% Generated solely from DER_EN_annotated_master.tex: DERform=1
\documentclass[11pt,oneside]{article}
\usepackage{fontspec}
\usepackage[french,english]{babel}
\usepackage{amsmath}
\usepackage{mathtools}
\usepackage{amssymb}
\usepackage{amsthm}
\usepackage{mathrsfs}
\usepackage{tikz-cd}
\usepackage{enumitem}
\usepackage{array}
\usepackage{ragged2e}
\usepackage{marginnote}
\usepackage{marginfix}
\usepackage{lmodern}
\usetikzlibrary{babel,arrows.meta,decorations.markings,decorations.pathreplacing}
\usepackage[a4paper,left=27mm,right=35mm,top=28mm,bottom=27mm,headsep=9mm,marginparwidth=25mm,marginparsep=5mm]{geometry}
\usepackage{xcolor}
\usepackage{fancyhdr}
\usepackage{xurl}
\usepackage{adjustbox}
% a diagram wider than the line or taller than the page is scaled down to fit (the widest, a chart of chapter II,
% runs to more than twice the width of the line); a diagram that fits is set as it is
\AddToHook{env/tikzcd/before}{\begin{adjustbox}{max width=\linewidth,max totalheight=\textheight}}
\AddToHook{env/tikzcd/after}{\end{adjustbox}}
\usepackage[hyperfootnotes=false,bookmarksopen=true,bookmarksopenlevel=1]{hyperref}
\hypersetup{pdftitle={Les Dérivateurs (English translation)},pdfauthor={Alexander Grothendieck},
  pdfsubject={English translation of the edition by M. Künzer, J. Malgoire and G. Maltsiniotis},
  colorlinks=true,linkcolor=black,urlcolor=blue!45!black,pdfdisplaydoctitle=true}
\setlength{\emergencystretch}{2em}
\setlength{\headheight}{14pt}
\makeatletter
\newcommand{\derleqno}{\tagsleft@true}
\renewcommand*{\@pnumwidth}{2.4em}
\renewcommand*{\@tocrmarg}{3.4em}
% \mathllap, \mathrlap, \mathclap as mathtools defines them (one unit uses \mathllap)
\providecommand*{\mathllap}[1]{\mathpalette\der@mathllap{#1}}
\providecommand*{\mathrlap}[1]{\mathpalette\der@mathrlap{#1}}
\providecommand*{\mathclap}[1]{\mathpalette\der@mathclap{#1}}
\def\der@mathllap#1#2{\llap{$\m@th#1{#2}$}}
\def\der@mathrlap#1#2{\rlap{$\m@th#1{#2}$}}
\def\der@mathclap#1#2{\hb@xt@\z@{\hss$\m@th#1{#2}$\hss}}
% the margin marks: p. n where page n of the edition begins, [page n] where page n of the manuscript begins
\newcommand{\derpp}[1]{\marginpar{\raggedright\footnotesize\color{black!55}p.~#1}}
\newcommand{\dermsp}[1]{\marginpar{\raggedright\footnotesize\ttfamily[#1]}}
% the editors' note of an edition file, under the chapter heading
\newenvironment{dernote}{\par\begingroup\small\setlength{\parindent}{0pt}\setlength{\parskip}{0.6ex}}%
  {\par\endgroup\vspace{2mm}\noindent\rule{0.3\textwidth}{0.4pt}\par\vspace{5mm}\@afterheading}
% #1 heading line, #2 title, #3 running head, #4 TOC text
\newcommand{\derchapter}[4]{\clearpage\setcounter{footnote}{0}\phantomsection
  \addcontentsline{toc}{section}{#4}\markboth{#3}{#3}\thispagestyle{plain}%
  \vspace*{10mm}{\centering\ifx\relax#1\relax\else{\large\scshape #1\par}\vspace{3mm}\fi
  {\LARGE\bfseries #2\par}}\vspace{9mm}\@afterheading}
% a later edition file of chapter XII
\newcommand{\derpart}[4]{\clearpage\setcounter{footnote}{0}\phantomsection
  \addcontentsline{toc}{subsection}{#4}\markboth{#3}{#3}%
  \vspace*{4mm}{\centering{\large\scshape #1\par}\vspace{2mm}{\Large\bfseries #2\par}}\vspace{7mm}\@afterheading}
\renewcommand*{\l@section}[2]{\addpenalty{-\@highpenalty}\vskip 0.8em plus 1pt
  {\leftskip 0pt\rightskip\@tocrmarg\parfillskip -\rightskip\parindent 0pt\interlinepenalty\@M
   \leavevmode\normalsize #1\nobreak\leaders\hbox{$\m@th\mkern\@dotsep mu\hbox{.}\mkern\@dotsep mu$}\hfill
   \nobreak\hb@xt@\@pnumwidth{\hss #2}\par}}
\renewcommand*{\l@subsection}{\@dottedtocline{2}{1.5em}{0em}}
\makeatother
% two styles, each setting all its fields: the body (running head, page number at the right of the head) and plain
% (page number in the foot: front matter, chapter openings); a style set with \pagestyle{fancy} would keep the fields
% that \pagestyle{plain} cleared
\fancypagestyle{derbody}{\fancyhf{}\fancyhead[L]{\small\itshape\leftmark}\fancyhead[R]{\small\thepage}%
  \renewcommand{\headrulewidth}{0pt}}
\fancypagestyle{plain}{\fancyhf{}\fancyfoot[C]{\small\thepage}\renewcommand{\headrulewidth}{0pt}}
\pagestyle{derbody}

\usepackage{luacode}
\providecommand{\DERform}{0}
\newcommand{\DERmspin}[1]{\marginnote{\raggedright\scriptsize\ttfamily[#1]}}
\newcommand{\DERinlinepage}[1]{\marginpar{\raggedright\footnotesize\ttfamily[#1]}}
\newcommand{\DERlacuna}{\ifmmode\mathord{[\ldots]}\else[\ldots]\fi}
\newfontfamily\DERappfont{DejaVu Sans}[Scale=MatchLowercase]
\newfontfamily\DERcodefont{DejaVu Sans Mono}[Scale=MatchLowercase]
\DeclareRobustCommand{\DERmark}[1]{\ifmmode\text{\hyperlink{app-#1}{\textsuperscript{\sffamily\tiny[#1]}}}\else\hyperlink{app-#1}{\textsuperscript{\sffamily\tiny[#1]}}\fi\allowbreak}
\newcommand{\DERapparatusitem}[5]{\par\medskip{\DERappfont\noindent\hypertarget{app-#1}{}\textbf{#1}\quad #2\par
  {\footnotesize\noindent\textbf{Print / imprimé:}\ {\DERcodefont #3}\par
  \noindent\textbf{Correction:}\ {\DERcodefont #4}\par
  \noindent\textbf{Source:}\ #5\par}}}
\begin{document}
\pagestyle{plain}
\hypersetup{pageanchor=false}
\begin{titlepage}\centering\vspace*{0.16\textheight}
{\Large Alexander Grothendieck\par}\vspace{2em}
{\Huge\bfseries Les Dérivateurs\par}\vspace{2em}
{\Large Corrected text with apparatus\par}\vspace{1em}
{\large English translation\par}\vfill
{\normalsize Edition by M.~Künzer, J.~Malgoire et G.~Maltsiniotis\par}
\vspace{1em}{\small 12 November 2026\par}\end{titlepage}
\hypersetup{pageanchor=true}
\textbf{Credits.} Author: Alexander Grothendieck, \emph{Les Dérivateurs} (1990–1991).
The edition transcribed is the electronic edition by M.~Künzer, J.~Malgoire and G.~Maltsiniotis
(Universités de Montpellier~II et Paris~VII). The deciphering, edition and editors’ notes are theirs.
AI-transcribed by Claude Opus~5.5 (Anthropic). Companion commentary and AI-produced corrections by GPT-6 Astra
(OpenAI, ChatGPT Pro mode), checked and integrated by Claude Opus~5.5 (Anthropic).
Eight commentary blocks in II-a, III-V and XIIb-a, seven paragraphs of other blocks, nineteen rows of the claims
register (II-a, III-V, VII-c, VIII-X, XIIb-a) and three bibliography entries are AI-produced by Claude Fable~5.2
(Anthropic) from GPT-6 Astra’s answers that it reviewed; each relevant row names the answer and reviewer.
Macro-definition corrections and changes in the checked texts absent from their tables were compared against the
edition’s pages by Claude Opus~5.5 (\texttt{transcription/DECISIONS.md}).
AI-translated: translated by GPT-6.1 Sol (OpenAI, effort max). Glossary by Claude Fable~5.2 and Claude Opus~5.5;
prompt, checks, fixes, completeness comparison and the published English book by Claude Opus~5.5 (Anthropic).
Raw research answers AI-produced by GPT-6 Astra; reviews by Claude Fable~5.2.
This reader build, apparatus and audit are by GPT-6 Astra (OpenAI, ChatGPT Pro mode).

\medskip\noindent\url{https://webusers.imj-prg.fr/~georges.maltsiniotis/groth/Derivateurs.html}

\bigskip\noindent\textbf{Rights}\par
The rights in the original work belong to its author or the author's representatives or heirs. Where we hold rights, in the AI-produced transcription, translation, corrections and typesetting, they are dedicated to the public domain under CC0 1.0. This edition is published for educational, non-profit purposes, and we claim no rights we do not hold. If you hold rights in this work and object to its publication here, please open a GitHub issue and we will take it down.

\clearpage\textbf{Reading.} Margin marks \texttt{p. n} locate edition pages; \texttt{[page n]} locates manuscript
pages. Anchor arguments and order are retained. Typewriter words in square brackets are the original editors’
notes. Isolated question marks in square brackets also record the editors’ doubts, under the convention stated in their preface.
Our apparatus variants carry an FRV or ENV identifier linked to the final register, with the printed
reading, correction and source. To identify symbols and unit-specific macros exactly, the readings in that
register are quoted as literal \LaTeX{} source. The original notes are preserved verbatim. All three forms
come from one annotated master. Typesetting repairs common to the three forms are recorded separately from
textual variants.

\clearpage\tableofcontents\clearpage\pagestyle{derbody}
\clearpage\setcounter{footnote}{0}\phantomsection\addcontentsline{toc}{section}{The edition’s table of contents}\markboth{The edition’s table of contents}{The edition’s table of contents}
% BEGIN UNIT dertbmt/dertbmt_p001-005.tex
\begingroup\makeatletter%
\newcommand{\printedpage}[1]{\marginpar{\footnotesize p.~#1}}%
\newcommand{\msp}[1]{\marginpar{\footnotesize [#1]}}%
\newcommand{\tochead}[2]{\par\medskip\noindent #1\hfill #2\par\smallskip}%
\newcommand{\tocl}[4]{\par\noindent\hangindent=\dimexpr #1+3em\relax\hspace*{#1}\makebox[2.5em][r]{#2}\hspace{0.5em}#3\dotfill #4}%
\newcommand{\Ch}{\underline{\mathrm{Ch}}}%
\makeatother\let\printedpage\derpp\let\msp\dermsp%
\let\mspin\DERmspin%


\printedpage{1}
% head: (none)

\begin{center}
{\LARGE\textbf{The Derivators}}

\bigskip

{\large Alexandre Grothendieck}
\end{center}

\bigskip

\tochead{\textbf{I. Generalities on (pre)derivators.}}{[1]-[126]}
\tocl{0em}{1.}{Preliminaries on the inverse image functors for presheaves on index categories, and their multiple adjoints}{[1]}
\tocl{0em}{2.}{Heuristic preliminaries on derivators}{[4]}
\tocl{0em}{}{Homotopically proper, homotopically smooth morphisms}{[53]}
\tocl{0em}{}{Notions of $\mathbb{D}$-equivalence and of $\mathbb{D}$-contractibility}{[61]}
\tocl{0em}{}{$\mathbb{D}$-fibrations}{[76]}

\tochead{\textbf{II. Cofinality (on the right and on the left) (preliminary to cohomological cofinality).}}{[1]-[99]}
\tocl{0em}{1.}{Cohomological cofinality}{[1]}
\tocl{0em}{2.}{The locally filtering categories are the locally totally 0-connected categories}{[11]}
\tocl{0em}{3.}{Return to cofinality}{[37]}
\tocl{0em}{4.}{Biaspheric functors}{[49]}
\tocl{0em}{5.}{Formal duality}{[69]}
\tocl{0em}{6.}{Transporters, kernels, stabilizers relative to $W$}{[83]}
\tocl{0em}{7.}{Cofiltering and locally cofiltering morphisms}{[90]}

\tochead{\textbf{III. External Hom in derivators.}}{[1]-[12]}
\tocl{0em}{1.}{Relation of derivators with the category Hot}{[1]}

\tochead{\textbf{IV. Substantial diagrams.}}{[73]-[83]}
\tocl{0em}{}{Construction of substantial diagrams for a derivator}{[73]}

\tochead{\textbf{V. Categories of paths and localization.}}{[1]-[52]}
\tocl{0em}{1.}{The category of paths}{[1]}
\tocl{0em}{2.}{Calculation of a category of fractions $X\Sigma^{-1}$}{[5]}
\tocl{0em}{3.}{Passage to categories of diagrams}{[28]}

\tochead{\textbf{VI. HOT.}}{[1]-[204]}
\tocl{0em}{}{The derivator HOT}{[1]}
\tocl{0em}{}{New version of the theory of $\mathrm{HOT}_W$}{[171]}

\par
\printedpage{2}
% head: 2

\tochead{\textbf{VII. Categories of paths (2).}}{[1]-[171]}
\tocl{0em}{1.}{Intervals and their variants}{[1]} % [sic] leur variantes -> leurs variantes, accord du pluriel
\tocl{0em}{2.}{Full subcategories of an interval}{[10]}
\tocl{0em}{3.}{Paths. Category of paths of a given type $\tau$}{[13]}
\tocl{0em}{4.}{Category of paths $\Ch(X)$}{[24]}
\tocl{0em}{5.}{Properties of morphisms of intervals}{[31]}
\tocl{0em}{6.}{Composition in \underline{Ty}}{[48]}
\tocl{0em}{7.}{Explicitation of the categories of fractions $X\Sigma^{-1}$ in terms of the $\Ch(X;a,b)$}{[59]}
\tocl{0em}{8.}{The homotopy lemma for morphisms of intervals}{[83]}
\tocl{0em}{9.}{Homotopy theorem for categories of paths}{[94]}
\tocl{0em}{10.}{Comparison question for Cat-fibered categories satisfying the ``homotopy theorem'' (filtering base)}{[99]}
\tocl{0em}{11.}{Quick return to the $(\mathcal{M},\Sigma)$, from the viewpoint of derivators}{[129]}
\tocl{0em}{12.}{Return to the ``categories of paths'' in Cat, as a category of models for Hot}{[137]}
\tocl{0em}{13.}{Where we finally prove that $\Ch(X;-,-)$ is as good as (or even better than) $\Ch_\infty(X;-,-)$}{[157]}

\tochead{\textbf{VIII. Relative homotopy 1-types: their integration\,\dots}}{[1]-[9]}
\tocl{0em}{1.}{Integration of relative homotopy 1-types over a category}{[1]}
\tocl{0em}{2.}{Questions of base change}{[6]}

\tochead{\textbf{IX. Return to the categories $\mathcal{M}W^{-1}$, comparison with Quillen.}}{[1]-[25]}
\tocl{0em}{1.}{Functoriality of $\mathrm{Ch}^{\Sigma}_{\infty}(a,b)$ in $a$, $b$ in $\mathcal{M}\Sigma^{-1}$}{[1]}
\tocl{0em}{2.}{Comparison with Quillen}{[14]}
\tocl{0em}{3.}{Attempt at constructing the functor $\Omega$}{[24]}

\tochead{\textbf{X. Comparison of Cat with $\Delta^\wedge$.}}{[1]-[5]}
\tocl{0em}{1.}{Reminder of the notions of $\mathbb{D}$-asphericity, and link with the ``basic localizer $W_{\mathbb{D}}$''}{[1]}

\tochead{\textbf{XI. Hot-fibrations, proper functors and smooth functors etc.}}{[127]-[136]}
\noindent\hspace*{1.5em}(in Cat).\par\smallskip
\tocl{0em}{1.}{Hot-fibrations in Cat}{[127]}

\tochead{\textbf{XII. Characterization of $W_\infty$. $W$-proper, $W$-smooth functors etc. Amalgamated sums and $W$-cocartesian squares in Cat.}}{[1]-[275]}
\tocl{0em}{I}{Characterization of $W_\infty$}{[1]}
\tocl{0em}{II}{Perplexities about the axioms and the terminology}{[18]}
\tocl{0em}{III}{Remarkable classes of functors related to $W$}{[22]}
\tocl{0em}{IV}{On axiom W(4) and its variants}{[53]}
\tocl{0em}{V}{Amalgamated sums and $W$-cocartesian squares in Cat}{[59]}
\tocl{2em}{1.}{Generalities on $U\overset{i}{\hookrightarrow}X\overset{j}{\hookleftarrow}T$ and amalgamated sums}{[59]}
\tocl{4.5em}{1.1.}{Gluing data between an open and a closed}{[59]}
\tocl{4.5em}{1.2.}{Relation with the gluing of topoi}{[61]}

\par
\printedpage{3}
% head: 3

\tocl{4.5em}{1.3.}{Arrows in Cat, with source or target a glued object $(X,U,T)$}{[65]}
\tocl{4.5em}{1.4.}{Construction of amalgamated sums from an open where closed immersion}{[66]} % [sic] où fermée -> ou fermée, conjonction (cf. 2.2)
\tocl{2em}{2.}{Amalgamated sums and $W$-equivalences}{[75]}
\tocl{4.5em}{2.1.}{Cocartesian square axiom and cube axiom (or ``Five Lemma'')}{[75]}
\tocl{4.5em}{2.2.}{Application to cocartesian squares deduced from an open or closed immersion}{[80]}
\tocl{4.5em}{2.3.}{Totally $W$-aspheric objects of Cat, and totally $W$-aspheric morphisms}{[87]}
\tocl{4.5em}{2.4.}{Application to open Dwyerian immersions}{[107]}
\tocl{4.5em}{2.5.}{Complement: criterion for total asphericity of a $W$-proper category over another}{[115]}
\tocl{2em}{3.}{Amalgamated sums and integration}{[121]}
\tocl{4.5em}{3.1.}{Canonical reduction to the case of a diagram of closed immersions}{[121]}
\tocl{4.5em}{3.2.}{Comparison of the integral with an ordinary amalgamated sum. $W$-homotopically cocartesian squares}{[128]}
\tocl{4.5em}{3.3.}{Proper integration versus smooth integration. Statement of equivalence between the two points of view. Comments on the case of arbitrary $W$-homotopic limits}{[133]}
\tocl{4.5em}{3.4.}{Criteria for homotopically cocartesian squares}{[146]}
\tocl{4.5em}{3.5.}{Applications to amalgamated sums in a category $X^\wedge$ (with $X\in\mathrm{Ob}\,\mathrm{Cat}$)}{[158]}
\tocl{0em}{VI}{Criterion for $W_\infty$-asphericity of the $\Omega^n(X,x)$}{[171]}
\tocl{2em}{1.}{Cogitations on $W_n$}{[171]}
\tocl{2em}{2.}{Criterion for $W_n$-asphericity}{[187]}
\tocl{2em}{3.}{On the philosophy of homotopy fibers}{[191]}
\tocl{2em}{4.}{Axiomatization of ``$W$-homotopy fibers''}{[207]}
\tocl{2em}{5.}{Axiomatics of non-locally constant relative $W$-homotopy types $\mathrm{Ho}^{\mathcal{L}}_{W}(S)$}{[218]}
\tocl{2em}{6.}{$W$-homotopy cofiber in Cat}{[227]}
\tocl{2em}{7.}{Return to the axiomatics of the $W$-homotopy fiber}{[240]}
\tocl{2em}{8.}{Appendix: Return to axioms W(7), W(7 bis) (pages 75, 76)}{[265]}

\tochead{\textbf{XIII. Categories of models (1).}}{[1]-[98]}
\tocl{0em}{1.}{Quillen couples $(\Phi,\Psi)$}{[1]}
\tocl{0em}{2.}{Weak Quillen categories, and construction thereof (Conjectural theorem-scholie)}{[7]}
\tocl{0em}{3.}{Quillen categories (second version)}{[27]}
\tocl{0em}{4.}{The variances $u^*$, $u_*$, $u_!$ for $(\mathcal{M}(I),W(I))$ (I variable)}{[44]}
\tocl{0em}{5.}{Quillenizable categories of models}{[50]}
\tocl{0em}{6.}{Quillenizable categories of models, and functors $u_!$, $u_*$}{[57]}
\tocl{0em}{7.}{Cardinal filtrations and accessible subsets of a big category}{[77]}

\par
\printedpage{4}
% head: 4

\tochead{\textbf{XIV. h-cartesian and h-cocartesian squares.}}{[1]-[8]}
\tocl{0em}{1.}{Axioms for quadrangulated structures (an aborted attempt)}{[1]}
\tocl{0em}{2.}{Examples}{[7]}

\tochead{\textbf{XV. Factorization theorems. (Models (2)).}}{[1]-[137]}
\tocl{0em}{1.}{Fundamental construction: the transfinite factorization sequence of an arrow $f$}{[1]}
\tocl{0em}{2.}{Variation of the factorizations $f=p_\alpha i_\alpha$ with the generating subset $C_0$}{[12]}
\tocl{0em}{3.}{Bounding the cardinals}{[16]}
\tocl{0em}{4.}{Complements on Quillen structures}{[28]}
\tocl{0em}{5.}{Preliminary technical results on accessible subsets of $\mathrm{Ob}\,\mathcal{M}$, stable under filtering direct limits}{[43]} % [sic] téchniques -> techniques, accent
\tocl{0em}{6.}{Aborted attempt at a generation criterion for L-accessible subsets of $\mathrm{Fl}(\mathcal{M})$}{[73]}
\tocl{0em}{7.}{Construction of closed Quillen triples with $C=\mathrm{Mon}\subset\mathrm{Cof}_W$}{[84]}
\tocl{0em}{8.}{Application to the categories of models $(\widehat{A},W_A)$}{[100]}
\tocl{0em}{9.}{Revision of the notations, for the axioms of a basic localizer}{[116]}
\tocl{0em}{10.}{Application to the derivator $\mathrm{HOT}_W$}{[134]}

\tochead{\textbf{XVI. Basic localizers in Cat.}}{[1]-[124]}
\tocl{0em}{1.}{Axioms Loc\,(1) to Loc\,(3) (and variants): basic localizers}{[1]}
\tocl{0em}{2.}{Axioms Loc\,(4) and Loc\,(5) (of direct sums, of connectedness)}{[8]}
\tocl{0em}{3.}{Axiom Loc\,(6) of the cocartesian square}{[16$'$]}
\tocl{0em}{4.}{Axiom Loc\,(7) of $W$-fibrations}{[27]}
\tocl{0em}{5.}{Limit axiom Loc\,(8), and axiom Loc\,(9) of accessibility}{[71]}
\tocl{0em}{6.}{The condition Loc\,$(5,\mathbb{D}_{W_\omega})$ implies $W\subset\mathrm{W}_\infty$}{[83]}
\tocl{0em}{7.}{Comments on the axioms}{[97]}

\tochead{\textbf{XVII. Categories with fibrations and with cofibrations. (Models (3)).}}{[1]-[118]}
\tocl{0em}{1.}{Preliminaries}{[1]}
\tocl{0em}{2.}{Description of $C_I\subset\mathrm{Fl}(\mathcal{M}(I^\circ))$, for $I$ an ordered set}{[25]}
\tocl{0em}{3.}{Categories of models with cofibrations}{[39]}
\tocl{0em}{4.}{Special categories with cofibrations}{[71]}
\tocl{0em}{5.}{Relation of special categories with cofibrations to special Quillen categories}{[84]}
\tocl{0em}{6.}{Application to the construction of derivators (associated with special Quillen triples)}{[98]}

\tochead{\textbf{XVIII. Accessible categories and sets.}}{[50]-[492]}
\tocl{0em}{4.}{Paratopoi and the operation $\mathcal{M}\boxtimes\mathcal{N}$}{[50]}
\tocl{2em}{4.1.}{Pseudo-topoi. Product problem}{[50]}
\tocl{2em}{4.2.}{Muddles over $\mathcal{M}\boxtimes\mathcal{N}$}{[57]}
\tocl{2em}{4.3.}{The operation $\mathcal{M}\boxtimes\mathcal{N}$ and the accessibility axiom}{[64]}
\tocl{2em}{4.4.}{Study of $\mathrm{Ind}_\pi(C)$ (for $C$ a small category stable under $\varinjlim$ of cardinality $\leq\pi$)}{[79]}

\par
\printedpage{5}
% head: 5

\tocl{2em}{4.5.}{Canonical cardinal filtration of a $\mathcal{U}$-paratopos category}{[106]}
\tocl{2em}{4.6.}{Characterization of paratopos categories (proof of 4.3.6)}{[139]}
\tocl{2em}{4.7.}{$\mathrm{Ind}_\pi C$ ($C$ arbitrary) and generalized cardinal filtrations}{[161]}
\tocl{2em}{4.8.}{$\mathcal{U}$-accessible categories}{[187]}
\tocl{2em}{4.9.}{The indicial representation theorem for $\mathrm{Ind}_\pi(C)$}{[217]}
\tocl{2em}{4.10.}{The indicial representation theorem: proof}{[243]}
\tocl{2em}{4.11.}{Application to the existence of limits in a category of the form $\mathrm{Ind}_{\pi_0}(C)$}{[267]}
\tocl{2em}{4.12.}{Accessible functors between accessible categories}{[273]}
\tocl{2em}{4.13.}{Accessible sets of objects of an accessible category}{[281]}
\tocl{2em}{4.14.}{Accessibility criteria for big categories}{[370]}
\tocl{2em}{4.15.}{Accessible category over another}{[425]}
\tocl{2em}{4.16.}{Examples and questions}{[464]}
\tocl{2em}{4.17.}{Return to paratopoi, and to the operations $\mathcal{M}\boxtimes\mathcal{N}$, $\underline{\mathrm{Hom}}_!(\mathcal{M},\mathcal{N})$}{[492]}
\par\noindent\hspace*{1.5em}\textbf{Appendix}.\hfill[1]-[19]
\tocl{2em}{}{Ultraproducts of fields, and accessible categories}{[1]}
\tocl{2em}{}{Application to an example of a non-accessible subcategory}{[11]}

\tochead{\textbf{XIX. Models (4).}}{[1]-[50]}
\tocl{0em}{1.}{Reminders on Quillen couples (special Quillen couples)}{[1]}
\tocl{0em}{2.}{Reminders on closed Quillen triples and special Quillen triples}{[13]}
\tocl{0em}{3.}{Case of categories with cofibrations, of categories with fibrations}{[33]}

\bigskip
\noindent\rule{4cm}{0.4pt}

\par\noindent\textbf{Note by the editors of the transcribed edition.}\par
{\footnotesize\noindent\hangindent=2em \textsc{Note}. In this table of contents, as well as in the provisional electronic edition of the ``Derivators'', we follow the division of the work carried out a posteriori by Grothendieck. The page numbers are those of the original, indicated in brackets in the digitized text. Chapter XI, which begins on page 127, was originally the last section of Chapter I, which ends on page 126. Likewise, Chapter XVIII was section 4 of Chapter XIX.\par}


\endgroup
% END UNIT dertbmt/dertbmt_p001-005.tex
\clearpage\setcounter{footnote}{0}\phantomsection\addcontentsline{toc}{section}{Chapter I}\markboth{Chapter I}{Chapter I}
% BEGIN UNIT Der-ChI-234/Der-ChI-234_p001-005.tex
\begingroup\makeatletter%
\newcommand{\printedpage}[1]{\marginpar{\footnotesize p.~#1}}%
\newcommand{\msp}[1]{\marginpar{\footnotesize [#1]}}%
\newcommand{\AMSTeX}{$\mathcal{A}\kern-.1667em\lower.5ex\hbox{$\mathcal{M}$}\kern-.125em\mathcal{S}$-\TeX}%
\newcommand{\uHom}{\underline{\mathsf{Hom}}}%
\newcommand{\Ob}{\mathsf{Ob}}%
\newcommand{\Ab}{\mathsf{Ab}}%
\newcommand{\bDelta}{\boldsymbol{\Delta}}%
\makeatother\let\printedpage\derpp\let\msp\dermsp%
\let\mspin\DERmspin%


% ===================== page 1 =====================
\printedpage{--} % no page number printed (title page)
% head: (none)

\begin{center}
{\Large\textsf{THE DERIVATORS}}

\bigskip

\textbf{ALEXANDRE GROTHENDIECK}

\bigskip\bigskip

\textbf{Chapter I}

\medskip

\textbf{Generalities on (pre)derivators}
\end{center}

\par

% ===================== page 2 =====================
\printedpage{--} % no page number printed (editors' note)
% head: (none)

\par\noindent\textbf{Note by the editors of the transcribed edition.}\par
This text was deciphered and transcribed in \AMSTeX{} by G. Maltsiniotis. It was edited by M. Künzer, J. Malgoire, and G. Maltsiniotis. The transcription is as faithful as possible to the manuscript. For the few obvious corrections, or the rare comments by the editors, as well as for the original numbering of the manuscript pages, typewriter characters \texttt{[typewriter]} in brackets are used. A question mark in brackets means that one is not sure of the preceding word.

This edition is provisional. Remarks, comments and corrections are welcome. Send a message to:

\textbf{[e-mail address removed]}

\hfill G. Maltsiniotis

\par

% ===================== page 3 =====================
\printedpage{1}
% head: A. Grothendieck : Dérivateurs, Ch. I, Édition des Universités de Montpellier II et Paris VII
% foot: Version 8 juillet 2001

\msp{page 1}

\section*{Derivators}

\subsection*{\textmd{1.}\ Preliminaries on the inverse image functors of presheaves on index categories, and their multiple adjoints.}

$A$ ``category of coefficients'',\\
$I$ ``index category'',\\
$A(I)=\uHom(I^{\circ},A)$ presheaves on $I$ with values in $A$.

If $f : I \to J$, we get:
\[
f^{*} : A(J) \to A(I) \quad ,
\]
and we are interested in the existence of the successive adjoint functors on the left and on the right. Provisional notation:
\[
f^{?}\ f_{!}\ f^{*}\ f_{*}\ f^{!} \quad .
\]

Here are the main facts:

\textit{a}) Subject to the existence of the relevant $\varprojlim$ (resp. $\varinjlim$), $f_{*}$ (resp. $f_{!}$) exists. For example if there exist \texttt{[some]} $\varprojlim$ (resp. $\varinjlim$) that are \textit{finite}, then $f_{*}$ (resp. $f_{!}$) exists if $I$, $J$ are finite categories.

\textit{b}) Even if $A$ is an abelian category, with $\varprojlim$ and $\varinjlim$ (for example $\Ab$), the \textit{biadjoints} do not always exist. For example, in the particular case where $J=e$ is the one-point category, and $I$ the ordered set
\[
\begin{tikzcd}
 & a \\
0 \arrow[ur] \arrow[dr] & \\
 & b
\end{tikzcd}
\]
(in the case of $f^{!}$), or the opposite ordered set (in the case of $f^{?}$). Even if $f$ is an open immersion, $f^{!}$ does not necessarily exist (for example it does not exist if $I$ is as \texttt{[here-]}above, and
\[
J=\begin{tikzcd}
 & a \arrow[dr] & \\
0 \arrow[ur] \arrow[dr] & & c \\
 & b \arrow[ur] &
\end{tikzcd}
\qquad\qquad \text{),}
\]
and if $f$ is a closed immersion, $f^{?}$ does not necessarily exist.

\textit{c}) If $f$ is a \textit{closed immersion}, and $A$ is \textit{pointed} (\textit{i.e.} admits an object both initial and final), then $f^{!}$ exists. If $f$ is an open immersion, $f^{?}$ exists. (Always subject
\printedpage{2}
% head: A. Grothendieck : Dérivateurs, Ch. I, Édition des Universités de Montpellier II et Paris VII
% foot: Version 8 juillet 2001
to the existence in $A$ of the relevant limits). If $f$ is an immersion, $f_{*}$ and $f_{!}$ \texttt{[are]} fully faithful \textit{i.e.} $f^{*}$ \texttt{[is]} a localization functor.

\textit{d}) Let us assume only the existence of the final object $e$ in $A$ (plus if need be certain finite $\varprojlim$). If

\msp{page 2}
\[
i : Y \hookrightarrow S
\]
is a closed immersion, with complement $U$
\[
Y \overset{i}{\hookrightarrow} S \overset{j}{\hookleftarrow} U
\]
so that $j$ is an open immersion, then $A(Y)$ is identified by the functor $i_{*}$ with the full subcategory of $A(S)$ consisting of the $F \in \Ob\ A(S)$ such that $F|U \simeq e_{U}$ (final presheaf). Similarly, if the initial object $\emptyset$ exists in $A$, and the relevant $\varinjlim$, then $A(U)$ is identified by $j_{!}$ with the full subcategory of $A(S)$ consisting of the $F$ such that $F|Y=\emptyset_{Y}$ (initial presheaf, taking the value $\emptyset$ everywhere).

If $A$ is even \textit{pointed} ($e=\emptyset$), then we have, for every $F\in \Ob\ A(S)$, a square both cartesian and cocartesian, functorial in $F$
\[
\begin{tikzcd}
 & F \arrow[dr, "\text{adj}"] & \\
j_{!}j^{*}F \arrow[ur, "\text{adj}"] \arrow[dr] & & i_{*}i^{*}F \\
 & e \arrow[ur] &
\end{tikzcd}
\qquad\qquad \text{.}
\tag*{(1)}
\]
In the case $A$ abelian, this diagram takes the form of an exact sequence
\[
0 \longrightarrow j_{!}j^{*}F \longrightarrow F \longrightarrow i_{*}i^{*}F \longrightarrow 0 \quad .
\]

\textit{e}) Same conditions for $i$ and $j$. We assume that $A$ has the relevant $\varprojlim$, including final object. We then have a cartesian square functorial in $F\in \Ob\ A(S)$ \texttt{[one must assume the category} $A$ \texttt{pointed]}
\[
\begin{tikzcd}
F \arrow[d, "\text{adj}"'] & i_{*}i^{!}F \arrow[l] \arrow[d] \\
j_{*}j^{*}F & e \arrow[l]
\end{tikzcd}
\tag*{(2)}
\]
\textit{i.e.} $i^{!}F$ is identified by the fully faithful functor $i_{*}$ with the \textit{kernel} of $F\to j_{*}j^{*}F$. In the abelian case, this takes the form of an exact sequence
\[
0 \longrightarrow i_{*}i^{!}F \longrightarrow F \longrightarrow j_{*}j^{*}F \quad .
\]

\par

% ===================== page 5 =====================
\printedpage{3}
% head: A. Grothendieck : Dérivateurs, Ch. I, Édition des Universités de Montpellier II et Paris VII
% foot: Version 8 juillet 2001

\msp{page 3}

I dispense with writing down the dual cocartesian square, when $A$ has the relevant $\varinjlim$, including initial object. It seems that it has not been of use in practice. The classically useful adjoint and biadjoint functors are included below (where the ``not useful'' ones are in brackets)
\[
\begin{aligned}
{}[i_{!}]\ i^{*}\ i_{*}\ i^{!} \\
[j^{?}]\ j_{!}\ j^{*}\ j_{*}
\end{aligned}
\]
\textit{i.e.} the functors $i_{!}$ and $j^{?}$ have not been of use classically.

\textit{f}) Case of the categories $\bDelta_{i}$ (standard simplices, $i\in\mathbb{N}$), and of the functors between them (semi-simplicial maps). It would seem that for such an $f$, if $A$ is pointed, there are \textit{bi-adjoints} both on the left and on the right. Moreover, for $\bDelta_{i}$ and $\bDelta_{j}$ fixed when we look at all the $f:\bDelta_{i}\to\bDelta_{j}$ and all the $g:\bDelta_{j}\to\bDelta_{i}$, that the total sequences of multi-adjoint functors for each of the corresponding $f^{*}$, $g^{*}$ are all the same \texttt{[this is not true]}. That looks true at least if $j=i+1$ (and I have checked it for $i=0$, $i=1$), we should find a ``sequence of adjunctions'' of length $2i+7=(i+1)+(i+2)+2+2$. To be precise, the sequence has four ``extreme'' terms (the first two, the last two) and the $(i+1)+(i+2)$ intermediate terms corresponding to the functors of the form $f^{*}$ or $g^{*}$, where $f:\bDelta_{i+1}\to\bDelta_{i}$ is one of the $(i+1)$ degeneracy maps, and $g:\bDelta_{i}\to\bDelta_{i+1}$ one of the $i+2$ face maps \texttt{[this is not quite true; one must restrict oneself to the maps that are injective or surjective]}.


\endgroup
% END UNIT Der-ChI-234/Der-ChI-234_p001-005.tex
% BEGIN UNIT Der-ChI-234/Der-ChI-234_p006-010.tex
\begingroup\makeatletter%
\newcommand{\printedpage}[1]{\marginpar{\footnotesize p.~#1}}%
\newcommand{\msp}[1]{\marginpar{\footnotesize [#1]}}%
\newcommand{\Dia}{\mathsf{Dia}}%
\newcommand{\Cat}{\mathcal{C}\mathit{at}}%
\newcommand{\DD}{\mathbb{D}}%
\newcommand{\Hom}{\underline{\mathsf{Hom}}}%
\newcommand{\id}{\mathsf{id}}%
\newcommand{\Ob}{\mathsf{Ob}}%
\newcommand{\Del}{\boldsymbol{\Delta}}%
\makeatother\let\printedpage\derpp\let\msp\dermsp%
\let\mspin\DERmspin%


\printedpage{4}
% head: A. Grothendieck : Dérivateurs, Ch. I, Édition des Universités de Montpellier II et Paris VII
% foot: Version 8 juillet 2001

\msp{page 4}

\section*{2. Heuristic preliminaries on derivators.}

We take a full subcategory $\Dia$ (as in ``diagrams'') of $\Cat$. It can be all of $\Cat$, or smaller. But we will most often assume that $\Dia$ contains at least the ``ordered categories'' (corresponding to ordered sets, \textit{i.e.} such that $a$) every morphism is a mono and $b$) every iso is an identity \texttt{[this is not true: the category reduced to a double arrow satisfies these conditions, but is not an ordered set]}). By abuse of language, we identify an ordered set with the corresponding category (and a category with the topos it defines?). We regard $\Dia$ as a 2-\textit{category}.\footnote{N.B. We need stability conditions for $\Dia$, at least stability under $I \times J$ and $I \amalg J$, in fact, considerably more \dots}

A \textit{derivator} $\DD$ (on $\Dia$, as the ``\textit{domain}'' of the derivator) is a contravariant 2-functor from $\Dia$ to $\Cat$, satisfying certain conditions. Thus
\[ \DD : \Dia^{\circ} \longrightarrow \Cat \quad . \]
If $I \in \Ob_0 \Dia$ is a category (or a ``type of diagram'') in the domain of $\DD$, we must therefore have a category $\DD(I)$. The objects of $\DD(I)$ will be called the \textit{coefficient-objects} (or \textit{coefficients}) of type $\DD$, over $I$. For
\[ f : I \longrightarrow J \]
a 1-arrow of $\Dia$, we must have a functor
\[ f^*_{\DD} \text{ or } f^* : \DD(J) \longrightarrow \DD(I) \quad , \]
called the \textit{``inverse image functor''}. If $F$ is a coefficient of type $\DD$ on $J$ its ``\textit{inverse image}'' is therefore a coefficient of type $\DD$ on $I$. We assume (to simplify!) that the 2-functor $\DD$ is \textit{strict}, hence that we have strict transitivity
\[ (gf)^* = f^*g^* \ , \qquad \text{if} \qquad I \xrightarrow{f} J \xrightarrow{g} K \qquad \text{in} \ \ \Dia \ . \]

\msp{page 5}

Moreover, if we have two functors $f, g : I \rightrightarrows J$, and a morphism between them
\[
\begin{tikzcd}[column sep=large]
I \arrow[r, shift left=2, "f"] \arrow[r, shift right=2, "g"'] \arrow[r, phantom, "{\downarrow}u"] & J
\end{tikzcd}
\quad ,
\]
we take a morphism $u^*_{\DD}$ or $u^*$ between the functors $g^*$ and $f^*$, whence
\[
\begin{tikzcd}[column sep=large]
\DD(I) & \DD(J) \arrow[l, shift right=2, "f^*"'] \arrow[l, shift left=2, "g^*"] \arrow[l, phantom, "{\uparrow}u^*"]
\end{tikzcd}
\quad .
\]

\printedpage{5}
% head: A. Grothendieck : Dérivateurs, Ch. I, Édition des Universités de Montpellier II et Paris VII
% foot: Version 8 juillet 2001

We must also have transitivity for a composition of morphisms of functors:
\[
\begin{tikzcd}[column sep=large]
I \arrow[r, shift left=4, "f"] \arrow[r, "g"' near start] \arrow[r, shift right=4, "h"'] \arrow[r, phantom, shift left=2, "{\downarrow}u"] \arrow[r, phantom, shift right=2, "{\downarrow}v"] & J
\end{tikzcd}
\]
must give
\[ (vu)^* = u^*v^* \]
thus
\[
\begin{tikzcd}[column sep=large]
\DD(I) & \DD(J) \arrow[l, shift right=4, "f^*"'] \arrow[l, "g^*" near end] \arrow[l, shift left=4, "h^*"] \arrow[l, phantom, shift right=2, "{\uparrow}u^*"] \arrow[l, phantom, shift left=2, "{\uparrow}v^*"]
\end{tikzcd}
\quad .
\]
Moreover, if $f$ is an identity, $f^*$ too
\[ (\id_I)^* = \id_{\DD(I)} \quad , \]
and likewise if $u$ is an identity
\[ (\id_f)^* = \id_{f^*} \quad . \]

\msp{page 6}

I indicate right away the most important axiom on a derivator (beyond the datum of a strict 2-contra functor):

\textbf{Der}$_1$: If $f : I \to J$ in $\Dia$, then $f^*_{\DD} = f^* : \DD(J) \to \DD(I)$ has a right adjoint and a left adjoint.

These functors will be denoted (at least provisionally) $f^{\DD}_*$ or $f_*$ , and $f^{\DD}_!$ or $f_!$ , that is
\[ f_! \, , \ f_* : \DD(I) \rightrightarrows \DD(J) \quad , \]
so that we find a sequence of \textit{three} adjoint functors between $\DD(I)$ and $\DD(J)$
\[ f_! \, , \ f^*, \ f_* \quad . \]

N.B. For most of the developments we will not need the existence of both the $f_!$ and the $f_*$ . A \textit{left derivator} requires only the $f_*$ (right adjoints of the $f^*$), and the existence of $f_!$ only when $f$ is an \textit{open immersion}. Dually, a \textit{right derivator} requires only the $f_!$ (left adjoints of the $f^*$), and $f_*$ only if $f$ is a \textit{closed immersion}. It is probably the left derivators that are the most important.

\msp{page 7}

\textbf{Der}$_2$: If $f : I \to J$ is in $\Dia$, if it is a \textit{closed} immersion, the right biadjoint $f^!_{\DD}$ or $f^!$ of $f^*$ exists. If it is an open immersion, the left biadjoint $f^?$ exists.\footnote{The biadjoints should be required only for a pointed derivator. \textsf{HOT} does not satisfy this.} In
\printedpage{6}
% head: A. Grothendieck : Dérivateurs, Ch. I, Édition des Universités de Montpellier II et Paris VII
% foot: Version 8 juillet 2001
both cases, we want $f^*$ to be a localization functor, which is equivalent to saying that $f_*$ is fully faithful, or that $f_!$ is fully faithful.\footnote{We must in addition characterize the image of $f_*$ if $f$ is a closed immersion (the $F$ such that $g^*F \overset{\sim}{\to} e_{\DD(U)}$) and that of $f_!$ if it is an open immersion (the $F$ such that $g^*F \simeq \emptyset_{\DD(Y)}$), where $g$ is the complementary immersion.}

N.B. In the case of a left derivator, we assume only the existence of $f^!$ if $f$ is a closed immersion. Thus to summarize, for such a derivator, we assume the existence of
\[
\left\{
\begin{aligned}
& f_* \text{ for any } f : I \to J \text{ in } \Dia \\
& f_! \text{ if } f \text{ is an open immersion} \\
& f^! \text{ if } f \text{ is a closed immersion (pointed case only).}
\end{aligned}
\right.
\]
Dually, for a right derivator, we assume the existence of
\[
\left\{
\begin{aligned}
& f_! \text{ for any } f : I \to J \text{ in } \Dia \\
& f_* \text{ if } f \text{ is a closed immersion} \\
& f^? \text{ if } f \text{ is an open immersion (pointed case only).}
\end{aligned}
\right.
\]
I would like to formulate clever axioms, for the case of a pair
\[ Y \overset{i}{\hookrightarrow} I \overset{j}{\hookleftarrow} U \]
of an open immersion $j$ and a closed immersion $i$ in $\Dia$, complementary to one another. First a formulation if we assume $\DD$ ``\textit{pointed}'', \textit{i.e.} the categories $\DD(I)$ pointed, and the functors $f^*$ compatible with the pointings. We want to ``interpret'' the square (1) on page 2:

\msp{page 8}

\[
\begin{tikzcd}
& F \arrow[dr, "\mathrm{adj}"] & \\
j_!j^*F \arrow[ur, "\mathrm{adj}"] \arrow[dr] & & i_*i^*F \\
& e \arrow[ur] &
\end{tikzcd}
\quad ,
\tag*{(1)}
\]

\textbf{Der}$'_2$: For all $I, J \in \Dia$, we have
\[ \DD(I \amalg J) \xrightarrow[(\alpha^*, \beta^*)]{\sim} \DD(I) \times \DD(J) \quad , \]
where $\alpha, \beta$ are the canonical injections of $I$ and $J$ into $I \amalg J$.

\printedpage{7}
% head: A. Grothendieck : Dérivateurs, Ch. I, Édition des Universités de Montpellier II et Paris VII
% foot: Version 8 juillet 2001

square both cartesian and cocartesian. We are thus looking for a square in $\DD(I)$, but better than that. Let $J = \Del_1 \times \Del_1$ be the standard square
\[
\begin{tikzcd}
& (1,0) \arrow[dl] & \\
(1,1) & & (0,0) \arrow[ul] \arrow[dl] \\
& (0,1) \arrow[ul] &
\end{tikzcd}
\]
and let us consider the category $I \times J$. We have
\[
\begin{aligned}
& J \longrightarrow \Hom(I, I \times J) \\
& j \longmapsto \bigl(i \longmapsto (i,j)\bigr)
\end{aligned}
\]
On the other hand, by $\DD$
\[ \Hom(I, I \times J)^{\circ} \longrightarrow \Hom(\DD(I \times J), \DD(I)) \]
whence a composite
\[ J^{\circ} \longrightarrow \Hom(\DD(I \times J), \DD(I)) \]
or again
\[ \DD(I \times J) \longrightarrow \Hom(J^{\circ}, \DD(I)) \quad . \]
This holds for every $J \in \Ob \, \Dia$. We can therefore say that the datum of a ``coefficient'' $\Phi$ on a product category $I \times J$ gives rise to a diagram of type $J^{\circ}$ in the category $\DD(I)$ of coefficients on $I$. The philosophy
\msp{page 9}
of derivators consists in large part, conversely, whenever we have a diagram in a category of coefficients $\DD(I)$ (on a diagram $I$), say diagram of type $K$: $K \to \DD(I)$, in asking ourselves whether it does not come ``canonically'' from an object of $\DD(I \times K^{\circ})$ \textit{i.e.} from a coefficient on $I \times K^{\circ}$.

Thus for the square p. 8. We would like it to come from $\DD(I \times J)$, where $J = \Del_1 \times \Del_1$, and this functorially, so we want a ``canonical'' functor
\[ \DD(I) \longrightarrow \DD(I \times \Del_1 \times \Del_1) \]
which on composing with
\[ \DD(I \times \Del_1 \times \Del_1) \longrightarrow \Hom((\Del_1 \times \Del_1)^{\circ}, \DD(I)) \]
\printedpage{8}
% head: A. Grothendieck : Dérivateurs, Ch. I, Édition des Universités de Montpellier II et Paris VII
% foot: Version 8 juillet 2001
gives the functor described by square (1). Moreover, we must express in the present context the ``bicartesian'' condition on square (1), in the context of page 7. Analogous question for square (2) on page 2:
\[
\begin{tikzcd}
& F \arrow[dr] & \\
i_*i^!F \arrow[ur] \arrow[dr] & & j_*j^*F \\
& e \arrow[ur] &
\end{tikzcd}
\quad ,
\tag*{(2)}
\]
where we must in addition interpret the cartesian condition. We are in fact going to define a coefficient on the following open subcategory $W$ of $\Del_1 \times \Del_1$
\[
\begin{tikzcd}
(1,0) & \\
& (0,0) \arrow[ul] \arrow[dl] \\
(0,1) &
\end{tikzcd}
\quad ,
\]

\msp{page 10}

more precisely on $W \times I$, and if $\alpha$
\[ \alpha : W \times I \longrightarrow \Del_1 \times \Del_1 \times I \]
is the canonical inclusion, we will define the ``cartesian'' square as the image under $\alpha_*$ of this coefficient. This thus brings us back to the problem of first ``globalizing'' the partial diagram
\[
\begin{tikzcd}
F \arrow[dr] & \\
& j_*j^*F \\
e \arrow[ur] &
\end{tikzcd}
\quad .
\]
And likewise for square (1) except that we are at a loss if we choose to look at it mainly as cartesian (which suggests proceeding as just now) or as \underline{co}cartesian (which % verified transcription correction: the edition prints embarrassé, with double r


\endgroup
% END UNIT Der-ChI-234/Der-ChI-234_p006-010.tex
% BEGIN UNIT Der-ChI-234/Der-ChI-234_p011-015.tex
\begingroup\makeatletter%
\NoAutoSpacing%
\newcommand{\printedpage}[1]{\marginpar{\footnotesize p.~#1}}%
\newcommand{\msp}[1]{\marginpar{\footnotesize [#1]}}%
\newcommand{\DD}{\mathbb{D}}%
\newcommand{\uHom}{\underline{\mathsf{Hom}}}%
\newcommand{\Dia}{\mathsf{Dia}}%
\newcommand{\Ens}{\mathcal{E}ns}%
\newcommand{\bDelta}{\boldsymbol{\Delta}}%
\newcommand{\defeq}{\overset{\mathrm{def}}{=}}%
\makeatother\let\printedpage\derpp\let\msp\dermsp%
\let\mspin\DERmspin%

\DERmark{ENV0001}


\printedpage{9}
% head: A. Grothendieck : Dérivateurs, Ch. I, Édition des Universités de Montpellier II et Paris VII
\noindent suggests proceeding in the dual way). But the two coefficients one will find by these two procedures, products either of
\[
\begin{tikzcd}
F \arrow[dr] & \\
 & i_*i^*F \\
e \arrow[ur] &
\end{tikzcd}
\quad ,
\]
or of
\[
\begin{tikzcd}
 & F \\
j_!j^*F \arrow[ur] \arrow[dr] & \\
 & e
\end{tikzcd}
\quad ,
\]
should be \textit{canonically} isomorphic.

\noindent\msp{page 11}If $\DD$ is a derivator, the category
\[
\mathcal{A} = A_{\DD} = \DD(e) % edition reading verified: plain italic A_D is retained; this alphabet variation is not established to be an error
\]
plays an important role. It is called the \textit{base category,} or \textit{the fundamental category}, of the derivator, these objects are the \textit{fundamental coefficients.} For arbitrary $I$ in $\Dia$, we thus have for $F \in \DD(I)$, coefficient on $I$ % [sic] ces objets -> ses objets, possessive intended
\[
I^\circ \longrightarrow \mathcal{A} \quad , \qquad i \longmapsto F_i
\]
a diagram of type $I^\circ$ in $\mathcal{A}$, functorially, \textit{i.e.}
\[
\DD(I)\longrightarrow
\uHom(I^\circ,\underset{\substack{\Vert\\\DD(e)}}{\mathcal A})
\]
% the vertical equality sign is printed under the letter \mathcal{A} of \uHom(I^\circ,\mathcal{A}), with \DD(e) beneath it
(\textit{cf.} page 8).

\noindent\textbf{Der}$_3$ : For every $I \in \Dia$, this functor is \textit{conservative,}\footnote[4]{N.B. But this functor has no desire to be \textit{faithful.}} \textit{i.e.} for
\[
F \longrightarrow G
\]
\printedpage{10}%
% head: A. Grothendieck : Dérivateurs, Ch. I, Édition des Universités de Montpellier II et Paris VII
in $\DD(I)$ to be \texttt{[an]} isomorphism, it is necessary and sufficient that for every $i \in I$,
\[
F_i \longrightarrow G_i
\]
be one.

\noindent\textbf{Corollary.} \textit{The canonical functor (for $I$, $J \in \Dia$)}
\[
\DD(I \times J) \xrightarrow{\ \theta\ } \uHom(J^\circ, \DD(I))
\]
\textit{is conservative.}

Indeed, if $F$ is in $\DD(I \times J)$, we have for $(i,j) \in I \times J$
\[
F_{i,j} \simeq \theta(F)(j)_i \quad .
\]

\noindent\msp{page 12}Intuitively, one regards a derivator $\DD$ as ``being'' the fundamental category $\mathcal{A}$, with an additional ``structure'', \texttt{[to]} wit the datum of the ``extensions'' of $\mathcal{A}$ (as category of point coefficients) to the diagrams $I \in \Dia$, contravariantly in $I$. A coefficient $F$ in $\DD(I) \defeq \mathcal{A}_I$ is regarded as ``being'' the corresponding diagram $(F_i)_{i\in I}$ of fundamental objects, with an additional datum of ``globalization''. The $F_i$, or ``fibers'' of $F$ at the $i \in I$, are the local characteristics, of the ``global on $I$'' coefficient $F$. % [sic] caractéristiques locaux -> caractéristiques locales, agreement (feminine)

All known derivators are obtained in the following standard way. One starts from a ``category of models''
\[
\mathcal{M} \quad , \quad \Sigma
\]
equipped with a set of arrows, called ``equivalences'', and satisfying the relevant axioms for passing to a category of fractions
\[
\mathcal{A} = \mathcal{M}\Sigma^{-1}
\]
which will be the ``fundamental category''. One defines for every $I$ in $\Dia$
\[
\begin{aligned}
\mathcal{M}_I &= \uHom(I^\circ,\mathcal{M}) \\
\Sigma_I &= \{u \in \mathsf{Fl}(\mathcal{M}_I) \mid \forall i \in I, u_i \in \Sigma\} \\
\mathcal{A}_I &= \mathcal{M}_I\Sigma_I^{-1}
\end{aligned}
\]

\noindent\msp{page 13}It is clear that $\mathcal{A}_I \defeq \DD_{\mathcal{M},\Sigma}(I)$ depends on $I$ in a contravariant functorial way, and even in a 2-functorial way.

Of course, to really have a derivator, one must check the axioms. The strongest is the existence of the $f_!$, $f_*$. It is proved by Thomason when $(\mathcal{M},\Sigma)$ comes from a ``model category'' in Quillen's sense. It is not clear to me that then $(\mathcal{M}_I,\Sigma_I)$ also comes from a model category (in which for example the fibrations \printedpage{11}%
% head: A. Grothendieck : Dérivateurs, Ch. I, Édition des Universités de Montpellier II et Paris VII
and cofibrations would be defined from these data in $\mathcal{M}$, as earlier for the equivalences).

When one takes for $(\mathcal{M},\Sigma)$ a category of models for $\mathsf{Hot}$ (for example the semi-simplicial sets $\mathcal{M} = \uHom(\bDelta^\circ,\Ens)$, whence $\mathcal{M}_I \simeq \uHom(\bDelta^\circ \times I^\circ,\Ens) \simeq \uHom(\bDelta^\circ,\underbrace{\uHom(I^\circ,\Ens)}_{\widehat{I}})$, category of semi-simplicial objects in the topos $\widehat{I}$), one finds a description of a derivator (nonadditive), which will play a particularly important role. I denote it by % Resolved: exactly two closing parentheses in the formula; the source's outer prose parenthesis remains unclosed.
\DERmark{ENV0002}

\[
\mathsf{HOT} \quad ,
\]
its category of fundamental coefficients is $\mathsf{Hot}$, the standard homotopy category.

\noindent\msp{page 14}The notion of ``morphism of Derivators'' is clear, it is a morphism between two strict 2-functors
\[
\Dia^\circ \longrightarrow \mathcal{C}at
\]
Thus
\[
u : \DD \longrightarrow \DD'
\]
is defined by giving, for every $I \in \Dia$,
\[
u_I : \DD(I) \longrightarrow \DD'(I)
\]
with commutativity (strict?) for $f : I \to J$
\[
\begin{tikzcd}
\DD(J) \arrow[r, "f^*_{\DD}"] \arrow[d, "u_J"'] & \DD(I) \arrow[d, "u_I"] \\
\DD'(J) \arrow[r, "f^*_{\DD'}"] & \DD'(I)
\end{tikzcd}
\quad . \tag*{(1)}
\]
We have
\[
\begin{tikzcd}[column sep=large]
\uHom(I,J)^\circ \arrow[r, "\DD_{I,J}"] \arrow[d, "\DD'_{I,J}"'] & \uHom(\DD(J),\DD(I)) \arrow[d, "\phi \mapsto u_I \circ \phi"] \\
\uHom(\DD'(J),\DD'(I)) \arrow[r, "\psi \mapsto \psi \circ u_J"'] & \uHom(\DD(J),\DD'(I))
\end{tikzcd}
\tag*{(2)}
\]
and the commutativity of (1) means the commutativity of (2) \textit{for the objects}. But of course we want it for the arrows too, \textit{i.e.} we want the commutativity of diagram (2).\footnote[5]{Give the example of a morphism of derivators defined by a functor of categories of models $(\mathcal{M},\Sigma) \to (\mathcal{M}',\Sigma')$.}

Probably we want a little more from morphisms of derivators - we want commutation with the $f_*$ and (?) with the $f_!$ (this time only up to canonical isomorphism) \ldots\footnote[6]{Commutation with $f_*$ (resp. $f_!$) \textit{left exact} (resp. \textit{right exact}) morphism. The \textit{exact} morphisms are those that are both left exact and right exact.}

\printedpage{12}
% head: A. Grothendieck : Dérivateurs, Ch. I, Édition des Universités de Montpellier II et Paris VII
\noindent\msp{page 15}For example, consider the derivator
\[
\mathsf{HOTAB}
\]
whose fundamental category is $\mathsf{Hotab}$, the full subcategory of the derived category $\mathsf{Der}(\mathsf{Ab})$ of $\mathsf{Ab}$, consisting of the complexes with vanishing cohomology in degree $> 0$. One can describe $\mathsf{HOTAB}$ by the category of models $\uHom(\bDelta^\circ,\mathsf{Ab})$ of semi-simplicial abelian groups, with as equivalences the quasi-isomorphisms in the usual sense. The canonical ``abelianization functor''
\[
\mathsf{Hot} \longrightarrow \mathsf{Hotab}
\]
extends to a morphism of derivators
\[
\mathsf{HOT} \longrightarrow \mathsf{HOTAB} \quad .
\]
The question of its compatibility with the $f_*$ does not seem tautological. It would imply, for example, that if
\[
\begin{tikzcd}
X \arrow[d, "f"'] & X' \arrow[l, "g"'] \arrow[d, "f'"] \\
Y & Y' \arrow[l, "j"]
\end{tikzcd}
\]
is a cartesian diagram in $\mathcal{M} = \uHom(\bDelta^\circ,\Ens)$ with $f$ a fibration (of Thom in the setting of topological spaces, with $f$ for example a locally trivial fibration) then the diagram in $\mathsf{Hotab}$
\[
\begin{tikzcd}
\mathbb{R}\Gamma_!(X) \arrow[d] & \mathbb{R}\Gamma_!(X') \arrow[l] \arrow[d] \\
\mathbb{R}\Gamma_!(Y) & \mathbb{R}\Gamma_!(Y') \arrow[l]
\end{tikzcd}
\]
is $\mathsf{hot}$-cartesian\footnote[7]{This is false already if $Y$ is the one-point space : trivially the $\mathbb{R}\Gamma_!$ does not transform products into sums! This functor is surely \texttt{[not]} ``left exact'' (as a morphism of derivators).}. And likewise for a $\mathsf{hot}$-amalgamated sum ($\mathsf{hot}$-cocartesian square in $\mathsf{Hot}$). % [sic] carrée -> carré, agreement (masculine noun)

\noindent\msp{page 16}If this is true, and not known to boot, it would be scandalous all the same!

Proceeding as for the base change morphisms for the $Rf_*$, we find in any case a canonical morphism, for $f : I \to J$
\[
u_Jf_*^{\DD} \longrightarrow f_*^{\DD'}u_I
\]
\printedpage{13}%
% head: A. Grothendieck : Dérivateurs, Ch. I, Édition des Universités de Montpellier II et Paris VII
and likewise
\[
f_!^{\DD'}u_I \longrightarrow u_Jf_!^{\DD} \quad ,
\]
and the question is whether these are isomorphisms. Analogous questions for the $f^!$, for $f$ a closed immersion (or for $f^?$, for $f$ an open immersion). Since the notion of a morphism of derivators ought all the same to ``fit'' the well-known notion of a (triangulated) morphism of triangulated categories\footnote[8]{No! The case of triangulated categories is ``misleading'', because for these left and right exactness amount to the same thing.}, hence be exact in a suitable sense, this does show that one must require a morphism of derivators to commute with the $f_*$ and with the $f_!$.\footnote[9]{N.B. One can conceive of a functor between derivators being left exact or right exact or both.} I should ask Thomason whether he has looked, for a morphism
\[
\mathcal{M} \longrightarrow \mathcal{M}'
\]
between categories of models, which is say \textit{left exact}, and ``compatible'' with the equivalences, fibrations (and cofibrations (???)).

\noindent\textsc{Notation} : $\mathsf{DER}$ for the 2-category of \textit{derivators.}

\noindent\msp{page 17}\textbf{Duality}. If $\DD$ is a derivator, we are going to define a \textit{dual} derivator $\DD^\circ$ by
\[
\DD^\circ(I) \defeq \DD(I^\circ)^\circ \quad .
\]
The coefficients on $I$ of type $\DD^\circ$ are the coefficients on the category $I^\circ$ opposite to $I$ of type $\DD$, and as a category it is the category \textit{opposite} to that of the coefficients of type $\DD$ on $I^\circ$. To define the 2-functorial law, we note that for $I$, $J$ in $\Dia$ we have % [sic] les coefficient -> les coefficients, plural agreement
\[
\begin{tikzcd}[column sep=large]
\uHom(I,J)^\circ \arrow[r, "\DD_{I,J}"] \arrow[d, "\wr"', "\text{duality}"] & \uHom(\DD(J),\DD(I)) \arrow[d, "\wr"', "\text{duality}"] \\
\uHom(I^\circ,J^\circ) \arrow[r, "\theta"] & \uHom(\DD(J)^\circ,\DD(I)^\circ)^\circ
\end{tikzcd}
\quad ,
\]
where $\theta$ is the unique functor making the diagram commutative. Then $\theta^\circ$is a functor
\[
\uHom(I^\circ,J^\circ)^\circ \longrightarrow \uHom(\DD(J)^\circ,\DD(I)^\circ) \quad ,
\]
that is by definition $\DD^\circ_{I,J}$ \texttt{[rather} $\DD^\circ_{I^\circ,J^\circ}$\texttt{]}.

If $\DD$ is defined by $(\mathcal{M},\Sigma)$, then $\DD^\circ$ by $(\mathcal{M}^\circ,\Sigma^\circ)$, since
\[
(\mathcal{M}^\circ)_I \defeq \uHom(I^\circ,\mathcal{M}^\circ) \simeq \uHom(I,\mathcal{M})^\circ \simeq (\mathcal{M}_{I^\circ})^\circ
\]
% foot (pages 9--13): Version 8 juillet 2001


\endgroup
% END UNIT Der-ChI-234/Der-ChI-234_p011-015.tex
% BEGIN UNIT Der-ChI-234/Der-ChI-234_p016-020.tex
\begingroup\makeatletter%
\newcommand{\printedpage}[1]{\marginpar{\footnotesize p.~#1}}%
\newcommand{\msp}[1]{\marginpar{\footnotesize [#1]}}%
\newcommand{\D}{\mathbb{D}}%
\newcommand{\uHom}{\underline{\mathsf{Hom}}}%
\newcommand{\sHom}{\mathsf{Hom}}%
\newcommand{\Dia}{\mathsf{Dia}}%
\newcommand{\DER}{\mathsf{DER}}%
\newcommand{\id}{\mathsf{id}}%
\newcommand{\Cat}{\mathcal{C}at}%
\newcommand{\Ens}{\mathcal{E}ns}%
\newcommand{\sIm}{\mathsf{Im}}%
\makeatother\let\printedpage\derpp\let\msp\dermsp%
\let\mspin\DERmspin%


\printedpage{14}
% head: A. Grothendieck : Dérivateurs, Ch. I, Édition des Universités de Montpellier II et Paris VII
% foot: Version 8 juillet 2001

and $\Sigma^\circ{}_I$ is deduced from $\Sigma_{I^\circ}$, whence
\[
(\mathcal{M}^\circ)_I(\Sigma^\circ{}_I)^{-1} \simeq (\mathcal{M}_{I^\circ}\Sigma_{I^\circ}^{-1})^\circ = \D(I^\circ)^\circ \quad .
\]

\msp{page 18}

\textbf{Philosophy.} Each time we have a functor
\[
\mathcal{A} \xrightarrow{u_0} \mathcal{A}'
\]
between two categories which can be naturally interpreted as the fundamental categories of two derivators $\D$, $\D'$, the question arises of interpreting $u_0$ as coming from a morphism of derivators
\[
\D \xrightarrow{u} \D' \quad .
\]
The latter is, in practice, always canonically determined, up to canonical isomorphism. (But I nevertheless doubt that the restriction functor to the fundamental categories
\[
\uHom(\D,\D') \longrightarrow \uHom(\mathcal{A},\mathcal{A}')
\]
is fully faithful \dots\ )

We are going to apply this philosophy to the functors $f^*$, $f_!$, $f_*$ (and possibly $f^!$, $f^?$), for $f : I \to J$, thus to the functors between $\D(I)$ and $\D(J)$. This therefore leads us to regard $\D(I)$ and $\D(J)$ as fundamental categories of suitable derivators.

\textbf{Induced derivator} by a derivator $\D$, on an index category $I \in \Dia$. It is a derivator $\D_I$ defined by
\[
\D_I(J) = \D(I \times J) \quad .
\]
The functor associated with $\D_I$, for two categories $J$, $J' \in \Dia$

\msp{page 19}
\[
\uHom(J,J')^\circ \xrightarrow{(\D_I)_{J,J'}} \uHom(\D_I(J'),\D_I(J)) = \uHom(\D(I\times J'),\D(I\times J))
\]
is the composite
\[
\uHom(J,J')^\circ \longrightarrow \uHom(I\times J, I\times J')^\circ \xrightarrow{\D_{I\times J,I\times J'}} \uHom(\D(I\times J'),\D(I\times J)) \quad ,
\]
where the first arrow is the functor
\[
f \longmapsto \id_I \times f \quad .
\]
In other words, $\D_I$ is the composite
\[
\begin{aligned}
&\Dia^\circ \longrightarrow \Dia^\circ \xrightarrow{\D} \Cat \\
&J \longmapsto I \times J \qquad .
\end{aligned}
\]
We have
\[
f^*_{\D_I} = (\id \times f)^*_{\D} \quad .
\]
If the $f_*$ (resp. the $f_!$) exist for $\D$, they exist for $\D_I$. (Likewise for every kind of multiadjoints.) Since $J \mapsto I \times J$ takes closed immersion (resp. open) to likewise, axiom \textbf{Der}$_2$ (more precisely, each of these ingredients, \textit{cf.} marginal annotations p.~7) is stable under passage from $\D$ to $\D_I$ . Likewise for \textbf{Der}$_3$. (For each new axiom we bring out, we will have to check stability.)

\printedpage{15}
% head: A. Grothendieck : Dérivateurs, Ch. I, Édition des Universités de Montpellier II et Paris VII
% foot: Version 8 juillet 2001

\msp{page 20}

Let
\[
f : I \longrightarrow I' \quad .
\]
Consider between $\D(I)$, $\D(I')$ the functors
\[
\begin{tikzcd}[column sep=huge]
\D(I) \arrow[r, shift left=3ex, "f_!"] \arrow[r, shift right=3ex, "f_*"'] & \D(I') \arrow[l, "f^*"']
\end{tikzcd}
\]
(subject to the existence of the $f_*$ resp. $f_!$). Can we interpret them in terms of morphisms of derivators
\[
\begin{tikzcd}[column sep=huge]
\D_I \arrow[r, shift left=3ex, "\overline{f_!}"] \arrow[r, shift right=3ex, "\overline{f_*}"'] & \D_{I'} \arrow[l, "\overline{f^*}"']
\end{tikzcd}
\qquad ?
\]
So we must define, for $J$ in $\Dia$, corresponding functors between $\D_I(J)$ and $\D_{I'}(J)$, \textit{i.e.} between $\D(I \times J)$ and $\D(I' \times J)$. Now we do have
\[
f \times \id_J : I \times J \longrightarrow I' \times J
\]
and we will set (subject to the existence of the last two)
\[
\left\{
\begin{aligned}
&\overline{f^*}(J) = (f \times \id_J)^*_{\D} \\[1ex]
&\overline{f_!}(J) = (f \times \id_J)^{\D}_! \\[1ex]
&\overline{f_*}(J) = (f \times \id_J)^{\D}_* \quad .
\end{aligned}
\right.
\]

\msp{page 21}

In other words, we can regard
\[
\Dia \ni I \longmapsto \D_I \quad , \qquad \D \in \DER \quad ,
\]
as functorial in $I$ in three ways (one contravariant, the other two covariant), according as we use the $f^*$, the $f_!$ (if they exist) or the $f_*$ (if they exist). We will denote by
\[
f^{\D}_! \quad , \qquad f^*_{\D} \quad , \qquad f^{\D}_*
\]
indifferently the functors between $\D(I)$ and $\D(I')$ or (more finely) between $\D_I$ and $\D_{I'}$, from which the preceding ones are deduced by passing to the corresponding fundamental categories. Likewise for $f^!$ or for $f^?$ subject to existence. For example, if we assume that the $f^!$ exist for $\D$, for every $f$ a closed immersion, then for such an $f$ we find a morphism of derivators
\[
\D_I \xleftarrow{f^!_{\D}} \D_{I'} \quad .
\]
We must see whether all these morphisms of derivators are exact, \textit{i.e.} whether for
\[
g : J \longrightarrow J'
\]
in $\Dia$, they commute with $g_!$, $g_*$.

\printedpage{16}
% head: A. Grothendieck : Dérivateurs, Ch. I, Édition des Universités de Montpellier II et Paris VII
% foot: Version 8 juillet 2001

\msp{page 22 blank}

\msp{page 23}

To make this question explicit, we are led to look at the diagram of categories
\[
\begin{tikzcd}[column sep=huge, row sep=huge]
I \times J \arrow[r, "G = \id_I \times g"] \arrow[d, "F = f \times \id_J"'] & I \times J' \arrow[d, "F' = f \times \id_{J'}"] \\
I' \times J \arrow[r, "G' = \id_{I'} \times g"] & I' \times J'
\end{tikzcd}
\qquad ,
\]
which is moreover a cartesian square in $\Cat$. The commutativity of $f^*_{\D}$ with the $g_*$ is then made explicit by the commutativity in
\[
\begin{tikzcd}[column sep=huge, row sep=huge]
\D(I \times J) \arrow[r, "G_*", ""{name=T, below, near end}] & \D(I \times J') \\
\D(I' \times J) \arrow[u, "F^*"] \arrow[r, "G'_*"] & \D(I' \times J') \arrow[u, "{F'}^*"', ""{name=R, left}]
\arrow[dashed, from=R, to=T, "\text{iso?}"]
\end{tikzcd}
\qquad .
\]
This commutativity will be ensured by later axioms on derivators. (cf. further % [sic] ultérieures -> ultérieurs, accord avec « axiomes »
below on page 59). Thus, $f^*_{\D}$ is right exact \texttt{[rather left]}, and dually, it must also be left exact \texttt{[rather right]}. So $f^*_D : \D_{I'} \to \D_I$ is \textit{exact.} Consider $f^{\D}_*$, % [sic] f^*_D -> f^*_{\mathbb{D}}, indice imprimé en D italique au lieu de \mathbb{D}
this time it is the commutativity of
\[
\begin{tikzcd}[column sep=huge, row sep=huge]
\D(I \times J) \arrow[r, "G_*"] \arrow[d, "F_*"'] & \D(I \times J') \arrow[d, "F'_*"] \\
\D(I' \times J) \arrow[r, "G'_*"] & \D(I' \times J')
\end{tikzcd}
\qquad ,
\]
it is trivial (up to canonical isomorphism). Thus, we
\msp{page 24}
find that $f^{\D}_* : \D_I \to \D_{I'}$ is \textit{left exact}, and dually $f^{\D}_!$ is \textit{right exact.} But is $f^{\D}_*$ also \textit{right} exact? This therefore means commutativity of
\[
\begin{tikzcd}[column sep=huge, row sep=huge]
\D(I \times J) \arrow[r, "G_!"] \arrow[d, "F_*"'] & \D(I \times J') \arrow[d, "F'_*", ""{name=R2, left}] \\
\D(I' \times J) \arrow[r, "G'_!", ""{name=B, above, near end}] & \D(I' \times J')
\arrow[dashed, from=B, to=R2]
\end{tikzcd}
\qquad .
\]

\printedpage{17}
% head: A. Grothendieck : Dérivateurs, Ch. I, Édition des Universités de Montpellier II et Paris VII
% foot: Version 8 juillet 2001

We do manage to define
\[
G'_!F_* \xrightarrow{?} F'_*G_!
\]
or what comes to the same thing by adjunction
\[
F_* \xrightarrow{?} {G'}^*F'_*G_!
\]
using the base change morphism
\[
{G'}^*F'_* \xrightarrow{\overset{\alpha}{\sim}} F_*G^*
\]
which is \texttt{[an]} \textit{isomorphism} (we have seen this), whence
\[
\begin{tikzcd}[column sep=large, row sep=normal]
& {G'}^*F'_*G_! \arrow[dd, "\wr"', "\alpha * G_!"] \\
F_* \arrow[ur, "?"] \arrow[dr, "F_* * \beta"'] & \\
& F_*G^*G_!
\end{tikzcd}
\]
where $\beta : \id \to G^*G_!$ is the adjunction morphism, whence the desired morphism. (We can also look for ${F'}^*(G'_!F_*) \to G_!$, and use the commutation $G_!F^* \overset{\sim}{\to} {F'}^*G'_!$, \textit{i.e.} the \textit{right} exactness of $f^*_{\D}$.)

\msp{page 25}

But is this morphism an isomorphism? I think that in the additive case it may be reasonable, but in the nonadditive case there is no reason for it to be true. We see in any case that if for every $f : I \to I'$ in $\Dia$, $f^{\D}_* : \D_I \to \D_{I'}$ is (not only left exact but also) right exact, then for every $g : J \to J'$ in $\Dia$, $g^{\D}_!$ is (not only right exact but even) left exact. % [sic] exacte -> exact, accord avec « $g^{\D}_!$ est (non seulement exact\dots) »

Let us take the derivator of presheaves of sets on variable $I$'s, its fundamental category is $\Ens$. It is a nonpointed derivator. Consider $I \to I'$ a nonsurjective closed immersion, $f_* : \D(I) \to \D(I')$. If $f^{\D}_*$ were right exact $f_*$ would have to take initial object of $\D(I)$ to initial object of $\D(I')$, unless I am mistaken (\textit{cf.} below). But this is not so, since for every $F \in \sIm\bigl(f^{\D}_* : \D(I) \to \D(I')\bigr)$ the restriction of $F$ to $I' \setminus I$ is the final sheaf $e_U = (i \mapsto e)$, and not the initial sheaf $\emptyset_U = (i \mapsto \emptyset)$.

We are going to prove this:

\textbf{Proposition.} \textit{If $\D$ is a left (right) derivator, then finite products (finite sums) exist in the $\D(I)$. Every morphism $\varphi : \D \to \D'$ of derivators which is left (right) exact commutes with the said products (sums), each $I$ giving $\varphi(I) : \D(I) \to \D'(I)$. In particular, it takes final (initial) object to likewise.}

\printedpage{18}
% head: A. Grothendieck : Dérivateurs, Ch. I, Édition des Universités de Montpellier II et Paris VII
% foot: Version 8 juillet 2001

\msp{page 26}

Saying that finite products indexed by the finite set $I$ exist in the category $\mathcal{A} = \D(e)$ (say) is equivalent to saying that the diagonal functor
\[
\mathcal{A} \xrightarrow{\Delta} \mathcal{A}^I
\]
admits a right adjoint. Now if
\DERmark{ENV0003}

\[
f : I \to e
\]
is the projection of $I$ (regarded as a discrete ordered set) onto the one-point ordered set $e$, the functor $\Delta$ is identified with $f^* : \D(e) \to \D(I)$ (thanks to axiom \textbf{Der}$_2$ on the behavior of the derivator for a discrete sum of index categories). So the existence of the product amounts to that of the right adjoint $f_*$ of $f^*$, OK for left derivator. And a left exact homomorphism of left derivators, since it commutes with the $f_*$, commutes with finite products. Likewise for finite sums, and right derivators.

The preceding example uses the fact that the derivator under consideration is not pointed. But even for a pointed derivator $\D$, a morphism $f^{\D}_* : \D_I \to \D_{I'}$ is not necessarily right exact. If we take the derivator associated with an arbitrary category $\mathcal{M}$ (with $\Sigma = \emptyset$),
\[
\D_{\mathcal{M}}(I) = \uHom(I^\circ,\mathcal{M}) \overset{\mathrm{def}}{=} \mathcal{M}(I)
\]
\DERmark{ENV0004}

then the induced derivator
\[
(\D_{\mathcal{M}})_I = \{J \mapsto \D_{\mathcal{M}}(I \times J) = \uHom((I \times J)^\circ,\mathcal{M}) \simeq \uHom(J^\circ,\uHom(I^\circ,\mathcal{M}))\}
\]

\msp{page 27}

is identified with
\[
\D_{\uHom(I^\circ,\mathcal{M})} = \D_{\mathcal{M}(I)}
\]
\textit{i.e.}
\[
(\D_{\mathcal{M}})_I = \D_{\mathcal{M}(I)}
\]
N.B. Here $\mathcal{M}$ and $\D_{\mathcal{M}}$ determine each other and the functor
\[
\mathcal{M} \longmapsto \D_{\mathcal{M}}
\]
\[
\Cat \longrightarrow \DER
\]
is 2-faithful, \textit{i.e.} if $\mathcal{M}$, $\mathcal{M}'$ are two categories
\[
\uHom(\mathcal{M},\mathcal{M}') \xrightarrow{\simeq} \sHom(\D_{\mathcal{M}},\D_{\mathcal{M}'})
\]
is an equivalence of categories. On the other hand, we check at once that $u : \mathcal{M} \to \mathcal{M}'$ is left exact (or right exact) if and only if $\D_u : \D_{\mathcal{M}} \to \D_{\mathcal{M}'}$ is so.

Thus, to see whether a morphism of derivators of the form
\[
\begin{tikzcd}[column sep=large, row sep=normal]
f^{\D}_* : &[-2.5em] (\D_{\mathcal{M}})_I \arrow[r] \arrow[d, equal] & (\D_{\mathcal{M}})_{I'} \arrow[d, equal] \\
& \D_{\mathcal{M}(I)} & \D_{\mathcal{M}(I')}
\end{tikzcd}
\]
is \textit{right} exact, it comes to the same thing to say


\endgroup
% END UNIT Der-ChI-234/Der-ChI-234_p016-020.tex
% BEGIN UNIT Der-ChI-234/Der-ChI-234_p021-025.tex
\begingroup\makeatletter%
\newcommand{\printedpage}[1]{\marginpar{\footnotesize p.~#1}}%
\newcommand{\msp}[1]{\marginpar{\footnotesize [#1]}}%
\newcommand{\sDia}{\mathsf{Dia}}%
\newcommand{\sOb}{\mathsf{Ob}}%
\newcommand{\sid}{\mathsf{id}}%
\newcommand{\DD}{\mathbb{D}}%
\makeatother\let\printedpage\derpp\let\msp\dermsp%
\let\mspin\DERmspin%


\printedpage{19}
% head: A. Grothendieck : Dérivateurs, Ch. I, Édition des Universités de Montpellier II et Paris VII
% foot: Version 8 juillet 2001
\msp{page 28}
that the functor for the associated fundamental categories
\[ \mathcal{M}(I) \xrightarrow{f^{\mathcal{M}}_*} \mathcal{M}(I') \]
is right exact. Now here it is a direct image functor of ordinary presheaves, and this functor has no desire at all to be \textit{right} exact. (Even if $\mathcal{M}$ is \textit{abelian}, for that matter!) Let us take for example for $f : I \hookrightarrow I'$ an open immersion with $I' = \boldsymbol{\Delta}_1 \times \boldsymbol{\Delta}_1$, $I = I' - \{c\}$, $c$ a final object of $I'$
\[
I' = \begin{tikzcd}
c & b \arrow[l] \\
a \arrow[u] & 0 \arrow[l] \arrow[u]
\end{tikzcd}
\qquad\qquad
I = \begin{tikzcd}
 & b \\
a & 0 \arrow[l] \arrow[u]
\end{tikzcd}
\]
The functor $f_*$ is the one which associates to a diagram
\[
F = \left\{
\begin{tikzcd}
 & F_b \arrow[d] \\
F_a \arrow[r] & F_0
\end{tikzcd}
\right\}
\]
the completed cartesian square
\[
\overline{F} = \left\{
\begin{tikzcd}
F_a \times_{F_0} F_b \arrow[r] \arrow[d] & F_b \arrow[d, "\beta"] \\
F_a \arrow[r, "\alpha"] & F_0
\end{tikzcd}
\right\}
\]
Now if $F \mapsto \overline{F}$ were right exact, $F \mapsto F_a \times_{F_0} F_b$ would be so too. Now $F_a \times_{F_0} F_b = K(F) = \mathsf{Ker}(F_a \times F_b \xrightarrow{\alpha pr_1 - \beta pr_2} F_0)$, as an additive functor in $F$, is not right exact, for example it does not necessarily transform epimorphism into epimorphism. Let us take for example
\[
F' = \left\{
\begin{tikzcd}
 & F_b \arrow[d] \\
F_a \arrow[r] & 0
\end{tikzcd}
\right\}
\]
and the obvious homomorphism $F \to F'$ which is the identity on
\msp{page 29}
$F_a$ and on $F_b$. It is an epimorphism and the corresponding homomorphism $\overline{F} \to \overline{F'}$ is the inclusion
\[ F_a \times_{F_0} F_b \longrightarrow F_a \times F_b \]
which is an epimorphism only if $\alpha = \beta = 0$.

\printedpage{20}
% head: A. Grothendieck : Dérivateurs, Ch. I, Édition des Universités de Montpellier II et Paris VII
% foot: Version 8 juillet 2001
\textbf{Perplexity} : One should be able to define a notion of an (additive) ``\textit{triangulated}'' derivator (which would correspond to the intuition of triangulated categories; it is a derivator satisfying fairly strong axioms), so that for morphisms between these, left exactness and right exactness are \textit{equivalent.}

\textbf{``Base change theorem''} for derivators. Let $\DD$ be a left derivator (existence of $f_*$) and let us consider
\[ f : I \longrightarrow J \qquad \text{in } \sDia \quad , \]
let $j_0 \in \sOb\ J$, we propose to ``calculate'' the functor
\[ F \mapsto f_*(F)_{j_0} = \alpha^*_{j_0}(f_*(F)) \quad , \]
where
\[ \alpha_{j_0} : e \longrightarrow J \]
is the inclusion of the one-point category $e$ in $J$, with value $j_0$.

\msp{page 30}
Let us consider the diagram
\[
\begin{tikzcd}[row sep=large, column sep=large]
I/j_0 \arrow[r, "q"] \arrow[d, "f'"'] & I \arrow[d, "f"] \\
e \arrow[r, "p = \alpha_{j_0}"'] & J
\arrow[from=1-2, to=2-1, dashed, shorten <=14pt, shorten >=14pt, "\gamma"']
\end{tikzcd}
\quad ,
\tag*{(1)}
\]
where $I/j_0$ is the category $I \times_J J/j_0$ of pairs
\[ \{(i,\varphi) \mid i \in \sOb\ I,\ \varphi : f(i) \to j_0\} \]
and $q$ is the functor $(i,\varphi) \mapsto i$. This diagram is not commutative, but there is a canonical functor \texttt{[rather morphism of functors]}
\[ fq \xrightarrow{\gamma} pf' \quad , \]
where $fq$ is the functor $(i,\varphi) \mapsto f(i)$ and $pf'$ the constant functor $(i,\varphi) \mapsto j_0$ and where $\gamma(i,\varphi) : f(i) \to j_0$ is defined by $\gamma(i,\varphi) = \varphi$.

More generally, let us consider a diagram in $\sDia$
\[
\begin{tikzcd}[row sep=large, column sep=large]
I' \arrow[r, "q"] \arrow[d, "f'"'] & I \arrow[d, "f"] \\
J' \arrow[r, "p"'] & J
\arrow[from=1-2, to=2-1, dashed, shorten <=14pt, shorten >=14pt, "\gamma"']
\end{tikzcd}
\quad ,
\tag*{(2)}
\]
\printedpage{21}
% head: A. Grothendieck : Dérivateurs, Ch. I, Édition des Universités de Montpellier II et Paris VII
% foot: Version 8 juillet 2001
with $\gamma : fq \to pf'$. We then define a \textit{base change morphism for direct images} relative to this diagram
\[ p^* f_* \xrightarrow{c_\gamma} f'_* q^* \]
\msp{page 31}
in two ways, which give the same result (general story of adjoint functors to be made explicit once and for all, in terms of a diagram of functors
\[
\begin{tikzcd}[row sep=large, column sep=large]
X' & X \arrow[l, "q^*"'] \\
Y' \arrow[u, "{f'}^*"] & Y \arrow[l, "p^*"'] \arrow[u, "f^*"']
\arrow[from=2-1, to=1-2, dashed, shorten <=14pt, shorten >=6pt]
\end{tikzcd}
\]
assuming that the right adjoint functors $f'_*$, $f_*$ of the two vertical functors exist). An interesting case is the one where
\DERmark{ENV0005}

\[ pf' = fq \quad , \qquad \gamma = \sid_{fq} \]
and where we put $c_\gamma = c$.

Here
\[ p^* f_* \xrightarrow{?} f'_* q^* \]
amounts to
\[
\begin{tikzcd}[column sep=small]
f_* \arrow[r, "?"] & p_* f'_* q^* \simeq & (pf')_* q^* \\
 & & (fq)_* q^* \simeq (f_* q_*) q^* = f_*(q_* q^*) \quad . \arrow[u, "\gamma_* * q^*"']
\end{tikzcd}
\]
Now we have
\[ \sid \xrightarrow{\alpha} q_* q^* \qquad\qquad \text{(adjunction)} \quad , \]
whence
\[ f_* \xrightarrow{f_* * \alpha} f_*(q_* q^*) \]
and composing with $\gamma * q^*$, we find what we want. And instead of this we can look for
\[
\begin{tikzcd}
{f'}^*(p^* f_*) \arrow[r, "?"] \arrow[d, equal] & q^* \\
(pf')^* f_* \arrow[d, "\gamma^* * f_*"] & \\
(fq)^* f_* \simeq q^*(f^* f_*) \arrow[r, "q^* * \gamma'"] & q^*
\end{tikzcd}
\quad ,
\]
where $\gamma' : f^* f_* \to \sid$ is the adjunction.

\printedpage{22}
% head: A. Grothendieck : Dérivateurs, Ch. I, Édition des Universités de Montpellier II et Paris VII
% foot: Version 8 juillet 2001
\msp{page 32}
\textbf{Der}$_4$ : The base change morphism for direct images $f_*$ associated with a diagram (1) is always an isomorphism (left derivators). Dual condition for right derivators and base change for direct images $f_!$, for the diagram
\[
\begin{tikzcd}[row sep=large, column sep=large]
j_0 \backslash I \arrow[r] \arrow[d] & I \arrow[d, "f"] \\
e \arrow[r, "p = \alpha_{j_0}"'] & J
\arrow[from=2-1, to=1-2, dashed, shorten <=22pt, shorten >=8pt, "\gamma"]
\end{tikzcd}
\tag*{(1$'$)}
\]
associated with $f : I \to J$ in $\sDia$, and $j_0 \in \sOb\ J$.

\textbf{Der}$_5$\footnote[10]{\underline{Totally} botched, becomes reasonable only for constant coefficients on $F$.} : Let $f : I \to J$ be a \textit{strict hot-equivalence, i.e.} $\forall j_0 \in \sOb\ J$, $I/j_0$ is contractible. Then $f^*$ is fully faithful, \textit{i.e.} the adjunction morphism
\[ \sid_{\DD(J)} \longrightarrow f_* f^* \]
is an isomorphism ($\DD$ a left derivator). Dually, if $f$ is a co-strict equivalence, \textit{i.e.} $f^\circ : I^\circ \to J^\circ$ is a strict equivalence, \textit{i.e.} $\forall j_0 \in \sOb\ J$, $j_0 \backslash I$ is contractible, then (if $\DD$ is a right derivator) the functor $f^*$ is fully faithful, \textit{i.e.} the adjunction morphism % Resolved: the edition line-break gives co-stricte; retain that spelling.
\[ f_! f^* \longrightarrow \sid_{\DD(J)} \]
is an isomorphism.

N.B. Thus the functors $f_*$ (resp. $f_!$) are localization functors.

\textbf{N.B.} Pages 33 to 40 have been deleted, because botched. % [sic] suprimées -> supprimées, spelling

\msp{page 41}
I am going to take up again the terminology and the axioms relating to derivators.

A \textit{prederivator} is a 2-functor
\[ \sDia^\circ \xrightarrow{\DD} \mathcal{C}at \quad , \]
where $\sDia$ is a full 2-subcategory of $\mathcal{C}at$.\footnote[11]{$\sDia$ \textit{domain} of the derivator. The objects are the \textit{index categories} for $\DD$. The $\DD(I)$ are the categories of coefficients for $\DD$, the objects of $\DD(I)$ the ``coefficients of type $\DD$ on $I$''. % [sic] coeficients -> coefficients, spelling
} We assume at least that $\sDia$ contains the \textit{finite ordered sets,} that $\sDia$ is stable under finite $\varprojlim$, under finite sums. That with every category $C$ and every strictly full open (resp. closed) subcategory it contain the complementary subcategory.\footnote[12]{Notation $f^*$ : inverse image of coefficients.}

\printedpage{23}
% head: A. Grothendieck : Dérivateurs, Ch. I, Édition des Universités de Montpellier II et Paris VII
% foot: Version 8 juillet 2001
\msp{page 42}
We define the prederivator $\DD_{\mathcal{M},\Sigma}$ associated with a system
\[ (\mathcal{M},\Sigma) \qquad\qquad \mathcal{M} \text{ Category,} \quad \Sigma \subset \mathsf{Fl}(\mathcal{M}) \quad . \] % Resolved: Fl abbreviates fleches, with lowercase l; retain Fl.
If $\Sigma = \emptyset$, we write $\DD_{\mathcal{M}}$.

Prederivators form a 2- category $\mathsf{PREDER}$, and
\[ (\mathcal{M},\Sigma) \longmapsto \DD_{\mathcal{M},\Sigma} \]
is a 2-functor
\[ \mathsf{PREMOD}
\begin{pmatrix}
\text{``Category of categories} \\
\text{of models'' (weak)}
\end{pmatrix}
\longrightarrow \mathsf{PREDER} \quad . \]
For $\DD$ given, we define the induced derivator ($I \in \sDia$)
\[ \DD_I = \big(J \longmapsto \DD(I \times J)\big) \quad , \]
depending on $I$ contravariantly, whence a 2-functor
\[ \begin{aligned}
\sDia^\circ &\longrightarrow \mathsf{PREDER} \\
I &\longmapsto \DD_I
\end{aligned} \]
or again a bi-2-functor
\[ \begin{aligned}
\mathsf{PREDER} \times \sDia^\circ &\longrightarrow \mathsf{PREDER} \\
(\DD, I) &\longmapsto \DD_I \quad .
\end{aligned} \]
If $\DD$ is defined by $(\mathcal{M},\Sigma)$, then $\DD_I$ is defined by $(\mathcal{M}(I),\Sigma(I))$ :
\[ (\DD_{\mathcal{M},\Sigma})_I \simeq \DD_{\mathcal{M}(I),\Sigma(I)} \quad , \]
\msp{page 43}
where $\mathcal{M}(I) = \underline{\mathsf{Hom}}(I^\circ,\mathcal{M})$ (presheaves on $I$ with values in $\mathcal{M}$), and $\Sigma(I) \subset \mathsf{Fl}\ \mathcal{M}(I)$ is deduced from $\Sigma$ by looking fiber by fiber (now assuming $\Sigma$ saturated). % Resolved: the repeated arrow-class symbol is Fl, with lowercase l.

For any axioms we shall single out for certain prederivators, we shall have to check that they are stable under passage to the induced prederivators $\DD_I$, and also how they behave under duality (passage from $\DD$ to $\DD^\circ$). % [sic] qu'il sont -> qu'ils sont, agreement

Definition of the dual prederivator $\DD^\circ$
\[ \DD^\circ(I) \overset{\mathrm{def}}{=} \DD(I^\circ)^\circ \qquad\qquad \textit{etc.} \]
\DERmark{ENV0006}

\textbf{Der 1} (First localization axiom) : \textit{a}) If $I$, $J \in \sDia$, then
\[ \DD(I \amalg J) \longrightarrow \DD(I) \times \DD(J) \]
is an equivalence of categories.

\textit{b}) $\DD(\emptyset)$ is the final category.

Stable under passage to the induced prederivators, and passage to the dual.

\textbf{Corollary.} \textit{Let} $(I_\alpha)_{\alpha \in A}$ be \texttt{[a]} \textit{finite family of categories in} $\sDia$, $I = \coprod\limits_{\alpha \in A} I_\alpha$. \textit{Then}
\[ \DD(I) \longrightarrow \prod_{\alpha \in A} \DD(I_\alpha) \]
\textit{is an equivalence of categories.}


\endgroup
% END UNIT Der-ChI-234/Der-ChI-234_p021-025.tex
% BEGIN UNIT Der-ChI-234/Der-ChI-234_p026-030.tex
\begingroup\makeatletter%
\NoAutoSpacing%
\newcommand{\printedpage}[1]{\marginpar{\footnotesize p.~#1}}%
\newcommand{\msp}[1]{\marginpar{\footnotesize [#1]}}%
\newcommand{\D}{\mathbb{D}}%
\newcommand{\Dia}{\mathsf{Dia}}%
\newcommand{\HOT}{\mathsf{HOT}}%
\newcommand{\Ob}{\mathsf{Ob}}%
\makeatother\let\printedpage\derpp\let\msp\dermsp%
\let\mspin\DERmspin%


\printedpage{24}
% head: A. Grothendieck : Dérivateurs, Ch. I, Édition des Universités de Montpellier II et Paris VII

\noindent\textbf{Der 2} (Second localization axiom) : For every $I \in \Dia$, the canonical functor
\[ \D(I) \longrightarrow \underline{\mathsf{Hom}}(I^{\circ}, \mathcal{A}_{\D}) \quad , \]
$\mathcal{A}_{\D} = \D(e)$ fundamental category for the derivator $\D$, is conservative.\\
Stable under induction, and under duality.

\noindent\msp{page 44}\texttt{[page 44]}

\noindent A prederivator is called a \textit{left derivator} if it satisfies \textbf{Der 1} and \textbf{Der 2}, and in addition the following axioms (of which here is the first).

\noindent\textbf{Der 3} (g) (Existence of direct images) : For every $f : I \to J$ in $\Dia$, the functor $f^{*} : \D(J) \to \D(I)$ admits a \textit{right adjoint} $f_{*}$\,.\\
Stable under induction.

We call $f_{*}$ the direct image functor (cohomological) for $f$. The object $f_{*}(F) \in \D(J)$ (for $F \in \D(I)$) is also called the object of ``\textit{relative cohomology}'' of $I$ over $J$ (via $f$), with coefficients in $F$. We have canonical isomorphisms
\[ (gf)_{*} \simeq g_{*} f_{*} \quad . \]

\noindent\textbf{Corollary.} \textit{The categories $\D(I)$ have finite products, the functors $f^{*}$ and $f_{*}$ commute with them.}

Uses \textbf{Der 1} for $I = \emptyset$ and $I = e \amalg e$, and \textbf{Der 3} for $I \xrightarrow{f} e$.

We have a dual axiom :

\noindent\textbf{Der 3} (d) : existence of left adjoints $f_{!}$ to the $f^{*}$, the functor $f_{!}$ is called ``homological direct image'' via $f$, and we have a notion of \textit{relative homology} of $I$ over $J$ (for $f$), with coefficients in $F$.

In what follows, the indication (g) or (d) next to the designation \textbf{Der} $n$ of an axiom, % [sic] coté -> côté, missing circumflex
means that it applies to left derivators, respectively to right derivators. Likewise for statements, definitions \textit{etc.}

\noindent\msp{page 45}\texttt{[page 45]}

\noindent\textbf{Der 4}$_{0}$ (g)\footnote[13]{Unnecessary as an axiom, will become a theorem (follows from the other axioms).} : Let $U \overset{j}{\hookrightarrow} X \overset{i}{\hookleftarrow} Y$ be in $\Dia$, open and closed immersions complementary to one another. Then the adjunction morphism in $\D(Y)$
\[ i^{*} i_{*}(F_{Y}) \longrightarrow F_{Y} \]
is an isomorphism, hence
\[ i_{*} : \D(Y) \longrightarrow \D(X) \]
is fully faithful, and $i^{*} : \D(X) \to \D(Y)$ is a localization functor. The essential image of $i_{*}$ consists of the objects $F$ of $\D(X)$ such that $j^{*}(F)$ is the final object of $\D(U)$.

% foot: Version 8 juillet 2001

\printedpage{25}
% head: A. Grothendieck : Dérivateurs, Ch. I, Édition des Universités de Montpellier II et Paris VII

\noindent\textbf{Der 4}$_{0}$ (d) : With the preceding notation for $i$, $j$, the adjunction morphism in $\D(U)$
\[ F_{U} \longrightarrow j^{*} j_{!} F_{U} \]
is an isomorphism, - $j_{!}$ is fully faithful, $j^{*}$ a localization functor. The essential image of $j_{!}$ is the full subcategory of $\D(X)$ consisting of the $F$ such that $i^{*}(F)$ is an initial object of $\D(Y)$.

There are the two squares of page 8 and 9, but as for interpreting the cartesian or cocartesian condition it is premature here.

\noindent\msp{page 46}\texttt{[page 46]}

\noindent Base change axioms.

\noindent\textbf{Der 4} (g) (Fibers of a direct image) : Let $f : X \to Y$ be in $\Dia$, $y \in Y$, let us consider the known diagram
\[
\begin{tikzcd}[row sep=large, column sep=large]
X/y \arrow[r, "q"] \arrow[d, "f'"'] & X \arrow[d, "f"] \\
e \arrow[r, "p = \alpha_{y}"'] & Y
\end{tikzcd}
\]
and the corresponding base change morphism
\[ f_{*}(F)_{y} \overset{\mathrm{def}}{=} p^{*} f_{*}(F) \longrightarrow f'_{*}(q^{*}(F))\bigl(\overset{\mathrm{def}}{=} \Gamma_{\D}(X/y, q^{*}(F))\bigr) \quad . \]
\DERmark{ENV0007}\DERmark{ENV0008}


\noindent This is always an isomorphism.

\noindent\msp{page 47}\texttt{[page 47]}

\noindent\textbf{Der 5} (g) (Cofinality axiom for cohomology)\footnote[14]{It is a consequence of weaker statements.} : Let $\alpha : X' \to X$ in $\Dia$ be ``cohomologically cofinal'' \textit{i.e.} the $x\backslash X'$ are contractible ($x \in X$) \textit{i.e.} $\alpha^{\circ} : {X'}^{\circ} \to X^{\circ}$ is a hot-equivalence ``locally over $X^{\circ}$\,''. Then $\forall F \in \D(X)$, the homomorphism
\[ \Gamma_{\D}(X, F) \longrightarrow \Gamma_{\D}(X', \alpha^{*} F) \]
is an isomorphism (in $\mathcal{A} = \D(e)$).

\noindent N.B. If $f : X \to e$, we set $\Gamma_{\D}(X, F) = f_{*}(F)$ (``absolute'' cohomology of $X$, with coefficients in $F$).

\noindent Same thing\footnote[15]{False!} if $\alpha$ is ``homologically cofinal'', i.e. the $X'/x$ are contractible \texttt{[sentence crossed out]}.

% foot: Version 8 juillet 2001

\printedpage{26}
% head: A. Grothendieck : Dérivateurs, Ch. I, Édition des Universités de Montpellier II et Paris VII

\noindent\textbf{Remark.} We should be able to strengthen this axiom considerably, by requiring that
\[ \alpha_{!}(\varepsilon_{X'}) \xrightarrow{\sim} \varepsilon_{X} \qquad\qquad \text{in} \quad \HOT(X) \quad , \]
where $\varepsilon_{X'}$, $\varepsilon_{X}$ \texttt{[are]} the final objects of $\HOT(X')$, $\HOT(X)$, and where $\alpha_{!}(\varepsilon_{X'}) \to \varepsilon_{X}$ is the adjunction homomorphism, taking into account that $\varepsilon_{X'} \simeq \alpha^{*}(\varepsilon_{X})$. We say that $X'$ is \textit{homotopically trivial over $X$}. (And if the dual condition is satisfied in $\HOT(X^{\circ})$, we say that $X'$ is cohomotopically trivial over $X$.) This condition is probably \underline{not} self-dual. It is satisfied both if $\alpha$ is cohomologically cofinal, and if it is homologically cofinal (costrict or strict hot-equivalence).

\noindent N.B. For $f : X \to e$, saying that $f$ is cohomologically cofinal, homologically cofinal, or that $X$ is contractible, is equivalent.

\noindent\textbf{Corollary 1} (g) (Asphericity theorem for constant coefficients in cohomology). \textit{Let $f : X \to e$ be in} $\Dia$\textit{, $X$ contractible, $M \in \D(e)$. Then the adjunction homomorphism}
\[ M \xrightarrow{\sim} f_{*} f^{*}(M) \quad \bigl(\overset{\mathrm{def}}{=} \Gamma(X, f^{*}(M))\bigr) \]
\DERmark{ENV0009}

\textit{is an isomorphism.}

\noindent\textbf{Corollary 2} (g)\footnote[16]{\underline{False}, except for constant coefficients.} \textit{Let $f : X \to Y$ be} \texttt{[a]} \textit{strict hot-equivalence (i.e. $\forall y \in Y$, $X/y$ contractible). Then for $F \in \Ob\,\D(Y)$}
\[ F \overset{\text{adj}}{\xrightarrow{\sim}} f_{*} f^{*} F \]
\texttt{[is an]} \textit{isomorphism.} \texttt{[This corollary is crossed out, but not its proof.]}

\noindent\msp{page 48}\texttt{[page 48]}

\noindent We apply \textbf{Der 2}, \textbf{Der 4} (g), to reduce to looking, for $y_{0} \in Y$, at the homomorphism $\alpha$ of the following (commutative!) diagram :
\[
\begin{tikzcd}[row sep=large, column sep=huge]
\Gamma(\{y_{0}\}, F_{y_{0}}) = F_{y_{0}} \arrow[r, "\alpha"] & \Gamma_{\D}(X/y_{0}, f^{*}(F)) \\
\Gamma(Y/y_{0}, F) \arrow[u, "\varphi\,\wr"] \arrow[ur, "\alpha'"'] &
\end{tikzcd}
\quad ,
\]
where $\alpha'$ is deduced from the functor induced by $f$
\[ X/y_{0} \longrightarrow Y/y_{0} \]
and where $\varphi$ is deduced from the functor
\[ \{y_{0}\} \longrightarrow Y/y_{0} \qquad \text{(final object)}\,. \]

% foot: Version 8 juillet 2001

\printedpage{27}
% head: A. Grothendieck : Dérivateurs, Ch. I, Édition des Universités de Montpellier II et Paris VII

\noindent Thus $\varphi$ is an isomorphism by \textbf{Der 5} (g)\footnote[17]{\textit{cf.} corollary 3, page 50, which we should move up here.} (N.B. If $z \in Z$, whence $\{z\} \to Z$, this functor is cohomologically cofinal if and only if $z$ is a final object, homologically final \texttt{[rather cofinal]} if and only if $z$ is an initial object). Thus the conclusion of the corollary is equivalent to saying that $\alpha'$ is an isomorphism. But $X/y_{0} \to Y/y_{0}$ is also homologically cofinal, for if $u : y_{1} \to y_{0}$ is in $Y/y_{0}$, we have a canonical isomorphism
\[ (X/y_{0})/(y_{1}, u) \simeq X/y_{1} \]
which is contractible. So by the second part of \textbf{Der 5} (g),\footnote[18]{(which is \underline{false})} we win.

\noindent\msp{page 49}\texttt{[page 49]}

\noindent\textbf{Corollary} 4 (g) (base change criterion for cohomological direct images). \textit{Let}
\[
\begin{tikzcd}[row sep=huge, column sep=huge]
X' \arrow[r, "q"] \arrow[d, "f'"'] & X \arrow[d, "f"] \arrow[dl, dashed, "\gamma"', shorten <=4ex, shorten >=1ex] \\
Y' \arrow[r, "p"'] & Y
\end{tikzcd}
\]
\textit{be in} $\Dia$ (\textit{cf.} page 30), \textit{whence the base change homomorphism}
\[ p^{*} f_{*}(F) \longrightarrow f'_{*} q^{*}(F) \qquad (F \in \Ob\,\D(X)) \quad . \]
\textit{For this to be an isomorphism for all $F$, it is enough that for every $y' \in \Ob\,Y'$, the functor induced by} (p. 9)
\[ X'/y' \longrightarrow X/p(y') \]
\textit{be a strict hot-equivalence (or costrict);\footnote[19]{\underline{False} for ``strict'', works \underline{only} for ``\underline{\underline{co}}strict''.} this also means (in the strict case) that for}
\[ y'_{0} \in \Ob\,Y' \,,\ \ x_{0} \in \Ob\,X \,,\ \ u_{0} : f(x_{0}) \longrightarrow p(y'_{0}) \]
\textit{the category $K(y'_{0}, x_{0}, u_{0})$ of objects of the form}
\[ (x', v', u) \mid x' \in \Ob\,X',\ \ v' : f'(x') \to y'_{0},\ \ u : q(x') \overset{\dashleftarrow}{\longrightarrow} x_{0} \]
\textit{making the square in $Y$ commutative}
\DERmark{ENV0010}

\[
\begin{tikzcd}[row sep=huge, column sep=huge]
f(q(x')) \arrow[r, "\gamma(x')"] \arrow[d, shift right=0.5ex, "f(u)"'] & p(f'(x')) \arrow[d, "p(v')"] \\
f(x_{0}) \arrow[r, "u_{0}"'] \arrow[u, dashed, shift right=0.5ex, "\leftarrow\ \text{\scriptsize\shortstack[l]{case of\\costrict\\hot-equivalence}}"'] & p(y'_{0})
\end{tikzcd}
\tag*{$(*)$}
\]
\textit{is contractible.}

% foot: Version 8 juillet 2001

\printedpage{28}
% head: A. Grothendieck : Dérivateurs, Ch. I, Édition des Universités de Montpellier II et Paris VII

\noindent N.B. In the case of the \underline{co}strict hot-equivalence, it is the same description of
\[ K = (x_{0}, u_{0})\backslash(X'/y'_{0}) \quad , \]
except that the arrow $u$ is in the opposite direction $x_{0} \to q(x')$ and likewise in the commutative diagram $(*)$.\\
We apply \textbf{Der 2}, \textbf{Der 4} (g) and corollary 1 \texttt{[rather} \textbf{Der 5} (g)\texttt{]}.

\noindent\msp{page 50}\texttt{[page 50]}

\noindent\textbf{Corollary 5} (g). \textit{Let $f : X \to Y$ be fully faithful in} $\Dia$\textit{, and $F \in \D(X)$. Then the adjunction homomorphism}
\[ f^{*} f_{*}(F) \longrightarrow F \]
\textit{is an isomorphism.}

Indeed, it is enough to check it fiber by fiber (\textbf{Der 2}). Let $x \in X$, $y = f(x) \in Y$, % [sic] efffet -> effet, tripled f
then
\[ f^{*} f_{*}(F)_{x} \simeq f_{*}(F)_{y} \simeq \Gamma_{\D}(X/y, F) \]
now because of full faithfulness $X/y \simeq X/x$.

\noindent\textbf{Corollary 3} (move up to page 48).\footnote[20]{Will follow from the stronger \textbf{Der 6} (g) below.} \textit{Let $Z \in$} $\Dia$\textit{, suppose that $Z$ has a final object $e_{Z}$. Let $p : e \to Z$ be the corresponding functor. Then the canonical homomorphism (for $F$ in} $\D(Z)$\,)
\[
\begin{tikzcd}[row sep=huge]
Z \arrow[d, shift left=0.6ex, "f"] \\
e \arrow[u, shift left=0.6ex, "p"]
\end{tikzcd}
\qquad\qquad\qquad\qquad
f_{*}(F) \longrightarrow p^{*}(F)
\]
\textit{is an isomorphism.}\footnote[21]{Axiom of the index category with a final object (arbitrary coefficient).}

This homomorphism is \texttt{[a]} special case of this one, for
\[
\begin{tikzcd}[row sep=large, column sep=large]
Z' \arrow[rr, "p"] \arrow[dr, "f'"'] & & Z \arrow[dl, "f"] \\
& S &
\end{tikzcd}
\]
(taking $Z' = S = e$, $f' = \mathsf{id}_{e}$),
\[ f_{*}(F) \longrightarrow f'_{*}(p^{*}(F)) \quad , \qquad (\, F \in \Ob\,\D(Z)\,) \quad , \]
and it is a special case\footnote[22]{already considered in cor. 2, p. 47} of the morphisms of

% foot: Version 8 juillet 2001


\endgroup
% END UNIT Der-ChI-234/Der-ChI-234_p026-030.tex
% BEGIN UNIT Der-ChI-234/Der-ChI-234_p031-035.tex
\begingroup\makeatletter%
\newcommand{\printedpage}[1]{\marginpar{\footnotesize p.~#1}}%
\newcommand{\msp}[1]{\marginpar{\footnotesize [#1]}}%
\newcommand{\Ob}{\mathsf{Ob}}%
\newcommand{\id}{\mathsf{id}}%
\newcommand{\DD}{\mathbb{D}}%
\makeatother\let\printedpage\derpp\let\msp\dermsp%
\let\mspin\DERmspin%


\printedpage{29}
% head: A. Grothendieck : Dérivateurs, Ch. I, Édition des Universités de Montpellier II et Paris VII
% foot: Version 8 juillet 2001

\noindent\msp{page 51}
base change for
\[
\begin{tikzcd}
Z' \arrow[r, "p"] \arrow[d, "f'"'] & Z \arrow[d, "f"] \\
S \arrow[r, "\id"] & S
\end{tikzcd}
\qquad .
\]

\noindent\textsc{Remark} : We recover the fact that $i^*i_*(F_Y) \overset{\sim}{\to} F_Y$ is \texttt{[an]} isomorphism, for $i$ a closed immersion. Moreover, with the notation of \textbf{Der 4} (g), $j^*(i_*(F_Y))$ is computed by the base change theorem (corollary 2 \texttt{[rather 4]} above), which applies in every commutative situation
\[
\begin{tikzcd}[row sep=large, column sep=large]
X' \arrow[r, "q"] \arrow[d, "f'"'] \arrow[dr, phantom, "\text{\scriptsize cart}"] & X \arrow[d, "f"] \\
Y' \arrow[r, "p"] & Y
\end{tikzcd}
\]
when $p$ is an open immersion (and $\gamma = \id$) (then $X'/y' \to X/p(y')$ is even an isomorphism!), and we find, returning to the notation of \textbf{Der 4} (g),
\[
\begin{tikzcd}
\emptyset \arrow[r, "q"] \arrow[d, "f'"'] & Y \arrow[d, hook, "i"] \\
U \arrow[r, hook, "j"'] & X
\end{tikzcd}
\]
\textit{i.e.} $j^*(i_*(F_Y)) \simeq f'_*(\varphi)$, where $\varphi = q^*(F_Y)$ is the unique element of the category (equivalent to a one-point category, by \textbf{Der 1}) $\DD(\emptyset)$. On the other hand, $f'_*$ takes final object to final object, whence $j^*i_*(F_Y)$ is a final object.

Conversely, let $F$ be in $\DD(X)$, such that $j^*(F) = e_{\DD(U)}$. Consider
\[
F \xrightarrow{u} i_*i^*F \qquad\quad \text{(adjunction) ;}
\]
I say that this is an isomorphism. Indeed, by the axiom \msp{page 52} of localization \textbf{Der 2}, it suffices to prove that $i^*(u)$ and $j^*(u)$ are \texttt{[some]} isomorphisms. For $i^*(u)$ this follows from the fact that $i^*$ is a localization functor, and for $j^*$ from the fact that $j^*(F)$ and $j^*(i_*i^*(F))$ are final objects in $\DD(U)$.

\printedpage{30}
% head: A. Grothendieck : Dérivateurs, Ch. I, Édition des Universités de Montpellier II et Paris VII
% foot: Version 8 juillet 2001

\noindent\msp{page 53}

\subsection*{Homotopically proper, homotopically smooth morphisms.}

\noindent Let $f : X \to Y$ be \texttt{[a]} morphism of categories, consider $y_0 \in \Ob\, Y$, and the canonical functor
\[
X_{y_0} \xrightarrow{\varphi} X/y_0 \qquad\qquad\qquad \varphi(x) = (x, \id_{f(x)}) \tag*{(1)}
\]
($X_{y_0}$ fiber of $X$ at $y_0$). Let $(x_0, u_0) \in \Ob\, X/y_0$ hence $x_0 \in \Ob\, X$, $u_0 : f(x_0) \to y_0$. Consider the category coinduced by $\varphi$ over $(x_0, u_0)$
\[
(x_0, u_0)\backslash(X_{y_0}) \quad ,
\]
to test whether $\varphi$ is a costrict hot-equivalence \textit{i.e.} homologically \texttt{[rather co\-ho\-mo\-lo\-gi\-cal\-ly]} cofinal. We are going to interpret it via the functor induced by $f$
\[
x_0\backslash X \xrightarrow{\psi} f(x_0)\backslash Y
\]
going from the category of specializations of $x_0$ in $X$ to that of specializations of $f(x_0)$ in $Y$. We then find a canonical isomorphism between the category coinduced $(x_0, u_0)\backslash X_{y_0}$ by $X_{y_0}$ over the object $(x_0, u_0)$ of $X/y_0$ and the fiber category $(x_0\backslash X)_{u_0}$ of $x_0\backslash X$ over $f(x_0)\backslash Y$ at $u_0$. Thus

\noindent\textbf{Proposition.} \textit{For the functors $\varphi : X_{y_0} \to X/y_0$ ($y_0 \in \Ob\, Y$) to be costrict hot-equivalences, it is necessary and sufficient that the fibers of the functors $x_0\backslash X \to f(x_0)\backslash Y$ induced by $f$ be contractible.}

\noindent\msp{page 54}
We then say that $f$ is \textit{homotopically proper.} The dual notion is expressed, either by the condition that the functors induced by $f$
\[
X_{y_0} \longrightarrow y_0\backslash X
\]
\DERmark{ENV0011}

be strict hot-equivalences \textit{i.e.} that they be cohomologically \texttt{[rather ho\-mo\-lo\-gi\-cal\-ly]} cofinal, or that the fibers of the morphisms induced on the ``localized categories'' (categories of generizations)
\[
X/x_0 \longrightarrow Y/f(x_0)
\]
be contractible. (N.B. We use here that a category $Z$ is contractible if and only if $Z^\circ$ is.) We then say that $f$ is \textit{homotopically smooth.}

\noindent\textbf{Proposition.} \textit{The notion of proper (resp. smooth) functor is stable under base change, under composition. Isomorphisms, and more generally, closed immersions are proper. The image of a proper morphism is a closed subcategory.}\footnote[23]{Beware, an equivalence of categories is not proper! If $X' \hookrightarrow X$ \texttt{[is]} an inclusion of a subcategory of $X$, it is proper only if it is a closed immersion.}

\noindent N.B. The functor $X \to e$ is always proper and smooth, for the categories $x_0\backslash X$ and $X/x_0$ are contractible (having an initial object or a final object). Thus every projection functor $X \times Y \xrightarrow{pr_2} Y$ is both proper and smooth.

\printedpage{31}
% head: A. Grothendieck : Dérivateurs, Ch. I, Édition des Universités de Montpellier II et Paris VII
% foot: Version 8 juillet 2001

\noindent\msp{page 55}
\textsc{Proof}. Stability under base change :
\[
\begin{tikzcd}
X \arrow[d, "f"'] \arrow[dr, phantom, "\text{\scriptsize cart}"] & X' \arrow[l, "q"'] \arrow[d, "f'"] \\
Y & Y' \arrow[l, "p"]
\end{tikzcd}
\qquad ,
\]
whence for $x' \in \Ob\, X'$
\[
\begin{tikzcd}[row sep=large]
X/x \arrow[d] & X'/x' \arrow[l] \arrow[d] \\
Y/y & Y'/y' \arrow[l]
\end{tikzcd}
\qquad ,
\]
where $x$, $y'$, $y$ are the images of $x'$ in $X$, $Y'$, $Y$. We check that this diagram is \textit{cartesian}, hence the fibers of $X'/x' \longrightarrow Y'/y'$ are isomorphic to fibers of $X/x \longrightarrow Y/y$, hence contractible if these are.

In a case $X \xrightarrow{f} Y$ we check that for a bilocalized functor of $f$
\[
X/x \xrightarrow{\varphi=f/x} Y/y \quad ,
\]
every ``bilocalized''
\[
(X/x)/\xi \longrightarrow (Y/y)/\eta \qquad\qquad (\xi : x' \to x\,, \quad \eta : y' \to y)
\]
of $\varphi$ is again identified with a bilocalized of $f$
\[
X/x' \longrightarrow Y/y' \quad .
\]
Thus if $f$ is smooth, the bilocalized $f/x$ ($x \in \Ob\, X$) are so too. Let us note on the other hand.

\noindent\textbf{Proposition.}\footnote[24]{\textit{cf.} improvement in cor. 2, next page.} \textit{Every proper (or dually smooth) morphism with contractible fibers is a hot-equivalence.}

\noindent\msp{page 56}
This comes from the fact that for such a proper functor, if $F$ is a presheaf, or complex of abelian presheaves on $X$, the formation of the $Rf_*(F)$ commutes with passage to fibers, and likewise for $R^1f_*(F)$ if $F$ is a presheaf of groups (not even abelian).

\noindent\textbf{Corollary 1.} \textit{If $f : X \to Y$ is cohomologically proper with contractible fibers, then $Y$ is contractible if and only if $X$ is.}

Thanks to this, we find that the composite of two proper (smooth) functors is proper (smooth).

\printedpage{32}
% head: A. Grothendieck : Dérivateurs, Ch. I, Édition des Universités de Montpellier II et Paris VII
% foot: Version 8 juillet 2001

\noindent N.B. Contrary to what happens in topology for the notion of properness, if $gf$ is proper and $g$ proper, $f$ is not necessarily so, \textit{i.e.} an $S$-morphism of $S$-categories $X$, $Y$ proper over $S$ is not necessarily proper. Thus a section of $X/S$, for $X$ proper over $S$, is not necessarily proper, \textit{i.e.} is not necessarily a closed immersion. This is already false if $S = e$ (in which case $f : X \to e$ is always proper).

\noindent\textbf{Corollary 2.} \textit{A proper or smooth morphism, with contractible fibers is both a strict and a costrict hot-equivalence.}

Indeed, we apply corollary 1 to $X/y \to Y/y$ or $y\backslash X \to y\backslash Y$.

\noindent\textbf{Theorem} (g) \textit{Consider a base change situation in} \textsf{Dia}
\[
\begin{tikzcd}[row sep=large, column sep=huge]
X \arrow[d, "f"'] \arrow[dr, phantom, "\text{cart}"] & X' \arrow[l, "q"'] \arrow[d, "f'"] \\
Y & Y' \arrow[l, "p"]
\end{tikzcd}
\qquad , \qquad\qquad p^*f_*(F) \xrightarrow{c} f'_*q^*(F)
\]
\textit{($F \in \DD(X)$), $\DD$ a left exact derivator. If $f$ is (homotopically) proper or $p$ homotopically smooth, then $c$ is an isomorphism.}

\noindent\msp{page 57}
The question is whether for every $y' \in Y'$, putting $y = p(y')$ and considering the functor induced by $q$
\[
X/y \longleftarrow X'/y'
\]
the corresponding morphism
\[
\Gamma_{\DD}(X/y, F) \longrightarrow \Gamma_{\DD}(X'/y', q^*F)
\]
is \texttt{[an]} isomorphism. In the case $f$ proper, we consider
\[
\begin{tikzcd}[column sep=large]
X_y \arrow[d, "\sim\ \ \text{costrict}"'] & X'_{y'} \arrow[l, "\sim"'] \arrow[d, "\sim\ \ \text{costrict}"] \\
X/y & X'/y' \arrow[l]
\end{tikzcd}
\]
\DERmark{ENV0012}

whence the conclusion, since by \textbf{Der 5} (g), in
\[
\begin{tikzcd}[row sep=large, column sep=large]
\Gamma(X_y) \arrow[r] & \Gamma(X'_{y'}) \\
\Gamma(X/y) \arrow[u] \arrow[r] & \Gamma(X'/y') \arrow[u]
\end{tikzcd}
\]
\DERmark{ENV0013}

the vertical arrows and the upper horizontal one are \texttt{[some]} isomorphisms, hence also the lower horizontal one.

\printedpage{33}
% head: A. Grothendieck : Dérivateurs, Ch. I, Édition des Universités de Montpellier II et Paris VII
% foot: Version 8 juillet 2001

In the case $p$ smooth, let us consider on the $Y$, $Y'$ side the base change square % [sic] coté -> côté, missing circumflex (cf. "Du côté" below)
\[
\begin{tikzcd}[row sep=large, column sep=large]
Y & Y' \arrow[l] \\
Y/y \arrow[u] & Y'/y' \arrow[l] \arrow[u]
\end{tikzcd}
\]

\noindent\msp{page 58}
whence the adjoining diagram
\[
\begin{tikzcd}[column sep=small]
X \arrow[dd] & & X' \arrow[ll] \arrow[dd] & \\
& X/y \arrow[ul] & & X'/y' \arrow[ll, crossing over, "\kappa"'] \arrow[ul] \arrow[dd] \\
Y & & Y' \arrow[ll] & \\
& Y/y \arrow[ul] \arrow[from=uu, crossing over] & & Y'/y' \arrow[ll, "\pi"] \arrow[ul]
\end{tikzcd}
\qquad ,
\]
where the lateral (vertical) squares are cartesian. We note that $\pi$ is smooth and has contractible fibers (for $p$ is smooth) hence $\kappa$ also. Thus $\kappa$ is a costrict hot-equivalence, hence $\Gamma(X/y) \overset{\sim}{\longrightarrow} \Gamma(X'/y')$, we win again.

\noindent N.B. Unless I am mistaken, we have a converse of the theorem, \texttt{[to]} wit : If $f : X \to Y$ is such that the base change theorem for $f_*(F)$ holds for \textit{every} derivator (it probably suffices to take \textsf{HOTAB} and \textsf{HOT}) and for every $Y' \to Y$, then $f$ is homotopically proper. On the other hand, if $g : Y' \to Y$ is such that the base change theorem for $g$ holds for every $X \to Y$ and every $F \in \DD(X)$ ($\DD$ any left exact derivator), then $g$ is homotopically smooth. % [sic] pout tout -> pour tout, misprint
\DERmark{ENV0014}


\noindent On the homological side the smooth functors $f$ are exactly those for which $f_!$ commutes with every base change. The proper functors $g$ are exactly those for which base change by $g$ commutes with $f_!$, for every $f : X \to Y$ and every $\DD$ a right exact derivator.
\DERmark{ENV0015}


\noindent\msp{page 59}
Thus, the proper functors $f : X \to Y$ are characterized in \textit{four} remarkable ways
\begin{itemize}[label={}, leftmargin=0pt, itemindent=2\parindent, labelsep=0.3em]
\item[\textit{a})] $f_*$ commutes with every base change ($\forall$ left exact derivator $\DD$);
\item[\textit{b})] the base change $f^*$ commutes with every $g_!$ ($g : Y' \to Y$, $\DD$ a right exact derivator);
\item[\textit{c})] $\forall y \in \Ob\, Y$, $X_y \to X/y$ is a costrict hot-equivalence \textit{i.e.} is a homologically \texttt{[rather co\-ho\-mo\-lo\-gi\-cal\-ly]} cofinal functor;
\item[\textit{d})] $\forall x \in \Ob\, X$ the functor induced on the colocalized
\[
x\backslash X \longrightarrow f(x)\backslash Y
\]
\end{itemize}


\endgroup
% END UNIT Der-ChI-234/Der-ChI-234_p031-035.tex
% BEGIN UNIT Der-ChI-234/Der-ChI-234_p036-040.tex
\begingroup\makeatletter%
\derleqno%
\NoAutoSpacing%
\newcommand{\printedpage}[1]{\marginpar{\footnotesize p.~#1}}%
\newcommand{\msp}[1]{\marginpar{\footnotesize [#1]}}%
\newcommand{\DD}{\mathbb{D}}%
\newcommand{\cA}{\mathcal{A}}%
\newcommand{\Dia}{\mathsf{Dia}}%
\newcommand{\Ob}{\mathsf{Ob}}%
\newcommand{\Hom}{\mathsf{Hom}}%
\newcommand{\hot}{\mathsf{hot}}%
\newcommand{\Hot}{\mathsf{Hot}}%
\newcommand{\HOT}{\mathsf{HOT}}%
\newcommand{\isoto}{\overset{\sim}{\to}}%
\newcommand{\isolongto}{\overset{\sim}{\longrightarrow}}%
\newcommand{\sHom}{\mathsf{Hom}}%
\newcommand{\uHom}{\underline{\mathsf{Hom}}}%
\makeatother\let\printedpage\derpp\let\msp\dermsp%
\let\mspin\DERmspin%


\printedpage{34}
% head: A. Grothendieck : Dérivateurs, Ch. I, Édition des Universités de Montpellier II et Paris VII
% foot: Version 8 juillet 2001
\noindent has contractible fibers.

Another standard situation where base change works (but where the functors involved have no reason to be proper or smooth).

Let
\[
f : I \longrightarrow I' \quad , \qquad g : J \longrightarrow J'
\]
be in $\Dia$, whence the diagram (compare page 23)
\[
\begin{tikzcd}
X' = I \times J \arrow[r, "G'"] \arrow[d, "F'"'] & I \times J' = X \arrow[d, "F"] \\
Y' = I' \times J \arrow[r, "G"] & I' \times J' = Y
\end{tikzcd}
\quad .
\]
This diagram is such that for every left derivator

\noindent\msp{page 60}

\noindent we have base change
\[
G^*F_* \simeq F'_*{G'}^* \quad ,
\]
the reason here being that the induced functors
\[
X'/y' \longrightarrow X/y
\]
are \underline{\underline{co}}strict hot-equivalences, \textit{i.e.} that the categories $K$ (page 49) are contractible. Indeed, if $x_0 = (i_0, j'_0) \in \Ob X$, $y'_0 = (i'_0, j_0) \in \Ob Y'$, $u_0 = (u'_0, v'_0) : F(x_0) \to G(y'_0)$, where $u'_0 : f(i_0) \to i'_0$, $v'_0 : j'_0 \to g(j_0)$ are given, one checks that the category $K$ sought is simply isomorphic to $(i_0, u'_0)\backslash(I/i'_0) \times (j'_0\backslash J)/(j_0, v'_0)$ which is contractible, being a product of a category having an initial object and a category having a final object \texttt{[rather $(i_0, u'_0)\backslash(I/i'_0) \times (j'_0\backslash J)/(j_0, v'_0)$ which is contractible being the product of a category having an initial object and a category having a final object]}. Whew!
\DERmark{ENV0016}\DERmark{ENV0017}


I notice a very big blunder. I thought I remembered that for $f : X \to Y$ which is a strict hot-equivalence, one had for every $F \in \DD(Y)$, $F \isoto f_*f^*F$ (for example for $\DD = \mathsf{Der}(\mathit{Ab})$, or $\DD = \Hot$), \textit{i.e.} $f^*$ fully faithful, whence also that $\Gamma_{\DD}(Y, F) \to \Gamma_{\DD}(X, f^*F)$ is \texttt{[an]} isomorphism. But this is true only if $F$ is a \textit{constant} coefficient on $Y$, \textit{i.e.} of the form $q^*(M)$ ($q : Y \to e$, $M \in \DD(e) = \cA_{\DD}$). This is false, for general $F$, already in the case of the immersion $f : \{\varphi\} \to Y$, where $\varphi$ is an initial object. Then, $f$ is always a strict hot-equivalence, but the homomorphism $\Gamma(Y, F) \to \Gamma(X, f^*F)$ is in general not an isomorphism.

\par
\printedpage{35}
% head: A. Grothendieck : Dérivateurs, Ch. I, Édition des Universités de Montpellier II et Paris VII
% foot: Version 8 juillet 2001
\noindent\msp{page 61}

\subsection*{Notions of $\mathbb{D}$-equivalence and of $\mathbb{D}$-contractibility.}

So we are going to take up again the questions of acyclicity and base change.

\noindent\textbf{Definition 1.} Let $\DD$ be a prederivator, and let
\[
\begin{tikzcd}
X \arrow[rr, "f"] \arrow[dr, "p"'] & & Y \arrow[dl, "q"] \\
& e &
\end{tikzcd}
\]
be in $\Dia$. We say that $f$ is a $\DD$-\textit{equivalence} if the functor
\[
F \longmapsto f^*F \quad , \qquad \DD(Y) \longrightarrow \DD(X)
\]
induces a fully faithful functor on the subcategory of $\DD(Y)$ formed by the \textit{constant} $\DD$-coefficients, \textit{i.e.} if for all $M, N \in \Ob \cA = \DD(e)$ (``absolute coefficients''), we have
\[
\sHom(q^*M, q^*N) \isolongto \sHom(p^*M, p^*N) \quad . \tag*{(1)}
\]
Suppose that $p_*$ and $q_*$ exist. We shall put, for a sheaf $F$ on $X$
\[
H^*_{\DD}(X, F) \overset{\mathrm{def}}{=} p_*(F) \qquad \text{(and likewise } \ H^{\DD}_*(X, F) = p_!(F) \ \text{).}
\]
\DERmark{ENV0018}

N.B. And if $F = p^*M$, $M \in \cA$, we write simply $H^*_{\DD}(X, M)$, $H^{\DD}_*(X, M)$.

\noindent The homomorphism (1) is interpreted by adjunction as
\[
\begin{tikzcd}[row sep=small]
\sHom(M, q_*q^*N) \arrow[r] \arrow[d, equal] & \sHom(M, p_*p^*N) \arrow[d, equal] \\
\sHom(M, H^*_{\DD}(Y, q^*N)) \arrow[r] & \sHom(M, H^*_{\DD}(X, p^*N))
\end{tikzcd}
\]
deduced from the homomorphism
\[
q_*q^*N \longrightarrow p_*p^*N \quad , \qquad \textit{i.e.} \quad H^*_{\DD}(Y, N) \longrightarrow H^*_{\DD}(X, N) \quad . \tag*{(2)}
\]
Thus to say that (1) is an isomorphism for all $M, N$ means that (2) is an isomorphism for every $N$. Likewise, if $p_!$, $q_!$ exist this is equivalent to saying that
\[
p_!p^*N \longrightarrow q_!q^*N \quad , \qquad \textit{i.e.} \quad H^{\DD}_*(X, N) \longrightarrow H^{\DD}_*(Y, N) \tag*{$(2')$}
\]
is an isomorphism in $\cA$ for every $N \in \cA$.

\par
\printedpage{36}
% head: A. Grothendieck : Dérivateurs, Ch. I, Édition des Universités de Montpellier II et Paris VII
% foot: Version 8 juillet 2001
\noindent\msp{page 62}

\noindent In summary:

\noindent\textbf{Proposition 1.} \textit{If $p_*$ and $q_*$ exist, i.e. if $H^*_{\DD}(X, \bullet)$ and $H^*_{\DD}(Y, \bullet)$ exist (it even suffices that this exist for constant coefficients), then $f$ is a $\DD$-equivalence if and only if $f$ induces an isomorphism for the $\DD$-cohomologies with constant coefficients of $X$ and $Y$. If $p_!$ and $q_!$ exist, i.e. if $H^{\DD}_*(X, \bullet)$ and $H^{\DD}_*(Y, \bullet)$ exist, $f$ is a $\DD$-equivalence if and only if $f$ induces an isomorphism for the $\DD$-homologies with constant coefficients of $X$ and $Y$.}

\noindent\textbf{Remark 1.} For $f$ to be \texttt{[a]} $\DD$-equivalence, it \textit{suffices} that for every $F$ \textit{constant} on $Y$ the adjunction morphism $F \to f_*f^*F$ be \texttt{[an]} isomorphism.

\noindent\textbf{Remark 2.} If $f$, $f' : X \to Y$ are in the same connected component of $\uHom(X, Y)$, then $f$ is a $\DD$-equivalence $\Leftrightarrow$ $f'$ a $\DD$-equivalence.

\noindent\textbf{Proposition 2.} \textit{Isomorphisms are $\DD$-equivalences. If $X \overset{f}{\to} Y \overset{g}{\to} Z$, if two among $f$, $g$, $gf$ are $\DD$-equivalences, the third too.} % [sic] parmis -> parmi, spelling

\noindent\textbf{Definition 2.} We say that $X$ in $\Dia$ is $\DD$-\textit{aspheric} (or $\DD$-\textit{contractible}, or $\DD$-\textit{trivial}) if $X \overset{p}{\to} e$ is a $\DD$-equivalence, \textit{i.e.} if $p^*$ is fully faithful.

If $p_*$ exists, this therefore means that for $M \in \Ob(\cA)$
\[
M \overset{\mathrm{adj}}{\isolongto} p_*p^*M \quad , \qquad \textit{i.e.} \quad M \isolongto H^*_{\DD}(X, M) \quad ,
\]
and if $p_!$ exists, this means that for every $M$
\[
p_!p^*M \overset{\mathrm{adj}}{\isolongto} M \quad , \qquad \textit{i.e.} \quad H^{\DD}_*(X, M) \isolongto M \quad .
\]
\textbf{Corollary.} \textit{If $Y$} (resp. $X$) \textit{is contractible, then $f : X \to Y$ is a $\DD$-equivalence if and only if $X$} (resp. $Y$) \textit{is contractible.}

\noindent\textsc{Example}. Let us take $\DD = \HOT$. (We can take $\Dia$ = category of small categories). Then, for $M \in \cA = \Hot$ and $X$ a small category, we have
\[
H^*_{\HOT}(X, M) \simeq \sHom_{\cA}(\hot(X), M) \quad ,
\]
where
\[
\hot : \mathcal{C}\mathit{at} \longrightarrow \cA = \Hot
\]
is the canonical functor which associates to every $X$ its ``homotopy type''. We have
\[
H^{\HOT}_*(X, M) = \hot(X) \times M \quad .
\]

\noindent\msp{page 63}

\noindent We deduce from this that $f : X \to Y$ is a $\HOT$-equivalence (in the sense of definition 1) if and only if it is a $\hot$-equivalence (or ``weak equivalence'') \textit{i.e.} if and only if
\[
\hot(f) : \hot(X) \isolongto \hot(Y)
\]
\printedpage{37}
% head: A. Grothendieck : Dérivateurs, Ch. I, Édition des Universités de Montpellier II et Paris VII
% foot: Version 8 juillet 2001
is an isomorphism.

\noindent\textbf{Proposition 2.} \textit{Let $f : X \to Y$. Equivalent conditions (if $\DD$} \texttt{[is a]} left \textit{derivator}). % Resolved against official supplemental edition p37: a gauche). is upright.

a) \textit{$f$ is a $\DD$-equivalence, and remains so after every base change $Y' \to Y$ which is a local isomorphism (for example for $Y/F \to Y$, $F \in \Ob \widehat{Y}$).}

b) \textit{For every $y \in \Ob Y$, the induced functor $X/y \to Y/y$ is a $\DD$-equivalence.}

c) \textit{For every $y \in \Ob Y$, $X/y$ is $\DD$-contractible.}\footnote[25]{Other equivalent conditions:\par
\indent d) \textit{For every \underline{constant} sheaf $M_Y$ on $Y$, we have $M_Y \isoto f_*f^*M_Y$.}\par
\indent d$'$) \textit{As in} d) \textit{with \underline{locally} constant sheaf.}}

\noindent N.B. We say that $g : Z \to Y$ is a \textit{local isomorphism} if for every $z \in \Ob Z$, put $y = g(z)$, the induced morphism on the categories localized at $z$ and $y$
\[
Z/z \isolongto Y/y
\]
is an isomorphism. For example every open immersion, every induced functor $Y/y \to Y$ (localization functor) are local isomorphisms.

From this a) $\Rightarrow$ b) follows trivially.

b) $\Leftrightarrow$ c). Follows from the corollary on the preceding page, if one knows that $Y/y$ is $\DD$-contractible\footnote[26]{Cf. cor. 1 of \textbf{Der 5} (g)}. This will follow from \textbf{Der 5} below.

\noindent\msp{page 64}

c) $\Rightarrow$ a) ``Passage from local to global''. Since hypothesis c) is stable under base change by open immersion, it suffices to prove that c) implies that $f$ is a $\DD$-equivalence, \textit{i.e.} that for every $F = M_Y$ ($M \in \Ob \cA$), constant coefficient on $Y$, the adjunction homomorphism
\[
M_Y \longrightarrow f_*f^*M_Y = f_*M_X
\]
is \texttt{[an]} isomorphism. It suffices to prove that this is an isomorphism fiber by fiber (\textbf{Der 2}), hence that for every $y \in Y$,
\[
M_y \longrightarrow f_*(M_X)_y
\]
is an isomorphism. By \textbf{Der 4} (g), we have
\[
f_*(M_X)_y \isolongto \Gamma(X/y, M) \quad ,
\]
and the hypothesis that $X/y$ is $\DD$-contractible gives what we want.

\noindent\textbf{Proposition 2$'$.} \textit{Dually, if $\DD$ is a right derivator} \texttt{[rather right]}\textit{, equivalent conditions on $f : X \to Y$:}
\DERmark{ENV0019}


a) \textit{$f$ a $\DD$-equivalence,} \texttt{[and]} \textit{remains so under every base change $g : Y' \to Y$ which is a ``colocal isomorphism'' (i.e. for every $y' \in \Ob Y'$, the induced functor on the ``colocalized'' categories $y'\backslash Y' \isolongto y\backslash Y$ (where $y = g(y')$)} \texttt{[is an isomorphism]}\textit{).}

b) \textit{For every $y \in Y$, the induced functor $y\backslash X \to y\backslash Y$ is a $\DD$-equivalence.}

c) \textit{For every $y \in Y$, $y\backslash X$ is $\DD$-contractible.}
\DERmark{ENV0020}


\noindent To prove it, one must use that the $y\backslash Y$ are $\DD$-contractible (cf. cor. 1 of \textbf{Der 5} (d)).

\par
\printedpage{38}
% head: A. Grothendieck : Dérivateurs, Ch. I, Édition des Universités de Montpellier II et Paris VII
% foot: Version 8 juillet 2001
\noindent\msp{page 65}

\noindent\textbf{Definition 3.} Let $\DD$ be a prederivator (perhaps even a left or right derivator). We say that $f$ is a \textit{strict} $\DD$-equivalence if $f$ satisfies condition a) of \texttt{[the]} proposition 2. We say that it is a \textit{costrict} $\DD$-equivalence if $f$ satisfies that of \texttt{[the]} proposition $2'$; we also sometimes say that $f$ is $\DD$-cofinal. (The dual condition, \textit{i.e.} strict $\DD$-equivalence, would be called co-$\DD$-cofinal!)\footnote[27]{Better terminology: $f$ is $\DD$-\textit{aspheric} (strict $\DD$-equivalence), $\DD$-\textit{coaspheric} (co-strict $\DD$-equivalence).} % [sic] s'appelerait -> s'appellerait, spelling

\noindent\textbf{Proposition 3.} \textit{Let $X \overset{f}{\to} Y \overset{g}{\to} Z$. If $f, g$ are strict $\DD$-equivalences, then $gf$ too. If $gf$ and $f$ are strict $\DD$-equivalences, then $g$ too. If $\DD$} \texttt{[is a]} \textit{left or right derivator, every equivalence} \texttt{[of categories]} \textit{is a strict $\DD$-equivalence.} (For it has a right adjoint.)

We use the fact that the notion of local isomorphism is stable under base change.

\noindent\msp{page 66}

\noindent\textbf{Proposition 4.} \textit{Let $\DD$ be a right derivator, $f : X \to Y$ a functor in} $\Dia$\textit{, such that for $y \in Y$, $y\backslash X$ is $\DD$-contractible (hence $f$ a costrict $\DD$-equivalence, subject to $\DD$ being a right derivator \ldots\ ). Then for every $F$ in $\DD(Y)$ (not necessarily with constant coefficients) we have} % [sic] constant -> constants, agreement
\[
H^*_{\DD}(Y, F) \isolongto H^*_{\DD}(X, f^*(F))
\]
\textit{is an isomorphism. Conversely, if the condition is true for every $F$, then the $y\backslash X$ are $\DD$-contractible.}

\noindent\textsc{Proof}. It amounts to the same thing to say that for every $M \in \Ob \cA = \DD(e)$, we have
\[
\sHom(M, H^*_{\DD}(Y, F)) \overset{\mathrm{iso}}{\longrightarrow} \sHom(M, H^*_{\DD}(X, f^*(F))) \ .
\]
But by adjunction, this arrow is interpreted as
\[
\sHom(M_Y, F) \longrightarrow \sHom(M_X, f^*F) \tag*{$(*)$}
\]
On the other hand, if $f_!$ exists, we have
\[
\sHom(M_X, f^*F) = \sHom(f_!(M_X), F) \ ,
\]
and the homomorphism $(*)$ is deduced from the adjunction homomorphism
\[
f_!(M_X) = f_!(f^*(M_Y)) \longrightarrow M_Y \ , \tag*{$(**)$}
\]
by taking the $\sHom$ into $F$. Thus the condition is equivalent to saying that this homomorphism is \texttt{[an]} isomorphism, whatever the constant coefficient $M_Y$ on $Y$. To see whether this is an isomorphism, we look fiber by fiber (\textbf{Der 2}), now the fibers are interpreted as $H^{\DD}_*(y\backslash X, M_X)$ (by \textbf{Der 4} (d)) and as $M$.


\endgroup
% END UNIT Der-ChI-234/Der-ChI-234_p036-040.tex
% BEGIN UNIT Der-ChI-234/Der-ChI-234_p041-045.tex
\begingroup\makeatletter%
\newcommand{\printedpage}[1]{\marginpar{\footnotesize p.~#1}}%
\newcommand{\msp}[1]{\marginpar{\footnotesize [#1]}}%
\newcommand{\D}{\mathbb{D}}%
\newcommand{\Dia}{\mathsf{Dia}}%
\newcommand{\Ob}{\mathsf{Ob}}%
\newcommand{\Fl}{\mathsf{Fl}}%
\newcommand{\Hom}{\underline{\mathsf{Hom}}}%
\newcommand{\hot}{\mathsf{hot}}%
\newcommand{\Cat}{\mathcal{C}at}%
\newcommand{\isoto}{\overset{\sim}{\longrightarrow}}%
\makeatother\let\printedpage\derpp\let\msp\dermsp%
\let\mspin\DERmspin%


\printedpage{39}
% head: A. Grothendieck : Dérivateurs, Ch. I, Édition des Universités de Montpellier II et Paris VII
% foot: Version 8 juillet 2001
\msp{page 67}

\noindent Thus the arrow under consideration $(**)$ is interpreted as
\[
H_*^{\D}(y\backslash X, M_X) \longrightarrow M
\]
and this is the adjunction homomorphism $p_!(p^*(M)) \to M$.
\[
\begin{tikzcd}
X \arrow[rr, "f"] \arrow[dr, "p"{swap}] & & Y \arrow[dl, "q"] \\
& e &
\end{tikzcd}
\]

\noindent Thus it is an isomorphism for every $M$ if and only if $y\backslash X$ is $\D$-contractible.

\noindent We are going to state what we have proved.

\noindent\textbf{Theorem 1.} \textit{Let}
\[
\begin{tikzcd}
X \arrow[rr, "f"] \arrow[dr, "p"{swap}] & & Y \arrow[dl, "q"] \\
& e &
\end{tikzcd}
\]
\textit{be in} $\Dia$. \textit{We assume that} $\D$ \textit{is a prederivator satisfying} \textbf{Der 2}, \textit{that} $p_*, q_*$ \textit{and} $f_!$ \textit{exist, that} $f_!$ \textit{satisfies} \textbf{Der 4} (d) \textit{for the calculation of fibers. Equivalent conditions:}

a) $f$ \textit{is a costrict} $\D$\textit{-equivalence, i.e.} (prop. 2$'$) \textit{for every} $y \in \Ob\, Y$, $y\backslash X$ \textit{is} $\D$\textit{-contractible;}
\DERmark{ENV0021}


b) \textit{for every} $F \in \Ob\, \D(Y)$, \textit{the following homomorphism is} \texttt{[an]} \textit{isomorphism}
\[
H^*_{\D}(Y, F) \isoto H^*_{\D}(X, f^*(F)) \ ;
\]

c) \textit{for every} $M \in \Ob(\mathcal{A} = \D(e))$, \textit{the adjunction homomorphism is} \texttt{[an]} \textit{isomorphism}
\[
f_!(M_X) = f_!(f^*(M_Y)) \isoto M_Y \ .
\]

I am going to renumber the axioms on the prederivators. % [sic] le prédérivateurs -> les prédérivateurs, article singulier devant un pluriel

\noindent\textbf{Der 1}, \textbf{Der 2} as before (p. 45).

\noindent\textbf{Der 3}\footnote[28]{Axiom of aspheric indexing categories, \textit{cf.} variants \textbf{Der 3} a) b) c) page 97, and especially pages 111-112.} : Let $X$ be in $\Dia$, $p : X \to e$. Suppose $X$ contractible.\footnote[29]{N.B. It suffices to assume that $X$ has a final object. But this is not a self-dual axiom.} Then
\[
p^* : M \longmapsto M_X \quad , \qquad p^* : \mathcal{A} = \D(e) \longrightarrow \D(X)
\]
is fully faithful.

\printedpage{40}
% head: A. Grothendieck : Dérivateurs, Ch. I, Édition des Universités de Montpellier II et Paris VII
% foot: Version 8 juillet 2001
\msp{page 68}

\noindent\textbf{Der 4} (g) : existence of the $f_*$.

\noindent\textbf{Der 4} (d) : existence of the $f_!$.

\noindent\textbf{Der 5} (g) Calculation of the fibers of $f_*(F)$.

\noindent\textbf{Der 5} (d) Calculation of the fibers of $f_!(F)$.

\noindent The old axiom \textbf{Der 5} (g) follows from \textbf{Der 2}, \textbf{Der 3}, \textbf{Der 4} (g,d), \textbf{Der 5} (d), thanks to the preceding theorem.

An interesting aspect of the preceding theorem, if we view it as stating the equivalence of a) and b), is that this statement does not appeal to homological notions, but that its proof does appeal to them, in a relatively strong form (calculation of the fibers of $f_!(F)$). For the base change theorems for relative $\D$-cohomology, \textit{i.e.} for the $f_*^{\D}$, we also constantly need the calculation of the fibers of $f_*(F)$. Thus it seems that we cannot dispense with any of the preceding axioms, including \textbf{Der 4} (d) and \textbf{Der 5} (d), even if we only want to concern ourselves with the formalism of the $f^*, f_*$ (cohomological aspect). In what follows, we shall therefore assume these axioms satisfied.

\noindent\textbf{Theorem 2.} \textit{Let} $f : X \to Y$ \textit{be in} $\Dia$. \textit{Equivalent conditions:}

a) \textit{For every} $y_0 \in \Ob\, Y$, \textit{the canonical functor}
\[
X_{y_0} \overset{\varphi}{\longrightarrow} X/y_0 \quad (\simeq X \times_Y Y/y_0) \ ,
\]
\textit{where} $X_{y_0}$ \textit{denotes the fiber at} $y_0$, \textit{is a co-strict} $\D$\textit{-equivalence.}

b) \textit{For every} $x_0 \in \Ob\, X$, \textit{the induced functor on the categories colocalized at} $x_0$ \textit{and at} $y_0 = f(x_0)$
\[
x_0\backslash X \longrightarrow y_0\backslash Y
\]
\textit{has} $\D$\textit{-contractible fibers.}

\msp{page 69}

c) \textit{The formation of the} $f_*F$ \textit{(}$F$ \textit{in} $\D(X)$\textit{) commutes with every base change} $Y' \to Y$, \textit{and this remains true after every base change -- i.e. if we have} $Y'' \to Y' \to Y$, \textit{whence}
\[
\begin{tikzcd}
X'' \arrow[r] \arrow[d, "f''"{swap}] & X' \arrow[r] \arrow[d, "f'"] & X \arrow[d, "f"] \\
Y'' \arrow[r, "p'"{swap}] & Y' \arrow[r, "p"{swap}] & Y
\end{tikzcd}
\]
\textit{the ``base change theorem'' is true for} $f'_*$ \textit{and} $p'^{*}$.

c$'$) \textit{The formation of the} $f_*(F)$ \textit{(}$F \in \Ob\, \D(X)$\textit{) commutes with passage to the fibers and this remains true after every base change} $p : Y' \to Y$.

c$''$) \textit{As in} c$'$) \textit{with} $Y' \to Y$ \textit{a localization morphism} $Y/y \to Y$.

d) \textit{If} $Y'' \overset{p}{\to} Y' \to Y$, \textit{with} $p$ \textit{coaspheric, then} $q : X'' = X \times_Y Y'' \to X' = X \times_Y Y'$ \textit{is coaspheric.}

\printedpage{41}
% head: A. Grothendieck : Dérivateurs, Ch. I, Édition des Universités de Montpellier II et Paris VII
% foot: Version 8 juillet 2001

d$'$) \textit{As in} d) \textit{with} $Y'' = e$ (where $e\longrightarrow Y'$ selects a final object of $Y'$).
\DERmark{ENV0022}


d$''$) \textit{As in} d) \textit{with} $Y'' = \{1\}$, $Y' = \boldsymbol{\Delta}_1$.

.

\noindent\textsc{Proof}. a) $\Leftrightarrow$ b) cf. proposition page 53. Let us show that these conditions imply c). Since they are stable under base change (\textit{cf.} immediate proof page 55) therefore by \textbf{Der 2} we are reduced to proving that $H^*_{\D}(X/y, F) \overset{\sim}{\to} H^*_{\D}(X_y, F)$ $(F \in \Ob\, \D(X))$, which follows from theorem 1.\footnote[30]{Thus we use \textbf{Der 5} (d).} (Cf. proof of the theorem page 56, which works with more general hypotheses.)
\DERmark{ENV0023}\DERmark{ENV0024}


We have c) $\Leftrightarrow$ c$'$) $\Rightarrow$ c$''$) (the equivalence c) $\Leftrightarrow$ c$'$) by \textbf{Der 2}).

Let us prove c$''$) $\Rightarrow$ a).

\msp{page 70}

\noindent Consider the diagram with cartesian squares
\[
\begin{tikzcd}
X_{y_0} \arrow[r, "q"] \arrow[d, "p"{swap}] & X/y_0 \arrow[r] \arrow[d, "f_0"] & X \arrow[d, "f"] \\
\{y_0\} \arrow[r] & Y/y_0 \arrow[r] & Y
\end{tikzcd}
\qquad\qquad .
\]

\noindent Hypothesis c$''$) proves that for every $F$ in $\D(X/y_0)$, we have $H^*_{\D}(X/y_0, F) \overset{\sim}{\to} H^*_{\D}(X_{y_0}, F)$. By theorem 1, this implies that $q$ is a co-strict $\D$-equivalence. % [sic] un D-équivalence -> une D-équivalence, accord du genre
\DERmark{ENV0025}


\noindent\textbf{Corollary 1.} Dual statement of the theorem, for the functors $X_{y_0} \to y_0\backslash X$, $X/x_0 \to Y/y_0$, and base change for the $f_!(F)$ \ldots
\DERmark{ENV0026}


\noindent\textbf{Definition 4.} Under the equivalent conditions of theorem 2, we say that $f$ is $\D$-\textit{proper.} The dual condition is called $f$ $\D$-\textit{smooth.}
\DERmark{ENV0027}


\noindent\textbf{Corollary 2.} \textit{Let} $f : X \to Y$ \textit{be in} $\Dia$. \textit{If} $f$ \textit{is} $\D$\textit{-proper or} $\D$\textit{-smooth, and has} $\D$\textit{-contractible fibers, then} $f$ \textit{is a strict and costrict} $\D$\textit{-equivalence. Better,} $f$ \textit{is a} universal $\D$\textit{-equivalence} \textit{i.e. is a} $\D$\textit{-equivalence and remains so after every base change.}\footnote[31]{Generalizes: \textbf{Corollary 2$'$.} \textit{Let}
\[
\begin{tikzcd}[ampersand replacement=\&]
X \arrow[rr, "f"] \arrow[dr] \& \& Y \arrow[dl] \\
\& S \&
\end{tikzcd}
\]
\textit{with} $X, Y$, $\D$\textit{-proper over} $S$. \textit{If} $f$ \textit{induces a} $\D$\textit{-equivalence on the fibers it is a} $\D$\textit{-equivalence. Likewise in the smooth case.}\par
This is proposition 6 below.}

\printedpage{42}
% head: A. Grothendieck : Dérivateurs, Ch. I, Édition des Universités de Montpellier II et Paris VII
% foot: Version 8 juillet 2001
\msp{page 71}

\noindent It suffices to prove that $f$ is a $\D$-equivalence. Suppose for example $f$ $\D$-proper. It suffices to prove that if $Y$ is $\D$-contractible, then $X$ is so (if need be, applying it to the $Y/y_0$ after base change, to show that the $X/y_0$ are contractible). For this, it suffices to prove that $f$ is a $\D$-equivalence, and we know that this is the case if for every $F$ \textit{constant} on $Y$, the homomorphism $F \overset{\text{adj}}{\longrightarrow} f_*f^*F$ is \texttt{[an]} isomorphism. In fact, this is true for \textit{every} $F$.

\noindent\textbf{Corollary 3.} \textit{If} $f : X \to Y$ \textit{is} $\D$\textit{-proper or} $\D$\textit{-smooth, with} $\D$\textit{-contractible fibers, then} $f^* : \D(Y) \to \D(X)$ \textit{is fully faithful, i.e. the adjunction homomorphisms}
\[
F \longrightarrow f_*f^*F \quad , \qquad f_!f^*F \longrightarrow F
\]
\textit{in} $\D(Y)$ \textit{are isomorphisms.}

\noindent It suffices to see this for $f$ proper, and this follows from the calculation of the fibers of $f_*f^*F$ \ldots

\noindent\textbf{Proposition 5.} \textit{The notion of a} $\D$\textit{-proper functor} (resp. $\D$\textit{-smooth}) \textit{is stable under base change,} (trivial), \textit{and under composition} (use \texttt{[the]} corollary 2). \textit{Every closed immersion is} $\D$\textit{-proper, every open immersion} $\D$\textit{-smooth. More generally, every local isomorphism is} $\D$\textit{-smooth, every local co-isomorphism is} $\D$\textit{-proper.}

\noindent\textbf{Proposition 6.} \textit{Let} $X, Y$ \textit{be over} $S$, $f : X \to Y$ \textit{an} $S$\textit{-morphism. Suppose} $X, Y$ \textit{proper over} $S$, \textit{or} $X$ \textit{and} $Y$ \textit{smooth over} $S$, \textit{and that for every} $s \in \Ob\, S$, $f_s : X_s \to Y_s$ \textit{is a} $\D$\textit{-equivalence, then} $f$ \textit{also.}\footnote[32]{\textsc{Proof}. We shall even prove that on the subcategory of coefficients on $Y$ which come from $\D(S)$ \texttt{[rather from }$\D(S)$\texttt{]}, $f^* : \D(Y) \to \D(X)$ induces a fully faithful functor. If $p$ and $q$ are \texttt{[is]} proper, I interpret it as $q_*q^*F \overset{\sim}{\to} p_*p^*F$ for every $F$, and this follows from the calculation of $p_*$, $q_*$ by passage to the fibers (proper), and from the hypothesis on the $f_s$. OK.}
\DERmark{ENV0028}\DERmark{ENV0029}


\msp{page 72}

\noindent\textbf{Theorem 3.} \textit{Let}
\[
\begin{tikzcd}
X \arrow[d, "f"{swap}] & X' \arrow[l, "q"{swap}] \arrow[d, "f'"] \\
Y & Y' \arrow[l, "p"{swap}]
\end{tikzcd}
\]
\textit{be cartesian in} $\Dia$. \textit{Then the base change morphism}
\[
p^*f_* \longrightarrow f'_*q^*
\]
\textit{is an isomorphism if} $f$ \textit{is} $\D$\textit{-proper, or} $p$ $\D$\textit{-smooth.}

\noindent The $\D$-proper case already proved. The $\D$-smooth case: the proof p. 56 is OK.

\printedpage{43}
% head: A. Grothendieck : Dérivateurs, Ch. I, Édition des Universités de Montpellier II et Paris VII
% foot: Version 8 juillet 2001

\subsection*{Comparison of $X$ and $X^{\circ}$.}

Let $X$ be \texttt{[a]} category, we define a category $\Phi(X)$ and functors
\[
\begin{tikzcd}[column sep=small]
& & \Phi(X) \arrow[ddl, "p_X"{swap, name=P}] \arrow[ddr, "q_X"{name=Q}] & & \\
|[alias=L]| \text{\footnotesize\begin{tabular}{c}$\Cat$-cofibrant\\(hence proper)\end{tabular}} & & & & |[alias=R]| \text{\footnotesize\begin{tabular}{c}$\Cat$-cofibrant\\(hence proper)\\ \texttt{[rather }$\Cat$\texttt{-cofibrant}\\ \texttt{(hence proper)]}\end{tabular}} \\
& X & & X^{\circ} &
\arrow[from=P, to=L]
\arrow[from=Q, to=R]
\end{tikzcd}
\]
\DERmark{ENV0030}

which will be $\hot$-equivalences, and $\D$-equivalences if $\D$ is a derivator whose domain contains $X$ and $\Phi(X)$. We set
\[
\Ob\, \Phi(X) = \Fl(X) \quad ,
\]
so the same objects as $\mathcal{F}\ell(X) = \Hom(\boldsymbol{\Delta}_1, X)$, but the morphisms are different:

\msp{page 73}
\[
\begin{tikzcd}
Y_1 \arrow[r, "u"] & Y_0 \arrow[d, "v_0"] \\
Y'_1 \arrow[u, "v_1"] \arrow[r, "u'"{swap}] & Y'_0
\end{tikzcd}
\]
a morphism from $u$ to $u'$ is a pair $v_0 : Y_0 \to Y'_0$, $v_1 : Y'_1 \to Y_1$ such that $v_0uv_1 = u'$. Thus source and target $s(u)$ and $b(u)$ are the one contravariant, the other covariant in $u$. We take for $q$ the source functor, for $p$ the target functor. % [sic] l'autre covariant -> l'autre covariante, accord avec « l'une contravariante »

If we apply this to $X^{\circ}$, we find
\[
\begin{tikzcd}
& \Phi(X^{\circ}) \arrow[dl, "p_{X^{\circ}}"{swap}] \arrow[dr, "q_{X^{\circ}}"] & \\
X^{\circ} & & (X^{\circ})^{\circ} = X
\end{tikzcd}
\]
We have $\Ob\, \Phi(X^{\circ}) = \Fl(X^{\circ}) \simeq \Fl(X)$
\[
\begin{tikzcd}
Y_1 \arrow[r, "u"] & Y_0 \arrow[d, "v_1"] & \text{in}\ \ X \\
Y'_1 \arrow[u, "v_0"] \arrow[r, "u'"{swap}] & Y'_0 & \text{in}\ \ X
\end{tikzcd}
\]
and the same homomorphisms in the sense of $\Phi(X)$ and in the sense of $\Phi(X^{\circ})$. Thus
\[
\left\{
\begin{aligned}
&\Phi(X^{\circ}) = \Phi(X) \\
&p_{X^{\circ}} = q_X \\
&q_{X^{\circ}} = p_X \quad ,
\end{aligned}
\right.
\]


\endgroup
% END UNIT Der-ChI-234/Der-ChI-234_p041-045.tex
% BEGIN UNIT Der-ChI-234/Der-ChI-234_p046-050.tex
\begingroup\makeatletter%
\NoAutoSpacing%
\newcommand{\printedpage}[1]{\marginpar{\footnotesize p.~#1}}%
\newcommand{\msp}[1]{\marginpar{\footnotesize [#1]}}%
\newcommand{\D}{\mathbb{D}}%
\newcommand{\Ob}{\mathop{\mathsf{Ob}}}%
\newcommand{\Dia}{\textsf{Dia}}%
\newcommand{\bs}{\backslash}%
\makeatother\let\printedpage\derpp\let\msp\dermsp%
\let\mspin\DERmspin%


\printedpage{44}
% head: A. Grothendieck : Dérivateurs, Ch. I, Édition des Universités de Montpellier II et Paris VII
the passage from $X$ to $X^\circ$ amounts to interchanging $p_X$ and $q_X$. The fiber of $\Phi(X)$ at $x \in \Ob X$ is $\simeq (X/x)^\circ$, that at $y^\circ \in \Ob X^\circ$ is $\simeq y\bs X$
\[
\left\{
\begin{aligned}
&p_X^{-1}(\{x\}) \simeq (X/x)^\circ\\
&q_X^{-1}(\{y^\circ\}) \simeq y\bs X
\end{aligned}
\right.
\]

N.B. : $p_X$ is $W$-proper \texttt{[rather $W$-proper]} with $W$-aspheric fibers, hence in $W^u$ \texttt{[$W^u =$ set of morphisms which are in $W$ and remain so after any base change]}, $q_X$ is $W$-proper with $W$-aspheric fibers, hence in $W^u$.
\DERmark{ENV0031}


\msp{page 74}

One checks
\[
\left\{
\begin{aligned}
&\Phi(X)/x \simeq \Phi(X/x) &\qquad& x \in \Ob X\\
&\Phi(X)/x^\circ \simeq \Phi(x\bs X) &\qquad& x^\circ \in \Ob X^\circ
\end{aligned}
\right.
\]

I say that $p_X$, $q_X$ are $\D$-\textit{strict equivalences} and \textit{costrict.} By the preceding formulas, this means that the $\Phi(X/x)$, $\Phi(x\bs X)$ are $\D$-contractible. Thus that if $Z$ has a final object, $\Phi(Z)$ is $\D$-contractible. But let us consider
\[
\begin{tikzcd}
\Phi(Z) \arrow[dr, "q_Z"] & \\
& Z^\circ
\end{tikzcd}
\quad ;
\]
it suffices to prove that this is a strict $\D$-equivalence (since $Z^\circ$, having an initial object, is $\D$-contractible), hence that the $\Phi(Z)/z_0^{\mathrm{op}}$ are $\D$-contractible. We have by the preceding formulas % [sic] contractil -> contractile, adjectif ailleurs écrit « contractile »
\DERmark{ENV0032}

\[
\Phi(Z)/z_0^{\mathrm{op}} \simeq \Phi(z_0\bs Z) = \Phi(T) \quad ,
\]
\DERmark{ENV0033}

where this time $T$ has an initial object $\emptyset$ and a final object $e$. Let us then consider the object
\[
\emptyset \longrightarrow e
\]
of $\Phi(T)$, it is a final object, since for every $x \to y$ in $\Phi(T)$, there is a unique arrow
\[
\begin{tikzcd}
x \arrow[r] & y \arrow[d] \\
\emptyset \arrow[u] \arrow[r] & e
\end{tikzcd}
\]
% foot: Version 8 juillet 2001

\printedpage{45}
% head: A. Grothendieck : Dérivateurs, Ch. I, Édition des Universités de Montpellier II et Paris VII
\msp{page 75}

To summarize

\textbf{Proposition 1 :} \textit{$p_X$ and $q_X$ are strict} \textsf{hot}\textit{-equivalences,} i.e. \textsf{Hot}\textit{-aspheric (hence also strict $\D$-equivalences if $\Phi(X)$ is in the domain of $\D$) with fibers $(X/x)^\circ$ and $y\bs X$ respectively (hence with contractible fibers).}

\textbf{Corollary 1.} \textit{For $X$ to be $\D$-contractible, it is necessary and sufficient that $X^\circ$ be so.}

\hspace*{\parindent}More generally :

\textbf{Corollary 2.} \textit{Let $f : X \to Y$. For $f$ to be a $\D$-equivalence, it is necessary and sufficient that $f^\circ$ be so.}

\hspace*{\parindent}We shall use that the diagram
\[
\begin{tikzcd}
& \Phi(X) \arrow[dl, "p_X"'] \arrow[dr, "q_X"] & \\
X & & X^\circ
\end{tikzcd}
\]
is functorial in $X$. We therefore have
\[
\begin{tikzcd}
& \Phi(X) \arrow[rr, "\Phi(f)"] \arrow[dl, "p_X"'] & & \Phi(Y) \arrow[dl, "p_Y"] \arrow[dd, "q_Y"] \\
X \arrow[rr, "f" near end] \arrow[from=1-2, to=3-2, "q_X"' near end, crossing over] & & Y & \\
& X^\circ \arrow[rr, "f^\circ"'] & & Y^\circ
\end{tikzcd}
\]
the two squares commutative. Since $p_X$ and $p_Y$ are $\D$-equivalences, $f$ is a $\D$-equivalence if and only if $\Phi(f)$ is, and likewise for $f^\circ$. These arguments assume that if $X$ is in the domain of $\D$, $\Phi(X)$ is too. This is the case if \Dia{} consists of finite ordered categories.
% foot: Version 8 juillet 2001

\printedpage{46}
% head: A. Grothendieck : Dérivateurs, Ch. I, Édition des Universités de Montpellier II et Paris VII
\msp{page 76}

\subsection*{$\D$-fibrations}

\textbf{Proposition 1.} \textit{Let $f : X \to Y$ be in} \Dia\textit{, equivalent conditions :}

\hspace*{\parindent}a) \textit{For every $y \in Y$, the inclusion $X_y \to X/y$ is a $\D$-equivalence, and for every arrow $y_1 \to y_0$ of $Y$, $X/y_1 \to X/y_0$ is a $\D$-equivalence.}

\hspace*{\parindent}b) \textit{For every} constant \textit{coefficient $F$ on $Y$, $f_*f^*F$ is locally constant,}\footnote[33]{By this we mean that for every $u : y \to y'$ in $Y$, $G_{y'} \to G_y$ is \texttt{[an]} isomorphism. Without additional axioms on $\D$, it is not clear that this implies that the ``restrictions'' of $G$ to the $Y/y$ or $y\bs Y$, are isomorphic to constant coefficients. See however axiom \textbf{Der 3} a$'$), page 97, which ensures that it comes to the same thing.} \textit{and ``is computed fiber by fiber''} i.e. \textit{the canonical homomorphism}
\[
f_*f^*(F)_y \overset{\sim}{\longrightarrow} H^*_{\D}(X_y, F|X_y)
\]
\texttt{[is an isomorphism]}.

These conditions are equivalent, and imply the conditions of (b) for every $F$ \textit{locally} constant on $Y$.

The first condition in (a) means that for every $y \in Y$, and every $M \in \Ob \mathcal{A}$,
\[
H^*_{\D}(X/y, M) \longrightarrow H^*_{\D}(X_y, M)
\]
is an isomorphism. Now this arrow is interpreted as
\[
f_*(M_X)_y \longrightarrow H^*_{\D}(X_y, M)
\]
is an isomorphism, \textit{i.e.} that $f_*(M_X)$ is computed fiber by fiber. Since $M_X \simeq f^*(M_Y)$, this is the second condition (b). The second condition in (a), for $y_1 \to y_0$ means that for $M$ in $\mathcal{A}$,
\[
H^*_{\D}(X/y_0, M) \longrightarrow H^*_{\D}(X/y_1, M)
\]
is \texttt{[an]} isomorphism, now this arrow is interpreted as
\[
f_*(M_X)_{y_0} \longrightarrow f_*(M_X)_{y_1} \quad ,
\]
that this is always an isomorphism, this is the first condition (b). That (b) also holds for $F$ locally constant on $Y$ is immediate, since this is a condition of a local nature with respect to $Y$.

\msp{page 77}

The dual proposition calls on the functors
\[
y_0\bs X \longrightarrow y_1\bs X \qquad \text{and} \qquad X_y \longrightarrow y\bs X \quad ,
\]
and the sheaves $f_!f^*M$, for $M$ constant, or locally constant on $Y$.
% foot: Version 8 juillet 2001

\printedpage{47}
% head: A. Grothendieck : Dérivateurs, Ch. I, Édition des Universités de Montpellier II et Paris VII
\textbf{Definition 1.} If the equivalent conditions of proposition 1 are satisfied, we say that $f$ is a $\D$-\textit{cohomological fibration} (or $\D$-\textit{left fibration}). If the dual conditions are satisfied, we say that it is a $\D$-\textit{homological fibration} (or $\D$-\textit{right fibration}).

N.B. Cohomological $\D$-fibrations are stable under any base change which is a local isomorphism. Homological $\D$-fibrations under any base change which is a local co-isomorphism. But these are not notions stable under \textit{every} base change\footnote[34]{? give example!}. But the notion of a cohomological $\D$-fibration is stable under $\D$-smooth base change, and that of a homological $\D$-fibration under $\D$-proper base change.

\hspace*{\parindent}I am going to try to construct an example where base change gets botched, by setting $Y = \boldsymbol{\Delta}_1$. Giving $X$ over $\boldsymbol{\Delta}_1$ amounts to giving an open subcategory $X_0$, and its closed complement $X_1$
\[
\begin{tikzcd}
& X & \\
X_1 \arrow[ur, hook] \arrow[d, dash, dashed] & & X_0 \arrow[ul, hook'] \arrow[d, dash, dashed] \\
1 & & 0 \arrow[ll]
\end{tikzcd}
\quad .
\]
To say that $X \to Y = \boldsymbol{\Delta}_1$ is a cohomological $\D$-fibration simply means
\msp{page 78}
that the restrictions
\[
X_0 \overset{j}{\longrightarrow} X \quad , \qquad X_1 \overset{i}{\longrightarrow} X
\]
are $\D$-equivalences. Let us take a base change $Y' \to Y = \boldsymbol{\Delta}_1$, \textit{i.e.} a $Y'$ with $Y'_0, Y'_1$
\[
\begin{tikzcd}
& Y' & \\
Y'_1 \arrow[ur, hook] \arrow[d, dash, dashed] & & Y'_0 \arrow[ul, hook'] \arrow[d, dash, dashed] \\
1 & & 0 \arrow[ll]
\end{tikzcd}
\quad .
\]
We see that the condition $Y'$ smooth over $\boldsymbol{\Delta}_1$ means that for every $y'_1 \in \Ob Y'_1$, the category $Y'_0/y'_1$ of objects $y_0$ of $Y'_0$ over $y'_1$ (\textit{i.e.} equipped with $y_0 \to y'_1$) is contractible. We see that if $Y'_1 \neq \emptyset$ then $Y'_0 \neq \emptyset$ -- but an example with $Y'_0 = \emptyset$, or $Y'_1 = \emptyset$, is not interesting
\DERmark{ENV0034}

% foot: Version 8 juillet 2001
\printedpage{48}
% head: A. Grothendieck : Dérivateurs, Ch. I, Édition des Universités de Montpellier II et Paris VII
(for a projection $I \times J \to I$ is always a cohomological $\D$-fibration). Let us take for example for $Y'$
\[
\begin{tikzcd}
& a \arrow[dl] \\
b & \\
& a' \arrow[ul] \\
\underbrace{\qquad}_{\displaystyle Y'_1} & \underbrace{\qquad}_{\displaystyle Y'_0}
\end{tikzcd}
\quad ,
\]
an amalgamated sum of two copies of $Y = \boldsymbol{\Delta}_1$, hence $X' = X \times_Y Y'$ can be viewed as an amalgamated sum $X \amalg_{X_1} X$. But that seems \textit{not} to be a counterexample. We would need $X_1 \to X'$ \textit{not} to be a $\D$-equivalence.

\msp{page 79}

\textbf{Proposition 2.} \textit{Let $f : X \to Y$ be in} \Dia. \textit{Equivalent conditions :}

\hspace*{\parindent}a) \textit{For every $Y'' \overset{p}{\to} Y' \to Y$, whence}
\[
\begin{tikzcd}
X'' \arrow[r, "q"] \arrow[d] \arrow[dr, phantom, "\text{cart}"] & X' \arrow[r] \arrow[d] \arrow[dr, phantom, "\text{cart}"] & X \arrow[d, "f"] \\
Y'' \arrow[r, "p"'] & Y' \arrow[r] & Y
\end{tikzcd}
\quad ,
\]
\textit{if $p$ is a $\D$-equivalence,} \texttt{[then]} \textit{so is $q$.}\footnote[35]{One must assume $Y', Y''$ \textsf{Hot}-aspheric, cf. page 96, and assume axiom Der 3 $a'$) of page 97 to be satisfied.}

\hspace*{\parindent}b) \textit{$f$ is a ``universal'' cohomological $\D$-fibration,} i.e. \textit{is a cohomological fibration, and remains so after any base change.}

\hspace*{\parindent}b$'$) \textit{Condition dual to} (b) \textit{: $f$ is a universal homological $\D$-fibration.}\footnote[36]{b) $\Rightarrow$ a) \underline{botched}, \textit{cf.} marginal annotation,\footnotemark[37] \textit{cf.} counterexample page 89. But it works for the derivators $\D_i$ ($i \in \mathbb{N}$) associated with the ``basic localizers'' $W_i$, \textit{i.e.} with the theory of homotopy $i$-types.}\footnotetext[37]{The implication b) $\Rightarrow$ a) not proved, only in a weaker form :\par \hspace*{\parindent}a$'$) if $Y''$, $Y'$ are $\D$-aspheric and 1-connected, then $X'' \to X'$ is a $\D$-equivalence. (It even suffices to assume it for $Y'$, $Y''$ having final objects (to imply b) or having initial objects (to imply b$'$).)}

\textit{a}) $\Rightarrow$ \textit{b}) : Let us assume \textit{a}), let $Y' \to Y$ be a base change, whence
\[
\begin{tikzcd}
X' \arrow[r] \arrow[d, "f'"'] & X \arrow[d, "f"] \\
Y' \arrow[r] & Y
\end{tikzcd}
\quad .
\]
% foot: Version 8 juillet 2001


\endgroup
% END UNIT Der-ChI-234/Der-ChI-234_p046-050.tex
% BEGIN UNIT Der-ChI-234/Der-ChI-234_p051-055.tex
\begingroup\makeatletter%
\newcommand{\printedpage}[1]{\marginpar{\footnotesize p.~#1}}%
\newcommand{\msp}[1]{\marginpar{\footnotesize [#1]}}%
\newcommand{\DD}{\mathbb{D}}%
\newcommand{\Ob}{\mathsf{Ob}}%
\newcommand{\Cat}{\mathcal{C}at}%
\newcommand{\Dia}{\mathsf{Dia}}%
\makeatother\let\printedpage\derpp\let\msp\dermsp%
\let\mspin\DERmspin%


\printedpage{49}
% head: A. Grothendieck : Dérivateurs, Ch. I, Édition des Universités de Montpellier II et Paris VII
% foot: Version 8 juillet 2001

To prove that $f'$ is a cohomological $\DD$-fibration, we look over $Y'$ at the functors of the form $\{y'\} \to Y'/y'$ ($y' \in \Ob\, Y'$) and $Y'/y'_1 \to Y'/y'_0$ (if $u' : y'_1 \to y'_0$ an arrow of $Y'$), and we see that the arrows in $\Cat/Y'$ deduced from these by the base change $X' \xrightarrow{f'} Y'$ remain $\DD$-equivalences

\msp{page 80}

\[
\begin{tikzcd}[column sep=large, row sep=large]
X'_1 \arrow[r, "q"] \arrow[d, "f'_1"'] \arrow[dr, phantom, "\text{cart}"] & X'_0 \arrow[r] \arrow[d, "f'_0"] \arrow[dr, phantom, "\text{cart}"] & X' \arrow[r] \arrow[d, "f'"] \arrow[dr, phantom, "\text{cart}"] & X \arrow[d, "f"] \\
Y'_1 \arrow[r, "p"'] & Y'_0 \arrow[r] & Y' \arrow[r] & Y & .
\end{tikzcd}
\]
Since $p$ is a $\DD$-equivalence, this indeed follows from a).

\indent\textit{b}) $\Rightarrow$ \textit{a}): As condition \textit{b}) is obviously stable under base change, we are reduced to proving that if $f$ satisfies \textit{b}), and if we have
\[
\begin{tikzcd}[column sep=large, row sep=large]
X \arrow[d, "f"'] \arrow[dr, phantom, "\text{cart}"] & X' \arrow[l, "q"'] \arrow[d, "f'"] \\
Y & Y' \arrow[l, "p"]
\end{tikzcd}
\tag*{$(*)$}
\]
if $p$ is a $\DD$-equivalence, then $q$ too, \textit{i.e.}
\[
\begin{tikzcd}[row sep=small]
H^*_{\DD}(X,M) \arrow[r, "\sim"] & H^*_{\DD}(X',M) \\
\wr| & \wr| \\
H^*_{\DD}(Y, f_*(M_X)) & H^*_{\DD}(Y', f'_*(M_{X'}))
\end{tikzcd}
\]
for every $M$ in $\mathcal{A} = \DD(e)$. But $F = f_*(M_X)$ is locally constant and is computed ``fiber by fiber'', likewise for $F' = f'_*(M_{X'})$.

\msp{page 81}

The equivalence \textit{a}) $\Leftrightarrow$ \textit{b}$'$) follows from \textit{a}) $\Leftrightarrow$ \textit{b}) by duality (condition \textit{a}) being self-dual), or by repeating the same argument for \textit{a}) $\Rightarrow$ \textit{b}$'$), and for \textit{b}$'$) $\Rightarrow$ \textit{a}) too (replacing cohomology by homology).

\textbf{Definition 2.} If the equivalent conditions of proposition 2 are satisfied, we say that $f$ is a $\DD$-\textit{fibration}.

N.B. This is a notion \textit{both} local and ``colocal'' on $Y$ (whereas left fibration is local, right fibration colocal).

\printedpage{50}
% head: A. Grothendieck : Dérivateurs, Ch. I, Édition des Universités de Montpellier II et Paris VII
% foot: Version 8 juillet 2001

\textbf{Proposition 3.} a) \textit{The notion of $\DD$-fibration is stable under base change and under composition.}\footnote[38]{Stability under composition \textit{is} not clear.}

\indent b) \textit{Let}
\[
\begin{tikzcd}[column sep=normal, row sep=large]
X \arrow[rr, "f"] \arrow[dr, "p"'] & \phantom{X} & X' \arrow[dl, "p'"] \\
& Y &
\arrow[from=1-2, to=2-2, phantom, "\text{\footnotesize comm}"]
\end{tikzcd}
\]
\textit{with $p,p'$ $\DD$-fibrations (it suffices that $p,p'$ be cohomological $\DD$-fibrations, or that they be homological $\DD$-fibrations, or that $p,p'$ be proper, or $p,p'$ smooth). If $f$ induces $\DD$-equivalences in the fibers, then $f$ is a $\DD$-equivalence.}\footnote[39]{Can be made considerably more precise in the case where $p$, $p'$ are $\DD$-fibrations, \textit{cf.} prop. 6, page 84.}

\indent\textsc{Proof.} \textit{a}) is immediate under condition \textit{a}) (prop. 2) of the notion of $\DD$-fibration. For \textit{b}), suppose that $f$ induces $\DD$-equivalences on the fibers. We want to prove that % [sic] On veux -> On veut, verb agreement
\[
H^*_{\DD}(X',M) \longrightarrow H^*_{\DD}(X,M)
\]
\DERmark{ENV0035}

is an isomorphism, for every $M$ in $\mathcal{A}$, but this arrow is identified with
\[
H^*_{\DD}(Y, p'_*(M_{X'})) \longrightarrow H^*_{\DD}(Y,p_*(M_X))
\]
\DERmark{ENV0036}

deduced from
\[
p'_*(M_{X'}) \longrightarrow p_*(M_X) \qquad\qquad (\text{NB}\ \ M_X \simeq f^*M_{X'}\,)
\]
\DERmark{ENV0037}

\msp{page 82}
and it suffices to prove that this arrow is an isomorphism. For this, we look fiber by fiber, and we use the fact that ``direct images are computed fiber by fiber''.

\textbf{Question.} Let $f : X \to Y$ be a $\DD$-fibration, and let $F \in \DD(Y)$. Then is $f_*f^*F$ computed by passing to the fibers? Same question for $f_!f^*F$. For the first property, concerning $f_*$, we show that it works if $F$ is of the form $p_*(M_{Y'})$, where $p : Y' \to Y$ is $\DD$-\textit{proper,} and $M \in \Ob\, \mathcal{A}$. On the other hand, the full subcategory $C$ of $\DD(Y)$ consisting of the $F$ for which it works is (unless I am mistaken) ``stable under $\DD$-$\mathsf{hot}$-$\varprojlim$'' in the following sense: if we have a category $I$ with final object $e$ forming a closed subcategory (\textit{i.e.} $e \to x \Rightarrow x = e$), and if we have $\mathcal{F} \in \DD(Y \times I)$, whence $I^\circ \to \DD(Y)$, $i \mapsto \mathcal{F}_i$ (skeletal diagram), if the $\mathcal{F}_i$ for $i \neq e$ are in $C$, and if $\mathcal{F}$ is in the essential image of $j_*$, where $j : Y \times (I \setminus \{e\}) \hookrightarrow Y \times I$ is the inclusion, then $\mathcal{F}_e$ is also in $C$.
\DERmark{ENV0038}


\indent We may however consider that if we want $f : X \to Y$ to satisfy all the cohomological and homological properties that attach to the notion of a ``fibration'' (``locally trivial''), we should put among the conditions that $f$ is \textit{smooth}, and \textit{proper}, so that base changes by $f$ commute with the functors $g_*$, $g_!$. % [sic] on veux -> on veut, verb agreement

\printedpage{51}
% head: A. Grothendieck : Dérivateurs, Ch. I, Édition des Universités de Montpellier II et Paris VII
% foot: Version 8 juillet 2001

\msp{page 83}

Let us point out:

\textbf{Proposition 4.} \textit{Let $f : X \to Y$ be} proper. \textit{For $f$ to be a $\DD$-fibration, it is necessary and sufficient that for every $y_1 \to y_0$ in $Y$, the induced functor $X/y_1 \to X/y_0$ be a $\DD$-equivalence} (i.e. \textit{that $f$ be a} cohomological \textit{$\DD$-fibration}). \textit{Dually, if $f$ is $\DD$-smooth, for it to be a $\DD$-fibration, it is necessary and sufficient that for every $y_1 \to y_0$ in $Y$, $y_0\backslash X \to y_1\backslash X$ be a $\DD$-equivalence.}

\indent One could call \textit{perfect} $\DD$-fibrations the functors both proper and smooth that are cohomological $\DD$-fibrations (or what amounts to the same thing homological $\DD$-fibrations).

\textbf{Corollary.} \textit{Suppose $f$ a $\Cat$-cofibration (hence proper). For $f$ to be a $\DD$-fibration, it is necessary and sufficient that the functors $X_{y_0} \to X_{y_1}$ (for $y_0 \to y_1$ in $Y$) be $\DD$-equivalences. Dually, for $\Cat$-fibration (hence smooth).} % [sic] u_1 -> y_1, the arrow is y_0 -> y_1 ; [sic] Duallement -> Dualement, spelling
\DERmark{ENV0039}


\textbf{Proposition 5.} \textit{Let $f : X \to Y$ be in} $\Dia$. \textit{Equivalent conditions:}
\begin{enumerate}[label={}, leftmargin=2em, labelsep=0pt]
\item a) \textit{$f$ is a $\DD$-equivalence, and remains so under every base change;}
\item b) \textit{$f$ is a $\DD$-fibration with $\DD$-contractible fibers.}
\end{enumerate}

\textsc{Proof.} \textit{a}) $\Rightarrow$ \textit{b}): That \textit{a}) implies that $f$ is a $\DD$-fibration is immediate from its definition (condition \textit{a}) of proposition 2), and that the fibers are contractible is tautological.

\textit{b}) $\Rightarrow$ \textit{a}): It suffices to show that \textit{b}) implies that $f$ is a $\DD$-equivalence. But this is a special case of proposition 3 \textit{b}) (taking $X' = Y$ there).

\msp{page 84}

When these conditions are satisfied, we expect that the functor
\[
f^* : \DD(Y) \longrightarrow \DD(X)
\]
is fully faithful, \textit{i.e.} that the adjunction morphisms % [sic] les morphisme -> les morphismes, plural agreement
\[
F \longrightarrow f_*f^*F \quad , \qquad f_!f^*F \longrightarrow F % [sic] f_!f^* -> f_!f^*F, the object F is missing
\]
\DERmark{ENV0040}

in $\DD(Y)$ are isomorphisms. (This is true if $f$ is $\DD$-proper or $\DD$-smooth, cf. cor. 3, p.71.) This again amounts to saying that the $f_*f^*F$ are computed by passing to the fibers, or again likewise for the $f_!f^*F$. This brings us back to the question on page 82.

\textbf{Definition 3.} We then say that $f$ is a \textit{trivial} $\DD$-fibration.

\textbf{Proposition 6} (makes prop. 3 page 81 more precise). \textit{Let us consider a commutative triangle in} $\Dia$
\[
\begin{tikzcd}[column sep=small, row sep=large]
X \arrow[rr, "f"] \arrow[dr, "p"'] & & X' \arrow[dl, "p'"] \\
& S & & & .
\end{tikzcd}
\]

\printedpage{52}
% head: A. Grothendieck : Dérivateurs, Ch. I, Édition des Universités de Montpellier II et Paris VII
% foot: Version 8 juillet 2001

\indent a) \textit{Suppose that $p$ and $p'$ are left $\DD$-fibrations, or that $p$ and $p'$ are right $\DD$-fibrations. Let $S_0$ be a subset of $S$ that meets every connected component of $S$. For $f$ to be a $\DD$-equivalence, it} suffices \textit{that for every $s_0 \in S_0$, $f$ induce a $\DD$-equivalence $X_{s_0} \to X'_{s_0}$.}
\DERmark{ENV0041}


\indent b) \textit{This condition is also necessary}\footnote[40]{Not proved!\footnotemark[41] Even if $X' = S$, it is not clear, even if $\DD = \mathsf{HOTAB}$, $f$ proper and smooth. But OK if $S$ is completely $\DD$-aspheric.}\footnotetext[41]{\underline{False}, \textit{cf.} counterexample, p. 87.} \textit{if $p$ and $p'$ are $\DD$-fibrations. So in this case, $f$ is a $\DD$-equivalence if and only if $f$ induces $\DD$-equivalences fiber by fiber.}
\DERmark{ENV0042}


\msp{page 85}

\textsc{Proof.} a) It suffices, by virtue of proposition 3, to show that the hypothesis implies that for \textit{every} $s \in S$, $f$ induces a $\DD$-equivalence $f_s : X_s \to X'_s$. This follows at once from the lemma:

\textbf{Lemma.} \textit{If $p$ and $p'$ are both either left $\DD$-fibrations or right $\DD$-fibrations, and if $s_1 \to s_0$ is an arrow in $S$, then $f_{s_1}$ is a $\DD$-equivalence if and only if $f_{s_0}$ is. (Thus the set $S' \subset \Ob\, S$ of $s \in S$ such that $f_s$ is a $\DD$-equivalence is a union of connected components of $S$.)}

\indent Suppose for example that $p,p'$ are left $\DD$-fibrations. Let us consider, for $s \in S$, the diagram
\[
\begin{tikzcd}[row sep=large]
X_s \arrow[r, "f_s"] \arrow[d] & X'_s \arrow[d] \\
X/s \arrow[r, "f/s"'] & X'/s & .
\end{tikzcd}
\]
The vertical arrows are $\DD$-equivalences by the hypothesis on $p,p'$, so $f_s$ is a $\DD$-equivalence if and only if $f/s$ is. Let us now consider
\[
\begin{tikzcd}[row sep=large]
X/s_1 \arrow[r, "f/s_1"] \arrow[d] & X'/s_1 \arrow[d] \\
X/s_0 \arrow[r, "f/s_0"'] & X'/s_0 & .
\end{tikzcd}
\]
Here again, the hypothesis on $p,p'$ implies that the vertical arrows are $\DD$-equivalences, so $f_{s_1}$ is a $\DD$-equivalence if and only if $f_{s_0}$ is.

\msp{page 86}

If $X' = S$, the statement (not proved) says that if $f : X \to S$ is a $\DD$-fibration and a $\DD$-equivalence, then the fibers are $\DD$-contractible. This is not certain even for a projection $X = T \times S \to S$. If $\eta = \mathit{hot}_0(S)$, $\xi = \mathit{hot}_0(T)$, the question reduces to this one: If for every $M$ in $\mathcal{A}$, we have: the homomorphism
\DERmark{ENV0043}

\[
\alpha_M : \xi \times_{\DD} M \longrightarrow M
\]
\printedpage{53}
% head: A. Grothendieck : Dérivateurs, Ch. I, Édition des Universités de Montpellier II et Paris VII
% foot: Version 8 juillet 2001
satisfies
\[
\eta \times_{\DD} \alpha_M : \eta \times_{\DD} (\xi \times_{\DD} M) \xrightarrow{\sim} \eta \times_{\DD} M
\]
is an isomorphism, then $\alpha_M$ is itself an isomorphism.

\msp{page 87}

Here is a \textit{counterexample}\footnote[42]{See an even simpler counterexample on pages 93, 94. (But it has the drawback of not having connected fibers.) % [sic] inconvenient -> inconvénient, missing accent
} which I first give in the context of topological spaces, instead of $\Cat$. I am going to construct a locally trivial and nontrivial fibration
\[
\begin{tikzcd}[column sep=0.2em, row sep=large]
X \arrow[d, "f"'] & \supset S^1 & \\
S & = S^1 & \qquad ,
\end{tikzcd}
\]
with fiber $S^1$, such that $f$ is a $\DD$-equivalence for every $\DD = \mathsf{Der}(k\text{-}\mathsf{Mod})$, $k$ being a ring in which 2 is invertible. This means that
\[
H_*(X,k) \longrightarrow H_*(S,k)
\]
is an isomorphism modulo 2-torsion (kernel and cokernel are annihilated by a power of 2 - in fact by 2 here), or again $H_*(X,k) \xrightarrow{\sim} H_*(S,k)$ if 2 invertible in $k$. But the fiber $S^1$ is not $\DD$-contractible (unless $k = 0$!).\footnote[43]{Much simpler example still: we take $X = S^1$, and $X \to S$ the two-sheeted covering. If we take the $n$-sheeted covering, it is also OK if $k$ a ring in which $n$ is invertible.}

\indent The fibration is given up to isomorphism, by fixing a point $s \in S$, and giving oneself an automorphism $u$ of the fiber $X_s$ (``monodromy''). (The isomorphism class of the bundle depends only on the isotopy class of $u$). I take for $u$ the symmetry (reversing orientation) $z \mapsto z^{-1}$ of the standard circle $\{z \in \mathbb{C} \mid |z| = 1\}$. This fibration ``is realized'' by ordered sets $X \to S$, $S^1$ being represented (on the base, as on the fibers) by % [sic] Je prend -> Je prends, verb agreement ; [sic] ordonés -> ordonnés, spelling
\DERmark{ENV0044}

\[
\begin{tikzcd}[column sep=normal, row sep=normal]
& & a & & \\
S \ = & 0 \arrow[ur] \arrow[dr] & & 0' \arrow[ul] \arrow[dl] & = \ S' \amalg_R S'' \quad , \\
& & b & &
\end{tikzcd}
\]
where $R$ is the discrete category having two elements $a$, $b$. We define $X$ as $X' \amalg_Y X''$, where $X' = X|S'$, $X'' = X|S''$, $Y = X|R$.\footnote[44]{N.B. $S'$, $S''$, $R = S' \cap S''$ closed in $S$, $X'$, $X''$, $Y = X' \cap X''$ closed in $X$.}


\endgroup
% END UNIT Der-ChI-234/Der-ChI-234_p051-055.tex
% BEGIN UNIT Der-ChI-234/Der-ChI-234_p056-060.tex
\begingroup\makeatletter%
\newcommand{\printedpage}[1]{\marginpar{\footnotesize p.~#1}}%
\newcommand{\msp}[1]{\marginpar{\footnotesize [#1]}}%
\newcommand{\Cat}{\mathcal{C}\mathit{at}}%
\newcommand{\D}{\mathbb{D}}%
\newcommand{\RR}{\mathbb{R}}%
\makeatother\let\printedpage\derpp\let\msp\dermsp%
\let\mspin\DERmspin%


\printedpage{54}
% head: A. Grothendieck : Dérivateurs, Ch. I, Édition des Universités de Montpellier II et Paris VII

\msp{page 88}
We put
\[ X' = S' \times \bar S \qquad\qquad X'' = S'' \times \bar S \]
($\bar S$ an exact copy of $S$, cf. below) and, to glue $X'$ and $X''$, we must give ourselves an isomorphism between $X'|R \simeq R \times \bar S = \bar S \amalg \bar S$ and $X''|R \simeq R \times \bar S = \bar S \amalg \bar S$ which is an $R$-functor, \textit{i.e.} which comes from two isomorphisms
\[
\begin{aligned}
\bar S &\xrightarrow{u_a} \bar S \\
\bar S &\xrightarrow{u_b} \bar S \qquad .
\end{aligned}
\]
We take
\[
\begin{aligned}
u_a &= \mathsf{id}_{\bar S} \\
u_b &= \text{reflection about the axis } \bar 0, \bar 0' \ \textit{i.e.} % Doubt resolved: PDF leaf 18 (printed p.54) reads u_b, with plain b.
\\
&\phantom{=}\quad u_b \text{ leaves fixed } \bar 0, \bar 0' \text{ and interchanges } \bar a \text{ and } \bar b.
\end{aligned}
\]
(It is clearer if we disidentify the base and the fiber, taking fiber
\[
\bar S \quad = \quad
\begin{tikzcd}
 & \bar a & \\
\bar 0 \arrow[ur] \arrow[dr] & & \bar 0' \arrow[ul] \arrow[dl] \\
 & \bar b &
\end{tikzcd}
\qquad ).
\]

To see that $X$ is a $\D$-fibration over $S$ (whatever $\D$ is) let us prove that it is \textit{proper}\footnote[45]{N.B. $f : X \to S$ is also \textit{smooth}, for it is a trivial bundle (a product) over two opens $S - \{a\}$, $S - \{b\}$ which cover $S$. Better: $f$ is both $\Cat$-fibrant and $\Cat$-cofibrant (for this condition uses only the properties of the $X \times_S \boldsymbol{\Delta}_2 \to \boldsymbol{\Delta}_2$, for $\boldsymbol{\Delta}_2 \to S$).} over $S$ - we are reduced to seeing that $X' \to S'$ and $X'' \to S''$ are proper, which is trivial. To see that it is a $\D$-fibration, we are likewise reduced to the case of $X' \to S'$ and $X'' \to S''$, where this is also trivial. We have used the: % [sic] quelque soit -> quel que soit, spelling

\textbf{Lemma.} \textit{Let $f : X \to S$, $(S_i)_{i \in I}$ a family of closed subcategories covering $S$. For every $i$, let $X_i = X \times_S S_i \xrightarrow{f_i} S_i$. For $f$ to be proper, it is necessary and sufficient that the $f_i$ be so. For $f$ to be a proper $\D$-fibration it is necessary and sufficient that the $f_i$ be so.}

\msp{page 89}
Here is the little computation which shows that $f : X \to S$ is a $\D$-equivalence. We do the computation in cohomology.\footnote[46]{But in the case of the 2-sheeted covering of $S$, there is not even any computation to do!} We must prove that
\[ H^*(S,k) \longrightarrow H^*(X,k) \tag*{$(*)$} \]
\printedpage{55}
% head: A. Grothendieck : Dérivateurs, Ch. I, Édition des Universités de Montpellier II et Paris VII
is an isomorphism (and likewise for any constant coefficient $M$ on $S$). Now
\[ H^*(X,k) \overset{\mathrm{def}}{=} \RR\Gamma_X(k_X) = \RR\Gamma_S(\RR f_*(k_X)) \quad , \]
and $(*)$ is induced by
\[ i : k_S \longrightarrow \RR f_* k_X \quad . \]
Now we have
\[ \RR f_* k_X \simeq \underbrace{k_S}_{= \mathsf{Im}(i)} \quad \oplus \quad \underbrace{k_S[1]}_{k_S \text{ in degree } 1}{}^{\underline{\mathrm{tordu}}} \]
namely ``twisted'' by the monodromy action on $S = S^1$, which is the symmetry $\lambda \mapsto -\lambda$. We have
\[ \RR\Gamma_S(k_S[1]^{\mathrm{tordu}}) = H^1(\underbrace{\pi_1(S)}_{= \mathbb{Z}}\, , k_S^{\mathrm{tordu}}) = k_{(\mathbb{Z})} = k/(u - \mathsf{id})k = k/2k = 0 \ , \]
hence the result.

Let us note that if we look at
\[ X' = X \times_S X \xrightarrow{f'} S \quad , \]
we have by relative Künneth
\[ \RR f'_*(k_{X'}) \simeq \RR f_*(k_X) \otimes^{\mathbb{L}} \RR f_*(k_X) \simeq k_S \oplus k_S[1]^t \oplus k_S[1]^t \oplus \underbrace{k_S[1]^t \otimes k_S[1]^t}_{\simeq k_S[2] \text{ untwisted}} \ , \]
hence
\[ H^*_{\D}(X',k) \simeq \RR\Gamma_{X'}(k_{X'}) \simeq \RR\Gamma_S(\RR f'_*(k_{X'})) \simeq \RR\Gamma_S(k_S) \oplus \RR\Gamma_S(k_S)[2] \]
and
\[ H^*_{\D}(S,k) \longrightarrow H^*_{\D}(X',k) \]
is \textit{not} an isomorphism.

\msp{page 90}
Thus $f : X \to S$ which is a $\D$ fibration (proper and smooth) does satisfy the equivalent conditions $b$), $b'$) of proposition 2 (page 79), but not condition $a$), since the base change
\[
\begin{tikzcd}[column sep=large, row sep=large]
X' = X \times_S X \arrow[r, "{q = pr_1}"] \arrow[d, "{f' = pr_2}"'] & X \arrow[d, "f"] \\
S' = X \arrow[r, "{p = f}"'] & S
\end{tikzcd}
\]
is a $\D$-equivalence $S' \to S$, whereas $X' \to X$ is not.

Here is another counterexample, on the same principle but a little simpler. Let $G$ be a finite group (for example $G = \mathbb{Z}/2\mathbb{Z}$), let us make it act on an object $Z$ of $\Cat$, let $S = [G] = B_G$ be the category reduced to one object $a$, with $G = \mathsf{End}(a)$. Let $E_G$ be the category (isomorphic to $(B_G)/a$)
\[ \mathsf{Ob}\ E_G = G \quad , \qquad \mathsf{Hom}_{E_G}(g,g') = \{ h \in G \mid g = g'h \} \quad , \]
obvious composition - then $E_G$ is a trivial groupoid (1-connected). The group $G$ acts freely on it on the left by translations, and $B_G = E_G/G$. % [sic] goupe -> groupe, misprint

\printedpage{56}
% head: A. Grothendieck : Dérivateurs, Ch. I, Édition des Universités de Montpellier II et Paris VII

\msp{page 91}
If $G$ acts on $Z$ on the left, we can then use it to twist $Z$, by
\[ X = E_G \times_G Z = (E_G \times Z)/\text{``diagonal'' actions of G} \quad . \]
So we have a cartesian diagram
\[
\begin{tikzcd}
 & X' \arrow[d, equal] & X'' \arrow[d, equal] \\
X \arrow[d, "f"'] & Z \times E_G \arrow[l] \arrow[d, "pr_2"] & Z \arrow[l, dashed, "{\overset{\D}{\sim}}"'] \arrow[d] \\
B_G \arrow[d, equal] & E_G \arrow[l] \arrow[d, equal] & e \arrow[l, dashed, "{\overset{\D}{\sim}}"'] \arrow[d, equal] \\
S & S' & S''
\end{tikzcd}
\]
Since the base change $S' = E_G = (B_G)/a \to B_G = S$ transforms $X$ over $B_G$ into a product (which is a proper and smooth $\D$-bundle over $S'$), we see that $X$ over $B_G$ is a proper and smooth $\D$-bundle (for this is a local condition on $B_G$).

We now suppose that $G$ is finite of order $n > 1$, and let $k$ be a ring in which $n$ is invertible, $\D = \mathsf{Der}(k - \mathsf{Mod})$. Then $E_G \to B_G$ is a $\D$-equivalence, or, what amounts to the same thing, $e \to B_G$ is.\footnote[47]{In fact, $e$, $E_G$, $B_G$ are $\D$-aspheric.} Yet, after base change either by $S' = E_G$, or by $S'' = e$, $X' \to X$ (or what amounts to the same thing, $X'' \to X$) has hardly any reason to be one too. We have in cohomology % [sic] invesible -> inversible, misprint
\[ \RR\Gamma_X(k_X) \simeq \RR\Gamma_{BG}(f_*(k_X)) \simeq \RR\Gamma_G(G, H^*(Z,k)) \]
\msp{page 92}
therefore ($n$ being invertible in $k$):
\[ H^*(X,k) = \RR\Gamma_X(k_X) \simeq H^*(Z,k)^G \longrightarrow H^*(Z,k) \quad , \]
this is not in general an isomorphism. We can take for example $G = \mathbb{Z}/2\mathbb{Z}$, $Z \sim S^1$ as before, and $G$ acting by reflection of $S^1$ as previously.

This example is particularly convincing, for it is an example \textit{with $S$, $S'$ $\D$-aspheric,} and $S' \to S$ a local isomorphism (in fact, a principal covering of $S$). It is also an example of a smooth and proper $\D$-bundle $X \to S$, over an \textit{aspheric} base $S$, and such that the inclusion of a fiber in $X$ % [sic] convainquant -> convaincant, spelling of the adjective
\[ X_s \longrightarrow X \]
is \textit{not} a $\D$-equivalence.

The only complaint we can make about this example, in comparison with the preceding one, is that the categories $X, S$ are not ordered sets. And $S$ cannot be a finite ordered set, for the classifying space of a finite group has cohomology in all dimensions (unless I am mistaken).

\printedpage{57}
% head: A. Grothendieck : Dérivateurs, Ch. I, Édition des Universités de Montpellier II et Paris VII

\msp{page 93}
But we can take inspiration from this example to modify it so that it carries over to ordered sets. It suffices to realize, for given $G$:

\textit{a}) A \textit{free} action of $G$ on an object $E$ of $\Cat$, such that $H^*(E,k)^G \simeq k$, \textit{i.e.} $B = E/G$ is $\D$-aspheric.

\textit{b}) An action of $G$ on an object $Z$, such that the inclusion
\[ H^*(Z,k)^G \longrightarrow H^*(Z,k) \] % [sic] H^*(Z,k)^G \longrightarrow H^*(Z,k)^G -> H^*(Z,k)^G \longrightarrow H^*(Z,k), an inclusion of the invariants (cf. p. 56)
\DERmark{ENV0045}

is not \texttt{[an]} isomorphism.

\textit{c}) Choose for $E$, $Z$ ordered sets, and make sure that $B = E/G$ is also represented by an ordered category. Unless I am mistaken, the same will hold for $X = E \times Z/G$. We will have in any case % [sic] en tout les cas -> en tous les cas, agreement
\[
\begin{tikzcd}[column sep=large]
 & Z \times E \arrow[d, equal] & \\
X \arrow[d] & X' \arrow[l] \arrow[d] & Z \arrow[l] \arrow[d] \\
S \arrow[d, equal] & S' \arrow[l] \arrow[d, equal] & S'' = \{s\}\, , \quad s \in E \ . \arrow[l] \\
E/G & E & \text{\footnotesize asph} \\
\text{\footnotesize asph} & \substack{\text{asph} \\ \text{if we can}} &
\end{tikzcd}
\]
This gives another example with $X$ a proper and smooth $\D$-bundle, over a $\D$-aspheric base $S$, such that the inclusion of a fiber $X_s$ in $X$ is not a $\D$-equivalence. We can even manage to have $H^*(Z,k)^G \simeq k$ which implies
\msp{page 94}
that $X$ too is aspheric (but the fibers of $X \to S$ are not $\D$-aspheric). This therefore gives another counterexample to b) of prop. 6 (p. 84).

Here is the most economical way, I think, to realize the situation I have spoken of. Topologically, I take % [sic] je prend -> je prends, conjugation
\[
\begin{aligned}
G &= \{\pm 1\} \\
E &= S^2 &\qquad& \text{Euclidean sphere, on which } \{\pm 1\} \\
  &      &&      \text{acts by antipodism } x \mapsto -x \\
B &= S^2/\{\pm 1\} && \text{real projective plane} \\
Z &= S^1 && \text{where G acts as before by the reflection\,.}
\end{aligned}
\]
So this is an example very close to the one at the beginning (p. 87), where we could also consider the base ($B = S^1$) as the quotient of an $E$ ($= S^1$) by $G = \{\pm 1\}$ acting freely (by
\printedpage{58}
% head: A. Grothendieck : Dérivateurs, Ch. I, Édition des Universités de Montpellier II et Paris VII
the antipodism $z \mapsto -z$ in $\mathbb{C} \simeq \RR^2$). But in this example, the base was not $\D$-aspheric.

The most economical way of realizing the sphere $S^2$ by an ordered set is probably by its ``globular'' decomposition into three pairs of cells of dimension 0, 1, 2, giving the ordered set
\[
I \qquad = \qquad
\begin{tikzcd}[row sep=large, column sep=large]
S_2^- & S_2^+ \\
S_1^- \arrow[u] \arrow[ur] & S_1^+ \arrow[u] \arrow[ul] \\
S_0^- \arrow[u] \arrow[ur] & S_0^+ \arrow[u] \arrow[ul]
\end{tikzcd}
\]
and the antipodism exchanges $S_i^-$ and $S_i^+$, but the quotient is not ordered, it is the category (important for the theory of 2-categories):
\msp{page 95}
\[
\begin{tikzcd}[row sep=large]
S_2 \\
S_1 \arrow[u, shift left, "p_1^-"] \arrow[u, shift right, "p_1^+"'] \\
S_0 \arrow[u, shift left, "p_0^-"] \arrow[u, shift right, "p_0^+"']
\end{tikzcd}
\qquad\qquad \text{with} \qquad\qquad
\begin{cases}
p_1^- p_0^- = p_1^+ p_0^+ \\
p_1^- p_0^+ = p_1^+ p_0^-
\end{cases}
\]
\DERmark{ENV0046}\DERmark{ENV0047}

But let us take the barycentric subdivision of the preceding cellular decomposition of $S$ whose vertices correspond to the elements of $\mathsf{Ob}\ I$, the simplices to the totally ordered subsets of $\mathsf{Ob}\ I$. We find the octahedron, having 6 vertices, 12 edges, 8 faces, and $G$ acts freely on the set of simplices, hence the quotient: 3 vertices, 6 edges, 4 faces, a cellular decomposition of the projective plane. (Its decomposition by three lines.) We present $S^1$ by the ordered set (as at the beginning) % [sic] parties totalement ordonnés -> parties totalement ordonnées, agreement
\[
\begin{tikzcd}[row sep=large, column sep=large]
S_1^- & S_1^+ \\
S_0^- \arrow[u] \arrow[ur] & S_0^+ \arrow[u] \arrow[ul]
\end{tikzcd}
\qquad\qquad \textit{i.e.} \qquad\qquad
\begin{tikzcd}
 & S_1^+ & \\
0 = S_0^- \arrow[ur] \arrow[dr] & & S_0^+ = 0' \arrow[ul] \arrow[dl] \\
 & S_1^- &
\end{tikzcd}
\]
and $G$ acts on it by a reflection exchanging $S_1^-$ and $S_1^+$. So the ordered set $X$ will have $4 \times 13 = 52$ vertices and that is modest \ldots


\endgroup
% END UNIT Der-ChI-234/Der-ChI-234_p056-060.tex
% BEGIN UNIT Der-ChI-234/Der-ChI-234_p061-065.tex
\begingroup\makeatletter%
\NoAutoSpacing%
\newcommand{\printedpage}[1]{\marginpar{\footnotesize p.~#1}}%
\newcommand{\msp}[1]{\marginpar{\footnotesize [#1]}}%
\newcommand{\Hot}{\textsf{Hot}}%
\newcommand{\Dia}{\textsf{Dia}}%
\newcommand{\Ob}{\textsf{Ob}}%
\newcommand{\Hom}{\textsf{Hom}}%
\newcommand{\DD}{\mathbb{D}}%
\makeatother\let\printedpage\derpp\let\msp\dermsp%
\let\mspin\DERmspin%


\printedpage{59}
% head: A. Grothendieck : Dérivateurs, Ch. I, Édition des Universités de Montpellier II et Paris VII
% foot: Version 8 juillet 2001

\msp{page 96}
Finally, it seems that nothing remains of proposition 2 (p.79) introducing the notion of $\DD$-fibration – not even the equivalence of the two conditions $b$) and $b'$) (universal left $\DD$-fibration, universal right $\DD$-fibration). One thinks of their relation to a condition $a$) modified as follows (for $f : X \to Y$) :

$a$) For every $Y'' \to Y' \to Y$, with $Y'$ and $Y''$ \Hot-\underline{aspheric} (and not only $\DD$-aspheric) $X'' = X \times_Y Y'' \to X' = X \times_Y Y'$ is a $\DD$-equivalence. (When one assumes that $Y''$ and $Y'$ have final elements, this is the condition that $f$ be a universal left $\DD$-fibration; when one assumes that $Y''$ and $Y'$ have initial elements, this is the condition universal right $\DD$-fibration.)\footnote[48]{The reasoning must be made explicit. It suffices to test even on $Y'' \to Y'$, where $Y''$ is \underline{one-point}, and where $Y'$ has a final (resp. initial) object.}

We obviously have $a$) $\Rightarrow$ $b$) and $a$) $\Rightarrow$ $b'$) and we would like these to be equivalences. Here is an interesting axiom\footnote[49]{\textit{cf.} \textbf{Der 3} $a'$) on the following page.} on $\DD$, which seems to hold in all known cases, which ensures the desired equivalence. At the same time I am going to reformulate \textbf{Der 3} in three parts (\textit{cf.} p. 67).

\msp{page 97}
\textbf{Der} 3 $a$) If $Z \in \Dia$ is \Hot-aspheric, then $p : Z \to e$ is a $\DD$-equivalence, \textit{i.e.} for example $p^* : \mathcal{A} = \DD(e) \to \DD(Z)$ is fully faithful.

$b$) Let $f : Z \to Z'$ be an arrow in \Dia. If $f$ is a \Hot-equivalence, it is a $\DD$-equivalence (\textit{i.e.} $f^* : \DD(Z') \to \DD(Z)$ induces an \underline{equivalence} between the subcategories of $Z$, $Z'$ \texttt{[rather of} $\DD(Z)$\texttt{,} $\DD(Z')$\texttt{]} consisting of the constant objects).

$a'$) (Makes $a$ considerably more precise). If $Z \in \Dia$ is \Hot-aspheric, then we have $a$), and the essential image of $\DD(e) = \mathcal{A}$ in $\DD(Z)$ consists of the \underline{locally} \underline{constant} objects of $Z$, \textit{i.e.} of the objects such that the functor
\[
Z^{\circ} \longrightarrow \mathcal{A} \qquad z \longmapsto F_z
\]
is locally constant (\textit{i.e.} sends every arrow to an isomorphism) hence isomorphic to a constant functor (since $Z$ is \Hot-aspheric, hence 1-connected).

It must be clearly noted that in $a'$), it is \underline{not} enough to assume that $Z$ is $\DD$-aspheric, instead of \Hot-aspheric. As soon as $\DD$-equivalence is strictly weaker than \Hot, it seems that $a'$) always becomes false, as a condition on $\DD$, if one assumes there only that $Z$ is $\DD$-aspheric. For example, even if $\DD = \DD_i$ (theory of relative homotopy $i$-types), the category
\msp{page 98}
of locally constant objects on $Z$ is expressed in terms of the homotopy $(i+1)$-type, but not in terms of its homotopy $i$-type alone.

It does seem that it is axiom \textbf{Der} 3 ($a'$) which ensures the existence of a good notion (``bilateral'') of $\DD$-fibration, which is stable under base change, and which reduces either
\printedpage{60}
% head: A. Grothendieck : Dérivateurs, Ch. I, Édition des Universités de Montpellier II et Paris VII
% foot: Version 8 juillet 2001
to ``universal left $\DD$-fibration'' or to ``universal right $\DD$-fibration'' and thus which is in addition a \underline{verifiable} condition in practice. In the case where $f$ is proper or smooth, the criterion is particularly simple, since it is limited to the condition that the $f_* f^* a_Y$, or the $f_! f^* a_Y$ (for $a \in \mathcal{A}$), be locally constant objects on $Y$ (hence \underline{constant} if $Y$ is \underline{\underline{\Hot}}-aspheric, by virtue of \textbf{Der} 3 ($a'$)). % [sic] localements -> localement, plural ending on the adverb; % [sic] "(donc ... ($a'$)." -> closing parenthesis after "($a'$)" missing, unbalanced parenthesis
\DERmark{ENV0048}


I am also going to introduce a strengthened variant of \textbf{Der} 3 $b$) :

\textbf{Der} 3 $b'$) Let $Z' \to Z$ be a \Hot-equivalence in \Dia. Then $f^* : \DD(Z) \to \DD(Z')$ induces an equivalence between the categories of \underline{locally} \underline{constant} coefficients on $Z$ and on $Z'$.\footnote[50]{\textbf{Der} 3 $b'_0$) : the functor is faithful\\
\hspace*{6em}$b'_1$) : the functor is fully faithful\\
\hspace*{6em}$b'_2$) : the functor is an equivalence} % [sic] localements -> localement, plural ending on the adverb

\msp{page 99}
In certain questions it suffices to know that the said functor is fully faithful. This implies that
\[
H^*_{\DD}(Z,F) \overset{\sim}{\longrightarrow} H^*_{\DD}(Z', f^*F)
\]
when $F$ is \underline{locally} \underline{constant} on $Z$\footnote[51]{One may say that $f : Z' \to Z$ is a $\DD$-equivalence 0-\textit{complete on the left,} and 0-\textit{complete on the right} in the homological case.} (but it is not clear that this condition always implies full faithfulness\footnote[52]{But this is OK if one has internal \underline{\Hom} in $\DD(Z)_{\mathrm{loc.const.}}$ satisfying
\[
\Hom(F,G) \simeq \Hom(e_Z, \underline{\Hom}(F,G)) \quad ,
\]
where $e$ is the final element of $\mathcal{A}$, the \underline{\Hom} being compatible with inverse images.}) and similarly for homology
\[
H_*^{\DD}(Z', f^*F) \longrightarrow H_*^{\DD}(Z,F) \quad .
\]
Under this condition, we have the :

\textbf{Proposition :} \textit{We assume} \textbf{Der} 3 \textup{b}$'$) \textit{holds. Let $f : X \to Y$ in} \Dia. \textit{For $f$ to be a $\DD$-fibration, it is necessary and sufficient that for $Y'' \to Y'$ over $Y$ which is a} \Hot\textit{-equivalence, $X'' = X \times_Y Y'' \to X' = X \times_Y Y'$ be a $\DD$-equivalence.}

To tell the truth, one can state something reasonable, without assuming axiom \textbf{Der} 3 $b'$). One must introduce the :

\textbf{Definition :} Let $f : Z' \to Z$ in \Dia. We say that $f$ is a 1-\textit{complete} $\DD$-equivalence if the functor $f^* : \DD(Z) \to \DD(Z')$ is fully faithful on the full subcategory of $\DD(Z)$ consisting of the locally constant coefficients. We say that $f$ is a 2-\textit{complete} $\DD$-equivalence if it induces an equivalence between the categories of locally constant coefficients. % [sic] f et une -> f est une, verb "est" misprinted as "et"
\DERmark{ENV0049}


\msp{page 100}
We say that $Z$ is \textit{completely} $\DD$-\textit{aspheric}, if $Z \to e$ is a 2-complete $\DD$-equivalence.

\printedpage{61}
% head: A. Grothendieck : Dérivateurs, Ch. I, Édition des Universités de Montpellier II et Paris VII
% foot: Version 8 juillet 2001

Thus, axiom \textbf{Der} 3 $a'$) can be reformulated, by saying that if $Z$ is \Hot-aspheric, then $Z$ is completely $\DD$-aspheric. And likewise, \textbf{Der} 3 $b'$) by saying that if $f : X \to Y$ in \Dia{} is a \Hot-equivalence, then $f$ is a 2-complete $\DD$-equivalence. We are going to give a weakened variant, of \textbf{Der} 3 $a'$), which suffices to give a reasonable notion of $\DD$-fibration :

\textbf{Der} 3 $a''$) Every $Z$ in \Dia{} which has either a final object, or an initial object, is completely $\DD$-aspheric.

\textbf{Definition.} We say that $f : X \to Y$ in \Dia{} is a $\DD$-fibration, if for
\[
Y'' \longrightarrow Y' \longrightarrow Y
\]
with $Y'' \to Y'$ a 1-complete $\DD$-equivalence (for example $Y'$ and $Y''$ completely $\DD$-aspheric), $X'' = X \times_Y Y'' \longrightarrow X' = X \times_Y Y'$ is a $\DD$-equivalence. % [sic] asphèriques -> asphériques, grave accent instead of acute

This is equivalent to each of conditions $b$), $b'$) of prop. 2 (page 79), \textit{i.e.} that $f$ be universally a left $\DD$-fibration, or universally a right $\DD$-fibration (thanks to \textbf{Der 3} $b''$ \texttt{[rather} $a''$\texttt{]}). % [sic] $\DD$ fibration -> $\DD$-fibration, hyphen missing
\DERmark{ENV0050}


\msp{page 101}
\textbf{Definition 1.} $f : X \to Y$ in \Dia{} is called a \textit{cohomological $\DD$-equivalence} (understood : \textit{complete}) if for every $F \in \Ob\,\DD(Y)$, $F$ \underline{locally} \underline{constant}, we have
\[
H^*_{\DD}(Y,F) \overset{\sim}{\longrightarrow} H^*_{\DD}(X, f^*(F)) \quad .
\]
Dual notion \textit{homological $\DD$-equivalence}
\[
H_*^{\DD}(X, f^*(F)) \overset{\sim}{\longrightarrow} H_*^{\DD}(Y,F) \quad .
\]
If $Y$ is completely $\DD$-aspheric, then the locally constant coefficients on $Y$ are constant, and $f$ is a cohomological (resp. homological) $\DD$-equivalence if and only if it is a $\DD$-equivalence \textit{i.e.} if $X$ is $\DD$-aspheric.\footnote[53]{Axiom \textbf{Der} 3 $b'_1$) implies that every \Hot-equivalence is a cohomological $\DD$-equivalence.}

\textbf{Proposition 1} (Sorital). \textit{Let $X \overset{f}{\to} Y \overset{g}{\to} Z$. If $f$ and $g$ are cohomological $\DD$-equivalences, then $gf$ too. If $gf$ and $f$ are cohomological $\DD$-equivalences, then $g$ too. If $gf$ is a cohomological $\DD$-equivalence, and if $g$ is a \underline{2-complete} $\DD$-equivalence} (hence every $F$ locally constant on $Y$ comes from a $G$ locally constant on $Z$), \textit{then $f$ is a cohomological $\DD$-equivalence.}

\textbf{Definition 2.} 1°) $f : X \to Y$ in \Dia{} is called a \textit{complete cohomological $\DD$-fibration,} I do mean \underline{complete}, if

$a$) For every $y \in \Ob\, Y$, $X_y \to X/y$ is a cohomological $\DD$-equivalence (complete).

$b$) For every $y' \to y$ in $Y$, $X/y' \to X/y$ is a cohomological $\DD$-equivalence (complete).

\printedpage{62}
% head: A. Grothendieck : Dérivateurs, Ch. I, Édition des Universités de Montpellier II et Paris VII
% foot: Version 8 juillet 2001

\msp{page 102}
$c$) These conditions remain satisfied after every base change $Y' \to Y$.

2°) Dual notion of complete homological fibration.

\textsc{Example.} If $f$ is a \Hot-fibration and is \Hot-proper, the arrows of $a$) and $b$) are \Hot-equivalences, hence by virtue of \textbf{Der 3} $b'_1$) (which must be admitted), they are cohomological $\DD$-equivalences. This means that $a$) and $b$) hold, and $c$) too, for the notion of a \Hot-proper arrow is stable under base change.
\DERmark{ENV0051}


\textbf{Proposition 2.} \textit{Let $f : X \to Y$ be a complete cohomological $\DD$-fibration, $F \in \Ob\,\DD(X)$ locally constant on $X$. Then $f_*(F)$ is locally constant, its formation commutes with every base change, and in particular is computed fiber by fiber.}\footnote[54]{There is a converse \textit{cf.} prop. 4.}
\DERmark{ENV0052}


\textbf{Corollary 1.} \textit{Let there be a cartesian diagram}
\[
\begin{tikzcd}
X & X' \arrow[l, "q"'] \arrow[d, "f'"] \\
Y & Y' \arrow[l, "p"']
\arrow[from=1-1, to=2-1, "f"']
\end{tikzcd}
\quad .
\]
\textit{We assume $f$ a cohomological $\DD$-fibration and $p$ a cohomological $\DD$-equivalence. Then $q$ is a cohomological $\DD$-equivalence.}

\msp{page 103}
Indeed, we have for $F$ locally constant on $X$
\[
H^*_{\DD}(X,F) \simeq H^*_{\DD}(Y, f_*F) \overset{\sim}{\longrightarrow} H^*_{\DD}(Y', p^*f_*F) \simeq H^*_{\DD}(Y', f'_*q^*F) \simeq H^*_{\DD}(X', q^*F)
\]
qed.

\textbf{Proposition 3.} \textit{The notion of complete cohomological $\DD$-fibration is stable under base change} (trivial) \textit{and under composition.}

For stability under composition, it suffices to note the following corollary of proposition 2.

\textbf{Corollary 2} (of prop. 2). \textit{Let $f : X \to Y$. For $f$ to be a complete cohomological $\DD$-fibration, it is necessary and sufficient that for every $Y'' \overset{p}{\to} Y' \to Y$ over $Y$, with $p$ a cohomological $\DD$-equivalence, the corresponding arrow $X'' = X \times_Y Y'' \longrightarrow X' = X \times_Y Y'$ also be a cohomological $\DD$-equivalence.}

This is necessary by corollary 1, taking into account that $f' : X' = X \times_Y Y' \to Y'$ is also a cohomological $\DD$-fibration. It is sufficient, for it suffices to apply it to the cases $\{y\} \to Y/\{y\}$, $Y/y \to Y/y'$, which are cohomological $\DD$-equivalences by
\msp{page 104}
virtue of \textbf{Der} 3 $b'_1$) (or only \textbf{Der} 3 $a'$)).

\printedpage{63}
% head: A. Grothendieck : Dérivateurs, Ch. I, Édition des Universités de Montpellier II et Paris VII
% foot: Version 8 juillet 2001

\textbf{Proposition 4.} \textit{We assume that} \textbf{Der} 3 a$'$) \textit{is satisfied} (\textit{i.e.} every $Z$ in \Dia{} which is \Hot-aspheric, is completely $\DD$-aspheric). \textit{Let $f : X \to Y$ in} \Dia. \textit{The following conditions are equivalent}

1°) \textit{$f$ is a \underline{complete} cohomological $\DD$-fibration} (conditions $a$), $b$), $c$) of \texttt{[the]} definition 2, page 101).

2°) \textit{For every $F \in$} \Ob\,$\DD(X)$\textit{, locally constant, $f_*(F)$ is locally constant, is computed fiber by fiber; and that remains true after every base change $Y' \to Y$ (with $F'$ locally constant on $X' = X \times_Y Y'$). (Then the formation of $f_*(F)$ commutes with \underline{every} base change).}

3°) \textit{For every $Y'' \overset{p}{\to} Y'$ over $Y$, if $p$ is a cohomological $\DD$-equivalence, then $X'' = X \times_Y Y'' \longrightarrow X' = X \times_Y Y'$ too.}

3$'$) \textit{As in 3, assuming that $Y''$ is the one-point category, and that $Y'$ has a final element : $X'' \to X'$ must be a cohomological $\DD$-equivalence.}

\msp{page 105}
\textbf{Definition 3.} We say that $f : X \to Y$ in \Dia{} is a \textit{complete $\DD$-fibration} if it is a complete cohomological $\DD$-fibration and a complete homological $\DD$-fibration.

This is stable under base change, under composition.

Definition 3 is equivalent to this : for every $p : Y'' \to Y'$ over $Y$, with $Y''$ and $Y'$ \Hot-aspheric, the corresponding morphism $X'' \to X'$ is \underline{at} \underline{the same} \underline{time} a cohomological $\DD$-equivalence, and a homological $\DD$-equivalence. (And it even suffices to assume it for $Y'' = e$ and for $Y'$ with a final object or an initial object.) One can think of strengthening this condition still further by requiring (without looking to the left or to the right) that $X'' \to X'$ be a 1-complete $\DD$-equivalence. But it is not clear that this strengthened notion is stable under composition – unless one strengthens it still further, by requiring that if $Y'' \to Y'$ is a 1-complete $\DD$-equivalence, then so is $X'' \to X'$. The trouble with this strengthened definition is that for lack of a cohomological or homological description of the notion of
\msp{page 106}
1-complete $\DD$-equivalence, I do not feel capable of giving interesting (not totally trivial) examples of $\DD$-fibrations in this very strong sense. Except for the trivial \Hot-fibrations (\textit{i.e.} with \Hot-aspheric fibers) – but that is no fun!

But when one has associative and unital products $*$ and corresponding internal \underline{\Hom} in the $\DD(X)_{\mathrm{loc.\ const.}}$, compatible with \texttt{[the]} inverse images, so one will have on $X$ if $F$ and $G$ are locally constant in $\DD(X)$, and if $\varepsilon_X$ denotes the unit for $*$ in $\DD(X)$,
\[
\Hom(F,G) \simeq \Hom(\varepsilon_X * F, G) \simeq \Hom(\varepsilon_X, \underline{\Hom}(F,G)) \simeq \Hom_{\mathcal{A}}(\varepsilon_e, p_*\, \underline{\Hom}(F,G)) \ ,
\]
for $\varepsilon_X = p^*(\varepsilon_e)$, where $p : X \to e$, so in this case an arrow $f : X \to Y$ in \Dia{} is a 1-complete $\DD$-equivalence if (and only if) it is a cohomological $\DD$-equivalence. So in this case, if $f$ is a complete cohomological $\DD$-fibration, it is also a complete homological $\DD$-fibration, hence also a complete $\DD$-fibration (definition 3). (And even in the so-called strengthened sense of just now : if $Y'' \to Y'$ over $Y$ is a 1-complete $\DD$-equivalence ($\Leftrightarrow$ a cohomological $\DD$-equivalence!) then $X'' \to X'$ too). % [sic] soit-disant -> soi-disant, spelling; % [sic] "(Et même ... aussi." -> closing parenthesis missing, unbalanced parenthesis
\DERmark{ENV0053}



\endgroup
% END UNIT Der-ChI-234/Der-ChI-234_p061-065.tex
% BEGIN UNIT Der-ChI-234/Der-ChI-234_p066-070.tex
\begingroup\makeatletter%
\newcommand{\printedpage}[1]{\marginpar{\footnotesize p.~#1}}%
\newcommand{\msp}[1]{\marginpar{\footnotesize [#1]}}%
\newcommand{\HOT}{\textsf{HOT}}%
\newcommand{\Hot}{\textsf{Hot}}%
\newcommand{\Dia}{\textsf{Dia}}%
\newcommand{\Cat}{\mathcal{C}at}%
\newcommand{\DD}{\mathbb{D}}%
\makeatother\let\printedpage\derpp\let\msp\dermsp%
\let\mspin\DERmspin%


\printedpage{64}
% head: A. Grothendieck : Dérivateurs, Ch. I, Édition des Universités de Montpellier II et Paris VII
% foot: Version 8 juillet 2001

\msp{page 107}

The question arises of finding $f$ which are cohomological $\DD$-fibrations, or homological ones. I do not know whether I shall be able to find any, which are not already cohomological \HOT-fibrations, or homological ones (this now amounts to the same thing\footnote[55]{Cohomological \HOT-equivalences $=$ homological \HOT-equivalences $=$ \HOT-equivalences plain and simple.} ), for in practice I hardly know how to say that an arrow $Z' \to Z$ in $\Cat$ is a complete cohomological $\DD$-equivalence, except when this arrow is a \HOT-equivalence (and assuming axiom \textbf{Der} 3 $b'_1$), or its weakened cohomological variant).

I shall therefore look for the \HOT-fibrations, \textit{i.e.} the $f : X \to Y$ such that for every $Y'' \to Y' \to Y$, with $Y''$ and $Y'$ aspheric (understood: for \HOT), $X'' \to X'$ is a \Hot-equivalence. (It even suffices to suppose $Y''$ reduced to a point, $Y'$ with a final object $\ldots$ )

Suppose for example $f$ proper (\textit{i.e.} \HOT-proper), so condition $a$) of definition 2, p. 101 is satisfied and likewise after every base change. Does it suffice to check in addition condition $b$), \textit{i.e.} that
\[ X/y \longrightarrow X/y' \quad , \qquad \text{for } y \to y' \text{ in } Y \ , \]
are \Hot-equivalences? (I mean staying over $Y$, without also supposing we have the same thing after base change).
\msp{page 108}
But does this remain true after every base change? % [sic] reste-il -> reste-t-il, forme interrogative

The situation seems better here for smooth $f$ (instead of proper), interpreting conditions $a$) and $b$) in homological form, as meaning that $f_!(e_X)$ is a \underline{locally} \underline{constant} object of $\mathsf{HOT}(Y)$, which is computed fiber by fiber. Since smooth $f$ implies that $f_!$ commutes with every base change, this hypothesis on $f$ (that $f$ be a cohomological \HOT-fibration, perhaps not complete) is stable under base change, hence implies that it is a complete cohomological \HOT-fibration. But this argument uses the general properties of derivators for \HOT, which I have not yet established - especially not in the whole domain $\Cat$! So it is healthy for me to be wary. I prefer to feel out the difficulty while staying at first within the framework of the \underline{cohomological} formalism (noncommutative), of which I have some mastery. I therefore suppose $f$ proper again, and that the % [sic] meilleur -> meilleure, accord avec « situation »
\[ X_y \longrightarrow X/y \quad , \qquad X/y \longrightarrow X/y' \]
\msp{page 109}
are \HOT-equivalences, I want to prove the same thing after every base change, at least that in the situation
\[
\begin{tikzcd}
X \arrow[r, leftarrow, "q"] \arrow[d, "f"'] & X' \arrow[r, leftarrow, "q'"] \arrow[d, "f'"'] & X'' \arrow[d, "f''"] \\
Y \arrow[r, leftarrow, "p"'] & Y' \arrow[r, leftarrow, "p'"'] & Y''
\end{tikzcd}
\]
\printedpage{65}
% head: A. Grothendieck : Dérivateurs, Ch. I, Édition des Universités de Montpellier II et Paris VII
% foot: Version 8 juillet 2001
if $Y'$, $Y''$ are aspheric, then $q'$ is a \Hot-equivalence. If I want to establish this by a cohomological criterion, I must start by taking on $X'$ a locally constant sheaf of \textit{sets} $F$, and prove that $\Gamma(X',F) \overset{\sim}{\to} \Gamma(X'',{q'}^{*}F)$. Since this sheaf does not come from $X$, I am at a loss how to apply the hypothesis on $f$, which concerns a priori constant sheaves of sets on $X$, or on the $X/y$. But we should have this:

\textbf{Lemma.} \textit{Let $f : X'' \to X'$. Then $f$ is a} \HOT-\textit{equivalence} i.e.
\[ H^{*}_{\mathsf{HOT}}(X', a_{X'}) \xrightarrow{\ \sim\ } H^{*}_{\mathsf{HOT}}(X'', a_{X''}) \]
\textit{is} \texttt{[an]} \textit{isomorphism for every $a \in \mathsf{Hot} = \mathsf{HOT}(e)$, if and only if $f$ is a} \Hot-\textit{equivalence,}\footnote[56]{Note here \texttt{[illegible word]} the distinction between \HOT-equivalence and \Hot-equivalence.} \textit{and then $f$ is even a} $2$-\textit{complete} \HOT-\textit{equivalence} (I need only: complete cohomological and homological \HOT-equivalence) i.e. \textit{for every $\xi \in \DD(X')$ \underline{locally} \underline{constant} on $X'$, we have}\footnote[57]{Is this also true for the derivators $\DD_i$?}
\[ H^{*}_{\mathsf{HOT}}(X',\xi) \xrightarrow{\ \sim\ } H^{*}_{\mathsf{HOT}}(X'',{f}^{*}(\xi)) \quad . \]
\DERmark{ENV0054}

\msp{page 110}

The first assertion should follow from the formulas
\[ H^{*}_{\mathsf{HOT}}(Z,a_Z) \simeq \underline{\mathsf{Hom}}_{\mathsf{Hot}}(\mathsf{hot}(Z),a) \tag*{$(*)$} \]
\[ \text{(and also} \qquad\qquad H_{*}^{\mathsf{HOT}}(Z,a_Z) \simeq \mathsf{hot}(Z) \times a \qquad\qquad ) \]
hence if
\[ \mathsf{hot}(X'') \xrightarrow{\ \sim\ } \mathsf{hot}(X') \tag*{$(**)$} \]
is \texttt{[an]} isomorphism, then
\[ H^{*}_{\mathsf{HOT}}(X', a_{X'}) \longrightarrow H^{*}_{\mathsf{HOT}}(X'', a_{X''}) \]
is \texttt{[an]} isomorphism, since it is deduced from $(*)$ by taking $\underline{\mathsf{Hom}}_{\mathsf{Hot}}(\,.\,,a)$. Conversely, if $(*)$ is always \texttt{[an]} isomorphism, expressing the two sides \texttt{[illegible word]} of $\underline{\mathsf{Hom}}_{\mathsf{Hot}}(\,.\,,a)$ and taking the $\pi_0$, we find that
\[ \mathsf{Hom}_{\mathsf{Hot}}(\mathsf{hot}(X'),a) \xrightarrow{\ \sim\ } \mathsf{Hom}_{\mathsf{Hot}}(\mathsf{hot}(X''),a) \]
is an isomorphism for every $a$, hence $(**)$ is an isomorphism.

But the fact that this implies that $f$ is a complete cohomological \HOT-equivalence, \textit{i.e.} that the relation $(*)$ remains true for a $\xi \in \DD(X')$ \underline{locally} \underline{constant} on $X'$, seems to me a delicate property of the derivator \HOT, to which I shall have to return.
\DERmark{ENV0055}


So the lemma is far from established, for want of having so much as constructed the derivator \HOT, without even speaking of establishing its delicate special properties! But I have no reason to doubt the lemma, it is part of the litany of basic intuitions that I would like to put on a solid footing.

\printedpage{66}
% head: A. Grothendieck : Dérivateurs, Ch. I, Édition des Universités de Montpellier II et Paris VII
% foot: Version 8 juillet 2001

\msp{page 111}

The lemma implies that a $\DD = \mathsf{HOT}$-fibration (in the sense of definition 2 page 81, revised and corrected on page 96-97) is in fact already a complete cohomological and homological \HOT-fibration. Now we have seen that if $f$ is \underline{proper}, $f$ is a $\DD$-fibration if and only if the preceding conditions \underline{over $Y$} (without base change) are satisfied. And dually for smooth $f$.

\begin{flushright}
11.11.90
\end{flushright}

I feel the need to review the notions of $\DD$-fibrations of all kinds (p. 76 ff). First the various useful variants of axiom \textbf{Der} 3 on derivators:

\textbf{Der} 3 $a$)\footnote[58]{\textit{cf.} also the series \textbf{Der 3} $c \ldots$ , page 124 ff.} If $Z \in \mathsf{Dia}$ is \Hot-aspheric, it is $\DD$-aspheric \textit{i.e.} $\mathcal{A} = \DD(e) \to \DD(Z)$ is fully faithful. % Doubt resolved: attached PDF leaf 30 (printed p.66), footnote 58 has bold Der 3.

$a'$) As $a$), with in addition the characterization of the essential image of $\mathcal{A}$ in $\DD(Z)$, as consisting of the locally constant $\DD$-coefficients. % [sic] localements -> localement, adverbe invariable

\textbf{Der} 3 $b$) If $f : Z' \to Z$ in \Dia{} is a \Hot-equivalence, then it is a $\DD$-equivalence, \textit{i.e.} the functor $f^{*} : \DD(Z) \to \DD(Z')$ induces an equivalence between the subcategories consisting of the constant coefficients (or again a fully faithful functor $\DD(Z)_{\text{const.}} \longrightarrow \DD(Z')$).
\DERmark{ENV0056}


$b'_1$) If $f$ is as above, then $f$ is a 1-complete $\DD$-equivalence, \textit{i.e.} the functor $f^{*} : \DD(Z) \to \DD(Z')$ induces a fully faithful functor
\msp{page 112}
between the full subcategories consisting of the locally constant objects.

$b'_2$) As $b'_1$), but with 2-\underline{complete} $\DD$-equivalence instead of 1-complete, \textit{i.e.} $f^{*}$ induces an equivalence between the subcategories consisting of the locally constant coefficients.

$b''_{\mathrm{coh}}$) If $f : Z' \to Z$ in \Dia{} is a \Hot-equivalence, then it is a complete cohomological $\DD$-equivalence, \textit{i.e.} for every $F \in \mathsf{Ob}\ \DD(Z)$ which is \underline{locally} \underline{constant}, we have
\[ H^{*}_{\DD}(Z,F) \xrightarrow{\ \sim\ } H^{*}_{\DD}(Z',f^{*}F) \quad . \]

$b''_{\mathrm{hom}}$) If $f : Z' \to Z$ in \Dia{} is a \Hot-equivalence, it is a complete homological $\DD$-equivalence, \textit{i.e.} for every $F \in \mathsf{Ob}\ \DD(Z)$ which is \underline{locally} \underline{constant}, we have\footnote[59]{\textit{cf.} special $\pm$ variant for \HOT{} pages 124-125.}
\[ H_{*}^{\DD}(Z',f^{*}F) \xrightarrow{\ \sim\ } H_{*}^{\DD}(Z,F) \quad . \]

\printedpage{67}
% head: A. Grothendieck : Dérivateurs, Ch. I, Édition des Universités de Montpellier II et Paris VII
% foot: Version 8 juillet 2001

We have the diagram of tautological implications for the variants of \textbf{Der} 3:
\[
\begin{tikzcd}
 & b'_2 \arrow[d, Rightarrow] \arrow[drr, Rightarrow] & & & \\
 & b'_1 \arrow[dl, Rightarrow] \arrow[dr, Rightarrow] \arrow[dd, Rightarrow] & & a' \arrow[ddd, Rightarrow] & \\
b''_{\mathrm{coh}} \arrow[dr, Rightarrow] & & b''_{\mathrm{hom}} \arrow[dl, Rightarrow] & & \\
 & b \arrow[drr, Rightarrow] & & & \\
 & & & a & .
\end{tikzcd}
\]
In all cases known to me, derivators satisfy the strongest condition \textbf{Der} 3 $b'_2$.
\msp{page 113}

I also recall the:

\textbf{Definition 1:} Let $f : X \to Y$ be in \Dia. We say that $f$ is a \textit{cohomological $\DD$-equivalence} (understood: ``\textit{complete}'') if for every $\DD$-coefficient $F$ \underline{locally} \underline{constant} on $Y$,
\[ H^{*}_{\DD}(Y,F) \xrightarrow{\ \sim\ } H^{*}_{\DD}(X,f^{*}F) \quad ; \]
we say that $f$ is a \textit{homological $\DD$-equivalence} (``\textit{complete}'') if for every $F$ as \texttt{[here-]}above,
\[ H_{*}^{\DD}(X,f^{*}F) \xrightarrow{\ \sim\ } H_{*}^{\DD}(Y,F) \quad . \]
We say that $f$ is a 1-\textit{complete} $\DD$-\textit{equivalence} (resp. 2-\textit{complete}) if the functor
\[ f^{*} : \DD(Y)_{\text{loc. const.}} \longrightarrow \DD(X)_{\text{loc. const.}} \]
is fully faithful (resp. an equivalence).
\[
\begin{tikzcd}[column sep=small]
 & \text{2-complete $\DD$-equivalence} \arrow[d, Rightarrow] & & \\
 & \text{1-complete $\DD$-equivalence} \arrow[dl, Rightarrow] \arrow[dr, Rightarrow] & & \\
\text{cohomological $\DD$-equivalence} \arrow[dr, Rightarrow] & & \text{homological $\DD$-equivalence} \arrow[dl, Rightarrow] & \\
 & \text{$\DD$-equivalence} & & .
\end{tikzcd}
\]
If in the subcategories $\DD(X)_{\text{loc. const.}}$ there is an operation $*$ (associative, unital, commutative) and associated internal \underline{\textsf{Hom}}, preserved by inverse images, then cohomological $\DD$-equivalences are already 1-complete and are therefore homological $\DD$-equivalences, therefore in this particular case
\printedpage{68}
% head: A. Grothendieck : Dérivateurs, Ch. I, Édition des Universités de Montpellier II et Paris VII
% foot: Version 8 juillet 2001
\msp{page 114 blank}
\msp{page 115}
\[
\begin{tikzcd}
\text{2-complete $\DD$-equivalence} \arrow[d, Rightarrow] & & \\
\text{1-complete $\DD$-equivalence} \arrow[r, Leftrightarrow] & \text{cohomological $\DD$-equivalence} \arrow[d, Rightarrow] & \\
 & \text{homological $\DD$-equivalence} \arrow[d, Rightarrow] & \\
 & \text{$\DD$-equivalence} & .
\end{tikzcd}
\]
In $\DD = \mathsf{HOT}$, all these notions should be equivalent, and equivalent to the apparently still weaker condition, that $f : X \to Y$ be a \Hot-equivalence. But I know only \HOT{} as a derivator where the $\DD$-equivalence of an arrow $f : X \to Y$ in \Dia{} already implies cohomological $\DD$-equivalence (complete), or homological $\DD$-equivalence (complete), or even 1-complete $\DD$-equivalence. (I shall come back to this presently.)

\textbf{Proposition 1.} \textit{Let $f : X \to Y$ be in} \Dia. \textit{Equivalent conditions:}

a) \textit{For every $y \in \mathsf{Ob}\, Y$, the inclusion $X_y \to X/y$ is a $\DD$-equivalence, and for every arrow $y' \to y$ in $Y$, $X/y' \to X/y$ is a $\DD$-equivalence.}

b) \textit{For every \underline{constant} coefficient $F$ on $Y$, $f_{*}f^{*}F$ is locally constant, and is computed ``fiber by fiber''.}

(This is prop. 1, p. 76).

\textbf{Proposition 1$'$.} Dual.

\msp{page 116}

\textbf{Definition 2.} Under the equivalent conditions of proposition 1, we say that $f$ is a \textit{weak cohomological $\DD$-fibration.} Under the dual conditions, that $f$ is a \textit{weak homological $\DD$-fibration.}

This is what in definition 1, p. 77 I called ``cohomological $\DD$-fibration'' or ``left $\DD$-fibration'', and dually. Now the notion seems to me too weak to deserve this name, which I reserve for the ``complete'' case.

\textbf{Proposition 2.} \textit{Weak cohomological $\DD$-fibrations are stable under smooth base changes, weak homological $\DD$-fibrations stable under proper base changes.}

I should construct an example where $f$ is a weak cohomological $\DD$-fibration, and is no longer so after base change.


\endgroup
% END UNIT Der-ChI-234/Der-ChI-234_p066-070.tex
% BEGIN UNIT Der-ChI-234/Der-ChI-234_p071-075.tex
\begingroup\makeatletter%
\newcommand{\printedpage}[1]{\marginpar{\footnotesize p.~#1}}%
\newcommand{\msp}[1]{\marginpar{\footnotesize [#1]}}%
\newcommand{\D}{\mathbb{D}}%
\newcommand{\sfup}[1]{\textsf{\upshape #1}}%
\newcommand{\Dia}{\sfup{Dia}}%
\newcommand{\Hot}{\sfup{Hot}}%
\newcommand{\HOT}{\sfup{HOT}}%
\newcommand{\Ob}{\sfup{Ob}}%
\newcommand{\Hom}{\sfup{Hom}}%
\newcommand{\Cat}{$\mathcal{C}$\textit{at}}%
\makeatother\let\printedpage\derpp\let\msp\dermsp%
\let\mspin\DERmspin%


\printedpage{69}
% head: A. Grothendieck : Dérivateurs, Ch. I, Édition des Universités de Montpellier II et Paris VII

\textbf{Proposition 3.} \textit{Let}
\[
\begin{tikzcd}
X \arrow[rr, "f"] \arrow[dr, "p"'] & & X' \arrow[dl, "p'"] \\
& S &
\end{tikzcd}
\]
\textit{be in \Dia, with $p$ and $p'$ weak cohomological $\D$-fibrations} (resp. \textit{weak homological $\D$-fibrations}), \textit{let $S_0 \subset S$ be such that $\pi_0(S_0) \to \pi_0(S)$ is surjective. For $f$ to be a $\D$-equivalence, it suffices that the $f_s : X_s \to X'_s$, for $s \in S_0$, be so.}

\textbf{Proposition 4.} \textit{If $f : X \to Y$ in \Dia{} is $\D$-proper} (resp. $\D$-\textit{smooth}), \textit{then $f$ is a weak cohomological} (resp. \textit{homological}) \textit{$\D$-fibration if and only if for every $F$ constant on $Y$, $f_*f^*F$} (resp. $f_!f^*F$) \textit{is locally constant on $Y$. And if this is so, for every $Y' \to Y$ in \Dia, $f' : X' = X \times_Y Y' \to Y'$ is again a weak cohomological} (resp. \textit{homological}) \textit{$\D$-fibration.}

Immediate from criterion \textit{b}) of proposition 1.

\msp{page 117}\textbf{Proposition 5.} \textit{Let $f : X \to Y$ be in \Dia. Consider the following conditions:}

a) \textit{For every $Y'' \to Y'$ over $Y$ in \Dia, if $Y''$ and $Y'$ are \Hot-aspheric, then}
\[
X'' = X \times_Y Y'' \longrightarrow X' = X \times_Y Y'
\]
\textit{is a $\D$-equivalence.}
\[
\begin{tikzcd}
X & X' \arrow[l] & X'' \arrow[l] \\
Y \arrow[from=u] & Y' \arrow[l] \arrow[from=u] & Y'' \arrow[l, "p"] \arrow[from=u]
\end{tikzcd}
\]

a$'$) \textit{As in} a), \textit{but assuming $p : Y'' \to Y'$ is a \Hot-equivalence (instead of assuming $Y'$, $Y''$ \Hot-aspheric).}

a$_1$) \textit{As in} a), \textit{but assuming (instead of asphericity) that $Y'$ has a final object, and $Y''$ is the unit category.}

a$_2$) \textit{As in} a), \textit{but assuming that $Y'$ has an initial object, and $Y''$ is the unit category.}
\DERmark{ENV0057}


b$_1$) \textit{$f$ is a weak cohomological $\D$-fibration, and remains so after any base change.}

b$_2$) \textit{$f$ is a weak homological $\D$-fibration, and remains so after any base change.}

\printedpage{70}
% head: A. Grothendieck : Dérivateurs, Ch. I, Édition des Universités de Montpellier II et Paris VII
\textit{We then have the diagram of tautological implications}
\[
\begin{tikzcd}[column sep=small]
& & & a' \arrow[d, Rightarrow] & & & & \\
& & & a \arrow[dlll, Rightarrow] \arrow[drrr, Rightarrow] & & & & \\
a_1 & \Longleftrightarrow & b_1 & & b_2 & \Longleftrightarrow & a_2 & .
\end{tikzcd}
\]
\textit{When $\D$ satisfies axiom} \textbf{Der} 3 a$'$ (p. 111), \textit{then all these conditions, except} a$'$), \textit{are equivalent. When $\D$ satisfies axiom} \textbf{Der} 3 b$'_1$ \textit{(or only} \textbf{Der} 3 b$''_{\mathrm{coh}}$ or b$''_{\mathrm{hom}}$\textit{),} \textit{then all the conditions (including} a$'$\textit{) are equivalent.}

To simplify, we shall assume henceforth that axiom \textbf{Der} 3 $a'$) is satisfied.

\textbf{Definition 3.} Under the equivalent conditions $a$, $a_1$, $b_1$, $a_2$, $b_2$ of proposition 5, we say that $f$ is a $\D$-\textit{fibration}. % [sic] proposition 4 -> proposition 5, the conditions a, a_1, b_1, a_2, b_2 are those of proposition 5
\DERmark{ENV0058}


\msp{page 118}\textbf{Proposition 6.} \textit{Suppose that $f : X \to Y$ in \Dia{} is $\D$-proper} (resp. $\D$-\textit{smooth}). \textit{Equivalent conditions:}

a) \textit{$f$ is a $\D$-fibration.}

b) \textit{$f$ is a weak cohomological} (resp. \textit{homological}) \textit{$\D$-fibration.}

c) \textit{For every $y' \to y$ in $Y$, $X/y' \to X/y$} (resp. $y\backslash X \to y'\backslash X$) \textit{is a $\D$-equivalence.}

d) \textit{For every $F$ constant on $Y$, $f_*f^*F$} (resp. $f_!f^*F$) \textit{is locally constant on $Y$.}

d$'$) \textit{For every $F$ locally constant on $Y$, $f_*f^*F$} (resp. $f_!f^*F$) \textit{is locally constant on $Y$.}

This follows formally from Prop. 4 and Prop. 5.

\textbf{Corollary.} \textit{Suppose $f : X \to Y$ is} \Cat\textit{-fibrant} (resp. \Cat\textit{-cofibrant}). \textit{For $f$ to be a $\D$-fibration, it is necessary and sufficient that for every $u : y \to y'$ in $Y$, $u^* : X_{y'} \to X_y$} (resp. $u_* : X_y \to X_{y'}$) \textit{be a $\D$-equivalence.}

It is immediate that the notion of $\D$-fibration is stable under base change, but it is doubtful that it is so under composition.\footnote[60]{N.B. If the $\D$-equivalences are already \Hot-equivalences, then the notion coincides with that of \Hot-fibration, and it is stable under composition. Also, if the $\D$-equivalences are in the sense of cohomological, or homological, $\D$-equivalences, then the notion coincides with that of cohomological (resp. homological) $\D$-fibrations below, and is also stable under composition.} (Here again, I have omitted to make a counterexample, with derivators $\D = \sfup{Der}(k\text{-}\sfup{Mod})$, say.) Another shortcoming of the notion: it is not clear, if $f$ is a $\D$-fibration and $F$ a $\D$-coefficient on $Y$, whether
\[
f_*f^*F \qquad \text{and} \qquad f_!f^*F
\]
\printedpage{71}
% head: A. Grothendieck : Dérivateurs, Ch. I, Édition des Universités de Montpellier II et Paris VII
are computed ``fiber by fiber'', as one would expect in the case of a ``bundle''.

On the other hand, here is a nice result:

\msp{page 119}\textbf{Proposition 7.} \textit{Let $f : X \to Y$ be in \Dia. Equivalent conditions:}

a) \textit{$f$ is a ``universal'' $\D$-equivalence,} i.e. \textit{for every $Y' \to Y$ in \Dia, $f' : X' = X \times_Y Y' \to Y'$ is a $\D$-equivalence.}

b) \textit{$f$ is a $\D$-fibration with $\D$-aspheric fibers.}

\textbf{Definition 4.} Under the equivalent conditions of proposition 7, we say that $f$ is a \textit{trivial} $\D$-fibration.

But beware, $f$ can be a $\D$-fibration and a $\D$-equivalence, without the $\D$-fibration being trivial \textit{i.e.} without the fibers being $\D$-aspheric, cf. p. 87 and following. % [sic] peut-être -> peut être, verb "pouvoir" followed by infinitive
And also (should go above) even if $\D$ satisfies \textbf{Der} 3 $b'_2$ (the strongest of the variants of \textbf{Der} 3), thus when \underline{all} the conditions of proposition 5, including $a'$), are equivalent, these (\textit{i.e.} that $f$ be a $\D$-fibration) nevertheless do not imply, in general, that in condition $a'$) one can restrict oneself to $p : Y'' \to Y'$ a $\D$-equivalence (instead of a \Hot-equivalence). In other words, if we have a cartesian diagram in \Dia
\DERmark{ENV0059}

\[
\begin{tikzcd}
X \arrow[d, "f"'] & X' \arrow[l, "q"'] \arrow[d, "f'"] \\
Y & Y' \arrow[l, "p"] & ,
\end{tikzcd}
\]
with $f$ a $\D$-fibration, and $p$ a $\D$-equivalence, it does not necessarily follow that $q$ is a $\D$-equivalence. (But this will nevertheless be the case, ``of course'', if $\D = \HOT$.)

\msp{page 120}On the other hand, we have:

\textbf{Proposition 8.} \textit{Let $f : X \to Y$ be in \Dia. Equivalent conditions:}

a) \textit{$f$ is a $\D$-fibration.}

b) \textit{For every $p : Y'' \to Y'$ over $Y$, if $p$ is a} 1-\textit{complete $\D$-equivalence} (def. 1, p. 113), \textit{then $q : X'' \to X'$ is a $\D$-equivalence.}
\[
\begin{tikzcd}
X \arrow[d, "f"'] & X' \arrow[l] \arrow[d, "f'"'] & X'' \arrow[l, "q"'] \arrow[d, "f''"] \\
Y & Y' \arrow[l] & Y'' \arrow[l, "p"]
\end{tikzcd}
\]
\textbf{Corollary.} \textit{Suppose that in \Dia, every $\D$-equivalence is} 1-\textit{complete.} (N.B. I know only the derivator \HOT{} to satisfy this condition, which seems almost characteristic of \HOT.) \textit{Then $f : X \to Y$ is a $\D$-fibration if and only if for every $p : Y'' \to Y'$ over $Y$, if $p$ is a $\D$-equivalence, then $q : X'' \to X'$ too. In this exceptional case among derivators, the notion of $\D$-fibration is stable under composition} (which is tautological from the criterion I have just given.)

\printedpage{72}
% head: A. Grothendieck : Dérivateurs, Ch. I, Édition des Universités de Montpellier II et Paris VII

We are going to strengthen the notion of $\D$-fibration, to make it more flexible, and in particular stable under composition.

\textbf{Proposition 9.} \textit{Let $f : X \to Y$ be in \Dia. Equivalent conditions:}

a) \textit{For every $y \in Y$, $X_y \to X/y$ is a \underline{cohomological} $\D$-\underline{equivalence}} (def. 1, p. 113), \textit{and likewise for $X/y \to X/y'$ for every arrow $y \to y'$ in $Y$. And these conditions on $f$ remain valid after any base change} (N.B. it suffices that the first of these conditions remain so).

\msp{page 121}b) \textit{for every $F \in \Ob\ \D(X)$ \underline{locally} \underline{constant}, $f_*F$ is locally constant on $Y$ and is computed ``fiber by fiber'' -- and likewise after any base change.} % Doubt resolved: attached PDF leaf 36 (printed p.72) has an en dash, encoded --.

c) \textit{For every $p : Y'' \to Y'$ over $Y$, if $p$ is a cohomological $\D$-equivalence, then the same holds for $q : X'' \to X'$.}

c$'$) \textit{As in} c, \textit{but assuming as a hypothesis on $p$ that $Y''$ is the final category, and $Y'$ has a final object.}

\textsc{Proof}. In the order ($\pm$ tautological)
\[
\mathrm{a} \Longrightarrow \mathrm{b} \Longrightarrow \mathrm{c} \Longrightarrow \mathrm{c}' \Longrightarrow \mathrm{a} \quad .
\]

\textbf{Definition 5.} Under the equivalent conditions of proposition 9, we say that $f$ is a \textit{cohomological $\D$-fibration}. Under the dual conditions, we say that $f$ is a \textit{homological $\D$-fibration}. We say that $f$ is a \textit{strong $\D$-fibration}, if it is both a cohomological $\D$-fibration and a homological $\D$-fibration.

These notions strengthen the notion of $\D$-fibration. They have the further advantage of being \underline{stable} \underline{under} \underline{composition}, not only under base change.

\msp{page 122}The notion of strong $\D$-fibration was chosen to cover the two dual notions of cohomological and homological $\D$-fibration (for lack of a better notion) while remaining stable under composition. One could require that for every % [sic] meilleur -> meilleure, agreement with "notion"
\[
Y'' \longrightarrow Y'
\]
over $Y$, with $Y'$, $Y''$ \Hot-aspheric, $q : X'' \to X'$ be a 1-complete $\D$-equivalence (or even 2-complete), call these the 1-\textit{complete} (resp. 2-\textit{complete}) $\D$-\textit{fibrations}. This strengthens in a rather elegant way the notion of strong $\D$-fibration, but has the great fault of not being stable under composition (as far as I know!), except in the particular case where we have nice internal \underline{\Hom} (cf. bottom of page 113), in which case the notion of cohomological $\D$-equivalence already implies that of 1-complete $\D$-equivalence,
\DERmark{ENV0060}

\printedpage{73}
% head: A. Grothendieck : Dérivateurs, Ch. I, Édition des Universités de Montpellier II et Paris VII
\msp{page 123}thus also that of homological $\D$-equivalence. So in this case
\[
\begin{tikzcd}
\text{cohomological $\D$-fibration} \arrow[r, Leftrightarrow] \arrow[d, Rightarrow] & \text{1-complete $\D$-fibration} \\
\text{homological $\D$-fibration} \quad , &
\end{tikzcd}
\]
and there is no need to look further than cohomological $\D$-fibrations.

One could think, it is true, of defining (in the general case) a stronger notion, by requiring that for every $p : Y'' \to Y'$ over $Y$ which is a 1-complete $\D$-equivalence, $q : X'' \to X'$ also be a 1-complete $\D$-equivalence. This time, the notion is stable under composition, but it nevertheless seems unusable to me (outside the case we have just mentioned, where it coincides with the notion of cohomological $\D$-fibration, or again of strong $\D$-fibration), for, not being in a position to describe the 1-complete $\D$-equivalences, I am also unable to give nontrivial examples of ``strictly 1-complete'' $\D$-fibrations.

To check that an $f : X \to Y$ in \Dia{} is a cohomological $\D$-fibration, one must
\msp{page 124}know how to check that arrows $q : Z' \to Z$ in \Dia{} are cohomological $\D$-equivalences, hence that $H^*_{\D}(Z,F) \overset{\sim}{\rightarrow} H^*_{\D}(Z',q^*(F))$ is \texttt{[an]} isomorphism for every $F \in \Ob\ \D(Z)$ which is locally constant on $Z$. Now I know how to check a statement of this nature (outside the case where I know in advance that $Z$ is completely $\D$-aspheric, \textit{i.e.} that the locally constant coefficients on $Z$ are in fact constant) \underline{only} if $q$ is already a \Hot-equivalence. % [sic] constant) -> constants), agreement with "coefficients"
That is to say that I know how to exhibit cohomological $\D$-fibrations, only when they are already cohomological \HOT-fibrations. I want, in what follows, to give usable criteria for \HOT-fibrations. But we are going to axiomatize these developments, keeping a general $\D$, on which I am going to make an assumption so draconian that practically, I know only the derivator \HOT{} which satisfies it. This is $c_1$ below.

\textbf{Der} 3 $c_1$) Every $\D$-equivalence is a cohomological $\D$-equivalence.

$c_2$) Every $\D$-equivalence is a homological $\D$-equivalence.

\msp{page 125}$c'$) Every $\D$-equivalence is 1-complete.

$c''$) Every $\D$-equivalence is 2-complete.

N.B. We have
\[
\begin{tikzcd}
& c'' \arrow[d, Rightarrow] & \\
& c' \arrow[dl, Rightarrow] \arrow[dr, Rightarrow] & \\
c_1 & & c_2
\end{tikzcd}
\]


\endgroup
% END UNIT Der-ChI-234/Der-ChI-234_p071-075.tex
% BEGIN UNIT Der-ChI-234/Der-ChI-234_p076-076.tex
\begingroup\makeatletter%
\newcommand{\printedpage}[1]{\marginpar{\footnotesize p.~#1}}%
\newcommand{\msp}[1]{\marginpar{\footnotesize [#1]}}%
\newcommand{\HOT}{\textsf{HOT}}%
\newcommand{\Hot}{\textsf{Hot}}%
\newcommand{\Dia}{\textsf{Dia}}%
\newcommand{\Cat}{\mathcal{C}at}%
\newcommand{\D}{\mathbb{D}}%
\makeatother\let\printedpage\derpp\let\msp\dermsp%
\let\mspin\DERmspin%


\printedpage{74}
% head: A. Grothendieck : Dérivateurs, Ch. I, Édition des Universités de Montpellier II et Paris VII
% foot: Version 8 juillet 2001
tautologically. I know only \HOT{} as satisfying the two weakest conditions $c_1$, $c_2$ together. And in fact, subject to writing the theory of the derivator \HOT, it must be true that it satisfies even the strongest form $c''$. (And the \HOT-equivalences must be simply the \Hot-equivalences.) (Cf. the ``lemma'' p.~109.)

The following proposition is a tautology.

\noindent\textbf{Proposition 10.} \textit{Suppose that $\D$ satisfies condition \textup{\textbf{Der} 3 $\mathrm{c}_1$ (p.~124),} for example $\D = \text{\HOT}$. Then the notion of $\D$-fibration coincides with that of cohomological $\D$-fibration. It is therefore stable under composition. For $f : X \to Y$ to be a $\D$-fibration, it is necessary and sufficient that for $p : Y'' \to Y'$ over $Y$, if $p$ is a $\D$-equivalence, the same be true of $q : X'' \to X'$.}

\noindent\textbf{Corollary 1.} \textit{If $f : X \to Y$ in \textup{\Dia} is $\D$-proper, then $f$ is a cohomological $\D$-fibration if and only if for every arrow $y' \to y$ in $Y$, $X/y' \to X/y$ is a $\D$-equivalence \textup{(i.e.} if and only if $f$ is a weak cohomological $\D$-fibration\textup{).}}

Apply criterion $c$) of \texttt{[the]} prop.~6.

\msp{page 126}

\noindent\textbf{Corollary 2.} \textit{Suppose that $\D$ satisfies \textup{\textbf{Der} 3 $(\mathrm{c}_1, \mathrm{c}_2)$ (for example $\D = \text{\HOT}$).} Then for an arrow $f : X \to Y$ in \textup{\Dia}, it amounts to the same thing for $f$ to be a $\D$-fibration, a cohomological $\D$-fibration, or a homological $\D$-fibration. If $f$ is proper \textup{(resp.} smooth\textup{),} it amounts to the same thing to say that $f$ is a cohomological $\D$-fibration \textup{(resp.} homological\textup{),} weak.}

We therefore see that in $\Cat$, we must have a single reasonable notion of \HOT-fibration, which will coincide with that of cohomological or homological \HOT-fibration. But for this notion to be available as of now, the derivator \HOT{} would have to have already been constructed on the domain $\text{\Dia} = \Cat$ in its entirety. Now this has not yet been done (and there could therefore have been no question of proving the properties of \HOT{} just mentioned, say the strongest \textbf{Der} 3 $c''$).

Nevertheless, I would like to develop in a direct (``elementary'') way the notion of \Hot-fibration, without referring to the theory of the derivator \HOT, and to the preceding theoretical developments. % [sic] réferer -> référer, missing accent on the second e


\endgroup
% END UNIT Der-ChI-234/Der-ChI-234_p076-076.tex
\clearpage\setcounter{footnote}{0}\phantomsection\addcontentsline{toc}{section}{Chapter II}\markboth{Chapter II}{Chapter II}
% BEGIN UNIT Der-ChII-234/Der-ChII-234_p001-005.tex
\begingroup\makeatletter%
\newcommand{\printedpage}[1]{\marginpar{\footnotesize p.~#1}}%
\newcommand{\msp}[1]{\marginpar{\footnotesize [#1]}}%
\newcommand{\Vcat}[2]{\underset{#2}{\underset{\vee}{#1}}}%
\makeatother\let\printedpage\derpp\let\msp\dermsp%
\let\mspin\DERmspin%


% ===================== PDF page 1 =====================
\printedpage{--}
% no page number printed (title page)
% head: (none)

\begin{center}
\vspace*{4cm}
{\Large\textsf{THE DERIVATORS}}

\bigskip\bigskip

\textbf{ALEXANDRE GROTHENDIECK}

\bigskip\bigskip\bigskip

\textbf{Chapter II}

\textbf{Cofinality (on the right and on the left)}

\textbf{(preliminary to cohomological cofinality)}
\end{center}

\par

% ===================== PDF page 2 =====================
\printedpage{--}
% no page number printed (editorial note)
% head: (none)

\vspace*{6cm}

\noindent \par\noindent\textbf{Note by the editors of the transcribed edition.}\par
This text was deciphered and transcribed in \LaTeXe{} by M. Künzer. It was edited by M. Künzer, J. Malgoire, and G. Maltsiniotis. The transcription is as faithful as possible to the manuscript. For the few obvious corrections, or rare comments by the editors, as well as for the original numbering of the manuscript pages, typewriter type \texttt{[typewriter]} in brackets is used. A question mark in brackets means that one is not sure of the preceding word.

\noindent This edition is provisional. Remarks, comments and corrections are welcome. Send a message to:

\noindent\textbf{[e-mail address removed]}

\hfill G. Maltsiniotis\hspace*{2cm}

\par

% ===================== PDF page 3 =====================
\printedpage{1}
% head: A. Grothendieck : Dérivateurs, Ch. II, Édition des Universités de Montpellier II et Paris VII
% foot: Version 8 juillet 2001

\msp{page 1}

\section*{1\quad Cohomological cofinality}

Let
\[ \alpha : C' \longrightarrow C \]
be a functor of (small) categories, $\mathcal{M}$ a category having the relevant $\varprojlim$, $F$ a presheaf (thus contravariant in $i$) on $C$ with values in $\mathcal{M}$. We have a canonical homomorphism in $\mathcal{M}$
\[ \varprojlim_{C} F \longrightarrow \varprojlim_{C'} \underbrace{\alpha^*(F)}_{F\circ\alpha} \tag*{(1)} \]
and we want conditions on $\alpha$ for this always to be an isomorphism (\footnote{One then says that $\alpha$ is \textit{cofinal,} cf. SGA 4, I 8.1, or 3.3, p. 45.}). We are at once reduced to the case where $\mathcal{M} = (\mathrm{Ens})$. Let us say that a category $I$ is \textit{weakly left filtering} if $I \neq \emptyset$ and if $\forall\, i', i'' \in \operatorname{Ob} I$, $\exists\, i \in \operatorname{Ob} I$ and arrows
\[
\begin{tikzcd}[row sep=small, column sep=small]
 & i' \\
i \arrow[ur] \arrow[dr] & \\
 & i''.
\end{tikzcd}
\]
This also means that $\operatorname{Ob} I$ as an ordered set is left filtering.

\textbf{Proposition.} \textit{For} (1) \textit{always to be an isomorphism} ($\alpha$ \textit{``anodyne for} $\varprojlim$ \textit{''), it suffices that for every} $x_0 \in \operatorname{Ob} C$\textit{, the category} $x_0\backslash C'$ \textit{of pairs} $(x', u)$\textit{, with} $x' \in \operatorname{Ob} C'$\textit{,} $u : x_0 \longrightarrow \alpha(x')$\textit{, be weakly left filtering} (\footnote{Cf. improvement in the remarks p. 5.}, \footnote{Cf. 3.2, p. 43-44, for necessary and sufficient conditions for cofinality.}). % [sic] u : \alpha(x') \longrightarrow x -> x_0 in place of x, the category is x_0\backslash C'
\DERmark{ENV0061}


This does not imply the stronger conclusion
\[ F \underbrace{\xrightarrow{\;\sim\;}}_{\text{adj.}} \alpha_*\alpha^* F \]
that the canonical adjunction morphism is an isomorphism. For this to be so, it suffices that $\forall\, x_0 \in C$, the functor $\alpha/x_0 : C'/x_0 \longrightarrow C/x_0$ induced by $\alpha$ be ``anodyne for the $\varprojlim$ '', which is ensured when for every arrow $u : x_1 \longrightarrow x_0$ in $C$, the category $\Vcat{C'}{u}$ formed of triples $(x', \xi, \eta)$ with $x' \in \operatorname{Ob} C'$ and $\xi : x_1 \longrightarrow \alpha(x')$, $\eta : \alpha(x') \longrightarrow x_0$ rendering commutative
\DERmark{ENV0062}

% printed: C' inside a V-shaped wedge, with u printed below the wedge
\msp{page 2}
the diagram
\printedpage{2}
% head: A. Grothendieck : Dérivateurs, Ch. II, Édition des Universités de Montpellier II et Paris VII
% foot: Version 8 juillet 2001
\[
\begin{tikzcd}
x_1 \arrow[rr, "u"] \arrow[dr, "\xi"'] & & x_0 \\
 & \alpha(x'), \arrow[ur, "\eta"'] &
\end{tikzcd}
\]
is weakly left filtering.

More generally, let us consider
\[
\begin{tikzcd}
X' = C' \arrow[rr, "\alpha"] \arrow[dr, "f'"'] & & C = X \arrow[dl, "f"] \\
 & S &
\end{tikzcd}
\]
and let us ask ourselves the question when the functorial homomorphism, with respect to a variable presheaf $F$ on $X$,
\[ f_*(F) \longrightarrow f'_*(\alpha^*(F)) \]
is always an isomorphism. (The preceding case is the one where $S = X$, $f = \mathrm{id}_X$.) We find that it suffices that $\forall\, s_0 \in \operatorname{Ob} S$, the functor
\[ \alpha/s_0 : X'/s_0 \longrightarrow X/s_0 \]
be ``anodyne for the $\varprojlim$ '', and for this, it suffices that for \texttt{[every]} object $(x_0, u_0)$ of $X/s_0$ (where
\[ f(x_0) \xrightarrow{u_0} s_0) \]
the category $(x_0, u_0)\backslash(X'/s_0)$ be weakly left filtering.

\msp{page 3}

This category is interpreted as that of triples $(x', \xi, u)$ with
\[
\left\{
\begin{aligned}
& x' \in \operatorname{Ob} C' \\
& u : f'(x') = f(\alpha(x')) \longrightarrow s_0 \\
& \xi : x_0 \longrightarrow \alpha(x')
\end{aligned}
\right.
\]
\DERmark{ENV0063}

\texttt{[rather} $f'(x')$ \texttt{instead of} $g(x')$\texttt{]} so as to render commutative the diagram
\[
\begin{tikzcd}
f(x_0) \arrow[rr, "u_0"] \arrow[dr, "f(\xi)"'] & & s_0 \\
 & f(\alpha(x')). \arrow[ur, "u"'] &
\end{tikzcd}
\]
I would like to bring out variants of these statements in the cohomological setting, the functors $\varprojlim$ being replaced by $R\Gamma$'s (of abelian presheaves e.g. - but one also thinks of the ``noncommutative cases'', defined by derivators in particular.)

\par

% ===================== PDF page 5 =====================
\printedpage{3}
% head: A. Grothendieck : Dérivateurs, Ch. II, Édition des Universités de Montpellier II et Paris VII
% foot: Version 8 juillet 2001

If $I$ is a \textit{finite ordered} category, the \texttt{[weakly]} left filtering condition is equivalent to the existence of an initial object, which implies that $I$ is contractible. In the preceding statements, it might be that the right philosophy is to replace ``weakly left filtering'' everywhere by ``contractible''. If e.g. $C$ is the final category, then replace the condition that $C'$ be left filtering, by acyclicity. This is in any case in accord with the spirit of the axiom $\mathrm{Der}_5$ (case where $J$ is the final category). Indeed, in the case where $C = e$, the condition that $\alpha$ be anodyne for the $\varprojlim$ is visibly equivalent to saying that the functor $\alpha^* : \mathcal{M} \longrightarrow \underline{\mathrm{Hom}}(C'^{o}, \mathcal{M})$, associating to every $F \in \mathcal{M}$ the constant presheaf with value $F$ % [sic] précedents -> précédents, missing accent on the second e
\msp{page 4}
on $C'$, gives an adjunction morphism
\[ \varprojlim_{e} F = F \longrightarrow \alpha_*(\alpha^*(F)) = \varprojlim_{C'} \alpha^* F \]
which is an isomorphism. And the axiom $\mathrm{Der}_5$ ensures that it must be so if $C'$ is contractible.

The filtering condition in the proposition is moreover a \textit{sufficient} condition for anodyneness, but by no means necessary. If e.g. $C$ is the final category, then the condition of anodyneness for $C' \longrightarrow C$ is simply that $C'$ be connected and nonempty, i.e. $\pi_0(C') = \{0\}$. (Condition which visibly ignores the higher $\pi_i$.)

Let \fbox{$\alpha : X' \longrightarrow X$} be a functor of small categories. The condition suggested by the proposition, that the $x_0\backslash X'$ be contractible, means that $\alpha$ is a co-strict hot-equivalence, i.e. that $\alpha^{o} : X'^{o} \longrightarrow X^{o}$ is a hot-equivalence ``locally on $X^{o}$'', hence also globally. If $\alpha$
\msp{page 5}
itself were a strict hot-equivalence, instead of $\alpha^{o}$, then we would have by $\mathrm{Der}_5$, that for $F \in \mathbf{D}(X)$, $F \xrightarrow{\;\sim\;} \alpha_*\alpha^* F$ ($\mathbf{D}$ a left derivator), which is much stronger than $\Gamma_{\mathbf{D}}(X, F) \xrightarrow[\sim]{\alpha^*} \Gamma_{\mathbf{D}}(X', \alpha^* F)$. The dual condition implies (if $\mathbf{D}$ \texttt{[is a]} right derivator) $\alpha_!\alpha^* F \xrightarrow{\;\sim\;} F$. But is it reasonable to expect that it also implies the condition ``on the cohomology side, instead of homology'',
\[ \Gamma_{\mathbf{D}}(X, F) \xrightarrow{\;\sim\;} \Gamma_{\mathbf{D}}(X', \alpha^*(F)) \ \ ? \]

The first case, which should carry conviction, is that of the derivator DerAb, i.e. of the derived categories of abelian sheaves:

\textbf{Question.} \textit{Let} $\alpha : X' \longrightarrow X$ \textit{be a co-strict} hot-\textit{equivalence,} $F$ \textit{an abelian presheaf on} $X$\textit{. Do we then have} $H^i(X, F) \xrightarrow[\sim]{\alpha^*} H^i(X', \alpha^*(F))$ \textit{for every} $i \in \mathbf{Z}$ \textit{?}

If $X$ admits an initial object $\emptyset$, then $\emptyset\backslash X' \simeq X'$, so $X'$ contractible, so it's OK for $F$ constant.

\textbf{Remark.} I realize that in proposition 1, instead of the condition that the $x\backslash C'$ be weakly left filtering, it would suffice that they be nonempty and connected, i.e.


\endgroup
% END UNIT Der-ChII-234/Der-ChII-234_p001-005.tex
% BEGIN UNIT Der-ChII-234/Der-ChII-234_p006-010.tex
\begingroup\makeatletter%
\newcommand{\printedpage}[1]{\marginpar{\footnotesize p.~#1}}%
\newcommand{\msp}[1]{\marginpar{\footnotesize [#1]}}%
\newcommand{\Homint}{\operatorname{Hom\,int}}%
\makeatother\let\printedpage\derpp\let\msp\dermsp%
\let\mspin\DERmspin%


\printedpage{4}
% head: A. Grothendieck : Dérivateurs, Ch. II, Édition des Universités de Montpellier II et Paris VII
% foot: Version 8 juillet 2001
0-connected (\footnote[4]{This also means that $\alpha$ is $W_0$-coaspheric. For a treatment ab ovo, cf. prop. 3.2, page 43.}). The fact that the $x\backslash C'$ nonempty implies that $\varprojlim\limits_{C} \longrightarrow \varprojlim\limits_{C'}$ is \textit{injective.} For surjectivity, if $\xi' \in \varprojlim\limits_{C'}$ (set-theoretic case), we define $\xi \in \varprojlim\limits_{C}$ by the $\xi(x) \in F(x)$ with the help of % [sic] sujectivité -> surjectivité, lettre r omise
\msp{page 6}
\[ x \xrightarrow{u} \alpha(x') \]
by taking $u^*(\xi'(x'))$ (where $\xi'(x') \in (\alpha^*(F))(x') = F(\alpha(x')))$, the essential point is to show that the $u^*(\xi'(x'))$ depend only on $x$, not on the choice of $(x',u) \in x\backslash C'$. Now for $(x',u)$ variable in $x\backslash C'$, the map
\DERmark{ENV0064}\DERmark{ENV0065}\DERmark{ENV0066}

\[
\begin{aligned}
\varphi : \operatorname{Ob}(x\backslash C') &\quad\longrightarrow\quad F(x)\\
(x',u) &\quad\longmapsto\quad u^*(\xi'(x'))
\end{aligned}
\]
\DERmark{ENV0067}

is such that if there exists an arrow
\[ (x',u) \xrightarrow{u'} (\bar{x}',\bar{u}) \qquad \text{where} \qquad u' : x' \longrightarrow \bar{x}' \]
(thus the diagram
\[
\begin{tikzcd}
 & \alpha(x') \arrow[dd, "\alpha(u')"] \\
x \arrow[ur, "u"] \arrow[dr, "\bar{u}"'] & \\
 & \alpha(\bar{x}')
\end{tikzcd}
\]
is commutative), then $\varphi(x',u) = \varphi(\bar{x}',\bar{u})$. Thus $\varphi$ is constant on the connected components of $x\backslash C'$, thus constant if $x\backslash C'$ is 0-connected, whence $\xi(x)$. We then prove at once that for variable $x$,
\[ \xi \overset{\mathrm{def}}{=} (\xi(x))_{x\in\operatorname{Ob} C} \in \varprojlim F, \qquad \text{\texttt{[rather}}\ \ \in \varprojlim F, \text{\texttt{]}} \]
\DERmark{ENV0068}\DERmark{ENV0069}

and that $\alpha^*(\xi) = \xi'$, q.e.d.

We can say that for $x_0 \in \operatorname{Ob} C$ fixed, if we look at the \textit{constant} presheaf on $x_0\backslash C'$ with value $F(x_0)$, say $\Phi_{x_0}$, the map
\DERmark{ENV0070}

\[ (x',u) \longmapsto u^*(\xi'(x')) : x_0\backslash C' \longrightarrow F(x_0) \]
\DERmark{ENV0071}

can be regarded as a section of this presheaf on $x_0\backslash C'$, defined by $\xi' \in \varprojlim\limits_{C'} \alpha^*F$, so we have
\DERmark{ENV0072}

\[ \varprojlim_{C'} \alpha^*F \xrightarrow{c_{x_0}} \Gamma(x_0\backslash C', \Phi_{x_0}) \qquad (\xleftarrow{\ \sim\ } F(x_0) \text{ if } x_0\backslash C '\text{ 0-connected}). \]
\DERmark{ENV0073}

If we have \fbox{$x_0 \xrightarrow{v} x_1$} in $C$, we deduce from this
\DERmark{ENV0074}

\[ x_1\backslash C' \xrightarrow{V} x_0\backslash C' \]
(the $x\backslash C'$ are the fibers of a \textit{fibered category} $C\backslash C'$ over $C$ \dots) and we have $F(x_1) \longrightarrow F(x_0)$, whence \fbox{$\Phi_{x_1} \longrightarrow V^*\Phi_{x_0}$} for the corresponding constant sheaves on $x_1\backslash C'$ resp. $x_0\backslash C'$
\DERmark{ENV0075}

\printedpage{5}
% head: A. Grothendieck : Dérivateurs, Ch. II, Édition des Universités de Montpellier II et Paris VII
% foot: Version 8 juillet 2001
\msp{page 7}
and we have commutativity in
\[
\begin{tikzcd}[column sep=large]
\varprojlim\limits_{C'} \alpha^*F \arrow[r, "c_{x_0}"] \arrow[d, "c_{x_1}"'] & \Gamma(x_0\backslash C', F(x_0)) \simeq \operatorname{Hom}(\pi_0(x_0\backslash C'), F(x_0)) \arrow[d, "V^*"] & \\
\operatorname{Hom}(\pi_0(x_1\backslash C'), F(x_1)) \simeq \Gamma(x_1\backslash C', F(x_1)) \arrow[r, "\substack{v^*\\ (F(x_1) \longrightarrow F(x_0))}"'] & \Gamma(x_1\backslash C', F(x_0)) \simeq \operatorname{Hom}(\pi_0(x_1\backslash C'), F(x_0)) & \\
\pi_0(x_1\backslash C') \arrow[r, "\pi_0(V)"] \arrow[d, dotted] & \pi_0(x_0\backslash C') \arrow[d, dotted] & \text{commutative.} \\
F(x_1) \arrow[r, "F(v)"] & F(x_0) &
\end{tikzcd}
\]
\DERmark{ENV0076}\DERmark{ENV0077}

We consider the presheaf of sets
\[ x \longmapsto \pi_0(x\backslash C') \qquad \text{on } C, \]
let us call it $\pi_0(C'/C)$, we will have a canonical morphism
\[ \boxed{\varprojlim_{C'} \alpha^*(F) \xrightarrow{\ \sim\ } \operatorname{Hom}_{\substack{\text{presheaves}\\ \text{of sets on $C$}}} (\pi_0(C'/C), F),} \]
probably always \texttt{[an]} isomorphism.

If we work in the derived category of presheaves of $k$-modules ($k$ a ring), the preceding formula suggests conjecturally
\[ \boxed{\mathbf{R}\Gamma(C', \alpha^*F) \simeq \mathbf{R}\underbrace{\operatorname{Hom}}_{\text{internal Hom}}\ (\underbrace{\mathbf{L}\alpha_!(k_{C'})}_{\substack{\text{relative}\\ \text{homology of}\\ \text{$C'$ over $C$}}}, F).} \]

\msp{page 8}
The canonical augmentation
\[ \mathbf{L}\alpha_!(k_{C'}) \longrightarrow k_C \]
gives back, by taking the $\mathbf{R}\mathrm{Hom}$ into $F$,
\[ \mathbf{R}\Gamma(C,F) \longrightarrow \mathbf{R}\operatorname{Hom}(\mathbf{L}\alpha_!(k_{C'}), F) \xleftarrow{\ \sim\ } \mathbf{R}\Gamma(C', \alpha^*F) \]
the canonical homomorphism of ``contravariance'' of cohomology.

If on the other hand we work with the nonadditive derivator HOT, thus $F$ a ``relative homotopy type'' over $C$, we must have
\[ \boxed{\mathbf{R}\Gamma_{\mathrm{HOT}}(C', \alpha^*F) \xrightarrow{\ \sim\ } \mathbf{R}\Gamma(C\,,\ \underbrace{\operatorname{Hom}}_{\substack{\text{internal Hom}\\ \text{of the category of}\\ \text{relative homotopy}\\ \text{types over $C$}}}\ (\ \alpha_!\ (\ \underbrace{e_{C'}}_{\substack{\text{final object}\\ \text{of HOT($C'$),}\\ \text{final relative}\\ \text{homotopy type}\\ \text{over $C'$}}}\ ),\ F)).} \]

\printedpage{6}
% head: A. Grothendieck : Dérivateurs, Ch. II, Édition des Universités de Montpellier II et Paris VII
% foot: Version 8 juillet 2001
In fact, these formulas are quasi-tautological, if we admit the existence of the functors $\alpha_!$ as \textit{left} adjoints to the $\alpha^*$. In the abelian case, say,
\msp{page 9}
the isomorphy of the $H^0$ of the two sides is the adjunction formula
\[ \operatorname{Hom}_{C'}(k_{C'}, \alpha^*F) \simeq \operatorname{Hom}_C(\alpha_!(k_{C'}), F), \]
whence the isomorphy for the $H^i$ of the two sides, by shifting. Likewise in the case of HOT, the canonical isomorphy of the $\pi_0$ of the two sides is the adjunction formula. To have the finer isomorphism in $\mathbf{D}(e)$, the fundamental category of the derivator considered, in terms of internal Hom, we must admit the existence of unit objects (the $k_X$ in one case, the $e_X$ \texttt{[rather} $\varepsilon_X$\texttt{]} in the other), compatible with inverse images, in such a way \texttt{[that]} $\Homint(\varepsilon_X, \underbrace{F}_{\in\mathbf{D}_X}) \simeq F$. We will have
\[ \Gamma(C', \alpha^*F) \overset{\mathrm{def}}{=} f'_*(\alpha^*(F)) \simeq f_*(\alpha_*\alpha^*(F)), \]
\DERmark{ENV0078}

and the question is therefore to give an interpretation of $\alpha_*\alpha^*F$ as
\[ \alpha_*\alpha^*F \simeq \Homint(\alpha_!(\varepsilon_{C'}), F) \qquad \text{in } \mathbf{D}(C). \]
For this, we are going to write
\[ \alpha^*(F) \simeq \Homint_{C'}(\varepsilon_{C'},\alpha^* F), \]
\DERmark{ENV0079}

and the formula must follow from an ``internal adjunction formula''
\msp{page 10}
\[ ? \qquad\qquad \boxed{\alpha_*(\Homint_{C'}(G', \alpha^*F)) \simeq \Homint_C(\alpha_!G', F).} \]
(\textit{``Internal'' adjunction formula,} for $\alpha_!$ and $\alpha^*$, valid in principle for \textit{every} $\alpha : C' \longrightarrow C$ in Diag, every $G'$ in $\mathbf{D}(C')$ and every $F$ in $\mathbf{D}(C)$.)

How to prove this formula? By taking the product operation $\otimes$ associated with the operation Hom\,int, we are reduced (by taking the associated contravariant functors) to a formula
\[ \alpha_!(\alpha^*G \otimes G') \xrightarrow{\ \sim\ } G \otimes \alpha_!G'. \]
\DERmark{ENV0080}

We find by adjunction a canonical homomorphism, we must see whether it is always \texttt{[an]} isomorphism.

\printedpage{7}
% head: A. Grothendieck : Dérivateurs, Ch. II, Édition des Universités de Montpellier II et Paris VII
% foot: Version 8 juillet 2001
\msp{page 11}

\section*{2\quad Locally filtering categories are the locally totally $0$-connected categories}

Let $\mathcal{T}$ be a topos. We say that $\mathcal{T}$ is \textit{irreducible} if the ordered set $\operatorname{Ouv}(\mathcal{T})$ of subobjects of $e_{\mathcal{T}}$ has the following property: a) $\operatorname{Ouv}(\mathcal{T})$ is not reduced to one element, i.e. $\emptyset_{\mathcal{T}} \longrightarrow e_{\mathcal{T}}$ is not an isomorphism ($\mathcal{T}$ is not identical to its empty subtopos). b) If $U'$, $U''$ are two ``nonempty'' objects, i.e. noninitial, of $\operatorname{Ouv}(\mathcal{T})$, their intersection is ``nonempty''. It amounts to the same for b) to say that if $\mathcal{T}'$, $\mathcal{T}''$ are two closed subtopoi distinct from $\mathcal{T}$, then their supremum $\mathcal{T}' \cup \mathcal{T}''$ is distinct from $\mathcal{T}$, or again that every open of $\mathcal{T}$ is connected. We say that $\mathcal{T}$ is \textit{locally irreducible} if a) $\mathcal{T}$ is locally connected (i.e. each sheaf on $\mathcal{T}$ is a sum of 0-connected sheaves), and if b) for every 0-connected sheaf $U$ on $\mathcal{T}$, the induced topos $\mathcal{T}/U$ is irreducible (i.e. the ordered set $\operatorname{Ouv}(U)$ of subobjects of $U$ has property b) above ( a) simply means that $U \not\simeq \emptyset_{\mathcal{T}}$, i.e. that $U$ is not an initial object in $\mathcal{T}^\wedge$; or more simply): every subsheaf of a connected sheaf is connected (\footnote[5]{\textbf{NB} Irreducible does not imply locally irreducible, any more than connected does not imply locally connected etc.}) (\footnote[6]{$\mathcal{T}$ is locally irreducible if and only if it has a set of generators $(U_\alpha)$, with the $U_\alpha$ irreducible (hence 0-connected, hence $\mathcal{T}$ is locally connected).}).

When $\mathcal{T}$ is of the form $\operatorname{Top}(T)$, where $T$ is a topological space, then $\mathcal{T}$ irreducible means that $T$ is irreducible (for then $\operatorname{Ouv}(\mathcal{T}) \simeq \operatorname{Ouv}(T)$). This already implies trivially that $T$ is locally connected (every open is a disjoint sum of connected opens), thus $\operatorname{Top}({T})$ is locally connected. Moreover, this also implies that $\operatorname{Top}(T)$ is locally irreducible (\footnote[7]{make the proof explicit.}). (\textbf{NB} The 0-connected objects % [sic] \operatorname{Top}(\mathcal{T}) -> \operatorname{Top}(T), Top s'applique à l'espace T et non au topos \mathcal{T}
\DERmark{ENV0081}

\msp{page 12}
of $\operatorname{Top}(T)^\wedge$, i.e. the 0-connected sheaves of sets on $T$, are those whose ``generic fiber'' has exactly one point \dots) Thus $\mathcal{T} = \operatorname{Top}(T)$ is locally irreducible if and only if $T$ \texttt{[is]} locally connected, and the connected components of $T$ are irreducible, i.e. if and only if $T = \coprod_i T_i$, the $T_i$ irreducible.

If $\mathcal{T}$ is of the form $\operatorname{Top}(X)$, $X$ a category, we know that $\mathcal{T}$ is locally connected (and even locally $\infty$-connected). We see on the other hand that $\mathcal{T}$ is irreducible if and only if $\operatorname{Esp}(X)$ (the associated topological space, whose underlying set is $\operatorname{Ob} X$) is irreducible (\footnote[8]{for the opens of $\operatorname{Top}(X)$, those of $\operatorname{Esp}(X)$ and those of $X$ correspond to one another.}); and we then say that $X$ is irreducible (\footnote[9]{This also means that for $x$, $y$ in $X$, $\exists\, z$ with \begin{tikzcd}[ampersand replacement=\&, row sep=0.2em, column sep=1.2em, baseline=(\tikzcdmatrixname-2-1.base)] \& x \\ z \arrow[ur] \arrow[dr] \& \\ \& y \end{tikzcd}, i.e. $\operatorname{Ord}(X)$ is left filtering.}). As the connected components of $X$ and those of $\operatorname{Esp}(X)$ or $\operatorname{Top}(X) = \mathcal{T}$ correspond to one another, we see therefore that the following conditions are equivalent: a) the topos $\operatorname{Top}(X)$ is locally irreducible, b)
\printedpage{8}
% head: A. Grothendieck : Dérivateurs, Ch. II, Édition des Universités de Montpellier II et Paris VII
% foot: Version 8 juillet 2001
$\operatorname{Esp}(X)$ is locally irreducible, c) the connected components of $X$ are irreducible, d) the $X/x$ \texttt{[are]} irreducible. We then say that $X$ is locally irreducible (\footnote[10]{False, we have (a $\Longleftrightarrow$ d) $\Longrightarrow$ (b $\Longleftrightarrow$ c) $X$ irreducible does \textit{not} imply that the $X/x$ are so, example $X = (0 \overset{\alpha}{\underset{\beta}{\rightrightarrows}} 1)$ $(\alpha \neq \beta)$, we have $X/1 =$ \begin{tikzcd}[ampersand replacement=\&, row sep=0.2em, column sep=1.2em, baseline=(\tikzcdmatrixname-2-2.base)] 0' \arrow[dr] \& \\ \& 1 \\ 0'' \arrow[ur] \& \end{tikzcd}, which is \textit{not} irreducible, nor locally irreducible since it is connected. Yet it is étale over $X$.}).

This therefore implies that for every $X'$ étale over $X$, with $X'$ 0-connected, $X'$ is irreducible.

\textit{State equivalent conditions.}
\begin{enumerate}
\item[a)] $\operatorname{Top}(X)$ is locally irreducible.
\item[b)] The $X/x$ \texttt{[are]} irreducible.
\item[c)] Every 0-connected object of $X^\wedge$ is irreducible.
\item[d)] The connected components of $X$ are weakly cofiltering (i.e. weakly left filtering).
\item[e)] Condition PS $1_{\mathrm{g}}$ of SGA 4, I 2.
\end{enumerate}
If $X$ is \texttt{[a]} category, $X$ irreducible also means that $X \neq \emptyset$ and $\forall\, x, y \in \operatorname{Ob} X$, \texttt{[there]} exists \texttt{[an]} $s$ with \begin{tikzcd}[row sep=0.2em, column sep=1.2em, baseline=(\tikzcdmatrixname-2-1.base)] & x \\ s \arrow[ur] \arrow[dr] & \\ & y \end{tikzcd}. (This also means that $X$ is \textit{weakly filtering} (\footnote[11]{It is better to call it \textit{pseudo-filtering,} and let us reserve weakly filtering for: locally weakly filtering + connected.}) on the left, or weakly cofiltering.) $X$ locally irreducible means that for $x$, $y$ \textit{belonging to the same connected component,} \texttt{[there]} exists \texttt{[an]} $s$ with \begin{tikzcd}[row sep=0.2em, column sep=1.2em, baseline=(\tikzcdmatrixname-2-1.base)] & x \\ s \arrow[ur] \arrow[dr] & \\ & y \end{tikzcd}, % [sic] irreductible -> irréductible, accent omis
\msp{page 13}
i.e. $X/[x] \cap X/[y] \neq \emptyset$ (where $[x]$ is regarded as an element of the ordered set $\operatorname{Ob} X$). It again amounts to the same to say that the induced categories $X/x$ are irreducible (\footnote[12]{false, cf. marginal annotation above.}).

\textbf{Example.} A category \textit{having an initial object,} or merely a ``smallest object'', is irreducible. Conversely, if the \textit{ordered} set $\operatorname{Ord}(X)$ associated with $X$ (quotient of $\operatorname{Ob}(X)$ by the equivalence relation associated with its preorder) is finite, e.g. $X$ finite, or only locally finite, e.g. $X$ locally finite (i.e. the categories $X/x$ finite), then if $X$ is irreducible, $\operatorname{Ob} X$ has a smallest element, i.e. $X$ has a smallest object. This is true more generally still, if $\operatorname{Ord}(X)$ is artinian, i.e. every decreasing sequence in $\operatorname{Ord}(X)$ is stationary. % [sic] initiale -> initial, accord avec « objet »

But the interesting condition is the topos-theoretic condition of \textit{local} irreducibility, which is equivalent to: $I$ left pseudo-filtering, i.e. the connected components of $I$ filtering, i.e. conditions PS 1, PS 2

PS 1\enspace Every diagram \begin{tikzcd}[row sep=0.2em, column sep=1.2em, baseline=(\tikzcdmatrixname-2-2.base)] x \arrow[dr] & \\ & s \\ y \arrow[ur] & \end{tikzcd} can be completed to \begin{tikzcd}[row sep=0.2em, column sep=1.2em, baseline=(\tikzcdmatrixname-2-1.base)] & x \arrow[dr] & \\ z \arrow[ur] \arrow[dr] & & s \\ & y \arrow[ur] & \end{tikzcd} \textit{commutative.}

PS 2\enspace For every double arrow $x \rightrightarrows y$, \texttt{[there]} exists an equalizer $z \longrightarrow x \rightrightarrows y$.


\endgroup
% END UNIT Der-ChII-234/Der-ChII-234_p006-010.tex
% BEGIN UNIT Der-ChII-234/Der-ChII-234_p011-015.tex
\begingroup\makeatletter%
\newcommand{\printedpage}[1]{\marginpar{\footnotesize p.~#1}}%
\newcommand{\msp}[1]{\marginpar{\footnotesize [#1]}}%
\newcommand{\sX}{\mathcal{X}} % resolved against PDF: CMSY calligraphic X
\makeatother\let\printedpage\derpp\let\msp\dermsp%
\let\mspin\DERmspin%


\printedpage{9}
% head: A. Grothendieck : Dérivateurs, Ch. II, Édition des Universités de Montpellier II et Paris VII
% foot: Version 8 juillet 2001

\msp{page 14}

In summary, we have the

\noindent\textbf{Proposition 2.1.} \textit{Let $X$ be an object of} Cat. \textit{Equivalent conditions:}
\begin{enumerate}[label={},leftmargin=3em]
\item[\textit{a)}] \textit{The associated topos} $\operatorname{Top}(X)$ \textit{is locally irreducible, i.e. (being locally connected) for every object $F$ of $X^{\wedge}$ and every subobject $F'$ of $F$, if $F$ is connected, so is $F'$.}
\item[\textit{a$'$)}] \textit{For every $X'$ étale over $X$, if $X'$ is connected, every open of $X'$ is connected.}
\item[\textit{a$''$)}] \textit{Every open of a localized category $X/a$ ($a \in \operatorname{Ob} X$) is connected (i.e. $X/a$ is an irreducible category, i.e. the topos} $\operatorname{Top}(X/a)$ \textit{is).}
\item[\textit{b)}] \textit{For every object $a$ of $X$,} $\operatorname{Esp}(X/a)$ \textit{is an irreducible topological space.}
\item[\textit{c)}] \textit{The categories $X/a$ are weakly left filtering (or weakly cofiltering), i.e. for two objects $\xi$, $\eta$ of such a category,} \texttt{[there]} \textit{exists an object $\zeta$ and two arrows}
\begin{tikzcd}[ampersand replacement=\&, row sep=tiny, column sep=small, baseline=(\tikzcdmatrixname.center)]
 \& \xi \\
\zeta \arrow[ur] \arrow[dr] \& \\
 \& \eta
\end{tikzcd}
\textit{(i.e. the ordered set Ord($X/a$) associated with $X/a$ is left filtering)}
\item[\textit{c$'$)}] \textit{For every diagram in solid lines in $X$,}
\begin{tikzcd}[ampersand replacement=\&, row sep=tiny, column sep=small, baseline=(\tikzcdmatrixname.center)]
 \& x \arrow[dr] \& \\
z \arrow[ur, dotted] \arrow[dr, dotted] \& \& s \\
 \& y \arrow[ur] \&
\end{tikzcd}
\textit{, we can complete it to a} commutative \textit{square as indicated. (Condition} PS 1 \textit{for locally filtering categories.)}
\end{enumerate}

\noindent These conditions imply that every connected component $X_{\alpha}$ of $X$ is irreducible, or what amounts to the same thing, that $\operatorname{Esp}(X_{\alpha})$ is irreducible,

\msp{page 15}

\noindent but are not implied by the latter.

Counterexample, the 0-connected category $0 \underset{\beta}{\overset{\alpha}{\rightrightarrows}} 1$, with $\alpha \neq \beta$. We have $X/1 \simeq$
\begin{tikzcd}[ampersand replacement=\&, row sep=tiny, column sep=small, baseline=(\tikzcdmatrixname.center)]
0 \arrow[dr] \& \\
 \& 1 \\
0' \arrow[ur] \&
\end{tikzcd}, which is \textit{not} irreducible, since $\Phi \backslash \{1\} \simeq \varepsilon$ is a disconnected open \texttt{[}$\Phi = X/1$\texttt{,} $\varepsilon = \bullet\ \bullet$\texttt{]}.

I would like a ``topossic'' condition on $X$, equivalent to $X$ ``left pseudo-filtering'', i.e. the connected components filtering, i.e. PS 1\textsuperscript{opp} (condition c$'$) above) + PS 2\textsuperscript{opp} (double arrow condition), and which moreover implies that every 0-connected object of $X^{\wedge}$ (i.e. every object $X'$ of Cat étale over $X$ and 0-connected) is $W_{\infty}$-aspheric. And I am thinking of the notion of ``total asphericity'' of Pursuing Stacks.

\noindent\textbf{Definition 2.2.} (\footnote[13]{Cf. summary table pages 97, 98 and especially 99.})\ \ Let $X$ be a category.
\begin{enumerate}[label={},leftmargin=3em]
\item[a)] We say that $X$ is \textit{weakly filtering} (\footnote[14]{rather say \textit{pseudo-filtering}, reserve ``weakly filtering'' for ``locally weakly filtering'' (or \textit{irreducible}) + connected, condition PS $1'_{\mathrm{g}}$.}) if $X \neq \emptyset$, and if for $a, b \in \operatorname{Ob} X$, $\exists\, c$ and arrows $c$
\begin{tikzcd}[ampersand replacement=\&, row sep=tiny, column sep=small, baseline=(\tikzcdmatrixname.center)]
 \& a \\
{} \arrow[ur, phantom] \& \\
 \& b
\end{tikzcd}
\hspace{-2.6em}% layout aid only
\begin{tikzcd}[ampersand replacement=\&, row sep=tiny, column sep=small, baseline=(\tikzcdmatrixname.center)]
 \& {} \\
{} \arrow[ur] \arrow[dr] \& \\
 \& {}
\end{tikzcd}
, i.e. $X/[a] \cap X/[b] \neq \emptyset$. (\textbf{NB} $X$ is weakly left filtering (\footnote[15]{pseudo-filtering.}) if and only if $\operatorname{Ord}(X)$ is, i.e. is left filtering. \textbf{ReNB} $X$ weakly filtering $\Longrightarrow$ $X$ 0-connected.)

\item[b)] \printedpage{10}
% head: A. Grothendieck : Dérivateurs, Ch. II, Édition des Universités de Montpellier II et Paris VII
% foot: Version 8 juillet 2001
We say that $X$ is \textit{locally weakly left filtering} or more simply left prefiltering if the categories $X/x$ ($x \in \operatorname{Ob} X$) are weakly left filtering, i.e. if $X$ satisfies the condition: $\forall$ diagrams
\begin{tikzcd}[ampersand replacement=\&, row sep=tiny, column sep=small, baseline=(\tikzcdmatrixname.center)]
a \arrow[dr] \& \\
 \& x \\
b \arrow[ur] \&
\end{tikzcd}, $\exists\, c$ in $X$ and arrows from $c$ to $a$, $b$ making the square \textit{commutative}
\[
\begin{tikzcd}[ampersand replacement=\&, row sep=tiny, column sep=small]
 \& a \arrow[dr] \& \\
c \arrow[ur, dotted] \arrow[dr, dotted] \& \& x \\
 \& b \arrow[ur] \&
\end{tikzcd}
\]
(\footnote[16]{Condition PS $1_{\mathrm{g}}$.}). (\textbf{NB} This condition implies that $\operatorname{Ord}(X)$ is locally weakly filtering, but is not equivalent to the latter. Example, $X = 0 \underset{\beta}{\overset{\alpha}{\rightrightarrows}} 1$, $X/1 =$
\begin{tikzcd}[ampersand replacement=\&, row sep=tiny, column sep=small, baseline=(\tikzcdmatrixname.center)]
0' \arrow[dr] \& \\
 \& 1 \\
0'' \arrow[ur] \&
\end{tikzcd}
is not weakly left filtering.)

\msp{page 16}

\textbf{NB} If $X$ locally weakly left filtering, every $X'$ étale over $X$ likewise.

\textbf{NB} (\footnote[17]{\textbf{False.}}) $X$ locally weakly left filtering $\Longleftrightarrow$ the connected components are weakly left filtering, or again if $X$ is 0-connected, $X$ locally weakly left filtering $\Longleftrightarrow$ $X$ is weakly left filtering (\footnote[18]{pseudo-filtering.}). We have indeed $\Longleftarrow$ tautologically (\footnote[19]{\textbf{false,} thus $0 \underset{\beta}{\overset{\alpha}{\rightrightarrows}} 1$ is 0-connected and weakly left filtering, but not locally weakly left filtering because
\begin{tikzcd}[ampersand replacement=\&, row sep=tiny, column sep=small, baseline=(\tikzcdmatrixname.center)]
0 \arrow[dr, "\alpha"] \& \\
 \& 1 \\
0 \arrow[ur, "\beta"'] \&
\end{tikzcd}
cannot be completed to a \textit{commutative} square.}), and $\Longrightarrow$ by showing that if PS 1 \texttt{[is]} satisfied, then $PS(a,b)$ defined by $X/[a] \cap X/[b] \neq \emptyset$ is an equivalence relation, implied by the relation $C(a,b)$: $\{\operatorname{Hom}(a,b) \neq \emptyset$ \textit{or} $\operatorname{Hom}(b,a) \neq \emptyset\}$, hence the equivalence relation (homotopy) generated by $C(a,b)$ is finer than $PS(a,b)$.

\item[c)] We say that $X$ is \textit{left filtering}, or \textit{cofiltering,} if it is weakly left filtering (\footnote[20]{pseudo-filtering.}) (hence \textit{nonempty}), and satisfies (in addition to condition PS $1'_{\mathrm{g}}$ of a)) the following condition PS $2_{\mathrm{g}}$:

For every double arrow $a \underset{v}{\overset{u}{\rightrightarrows}} b$, \texttt{[there]} exists \texttt{[an]} arrow $c \xrightarrow{w} a$ such that $uw = vw$ (\footnote[21]{i.e. three conditions a) $X \neq \emptyset$ b) usual PS $1'$ c) PS 2 (double arrow).}).

It amounts to the same to say that $X$ satisfies the three conditions a$'$) $X$ is 0-connected (stronger than $X \neq \emptyset$), b$'$) PS 1, i.e. $X$ locally weakly left filtering and c$'$) PS 2 as before.

\item[d)] We say that $X$ is \textit{locally left filtering} if the connected components are left filtering, i.e. if $X$ satisfies conditions PS 1 and PS 2 (``pseudo-cofiltering'' in the terminology of SGA 4). \textbf{NB} If $X$ is locally left filtering, the same holds for every $X'$ étale over $X$, for this is obvious for PS 1 which concerns the $X/x$, and for PS 2 this follows from the fact that $X' \longrightarrow X$ is faithful. So if $X$ is locally cofiltering, the $X/a$ are filtering, but the converse is false.
\end{enumerate}

\printedpage{11}
% head: A. Grothendieck : Dérivateurs, Ch. II, Édition des Universités de Montpellier II et Paris VII
% foot: Version 8 juillet 2001

\noindent We have the dual notions: \textit{weakly right filtering} (or simply weakly filtering), and \textit{right filtering} or simply filtering. (Dual notions of a), c) above.) For the dual notion of b), namely that the $x\backslash X$ be weakly filtering, we can no longer call it ``locally weakly filtering'', but it would be \textit{co}locally weakly filtering.

\noindent\textbf{Example.} $X = \Bigl(0 \xrightarrow{\gamma} 1 \underset{\beta'}{\overset{\alpha'}{\rightrightarrows}} 2 \underset{\beta}{\overset{\alpha}{\rightrightarrows}} 3\Bigr)$ with $(\alpha\alpha' = \beta\beta') \neq (\alpha\beta' = \beta\alpha')$.

\msp{page 17}

\noindent Only the notion a) (weakly left filtering (\footnote[22]{pseudo-filtering.})) depends only on $\operatorname{Ord} X$ - but of course if $X$ is \textit{pre}ordered, all three depend only on $\operatorname{Ord} X$. Let us note that saying it depends only on $\operatorname{Ord} X$, or only on $\mathrm{Esp}\,\mathrm{sep}X$, amounts to the same thing, since the two determine each other. Condition a) means that $\operatorname{Esp} X$, or also $\operatorname{Top}(X)$, is \textit{irreducible.} We shall say that $X$ is \textit{irreducible} (instead of pseudo-filtering). The other three notions cannot be expressed in terms of $\operatorname{Esp} X$ or $\operatorname{Ord} X$ alone, but they can be expressed in terms of the $\mathrm{Topos}\, X$. Condition b) ($X$ locally weakly filtering) means that $\operatorname{Top}(X)$ is \textit{locally irreducible,} and we shall also say that $X$ is \textit{locally irreducible} (instead of locally weakly filtering). This means that for every $X'$ étale over $X$ and 0-connected, $X'$ is weakly left filtering, i.e. is irreducible.

It remains to express the filtering condition itself in terms of $\operatorname{Top} X$. I recall in this connection:

\noindent\textbf{Proposition 2.3.} \textit{The following conditions on a topos $\sX$ are equivalent (when $S$ (as a ``site'') denotes a generating subset of} $\mathrm{Fais}(\sX) \simeq S^{\sim}$, consisting of 0-connected objects (\footnote[23]{thus $\sX$ is locally 0-connected.}).\textit{)}
\DERmark{ENV0082}

\begin{enumerate}[label={},leftmargin=3em]
\item[\textit{a)}] \textit{If $U$, $V$ are two $0$-connected objects of $\operatorname{Ob} \sX$, $U \times V$ is $0$-connected.}
\item[\textit{b)}] \textit{If $U$ is a $0$-connected object of $\sX^{\wedge}$, then $U \longrightarrow e$ is ``$W_0$-aspheric'', i.e. for every $V$, the projection $V \times U \longrightarrow V$ induces $\pi_0(U \times V) \xrightarrow{\sim} \pi_0(V)$.}
\item[\textit{c)}] \textit{If $U$, $V$} \texttt{[are]} \textit{in $\sX^{\wedge}$, then}
\[
\pi_0(U \times V) \xrightarrow{\sim} \pi_0(U) \times \pi_0(V).
\]
\DERmark{ENV0083}

\item[\textit{d)}] \textit{As in a), with $U$, $V$ in $S$. (But we assume the objects of $S$ to be $0$-connected.)}
\end{enumerate}
\noindent(\footnote[24]{The proof should be made explicit together with that of 2.9 below (p.~25).}). % [sic] démonstation -> démonstration, letter r missing

\msp{page 18}

\noindent Moreover, these conditions imply that $\sX$ is connected, hence 0-connected if $\sX \neq \emptyset$. This moreover implies that $\sX$ is $W_{\infty}$-aspheric, cf. Pursuing Stacks §35 (Corollary p.~52) (\footnote[25]{\textbf{False,} loc. cit. says nothing of the kind.}).

It is essentially Proposition 1 of Pursuing Stacks §65, cf. loc. cit. p.~170-171 (\footnote[26]{But there I restrict myself to topoi of the form $\operatorname{Top}(A)$, with $A$ in Cat.}). Cf. also Pursuing Stacks, §35, pages 50-55.

\printedpage{12}
% head: A. Grothendieck : Dérivateurs, Ch. II, Édition des Universités de Montpellier II et Paris VII
% foot: Version 8 juillet 2001

\noindent\textbf{Definition 2.4.} We say that $\sX$ is \textit{totally $0$-connected} if it is locally connected, nonempty and satisfies the equivalent conditions of 2.3. (NB this is a special case of the notion ``totally $W$-aspheric'' of loc. cit. Thus we have the notion of totally $n$-connected.)

Let us note that for an object $X$ of Cat, saying that $\operatorname{Top}(X)$ is totally 0-connected amounts to saying that $X \neq \emptyset$, and that for $x, y \in \operatorname{Ob} X$, $x \times y$ in $X^{\wedge}$ is 0-connected, i.e. $X/x \times y$ is 0-connected (\footnote[27]{This also means that $\mathrm{diag}_X : X \longrightarrow X \times X$ is co-cofinal, i.e. $X^o \longrightarrow X^o \times X^o$ is cofinal, and implies that the same holds for $X \longrightarrow X^E$ for every \textit{finite} set $E$.}). It is a very strong condition. We shall also say that $X$ is \textit{totally $0$-connected.} Let us note here

\noindent\textbf{Proposition 2.5.} \textit{Let $\sX$ be a topos. If $\sX$ is totally $0$-connected, then $\sX$ is irreducible. (In fact the conclusion holds, tautologically, assuming only that the product of two $0$-connected objects of $\sX^{\wedge}$ is nonempty, and it even suffices to assume this for subobjects of the final object $e_{\sX}$.) In particular, if $\sX = \mathit{Top}(X)$, where $X$ is in} Cat\textit{, then $X$ is weakly left filtering.}

But of course, this does \textit{not} imply that $\sX$ is locally irreducible.

\msp{page 19}

\noindent For example the category $\boldsymbol{\Delta}$, or any small category stable under binary products, is totally 0-connected, but $\boldsymbol{\Delta}$ is not locally irreducible (the induced categories $\boldsymbol{\Delta}/\boldsymbol{\Delta}^n$, for $n > 0$, are visibly \textit{not} irreducible). (But let us recall that if $\sX = \operatorname{Top}(T)$, $T$ a topological space, then $T$ irreducible implies that $\sX$ is also \textit{locally} irreducible. So $T$ totally 0-connected implies $\operatorname{Top}(T)$ globally \textit{and} locally irreducible, i.e. for every $T'$ étale over $T$, and 0-connected, $T'$ is irreducible.)

\noindent\textbf{Theorem 2.6.} \textit{Let $\sX$ be a locally connected topos, $S \subseteq \sX^{\wedge}$ a small generating subset consisting of $0$-connected objects. Equivalent conditions (}except \textit{c, c$'$, which seem} (\footnote[28]{\textbf{are.}}) \textit{weaker, cf. proof).)}
\begin{enumerate}[label={},leftmargin=3em]
\item[\textit{a)}] \textit{Every induced topos $\sX/U$ over $U$ $0$-connected in $\sX^{\wedge}$ is totally $0$-connected.}
\item[\textit{a$'$)}] \textit{Every fiber product $U \times_W V$ of $0$-connected objects in $\sX^{\wedge}$ is $0$-connected.}
\item[\textit{b)}] \textit{$\sX$ is locally irreducible, and its connected components are totally $0$-connected.}
\item[\textit{c)}] \textit{As in a), with $U$ in $S$, i.e. the $\sX/U$ ($U \in \operatorname{Ob} S$) are totally $0$-connected.}
\item[\textit{c$'$)}] \textit{As in a$'$), with $U$, $V$, $W$ in $S$.}
\item[\textit{d)}] \textit{$S$ is locally left filtering, i.e. $S$ satisfies the two conditions}

PS$^*$ 1\ \ \textit{Every diagram}
\begin{tikzcd}[ampersand replacement=\&, row sep=tiny, column sep=small, baseline=(\tikzcdmatrixname.center)]
a \arrow[dr] \& \\
 \& x \\
b \arrow[ur] \&
\end{tikzcd}
\textit{in $S$ fits into}
\begin{tikzcd}[ampersand replacement=\&, row sep=tiny, column sep=small, baseline=(\tikzcdmatrixname.center)]
 \& a \arrow[dr] \& \\
c \arrow[ur, dotted] \arrow[dr, dotted] \& \& x \\
 \& b \arrow[ur] \&
\end{tikzcd}
\textit{which is commutative.}

PS$^*$ 2\ \ \textit{Every diagram $a \rightrightarrows b$ fits into a commutative diagram $c \longrightarrow a \rightrightarrows b$.}
\end{enumerate}
\noindent(\footnote[29]{NB We have equivalence of \tikz[baseline=(n.base)]\node[draw, rounded corners, inner sep=1.5pt](n){a), a$'$), b), d),}; which imply c), c$'$), which are equivalent. But it seems (\footnotemark[30]) that c), c$'$) are strictly weaker. Cf. however the variants e), e$'$) in 2.12 (p.~28, 29).}). % [sic] plus faible -> plus faibles, adjective does not agree with the plural subject c), c')
\footnotetext[30]{It turns out. It is indeed true, c) and c$'$) do \textit{not} imply the other conditions, cf. marginal annotation
% footnote 30 continues at the foot of printed page 13:
p.~28 (case of $X = B_G$).}

\printedpage{13}
% head: A. Grothendieck : Dérivateurs, Ch. II, Édition des Universités de Montpellier II et Paris VII
% foot: Version 8 juillet 2001

\msp{page 20}

\noindent These conditions imply that the 0-connected objects of $\sX^{\wedge}$ (in particular the elements of $S$) are $\infty$-connected.

\noindent\textbf{Corollary 2.7} (of the equivalence b $\Longleftrightarrow$ d). \textit{Let $X$ be an object of} Cat. \textit{For $X$ to be left filtering it is necessary and sufficient that $\mathit{Top}(X)$ be totally $0$-connected and locally (and therefore also globally) irreducible (or again that $X$ be totally $0$-connected and that for every $x \in \operatorname{Ob} X$, the localized category $X/x$ be irreducible (it is even totally $0$-connected), or again that $X$ be nonempty and that for two objects $a$, $b$ in $X$, $a \times b \in X^{\wedge}$ be $0$-connected, and for}
\DERmark{ENV0084}

\begin{tikzcd}[ampersand replacement=\&, row sep=tiny, column sep=small, baseline=(\tikzcdmatrixname.center)]
a \arrow[dr] \& \\
 \& x \\
b \arrow[ur] \&
\end{tikzcd}
\textit{in $X$, $a \times_x b$ (in $X^{\wedge}$) be $0$-connected)} (\footnote[31]{Example of $X$ totally 0-connected, but not cofiltering, i.e. not locally irreducible? The category $\boldsymbol{\Delta}$ of standard simplices is totally 0-connected, but is not cofiltering because $\boldsymbol{\Delta}^0 = e \underset{\partial_1}{\overset{\partial_0}{\rightrightarrows}} \boldsymbol{\Delta}^1$ cannot be equalized! So \texttt{[the opposite category]} $\boldsymbol{\Delta}^o = I$ is such that $I \longrightarrow I^E$ is cofinal for every finite set $E$, but $I$ not filtering.}).

\noindent\textsc{Proof} of 2.6. We have the diagram of $\pm$ tautological implications
\[
\begin{tikzcd}[ampersand replacement=\&, column sep=large, row sep=large]
 \& \mathrm{a} \arrow[r, Leftrightarrow] \arrow[dl, Rightarrow, "{\substack{\text{by}\\ \text{2.5}}}"'] \arrow[d, Rightarrow] \& \mathrm{a}' \arrow[d, Rightarrow] \\
\mathrm{b} \arrow[d, Rightarrow] \& \mathrm{c} \arrow[r, Leftrightarrow, "{(\Leftarrow \text{ by 2.3.})}"] \& \mathrm{c}' \arrow[dll, Rightarrow] \\
\text{PS 1.} \& \&
\end{tikzcd}
\]
Let us nevertheless show the two implications c$'$) $\Longrightarrow$ PS 1 and b $\Longrightarrow$ PS 1.

\msp{page 21}

\noindent This is contained in the

\noindent\textbf{Corollary 2.8.} \textit{Let $\sX$, $S$ be as in 2.6. Equivalent conditions:}
\begin{enumerate}[label={},leftmargin=3em]
\item[\textit{a)}] \textit{$\sX$ is locally irreducible, i.e. the induced topoi $\sX/U$, where $U$ is $0$-connected, are irreducible.}
\item[\textit{b)}] \textit{Every fiber product $U \times_W V$ of $0$-connected objects $U$, $V,W$ is ``nonempty''.}
\item[\textit{b$'$)}] \textit{Every fiber product $U \times_W V$, with $U$, $V$ ``nonempty'', $W$ $0$-connected, is nonempty.}
\item[\textit{c)}] \textit{The $\sX/x$ with $x \in S$ are irreducible.}
\item[\textit{d)}] \textit{As in b), with $U$, $V$, $W$ in $S$.}
\item[\textit{e)}] \textit{$S$ is locally weakly left filtering, i.e. satisfies PS$^*$ 1.}
\end{enumerate}
\DERmark{ENV0085}

\noindent Here the $\pm$ obvious implications are


\endgroup
% END UNIT Der-ChII-234/Der-ChII-234_p011-015.tex
% BEGIN UNIT Der-ChII-234/Der-ChII-234_p016-020.tex
\begingroup\makeatletter%
\newcommand{\printedpage}[1]{\marginpar{\footnotesize p.~#1}}%
\newcommand{\msp}[1]{\marginpar{\footnotesize [#1]}}%
\newcommand{\X}{\mathcal{X}}%
\makeatother\let\printedpage\derpp\let\msp\dermsp%
\let\mspin\DERmspin%


\printedpage{14}
% head: A. Grothendieck : Dérivateurs, Ch. II, Édition des Universités de Montpellier II et Paris VII
% foot: Version 8 juillet 2001

\[
\begin{tikzcd}
\mathrm{b}' \arrow[r, Rightarrow] & \mathrm{b} \arrow[rr, Rightarrow] \arrow[d, Rightarrow] & & \mathrm{a} \arrow[d, Rightarrow] \\
 & \mathrm{d} \arrow[r, Leftrightarrow] & \mathrm{e} \arrow[r, Leftrightarrow] & \mathrm{c}.
\end{tikzcd}
\]

Let us prove the implications on the bottom line. Let us consider d). If $\begin{tikzcd}[cramped, row sep=0.1em, column sep=0.8em, baseline=(current bounding box.center)] a \arrow[dr] & \\ & x \\ b \arrow[ur] & \end{tikzcd}$ \texttt{[is]} in $S$, to say that the fiber product in $\X$ is ``nonempty'' means, since $S$ is a system of generators and in particular \textit{nonempty,} that there exists $c$ in $S$ and $c\longrightarrow a\times_x b$, or what amounts to the same, $\begin{tikzcd}[cramped, row sep=0.1em, column sep=0.8em, baseline=(current bounding box.center)] & a \\ c \arrow[ur, dotted] \arrow[dr, dotted] & \\ & b \end{tikzcd}$, making the diagram $\begin{tikzcd}[cramped, row sep=0.1em, column sep=0.8em, baseline=(current bounding box.center)] & a \arrow[dr] & \\ c \arrow[ur, dotted] \arrow[dr, dotted] & & x \\ & b \arrow[ur] & \end{tikzcd}$ commutative. Thus d) is equivalent to PS 1, i.e. to e) ($\X$ locally weakly left filtering). Let us consider condition c): the $\X/a=(S/a)^\sim$ are irreducible. This means that for two \textit{nonempty} subobjects $U,V\subseteq a$ in $\X^\wedge$, we have $U\cap V\neq\emptyset$, or else that for two sieves

\msp{page 22}$\underline{U}$, $\underline{V}$ of $S/x$, if $\underline{U}$, $\underline{V}$ nonempty (which means that the subobjects of $x$ in $\X^\wedge$ which they define are ``nonempty'', or that the objects of $S$ are ``nonempty'' in $\X^\wedge$), then $\underline{U}\cap\underline{V}$ is nonempty. We may suppose that $\underline{U}$, $\underline{V}$ is each generated by a single arrow $a\xrightarrow{u}x$, $b\xrightarrow{v}x$. The intersection sieve consists of the objects $c\longrightarrow x$ of $S/x$ such that $c\longrightarrow x$ factors through a $\begin{tikzcd}[cramped, row sep=0.1em, column sep=0.8em, baseline=(current bounding box.center)] & a \\ c \arrow[ur, dotted] \arrow[dr, dotted] & \\ & b \end{tikzcd}$, so to say that this sieve is not empty means that PS 1 is satisfied for the two given arrows $u$, $v$. Thus c) is equivalent to condition PS 1\texttt{[, i.e. to]} e.

Thus we have established
\[
\mathrm{b}'\Longrightarrow\mathrm{b}\Longrightarrow\mathrm{a}\Longrightarrow(\mathrm{c}\Longleftrightarrow\mathrm{d}\Longleftrightarrow\mathrm{e}),
\]
it remains to prove that c $\Longrightarrow$ b$'$. So let $\begin{tikzcd}[cramped, row sep=0.1em, column sep=0.8em, baseline=(current bounding box.center)] U \arrow[dr] & & V \arrow[dl] \\ & W & \end{tikzcd}$ be a diagram of ``nonempty'' objects, with $W$ 0-connected, we must prove that the fiber product is nonempty. Considering $a\longrightarrow U$ and $b\longrightarrow V$ with $a,b\in S$, we are reduced to the case where $U=a$, $V=b$ in $S$. Let us prove that if $a\times_W b$ is ``empty'', then $W$ is disconnected. For this, consider the subcategory $C$ of $S/W$ consisting of the $x\longrightarrow W$ such that $a\times_W x$ is ``nonempty''. Denoting by $w(a)$, $w(x)$ the image of $a$, $x$ in $W$, this means that $w(a)\cap w(x)\neq\emptyset$.

\msp{page 23}Note that any quotient in $\X^\wedge$ of an irreducible object is irreducible, so $w(a)$, $w(x)$ are irreducible objects. This shows that for $x$ in $C$ and $x'$ in $S/W$, we have $x'$ in $C$, i.e. $w(a)\cap w(x')\neq\emptyset$ if and only if $w(x)\cap w(x')\neq\emptyset$. (Direct lemma: in an object $W$ of a topos $\X$, if we have \texttt{[some]} irreducible subobjects $W_1$, $W_2$, $W_3$, if one of the two \texttt{[rather three]} meets the other two, these meet.) This shows that $C$ is a sieve, since $x'\longrightarrow x$ implies $w(x')\subseteq w(x)$, so $w(x')$ meets $w(x)$, hence also $w(a)$. But then the cosieve $C'=(S/W)\backslash C$ is also a sieve, for if $c$
\printedpage{15}
% head: A. Grothendieck : Dérivateurs, Ch. II, Édition des Universités de Montpellier II et Paris VII
% foot: Version 8 juillet 2001
is in $C'$, i.e. \texttt{[if]} $w(c)$ does not meet $w(x)$, a fortiori not any $x'$ over $x$, for which we will have $w(x')\subseteq w(x)$. Thus, $S/W$ is decomposed as a sum of two opens $C$, $C'$, nonempty since one contains $a$, the other $b$. Consequently, $W$ is decomposed as a sum of two nonempty subobjects, so $W$ is not 0-connected.

We have proved 2.8, hence also (returning to 2.6) the implications toward PS 1 in the diagram on page 20. It remains to prove
\[
\left\{
\begin{aligned}
&\mathrm{b}\Longrightarrow\mathrm{PS}\ 2,\ \mathrm{c}'\Longrightarrow\mathrm{PS}\ 2\\
&\mathrm{PS}\ 1+\mathrm{PS}\ 2\ (\text{i.e. d})\Longrightarrow\mathrm{a}'
\end{aligned}
\right.
\]
(\footnote[32]{NB in the end we do not prove c$'\Longrightarrow$ PS 2.}).

\msp{page 24}We therefore place ourselves under the hypothesis where PS 1 is satisfied (since b), c$'$) and d) imply PS 1), thus where we are under the equivalent conditions of 2.8 - $\X$ is locally irreducible. Let us then consider condition PS 2, and prove that it is implied both by b), and by c$'$), and that it implies a$'$). It means that for a double arrow $a\overset{u}{\underset{v}{\rightrightarrows}}b$ in $S$, the cokernel \texttt{[rather kernel]} in $S^\wedge$, or, what amounts to the same, in $\X^\wedge=S^\sim\subseteq S^\wedge$, is ``nonempty''. We can write it in the form of a fiber product
\[
\begin{tikzcd}
a \arrow[r, "{\scriptstyle (u,v)}"] & b\times b \\
\operatorname{Ker}(u,v) \arrow[u] \arrow[r] & b. \arrow[u, "\mathrm{diag}"']
\end{tikzcd}
\]
Here by hypothesis $a$ and $b$ are 0-connected objects of $\X^\wedge$ (those of $S$ being so), b) holds (($\X$ locally irreducible and) its connected components \texttt{[are]} totally 0-connected), then $b\times b$ is 0-connected as a product of two 0-connected objects over the \textit{same connected component,} which is totally 0-connected. Since $\X$ is locally irreducible, it follows that the fiber product is nonempty (2.8), hence b) $\Longrightarrow$ d). Let us prove d) $\Longrightarrow$ a), i.e. if
\DERmark{ENV0086}\DERmark{ENV0087}


\msp{page 25}$S$ is locally cofiltering, then the induced topoi $\X/W$, where $W$ is 0-connected, are totally 0-connected. This induced topos is described by the site $S/W$, i.e. $S/W$ is a system of generators. Now if $S$ is locally filtering, $S/W$ (étale over $S$) is so too. So we may suppose $W=e$, and to prove that if $S$ is locally cofiltering and $\X$ 0-connected, then $\X$ is totally 0-connected. Note that $\X$ 0-connected is equivalent to $S$ 0-connected - one should make this explicit all the same:

\textbf{Lemma 2.9.} \textit{Let $\X$ be a topos, $S\subseteq\X^\wedge$ a full generating subcategory consisting of objects} that are 0-connected. \textit{Then $\X$ is $0$-connected if and only if $S$ is - more generally, the ordered set of both open and closed subobjects of $S$ in} Cat \textit{(or else of the
\printedpage{16}
% head: A. Grothendieck : Dérivateurs, Ch. II, Édition des Universités de Montpellier II et Paris VII
% foot: Version 8 juillet 2001
open subobjects $=$ sieves of $S$ that are also closed, i.e. cosieves) is isomorphic to that of the direct summands of $\X$.}

I used it tacitly in the proof of 2.8, for the induced topos $\X/W$. I do not spell out the proof.

Thus (returning to the proof of d) $\Longrightarrow$ a) in 2.6.) $S$ being 0-connected and locally cofiltering,
\DERmark{ENV0088}


\msp{page 26}\texttt{[it]} is in fact cofiltering. So for two objects $a$, $b$ of $S$, $\exists\ c$ with $c\longrightarrow a$, $c\longrightarrow b$ (condition PS 1$'$, stronger than PS 1). To prove that $\X$ is totally 0-connected, we are reduced to proving by 2.3 that the $a\times b$ in $\X^\wedge$ are 0-connected. By 2.9, it amounts to the same to prove that they are 0-connected in $S^\wedge$. So we are reduced to proving the

\textbf{Lemma 2.10.} \textit{Let $S$ be a cofiltering category. If $a$, $b$ are two objects of $S$, $a\times b$ is an object that is} $0$\textit{-connected in $S^\wedge$.}

We see at once that it suffices to prove (and in fact it amounts to the same, taking into account that $a\times b$ is locally irreducible) that if we have two arrows
\[
\begin{tikzcd}
x \arrow[r] & a\times b \\
 & y, \arrow[u]
\end{tikzcd}
\]
then $x\times_{a\times b}y\neq\emptyset$, i.e. $\exists\ z$ in $S$ making commutative
\[
\begin{tikzcd}
x \arrow[r] & a\times b \\
z \arrow[u] \arrow[r] & y. \arrow[u]
\end{tikzcd}
\]
Now by PS 1$'$ there does indeed exist a $z$ mapping to $x$ and to $y$, but the two composite arrows
\[
z\rightrightarrows a\times b,
\]
i.e. each of the double arrows
\[
z\rightrightarrows a \tag*{$(*)$}
\]
\[
z\rightrightarrows b \tag*{$(*')$}
\]
may not be equal. But by PS 2,

\msp{page 27}we can find
\[
z'\longrightarrow z
\]
\printedpage{17}
% head: A. Grothendieck : Dérivateurs, Ch. II, Édition des Universités de Montpellier II et Paris VII
% foot: Version 8 juillet 2001
equalizing the double arrow $(*)$, so we may suppose that this one is already trivial. Applying PS 2 once again, by a $z''\longrightarrow z'$, we similarly equalize the double arrow $(*')$, and we win. We have therefore proved, finally,
\[
\begin{tikzcd}
\mathrm{a} \arrow[r, Leftrightarrow] \arrow[d, Rightarrow] & \mathrm{a}' \arrow[r, Rightarrow] & \mathrm{b} \arrow[r, Leftrightarrow] & \mathrm{d}\ (=\mathrm{PS}\ 1+\mathrm{PS}\ 2) \\
\mathrm{c} \arrow[r, Leftrightarrow] & \mathrm{c}',
\end{tikzcd}
\]
but not the reverse implication. In other words, if the induced topoi $\X/a$, for $a$ in $S$, are totally 0-connected, I have not proved that the same holds for the $\X/W$, $W$ a 0-connected object \textit{whatsoever} of $\X^\wedge$. Thus, we are at the following point

\textbf{Perplexity 2.11.} \textit{Let $X$ be in} Cat. \textit{If $X$ is locally left filtering, then the fiber products $a\times_x b$ of objects in $X$ are objects that are} $0$\textit{-connected in $X^\wedge$, and this condition in turn implies PS 1, i.e. that $X$ is locally irreducible. But in the presence of this hypothesis PS 1, or $X$ locally irreducible, the condition on the fiber products $a\times_x b$ is equivalent to the following: for every diagram} \texttt{[in]} \textit{solid lines}

\msp{page 28}
\[
\begin{tikzcd}
c \arrow[r] & a\times_x b \\
e \arrow[u, dotted] \arrow[r, dotted] & d, \arrow[u]
\end{tikzcd}
\qquad\qquad
\textit{with $x$, $a$, $b$, $c$, $d$ in $X$,}
\]
\textit{we can complete it to a diagram that is} commutative, \textit{with $e$ in $X$. In other words, every commutative diagram in solid lines below can be completed by adding to it an ``initial object'' $e$} % [sic] d'autre termes -> d'autres termes, missing plural s
\[
\begin{tikzcd}[column sep=small, row sep=small]
 & & a \arrow[ddr] & \\
 & c \arrow[ur] \arrow[dddr] & & \\
e \arrow[ur, dotted] \arrow[dr, dotted] & & & x \\
 & d \arrow[uuur] \arrow[dr] & & \\
 & & b \arrow[uur] &
\end{tikzcd}
\]
(\footnote[33]{It suffices to require for $e$ the commutativity of the two squares $(e\ c\ b\ d)$ and $(e\ d\ a\ c)$.}).

This condition is thus implied by $X$ locally cofiltering, and from what precedes it is also satisfied more generally if and only if (\footnote[34]{if and only if?}) the $X/x$ are cofiltering (\footnote[35]{It means that $X$ satisfies PS 1, and the condition

\textbf{PS$'$ 2} Every commutative diagram
\[
a\rightrightarrows b\longrightarrow x
\]
% footnote continued on printed page 18
fits into a commutative diagram
\[
y\longrightarrow a\rightrightarrows b\longrightarrow x.
\]
}). But is the converse true - does this condition of total 0-connectedness of the $X/x$ imply that $X$ is totally 0-connected? But I have not constructed an example where the $X/x$
\printedpage{18}
% head: A. Grothendieck : Dérivateurs, Ch. II, Édition des Universités de Montpellier II et Paris VII
% foot: Version 8 juillet 2001
are cofiltering without $X$ being locally cofiltering (or what amounts to the same, of $X$ 0-connected with the $X/x$ cofiltering without $X$ thereby being cofiltering)(\footnote[36]{But there is an obvious \textit{counterexample}, if $X$ is a 0-connected groupoid, the $X/x$ are $\approx e$, hence cofiltering, but $X$ is cofiltering (i.e. \texttt{[illegible word(s)]}) only if $X$ is 1-connected. E.g. if $G$ is a group $\neq 1$, then $X=B_G$ is such that the $X/x$ are cofiltering, but $X$ not cofiltering. % [sic] si $C$ est un groupoïde -> si $G$ est un groupoïde, the groupoid is then written $B_G$
}) ; nor an example where the $X/x$ are totally 0-connected, without the $X/x$ being cofiltering. But let us note the
\DERmark{ENV0089}


\textbf{Corollary 2.12:} \textit{The equivalent conditions a), a$'$), b), d) (excluding c), c$'$), apparently weaker) of theorem 2.6}

\msp{page 29}\textit{are also equivalent to the following}
\begin{itemize}[leftmargin=2.5em]
\item[\textit{e)}] \textit{$\forall\ x$ in $S$, the induced topos $\X/x$ and the induced topos $\X/x\times x$ is irreducible} (\footnote[37]{NB $\X/x$ unnecessary.}),
\end{itemize}
\textit{or else to the following} % [sic] où encore -> ou encore, conjunction "ou" without accent
\begin{itemize}[leftmargin=2.5em]
\item[\textit{e$'$)}] \textit{$\forall\ x$ in $S$, $S^\wedge/x$ and $S^\wedge/x{\times}x$ are irreducible, i.e. $S/x$ and $S/x{\times}x$ are irreducible} (\footnote[38]{$S/x$ unnecessary, for $x\xrightarrow{\mathrm{diag}}x\times x$ is a nonempty open immersion, so if $x\times x$ is irreducible, $x$ is too. But it amounts to the same to say that the $S/x$ are irreducible (condition PS$^*$ 1, cf. 2.8, p. 21) \textit{and} that the $x\xrightarrow{\mathrm{diag}}x\times x$ are dominant (which indeed implies that $x\times x$ is irreducible). This immediately gives back the form PS$^*$ 2 of the cofiltering condition. % [sic] ce que implique -> ce qui implique, relative pronoun subject
}).
\end{itemize}
This is the apparently weakest ``topological'' form that I can find for the condition $S$ locally cofiltering, i.e. $\X$ (or Top$(S)$, at will) locally totally 0-connected.

\textbf{Definition 2.13.} We say that a topos $\X$ is \textit{locally totally} $0$\textit{-connected} if it satisfies the equivalent conditions a, a$'$, b, d of theorem 2.6, i.e. if $\X$ is locally connected and if in addition every 0-connected induced topos is totally 0-connected - the fiber product of two 0-connected objects over a third is 0-connected. It amounts to the same to say that every (or that one) full generating subcategory $S$ of $\X^\wedge$ consisting of 0-connected objects is locally cofiltering (still assuming $\X$ locally connected, or what amounts to the same, the existence of an $S$ consisting of 0-connected objects).

We have not yet finished the proof of 2.6 - it remains to prove that if $\X$ (assumed locally totally 0-connected) is 0-connected, it is $\infty$-connected - or else that the connected components of $\X$, assumed locally totally 0-connected, are $\infty$-connected, i.e. $W_\infty$-aspheric.

\msp{page 30}\textbf{Corollary 2.14.} \textit{Let us note all the same that if $\X$ is of the form Top$(T)$, $T$ a topological space, then $\X$ is locally totally} $0$\textit{-connected if and only if it is locally irreducible, i.e. if and only if $T$ is locally irreducible.}


\endgroup
% END UNIT Der-ChII-234/Der-ChII-234_p016-020.tex
% BEGIN UNIT Der-ChII-234/Der-ChII-234_p021-025.tex
\begingroup\makeatletter%
\newcommand{\printedpage}[1]{\marginpar{\footnotesize p.~#1}}%
\newcommand{\msp}[1]{\marginpar{\footnotesize [#1]}}%
\newcommand{\Hom}{\operatorname{Hom}}%
\newcommand{\Ker}{\operatorname{Ker}}%
\newcommand{\Ob}{\operatorname{Ob}}%
\makeatother\let\printedpage\derpp\let\msp\dermsp%
\let\mspin\DERmspin%


\printedpage{19}
% head: A. Grothendieck : Dérivateurs, Ch. II, Édition des Universités de Montpellier II et Paris VII
% foot: Version 8 juillet 2001

We are reduced to proving that this last condition implies condition d) of 2.6, and we can take for $S$ the ordered set of nonempty connected open sets of $T$, this then becomes obvious.
\DERmark{ENV0090}


That said, we must remember that the cohomological theory of irreducible topological spaces is by no means trivial - as \textsc{Serre}'s FAC shows us\,! On the other hand, it was \textsc{Serre} who had pointed out to me that for coefficients that are \textit{constant,} the cohomology of an irreducible space is trivial.

In the present case, let us already prove, if $\mathcal{X}$ is 0-connected and locally totally 0-connected, that every covering of the final object of $\mathcal{X}$ is trivial. It suffices to see that if $U$ is 0-connected, $U \longrightarrow e$ is an isomorphism. But let us consider $U \hookrightarrow U \times U$, we know that it is a direct summand, and as $U \times U$ is connected, it is an isomorphism. So $U \longrightarrow e$ \texttt{[is a]} monomorphism, hence a direct summand, and as $U \neq \emptyset$, we have $U \longrightarrow e$ an isomorphism.

It remains to prove that for every constant abelian sheaf \msp{page 31} $F$ on $\mathcal{X}$, we have $H^i(\mathcal X,F)=0$ for $i>0$. We proceed by induction on $i$, it is already done for $i=1$. Suppose this has been proved for the $j<i$ ($i$ fixed) for every topos that is \textit{locally totally} 0\textit{-connected.} Let $\xi \in H^i(\mathcal{X},F)$, let $U \longrightarrow e$ be \texttt{[an]} epimorphism such that $\xi|U=0$, and consider the spectral sequence % [sic] H^i(X,F) -> H^i(\mathcal{X},F), plain X printed where the topos \mathcal{X} is meant elsewhere
\DERmark{ENV0091}\DERmark{ENV0092}

\[
E_2^{pq} = H^p(C(U^*,\mathcal{H}^q(F))) \Longrightarrow H^*(\mathcal{X},F).
\]
We have by the totally 0-connected hypothesis
\[
\pi_0(U^q) \simeq \pi_0(U)^q, \tag*{$(*)$}
\]
and the connected components of $U^q$ are totally 0-aspheric,
\begin{center}
\begin{tikzpicture}[scale=0.8]
\draw[->] (0,0) -- (0,5);
\node[left] at (0,5) {$q$};
\draw[->] (0,0) -- (5,0);
\node at (4.3,-0.6) {$p$};
\node[left] at (0,4) {$i$};
\foreach \x in {0,1,2,3,4} {\fill (\x,4) circle (2pt); \fill (\x,0) circle (2pt);}
\end{tikzpicture}
\end{center}
so by the induction hypothesis
\[
E_1^{pq} = 0 \qquad \text{if } 0<q<i,
\]
so we find an exact sequence
\[
0 \longrightarrow H^i(C(U^*,F)) \longrightarrow H^i(\mathcal{X},F) \longrightarrow H^0(\ \underbrace{H^i(U^*,F)}_{\substack{\text{complex}\\ \text{co-semi-simplicial}\\ \text{for variable \v{C}ech degree}\\ \text{at fixed }i}}\ ) = \Ker(H^i(U,F) \rightrightarrows H^i(U^2,F)).
\]
\DERmark{ENV0093}


\printedpage{20}
% head: A. Grothendieck : Dérivateurs, Ch. II, Édition des Universités de Montpellier II et Paris VII
% foot: Version 8 juillet 2001

By the hypothesis on $U$, the image of $\xi$ in the third term is zero, so $\xi$ belongs to the first term. But for $(*)$, the complex $C(U^*,F)$ is the \textit{trivial} complex of the powers of $I=\pi_0(U)$, with coefficients in $M$ (if $F=M_{\mathcal{X}}$). Its cohomology \msp{page 32} in degrees $>0$ is zero, so $\xi=0$, q.e.d.

\noindent\textbf{Proposition 2.15.} (\footnote[39]{This is already in SGA 4, I 8.}) \textit{Let $X$ be an object of} Cat\textit{. Equivalent conditions:}
\begin{enumerate}[label={}, leftmargin=2em, itemsep=0pt]
\item \textit{a) $X$ is filtering, i.e. $X$ is $0$-connected and $X^o$ satisfies the equivalent conditions of theorem 2.6 ($X^{o,\wedge}$ is locally totally $0$-connected).}
\item \textit{b) The final object of $X^\wedge$ is ind-representable.}
\end{enumerate}
\noindent\textsc{Proof}. We know that we have
\[
e_{X^\wedge} \simeq \varinjlim_{x \in X} x, \tag*{(2.15.1)}
\]
direct limit taken in $X^\wedge$. Indeed, this also means that for every object $a$ in $X$, we have
\[
\varinjlim_{x \in X} \ \underbrace{x(a)}_{\Hom_X(a,x)} \ \xrightarrow{\ \sim\ } e \qquad \text{(one-point set)}
\]
But let us consider the object $X' = a\backslash X \longrightarrow X$ coétale over $X$, which is therefore Cat-cofibered over $X$, we know that we have
\[
\pi_0(X') \simeq \varinjlim_{x \in X} \pi_0(X'_x),
\]
where $X'_x = \Hom(a,x)$, so we find
\[
\varinjlim_{x \in X} \Hom(a,x) = \pi_0(a\backslash X) = e
\]
since $a\backslash X$ is 0-connected.

Formula (2.15.1) shows that if \msp{page 33} $X$ is filtering, $e_{X^\wedge}$ is a filtering direct limit of representable functors. We know moreover that we can reindex by a filtering \textit{ordered} category (cf. SGA 4, I 8.1.6 (this is due to P. \textsc{Deligne}\,!)).

We have proved a $\Longrightarrow$ b, let us prove b $\Longrightarrow$ a. So suppose
\[
e_{X^\wedge} = \varinjlim_{i \in I} x_i,
\]
where $I$ is a (small) filtering ordered set, and $I \longrightarrow X$ a functor. Let us prove PS 1$'$ (right-hand form), i.e. for $a, b$ in $X$, \texttt{[there]} exists $\vcenter{\hbox{\begin{tikzcd}[row sep=0.1em, column sep=1.2em]
a \arrow[dr] & \\
& x \\
b \arrow[ur] &
\end{tikzcd}}}$. Indeed, since $\varinjlim\limits_{I} \Hom(a,x_i) = e$, it follows that $\Hom(a,x_i) \neq \emptyset$ for $i$ large, likewise $\Hom(b,x_i) \neq \emptyset$, so it suffices to take $\vcenter{\hbox{\begin{tikzcd}[row sep=0.1em, column sep=1.2em]
a \arrow[dr] & \\
& x_i \\
b \arrow[ur] &
\end{tikzcd}}}$. Let us prove PS 2$_{\mathrm{d}}$, i.e. that for \texttt{[a]} double arrow $a \overset{u}{\underset{v}{\rightrightarrows}} b$, \texttt{[there]} exists $b \longrightarrow x$ which equalizes.

\printedpage{21}
% head: A. Grothendieck : Dérivateurs, Ch. II, Édition des Universités de Montpellier II et Paris VII
% foot: Version 8 juillet 2001

\msp{page 34}
Let $\xi_{i_0}$ be in $\Hom(b,x_{i_0})$, its two images in $\Hom(a,x_{i_0})$ are not necessarily equal, but become so in the direct limit, reduced to $e$, so become equal by composing with $x_{i_0} \longrightarrow x_i$ for $i$ large enough, so the composite $b \longrightarrow x_{i_0} \longrightarrow x_i$
\[
a \rightrightarrows b \longrightarrow x_{i_0} \longrightarrow x_i
\]
equalizes the double arrow.

\noindent\textbf{Corollary 2.16.} \textit{Let $W$ be a basic localizer satisfying the limit axiom} Loc (8)\textit{. If $X$ in} Cat \textit{is a filtering or cofiltering category, then $X$ is $W$-aspheric.}

More generally and pleasantly:

\noindent\textbf{Corollary 2.17.} \textit{Let $W$ be as in 2.16. Then every filtering category is} totally $W$-aspheric (\footnote[40]{cf. XII, §V for this notion, which I have struggled a great deal to bring out and to understand\,!}), \textit{i.e. for every arrow $f : X \longrightarrow Y$ in} Cat\textit{, the object $f_!(e_X)$ of $Y^\wedge$ (also defined as $y \longmapsto \pi_0(y\backslash X)$, and occurring in the canonical factorization of $f$ as $ip$, with $i$ étale and the cofibers} \texttt{[?]} \textit{of $p$ $0$-connected)} \msp{page 35} \textit{is $W$-aspheric. In fact, $Y/f_!(e_{X^\wedge})$ is even a filtering category.}
\DERmark{ENV0094}\DERmark{ENV0095}


Indeed, the functor $f_!$ commutes with direct limits, so $f_! e_{X^\wedge}$ is a filtering direct limit of the $f_!(x_i) = f(x_i)$. So $Y/f_!(e_{X^\wedge})$ is a filtering direct limit of the $Y/f(x_i)$. So it is even a filtering category, a fortiori it is $W$-aspheric.

So finally, we find nontrivial examples of totally $W$-aspheric categories, i.e. examples of categories without a final object. % [sic] donc -> dont (?), the conjunction "donc" makes no sense here
\DERmark{ENV0096}


\noindent\textbf{Corollary 2.18.} (\footnote[41]{already in SGA 4, I 8.}) \textit{Let $X$ be in} Cat\textit{, $F$ in $X^\wedge$. For $F$ to be ind-representable it is necessary and sufficient that $X/F$ be filtering.}

Indeed, by (2.15), $X/F$ is filtering if and only if the final object of $(X/F)^\wedge$, corresponding to $F$ in $X^\wedge/F$, comes from an ind-object of $X/F$, but this also means that $F$ is the direct limit in $X^\wedge$ of an ind-object of $X$, i.e. that $F$ is ind-representable.
\DERmark{ENV0097}


\msp{page 36}
\noindent\textbf{Question.} \textit{If $X$ in} Cat \textit{is totally $W_\infty$-aspheric, i.e. the $f_!(e_{X^\wedge})$ are $W_\infty$-aspheric for every arrow $f : X \longrightarrow Y$ in} Cat\textit{, or again for every such arrow, with the $y\backslash X$ $0$-connected ($y \in \Ob Y$), $Y$ is $W_\infty$-aspheric - does this imply that $X$ is filtering?} That wouldn't surprise me - and even assuming only the $f_!(e_{X^\wedge})$ 1-connected (\footnote[42]{They always are (cf. XII, p. 101), it is 2-\textit{connected} that should be read here\,!}) (they are 0-connected, as soon as $X$ is). But I do not want to linger at present on this question.
\DERmark{ENV0098}


The answer to the preceding question is negative, as we see from the:

\printedpage{22}
% head: A. Grothendieck : Dérivateurs, Ch. II, Édition des Universités de Montpellier II et Paris VII
% foot: Version 8 juillet 2001

\noindent\textbf{Proposition 2.19.} (\footnote[43]{\textbf{False.}}) \textit{Let $X$ be in} Cat \textit{such that $\mathit{diag}_X : X \longrightarrow X \times X$ is cofinal} (\footnote[44]{cf. definition 3.3, p. 45.})\textit{, i.e. such that $X^o$ is totally $0$-connected. Then $X$ is totally $W_\infty$-aspheric (hence also totally $W$-aspheric for every basic localizer $W$ such that $W \supseteq W_\infty$).}

Indeed, let $f : X \longrightarrow Y$ be a cofinal functor, we must prove that $Y$ is $W_\infty$-aspheric. Now we know that totally 0-connected implies $W_\infty$-aspheric (\footnote[45]{This is \textbf{false.}}) (cf. proposition 2.3, p. 17-18), so it suffices to see that $Y^o$ is totally 0-connected, i.e. $Y \xrightarrow{d_Y} Y \times Y$ cofinal. Now we have the commutative diagram
\[
\begin{tikzcd}[row sep=large, column sep=large]
X \arrow[r, "d_X"] \arrow[d, "f"'] & X \times X \arrow[d, "f \times f"] \\
Y \arrow[r, "d_Y"] & Y \times Y,
\end{tikzcd}
\]
where by hypothesis $f$, hence also $f \times f$, and $d_X$ are cofinal. So $\varphi = (f \times f) \circ d_X$ \texttt{[rather} $\varphi = (f \times f) \circ d_X$\texttt{]} is cofinal, and since $\varphi = d_Y f$ and $f$ is also cofinal, we conclude that $d_Y$ is, q.e.d.
\DERmark{ENV0099}


\noindent\textbf{Corollary.} $\boldsymbol{\Delta}$ \textit{being totally $0$-connected, $\boldsymbol{\Delta}^o$ is totally $W_\infty$-aspheric, yet it is not filtering.}

\noindent\textbf{Question.} \textit{$X$ totally $W_\infty$-aspheric $\Longrightarrow$ $X^o$ totally $0$-connected? Or at least implies PS 1$'$?}

\msp{page 36bis}
\noindent\textbf{Counterexample 2.20} to the false proposition 2.19.

Let $X$ be an object of Cat, let us adjoin to it the final category $e$ as a closed subobject, i.e. let us consider an object $X'$ of Cat in which $X$ is an open subobject, $e$ the complementary closed subobject:
\[
X \overset{i}{\hookrightarrow} X' \overset{j}{\hookleftarrow} e.
\]
We know (cf. XII) that such an operation is determined, up to isomorphism, by the datum of an object
\[
H \in \Ob X^\wedge
\]
by putting
\[
H(x) = \Hom_{X'}(x,e),
\]
the category $X/e$ being identified (up to isomorphism) with $X/H$. It is therefore 0-connected if and only if the object $H$ of $X^\wedge$ is 0-connected. This also means that the morphism $i$ is co-cofinal, i.e. that $i^o : X^o \longrightarrow X'^o$ is cofinal. Suppose now $X$ $W_\infty$-aspheric and totally 0-connected, i.e. the current Loc (6), cf. XVI. Proposition 2.19 asserts that, under these conditions, $X'$ (which is itself also totally 0-connected) is even $W$-aspheric. It would therefore follow that $i \in W_\infty$. But by XII, th. 4, cor. 4 (p. 149), it would follow that
\printedpage{23}
% head: A. Grothendieck : Dérivateurs, Ch. II, Édition des Universités de Montpellier II et Paris VII
% foot: Version 8 juillet 2001
\msp{page 36ter}
$X/e$ is $W_\infty$-aspheric, i.e. $H$ $W_\infty$-aspheric. Thus, if $X$ is totally 0-aspheric, every 0-connected object of $X^\wedge$ would be $\infty$-connected. But this is manifestly false, e.g. if we take for $X$ a test category for $W_\infty$, e.g. $X = \boldsymbol{\Delta}$.

This makes me suspect that the conjecture
\[
X \text{ totally } W_\infty\text{-aspheric} \Longrightarrow X \text{ filtering}
\]
might be valid.


\endgroup
% END UNIT Der-ChII-234/Der-ChII-234_p021-025.tex
% BEGIN UNIT Der-ChII-234/Der-ChII-234_p026-030.tex
\begingroup\makeatletter%
\newcommand{\printedpage}[1]{\marginpar{\footnotesize p.~#1}}%
\newcommand{\msp}[1]{\marginpar{\footnotesize [#1]}}%
\newcommand{\mspin}[1]{\texttt{[#1]}}%
\newcommand{\HOT}{\mathrm{HOT}}%
\newcommand{\Cat}{\mathrm{Cat}}%
\newcommand{\Fl}{\mathrm{Fl}}%
\newcommand{\Ob}{\mathrm{Ob}}%
\newcommand{\Ens}{\mathrm{Ens}}%
\newcommand{\Hom}{\underline{\mathrm{Hom}}}%
\newcommand{\defeq}{\overset{\mathrm{def}}{=}}%
\newcommand{\pl}{\texttt{[rather}}%
\newcommand{\uHom}{\underline{\operatorname{Hom}}}%
\makeatother\let\printedpage\derpp\let\msp\dermsp%
\let\mspin\DERmspin%

\DERmark{ENV0100}


\printedpage{24}
% head: A. Grothendieck : Dérivateurs, Ch. II, Édition des Universités de Montpellier II et Paris VII
% foot: Version 8 juillet 2001
\msp{page 37}

\section*{3\quad Return to cofinality}

It is instructive to relate questions of cofinality (for an arrow $u : I' \longrightarrow I$ in Cat) to the derivator $\HOT_{W_0}$, where $W_0 \subseteq \Fl\,\Cat$ is the coarsest of the localizers in Cat satisfying the connectedness axiom Loc 5, i.e.
\[
W_0 = \{ f \in \Fl\,\Cat \mid \pi_0(f) \text{ bijective}\}. \tag*{(3.1)}
\]
One then finds, for $I$ in Cat, a canonical equivalence, 2-functorial in $I$,
\[
\HOT_{W_0}(I) \xrightarrow{\;\sim\;} I^{\wedge} \defeq \uHom(I^o,(\Ens)), \quad \text{\pl}\ \uHom(I^o,\Ens)\,\text{\texttt{]}} \tag*{(3.2)}
\]
\DERmark{ENV0101}

in other words, one has an equivalence
\[
\HOT_{W_0} \simeq \mathbf{D}_{(\Ens)}. \tag*{(3.3)}
\]
When $\HOT_{W_0}$ is defined by localization starting from the category $\Cat(I) \defeq \uHom(I^o,\Cat)$, the arrow (3.2) is the one deduced, by localization, from the arrow
\[
\begin{array}{rcl}
\pi_0^I : \uHom(I^o,\Cat) & \longrightarrow & \uHom(I^o,(\Ens)) \\[1ex]
X & \longmapsto & \pi_0 \circ X,
\end{array} \tag*{(3.4)}
\]
where
\[
\pi_0 : \Cat \longrightarrow (\Ens) \tag*{(3.5)}
\]
is the ``connected components'' functor. If one describes $\HOT_{W_0}(I)$ as a localized category of $\Cat/I$, (3.2) is deduced from the arrow
\[
\left\{
\begin{array}{rcl}
\Cat/I & \longrightarrow & I^{\wedge} = \uHom(I^o,(\Ens)) \\[1ex]
X & \longmapsto & \underline{\pi_0}(X/I) = (i \longmapsto \pi_0(i\backslash X)).
\end{array}
\right. \tag*{(3.6)}
\]

\msp{page 38}
Thus, the familiar functors, for an arrow
\[
u : I' \longrightarrow I
\]
in Cat
\[
\left\{
\; I'^{\wedge} \quad
\begin{array}{c}
\xrightarrow{\;u_!\;} \\[1.5ex]
\xleftarrow{\;u^*\;} \\[1.5ex]
\xrightarrow{\;u_*\;}
\end{array}
\quad I^{\wedge},
\right. \tag*{(3.7)}
\]
\printedpage{25}
% head: A. Grothendieck : Dérivateurs, Ch. II, Édition des Universités de Montpellier II et Paris VII
% foot: Version 8 juillet 2001
are interpreted as the functors of the same name for the derivator $\HOT_{W_0}$, or more directly, for the derivator (which one might mistakenly regard as hoax-like) $\mathbf{D}_{(\Ens)}$.

This gives a meaning to the cohort of cohomological notions associated with a derivator, or more particularly, with a basic localizer $W$ on Cat (giving rise to the derivator $\HOT_W$). We shall designate them by the name of $W_0$-\textit{fourbi} (gear) (instead of $\HOT_{W_0}$-\textit{fourbi}, or $\mathbf{D}_{(\Ens)}$-\textit{fourbi}), in particular for an arrow $u : I' \longrightarrow I$ of Cat such as the notion
\[
\left\{
\begin{array}{lll}
\text{$W_0$-equivalence} & \Bigl(\iff & \text{$\pi_0(u)$ bijective.}\Bigr) \\
\text{$W_0$-aspheric} & \Bigl(\iff & \text{the $I'/i$ for $i \in \Ob I$ are} \\
 & & \text{$0$-connected.}\Bigr) \\
\text{$W_0$-coaspheric} & \Bigl(\iff & \text{the $i\backslash I'$ for $i \in \Ob I$ are} \\
 & & \text{$0$-connected.}\Bigr) \\
\text{$W_0$-proper (\footnotemark[46])} & \Bigl(\iff & \text{the inclusions $I'_i \longrightarrow I'/i$ ($i \in \Ob I$)} \\
 & & \text{are $W_0$-coaspheric, i.e. the fibers} \\
 & & \text{of the colocalized $i'\backslash I' \longrightarrow i\backslash I$} \\
 & & \text{($i = u(i')$) \pl\ ($i = u(i')$)]} \\
 & & \text{are $0$-connected.}\Bigr) \\
\text{\mspin{page 39}} & & \\
\text{$W_0$-smooth (\footnotemark[47])} & \Bigl(\iff & \text{the inclusions $I'_i \longrightarrow i\backslash I'$} \\
 & & \text{are $W_0$-aspheric, i.e. the fibers} \\
 & & \text{of the localized arrows $I'/i' \longrightarrow I/i$} \\
 & & \text{($i = u(i')$) are $0$-connected.}\Bigr) \\
\text{$W_0$-perfect} & \Bigl(\iff & \text{$W_0$-proper and $W_0$-smooth.}\Bigr) \\
\text{universal $W_0$-equivalence (\footnotemark[48])} & \Bigl(\iff & \text{$\pi_0$ bijective and remains so under any} \\
 & & \text{base change. These are also} \\
 & & \text{the $W_0$-fibrations (or weak $W_0$-fibrations)} \\
 & & \text{with $W_0$-aspheric fibers,} \\
 & & \text{i.e. $0$-connected.}\Bigr) \\
\text{$W_0$-fibration} & \Bigl(\iff & \text{$u^* : \Cat/I \longrightarrow \Cat/I'$ transforms} \\
 & & \text{$W_0$-equivalences into $W_0$-equivalences.}\Bigr) \\
\text{weak $W_0$-fibration (\footnotemark[49])} & \Bigl(\iff & \text{the base change} \\
 & & \text{$u^* : \Cat/I \longrightarrow \Cat/I'$ transforms} \\
 & & \text{homotopisms $J_1 \longrightarrow J_2$ into} \\
 & & \text{$W_0$-equivalences - and it suffices to} \\
 & & \text{require it for the inclusion of a} \\
 & & \text{final object (or at one's choice initial) and for} \\
 & & \text{the inclusion of a point in $\boldsymbol{\Delta}^1$.}\Bigr)
\end{array}
\right. \tag*{(3.8)}
\]
\DERmark{ENV0102}\DERmark{ENV0103}

\footnotetext[46]{Example: the Cat-cofibrations.}
\footnotetext[47]{Example: the Cat-fibrations.}
\footnotetext[48]{Example: the $W_0$-smooth or $W_0$-proper arrows, with 0-connected fibers.}
\footnotetext[49]{Example: the $W_0$-smooth arrows such that $\pi_0(I'/I) = (i \longmapsto \pi_0(i\backslash I')) = u_!(e_{I'^{\wedge}})$ is locally constant on $I$, i.e. $\underline{\pi_0}(I'/I)$ is an étale \textit{covering} of $I$, or if $u^{\mathrm{opp}}$ satisfies this condition.}

\printedpage{26}
% head: A. Grothendieck : Dérivateurs, Ch. II, Édition des Universités de Montpellier II et Paris VII
% foot: Version 8 juillet 2001
\texttt{[We have a]} \textit{strict} implication (and how)
\[
\text{universal $W_0$-equivalence} \Longrightarrow \text{$W_0$-fibration}
\]
\texttt{[and a]} \textit{strict} implication (and how\,!) (\footnote[50]{But if the fibers are 0-connected, the two notions of $W_0$-fibrations are equivalent and mean that one has a universal $W_0$-equivalence \dots})
\[
\text{$W_0$-fibration} \Longrightarrow \text{weak $W_0$-fibration.}
\]

Let us point out, as an example of the use of these notions, that for a cartesian diagram in Cat
\[
\begin{tikzcd}
I' \arrow[d, "u"'] & J' \arrow[l, "p'"'] \arrow[d, "v"] \\
I & J, \arrow[l, "p"']
\end{tikzcd} \tag*{(3.9)}
\]
the base change arrow for $u_*$, $v_*$ for presheaves is \texttt{[an]} isomorphism if $u$ is $W_0$-proper, or $p$ $W_0$-smooth, and the base change arrow for $u_!$, $v_!$ is \texttt{[an]} isomorphism if $u$ is $W_0$-smooth or $p$ $W_0$ proper (\footnote[51]{Unless I am mistaken, this gives \textit{two characterizations} of $W_0$-proper morphisms, or of $W_0$-smooth morphisms.}). I point out in this connection that these statements, which a priori are formulated for presheaves with values in (Ens), remain valid % [sic] ca -> ça (in footnote 51), missing cedilla
\msp{page 40}
for presheaves with values in any category $\mathcal{M}$ (subject to the existence of the cohomological or homological direct image operations under consideration). This can be seen e.g. by noting that % [sic] operations -> opérations, missing accent
\[
\text{$W_0$-equivalence} \;\Longrightarrow\; \text{$\mathbf{D}_{\mathcal{M}}$-equivalence} \tag*{(3.10)}
\]
for every $\mathcal{U}$-category $\mathcal{M}$, from which it follows that every (\footnote[52]{Not quite correct for \textit{every} ``$W_0$-notion'', one must except the two notions of strong, and weak $W_0$-fibration.}) ``$W_0$-notion'' implies the corresponding $\mathbf{D}_{\mathcal{M}}$-notion, and thus allows one to apply to $\mathbf{D}_{\mathcal{M}}$ the results concerning it. (There is of course a direct, less learned argument, consisting in reducing to the case where $\mathcal{M}$ is small, and then using the fully faithful embedding $\mathcal{M} \hookrightarrow \mathcal{M}^{\wedge}$ resp. $\mathcal{M} \hookrightarrow ((\mathcal{M}^o)^{\wedge})^o$.)

I recall for the record a (third\,?) characterization of the arrows $u : I' \longrightarrow I$ which are $W_0$-proper (resp. $W_0$-smooth): they are those for which the base change functor $u^* : \Cat/I \longrightarrow \Cat/I'$ transforms $W_0$-coaspheric (resp. $W_0$-aspheric) arrows into likewise. This gains in interest, from the fact that the notion of a $W_0$-coaspheric (resp. $W_0$-aspheric) functor is interpreted as that of a cofinal (resp. co-cofinal) functor, cf. below.

\printedpage{27}
% head: A. Grothendieck : Dérivateurs, Ch. II, Édition des Universités de Montpellier II et Paris VII
% foot: Version 8 juillet 2001
\msp{page 41}
The notion of $W_0$-fibration appears extremely strong - much stronger in certain respects than that (say) of $W_\infty$-fibration or of $W_\omega$-fibration. It is equivalent to this (as one sees at once): for every connected component $I'_\alpha$ of $I'$, denoting by $I_\alpha$ the corresponding connected component of $I$, and for every $J_\alpha$ 0-connected and over $I_\alpha$, the fiber product $I'_\alpha \times_{I_\alpha} J_\alpha$ is again 0-connected. (That looks like those stories of total 0-connectedness of a topos\,!) But this also means that $I'_\alpha \longrightarrow I_\alpha$ is a universal $W_0$-equivalence. This implies that its fibers are 0-connected, and the converse is true if $u : I' \longrightarrow I$ is $W_0$-proper or $W_0$-smooth. Thus, the essential condition for a $W_0$-fibration is that the fibers of the induced arrows $I'_\alpha \longrightarrow I_{\alpha}$ \pl\ $I'_\alpha \longrightarrow I_\alpha$\texttt{]} be 0-connected. In fact, one has this
\DERmark{ENV0104}


\noindent\textbf{Proposition 3.1.} \textit{Let $u : I' \longrightarrow I$ be an arrow in} Cat.
\begin{itemize}
\item[\textit{a)}] \textit{For it to be a universal $W_0$-equivalence, it is necessary and sufficient that for every $\boldsymbol{\Delta}^1 \longrightarrow I$, i.e. every arrow $\alpha$ of $I$, the category $I' \times_I \boldsymbol{\Delta}^1 \defeq I'|\alpha$ be $0$-connected, or again that the fibers of $u$ be $0$-connected, and that for every} nonidentity \textit{arrow $\alpha$ of $I$, $I'|\alpha$ be $0$-connected.}
\item[\textit{b)}] \textit{For $u : I' \longrightarrow I$ to be a $W_0$-fibration, it is necessary and sufficient that for every connected component $I'_{\alpha'}$ of}
\msp{page 42}
\textit{$I'$, putting $\alpha = u(\alpha')$, the induced arrow $I'_{\alpha'} \longrightarrow I_\alpha$ be a universal $W_0$-equivalence (cf. criterion a) for this notion).}
\DERmark{ENV0105}

\item[\textit{c)}] \textit{For it to be a weak $W_0$-fibration, it is} necessary \textit{that $\forall\, i \in I$, the inclusion $I'_i \hookrightarrow i\backslash I'$ be a $W_0$-equivalence, i.e. $\pi_0(I'_i) \xrightarrow{\;\sim\;} \pi_0(i\backslash I')$ be bijective,} and \textit{that $i \longmapsto \pi_0(i\backslash I')$ be a locally constant presheaf on $I$ (i.e. $\underline{\pi_0}(I'/I)$ be an étale} covering \textit{of $I$). This condition is also sufficient, if $u$ is $W_0$-smooth. Moreover (for the record) $u$ is a weak $W_0$-fibration if and only if $u^{\mathrm{opp}}$ is.}
\end{itemize}

The greatest difference between the notion of $W_0$-fibration, and that of weak $W_0$-fibration, seems to me that $W_0$-perfect morphisms (i.e. $W_0$-smooth and $W_0$-proper) are necessarily weak $W_0$-fibrations, but are not necessarily $W_0$-fibrations. The $W$-perfect morphisms, for a basic localizer $W \subseteq W_0$, are weak $W$-fibrations. The most striking example in this respect is probably that of étale coverings, which are $W_\omega$-perfect (hence
\DERmark{ENV0106}

\msp{page 43}
weak $W_\omega$-fibrations and a fortiori weak $W_0$-fibrations), but not $W_0$-fibrations\,; for in the case where $I'$, $I$ are 0-connected, the fibers would have to be 0-connected (cf. b) above), i.e. $u$ would have to be an isomorphism\,!

A nice property of $W_0$-smooth (resp. $W_0$-proper) morphisms, is that they are \textit{open} (resp. closed) maps for the underlying topological spaces, which indeed means that the local fibers, i.e. of the arrows localized above $I'/i' \longrightarrow I/i$ (resp. colocal, i.e. the arrows colocalized below $i'\backslash I' \longrightarrow i\backslash I$), are surjective, i.e. with nonempty fibers (which is implied by the condition that these be 0-connected).

\printedpage{28}
% head: A. Grothendieck : Dérivateurs, Ch. II, Édition des Universités de Montpellier II et Paris VII
% foot: Version 8 juillet 2001
In all this I have still hardly spoken of $W_0$-aspheric resp. coaspheric arrows, which nevertheless are here a main motivation for introducing $W_0$\,!

\noindent\textbf{Proposition 3.2.} \textit{Equivalent conditions on the arrow $u : I' \longrightarrow I$ in} Cat:
\begin{itemize}
\item[a)] \textit{For every $\mathcal{U}$-category $\mathcal{M}$ and every functor $F : I \longrightarrow \mathcal{M}$, the canonical arrow (provided that small direct limits are representable in $\mathcal{M}$),}
\[
\varinjlim_{I'} u^*(F) \longrightarrow \varinjlim_{I} F \tag*{(3.2.1)}
\]
\msp{page 44}
\textit{is} \texttt{[an]} \textit{isomorphism.}
\item[a$'$)] \textit{For every $\mathcal{U}$-category $\mathcal{M}$ in which small inverse limits are representable and every $F \in \mathcal{M}(I) = \uHom(I^o,\mathcal{M})$, i.e. every presheaf on $I$ with values in $\mathcal{M}$, the canonical arrow}
\[
\varprojlim_{I} F \longrightarrow \varprojlim_{I'} u^*(F) \tag*{(3.2.2)}
\]
\textit{is} \texttt{[an]} \textit{isomorphism.}
\item[a$''$)] (\footnote[53]{This is a condition on morphisms of topoi, namely that if $F$ \texttt{[is]} in $\mathcal{X}^{\wedge}$, one has
\[
H^0(\mathcal{X},F) \xrightarrow{\;\sim\;} H^0(\mathcal{X}',u^*F).
\]
There is \texttt{[a]} similar notion for the $H^i$ ($i \leq n$), to be brought out. Notion of $W_0$-coaspheric morphisms of topoi\,! It is in the nature of a relative 0-connectedness condition, dual to the $W_0$-aspheric notion, which simply means that for $U$ 0-connected in $\mathcal{X}^{\wedge}$ \pl\ $\mathcal{X}^{\wedge}$\texttt{]}, $u^*(U)$ \pl\ $u^*(U)$\texttt{]} is 0-connected. Likewise, one probably has notions (``dual'' to the preceding ones) of $W_n$-coasphericity (well understood in principle if $\mathcal{X}$ is locally $n$-aspheric) (The $W_n$-coasphericity, being that mentioned above.)}) \textit{As in a$'$, with $\mathcal{M} = (\Ens)$, i.e. $F \in \Ob I^{\wedge}$.}~(\footnote[54]{\textbf{NB} One can introduce
\begin{itemize}
\item[a$'''$)] As in a), but with $\mathcal{M} = (\Ens)$.
\item[a$^{IV}$)] As in a), with the canonical functor $I\longrightarrow I^{\wedge}$.
\end{itemize}}).
\DERmark{ENV0107}\DERmark{ENV0108}\DERmark{ENV0109}

\item[b)] \textit{The arrow $u$ is $W_0$-coaspheric, i.e. for every $i \in I$, $i\backslash I'$ is $0$-connected.}
\end{itemize}
The equivalence of a, a$'$, a$''$ is $\pm$ tautological in terms of the definitions (cf. SGA 4, I, 8\dots). The form a$''$ reads, putting $\mathbf{D} = \HOT_{W_0}$, that $\forall\, F \in \mathbf{D}(I)$, one wants
\[
\underbrace{H^{\bullet}_{\mathbf{D}}(I,F)}_{\defeq p_{I,*}(F)} \longrightarrow \underbrace{H^{\bullet}_{\mathbf{D}}(I',u^*F)}_{\defeq p_{I',*}(u^*F)} \qquad \text{\pl}\quad \underbrace{H^{\bullet}_{\mathbf{D}}(I',u^*F)}_{\defeq p_{I',*}(u^*F)}\ \text{\texttt{]}} \tag*{(3.2.3)}
\]
\DERmark{ENV0110}

($p_I : I \longrightarrow e$ and $p_{I'} : I' \longrightarrow e$ the canonical functors) to be \texttt{[an]} isomorphism. One knows that this is equivalent to $u$ $\mathbf{D}$-coaspheric. The proof is an amusing game between $u_!$, $u^*$, $u_*$ all at once. One must express that the Hom from an object $M$ of $\mathbf{D}(e)$ ($=$ (Ens) in this case) into the arrow (3.2.3) is bijective. Here, it % [sic] occurence -> occurrence, French spelling


\endgroup
% END UNIT Der-ChII-234/Der-ChII-234_p026-030.tex
% BEGIN UNIT Der-ChII-234/Der-ChII-234_p031-035.tex
\begingroup\makeatletter%
\newcommand{\printedpage}[1]{\marginpar{\footnotesize p.~#1}}%
\newcommand{\msp}[1]{\marginpar{\footnotesize [#1]}}%
\newcommand{\Hom}{\operatorname{Hom}}%
\newcommand{\uHom}{\underline{\operatorname{Hom}}}%
\newcommand{\Ob}{\operatorname{Ob}}%
\newcommand{\bs}{\backslash}%
\newcommand{\cM}{\mathcal{M}}%
\newcommand{\dbl}[2]{\mathrel{\overset{#1}{\underset{#2}{\rightrightarrows}}}}%
\newcommand{\defeq}{\overset{\mathrm{def}}{=}}%
\makeatother\let\printedpage\derpp\let\msp\dermsp%
\let\mspin\DERmspin%

\DERmark{ENV0111}


\printedpage{29}
% head: A. Grothendieck : Dérivateurs, Ch. II, Édition des Universités de Montpellier II et Paris VII
% foot: Version 8 juillet 2001

\msp{page 45}
is more natural to take $M = e$ (one-point set) and to interpret the arrow (3.2.3) as
\[
\Hom(e_I, F) \longrightarrow \Hom(e_{I'}, u^*F)
\tag*{(3.2.4)}
\]
and by adjunction of $u_!$, $u^*$ the second member is interpreted as $\Hom(u_!(e_{I'}), F)$, and the arrow is deduced from
\[
u_!e_{I'} = u_!u^*(e_I) \xrightarrow{\text{adj.}} e_I
\tag*{(3.2.5)}
\]
by taking Hom into $F$. Thus to say that (3.2.4) is \texttt{[an]} isomorphism for every $F$ means that (3.2.5) is \texttt{[an]} isomorphism. Passing to the values of the two members at $i$, we find $\pi_0(i\bs I') \xrightarrow{\sim} 1$, i.e. that $i\bs I'$ is 0-connected, q.e.d.

\noindent\textbf{Definition 3.3.}
\begin{itemize}[label={}]
\item[a)] When the equivalent conditions of 3.2 are satisfied, thus when the $i\bs I'$ ($i \in \Ob I$) are 0-connected, we say that $u : I' \longrightarrow I$ is a \textit{cofinal} functor. We say that $u$ is \textit{co-cofinal} if $u^o : I'^o \longrightarrow I^o$ is cofinal, i.e. if the localized categories $I'/i$ of $I'$ over $I$ ($i \in \Ob I$) are 0-connected.
\item[b)] If $I'$ is a \textit{subcategory} of $I$, we say that $I'$ \textit{is cofinal in} $I$ if the inclusion functor is cofinal, \textit{co-cofinal} if this functor is co-cofinal.
\end{itemize}

\noindent\textbf{Corollary 3.4.} (\footnote[55]{\textbf{NB} For a complementary cofinality statement, cf. XVIII prop. 4.9.6.3 (p. 224-225, 228-242). A statement covering both would be needed.})
\begin{itemize}[label={}]
\item[\textit{a)}] \textit{For $u$ to be cofinal, it is} necessary \textit{that the colocalized categories $i\bs I'$ be nonempty (condition F 1); F 0 is also required. Thus we must have}

\msp{page 46}
\begin{itemize}[label={}]
\item[\textit{F 0)}] $\pi_0(u)$ \texttt{[rather $\pi_0(u)$]} \textit{is bijective.}
\item[\textit{F 1)}] \textit{For every} $i \in \Ob I$, \texttt{[there]} \textit{exists} \texttt{[an]} $i' \in I'$ \textit{and an arrow} $i \longrightarrow u(i')$ \textit{in} $I$.
\end{itemize}
\DERmark{ENV0112}\DERmark{ENV0113}

(\footnote[56]{\textbf{NB} F 1 already implies \texttt{[the]} surjectivity of $\pi_0$, hence its bijectivity if $I'$ is 0-connected - in this case F 0 follows from F 1.}).
\item[\textit{b)}] \textit{If $I'$ is weakly filtering, for $u$ to be cofinal, it is} sufficient \textit{that it satisfy F 1, as well as the condition:}
\begin{itemize}[label={}]
\item[\textit{F 2)}] \textit{For every double arrow} $i \dbl{\beta}{\gamma} u(i')$ \textit{(with} $i \in \Ob I$, $i' \in \Ob I'$\textit{),} \texttt{[there]} \textit{exists an arrow} $i' \xrightarrow{\alpha'} j'$ \textit{in} $I'$\textit{, in such a way that} $u(\alpha')$ \textit{equalizes} $(\beta, \gamma)$\textit{, i.e.} $u(\alpha')\beta = u(\alpha')\gamma$.
\end{itemize}
\textit{and these conditions F 1, F 2 imply that the $i\bs I'$ are weakly filtering, and $I$ as well} (\footnote[57]{If $I'$ satisfies right PS 1, then for the cofinality of $u$, it suffices that we have F 0, F 1, F 2, and then the $i\bs I'$ are weakly filtering, and $I$ is prefiltering just like $I'$.}).
\DERmark{ENV0114}


\printedpage{30}
% head: A. Grothendieck : Dérivateurs, Ch. II, Édition des Universités de Montpellier II et Paris VII
% foot: Version 8 juillet 2001
\item[\textit{c)}] \textit{If $I'$ is filtering, then $u$ is cofinal} if and only if \textit{it satisfies the two conditions F1, F2, or again if and only if the $i\bs I'$ are} filtering \textit{(and not merely} 0-\textit{connected). And this implies that $I$ is also filtering} (\footnote[58]{If $I'$ is locally filtering, $u$ is cofinal if and only if it satisfies F 0, F 1, F 2, i.e. if and only if the $i\bs I'$ are filtering, and then $I$ is locally filtering.}).
\item[\textit{d)}] \textit{If $I$ is right prefiltering (i.e. satisfying right PS 1), and $u$} fully faithful, \textit{for $u$ to be cofinal, it is} necessary \textit{and it is} sufficient \textit{that $u$ satisfy condition F 1 above. And then $I'$ is also prefiltering, and the $i\bs I'$ are weakly filtering. If moreover $I$ is filtering, the $i\bs I'$ and $I'$ are as well. (Compare SGA 4, I 8.1.3.)}
\end{itemize}
\DERmark{ENV0115}\DERmark{ENV0116}


\msp{page 47}
\textsc{Proof.} a) F 1 is obvious, since 0-connected $\Longrightarrow$ nonempty. F 0 because $W_0$-coaspheric $\Longrightarrow$ $W_0$-equivalence, i.e. $\pi_0(u)$ bijective.

To \textit{prove} b), it suffices to prove that F 1, F 2 (if $I'$ is weakly filtering) imply that the $i\bs I'$ are weakly filtering (a fortiori 0-connected) and moreover that $I$ is as well. We already know that they are nonempty by F 1, remains to see that if $(i', \alpha)$, $(j', \beta)$ are in $i\bs I'$, then they are ``dominated'' by a third object of $i\bs I'$. Since $I'$ is weakly filtering, we find
\begin{tikzcd}[row sep=0pt, column sep=small, baseline=(k.base)]
i' \arrow[dr, "r"] & \\
& |[alias=k]| k' \\
j' \arrow[ur, "s"'] &
\end{tikzcd}, whence $u(r)$, $u(s)$ in $I$ (cf. diagram)
\[
\begin{tikzcd}[row sep=small, column sep=small]
& i' \arrow[rr, "r"] \arrow[dd, dash] & & k' \arrow[r, dotted, "t"] \arrow[dd, dash] & l' \arrow[dd, dash] \\
& & j' \arrow[ur, "s"] \arrow[dd, dash] & & \\
& u(i') \arrow[rr, "u(r)"] & & u(k') \arrow[r, dotted, "u(t)"] & u(l') \\
i \arrow[ur, "\alpha"] \arrow[rr, "\beta"'] & & u(j'), \arrow[ur, "u(s)"'] & &
\end{tikzcd}
\]
applying F 2 to the double arrow $u(r)\alpha, u(s)\beta : i \rightrightarrows u(k')$, we find $k' \xrightarrow{t} l'$ (dotted on \texttt{[the]} diagram) such that $(u(t)u(r))\alpha = (u(t)u(s))\beta$, say $\gamma$, which shows that the object $(l', \gamma : i \longrightarrow u(l'))$ dominates $(i', \alpha)$ and $(j', \beta)$.

The proof that if $I'$ is weakly filtering and $u$ satisfies F 1, then $I$ is weakly filtering, is immediate.

\textit{Let us prove} c). It remains to prove that if $I'$ is \textit{filtering} (not merely weakly filtering), then if $u$ \texttt{[is]} cofinal, then F 2 is satisfied, and moreover the $i\bs I'$ and $I$ are filtering (not merely weakly filtering). But let us note at once that $i\bs I' \longrightarrow I'$ is coétale, hence $I'$ filtering, hence locally filtering implies that $i\bs I'$ is also locally filtering\,; since by hypothesis the $i\bs I'$ are 0-connected, they are therefore filtering. Condition F 2 comes out easily now:

\printedpage{31}
% head: A. Grothendieck : Dérivateurs, Ch. II, Édition des Universités de Montpellier II et Paris VII
% foot: Version 8 juillet 2001
\msp{page 48}
\[
\begin{tikzcd}[row sep=small, column sep=small]
& i' \arrow[rr, "r"] \arrow[dd, dash] & & j' \arrow[r, dotted, "t"] \arrow[dd, dash] & k' \arrow[dd, dash] \\
& & i' \arrow[ur, "s"] \arrow[dd, dash] & & \\
& u(i') \arrow[rr, "u(r)"] & & u(j') \arrow[r, dotted, "u(t)"] & u(k') \\
i \arrow[ur, "\alpha"] \arrow[rr, "\beta"'] & & u(i'). \arrow[ur, "u(s)"'] & &
\end{tikzcd}
\]
Since $i\bs I'$ is filtering, hence weakly filtering (PS 1), we can find $i' \dbl{r}{s} j'$ such that
\[
u(r)\alpha = u(s)\beta,
\]
and since $I'$ is filtering, we can find $t : j' \longrightarrow k'$ which equalizes the double arrow $(r, s)$:
\[
tr = ts(\defeq \tau),
\]
then
\[
\underbrace{u(t)u(r)}_{u(tr)}\alpha = \underbrace{u(t)u(s)}_{u(ts)}\beta \qquad\quad \text{i.e.} \qquad\quad u(\tau)\alpha = u(\tau)\beta,
\]
i.e. we have verified F 2.

Remains to prove that $I$ is filtering (if $I'$ is and \texttt{[if]} $u : I' \longrightarrow I$ is cofinal). We already know that $I$ is weakly filtering by b), remains to verify the double arrow condition
\[
i \dbl{\alpha}{\beta} j
\]
(PS 2). Let $i'$ be in $I'$ and $j \xrightarrow{\gamma} u(i')$, and consider
\DERmark{ENV0117}

\[
i \dbl{\gamma\alpha}{\gamma\beta} u(i'),
\]
by hypothesis F 2 $\exists\, r : i' \longrightarrow j'$ in $I'$ such that
\msp{page 49}
$u(r)(\gamma\alpha) = u(r)(\gamma\beta)$, hence $u(r)\gamma$ equalizes $(\alpha, \beta)$, q.e.d.

The verification of d) by AQT \texttt{[abbreviation whose meaning eludes the editors]}.

\printedpage{32}
% head: A. Grothendieck : Dérivateurs, Ch. II, Édition des Universités de Montpellier II et Paris VII
% foot: Version 8 juillet 2001
\section*{4\quad Biaspheric functors}

Let again
\[
u : I' \longrightarrow I
\]
be an arrow of Cat, let $\cM$ be a category, whence
\[
u^* : \cM(I) \longrightarrow \cM(I');
\tag*{(4.1)}
\]
we seek conditions on $u$, $\cM$ for this functor to be fully faithful. If the left adjoint functor $u_!$ (resp. right adjoint $u_*$) exists, we know that it amounts to the same to say that the adjunction morphism
\[
u_!u^*F \longrightarrow F
\tag*{(4.2)}
\]
resp.
\[
u_*u^*F \longleftarrow F
\tag*{(4.2 bis)}
\]
is \texttt{[an]} isomorphism for every $F$ in $\cM(I)$. Let us then say that $u$ is $\cM$-connected, or $\cM$-biaspheric, this terminology implying that this is a property which refers to the derivator $\mathbf{D}_{\cM} = (S \longmapsto \cM(S))$ on Cat defined by $\cM$. Let us note that
\[
\begin{aligned}
&\text{If $u$ is $\cM$-connected, then it is also $\cM(S)$-connected, for every $S$ in Cat,}\\
&\text{since $\cM(S)(I) \simeq \cM(I)(S) \defeq \uHom(S^o, \cM(I))$,}
\end{aligned}
\tag*{(4.3)}
\]

\msp{page 50}
and
\[
\text{If $\cM' \hookrightarrow \cM$ is fully faithful, then $\cM$-connected implies $\cM'$-connected,}
\tag*{(4.4)}
\]
finally, conversely,
\[
\begin{aligned}
&\text{for $u$ to be $\cM$-connected, it is necessary and sufficient that for every \textit{small}}\\
&\text{full subcategory $\cM'$ of $\cM$, $u$ be $\cM'$-connected.}
\end{aligned}
\tag*{(4.5)}
\]
Since every small $\cM'$ embeds in $\cM'^\wedge = \operatorname{Ens}(\cM')$, it follows from (4.3), (4.4), (4.5) that
\[
\text{If $u$ is Ens-connected, it is $\cM$-connected for every category $\cM$,}
\tag*{(4.6)}
\]
for this is the case if $\cM$ is small by (4.3), hence for every $\cM$ by (4.5).

Let us say that a functor $u$ in Cat is $\cM$-\textit{cofinal} or $\cM$-coaspheric if the functorial homomorphism
\[
\varprojlim_{I} F \longrightarrow \varprojlim_{I'} u^*(F)
\tag*{(4.7)}
\]
(arrow in $\uHom(\cM(I), \cM)$ if the $\varprojlim$ of types $I$, $I'$ exist in $\cM$, otherwise arrow in $\uHom(\cM(I), \cM^\wedge)$) is an isomorphism. Reducing to the case of small categories $\cM$, each embedded in $\cM^\wedge$, we see by criterion (4.2 bis) that
\[
\text{$u$ is $\cM$-connected if for every $i \in \Ob I$,}
\tag*{(4.8)}
\]

\printedpage{33}
% head: A. Grothendieck : Dérivateurs, Ch. II, Édition des Universités de Montpellier II et Paris VII
% foot: Version 8 juillet 2001
\msp{page 51}
the induced functor $I'/i \longrightarrow I/i$ is $\cM$-cofinal. The converse is not clear, for $u$ is $\cM$-connected if and only if for every $i \in \Ob I$ and every $\cM$-presheaf $F_i$ on $I/i$ \textit{induced by a presheaf $F$ on $I$,} the homomorphism $\varprojlim\limits_{I/i} F_i \longrightarrow \varprojlim\limits_{I'/i} (u/i)^*(F/i)$ is \texttt{[an]} isomorphism. We can therefore assert straightaway that this implies likewise for \textit{every} $F_i$ on $I/i$, only if we know that $\alpha_i^* : \cM(I) \longrightarrow \cM(I/i)$ is essentially surjective. This is so for every $i$ when $I$ is \textit{ordered,} for then $I/i \longrightarrow I$ is fully faithful, hence $\cM(I/i) \longrightarrow \cM(I)$ \texttt{[is]} fully faithful, hence $\cM(I) \longrightarrow \cM(I/i)$ \texttt{[is]} a localization functor, in particular essentially surjective.
\DERmark{ENV0118}


\msp{page 52}
But when $\cM = (\text{Ens})$, this perplexity disappears. We find the

\noindent\textbf{Proposition 4.1.} \textit{Let $u : I' \longrightarrow I$ be an arrow in} Cat. \textit{Equivalent conditions:}
\begin{itemize}[label={}]
\item[\textit{a)}] \textit{For every category $\cM$, $u^*_{\cM} : \cM(I) \longrightarrow \cM(I')$ is fully faithful (i.e. $u$ is $\cM$-connected).}
\item[\textit{a$'$)}] \textit{For every category $\cM$ in which direct limits of type $I$, $I'$ exist, the adjunction arrow for $F$ in $\cM(I)$} % Transcription corrected: official printed33 reads a prime.
\[
u_!u^*F \xrightarrow{\sim} F
\]
\textit{is} \texttt{[an]} \textit{isomorphism.}
\item[\textit{a$''$)}] \textit{Dual of a$'$), with} \texttt{[the]} \textit{adjunction arrow}
\[
u_*u^* F\xleftarrow{\sim} F. % [sic] u_*u^* -> u_*u^*F, the object F is missing on the left side (compare a') and (4.2 bis))
\]
\DERmark{ENV0119}

\item[\textit{b) b$'$) b$''$)}] \textit{As in a) a$'$) a$''$), with $\cM =$} (Ens).
\item[\textit{c)}] \textit{As in b$'$), with $F$ representable, i.e. of the form $j \longmapsto$} $\Hom_I(j, i)$.
\item[\textit{d)}] \textit{For every $i$ in $I$, $I'/i \longrightarrow I/i$ is cofinal.} % [sic] Pour tout I dans I -> Pour tout i dans I, the variable of the statement is i (compare d'))
\item[\textit{d$'$)}] \textit{For every $i$ in $I$, $i\bs I' \longrightarrow i\bs I$ is co-cofinal.}
\item[\textit{d$''$)}] \textit{For every arrow $i_0 \xrightarrow{\alpha} i_1$ in $I$, the category} $i_0\bs \underset{\alpha}{I'}/i_1 \defeq I' \times_I \underbrace{i_0\bs \underset{\alpha}{I}/i_1}_{\substack{\text{biinduced category}\\ (i_0, \alpha)\bs(I/i_1) \simeq (i_0\bs I)/(i_1, \alpha)}}$ \qquad \textit{is}

0-\textit{connected.}
\DERmark{ENV0120}\DERmark{ENV0121}

\end{itemize}

\msp{page 53}
\textsc{Proof.} We already know, from the preliminaries, the equivalence of a, a$'$, a$''$, b, b$'$, b$''$, let A be this condition - we shall say that $u$ is $W_0$-\textit{connected} or $W_0$-\textit{biaspheric,} if these equivalent conditions are satisfied. We have moreover the self-duality
\[
\text{$u$ is $W_0$-connected if and only if $u^o$ is,}
\tag*{(4.8)} % [sic] (4.8) -> another number, the label (4.8) is already used for a display on p. 32
\]
for we obviously have
\[
\text{$u$ is $\cM$-connected if and only if $u^o$ is $\cM^o$-connected.}
\tag*{(4.9)}
\]


\endgroup
% END UNIT Der-ChII-234/Der-ChII-234_p031-035.tex
% BEGIN UNIT Der-ChII-234/Der-ChII-234_p036-040.tex
\begingroup\makeatletter%
\newcommand{\printedpage}[1]{\marginpar{\footnotesize p.~#1}}%
\newcommand{\msp}[1]{\marginpar{\footnotesize [#1]}}%
\newcommand{\Hom}{\operatorname{Hom}}%
\newcommand{\Ob}{\operatorname{Ob}}%
\newcommand{\bD}{\mathbf{D}}%
\newcommand{\bDelta}{\boldsymbol{\Delta}}%
\makeatother\let\printedpage\derpp\let\msp\dermsp%
\let\mspin\DERmspin%


\printedpage{34}
% head: A. Grothendieck : Dérivateurs, Ch. II, Édition des Universités de Montpellier II et Paris VII
This shows the equivalence of d, d$'$, which can also be seen from the equivalence of each of these two conditions with d$''$, by the cofinality criterion of the preceding paragraph. Let D be this condition on $u$. By (4.8) we have D $\Longrightarrow$ A, hence % [sic] l'equivalence -> l'équivalence, accent missing; précedent -> précédent, accent missing
\[
\mathrm{D} \Longrightarrow \mathrm{A} \Longrightarrow \mathrm{c},
\]
it remains to prove that c $\Longrightarrow$ D, more precisely c $\Longrightarrow$ d$''$ and c$'$ $\Longrightarrow$ d$''$ \texttt{[c$'$?]}. Let us make explicit the adjunction morphism (for $i$ in $I$)
\[
u_!\underbrace{u^*i} \longrightarrow i,
\]
i.e. for $j$ in $I$
\[
\underbrace{(u_!u^*i)(j)}_{\pi_0(j\backslash(I'/u^*(i)))} \longrightarrow i(j) = \Hom(j,i).
\]

\msp{page 54}Now $I'/u^*(i) \simeq I'/i$, so we have the arrow
\[
(u_!u^*(i))(j) \simeq \pi_0(j\backslash(I'/i)) \longrightarrow \Hom_I(j,i). \tag*{(4.10)}
\]
Now we see at once that for $\alpha : j \longrightarrow i$ in $\Hom(j,i)$, the fiber of (4.10) at $\alpha$ identifies with $\pi_0(j\backslash\underset{\alpha}{I'}/i\ )$, more precisely, $j\backslash\underset{\alpha}{I'}/i$\ \ is the open and closed subcategory of $J = j\backslash(I'/i)$ corresponding to the part of $\pi_0(J)$ which lies over $\alpha$ in (4.10). This completes the proof of 4.1.
\DERmark{ENV0122}


NB Another way to prove that c $\Longrightarrow$ A is to note that the two sides in the adjunction morphism $u_!u^*F \longrightarrow F$ in $I^{\wedge}$ are functors $I^{\wedge} \longrightarrow I^{\wedge}$ which commute with arbitrary $\varinjlim$. Thus to check that it is \texttt{[an]} isomorphism, it suffices to check it for $F$ in $I$, since $I$ generates $I^{\wedge}$ by $\varinjlim$ . On the other hand, an argument of this type does not seem to prove that to check that we have an \textit{isomorphism} $F \longrightarrow u_*u^*F$, it suffices to do so for representable $F$, i.e. to check b$''$ for representable $F$.

It would remain to see whether c$'$ $\Longrightarrow$ d$''$ \texttt{[c$'$?]}.

\msp{page 55}Let us make explicit the adjunction morphism
\[
i \longrightarrow u_*u^*i
\]
(for $i \in \Ob I$) by its values on $j$
\[
\underbrace{i(j)}_{\Hom(j,i)} \longrightarrow \Gamma(I'/j, u^*i). \tag*{$(\ast)$}
\]
So let
\[
u_j : I'/j \longrightarrow I,
\]
% foot: Version 8 juillet 2001
\printedpage{35}
% head: A. Grothendieck : Dérivateurs, Ch. II, Édition des Universités de Montpellier II et Paris VII
we have
\[
\begin{aligned}
\Gamma(I'/j, \underbrace{u^*(i)|I'/j}_{u_j^*(i)}) &\simeq \Hom(e_{I'/j}, u^*(i)) \\ % ?? uncertain: extent of the underbrace -- over "u^*(i)|I'/j" or including the closing parenthesis
&\overset{\mathrm{adj}}{\simeq} \Hom(u_{j,!}e_{I'/j}, i),
\end{aligned}
\]
and the arrow $(\ast)$, functorial in $i$, is deduced from an arrow
\[
u_{j,!}(e_{I'/j}) \longrightarrow j. \tag*{$(\ast\ast)$}
\]
We see at once, as above, that $(\ast\ast)$ is \texttt{[an]} isomorphism if and only if for every $\alpha : k \longrightarrow j$, the category $k\backslash\underset{\alpha}{I'}/j$ is 0-connected, and if we require this for every $j$, $k$, $\alpha$, this is condition d$''$. But it is not clear that for $(\ast\ast)$ to be an isomorphism, it suffices that the Hom of the two sides with values in every $i$ in $I$ (and not every $F$ in $I^{\wedge}$) be \texttt{[an]} isomorphism. The difficulty comes from the fact that $u_{j,!}(e_{I'/j})$ is not necessarily representable, and that an arbitrary $F$ in $I^{\wedge}$ \msp{page 56}is not necessarily $\varprojlim$ of representable functors. So it is doubtful whether c$'$ \texttt{[c$'$?]} suffices to imply that $u$ is $W_0$-connected.

\textbf{Corollary 4.2.} \textit{For $u$ to be $W_0$-connected, it} suffices \textit{that it be a universal $W_0$-equivalence, i.e. (cf. 3.1) that the fibers be \textup{0}-connected, and that for every nonidentity arrow in $I$, defining $\bDelta^1 \longrightarrow I$, $I' \times_I \bDelta^1$ be \textup{0}-connected.}

4.3. This condition of universal $W_0$-equivalence is however \textit{strictly} stronger than $W_0$-connectedness, as we see from the example where $u$ is an \textit{equivalence of categories} not surjective on objects (e.g. the inclusion of a vertex in a 1-connected groupoid having more than one object). More generally, every localization functor is $W_0$-connected, whereas its fibers are not necessarily connected, and a fortiori it is not a universal $W_0$-equivalence.

However, if every endomorphism $p$ of an $i$ in $I$, such that $p^2 = p$ (projector) is the identity, and if \msp{page 57}$I' \longrightarrow I$ is ``transportable'', e.g. if in $I$ every isomorphism is an identity, then $u$ $W_0$-connected implies that the fibers of $u$ are 0-connected, as we see by expressing d$''$ for $\alpha$ an identity arrow. The conditions envisaged on $I$ are satisfied if $I$ is ordered. But even if $I'$ and $I$ are both ordered, $u$ can be $W_0$-connected, without being a universal $W_0$-equivalence. Here is an example with $I = \bDelta^2$:
\[
\begin{array}{r@{\;}l@{\quad}c}
I' = & \left\{\begin{tikzcd}[ampersand replacement=\&]
 \& 1' \arrow[r] \arrow[d] \& 2 \\
0 \arrow[r] \& 1 \&
\end{tikzcd}\right\} & \\[1ex]
 & & \Big\downarrow \\[1ex]
I = & \left\{\begin{tikzcd}[ampersand replacement=\&]
0 \arrow[r] \& 1 \arrow[r] \& 2
\end{tikzcd}\right\}. &
\end{array}
\]
% foot: Version 8 juillet 2001

\printedpage{36}
% head: A. Grothendieck : Dérivateurs, Ch. II, Édition des Universités de Montpellier II et Paris VII
The bilocalized categories of $I$ are the three vertices $\{0\}$, $\{1\}$, $\{2\}$, and the three segments $[0,1]$, $[1,2]$ and $[0,2] = I$, the inverse images are respectively
\[
\bullet\ ,\quad
\begin{tikzcd}[row sep=normal] {} \arrow[d] \\ {} \end{tikzcd}\ ,\quad
\bullet\ ,\quad
\begin{tikzcd}[row sep=normal] & {} \arrow[d] \\ {} \arrow[r] & {} \end{tikzcd}\ ,\quad
\begin{tikzcd}[row sep=normal] {} \arrow[r] \arrow[d] & {} \\ {} & \end{tikzcd}\ ,\quad
I'\quad ,
\]
they are all 0-connected, i.e. $W_0$-aspheric, and even contractible hence $W_\omega$-aspheric (so $u$ is even $W_\omega$-connected, i.e. $W_\omega$-biaspheric). But it is not a universal $W_0$-equivalence, since the inverse image of $I'$ under $\bDelta^1 \xrightarrow{\partial} \bDelta^2$, $\partial(0) = 0$, $\partial(1) = 2$, \msp{page 58}is the disconnected category $\varepsilon = \{\bullet\ \ \bullet\}$.

I had hoped for a moment that the $W_0$-connected functors $I' \xrightarrow{u} I$ would be those which factor as $I' \xrightarrow{u_0} I'' \xrightarrow{v} I$, with $u_0$ an equivalence, and $v$ a universal $W_0$-equivalence. But the preceding example defeats this, with $u$ even $W_\omega$-connected. Indeed, $I''$ would be derived from a preordered set, and $v$ would factor through the corresponding ordered set, namely $I'$, hence would be of the form
\[
I'' \longrightarrow I' \xrightarrow{u} I,
\]
where $I'' \longrightarrow I'$ is a passage to the quotient, and the inverse image of $\bDelta^1 \xrightarrow{\partial} I$ would again be disconnected (a $\pi_0$ of cardinality 2). So it is really hopeless.

However

\textbf{Corollary 4.3.} \textit{Suppose that $u$ is $W_0$-smooth, or $W_0$-proper, or a weak $W_0$-fibration. Then $u$ is $W_0$-connected if and only if it is a universal $W_0$-equivalence, i.e. if and only if a) its fibers are \textup{0}-connected and b) $I' \times_I \bDelta^1$} \texttt{[is]} \textit{\textup{0}-connected for every $\bDelta^1 \longrightarrow I$ (corresponding to a nonidentity arrow).}

\msp{page 59}In all three cases we know that $u$ is a universal $W_0$-equivalence if and only if its fibers are $W_0$-connected. Now $W_0$-smoothness implies that $\forall i \in \Ob I$, the inclusion
\[
I'_i \longrightarrow i\backslash I'
\]
is a $W_0$-equivalence, on the other hand $i\backslash I' \longrightarrow i\backslash I$ being $W_0$-aspheric, hence a $W_0$-equivalence by hypothesis, $i\backslash I'$ is 0-connected, hence $I'_i$ also. The case $u$ $W_0$-proper is dual. In the case \texttt{[where]} $u$ \texttt{[is]} a weak $W_0$-fibration, we again know that $I'_i \longrightarrow i\backslash I'$ is a $W_0$-equivalence, and we conclude as in the smooth case (or, as one chooses, via $I'_i \longrightarrow I'/i$).

But it occurs to me that the $W_0$-connected functors might be exactly the localization functors? It amounts to the same thing to say that a \textit{conservative} functor $u : I' \longrightarrow I$ (i.e. $u(\alpha)$ isomorphism $\Longrightarrow$ $\alpha$ isomorphism) which is $W_0$-connected, is an equivalence. But this is false, as we see by taking $I = e$, $I'$ a 0-connected and not 1-connected groupoid: $I' \longrightarrow e$ is $W_0$-connected (and even a universal $W_0$-equivalence), but it is not an equivalence (and not a localization functor).
% foot: Version 8 juillet 2001

\printedpage{37}
% head: A. Grothendieck : Dérivateurs, Ch. II, Édition des Universités de Montpellier II et Paris VII
\msp{page 60}I would like to generalize the notion of a $W_0$-biaspheric arrow to the case of an arbitrary derivator $\bD$.

\textbf{Proposition 4.4.} \textit{Let $\bD$ be a derivator on} Diag, \textit{and $u : I' \longrightarrow I$ an arrow in} Diag. \textit{Suppose that $u_*$ and $u_!$ exist, and are computed in the standard way, for every arrow $u$ in} Diag. \textit{(In fact, it suffices that this be the case for $u$ and its localized arrows $u/i : I'/i \longrightarrow I/i$, or $u$ and its colocalized arrows $i\backslash u : i\backslash I' \longrightarrow i\backslash I$.) Consider the conditions:} % [sic] pour tout flèche -> pour toute flèche, gender agreement
\begin{itemize}[leftmargin=4em]
\item[\textit{a)}] $u^* : \bD(I) \longrightarrow \bD(I')$ \texttt{[is]} \textit{fully faithful.}
\item[\textit{a$'$)}] \textit{$\forall$ $F$ in $\bD(I)$, the arrow}
\[
u_!u^*F \longrightarrow F
\]
\textit{is an isomorphism.}
\item[\textit{a$''$)}] \textit{$\forall$ $F$ in $\bD(I)$, the arrow}
\[
u_*u^*F \longleftarrow F
\]
\textit{is an isomorphism.}
\item[\textit{b)}] \textit{$\forall$ $i$ in $I$, the arrow $i\backslash I' \longrightarrow i\backslash I$ is $\bD$-aspheric.}
\item[\textit{b$'$)}] \textit{$\forall$ $i$ in $I$, the arrow $I'/i \longrightarrow I/i$ is $\bD$-coaspheric.}
\item[\textit{b$''$)}] \textit{$\forall$ $\alpha : i_0 \longrightarrow i_1$ in $I$, the category $i_0\backslash\underset{\alpha}{I'}/i_1$ is $\bD$-aspheric.}
\end{itemize}
\textit{We then have the implications}
\[
(b \Longleftrightarrow b' \Longleftrightarrow b'') \Longrightarrow (a \Longleftrightarrow a' \Longleftrightarrow a'').
\]
\textit{We have $a \Longrightarrow b$ (i.e. all the legitimate conditions envisaged are equivalent)} \msp{page 61}\textit{in each of the following cases.}
\begin{itemize}[leftmargin=4em]
\item[$1^o$] \textit{$I$ is preordered.}
\item[$2^o$] \textit{$\forall$ $i$ in $I$, $\alpha_i^* : \bD(I) \longrightarrow \bD(I/i)$ is essentially surjective.}
\item[$3^o$] \textit{$\forall$ $i$ in $I$, considering $\alpha_i : I/i \longrightarrow I$, the functor $\alpha_{i,!} : \bD(I/i) \longrightarrow \bD(I)$ is conservative.}
\item[$4^o$] $\mathcal{D} = \mathrm{HOT}_{W_0} \approx \bD_{\mathrm{Ens}}$.
\end{itemize}
\textsc{Proof}. The equivalence of a, a$'$, a$''$ is clear, say $A$. Likewise the equivalence of b, b$'$, b$''$, say $B$. % [sic] l'equivalence -> l'équivalence, accent missing

Let us prove
\[
B \Longrightarrow A
\]
assuming e.g. the existence of $u_*$. We tacitly assume, along with the existence of $u_*$ : $\bD(I') \longrightarrow \bD(I)$, the method of computing the fibers of $u_*(-)$. Thus $A$, via a$''$, is equivalent to
\[
H^{\bullet}_{\bD}(I'/i, u^*F) \xleftarrow{\mathit{iso}} H^{\bullet}_{\bD}(I/i, F) \simeq F(i)
\]
% foot: Version 8 juillet 2001
\printedpage{38}
% head: A. Grothendieck : Dérivateurs, Ch. II, Édition des Universités de Montpellier II et Paris VII
for every $F$ in $\bD(I)$, every $i$ in $I$.

Let
\[
\alpha_i : I/i \longrightarrow I, \qquad F_i = \alpha_i^*(F), % [sic] I/i \longrightarrow F -> I/i \longrightarrow I, alpha_i is defined as I/i -> I in 3^o
\]
\DERmark{ENV0123}

we must therefore express that for every $F_i$ coming from an $F$, we have
\[
H^{\bullet}_{\bD}(I/i, F_i) \xrightarrow{\mathrm{iso}} H^{\bullet}_{\bD}(I'/i, (u/i)^*(F_i)). \tag*{$(\ast)$}
\]
Now subject to the existence of $(u/i)_!$, to say that $(\ast)$ is \texttt{[an]} isomorphism for \textit{every} $F_i$ in $\bD(I/i)$, means that
\[
u/i : I'/i \longrightarrow I/i
\]
\msp{page 62}is $\bD$-coaspheric, whence b$'$ $\Longrightarrow$ a$''$. It is clear that the reverse implication holds if $\alpha_i^*$ is essentially surjective, this is case $2^o$. Now $2^o$ is satisfied if $I$ is preordered, whence case $1^o$. Case $4^o$ has already been treated (4.1). We must still treat case $3^o$, which moreover includes case $2^o$, hence case $1^o$ (and also case $4^o$, cf. below).

We must express the isomorphy of the arrow $(\ast)$ for $F_i$ of the form $\alpha_i^*(F)$. Let $\xi$ be \texttt{[an]} object of $\mathcal{A} = \mathcal{D}(e)$, we must express the isomorphy of $\Hom(\xi, -)$ in the two sides of $(\ast)$, i.e. of
\[
\begin{aligned}
\Hom(\xi_{I/i}, F_i) \quad &\longrightarrow \quad \Hom(\xi_{I'/i}, (u/i)^*(F_i)) \\
&\overset{\text{adj.}}{\simeq} \quad \Hom((u/i)_!(\xi_{I'/i}), F_i),
\end{aligned}
\tag*{$(\ast\ast)$}
\]
\DERmark{ENV0124}\DERmark{ENV0125}\DERmark{ENV0126}

\texttt{[rather} $\xi_{I'/i}$\texttt{]} arrow which is deduced from
\[
(u/i)_!(u/i)^*(\xi_{I/i}) \xrightarrow{\text{adj.}} \xi_{I/i} \tag*{$(\ast\ast\ast)$}
\]
\DERmark{ENV0127}

by taking the Hom into $F_i$. The condition $u/i$ coaspheric (i.e. b$'$) means that $(\ast\ast\ast)$ is an isomorphism for every $\xi$ in $\mathcal{A}$. Does it \textit{suffice} for this that its Hom into an $F_i$ be bijective, when we restrict $F_i$ in $\bD(I/i)$ to be in the image of $\bD(I)$ under $\alpha_i^*$? Now this arrow $(\ast\ast)$, for $F_i = \alpha_i^*F$, is made explicit by the adjunction formula relative \msp{page 63}to the functor $\alpha_i$, as
\[
\Hom(\alpha_{i,!}(\xi_{I/i}), F) \longrightarrow \Hom(\alpha_{i,!}(\ldots), F).
\]
To say that it is an isomorphism for every $F$, thus means \textit{exactly} that the transform of the arrow $(\ast\ast\ast)$ under $\alpha_!$ is an isomorphism. To conclude from this that $(\ast\ast\ast)$ itself is \texttt{[an]} isomorphism, the natural condition to impose is therefore that $\alpha_{i,!} : \bD(I/i) \longrightarrow \bD(I)$ be conservative. This is case $3^o$. This completes the proof of 4.4.

4.5. Let us show that condition $3^o$ on $I$ is not always satisfied, in the case of $\bD = (\mathrm{HOT}_{W_0})^{\mathrm{opp}} \approx (\bD_{\mathrm{Ens}})^{\mathrm{opp}}$, although in this case, it is nevertheless true that for every $I'$ over
% foot: Version 8 juillet 2001


\endgroup
% END UNIT Der-ChII-234/Der-ChII-234_p036-040.tex
% BEGIN UNIT Der-ChII-234/Der-ChII-234_p041-045.tex
\begingroup\makeatletter%
\newcommand{\printedpage}[1]{\marginpar{\footnotesize p.~#1}}%
\newcommand{\msp}[1]{\marginpar{\footnotesize [#1]}}%
\newcommand{\D}{\mathbf{D}}%
\makeatother\let\printedpage\derpp\let\msp\dermsp%
\let\mspin\DERmspin%


\printedpage{39}
% head: A. Grothendieck : Dérivateurs, Ch. II, Édition des Universités de Montpellier II et Paris VII

$I$ satisfying condition a$''$ (self-dual), condition b$'$ (also self-dual) is also satisfied (by case $4^o$ above). I say therefore that there exist $I$ and $i$ in $I$, such that
\[
\beta_{i,*} : (i\backslash I)^\wedge \longrightarrow I^\wedge
\]
is not conservative. (\textbf{NB} On the other hand $\alpha_{i,!} : (I/i)^\wedge \longrightarrow I^\wedge$ is always conservative, which gives back case $4^o$ as a special case of $3^o$). As $\beta_i : i\backslash I \longrightarrow I$ is $W_0$-proper (and even $W_\omega$-proper!), $\beta_{i,*}$ ``is computed fiber by fiber'', so we have
\msp{page 64}
\[
\beta_{i,*}(G)(j) \simeq \prod_{\alpha\in\mathrm{Hom}(i,j)} G(\alpha),
\]
the fiber in question at $j$ being discrete. Let us take
\[
I = 0 \underset{\beta}{\overset{\alpha}{\rightrightarrows}} 1, \qquad\quad \text{with } \alpha \neq \beta,
\]
then $0\backslash I$ is $\Psi = {}$\begin{tikzcd}[ampersand replacement=\&,row sep=0.1em, column sep=1.2em, baseline=(\tikzcdmatrixname-2-1.base)]
 \& 1 \\
0 \ar[ur, "\alpha'"] \ar[dr, "\beta'"'] \& \\
 \& 1'
\end{tikzcd}, and $\Psi^\wedge$ is the category of diagrams of sets $G = {}$\begin{tikzcd}[ampersand replacement=\&,row sep=0.1em, column sep=1.2em, baseline=(\tikzcdmatrixname-2-1.base)]
 \& G_1 \ar[dl, "\lambda"'] \\
G_0 \& \\
 \& G_{1'} \ar[ul, "\mu"]
\end{tikzcd}. The functor $\beta_*$ deduced from $\beta : i\backslash I \longrightarrow I$ transforms this diagram into
\[
G_1 \times G_{1'} \underset{\mu\,\mathrm{pr}_2}{\overset{\lambda\,\mathrm{pr}_1}{\rightrightarrows}} G_0.
\]
\DERmark{ENV0128}

\texttt{[rather} $\mu\,\mathrm{pr}_2$\texttt{]}. If we restrict ourselves to the diagrams $G$ for which $G_{1'} = \emptyset$, $\beta_*$ on these gives
\[
\emptyset \rightrightarrows G_0,
\]
so the restriction of $\beta_*$ to this subcategory (equivalent to $(\boldsymbol{\Delta}^1)^\wedge$) is not faithful, nor conservative, taking account only of the component $G_0$, and not of $G_1$.

\msp{page 63bis}

\textbf{Definition 4.6.} \textit{We say that $u : I' \longrightarrow I$ in} Diag \textit{is} $\D$-biaspheric \textit{if it satisfies the equivalent conditions b, b$'$, b$''$ of 4.4,} $\D$-connected \textit{if it satisfies condition a, or again (subject to the existence of $u_!$ resp. of $u_*$) the equivalent conditions a$'$, a$''$.} % [sic] Definition -> Définition, accent missing in print

Thus, $\D$-biaspheric implies $\D$-connected, and the converse is true when $I$ satisfies (with respect to $\D$) condition $3^o$ of 4.4 (thus also, if it satisfies one of the stronger conditions $2^o$, $1^o$), or the condition dual to this one (since the notions $\D$-biaspheric, $\D$-connected are self-dual). In this regard, let us note the
\DERmark{ENV0129}


\textbf{Corollary 4.7.} \textit{Suppose that $\mathcal{A} = \D(e)$ satisfies the following condition: the functor}
\[
\left\{
\begin{array}{lcccl}
\mathcal{A} & \times & \mathcal{A} & \longrightarrow & \mathcal{A} \\[1ex]
(\xi & , & \eta) & \longrightarrow & \xi \amalg \eta
\end{array}
\right.
\]
\textit{is conservative (resp. faithful). Then for every étale arrow $u : J \longrightarrow I$ in} Diag \textit{(and in particular for $I/i \longrightarrow I$), the functor $u_! : \D(J) \longrightarrow \D(I)$ is conservative (resp. faithful)} (\footnote{\textbf{NB} This condition is satisfied if $\forall\, \xi, \eta \in \mathcal{A}$ we have $\mathrm{Hom}(\xi,\eta) \neq \emptyset$, (for it reduces to this: if $u_1$,
% footnote continues at the foot of p. 40
$u_2$ are two arrows in Ens between sets \textit{that are nonempty,} and if $u_1 \times u_2$ is \texttt{[an]} isomorphism, $u_1$, $u_2$ are \texttt{[some]} isomorphisms). This applies in particular if $\mathcal{A}$ admits a zero object, e.g. if $\mathcal{A}$ is additive.}). \textit{If this condition, or the dual condition,}
\printedpage{40}
% head: A. Grothendieck : Dérivateurs, Ch. II, Édition des Universités de Montpellier II et Paris VII
\msp{page 64bis}
\textit{involving the products $\xi\times\eta$ ($\xi$, $\eta$ in $\mathcal{A}$), is satisfied, then for every arrow $u : I' \longrightarrow I$ in} Diag, \textit{$u$ is $\D$-connected if and only if $u$ is $\D$-biaspheric.}

\textsc{Proof.} The hypothesis of conservativity, or of faithfulness carries over to sums indexed by an arbitrary set (subject to existence), as one sees by writing $\coprod\limits_{\lambda} \xi_\lambda = \xi_{\lambda_0} \amalg \coprod\limits_{\lambda\neq\lambda_0} \xi_\lambda$. The assertion made about $u_!$ then follows from the computation of $u_!$ fiber by fiber (which is valid since $u$ is étale hence smooth), taking into account that the fibers are discrete.

I know of no example of a pair $(\D, u)$ of a derivator $\D$ and an arrow $u$ in Diag, for which $u$ is $\D$-connected without being $\D$-biaspheric. We have given conditions on $\D$, which ensure equivalence of the two notions. If $\D$ is of the form $\D_{\mathcal{M}}$, and if $\mathcal{M}$ satisfies the condition stated in 4.7 or its dual, then by 4.7 it is OK. So a counterexample with a $\D_{\mathcal{M}}$ would have to take this into account. I am not going to persist in trying to clear up how things stand.
\DERmark{ENV0130}


\msp{page 65}

\textbf{Corollary 4.8.} \textit{An arrow $u$ of} Diag \textit{which is a universal $\D$-equivalence is $\D$-biaspheric, and \textup{a fortiori} $\D$-connected. Conversely, if $u$ is either $\D$-smooth, or $\D$-proper, or a weak $\D$-fibration, and if $u$ is $\D$-biaspheric, it is a universal $\D$-equivalence.}

The proof is that of the special case 4.3. We thus have a diagram of implications:
\[
\begin{tikzcd}[column sep=small, row sep=large]
\text{Cat-equivalence} \ar[dr, Rightarrow] & \D\text{-universal equivalence } (W_{\mathrm{univ}}) \ar[d, Rightarrow, "{\text{\footnotesize\begin{tabular}{@{}l@{}}strict implication ; equivalence\\ in the case where $u$ is either $\mathbf{D}$-proper,\\ or $\mathbf{D}$-smooth, or a weak $\mathbf{D}$-fibration.\end{tabular}}}"] & \\
 & \D\text{-biaspheric} \ar[dl, Rightarrow] \ar[d, Rightarrow, "{\text{\footnotesize\begin{tabular}{@{}l@{}}equivalence\\ in all\\ known cases\end{tabular}}}"] \ar[dr, Rightarrow] & \\
\D\text{-aspheric} \ar[dr, Rightarrow] & \D\text{-connected} \ar[d, Rightarrow] & \D\text{-coaspheric} \ar[dl, Rightarrow] \\
 & \D\text{-equivalences } (W) &
\end{tikzcd}
\]
(\footnote{\textbf{NB} These seven classes of morphisms in Cat are stable under composition including the notion of $\D$-connectedness, all these notions depend only on the set $W = W_{\D}$ of $\D$-equivalences.}). Recall that for a given $u : I' \longrightarrow I$
\DERmark{ENV0131}

\begin{enumerate}[leftmargin=*, labelindent=1.5em]
\item[a)] $u$ is $\D$-aspheric if and only if $\forall\, F$ \textit{constant} in $\D(I)$
\[
F \xrightarrow[\sim]{\text{adj.}} u_*u^*F \qquad\quad \text{is \texttt{[an]} isomorphism.} \tag*{$(\ast)$}
\]
\printedpage{41}
% head: A. Grothendieck : Dérivateurs, Ch. II, Édition des Universités de Montpellier II et Paris VII
\item[b)] $u$ is $\D$-coaspheric if and only if $\forall\, F$ \textit{constant} in $\D(I)$
\[
F \xleftarrow[\sim]{\text{adj.}} u_!u^*F \qquad\quad \text{is \texttt{[an]} isomorphism.} \tag*{$(\ast\ast)$}
\]
\item[c)] $u$ is $\D$-connected if and only if for \textit{every} $F$ in $\D(I)$ (not necessarily constant, nor locally constant), $(\ast)$ is \texttt{[an]} isomorphism, or what amounts to the same thing, for every $F$ $(\ast\ast)$ is \texttt{[an]} isomorphism.
\end{enumerate}

\msp{page 66}

Recall that the class $W_{\mathrm{asph}}$ \texttt{[of]} $W$-aspheric arrows and that, $\mathrm{Liss}_W$, of $W$-smooth arrows, determine each other, thus
\[
\left\{
\begin{array}{lcl}
(f : X \longrightarrow Y) \in W_{\mathrm{asph}} & \Longleftrightarrow & \forall\, (Y' \xrightarrow{\ u\ } Y) \in \mathrm{Liss}_W, \\
 & & X' = X \times_Y Y' \longrightarrow Y' \text{ is } \in W. % [sic] X' = X \times_Y Y' \longrightarrow X' -> X' = X \times_Y Y' \longrightarrow Y', the target must be Y' (compare (4.12))
\\
(f : Y' \longrightarrow Y) \in \mathrm{Liss}_W & \Longleftrightarrow & \text{The functor } X \longmapsto X' = X \times_Y Y' \\
 & & \text{from } \mathrm{Cat}/Y \text{ to } \mathrm{Cat}/Y' \text{ transforms} \\
 & & \text{arrows in } W_{\mathrm{asph}} \text{ in \texttt{[rather into]}} \\
 & & \text{arrows in } W_{\mathrm{asph}}.
\end{array}
\right.
\tag*{(4.10)}
\]
\DERmark{ENV0132}


To every subset $A$ of Fl(Cat) stable under isomorphisms, one can associate a subset $\mathbf{L}(A)$ of Fl(Cat), by the relation
\[
\begin{array}{rcl}
(u : Y' \longrightarrow Y) \in \mathbf{L}(A) & \Longleftrightarrow & u^* : \mathrm{Cat}/Y \longrightarrow \mathrm{Cat}/Y' \\
 & & \text{transforms arrows} \\
 & & \in A \text{ into arrows } \in A.
\end{array}
\tag*{(4.11)}
\]

We see that $\mathbf{L}(A)$ contains the isomorphisms, is stable under composition, and is also stable under base change (as is the case for $\mathrm{Liss}_W$). Conversely, to a subset $L$ of Fl(Cat), stable under isomorphisms, we associate a subset $\mathbf{A}_W(L)$ by
\[
\begin{array}{rcl}
(f : X \longrightarrow Y) \in \mathbf{A}_W(L) & \Longleftrightarrow & \forall\, (Y' \longrightarrow Y) \text{ in } L, \\
 & & X' = X \times_Y Y' \longrightarrow Y' \\
 & & \text{is in } W.
\end{array}
\tag*{(4.12)}
\]

If $L$ is stable under base change, then $\mathbf{A}_W(L)$ is stable under composition. It contains the isomorphisms (since $W$ contains them), is contained in $W$ if $L$ contains the identities. But
\msp{page 67}
no reason for it to be stable under arbitrary base change (and in fact, $W_{\mathrm{asph}}$ is not).

(\footnote{Botched \texttt{[up to the end of} \S4\texttt{]}, cf. \S6.}). We have, for a given $A \subseteq W \subseteq \mathrm{Fl(Cat)}$, stable under isomorphisms,
\[
A \subseteq \mathbf{A}_W(\mathbf{L}(A)) \overset{\mathrm{def}}{=} A^\natural \subseteq W, \qquad \text{\footnotesize\begin{tabular}{@{}l@{}}subset stable under composition,\\ contains \texttt{[the]} isomorphisms,\end{tabular}}
\tag*{(4.13)}
\]
\DERmark{ENV0133}

\printedpage{42}
% head: A. Grothendieck : Dérivateurs, Ch. II, Édition des Universités de Montpellier II et Paris VII
and for every subset $L$ of Fl(Cat) which is stable under base change, and under composition
\[
L \subseteq \mathbf{L}(\mathbf{A}_W(L)) \overset{\mathrm{def}}{=} L^\flat, \qquad \text{\footnotesize\begin{tabular}{@{}l@{}}subset of Fl\,Cat stable\\ under composition and under base\\ change, containing the isomorphisms.\end{tabular}}
\tag*{(4.14)}
\]
\DERmark{ENV0134}

We have
\[
\left\{
\begin{array}{lcl}
A^{\natural\natural} & = & A^\natural, \quad \text{since } \mathbf{A}_W(L)^\natural = \mathbf{A}_W(L) \\[1ex]
L^{\flat\flat} & = & L^\flat, \quad \text{since } \mathbf{L}(A)^\flat = \mathbf{L}(A).
\end{array}
\right.
\]
Examples:
\begin{enumerate}[leftmargin=*, labelindent=1em]
\item[$1^o$)] $L = \text{Ét}$, $\mathbf{A}_W(L) = W_{\mathrm{asph}}$, $\mathbf{L}(\mathbf{A}_W(L)) = L^\flat = \mathrm{Liss}_W$.
\item[$2^o$)] $L = \text{Coét}$, $\mathbf{A}_W(L) = W_{\mathrm{coas}}$, $\mathbf{L}(\mathbf{A}_W(L)) = L^\flat = \mathrm{Prop}_W$.
\item[$3^o$)] $A = W$, $\mathbf{L}(A) = \mathrm{Fib}_W$, $A^\natural = W$.
\item[$4^o$)] $L = \mathrm{Fl\,Cat}$, $\mathbf{A}_W(L) = W_{\mathrm{univ}}$, $L^\flat = L$.
\end{enumerate}
\msp{page 68 blank}

\par
\printedpage{43}
% head: A. Grothendieck : Dérivateurs, Ch. II, Édition des Universités de Montpellier II et Paris VII
\msp{page 69}

\section*{5\quad Formal duality}

Let $\mathcal{P}$, $\mathcal{Q}$ be two ordered sets. I am trying to isolate a notion of ``duality'' (formal) linking $\mathcal{P}$ and $\mathcal{Q}$. This will be a relation between $\mathcal{P}$ and $\mathcal{Q}$
\[
\mathcal{R} \subseteq \mathcal{P} \times \mathcal{Q}
\tag*{(5.1)}
\]
having the following properties
\begin{enumerate}[leftmargin=*, labelindent=0.5em]
\item[(D)] For every subset $A$ of $\mathcal{P}$, the set
\end{enumerate}
\[
A^\cup \overset{\mathrm{def}}{=} \{ q \in \mathcal{Q} \mid (p,q) \in \mathcal{R} \text{ for every } p \in A,\ \text{ i.e. } A \times \{q\} \subseteq \mathcal{R} \} % [sic] \mathbf{R} -> \mathcal{R}, the relation is written \mathcal{R} everywhere else
\]
\DERmark{ENV0135}\DERmark{ENV0136}

is an open subset of $\mathcal{Q}$ having a greatest element $\beta(A) : A^\cup = \mathcal{Q}/\beta(A)$. Dually, for every subset $B$ of $\mathcal{Q}$, the set
\[
B^\cap \overset{\mathrm{def}}{=} \{ p \in \mathcal{P} \mid (p,q) \in \mathcal{R} \text{ for every } q \in B,\ \text{ i.e. } \{p\} \times B \subseteq \mathcal{R} \}
\]
\DERmark{ENV0137}

is an open subset of $\mathcal{P}$ having a greatest element $\alpha(B)$, i.e. $B^\cap = \mathcal{P}/\alpha(B)$.

Most often, $\mathcal{P}$ and $\mathcal{Q}$ will be stable under arbitrary suprema (hence also under arbitrary infima). In this case, (D) is equivalent, on the side of the $A^\cup$ (for $A \in \mathfrak{P}(\mathcal{P})$) to this:
\begin{enumerate}[leftmargin=*, labelindent=0em]
\item[(D$_0$)] For every $p \in \mathcal{P}$, $\mathcal{R}(p) = \{ q \in \mathcal{Q} \mid (p,q) \in \mathcal{R} \}$ is an open \textit{filtering} subset of $\mathcal{Q}$, or what amounts to the same thing, a nonempty subset ($\Longleftrightarrow$ containing the least element $\emptyset_{\mathcal{Q}}$ of $\mathcal{Q}$) stable under suprema of two elements. Moreover, $\mathcal{R}(p)$ is stable under filtering suprema, or what amounts to the same thing (being itself filtering), it contains a greatest element. It amounts % [sic] Q -> \mathcal{Q}, plain italic Q printed in "q \in Q" instead of the script \mathcal{Q}
\msp{page 70}
to the same thing to say that $\mathcal{R}(p)$ is an \textit{open subset} of $\mathcal{Q}$ stable under arbitrary suprema (which implies that it contains $\emptyset_{\mathcal{Q}}$). Dually, for every $q \in \mathcal{Q}$, ${}^{s}\mathcal{R}(q) = \{ p \in \mathcal{P} \mid (p,q) \in \mathcal{R} \}$ is an \textit{open} subset of $\mathcal{P}$, stable under arbitrary suprema, i.e. an open subset having a greatest element. % doubt resolved against PDF leaf 15: raised s, transpose relation
\DERmark{ENV0138}

\end{enumerate}

In other words, when one postulates condition (D) for $A$ resp. $B$ reduced to one element, then it remains satisfied for every subset. This comes from the formulas
\[
\left\{
\begin{array}{lcl}
(\bigcup_i A_i)^\cup & = & \bigcap_i A_i^\cup \\[1ex]
(\bigcup_j B_j)^\cap & = & \bigcap_j B_j^\cap,
\end{array}
\right.
\tag*{(5.2)}
\]
and from the fact that the intersection of open subsets of the form $\mathcal{P}/\alpha_i$ in $\mathcal{P}$ is of the form $\mathcal{P}/\alpha$, where $\alpha = \inf_i \alpha_i$, and likewise in $\mathcal{Q}$. Thus we find
\[
\left\{
\begin{array}{lcl}
\beta(\bigcup_i A_i) & = & \inf_i \beta(A_i) \\[1ex]
\alpha(\bigcup_j B_j) & = & \inf_j \alpha(B_j).
\end{array}
\right.
\tag*{(5.3)}
\]


\endgroup
% END UNIT Der-ChII-234/Der-ChII-234_p041-045.tex
% BEGIN UNIT Der-ChII-234/Der-ChII-234_p046-050.tex
\begingroup\makeatletter%
\newcommand{\printedpage}[1]{\marginpar{\footnotesize p.~#1}}%
\newcommand{\msp}[1]{\marginpar{\footnotesize [#1]}}%
\newcommand{\cP}{\mathcal{P}}%
\newcommand{\cQ}{\mathcal{Q}}%
\newcommand{\cR}{\mathcal{R}}%
\newcommand{\fP}{\mathfrak{P}}%
\newcommand{\sR}{{}^{s}\mathcal{R}}%
\newcommand{\circled}[1]{\tikz[baseline=(c.base)]\node[draw,rounded corners=2pt,inner sep=1.5pt](c){#1};}%
\makeatother\let\printedpage\derpp\let\msp\dermsp%
\let\mspin\DERmspin%


\printedpage{44}
% head: A. Grothendieck : Dérivateurs, Ch. II, Édition des Universités de Montpellier II et Paris VII

We shall put, for a subset $A$ of $\cP$ resp. $B$ of $\cQ$
\[
A^{\natural} = (A^{\cup})^{\cap}, \qquad B^{\flat} = (B^{\cap})^{\cup}
\tag*{(5.4)}
\]
so that we shall have
\[
A \subseteq A^{\natural}, \qquad B \subseteq B^{\flat}
\tag*{(5.5)}
\]
We shall say that $A$ is $\cR$-\textit{closed} if $A = A^{\natural}$ (or again if $A$ is of the form $B'^{\cap}$ for $B' \subseteq \cQ$, cf. 5.6), that
\msp{71}
$B$ is $\cR$-\textit{closed} if $B = B^{\flat}$. We shall have for every $A \subseteq \cP$, $B \subseteq \cQ$
\[
\left\{
\begin{array}{rcll}
(A^{\cup})^{\flat} & = & A^{\cup}, & \text{i.e. } A^{\cup} \text{ is } \cR\text{-closed,}\\[1ex]
(B^{\cap})^{\natural} & = & B^{\cap}, & \text{i.e. } B^{\cap} \text{ is } \cR\text{-closed.}
\end{array}
\right.
\tag*{(5.6)}
\]
In particular
\[
(A^{\natural})^{\natural} = A^{\natural}, \qquad (B^{\flat})^{\flat} = B^{\flat},
\tag*{(5.7)}
\]
i.e. the $\cR$-closure $A^{\natural}$ of $A \subseteq \cP$ is $\cR$-closed, and likewise the $\cR$-closure $B^{\flat}$ of $B \subseteq \cQ$. % [sic] cloture -> clôture, accent circonflexe manquant (cf. « clôture » dans la même ligne)

In fact, I realize that I am mixing two different kinds of facts. The relations (5.2), and (5.4) to (5.7), do not depend on the datum of order relations on $\cP$, $\cQ$, nor on axioms connecting these with the given relation $\cR$, but only on this relation, which need not be subject to any condition. Let us note that % [sic] la donné -> la donnée, accord manquant
\[
\emptyset^{\cup} = \cQ, \qquad \emptyset^{\cap} = \cP,
\tag*{(5.8)}
\]
whereas
\[
\emptyset^{\natural} = \cQ^{\cap}, \qquad \emptyset^{\flat} = \cP^{\cup}
\tag*{(5.9)}
\]
are respectively the smallest $\cR$-closed subset of $\cP$ and of $\cQ$.

The map
\[
A \longmapsto A^{\natural}
\tag*{(5.10)}
\]
is an \textit{increasing} (and \textit{idempotent}) map from $\fP(\cP)$ to itself, thus having the property
\msp{72}
\[
(\bigcup_i A_i)^{\natural} = (\bigcup_i A_i^{\natural})^{\natural}.
\]

\printedpage{45}
% head: A. Grothendieck : Dérivateurs, Ch. II, Édition des Universités de Montpellier II et Paris VII

For every $A \subseteq \cP$, $A^{\natural}$ is the smallest $\cR$-closed subset of $\cP$ containing $A$, and likewise for every $B \subseteq \cQ$, $B^{\flat}$ is the smallest $\cR$-closed subset of $\mathcal Q$ containing $B$. We have moreover (this should even have been said before saying that $A \longmapsto A^{\natural}$ is increasing) % [sic] de $Q$ -> de $\mathcal{Q}$, lettre droite au lieu de la lettre calligraphique ; [sic] contentant -> contenant, faute de frappe
\DERmark{ENV0139}

\[
\left\{
\begin{array}{rcl}
A \subseteq A' & \Longrightarrow & A'^{\cup} \subseteq A^{\cup}\\[1ex]
B \subseteq B' & \Longrightarrow & B'^{\cap} \subseteq B^{\cap},
\end{array}
\right.
\tag*{(5.11)}
\]
i.e. the maps
\[
\begin{tikzcd}
\fP(\cP) \arrow[r, shift left, "\beta"] & \fP(\cQ) \arrow[l, shift left, "\alpha"]
\end{tikzcd}
\qquad
\begin{array}{l}
A \longmapsto A^{\cup}\\
B \longmapsto B^{\cap}
\end{array}
\tag*{(5.12)}
\]
are decreasing (which implies that the two composites $\alpha\beta = (A \longmapsto A^{\natural})$ and $\beta\alpha = (B \longmapsto B^{\flat})$ are increasing). Their restrictions to the subsets $\fP_{\cR}(\cP)$ and $\fP_{\sR}(\cQ)$ of the $\cR$-closed subsets of $\cP$ resp. of $\cQ$, are, by definition of these as
\[
\begin{array}{rcl}
\fP_{\cR}(\cP) & = & \{A \subseteq \cP \mid A^{\natural} \ (\overset{\mathrm{def}}{=} \alpha\beta(A)) = A\}\\
\fP_{\sR}(\cQ) & = & \{B \subseteq \cQ \mid B^{\flat} \ (\overset{\mathrm{def}}{=} \beta\alpha(B)) = B\},
\end{array}
\]
\DERmark{ENV0140}\DERmark{ENV0141}

inverse to one another. We thus find

\textbf{Proposition 5.1.} \textit{Let $\cP$, $\cQ$ be two sets, $\cR \subseteq \cP \times \cQ$ a relation. Then the maps $\alpha$, $\beta$ of (5.12) induce bijections inverse to one}
\msp{73}
\textit{another, between the set $\mathit{Cl}_{\cR}(\cP)$ of the $\cR$-closed subsets of $\cP$, and the set $\mathit{Cl}_{\sR}(\cQ)$ of the ${}^{s}\cR$-closed subsets (we shall say $\cR$-closed by abuse of language) of $\cQ$. These bijections are decreasing, hence antiisomorphisms for the natural order relations (of inclusion).}

\textbf{Corollary 5.2.} \textit{Let $(A_i)_i$ be a family of $\cR$-closed subsets of $\cP$. Then $\bigcap A_i$ is a $\cR$-closed subset of $\cP$, and}
\[
(\bigcap_i A_i)^{\cup} = \underbrace{(\bigcup_i A_i^{\cup})^{\flat}}_{\substack{\text{sup of the}\\ \text{subsets } \cR\text{-closed}\\ A_i^{\cup} \text{ of } \cQ}} \quad .
\tag*{(5.2.1)}
\]

But the relation $\cR$ ``extends'' to a relation $\bar{\cR}$ between
\[
\bar{\cP} = \fP(\cP), \qquad \bar{\cQ} = \fP(\cQ)
\tag*{(5.13)}
\]
by the prescription
\[
\underbrace{\bar{\cR}}_{\subseteq \bar{\cP} \times \bar{\cQ}} = \{(A,B) \in \bar{\cP} \times \bar{\cQ} \mid B \subseteq A^{\cup},\ \text{i.e. } A \subseteq B^{\cap},\ \text{i.e. } A \times B \subseteq \cR\}.
\tag*{(5.14)}
\]
The sets $\bar{\cP}$, $\bar{\cQ}$ are stable under arbitrary sups (union of subsets of $\cP$ resp. $\cQ$), and the relation $\bar{\cR}$ on these ordered sets satisfies condition (D$_0$)
\printedpage{46}
% head: A. Grothendieck : Dérivateurs, Ch. II, Édition des Universités de Montpellier II et Paris VII
\msp{74}
above (p. 69), in an obvious way: for fixed $A$, the $B$ such that $(A,B) \in \bar{\cR}$ are those which are contained in $A^{\cup}$, so the set $\bar{\cR}(A)$ is nothing other than $\bar{\cQ}/A^{\cup}$, and dually ${}^{s}\bar{\cR}(B)$ is nothing other than $\bar{\cP}/B^{\cap}$. Thus
\[
\beta(A) = A^{\cup}, \qquad \alpha(B) = B^{\cap}.
\tag*{(5.14)} % [sic] (5.14) -> numéro distinct, numéro déjà employé pour la formule (5.14) de la page précédente
\]
If I thought of introducing axiom (D$_0$), for a pair of arbitrary ordered sets, it is because of the cases where these are \textit{not} of the form $\fP(\cP)$, $\fP(\cQ)$, and (even if they were) where the basic relation is not deduced as above from an arbitrary relation between sets $\cP$, $\cQ$.

I therefore return to the axiomatics of the beginning, with $\cP$, $\cQ$ ordered, and $\cR$ satisfying (D). In all the cases I can think of, $\cP$ and $\cQ$ are stable under arbitrary sups, so (D) is equivalent to (D$_0$). Let us note that the map
\[
\cP \longrightarrow \fP(\cP), \qquad p \longmapsto \cP/p
\]
is injective, and likewise for $\cQ \longrightarrow \fP(\cQ)$. Thus, knowledge of $A \longmapsto A^{\cup}$ is equivalent to that of % [sic] connaisance -> connaissance, faute d'orthographe
\[
\begin{array}{rcl}
A & \longmapsto & \beta(A)\\[0.5ex]
\beta : \fP(\cP) & \longrightarrow & \cQ,
\end{array}
\]
and likewise that of $B \longmapsto B^{\cap}$ is equivalent to
\msp{75}
that of
\[
\alpha : \fP(\cQ) \longrightarrow \cP.
\]
In fact, it suffices to give oneself
\[
\begin{array}{rcccl}
\beta_0 : & \cP & \longrightarrow & \cQ & \qquad \text{maps}\\
\alpha_0 : & \cQ & \longrightarrow & \cP & \qquad \textit{decreasing}
\end{array}
\tag*{(5.15)}
\]
defined by
\[
\left\{
\begin{array}{rcl}
\beta_0(p) & = & \beta(\{p\})\\[1ex]
\alpha_0(q) & = & \alpha(\{q\}),
\end{array}
\right.
\]
since we deduce $\beta$, $\alpha$ from them by
\[
\left\{
\begin{array}{rcl}
\beta(A) & = & \inf\nolimits_{p \in A} \beta_0(p)\\[1ex]
\alpha(B) & = & \inf\nolimits_{q \in B} \alpha_0(q).
\end{array}
\right.
\tag*{(5.16)}
\]
The relation $\cR$ is deduced from the datum of $\beta_0$, or at one's choice of $\alpha_0$, by
\[
(p,q) \in \cR \iff q \leq \beta_0(p) \iff p \leq \alpha_0(q).
\tag*{(5.17)}
\]

\printedpage{47}
% head: A. Grothendieck : Dérivateurs, Ch. II, Édition des Universités de Montpellier II et Paris VII

This therefore implies that $\alpha_0$, $\beta_0$ determine one another, by the formulas
\[
\left\{
\begin{array}{rcl}
\alpha_0(q) & = & \sup \underbrace{\{p \in \cP \mid q \leq \beta_0(p)\}}_{\{q\}^{\cap}}\\[3ex]
\beta_0(p) & = & \sup \underbrace{\{q \in \cQ \mid p \leq \alpha_0(q)\}}_{\{p\}^{\cup}}.
\end{array}
\right.
\tag*{(5.18)}
\]
It remains to ask which maps $\beta_0 : \cP \longrightarrow \cQ$, or $\alpha_0 : \cQ \longrightarrow \cP$, correspond to a \textit{``dualizing relation''} % [sic] \alpha_0 : \cP \to \cQ, \beta_0 : \cQ \to \cP -> \beta_0 : \cP \to \cQ, \alpha_0 : \cQ \to \cP, les rôles de \alpha_0 et \beta_0 sont échangés par rapport à (5.15)
\DERmark{ENV0142}

\msp{76}
$\cR$ on $\cP \times \cQ$. Here are the conditions on $\beta_0 : \cP \longrightarrow \cQ$:
\begin{itemize}[leftmargin=4em, labelsep=0.8em]
\item[\circled{a}] $\beta_0$ is \textit{decreasing},
\end{itemize}
which implies that the relation
\[
\underbrace{\cR_{\beta_0}}_{\subseteq \cP \times \cQ} = \{(p,q) \in \cP \times \cQ \mid q \leq \beta_0(p)\}
\]
satisfies condition (D$_0$) on the left side, and we still have to check it on the right side, i.e. look, for fixed $q$, at the set
\[
\{q\}^{\cap} = \{p \in \cP \mid q \leq \beta_0(p)\}.
\]
As $\beta_0$ is decreasing, we see that $p \in \{q\}^{\cap}$, and $p' \leq p \Longrightarrow p' \in \{q\}^{\cap}$, i.e. $\{q\}^{\cap}$ is \textit{open.} We must still check (assuming the existence of arbitrary sups in $\cP$) that this set is stable under arbitrary sups - it will follow that it is of the form $\cP/\alpha(q)$. This is equivalent to the condition
\begin{itemize}[leftmargin=4em, labelsep=0.8em]
\item[\circled{b}] $\beta_0$ transforms arbitrary sups in $\cP$ into infs in $\cQ$.
\end{itemize}
If we regard $\beta_0$ as a contravariant functor $\cP^{o} \longrightarrow \cQ$, this means that it commutes with inverse limits, or again that it has a left adjoint $\cQ \longrightarrow \cP^{o}$, which commutes with direct limits, and corresponds
\msp{77}
to the decreasing map $\alpha_0 : \cQ \longrightarrow \cP$, transforming sups into infs.

Here it is convenient to consider $\natural$ and $\flat$ as operations on $\cP$ resp. $\cQ$, by
\[
p^{\natural} = \alpha_0\beta_0(p), \qquad q^{\flat} = \beta_0\alpha_0(q).
\tag*{(5.19)}
\]
The relation between these endomorphisms, and the operation of the same name on $\fP(\cP)$, $\fP(\cQ)$, is given by
\[
\left\{
\begin{array}{rcl}
p^{\natural} & = & \text{Greatest element of } \{p\}^{\natural}\\[1ex]
A^{\natural} & = & \{p\in\cP\mid p\leq (\sup\nolimits_{a\in A}a)^{\natural}\} = \{p\in\cP\mid p\leq (\sup\nolimits_{a\in A}a^{\natural})^{\natural}\},
\end{array}
\right.
\tag*{(5.20)}
\]
\DERmark{ENV0143}

\printedpage{48}
% head: A. Grothendieck : Dérivateurs, Ch. II, Édition des Universités de Montpellier II et Paris VII
and symmetrically on the side of $\cQ$.

We say that $p$ is $\cR$-\textit{closed,} or $\beta_0$-closed, or $\alpha_0$-closed, if $p = p^{\natural}$, and likewise $q \in \cQ$ is said to be $\cR$-\textit{closed} etc. if $q = q^{\flat}$. The maps
\[
\left\{
\begin{array}{ccccccc}
p & \longmapsto & p^{\natural} & : & \cP & \longrightarrow & \cP\\[1ex]
q & \longmapsto & q^{\flat} & : & \cQ & \longrightarrow & \cQ
\end{array}
\right.
\tag*{(5.21)}
\]
are \textit{increasing} and \textit{idempotent.} We can characterize $p^{\natural}$ as the smallest closed element of $\cP$ lying above $p$, and likewise for $q^{\flat}$. We then find, by (5.19), that $\alpha_0$ and $\beta_0$ induce bijections inverse to one another,
\msp{78}
between the set of closed elements of $\cP$, and the set of closed elements of $\cQ$
\[
\begin{tikzcd}
\cP^{\natural} \arrow[r, shift left, "\bar{\beta}_0"] \arrow[r, phantom, "\scriptstyle\sim"] & \cQ^{\flat}. \arrow[l, shift left, "\bar{\alpha}_0"]
\end{tikzcd}
\tag*{(5.22)}
\]
But in practice, the maps $\alpha_0$, $\beta_0$ are not obtained directly, but we have at our disposal the relation $\cR$, satisfying condition (D$_0$), from which we deduce $\alpha_0$, $\beta_0$. In all the cases I am thinking of here, either $\cP$, $\cQ$ are of the form $\fP(P_0)$, $\fP(Q_0)$, and $\cR$ is deduced from a relation $R_0 \subseteq P_0 \times Q_0$ as above (by far the most frequent, it seems to me), or else $\cP$ and $\cQ$ appear as suitable subsets of sets $\fP(P_0)$, $\fP(Q_0)$, equipped with the induced order, and the relation $\cR$ on $\cP \times \cQ$ is obtained in a more complex way.

\noindent\textbf{Examples.}
\begin{enumerate}[label={\arabic*)}]
\item Relations of \textit{orthogonality} between subsets of two modules over a ring, equipped with a duality pairing. In the case of vector \texttt{[spaces]} of finite dimension and of a true duality, the closed subsets are the vector subspaces, in the infinite case, they are the vector subspaces closed for the ``weak topology'' defined by the pairing.
\item Relations of \textit{orthogonality} between subsets of two locally compact commutative topological groups dual to one another, in the sense of Pontrjagin.
\msp{79}

The closed subsets are the closed subgroups.
\item For two vector spaces ``in duality'' over a complete valued field, such as $\mathbf{R}$ or $\mathbf{C}$, the relation % [sic] où -> ou, accent grave fautif
\[
|\langle x, x' \rangle| \leq 1.
\]
In the real or complex case, the closed subsets are the nonempty convex circled [\,?] weakly closed subsets.
\DERmark{ENV0144}

\item For real vector spaces, the relation of semi-polarity
\[
\langle x, x' \rangle \geq -1.
\]
The closed subsets are the convex weakly closed subsets containing the origin.
\end{enumerate}


\endgroup
% END UNIT Der-ChII-234/Der-ChII-234_p046-050.tex
% BEGIN UNIT Der-ChII-234/Der-ChII-234_p051-055.tex
\begingroup\makeatletter%
\newcommand{\printedpage}[1]{\marginpar{\footnotesize p.~#1}}%
\newcommand{\msp}[1]{\marginpar{\footnotesize [#1]}}%
\newcommand{\Fl}{\mathrm{Fl}}%
\newcommand{\Pf}{\mathfrak{P}}%
\setcounter{footnote}{61}%
\makeatother\let\printedpage\derpp\let\msp\dermsp%
\let\mspin\DERmspin%


\printedpage{49}
% head: A. Grothendieck : Dérivateurs, Ch. II, Édition des Universités de Montpellier II et Paris VII
% foot: Version 8 juillet 2001

\begin{itemize}[leftmargin=2.5em]
\item[5)] Let $K$ be a field (for simplicity), $L$ a Galois extension with group $G$. Relation between $L$ and $G$
\[
\gamma\cdot x=x,\qquad \text{i.e. } x\in L^{\gamma}.
\]
The closed subsets of $L$ are the subextensions, the closed subsets of $G$ the subgroups of $G$ closed for the Krull topology.
\DERmark{ENV0145}

\end{itemize}

There are various analogues in inseparable cases (by dint of $p$-Lie algebras for example), and descent statements of all kinds.

\msp{page 80}All these examples are of the 1\textsuperscript{st} kind, the ordered sets that occur are of the form $\Pf(P_0)$, $\Pf(Q_0)$, and the relation between them is deduced from a relation $R_0$ between $P_0$ and $Q_0$. Here now are more complex examples.

\begin{itemize}[leftmargin=2.5em]
\item[6)] Let $\mathcal{M}$ be a category (``big'', most often \dots). We take
\[
\mathcal{P}=\mathcal{Q}=\Pf(\Fl(\mathcal{M})).
\]
The relation $\mathcal{R}\subseteq\mathcal{P}\times\mathcal{Q}$ is Quillen's relation, which I write, for $\Phi,\Psi\in\Pf$
\[
\Psi\subseteq\Phi_{*},
\]
or again equivalently
\[
\Phi\subseteq\Psi^{*},
\]
($\Psi$ has the LLP \texttt{[left lifting property]} with respect to $\Phi$, $\Phi$ the RLP \texttt{[right lifting property]} with respect to $\Psi$). Thus the maps $\alpha_0$, $\beta_0$ are given by % [sic] Psi a la LLP pour Psi -> Psi a la LLP pour Phi, the relation is between Psi and Phi (cf. "Phi la RLP pour Psi")
\DERmark{ENV0146}

\[
\beta_0(\Phi)=\Phi_{*},\qquad \alpha_0(\Psi)=\Psi^{*}.
\]
Here it would occur to no one to introduce $\Pf(\mathcal{P})$, $\Pf(\mathcal{Q})$ - it is already big enough as it is! The notation for $\Phi^{\natural}$ is here $\tilde{\Phi}$, for $\Psi^{\flat}$ it is $\bar{\Psi}$. The closed subsets are those I call Q-closed on the left and on the right.

The relation $\mathcal{R}$, which I write
\[
\Phi\longleftrightarrow\Psi,
\]
\msp{page 81}is not deduced from a relation between arrows of $\mathcal{M}$, i.e. from a subset of $\Fl(\mathcal{M})\times\Fl(\mathcal{M})$, but is of a more complex nature. To describe it, one must introduce the set $\mathcal{D}$ of commutative square diagrams (solid lines) of the form
\[
\begin{tikzcd}
A \arrow[r] \arrow[d, "i"] & X \arrow[d, "p"] \\
B \arrow[r] \arrow[ur, dotted, "f"] & Y
\end{tikzcd}
\tag*{$(*)$}
\]
\printedpage{50}
% head: A. Grothendieck : Dérivateurs, Ch. II, Édition des Universités de Montpellier II et Paris VII
% foot: Version 8 juillet 2001
(which can be interpreted as $\Fl(\underline{\Fl}(\mathcal{M}))$, the set of arrows in $\underline{\Fl}(\mathcal{M})$). We have a subset $\mathcal{D}_0$ of $\mathcal{D}$, consisting of the diagrams for which \texttt{[there]} exists \texttt{[an]} $f$ with two commutative triangles (it is thus the image of $\underbrace{\Fl_3(\mathcal{M})}_{\mathrm{Hom}(\boldsymbol{\Delta}^3,\mathcal{M})}$ in $\mathrm{Hom}(\boldsymbol{\Delta}^1\times\boldsymbol{\Delta}^1,\mathcal{M})$, transpose of $\boldsymbol{\Delta}^1\times\boldsymbol{\Delta}^1\longrightarrow\boldsymbol{\Delta}^3$ given by $(0,0)\longmapsto 0$, $(0,1)\longmapsto 1$, $(1,0)\longmapsto 2$, $(1,1)\longmapsto 3$), and a map
\[
\mathcal{D}_0\subseteq\mathcal{D}\xrightarrow{\ \varphi\ }\Fl(\mathcal{M})\times\Fl(\mathcal{M}),
\]
associating to $(*)$ the pair $(i,j)$ \texttt{[rather} $(i,p)$\texttt{]} (source and target of $(*)$, considered as an arrow in $\underline{\Fl}(\mathcal{M})$ from $i$ to $p$). The relation $\mathcal{R}$ is then deduced from $\varphi$, $\mathcal{D}_0$ by the condition
\msp{page 82}
\[
\Phi\longleftrightarrow\Psi\qquad\Longleftrightarrow\qquad\varphi^{-1}(\Psi\times\Phi)\subseteq\mathcal{D}_0.
\]
\DERmark{ENV0147}

But to tell the truth, I realize that it is after all deduced from a relation between elements of $\Fl(\mathcal{M})$, and this relation between $i : A\longrightarrow B$ and $p : X\longrightarrow Y$ can be written as follows:
\begin{itemize}[leftmargin=3em]
\item[$(**)$] for all arrows $A\longrightarrow X$, $B\longrightarrow Y$ making the diagram $(*)$ commutative, the dotted arrow $f$ in $(*)$ exists. % [sic] la flèche $\varphi$ -> la flèche $f$, the dotted arrow in $(*)$ is labelled $f$
\end{itemize}
\DERmark{ENV0148}

\end{itemize}

This makes me think that if we have a relation $\mathcal{R}$ between $\Pf(P)$ and $\Pf(Q)$, satisfying $(\mathrm{D}_0)$, hence described by a decreasing map
\[
\Pf(P)\longrightarrow\Pf(Q)
\]
transforming sup into inf, the datum of such a map is equivalent to that of an arbitrary map % [sic] la donné -> la donnée, feminine noun
\[
P\longrightarrow\Pf(Q),
\]
or what amounts to the same thing, of a relation $\mathcal{R}$ on $P\times Q$. Thus this structure must be regarded as ``of the first kind''. The situation changes from the moment when, to describe the relation $\mathcal{R}$ with the desired properties, we find ourselves compelled to replace $\Pf(P)$, $\Pf(Q)$ by suitable subsets - when we look only at certain subsets of $P$, $Q$. This is what we
\msp{page 83}are going to do in the next §. The case I am going to treat there is the only one of its kind that I have encountered - a case of ``formal duality'' which comes in naturally, and which is not ``of the 1\textsuperscript{st} kind''.

\printedpage{51}
% head: A. Grothendieck : Dérivateurs, Ch. II, Édition des Universités de Montpellier II et Paris VII
% foot: Version 8 juillet 2001
\section*{6\quad Transporters, kernels, stabilizers relative to $W$}

Let $\mathcal{M}$ be a category, equipped with a subset
\[
W\subseteq\Fl(\mathcal{M}).
\tag*{(6.1)}
\]
I am thinking of the case of a localizer, but will introduce an assumption on $W$ only as needed. For simplicity, I assume $\mathcal{M}$ stable under fiber products.

I set
\[
\left\{
\begin{array}{rcl}
\mbox{(\footnotemark)}\mathcal{P} & = & \Pf(W)=\{A\in\Pf(\Fl(\mathcal{M}))\mid A\subseteq W\}\\[1ex]
\mathcal{Q} & = & \{L\in\Pf(\Fl(\mathcal{M}))\mid\
\parbox[c]{17em}{$L$ \textit{stable} under base\\
change (by which I also mean,\\
stable under isomorphism)\\
and under composition. $\}$.}
\end{array}
\right.
\tag*{(6.2)}
\]
\footnotetext{cf. correction p. 86.}
I introduce a relation
\[
\mathcal{R}\subseteq\mathcal{P}\times\mathcal{Q}
\]
in the following way. For $A\in\mathcal{P}$, $L\in\mathcal{Q}$, we set
\[
(A,L)\in\mathcal{R}\iff
\begin{array}[t]{l}
\text{For every cartesian diagram}\\
\begin{tikzcd}[ampersand replacement=\&]
X \arrow[d, "f"] \& X' \arrow[l] \arrow[d, "f'"] \\
Y \& Y' \arrow[l, "p"]
\end{tikzcd}\\
\text{with $f\in A$, $p\in L$, we have $f'\in W$}\\
\text{(i.e. the base changes of arrows of $A$ by $p\in L$ belong to $W$).}
\end{array}
\tag*{(6.3)}
\]
\DERmark{ENV0149}


\msp{page 84}When I order $\mathcal{P},\mathcal{Q}\subseteq\Pf(\Fl(\mathcal{M}))$ by the order relation induced by inclusion of subsets of $\Fl(\mathcal{M})$, we find ordered subsets of $\Pf(\Fl(\mathcal{M}))$, stable under arbitrary unions, and also under arbitrary intersections, with the sole proviso that in the case of $\mathcal{P}$, the family of indices must be $\neq\emptyset$ (since then we find as intersection the whole of $\Fl(\mathcal{M})$, which is not in general contained in $W$). % [sic] la seul réserve -> la seule réserve, feminine agreement

We are going to verify condition $(\mathrm{D}_0)$ on page 69, for $\mathcal{R}$ and the order relations. There are two checks to make.

1$^{o}$) \texttt{[where is ``2$^{o}$''?]} Let us give ourselves $A\in\mathcal{P}$, i.e. $A\subseteq W$, and consider the set $\mathcal{R}(A)$ of $L\in\mathcal{Q}$ such that we have $(A,L)\in\mathcal{R}$. It is clear that this is an open set (and likewise, that ${}^{s}\mathcal{R}(L)$ is an open set of $\mathcal{P}$ for every $L$ in $\mathcal{Q}$) - i.e. $L\in\mathcal{R}(A)$, $L'\subseteq L$ implies % ?? uncertain: left superscript before R(L) (also on p. 53) -- italic s or upright s % [sic] ouvert de $\mathcal{Q}$ -> ouvert de $\mathcal{P}$, cf. "l'ouvert ${}^{s}\mathcal{R}(L)$ de $\mathcal{P}$" p. 53
\DERmark{ENV0150}

\printedpage{52}
% head: A. Grothendieck : Dérivateurs, Ch. II, Édition des Universités de Montpellier II et Paris VII
% foot: Version 8 juillet 2001
$L'\in\mathcal{R}(A)$. It remains to find a greatest element of $\mathcal{R}(A)$. Let us note that if $L\in\mathcal{R}(A)$, then for every $p\in L$, $p : Y'\longrightarrow Y$, the base change
\[
p^{*} : \mathcal M/Y\longrightarrow\mathcal M/Y' : X\longmapsto X\times_{Y}Y'
\]
\DERmark{ENV0151}\DERmark{ENV0152}

must transform arrows $f\in A$ into arrows $g\in W$, as follows from the fact that $L$ is stable under base change.
\[
\begin{tikzcd}
X_1 \arrow[d, "f"'] & X'_1 \arrow[l, "q_1"'] \arrow[d, "g"] \\
X_2 \arrow[d] & X'_2 \arrow[l, "q_2"'] \arrow[d] \\
Y & Y'. \arrow[l, "p"']
\end{tikzcd}
\]
Thus $L\subseteq\mathbf{L}_W(A)$, where we set
\msp{page 85}
\[
\begin{aligned}
\mathbf{L}_W(A)=\{(p : Y'\longrightarrow Y)\in\Fl(\mathcal{M})\quad\mid\quad & p^{*} : \mathcal M/Y\longrightarrow\mathcal M/Y'\ \text{transforms}\\
& \text{arrows of $A$ into arrows \textit{of} $W$.}\}
\end{aligned}
\]
\DERmark{ENV0153}\DERmark{ENV0154}

(\footnote{NB say it rather in terms of cartesian diagrams, since we have made no assumption of stability of $W$ under isomorphisms.}).
\[
\begin{tikzcd}
X_1 \arrow[d, "f"'] \arrow[dr, phantom, "\text{cart.}"] & X'_1 \arrow[l] \arrow[d, "f'"] \\
X_2 \arrow[d] \arrow[dr, phantom, "\text{cart.}"] & X'_2 \arrow[l] \arrow[d] \\
Y & Y' \arrow[l, "p"']
\end{tikzcd}
\]
\[
f\in A\implies f'\in W.
\]
It is clear on the other hand that this subset of $\Fl(\mathcal{M})$ is stable under base change, hence $\mathbf{L}_W(A)\in\mathcal{Q}$, and that $(A,\mathbf{L}_W(A))\in\mathcal{R}$. Thus $\mathbf{L}_W(A)$ is a greatest element of $\mathcal{R}(A)$, so we find
\[
\mathcal{R}=\{(A,L)\in\mathcal{P}\times\mathcal{Q}\mid L\subseteq\mathbf{L}_W(A)\}.
\tag*{(6.5)}
\]
\printedpage{53}
% head: A. Grothendieck : Dérivateurs, Ch. II, Édition des Universités de Montpellier II et Paris VII
% foot: Version 8 juillet 2001
\[
\begin{tikzcd}
X_1 \arrow[d, "f"'] \arrow[dr, phantom, "\text{cart.}"] & X'_1 \arrow[l, "q_1"'] \arrow[d, "f'"] \\
X_2 \arrow[d] \arrow[dr, phantom, "\text{cart.}"] & X'_2 \arrow[l, "q_2"'] \arrow[d] \\
Y_1 \arrow[d] \arrow[dr, phantom, "\text{cart.}"] & Y'_1 \arrow[l, "p_1"'] \arrow[d] \\
Y & Y'. \arrow[l, "p"']
\end{tikzcd}
\]
Conversely, let us start with $L\in\mathcal{Q}$, and show that the open set ${}^{s}\mathcal{R}(L)$ of $\mathcal{P}$ has a greatest element. For this, let us set
\[
\begin{aligned}
\mathbf{A}_W(L)=\{(f : X\longrightarrow Y)\in W\subseteq\Fl(\mathcal{M})\quad\mid\quad\forall\qquad
&\begin{tikzcd}[ampersand replacement=\&]
X \arrow[d, "f"'] \arrow[dr, phantom, "\textit{cart.}"] \& X' \arrow[l, "q"'] \arrow[d, "f'"] \\
Y \& Y', \arrow[l, "p"]
\end{tikzcd}\\
&\text{with $p\in L$, we have $f'\in W$}\}.
\end{aligned}
\tag*{(6.6)}
\]
(NB When $L$ contains the identities, as is the case in particular if it is of the form $\mathbf{L}_W(A)$ - then it contains the isomorphisms, at least if $A$ contains them - then $\mathbf{A}_W(L)$ is the set of $f\in\Fl(\mathcal{M})$ satisfying the condition stated after $|$, without having to specify $f\in W$.) We have $\mathbf{A}_W(L)\in\mathcal{P}$, and we verify that $(\mathbf{A}_W(L),L)\in\mathcal{R}$ (using the stability of $L$ under composition), finally it is clear that $(A,L)\in\mathcal{R}$ implies $A\subseteq\mathbf{A}_W(L)$ (using $A\subseteq W$), so $\mathbf{A}_W(L)$ is a greatest element of ${}^{s}\mathcal{R}(L)$.

\msp{page 86}Let us note that if $W$ \textit{is stable under composition} (which we are going to assume from now on), the same is true of $\mathbf{A}_W(L)$ (which again uses the fact that $L$ is stable under base change). Moreover, if $W$ \textit{contains the isomorphisms} (which we shall also assume), $\mathbf{A}_W(L)$ contains them too. As for $\mathbf{L}_W(A)$, it contains the isomorphisms \textit{provided that} $A$ is stable under isomorphisms; when $A$ is stable under composition and contains the isomorphisms, this condition will automatically be satisfied. We insist on this condition, that is why on the side of $A$, in the definition of $\mathcal{P}$, I am going to introduce this additional condition that $A$ be \textit{stable under isomorphisms.} So I set, correcting (6.2),
\[
\left\{
\begin{array}{rcl}
\mathcal{P} & = & \{A\in\Pf(\Fl(\mathcal{M}))\quad\mid\quad\mathrm{Is}(\mathcal{M})\subseteq A\subseteq W,\quad A\text{ stable under composition}\}\\[1ex]
\mathcal{Q} & = & \{L\in\Pf(\Fl(\mathcal{M}))\quad\mid\quad\mathrm{Is}(\mathcal{M})\subseteq L,\quad L\text{ stable under base change}\}.
\end{array}
\right.
\tag*{(6.7)}
\]


\endgroup
% END UNIT Der-ChII-234/Der-ChII-234_p051-055.tex
% BEGIN UNIT Der-ChII-234/Der-ChII-234_p056-060.tex
\begingroup\makeatletter%
\newcommand{\printedpage}[1]{\marginpar{\footnotesize p.~#1}}%
\newcommand{\msp}[1]{\marginpar{\footnotesize [#1]}}%
\newcommand{\Fl}{\mathrm{Fl}}%
\newcommand{\Is}{\mathrm{Is}}%
\newcommand{\Fib}{\mathrm{Fib}}%
\makeatother\let\printedpage\derpp\let\msp\dermsp%
\let\mspin\DERmspin%


\printedpage{54}
% head: A. Grothendieck : Dérivateurs, Ch. II, Édition des Universités de Montpellier II et Paris VII
We are of course interested above all in the closed elements of $\mathcal{P}$, $\mathcal{Q}$, i.e. the subsets $A$, $L$ of $\Fl(\mathcal{M})$, containing the isomorphisms, stable under composition (and $L$ also stable under base change, $A \subseteq W$) such that we have
\[
A = \mathbf{A}_W(L), \qquad L = \mathbf{L}_W(A) \tag*{(6.8)}
\]
\DERmark{ENV0155}

\texttt{[rather }$\mathbf{L}_W(A)$\texttt{]}.

\msp{page 87}

I shall keep the notations $\natural$, $\flat$ of the preceding paragraph. What are the two extreme pairs?

Starting from $A_0 = \Is(\mathcal{M})$ (smallest of the elements of $\mathcal{P}$), we find $\mathbf{L}_W(A_0) = \Fl(\mathcal{M})$ (largest element of $\mathcal{Q}$), and $A_0^{\natural} = \mathbf{A}_W(\Fl(\mathcal{M})) = W_{\mathrm{univ}}$, the set of arrows which are in $W$ and remain so under any base change:
\[
A_0^{\natural} = W_{\mathrm{univ}}, \qquad \beta_0(A_0) \overset{\mathrm{def}}{=} \mathbf{L}_W(A_0) = \Fl(\mathcal{M}), \qquad \text{\scriptsize where $A_0 = \Is(\mathcal{M})$ smallest element of $\mathcal{P}$.} \tag*{(6.9)}
\]
\DERmark{ENV0156}

Conversely, let us start from $L_0 = \Is(\mathcal{M})$, the smallest element of $\mathcal{Q}$. We have $\mathbf{A}_W(L_0) = W$ ($W$ being stable under composition and containing the isomorphisms), and $L_0^{\flat} = \mathbf{L}_W(W) = \Fib_W$, the $W$-fibrations:
\[
L_0^{\flat} = \Fib_W, \qquad \alpha_0(L_0) = W, \qquad \text{\scriptsize where $L_0 = \Is(\mathcal{M})$, smallest element of $\mathcal{Q}$.} \tag*{(6.10)}
\]
So the left $\mathcal{R}$-closed subsets of $\Fl(\mathcal{M})$ wander between $W_{\mathrm{univ}}$ and $W$, the right $\mathcal{R}$-closed subsets between $\Fib_W$ and $\Fl(\mathcal{M})$.

\textbf{Proposition 6.1.} \textit{Let $A \in \mathcal{P}$, $L \in \mathcal{Q}$ be such that $A = \alpha_0(L)$, $L = \beta_0(A)$ (i.e. $A = \mathbf{A}_W(L)$, $L = \mathbf{L}_W(A)$). For $L$ to be stable under composition, it is necessary and sufficient that $A$ be stable under base change by the arrows $f \in L$. (More generally, if $L \in \mathcal{Q}$ is any element of $\mathcal{Q}$, and if we assume in} \msp{page 88} \textit{addition that $L$ is stable under composition, then $A = \mathbf{A}_W(L)$ is stable under base change by $L$.)}

Terminology: for $A \in \mathcal{P}$, i.e. $A$ a subset of $\Fl(\mathcal{M})$ contained in $W$ containing the isomorphisms and stable under composition, $\mathbf{L}_W(A)$ is called the \textit{transporter of $A$ into $W$}. (There is a dual notion of \textit{cotransporter.}) It is an element of $\mathcal{Q}$, thus a subset of $\Fl(\mathcal{M})$ containing the isomorphisms and stable under base change. On the other hand, for $L$ in $\mathcal{Q}$, the arrows $f \in \mathbf{A}_W(L)$ are called \textit{$W$-$L$-aspheric,} or the \textit{$W$-$L$-equivalences,} or the \textit{$W$-equivalences that are $L$-universal,} and $\mathbf{A}_W(L)$ is called the \textit{$L$-kernel of $W$}. We could denote it more simply by $W_L$ (when this notation does not conflict with $W_S \subseteq \Fl(\mathcal{M}/S)$ for $S \in \mathrm{Ob}\,\mathcal{M}$). It is a subset of $W$ which (like $W$ itself) contains the isomorphisms and is stable under composition. % [sic] élement -> élément, missing accent

Finally, there is also reason to introduce the \textit{stabilizer} $\mathrm{Stab}(A)$, or transporter of $A$ into $A$, which does not depend on the choice of $W$, but only on the subset $A$ of $\Fl(\mathcal{M})$, containing the isomorphisms and stable under composition. We could also denote it by $\Fib_A$, and call its elements \textit{$A$-fibrations} ($A$ now playing the role of $W$). We have $\mathrm{Stab}(A) \subseteq \mathbf{L}_W(A)$, and it would seem that in % [sic] donné -> donnée, feminine noun
% foot: Version 8 juillet 2001
\printedpage{55}
% head: A. Grothendieck : Dérivateurs, Ch. II, Édition des Universités de Montpellier II et Paris VII
\msp{page 89}
practice, it is these stabilizers, more than the transporters $\mathbf{L}_W(A)$, that are important (when the two do not coincide). Thus in Cat, if $W$ is a basic localizer, the $W$-smooth (resp. $W$-proper) morphisms $\mathrm{Liss}_W$ ($\mathrm{Prop}_W$), which are the stabilizers of the sets $W_{\mathrm{as}}$ ($W_{\mathrm{coas}}$) of the sets of $W$-aspheric (resp. $W$-coaspheric) morphisms, playing an essential role, whereas I do not know how to characterize in a meaningful way the arrows of the transporters $\mathbf{L}_W(W_{\mathrm{as}})$, $\mathbf{L}_W(W_{\mathrm{coas}})$, which contain, in addition to the $W$-smooth (resp. $W$-proper) morphisms the $W$-fibrations $\Fib_W$. I do not know how to characterize a class of morphisms which looks so motley! Likewise, there is reason to introduce and characterize the stabilizer of $W_{\mathrm{bias}}$ (set of $W$-biaspheric arrows), as the ``$W$-neat'' morphisms $\mathrm{Net}_W$, by a description similar to that of the $\mathrm{Liss}_W$, $\mathrm{Prop}_W$, using the bilocalized categories $\backslash\underset{u}{X}/$, $\backslash\underset{v}{Y}/$ , instead of the localized categories $X/x$, $Y/y$ or the colocalized categories $x\backslash X$, $y\backslash Y$. Finally there is reason to determine the stabilizer of the set Equ of Cat-equivalences (which does not depend on the choice of a $W$ - but we have $\mathrm{Equ} \in \mathcal{P}(W)$ for every basic $W$). % [sic] donné -> donnée, feminine noun
\DERmark{ENV0157}

% foot: Version 8 juillet 2001

\printedpage{56}
% head: A. Grothendieck : Dérivateurs, Ch. II, Édition des Universités de Montpellier II et Paris VII
\msp{page 90}

\section*{7\quad Cofiltering and locally cofiltering morphisms}

But I do not want to dwell now on these developments, which do not seem urgent to me, and would like to return to the questions connected with the notion of filtering or cofiltering categories.

\textbf{Definition 7.1.} Let $f : X \longrightarrow Y$ be an arrow in Cat. We say that $f$ is \textit{cofiltering} (resp. \textit{locally cofiltering}, resp. \textit{weakly cofiltering} (\footnote[64]{I have changed the terminology: now weakly cofiltering $=$ 0-connected $+$ locally weakly cofiltering, \textit{stable under passage} to $X'$ over $X$, $X'$ étale and $\infty$-connected. Review the terminology in the weakly cofiltering case, which is a treacherous notion, for lack of being stable under passage \texttt{[to]} an $X'$ étale over $X$.}), resp. \textit{locally weakly cofiltering} or again \textit{precofiltering}) if for every $y$ in $Y$, the colocalized category $y\backslash X$ is cofiltering (resp. \dots). (Dually, we say that $f$ is \textit{filtering}, resp. \dots, if the categories $X/y$ are filtering, resp. \dots, i.e. if $f^o : X^o \longrightarrow Y^o$ is cofiltering resp. \dots.) % [sic] Definition -> Définition, missing accent; colocalisé -> colocalisée, agreement with catégorie

\textbf{Remark 7.2.} Notions stable under coétale base change. Contains the isomorphisms and even the Cat-equivalences. The \textit{étale} morphisms are locally cofiltering.
\DERmark{ENV0158}


\textbf{Proposition 7.3.} \textit{Let $f : X \longrightarrow Y$ be a cofiltering morphism (resp. locally cofiltering, resp. weakly cofiltering, resp. precofiltering), and suppose that $Y$ is cofiltering (resp. locally cofiltering, resp. weakly cofiltering, resp. precofiltering). Then $X$ itself is cofiltering (resp. \dots).}

\textsc{Proof}. It is a matter, in the four cases considered, of checking for $X$ the condition $\underbrace{\mathrm{PS}_g1' \;/\; \mathrm{PS}_g1}_{\text{last two cases}} \;/\; \mathrm{PS}_g1' + \mathrm{PS}_g2 \;/\; \mathrm{PS}_g1 + \mathrm{PS}_g2$. But \msp{page 91} for each of the three conditions $\mathrm{PS}_g1$, $\mathrm{PS}_g1'$, $\mathrm{PS}_g2$ separately, we see immediately that if it is satisfied for $Y$ and for the $y\backslash X$, it is satisfied for $X$, whence the proposition.

\textbf{Corollary 7.3.} \textit{The four notions considered in (7.1) are stable under composition.}

\textsc{Proof} immediate, using remark (7.2).
\[
\begin{tikzcd}
X \arrow[d] & z\backslash X \arrow[l] \arrow[d] \\
Y \arrow[d] & z\backslash Y \arrow[l] \arrow[d] \\
Z & z\backslash Z. \arrow[l]
\end{tikzcd}
\]
% foot: Version 8 juillet 2001

\printedpage{57}
% head: A. Grothendieck : Dérivateurs, Ch. II, Édition des Universités de Montpellier II et Paris VII
\textbf{Remark 7.4.}
\begin{enumerate}
\item[1)] Note that the notion of a cofiltering or weakly cofiltering morphism is \textit{not} stable under arbitrary base change, nor even under an open immersion. Example: $e \longrightarrow \boldsymbol{\Delta}^1$, mapping $e$ to 1, is cofiltering,
\[
\begin{tikzcd}
 & e \arrow[d, mapsto] \\
0 \arrow[r] & 1
\end{tikzcd}
\]
but its fiber at 0 is empty, hence is neither cofiltering nor weakly cofiltering. I have not seen an example showing that the notion of a \textit{locally} cofiltering resp. precofiltering arrow is not stable under base change, and would not want to dwell on it, having no need to know the answer.
\item[2)] For $X$ to be cofiltering (resp. \dots), it is necessary and sufficient that $X \longrightarrow e$ be so. Thus the absolute notion is reduced to a relative notion.
\item[3)] If $f : X \longrightarrow Y$ is cofiltering, it is \msp{page 92} $W$-coaspheric (thus $\in W$) for every $W$ satisfying the direct limit axiom Loc (8). In particular, it is $W_0$-coaspheric, i.e. a cofiltering functor (and even merely weakly cofiltering) is \textit{cofinal.}
\end{enumerate}
\DERmark{ENV0159}

The main reason for my interest in cofiltering arrows is contained in the

\textbf{Proposition 7.4.} \textit{Let $f : X \longrightarrow Y$ in} Cat \textit{be locally cofiltering (resp. locally weakly cofiltering). Then:}
\begin{enumerate}
\item[\textit{a)}] \textit{For every $X'$ étale over $X$, the composite $f' : X' \longrightarrow X \longrightarrow Y$ is locally cofiltering (resp. \dots), and for every factorization of $f'$ as $X' \xrightarrow{g'} Y' \longrightarrow Y$ with $Y' \longrightarrow Y$ étale, $g'$ is locally cofiltering (resp. \dots).}
\item[\textit{b)}] \textit{$\forall\ x$ in $X$, the morphism on the localized categories}
\[
f/x : X/x \longrightarrow Y/y \qquad (y = f(x)) % [sic] Y/x -> Y/y, since y = f(x)
\]
\DERmark{ENV0160}

\textit{is cofiltering (resp. weakly cofiltering).}
\item[\textit{c)}] \textit{If $x$, $y$ are as in b), and $U$ is an open subset of $X/x$, denoting by $V$ the open subset of $Y/y$ generated by the image of $(f/x)|_U$ (i.e. the open subset consisting of the $z$ in $Y/y$ such that $z\backslash U$ is nonempty), the induced arrow $g : U \longrightarrow V$} \texttt{[is]} \textit{cofiltering (resp. weakly cofiltering), a fortiori cofinal (and even $W$-coaspheric, for every $W$ satisfying the limit axiom Loc (8))} (\footnote[65]{It suffices here to be under the resp'ed hypothesis ($f$ \textit{weakly} locally cofiltering) except for the NB here \texttt{[``the NB here'' is ``(and even ...Loc (8))'']}. % [sic] respée -> intended word unclear (perhaps « respectée » or « resp. »), not a French word
}). % doubt resolved: PDF prints q here; later g is an edition naming inconsistency (recorded in corrections.csv).
\DERmark{ENV0161}\DERmark{ENV0162}\DERmark{ENV0163}


\msp{page 93}
\textit{Thus for every} cosheaf \textit{$G$ on $V$, with values in a category $\mathcal{M}$ with $\varinjlim$ , we have} \texttt{[an]} \textit{isomorphism} % [sic] pout -> pour, misprint
\[
\varinjlim_U g^*(G) \xrightarrow{\ \sim\ } \varinjlim_V G,
\]
% foot: Version 8 juillet 2001
\printedpage{58}
% head: A. Grothendieck : Dérivateurs, Ch. II, Édition des Universités de Montpellier II et Paris VII
\textit{and for every presheaf $F$ on $V$, with values in a category $\mathcal{M}$ with $\varprojlim$ , we have} \texttt{[an]} \textit{isomorphism}
\[
\varprojlim_U g^*(F) \xleftarrow{\ \sim\ } \varprojlim_V F.
\]
\end{enumerate}

\textsc{Proof}. I shall first spell out the

\textbf{Lemma 7.5.} \textit{Let $h : X' \longrightarrow X$ be an étale morphism in} Cat. \textit{Then $h$ is locally cofiltering. More precisely, for every $x$ in $X$, we have}
\[
\pi_0(x\backslash X') \xrightarrow{\ \sim\ } \mathrm{Ob}\,X'_x,
\]
\textit{every connected component of $x\backslash X'$ contains one and only one element of the discrete category $X'_x \subseteq x\backslash X'$. If $x'$ is in $X'_x$ (i.e. $x'$ an object of $X'$ such that $h(x') = x$), it is an} initial \textit{object of the connected component which contains it. (Thus this component is cofiltering.)}

This proves (and beyond) the first assertion of 7.4, in view of 7.3 (stability under composition), or because an \msp{page 94} étale category (here $y\backslash X'$) over a locally cofiltering (resp. locally weakly cofiltering) category (here $y\backslash X$) is likewise.
\[
\begin{tikzcd}[row sep=large]
y\backslash X' \arrow[r] \arrow[d, "\text{ét.}"'] \arrow[dr, phantom, "\text{cart.}"] & X' \arrow[d, "\text{ét.}"] \\
y\backslash X \arrow[r] \arrow[d] \arrow[dr, phantom, "\text{cart.}"] & X \arrow[d] \\
y\backslash Y \arrow[r] & Y.
\end{tikzcd}
\]
For the second assertion in a), we can now assume $X' = X$. We must prove that for every $y'$ in $Y'$, $y'\backslash X \overset{\mathrm{def}}{=} X \times_{Y'} y'\backslash Y'$ is locally cofiltering (resp. locally weakly cofiltering). Let $y$ in $Y$ be the image of $y'$, and consider the diagram of three cartesian squares opposite.
\DERmark{ENV0164}

\[
\begin{tikzcd}[row sep=large]
y'\backslash X \arrow[r] \arrow[d] \arrow[dr, phantom, "\textit{cart.}"] & y\backslash X \arrow[r] \arrow[d] \arrow[dr, phantom, "\textit{cart.}"] & X \arrow[d] \\
y'\backslash Y' \arrow[r] & y\backslash Y' \arrow[r] \arrow[d] \arrow[dr, phantom, "\textit{cart.}"] & Y' \arrow[d] \\
 & y\backslash Y \arrow[r] & Y
\end{tikzcd}
\]
\DERmark{ENV0165}

Lemma 7.5 implies that $y'\backslash Y' \longrightarrow y\backslash Y'$ is the inclusion into $y\backslash Y'$ of a connected component of $y\backslash Y'$, so it is an open immersion. So $y'\backslash X \longrightarrow y\backslash X$ is an immersion
% foot: Version 8 juillet 2001


\endgroup
% END UNIT Der-ChII-234/Der-ChII-234_p056-060.tex
% BEGIN UNIT Der-ChII-234/Der-ChII-234_p061-065.tex
\begingroup\makeatletter%
\newcommand{\printedpage}[1]{\marginpar{\footnotesize p.~#1}}%
\newcommand{\msp}[1]{\marginpar{\footnotesize [#1]}}%
\newcommand{\margmark}[1]{\marginpar{\footnotesize #1}}%
\newcommand{\PS}{\mathrm{PS}}%
\newcommand{\Ob}{\mathrm{Ob}}%
\newcommand{\tl}[1]{\begin{tabular}{@{}l@{}}#1\end{tabular}}%
\newcommand{\tls}[1]{\text{\footnotesize\begin{tabular}{@{}l@{}}#1\end{tabular}}}%
\newcommand{\circled}[1]{\tikz[baseline=(c.base)]\node[draw,circle,inner sep=1pt](c){#1};}%
\makeatother\let\printedpage\derpp\let\msp\dermsp%
\let\mspin\DERmspin%


%======================================================================
\printedpage{59}
% head: A. Grothendieck : Dérivateurs, Ch. II, Édition des Universités de Montpellier II et Paris VII

open. Since $y\backslash X$ is locally cofiltering (resp. locally weakly cofiltering) by hypothesis, the same is therefore true of $y'\backslash X$. This proves a). Applying a) to the situation $X/x \longrightarrow Y/y$ in b), we see that the colocalized categories $z\backslash(X/x)$ of
\DERmark{ENV0166}\DERmark{ENV0167}

\margmark{!} % margin: "!" with a triple vertical rule alongside the lines "X/x sur Y/y ... donc est 0-connexe."
$X/x$ over $Y/y$ ($z \in \Ob\,(Y/y)$) are locally cofiltering (resp. weakly). But since $X/x$ has a final object over the final object of $Y/y$, it follows that $z\backslash(X/x)$ has a final object, hence is $0$-\textit{connected.}

\msp{page 95}
Thus these categories, being locally cofiltering (resp. weakly), are in fact cofiltering (resp. weakly cofiltering).

Let us prove c). Note that for $z$ in $V$, $z\backslash U$ is identified with the inverse image, under $z\backslash(X/x) \longrightarrow X/x$, of the open subset $U$ of $X/x$, it is therefore an \textit{open subset} of $z\backslash(X/x)$, nonempty by construction of $V$. But a nonempty open subset of a cofiltering (resp. weakly cofiltering) category is again cofiltering (resp. weakly cofiltering). This means that $g : U \longrightarrow V$ is cofiltering (resp. weakly cofiltering), a fortiori cofinal (and $W$-coaspheric, for every $W$ satisfying Loc (8)). This completes the proof.

\noindent\textbf{Remark 7.5.}
\begin{itemize}
\item[1)] The same argument as in c) shows that if $f : X \longrightarrow Y$ is cofiltering (resp. weakly cofiltering), then for every open subset $U$ in $X$, if we denote by $V$ the open subset of $Y$ generated by $f(U)$, i.e. the open subset consisting of the $y$ in $Y$ such that $y\backslash U \neq \emptyset$, then the induced arrow $g : U \longrightarrow V$ is cofiltering (resp. weakly cofiltering), and a fortiori $g$ is cofinal (and more generally, if $f$ is cofiltering, $g$ is $W$-coaspheric for every $W$ satisfying condition Loc (8)).

\msp{page 96}
\item[2)] An interesting special case of 7.4 is the one where we assume that the $y\backslash X$ ($y \in \Ob\,Y$) are direct sums of categories each having an \textit{initial object.} Then we see that this condition is stable under the operations of 7.4 a) [\,?], and consequently, the colocalized categories $z\backslash(X/x)$ of the morphisms localized upstairs $X/x \longrightarrow Y/y$ are categories with an initial object. This then implies that the same is true of the colocalized categories $z\backslash U$ ($z \in \Ob\,V$) in $g : U \longrightarrow V$. It follows that $g$ is $W$-coaspheric for \textit{every} basic localizer $W$ (without having to worry about the axiom on limits).

The condition considered on $f : X \longrightarrow Y$ is satisfied in particular in the following two cases (which seem to me by far the most important): a) $f$ is étale and b) $X$ is discrete.
\DERmark{ENV0168}

\end{itemize}
% foot: Version 8 juillet 2001

%======================================================================
\par
\printedpage{60}
% head: A. Grothendieck : Dérivateurs, Ch. II, Édition des Universités de Montpellier II et Paris VII

\msp{page 97}

\begin{center}
{\small
\begin{tikzcd}[ampersand replacement=\&, column sep=small, row sep=large]
\left\{\begin{tabular}{@{}c@{\ }l@{}}
 & Top $X$\\
 & locally totally 0-connected\\
 & and 0-connected\\
$\Longleftrightarrow$ & totally 0-connected\\
 & and locally irreducible\\
$\Longleftrightarrow$ & $X$ \textit{cofiltering,}\\
 & i.e. $\PS^*1' + \PS^*2$\\
$\Longleftrightarrow$ & $X$ 0-connected\\
 & $+ \PS^*1 + \PS^*2$\\
$\Longleftrightarrow$ & $\left\{\begin{tabular}{@{}l@{}}
$\forall\, x \in \Ob\,X$, $X/x$ and\\
$X/x \times x$ are irreducible,\\ % [sic] irréductible -> irréductibles, agreement with "sont"
$X$ 0-connected
\end{tabular}\right.$\\
$\Longleftrightarrow$ & $\left\{\begin{tabular}{@{}l@{}}
$X$ 0-connected,\\
the $X/a \times_x b$ nonempty\\
(and even 0-connected),\\
the $X/a \times_{x\times x} b$ 0-connected
\end{tabular}\right.$
\end{tabular}\right.
\;
\left[\begin{tabular}{@{}l@{}}
stable under\\
passage\\
to $X' \longrightarrow X$\\
étale over $X$, if\\
$X'$ 0-connected,\\
\textit{but} if the\\
$X/x$ are cofiltering\\
and $X$ 0-connected,\\
this does\\
\textit{not} imply $X$ cofiltering\\
Ex. $B_G$
\end{tabular}\right]
\arrow[rrrd, Rightarrow, "{\tl{strict\\(e.g. $\bullet\ \bullet$)}}"]
\arrow[rrdd, Rightarrow, "{\tl{strict\\(e.g. $\boldsymbol{\Delta}$)}}"]
\arrow[ddd, Rightarrow, "{\tl{strict (e.g.\\$B_G$)}}"']
\& \& \& |[draw,rounded corners=6pt]| \varepsilon = \bullet\ \bullet \\
\& \& \&
\begin{tabular}{@{}l@{}}
\begin{tabular}{@{}c@{\ }l@{}}
 & Top $X$ locally\\
 & totally 0-connected\\
$\Longleftrightarrow$ & the connected components of $X$\\
 & are cofiltering\\
$\Longleftrightarrow$ & $\PS^*1 + \PS^*2$\\
 & ($X$ \textit{locally cofiltering}) % [sic] X (upright roman) -> $X$ (italic), inconsistent with all other occurrences
\end{tabular}\\[2ex]
$\left[\begin{tabular}{@{}l@{}}
Stable under passage to $X' \longrightarrow X$\\
étale, but if the $X'/x$\\
satisfy it, not necessarily $X'$!
\end{tabular}\right]$
\end{tabular}
\arrow[d, Rightarrow] \\
\& \&
\begin{tabular}{@{}l@{}}
\begin{tabular}{@{}c@{\ }l@{}}
 & Top $X$ totally 0-connected\\
 & (we say $X$ \textit{totally} 0-\textit{connected})\\
$\Longleftrightarrow$ & $X^o \longrightarrow X^o \times X^o$ is cofinal\\
 & \textit{and} $X \neq \emptyset$\\
$\Longleftrightarrow$ & $\forall\, x, y \in \Ob\,X$, $x \times y$ 0-connected in $X^\wedge$\\
 & i.e. $X/x \times y$ 0-connected\\
 & \textit{and} $X \neq \emptyset$
\end{tabular}\\[2ex]
$\left[\begin{tabular}{@{}l@{}}
NB is not stable under passage to the $X/x$ (\footnotemark[66]), and if\\
the $X/x$ are totally 0-connected and $X \neq \emptyset$,\\
it does not follow that\\
$X$ is so (\footnotemark[67]).
\end{tabular}\right]$
\end{tabular}
\arrow[dddl, Rightarrow, "{\tl{strict (e.g. $B_G$, or\\$0 \rightrightarrows 1$)}}"]
\&
\tl{Top $X$\\locally irreducible} \\
\begin{tabular}{@{}c@{\ }l@{}}
 & Top $X$ locally\\
 & irreducible\\
 & $+$ 0-connected\\
$\Longleftrightarrow$ & $X$ \textit{pseudo-filtering}\\
 & $+$ 0-\textit{connected,}\\
 & i.e. $\PS^*1$\\
 & $+$ 0-connected
\end{tabular}
\arrow[ddr, Rightarrow, "{\tl{strict\\(e.g. $0 \rightrightarrows 1$)}}"']
\&
\left[\begin{tabular}{@{}l@{}}
$\PS^*1$ depends\\
only on the induced\\
categories $X/x$,\\
hence stable under passage\\
to étale $X' \longrightarrow X$,\\
hence the condition\\
$\PS^*1 +{}$0-connected stable\\
under passage to\\
coverings\\
$X' \longrightarrow X$ 0-\textit{connected.}
\end{tabular}\right]
\& \& \\
|[draw,rounded corners=6pt]| B_G \& \& \& \\
\&
\begin{tabular}{@{}c@{\ }l@{}}
 & Top $X$ irreducible\\
 & ($X$ \textit{irreducible})\\
$\Longleftrightarrow$ & $\PS^*1'$
\end{tabular}
\& \& \\
|[draw,rounded corners=6pt]| 0 \rightrightarrows 1
\&
\left[\begin{tabular}{@{}l@{}}
depends only \texttt{[on]} $\mathrm{Ord}\,X$ or on $\mathrm{Esp}(X)$,\\
means that $\mathrm{Ord}\,X$ cofiltering,\\
or $\mathrm{Esp}(X)$ irreducible.\\
Is not stable under passage to\\
$X' \longrightarrow X$ étale 0-connected,\\
e.g. by passage to an $X/x$,\\
e.g. $X = (0 \rightrightarrows 1)$ is irreducible,\\
but $X/1 = \left(\begin{tabular}{@{}l@{}}$\cdot$\\ $\;\searrow$\\ $\qquad\cdot$\\ $\;\nearrow$\\ $\cdot$\end{tabular}\right)$ is\\
not.
\end{tabular}\right]
\& \& \\
\end{tikzcd}
}
\DERmark{ENV0169}

\end{center}
% footnote marks (66), (67) are set inside the diagram; texts below
\footnotetext[66]{\textbf{Ex.} \circled{$\boldsymbol{\Delta}$}.}
\footnotetext[67]{\textbf{Ex.} $B_G$.}
% foot: Version 8 juillet 2001

%======================================================================
\par
\printedpage{61}
% head: A. Grothendieck : Dérivateurs, Ch. II, Édition des Universités de Montpellier II et Paris VII

\msp{page 98}

\begin{center}
{\small
\begin{tikzcd}[ampersand replacement=\&, column sep=small, row sep=large]
\& \tl{Top $X$ locally totally 0-connected\\$+$ 0-connected\\($X$ \textit{cofiltering})}
\arrow[ddl, Rightarrow] \arrow[ddr, Rightarrow] \arrow[rrdd, Rightarrow]
\& \& \\
\& \& \& |[draw,rounded corners=6pt]| \varepsilon = \bullet\ \bullet \\
\tl{Top $X$ locally irreducible\\$+$ 0-connected\\($X$ \textit{pseudo-filtering} and 0-\textit{connected})\\$\PS^*1 +{}$ 0-connected}
\arrow[ddr, Rightarrow]
\& \&
\tl{Top $X$ totally 0-connected\\($X$ \textit{totally} 0-\textit{connected})\\$X\neq\emptyset$ and $X^o \longrightarrow X^o \times X^o$ cofinal}
\arrow[ddl, Rightarrow]
\&
\tl{Top $X$ locally totally\\0-connected\\($X$ \textit{locally cofiltering})\\$\PS^*1 + \PS^*2$}
\arrow[dd, Rightarrow] \\
|[draw,rounded corners=6pt]| B_G \& \& |[draw,circle,inner sep=1pt]| \boldsymbol{\Delta} \& \\
\& \tl{Top $X$ irreducible\\($X$ \textit{irreducible})\\$\PS^*1'$}
\& \&
\tl{Top $X$ locally irreducible\\($X$ \textit{pseudo-cofiltering} [\,?])\\$\PS^*1$} \\
\& |[draw,rounded corners=6pt]| 0 \rightrightarrows 1 \& \&
|[draw,rounded corners=12pt]| \begin{tabular}{@{}c@{}}$B_G \amalg e$\\with $G \neq 1$\end{tabular} \\
\end{tikzcd}
}
\DERmark{ENV0170}

\end{center}

\begin{center}
{\small
\begin{tikzcd}[ampersand replacement=\&, column sep=small, row sep=large]
\text{(\footnotemark[68]).}
\& \tl{$X$ \textit{cofiltering}\\$\PS^*1' + \PS^*2$\\$= \PS^*1 + \PS^*2 +{}$ 0-connected}
\arrow[ddl, Rightarrow] \arrow[ddr, Rightarrow] \arrow[rrdd, Rightarrow]
\& \& \\
\& \& \& |[draw,rounded corners=6pt]| \bullet\ \bullet \\
\tl{$X$ \textit{pseudo-cofiltering} and 0-\textit{connected}\\($\PS^*1 +{}$ 0-connected)}
\arrow[ddr, Rightarrow] \arrow[ddrrr, Rightarrow]
\& \&
\tl{$X$ \textit{totally} 0-\textit{connected}\\($X\neq\emptyset$ and $X^o \longrightarrow X^o \times X^o$ cofinal)}
\arrow[ddl, Rightarrow]
\&
\tl{$X$ \textit{locally cofiltering}\\$\PS^*1 + \PS^*2$}
\arrow[dd, Rightarrow] \\
|[draw,rounded corners=6pt]| B_G \& \& |[draw,circle,inner sep=1pt]| \boldsymbol{\Delta} \& \\
\& \tl{$X$ \textit{irreducible}\\$\PS^*1'$}
\& \&
\tl{$X$ \textit{pseudo-cofiltering}\\$\PS^*1$} \\
\& |[draw,rounded corners=6pt]| \bullet \rightrightarrows \bullet \& \&
|[draw,rounded corners=6pt]| B_G \amalg e \\
\end{tikzcd}
}
\DERmark{ENV0171}

\end{center}
\footnotetext[68]{\textbf{NB} The five typical categories $B_G$ ($G \neq 1$), $\boldsymbol{\Delta}$, $0 \rightrightarrows 1$, $\varepsilon = \bullet\ \bullet$, $B_G \amalg e$ ($G \neq 1$) are such that they belong to the class $C$ of objects of Cat under consideration, but to none of the five other classes except those which are linked to $C$ by $C \Longrightarrow C'$ on the diagram of implications (or by a composite implication).}
% foot: Version 8 juillet 2001

%======================================================================
\par
\printedpage{62}
% head: A. Grothendieck : Dérivateurs, Ch. II, Édition des Universités de Montpellier II et Paris VII

\noindent\begin{tabular}{@{}p{2.8cm}@{}p{11.5cm}@{}}
$\bullet \rightrightarrows \bullet$ & is neither pseudo-cofiltering, nor totally 0-connected.\\[1ex]
$B_G$ & is neither locally cofiltering, nor totally 0-connected.\\[1ex]
$\boldsymbol{\Delta}$ & is not pseudo-filtering.\\[1ex]
$B_G \amalg e$ & is not locally cofiltering, for $B_G$ is not, nor belongs to one of the 4 classes of connected objects.\\[1ex] % [sic] cofiltante -> cofiltrante, letter r missing
The same exclusion from the four connected classes for & $\varepsilon = \bullet\ \bullet$
\end{tabular}
\DERmark{ENV0172}


\msp{page 99}

\begin{center}
{\small
\begin{tikzcd}[ampersand replacement=\&, column sep=small, row sep=large]
\& \tl{$X$ \textit{cofiltering}\\($\PS^*1' + \PS^*2$\\$\Longleftrightarrow \PS^*1 + \PS^*2$\\$+$ 0-connected)}
\arrow[ddl, Rightarrow, "{\tls{$\PS^*1'$\\($\Longrightarrow$ 0-connected)}}"]
\arrow[dd, Rightarrow, "{\tls{$\PS^*1$\\(the $X/x$ satisfy\\$\PS^*1'$)}}"]
\arrow[ddr, Rightarrow]
\& \\
% layout: the two small labels stand between the arrows; attribution to the arrows as in this layout
\& \& |[draw,rounded corners=6pt]| \varepsilon = \bullet\ \bullet \\
\tl{$X$ totally 0-connected\\($X\neq\emptyset$ and $X^o \longrightarrow X^o \times X^o$ cofinal)}
\arrow[ddd, Rightarrow]
\&
\tl{$X$ \textit{pseudo-cofiltering}\\0-connected\\($\PS^*1 +{}$ 0-connected)\\(missing $\PS^*2$)}
\arrow[dddl, Rightarrow] \arrow[dddr, Rightarrow]
\&
\tl{$X$ \textit{locally cofiltering}\\(connected components cofiltering)\\($\PS^*1 + \PS^*2$)}
\arrow[ddd, Rightarrow] \\
|[draw,circle,inner sep=1pt]| \boldsymbol{\Delta} \& |[draw,rounded corners=6pt]| B_G \& \\
\& G \neq \{1\} \& \\
\tl{$X$ irreducible\\($\PS^*1'$)}
\& \&
\tl{$X$ \textit{pseudo-cofiltering}\\or again \textit{locally irreducible}\\(connected components\\pseudo-cofiltering)\\($\PS^*1$)} \\
|[draw,rounded corners=6pt]| \bullet \rightrightarrows \bullet \& \& |[draw,rounded corners=6pt]| B_G \amalg e \\
\& \& G \neq \{1\} \\
\end{tikzcd}
}
\DERmark{ENV0173}\DERmark{ENV0174}\DERmark{ENV0175}\DERmark{ENV0176}

\end{center}

and dually we have

\begin{center}
{\small
\begin{tikzcd}[ampersand replacement=\&, column sep=small, row sep=huge]
\& \text{$X$ filtering}
\arrow[dl, Rightarrow] \arrow[d, Rightarrow] \arrow[dr, Rightarrow]
\& \\
\text{$X^o$ totally 0-connected}
\arrow[d, Rightarrow]
\&
\tl{$X$ pseudo-filtering\\0-connected}
\arrow[dl, Rightarrow] \arrow[dr, Rightarrow]
\&
\text{$X$ locally filtering}
\arrow[d, Rightarrow] \\
\text{$X$ coirreducible} \& \& \text{$X$ pseudo-filtering} \\
\end{tikzcd}
}
\DERmark{ENV0177}

\end{center}
% foot: Version 8 juillet 2001

%======================================================================
\par
\printedpage{63}
% head: A. Grothendieck : Dérivateurs, Ch. II, Édition des Universités de Montpellier II et Paris VII

so we have a direct terminology except for ``$X^o$ totally 0-connected'' - I don't dare say ``$X$ totally co-connected''!!
% foot: Version 8 juillet 2001


\endgroup
% END UNIT Der-ChII-234/Der-ChII-234_p061-065.tex
\clearpage\setcounter{footnote}{0}\phantomsection\addcontentsline{toc}{section}{Chapter III}\markboth{Chapter III}{Chapter III}
% BEGIN UNIT Der-ChIII-567/Der-ChIII-567_p001-005.tex
\begingroup\makeatletter%
\newcommand{\printedpage}[1]{\marginpar{\footnotesize p.~#1}}%
\newcommand{\msp}[1]{\marginpar{\footnotesize [#1]}}%
\newcommand{\Hot}{\mathrm{Hot}}%
\newcommand{\D}{\mathbf{D}}%
\newcommand{\uHom}{\underline{\mathrm{Hom}}}%
\newcommand{\hm}{\mathrm{hom}}%
\makeatother\let\printedpage\derpp\let\msp\dermsp%
\let\mspin\DERmspin%


% ===================== page 1 (title page) =====================
\printedpage{--}
% no page number printed on this page
% head: (none)

\begin{center}
{\LARGE\sffamily THE DERIVATORS}

\bigskip

\textbf{ALEXANDRE GROTHENDIECK}

\bigskip\bigskip

\textbf{Chapter III}

\medskip

\textbf{External Hom in derivators}
\end{center}

\par

% ===================== page 2 (editorial note) =====================
\printedpage{--}
% no page number printed on this page
% head: (none)

\par\noindent\textbf{Note by the editors of the transcribed edition.}\par
This text was deciphered and transcribed in \LaTeXe{} by M. Künzer. It was edited by M. Künzer, J. Malgoire, and G. Maltsiniotis. The transcription is as faithful as possible to the manuscript. For the few obvious corrections, or rare comments by the editors, as well as for the original numbering of the manuscript pages, typewriter characters \texttt{[typewriter]} in brackets are used. A question mark in brackets means that one is not sure of the preceding word.

This edition is provisional. Remarks, comments and corrections are welcome. Send a message to:

\textbf{[e-mail address removed]}
\begin{flushright}
G. Maltsiniotis\qquad\qquad
\end{flushright}

\par

% ===================== page 3 (printed 1) =====================
\printedpage{1}
% head: A. Grothendieck : Dérivateurs, Ch. III, Édition des Universités de Montpellier II et Paris VII
% foot: Version 2 février 2002

\msp{page 1}

\section*{1\quad Relation of derivators with the category $\Hot$}

Let $\D$ be a derivator, $\mathcal{A}$ its fundamental category,
\[
\mathcal{A} = \D(e).
\]
I propose to define, if it can be done, functors
\[
\left\{
\begin{array}{rccccccl}
\Hot & \times & \mathcal{A} & \xrightarrow{\pi_{\D}} & \mathcal{A} & (\xi, M) & \longmapsto & \xi \times M \ \text{ or } \ \xi \otimes M \quad \text{(depending on the case)}\\
\Hot^o & \times & \mathcal{A} & \xrightarrow{\varpi_{\D}} & \mathcal{A} & (\xi, M) & \longmapsto & \hm(\xi, M).
\end{array}
\right.
\tag*{(1)}
\]
These functors must be interchanged by duality: if one takes the functors $\pi^o$, $\varpi^o$ on the opposite categories, one finds
\[
\left\{
\begin{array}{rcccc}
\Hot & \times & \mathcal{A}^o & \xrightarrow{\varpi^o_{\D}} & \mathcal{A}^o\\
\Hot^o & \times & \mathcal{A}^o & \xrightarrow{\pi^o_{\D}} & \mathcal{A}^o,
\end{array}
\right.
\tag*{(1$'$)}
\]
and one must have
\[
\left\{
\begin{array}{rcl}
\varpi^o_{\D} & = & \pi_{\D^o}\\
\pi^o_{\D} & = & \varpi_{\D^o}.
\end{array}
\right.
\]
To tell the truth, considering (1) as functors
\[
\begin{array}{rcl}
\Hot & \xrightarrow{\pi_{\D}} & \uHom(\mathcal{A}, \mathcal{A})\\
\Hot^o & \xrightarrow{\varpi_{\D}} & \uHom(\mathcal{A}, \mathcal{A}),
\end{array}
\]
one expects these functors to extend to
\[
\left\{
\begin{array}{rcl}
\Hot & \xrightarrow{\pi_{\D}} & \uHom(\D, \D)\\
\Hot^o & \xrightarrow{\varpi_{\D}} & \uHom(\D, \D)
\end{array}
\right.
\tag*{(2)}
\]
\DERmark{ENV0178}

\texttt{[rather twice} $\uHom(\D, \D)$\texttt{]}. This means, essentially, that for a given $\xi \in \mathrm{Ob}\,\Hot$, the operations $\xi \times_{\D} M$ and $\hm_{\D}(\xi, M)$ \texttt{[rather} $\hm_{\D}(\xi, M)$\texttt{]} are defined not only in $\mathcal{A} = \D(e)$, but in all the $\D(I)$, and this in such a way as to commute with the functors $f^*$:
\DERmark{ENV0179}


\msp{page 2}
If
\[
f : I' \longrightarrow I,
\]
\printedpage{2}
% head: A. Grothendieck : Dérivateurs, Ch. III, Édition des Universités de Montpellier II et Paris VII
% foot: Version 2 février 2002
then for $M \in \mathrm{Ob}\,\D(I)$, one must have canonical isomorphisms
\DERmark{ENV0180}

\[
\left\{
\begin{array}{rcl}
f^*(\xi \times_{\D} M) & \simeq & \xi \times_{\D} f^*(M)\\
f^*\hm_{\D}(\xi, M) & \simeq & \hm_{\D}(\xi, f^*(M)),
\end{array}
\right.
\tag*{(3)}
\]
with the usual compatibilities for a composite $fg$, relative to $I'' \xrightarrow{g} I' \xrightarrow{f} I$.

But I prefer to work in $\mathcal{A}$, which should imply the similar constructions in all the $\D(I)$, by applying what has been said to the derivator $\D_I$ if need be.

The fundamental algebraic intuition is this, denoting by
\[
X \longmapsto \mathrm{hot}(X) : (\mathrm{Cat}) \longrightarrow (\Hot)
\]
the localization functor from (Cat) to (Hot), obtained by inverting the arrows in (Cat) which are ``weak equivalences''. Suppose $X$ is in Diag, let
\[
p : X \longrightarrow e
\]
be the functor with values in the one-point category. We want
\[
\left\{
\begin{array}{rcccl}
\mathrm{hot}(X) \times_{\D} M & \simeq & p_!(p^*(M)) & = & H^{\D}_{\bullet}(X, M)\\
\hm_{\D}(\mathrm{hot}(X), M) & \simeq & p_*(p^*(M)) & = & H^{\bullet}_{\D}(X, M)
\end{array}
\right.
\tag*{(4)}
\]
(canonical isomorphisms).

\msp{page 3}
This would give a principle of definition for the operations sought. A difficulty comes from the fact that in general, Diag is not equal to (Cat), and that the homotopy types coming from Diag do not give (even up to isomorphism) \textit{all} homotopy types. Intuitively, one finds only ``small enough'' homotopy types. But one can localize Diag with respect to weak equivalences (it is better to say ``quasi-isomorphisms''), one then finds a category $(\Hot)_0$, and a canonical functor
\[
\Hot_0 \longrightarrow \Hot,
\]
if we are lucky, this functor will be fully faithful, and we shall try to describe the functors (1) or (2) by replacing Hot in them by the full subcategory of ``small'' homotopy types. One can also forget (temporarily!) Hot, and work with $\Hot_0$. This is what we are going to do for the moment.

In order to be able to define the functors
\[
\left\{
\begin{array}{lccccccl}
\Hot_0 & \times & \mathcal{A} & \longrightarrow & \mathcal{A} & (\xi, F) & \longmapsto & \xi \times_{\D} F\\
\Hot_0^o & \times & \mathcal{A} & \longrightarrow & \mathcal{A} & (\xi, F) & \longmapsto & \hm_{\D}(\xi, F),
\end{array}
\right.
\tag*{(5)}
\]
I note that for $M$ \texttt{[rather} $M$\texttt{]} fixed, $X$ variable in Diag, the
\DERmark{ENV0181}

\printedpage{3}
% head: A. Grothendieck : Dérivateurs, Ch. III, Édition des Universités de Montpellier II et Paris VII
% foot: Version 2 février 2002
\msp{page 4}
last members in (4)
\[
H^{\D}_{\bullet}(X, M), \qquad H^{\bullet}_{\D}(X, M)
\tag*{(6)}
\]
are indeed, the one covariant, the other contravariant in $X$. Thus there is, up to unique isomorphism, a single pair of bifunctors (5) satisfying (bifunctorially) the isomorphisms (4), if and only if the expressions (6), considered as functorial in $X$, transform quasi-isomorphisms into isomorphisms. This is therefore equivalent to the validity of the % [sic] "il y a ," -> "il y a,", espace imprimée avant la virgule
\DERmark{ENV0182}


\textbf{Theorem} (?) (\footnote{New axiom?}). \textit{Let $f : X' \longrightarrow X$ be in Diag. If $f$ is a quasi-isomorphism, it is a $\D$-equivalence.}

When $X$ is the one-point category, this is nothing other than axiom Der 3 (p. 67). I do not know whether it is a consequence of axioms Der 1 to Der 5 (d). If this is not the case, one would have to add it to the axioms, I think. I have not even proved that an equivalence of categories in Diag (in the sense of categories!) is a $\D$-equivalence!

Let us note the \textit{adjunction formula}
\[
\mathrm{Hom}(\xi \times_{\D} M, N) \xrightarrow{\sim} \mathrm{Hom}(M, \hm_\D(\xi, N)).
\tag*{(7)}
\] % [sic] \hom_D (D italique maigre) -> \hom_{\mathbf{D}}, le dérivateur est noté en gras partout ailleurs
\DERmark{ENV0183}

If $\xi$ comes from $X \in \mathrm{Diag}$, then this reads
\[
\mathrm{Hom}(p_!p^*M, N) \xrightarrow{\sim} \mathrm{Hom}(M, p_*p^*N),
\tag*{(8)}
\]
\msp{page 5}
now by adjunction both sides are canonically isomorphic to
\[
\mathrm{Hom}_{\D(X)}(p^*M, p^*N).
\]

\textbf{Example.} Taking for $\xi$ the homotopy type of the standard circle (represented e.g. by the ordered set
\[
\begin{tikzcd}[row sep=small, column sep=small]
 & b & \\
0 \arrow[ur] \arrow[dr] & & 0' \arrow[ul] \arrow[dl] \ , \\
 & a &
\end{tikzcd}
\]
where we agree to take 0 as origin, and to orient from 0 toward $a$), the two functors in question are (unless I am mistaken) the functors suspension $\Sigma$ and ``cosuspension'' or ``loop space'' functor $\Omega$. Subject to checking that $\Omega F$ does indeed coincide with what I have described elsewhere, concerning substantial cartesian squares in $\mathcal{A}$.

In $\Hot_0$ we have a product operation, in fact this must be the categorical product in $\Hot_0$. But on the models it is given by
\[
\mathrm{hot}_0(X) \times \mathrm{hot}_0(Y) \ \simeq \ \mathrm{hot}_0(X \times Y)
\tag*{(8)}
\] % [sic] (8) -> (9), le numéro (8) est déjà employé plus haut sur cette page


\endgroup
% END UNIT Der-ChIII-567/Der-ChIII-567_p001-005.tex
% BEGIN UNIT Der-ChIII-567/Der-ChIII-567_p006-010.tex
\begingroup\makeatletter%
\newcommand{\printedpage}[1]{\marginpar{\footnotesize p.~#1}}%
\newcommand{\msp}[1]{\marginpar{\footnotesize [#1]}}%
\newcommand{\D}{\mathbf{D}}%
\newcommand{\Hot}{\mathrm{Hot}}%
\newcommand{\Hom}{\mathrm{Hom}}%
\newcommand{\hm}{\mathrm{hom}}%
\newcommand{\homext}{\mathrm{hom}\,\mathrm{ext}}%
\newcommand{\HOT}{\mathrm{HOT}}%
\newcommand{\Ob}{\mathrm{Ob}\,}%
\newcommand{\Ltens}{\overset{\mathbf{L}}{\otimes}}%
\newcommand{\RHom}{\mathbf{R}\mathrm{Hom}}%
\newcommand{\ab}{\mathrm{ab}}%
\newcommand{\defeq}{\overset{\mathrm{def}}{=}}%
\makeatother\let\printedpage\derpp\let\msp\dermsp%
\let\mspin\DERmspin%

\DERmark{ENV0184}


\printedpage{4}
% head: A. Grothendieck : Dérivateurs, Ch. III, Édition des Universités de Montpellier II et Paris VII
% foot: Version 2 février 2002

\texttt{[twice} (8)\texttt{]}. This being so, we have the canonical isomorphisms, for $\xi,\eta\in\Ob\Hot_0$,
\[
\left\{
\begin{aligned}
\xi\times_{\D}(\eta\times_{\D}M) &\;\simeq\; (\xi\times\eta)\times_{\D}M\\
\hm_{\D}(\xi,\hm_{\D}(\eta,M)) &\;\simeq\; \hm_{\D}(\xi\times\eta,M).
\end{aligned}
\right.
\tag*{(9)}
\]

If $\xi$, $\eta$ are represented by $X$, $Y$, the first formula e.g.
\msp{page 6}
is established by computing the second member as
\[
\varphi_!(M_{X\times Y})\simeq p_!(\;\underbrace{q'_!(M_{X\times Y})}_{\substack{\simeq\ q_!(M_Y)_X\\ \text{since $q$ is $\mathbf{D}$-smooth}}}\;)\simeq\xi\times_{\D}(\eta\times_{\D}M),
\]
O.K.
\[
\begin{tikzcd}
X \ar[d,"p"'] & X\times_e Y \ar[l,"q'"'] \ar[d,"p'"] \ar[dl,dotted,"\varphi"'] \\
e & Y \ar[l,"q"]
\end{tikzcd}
\]

I also feel like proving:

\textbf{Proposition 1.} \textit{Let $\xi\in\Ob\Hot_0$. Then $F\longmapsto\xi\times_{\D}F$, as a morphism of derivators $\D\longrightarrow\D$, is} right exact, \textit{and $F\longmapsto\hm_{\D}(\xi,F)$ is left exact.}

We may assume that $\xi$ is defined by $X$ in Diag. Let $f:I\longrightarrow J$, let $F$ be in $\D(I)$, we have to prove
\[
f_*(\hm_{\D}(\xi,F))\;\simeq\;\hm_{\D}(\xi,f_*(F)).
\]
So let us consider
\[
\begin{tikzcd}
X\times I \ar[r,"f'"] \ar[d,"{p=(x,i)\longmapsto i}"'] & X\times J \ar[d,"{q=(x,j)\longmapsto j}"] \\
I \ar[r,"f"] & J.
\end{tikzcd}
\]

We have
\[
\hm_{\D}(\xi,F)=p_*(p^*F),
\]
so
\[
f_*(\hm_{\D}(\xi,F))=\underbrace{f_*p_*}_{(fp)_*}(p^*(F))=(qf')_*p^*(F)\simeq q_*(f'_*p^*(F)).
\]
But
\[
f'_*p^*(F)\;\simeq\;q^*f_*(F),
\]
the base change formula for the $f_*$, for the base change $q$, which is smooth, since obtained from $X\longrightarrow e$, which is smooth, by base change $J\longrightarrow e$. We find
\printedpage{5}
% head: A. Grothendieck : Dérivateurs, Ch. III, Édition des Universités de Montpellier II et Paris VII
% foot: Version 2 février 2002
\msp{page 7}
therefore
\[
f_*(\hm_{\D}(\xi,F))\simeq q_*q^*(f_*(F))\defeq\hm_{\D}(\xi,f_*(F)),
\]
q.e.d. The statement for $F\longmapsto\xi\times_{\D}F$ is dual to the one we have just proved.

Let us take up again, for $\xi\in\Hot_0$ \texttt{[}$=\Hot_0$\texttt{]}, the adjunction formula
\DERmark{ENV0185}

\[
\Hom(\xi\times_{\D}M,N)\simeq\Hom(M,\hm_{\D}(\xi,N))\defeq\Hom_{\xi}(M,N).
\tag*{$(*)$}
\]

If $\xi$ is defined by $X$ in Diag, the two members of $(*)$ are identified with
\[
\Hom_{\D(X)}(p^*M,p^*N),
\]
where $p:X\longrightarrow e$ is the structural functor.

Let us fix $M$, $N$, and consider the expression $(*)$ as a contrafunctor
\[
(\Hot_0)^o\;\longrightarrow\;\mathrm{Ens}.
\]

It would seem that in all known cases, this functor is ``representable'' by an object, if not of $\Hot_0$ itself, at least an object in Hot. If we cannot define it as an object $H$ of $\Hot_0$, the fact that we have a canonical isomorphism
\[
\Hom_{\xi}(M,N)\;\simeq\Hom_{\Hot}(\xi,H),
\tag*{(10)}
\]
functorial in $\xi\in\Ob\Hot_0$, does not necessarily define \texttt{[it]}
\msp{page 8}
up to unique isomorphism. In fact, when
\[
\D=\D_{\mathcal{M},\Sigma}
\]
is defined by a category of models $(\mathcal{M},\Sigma)$, then there is a natural definition of an operation $\homext$ (``external'' Hom) with values in Hot, namely a functor
\[
\left\{
\begin{aligned}
&\mathcal{A}^o\times\mathcal{A} &&\longrightarrow &&\mathrm{Hot}\\
&(M,N) &&\longmapsto &&\homext(M,N),
\end{aligned}
\right.
\tag*{(11)}
\]
where $\mathcal{A}=\mathcal{M}\Sigma^{-1}$ is the fundamental category of $\D=\D_{\mathcal{M},\Sigma}$. It is true that I have not checked, if e.g. $(\mathcal{M},\Sigma)$ underlies a Quillen model category, so that (by \textsc{Thomason}) $\D$ is indeed a derivator, that this notion of $\homext$ does indeed satisfy a formula of the form (10). Here I am going to posit it [?], if need be, these additional data: specification of (11), and of a functorial isomorphism
\[
\underbrace{\Hom_{\xi}(M,N)}_{\begin{subarray}{l}=\ \Hom(\xi\times_{\D}M,N)\\ =\ \Hom(M,\hm_{\D}(\xi,N))\end{subarray}}
\quad\simeq\;\Hom_{\Hot}(\xi,\homext(M,N)).
\tag*{(12)}
\]

\printedpage{6}
% head: A. Grothendieck : Dérivateurs, Ch. III, Édition des Universités de Montpellier II et Paris VII
% foot: Version 2 février 2002
\msp{page 9}
If now $\eta$ is a second element of $\Hot_0$, we feel like writing
\[
\left\{
\begin{aligned}
\homext(\eta\times_{\D}M,N) &\;\simeq\;\homext(M,\hm_{\D}(\eta,N))\\
&\;\simeq\;\hm_{\HOT}(\eta,\homext(M,N)),
\end{aligned}
\right.
\tag*{(13)}
\]
where in the second member, $\hm_{\HOT}$ is the operation (1) in the case where $\D=\mathrm{HOT}$, and which is also defined in terms of the cartesian product $\xi\times\eta$ in Hot, by
\[
\Hom_{\Hot}(\xi,\hm_{\HOT}(\eta,\zeta))\;\simeq\;\Hom_{\Hot}(\xi\times\eta,\zeta),
\]
in other words, for $\eta$ fixed, the functor
\[
\zeta\;\longmapsto\;\hm_{\HOT}(\eta,\zeta)
\]
is defined as the functor right adjoint to the functor
\[
\xi\;\longmapsto\;\xi\times\eta.
\]

To verify formula (13) heuristically (a verification valid ``in due form'' if the three objects under consideration are in fact defined as objects of $\Hot_0\ \ldots$), we take $\Hom_{\Hot}(\xi,-)$ with values in the three members, for $\xi$ given in $\Hot_0$. For the first, we find
\[
\begin{aligned}
\Hom_{\xi}(\eta\times_{\D}M,N) &\;=\;\Hom(\xi\times_{\D}(\eta\times_{\D}M),N)\\
&\;=\;\Hom((\xi\times\eta)\times_{\D}M,N).
\end{aligned}
\]
For the second
\[
\begin{aligned}
\Hom_{\xi}(M,\hm_{\D}(\eta,N)) &\;=\;\Hom_{\mathcal{A}}(\xi\times_{\D}M,\hm_{\D}(\eta,N))\\
&\;\simeq\;\Hom_{\mathcal{A}}(\eta\times_{\D}(\xi\times_{\D}M),N)\\
&\;\simeq\;\Hom_{\mathcal{A}}((\eta\times\xi)\times_{\D}M,N),
\end{aligned}
\]
\msp{page 10}
and finally for the third
\[
\begin{aligned}
\Hom_{\Hot}(\xi\times\eta,\homext(M,N)) &\;\simeq\;\Hom_{\xi\times\eta}(M,N)\\
&\;\simeq\;\Hom((\xi\times\eta)\times_{\D}M,N),
\end{aligned}
\]
so always the same thing! No doubt about it, this must work.

If we have to give ourselves the $\homext$ (11) as an additional datum, with in addition a trifunctorial isomorphism (12), we must also give ourselves the trifunctorial isomorphisms (13), \textit{``compatible''} with the isomorphism (12) in the sense of the preceding calculation.

\printedpage{7}
% head: A. Grothendieck : Dérivateurs, Ch. III, Édition des Universités de Montpellier II et Paris VII
% foot: Version 2 février 2002

For example, taking $\eta=$ homotopy type of the circle $S^1$, formulas (13) specialize to
\[
\begin{aligned}
\homext(\Sigma M,N) &\;\simeq\;\homext(M,\Omega N)\\
&\;\simeq\;\Omega\homext(M,N)\\
&\;\defeq\;\hm_{\HOT}(S^1,\homext(M,N)).
\end{aligned}
\tag*{(14)}
\]

But this formula in Hot is probably not quite the right one, for the right formalism of $\Omega$, $\Sigma$ is carried out in the categories that are \textit{pointed,} and in particular in the pointed homotopy category $\Hot_{\bullet}$. % [sic] ponctué -> ponctuée, agreement with « catégorie »

I am going to spell out an example, taking the derivator
\[
\mathrm{Dér}(A)
\]
``derived'' from an abelian category $A$. Thus $\D(e)=$ derived category of $A$. To avoid
\msp{page 11}
incidental difficulties here, I am going to assume that $A$ is the category of $k$-modules, $k$ an arbitrary ring.

To every $\xi\in\Ob\Hot$ is associated an ``abelianized''
\[
\xi_{\ab}\defeq H_{\bullet}(\xi)\in\underbrace{D(\mathrm{Ab})}_{\text{or }\mathbf{D}(\mathbf{Z}\text{-Mod})}\;,\qquad\quad\text{derived category of abelian groups,}
\]
an object which has cohomology only in degree $\le 0$. That said, if $M$ is in $D(k\text{-Mod})$, we shall have
\[
\left\{
\begin{aligned}
\xi\times_{\D}M &\;\simeq\;\xi_{\ab}\Ltens_{\mathbf{Z}}M\\
\hm_{\D}(\xi,M) &\;\simeq\;\RHom_{\mathbf{Z}}(\xi_{\ab},M),
\end{aligned}
\right.
\tag*{(15)}
\]
a formula which is in cahoots with the Künneth formula for homology
\[
(\xi\times\eta)_{\ab}\simeq\xi_{\ab}\Ltens_{\mathbf{Z}}\eta_{\ab},
\tag*{(16)}
\]
which again gives a familiar look to formulas (9), which reduce to
\[
\left\{
\begin{aligned}
\xi_{\ab}\Ltens_{\mathbf{Z}}(\eta_{\ab}\Ltens_{\mathbf{Z}}M) &\;\simeq\;(\xi_{\ab}\Ltens_{\mathbf{Z}}\eta_{\ab})\Ltens_{\mathbf{Z}}M\\
\RHom_{\mathbf{Z}}(\xi_{\ab},\RHom_{\mathbf{Z}}(\eta_{\ab},M)) &\;\simeq\;\RHom_{\mathbf{Z}}(\xi_{\ab}\Ltens_{\mathbf{Z}}\eta_{\ab},M).\qquad\text{(\footnotemark[2])}
\end{aligned}
\right.
\]
\footnotetext[2]{Also spell out the adjunction formula (7) (p. 4).}

On the other hand, does a $\homext(M,N)$ indeed exist, if $M$, $N$ in $D(k\text{-Mod})$? We think of $\RHom_k(M,N)$, but to make a homotopy type out of it, we have to ``truncate in degrees $>0$''. We could therefore take
\[
\homext(M,N)\;\simeq\;\tau_{\le 0}(\RHom_k(M,N)).
\]

\printedpage{8}
% head: A. Grothendieck : Dérivateurs, Ch. III, Édition des Universités de Montpellier II et Paris VII
% foot: Version 2 février 2002
\msp{page 12}
Is the fundamental formula (10) then valid? The second member is indeed
\[
\Hom_{\Hot}(\xi,H)\;\simeq\;\Hom_{\mathrm{Hotab}}(\xi_{\ab},H)\qquad\quad\text{(\footnotemark[3])}
\]
\footnotetext[3]{? ? is this true ? ? ?}
for $H$ is an abelian object of Hot (cf. the carpet on abelianization \ldots), from there it is in the bag:
\[
\begin{aligned}
\Hom_{\underbrace{\scriptstyle\mathrm{Hotab}}_{\begin{subarray}{l}\approx\ D(\mathbf{Z}\text{-Mod})^{\le 0},\\ \text{\textit{full} subcategory}\\ \text{of }D(\mathbf{Z}\text{-Mod})\end{subarray}}}(\xi_{\ab},\tau_{\le 0}\RHom_k(M,N)) &\;\simeq\;\Hom_{D(\mathbf{Z}\text{-Mod})}(\xi_{\ab},\RHom_k(M,N))\\
&\;\simeq\;\Hom_{D(k\text{-Mod})}(\xi_{\ab}\times M,N),
\end{aligned}
\]
O.K.

On the other hand, the bifunctor
\[
(M,N)\;\longmapsto\;\homext(M,N)\;=\;\tau_{\le 0}\RHom_k(M,N)
\]
factors canonically through Hotab and a fortiori through $\Hot_{\bullet}$. The ``right'' $\homext(M,N)$ would here rather be the one in Hotab, or even in the category $D^{-}(\mathbf{Z}\text{-Mod})$, obtained from that one ``by making the translation functor invertible''. It is evident, moreover, that here the $\xi_{\ab}\times_{\D}M$, where $\xi=S^1$, is not ``the right'' suspension, which must be the shift of degrees – $\xi_{\ab}$ should have cohomology only in degree $-1$, and not also in degree $0$.


\endgroup
% END UNIT Der-ChIII-567/Der-ChIII-567_p006-010.tex
\clearpage\setcounter{footnote}{0}\phantomsection\addcontentsline{toc}{section}{Chapter IV}\markboth{Chapter IV}{Chapter IV}
% BEGIN UNIT Der-ChIV-567/Der-ChIV-567_p001-005.tex
\begingroup\makeatletter%
\newcommand{\printedpage}[1]{\marginpar{\footnotesize p.~#1}}%
\newcommand{\msp}[1]{\marginpar{\footnotesize [#1]}}%
\newcommand{\D}{\mathbf{D}}%
\newcommand{\cA}{\mathcal{A}}%
\newcommand{\bDelta}{\boldsymbol{\Delta}}%
\newcommand{\uHom}{\underline{\mathrm{Hom}}}%
\makeatother\let\printedpage\derpp\let\msp\dermsp%
\let\mspin\DERmspin%


\printedpage{}% no page number printed on this page
% head: (none)

\begin{center}
{\LARGE\textsf{THE DERIVATORS}}

\bigskip

\textbf{ALEXANDRE GROTHENDIECK}

\bigskip\bigskip

\textbf{Chapter IV}

\medskip

\textbf{Substantial diagrams}
\end{center}

\par

\printedpage{}% no page number printed on this page
% head: (none)

\noindent \par\noindent\textbf{Note by the editors of the transcribed edition.}\par
This text was deciphered and transcribed in \LaTeXe\ by M. Künzer. It was edited by M. Künzer, J. Malgoire, and G. Maltsiniotis. The transcription is as faithful as possible to the manuscript. For the few obvious corrections, or rare comments by the editors, as well as for the original numbering of the manuscript pages, typewriter characters \texttt{[typewriter]} in brackets are used. A question mark in brackets means that one is not sure of the preceding word.

\noindent This edition is provisional. Remarks, comments and corrections are welcome. Send a message to:

\noindent\textbf{[e-mail address removed]}

\hfill G. Maltsiniotis\hspace*{2cm}

\par

\printedpage{1}
% head: A. Grothendieck : Dérivateurs, Ch. IV, Édition des Universités de Montpellier II et Paris VII
% foot: Version 2 février 2002

\noindent\texttt{[There are no pages 1-72]}

\msp{page 73}

\subsection*{Construction of substantial diagrams for a derivator.}

\noindent Let $\D$ be a prederivator, $I$ in Diag, let us consider

\noindent $\varphi_I : \D(I) \longrightarrow \uHom(I^o,\cA)$,\ \ where $\cA=\D(e)$ is the fundamental category of coefficients. % [sic] \mathcal{D}(I) -> \mathbf{D}(I), elsewhere the prederivator is the bold D
\DERmark{ENV0186}


The objects of $\uHom(I^o,\cA)$ are interpreted as \textit{diagrams of type $I^o$} in $\cA$. This language is especially convenient if $I$ is ordered. Let us denote, for $F\in\mathrm{Ob}\,\D(I)$, by $(F_i)_{i\in\mathrm{Ob}\,I}$ the associated $I^o$-diagram in $\cA$. By axiom Der 2, the functor
\[
\varphi : F \longmapsto (F_i)_{i\in\mathrm{Ob}\,I}
\]
is conservative. But it has no desire to be faithful, still less fully faithful. When one knows the diagram $(F_i)_i$, one is far from knowing $F$, not even up to nonunique isomorphism, I presume. It is by no means clear that if $F$, $F'$ in $\D(I)$ are such that the associated diagrams are isomorphic, they are themselves isomorphic.

To insist on this point, the objects of $\D(I)$ will be called the ``\textit{substantial diagrams} of type $I^o$ in $\cA$'' (or $\D$-substantial, if there is a risk of confusion). To support intuition and simplify the explanations, an object $F$ of $\D(I)$ ($\D$-coefficients on $I$) will be represented graphically by the corresponding ``weak'' or ``skeletal'' diagram -- one regards it
\msp{page 74}
as a sort of (pale) ``skeleton'' of the substantial diagram from which it comes. E.g. if $I=\bDelta_1 : 0\longrightarrow 1$ \texttt{[rather $\bDelta_1$]}, one denotes a substantial diagram of type $(\bDelta_1)^o$ \texttt{[rather $(\bDelta_1)^o$]} by the same graphical notation, as the ordinary diagram:
\DERmark{ENV0187}\DERmark{ENV0188}

\[
F_0 \longleftarrow F_1.
\]
If
\[
I = \begin{tikzcd}
a \arrow[r] & c \\
0 \arrow[u] \arrow[r] & b, \arrow[u]
\end{tikzcd}
\]
the substantial diagrams of type $I^o$ are written
\DERmark{ENV0189}

\[
\begin{tikzcd}
F_a \arrow[d] & F_c \arrow[l] \arrow[d] \\
F_0 & F_b. \arrow[l]
\end{tikzcd}
\]

\par

\printedpage{2}
% head: A. Grothendieck : Dérivateurs, Ch. IV, Édition des Universités de Montpellier II et Paris VII
% foot: Version 2 février 2002

It will of course be necessary to specify each time whether this graphical notation designates a skeletal diagram, or a substantial one. Often one starts from a skeleton, and tries to ``flesh it out'' with a little flesh, to make it into a substantial diagram -- if possible, in a canonical way. This is a point of derivatic philosophy essential for practice. One will often proceed step by step, by ``substantifying'' subdiagrams corresponding to a strictly increasing sequence $I_0\subset I_1\subset\cdots\subset I_n=I$ of $I$, as best suits the particular situation.

\msp{page 75}

Here is a typical example, assuming $\D$ a \textit{pointed} derivator (the $\D(I)$ are pointed, and the $f^*$ take zero objects to zero objects). It is a matter of defining the ``cosuspension'' or ``loop space'' functor in $\cA$,
\[
\Omega : \cA \longrightarrow \cA.
\]
It is understood that for $M$ in $\cA$, $\Omega(M)$ must fit into a commutative square
\[
\begin{tikzcd}
e \arrow[d] & \Omega M \arrow[l] \arrow[d] \\
M & e, \arrow[l]
\end{tikzcd}
\tag*{(C)}
\]
where $e$ denotes the unit elements. This square however should be substantial, more precisely, be a substantial ``$\D$-\textit{cartesian}'' square. It being understood that an object of $\D(I)$ (substantial square in $\cA$) is called $\D$-cartesian, if it is in the essential image of $\D(I_1)$ under $i_{1*}$, where
\[
i_1 : I_1 \hookrightarrow I,
\]
$I_1$ being the open subcategory of $I$ defined by
\[
I_1 = \begin{tikzcd}
a \\
0 \arrow[u] \arrow[r] & b.
\end{tikzcd}
\]

\msp{page 76}

The restriction of the substantial square sought (C) to the subcategory $I_1$ is therefore a substantial diagram with skeleton
\[
\begin{tikzcd}
e \arrow[d] \\
M & e, \arrow[l]
\end{tikzcd}
\tag*{(Sq)}
\]
\printedpage{3}
% head: A. Grothendieck : Dérivateurs, Ch. IV, Édition des Universités de Montpellier II et Paris VII
% foot: Version 2 février 2002
from which it is obtained (up to canonical isomorphism) by applying $i_*$ to the latter, so as to complete it to $\begin{tikzcd}
e \arrow[d] & \Omega M \arrow[l, dotted] \arrow[d, dotted] \\
M & e \arrow[l]
\end{tikzcd}$. It therefore remains to describe a canonical substantiation of the skeleton (Sq), in terms of $M$. This leads us to look at the open subcategory
\[
I_0 = \{0\} :\quad
\begin{tikzcd}
{} \arrow[r, dotted] & {} \\
0 \arrow[u, dotted] \arrow[r, dotted] & {} \arrow[u, dotted]
\end{tikzcd}
\]
of $I_1$ (hence also of $I$) reduced to the initial object, and the inclusion
\[
I_0 \overset{i_0}{\hookrightarrow} I_1,
\]
and to substantify the skeleton (Sq) by $(i_0)_!(M)$.

One can say a posteriori that, if one knew the substantial square, there is, up to unique isomorphism, only one way of obtaining it as an $i_{1*}i_{0!}(?)$, with ? in $\D(I_0)\simeq\cA$, for $i_{0!}$ and $i_{1*}$ being fully faithful, the same holds for their composite.

\msp{page 77}

In fact, one sees that the functor
\[
i_{1*}i_{0!} :\ \cA\simeq\D(I_0) \longrightarrow \D(I)
\]
induces an equivalence of $\cA$ with the full subcategory of $\D(I)$ formed by the $\D$-cartesian squares whose fibers at $a$ and $b$ are zero. (The obvious quasi-inverse functor is $i^*$, where $i : I_0 \hookrightarrow I$ \texttt{[rather $I_0 \hookrightarrow I$]}.) This therefore means that in this case, the datum of the skeleton of a very particular nature ($F_a\simeq F_b\simeq 0_{\cA}$), plus the additional $\D$-cartesian constraint, on the square sought, determines the latter up to unique isomorphism in terms of its skeleton: there is a substantiation of the latter not only ``canonical'', but even unique up to unique isomorphism, because of the $\D$-cartesian constraint.
\DERmark{ENV0190}


I would like to begin by looking at the still simpler case where $I=\bDelta_1=(0\longrightarrow 1)$, so the skeletons envisaged are simply the arrows
\[
F_1 \longrightarrow F_0 \qquad \text{in } \cA.
\]
The philosophy obliges us, for the most fundamental arrows in the theory of derivators $\D$ (arrows in $\cA$, or by extension in a $\D(I)\simeq\cA_{\D_I}$, fundamental category of the induced derivator $\D_I$) to ask ourselves whether there is not


\endgroup
% END UNIT Der-ChIV-567/Der-ChIV-567_p001-005.tex
% BEGIN UNIT Der-ChIV-567/Der-ChIV-567_p006-009.tex
\begingroup\makeatletter%
\newcommand{\printedpage}[1]{\marginpar{\footnotesize p.~#1}}%
\newcommand{\msp}[1]{\marginpar{\footnotesize [#1]}}%
\newcommand{\Hom}{\operatorname{Hom}}%
\newcommand{\uHom}{\underline{\operatorname{Hom}}}%
\newcommand{\Ob}{\operatorname{Ob}}%
\newcommand{\id}{\mathrm{id}}%
\newcommand{\bDelta}{\boldsymbol{\Delta}}%
\makeatother\let\printedpage\derpp\let\msp\dermsp%
\let\mspin\DERmspin%


\printedpage{4}
% head: A. Grothendieck : Dérivateurs, Ch. IV, Édition des Universités de Montpellier II et Paris VII
% foot: Version 2 février 2002
\msp{page 78}
a ``canonical'', ``natural'', ``willed by God'' way of ``fleshing it out'' into a substantial diagram. The most obvious arrows that present themselves are relative to a functor
\[ f : I \longrightarrow J \]
in Diag, and to the two adjunction morphisms for the pair $(f^*, f_*)$, namely
\[
\begin{aligned}
G &\longrightarrow f_*f^*G && \text{in } D(J)\\
f^*f_*F &\longrightarrow F && \text{in } D(I). \quad (\text{\footnotemark})
\end{aligned}
\]
\footnotetext{Of course, dual problem for $(f_!, f^*)$.}

I am going to try to flesh out these two arrows, and am going to generalize by considering the situation below \texttt{[below]}
\DERmark{ENV0191}

\[
\begin{tikzcd}
I \arrow[d, "f"'] & F \\
J & G, \arrow[u, "u"']
\end{tikzcd}
\]
where
\[ \Hom^f(G,F) \overset{\mathrm{def}}{=} \Hom(f^*G,F) \overset{\text{adj.}}{\simeq} \Hom(G,f_*F). \]
\DERmark{ENV0192}

Thus the same $u$ is interpreted by \textit{two} skeletal arrows
\[
\left\{
\begin{aligned}
f^*(G) &\longrightarrow F && \text{in } \mathbf{D}(I)\\
G &\longrightarrow f_*(F) && \text{in } \mathbf{D}(J).
\end{aligned}
\right.
\tag*{(1)}
\]
To flesh them out is to interpret them (if it can be done) in terms of two objects (canonical if possible) in $\mathbf{D}(I \times \bDelta_1)$ and $\mathbf{D}(J \times \bDelta_1)$ respectively.

I am going to be guided by the special case of a trivial derivator
\[ \mathbf{D} = \mathbf{D}_{\mathcal{M}}, \]
where $\mathcal{M}$ is a ``naked'' category of models $(\Sigma = \emptyset)$.

\msp{page 79}
For fixed $f$, we are going to construct an auxiliary category $K(f)$, in such a way that for every category $\mathcal{M}$, the systems $(G,F,u)$ consisting of an $\mathcal{M}$-presheaf $G$ on $J$, another $F$ on $I$, and an $f$-homomorphism $G \longrightarrow F$, are identified (functorially) with the $\mathcal{M}$-presheaves on
\printedpage{5}
% head: A. Grothendieck : Dérivateurs, Ch. IV, Édition des Universités de Montpellier II et Paris VII
% foot: Version 2 février 2002
$K(f)$. That said, we are going to try to interpret the skeletal arrows (1), but no matter here -- all the $\mathbf{D}$-skeletons are now fleshed out in terms of the geometry of $K(f)$. % [sic] remplumés) -> remplumés, closing parenthesis printed without any opening parenthesis
\DERmark{ENV0193}


\[
\left\{
\begin{aligned}
&\Ob K(f) = \Ob I \amalg \Ob J \qquad \text{\texttt{[rather} $\Ob I \amalg \Ob J$\texttt{ ]}}\\
&\text{$I$, $J$ being identified with full subcategories of $K(f)$,}\\
&\text{there are no arrows from a $j \in J$ to an $i \in I$.}\\
&\text{(Thus $I$ is open in $K(f)$, $J$ its closed complement.)}\\
&\text{For $i \in \Ob I$, $j \in \Ob J$, we set $\Hom_{K(f)}(i,j) = \Hom_J(f(i),j)$.}
\end{aligned}
\right.
\]
\DERmark{ENV0194}


The composition of arrows in $K(f)$ is obvious. We are going to define a canonical retraction
\[ \bar{f} : K(f) \longrightarrow J \]
of $K(f)$ onto its closed subcategory $J$. Thus $\bar{f}$ is known on $\Ob J$, and on $\Ob I$ it will be $f$:
\[
\left\{
\begin{aligned}
\bar{f}(j) &= j && \text{if } j \in \Ob J\\
\bar{f}(i) &= f(i) && \text{if } i \in \Ob I.
\end{aligned}
\right.
\]
It remains to define, for $x, y \in \Ob K(f)$,
\msp{page 80}
\[ \Hom_{K(f)}(x,y) \longrightarrow \Hom_J(\bar{f}(x), \bar{f}(y)). \]
If $x, y \in \Ob J$, it is the identity, if $x, y \in I$, it is $f$. If $x \in I$, $y \in J$, we must define an arrow
\[ \Hom(f(x),y) \longrightarrow \Hom(f(x),y), \]
and we again take the identity. We have
\[
\begin{tikzcd}
J \arrow[r, hook, "j"] \arrow[dr, equal] & K(f) \arrow[d, "\bar{f}"] & I \arrow[l, hook', "i"'] \\
& J &
\end{tikzcd}
\]
\[
\left\{
\begin{aligned}
\bar{f} \circ j &= \id_J\\
j \circ \bar{f} &= \left\{
\begin{aligned}
&j && \text{on } J\\
&jf && \text{on } I,
\end{aligned}
\right.
\end{aligned}
\right.
\qquad \text{i.e.} \quad
\boxed{
\begin{aligned}
j\bar{f}j &= j && \text{on } J\\
j\bar{f}i &= jf && \text{on } I
\end{aligned}
}\ ,
\]
and we have \texttt{[}$K = K(f)$\texttt{]}
\[
\begin{aligned}
\alpha :\ \id_K &\longrightarrow j\bar{f} && \text{i.e. functorial homomorphism}\\
x &\xrightarrow{\alpha(x)} j(\bar{f}(x)), && x \in \Ob K,
\end{aligned}
\]
\printedpage{6}
% head: A. Grothendieck : Dérivateurs, Ch. IV, Édition des Universités de Montpellier II et Paris VII
% foot: Version 2 février 2002
defined as follows:
\[
\left\{
\begin{aligned}
&\text{If } x \in J,\ \alpha(x) = \id_x\ :\ x \xrightarrow{\sim} j(\bar{f}(x)) = j(x) = x\\
&\text{If } x \in I,\ \alpha(x) : \underbrace{x \longrightarrow f(x)}_{\text{in } K(f)},\ \text{or } \underbrace{f(x) \longrightarrow f(x)}_{\text{in } J},\ \alpha(x) = \id_{f(x)}.
\end{aligned}
\right.
\]
We must check functoriality for $x$ varying in $K$ \dots

Let $\Phi$ be an $\mathcal{M}$-presheaf on $K$. It defines
\[
\begin{aligned}
F &= i^*(\Phi) && \in \Ob \mathbf{D}_{\mathcal{M}}(I)\\
G &= j^*(\Phi) && \in \Ob \mathbf{D}_{\mathcal{M}}(J).
\end{aligned}
\]
On the other hand $f^*(G)$ is interpreted as $(j\bar{f})^*(\Phi)|I$, since
\[
\begin{aligned}
f^*(G)(i) &= G(\ \underbrace{f(i)}_{= \bar{f}(i) \,\in\, J}\ )\\
&= j^*(\Phi)(\bar{f}(i))\\
&= \Phi(j\bar{f}(i))\\
&= (j\bar{f})^*(\Phi)(i).
\end{aligned}
\]

\msp{page 81}
Thus
\[ f^*(G) \simeq (j\bar{f})^*(\Phi)|I. \]
Now the arrow
\[ \alpha : \id_K \longrightarrow j\bar{f} \]
gives
\[ \underbrace{\id_K(\Phi)}_{= \Phi} \xleftarrow{\Phi(\alpha)} (j\bar{f})^*(\Phi), \qquad \text{functorial in } \Phi, \]
whence, taking the restrictions to $I$,
\[ i^*(\Phi) \overset{\mathrm{def}}{=} F \xleftarrow{\ u\ } f^*G. \]
\DERmark{ENV0195}


In fact, without too much twisting and turning, we should be able to show that for every category $\mathcal{M}$, the category $\uHom(K^o, \mathcal{M})$ of $\mathcal{M}$-presheaves on $K$ is isomorphic to that of triples $(F,G,u)$, with $F \in \Ob \uHom(I^o, \mathcal{M})$, $G \in \uHom(J^o, \mathcal{M})$, and $u : f^*(G) \longrightarrow F$. Thus, $K$ appears as a solution of a universal problem
\[ \uHom(K^o, \mathcal{M}) \longrightarrow \{F,\ G,\ u \text{ as above}\}(\mathcal{M}). \]

But to specify the arrow $\longrightarrow$ here in a simple conceptual way, it is good to give ourselves $K$ with the additional structures \begin{tikzcd}[cramped, column sep=normal] J \arrow[r, hook, shift left, "j"] & K \arrow[l, shift left, "\bar{f}"] & I \arrow[l, hook', "i"'] \end{tikzcd}, $\bar{f}j = \id_J$, $\id_K \xrightarrow{\alpha} j\bar{f}$. And this system % [sic] I (first object, source of j and target of f-bar) -> J, since j : J -> K(f) and f-bar : K(f) -> J
\DERmark{ENV0196}

\printedpage{7}
% head: A. Grothendieck : Dérivateurs, Ch. IV, Édition des Universités de Montpellier II et Paris VII
% foot: Version 2 février 2002
\msp{page 82}
$\{K,\ i,\ j,\ \bar{f},\ \alpha\}$ is (unless I am mistaken) defined up to canonical isomorphism by the condition of representing the functor $\mathcal{M} \longmapsto \{F,\ G,\ u\}(\mathcal{M})$ \dots
\DERmark{ENV0197}


In particular, the system defined by $G \in \Ob \mathbf{D}_{\mathcal{M}}(J)$
\[ (f^*(G),\ G,\ \id : f^*(G) \longrightarrow f^*(G)) \]
is defined unless I am mistaken by $\bar{f}^*(G)$ on $K(f)$. We propose [?] to find a similar interpretation for
\[ (F,\ f_*(F),\ f^*f_*F \longrightarrow F) \]
$(F \in \Ob \mathbf{D}(I))$. The natural candidate here would be $\Phi = i_*(F)$. Let us grant this point (i.e. especially the verification of the fact that $j^*(i_*(F)) \simeq f_*(F)$). All this shows that if $\mathbf{D}$ is once again an arbitrary derivator, we can, by the preceding formulas, find a canonical object $\Phi$ of $\mathbf{D}(K)$, in terms either of a $\mathbf{D}$-coefficient $F$ on $I$ (taking $\Phi = i_*(F)$), or of a $\mathbf{D}$-coefficient $G$ on $J$ (taking $\Phi = \bar{f}^*(G)$). On the other hand, we can hope to define canonically
\[
\begin{aligned}
\mathbf{D}(K) &\longrightarrow \mathbf{D}(I \times \bDelta_1)\\
\mathbf{D}(K) &\longrightarrow \mathbf{D}(J \times \bDelta_1)
\end{aligned}
\]
whose composites with the ``skeleton'' functors
\[
\begin{aligned}
\mathbf{D}(I \times \bDelta_1) &\longrightarrow \uHom(\bDelta_1^o, \mathbf{D}(I))\\
\mathbf{D}(J \times \bDelta_1) &\longrightarrow \uHom(\bDelta_1^o, \mathbf{D}(J))
\end{aligned}
\]
\msp{page 83}
give back (for $\Phi \in \Ob \mathbf{D}(K)$) the skeletal arrows
\[
\left\{
\begin{aligned}
f^*(G) &\xrightarrow{\ u\ } F\\
G &\xrightarrow{\ v\ } f_*F
\end{aligned}
\right.
\]
associated with $(F,G,u)$. To be precise, $f^*(G) \longrightarrow F$ is defined in terms of $\mathbf{D}$, and $\Phi \in \Ob \mathbf{D}(K)$, as above, and $v$ is deduced from it by adjunction.

The natural idea here, would be to reduce the problem of fleshing out the skeletal adjunction arrows, to the case of open or closed immersions. Note that the open subcategories of $K$ correspond precisely to the functors $K \longrightarrow \bDelta_1$, by taking the inverse images of the generic open. There is therefore a unique % [sic] sqelettiques -> squelettiques, missing u (spelled squelettiques elsewhere in the text)
\[ \varphi :\ K(f) \longrightarrow \bDelta_1 \]
with $I = \varphi^{-1}(\{0\})$. Thus \texttt{[end of the chapter]}


\endgroup
% END UNIT Der-ChIV-567/Der-ChIV-567_p006-009.tex
\clearpage\setcounter{footnote}{0}\phantomsection\addcontentsline{toc}{section}{Chapter V}\markboth{Chapter V}{Chapter V}
% BEGIN UNIT Der-ChV-234/Der-ChV-234_p001-005.tex
\begingroup\makeatletter%
\NoAutoSpacing%
\newcommand{\printedpage}[1]{\marginpar{\footnotesize p.~#1}}%
\newcommand{\msp}[1]{\marginpar{\footnotesize [#1]}}%
\newcommand{\Ob}{\operatorname{Ob}}%
\newcommand{\Fl}{\operatorname{Fl}}%
\newcommand{\uCh}{\underline{\mathrm{Ch}}}%
\newcommand{\uHom}{\underline{\mathrm{Hom}}}%
\newcommand{\bPhi}{\boldsymbol{\Phi}}%
\makeatother\let\printedpage\derpp\let\msp\dermsp%
\let\mspin\DERmspin%


\printedpage{}
% head: (none); no page number printed on this page

\begin{center}
\vspace*{6cm}

{\LARGE\textsf{THE DERIVATORS}}

\vspace{1.5cm}

{\large\textbf{ALEXANDRE GROTHENDIECK}}

\vspace{2.5cm}

{\large\textbf{Chapter V}}

\vspace{0.4cm}

{\large\textbf{Categories of paths and localization}}
\end{center}

\par

\printedpage{}
% head: (none); no page number printed on this page

\vspace*{6cm}

\par\noindent\textbf{Note by the editors of the transcribed edition.}\par
This text was deciphered and transcribed in \LaTeXe{} by M. Künzer. It was edited by M. Künzer, J. Malgoire, and G. Maltsiniotis. The transcription is as faithful as possible to the manuscript. For the few obvious corrections, or rare comments by the editors, as well as for the original numbering of the manuscript pages, the typewriter characters \texttt{[typewriter]} in brackets are used. A question mark in brackets means that one is not sure of the preceding word.

This edition is provisional. Remarks, comments and corrections are welcome. Send a message to:

\textbf{[e-mail address removed]}\\
\hspace*{\fill}G. Maltsiniotis\hspace*{2cm}

\par

\printedpage{1}
% head: \textbf{A. Grothendieck :} \textit{Dérivateurs,} Ch. V, Édition des Universités de Montpellier II et Paris VII
% foot: Version 9 septembre 2001

\msp{page 1}

\section*{1\quad The category of paths}

Let $\Phi$ be the following category:
\[
\left\{
\begin{array}{rcll}
\Ob\Phi & = & \mathbf{Z} & \\
\Fl\Phi & = & \Ob\Phi\sqcup\mathbf{Z} & \qquad (\Ob\Phi\leftrightarrow\text{identities})
\end{array}
\right.
\]
We denote by $\{\alpha\}$, or simply $\alpha$, the object corresponding to $\alpha\in\mathbf{Z}$. We denote by $u_\alpha$ the arrow corresponding to the integer $\alpha\in\mathbf{Z}$.
\[
\begin{array}{rcl@{\qquad}l@{\quad}l@{\quad}l}
s(u_\alpha) & = & \alpha & \text{if } \alpha \text{ even}, & \alpha+1 & \text{if } \alpha \text{ odd}\\[1ex]
b(u_\alpha) & = & \alpha+1 & \text{if } \alpha \text{ even}, & \alpha & \text{if } \alpha \text{ odd}
\end{array}
\]
\DERmark{ENV0222}

\[
\texttt{[rather}\qquad b(u_\alpha) \;=\; \alpha+1 \quad \texttt{if}\ \alpha\ \texttt{even,}\quad \alpha \quad \texttt{if}\ \alpha\ \texttt{odd]}
\]

No compositions to define, for if two arrows are composable, one of them is the identity. In fact, $\Phi$ is an \textit{ordered} category, $\{\alpha\}<\{\beta\}$ if and only if $\alpha$ even, $\beta=\alpha+1$ or $\beta=\alpha-1$ \texttt{[rather} $\beta=\alpha+1$ \texttt{or} $\beta=\alpha-1$\texttt{]}. % [sic] Il n'y pas -> Il n'y a pas, verb "a" missing
\DERmark{ENV0223}\DERmark{ENV0224}


We denote, for $n\in\mathbf{N}$, by
\[
\Phi_n\subseteq\Phi,\quad \Ob\Phi_n=\{\alpha\in\Ob\Phi\mid -n\le\alpha\le+n\},
\]
the full subcategory of $\Phi$ whose objects are the $\{\alpha\}$ with $-n\le\alpha\le+n$. Thus
\[
\Phi_1=\{-1\longleftarrow 0\longrightarrow 1\},
\]
the dual is
\[
\Psi_1=\Phi_1^{o}=\{-1\longrightarrow 0\longleftarrow 1\}.
\]
We can regard $\Phi$ as an amalgamated sum of an infinite ordered sequence (on both sides) of copies of $\Phi_1$, or also of a similar sequence of copies of $\Psi_1$.
\[
\begin{tikzcd}[column sep=small, execute at end picture={
\draw[decorate,decoration={brace,amplitude=4pt}] ([yshift=10pt]\tikzcdmatrixname-1-2.center) -- ([yshift=10pt]\tikzcdmatrixname-1-4.center) node[midway,above=4pt]{$\scriptstyle\Phi_1$};
\draw[decorate,decoration={brace,amplitude=4pt}] ([yshift=10pt]\tikzcdmatrixname-1-4.center) -- ([yshift=10pt]\tikzcdmatrixname-1-6.center) node[midway,above=4pt]{$\scriptstyle\Phi_1$};
\draw[decorate,decoration={brace,amplitude=4pt}] ([yshift=10pt]\tikzcdmatrixname-1-6.center) -- ([yshift=10pt]\tikzcdmatrixname-1-8.center) node[midway,above=4pt]{$\scriptstyle\Phi_1$};
\draw[decorate,decoration={brace,amplitude=4pt}] ([yshift=10pt]\tikzcdmatrixname-1-8.center) -- ([yshift=10pt]\tikzcdmatrixname-1-10.center) node[midway,above=4pt]{$\scriptstyle\Phi_1$};
\draw[decorate,decoration={brace,amplitude=4pt}] ([yshift=10pt]\tikzcdmatrixname-1-10.center) -- ([yshift=10pt]\tikzcdmatrixname-1-12.center) node[midway,above=4pt]{$\scriptstyle\Phi_1$};
\draw[decorate,decoration={brace,mirror,amplitude=4pt}] ([yshift=-10pt]\tikzcdmatrixname-1-3.center) -- ([yshift=-10pt]\tikzcdmatrixname-1-5.center) node[midway,below=4pt]{$\scriptstyle\Psi_1$};
\draw[decorate,decoration={brace,mirror,amplitude=4pt}] ([yshift=-10pt]\tikzcdmatrixname-1-5.center) -- ([yshift=-10pt]\tikzcdmatrixname-1-7.center) node[midway,below=4pt]{$\scriptstyle\Psi_1$};
\draw[decorate,decoration={brace,mirror,amplitude=4pt}] ([yshift=-10pt]\tikzcdmatrixname-1-7.center) -- ([yshift=-10pt]\tikzcdmatrixname-1-9.center) node[midway,below=4pt]{$\scriptstyle\Psi_1$};
\draw[decorate,decoration={brace,mirror,amplitude=4pt}] ([yshift=-10pt]\tikzcdmatrixname-1-9.center) -- ([yshift=-10pt]\tikzcdmatrixname-1-11.center) node[midway,below=4pt]{$\scriptstyle\Psi_1$};
}]
\cdots & {} & -4 \arrow[l] \arrow[r] & -3 & -2 \arrow[l] \arrow[r] & -1 & 0 \arrow[l] \arrow[r] & 1 & 2 \arrow[l] \arrow[r] & 3 & 4 \arrow[l] \arrow[r] & {} & \cdots
\end{tikzcd}
\]

The $\Phi_n$, for $n$ odd, are likewise sums of copies of $\Phi_1$ put end to end, and the $\Phi_n$, for $n$ even $>0$, sums of copies of $\Psi_1$.
\DERmark{ENV0225}


\par

\printedpage{2}
% head: \textbf{A. Grothendieck :} \textit{Dérivateurs,} Ch. V, Édition des Universités de Montpellier II et Paris VII
% foot: Version 9 septembre 2001

The $\Phi_n$ form a direct system with fully faithful transition functors, of which $\Phi$ is the direct limit
\[
\Phi=\varinjlim_n\Phi_n.
\]
But for what follows we are going to make the $\Phi_n$ into an \textit{inverse} system, by defining for every $n\le m$ a \textit{retraction} $\varphi_{n,m}$ of $\Phi_m$ onto $\Phi_n$, by

\msp{page 2}
\[
\begin{array}{rcl}
\varphi_{n,m}:\Phi_m & \longrightarrow & \Phi_n\subset\Phi_m\\[1ex]
\varphi_{n,m}(\{i\}) & = &
\left\{
\begin{array}{l@{\qquad}l@{\quad}l@{\ \ }r}
\{i\} & \text{if}\ i\in\Ob\Phi_n, & \text{i.e} & -n\le i\le n\\[1ex]
\{n\} & & \text{if} & n\le i\le m\\[1ex]
\{-n\} & & \text{if} & -m\le i\le -n
\end{array}
\right.
\end{array}
\]
\DERmark{ENV0226}


\texttt{[rather}
\[
\begin{array}{rcl}
\varphi_{n,m}:\Phi_m & \longrightarrow & \Phi_n\subset\Phi_m\\[1ex]
\varphi_{n,m}(\{i\}) & = &
\left\{
\begin{array}{l@{\qquad}l@{\quad}l@{\ \ }r}
\{i\} & \texttt{if}\ i\in\Ob\Phi_n, & \texttt{i.e} & -n\le i\le n\\[1ex]
\{n\} & & \texttt{if} & n\le i\le m\\[1ex]
\{-n\} & & \texttt{if} & -m\le i\le -n\ \texttt{]}
\end{array}
\right.
\end{array}
\]
We denote the inverse system of ordered sets thus obtained by
\[
\bPhi=(\Phi_n)_{n\in\mathbf{N}}\in\mathrm{Pro}(\mathrm{Cat}).
\]
Let $X$ be a category, then for $n\in\mathbf{N}$
\[
\uCh_n(X)\overset{\mathrm{def}}{=}\uHom(\Phi_n,X)
\]
\DERmark{ENV0227}

is the category of diagrams in $X$ of the form
\[
A_{-n}\cdots\longleftarrow A_{-2}\xrightarrow{u_{-2}}A_{-1}\xleftarrow{u_{-1}}A_0\xrightarrow{u_0}A_1\xleftarrow{u_1}A_2\xrightarrow{u_2}A_3\longleftarrow\cdots A_n
\tag*{$(*)$}
\]
with $2n$ arrows end to end $u_i$ $(-n\le i\le n-1)$ \texttt{[rather (}$-n\le i\le n-1$\texttt{)]}, two consecutive arrows always going in opposite directions,
\DERmark{ENV0228}

\[
\text{either}\ \ \cdot\longrightarrow\cdot\longleftarrow\cdot,\ \ \text{or}\ \ \cdot\longleftarrow\cdot\longrightarrow\cdot\,.
\]
The canonical functor deduced from $\varphi_{n,m}:\Phi_m\longrightarrow\Phi_n$
\[
\varphi_X^{m,n}:\uCh_n(X)\longrightarrow\uCh_m(X),\qquad n\le m,
\]
\DERmark{ENV0229}

\printedpage{3}
% head: \textbf{A. Grothendieck :} \textit{Dérivateurs,} Ch. V, Édition des Universités de Montpellier II et Paris VII
% foot: Version 9 septembre 2001
\texttt{[rather} $\uCh_m(X)$\texttt{]} is the ``degeneracy functor'', consisting in completing a diagram $(*)$, with $m-n$ terms $A_{-n}$ with identities on the left, and with $m-n$ terms \texttt{[}$A_n$\texttt{]} with identities on the right, so for $n$ even we find
\[
\underbrace{A_{-n}\cdots\xrightarrow{\text{id}}A_{-n}}_{m-n\text{ terms}}\xleftarrow{\text{id}}A_{-n}\longrightarrow A_{-(n-1)}\cdots\cdots A_{n-1}\longleftarrow A_n\xrightarrow{\text{id}}\underbrace{A_n\xleftarrow{\text{id}}A_n\longrightarrow\cdots A_n}_{m-n\text{ terms}}.
\]
This is a fully faithful functor (which corresponds to the fact that $\varphi_{n,m}$ is surjective on objects and on arrows). % [sic] flèches. -> flèches)., opening parenthesis never closed
\DERmark{ENV0230}


\msp{page 3}
We set
\[
\underbrace{\uCh_\infty(X)}_{\text{or simply }\uCh(X)}\qquad=\uHom(\bPhi,X)\overset{\mathrm{def}}{=}\varinjlim_n\uCh_n(X)
\]
\DERmark{ENV0231}

(\footnote{The set of objects of this category will be denoted by $\mathrm{Ch}_\infty(X)$ or $\mathrm{Ch}(X)$, likewise we define the $\mathrm{Ch}_n(X)$.}) (\footnote{in this formula: $\bPhi$ procategory}). We have a fully faithful inclusion
\[
\uCh_\infty(X)\lhook\joinrel\longrightarrow\uHom(\Phi,X),
\]
where the second category is that of infinite diagrams (in both directions)
\[
\cdots\longleftarrow A_{-2}\xrightarrow{u_{-2}}A_{-1}\xleftarrow{u_{-1}}A_0\xrightarrow{u_0}A_1\xleftarrow{u_1}A_2\xrightarrow{u_2}\cdots
\]
and where $\uCh_\infty(X)$ is identified with the full subcategory formed by the objects for which the arrows $u_i$ are identities, for $|i|$ sufficiently large. % [sic] formées -> formée, agreement with "la sous-catégorie"

We shall call $\uCh_\infty(X)$ the \textit{category of paths} in $X$, its objects will be the paths in $X$. Those which come from $\uCh_n(X)$ will be called \textit{standard paths of length} $2n$. (The term ``standard'' applies here to our notational conventions. Thus we consider only paths of even length, and for a given length, we choose one of the two reasonable kinds, for example for paths of length 2 we have taken only
\[
A_{-1}\longleftarrow A_0\longrightarrow A_1
\]
and not
\[
A_{-1}\longrightarrow A_0\longleftarrow A_1.)
\]

\msp{page 4}
I realize that my notations $u_i$ are not convenient for duality, I am going to change them, to $u_i$ for the arrows between objects $\{i\}$ with $i\ge0$, and $v_i$ the arrow symmetric to $u_i$ with respect to the origin. (The group of automorphisms of the $\Phi_n$, for $n>0$, is reduced to the group with two elements, the only nontrivial automorphism being the symmetry $\{i\}\longrightarrow\{-i\}$). So our paths will henceforth be denoted % [sic] {-i}. -> {-i}). , opening parenthesis "(Le groupe" not closed on this page
\DERmark{ENV0232}\DERmark{ENV0233}

\[
A_{-2}\xrightarrow{v_1}A_{-1}\xleftarrow{v_0}A_0\xrightarrow{u_0}A_1\xleftarrow{u_1}A_2\longrightarrow\cdots.
\tag*{$(*)$}
\]


\endgroup
% END UNIT Der-ChV-234/Der-ChV-234_p001-005.tex
% BEGIN UNIT Der-ChV-234/Der-ChV-234_p006-010.tex
\begingroup\makeatletter%
\newcommand{\printedpage}[1]{\marginpar{\footnotesize p.~#1}}%
\newcommand{\msp}[1]{\marginpar{\footnotesize [#1]}}%
\newcommand{\uCh}{\underline{\mathrm{Ch}}}%
\newcommand{\uHom}{\underline{\mathrm{Hom}}}%
\newcommand{\uFl}{\underline{\mathrm{Fl}}}%
\newcommand{\Hom}{\mathrm{Hom}}%
\makeatother\let\printedpage\derpp\let\msp\dermsp%
\let\mspin\DERmspin%


\printedpage{4}
% head: A. Grothendieck : Dérivateurs, Ch. V, Édition des Universités de Montpellier II et Paris VII
% foot: Version 9 septembre 2001

(This is not exciting either \dots)

Let
\[
\varepsilon=\{-1,+1\}
\]
be the ordered set, \textit{discrete,} reduced to the two elements $-1$, $+1$, and let us define functors (inclusions and fully faithful for $n>0$)
\DERmark{ENV0234}

\[
\begin{aligned}
\alpha_n:\varepsilon\longrightarrow\Phi_n\quad & \alpha_n(-1)=-n\\
& \alpha_n(1)=n,
\end{aligned}
\]
these functors are compatible with the $\varphi_{n,m}$, whence % [sic] compatible -> compatibles, accord avec « ces foncteurs »
\[
\alpha_\infty:\varepsilon\longrightarrow\boldsymbol{\Phi}\in\mathrm{ProCat},
\]
whence for every $X$ a functor
\[
\kappa:\uHom(\boldsymbol{\Phi},X)\longrightarrow\uHom(\varepsilon,X)\simeq X\times X
\]
i.e.

\msp{page 5}
\[
\begin{aligned}
\uCh(X) &\longrightarrow X\times X\\
c &\longrightarrow (\sigma(c),\beta(c)).
\end{aligned}
\]
\DERmark{ENV0235}


\texttt{[rather} $(\sigma(c),\beta(c))$\texttt{]} For every path $c$ in $X$ i.e. a diagram $(*)$ (satisfying the condition of degeneracy at infinity), we have a source or origin of the path, and a target or an \textit{endpoint}
\[
\left\{
\begin{aligned}
\sigma(c) &= A_i\quad\text{for } i\leq -\mathrm{long}(c)\\
\beta(c) &= A_i\quad\text{for } i\geq \mathrm{long}(c)
\end{aligned}
\right.
\]
\DERmark{ENV0236}\DERmark{ENV0237}

\texttt{[replace} $x$ \texttt{by} $c$\texttt{; for long}$(c)$ \texttt{suitably defined]}. The fiber of the preceding functor at $x$, $y$ will be denoted
\[
\uCh(X;x,y)\qquad\text{or simply }\uCh(x,y),
\]
\textit{category of paths from $x$} to \textit{$y$.}

\printedpage{5}
% head: A. Grothendieck : Dérivateurs, Ch. V, Édition des Universités de Montpellier II et Paris VII
% foot: Version 9 septembre 2001

\section*{2\quad Calculation of a category of fractions $X\Sigma^{-1}$}

Before continuing the formalism of paths (composition etc.), I am going to test the usefulness of the notion of category of paths, for the ``calculation'' of the $X\Sigma^{-1}$, where
\[
\Sigma\subset\mathrm{Fl}(X),\qquad\Sigma\text{ contains the identities}
\]
is a subset, otherwise \textit{arbitrary}, of $\mathrm{Fl}(X)$. We know that $\mathrm{Ob}\,X\Sigma^{-1}=\mathrm{Ob}\,X$, and everything comes down to defining the
\[
\Hom_{X\Sigma^{-1}}(x,y),\qquad x,y\in\mathrm{Ob}\,X,
\]
\msp{page 6}
and for $x,y,z\in\mathrm{Ob}\,X$ the composition
\[
\Hom_\Sigma(x,y)\times\Hom_\Sigma(y,z)\xrightarrow{(u,v)\mapsto vu}\Hom_\Sigma(x,z),
\]
(where we write $\Hom_\Sigma$ instead of $\Hom_{X\Sigma^{-1}}$, to lighten the notation).

Using $\Sigma$, we define a full subcategory
\[
\uCh_\Sigma(X)\hookrightarrow\uCh(X),
\]
(\footnote[3]{The set of its objects is denoted $\mathrm{Ch}_\Sigma(X)$. Its elements are the $\Sigma$-paths \texttt{[from} $x$ \texttt{to} $y$\texttt{?]}.}) by the condition on $c\in\mathrm{Ch}(X)$: For $i$ even, \texttt{[rather even]} i.e. $i-1\xleftarrow{u_i}i$ ($u_i$ is called ``retrograde''), $c(u_i)\in\Sigma$, i.e. the retrograde ``segments'' in the path $c$ must be in $\Sigma$. We similarly define the full subcategories
\DERmark{ENV0238}\DERmark{ENV0239}\DERmark{ENV0240}

\[
\uCh_\Sigma(X;x,y)=\{c\in\uCh(X;x,y)\mid c\in\uCh_\Sigma(X)\},
\]
or again the fiber at $(x,y)$ of the source-target functor for $\uCh_\Sigma(X)$
\[
\uCh_\Sigma(X)\longrightarrow X\times X.
\]

We define a canonical map
\[
\mathrm{Ch}_\Sigma(X,x,y)\longrightarrow\Hom_\Sigma(x,y),
\]
\msp{page 7}
by associating to every path from $x$ to $y$ the $\Sigma$-homomorphism from $x$ to $y$ obtained by \textit{inverting} (in $X\Sigma^{-1}$) the retrograde segments (which is possible, since they are in $\Sigma$), so that we find in $X\Sigma^{-1}$ a ``direct path'' % [sic] `a -> à, accent imprimé comme guillemet ouvrant
\[
x=A_{-n}\longrightarrow\cdots\longrightarrow A_{-1}\longrightarrow A_0\longrightarrow A_1\longrightarrow A_2\cdots A_n=A_{n+1}=\cdots=y.
\]
and we compose all these morphisms.

In fact, we find at once that we thus define a \textit{functor}
\[
\uCh_\Sigma(X)\longrightarrow\uFl(X\Sigma^{-1}),
\]
\texttt{[}$\uFl$ \texttt{category of arrows]} from the fact that if we have a diagram

\printedpage{6}
% head: A. Grothendieck : Dérivateurs, Ch. V, Édition des Universités de Montpellier II et Paris VII
% foot: Version 9 septembre 2001
\[
\begin{tikzcd}
A_{i-1} \arrow[d, "f_{i-1}"] & A_i \arrow[l, "u_i"'] \arrow[d, "f_i"] \\
A'_{i-1} & A'_i \arrow[l, "u'_i"']
\end{tikzcd}
\qquad\qquad (i\text{ even})
\]
with $u_i$, $u'_i$ in $\Sigma$, this diagram, as a diagram in $X\Sigma^{-1}$, is commutative with horizontal arrows \textit{isomorphisms,} and therefore remains commutative on inverting the said arrows. Whence an induced functor
\[
\uCh_\Sigma(X;x,y)\longrightarrow\uFl(\underbrace{X\Sigma^{-1}}_{X'};x,y).
\]

But for any category $X'$, and objects $x$, $y$ in $X'$, the category $\uFl(X',x,y)$ is
\msp{page 8}
a category, \textit{discrete,} having as its set of objects $\Hom_{X'}(x,y)$. Thus the preceding functor in fact defines a canonical map % [sic] précédant -> précédent, adjectif
\[
\boxed{\pi_0\uCh_\Sigma(X;x,y)\longrightarrow\Hom_\Sigma(x,y).}
\]
I say that this map is always bijective. To see this, we are going to \textit{construct} ``from scratch'' the category $X\Sigma^{-1}$, taking as objects those of $X$, and \textit{defining} % [sic] constuire -> construire, lettre omise
\[
\boxed{\Hom_\Sigma(x,y)=\pi_0\uCh_\Sigma(X;x,y).}
\]
We must therefore define the composition
\[
\Hom_\Sigma(x,y)\times\Hom_\Sigma(y,z)\longrightarrow\Hom_\Sigma(x,z).
\]
For this, we are going to define functors ``composition of paths''
\[
\begin{aligned}
\uCh_m(X;x,y)\times\uCh_n(X;y,z) &\longrightarrow \uCh_{m+n}(X;x,z)\\
(c,c') &\longmapsto c'\circ c,
\end{aligned}
\]
but restricting ourselves for simplicity to the case where $m$, $n$ are \textit{even,} so that by ``putting end to end'' $c$ and $c'$ (in this order), we do indeed find a path of type $n+m$.
\DERmark{ENV0241}


\msp{page 9}
This operation is ``transposed'' from the pair of functors
\[
\begin{aligned}
\Phi_m &\xrightarrow{\alpha} \Phi_{m+n}\quad i\longmapsto (i-n)\\
\Phi_n &\xrightarrow{\beta} \Phi_{m+n}\quad i\longmapsto (i+m),
\end{aligned}
\]
\DERmark{ENV0242}

\printedpage{7}
% head: A. Grothendieck : Dérivateurs, Ch. V, Édition des Universités de Montpellier II et Paris VII
% foot: Version 9 septembre 2001
where we will have
\[
\alpha(\{m\})=\beta(\{-n\})=m-n,
\]
\DERmark{ENV0243}

and we will have
\[
\Phi_{m+n}\simeq\Phi_m\sqcup_e\Phi_n,
\]
\DERmark{ENV0244}

i.e. $\Phi_{m+n}$ is the ``amalgamated sum'' in (Cat), obtained by ``gluing'' $\Phi_m$ and $\Phi_n$, by identification of the ``target element'' of $\Phi_m$ with the ``source element'' of $\Phi_n$. The composition of paths functor induces \texttt{[in what precedes, from the beginning of page 9, one must interchange} $m$ \texttt{and} $n$\texttt{]}
\DERmark{ENV0245}

\[
\uCh_{m,\Sigma}(x,y)\times\uCh_{n,\Sigma}(y,z)\longrightarrow\uCh_{m+n,\Sigma}(x,z),
\]
whence on passing to the $\pi_0$
\[
\Hom_{\Sigma,m}(x,y)\times\Hom_{\Sigma,n}(y,z)\longrightarrow\Hom_{\Sigma,m+n}(x,z),
\]
where we denote, for $a,b\in\mathrm{Ob}\,X$, by
\[
\Hom_{\Sigma,n}(a,b)\overset{\mathrm{def}}{=}\pi_0(\uCh_{n,\Sigma}(a,b)).
\]
\DERmark{ENV0246}\DERmark{ENV0247}


\texttt{[rather} $\uCh_{n,\Sigma}$\texttt{]}

\msp{page 10}
To deduce from this a map on the $\varinjlim$
\[
\underbrace{\Hom_\Sigma(x,y)}_{=\ \varinjlim\ \Hom_{\Sigma,m}(x,y)}\quad\times\quad\underbrace{\Hom_\Sigma(y,z)}_{=\ \varinjlim\ \Hom_{\Sigma,n}(y,z)}\quad\longrightarrow\quad\underbrace{\Hom_\Sigma(x,z)}_{=\ \varinjlim\ \Hom_{\Sigma,n+m}(x,z)}
\]
we must prove, for $n\leq n'$, $m\leq m'$, commutativity in
\[
\begin{tikzcd}
\Hom_{\Sigma,m}(x,y)\times\Hom_{\Sigma,n}(y,z) \arrow[r] \arrow[d] & \Hom_{\Sigma,m+n}(x,z) \arrow[d] \\
\Hom_{\Sigma,m'}(x,y)\times\Hom_{\Sigma,n'}(y,z) \arrow[r] & \Hom_{\Sigma,m'+n'}(x,z)
\end{tikzcd}
\]
\DERmark{ENV0248}\DERmark{ENV0249}

\texttt{[rather} $\Hom_{\Sigma,m+n}(x,z)$\texttt{,} $\Hom_{\Sigma,m'+n'}(x,z)$\texttt{]} In other words, that for $c:x\longrightarrow y$ and $c':y\longrightarrow z$ two composable $\Sigma$-paths of length $2m$ and $2n$ respectively, the composite $c'c:x\longrightarrow z$ does not change, ``up to homotopy'' in $\uCh_\Sigma(x,z)$, when we replace $c$, $c'$ by paths $\bar{c}$, $\bar{c}'$ of length $2m'$, $2n'$ deduced from $c$ and $c'$ by ``degeneracy at the ends''.
\DERmark{ENV0250}\DERmark{ENV0251}


More precisely, % [sic] précisement -> précisément, accent omis
\[
\bar{c}'\bar{c}\sim\overline{c'c},
\]
the composite path is homotopic to the one deduced from $c'c$ by ``degeneracy at the ends''.

I have just looked at $n=m=1$, $n'=m'=2$, and found that there is no natural arrow between $\bar{c}'\bar{c}$ and $\overline{c'c}$ - the homotopy sought must
\printedpage{8}
% head: A. Grothendieck : Dérivateurs, Ch. V, Édition des Universités de Montpellier II et Paris VII
% foot: Version 9 septembre 2001
\msp{page 11}
be defined with the help of a ``path'' in $\uCh_{\Sigma,m'+n'}(X;x,z)$. In the general case, we are in fact going to find a homotopy
\DERmark{ENV0252}

\[
h=h_{(n,m),(n',m')}
\]
``making commutative'' the diagram of functors
\[
\begin{tikzcd}[column sep=large]
\mathcal{A}=\uCh_{\Sigma,m}(X;x,y)\times\uCh_{\Sigma,n}(X;y,z) \arrow[r, "*_{m,n}"] \arrow[d, "i_{m,m'}\times j_{n,n'}"] & \uCh_{\Sigma,m+n}(X;x,z) \arrow[d, "k_{m+n,m'+n'}", ""'{name=K}] \\
\uCh_{\Sigma,m'}(X;x,y)\times\uCh_{\Sigma,n'}(X;y,z) \arrow[r, "*_{m',n'}", ""{name=S}] & \uCh_{\Sigma,m'+n'}(X;x,z)=\mathcal{B},
\arrow[from=K, to=S, bend right, "h"']
\end{tikzcd}
\]
where the $i_{\alpha,\beta}$, $j_{\alpha,\beta}$, $k_{\alpha,\beta}$ are the transition functors
\[
\uCh_{\Sigma,\alpha}(X;a,b)\longrightarrow\uCh_{\Sigma,\beta}(X;a,b)
\]
$(\alpha\leq\beta)$ for the $(x,y)$, $(y,z)$ and $(x,z)$ (in place of $(a,b)$) respectively. The homotopy $h$ will be a path in the category $\uHom(\mathcal{A},\mathcal{B})$, going from
\[
k_{m+n,m'+n'}\circ{*}_{m,n}\longrightarrow{*}_{m',n'}\circ(i_{m,m'}\times j_{n,n'}).
\]
\DERmark{ENV0253}\DERmark{ENV0254}

The existence of such a homotopy means neither more nor less than that the square diagram above is commutative, as a diagram in the category (Cat\,hot) of ``categories up to homotopy,'' whose objects are
\msp{page 12}
the (small) categories, and where we define
\[
\Hom_{\mathrm{Cat\,hot}}(\mathcal{A},\mathcal{B})\overset{\mathrm{def}}{=}\underbrace{\pi_0(\uHom(\mathcal{A},\mathcal{B}))}_{\begin{array}{@{}l@{}}\text{homotopy classes}\\ \text{of functors from }\mathcal{A}\\ \text{to }\mathcal{B}\end{array}}\,.
\]
\DERmark{ENV0255}

It is this conclusion that will be useful to us, and not the explicit description of the path $h$. This existence granted (cf. the proof below), we see therefore that if we regard, for $a,b\in\mathrm{Ob}\,X$, the system of categories $\uCh_{\Sigma,r}(X;a,b)$ $(r\in\mathbf{N})$ as an \textit{ind-object of} Cat\,hot\textit{,} which we denote by
\[
\uCh^{\mathrm{ind}}_{\Sigma}\,(X;a,b)\in\mathrm{Ob}\,\mathrm{Ind}(\mathrm{Cat\,hot}),
\]
we will have a canonical morphism in Ind(Cat\,hot)
\[
\uCh^{\mathrm{ind}}_{\Sigma}\,(X;x,y)\times\uCh^{\mathrm{ind}}_{\Sigma}\,(X;y,z)\xrightarrow{\ *\ }\uCh^{\mathrm{ind}}_{\Sigma}\,(X;x,z),
\]
called \textit{``ind-morphism of composition of paths.''} We would have liked to deduce it from an honest \textit{functor} ``composition of paths''
\[
\begin{aligned}
\uCh_\Sigma(X;x,y)\times\uCh_\Sigma(X;y,z) &\xrightarrow{\ ?\ } \uCh_\Sigma(X;x,z),\\
\text{which we would write }(c,c') &\longmapsto c'\circ c\text{ or }c'\cdot c,
\end{aligned}
\]
but I doubt that such a functor exists,


\endgroup
% END UNIT Der-ChV-234/Der-ChV-234_p006-010.tex
% BEGIN UNIT Der-ChV-234/Der-ChV-234_p011-015.tex
\begingroup\makeatletter%
\newcommand{\printedpage}[1]{\marginpar{\footnotesize p.~#1}}%
\newcommand{\msp}[1]{\marginpar{\footnotesize [#1]}}%
\newcommand{\Ch}{\underline{\mathrm{Ch}}}%
\newcommand{\Hom}{\underline{\mathrm{Hom}}}%
\newcommand{\Fl}{\mathrm{Fl}}%
\newcommand{\card}{\mathrm{card}}%
\newcommand{\uHom}{\underline{\mathrm{Hom}}}%
\makeatother\let\printedpage\derpp\let\msp\dermsp%
\let\mspin\DERmspin%


\printedpage{9}
% head: A. Grothendieck : Dérivateurs, Ch. V, Édition des Universités de Montpellier II et Paris VII

\msp{page 13}
which on passing to ind-objects gives back the preceding morphism, or, what amounts to the same thing, such that for every $m,n\in\mathbf{N}$ (of the same parity), there exists an $N\geq m+n$ large enough that we have a diagram
\[
\begin{tikzcd}
\Ch_{\Sigma,m}(X;x,y)\times\Ch_{\Sigma,n}(X;y,z) \arrow[r,"*_{m,n}"] \arrow[dd,hook] \arrow[dr] & \Ch_{\Sigma,m+n}(X;x,z) \arrow[d,hook] \\
 & \Ch_{\Sigma,N}(X;x,z) \arrow[d,hook] \\
\Ch_{\Sigma}(X;x,y)\times\Ch_{\Sigma}(X;y,z) \arrow[r,"?"] & \Ch_{\Sigma}(X;x,z)
\end{tikzcd}
\]
and where the upper triangle is \textit{commutative in} Cat\,hot. (One must suppose that for $N$ large enough, the composite
\begin{tikzcd}[cramped, sep=small]
{} \arrow[d,hook] & \\
{} \arrow[r,"?"] & {}
\end{tikzcd}
factors through $\Ch_{\Sigma,N}(X;x,z)$ \dots)
To see that the composition functors $*_{m,n}$ define a homomorphism of the ind-objects in (Cat\,hot), it suffices to prove it in the particular cases $(m',n')=(m+2,n)$ ou $(m',n')=(m,n+2)$. We are going to draw the figure assuming $m$, $n$ \textit{odd.}
\DERmark{ENV0256}


I realize here that I have stupidly botched it, the composition of two paths of lengths \textit{$m$, $n$ odd is}

\msp{page 14}
\textit{not a path} of length $m+n$, i.e. a diagram of type $\Phi_{m+n}$, but a diagram of \textit{dual} type $\Phi_{m+n}^o=\Psi_{m+n}$. For example, for $m=n=1$
\DERmark{ENV0257}

\[
\underbrace{A_0\longleftarrow A_1\longrightarrow A_2}_{c}=\underbrace{B_0\longleftarrow B_1\longrightarrow B_2}_{c'}
\]
is of type $\Phi_2^o$! So as not to complicate life too much, we are going to restrict ourselves in the preceding developments to $m$, $n$ \textit{even,} hence paths of length at least 4. It is a little ridiculous! We are reduced to proving commutativity up to homotopy in the case where we leave $c'$ fixed and where we send $\Ch_{\Sigma,m}(X;x,y)$ into $\Ch_{\Sigma,m+2}(X;x,y)$, or conversely $c$ fixed and $\Ch_{\Sigma,n}(X;y,z)$ into $\Ch_{\Sigma,n+2}(X;y,z)$ \texttt{[read:}$\Ch_{\Sigma,n+2}(X;y,z)$\texttt{]}. Let us take the first case, and write a path $c$ without negative indices (this is not a very happy notation!)
\DERmark{ENV0258}

\[
A_0\longrightarrow A_1\longleftarrow A_2\cdots\longleftarrow A_{2m}
\]
(NB because of the condition $m$ even, all our paths now begin with an arrow from $A_0$ to $A_1$, never in the opposite direction).
\DERmark{ENV0259}


I see that making the homotopy explicit, while offering no difficulty, is very tedious to explain. I am going to try to give a conceptual proof, without tedious ``calculation''.

% foot: Version 9 septembre 2001

\printedpage{10}
% head: A. Grothendieck : Dérivateurs, Ch. V, Édition des Universités de Montpellier II et Paris VII

\msp{page 15}
The two composite functors
\[
f=k\circ *_{m,n}\ ,\quad g=*_{m',n'}\circ(i\times j)
\]
in the square on page 11, which we have to prove are homotopic in $\uHom(\mathcal{A},\mathcal{B})$, can be viewed as obtained in the following way, in terms of two functors
\[
\varphi,\psi:\Phi_{m'+n'}\rightrightarrows\Phi_{m+n}, \tag*{(1)}
\]
whence two ``transposed'' functors
\[
\begin{tikzcd}
\uHom(\Phi_{m+n},X) \arrow[r,shift left,"F"] \arrow[r,shift right,"G"'] & \uHom(\Phi_{m'+n'},X)
\end{tikzcd} \tag*{(2)}
\]
\DERmark{ENV0260}\DERmark{ENV0261}

\texttt{[rather} $\Phi_{m+n}$\texttt{,} $\Phi_{m'+n'}$\texttt{]} which functors ``induce'' $f$ and $g$. To be precise, denoting by $a$, $b$ the source and target objects of $\Phi_{m+n}$, by $c$ the object which corresponds to the target of the path $c$ of length $m$ \texttt{[too many} $c$\texttt{'s]}, to the source of the path $c'$ of length $n$ which is to be composed, finally by $a_0$, $b_0$ the source and target elements of $\Phi_{m'+n'}$, the functors under consideration (1) send $a_0$ to $a$, $b_0$ to $b$, hence define a commutative diagram
\[
\begin{tikzcd}
\Phi_{m'+n'} \arrow[r,shift left,"\varphi"] \arrow[r,shift right,"\psi"'] & \Phi_{m+n} \\
\underset{\text{\footnotesize discrete category}}{\varepsilon=\{a_0,b_0\}.} \arrow[u,hook,"\alpha'"] \arrow[ur,hook,"\alpha"] &
\end{tikzcd} \tag*{(3)}
\]
This said, the target category $\mathcal{B}$ for $f$, $g$, namely $\Ch_{\Sigma,m'+n'}(X;x,z)$, is interpreted as the full subcategory of $\uHom(\Phi_{m'+n'},X)$ consisting
\DERmark{ENV0262}


\msp{page 16}
of the functors $c':\Phi_{m'+n'}\longrightarrow X$ whose composite with $\alpha'$ is the functor
\[
\beta:\varepsilon\longrightarrow X\qquad\qquad
\begin{aligned}
a_0&\longrightarrow x\\
b_0&\longrightarrow z,
\end{aligned} \tag*{$(*)$}
\]
and likewise the source category $\mathcal{A}$ is interpreted as the full subcategory of $\uHom(\Phi_{m+n},X)$ consisting of the $c:\Phi_{m+n}\longrightarrow X$ whose restriction to $\{a,b,c\}$ is the functor
\[
a\longrightarrow x,\ b\longrightarrow z,\ c\longrightarrow y.
\]
\DERmark{ENV0263}

\texttt{[rather} $b\longrightarrow z$\texttt{,} $c\longrightarrow y$\texttt{]} It is contained in the full subcategory $\mathcal{A}'$ of functors whose restriction to $\{a,b\}$ is $a\longrightarrow x,\ b\longrightarrow z$, i.e. whose composite with $\alpha$ is the functor $\beta:\varepsilon\longrightarrow X$ already considered above. Thus the functors $f,g:\mathcal{A}\rightrightarrows\mathcal{B}$ factor through $\mathcal{A}'$
% foot: Version 9 septembre 2001
\printedpage{11}
% head: A. Grothendieck : Dérivateurs, Ch. V, Édition des Universités de Montpellier II et Paris VII
\[
\begin{tikzcd}
\mathcal{A} \arrow[r,hook] \arrow[dr,shift left,"f"] \arrow[dr,shift right,"g"'] & \mathcal{A}' \arrow[d,shift right,"g'"'] \arrow[d,shift left,"f'"] \\
 & \mathcal{B},
\end{tikzcd}
\]
functors deduced from $\varphi$, $\psi$ above, giving rise to the commutative triangles (3), and from the functor $\beta:\varepsilon\longrightarrow X$, in the said way. To prove that $f$ and $g$ are homotopic, it suffices to prove it for $f'$, $g'$. And for this, it suffices to show that $\varphi$ and $\psi$ are homotopic as objects of the category $\uHom(\Phi_{m'+n'},\Phi_{m+n};\alpha')$ consisting of the functors whose restriction to $\varepsilon$ is $\alpha$ \texttt{[rather of the category} $\uHom(\Phi_{m'+n'},\Phi_{m+n};\alpha')$ \texttt{consisting of the functors whose restriction to} $\varepsilon$ \texttt{is} $\alpha'$ \texttt{]}.
\DERmark{ENV0264}


\msp{page 17}
Moreover, the two assertions (on $f'$, $g'$, for arbitrary $X$, $x$, $z$, \texttt{[rather} $x$\texttt{,} $z$\texttt{]} and on $\varphi$, $\psi$, $\alpha'$ \texttt{[rather} $\alpha'$\texttt{]}) are in fact equivalent as one sees by taking $X=\Phi_{m+n}$, $x=a$, $z=b$ \texttt{[rather} $x=a$\texttt{,} $z=b$\texttt{]}. This will be contained in the more general lemma. (\footnote[4]{In the considerations that follow, we in fact tacitly take the functors under consideration to be \textit{increasing} on the vertices with respect to the order of their numbering. Thus one should write $\uHom_{cr}$ instead of $\uHom$, in the statements that follow, including in the connectedness lemma opposite.})
\DERmark{ENV0265}\DERmark{ENV0266}\DERmark{ENV0267}


\textbf{Lemma.} \textit{Let $k$ and $l$ be two integers. Then the full subcategory of $\uHom_{cr}(\Phi_k,\Phi_l)$ consisting of the functors $\Phi_k\longrightarrow\Phi_l$ sending source to source, target to target, is connected.} (NB It is nonempty if and only if $k\geq l$.)
\DERmark{ENV0268}


We are going to describe these functors completely (= increasing maps) $\Phi_k\xrightarrow{\varphi}\Phi_l$, identifying $\Fl(\Phi_k)$ \texttt{[rather the set of nonidentity arrows of} $\Phi_k$\texttt{]} with the interval of integers $[1,2k]$, and associating to $\varphi$ the subset
\[
E_\varphi\subset\Fl(\Phi_k)=[1,2k]\qquad\qquad E_\varphi=\{u\in\Fl(\Phi_k)\mid\varphi(u)\text{ an identity}\}
\]
consisting of the $u\in\Fl(\Phi_k)$ which are ``contracted'' by $\varphi$ to an identity. If $E$ is a set of integers, we introduce on it the notion of connected components (by the equivalence relation in $E$ generated by the successor relation $y=x+1$). One sees at once that the connected components of $E_\varphi$ which contain neither the first nor the last arrow of $\Phi_k$ are even, and that % [sic] connnexes -> connexes, triple n
\[
\underbrace{\card\Fl(\Phi_k)}_{2k}=\card\,(E_\varphi)+\underbrace{\card\Fl(\Phi_l)}_{2l},
\]

\msp{page 18}
i.e.
\[
\card\,(E_\varphi)=2(k-l).
\]

% foot: Version 9 septembre 2001

\printedpage{12}
% head: A. Grothendieck : Dérivateurs, Ch. V, Édition des Universités de Montpellier II et Paris VII

It is immediate that a functor $\varphi$ is surjective on objects and on arrows, hence identifies $\Phi_l$ with a quotient category of $\Phi_k$ - and the equivalence relation in $\Phi_k$ defined by $\varphi$ is the one obtained by identifying all the objects that are source or target of arrows of one and the same connected component $I$ of $E_\varphi$
\[
\cdots\longleftarrow\ \underbrace{\longrightarrow\ \longleftarrow\ \longrightarrow\ \longleftarrow}_{I}\ \longrightarrow\cdots
\]
and the arrows in this component with the corresponding identity. The conditions stated on $E$ (independently of a $\varphi$ which would give rise to it) ensure that the quotient category is isomorphic (with \textit{preservation} of source and target) either to $\Phi_l$, or to the dual $\Psi_l=\Phi_l^o$, according to the direction of the first arrow of $\bar{\Phi}_k$ (quotient of $\Phi_k$). This ``initial direction'' is the same as that for $\Phi_k$ if the first arrow of $\Phi_k$ does not belong to $E$, or else if it belongs to it and its connected component in $E$ has even cardinality. In the opposite case ($u\in E$ and its connected component in $E$ has odd cardinality), the initial direction of $\bar{\Phi}_k$ is opposite to that of $\Phi_k$. % [sic] "(indépendamment d'un $\varphi$ qui lui donnerait naissance" -> parenthèse fermante manquante, the opening parenthesis is never closed; [sic] connnexe -> connexe, triple n
\DERmark{ENV0269}


\msp{page 19}
One sees therefore that $E$ must satisfy yet another condition, for the corresponding $\bar{\Phi}_k$ to be isomorphic to $\Phi_l$, i.e. for $E$ to come from a functor $\varphi$. I summarize all the conditions

\begin{itemize}
\item[a)] $\card E=2(k-l)$.
\item[b)] The connected components of $E$ which contain no initial or final arrow have even cardinality.
\item[c)] If $\Phi_k$ and $\Phi_l$ have the same initial direction ($k$ and $l$ of the same parity), then the possible component of $E$ containing the initial arrow must itself also have even cardinality (which, because of a) and b) \texttt{[rather a) and b)]}, is equivalent to the same condition for the possible connected component containing the final arrow of $\Phi_k$). If $\Phi_k$ and $\Phi_l$ have different initial directions, the initial arrow of $\Phi_k$ must belong to a connected component of \textit{odd} cardinality in $E$ (which is equivalent to likewise for the final arrow).
\end{itemize}
\DERmark{ENV0270}


On the other hand, the functor $\varphi$ is uniquely determined by the knowledge of $E$, for there is a \textit{unique} isomorphism $\bar{\Phi}_k\xrightarrow{\sim}\Phi_l$ compatible with sources and targets, because of the

\textbf{Lemma} (immediate) : \textit{An automorphism of $\Phi_l$ which is the identity on the source object (or on the target object) is the identity everywhere.}

The functors $\varphi$, $\psi$ from earlier corresponded to the choice, one of $E_\varphi=\{1,2,2k-1,2k\}$ \texttt{[rather} $E_\varphi=\{1,2,2k-1,2k\}$\texttt{]}, the other of $E_\psi=\{1,2,i,i+1\}$, 2 connected components
\DERmark{ENV0271}


\msp{page 20}
in one case, two \texttt{[rather two]} in the other. % [sic] "l'autre.)" -> parenthèse fermante sans parenthèse ouvrante, no matching opening parenthesis
\DERmark{ENV0272}\DERmark{ENV0273}


% foot: Version 9 septembre 2001

\printedpage{13}
% head: A. Grothendieck : Dérivateurs, Ch. V, Édition des Universités de Montpellier II et Paris VII

In the general case, one should make explicit the order relation on the set $\uHom^{!}(\Phi_k,\Phi_l)$ (! means : compatible with source and target), in terms of the representative sets $E_\varphi$, $E_\psi$. I am not going to make this order relation explicit, but only a property of it which will suffice for us. Let $\varphi_E$ be the functor defined by a subset $E\subset[1,2k]$. First of all :

\textbf{Lemma 1} (immediate) : \textit{Let $E,E'\subset[1,2k]$ be admissible, and let $j\in[1,2k]$ be such that $E\cap[1,j]=E'\cap[1,j]$. Then $\varphi_E$ and $\varphi_{E'}$ agree on the target of $u_j$ ($j$ an arrow of $\Phi_k$) and on all the objects of $\Phi_k$ preceding it (in the sense of the numbering of the objects, and not of the order relation of the category $\Phi_k$ of course.) Dually, if $E\cap[j,2k]=E'\cap[j,2k]$, then $\varphi_E$ and $\varphi_{E'}$ agree on the source of $u_j$, and on all the objects of $\Phi_k$ following it.} % [sic] la source de $y$ -> la source de $u_j$, apparently intended by duality with "le but de $u_j$"
\DERmark{ENV0274}


Let us now consider a connected subset $I$ of $E$ (not necessarily a connected component), $I=[r,s]\subset[1,2k]$ ($r$ the smallest element of $I$, $s$ its largest element) and suppose that

\begin{itemize}
\item[a)] $\card I$ is even (i.e. $s-r$ is odd) \texttt{[rather} $s-r$ \texttt{odd]}
\item[b)] $r\neq1$, and $r-1\notin E$ (i.e. $r-1\in[1,2k]\backslash E$)
\end{itemize}
\DERmark{ENV0275}


and let
\[
E'=(E\backslash\{s\})\cup\{r-1\}=E\backslash I\cup\overbrace{[r-1,s-1]}^{I'}
\]
obtained from $I$ by ``translating'' one notch to the left.

\msp{page 21}
It is immediate that if $E$ is admissible, $E'$ is too (and conversely). This said, we have the

\textbf{Lemma 2.} \textit{$E$ and $E'$ being as above, there exists a homomorphism of functors between $\varphi_E$ and $\varphi_{E'}$} (I am not saying in which direction the arrow goes!). % [sic] une homomorphisme -> un homomorphisme, wrong gender

\textbf{Corollary.} Dual statement, when $E'$ is obtained from $E$ by translating to the right by one notch a connected interval of even length $I=[r,s]$ of $E$, $s\neq2k$ and $s+1\notin E$.

Indeed, in this case $E$ is obtained from $E'$ by an operation of the type made explicit in lemma 2.

The connectedness of the category $\uHom^{!}(\Phi_k,\Phi_l)$ is now clear. If $k$ and $l$ have the same parity, by pushing the connected components of $E$ to the left, we write $E$ in the ``normal'' form
\[
E_0=[1,2(k-l)].
\]
If $k$ and $l$ are of opposite parity, we cannot push the whole connected component of $E$ which contains the last element $2k$, for it has odd cardinality. One must leave $2k$ in $E$ and push the rest, and we reduce $E$ to the normal form
\[
E_0=[1,2(k-l)-1]\cup\{2k\}.
\]

\hfill q.e.d.!

% foot: Version 9 septembre 2001


\endgroup
% END UNIT Der-ChV-234/Der-ChV-234_p011-015.tex
% BEGIN UNIT Der-ChV-234/Der-ChV-234_p016-020.tex
\begingroup\makeatletter%
\newcommand{\printedpage}[1]{\marginpar{\footnotesize p.~#1}}%
\newcommand{\msp}[1]{\marginpar{\footnotesize [#1]}}%
\newcommand{\uCh}{\underline{\mathrm{Ch}}}%
\newcommand{\uHom}{\underline{\mathrm{Hom}}}%
\newcommand{\uFl}{\underline{\mathrm{Fl}}}%
\newcommand{\bDelta}{\boldsymbol{\Delta}}%
\newcommand{\Hom}{\mathrm{Hom}}%
\makeatother\let\printedpage\derpp\let\msp\dermsp%
\let\mspin\DERmspin%


\printedpage{14}
% head: A. Grothendieck : Dérivateurs, Ch. V, Édition des Universités de Montpellier II et Paris VII
What precedes justifies our reasoning, which gave by composition of paths well-defined composition maps
\[
\underbrace{\Hom_\Sigma(x,y)}_{=\pi_0\uCh_\Sigma(x,y)} \times \underbrace{\Hom_\Sigma(y,z)}_{=\pi_0\uCh_\Sigma(y,z)} \longrightarrow \underbrace{\Hom_\Sigma(x,z)}_{=\pi_0\underline{\textstyle\mathrm{Ch}}_\Sigma(x,z)}.
\]
% (the third "Ch" under the brace is printed in a larger size than the two others)

\msp{page 22}I say that this composition is associative. This follows from the fact that if we have three paths (of \textit{even} indices!)
\[
x \xrightarrow{c} y \xrightarrow{c'} z \xrightarrow{c''} t,
\]
we have (without having to carry out homotopies \dots) strict associativity
\[
c''(c'c) = (c''c')c.
\]

I say that there are identities. We define
\[
1_x \in \Hom_\Sigma(x,x)
\]
by the path of index 0 (which is even!), i.e. by the diagram reduced to $x$, with no arrow (or if we want to interpret the diagram as a functor $e \longrightarrow X$, there is the implicit identity arrow). Our arguments did not exclude the case $n = 0$ or $m = 0$, and it is tautological that for a path
\[
x \xrightarrow{c} y
\]
we have
\[
c \cdot 1_x = 1_y \cdot c = c.
\]
So we have indeed defined a category, which I am going to call $X_\Sigma$. I am going to define a canonical functor % [sic] Ja -> Je, misprint of the pronoun
\[
\theta : X \longrightarrow X_\Sigma
\]
which will be the identity on objects, we must define it on arrows
\[
\Hom_X(x,y) \longrightarrow \Hom_{X_\Sigma}(x,y) = \pi_0\uCh_\Sigma(x,y)
\]
\DERmark{ENV0276}

\texttt{[rather} $\pi_0\uCh_\Sigma(x,y)$\texttt{]}. Unfortunately, an arrow
\[
x \xrightarrow{u} y
\]
\msp{page 23}cannot be interpreted directly as a path in our sense, we have to pay for our evasions! But we are going to associate to such an arrow the path
\[
x \longrightarrow y \xleftarrow{\mathrm{id}} y \xrightarrow{\mathrm{id}} y \xleftarrow{\mathrm{id}} y,
\]
% foot: Version 9 septembre 2001
\printedpage{15}
% head: A. Grothendieck : Dérivateurs, Ch. V, Édition des Universités de Montpellier II et Paris VII
an element of $\uCh_{\Sigma,2}(x,y)$, which defines the desired map. We must check compatibility with identities and under composition. To the arrow $1_x : x \longrightarrow x$ is associated the element of $\uCh_2(X;x,x)$
\DERmark{ENV0277}

\[
x \xrightarrow{\mathrm{id}} x \xleftarrow{\mathrm{id}} x \xrightarrow{\mathrm{id}} x \xleftarrow{\mathrm{id}} x,
\]
we must prove that in $\pi_0(\uCh_{\Sigma,2}(x,x))$ it becomes equal to the element deduced from $1_x \in \uCh_{\Sigma,0}(X;x,x)$, by the transition functor. But this element is precisely none other than
\DERmark{ENV0278}

\[
x \xrightarrow{\mathrm{id}} x \xleftarrow{\mathrm{id}} x \xrightarrow{\mathrm{id}} x \xleftarrow{\mathrm{id}} x,
\]
it is the same (without even having to pass to $\pi_0$)!

For compatibility with composition
\[
x \xrightarrow{u} y \xrightarrow{v} z,
\]
we have on the one hand
\[
x \xrightarrow{vu} z \xleftarrow{\mathrm{id}} z \xrightarrow{\mathrm{id}} z \xleftarrow{\mathrm{id}} z,
\]
on the other hand
\[
x \xrightarrow{u} y \xleftarrow{\mathrm{id}} y \xrightarrow{\mathrm{id}} y \xleftarrow{\mathrm{id}} y \xrightarrow{v} z \xleftarrow{\mathrm{id}} z \xrightarrow{\mathrm{id}} z \xleftarrow{\mathrm{id}} z.
\]
We complete the first path by degeneracy at both ends, to compare it with the second as an element of $\uCh_4(x,z)$,
\[
x \xrightarrow{\mathrm{id}} x \xleftarrow{\mathrm{id}} x \xrightarrow{vu} z \xleftarrow{\mathrm{id}} z \xrightarrow{\mathrm{id}} z \xleftarrow{\mathrm{id}} z \xrightarrow{\mathrm{id}} z \xleftarrow{\mathrm{id}} z.
\]

\msp{page 24}We must prove that the two are homotopic. We thus have two functors
\[
f,g : \underbrace{\overbrace{\uFl(X;x,y)}^{\text{discrete category}} \times \overbrace{\uFl(X;y,z)}^{\text{discrete category}}}_{\mathcal{A}} \underset{g}{\overset{f}{\rightrightarrows}} \underbrace{\uCh_{4,\Sigma}(X;x,z)}_{\mathcal{B}}, \tag*{$(\ast)$}
\]
we want to prove that they are homotopic. For this, we regard them (as a moment ago) as ``transposes'' of the two functors
\[
\begin{tikzcd}[column sep=large]
\bDelta_1 \sqcup_{\bDelta_0} \bDelta_1 \simeq \bDelta_2 & \Phi_4 \arrow[l, shift right, "\varphi"'] \arrow[l, shift left, "\psi"] \\
& \varepsilon, \arrow[ul, "\alpha"] \arrow[u, "\alpha'"']
\end{tikzcd}
\]
\DERmark{ENV0279}

since the categories $\mathcal{A}$, $\mathcal{B}$ of the two sides of (*) can be interpreted as full subcategories of $\uHom(\bDelta_2, X)$ and $\uHom(\Phi_4, X)$ respectively; namely the full subcategories consisting of the functors whose composite with $\alpha$ resp. $\alpha'$ is given by $x$ and by $z$, and moreover (for $\mathcal{A}$) whose value at the middle object of $\bDelta_2$ is $y$. We can also replace $\mathcal{A}$ by the larger full subcategory $\mathcal{A}'$ of $\uHom(\bDelta_2, X)$, obtained by ignoring this last condition, and it suffices to prove that the two ``bigger'' functors
% foot: Version 9 septembre 2001
\printedpage{16}
% head: A. Grothendieck : Dérivateurs, Ch. V, Édition des Universités de Montpellier II et Paris VII
\[
\mathcal{A}' \underset{g'}{\overset{f'}{\rightrightarrows}} \mathcal{B}'
\]
are homotopic. And for this, it suffices to
\msp{page 25}see again that $\varphi$ and $\psi$ are homotopic in the category of functors $\Phi_4 \longrightarrow \bDelta_2$ which are compatible with the source, and the target. We thus see that our homotopy lemma of a moment ago was not general enough. We would have to define the \textit{generalized interval categories,} of arbitrary length $n$ ($n =$ number of arrows end to end), a \textit{standard} generalized interval category of length $n$ has as its set of objects the interval of integers % [sic] la catégories -> la catégorie, singular article with plural noun
\[
[0,n],
\]
and for every pair of consecutive integers $i-1$, $i$ among the objects ($1 \leq i \leq n$), there is exactly one arrow connecting $i-1$, $i$ \texttt{[rather} $i-1$\texttt{,} $i$\texttt{]}. We assign to the interval $[i,i+1]$ the sign $+$ if $i \longrightarrow i+1$ (arrow in the direct direction), the sign $-$ \texttt{[rather} $-$\texttt{]} if $i \longleftarrow i+1$ (arrow in the retrograde direction) - thus the standard categories are defined, for fixed $n$, by a system of signs $+$ or $-$, which we assign to the $n$ consecutive ``segments'' $A_i := [i-1,i]$ ($1 \leq i \leq n$) of the interval of integers $[0,n]$. There are thus exactly $2^n$ of them.
\DERmark{ENV0280}\DERmark{ENV0281}


We expect the following statement to be true:

\textbf{Lemma.} \textit{Let $I$, $I'$ be two generalized intervals, equipped with source objects and target object (so that $I$, $I'$ are} canonically \textit{isomorphic to the corresponding standard generalized intervals). Let} $\uHom^!(I,I')$ % [sic] genéralisés -> généralisés, missing accent on the first e
\msp{page 26}\textit{be the category of functors from $I$ to $I'$, compatible with source and target. If this category is $\neq \emptyset$, it is} 0\textit{-connected} (and, I presume, even contractible?).

I am going to assume provisionally that this is so, at least in the particular case of $\Phi_4$ (generalized interval of length 8) and $\bDelta_2$. This therefore gives a functor
\[
\theta : X \longrightarrow X_\Sigma.
\]
We have won, for the universal property of this functor is practically obvious. But I forgot to check that $\theta$ transforms elements of $\Sigma$ into invertible arrows. So let
\[
x \xrightarrow{u} y
\]
with $u \in \Sigma$, we are going to define a path in the reverse direction
\[
y \xrightarrow{c} x
\]
which will define in $X_\Sigma$ an arrow inverse to $u$. The natural idea is to take the ``path''
\[
y \xleftarrow{u} x
\]
% foot: Version 9 septembre 2001
\printedpage{17}
% head: A. Grothendieck : Dérivateurs, Ch. V, Édition des Universités de Montpellier II et Paris VII
from $y$ to $x$, but this is not a path in the sense of our definitions, we are going to complete it to
\[
c : y \xrightarrow{\mathrm{id}} y \xleftarrow{u} x \xrightarrow{\mathrm{id}} x \xleftarrow{\mathrm{id}} x.
\]
I say that $c$ does the job. We must prove two homotopy relations, of $c \cdot \bar{u}$ and of $\bar{u} \cdot c$ with the identity paths at $x$ \texttt{[resp. at} $y$\texttt{]} of length 8 (index 4, i.e. in $\uCh_{\Sigma,4}(X;x,x)$ resp. $\uCh_{\Sigma,4}(X;y,y)$),
\msp{page 27}where $\bar{u}$ is the path associated to $u$
\[
x \xrightarrow{u} y \xleftarrow{\mathrm{id}} y \xrightarrow{\mathrm{id}} y \xleftarrow{\mathrm{id}} y.
\]

We thus have to look at
\[
x \xrightarrow{u} y \xleftarrow{\mathrm{id}} y \xrightarrow{\mathrm{id}} y \xleftarrow{\mathrm{id}} y \xrightarrow{\mathrm{id}} y \xleftarrow{u} x \xrightarrow{\mathrm{id}} x \xleftarrow{\mathrm{id}} x
\]
and see that it is homotopic to the trivial path
\[
x \xrightarrow{\mathrm{id}} x \xleftarrow{\mathrm{id}} \cdots x \xrightarrow{\mathrm{id}} x \xleftarrow{\mathrm{id}} x
\]
(identities everywhere), and likewise for
\[
y \xrightarrow{\mathrm{id}} y \xleftarrow{u} x \xrightarrow{\mathrm{id}} x \xleftarrow{\mathrm{id}} x \xrightarrow{u} y \xleftarrow{\mathrm{id}} y \xrightarrow{\mathrm{id}} y \xleftarrow{\mathrm{id}} y
\]
and the trivial path at $y$. This time, we again have two functors
\[
\begin{tikzcd}[column sep=large]
\bDelta_1 & \Phi_4 \arrow[l, shift right, "\varphi"'] \arrow[l, shift left, "\psi"] \\
& \varepsilon, \arrow[ul, "\alpha'"] \arrow[u, "\alpha"']
\end{tikzcd}
\]
compatible with source objects and target objects, and we must prove that they are homotopic in $\uHom^!(\Phi_4, \bDelta_1)$. This must therefore follow from the same general lemma. We have won - by a simple interpretation of $X\Sigma^{-1}$, and in particular of the $\Hom_{X\Sigma^{-1}}(x,y)$ and of their composition, in terms of the geometry of path spaces:
\[
\left\{
\begin{aligned}
&\Hom_{X\Sigma^{-1}}(x,y) \simeq \pi_0(\uCh_\Sigma(X;x,y)) \\ % [sic] Hom_{XΣ^{-1}} -> Hom_{XΣ^{-1}}(x,y), arguments (x,y) missing
&\text{Composition of homomorphisms in } X\Sigma^{-1} \text{ corresponds to} \\
&\text{the composition of } \Sigma\text{-paths in } X
\end{aligned}
\right.
\]
\DERmark{ENV0282}

% foot: Version 9 septembre 2001

\printedpage{18}
% head: A. Grothendieck : Dérivateurs, Ch. V, Édition des Universités de Montpellier II et Paris VII
\noindent\msp{page 28}

\section*{3\quad Passage to categories of diagrams}

We should have denoted the category $X$ by $\mathcal{M}$, for later on it will play the role of a ``category of models''. On the other hand, we shall denote by $X$ the ``index categories'' on which we want to construct the ``categories of diagrams'' of type $X$
\[
\mathcal{M}(X) = \uHom(X,\mathcal{M}).
\]
We thus take
\[
\underbrace{\Sigma}_{\text{quasi-isomorphisms}} \subseteq \mathrm{Fl}(\mathcal{M}), \qquad \Sigma \supset \text{set of identity arrows},
\]
and for every category $X$ we define
\[
\Sigma_X \subset \mathrm{Fl}\,\mathcal{M}(X)
\]
as follows: if $F, G \in \mathrm{Ob}\,\mathcal{M}(X)$ ($F, G : X \rightrightarrows \mathcal{M}$), an arrow $u : F \longrightarrow G$ (homomorphism of functors) is in $\Sigma_X$ if and only if $\forall x \in X$, $u(x) : F(x) \longrightarrow G(x)$ is in $\Sigma$ (``criterion for quasi-isomorphy fiber by fiber'', for the homomorphism between the copresheaves $F$, $G$ under consideration). % [sic] homomophisme -> homomorphisme, letter r missing

Let
\[
a, b \in \mathrm{Ob}\,\mathcal{M}
\]
and consider the corresponding constant functors on $X$
\[
a_X, b_X \in \mathrm{Ob}\,\mathcal{M}(X).
\]

\msp{page 29}We propose to apply the theory of \S2 to the ``calculation'' of
\[
\Hom_{\Sigma_X}(a_X, b_X),
\]
we thus have a canonical bijection
\[
\boxed{\Hom_{\Sigma_X}(a_X, b_X) \simeq \pi_0\uCh_{\Sigma_X}(\mathcal{M}(X); a_X, b_X).} % [sic] a_x, b_x -> a_X, b_X, lowercase index against a_X, b_X elsewhere
\]
\DERmark{ENV0283}

We propose to make explicit the category of paths on the right-hand side. We have, leaving aside the data of $\Sigma$, $a$, $b$ in $\mathcal{M}$, the expression for the category of \textit{all} paths
% foot: Version 9 septembre 2001


\endgroup
% END UNIT Der-ChV-234/Der-ChV-234_p016-020.tex
% BEGIN UNIT Der-ChV-234/Der-ChV-234_p021-025.tex
\begingroup\makeatletter%
\NoAutoSpacing%
\newcommand{\printedpage}[1]{\marginpar{\footnotesize p.~#1}}%
\newcommand{\msp}[1]{\marginpar{\footnotesize [#1]}}%
\newcommand{\cM}{\mathcal{M}}%
\newcommand{\Ch}{\underline{\mathrm{Ch}}}%
\newcommand{\Chh}{\underline{\underline{\mathrm{Ch}}}}%
\newcommand{\Hom}{\underline{\mathrm{Hom}}}%
\newcommand{\KK}{\underline{\underline{K}}}%
\newcommand{\dfeq}{\overset{\mathrm{def}}{=}}%
\newcommand{\Cathot}{\mathrm{Cat\,hot}}%
\newcommand{\uHom}{\underline{\mathrm{Hom}}}%
\makeatother\let\printedpage\derpp\let\msp\dermsp%
\let\mspin\DERmspin%

\DERmark{ENV0284}


\printedpage{19}
% head: A. Grothendieck : Dérivateurs, Ch. V, Édition des Universités de Montpellier II et Paris VII

in $\cM(X)$
\[
\begin{aligned}
\Ch(\cM(X)) \quad &\simeq \quad \uHom(\ \underbrace{\boldsymbol{\Phi}}_{\substack{\text{pro-category}\\ (\Phi_n)_{n\in\mathbf{N}}}}\ ,\ \underbrace{\cM(X)}_{\uHom(X,M)}\ ) \\ % reading resolved: italic M is printed; category-valued Hom underlines restored.
&\dfeq \quad \varinjlim_n \uHom(\Phi_n,\cM(X)) \\
&\simeq \quad \varinjlim_n \uHom(\Phi_n\times X,\cM) \\
&\simeq \quad \varinjlim_n \uHom(X,\uHom(\Phi_n,\cM)) \\
&\dfeq \quad \varinjlim_n \uHom(X,\Ch_n(\cM)) \\
&\simeq \quad \uHom^{!}(X,\Ch(\cM)),
\end{aligned}
\]
where the sign ! means that we take the full subcategory of functors
\[
X\longrightarrow \Ch(\cM),
\]
\msp{page 30}
which factor through one of the full subcategories $\Ch_n(\cM)$. We can also say that if we consider the direct system of categories $\Ch_n(\cM)$ as an ind-object of Cat, which we are going to denote
\[
\Chh(\cM) \qquad\qquad\qquad\qquad \text{ind-object of Cat,}
\]
we have
\[
\Ch(\cM(X)) \xrightarrow{\ \sim\ } \uHom(X,\Chh(\cM)).
\]
If we turn to $\Sigma$, $\Sigma_X$, we find at once that a path in $\cM(X)$ is a $\Sigma_X$-path, if and only if the corresponding functor $X\longrightarrow\Ch(\cM)$ factors through $\Ch_{\Sigma}(\cM)$. Thus we find
\[
\Ch_{\Sigma_X}(\cM(X))\simeq\uHom^{!}(X,\Ch_{\Sigma}(\cM)),
\]
where the sign ! has the same meaning as above. If we introduce the ind-object $\Chh_{\Sigma}(\cM)$ defined by the $\Ch_{\Sigma,n}(\cM)$, we thus also find
\[
\Ch_{\Sigma_X}(\cM(X))\xrightarrow{\ \sim\ }\uHom(X,\Chh_{\Sigma}(\cM)).
\]

\msp{page 31}
Let $c$ be a path in $\cM(X)$, and
\[
X\longrightarrow\Ch(\cM)
\]
the functor it defines, let us compose it with the source-target functors
\[
\Ch(\cM)\longrightarrow\cM\times\cM,
\]
we find a functor
\[
X\longrightarrow\cM\times\cM,
\]
i.e. an object $(F_0,F_1)$ of $\cM(X)\times\cM(X)$. I say that $F_0$ and $F_1$ are respectively the source and the target of the path $c$. A little more precisely, we have commutativity (strict) in the diagram of functors
% foot: Version 9 septembre 2001
\printedpage{20}
% head: A. Grothendieck : Dérivateurs, Ch. V, Édition des Universités de Montpellier II et Paris VII
\[
\begin{tikzcd}[column sep=large, row sep=large]
\Ch_{\Sigma_X}(\cM(X)) \arrow[r, "\sim"] \arrow[d, "{\sigma_{\cM(X)},\,\beta_{\cM(X)}}"] & \uHom^{!}(X,\Ch_{\Sigma}(\cM)) \arrow[d, "{\text{deduced from }\sigma_{\cM},\,\beta_{\cM}}"] \\
\cM(X)\times\cM(X) \arrow[r, "\sim"] & \uHom(X,\cM\times\cM),
\end{tikzcd}
\]
where the functors $\sigma$ are the source functors, the functors $\beta$ the target functors.

If we give ourselves arbitrary objects $F_0$, $F_1$ in $\cM(X)$, hence a functor
\[
X\longrightarrow\cM\times\cM,
\]
$\Ch_{\Sigma_X}(F_0,F_1)$ thus identifies with the category of ``admissible'' functors (i.e. factoring through a $\Ch_{\Sigma,n}(\cM)$) $X\xrightarrow{\ C\ }\Ch_{\Sigma}(\cM)$ which \textit{lift}
\msp{page 32 blank}
\msp{page 33}
the functor $(F_0,F_1)$ \texttt{[rather} $(F_0,F_1)$\texttt{]} above, i.e. which give rise to a commutative triangle (strict)
\DERmark{ENV0285}

\[
\begin{tikzcd}[row sep=large]
 & \Ch_{\Sigma}(\cM) \arrow[d, "{(\sigma_{\cM},\,\beta_{\cM})}"] \\
X \arrow[ur, "C"] \arrow[r, "{(F_0,F_1)}"] & \cM\times\cM.
\end{tikzcd}
\]
But if the functors $F_0$, $F_1$ are constant, with value $a$, $b$ as before, i.e. $(F_0,F_1)$ constant with value $(a,b)$, then the category of functors $C$ in question is tautologically isomorphic to that of functors from $X$ into the fiber of $(\sigma_{\cM},\beta_{\cM})$ at $(a,b)$, whence the fundamental isomorphism of categories
\[
\boxed{\Ch_{\Sigma_X}(\cM(X);a_X,b_X)\simeq\uHom^{!}(X,\Ch_{\Sigma}(\cM;a,b)).}
\]
\DERmark{ENV0286}


We are going to make explicit the essential content of this formula in a

\textbf{Scholie.} Let $\cM$ be a category with a set $\Sigma$ of localization arrows, $a,b\in\mathrm{Ob}\,\cM$. Then we can find an object $\KK=K(a,b)$, $K=(K_n)$, of Ind\,Cat (indexed by $\mathbf{N}$ but this fact matters little), and a canonical isomorphism for every category $X$ (isomorphism functorial in $X$)
\[
\begin{aligned}
\mathrm{Hom}_{\Sigma_X}(a_X,b_X) \quad &\simeq \quad \pi_0(\uHom(X,\KK)) \\
&\dfeq \quad \pi_0(\ \varinjlim \uHom(X,K_n)) \\
&\simeq \quad \varinjlim \pi_0(\uHom(X,K_n)).
\end{aligned}
\]
% foot: Version 9 septembre 2001

\printedpage{21}
% head: A. Grothendieck : Dérivateurs, Ch. V, Édition des Universités de Montpellier II et Paris VII

\msp{page 34}
I am going to dwell on this result, without worrying for the moment about the particular way in which the ind-category $\KK$ \texttt{[rather} $\KK$\texttt{]} was obtained in terms of $\cM$, $\Sigma$, $a$, $b$.
\DERmark{ENV0287}


Since the $\pi_0\uHom(X,K_n)=\mathrm{Hom}_{\Cathot}(X,K_n)$ depend only on $X$ and on $K_n$ as objects of Cat\,hot, I am going to interpret $\KK$ as an ind-object of Cat\,hot, and write the fundamental formula in the form
\[
\mathrm{Hom}_{\Sigma_X}(a_X,b_X)\simeq\mathrm{Hom}_{\mathrm{Ind\,Cat\,hot}}(X,\KK).
\]

Thus, the statement of the scholie implies two things:

\begin{enumerate}
\item[a)] The contrafunctor $X\longmapsto\mathrm{Hom}_{\Sigma_X}(a_X,b_X)$ on Cat factors canonically through Cat\,hot, via a canonical functor
\[
\Cathot^{o}\ \xrightarrow{\ H=H^{\cM,\Sigma}_{a,b}\ }\mathrm{Ens}\ ;
\]
\item[b)] This latter contrafunctor on Cat\,hot is \textit{ind-representable} by an ind-object $\KK$ of Cat\,hot.
\end{enumerate}

It is this ind-object of Cat\,hot which is defined up to unique isomorphism by $\cM$, $\Sigma$, $a$, $b$, independently of any computational artifice for making the $\mathrm{Hom}_{\Sigma_X}(a_X,b_X)$ explicit.

Let us also note the following conclusion, by no means obvious a priori:

\msp{page 35}
Let us first note (this being implicit in the scholie) that every arrow in Cat
\[
X\xrightarrow{\ f\ }Y
\]
defines
\[
\cM(Y)\xrightarrow{\cM(f)}\cM(X) \qquad\qquad\qquad \text{mapping } \Sigma_Y \text{ into } \Sigma_X,
\]
whence
\[
\cM(Y)\Sigma_Y^{-1}\dfeq\mathbf{D}_{\cM,\Sigma}(Y)\longrightarrow\cM(X)\Sigma_X^{-1}=\mathbf{D}_{\cM,\Sigma}(X),
\]
\DERmark{ENV0288}

\texttt{[rather} $\cM(X)$\texttt{]} whence, for every $a,b\in\mathrm{Ob}\,\cM$, since $\cM(f)$ transforms $a_Y$, $b_Y$ into $a_X$, $b_X$:
\[
\boxed{\mathrm{Hom}_{\Sigma_Y}(a_Y,b_Y)\longrightarrow\mathrm{Hom}_{\Sigma_X}(a_X,b_X).}
\]

This being so, we have:

\textbf{Corollary:} \textit{If $f$ is a homotopism, then for any $a,b\in\cM$, the preceding arrow is an isomorphism.} % [sic] quelque soient -> quels que soient, grammatical form of the locution

The natural question is to know whether the same conclusion is still valid, assuming only that $f$ is a ``\textit{quasi-isomorphism}'' (or ``weak equivalence''), here that $f$ defines an isomorphism in the category of fractions Hot of Cat. (Not to be confused with Cat\,hot,
% foot: Version 9 septembre 2001
\printedpage{22}
% head: A. Grothendieck : Dérivateurs, Ch. V, Édition des Universités de Montpellier II et Paris VII
\msp{page 36}
distinctly finer). If this were so, the contrafunctor
\[
X\longmapsto\mathrm{Hom}_{\Sigma_X}(a_X,b_X)
\]
would factor not only through Cat\,hot, but even through the coarser category of fractions Hot. And then the question would arise whether this hypothetical contrafunctor
\[
\mathrm{Hot}^{o}\longrightarrow\mathrm{Ens}
\]
is still ind-representable; and preferably (\footnote[5]{this will be automatic if we have ind-representability as a functor on Hot, cf. below.}) ind-representable by the ind-object image of $\KK$ in Ind\,Hot under
\[
\mathrm{Ind}(\Cathot)\longrightarrow\mathrm{Ind}(\mathrm{Hot}),
\]
deduced from the localization functor $\Cathot\longrightarrow\mathrm{Hot}$. (For this ind-object is a ready-made candidate.)

\msp{page 37}
I am going to give a few categorical reminders, generalizing the situation of the canonical localization functor
\[
\Cathot\xrightarrow{\ \varphi\ }\mathrm{Hot},
\]
which we know admits a right adjoint \texttt{[if a reader has a proof of this fact, please contact G. Maltsiniotis: [e-mail address removed]]}
\[
\Cathot\xleftarrow{\ \psi\ }\mathrm{Hot}, \qquad\qquad\qquad \mathrm{Hom}_{\mathrm{Hot}}(\varphi X,Y)\simeq\mathrm{Hom}_{\Cathot}(X,\psi Y),
\]
where $\psi$ is the ``Kan envelope'' or ``Kan representative'' functor. We are therefore going to consider the more general situation
\[
\begin{tikzcd}[column sep=large]
C \arrow[r, shift right=1, "\varphi"'] & C' \arrow[l, shift right=1, "\psi"']
\end{tikzcd}
\qquad\qquad \mathrm{Hom}_{C'}(\varphi X,X')\simeq\mathrm{Hom}_{C}(X,\psi X')
\]
i.e.
\[
\varphi,\psi \text{ pair of adjoint functors.}
\]
We assume moreover
\[
\left\{
\begin{aligned}
&\varphi \text{ a localization functor, i.e. } \psi \text{ fully faithful,}\\
&\text{i.e. } \varphi\psi\xrightarrow[\sim]{\text{adj}}\mathrm{id}_{C'} \quad .
\end{aligned}
\right.
\]
Let $W=W_{\varphi}\subset\mathrm{Fl}\,C$, defined by $\varphi$: $W_{\varphi}=\{u\in\mathrm{Fl}\,C\mid\varphi(u)\text{ isomorphism}\}$. Here is the proposition I had forgotten, and which I had trouble reconstructing, alas!

\textbf{Proposition 1.} \textit{Let $K\in \mathit{Ob}C$. The following conditions on $K$ are equivalent:}
% foot: Version 9 septembre 2001

\printedpage{23}
% head: A. Grothendieck : Dérivateurs, Ch. V, Édition des Universités de Montpellier II et Paris VII

\begin{enumerate}
\item[(i)] \textit{The functor}
\[
X\longmapsto K(X)\dfeq\mathrm{Hom}_C(X,K) \qquad\quad C^{o}\longrightarrow\mathrm{Ens}
\]
\textit{transforms arrows of $W$ (``quasi-isomorphisms'') into bijections.}
\item[(i$'$)] \textit{Same condition for the adjunction arrows}
\[
X\longrightarrow\psi\varphi X \qquad\quad (\forall X\in\mathrm{Ob}\ \mathrm{C})
\]
\item[(i$''$)] \textit{Same condition for the adjunction arrow} (\footnote[6]{N.B. It even suffices that the element $\mathrm{id}_K$ of $F_K(K)$ come from $F_K(\psi\varphi K)$ \texttt{[where} $F_K=K(-)=\mathrm{Hom}(-,K)$\texttt{]}, i.e.\par (i$'''$) There exists a retraction of $\psi\varphi K$ onto $K$ \texttt{[rather} $\psi\varphi K$ \texttt{onto} $K$\texttt{]} (NB This will even be an isomorphism).})
\[
K\longrightarrow\psi\varphi K.
\]
\DERmark{ENV0289}

\item[(ii)] \textit{$K$ belongs to the essential image of $\psi$}
\item[(ii$'$)] \textit{$K\longrightarrow\psi\varphi K$ (adjunction) is an isomorphism.}
\item[(ii$''$)] \textit{The homomorphism functorial in $X$}
\[
\mathrm{Hom}_C(X,K)\longrightarrow\mathrm{Hom}_{C'}(\varphi X,\varphi K)
\]
\textit{is an isomorphism.}
\end{enumerate}

\textsc{Proof.} We first have as quasi-tautologies (ii$'$)$\Longrightarrow$ (ii)$\Longrightarrow$ (ii$''$)$\Longrightarrow$ (ii$'$), i.e. the equivalence of the last three conditions. Indeed, (ii$'$)$\Longrightarrow$ (ii) tautological, (ii)$\Longrightarrow$ (ii$''$) since if $K=\psi(K')$ we have % reading resolved: single and double mathematical primes as printed.
\msp{page 38}
by adjunction
\[
\mathrm{Hom}_C(X,K)=\mathrm{Hom}_C(X,\psi(K'))\overset{\text{adj}}{\simeq}\mathrm{Hom}_{C'}(\varphi X,K'),
\]
now $\varphi(K)=\varphi\psi(K')\xrightarrow{\ \sim\ }K'$ (adjunction, and $\psi$ fully faithful).
\[
\begin{tikzcd}[column sep=small, row sep=large]
\mathrm{Hom}_C(X,K) \arrow[rr, "\sim"] \arrow[dr, "\varphi_{X,K}"] & & \mathrm{Hom}_{C'}(\varphi X,K') & \text{\footnotesize which is completed to} \\ % [sic] Hom_C(\varphi X,K') -> Hom_{C'}(\varphi X,K'), \varphi X and K' are objects of C' (cf. the display above)
& \mathrm{Hom}_{C'}(\varphi X,\ \underbrace{\varphi K}_{=\varphi\psi(K')}\ ) \arrow[ur, "\sim", "\text{adj}"'] & &
\end{tikzcd}
\]
\DERmark{ENV0290}

which proves that $\varphi_{X,K}$ is an isomorphism.

(ii$''$) $\Longrightarrow$ (ii$'$). For the arrow
% foot: Version 9 septembre 2001


\endgroup
% END UNIT Der-ChV-234/Der-ChV-234_p021-025.tex
% BEGIN UNIT Der-ChV-234/Der-ChV-234_p026-030.tex
\begingroup\makeatletter%
\newcommand{\printedpage}[1]{\marginpar{\footnotesize p.~#1}}%
\newcommand{\msp}[1]{\marginpar{\footnotesize [#1]}}%
\newcommand{\Hom}{\operatorname{Hom}}%
\newcommand{\Ind}{\operatorname{Ind}}%
\newcommand{\Ens}{\mathrm{Ens}}%
\newcommand{\Cat}{\mathrm{Cat}}%
\newcommand{\Cathot}{\mathrm{Cat\,hot}}%
\newcommand{\Hot}{\mathrm{Hot}}%
\newcommand{\Ob}{\mathrm{Ob}}%
\newcommand{\id}{\mathrm{id}}%
\newcommand{\Ch}{\mathrm{Ch}}%
\newcommand{\Chh}{\underline{\underline{\mathrm{Ch}}}}%
\makeatother\let\printedpage\derpp\let\msp\dermsp%
\let\mspin\DERmspin%


\printedpage{24}
% head: A. Grothendieck : Dérivateurs, Ch. V, Édition des Universités de Montpellier II et Paris VII
% foot: Version 9 septembre 2001

\[
\begin{tikzcd}
\Hom_C(X,K) \arrow[r,"\alpha_X"] \arrow[dr,dashed] & \Hom_{C'}(\varphi X,\varphi K) \arrow[d,dash,"\wr"',"\mathrm{adj}"] \\
 & \Hom_C(X,\psi\varphi K)
\end{tikzcd}
\]
is deduced from the adjunction arrow
\[
K \longrightarrow \psi\varphi K
\]
by taking $\Hom_C(X,-)$, so $\alpha_X$ is an isomorphism for every $X$ if and only if this adjunction arrow is an isomorphism. So
\[
\text{(ii)} \Longleftrightarrow \text{(ii$'$)} \Longleftrightarrow \text{(ii$''$)}.
\]
On the other hand,
\[
\text{(ii$''$)} \Longrightarrow \text{(i)} \Longrightarrow \text{(i$'$)} \Longrightarrow \text{(i$''$)}
\]
are tautological, bearing in mind for (i) $\Longrightarrow$ (i$'$) \texttt{[probably (i) $\Longrightarrow$ (i$'$)]} that the adjunction arrows
\DERmark{ENV0291}

\[
X \xrightarrow{u} \psi\varphi X \qquad\qquad\qquad \text{in } C
\]
are indeed in $W$, for we have a commutative triangle
\[
\begin{tikzcd}[column sep=large]
\varphi(X) \arrow[r,"\varphi(u)"] \arrow[dr,"\id_{\varphi(X)}"] & \varphi(\psi\varphi X) = (\varphi\psi)(\varphi X) \arrow[d,"{\begin{array}{l}\wr\ \mathrm{adj}\ \varphi\psi \longrightarrow 1,\ \text{which is an}\\ \text{isomorphism since } \psi \text{ is fully faithful}\end{array}}"] \\
 & \varphi(X),
\end{tikzcd}
\]
so $\varphi(u)$ is indeed an isomorphism.

It remains to prove (i$''$) $\Longrightarrow$ (ii$'$), i.e. we suppose

\noindent\msp{page 39}%
that the functor $X \longrightarrow \Hom(X,K)$ transforms the particular adjunction morphism
\[
K \xrightarrow{i} \psi\varphi K
\]
into
\[
\begin{aligned}
u \quad &\longrightarrow \quad u\circ i\\
\Hom(\psi\varphi(K),K) \quad &\longrightarrow \quad \Hom(K,K),
\end{aligned}
\]
which is a bijection. So $\id_K$ is in the image, let $p$ be an element of the first member which gives rise to it

\printedpage{25}
% head: A. Grothendieck : Dérivateurs, Ch. V, Édition des Universités de Montpellier II et Paris VII
% foot: Version 9 septembre 2001

\[
\begin{tikzcd}
K \arrow[r,hook,shift left,"i"] & \psi\,\underbrace{\varphi(K)}_{K'} \arrow[l,shift left,"p"]
\end{tikzcd}
\qquad\qquad {\scriptstyle \text{so } p\circ i = \id_K}
\]
Let us consider
\[
ip = q : \psi(K') \longrightarrow \psi(K') \qquad\qquad\qquad (\text{NB } q^2 = q),
\]
since $\psi$ is fully faithful, there exists a unique
\[
q_0 : K' \longrightarrow K' \qquad\qquad\qquad (K' \overset{\mathrm{def}}{=} \varphi(K))
\]
\DERmark{ENV0292}

such that we have
\[
q = \psi(q_0).
\]
The relation $pi = \id_K$ gives, on applying $\varphi$,
\[
\varphi(p)\varphi(i) = \id,
\]
so since $\varphi(i)$ is an isomorphism (it is
\[
\begin{tikzcd}[column sep=large]
\varphi(K) \arrow[r,"\sim"] \arrow[dr,"\sim"{sloped,pos=0.45},"\id"{pos=0.6}] & \varphi(\psi\varphi)(K) = (\varphi\psi)(\varphi K) \arrow[d,"\wr"',"\mathrm{adj}"] \\
 & \varphi(K),
\end{tikzcd}
\]
making the triangle opposite commutative), $\varphi(p)$ is one. So $\varphi(q) = \varphi(i)\varphi(p)$ is one, since $\varphi(i)$ and $\varphi(p)$ \texttt{[rather $\varphi(p)$]} are. But
\DERmark{ENV0293}

\[
\varphi(q) = \varphi\psi(q_0) \quad \underbrace{\xrightarrow{\ \sim\ }}_{\mathrm{adj}\ \varphi\psi\ \xrightarrow{\sim}\ \id} \quad q_0,
\]
so $q_0$ is an isomorphism, hence also $q = \psi q_0$. Since $q^2 = q$, this implies $q = \id$. So $i$ and $p$ are isomorphisms inverse to one another, which establishes (ii$'$) \texttt{[rather (ii$'$)]} and completes the proof.
\DERmark{ENV0294}


\textbf{Corollary:} \textit{Let us consider functors $C^o \xrightarrow{H} \Ens$, and $C'^o \xrightarrow{H'} \Ens$, such that $H \simeq H'\circ\varphi$. For $H$ to be representable, it is necessary and sufficient that $H'$ be so, and in this case, the object $K$ of $C$}

\noindent\msp{page 40}%
\textit{which represents $H$ satisfies the equivalent conditions of the proposition. If $K$, $K'$ represent them, we have canonical isomorphisms $K' \simeq \varphi K$, $K \simeq \psi K'$.}

If $H'$ is representable by $K' \in \Ob C'$, we will therefore have
\[
H'\circ\varphi(X) \simeq \Hom_{C'}(\varphi X,K') \overset{\mathrm{adj}}{\simeq} \Hom_C(X,\psi(K'))
\]
\printedpage{26}
% head: A. Grothendieck : Dérivateurs, Ch. V, Édition des Universités de Montpellier II et Paris VII
% foot: Version 9 septembre 2001
so $H$ is representable by $K = \psi(K')$, satisfying condition (ii) of the proposition. Conversely, if $H$ is representable by an object $K$, then $K$ satisfies condition (i), so by (ii$''$) we have
\[
\Hom_C(X,K) \xrightarrow{\sim} \Hom_{C'}(\varphi X,\varphi K),
\]
so the functor $H' : C'^o \longrightarrow \Ens$ which factors $H$ is given explicitly as $X' \longmapsto \Hom_{C'}(X',\varphi K)$, hence is representable by $\varphi K$. % [sic] C^o -> C'^o, H' is a functor on C'^o (cf. the Corollaire)
\DERmark{ENV0295}


Let us pass to ind-objects. The functors $\varphi$, $\psi$ pass to ind-objects, whence functors
\[
\begin{tikzcd}[column sep=large]
\Ind C \arrow[r,shift left,leftarrow,"\Psi = \Ind\psi"] \arrow[r,shift right,"\Phi = \Ind\varphi"'] & \Ind C'.
\end{tikzcd} % [sic] both arrows printed pointing to Ind C' -> the arrow Psi = Ind psi apparently intended from Ind C' to Ind C, psi goes from C' to C
\]
\DERmark{ENV0296}

We will again have the adjunction formula for $X = (X_i)_{i\in I}$ in $\Ind C$, $X' = (X'_{i'})_{i'\in I'}$ in $\Ind C'$, for
\[
\begin{aligned}
\Hom(X,\Psi X') &= \varprojlim_i \varinjlim_{i'} \Hom(X_i,\psi(X'_{i'}))\\
&= \varprojlim_i \varinjlim_{i'} \Hom(\varphi X_i,X'_{i'})\\
&= \Hom(\Phi X,X').
\end{aligned}
\]

\noindent\msp{page 41}%
We prove in the same way that $\Psi$ is fully faithful, for if $X' = (X'_i)_{i\in I}$, $Y' = (Y'_j)_{j\in J}$ are objects of $\Ind C'$, we will have
\[
\begin{tikzcd}[column sep=small]
\Hom(X',Y') \arrow[d,"\Psi_{X',Y'}"] & \simeq \varprojlim\limits_i \varinjlim\limits_j & \Hom_{C'}(X'_i,Y'_j) \arrow[d,"\psi_{X'_i,Y'_j}"] \\
\Hom(\Psi X',\Psi Y') & \simeq \varprojlim\limits_i \varinjlim\limits_j & \Hom_C(\psi(X'_i),\psi(Y'_j)),
\end{tikzcd}
\]
so $\Psi_{X',Y'}$ is an isomorphism if the $\psi_{X'_i,Y'_j}$ are.

So proposition 1 can be applied to the situation $\Phi$, $\Psi$ between ind-objects. I am going to spell it out in the most convenient form:

\textbf{Proposition 2.} \textit{Let $K = (K_i)_{i\in I}$ be an ind-object of $C$, and let $H$ be the functor}
\[
\left\{
\begin{aligned}
H : X \quad &\longmapsto \quad \underbrace{\Hom_C(X,K)}_{K(X)} \overset{\mathrm{def}}{=} \varinjlim_i \Hom_C(X,K_i)\\
C^o \quad &\longrightarrow \quad \Ens
\end{aligned}
\right.
\]
\DERmark{ENV0297}

\textit{which it ind-represents. Equivalent conditions}

\begin{enumerate}[label={}]
\item[(i)] \textit{The functor $H$ factors (up to isomorphism) through $C'$ via $\varphi$, i.e. transforms the $u \in W_\varphi$ into isomorphisms.}
\printedpage{27}
% head: A. Grothendieck : Dérivateurs, Ch. V, Édition des Universités de Montpellier II et Paris VII
% foot: Version 9 septembre 2001
\item[(i$'$)] \textit{Same condition for the adjunction arrows $X \longrightarrow \psi\varphi(X)$ (which are in $W$).}
\item[(i$''$)] \textit{Same condition for the adjunction arrow $K \longrightarrow \psi\varphi(K)$.}
\item[(ii)] \textit{$K$ belongs to the essential image of} $\Ind C'$ \textit{under $\Psi$.}
\item[(ii$'$)] $K \longrightarrow \Psi\Phi(K)$ \textit{is an} isomorphism \textit{in} $\Ind C$. % [sic] ".)" -> ".", closing parenthesis without opening parenthesis
\item[(ii$''$)] \textit{The functorial homomorphism}
\[
K(X) = \Hom_{\mathrm{Ind}(C)}(X,K) \longrightarrow \Hom_{\mathrm{Ind}(C')}(\Phi(X),\Phi(K))
\]
\DERmark{ENV0298}\DERmark{ENV0299}

\texttt{[rather} $\Hom_{\mathrm{Ind}(C)}(X,K) \longrightarrow \Hom_{\mathrm{Ind}(C')}(\Phi(X),\Phi(K))$\texttt{, identifying the object} $X$ \texttt{of} $C$ \texttt{with its image in} $\mathrm{Ind}(C)$\texttt{]} \textit{is an isomorphism.}
\end{enumerate}

\textsc{Proof.} The equivalence of (ii), (ii$'$), (ii$''$) is proposition 1 applied to $\Phi$ and $\Psi$, except that we must note that (ii$''$) obviously remains valid if we replace $X$ in it by an arbitrary ind-object $(X_j)_{j\in J}$. We have

\noindent\msp{page 42}%
\[
\text{(ii$''$)} \Longrightarrow \text{(i)} \Longrightarrow \text{(i$'$)}
\]
tautologically, it remains to prove that (i$'$) implies (i$''$) and (i$''$) $\Longrightarrow$ (ii$'$). We cannot apply proposition 1 directly, for hypothesis (i$'$) here concerns only the objects of $\Ind C$ which come from $C$, and it is not obvious a priori that it passes to arbitrary objects. But it is in the case of a homomorphism of ind-objects $(X_i)$, $(Y_i)$ on the same index set, given by an ordinary (term-by-term) homomorphism of direct systems on $I$, $X_i \longrightarrow Y_i$: if the $X_i \longrightarrow Y_i$ are in $W$, hypothesis (i) implies that $H(Y) \longrightarrow H(X)$ is an isomorphism, and likewise for (i$'$) for $\underbrace{X}_{=(X_i)_i} \;\longrightarrow\; \underbrace{\Psi\Phi(X)}_{=(\psi\varphi(X_i))_i}$\,, where the adjunction homomorphism is given term by term by the adjunction homomorphisms $X_i \longrightarrow \psi\varphi(X_i)$. So hypothesis (i$'$) here implies hypothesis (i$'$) of proposition 1 for $\Phi$ and $\Psi$, hence also (ii$'$) (by proposition 1), and we win. % [sic] "et (i'') (ii')" -> "et (i'') => (ii')", the implication sign or verb is missing
\DERmark{ENV0300}


Let us return to the situation where $C = \Cathot$, $C' = \Hot$, $\varphi$ the functor \texttt{[of]} localization. If $K \in \Ob\Cathot$, then $\psi\varphi(K)$ is called \textit{the Thomason-Kan envelope} of $K$ - if $K$ is a category, its envelope $\bar{K}$ and the arrow $K \longrightarrow \bar{K}$ are only defined in the category Cat\,hot, i.e. $\bar{K}$ is only defined ``up to homotopism''. % reading resolved: no underline is printed under psi phi(K).

\noindent\msp{page 43}%
But we can always \textit{choose} an object $\bar{K}$ of Cat, and an arrow
\[
i : K \longrightarrow \bar{K} \qquad\qquad\qquad \textit{in } \Cat,
\]
\printedpage{28}
% head: A. Grothendieck : Dérivateurs, Ch. V, Édition des Universités de Montpellier II et Paris VII
% foot: Version 9 septembre 2001
not only in Cat\,hot, in such a way moreover that $i$ is a ``Thomason \textit{cofibration}'' (and in particular an open immersion?). It follows that if we have a direct system
\[
K = (K_n)_{n\in\mathbf{N}}
\]
indexed by $\mathbf{N}$, and if we choose for each $n$ a Thomason-Kan envelope
\[
K_n \overset{\alpha_n}{\hookrightarrow} \bar{K}_n
\]
in this way, we can find transition functors
\[
\bar{K}_n \xrightarrow{\bar{i}_n} \bar{K}_{n+1}
\]
such that the squares
\[
\begin{tikzcd}
K_n \arrow[r,hook,"\alpha_n"] \arrow[d,"i_n"'] & \bar{K}_n \arrow[d,"\bar{i}_n"] \\
K_{n+1} \arrow[r,hook,"\alpha_{n+1}"] & \bar{K}_{n+1}
\end{tikzcd}
\]
are commutative, so that we thus have a homomorphism of \textit{direct systems of} Cat
\[
K \longrightarrow \bar{K}
\]
and a fortiori a homomorphism of ind-objects of Cat. This being so, conditions (i), (ii$'$) of proposition 2 show that the functor ind-represented by $K$ \textit{in} $C = \Cathot$ transforms ordinary

\noindent\msp{page 44}%
quasi-isomorphisms (weak equivalences) into bijections, \textit{if and only if} the preceding homomorphism is an isomorphism of ind-objects of Ind\,Cat\,hot.

So that will be the natural condition to impose on the ind-categories of paths
\[
\Chh_{\Sigma}(\mathcal{M};a,b) \qquad\qquad\qquad (a,b \in \mathrm{ob}\,\mathcal{M})
\]
coming from a ``category of models'' $\mathcal{M}$, $\Sigma$.

Of course, having at our disposal an explicit direct system of categories (not only of objects of Cat\,hot)
\[
K_n = \underline{\Ch}_{\Sigma,n}(\mathcal{M};a,b),
\]
it is tempting to pass to the direct limit
\[
K_\infty = \varinjlim_n \underline{\Ch}_{\Sigma,n}(\mathcal{M};a,b) \overset{\mathrm{def}}{=} \underline{\Ch}_{\Sigma}(\mathcal{M};a,b),
\]
\DERmark{ENV0301}

and this all the more so since the transition functors are monomorphisms and fully faithful, so that the $K_n$ are identified with a sequence of full subcategories of


\endgroup
% END UNIT Der-ChV-234/Der-ChV-234_p026-030.tex
% BEGIN UNIT Der-ChV-234/Der-ChV-234_p031-034.tex
\begingroup\makeatletter%
\newcommand{\printedpage}[1]{\marginpar{\footnotesize p.~#1}}%
\newcommand{\msp}[1]{\marginpar{\footnotesize [#1]}}%
\newcommand{\Cathot}{\mathrm{Cat\,hot}}%
\newcommand{\uHom}{\underline{\mathrm{Hom}}}%
\newcommand{\uuK}{\underline{\underline{K}}}%
\newcommand{\uCh}{\underline{\mathrm{Ch}}}%
\newcommand{\Hom}{\mathrm{Hom}}%
\newcommand{\Hot}{\mathrm{Hot}}%
\newcommand{\Ens}{\mathrm{Ens}}%
\newcommand{\Ch}{\mathrm{Ch}}%
\newcommand{\defeq}{\overset{\mathrm{def}}{=}}%
\newcommand{\Chh}{\underline{\underline{\mathrm{Ch}}}}%
\makeatother\let\printedpage\derpp\let\msp\dermsp%
\let\mspin\DERmspin%

\DERmark{ENV0302}


\printedpage{29}
% head: A. Grothendieck : Dérivateurs, Ch. V, Édition des Universités de Montpellier II et Paris VII
% foot: Version 9 septembre 2001
$K_\infty$, whose union is $K_\infty$. But on the other hand I do not know how to give any convincing interpretation of $K_\infty$, like that of the ind-object of Cat\,hot defined by $\uuK = (K_n)_{n\in\mathbf{N}}$, which ind-represents the remarkable functor % [sic] defini -> défini, missing accent
\[
H = H_{\mathcal{M};x,y} : \Cathot^{o} \longrightarrow \Ens
\]
factoring the functor
\msp{page 45}
\[
X \longmapsto H_{\mathcal{M};x,y}(X) \defeq \Hom_{\mathcal{M}(X)\Sigma_X^{-1}}(x_X, y_X),
\]
by the formula, then
\[
H(X) = \varinjlim_{n} \underbrace{\pi_0\uHom(X, K_n)}_{\Hom_{\Cathot}(X,K_n)} .
\]
I do not know how to give an interpretation of this type either of the category $K_\infty$ itself, or of its image in Cat\,hot, or of its image in Hot. It is true that if I restrict the functor $H$ to the full subcategory of (Cat) consisting of \textit{finite} categories (Ob\,$X$ and Fl\,$X$ finite), then we will have
\[
\uHom^{!}(X, K_\infty) = \uHom(X, K_\infty),
\]
and the preceding isomorphism is interpreted as % [sic] précédant -> précédent, adjective intended
\[
\Hom_{\mathcal{M}(X)\Sigma_X^{-1}}(a_X, b_X) \defeq H(X) \simeq \pi_0\uHom(X, K_\infty) \defeq \Hom_{\Cathot}(X, K_\infty).
\]
\DERmark{ENV0303}

\texttt{[rather} $(a_X, b_X)$\texttt{]} This proves at least that knowledge of $K_\infty$ is useful. But this formula, with $X$ restricted to categories that are \textit{finite,} will not define $K_\infty$ as an object of Cat\,hot, up to unique isomorphism in this category - \textit{except} in the case where $K_\infty$ is itself homotopic to a
\msp{page 46}
finite category. (I do not quite see whether this happens very often \ldots\ But one cannot found a general theory of derivators and ``derivable'' categories of models on an assumption of this sort.)

But suppose that the functor $H$ under consideration even factors through Hot, i.e. that it transforms quasi-isomorphisms into bijections. We have seen that this is equivalent to the condition that it can be represented by a direct system over $\mathbf{N}$ made up of Thomason-Kan categories. It must moreover be easy, if one wants to, to arrange for the transition functors to be monomorphisms and fully faithful, so as to have a nice $\varinjlim$. It will of course be a Thomason-Kan category. Is this category determined by the functor $H : \Cathot^{o} \longrightarrow \Ens$ up to unique isomorphism in Cat\,hot? Of course, there is no question of its representing the functor $H$, except \texttt{[if]} the latter is really and truly representable, i.e. if the ind-object $\uuK$ of Cat\,hot is essentially constant (which must practically never be the case). Let us already note
\printedpage{30}
% head: A. Grothendieck : Dérivateurs, Ch. V, Édition des Universités de Montpellier II et Paris VII
% foot: Version 9 septembre 2001
\msp{page 47}
that if $u$ is a homomorphism of ind-objects \textit{in} Cat, or, what now amounts to the same thing, a homomorphism in Cat % [sic] homomorphimse -> homomorphisme, transposed letters
\[
K_\infty \xrightarrow{u_\infty} K'_\infty
\]
such that $\forall n$, $u_\infty(K_n) \subseteq K'_{n'}$ for a suitable $n' = \tau(n)$, and if we assume that the homomorphism of ind-objects induces an isomorphism for the ind-representable functors \textit{on} Cat\,hot (for example, if $\uuK$ and $\uuK'$ ind-represent the same functor $H$, and the homomorphism of ind-objects is the one corresponding to $\mathrm{id}_H$), then $u_\infty$ is necessarily an isomorphism \textit{in} Hot, hence also (unless I am mistaken) in Cat\,hot since $K_\infty$, $K'_\infty$ are Thomason-Kan categories. But what is not evident is the existence of a homomorphism $u_\infty$ as \texttt{[here-]}above, and if it exists, that it is unique (or at least canonical?) as an arrow of Cat\,hot.

For the existence of a $u_\infty$, we are reduced to proving that if we have a diagram with solid arrows
\[
\begin{tikzcd}
K_n \arrow[r, "i_n", "\text{cofib.}"'] \arrow[d, "u_n"] & K_{n+1} \arrow[dr, dashed] & \\
K'_{\tau(n)} \arrow[rr, "\text{transition}"'] & & K'_N ,
\end{tikzcd}
\]
and if an arrow $u_{n+1}$ is given ``in dotted lines'' \textit{in} Cat\,hot, so as to make the square commutative,
\msp{page 48}
then we can realize it by an arrow \textit{in} Cat that makes the diagram commute for real (in Cat). We are reduced to the same problem in a situation
\[
\begin{tikzcd}
X \arrow[r, hook, shift right, "i"'] & Y , \arrow[l, dashed, shift right, "p"']
\end{tikzcd}
\]
where $X$ is a Thomason category, $i$ a Thomason cofibration, $p$ a retraction from $Y$ to $X$ \textit{in} Cat\,hot, we want to realize it by a genuine retraction (\footnote[7]{\texttt{[a ?\ \ in the margin]}}).

Admitting that this is the case, hence that every homomorphism of the ind-objects $\uuK$, $\uuK'$ \textit{in} Cat\,hot can be represented by a homomorphism between the ind-objects \textit{in} Cat, hence also by a functor $K_\infty \longrightarrow K'_\infty$, there remains the question whether the corresponding homomorphism in Cat\,hot is unique. So if we have two functors $f$, $g$ which induce the same homomorphism on the ind-represented functors on Cat\,hot, whether $f$ and $g$ are homotopic. So

\textbf{Question}. Two functors $f, g : K_\infty \longrightarrow K'_\infty$ such that for every $n \in \mathbf{N}$, $f_n$, $g_n$ obtained from them by restriction to $K_n$ are homotopic (and even homotopic as functors from $K_n$ to a suitable $K'_N$), are they homotopic? (\footnote[8]{$K$, $K'$ Thomason categories, the $K_n$, $K'_n$ too.})

\printedpage{31}
% head: A. Grothendieck : Dérivateurs, Ch. V, Édition des Universités de Montpellier II et Paris VII
% foot: Version 9 septembre 2001
\msp{page 49}
I presume not - that it can happen that the minimum required length $l(n)$ for a path in $\uHom(K_n, K'_{\tau(n)})$ joining $f_n$ and $g_n$, tends to infinity if $n \longrightarrow \infty$. Does the hypothesis that the $K'_n$ are Thomason categories change anything here (are two homotopic objects in $\uHom(X, K')$, if $K'$ is Thomason, joined by an arrow?)

Of course, one can claim that the categories $\uCh_{\Sigma,n}(\mathcal{M}; a, b)$ have been defined ``canonically'', and their direct limit $\uCh_{\Sigma}(\mathcal{M}; a, b)$ too. But the definition was rather arbitrary and outlandish - I am convinced that ten mathematicians, proposing to develop this same idea of categories of paths, with the same aim as me, would arrive at ten different definitions! Yet it is not ruled out that we will end up finding a transitive system of isomorphisms in Cat\,hot, between the categories of ``infinite'' paths (i.e. of unspecified length) $K_\infty = K_\infty(\mathcal{M}; a, b)$ that we will have constructed, with the
\msp{page 50}
main aim (let us say) of representing the functor $H$ \textit{on finite categories.} And that the isomorphisms obtained will, of course, be homomorphisms in Ind(Cat\,hot). I have the impression that I have not yet thoroughly axiomatized the very notion of paths, which after all has a fairly obvious and important geometric meaning in its own right. So by what is to be understood by ``\textit{a} theory of paths in a category''. Such a theory $\mathcal{T}$ would thus associate to $\mathcal{M}$, $\Sigma$, $a$, $b$ a direct system of categories $\Ch^{\mathcal{T}}_{n,\Sigma}(\mathcal{M}; a, b)$. And the natural geometric constraints to be imposed on such theories, also with a view to their use in making explicit
\[
\Hom_{\mathcal{M}(X)\Sigma_X^{-1}}(a_X, b_X) \xrightarrow{\ \sim\ } \pi_0\uHom(X, \Chh^{\mathcal{T}}(\mathcal{M}; a, b)) % reading resolved: two underlines, the ind-path object.
\]
would give a transitive system of isomorphisms not only between the ind-objects of Cat\,hot, but perhaps also between the $\Ch^{\mathcal{T}}_{\infty,\Sigma}(\mathcal{M}; a, b)$, for different theories. If we have such a statement at our disposal, the category $\Ch_{\infty,\Sigma}(\mathcal{M}; a, b)$, or more exactly its images in Cat\,hot and Hot, % [sic] differentes -> différentes, missing accent
\msp{page 51}
would take on a convincing meaning.

On the other hand there is another point of view, which perhaps could ``canonify'' the categories $\uCh_\infty$ as an object of Cat\,hot and especially of Hot (one would pass to Cat\,hot by applying the functor $\psi$, Thomason-Kan envelope \ldots). Indeed, we are interested in constructing, in terms of $\mathcal{M}$, $\Sigma$, a ``derivator'' (or a prederivator) $\mathbf{D}_{\mathcal{M},\Sigma}$ especially on the ``domain'' of finite categories. Extending it to a larger domain is perhaps important in a ``second wind'', but can be considered a luxury in the first movement of founding a theory of derivators. This said, we can ask ourselves questions of the following type.

Let
\[
\Hot_0 \subset \Hot
\]
\printedpage{32}
% head: A. Grothendieck : Dérivateurs, Ch. V, Édition des Universités de Montpellier II et Paris VII
% foot: Version 9 septembre 2001
be the full subcategory of Hot, defined by finite categories (\footnote[9]{\textit{NB} Hot$_0$ is equivalent to a small category.}) - let us call the homotopy types ``of finite type'' (or of finite \textit{presentation}?). We have an obvious functor
\msp{page 52}
\[
\begin{aligned}
\Hot \quad &\longrightarrow \quad \Hot_0^{\wedge} = \uHom(\Hot_0^{o}, (\Ens)) \, ,\\
X \quad &\longmapsto \quad \tilde{X} = (\xi \longmapsto \Hom_{\Hot}(\xi, X)).
\end{aligned}
\]
This functor is \textit{conservative} \texttt{[this seems false]}, since a homomorphism $X \longrightarrow X'$ in Hot is an isomorphism if it induces isomorphisms for the functors $\tilde{X} \longrightarrow \tilde{X}'$, and even just for their values on the spheres $S^n$ ($n \geq 0$). (Unless I am mistaken. Are there any blunders, because the criterion for weak equivalence in terms of the $\pi_i$ does not apply directly, in the absence of base points. That should not be very serious \ldots\ \texttt{[this in fact seems serious]}). But is it perhaps even \textit{fully faithful} (\footnote[10]{\texttt{[a ?\ \ in the margin]}})? So that Hot would be realized as a full subcategory of $\Hot_0^{\wedge}$, containing the representable functors?


\endgroup
% END UNIT Der-ChV-234/Der-ChV-234_p031-034.tex
\clearpage\setcounter{footnote}{0}\phantomsection\addcontentsline{toc}{section}{Chapter VI}\markboth{Chapter VI}{Chapter VI}
% BEGIN UNIT Der-ChVI-345/Der-ChVI-345_p001-005.tex
\begingroup\makeatletter%
\NoAutoSpacing%
\newcommand{\printedpage}[1]{\marginpar{\footnotesize p.~#1}}%
\newcommand{\msp}[1]{\marginpar{\footnotesize [#1]}}%
\newcommand{\deq}{\overset{\mathrm{def}}{=}}%
\makeatother\let\printedpage\derpp\let\msp\dermsp%
\let\mspin\DERmspin%

\DERmark{ENV0304}


% ===================== page 1 of the PDF (title page) =====================
\begin{center}
\printedpage{[no number]}% no page number printed on this page
% head: (none)

\vspace*{6cm}

{\Large\textsf{THE DERIVATORS}}

\bigskip

\textbf{ALEXANDRE GROTHENDIECK}

\vspace{2cm}

\textbf{Chapter VI}

\medskip

\textbf{HOT}
\end{center}

\par

% ===================== page 2 of the PDF (editorial note) =====================
\vspace*{6cm}

\noindent\printedpage{[no number]}% no page number printed on this page
% head: (none)
\par\noindent\textbf{Note by the editors of the transcribed edition.}\par
This text was deciphered and transcribed in \LaTeXe{} by M. Künzer. It was edited by M. Künzer, J. Malgoire, and G. Maltsiniotis. The transcription is as faithful as possible to the manuscript. For the few obvious corrections, or rare comments by the editors, as well as for the original numbering of the manuscript pages, typewriter characters \texttt{[typewriter]} in brackets are used. A question mark in brackets means that one is not sure of the preceding word.

\noindent This edition is provisional. Remarks, comments and corrections are welcome. Send a message to:

\noindent\textbf{[e-mail address removed]}

\hfill G. Maltsiniotis\hspace*{2cm}

\par

% ===================== printed page 1 =====================
\noindent\printedpage{1}
% head: A. Grothendieck : Dérivateurs, Ch. VI, Édition des Universités de Montpellier II et Paris VII
% foot: Version 9 novembre 2001
\msp{page 1}

\section*{\normalsize The derivator \textmd{HOT}}

\begin{flushright}
14.11.90
\end{flushright}

\begin{flushright}
(The first sheets, from two weeks ago, have been ditched, overtaken by VII)
\end{flushright}

\noindent I define Hot starting from the category of models
\[
\mathcal{M} = \mathrm{Cat},\ W \subseteq \mathrm{Fl}(\mathrm{Cat}),
\]
$W$ being in principle the Hot-equivalences. But as far as possible I am going to axiomatize more, assuming given a more general set $W$, called ``basic localizer'' (basic localizer), e.g. $W = W_{\mathbf{D}}$, where $\mathbf{D}$ is an already given derivator (but assuming that its domain is all of (Cat)). We are going to note as we go along which properties of $W$ we will need. For purposes of later reference, I am going to recall here those I isolated in Pursuing Stacks.

\noindent\textbf{L1}\ \ Saturation of $W$

\begin{itemize}[leftmargin=2.5em]
\item[a)] \textit{The identities} \texttt{[are]} \textit{in $W$.}
\item[b)] \textit{If two among $gf$, $f$, $g$} \texttt{[are]} \textit{in $W$, the third also, $X \xrightarrow{f} Y \xrightarrow{g} Z$.}
\item[c)] \textit{If $gf$ and $fg$ are in $W$, then $f$ and $g$ too, $X \overset{f}{\underset{g}{\rightleftarrows}} Y$ .}
\end{itemize}

\noindent\textbf{L2}\ \ $\boldsymbol{\Delta}^1$ \textit{is universally $W$-aspheric, i.e. $\boldsymbol{\Delta}^1 \times X \longrightarrow X$ is in $W$ for every $X$ in} Cat. This implies

\noindent\textbf{L2 bis}\ \ $\boldsymbol{\Delta}^1$ \textit{is $W$-aspheric.} (This was L3$'$ in Pursuing Stacks.)

\noindent\textbf{L3}\ \ \textit{If $X$ has a final object, $X$ is $W$-aspheric.} (\footnote[1]{There is an L3$'$ dual to L3.})

\noindent I note here \texttt{[some]} strengthened variants.

\noindent\textbf{L3 bis}\ \ \textit{If $X$ is} Hot-\textit{aspheric, it is $W$-aspheric.}

\noindent\textbf{L3 ter}\ \ \textit{If $X$ is} Hot-\textit{aspheric, it is universally $W$-aspheric, i.e. $X \times Y \xrightarrow{\mathrm{pr}_2} Y$ is in $W$ for every $Y$ in} Cat. (\footnote[2]{Suffices to imply L2, L3.})

\par

% ===================== printed page 2 =====================
\noindent\printedpage{2}
% head: A. Grothendieck : Dérivateurs, Ch. VI, Édition des Universités de Montpellier II et Paris VII
% foot: Version 9 novembre 2001
\textbf{L3 quater}\ \ \textit{Every} Hot-\textit{equivalence is in $W$.}

\noindent L3 quater $\Longrightarrow$ L3 ter $\Longrightarrow$ L3 bis $\Longrightarrow$ L3.

\noindent L3 ter $\Longrightarrow$ L2 $\Longrightarrow$ L2 bis.

\noindent\msp{page 2}

\begin{itemize}[leftmargin=2.5em]
\item[\textbf{L4}] \textit{If $f : X \longrightarrow Y$ is $W$-aspheric (i.e. the $X/y \longrightarrow Y/y$ in $W$, i.e. (if we have L3) if the $X/y$ are $W$-aspheric), then $f$ too.} (Localization axiom.) (\footnote[3]{there is an L4$'$ dual to L4.}) \texttt{[``also''} \texttt{= in} $W$\texttt{]}

\item[\textbf{L5}] (Fibration axiom). \textit{If}
\[
\begin{tikzcd}[column sep=small]
X \arrow[rr, "f"] \arrow[dr] & & Y \arrow[dl] \\
& S &
\end{tikzcd}
\]
\textit{with $X$ and $Y$ fibered over $S$ and $f$ cartesian, if for every $s \in S$ $f_s : X_s \longrightarrow Y_s$ is in $W$, $f$ too.}
\end{itemize}

\noindent In Pursuing Stacks I always assumed \fbox{L1, L3, L4} (implying L2) - it was these conditions that defined the notion of ``localizer'' (basic) in Cat. (\footnote[4]{If $W$ \texttt{[is]} a localizer, then contractible categories are $W$-aspheric, homotopisms are in $W$, homotopic maps are equal in (Cat)$W^{-1} \deq \mathrm{Hot}_W$.}) On the other hand, \fbox{L1, L2 bis, L5} imply L1,2,3,4,5.

\noindent\textit{We are going to assume here too that we have at least} L1, L3, L4 (i.e. $W$ a localizer). Applying the construction of I, p.~72-\textbf{75,} we see that $f : X \longrightarrow Y$ is in $W$ if and only if $f^o : X^o \longrightarrow Y^o$ is. So L1, L3, L4 imply L3$'$, L4$'$ (axioms dual to L3, L4).

\noindent There were supplementary axioms

\begin{itemize}[leftmargin=2.5em]
\item[\textbf{La}] $f \in W \Longrightarrow \pi_0(f)$ \textit{bijective.}

\item[\textbf{Lb}] \textit{The canonical localization functor}
\[
\mathrm{Cat} \longrightarrow (\mathrm{Cat})W^{-1} \deq \mathrm{Hot}_W
\]
\textit{commutes with finite sums.}

\item[\textbf{Lc}] \textit{Likewise for finite products.}
\end{itemize}

\par

% ===================== printed page 3 =====================
\noindent\printedpage{3}
% head: A. Grothendieck : Dérivateurs, Ch. VI, Édition des Universités de Montpellier II et Paris VII
% foot: Version 9 novembre 2001
One should all the same isolate concrete conditions on $W$ for the validity of Lb, Lc. The conditions L3 quater and La are written respectively
\[
W_\infty \subseteq W,\ \ W \subseteq W_0.
\]

\noindent\msp{page 3}

\noindent Their conjunction is the very reasonable axiom
\[
W_\infty \subseteq W \subseteq W_0.
\]
Of course, axiom L5 now looks a bit cramped, i.e. feeble. We are going to see very soon what we will need!

\noindent We thus want to define a derivator
\[
\mathrm{HOT}_W
\]
starting from the category of models
\[
(\mathrm{Cat}, W \subseteq \mathrm{Fl}(\mathrm{Cat})).
\]
So to define
\[
\mathrm{HOT}_W(S) \qquad\qquad\qquad (S \in \mathrm{Ob}\,\mathrm{Cat}),
\]
the general recipe is to take the category % [sic] récette -> recette, misspelling of "recette"
\[
\mathcal{M}(S) = \mathrm{Cat}(S) = \underline{\mathrm{Hom}}(S^o, (\mathrm{Cat})),
\]
to introduce in it the set of arrows
\[
\left\{
\begin{aligned}
& W_S \subseteq \mathrm{Fl}(\mathcal{M}(S)) \\
& u \in W_S \Longleftrightarrow u_s \in W \ \ \forall\, s \in \mathrm{Ob}\,S.
\end{aligned}
\right.
\]
We must set
\[
\mathrm{HOT}_W(S) \deq \underline{\mathrm{Hom}}(S^o, \mathrm{Cat})W_S^{-1}.
\]
But the category $\underline{\mathrm{Hom}}(S^o, \mathrm{Cat})$ is interpreted, up to canonical isomorphism, as that
\[
\underline{\mathrm{Fibsc}}\,S
\]
of categories $X$ Cat-fibered over $S$, equipped with a ``splitting'', i.e. with a choice of base change functors $u^* : X_t \longrightarrow X_s$ (for $t \xrightarrow{u} s$ in $S$) \texttt{[rather} $u^* : X_s \longrightarrow X_t$\texttt{]}, with moreover \textit{strict transitivity}
\[
(vu)^* = u^* v^* \qquad\qquad\qquad \text{if}\ \ t \xrightarrow{u} s \xrightarrow{v} r\ \ \text{in}\ S
\]

\noindent\msp{page 4}


\endgroup
% END UNIT Der-ChVI-345/Der-ChVI-345_p001-005.tex
% BEGIN UNIT Der-ChVI-345/Der-ChVI-345_p006-010.tex
\begingroup\makeatletter%
\newcommand{\printedpage}[1]{\marginpar{\footnotesize p.~#1}}%
\newcommand{\msp}[1]{\marginpar{\footnotesize [#1]}}%
\newcommand{\ul}[1]{\underline{\mathrm{#1}}}%
\newcommand{\xline}[1]{\mathrel{\overset{#1}{\rule[0.6ex]{2em}{0.4pt}}}}%
\makeatother\let\printedpage\derpp\let\msp\dermsp%
\let\mspin\DERmspin%


\printedpage{4}
% head: A. Grothendieck : Dérivateurs, Ch. VI, Édition des Universités de Montpellier II et Paris VII
and not only transitivity up to canonical isomorphism. Thus one can also write
\[
\mathrm{HOT}_W(S) \xrightarrow{\ \sim\ } (\ul{Fibsc}\,S)W_S^{-1}, \tag*{$(*)$}
\]
where this time $W_S$ is interpreted as
\[
W_S \subseteq \mathrm{Fl}(\ul{Fibsc}\,S),
\]
namely
\[
W_S = \{u \in \mathrm{Fl}(\ul{Fibsc}\,S) \mid u_s \in W \text{ for all } s \in \mathrm{Ob}\,S\}
\]
(strictly cartesian functors, i.e. commuting with the splittings, which induce $W$-equivalences between the fiber categories).

The arrow in $(*)$ is indeed an isomorphism of categories, not just an equivalence. But we are going to define other categories, more convenient in certain respects, which will be \textit{equivalent} (we hope), but certainly not isomorphic, to $\mathrm{HOT}_W(S)$. (We shall deal later with the variance of $\mathrm{HOT}_W(S)$ in $S$.)

The categories $\ul{Hom}(S^o,\mathrm{Cat})$, $\ul{Fibsc}\,S$, as well as the variants we are now going to introduce, are in fact 2-categories. But to describe $\mathrm{HOT}_W(S)$, only the 1-category structures are used. One has the impression that this is quite % [sic] seule les structures -> seules les structures, accord du pluriel
\msp{page 5}
crude, that one loses something in doing so. But what is lost will be recovered, and superabundantly, in the structure of the internal \underline{Hom} in the $\mathrm{HOT}(S)$, if one ``remembers'' that these internal \underline{Hom}, giving rise to ``relative homotopy types'' $\ul{Hom}_{\mathrm{HOT}_W}(S;\xi,\eta)$ (if $\xi$, $\eta$ belong to them), can be interpreted in terms of $\infty$-stacks in groupoids over $S$, with all the richness of structure that belongs to them. But this is another story\dots

In fact, I see a diagram of 13 categories of ``categories over $S$\,'', to be examined now % [sic] categories -> catégories, accent manquant
\[
\begin{tikzcd}[column sep=small]
 & \ul{Fibsc}_0S \arrow[dl, hook] \arrow[d, "\mathrm{fid}"] & & \ul{Cofibsc}_0S \arrow[d] \arrow[dr, hook] & \\
\ul{Fibsc}\,S \arrow[d, "\mathrm{fid}"] & \ul{Fib}_0S \arrow[dl, hook] & & \ul{Cofib}_0S \arrow[dr, hook] & \ul{Cofibsc}\,S \arrow[d, "\mathrm{fid}"] \\
\ul{Fib}\,S \arrow[dr, hook] & & \ul{Parf}\,S \arrow[dl, hook] \arrow[d, hook] \arrow[dr, hook] & & \ul{Cofib}\,S \arrow[dl, hook] \\
 & \ul{Liss}\,S \arrow[dr, hook] & \ul{Fibhot}\,S \arrow[d, hook] & \ul{Prop}\,S \arrow[dl, hook] & \\
 & & \mathrm{Cat}/S & &
\end{tikzcd}
\]
% foot: Version 9 novembre 2001

\printedpage{5}
% head: A. Grothendieck : Dérivateurs, Ch. VI, Édition des Universités de Montpellier II et Paris VII
where \underline{Fib}, \underline{Cofib} denote the Cat-fibered, Cat-cofibered categories over $S$, \underline{Liss} and \underline{Prop} the categories
\msp{page 6}
smooth resp. proper over $S$, \underline{Parf} those that are both smooth and proper (perfect fibrations), finally \underline{Fibhot} the Hot-fibered categories over $S$. All the arrows are strictly fully faithful functors, except for the two (rather four) upper vertical arrows, which are faithful and no more. I have just added another 4 of them, categories: $\ul{Fibsc}_0$, $\ul{Fib}_0$, $\ul{Cofibsc}_0$, $\ul{Cofib}_0$ are the full subcategories of the corresponding categories without subscript, consisting of the fibered resp. cofibered categories $X$ over $S$, for which the base change resp. cobase change functors
\[
u^* : X_t \longrightarrow X_s, \qquad\quad \text{resp.} \qquad\quad u_* : X_s \longrightarrow X_t,
\]
for $s \xrightarrow{u} t$ in $S$, are Hot-equivalences. In terms of functors
\[
S^o \xrightarrow{F} \mathrm{Cat} \qquad\quad \text{resp.} \qquad\quad S \xrightarrow{F'} \mathrm{Cat},
\]
this means that this functor transforms arbitrary arrows in $S$ into elements of $W$ in Cat, i.e. that we have a commutative diagram
\[
\begin{tikzcd}[row sep=large]
S^o \arrow[r, "F"] \arrow[d, "\mathrm{can}"'] & \mathrm{Cat} \arrow[d, "\mathrm{can}"] \\
\Pi(S^o) \simeq (\Pi S)^o = S^o(\mathrm{Fl}\,S^o)^{-1} \arrow[r, "\bar F"'] & \mathrm{Hot}_W = (\mathrm{Cat})W^{-1}
\end{tikzcd}
\]
(and likewise for $F'$), i.e. that we have a ``local
\msp{page 7}
system'' $\bar F$ on $S$, with values in $\mathrm{Hot}_W$, defined using $F$. I recall that one also finds such a local system in the case of a Hot-fibered category over $S$ (and even a local system with values in Hot itself, instead of $\mathrm{Hot}_W$ - we tacitly use here that $W_\infty \subseteq W$, i.e. \fbox{L3 quater} (\footnote[5]{We shall assume it true throughout what follows.})). And the same holds a fortiori for the objects of $\ul{Parf}\,S$. For $\ul{Liss}\,S$ resp. $\ul{Prop}\,S$, one simply finds a functor
\[
S^o \xrightarrow{\bar F} \mathrm{Hot} \qquad\quad \text{resp.} \qquad\quad S \xrightarrow{\bar F'} \mathrm{Hot},
\]
whence, by composing with $\mathrm{Hot} \longrightarrow \mathrm{Hot}_W$, functors with values in $\mathrm{Hot}_W$.

In each of these 14 \texttt{[13?]} categories, we introduce a set of arrows, which will serve to localize, and we shall always denote it by the same notation $W_S$. In all these categories except only $\mathrm{Cat}/S$, $W_S$ denotes $W$-equivalence ``fiber by fiber''. In $\mathrm{Cat}/S$ on the other hand, it will be just $W$-equivalence. All the arrows of the diagram are compatible with the localizing sets of arrows.
% foot: Version 9 novembre 2001

\printedpage{6}
% head: A. Grothendieck : Dérivateurs, Ch. VI, Édition des Universités de Montpellier II et Paris VII
\msp{page 8}
\[
\begin{tikzcd}[column sep=small, row sep=large]
 & \ul{Fibsc}_0S \arrow[dl, hook, "\alpha'"'] \arrow[d] & & \ul{Cofibsc}_0S \arrow[d] \arrow[dr, hook, "\beta'"] & \\
\ul{Fibsc}S \arrow[d] & \ul{Fib}_0S \arrow[dl, hook, "\alpha"'] \arrow[ddr, dotted, bend right=20] & & \ul{Cofib}_0S \arrow[dr, hook, "\beta"] \arrow[ddl, dotted, bend left=20] & \ul{Cofibsc}\,S \arrow[d] \\
\ul{Fib}S \arrow[dr, hook] & & \ul{Parf}\,S \arrow[dl, hook] \arrow[d, hook] \arrow[dr, hook] & & \ul{Cofib}\,S \arrow[dl, hook] \\
 & \ul{Liss}\,S \arrow[dr, hook] & \ul{Fibhot}\,S \arrow[d, hook] \arrow[uuul, dotted, bend left=10, "\Phi_0"] \arrow[uuur, dotted, bend right=10, "\Psi_0"'] & \ul{Prop}\,S \arrow[dl, hook] & \\
 & & \mathrm{Cat}/S \arrow[uuull, dotted, bend left=60, "\Phi"] \arrow[uuurr, dotted, bend right=60, "\Psi"'] & &
\end{tikzcd}
\]
% ?? uncertain: starting point of the dotted arrows labelled Phi_0 and Psi_0 (the dotted lines run alongside Parf S down to Fibhot S) -- Fibhot S or Parf S

\texttt{[In this diagram, there are also} $\ul{Liss}_0$ \texttt{and} $\ul{Prop}_0$ \texttt{but without connection to the other objects.]}

I am almost sure of the

\textbf{Theorem.} (\footnote[6]{This is totally botched, cf. comments page 32 ff., solution of the perplexities pages 41 ff. See especially the recapitulatory Theorem-Scholie p. 70-81.}) \textit{All the preceding functors induce equivalences between the localized categories.} % [sic] Theorème -> Théorème, accent manquant

As there are 16 functors, it looks like an exercise in patience. The general principle is that one manages to define functors in the opposite direction, compatible with the $W_S$, in such a way that the composites are linked to the identity functors by arrows which are in $W_S$, so that by passage to the localizations these functors induce (equivalences) quasi-inverses of the functors deduced from those of the diagram. All these functors in the opposite direction arise essentially by composition, from three very familiar functors. Two of these are indicated by dotted lines on the diagram, as $\Phi$, $\Psi$. % [sic] exercise -> exercice, orthographe française ; % [sic] reliées -> reliés, accord avec « les composés »

\msp{page 9}
The third is a canonical functor
\[
C : \mathrm{Cat}/S \longrightarrow \ul{Parf}\,S,
\]
which we are also going to make explicit.
% foot: Version 9 novembre 2001

\printedpage{7}
% head: A. Grothendieck : Dérivateurs, Ch. VI, Édition des Universités de Montpellier II et Paris VII
The functor $\Phi$ (resp. $\Psi$) applied to $\ul{Fib}_0S$ (resp. $\ul{Cofib}_0S$), induces a functor
\[
\ul{Fib}_0S \longrightarrow \ul{Fibsc}_0S \qquad\quad \text{resp.} \qquad\quad \ul{Cofib}_0S \longrightarrow \ul{Cofibsc}_0S \tag*{$(*)$}
\]
(which will have to be checked). This provides all the ``inverse'' arrows we need, with the sole exception of $\alpha$, $\beta$, $\alpha'$, $\beta'$. But if we have proved, by the indicated method, that all the arrows of the diagram induce equivalences on the localized categories, to prove that the same is still true for $\alpha$, $\beta$, $\alpha'$, $\beta'$, it suffices to prove it for the two functors
\[
\ul{Fibsc}_0S \longrightarrow \ul{Fib}\,S, \qquad\qquad \ul{Cofibsc}_0S \longrightarrow \ul{Cofib}\,S,
\]
so it remains to define the functors in the opposite direction
\[
\ul{Fib}\,S \longrightarrow \ul{Fibsc}_0S, \qquad\qquad \ul{Cofib}\,S \longrightarrow \ul{Cofibsc}_0S,
\]
and as the sources are full subcategories of $\mathrm{Cat}/S$, it will suffice to define
\[
\Phi_0 : \mathrm{Cat}/S \longrightarrow \ul{Fibsc}_0S, \qquad\qquad \Psi_0 : \mathrm{Cat}/S \longrightarrow \ul{Cofibsc}_0S.
\]
For this, we shall set
\[
\Phi_0 = \Phi r C, \qquad\qquad \Psi_0 = \Psi r C, \tag*{$(**)$}
\]
\msp{page 10}
where $C : \mathrm{Cat}/S \longrightarrow \ul{Parf}\,S$, $\Phi$ and $\Psi$ are the 3 ``retrograde'' functors inverted above, and $r : \ul{Parf}(S) \hookrightarrow \mathrm{Cat}/S$ is the inclusion. We must therefore check that $\Phi_0$, $\Psi_0$ defined by $(**)$, which a priori have as target categories $\ul{Fibsc}\,S$ and $\ul{Cofibsc}\,S$, take their value in the subcategories with subscript 0.

I am going to give a single equivalence statement, which will cover (almost) all 16 cases at once:

\textbf{Corollary.} \textit{Let us consider the diagram}
\[
\begin{tikzcd}[column sep=small, row sep=small]
 & \ul{Parf}\,S \arrow[dd, "r"'] & \\
\ul{Fibsc}_0S \arrow[dr, "p"'] & & \ul{Cofibsc}_0S \arrow[dl, "q"] \\
 & \mathrm{Cat}/S &
\end{tikzcd}
\]
\textit{and let $\mathcal{F}$ be a full subcategory of $\mathrm{Cat}/S$, such that one of the functors $p$, $q$, $r$ factors through $\mathcal{F}$. Let $W_{\mathcal{F}}$ be induced on $\mathcal{F}$ by $W_S$ in $\mathrm{Cat}/S$. Then $\mathcal{F} \hookrightarrow \mathrm{Cat}/S$ induces an equivalence on the localized categories.}

\msp{page 11}
\textsc{Proof}. I say that it suffices to prove, for the three functors $p$, $q$, $r$, and for the functors in the opposite direction
% foot: Version 9 novembre 2001
\printedpage{8}
% head: A. Grothendieck : Dérivateurs, Ch. VI, Édition des Universités de Montpellier II et Paris VII
\[
\begin{tikzcd}[column sep=small, row sep=small]
 & \ul{Parf}\,S & \\
\ul{Fibsc}_0S & & \ul{Cofibsc}_0S \\
 & \mathrm{Cat}/S \arrow[uu, "C"] \arrow[ul, "\Phi_0"] \arrow[ur, "\Psi_0"'] &
\end{tikzcd}
\]
that the six composites
\[
p\Phi_0, \Phi_0p, q\Psi_0, \Psi_0q, rC, Cr
\]
are linked to the identity functors (in the relevant category) by an arrow which is in the corresponding localizer. For each of these three pairs of functors $(p,\Phi_0)$, $(q,\Psi_0)$, $(r,C)$, this will already prove that on passing to the localizations, they induce equivalences quasi-inverse to one another. But we can do better with the help of the following formal lemma.

\textbf{Lemma.} \textit{Let $(C,W)$, $(C',W')$ be two categories with localizers $W \subseteq \mathrm{Fl}\,C$, $W' \subseteq \mathrm{Fl}\,C'$. Let}
\[
p : C \longrightarrow C'
\]
\textit{be a functor compatible with the localizers, and suppose there exists a functor $\varphi : C' \longrightarrow C$ likewise compatible, and such that we have arrows (of unspecified direction)}
\[
\varphi p \xline{u} id_C, \qquad\qquad p\varphi \xline{v} id_{C'},
\]
\textit{which are in $W$ resp. $W'$ ``argument by argument''.}

\msp{page 12}
\textit{Then $p$ and $\varphi$ induce equivalences quasi-inverse to one another on passing to the localized categories (they have no merit). Moreover, for any full subcategory $C'_0$ of $C'$, such that $p$ factors (as $p_0 : C \longrightarrow C'_0$) through $C'_0$, if we denote by $W'_0$ the localizer induced in $C'_0$ by $C'$, then $p_0 : C \longrightarrow C'_0$ and $\varphi_0 = \varphi|C'_0 : C' \longrightarrow C$} \texttt{[rather} $C'_0 \longrightarrow C$\texttt{]} \textit{also induce equivalences quasi-inverse to one another on the localizations.} % [sic] merite -> mérite, accent manquant

Indeed, $\varphi_0p_0 = (\varphi i)p_0 = \varphi(ip_0) = \varphi p \xline{u} \mathrm{id}_C$ (here $i : C'_0 \hookrightarrow C'$ denotes the inclusion), and $F \overset{\mathrm{def}}{=} p_0\varphi_0 = p_0\varphi i$ satisfies
\DERmark{ENV0305}

\[
iF = \underbrace{ip_0}_{p}\varphi i = (p\varphi)i.
\]
To define $F \xline{v_0} \mathrm{id}_{C'_0}$, which is to be in $W'_0$ argument by argument,
\msp{page 13}
it suffices to define
\[
\underbrace{iF}_{=\,(p\varphi)i} \xline{\bar v_0} i,
\]
which is to be in $W'$ argument by argument, since $W'_0$ induced by $W'$ and since $F \longmapsto iF$
\[
\ul{Hom}(C'_0,C'_0) \longrightarrow \ul{Hom}(C'_0,C')
\]
% foot: Version 9 novembre 2001


\endgroup
% END UNIT Der-ChVI-345/Der-ChVI-345_p006-010.tex
% BEGIN UNIT Der-ChVI-345/Der-ChVI-345_p011-015.tex
\begingroup\makeatletter%
\newcommand{\printedpage}[1]{\marginpar{\footnotesize p.~#1}}%
\newcommand{\msp}[1]{\marginpar{\footnotesize [#1]}}%
\newcommand{\cat}[1]{\underline{\mathrm{#1}}}%
\newcommand{\circled}[1]{\tikz[baseline=(c.base)]\node[draw,circle,inner sep=1pt](c){#1};}%
\makeatother\let\printedpage\derpp\let\msp\dermsp%
\let\mspin\DERmspin%


\printedpage{9}
% head: A. Grothendieck : Dérivateurs, Ch. VI, Édition des Universités de Montpellier II et Paris VII
is fully faithful ($i$ being so). Now it suffices to take
\[
\bar{v}_0 = v * i : (p\varphi)\circ i \mathrel{\relbar\joinrel\relbar} \mathrm{id}_{C'}\circ i.
\] % ?? uncertain: the arrow between (p\varphi)\circ i and id_{C'}\circ i is printed as a plain line with no visible head -- plain line / arrow

To sum up, we have to show that the composite of the pairs of quasi-inverse functors works in the following five cases
\[
\begin{tikzcd}[column sep=large, row sep=huge]
\cat{Fibsc} \arrow[rrd, shift left] & \cat{Fibsc}_0 \arrow[l, hook'] \arrow[rd, shift left] & \cat{Parf} \arrow[d, shift left] & \cat{Cofibsc}_0 \arrow[r, hook] \arrow[ld, shift left] & \cat{Cofibsc} \arrow[lld, shift left] \\
& & \mathrm{Cat}/S \arrow[llu, shift left, "\Phi"] \arrow[lu, shift left, "\Phi_0"] \arrow[u, shift left, "C"] \arrow[ru, shift left, "\Psi_0"] \arrow[rru, shift left, "\Psi"] & &
\end{tikzcd}
\]

Besides, by duality, we are reduced to the three cases of $\cat{Fibsc}$, $\cat{Fibsc}_0$, $\cat{Parf}$. (Perhaps the case $\cat{Fibsc}_0$, with $\Phi_0$ defined in terms of $\Phi$ and $C$, will reduce to the two cases surrounding it, of $\cat{Fibsc}$ and $\cat{Parf}$.)

I see that I have to adjust my aim, first of all because I forgot that the functors
\msp{page 14}
$\Phi$, $\Psi$ I had in mind are \textit{not} compatible with the localizers \texttt{[= localizers]} - an arrow $f : X \longrightarrow Y$ of $\mathrm{Cat}/S$ such that the corresponding arrow in Cat is in $W$, is not going to induce an arrow in $\cat{Fibsc}\,S$ or $\cat{Cofibsc}\,S$ which is in $W$ on each fiber. This is not serious, for we can replace the functors by $\Phi' = \alpha'\Phi_0$, $\Psi' = \beta'\Psi_0$, which, unless I am mistaken, should work. But in view of the preceding Lemma, there is a perplexity, for I seem to have proved that if $\mathcal{F}$ is a full subcategory of $\mathrm{Cat}/S$ of which I assume only that it contains the essential image of one of the three categories $\cat{Fib}_0S$, $\cat{Cofib}_0S$, or $\cat{Parf}\,S$, then the inclusion into $\mathrm{Cat}/S$ induces an equivalence of the localized categories, taking in $\mathcal{F}$ as localizer the one induced by $\mathrm{Cat}/S$, i.e.
\begin{tikzcd}[column sep=small, row sep=small, baseline=(\tikzcdmatrixname-2-2.base)]
f : X \arrow[rr] \arrow[dr] & & Y \arrow[dl] \\
& S &
\end{tikzcd}
is in $W_{\mathcal{F}}$ if and only if as an arrow of Cat, $f$ is in $W$. Now this does indeed coincide (via a nontrivial theorem to which I shall have to return in the setting of a general basic localizer $W$) with the set of localization arrows ``fiber by fiber'', in the case of $\mathcal{F} = \cat{Parf}\,S$, $\cat{Fibhot}\,S$, $\cat{Fib}_0S$, $\cat{Cofib}_0S$.

\msp{page 15}
On the other hand, this is not at all the case on taking $\mathcal{F} = \cat{Fib}\,S, \cat{Liss}\,S$ or $\cat{Cofib}\,S$. Having chosen in these categories as localizers the notion of $W$-equivalences ``fiber by fiber'', I thus seem to be saying that I find for these four categories the \textit{same} localized ones (up to equivalence) as if I took the much coarser localizer inherited by $\mathrm{Cat}/S$. This seems hard to believe, and I shall have to be very careful in proving the key statements.

To see that the solid-line arrows of the big diagram p.~6 are all compatible with the localiz(at)ers \texttt{[sic]}, hence also the arrows composed by means of these, I already have to use this (for the case of the arrows $\cat{Liss} \longrightarrow \mathrm{Cat}/S$, $\cat{Prop} \longrightarrow \mathrm{Cat}/S$):
% foot: Version 9 novembre 2001

\printedpage{10}
% head: A. Grothendieck : Dérivateurs, Ch. VI, Édition des Universités de Montpellier II et Paris VII
\circled{1}\ \ If $X, Y$ smooth or proper over $S$, and $f : X \longrightarrow Y$ an $S$-functor, if $f_s : X_s \longrightarrow Y_s$ is in $W$ for every $s \in S$, then $f \in W$. (\footnote[7]{OK if we assume L3 quater and L4.}).

This considerably strengthens condition L5 on $W$, where we assumed that $f$ is a \textit{cartesian} functor between categories \textit{fibered} over $S$. (To sort out, for the various properties we need for $W$, to what standard properties we can reduce them. Everything will be OK if $W = W_\infty$.)

\msp{page 16}
I) \textbf{Definition of} $\Phi$, $\Psi$.

We are going to copy the canonical construction of $\cat{Ch}(X) \longrightarrow X \times X$ (or $\cat{Ch}_\infty(X) \longrightarrow X \times X$) factoring $X \xrightarrow{\mathrm{diag}} X \times X$, and the corresponding canonical factorization of an $f : X' \longrightarrow X$ as $X' \xrightarrow{i} \bar{X}' \xrightarrow{\bar{f}} Y$ \texttt{[rather} $\bar{X}' \xrightarrow{\bar{f}} X$\texttt{]}. We shall simply replace $\cat{Ch}(X)$ by the subcategory $\cat{Ch}_+(X)$, which is none other than the category of arrows of $X$
\[
\cat{Fl}(X) = \cat{Hom}(\boldsymbol{\Delta}_1, X),
\]
equipped with the functors
\[
X \xrightarrow{i_X} \cat{Fl}(X) \xrightarrow{p_X=(\sigma_X,\beta_X)} X \times X,
\]
where $i_X$ is the arrow $x \longmapsto \mathrm{id}_x$, and $\sigma_X$, $\beta_X$ are the source and target functors. The fiber of $p_X$ at $(a,b) \in \mathrm{Ob}(X \times X)$ is the discrete category defined by the set $\mathrm{Hom}(a,b)$, that depends a great deal on $(a,b)$ and not only on its connected component in $X \times X$ - so $p_X$ is nothing like a Hot-fibration. It is in general neither proper nor smooth nor anything, on the other hand the two arrows composed with $\mathrm{pr}_1$ and $\mathrm{pr}_2$ % [sic] géneral -> général, missing accent
\[
\begin{tikzcd}[column sep=small]
& \cat{Fl}(X) \arrow[ld, "\sigma_X"'] \arrow[rd, "\beta_X"] & \\
X & & X
\end{tikzcd}
\]
are nice, we have in fact

\msp{page 17}
\textbf{Lemma.} \textit{We have}
\[
\sigma_X^{-1}(\{x\}) \simeq x\backslash X, \qquad\qquad \beta_X^{-1}(\{x\}) \simeq X/x.
\]
\textit{The functor} $\sigma_X$ \textit{is} Cat-\textit{fibrant with contractible fibers (hence aspheric, hence} $W$\textit{-aspheric), the functor} $\beta_X$ \textit{cofibrant with contractible fibers (hence aspheric, hence} $W$\textit{-aspheric). They are thus also} Hot-\textit{equivalences, a fortiori} $W$\textit{-equivalences} (\footnote[8]{In fact, they are even homotopisms, cf. marginal remark below (more general - we recover the case of $\sigma_X$ by taking $X' = X$.)}).

\textbf{Corollary 1.} \textit{The functor} $i_X : X \longrightarrow \cat{Fl}(X)$ \textit{is a} Hot-\textit{equivalence} (since its composite with the Hot-equivalence $\sigma_X$ is the identity).
% foot: Version 9 novembre 2001

\printedpage{11}
% head: A. Grothendieck : Dérivateurs, Ch. VI, Édition des Universités de Montpellier II et Paris VII
Here is the general construction of the factorization of $f : X' \longrightarrow X$, corresponding to the preceding choice of a ``category of paths'' $\mathcal{C}(X) = \cat{Fl}(X)$
\[
\begin{tikzcd}
\bar{X}' \arrow[d, "\pi_f"] \arrow[rd] \arrow[rrd, "\bar{f}"] & & \\
X' \arrow[u, dotted, bend left=60, "i_f"] \arrow[rd, "f"'] \arrow[r, phantom, "\text{\footnotesize Cart.}"] & \cat{Fl}(X) \arrow[r, "\beta_X"'] \arrow[d, "\sigma_X"'] & X \\
& X &
\end{tikzcd}
\]
where $i_f$ is deduced from the section $i$ \texttt{[rather} $i_X$\texttt{]} of $\cat{Fl}(X)$ over $X$ (\footnote[9]{\textbf{NB} $\pi_f i_f = \mathrm{id}$, and $i_f\pi_f \longrightarrow \mathrm{id}$, hence $i_f$, $\pi_f$ are homotopisms inverse to each other in Cat\,hot. This makes more precise \texttt{[the]} cor.~3, page 18.}). Thus, $\bar{X}'$ is the category of pairs
\[
\{(x',u) \mid x' \in \mathrm{Ob}\,X',\ \underbrace{u}_{\in\mathrm{Fl}\,(X)} : f(x') \longrightarrow y\},
\]
the functor $\bar{f}$ is defined by
\[
\bar{f}(x',u) = \mathrm{but}(u) = y
\]
and $i_f$ by
\[
i_f(x') = (x', \mathrm{id}_{f(x')} : f(x') \xrightarrow{\mathrm{id}} f(x')),
\]
so we do indeed have
\[
f(x') = \bar{f} i_f(x').
\]

\msp{page 18}
\textbf{Corollary 2.} \textit{The functor}
\[
\bar{f} : \bar{X}' \longrightarrow X
\]
\textit{has as fibers the categories $X'/x$, it is cofibrant, the base change functors are the functors $X'/x \longrightarrow X'/y$ associated with the arrows $x \longrightarrow y$ in $X$.}

In fact, we see that this category is not only cofibered, but canonically \textit{split.}

(\footnote[10]{\textbf{Corollary 3.} \textit{The arrow $i_f : X' \longrightarrow \bar{X}'$ is an $X$-morphism (if $X'$ over $X$ by $f$, $\bar{X}'$ by $\bar{f}$), i.e. $f = \bar{f}\circ i_f$, and it is a} Hot-\textit{equivalence} (for $\pi_f : \bar{X}' \longrightarrow X'$ is Cat-fibrant, with contractible fibers, hence a Hot-equivalence).}).

Applying this to $f : X \longrightarrow S$, we find a construction
\[
\mathrm{Cat}/S \xrightarrow{\Psi} \cat{Cofibsc}\,S,
\]
it is well and truly functorial in $X$ varying in $\mathrm{Cat}/S$. If
\[
\psi : X \longrightarrow Y
\]
% foot: Version 9 novembre 2001
\printedpage{12}
% head: A. Grothendieck : Dérivateurs, Ch. VI, Édition des Universités de Montpellier II et Paris VII
in $\mathrm{Cat}/S$, then $\psi$ induces
\[
\psi/s : X/s \longrightarrow Y/s\ ,
\]
for $s \in S$, and this in a way compatible with the
\msp{page 19}
``cobase change'' functors associated with the $u : s \longrightarrow t$ in $S$
\[
\begin{tikzcd}[column sep=large, row sep=large]
X/s \arrow[r, "u_*^X"] \arrow[d, "\psi/s"'] & X/t \arrow[d, "\psi/t"] \\
Y/s \arrow[r, "u_*^Y"] & Y/t \quad .
\end{tikzcd}
\]
This therefore defines a morphism of split categories over $S$
\[
\Psi(X) \longrightarrow \Psi(Y),
\]
that's the one!

We have a morphism functorial in $X$
\[
X \xrightarrow{i_X} \Psi(X)
\]
in $\mathrm{Cat}/X$, and it is a Hot-equivalence, hence $\in W$. On the fibers, the morphism $i_X$ induces
\[
X_s \xrightarrow{i_{X,s}} X/s,
\]
the canonical inclusion.

\msp{page 20}
\textbf{Corollary 1.} \textit{If $X \xrightarrow{f} Y$ in} $\mathrm{Cat}/S$ \textit{is in $W$, then $\Psi(X) \xrightarrow{\Psi(f)} \Psi(Y)$ too.} But we shall beware of thinking that it is in $W_S$ in the sense of $\cat{Cofib}\,S$ or $\cat{Cofibsc}\,S$.

For on the fibers of $\Psi(X)$, $\Psi(Y)$ we find the arrows
\[
f/s : X/s \longrightarrow Y/s
\]
induced by $f$, which have no reason to be in $W$. Their being so means that $f : X \longrightarrow Y$ is in $W$ ``locally over S'' (\footnote[11]{cf. \circled{1} page 15.}), which (for the $W$ that interest us) is much stronger than being in $W$. (If $Y = S$, this means in fact that $p : X \longrightarrow S$ is $W$-aspheric, much stronger than $p \in W$.)

So to prove that
\[
\cat{Cofib}\,S \longrightarrow \mathrm{Cat}/S \qquad\quad \text{or} \qquad\quad \cat{Cofibsc}\,S \longrightarrow \mathrm{Cat}/S
\]
% foot: Version 9 novembre 2001
\printedpage{13}
% head: A. Grothendieck : Dérivateurs, Ch. VI, Édition des Universités de Montpellier II et Paris VII
induce equivalences for the localized categories, the functor $\Psi$ directly will not be usable. However:

\textbf{Corollary 2.} \textit{Suppose $X$, $Y$ in} $\cat{Prop}\,S$ \textit{or in} $\cat{Fibhot}\,S$\textit{, and consider}
\[
\Psi(f) : \Psi(X) \longrightarrow \Psi(Y),
\]
\textit{for this to be a $W$-equivalence in} $\cat{Cofibsc}\,S$\textit{,}
\msp{page 21}
\textit{i.e. a $W$-equivalence fiber by fiber, it is necessary and sufficient that $f$ be a $W$-equivalence fiber by fiber} (and this must imply (\footnote[12]{cf. \circled{1} page 15.}) that $f$ is a $W$-equivalence).

Indeed, in the commutative diagram
\[
\begin{tikzcd}[column sep=large, row sep=huge]
X_s \arrow[r, "f_s"] \arrow[d] & Y_s \arrow[d] \\
X/s \arrow[r, "f/s"] & Y/s
\end{tikzcd}
\]
the vertical arrows are in $W$, for $X$, $Y$ \texttt{[are]} proper over $S$ (or $X$, $Y$ in $\cat{Fibhot}\,S$), so $f_s$ is in $W$ if and only if $f/s$ is.

\texttt{[The functor]} $\Phi$ is defined dually to $\Psi$, which amounts to interchanging the roles of $\sigma_X$, $\beta_X$ in the diagram page 17.

(\footnote[13]{\textbf{Corollary 3.} \textit{If $X \in \mathit{Ob}\,$}$\cat{Fibhot}\,S$\textit{, then $\Psi(X)$ is in} $\cat{Fibsc}_0S$.\par This means in fact that the $X/s \longrightarrow X/t$ are in $W$, now the hypothesis on $X$ implies that it is in $W_\infty$. OK.})

\textbf{II) Definition of} $C : \mathrm{Cat}/S \longrightarrow \cat{Parf}\,S$

This time, it is the factorization of the arrows $p : X \longrightarrow S$ in Cat, as
\[
X \xrightarrow{j_{X/S}} \bar{X} \xrightarrow{\bar{p}} S,
\]
where $i_X$ \texttt{[rather} $j_{X/S}$\texttt{]} is a Hot-equivalence (hence in $W$), and $\bar{p}$ perfect. This is constructed formally as before, using $\cat{Ch}_\infty(S)$ instead of $\cat{Fl}(S)$, and the factorization
\[
S \xrightarrow{j_S} \cat{Ch}_\infty(S) \xrightarrow{p_S=(\sigma_S,\beta_S)} S \times S
\]
of $\mathrm{diag}_S : S \longrightarrow S \times S$. The situation is prettier, for $p_S$ is already ``perfect'', which formally implies that the $\bar{p}$ deduced from it, to factor $p : X \longrightarrow S$, is itself perfect.
% foot: Version 9 novembre 2001


\endgroup
% END UNIT Der-ChVI-345/Der-ChVI-345_p011-015.tex
% BEGIN UNIT Der-ChVI-345/Der-ChVI-345_p016-020.tex
\begingroup\makeatletter%
\NoAutoSpacing%
\newcommand{\printedpage}[1]{\marginpar{\footnotesize p.~#1}}%
\newcommand{\msp}[1]{\marginpar{\footnotesize [#1]}}%
\newcommand{\und}[1]{\underline{\mathrm{#1}}}%
\newcommand{\Cat}{\mathrm{Cat}}%
\newcommand{\circled}[1]{\tikz[baseline=(c.base)]\node[draw,circle,inner sep=1pt](c){#1};}%
\makeatother\let\printedpage\derpp\let\msp\dermsp%
\let\mspin\DERmspin%


\printedpage{14}
% head: A. Grothendieck : Dérivateurs, Ch. VI, Édition des Universités de Montpellier II et Paris VII
% foot: Version 9 novembre 2001

\msp{page 22}
Here again, the construction is functorial, we get the functor
\[ C : \Cat/S \longrightarrow \und{Parf}\,S \]
which must be composed with the functor in the opposite direction
\[ r : \Cat/S \longleftarrow \und{Parf}\,S. \]
First of all, let us check that $C$, just like $r$, is compatible with the localizers. Let $f : X \longrightarrow Y$ be in $\Cat/S$, suppose $f$ is in $W$, let us prove $C(f)$ is in $W_S$. We have commutativity in
\[
\begin{tikzcd}[row sep=large, column sep=large]
X \arrow[r, "f"] \arrow[d, "j_X"'] & Y \arrow[d, "j_Y"] \\
CX \arrow[r, "Cf"] & CY
\end{tikzcd}
\]
and $j_X$, $j_Y$ are in $W$, so $f$ is in $W$ if and only if $Cf$ is. But in the case $W = W_\infty$ we know that this is equivalent to $Cf$ being in $W_\infty$ fiber by fiber. Admitting \circled{1} p. 15 for $W$, it remains to make sure of a converse in \underline{Parf}, i.e.

\circled{2} Let
\begin{tikzcd}[baseline=(\tikzcdmatrixname-2-2.base), row sep=small, column sep=small]
X \arrow[rr, "f"] \arrow[dr] & & Y \arrow[dl] \\
& S &
\end{tikzcd}
be in $\und{Parf}\,S$ (or merely in $\und{Fibhot}\,S$). If $f$ is in $W$, then so are the $f_s : X_s \longrightarrow Y_s$.

\msp{page 23}
Let us try to prove it, by copying the proof in the case $C = C_\infty$ (XI, th.~2). We can suppose $S$ connected (\footnote[14]{serves no purpose}), using La (i.e. $W \subseteq W_0$) and admitting the

\textbf{Lemma.} \textit{Let $f : X \longrightarrow Y$ be in} Cat. \textit{Then $f \in W$ if and only if $\pi_0(f)$ is bijective, and for every connected component $X_i$ of $X$, if $Y_i$ is the corresponding connected component of $Y$, $f_i : X_i \longrightarrow Y_i$ induced by $f$ is in $W$.}

Let us then use the existence of
\[ S' \longrightarrow S \]
a Hot-fibration with $S'$ Hot-aspheric (\footnote[15]{\textbf{NB} It suffices to have a Hot-fibration $S' \longrightarrow S_0 \hookrightarrow S$, where $S_0$ is the connected component of $S$ containing the element $s \in \mathrm{Ob}\,S$ under consideration.}). We can even suppose that $S' \longrightarrow S$ is perfect. Suppose we know that base change by such perfect fibrations transforms $W$-equivalences into likewise. (That must be true for the $W_i$, but we have seen that it was false for $k$-homological equivalence, cf. I p.~87 ff.) Then $X' \xrightarrow{f'} Y'$ over $S'$ is
\printedpage{15}
% head: A. Grothendieck : Dérivateurs, Ch. VI, Édition des Universités de Montpellier II et Paris VII
% foot: Version 9 novembre 2001
a $W$-equivalence. Now for $s' \in S'$, since $\{s'\} \longrightarrow S'$ is a Hot-equivalence ($S'$ Hot-aspheric), and $X'$, $Y'$ being Hot-fibered over $S$, it follows that $X'_{s'} \longrightarrow X'$, $Y'_{s'} \longrightarrow Y'$ are Hot-equivalences, hence also $X'_{s'} \xrightarrow{f'_{s'}} Y'_{s'}$, as we see from the commutative square
\[
\begin{tikzcd}[row sep=large, column sep=large]
X'_{s'} \arrow[r, "f'_{s'}"] \arrow[d, "\wr"'] & Y'_{s'} \arrow[d, "\wr"'] \\
X' \arrow[r, "{f'}"] & Y'.
\end{tikzcd}
\]

\msp{page 24}
Since $S' \longrightarrow S$ is surjective, this shows that the $f_s : X_s \longrightarrow Y_s$ are in $W$, qed.

Thus, (2) will be a consequence of \circled{2.} Let $p : Y' \longrightarrow Y$ be a Hot-fibration (or even a perfect functor, if that can help \dots). Let $f : X \longrightarrow Y$ be an $S$-arrow (in $\Cat/S$), then $f' : X' \longrightarrow Y'$ by base change. If $f \in W$, $f' \in W'$
\[
\begin{tikzcd}[row sep=large]
X' \arrow[r, "q"] \arrow[d, "{f'}"'] \arrow[dr, phantom, "\scriptstyle\text{cart.}"] & X \arrow[d, "f"] \\
Y' \arrow[r, "p"] & Y
\end{tikzcd}
\qquad\qquad
f \in W \text{ and } p \text{ perfect} \Longrightarrow f' \in W
\]

In any case, we absolutely need \circled{2} to be able to conclude that $C$ is compatible with the localizers.

Let us consider $rC : \Cat/S \longrightarrow \Cat/S$, we have an arrow
\[
\left\{
\begin{array}{l}
\mathrm{id}_{\Cat/S} \xrightarrow{j} rC \\[1.5ex]
j_X : X \longrightarrow CX \quad \text{the canonical arrow } i_{X/S}\ \text{\texttt{[rather} } j_{X/S}\text{\texttt{]}} \text{ (p.~21),}
\end{array}
\right.
\]
and $j_X$ is always in $W$.

Let us consider $Cr : \und{Parf}\,S \longrightarrow \und{Parf}\,S$, i.e. let us apply the construction $C$ to an $X$ perfect over $S$. Then we again have
\[
\left\{
\begin{array}{l}
\mathrm{id}_{\und{Parf}\,S} \xrightarrow{j'} Cr \\[1ex]
j'_X = j_{r(X)} : X = r(X) \xrightarrow{j_X} CX.
\end{array}
\right.
\]
This morphism is in $W$, hence also a $W$-equivalence fiber by fiber, since source and target are perfect, hence Hot-fibered over $S$.

\msp{page 25}
This already proves that the functors $r$ and $C$ induce equivalences quasi-inverse to one another on the localized categories. By the lemma p.~11, this shows that for every full subcategory $\mathcal{F}$ of $\Cat/S$ containing $\und{Parf}\,S$, the inclusion of $\mathcal{F}$ in $\Cat/S$ induces an
\printedpage{16}
% head: A. Grothendieck : Dérivateurs, Ch. VI, Édition des Universités de Montpellier II et Paris VII
% foot: Version 9 novembre 2001
equivalence on the localized categories, taking as localizer in $\mathcal{F}$ the ``coarse'' localizer induced by $\Cat/S$. This applies in particular to the full subcategories
\[ \und{Liss}\,S, \quad \und{Fibhot}\,S, \quad \und{Prop}\,S, \]
but taking in $\und{Liss}\,S$ and $\und{Prop}\,S$ not the fine ``fiber by fiber'' localizer, as I had wrongly announced, but the coarse localizer induced by $\Cat/S$. (For $\und{Fibhot}\,S$, there is no risk of ambiguity). I ought to convince myself that if I take the fine $W$-equivalence (in the case $W = W_\infty$, say), I do \textit{not} find the ``good'' localized categories, i.e. the functor of localized categories, to the localized category of $\Cat/S$, is not an equivalence. This in fact supposes that in $\und{Fib}\,S$ and $\und{Cofib}\,S$ as well, it is the \textit{coarse} equivalence
\msp{page 26}
that should be taken. (In $\und{Fib}_0S$ and $\und{Cofib}_0S$, or $\und{Fibsc}_0S$ and $\und{Cofibsc}_0S$, the two equivalences are identical, since these four categories are contained in $\und{Fibhot}\,S$.) The sequel will show us whether this is indeed so. On the other hand, in $\und{Fibsc}\,S$ and $\und{Cofibsc}\,S$, there is no doubt, according to the general principles of constructing a derivator from an $(\mathcal{M}, W)$, that we must take the fine equivalence, fiber by fiber. In this connection, we must be aware that the vertical functors of the diagram p.~8
\[ \und{Fibsc}\,S \longrightarrow \und{Fib}\,S, \quad \und{Cofibsc}\,S \longrightarrow \und{Cofib}\,S, \]
though essentially surjective and faithful, are not fully faithful. They would be, and even equivalences of categories, if in $\und{Fib}\,S$ resp. $\und{Cofib}\,S$ I confined myself to the $S$-functors that are \textit{cartesian} resp. \textit{cocartesian}. So we can interpret, up to equivalence, $\und{Fibsc}\,S$ and $\und{Cofibsc}\,S$ as $\und{Fib}\,S$ and $\und{Cofib}\,S$, where we would have restricted the morphisms to just those $S$-functors that are \textit{cartesian} resp. \textit{cocartesian.} It would thus seem that by restricting the arrows in this way, but by restricting the localizer in a still more draconian way, we find equivalent localized categories. % [sic] et mêmes des équivalences -> et même des équivalences, adverb "même" printed with plural s

\msp{page 27}
That's rather strange - we'll see if it's really true!

\subsection*{III) Construction of $\Phi_0, \Psi_0 :$ \begin{tikzcd}[ampersand replacement=\&, baseline=(\tikzcdmatrixname-2-1.base), row sep=small, column sep=normal]
\& \und{Fibsc}_0S \\
\Cat/S \arrow[ur] \arrow[dr] \& \\
\& \und{Cofibsc}_0S.
\end{tikzcd}}

We must have
\[ \Phi_0(X) = \Phi r C(X) \qquad\qquad\qquad \Psi_0(X) = \Psi r C(X), \]
where
\[ r : \und{Parf}\,S \longrightarrow \Cat/S \]
is the canonical inclusion. As functors with values in \underline{Fibsc} resp. \underline{Cofibsc}, this already defines $\Phi_0$, $\Psi_0$ by
\[ \Phi_0 = \Phi r C \qquad\qquad\qquad \Psi_0 = \Psi r C, \]
\printedpage{17}
% head: A. Grothendieck : Dérivateurs, Ch. VI, Édition des Universités de Montpellier II et Paris VII
% foot: Version 9 novembre 2001
but we must see that $\Phi_0$, $\Psi_0$ take their values in the full subcategories with index 0. This is therefore equivalent to the following Lemma, applied to $Y = rC(X)$:

\textbf{Lemma.} \textit{Let $Y$ be perfect over $S$ (or merely a \textup{Hot}-fibered object over $S$). Then $\Phi(Y)$ is in $\und{Fibsc}_0S$, $\Psi(Y)$ in $\und{Cofibsc}_0S$.}

We must prove (for $\Psi(Y)$) that if $s \longrightarrow t$ is an arrow in $S$, then
\[ \Psi(Y)_s \longrightarrow \Psi(Y)_t \]
is a $W$-equivalence. Now this arrow identifies with
\[ Y/s \longrightarrow Y/t \]
and we know well that it is a Hot-equi\msp{page 28}valence, by definition of Hot-fibered objects, hence it is a $W$-equivalence.

We thus have the functors
\[
\begin{tikzcd}[column sep=large]
\und{Fibsc}_0S \arrow[r, shift left, "p"] & \Cat/S, \arrow[l, shift left, "\Phi_0"]
\end{tikzcd}
\]
we must see that they satisfy the general conditions of the lemma p.~11 (and it will be the same for $q$, $\Psi_0$). We already know that $p$ is compatible with the localizers. Let us prove it for $\Psi_0$. Let
\[ f : X \longrightarrow Y \]
in $\mathrm{Cat}/S$ be a $W$-equivalence, then $rC(f)$ is one too. It suffices to prove this: % [sic] Cat/S (upright S) -> Cat/$S$, the S printed upright unlike elsewhere
\DERmark{ENV0306}


Let
\[ q : X' \longrightarrow Y' \]
be a morphism in $\und{Fibhot}\,S$ (for example a morphism of $\und{Parf}\,S$), which is a $W$-equivalence, hence \texttt{[which induces]} $W$-equivalences on the fibers. Then $\Psi(q) : \Psi(X') \longrightarrow \Psi(Y')$ is a $W$-equivalence fiber by fiber, i.e.
\[ X'/s \longrightarrow Y'/s \]
a $W$-equivalence for every $s \in S$ (was seen in corollary 2, page 20).

\msp{page 29}
Let us consider
\[ p\Phi_0 : \Cat/S \longrightarrow \Cat/S, \]
this thus transforms $X$ into $\Phi C(X)$, regarded as an object of $\Cat/S$ and nothing more. We have a functorial morphism
\[ X \xrightarrow{j_X} C(X) \qquad\qquad\qquad j : \mathrm{id}_{\Cat/S} \longrightarrow C, \]
\printedpage{18}
% head: A. Grothendieck : Dérivateurs, Ch. VI, Édition des Universités de Montpellier II et Paris VII
% foot: Version 9 novembre 2001
hence a functorial morphism
\[ \Phi(X) \xrightarrow{\Phi(j_X)} \Phi C(X) \qquad\qquad\qquad \Phi * j : \Phi \longrightarrow \Phi_0 \]
and also a functorial morphism
\[ X \xrightarrow{i_X} \Phi(X), \]
hence on composing
\[
\begin{tikzcd}
X \arrow[r, "i_X"] \arrow[rr, bend right, "\alpha_X"'] & \Phi(X) \arrow[r, "\Phi(j_X)"] & \Phi C(X),
\end{tikzcd}
\]
i.e.
\[ \alpha : \mathrm{id}_{\Cat/S} \longrightarrow p\Phi_0. \]
For $X$ fixed in $\Cat/S$, we know that $j_X$ is in $W$, hence so is $\Phi(j_X)$ (though not fiber by fiber), cf. cor.~1 page 20. We also know that $i_X$ is in $W$, hence so is the composite $\alpha_X$.

Let us consider
\[ \Phi_0 p : \und{Fibsc}_0\,S \longrightarrow \und{Fibsc}_0\,S\ \ . \]
We start with $X$ in $\und{Fibsc}_0\,S$, we forget the splitting and the very fact that it is Cat-fibrant
\msp{page 30}
(we have applied $p$), we apply $\Phi_0$ to it, i.e. we regard $\Phi CX$ as an object of $\und{Fibsc}\,S$ (and even of $\und{Fibsc}_0\,S$). We want again to define
\[ X \xrightarrow{\beta_X} \Phi_0 p(X) = \Phi C(X) \qquad\quad \text{\textbf{in} } \und{Fibsc}_0\text{\textbf{,}} \] % ?? uncertain: weight of "Fibsc_0" after bold "dans" -- regular or bold
functorial in $X$. Now we have as before
\[ X \xrightarrow{j_X} C(X), \]
hence on applying $\Phi$
\[ \Phi(X) \xrightarrow{\Phi(j_X)} \Phi C(X), \]
but regarded this time as an arrow in $\und{Fibsc}_0\,S$ (N.B. we do have $\Phi(X)$ in $\und{Fibsc}_0\,S$, for $X$, coming from $\und{Fibsc}_0\,S$, is in $\und{Fibhot}\,S$, and we apply the remark ``Cor.~3'' page 21.) On the other hand, we have an arrow in $\Cat/S$
\[ X \xrightarrow{i_X} \Phi(X), \]
is it an arrow in $\und{Fibsc}_0(S)$, i.e. is it cartesian? We see at once that no, that we have only, for an arrow $u : s \longrightarrow t$ in $S$, a commutation homomorphism (not an isomorphism) in


\endgroup
% END UNIT Der-ChVI-345/Der-ChVI-345_p016-020.tex
% BEGIN UNIT Der-ChVI-345/Der-ChVI-345_p021-025.tex
\begingroup\makeatletter%
\newcommand{\printedpage}[1]{\marginpar{\footnotesize p.~#1}}%
\newcommand{\msp}[1]{\marginpar{\footnotesize [#1]}}%
\newcommand{\uu}[1]{\underline{\mathrm{#1}}}%
\newcommand{\circled}[1]{\tikz[baseline=(c.base)]\node[draw,circle,inner sep=1pt](c){#1};}%
\newcommand{\lline}{\mathrel{\relbar\joinrel\relbar\joinrel\relbar}}%
\makeatother\let\printedpage\derpp\let\msp\dermsp%
\let\mspin\DERmspin%


\printedpage{19}
% head: A. Grothendieck : Dérivateurs, Ch. VI, Édition des Universités de Montpellier II et Paris VII
% foot: Version 9 novembre 2001

\[
\begin{tikzcd}[row sep=large, column sep=large]
X_s \arrow[d, "i_s"] & X_t \arrow[l, "u_X^*"'] \arrow[d, "i_t"] \\
s\backslash X & t\backslash X \arrow[l, "u\backslash\mathrm{id}_X = u_{\Phi(X)}^*"']
\arrow[from=1-2, to=2-1, phantom, "{\Downarrow\alpha}" description]
\end{tikzcd}
\]
\DERmark{ENV0307}


\msp{page 31}
So I have to pack up! One can regard $i_X$ at best as an arrow of $\uu{Fib}_0\, S$, and taking the composites one finds an arrow in $\uu{Fib}_0\, S$:
\[
\begin{tikzcd}
X \arrow[r, "i_X"] \arrow[rr, bend right=25, "\beta_X(=\alpha_X)"'] & \Phi(X) \arrow[r, "\Phi(j_X)"] & \Phi C(X),
\end{tikzcd}
\]
i.e. (now regarding $\Phi_0$ as a functor with values in $\uu{Fib}_0\, S$) a homomorphism of functors
\[
\beta : \mathrm{id}_{\uu{Fib}_0 S} \longrightarrow \Phi_0 p_0, % [sic] \mathrm{id}_{\uu{Fib}_0} S -> \mathrm{id}_{\uu{Fib}_0 S}, the S is printed on the baseline, not in the subscript
\]
\DERmark{ENV0308}

where $p_0 : \uu{Fib}_0 S \hookrightarrow \mathrm{Cat}/S$ is the canonical inclusion. This $\beta$ is in $W$ argument by argument, and likewise for $\alpha$ page 29, interpreted this time as
\[
\mathrm{id}_{\mathrm{Cat}/S} \longrightarrow p_0\Phi_0.
\]
What we have proved in the end is this:

\textbf{Packing-up theorem.} \textit{Let $\mathcal{F}$ be a full subcategory of $\mathrm{Cat}/S$ which contains one of the three full subcategories $\uu{Fib}_0 S$, $\uu{Parf}\, S$, $\uu{Cofib}_0 S$, and equipping $\mathcal{F}$ with the ``coarse'' localizer $W$ (induced by that of $\mathrm{Cat}/S$). Then the inclusion}
\[
\mathcal{F} \longrightarrow \mathrm{Cat}/S
\]
\textit{induces an equivalence on the localizers.}

\msp{page 32}
On the other hand, the localized categories of the four categories of split fibrations
\[
\uu{Fibsc}\, S, \quad \uu{Cofibsc}\, S, \quad \uu{Fibsc}_0 S, \quad \uu{Cofibsc}_0 S
\]
escape our arguments. \textit{And for good reason!} It is completely open, over $S$ connected say (e.g. $S = \boldsymbol{\Delta}^1$), whether for a relative homotopy type $\xi$ over $S$, the induced types $\xi_s$ ($s \in S$) form a local system on $S$, with values in $\mathrm{Hot}_W$, as would be the case if the theorem announced at the beginning were true. On the contrary, \textit{every} arrow of $\mathrm{Hot}_W$ must come from an object of $\mathrm{Hot}_W(\boldsymbol{\Delta}^1)$. The localizations of the 9 other categories of the diagram p. 8, and of all those discussed in the theorem we have just proved, must correspond ``visibly'' to the full subcategories of $\mathrm{HOT}_W(S)$ formed by the locally constant HOT-coefficients on $S$.

Also, I realize that it was absurd to expect $\uu{Fibsc}\, S$ and $\uu{Cofibsc}\, S$ to give canonically equivalent localized categories. For the elements $\xi$ of the first
\printedpage{20}
% head: A. Grothendieck : Dérivateurs, Ch. VI, Édition des Universités de Montpellier II et Paris VII
% foot: Version 9 novembre 2001
\msp{page 33}
of the two categories each give rise to a contravariant functor on $S$
\[
\bar{\xi} : S^o \longrightarrow \mathrm{Hot}_W,
\]
those $\eta$ of the second (which is none other than $\uu{Fibsc}(S^o)$) each give rise to a \textit{covariant} functor
\DERmark{ENV0309}

\[
\bar{\eta} : S \longrightarrow \mathrm{Hot}_W,
\]
which is not the same! It is only ``the same'' if these are local coefficients, which means exactly that $\xi$ is in $\uu{Fibsc}_0 S$ resp. $\eta$ in $\uu{Cofibsc}_0 S$. So one does indeed expect \circled{$1^o$} that the arrows
\[
\uu{Fibsc}_0 S \longrightarrow \uu{Fib}_0 S, \quad \uu{Cofibsc}_0 S \longrightarrow \uu{Cofib}_0 S
\]
induce equivalences on the localizations (something I was unable to prove a little while ago) (\setcounter{footnote}{15}\footnote{The quasi-inverses being well and truly given by $\Phi_0$, $\Psi_0$ - for [?] by what else then ???}). On the other hand, it is worse than unreasonable to expect the same to hold for the functors
\[
\uu{Fibsc}_0 S \longrightarrow \uu{Fibsc}\, S, \quad \uu{Cofibsc}_0 S \longrightarrow \uu{Cofibsc}\, S,
\]
that should give by localization the fully faithful inclusion functors, of the categories of locally constant objects into $\mathrm{HOT}_W(S)$ resp. into $\mathrm{HOT}_W(S^o)$. And likewise for
\DERmark{ENV0310}

\msp{page 34}
the functors induced on the localizations by
\[
(*) \qquad \uu{Fibsc}\, S \longrightarrow \uu{Fib}\, S, \quad \uu{Cofibsc}\, S \longrightarrow \uu{Cofib}\, S.
\]
Since the localizations of the target categories of these two functors look very much as though they identify with the preceding subcategories of locally constant objects, it would therefore seem that: \circled{$2^o$} The preceding functors define retractions of the category $\mathrm{HOT}_W(S)$ resp. $\mathrm{HOT}_W(S^o)$ onto the subcategory of locally constant objects on $S$, the localized functors of the inclusion functors
\DERmark{ENV0311}

\[
\uu{Fibsc}_0 S \longrightarrow \uu{Fibsc}\, S, \quad \uu{Cofibsc}_0 S \longrightarrow \uu{Cofibsc}\, S
\]
being fully faithful and having as essential images the subcategories in question.

This retraction seems to me a wholly mysterious thing. This is tied to the fact that for the very special derivator $\mathbf{D} = \mathrm{HOT}$, when one has a $\xi \in \mathbf{D}(S)$, one deduces from it (up to $\mathrm{HOT}/S$-equivalence) a type of diagram $X$ over $S$,
\msp{page 35}
which one can then regard (via the construction $C$) as defining a relative homotopy type \textit{locally constant} over $S$.

\printedpage{21}
% head: A. Grothendieck : Dérivateurs, Ch. VI, Édition des Universités de Montpellier II et Paris VII
% foot: Version 9 novembre 2001

So to elucidate the situation, i.e. the relations between the 13 localized categories coming from the diagram on page 8, it remains to see whether the hypothetical assertions I have just made are indeed valid, namely

\begin{enumerate}[leftmargin=*]
\item[\circled{$1^o$}] The functor deduced from
\[
\uu{Fibsc}_0 S \longrightarrow \uu{Fib}_0 S
\]
by passing to the localizations is an equivalence of categories (and in this question one can replace $\uu{Fib}_0 S$ by any other full subcategory of $\mathrm{Cat}/S$ containing it - it amounts to the same thing), and

\item[\circled{$2^o$}] The functor deduced from
\[
\uu{Fibsc}_0 S \longrightarrow \uu{Fibsc}\, S
\]
by localization is fully faithful. (That its image is exactly made up of the locally constant objects of the first localization, interpreted as $\mathrm{HOT}_W(S)$, is tautological.)
\end{enumerate}

There is no longer any need to worry about the side of
\msp{page 36}
cofibered categories, since all one will be able to say about it comes down to what I have just spelled out, for $X^o$ instead of $X$.

In the direction of \circled{$1^o$}, our preceding arguments by dint of $\Phi_0$, $p$ already give us a useful piece of information. But it will be more convenient now to drop $\uu{Fib}_0 S$ or rather, to replace it by $\uu{Fib}\, S$ (which is legitimate, we said so in \circled{$1^o$}), so that we now have a commutative diagram (solid lines)
\[
\begin{tikzcd}
\uu{Fibsc}\, S \arrow[dr, "j"'] & & \uu{Fibsc}_0 S \arrow[ll, hook', "i"'] \arrow[dl, "k"'] \\
& \uu{Fib}\, S, \arrow[ur, dotted, bend right=40, "\Phi_0=\psi"'] &
\end{tikzcd}
\]
where $i$ is a (strictly) fully faithful inclusion, $j$ (hence also $k$) faithful. Finally, we have a canonical functor $\psi$ ($= \Phi_0$) (dotted). The localizers on the top line are the fine localizers (for $\uu{Fibsc}_0 S$ this is the same as coarse $W$-equivalence), on $\uu{Fib}\, S$ % [sic] les localiseur fins -> les localiseurs fins, missing plural s
on the other hand we take coarse $W$-equivalence. (We know that the localization is the same, up to equivalence, as if one restricts to the full subcategory $\uu{Fib}_0 S$, which is quite remarkable). % [sic] "(On sait ... remarquable." -> "(On sait ... remarquable.)", the opening parenthesis is never closed
\DERmark{ENV0312}


\msp{page 37}
All four functors $i, j, k, \Phi_0$ are compatible with the localizers, hence pass to the localizations
\printedpage{22}
% head: A. Grothendieck : Dérivateurs, Ch. VI, Édition des Universités de Montpellier II et Paris VII
% foot: Version 9 novembre 2001
\[
\begin{tikzcd}
\mathrm{HOT}_W(S) \arrow[dr, "\bar{j}"'] & & \mathrm{HOT}^{?}_W(S) \arrow[ll, "\bar{\imath}"'] \arrow[dl, "\bar{k}"'] \\
& \mathrm{HOT}^{\mathrm{lc}}_W(S), \arrow[ur, bend right=40, "\bar{\psi}"'] &
\end{tikzcd}
\]
where the superscript ? in the notation means that I do not know how to interpret this category, unless I prove that $\bar{k}$ is an equivalence, and the superscript lc in the bottom category should suggest the intuition that in any case we are dealing with a sort of ``relative homotopy types'' (in a modified sense) \textit{locally constant} on $S$. The problem is to understand this diagram. The natural conjecture is % [sic] dois suggerer -> doit suggérer, verb agreement and missing accent

\begin{enumerate}[leftmargin=*]
\item[$1^o$)] $\bar{k}$ is an equivalence. (Which implies that $\bar{\psi}$ is the quasi-inverse equivalence, that $\bar{k}$, $\bar{\psi}$ essentially allow one to identify $\mathrm{HOT}^{?}_W$ and $\mathrm{HOT}^{\mathrm{lc}}_W$, so that $\bar{\imath}$ and $\bar{\jmath}$ become functors between $\mathrm{HOT}_W$ and $\mathrm{HOT}^{\mathrm{lc}}_W$.)

\item[$2^o$)] $\bar{\imath}$ is fully faithful (NB the essential image is well understood in any case.)
\end{enumerate}

\msp{page 38}
In the direction of $1^o$), I know for now that one has
\[
\bar{k}\bar{\psi} \simeq \mathrm{id}_{\mathrm{HOT}^{\mathrm{lc}}_W(S)},
\]
so $\bar{k}$ is essentially surjective (on objects), and induces \textit{surjections} on the $\mathrm{Hom}(-,-)$ between objects in the image of $\bar{\psi}$, while $\bar{\psi}$ is injective on the isomorphism classes of objects, and on the $\mathrm{Hom}(-,-)$ (i.e. \textit{faithful}).
\DERmark{ENV0313}


If the relations between $\mathrm{HOT}^{?}_W$ and $\mathrm{HOT}^{\mathrm{lc}}_W$ turn out to be impossible to make head or tail of (e.g. $\bar{\psi}$ not even fully faithful, which implies that $\bar{k}$ is not even faithful), then the category $\mathrm{HOT}^{?}_W$ deserves to be dumped in the wastebasket, and we will keep only
\DERmark{ENV0314}

\[
\begin{tikzcd}[row sep=large]
\mathrm{HOT}_W(S) \arrow[d, shift right=1.5, "\bar{j}"'] \\
\mathrm{HOT}^{\mathrm{lc}}_W(S), \arrow[u, shift right=1.5, "\bar{\psi}'"']
\end{tikzcd}
\]
where we define $\bar{\psi}'$ as
\[
\bar{\psi}' = \bar{\imath}\bar{\psi} = \overline{i\psi}.
\]
And the question (more modest) now becomes:

\begin{enumerate}[leftmargin=*]
\item[$3^o$)] Is $\bar{\psi}'$ fully faithful? (The \textit{essential} image is the same as that of $\uu{Fib}_0$ ($\approx$ $\uu{Fib}$), it is made up, unless I am mistaken, of the locally constant objects - to be checked if need be \dots)
\end{enumerate}

\printedpage{23}
% head: A. Grothendieck : Dérivateurs, Ch. VI, Édition des Universités de Montpellier II et Paris VII
% foot: Version 9 novembre 2001

\msp{page 39}
It is hard to imagine that $\bar{k}$, or what amounts to the same thing $\bar{\psi}$, could be an equivalence, without there being at least a ``$W_S$-homotopy''
\[
\mathrm{id}_{\uu{Fibsc}_0 S} \ \leftrightarrow\ \leftrightarrow^{h}\!\leftrightarrow\ \leftrightarrow\ \psi k
\]
in the category of endofunctors of $\uu{Fibsc}_0 S$. To tell the truth, I now realize that the description of $\bar{\psi}$ starting from $\psi : X \longmapsto \Phi_0 X = \Phi C X$ is needlessly complicated, for we had made the detour through $C(X)$ solely to be sure of landing in $\uu{Fibsc}_0 S$, not just in $\uu{Fibsc}\, S$. But if we drop $\uu{Fibsc}_0 S$, we don't give a damn, and the functor $\Phi : X \longmapsto \Phi(X)$ is much more natural. And if we do not drop it, there is no harm in replacing the source $\uu{Fib}\, S$ by $\uu{Fib}_0 S$ (which has the same localization, up to equivalence), and on $\uu{Fib}_0 S$ again there is no need to burden ourselves with $CX$, for the pleasure of applying $\Phi$ to it, since $\Phi(X)$ itself is already in $\uu{Fibsc}_0 S$.

So, the question is this: If $X$ varies
\msp{page 40}
% layout: the arrows nearrow, uparrow, nwarrow are printed under each "Fibsc_0 S", pointing up to it; the arrow labelled c from X to Phi(X) is a long line with its arrowhead printed in the middle
\[
\begin{array}{l}
\text{in } \underset{\nearrow\uparrow\nwarrow}{\uu{Fibsc}_0 S}, \quad \text{is there a homotopy in } \underset{\nearrow\uparrow\nwarrow}{\uu{Fibsc}_0 S}, \quad \text{``natural'', i.e. functorial in } X, \\[2ex]
\qquad\qquad X \xrightarrow{\qquad c \qquad} \Phi(X) = \ \text{Split fibered category with fibers the } s\backslash X
\end{array}
\]
But a path which is made up of arrows
\[
X = X_0 \lline X_1 \lline X_2 \cdots \lline X_n = \Phi(X)
\]
which are all cartesian (and moreover, of course, in $W$!). This therefore means that we would have split fibered categories over $S$, $X_1$, \dots, $X_{n-1}$, depending \textit{functorially} on $X$, which connect
\[
X = X_0 \qquad \text{to} \qquad \Phi(X) = X_n
\]
i.e.
\[
X_s \qquad \text{to} \qquad s\backslash X \qquad \text{???}
\]
I see nothing on the horizon.


\endgroup
% END UNIT Der-ChVI-345/Der-ChVI-345_p021-025.tex
% BEGIN UNIT Der-ChVI-345/Der-ChVI-345_p026-030.tex
\begingroup\makeatletter%
\newcommand{\printedpage}[1]{\marginpar{\footnotesize p.~#1}}%
\newcommand{\msp}[1]{\marginpar{\footnotesize [#1]}}%
\makeatother\let\printedpage\derpp\let\msp\dermsp%
\let\mspin\DERmspin%


\printedpage{24}
% head: A. Grothendieck : Dérivateurs, Ch. VI, Édition des Universités de Montpellier II et Paris VII

\msp{page 41}

\begin{flushright}
15.11.90
\end{flushright}

I think I finally understood, last night, the exact situation. There is reason, in Cat$/S$ and in the other categories of the diagram p. 8 (all mapping to Cat$/S$) to consider four sorts of ``$W$-equivalences''. % [sic] qu j'ai -> que j'ai, letter missing

\begin{itemize}
\item[a)] The \textit{``coarse'' $W$-equivalence} : \begin{tikzcd}[baseline=(s.base), row sep=small, column sep=small]
f : X \arrow[rr] \arrow[dr] & & Y \arrow[dl] \\
& |[alias=s]| S &
\end{tikzcd} in Cat$/S$ is a coarse $W$-equivalence if as an arrow of Cat, it is a $W$-equivalence. (\footnote[17]{stable under base change both smooth and proper, more generally if $S' \longrightarrow S$ is a Hot-fibration.}).

\item[b)] The \textit{``fiberwise'' $W$-equivalence} (over $S$) : $\forall\ s \in S$, $f_s : X_s \longrightarrow Y_s$ is a $W$-equivalence. (\footnote[18]{stable under any base change $S' \longrightarrow S$}). % [sic] chagement -> changement, letter missing

\item[c)] The \textit{$W$-equivalence ``local over $S$''} : $\forall\ s \in S$, $f/s : X/s \longrightarrow Y/s$ is a $W$-equivalence. Means that it is a $W$-equivalence, and remains so under any base change $S' \longrightarrow S$ which is a local isomorphism. (\footnote[19]{should be stable under \textit{smooth} base change $S' \longrightarrow S$}).

\item[c$'$)] The \textit{$W$-equivalence ``colocal over $S$''} : $\forall\ s \in S$, $s\backslash f : s\backslash X \longrightarrow s\backslash Y$ is a $W$-equivalence. Means that it is a $W$-equivalence, and remains so under any base change $S' \longrightarrow S$ which is a local co-isomorphism. (\footnote[20]{should be stable under \textit{proper} base change over $S$}).
\end{itemize}

I am going to summarize in a recapitulative statement the relations between these four notions of $W$-equivalence, thinking of the case $W = W_\infty$. Once this is done, I shall concern myself with isolating the minimal conditions on a general $W$ which ensure that it is true for $W$.

\textbf{Proposition 1.} \textit{Let} \begin{tikzcd}[baseline=(s.base), row sep=small, column sep=small]
f : X \arrow[rr] \arrow[dr] & & Y \arrow[dl] \\
& |[alias=s]| S &
\end{tikzcd} \textit{be in} Cat$/S$.

\begin{itemize}
\item[1)] \textit{Local $W$-equivalence over $S$, and colocal $W$-equivalence over $S$, imply coarse $W$-equivalence :}
\[
\begin{tikzcd}
\text{c} \arrow[dr, Rightarrow] & & \text{c}' \arrow[dl, Rightarrow] \\
& \text{a} &
\end{tikzcd}
\]

\printedpage{25}
% head: A. Grothendieck : Dérivateurs, Ch. VI, Édition des Universités de Montpellier II et Paris VII

\item[2)] \textit{If $X$, $Y$ are proper over $S$, or $X$, $Y$ smooth over $S$, or $X$, $Y$} Hot-\textit{fibered over $S$, then fiberwise $W$-equivalence implies that it is a} universal \textit{$W$-equivalence over $S$ ($W$-equivalence, and remains so under} any \textit{base change $S' \longrightarrow S$), hence in these three cases, we have}
\[
\begin{tikzcd}
& \text{b} \arrow[dl, Rightarrow] \arrow[dr, Rightarrow] & \\
\text{c} \arrow[dr, Rightarrow] & & \text{c}' \arrow[dl, Rightarrow] \\
& \text{a} &
\end{tikzcd}
\]
(\footnote[21]{It suffices $X$, $Y$ $W$-proper over $S$, or $W$-smooth over $S$, or $W$-fibered over $S$.}).

But we can be more precise about each of these three cases :

\item[(2 fib)] \textit{If $X$, $Y$ are} Hot-\textit{fibered over $S$, then coarse $W$-equivalence of $f$ implies that it is a fiberwise $W$-equivalence, hence the four notions of $W$-equivalence are equivalent for $f$.}
\[
\text{a} \Longleftrightarrow \text{c} \Longleftrightarrow \text{c}' \Longleftrightarrow \text{b}
\]
(\footnote[22]{a $\Longrightarrow$ b only assuming L5, see below p. 43.})

\item[(2 smooth)] \msp{page 42}\textit{If $X$, $Y$ are smooth over $S$, then fiberwise $W$-equivalence over $S$ is equivalent to colocal $W$-equivalence over $S$ and implies} universal \textit{$W$-equivalence over $S$, hence in this case} % [sic] lisse -> lisses, plural agreement with X, Y
\[
\begin{tikzcd}
\text{b} \arrow[r, Leftrightarrow] \arrow[d, Rightarrow] & \text{c}' \\
\text{c} \arrow[d, Rightarrow] & \\
\text{a} &
\end{tikzcd}
\]

\printedpage{26}
% head: A. Grothendieck : Dérivateurs, Ch. VI, Édition des Universités de Montpellier II et Paris VII
(\footnote[23]{it even suffices $X$, $Y$ $W$-smooth over $S$.}).

\item[(2 prop)] \textit{If $X$, $Y$ are proper over $S$, then fiberwise $W$-equivalence over $S$ is equivalent to local $W$-equivalence over $S$ :}
\[
\begin{tikzcd}
\text{b} \arrow[r, Leftrightarrow] \arrow[d, Rightarrow] & \text{c} \\
\text{c}' \arrow[d, Rightarrow] & \\
\text{a} &
\end{tikzcd}
\]
(\footnote[24]{It even suffices $X$, $Y$ $W$-proper over $S$.}).
\end{itemize}

\textsc{Proof} 1) is none other than L4 (\footnote[25]{No, there is confusion, L4 is weaker, cf. p. 44 the axiom L4 bis.}) and its dual L4$'$ for a basic localizer $W$, and we have seen that if we assume L1, L3, then L4 is equivalent to L4$'$, and in the notion of basic localizer we shall always assume that this is satisfied. Also, we assume that L3 is true even in the stronger form L3 quater (p. 1) : $W_\infty \subseteq W$, i.e. every Hot-equivalence is a $W$-equivalence.

(2 smooth) If $X$, $Y$ \texttt{[are]} smooth over $S$, then the $X_s \longrightarrow s\backslash X$, $Y_s \longrightarrow s\backslash Y$ are $W_\infty$-equivalences, hence $W$-equivalences (it would even suffice that $X$, $Y$ be $W$-smooth$/S$), hence the commutative diagram
\[
\begin{tikzcd}
X_s \arrow[r, "f_s"] \arrow[d, "i_s^X"'] & Y_s \arrow[d, "i_s^Y"] \\
s\backslash X \arrow[r, "s\backslash f"] & s\backslash Y
\end{tikzcd}
\]
shows us that $f_s$ is a $W$-equivalence if and only if $s\backslash f$ is, whence b $\Longleftrightarrow$ c$'$. To see that fiberwise $W$-equivalence is even a \textit{universal} $W$-equivalence over $S$, we are reduced to showing that it is a coarse $W$-equivalence (since smoothness \texttt{[is]} stable under base change over $S$), and we win since c$'$ $\Longrightarrow$ a by 1).

The case (2 proper) is dual.

(2 fib) If $X$, $Y$ are Hot-fibered over $S$, then
\printedpage{27}
% head: A. Grothendieck : Dérivateurs, Ch. VI, Édition des Universités de Montpellier II et Paris VII
\msp{page 43}
the inclusions
\[
X_s \longrightarrow s\backslash X, \qquad X_s \longrightarrow X/s,
\]
and likewise for $Y$, are $W_\infty$-equivalences, hence $W$-equivalences (it would even suffice that $X$, $Y$ be $W$-fibrations over $S$ (\footnote[26]{compare def. 3, p. 117, in (I).}), i.e. by definition that the preceding functors be $W$-equivalences, and that the same condition be satisfied after any base change over $S$), the argument in the smooth resp. proper case already shows that
\[
\begin{tikzcd}
\text{c} & \text{b} \arrow[l, Rightarrow] \arrow[r, Leftrightarrow] \arrow[d, Rightarrow] & \text{c}' \\
& \text{a,} &
\end{tikzcd}
\]
and that fiberwise $W$-equivalence (or local $W$-equivalence over $S$, resp. colocal over $S$) already implies \textit{universal} $W$-equivalence over $S$. (\footnote[27]{valid assuming only $X$, $Y$ $W$-fibered over $S$.})

The question a $\Longrightarrow$ b is on the other hand distinctly more delicate, and we have seen that this is by no means the case for every basic localizer - probably not for the one associated to the derivator Der($\mathbf{Z}$-Mod), in any case not for the one associated to Der($k$-Mod), when 2 is invertible in $k$. (It would even suffice that $\exists\ n \in \mathbf{N}$, $n \geq 2$, such that $n$ invertible in $k$ - the typical counterexample is $X = S^1 \xrightarrow{\ n\ } S^1 = Y = S$, the étale covering of order $n$ of the circle.)

The implication a $\Longrightarrow$ b is thus a very strong condition on $W$, I am going to denote it L5 (the old L5, cf. page 2, not deserving to survive, for it is a botched compromise).

\textbf{L5} \textit{Let} \begin{tikzcd}[baseline=(s.base), row sep=small, column sep=small]
X \arrow[rr, "f"] \arrow[dr] & & Y \arrow[dl] \\
& |[alias=s]| S &
\end{tikzcd} \textit{be in} Cat\textit{, with $X$, $Y$} Hot-\textit{fibered over $S$. If $f$ is a $W$-equivalence (``coarse'' relative to $S$), then it is a fiberwise $W$-equivalence} (\footnote[28]{In XII (p. 12) we see that if $W$ satisfies L1, L3 quater, L4 bis (cf. next page), L5, then \fbox{$W = W_\infty$}. So it is perfectly illusory, if we assume L5 to go on dragging a $W$ along.}).

We have seen that up to an $\varepsilon$ (``lemma'' p. 23), L5 is a consequence of

\textbf{L6} \textit{Let $f : X \longrightarrow Y$ be a $W$-equivalence, $Y' \longrightarrow Y$ a perfect morphism (i.e. smooth and proper), then $f' : X' = X \times_Y Y' \longrightarrow Y'$ is a $W$-equivalence.}

\msp{page 44}
Let us note the quite natural variant

\textbf{L6 bis} \textit{As in L6, but assuming only $Y' \longrightarrow Y$ a} Hot-\textit{fibration.}

\printedpage{28}
% head: A. Grothendieck : Dérivateurs, Ch. VI, Édition des Universités de Montpellier II et Paris VII

There is a chance that in practice, when L6 (or only L5) is satisfied, L6 bis is too. % [sic] verifié -> vérifié, missing accent

We are going to introduce an axiom

\textbf{L7} \ a) \textit{Let $(f_i : X_i \longrightarrow Y_i)_{i \in I}$ be a family of arrows in} Cat\textit{, let $f : \coprod X_i \longrightarrow \coprod Y_i$ be the corresponding arrow on the sum categories. For $f \in W$, it is necessary and sufficient that $f_i \in W$ for every $i \in I$.}

\hspace*{2em}b) \textit{Let $f : X \longrightarrow Y$ be in $W$. If $Y$ is} $0$\textit{-connected, $X$ too.}

This axiom is equivalent to the conjunction of La ($f \in W \Longrightarrow \pi_0 f$ bijective) and the validity of the ``Lemma'' page 23, which I put here again as a reminder :

\textbf{Lemma.} \textit{The preceding axiom L7 on the basic localizer $W$ is equivalent to the following condition : $f : X \longrightarrow Y$ is in $W$ if and only if $\pi_0(f)$ is bijective, and $f$ induces $W$-equivalences $f_i : X_i \longrightarrow Y_i$ on the connected components.}

I must come back once more to 1) of the proposition, it is not true that ``it is'' the condition L4 on $W$ (which I assume, I said), it is L4 in the particular case where $Y = S$. On reflection, I see it as a strengthened variant of L4, which does not seem likely to follow from L4 :

\textbf{L4 bis} \textit{Let} \begin{tikzcd}[baseline=(s.base), row sep=small, column sep=small]
X \arrow[rr, "f"] \arrow[dr] & & Y \arrow[dl] \\
& |[alias=s]| S &
\end{tikzcd} \textit{be in} Cat\textit{, let us suppose that $f$ is a local $W$-equivalence over $S$ (i.e. $f/s : X/s \longrightarrow Y/s$ a $W$-equivalence $\forall\ s \in S$), then $f$ is a $W$-equivalence} (\footnote[29]{equivalent to the dual axiom, in terms of the $s\backslash f$}).

\msp{page 45}
In what follows, by a ``basic localizer'', we shall mean a set
\[
W \subseteq \mathrm{Fl}\,\mathrm{Cat}
\]
satisfying (at least) the axioms % [sic] axioms -> axiomes, French spelling
\[
\text{L1, L3 quater, L4 bis, L7,}
\]
i.e.

\begin{itemize}
\item[a)] \textit{Saturation} of $W$ (conditions a) b) c) p. 1)

\item[b)] $W_\infty \subseteq W$ $[\subseteq W_0]$ (\footnote[30]{I am no longer sure whether it is really useful to impose $W_\infty \subseteq W$. Take instead the final object axiom, namely b$'$). If we want a) b$'$) c) + L5, then we shall have $W \subseteq W_\infty$. \textbf{NB} a) b) c) + L5 implies $W = W_\infty$, cf. annotation page 43.}) (\footnote[31]{cf. d)}).

\item[c)] If $f : X \longrightarrow Y$ over $S$ is locally over $S$ a $W$-equivalence, then $f \in W$ (L4 bis)
\end{itemize}


\endgroup
% END UNIT Der-ChVI-345/Der-ChVI-345_p026-030.tex
% BEGIN UNIT Der-ChVI-345/Der-ChVI-345_p031-035.tex
\begingroup\makeatletter%
\newcommand{\printedpage}[1]{\marginpar{\footnotesize p.~#1}}%
\newcommand{\msp}[1]{\marginpar{\footnotesize [#1]}}%
\newcommand{\uc}[1]{\underline{\mathrm{#1}}}%
\newcommand{\HOT}{\mathrm{HOT}}%
\newcommand{\Cat}{\mathrm{Cat}}%
\newcommand{\defeq}{\overset{\mathrm{def}}{=}}%
\makeatother\let\printedpage\derpp\let\msp\dermsp%
\let\mspin\DERmspin%

\DERmark{ENV0315}


\printedpage{29}
% head: A. Grothendieck : Dérivateurs, Ch. VI, Édition des Universités de Montpellier II et Paris VII
% foot: Version 9 novembre 2001

\begin{enumerate}[leftmargin=*]
\item[d)] For $f : X \longrightarrow Y$ to be in $W$, it is necessary and sufficient that $\pi_0(f)$ be bijective, and that the $f_i : X_i \longrightarrow Y_i$ induced on the connected components be in $W$.
\end{enumerate}

These axioms are satisfied by all the $W$ that I know, and that one would feel like regarding as localizers. On the other hand, L5, L6, L6 bis (p. 43, 44) are axioms of a much more special nature. I feel like saying that they distinguish localizers of a strictly \textit{``homotopical''} nature from those of a \textit{``homological''} (or ``cohomological'') nature. I should make sure that they are satisfied for the $W_i$ ($i$-homotopy equivalence) (\footnote[32]{This is \textit{false}, cf. annotation above.}). % [sic] qui je connais -> que je connais, relative pronoun object requires "que"

The fact that L4 bis, i.e. c) above, implies the dual axiom, follows from the fact that L4 bis implies this

\textbf{Proposition 2.} \textit{Let $f : X \longrightarrow Y$ be in} $\Cat$\textit{, then $f \in W$ if and only if $f^o : X^o \longrightarrow Y^o$ is in $W$.}

Seen by the argument of I, page 75 (cor. 2) - uses only L4, not even L1, L3.

\msp{page 46}

I am going to try to summarize what I think I have understood:

\begin{enumerate}[leftmargin=*]
\item[1)] $W$-equivalence ``by fibers'' is of interest (for constructing localized categories) only in the case where one knows that it already implies, either local $W$-equivalence, or colocal $W$-equivalence (or even universal $W$-equivalence) over $S$, hence when one assumes $X$, $Y$ to be at the very least both either $W$-smooth, or $W$-proper, or $W$-fibered over $S$.

\item[2)] One could call local $W$-equivalence over $S$ \textit{``left fine $W$-equivalence''} over $S$, and colocal $W$-equivalence over $S$ \textit{``right fine $W$-equivalence''} over $S$. They coincide in $\uc{Fibhot}\,S$, a fortiori in $\uc{Parf}\,S$ and in the four categories (cf. diagram p. 8) $\uc{Fib}_0S$, $\uc{Fibsc}_0S$, $\uc{Cofib}_0S$, $\uc{Cofibsc}_0S$ (\footnote[33]{One could call them the six \textit{inner} categories of the diagram.}). On the other hand, they do not coincide on the \textit{bordering} categories of the diagram, in particular on $\uc{Liss}\,S$, $\uc{Prop}\,S$, $\uc{Fib}\,S$, $\uc{Cofib}\,S$, $\uc{Fibsc}\,S$, $\uc{Cofibsc}\,S$. In this regard, it is the \textit{fine $W$-equivalence} \textbf{on the right} (in terms of the $s\backslash X$) that is of interest on the (alas) left side of the diagram, i.e. for $\uc{Liss}\,S$, $\uc{Fib}\,S$, $\uc{Fibsc}\,S$, whereas it is the \textit{fine $W$-equivalence} \textbf{on the left} (in terms of the $X/s$) that is relevant on the right side, for $\uc{Prop}\,S$, $\uc{Cofib}\,S$, $\uc{Cofibsc}\,S$. One will never have to look simultaneously at the two fine equivalences on one and the same category of the diagram, with the sole exception of $\Cat/S$ as a whole.
\end{enumerate}

\msp{page 47}

\begin{enumerate}[leftmargin=*]
\item[3)] Coarse $W$-equivalence, for its part, is ``good everywhere'' - but that depends for whom. For the 6 categories that are above $\uc{Fibhot}\,S$, i.e.
\[
\uc{Fibhot}\,S,\quad \uc{Parf}\,S,\quad \uc{Fib}_0S,\quad \uc{Fibsc}_0S,\quad \uc{Cofib}_0S,\quad \uc{Cofibsc}_0S,
\]
\printedpage{30}
% head: A. Grothendieck : Dérivateurs, Ch. VI, Édition des Universités de Montpellier II et Paris VII
% foot: Version 9 novembre 2001
it coincides with all the other equivalences (assuming axiom L5), for the three upper left bordering categories % [sic] moyennnant -> moyennant, triple n
\[
\uc{Liss}\,S,\quad \uc{Fib}\,S,\quad \uc{Fibsc}\,S
\]
of the diagram, fiberwise $W$-equivalence coincides with right fine $W$-equivalence, and dually on
\DERmark{ENV0316}

\[
\uc{Prop}\,S,\quad \uc{Cofib}\,S,\quad \uc{Cofibsc}\,S
\]
fiberwise $W$-equivalence coincides with left fine $W$-equivalence. As for $\Cat/S$ as a whole, it is the only category of the diagram where there are \textit{three} different interesting notions of $W$-equivalence, namely: the coarse one, and the two fine ones (left, right). On the other six bordering categories of the diagram, there are in principle two: the coarse one, and a specified one of the two fine ones (forbidden to look at the other - at one's own risk and peril \ldots): the right one on the left side, the left one on the right side. (NB I ought to redo the diagram, really, putting the fibrations on the right, not on the left.) As for the six inner categories of the diagram, in them there is only a single notion of $W$-equivalence left (since on them the four notions coincide).
\DERmark{ENV0317}

\end{enumerate}

\msp{page 48 (blank)}

\msp{page 49} (\textbf{NB} there is no page 48)

I now summarize what I think I have understood about the corresponding localized categories. There are, up to equivalence, only \textit{three} different ones all told, when one takes the ``admissible'' localizers in the categories, as I have just said (that being in short the objective interest of the prescriptions!) To be precise:

\begin{enumerate}[leftmargin=*]
\item[$1^o$)] The \textit{coarse $W$-equivalence} gives ``the same'' localized category for the 13 categories of the diagram, i.e. all the transition arrows of the diagram induce equivalences for this coarse $W$-equivalence. The localized category will be denoted \fbox{$\HOT^{\mathrm{lc}}_W(S)$}, where the exponent lc means: (relative $W$-homotopy types over $S$) \textit{``locally constant.''} (I have only a small doubt left for $\uc{Fibsc}\,S$ and $\uc{Cofibsc(S)}$ (\footnote[34]{This doubt is unfounded, everything works perfectly!}), whether coarse equivalence there is really and truly reasonable. We will see when doing the proof.)

\item[$2^o$)] The \textit{right fine $W$-equivalence} gives the same result for all the categories of the diagram where it is considered as admitted, \textit{and where it does not coincide with the coarse $W$-equivalence} already considered, hence on the four left bordering categories
\[
\Cat/S,\quad \uc{Liss}\,S,\quad \uc{Fib}\,S,\quad \uc{Fibsc}\,S
\]
(\footnote[35]{\textbf{NB} On the last three this equivalence is $W$-equivalence by fibers.}). The localized category is denoted \fbox{$\HOT_W(S)$}.
\end{enumerate}

\msp{page 50}

\printedpage{31}
% head: A. Grothendieck : Dérivateurs, Ch. VI, Édition des Universités de Montpellier II et Paris VII
% foot: Version 9 novembre 2001

\begin{enumerate}[leftmargin=*]
\item[$2'^o$)] Dually, the \textit{left fine $W$-equivalence} gives the same result on all the categories of the diagram where it is admitted, and where it does not coincide with the coarse one, i.e. on the 4 categories
\[
\Cat/S,\quad \uc{Prop}\,S,\quad \uc{Cofib}\,S,\quad \uc{Cofibsc}\,S.
\]
It is appropriate to denote the common category of fractions by the symbol
\[
\HOT_W(S^o) \qquad \text{(or also $\HOT^o_W(S)$)},
\]
\DERmark{ENV0318}

since it is isomorphic to the one denoted by this symbol.
\end{enumerate}

If one wants unambiguous notations, designating well-defined categories, and not just ``up to equivalence'', I will choose
\[
\HOT_W(S) \defeq \uc{Fibsc}(S)(W^{\mathrm{d}}_S)^{-1}
\]
whence
\[
\HOT_W(S^o) \qquad \underbrace{\simeq}_{\text{\footnotesize isomorphism of categories}} \qquad \uc{Cofibsc}(S)(W^{\mathrm{g}}_S)^{-1},
\]
where $W^{\mathrm{d}}_S$, $W^{\mathrm{g}}_S$ denote the sets of arrows corresponding to right resp. left fine $W$-equivalence. For $\HOT^{\mathrm{lc}}_W(S)$, I prefer a definition as symmetric as possible in $S$, $S^o$, and moreover a category as small as possible, which leads me to choose
\[
\HOT^{\mathrm{lc}}_W(S) \defeq (\uc{Parf}\,S)W^{-1}_S,
\]
where we have denoted by $W_S$ the $W$-equivalences in $\uc{Parf}\,S$. (No need this time to specify which ones - all four \texttt{[notions]} coincide in $\uc{Parf}\,S$!)

\msp{page 51}

Let us note the isomorphism
\[
\uc{Parf}(S^o) \simeq \uc{Parf}\,S
\]
(coming from the fact that $f : X \longrightarrow Y$ in $\Cat$ is a perfect fibration if and only if $f^o$ is). This isomorphism $X \longmapsto X^o$ is compatible with the localizers (since $f \in W \Longleftrightarrow f^o \in W$), whence
\[
\HOT^{\mathrm{lc}}_W(S^o) \simeq \HOT^{\mathrm{lc}}_W(S)
\]
(isomorphism of categories). We will have two canonical functors between the three categories of relative $W$-homotopy types over $S$
\[
\begin{tikzcd}[column sep=small]
(\uc{Fibsc}_0S)W_S^{-1} \arrow[rr, "\approx"] \arrow[dr, dotted, "i'"'] & & \HOT^{\mathrm{lc}}_W(S) \arrow[dl, "i"] \arrow[dr, "j"'] & & (\uc{Cofibsc}_0S)W_S^{-1} \arrow[ll, "\approx"'] \arrow[dl, dotted, "j'"] \\
& \HOT_W(S) & & \HOT^o_W(S) \defeq \HOT_W(S^o). &
\end{tikzcd}
\]
\printedpage{32}
% head: A. Grothendieck : Dérivateurs, Ch. VI, Édition des Universités de Montpellier II et Paris VII
% foot: Version 9 novembre 2001
\texttt{[rather} $\defeq \HOT_W(S^o)^o$\texttt{]} These arrows $i$, $j$ are suggested by the dotted arrows $i'$, $j'$, coming from the inclusions % [sic] suggerées -> suggérées, missing accent
\[
\uc{Fibsc}_0S \lhook\joinrel\longrightarrow \uc{Fibsc}\,S,\quad \uc{Cofibsc}_0S \lhook\joinrel\longrightarrow \uc{Cofibsc}\,S,
\]
which would define $i$, $j$ with the help of a choice of quasi-inverse functors for the horizontal arrows. But $i$, $j$ would be defined only up to (unique) isomorphism. It will be prettier to choose the quasi-inverses in a canonical way, with the help of the arrows
\[
\uc{Fibsc}_0S \xleftarrow{\Phi_0} \uc{Parf}\,S \xrightarrow{\Psi_0} \uc{Cofibsc}_0S
\]
(where $\Phi_0 = \Phi|\uc{Parf}\,S$, $\Psi_0 = \Psi|\uc{Parf}\,S$) of the diagram p. 8. The main question
\msp{page 52}
which remains unelucidated, is whether these functors $i$, $j$ really are \textit{fully faithful,} as I certainly expect.

I now want to prove the preceding assertions about the localized categories.

\textcircled{1} The localized category for coarse $W$-equivalence. It will suffice, up to an $\varepsilon$, to establish the conditions of the lemma p. 11, in the following three cases.
\begin{tikzcd}[baseline=(\tikzcdmatrixname-1-1.base)]
\uc{Parf}\,S \arrow[r, "r"] & \Cat/S \arrow[l, dotted, bend left, "C"]
\end{tikzcd}
(that will give the case of $\uc{Fibhot}\,S$, $\uc{Liss}\,S$ and $\uc{Prop}\,S$ over $\Cat/S$) and
\begin{tikzcd}[baseline=(\tikzcdmatrixname-1-1.base)]
\uc{Fibsc}_0S \arrow[r, "p"] & \Cat/S \arrow[l, dotted, bend left]
\end{tikzcd},
\begin{tikzcd}[baseline=(\tikzcdmatrixname-1-1.base)]
\uc{Cofibsc}_0S \arrow[r, "q"] & \Cat/S \arrow[l, dotted, bend left]
\end{tikzcd},
which will give respectively the cases of
\[
\uc{Fib}_0S,\ \uc{Fib}\,S,\ \uc{Liss}\,S \qquad \text{and} \qquad \uc{Cofib}_0S,\ \uc{Cofib}\,S,\ \uc{Prop}\,S \qquad \text{over } S.
\]
But the trouble for $p$ and $q$, is that the functors in the reverse direction that I had constructed (and again denoted $\Phi_0$, $\Psi_0$ (but they were essentially $\Phi\circ r\circ C$ and $\Psi\circ r\circ c$, not to be confused with the functors $\Phi_0$, $\Psi_0$ at the bottom of the preceding page, and also marked in the diagram p. 8), let us rather denote them $\Phi'$, $\Psi'$, work well on the $\Cat/S$ side, i.e. $p\Phi'$, $q\Psi'$ receive the functor $\mathrm{id}_{\Cat/S}$ by a $W$-arrow, but not on the side of the target categories - one cannot manage to find a natural arrow between $\Phi'p$ resp. $\Psi'q$, and the identity functors of $\uc{Fibsc}_0S$ resp. $\uc{Cofibsc}_0S$. I do
\msp{page 53}
not see, indeed, that there are any - the difficulty, by this route, seems irremediable. On the other hand, if one regards $\Phi'$, $\Psi'$ as functors with values in $\uc{Fib}_0S$, $\uc{Cofib}_0S$ respectively (thus in the categories where one does not require the arrows to be cartesian resp. cocartesian), one again finds $W$-morphisms % [sic] catésiennes -> cartésiennes, missing r
\[
\mathrm{id} \longrightarrow \Phi'p, \qquad \mathrm{id} \longrightarrow \Psi'q, % [sic] \psi'q -> \Psi'q, lowercase psi where \Psi' is used everywhere else
\]
\DERmark{ENV0319}

which at least already settles the case of all the categories of the diagrams p. 8 that are
\printedpage{33}
% head: A. Grothendieck : Dérivateurs, Ch. VI, Édition des Universités de Montpellier II et Paris VII
% foot: Version 9 novembre 2001
full subcategories of $\Cat/S$ (\footnote[36]{But it is idiotic to take such complicated functors, because one wants to reduce everything to $\Cat/S$. It is moreover simple to compare $\uc{Fibhot}\,S$ directly with $\uc{Fibsc}_0S$ and with $\uc{Cofibsc}_0S$. $\Phi_0$ and $\Psi_0$ (which are in fact defined on $\uc{Fibhot}$, not just on $\uc{Parf}$), and one finds as below $\mathrm{id} \xrightarrow{W} \pi_0\Phi_0$ and $\Phi_0\pi_0 \xrightarrow{W} \mathrm{id}$ (likewise for $\Psi_0$ and $\uc{Cofibsc}_0S$), where $\pi_0 : \uc{Fibsc}_0S \longrightarrow \uc{Fibhot}\,S$.}). It remains only to treat the case of the four categories of ``split'' structures
\[
\uc{Fibsc}\,S,\quad \uc{Fibsc}_0S \qquad \text{and} \qquad \uc{Cofibsc}\,S,\quad \uc{Cofibsc}_0S.
\]
Of course it suffices to do it on one side, say on the fibered side, by duality ($S \longmapsto S^o$) this will establish the symmetric result. So we are going to look a little more closely at the square of functors
\[
\begin{tikzcd}[column sep=huge, row sep=large]
\uc{Fibsc}\,S \arrow[d, "p"] & \uc{Fibsc}_0S \arrow[l, hook, "\alpha'"'] \arrow[d, "p_0"'] \\
\uc{Fib}\,S \arrow[u, dotted, bend left, "\varphi"] & \uc{Fib}_0S \arrow[l, hook, "\alpha"'] \arrow[u, dotted, bend right, "\varphi_0"']
\end{tikzcd}
\]
(\footnote[37]{NB It is a matter of proving only that $p$, $p_0$ induce equivalences for the localized categories (for coarse $W$-equivalence). We already know that this is so for $\alpha$, and it will follow for $\alpha'$.}, \footnote[38]{needless to look at this square, cf. marginal annotations above (top of the page)}). The functor
\[
\Phi : \Cat/S \longrightarrow \uc{Fibsc}\,S
\]
induces functors in the reverse direction to $p$, $p_0$ on the full subcategories $\uc{Fib}\,S$, $\uc{Fib}_0S$ of $\Cat/S$, functors marked by dotted lines as $\varphi$, $\varphi_0$.

\msp{page 54}

It is only necessary to check that if $X/S$ is in $\uc{Fib}_0S$, then $\varphi(X) = \Phi(X)$ is in the full subcategory $\uc{Fibsc}_0S$ of $\uc{Fibsc}\,S$. But we have seen that this is the case, provided only that $X/S$ is in $\uc{Fibhot}\,S$, for the base change morphisms in $\Phi(X)$, associated to the arrows $u : s \longrightarrow t$ in $S$
\[
u^*_{\Phi_X} : t\backslash X \longrightarrow s\backslash X
\]
are $W_\infty$-equivalences if $X$ is Hot-fibered over $S$. On the other hand, for arbitrary $X$ in $\Cat/S$, one has an arrow, functorial in $X$,
\[
X \longrightarrow \Phi(X),
\]
i.e. $\mathrm{id}_{\Cat/S} \longrightarrow \pi\Phi$ where $\pi : \uc{Fibsc}\,S \longrightarrow \Cat/S$ is the canonical functor, which is a $W_\infty$-equivalence (even a homotopism), hence a $W$-equivalence. A fortiori the functors $p\varphi$, $p_0\varphi_0$ give rise ($\uc{Fib}\,S$ and $\uc{Fib}_0S$ being full subcategories of $\Cat/S$) to arrows
\[
\mathrm{id} \longrightarrow p\varphi, \qquad \mathrm{id} \longrightarrow p_0\varphi_0.
\]


\endgroup
% END UNIT Der-ChVI-345/Der-ChVI-345_p031-035.tex
% BEGIN UNIT Der-ChVI-345/Der-ChVI-345_p036-040.tex
\begingroup\makeatletter%
\newcommand{\printedpage}[1]{\marginpar{\footnotesize p.~#1}}%
\newcommand{\msp}[1]{\marginpar{\footnotesize [#1]}}%
\newcommand{\ul}[1]{\underline{\mathrm{#1}}}%
\newcommand{\cercle}[1]{\tikz[baseline=(c.base)]\node[draw,circle,inner sep=1pt](c){#1};}%
\makeatother\let\printedpage\derpp\let\msp\dermsp%
\let\mspin\DERmspin%


\printedpage{34}
% head: A. Grothendieck : Dérivateurs, Ch. VI, Édition des Universités de Montpellier II et Paris VII

It remains to find an arrow (and coarse $W$-equivalence)
\[\varphi p \longrightarrow \mathrm{id},\]
\msp{page 55}
which will automatically induce
\[\varphi_0 p_0 \longrightarrow \mathrm{id}\]
since $\underline{\mathrm{Fibsc}}_0 S$ is \texttt{[a]} full subcategory of $\ul{Fibsc}\,S$. In fact, we are going to define $\varphi p \longrightarrow \mathrm{id}$, hence
\[\boxed{\gamma_X : \Phi(X) \longrightarrow X}\]
when $X$ is split fibered over $S$ - an arrow that is \textit{cartesian,} hence (if one likes) compatible with the splittings of $\Phi(X)$, and of $X$. In fact, we will have such an arrow, functorial in $X$, for $X$ varying in $\ul{Fib}\,S$ (thus without restricting ourselves to cartesian $S$-morphisms between fibered categories over $S$), but the arrow $\longrightarrow$ being of course \textit{cartesian.} (This makes sense, without choosing a splitting on $X$, and without there necessarily being one.) The definition of this arrow $\gamma$ makes essential use of the hypothesis that $X$ be Cat-fibrant over $S$. Of course, composed with the obvious inclusion arrow % [sic] fibrée scindée sur X -> fibrée scindée sur S, the base category is S
\DERmark{ENV0320}

\[X \xrightarrow{i_X} \Phi(X)\]
we are going to have
\[\gamma_X i_X = \mathrm{id}_X\]
(\footnote[39]{NB Since we know that $i_X \in W$, it will indeed follow that $\gamma_X \in W$ as we want.}). We define $\gamma$ fiber by fiber,
\[\gamma_{X,s} : \Phi(X)_s \overset{\mathrm{def}}{=} s\backslash X \longrightarrow X_s,\]
\DERmark{ENV0321}

thus by a \textit{retraction} of $s\backslash X$ onto the full subcategory $X_s$. An object of $s\backslash X$ is a pair $(x,u)$, $x \in \mathrm{Ob}\, X$, $u : s \longrightarrow p_X(x)$ ($p_X : X \longrightarrow S$ the structural morphism of $X$). We put
\[\gamma(x,u) = u_X^*(x) \in \mathrm{Ob}\, X_s,\]
\msp{page 56}
which thus supposes that we have chosen inverse image functors $u_X^* : X_t \longrightarrow X_s$, for the Cat-fibered category $X/S$. (But up to unique isomorphism, $\gamma$ does not depend on these choices.) I dispense with the tedious verification that this rule is indeed ``functorial in $s$'', i.e. does indeed define a cartesian arrow
\[\gamma_X : \Phi(X) \longrightarrow X,\]
and with verifying that the latter is functorial in $X$.

\printedpage{35}
% head: A. Grothendieck : Dérivateurs, Ch. VI, Édition des Universités de Montpellier II et Paris VII

Perhaps it is more elegant to show that $\pi : \ul{Fibsc}\,S \longrightarrow \mathrm{Cat}/S$ and $\Phi : \mathrm{Cat}/S \longrightarrow \ul{Fibsc}\,S$ are adjoint to one another:
\[\boxed{\mathrm{Homcart}_S(\Phi(\underbrace{\;X\;}_{\in\mathrm{Cat}/S}),\underbrace{\;Y\;}_{\in\ul{Fibsc}}) \simeq \mathrm{Hom}_S(X,\pi Y),}\]
from which there will follow adjunction arrows
\[\Phi\pi \longrightarrow \mathrm{id}_{\ul{Fibsc}}, \qquad\qquad \pi\Phi \longleftarrow \mathrm{id}_{\ul{Fib}},\]
\texttt{[rather} $\pi\Phi \longleftarrow \mathrm{id}_{\mathrm{Cat}/S}$\texttt{]} and we will have won. But to prove this adjunction formula, I have the impression that one has to use $\varphi p \longrightarrow \mathrm{id}$ \ldots\ There are tedious verifications \ldots

Finally, I now see that we have five cases to deal with, and three if we play on the symmetries:

\msp{page 57}

\[
\begin{tikzcd}[row sep=large]
\ul{Parf}\,S \arrow[d, "r"'] \\
\mathrm{Cat}/S \arrow[u, dotted, bend right, "C"']
\end{tikzcd}
\qquad\qquad\qquad
\left\{\begin{array}{lcl}
rC & \longleftarrow & \mathrm{id} \\
Cr & \longleftarrow & \mathrm{id}
\end{array}\right.
\]

\[
\begin{tikzcd}
\ul{Fibsc}\,S \arrow[dr, "\lambda"] & & \ul{Cofibsc}\,S \arrow[dl, "\mu"'] \\
& \mathrm{Cat}/S \arrow[ul, dotted, bend left, "\Phi"] \arrow[ur, dotted, bend right, "\Psi"'] &
\end{tikzcd}
\qquad
\begin{array}{l}
\left\{\begin{array}{lcl}
\lambda\Phi & \longleftarrow & \mathrm{id} \\
\Phi\lambda & \longrightarrow & \mathrm{id}
\end{array}\right. \\[2ex]
\left\{\begin{array}{lcl}
\mu\Psi & \longleftarrow & \mathrm{id} \\
\Psi\mu & \longrightarrow & \mathrm{id}
\end{array}\right.
\end{array}
\]

\[
\begin{tikzcd}
\underline{\mathrm{Fibsc}}_0 S \arrow[dr, "\lambda_0"] & & \underline{\mathrm{Cofibsc}}_0 S \arrow[dl, "\mu_0"'] \\
& \ul{Fibhot}\,S \arrow[ul, dotted, bend left, "\Phi_0"] \arrow[ur, dotted, bend right, "\Psi_0"'] &
\end{tikzcd}
\qquad
\begin{array}{l}
\left\{\begin{array}{lcl}
\lambda_0\Phi_0 & \longleftarrow & \mathrm{id} \\
\Phi_0\lambda_0 & \longrightarrow & \mathrm{id}
\end{array}\right. \\[2ex]
\left\{\begin{array}{lcl}
\mu_0\Psi_0 & \longleftarrow & \mathrm{id} \\
\Psi_0\mu_0 & \longrightarrow & \mathrm{id}
\end{array}\right.
\end{array}
\]

(\footnote[40]{\textbf{NB} All the marked functors are compatible with coarse $W$-equivalence, and of the ten functorial homomorphisms are coarse $W$-equivalences.})

NB The last four homomorphisms of functors, relating to the third diagram (with the subscripts 0) are simply induced by those relating to the diagram preceding it. Thus the key verifications relate to the first two diagrams, and by symmetry, in the end we only have to look at the \textit{two} pairs of functors $(r,C)$, $(\lambda,\Phi)$, and for each of these pairs, exhibit the two functorial homomorphisms, of the two composites with the identity functors.

\printedpage{36}
% head: A. Grothendieck : Dérivateurs, Ch. VI, Édition des Universités de Montpellier II et Paris VII

\msp{page 58}

\cercle{2} The localized categories for \textit{right fine $W$-equivalence.} There is no longer any reason, as we have said, to look at the inner categories of the diagram p. 8, where the said equivalence coincides with coarse $W$-equivalence, for which we already know the localized categories. It remains to look at the categories
\[\mathrm{Cat}/S,\ \ \ul{Liss}\,S,\ \ \ul{Fib}\,S,\ \ \ul{Fibsc}\,S.\]
The lemma on page 11 tells us that it suffices to show that in the situation
\[\ul{Fibsc}\,S \xrightarrow{\lambda} \mathrm{Cat}/S,\]
there is a functor in the opposite direction, with the specified properties. It will follow that any full subcategory $\mathcal{F}$ of $\mathrm{Cat}/S$ which contains $\ul{Fib}\,S$, when we equip it with the localizer of right fine $W$-equivalence, gives a localized category equivalent, by the functor induced by the inclusion, to the corresponding localized category of $\mathrm{Cat}/S$, i.e. essentially to $\mathrm{HOT}_W(S)$. % [sic] specifiées -> spécifiées, missing accent

As functor in the opposite direction, it is natural to take $\Phi$ again. There are three things to verify:

a) $\Phi$ \texttt{[is]} compatible with the localizers, i.e. if
\begin{tikzcd}[baseline=(\tikzcdmatrixname-2-2.base), row sep=small, column sep=small]
X \arrow[rr, "f"] \arrow[dr] & & Y \arrow[dl] \\
& S &
\end{tikzcd}
over $S$ is a right fine $W$-equivalence, i.e. the $s\backslash X \longrightarrow s\backslash Y$ are $W$-equivalences, then so is $\Phi(f)$. But the said equivalence coincides
\msp{page 59}
in $\ul{Fibsc}\,S$ (as also in $\ul{Fib}\,S$ and $\ul{Liss}\,S$) with $W$-equivalence on fibers, and on the other hand $\Phi(f)_s$ is none other than the preceding arrow $s\backslash X \longrightarrow s\backslash Y$. Thus it is obvious that
\[f \in \underbrace{W_S^{\mathrm{d}}}_{W\text{-equivalence on the right}} \iff \Phi(f) \in \underbrace{W_S^{\mathrm{f}}}_{W\text{-equivalence on fibers}} .\]

b) That $X \xrightarrow{i_X} \lambda\Phi(X)$ is in $W_S^{\mathrm{d}}$, for every $X$ in $\mathrm{Cat}/S$. This thus means that $i_X$ induces
\[s\backslash X \longrightarrow s\backslash\Phi(X)\]
which is a $W$-equivalence. This would be O.K. if it were clear that $X \longmapsto \Phi(X)$ commutes with every proper base change over $S$ (hence in particular with $s\backslash S \longrightarrow S$), since we know that $X \longrightarrow \Phi(X)$ is a $W$-equivalence. This is probably asking too much, but at least that it commutes with base change which is a local co-isomorphism (such as the $s\backslash S \longrightarrow S$).

For such a base change $S' \xrightarrow{p} S$, verifying that the fiber of $\Phi(X')$ at $s'$ is isomorphic to that of $\Phi(X)$ at $s = p(s')$, I will be content with that. Now this fiber is $s'\backslash X'$, obtained from $X'$ over $S'$ by the base change $s'\backslash S' \longrightarrow S'$, hence obtained from $X$ by $s'\backslash S' \longrightarrow S$. Now $s'\backslash S' \longrightarrow s\backslash S$ is an isomorphism by the hypothesis
\printedpage{37}
% head: A. Grothendieck : Dérivateurs, Ch. VI, Édition des Universités de Montpellier II et Paris VII
\msp{page 60}
on $p$ (local co-isomorphism), hence the base change $s'\backslash S' \longrightarrow S$ is identified with $s\backslash S \longrightarrow S$, and the inverse image is indeed isomorphic to $s\backslash X$, OK.

c) That likewise
\[\Phi\lambda \longrightarrow \mathrm{id}\]
is not only in $W_S$ (coarse), but that it is a $W$-equivalence on fibers. So we start with $X$ fibered over $S$. We look at
\[\Phi(X) \longrightarrow X,\]
to see that it is a $W$-equivalence on fibers, it suffices to see that this is so for $X \longrightarrow \Phi(X)$,
\[X_s \longrightarrow s\backslash X\]
is a $W$-equivalence for every $s \in S$. But this is already true as soon as $X$ is smooth over $S$.

Since the case of left fine $W$-equivalence reduces to the present case by $S \longmapsto S^{o}$, we have thus proved all that was expected, with the sole exception of the full faithfulness of the functors $i$, $j$ on page 51. % [sic] ce qu'était -> ce qui était, grammar; fidèlité -> fidélité, accent

But in the prospective summary of the situation for the localizers, p. 49-51, I forgot other functors besides $i$, $j$, of a nature if possible even more trivial. They are
\msp{page 61}
those which one can describe starting from the identity functor, on one of the categories of the diagram where two, or even three, different notions of $W$-equivalence are in use. These are, I recall

\begin{enumerate}[label={}]
\item[$1^{o}$)] The four categories
\[(*) \qquad\qquad \ul{Fibsc}\,S \longrightarrow \ul{Fib}\,S \longrightarrow \ul{Liss}\,S \longrightarrow \mathrm{Cat}/S, \qquad\qquad W_S^{\mathrm{d}} \subseteq W\]
with right fine $W$-equivalence on the one hand, coarse on the other, and dually
\[(*') \qquad\qquad \ul{Cofibsc}\,S \longrightarrow \ul{Cofib}\,S \longrightarrow \ul{Prop}\,S \longrightarrow \mathrm{Cat}/S, \qquad\qquad W_S^{\mathrm{g}} \subseteq W\]
for left fine $W$-equivalence on the one hand, coarse on the other, and moreover
\item[$2^{o}$)] $\mathrm{Cat}/S$, which was included in the two preceding sequences, and which is the only category of the diagram giving rise to three different notions of $W$-equivalence % [sic] inclue -> incluse, past participle of inclure
\end{enumerate}

\[
\begin{tikzcd}[row sep=small, column sep=large]
W_S^{\mathrm{d}} \arrow[dr, hook] & \\
& W_S \\
W_S^{\mathrm{g}} \arrow[ur, hook] &
\end{tikzcd}
\qquad\qquad .
\]

\printedpage{38}
% head: A. Grothendieck : Dérivateurs, Ch. VI, Édition des Universités de Montpellier II et Paris VII

In each of the four cases ($C$, equipped with two localizers $\Sigma \subseteq \Sigma'$) of the sequence $(*)$, the functor
\[C\Sigma^{-1} \longrightarrow C\Sigma'^{-1}\]
is essentially ``the same'', because of the comparison theorems. It suffices to look at this functor in a single case, e.g. one of the two extremes $\ul{Fibsc}\,S$ or $\mathrm{Cat}/S$. Same remark for the sequence $(*')$. So finally there remain two remarkable functors % [sic] das -> cas, misprint
\msp{page 62}
to examine, which do not at all look as though they come from the functors of the diagram p. 8 or from some composite (that only gives canonical equivalences or else the functors $i$, $j$ on page 51, when we go from the inner categories of the diagram to the bordering categories). The prettiest thing is probably to look at them starting from $\mathrm{Cat}/S$, since it is only there that we have both at once
\[
\begin{tikzcd}[row sep=small, column sep=tiny]
(\mathrm{Cat}/S)(W_S^{\mathrm{d}})^{-1} \approx \mathrm{HOT}_W(S) \arrow[dr, "\rho"] & \\
& (\mathrm{Cat}/S)W_S^{-1} \approx \mathrm{HOT}_W^{\mathrm{lc}}(S). \\
(\mathrm{Cat}/S)(W_S^{\mathrm{g}})^{-1} \approx \mathrm{HOT}_W^{o}(S) \arrow[ur, "\sigma"'] &
\end{tikzcd}
\]
On the other hand, the functors (which we feel like calling ``inclusion functors'') $i$, $j$ on page 51, in the opposite direction to $\rho$, $\sigma$, are those which come, by localization at $W_S^{\mathrm{d}}$ or $W_S^{\mathrm{g}}$ as the case may be, from one of the functors of the diagram page 8 (or from a composite of such functors), going from one of the six inner categories (where all the $W$-equivalences considered coincide) to one of the seven bordering categories, except when the latter is $\mathrm{Cat}/S$, is already determined with which $W$-equivalence, $W_S^{\mathrm{d}}$ or $W_S^{\mathrm{g}}$, we are going to work \texttt{[structure of this sentence?]}. The determi-\msp{page 63}% hyphen kept as printed: the word is divided at the manuscript page marker
nation that is perhaps the prettiest is by
\[(*) \qquad\qquad \ul{Parf}\,S \longrightarrow \mathrm{Cat}/S\]
which gives the functor $i$ when we work with $W_S^{\mathrm{d}}$, the functor $j$ when we work with $W_S^{\mathrm{g}}$. We then see at once that if we compose with the relevant functor, either $\rho$ or $\sigma$, defined by the inclusion $W_S^{\mathrm{d}} \subseteq W_S$ resp. $W_S^{\mathrm{g}} \subseteq W_S$ (as subsets of $\mathrm{Fl}(\mathrm{Cat}/S)$), we find the functor
\[(**) \qquad\qquad (\ul{Parf}\,S)W_S^{-1} \xrightarrow{\approx} (\mathrm{Cat}/S)W_S^{-1}\]
which is the \textit{canonical equivalence} between these two categories. Thus the functor corresponding to $i$, $j$ page 51, expressed in terms of $\mathrm{Cat}/S$ as a functor in the opposite direction to $\rho$, $\sigma$ (p. 62), is the composite of the equivalence quasi-inverse to $(**)$, which is made explicit by the functor $C$, and of the functor on the category localized at $W_S^{\mathrm{d}}$ or $W_S^{\mathrm{g}}$ obtained from $(*)$, i.e.


\endgroup
% END UNIT Der-ChVI-345/Der-ChVI-345_p036-040.tex
% BEGIN UNIT Der-ChVI-345/Der-ChVI-345_p041-045.tex
\begingroup\makeatletter%
\newcommand{\printedpage}[1]{\marginpar{\footnotesize p.~#1}}%
\newcommand{\msp}[1]{\marginpar{\footnotesize [#1]}}%
\newcommand{\HOT}{\mathrm{HOT}}%
\newcommand{\Cat}{\mathrm{Cat}}%
\newcommand{\Ob}{\mathrm{Ob}}%
\newcommand{\Fl}{\mathrm{Fl}}%
\newcommand{\id}{\mathrm{id}}%
\newcommand{\cA}{\mathcal{A}}%
\newcommand{\ul}[1]{\underline{\mathrm{#1}}}%
\newcommand{\circled}[1]{\tikz[baseline=(c.base)]\node[draw,circle,inner sep=1pt](c){#1};}%
\makeatother\let\printedpage\derpp\let\msp\dermsp%
\let\mspin\DERmspin%


\printedpage{39}
% head: A. Grothendieck : Dérivateurs, Ch. VI, Édition des Universités de Montpellier II et Paris VII

\[
\begin{tikzcd}[column sep=large, row sep=large]
 & \HOT_W(S)\approx(\Cat/S)(W_S^{\mathrm{d}})^{-1} \\
\HOT^{\mathrm{lc}}_W(S)\approx(\Cat/S)W_S^{-1} \arrow[ur, "i"] \arrow[r, "C"] \arrow[dr, "j"'] & (\ul{Parf}\,S)W_S^{-1} \arrow[u, "\text{\footnotesize since $W_S=W_S^{\mathrm{d}}$ on $\ul{Parf}$}"'] \arrow[d, "\text{\footnotesize since $W_S=W_S^{\mathrm{g}}$ on $\ul{Parf}$}"] \\
 & \HOT^{o}_W(S)\approx(\Cat/S)(W_S^{\mathrm{g}})^{-1}.
\end{tikzcd}
\]

\msp{page 64}
We see from this that we have canonical isomorphisms
\[
\rho i\simeq\id,\qquad\sigma j\simeq\id
\]
in the category of endofunctors of $(\Cat/S)W_S^{-1}$ ($\approx\HOT^{\mathrm{lc}}_W(S)$), an isomorphism which is made explicit here by the fact that $\rho i$, $\sigma j$ are deduced from the same functor $C$ by passage to the localized categories $W_S^{-1}$ (so in fact $\rho i=\sigma j$ !), and that we have an arrow functorial in $X$
\[
X\xrightarrow{i_X}C(X),
\]
which is in $W$ for every $X$. Thus, $\rho$ and $\sigma$ are, up to isomorphism, \textit{retractions} of the categories $\HOT_W(S)$, $\HOT^{o}_W(S)$ onto the category $\HOT^{\mathrm{lc}}_W(S)$. Moreover, these retractions are what are called localization functors (as we see from the diagram on page 62). (We say that $C\xrightarrow{F}C'$ is a \textit{localization functor} if, denoting by $\Sigma\subseteq\Fl\,C$ the set of arrows of $C$ transformed into isomorphisms by $F$, the functor
\[
C\Sigma^{-1}\xrightarrow{\bar{F}}C'
\]
which factors $F$ is an equivalence of categories.)

\msp{page 65}
We know, moreover, that if a functor $f$ has an adjoint $g$ (left or right), then $f$ is a functor \texttt{[of]} localization if and only if the adjoint $g$ is fully faithful. To prove that $i$, $j$ are fully faithful, it would therefore be tempting to prove that they are adjoint to the functors $\rho$, $\sigma$.

To get our bearings, I shall set
\[
\left\{
\begin{aligned}
&\cA=\Cat/S,\ \ \Sigma_0=W_S^{\mathrm{d}},\ \ \Sigma=W_S\qquad\quad(\subseteq\Fl(\cA))\\
&C:\cA\longrightarrow\cA\ \text{ the familiar functor, which has the virtue}\\
&\qquad\qquad\qquad\quad C(\Sigma)\subseteq\Sigma_0\quad.
\end{aligned}
\right.
\]

We thus have the situation
\printedpage{40}
% head: A. Grothendieck : Dérivateurs, Ch. VI, Édition des Universités de Montpellier II et Paris VII
\[
(*)\qquad\qquad(\cA,\Sigma_0)\ \underset{C}{\overset{\rho_0\,=\,\id_{\cA}}{\rightleftarrows}}\ (\cA,\Sigma),
\]
inducing by passage to the localized categories
\[
\cA\Sigma_0^{-1}\quad\underset{\bar{C}}{\overset{\rho\,=\,\bar{\rho}_0}{\rightleftarrows}}\ \cA\Sigma^{-1},
\]
we want to check that $\rho$ and $C$ \texttt{[rather $\bar{C}$]} are adjoint to one another, $C$ \texttt{[rather $\bar{C}$]} \textit{surely on the right,} which means that we must have functorial homomorphisms
\[
\left\{
\begin{aligned}
&\id_{\cA\Sigma^{-1}}&&\rightleftarrows&&\rho\,\bar{C}\quad(=\overline{\rho_0C})\\
&\id_{\cA\Sigma_0^{-1}}&&\rightleftarrows&&\bar{C}\rho\quad(=\overline{C\rho_0})
\end{aligned}
\right.
\]
(\footnote[41]{but if $C$ is on the \textit{right,} the arrows here must go in the opposite direction and this is indeed what I must resign myself to on p.~67.}) \texttt{[the arrows at the bottom go in the corrected direction, but the arrows at the top are not erased]} satisfying two compatibility relations to which we shall return. We may hope to find these functorial homomorphisms starting from
\[
\left\{
\begin{aligned}
&\id_{\cA}&&\xrightarrow{\ i\ }&&\rho_0C&&=&&C\\
&\id_{\cA}&&\xleftarrow{\ p\ }&&C\rho_0&&=&&C
\end{aligned}
\right.
\]

\msp{page 66}
and in our case we do indeed have
\[
i:\id_{\cA}\longrightarrow C\qquad\text{i.e. }\ i_X:X\longrightarrow C(X)
\]
functorial in $X\in\Ob\,\cA$ (and moreover $i_X$ is in $\Sigma$). There is however no natural arrow in the opposite direction $p_X$
\[
C(X)\xrightarrow{\ ?\ }X.
\]
I had thought of the projection $\pi_X:C(X)\longrightarrow X$, deduced from the construction $C(X)\simeq X\times_S(\ul{Ch}_\infty(S),\sigma)$, but it is unfortunately not an $S$-morphism. As we are looking for an arrow in $\cA\Sigma_0^{-1}$, we could make do with having a functorial arrow in the opposite direction, provided it is in $\Sigma_0$, hence can be inverted when we work in $\cA\Sigma_0^{-1}$. The obvious candidate is $i_X:X\longrightarrow C(X)$ which we have just spoken of. The trouble is that $i_X$ is only in $\Sigma$, not in $\Sigma_0$ (\footnote[42]{If this were not so, the relative homotopy type over $S$ defined by $X/S$ would be locally constant (since this is so for the one defined by $C(X)$), which is not always the case.}). Yet, when $X$ is itself of the form $C(Y)$ (more generally, if it is Hot-fibered over $S$), then the fact that $X\longrightarrow C(X)$ is in $\Sigma$ already implies that it is in $\Sigma_0$. Let us therefore note this % [sic] evident -> évident, accent missing
\[
\left\{
\begin{array}{lrl}
i:\id_{\cA}\longrightarrow C\ \text{such that} & 1^o) & i_X:X\longrightarrow C(X)\ \textit{in}\ \Sigma\\[1ex]
 & & \text{for every }X\in\Ob\,\cA\\[1ex]
 & 2^o) & i_X\ \text{in }\Sigma_0\text{ if }X\text{ of the}\\[1ex]
 & & \text{form }C(Y)
\end{array}
\right.
\]

\printedpage{41}
% head: A. Grothendieck : Dérivateurs, Ch. VI, Édition des Universités de Montpellier II et Paris VII

\msp{page 67}
The idea comes to me, after some shitty attempts \texttt{[sic]}, that $\bar{C}$ might be the right adjoint, and not the left \texttt{[this is the corrected version (right and left interchanged) cf.~p.~65]}, i.e. that the arrows sought might go in the opposite direction
\[
\left\{
\begin{aligned}
&\id_{\cA\Sigma^{-1}}&&\xleftarrow{\ p\ }&&\rho\,\bar{C}\quad(=\overline{\rho_0C})\\
&\id_{\cA\Sigma_0^{-1}}&&\xrightarrow{\ i\ }&&\bar{C}\rho\quad(=\overline{C\rho_0})
\end{aligned}
\right.
\]
(so I reverse the notations for $p$ and $i$), thus we are led if possible to look for arrows
\[
\begin{aligned}
&\id_{\cA}&&\xleftarrow{\ p\ }&&\rho_0C&&=&&C\\
&\id_{\cA}&&\xrightarrow{\ i\ }&&C\rho_0&&=&&C,
\end{aligned}
\]
thus arrows functorial in $X$
\[
X\xrightarrow{i_X}C(X)\xrightarrow{p_X}X.
\]
We do indeed have a functorial $i_X$ - always the same one! For $p_X$, we have no more than before, but recalling that it suffices to find it in $\cA\Sigma^{-1}$ (and no longer $\cA\Sigma_0^{-1}$!), it suffices to invert the arrow $i_X$, which is indeed an \textit{isomorphism} in $\cA\Sigma^{-1}$ (but not in $\cA\Sigma_0^{-1}$). We shall therefore set
\[
p_X=i_X^{-1}:C(X)\longrightarrow X\qquad\quad\text{in }\cA\Sigma^{-1}.
\]
We have to prove just two compatibilities, for the functors $\bar{C}$ and $\rho$ between $\cA\Sigma_0^{-1}$ and $\cA\Sigma^{-1}$ (and I shall write $C$ instead of $\bar{C}$): the two composites
\[
\left\{
\begin{array}{l}
\begin{tikzcd}[ampersand replacement=\&]
\rho \arrow[r, "\rho*i"] \arrow[rr, bend left=20, dotted, "\id"] \& \rho\,C\rho \arrow[r, "p*\rho"] \& \rho
\end{tikzcd}
\qquad\quad\text{\footnotesize (relation for functors with values in $\cA\Sigma^{-1}$)}\\[2ex]
\begin{tikzcd}[ampersand replacement=\&]
C \arrow[r, "i*C"] \arrow[rr, bend right=20, dotted, "\id"'] \& C\rho\,C \arrow[r, "C*p"] \& C
\end{tikzcd}
\qquad\quad\text{\footnotesize (relation for functors with values in $\cA\Sigma_0^{-1}$)}
\end{array}
\right.
\]
must be the identity $\id_\rho$ resp. $\id_C$.

\msp{page 68}
The first relation is checked tautologically, it amounts to saying that in $\cA\Sigma^{-1}$, we have an identity composite of the
\[
X\xrightarrow{i_X}C(X)\xrightarrow{i_X^{-1}}X\qquad\quad!
\]
The second means that for every $X$, we have an identity composite in $\cA\Sigma_0^{-1}$
\[
C(X)\xrightarrow{i_{C(X)}}C(C(X))\xrightarrow{C(p_X)}C(X).
\]
Now $p_X:C(X)\longrightarrow X$ had been defined in $\mathcal{A}\Sigma^{-1}$ as the inverse $i_X^{-1}$ of $i_X$, on the other hand the functor $C$ on $\cA$ defines by hypothesis $\cA\Sigma^{-1}\longrightarrow\cA\Sigma_0^{-1}$, more precisely $C(\Sigma)\subseteq\Sigma_0$ % [sic] $A\Sigma^{-1}$ (italic A) -> $\mathcal{A}\Sigma^{-1}$, elsewhere script A
\DERmark{ENV0322}

\printedpage{42}
% head: A. Grothendieck : Dérivateurs, Ch. VI, Édition des Universités de Montpellier II et Paris VII
thus $C(i_X)$ is invertible in $\cA\Sigma_0^{-1}$, and $C(p_X)=C(i_X)^{-1}$. So the compatibility amounts to the relation
\[
C(i_X)^{-1}i_{C(X)}=\id\qquad\quad\textit{in}\ \ \cA\Sigma_0^{-1},
\]
i.e.
\[
\boxed{\,i_{C(X)}=C(i_X)\,}:C(X)\longrightarrow C^2(X).
\]
In other words, we must prove the: % [sic] d'autre termes -> d'autres termes, plural

\textbf{Lemma.} \textit{Let $X$ be in $\Cat/S$, let us consider the two arrows}
\[
i_{C(X)},\ C(i_X):C(X)\rightrightarrows C(C(X)).
\]
\textit{These arrows are equal, if not in the absolute sense, at least in $(\Cat/S)(W_S^{\mathrm{d}})^{-1}$.}

Here $f_X:X\longrightarrow S$ is the standard projection and
\[
\begin{aligned}
C(X)\ &=\ \quad\{(x,c)\ \mid\ x\in\Ob\,X,\ c:f_X(x)\ \underbrace{\longrightarrow}_{\text{path in }S}\ t\}\\
C^2(X)\ &=\ \{(x,c,c')\ \mid\ x\in\Ob\,X,\ \underbrace{f_X(x)\xrightarrow{c}t\xrightarrow{c'}t'}_{c,\ c'\text{ paths in }S}\},
\end{aligned}
\]
the projections to $S$ are given by
\[
\begin{aligned}
(x,c)&&\longmapsto&&t&&=&&\text{target }c\\
(x,c,c')&&\longmapsto&&t'&&=&&\text{target }c'
\end{aligned}
\]

\msp{page 69}
We find
\[
i_{C(X)}(x,c)=(x,c,1_t):
\begin{tikzcd}
x \arrow[d] & & \\
f_X(x) \arrow[r, "c"] & t \arrow[r, "1_t", "\text{empty path}"'] & t
\end{tikzcd}
\]
\[
C(i_X)(x,c)=(x,1_{f_X(x)},c):
\begin{tikzcd}
x \arrow[d] & & \\
f_X(x) \arrow[r, "1_{f_X(x)}", "\text{empty path}"'] & f_X(x) \arrow[r, "c"] & t
\end{tikzcd}
\]
But we have a choice between two theories of paths, to construct $C(X)$: either with the $\ul{Ch}_\infty(S)$, or with the $\ul{Ch}(S)$. ``It gives the same results'', but depending on the situation one is more convenient than the other. I shall work with the $\ul{Ch}(S)$, in which case the $1_S$ are left and right units, and we have a composition law for paths. This law allows us to define
\[
\begin{aligned}
C^2(X)&&\xrightarrow{\ \gamma\ }&&C(X)\\
(x,c,c')&&\longmapsto&&(x,c'\circ c),
\end{aligned}
\]
and the composites of $i_{C(X)}$ and of $C(i_X)$ with this arrow are $\id_{C(X)}$ in both cases % [sic] $C_{i(X)}$ -> $C(i_X)$, the arrow considered is C(i_X)
\DERmark{ENV0323}

\printedpage{43}
% head: A. Grothendieck : Dérivateurs, Ch. VI, Édition des Universités de Montpellier II et Paris VII
\[
\begin{tikzcd}
C(X) \arrow[r, shift left] \arrow[r, shift right] \arrow[rr, bend right=15, "\id"'] & C^2(X) \arrow[r, "\gamma"] & C(X).
\end{tikzcd}
\]

So to show that $i_{C(X)}=C(i_X)$ in $(\Cat/S)(W_S^{\mathrm{d}})^{-1}$, it suffices to prove that $\gamma$ is in $W_S^{\mathrm{d}}$. But we know that the functor $C$ transforms $W_S$ into $W_S^{\mathrm{d}}$, so $C(i_X)$ is in $W_S^{\mathrm{d}}$, and since $\gamma\circ C(i_X)=\id$, $\gamma$ is in $W_S^{\mathrm{d}}$, OK! % [sic] foncteur C (upright roman C) -> foncteur $C$, math italic elsewhere
\DERmark{ENV0324}


\msp{page 70}
I feel like grouping what has been proved into a summary statement.

\textbf{Theorem-Scholie.} We consider the diagram of categories on page 8 (\footnote[43]{We assume for $W$ the conditions a) b) c) of p.~45, in addition L5 (p.~43). Cf. comments p.~81.}).

\circled{I} When we equip all 13 categories of the diagram with the localizer $W_S$ (coarse $W$-equivalence, i.e. $W$-equivalence in Cat, nothing more), \textit{all} the functors of the diagram induce equivalences on the localized categories (\footnote[44]{Moreover, if $\mathcal{F}$ is a full subcategory of $\Cat/S$ which contains one of the two subcategories $\ul{Fib}_0S$, $\ul{Cofib}_0S$ (which includes the case of the 8 categories of the diagrams, other than $\Cat/S$ itself), the inclusion $(\mathcal{F},W_S|\mathcal{F})$ into $(\Cat/S,W_S)$ induces an equivalence for the localized categories.}). Let us denote by
\[
\HOT^{\mathrm{lc}}_W(S)
\]
``the'' common localized category.

\circled{II} Let us equip the four left bordering categories
\[
(*)\qquad\qquad
\begin{tikzcd}
\Cat/S & \ul{Liss}\,S \arrow[l, hook'] & \ul{Fib}\,S \arrow[l, hook'] & \ul{Fibsc}\,S \arrow[l, hook']
\end{tikzcd}
\]
with the localizer $W_S^{\mathrm{d}}$ of \textit{right} fine $W$-equivalence over $S$ ( $\begin{tikzcd}[row sep=small, column sep=small]
X \arrow[rr, "f"] \arrow[dr] & & Y \arrow[dl]\\
 & S &
\end{tikzcd}$ is in $W_S^{\mathrm{d}}$ if and only if the $s\backslash X\longrightarrow s\backslash Y$ are \texttt{[rather, iff they are in $W$]}, for every $s\in\Ob\,S$). The functors of the diagram, i.e. those of $(*)$, between these categories, hence also their composites, induce equivalences between the localized categories (\footnote[45]{More generally, this works for every full subcategory of $\Cat/S$ which contains $\ul{Fib}\,S$ (and dually, with $W_S^{\mathrm{g}}$ instead of $W_S^{\mathrm{d}}$, if $\mathcal{F}$ contains $\ul{Cofib}\,S$).}). Let us denote by
\[
\HOT_W(S)
\]
``the'' common localized category.

Dual statement for
\[
(*')\qquad\qquad
\begin{tikzcd}
\Cat/S & \ul{Prop}\,S \arrow[l, hook'] & \ul{Cofib}\,S \arrow[l, hook'] & \ul{Cofibsc}\,S, \arrow[l, hook']
\end{tikzcd}
\]

% foot (each page): Version 9 novembre 2001


\endgroup
% END UNIT Der-ChVI-345/Der-ChVI-345_p041-045.tex
% BEGIN UNIT Der-ChVI-345/Der-ChVI-345_p046-050.tex
\begingroup\makeatletter%
\newcommand{\printedpage}[1]{\marginpar{\footnotesize p.~#1}}%
\newcommand{\msp}[1]{\marginpar{\footnotesize [#1]}}%
\newcommand{\HOT}{\mathrm{HOT}}%
\newcommand{\Cat}{\mathrm{Cat}}%
\newcommand{\Hot}{\mathrm{Hot}}%
\newcommand{\hot}{\mathrm{hot}}%
\newcommand{\Ob}{\operatorname{Ob}}%
\newcommand{\Hom}{\operatorname{Hom}}%
\newcommand{\ul}[1]{\underline{\mathrm{#1}}}%
\newcommand{\cercle}[1]{\tikz[baseline=(c.base)]\node[draw,circle,inner sep=1pt](c){#1};}%
\newcommand{\ovale}[1]{\tikz[baseline=(c.base)]\node[draw,rounded corners=0.8ex,inner sep=2pt](c){#1};}%
\makeatother\let\printedpage\derpp\let\msp\dermsp%
\let\mspin\DERmspin%


\setcounter{footnote}{45}

\printedpage{44}
% head: A. Grothendieck : Dérivateurs, Ch. VI, Édition des Universités de Montpellier II et Paris VII
% foot: Version 9 novembre 2001
and the localizers $W_S^{\mathrm{g}}$ in these categories ($f\in W_S^{\mathrm{g}}$ if and only if the $f/s : X/s\longrightarrow Y/s$ are in $W$). We find a ``same localized category'', denoted
\[ \HOT^o_W(S), \]
and we have tautologically
\msp{page 71}
\[ \HOT^o_W(S)\approx\HOT_W(S^o). \]
\texttt{[rather} $(\HOT_W(S^o))^o$ \texttt{cf. p. 50]}

\ovale{III}\quad On the six inner categories of the diagram, the three notions of $W$-equivalence, in terms of $W_S$, $W_S^{\mathrm{d}}$, $W_S^{\mathrm{g}}$ coincide, and coincide with fiberwise $W$-equivalence $W_S^{\mathrm{f}}$. On the three categories of $(*)$ other than $\Cat/S$, the localizers $W_S^{\mathrm{f}}$, $W_S^{\mathrm{d}}$ coincide; on the three categories of $(*')$ other than $\Cat/S$, the localizers $W_S^{\mathrm{f}}$, $W_S^{\mathrm{g}}$ coincide. But in general (\footnote{probably in all cases where $S$ is not equivalent to a discrete category.% [sic] dicrète -> discrète, faute d'orthographe
}), on these six categories, these localizers $W_S^{\mathrm{f}}$ are not equal to $W_S$, more precisely give nonequivalent localized categories, they are therefore each equipped with \textit{two} essentially distinct localizers, $W_S^{\mathrm{f}}$ and $W_S$. The category $\Cat/S$ is the only category in the diagram on which the localizers $W_S$, $W_S^{\mathrm{d}}$ and $W_S^{\mathrm{g}}$ are all three essentially distinct. % [sic] essentiellment -> essentiellement, lettre omise
It is the only one that allows one to obtain the three localized categories simultaneously, as
\[
\begin{tikzcd}[column sep=small]
 & \HOT^{\mathrm{lc}}_W(S)\simeq(\Cat/S)W_S^{-1} & \\
\HOT_W(S)\simeq(\Cat/S)(W_S^{\mathrm{d}})^{-1} \arrow[ur,"\rho"] & & \HOT^o_W(S)\simeq(\Cat/S)(W_S^{\mathrm{g}})^{-1} \arrow[ul,"\sigma"']
\end{tikzcd}
\]

Here, $\rho$ and $\sigma$ are localization functors (as is evident from the definition).

Any functor or composite functor in the diagram on page 8, going from an inner category of the diagram to a bordering category, induces a functor on the localized categories for $W_S^{\mathrm{d}}$ (when the target category is one of the four on the left,
\msp{page 72}
recalled in $(*)$, p. 70), or for $W_S^{\mathrm{g}}$ (when these categories are among the four on the right-hand side, recalled in $(*')$), which localizers also coincide with $W_S^{\mathrm{f}}$ except in the sole case of the category $\Cat/S$. All the functors thus obtained reduce, up to equivalence, to two, according as the target localized category is $\HOT_W(S)$ (case of localization $W_S^{\mathrm{d}}$), or is $\HOT^o_W(S)$ (case of localization $W_S^{\mathrm{g}}$). We thus find functors (in the direction opposite to $\rho$ and to $\sigma$)
\printedpage{45}
% head: A. Grothendieck : Dérivateurs, Ch. VI, Édition des Universités de Montpellier II et Paris VII
% foot: Version 9 novembre 2001
\[
\begin{tikzcd}
 & \HOT^{\mathrm{lc}}_W(S) \arrow[dl,"i"'] \arrow[dr,"j"] & \\
\HOT_W(S) & & \HOT^o_W(S).
\end{tikzcd}
\]

The functors $\rho$, $\sigma$ are interchanged by duality $S\longmapsto S^o$, and likewise for $i$, $j$.

\ovale{IV}\quad The functors $\rho$, $i$ are adjoint to one another in that order, i.e. we have for two objects $\xi\in\Ob\HOT_W(S)$, $\eta\in\Ob\HOT^{\mathrm{lc}}_W(S)$
\[ \Hom(\rho(\xi),\eta)\simeq\Hom(\xi,i(\eta)) \]
(isomorphism of bifunctors). And dually the functors $\sigma$, $j$ are adjoint to one another % [sic] adjoint -> adjoints, accord du pluriel
in that order: for $\xi\in\Ob\HOT^o_W(S)$, $\eta\in\Ob\HOT^{\mathrm{lc}}_W(S)$
\[ \Hom(\xi,j(\eta))\simeq\Hom(\sigma(\xi),\eta). \]

\msp{page 73}
It follows in particular that \textit{the functors $i$, $j$ are fully faithful.}

\cercle{V}\quad We have obvious functors
\[
\begin{aligned}
&(*) &\qquad& \HOT_W(S) &\quad& \longrightarrow &\quad& \underline{\mathrm{Hom}}(S^o,\Hot_W)\\
&(*') && \HOT^o_W(S) && \longrightarrow && \underline{\mathrm{Hom}}(S,\Hot_W).
\end{aligned}
\]
The first is interpreted on the four categories $(*)$ (bordering on the left) on page 70, which serve indifferently to describe $\HOT_W(S)$, by % [sic] indifférement -> indifféremment, faute d'orthographe
\[ s\longmapsto\hot_W(X_s) \]
in the case of the three categories distinct from $\Cat/S$, and by
\[ s\longmapsto\hot_W(s\backslash X) \]
in the case of $\Cat/S$. The functorial dependence on $s$ is obvious in the case % [sic] dépendence -> dépendance, faute d'orthographe
of $\Cat/S$ (given by $s\longmapsto\hot_W(s\backslash X)$), and in the case of $\ul{Fibsc}\,S$, $\ul{Fib}\,S$, where we have, for $u : s\longrightarrow t$ in $S$, the base change
\[ u^* : X_t\longrightarrow X_s. \]
In the case of $\ul{Liss}\,S$, the functorial dependence is deduced from that of $\hot_W(s\backslash X)$, thanks % [sic] dépendence -> dépendance, faute d'orthographe
to the fact that the inclusion
\[ X_s\longrightarrow s\backslash X \]
is then a $W$-equivalence.

We have a dual interpretation of the map
\printedpage{46}
% head: A. Grothendieck : Dérivateurs, Ch. VI, Édition des Universités de Montpellier II et Paris VII
% foot: Version 9 novembre 2001
\msp{page 74}
$(*)'$, \texttt{[illegible word]}
\[ s\longmapsto\hot_W(X_s)\qquad\text{or}\qquad s\longmapsto\hot_W(X/s), \]
\DERmark{ENV0325}

with the functorial law in $s$ deduced from the functor $t\backslash X\longrightarrow s\backslash X$ \texttt{[rather} $X/s\longrightarrow X/t$\texttt{]} (for $u : s\longrightarrow t$) in the case of $\Cat/S$, from the cobase change functor in the cases of $\ul{Cofib}\,S$ and $\ul{Cofibsc}\,S$, and from the variant for $\Cat/S$ and from the $W$-equivalence
\[ X_s\longrightarrow X/s \]
in the case of $\ul{Prop}\,S$.

We then find this: for an object $\xi$ of $\HOT_W(S)$ (resp. $\HOT^o_W(S)$) to be in the essential image of $\HOT^{\mathrm{lc}}_W(S)$, it is necessary and sufficient that the associated functor $S^o\longrightarrow\Hot_W$ (resp. $S\longrightarrow\Hot_W$) be a local system (\footnote{the verification will still have to be written}), i.e. transform every arrow of $S$ into an isomorphism, i.e. factor through $(\Pi S)^o\longrightarrow\Hot_W$ resp. $\Pi S\longrightarrow\Hot_W$ (\footnote{where $\Pi S$ is the fundamental groupoid of $S$}). Thus for an $X$ in $\Cat/S$, this means either that the functors
\[ t\backslash X\longrightarrow s\backslash X \]
associated to $u : s\longrightarrow t$ in $S$ are all $W$-equivalences (case of $\HOT_W(S)$), or that the functors
\[ X/s\longrightarrow X/t \]
are (case of $\HOT^o_W(S)$). (NB. these conditions are by no means equivalent. But if $X$ \texttt{[is]} smooth \texttt{[over]} $S$,
\msp{page 75}
the first implies the second, and the converse is true if $X$ is proper over $S$, and in this case, the condition is equivalent to $X$ being in $\ul{Fibhot}$. Finally, if $X$ is already in $\ul{Fibhot}\,S$, then the two conditions are automatically satisfied (\footnote{We should also have introduced in the diagram $\ul{Liss}_0S$, $\ul{Prop}_0S$!}).
\[
\begin{tikzcd}[column sep=small, row sep=small]
\ul{Fibsc}_0S \arrow[dd,hook] \arrow[drr] & & & & & & & & \ul{Cofibsc}_0S \arrow[dll] \arrow[dd,hook] \\
 & & \ul{Fib}_0S \arrow[dd,hook] \arrow[dr,hook] & & \ul{Parf}\,S \arrow[dl,hook] \arrow[dd,hook] \arrow[dr,hook] & & \ul{Cofib}_0S \arrow[dl,hook] \arrow[dd,hook] & & \\
\ul{Fibsc}\,S \arrow[drr] & & & \ul{Liss}_0S \arrow[dr,hook] \arrow[dd,hook] & & \ul{Prop}_0S \arrow[dl,hook] \arrow[dd,hook] & & & \ul{Cofibsc}\,S \arrow[dll] \\
 & & \ul{Fib}\,S \arrow[dr,hook] & & \ul{Fibhot}\,S \arrow[dd,hook] \arrow[uuullll,dotted,bend right,"\Phi_0"'] \arrow[uuurrrr,dotted,bend left,"\Psi_0"'] & & \ul{Cofib}\,S \arrow[dl,hook] & & \\
 & & & \ul{Liss}\,S \arrow[dr,hook] & & \ul{Prop}\,S \arrow[dl,hook] & & & \\
 & & & & \Cat/S \arrow[uuullll,dotted,bend left,"\Phi"] \arrow[uuurrrr,dotted,bend right,"\Psi"'] & & & &
\end{tikzcd}
\]
\printedpage{47}
% head: A. Grothendieck : Dérivateurs, Ch. VI, Édition des Universités de Montpellier II et Paris VII
% foot: Version 9 novembre 2001
(\footnote{The diagram of \textit{fifteen} remarkable categories.}). Moreover, if two elements $\xi\in\Ob\HOT_W(S)$, $\xi'\in\Ob\HOT^o_W(S)$ come from the same $\eta\in\Ob\HOT^{\mathrm{lc}}_W(S)$, then the two corresponding functors
\[ (\Pi S)^o\xrightarrow{\ \varphi\ }\Hot_W\xleftarrow{\ \varphi'\ }\Pi S \]
are \textit{associated,} i.e. each is isomorphic to ``the inverse'' of the other.
\[
\left\{
\begin{array}{lcl}
\varphi^o(s) & = & \varphi(s)\\
\varphi^o(u) & = & \varphi(u)^{-1}.
\end{array}
\right.
\]

\ovale{VI}\quad NB In the new version of the diagram, where I have added the categories $\ul{Liss}_0S$ and $\ul{Prop}_0S$ to make the drawing more transparent, the categories called ``bordering'' become % [sic] appelés -> appelées, accord du participe
the ``ground floor'', and the ``inner'' categories, the ``first upper floor'' or ``upper floor''. All % [sic] rez-de-chausseée -> rez-de-chaussée, lettre en trop
7 squares of the diagram are cartesian. All the arrows are fully % [sic] catésiens -> cartésiens, lettre omise
faithful inclusions, except the four outer horizontal arrows from the categories of Cat-fibered or Cat-cofibered categories that are \textit{split}, with morphisms exclusively cartesian or cocartesian and even compatible with the splittings, to the corresponding categories of nonsplit structures.

\msp{page 76}
\[
\begin{tikzcd}[column sep=small, row sep=small]
 & & & & \ul{Fibhot}\,S \arrow[dd,hook,"\delta"] \arrow[dddllll,dotted,bend right,"\Phi_0"'] \arrow[dddrrrr,dotted,bend left,"\Psi_0"] & & & & \\
 & & & \ul{Liss}_0S \arrow[ur,hook] \arrow[dd,hook] & & \ul{Prop}_0S \arrow[ul,hook] \arrow[dd,hook] & & & \\
 & \ul{Fib}_0S \arrow[urr,hook] \arrow[dd,hook] \arrow[ddrrr,dotted,bend left] & & & \Cat/S \arrow[dd,"C"] \arrow[dddllll,dotted,bend right,"\Phi"'] \arrow[dddrrrr,dotted,bend left,"\Psi"] & & & \ul{Cofib}_0S \arrow[ull,hook] \arrow[dd,hook] \arrow[ddlll,dotted,bend right] & \\ % Reading verified against PDF leaf 9: sources Fib_0 S and Cofib_0 S, target middle Parf S.
\ul{Fibsc}_0S \arrow[ur] \arrow[dd,hook] & & & \ul{Liss}\,S \arrow[ur,hook] & & \ul{Prop}\,S \arrow[ul,hook] & & & \ul{Cofibsc}_0S \arrow[ul] \arrow[dd,hook] \\
 & \ul{Fib}\,S \arrow[urr,hook] & & & \ul{Parf}\,S \arrow[uuul] \arrow[uuur] \arrow[dd,dotted,"\mathrm{id}"] & & & \ul{Cofib}\,S \arrow[ull,hook] & \\
\ul{Fibsc}\,S \arrow[ur] & & & & & & & & \ul{Cofibsc}\,S \arrow[ul] \\
 & & & & \ul{Parf}\,S \arrow[uuul,dotted,hook] \arrow[uuur,dotted,hook] & & & &
\end{tikzcd}
\]

I have drawn in green dotted lines \texttt{[color in the original only]} the five functors $\Phi$, $\Psi$, $\Phi_0$, $\Psi_0$, $C$ which by passage to the localized categories give quasi-inverses for all the equivalences of categories stated in \cercle{I} \cercle{II} of the present Scholie. To be precise, the functor
\[ \Phi : \Cat/S\longrightarrow\ul{Fibsc}\,S \]
associates to every $X$ in $\Cat/S$ the split category over $S$
\[ s\longmapsto s\backslash X \]
\printedpage{48}
% head: A. Grothendieck : Dérivateurs, Ch. VI, Édition des Universités de Montpellier II et Paris VII
% foot: Version 9 novembre 2001
over $S$. This functor, and its restrictions to $\ul{Liss}\,S$, $\ul{Fib}\,S$ and, more generally, to any full subcategory $\mathcal{F}$ of $\Cat/S$ containing $\ul{Fib}\,S$, is compatible with the localizers $W_S$ on the one hand (on $\mathcal{F}$ and $\ul{Fibsc}\,S$), $W_S^{\mathrm{d}}$ on the other hand, and therefore passes to the corresponding localized categories, to provide the desired quasi-inverse of
\[ (\ul{Fibsc}\,S)\Sigma^{-1}\longrightarrow\mathcal{F}\Sigma^{-1}\qquad(\Sigma=W_S\text{ or }W_S^{\mathrm{d}}) \]

We can also, for any intermediate full subcategory $\mathcal{F}'$
\[ \ul{Fib}\,S\subseteq\mathcal{F}'\subseteq\mathcal{F}\subseteq\Cat/S, \]
regard $\Phi$ as a functor $\mathcal{F}\longrightarrow\mathcal{F}'$ (more precisely, take the composite
\[ \mathcal{F}\xrightarrow{\Phi|\mathcal{F}}\ul{Fibsc}\,S\longrightarrow\ul{Fib}\,S\hookrightarrow\ \mathcal{F}'\ ), \]
say $\Phi_{\mathcal{F}',\mathcal{F}}$, and from what has just been said it follows that the two functors
\[
\begin{tikzcd}
\mathcal{F}' \arrow[r,hook] & \mathcal{F} \arrow[l,bend left,"\Phi_{\mathcal{F}',\mathcal{F}}"]
\end{tikzcd}
\]
\msp{page 77}
establish equivalences quasi-inverse to one another
\[
\begin{tikzcd}
\mathcal{F}'\Sigma^{-1} \arrow[r,"\approx"] & \mathcal{F}\Sigma^{-1} \arrow[l,bend left,"\approx"',"\Phi_{\mathcal{F}',\mathcal{F}}"]
\end{tikzcd}
\qquad\qquad(\Sigma=W_S\text{ or }W_S^{\mathrm{d}})
\]

Dual statement for the functor
\[ \Psi : \Cat/S\longrightarrow\ul{Cofibsc}\,S, \]
associating to $X$ in $\Cat/S$ the split cofibered category over $S$
\[ s\longmapsto X/s. \]
It gives pairs of quasi-inverse functors
\[
\begin{tikzcd}
(\ul{Cofibsc}\,S)\Sigma^{-1} \arrow[r] & \mathcal{F}\Sigma^{-1} \arrow[l,bend left,"\Psi_{\mathcal{F}}"]
\end{tikzcd}
\qquad\qquad(\Sigma=W_S\text{ or }W_S^{\mathrm{g}})
\]
for any full subcategory $\mathcal{F}$ of $\Cat/S$ such that % [sic] tout -> toute, accord avec sous-catégorie
\[ \ul{Cofib}\,S\subseteq\mathcal{F}\subseteq\Cat/S, \]
and pairs of quasi-inverse functors


\endgroup
% END UNIT Der-ChVI-345/Der-ChVI-345_p046-050.tex
% BEGIN UNIT Der-ChVI-345/Der-ChVI-345_p051-055.tex
\begingroup\makeatletter%
\newcommand{\printedpage}[1]{\marginpar{\footnotesize p.~#1}}%
\newcommand{\msp}[1]{\marginpar{\footnotesize [#1]}}%
\newcommand{\usl}[1]{\underline{\mathrm{#1}}}%
\newcommand{\encadre}[1]{\tikz[baseline=(n.base)]\node[draw,rounded corners=3pt,inner sep=1.5pt](n){#1};}%
\makeatother\let\printedpage\derpp\let\msp\dermsp%
\let\mspin\DERmspin%


\printedpage{49}
% head: A. Grothendieck : Dérivateurs, Ch. VI, Édition des Universités de Montpellier II et Paris VII
% foot: Version 9 novembre 2001

\[
\begin{tikzcd}
\mathcal{F}'\Sigma^{-1} \arrow[r, "\approx"] & \mathcal{F}\Sigma^{-1} \arrow[l, bend left=40, "\approx"', "\Psi_{\mathcal{F}',\mathcal{F}}"]
\end{tikzcd}
\qquad\qquad (\Sigma = W_S \text{ or } W_S^{\mathrm{g}})
\]

for two full subcategories $\mathcal{F}'$, $\mathcal{F}$
\[
\usl{Cofib}\,S \subseteq \mathcal{F}' \subseteq \mathcal{F} \subseteq \mathrm{Cat}/S.
\]
The functors $\Phi$, $\Psi$ send the full subcategory $\usl{Fibhot}\,S$ of ``the upper floor'', into the subcategories $\usl{Fibsc}_0S$, $\usl{Cofibsc}_0S$, respectively, hence induce on the upper floor functors $\Phi_0$, $\Psi_0$
\[
\usl{Fibsc}_0S \xleftarrow{\Phi_0} \usl{Fibhot}\,S \xrightarrow{\Psi_0} \usl{Cofibsc}_0S,
\]
playing on the upper floor exactly the same role as $\Phi$, $\Psi$ on the ground floor. The only difference is that on the upper floor we have $W_S = W_S^{\mathrm{d}} = W_S^{\mathrm{g}}$ $(= W_S^{\mathrm{f}})$,
\msp{page 78}
so there is only one sequence of localized categories to examine. We thus find
\[
\begin{tikzcd}
(\usl{Fibsc}_0S)W_S^{-1} \arrow[r, "\approx"] & \mathcal{F}W_S^{-1} \arrow[l, bend left=40, "\approx"', "\Phi_{0,\mathcal{F}}"]
\end{tikzcd}
\]
\[
\begin{tikzcd}
\mathcal{F}'W_S^{-1} \arrow[r, "\approx"] & \mathcal{F}W_S^{-1} \arrow[l, bend left=40, "\approx"', "\Phi_{0,\mathcal{F}',\mathcal{F}}"]
\end{tikzcd}
\]
for $\mathcal{F}$ (or $\mathcal{F}$, $\mathcal{F}'$) full subcategory(ies) of $\usl{Fibhot}\,S$ such that
\[
\usl{Fib}_0S \subseteq \mathcal{F} \subseteq \usl{Fibhot}\,S \qquad\text{resp.}\qquad \usl{Fib}_0S \subseteq \mathcal{F}' \subseteq \mathcal{F} \subseteq \usl{Fibhot}\,S,
\]
and likewise on the right-hand side for the full subcategories $\mathcal{F}$ or $\mathcal{F}$, $\mathcal{F}'$ of $\usl{Fibhot}\,S$ satisfying
\[
\usl{Cofib}_0S \subseteq \mathcal{F} \subseteq \usl{Fibhot}\,S \qquad\text{resp.}\qquad \usl{Cofib}_0S \subseteq \mathcal{F}' \subseteq \mathcal{F} \subseteq \usl{Fibhot}\,S,
\]
NB. The category $\usl{Parf}\,S$ contains both $\usl{Fib}_0S$ and $\usl{Cofib}_0S$, and is therefore among the admissible categories just like $\mathcal{F}$, or [?] $\mathcal{F}'$.

Thus, the four functors $\Phi$, $\Psi$, $\Phi_0$, $\Psi_0$ give all the quasi-inverses one can wish for on the ground floor on the one hand ($\Phi$ and $\Psi$), on the upper floor on the other ($\Phi_0$, $\Psi_0$). They therefore reduce all the relevant localized categories to the following four
\[
\begin{gathered}
(\mathrm{Cat}/S)W_S^{-1},\ \ (\mathrm{Cat}/S)(W_S^{\mathrm{d}})^{-1},\ \ (\mathrm{Cat}/S)(W_S^{\mathrm{g}})^{-1} \qquad \text{(on the ground floor)}\\
(\usl{Fibhot}\,S)W_S^{-1} \qquad \text{(on the upper floor).}
\end{gathered}
\]

\encadre{VII} It remains to compare the localized category on the upper floor with the localized categories on the ground floor. The vertical arrows descending from the upper floor to the ground floor, are compatible with $W_S^{\mathrm{d}}$ on the left-hand side of the diagram, with $W_S^{\mathrm{g}}$ on the right-hand side, so they induce essentially \textit{two} different functors
\printedpage{50}
% head: A. Grothendieck : Dérivateurs, Ch. VI, Édition des Universités de Montpellier II et Paris VII
% foot: Version 9 novembre 2001
\msp{page 79}
$i$, $j$ on $(\usl{Fibhot}\,S)W_S^{-1}$, into $(\mathrm{Cat}/S)(W_S^{\mathrm{d}})^{-1}$ and $(\mathrm{Cat}/S)(W_S^{\mathrm{g}})^{-1}$:
\[
\begin{tikzcd}[column sep=small, row sep=large]
& (\usl{Fibhot}\,S)W_S^{-1} \arrow[dl, "i"'] \arrow[dr, "j"] \arrow[dd, dotted, "\approx"', "\delta"] & \\
\mathrm{HOT}_W(S) \overset{\mathrm{def}}{\approx} (\mathrm{Cat}/S)(W_S^{\mathrm{d}})^{-1} \arrow[dr, dotted, "\rho"'] & & (\mathrm{Cat}/S)(W_S^{\mathrm{g}})^{-1} \overset{\mathrm{def}}{\approx} \mathrm{HOT}^{o}_{W}(S) \arrow[dl, dotted, "\sigma"] \\
& \begin{array}{l} (\mathrm{Cat}/S)W_S^{-1} \\ \overset{\mathrm{def}}{\approx} \mathrm{HOT}^{\mathrm{lc}}_{W}(S) \end{array} &
\end{tikzcd}
\qquad ,
\]
\DERmark{ENV0326}\DERmark{ENV0327}\DERmark{ENV0328}

so as to make the square completed by the dotted lines commute, representing the localization functors $\rho$, $\sigma$. We have of course
\[
\begin{aligned}
(*) \qquad \rho i = \sigma j : (\usl{Fibhot}\,S)W_S^{-1} \longrightarrow (\mathrm{Cat}/S)W_S^{-1}\\
\text{induced by the inclusion}\\
\delta : \usl{Fibhot}\,S \hookrightarrow \mathrm{Cat}/S
\end{aligned}
\]
We have stated in \encadre{I} that this arrow is an equivalence, so that $\rho$ and $\sigma$ can be interpreted, up to isomorphism of functors, as retractions of $\mathrm{HOT}_W(S)$ and of $\mathrm{HOT}^{o}_{W}(S)$ onto a common subcategory, interpreted either as $(\usl{Fibhot}\,S)W_S^{-1}$, or as $(\mathrm{Cat}/S)W_S^{-1}$. It remains to describe the quasi-inverse functor of $(*)$, which will at the same time give the quasi-inverses of all the vertical functors (from the upper floor to the ground floor) in the diagram p. 76, more precisely for their localized versions, for $W_S$.

\msp{page 80}
This is what is obtained with the help of the ``ascending'' functor (from the ground floor to the upper floor) $C$, which one can regard as a functor $\mathrm{Cat}/S \longrightarrow \usl{Fibhot}\,S$, but which in fact takes its values in $\usl{Parf}\,S$. (This is a detail of no great importance, it seems to me.)
\[
\mathrm{Cat}/S \xrightarrow{C} \usl{Parf}\,S \hookrightarrow \usl{Fibhot}\,S.
\]
This functor is compatible with the localizers $W_S$ on both sides, and establishes equivalences (quasi-inverses of the equivalences
\[
\begin{gathered}
(\usl{Parf}\,S)W_S^{-1} \xrightarrow{\sim} (\mathrm{Cat}/S)W_S^{-1} \qquad\text{and}\qquad (\usl{Fibhot}\,S)W_S^{-1} \xrightarrow{\sim} (\mathrm{Cat}/S)W_S^{-1}),\\
(\mathrm{Cat}/S)W_S^{-1} \xrightarrow{\sim} (\usl{Parf}\,S)W_S^{-1} \left( \xrightarrow{\sim} (\usl{Fibhot}\,S)W_S^{-1} \right).
\end{gathered}
\]

\printedpage{51}
% head: A. Grothendieck : Dérivateurs, Ch. VI, Édition des Universités de Montpellier II et Paris VII
% foot: Version 9 novembre 2001
So, when the three HOT categories for $S$ are interpreted in terms of localized categories of $\mathrm{Cat}/S$, it is the functor $C$ (more precisely the composite $\mathrm{Cat}/S \xrightarrow{C} \usl{Parf}\,S \hookrightarrow \mathrm{Cat}/S$) that provides both the inclusion functor (!) $i$, and the inclusion functor $j$ into $\mathrm{HOT}_W(S)$ and $\mathrm{HOT}^{o}_{W}(S)$ respectively. The functor in question
\[
C' : \mathrm{Cat}/S \longrightarrow \mathrm{Cat}/S
\]
does indeed satisfy
\[
C'(W_S) \subseteq W_S^{\mathrm{d}} \cap W_S^{\mathrm{g}}
\]
and defines the two functors
\[
\begin{tikzcd}[column sep=large]
& (\mathrm{Cat}/S)(W_S^{\mathrm{d}})^{-1} \approx \mathrm{HOT}_W(S) \\
\begin{array}{l} (\mathrm{Cat}/S)W_S^{-1} \\ \approx \mathrm{HOT}^{\mathrm{lc}}_{W}(S) \end{array} \arrow[ur, "i"] \arrow[dr, "j"'] & \\
& (\mathrm{Cat}/S)(W_S^{\mathrm{g}})^{-1} \approx \mathrm{HOT}^{o}_{W}(S)
\end{tikzcd}
\]
This allows us to make explicit the adjunction formulas of \encadre{IV}, when we interpret $\xi$, $\eta$ there
\msp{page 81}
as coming from $X$, $Y$ in $\mathrm{Cat}/S$:
\[
\boxed{\mathrm{Hom}_{W_S^{\mathrm{d}}}(X,CY) \simeq \mathrm{Hom}_{W_S^{\mathrm{g}}}(X,CY) \simeq \mathrm{Hom}_{W_S}(X,Y),}
\]
where $\mathrm{Hom}_{W_S^{\mathrm{d}}}$, $\mathrm{Hom}_{W_S^{\mathrm{g}}}$ and $\mathrm{Hom}_{W_S}$ denote respectively the Hom in the localized categories for $W_S^{\mathrm{d}}$, $W_S^{\mathrm{g}}$ and $W_S$, i.e. in $\mathrm{HOT}_W(S)$, $\mathrm{HOT}^{o}_{W}(S)$ and $\mathrm{HOT}^{\mathrm{lc}}_{W}(S)$ respectively (\footnote[51]{This formula suggests that if $X \in \mathrm{Ob}\,\mathrm{Cat}/S$, $Y \in \mathrm{Ob}\,\usl{Fibhot}\,S$, then
\[
\mathrm{Hom}_{W_S^{\mathrm{d}}}(X,Y) \simeq \mathrm{Hom}_{W_S^{\mathrm{g}}}(X,Y) \simeq \mathrm{Hom}_{W_S}(X,Y),
\]
but this follows from the adjunction formula and from $Y \longrightarrow C(Y)$ in $W_S^{\mathrm{d}} \cap W_S^{\mathrm{g}}$ for $Y$ in $\usl{Fibhot}\,S$.}). % [sic] adjoinction -> adjonction, spelling (elsewhere "adjonction")

\textbf{Comment.} All the constructions and statements from \encadre{I} to \encadre{VI}, with the exception of the fact that the descending functors induce equivalences for the \textit{coarse} localized categories, and of the adjunction formulas of \encadre{IV}, are of a virtually sorital nature, and require no conditions on the basic localizer $W \subseteq \mathrm{Fl}\,\mathrm{Cat}$ other than L1, L3 quater, L4 bis, summarized as a) b) c) on page 45. (We do not even have to use d).) It is true that in the absence of condition L5 (page 43), we do not know whether on the upper floor the two relevant localizers, $W_S$ and $W_S^{\mathrm{f}}$ (equal to $W_S^{\mathrm{d}}$ on the left-hand side, to $W_S^{\mathrm{g}}$ on the right-hand side) are equal. Thus a priori we find on the upper floor \textit{two} essentially distinct localized categories, and not just one, namely $(\usl{Fibhot}\,S)\Sigma^{-1}$ with $\Sigma = W_S$ or $\Sigma = W_S^{\mathrm{f}}$ $(= W_S^{\mathrm{d}} = W_S^{\mathrm{g}})$.

\printedpage{52}
% head: A. Grothendieck : Dérivateurs, Ch. VI, Édition des Universités de Montpellier II et Paris VII
% foot: Version 9 novembre 2001
On the other hand, the fact that the descent functors induce equivalences between the coarse localized categories is deeper, it requires (it seems to me) that we know an ``ascending'' functor $C$, say from $\mathrm{Cat}/S$ to $\usl{Fibhot}\,S$, compatible with $W_S$, and equipped with
\[
X \xrightarrow{i_X} C(X),
\]
which is in $W_S$ for arbitrary $X$. When we know that Cat is a Quillen category, even without having the properties of the functors $\usl{Ch}_\infty$ or $\usl{Ch}$ (\footnote[52]{nor even being sure of the existence of a \textit{functorial} factorization of $X \xrightarrow{f} S$ into $X \xrightarrow{i_X} C(X) \xrightarrow{\bar{f}} S$ with $i_X \in W$ and $\bar{f}$ perfect.}), we should probably
\msp{page 82}
be able to manage. But the adjunction formula seems to me to remain unexpected and deep, even admitting the results of \textsc{Quillen} and \textsc{Thomason} (\footnote[53]{See however theorem-scholie 2, p. 108-113. It is after all not as deep as it seemed to me, but virtually sorital!}). It seems linked here to a rather special property of the functor $C$, namely the lemma p. 68
\[
\begin{gathered}
i_{C(X)} \underset{\big\uparrow}{=} C(i_X) : C(X) \longrightarrow C(C(X)),\\
\scriptstyle \text{in } \underbrace{(\mathrm{Cat}/S)(W_S^{\mathrm{d}})^{-1}}_{\mathrm{HOT}_W(S)} \text{ and in } \underbrace{(\mathrm{Cat}/S)(W_S^{\mathrm{g}})^{-1}}_{\mathrm{HOT}^{o}_{W}(S)}
\end{gathered}
\]
which does not look purely formal. It would be interesting to examine to what extent this formula still makes sense in a more or less arbitrary ``category of models'' $(\mathcal{M},W)$, or coming from a Quillen triple $(W,\mathrm{Cof},\mathrm{Fib})$ in $\mathcal{M}$. It is clear that the factorization $X \xrightarrow{f} S$ into $X \xrightarrow{i_X} C(X) \xrightarrow{\bar{f}} S$ $(i_X \in W,\ \bar{f} \in \mathrm{Parf})$, playing the role of ``fibrations'' in the Quillen context, has an intrinsic meaning independent of the choice of a particular construction, something we can make perfectly precise. It is the description of the localizers $W_S^{\mathrm{d}}$, $W_S^{\mathrm{g}}$ in $\mathcal{M}/S$ (for $S \in \mathrm{Ob}\,\mathcal{M}$) that is not understood in a general way, outside the case $\mathcal{M} = \mathrm{Cat}$. In other words, I do not understand how to construct categories $\mathrm{HOT}_{\mathcal{M},W}(S)$, $\mathrm{HOT}^{o}_{\mathcal{M},W}(S)$, even if $(\mathcal{M},W)$ comes from a Quillen triple, I see only the description of a category $\mathrm{HOT}^{\mathrm{lc}}_{\mathcal{M},W}(S) = (\mathcal{M}/S)(W_S^{\text{coarse}})^{-1}$, and a localization functor $\mathcal{M}/S \longrightarrow \mathrm{HOT}^{\mathrm{lc}}_{\mathcal{M},W}(S)$. We have to manage to describe in $\mathcal{M}/S$ $W$-equivalences (\textit{``relative to $S$''}) finer than the coarse one given by $W$ on the source $X$ and nothing more. In Cat I manage it by using either the family of local isomorphisms (to describe $W_S^{\mathrm{g}}$), or that of colocal isomorphisms (to describe $W_S^{\mathrm{d}}$) - in any case, there is a notion of ``localization'' (either on $S$, or on $S^{o}$ \dots), with % [sic] En d’autre termes -> En d’autres termes, fixed expression requires plural
\DERmark{ENV0329}

\msp{page 83}
all the richness of the intuitions linked to such a notion. Perhaps, without knowing it I have always treated Cat not as a ``bare'' category, but as a site, or rather (since there
\printedpage{53}
% head: A. Grothendieck : Dérivateurs, Ch. VI, Édition des Universités de Montpellier II et Paris VII
% foot: Version 9 novembre 2001
is the involution $X \longmapsto X^{o}$ which cannot be evaded), as a ``bi-site'' - we have \textit{two} ways of localizing there, either by the $X/x$ (Zariski topology!) or by the $x\backslash X$ (co-Zariski topology). There is also the smooth topology (probably very close to the intuition of étale cohomology), and the proper topology - and we expect the smooth one to give essentially the same results as the Zariski one (since the étale topology is equal to it - there is no need to make a fuss about it \dots); and the proper topology the same results (e.g. the same cohomology?) as the co-Zariski topology. % [sic] vue que -> vu que, conjunctional locution

It does seem that in all the categories of models used up to now to describe Hot as a localized category, we have notions of localization: (Top), $\boldsymbol{\Delta}^\wedge$ (or $A^\wedge$, with $A$ an arbitrary test category), (Cat); and in Cat we even have \textit{two} such notions. On the other hand I am thinking of $\infty$-groupoids, it is perhaps the only category of models already so close to Hot itself, that we discern no notion of localization in it, and that the very idea of wanting to localize in it (except by ``localizing'' Hot itself) seems almost
\msp{page 84}
preposterous.

This quick digression now convinces me that the said compatibility for the construction of $C(X)$ and $i_X : X \longrightarrow C(X)$, can in no case be purely formal, a mere matter of Quillen categories say. That will at best give the compatibility in $(\mathcal{M}/S)W_S^{-1}$, which is ridiculous! But if we take as localizer $\Sigma$ in $\mathcal{M}/S$ the isomorphisms and nothing more, the compatibility relation will become false every time! As a slightly more reasonable set of arrows $\Sigma$, which is still very small, we can think e.g. of \textit{universal} $W$-equivalence over $S$ (distinctly finer still than $W_S^{\mathrm{d}}$, $W_S^{\mathrm{g}}$ in the special case of Cat). Would we still have compatibility for this set of arrows? I am also a little perplexed by the fact (not well understood) that the adjunction formulas for the Zariski and co-Zariski topologies go in opposite directions. (Or might I have made a mistake in dualizing in a hurry?) But to tell the truth, this compatibility by itself, even when we use a site structure on $\mathcal{M}$, is not enough to write down an adjunction formula. We have two localized categories, % [sic] ne convainc -> me convainc, the verb needs an object pronoun
\DERmark{ENV0330}

\[
\mathrm{HOT}^{\mathrm{lc}}_{\mathcal{M}} = (\mathcal{M}/S)W_S^{-1}, \qquad \mathrm{HOT}_W(S) = (\mathcal{M}/S)(W_S^{\mathrm{f}})^{-1},
\]
where $W_S^{\mathrm{f}}$ denotes a ``fine'' $W$-equivalence, corresponding to the notion of localization, more precisely to the ``admissible'' base changes we have chosen, and we have a localization functor
\msp{page 85}
\[
\mathrm{Hot}_W(S) \longrightarrow \mathrm{Hot}^{\mathrm{lc}}_{W}(S),
\]
but no functor in the reverse direction. But yes, but yes, working with the similar category $\usl{Fibhot}\,S$, this must still make sense at least when we have a good notion of ``fibration''. I see a little work of conceptual roughing-out to do, with perhaps an adjunction formula at the end of it, but I am not going to launch into it now.


\endgroup
% END UNIT Der-ChVI-345/Der-ChVI-345_p051-055.tex
% BEGIN UNIT Der-ChVI-345/Der-ChVI-345_p056-060.tex
\begingroup\makeatletter%
\newcommand{\printedpage}[1]{\marginpar{\footnotesize p.~#1}}%
\newcommand{\msp}[1]{\marginpar{\footnotesize [#1]}}%
\newcommand{\cM}{\mathcal{M}}%
\newcommand{\cF}{\mathcal{F}}%
\newcommand{\Fib}{\underline{\mathrm{Fib}}\,}%
\newcommand{\Fibhot}{\underline{\mathrm{Fibhot}}\,}%
\newcommand{\Ho}{\mathrm{Ho}}%
\newcommand{\Fl}{\mathrm{Fl}}%
\newcommand{\Ob}{\mathrm{Ob}}%
\makeatother\let\printedpage\derpp\let\msp\dermsp%
\let\mspin\DERmspin%


\printedpage{54}
% head: A. Grothendieck : Dérivateurs, Ch. VI, Édition des Universités de Montpellier II et Paris VII

\begin{flushright}
16.11.90
\end{flushright}

I am going to try all the same to do some ``roughing-out'' work, for the diagram of the
15 categories looks a bit bushy, % [sic] peut -> peu, adverb "un peu" intended
after all! I give myself
\[
(\cM,W) \qquad W\subseteq \Fl\,\cM,
\]
an $S\in\Ob\,\cM$, whence $\cM/S$, with $W_S\subseteq\Fl(\cM/S)$ induced by $W$. We are going to put
\[
\Ho^{\mathrm{lc}}(S)=(\cM/S)W_S^{-1},
\]
this will play the role of relatively constant homotopy types. I suppose on the other hand
that we have singled out a class $F\subseteq\Fl(\cM)$ of arrows, playing the role of ``fibrations'' (these
will be the Hot-fibrant morphisms in Cat, giving rise to the subcategory $\Fibhot S$
of $\mathrm{Cat}/S$). I foresee that we will have to
\msp{page 86}suppose, besides $\cM$ being stable under finite $\varprojlim$, that $F$ is stable under arbitrary base change,
and stable under composition, that it contains the isomorphisms. And above all that we
have a factorization theorem
\[
\begin{tikzcd}
X \arrow[rr, "f"] \arrow[dr, "i_X"'] & & S \\
& \bar{X} \arrow[ur, "\bar{f}"'] &
\end{tikzcd}
\qquad\qquad\qquad
i_X\in W,\ \bar{f}\in F
\]
for % [sic] pout -> pour, misprint
every arrow $f$ in $\cM$. With \textsc{Quillen} one even requires that $i_X$ have the left lifting property with respect to $F$, but that should not be essential, since I did not
\DERmark{ENV0331}

have to use this fact in Cat. When I work in the category $\cM/S$, where $f$ is
understood for an object $X$ of the said category, I will write (by abuse of language) $C(X)$ for $\bar{X}$,
although it is not said that we have a canonical choice of $C(X)$, still less that $C(X)$
depends functorially on $X$.

Let us denote by
\[
\Fib S\subseteq\cM/S \tag*{$(*)$}
\]
the subcategory of $\cM/S$ formed by the $X$ whose structural morphism $f_X:X\longrightarrow S$ is
in $F$. Let us denote (by abuse of notation) by $W_S$ the localizer induced on $\Fib S$ by $W_S$,
then the inclusion $(*)$ defines
\[
(\Fib S)W_S^{-1}\xrightarrow{\ \delta\ }(\cM/S)W_S^{-1}\overset{\mathrm{def}}{=}\Ho^{\mathrm{lc}}_W(S).
\]
\DERmark{ENV0332}


\msp{page 87}A first question is to see that this is an equivalence. We see that it is essentially
surjective thanks to the existence of the $C(X)$, $i_X$. That is not much! We are going to try to define
an arrow in the opposite direction
\[
\bar{C}:(\cM/S)W_S^{-1}\longrightarrow(\Fib S)W_S^{-1},
\]
% foot: Version 9 novembre 2001
\printedpage{55}
% head: A. Grothendieck : Dérivateurs, Ch. VI, Édition des Universités de Montpellier II et Paris VII
i.e. a \textit{functor}
\[
\gamma:\cM/S\longrightarrow(\Fib S)W_S^{-1}
\]
which transforms $W_S\subseteq\Fl(\cM/S)$ into isomorphisms. We will choose, for $X$ in $\cM/S$, a
$C(X)$ in $\Fib S$ and an
\[
i_X:X\longrightarrow C(X)\ , \qquad i_X\in W_S,
\]
and we will put
\[
\gamma(X)=C(X).
\]
The difficulty is to show that this does not depend on the choices made, up to canonical
isomorphism, and that it is functorial in $X$ - and the second implies the first.

To see that it works, clearly we have to suppose something more. I see two
hypotheses which, either of them, could make it work: one is that we already give ourselves
$C$ as a functor and $i_X$ functorial in $X$ (this was the case in Cat just now).
The other is that
\msp{page 88}the $i_X$ have, with respect to $F$, the lifting property in the Quillen style. % [sic] à là -> à la, accent on article "la"

I begin with the hypothesis that $C$ is already a functor. The diagram
\[
\begin{tikzcd}[row sep=large]
X \arrow[r, "i_X"] \arrow[d, "u"'] & CX \arrow[d, "C(u)"] \\
X' \arrow[r, "i_{X'}"] & CX',
\end{tikzcd}
\]
where $i_X,\ i_{X'}\in W$, shows that $C$ transforms $W_S$ into $W_S$, hence into isomorphisms of
$\Ho^{\mathrm{lc}}_W(S)$. So we have got the sought-for functor $\gamma$. In fact we have functors
\[
\begin{tikzcd}
\Fib S \arrow[r, hook, "\overset{\text{incl.}}{D}"] & \cM/S \arrow[l, bend left=50, "C"]
\end{tikzcd}
\]
compatible with the localizers $W_S$, inducing $\delta$ and $\gamma$, and we have
\[
\mathrm{id}_{\cM/S}\longrightarrow DC, \qquad \text{in $W_S$ argument by argument,}
\]
for we have
\[
X\xrightarrow{\ i_X\ }DC(X)=C(X)
\]
with $i_X\in W_S$. Likewise, we have
\[
\mathrm{id}_{\Fib S}\longrightarrow CD \qquad \text{in $W_S$ argument by argument,}
\]
% foot: Version 9 novembre 2001
\printedpage{56}
% head: A. Grothendieck : Dérivateurs, Ch. VI, Édition des Universités de Montpellier II et Paris VII
with the same
\[
X\xrightarrow{\ i_X\ }CX \qquad \text{(this time $X$ in $\Fib S$).}
\]

\msp{page 89}So we have without tears the two functors $\delta$, $\gamma$ as quasi-inverses of each other. And the
lemma on page 11 even shows that for \textit{every} full subcategory $\cF$ of $\cM/S$, such that
\[
\Fib S\subseteq\cF\subseteq\cM/S,
\]
the inclusions induce equivalences on the $W_S$-localized categories
\[
(\Fib S)W_S^{-1}\xrightarrow{\ \sim\ }\cF W_S^{-1}\xrightarrow{\ \sim\ }\underbrace{(\cM/S)W_S^{-1}}_{\overset{\mathrm{def}}{=}\Ho^{\mathrm{lc}}_W(S)}.
\]
\DERmark{ENV0333}

I am now going to treat the case of Quillen type, to convince myself that the existence of the
functor $C$, and the particular choice one can make of it, is irrelevant. (This last point
is moreover obvious a posteriori, when we know that $\delta$ is an equivalence, for the quasi-inverse
$\gamma$ will then be unique up to unique isomorphism, and the choice of a functor $C$, and
of a functorial homomorphism $i:\mathrm{id}\longrightarrow C$, is only one way among others of making
the quasi-inverse explicit.)

\msp{page 90}We therefore place ourselves under the Quillenian hypothesis on the $i_X$. We choose $C(X)$, $i_X$ for
each $X$ in $\cM/S$, and we put as before
\[
\gamma(X)=C(X) \qquad \text{as an object of $(\cM/S)W_S^{-1}$.}
\]
Let us show that this is functorial in $X$. Let
\[
u:X\longrightarrow X'
\]
be in $\cM/S$, let us consider the harmless (commutative) diagram
\[
\begin{tikzcd}[row sep=large]
X \arrow[rr, "u"] \arrow[d, "i_X"'] \arrow[drr, loosely dotted, "v=i_{X'}u"] & & X' \arrow[d, "i_{X'}"] \\
C(X) \arrow[rr, dotted, "\bar{u}"] \arrow[dr] & & C(X') \arrow[dl, "\in F"] \\
& S, &
\end{tikzcd}
\]
let us complete it with the coarsely dotted arrow, $v=i_{X'}u$ (making the upper triangle
commutative). (\footnote[54]{Lemma: let
% footnote continues on printed p. 57
\[
\begin{tikzcd}[ampersand replacement=\&, row sep=large]
X \arrow[r, "u"] \arrow[d, "i"'] \& X' \arrow[d] \\
C \arrow[r, dotted] \& C'
\end{tikzcd}
\]
be in $\cM/S$ with $i\in\mathrm{Cof}_0$, $(C'\longrightarrow S)\in F=\mathrm{Fib}$, then the dotted arrow exists in $\cM/S$, so as
to make the square commutative. \texttt{[cf. Lemma p. 91]}}). The hypothesis on $i_X$ implies that the finely dotted arrow $\bar{u}$ exists,
so as to have two more commutative triangles. But this means exactly that there is
an $S$-\textit{morphism} $\bar{u}$, i.e. an arrow in $\cM/S$, making the square commutative
% foot: Version 9 novembre 2001
\printedpage{57}
% head: A. Grothendieck : Dérivateurs, Ch. VI, Édition des Universités de Montpellier II et Paris VII
\[
\begin{tikzcd}[row sep=large]
X \arrow[r, "u"] \arrow[d, "i_X"'] & X' \arrow[d, "i_{X'}"] \\
C(X) \arrow[r, "\bar{u}"] & C(X'),
\end{tikzcd}
\]
\msp{page 91}and we must take
\[
\gamma(u)=\ \text{image of $\bar{u}$ in $\Fl((\Fib S)W_S^{-1})$.}
\]
We have to prove that this arrow does not depend on the choice of $\bar{u}$. We must isolate a Lemma,
where we denote by $\mathrm{Cof}_0$ the set of arrows of $\cM$ which have % [sic] qu'ont -> qui ont, relative pronoun
the lifting property with
respect to the elements of $\mathrm{Fib}=F$. ($\mathrm{Cof}_0$ plays the role of trivial cofibrations.)

\textbf{Lemma.} \textit{Let there be a solid-line diagram in $\cM/S$}
\[
\begin{tikzcd}[row sep=large]
X \arrow[r, "u"] \arrow[d, "i"'] & X' \arrow[d, "i'"] \\
C \arrow[r, dotted, "v"] & C'
\end{tikzcd}
\]
\textit{with $i\in$ Cof$_0$, and $C'\in$ Ob$\,\Fib S$ i.e. $(C'\longrightarrow S)\in$ Fib. Then $v$ exists, so as to make
the square commutative. Moreover, if $i\in W$ and $C\in\Ob\,\Fib S$, the class of $v$ as an arrow of
$(\Fib S)W_S^{-1}$ is independent of the choice made.}

It remains to prove this independence,
\msp{page 92}which will follow at once from the corollary:

\textbf{Corollary.} \textit{Let there be a diagram in $\cM/S$}
\[
\begin{tikzcd}
X \arrow[r, "i"] & C \arrow[r, shift left, "v_1"] \arrow[r, shift right, "v_2"'] & C',
\end{tikzcd}
\]
\textit{with $C$ and $C'$ in $\Fib S$ (i.e. $C\longrightarrow S$ and $C'\longrightarrow S$ in Fib), and $i\in$ Cof$_0$ and $i\in W$. If
$v_1i=v_2i$, then $v_1$ and $v_2$ define the same arrow of $(\Fib S)W_S^{-1}$.}
% foot: Version 9 novembre 2001

\printedpage{58}
% head: A. Grothendieck : Dérivateurs, Ch. VI, Édition des Universités de Montpellier II et Paris VII
\textsc{Proof}. Let
\[
v=(v_1,v_2):C\longrightarrow C'\times_S C',
\]
\DERmark{ENV0334}

\texttt{[rather} $C'\times_S C'$\texttt{]} which by hypothesis makes the diagram commutative
\[
\begin{tikzcd}[row sep=large]
X \arrow[r, "w"] \arrow[d, "i"'] & C' \arrow[d, "\mathrm{diag}_{C'}"] \\
C \arrow[r, "v"] & C'\times_S C',
\end{tikzcd}
\]
where $w=v_1i=v_2i$. Let us factor $\mathrm{diag}_{C'}$ as
\[
C'\xrightarrow{\ i'\ }\bar{C}'\xrightarrow{\ \delta'\ }C'\times_S C', \qquad \text{with $\delta'\in\mathrm{Fib}$,}
\]
and $i'\in W$. We find a commutative square
\[
\begin{tikzcd}[row sep=large]
X \arrow[r, "w'=i'w"] \arrow[d, "i"'] & \bar{C}' \arrow[d, "\delta'"] & \\
C \arrow[r, "v"] \arrow[ur, dotted, "\bar{v}"] & C'\times_S C' \arrow[r, shift left, "\mathrm{pr}_1"] \arrow[r, shift right, "\mathrm{pr}_2"'] & C'
\end{tikzcd}
\]
with $i\in\mathrm{Cof}_0$, $\delta'\in\mathrm{Fib}$, whence the existence of the dotted arrow, making the two triangles
commutative. We want to prove $\mathrm{pr}_1\circ v=\mathrm{pr}_2\circ v$ in $(\Fib S)W_S^{-1}$, as $v=\delta'\bar{v}$ it suffices
to prove that we have $\mathrm{pr}_1\delta'=\mathrm{pr}_2\delta'$ in $(\Fib S)W_S^{-1}$.

\msp{page 93}We only have to check first (to justify the argument) that this assertion makes sense,
i.e. that $\bar{C}'$ is itself also in $\Fib S$. But the hypotheses of stability of $F=\mathrm{Fib}$ under
base change and composition ensure that $1^o$) $C'\times_S C'$ is in $\Fib S$ (which is
stable under products), and consequently, $2^o$) that $\bar{C}'$ is also, since $\delta'\in\mathrm{Fib}$. As $C'$ is
also in $\Fib S$ and $C'\xrightarrow{\ i'\ }\bar{C}'$ is in $W_S$, to check $\mathrm{pr}_1\delta'=\mathrm{pr}_2\delta'$ in the localized category of
$\Fib S$, it suffices to see that $(\mathrm{pr}_1\delta')i'=(\mathrm{pr}_2\delta')i'$, now since $\delta'i'=\mathrm{diag}_{C'}$, we have the equality even
in $\Fib S$ (without localizing).

We have therefore defined $\gamma(X)$, $\gamma(u)$ for objects and arrows of $\cM/S$. Transitivity for $\gamma(u)$ is
clear, and also $\gamma(\mathrm{id})=\mathrm{id}$. So we have got the functor
\[
\gamma:\cM/S\longrightarrow(\Fib S)W_S^{-1}.
\]
Let now $\cF$ be a full subcategory of $\cM/S$, with
\[
\Fib S\subseteq\cF\subseteq\cM/S,
\]
let us consider $\gamma$ as a functor with values in $\cF W_S^{-1}$. (I forgot to check that $\gamma$ transforms
$W_S$ into $W_S$, but this is tautological.) Thus we have functors
% foot: Version 9 novembre 2001


\endgroup
% END UNIT Der-ChVI-345/Der-ChVI-345_p056-060.tex
% BEGIN UNIT Der-ChVI-345/Der-ChVI-345_p061-065.tex
\begingroup\makeatletter%
\NoAutoSpacing%
\newcommand{\printedpage}[1]{\marginpar{\footnotesize p.~#1}}%
\newcommand{\msp}[1]{\marginpar{\footnotesize [#1]}}%
\newcommand{\Fibu}{\underline{\mathrm{Fib}}\,}%
\newcommand{\Ho}{\mathrm{Ho}}%
\newcommand{\WL}{W_S^{\mathcal{L}}}%
\newcommand{\defeq}{\overset{\mathrm{def}}{=}}%
\makeatother\let\printedpage\derpp\let\msp\dermsp%
\let\mspin\DERmspin%

\DERmark{ENV0335}


\printedpage{59}
% head: A. Grothendieck : Dérivateurs, Ch. VI, Édition des Universités de Montpellier II et Paris VII
\[
\begin{tikzcd}
\mathcal{F}W_S^{-1} \arrow[r, shift left, "\delta"] & (\mathcal{M}/S)W_S^{-1}, \arrow[l, shift left, "\gamma"]
\end{tikzcd}
\]
we must see that they are quasi-inverse to one \msp{page 94} another.

On the $\mathcal{M}/S$ side: we must compare $X$ and $\delta\gamma(X) = C(X)$, this is done by
\[
i_X : X \longrightarrow C(X)
\]
which is in $W_S$, hence the isomorphism in the localized category
\[
X \simeq \delta\gamma(X) \qquad \text{in } (\mathcal{M}/S)W_S^{-1}.
\]
We must check that this is functorial in $X$.

On the $\mathcal{F}$ side, compare $X$ and $\gamma\delta(X)$, i.e. $X$ and $C(X)$. Same story, we choose $i_X : X \simeq C(X) = \gamma\delta(X)$. Functoriality to be checked.

It is now a matter of introducing a finer localized category of $\mathcal{M}/S$. For this, we suppose given another class of morphisms, I call it $\mathcal{L}$ (as in ``localization'')
\[
\mathcal{L} \subseteq \mathrm{Fl}(\mathcal{M}),
\]
e.g. in the case of $\mathcal{M} = \mathrm{Cat}$, $\mathcal{L}$ will consist of the local isomorphisms (\footnote[55]{Or also, \textit{smooth} morphisms. That should give the same result.}) (if it is a matter of defining $W_S^{\mathrm{g}}$), or the colocal isomorphisms (\footnote[56]{Or also proper morphisms.}) (if it is a matter of defining $W_S^{\mathrm{d}}$). We shall assume that $\mathcal{L}$ contains the isomorphisms, and $\mathcal{L}$ stable under composition (\footnote[57]{Perhaps this is not needed.}), but not under base change. We then define
\[
\WL \subseteq \mathrm{Fl}(\mathcal{M}/S)
\]
by
\[
f \in \WL \overset{\mathrm{def}}{\Longleftrightarrow} (f \times_S S') \in W,\ \ \forall (S' \longrightarrow S) \in \mathcal{L}.
\]
\DERmark{ENV0336}

\msp{page 95}

Since $\mathcal{L}$ contains $\mathrm{id}_S$, we shall have
\[
\WL \subseteq W_S,
\]
hence a canonical functor
\[
\underbrace{(\mathcal{M}/S)(\WL)^{-1}}_{\defeq \Ho_W^{\mathcal{L}}(S)} \xrightarrow{\ \rho\ } \underbrace{(\mathcal{M}/S)(W_S)^{-1}}_{\defeq \Ho_W^{\mathrm{lc}}(S)} % [sic] \Ho_W^{\mathrm{lc}}\mathcal{L}(S) -> \Ho_W^{\mathrm{lc}}(S), cf. the definition \Ho_W^{\mathrm{lc}}(S) \defeq (\mathcal{M}/S)W_S^{-1} on p. 61
\]
\DERmark{ENV0337}

% foot: Version 9 novembre 2001
\printedpage{60}
% head: A. Grothendieck : Dérivateurs, Ch. VI, Édition des Universités de Montpellier II et Paris VII
which is a localization functor. I would like to define a functor in the reverse direction. For this, it is natural to assume (by analogy with the situation in Cat)

\textbf{Hypothesis 1.} \textit{On $\Fibu S$, $W_S = \WL$, i.e. if $f : X \longrightarrow Y$ is in $W$ and $X$, $Y$ in $\Fibu S$, then for every base change $S' \longrightarrow S$ which is in $\mathcal{L}$, $f' : X_{S'} \longrightarrow Y_{S'}$ is still in $W$.}

In the Quillen context, this hypothesis has a definite tendency to be satisfied whatever $\mathcal{L}$ may be, i.e. for \textit{every} base change $S' \longrightarrow S$ without restriction. In fact it suffices that the Quillen structure be ``proper'', more precisely (and less restrictively), that a \texttt{[morphism]} $\in W$ remain so after base change by a morphism $\in$ Fib. Let us single this out as a Lemma. % [sic] pout -> pour, misprint

\textbf{Lemma.} \textit{The hypothesis is satisfied if we assume} \msp{page 96} \textit{that $F =$ Fib and $W$ are related by the additional condition (besides the factorization $i_X$): for every cartesian square}
\[
\begin{tikzcd}
X \arrow[d, "f"'] & X' \arrow[l, "q"'] \arrow[d, "f'"] \\
S & S', \arrow[l, "p"']
\end{tikzcd}
\]
\textit{if $f$ is in Fib and $p$ in $W$, then $q \in W$.}

\textsc{Proof}. Let
\[
\begin{tikzcd}
X \arrow[rr, "f"] \arrow[dr, "p"'] & & Y \arrow[dl, "q"] \\
& S &
\end{tikzcd}
\]
with $p$, $q \in \mathrm{Fib}$, $f \in W$, let us show that for \textit{every} base change $g : S' \longrightarrow S$, $f_{S'} : X_{S'} \longrightarrow Y_{S'}$ is still in $W$. We factor $g$ as
\[
\begin{tikzcd}
S' \arrow[r, "i"] \arrow[rr, bend right=15, "g"'] & \bar{S}' \arrow[r, "\bar{g}"] & S,
\end{tikzcd}
\qquad\qquad i \in W,\ \bar{g} \in \mathrm{Fib},
\]
\DERmark{ENV0338}

\texttt{[rather} $\bar{g} \in$ \texttt{Fib]} and we are reduced to proving that the assertion is true in the two special cases where $g \in \mathrm{Fib}$, or $g \in W$. The case $g \in \mathrm{Fib}$ follows at once from the hypothesis of the Lemma (applied to the cartesian square
\[
\begin{tikzcd}
X' \arrow[r, "f'"] \arrow[d, "q"'] & Y' \arrow[d, "p"] \\
X \arrow[r, "f"] & Y,
\end{tikzcd}
\]
% foot: Version 9 novembre 2001
\printedpage{61}
% head: A. Grothendieck : Dérivateurs, Ch. VI, Édition des Universités de Montpellier II et Paris VII
where $p \in \mathrm{Fib}$.)

In the case where $S' \longrightarrow S$ is in $W$, it follows from the hypothesis of the Lemma applied to
\[
\begin{tikzcd}
Y \arrow[d, "\in\mathrm{Fib}"] & Y' \arrow[l, "p"'] \arrow[d] \\
S & S' \arrow[l]
\end{tikzcd}
\]
that \msp{page 97} $p : Y' \longrightarrow Y$ is $\in W$, and likewise for $q : X' \longrightarrow X$, and $f \in W$ again implies $f' \in W$.

Let us return to the hypothesis p. 95. We shall therefore have a functor
\[
(\Fibu S, W_S = \WL) \longrightarrow (\mathcal{M}/S, \WL),
\]
hence for the localized categories a horizontal functor
\[
\begin{tikzcd}
(\Fibu S)W_S^{-1} \arrow[r, "i_0"] \arrow[d, "\approx"] & (\mathcal{M}/S)(\WL)^{-1} \defeq \Ho_W^{\mathcal{L}}(S) \arrow[dl] \\
\Ho_W^{\mathrm{lc}}(S) \defeq (\mathcal{M}/S)W_S^{-1} &
\end{tikzcd}
\]
which factors the vertical functor induced by the inclusion $(\Fibu S, W_S) \longrightarrow (\mathcal{M}/S, W_S)$ - a functor which we have seen (in the two cases considered) is an equivalence. By inverting the latter and composing it with $i_0$, we thus find a functor
\[
\underbrace{\Ho_W^{\mathrm{lc}}(S)}_{\defeq (\mathcal{M}/S)W_S^{-1}} \quad \xrightarrow{\ i\ } \quad \underbrace{\Ho_W^{\mathcal{L}}(S)}_{\defeq (\mathcal{M}/S)(\WL)^{-1}} \qquad ,
\]
such that
\[
\rho i \simeq \mathrm{id},
\]
hence a \textit{retraction} of $\Ho_W^{\mathcal{L}}(S)$ onto $\Ho_W^{\mathrm{lc}}(S)$. We expect $i$ to be fully faithful, and more precisely still, \msp{page 98} that $\rho$ and $i$ be adjoint, $i$ being the right adjoint of $\rho$:
\[
\begin{tikzcd}
\Ho_W^{\mathrm{lc}}(S) \arrow[r, shift right, "i"'] & \Ho_W^{\mathcal{L}}(S), \arrow[l, shift right, "\rho"']
\end{tikzcd}
\]
% foot: Version 9 novembre 2001
\printedpage{62}
% head: A. Grothendieck : Dérivateurs, Ch. VI, Édition des Universités de Montpellier II et Paris VII
i.e. we seek
\[
\begin{aligned}
i\rho &\quad \underbrace{\xleftarrow{\ i\ }}_{\textit{not}\ \text{iso}} \quad \mathrm{id} \quad \text{in } \Ho_W^{\mathcal{L}} \quad (\text{localization by } \WL)\\
\rho i &\quad \underbrace{\xrightarrow{\ p\ }}_{\text{iso}} \quad \mathrm{id} \quad \text{in } \Ho_W^{\mathrm{lc}} \quad (\text{localization by } W_S),
\end{aligned}
\]
\texttt{[too many} $i$\texttt{'s]} so for $X$ in $\mathcal{M}/S$, we seek
\[
\begin{aligned}
i_X :&\quad X \quad \longrightarrow \quad C(X) \quad \text{at least in } (\mathcal{M}/S)(\WL)^{-1}\\
p_X :&\quad X \quad \longleftarrow \quad C(X) \quad \text{at least in } (\mathcal{M}/S)(W_S)^{-1}.
\end{aligned}
\]
We shall choose for $i_X$ the ``canonical'' arrow from $X$ to $C(X)$, defined in $\mathcal{M}/S$ itself, without having to localize. But in the ``Quillen case'' (when $C$ is not given as a functor $\mathcal{M}/S \longrightarrow \Fibu S$), we must see that this $i_X$ is functorial in $X$, when we look at it in the localized category $(\mathcal{M}/S)W_S^{-1}$. But this is contained in the construction of $C$ as a \textit{functor}
\[
\mathcal{M}/S \longrightarrow (\Fibu S)W_S^{-1}.
\]
We shall note that $i_X$ is in $W_S$ and no more, not in $\WL$, hence does not in general define an isomorphism of $\Ho_W^{\mathcal{L}}$.

We define $p_X$ with the help of $i_X$, as the inverse of $i_X$ \textit{in} $(\mathcal{M}/S)W_S^{-1}$, which exists since \msp{page 99} $i_X$ is in $W_S$, hence \texttt{[is an]} isomorphism in $(\mathcal{M}/S)W_S^{-1} = \Ho_W^{\mathrm{lc}}$, and $p$ is therefore \texttt{[an]} isomorphism. We thus recover $\rho i \simeq \mathrm{id}$, hence that $\rho$ is a \textit{retraction} of $\Ho_W^{\mathcal{L}}$ onto $\Ho_W^{\mathrm{lc}}$.

To prove the adjunction, proceeding as in the case of Cat, we are reduced to proving that for every $X$,
\[
\left\{
\begin{aligned}
&C(i_X) = i_{C(X)} : C(X) \rightrightarrows C(C(X))\\
&\textit{in } \Ho_W^{\mathcal{L}} \ (\text{localization by } \WL).
\end{aligned}
\right.
\]
I had managed this in Cat, using a morphism
\[
C(C(X)) \xrightarrow{\ c\ } C(X)
\]
(essentially, a path composition morphism) such that $c \in \WL$, and that $c \circ C(i_X) = c \circ i_{C(X)}$. That had looked very delicate, for it was with great difficulty that I had managed to construct a reasonable functor \underline{Ch} (``paths''), where we could really and truly compose paths! For the moment I give up axiomatizing the properties of the functor \underline{Ch} that I need to make the machine work (I should come back to this shortly). I am going to begin instead with the ``Quillen case'', with $i_X \in \mathrm{Cof}_0$.
% foot: Version 9 novembre 2001

\printedpage{63}
% head: A. Grothendieck : Dérivateurs, Ch. VI, Édition des Universités de Montpellier II et Paris VII
\msp{page 100}

We take up again the description of the ``functor'' $C$ (with values in $(\Fibu S)W_S^{-1}$, and then without quotation marks), we find that $C(i_X)$, where
\[
X \xrightarrow{u = i_X} X' = C(X)
\]
is defined simply (in $(\Fibu S)W_S^{-1}$) by ``the'' arrow $C(u)$ making the diagram commute (in $\mathcal M/ S$)
\DERmark{ENV0339}

\[
\begin{tikzcd}
X \arrow[r, "u\ (=i_X)"] \arrow[d, "i_X"'] & X'\ (= C(X)) \arrow[d, "i_{X'}\ (=i_{C(X)})"] \\
C(X) \arrow[r, "C(u)\ (=C(i_X))"] & C(X')\ (= C(C(X))),
\end{tikzcd}
\]
so that in $\mathcal{M}/S$ we have the diagram
\[
\begin{tikzcd}
X \arrow[r, "\overbrace{i_X}^{\in\mathrm{Cof}_0\cap W}"] & \overbrace{C(X)}^{\in\mathrm{Ob}\,\Fibu S} \arrow[r, shift left, "i_{C(X)}"] \arrow[r, shift right, "C(i_X)"'] & \overbrace{C(C(X))}^{\in\mathrm{Ob}\,\Fibu S},
\end{tikzcd}
\]
commutative, i.e. $i_{C(X)}i_X = C(i_X)i_X$. The corollary on page 92 tells us that the arrows $i_{C(X)}$ and $C(i_X)$ are equal in $(\Fibu S)W_S^{-1}$. But by hypothesis (p. 95), in $\Fibu S$ we have $W_S = \WL$, so the arrows are equal in $(\Fibu S)(\WL)^{-1}$, and a fortiori in $(\mathcal{M}/S)(\WL)^{-1}$, and we win.

Let us see what we can do in the case where we have a \textit{functor} $C$ at our disposal. We assume of course that this functor in $\mathcal{M}/S$ \msp{page 101} is not tied to the particular $S$ in $\mathcal{M}$ that we have chosen as base, that its construction comes in \textsc{Cartan-Serre} fashion from a functor
\[
X \longmapsto \mathrm{Ch}(X) : \mathcal{M} \longrightarrow \mathcal{M},
\]
(``path object'') and functorial morphisms
\[
X \xrightarrow{\alpha_X} \mathrm{Ch}(X) \xrightarrow{p_X} X \times X, \qquad\qquad \boxed{p_X\alpha_X = \mathrm{diag}_X}
\]
(\footnote[58]{\textbf{NB} We put $\mathrm{pr}_1 \circ p_X = \sigma_X$, $\mathrm{pr}_2 \circ p_X = \beta_X$ (``source'' and ``target'' arrows for the object $X$).}). I presume we shall have to assume above all the following

\textbf{Hypothesis 2.} \textit{$\alpha_X \in W$, $p_X \in$} Fib.

For an $X$ over $S$, we shall construct $C_S(X)$ or $C(X)$ as follows
% foot: Version 9 novembre 2001


\endgroup
% END UNIT Der-ChVI-345/Der-ChVI-345_p061-065.tex
% BEGIN UNIT Der-ChVI-345/Der-ChVI-345_p066-070.tex
\begingroup\makeatletter%
\newcommand{\printedpage}[1]{\marginpar{\footnotesize p.~#1}}%
\newcommand{\msp}[1]{\marginpar{\footnotesize [#1]}}%
\newcommand{\Fib}{\underline{\mathrm{Fib}}}%
\newcommand{\Ch}{\mathrm{Ch}}%
\newcommand{\pr}{\mathrm{pr}}%
\newcommand{\circled}[1]{\tikz[baseline=(c.base)]\node[draw,circle,inner sep=1pt](c){#1};}%
\makeatother\let\printedpage\derpp\let\msp\dermsp%
\let\mspin\DERmspin%


\printedpage{64}
% head: A. Grothendieck : Dérivateurs, Ch. VI, Édition des Universités de Montpellier II et Paris VII

\[
\begin{tikzcd}[column sep=large]
 & \Ch(S) \arrow[d, "p_S"] \arrow[dl, dotted, "\beta_S"'] \arrow[dd, dotted, bend right=40, "\sigma_S"'] \arrow[dr, phantom, "\text{cart.}"]
 & C_S(X)=C(X) \arrow[l, "\pi"'] \arrow[d, "p'"] \arrow[dll, bend right=40, "\bar f_X=f_{C(X)}"'] \\
S & S\times S \arrow[l, "\pr_2^S"'] \arrow[d, "\pr_1^S"] \arrow[dr, phantom, "\text{cart.}"]
 & X\times S \arrow[l, "\pi_1=f_X\times\mathrm{id}_S"] \arrow[d, "\pr_1"] \\
 & S & X(=S'). \arrow[l, "f_X"'] \arrow[uu, dotted, bend right=40, "i_X"']
\end{tikzcd}
\]

It is the inverse image of $(\Ch(S),\pr_1^Sp_S=\sigma_S)$ under $X\xrightarrow{f_X}S$, but we have made the factorization of $\sigma_S$ explicit in the diagram to reinterpret the structural arrow $\bar f_X=f_{C(X)}$, defined as
\[
f_{C(X)}=\beta_S\pi(=\pr_2^S\underbrace{p_S\pi}_{=\pi_1p'})
\]
by the formula
\[
f_{C(X)}=\underbrace{(\pr_2^S\pi_1)}_{\text{projection } \pr_2 \text{ from } X\times S \text{ to } S}p'.
\]

\msp{page 102}
Now
\[
\begin{aligned}
&p'\in\mathrm{Fib}, && \text{since obtained from } p_S\in\mathrm{Fib} \text{ by base change } \pi_1,\\
&(\pr_2 : X\times S\longrightarrow S)\in\mathrm{Fib} && \text{provided that } X\longrightarrow\underbrace{e}_{\text{final object of } \mathcal{M}}\quad \text{is in Fib (same reason),}
\end{aligned}
\]
so if $X\in\Fib\,e$, we have $f_{C(X)}\in\mathrm{Fib}$. Since this is an essential property that we want to have for \textit{every} $X$ in $\mathcal{M}/S$, we are going to make the additional hypothesis

\textbf{Hypothesis 3.} \textit{$\forall\,X\in\mathit{Ob}\,\mathcal{M}$, $X\longrightarrow e$ is in Fib.}

As for
\[
i_X : X\longrightarrow C(X),
\]
it is obtained from the section $\alpha_S$ of $(\Ch(S),\sigma_S)$ over $S$ by base change $X\longrightarrow S$. Let us now note that
\[
\begin{tikzcd}
\Ch(S) \arrow[r, shift left, "\sigma_S"] \arrow[r, shift right, "\beta_S"'] & S
\end{tikzcd}
\qquad\qquad\qquad \text{are in } W,
\]
since their composites with $S\xrightarrow{\alpha_S}\Ch(S)$ are the identity, and $\alpha_S\in W$. On the other hand, $\sigma_S=\pr_1^S\circ p_S$ is in Fib, as a composite of two elements in Fib (NB $S\longrightarrow e$ in Fib by Hypothesis 3). Thus $\sigma_S$, $\beta_S$ are in $\mathrm{Fib}\cap W$. We shall need % [sic] \pr_1^S\circ p -> \pr_1^S\circ p_S, la flèche est notée p_S dans le diagramme et plus haut (\pr_1^Sp_S=\sigma_S)
\DERmark{ENV0340}


\textbf{Hypothesis 4.} \textit{The set of arrows $\mathit{Fib}\cap W$ is stable under every base change.}

% foot: Version 9 novembre 2001

\printedpage{65}
% head: A. Grothendieck : Dérivateurs, Ch. VI, Édition des Universités de Montpellier II et Paris VII

\msp{page 103}
This therefore implies that
\[
C(X)\xrightarrow{\pr_1p'}X
\]
\DERmark{ENV0341}

\texttt{[rather} $\pr_1p'$\texttt{]} is $\in W$ (being obtained from $\sigma_S : \Ch(S)\longrightarrow S$ by base change), so every section $i_X : X\longrightarrow C(X)$ of $C(X)$ over $X$ is in $W$ \texttt{[rather in} $W$\texttt{]}. We thus do indeed have a factorization, functorial in $X$
\DERmark{ENV0342}

\[
\begin{tikzcd}
 & C(X) \arrow[dr, "\bar f_X=f_{C(X)}"] & \\
X \arrow[ur, "i_X"] \arrow[rr, "f_X"] & & S
\end{tikzcd}
\]
with
\[
i_X\in W,\qquad\qquad f_{C(X)}\in\mathrm{Fib}.
\]
This is all we had used to establish
\[
(\Fib\,S)W_S^{-1}\xrightarrow{\sim}(\mathcal{M}/S)W_S^{-1}.
\]
But it is now a matter, under the hypothesis where we introduce an $\mathcal{L}$, of proving that for every $X$ in $\mathcal{M}/S$, the two arrows
\[
\begin{tikzcd}
C(X) \arrow[r, shift left, "C(i_X)"] \arrow[r, shift right, "i_{C(X)}"'] & C(C(X))
\end{tikzcd}
\]
are equal in $(\mathcal{M}/S)(W_S^{\mathcal{L}})^{-1}$. Since by hypothesis we have $W_S^{\mathcal{L}}=W_S$ in $\Fib\,S$ (Hypothesis 1), and $C(X),\ C(C(X))\in\mathrm{Ob}\,\Fib\,S$,
\msp{page 104}
it suffices to prove that the two arrows are equal \textit{in} $(\Fib\,S)W_S^{-1}$. This statement is independent of the datum of $\mathcal{L}$, and we feel like proving it.

Perhaps there is a way to make a geometric language explicit, to find our way around. Let us identify the objects $X$ of $\mathcal{M}$ with the functors $\mathcal{M}^o\longrightarrow(\mathrm{Ens})$ that they define,
\[
X\longmapsto(T\longmapsto X(T)=\mathrm{Hom}_{\mathcal{M}}(T,X))\,.
\]
Thus, for every $T$ % [sic] pout -> pour, faute de frappe
\[
\Ch(X)(T)
\]
is a set, called \textit{the set of paths in $X(T)$.} It maps (by $p_X$) functorially (in $T$) to $(X\times X)(T)=X(T)\times X(T)$, i.e. every path $c$ in $X(T)$ has an origin $\sigma(c)$, and a target $\beta(c)$. Moreover, by $\alpha_X$, $X(T)$ maps (functorially in $T$) to the set of paths of $X(T)$: for every element $x$ in $X(T)$, there is therefore a well-defined path (denoted $1_x$ or $\alpha_x$) with source and target $x$, which we call the trivial, or empty, path. This said,

% foot: Version 9 novembre 2001

\printedpage{66}
% head: A. Grothendieck : Dérivateurs, Ch. VI, Édition des Universités de Montpellier II et Paris VII

\msp{page 105}
for $X$ in $\mathcal{M}/S$, $C(X)$ represents the following contrafunctor in $\mathcal{M}$: for every $T$ in $\mathcal{M}$,
\[
C(X)(T)=\Bigl\{(x,c)\Bigm| x\in X(T),\ c : f_X(x)\xrightarrow{c}s\ \text{ a path in } S(T), \text{ with source } f_X(x)\Bigr\}\,,
\]
where $f_X : X\longrightarrow S$ is the structural morphism.

The ``calculations'' done, two (?) days ago show that we likewise have an interpretation
\[
\begin{aligned}
C(C(X))(T)=\Bigl\{(x,c,c')\Bigm| x\in X(T),\ f_X(x)\xrightarrow{c}s\xrightarrow{c'}t,\\
c\text{ a path with source } f_X(x),\\
c'\text{ a path with source equal to } \beta(c)\Bigr\}
\end{aligned}
\]
and the maps $C(i_X)$ and $i_{C(X)}$ are identified respectively with
\DERmark{ENV0343}

\[
\begin{aligned}
(x,c)\quad&\longmapsto\quad(x,1_{f(x)},c)\quad f(x)\xrightarrow{1_f(x)}f(x)\xrightarrow{c}s\\
(x,c)\quad&\longmapsto\quad(x,c,1_{\beta(c)})\quad f(x)\longrightarrow s\xrightarrow{1_s}s
\end{aligned}
\]
\DERmark{ENV0344}

\texttt{[rather} $f(x)\xrightarrow{1_{f(x)}}f(x)\xrightarrow{c}s$\texttt{]}. To compare them, we have an irresistible urge to write
\[
c\circ 1_{f(x)}=1_s\circ c=c,
\]
\DERmark{ENV0345}

but do we have the right? We would have to be able to compose the paths, and that is very, very difficult!

But we can always axiomatize. We therefore suppose that for $X$ in $\mathcal{M}$, $\Ch(X)$
\msp{page 106}
is not just an amorphous structure of (so-called) ``paths'' in $X$, just the datum of the source and target of the paths and of the trivial paths (these structures on $\Ch(X)$ depending functorially on $X$), but that we have a composition of paths when they are ``composable'', i.e. an arrow from a fiber product
\[
(c,c')\longmapsto c'\circ c : (\Ch(X),\beta_X)\times_X(\Ch(X),\sigma_X)\xrightarrow{\pi_X}\Ch(X),
\]
i.e. we can compose paths in $X(T)$
\[
x\xrightarrow{c}y\xrightarrow{c'}z
\]
if and only if $\text{target}\,c=\text{source}\,c'$, and moreover we of course require that $\text{source}(c'\circ c)=\text{source}\,c$, $\text{target}(c'\circ c)=\text{target}\,c'$, which is expressed by a commutative diagram
\[
\begin{tikzcd}
(\Ch(X),\beta_X)\times_X(\Ch(X),\sigma_X) \arrow[rr] \arrow[dr] & & \Ch(X) \arrow[dl, "p_X"] \\
 & X\times X. &
\end{tikzcd}
\]

% foot: Version 9 novembre 2001

\printedpage{67}
% head: A. Grothendieck : Dérivateurs, Ch. VI, Édition des Universités de Montpellier II et Paris VII

The natural thought is that with this composition and these identities, the $X(T)$, $\Ch(X)(T)$ become the objects and the arrows of a \textit{category} - i.e. that $(X,\Ch(X),\alpha_X,p_X,\pi_X)$ define a structure of ``category in $\mathcal{M}$''. But for what we want to do now, we probably do not need everything, i.e.
\DERmark{ENV0346}

\msp{page 107}
the relevant axiom of associativity of composition and the one on units, but only that the $1_x$ be two-sided units for composition.

Supposing this additional structure on the functor Ch given, and returning to the consideration of the associated functors % [sic] supossant -> supposant, faute de frappe ; foncteurs -> foncteur, accord avec « du » et « associé »
\[
C : \mathcal{M}/S\longrightarrow\Fib\,S,
\]
we find a canonical morphism
\[
C(C(X))\xrightarrow{\gamma_X}C(X)
\]
by
\[
(x,c,c')\longmapsto(x,c'\circ c),
\]
and the axiom of two-sided units assures us at once that
\[
(\ast)\qquad\qquad\gamma_X\circ i_{C(X)}=\gamma_X\circ C(i_X)=\mathrm{id}_{C(X)}.
\]
Since $i_{C(X)}\in W$, this already proves that $\gamma_X\in W$, so $\gamma_X\in W_S$. So $\gamma_X$ becomes an isomorphism in $(\Fib\,S)W_S^{-1}$, and the relation $(\ast)$ does indeed imply $i_{C(X)}=C(i_X)$ in the localized category, q.e.d.

I am going to summarize again what I have worked out.

\msp{page 108}
\textbf{Theorem-Scholie 2.} Let $\mathcal{M}$ be a category, equipped with two sets of arrows
\[
W,\ \mathrm{Fib}\subseteq\mathrm{Fl}(\mathcal{M}).
\]
We make the following hypotheses.

\begin{enumerate}[label={}]
\item[$1^o$] $W$ contains the isomorphisms, if $f$, $g$ are composable, if two of $f$, $g$, $gf$ are in $W$, the third also.
\item[$2^o$] $\mathcal{M}$ is stable under finite $\varprojlim$, Fib contains the isomorphisms, is stable under composition and under base change (\footnote[59]{These stability hypotheses on Fib are also used in I under hypothesis B, pp.~92--93, and in II, in case A), taken up again in toto in the form A$'$ below (p.~112).}).
\end{enumerate}
\DERmark{ENV0347}


For $S$ in $\mathcal{M}$, \fbox{$\Fib\,S$} denotes the full subcategory of $\mathcal{M}/S$ formed by the $X\longrightarrow S$ which are in Fib. We distinguish two different versions of an axiom of factorization of an arbitrary $f : X\longrightarrow S$ in $\mathcal{M}$ as $X\xrightarrow{i}\bar X\xrightarrow{\bar f}S$, with $i\in W$, $\bar f\in\mathrm{Fib}$.

% foot: Version 9 novembre 2001

\printedpage{68}
% head: A. Grothendieck : Dérivateurs, Ch. VI, Édition des Universités de Montpellier II et Paris VII

\begin{enumerate}[label={}]
\item[A)] We are given, for every $S$ in $\mathcal{M}$, a functor
\[
\mathcal{M}/S\longrightarrow\mathcal{M}/S\qquad\qquad X\longmapsto C(X)\ \ (=\bar X)
\]
and a morphism of functors in $\mathrm{X}$
\[
X\xrightarrow{i_X}C(X)
\]
with
\[
i_X\in W,\qquad\qquad C(X)\in\mathrm{Ob}\,\Fib\,S.
\]
\item[B)] Denoting by
\[
\mathrm{Cof}_0\subseteq\mathrm{Fl}\,\mathcal{M}
\]
the set of arrows of $\mathcal{M}$ which have the left lifting property \texttt{[rather left]} with respect to Fib, we suppose ``in the style of \textsc{Quillen}'', that every arrow $f : X\longrightarrow S$ factors as
\DERmark{ENV0348}

\[
X\xrightarrow{i_X}C(X)\xrightarrow{\bar f}S,
\]
with
\[
i\in\mathrm{Cofib}_0\cap W,\qquad\qquad\bar f\in\mathrm{Fib}.
\]
\texttt{[where} $\mathrm{Cof}_0=\mathrm{Cofib}_0$\texttt{]}.
\end{enumerate}

\msp{page 109}
Let us take $S$ in $\mathcal{M}$, consider
\[
W_S\subseteq\mathrm{Fl}(\mathcal{M}/S)
\]
the set of arrows induced by $W$ (
\begin{tikzcd}[column sep=small, row sep=small, baseline={([yshift=.5ex]current bounding box.south)}]
f : X \arrow[rr] \arrow[dr] & & Y \arrow[dl] \\
 & S &
\end{tikzcd}
is in $W_S$ if and only if it is on forgetting $S$), and let us denote by $W_S$ (by abuse of notation) the set of arrows induced by $W_S$ on $\Fib\,S$, or more generally on any full subcategory $\mathcal{F}$ of $\mathcal{M}/S$, with
\[
\Fib\,S\subseteq\mathcal{F}\subseteq\mathcal{M}/S.
\]

\circled{I} The functor
\[
(\ast)\qquad\qquad(\Fib\,S)W_S^{-1}\xrightarrow{\sim}(\mathcal{M}/S)W_S^{-1}
\]
induced by the inclusion is an equivalence of categories, more generally for every $\mathcal{F}$ as above
\[
(\Fib\,S)W_S^{-1}\longrightarrow\mathcal{F}W_S^{-1}\longrightarrow(\mathcal{M}/S)W_S^{-1}
\]
are equivalences of categories. A quasi-inverse functor is obtained (in both cases A) and B)) with the help of the functor
\[
\gamma : \mathcal{M}/S\longrightarrow(\Fib\,S)W_S^{-1}
\]

% foot: Version 9 novembre 2001


\endgroup
% END UNIT Der-ChVI-345/Der-ChVI-345_p066-070.tex
% BEGIN UNIT Der-ChVI-345/Der-ChVI-345_p071-075.tex
\begingroup\makeatletter%
\newcommand{\printedpage}[1]{\marginpar{\footnotesize p.~#1}}%
\newcommand{\msp}[1]{\marginpar{\footnotesize [#1]}}%
\newcommand{\circled}[1]{\tikz[baseline=(c.base)]\node[draw,circle,inner sep=1pt](c){#1};}%
\newcommand{\Fib}{\underline{\mathrm{Fib}}\,}%
\newcommand{\Fibsc}{\underline{\mathrm{Fibsc}}\,}%
\newcommand{\Cofib}{\underline{\mathrm{Cofib}}\,}%
\newcommand{\Cofibsc}{\underline{\mathrm{Cofibsc}}\,}%
\newcommand{\Fibhot}{\underline{\mathrm{Fibhot}}\,}%
\newcommand{\Parf}{\underline{\mathrm{Parf}}\,}%
\newcommand{\Liss}{\underline{\mathrm{Liss}}}%
\newcommand{\Prop}{\underline{\mathrm{Prop}}}%
\newcommand{\Ho}{\mathrm{Ho}}%
\newcommand{\Hom}{\mathrm{Hom}}%
\newcommand{\HOT}{\mathrm{HOT}}%
\newcommand{\hot}{\mathrm{hot}}%
\newcommand{\Fl}{\mathrm{Fl}}%
\newcommand{\Cat}{\mathrm{Cat}}%
\newcommand{\M}{\mathcal{M}}%
\makeatother\let\printedpage\derpp\let\msp\dermsp%
\let\mspin\DERmspin%


\setcounter{footnote}{59}

\printedpage{69}
% head: A. Grothendieck : Dérivateurs, Ch. VI, Édition des Universités de Montpellier II et Paris VII
which is given on objects by
\[\gamma(X) = C(X),\]
and which on arrows is defined in the obvious way in case A) (functoriality of $C(X)$ in $X$), and in the case of B) is constructed ad hoc, as the only functorial law for which the homomorphisms
\[i_X : X \longrightarrow C(X)\]
\msp{page 110}
(interpreted as arrows in $(\M/S)W_S^{-1}$) are themselves functorial in $X$.

We denote by
\[\Ho^{\mathrm{lc}}_W(S) \qquad\qquad \text{(lc as in ``locally constant'')}\]
the category of fractions $(\M/S)W_S^{-1}$, equivalent to $(\Fib S)W_S^{-1}$.

\circled{II} Let us take a localizer $W_S^0$ in $\M/S$
\[W_S^0 \subseteq W_S \subseteq \Fl(\M/S),\]
with the formal hypothesis similar to that on $W$ (cf. 1$^o$). We assume that

\begin{itemize}[leftmargin=4em]
\item[3$^o$)] $W_S^0 \cap \Fl(\Fib S) = W_S|\Fib S$.
\end{itemize}

(Cf. pages 94-96 for a general way of obtaining such localizers. One must clearly see that to define relative ``homotopy types'' over $S$ the localizer $W_S$ is much too coarse, morally it gives only a category of locally constant relative homotopy types, and this is precisely what is shown by the equivalence of \circled{I} , where $\Fib S$ is visualized as made up of the ``locally trivial fibrations'' over $S$.) We then have a second localized category % [sic] l'equivalence -> l'équivalence, accent manquant
\[\Ho_{W_S^0}(S) \overset{\mathrm{def}}{=} (\M/S)(W_S^0)^{-1},\]
\DERmark{ENV0349}

and two functors
\[
\begin{tikzcd}[column sep=large]
\Ho_{W_S^0}(S) \arrow[r, shift left, "\rho"] & \Ho^{\mathrm{lc}}_W(S), \arrow[l, shift left, "i"]
\end{tikzcd}
\tag*{$(*)$}
\]
where $\rho$ is the obvious (strict) localization functor, due to the inclusion $W_S^0 \subseteq W_S$, and where $i$ is
\msp{page 111}
essentially the functor deduced from the inclusion
\[(\Fib S, \underbrace{W_S|\Fib S}_{=W_S^0|\Fib S}) \longrightarrow (\M/S, W_S^0)\]
by localization \textit{by means of $W_S^0$}, hence the functor
\[(\Fib S)(W_S^0)^{-1} \longrightarrow (\M/S)(W_S^0)^{-1}, \tag*{$(**)$}\]
% foot: Version 9 novembre 2001
\printedpage{70}
% head: A. Grothendieck : Dérivateurs, Ch. VI, Édition des Universités de Montpellier II et Paris VII
taking into account the equivalence $(*)$ of \circled{I} , i.e. $i$ is the composite of the quasi-inverse of $(*)$, with the preceding functor $(**)$. Explicitly, it is defined by
\[i(X) = \text{class of C(X) in } (\M/S)(W_S^0)^{-1},\]
taking into account the fact that $X \longmapsto C(X)$ transforms $W_S$ into $W_S^0$, thanks to condition 3$^o$) on $W_S^0$. % [sic] W_S^o -> W_S^0, lettre o au lieu du chiffre 0 utilisé partout ailleurs
\DERmark{ENV0350}


This being so, \textit{under hypothesis} B (``Quillen''), we have: in the pair of functors $(\rho, i)$ of $(*)$, page 110, \textit{$i$ is right adjoint to $\rho$,} i.e. we have for $\xi \in \mathrm{Ob}\,\Ho_{W_S^0}(S)$, $\eta \in \mathrm{Ob}\,\Ho^{\mathrm{lc}}_W(S)$
\[\Hom_{W_S^0}(\xi, i(\eta)) \simeq \Hom_{W_S}(\rho(\xi), \eta),\]
which is made explicit, for $X$, $Y$ in $\M/S$, by
\[\boxed{\Hom_{W_S^0}(X, C(Y)) \simeq \Hom_{W_S}(X, Y)}\]
(\footnote{\textbf{NB} The first member is independent of the choice of $W_S^0$! (Provided only that the hypothesis $W_S^0 \cap \Fl\,\Fib S = W_S \cap \Fl\,\Fib S$, in addition to $W_S^0 \subseteq W_S$, is satisfied.)}). It follows that \textit{the functor $i$ is fully faithful.}

The same adjunction formula, and the
\msp{page 112}
full faithfulness of $i$, hold in case A, when one strengthens hypothesis A (relative to a particular $S$) to a more restrictive and more comprehensive hypothesis, as follows.

\begin{itemize}[leftmargin=3em]
\item[A$'$)] The functor $C : \M/S \longrightarrow \M/S$ and the functorial morphism $i_X : X \longrightarrow C(X)$ are constructed ``in the manner of \textsc{Cartan-Serre}'' from a functor defined on the whole of $\M$
\[\M \xrightarrow{\mathrm{Ch}} \M,\]
of the ``paths'' type, i.e. equipped in addition with the following data
\begin{itemize}[leftmargin=3em]
\item[a)] \[X \xrightarrow{\alpha_X} \mathrm{Ch}(X) \xrightarrow{p_X = (\sigma_X, \beta_X)} X \times X\]
(structure of source and target of paths, and of trivial paths), and
\item[b)] A composition of paths, without requiring here the associativity of composition but only that trivial paths be two-sided units for composition. Moreover, we assume that for every $X$ % [sic] pout -> pour, faute de frappe
\[\begin{cases} \alpha_X \in W,\ p_X \in \mathrm{Fib} \\ (X \longrightarrow e) \in \mathrm{Fib}. \end{cases}\]
\end{itemize}
\end{itemize}
% foot: Version 9 novembre 2001

\printedpage{71}
% head: A. Grothendieck : Dérivateurs, Ch. VI, Édition des Universités de Montpellier II et Paris VII
\msp{page 113}
\textbf{Corollary.} Let $X$ be in $\M/S$, $Y$ in $\Fib S$, then
\[\Hom_{W_S^0}(X, Y) \xrightarrow{\sim} \Hom_{W_S}(X, Y).\]
\textsc{Proof}. By the adjunction formula (p. 111), the second member is $\Hom_{W_S^0}(X, CY)$. But as $Y$ and $C(Y)$ are in $\Fib S$, where $W_S$ and $W_S^0$ coincide, $Y \xrightarrow{i_Y} C(Y)$ is not only in $W_S$, but in $W_S^0$, hence $\Hom_{W_S^0}(X, Y) \xrightarrow{\sim} \Hom_{W_S^0}(X, CY)$, whence the result.

Let us return to Cat, Cat$/S$. The preceding digression on $(\M, W, \mathrm{Fib})$ has made it possible to better situate what concerns the functor $C$, and the relations between the localized categories of
\[\Cat/S \ \ (\text{for } W_S,\ W_S^{\mathrm{d}},\ W_S^{\mathrm{g}}) \qquad\quad \text{and} \qquad\quad \Parf S \ \ (\text{for } W_S^{\mathrm{f}} = W_S = W_S^{\mathrm{d}} = W_S^{\mathrm{g}})\]
as well as for the intermediate categories
\[\Parf S \subseteq \mathcal{F} \subseteq \Cat/S\]
full in Cat$/S$ (in the diagram p. 76, there are five others, namely $\Fibhot S$, $\Liss\,S$, $\Liss_0 S$, $\Prop\,S$, $\Prop_0\,S$). As for the relations between the localized categories of these seven categories, which are at the center of the diagram, the statements of theorem-scholie 1 (p. 70-81) give nothing more than the content of theorem-scholie 2 (p. 108-113). To establish them, we therefore have the choice between the two hypotheses A) or B) of the said theorem. Either use an explicit functor $C_S$, with the help of a
\msp{page 114}
functor $X \longmapsto \underline{\mathrm{Ch}}(X)$ (``paths in $X$'') on Cat, hence the structures of source, target and composition for paths, and trivial paths (acting as two-sided units for composition). The construction of such a functor $C$ does not look trivial and seems not to be known - I struggled to get there! (With, in addition, the associativity of composition \ldots). Or else use, in the manner of \textsc{Quillen}, the existence of a class of morphisms called Q-fibrant, with good factorization properties. This is not obvious in Cat either, and it is established by \textsc{Thomason}, by a somewhat complicated route. Perhaps I shall be led to come back to this - but only if necessary! % [sic] n'a par l'air -> n'a pas l'air, faute de frappe

On the other hand, the assertions concerning the categories $\Fib S$, $\Fibsc S$, or $\Cofib S$, $\Cofibsc S$, and their variants with indices 0 (thus the eight wing-categories of the diagram), and their relations to Cat$/S$ on the one hand, to $\Fibhot S$ on the other (for the categories with index 0), cannot be contained in a general carpet on categories of models $(\M, W)$, it seems to me, in view of the very special nature of the four categories designated above, of their very particular role in the categorical formalism. Now at least the ``split'' categories $\Fibsc S$, $\Cofibsc S$, and consequently also their variants with index 0, play an essential role for the theory of the HOT$(S)$, since it is they that
\msp{page 115}
intervene dually when we apply the general recipe for constructing a derivator from a category of models $(\M, W)$. So I see no simplification
% foot: Version 9 novembre 2001
\printedpage{72}
% head: A. Grothendieck : Dérivateurs, Ch. VI, Édition des Universités de Montpellier II et Paris VII
of substance to be made to theorem-scholie 1. At most one could prune the diagram of categories under consideration, by dropping $\Liss\,S$, $\Liss_0 S$, $\Prop\,S$, $\Prop_0 S$, which do not play a really useful role in it. As for the stated results, what concerns them is covered by what is said for the full subcategories $\mathcal{F}$ of Cat$/S$ which contain (say) $\Parf S$ (or at one's choice those which contain $\Fib S$ or $\Cofib S$).

The description of categories of the type HOT$(S)$, for fixed $S$, now seems to me perfectly well understood. It is time to let $S$ vary, and to understand the contravariant character in $S$. We thus take
\[S' \xrightarrow{f} S,\]
and we propose to describe the inverse image functors for the three types of categories HOT$(S)$. But we shall now restrict ourselves to the $\HOT_W(S)$ and $\HOT^{\mathrm{lc}}_W(S)$, since everything concerning $\HOT^o_W(S) = \HOT_W(S^o)$ \texttt{[rather} $\HOT_W(S^o)^o$\texttt{, cf. p. 50]} will follow from this by the duality $S \longmapsto S^o$. So it is a matter of describing
\msp{page 116}
\[f^* : \HOT_W(S) \longrightarrow \HOT_W(S') \tag*{$(*)$}\]
and the induced functor
\[f^* : \HOT^{\mathrm{lc}}_W(S) \longrightarrow \HOT^{\mathrm{lc}}_W(S').\]
When we describe the categories HOT$(S)$, in accordance with the general recipe, by means of functors
\[F : S^o \longrightarrow \Cat\]
(presheaves of categories on $S$), or, what amounts to the same thing, of split categories, the description of the arrow $(*)$ is obvious: it comes down to the law
\[\left|\begin{array}{rcl}
F & \longmapsto & F \circ f^o \\[1ex]
\underline{\mathrm{Hom}}(S^o, \Cat) & \longrightarrow & \underline{\mathrm{Hom}}(S'^o, \Cat),
\end{array}\right.\]
by passing to the localized categories for $W$-equivalence argument by argument. In the language of fibered categories, it is well known that this trivial operation on presheaves of categories corresponds to the \textit{inverse image} operation on categories over $S$ - which operation passes to splittings. This therefore gives us
\[f^*(\hot_{W,S}(X)) = \hot_{W,S'}(\underbrace{X \times_S S'}_{X'}),\]
for every $X$ split fibered over $S$. Now we know that the splitting of $X$ plays no role, in the end, in the description of the class
\[\hot_{W,S}(X) \in \HOT_W(S)\]
% foot: Version 9 novembre 2001
\printedpage{73}
% head: A. Grothendieck : Dérivateurs, Ch. VI, Édition des Universités de Montpellier II et Paris VII
\msp{page 117}
associated to $X$. We can forget the splitting, and even (to study the $\Hom_{\mathrm{Hot}_W(S)}(X, Y)$) work with $S$-functors \textit{of any kind} between categories fibered over $S$ - no need even to assume them cartesian, or even compatible with splittings. This seems to me a very important thing for the formalism of the $\HOT_W(S)$, and very special to the case of the category of models Cat. This is somewhat tacit in the theorem-scholie - but my initial doubts concerning the category $\underline{\mathrm{Fibsc}}$ itself show clearly that I felt there was something there that was not obvious, nor even (almost) \texttt{[?]} plausible. And yet the existence of the functor $\Phi$ shows that it is indeed so.

We can therefore interpret $f^*$ in terms of $\Fib S$, $\Fib S'$, by the base change functor
\[f^* : \Fib S \longrightarrow \Fib S', \qquad\quad X \longmapsto X \times_S S'/S',\]
which functor obviously preserves the fiberwise $W$-equivalence $W_S^{\mathrm{f}}$, $W_{S'}^{\mathrm{f}}$, which is the relevant one here
\[f^*(W_S^{\mathrm{f}}) \subseteq W_{S'}^{\mathrm{f}},\]
so that $f^*$ does indeed induce
\[(\Fib S)(W_S^{\mathrm{f}})^{-1} \longrightarrow (\Fib S')(W_S^{\mathrm{f}})^{-1}\]
\texttt{[rather} $(\Fib S')(W_{S'}^{\mathrm{f}})^{-1}$\texttt{]},
\msp{page 118}
i.e. the functor $(*)$
\[f^* : \HOT_W(S) \longrightarrow \HOT_W(S').\]
On the other hand, $\HOT_W(S)$ can also be described as a localized category of Cat$/S$ as a whole, with respect to the localizer $W_S^{\mathrm{d}}$, which on $\Fib S$, and even on $\Liss\,S$, coincides with $W_S^{\mathrm{f}}$. But here we run up against the difficulty that the inverse image functor
\[f^* : \Cat/S \longrightarrow \Cat/S', \qquad\quad X \longmapsto X \times_S S'/S',\]
is \textit{not} compatible in general (if $S$ is not isomorphic to the one-point category, for arbitrary $S'$) with the localizers $W_S^{\mathrm{d}}$, $W_{S'}^{\mathrm{d}}$. We see this e.g. by taking $S' = \{s\}$, and recalling that $W_S^{\mathrm{d}}$ does \textit{not} imply, on the whole of Cat$/S$, fiberwise $W$-equivalence.

These troubles do not occur on $\Liss\,S$, where $W_S^{\mathrm{d}}$ and $W_S^{\mathrm{f}}$ coincide. We can therefore describe $f^*$ on the HOT also with the help of % [sic] sûr -> sur, accent fautif
\[f^* : \Liss\,S \longrightarrow \Liss\,S', \qquad\quad X \longmapsto X \times_S S'/S',\]
by passing to the localized categories
\[f^* : (\Liss\,S)(W_S^{\mathrm{f}})^{-1} \longrightarrow (\Liss\,S')(W_{S'}^{\mathrm{f}})^{-1}.\]
In other words:

\textbf{Proposition 1.} \textit{If $X$ in} Cat$/S$ \textit{is smooth over $S$, then we have a canonical isomorphism}
\[f^*(\hot_{W,S}(X)) \simeq \hot_{W,S'}(X'), \qquad\quad \text{\textit{where} } X' = X \times_S S'/S'.\]
% foot: Version 9 novembre 2001


\endgroup
% END UNIT Der-ChVI-345/Der-ChVI-345_p071-075.tex
% BEGIN UNIT Der-ChVI-345/Der-ChVI-345_p076-080.tex
\begingroup\makeatletter%
\NoAutoSpacing%
\newcommand{\printedpage}[1]{\marginpar{\footnotesize p.~#1}}%
\newcommand{\msp}[1]{\marginpar{\footnotesize [#1]}}%
\newcommand{\und}[1]{\underline{\mathrm{#1}}}%
\newcommand{\HOT}{\mathrm{HOT}}%
\newcommand{\hot}{\mathrm{hot}}%
\newcommand{\Cat}{\mathrm{Cat}}%
\makeatother\let\printedpage\derpp\let\msp\dermsp%
\let\mspin\DERmspin%


\printedpage{74}
% head: A. Grothendieck : Dérivateurs, Ch. VI, Édition des Universités de Montpellier II et Paris VII
% foot: Version 9 novembre 2001

\msp{page 119}
A similar remark applies to $\und{Fibhot}\,S$, where $W_S^{\mathrm{d}} = W_S^{\mathrm{f}}$ as well. The base change functor induces
\[
\und{Fibhot}\,S \longrightarrow \und{Fibhot}\,S', \qquad X \longmapsto X \times_S S'/S',
\]
compatible with the relevant relative localizers (all those considered so far coincide on these categories, under the hypotheses made on $W$), whence
\[
f^* : (\und{Fibhot}\,S)W_S^{-1} \longrightarrow (\und{Fibhot}\,S')W_{S'}^{-1}.
\]
The localized categories this time are not the $\HOT_W$, but the $\HOT_W^{\mathrm{lc}}$, so we find a nice description of
\[
f^* : \HOT_W^{\mathrm{lc}}(S) \longrightarrow \HOT_W^{\mathrm{lc}}(S'),
\]
still more advantageous than in terms of $\und{Liss}_0S$, $\und{Liss}_0S'$, since applying to somewhat larger categories (NB $\und{Liss}_0 \subseteq \und{Fibhot}$). Let us therefore note the

\textbf{Corollary 1.} \textit{The formula of proposition 1 (p. 118) remains valid if $X$ is assumed} Hot-\textit{fibrant over $S$ (instead of smooth over $S$).} (One then finds relative homotopy types \textit{locally constant} over $S$, $S'$.)

We thus have a description of $f^*$ which applies to all the categories on the left-hand side
\msp{page 120}
of the diagram \texttt{[page]} 76, with the sole exception of $\Cat/S$ itself. If $X$ is in $\Cat/S$, without being either smooth or Hot-fibered over $S$, how to find the inverse image in $\HOT_W(S')$? We can use $\Phi$ to improve the situation! We know that $X$ defines the same class in $\HOT_W(S)$ as $\Phi(X)$, while $C(X)$ represents the image of this class under $i\rho$. These two objects are, the one smooth (and even split fibered), the other Hot-fibered (and even perfect) over $S$. Thus
\DERmark{ENV0351}


\textbf{Corollary 2.} \textit{If $X$ in} $\Cat/S$ \textit{is arbitrary, then}
\[
f^*(\hot_{W,S}(X)) \simeq \hot_{W,S'}(\Phi(X) \times_S S'),\qquad f^*i\rho(\hot_{W,S}(X)) \simeq \hot_{W,S'}(C(X) \times_S S').
\]
\DERmark{ENV0352}


\bigskip

One may expect, on the other hand, that for certain particular base changes, the formula of proposition 1 is valid without restriction on $X$ in $\Cat/S$. It must come to $\pm$ the same thing that the inverse image functor
\[
X \longmapsto X \times_S S'
\]
be compatible with the localizers $W_S^{\mathrm{d}}$, $W_{S'}^{\mathrm{d}}$. The condition of $W$-equivalence ``colocal over $S$\,'' must be viewed as a condition of a ``homological'' nature, for the derivator HOT. It is associated, for an arbitrary derivator $\mathbf{D}$, with the condition (for $f : X \longrightarrow Y$ in $\Cat/S$) that the $s\backslash X \longrightarrow s\backslash Y$ are $\mathbf{D}$-equivalences,
\printedpage{75}
% head: A. Grothendieck : Dérivateurs, Ch. VI, Édition des Universités de Montpellier II et Paris VII
% foot: Version 9 novembre 2001
\msp{page 121}
which is written here elegantly in terms of relative homology, by the condition that $\forall\ a \in \mathcal{A} = \mathbf{D}(e)$, the arrow induced by $f$ % [sic] induit -> induite, accord avec « la flèche »
\[
\underbrace{H_\bullet^{\mathbf{D}}(X/S, a_X)}_{\overset{\mathrm{def}}{=}(p_X)_!(a_X)} \longrightarrow \underbrace{H_\bullet^{\mathbf{D}}(Y/S, a_Y)}_{\overset{\mathrm{def}}{=}(p_Y)_!(a_Y)}
\]
\DERmark{ENV0353}\DERmark{ENV0354}

be \texttt{[an]} isomorphism. (Use the ``calculation'' of the relative homology sheaves $(p_X)_!(a_X)$, $(p_Y)_!(a_Y)$, by the characterization of the fibers in Der 4 bis). Now, the $\mathbf{D}$-\textit{proper} base changes $S' \longrightarrow S$ commute with the $p_!$, so that we find the

\textbf{Lemma.} \textit{The notion of} $\mathbf{D}$\textit{-colocal equivalence over $S$ in} $\Cat/S$ \textit{is stable under every base change $S' \longrightarrow S$ that is} $\mathbf{D}$\textit{-proper. (Whatever the derivator} $\mathbf{D}$\textit{.)}

It will be recalled that the $s\backslash S \longrightarrow S$ are precisely proper base changes! One % [sic] rappelera -> rappellera, orthographe; des base -> de base, accord
can say that $\mathbf{D}$-colocal equivalence can be described by requiring that $f : X \longrightarrow Y$ be transformed into $W$-equivalences by certain types of base change $S' \longrightarrow S$, choosing at will one of the four types
\DERmark{ENV0355}


\begin{itemize}
\item[a)] The base changes $s\backslash S \longrightarrow S$ (which gives back the definition we had adopted).
\item[b)] The base changes that are local co-isomorphisms.
\item[c)] Proper base changes.
\item[d)] $\mathbf{D}$-proper base changes.
\end{itemize}

We cannot yet apply these results to the derivator
\DERmark{ENV0356}

\msp{page 122}
HOT, for lack of having yet constructed it, and verified, the relevant axioms. We can simply state % [sic] verifié -> vérifié, accent
the

\textbf{Conjecture} (sure!) : \textit{Let $p : S' \longrightarrow S$ be a proper arrow (understood:} HOT\textit{-proper) in} Cat. \textit{Then the base change functor}
\[
\Cat/S \longrightarrow \Cat/S', \qquad X \longmapsto X \times_S S'/S'
\]
\textit{is compatible with the localizers $W_S^{\mathrm{d}}$, $W_{S'}^{\mathrm{d}}$. (For a smooth functor $S' \longrightarrow S$, base change would be compatible with $W_S^{\mathrm{g}}$, $W_{S'}^{\mathrm{g}}$.)} (\footnote[61]{NB This ``conjecture'' is proved without a lemma on page 168-169, without needing anything that is not % [sic] prouvé -> prouvée, accord avec « conjecture »
already known here. Cf. theorem 4, case b) p. 162.}).

But for every base change $S' \longrightarrow S$ (proper or not) which has the preceding property, % [sic] pout -> pour, faute de frappe
hence which induces a functor
\[
\underbrace{(\Cat/S)(W_S^{\mathrm{d}})^{-1}}_{\approx\,\HOT_W(S)} \longrightarrow \underbrace{(\Cat/S')(W_{S'}^{\mathrm{d}})^{-1}}_{\approx\,\HOT_W(S')}\ ,
\]
\printedpage{76}
% head: A. Grothendieck : Dérivateurs, Ch. VI, Édition des Universités de Montpellier II et Paris VII
% foot: Version 9 novembre 2001
it is immediate, using the comparison theorem (e.g. with $\und{Liss}\,S$, $\und{Liss}\,S'$) over $S$ and over $S'$, that this functor is none other (up to canonical isomorphism) than the inverse image functor $f^*$ for relative homotopy types. So let us state the % [sic] énoncons -> énonçons, cédille

\textbf{Corollary 3.} \textit{The formula of proposition 1 remains valid without any hypothesis whatever on $X$ in} $\Cat/S$\textit{, provided that the base change arrow $f : S' \longrightarrow S$ is such that the functor $X \longmapsto X \times_S S'$ from} $\Cat/S$ \textit{to} $\Cat/S'$
\msp{page 123}
\textit{takes colocal $W$-equivalences over $S$ to colocal $W$-equivalences over $S'$.} (Which should be the case if (and only if?) $f$ is $W$-proper, hence also if $f$ is (``absolutely'') proper (\footnote[62]{cf. I, pages 55-59.}).)

So far we have worked only with the localizers $W_S^{\mathrm{d}}$, $W_{S'}^{\mathrm{d}}$, to the exclusion of $W_S$, $W_{S'}$. (Except on the categories contained in $\und{Fibhot}\,S$, where the two coincide.) When one makes no draconian hypothesis on $p : S' \longrightarrow S$, of the Hot-fibration type, one has no chance of $X \longmapsto X \times_S S'$ taking coarse $W$-equivalences into likewise, hence no chance of catching the inverse image of the ``coarse'' relative homotopy type of $X$, by taking that of $X' = X \times_S S'$. This is tied to the fact that for a general $S' \longrightarrow S$, the localization functors
\[
\HOT_W(S) \xrightarrow{\rho_S} \HOT_W^{\mathrm{lc}}(S), \qquad \HOT_W(S') \xrightarrow{\rho_{S'}} \HOT_W^{\mathrm{lc}}(S')
\]
have no desire to commute with the inverse image functor $f^*$ on the $\HOT_W$. We may hazard the

\textbf{Proposition 2.} \textit{Equivalent conditions on the functor $S' \longrightarrow S$ :}

\begin{itemize}
\item[\textit{a)}] \textit{The diagram of functors}
\[
\begin{tikzcd}
\HOT_W(S) \arrow[r, "\rho_S"] \arrow[d, "f^*"] & \HOT_W^{\mathrm{lc}}(S) \arrow[d, "f^*_{\mathrm{lc}}"] \\
\HOT_W(S') \arrow[r, "\rho_{S'}"] & \HOT_W^{\mathrm{lc}}(S')
\end{tikzcd}
\]
\textit{commutes (up to isomorphism of composite functors).}

\msp{page 124}

\item[\textit{b)}] \textit{The base change functor} % [sic] changment -> changement, faute de frappe
\[
\und{Liss}\,S \longrightarrow \und{Liss}\,S', \qquad X \longmapsto X \times_S S'/S'
\]
\textit{takes $W_S$ into $W_{S'}$.}
\printedpage{77}
% head: A. Grothendieck : Dérivateurs, Ch. VI, Édition des Universités de Montpellier II et Paris VII
% foot: Version 9 novembre 2001
\item[\textit{b$'$)}] \textit{Analogous statement for} $\und{Fib}\,S \longrightarrow \und{Fib}\,S'$\textit{.}
\item[\textit{b$''$)}] \textit{Analogous statement for} $\und{Fibsc}\,S \longrightarrow \und{Fibsc}\,S'$\textit{.}
\end{itemize}

(\footnote[63]{Not proved, only that b, b$'$, b$''$ imply a). For the reverse implication, one would need to know that every arrow of $\mathcal{F}'$ which becomes \texttt{[an]} isomorphism in $\mathcal{F}'W_{S'}^{-1}$, is in $W_{S'}$ (strong saturation of $W_{S'}$ in $\mathcal{F}'$), something I do not know. But I believe this must be true.}, \footnote[64]{\textbf{NB} If we assume $p$ proper, we should also be able to take $\mathcal{F} = \Cat/S$.}).

\textsc{Proof.} If $\mathcal{F}$ is one of the three categories $\und{Liss}\,S$, $\und{Fib}\,S$ or $\und{Fibsc}\,S$, the functor $\rho_S$ can be described as the canonical localization functor
\[
\mathcal{F}(W_S^{\mathrm{d}})^{-1} \longrightarrow \mathcal{F}W_S^{-1}
\]
(that would work for any full subcategory $\mathcal{F}$ of $\Cat/S$ which contains $\und{Fib}\,S$). Now if b), b$'$) or b$''$) is satisfied, we find a commutative diagram of functors induced by base change
\[
\begin{tikzcd}
\mathcal{F}(W_S^{\mathrm{d}})^{-1} \arrow[r, "\rho_S"] \arrow[d, "f^*_{W_S^{\mathrm{d}}}"'] & \mathcal{F}(W_S)^{-1} \arrow[d, "f^*_{W_S}"] \\
\mathcal{F}'(W_{S'}^{\mathrm{d}})^{-1} \arrow[r, "\rho_{S'}"] & \mathcal{F}'(W_{S'})^{-1},
\end{tikzcd}
\]
with the proviso, it is true, in the case $\mathcal{F} = \Cat/S$, that the inverse image functor is also compatible with $W_S^{\mathrm{d}}$, $W_{S'}^{\mathrm{d}}$ (which is automatic if $\mathcal{F} = \und{Fib}\,S$ or $\und{Fibsc}\,S$). We know that $f^*_{W_S^{\mathrm{d}}}$ describes the functor $f^*$ on relative homotopy types, and to prove a) as a consequence
\msp{page 125}
of b) (strengthened), b$'$ or b$''$, it remains to make sure that the functor $f^*_{W_S}$ likewise describes the functor induced by the preceding one on locally constant relative homotopy types. Now we have the canonical inclusions $\mathcal{F}_0 \subseteq \mathcal{F}$, $\mathcal{F}_0' \subseteq \mathcal{F}'$, and on $\mathcal{F}_0$, $\mathcal{F}_0'$ \texttt{[we have]} $W_S = W_S^{\mathrm{d}} = W_S^{\mathrm{f}}$, and the categories are OK for describing $\HOT_W^{\mathrm{lc}}(S)$ and $\HOT_W^{\mathrm{lc}}(S')$, and the functor $f^*$ induced between them by inverse images (\footnote[65]{NB $(\Cat/S)_0 = \und{Fibhot}\,S$, for the other cases, it is clear.}) (\textbf{NB} The base change functor $X \longmapsto X \times_S S'$ is compatible with the localizers on these categories, for it is trivial in the form $W_S^{\mathrm{f}}$.) We then have, by inverse images $X \longmapsto X \times_S S'$, a commutative diagram
\[
\begin{tikzcd}
\mathcal{F}_0(W_S)^{-1} \arrow[r, "\approx"] \arrow[d, "(f^*_{W_S})_0"'] & \mathcal{F}(W_S)^{-1} \arrow[d, "f^*_{W_S}"] \\
\mathcal{F}_0'(W_{S'})^{-1} \arrow[r, "\approx"] & \mathcal{F}'(W_{S'})^{-1},
\end{tikzcd}
\]
\printedpage{78}
% head: A. Grothendieck : Dérivateurs, Ch. VI, Édition des Universités de Montpellier II et Paris VII
% foot: Version 9 novembre 2001
where the horizontal arrows are \textit{equivalences} of categories, by theorem-scholie 1 \textcircled{I}~. This proves that $f^*_{W_S}$ too is good for calculating the inverse images of relative homotopy types.

We have proved that b, b$'$, b$''$ each imply a), let us prove that a) implies them. Now a) says exactly that the inverse images $f^*$ of the elements of $\HOT_W^{\mathrm{lc}}(S)$
\msp{page 126}
are, up to isomorphism, obtained by the inverse image functor on the whole of $\mathcal{F}$, interpreting $\HOT_W^{\mathrm{lc}}(S)$ as $\mathcal{F}W_S^{-1}$. But this does indeed imply that $\mathcal{F} \longrightarrow \mathcal{F}'$ must be compatible with the localizers. But I'm getting carried away - one only finds that if $u : X \longrightarrow Y$ in $\mathcal{F}$ is in $W_S$, then the arrow $u' : X' \longrightarrow Y'$ gives an isomorphism in $\mathcal{F}'(W_{S'})^{-1}$. I do not know, on the other hand, whether this implies that $u'$ is in $W_{S'}$.

\textbf{Corollary 1.} \textit{Suppose that $p : S' \longrightarrow S$ is a} Hot\textit{-fibration. Then the inverse image functor $X \longmapsto X \times_S S'$ from} $\Cat/S$ \textit{to $\mathit{Cat}/S'$ maps $W_S$ into $W_{S'}$, hence the functor $p^* : \HOT_W(S) \longrightarrow \HOT_W(S')$ commutes with the localization functors $\rho_S$, $\rho_{S'}$ in these categories.} (\textbf{NB} One must make the appropriate hypothesis on $W$, namely L6 bis page 44. It is OK for $W_\infty$.)

The fact that $X \longmapsto X \times_S S'$ is compatible with $W_S$, $W_{S'}$ is a tautology, in view of hypothesis L6 bis. It does not seem that it is also compatible with $W_S^{\mathrm{d}}$, $W_{S'}^{\mathrm{d}}$ (I should make a counterexample, with of course $S'$ not proper over $S$). So we are under the conditions of b), which implies a).

\msp{page 127}

I do not feel capable of exhibiting a functor other than a Hot-fibration (or at least a $W$-fibration), for which the inverse image functor on $\HOT_W$ is compatible with the localization functors $\rho_S$, $\rho_{S'}$. Of course, if $X$ is arbitrary in $\Cat/S$, nothing prevents us from looking at its localized object for $W_S$, i.e. $\rho_S(\hot_{W,S}(X))$, and asking how to find its inverse image. If $p : S' \longrightarrow S$ is a Hot-fibration, then it suffices to take $\hot_{W,S'}(X')$, $X' = X \times_S S'$. In any other case (I mean, without a hypothesis on $X/S$, more % [sic] X' = S \times_S S' -> X' = X \times_S S', cf. la définition de X' plus haut
\DERmark{ENV0357}

precisely without assuming that $X$ is itself Hot-fibered over $S$), this formula is hopelessly botched. We must, to calculate $f^*\rho_S(\hot_{W,S}(X))$, ``calculate'' $\rho_S(\hot_{W,S}(X))$ by ``going up to the upper floor'', as $\hot_{W,S}(C(X))$. We then find the inverse image $\hot_{W,S'}(C(X) \times_S S')$ (by \texttt{[the]} corollary 1, p. 119).

Now that the contravariant law of $\HOT_W(S)$ in $S \in \mathrm{Ob}\,\Cat$ seems well understood, it is a matter of seeing that we do indeed have a derivator HOT. I do not expect the $\HOT_W^{\mathrm{lc}}(S)$
\msp{page 128}
to define a derivator as well, but I would have to make sure of this expressly. (It is for the existence of $f_!$, $f_*$ that I must look.) % [sic] je doit -> je dois, conjugaison

To tell the truth, one would have to, to do things properly, look also at the 2-\textit{functor} structure of
\[
S \longmapsto \HOT_W(S) \qquad \Cat^o \longrightarrow \Cat,
\]


\endgroup
% END UNIT Der-ChVI-345/Der-ChVI-345_p076-080.tex
% BEGIN UNIT Der-ChVI-345/Der-ChVI-345_p081-085.tex
\begingroup\makeatletter%
\NoAutoSpacing%
\newcommand{\printedpage}[1]{\marginpar{\footnotesize p.~#1}}%
\newcommand{\msp}[1]{\marginpar{\footnotesize [#1]}}%
\newcommand{\HOT}{\mathrm{HOT}}%
\newcommand{\Hot}{\mathrm{Hot}}%
\newcommand{\Cat}{\mathrm{Cat}}%
\newcommand{\Ob}{\mathrm{Ob}}%
\newcommand{\hot}{\mathrm{hot}}%
\newcommand{\uHom}{\underline{\mathrm{Hom}}}%
\newcommand{\uFib}{\underline{\mathrm{Fib}}}%
\newcommand{\uFibsc}{\underline{\mathrm{Fibsc}}}%
\newcommand{\uCh}{\underline{\mathrm{Ch}}}%
\newcommand{\circled}[1]{\tikz[baseline=(c.base)]\node[draw,rounded corners=4pt,inner sep=1.5pt](c){#1};}%
\makeatother\let\printedpage\derpp\let\msp\dermsp%
\let\mspin\DERmspin%


\printedpage{79}
% head: A. Grothendieck : Dérivateurs, Ch. VI, Édition des Universités de Montpellier II et Paris VII
% foot: Version 9 novembre 2001

so for $S$, $S'$ given in Cat, define
\[ \uHom(S',S)^o \longrightarrow \uHom(\HOT_W(S),\HOT_W(S')). \]
This comes down above all, for
\[ f,g : S' \rightrightarrows S \]
two arrows from $S'$ to $S$, and a homomorphism of functors
\[ u : f \longrightarrow g, \]
to describing a homomorphism of functors
\[ u^* : g^* \longrightarrow f^* \]
between $\HOT_W(S)$ and $\HOT_W(S')$. This is moreover already done, in the general setting of $(\mathcal{M},W)$, when one describes $\HOT_W(S)$ and $\HOT_W(S')$ by localization of $\mathrm{Hom}(S^o,\Cat)$ and $\mathrm{Hom}(S'^o,\Cat)$. These categories are interpreted respectively as $\uFibsc\,S$ and $\uFibsc\,S'$, hence as categories of categories over $S$ resp. $S'$, so one would like to interpret $u^*$ in terms of such categories, and while one is at it, for categories $X$ over $S$ as general as possible. We shall
\msp{page 129}
then expect to have to make assumptions, either on $X$, or on $f$ and on $g$, which ensure that the images of $\hot_{W,S}(X)$ under $f^*$ and $g^*$ are none other than $\hot_{W,S}(f^*X)$, $\hot_{W,S}(g^*X)$ respectively, where we put
\[ f^*(X) = X \times_S (S',f), \qquad g^*(X) = X \times_S (S',g). \]
So we shall probably have to place ourselves under one of the following assumptions:
\begin{enumerate}
\item[a)] $X$ is smooth over $S$.
\item[b)] $X$ is Hot-fibered over $S$.
\item[c)] $f$ and $g$ are proper (and subject to the validity of the conjecture on page 122).
\end{enumerate}
One then expects $u^*$ to be associated with a homomorphism
\[ u^* : g^*(X) \longrightarrow f^*(X) \qquad \textit{in } \Cat/S' \]
defined in terms of $u$. Is there a way of actually defining it, without assuming $X$ at least fibered over $S$? Let us take $S'=e$, so $f$, $g$ corresponding to objects $s,t\in\Ob\,S$, and $u$ to an arrow % [sic] fléche -> flèche, wrong accent
\[ u : s \longrightarrow t, \]
the inverse images of $X/S$ under $f$, $g$ are the fibers $X_s$, $X_t$, and we seek
\[ u^* : X_t \longrightarrow X_s. \]
We have no chance of finding one canonically, if we do not assume $X$ fibered over $S$! So it is not a), b) or c) that can serve us, but on%
\printedpage{80}%
% head: A. Grothendieck : Dérivateurs, Ch. VI, Édition des Universités de Montpellier II et Paris VII
% foot: Version 9 novembre 2001
\msp{page 130}%
ly the hypothesis $X\in\Ob\,\uFib\,S$. Thus one can hope to find a description of $u^*$ between $\HOT_W(S)$ and $\HOT_W(S')$ in terms of operations on the categories over $S$, $S'$, as just suggested, only by restricting oneself to ``models'' $X$ which are in $\uFib\,S$, at least. And I do not doubt indeed that this works on $\uFib\,S$ no less well than directly on $\uFibsc\,S$, in terms of $\uHom(S^o,\Cat)$.

For the essential axioms for derivators, I refer to I p. 67-68 and p. 111-112. We have seen that Der 1, \circled{Der 2} (\footnote[66]{Botched for \underline{Der 2}, I do not know whether it is true even assuming $W$ \textit{strongly} saturated, i.e. such that $\hot_W(f)$ isomorphism $\Longrightarrow f\in W$.}) are satisfied for every prederivator defined by a pair $(\mathcal{M},W\subseteq\mathrm{Fl}\,\mathcal{M})$, so it is OK here for $\HOT_W$. It follows that it is OK also for $\HOT_W^{\mathrm{lc}}$, in view of the characterization of the $\HOT_W^{\mathrm{lc}}(S)$ as subcategories of $\HOT_W(S)$. For Der 1, does this work even for arbitrary sums of categories, i.e. do we have % [sic] quelconque -> quelconques, agreement with sommes; Canulé (footnote 66) is verified authorial wording; see corrections.csv
\[ \HOT_W(\coprod_i S_i) \xrightarrow{\sim} \prod_i \HOT_W(S_i)\,? \]
We have in any case
\[ \uFib(\overbrace{\coprod\limits_i S_i}^{S}) \xrightarrow{\sim} \prod_i \uFib(S_i), \]
and the localizer $W_S^{\mathrm{f}}$ in $\uFib(S)$ is
\msp{page 131}
indeed identified with the product of the localizers $W_{S_i}^{\mathrm{f}}$ in the $\uFib\,S_i$. But in the case of a perhaps infinite family of categories $\mathcal{F}_i$, equipped with localizers $W_i$, if we put
\[ \mathcal{F} = \prod\mathcal{F}_i, \qquad W = \prod W_i \subseteq \mathrm{Fl}(\mathcal{F}), \]
is the canonical functor
\[ \mathcal{F}W^{-1} \longrightarrow \prod_i(\mathcal{F}_iW_i^{-1}) \]
always an equivalence? Even in \texttt{[``Even in'' crossed out]} the case $I$ finite one needs (it seems to me) a rather careful argument, either by the universal property by induction on card\,$I$ (one is reduced to the case of two functors), or by dint of categories of paths, noting that
\[ \uCh_\infty^W(\mathcal{F};(s_i),(t_i)) \xrightarrow{\sim} \prod_i \uCh_\infty^{W_i}(\mathcal{F}_i;s_i,t_i), \tag*{$(*)$} \]
and passing to the $\pi_0$, which gives
\[ \mathrm{Hom}_W((s_i),(t_i)) \xrightarrow{\sim} \prod_i \mathrm{Hom}_{W_i}(s_i,t_i). \]
But formula $(*)$, visibly, only works for $I$ finite. In the general case, one finds the subcategory $\prod_i^!$ of $\prod_i$ formed by the systems $(c_i)_{i\in I}$ of paths which are of bounded length.

So the question of the validity of Der 1 without a finiteness assumption on the indexing set $I$
\printedpage{81}
% head: A. Grothendieck : Dérivateurs, Ch. VI, Édition des Universités de Montpellier II et Paris VII
% foot: Version 9 novembre 2001
\msp{page 132}
remains open for the moment. (I suspect that it is true, but one would need, to prove it, to have in hand a description of the arrows in $\HOT_W(S)$ by paths of bounded length, say of length two or three. This is an interesting question, but one that I reserve for later, if need forces me to it $\ldots\,$)

More important is the question of the validity of Der 3 and of its variants (I, p. 111-112). The minimum one asks of a derivator $\mathbf{D}$ is that if $S'\longrightarrow S$ in Diag is a Hot-equivalence, it should be a $\mathbf{D}$-equivalence, i.e. (in the present case) that the functor
\[ f^* : \HOT_W(S) \longrightarrow \HOT_W(S') \]
be fully faithful on the subcategory of \textit{constant} objects in $\HOT_W(S)$. This condition becomes more tangible when one already has either the cohomology functors $H^\bullet_{\mathbf{D}}(S,-)$, $H^\bullet_{\mathbf{D}}(S',-)$, or the homology $H_\bullet^{\mathbf{D}}(S,-)$, $H_\bullet^{\mathbf{D}}(S',-)$. In terms of the latter, it means that
\[ p_!p^*(a) \xrightarrow{\sim} p'_!p'^*(a) \qquad \text{\texttt{[is an]} isomorphism,} \]
\msp{page 133}
for every $a$ in $\mathcal{A}=\mathbf{D}(e)$. (Where $p:S\longrightarrow e$, $p':S'\longrightarrow e$ are the structural morphisms to the one-point category.) We have seen the general heuristic formula
\[ p_!\underbrace{p^*(a)}_{a_S} = \hot(S)\times a, \]
where $\hot\,S$ is the class of $S\in\Ob\,\Cat$ in $\Hot=(\Cat)W_\infty^{-1}$, and where the product $\times$ is an external operation of Hot on $\mathcal{A}=\mathbf{D}(e)$
\[ \Hot\times\mathcal{A} \longrightarrow \mathcal{A}\ . \]
In the present case, $\mathcal{A}=\Hot$, and the operation
\[ \Hot\times\Hot \longrightarrow \Hot \]
should be none other than the cartesian product functor, which is deduced from the similar functor in Cat. When one has this formalism, the desired result for an $f:S'\longrightarrow S$ simply reduces to saying that if $f$ is in $W$, then so is $f\times\mathrm{id}_X$ for every $X$ in Cat (taking $X$ such that $\hot(X)\simeq a$). Now here this is a well-known fact, resulting from the axioms imposed on $W$.

So I have the impression that I have no chance of proving Der 3a$'$ (p. 111), nor even Der 3a) (the case where $\mathrm{S}=\mathrm{e}$), without having a mastery of the formalism of the $p_!$. And I have
\printedpage{82}
% head: A. Grothendieck : Dérivateurs, Ch. VI, Édition des Universités de Montpellier II et Paris VII
% foot: Version 9 novembre 2001
\msp{page 134}
good hope that this mastery will be obtained at little cost, whereas the question of the $p_*$ seems much more delicate. % [sic] cetter -> cette, misspelling

I am therefore going to turn first to the axiom Der 4$'$ of existence of
\[ f_! : \HOT_W(S') \longrightarrow \HOT_W(S), \]
as a left adjoint to $f^*$. I am going to try to construct ad-hoc a functor $f_!$, then the two adjunction arrows
\[ \begin{aligned} \mathrm{id} &\quad\longrightarrow\quad f^*f_! \\ \mathrm{id} &\quad\longleftarrow\quad f_!f^*, \end{aligned} \]
and finally, to check the two relevant compatibilities.

The ``obvious'' geometric idea is, for $X'\in\Cat/S'$, representing $\xi'\in\Ob\,\HOT_W(S')$, to define $f_!(\xi')$ by
\[ f_!(\hot_{S',W}(X')) = \hot_{S,W}(X'), \]
where on the right-hand side, $X'$ is regarded as an object of $\Cat/S$, thanks to $f:S'\longrightarrow S$ and to the functor
\[ f_! : \Cat/S' \longrightarrow \Cat/S \]
which it induces. The only problem is to see whether this functor is compatible with the
\msp{page 135}
localizers $W_{S'}^{\mathrm{d}}$, $W_S^{\mathrm{d}}$, and if not, whether it is so at least on the subcategory $\uFib\,S'$ or $\uFibsc\,S'$, which suffices to compute $\HOT_W(S')$ by localization. We must therefore see whether an $X'_1\longrightarrow X'_2$ in $\Cat/S'$, which is in $W_{S'}^{\mathrm{d}}$, i.e. which is in $W$ and remains so under every base change of the form $s'\backslash S'\longrightarrow S'$, more generally for every local co-isomorphism, is also in $W_S^{\mathrm{d}}$, i.e. whether base changes of the form $s\backslash S\longrightarrow S$ turn it into $W$-equivalences (coarse). Now local co-isomorphisms are stable under \textit{every base change} (cf. I, p. 55 in the dual case of local isomorphisms), so we win!
\[
\begin{tikzcd}
\overline{S'} = s\backslash S' \arrow[r] \arrow[d] & S' \arrow[d] \\
\overline{S} = s\backslash S \arrow[r] & S
\end{tikzcd}
\quad .
\]
I am going to check it again all the same. If
\[
\begin{tikzcd}
\overline{S'} \arrow[r] \arrow[d] & S' \arrow[d] \\
\overline{S} \arrow[r] & S
\end{tikzcd}
\]
\printedpage{83}
% head: A. Grothendieck : Dérivateurs, Ch. VI, Édition des Universités de Montpellier II et Paris VII
% foot: Version 9 novembre 2001
is a cartesian square in Cat, we said in loc. cit., essentially, that for $\overline{s'}$ in $\overline{S'}$, having images $s'$, $\overline{s}$, $s$ in $S'$, $\overline{S}$, $S$ respectively, the square of the colocalized categories
\[
\begin{tikzcd}
\overline{s'}\backslash\overline{S'} \arrow[r] \arrow[d] & s'\backslash S' \arrow[d] \\
\overline{s}\backslash\overline{S} \arrow[r] & s\backslash S
\end{tikzcd}
\]
is again
\msp{page 136}
cartesian. I have just rechecked it again this very instant, for the objects at least. To say that $\overline{S}\longrightarrow S$ is a local co-isomorphism, means that the $\overline{s}\backslash\overline{S}\longrightarrow s\backslash S$ are \texttt{[some]} isomorphisms. But the preceding cartesian square shows that then the $\overline{s'}\backslash\overline{S'}\longrightarrow s'\backslash S'$ are also \texttt{[some]} isomorphisms, we win! % [sic] le de revérifier -> de le revérifier, word order
\DERmark{ENV0358}


\textbf{Theorem 1.} \textit{For every arrow $f:S'\longrightarrow S$ in} Cat\textit{, the induced functor}
\[ f_! : \Cat/S' \longrightarrow \Cat/S \]
\textit{is compatible with the localizers $(W_{S'}^{\mathrm{d}},W_S^{\mathrm{d}})$, hence dually with the localizers $(W_{S'}^{\mathrm{g}},W_S^{\mathrm{g}})$, and trivially with the localizers $(W_{S'},W_S)$. It therefore induces a commutative diagram of functors on the localized categories}
\[
\begin{tikzcd}
\HOT_W(S') \arrow[r, "\rho_{S'}"] \arrow[d, "f_!"] & \HOT_W^{\mathrm{lc}}(S') \arrow[d, "f_!^{\mathrm{lc}}"] & \HOT_W^o(S') \arrow[l, "\sigma_{S'}"'] \arrow[d, "f_!^o"'] \\
\HOT_W(S) \arrow[r, "\rho_S"] & \HOT_W^{\mathrm{lc}}(S) & \HOT_W^o(S). \arrow[l, "\sigma_S"']
\end{tikzcd}
\]

\textbf{Definition 1.} \textit{We have thus defined the vertical functors $f_!$, $f_!^{\mathrm{lc}}$, $f_!^o=(f^o)_!$, called homological direct image functors, or integration functors for $f$ (and relative homotopy types).}

\msp{page 137}
We have obtained at the same time, for free, the compatibility of these functors with the functors \texttt{[of]} localization $\rho_{S'}$, $\sigma_{S'}$ for $\HOT(S')$, and $\rho_S$, $\sigma_S$ for $\HOT(S)$. On the other hand, one will not expect them to commute with the fully faithful inclusion functors of the $\HOT_W^{\mathrm{lc}}$ into $\HOT_W$. Indeed, if one does not assume $f$ in \underline{Fibhot} (\footnote[67]{This case must be looked at carefully later.}), it is quite immoral to expect that $f_!$ in $\HOT_W(S')$ should transform locally constant types into likewise. The situation is therefore exactly the reverse of that for the functor $f^*$, which was compatible with the inclusion functors $i_S$, $i_{S'}$ and $j_S$, $j_{S'}$ - this was so trivial that there was no question even of saying it, let alone making a proposition of it! The situation is


\endgroup
% END UNIT Der-ChVI-345/Der-ChVI-345_p081-085.tex
% BEGIN UNIT Der-ChVI-345/Der-ChVI-345_p086-090.tex
\begingroup\makeatletter%
\newcommand{\printedpage}[1]{\marginpar{\footnotesize p.~#1}}%
\newcommand{\msp}[1]{\marginpar{\footnotesize [#1]}}%
\newcommand{\HOT}{\mathrm{HOT}}%
\newcommand{\hot}{\mathrm{hot}}%
\newcommand{\Cat}{\mathrm{Cat}}%
\newcommand{\id}{\mathrm{id}}%
\newcommand{\pr}{\mathrm{pr}}%
\newcommand{\Ob}{\mathrm{Ob}}%
\newcommand{\Wd}{W^{\mathrm{d}}}%
\makeatother\let\printedpage\derpp\let\msp\dermsp%
\let\mspin\DERmspin%


\printedpage{84}
% head: A. Grothendieck : Dérivateurs, Ch. VI, Édition des Universités de Montpellier II et Paris VII
even better still for the functor $f_!$ between the HOT, since the description of this functor can be given directly on the whole of $\Cat/S'$, without having to restrict oneself to more or less \msp{page 138}clever subcategories, where the induced functor would be compatible with the localizers ($\underline{\mathrm{Liss}}\,S$ and $\underline{\mathrm{Fibhot}}\,S$ in the case of $f^*$).

We have immediately the transitivity of $f_!$, $f_!^{\mathrm{lc}}$, $f_!^{o}$, i.e.
\[
(gf)_! = g_!f_!, \qquad (gf)_!^{\mathrm{lc}} = g_!^{\mathrm{lc}}f_!^{\mathrm{lc}}, \qquad (gf)_!^{o} = g_!^{o}f_!^{o}.
\]

But we are not going to move $S$, $S'$, $f$ for the moment, for we must now, before anything else, try to establish the adjunction of $f_!$ and $f^*$ (in that order).

For the moment, let us deal with $\HOT_W$ only - the case of $\HOT_W^{o}$ will follow from this by applying what we shall know to $f^{o}$. And we shall deal with $\HOT_W^{\mathrm{lc}}$ later.

1) \textit{The homomorphism}
\[
\alpha : \id_{\HOT_W(S')} \longrightarrow f^*f_!.
\]

If $X'$ is in $\Cat/S'$, $f_!(\hot_{S',W}(X'))$ is $\hot_{S,W}(X')$, but we cannot find its inverse image $f^*$ over $S'$ by taking the categorical inverse image. For even if we make draconian hypotheses on $X'$ as an object of $\Cat/S'$ (e.g. $X'$ fibered over $S'$), these will be reduced to nothing when we have composed with $f$, on which we do not want to make any hypothesis here. We \msp{page 139}must therefore first apply the functor $\Phi_S$ to $X'$ over $S$, then take its inverse image over $S'$:
\[
f^*f_!(\hot_{S',W}(X')) = \hot_{S',W}(\Phi_S(X'/S)\times_S S').
\]
We must therefore find an arrow
\[
\alpha_{X'} : X' \longrightarrow \Phi_S(X'/S)\times_S S'
\]
in the localized $(\Cat/S')(\Wd_{S'})^{-1}$, or even an arrow in $\Cat/S'$ itself. But by the adjunction properties of the functors $f_!$ and $f^*$ at the \textit{categorical level,} such an arrow amounts to an arrow
\[
f_!(X') \longrightarrow \Phi_S(X'/S)
\]
in $\Cat/S$, i.e. to an arrow which makes the square commute
\[
\begin{tikzcd}
X' \arrow[r] \arrow[d] & \Phi_S(X'/S) \arrow[d] \\ % [sic] \Phi_S(X'/S') -> \Phi_S(X'/S), the arrow f_!(X') -> \Phi_S(X'/S) just above has X'/S
S' \arrow[r] & S.
\end{tikzcd}
\]
\DERmark{ENV0359}

% foot: Version 9 novembre 2001

\printedpage{85}
% head: A. Grothendieck : Dérivateurs, Ch. VI, Édition des Universités de Montpellier II et Paris VII
Now we have such a canonical arrow, namely $i_{X'/S}$. It is functorial in $X'$ over $S$ (without even having to worry whether $X'$ over $S$ comes from an $X'$ over $S'$, i.e. whether the structural arrow $X'\longrightarrow S$ factors through $S'\longrightarrow S$).

\leavevmode\msp{page 140}2) \textit{The homomorphism}
\[
\beta : \id_{\HOT_W(S)} \longleftarrow f_!f^*.
\]

We can describe $\HOT_W(S)$ by $(\underline{\mathrm{Fib}}\,S)(\Wd_S)^{-1}$ or $(\underline{\mathrm{Liss}}\,S)(\Wd_S)^{-1}$ for example, so that $f^*$ has a simple description by
\[
f^*(\hot_{S,W}(X)) = \hot_{S',W}(X\times_S S')
\]
(for $X$ smooth over $S$). We then have
\[
f_!f^*(\hot_{S,W}(X)) = \hot_{S,W}(X\times_S S'),
\]
but watch out, $X\times_S S'$ is not a priori in $\underline{\mathrm{Liss}}\,S$, only in $\Cat/S$ and nothing more. But that does not prevent it from giving what is needed, by localization $(\Wd_S)^{-1}$ in $\Cat/S$. To define $\beta_X$, it therefore suffices to find an arrow
\[
\beta_X : X \longleftarrow X\times_S S'
\]
in $\Cat/S$, and we do see $\pr_1^X$ (\footnote[68]{when $\xi\in\HOT_W(S)$ is defined by $X$ in $\Cat/S$ not necessarily smooth, we must find $\Phi_S(X)\times_S S'\longrightarrow X$ in $(\Cat/S)(\Wd_S)^{-1}$, we take $(i_X)^{-1}\pr_1^{\Phi_S(X)}$ (cf. page 143).}). We thus have
\[
\left\{
\begin{array}{lcl}
\alpha : \id & \longrightarrow & f^*f_! \\
\beta : \id & \longleftarrow & f_!f^*
\end{array}
\right.
\]

3) \textit{The compatibility}
\[
\begin{tikzcd}
f_! \arrow[r, "f_!*\alpha"] \arrow[rr, bend right, "\id"'] & f_!f^*f_! \arrow[r, "\beta*f_!"] & f_!.
\end{tikzcd}
\]

We must therefore take $X'$ over $S'$, and make the arrows explicit starting from the object
\[
\xi' = \hot_{S',W}(X').
\]
% foot: Version 9 novembre 2001

\printedpage{86}
% head: A. Grothendieck : Dérivateurs, Ch. VI, Édition des Universités de Montpellier II et Paris VII
\leavevmode\msp{page 141}We have
\[
f_!(\xi') = \hot_{S,W}(X'/S),
\]
and to calculate $f^*f_!(\xi')$, we must as before take $\Phi_S(X')$ and we find
\[
f^*f_!(\xi') = \hot_{S',W}(\Phi_S(X')\times_S S').
\]
Finally,
\[
f_!f^*f_!(\xi') = \hot_{S,W}(\Phi_S(X')\times_S S').
\]
We must make the arrows explicit
\[
\begin{tikzcd}[row sep=small]
\underbrace{f_!(\xi')} \arrow[r, "{f_!(\alpha_{\xi'})}"] & \underbrace{f_!f^*f_!(\xi')} \arrow[r, "{\beta_{f_!(\xi')}}"] & \underbrace{f_!(\xi')}. \\
X'/S \arrow[r, "{\alpha_{X'}}", "{\text{deduced from } i_{X'}}"'] & \Phi_S(X')\times_S S' \arrow[r, dotted, "{\beta_{X'}}"] & X'/S \\
 & \text{over } S &
\end{tikzcd}
\tag*{$(*)$}
\]

$\alpha_{X'}$ is deduced from $i_{X'} : X'\longrightarrow\Phi_S(X')$ as the unique $S'$-morphism making the diagram commute
\[
\begin{tikzcd}
X' \arrow[r, "{\alpha_{X'}}"] \arrow[d] \arrow[rr, bend left, "{i_{X'}}"] & \Phi_S(X')\times_S S' \arrow[r, "\pr_1"] \arrow[dl, dotted] & \Phi_S(X') \arrow[d] \\ % Reading verified against PDF leaf 48: the curved i_X' arrow is solid.
S' \arrow[rr] & & S.
\end{tikzcd}
\tag*{$(**)$}
\]

To describe $\beta_{f_!(\xi')}$, where $f_!(\xi') = \hot_{S,W}(X')$, we must give a general description of $\beta_\xi$ that is valid for \textit{every} representative $X$ of $\xi$, not only in the case $X$ smooth over $S$. To calculate $f^*(\xi)$, we must therefore take the $\hot_{S',W}$ of $\Phi(X)\times_S S'$, then for $f_!f^*(\xi)$, we \msp{page 142}must take the $\hot_{S,W}$ of the same $\Phi_S(X)\times_S S'$. We must define, if we can, an arrow
\[
\Phi_S(X)\times_S S' \longrightarrow X \qquad\qquad \text{in } \Cat/S,
\]
or at least in $(\Cat/S)(\Wd_S)^{-1}$. But we have
\[
i_X : X \longrightarrow \Phi_S(X)
\]
which is in $\Wd_S$, and which by Cartesian product defines
\[
i_X\times_S\id_{S'} : X\times_S S' \longrightarrow \Phi_S(X)\times_S S'
\]
(in other words we have taken the inverse image of $i_X$ under $S'\longrightarrow S$). Is this arrow still in $\Wd_S$? (\footnote[69]{No, it is not, for if $S'=\{s\}$, this would mean that $X'_s\longrightarrow s\backslash X$ is in $W$, which is false in general.}). If so, we have the diagram
% foot: Version 9 novembre 2001
\printedpage{87}
% head: A. Grothendieck : Dérivateurs, Ch. VI, Édition des Universités de Montpellier II et Paris VII
\[
\begin{tikzcd}
\Phi_S(X)\times_S S' \arrow[r, dotted, "\beta_X ?"] & X \\
X\times_S S' \arrow[u, "{i_X\times_S\id_{S'}}", "{=\pi_X}"'] \arrow[ur, "\pr_1^X"] &
\end{tikzcd}
\]
and we define the desired arrow $\beta_X$ \textit{in} $(\Cat/S)(\Wd_S)^{-1}$ by the commutativity of this triangle, thus as
\[
\beta_X = \pr_1^X\circ\underbrace{\pi_X^{-1}}_{\text{inverse in }(\Cat/S)(\Wd_S)^{-1}} \qquad\qquad .
\]

\leavevmode\msp{page 143}But that doesn't work - cf. annotation! We shall then do the reverse
\[
\begin{tikzcd}
\Phi_S(X)\times_S S' \arrow[r, "\pr_1^{\Phi_S(X)}"] \arrow[dr, dotted, "\beta_X ?"'] & \Phi_S(X) \\
 & X \arrow[u, "{i_X \qquad \text{in } \Wd_S}"']
\end{tikzcd}
\]

(\footnote[70]{Description of $\beta_\xi$ for $\xi=\hot_{S,W}(X)$, with $X$ arbitrary in $\Cat/S$, to be made explicit in definition in due form.}), we define $\beta_X$ as $(i_X)^{-1}\pr_1^{\Phi_S(X)}$, where $(i_X)^{-1}$ is defined in $(\Cat/S)(\Wd_S)^{-1}$.

To prove the commutativity of $(*)$ (p. 141), since $i_{X'} : X'\longrightarrow\Phi_S(X')$ is in $\Wd_S$, it suffices to prove, instead of $\beta_{X'}\alpha_{X'} = \id_{X'}$, the relation deduced by composing with $i_{X'} : X'\longrightarrow\Phi_S(X')$
\[
\underbrace{i_{X'}\beta_{X'}}_{\pr_1^{\Phi_S(X')}}\alpha_{X'} = i_{X'},
\]
i.e.
\[
\pr_1^{\Phi_S(X')}\circ\alpha_{X'} = i_{X'},
\]
which is nothing other than the defining relation of $\alpha_{X'}$ (cf. $(**)$ p. 141), we win!

4) \textit{Compatibility}
\[
\begin{tikzcd}
f^* \arrow[r, "\alpha*f^*"] \arrow[rr, bend right, "\id"'] & f^*f_!f^* \arrow[r, "f^**\beta"] & f^*.
\end{tikzcd}
\]
% foot: Version 9 novembre 2001

\printedpage{88}
% head: A. Grothendieck : Dérivateurs, Ch. VI, Édition des Universités de Montpellier II et Paris VII
\leavevmode\msp{page 144}Let us therefore start with a $\xi\in\Ob\HOT_W(S)$, which we may suppose defined by $X\in\Ob\underline{\mathrm{Fib}}\,S$. Then
\[
\begin{aligned}
f^*(\xi) &\;=\; \hot_{S',W}(X\times_S S'),\\
f_!f^*(\xi) &\;=\; \hot_{S,W}(X\times_S S'),\\
f^*f_!f^*(\xi) &\;=\; \hot_{S',W}(\Phi_S(X\times_S S')\times_S S').
\end{aligned}
\]
We must now make explicit
\[
\underbrace{f^*(\xi)}_{\hot_{S',W}(X')} \xrightarrow{\alpha_{f^*(\xi)}} \underbrace{f^*f_!f^*(\xi)}_{\hot_{S',W}(\Phi_S(X')\times_S S')} \xrightarrow{f^*(\beta_\xi)} \underbrace{f^*(\xi)}_{\hot_{S',W}(X')} \quad ,
\]
and check that the composite is the identity. We have set
\[
X' = X\times_S S',
\]
and the diagram is therefore also
\[
\begin{tikzcd}
X' \arrow[r, "{\alpha_{X'}}"] \arrow[rr, bend left, "{\id_{X'}}"] & \Phi_S(X')\times_S S' \arrow[r, "f^*(\beta_\xi)"] & X' % Reading verified against PDF leaf 50: the curved identity arrow is solid.
\end{tikzcd}
\]
in $(\Cat/S')(\Wd_{S'})^{-1}$. To check commutativity, we can compose with
\[
X\times_S S' = X' \xrightarrow{i_X\times_S\id_{S'}} \Phi_S(X)\times_S S' \qquad\qquad \text{(where } i_X : X\longrightarrow\Phi_S(X)\text{)}
\]
\DERmark{ENV0360}

which is in $\Wd_{S'}$ (\footnote[71]{\textbf{NB} Since $X$ is smooth$/S$, and $\Phi_S(X)$ too, $X\longrightarrow\Phi_S(X)$ remains a colocal $W$-equivalence under any base change.}), now we check, from the general construction of $\beta_\xi=\beta_X$ (preceding page), that we have commutativity in
\[
\begin{tikzcd}
\Phi_S(X')\times_S S' \arrow[r, "f^*(\beta_\xi)"] \arrow[dr, "\Phi_S(\pr_1^X)\times_S S'"'] & X' \arrow[d, "i_X\times_S S'"] \\
 & \Phi_S(X)\times_S S'.
\end{tikzcd}
\]

\leavevmode\msp{page 145}We are reduced to the commutativity of
\[
\begin{tikzcd}
X\times_S S' = X' \arrow[r, "{\alpha_{X'}}"] \arrow[dr, "i_X\times_S S'"', "?"] & \Phi_S(X')\times_S S' \arrow[d, "\Phi_S(\pr_1^X)\times_S S'"] \\
 & \Phi_S(X)\times_S S'
\end{tikzcd}
\tag*{$(*)$}
\]
% foot: Version 9 novembre 2001


\endgroup
% END UNIT Der-ChVI-345/Der-ChVI-345_p086-090.tex
% BEGIN UNIT Der-ChVI-345/Der-ChVI-345_p091-095.tex
\begingroup\makeatletter%
\newcommand{\printedpage}[1]{\marginpar{\footnotesize p.~#1}}%
\newcommand{\msp}[1]{\marginpar{\footnotesize [#1]}}%
\newcommand{\Cat}{\mathrm{Cat}}%
\newcommand{\HOT}{\mathrm{HOT}}%
\newcommand{\hot}{\mathrm{hot}}%
\newcommand{\pr}{\mathrm{pr}}%
\newcommand{\id}{\mathrm{id}}%
\newcommand{\Ob}{\mathrm{Ob}\,}%
\newcommand{\Hom}{\mathrm{Hom}}%
\newcommand{\uHom}{\underline{\mathrm{Hom}}}%
\newcommand{\Liss}{\underline{\mathrm{Liss}}\,}%
\newcommand{\cA}{\mathcal{A}}%
\newcommand{\dd}{\mathrm{d}}%
\newcommand{\lc}{\mathrm{lc}}%
\newcommand{\adj}[2]{\mathrel{\underset{#2}{\overset{#1}{\rightleftarrows}}}}%
\makeatother\let\printedpage\derpp\let\msp\dermsp%
\let\mspin\DERmspin%


\printedpage{89}
% head: A. Grothendieck : Dérivateurs, Ch. VI, Édition des Universités de Montpellier II et Paris VII

in $(\Cat/S')(W^{\dd}_{S'})^{-1}$. Let us try to prove that this is commutative in $\Cat/S'$, not just in the category localized by $W^{\dd}_{S'}$. We can look at it as a diagram in $\Cat/S$ and prove its commutativity as such. This commutativity relation $vu = w$ then splits into two, according to the two factors $\Phi_S(X)$ and $S'$ of the target object $\Phi_S(X) \times_S S'$. For the factor $S'$, the relation simply means that $\alpha_{X'}$ is an $S'$-morphism, that's OK. For the factor $\Phi_S(X)$, we extend the diagram $(*)$ to the right by the square
\[
\begin{tikzcd}
X' \arrow[r, dotted, "\alpha_{X'}"] \arrow[dr, dotted] & \Phi_S(X') \times_S S' \arrow[r, "\pr_1^{\Phi_S(X')}"] \arrow[d, "\Phi_S(\pr_1 ^X)\times_S S'"] & \Phi_S(X') \arrow[d, "\Phi_S(\pr_1^X)"] \\
 & \Phi_S(X) \times_S S' \arrow[r, "\pr_1^{\Phi_S(X)}"] & \Phi_S(X),
\end{tikzcd}
\] % [sic] \Phi_S(\pr_1 X)\times_S S' -> \Phi_S(\pr_1^X)\times_S S', the X is printed on the line instead of as superscript of \pr_1
\DERmark{ENV0361}

\msp{page 146}
and taking into account that the composite of the first horizontal row is $i_{X'}$, by construction of $\alpha_{X'}$, we are reduced to the commutativity of
\[
\begin{tikzcd}
X \times_S S' = X' \arrow[r, "i_{X'}"] \arrow[d, "i_X\times_S S'"'] & \Phi_S(X') \arrow[d, "\Phi_S(\pr_1^X)"] \\
\Phi_S(X) \times_S S' \arrow[r, "\pr_1^{\Phi_S(X)}"] & \Phi_S(X).
\end{tikzcd}
\tag*{$(*)$}
\]

Now the morphism $i_X : X \longrightarrow \Phi_S(X)$ in $\Cat/S$ gives rise to the commutative square (as does any arrow $i : X \longrightarrow Y$)
\[
\begin{tikzcd}
X \times_S S' \arrow[r, "\pr_1^X"] \arrow[d, "i_X\times_S S'"'] & X \arrow[d, "i_X"] \\
\Phi_S(X) \times_S S' \arrow[r, "\pr_1^{\Phi_S(X)}"] & \Phi_S(X),
\end{tikzcd}
\]
so that the lower composite of $(*)$ is none other than that of $(**)$, hence equal to the upper composite of the latter, and the commutativity of $(*)$ comes down to that of
\[
\begin{tikzcd}
X \times_S S' = X' \arrow[r, "i_{X'}"] \arrow[d, "\pr_1^X"'] & \Phi_S(X') \arrow[d, "\Phi_S(\pr_1^X)"] \\
X \arrow[r, "i_X"] & \Phi_S(X),
\end{tikzcd}
\]
which is obvious by functoriality of $Z \xrightarrow{i_Z} \Phi_S(Z)$, applied to the arrow $\pr_1^X : X' \longrightarrow X$ of $\Cat/S$.

% foot: Version 9 novembre 2001

\printedpage{90}
% head: A. Grothendieck : Dérivateurs, Ch. VI, Édition des Universités de Montpellier II et Paris VII

\msp{page 147}
We have thus proved that the functors
\[
\HOT_W(S') \adj{f_!}{f^*} \HOT_W(S)
\tag*{$(*)$}
\]
are adjoint to each other, with adjunction morphisms
\[
\left\{
\begin{array}{lcl}
\id & \xrightarrow{\ \alpha\ } & f^*f_! \\[1ex]
\id & \xleftarrow{\ \beta\ } & f_!f^*
\end{array}
\right.
\]
specified above. I presume that the same proof, working with the coarse localizers $W_S$, $W_{S'}$, instead of $W^{\dd}_S$, $W^{\dd}_{S'}$, would give that the corresponding functors
\[
\HOT^{\lc}_W(S') \adj{f_!^{\lc}}{f^*_{\lc}} \HOT^{\lc}_W(S)
\tag*{$(**)$}
\]
are adjoint to each other, with essentially ``the same'' adjunction morphisms $\alpha$, $\beta$. I do not feel like going over the preceding arguments again and prefer to \textit{deduce} it from the adjunction in the situation $(*)$. This will give me the opportunity to better understand the essential formal features of the situation. I shall put for simplicity
\msp{page 148}
\[
\begin{array}{rcl@{\qquad\qquad}rcl}
\HOT_W(S) & = & \cA, & \HOT_W(S') & = & \cA', \\
\HOT^{\lc}_W(S) & = & \cA_0, & \HOT^{\lc}_W(S') & = & \cA'_0,
\end{array}
\]
and I use the pairs of adjoint functors
\[
\cA \adj{\rho}{i} \cA_0, \qquad\qquad \cA' \adj{\rho'}{i'} \cA'_0,
\]
where $\rho$, $\rho'$ are localization functors, i.e. $i$, $i'$ fully faithful functors. I have a pair of adjoint functors
\[
\cA' \adj{f_!}{f^*} \cA,
\]
and I assume moreover

\begin{enumerate}
\item[a)] $f_!$ is compatible with the localizers of $\cA'$, $\cA$ defined by $\rho'$, $\rho$ respectively, hence defines a functor (well determined up to unique isomorphism)
\[
f_!^0 : \cA'_0 \longrightarrow \cA_0,
\]
thus giving rise to a commutative diagram (up to functorial isomorphism)

% foot: Version 9 novembre 2001

\printedpage{91}
% head: A. Grothendieck : Dérivateurs, Ch. VI, Édition des Universités de Montpellier II et Paris VII
\[
\begin{tikzcd}
\cA' \arrow[r, "\rho'"] \arrow[d, "f_!"'] & \cA'_0 \arrow[d, "f_!^0"] \\
\cA \arrow[r, "\rho"] & \cA_0.
\end{tikzcd}
\tag*{(I)}
\]
\DERmark{ENV0362}

\texttt{[rather} $\cA \xrightarrow{\rho} \cA_0$\texttt{]}

\item[b)] ``Dually'', $f^* : \cA \longrightarrow \cA'$ is compatible with the full subcategories, essential images of the fully faithful functors $i$, $i'$, and then induces
\msp{page 149}
a functor, defined up to unique isomorphism:
\[
f_0^* : \cA_0 \longrightarrow \cA'_0,
\]
giving rise to a square of functors, commutative up to canonical functorial isomorphism
\[
\begin{tikzcd}
\cA_0 \arrow[r, "i"] \arrow[d, "f_0^*"'] & \cA \arrow[d, "f^*"] \\
\cA'_0 \arrow[r, "i'"] & \cA'.
\end{tikzcd}
\tag*{(II)}
\]
\end{enumerate}

\textit{Under these conditions, the pair of functors}
\[
\cA'_0 \adj{f_!^0}{f_0^*} \cA_0, \qquad \text{\texttt{[rather} } \cA'_0 \adj{f_!^0}{f_0^*} \cA_0, \text{\texttt{]}}
\]
\DERmark{ENV0363}

\textit{is a pair of adjoint functors.} (NB when I speak of pairs of adjoint functors, the one written above the other is the left adjoint.)

\textsc{Proof}. Let
\[
\eta_0 \in \Ob \cA_0, \qquad\qquad \xi'_0 \in \Ob \cA'_0,
\]
we must define a bifunctorial bijection
\[
\Hom_{\cA_0}(f_!^0(\xi'_0), \eta_0) \simeq \Hom_{\cA'_0}(\xi'_0, f_0^*(\eta_0))
\]
For this, we are going to write $\xi'_0$ in the form
\[
\xi'_0 = \rho'(\xi')\,, \qquad\qquad \xi' \in \Ob \cA',
\]
it will suffice to establish the bifunctorial bijection by replacing $\xi'_0$ by $\rho'(\xi')$, i.e.
\[
\Hom_{\cA_0}(f_!^0(\rho'(\xi')), \eta_0) \simeq \Hom_{\cA'_0}(\rho'(\xi'), f_0^*(\eta_0))
\]

% foot: Version 9 novembre 2001

\printedpage{92}
% head: A. Grothendieck : Dérivateurs, Ch. VI, Édition des Universités de Montpellier II et Paris VII

\msp{page 150}
(Indeed, the adjunction formula is interpreted as an isomorphism of bifunctors, or, what amounts to the same thing, of functors $\cA'_0 \mathrel{\underset{\psi}{\overset{\varphi}{\rightrightarrows}}} \cA_0^{\wedge}$, now
\[
\begin{array}{ccc}
\uHom(\cA'_0, \cA_0^{\wedge}) & \longrightarrow & \uHom(\cA', \cA_0^{\wedge}) \\[1ex]
\varphi & \longmapsto & \varphi \circ \rho'
\end{array}
\]
is fully faithful, since $\rho'$ is a localization functor, so it amounts to the same thing to give oneself an isomorphism $\varphi \simeq \psi$, or $\varphi \circ \rho' \simeq \psi \circ \rho'$.) We must therefore transform the expression
\[
\begin{array}{cccc}
 & \Hom_{\cA_0}(\underbrace{f_!^0(\rho'(\xi'))}_{\rho f_!(\xi')}, \eta_0) & \underbrace{\simeq}_{\text{\footnotesize commutativity of (I)}} & \Hom_{\cA_0}(\rho f_!(\xi'), \eta_0) \\[3ex]
\underbrace{\simeq}_{\text{\footnotesize adjunction $\rho$, $i$}} & \Hom_{\cA}(f_!(\xi'), i(\eta_0)) & \underbrace{\simeq}_{\text{\footnotesize adjunction $f_!$, $f^*$}} & \Hom_{\cA'}(\xi', \underbrace{f^*i(\eta_0)}_{i'(f_0^*(\eta_0))}) \\[3ex]
\underbrace{\simeq}_{\text{\footnotesize commutativity of (II)}} & \Hom_{\cA'}(\xi', i'(f_0^*(\eta_0))) & \underbrace{\simeq}_{\text{\footnotesize adjunction $\rho'$, $i'$}} & \Hom_{\cA'_0}(\underbrace{\rho'(\xi')}_{\xi'_0}, f_0^*(\eta_0)),
\end{array}
\]
which is the announced bijection.

I summarize the essentials of the results obtained on $f_!$, $f_!^{\lc}$:

\textbf{Theorem 2.} \textit{Let $f : S' \longrightarrow S$ be an arrow in} $\Cat$\textit{, let us consider the functors}
\[
\begin{aligned}
\HOT_W(S') &\adj{f_!}{f^*} \HOT_W(S) \\
\HOT^{\lc}_W(S') &\adj{f_!^{\lc}}{f^*_{\lc}} \HOT^{\lc}_W(S),
\end{aligned}
\]
\msp{page 151}
\textit{where the functors $f_!$, $f_!^{\lc}$ are described in theorem} 1 \textit{(p. 136). The adjunction morphisms}
\[
\left\{
\begin{array}{lcl}
\id & \longrightarrow & f^*f_! \\[1ex]
\id & \longleftarrow & f_!f^*
\end{array}
\right.
\]
\textit{are described on pp. 138 to 140.}

We would like to make explicit the nice formula, obtained in passing
\[
\boxed{\;f_!f^*(\xi) = \hot_{S,W}(S'/S) \times_S \xi\;}\ ,
\]
for this we must make the right-hand side explicit. Let us note the

\textbf{Proposition 4.} \textit{There exists a unique functor}
\[
(\xi, \eta) \longmapsto \xi \times_S \eta : \HOT_W(S) \times \HOT_W(S) \longrightarrow \HOT_W(S)
\]

% foot: Version 9 novembre 2001

\printedpage{93}
% head: A. Grothendieck : Dérivateurs, Ch. VI, Édition des Universités de Montpellier II et Paris VII

\textit{with the property that for $X$, $Y$ varying in} $\Liss S$\textit{, we have an isomorphism bifunctorial in $X$, $Y$}
\[
\hot_{S,W}(X) \times_S \hot_{S,W}(Y) \simeq \hot_{S,W}(X \times_S Y)
\tag*{$(*)$}
\]
(\footnote[72]{\textbf{NB} The multiplication $\times_S$ in $\HOT_W(S)$ is associative, commutative, unital \textit{etc., etc.}, since this is the case in $\Liss S$.}).

This is formal, using only that on $\Liss S$, the bifunctor
\[
(X, Y) \longrightarrow X \times_S Y
\]
is compatible with the localizers $W^{\dd}_S \times W^{\dd}_S$ and $W^{\dd}_S$. To see this, we are reduced to seeing that if $X \xrightarrow{u} X'$ in $\Liss S$ is in $W^{\dd}_S$, then the same holds for $X \times_S Y \xrightarrow{\ u\times_S \id_Y\ } X' \times_S Y$, for every
\msp{page 152}
$Y$ in $\Liss S$. In fact, we use only the smoothness of $X$, $X'$, not that of $Y$, and this is immediate, for if $X \longrightarrow X'$ is in $W^{\dd}_S$, $X \times_S Y \longrightarrow X' \times_S Y$ is in $W^{\dd}_Y$, hence also in $W^{\dd}_S$ (cf. theorem 1 applied to $Y \longrightarrow S$). This remark suggests at the same time:

\textbf{Corollary 1. (False)} \textit{The bifunctorial formula} $(*)$ \textit{remains valid for $X$ in} $\Liss S$\textit{, $Y$ arbitrary in} $\Cat/S$\textit{, and also in the reverse case.} (\footnote[73]{This statement is surely \textit{false!} Counterexample p. 153, cf. corollary 1 corrected on page 154-156.})

In the first case, let us consider
\[
Y \xrightarrow{i_Y} \Phi(Y)
\]
in $W^{\dd}_S$, hence an isomorphism in $(\Cat/S)(W^{\dd}_S)^{-1}$, we shall therefore have
\[
\begin{array}{lcl}
\hot_{S,W}(X) \times_S\ \hot_{S,W}(Y) & \xrightarrow{\ \sim\ } & \hot_{S,W}(X) \times_S\ \hot_{S,W}(\Phi(Y)) \\[1ex]
 & \underbrace{\xrightarrow{\ \sim\ }}_{\text{\footnotesize since $X$ \textit{and} $\Phi(Y)$ smooth over $S$}} & \hot_{S,W}(X \times_S \Phi(Y)),
\end{array}
\]
now I say that
\[
X \times_S Y \xrightarrow{\ X\times_S i_Y\ } X \times_S \Phi(Y)
\]
is again in $W^{\dd}_S$. It suffices to see that it is in $W^{\dd}_X$ (cf. theorem 1), and for this, that the smooth base change $X \longrightarrow S$ induces $\Cat/S \longrightarrow \Cat/X$ compatible with the localizers $W^{\dd}_S$, $W^{\dd}_X$. But this is \textit{false} (take $X = \{s\}$).

\msp{page 153}
Let us show that the corollary is false. I am going to take $Y = \{s\}$, the one-point category at $s$, then $\Phi(Y)$ is the fibered category with discrete fibers defined by the representable functor $F$ defined by $s$, so it is $S/F$ and, in this case $S/s$: % [sic] les corollaire -> le corollaire, agreement; occurence -> occurrence, spelling
\[
\Phi(\{s\}) \simeq S/s \qquad\quad \text{as a category over } S.
\]

% foot: Version 9 novembre 2001


\endgroup
% END UNIT Der-ChVI-345/Der-ChVI-345_p091-095.tex
% BEGIN UNIT Der-ChVI-345/Der-ChVI-345_p096-100.tex
\begingroup\makeatletter%
\newcommand{\printedpage}[1]{\marginpar{\footnotesize p.~#1}}%
\newcommand{\msp}[1]{\marginpar{\footnotesize [#1]}}%
\newcommand{\hot}{\mathrm{hot}}%
\newcommand{\HOT}{\mathrm{HOT}}%
\newcommand{\Cat}{\mathrm{Cat}}%
\newcommand{\Liss}{\underline{\mathrm{Liss}}}%
\newcommand{\multhot}{\mathrm{multhot}}%
\newcommand{\mult}{\mathrm{mult}}%
\makeatother\let\printedpage\derpp\let\msp\dermsp%
\let\mspin\DERmspin%


\printedpage{94}
% head: A. Grothendieck : Dérivateurs, Ch. VI, Édition des Universités de Montpellier II et Paris VII
If e.g. $s$ is the final element $e$ of $S$, we find
\[\Phi(\underset{S}{e}) \simeq S,\]
hence
\[\hot_{S,W}(X)\times_S \hot_{S,W}(e)=\hot_{S,W}(X),\]
on the other hand
\[\hot_{S,W}(X\times_S e)=\hot_{S,W}(X_e).\]

\msp{154}
If the two were isomorphic in $\HOT_W(S)$ by the canonical arrow
\[X_e\longrightarrow X,\]
applying the functor $p_!$ (where $p : S\longrightarrow e$), it would follow that $X_e\longrightarrow X$ is in $W$. If for example $S=\boldsymbol{\Delta}^1$, this would imply that for every $X$ smooth over $S$, $X_1\longrightarrow X$ is in $W$, whereas normally smoothness implies that $X_0\longrightarrow X$ is in $W$. So that would amount to saying that $\Liss\,\boldsymbol{\Delta}^1=\Liss_0\boldsymbol{\Delta}^1$, whence at once $\Liss\,S=\Liss_0S$, for every $S$ in Cat, we'd know about it! The most trivial counterexample is to take $X=\{0\}\hookrightarrow S=\boldsymbol{\Delta}^1$, the inclusion of the ``generic'' open, an open immersion, hence smooth, the fiber $X_1$ is empty! % [sic] s'ensuiverait -> s'ensuivrait, nonstandard verb form

% marginal mark: handwritten sign in the left margin beside the following statement
\textbf{Corollary 1 corrected.} \textit{If $X$ is in} $\Cat/S$\textit{, not necessarily smooth over $S$, then the functor}
\[\mult_{X/S} : Y\longmapsto Y\times_S X : \Cat/S\longrightarrow \Cat/S\]
\textit{induces a functor}
\[\Liss\,S\longrightarrow \Cat/S\]
\textit{compatible with the localizers $W_S^{\mathrm{d}}$, $W_S^{\mathrm{d}}$, hence induces a functor}
\[(\Liss\,S)(W_S^{\mathrm{d}})^{-1}\longrightarrow(\Cat/S)(W_S^{\mathrm{d}})^{-1},\]

\msp{155}
\textit{or again (taking into account the equivalence induced by the inclusion}
\[(\underline{\mathrm{Liss}}\,S)(W_S^{\mathrm{d}})^{-1}\xrightarrow{\ \sim\ }(\Cat/S)(W_S^{\mathrm{d}})^{-1}\quad\textit{),}\] % doubt resolved in DER_VI-c_corrections.csv; input note: underline of Liss -- double underline or single underline
\textit{a functor, again denoted} $\multhot_{X/S}$,
\[\multhot_{X/S} : \HOT_W(S)\longrightarrow\HOT_W(S).\]

This functor therefore satisfies the relation
\[\multhot_{X/S}(\hot_{S,W}(Y))\qquad\underbrace{\ \simeq\ }_{\text{canonical isomorphism}}\qquad\hot_{S,W}(X\times_S Y)\]
% foot: Version 9 novembre 2001
\printedpage{95}
% head: A. Grothendieck : Dérivateurs, Ch. VI, Édition des Universités de Montpellier II et Paris VII
% marginal mark: handwritten sign in the left margin beside the following two lines
for $Y$ \textit{smooth over $S$,} and for arbitrary $Y$ in $\Cat/S$,
\[\multhot_{X/S}(\hot_{S,W}(Y))\qquad\underbrace{\ \simeq\ }_{\text{canonical isomorphism}}\qquad\hot_{S,W}(X\times_S \Phi_S(Y)).\]

When $X$ is itself smooth over $S$, this functor is none other than multiplication by $\hot_{S,W}(X)$:
\[\multhot_{X/S}=\mult_{\hot_{S,W}(X)}\qquad\text{if $X$ \textit{smooth over $S$.}}\]

% marginal mark: handwritten sign in the left margin beside the following paragraph
But its value for $\hot_{S,W}(Y)$, when $Y$ is \textit{not} smooth over $S$, must not be confused with $\hot_{S,W}(X\times_S Y)$.

Thus, if $X$ and $Y$ are in $\Cat/S$, let us consider
\[X\times_S Y\qquad\text{and}\qquad\boxed{\hot_{S,W}(X\times_S Y).}\] % the frame is printed with rounded corners

\begin{enumerate}
\item[a)] If neither $X$ nor $Y$ is assumed smooth over $S$, this element ``has no natural interpretation in terms of operations on the categories $\HOT_W(S)$'' (\footnote[74]{Except in the case where one of the two is in $\underline{\mathrm{Fibhot}}\,S$, cf. annotation on the next page.}).

\msp{156}

\item[b)] If $X$ and $Y$ are smooth over $S$, this element is interpreted as
\[\hot_{S,W}(X)\times_S\hot_{S,W}(Y)\]
(\footnote[75]{But also if \textit{one} of the two elements is in $\underline{\mathrm{Fibhot}}\,S$ (and assuming the axiom L6 bis on $W$).}).

\item[c)] If one of the two $X$, $Y$ is smooth over $S$, without assuming anything about the other, we have an interpretation in terms of $\mult_X$ (if $Y$ smooth), or $\mult_Y$ (if $X$ smooth), in particular
\[\begin{aligned}
\hot_{S,W}(X\times_S Y)\ &\simeq\ \multhot_X(\hot_{S,W}(Y))\quad\text{if $Y$ smooth over $S$}\\
&\simeq\ \multhot_Y(\hot_{S,W}(X))\quad\text{if $X$ smooth over $S$.}
\end{aligned}\]
\end{enumerate}

\textbf{Corollary 2.} \textit{Let $f : S'\longrightarrow S$ be an arrow in} Cat\textit{, whence $f_!$, $f^*$ between} $\HOT_W(S')$ \textit{and} $\HOT_W(S)$\textit{. We have the relation}
\[f_!f^*(\xi)\simeq\multhot_{S'}(\xi),\]
\textit{hence if $S'$ is smooth over $S$,}
\[f_!f^*(\xi)\simeq\hot_{S,W}(S')\times_S\xi.\]

We are now in a position to verify the axioms Der 4$'$b, Der 5$'$ relating to direct images (I, p.~67, 68), and moreover Der 3b (I, p.~111). The axiom Der 4$'$ is none other than the existence of the left adjoints $f_!$ of the $f^*$. The axiom Der 5$'$ allows the ``calculation'' of the fibers of an object $f_!(\xi)$. It says that a certain
% foot: Version 9 novembre 2001
\printedpage{96}
% head: A. Grothendieck : Dérivateurs, Ch. VI, Édition des Universités de Montpellier II et Paris VII
\msp{157}
base change arrow, relative to the diagram in Cat equipped with the natural transformation $pf'\Longrightarrow fq$
\DERmark{ENV0364}

\[\begin{tikzcd}[row sep=large]
S' \arrow[d,"f"'] & s\backslash S' \arrow[l,"q"'] \arrow[d,"f'"] \\
S & e\simeq\{s\} \arrow[l,"p"']
\end{tikzcd}\qquad(s\in\mathrm{Ob}\,S),\]
namely
\[f'_!(q^*(\xi'))\longrightarrow p^*(f_!(\xi'))\]
\DERmark{ENV0365}

is an isomorphism (for $\xi'\in\mathbf{D}(S')$). So let
\[\xi'\in\mathrm{Ob}\,\HOT_W(S'),\qquad\xi'=\hot_{S',W}(X')\]
with $X'$ in $\Cat/S'$. Since
\[f_!(\xi')=\hot_{S,W}(X'),\]
and since $X'$ has no reason to be smooth over $S$ (if we do not assume $S'$ itself smooth over $S$, in which case it might be useful to take $X'$ smooth over $S'$, hence also over $S$), to calculate $p^*(f_!(\xi'))$ there is no question of taking the fiber $X'_s$, but we must take $s\backslash X'=\Phi_S(X')_s$:
\[p^*(f_!(\xi'))\simeq\hot_W(s\backslash X').\]
\[\begin{tikzcd}[row sep=large]
X' \arrow[d] & s\backslash X' \arrow[l] \arrow[d] \\
S' \arrow[d] & s\backslash S' \arrow[l,"q"'] \arrow[d] \\
S & s\backslash S \arrow[l]
\end{tikzcd}\]

Consider the base change diagram along $s\backslash S\longrightarrow S$ opposite \texttt{[i.e. above]}, with cartesian squares. We first see that $q^*(\xi')$ is none other than the class in $\HOT_W(s\backslash S')$ of the inverse image of $X'$ by the base change $s\backslash S'\longrightarrow S'$. To ensure this, it would suffice to take $X'$ smooth over $S'$ (which
\msp{158}
is permissible) - then the description of the inverse image is OK for every base change. But we have seen that in the case of a local co-isomorphism such as $s\backslash S'\longrightarrow S'$, there is not even any need for this artifice - the formula is OK without any smoothness restriction. Thus
\[q^*(\xi')=\hot_{s\backslash S',W}(s\backslash X'),\]
% foot: Version 9 novembre 2001
\printedpage{97}
% head: A. Grothendieck : Dérivateurs, Ch. VI, Édition des Universités de Montpellier II et Paris VII
whence
\[f'_!(q^*(\xi'))\simeq\hot_{e,W}(s\backslash X')=\hot_W(s\backslash X').\]
\DERmark{ENV0366}

So we have indeed found
\[p^*f_!\simeq f'_!q^*.\]

We can now also prove Der 3b. We have seen that once we have the left adjoints $f_!$ of the $f^*$, this axiom is stated in \textit{homological form} as follows: If
\[f : S'\longrightarrow S\]
is in $W$, we must prove that for every
\[\xi\in\HOT_W(e)\overset{\mathrm{def}}{=}\mathrm{Hot}_W,\]
\DERmark{ENV0367}

the corresponding arrow
\[(*)\qquad p'_!p'^*(\xi)\longrightarrow p_!p^*(\xi)\]
is \texttt{[an]} isomorphism in $\mathrm{Hot}_W$, where $p : S\longrightarrow e$, $p' : S'\longrightarrow e$ are the structural arrows. Now since every object of Cat is smooth over $e$, the preceding corollary 2 allows us to identify this arrow $(*)$ with
\[\eta'\times\xi\longrightarrow\eta\times\xi\]
\DERmark{ENV0368}

\msp{159}
deduced from
\[\hot_W(f) :\quad\underbrace{\ \eta'\ }_{\overset{\mathrm{def}}{=}\,\hot_W(S')}\quad\longrightarrow\quad\underbrace{\ \eta\ }_{\overset{\mathrm{def}}{=}\,\hot_W(S)}\quad.\]
\DERmark{ENV0369}\DERmark{ENV0370}

Since the latter is \texttt{[an]} isomorphism by hypothesis, we win.

So we have proved

\textbf{Theorem 3.} \textit{The prederivator} HOT \textit{satisfies the axioms} Der 1, Der 2 (\footnote[76]{\textit{not proved}, the ``result'' I had in mind is botched, cf. supplement proposition 5, page 164.}) \textit{(for the record), and} Der 3b (\footnote[77]{And even the strongest variant Der 3b$'_2$ of Der 3, cf. proof on the next page. \textbf{NB} For Der 3b$'_2$, we used the axiom L6 on $W$ - and this is essential. We see that the validity of Der 3b$'_2$ implies L6, and even the stronger form L6 bis, and even if we assume only Der 3b$'_1$, i.e. $f_0^*$ fully faithful.}), Der 4$'$, Der 5$'$\textit{.} % [sic] canulé (footnote 76) -> annulé (?), not a French word

We also feel like proving the strengthened variants of Der 3, the strongest form being this one (Der 3b$'_2$): if $f\in W$, then $f^*$ induces an \textit{equivalence} of categories
\[f_0^* : \HOT^{\mathrm{lc}}_W(S)\longrightarrow\HOT^{\mathrm{lc}}_W(S').\]

The full faithfulness of the functor $f_0^*$ must follow from the natural variant Der 3b$''_{\mathrm{coh}}$ of Der 3b, once we have developed the formalism of internal Hom in the $\HOT_W(S)$. But to give a meaning to Der 3b$''_{\mathrm{coh}}$, we need the functors \textit{cohomology,} i.e. $p_*$ and $p'_*$ (instead of $p_!$, $p'_!$) for HOT, so this is still premature for the moment. The construction of the $f_*$ and that of internal Hom will moreover surely be done $\pm$ jointly.
% foot: Version 9 novembre 2001

\printedpage{98}
% head: A. Grothendieck : Dérivateurs, Ch. VI, Édition des Universités de Montpellier II et Paris VII
\msp{160}
But I realize just now that we find almost ``by eye'' the fact that $f_0^*$ is indeed an \textit{equivalence} of categories, and that $f_!^0$ is the quasi-inverse equivalence, by using the adjunction of the functors $f_!^0$, $f_0^*$, and recalling that to prove that they are equivalences quasi-inverse to one another, it amounts to the same thing to show that the adjunction morphisms
\[\begin{aligned}
\mathrm{id}&\quad\xrightarrow{\ \alpha_0\ }\quad f_0^*f_!^0\\
\mathrm{id}&\quad\xleftarrow{\ \beta_0\ }\quad f_!^0f_0^*
\end{aligned}\]
are \texttt{[some]} isomorphisms. I am going to try to describe these adjunction arrows in the general case (without assuming $f\in W$), then to show that they are in $W$ if $f$ is.

\begin{enumerate}
\item[\texttt{[1)}] \texttt{Let]} $\xi=\hot^{\mathrm{lc}}_{S,W}(X)$ in $\HOT^{\mathrm{lc}}_W(S)$. We may assume $X\in\mathrm{Ob}\,\underline{\mathrm{Parf}}\,S$. Since \texttt{[``Since'' crossed out]} $X$ smooth over $S$, we have
\[f_0^*(\xi)=\hot^{\mathrm{lc}}_{S',W}(X\times_S S'),\qquad f_!^0f_0^*(\xi)=\hot^{\mathrm{lc}}_{S,W}(X\times_S S').\]
\DERmark{ENV0371}

So
\[\beta_{0,\xi} : f_!^0f_0^*(\xi)\longrightarrow\xi\]
is an arrow
\[X\times_S S'\longrightarrow X\qquad\text{in }(\Cat/S)(W_S)^{-1}\]
and I presume that it is $\mathrm{pr}_1$ (\footnote[78]{\textbf{NB} We see that the axiom L6 is $\pm$ \textit{necessary} for the result on $f_0^*$, $f_!^0$ to be valid.}). Is it in $W$ if $f : S'\longrightarrow S$ is? Yes (for $X$ is perfect

\msp{161}

over $S$), provided that $W$ satisfies the axiom L6 (cf. page 43).

\item[2)] Let
\[\xi'=\hot^{\mathrm{lc}}_{S',W}(X')\qquad\text{in }\HOT^{\mathrm{lc}}_W(S'),\]
\DERmark{ENV0372}

we have
\[f_!^0(\xi')=\hot^{\mathrm{lc}}_{S,W}(X'),\]
but to make $f_0^*$ of this element explicit, since $X'$ has no reason to be Hot-fibered over $S$ (in which case it would suffice to take $f^*(X')\overset{\mathrm{def}}{=}X'\times_S S'$), we must first take $C_S(X')$, then apply $f^*$:
\DERmark{ENV0373}\DERmark{ENV0374}

\[f_0^*f_!^0(X')=\hot^{\mathrm{lc}}_{S',W}(C_S(X')\times_S S').\]

We must therefore define
\[X'\xrightarrow{\alpha_{0,X'}}C_S(X')\times_S S'\qquad\text{in }(\Cat/S')W_{S'}^{-1}.\]
\DERmark{ENV0375}


Now we find a canonical arrow in $\Cat/S'$, deduced from the $S$-morphism $i_{X'}$: $X'\longrightarrow C_S(X')$, i.e. from the commutative square in solid lines
\end{enumerate}
% foot: Version 9 novembre 2001


\endgroup
% END UNIT Der-ChVI-345/Der-ChVI-345_p096-100.tex
% BEGIN UNIT Der-ChVI-345/Der-ChVI-345_p101-105.tex
\begingroup\makeatletter%
\newcommand{\printedpage}[1]{\marginpar{\footnotesize p.~#1}}%
\newcommand{\msp}[1]{\marginpar{\footnotesize [#1]}}%
\newcommand{\HOT}{\mathrm{HOT}}%
\newcommand{\hot}{\mathrm{hot}}%
\newcommand{\Hot}{\mathrm{Hot}}%
\newcommand{\Cat}{\mathrm{Cat}}%
\newcommand{\Hom}{\mathrm{Hom}}%
\newcommand{\pr}{\mathrm{pr}}%
\newcommand{\D}{\mathbf{D}}%
\newcommand{\bs}{\backslash}%
\makeatother\let\printedpage\derpp\let\msp\dermsp%
\let\mspin\DERmspin%


\noindent\printedpage{99}%
% head: A. Grothendieck : Dérivateurs, Ch. VI, Édition des Universités de Montpellier II et Paris VII
\[
\begin{tikzcd}
X' \arrow[rr, "i_{X'}"] \arrow[ddd, "p_{X'}"'] \arrow[dr, dotted, "\beta_{0,X'}"'] & & C_S(X') \arrow[ddd] \\
& C_S(X') \times_S S' \arrow[ur, dotted, "\pr_1"'] \arrow[ddl, dotted] & \\
& \text{\footnotesize commutative} & \\
S' \arrow[rr, "f"] & & S,
\end{tikzcd}
\]

\begin{quote}
as the unique dotted arrow $\beta_{0,X'}$ which is an $S'$-morphism (left triangle commutative) and which makes the upper triangle commute.
\end{quote}

If $f \in W$, then $\pr_1$, obtained from $f$ by base change $C_S(X') \longrightarrow S$, which is perfect, is in $W$ by L6, we know that the same holds for $i_{X'}$, hence also for $\beta_{0,X'}$, we win.

\msp{page 162}To finish with the general properties of the functors $f_!$ for $\HOT_W$, I would like to give the reasonable base change statements. This would be a

\noindent\textbf{Theorem 4.} \textit{Let}
\[
\begin{tikzcd}
X \arrow[d, "f"'] & X' \arrow[l, "q"'] \arrow[d, "f'"] \\
S & S' \arrow[l, "p"']
\end{tikzcd}
\]
\textit{be a cartesian diagram in} $\Cat$\textit{, let us consider the base change arrow (in} $\underline{\Hom}(\HOT_W(X), \HOT_W(S'))$\textit{):}
\[
p^* f_! \longleftarrow f'_! q^*.
\]
\textit{This is an isomorphism in each of the following two cases.}

\begin{enumerate}
\item[\textit{a)}] \textit{$f$ is smooth.}
\item[\textit{b)}] \textit{$p$ is proper.}
\end{enumerate}

\textsc{Proof.} \textit{Case a).} We test with $\xi$ in $\HOT_W(X)$, we may suppose $\xi$ to be of the form
\[
\xi = \hot_{X,W}(Z) \qquad \text{with } Z \in \mathrm{Ob}\, \underline{\mathrm{Liss}}\, X.
\]
Let us note that $Z$ is then also smooth over $S$, since $f$ is smooth.
\[
\begin{aligned}
f_!(\xi) &= \hot_{S,W}(Z), &\qquad p^* f_!(\xi) &= \hot_{S',W}(\quad \underbrace{Z \times_S S'}_{\text{$Z'$, cf. diagram}} \quad) \\
q^*(\xi) &= \hot_{X',W}(Z'), & f'_! q^*(\xi) &= \hot_{S',W}(Z'),
\end{aligned}
\]
OK!

\noindent\printedpage{100}%
% head: A. Grothendieck : Dérivateurs, Ch. VI, Édition des Universités de Montpellier II et Paris VII
\[
\begin{tikzcd}
Z \arrow[d] & Z' \arrow[l] \arrow[d] \\
X \arrow[d, "f"'] & X' \arrow[l, "q"] \arrow[d, "f'"] \\
S & S' \arrow[l, "p"]
\end{tikzcd}
\]
\textit{Case b).} We have to admit that \msp{page 163}the base change functor by $p$:
\[
Z \longmapsto Z \times_S S' : \Cat/S \longrightarrow \Cat/S'
\]
is compatible with the localizers $W^{\mathrm{d}}_{S}$, $W^{\mathrm{d}}_{S'}$ (cf. below) - that's, formally, the relevant property for $p$ (\footnote[79]{This must be \textit{equivalent} to $p$ proper. Compare with corollary 1, page 166.}). It will follow that % [sic] formellemet -> formellement, letter omitted
\[
p^* : \HOT_W(S) \longrightarrow \HOT_W(S')
\]
is computed, for an element in $\HOT_W(S)$ given by any $Z$ \textit{whatsoever} in $\Cat/S$ (not necessarily smooth over $S$), by
\[
p^*(\hot_{S,W}(Z)) = \hot_{S',W}(Z \times_S S').
\]
Let us then start, as in a), from an element $\xi$ of $\HOT_W(X)$ given by a $Z$ in $\Cat/X$ (no need this time to suppose $Z$ smooth), the preceding immediate calculation again establishes the commutation formula.

But it does not seem that the compatibility property of $X \longmapsto X \times_S S'$, with colocal equivalence over $S$ and over $S'$, follows in a $\pm$ obvious way from the condition that $p : S' \longrightarrow S$ be proper (\footnote[80]{But yes, cf. proof page 168. The developments of the four pages that follow are unnecessary by this proof.}). I feel like deducing it from general \textit{fourbi} on the formalism of derivators, namely from the base change theorem for the $f_!^{\D}$ for a $\D$-proper morphism.
\DERmark{ENV0376}


\msp{page 164}In the dual form of a base change theorem for $f_*^{\D}$, by $\D$-smooth base change, this is done in I, p. 56 to 58, and we should be able to apply it here. Let us note that in the proof of the theorem loc. cit., we use only L4 \texttt{[rather} Der 4\texttt{]}, i.e. the existence of the $f_*^{\D}$, to the exclusion of the $f_!^{\D}$ (L4$'$ \texttt{[rather} Der 4$'$\texttt{]}), plus L5 \texttt{[rather} Der 5\texttt{]} of course, so the dual form does not call on the existence of the $f_*^{\D}$, only of the $f_!^{\D}$. We thus have a good chance of being able to apply it to $\HOT_W$ (where I have not yet established the existence of the $f_*^{\mathrm{HOT}}$). % [sic] ne fait appel -> ne fait pas appel, "pas" apparently omitted (sense and "seulement des")

\noindent\printedpage{101}%
% head: A. Grothendieck : Dérivateurs, Ch. VI, Édition des Universités de Montpellier II et Paris VII
Let us first note

\noindent\textbf{Proposition 5.} \textit{Let} $f : S' \longrightarrow S$ \textit{be an arrow in} $\Cat$. \textit{Then $f$ is a} $\HOT_W$\textit{-equivalence} if and only if $\hot_W(f)$ \textit{is} \texttt{[an]} \textit{isomorphism. If $W$ is strongly saturated, this therefore means that $f \in W$ (and in any case $f \in W \Longrightarrow f$} \texttt{[is]} \textit{a} $\HOT_W$\textit{-equivalence)} (\footnote[81]{should be moved up toward theorem 3, p. 159.}).

This is contained in the proof of Der 3b earlier. I recall that $f$ is a HOT-equivalence, by definition, if and only if for $a, b \in \Hot_W (= \HOT_W(e))$,
\[
\Hom_{\HOT_W(S)}(a_S, b_S) \longrightarrow \Hom_{\HOT_W(S')}(a_{S'}, b_{S'}) \tag*{$(*)$}
\]
is \texttt{[an]} isomorphism, where $a_S$ etc. denotes $p_S^*(a)$ etc. ($p_S : S \longrightarrow e$). The adjunction formulas for $p_{S!}$, $p_S^*$ and for $p_{S'!}$, $p_{S'}^*$ allow us to interpret the arrow $(*)$ as
\[
\Hom_{\Hot_W}(p_{S!}(a_S),b)\longrightarrow\Hom_{\Hot_W}(p_{S'!}(a_{S'}), b),
\]
\DERmark{ENV0377}

\msp{page 165}deduced by ``transposition'' $\Hom(-, b)$ from the arrow
\[
p_{S'!}(a_{S'}) \longrightarrow p_{S!}(a_S)
\]
induced by $f$ - so we have to prove that this is an isomorphism (this is the form we had used earlier). Now this arrow is identified with
\[
\eta' \times a \longrightarrow \eta \times a
\]
induced by
\[
\hot_W(f) : \underbrace{\eta'}_{= \hot_W(S')} \quad \longrightarrow \quad \underbrace{\eta}_{= \hot_W(S)} \quad ,
\]
and to say that this is \texttt{[an]} isomorphism for any $a$ in $\Hot_W$ amounts to saying that $\eta' \longrightarrow \eta$ is (since multiplication in $\Hot_W$ has a unit element \dots). % [sic] une élément -> un élément, gender agreement

Now I remember that to have axiom Der 2 of derivators, I should also suppose $W$ \textit{strongly} saturated (i.e. equal to the localizer defined by the canonical functor $\Cat \longrightarrow (\Cat)W^{-1} \overset{\mathrm{def}}{=} \Hot_W$). (This is what I assumed from the outset, starting from a general $(\mathcal{M}, W)$ to define a prederivator, to have axiom Der 2.)
\DERmark{ENV0378}


Applying the dual of the theorem in I, p. 56, we thus find the base change theorem for $f_!$, for base change $f : S' \longrightarrow S$ \textit{$W$-proper.}

% the following paragraph is marked in the left margin by a double vertical rule
\msp{page 166}\marginpar{$\Vert$}But the proof surely uses Der 2 (fiberwise isomorphism criterion, in $\HOT_W(S)$), which I have not proved yet. It is incomplete for the moment. But I can already note the

\noindent\printedpage{102}%
% head: A. Grothendieck : Dérivateurs, Ch. VI, Édition des Universités de Montpellier II et Paris VII
\textbf{Corollary 1} (of theorem 4, p. 162). \textit{If the arrow} $S' \xrightarrow{p} S$ \textit{in} $\Cat$ \textit{gives rise (for every} $f : X \longrightarrow S$ \textit{in} $\Cat/S$\textit{) to the base change formula for}
\[
f_! : \HOT_W(X) \longrightarrow \HOT_W(S),
\]
\textit{then the functor} $p^* : X \longmapsto X \times_S S'$ \textit{from} $\Cat/S$ \textit{to} $\Cat/S'$ \textit{is compatible with the localizers} $W^{\mathrm{d}}_{S}$\textit{,} $W^{\mathrm{d}}_{S'}$\textit{.} (The converse implication has already been seen, cf. pages 161-163.)

This will follow from the

\noindent\textbf{Lemma.} \textit{Let} $u : X' \longrightarrow X$ \textit{be an arrow in} $\Cat/S$\textit{. For} $u \in W^{\mathrm{d}}_{S}$\textit{, it is necessary and sufficient (if $W$} \texttt{[is]} \textit{strongly saturated) that for every} $\xi \in \Hot_W$\textit{, the corresponding arrow}
\[
u_* : p'_!(\xi_{X'}) \longrightarrow p_!(\xi_X)
\]
\textit{be} \texttt{[an]} \textit{isomorphism in} $\HOT_W(S)$\textit{.}

This is sufficient, subject to strong saturation of $W$, for if $u_*$ is \texttt{[an]} isomorphism, passing to the fibers, we find by Der 5$'$
\[
H_\bullet^{\HOT_W}(s\bs X', \xi) \longrightarrow H_\bullet^{\HOT_W}(s\bs X, \xi),
\]
if this is \texttt{[an]} isomorphism, then $s\bs X' \longrightarrow s\bs X$ is by definition a HOT-equivalence, which implies that it is in $W$ (proposition 5).

\msp{page 167}This is necessary, for if $\xi = \hot_W(T)$, hence $\xi_S = \hot_{S,W}(T_S)$ (where $T_S = T \times S/S$), and we therefore have
\[
p'_!(\xi_{X'}) = \hot_{S,W}(X' \times T), \qquad p_!(\xi_X) = \hot_{S,W}(X \times T),
\]
and we have to prove that if $u : X' \longrightarrow X$ is in $W^{\mathrm{d}}_{S}$, then $\underbrace{X' \times T}_{X' \times_S T_S} \longrightarrow \underbrace{X \times T}_{X \times_S T_S}$ also, which simply means that
\[
s\bs(X' \times T) \longrightarrow s\bs(X \times T) \tag*{$(*)$}
\]
is in $W$ for every $s \in S$. But the base change $s\bs S \longrightarrow S$ commutes with fiber products, so transforms $X \times T = X \times_S T_S$ into $(s\bs X) \times_{s\bs S} T_{(s\bs S)} = s\bs X \times T$, and likewise for $X'$. Thus the arrow is identified with
\[
(s\bs X') \times T \longrightarrow (s\bs X) \times T,
\]
obtained from
\[
s\bs X' \longrightarrow s\bs X
\]
by the functor $? \times T$. Since this functor transforms $W$ into $W$, we win.

\noindent\printedpage{103}%
% head: A. Grothendieck : Dérivateurs, Ch. VI, Édition des Universités de Montpellier II et Paris VII
\msp{page 168}Finally, I realize that I had a stupid mental block, over the question of proper base change being compatible with colocal $W$-equivalence. Let
\[
\begin{tikzcd}[column sep=small]
X \arrow[rr, "f"] \arrow[dr] & & Y \arrow[dl] & \\
& S & & S', \arrow[ll, "g"']
\end{tikzcd}
\]
with $f \in W^{\mathrm{d}}_{S}$ and $g$ proper, we have to see that $f' : X' \longrightarrow Y'$ obtained from $f$ by base change $S' \longrightarrow S$ is still in $W^{\mathrm{d}}_{S'}$. This means that every base change of the form $s'\bs S' \longrightarrow S'$ transforms it into $\in W$. But let $s = g(s')$, we have the commutative square
\[
\begin{tikzcd}
S & S' \arrow[l] \\
\underbrace{s\bs S}_{= Z} \arrow[u] & \underbrace{s'\bs S'}_{= Z'} \arrow[l] \arrow[u] \ ,
\end{tikzcd}
\]
and it suffices to see that for every arrow $\varphi : \mathcal{X} \longrightarrow \mathcal{Y}$ in $\Cat/Z$ which is in $W$, its inverse image by base change $Z' \longrightarrow Z$ is likewise. (Applying this to the case $\mathcal{X} = X \times_S (s\bs S)$, $\mathcal{Y} = Y \times_S (s\bs S)$, $\varphi = s\bs f$.) Now this follows from the fact that, by \msp{page 169}the properness assumption, $Z' \longrightarrow Z$ is proper and has Hot-aspheric fibers, hence is a universal Hot-equivalence. It would even suffice for $g$ to be $W$-proper, to imply that $Z' \longrightarrow Z$ is a universal $W$-equivalence, i.e. which remains so after every base change. Now we have the

\noindent\textbf{Lemma.} If $g : Z' \longrightarrow Z$ is a universal $W$-equivalence, then the functor
\[
g^* : X \longmapsto X \times_Z Z' : \Cat/Z \longrightarrow \Cat/Z'
\]
is compatible with the localizers $W_Z$, $W_{Z'}$, i.e. if $f : X \longrightarrow Y$ in $\Cat/Z$ is in $W$, likewise $f' : X' \longrightarrow Y'$ in $\Cat/Z'$.
\[
\begin{tikzcd}
X \arrow[dr, "f"'] \arrow[ddr] & & X' \arrow[ll, "g_X"'] \arrow[dr, "f'"] \arrow[ddr] & \\
& Y \arrow[d] & & Y' \arrow[ll, "g_Y"', crossing over] \arrow[d] \\
& Z & & Z' \arrow[ll, "g"']
\end{tikzcd}
\]
This is obvious from the upper square of the diagram above, where $g_X$, $g_Y$ are in $W$ by the hypothesis on $g$, so $f$ is in $W$ if and only if $f'$ is, q.e.d. % [sic] hypthèse -> hypothèse, letter omitted

\msp{page 170 missing}\texttt{[page 170 missing]}


\endgroup
% END UNIT Der-ChVI-345/Der-ChVI-345_p101-105.tex
% BEGIN UNIT Der-ChVI-345/Der-ChVI-345_p106-110.tex
\begingroup\makeatletter%
\newcommand{\printedpage}[1]{\marginpar{\footnotesize p.~#1}}%
\newcommand{\msp}[1]{\marginpar{\footnotesize [#1]}}%
\newcommand{\ul}[1]{\underline{\mathrm{#1}}}%
\newcommand{\circled}[1]{\tikz[baseline=(c.base)]\node[draw,rounded corners,inner sep=2pt](c){#1};}%
\makeatother\let\printedpage\derpp\let\msp\dermsp%
\let\mspin\DERmspin%


\setcounter{footnote}{81}

\printedpage{104}
% head: A. Grothendieck : Dérivateurs, Ch. VI, Édition des Universités de Montpellier II et Paris VII
% foot: Version 9 novembre 2001

\msp{page 171}

\subsection*{New version of the theory of $\mathrm{HOT}_W$.}

I realized that the hypothesis $W \supseteq W_\infty$ (L3 quater, p. 1), which I had ended up putting on $W$, is artificial, and that the hypothesis L5 (p. 43), or even L6, L6 bis (p. 43-44), is unnecessary for a large part of the theory of $\mathrm{HOT}_W$. I have thus been led to review the theory (i.e. what I have done with it up to now) under the following hypotheses on $W$ (cf. XII):

\begin{itemize}[leftmargin=4em]
\item[W(1)] (\textit{Saturation} of $W$): as before (L1 page 1)
\item[W(2)] (localization): this is the old L4 bis (\footnote{i.e. L5 of page 2. The axiom is taken up again (in a similar and equivalent form) at the bottom of page 43.}).
\item[W(3)] (final object): If $X$ has a final object, $X$ is $W$-aspheric (old L3).
\end{itemize}

These three axioms are sufficient for everything I had done in Pursuing Stacks, and in a large part of what has been done here on HOT. Sometimes however one needs the old L 5, which I restate here (compare XII, theorem \texttt{[theorem number missing]}): % [sic] pout -> pour, misprint

\begin{itemize}[leftmargin=4em]
\item[W(4b)] Let $X$, $Y$ be $W$-fibered over $S$, $f : X \longrightarrow Y$ an $S$-morphism. If $f \in W$, then $\forall s \in \mathrm{Ob}\, S$, $f_s : X_s \longrightarrow Y_s$ is in $W$.
\end{itemize}

For the notion of $W$-fibered (\textit{weak}), I took a long time finally to find \textit{the} right notion, right for a $W$ satisfying only $\mathrm{W}(1,2,3)$ (and even only $\mathrm{W}(1, 2 \text{ bis}, 3)$, where W(2 bis) is the old L4). For this I refer to (XII, theorem 3, p. 51), where the notion is well pinned down by means of equivalent conditions. In fact, for the main applications to $\mathrm{HOT}_W$, seen up to now, one needs only (\footnote{This is ultimately an error, it seems now. I think that the axiom W(4b$'$) ultimately serves strictly no purpose (except in fact to keep false ideas alive.)}) a weakened variant of W(4b):

\begin{itemize}[leftmargin=4em]
\item[W(4b$'$)] Like W(4b), but requiring in addition that $X$, $Y$ be in $\ul{Parf}_\omega\, S$.
\end{itemize}

\msp{page 172}

For the notion of $W$-proper, $W$-smooth, $W$-perfect morphisms in Cat, I refer to XII, part III (p. 22-52), and the notations $\ul{Liss}_W S$, $\ul{Prop}_W S$, $\ul{Parf}_W S$ (subcategories of $\mathrm{Cat}/S$), and $\mathrm{Parf}_W \subseteq \mathrm{Fl}(\mathrm{Cat})$. In the statement of W(4b$'$) we have set
\[
\ul{Parf}_\omega S = \ul{Parf}_{W_\omega} S, \quad
\begin{array}[t]{l}
\text{where } W_\omega \subseteq \mathrm{Fl}(\mathrm{Cat}) \text{ is the smallest subset of } \mathrm{Fl}(\mathrm{Cat})\\
\text{satisfying } \mathrm{W}(1, 2\mathrm{bis}, 3).
\end{array}
\]

I similarly write
\[
W'_\omega \subseteq \mathrm{Fl}(\mathrm{Cat}), \qquad \text{smallest subset of } \mathrm{Fl}(\mathrm{Cat}) \text{ satisfying } \mathrm{W}(1, 2, 3, 4\mathrm{b}). % [sic] satifaisant -> satisfaisant, misprint
\]

\printedpage{105}
% head: A. Grothendieck : Dérivateurs, Ch. VI, Édition des Universités de Montpellier II et Paris VII
% foot: Version 9 novembre 2001

The hypothesis $W \supseteq W_\infty$ which we had ended up imposing on $W$, and which appeared artificial in the course of the reflection in XII, is better replaced by % [sic] apparu -> apparue, agreement with « l'hypothèse »
\[
W \supseteq W'_\omega \qquad (\text{or by } W \supseteq \overline{W'_\omega}, \text{ where } \overline{W'_\omega} \text{ is the strong saturation of } W'_\omega)
\]
(\footnote{$\overline{W'_\omega}$ is the smallest \textit{strongly saturated} subset of $\mathrm{Fl}(\mathrm{Cat})$, satisfying W(4b), and it even satisfies condition W$'$(4a), the strongest variant of the axioms of type W(4).}). But we do not impose it as a mandatory condition, on the same footing as $\mathrm{W}(1,2,3)$, any more than W(4b). The role I have had $W_\infty$ play in the past weeks, seems to belong more naturally to $W'_\omega$ or $\overline{W'_\omega}$. This is so in particular also, it seems to me, in the statement of the conditions of type Der 3 on a derivator (I, p. 111-112), which ought to be reviewed in this light. But everything depends, whether (as is my impression) one can really do substantial homotopical algebra, with in particular the long exact sequences of fibration and suspension, for $\mathrm{HOT}_W$, with $W$'s which are \textit{strictly} contained in $W_\infty$. Before even being sure, it is still too early for me to pronounce on the ``right'' formulation of Der 3,

\msp{page 173}

which is the axiom relating the derivator $\mathbf{D}$ to $W_\infty$, or to some other basic localizer.

I am now going to give a summary of the theory of $\mathrm{HOT}_W$, reviewed from the viewpoint I have just stated: we assume for $W$ \textit{the conditions} $\mathrm{W}(1,2,3)$ \textit{and no more,} unless explicitly mentioned.

\circled{1} \textbf{Main localizers} on $\mathrm{Cat}/S$ and on the main categories of $S$-categories (cf. diagram p. 76).

They will all be induced by the three familiar localizers
\[
\begin{tikzcd}[column sep=small]
W_S^{\mathrm{d}} \arrow[dr, hook] & & W_S^{\mathrm{g}} \arrow[dl, hook'] \\
& W_S &
\end{tikzcd}
\]
on $\mathrm{Cat}/S$. There is also reason to look at $W_S^{\mathrm{f}}$, but never on the whole of $\mathrm{Cat}/S$, but on other categories of the diagram, and always in cases where $W_S^{\mathrm{f}}$ coincides either with $W_S^{\mathrm{d}}$ or with $W_S^{\mathrm{g}}$. But the $W_S^{\mathrm{f}}$ point of view is often more directly convenient, and moreover has the advantage of being symmetric.

\circled{a} In $\ul{Liss}_W S$, we have

\printedpage{106}
% head: A. Grothendieck : Dérivateurs, Ch. VI, Édition des Universités de Montpellier II et Paris VII
% foot: Version 9 novembre 2001

\[
\begin{tikzcd}[row sep=large, column sep=0pt]
W_S^{\mathrm{f}} \arrow[d, Rightarrow] & {} = W_S^{\mathrm{d}} = W_S^{\mathrm{univ}} \\
W_S^{\mathrm{g}} \arrow[d, Rightarrow] & \\
W_S &
\end{tikzcd}
\]

Dually,

\circled{b} in $\ul{Prop}_W S$, we have
\[
\begin{tikzcd}[row sep=large, column sep=0pt]
W_S^{\mathrm{f}} \arrow[d, Rightarrow] & {} = W_S^{\mathrm{g}} = W_S^{\mathrm{univ}} \\
W_S^{\mathrm{d}} \arrow[d, Rightarrow] & \\
W_S, &
\end{tikzcd}
\]

\msp{page 174}

finally,

\circled{c} in $\ul{Fib}_W S$, we have
\[
\begin{tikzcd}[row sep=large]
W_S^{\mathrm{f}} \arrow[r, Leftrightarrow] \arrow[d, Rightarrow] & W_S^{\mathrm{d}} \arrow[r, Leftrightarrow] & W_S^{\mathrm{g}} \arrow[r, Leftrightarrow] & W_S^{\mathrm{univ}} \\
W_S. & & &
\end{tikzcd}
\]

This implication $W_S^{\mathrm{f}} \Longrightarrow W_S$ is strict if one does not impose on $W$ axiom W(4b), whose usefulness here is exactly to ensure
\[
\left|\;
\begin{array}{l}
W_S^{\mathrm{f}} \Longleftrightarrow W_S \ \text{ in } \ul{Fib}_W S \ \textit{if}\ \mathrm{W}(4\mathrm{b})\ \texttt{[is]}\ \text{satisfied}\\[1.5ex]
\text{and then the five localizers considered above}\\[1.5ex]
\text{in } \ul{Fib}_W S \text{ are all equal.}
\end{array}
\right.
\]

The relations between [?] localizers in $\ul{Liss}_W S$, $\ul{Prop}_W S$, $\ul{Fib}_W S$ carry over automatically to subcategories of these, more generally to all those which map
\printedpage{107}
% head: A. Grothendieck : Dérivateurs, Ch. VI, Édition des Universités de Montpellier II et Paris VII
% foot: Version 9 novembre 2001
by an inclusion functor into one of these. In the diagram p. 76, all the categories, with the sole exception of the whole of $\mathrm{Cat}/S$, are in this case.

For the main results developed up to now on $\mathrm{HOT}_W$, when axiom W(4b) is used, the weak form W(4b$'$) suffices. We have
\[
\left|\;
\begin{array}{l}
W_S^{\mathrm{f}} \Longleftrightarrow W_S \ \text{ in } \ul{Parf}_\omega S \ \textit{if}\ \mathrm{W}(4\mathrm{b}')\ \texttt{[is]}\ \text{satisfied}\\[1.5ex]
\text{and then the five localizers } W_S^{\mathrm{f}},\, W_S^{\mathrm{d}},\, W_S^{\mathrm{g}},\, W_S^{\mathrm{univ}},\, W_S\\[1.5ex]
\text{are equal in } \ul{Parf}_\omega S.
\end{array}
\right.
\]

\circled{2} \textbf{Main categories of categories over $S$.}

These are, up to changes in lighting, those of the diagram p. 76. The categories $\ul{Fibsc}\,S$, $\ul{Fib}\,S$, $\ul{Cofibsc}\,S$, $\ul{Cofib}\,S$ and their variants with subscript 0, are unchanged - their role is to be as small

\msp{page 175}

as possible to fulfill their function (of describing $\mathrm{HOT}_W(S)$, $\mathrm{HOT}^{o}_W(S)$, $\mathrm{HOT}^{\mathrm{lc}}_W(S)$, or $\mathrm{HOT}^{\mathrm{lc}!}_W(S)$ below). The category $\ul{Parf}\,S$ of the diagram fulfilled a function of the same type, which is why we now replace it (bidding higher) by $\ul{Parf}_\omega S$. The other categories of the diagram are, in addition to $\mathrm{Cat}/S$, $\ul{Liss}\,S$ and $\ul{Prop}\,S$ and their variants with subscript 0, as well as $\ul{Fibhot}\,S$. (All these categories, as well as $\ul{Parf}\,S$, were thus constructed with reference to $W_\infty$, not to $W$.) Their role is to be as large as possible to give rise to certain properties and relations. Now it turns out that here the reference to $W_\infty$ was artificial, and one should put instead the categories $\ul{Liss}_W S$, $\ul{Prop}_W S$, $\ul{Fib}_W S$, and in addition the variants $\ul{Liss}_{0,W} S$, $\ul{Prop}_{0,W} S$. Thus the diagram becomes % [sic] propiétés -> propriétés, misprint
\[
\begin{tikzcd}[column sep=tiny, row sep=small]
 & & & & \ul{Fib}_W S \arrow[dd, hook, "\delta"] \arrow[dddllll, dotted, bend right=35, "\Phi_0"'] \arrow[dddrrrr, dotted, bend left=35, "\Psi_0"] & & & & \\
 & & \ul{Liss}_{0,W} S \arrow[urr, hook] \arrow[dd, hook] & & & & \ul{Prop}_{0,W} S \arrow[ull, hook'] \arrow[dd, hook] & & \\
 & \ul{Fib}_{0,W} S \arrow[ur, hook] \arrow[dd, hook] & & & \mathrm{Cat}/S \arrow[dd, "C"] \arrow[dddllll, dotted, bend right=30, "\Phi"] \arrow[dddrrrr, dotted, bend left=30, "\Psi"'] & & & \ul{Cofib}_{0,W} S \arrow[ul, hook'] \arrow[dd, hook] & \\
\ul{Fibsc}_{0,W} S \arrow[ur] \arrow[dd, hook] & & \ul{Liss}_W S \arrow[urr, hook] & & & & \ul{Prop}_W S \arrow[ull, hook'] & & \ul{Cofibsc}_{0,W} S \arrow[ul] \arrow[dd, hook] \\
 & \ul{Fib}\,S \arrow[ur, hook] & & & \ul{Parf}_\omega S \arrow[uuull, hook] \arrow[uuurr, hook'] & & & \ul{Cofib}\,S \arrow[ul, hook'] & \\
\ul{Fibsc}\,S \arrow[ur] & & & & & & & & \ul{Cofibsc}\,S \arrow[ul]
\end{tikzcd}
\]
(\footnote{NB the new notation would rather be $\ul{Fibf}_W S$, for these are $W$-fibrations which I now call ``weak'' (and defined only in Cat, as opposed to $W$-fibrations, definable in any category of models).}) (all the squares \texttt{[are]} cartesian, the functors denoted $\hookrightarrow$ are fully faithful),

\printedpage{108}
% head: A. Grothendieck : Dérivateurs, Ch. VI, Édition des Universités de Montpellier II et Paris VII
% foot: Version 9 novembre 2001

\msp{page 176}

where I have again marked in green \texttt{[in the original only]} the five \textit{``inversion functors''}
\[
\Phi,\ \Psi,\ \Phi_0,\ \Psi_0,\ C,
\]
which provide quasi-inverses of the functors on the localized categories induced by the \textit{inclusion functors.}

\circled{3} \textbf{The four types of categories} $\mathrm{HOT}_W$\textbf{.}

On each of the categories of the upper floor, there are \textit{two} localizers to look at, % [sic] l'etage -> l'étage, missing accent (printed « l'étage » below)
\[
\begin{tikzcd}[row sep=huge, column sep=0pt]
W_S^{\mathrm{f}} \arrow[d, Rightarrow] & {} = W_S^{\mathrm{g}} = W_S^{\mathrm{d}} = W_S^{\mathrm{univ}} \\
W_S, &
\end{tikzcd}
\]
giving rise to two distinct localized categories, denoted
\[
\mathrm{HOT}^{\mathrm{lc}}_W(S) \qquad \text{and} \qquad \mathrm{HOT}^{\mathrm{lc}!}_W(S),
\]
it turns out that (thanks to $\Phi_0$ and $\Psi_0$) these localized categories ``are the same'' for the 8 categories of the upper floor, with the exception of $\ul{Parf}_\omega S$; more exactly, the inclusion functors (with the exception perhaps of those having as source $\ul{Parf}_\omega S$) induce equivalences on the two types of localized categories, the quasi-inverse equivalences being obtained by $\Phi_0$ (left wing) or $\Psi_0$ (right wing). When one localizes \textit{for} $W_S$ (to the exclusion of $W_S^{\mathrm{f}}$), one also finds quasi-inverses by $C|\ul{Liss}_{0,W} S$ resp. $C|\ul{Prop}_{0,W} S$, for the two inclusions
\[
\begin{tikzcd}
\ul{Liss}_{0,W} S & \ul{Parf}_\omega S \arrow[l, hook'] \arrow[r, hook] & \ul{Prop}_{0,W} S
\end{tikzcd}
\]
(\footnote{cf. below, what one can say for the localized categories for $W_S^{\mathrm{f}}$.}). For reasons of symmetry, the prettiest seems to be to define
\[
\left\{
\begin{array}{lcl}
\mathrm{HOT}^{\mathrm{lc}}_W(S) & = & (\ul{Fib}_W S)(W_S^{\mathrm{f}})^{-1}\\
\mathrm{HOT}^{\mathrm{lc}!}_W(S) & = & (\ul{Fib}_W S)(W_S)^{-1},
\end{array}
\right.
\]
and we have a localization functor
\[
\mathrm{HOT}^{\mathrm{lc}}_W(S) \xrightarrow{\ \lambda\ } \mathrm{HOT}^{\mathrm{lc}!}_W(S)
\]
induced by the inclusion
\[
W_S^{\mathrm{f}} \subseteq W_S.
\]


\endgroup
% END UNIT Der-ChVI-345/Der-ChVI-345_p106-110.tex
% BEGIN UNIT Der-ChVI-345/Der-ChVI-345_p111-115.tex
\begingroup\makeatletter%
\NoAutoSpacing%
\newcommand{\printedpage}[1]{\marginpar{\footnotesize p.~#1}}%
\newcommand{\msp}[1]{\marginpar{\footnotesize [#1]}}%
\newcommand{\Fib}{\underline{\mathrm{Fib}}}%
\newcommand{\Parf}{\underline{\mathrm{Parf}}}%
\newcommand{\Prop}{\underline{\mathrm{Prop}}}%
\newcommand{\Liss}{\underline{\mathrm{Liss}}}%
\newcommand{\Fibsc}{\underline{\mathrm{Fibsc}}}%
\newcommand{\Cofibsc}{\underline{\mathrm{Cofibsc}}}%
\newcommand{\HOT}{\mathrm{HOT}}%
\newcommand{\Cat}{\mathrm{Cat}}%
\newcommand{\Hom}{\mathrm{Hom}}%
\newcommand{\Ob}{\mathrm{Ob}}%
\newcommand{\lcs}{\mathrm{lc}\,!}%
\newcommand{\cercle}[1]{\tikz[baseline=(c.base)]\node[draw,circle,inner sep=1pt](c){#1};}%
\newcommand{\troisbarres}{\leavevmode\llap{$|||$\quad}}%
\makeatother\let\printedpage\derpp\let\msp\dermsp%
\let\mspin\DERmspin%


\printedpage{109}
% head: A. Grothendieck : Dérivateurs, Ch. VI, Édition des Universités de Montpellier II et Paris VII

\noindent\msp{page 177}
For every full subcategory
\[
\mathcal{F} \subseteq \Fib_W S
\]
which contains one of the three subcategories
\[
\Fib_{0,W}S,\ \Parf_\omega S,\ \Prop_{0,W}S,
\]
the inclusion of $\mathcal{F}$ into $\Fib_W S$ defines equivalences
\[
\mathcal{F}(W_S)^{-1} \approx \HOT^{\lcs}_W(S),
\]
and if $\mathcal{F}$ contains one of the two subcategories $\Fib_{0,W}S$, $\Prop_{0,W}S$ (to the exclusion of $\Parf_\omega S$), one also finds
\[
\mathcal{F}(W^{\mathrm{f}}_S)^{-1} \approx \HOT^{\mathrm{lc}}_W(S).
\]
The ground floor contributes two other variants of HOT, by localizing either by $W^{\mathrm{d}}_S$ (left wing), or by $W^{\mathrm{g}}_S$ (right wing). The inclusion functors define equivalences for the localized categories either for $W^{\mathrm{d}}_S$ (left wing), or for $W^{\mathrm{g}}_S$ (right wing). The prettiest (symmetry) seems to be to define these two localized categories by $\Cat/S$
\[
\HOT_W(S) = (\Cat/S)(W^{\mathrm{d}}_S)^{-1}, \qquad \HOT^{o}_W(S) = (\Cat/S)(W^{\mathrm{g}}_S)^{-1}.
\]
The quasi-inverses for the equivalences of localized categories, induced by the inclusion functors, are given either by $\Phi$ (left wing), or by $\Psi$ (right wing).

But $\Phi$ and $\Psi$ supply quasi-inverses also when one localizes for $W_S$, the coarse $W$-equivalence instead of the fine - just as on the upper floor. One therefore finds on the ground floor a (unique) third localized category, which one can describe more prettily by
\[
(\Cat/S)(W_S)^{-1},
\]
but it is equivalent to $(\Fib_W S)W_S^{-1}$ by the
\msp{page 178}
functor induced by the inclusion. (A quasi-inverse being supplied by $C$.) It is therefore $\HOT^{\lcs}_W(S)$. One thus finds, from the commutative diagram of categories with localizers
\[
\begin{tikzcd}[column sep=small]
 & (\Fib_W S, W^{\mathrm{f}}_S) \arrow[dl] \arrow[dd] \arrow[dr] & \\
(\Cat/S, W^{\mathrm{d}}_S) \arrow[ddr] & & (\Cat/S, W^{\mathrm{g}}_S) \arrow[ddl] \\
 & (\Fib_W S, W_S) \arrow[d] & \\
 & (\Cat/S, W_S) &
\end{tikzcd}
\]
% foot: Version 9 novembre 2001
\printedpage{110}
% head: A. Grothendieck : Dérivateurs, Ch. VI, Édition des Universités de Montpellier II et Paris VII
(where the short vertical arrow induces an equivalence on the localized categories), the commutative diagram between the $\HOT_W$ categories linked to $S$:
\[
\begin{tikzcd}[column sep=small]
 & \HOT^{\mathrm{lc}}_W(S) \arrow[dl, "i"'] \arrow[dd, "\lambda"'] \arrow[dr, "j"] & \\
\HOT_W(S) \arrow[dr, "\rho"'] & & \HOT^{o}_W(S)(\simeq \HOT_W(S^{o})) \arrow[dl, "\sigma"] \\
 & \HOT^{\lcs}_W(S), &
\end{tikzcd}
\tag*{$(\ast)$}
\]
where $\rho$, $\sigma$, $\lambda$ are localization functors, and where $i$, $j$ are felt to be inclusion functors. This diagram $(\ast)$ seems to exhaust the natural functors which link the four categories considered to one another.

\noindent\msp{page 179}
The properties of the functors $\Phi$, $\Psi$, $C$ which have been used, to supply all the desired quasi-inverses of functors (between localized categories) induced by the inclusion functors of the diagram, are the following.

\cercle{a} We have a functorial homomorphism in $\Cat/S$
\[
X \xrightarrow{i_X} \Phi(X), \qquad \underline{i_X \in W^{\mathrm{d}}_S}.
\]
When $X$ is in $\Fib_W S$, $\Phi(X)$ is in $\Fibsc_0 S$ \texttt{[rather }$\Fibsc_{0,W}S$\texttt{, cf. diag. p. 175]}, so $i_X \in W^{\mathrm{f}}_S$ ($= W^{\mathrm{d}}_S$ on $\Fib_W S$). Thus $\Phi$ induces the functor $\Phi_0$ of the diagram.

Dually, we have
\[
X \xrightarrow{i_X} \Psi(X), \qquad \underline{i_X \in W^{\mathrm{g}}_S},
\]
and when $X$ is in $\Fib_W S$, $\Psi(X)$ is in $\Cofibsc_0 S$ \texttt{[rather }$\Cofibsc_{0,W}S$\texttt{]}, hence the functor $\Psi_0$.

\cercle{b} We have a functorial homomorphism
\[
\left|
\begin{aligned}
& X \xrightarrow{i_X} C(X), \qquad \underline{i_X \in W_S} \\
& \text{and } C(X) \in \Ob\,\Parf_\omega S.
\end{aligned}
\right.
\]

\troisbarres But one must bear in mind that in the absence of axiom W(4b) on $W$, one can \textit{not} assert $i_X \in W^{\mathrm{f}}_S$, even if $X$ is in $\Fib_W S$ (or even, in $\Parf_\omega S$), the only case (at any rate) in which one is entitled to expect it. I note in this connection the diagram, \textit{cartesian,} for $f : X \longrightarrow Y$ in $\Cat/S$
% foot: Version 9 novembre 2001
\printedpage{111}
% head: A. Grothendieck : Dérivateurs, Ch. VI, Édition des Universités de Montpellier II et Paris VII
\[
\begin{tikzcd}
C(Y)_s \arrow[d, "\pi^{s}_{Y}"'] & C(X)_s \arrow[l, "C(f)_s"'] \arrow[d, "\pi^{s}_{X}"] \\
Y & X, \arrow[l, "f"']
\end{tikzcd}
\qquad \text{\footnotesize vertical arrows in $\Parf_\omega \subseteq \Fib_W$}
\]
which shows for example that
\msp{page 180}
one cannot expect that $f \in W_S \implies C(f) \in W^{\mathrm{f}}_S$ without having axiom W(4a$'$) \texttt{[what is this?]} (or an axiom of comparable strength, such as W(4b)).

It is these negative facts (in the absence of W(4b)) that are the reason why I cannot say anything in particular about the functors induced by \textit{fine localization} $(W^{\mathrm{f}}_S)^{-1}$ of the two inclusion functors
\[
\Liss_0 S \longleftarrow \Parf_\omega S \longrightarrow \Prop_0 S,
\]
or, what amounts to the same thing, about the one induced by the common composite of those with the inclusions into $\Fib_W S$
\[
\begin{tikzcd}[column sep=tiny, row sep=small]
 & \Parf_\omega S \arrow[dl] \arrow[dd] \arrow[dr] & \\
\Liss_0 S \arrow[dr] & & \Prop_0 S \arrow[dl] \\
 & \Fib_W S, &
\end{tikzcd}
\]
in other words, I can say nothing about the faithfulness properties of the functor
\[
(\Parf_\omega S)(W^{\mathrm{f}}_S)^{-1} \longrightarrow (\Fib_W S)(W^{\mathrm{f}}_S)^{-1},
\]
and still less that it is an equivalence of categories. Thus, in the absence of W(4b), for the full subcategories $\mathcal{F}$ of $\Fib_W S$ which contain $\Parf_\omega S$, without containing either $\Liss_0 S$ or $\Prop_0 S$, it could a priori be that the localized categories
\[
\mathcal{F}(W^{\mathrm{f}}_S)^{-1}
\]
are all different. I am thinking in particular of $\Parf_W S$, as interesting in itself, after all, as $\Parf_\omega S$. So in addition to the four localized categories of diagram $(\ast)$, p. 178, there is a whole infinite ``teratological'' series of localized categories
\msp{page 181}
which refuse to let themselves be ``closed up'' \texttt{[?]} by $\Phi_0$ or $\Psi_0$, and for which even $C$ is of no help, for lack of axiom W(4b). % [sic] axime -> axiome, missing letter
\DERmark{ENV0379}


\subsection*{\cercle{4} Case where axiom \textmd{W(4b)} is satisfied.}
% foot: Version 9 novembre 2001

\printedpage{112}
% head: A. Grothendieck : Dérivateurs, Ch. VI, Édition des Universités de Montpellier II et Paris VII
Then the perplexities I have just spoken of disappear. We have at once, by virtue of W(4b):
\[
|||\qquad
\begin{array}{@{}r@{}}
\text{If } X \in \Ob\,\Fib_W S, \text{ then } i_X : X \longrightarrow C(X) \text{ is in } W^{\mathrm{f}}_S \\
(\textit{assuming axiom } \text{W(4b)})
\end{array}
\]
and $C$ supplies quasi-inverses for the ``inclusion'' functors between fine localized categories (for $W^{\mathrm{f}}_S$) considered on the preceding page. It supplies \texttt{[``It supplies'' crossed out]} one as well for the functor $\lambda$ considered above ($(\ast)$ page 178)
\[
\underbrace{(\Fib_W S)(W^{\mathrm{f}}_S)^{-1}}_{\overset{\mathrm{def}}{=}\HOT^{\mathrm{lc}}_W(S)} \longrightarrow \underbrace{(\Cat/S)(W_S)^{-1}}_{\simeq \HOT^{\lcs}_W(S)},
\]
\DERmark{ENV0380}

which is therefore an equivalence. Another, more obvious, way of saying the same thing is that on each of the categories of the upper floor the two localizers $W^{\mathrm{f}}_S$ and $W_S$ are now identical (since they are all contained in $\Fib_W S$, where they \texttt{[rather they]} are so precisely by axiom W(4b)), so the functor $\lambda$ between these two localized categories (original definition of $\lambda$, in the case $\mathcal{F} = \Fib_W S$, cf. page 176) is even the \textit{identity functor,} the two localized categories being \textit{equal:}
\[
\HOT^{\mathrm{lc}}_W(S) = \HOT^{\lcs}_W(S) \qquad\qquad !
\]
Thus the diagram $(\ast)$ p. 178 takes the form
\[
\begin{tikzcd}[column sep=small]
 & \HOT^{\mathrm{lc}}_W(S) \arrow[dl, shift left=1.2, "i"] \arrow[dr, shift right=1.2, "j"'] & \\
\HOT_W(S) \arrow[ur, shift left=1.2, "\rho"] & & \HOT^{o}_W(S). \arrow[ul, shift right=1.2, "\sigma"']
\end{tikzcd}
\tag*{$(\ast)$}
\]

\noindent\msp{page 182}
When one interprets not only the two lateral categories $\HOT_W(S)$, $\HOT^{o}_W(S)$, but also the one in the middle, as localized categories of $\Cat/S$, for the localizers $W^{\mathrm{d}}_S$, $W^{\mathrm{g}}_S$ and $W_S$ respectively, the functors $\rho$, $\sigma$ are those deduced from the two inclusions
\[
\begin{tikzcd}[column sep=small, row sep=small]
W^{\mathrm{d}}_S \arrow[dr, hook] & & W^{\mathrm{g}}_S \arrow[dl, hook'] \\
 & W_S &
\end{tikzcd}
\qquad \text{\footnotesize as subsets of Fl(Cat/$S$),}
\]
whereas the functors $i$, $j$ are deduced, by passage to the localized categories, from one and the same functor
\[
\Cat/S \xrightarrow{C} \Cat/S,
\]
this functor having the virtue of mapping $W_S$ into $W^{\mathrm{d}}_S \cap W^{\mathrm{g}}_S$,
\[
f \in W_S \overset{\text{assuming W(4b)}}{\Longrightarrow} C(f) \in W^{\mathrm{d}}_S \cap W^{\mathrm{g}}_S \qquad (\text{and even } \in W^{\mathrm{univ}}_S).
\]
% foot: Version 9 novembre 2001

\printedpage{113}
% head: A. Grothendieck : Dérivateurs, Ch. VI, Édition des Universités de Montpellier II et Paris VII
(And for good reason, since $C$ factors through $\Fib_W S$ (and even through $\Parf_\omega S$), where $W_S$ is identical to $W^{\mathrm{d}}_S$, $W^{\mathrm{g}}_S$, $W^{\mathrm{univ}}_S$. That $f \in W_S$ implies $C(f) \in W$ is independent of axiom W(4b), and is seen from the diagram
\[
\begin{tikzcd}
X \arrow[r, "f"] \arrow[d, "i_X"'] & Y \arrow[d, "i_Y"'] \\
C(X) \arrow[r, "C(f)"] & C(Y),
\end{tikzcd}
\]
where the arrows $i_X$, $i_Y$ are in $W_S$.)

That said, it is true that the pairs
\[
(\rho, i),\ (\sigma, j)
\]
are pairs of adjoint functors, hence ($\rho$ and $\sigma$ being localization functors) that $i$ and $j$ are \textit{fully faithful.} The adjunction formula can be read as
\[
\boxed{\Hom_{W^{\mathrm{g}}_S}(X, CY) \simeq \Hom_{W^{\mathrm{d}}_S}(X, CY) \simeq \Hom_{W_S}(X, Y)}
\tag*{$(\ast)$}
\]
\msp{page 183}
for $X$, $Y$ in $\Cat/S$ with nothing more. A $\Hom_\Sigma$ indicates the Hom in $(\Cat/S)\Sigma^{-1}$.

\noindent\textbf{Corollary.} \textit{If $X$ in} $\Cat/S$\textit{, $Y$ in} $\Fib_W S$\textit{, then}
\[
\begin{tikzcd}[column sep=small]
\Hom_{W^{\mathrm{g}}_S}(X, Y) \arrow[dr, "\sim"] & & \Hom_{W^{\mathrm{d}}_S}(X, Y) \arrow[dl, "\sim"'] \\
 & \Hom_{W_S}(X, Y) &
\end{tikzcd}
\]
\textit{are bijective.}

Indeed, $Y$ in $\Fib_W(S)$ implies that $Y \xrightarrow{i_Y} C(Y)$ is in $W^{\mathrm{univ}}_S \subseteq W^{\mathrm{d}}_S \cap W^{\mathrm{g}}_S$, so
\[
\Hom_{W^{\mathrm{g}}_S}(X, Y) \xrightarrow{\sim} \Hom_{W^{\mathrm{g}}_S}(X, CY) \underbrace{\xrightarrow{\sim}}_{\text{adjunction}} \Hom_{W_S}(X, Y),
\]
and likewise for $\Hom_{W^{\mathrm{d}}_S}(X, Y)$.

\subsection*{\cercle{5} Perplexities about the weakened variants of axiom \textmd{W(4b)}.}

What remains of all these beautiful things (only three HOT categories left, instead of $4+\infty$!) when one now assumes only W(4b$'$) instead of W(4b), i.e. the equality of $W_S$ and $W^{\mathrm{f}}_S$ in
% foot: Version 9 novembre 2001


\endgroup
% END UNIT Der-ChVI-345/Der-ChVI-345_p111-115.tex
% BEGIN UNIT Der-ChVI-345/Der-ChVI-345_p116-120.tex
\begingroup\makeatletter%
\newcommand{\printedpage}[1]{\marginpar{\footnotesize p.~#1}}%
\newcommand{\msp}[1]{\marginpar{\footnotesize [#1]}}%
\newcommand{\Parf}{\underline{\mathrm{Parf}}}%
\newcommand{\Fib}{\underline{\mathrm{Fib}}}%
\newcommand{\Liss}{\underline{\mathrm{Liss}}}%
\newcommand{\Prop}{\underline{\mathrm{Prop}}}%
\newcommand{\Bifib}{\underline{\mathrm{Bifib}}}%
\newcommand{\Cofibsc}{\underline{\mathrm{Cofibsc}}}%
\newcommand{\Fibsc}{\underline{\mathrm{Fibsc}}}%
\newcommand{\HOT}{\mathrm{HOT}}%
\newcommand{\margbar}{\marginpar{\(\|\)}}%
\makeatother\let\printedpage\derpp\let\msp\dermsp%
\let\mspin\DERmspin%


\printedpage{114}
% head: A. Grothendieck : Dérivateurs, Ch. VI, Édition des Universités de Montpellier II et Paris VII
% foot: Version 9 novembre 2001

$\Parf_\omega S$? Finally, it seems to me (contrary to what I had thought I saw) that almost nothing remains (\footnote{\textbf{NB} If $W$ satisfies the axiom of limits, then W(4b$'$) implies W(4b), cf. the summary theorem in XVI 4.4.13 (p. 58-62), in particular $2^{o}$). Condition W(4b) is condition B of this statement, and W(4b$'$) condition D. (It is implied by W(4a) $=$ C, which I did not suspect in the least when writing XII \dots)}). We still do not then know, for $X$ in $\Fib_W S$, whether
\[ i_X : X \longrightarrow C(X) \]
might not be in $W^{\mathrm{f}}_S$, and not only in $W_S$; we only know it for $X \in \Parf_\omega S$, but in this case we do not need it, for the point of the functor $C(X)$ and of the inclusion $i_X : X \longrightarrow C(X)$ is ``to improve'' $X$ and to make it perfect, or even $W$-perfect. But if it already is, it is no longer worth the trouble. There will be, as
\msp{page 184}
before, 4 essentially distinct localized categories $\HOT_W$, $\HOT^{o}_W$, $\HOT^{\mathrm{lc}}_W$, $\HOT^{\mathrm{lc}\,!}_W$, plus infinitely many others $\mathcal{F}(W^{\mathrm{f}}_S)^{-1}$ for $\mathcal{F}$ between $\Parf_\omega S$ and $\Fib_W S$ (cf. page 180), with the sole benefit that we can now exclude the case $\mathcal{F} = \Parf_\omega S$, for in this case $\mathcal{F}(W^{\mathrm{f}}_S)^{-1} = \mathcal{F}(W_S)^{-1} = \HOT^{\mathrm{lc}\,!}_W(S)$ is one of the four canonical ones. The decisive thing is that we still have no way of proving that
\[ \lambda : \HOT^{\mathrm{lc}}_W(S) \longrightarrow \HOT^{\mathrm{lc}\,!}_W(S) \]
is an equivalence, and I doubt that this is the case. We would have to look with $W = W_{\mathbf{D}}$, the $\mathbf{D}$-equivalence, where $\mathbf{D}$ is a derivator Der($k$-Mod), $k$ a given ring (e.g. $k = \mathbf{Z}$, or $k$ a field $\mathbf{Z}/p\mathbf{Z}$ or $\mathbf{Q}$ \dots). I am not going to do it now, however. (It is true that the axiom W(4b$'$) is not satisfied in this case either \dots)

If in the statement of the axiom, we replaced $\Parf_\omega S$ by $\Parf_W S$, nothing essential would be changed - at most I would replace $\Parf_\omega S$ by $\Parf_W S$ in the brilliant diagram p. 175. Finally, I am now convinced that it has to be all or nothing: either \textit{four} categories $\HOT(S)$ for $W$, being content with W$(1,2,3)$, or else only three, i.e. no $\HOT^{\mathrm{lc}\,!}_W(S)$ - and for that, it does not seem that we can make do with less than W(4b).

\msp{page 185}
It would suffice, it is true, to assume that $W_S = W^{\mathrm{f}}_S$ on any one of the seven categories of the upper floor other than $\Parf_\omega S$. This would then imply that in the other six, we will have $\overline{W_S} = \overline{W^{\mathrm{f}}_S}$ (same strong saturations in $\mathrm{Fl}(\mathcal{F})$). One of the two weakest of the axioms in question would be: % doubt resolved in DER_VI-c_corrections.csv; input note: line under "étage" -- read as the overline of $\overline{W^{\mathrm{f}}_S}$ on the next line, not an underline of "étage"

\begin{quote}
If $f : X \longrightarrow Y$ is a \textit{cartesian} functor of Cat-fibered and $W$-fibered categories (\footnote{``$W$-fibered'' meaning here that for $u \in \mathrm{Fl}\,S$, the base change functors $u^*$ $X_{s'} \longrightarrow X_s$, $Y_{s'} \longrightarrow Y_s$ are in $W$.}) over another $S$, if $f \in W$, then $f_s \in W$ for every $s \in S$,
\end{quote}

(and the other is the dual version, for cocartesian functors of cofibered categories). Unless I am mistaken, this axiom must already imply that $\overline{W}$ satisfies the strongest version W(4a)
\printedpage{115}
% head: A. Grothendieck : Dérivateurs, Ch. VI, Édition des Universités de Montpellier II et Paris VII
% foot: Version 9 novembre 2001
of all (a fortiori W(4b)) of the axioms W(4). This axiom therefore now seems to me much more reasonable than W(4b$'$), and this all the more so as it refers directly to the description of $\HOT^{\mathrm{lc}}_W(S)$ in terms of functors (or of ``pseudofunctors'', if we do not want to assume the categories split and the functors compatible with the splittings) $S^{o} \longrightarrow \mathrm{Cat}$. % [sic] directemment -> directement, spelling
Besides, although it does not look it, this axiom \textit{is} self-dual, since the dual version, for fixed $S$, is the direct version for $S^{o}$! This makes this axiom still more appetizing. So to formulate W(4a$'$), I now feel like dropping $\Parf_\omega$, and doing the same. If we keep two versions W(4a$'$) (the old one) and W(4a$''$) (the new one), both have the virtue of implying W(4b). And W(4b$'$) (new version) must imply W(4a) (ultrastrong version), if not for $W$, at least for $\overline{W}$.

My ideas were tenaciously [?] muddled
\msp{page 186}
until this very day, for I had the idea that, thanks to the wondrous functor $C$, the category $\Parf_\omega S$, like the seven other categories of the upper floor, could serve to define $\HOT^{\mathrm{lc}}_W(S)$. But this is not so, it seems, unless we assume the axiom W(4b), or the variant W(4b$'$) (new version). The old version of W(4b$'$) turns out to be perfectly unusable. % [sic] supposeer -> supposer, spelling

On this subject, I have still not given a clear answer, however, not a ``yes'' or a ``no'', to the question:

\begin{quote}
\textbf{Question:} \textit{do the inclusions}
\[
\begin{tikzcd}
 & \Liss_{0,W}S \\
\Parf_W S \arrow[ur] \arrow[dr] & \\
 & \Prop_{0,W}S
\end{tikzcd}
\]
\textit{induce, even in the absence of W(4b) (or of its weakened variants), equivalences for the fine localized categories (by $W^{\mathrm{f}}_S$)? Same question for} $\Parf_\omega S$ \textit{instead of} $\Parf_W S$ (\footnote{OK in the presence of the axiom W(6bis) on the $\varinjlim$ , cf. p. 192 ff.}). \textit{Same question for}
\[ \Bifib\, S \subseteq \Parf_\omega S, \]
\textit{the category of categories over $S$ which are both} Cat-\textit{fibered and} Cat-\textit{cofibered.}
\end{quote}

I cannot tell, even granting (this time) the axiom W(4a), whether $\Bifib\, S$ is big enough to be able to describe all the locally constant relative homotopy types over $S$, i.e. whether % [sic] tous le types -> tous les types, missing letter
\[ (\Bifib\, S)(W^{\mathrm{f}}_S)^{-1} \longrightarrow (\Fib_W S)(W^{\mathrm{f}}_S)^{-1} \]
is at least essentially surjective, if not an equivalence of categories. What is missing is a functor
\[ \mathrm{Cat}/S \longrightarrow \Bifib\, S \qquad \text{or at least} \qquad \underbrace{\Fib_W S}_{\text{or $\Fib_\omega S$\,?}} \longrightarrow \Bifib\, S, \]
\printedpage{116}
% head: A. Grothendieck : Dérivateurs, Ch. VI, Édition des Universités de Montpellier II et Paris VII
% foot: Version 9 novembre 2001
which could play the role of $\Phi$ and of $\Psi$ \textit{simultaneous\msp{page 187}ly.} I have not even checked that $\Phi$ does \textit{not} transform objects of $\Prop_W S$ or of $\Prop_{0,W}(S)$ into $W$-proper objects (hence $W$-perfect, being in any case $W$-smooth and even $W_\omega$-smooth) over $S$ (\footnote{See the answer to this interesting question on page 192.}). If this were nevertheless true on $\Prop_W S$, that would show that the essential image of $\Prop_W S$ in $\HOT_W(S)$ ($= (\mathrm{Cat}/S)(W^{\mathrm{d}}_S)^{-1}$) is contained in $\HOT^{\mathrm{lc}}_W(S)$. % [sic] contenu -> contenue, agreement with "l'image"
Taking $S = \boldsymbol{\Delta}^1$, this would mean that if in $X$ we have a closed subcategory $X_1$ such that $X_1 \overset{i}{\hookrightarrow} X$ is $W_\omega$-aspheric, then $1\backslash X \longrightarrow 0\backslash X$ is in $W$ - but $0\backslash X \simeq X$, $1\backslash X \simeq X_1$, the arrow is $X_1 \longrightarrow X$, which is indeed in $W$ being $W$-aspheric! This is strange and almost immoral (\footnote{We observe below that this is so whenever $S$ has a final object.}). But by dint of scribbling I have just convinced myself that even if $X$ is cofibered over $S$ (hence also $W$-proper), $\Phi(X)$ is not $W$-proper (which would imply $W$-perfect), nor
\msp{page 188}
even a $W$-fibered one, which would mean that
\[ s'\backslash X \longrightarrow s\backslash X \qquad \text{associated to } s \longrightarrow s' \text{ in } S \]
would always be in $W$. Setting $Z' = s'\backslash S$, $Z = s\backslash S$, this would mean that if we have a \textit{local co-isomorphism}
\[ i : \underbrace{Z'}_{\simeq z_0\backslash Z} \longrightarrow Z, \]
where $Z'$, $Z$ have initial elements, and if $X$ is cofibered over $Z$, $X'$ its inverse image over $Z'$, then $X' \longrightarrow X$ is in $W$. In fact, if $Z$ also has a \textit{final} element $e$ (as was the case just now for $\boldsymbol{\Delta}^1$), then $Z' \simeq z_0\backslash Z$ also has one at $e'$, which maps to $e$, and we will have the commutative diagram
\[
\begin{tikzcd}
X' \arrow[r, "j"] & X \\
X'_{e'} \arrow[u, "W"] \arrow[r, "\sim"] & X_e, \arrow[u, "W"']
\end{tikzcd}
\]
whence $j \in W$. So to make a counterexample,
\msp{page 189}
we must take $Z$ without a final element.

But let us take $Z$ \texttt{[an]} ordered set, with an initial object $0$, but without a final object, and let us take as $z_0$ a \textit{maximal} element of $Z$, so that $Z' = z_0\backslash Z$ is isomorphic to $\{z_0\}$. We must give in Cat an example where the inclusion of $X_{z_0}$ in $X$ is not in $W$. We are going to take
\printedpage{117}
% head: A. Grothendieck : Dérivateurs, Ch. VI, Édition des Universités de Montpellier II et Paris VII
% foot: Version 9 novembre 2001
\[
\begin{tikzcd}
 & & a \\
Z = & 0 \arrow[ur] \arrow[dr] & \\
 & & b = z_0,
\end{tikzcd}
\]
so giving $X$ cofibered cosplit over $Z$ is equivalent to a diagram in Cat
\[
\begin{tikzcd}
 & X_a \\
X_0 \arrow[ur, "\alpha"] \arrow[dr, "\beta"] & \\
 & X_b.
\end{tikzcd}
\]
I think that under strong enough conditions (e.g. if $\alpha$, $\beta$ are open immersions, or closed immersions - to be checked!), the category $X$ is isomorphic in Hot to the amalgamated sum in Cat
\[ X_a{\textstyle\coprod_{X_0}}X_b, \]
and there \texttt{[is]} no reason (unless we assume $\alpha \in W$, and even then \dots) for $X_b \longrightarrow X_a \coprod_{X_0} X_b$ to be in $W$. Let us take $X_0 = X_b = e$, $\beta = \mathrm{id}_e$, so $\alpha$ is given by the datum of an element $t$ of $\mathrm{Ob}(X_a = T)$. The amalgamated sum is $T$, which is not isomorphic to $e$ in $\mathrm{Hot}_W$,
\msp{page 190}
if $T$ has not been taken $W$-aspheric. But I am going to try to show directly that the inclusion
\[ T = X_a \longrightarrow X \]
is $\in W$. For this I look at the open subcategory $Z''$ of $Z$, complementary to the full closed subcategory $Z'$, and I note that $X'' = X|Z''$ being proper over $Z''$ and $a$ being a final object of $Z''$, we have
\[ X_a \longrightarrow X'' \qquad \text{in } W, \]
it remains to see that the same holds for $X'' \longrightarrow X$. Now $X$ is obtained from $X''$ by adding a single element $\xi$ (the one of $X_b$), which does not map to any element of $X''$, and which receives no element of $X_a$ - on the other hand for the unique element $\xi_0$ of $X_0$ there is a unique arrow from it to $\xi$. We thus have a retraction of $X$ onto $X''$, sending $\xi$ to $\xi_0$ and the preceding arrow to $\mathrm{id}_{\xi_0}$. Is this a homotopism? It suffices to \texttt{[show]} that as an endomorphism $\varphi$ of $X$, it is $\sim \mathrm{id}_X$. Now we do have $\mathrm{id}_X \longleftarrow \varphi$, whose value on $X''$ is the identity, and the value on $\xi$ is the unique $\varphi(\xi) = \xi_0 \longrightarrow \xi$. This is a little brief, but it must be true. (In any case, $X$ is the union of two opens $X''$, and of the open $X|\{0,b\} = X'''$, with $X''' \simeq \boldsymbol{\Delta}^1$, and $X'' \cap X''' \simeq e$, so at least for the cohomology with locally constant coefficients of $X$ I am at ease by Mayer-Vietoris, and know full well that it is that of the piece $X''$, hence that of $X_a$ \texttt{[or $X_a$, illegible subscript]} not at all that of $e$!)
\DERmark{ENV0381}\DERmark{ENV0382}


\printedpage{118}
% head: A. Grothendieck : Dérivateurs, Ch. VI, Édition des Universités de Montpellier II et Paris VII
% foot: Version 9 novembre 2001
\msp{page 191}
\margbar We have thus seen at least this: If $X$ is cofibered hence $W$-proper over $S$, it by no means follows that $\Phi(X)$ is still $W$-proper over $S$ (\footnote{see a more precise result on the next page!}), nor that it is a $W$-fibered one. But what happens if $X \in \Prop_{0,W}S$, for example if it comes from $\Cofibsc_0 S$? Then $X$ is in $\Fib_W S$, so we already know that $\Phi(X)$ is in $\Fibsc_{0,W}S$, so it is not that way that we will make a counterexample. But let us again consider $X$ cofibered over
\[
\begin{tikzcd}
 & & a \\
S = & 0 \arrow[ur] \arrow[dr] & \\
 & & b,
\end{tikzcd}
\]
given by the diagram
\[
\begin{tikzcd}
 & X_a \\
X_0 \arrow[ur, "u"] \arrow[dr, "v"] & \\
 & X_b
\end{tikzcd}
\]
in Cat, assuming this time $u, v \in W$ ($X \in \Cofibsc_{0,W}S$). Then $\Phi(X)$ is the split fibered category over $S$ defined by the contravariant diagram on $S$
\[
\begin{tikzcd}
 & X_a \arrow[dl, hook, "i"'] \\
X & \\
 & X_b, \arrow[ul, hook, "j"']
\end{tikzcd}
\]
\DERmark{ENV0383}

where $i$, $j$ are the inclusions (which are closed immersions). To say that $\mathcal{X} = \Phi(X)$ is proper over $S$ also means that its restrictions over the open subcategories $\{0,a\}$ and $\{0,b\}$ of $S$ are so (for every arrow $\boldsymbol{\Delta}^1 \longrightarrow S$ factors through one or the other). This
\msp{page 192}
therefore also means that the inclusion of the closed fibers
\[ X_a \longrightarrow X, \qquad X_b \longrightarrow X \]
are coaspheric. But $X_a \longrightarrow X$ can be coaspheric only if $X_b = \emptyset$, for $x \in X_b \subseteq X$, $x\backslash X_a$ (as well as $X_a/x$) is empty, hence \textit{not} $W$-aspheric (unless we assume $W = \mathrm{Fl}(\mathrm{Cat})$, a case which we exclude). So $\Phi(X)$ is proper over $S$ \textit{only} if $X_a$ and $X_b = \emptyset$ (a result valid for \textit{every} $X$ over $S$, whether it is cofibered over $S$ or $W$-proper, or not). If we want $u, v \in W$, this therefore implies that $X_0 = \emptyset$, so we have proved:


\endgroup
% END UNIT Der-ChVI-345/Der-ChVI-345_p116-120.tex
% BEGIN UNIT Der-ChVI-345/Der-ChVI-345_p121-125.tex
\begingroup\makeatletter%
\NoAutoSpacing%
\newcommand{\printedpage}[1]{\marginpar{\footnotesize p.~#1}}%
\newcommand{\msp}[1]{\marginpar{\footnotesize [#1]}}%
\newcommand{\und}[1]{\underline{\mathrm{#1}}}%
\newcommand{\circnum}[1]{\tikz[baseline=(c.base)]\node[draw,circle,inner sep=1pt](c){#1};}%
\makeatother\let\printedpage\derpp\let\msp\dermsp%
\let\mspin\DERmspin%


\printedpage{119}
% head: A. Grothendieck : Dérivateurs, Ch. VI, Édition des Universités de Montpellier II et Paris VII
% foot: Version 9 novembre 2001

% the following paragraph is marked by a triple vertical rule in the left margin
\noindent\llap{$|||$\quad}Let $X \in \mathrm{Cat}/S$, $S = \begin{tikzcd}[row sep=tiny, column sep=small] & a \\ 0 \arrow[ur] \arrow[dr] & \\ & b \end{tikzcd}$. If $\Phi(X)$ is proper over $S$, then $X = \emptyset$ (if we assume $W \neq \mathrm{Fl}(\mathrm{Cat})$).

But along the way a bizarre idea appeared. Consider the infinite diagram, for $X$ in $\mathrm{Cat}/S$:
\[
X \longrightarrow \Phi(X) \longrightarrow \Psi\Phi(X) \longrightarrow \Phi\Psi\Phi(X) \longrightarrow \cdots
\]
If $X$ is arbitrary, these arrows are alternately in $W_S^{\mathrm{d}}$ and $W_S^{\mathrm{g}}$, so their composites are in $W_S$ and no more. But if $X$ is in $\und{Fib}_W S$, since this category is \textit{stable} under $\Phi$, $\Psi$, and since on $\und{Fib}_W S$ $W_S^{\mathrm{d}} = W_S^{\mathrm{g}} = W_S^{\mathrm{f}} = W_S^{\mathrm{univ}}$, we see that in this case the arrows are all in $W_S^{\mathrm{f}}$. We now feel like setting
\[
\Phi_\infty(X) = \varinjlim \text{ of all that}
\]
\msp{page 193}
and we would define symmetrically
\[
\Psi_\infty(X) = \varinjlim (X \longrightarrow \Psi(X) \longrightarrow \Phi\Psi(X) \longrightarrow \cdots).
\]

For this to be usable, we must postulate the axiom \underline{W(6 bis)} (compare XII, pages 14, 15). The direct limit of a sequence of morphisms $f_i : X_i \longrightarrow Y_i$, for a morphism of direct systems in Cat indexed by $\mathbf{N}$, is in $W$ if the $f_i$ are.

This axiom is close to W(6) of loc. cit., but is implied by it only with W(4b). It is stronger in the case of direct systems indexed by $\mathbf{N}$ (so also, more generally, of direct systems on filtering ordered sets having a countable cofinal set), but says nothing about direct systems on more general ordered index sets.

This axiom also implies that a sequential filtering direct limit of morphisms $f_i$ is $W$-smooth, resp. $W$-proper, if the $f_i$ are. We then find that $\Phi_\infty(X)$ is both $W$-smooth (as a direct limit of $\Phi(X)$, $\Phi\Psi\Phi(X)$, \ldots, $\Phi(\Psi\Phi)^n(X)$, which are $W$-smooth over $S$) and $W$-proper (as a direct limit of $\Psi\Phi(X)$, \ldots, $(\Psi\Phi)^n(X)$ \ldots). On the other hand, that $X \longrightarrow \Phi_\infty(X)$ is in $W_S^{\mathrm{f}}$ - for $W_S^{\mathrm{f}}$ is stable under countable filtering direct limits in $\mathrm{Cat}/S$. We have thus found a functor
\DERmark{ENV0384}

\msp{page 194}
\[
\Phi_\infty : \und{Fib}_W S \longrightarrow \und{Parf}_W S \subseteq \und{Fib}_W S
\]
and a functorial morphism
\[
i_X : X \longrightarrow \Phi_\infty(X), \qquad \boxed{i_X \in W_S^{\mathrm{f}}} \qquad \text{if } X \in \und{Fib}_W S.
\]
% the box around i_X \in W_S^f is printed with rounded corners

This functor then allows us, in the absence of W(4b), to do what the functor $C$ absolutely refused. It is true that $C$ has the advantage of being defined on all of $\mathrm{Cat}/S$, and of being
\printedpage{120}
% head: A. Grothendieck : Dérivateurs, Ch. VI, Édition des Universités de Montpellier II et Paris VII
% foot: Version 9 novembre 2001
reasonable and even very useful there, when we have W(4b). But $\Phi_\infty$ too is defined on all of $\mathrm{Cat}/S$, as is $i_X$
\[
\begin{aligned}
&\Phi_\infty : \mathrm{Cat}/S \longrightarrow \und{Parf}_W S\\
&i : \mathrm{id} \longrightarrow \Phi_\infty \qquad \text{i.e.}\ \begin{cases} i_X : X \longrightarrow \Phi_\infty(X)\\ i_X \in W_S, \end{cases}
\end{aligned}
\]
with the sole handicap (inevitable) that if $X$ is not in $\und{Fib}_W S$, we can say no better than $i_X \in W_S$, to the exclusion of $i_X \in W_S^{\mathrm{f}}$. But for $C(X)$ it was necessarily the same. So we get the impression that $\Phi_\infty$, take your pick, $\Psi_\infty$, six of one, half a dozen of the other, plays in the presence of W(6 bis) the same role as C in the presence of W(4b). At most (in the presence of both axioms at once) the functor $C$ has the aesthetic advantage of taking its values in $\und{Parf}_\omega S$, instead of merely $\und{Parf}_W S$ (\footnote[93]{\textit{In fact} $\Phi_\infty = C$, if $C$ is defined simply via one-sided infinite paths. If we define $C$ via \textit{finite} paths of all types $c$, then it does indeed take its values in $\und{Parf}_\omega S$.}). But it is imperative to % [sic] "que C en présence" (C in roman type) -> $C$, italic everywhere else
\msp{page 195}
realize now that $\Phi_\infty$ has an advantage of exactly the same order over $C$, by introducing
\[
W''_\omega = \text{smallest subset } \Sigma \text{ of } \mathrm{Fl}(\mathrm{Cat}) \text{ satisfying } \mathrm{W}(1, 2\text{ bis}, 3, 6\text{ bis})
\]
and note that
\[
\Phi_\infty(X) \in \mathrm{Ob}\,\und{Parf}_{W''_\omega} S, \qquad i_X \in W''_\omega % [sic] i_x -> i_X, the morphism is i_X everywhere else
\]
\DERmark{ENV0385}

- but here I notice that we have
\[
W_\omega \subseteq W''_\omega,
\]
so $C$ is still a bit better in this respect \ldots

When we have $\Phi_\infty$ with the properties we have just made explicit, we find that the inclusion
\[
\und{Parf}_W S \longrightarrow \und{Fib}_W(S)
\]
induces an equivalence not only for the categories localized for $W_S$, but also for the fine localized categories
\[
(\und{Parf}_W S)(W_S^{\mathrm{f}})^{-1} \longrightarrow (\und{Fib}_W(S))(W_S^{\mathrm{f}})^{-1} \overset{\mathrm{def}}{=} \mathrm{HOT}^{\mathrm{lc}}_W(S)
\]
\DERmark{ENV0386}

(\footnote[94]{We can even replace $\und{Parf}_W S$ there by $\und{Parf}_{W''_\omega} S$, \textit{or} by any other full subcategory of $\und{Fib}_W S$ containing $\und{Parf}_{W''_\omega} S$.}). But we cannot for all that reduce the number of useful localized categories to three, in the absence of W(4b). At most, having a functor $\Phi_\infty$ which is simpler than $C$, in the absence of W(4b), because we know in any case that
\[
i_X : X \longrightarrow \Phi_\infty(X)
\]
is in $W_S^{\mathrm{f}}$ if $X$ is in $\und{Fib}_W S$, can we
\printedpage{121}
% head: A. Grothendieck : Dérivateurs, Ch. VI, Édition des Universités de Montpellier II et Paris VII
% foot: Version 9 novembre 2001
\msp{page 196}
hope to get by with the weakened variant of W(4b) obtained by restricting ourselves in the statement of W(4b) (p. 171) to the case where $X$, $Y$ are in $\und{Parf}_{W''_\omega} S$. This time, that does indeed suffice to imply that the functor
\[
\lambda : (\und{Fib}_W S)(W_S^{\mathrm{f}})^{-1} \longrightarrow (\und{Fib}_W S)(W_S)^{-1}
\]
is an equivalence of categories, since we can ``compute'' it just as well using $\und{Parf}_{W''_\omega} S$ instead of $\und{Fib}_W S$, now the hypothesis W(4b($\mathrm{W}''_\omega$)) implies that on $\und{Parf}_{W''_\omega} S$ the two localizers $W_S^{\mathrm{f}}$, $W_S$ are \textit{equal.}

We need a usable practical conclusion from all these perplexities and observations. Let us try

\textbf{Proposition.} \textit{Suppose that we are under one of the following conditions.}

\begin{enumerate}[label={}]
\item[$1^o$)] \textit{The axiom} W(4b) \textit{is satisfied in the case of categories $X$, $Y$ which are in} $\und{Fibsc}_{0,W} S$ (\footnote[95]{but this is \textit{equivalent} to W(4b)! Cf. XVI 4.9. loc. cit.}) \textit{(this condition on $W$ is in fact self-dual).}

\item[$2^o$)] \textit{The axiom} W(4b) \textit{is satisfied in the case of categories $X$, $Y$ which are in} $\und{Parf}_{W''_\omega} S$\textit{, where $W''_\omega$ is the smallest set $\Sigma \subseteq$ Fl}(Cat) \textit{satisfying the conditions} $\mathrm{W}(1, 2\text{ bis}, 3, 6\text{ bis})$. \textit{(We can also make do with the condition, stronger but more understandable, that} W(4b) \textit{is OK if $X$, $Y$ are in} $\und{Parf}_W S$ (\footnote[96]{But this also implies W(4b), cf. loc. cit.}).\textit{) Moreover, we assume that $W$ satisfies} W(6 bis).
\end{enumerate}

\textit{Under these conditions, the corresponding [?] functor $\lambda$}
\[
\mathrm{HOT}^{\mathrm{lc}}_W(S) \longrightarrow \mathrm{HOT}^{\mathrm{lc!}}_W(S)
\]
\textit{is an equivalence of categories (cf. diagram p. 178), so $(\rho, i)$, $(\sigma, j)$ can be regarded as pairs of functors between two categories (cf. $(*)$ p. 181).} These are pairs of adjoint functors, \textit{so $i$, $j$ are fully faithful, and we have the adjunction formula $(*)$, page 182.}

\msp{page 197}
To prove the adjunction of $(\rho, i)$ (or dually of $(\sigma, j)$), we again express the two categories at hand $\mathrm{HOT}_W(S)$, $\mathrm{HOT}^{\mathrm{lc}}_W(S)$ as $(\mathrm{Cat}/S)(W_S^{\mathrm{d}})^{-1}$, $(\mathrm{Cat}/S)(W_S)^{-1}$, the functor $\rho$ being the canonical localization going from the first to the second, and the functor $i$ being given by an inversion functor
\[
\gamma : (\mathrm{Cat}/S, W_S) \longrightarrow (\mathrm{Cat}/S, \overline{W_S^{\mathrm{d}}})
\]
which is $C$ in case $1^o$, $\Phi_\infty$ in case $2^o$. The adjunction morphisms
\[
\rho i \longrightarrow \mathrm{id}, \qquad i\rho \longleftarrow \mathrm{id}
\]
should then be described as in the case where we simply assume W(4b), and the compatibilities are seen in the same way: we are reduced to proving that two arrows
\printedpage{122}
% head: A. Grothendieck : Dérivateurs, Ch. VI, Édition des Universités de Montpellier II et Paris VII
% foot: Version 9 novembre 2001
\[
\begin{tikzcd}
\gamma(X) \arrow[r, shift left, "i_{\gamma(X)}"] \arrow[r, shift right, "\gamma(i_X)"'] & \gamma(\gamma(X))
\end{tikzcd}
\]
are equal in $(\mathrm{Cat}/S)(W_S^{\mathrm{d}})^{-1}$. The verification in case $1^o$), where $\gamma = C$, should not be changed. In case $2^o$) where $\gamma = \Phi_\infty$, I have not done it. I will do it only under the pressure of need.

\subsection*{\circnum{6} Contravariant law of $\mathrm{HOT}_W(S)$, $\mathrm{HOT}^{\mathrm{lc}}_W(S)$.}

There is for
\[
f : S' \longrightarrow S
\]
a homomorphism from the diagram
\[
\begin{tikzcd}[column sep=small]
& \mathrm{HOT}^{\mathrm{lc}}_W(S') \arrow[dl] \arrow[dr] & & \text{\texttt{[Rather}} & & \mathrm{HOT}^{\mathrm{lc}}_W(S) \arrow[dl] \arrow[dr] & \\
\mathrm{HOT}_W(S') & & \mathrm{HOT}^{o}_W(S') & & \mathrm{HOT}_W(S) & & \mathrm{HOT}^{o}_W(S)\ \text{\texttt{]}}
\end{tikzcd}
\]
\msp{page 198}
into the similar diagram for $S$ \texttt{[rather} $S'$\texttt{]}, by functors $f^*$ which are defined respectively on
\[
\begin{aligned}
\und{Liss}_W S &\quad \text{for}\quad \mathrm{HOT}_W(S) & \quad\longrightarrow\quad & \mathrm{HOT}_W(S')\\
\und{Prop}_W S &\quad \text{for}\quad \mathrm{HOT}^{o}_W(S) & \quad\longrightarrow\quad & \mathrm{HOT}^{o}_W(S')\\
\und{Fib}_W S &\quad \text{for}\quad \mathrm{HOT}^{\mathrm{lc}}_W(S) & \quad\longrightarrow\quad & \mathrm{HOT}^{\mathrm{lc}}_W(S'),
\end{aligned}
\]
through the functors % [sic] intermédiare -> intermédiaire, missing i
\[
f^* : X \longmapsto X \times_S S' : \left\{
\begin{aligned}
\und{Liss}_W S &\quad\longrightarrow\quad \und{Liss}_W S'\\
\und{Prop}_W S &\quad\longrightarrow\quad \und{Prop}_W S'\\
\und{Fib}_W S &\quad\longrightarrow\quad \und{Fib}_W S',
\end{aligned}
\right.
\]
which functor is compatible with the relevant localizers, which are $W_S^{\mathrm{f}}$ in all cases. On the other hand, in the absence of draconian hypotheses on $W$ or on $S$, we cannot say that $f^*$ is compatible with $W_S$, $W_{S'}$ on any of the categories of the diagram p. 175, allowing us to define $\mathrm{HOT}^{\mathrm{lc!}}_W(S)$, so as to give us a functor % [sic] pretinents -> pertinents, transposed letters
\[
\mathrm{HOT}^{\mathrm{lc!}}_W(S) \longrightarrow \mathrm{HOT}^{\mathrm{lc!}}_W(S').
\]
In the absence of axioms on $W$, the only case known to me is that where $f$ is \textit{a universal $W$-equivalence,} i.e. a $W$-fibration with $W$-aspheric fibers, in which case it is obvious that for \textit{every} arrow $u : X \longrightarrow Y$ in $\mathrm{Cat}/S$, $u \in W_S$ if and only if $(u' : X' \longrightarrow Y') \in W_{S'}$. We then find a homomorphism from the complete diagram $(*)$ p. 176 \texttt{[rather p. 178]} for $S$, into the corresponding diagram for $S'$.

If we assume W(4b), then $W_S = W_S^{\mathrm{f}}$ on $\und{Fib}_W S$, so on that category the functor $f^*$ is compatible (no merit!) with the localizers $W_S$, $W_{S'}$, but this is not really new
\printedpage{123}
% head: A. Grothendieck : Dérivateurs, Ch. VI, Édition des Universités de Montpellier II et Paris VII
% foot: Version 9 novembre 2001
\msp{page 199}
compared with what was said above, for $\und{Fib}_W S$, $\und{Fib}_W S'$, and for the $\mathrm{HOT}^{\mathrm{lc}}_W$.

But if we assume W(4a), and $f \in \und{Fib}_W$, then $f^*$ on all of $\mathrm{Cat}/S$ is compatible with the localizers $W_S$, $W_{S'}$, and we find a computation of $f^*$ on $\mathrm{HOT}^{\mathrm{lc!}}_W(S) \approx \mathrm{HOT}^{\mathrm{lc}}_W(S)$ by
\[
f^*(\mathrm{hot}^{\mathrm{lc}}_{S,W}(X)) = \mathrm{hot}^{\mathrm{lc}}_{S',W}(X \times_S S').
\]

So this is a case, exceptionally, where the functors $f^*$ for the $\mathrm{HOT}_W$ are compatible not only with the inclusion functors $i$, $j$ (of page 176 \texttt{[rather p. 178]}), but also with the localization functors $\rho$, $\sigma$. % [sic] exceptionellement -> exceptionnellement, French spelling has double n

\textbf{Omission:} If $p : S' \longrightarrow S$ is $W$-proper, then $p^* : \mathrm{Cat}/S \longrightarrow \mathrm{Cat}/S'$ is compatible with $W_S^{\mathrm{d}}$, $W_{S'}^{\mathrm{d}}$, and $p^* : \mathrm{HOT}_W(S) \longrightarrow \mathrm{HOT}_W(S')$ can be ``computed'' by the formula
\[
p^*(\mathrm{hot}_{S,W}(X)) = \mathrm{hot}_{S',W}(X \times_S S')
\]
for \textit{every} $X$ in $\mathrm{Cat}/S$.

\subsection*{\circnum{7} Product in $\mathrm{HOT}_W(S)$.}

The functor
\[
(X, Y) \longmapsto X \times_S Y : \und{Liss}_W S \times \und{Liss}_W S \longrightarrow \und{Liss}_W S
\]
is compatible with the localizers $W_S^{\mathrm{f}} \times W_S^{\mathrm{f}}$ and $W_S^{\mathrm{f}}$, so passes to the localized categories as
\[
\left\{
\begin{aligned}
\mathrm{HOT}_W(S) \times \mathrm{HOT}_W(S) &\quad\longrightarrow\quad \mathrm{HOT}_W(S)\\
(\xi, \eta) &\quad\longmapsto\quad \xi \times_S \eta
\end{aligned}
\right.
\]
(\footnote[97]{\textbf{NB} If $X$, $Y$ are \textit{arbitrary} in $\mathrm{Cat}/S$, we have $\mathrm{hot}_{S,W}(X) \times_S \mathrm{hot}_{S,W}(Y) = \mathrm{hot}_{S,W}(\Phi_S(X) \times_S \Phi_S(Y))$. If $X$ or $Y$ is smooth over $S$, we can replace $\Phi(X)$ resp. $\Phi(Y)$ there by $X$ resp. $Y$.}). This product is associative, unital, commutes with all the usual compatibilities, as follows from the fact that that's how it is in $\und{Liss}_S$. The unit element is the class of $S \xrightarrow{\mathrm{id}} S$. The subcategory $\mathrm{HOT}^{\mathrm{lc}}_W(S)$ is stable under the product, which induces
\msp{page 200}
\[
\mathrm{HOT}^{\mathrm{lc}}_W(S) \times \mathrm{HOT}^{\mathrm{lc}}_W(S) \longrightarrow \mathrm{HOT}^{\mathrm{lc}}_W(S),
\]
itself induced by the product functor $(X, Y) \longrightarrow X \times_S Y$ in $\und{Fib}_W S$.

If $X$ is arbitrary in $\mathrm{Cat}/S$, the functor
\[
\mathrm{mult}_{X/S} : Y \longmapsto X \times_S Y : \und{Liss}_W S \longrightarrow \mathrm{Cat}/S
\]
is compatible with the localizers $W_S^{\mathrm{d}}$, for this functor can be written as the composite of the corresponding functor \textit{with values in $\und{Liss}_W X$,} which is compatible with the localizers % [sic] correspondent -> correspondant, adjective intended


\endgroup
% END UNIT Der-ChVI-345/Der-ChVI-345_p121-125.tex
% BEGIN UNIT Der-ChVI-345/Der-ChVI-345_p126-129.tex
\begingroup\makeatletter%
\newcommand{\printedpage}[1]{\marginpar{\footnotesize p.~#1}}%
\newcommand{\msp}[1]{\marginpar{\footnotesize [#1]}}%
\newcommand{\cercle}[1]{\tikz[baseline=(c.base)]\node[draw,circle,inner sep=1pt](c){#1};}%
\makeatother\let\printedpage\derpp\let\msp\dermsp%
\let\mspin\DERmspin%


\setcounter{footnote}{97}

\printedpage{124}
% head: A. Grothendieck : Dérivateurs, Ch. VI, Édition des Universités de Montpellier II et Paris VII
% foot: Version 9 novembre 2001

$W_S^{\mathrm{d}}$, $W_X^{\mathrm{d}}$, and the ``forget $X$'' functor from $\underline{\mathrm{Liss}}_W X$ (even from $\mathrm{Cat}/X$) to $\mathrm{Cat}/S$, which is compatible with $W_X^{\mathrm{d}}$ and $W_S^{\mathrm{d}}$. We thus get a functor % [sic] W_S^d, W_S^d -> apparently W_S^d, W_X^d (or W_X^d, W_S^d), the same symbol is printed twice
\DERmark{ENV0387}

\[
\left\{\begin{aligned}
&\mathrm{mult}_{X/S} : \mathrm{HOT}_W(S) \longrightarrow \mathrm{HOT}_W(S)\\
&\mathrm{mult}_{X/S}(\mathrm{hot}_{S,W}(Y)) \overset{\mathrm{def}}{=} \mathrm{hot}_{S,W}(X\times_S Y) \qquad \text{if $Y$ \textit{smooth} over $S$}
\end{aligned}\right.
\]
\DERmark{ENV0388}

(\footnote{This formula is also valid if $Y$ is in $\underline{\mathrm{Fib}}_W\, S$, cf. $(*)$, p. 202.}), since $\mathrm{HOT}_W(S)$ can be obtained as the localized category of both $\underline{\mathrm{Liss}}_W S$ and $\mathrm{Cat}/S$, for $W_S^{\mathrm{d}}$. When $X$ is itself $W$-smooth over $S$, this functor is none other than
\[
\left\{\begin{aligned}
&\mathrm{mult}_\xi : \eta \longmapsto \xi\times_S \eta,\\
&\text{where }\ \xi = \mathrm{hot}_{S,W}(X).
\end{aligned}\right.
\]

But this is not so if $X$ is not smooth, i.e. we must distinguish in this case, even for $Y$ $W$-smooth over $S$,
\[
\mathrm{hot}_{S,W}(X\times_S Y) \qquad\text{and}\qquad \mathrm{hot}_{S,W}(X)\times_S \mathrm{hot}_{S,W}(Y).
\]

The second term can be made explicit as $\mathrm{hot}_{S,W}(\Phi(X)\times_S Y)$ which receives $\mathrm{hot}_{S,W}(X\times_S Y)$, i.e.
\msp{page 201}
we have a canonical homomorphism
\[
\underbrace{\mathrm{hot}_{S,W}(X\times_S Y)}_{=\,\mathrm{mult}_{X/S}(Y)} \longrightarrow \underbrace{\mathrm{hot}_{S,W}(X)\times_S \mathrm{hot}_{S,W}(Y)}_{=\,\mathrm{hot}_{S,W}(\Phi(X)\times_S Y)},
\]
induced by
\[
X\times_S Y \longrightarrow \Phi(X)\times_S Y,
\]
but this last arrow has no reason to be in $W_S^{\mathrm{d}}$, and thus does not necessarily induce an isomorphism in $\mathrm{HOT}_W(S)$. (Counterexample page 154.)

We must look separately at the case \textit{where $X$ is in} $\underline{\mathrm{Fib}}_W S$, I feel like saying that in this case we have
\[
\mathrm{hot}_{S,W}(X\times_S Y) \xrightarrow{\ \sim\ } \underbrace{\mathrm{hot}_{S,W}(X)\times_S \mathrm{hot}_{S,W}(Y)}_{=\,\mathrm{hot}_{S,W}(\Phi(X)\times_S\, \Phi(Y))} \tag*{$(*)$} % doubt resolved in DER_VI-c_corrections.csv; input note: under the brace "hot" is printed larger than "(Phi(X)x_S Phi(Y))", which stands at subscript level -- argument of hot_{S,W} or part of its subscript
\]
(\footnote{This formula is established only by assuming either $Y$ $W$-smooth (in addition to $X$ $W$-fibered over $S$), or $X$ $W$-proper, e.g. $X$ $W$-perfect.}) for \textit{every} $Y$ in $\mathrm{Cat}/S$ (not necessarily $W$-smooth), a fortiori that in this case
\[
\left\{\begin{aligned}
&\mathrm{mult}_{X/S} = \mathrm{mult}_\xi \qquad \text{in } \mathrm{HOT}_W(S)\\
&\text{where } \xi = \mathrm{hot}_{S,W}(X/S) \text{ (if $X$ in $\underline{\mathrm{Fib}}_W S$).} % [sic] hot_{S/W} -> hot_{S,W}, the notation used everywhere else
\end{aligned}\right. \tag*{$(**)$}
\]
\DERmark{ENV0389}


\printedpage{125}
% head: A. Grothendieck : Dérivateurs, Ch. VI, Édition des Universités de Montpellier II et Paris VII
% foot: Version 9 novembre 2001

In the case where $Y$ is smooth over $S$ (which suffices to establish $(**)$), it suffices to recall that $X\longrightarrow \Phi(X)$ is in $W_S^{\mathrm{f}}$ (follows from the fact that $X$ is $W$-fibered over $S$ (\footnote{even $W$-presmooth over $S$ suffices}), hence $X_s \longrightarrow s\backslash X$ is in $W$, $\forall\ s\in S$), hence $\underbrace{X\times_S Y}_{X_Y} \longrightarrow \underbrace{\Phi(X)\times_S Y}_{\Phi(X)_Y}$ is also in $W_Y^{\mathrm{f}}$, hence (since $X_Y$, $\Phi(X)_Y$ are $W$-fibered over $Y$) in $W_Y^{\mathrm{d}}$, hence also in $W_S^{\mathrm{d}}$. In fact, even without assuming $Y$ smooth, this shows that we have
\msp{page 202}
\[
\mathrm{hot}_{S,W}(X\times_S Y) \xrightarrow{\ \sim\ } \mathrm{mult}_{Y/S}(\mathrm{hot}_{S,W}(X)) \qquad \text{if $X$ is in $\underline{\mathrm{Fib}}_W S$.} \tag*{$(*)$}
\]

To have $(*)$ for arbitrary $Y$, we can by the foregoing assume $X$ smooth (replacing it by $\Phi(X)$ if need be), hence $X$ in $\underline{\mathrm{Liss}}_{0,W}S$. We must see that
\[
X\times_S Y \longrightarrow X\times_S \Phi(Y)
\]
is still in $W_S^{\mathrm{d}}$. Now if $X$ is \textit{$W$-proper} over $S$, the base change $X\longrightarrow S$ is compatible with $W_S^{\mathrm{d}}$, $W_X^{\mathrm{d}}$ (on all of $\mathrm{Cat}/S$), and we win. (This must be false without an assumption of $W$-properness on $X$, \textit{or} of $W$-smoothness on $Y$ \dots)

When we have
\[
f : S' \longrightarrow S,
\]
then
\[
f^* : \mathrm{HOT}_W(S) \longrightarrow \mathrm{HOT}_W(S')
\]
is compatible with the product structures.

\subsection*{\protect\cercle{8} The functors $f_!$: adjunction theorem, properties of the derivator $\mathrm{HOT}_W$.}

Let
\[
f : S' \longrightarrow S
\]
be in Cat, then the induced functor
\[
f_! : \mathrm{Cat}/S' \longrightarrow \mathrm{Cat}/S
\]
is compatible with the localizers $(W_{S'}^{\mathrm{d}}, W_S^{\mathrm{d}})$, or $(W_{S'}^{\mathrm{g}}, W_S^{\mathrm{g}})$, or $(W_{S'}, W_S)$. We thus get an induced functor for the localized categories, and \texttt{[a]} commutative diagram
\msp{page 203}
\[
\begin{tikzcd}[column sep=small]
\mathrm{HOT}_W(S') \arrow[rrd, "\rho'"] \arrow[dd, "f_!"'] & & & & \mathrm{HOT}^{o}_W(S') \arrow[lld, "\sigma'"'] \arrow[dd, "f^{o}_!"] \\
& & \mathrm{HOT}^{\mathrm{lc}\,!}_W(S') \arrow[dd, "f^{\mathrm{lc}\,!}_!"] & & \\
\mathrm{HOT}_W(S) \arrow[rrd, "\rho"] & & & & \mathrm{HOT}^{o}_W(S). \arrow[lld, "\sigma"'] \\
& & \mathrm{HOT}^{\mathrm{lc}\,!}_W(S) & &
\end{tikzcd}
\]

\printedpage{126}
% head: A. Grothendieck : Dérivateurs, Ch. VI, Édition des Universités de Montpellier II et Paris VII
% foot: Version 9 novembre 2001

We do not have a functor $f_!^{\mathrm{lc}}$ at the level of the $\mathrm{HOT}_W^{\mathrm{lc}}$, unless we make an assumption, such as \underline{W(4b)}, which ensures that $\mathrm{HOT}_W^{\mathrm{lc}\,!} \approx \mathrm{HOT}_W^{\mathrm{lc}}$.

The main theorem, that the functor $f_!$ is left adjoint to $f^*$, is valid \textit{under the sole assumption} W$(1,2,3)$ (\footnote{i.e. with these new notations Loc$(1,2,3\ \text{bis})$ (cf. XVI).}). Thus dually for $f_!^{o}$, $f_{o}^*$. But since we do not have a functor $f_{\mathrm{lc}\,!}^*$, the question of adjunction does not arise for general $f$, and without an assumption W(4b) on $W$. It only arises, apparently, for $f$ a trivial $W$-fibration (cf. end of \cercle{6}, p. 198), and in this case the adjunction should not present any difficulties. But as long as I do not have a very clear feeling about the meaning of the $\mathrm{HOT}_W^{\mathrm{lc}\,!}$, this is not a result to get much of a hard-on over \texttt{[sic]}.

The adjunction of $f_!$, $f^*$ for $\mathrm{HOT}_W$ already allows us to show that the prederivator $\mathrm{HOT}_W$
\msp{page 204}
satisfies the conditions Der 1, Der $4'$, Der $5'$ for derivators. Moreover, we have what we can call the property Der 3b for $\mathrm{HOT}_W$ as a $W$-derivator, i.e.

\begin{quote}
\textit{If $f\in W$, then $f$ is a} $\mathrm{HOT}_W$\textit{-equivalence} (\footnote{We prove this by interpreting the notion of a $\mathrm{HOT}_W$-equivalence in terms of $\mathrm{HOT}_W$\textit{-homology.} This comes down in the end to this: If $f : S' \longrightarrow S$ is in $W$, then $\forall\, T$, $S'\times T \longrightarrow S\times T$ too!}) (assuming only the axioms W$(1,2,3)$ on $W$). More precisely, $f$ is a $\mathrm{HOT}_W$-equivalence if and only if $\mathrm{hot}_W(f)$ is \texttt{[an]} isomorphism, i.e. $f\in \overline{W}$, the strong saturation of $W$.
\end{quote}

If $W$ \textit{satisfies} W(4b)\textit{,} then we also have a pair of adjoint functors
\[
\begin{tikzcd}
\mathrm{HOT}_W^{\mathrm{lc}}(S') \arrow[r, shift left, "f_!^{\mathrm{lc}}"] & \mathrm{HOT}_W^{\mathrm{lc}}(S) \arrow[l, shift left, "f^*_{\mathrm{lc}}"]
\end{tikzcd}
\]
for every arrow $S'\longrightarrow S$ in Cat. We also see in this case that % [sic] pout -> pour, misprint

\begin{quote}
% printed with a vertical rule at the left of the block
\begin{tabular}{|p{0.9\linewidth}}
if $f : S' \longrightarrow S$ is in $W$, then $f^*_{\mathrm{lc}}$ and $f_!^{\mathrm{lc}}$ between $\mathrm{HOT}_W^{\mathrm{lc}}(S)$ and $\mathrm{HOT}_W^{\mathrm{lc}}(S')$ are equivalences quasi-inverse to one another (assuming W(4b)),
\end{tabular}
\end{quote}

which is axiom Der 3b$_2$ ``$W$-derivators version''. We cannot expect any better, on the Der 3 side!

\textbf{Omission:} Let $f : S' \longrightarrow S$, we then have $\boxed{f_!f^*} : \mathrm{HOT}_W(S) \longrightarrow \mathrm{HOT}_W(S)$ is equal to $\mathrm{mult}_{S'/S}$ (cf. p. 200). It is therefore equal to $\mathrm{mult}_\xi$, where $\xi = \mathrm{hot}_{S,W}(S')$, if we assume either $S'$ $W$-smooth over $S$, or $S'$ $W$-fibered over $S$ (cf. $(**)$, p. 201).

\subsection*{\protect\cercle{9} Base change theorem for $f_!$ in $\mathrm{HOT}_W$.}

If we have
\printedpage{127}
% head: A. Grothendieck : Dérivateurs, Ch. VI, Édition des Universités de Montpellier II et Paris VII
% foot: Version 9 novembre 2001
\[
\begin{tikzcd}
X \arrow[d, "f"'] & X' \arrow[l, "q"'] \arrow[d, "f'"] \\
S & S', \arrow[l, "g"']
\end{tikzcd}
\]
with $f$ $W$-smooth \textit{or} $g$ $W$-proper \texttt{[the continuation is missing]}


\endgroup
% END UNIT Der-ChVI-345/Der-ChVI-345_p126-129.tex
\clearpage\setcounter{footnote}{0}\phantomsection\addcontentsline{toc}{section}{Chapter VII}\markboth{Chapter VII}{Chapter VII}
% BEGIN UNIT Der-ChVII-456/Der-ChVII-456_p001-005.tex
\begingroup\makeatletter%
\newcommand{\printedpage}[1]{\marginpar{\footnotesize p.~#1}}%
\newcommand{\msp}[1]{\marginpar{\footnotesize [#1]}}%
\newcommand{\Fltr}{\operatorname{Fl}\operatorname{tr}}%
\newcommand{\Ob}{\operatorname{Ob}}%
\newcommand{\bDelta}{\boldsymbol{\Delta}}%
\makeatother\let\printedpage\derpp\let\msp\dermsp%
\let\mspin\DERmspin%


% ===================== page 1 =====================
\printedpage{}
% head: (none; title page, no page number printed)

\begin{center}
{\Large\textsf{THE DERIVATORS}}

\bigskip

\textbf{ALEXANDRE GROTHENDIECK}

\bigskip\bigskip

\textbf{Chapter VII}

\medskip

\textbf{Categories of paths (2)}
\end{center}

\par

% ===================== page 2 =====================
\printedpage{}
% head: (none; no page number printed)

\par\noindent\textbf{Note by the editors of the transcribed edition.}\par
This text was deciphered and transcribed in \LaTeXe\ by M. Künzer. It was edited by M. Künzer, J. Malgoire, and G. Maltsiniotis. The transcription is as faithful as possible to the manuscript. For the few obvious corrections, or rare comments by the editors, as well as for the original numbering of the manuscript pages, typewriter characters \texttt{[typewriter]} in brackets are used. A question mark in brackets means that one is not sure of the preceding word.

This edition is provisional. Remarks, comments and corrections are welcome. Send a message to:

\textbf{[e-mail address removed]}

\hfill G. Maltsiniotis\hspace*{3em}

\par

% ===================== page 3 =====================
\printedpage{1}
% head: A. Grothendieck : Dérivateurs, Ch. VII, Édition des Universités de Montpellier II et Paris VII
% foot: Version 22 décembre 2001

\msp{page 1}
Since the ``simplified'' formalism of paths that I wanted to work out in a hurry has turned out to be unbearably cumbersome, enough to make it virtually unusable, I am going to try to work out a flexible formalism, sticking to intuition, at the cost only of an additional conceptual effort.

\section*{1\quad Intervals and their variants} % [sic] leur variantes -> leurs variantes, plural agreement

We call a \textit{standard interval category} a category defined by an ordered set, whose set of elements is an interval of integers $[0,n]$ ($n\in\mathbf{N}$, we do not exclude the case $n=0$, i.e. of the one-point category $\{0\}$): % [sic] élements -> éléments, missing accent
\[
\Ob I=[0,n]\subseteq\mathbf{N},
\]
and satisfying the following hypothesis:

\begin{enumerate}
\item[a)] for every pair of two consecutive vertices $i-1$, $i$ of $I$ ($1\leq i\leq n$), there is an arrow joining them, thus either $i\longrightarrow i+1$, or $i\longleftarrow i+1$ \texttt{[rather} $i-1\longrightarrow i$\texttt{,} $i-1\longleftarrow i$\texttt{]} (this arrow is denoted by $\underline{u}_i$,
\[
(i-1)\xrightarrow{\underline{u}_i}i\qquad\quad\text{or}\qquad\quad(i-1)\xleftarrow{\underline{u}_i}i\ \ ).
\]
\item[b)] The category $I$ is generated by the set of its arrows $\underline{u}_i$ (called \textit{transition arrows}).
\end{enumerate}

The transition arrow is called \textit{direct} if it goes from $i-1$ to $i$, \textit{retrograde} in the opposite direction. If we assign to the interval $[i-1,i]$ the sign $+$ or $-$ according as $\underline{u}_i$ is direct or retrograde, we thus see that the
\msp{page 2}
standard intervals of given length $n$ correspond one-to-one to the systems of signs $+$ or $-$ on the $n$ transition intervals $[i-1,i]$ ($1\leq i\leq n$), so there are $2^n$ of them. The vertex $0$ is the \textit{origin} or the \textit{source} of the interval, the vertex $n$ the \textit{endpoint} or the \textit{target} of the interval.

A standard interval is \textit{oriented,} in that we have the natural \textit{total} order relation on the set of its vertices, namely that of the integers in $[0,n]$. When we speak of the order relation on $I$, that a vertex $i$ \textit{``precedes''} or \textit{``follows''} another one $j$ ($i\leq j$ or $i\geq j$ respectively), it is this order relation that we will mean. The one associated with the category structure will never be regarded as an order relation. There is also a natural total order relation on the set of transition arrows $\underline{u}_i$ ($1\leq i\leq n$), given by the order relation of the integers $i$, $1\leq i\leq n$, that index them. We will say that a transition arrow $i$ \texttt{[rather} $\underline{u}_i$\texttt{]} \textit{precedes} or \textit{follows} another \texttt{[}$\underline{u}_j$\texttt{]}, or that it is the \textit{immediate successor} or \textit{immediate predecessor} of another (if $i=j+1$ or $i=j-1$). % [sic] qu'on à -> qu'on a, verb; [sic] un sommet i en ‘précède’ -> un sommet i ‘précède’, superfluous "en"

% Reader layout: continuous text; edition coordinate retained below.
% ===================== page 4 =====================
\printedpage{2}
% head: A. Grothendieck : Dérivateurs, Ch. VII, Édition des Universités de Montpellier II et Paris VII
% foot: Version 22 décembre 2001

If we have a set of integers
\[
E\subseteq\mathbf{Z},
\]
we will consider in $E$ the equivalence relation generated by the relation between two elements $(i,j)$ of $E$: $j$ is the immediate successor of $i$ \textit{in} $\mathbf{Z}$ (i.e. $j=i+1$). Two distinct elements $i$, $j$ of $E$ are congruent for this relation if and only if the interval of integers that they determine is $\subseteq E$, thus % [sic] distinct -> distincts, plural agreement
\[
\text{if } i<j,\ i \text{ and } j \text{ are \textit{``connected''} in } E \text{ if and only if } [i,j]\subseteq E.
\]

\msp{page 3}
When we speak of the ``connected components'' of $E$, it will always be a matter of the classes under this equivalence relation.

We will identify the set of transition arrows $\Fltr(I)$ of $I$ with the interval of integers $[1,n]$, hence with a subset of $\mathbf{Z}$. If we have a subset $E\subseteq\Fltr(I)$, we thus have in it the notion of transition arrows ``connected'' in $E$, and the notion of connectedness for $E$. If $i<j$, $\underline{u}_i$ and $\underline{u}_j$ are connected if and only if all the intermediate transition arrows between $\underline{u}_i$ and $\underline{u}_j$ are in $E$. And $E$ is connected if and only if it is either empty, or formed of a sequence of transition arrows each of which is the immediate successor of the one preceding it, % [sic] soit formée -> soit formé, agreement with "il"
\[
E=\{\underline{u}_i,\underline{u}_{i+1},\ldots,\underline{u}_{i+l-1}\}, % [sic] \underline{u}_{i+1}\ldots -> \underline{u}_{i+1},\ldots, missing comma
\]
\DERmark{ENV0390}

where $l$ is the cardinal of $E$. Such an $E$ is called a \textit{chain} of transition arrows, and $l$ is its length. A chain is called \textit{composable} if the composite $\underline{u}_{i+l-1}\circ\cdots\circ\underline{u}_{i+1}\circ\underline{u}_i$ \texttt{[or} $\underline{u}_i\circ\underline{u}_{i+1}\circ\ldots\circ\underline{u}_{i+l-1}$\texttt{]} makes sense, i.e. if they are all in the same direction, either direct, or retrograde.

We denote by $\Fltr^+(I)$, $\Fltr^-(I)$ the set of direct, or retrograde, transition arrows, respectively. Their connected components are nothing other than the maximal composable chains. Thus, the maximal composable chains form a partition of the set of transition arrows, into mutually disjoint intervals of integers. The chains themselves form a totally ordered set
\[
C_1,\ C_2,\ \ldots\ C_r\qquad\quad(1\leq r\leq n).
\]
The integer $r$ is called the \textit{``useful length''} of the interval
\msp{page 4}
$I$, it is thus the cardinal of the set of its maximal composable chains.

Two maximal composable chains \texttt{[consecutive]} are always of opposite sign, if one is direct the other is retrograde, and conversely. The first maximal composable chain $C_1$ (let us say simply \textit{``structural chain''}) has the same sign as $\underline{u}_1$; \texttt{[we]} say that the \textit{initial direction} of $I$ is direct, resp. retrograde, if this is so for $\underline{u}_1$, or what amounts to the same thing, for $C_1$. Likewise we define the \textit{final direction} of $I$ as that of $\underline{u}_n$, or again that of $C_r$. Initial and final directions are equal if the useful length $r$ is
\printedpage{3}
% head: A. Grothendieck : Dérivateurs, Ch. VII, Édition des Universités de Montpellier II et Paris VII
% foot: Version 22 décembre 2001
odd, unequal if it is even. (\textbf{NB} It is only for $n=0$, i.e. when $I=\{0\}$ is the one-point category, that initial direction and final direction are not defined.)

The composable chains correspond to the full subcategories of $I$ whose vertices are consecutive in the total order and that are isomorphic to a category $\bDelta_i$ (where $i$ is the length of the chain). Let us note that $\bDelta_i$ has no automorphism other than the identity, so the maximal chains of given length $i\geq 1$ correspond exactly to the functors
\DERmark{ENV0391}

\[
\bDelta_i\lhook\joinrel\longrightarrow\ I
\]
that are injective on the vertices and whose image is an entire structural interval. Such a functor is fully faithful, and a monomorphism. This functor defines a strictly increasing or strictly decreasing map on
\DERmark{ENV0392}

\msp{page 5}
the vertices, according as the chain is direct or retrograde.

Let us note that every standard simplex $\bDelta_i$, $i\geq 0$, can be viewed as an interval, of length $i$, of useful length $1$ (if $i\geq 1$, $0$ if $i=0$). Its origin is $0$, its endpoint is $i$. But let us call an arbitrary category an \textit{interval category} if it is isomorphic to a standard interval. Let us note here

\textbf{Proposition 1.} \textit{Let $I$ be a standard interval of length $n\geq 1$.}

\begin{enumerate}
\item[\textit{a)}] \textit{Let $i\in I$. For $i\in\{0,n\}$, i.e. for $i$ to be one of the two} ``ends'' \textit{(or} ``endpoints''\textit{) (source or target) of $I$, it is necessary and sufficient that one of the two categories $I/i$ or $i\backslash I$ be the one-point category and that the other be isomorphic to a} $\bDelta_j$ (\footnote{where $j$ is the length of the unique structural chain that ``touches'' $i$ (i.e. such that $i$ is in the full subcategory defined by the said chain).}). \textit{(Intrinsic characterization of the ends of an interval.) For $i$ to be the source, it is necessary and sufficient that $i\backslash I$} \texttt{[incomplete sentence]}
\item[\textit{b)}] \textit{Every automorphism $I\longrightarrow I$ that fixes the two endpoints of $I$ is the identity.}
\end{enumerate}
\DERmark{ENV0393}


\textbf{Corollary.} \textit{The group of automorphisms of $I$ has at most two elements. If it has two, the nontrivial automorphism interchanges the two endpoints. (We then say that the interval is} symmetric.\textit{)}

\msp{page 6}
NB If $i$ is not an end, and if nevertheless one of the two categories $i\backslash I$, $I/i$ is one-point, say $I/i$, then in the immediate neighborhood of $i$, $I$ has the form
\[
\begin{tikzcd}[row sep=0pt]
\bullet & \bullet \arrow[l, "\underline{u}_i"'] \arrow[r, "\underline{u}_{i+1}"] & \bullet \\
i-1 & i & i+1
\end{tikzcd}
\]
($\underline{u}_i$ retrograde, $\underline{u}_{i+1}$ direct - if $i\backslash I$ were the one-point one, then $\underline{u}_i$ would be direct, $\underline{u}_{i+1}$ retrograde - in any case these two arrows have opposite directions, i.e. not composable). We thus find that $i\backslash I$ is a category of the form


\endgroup
% END UNIT Der-ChVII-456/Der-ChVII-456_p001-005.tex
% BEGIN UNIT Der-ChVII-456/Der-ChVII-456_p006-010.tex
\begingroup\makeatletter%
\newcommand{\printedpage}[1]{\marginpar{\footnotesize p.~#1}}%
\newcommand{\msp}[1]{\marginpar{\footnotesize [#1]}}%
\newcommand{\bD}{\boldsymbol{\Delta}}%
\newcommand{\iddots}{\mathinner{\mkern1mu\raise1pt\hbox{.}\mkern2mu\raise4pt\hbox{.}\mkern2mu\raise7pt\hbox{.}\mkern1mu}}%
\makeatother\let\printedpage\derpp\let\msp\dermsp%
\let\mspin\DERmspin%


\printedpage{4}
% head: A. Grothendieck : Dérivateurs, Ch. VII, Édition des Universités de Montpellier II et Paris VII
% foot: Version 22 décembre 2001

\[
\begin{tikzcd}[column sep=small, row sep=small]
 & & & i \arrow[dl] \arrow[dr] & & & \\
 & & i-1 & & i+1 & & \\
 & \iddots \arrow[dl] & & & & \ddots \arrow[dr] & \\
 i-r & & & & & & i+s, % [sic] i+r -> i+s, s est la longueur de la chaîne contenant u_{i+1}
\end{tikzcd}
\]
\DERmark{ENV0394}


where $r$ is the length of the structural chain containing $\underline{u}_i$, $s$ that of the structural chain containing $\underline{u}_{i+1}$. This category is an amalgamated sum $\bD_r \amalg \bD_s$ (glued together along their origin), but is not isomorphic to a $\bD_j$, e.g. for lack of a final object.

This observation proves a). To prove b), it suffices to prove that the total order on $I$ is determined intrinsically, by the choice of \textit{the origin} $s_0$ of $I$ (among its % [sic] ces -> ses, possessif
\DERmark{ENV0395}

two endpoints). Suppose that we have already determined intrinsically the vertices $s_0$, $s_1$, $\ldots$, $s_i$ ($i \geq 0$). How to determine $s_{i+1}$? (When $i < n$.) Let $I'$ be the full subcategory of $I$
\msp{page 7}
formed by the $j \in I$ which are not among the $s_j$ \texttt{[rather} $s_k$\texttt{]} already constructed, except the last $s_i$ (thus $s_i \in I'$, $\operatorname{card} I' \geq 2$). Then $I'$ is isomorphic to a standard interval, with $s_i$ corresponding to the origin. One of the categories $s_i\backslash I'$, $I'/s_i$ is one-point, the other isomorphic to a category $\bD_j$, and moreover, $i$ is an initial or final object of $\bD_j$. If it is the initial object $0$, we take for $s_{i+1}$ the immediate successor of $i$ in $\bD_j$. If it is the final object $j$, we take for $s_{i+1}$ the immediate predecessor $j-1$.
\DERmark{ENV0396}


If we had started from the target $n$ of $I$, instead of starting from the source, we would have found the opposite total order relation. Assigning the target element the order number $0$, we find a renumbering of $I$, $i$ becoming $n-i$, and the transition arrows which were direct for the initial order relation, become retrograde in the opposite order relation, and conversely. We thus see:

\textbf{Proposition 2.} \textit{If $I$ is a standard interval, there exists a unique pair $(I',\theta)$, where $I'$ is a standard interval, and $\theta$ an isomorphism $I \xrightarrow{\sim} I'$ which takes the target of $I$ into the source of $I'$ (and the source of $I$ into the target of $I'$).}

This interval $I'$ is called the \textit{opposite} interval of $I$. We must see it, intrinsically, as the ``same'' category $I$, % [sic] categorie -> catégorie, accent manquant
but with an opposite ``orientation''.

\msp{page 8}

A clearer way to see this is the following. Let $J$ be an interval category, i.e. isomorphic to a standard interval. Then there are exactly (if $I \not\simeq e$) \textit{two} pairs $(I,\theta)$ of a standard interval $I$ and an isomorphism $\theta : I \xrightarrow{\sim} J$. These two isomorphisms induce on $J$ two total orders, opposite to one another. Each pair is uniquely determined by the said total order, which we shall call % [sic] appelera -> appellera, orthographe
an \textit{orientation} of the interval $J$. Thus $J$ has two orientations, and each determines an \textit{oriented isomorphism} (i.e. respecting
\printedpage{5}
% head: A. Grothendieck : Dérivateurs, Ch. VII, Édition des Universités de Montpellier II et Paris VII
% foot: Version 22 décembre 2001
the orientations) of $J$ with a standard interval (\footnote[2]{\textbf{NB} We shall speak of the \textit{oppositely oriented} interval $I^{\mathrm{opp}}$ of an \textit{oriented} interval $I$, which we must take care not to confuse with the opposite category $I^o$. Here $I^{\mathrm{opp}}$ has the same underlying category as $I$, but with an opposite orientation.}). These two intervals are the same if and only if there exists an automorphism of $J$ (necessarily unique) which interchanges its two orientations, or what amounts to the same thing, which interchanges the two endpoints. (Thus the categories opposite to one another
\[
\begin{tikzcd}[column sep=small, row sep=small]
 & \bullet \arrow[dl] \arrow[dr] & & & & \bullet & \\
 \bullet & & \bullet & & \bullet \arrow[ur] & & \bullet \arrow[ul]
\end{tikzcd}
\]
are each % [sic] chacunes -> chacune, orthographe
symmetric intervals.) Choosing an orientation also amounts to % [sic] `a -> à, accent imprimé comme guillemet
choosing the origin among the two endpoints.

Let us note that the \textit{opposite category} of an interval is an interval, and the opposite category of a standard interval is the standard interval deduced from it by changing all the signs on the transition intervals $[i-1,i]$ ($1 \leq i \leq n$). We shall take care not to confuse the property $I \simeq I^o$, with the property of symmetry.

\msp{page 9}

Thus the two categories ($\Psi_1$ and $\Phi_1$) above are symmetric, but opposite to one another without being isomorphic. On the other hand, the $\bD_i$ are not symmetric, but each isomorphic to the opposite category.

If $I$ is isomorphic to the opposite category $I^o$, and if $\theta$ is such an isomorphism (necessarily unique), then it is immediate that $\theta$ reverses the orientation (otherwise it would be the identity on the vertices, and the transition arrows would be both direct and indirect), hence $\theta(0) = n$, $\theta(1) = n-1$, and we want $\underline{u}_1$ and $\underline{u}_n$ to be of the same sign. But if $I$ is symmetric, $\underline{u}_1$ and $\underline{u}_n$ are of opposite sign. Thus $I$ cannot be both symmetric and isomorphic to the opposite interval \texttt{[rather to the opposite category; exception:} $I = \{0\}$\texttt{]}.

For any category $\bD_i$, this being an interval, there are (if $i > 0$) \textit{two orientations.} As oriented intervals, there are $\bD_i^+$ (with the ordinary orientation, given by the ordinary order on $\operatorname{Ob} \bD_i = [0,i]$), and $\bD_i^-$ with the opposite orientation.

\printedpage{6}
% head: A. Grothendieck : Dérivateurs, Ch. VII, Édition des Universités de Montpellier II et Paris VII
% foot: Version 22 décembre 2001

\msp{page 10}

\section*{2\quad Full subcategories of an interval}

Let $I$ be an oriented interval category, of length $n \geq 1$. We shall identify $\operatorname{Ob} I$ with the set of integers $[0,n]$. For a chain $\underline{u}_i$, $\underline{u}_{i+1}$, $\ldots$, $\underline{u}_{i+l-1}$ of transition arrows in $I$, we call the \textit{domain} of the chain, the set of sources and targets of these elements. This is thus an interval of integers $[i-1, i+l-1]$, of length $l$ equal to that of the chain, whose origin is determined % [sic] déterminé -> déterminée, accord avec l'origine
by the first arrow $\underline{u}_i$ of the chain (it is its source, if $\underline{u}_i$ is direct, and its target if $\underline{u}_i$ is retrograde), and whose endpoint is determined by the last $\underline{u}_{i+l-1}$. The intervals of integers (viewed as vertices of $I$) corresponding to the structural chains are called the \textit{structural intervals} or \textit{critical intervals} of $I$. Unlike the structural chains, these intervals do not form a partition of $I$, since two structural intervals may have a common endpoint. If the structural chains are
\[
C_1,\ C_2,\ \ldots\ ,\ C_r \qquad (r = \text{ useful length of } I)
\]
of lengths
\[
l_1,\ l_2,\ \ldots\ ,\ l_r,
\]
then the structural intervals are
\[
\underbrace{[0,l_1]}_{I_1},\ \underbrace{[l_1,l_1+l_2]}_{I_2},\ \ldots\ ,\ \underbrace{[l_1+\cdots+l_{r-1},l_1+\cdots+l_{r-1}+l_r]}_{I_r},
\]
a sequence of intervals of integers of which two consecutive ones have a common endpoint, and such that if $i < j$, each element of $I_i$ % [sic] I -> I_i, comparaison des éléments de I_i et de I_j
\DERmark{ENV0397}

is $\leq$ each element of $I_j$. We call \textit{critical vertex} of $I$ the endpoints of the structural intervals.

\textbf{Proposition 3.} \textit{Let $s \in I$, let $I' = I/s$ (full subcategory of $I$ of the $s'$ such that there exists $s' \longrightarrow s$), $I'' = I/s$ \textup{\texttt{[rather} $I'' = s\backslash I$\texttt{]}}, $I'_0 = (I/s)\backslash\{s\}$, $I''_0 = (s\backslash I)\backslash\{s\}$} (\footnote[3]{i.e. that we have: there are no arrows between an element of $I'_0$ and an element of $I''_0$.}). For $s$ to be critical, it is necessary and sufficient that $((I/s)\cup(s\backslash I))\backslash\{s\}$ be the sum category of $I'_0$ and $I''_0$.
\DERmark{ENV0398}\DERmark{ENV0399}


\msp{page 11}

\textbf{Proposition 4.} \textit{Let $i,j$ be in $I$, with $i < j$. Equivalent conditions:}
\begin{enumerate}
\item[\textit{a)}] \textit{There exists an arrow between $i$ and $j$ (either $i \longrightarrow j$, or $j \longrightarrow i$).}
\item[\textit{b)}] \textit{$i$ and $j$ belong to the same structural interval (necessarily unique).}
\item[\textit{c)}] \textit{The full subcategory of $I$ defined by the set of objects $[i,j]$ is isomorphic to a $\bD_\alpha$ (we shall have $\alpha = j-i$).}
\end{enumerate}

\printedpage{7}
% head: A. Grothendieck : Dérivateurs, Ch. VII, Édition des Universités de Montpellier II et Paris VII
% foot: Version 22 décembre 2001
We shall then say that $i$ and $j$ are \textit{directly connected} in $I$. (And we shall use the same terminology in an arbitrary category: $i$ and $j$ are directly connected if and only if $\operatorname{Hom}_I(i,j)$ \textit{or} $\operatorname{Hom}_I(j,i)$ is nonempty.)

\textbf{Proposition 5.} \textit{Let $J \subseteq I$ be a subset of $I$, let us consider the full subcategory of $I$ which it defines. Then:}
\begin{enumerate}
\item[\textit{a)}] \textit{The connected components $J_\alpha$ of $J$ are intervals and the order induced on a $J_\alpha$ by that of $I$ is an orientation (called the \textup{induced orientation}) of $J_\alpha$.}
\item[\textit{b)}] \textit{Suppose $\operatorname{card} J > 1$. Let $i$ be the smallest element of $J$, $j$ the largest element, thus $i < j$. For $J$ itself to be an interval, it is necessary and sufficient that every critical vertex in $[i,j]$ be an element of $J$.}
\end{enumerate}

(\footnote[4]{We call a \textit{subinterval} of $I$ a full subcategory % [sic] sous-categorie -> sous-catégorie, accent manquant
which is connected, i.e. which is itself an interval. This fullness condition is part of the definition: a subcategory $I'$ of $I$ which is an interval is not necessarily full, cf. below \texttt{[paragraph]} on functors between intervals.}).
\DERmark{ENV0400}


\textbf{NB} If $s \in [i,j]$ is a critical vertex of $I$ which is not in $J$, then by proposition 3 it would disconnect $J$, i.e. $J$ would be a direct sum of two nonempty subcategories $J'$ and $J''$ ($i \in J'$, $j \in J''$), hence $J$ would not be an interval.

\msp{page 12}

\textbf{Corollary.} \textit{The set of critical vertices of $I$ is also the set of vertices of the smallest subinterval of $I$ containing the endpoints. More generally, if $i < j$ are two vertices of $I$, the smallest subinterval of $I$ containing $i$, $j$ is formed by $i$, $j$ and the critical vertices $s$ such that $i < s < j$.}

\textbf{Definition.} An interval is said to be \textit{reduced} (or \textit{irredundant}) if all its vertices are critical, or what amounts to the same thing, if its structural chains are of length 1.

A \textit{reduced} oriented interval is known, up to isomorphism, when we know
\begin{enumerate}
\item[a)] its length $\mathrm{n}$,
\item[b)] its initial direction $+$ (direct) or $-$ (retrograde).
\end{enumerate}

\textbf{NB} This direction is the same as the final direction if $n$ is odd, it is opposite if $n$ is even.

Thus, according to the two cases of initial direction, we find
\[
\begin{aligned}
0 \longrightarrow 1 \longleftarrow 2 \longrightarrow \cdots \qquad\qquad \cdots\ n & \qquad (\text{With } & n-1 \longrightarrow n & \quad \text{if } n \text{ odd,}\\
& & n-1 \longleftarrow n & \quad \text{if } n \text{ even.)}\\
0 \longleftarrow 1 \longrightarrow 2 \longleftarrow \cdots \qquad\qquad \cdots\ n & \qquad (\text{With } & n-1 \longleftarrow n & \quad \text{if } n \text{ odd,}\\
& & n-1 \longrightarrow n & \quad \text{if } n \text{ even.)}
\end{aligned}
\]

\printedpage{8}
% head: A. Grothendieck : Dérivateurs, Ch. VII, Édition des Universités de Montpellier II et Paris VII
% foot: Version 22 décembre 2001

\textbf{Proposition 6.} \textit{A subinterval $J$, having $i$ and $j$ ($i < j$) as smallest and largest vertex, is reduced if and only if every $s \in J$ distinct from $i$ and from $j$ is critical.}

\textbf{Corollary.} \textit{The full subcategory of $I$ formed by the critical vertices is a subinterval. We call it the reduced interval associated to $I$, denoted $I_{\textup{red}}$.}


\endgroup
% END UNIT Der-ChVII-456/Der-ChVII-456_p006-010.tex
% BEGIN UNIT Der-ChVII-456/Der-ChVII-456_p011-015.tex
\begingroup\makeatletter%
\newcommand{\printedpage}[1]{\marginpar{\footnotesize p.~#1}}%
\newcommand{\msp}[1]{\marginpar{\footnotesize [#1]}}%
\newcommand{\Ch}{\underline{\mathrm{Ch}}}%
\newcommand{\uHom}{\underline{\mathrm{Hom}}}%
\newcommand{\uFl}{\underline{\mathrm{Fl}}}%
\newcommand{\bDelta}{\boldsymbol{\Delta}}%
\newcommand{\opp}{\mathrm{opp}}%
\makeatother\let\printedpage\derpp\let\msp\dermsp%
\let\mspin\DERmspin%


\printedpage{9}
% head: A. Grothendieck : Dérivateurs, Ch. VII, Édition des Universités de Montpellier II et Paris VII
% foot: Version 22 décembre 2001

\msp{page 13}

\section*{3\quad Paths. Category of paths of a given type $\tau$.}

Let $I$ be an oriented interval, and $X$ a category. We call a \textit{path of type $I$ in $X$} any functor $I \longrightarrow X$, and the \textit{category of paths of type $I$ in $X$} the category
\[
\Ch_I(X) = \uHom(I,X).
\]
Since $I$ is \textit{canonically} isomorphic to a standard interval category $I_0$, we can identify the paths of type $I$, with those of type $I_0$. The type $\tau$ of a path of length $n$ $(= \operatorname{card}(\operatorname{Ob} I)-1)$ in $X$ can therefore be described by a sequence of $n$ signs $+$ or $-$ % [sic] au de type -> à ceux de type, word(s) missing; [sic] le longueur -> de longueur, wrong preposition
\DERmark{ENV0401}\DERmark{ENV0402}\DERmark{ENV0403}

\[
\tau_1,\ \tau_2,\ \dots\ ,\ \tau_n \qquad (\tau_i \in \{-1,+1\} \text{ for } 1 \leq i \leq n).
\]
Because of the fact that $I$ is generated by its transition arrows, a functor $c : I \longrightarrow X$ is determined when we know % [sic] pas ses flèches -> par ses flèches, misprint

\begin{enumerate}
\item[a)] The $c(i)$ $(i \in I)$, i.e. a sequence of $n+1$ objects of $X$:
\[
a_0,\ a_1,\ \dots\ ,\ a_n\ (\text{\footnotemark[5]});
\]
\footnotetext[5]{$a_0$ is called the \textit{source} of the path, $a_n$ its target.}

\item[b)] The $c(\underline{u}_i)$ (where the $\underline{u}_i$ are the transition arrows in $I$), i.e. for every $i \in [1,n]$ an arrow $u_i$ in $X$ between $a_{i-1}$ and $a_i$, going from $a_{i-1}$ to $a_i$ if $\underline{u}_i$ is direct $(\tau_i = +1)$, and from $a_i$ to $a_{i-1}$ if $\underline{u}_i$ is retrograde $(\tau_i = -1)$,
\end{enumerate}

and to any such datum a), b) there corresponds a functor. Thus we shall be able to identify the paths of type $I$ \msp{page 14} in $X$ with the diagrams of the specified type in $X$. And the morphisms between paths likewise correspond to the morphisms between diagrams $a$ and $a'$, an arrow $a \xrightarrow{\alpha} a'$ thus being a sequence of $n+1$ arrows
\[
\alpha_i : a_i \longrightarrow a'_i, \qquad i \in [0,n],
\]
satisfying the commutativity condition for the $n$ squares
\[
\begin{tikzcd}
a_{i-1} \arrow[r, "u_i"] \arrow[d, "\alpha_{i-1}"'] & a_i \arrow[d, "\alpha_i"] \\
a'_{i-1} \arrow[r, "u'_i"] & a'_i
\end{tikzcd}
\qquad\qquad \text{(case $\underline{u}_i$ direct)}
\]

\printedpage{10}
% head: A. Grothendieck : Dérivateurs, Ch. VII, Édition des Universités de Montpellier II et Paris VII
% foot: Version 22 décembre 2001
or
\[
\begin{tikzcd}
a_{i-1} \arrow[d, "\alpha_{i-1}"'] & a_i \arrow[l, "u_i"'] \arrow[d, "\alpha_i"] \\
a'_{i-1} & a'_i. \arrow[l, "u'_i"']
\end{tikzcd}
\qquad\qquad \text{(case $\underline{u}_i$ retrograde)}
\]
If we have an arrow between any two vertices $i$, $j$ in $I$, with $i < j$ and $j$ not the immediate successor of $i$, then $i$, $j$ belong to the same critical interval, and in $I$ we shall have % [sic] quelconque -> quelconques, agreement
\[
i \xrightarrow{\underline{u}_{i+1}} i+1 \xrightarrow{\underline{u}_{i+2}}\ \cdots\ j-1 \xrightarrow{\underline{u}_j} j \qquad \text{a chain of length } j-i \geq 2 % doubt resolved: PDF p. 12 (edition p. 10) clearly underlines u; the existing reading is retained.
\]
or
\[
i \longleftarrow i+1 \longleftarrow\ \ \cdots\ j-1 \longleftarrow j,
\]
and the corresponding arrow between $a_i$ and $a_j$ is obtained by composing the corresponding arrows $u_i$.

\msp{page 15}

\textbf{Examples.}

\begin{enumerate}
\item[1.] $I = \bDelta_0 = e$, $\Ch_I(X) \simeq X$. The paths of length 0 are identified with the objects $a$ of $X$. The unique element $\mathrm{id}_e$ of $\mathrm{Fl}(I)$ gives $\mathrm{id}_a$ in $X$.

\item[2.] $I = \bDelta_1^+$, $\Ch_I(X) \simeq \uFl(X)$. The paths of type $\bDelta_1^+$ are those of length 1 and \textit{direct,} i.e. given by an arrow
\[
a_0 \longrightarrow a_1
\]
whose source is also the source $a_0$ of the path. On the other hand, the paths of type $\bDelta_1^-$ are given by diagrams
\[
a_0 \longleftarrow a_1,
\]
i.e. by arrows whose target is regarded as the source of the path. Of course, the categories $\Ch_{\bDelta_1^+}(X)$ and $\Ch_{\bDelta_1^-}(X)$ are canonically isomorphic, but by an isomorphism which \textit{exchanges source and target of the path.} We shall therefore take good care not to confuse the two preceding paths.

\msp{page 16}

\item[3.] $\bDelta_2^+$, $\bDelta_2^-$, corresponding to the diagrams
\[
a_0 \longrightarrow a_1 \longrightarrow a_2
\]
resp.
\[
a_0 \longleftarrow a_1 \longleftarrow a_2,
\]
we can make a similar remark as for $\bDelta_1^+$, $\bDelta_1^-$.

\printedpage{11}
% head: A. Grothendieck : Dérivateurs, Ch. VII, Édition des Universités de Montpellier II et Paris VII
% foot: Version 22 décembre 2001

\item[4.] $\Phi_1$, $\Psi_1$, these are respectively the intervals
\[
\Phi_1 :\ 0 \xleftarrow{\underline{u}_0} 1 \xrightarrow{\underline{u}_1} 2 \qquad\qquad \text{or} \qquad\qquad
\begin{tikzcd}[row sep=small, column sep=small]
 & 0 \\
1 \arrow[ur] \arrow[dr] & \\
 & 2
\end{tikzcd}
\qquad , % doubt resolved: PDF p. 13 (edition p. 11) clearly underlines u in both labels; the existing reading is retained.
\]
and
\[
\Psi_1 :\ 0 \xrightarrow{\underline{u}_0} 1 \xleftarrow{\underline{u}_1} 2 \qquad\qquad \text{or} \qquad\qquad
\begin{tikzcd}[row sep=small, column sep=small]
 & 0 \arrow[dl] \\
1 & \\
 & 2 \arrow[ul]
\end{tikzcd}
\qquad .
\]
They correspond respectively to the diagrams
\[
a_0 \xleftarrow{u_0} a_1 \xrightarrow{u_1} a_2
\]
and
\[
a_0 \longrightarrow a_1 \longleftarrow a_2.
\]
The categories $\Phi_1$, $\Psi_1$ are opposite to one another, and are not isomorphic. So there is in general no isomorphism between $\Ch_{\Phi_1}(X)$ and $\Ch_{\Psi_1}(X)$.
\end{enumerate}

Let $I$ be an oriented interval, and let us consider the \textit{opposite interval $I^{\opp}$}, having the same underlying category but the opposite orientation. As the definition of the category of paths of type $I$ does not involve the orientation of $I$, we thus have a canonical isomorphism % [sic] l'orientation opposé -> l'orientation opposée, agreement
\[
\Ch_I(X) \xrightarrow{\sim} \Ch_{I^{\opp}}(X).
\]

\msp{page 17}

Nevertheless we shall take care not to confuse a path of type $I$ with the opposite path of type $I^{\opp}$, the source of the one being the target of the other and conversely. This is to say that a path must be seen as \textit{more} than a mere functor
\[
I \longrightarrow X,
\]
we must in addition give ourselves the \textit{orientation of $I$} as forming part of the path structure. So the reasonable formal definition, which fits the intuition of a path as going from a well-defined \textit{source} toward an equally well-defined \textit{target}, is that a path (defined in terms of an interval category $I$) is the datum of a pair $(\omega,c)$, where $\omega$ is an orientation of $I$, and $c$ a functor $c : I \longrightarrow X$. When we look at the category of paths of a fixed ``type'', this \textit{type} is $(I,\omega)$, and not $I$. It is when $\omega$ is fixed once and for all that this category can with impunity be identified with $\uHom(I,X)$ without more. % [sic] la donné -> la donnée, missing final e

We must view this category of paths as equipped with the essential structure
\[
\begin{tikzcd}
\Ch_{(I,\omega)}(X) \arrow[d, "{(\sigma,\beta)}"] \\
X \times X
\end{tikzcd}
\qquad
\begin{array}{l}
c \longmapsto (\sigma(\varphi,v),\beta(\varphi,v)) \\
\text{where $\sigma(c)$ is the \textit{source,}} \\
\qquad \text{$\beta(c)$ the \textit{target} of $c$.}
\end{array} % [sic] ou sigma(c) -> où sigma(c), missing accent
\]
\DERmark{ENV0404}


\printedpage{12}
% head: A. Grothendieck : Dérivateurs, Ch. VII, Édition des Universités de Montpellier II et Paris VII
% foot: Version 22 décembre 2001

\msp{page 18}

This functor is identified with the composition functor
\[
c \longmapsto c \circ \alpha
\]
where $\alpha$ is the endpoint functor (an inclusion when the length is positive)
\DERmark{ENV0405}

\[
\left\{
\begin{array}{lll}
\varepsilon \xrightarrow{\alpha} I & & \\[1ex]
\varepsilon = \{0,1\}, & \multicolumn{2}{l}{\text{discrete category with two elements,}} \\[1ex]
\alpha \text{ given by} & \alpha(0) = & \text{source of } I \\[1ex]
 & \alpha(1) = & \text{target of } I.
\end{array}
\right.
\]
If $\omega'$ is the opposite orientation of $I$, we have a commutative diagram
\[
\begin{tikzcd}
\Ch_{(I,\omega)}(X) \arrow[r, "\varphi", "\sim"'] \arrow[d, "{(\sigma,\beta)}"] & \Ch_{(I,\omega')}(X) \arrow[d, "{(\sigma',\beta')}"] \\
X \times X \arrow[r, "s"] & X \times X,
\end{tikzcd}
\]
where $s$ is the symmetry functor
\[
s(x,y) = (y,x)
\]
(\footnote[6]{The functor $\varphi$ is the one which associates to every path $c$, of type $\tau$, the \textit{opposite path $c^{\opp}$} or $c^{-1}$, of type $\tau^{\opp}$. If $\tau = \tau^{\opp}$, this is therefore a path \textit{of the same type,} and $\varphi$ is an \textit{automorphism} of $\Ch_\tau(X)$. This automorphism (functorial in $X$) is deduced from an automorphism $\theta$ of $I$, cf. below.}). Thus the isomorphism $\varphi$ does not respect the structural ``source-target'' arrows $(\sigma,\beta)$.

There is however a difficulty when there exists an isomorphism of \textit{oriented} intervals
\[
\theta : (I,\omega) \xrightarrow{\sim} (I,\omega'),
\]
i.e. when $I$ is symmetric. This is not the case for the $\bDelta_i$ $(i \geq 1)$ (these categories have no automorphisms other than the identity), but it is the case for the categories $\Phi_1$, $\Psi_1$ of example $4^o$, more generally \msp{page 19} for the reduced categories whose length is \textit{even.} As every time we have an isomorphism $\theta$ of intervals that are \textit{oriented,} we find an isomorphism between the categories of paths of type $(I,\omega)$ and $(I,\omega)^{\opp} = (I,\omega')$, \textit{compatible} this time with the functors $(\sigma,\beta)$:
\[
\begin{tikzcd}
\Ch_{(I,\omega)}(X) \arrow[rr, "\Psi_\theta", "\sim"'] \arrow[dr, "{(\sigma,\beta)}"'] & & \Ch_{(I,\omega')}(X) \arrow[dl, "{(\sigma',\beta')}"] \\
 & X \times X. &
\end{tikzcd}
\]

\printedpage{13}
% head: A. Grothendieck : Dérivateurs, Ch. VII, Édition des Universités de Montpellier II et Paris VII
% foot: Version 22 décembre 2001

\msp{page 20}

It is good to regard a \textit{type} of paths as a noncategorical, but combinatorial, datum, a simple sequence $\tau$ of $n$ signs $\tau_i$ $(1 \leq i \leq n)$ $+$ or $-$ ($n =$ length of the paths under consideration). This combinatorial datum defines a standard oriented interval, $I_\tau$. When we have an arbitrary oriented interval $(I,\omega)$, it defines a type $\tau$, and we identify it canonically with the type $I_\tau$. It is better to write $\Ch_\tau(X)$, rather than $\Ch_{(I,\omega)}(X)$, to make it quite clear that for given $X$, this is a totally ``pinned down'' category, without automorphisms tied to the particular nature of $\tau$ (to automorphisms of $\tau$!). % [sic] une donné -> une donnée, missing final e; [sic] Cette donné -> Cette donnée, missing final e

We can make this point precise as follows. Let $\tau'$ be another type of intervals, and let us set out to determine ``laws'' which allow us to associate to any path of type $\tau$ in an arbitrary category $X$, a path of type $\tau'$ of this same category, and this in a ``functorial'' way, i.e. compatible with the horizontal maps
\[
\begin{tikzcd}
\mathrm{Ch}_\tau(X) \arrow[r, "\mathrm{Ch}_\tau(f)"] \arrow[d, "L_X"'] & \mathrm{Ch}_\tau(X') \arrow[d, "L_{X'}"] \\
\mathrm{Ch}_{\tau'}(X) \arrow[r, "\mathrm{Ch}_{\tau'}(f)"] & \mathrm{Ch}_{\tau'}(X')
\end{tikzcd}
\]
associated with a functor $f : X \longrightarrow X'$.

\msp{page 21}

For example, taking the type $\bDelta_2^+$ (i.e. $+1,+1$), corresponding to the paths
\[
a_0 \xrightarrow{u_1} a_1 \xrightarrow{u_2} a_2,
\]
we can associate to such a path the path of type \texttt{[}$\bDelta_1^+$\texttt{]}
\[
a_0 \xrightarrow{u_2u_1} a_2,
\]
or also the path of type $\bDelta_2^+$
\[
a_0 \xrightarrow{u_2u_1} a_2 \xrightarrow{\mathrm{id}} a_2,
\]
etc. Since the functor
\[
\left\{
\begin{array}{cccl}
X & \longmapsto & \mathrm{Ch}_\tau(X) & \qquad (= \mathrm{Ob}\, \Ch_\tau(X)) \\[1ex]
\mathrm{Cat} & \longrightarrow & (\mathrm{Ens}) &
\end{array}
\right.
\]
is \textit{represented} by the object $I_\tau$ of $X$ \texttt{[rather of} Cat\texttt{]}, the sorites of representable functors tells us that such laws of formation of paths of type $\tau$ from paths of type $\tau'$ \texttt{[rather} $\tau'$ \texttt{from paths of type} $\tau$\texttt{]} correspond bijectively to the arrows in Cat
\[
I_{\tau'} \xrightarrow{\varphi} I_\tau\ ,
\]


\endgroup
% END UNIT Der-ChVII-456/Der-ChVII-456_p011-015.tex
% BEGIN UNIT Der-ChVII-456/Der-ChVII-456_p016-020.tex
\begingroup\makeatletter%
\newcommand{\printedpage}[1]{\marginpar{\footnotesize p.~#1}}%
\newcommand{\msp}[1]{\marginpar{\footnotesize [#1]}}%
\newcommand{\Hom}{\mathrm{Hom}}%
\newcommand{\uHom}{\underline{\mathrm{Hom}}}%
\newcommand{\Ch}{\mathrm{Ch}}%
\newcommand{\uCh}{\underline{\mathrm{Ch}}}%
\newcommand{\Ty}{\mathrm{Ty}}%
\newcommand{\uTy}{\underline{\mathrm{Ty}}}%
\newcommand{\Cat}{\mathrm{Cat}}%
\newcommand{\bDelta}{\boldsymbol{\Delta}}%
\newcommand{\circled}[1]{\text{\textcircled{$\scriptstyle #1$}}}%
\makeatother\let\printedpage\derpp\let\msp\dermsp%
\let\mspin\DERmspin%


\printedpage{14}
% head: A. Grothendieck : Dérivateurs, Ch. VII, Édition des Universités de Montpellier II et Paris VII
% foot: Version 22 décembre 2001

by simply taking the ``law'' corresponding to $\varphi$
\[
c \longmapsto c\circ\varphi \qquad \text{for } c : I_\tau \longrightarrow X.
\]
More precisely, the natural map
\[
\underbrace{\Hom_{\Cat}(I_{\tau'},I_\tau)}_{=\mathrm{Ob}\,\uHom(I_{\tau'},I_\tau)} \longrightarrow \Hom(\mathrm{Ch}_\tau,\mathrm{Ch}_{\tau'}),
\]
\msp{page 22}
where the second Hom is regarded as a Hom of objects in $\Cat^{\vee}=\uHom(\Cat,\mathrm{Ens})$. Thus, in the two examples above, where $\tau=\bDelta_2^+$ and $\tau'=\bDelta_1^+$ or $\bDelta_2^+$, the functors in the reverse direction are
\[
\begin{tikzcd}
\bDelta_2^+\ \Big\{ & 0 \arrow[r] & 1 \arrow[r] & 2 \\
\bDelta_1^+\ \Big\{ & 0 \arrow[u, mapsto] \arrow[rr] & & 1 \arrow[u, mapsto]
\end{tikzcd}
\]
and
\[
\begin{tikzcd}
\bDelta_2^+\ \Big\{ & 0 \arrow[r] & 1 \arrow[r] & 2 \\
\bDelta_2^+\ \Big\{ & 0 \arrow[u, mapsto] \arrow[r] & 1 \arrow[ur, mapsto] \arrow[r] & 2\,. \arrow[u, mapsto]
\end{tikzcd}
\]
This is a general principle for reducing all the $\pm$ canonical ``operations'' on paths in arbitrary categories, and statements concerning them, to questions where $X$ disappears and where we work in the interval categories themselves. One of the technical advantages is that these are very small categories, giving rise moreover to a sure topological intuition (which risks evaporating at times when working in categories of paths), and also that in such categories, hence also in the categories of functors between them, there is always at most
\msp{page 23}
\textit{one single arrow} going from a given object to another. In other words, in such categories all diagrams are commutative.

Of course, we also see that if we have a law of formation $L$ of paths of type $\tau'$ from paths of type $\tau$ (with the condition of functoriality with respect to the category $X$ in which we take the paths), this law is necessarily \textit{functorial with respect to paths:} if we have a homomorphism
\[
c \xrightarrow{v} c' \qquad \text{in } \uCh_\tau(X),
\]
\printedpage{15}
% head: A. Grothendieck : Dérivateurs, Ch. VII, Édition des Universités de Montpellier II et Paris VII
% foot: Version 22 décembre 2001
then the law $L$ ``extends'' to such $v$ to give
\[
L(c) \xrightarrow{L(v)} L(c'),
\]
with the usual relations
\[
L(\mathrm{id}_c)=\mathrm{id}_{L(c)}, \qquad L(v'v)=L(v')L(v). % [sic] \mathrm{id}_C -> \mathrm{id}_c, capital C printed while the chemin is c (cf. L(c))
\]
\DERmark{ENV0406}

This follows from the description of $L$ as an operation $L_\varphi$ in terms of a well-determined functor
\[
I_{\tau'} \xrightarrow{\varphi} I_\tau,
\]
which in turn gives for every $X$ a \textit{functor}
\[
\left\{
\begin{array}{rcl}
\uCh_\tau(X) & \longrightarrow & \uCh_{\tau'}(X)\\[1ex]
c & \longmapsto & c\circ\varphi
\end{array}
\right.
\]
more precise than simply a \textit{map} $\mathrm{Ch}_\tau(X)\longrightarrow\mathrm{Ch}_{\tau'}(X)$ between the corresponding sets of objects. % [sic] un application -> une application, gender agreement

\msp{page 24}
We now have the choice between

\begin{enumerate}[label=\alph*)]
\item A more or less systematic study of functors between intervals. I prefer to wait for this until I need this and that, and to do only the strict minimum for the constructions I have in view.
\item A description of the composition of paths. % [sic] la compositions -> la composition, number agreement
\item A formalism of categories of paths \textit{of unspecified type.}
\end{enumerate}
\DERmark{ENV0407}


Since my aim is above all to describe such categories $\uCh(X)$ (without a subscript $\tau$!), with good properties, I will start with c), then look at b) as a crucial test for the usability of the notions obtained, and finally take up a) when I know what I will need of it.

\printedpage{16}
% head: A. Grothendieck : Dérivateurs, Ch. VII, Édition des Universités de Montpellier II et Paris VII
% foot: Version 22 décembre 2001

\section*{4\quad Category of paths $\uCh(X)$}

I want to ``include'' all the categories $\uCh_\tau(X)$ of paths of all the types one wants, in one and the same category $\uCh(X)$. Recall
\[
\uCh_\tau(X)=\uHom(I_\tau,X),
\]
where $I_\tau$ is the standard (oriented) interval of type $\tau$.

The set of all types of intervals
\msp{page 25}
can be identified with the free monoid Ty on two generators $\Delta^+$ and $\Delta^-$, whose elements are arbitrary finite sequences (possibly empty) of objects alternately equal to $\Delta^+$ and to $\Delta^-$. If the first element is $\Delta^+$, we have a type of path with direct initial direction, retrograde otherwise. The decomposition of $I_\tau$, or of $\mathrm{Fl}\,(I_\tau)$, into critical intervals corresponds to the reduced expression of a word as % [sic] eventuellement -> éventuellement, missing accent
\[
(\Delta^+)^{l_1}(\Delta^-)^{l_2}\ \cdots\ (\Delta^?)^{l_r} \qquad (\Delta^?=\Delta^+ \text{ if } r \text{ odd, } \Delta^?=\Delta^- \text{ if } r \text{ even}),
\]
or as
\[
(\Delta^-)^{l_1}(\Delta^+)^{l_2}\ \cdots\ (\Delta^?)^{l_r} \qquad (\Delta^?=\Delta^- \text{ if } r \text{ odd, } \Delta^?=\Delta^+ \text{ if } r \text{ even}).
\]
The $l_i$ are the lengths of the successive critical intervals (or of the successive structural chains), their number is the reduced length of the interval.

We are going to regard Ty as the set of objects of a category $\uTy$:
\[
\left\{
\begin{array}{l}
\mathrm{Ob}\,\uTy=\Ty=\ \text{free monoid on the two generators } \Delta^+,\ \Delta^- \text{ above.}\\[1ex]
\text{If } \tau,\tau'\in\Ty,\ \Hom(\tau,\tau')=\Hom_{\Cat}(I_\tau,I_{\tau'}).
\end{array}
\right.
\]

(\footnote[7]{A somewhat botched definition, one must restrict oneself to homomorphisms from $I_\tau$ to $I_{\tau'}$ compatible with source and target - cf. p.~29.}). We will note that we have not only a set $\Hom_{\Cat}(I_\tau,I_{\tau'})$, but even a category $\uHom(I_\tau,I_{\tau'})$, which here amounts to saying that $\Hom(I_\tau,I_{\tau'})$ (i.e. $\Hom(\tau,\tau')$) can be viewed as an \textit{ordered} set. But for the moment we are going to ignore this order relation.

\msp{page 26}
The composition of morphisms is the composition of functors between categories $I_\tau$. Thus we have a \textit{fully faithful} (\footnote[8]{false, cf. marginal annotation on the preceding page - we do not take \textit{all} the functors in Cat between $I_\tau$ and $I_{\tau'}$.}) canonical functor
\[
\tau \longmapsto I_\tau : \uTy \xrightarrow{\rho} (\Cat) \qquad \text{(\footnotemark[9]),}
\]
\footnotetext[9]{\textit{realization} functor of a (combinatorial) type of intervals.}% footnote text printed at the foot of p. 17
\printedpage{17}
% head: A. Grothendieck : Dérivateurs, Ch. VII, Édition des Universités de Montpellier II et Paris VII
% foot: Version 22 décembre 2001
by which the category $I_\tau$ \texttt{[rather} $\uTy$\texttt{]} could be identified with a full subcategory of (Cat) (that of the $I_\tau$). But I prefer not to identify $\tau\in\Ty$ (combinatorial datum) with $I_\tau$ (categorical datum).

Let $X$ be a (small?) category. We have a contrafunctor
\[
\left\{
\begin{array}{rcl}
S & \longmapsto & \uHom(S,X)\\[1ex]
(\Cat)^o & \longrightarrow & (\Cat)
\end{array}
\right.
\]
defined by $X$, whence by composing with $\rho$ a contrafunctor
\[
\uTy^{o} \xrightarrow{\mathrm{CH}_X} \Cat.
\]
This is the contrafunctor
\[
\left\{
\begin{array}{l}
\tau \longmapsto \uHom(I_\tau,X) \overset{\mathrm{def}}{=} \uCh_\tau(X),\\[1ex]
\text{i.e. } \mathrm{CH}_X(\tau)=\uCh_\tau(X).
\end{array}
\right.
\]
\DERmark{ENV0408}

But a contrafunctor on a small category $\uTy$ gives rise to a \textit{fibered category} (and even split) over it.

\msp{page 27}
I am going to denote this fibered category by $\uCh(X)$. (Provisionally - if it turns out to be too big, I will restrict it \dots) Thus we have
\[
\begin{tikzcd}[row sep=large]
\uCh(X) \arrow[d, "\pi_X"] \\
\uTy\ .
\end{tikzcd}
\]
The fiber of $\uCh(X)$ at $\tau\in\mathrm{Ob}\,\uTy$ is the category $\uCh_\tau(X)$ of paths of type $\tau$ in $X$. This therefore tells us what the arrows in $\uCh(X)$ over the identity arrows in $\uTy$ are. But what are the more general arrows? We give ourselves an arrow in $\uTy$
\[
\varphi : \tau' \longrightarrow \tau, \qquad \text{i.e. } I_{\tau'} \xrightarrow{\varphi} I_\tau,
\]
what are the arrows of $\uCh(X)$ over $\varphi$? By the general sorites of fibered categories, if
\[
\begin{array}{lclcl}
c\in\uCh(X)_\tau & = & \uCh_\tau(X) & \overset{\mathrm{def}}{=} & \uHom(I_\tau,X)\\[1ex]
c'\in\uCh(X)_{\tau'} & = & \uCh_{\tau'}(X) & \overset{\mathrm{def}}{=} & \uHom(I_{\tau'},X)
\end{array}
\]
\DERmark{ENV0409}\DERmark{ENV0410}

we have
\[
\boxed{\Hom_\varphi(c',c)\simeq\Hom(c',\underbrace{\varphi^*(c)}_{=c\circ\varphi}).}
\]
Thus: the homomorphisms from the
\printedpage{18}
% head: A. Grothendieck : Dérivateurs, Ch. VII, Édition des Universités de Montpellier II et Paris VII
% foot: Version 22 décembre 2001
\msp{page 28}
path $c'$ of type $\tau'$ in $X$, to the path $c$ of type $\tau$, ``over'' the morphism $\varphi$ from $\tau'$ to $\tau$, i.e. the functor
\[
\varphi : I_{\tau'} \longrightarrow I_\tau,
\]
is a homomorphism $v$ from $c'$ to $c\circ\varphi$, i.e.
\[
\begin{tikzcd}[row sep=large]
I_{\tau'} \arrow[r, "\varphi"] \arrow[dr, "c'"', ""{name=A}] & I_\tau \arrow[d, "c", ""'{name=B}] \\
& X.
\arrow[from=A, to=B, bend left=20, shorten <=6pt, shorten >=2pt, "\circled{v}"]
\end{tikzcd}
\]
The composition of arrows is obvious:
\[
\begin{tikzcd}[row sep=large]
I_{\tau''} \arrow[r, "\varphi'"] \arrow[drr, "c''"', ""{name=C}] & I_{\tau'} \arrow[r, "\varphi"] \arrow[dr, "c'"', ""{name=A}] & I_\tau \arrow[d, "c", ""'{name=B}] \\
& & X,
\arrow[from=A, to=B, bend left=20, shorten <=6pt, shorten >=2pt, "v"]
\arrow[from=C, to=A, shorten <=4pt, shorten >=4pt, "v'"]
\end{tikzcd}
\]
we have
\[
\begin{array}{rccl}
v : & c' & \longrightarrow & c\circ\varphi, \quad \text{whence } c'\circ\varphi' \xrightarrow{v*\varphi'} \underbrace{(c\circ\varphi)\circ\varphi'}_{=c\circ(\varphi\varphi')}\\[3ex]
v' : & c'' & \longrightarrow & c'\circ\varphi'
\end{array}
\]
and
\[
v''=(v*\varphi')\circ v', \qquad \text{thus} \qquad
\begin{tikzcd}
c'' \arrow[r, "v'"] \arrow[rr, bend right=12, "\circled{v''}"'] & c'\circ\varphi' \arrow[r, "v*\varphi'"] & c\circ(\varphi\varphi').
\end{tikzcd}
\]
But do we have source and target functors
\[
\uCh(X) \overset{\sigma}{\underset{\beta}{\rightrightarrows}} X \qquad
\begin{array}{ccl}
c & \longmapsto & \text{source}(c)\\[1ex]
c & \longmapsto & \text{target}(c)
\end{array}
\qquad ?
\]
It would therefore be necessary for a homomorphism $v$ as above
\[
v : c' \longrightarrow c\circ\varphi
\]
to define
\[
\begin{array}{lrcl}
\sigma(c) : & \text{source}(c') & \longrightarrow & \text{source}(c)\\[1ex]
\beta(c) : & \text{target}(c') & \longrightarrow & \text{target}(c).
\end{array}
\]


\endgroup
% END UNIT Der-ChVII-456/Der-ChVII-456_p016-020.tex
% BEGIN UNIT Der-ChVII-456/Der-ChVII-456_p021-025.tex
\begingroup\makeatletter%
\newcommand{\printedpage}[1]{\marginpar{\footnotesize p.~#1}}%
\newcommand{\msp}[1]{\marginpar{\footnotesize [#1]}}%
\newcommand{\ul}[1]{\underline{\mathrm{#1}}}%
\newcommand{\bDelta}{\boldsymbol{\Delta}}%
\newcommand{\red}{\text{\upshape red}}%
\newcommand{\circled}[1]{\tikz[baseline=(c.base)]\node[draw,circle,inner sep=0.4pt](c){\scriptsize #1};}%
\makeatother\let\printedpage\derpp\let\msp\dermsp%
\let\mspin\DERmspin%


\printedpage{19}
% head: A. Grothendieck : Dérivateurs, Ch. VII, Édition des Universités de Montpellier II et Paris VII
% foot: Version 22 décembre 2001

\msp{page 29}
But since
\[
\left\{
\begin{array}{l}
\sigma(c)=c(0_\tau),\ \beta(c)=c(1_\tau)\\[1ex]
\text{where } 0_\tau,\ 1_\tau \text{ are the source and target objects of the standard oriented interval } I_\tau\ ,
\end{array}
\right.
\]
and similarly for $\tau'$, $v$ will indeed induce
\[
\underbrace{\sigma(c')}_{=c'(0_{\tau'})} \longrightarrow c(\varphi(0_{\tau'})) \overset{?}{=} \underbrace{\sigma(c)}_{c(0_\tau)},
\]
and likewise for $\beta(c')$ and $\beta(c)$, only if we impose
\[
\boxed{
\left\{
\begin{array}{lcl}
\varphi(0_{\tau'}) & = & 0_\tau\\
\varphi(1_{\tau'}) & = & 1_\tau,
\end{array}
\right.
}
\]
i.e. that the functor $\varphi$ under consideration be \textit{compatible with the source and the target.} This is what visibly had to be imposed in the definition of % [sic] C'est ce visiblement -> C'est visiblement ce, word order
\[
\mathrm{Hom}_{\ul{Ty}}(\tau',\tau)=\underline{\mathrm{Hom}}^{\circled{int}}(I_{\tau'},I_\tau)\quad\text{-}
\]
one must take the functors $\varphi$ respecting the interval structures, i.e. making commutative the diagram
\[
\begin{tikzcd}[column sep=small]
I_{\tau'} \arrow[rr] & & I_\tau \\
& \varepsilon=\{0,1\} \arrow[ul,  "i_{\tau'}"] \arrow[ur,  "i_\tau"'] &
\end{tikzcd}
\]
\DERmark{ENV0411}\DERmark{ENV0412}

defining the orientation of $I_{\tau'}$ and of $I_\tau$ (\footnote[10]{In fact, one must probably restrict the morphisms still further by requiring that they be \textit{increasing} on objects - cf. page 31.}). This rectification made, I have good hope that the category $\ul{Ch}(X)$ we have just constructed is ``the right one''. One could denote by $\ul{CH}(X)$ the larger category, over the category $\ul{Typ}$ having the same objects as $\ul{Ty}$, but more morphisms: namely % [sic] (in footnote 10) qu'il soient -> qu'ils soient, missing plural s
\msp{page 30}
\[
\begin{array}{lcl}
\mathrm{Hom}_{\ul{Typ}}(\tau,\tau') & = & \mathrm{Hom}_{\mathrm{Cat}}(I_\tau,I_{\tau'})\\[0.5ex]
\mathrm{Hom}_{\ul{Ty}}(\tau,\tau') & = & \mathrm{Hom}^{\mathrm{int}}_{\mathrm{Cat}}(I_\tau,I_{\tau'})\\[0.5ex]
 & & \text{(homomorphisms respecting the interval structure).}
\end{array}
\]
We will have
\[
\ul{Ch}(X)\simeq\ul{CH}(X)\times_{\ul{Typ}}\ul{Ty},
\]
\printedpage{20}
% head: A. Grothendieck : Dérivateurs, Ch. VII, Édition des Universités de Montpellier II et Paris VII
% foot: Version 22 décembre 2001
i.e. $\ul{Ch}(X)$ is the \textit{restriction} of the fibered category $\ul{CH}(X)$ over $\ul{Typ}$, to the (non-full) subcategory $\ul{Ty}$ of $\ul{Typ}$.

\textbf{NB} The fibered categories $\ul{CH}(X)$, $\ul{Ch}(X)$ over $\ul{Typ}$, $\ul{Ty}$ depend functorially on $X$, varying in (Cat), as it should be. As is only right, we are interested in them above all when $X$ \textit{is not} in (Cat), i.e. when $X$ is a ``big'' category (such as (Cat) itself). I dispense with making explicit here the pertinent Universe questions.

Two essential tests, to ascertain whether the notion is good
\begin{enumerate}
\item[a)] The theory of the $\mathcal{M}\Sigma^{-1}$.
\item[b)] The composition of paths must work well.
\end{enumerate}

\msp{page 31}
We have the diagram of canonical functors
\[
\begin{tikzcd}
\ul{Ty} \arrow[r, hook] \arrow[d, hook] & \varepsilon\backslash\mathrm{Cat} \arrow[d] \\
\ul{Typ} \arrow[r, hook] & \mathrm{Cat},
\end{tikzcd}
\]
where the horizontal arrows are fully faithful.

\printedpage{21}
% head: A. Grothendieck : Dérivateurs, Ch. VII, Édition des Universités de Montpellier II et Paris VII
% foot: Version 22 décembre 2001

\section*{5\quad Properties of morphisms of intervals}

I have more or less convinced myself that to lend itself to calculations and to simple arguments, one must restrict the morphisms still further in the definition of the category $\ul{Ty}$. I will require that the morphisms $\tau\longrightarrow\tau'$ be functors
\[
w : I_\tau\longrightarrow I_{\tau'}
\]
which are compatible with the \textit{oriented interval} structure of $I_\tau$, $I_{\tau'}$, i.e. that they take source to source, target to target, but moreover that they be \textit{increasing} on objects (for the order relation of progression of paths). This is a restriction which geometrically seems quite natural to me. (It is similar to the analogous restriction, in the pre-categorical definition of morphisms between $\bDelta_i$ \ldots\,) It is this hypothesis that allows one to make a fairly simple structure theory of morphisms of intervals. But for the moment I am going to distinguish ``morphisms'' and ``increasing morphisms''.

\msp{page 32}
Let
\[
f : I_\tau\longrightarrow I_{\tau'}
\]
be a morphism of oriented intervals. Consider the full subcategories, isomorphic to $\Delta_i^+$ or $\Delta_i^-$ as oriented intervals, defined by the critical intervals of $I_\tau$. They form a sequence $\Sigma_\tau$ of subcategories of $I_\tau$
\[
\Sigma_\tau=(\bDelta^{\varepsilon_1}_{l_1}(\tau),\bDelta^{\varepsilon_2}_{l_2}(\tau),\ldots,\bDelta^{\varepsilon_r}_{l_r}(\tau))
\]
of which $I_\tau$ is the amalgamated sum (we will come back to this point further \texttt{[down]}), $r$ being the useful length of $\tau$. We call them the \textit{critical subintervals} of $I_\tau$. Thus $f$ decomposes into a sequence of morphisms
\[
f_i : \bDelta^{\varepsilon_i}_{l_i}\longrightarrow I_{\tau'},\qquad i\in[1,r],
\]
such that
\begin{enumerate}
\item[a)] for $i\in[1,r-1]$, the target of $f_i$ is equal to the source of $f_{i+1}$.

Moreover,
\item[b)] $f_1$ takes source to source, and $f_r$ takes target to target.
\end{enumerate}

Conversely, when we are given $f_i$ satisfying these two conditions, they define a unique morphism of intervals $f$. (And if condition b) is not required, we find a description of the morphisms of categories $I_\tau\longrightarrow I_{\tau'}$.)

\msp{page 33}
We are first going to look at an arbitrary functor
\[
f : \bDelta^{\varepsilon}_i\longrightarrow I_{\tau'}\qquad(\varepsilon\in\{\pm1\})\quad .
\]

\printedpage{22}
% head: A. Grothendieck : Dérivateurs, Ch. VII, Édition des Universités de Montpellier II et Paris VII
% foot: Version 22 décembre 2001
Suppose to fix ideas that $\varepsilon=+1$, i.e. that the first member is a direct critical interval \ldots{} We find that $f$ factors through a critical subinterval of $I_{\tau'}$, say $\bDelta^{\varepsilon'}_j$. It is therefore defined by an increasing map from $\bDelta^{\varepsilon}_i$ to $\bDelta^{\varepsilon'}_j$ (if $\varepsilon$, $\varepsilon'$ \texttt{[are]} of the same sign), a decreasing map in the opposite direction (\footnote[11]{\textbf{Lemma.} \textit{If $I_{\tau'}$ is an interval of strictly positive length, every subset $J$ of $I_{\tau'}$ such that for two elements of $J$, there is always an arrow from one to the other, is contained in a critical interval (unique if} $\operatorname{card} J\geq 2$\textit{).}}). But if we require that $f$ be increasing, then it is necessary that \textit{either} $f$ be constant (in which case, if its value is an interior critical vertex of $I_{\tau'}$, there are exactly \textit{two} critical intervals of $I_{\tau'}$ through which $f$ factors) \textit{or} $f$ nonconstant and $\varepsilon=\varepsilon'$ (and in any case, if $f$ nonconstant, then the critical interval through which $f$ factors is unique). To sum up:
\DERmark{ENV0413}\DERmark{ENV0414}


\textbf{Proposition 7.} \textit{Let $I_\tau$, $I_{\tau'}$ be two oriented intervals,}
\[
f : I_\tau\longrightarrow I_{\tau'}
\]
\textit{a functor. Then for every critical interval $J\subseteq I_\tau$ such that $f|J$ is nonconstant, there exists a unique critical interval $J'\subseteq I_{\tau'}$ such that $f|J$ factors through $J'$, i.e. such that $f(J)\subseteq J'$. If $f$ is increasing, then $J$ and $J'$ are of the same sign,}
\msp{page 34}
\textit{i.e. both direct, or both retrograde.}

Thus, if we look at the sequence $s_0,s_1,\ldots,s_n$ of critical vertices of $\tau$, then
\[
t_0=f(s_0),\ t_1=f(s_1),\ \ldots\ ,\ t_n=f(s_n)
\]
form a sequence of vertices of $I_{\tau'}$, such that two consecutive ones are joined by an arrow of $I_{\tau'}$. Thus the full subcategory of $I_{\tau'}$ defined by these $t_i$ is connected, it is therefore a subinterval. If $f$ is increasing, this subinterval has endpoints $t_0,t_n$ and convex hull $[t_0,t_n]$. If $f$ is a morphism of intervals, it contains all the critical vertices of $I_{\tau'}$. In all cases, taking $t_i$, $t_j$ $(i<j)$ and applying the above to $f|J$, where $J$ is the subinterval of $I$ with origin $s_i$, endpoint $s_j$, we see that every critical vertex of $I_{\tau'}$ lying between $t_i$ and $t_j$ is among the $t_k$ $(i\leq k\leq j)$. To sum up:
\DERmark{ENV0415}


\textbf{Corollary.} \textit{Let $s_i\leq s_j$ be two critical vertices of $I_{\tau}$. Then every critical vertex of $I_{\tau'}$ lying between $f(s_i)=t_i$ and $f(s_j)=t_j$, is the image of a critical vertex $s_k$ lying between $s_i$ and $s_j$ (i.e. $i\leq k\leq j$). In particular, if $f$ is a morphism of intervals,} % [sic] sommets critiques de I_{\tau'} (first occurrence) -> de I_\tau, the s_i are the critical vertices of \tau
\DERmark{ENV0416}

\[
f(I_{\tau\,\red})\supseteq I_{\tau'\,\red}
\]
\textit{(where $I_{\tau\,\red}$ denotes the set of critical vertices of $I_\tau$ or also the reduced subinterval it defines).}

\msp{page 35}
One should be aware that even if $f$ is increasing, one \textit{does not have} in general $f(I_{\tau\,\red})\subseteq I_{\tau'\,\red}$, i.e. $f$ does not necessarily take critical vertices to critical vertices, Minimal example:
\printedpage{23}
% head: A. Grothendieck : Dérivateurs, Ch. VII, Édition des Universités de Montpellier II et Paris VII
% foot: Version 22 décembre 2001

\noindent
\begin{minipage}[c]{0.6\textwidth}
\[
\begin{tikzcd}
\Phi_3^+\ : & s_0 \arrow[r] \arrow[d, mapsto] & s_1 \arrow[d, mapsto] & s_2 \arrow[l] \arrow[r] \arrow[dl, mapsto] & s_3 \arrow[dl, mapsto] \\
\bDelta_2^+\ : & {} \arrow[r] & {} \arrow[r] & {} &
\end{tikzcd}
\]
\end{minipage}\hfill
\begin{minipage}[c]{0.35\textwidth}
$\Phi_3^+$ is reduced, the two critical vertices taken to the noncritical vertex 1.
\end{minipage}

Let us denote by $\mathrm{Cr}(\tau)$ the set of critical intervals of $I_\tau$, in one-to-one correspondence (respecting the total order) with the set of transition arrows of $I_{\tau\,\red}$, % [sic] l'ensembles -> l'ensemble, spurious plural s
\[
\mathrm{Cr}(\tau)\simeq\mathrm{Fl}\,\mathrm{tr}\,(I_{\tau\,\red}).
\]
Let, for a functor
\[
f : I_\tau\longrightarrow I_{\tau'},
\]
$\mathrm{Cr}(\tau,f)$ be the set of critical intervals of $I_\tau$ on which $f$ is not constant. We then have a map
\[
\mathrm{Cr}(\tau,f)\longrightarrow\mathrm{Cr}(\tau') \tag*{$(*)$} % label (*) printed at the left margin
\]
induced by $f$ via proposition 7, associating to each $J\in\mathrm{Cr}(\tau,f)$ the unique $J'\in\mathrm{Cr}(\tau')$ such that $f(J)\subseteq J'$. We have moreover

\textbf{Corollary 2.}
\begin{enumerate}
\item[\textit{a)}] \textit{The image of $(*)$ is a segment of the totally ordered set} $\mathrm{Cr}(\tau')$\textit{, i.e. if two critical intervals $J'$, $J''$ in $I_{\tau'}$ are in the image, the same is true of all} % [sic] I'_{\tau'} -> I_{\tau'}, spurious prime
\msp{page 36}
\textit{those lying between them.}
\DERmark{ENV0417}

\item[\textit{b)}] \textit{This image also contains the critical interval of $I_{\tau'}$ (there is at most one) which contains $f(0_\tau)$ (resp. $f(1_\tau)$) and whose intersection with the segment $[f(0_\tau),f(1_\tau)]$ has cardinality $\geq 2$, as well as all the critical intervals contained in the said segment.}
\item[\textit{c)}] \textit{If $f$ is a morphism of oriented intervals, then the induced map}
\[
\mathrm{Cr}(f):\mathrm{Cr}(\tau,f)\longrightarrow\mathrm{Cr}(\tau')
\]
\textit{is surjective. More precisely, let $J'\in$} $\mathrm{Cr}(\tau')$\textit{, and let $J=f^{-1}(J')$. Then} $\mathrm{Cr}(f)^{-1}(\{J'\})$ \textit{consists of the critical intervals in} $\mathrm{Cr}(\tau,f)$ contained in $J$. \textit{We have moreover,} if $f$ is increasing, % [sic] J' \subseteq Cr(\tau') -> J' \in Cr(\tau'), J' is an element of Cr(\tau')
\DERmark{ENV0418}

\[
\mathrm{long}(J')=\sum_{K\in\mathrm{Cr}(f)^{-1}(\{J'\})}\pi_{\mathrm{n}}(K,f),
\]
\textit{where $\pi_{\mathrm{n}}(K,f)$ is the ``part'' (numerical) of $f|K$, i.e. the length of the interval $[f(0_K),f(1_K)]$. (}\textbf{NB} \textit{If $K$ is a critical interval of $I$ such that $K\notin Cr(\tau,f)$, i.e. $f|K$ is constant, we will put $\pi_{\mathrm{n}}(K,f)=0$. Thus the preceding formula also reads}
\[
\mathrm{long}(J')=\sum\limits_{K\in\mathrm{Cr}(I_\tau),\ K\subseteq J}\pi_{\mathrm{n}}(K,f).\ \textit{)}
\]
\DERmark{ENV0419}

\end{enumerate}


\endgroup
% END UNIT Der-ChVII-456/Der-ChVII-456_p021-025.tex
% BEGIN UNIT Der-ChVII-456/Der-ChVII-456_p026-030.tex
\begingroup\makeatletter%
\newcommand{\printedpage}[1]{\marginpar{\footnotesize p.~#1}}%
\newcommand{\msp}[1]{\marginpar{\footnotesize [#1]}}%
\newcommand{\Tr}{\mathrm{Tr}}%
\newcommand{\Cr}{\mathrm{Cr}}%
\newcommand{\Hom}{\mathrm{Hom}}%
\newcommand{\Ty}{\underline{\mathrm{Ty}}}%
\newcommand{\Ob}{\mathrm{Ob}}%
\newcommand{\lng}{\operatorname{long}}%
\newcommand{\card}{\operatorname{card}}%
\newcommand{\bD}{\boldsymbol{\Delta}}%
\newcommand{\uTy}{\underline{\mathrm{Ty}}}%
\makeatother\let\printedpage\derpp\let\msp\dermsp%
\let\mspin\DERmspin%


\printedpage{24}
% head: A. Grothendieck : Dérivateurs, Ch. VII, Édition des Universités de Montpellier II et Paris VII

\textbf{Corollary 3.} \textit{If $f$ is an} increasing \textit{morphism of oriented intervals, then for every $J' \in \Cr(\tau')$, we have}
\[
(0 \leq)\ \card \Cr(f)^{-1}(\{J'\}) \leq \lng(J'),
\] % [sic] \longrightarrow -> \leq, the next sentence speaks of equality; and {J} -> {J'}, as in Cr(f)^{-1}({J'}) below
\DERmark{ENV0420}\DERmark{ENV0421}

\msp{page 37}
\textit{equality being obtained only if the parts of $f$ on the critical intervals of $I_\tau$ which lie over $J'$ are all equal to $0$ or $1$ (hence if and only if on the $J \in \Cr(f)^{-1}(\{J'\})$, they are equal to $1$).}

\textbf{Corollary 4.} \textit{Suppose moreover that $I'$} \texttt{[}$= I_{\tau'}$\texttt{]} \textit{is reduced. Then}
\[
\mathit{Cr}(f) : \mathit{Cr}(\tau, f) \longrightarrow \mathit{Cr}(\tau')\ (\simeq \mathrm{Fl}\,\mathrm{tr}(\tau'))
\]
\textit{is} bijective.

Thus in this case, $\tau'$ reduced, we obtain a map in the opposite direction
\[
\begin{tikzcd}[column sep=small]
\Tr(\tau') \simeq \Cr(\tau') \arrow[rr, "\Cr^*(f)"] \arrow[dr, "\Cr(f)^{-1}"'] & & \Cr(\tau) \\
 & \Cr(\tau, f), \arrow[ur, hook, "\text{inclusion}"'] &
\end{tikzcd}
\]
a map which transforms the critical interval $J'$ $(= \{i - 1, i\})$ into the \textit{unique} critical interval $J$ of $I_\tau$ such that
\[
f(J) = J'.
\]

\msp{page 38}
But $J \simeq \bD_i^{\varepsilon}$, and we therefore have a \textit{surjective} morphism
\[
f_0 : \bD_i^{\varepsilon} \longrightarrow J' \simeq \bD_1^{\varepsilon}.
\]
We see right away that there then exists a \textit{unique} transition arrow $u$ of $\bD_i^{\varepsilon}$, such that $f(u) \neq \mathrm{id}$. Thus we find the following corollary, more precise than corollary 4:

\textbf{Corollary 5.} \textit{Let $f : I_\tau \longrightarrow I_{\tau'}$ be an increasing morphism of oriented intervals, with $I_{\tau'}$ reduced, and let $\Tr(\tau, f)$ be the set of transition arrows $\underline{u}_i$ of $I_\tau$ such that $f(\underline{u}_i) \neq \mathrm{id}$, i.e. such that $\pi(u_i, f) = 1$. Then $f$ induces a map} % [sic] telle -> telles, agreement with "flèches" (cf. "telles que" just after)
\[
\Tr(\tau, f) \longrightarrow \Tr(\tau')
\]
\textit{which is bijective, and (of course) increasing, hence an isomorphism of ordered sets.}

Thus the map $\Tr(\tau') \longrightarrow \Cr(\tau, f)$ from earlier is made more precise by a bijective map
\[
\Tr(f)^* : \Tr(\tau') \longrightarrow \Tr(\tau, f),
\]
related to $\Cr(f)^*$ above by the commutative triangle
\[
\begin{tikzcd}
 & \Tr(\tau, f) \arrow[dd, "\wr\ \alpha"] \\
\Tr(\tau') \arrow[ur, "\Tr(f)^*"] \arrow[dr, "\Cr(f)^*"'] & \\
 & \Cr(\tau, f),
\end{tikzcd}
\]
% foot: Version 22 décembre 2001
\printedpage{25}
% head: A. Grothendieck : Dérivateurs, Ch. VII, Édition des Universités de Montpellier II et Paris VII
\msp{page 39}
the map $\alpha$ associating to every $u_i \in \Tr(\tau, f)$ the unique critical interval which ``contains'' $u_i$ (i.e. such that $u_i$ is an arrow of the said interval, or again, is part of the associated ``critical chain'' \dots).

We have better still:

\textbf{Proposition 8.} \textit{Let $\tau, \tau' \in \mathrm{Ty}$, $\tau'$} reduced. \textit{For every increasing morphism of oriented intervals}
\[
f : I_\tau \longrightarrow I_{\tau'}
\]
\textit{(we shall simply say that $f$ is an arrow in $\uTy$, and write $f \in \Hom_{\uTy}(I_\tau, I_{\tau^\prime})$), let us associate to it the subset $\Tr(\tau, f)$ of transition arrows of $I_\tau$ nondegenerate for $f$. These arrows form a totally ordered set equipped with signs $+$ or $-$ (direct or retrograde character of the arrows), and $f$ induces a bijection of totally ordered sets} % [sic] I'_\tau -> I_{\tau'}, the target of f is I_{\tau'}
\DERmark{ENV0422}

\[
\underbrace{E_f}_{\subseteq \Tr(\tau)} = \Tr(\tau, f) \simeq \Tr(\tau'),
\]
\textit{and even an isomorphism of totally ordered sets ``equipped with signs''. Thus, to every $f$ we associate a subset $E_f$ of $\Tr(\tau)$, and such that its ``type'' is equal to $\tau'$. Let us denote by}
\[
\mathfrak{P}_{\tau'}(\Tr(\tau)) \subseteq \mathfrak{P}(\Tr(\tau))
\]
\textit{the set of subsets $E$ of $\Tr(\tau)$ whose type is $\tau'$.}

\msp{page 40}
\textit{Then the map}
\[
f \longmapsto \Tr(\tau, f) : \Hom_{\uTy}(\tau, \tau') \longrightarrow \mathfrak{P}_{\tau'}(\Tr(\tau))
\]
\textit{is} bijective \textit{: giving oneself a morphism $f : \tau \longrightarrow \tau'$ in $\uTy$,} amounts to the same thing \textit{as giving oneself a strictly increasing sequence of transition arrows}
\[
\underline{u}_{i_1},\ \underline{u}_{i_2},\ \dots\ ,\underline{u}_{i_{n'}} \qquad (n' = \lng \tau')
\]
\textit{in $I_\tau$ $(i_1 < i_2 < \cdots < i_{n'})$, such that the corresponding sequence of signatures}
\[
(\varepsilon_{i_1}, \varepsilon_{i_2}, \dots, \varepsilon_{i_{n'}}) = \tau'
\]
\textit{is equal to $\tau'$.}

\textbf{Corollary 1.} \textit{For every $\tau$, we have canonical bijections}
\[
\left\{
\begin{array}{lcl}
\Tr^+(\tau) & \simeq & \Hom_{\uTy}(\tau, \Delta^+) \\
\Tr^-(\tau) & \simeq & \Hom_{\uTy}(\tau, \Delta^-),
\end{array}
\right.
\]
functorial in $\tau$ \textit{(varying in} $\uTy$\textit{).}
% foot: Version 22 décembre 2001

\printedpage{26}
% head: A. Grothendieck : Dérivateurs, Ch. VII, Édition des Universités de Montpellier II et Paris VII

To give the corollary a meaning, we have to define
\[
\tau \longmapsto \Tr^+(\tau), \qquad \tau \longmapsto \Tr^-(\tau)
\]
as \textit{functors} contravariant in $\tau$, thus for a morphism
\[
f : I_{\tau_2} \longrightarrow I_{\tau_1}
\]
in $\uTy$, define maps
\msp{page 41}
\[
\begin{array}{rcl}
\Tr^+(f)^* : \Tr^+(\tau_1) & \longrightarrow & \Tr^+(\tau_2) \\
\Tr^-(f)^* : \Tr^-(\tau_1) & \longrightarrow & \Tr^-(\tau_2),
\end{array}
\]
in other words, a map
\[
\Tr(f)^* : \Tr(\tau_1) \quad \longrightarrow \quad \Tr(\tau_2)
\]
``compatible with the signatures''. In corollary 5 above, we defined such a map only for $\tau'$ reduced \texttt{[}$\tau_1 = \tau'$\texttt{,} $\tau_2 = \tau$\texttt{]} (\footnote[12]{What follows does not involve the signatures - it is a sorites on increasing maps between finite totally ordered sets.}). For an arrow $u$ of $I_\tau$, let us call \textit{set-theoretic part} of $f$ on $u$, or of $u$ for $f$, the set-theoretic part of the arrow $f(u)$ in $I_{\tau'}$, i.e. the segment $[f(i-1), f(i)] \subseteq I_{\tau'}$ (here $u = \underline{u}_i$). Let % [sic] où -> ou, alternative "de f sur u, ou de u pour f"
\DERmark{ENV0423}

\[
\underline{u}_{i_1},\ \dots\ ,\underline{u}_{i_N} \qquad \text{or } (v_1, v_2, \dots, v_N)
\] % [sic] v_n -> v_N, the sequence has N elements (cf. v_N below)
\DERmark{ENV0424}

be the increasing sequence of elements of $\Tr(\tau, f)$, we then see by induction on $i$ that for $i \in [1, N-1]$, we have
\[
\begin{array}{rcl}
f(\mathrm{but}(v_i)) & = & f(\mathrm{source}(v_{i+1}))\ \ (= t_i), \\
f(\mathrm{source}(v_1)) & = & 0_{\tau'}, \\
f(\mathrm{but}(v_N)) & = & 1_{\tau'}.
\end{array}
\]
Let us consider the sequence
\[
\{\overbrace{0_{\tau'}}^{= t_0}, t_1, t_2, \dots, t_{N-1}, \overbrace{1_{\tau'}}^{= t_N}\},
\]
it is defined by the corresponding subset of $I_{\tau'}$, which we shall call the set of \textit{critical vertices of $I_{\tau'}$ for $f$.} It is a sequence containing $0_{\tau'}$, $1_{\tau'}$, and such that two consecutive elements are connected in $I_{\tau'}$ by an arrow, which is equivalent to saying that it contains the critical vertices (intrinsic) of $I_{\tau'}$. On the other hand, for every transition interval $[i-1, i]$ of $I_{\tau'}$ $(i \in [1, n'])$, % [sic] appelera -> appellera, French spelling
% foot: Version 22 décembre 2001
\printedpage{27}
% head: A. Grothendieck : Dérivateurs, Ch. VII, Édition des Universités de Montpellier II et Paris VII
\msp{page 42}
there exists one and only one transition interval $[t_{j-1}, t_j]$ for this sequence $(j \in [1, N])$, such that we have $[i-1, i] \subseteq [t_{j-1}, t_j]$. This says that there exists one and only one $v_j \in \Tr(\tau, f)$ such that the part of $f$ on $v_j$ (namely precisely $[t_{j-1}, t_j]$) contains (that of) $\underline{u}_i$. To sum up:
\DERmark{ENV0425}\DERmark{ENV0426}


\textbf{Corollary 2.} \textit{Let $f : \tau \longrightarrow \tau'$ be in $\uTy$. Then for every $u \in \Tr(\tau') = \Tr(I_{\tau'})$, there exists one and only one $v \in \Tr(\tau)$ such that the part of $f$ on $v$ contains $u$. Moreover, the signatures of $u$ and $v$ are equal.}

It is in this way that we therefore obtain a map
\[
\Tr(f)^* : \Tr(\tau') \quad \longrightarrow \quad \Tr(\tau) \qquad (\text{factoring through } \Tr(\tau, f))
\]
inducing
\[
\begin{array}{rcl}
\Tr^+(f)^* : \Tr^+(\tau') & \longrightarrow & \Tr^+(\tau) \\
\Tr^-(f)^* : \Tr^-(\tau') & \longrightarrow & \Tr^-(\tau).
\end{array}
\]
Transitivity for a composite
\[
\tau \xrightarrow{f} \tau' \xrightarrow{g} \tau''
\]
or
\[
I_\tau \xrightarrow{f} I_{\tau'} \xrightarrow{g} I_{\tau''}
\]
is immediate. If $u$, $u'$, $u''$ are transition arrows in $I_\tau$, $I_{\tau'}$, $I_{\tau''}$ respectively, such that
\[
\begin{array}{rcll}
\pi(g, u') & \supseteq & \pi(u''), & \qquad \text{i.e. } u' = \Tr(g)^*(u'') \\
\pi(f, u) & \supseteq & \pi(u'), & \qquad \text{i.e. } u = \Tr(f)^*(u'),
\end{array}
\]
then we have
\[
\pi(gf, u) \ \supseteq\ \pi(u''), \qquad \text{i.e. } u = \Tr(gf)^*(u'').
\]

\msp{page 43}
If the arrows in question have as intrinsic part $[i-1, i]$, $[i'-1, i']$, $[i''-1, i'']$ respectively, the hypothesis is
\[
\begin{array}{rclrcl}
g(i'-1) & \leq & i''-1, & g(i') & \geq & i'' \\
f(i-1) & \leq & i'-1, & f(i) & \geq & i'
\end{array}
\]
and since $g$ is increasing, we have deduced
\[
\begin{array}{rcll}
gf(i-1) & \leq & g(i'-1) & (\leq i''-1) \\
gf(i) & \geq & g(i') & (\geq i''),
\end{array}
\]
\begin{flushright}
q.e.d.
\end{flushright}
% foot: Version 22 décembre 2001

\printedpage{28}
% head: A. Grothendieck : Dérivateurs, Ch. VII, Édition des Universités de Montpellier II et Paris VII

Thus, corollary 1 above acquires a precise meaning and it is proved at the same time, since the correspondences of corollary 1 are given by
\[
\underbrace{f}_{\in \Hom_{\uTy}(\tau, \Delta^{\varepsilon})} \qquad \longmapsto \Tr(f)^*(\underline{u}),
\]
where $\underline{u}$ is the unique transition arrow of $\Delta^{\varepsilon}$.

We are going to use these sorites to give, for every $\tau \in \uTy$, a simple combinatorial description of the categories $\uTy/\tau$ and $\tau\backslash\uTy$, or at least of their sets of objects.

\textbf{1$^o$ Description of} $\Ob\,(\tau\backslash\uTy)$\textbf{.}

To $f \in \Ob(\tau\backslash\uTy)$, i.e.
\[
f : \tau \longrightarrow \tau' \qquad \text{i.e. } f : I_\tau \longrightarrow I_{\tau'}
\]
in $\uTy$, is associated $E = E_f = \Tr(\tau, f) \subseteq \Tr(\tau)$, the set of critical transition arrows (i.e. nondegenerate) for $f$. To each $u \in E$, let us associate the integer $\lng(\pi(f, u))$,
\msp{page 44}
the \textit{part} of $f$ on $u$. Then $(\tau', f)$ is uniquely determined by the pair
\[
(E, \pi), \qquad E \in \mathfrak{P}(\Tr(\tau)), \qquad \underbrace{\pi}_{\text{weighting}} : E \ \longrightarrow\ \underbrace{\mathbf{N}^*}_{\substack{= \text{set of} \\ \text{integers} \geq 1}} \quad .
\]
Or if one prefers, by the map $\pi = \pi_f$
\[
\Tr(\tau) \xrightarrow{\pi} \mathbf{N}, \qquad u \longmapsto \pi(u) = \pi(u, f) = \lng(\pi_{\tau'}(f(u)))
\]
(which gives back $E$ as
\[
E = \pi^{-1}(\mathbf{N}^*) \qquad ).
\]
This map is subject to no restriction,
\[
\Ob\,(\tau\backslash\uTy) \ \xrightarrow[\sim]{\text{bij.}}\ \Hom_{\mathrm{Ens}}(\Tr(\tau), \mathbf{N}).
\]
We recover $\tau'$ from $\pi$ as follows. We take $E = \pi^{-1}(\mathbf{N}^*)$ (transition arrows nondegenerate for $\pi$), $E$ is totally ordered by the order induced by that of $\Tr(\tau)$, it is moreover signatured \texttt{[sic]} by the signaturation \texttt{[sic]} of $\Tr(\tau)$, let
\[
(\varepsilon_1, n_1),\ \dots\ ,\ (\varepsilon_r, n_r)
\]
be the sequence of signatures and weights, of the sequence of elements of $E$ $(\varepsilon_i \in \{\pm 1\}$, $n_i \in \mathbf{N}^*$ for $i \in [1, r])$. Then,
\[
\tau' = \Delta_{\varepsilon_1}^{n_1} \Delta_{\varepsilon_2}^{n_2} \cdots \Delta_{\varepsilon_r}^{n_r}, \qquad r = \card E,
\]
in the free monoid $\mathrm{Ty}$, where we have denoted by $\Delta_+$, $\Delta_-$ (instead of $\Delta^+$, $\Delta^-$) the two generators. The map
\[
I_\tau \ \longrightarrow\ I_{\tau'}
\]
% foot: Version 22 décembre 2001


\endgroup
% END UNIT Der-ChVII-456/Der-ChVII-456_p026-030.tex
% BEGIN UNIT Der-ChVII-456/Der-ChVII-456_p031-035.tex
\begingroup\makeatletter%
\newcommand{\printedpage}[1]{\marginpar{\footnotesize p.~#1}}%
\newcommand{\msp}[1]{\marginpar{\footnotesize [#1]}}%
\newcommand{\Ty}{\mathrm{Ty}}%
\newcommand{\uTy}{\underline{\mathrm{Ty}}}%
\newcommand{\uTyp}{\underline{\mathrm{Typ}}}%
\newcommand{\uCh}{\underline{\mathrm{Ch}}}%
\newcommand{\ovalbox}[1]{\tikz[baseline=(X.base)]\node[draw,rounded corners=6pt,inner sep=2pt](X){#1};}%
\makeatother\let\printedpage\derpp\let\msp\dermsp%
\let\mspin\DERmspin%


\printedpage{29}
% head: A. Grothendieck : Dérivateurs, Ch. VII, Édition des Universités de Montpellier II et Paris VII
\msp{page 45}is described uniquely by saying that it is \textit{constant} on each of the connected components of the graph $E' = E'(f)$ whose vertices are all the vertices of $I_\tau = [0,n]$ and whose edges are the transition arrows degenerate for $f$, i.e. in $\mathrm{Tr}(\tau)\backslash E$. And on the $\alpha$-th element of $E$ ($\alpha \in [1, r = \operatorname{card} E]$), of index $i_\alpha$ in $\mathrm{Tr}(\tau)$, $f$ induces the functor
\[
\boldsymbol{\Delta}_{\varepsilon_\alpha} \simeq \underbrace{J_{i_\alpha}}_{\substack{\text{transition}\\ \text{interval}\\ \text{of index } i_\alpha\\ \text{in } I_\tau}} \longrightarrow I_{\tau'}
\]
\DERmark{ENV0427}

given by \textit{the unique} \underline{Ty}-\textit{arrow} in $I_{\tau'}$ from $J_{i_\alpha}$ to the $\alpha$-th factor, of type $\Delta^{n_\alpha}_{\varepsilon_\alpha}$, of $\tau'$, given by $(*)$
\DERmark{ENV0428}\DERmark{ENV0429}

\[
\Delta^{n_\alpha}_{\varepsilon_\alpha} \simeq \text{type of } \boldsymbol{\Delta}^{\varepsilon_\alpha}_{n_\alpha}\ \ \ldots
\]

\subsection*{2$^o$ Determination of $\mathrm{Ob}\,(\uTy/\tau')$.}

If $(\tau, f) \in \mathrm{Ob}\,(\uTy/\tau')$, thus
\[
f : I_\tau \longrightarrow I_{\tau'},
\]
we first of all associate with it the subset
\[
\mathrm{Cr}(\tau', f) = E'_f = E' \subseteq I_{\tau'}
\]
consisting of the vertices (called ``critical for $f$'', cf. page 41) of $I_{\tau'}$ which are source or target of an arrow $f(u)$, where $u \in \mathrm{Tr}(\tau, f)$, i.e. $f(u)$ is not degenerate (\footnote[13]{isolate a lemma $f$ increasing unnecessary.}). When $I_{\tau'}$ is a point, we set $E'=\{0\}$, so that in every case $E'=f(I_\tau)$. $E'$ contains $f(I_{\tau\,\text{red}})$ (i.e. the images under $f$ of the critical vertices of $I_\tau$),
\DERmark{ENV0430}

\msp{page 46}but can be larger. (E.g. if $\tau = \tau'$, $f = \mathrm{id}$, $E' = I_{\tau'}$, which is thus strictly larger than $I_{\tau'\,\text{red}}$ if $\tau'$ is not reduced.) For every $i \in E'$, let us consider $f^{-1}(\{i\}) \subseteq I_\tau$. Since $f$ \texttt{[is]} increasing, it is an interval of the totally ordered set $I_\tau$, so it is a subinterval of the interval category $I_\tau$. As such, it has a signature $\tau_i \in \Ty$. We thus have a map
\[
\eta_f \text{ or } \eta : E' \longrightarrow \Ty.
\]
I say that $(\tau, f) \in \mathrm{Ob}\,(\uTy/\tau')$ is uniquely determined by this pair
\[
(E', \eta) \begin{cases} E' \in \mathfrak{P}(I_{\tau'}) = \mathfrak{P}([0,n']) \\[1ex] \eta : E' \longrightarrow \Ty\ . \end{cases}
\]
Moreover, the latter is subject to the sole restriction that
\[
E' \supseteq \mathrm{Cr}(\tau') = \text{set of critical vertices of } I_{\tau'}
\]
\printedpage{30}
% head: A. Grothendieck : Dérivateurs, Ch. VII, Édition des Universités de Montpellier II et Paris VII
(or what amounts to the same thing, that $E'$ contains $0_{\tau'}$, $1_{\tau'}$ and that two consecutive elements of $E'$ are joined by an arrow of $I_{\tau'}$): % [sic] où -> ou, « ou ce qui revient au même »
\DERmark{ENV0431}

\[
\left\{
\begin{array}{lcl}
\mathrm{Ob}\,(\Ty/\tau') & \simeq & \text{set of pairs } (E',\eta), \text{ with} \\[1ex]
 & & E' \subseteq [0, \underbrace{\ n'\ }_{=\,\mathrm{long}\,\tau'}], \quad E' \supseteq \mathrm{Cr}(\tau'), \\[1ex]
 & & \eta : E' \longrightarrow \Ty\ .
\end{array}
\right.
\]
Let
\[
0 = t_0 < t_1 < \ \cdots\ < t_N = n'
\]
\DERmark{ENV0432}

be the sequence of elements of $E'$, let for \ovalbox{$i \in [1,N]$}
\[
\left\|
\begin{array}{l}
l_i = t_i - t_{i-1} \\[1ex]
\varepsilon_i \quad \text{signature of } [t_{i-1}, t_i] \text{ in } I_{\tau'}\ ,
\end{array}
\right.
\]
\msp{page 47}and for \ovalbox{$i \in [0,N]$}\,, let
\[
\|\ \ \tau_i = \eta(t_i).
\]
Then the type $\tau$ is given by
\[
\tau = \tau_0 \Delta_{\varepsilon_1} \tau_1 \Delta_{\varepsilon_2} \cdots \Delta_{\varepsilon_N} \tau_N, \tag*{$(*)$}
\]
\texttt{[of]} length $N + \sum_{i=0}^N \underbrace{\mathrm{long}(\tau_i)}_{\lambda_i \geq 0}$.

\subsection*{3$^o$ Description of $\mathrm{Fl}(\uTy)$.}

If
\[
f : \tau \longrightarrow \tau' \qquad \texttt{[is]} \ \text{in } \uTy\ ,
\]
we associate with it as before an integer \underline{$N \geq 0$} ($N = \operatorname{card}\mathrm{Cr}(\tau', f){-}1$, $N = 0$ if and only if $\tau' = 1$, type of the one-point interval), maps
\[
\left\{
\begin{array}{lcl@{\qquad\qquad}ccl}
{[1,N]} & \longrightarrow & \{+1,-1\} \times \mathbf{N}^*, & i & \longmapsto & (\varepsilon_i, l_i) \\[1ex]
{[0,N]} & \longrightarrow & \Ty & i & \longmapsto & \tau_i
\end{array}
\right.
\]
\DERmark{ENV0433}

(\footnote[14]{\textbf{NB} If $N = 0$, $[1,N] = \emptyset$, and there are no $\varepsilon_i$, $l_i$; the map $\underbrace{[0,N]}_{=\,\{0\}} \longrightarrow \Ty$ reduces to the datum of $\tau_0 \in \Ty$.}). The arrow $f$ in $\uTy$ is determined by these data. The type $\tau$ is given by the formula $(*)$ above, the type $\tau'$ by % [sic] determinée -> déterminée, accent manquant
\[
\tau' = \Delta^{l_1}_{\varepsilon_1} \Delta^{l_2}_{\varepsilon_2}\ \cdots\ \Delta^{l_N}_{\varepsilon_N}, \qquad \text{length } n' = \textstyle\sum l_i, \tag*{$(*')$}
\]
\printedpage{31}
% head: A. Grothendieck : Dérivateurs, Ch. VII, Édition des Universités de Montpellier II et Paris VII
and the increasing map from $I_\tau = [0, N + \sum \mathrm{long}(\tau_i)]$ to $I_{\tau'} = [0, \sum l_i]$ is given thus:
\[
\begin{array}{l@{\qquad}cl}
\text{constant on } [0, 0+\lambda_0] & \longmapsto & 0 \\[1ex]
\text{constant on } [\lambda_0+1, \lambda_0+1+\lambda_1] & \longmapsto & l_1 \\[1ex]
\text{constant on } [\lambda_0+\lambda_1+2, \lambda_0+\lambda_1+2+\lambda_2] & \longmapsto & l_1+l_2 \\[1ex]
\ldots & & \\[1ex]
\text{constant on } [\lambda_0+\lambda_1+\cdots+\lambda_i+(i+1), \lambda_0+\cdots+\lambda_i+\lambda_{i+1}+(i+1)] & \longmapsto & \sum_{h=1}^{i+1}l_h.
\end{array}
\]
\DERmark{ENV0434}\DERmark{ENV0435}\DERmark{ENV0436}


\printedpage{32}
% head: A. Grothendieck : Dérivateurs, Ch. VII, Édition des Universités de Montpellier II et Paris VII
\msp{page 48}

\section*{6\quad Composition in $\uTy$}

We have a composition law (everywhere defined) in $\mathrm{Ob}\,(\uTy) = \Ty$, it is that of the free monoid on two generators $\Delta_+$, $\Delta_-$. This composition law is interpreted, at the categorical level (and no longer merely combinatorial) as an amalgamated sum
\[
I_{\tau\tau'} = I_\tau \amalg_{(e, \underline{1}_\tau, \underline{0}_{\tau'})} I_{\tau'},
\]
where $e$ is the unit category, and where $\underline{1}_\tau$, $\underline{0}_{\tau'}$ now represent the morphisms $e \longrightarrow I_\tau$, $e \longrightarrow I_{\tau'}$ , having as image the target element and the source element respectively of the oriented interval in question. We thus have a cocartesian square in (Cat)
\[
\begin{tikzcd}[column sep=small, row sep=small]
 & e \arrow[dl, hook', "\underline{1}_\tau"'] \arrow[dr, hook, "\underline{0}_{\tau'}"] & \\
I_\tau \arrow[dr, hook, "i_{\tau,\tau'}"'] & \text{\scriptsize cocart.} & I_{\tau'} \arrow[dl, hook', "j_{\tau,\tau'}"] \\
 & I_{\tau\tau'}, &
\end{tikzcd}
\]
where all the arrows are monic and fully faithful functors. That this diagram (obviously commutative) is indeed cocartesian in (Cat) amounts, by definition, to saying that for every $X$ in Cat, the transposed diagram
\[
\begin{tikzcd}[column sep=small, row sep=small]
 & X & \\
\uCh_\tau(X) \arrow[ur, "\beta_\tau" near end, "\text{\scriptsize (target)}" near start] & & \uCh_{\tau'}(X) \arrow[ul, "\sigma_{\tau'}"' near end, "\text{\scriptsize (source)}"' near start] \\
 & \uCh_{\tau\tau'}(X), \arrow[ul] \arrow[ur] &
\end{tikzcd}
\]

\texttt{[is cartesian]}

\msp{page 49}which amounts to the following tautology that a path $C$ in $X$ of type $\tau\tau'$ ``is nothing other'' than a pair $(c,c')$ of a path $c$ of type $\tau$ and a path $c'$ of type $\tau'$ ``put end to end'' (in this order), i.e. such that
\[
\mathrm{but}(c) = \mathrm{source}(c')
\]
(and we then have
\[
\begin{array}{rcll}
\mathrm{source}(C) & = & \mathrm{source}(c) & \\[1ex]
\mathrm{but}(C) & = & \mathrm{but}(c') & \qquad\quad ).
\end{array}
\]

\printedpage{33}
% head: A. Grothendieck : Dérivateurs, Ch. VII, Édition des Universités de Montpellier II et Paris VII
Here arises the everlasting problem whether we are going to denote this path by $c{\cdot}c'$, which at first sight seems to impose itself, or rather $c' \circ c$, regarding the path $c$ from $x = \mathrm{source}(c)$ to $y = \mathrm{but}(c)$ as a sort of morphism from $x$ to $y$ (cf. below the formalism of categories of fractions), $c'$ as a morphism from $y = \mathrm{source}(c')$ to $z = \mathrm{but}(c')$, and interpreting the composite $C$ as the composite of these two morphisms $c$, $c'$ (in this order from the point of view of the operations to be carried out), which is traditionally denoted $c' \circ c$ or $c' \cdot c$, and not $c \circ c'$ or $c \cdot c'$, alas! It is this notation that we are going to follow below. % [sic] traditionellement -> traditionnellement, orthographe

But for the moment we are not going to concern ourselves with paths in arbitrary categories $X$, but return to the formalism of the $\tau$ and the $I_\tau$ themselves. The categorical interpretation of $\tau\tau'$ as an amalgamated sum assures us that this composition operation is functorial:

\msp{page 50}
\[
\tau_1 \xrightarrow{\ f\ } \tau_2, \quad \tau'_1 \xrightarrow{\ f'\ } \tau'_2 \qquad\quad \text{gives } \tau_1\tau'_1 \xrightarrow{f*f'} \tau_2\tau'_2
\]
(we shall not write $f \cdot f'$ for the arrow denoted $f * f'$, because of the possible confusion with the composition of arrows in $\uTy$). We have transitivity in an obvious way: if moreover $\tau_2 \xrightarrow{\ g\ } \tau_3$, $\tau'_2 \xrightarrow{\ g'\ } \tau'_3$, \texttt{[we have]}
\[
(gf) * (g'f') = (g * g') \circ (f * f').
\]
Thus we have a composition bifunctor
\[
\uTy \times \uTy \longrightarrow \uTy \qquad\quad (c,c') \longmapsto c \cdot c' \text{ or } c * c'.
\]
(In fact, restricting here the morphisms of $\uTyp$ to functors preserving the two endpoints, we also have a bifunctor
\DERmark{ENV0437}

\[
\uTyp \times \uTyp \longrightarrow \uTyp,
\]
but we prefer now to forget the category \underline{Typ}, which gives a distinctly less simple formalism.)

This composition is \textit{strictly} associative and \textit{strictly unital} (the unit being of course the unit element 1 of the free monoid Ty, corresponding to the one-point category $\{0\}$), and not just ``up to isomorphism''. And with good reason, since in $\uTy$ the isomorphisms are identities. \textit{Thus} $\uTy$ \textit{is a monoid in the category} (Cat) of small categories!

\msp{page 49$'$ (the numbering is shifted)}This implies that not only $\mathrm{Ob}\,\uTy = \Ty$, but also $\mathrm{Fl}(\uTy)$, and more generally all the sets of the form $\mathrm{Hom}_{\mathrm{Cat}}(X, \uTy)$ are monoids, and this functorially in $X$, for example the components $\mathrm{Hom}_{\mathrm{Cat}}(\boldsymbol{\Delta}_i, \uTy)$ of the semi-simplicial set ``nerve'' of $\uTy$,
\[
\mathrm{Nerf}_i(\uTy) = \mathrm{Hom}_{\mathrm{Cat}}(\boldsymbol{\Delta}_i, \uTy),
\]
are monoids, and the semi-simplicial maps between them are monoid homomorphisms: the nerve of $\uTy$ is a \textit{semi-simplicial monoid.}


\endgroup
% END UNIT Der-ChVII-456/Der-ChVII-456_p031-035.tex
% BEGIN UNIT Der-ChVII-456/Der-ChVII-456_p036-040.tex
\begingroup\makeatletter%
\newcommand{\printedpage}[1]{\marginpar{\footnotesize p.~#1}}%
\newcommand{\msp}[1]{\marginpar{\footnotesize [#1]}}%
\newcommand{\uTy}{\underline{\mathrm{Ty}}}%
\makeatother\let\printedpage\derpp\let\msp\dermsp%
\let\mspin\DERmspin%


\printedpage{34}
% head: A. Grothendieck : Dérivateurs, Ch. VII, Édition des Universités de Montpellier II et Paris VII
We are going to describe these monoids $\mathrm{Nerf}_i(\uTy)$, and observe that they are all monoids that are \textit{free,} just like $\mathrm{Nerf}_0(\uTy) = \mathrm{Ob}\,\uTy = \mathrm{Ty}$. The semi-simplicial maps between these monoids are therefore known when we know them on the \textit{generators} of these free monoids (which generators are also the indecomposable elements, hence defined intrinsically in terms of the multiplicative structure). % [sic] connus -> connues, agreement with « applications »

For $\mathrm{Nerf}_0(\uTy) = \mathrm{Ty}$, this is something we have seen, the generators are
\[
\Delta_+,\Delta_-,
\]
corresponding respectively to the oriented intervals
\[
\begin{array}{rccc}
\boldsymbol{\Delta}_+ : & 0 & \longrightarrow & 1\\[4pt]
\boldsymbol{\Delta}_- : & 0 & \longleftarrow & 1.
\end{array}
\]

\msp{page 50$'$}
For $\mathrm{Nerf}_1(\uTy) = \mathrm{Fl}(\uTy)$, i.e. for the arrows, result $3^o$) of the preceding paragraph (determination of $\mathrm{Fl}(\uTy)$) is also interpreted by saying that $\mathrm{Fl}(\uTy)$ is the free monoid generated by an infinite system of generators
\[
\mathfrak{f}_{\varepsilon,n}\ , \qquad\quad (\varepsilon,n)\in\{+1,-1\}\times\mathbf{N},
\]
indexed by the product set $\{+1,-1\}\times\mathbf{N}$. Here $\mathfrak{f}_{\varepsilon,n}$ denotes the unique arrow in $\uTy$
\[
\mathfrak{f}_{\varepsilon,l} : \Delta_\varepsilon \longrightarrow \Delta^l_\varepsilon
\]
(where of course the exponent $l$ denotes an exponentiation for the composition law in $\uTy$
\[
\Delta^l_\varepsilon = \underbrace{\Delta_\varepsilon\circ\Delta_\varepsilon\circ\ \cdots\ \circ\Delta_\varepsilon}_{l\text{ copies}},
\]
and not the Cartesian power in Cat, or in $\uTy$ (where it exists only if $l\in\{0,1\}$)). % [sic] où dans -> ou dans, conjunction intended
\DERmark{ENV0438}

In the description of $\mathrm{Fl}(\uTy)$ in the preceding paragraph, we wrote the ``general term'' of the free monoid by grouping together the groups of successive functors of the form $f_{+,0}$, $f_{-,0}$, i.e. corresponding to the two indecomposable degeneracy arrows (for $*$) % [sic] corréspondant -> correspondant, accent
\[
\left\{
\begin{array}{lccc}
f_{+,0} : & \Delta_+ & \longrightarrow & e\\[4pt]
f_{-,0} : & \Delta_- & \longrightarrow & e
\end{array}
\right.
\]
(\footnote[15]{\textbf{NB} The empty ``word'' of $\mathrm{Fl}(\uTy)$ corresponds to the arrow $\mathrm{id}_1$, where $1$ is the unit element of Ty, realizing the one-point ``interval''.}). The free submonoid of $\mathrm{Fl}(\uTy)$ generated by these two elements is the monoid of arbitrary (pure) degeneracy arrows
% foot: Version 22 décembre 2001
\printedpage{35}
% head: A. Grothendieck : Dérivateurs, Ch. VII, Édition des Universités de Montpellier II et Paris VII
\msp{page 51}
in $\uTy$:
\[
f_\tau : I_\tau \longrightarrow e \qquad\quad (\tau\in\mathrm{Ty}),
\]
and we have
\[
f_\tau * f_{\tau'} = f_{\tau\tau'},
\]
i.e. the map
\[
\tau \longmapsto f_\tau :\quad \underbrace{\mathrm{Ob}\,(\uTy)}_{=\,\mathrm{Ty}\,=\,\mathrm{Nerf}_0(\uTy)} \quad\longrightarrow\quad \underbrace{\mathrm{Fl}(\uTy)}_{\mathrm{Nerf}_1(\uTy)}
\]
is a monoid homomorphism, taking $\Delta_\varepsilon$ ($\varepsilon\in\{+1,-1\}$) to $\mathfrak{f}_{\varepsilon,0}$. This is not the degeneracy map from $\mathrm{Nerf}_0$ to $\mathrm{Nerf}_1$, i.e. the map $\tau\longmapsto\mathrm{id}_\tau$. The latter is given by % [sic] dégénérenscence -> dégénérescence, spelling
\[
\begin{array}{ccc}
\mathrm{Ob}\,(\uTy) & \longrightarrow & \mathrm{Fl}(\uTy)\\[4pt]
\tau & \longmapsto & \mathrm{id}_\tau
\end{array}
\ :\
\left|\
\begin{array}{rcccl}
 & \Delta_\varepsilon & \longmapsto & \mathfrak{f}_{\varepsilon,1} & \qquad(\varepsilon\in\{+1,-1\})\\[4pt]
\text{i.e.} & \Delta_+ & \longmapsto & \mathfrak{f}_{+,1} & \\[4pt]
 & \Delta_- & \longmapsto & \mathfrak{f}_{-,1}\ . &
\end{array}
\right.
\]
As for the source and target homomorphisms
\[
\begin{tikzcd}[column sep=large]
\mathrm{Fl}(\uTy) \arrow[r, shift left, "\sigma\text{ (source)}"] \arrow[r, shift right, "\beta\text{ (target)}"'] & \mathrm{Ob}\,(\uTy),
\end{tikzcd}
\]
they are obvious on the generators
\[
\left\{
\begin{array}{lcl}
\sigma(\mathfrak{f}_{\varepsilon,l}) & = & \Delta_\varepsilon\\[4pt]
\beta(\mathfrak{f}_{\varepsilon,l}) & = & \Delta^l_\varepsilon.
\end{array}
\right.
\]

I am now going to describe $\mathrm{Nerf}_2(\uTy)$, i.e. the systems
\[
\tau \xrightarrow{\ f\ } \tau' \xrightarrow{\ g\ } \tau''.
\]

\msp{page 52}
Giving $f\in\mathrm{Fl}(\uTy)$ amounts to giving an ordered sequence of generators $f_{\varepsilon,l}$, hence a sequence
\[
\begin{aligned}
&\mathfrak{f}_{\varepsilon_1,l_1},\ \mathfrak{f}_{\varepsilon_2,l_2},\ \ \dots\ \ ,\mathfrak{f}_{\varepsilon_n,l_n}\\
&((\varepsilon_i,l_i))_{1\leq i\leq n}\in[\{\pm1\}\times\mathbf{N}]^n.
\end{aligned}
\]
We then have
\[
\begin{aligned}
x &= \Delta_{\varepsilon_1}\Delta_{\varepsilon_2}\cdots\Delta_{\varepsilon_n} &\qquad\qquad& \mathrm{long}\,x = n\\
y &= \Delta^{l_1}_{\varepsilon_1}\Delta^{l_2}_{\varepsilon_2}\cdots\Delta^{l_n}_{\varepsilon_n} && \mathrm{long}\,y = \textstyle\sum_{i\in[1,n]} l_i.
\end{aligned}
\]
% foot: Version 22 décembre 2001

\printedpage{36}
% head: A. Grothendieck : Dérivateurs, Ch. VII, Édition des Universités de Montpellier II et Paris VII
Giving $g : y\longrightarrow z$, i.e. an object of $y\backslash\uTy$, amounts to giving a map from the set of the $m=\sum l_i$ transition arrows of $y$ to $\mathbf{N}$. In fact, we have % [sic] q -> g, the map is called g below (« la donnée de g »)
\DERmark{ENV0439}

\[
\mathrm{Tr}(y) = \mathrm{Tr}(\Delta^{l_1}_{\varepsilon_1})\amalg\mathrm{Tr}(\Delta^{l_2}_{\varepsilon_2})\amalg\ \cdots\ \amalg\mathrm{Tr}(\Delta^{l_n}_{\varepsilon_n})
\]
(``sum set'', i.e. disjoint union), hence
\[
\mathrm{Hom}_{\mathrm{Ens}}(\mathrm{Tr}(y),\mathbf N) \simeq \prod_{i\in[1,n]}\mathrm{Hom}_{\mathrm{Ens}}(\underbrace{\mathrm{Tr}(\Delta^{l_i}_{\varepsilon_i})}_{[1,l_i]},\mathbf{N}), % [sic] Hom_Ens(Tr(y)) -> Hom_Ens(Tr(y),N), second argument missing
\]
\DERmark{ENV0440}

hence giving $g$ is equivalent to giving a system of $n$ maps
\[
\left\{
\begin{array}{l}
\pi_i : [1,l_i]\longrightarrow\mathbf{N} \qquad\quad (\text{weighting of }\Delta^{l_i}_{\varepsilon_i})\\[6pt]
\text{for } i\in[1,n].
\end{array}
\right.
\]
But for $i\in[1,n]$, let us consider the element of $\mathrm{Nerf}_2$
\[
\Delta_{\varepsilon_i}\xrightarrow{\mathfrak{f}_{\varepsilon_i,l_i}}\Delta^{l_i}_{\varepsilon_i}\xrightarrow{g_{\varepsilon_i,\pi_i}}\underbrace{\Delta^{\pi_i(1)}_{\varepsilon_i}\Delta^{\pi_i(2)}_{\varepsilon_i}\cdots\Delta^{\pi_i(l_i)}_{\varepsilon_i}}_{\Delta_{\varepsilon_i}^{\Sigma_j\pi_i(j)\overset{\mathrm{def}}{=}\mathrm{long}(\pi_i)}}, % [sic] \Delta_{\varepsilon_1} -> \Delta_{\varepsilon_i}, and g_{\varepsilon,\pi_i} -> g_{\varepsilon_i,\pi_i}, index i intended
\tag*{(1)}
\]
\DERmark{ENV0441}\DERmark{ENV0442}\DERmark{ENV0443}


\msp{page 53}
where $g_{\varepsilon_i,\pi_i}$ is the arrow defined in terms of $\pi_i$ in the preceding paragraph. But this is also nothing other than
\[
g_{\varepsilon_i,\pi_i} = \mathfrak{f}_{\varepsilon_i,\pi_i(1)}*\mathfrak{f}_{\varepsilon_i,\pi_i(2)}*\cdots*\mathfrak{f}_{\varepsilon_i,\pi_i(l_i)}.
\tag*{(2)}
\]
The element of $\mathrm{Nerf}_2$ under consideration depends only on
\[
(\varepsilon_i,l_i,\pi_i), \qquad\quad \text{where } (\varepsilon_i,l_i)\in\{\pm1\}\times\mathbf{N},\ \pi_i : [1,l_i]\longrightarrow\mathbf{N}.
\]
(\textbf{NB} If $l_i=0$, $[1,l_i]=[1,0]=\emptyset$ by convention, and $\pi_i$ is the unique map $\emptyset\longrightarrow\mathbf{N}$, so the datum $(\varepsilon_i,l_i,\pi_i)$ then amounts to $(\varepsilon_i,l_i)$.)

If
\[
(\varepsilon,l,\pi)\ \text{ as above,} \qquad\quad (\varepsilon,l)\in\{\pm1\}\times\mathbf{N},\ \pi : [1,l]\longrightarrow\mathbf{N},
\]
let us denote by
\[
\mathfrak{f}_{\varepsilon,l,\pi}
\]
the element of $\mathrm{Nerf}_2$ given by
\[
\Delta_\varepsilon\xrightarrow{\mathfrak{f}_{\varepsilon,l}}\Delta^l_\varepsilon\ \xrightarrow{\quad\mathfrak{f}_{\varepsilon,\pi(1)}\ \ *\ \cdots\ *\ \ \mathfrak{f}_{\varepsilon,\pi(l)}\quad}\Delta^{\pi(1)}_\varepsilon\Delta^{\pi(2)}_\varepsilon\cdots\Delta^{\pi(l)}_\varepsilon.
\]
We can formulate the preceding construction by saying that every element of $\mathrm{Nerf}_2$ can be decomposed uniquely into a product
\[
(f,g) = \mathfrak{f}_{\varepsilon_1,l_1,\pi_1}*\mathfrak{f}_{\varepsilon_2,l_2,\pi_2}*\cdots*\mathfrak{f}_{\varepsilon_n,l_n,\pi_n}, % [sic] l_n\pi_n -> l_n,\pi_n, comma missing in index
\]
\DERmark{ENV0444}

which proves that $\mathrm{Nerf}_2(\uTy)$ is the free monoid on the preceding generators $\mathfrak{f}_{\varepsilon,l,\pi}$. (\textbf{NB} to the empty sequence corresponds
\[
1\xrightarrow{\mathrm{id}_1}1\xrightarrow{\mathrm{id}_1}1 \qquad\qquad .)
\]
% foot: Version 22 décembre 2001

\printedpage{37}
% head: A. Grothendieck : Dérivateurs, Ch. VII, Édition des Universités de Montpellier II et Paris VII
\msp{page 54}
If the object of $\mathrm{Nerf}_2$ is
\[
x\xrightarrow{\ f\ }y\xrightarrow{\ g\ }z,
\]
the preceding construction tells us, in terms of the combinatorial datum, i.e. of the sequence of data
\[
(\varepsilon_1,l_1,\pi_1),\ \dots\ ,\ (\varepsilon_n,l_n,\pi_n),
\]
what $x$, $y$, $z$ are: each of them is the product in $\mathrm{Nerf}_0=\mathrm{Ty}$ of the $n$ factors $x_i$, $y_i$ or $z_i$ respectively, for $i\in[1,n]$, given in the definition of the $\mathfrak{f}_{\varepsilon_i,l_i,\pi_i}$, i.e. of
\[
\Delta_{\varepsilon_i},\ \Delta^{l_i}_{\varepsilon_i},\ \text{ and }\Delta^{\pi_i(1)}_{\varepsilon_i}\Delta^{\pi_i(2)}_{\varepsilon_i}\cdots\Delta^{\pi_i(l_i)}_{\varepsilon_i}=\Delta^{\mathrm{long}(\pi_i)}_{\varepsilon_i}
\]
(where $\mathrm{long}(\pi_i)=\sum_{j\in[1,l_i]}\pi_i(j)$). Thus
\[
\left\{
\begin{array}{lcl}
x & = & \prod^{\mathrm{Ty}}_{i\in[1,n]}\Delta_{\varepsilon_i}\\[4pt]
y & = & \prod^{\mathrm{Ty}}_{i\in[1,n]}\Delta^{l_i}_{\varepsilon_i}\\[4pt]
z & = & \prod^{\mathrm{Ty}}_{i\in[1,n]}\Delta^{\mathrm{long}(\pi_i)}_{\varepsilon_i}
\end{array}
\right.
\]
(\footnote[16]{careful, the products must be taken respecting the order of the factors given by the index notation in $i$.}), and we similarly find $f$, $g$ as products (in the sense $*$) of $n$ factors $f_i$ resp. $g_i$ ($i\in[1,n]$) made explicit in formulas (1), (2) (p. 52, 53):
\[
\begin{array}{lclccl}
f_i & = & \mathfrak{f}_{\varepsilon_i,l_i} : & \Delta_{\varepsilon_i} & \longrightarrow & \Delta^{l_i}_{\varepsilon_i}\\[4pt]
g_i & = & g_{\varepsilon_i,\pi_i} : & \Delta^{l_i}_{\varepsilon_i} & \longrightarrow & \Delta^{\mathrm{long}(\pi_i)}_{\varepsilon_i}, % [sic] g_{\varepsilon_i,l_i} -> g_{\varepsilon_i,\pi_i}, notation of (1), (2)
\end{array}
\]
\DERmark{ENV0445}

$g_i$ itself being represented as a $*$-product of $l_i$ indecomposable factors, so that $f$ and $g$ appear in the form
\msp{page 55}
of $*$-products of $n$, resp. of $\sum_{i\in[1,n]}l_i$ factors of the form $\mathfrak{f}_{\varepsilon,\nu}$. To finish getting the truncated semi-simplicial object $(\mathrm{Nerf}_i(\uTy))_{0\leq i\leq 2}$ in hand, it remains to describe the two degeneracy maps $\mathrm{Nerf}_1\longrightarrow\mathrm{Nerf}_2$ (the case where $f$ or $g$ is an identity arrow) and above all the \textit{composition} map for arrows
\[
\mathrm{Nerf}_2(\uTy)\xrightarrow{\ \gamma\ }\mathrm{Nerf}_1(\uTy)=\mathrm{Fl}(\uTy).
\]
It suffices to give $\gamma$ on the generators, and we have tautologically, by definition of the $\mathfrak{f}_{\varepsilon,l,\pi}$
\[
\gamma(\mathfrak{f}_{\varepsilon,l,\pi})=\mathfrak{f}_{\varepsilon,\mathrm{long}(\pi)}.
\]
(\textbf{NB} There is at most one morphism from $\Delta_\varepsilon$ (from $\Delta_+$ or from $\Delta_-$) to an interval $I_\tau$. It exists if and only if $\tau$ is of the form $\Delta^l_\varepsilon$, $l\in\mathbf{N}$.)
% foot: Version 22 décembre 2001

\printedpage{38}
% head: A. Grothendieck : Dérivateurs, Ch. VII, Édition des Universités de Montpellier II et Paris VII
The case $g=\mathrm{id}$ corresponds to the case where
\[
\pi_i(j)=1 \qquad\quad \forall j\in[1,l_i],
\]
the $\mathfrak{f}_{\varepsilon_i,l_i,\pi_i}$ are then determined in terms of the $\mathfrak{f}_{\varepsilon_i,l_i}$, and the degeneracy map is % [sic] détérminés -> déterminés, accent
\[
\mathfrak{f}_{\varepsilon,l}\longmapsto\mathfrak{f}_{\varepsilon,l,1_l} \qquad\quad (\text{where } 1_l \text{ is the constant function with value 1 on } [1,l]).
\]
The case $f=\mathrm{id}$ is the one where the $l_i$ are equal to 1. Then each $\pi_i : [1,1]\longrightarrow\mathbf{N}$ can be viewed as an element of $\mathbf{N}$, giving $(\varepsilon_i,l_i,\pi_i)$ amounts to giving the $(\varepsilon_i,n_i)\in\{\pm1\}\times\mathbf{N}$, i.e. giving % [sic] Le f = id -> Le cas f = id, word « cas » missing
\DERmark{ENV0446}

\msp{page 56}
a sequence of $n$ generators among the $\mathfrak{f}_{\varepsilon,l}$. The degeneracy map $\mathrm{Nerf}_1\longrightarrow\mathrm{Nerf}_2$ is here given on the generators by
\[
\mathfrak{f}_{\varepsilon,l}\longmapsto\mathfrak{f}_{\varepsilon,1,l}.
\]
Fortunately, we no longer have to prove the associativity of composition.

\textbf{NB} The simplicity of these descriptions reinforces my feeling that the category $\uTy$ now defined, with the \textit{increasing morphisms of oriented intervals} $I_\tau\longrightarrow I_{\tau'}$ as arrows of $\uTy$, is indeed ``the right one''. We shall see for a start whether it gives a nice description of the categories of fractions $X\Sigma^{-1}$, without contortions as before.

But we have promised to determine $\mathrm{Nerf}_r(\uTy)$, for every $r\in\mathbf{N}$. The general generator of $\mathrm{Nerf}_r$ is of the form (for $\varepsilon\in\{\pm1\}$ fixed)
\[
\underbrace{\Delta_\varepsilon}_{=\,\Delta^{l_0}_\varepsilon\text{with } l_0=1}\xrightarrow{\ f_1\ }\Delta^{l_1}_\varepsilon\xrightarrow{\ f_2\ }\Delta^{l_2}_\varepsilon\longrightarrow\ \cdots\ \xrightarrow{\ f_r\ }\Delta^{l_r}_\varepsilon,
\]
where the $f_i : \Delta^{l_{i-1}}_\varepsilon\longrightarrow\Delta^{l_i}_\varepsilon$, $i\in[1,r]$, are given by maps
\[
\pi_{i-1} : [1,l_{i-1}]\longrightarrow\mathbf{N} \qquad\quad \text{such that }\ \underbrace{\mathrm{long}(\pi_{i-1})}_{\overset{\mathrm{def}}{=}\ \Sigma_{j\in[1,l_{i-1}]}\pi_{i-1}(j)}\ = l_i. % [sic] \Sigma_{j\in[1,l_i]}\pi_i(j) -> \Sigma_{j\in[1,l_{i-1}]}\pi_{i-1}(j), indices shifted with respect to long(\pi_{i-1})
\]
\DERmark{ENV0447}\DERmark{ENV0448}\DERmark{ENV0449}\DERmark{ENV0450}\DERmark{ENV0451}

(NB $\pi_0 : [1,1]=\{1\}\longrightarrow\mathbf{N}$ is identified with $\mathrm{long}(\pi_0)=\pi_0(1)=l_1$.)

\msp{page 57}
To avoid making the definition of the domain of the $\pi_i$ ($i\in[0,r-1]$) depend on the integers $l_i$, we can also consider the $\pi_i$ as maps
\[
\mathbf{N}^*\xrightarrow{\ \pi\ }\mathbf{N} \qquad\quad (\mathbf{N}^*=\{i\in\mathbf{N}\mid i\geq 1\})
\]
with finite supports (i.e. such that $\pi(i)=0$ for $i$ large), and require precisely
\[
\begin{aligned}
&\mathrm{supp}\,\pi_0\ \subseteq\ \{1\}\\
&\mathrm{supp}\,\pi_1\ \subseteq\ [1,\mathrm{long}(\pi_0)]\\
&\qquad\ \cdots\\
&\boxed{\mathrm{supp}\,\pi_i\ \subseteq\ [1,\mathrm{long}(\pi_{i-1})]}, \qquad\quad \text{for }\ \boxed{i\in[1,r-1]}\quad .
\end{aligned}
\tag*{$(**)$}
\]
\DERmark{ENV0452}

% foot: Version 22 décembre 2001


\endgroup
% END UNIT Der-ChVII-456/Der-ChVII-456_p036-040.tex
% BEGIN UNIT Der-ChVII-456/Der-ChVII-456_p041-045.tex
\begingroup\makeatletter%
\newcommand{\printedpage}[1]{\marginpar{\footnotesize p.~#1}}%
\newcommand{\msp}[1]{\marginpar{\footnotesize [#1]}}%
\newcommand{\Ch}{\underline{\mathrm{Ch}}}%
\newcommand{\Ty}{\underline{\mathrm{Ty}}}%
\newcommand{\Nerf}{\mathrm{Nerf}}%
\makeatother\let\printedpage\derpp\let\msp\dermsp%
\let\mspin\DERmspin%


\printedpage{39}
% head: A. Grothendieck : Dérivateurs, Ch. VII, Édition des Universités de Montpellier II et Paris VII
On the other hand, setting $l_i = \mathrm{long}(\pi_{i-1})$ (for $i \in [1,r]$),
\[
f_i : \Delta_\varepsilon^{l_{i-1}} \longrightarrow \Delta_\varepsilon^{l_i\ (=\mathrm{long}(\pi_{i-1}))}
\]
is defined as
\[
f_i = \mathfrak{f}_{\varepsilon,\pi_{i-1}(1)} * \mathfrak{f}_{\varepsilon,\pi_{i-1}(2)} * \cdots * \mathfrak{f}_{\varepsilon,\pi_{i-1}(l_{i-1})},
\]
so that the element $(*)$ of $\Nerf_r(\Ty)$ is uniquely determined by the data of the $\varepsilon \in \{\pm 1\}$ and of
\[
\pi_\bullet = (\pi_i)_{i\in[0,r-1]} \in (\mathrm{Hom}_{\mathrm{Ens}}(\mathbf{N}^*,\mathbf{N}))^r
\]
satisfying the conditions $(**)$. Let $F_{\varepsilon,\pi_\bullet}$ be this element of $\Nerf_r(\Ty)$, and let $N_r$ be the set of $\pi_\bullet$ we have just made explicit. Then
\[
(F_{\varepsilon,\pi_\bullet})_{(\varepsilon,\pi_\bullet)\in\{\pm1\}\times N_r}
\]
is a free system of generators of the monoid $\Nerf_r(\Ty)$.

If $l_i = 0$ (i.e. \fbox{$\pi_{i-1} = 0$}), for $i<r$, we must have $\mathrm{supp}(\pi_i) = \emptyset$, i.e. $\pi_i = 0$, whence $l_{i+1} = 0$, and likewise for the following $l_j$ and $\pi_j$. Geometrically the hypothesis is expressed by the condition that the target
\DERmark{ENV0453}

\msp{page 58}
of $f_i$ is the element 1 of Ty (corresponding to the one-point interval), and since $\mathrm{Hom}(1,\tau) = \emptyset$ for $\tau \neq 1$, this means that the arrows following $f_i$ must all be equal to the arrow $\mathrm{id}_1 : 1 \longrightarrow 1$.

\textbf{Remark.} In the combinatorial description of $\Ty$ and of $\Nerf_\bullet(\Ty)$, the set of ``signatures'' $\{\pm1\}$ plays a rather phony role. It can be replaced by any set $\Sigma$ (called the ``set of signatures''), and we then construct a monoidal category (or a categorical monoid) having as its set of objects the free monoid generated by $\Sigma$, interpreted as intervals of integers $[0,n]$, equipped with ``signatures'' in $\Sigma$ for each transition interval, or what amounts to the same thing, $\Sigma$-weighted directed ``graphs'' of combinatorial type an interval $\boldsymbol{\Delta}_i$, which amounts to giving an (ordered) sequence of $n$ elements $\varepsilon_1,\ \varepsilon_2,\ \ldots\ ,\ \varepsilon_n \in \Sigma$, i.e. a ``word'' $\varepsilon_1\varepsilon_2\ldots\varepsilon_n$, which we prefer to write $\Delta_{\varepsilon_1}\Delta_{\varepsilon_2}\cdots\Delta_{\varepsilon_n}$. We can define in the obvious way the \textit{increasing} homomorphisms of such weighted directed graphs (understood: with the source element and the target element in the structure, to be respected by the morphisms), without needing a categorical interpretation. The functoriality of the operation of composition of weighted intervals, for this notion of morphism, is obvious. The
\msp{page 59}
choice $\Sigma = \{\pm1\}$ is to my knowledge the only useful one, because of its link with the notion of paths and the calculation of $\pi_0$ of an object $X$ of Cat \texttt{[rather of an object} $X$ \texttt{of} Cat\texttt{]}, etc.
\DERmark{ENV0454}

% foot: Version 22 décembre 2001

\printedpage{40}
% head: A. Grothendieck : Dérivateurs, Ch. VII, Édition des Universités de Montpellier II et Paris VII
\section*{7\quad Explicit description of the categories of fractions $X\Sigma^{-1}$ in terms of the $\Ch(X;a,b)$}

If $X$ is a category, we denote by $\Ch(X)$ the category of paths in $X$ (cf. §4), and for $a,b \in \mathrm{Ob}\,X$, by
\[
\Ch(X;a,b) \qquad \text{or simply } \Ch(a,b)
\]
that of paths from $a$ to $b$, i.e. with source $a$, target $b$. The geometric intuition (which we will make explicit) is that such a path $c$ defines a ``direct'' path from $a$ to $b$ (i.e. a path of type $\tau'$ ``direct'', i.e. all the transition arrows are direct, i.e. $\tau'$ of the form $\Delta_+^n$), of the same length as the initial path $c$, obtained by inverting all the retrograde arrows in the path $c$. This operation is actually possible in $X$ itself if and only if the said arrows are invertible. We then obtain a composable sequence of arrows, whose total composite (which we can call ``the useful effect'' of the path)
\msp{page 60}
is an arrow from $a$ to $b$. In the general case, paths from $a$ to $b$ can be seen as a sort of ``generalized arrows'' going from $a$ to $b$, whose composition would be the obvious composition, and the identities the ``empty paths'' relative to the $a \in \mathrm{Ob}\,X$. We do moreover find in this way a category, having the same objects as $X$, but where the
\[
\mathrm{Hom}_{\mathrm{Ch}}(a,b) = \mathrm{Ch}(a,b) = \mathrm{Ob}\,\Ch(a,b)
\]
are prohibitively large. This is expressed in particular by the fact that the only invertible arrows in $X_{\mathrm{Ch}}$ are the identities! We have no natural functor between $X_{\mathrm{Ch}}$ and $X$ that is the identity on objects. For to a path $a \xrightarrow{c} b$ we cannot associate any $X$-morphisms $a \longrightarrow b$ (unless $c$ \texttt{[is]} of type $\Delta_+^n$, or more generally if its retrograde arrows are invertible), and to an $X$-morphism $a \xrightarrow{u} b$ we can associate a path $c(u)$ of type $\Delta_+$, but if $a \xrightarrow{u} b \xrightarrow{v} c$, $c(vu)$ is still of type $\Delta_+$ whereas $c(v)c(u)$ is of type $\Delta_+^2$, they are not equal. On the other hand, there is a natural arrow between these two elements of $\Ch(a,c)$, an arrow over the unique homomorphism
\DERmark{ENV0455}

\[
\Delta_+ \xrightarrow{\mathfrak{f}_{+,2}} \Delta_+^2.
\]
\msp{page 61}
We then have
\[
\mathfrak{f}_{+,2}^*(\underbrace{a \xrightarrow{u} b \xrightarrow{v} c}_{c(v)c(u)}) = (\underbrace{a \xrightarrow{vu} c}_{c(vu)}),
\]
whence a canonical arrow
\[
c(vu) \longrightarrow c(v)c(u) \qquad \text{in } \Ch(a,c).
\]
This shows that $c(vu)$ and $c(v)c(u)$ are at least in the same connected component of $\Ch(a,c)$, or (as we also say) that they are ``homotopic''. This suggests here modifying the definition of $\mathrm{Hom}_{\mathrm{Ch}}(a,b)$ to
\[
\pi(a,b) = \pi_0(\Ch(a,b)).
\]
% foot: Version 22 décembre 2001

\printedpage{41}
% head: A. Grothendieck : Dérivateurs, Ch. VII, Édition des Universités de Montpellier II et Paris VII
The path composition functor
\[
\begin{array}{ccccc}
\Ch(a,b) & \times & \Ch(b,c) & \longrightarrow & \Ch(a,c)\\
(c & , & c') & \longmapsto & c' \circ c
\end{array}
\]
\texttt{[there are too many} $c$\texttt{]} (\footnote[17]{This functor is deduced from the functors
\[
\Ch_\tau(a,b) \times \Ch_{\tau'}(b,c) \longrightarrow \Ch_{\tau\tau'}(a,c)
\]
which are obvious, cf. page 77.}) induces a map on the $\pi_0$, which gives a composition law
\[
\pi(a,b) \times \pi(b,c) \longrightarrow \pi(a,c)
\]
on arrows. It is immediate that it is associative (the composition of paths being so), and unital. It admits as identity arrow at $x$ either the ``empty loop at $x$'', which is (directly) homotopic in $\Ch(x,x)$ to the ``degenerate'' path $x \xrightarrow{\mathrm{id}} x$
\msp{page 62}
of type $\Delta_+$, or also to the degenerate path of any type $\tau$, say
\[
x \xleftarrow{\mathrm{id}} x \xrightarrow{\mathrm{id}} x \xrightarrow{\mathrm{id}} x \xleftarrow{\mathrm{id}} x \xleftarrow{\mathrm{id}} x,
\]
such a path being none other than $p_\tau^*(x)$, where
\[
p_\tau : \tau \longrightarrow 1
\]
is the degeneracy morphism (the unique morphism in $\Ty$ to 1, which is a final object), whence
\[
p_\tau^* : \Ch_1(x,x) = \{x\} \longrightarrow \Ch_\tau(x,x) \subseteq \Ch(x,x).
\]
The category $\Pi(X)$ thus obtained must be called the \textit{fundamental groupoid} of $X$. Indeed, % [sic] obtenu -> obtenue, agreement with « La catégorie »

\begin{enumerate}[label=\alph*)]
\item all the arrows of $\Pi(X)$ are invertible, i.e. $\Pi(X)$ is indeed a groupoid, and
\item the canonical functor
\[
X \longrightarrow \Pi(X)
\]
which is the identity on objects, and which is defined on arrows by $u \longmapsto c(u)$ as above, is \textit{universal} (in Cat) for functors (i.e. Cat-morphisms) from $X$ to \textit{groupoids.} (More generally still, it is universal for functors $X \xrightarrow{F} Y$ such that $F$ transforms every arrow of $X$ into \texttt{[an]} invertible arrow.)
\end{enumerate}
% foot: Version 22 décembre 2001

\printedpage{42}
% head: A. Grothendieck : Dérivateurs, Ch. VII, Édition des Universités de Montpellier II et Paris VII
\msp{page 63}
To prove that every arrow in $\Pi(X)$ is invertible, since every nonempty path is a composite of ``elementary'' paths of type $\Delta_+$ or $\Delta_-$, it suffices to prove that these paths are invertible. In fact, they correspond bijectively,
\[
\Ch_{\Delta_+}(a,b) \simeq \Ch_{\Delta_-}(b,a),
\]
by associating to the path from $a$ to $b$
\[
c(u) : a \xrightarrow{u} b \qquad \text{of type } \Delta_+
\]
the path
\[
c(u)^o : b \xleftarrow{u} a \qquad \text{of type } \Delta_-
\]
from $b$ to $a$. I say that
\[
\begin{array}{rcl}
c(u)^o c(u) & \sim & 1_a\\
c(u) c(u)^o & \sim & 1_b
\end{array}
\]
\[
\left.
\begin{tikzcd}
 & a \arrow[r, "u"] & b & a \arrow[l, "u"'] \\
 {} & a \arrow[u, "\mathrm{id}"'] \arrow[r, "\mathrm{id}"] & a \arrow[u, "u"'] & a \arrow[u, "\mathrm{id}"'] \arrow[l, "\mathrm{id}"'] \\
 {} \arrow[u] & & a &
\end{tikzcd}
\right\}\
\begin{array}{l}
\text{we have a homomorphism in } \Ch_{\Delta_+\Delta_-}(a,a)\\
\text{from } 1_a \text{ to } c(u)^o c(u).
\end{array}
\]
\[
\left.
\begin{tikzcd}
b \arrow[d, "\mathrm{id}"] & a \arrow[l, "u"'] \arrow[r, "u"] \arrow[d, "u"] & b \arrow[d, "\mathrm{id}"] \\
b & b \arrow[l, "\mathrm{id}"'] \arrow[r, "\mathrm{id}"] & b
\end{tikzcd}
\right\}\
\begin{array}{l}
\text{likewise in } \Ch_{\Delta_-\Delta_+}(b,b),\\
\text{from } 1_b \longleftarrow c(u)c(u)^o % [sic] c(u)^0 -> c(u)^o, digit 0 printed instead of the letter o used elsewhere
\end{array}
\]
\DERmark{ENV0456}\DERmark{ENV0457}\DERmark{ENV0458}

\texttt{[rather} $1_b \longleftarrow c(u)c(u)^0$\texttt{, viz.} % [sic] c(u)^0 -> c(u)^o, digit 0 printed instead of the letter o used elsewhere
\[
\begin{tikzcd}
b \arrow[d, "\mathrm{id}"] & a \arrow[l, "u"'] \arrow[r, "u"] \arrow[d, "u"] & b \arrow[d, "\mathrm{id}"] \\
b & b \arrow[l, "\mathrm{id}"'] \arrow[r, "\mathrm{id}"] & b \ \texttt{].}
\end{tikzcd}
\]
% foot: Version 22 décembre 2001

\printedpage{43}
% head: A. Grothendieck : Dérivateurs, Ch. VII, Édition des Universités de Montpellier II et Paris VII
\msp{page 64}
Let us note that, more generally, for any path
\[
c : a \longrightarrow b, \qquad \text{of type } \tau,
\]
we define an \textit{opposite path}
\[
c^o : b \longrightarrow a \qquad \text{of type } \tau^o,
\]
where
\[
\tau \longmapsto \tau^o : \mathrm{Ty} \longrightarrow \mathrm{Ty}
\]
is the anti-involution of the free monoid Ty which exchanges the two generators $\Delta_-$, $\Delta_+$, and which is therefore given on the general element by
\[
(\Delta_{\varepsilon_1}\Delta_{\varepsilon_2}\cdots\Delta_{\varepsilon_n})^o = \Delta_{-\varepsilon_n}\Delta_{-\varepsilon_{n-1}}\cdots\Delta_{-\varepsilon_1}.
\]
By induction on the length of $c$ (\footnote[18]{unnecessary, follows directly from formula $(*)$ \texttt{[on the]} following page, by decomposing $c$ into composites of ``elementary'' paths of type $\Delta_+$ or $\Delta_-$.}), using the preceding result, that $c(u)$, $c(u)^o$ are inverses of one another in $\Pi(X)$, we find
\[
\left.
\begin{array}{rcl}
c^o c & = & \mathrm{id}_a\\
c c^o & = & \mathrm{id}_b
\end{array}
\right\}\text{in } \Pi(X).
\]

\textbf{Example.} If
\[
\begin{array}{rcl}
\tau & = & \Delta_+\Delta_-\Delta_+^2\Delta_-\\
\tau^o & = & \Delta_+\Delta_-^2\Delta_+\Delta_-,
\end{array}
\]
and if we have a path
\[
c : a_0 \xrightarrow{u_1} a_1 \xleftarrow{u_2} a_2 \xrightarrow{u_3} a_3 \xrightarrow{u_4} a_4 \xleftarrow{u_5} a_5
\]
of type $\tau$, we deduce from it the path
\[
c^o : a_5 \xrightarrow{u_5} a_4 \xleftarrow{u_4} a_3 \xleftarrow{u_3} a_2 \xrightarrow{u_2} a_1 \xleftarrow{u_1} a_0
\]
of type $\tau^o$. The operation $c \longmapsto c^o$ satisfies in
\msp{page 65}
fact \texttt{[for]}
\[
\begin{array}{ccccc}
a & \xrightarrow{\gamma} & b & \xrightarrow{\gamma'} & c\\
c & \xrightarrow{\gamma'^o} & b & \xrightarrow{\gamma^o} & a
\end{array}
\]
the relation
\[
(\gamma'\gamma)^o = \gamma^o\gamma'^o, \tag*{$(*)$}
\]
% foot: Version 22 décembre 2001


\endgroup
% END UNIT Der-ChVII-456/Der-ChVII-456_p041-045.tex
% BEGIN UNIT Der-ChVII-456/Der-ChVII-456_p046-050.tex
\begingroup\makeatletter%
\newcommand{\printedpage}[1]{\marginpar{\footnotesize p.~#1}}%
\newcommand{\msp}[1]{\marginpar{\footnotesize [#1]}}%
\newcommand{\Ch}{\underline{\mathrm{Ch}}}%
\newcommand{\Ty}{\mathrm{Ty}}%
\newcommand{\uTy}{\underline{\mathrm{Ty}}}%
\newcommand{\Fl}{\mathrm{Fl}}%
\newcommand{\Hom}{\mathrm{Hom}}%
\newcommand{\ch}{\mathrm{ch}}%
\newcommand{\hot}{\mathrm{hot}}%
\newcommand{\id}{\mathrm{id}}%
\makeatother\let\printedpage\derpp\let\msp\dermsp%
\let\mspin\DERmspin%


\printedpage{44}
% head: A. Grothendieck : Dérivateurs, Ch. VII, Édition des Universités de Montpellier II et Paris VII
% foot: Version 22 décembre 2001

which is in accord with the intuition that the operation $\gamma \longmapsto \gamma^o$ plays the role of passing to the inverse (``up to homotopy''), and we have the formula
\[
(\gamma'\gamma)^{-1} = \gamma^{-1}\gamma'^{-1}
\]
for two invertible arrows $\gamma$, $\gamma'$ composable (in the order $\gamma'$, $\gamma$).

As for the universal property of
\[
X \longrightarrow \Pi(X)
\]
(which will allow us to write
\[
\Pi(X) \simeq X\Sigma^{-1}, \qquad \text{with } \Sigma = \Fl(X) \ \ ),
\]
we may as well establish it right away in the general setting where we give ourselves a subset
\[
\Sigma \subseteq \Fl(X),
\]
and we define, for $\tau \in \Ty$,
\[
\Ch_{\Sigma,\tau}(X) \subseteq \Ch_{\tau}(X)
\]
as the full subcategory of $\Ch_{\tau}(X)$ consisting of the paths in $X$ for which \textit{all the retrograde arrows}

\msp{page 66}

are in $\Sigma$ \texttt{[rather in $\Sigma$]}. We shall note that for a path $c$ and its opposite $c^o$ both to be $\Sigma$-\textit{paths,} it is necessary and sufficient that \textit{all} the arrows of the path $c$ be in $\Sigma$ (we shall then say that $c$ is a $\Sigma$-\textit{strict path}), since the retrograde arrows of $c^o$ are the direct arrows of $c$.
\DERmark{ENV0459}


For the full subcategories $\Ch_{\tau,\Sigma}(X)$ to be stable under the inverse image functors
\[
f^* : \Ch_{\tau'}(X) \longrightarrow \Ch_{\tau}(X)
\]
associated with the arrows
\[
f : \tau \longrightarrow \tau'
\]
in $\uTy$, or simply for there to be stability for the two functors $f^*$ associated with an $f : \Delta_- \longrightarrow 1$ and $\Delta_- \longrightarrow \Delta_-^2$, it is necessary and sufficient that $\Sigma$ be stable under composition and contain all the identity arrows % [sic] qu'y ait -> qu'il y ait, missing word
\[
\left\{
\begin{aligned}
& x \in \mathrm{Ob}\,X \ \Longrightarrow \ \id_x \in \Sigma \\
& u, v \in \Sigma, \ \ u,\, v \text{ composable} \Longrightarrow uv \in \Sigma.
\end{aligned}
\right.
\]
Indeed, the arrows in a path $f^*(c') \in \Ch_{\tau}(X)$ are either identity arrows (which correspond to the transition arrows of $I_\tau$ degenerate for $f$), or arrows among those of the path $c'$ of type $\tau'$, or composites of a sequence of consecutive composable arrows of this path (which correspond to the transition arrows in $I_{\tau'}$ included in the part of a transition arrow of $I_\tau$), with preservation of signatures.
\DERmark{ENV0460}


\printedpage{45}
% head: A. Grothendieck : Dérivateurs, Ch. VII, Édition des Universités de Montpellier II et Paris VII
% foot: Version 22 décembre 2001

\msp{page 67}

We assume these conditions satisfied, so that the $\Ch_{\tau,\Sigma}(X)$ form the fibers of a full fibered subcategory of $\Ch(X)$ over $\uTy$, which we shall denote by
\[
\Ch_{\Sigma}(X) \hookrightarrow \Ch(X).
\]
It is obvious that this subcategory is stable under the operation of composition of paths (but not, in general, as we have said, under passage from $c$ to $c^o$ (for this $\Sigma$ would have to contain \textit{all} the arrows of $X$, since for every $u : a \longrightarrow b$, which is thus in $\Ch_{\Delta_+,\Sigma}(X)$, $u^o$ given by $b \xleftarrow{\ u\ } a$ in $\Ch_{\Delta_-}(X)$ would have to be a $\Sigma$-path, whence $u \in \Sigma$)).

We are going to define a canonical category
\[
\Pi_\Sigma(X),
\]
having the same objects as $X$, and a canonical functor
\[
X \longrightarrow \Pi_\Sigma(X)
\]
which is the identity on objects, and which transforms the $u \in \Sigma$ into invertible arrows. For this we copy the description of the fundamental groupoid (but $\Pi_\Sigma(X)$ will not in general be a groupoid), by putting
\[
\Hom_\Sigma(x,y) = \pi_0(\Ch_\Sigma(x,y)),
\]

\msp{page 68}

and defining the composition of ``$\Sigma$-arrows'' in terms of the composition of paths. The argument which proved, in the case of $\Pi(X)$, that \textit{every} arrow of $X$, viewed as a path $c$ of type $\Delta_+$, was invertible in $\Pi(X)$, having as inverse $c^o$ of type $\Delta_-$, is still valid, provided that $c^o$ is indeed still a $\Sigma$-path, which requires that $u \in \Sigma$.

We are going to interpret it in terms of a canonical functor
\[
\kappa : X \longrightarrow \Pi_\Sigma(X) \qquad \text{(\footnotemark[19])}
\tag*{$(*)$}
\]
\footnotetext[19]{On arrows, $\kappa$ is defined, for $u : a \longrightarrow b$, by $\kappa(u) =$ homotopy class of $\ch(u)$ in $\Ch_\Sigma(a,b)$, where $\ch(u)$ is the path of type $\Delta_+$ from $a$ to $b$ defined by $u$.}
which is again described as in the case of the fundamental groupoid, by associating to every arrow $u$ of $X$ the homotopy class of the path $c(u)$ it defines. We thus see that this canonical functor transforms the $u \in \Sigma$ into invertible arrows.

When the arrows of $\Sigma$ are already invertible in $X$, I say that the functor $(*)$ is an isomorphism, i.e. that for every $a, b \in \mathrm{Ob}\,X$,
\[
\alpha : \Hom_X(a,b) \longrightarrow \pi_0\Ch_\Sigma(a,b), \qquad u \longmapsto \underbrace{\ch_\Sigma(u)}_{= \ch(u)}
\]
is bijective. To see this, we are going to define
\[
\beta : \pi_0(\Ch_\Sigma(a,b)) \longrightarrow \Hom_X(a,b),
\]
which will be inverse to $\alpha$. We define % [sic] inverse de $\beta$ -> inverse de $\alpha$, $\beta$ is meant to be the inverse of $\alpha$
\DERmark{ENV0461}


\printedpage{46}
% head: A. Grothendieck : Dérivateurs, Ch. VII, Édition des Universités de Montpellier II et Paris VII
% foot: Version 22 décembre 2001

\msp{page 69}

$\beta(c)$, for a $\Sigma$-arrow of type $\tau$, the length $n$, by inverting all the retrograde arrows (which is possible, since they are in $\Sigma$), so as to obtain a path $c^\natural$ of type $\Delta_+^n$, and by taking the composite (in $X$) of the arrows of this path $c^\natural$. It is then obvious that $\beta\alpha = \id$, we must prove that we also have $\alpha\beta = \id$, i.e. that in $\Pi_\Sigma(X)$, the arrows $u = \int c^\natural$ and $c$ are equal, i.e. $u$ and $c$ are homotopic in $\Ch_\Sigma(a,b)$. Let us consider the path $c^\natural$, of type $\Delta_+^n$: % [sic] le longueur -> la longueur, gender; possibles -> possible, number
\[
a = a_0 \xrightarrow{u_1} a_1 \xrightarrow{u_2} a_2 \longrightarrow \ \cdots \ \longrightarrow a_{n-1} \xrightarrow{u_n} a_n = b,
\]
and let $u_i$ ($i \in [1,n]$) be an arrow of this path, corresponding to a transition arrow $\underline{u}_i$ that is \textit{retrograde} in $\tau$, hence by construction
\[
u_i = v_i^{-1} \qquad \text{in } X\text{, with } v_i \in \Sigma
\]
the arrow with index $i$ in the initial path $c$ of type $\tau$. If $\underline{u}_i$, on the other hand, is direct, we will have
\[
u_i = v_i.
\]
(NB We denote by $v_i$ the $i$-arrow in the path $c$.) Now we have
\[
c = c_n * c_{n-1} * \cdots * c_1 \qquad \text{in } \Pi_\Sigma(X)
\tag*{$(*)$}
\]
and the factor $c_i$ is either the homotopy class of $\ch(v_i) = \ch(u_i)$ (the case where the transition arrow $\underline{u}_i$ of $\tau$ is direct), or that of the path \textit{opposite} to the path $\ch(v_i) : a_i \xrightarrow{v_i} a_{i-1}$ of type $\Delta_+$, hence the path $\ch(v_i)^o : a_{i-1} \xleftarrow{v_i} a_i$ of type $\Delta_-$. I say that
\[
\boxed{\ \ch(v_i)^o = \ch(u_i) \qquad \text{\textit{in} } \Pi_\Sigma(X),\ }
\tag*{(1)}
\]
\DERmark{ENV0462}

which proves that in $\Pi_\Sigma(X)$, $c$ is equal to the composite of the $\ch(u_i)$ ($i \in [1,n]$), hence to $c^\natural$. We must prove moreover that
\[
c^\natural = \ch(u) \qquad (u = u_nu_{n-1}\cdots u_1) \qquad \text{in } \Pi_\Sigma(X).
\tag*{(2)}
\]
\DERmark{ENV0463}


\msp{page 70}

So there are two verifications to make (where I again assume $\Sigma$ arbitrary):

\begin{enumerate}[label={}]
\item[$1^o$)] Let $u \in \Sigma$, $u : a \longrightarrow b$, suppose $\Sigma$ invertible in $X$. Then
\[
\ch(u)^o = \ch(u^{-1}) \qquad \text{in } \Pi_\Sigma(X),
\]
i.e. the two paths from $b$ to
\[
\begin{array}{cccl}
b & \xleftarrow{\ u\ } & a & \text{of type } \Delta_- \\
b & \xrightarrow{u^{-1}} & a & \text{of type } \Delta_+
\end{array}
\]
\printedpage{47}
% head: A. Grothendieck : Dérivateurs, Ch. VII, Édition des Universités de Montpellier II et Paris VII
% foot: Version 22 décembre 2001
are homotopic in $\Ch_\Sigma(b,a)$. Now we already know, since $\kappa : X \longrightarrow \Pi_\Sigma(X)$ is a functor, that
\[
\hot\,\ch(u^{-1}) = (\hot\,\ch(u))^{-1} \qquad \text{in } \Pi_\Sigma(X)
\]
where $\hot(\bullet)$ denotes the homotopy class of a chain (here a class in $\Ch_\Sigma(b,a)$ and in $\Ch_\Sigma(a,b)$ respectively). So the statement to prove is that $\hot(\ch(u)^o)$ is equal to the \textit{inverse} of $\hot(\ch(u))$ in $\Pi_\Sigma(X)$. This is thus what we proved for \textit{every} $u \in \Fl(X)$, in the case of $\Pi(X)$. The argument must work, so long as it makes sense, i.e. that $\ch(u)^o$ really is a $\Sigma$-path, i.e. $u \in \Sigma$.

\item[$2^o$)] Formula (2) means that for every path $c$ of type $\Delta_+^n$ (thus tautologically a \msp{page 71} $\Sigma$-path), going from $a$ to $b$, its class in $\Hom_\Sigma(a,b)$ is none other than that of $\ch(u)$, where $u$ is the composite $u_nu_{n-1}\cdots u_1$. But we have
\[
c = c_1 * c_2 * \cdots * c_n,
\]
whence in $\Pi_\Sigma(X)$
\[
\hot(c) = \hot(c_n)\hot(c_{n-1})\cdots\hot(c_1),
\]
now $\hot(c_i) = \hot\,\ch(u_i) = \kappa(u_i)$, hence
\[
\begin{aligned}
\hot(c) \ &= \ \kappa(u_n)\kappa(u_{n-1})\cdots\kappa(u_1) \\
&= \ \kappa(\underbrace{u_nu_{n-1}\cdots u_1}_{u}),
\end{aligned}
\]
since $\kappa$ is a functor, hence
\[
\hot(c) = \hot\,\ch(u),
\]

\hfill q.e.d.
\end{enumerate}

Let us now note that if $(X', \Sigma')$ is a second pair consisting of a category $X'$ and an \textit{``admissible''} set of arrows of $X'$, and if
\[
f : X \longrightarrow X'
\]
is a functor such that
\[
f(\Sigma) \subseteq \Sigma',
\]
then there exists a \textit{unique} functor
\[
\bar{f} = \Pi_\Sigma(f) : \Pi_\Sigma(X) \longrightarrow \Pi_{\Sigma'}(X')
\]
which ``extends'' $f$, i.e. such that the square
\printedpage{48}
% head: A. Grothendieck : Dérivateurs, Ch. VII, Édition des Universités de Montpellier II et Paris VII
% foot: Version 22 décembre 2001
\msp{page 72}
\[
\begin{tikzcd}
X \arrow[r, "\kappa_X"] \arrow[d, "{\scriptstyle f}"] & \Pi_\Sigma(X) \arrow[d, "{\scriptstyle \bar{f} = \Pi_\Sigma(f)}"] \\
X' \arrow[r, "\kappa_{X'}"] & \Pi_{\Sigma'}(X')
\end{tikzcd}
\]
is commutative. Uniqueness follows from the fact that on objects $\bar{f}$ must be given by $f$, and that on arrows the commutativity of the square implies that $\bar{f}$ is known on arrows of the form $\kappa_X(u)$, and thus also on the inverses of such arrows, when the arrows under consideration are invertible, hence in particular on the $\kappa_X(u)^{-1}$ for $u \in \Sigma$. Now every arrow of $\Pi_\Sigma(X)$, coming from a path, is a composite of a sequence of arrows of this form. For the existence of $\bar{f}$, we note that $f$ induces functors % [sic] les flèche -> les flèches, plural
\DERmark{ENV0464}

\[
\Ch_\tau(X; a, b) \longrightarrow \Ch_\tau(X'; f(a), f(b))
\]
which commute with the operations $\varphi^*$ on paths associated with the functors
\[
\varphi : \tau \longrightarrow \tau',
\]
whence a homomorphism of fibered categories over $\uTy$
\[
f_{a,b} : \Ch(X; a, b) \longrightarrow \Ch(X'; f(a), f(b)).
\]

\msp{page 73}

Whence a map on the $\pi_0$, i.e.
\[
\bar{f}_{a,b} : \Hom_\Sigma(a,b) \longrightarrow \Hom_\Sigma'(f(a), f(b)),
\]
\DERmark{ENV0465}

given by
\[
\bar{f}_{a,b}(\hot(c)) = \hot(f_{a,b}(c)).
\]
We obviously have, for two paths $\gamma : a \longrightarrow b$ and $\gamma' : b \longrightarrow c$,
\[
f_{a,c}(\gamma'\gamma) = f_{b,c}(\gamma')f_{a,b}(\gamma),
\]
\DERmark{ENV0466}\DERmark{ENV0467}

which proves that $\bar{f}$ is compatible with composition. It is trivial that it is compatible with identities (for $f_{a,a}$ transforms the empty path at $a$ into the empty path at $f(a)$), so we have a functor $\bar{f}$, and it really does make the square commutative.

We can now prove the universal property of $\Pi_\Sigma(X)$ and of
\[
X \xrightarrow{\ \kappa\ } \Pi_\Sigma(X),
\]
namely that every functor $f : X \longrightarrow X'$ such that $u \in \Sigma \Longrightarrow f(u)$ invertible, factors \textit{uniquely} through $\Pi_\Sigma(X)$. Uniqueness is seen as before (\footnote[20]{This is unnecessary, cf. below.}). For existence, we note that $f$ maps $\Sigma$ into the set $\Sigma'$


\endgroup
% END UNIT Der-ChVII-456/Der-ChVII-456_p046-050.tex
% BEGIN UNIT Der-ChVII-456/Der-ChVII-456_p051-055.tex
\begingroup\makeatletter%
\newcommand{\printedpage}[1]{\marginpar{\footnotesize p.~#1}}%
\newcommand{\msp}[1]{\marginpar{\footnotesize [#1]}}%
\newcommand{\Ch}{\underline{\mathrm{Ch}}}%
\newcommand{\CH}{\underline{\mathrm{CH}}}%
\newcommand{\Ty}{\underline{\mathrm{Ty}}}%
\newcommand{\Hom}{\underline{\mathrm{Hom}}}%
\newcommand{\red}{\mathrm{red}}%
\newcommand{\uHom}{\underline{\mathrm{Hom}}}%
\makeatother\let\printedpage\derpp\let\msp\dermsp%
\let\mspin\DERmspin%


\setcounter{footnote}{20}

\printedpage{49}
% head: A. Grothendieck : Dérivateurs, Ch. VII, Édition des Universités de Montpellier II et Paris VII
\msp{page 74}
of the invertible arrows of $X'$ (which set does indeed satisfy the two admissibility conditions). Thus by what precedes we have a commutative diagram
\[
\begin{tikzcd}
X \ar[r,"\kappa"] \ar[d,"f"'] & \Pi_{\Sigma}(X) \ar[d,"\bar f = \Pi_{\Sigma}(f)"] \ar[dl,dotted,"\tilde f"'] \\
X' \ar[r,"\kappa'","\sim"'] & \Pi_{\Sigma'}(X'),
\end{tikzcd}
\]
where $\bar f$ is uniquely determined. Since $\kappa'$ is an isomorphism (as we saw earlier, since $\Sigma'$ \texttt{[is]} made up of the invertible arrows), the existence and uniqueness of $\bar f$ is equivalent to the announced conclusion,

\hfill q.e.d.

\textbf{One more sorites} on the homotopy classes of elements of a category $X$. In general, one sees that $X$ and $\Pi_{\Sigma}(X) = X\Sigma^{-1}$ have the same connected components, for if
\msp{page 75}
$a$ and $b$ in $X$ are joined by an arrow of $X\Sigma^{-1}$, i.e. by a $\Sigma$-path, then they are in the same connected component of $X$. But the converse is also true, as one sees e.g. by taking $\Sigma = \mathrm{Fl}(X)$, and by saying that in the \textit{groupoid} $\Pi(X)$, where two vertices are in the same \texttt{[connected]} component if and only if $\exists\ x \longrightarrow y$ (a relation which is not only transitive and reflexive, but also symmetric in a groupoid). A path
\[
a \xrightarrow{\quad c \quad} b
\]
will also be called a \textit{homotopy} from $a$ to $b$. We say that $a$ and $b$ are \textit{homotopic} if there exists a homotopy $c$ from $a$ to $b$, in which case there also exists an ``inverse'' homotopy $c^{o}$ from $b$ to $a$. This also means that $a$ and $b$ are in the same connected component of $X$, i.e. define the same element of $\pi_0(X)$.

The homotopy will be called \textit{simple} if $c$ is of type $\Delta_{+}$ or $\Delta_{-}$, i.e. if it is given by an arrow $u : a \longrightarrow b$ or $v : b \longrightarrow a$. If one is looking for a homotopy between $a$ and $b$, one will try to find it as ``short'' as possible. If one has a homotopy $c$ of length $l$ from $a$ to $b$, then on replacing $\tau$ by $\tau_{\red}$, hence $I_{\tau}$ by $I_{\tau\,\red}$, the corresponding reduced
\msp{page 76}
subinterval, one finds a homotopy from $a$ to $b$ of type $\tau_{\red}$, whose length is equal to the \textit{``useful length''} of the type $\tau$ of $c$ (which justifies the terminology). Let us note that if
\[
\tau = \Delta_{\varepsilon}^{l_1}\Delta_{-\varepsilon}^{l_2}\Delta_{\varepsilon}^{l_3}\cdots\Delta_{(-1)^{r-1}\varepsilon}^{l_r},
\]
\DERmark{ENV0468}

then
\[
\tau_{\red} = \Delta_{\varepsilon}\Delta_{-\varepsilon}\Delta_{\varepsilon}\cdots\Delta_{(-1)^{r-1}\varepsilon},
\]
\DERmark{ENV0469}

% foot: Version 22 décembre 2001
\printedpage{50}
% head: A. Grothendieck : Dérivateurs, Ch. VII, Édition des Universités de Montpellier II et Paris VII
\texttt{[rather twice} $(-1)^{r-1}$\texttt{]} where $r$ is the useful length of $\tau$ (the not useful length being $n = \sum\limits_{i\in[1,r]} l_i$).

When one starts with a homotopy between \textit{functors}
\[
X \longrightarrow Y,
\]
it is understood that we are dealing with homotopies in the category $\uHom(X,Y)$. (Or, depending on the context, in a specified subcategory of $\uHom(X,Y)$.)

\msp{page 77}
On page 61, I invoked the existence of a \textit{functor} of \textit{``composition of paths''}
\[
\Ch_{\Sigma}(X;a,b)\times\Ch_{\Sigma}(X;b,c) \ \longrightarrow\ \Ch_{\Sigma}(X;a,c),
\]
which on the fiber of the first member over $(\tau,\tau') \in \operatorname{Ob}(\Ty\times\Ty)$, would be given by
\[
\Ch_{\Sigma,\tau}(X;a,b)\times\Ch_{\Sigma,\tau'}(X;b,c) \ \longrightarrow\ \Ch_{\Sigma,\tau\tau'}(X;a,c), \tag*{$(*)$}
\]
the composition of paths of specified types $\tau$, $\tau'$. That we have here a \textit{functor,} not just a map on the sets of objects, is fairly obvious, and can be expressed for example with the help of the amalga\texttt{[ma]}tion diagram
\[
\begin{tikzcd}[column sep=small]
 &  \{1_{\tau}=0_{\tau'}\} \ar[dl,hook'] \ar[dr,hook] & \\
I_{\tau} \ar[dr,hook] & & I_{\tau'} \ar[dl,hook'] \\
 & I_{\tau\tau'} &
\end{tikzcd}
\]
\DERmark{ENV0470}\DERmark{ENV0471}\DERmark{ENV0472}

by transforming it by $\uHom(\cdot,X)$ (instead of transforming it by $\mathrm{Hom}_{\mathrm{Cat}}(\cdot,X)$ only), at least in the particular case where $\Sigma = \mathrm{Fl}(X)$ (which was the one on page 61). But the passage to arrows is fairly obvious - for two composites $c'_1\circ c_1$ and $c'_2\circ c_2$, from
\[
\begin{tikzcd}[column sep=small]
a \ar[rr,"c_1"] \ar[d,equal] & {} \ar[d,"w"'] & b \ar[rr,"c'_1"] \ar[d,equal] & {} \ar[d,"w'"] & c \ar[d,equal] \\
a \ar[rr,"c_2"] & {} & b \ar[rr,"c'_2"] & {} & c
\end{tikzcd}
\]
\texttt{[one]} sees that the homomorphisms $w''$ from the first to the second, in $\Ch_{\Sigma}(X;a,c)$, which are the identity on the vertex of index $n$ of the two paths (where $n = \operatorname{long}\tau$), corresponding to $1_{\tau}$ and to $0_{\tau'}$, ``are'' exactly the pairs
\DERmark{ENV0473}

% foot: Version 22 décembre 2001
\printedpage{51}
% head: A. Grothendieck : Dérivateurs, Ch. VII, Édition des Universités de Montpellier II et Paris VII
\msp{page 78}
formed of a homomorphism $w : c_1 \longrightarrow c_2$ and $w' : c'_1 \longrightarrow c'_2$ in $\Ch_{\Sigma,\tau}(a,b)$ and $\Ch_{\Sigma,\tau'}(b,c)$ respectively, and if $w''$, $w$, $w'$ correspond in this way, we shall set
\[
w'' = w * w'.
\]
There is an obvious transitivity to check to see that this defines a \textit{functor} $(*)$. We can express the situation with the help of the diagram
\[
\begin{tikzcd}[column sep=small]
 & \Ty\times\Ty \ar[rrr,"\pi\ (\text{composition})"] \ar[dl,"p_1"'] \ar[dr,"p_2"] & & & \Ty \\
\Ty & & \Ty, & &
\end{tikzcd}
\]
where $p_1$, $p_2$ are the two projections. We consider over $\Ty\times\Ty$ the three fibered categories % [sic] fibrés -> fibrées, agreement with "catégories"
\[
p_1^*(\CH_{\Sigma,\tau}(a,b)),\ \ p_2^*(\CH_{\Sigma,\tau'}(b,c)),\ \ \pi^*(\Ch_{\Sigma,\tau\tau'}(a,c))
\]
\texttt{[rather} $\pi^*(\CH_{\Sigma,\tau\tau'}(a,c)$\texttt{]}, whose fibers at $(\tau,\tau')$ are respectively % [sic] \pi^*(\CH_{\Sigma,\tau\tau'}(a,c) -> \pi^*(\CH_{\Sigma,\tau\tau'}(a,c)), one closing parenthesis missing
\[
\Ch_{\Sigma,\tau}(a,b),\ \ \Ch_{\Sigma,\tau'}(b,c),\ \ \Ch_{\Sigma,\tau\tau'}(a,c).
\]
(We write $\CH$ instead of $\Ch$, to indicate that we regard the categories under consideration as (fibered) categories \textit{over} $(\Ty)$ (\footnote{But for $p_1^*$, $p_2^*$ and $\pi^*$ these are ordinary Cartesian products in Cat, there is no need to be afraid and to write $\CH$'s \ldots}).) Over $\Ty\times\Ty$, we have a (Cartesian) homomorphism of fibered categories
\[
p_1^*(\CH_{\tau}(a,b))\times_{\Ty\times\Ty} p_2^*(\CH_{\tau'}(b,c)) \ \longrightarrow\ \pi^*(\CH_{\tau\tau'}(a,c)) \tag*{$(**)$}
\]
\DERmark{ENV0474}\DERmark{ENV0475}

\texttt{[rather} $p_1^*(\CH_{\tau}(a,b))\times_{\Ty\times\Ty} p_2^*(\CH_{\tau'}(b,c)) \ \longrightarrow\ \pi^*(\CH_{\tau\tau'}(a,c))$\texttt{]},
\msp{page 79}
for it is OK on the fibers ($(*)$ page 77), and if
\[
(\varphi,\varphi') : (\tau_1,\tau'_1) \longrightarrow (\tau,\tau') \qquad (\varphi : \tau_1 \longrightarrow \tau,\ \varphi' : \tau'_1 \longrightarrow \tau')
\]
is an arrow in $\Ty\times\Ty$, we deduce from it a commutative diagram
\[
\begin{tikzcd}
\Ch_{\Sigma,\tau}(a,b)\times\Ch_{\Sigma,\tau'}(b,c) \ar[r,"\text{comp.}"] \ar[d,"\varphi^*_{a,b}\times(\varphi'_{b,c})^*"] & \Ch_{\Sigma,\tau\tau'}(a,c) \ar[d,"(\varphi*\varphi')^*_{a,c}"] \\
\Ch_{\Sigma,\tau_1}(a,b)\times\Ch_{\Sigma,\tau'_1}(b,c) \ar[r,"\text{comp.}"] & \Ch_{\Sigma,\tau_1\tau'_1}(a,c)
\end{tikzcd}
\]
\DERmark{ENV0476}

% foot: Version 22 décembre 2001
\printedpage{52}
% head: A. Grothendieck : Dérivateurs, Ch. VII, Édition des Universités de Montpellier II et Paris VII
\texttt{[rather} $\Ch_{\Sigma,\tau\tau'}(a,c)$\texttt{]}, i.e. the composition of paths with fixed endpoints is compatible with the inverse images of paths under arrows between the types of paths.

But the Cartesian homomorphism $(**)$ of fibered categories over $\Ty\times\Ty$ is not quite the same thing as an innocuous functor
\[
\Ch_{\Sigma}(a,b)\times\Ch_{\Sigma}(b,c) \ \longrightarrow\ \Ch_{\Sigma}(a,c). \tag*{$(*)$}
\]
The formula $(*)$ defines in any case a map on the sets of objects, since the sets of objects of each of the three categories under consideration decompose as a sum of sets $\Ch_{\Sigma,\tau}(\bullet,\bullet)$, where $\tau$ runs through $\mathrm{Ty} = \operatorname{Ob}\Ty$. But does this extend in a natural way to morphisms?

\msp{page 80}
But now I remember that this famous Cartesian product $p_1^*(\text{\textit{fourbi}})\times p_2^*(\text{thingamajig})$ is none other than the plain old Cartesian product
\DERmark{ENV0477}

\[
\Ch_{\Sigma}(a,b)\times\Ch_{\Sigma}(b,c),
\]
sent to $\Ty\times\Ty$ by the Cartesian product of the two projection functors
\[
\begin{aligned}
p_{a,b} : \Ch_{\Sigma}(a,b) &\ \longrightarrow\ \ \Ty\\
p_{b,c} : \Ch_{\Sigma}(b,c) &\ \longrightarrow\ \ \Ty,
\end{aligned}
\]
hence
\[
\begin{tikzcd}
\Ch_{\Sigma}(a,b)\times\Ch_{\Sigma}(b,c) \ar[d,"p_{a,b}\times p_{b,c}"] \\
\Ty\times\Ty
\end{tikzcd}
\]
\msp{page 81}
which is well and truly fibrant (as a Cartesian product of two fibrant functors) and has the right fibers! So we do find, by the general sorites on fibrant functors (here over $\Ty\times\Ty$) a functor of fibered categories over $\Ty\times\Ty$
\[
\Ch_{\Sigma}(a,b)\times\Ch_{\Sigma}(b,c) \ \longrightarrow\ \qquad \underbrace{\pi^*(\CH_{\Sigma}(a,c))}_{=\,(\Ty\times\Ty)\times_{\Ty}\Ch_{\Sigma}(a,c)\ \ \text{(\footnotemark)}} \qquad .
\]
\footnotetext{fiber product along $\pi : \Ty\times\Ty \longrightarrow \Ty$.}
Since on the other hand we have a functor from the last member (fiber product) to the factor $\Ch_{\Sigma}(a,c)$, we win. (We would have found it again by the direct argument I was about to make, and which I crossed out - but I am happier this way!)

% foot: Version 22 décembre 2001
\printedpage{53}
% head: A. Grothendieck : Dérivateurs, Ch. VII, Édition des Universités de Montpellier II et Paris VII
To do things properly, we would have to prove the \textit{associativity} of the composition functors, not just at the level of objects, by arguments like this, but it comes down simply to a relation of the type
\[
(w''*w')*w = w''*(w'*w)
\]
for a system of three arrows in
\[
\Ch_{\Sigma}(a,b),\ \Ch_{\Sigma}(b,c)\ \ \text{and}\ \ \Ch_{\Sigma}(c,d)
\]
respectively. It can only be trivially
\msp{page 82}
true, but tedious to spell out.

The cold sweats are over!

In summary, the formalism of the $\Ch_{\Sigma}(X)$ and of the $\Ch_{\Sigma}(X;a,b)$ developed up to now gives complete satisfaction on the following points:

\begin{enumerate}
\item[a)] Good formalism for the composition of paths (which was disgusting before).

\item[b)] Elegant treatment of the construction of categories of fractions $\mathcal{M}\Sigma^{-1}$. But I expect that it will be less pretty than before, when it comes to replacing an $(\mathcal{M},\Sigma)$ by $(\ \underbrace{\mathcal{M}(X)}_{=\,\uHom(X^{o},\mathcal{M})}\ ,\Sigma_X)$'s ($X$ an index category) \texttt{[rather} $\uHom(X^{o},\mathcal{M})$\texttt{]}. It is to remedy this that I will have to develop comparison statements, between the category $\Ch_{\Sigma}$ constructed here, and others whose type in Hot lends itself to computations in a more direct way, as in the original state of the question. Moreover, we will surely have to turn to Ind-categories, in order to be able to follow the variance of $\mathcal{M}(X)\Sigma_X^{-1}$ for variable $X$.

\item[c)] Good formalism for ``homotopies'' and for compositions of homotopies (which is another way of speaking of paths and of their composition).
\end{enumerate}
\DERmark{ENV0478}

% foot: Version 22 décembre 2001


\endgroup
% END UNIT Der-ChVII-456/Der-ChVII-456_p051-055.tex
% BEGIN UNIT Der-ChVII-456/Der-ChVII-456_p056-060.tex
\begingroup\makeatletter%
\newcommand{\printedpage}[1]{\marginpar{\footnotesize p.~#1}}%
\newcommand{\msp}[1]{\marginpar{\footnotesize [#1]}}%
\newcommand{\uHom}{\underline{\mathrm{Hom}}}%
\newcommand{\Ty}{\mathrm{Ty}}%
\newcommand{\uTy}{\underline{\mathrm{Ty}}}%
\newcommand{\ul}[1]{\underline{#1}}%
\newcommand{\red}{\text{\upshape red}}%
\newcommand{\triv}{\mathrm{triv}}%
\makeatother\let\printedpage\derpp\let\msp\dermsp%
\let\mspin\DERmspin%


\printedpage{54}
% head: A. Grothendieck : Dérivateurs, Ch. VII, Édition des Universités de Montpellier II et Paris VII
% foot: Version 22 décembre 2001

\msp{page 83}

\section*{8\quad The homotopy lemma for morphisms of intervals}

It is the following:

\textbf{Lemma 1.} \textit{Let $\tau,\tau'\in\mathrm{Ty}$, and let $f,g:\tau\rightrightarrows\tau'$ be two morphisms in $\uTy$. Then $f$, $g$ are homotopic in the category $\uHom\,\underline{\mathrm{int}}(I_\tau,I_{\tau'})$ of increasing morphisms of oriented intervals from $I_\tau$ to $I_{\tau'}$.}

This category could legitimately be denoted by $\uHom_{\uTy}(\tau,\tau')$, or $\uHom(\tau,\tau')$. It is a subcategory of the category $\uHom(I_\tau,I_{\tau'})$ of \textit{all} functors from $I_\tau$ to $I_{\tau'}$, and a full subcategory of the category $\uHom^\varepsilon(I_\tau,I_{\tau'})$ of functors $f$ that are compatible with the datum of the sources and targets, i.e. that make the diagram commute
\[
\begin{tikzcd}[column sep=small]
 & \varepsilon=\{0,1\} \arrow[dl, "{(\sigma_\tau,\beta_\tau)}"'] \arrow[dr,  "{(\sigma_{\tau'},\beta_{\tau'})}"] & \\
I_\tau \arrow[rr, "f"] & & I_{\tau'}.
\end{tikzcd}
\]
\DERmark{ENV0479}\DERmark{ENV0480}


If $f$, $g$ are in this category (not necessarily increasing morphisms, by the way), an arrow from $f$ to $g$ is thus the datum, for every $i\in I_\tau$, of an arrow $f(i)\xrightarrow{w_i}g(i)$ in $I_{\tau'}$, with the two conditions:

\begin{enumerate}[label={}, leftmargin=*]
\item[a)] If $i\in\{0_\tau,1_\tau\}$, $w_i=\mathrm{id}$, and
\item[b)] the usual commutative square for an arrow $i\longrightarrow j$ in $I_\tau$. (It suffices to have this compatibility for the transition arrows $\ul{u}_i:i-1\longrightarrow i$ if $\tau_i=+$, or $\ul{u}_i:i\longrightarrow i-1$ if $\tau_i=-$, in $I_\tau$ ($i\in[1,\operatorname{long}\tau]$).)
\end{enumerate}
\DERmark{ENV0481}


Another way of describing this category $\uHom^\varepsilon$ is to say that it is the category
\msp{page 84}
\[
\underline{\mathrm{Ch}}_\tau(I_{\tau'};0_{\tau'},1_{\tau'})\qquad(\simeq\uHom^\varepsilon(I_\tau,I_{\tau'})),
\]
but $\uHom_{\uTy}(\tau,\tau')$ is the full subcategory of this category consisting of the \textit{increasing paths} in $I_{\tau'}$. The homotopy lemma also means that this category $\uHom_{\uTy}(\tau,\tau')$ is always \textit{connected.} It is not necessarily 0-connected, however, for it can be empty. Let us note here, for later, the

\textbf{Proposition 1:} \textit{Let $\tau,\tau'\in\mathrm{Ty}$, let $r$, $r'$ be their useful lengths. For $\mathrm{Hom}_{\uTy}(\tau,\tau')$ to be nonempty, it is necessary and sufficient that one or the other of the following two conditions be satisfied:}

\begin{enumerate}[label={}, leftmargin=*]
\item[\textit{a)}] $r\geq r'+1$

\printedpage{55}
% head: A. Grothendieck : Dérivateurs, Ch. VII, Édition des Universités de Montpellier II et Paris VII
% foot: Version 22 décembre 2001

\item[\textit{b)}] \textit{$r=r'$, and (if $r\neq0$) $\tau$ and $\tau'$ have the same initial direction (so that $\tau_{\red}=\tau'_{\red}$).}
\end{enumerate}

(Proof $\pm$ immediate.)

Let us note that if $\tau$, $\tau'$ do not have the same initial direction, then every morphism $\tau\longrightarrow\tau'$ is necessarily degenerate on the first critical interval of $I_\tau$. We will thus have $\tau=\Delta_\varepsilon^{l_0}\cdot\tau_0$ (where $l_0$ is the length of the first critical interval, and where $\tau_0$ has \texttt{[the]} initial direction opposite to that of $\tau$), and an isomorphism
\[
\uHom(\tau,\tau')\simeq\uHom(\tau_0,\tau'),
\]
where this time $\tau_0$, $\tau'$ have the same initial direction. We will therefore assume henceforth that \textit{the initial direction of $\tau$, $\tau'$ is the same.} The condition $\uHom(\tau,\tau')\neq\emptyset$ is now equivalent to $\underline{r\geq r'}$, which we can assume. I will then describe a canonical element
\msp{page 85}
$F_0$ of $\uHom(\tau,\tau')$, as follows. Let
\[
I_1,I_2,\dots,I_r
\]
be the critical intervals of $I_\tau$, and
\[
I'_1,I'_2,\dots,I'_{r'}\qquad(r'\leq r)
\]
those of $I_{\tau'}$. The functor $F_0$ will come from morphisms of oriented intervals (such as
\[
\left\{
\begin{array}{ccc}
I_1 & \xrightarrow{f_1} & I'_1\\[4pt]
I_2 & \xrightarrow{f_2} & I'_2\\[4pt]
 & \dots & \\[4pt]
I_{r'} & \xrightarrow{f_{r'}} & I'_{r'}
\end{array}
\right.\qquad)
\]
and will therefore necessarily be trivial on the critical intervals $I_\alpha$ for $r'<\alpha\leq r$ (there are some if and only if $r>r'$). Since $I_\tau$, $I_{\tau'}$ have the same initial direction, it follows that for $\alpha\in[1,r']$, $I_\alpha$ and $I'_\alpha$ have the same signature (namely $\varepsilon_\alpha=(-1)^{\alpha+1}\varepsilon_1$, where $\varepsilon_1$ is the common signature of $I_1$ and $I'_1$). Thus
\[
I_\alpha=\Delta^{l_\alpha}_{\varepsilon_\alpha},\qquad I'_\alpha=\Delta^{l'_\alpha}_{\varepsilon_\alpha},
\]
and we define $f_\alpha$ as
\[
f_\alpha=\mathfrak{f}_{\varepsilon_\alpha,l'_\alpha}*\triv_\alpha,
\]
where
\[
\mathfrak{f}_{\varepsilon_\alpha,l'_\alpha}:\Delta_{\varepsilon_\alpha}\longrightarrow\Delta^{l'_\alpha}_{\varepsilon_\alpha}
\]
and where
\[
\triv_\alpha:\Delta^{l_\alpha-1}_{\varepsilon_\alpha}\longrightarrow1. % [sic] l_{\alpha-1} -> l_\alpha-1, the exponent is printed as l with subscript \alpha-1, while I_\alpha = \Delta^{l_\alpha} = \Delta * \Delta^{l_\alpha-1}
\]
\DERmark{ENV0482}\DERmark{ENV0483}

In other words, for $\varepsilon=+1$ (e.g.) we have to define a morphism of the standard intervals $\underbrace{\boldsymbol{\Delta}^+_{l_\alpha}}_{=[0,l_\alpha]}\longrightarrow\underbrace{\boldsymbol{\Delta}^+_{l'_\alpha}}_{=[0,l'_\alpha]}$, and we set (with the numberings internal to $\boldsymbol{\Delta}_{l_\alpha}$, $\boldsymbol{\Delta}_{l'_\alpha}$)
\[
f_\alpha(0)=0,\qquad f_\alpha(i)=l'_\alpha\ \text{ for }i\in[1,l_\alpha].
\]

\printedpage{56}
% head: A. Grothendieck : Dérivateurs, Ch. VII, Édition des Universités de Montpellier II et Paris VII
% foot: Version 22 décembre 2001

\msp{page 86}
The idea then, is to be in the frame of mind of a traveler who, bound to a travel plan in stages described by $\tau$, would like to arrive as quickly as possible at $1_{\tau'}$, starting from $0_{\tau'}$. He will, from the very first stage $\Delta_\varepsilon\subseteq I_1$, go to the end of $I'_1$, regretting not being able to go farther (because of the travel plan). He has to remain inactive until he has reached the moment marked by the final endpoint of $I_1$ (if this is beyond that of his first transition arrow, i.e. if $l_1>1$), then he takes his second big step to the terminal endpoint of $I'_2$ etc. He can hope, at each moment of his journey, to be ahead of (or at least not behind) the other travelers, who might have followed different tactics for advancing. If this hope were justified, that would mean that $F_0$ would be a greatest element of $\uHom(\tau,\tau')$ (\textbf{NB} all these categories $\uHom$ between intervals are ``ordered'', i.e. there is at most one arrow between two objects), i.e. a final object. This would imply that $\uHom(\tau,\tau')$ would even be contractible - something perhaps true, but surely not in so simplistic a way I fear. I have botched it, by the way: $F_0$ probably is indeed a ``greatest'' element, but for the order defined by the numbering of the vertices of $I_{\tau'}$ (so $F\leq G$ if and only if $F(i)\leq G(i)$ $\forall i\in I_\tau$), and not at all for the order defined by the category structure of $\uHom(\tau,\tau')$.

\msp{page 87}
The idea for proving the connectedness of $\uHom(\tau,\tau')$ is to bring any $F\in\mathrm{Hom}(\tau,\tau')$, by a sequence of simple homotopies, to the normal form. Let us recall that $F$ is given by an increasing sequence of transition arrows
\[
\ul{u}_{i_1},\ \ul{u}_{i_2},\ \dots\ ,\ \ul{u}_{i_s},\qquad i_1<i_2<\cdots<i_s,
\]
and a sequence of integers $\geq1$
\[
\lambda_1,\ \lambda_2,\ \dots\ ,\ \lambda_s
\]
with the condition that
\[
\Delta^{\lambda_1}_{\varepsilon_{i_1}}\Delta^{\lambda_2}_{\varepsilon_{i_2}}\cdots\Delta^{\lambda_s}_{\varepsilon_{i_s}}=\tau'. \tag*{$(*)$}
\]
The $\ul{u}_{i_\alpha}$ are the transition arrows of $I_\tau$ \textit{critical} for $F$, i.e. nondegenerate for $F$, and $\lambda_\alpha$ is the \textit{part} of $\ul{u}_{i_\alpha}$ for $F$. The normalized (``hurried'') functor $F_0$ is distinguished by the following conditions a) and b) \texttt{[and c)]}. First of all:

\begin{enumerate}[label={}, leftmargin=*]
\item[a)] The sequence of signs $\varepsilon_{i_1},\varepsilon_{i_2},\dots,\varepsilon_{i_s}$ is ``alternating'', i.e. $\varepsilon_{i_\alpha}\neq\varepsilon_{i_{\alpha+1}}$ ($\alpha\in[1,s-1]$). (This expresses the fact that the hurried traveler does not do in several steps what he could do in a single one.)

Since the first of these signs is necessarily $\varepsilon'_1=\varepsilon_1$, this means that the sequence of signs is the same as that of the $\varepsilon_1$, $\varepsilon_2$, \dots, $\varepsilon_{r'}$ (apart from the length). But better still, that the expression $(*)$ for $\tau'$ is its unique reduced expression, hence that $s=r'$ (useful length of $\tau'$), and that the $\lambda_i$ are the lengths $l'_i$ of the critical intervals of $I_{\tau'}$.

\printedpage{57}
% head: A. Grothendieck : Dérivateurs, Ch. VII, Édition des Universités de Montpellier II et Paris VII
% foot: Version 22 décembre 2001

\msp{page 88}
Let then $\alpha\in[1,s=r']$, and let us consider the $F$-critical transition arrow $\ul{u}_{i_\alpha}$, let $j_\alpha$ be the index marking the critical interval of $I_\tau$ containing the part of $\ul{u}_{i_\alpha}$ in $I_\tau$,
\[
\ul{u}_{i_\alpha}\in\mathrm{Fl}(I_{j_\alpha}),
\]
then $F$ induces (increasing) morphisms of oriented intervals
\[
F_\alpha=F|I_{j_\alpha}:I_{j_\alpha}\longrightarrow I'_\alpha
\]
and is trivial on the other critical intervals (those not of the form $I_{j_\alpha}$). The datum of the $F_\alpha$ thus determines $F$ on the $I_{j_\alpha}$, hence on \textit{all} the $I_j$, hence everywhere (by the operation $*$ of composition by gluing). Moreover condition a) ensures that each $F_\alpha$ is of a rather special nature: a single critical transition arrow $\ul{u}_{i_\alpha}$. Moreover, for us to have $F=F_0$, it will be necessary:

\item[b)] $\forall\alpha\in[1,r']$, $\ul{u}_{i_\alpha}$ is the \textit{first} transition arrow of the critical interval $I_{j_\alpha}$ that contains it.

And finally:

\item[c)] The sequence of critical intervals of $I_\tau$
\[
I_{j_1},\ I_{j_2},\ \dots\ ,\ I_{j_{r'}}\qquad(\text{of alternating signs }\varepsilon_1,\ -\varepsilon_1,\ \varepsilon_1,\ \ldots)
\]
associated with the $F$-critical arrows, is none other than that
\[
I_1,\ I_2,\ \dots\ ,\ I_{r'}
\]
of the first $r'$ critical intervals of $I_\tau$, i.e.
\[
(j_1,j_2,\dots,j_{r'})=(1,2,\dots,r').
\]
\end{enumerate}

\msp{page 89}
This being seen, the natural idea for ``deforming'' $F$ into $F_0$, is to find a sequence of homotopies that transform $F$ so as to make it acquire successively the characteristics a), b), c) of the hurried traveler. If need be, to get to the goal faster, we will use the induction hypothesis on $n=\operatorname{long}\tau$, that the connectedness theorem is already proved when $I$ is of length $\leq n-1$. (NB This theorem is trivial for $n\leq1$, for then the categories $\uHom(\tau,\tau')$ are all either empty, or one-point.)

\msp{page 90}
1) We want to ensure a). For this, we use the induction hypothesis, to replace a group of \texttt{[illegible word]} consecutive critical arrows of the same sign, by a single critical arrow. This method only breaks down if the critical arrows are all of the same sign $\varepsilon$ and include $\ul{u}_1$ and $\ul{u}_n$. But the induction hypothesis applied to $I=[0,n-1]\subseteq I_\tau$ reduces us to the case where there are only two critical arrows $\ul{u}_1$ and $\ul{u}_n$. If the first critical interval is not reduced to $\Delta_\varepsilon$, i.e. is of length at least 2, then the induction hypothesis applied to $I=[1,n]\subseteq I_\tau$ reduces us to the case where the two critical arrows are $\ul{u}_1$ and $\ul{u}_2$, and we are
\DERmark{ENV0484}\DERmark{ENV0485}

\printedpage{58}
% head: A. Grothendieck : Dérivateurs, Ch. VII, Édition des Universités de Montpellier II et Paris VII
% foot: Version 22 décembre 2001
\msp{page 91}
then reduced to the case where $\boxed{I=\Delta^2_\varepsilon,\ \tau'=\Delta^2_\varepsilon.}$ Otherwise, to the case where $\boxed{\tau=\Delta_\varepsilon\Delta^{l_2}_{-\varepsilon}\Delta_\varepsilon,\ \tau'=\Delta^2_\varepsilon,}$ and where the two critical arrows are again the first and the last, and $\tau'=\Delta^2_\varepsilon$. Those are the two irreducible cases, which will have to be treated separately.

These cases granted, we thus have a way to deform $F$ so as to satisfy a). Next we want:
\DERmark{ENV0486}


2) To ensure that the sequence of critical intervals $I_{j_1}$, $I_{j_2}$, \dots, $I_{j_{r'}}$ which contains the critical arrows $\ul{u}_{i_1}$, \dots, $\ul{u}_{i_{r'}}$, is the sequence $I_1$, \dots, $I_{r'}$. This is again immediate by the induction hypothesis, the only irreducible case being that where $r'=1$ (a single critical arrow), and where it is a matter of bringing this critical arrow down to $\ul{u}_1$. If $I_1$ is of length $>1$, then we bring it down to $\ul{u}_2$ by the induction hypothesis, and the irreducible case is that where $\boxed{\tau=\Delta^2_\varepsilon,\ \tau'=\Delta_\varepsilon,}$ a critical arrow $\ul{u}_2$, bring it down to $\ul{u}_1$. If on the other hand $I_1$ \texttt{[is]} of length 1, the irreducible case is $\boxed{\tau=\Delta_\varepsilon\Delta^{l_2}_{-\varepsilon}\Delta_\varepsilon,\ \tau'=\Delta_\varepsilon,}$ a critical arrow $\ul{u}_{l_2+2}$, bring it down to $\ul{u}_1$.

Subject to treating the two new irreducible cases, we can deform $F$ into a functor that satisfies a) and c). It remains to normalize each of the $F_\alpha:I_\alpha\longrightarrow I'_\alpha$ ($\alpha\in[1,r']$). Thus we are reduced to $\Delta^n_\varepsilon\longrightarrow\Delta_\varepsilon$: a single critical arrow $\ul{u}_\alpha$, $\alpha\in[1,n]$, to be brought down to $\ul{u}_1$. If $n\geq2$, by the induction hypothesis it comes down to $\ul{u}_2$, case assumed treated in $2^o$). So
\msp{page 92}
we win, subject to treating the 4 irreducible cases.

A) Case $\tau=\Delta^2_\varepsilon$, $\tau'=\Delta^2_\varepsilon$ or $\Delta_\varepsilon$.

The category $\uHom(\Delta^2_\varepsilon,\Delta_\varepsilon)$ has two elements, the morphisms $\Delta^2_\varepsilon\longrightarrow\Delta_\varepsilon$ correspond to the transition arrows of $\Delta^2_\varepsilon$ ($=\boldsymbol{\Delta}^+_2$ or $\boldsymbol{\Delta}^-_2$). In the case $\varepsilon=+1$, these are the two maps from $\{0,1,2\}$ to $\{0,1\}$, which map 0 to 0, 2 to 1, and 1 either to 1 ($F_0$), or to 0 \texttt{[rather to 0]} ($F_1$).
\DERmark{ENV0487}

\[
\begin{tikzcd}
0 \arrow[r] & 1 \arrow[r] & 2
\end{tikzcd}
\]
\[
\begin{tikzcd}
0 \arrow[r] & 1
\end{tikzcd}
\]

We obviously have an arrow $F_1\longrightarrow F_0$ (the hurried traveler arrives faster \ldots). Similar for $\varepsilon=-1$,
\[
\uHom(\Delta^2_\varepsilon,\Delta_\varepsilon)\simeq\boldsymbol{\Delta}^1,\qquad\text{0-connected category.}
\]

The category $\uHom(\Delta^2_\varepsilon,\Delta^2_\varepsilon)$ has three elements, identified with the set of nonempty subsets of the two-element set $\mathrm{Tr}(\tau)$ of transition arrows $\ul{u}_1$, $\ul{u}_2$ of $\tau$: there are $F_{\ul{u}_1}=F_0$, $F_{\ul{u}_2}$, $F_{\ul{u}_1,\ul{u}_2}$. If $\varepsilon=+1$, these are the increasing maps (nothing more) from $\{0,1,2\}$ to $\{0,1,2\}$, mapping 0 to 0, 2 to 2, and 1 to whatever it likes. % [sic] parties non vide -> parties non vides, agreement


\endgroup
% END UNIT Der-ChVII-456/Der-ChVII-456_p056-060.tex
% BEGIN UNIT Der-ChVII-456/Der-ChVII-456_p061-065.tex
\begingroup\makeatletter%
\newcommand{\printedpage}[1]{\marginpar{\footnotesize p.~#1}}%
\newcommand{\msp}[1]{\marginpar{\footnotesize [#1]}}%
\newcommand{\uHom}{\underline{\mathrm{Hom}}}%
\newcommand{\uCh}{\underline{\mathrm{Ch}}}%
\newcommand{\uTy}{\underline{\mathrm{Ty}}}%
\makeatother\let\printedpage\derpp\let\msp\dermsp%
\let\mspin\DERmspin%


\printedpage{59}
% head: A. Grothendieck : Dérivateurs, Ch. VII, Édition des Universités de Montpellier II et Paris VII
% foot: Version 22 décembre 2001

\begin{tabular}{lll}
Thus & $1\longmapsto 2$ & corresponds to \quad $F_0$\\[1ex]
 & $1\longmapsto 1$ & corresponds to \quad $F_1$\\[1ex]
 & $1\longmapsto 0$ & corresponds to \quad $F_2$
\end{tabular}

and of course we have $F_2\longrightarrow F_1\longrightarrow F_0$, % [sic] bien sur -> bien sûr, accent missing
\[
\uHom(\Delta^2_\varepsilon,\Delta^2_\varepsilon)\simeq\boldsymbol{\Delta}^2,\qquad\quad\text{0-connected category.}
\]

In these two cases, two elements of $\mathrm{Hom}(\tau,\tau')$ are joined by a homotopy which is \textit{simple.}

\msp{page 93}

B) Case $\tau=\Delta_\varepsilon\Delta^l_{-\varepsilon}\Delta_\varepsilon$, $\tau'=\Delta^2_\varepsilon$ or $\Delta_\varepsilon$.

Let us note that we have, for arbitrary $\tau,\tau'\in\uTy$, an isomorphism of categories
\[
\uHom(\tau,\tau')\longrightarrow\uHom(\tau^o,\tau'^o),
\]
\DERmark{ENV0488}\DERmark{ENV0489}

simply expressing that
\[
\tau\longmapsto\tau^o
\]
is a \texttt{[2]}-functor from $\uTy$ to $\uTy$, and even an automorphism of this \texttt{[2]}-category. Thus to say that $\uHom(\tau,\tau')$ is connected amounts to saying that $\uHom(\tau^o,\tau'^o)$ is. Thus we can suppose above that $\varepsilon=+1$ (and the same in A!).
\[
\begin{aligned}
\tau :\quad & \longrightarrow\ \underbrace{\longleftarrow\ \longleftarrow\ \cdots\ \longleftarrow}_{l}\ \longrightarrow\\
\tau' :\quad & \longrightarrow\ \longrightarrow\quad\text{or}\quad\longrightarrow\ .
\end{aligned}
\]

But all the morphisms $\tau\longrightarrow\tau'$ are trivial on the critical interval $I_2$ (retrograde), which implies that
\[
\uHom(\tau,\tau')\xleftarrow{\ \sim\ }\uHom(\Delta^2_\varepsilon,\tau'),
\]
via the canonical contraction functor
\[
I_\tau\longrightarrow I_{\Delta^2_\varepsilon}=\boldsymbol{\Delta}^2.
\]

Now in A we saw that $\uHom(\Delta^2_\varepsilon,\tau')$ is connected (for $\tau'$ of the form $\Delta_\varepsilon$ or $\Delta^2_\varepsilon$).

We've won!

\par
\printedpage{60}
% head: A. Grothendieck : Dérivateurs, Ch. VII, Édition des Universités de Montpellier II et Paris VII
% foot: Version 22 décembre 2001

\msp{page 94}

\section*{9\quad Homotopy theorem for categories of paths}

Let $X$ be a category, $\Sigma\subseteq\mathrm{Fl}(X)$ \texttt{[an]} admissible set of arrows in $X$, $a,b\in\mathrm{Ob}\,X$. Let us consider for variable $\tau\in\mathrm{Ty}$, the categories
\[
\uCh_{\Sigma,\tau}(X),\qquad\quad\uCh_{\Sigma,\tau}(X;a,b),
\]
which form the fibers of the Cat-fibered categories
\[
\begin{tikzcd}[column sep=small]
\uCh_\Sigma(X) \arrow[dr] & & \uCh_\Sigma(X;a,b) \arrow[ll, "\text{incl.}"'] \arrow[dl]\\
 & \uTy &
\end{tikzcd}
\]
over the category $\uTy$. Thus, for $\tau,\tau'\in\mathrm{Ty}$, and for any
\[
\varphi:\tau\longrightarrow\tau',
\]
we have a commutative square of homomorphisms
\[
\begin{tikzcd}
\uCh_{\Sigma,\tau'}(X) \arrow[d, "\varphi^*_X"'] & \uCh_{\Sigma,\tau'}(X;a,b) \arrow[l, "\text{incl.}"'] \arrow[d, "\varphi^*_{a,b}"]\\
\uCh_{\Sigma,\tau}(X) & \uCh_{\Sigma,\tau}(X;a,b), \arrow[l, "\text{incl.}"']
\end{tikzcd}
\]
with transitivity for variable $\varphi$ (composition $\psi\circ\varphi$ of arrows in $\uTy$).

In fact, for $\varphi$ variable in $\uHom(\tau,\tau')$, the inverse image functors $\varphi^*_X$, $\varphi^*_{a,b}$ depend \textit{functorially} on $\varphi$; i.e. we have \textit{functors}
\[
\begin{tikzcd}[row sep=small]
 & \uHom(\uCh_{\Sigma,\tau'}(X),\uCh_{\Sigma,\tau}(X))\\
\uHom(\tau,\tau') \arrow[ur] \arrow[dr] & \\
 & \uHom(\uCh_{\Sigma,\tau'}(X;a,b),\uCh_{\Sigma,\tau}(X;a,b))
\end{tikzcd}
\tag*{$(*)$}
\]

\par
\printedpage{61}
% head: A. Grothendieck : Dérivateurs, Ch. VII, Édition des Universités de Montpellier II et Paris VII
% foot: Version 22 décembre 2001

\msp{page 95}

This comes from the inclusion
\[
\uHom(\tau,\tau')\hookrightarrow\uHom^{\varepsilon}(I_\tau,I_{\tau'})
\]
and from the description of the categories $\uCh_\tau$, $\uCh_{\tau'}$ as categories of functors, such as
\[
\uCh_{\Sigma,\tau}(X)\simeq\uHom_{\Sigma_\tau,\Sigma}(I_\tau,X),
\]
where $\Sigma_\tau$ denotes the set of arrows in $I_\tau$ which are either invertible or \textit{retrograde.} (This set of arrows is admissible, since stable under composition.) We denote by $\uHom_{\Sigma_\tau,\Sigma}$ the full subcategory of $\uHom$ formed by the functors which map $\Sigma_\tau$ into $\Sigma$. Now in general, if
\[
(I,\Sigma_I),\ \ (I',\Sigma_{I'}),\ \ (X,\Sigma)
\]
are three categories, equipped with admissible sets of arrows, considering the functors $\varphi:I\longrightarrow I'$ with $\varphi(\Sigma_I)\subseteq\Sigma_{I'}$, then the familiar functor
\[
\uHom(I,I')\longrightarrow\uHom(\uHom(I',X),\uHom(I,X))
\]
(transpose of the composition functor
\[
\begin{matrix}
\uHom(I,I') & \times & \uHom(I',X) & \longrightarrow & \uHom(I,X) & \\
(\varphi & , & F') & \longmapsto & F'\circ\varphi & ),
\end{matrix}
\]
generalizes to a functor
\[
\uHom_{\Sigma_I,\Sigma_{I'}}(I,I')\longrightarrow\uHom(\uHom_{\Sigma_{I'},\Sigma}(I',X),\uHom_{\Sigma_I,\Sigma}(I,X)),
\tag*{$(*)$}
\]
transpose of the composition functor
\[
\uHom_{\Sigma_I,\Sigma_{I'}}(I,I')\times\uHom_{\Sigma_{I'},\Sigma}(I',X)\longrightarrow\uHom_{\Sigma_I,\Sigma}(I,X)
\]
induced by $(*)$.

\msp{page 96}

Since every
\[
\varphi:\tau\longrightarrow\tau'
\]
induces
\[
\varphi:I_\tau\longrightarrow I_{\tau'},
\]
transforming retrograde arrows into identity or retrograde arrows, i.e.
\[
\varphi(\Sigma_\tau)\subseteq\Sigma_{\tau'},
\]
\printedpage{62}
% head: A. Grothendieck : Dérivateurs, Ch. VII, Édition des Universités de Montpellier II et Paris VII
% foot: Version 22 décembre 2001
we find the first of the two functors $(*)$ announced on page 94, which is associated by transposition with the functor $\gamma$
\[
\begin{tikzcd}[column sep=large]
\uHom(\tau,\tau')\times\ \underbrace{\uCh_{\Sigma,\tau'}(X)}_{\simeq\,\uHom_{\Sigma_{\tau'},\Sigma}(I_{\tau'},X)} \arrow[r, "\gamma"] \arrow[d, "\text{can.}"'] & \uCh_{\Sigma,\tau}(X) \arrow[d, equal]\\
\uHom_{\Sigma_\tau,\Sigma_{\tau'}}(I_\tau,I_{\tau'})\times\uHom_{\Sigma_{\tau'},\Sigma}(I_{\tau'},X) \arrow[r, dotted, "\gamma_0"] & \uHom_{\Sigma_\tau,\Sigma}(I_\tau,X)
\end{tikzcd}
\tag*{$(*)$}
\]
deduced from the composition functor $\gamma_0$ in this diagram: $\gamma=\gamma_0\circ\mathrm{can}.$ On the other hand, this composition functor $\gamma$ visibly induces a functor on the subcategories of paths with fixed source and target
\[
\uHom(\tau,\tau')\times\uCh_{\Sigma,\tau'}(X;a,b)\longrightarrow\uCh_{\Sigma,\tau}(X;a,b),
\]
whence the second of the two arrows announced in $(*)$ p. 94.

This therefore shows that if
\[
\varphi_1,\varphi_2:\tau\rightrightarrows\tau'
\]
are two morphisms from $\tau$ to $\tau'$ homotopic in $\uHom(\tau,\tau')$, then the induced functors between categories of $\Sigma$-paths or between categories of $\Sigma$-paths with specified endpoints: % [sic] précisés -> précisées, agreement with extrémités

\msp{page 97}
\[
\begin{aligned}
\varphi^*_{1\,X},\varphi^*_{2\,X} : && \uCh_{\Sigma,\tau'}(X) &\quad\rightrightarrows\quad\uCh_{\Sigma,\tau}(X),\\
(\varphi_1)^*_{a,b},(\varphi_2)^*_{a,b} : && \uCh_{\Sigma,\tau'}(X;a,b) &\quad\rightrightarrows\quad\uCh_{\Sigma,\tau}(X;a,b)
\end{aligned}
\]
are also homotopic.

Now the homotopy lemma assures us that the hypothesis $\varphi_1\sim\varphi_2$ is \textit{always} satisfied. Thus:

\textbf{Homotopy theorem on the universal operations on paths.} \textit{Let $X$, $\Sigma$, $a$, $b$ be as above, and $\tau,\tau'\in\mathrm{Ty}$ such that $\mathrm{Hom}(\tau,\tau')\neq\emptyset$. Then the functors $\varphi^*_X$, $\varphi^*_{X;a,b}$ induced by a $\varphi\in\mathrm{Hom}(\tau,\tau')$}
\[
\left\{
\begin{aligned}
\varphi^*_X : && \uCh_{\Sigma,\tau'}(X) &\quad\longrightarrow\quad\uCh_{\Sigma,\tau}(X)\\
\varphi^*_{X;a,b} : && \uCh_{\Sigma,\tau'}(X;a,b) &\quad\longrightarrow\quad\uCh_{\Sigma,\tau}(X;a,b)
\end{aligned}
\right.
\]
\textit{do not depend on $\varphi$,} up to homotopy.

We can make this statement more precise, in such a way that the statement for the $\varphi^*_{X;a,b}$ appears, as it should, as a corollary to the one on $\varphi^*_X$. For this, we are going to look at the categories of paths in $X$ as categories \textit{over $X\times X$,} by the ``source and target'' functors:
\DERmark{ENV0490}

\[
\begin{tikzcd}[column sep=small]
\uCh_{\Sigma,\tau}(X) \arrow[dr, "(\sigma_\tau{,}\beta_\tau)"'] & & \uCh_{\Sigma,\tau'}(X) \arrow[dl, "(\sigma_{\tau'}{,}\beta_{\tau'})"]\\
 & X\times X &
\end{tikzcd}
\]

\par
\printedpage{63}
% head: A. Grothendieck : Dérivateurs, Ch. VII, Édition des Universités de Montpellier II et Paris VII
% foot: Version 22 décembre 2001

\msp{page 98}

Then the functor $(*)$ (p. 94) factors through the subcategory $\uHom_{\mathrm{Cat}/X\times X}(\cdot,\cdot)$ of $\uHom_{\mathrm{Cat}}(\cdot,\cdot)$ formed by the functors which commute with the projection functors to $X\times X$, and the homomorphisms between such functors which are effected by ``vertical'' arrows of the target category (i.e. by arrows over identity arrows of the base category $X\times X$):
\[
\uHom(\tau,\tau')\longrightarrow\uHom_{\mathrm{Cat}/X\times X}(\uCh_{\Sigma,\tau'}(X),\uCh_{\Sigma,\tau}(X)).
\]

Thus in the statement of the homotopy theorem, we can specify that the homotopy takes place in this subcategory, i.e. is effected by a sequence of simple homotopies which are morphisms of functors between categories over $X\times X$. This then implies the homotopy statement on the $\varphi^*_{X;a,b}$.


\endgroup
% END UNIT Der-ChVII-456/Der-ChVII-456_p061-065.tex
% BEGIN UNIT Der-ChVII-456/Der-ChVII-456_p066-070.tex
\begingroup\makeatletter%
\NoAutoSpacing%
\newcommand{\printedpage}[1]{\marginpar{\footnotesize p.~#1}}%
\newcommand{\msp}[1]{\marginpar{\footnotesize [#1]}}%
\newcommand{\Cat}{\mathrm{Cat}}%
\newcommand{\Cathot}{\mathrm{Cat\,hot}}%
\newcommand{\Hot}{\mathrm{Hot}}%
\newcommand{\Hom}{\mathrm{Hom}}%
\newcommand{\Ob}{\mathrm{Ob}\,}%
\newcommand{\uTy}{\underline{\mathrm{Ty}}}%
\newcommand{\uCh}{\underline{\mathrm{Ch}}}%
\newcommand{\uHom}{\underline{\mathrm{Hom}}}%
\newcommand{\sHom}{\underline{\mathcal{H}\mathit{om}}}%
\newcommand{\To}{T_{\mathrm{o}}}%
\newcommand{\circled}[1]{\tikz[baseline=(c.base)]\node[draw,circle,inner sep=1pt](c){#1};}%
\makeatother\let\printedpage\derpp\let\msp\dermsp%
\let\mspin\DERmspin%


\leavevmode\printedpage{64}
% head: A. Grothendieck : Dérivateurs, Ch. VII, Édition des Universités de Montpellier II et Paris VII
\msp{page 99}

\section*{10\quad Comparison question for \textmd{Cat}-fibered categories satisfying the ``homotopy theorem'' (filtering base)}

We are now going to generalize to the case of a (Cat)-fibered category over another
\[
\begin{tikzcd}
C \arrow[d, "\pi"'] \\
T,
\end{tikzcd}
\]
thinking mainly of the case $T = \uTy$, $C = \uCh_{\Sigma}(X;a,b)$. The essential hypothesis is the following:

\noindent\begin{tabular}{@{}p{0.3\textwidth}p{0.64\textwidth}@{}}
(Homotopy hypothesis) & If $t,t' \in \Ob T$ are such that $\Hom(t,t') \neq \emptyset$, then the functors $\varphi^{*} : C_{t'} \longrightarrow C_{t}$ associated to the $\varphi \in \Hom(t,t')$ are homotopic.
\end{tabular}

They therefore define one and the same arrow in Cat\,hot. I denote by $T_{\mathrm{ord}}$ or $\To$ the set $\Ob(T)$, with as order relation the one deduced from the category structure:
\[
t \leq t' \qquad\Longleftrightarrow\qquad \Hom_{T}(t,t') \neq \emptyset,
\]
so that we have a functor
\[
T \longrightarrow \To
\]
which is \textit{universal} for functors from $T$ into ``ordered'' categories (i.e. where always $\#\Hom(a,b)\le1$ \texttt{[no, cf. }$\bullet\rightrightarrows\bullet$\texttt{; rather, where always }$\#\Hom(a,b) \leq 1$ \texttt{(admitting nonidentity isomorphisms)]}) or for functors $F$ on $T$ such that $F$ is constant on every $\Hom(t,t')$. The homotopy hypothesis can be expressed by saying that the composite functor
\DERmark{ENV0491}

\[
\begin{tikzcd}
T^{o} \arrow[r] & \Cat \arrow[r] & \Cathot \\
t \arrow[r, mapsto] & C_{t} &
\end{tikzcd}
\]
\msp{page 100}
factors (obviously in a unique way) as a functor
\[
H_{C} : \To^{o} \longrightarrow \Cathot,
\]
i.e. that we have a commutative diagram
\printedpage{65}
% head: A. Grothendieck : Dérivateurs, Ch. VII, Édition des Universités de Montpellier II et Paris VII
\[
\begin{tikzcd}
T^{o} \arrow[r, "t\,\longmapsto\, C_{t}"] \arrow[d] & \Cat \arrow[d] \\
\To^{o} \arrow[r, "H_{C}"] & \Cathot.
\end{tikzcd}
\]

We are going to assume in addition:

\noindent\begin{tabular}{@{}p{0.36\textwidth}p{0.58\textwidth}@{}}
(Filtering hypothesis) & The ordered set $\To$ associated to $T$ is downward \textit{filtering} (left filtering), i.e. if $t,t' \in \Ob T$, there exist $t'' \in \Ob T$ and arrows $t'' \longrightarrow t$, $t'' \longrightarrow t'$.
\end{tabular}

This hypothesis is indeed satisfied in $\uTy$, since for $\tau,\tau' \in \Ob \uTy = \mathrm{Ty}$, it suffices to take $\tau\tau'$, which maps to $\tau$ by $\mathrm{id}_{\tau} * p_{\tau'}$, to $\tau'$ by $p_{\tau} * \mathrm{id}_{\tau'}$, where $p_{\tau} : \tau \longrightarrow 1$ is the morphism from $\tau$ to the final object.

Thus, $H_{C}$ defines an ind-object of Cat\,hot, which we are going to denote by $\underline{\underline{C}}$. If we assume in addition:

\noindent\begin{tabular}{@{}p{0.36\textwidth}p{0.58\textwidth}@{}}
(Countability hypothesis, always satisfied if $\mathrm{card}\,\To \leq \aleph_{0}$) & $\To$ has a \textit{countable} subset cofinal on the left.
\end{tabular}

then this ind-object of Cat\,hot can be defined, in many ways, in terms
\msp{page 101}
of a direct system in Cat (not just in Cat\,hot)
\[
(C_{n})_{n\in\mathbf{N}},
\]
from which we recover the ind-object $\underline{\underline{C}}$ in Cat\,hot by passing to the corresponding direct system in Cat\,hot, then to the associated ind-object.

To construct the direct system $C_{n}$, it suffices to take an infinite chain
\[
t_{0} \longleftarrow t_{1} \longleftarrow t_{2} \longleftarrow \;\cdots\; \longleftarrow t_{n} \longleftarrow t_{n+1} \longleftarrow \;\cdots
\]
of morphisms in $T$, such that the set of the $t_{i}$ is cofinal on the left in $\To$. We are going to assume in addition that $C$ is \textit{strictly fibered} over $T$ (i.e. coming from a genuine functor $T^{o} \longrightarrow \Cat$, i.e. we can choose the inverse image functors
\[
\varphi^{*} : C_{t'} \longrightarrow C_{t}
\]
in a strictly transitive way, not just transitive up to isomorphism) - a case to which we can always reduce, moreover, by the trick of \textsc{Giraud}. This hypothesis is moreover verified in the case $\uCh_{\Sigma}(X;a,b)$ over $\uTy$. Thus, the
\[
(C_{t_{i}})_{i\in\mathbf{N}}
\]
form a direct system of categories indexed by $\mathbf{N}$, the transition morphism
\[
C_{n} \longrightarrow C_{n'} \qquad \text{i.e. } C_{t_{n}} \longrightarrow C_{t_{n'}}
\]
\printedpage{66}
% head: A. Grothendieck : Dérivateurs, Ch. VII, Édition des Universités de Montpellier II et Paris VII
\msp{page 102}
being the inverse image functor $\varphi^{*}_{n,n'}$, where
\[
\varphi_{n,n'} : t_{n'} \longrightarrow t_{n}
\]
is the composite of the $t_{n'} \longrightarrow t_{n'-1} \longrightarrow \;\cdots\; \longrightarrow t_{n}$. (If we did not have strict transitivity, we would define $\varphi^{*}_{n,n'}$ as the composite of the inverse image functors along successive transition arrows $t_{n'} \longrightarrow t_{n'-1} \longrightarrow \;\cdots\; \longrightarrow t_{n+1} \longrightarrow t_{n}$, this would be $\varphi^{*}_{n,n'}$ only up to canonical isomorphism.)

Let
\[
C_{\infty} = \varinjlim_{n} C_{n} = \varinjlim_{n} C_{t_{n}}
\]
be the category $\varinjlim$ . (It will be especially nice if we have been able to choose the system $(t_{n})_{n\in\mathbf{N}}$ in such a way that the transition functors
\[
C_{t_{n-1}} \longrightarrow C_{t_{n}}
\]
\DERmark{ENV0492}\DERmark{ENV0493}

are fully faithful, and perhaps even \texttt{[some]} monomorphisms, in such a way that the $C_{n}$ identify with full subcategories of their direct limit $C_{\infty}$.) I would like to understand the precise relations between the following three kinds of objects:

\begin{itemize}[leftmargin=2.5em]
\item[\circled{a}] $C$ itself (as an object in Cat, not of $\Cat/T$), and especially the objects of Cat\,hot and of Hot which correspond to it.

\item[\circled{b}] \msp{page 103}$C_{\infty}$, and the objects of Cat\,hot and of Hot which correspond to it. To what extent are these latter % [sic] deniers -> derniers, missing letter
independent % [sic] indépendant -> indépendants, agreement with the plural subject
of the choice of the inverse system cofinal on the left $(t_{n})_{n\in\mathbf{N}}$ in $T$? What is the role of $C_{\infty}$? Is it canonically isomorphic to $C$ in Cat\,hot, or in Hot? (In which case the question of its ``role'' would seem to me resolved!)

\item[\circled{c}] The ind-object % [sic] Le ind-objet -> L'ind-objet, missing elision
$H_{C}$ in Cat\,hot, and the ind-object of Hot which it defines.

I note that the ind-object $H_{C}$, and also the categories $C_{\infty}$ constructed as above, do not change when we replace $T$ by a subcategory $T'$ cofinal on the left, or also by a category of the form $T/t_{0}$ ($t_{0} \in \Ob T$), which gives rise to the associated ordered set % [sic] ordonnée -> ordonné, agreement with "ensemble"
\[
(T/t_{0})_{\mathrm{o}} \longrightarrow \To/t_{0} = \{t \in \To \mid t \leq t_{0}\}
\]
\DERmark{ENV0494}\DERmark{ENV0495}\DERmark{ENV0496}

\texttt{[rather} $(T/t_{0})_{\mathrm{o}} = \To/t_{0} = \{t \in \To \mid t \leq t_{0}\}$\texttt{]}. If the relation between the objects in a), and those in b) and in c), is very close, we would conclude from this % [sic] concluerait -> conclurait, spelling of the conditional
that the ``type'' of $C$ in Cat\,hot, or in Hot as the case may be, does not change on replacing $C$ by
\printedpage{67}
% head: A. Grothendieck : Dérivateurs, Ch. VII, Édition des Universités de Montpellier II et Paris VII
\msp{page 104}
$C \times_{T} T'$ or by $C/t_{0} \simeq C \times_{T} T/t_{0}$. Better still, it would not change on replacing $T$ by $(\mathbf{N})^{o}$, the (ordered) category opposite to $\mathbf{N}$ (thus having a unique final object), and the fibered category $C$ over $T$ by the category $C_{\underline{t}}$ induced by $C$ over $(\mathbf{N})^{o}$, via the functor
\[
\underline{t} : (\mathbf{N})^{o} \longrightarrow T
\]
defined by the inverse system $\underline{t} = (t_{n})_{n\in\mathbf{N}}$ in $T$. This would thus mean that to understand the homotopy type (in the strict sense, i.e. in Cat\,hot, or broad, i.e. in Hot) of $C$, we can ``clean'' it in a draconian way, the base $T$ being viewed as a sort of somewhat muddy gangue, or at least plethoric.

\item[\circled{d}] In any case, if $T'$ is another left filtering category, and
\[
\alpha : T' \longrightarrow T
\]
a functor cofinal on the left, so that $C$ over $T$, and $C' = C \times_{T} T'$ over $T'$, define essentially the ``same'' object of Ind\,Cat\,hot, and the category $C'_{\infty}$ for $C'$ over $T'$ is among the categories $C_{\infty}$ for $C$ over $T$ - we may ask
\msp{page 105}
whether the canonical functor
\[
C' \;\longrightarrow\; C
\]
is also always an isomorphism in Cat\,hot, or if not, at least in Hot.
\end{itemize}

A particular interest attaches to $\pi_{0}(C)$, and moreover to its variance for $C$ over $T$ varying in the following way. If $S$ is in Cat, we look at the categories
\[
\uHom_{\Cat}(S,C_{t}) \qquad (t \in T),
\]
we find a functor
\[
t \longmapsto \uHom_{\Cat}(S,C_{t}), \qquad T^{o} \longrightarrow (\Cat),
\]
whence a fibered category over $T$, which I am going to denote by
\[
\sHom_{T}(S_{T},C)
\]
(Unless I am mistaken, it is interpreted as the fibered category over $T$ corresponding to the stack on $\mathrm{Top}(T) = T^{\wedge}$ of the morphism of $T$-stacks from $S_{T}$ to $C$, where $S_{T} = S \times T$ (considered as a fibered category over $T$ with essentially constant fiber $S$) to $C$. Here I identify fibered category over $T$, and stacks on the topos $T^{\wedge}$.) It is immediate that the homotopy property (p. 99) is stable under passage from $C$ to $\sHom_{T}(S_{T},C)$. We are interested in the variation of
\[
\pi_{0}(\sHom_{T}(S_{T},C))
\]
\printedpage{68}
% head: A. Grothendieck : Dérivateurs, Ch. VII, Édition des Universités de Montpellier II et Paris VII
\msp{page 106}
as a contrafunctor in $S$,
\[
H : S \longmapsto \pi_{0}\sHom_{T}(S_{T},C) \qquad (\Cat)^{o} \longrightarrow (\mathrm{Ens}),
\]
and in a possible expression of this contrafunctor $H$ in terms of the ind-object $H_{C}$ of Cat\,hot, by the hypothetical formula
\[
H(S) \simeq \Hom_{\mathrm{Ind}(\Cathot)}(\gamma(S),H_{C}),
\]
where
\[
\gamma : \Cat \longrightarrow \Cathot
\]
is the canonical functor. If we had such a canonical isomorphism functorial in $S$, this would give a description (up to unique isomorphism) of the ind-object $H_{C}$ in terms of the preceding functor $H$ on Cat.

A doubt comes to me here, in view of the example $T = \uTy$, $C = \uCh_{\Sigma}(X;a,b)$. The three types of objects a) b) c) have the same $\pi_{0}$, and it is certain that ``it's the right one''. On the other hand, for a variable parameter category $S$, it is also certain that the right functor in $S$ is
\[
S \longmapsto \pi_{0}(\uHom_{T}(S_{T},C)) \simeq \Hom_{\mathrm{Ind}\,\Cathot}(S,H_{C}),
\]
also isomorphic, \textit{if $S$ is finite,} to
\[
\Hom_{\Cathot}(S,C_{\infty}),
\]
which shows that, at least for arguments $S \in (\Cat)$ which are \textit{finite} categories, the objects
\msp{page 107}
$C_{\infty}$ and $H_{C}$ are ``good'', they allow us % [sic] il nous permettent -> ils nous permettent, agreement with the plural subject
to describe an important functor which dominates the situation under consideration. Under a very strong hypothesis, but one satisfied in the important cases, we can in these two formulas replace Cat\,hot by Hot.

By contrast, except for the case where $S$ is reduced to a point and where $\Hom_{\Cathot}$ or $\Hom_{\Hot}$ from $S$ to $C$ is none other than $\pi_{0}(C)$ ($\simeq \pi_{0}(C_{\infty}) \simeq \pi_{0}(H_{C})$), thus indeed ``good'', I do not know whether
\[
\underbrace{\Hom_{\Cathot}(S,C)}_{=\,\pi_{0}\uHom(S,C)} \qquad \text{or} \qquad \Hom_{\Hot}(S,C)
\]
is related in a reasonable way to the preceding sets (determined by $H_{C}$, resp. by $C_{\infty}$), e.g. whether they are isomorphic to it in a natural way, and whether they have a reasonable geometrico-topological meaning. I am beginning to suspect \textit{not} - hence that we would have here a case where a relation of a very general type concerning a $\pi_{0}$, here thus
\[
\pi_{0}(C) \simeq \pi_{0}(C_{\infty}) \simeq \pi_{0}(H_{C}),
\]
does not extend to an isomorphism on the higher $\pi_{i}$, and to a property of Hot-equivalence.


\endgroup
% END UNIT Der-ChVII-456/Der-ChVII-456_p066-070.tex
% BEGIN UNIT Der-ChVII-456/Der-ChVII-456_p071-075.tex
\begingroup\makeatletter%
\newcommand{\printedpage}[1]{\marginpar{\footnotesize p.~#1}}%
\newcommand{\msp}[1]{\marginpar{\footnotesize [#1]}}%
\newcommand{\Hotz}{\mathrm{Hot}^{\wedge}_{0}}%
\newcommand{\Ch}{\underline{\mathrm{Ch}}}%
\newcommand{\Ty}{\underline{\mathrm{Ty}}}%
\newcommand{\CNo}{C_{\mathbf{N}^{o}}}%
\newcommand{\lab}[2][l]{\text{\footnotesize\begin{tabular}{@{}#1@{}}#2\end{tabular}}}%
\makeatother\let\printedpage\derpp\let\msp\dermsp%
\let\mspin\DERmspin%


\printedpage{69}
% head: A. Grothendieck : Dérivateurs, Ch. VII, Édition des Universités de Montpellier II et Paris VII

\msp{page 108}
I find no direct arrow, for that matter, between $C$ and either of the objects $C_{\infty}$, $H_{C}$; all I see in this direction is
\[
\begin{tikzcd}[column sep=large, row sep=huge]
\CNo \arrow[r, "{\substack{\text{localization}\\ \text{functor}}}"] \arrow[d] & C_{\infty} \\
C, &
\end{tikzcd}
\]
where the vertical arrow is the projection arrow of $\CNo = C \times_{T} \mathbf{N}^{o}$ (with $\mathbf{N}^{o} \xrightarrow{\ t\ } T$). This establishes a close topological relation between $C$ (as an object of Cat\,hot or simply of Hot) and $C_{\infty}$, only if the two preceding arrows are at least Hot-equivalences. And the typical (problematic) case where the question of the relation between $C$ and $C_{\infty}$ arises, is therefore the one where $T = \mathbf{N}^{o}$, i.e. that of an arbitrary direct system of (small) categories, indexed by $\mathbf{N}$, and interpreted as defined by a fibered category $C = \CNo$ over $\mathbf{N}^{o}$. In this particular case, one has at least a canonical arrow
\[
C \longrightarrow C_{\infty}
\]
about which one can ask oneself the question whether it is
\msp{page 109}
a Hot-equivalence. If not, this will prove that in the general case, the Hot-type of $C$ (and a fortiori its isomorphism class in Cat\,hot) has no interesting geometric meaning. That is to say that $C$ is interesting only \textit{as a fibered category over $T$}, making it possible in particular to construct the ind-object $H_{T}$ of Cat\,hot or of Hot, but not as a ``bare category''.

But I have just convinced myself that $\CNo \longrightarrow C_{\infty}$ is indeed a Hot-equivalence. So the arrow $\CNo \longrightarrow C$ defines an arrow \textit{in} Hot
\[
C_{\infty} \longrightarrow C
\]
(by composing with the inverse in Hot of the preceding one). Is it an isomorphism? So for this question, one is reduced to the question for $\CNo \longrightarrow C$. But if one takes $C = T$, this already poses the same question for $\mathbf{N}^{o} \longrightarrow T$, hence the question whether $T$ is aspheric, under the sole hypothesis that $T_{\mathrm{o}}$ is left filtering and countable at infinity (on the left). But of course this is false (e.g. if $T$ has only one object, it is not aspheric for all that!).

This therefore convinces me that in the absence
% foot: Version 22 décembre 2001
\printedpage{70}
% head: A. Grothendieck : Dérivateurs, Ch. VII, Édition des Universités de Montpellier II et Paris VII
\msp{page 110}
of an additional draconian hypothesis on $T$, the ``bare'' category $C$ is of no particular interest, apart from its $\pi_{0}$ (which is also given by $C_{\infty}$, and by $H_{C}$). It is not certain, to be sure, that \underline{Ty} does not satisfy these ``draconian conditions''. But since ``bare'' $C$ does not fit into a congenial formalism (e.g. one has no geometric interpretation of $\mathrm{Hom}_{\mathrm{Cat\,hot}}(S,C)$, or of $\mathrm{Hom}_{\mathrm{Hot}}(S,C)$), I think that in what follows I can leave it out of account. Exit a)!

\begin{flushright}
{\footnotesize 5.11.90}
\end{flushright}

All the same, I spent the whole day yesterday trying to understand the relation of $C$ and of its type in Hot, to that of $C_{\infty}$. Indeed, since one finally has an arrow in Hot
\[
C_{\infty} \longrightarrow C,
\]
deduced from the diagram in (Hot)
\[
\begin{tikzcd}[column sep=huge, row sep=huge]
\CNo \arrow[r, "\sim"] \arrow[d] & C_{\infty} \\
C, &
\end{tikzcd}
\]
to prove that $C_{\infty} \longrightarrow C$ is an isomorphism in Hot
\msp{page 111}
seems the simplest, and perhaps the only, way of proving that the image of $C_{\infty}$ in Hot does not depend, up to unique isomorphism, on the choices made to construct $C_{\infty}$. Not to mention that in the motivating case, $C = \Ch_{\Sigma}(X)$ or $\Ch_{\Sigma}(X;a,b)$, $T = \Ty$, it is the category of paths $C$ that gives rise to an impeccable composition law for paths, functorial, associative and all, whereas there is nothing of the sort for $C_{\infty}$. One does not find a composition law
\[
C_{\infty} \times C_{\infty} \longrightarrow C_{\infty}
\]
in Cat, nor even in Cat\,hot or in Hot - or rather, if one does nevertheless manage to construct an arrow $C_{\infty} \times C_{\infty} \longrightarrow C_{\infty}$ in Cat, it is still laboriously by dint of $\pm$ arbitrary choices, notch by notch, and without knowing whether what one obtains has a uniqueness property, even if only in Hot; one finds a tangible geometric interpretation of this composition only in $\Hotz$ (where one has a genuine groupoid law), or in $H_{T}$ as an object of Ind\,Cat\,hot
% foot: Version 22 décembre 2001
\printedpage{71}
% head: A. Grothendieck : Dérivateurs, Ch. VII, Édition des Universités de Montpellier II et Paris VII
\msp{page 112}
or of Ind\,Hot. (So the direct limit in $\Hotz$ is none other than the element of $\Hotz$ defined by $C_{\infty}$.)

My efforts to prove that
\[
C_{\infty} \longrightarrow C
\]
was an isomorphism in Hot, or even just a 1-isomorphism (i.e. an isomorphism not only for the $\pi_{0}$, which is always true, but also for the $\pi_{1}$ of all the components), have failed. I found that it was enough for the functor % [sic] ont échoués -> ont échoué, accord fautif du participe
\[
T \longrightarrow T_{\mathrm{o}}
\]
to be either smooth and coaspheric, or proper and aspheric - so everything is OK if $T = T_{\mathrm{o}}$, i.e. if $T$ is ordered. But in the case $T = \Ty$, the functor $T \longrightarrow T_{\mathrm{o}}$ is neither proper, nor smooth (not even 0-proper or 0-smooth).

Finally I have come to the conviction that in the case we are concerned with ($T = \Ty$, $C = \Ch_{\Sigma}(X;a,b)$), the arrow $C_{\infty} \longrightarrow C$ is in general \textit{not} an isomorphism in Hot. Indeed, in the case $a = b$, if this \texttt{[were true]},
\msp{page 113}
one would have a description of the Hot-group $\Omega(X,a) = \Ch(X,a)$ (\footnote[23]{Here we take $\Sigma = \mathrm{Fl}(X)$.}) in terms of a category $C$ equipped with an \textit{associative} (and unital) composition law \textit{in} Cat, giving back the group law of $\Omega$ on passing from (Cat) to Hot. Taking the nerve of $C$, one would find a semi-simplicial monoid which, in Hot, would give the group law of $\Omega(X,a)$. Now I remember (\footnote[24]{This is an error - it must be for \textit{other} group objects of Hot than loop groups. I saw in Quillen that $\Omega(X,x)$ must be describable by a semi-simplicial group (when $X$ is a semi-simplicial set, $x$ a vertex). So the problem whether $C_{\infty} \simeq C$ in Hot remains completely open.}) that it was explained to me that in general, the group law of $\Omega(X,a)$ could not be obtained by an \textit{associative} law on a semi-simplicial set. Thus, the failure of my efforts seems to be due, not to my clumsiness, but to the nature of things. % [sic] C'est un erreur (note 24) -> C'est une erreur, genre

I am going to summarize the positive facts I have arrived at in the following diagram, which takes place in the general situation $C \longrightarrow T$ considered at the beginning of the present paragraph:

\msp{page 114}
% foot: Version 22 décembre 2001
\printedpage{72}
% head: A. Grothendieck : Dérivateurs, Ch. VII, Édition des Universités de Montpellier II et Paris VII
\[
\begin{tikzcd}[column sep=small, row sep=huge]
 & & C_{t'} \arrow[d, "{u^{*}\quad\lab{(for $u : t \longrightarrow t'$\\ in $T$)}}"] \arrow[ldd, dotted] \arrow[rdd, dotted] & \\
 & & C_{t} \arrow[ld, swap, "{\lab{arrow defined in\\ Ind\,Cat\,hot\\ (\texttt{[or in]}\\ IndHot,\\ Ind\,($\Hotz$))}}"] \arrow[rd, "{\lab[r]{arrow in Cat\\ (hence in\\ Cat\,hot,\\ Hot,\\ $\Hotz$)}}"] & \\
\varinjlim C_{t} = & |[draw, circle]| H_{T} \arrow[rd, swap, "{\lab{arrow\\ defined in\\ Ind\,Cat\,hot\\ or in Ind\,Hot\\ or in Ind\,$\Hotz$}}"] & \lab{commutative\\ diagram\\ in Ind\,Hot\\ or Ind\,($\Hotz$)} & C \\
\lab{in Ind\,Cat\,hot\\ (or in Ind\,Hot\\ or in Ind\,($\Hotz$))} & & |[draw, circle]| C_{\infty} \arrow[ru, swap, "{\lab[r]{arrow\\ defined\\ in Hot,\\ hence in $\Hotz$,\\ in general not an\\ isomorphism, even\\ in $\Hotz$}}"] & \\
 & & = \varinjlim H_{T} \text{ in } \Hotz &
\end{tikzcd}
\]

\texttt{[The annotations in this diagram concern the solid arrows.]}

The important objects are $H_{T}$ and $C_{\infty}$, whereas $C$ (over $T$) and the $C_{t}$ take on an intermediate technical role, to construct $H_{T}$ and $C_{\infty}$. One can consider $C_{\infty}$ as a category (object of (Cat)), but whose construction depends on many arbitrary choices, unlike $H_{T}$ in Ind\,Cat\,hot. It is the object of $\Hotz$ deduced from $C_{\infty}$ that has a canonical meaning, and can be described as the direct limit of $H_{T}$ as an object of Ind\,$\Hotz$, i.e. as representing the same functor on $\mathrm{Hot}_{0}$ as $H_{T}$.

\msp{page 115}
In the geometric case $T = \Ty$, and $C$ of the form $\Ch_{\Sigma}(X;a,b)$, giving rise to $C_{\infty}(a,b)$, one \textit{does not have} a natural definition of a bifunctor of composition of paths
\[
C_{\infty}(a,b) \times C_{\infty}(b,c) \ \longrightarrow\ C_{\infty}(a,c).
\]
One can define, to be sure, pairings
\[
\Ch_{\Sigma,\tau}(a,b) \times \Ch_{\Sigma,\tau'}(b,c) \ \longrightarrow\ \Ch_{\Sigma,\tau\tau'}(a,c),
\]
or, with the notation relating to the datum of
\[
\begin{aligned}
&\underline{t} : \mathbf{N}^{o} \longrightarrow \Ty,\\
&\Ch_{n}(a,b) \times \Ch_{m}(b,c) \longrightarrow \Ch_{n+m}(a,c),
\end{aligned}
\]
% foot: Version 22 décembre 2001
\printedpage{73}
% head: A. Grothendieck : Dérivateurs, Ch. VII, Édition des Universités de Montpellier II et Paris VII
but these pairings \textit{are not} compatible with the transition morphisms of the three direct systems of categories $(\Ch_{n})_{n\in\mathbf{N}}$, they are so \textit{only up to homotopy.} This is enough, however, to have perfect pairings
\[
H_{T}(a,b) \times H_{T}(b,c) \ \longrightarrow\ H_{T}(a,c),
\]
with associativity, and (if $\Sigma = \mathrm{Fl}(X)$ in its entirety) existence of \textit{true} inverses, i.e. of arrows (in Ind\,Cat\,hot) of ``passage to the opposite path''

\msp{page 116}
\[
H_{T}(a,b) \ \longrightarrow\ H_{T}(b,a), \qquad c \longmapsto c^{o}, \tag*{$(*)$}
\]
giving rise to the commutative diagram expressing that this is indeed an inverse for composition. Another way of saying it is that for every object $Y$ of Cat\,hot (or of Ind\,Cat\,hot), the sets
\[
\mathrm{Hom}_{\mathrm{{Ind}(\mathrm{Cat\,hot})}}(Y, H_{T}(a,b)) \overset{\mathrm{def}}{=} H_{T}^{Y}(a,b)
\]
\DERmark{ENV0497}\DERmark{ENV0498}

form, for variable $a,b \in X$, the arrows of a \textit{groupoid} (\footnote[25]{If $\Sigma = \mathrm{Fl}(X)$. If $\Sigma$ is smaller, instead of a groupoid, one finds a \textit{category}, nothing more.}) $\mathrm{Ch}_{\Sigma}^{Y}(X)$ having as set of vertices $\mathrm{Ob}\,X$, and $H_{T}^{Y}(a,b)$ being made up of the arrows from $a$ to $b$. The map
\[
H_{T}^{Y}(a,b) \longrightarrow H_{T}^{Y}(b,a), \qquad c \longmapsto c^{o},
\]
deduced from the functor $(*)$ above is none other than the map of passage to the inverse. This category is moreover none other than the one deduced from % [sic] celle déduit -> celle déduite, accord
\[
X(Y)\Sigma_{Y}^{-1}, \qquad \text{where } \left\{
\begin{aligned}
&X(Y) = \underline{\mathrm{Hom}}_{\mathrm{Cat}}(Y,X)\\
&\Sigma_{Y} \subseteq \mathrm{Fl}(X(Y)),\\
&\text{defined in terms of } \Sigma \text{ ``argument by argument''},
\end{aligned}
\right.
\]
by means of the functor
\[
X \longrightarrow X(Y), \qquad a \longmapsto a_{Y},
\]
\msp{page 117}
associating to every $a \in X$ the \textit{constant} functor on $Y$ with value $a$, taking
\[
H_{T}^{Y}(a,b) \simeq \mathrm{Hom}_{X(Y)\Sigma_{Y}^{-1}}(a_{Y}, b_{Y}) \ .
\]
\DERmark{ENV0499}


In the same geometric case, I have convinced myself, moreover, that $C_{\infty}$ was, as an object of Hot (not only of $\Hotz$) \textit{unique up to} Hot-\textit{isomorphism,} but (strange thing) perhaps not up to \textit{unique} Hot-isomorphism. In general, returning to the
% foot: Version 22 décembre 2001


\endgroup
% END UNIT Der-ChVII-456/Der-ChVII-456_p071-075.tex
% BEGIN UNIT Der-ChVII-456/Der-ChVII-456_p076-080.tex
\begingroup\makeatletter%
\NoAutoSpacing%
\newcommand{\printedpage}[1]{\marginpar{\footnotesize p.~#1}}%
\newcommand{\msp}[1]{\marginpar{\footnotesize [#1]}}%
\newcommand{\Hot}{\mathrm{Hot}}%
\newcommand{\Hotz}{\mathrm{Hot}_0^{\wedge}}%
\newcommand{\Hc}{\underline{\mathrm{Hom\,cof}}}%
\newcommand{\Nop}{\mathbf{N}^o}%
\makeatother\let\printedpage\derpp\let\msp\dermsp%
\let\mspin\DERmspin%


\printedpage{74}
% head: A. Grothendieck : Dérivateurs, Ch. VII, Édition des Universités de Montpellier II et Paris VII
conditions at the beginning of the paragraph, $C_\infty$ must be written $C_\infty(\underline{t})$, it depends on the choice of the cofinal functor
\[
\underline{t} : \Nop \longrightarrow T.
\]
When $\underline{t}$ varies in the category
\[
\Hc(\Nop, T),
\]
the $C_\infty(\underline{t})$ vary contravariantly in $\underline{t}$ (thus define a fibered category over $\Hc(\Nop, T)$). Moreover, the functor
\[
C_\infty(\underline{t'}) \longrightarrow C_\infty(\underline{t})
\]
associated to a morphism of inverse systems
\[
\underline{t} \longrightarrow \underline{t'}
\]
is an \textit{isomorphism} in Hot. So the
\msp{page 118}
functor
\[
\underline{t} \longmapsto C_\infty(\underline{t}), \qquad \Hc(\Nop, T) \longrightarrow (\Hot)
\]
factors through the fundamental groupoid
\[
\Pi(\Hc(\Nop, T)) \longrightarrow (\Hot).
\]
When this groupoid, i.e. when the category $\Hc$ itself is connected, we can therefore conclude that all the $C_\infty(\underline{t})$ are isomorphic to one another in Hot. To conclude even that there is a \textit{transitive} system of isomorphisms between the $C_\infty(\underline{t})$, i.e. that $C_\infty(\underline{t})$ is ``unique up to \textit{unique} isomorphism'', one would have to know in addition that the fundamental groupoid is equivalent to the one-point category (trivial fundamental group), or, as one says, that $\Hc$ is \textit{simply connected,} or 1-connected.

The connectedness of $\Hc$ is assured when the category $T$ has (as is the case for $\underline{\mathrm{Ty}}$) a law of composition $(t,t') \longmapsto t * t'$, with morphisms functorial in $t, t'$
\[
\begin{tikzcd}[row sep=small]
 & t \\
t * t' \arrow[ru] \arrow[rd] & \\
 & t'
\end{tikzcd}
\]
(this last condition being automatic
% foot: Version 22 décembre 2001
\printedpage{75}
% head: A. Grothendieck : Dérivateurs, Ch. VII, Édition des Universités de Montpellier II et Paris VII
\msp{page 119}
if the law of composition admits the final object of $T$, assumed to exist, as a two-sided unit). In this case, $\underline{t} = (t_n)_n$ and $\underline{t'} = (t'_n)_n$ are linked to one another by $\underline{t} * \underline{t'} = (t_n * t'_n)_n$, since we have a morphism of inverse systems
\[
\begin{tikzcd}[column sep=small]
 & \underline{t} * \underline{t'} \arrow[ld] \arrow[rd] & \\
\underline{t} & & \underline{t'}.
\end{tikzcd}
\]
I do not know, for that matter, whether the arrow in Hot
\[
C_\infty(\underline{t}) \longrightarrow C_\infty(\underline{t'})
\]
associated to $\underline{t'} \longrightarrow \underline{t}$ in $\Hc$ depends only on $\underline{t}$, $\underline{t'}$, and not on the choice of $\underline{t'} \longrightarrow \underline{t}$. What is clear is that its image in $\Hotz$ is unique, it is the identity when we identify $C_\infty(\underline{t})$ and $C_\infty(\underline{t'})$ with one and the same element of $\Hotz$ defined by $\varinjlim H_T$.

And yet, I realize just now that, under the hypothesis of the existence of the law of composition and of its two-sided unit, and assuming in addition that the law in question is associative up to isomorphism, \textit{we have transitivity} in the isomorphisms (canonical!) between the $C_\infty(\underline{t})$ which we have just indicated!

\msp{page 120}
(Without having for all that to compute the fundamental group of $\Hc(\Nop, T)$!) If we have three elements $\underline{t}$, $\underline{t'}$, $\underline{t''}$ of $\Hc$, it suffices to look at the diagram deduced from the following commutative diagram
\[
\begin{tikzcd}[column sep=small]
 & & \underline{t} * \underline{t'} * \underline{t''} \arrow[ld] \arrow[rd] & & \\
 & \underline{t} * \underline{t'} \arrow[ld] \arrow[rd] & & \underline{t'} * \underline{t''} \arrow[ld] \arrow[rd] & \\
\underline{t} & & \underline{t'} & & \underline{t''}
\end{tikzcd}
\]
by applying to it the contrafunctor $\underline{s} \longrightarrow C_\infty(\underline{s})$ from $\Hc$ to Hot.

So in the case of the $C(a,b) = \underline{\mathrm{Ch}}_\Sigma(X; a, b)$, I no longer have any scruples about regarding the category $C_\infty(a,b)$ as a \textit{canonical} object of Hot, linked to the canonical ind-object $H_T(a,b)$ in Hot by a canonical homomorphism
\[
C_\infty(a,b) \longleftarrow H_T(a,b), \tag*{$(*)$}
\]
which in $\Hotz$ gives an isomorphism
\[
C_\infty(a,b) \simeq \varinjlim H_T(a,b), \tag*{$(**)$}
\]
i.e. the arrow $(*)$ induces, for $\xi \in \mathrm{Ob}\,\Hot$, a set-theoretic map
\[
\mathrm{Hom}_{\Hot}(\xi, C_\infty(a,b)) \longleftarrow \underbrace{\mathrm{Hom}_{\mathrm{Ind}\,\Hot}(\xi, H_T(a,b))}_{\underset{n}{\varinjlim}\, \mathrm{Hom}_{\Hot}(\xi, C_n(a,b))},
\]
which is \textit{bijective} if $\xi$ is in Hot.
% foot: Version 22 décembre 2001

\printedpage{76}
% head: A. Grothendieck : Dérivateurs, Ch. VII, Édition des Universités de Montpellier II et Paris VII
\msp{page 121}
Of course, the arrows
\[
C_\infty(\underline{t'}) \longrightarrow C_\infty(\underline{t}) \tag*{$(*)$}
\]
in Hot, associated to the arrows $\underline{t} \longrightarrow \underline{t'}$ in $\Hc(\Nop, T)$, are in fact arrows in % [sic] associés -> associées, accord avec « les flèches »
(Cat) itself - and even fully faithful functors, in the ``geometric'' case and for the two canonical arrows
\[
\begin{tikzcd}[column sep=small]
 & \underline{t} * \underline{t'} \arrow[ld] \arrow[rd] & \\
\underline{t} & & \underline{t'}.
\end{tikzcd}
\]
They are isomorphic \textit{in} Hot, because they define an isomorphism on the associated ind-objects of Cat\,hot, hence also for the associated ind-objects in Hot and in $\Hotz$, from which it follows that
\DERmark{ENV0500}

\[
\mathrm{Hom}_{\Hot}(Y, C_\infty(\underline{t'})) \longrightarrow \mathrm{Hom}_{\Hot}(Y, C_\infty(\underline{t}))
\]
associated to $(*)$ is a bijection if $Y$ \textit{is finite} (\footnote[26]{Which suffices to ensure that $(*)$ is \texttt{[an]} isomorphism in Hot: $\Hot \longrightarrow \Hotz$ is conservative!}), this arrow then identifies with
\[
\underset{n}{\varinjlim}\, \mathrm{Hom}_{\Hot}(Y, C_{t'_n}) \longrightarrow \underset{n}{\varinjlim}\, \mathrm{Hom}_{\Hot}(Y, C_{t_n}).
\]
(For $Z \longmapsto \mathrm{Hom}_{\Hot}(Y, Z)$ commutes with filtering direct limits, if $Y$ is a category which is \textit{finite.} This comes from the fact that a filtering $\varinjlim$ of Thomason-Kan categories is likewise \dots) But I do not know, on the other hand, whether $(*)$ is \texttt{[an]} isomorphism in Cat\,hot. Only that for a \textit{finite} category $Y$, it induces
\[
\mathrm{Hom}_{\mathrm{Cat\,hot}}(Y, C_\infty(\underline{t'})) \xrightarrow{\ \sim\ } \mathrm{Hom}_{\mathrm{Cat\,hot}}(Y, C_\infty(\underline{t})),
\]
the two sides then identifying with
\msp{page 122}
\[
H_T(Y)_{\mathrm{Cat\,hot}} = \mathrm{Hom}_{\mathrm{Ind\,Cat\,hot}}(Y, H_T).
\]
To sum up these insights and perplexities about the $C_\infty(\underline{t})$:

\begin{enumerate}
\item[a)] To every $\underline{t} = (t_n)_{n \in \Nop} \in \Hc(\Nop, T)$ is associated a direct system, indexed by $\mathbf{N}$, of categories $C_n(\underline{t})$ $(= C_{t_n})$, and their direct limit
\[
C_\infty(\underline{t}).
\]
% foot: Version 22 décembre 2001

\printedpage{77}
% head: A. Grothendieck : Dérivateurs, Ch. VII, Édition des Universités de Montpellier II et Paris VII
\item[b)] This direct system, and also its limit is \texttt{[a]} contravariant functor in $\underline{t}$. If $\underline{t}, \underline{t'} \in \mathrm{Ob}\,\Hc$, then we can compare $C_\infty(\underline{t})$, $C_\infty(\underline{t'})$ and the corresponding direct systems, thanks to
\[
\begin{tikzcd}[column sep=small]
 & \underline{t} * \underline{t'} \arrow[ld] \arrow[rd] & \\
\underline{t} & & \underline{t'}.
\end{tikzcd}
\]

\item[c)] The object
\[
(C_n(\underline{t}))_{n \in \mathbf{N}} \in \mathrm{Ob\,Ind\,Cat\,hot}
\]
defined by the direct system $(C_n(\underline{t}))_{n \in \mathbf{N}}$ in Cat is (up to unique isomorphism) independent of the choice of $\underline{t}$, say $H_C$ or $H$. The homomorphisms
\[
C_\infty(\underline{t'}) \xrightarrow{u^*} C_\infty(\underline{t}) \qquad \text{associated to } u : \underline{t} \longrightarrow \underline{t'} \tag*{$(*)$}
\]
are compatible with this datum, i.e. we have commutativity in
\[
\begin{tikzcd}
\underset{\mathrm{Ind\,Cat\,hot}}{\varinjlim} C_*(\underline{t'}) \arrow[rr] & & \underset{\mathrm{Ind\,Cat\,hot}}{\varinjlim} C_*(\underline{t}) \\
 & H \arrow[lu, "\sim"] \arrow[ru, "\sim"'] &
\end{tikzcd}
\]

\msp{page 123}
\item[d)] It follows that the homomorphisms $(*)$ are \textit{isomorphisms} in Hot, hence all the $C_\infty(\underline{t})$ are isomorphic \textit{in} Hot. The trick with the $t * t'$ (in the case where we have a law $t * t'$ which is associative and admits the final object as unit object) even gives \textit{canonical} isomorphisms
\[
C_\infty(\underline{t'}) \longrightarrow C_\infty(\underline{t}) \qquad \text{in Hot,}
\]
and this is a transitive system of isomorphisms. So we have in Hot a canonical object
\[
C_\infty \in \mathrm{Ob}\,\Hot,
\]
whose relation to $H \in \mathrm{Ob\,Ind\,Hot}$ is deduced from c), it is a canonical arrow
\[
H \longrightarrow C_\infty \qquad \text{in Ind\,Hot (\footnotemark[27]).}
\]
\footnotetext[27]{\textbf{NB} $H$ is practically never essentially constant, so the arrow opposite \texttt{[is]} not an isomorphism in Ind\,Hot.}

This arrow induces an isomorphism for the corresponding objects of $\Hotz$; but viewing $H$ not as an ind-object of $\Hotz$, but taking the functor $H_0$ induced on $\Hot_0$ by the ind-representable functor $H$ on Hot. This induced functor is ``representable'' by the object $C_\infty$ of the larger category
\[
H_0 \simeq C_\infty \qquad \text{in } \Hotz.
\]

(If the canonical functor $\Hot \longrightarrow \Hotz$ is fully faithful, then this formula describes $C_\infty$ as an object of Hot, up to unique isomorphism, in terms of $H_0$, hence also in terms of $H \in \mathrm{Ob\,Ind\,Cat\,hot}$.)
% foot: Version 22 décembre 2001

\printedpage{78}
% head: A. Grothendieck : Dérivateurs, Ch. VII, Édition des Universités de Montpellier II et Paris VII
\msp{page 124}
\item[e)] I do not know, on the other hand, whether the arrows
\[
C_\infty(\underline{t'}) \longrightarrow C_\infty(\underline{t})
\]
are isomorphisms in Cat\,hot, although they are defined in terms of homomorphisms of direct systems of categories, which are isomorphisms in Ind\,Cat\,hot. This is so, even for the arrows $C_\infty(\underline{t}), C_\infty(\underline{t'}) \rightrightarrows C_\infty(\underline{t} * \underline{t'})$ used to compare $C_\infty(\underline{t})$ and $C_\infty(\underline{t'})$ canonically. Thus, I do not see how to define arrows in the reverse direction in Cat, which would be candidates to be inverse isomorphisms in Cat\,hot.

\item[f)] In the geometric case ($T = \underline{\mathrm{Ty}}$, situation described in terms of $X$, $\Sigma$, and $a, b \in \mathrm{Ob}\,X$, hence $C_\infty(\underline{t}; a, b)$, $C_\infty(\overline{a, b})$, $H(a,b)$ etc.), we \textit{do not have} a path composition operation at the level of the \textit{categories} $C_\infty(\underline{t}; a, b)$ and $C_\infty(\underline{t}; b, c)$, but only at the level of the objects $H(a,b)$, $H(b,c)$ of Ind\,Cat\,hot or of Ind\,Hot. Am I even sure I know how to define \textit{in} Hot
\[
C_\infty(a,b) \times C_\infty(b,c) \longrightarrow C_\infty(a,c) \ \ ?
\]

Surely, there are no problems for the images in $\Hotz$, since there they are the same as those of $H(a,b)$, $H(b,c)$, $H(a,c)$ respectively, and we have perfect, associative pairings, with unit, and with inverses where there ought to % [sic] là ou -> là où, accent omis sur le relatif
\msp{page 125}
be some. But not for the $C_\infty(a,b)$ themselves?
\end{enumerate}

It is a bit vexing that with all this work it would seem that we haven't even been able to damn well define the analogue in Cat of the classical loop space $\Omega(X, x)$ as a group \textit{in} Hot. (What does \textsc{Thomason} do?) But here is a thinkable idea [?]. To define \textit{all the same}
\[
C_\infty(a,b) \times C_\infty(b,c) \longrightarrow C_\infty(a,c)
\]
in Hot, we are going to describe the two functors $C_\infty(a,b)$ and $C_\infty(b,c)$ by different inverse systems $\underline{t}$, $\underline{t'}$. We can make it explicit by starting from a $\tau_0$ of useful length $\geq 2$, and taking in both cases
\[
t_n = t'_n = \tau_0^n,
\]
but different transition morphisms
\[
\begin{aligned}
(t_{n+p} \longrightarrow t_n) &= \pi_p * \mathrm{id}_{\tau_0^n} \\
(t'_{n+p} \longrightarrow t'_n) &= \mathrm{id}_{\tau_0^n} * \pi_p,
\end{aligned}
\]
where $\pi_p : \tau_0^p \longrightarrow 1$ is the total degeneracy. We will define $C_\infty(a,c)$ in terms of $\underline{t''} = \underline{t} * \underline{t'}$, hence
\[
\left\{
\begin{aligned}
t''_n &= \tau_0^{2n} \\
(t''_{n+p} \longrightarrow t''_n) &= \pi_p * \mathrm{id}_{\tau_0^{2n}} * \pi_p.
\end{aligned}
\right.
\]
% foot: Version 22 décembre 2001


\endgroup
% END UNIT Der-ChVII-456/Der-ChVII-456_p076-080.tex
% BEGIN UNIT Der-ChVII-456/Der-ChVII-456_p081-085.tex
\begingroup\makeatletter%
\newcommand{\printedpage}[1]{\marginpar{\footnotesize p.~#1}}%
\newcommand{\msp}[1]{\marginpar{\footnotesize [#1]}}%
\newcommand{\Ch}{\underline{\mathrm{Ch}}}%
\newcommand{\uHom}{\underline{\mathrm{Hom}}}%
\newcommand{\comp}{\operatorname{comp}}%
\newcommand{\Hot}{\mathrm{Hot}}%
\makeatother\let\printedpage\derpp\let\msp\dermsp%
\let\mspin\DERmspin%


\printedpage{79}
% head: A. Grothendieck : Dérivateurs, Ch. VII, Édition des Universités de Montpellier II et Paris VII
% foot: Version 22 décembre 2001

This time, the composition functors
\[ C_n(\underline{t};a,b)\times C_n(\underline{t'};b,c) \longrightarrow C_n(\underline{t''};a,c), \]
are compatible with the transition functors for varying $n$, and one finds
\msp{page 126}
\[ C_\infty(\underline{t};a,b)\times C_\infty(\underline{t'};b,c) \longrightarrow C_\infty(\underline{t''};a,c), \]
hence an arrow in Hot
\[ \comp_{\tau_0} : C_\infty\times C_\infty \longrightarrow C_\infty. \]
We must see that it does not depend on the choice of $\tau_0$, but if we have $\tau_0$ and $\tau'_0$, we compare them by $\tau_0*\tau'_0 \begin{array}{l}\longrightarrow \tau_0\\ \longrightarrow \tau'_0\end{array}$, so it suffices to see that if there exists an arrow $\tau_0\longrightarrow\tau'_0$, then $\comp_{\tau_0} = \comp_{\tau'_0}$. But such an arrow induces arrows $\underline{t}_{\tau_0}\longrightarrow\underline{t}_{\tau'_0}$, $\underline{t'}_{\tau_0}\longrightarrow\underline{t'}_{\tau'_0}$, $\underline{t''}_{\tau_0}\longrightarrow\underline{t''}_{\tau'_0}$, whence a diagram in Cat
\DERmark{ENV0501}

\[
\begin{tikzcd}
C_\infty(\underline{t}_{\tau'_0};a,b)\times C_\infty(\underline{t'}_{\tau'_0};b,c) \arrow[r, "{\comp_{\tau'_0}}"] \arrow[d] & C_\infty(\underline{t''}_{\tau'_0};a,c) \arrow[d] \\
C_\infty(\underline{t}_{\tau_0};a,b)\times C_\infty(\underline{t'}_{\tau_0};b,c) \arrow[r, "{\comp_{\tau_0}}"] & C_\infty(\underline{t''}_{\tau_0};a,c),
\end{tikzcd}
\]
which one checks at once is commutative. Now this diagram can be rewritten in Hot as
\[
\begin{tikzcd}
C_\infty(a,b)\times C_\infty(b,c) \arrow[r, "{\comp_{\tau'_0}}"] \arrow[d, "{\mathrm{id}}"] & C_\infty(a,c) \arrow[d, "{\mathrm{id}}"] \\
C_\infty(a,b)\times C_\infty(b,c) \arrow[r, "{\comp_{\tau_0}}"] & C_\infty(a,c),
\end{tikzcd}
\]
\msp{page 127}
whence $\comp_{\tau_0} = \comp_{\tau'_0}$.

I shrink from the real difficulty - that of proving the associativity of this composition law. It becomes associative in $\Hot^{\wedge}_{0}$ - but for two arrows in Hot to be equal (here then
\[ C_\infty(a,b)\times C_\infty(b,c)\times C_\infty(c,d) \rightrightarrows C_\infty(a,d)\quad ), \]
is it enough that this be so for their images in $\Hot^{\wedge}_{0}$ - in other words, is the functor
\[ \Hot\longrightarrow\Hot^{\wedge}_{0}, \]
(besides being conservative) at least faithful (if not fully faithful)?

\par
\printedpage{80}
% head: A. Grothendieck : Dérivateurs, Ch. VII, Édition des Universités de Montpellier II et Paris VII
% foot: Version 22 décembre 2001

There is also the fact that this composition law is compatible with that of the $H(a,b)$ in Ind\,Hot, for the canonical arrows
\[ H(a,b)\longrightarrow C_\infty(a,b) \qquad \text{in Ind\,Hot}, \]
if we would know that these are epimorphisms, and remain so under cartesian products, we would have won.

Of course, we will know that it is associative, once we know that the composition defined here is compatible with the composition of paths in the semi-simplicial framework. But how does one proceed in this framework, to compose paths of unbounded length? Is there not exactly the same difficulty here? Unless one gets % [sic] composititon -> composition, misspelling % Reading resolved against the supplied PDF; see corrections.csv.
\DERmark{ENV0502}

\msp{page 128}
out of it by reducing at once, for the theory of $\Omega$ operations, to Kan complexes, and to paths of type $\Delta_+$. But to compose them \textit{functorially,} there still seems to me to be the same problem - one will find paths of type $\Delta_+\Delta_+ = \boldsymbol{\Delta}_2^+$, and not $\Delta_+$, but one can then compose the direct homotopies (amounts to restricting to $\Delta_+\simeq\{0,2\}\subseteq\boldsymbol{\Delta}_2^+$). Now I worked out the formalism of categories of paths without worrying about conditions of Kan type. I presume that in the semi-simplicial framework, the $\uHom(\boldsymbol{\Delta}_1,X)$ do \textit{not} give the right path object, for $X$ not Kan. (For otherwise, it would seem that for $X=X_0$ whose components are \textit{finite,} the calculation of the $\pi_i$ ought to be finite, which is nevertheless not the case, I have always been assured \dots)

I presume nevertheless, when one takes $\Sigma=\mathrm{Fl}(X)$, that the objects $\Ch_\infty(X;a,b)$ and $\Ch_\infty(\bar{X};a,b)$ are isomorphic in Hot, when one embeds $X$ in a Thomason-Kan envelope $\bar{X}$.

\par
\printedpage{81}
% head: A. Grothendieck : Dérivateurs, Ch. VII, Édition des Universités de Montpellier II et Paris VII
% foot: Version 22 décembre 2001

\msp{page 129}

\section*{11\quad Quick return to the $(\mathcal{M},\Sigma)$, from the perspective of derivators}

When one is interested in the formalism of the $\Ch_\Sigma(X;a,b)$ with a view to a \textit{derivator} $\mathbf{D}_{X,\Sigma}$ that one would like to define in terms of $X$, $\Sigma$, it would seem that the prederivator one does indeed obtain only deserves the name of derivator when one assumes that for every arrow
\[ Y\longrightarrow Y' \]
in Cat which is a Hot-equivalence, the corresponding arrows
\[ \mathrm{Hom}_{\mathrm{Ind\,Cat\,hot}}(Y',H(a,b)) \longrightarrow \mathrm{Hom}_{\mathrm{Ind\,Cat\,hot}}(Y,H(a,b)), \]
or again
\[ \varinjlim_n \mathrm{Hom}_{\mathrm{Cat\,hot}}(Y',\Ch_{n,\Sigma}(a,b)) \longrightarrow \varinjlim_n \mathrm{Hom}_{\mathrm{Cat\,hot}}(Y,\Ch_{n,\Sigma}(a,b)) \]
are bijections. I looked in V at what this means for the object $H(a,b) = (\Ch_{n,\Sigma}(a,b))_n$ of Ind\,Cat\,hot. It is this condition alone that ensures, from the point of view of derivators, a convincing interpretation of $\Ch_{\infty,\Sigma}(a,b)$,
\[ \Ch_{\infty,\Sigma}(a,b) \;=\; \varinjlim_n \Ch_{n,\Sigma}(a,b) \]
as an object of Hot, or more exactly of $\Hot^{\wedge}_{0}$, for then one has for $Y$ a \textit{finite} category
\[ \underbrace{\mathrm{Hom}_{\mathrm{IndCat\,hot}}(Y,H(a,b))}_{\simeq\,\mathrm{Hom}_{X(Y)\Sigma_Y^{-1}}(a_Y,b_Y)} \;\simeq\; \mathrm{Hom}_{\Hot}(Y,\Ch_{\infty,\Sigma}(a,b)), \]
\DERmark{ENV0503}\DERmark{ENV0504}\DERmark{ENV0505}

\msp{page 130}
for I do want the type in Hot to be the right one, that it does correspond to the homotopy of path spaces. So one would like a morphism
\[ X\xrightarrow{\;f\;}X' \]
in (Cat), which is a Hot-equivalence, to define
\[ \Ch_\infty(X;a,b) \longrightarrow \Ch_\infty(X';f(a),f(b)), \]
which are themselves also Hot-equivalences. On the other hand, this is obviously compatible with the brilliant composition of paths, which I managed to define between these objects of Hot. So (subject to having checked what I have just said) the verification of associativity would reduce to the case where $X$ is Thomason-Kan. But in this case, there is a chance that, just as in the semi-simplicial case, the $\Ch_\infty(X;a,b)$ can be defined (as objects of (Hot)) as being simply the paths of type $\Delta_+$ from $a$ to $b$, which one then composes in the trivial way. But no, these categories are discrete! % [sic] ramenerait -> ramènerait, missing accent

I give up for the moment pursuing the question of associativity any further. Ask \textsc{Thomason} about the properties of the functor $\Hot\longrightarrow\Hot^{\wedge}_{0}$.

\par
\printedpage{82}
% head: A. Grothendieck : Dérivateurs, Ch. VII, Édition des Universités de Montpellier II et Paris VII
% foot: Version 22 décembre 2001

\msp{page 131}
This means, very nearly, that one has an interpretation of $\Ch_{\infty,\Sigma}(a,b)$, as an object of Hot, as an ``external Hom'' $\mathrm{Hom\,ext}(a,b)$ with values in Hot, defined on the category $\mathcal{M}=X$ using just the admissible set of arrows $\Sigma\subseteq\mathrm{Fl}(\mathcal{M})$, an operation satisfying the condition
\[ \boxed{\mathrm{Hom}_{\Hot}(Y,\mathrm{Hom\,ext}(a,b)) \;\simeq\; \mathrm{Hom}_{\mathcal{M}(Y)\Sigma_Y^{-1}}(a_Y,b_Y)} \tag*{$(*)$} \]
\DERmark{ENV0506}\DERmark{ENV0507}

(which can also be interpreted, either as
\[ \mathrm{Hom}_{\mathbf{D}(*)}(p_!(a_Y),b) \]
\DERmark{ENV0508}\DERmark{ENV0509}\DERmark{ENV0510}

if $p_! : \underbrace{\mathcal{M}(Y)\Sigma_Y^{-1}}_{\mathbf{D}(Y)} \longrightarrow \mathcal{M}\Sigma^{-1}$ exists, or as
\DERmark{ENV0511}

\[ \mathrm{Hom}_{\mathbf{D}(*)}(a,p_*(b_Y)) \]
\DERmark{ENV0512}\DERmark{ENV0513}\DERmark{ENV0514}

if $p_* : \underbrace{\mathcal{M}(Y)\Sigma_Y^{-1}}_{\mathbf{D}(Y)} \longrightarrow \underbrace{\mathcal{M}\Sigma^{-1}}_{=\,\mathbf{D}(*)}$ exists); with the one small nuisance that this formula is valid only for $Y$ \textit{finite.} If one wants to express the right-hand side for $Y$ varying arbitrarily in Cat, an object of Hot ($\mathrm{Hom\,ext}(a,b)$) on the left-hand side does not suffice, but one needs the ind-object $H(a,b) = (\Ch_n(a,b))_{n\in\mathbf{N}}$ in Hot. So if one insists on working in a ``domain'' Diag of the derivator $\mathbf{D}=\mathbf{D}_{\mathcal{M},\Sigma}$ as large as possible, while restricting oneself to finite categories,
\DERmark{ENV0515}

\msp{page 132}
and one wants a formula $(*)$ to remain valid for every $Y$ in the domain, one must ``in all rigor'' define the $\mathrm{Hom\,ext}(a,b)$ not as an object of Hot, but of Ind\,Hot. It would therefore then be described as
\[ \mathrm{Hom\,ext}(a,b) = (\Ch_{n,\Sigma}(a,b))_n \qquad \text{(ind-object in Hot)}\ , \]
and the object $\Ch_{\infty,\Sigma}$ in Hot is a sort of direct limit of it.

I wondered whether the preliminary condition, namely that the functors $H(a,b)$ on Cat turn Hot-equivalences $Y'\longrightarrow Y$ into bijections, was really a (strong) restriction. I convinced myself that this is indeed the case, already in the special case where $\Sigma=\mathrm{Fl}(X)$. It seems to me that when the connected components of $\Omega(X,a)$ are not contractible, the said condition is practically never satisfied. I thought e.g. of the case where $X$ has the homotopy type (in Hot) of a 2-sphere, so that its $\Omega(X,a)=\Ch_\infty(X,a,a)$ must have a $\pi_1\simeq\mathbf{Z}$. Let us take $Y$ having the type in Hot of a circle, say
\[
\begin{tikzcd}
 & & a & & \\
Y\colon & 0 \arrow[ur] \arrow[dr] & & 0' \arrow[ul] \arrow[dl] & , \\
 & & b & &
\end{tikzcd}
\]

\par
\printedpage{83}
% head: A. Grothendieck : Dérivateurs, Ch. VII, Édition des Universités de Montpellier II et Paris VII
% foot: Version 22 décembre 2001

\msp{page 133}
so the functors
\[ f : Y\longrightarrow \Ch_\infty(X;\bullet,\bullet) \overset{\mathrm{def}}{=} \Omega(X,\bullet) \]
\DERmark{ENV0516}

correspond to \textit{certain} ``loops'' in $\Omega(X,\bullet)$ (but loops with unspecified origin), they have a homotopy type in the topological sense (``number of turns'' around the sphere $X$) given by an integer $n\in\mathbf{Z}$ (the ``degree'' of the loop $\lambda$). The obvious intuition is that the set of degrees $\delta$ which one can obtain in this way, with a given \textit{finite} category $Y$ which represents the circle $S^1\in(\Hot)$, is finite - $|\delta|$ surely bounded above by the number of these arrows (hence at most by 4 in the case under consideration, each arrow of $Y$ making one go \textit{at most} once around the pseudo-circle $\Omega(S^2,\bullet)$) (\footnote[28]{probably one finds at most only the degrees $\delta$ between $-2$ and $+2$.}). On the other hand, when one represents $S^1$ by arbitrarily large categories $Y_n$ (say obtained from the intervals $\Phi_n=(\Delta_+\Delta_-)^n$ by identifying source and target), one must find arbitrarily high degrees. Now every $Y_n$ maps to $Y(=Y_2)$ by an arrow
\DERmark{ENV0517}

\[ Y_n = Y' \longrightarrow Y = Y_2 \]
\msp{page 134}
which is of degree 1, hence which is a Hot-equivalence. But the corresponding map
\[ \pi_0\uHom(Y,\Omega(X,\bullet)) \longrightarrow \pi_0\uHom(Y',\Omega(X,\bullet)) \]
is not bijective. This kind of argument seems to me bound to work whenever the set of homotopy groups $\pi_j(X,x)$, for $j\geq 2$ (hence the set of homotopy groups $\pi_i$ of $\Omega(X,x)$, for $i\geq 1$) is \textit{infinite,} taking for $Y$, $Y'$ different representatives $Y_n$ of the sphere $S^i$, say the $Y_n$ obtained from $\Phi_n$ $(n\geq 2)$ by taking $\Phi_n^i = \Phi_n\times\Phi_n\times\cdots\times\Phi_n$ (which is a sort of combinatorial $i$-cube), and identifying its ``boundary'' with a single point. I presume that the number of elements in $\pi_i(\Omega(X,x))$ that one can obtain from arrows $Y_n\longrightarrow\Omega(X,x)$ is limited by the size of $Y_n$, but one can surely obtain each of them, by taking $n$ large enough. So if $\xi\in\pi_i(\Omega)$ escapes $Y_{n_0}$ ($n_0\geq2$, say), one takes an $n'$ large enough that $\xi$ does not escape $Y_{n'}$, and one chooses % [sic] d'élements -> d'éléments, missing accent
\DERmark{ENV0518}

\[ Y_{n'} \longrightarrow Y_n\ , \]
\texttt{[}$n=n_0$\texttt{?]}

\msp{page 135}
which is a Hot-equivalence (e.g. with the aid of a retraction of $\Phi_{n'}$ onto $\Phi_n$, compatible with sources and targets).

This draws strong attention to the fact that if $Y$, $C$ are two categories, the category $\uHom(Y,C)$ \textit{is not} a representative of $\uHom_{\Hot}(\mathrm{hot}(Y),\mathrm{hot}(C))$, one has only a canonical arrow
\[ \mathrm{hot}(\uHom(Y,C)) \longrightarrow \uHom_{\Hot}(\mathrm{hot}(Y),\mathrm{hot}(C)), \tag*{$(*)$} \]


\endgroup
% END UNIT Der-ChVII-456/Der-ChVII-456_p081-085.tex
% BEGIN UNIT Der-ChVII-456/Der-ChVII-456_p086-090.tex
\begingroup\makeatletter%
\newcommand{\printedpage}[1]{\marginpar{\footnotesize p.~#1}}%
\newcommand{\msp}[1]{\marginpar{\footnotesize [#1]}}%
\newcommand{\uHom}{\underline{\mathrm{Hom}}}%
\newcommand{\uCh}{\underline{\mathrm{Ch}}}%
\newcommand{\hot}{\mathrm{hot}}%
\makeatother\let\printedpage\derpp\let\msp\dermsp%
\let\mspin\DERmspin%


\printedpage{84}
% head: A. Grothendieck : Dérivateurs, Ch. VII, Édition des Universités de Montpellier II et Paris VII
% foot: Version 22 décembre 2001

i.e.
\[
\underbrace{\hot(\uHom(Y,C))\times\hot(Y)}_{\hot(\uHom(Y,C)\times Y)}\longrightarrow\hot(C),
\]
which arrow is simply deduced from
\[
\uHom(Y,C)\times Y\longrightarrow C.
\]
But the arrow $(*)$ has no desire to be \texttt{[an]} isomorphism. It is because it is already not one in the special case $Y=\boldsymbol{\Delta}^1$, essentially, which forces us, in order to make a theory of paths in categories that is topologically reasonable, not to confine ourselves to $\boldsymbol{\Delta}^1$ as a type of paths, but to take composite intervals $\tau\in\mathrm{Ty}$.

\msp{page 136}

On the other hand, $(*)$ is \texttt{[an]} isomorphism when $C$ is a Thomason-Kan category. It suffices to see that for every $Z$ in (Cat), the Hom from $\hot(Z)$ into the two sides is isomorphic via the arrow deduced from $(*)$. Now if $C$ is Thomason-Kan, so is $\uHom(Y,C)$ (immediate from the definition of Thomason-Kan in terms of extensions of
\begin{tikzcd}[baseline=(\tikzcdmatrixname-1-1.base), row sep=small]
A_0 \arrow[r, hook, "i"] \arrow[dr] & A \arrow[d, dotted] \\
& C
\end{tikzcd}
$\ldots$(\footnote[29]{a bit hasty! One has to know that the cofibrant maps $i$ of \textsc{Thomason} remain so after cartesian multiplication by $Y$ \ldots})), hence
\[
\mathrm{Hom}_{\mathrm{Hot}}(\hot(Z),\hot(\uHom(Y,C)))\xleftarrow{\ \sim\ }\pi_0(\underbrace{\uHom_{\mathrm{Cat}}(Z,\uHom(Y,C))}_{\uHom_{\mathrm{Cat}}(Z\times Y,C)})\simeq\pi_0\uHom(Z\times Y,C), % Reading resolved against the supplied PDF; see corrections.csv.
\]
and
\[
\begin{aligned}
&\mathrm{Hom}_{\mathrm{Hot}}(\hot(Z),\uHom_{\mathrm{Hot}}(\hot(Y),\hot(C)))\simeq\mathrm{Hom}_{\mathrm{Hot}}(\underbrace{\hot(Z)\times\hot(Y)}_{\hot(Z\times Y)},\hot(C))\\
&\qquad\underbrace{\xleftarrow{\ \sim\ }}_{\text{since $C$ is Thomason-Kan}}\qquad\quad\pi_0(\uHom_{\mathrm{Cat}}(Z\times Y,C)),
\end{aligned}
\]
O.K.

We thus see that the categories $\uCh_{\infty,\Sigma}(X;a,b)$ have no tendency to be Thomason-Kan categories. Yet, its type in Hot must be ``the right one'', without first having to replace $X$ by a Thomason-Kan envelope $\bar{X}$.

\printedpage{85}
% head: A. Grothendieck : Dérivateurs, Ch. VII, Édition des Universités de Montpellier II et Paris VII
% foot: Version 22 décembre 2001

\msp{page 137}

\section*{12\quad Back to the ``categories of paths'' in \textmd{Cat}, as a category of models for \textmd{Hot}}

I see exactly two interesting applications of the formalism of paths in categories. One, which above all has forced my hand, is to be able to construct Hom\,ext, \textit{external} Hom with values in Hot, for every prederivator $\mathbf{D}$ defined by a pair $(\mathcal{M},\Sigma)$ (and I know of no prederivator that is not defined this way).

The other application is to have in (Cat) a \textit{functorial construction,} in terms of an arbitrary object $X$ or an arrow $X\longrightarrow Y$, of the analog of the formalism of \textsc{Cartan-Serre} on path spaces. In this setting, it is understood that one always takes $\Sigma_X=\mathrm{Fl}(X)$, the set of \textit{all} arrows of $X$. To construct functorially in $X$ in (Cat)
\[
\begin{tikzcd}[row sep=huge]
\uCh_\infty(X) \arrow[d] \\
X\times X,
\end{tikzcd}
\]
one must have fixed an inverse system
\[
\underline{t}=(t_n)_{n\in\mathbf{N}}
\]
\msp{page 138}
in \underline{Ty} which is cofinal in \underline{Ty}, which simply means that among the $t_n$ there are some of arbitrarily large useful length (and since
\[
n\longmapsto\mathrm{long.ut.}(t_n)
\]
is increasing, this implies that long.ut.$(t_n)$ tends to $+\infty$ for $n\longrightarrow+\infty$). The main facts to check next (if it can be done) are the essential properties of the arrows in the diagram
\printedpage{86}
% head: A. Grothendieck : Dérivateurs, Ch. VII, Édition des Universités de Montpellier II et Paris VII
% foot: Version 22 décembre 2001
\[
\begin{tikzcd}[column sep=large, row sep=large]
& X \arrow[d, "\tau"] & \\
& \uCh_\infty(X) \arrow[d, "p"] \arrow[ddl, "\sigma"'] \arrow[ddr, "\beta"] & \\
& X\times X \arrow[dl, "\mathrm{pr}_1"] \arrow[dr, "\mathrm{pr}_2"'] & \\
X & & X\ ,
\end{tikzcd}
\qquad\qquad
\left|\begin{array}{l}
\sigma\tau=\beta\tau=\mathrm{id}_X\\
p\tau=\mathrm{diag}_X
\end{array}\right.
\]
and more generally, for an arrow
\[
X'\xrightarrow{\ f\ }X
\]
in Cat, of the arrows in the corresponding diagram
\[
\begin{tikzcd}
& \overline{X'}=C(f) \arrow[dl] \arrow[dr, "g"] \arrow[ddrr, bend left=30, "\bar{f}"] & & \\
X' \arrow[ur, dotted, bend left=40, "\alpha"] \arrow[dr, "f"'] & |[draw, rounded corners]| \text{cart.} & \uCh_\infty(X) \arrow[dl, "\sigma"] \arrow[dr, "\beta"] & \\
& X & & X,
\end{tikzcd}
\qquad
\bar{f}\alpha=(\beta g)\alpha=f,
\]
and more particularly of $\beta g$. The arrow
\msp{page 139}
in dotted lines $\alpha$ is the section of $\overline{X'}$ over $X'$ deduced from the section $\tau$ of $\uCh_\infty(X)$ over $X$ which realizes the decomposition of $f$ into
\[
X'\xrightarrow{\ \alpha\ }\overline{X'}\xrightarrow{\ \bar{f}\ }X,
\]
where $\bar{f}$ should be a \textit{``fibrant''} arrow, $\alpha$ a Hot-equivalence and a \textit{``cofibrant''} arrow in the sense of \textsc{Thomason}. (But whether this is so, remains to be seen!) Here for the moment are the essential properties I would like to check.
\DERmark{ENV0519}


\begin{enumerate}[label={}, leftmargin=*]
\item[1$^o$)] $p$ is smooth and proper, and is a HOT-fibration.

\item[2$^o$)] So the same holds for $\sigma=\mathrm{pr}_1p$ and $\beta=\mathrm{pr}_2p$. Moreover, these HOT-fibrations are \textit{trivial,} i.e. with aspheric fibers.

\printedpage{87}
% head: A. Grothendieck : Dérivateurs, Ch. VII, Édition des Universités de Montpellier II et Paris VII
% foot: Version 22 décembre 2001

\item[3$^o$)] More generally (as regards $\beta$), $\overline{X'}\xrightarrow{\ \bar{f}\ }X$ should still be proper and smooth and a HOT-fibration, but of course in general not trivial. (It should be trivial if and only if $f$ is a Hot-equivalence.) (\footnote[30]{follows formally from 1$^o$ and 2$^o$, cf. page 170.}).
\end{enumerate}

A word on ``duality''. The statements about $\sigma$ (resp. $\beta$) translate into statements about $\beta$ (resp. $\sigma$) by passing from $\underline{t}$ to $\underline{t}^o=(t^o_n)$, where $\tau\longmapsto\tau^o$ is the functor in \underline{Ty} ``passage to the opposite interval''. We have in fact % Reading resolved against the supplied PDF; see corrections.csv.

\msp{page 140}

an isomorphism functorial in $X$ (passage to the opposite path)
\[
\boxed{\begin{array}{ccc}
c & \longmapsto & c^o\\[4pt]
\uCh^{\underline{t}}_\infty(X) & \longrightarrow & \uCh^{\underline{t}^o}_\infty(X)
\end{array}}
\qquad\qquad
\text{(\footnotemark[31])}
\]
\footnotetext[31]{cf. page 143 for the passage from $\underline{t}$ to $\underline{t}^-$.}
which ``exchanges source and target''. Of course, we do well to take $\underline{t}$ such that we have $\underline{t}=\underline{t}^o$ (this is what I had done, by a sound instinct, right from the start in V), so that we have a canonical involutive functor
\[
\begin{array}{ccc}
c & \longmapsto & c^o\\[4pt]
\uCh_\infty(X) & \longrightarrow & \uCh_\infty(X)
\end{array}
\]
(I again drop the superscript $\underline{t}$), which ``exchanges source and target''. So we find, up to canonical isomorphism, the ``same'' factorization of $f:X'\longrightarrow X$ into $\bar{f}\circ\alpha$, if we do as on page 138 to define $\overline{X'}$ and $\alpha$, or if in this diagram we interchange the roles of $\sigma$ and $\beta$. % [sic] page 38 -> page 138, the diagram defining \overline{X'} and \alpha stands on manuscript page 138 (apparently)
\DERmark{ENV0520}


Finally:

\begin{enumerate}[label={}, leftmargin=*]
\item[4$^o$)] If
\[
c:a\longrightarrow b
\]
is an arbitrary path in $X$, we would like to deduce from this an isomorphism in Hot (of ``transport by $c$'') % [sic] isommorphisme -> isomorphisme, doubled m
\[
\overline{X'}_a\xrightarrow[\sim]{\bar{c}}\overline{X'}_b,
\]
with of course transitivity for composition,

\msp{page 141}

which will also give
\[
\overline{c^{\mathrm{opp}}}=\overline{c}^{-1}
\]
(assuming that the trivial paths $a\xrightarrow{\mathrm{id}}a$ give the identity $\overline{X'}_a\xrightarrow[\sim]{\mathrm{id}}\overline{X'}_a$!) (\footnotemark[32]\textsuperscript{,}\footnotemark[33]). Of course, $\overline{c}$ must depend only on the \textit{homotopy class} of ${c}$ (in % [sic] de $\overline{c}$ (dans -> de $c$ (dans, the homotopy class is apparently that of the path c
\DERmark{ENV0521}

\footnotetext[32]{Above all, not to confine oneself to the case where $c$ would be one of the $t_n$!}
\footnotetext[33]{describe the effect on the $\pi_0$ all the same!}
\printedpage{88}
% head: A. Grothendieck : Dérivateurs, Ch. VII, Édition des Universités de Montpellier II et Paris VII
% foot: Version 22 décembre 2001
$\uCh_\tau(X;a,b)$, or even in $C=\uCh(X;a,b)$, category of all paths in $X$, from $a$ to $b$, of unspecified type). In other words, we must have a functor
\[
\begin{array}{rcl}
\Pi(X) & \longrightarrow & \mathrm{Hot}\\[4pt]
a & \longmapsto & \hot(\overline{X'}_a)
\end{array}
\]
from the fundamental groupoid of $X$ to Hot. To tell the truth, this has nothing to do with the particular situation studied here - it must work whenever we have a HOT-fibration, here $\overline{X'}\longrightarrow X$, but also $\uCh_\infty(X)\longrightarrow X\times X$ (\footnote[34]{Here too, one should describe the effect of the transport $\overline{c}$ on the $\pi_0$ of the fibers.}). And for a $\mathbf{D}$-fibration $X'\longrightarrow X$ ($\mathbf{D}$ an arbitrary derivator), and $a\in\mathbf{D}(e)$, we must likewise have two functors
\[
\Pi(X)\rightrightarrows\mathbf{D}(e)
\]
corresponding to
\[
\left\{\begin{array}{lcl}
x & \longmapsto & H^\bullet_{\mathbf{D}}(X'_x,\bullet)\\[4pt]
x & \longmapsto & H_\bullet^{\mathbf{D}}(X'_x,\bullet).
\end{array}\right.
\]

\msp{page 142}

The case we are concerned with here would be the one where $\mathbf{D}=\mathrm{HOT}$, where we take $H^{\mathbf{D}}_\bullet$, and $a=$ final object of $\mathbf{D}(e)=\mathrm{Hot}$.)

\item[5$^o$)] If 1$^o$, 2$^o$, 3$^o$ are verified, we are morally sure that the construction of the $\uCh_\infty(X)$ is indeed the exact equivalent of the formalism of \textsc{Cartan-Serre}. To be completely sure of it, i.e. to link this theory explicitly to the classical constructions, one would have to prove a \textit{comparison statement} with semi-simplicial theory, which would essentially say this: If we denote, for $X$ in Cat, by $\tilde{X}$ the associated object in $\boldsymbol{\Delta}^\wedge$ (the ``nerve'' of $X$), then the arrow
\[
\uCh_\infty(X)\tilde{}\longrightarrow\underbrace{(X\times X)\tilde{}}_{\simeq\,\tilde{X}\times\tilde{X}}
\]
in $\boldsymbol{\Delta}^\wedge/(X\times X)\tilde{}$, is $W$-equivalent over $(X\times X)\tilde{}$ to the arrow
\[
\uCh(\tilde{X})\longrightarrow\tilde{X}\times\tilde{X}
\]
of the semi-simplicial theory of paths (assuming that there is such a classical arrow defined for \textit{every} object $X$ of $\boldsymbol{\Delta}^\wedge$). Or that we have at least canonical isomorphisms in Hot
\[
\hot_{\mathrm{Cat}}(\uCh_\infty(X;a,b))\ \simeq\ \hot_{\boldsymbol{\Delta}^\wedge}(\uCh(\tilde{X};\tilde{a},\tilde{b})).
\]
\end{enumerate}


\endgroup
% END UNIT Der-ChVII-456/Der-ChVII-456_p086-090.tex
% BEGIN UNIT Der-ChVII-456/Der-ChVII-456_p091-095.tex
\begingroup\makeatletter%
\newcommand{\printedpage}[1]{\marginpar{\footnotesize p.~#1}}%
\newcommand{\msp}[1]{\marginpar{\footnotesize [#1]}}%
\newcommand{\Ch}{\underline{\mathrm{Ch}}}%
\newcommand{\Ty}{\underline{\mathrm{Ty}}}%
\newcommand{\dashto}[1]{\mathrel{\overset{#1}{\rule[0.6ex]{2em}{0.4pt}}}}%
\makeatother\let\printedpage\derpp\let\msp\dermsp%
\let\mspin\DERmspin%


\printedpage{89}
% head: A. Grothendieck : Dérivateurs, Ch. VII, Édition des Universités de Montpellier II et Paris VII
% foot: Version 22 décembre 2001

\msp{page 143}
In $1^o$ and $3^o$, we want to prove that certain arrows ($p$ and $\bar{f}$) are both \textit{proper} and \textit{smooth} - which is really ideal! I say that each time it suffices to prove one of the two, that the other follows from it (as a ``special case'') by duality. For this, let us denote by
\[
\tau \longmapsto \tau^{-}
\]
the automorphism (\textit{not anti}-automorphism) of $\Ty$ which consists in changing the sign of all the transition arrows - it is thus, at the level of objects, the unique automorphism of the free monoid Ty, which interchanges $\Delta_+$ and $\Delta_-$. To the inverse system $\underline{t}$ in $\Ty$ corresponds the inverse system $\underline{t}^-$. That said, we have an isomorphism of categories
\[
\boxed{\Ch^{\underline{t}}_{\infty}(X^o) \ \simeq\ (\Ch^{\underline{t}^-}_{\infty}(X))^o} \tag*{$(^{35})$} % Reading resolved against the supplied PDF; see corrections.csv.
\]
\footnotetext[35]{cf. page 140 for the passage from $\underline{t}$ to $\underline{t}^o$.}
and the commutative diagram
\[
\begin{tikzcd}
\Ch^{\underline{t}}_{\infty}(X^o) \arrow[r, "\sim"] \arrow[d, "p^{\underline{t}}_{X^o}"] & \Ch^{\underline{t}^-}_{\infty}(X)^o \arrow[d, "(p^{\underline{t}^-}_{X})^o"] \\
X^o \times X^o \arrow[r, equal] & (X \times X)^o.
\end{tikzcd}
\]
Thus, to say that $p^{\underline{t}}_{X}$ \texttt{[rather} $p^{\underline{t}^-}_{X}$\texttt{]} is smooth, i.e. that $(p^{\underline{t}}_{X})^o$ \texttt{[rather} $(p^{\underline{t}^-}_{X})^o$\texttt{]} is proper, is equivalent to saying that $p^{\underline{t}}_{X^o}$ is proper.

\msp{page 144}
So we are left, on this subject, with only two facts to check: the smoothness of $p : \Ch_{\infty}(X) \longrightarrow X \times X$, and that of $\bar{f} : \overline{X'} \longrightarrow X$.

a) Smoothness of $p : \Ch_{\infty}(X) \longrightarrow X \times X$. We take an arrow
\[
w : (a,b) \longrightarrow (a',b'), \qquad \text{i.e.} \quad \left\{ \begin{array}{lccc} u : & a & \longrightarrow & a' \\[1ex] v : & b & \longrightarrow & b', \end{array} \right.
\]
in the base, an object $c'$ over the target, i.e.
\[
c' \in \Ch_{\infty}(X; a', b'),
\]
and we look at the category of arrows in $\Ch_{\infty}(X)$, with target $c'$, that lift $w$,
\[
\begin{tikzcd}
\underset{\text{\scriptsize variable}}{c} \arrow[r, "\varphi"] \arrow[d, dash] & \underset{\text{\scriptsize fixed}}{c'} \arrow[d, dash] \\
(a,b) \arrow[r, "w"] & (a',b'),
\end{tikzcd}
\]
\printedpage{90}
% head: A. Grothendieck : Dérivateurs, Ch. VII, Édition des Universités de Montpellier II et Paris VII
% foot: Version 22 décembre 2001
i.e. the category of pairs $(c,\varphi)$, where
\[
\left\{ \begin{array}{l}
c \in \Ch_{\infty}(X; a, b) \\[1ex]
\varphi : c \longrightarrow c' \ \ (\in \Ch_{\infty}(X)) \\[1ex]
\text{an arrow inducing } \begin{array}[t]{l} u \ \text{ on the sources} \\ v \ \text{ on the targets,} \end{array}
\end{array} \right.
\]
\DERmark{ENV0522}

\[
\begin{tikzcd}
a \arrow[rrrr, dotted, dash, "c"] \arrow[d] & {} \arrow[d, dotted] & {} \arrow[d, dotted] & {} \arrow[d, dotted] & b \arrow[d] \\
a' \arrow[rrrr, dash, "c'"] & {} & {} & {} & b'.
\end{tikzcd}
\]

\msp{page 145}
\begin{flushright}
{\footnotesize 9.11.90}
\end{flushright}

For the preceding question, concerning the arrows in the diagram
\[
\begin{tikzcd}
& X \arrow[d, "i"] & \\
& \Ch_{\infty}(X) \arrow[d, "p"] \arrow[ddl, "\sigma"'] \arrow[ddr, "\beta"] & \\
& X \times X \arrow[dl, "\mathrm{pr}_1"] \arrow[dr, "\mathrm{pr}_2"'] & \\
X & & X,
\end{tikzcd}
\]
we must absolutely specify the description of $\Ch_{\infty}(X)$ \textit{in} Cat, so we must specify the inverse system $\underline{t} = (t_n)_{n \in \mathbf{N}}$. We are in fact going to take % [sic] déscription -> description, misspelling
\[
t_n = (\underbrace{\Delta_+ \Delta_-}_{\tau_0})^{2n},
\]
with the transition morphisms described elsewhere
\[
\underbrace{t_{n+p}}_{=\tau_0^{2n+2p}} \ \longrightarrow \ \underbrace{t_n}_{\tau_0^{2n}}, % [sic] \tau_0^{n} -> \tau_0^{2n}, since t_n = (\Delta_+\Delta_-)^{2n} = \tau_0^{2n}
\]
\DERmark{ENV0523}

\printedpage{91}
% head: A. Grothendieck : Dérivateurs, Ch. VII, Édition des Universités de Montpellier II et Paris VII
% foot: Version 22 décembre 2001
being given by $\delta_p * \mathrm{id}_{\tau_0^{2n}} * \delta_p$, where $\delta_p : \tau_0^p \longrightarrow 1$ is the degeneracy of $\tau_0^p$. Thus we can regard $\Ch_{\infty}(X)$ as being the category of bi-infinite paths in $X$ % [sic] \tau_o^{2n} -> \tau_0^{2n}, letter o printed in the index instead of 0
\DERmark{ENV0524}

\[
\cdots \longrightarrow a_{-5} \longleftarrow \overbrace{a_{-4} \longrightarrow a_{-3} \xleftarrow{u_{-2}} \overbrace{a_{-2} \xrightarrow{u_{-1}} a_{-1} \xleftarrow{u_0} \overbrace{a_0}^{\Ch_0} \xrightarrow{u_1} a_1 \xleftarrow{u_2} a_2}^{\Ch_1} \xrightarrow{u_3} a_3 \longleftarrow a_4}^{\Ch_2} \longrightarrow a_5 \longleftarrow \cdots,
\]
with the $u_i$ identities for $|i|$ large enough.

The (full) subcategory
\[
\Ch^{W}_{t_n}(X) \ = \ \Ch^{W}_{(\Delta_+\Delta_-)^{2n}}
\]
is identified with the one for which we have
\begin{quote}
$u_i$ an identity for $i \notin [-2n+1, 2n]$, i.e. the only ``significant'' $u_i$ are the $4n$ arrows $u_{-2n+1}$, $u_{-2n+2}$, $\ldots$, $u_{2n}$. % [sic] i (in [-2i+1,2i], 4i, u_{-2i+1}, u_{-2i+2}, u_{2i}) -> n, the index of t_n is n
\DERmark{ENV0525}\DERmark{ENV0526}\DERmark{ENV0527}\DERmark{ENV0528}\DERmark{ENV0529}\DERmark{ENV0530}

\end{quote}

But it will be more convenient to regard $\Ch_{\infty}(X)$ as the direct limit of a
\msp{page 146}
doubly indexed family of full subcategories
\[
\Ch_{n,m}(X), \qquad n, m \in \mathbf{Z},
\]
defining this category as \textit{empty} if $n > m$, and for $n \leq m$ as consisting of the objects
\[
\left\{ \begin{array}{l}
(a_i, u_i)_{i \in \mathbf{Z}} \in \Ch^{W}_{n,m}(X) \iff u_i \text{ an identity for } i \notin [n+1, m] \\[1ex]
\overset{\mathrm{id}}{\cdots}\, a_{n-1} \dashto{\mathrm{id}} \underbrace{a_n \dashto{u_{n+1}} a_{n+1} \cdots \dashto{} a_{m-1} \dashto{u_m} a_m}_{\text{\scriptsize ``useful part'' of the path}} \dashto{\mathrm{id}} a_{m+1} \overset{\mathrm{id}}{\cdots}\,.
\end{array} \right.
\]
\DERmark{ENV0531}\DERmark{ENV0532}\DERmark{ENV0533}\DERmark{ENV0534}\DERmark{ENV0535}\DERmark{ENV0536}

Thus for $n = m$, we find the full subcategory isomorphic to $X$, consisting of the ``totally degenerate'' infinite paths, i.e. such that \textit{all} the $u_i$ are identities. (This subcategory of $\Ch_{\infty}$ therefore does not depend on the choice of $n = m$.)
\[
\xrightarrow{\mathrm{id}} \underbrace{a_{-1}}_{=\, a} \xleftarrow{\mathrm{id}} \underbrace{a_0}_{=\, a} \xrightarrow{\mathrm{id}} \underbrace{a_1}_{=\, a} \xleftarrow{\mathrm{id}}
\]
All the vertices of such a path are equal to one and the same object of $X$, and the transition arrows between them are identical.

The structure of $\Ch^{W}_{n,m}(X) \ \simeq \ \Ch^{W}_{\tau_{n,m}}(X)$ depends only on the length $m-n$ of $\tau_{n,m}$ and on the \textit{parity} of $n$. If $n$ is even, the initial direction of $\tau_{n,m}$ is positive, and
\[
\tau_{n,m} = \underbrace{\Delta_+ \Delta_- \cdots}_{m-n \text{ factors in alternation}} = \left\{ \begin{array}{ll}
(\Delta_+ \Delta_-)^{\frac{m-n}{2}} & \text{if } m \text{ even} \\[1ex]
(\Delta_+ \Delta_-)^{\frac{m-n-1}{2}} \Delta_+ & \text{if } m \text{ odd}
\end{array} \right.
\]
\DERmark{ENV0537}

\texttt{[rather} $\frac{m-n-1}{2}$ \texttt{instead of} $\frac{m-n}{2} - 1$\texttt{]}, whereas for $n$ odd it is just the reverse
\printedpage{92}
% head: A. Grothendieck : Dérivateurs, Ch. VII, Édition des Universités de Montpellier II et Paris VII
% foot: Version 22 décembre 2001
\msp{page 147}
\[
\tau_{n,m} = \underbrace{\Delta_- \Delta_+ \cdots}_{m-n \text{ alternating factors}} = \left\{ \begin{array}{ll}
(\Delta_- \Delta_+)^{\frac{m-n}{2}} & \text{if } m \text{ odd} \\[1ex]
(\Delta_- \Delta_+)^{\frac{m-n-1}{2}} \Delta_- & \text{if } m \text{ even}
\end{array} \right.
\]
\DERmark{ENV0538}

\texttt{[rather} $\frac{m-n-1}{2}$ \texttt{instead of} $\frac{m-n}{2} - 1$\texttt{]}. The initial direction agrees with the parity of $n$, i.e. equals $(-1)^n$, whereas the terminal direction is opposite to the parity of $m$, i.e. equals $(-1)^{m+1}$ \texttt{[rather} $(-1)^{m+1}$\texttt{]}.
\DERmark{ENV0539}


I note that we have $\tau_0 = (\tau_0)^o$, whence
\[
\underline{t} = \underline{t}^o, \qquad \text{but } \underline{t} \neq \underline{t}^-,
\]
and yet we have isomorphisms functorial in $X$
\[
\Ch_{\underline{t}}(X) \ \simeq \ \Ch_{\underline{t}^-}(X)
\]
respecting source and target, because of an isomorphism
\[
\underline{t} \ \simeq \ \underline{t}^- \qquad \text{in Pro}(\Ty).
\]
We can say e.g. that the \textit{translation functor} $T^p$ of amplitude $p$ (for $p \in \mathbf{Z}$) is defined on $\Ch^{W}_{\underline{t}}(X)$, by
\[
\left\{ \begin{array}{l}
T^p(((a_i),(u_i))_{i \in \mathbf{Z}}) = ((b_j),(v_j))_{j \in \mathbf{Z}} \\[1ex]
\text{with} \left\{ \begin{array}{lcl} b_j & = & a_{j-p} \\[1ex] v_j & = & u_{j-p} \end{array} \right.
\end{array} \right.
\]
and that
\[
\begin{aligned}
T^p &: \left\{ \begin{array}{lcl} \Ch_{t}(X) & \xrightarrow{\sim} & \Ch_{-t}(X) \\[1ex] \Ch_{-t}(X) & \xrightarrow{\sim} & \Ch_{t}(X) \end{array} \right\} \ \text{if } p \text{ odd} \\
T^p &: \left\{ \begin{array}{lcl} \Ch_{t}(X) & \xrightarrow{\sim} & \Ch_{t}(X) \\[1ex] \Ch_{-t}(X) & \xrightarrow{\sim} & \Ch_{-t}(X) \end{array} \right\} \ \text{if } p \text{ even.}
\end{aligned} % Reading resolved against the supplied PDF; see corrections.csv.
\]

\msp{page 148}
In other words, we can formally regard the sum category
\[
\Ch_{\pm}(X) = \Ch_{+}(X) \amalg \Ch_{-}(X)
\]
and the translation operation $T$ as an automorphism which interchanges the two pieces, and $T^p$ is nothing other than its $p$-th iterate, which interchanges the pieces $+$ and $-$ if $p$ odd, which stabilizes them if $p$ even.

\printedpage{93}
% head: A. Grothendieck : Dérivateurs, Ch. VII, Édition des Universités de Montpellier II et Paris VII
% foot: Version 22 décembre 2001

There is also the involutive automorphism of passage to the opposite path, which interchanges source and target, and which I define (in accordance with the canonical identity isomorphism
\[
\underline{t} \xrightarrow{\sim} \underline{t}^o \quad )
\]
as the symmetry $S$
\[
\left\{ \begin{array}{l}
S((a_i),(u_i)) = ((b_j),(u_j)) \ \ \text{with} \\[1ex]
\begin{array}{lcl} b_j & = & a_{-j} \\[1ex] u_j & = & u_{-j+1}. \end{array}
\end{array} \right.
\]
If I regard it as acting on $\Ch^{W}_{\pm}(X)$, it is related to the functor $T$ by
\[
\left\{ \begin{array}{rcl}
S^2 & = & \mathrm{id} \\[1ex]
STS^{-1} & = & T^{-1},
\end{array} \right.
\]
so that it is the semidirect group of $\mathbf{Z}/2\mathbf{Z}$ by $\mathbf{Z}$ (affine group of dimension 1 of the ring $\mathbf{Z}$) that acts on $\Ch_{\pm}$.

\msp{page 149}
We are embarrassed to find an automorphism $T^2$ of the situation $\Ch_{\infty}(X) \longrightarrow X \times X$, without it being obvious that in Hot, this functor induces the identity! I see no way to connect a path $c$ and $T^2 c$ by an arrow in $\Ch_{\infty}$, nor even a ``canonical'' path. But I tell myself that this automorphism is described by an automorphism, not of the inverse system % [sic] embarassé -> embarrassé, misspelling
\[
(\tau_{n,m})_{n,m \in \mathbf{Z}^o \times \mathbf{Z},\ n \leq m}
\]
(where $\mathbf{Z}^o$ denotes the ordered set \textit{opposite} to $\mathbf{Z}$), but of the \textit{pro-object} of $\Ty$ which it defines (which is of course ``the same'' for all the inverse systems cofinal on the left in $\Ty$). It already follows from this that $T^2$ induces the identity on the ind-object of Cat\,hot, $\underline{\underline{\mathrm{Ch}}}_{\infty}(X)$, and a fortiori on the corresponding ind-object of Hot, or of $\mathrm{Hot}^{\wedge}_{0}$, or finally for the object of $\mathrm{Hot}^{\wedge}_{0}$ itself. Likewise for the $\Ch_{\infty}(X; a, b)$. It follows, returning to $\Ch_{\infty}(X)$ itself and to the $\Ch_{\infty}(X; a, b)$, that the automorphism induced by $T^2$ is the identity in $\mathrm{Hot}^{\wedge}_{0}$. I hope that $\mathrm{Hot} \longrightarrow \mathrm{Hot}^{\wedge}_{0}$ is
\msp{page 150}
at least faithful, so that this will prove that $T^2$ induces the identity in Hot as well.

I return to questions $1^o$ to $4^o$ of a few days ago.

For $1^o$), we must prove that
\[
p : \Ch_{\infty}(X) \longrightarrow X \times X
\]
is smooth and proper and a Hot-fibration. By duality, $p$ proper will follow from $p$ smooth. But $\Ch_{\infty}(X)$ is $\varinjlim$ of the $\Ch_{n,m}(X)$, with $n \longrightarrow -\infty$, $m \longrightarrow +\infty$, and where we can restrict ourselves to $n$ and $m$ \textit{even.} Now we see that for $n$, $m$ even, i.e. $\tau_{n,m}$ with forward initial arrow, with final arrow


\endgroup
% END UNIT Der-ChVII-456/Der-ChVII-456_p091-095.tex
% BEGIN UNIT Der-ChVII-456/Der-ChVII-456_p096-100.tex
\begingroup\makeatletter%
\newcommand{\printedpage}[1]{\marginpar{\footnotesize p.~#1}}%
\newcommand{\msp}[1]{\marginpar{\footnotesize [#1]}}%
\newcommand{\Ch}{\underline{\mathrm{Ch}}}%
\newcommand{\CCh}{\underline{\underline{\mathrm{Ch}}}}%
\newcommand{\Xb}{\overline{X'}}%
\newcommand{\Xbb}{\overline{\overline{X'}}}%
\newcommand{\Hotz}{\mathrm{Hot}^{\wedge}_{0}}%
\newcommand{\Ty}{\underline{\mathrm{Ty}}}%
\makeatother\let\printedpage\derpp\let\msp\dermsp%
\let\mspin\DERmspin%


\printedpage{94}
% head: A. Grothendieck : Dérivateurs, Ch. VII, Édition des Universités de Montpellier II et Paris VII
retrograde, $\Ch_{n,m}(X) \simeq \Ch_{\tau_{n,m}}(X)$ is Cat-\textit{fibered} over $X \times X$, hence smooth (i.e. HOT-smooth), and this passes to $\varinjlim$ (\footnote[36]{\textbf{NB} $\Ch_\infty(X)$ is neither fibered nor cofibered over $X \times X$ - for the inclusion functors between the $\Ch_{n,m}$ are not cartesian resp. cocartesian.}).

Thus $\Ch_\infty(X) \xrightarrow{p} X \times X$ is both proper and smooth (\footnote[37]{12.11. But that already implies that it is a fibration!}). To say that it is a HOT-fibration, amounts e.g. to saying that for
\[
(a,b) \xrightarrow{(u,v)} (a',b')
\]
in $X \times X$, the corresponding arrow
\[
\Ch_\infty(X)/(a,b) \longrightarrow \Ch_\infty(X)/(a',b')
\]
is a Hot-equivalence. Now, since $p$ is proper, the inclusions of the fibers at $(a,b)$ and at $(a',b')$ into the respective terms are Hot-equivalences:
\msp{page 151}
\[
\begin{tikzcd}[row sep=large]
\Ch_\infty(X)/(a,b) \arrow[r, "\rho"] & \Ch_\infty(X)/(a',b') \\
\Ch_\infty(X;a,b) \arrow[u, hook, "\wr"] \arrow[r, dotted, "\varphi"] & \Ch_\infty(X;a',b') \arrow[u, hook, "\wr"]
\end{tikzcd}
\tag*{$(*)$}
\]
Thus we find a dotted arrow in Hot, which we have to prove is a Hot-equivalence. But by considering
\[
u : a \longrightarrow a' \qquad v : b \longrightarrow b'
\]
as \textit{paths} in $X$, and
\[
u^o : a' \longleftarrow a
\]
as a retrograde path from $a'$ to $a$, the general variance theory (in IX) does indeed give us an isomorphism in (Hot):
\[
\Ch_\infty(X;a,b) \longrightarrow \Ch_\infty(X;a',b'),
\]
we feel like saying that it is $\varphi$, i.e. that by inserting it in place of $\varphi$, the square above becomes commutative (which would prove what we want). I think that if we work with the ind-objects $\CCh_\infty(X)/(a,b)$, $\CCh_\infty(X;a,b)$ etc., for which we have pairings which can be spelled out nicely at the finite stages, and give rise to the $\varphi$, this should be easy (though a little tedious). This would not for all that prove the commutativity of $(*)$ in Hot, but
% foot: Version 22 décembre 2001
\printedpage{95}
% head: A. Grothendieck : Dérivateurs, Ch. VII, Édition des Universités de Montpellier II et Paris VII
\msp{page 152}
only in $\Hotz$, which is perhaps not sufficient. But this suffices to prove that the arrow $\rho$ in $(*)$ becomes \texttt{[an]} isomorphism in $\Hotz$, hence it is \texttt{[an]} isomorphism.

2$^o$) We have to show that the HOT-fibrations
\[
\Ch_\infty(X) \overset{\sigma}{\underset{\beta}{\rightrightarrows}} X
\]
are trivial, or what amounts to the same thing, that they have aspheric fibers. It suffices to see this for $\sigma$, hence that
\[
\Ch_\infty(X;a,?)
\]
is aspheric. But this category is the direct limit of the $\Ch_n(X;a,?)$, i.e. of the
\[
\underline{\mathrm{Hom}}^{\text{fixed source}}(\tau^n, X).
\]
But since the interval $I_{\tau^n}$ retracts by deformation onto its origin, it is immediate that this \underline{Hom} category is even contractible.

3$^o$) Let us consider
\[
\begin{tikzcd}[column sep=small, row sep=large]
 & \Xb = C(f) \arrow[dl] \arrow[dr, "g"] \arrow[ddrr, bend left, "\bar f"] & & \\
X' \arrow[ur, dotted, bend left, "\alpha"] \arrow[dr, "f"'] & \underline{\text{cart.}} & \Ch_\infty(X) \arrow[dl, "\sigma"'] \arrow[dr, "\beta"] & \\
 & X & & X,
\end{tikzcd}
\]
\msp{page 153}
I would like to prove that $\bar f$ is a proper and smooth \textbf{D}-fibration (\footnote[38]{That follows formally from what has already been seen, cf. pages 170-171.}).

$\Xb$ is the category of pairs
\[
(x',c) \ \in \ \mathrm{Ob}\,X' \times \mathrm{Ch}_\infty(X)
\]
such that
\[
\sigma(c) = f(x')
\]
and
\[
\bar f(x',c) = \beta(c), \qquad \text{target of } c.
\]
The fiber of $\Xb$, $\bar f$ at $x_0 \in X$ is thus the category of pairs $(x',c) \in \mathrm{Ob}\,X' \times \mathrm{Ch}_\infty(X)$, such that
\[
\sigma(c) = f(x'), \qquad \beta(c) = x_0.
\]
% foot: Version 22 décembre 2001

\printedpage{96}
% head: A. Grothendieck : Dérivateurs, Ch. VII, Édition des Universités de Montpellier II et Paris VII
We can also regard $\Xb$ as the direct limit of the subcategories $\Xb_{n,m}$, corresponding to the pairs $(x',c)$ with $c \in \mathrm{Ch}_{n,m}(X)$. We find that for $m$ even (terminal direction of $\tau_{n,m}$ retrograde) the $\Xb_{n,m}$ are Cat-fibered over $X$, for $m$ odd (terminal direction of $\tau_{n,m}$ direct) they are Cat-cofibered. So we conclude again as before, $\Xb$ being (at will) a filtering direct limit of subcategories smooth over $X$, or proper over $X$. % [sic] correspondante -> correspondantes, agreement with « sous-catégories »

\msp{page 154}
Finally we have to prove that $\Xb \xrightarrow{\bar f} X$ is a \textbf{D}-fibration (\footnote[39]{This already follows from the fact that it is proper and smooth!}), which amounts e.g. to saying that for an arrow
\[
a \xrightarrow{v} b
\]
in $X$, the induced arrow
\[
\Xb/a \xrightarrow{\bar v} \Xb/b
\]
is a Hot-equivalence. We already know as before that we have a diagram with Hot-equivalences as vertical arrows
\[
\begin{tikzcd}[row sep=large]
\Xb/a \arrow[r] & \Xb/b \\
\Xb_a \arrow[u, hook, "\wr"] \arrow[r, dotted] & \Xb_b, \arrow[u, hook, "\wr"]
\end{tikzcd}
\]
and it all comes down to showing that a certain arrow between $\Xb_a$ and $\Xb_b$ in Hot is \texttt{[an]} isomorphism. I feel like defining such an arrow directly in Hot, or even in Cat. Now for odd $m$, the $\Xb_{n,m}$ are Cat-cofibered over $X$. For $\Xb_{-m,m}$ ($m$ odd) we therefore find
\[
(\Xb_{-m,m})_a \xrightarrow{v_*} (\Xb_{-m,m})_b,
\]
\msp{page 155}
the trouble is that they are not compatible with the transition morphisms for variable $m$. They surely are, however, \textit{up to homotopy}, so we will find at least an arrow of ind-objects in Cat\,hot
\[
\Xbb_a \xrightarrow{v_*} \Xbb_b,
\]
which by passing to $\Hotz$ gives an arrow all the same
\[
\Xb_a \xrightarrow{v_*} \Xb_b.
\]
Working with even $m$, we likewise find functors in the opposite directions % [sic] m impairs -> m pairs, the Cat-fibrées (giving v^*) are stated above for m pair
\DERmark{ENV0540}

\[
(\Xb_{-m,m})_a \xleftarrow{v^*} (\Xb_{-m,m})_b,
\]
% foot: Version 22 décembre 2001
\printedpage{97}
% head: A. Grothendieck : Dérivateurs, Ch. VII, Édition des Universités de Montpellier II et Paris VII
which do not glue together for variable $m$, but which glue together up to homotopy, to give an arrow in Ind\,Cat\,hot
\[
\Xbb_a \xleftarrow{v^*} \Xbb_b,
\]
whence also
\[
\Xb_a \xleftarrow{v^*} \Xb_b \qquad \text{in } \Hotz.
\]
On the other hand, we should be able to show soritally,
\msp{page 156}
stage by stage, that $v_*v^*$ and $v^*v_*$ in Ind\,Cat\,hot are equal to the identity, hence $v_*$ and $v^*$ are isomorphisms inverse to each other. We must check on the other hand soritally that
\[
\begin{tikzcd}[row sep=large]
\Xbb/a \arrow[r, "\bar v"] & \Xbb/b \\
\Xbb_a \arrow[u] \arrow[r, "v_*"] & \Xbb_b, \arrow[u]
\end{tikzcd}
\]
in Ind\,Cat\,hot is commutative. Passing to the corresponding diagram in $\Hotz$, by taking direct limits, we find a commutative square in $\Hotz$
\[
\begin{tikzcd}[row sep=large]
\Xb/a \arrow[r, "\bar v"] & \Xb/b \\
\Xb_a \arrow[u, "\wr"'] \arrow[r, "v_*", "\sim"'] & \Xb_b, \arrow[u, "\wr"']
\end{tikzcd}
\]
whence it follows that $\bar v$ is \texttt{[an]} isomorphism in $\Hotz$, hence also in Hot.

To check that the HOT-fibration $\bar f$ is trivial if and only if $f$ is a HOT-equivalence, we would have to have developed in Cat the formalism of the exact sequence of loops - we are very close, evidently!
% foot: Version 22 décembre 2001

\printedpage{98}
% head: A. Grothendieck : Dérivateurs, Ch. VII, Édition des Universités de Montpellier II et Paris VII
\msp{page 157}

\begin{flushright}
\footnotesize 13.11.90
\end{flushright}

\section*{13\quad Where we finally prove that $\Ch(X;-,-)$ is as good as (or even better than) $\Ch_\infty(X;-,-)$.}

I believe after all that the categories $\Ch(X)$, $\Ch(X;a,b)$ (without the subscript $\infty$!) are ``good'', that the canonical arrows in Hot
\[
\left\{
\begin{aligned}
\Ch_\infty(X) \quad &\xrightarrow{\sim} \quad \Ch(X)\\
\Ch_\infty(X;a,b) \quad &\xrightarrow{\sim} \quad \Ch(X;a,b)
\end{aligned}
\right.
\]
are isomorphisms. (The first is even \texttt{[an]} isomorphism in $\mathrm{HOT}(X \times X)$, which implies the second relation, on passing to the fibers \dots) I believe that I am now (after XI) in a position to prove it. This will follow from the following two key results:

\begin{enumerate}
\item[1)] The canonical functor
\[
X \longrightarrow \Ch(X)
\]
(associating to every $x \in \mathrm{Ob}\,X$ the empty path with source and target $x$) is a Hot-equivalence, and: % [sic] but $X$ -> but $x$, the object is $x \in \mathrm{Ob}\,X$
\DERmark{ENV0541}

\item[2)] The source-target functor
\[
\Ch(X) \longrightarrow X \times X
\]
is \textit{smooth} and \textit{proper} (hence a ``perfect Hot-fibration'').
\end{enumerate}

This should give what we are looking for. Since the $X \times X$-arrow $\Ch_\infty(X) \longrightarrow \Ch(X)$ (and the deduced arrows $\Ch_\infty(X;a,b) \longrightarrow \Ch(X;a,b)$) was defined by way of the category
\[
\mathcal{C}_{\mathbf{N}^o} = \mathbf{N}^o \times_{\Ty} \mathcal{C}
\qquad
\begin{aligned}
&\text{(where $\mathcal{C} = \Ch(X)$ \texttt{[is]} considered as a fibered category over $\Ty$,}\\
&\text{and $\mathbf{N}^o \longrightarrow \Ty$ is the inverse system of the $\underbrace{(\Delta_+\Delta_-)^{2n}}_{t_n} \dots$)}
\end{aligned}
\]
\DERmark{ENV0542}

\msp{page 158}
with the help of the two canonical arrows
\[
\begin{tikzcd}[row sep=huge, column sep=large]
\mathcal{C}_{\mathbf{N}^o} \arrow[r] \arrow[d] & \Ch_\infty(X) \\
\Ch(X), &
\end{tikzcd}
\]
% foot: Version 22 décembre 2001


\endgroup
% END UNIT Der-ChVII-456/Der-ChVII-456_p096-100.tex
% BEGIN UNIT Der-ChVII-456/Der-ChVII-456_p101-105.tex
\begingroup\makeatletter%
\newcommand{\printedpage}[1]{\marginpar{\footnotesize p.~#1}}%
\newcommand{\msp}[1]{\marginpar{\footnotesize [#1]}}%
\newcommand{\Ch}{\underline{\mathrm{Ch}}}%
\newcommand{\Ty}{\underline{\mathrm{Ty}}}%
\newcommand{\cC}{\mathcal{C}}%
\newcommand{\cX}{\mathcal{X}}%
\newcommand{\NN}{\mathbf{N}^{o}}%
\newcommand{\Ob}{\operatorname{Ob}}%
\makeatother\let\printedpage\derpp\let\msp\dermsp%
\let\mspin\DERmspin%


\printedpage{99}
% head: A. Grothendieck : Dérivateurs, Ch. VII, Édition des Universités de Montpellier II et Paris VII
we will also have to prove the results corresponding to 1) 2) for $\cC_{\NN}$, i.e.

\begin{itemize}[label={}]
\item[1$'$)]
\[
\begin{tikzcd}
X \ar[r] \ar[d,"\wr"'] & \cC_{\NN} \\
\cC_{t_0} = \cC_0 &
\end{tikzcd}
\]
(the canonical inclusion of $\cC_0$, fiber of $\cC_{\NN}$ (fibered over $\mathbf{N}^o$) at $0 \in \mathbf{N}^o$) is a Hot-equivalence, and:

\item[2$'$)] $\cC_{\NN} \longrightarrow X \times X$ is proper and smooth.
\end{itemize}

With these four technical results, which I am going to prove shortly (and $1'$ and $2'$ come as no surprise to me), we win. Let us indeed consider the diagram
\[
\begin{tikzcd}
& X \ar[dl,"\alpha_1"'] \ar[d,"\alpha_0"] \ar[dr,"\alpha_2"] & \\
\Ch(X) \ar[dr,"p_1"'] & \cC_{\NN} \ar[l,"\pi_1"'] \ar[r,"\pi_2"] \ar[d,"p_0"] & \Ch_{\infty}(X) \ar[dl,"p_2"] \\
& X \times X. &
\end{tikzcd}
\]

By 1) and $1'$) (concerning $\alpha_1$ and $\alpha_0$) and the similar result
\msp{page 159}
already known for $\alpha_2$, these three functors are Hot-equivalences. So the same holds for $\pi_1$, $\pi_2$. Since $p_1$, $p_0$, $p_2$ are Hot-fibrations (being proper and smooth), it follows that $\pi_1$ and $\pi_2$ induce Hot-equivalences on the fibers
\[
\Ch(X;a,b) \quad \xleftarrow{\ \pi_1(a,b)\ } \quad \cC_{\NN}(a,b) \quad \xrightarrow{\ \pi_2(a,b)\ } \quad \Ch_{\infty}(X;a,b).
\]
I had checked this (without difficulty) already for $\pi_2$ and the $\pi_2(a,b)$, which had made it possible, by inverting $\pi_2$ in Hot, to define \textit{in} Hot
\[
\left\{
\begin{aligned}
\Ch_{\infty}(X) &\quad\longrightarrow\quad \Ch(X) \\
\Ch_{\infty}(X;a,b) &\quad\longrightarrow\quad \Ch(X;a,b),
\end{aligned}
\right.
\]
and to say that these are isomorphisms means precisely that $\pi_1$ and the $\pi_1(a,b)$ are isomorphisms.

I do not think that $\cC_{\NN}$ is in itself of any other interest, than to make it possible to compare % [sic] interêt -> intérêt, missing acute accent
the $X \times X$-categories $\Ch_{\infty}(X)$ and $\Ch(X)$, with the help of the preceding diagram. On the other hand, $\Ch_{\infty}(X)$ and $\Ch(X)$ appear to me as the two essential ingredients of the formalism of % [sic] m'aparaissent -> m'apparaissent, spelling
paths in Cat.
% foot: Version 22 décembre 2001

\printedpage{100}
% head: A. Grothendieck : Dérivateurs, Ch. VII, Édition des Universités de Montpellier II et Paris VII
\msp{page 160}
I am going to prove 1, 2, $1'$, $2'$.

1) Let us consider the diagram
\[
\begin{tikzcd}
X = (X_{\Ty})_e \ar[rr,hook,"i"] \ar[dr,dotted] & & X_{\Ty} = X \times \Ty \ar[rr,"j"] \ar[dr] & & \Ch(X) \ar[dl] \\
& e \ar[rr,dotted] & & \Ty\ , &
\end{tikzcd}
\]
where $X_{\Ty}$ and $\Ch(X)$ are considered as Cat-fibered categories over $\Ty$, and the functor $j$ is deduced from the canonical functors
\[
X \longrightarrow \Ch(X)_{\tau} \overset{\mathrm{def}}{=} \Ch_{\tau}(X),
\]
\DERmark{ENV0543}

associating to $x \in \Ob X$ the path $p_{\tau}^{*}(x)$, where $p_{\tau} : \tau \longrightarrow e$ is the projection from $\tau$ to the unit element (and final element) $e = 1_{\Ty}$ of $\Ty$, and $\Ch_{e}(X)$ is identified with $X$. We have to check the

\textbf{Lemma.} \textit{For every $\tau \in \Ty$, the preceding functor is a} Hot-\textit{equivalence (and even a homotopism).}

But we know that for two categories $\cX$, $\cC$ smooth over another, $T$, every $T$-morphism from $\cX$ to $\cC$ that induces Hot-equivalences on the fibers, is a Hot-equivalence. Thus $j$ is
\msp{page 161}
a Hot-equivalence. On the other hand, $\Ty$, having a final object, is aspheric, hence % [sic] est ... est -> est, word repeated across the manuscript page break
\DERmark{ENV0544}

$X \times \Ty \xrightarrow{\mathrm{pr}_1} X$ is a Hot-equivalence, whence it follows at once that $i$ is a Hot-equivalence. So likewise $ji$, now $ji = \alpha_1$, q.e.d.

2) Let us prove e.g. the smoothness of $\Ch(X) \longrightarrow X \times X$, properness will follow from it by applying this result to $X^{\mathrm{opp}}$. So let
\DERmark{ENV0545}

\[
(a',b') \ \longrightarrow\ (a,b), \qquad \text{i.e. } a' \xrightarrow{u} a,\ \ b' \xrightarrow{v} b,
\]
be an arrow in $X \times X$, and $c$ an object of $\Ch(X)$ over $(a,b)$:
\[
\begin{tikzcd}[column sep=large]
a' \ar[r,dotted,"c'","\scriptstyle\rightarrow"'] \ar[d,"u"'] & b' \ar[d,"v"] \\
a \ar[r,"c","\scriptstyle\rightarrow"'] & b.
\end{tikzcd}
\]

We look at the category $\cC(c;u,v)$ of pairs $(c',w)$ that lift $(c,(u,v))$, where therefore $c'$ is a path (of unspecified type) from $a'$ to $b'$, and
\[
w : c' \longrightarrow c
\]
an arrow in $\Ch(X)$ which on the source objects induces $u$, and on the target objects induces $v$ (thus $c'$ is a path from $a'$ to $b'$). We have to
% foot: Version 22 décembre 2001
\printedpage{101}
% head: A. Grothendieck : Dérivateurs, Ch. VII, Édition des Universités de Montpellier II et Paris VII
\msp{page 162}
prove that this category $\cC(c;u,v)$ is aspheric. For this, we look at the composite path
\[
\tilde{c} = \mathrm{ch}(u)\ c\ \mathrm{ch}(v)^{o},
\]
where $\mathrm{ch}(u)$ is the (direct) path from $a'$ to $a$ defined by $u : a' \longrightarrow a$, and $\mathrm{ch}(v)^{o}$ therefore denotes the (retrograde) path from $b$ to $b'$ defined by $v : b' \longrightarrow b$. If $c$ is of type $\tau$, $\tilde{c}$ is therefore of type $\Delta_{+}\tau\Delta_{-}$. We are going to compare $\cC(c;u,v)$ and the category $\Ch(X)/\tilde{c}$, with the help of two functors
\[
\begin{tikzcd}
\cC(c;u,v) \ar[r,shift left,"\varphi"] & \Ch(X;a',b')/\tilde{c}\ , \ar[l,shift left,"\psi"]
\end{tikzcd}
\]
which we will prove define arrows inverse to one another in Hot. Since $\Ch(X)/\tilde{c}$ is contractible (having a final object), hence aspheric, the same holds for $\cC(c;u,v)$, and we will win.

We define $\varphi$ by $\boxed{\varphi(c',w') = (\tilde{c}',\tilde{w}'),}$ where ($c'$ being of type $\tau'$) $\tilde{c}'$ is the path of type $\Delta_{+}\tau'\Delta_{-}$ obtained from $c'$ as inverse image under the degeneracy $\Delta_{+}\tau'\Delta_{-} \longrightarrow \tau'$, i.e. we extend $c'$ by a ``$\Delta_{+}$ step'' $a' \xrightarrow{\mathrm{id}} a'$ on the left and by a ``$\Delta_{-}$ step'' $b' \xleftarrow{\mathrm{id}} b'$ on the right, cf. figure.
\[
\overbrace{\underbrace{\begin{tikzcd}[ampersand replacement=\&]
a' \ar[r,"\mathrm{id}"] \ar[d,"\mathrm{id}"'] \ar[dr,phantom,"\ast"] \& a' \ar[rr,dotted,no head,"c'"] \ar[d,"u"'] \& {} \ar[d,dotted,"w'"] \& b' \ar[d,"v"'] \ar[dr,phantom,"\ast"] \& b' \ar[l,"\mathrm{id}"'] \ar[d,"\mathrm{id}"] \\
a' \ar[r,"u"] \& a \ar[rr,no head,"c"'] \& {} \& b \& b' \ar[l,"v"'] % [sic] a' ... b' with label c' in the bottom row -> a ... b with label c, the bottom row of the figure is the path c~ = ch(u) c ch(v)^o
\end{tikzcd}}_{\tilde{c}}}^{\tilde{c}'}
\]
\DERmark{ENV0546}\DERmark{ENV0547}\DERmark{ENV0548}


\msp{page 163}
We can say that
\[
\tilde{c}' = a'_{\Delta_{+}} \circ c' \circ b'_{\Delta_{-}},
\]
where in general, if $x \in X$ and $\tau \in \mathrm{Ty}$, $x_{\tau}$ denotes the ``constant path'' of value $x$, type $\tau$. On the other hand we have obvious morphisms
\[
a'_{\Delta_{+}} \longrightarrow \mathrm{ch}(u), \qquad b'_{\Delta_{-}} \longrightarrow \mathrm{ch}(v)^{\mathrm{opp}},
\]
as shown [?] by the vertical arrows $(\mathrm{id}_{a'},u)$ and $(v,\mathrm{id}_{b'})$ in the two squares marked $\ast$ in the diagram, arrows which lie one in $\Ch_{\Delta_{+}}(X)$, the other in $\Ch_{\Delta_{-}}(X)$. This said, we set
\[
\tilde{w}' : \tilde{c}' \longrightarrow \tilde{c}, \qquad \tilde{w}' = \lambda(u) \ast w' \ast \mu(v), % [sic] \tilde{w} -> \tilde{w}', the morphism just named is \tilde{w}'
\]
\DERmark{ENV0549}

where
\[
\left\{
\begin{aligned}
\lambda(u) :\ \ a'_{\Delta_{+}} &\quad\longrightarrow\quad \mathrm{ch}(u) \\
\mu(v) :\ \ b'_{\Delta_{-}} &\quad\longrightarrow\quad \mathrm{ch}(v)^{\mathrm{opp}}
\end{aligned}
\right.
\]
% foot: Version 22 décembre 2001
\printedpage{102}
% head: A. Grothendieck : Dérivateurs, Ch. VII, Édition des Universités de Montpellier II et Paris VII
(cf. figure). The pair
\[
(\tilde{c}',\tilde{w}') \in \Ob \Ch(X)/\tilde{c}
\]
thus defined depends functorially on $(c',w')$ varying in $\cC(c;u,v)$: If $(c'_1,w'_1)$ is a second object of this category, an arrow from the first to the second is an arrow
\[
c' \xrightarrow{\alpha} c'_1
\]
in $\Ch(X;a',b')$, such that we have commutativity in
\[
\begin{tikzcd}
c' \ar[rr,"\alpha"] \ar[dr,"w'"'] & & c'_1 \ar[dl,"w'_1"] \\
& c, &
\end{tikzcd}
\qquad\qquad w' = w'_1\alpha.
\]

Such an arrow $\alpha$ defines
\[
\tilde{\alpha} = a'_{\Delta_{+}} \ast \alpha \ast b'_{\Delta_{-}} : \tilde{c}' \longrightarrow \tilde{c}'_1
\]
in $\Ch(X;a',b')$, and we check
\msp{page 164}
at once that we have commutativity in
\[
\begin{tikzcd}
\tilde{c}' \ar[rr,"\tilde{\alpha}"] \ar[dr,"\tilde{w}'"'] & & \tilde{c}'_1 \ar[dl,"\tilde{w}'_1"] \\
& \tilde{c}, &
\end{tikzcd}
\qquad\qquad \text{i.e. } \tilde{w}' = \tilde{w}'_1\tilde{\alpha},
\]
i.e. that $\tilde{\alpha}$ is indeed an arrow in $\Ch(X;a',b')/\tilde{c}$. Likewise, that $\alpha \longmapsto \tilde{\alpha}$ defines a functorial dependence, i.e. $(\beta\alpha)\tilde{} = \tilde{\beta}\tilde{\alpha}$, is obvious from the definition of $\tilde{\alpha}$. So, we have indeed defined % Reading resolved against the supplied PDF; see corrections.csv.
\[
\cC(c;u,v) \ \xrightarrow{\ \varphi\ }\ \Ch(X;a',b')/\tilde{c},
\]
let us now define
\[
\cC(c;u,v) \ \xleftarrow{\ \psi\ }\ \Ch(X;a',b')/\tilde{c}.
\]
Let $(c'',w'')$ be an object of $\Ch(X;a',b')/\tilde{c}$,
\[
\overbrace{\underbrace{\begin{tikzcd}[ampersand replacement=\&]
a'_0 = a' \ar[r] \ar[d,"\mathrm{id}"'] \& a'_1 \ar[rr,no head,"\rightarrow"] \ar[d,"w''_1"'] \& {} \ar[d,dotted,"w''"] \& a'_n \ar[d,"w''_n"'] \& b' \ar[l] \ar[d,equal] \\
a' \ar[r,"u"] \& a \ar[rr,no head] \& {} \& b \& b', \ar[l,"v"']
\end{tikzcd}}_{\tilde{c}}}^{c''}
\]
\DERmark{ENV0550}\DERmark{ENV0551}

% foot: Version 22 décembre 2001

\printedpage{103}
% head: A. Grothendieck : Dérivateurs, Ch. VII, Édition des Universités de Montpellier II et Paris VII
thus
\[
w'' : c'' \longrightarrow \tilde{c}, \qquad\quad \text{arrow in } \Ch(X;a',b'),
\]
let us note that we have an obvious retraction arrow
\[
\tilde{c} \xrightarrow{\rho_c} c \qquad\quad \text{associated with } \Delta_{+}\tau\Delta_{-} \longrightarrow \tau \text{ (retraction),}
\]
inducing $u$ and $v$ on the source and the target:
\[
\begin{tikzcd}
a' \ar[r,"u"] \ar[d,"u"] & a \ar[rrr,no head,"c"] \ar[d,"\mathrm{id}"] & {} \ar[d,"\mathrm{id}"] & {} \ar[d,"\mathrm{id}"] & b \ar[d,"\mathrm{id}"] & b' \ar[l,"v"'] \ar[d,"v"] & \Big\}\ \tilde{c} & \in \Ch_{\Delta_{+}\tau\Delta_{-}}(a',b') \\
a \ar[r,"\mathrm{id}"] & a \ar[rrr,no head,"c"] & {} & {} & b & b \ar[l,"\mathrm{id}"'] & \Big\}\ \delta^{*}_{\tau}(c) & \in \Ch_{\Delta_{+}\tau\Delta_{-}}(a,b) \\
& a \ar[rrr,no head,"c"] & & & b & & \Big\}\ c & \in \Ch_{\tau}(a,b).
\end{tikzcd}
\]
\DERmark{ENV0552}\DERmark{ENV0553}


\msp{page 165}
So we define an object
\[
(c',w') = \psi(c'',w'')\ \in \Ob \cC(c;u,v)
\]
by
\[
c' = c'', \qquad\quad w' = \rho_c w'',
\]
and this is functorial in $(c'',w'')$.

The composite $\varphi\psi$ being an endomorphism of the contractible category $\Ch(X;a',b')/\tilde{c}$, is homotopic to the identity. We have to check that the same is true of $\psi\varphi$, an endomorphism of $\cC(c;u,v)$. This functor associates to $(c',w')$, $(\tilde{c}',\overline{w'})$, where $\tilde{c}'$ is defined as above, $\tilde{c}' = a'_{\Delta_{+}} \circ c' \circ b'_{\Delta_{-}}$, and where $\overline{w'}$ is the composite
\[
\underbrace{\tilde{c}'}_{= \delta^{\prime *}_{\tau'}(c')}\ \ \underbrace{\xrightarrow{\ \rho_{c'}\ }}_{\substack{\text{canonical}\\ \text{arrow}}}\ \ c' \xrightarrow{w'} c.
\]

So we have
\[
(\psi\varphi)(c',w') = (\tilde{c}',\overline{w'}) \xrightarrow{\rho_{c'}} (c',w'),
\]
this is indeed a morphism in $\cC(c;u,v)$, i.e. we have commutativity in
\[
\begin{tikzcd}
\tilde{c}' \ar[rr,"\rho_{c'}"] \ar[dr,"\overline{w'}"'] & & c' \ar[dl,"w'"] \\
& c &
\end{tikzcd}
\qquad\qquad \text{(arrow in } \Ch(X;a',b')\text{).}
\]
% foot: Version 22 décembre 2001


\endgroup
% END UNIT Der-ChVII-456/Der-ChVII-456_p101-105.tex
% BEGIN UNIT Der-ChVII-456/Der-ChVII-456_p106-109.tex
\begingroup\makeatletter%
\newcommand{\printedpage}[1]{\marginpar{\footnotesize p.~#1}}%
\newcommand{\msp}[1]{\marginpar{\footnotesize [#1]}}%
\newcommand{\CN}{\mathcal{C}_{\mathbf{N}^o}}%
\tikzset{zigzag/.style={-, postaction={decorate},%
  decoration={markings, mark=at position 0.12 with {\arrow{>}},%
  mark=at position 0.88 with {\arrow{<}}}}}%
\makeatother\let\printedpage\derpp\let\msp\dermsp%
\let\mspin\DERmspin%


\printedpage{104}
% head: A. Grothendieck : Dérivateurs, Ch. VII, Édition des Universités de Montpellier II et Paris VII
\msp{page 166}

Thus, the functor $\psi\varphi$ is indeed homotopic to the identity functor. We have won!

Let us pass to $1'$, $2'$) for $\CN$. The assertion $1'$) is proved in exactly the same way as 1) for \underline{Ch}$(X)$, by regarding $\CN$ as fibered over $\mathbf{N}^o$ (which has a final object 0), and using the diagram
\[
\begin{tikzcd}
X = (X_{\mathbf{N}^o})_e \arrow[rr, "i'"] & & X_{\mathbf{N}^o} = X \times \mathbf{N}^o \arrow[rr, "j'"] \arrow[dr] & & \mathcal{C}_{\mathbf{N}^o} \arrow[dl] \\
 & & & \mathbf{N}^o \ , &
\end{tikzcd}
\]
where $j'$ is (cartesian and) a Hot-equivalence on each fiber by the lemma, hence a Hot-equivalence globally, and $i'$ is a Hot-equivalence since $\mathbf{N}^o$ is aspheric (having a final object). Thus $j'i'$ is a Hot-equivalence, q.e.d.
\DERmark{ENV0554}


It remains to prove that
\[ \CN \longrightarrow X \times X \]
is proper and smooth. Here again,
\msp{page 167}
it is exactly the same argument as before, I presume. There is, however, a technical complication. We had taken, to define
\[ \underline{t} : \mathbf{N}^o \longrightarrow \underline{\mathrm{Ty}}, \qquad \underline{t} = (t_n)_n, \]
the choice
\[ t_n = (\Delta_+\Delta_-)^{2n}, \]
or if one prefers
\[ t_n = (\Delta_+\Delta_-)^n, \]
in any case the $t_n$ have initial direction $+$, terminal $-$, and they are reduced types. If one takes $\Delta_+t_n\Delta_-$, one finds a nonreduced type, which is not among the $t_n$. It would be of no use to change the definition of $t_n$ to $t_n = (\Delta_-\Delta_+)^{\text{something}}$, for this time $\Delta_+t_n\Delta_-$ will indeed be reduced, but will not have the right initial and terminal signs. Thus to define $\tilde{c}$ in terms of $c$, as an object of $\CN$, it will be necessary, if $c$ is of type $\tau = t_n$, to take $\tilde{c}$ of type $\Delta_+\Delta_-\tau\Delta_+\Delta_- = t_{n+1}$ (taking $t_{n+1} = (\Delta_+\Delta_-)^{2n+2}$), and likewise to define $\tilde{c}'$ in terms of $c'$. The relevant definitions of the passage from $(c',w')$ to $(\tilde{c}',\tilde{w}')$ are described in the figure: % [sic] t_{n+1} = (\Delta_+\Delta_-)^{2n} -> t_{n+1} = (\Delta_+\Delta_-)^{2n+2}, exponent inconsistent with t_n = (\Delta_+\Delta_-)^{2n}
\DERmark{ENV0555}


\msp{page 168}
% foot: Version 22 décembre 2001

\printedpage{105}
% head: A. Grothendieck : Dérivateurs, Ch. VII, Édition des Universités de Montpellier II et Paris VII
\[
\tilde{w}'\Big\downarrow
\left\{
\overbrace{\underbrace{
\begin{tikzcd}[ampersand replacement=\&]
a' \arrow[r, "\mathrm{id}"] \arrow[d, "\mathrm{id}"'] \& a' \arrow[d, "u"'] \& a' \arrow[l, "\mathrm{id}"'] \arrow[d, "u"'] \arrow[rrrr, dotted, no head, "c'", "\longrightarrow"'] \& {} \arrow[d, "w'"] \& {} \arrow[d, "w'"] \& {} \arrow[d, "w'"] \& b' \arrow[r, "\mathrm{id}"] \arrow[d, "v"'] \& b' \arrow[d, "v"] \& b' \arrow[l, "\mathrm{id}"'] \arrow[d, "\mathrm{id}"] \\
a' \arrow[r, "u"] \& a \& a \arrow[l, "\mathrm{id}"'] \arrow[rrrr, no head, "c", "\longrightarrow"'] \& {} \& {} \& {} \& b \arrow[r, "\mathrm{id}"] \& b \& b'. \arrow[l, "v"']
\end{tikzcd}
}_{\tilde{c}}}^{\tilde{c}'}
\right\}
\tilde{w}'
\]

Thus
\[
\left\{
\begin{array}{lcl ll}
\tilde{c} & = & \xi_u \circ c \circ \eta_v & \text{where} &
\left\{
\begin{array}{lccccc}
\xi_u : & a' & \xrightarrow{u} & a & \xleftarrow{\mathrm{id}} & a \\[1ex]
\eta_v : & b & \xrightarrow{\mathrm{id}} & b & \xleftarrow{v} & b'
\end{array}
\right\}
\text{ type } \Delta_+\Delta_- \\[4ex]
\tilde{c}' & = & a'_{\Delta_+\Delta_-} \circ c' \circ b'_{\Delta_+\Delta_-} & & \\[4ex]
\tilde{w}' & = & \lambda(u) * w' * \mu(v) & \text{with} & % [sic] w -> w', the homomorphism of the pair (c',w')
\left\{
\begin{array}{lccc}
\lambda(u) : & a'_{\Delta_+\Delta_-} & \longrightarrow & \xi_u \\[1ex]
\mu(v) : & b'_{\Delta_+\Delta_-} & \longrightarrow & \eta_v % [sic] \lambda(v) -> \mu(v), the formula above has \mu(v)
\end{array}
\right\}
\text{ as in \texttt{[the]} figure.}
\end{array}
\right.
\]
\DERmark{ENV0556}\DERmark{ENV0557}\DERmark{ENV0558}

One thus finds, mutatis mutandis,
\[ \CN(c;u,v) \xrightarrow{\varphi} \CN(a',b')/\tilde{c}, \qquad \varphi(c',w') = (\tilde{c}',\tilde{w}'), \]
and one defines the functor
\[ \CN(c;u,v) \xleftarrow{\psi} \CN(a',b')/\tilde{c} \]
as follows: to $(c'',w'')$ in $\CN(a',b')/\tilde{c}$, thus a path $c''$ of type $t_n = (\Delta_+\Delta_-)^{2n}$ from $a'$ to $b'$ and a homomorphism $w'' : c'' \longrightarrow \tilde{c}$ with fixed endpoints, one associates $(c',w')$ in $\CN(c;u,v)$, with
\[ c' = c'', \qquad w' = \rho_c w'', \]
where
\[ \rho_c : \tilde{c} \longrightarrow c \]
is defined in a similar way as before, by
\msp{page 169}
\[
\begin{tikzcd}
a' \arrow[r, "u"] \arrow[d, "u"'] & a \arrow[d, "\mathrm{id}"] & a \arrow[l, "\mathrm{id}"'] \arrow[d, "\mathrm{id}"] \arrow[rr, zigzag] & {} \arrow[d, "\mathrm{id}"'] & b \arrow[r, "\mathrm{id}"] \arrow[d, "\mathrm{id}"] & b \arrow[d, "\mathrm{id}"] & b' \arrow[l, "v"'] \arrow[d, "v"] & \tilde{c},\ \text{ type } t_{n+1}. \\
a \arrow[r, "\mathrm{id}"] & a & a \arrow[l, "\mathrm{id}"'] \arrow[rr, zigzag, "c"] & {} & b \arrow[r, "\mathrm{id}"] & b & b \arrow[l, "\mathrm{id}"'] & \varphi^*_{n,n+1}(c).
\end{tikzcd}
\]
\DERmark{ENV0559}


One need only check that the endofunctor $\psi\varphi$ of $\CN(c;u,v)$ is homotopic to the identity. But this functor transforms $(c',w')$ into $(\tilde{c}' = a'_{\Delta_+\Delta_-}c'b'_{\Delta_+\Delta_-}, \overline{w'})$, where $\overline{w'}$ is again the composite
% foot: Version 22 décembre 2001
\printedpage{106}
% head: A. Grothendieck : Dérivateurs, Ch. VII, Édition des Universités de Montpellier II et Paris VII
\[
\begin{tikzcd}
\tilde{c}' \arrow[r, "\text{can.}"] \arrow[rr, bend right=20, "\overline{w'}"'] & c' \arrow[r, "w'"] & c.
\end{tikzcd}
\]

We therefore have
\[ \psi\varphi \longrightarrow \mathrm{id}, \]
$\psi\varphi$ is homotopic to the identity functor.

This proves the smoothness of $\CN \longrightarrow X \times X$. For properness, one must prove the smoothness in the case of $X^o$, and of the inverse system $\underline{t}^{\mathrm{opp}} \simeq (t_n^{\mathrm{opp}})_n$. But we have $\underline{t}^{\mathrm{opp}} = \underline{t}$, and it suffices to apply what we have already proved to the case of $X^o$, to see that $\CN \longrightarrow X \times X$ is proper. % [sic] de le cas -> dans le cas, grammatical
\DERmark{ENV0560}


Only the lemma remains to be proved. It must be essentially
\msp{page 170}
trivial, and I give up writing out a proof.

I am now going to show ``formally'' that the ``canonical'' factorization of an arrow in Cat $f : X' \longrightarrow X$ as
\[ X' \xrightarrow{i_f} \overline{X'} \xrightarrow{\bar{f}} X, \qquad i_f \text{ Hot-equivalence,} \]
with (we hoped) $\bar{f}$ proper and smooth, starting from a similar factorization for the diagonal arrow $X \xrightarrow{\ \ \mathrm{diag}_X\ \ } X \times X$ as
\[
\begin{tikzcd}
X \arrow[d, "i_X\ \ \text{\footnotesize Hot-equivalence}"] \\
\mathcal{C}_X\ \ \text{``paths''} \arrow[d, "p_X"] \\
X \times X
\end{tikzcd}
\]
is a formal matter. (Cf. page 139 for the notations and \texttt{[the]} problematics.) We look at the diagram
\[
\begin{tikzcd}
 & \mathcal{C}_X \arrow[d, "p_X"] \arrow[dl, "\beta_X"'] \arrow[dr, phantom, "\text{cart.}"] & C(f) \arrow[l, "\mathrm{pr}'_2"'] \arrow[d, "p_{X,f}"] \arrow[dll, bend right=50, "\bar{f}"'] & [-1.5em] \simeq X' \times_X (C_X, \sigma = \mathrm{pr}_1 \circ p_X) \\ % Reading resolved against the supplied PDF; see corrections.csv.
X & X \times X \arrow[l, "\mathrm{pr}_2^X"] \arrow[d, "\mathrm{pr}_1^X"'] \arrow[dr, phantom, "\text{cart.}"] & X' \times X \arrow[l, "f \times \mathrm{id}_X = f'"] \arrow[d, "\mathrm{pr}_1"] & \\
 & X & X' \arrow[l, "f"'] \arrow[uu, bend right=60, dotted, "i_f"'] &
\end{tikzcd}
\]
% foot: Version 22 décembre 2001
\printedpage{107}
% head: A. Grothendieck : Dérivateurs, Ch. VII, Édition des Universités de Montpellier II et Paris VII
(\footnote[40]{NB $C(f)$ is defined as the inverse image over $X'$ of the bundle $C_X$ over $X$ (by $\sigma_X = \mathrm{pr}_1^X \circ p_X$) by $f : X' \longrightarrow X$.}). Since $p_X$ is proper and smooth, i.e. ``perfect'', its inverse image over $X' \times X$, namely $p_{X,f}$, is so too; on the other hand $\mathrm{pr}_2^X \circ f'$ is the
\msp{page 171}
projection $\mathrm{pr}_2$ of $X' \times X$ onto $X$, hence it is also a perfect morphism (since $X' \longrightarrow e$ is perfect, hence the same holds for the composite
\[ (\mathrm{pr}_2^X \underbrace{f') \circ p_{X,f}} = \underbrace{\mathrm{pr}_2^X(p_X}\, \mathrm{pr}'_2) = \beta_X \mathrm{pr}'_2 \overset{\mathrm{def}}{=} \bar{f} \ \ ). \]
\DERmark{ENV0561}


I recall that the section $i_f$ of $C(f)$ over $X'$ is defined from the section $i_X$ of $(\mathcal{C}_X,\sigma_X)$ over $X$. It is a Hot-equivalence, for $\sigma_X = \mathrm{pr}_1^X p_X$ is perfect with aspheric fibers, hence universally a Hot-equivalence, hence $\mathrm{pr}_1 p_{X,f} : C(f) \longrightarrow X'$ is also a Hot-equivalence, and any section is therefore one too.
% foot: Version 22 décembre 2001


\endgroup
% END UNIT Der-ChVII-456/Der-ChVII-456_p106-109.tex
\clearpage\setcounter{footnote}{0}\phantomsection\addcontentsline{toc}{section}{Chapter VIII}\markboth{Chapter VIII}{Chapter VIII}
% BEGIN UNIT Der-ChVIII-567/Der-ChVIII-567_p001-005.tex
\begingroup\makeatletter%
\newcommand{\printedpage}[1]{\marginpar{\footnotesize p.~#1}}%
\newcommand{\msp}[1]{\marginpar{\footnotesize [#1]}}%
\newcommand{\Grp}{\underline{\mathrm{Grp}}}%
\makeatother\let\printedpage\derpp\let\msp\dermsp%
\let\mspin\DERmspin%


% ===================== page 1 =====================
\printedpage{}
% head: (aucun ; page de titre non numérotée)

\begin{center}
{\Large\textsf{THE DERIVATORS}}

\bigskip

\textbf{ALEXANDRE GROTHENDIECK}

\bigskip\bigskip

\textbf{Chapter VIII}

\medskip

\textbf{Relative homotopy 1-types: their integration \dots}
\end{center}

\par

% ===================== page 2 =====================
\printedpage{}
% head: (aucun ; page non numérotée)

\par\noindent\textbf{Note by the editors of the transcribed edition.}\par
This text was deciphered and transcribed in \LaTeXe{} by M. Künzer. It was edited by M. Künzer, J. Malgoire, and G. Maltsiniotis. The transcription is as faithful as possible to the manuscript. For the few obvious corrections, or rare comments by the editors, as well as for the original numbering of the manuscript pages, typewriter characters \texttt{[typewriter]} in brackets are used. A question mark in brackets means that one is not sure of the preceding word.

This edition is provisional. Remarks, comments and corrections are welcome. Send a message to:

\textbf{[e-mail address removed]}

\hfill G. Maltsiniotis\hspace*{2cm}

\par

% ===================== page 3 =====================
\printedpage{1}
% head: A. Grothendieck : Dérivateurs, Ch. VIII, Édition des Universités de Montpellier II et Paris VII
% foot: Version 2 février 2002

\msp{page 1}

\section*{1\quad Integration of relative homotopy $1$-types over a category.}

Relative homotopy types over a (small) category $S$ are defined by \textit{stacks} (in categories) on $S^{\wedge}$, which correspond essentially to contrafunctors $S^o \longrightarrow (\mathrm{Cat})$, or again to categories $X$ over $S$ which are Cat-fibered over $S$. (In $(\mathrm{Cat}/S)_{\mathrm{fib}}$, one inverts the arrows which are Hot-equivalences on the fibers.) On the other hand, relative homotopy 1-types over $S$ should be described by \textit{stacks in groupoids} on $S$. If one has a relative homotopy type, defined by
\[ X \longrightarrow S, \qquad\qquad \text{Cat-fibered category over } S, \]
or by
\[
\begin{aligned}
S^o \quad &\longrightarrow \quad \mathrm{Cat}\\
s \quad &\longmapsto \quad X_s\ ,
\end{aligned}
\]
the associated homotopy 1-type is defined by
\[ s \longmapsto \Pi(X_s), \qquad\qquad S^o \longrightarrow (\mathrm{Grp}), \qquad\qquad \text{category of groupoids,} \]
where $\Pi(X_s)$ is the fundamental groupoid of $X_s$. We denote the corresponding fibered category over $S$ by \fbox{$\Pi(X/S)$}. It is true that if
\[ X \ \longrightarrow\ X' \]
is a homomorphism of fibered categories over $S$ which induces a Hot-equivalence on the fibers, it induces a fibered equivalence
\[ \Pi(X/S) \ \longrightarrow\ \Pi(X'/S), \]
but not an isomorphism. There would thus be work to do to define a functor
\[
\underbrace{\mathrm{Hot}(S)}_{\text{\footnotesize\begin{tabular}{c}category of\\ relative homotopy\\ types over $S$\end{tabular}}}
\quad\xrightarrow{\ \xi\,\longmapsto\,\Pi(\xi/S)\ }\quad
\underbrace{\Grp(S)}_{\text{\footnotesize\begin{tabular}{c}category of\\ groupoids\\ fibered over $S$\\ (or categories\\ of stacks in\\ groupoids\\ on $S^{\wedge}$)\end{tabular}}}
\quad .
\]

But this is not my purpose for the moment. Rather, to give a transitivity formula in noncommutative homology, analogous to $R(gf)_* \simeq Rg_*Rf_*$ in cohomology.

\par

% ===================== page 4 =====================
\printedpage{2}
% head: A. Grothendieck : Dérivateurs, Ch. VIII, Édition des Universités de Montpellier II et Paris VII
% foot: Version 2 février 2002

\msp{page 2}

If
\[ f : S \longrightarrow T \]
is a functor, we are going to define a functor
\[ f_! :\ \Grp(S) \ \longrightarrow\ \Grp(T) \]
(direct image of relative homotopy 1-types). But let us first note that if
\[ p :\ X \ \longrightarrow\ S \]
is any $S$-category \textit{whatsoever} (not necessarily Cat-fibered), one nevertheless associates to it a relative homotopy type, via the fibered category $\tilde{X}^{(S)}$, or $\tilde{X}$, having as fibers the $s\backslash X$ (noting that $s' \longrightarrow s$ in $S$ does indeed define $s\backslash X \longrightarrow s'\backslash X$) (\footnote{Having $\tilde{X}$, one defines $\Pi(X/S) \overset{\mathrm{def}}{=} \Pi(\tilde{X}/S)$.}). We have a canonical $S$-functor
\DERmark{ENV0562}

\[ X \xrightarrow{\ i\ } \tilde{X}, \qquad\qquad x \longmapsto (x, p(x), \underbrace{p(x)}_{=\,\text{‘}\,s\,\text{’}} \xrightarrow{\ \mathrm{id}\ } p(x)). \]
This is an $S$-Hot-equivalence if $X$ was already fibered. In fact, on each fiber it is a homotopism, for one then defines $\tilde{X}_{s_0} \xrightarrow{\ j\ } X_{s_0}$ in an obvious way, in such a way that
\[
X_{s_0} \underset{j}{\overset{i}{\rightleftarrows}} \tilde{X}_{s_0}, \qquad ji \simeq \mathrm{id}, \qquad ij \underbrace{\xrightarrow{\text{simple homotopy}}}_{\text{\footnotesize\begin{tabular}{l}NB this arrow\\ is not \texttt{[an]}\\ isomorphism\\ in general.\end{tabular}}} \mathrm{id}.
\]
Let us then return to $f : S \longrightarrow T$, and consider a fibered groupoid over $S$,
\[ G \ \longrightarrow\ S, \]
this is not a fibered category over $T$, but one can associate to it
\[ t \longmapsto t\backslash G, \qquad\qquad T^o \longrightarrow (\mathrm{Cat}), \]
hence a fibered category over $T$, of which one takes the relative $\Pi$. But it is simpler to take directly
\[ t \ \longmapsto\ \Pi(t\backslash G), \]
which defines a category fibered in groupoids over $T$.

\par

% ===================== page 5 =====================
\printedpage{3}
% head: A. Grothendieck : Dérivateurs, Ch. VIII, Édition des Universités de Montpellier II et Paris VII
% foot: Version 2 février 2002

\msp{page 3}

Thus we do indeed find a functor (in fact, even a 2-functor of 2-categories)
\[ f_! :\ \Grp(S) \ \longrightarrow\ \Grp(T). \]
If
\[ u :\ G \ \longrightarrow\ G' \]
is a homomorphism of fibered groupoids over $S$, which induces Cat-equivalences on the fibers, then
\[ f_!(u) :\ f_!(G) \ \longrightarrow\ f_!(G') \]
is also a Cat-equivalence on the fibers. (This is formal, from the fact that $f_!$ is a 2-functor, taking into account that if one has an equivalence on the fibers, then there is a quasi-inverse -- this is a general story about fibered categories \dots)

Now let
\[ p :\ X \ \longrightarrow\ S, \]
and let us take the case $G = \Pi(X/S)$. I say that we then have a $T$-equivalence
\[ \boxed{\ f_!(\Pi(X/S)) \ \simeq\ \Pi(X/T),\ } \tag*{$(*)$} \]
in other words, to know $\Pi(X/T)$ (up to $T$-equivalence), it suffices to know $\Pi(X/S)$, and to apply to it the functor $f_!$ of integration of relative homotopy 1-types, for $f : S \longrightarrow T$.

Let us note that we trivially have
\[ \Pi(X/S) \ =\ p_!(e_X), \]
where $e_X$ denotes the trivial fibered $X$-groupoid (all fibers one-point). Thus the preceding formula is written as a transitivity formula
\[ f_!(p_!(\Gamma)) \ \simeq\ (fp)_!(\Gamma), \]
\msp{page 4}
where $\Gamma$ is the trivial fibered groupoid over $X$. This suggests generalizing the formula $(*)$ to
\[ \boxed{\ (fp)_!(\Gamma) \ \simeq\ f_!(p_!(\Gamma))\ } \]
for \textit{every} fibered groupoid over $X$, in a situation
\[ X \xrightarrow{\ p\ } S \xrightarrow{\ f\ } T. \]

This will be formal, if we establish a 2-adjunction property for the functor $f_!$:


\endgroup
% END UNIT Der-ChVIII-567/Der-ChVIII-567_p001-005.tex
% BEGIN UNIT Der-ChVIII-567/Der-ChVIII-567_p006-010.tex
\begingroup\makeatletter%
\newcommand{\printedpage}[1]{\marginpar{\footnotesize p.~#1}}%
\newcommand{\msp}[1]{\marginpar{\footnotesize [#1]}}%
\newcommand{\uHom}{\underline{\mathrm{Hom}}}%
\newcommand{\uGrp}{\underline{\mathrm{Grp}}}%
\newcommand{\uuGrp}{\underline{\underline{\mathrm{Grp}}}}%
\newcommand{\fl}{\longrightarrow}%
\makeatother\let\printedpage\derpp\let\msp\dermsp%
\let\mspin\DERmspin%


\printedpage{4}
% head: A. Grothendieck : Dérivateurs, Ch. VIII, Édition des Universités de Montpellier II et Paris VII
% foot: Version 2 février 2002

\textbf{Theorem.} \textit{Let $f : S \fl T$ be an arrow in} (Cat)\textit{, $G$ a stack in groupoids over $S$, $H$ a stack in groupoids over $T$. We then have a canonical equivalence of categories}
\[
\boxed{\ \uHom_{\uGrp(S)}(G, f^*H) \ \simeq\ \uHom_{\uGrp(T)}(f_!G, H).\ }
\]

The adjunction arrows
\[
\left\{
\begin{array}{lclcrcl}
f_!f^*H & \fl & H & \qquad\qquad & H & \in & \mathrm{Ob}\,\uGrp(T)\\[4pt]
f^*f_!G & \longleftarrow & G & \qquad\qquad & G & \in & \mathrm{Ob}\,\uGrp(S)
\end{array}
\right.
\]
should be definable without any problem. For the second, we have to define for $s \in S$
\[
G_s \ \fl\ (f^*(f_!(G)))_s = (f_!(G))_{f(s)} \overset{\mathrm{def}}{=} \Pi(f(s)\backslash_T G),
\]
\DERmark{ENV0563}\DERmark{ENV0564}

now we have
\[
G_s \fl f(s)\backslash_T G,
\]
\DERmark{ENV0565}

hence
\[
\underbrace{\Pi(G_s)}_{\begin{subarray}{l} =\ G_s\\ \text{since } G_s\\ \text{is a}\\ \text{groupoid}\end{subarray}} \ \fl \Pi(f(s)\backslash_T G),
\]
\DERmark{ENV0566}

O.K. This is functorial in $s$ \dots

\msp{page 5}
For the arrow $f_!f^*H \fl H$ over $T$, we have to find for every $t \in \mathrm{Ob}\,T$, and functorially in $t$,
\[
\underbrace{(f_!f^*H)_t}_{\overset{\mathrm{def}}{=}\Pi(t\backslash f^*(H))} \quad \fl \quad \underbrace{H_t}_{=\,\Pi(H_t)} \ ,
\]
\DERmark{ENV0567}

it would therefore suffice to define
\[
t\backslash f^*(H) \ \fl\ H_t,
\]
but we have the cartesian diagram
\[
\begin{tikzcd}[column sep=large, row sep=large]
f^*H = H \times_T S \arrow[r] \arrow[d] & H \arrow[d] \\
S \arrow[r, "f"] & T \ni t
\end{tikzcd}
\]
\DERmark{ENV0568}

\texttt{[rather} $H \times_T S$\texttt{]} inducing
\[
f^*H \ \fl\ H,
\]
\printedpage{5}
% head: A. Grothendieck : Dérivateurs, Ch. VIII, Édition des Universités de Montpellier II et Paris VII
% foot: Version 2 février 2002
hence
\[
t\backslash f^*H \ \fl\ t\backslash H \qquad\quad \text{(instead of mapping into $H_t$),}
\]
but on the other hand, since $H$ is fibered over $T$, hence smooth over $T$, the inclusion
\[
H_t \fl t\backslash H
\]
is a Hot-equivalence. Or rather, let us say: $f^*H \fl H$ is a morphism of categories over $T$, inducing
\[
\underbrace{\Pi(f^*H/T)}_{\overset{\mathrm{def}}{=}\ f_!(f^*H)} \ \fl\ \Pi(H/T),
\]
\DERmark{ENV0569}


\msp{page 6}
now the canonical $T$-functor
\[
H \fl \Pi(H/T) \qquad\quad (= \text{category over } T \text{ defined by } t \longmapsto \Pi(t\backslash H))
\]
is an equivalence, because $H$ is a groupoid fibration over $T$. So we can find a $T$-quasi-inverse
\[
\Pi(H/T) \fl H
\]
and by composing we find
\[
f_!(f^*H) \ \fl\ \Pi(H/T) \ \fl\ H
\]
the arrow sought.

Having the adjunction 2-arrows, I stop there with the verifications!

\printedpage{6}
% head: A. Grothendieck : Dérivateurs, Ch. VIII, Édition des Universités de Montpellier II et Paris VII
% foot: Version 2 février 2002

\section*{2\quad Base change questions}

\[
\begin{tikzcd}[column sep=large, row sep=large]
S & S' \arrow[l, "q"'] \arrow[d, "f'"] \\
T \arrow[from=1-1, "f"'] & T' \arrow[l, "p"]
\arrow[from=2-2, to=1-1, bend right=50, shorten <=18pt, shorten >=18pt, "\gamma"']
\end{tikzcd}
\qquad\qquad\qquad
p^*f_!(G) \longleftarrow f'_!q^*(G)
\tag*{$(*)$}
\]

The arrow is defined, provided we have commutativity or a commutativity arrow $\gamma : pf' \fl fq$ (not necessarily \texttt{[an]} isomorphism). By adjunction, it amounts to the same thing to give
\[
\begin{tikzcd}[row sep=large]
p_!f'_!q^*G \arrow[r] \arrow[d, dash, "\wr"', "\text{transitivity}"] & f_!(G) \\
(pf')_!q^*G \arrow[d, "\gamma"] & \\
(fq)_!q^*G \arrow[d, dash, "\text{transitivity}\ \wr"'] & \\
f_!(q_!q^*)G. \arrow[uuur, dashed, "\text{since } q_!q^*G \longrightarrow G"'] &
\end{tikzcd}
\]

We wonder when the arrow is \texttt{[an]} isomorphism.

\msp{page 7}
The following theorem must be true, and follows from my general \textit{fourbis}.
\DERmark{ENV0570}


% vertical bar printed in the left margin beside the following theorem
\textbf{Theorem.} \textit{If the square $(*)$ is} cartesian \textit{in} (Cat)\textit{, $\gamma$ an isomorphism, and if we assume that $f$ is} 1\textit{-smooth, or $p$} 1\textit{-proper, then the $T'$-functor of base change is an equivalence of $T'$-categories.}

The general ``\textit{fourbis}'' would be in terms of an ad hoc derivator $\mathbf{D}$, that of relative homotopy 1-types. I will return to this later.

We ask the question when the 2-functor associated to $f : S \fl T$
\[
f^* :\ \uuGrp(T) \ \fl\ \uuGrp(S)
\]
\printedpage{7}
% head: A. Grothendieck : Dérivateurs, Ch. VIII, Édition des Universités de Montpellier II et Paris VII
% foot: Version 2 février 2002
is 0-faithful, 1-faithful, 2-faithful (\footnote[2]{The $i$-faithfulness of $f^*$ ($i \in \{0,1,2\}$) is equivalent to that of the functors (obtained by taking $G = e_T$)
\[
\underline{\Gamma}(H/T) \ \fl\ \underline{\Gamma}(f^*H/S).
\]}). If $G$, $H$ are groupoid fibrations over $T$, we consider
\[
\uHom_{\uuGrp(T)}(G,H) \ \fl\ \uHom_{\uuGrp(S)}(f^*G, f^*H),
\]
and according as this functor is always faithful, or fully faithful, or an equivalence of categories, we say that $f^*$ is 0-, 1-, or 2-faithful. In the 2-faithful case, this implies that if $K$ is a groupoid fibration over $S$, if it is equivalent to an $f^*G$, then $G$ is determined up to $T$-equivalence, itself determined up to unique isomorphism.

\msp{page 8}
Let $i \in \mathbf{N}$. We say that $f$ is $i$-aspheric if $\forall t \in T$, $Z = S/t$ is $i$-aspheric, i.e.
\[
\pi_0(Z),\ \dots,\ \pi_i(Z) \qquad\quad \text{trivial.}
\tag*{$(*)$}
\]

We also define $(-1)$-aspheric: the $S/t$ are nonempty. Likewise for the notion of $i$-coaspheric.

In fact, for $i \in [-1,+\infty[ \ \subseteq \mathbf{Z}$, we define a derivator $\mathbf{D}_i$ of relative homotopy $i$-types (cf. below). Hence the notions of a $\mathbf{D}_i$-aspheric object ($=$ $i$-aspheric), of a $\mathbf{D}_i$-equivalence, of a $\mathbf{D}_i$-aspheric and $\mathbf{D}_i$-coaspheric functor, of a $\mathbf{D}_i$-smooth and $\mathbf{D}_i$-proper functor, of a $\mathbf{D}_i$-fibration and of a $\mathbf{D}_i$-fibration that is \textit{trivial} (i.e. with $\mathbf{D}_i$-aspheric fibers). We should have the

\textbf{Theorem 3.} \textit{Let $f : S \fl T$ be in} (Cat)\textit{, $i \in \{0,1,2\}$. Suppose one of the following conditions is satisfied:}

\begin{itemize}[leftmargin=3em]
\item[\textit{a)}] \textit{$f$ is $(i-1)$-proper, and $(i-1)$-aspheric (is expressed by the condition that certain categories, such as the $S/t$, be $(i-1)$-aspheric).}
\item[\textit{b)}] \textit{$f$ is $(i-1)$-smooth, and $(i-1)$-coaspheric.}
\end{itemize}

\textit{Then:}

\begin{itemize}[leftmargin=3em]
\item[\textit{1\textsuperscript{o})}] \textit{The inverse image functor $f^*$}
\[
f^* : \uGrp(T) \ \fl\ \uGrp(S)
\]
\textit{is $i$-faithful.}
\item[\textit{2\textsuperscript{o})}] \textit{In the case $i = 2$ ($f^*$ 2-faithful), we}
\printedpage{8}
% head: A. Grothendieck : Dérivateurs, Ch. VIII, Édition des Universités de Montpellier II et Paris VII
% foot: Version 2 février 2002
\msp{page 9}
\textit{also have this: if $G \in \mathrm{Ob}\,\uGrp(S)$ is a groupoid fibration over $S$, then $G$ is equivalent to a groupoid fibration of the form $f^*(H)$ if and only if for every $t \in T$, the restriction $G_t$ of $G$ to $S_t$ is equivalent to a constant groupoid fibration over $S_t$. (}\textbf{NB} \textit{Condition a) or b) already implies that $S_t$ is} 1\textit{-connected, i.e. that its fundamental groupoid is equivalent to $e$.)}
\DERmark{ENV0571}

\item[\textit{3\textsuperscript{o})}] \textit{Let $Z$ be an object} 1\textit{-connected in} (Cat)\textit{, $G$ a groupoid fibration over $Z$. For $G$ to be $Z$-equivalent to a constant groupoid over $Z$, it is necessary and sufficient that for every arrow $u$ in $Z$, the morphism of groupoids $u^*$ induced on the fibers of $G$ be an equivalence of groupoids} (\footnote[3]{We then say that $G$ is \textit{essentially locally constant} over $Z$.}).
\end{itemize}
\DERmark{ENV0572}


\textbf{Corollary.} \textit{The necessary and sufficient condition of 2\textsuperscript{o}) can also be written as follows: for every arrow $u$ of $S$ ``vertical'' over $T$, i.e. such that $f(u)$ is an identity arrow, the functor $u^*$ induced by $u$ on the fibers of $G$ be an equivalence of groupoids.}


\endgroup
% END UNIT Der-ChVIII-567/Der-ChVIII-567_p006-010.tex
\clearpage\setcounter{footnote}{0}\phantomsection\addcontentsline{toc}{section}{Chapter IX}\markboth{Chapter IX}{Chapter IX}
% BEGIN UNIT Der-ChIX-678/Der-ChIX-678_p001-005.tex
\begingroup\makeatletter%
\newcommand{\printedpage}[1]{\marginpar{\footnotesize p.~#1}}%
\newcommand{\msp}[1]{\marginpar{\footnotesize [#1]}}%
\newcommand{\Ch}{\mathrm{Ch}}%
\newcommand{\Hot}{\mathrm{Hot}}%
\newcommand{\Hom}{\mathrm{Hom}}%
\makeatother\let\printedpage\derpp\let\msp\dermsp%
\let\mspin\DERmspin%


\printedpage{[unnumbered]}
% head: (aucun titre courant imprimé)

\begin{center}
{\LARGE\textsf{THE DERIVATORS}}

\bigskip

\textbf{ALEXANDRE GROTHENDIECK}

\bigskip\bigskip

\textbf{Chapter IX}

\textbf{Return to the categories of fractions $\boldsymbol{\mathcal{M}W^{-1}}$,\\
comparison with Quillen}
\end{center}

\par

\printedpage{[unnumbered]}
% head: (aucun titre courant imprimé)

\par\noindent\textbf{Note by the editors of the transcribed edition.}\par
This text was deciphered and transcribed in \LaTeXe{} by M. Künzer. It was edited by M. Künzer, J. Malgoire, and G. Maltsiniotis. The transcription is as faithful as possible to the manuscript. For the few obvious corrections, or rare comments by the editors, as well as for the original numbering of the manuscript pages, typewriter characters \texttt{[typewriter]} in square brackets are used. A question mark in square brackets means that one is not sure of the preceding word.

This edition is provisional. Remarks, comments and corrections are welcome. Send a message to:

\textbf{[e-mail address removed]}

\hfill G. Maltsiniotis\hspace*{2cm}

\par

\printedpage{1}
% head: A. Grothendieck : Dérivateurs, Ch. IX, Édition des Universités de Montpellier II et Paris VII
% foot: Version 12 mars 2002

\msp{page 1}

\section*{1\quad Functoriality of $\Ch^{\Sigma}_{\infty}(a,b)$ in $a$, $b$ in $\mathcal{M}\Sigma^{-1}$}

We are given a category $\mathcal{M}$ with \texttt{[an]} admissible set of arrows $\Sigma \subseteq \mathrm{Fl}(\mathcal{M})$, whence (taking e.g. $\underline{t} = (t_n)$, with $t_n = (\underbrace{\Delta_+\Delta_-}_{\tau_0})^{2n}$, and $\underbrace{t_{n+1}}_{=\,\tau_0^{2n+2}} \longrightarrow \underbrace{t_n}_{=\,\tau_0^{2n}}$ defined as $p * \mathrm{id}_{\tau_0^{2n}} * p$, where $p : \tau_0 \longrightarrow 1$ the degeneracy.) We then have a theory of the $\underline{\Ch}_{\Sigma,\infty}(\mathcal{M},a,b)$, which I denote simply $\Ch_\infty(a,b)$ as an object of Hot. Pairings
\[
\Ch_\infty(a,b) \times \Ch_\infty(b,c) \longrightarrow \Ch_\infty(a,c) \qquad \text{in Hot}
\]
with (we shall admit it) \textit{associativity} in Hot. This therefore defines pairings on the $\pi_0(\Ch_\infty(a,b)) \overset{\mathrm{def}}{=} \pi_0(a,b)$, which is none other than the composition of arrows in $\mathcal{M}\Sigma^{-1}$. On the other hand, we have for $a,b \in \mathrm{Ob}\,\mathcal{M}$
\DERmark{ENV0198}

\[
\Hom_{\mathcal{M}}(a,b) \longrightarrow \pi_0(\Ch_\infty(a,b)) = \pi_0(a,b),
\]
compatible with compositions. But we can ignore it for the moment, and look at $\mathrm{Ob}\,\mathcal{M}$, or simply $\mathcal{M}_0$, with only the structure
\[
\left\{
\begin{array}{ccccccc}
\mathcal{M}_0 \times \mathcal{M}_0 & \longrightarrow & \Hot & \qquad & (a,b) & \longmapsto & \Ch_\infty(a,b) \\[1ex]
\mathcal{M}_0 \times \mathcal{M}_0 \times \mathcal{M}_0 & \longrightarrow & \mathrm{Fl}\,\Hot & \qquad & (a,b,c) & \longmapsto & (\Ch_\infty(a,b) \times \Ch_\infty(b,c) \underbrace{\longrightarrow}_{\substack{\text{arrow}\\ \text{in Hot}}} \Ch_\infty(a,c))
\end{array}
\right.
\]
with condition \texttt{[of]} associativity, and also a condition on units:
\[
\Ch_\infty(a,a) \text{ is equipped with a } 1_a \in \pi_0\Ch_\infty(a,a) = \pi_0(a,a),
\]
playing the role of a ``two-sided identity'', in the following sense. Let us note that the pairings
\[
\begin{array}{ccccc}
\Ch_\infty(a,b) & \times & \Ch_\infty(b,c) & \longrightarrow & \Ch_\infty(a,c) \\[1ex]
P & \times & Q & \longrightarrow & R
\end{array}
\]
\msp{page 2}
define
\[
\begin{array}{ccc}
P & \longrightarrow & \underline{\Hom}_{\Hot}(Q,R) \\[1ex]
Q & \longrightarrow & \underline{\Hom}_{\Hot}(P,R),
\end{array}
\]
whence on passing to the $\pi_0$
\[
\begin{array}{ccc}
\pi_0(P) & \longrightarrow & \Hom_{\Hot}(Q,R) \\[1ex]
\pi_0(Q) & \longrightarrow & \Hom_{\Hot}(P,R),
\end{array}
\]
\printedpage{2}
% head: A. Grothendieck : Dérivateurs, Ch. IX, Édition des Universités de Montpellier II et Paris VII
% foot: Version 12 mars 2002
in other words,
\[
\begin{array}{ccc}
\pi_0(a,b) & \longrightarrow & \Hom_{\Hot}(\Ch_\infty(b,c),\Ch_\infty(a,c)) \\[1ex]
\pi_0(b,c) & \longrightarrow & \Hom_{\Hot}(\Ch_\infty(a,b),\Ch_\infty(a,c)).
\end{array}
\]

We deduce from this in particular for $a,b \in \mathcal{M}_0$
\[
\left\{
\begin{array}{cccc}
1) & \pi_0(a,a) & \longrightarrow & \Hom_{\Hot}(\Ch_\infty(a,b),\Ch_\infty(a,b)) \\[1ex]
2) & \pi_0(a,a) & \longrightarrow & \Hom_{\Hot}(\Ch_\infty(b,a),\Ch_\infty(b,a)),
\end{array}
\right.
\tag*{$(*)$}
\]
and an object $\alpha_a$ of $\pi_0(a,a)$ is a two-sided unit if its two images are the identity homomorphisms in $\Ch_\infty(a,b)$ resp. $\Ch_\infty(b,a)$. Passing to the $\pi_0$ of these, this implies that $\alpha_a$ is a two-sided unit for the composition of the $\pi_0(x,y)$, relative to $b$. If we want this to be the case for \textit{every $b$,} this means (taking $b = a$) that $\alpha_a$ is a two-sided unit object for the structure of associative composition (hence of a monoid) in $\pi_0(a,a)$, hence that $\pi_0(a,a)$ is a monoid \textit{with unit,} having $\alpha_a$ as unit object. The element $\alpha_a$ is therefore uniquely determined. And we want $\alpha_a$ to exist
\msp{page 3}
for every $a \in \mathcal{M}_0$, so $\mathcal{M}_0$, equipped with the $\pi_0(a,b)$ and their composition, becomes a \textit{category,} say $\mathcal{M}_\Pi$; and its identity arrows must play the role of a two-sided unit also for the ``external Hom'' $\Ch_\infty(a,b)$, with values in Hot.

We can say that for fixed $a = a_0$, the $\Ch_\infty(a_0,b)$ ($b$ variable) form the values of a covariant \textit{functor}
\[
\mathcal{M}_\Pi \longrightarrow \Hot \qquad b \longmapsto \Ch_\infty(a_0,b),
\]
by associating to
\[
b \xrightarrow{\ u\ } b', \qquad u \in \pi_0(b,b'),
\]
the corresponding arrow
\[
\Ch_\infty(a_0,b) \longrightarrow \Ch_\infty(a_0,b') % [sic] \Ch_\infty(b') -> \Ch_\infty(a_0,b'), argument manquant (cf. la note des éditeurs qui suit)
\]
\DERmark{ENV0199}

\texttt{[rather} $\Ch_\infty(a_0,b) \longrightarrow \Ch_\infty(a_0,b')$\texttt{]} $*_1$ [?]. That this is indeed a functorial law should follow formally from the associativity of the pairings on the $\Ch_\infty$. Likewise, we have a functorial law, for fixed $b$,
\[
\mathcal{M}^o_\Pi \longrightarrow \Hot, \qquad a \longmapsto \Ch_\infty(a,b),
\]
thanks to associativity. In fact, associativity will also imply that we have a \textit{bifunctor}
\[
\begin{array}{ccc}
\mathcal{M}^o_\Pi \times \mathcal{M}_\Pi & \longrightarrow & \Hot \\[1ex]
(a\,,\,b) & \longmapsto & \Ch_\infty(a,b),
\end{array}
\]
\printedpage{3}
% head: A. Grothendieck : Dérivateurs, Ch. IX, Édition des Universités de Montpellier II et Paris VII
% foot: Version 12 mars 2002
\msp{page 4}
for if
\[
a' \xrightarrow{\ u\ } a, \qquad b \xrightarrow{\ v\ } b',
\]
we have the diagram
\[
\begin{tikzcd}[column sep=large, row sep=large]
\Ch_\infty(a,b) \arrow[r, "\substack{\text{comp. with}\\ v}"] \arrow[d, "\text{comp. with } u"'] & \Ch_\infty(a,b') \arrow[d, "\text{comp. with } u"] \\
\Ch_\infty(a',b) \arrow[r, "\substack{\text{comp. with}\\ v}"] & \Ch_\infty(a',b'),
\end{tikzcd}
\]
which is commutative thanks to associativity for
\[
\Ch_\infty(a',a),\ \Ch_\infty(a,b),\ \Ch_\infty(b,b').
\]

\textbf{Corollary} of this bifunctorial dependence: \textit{If $u : a' \longrightarrow a$, $v : b \longrightarrow b'$ are invertible in $\mathcal{M}_\Pi$, then the corresponding arrow}
\DERmark{ENV0200}

\[
\Ch_\infty(a,b) \longrightarrow \Ch_\infty(a',b') \qquad \textit{in } \Hot
\]
\textit{is an} isomorphism.

The same arguments should work if instead of the $\Hom\,\mathrm{ext}$ with values in Hot, we took them in Ind\,Hot, or in Cat\,hot, or in Ind\,Cat\,hot -- I am thinking of course of the $\mathrm{CH}_\infty(a,b)$ (in the theory of localization) in Ind\,Cat\,hot, which can be regarded as coming from a bifunctor (where $\mathcal{M}_\Pi = \mathcal{M}\Sigma^{-1}$)
\[
\mathcal{M}^o_\Pi \times \mathcal{M}_\Pi \longrightarrow \mathrm{Ind}\,\mathrm{Cat}\,\mathrm{hot}.
\]
\DERmark{ENV0201}


\msp{page 5}
Here is the definitive axiomatization (?). We replace Hot by a category

\begin{itemize}[label={$\mathcal{H}$,}]
\item category with product $*$ associative, commutative (?) (\footnote{We should not need commutativity.}), unital, and all the usual compatibilities as needed,
\end{itemize}

the unit object for $*$ will be denoted 1, and we shall write
\[
\left\{
\begin{array}{ccc}
\pi_0(\xi) & \overset{\mathrm{def}}{=} & \Hom_{\mathcal{H}}(1,\xi) \\[1ex]
\pi_0 : \mathcal{H} & \longrightarrow & \mathrm{Ens}.
\end{array}
\right.
\]
\DERmark{ENV0202}


If we have
\[
P * Q \xrightarrow{\ u\ } R,
\]


\endgroup
% END UNIT Der-ChIX-678/Der-ChIX-678_p001-005.tex
% BEGIN UNIT Der-ChIX-678/Der-ChIX-678_p006-010.tex
\begingroup\makeatletter%
\NoAutoSpacing%
\newcommand{\printedpage}[1]{\marginpar{\footnotesize p.~#1}}%
\newcommand{\msp}[1]{\marginpar{\footnotesize [#1]}}%
\newcommand{\Hot}{\mathrm{Hot}}%
\newcommand{\HotO}{\mathrm{Hot}_0^{\wedge}}%
\newcommand{\MOP}{\mathcal{M}_{0\,\Pi}}%
\makeatother\let\printedpage\derpp\let\msp\dermsp%
\let\mspin\DERmspin%


\printedpage{4}
% head: A. Grothendieck : Dérivateurs, Ch. IX, Édition des Universités de Montpellier II et Paris VII
% foot: Version 12 mars 2002

we deduce from this
\[
\left\{
\begin{array}{lcl}
\pi_0(P) & \xrightarrow{c_u^{\mathrm{l}}} & \mathrm{Hom}_{\mathcal{H}}(Q,R)\\
\pi_0(Q) & \xrightarrow{c_u^{\mathrm{r}}} & \mathrm{Hom}_{\mathcal{H}}(P,R)
\end{array}
\right.
\]
as follows:
\[
\begin{array}{rlcl}
\text{If } \alpha\in\pi_0(P),\ \text{i.e. } 1\xrightarrow{\ \alpha\ }P,\ \text{then} & \multicolumn{3}{l}{\alpha*\mathrm{id}_Q : Q\longrightarrow P*Q}\\
\text{and} & c^{\mathrm{l}}_u(\alpha) & = & u\circ(\alpha*\mathrm{id}_Q)\\
\text{and dually} & c^{\mathrm{r}}_u(\beta) & = & u\circ(\mathrm{id}_P*\beta).
\end{array}
\]
If we have three objects $P_0$, $P_1$, $P_2$ in $\mathcal{H}$, pairings
\[
\left\{
\begin{array}{rcl}
P_0*P_1 & \xrightarrow{u_1} & Q_0\\
P_1*P_2 & \xrightarrow{u_2} & Q_1\\
P_0*Q_1 & \xrightarrow{u_3} & R\\
Q_0*P_2 & \xrightarrow{u_4} & R
\end{array}
\right.
\]
\[
\begin{tikzcd}
a \arrow[r, "P_0"'] \arrow[rr, bend left=20, "Q_0"] \arrow[rrr, bend left=35, "R"] & b \arrow[r, "P_1"'] \arrow[rr, bend right=20, "Q_1"'] & c \arrow[r, "P_2"'] & d
\end{tikzcd}
\]

\msp{page 6}
with an associativity relation, saying that
\[
\begin{tikzcd}
P_0*P_1*P_2 \arrow[r] \arrow[d] & Q_0*P_2 \arrow[d]\\
P_0*Q_1 \arrow[r] & R
\end{tikzcd}
\]
is commutative, then we deduce from this two transitivity relations and one commutation relation, for the operations associated with the elements of $\pi_0(P_0)$, $\pi_0(P_1)$, $\pi_0(P_2)$, which we will not spell out, except for one:
\[
\begin{array}{rrrcl}
\multicolumn{5}{l}{\text{Let } \beta\in\pi_0(P_1),\ \gamma\in\pi_0(P_2),}\\
\text{whence by } \beta : & c_{u_1}(\beta) : & P_0 & \longrightarrow & Q_0,\\
\text{and by } \gamma : & c_{u_4}(\gamma) : & Q_0 & \longrightarrow & R,\\
\text{and by } \beta *_{u_2}\gamma\in\pi_0(Q_1) : & c_{u_3}(\beta\,\texttt{[}{*_{u_2}}\texttt{]}\gamma) : & P_0 & \longrightarrow & R,
\end{array}
\]
\DERmark{ENV0203}


\printedpage{5}
% head: A. Grothendieck : Dérivateurs, Ch. IX, Édition des Universités de Montpellier II et Paris VII
% foot: Version 12 mars 2002
\texttt{[rather} $\beta\in\pi_0(P_1)$\texttt{,} $\gamma\in\pi_0(P_2)$\texttt{]}, well, we have
\[
\boxed{\;c(\beta *_{u_2}\gamma) \;=\; c(\gamma)c(\beta).\;}
\]
We still have to define $\beta *_{u_2}\gamma$, via the pairing $P_1*P_2\xrightarrow{u_2}Q_1$, we have in fact
\[
\begin{array}{ccc}
\pi_0(P_1)\times\pi_0(P_2) & \longrightarrow & \pi_0(P_1*P_2)\\
(\beta\,,\,\gamma) & \longmapsto & \beta*\gamma,
\end{array}
\]
and we compose with
\[
\pi_0(u_2) : \pi_0(P_1*P_2) \longrightarrow \pi_0(Q_1)
\]
\DERmark{ENV0204}

\texttt{[rather} $\pi_0(u_2) : \pi_0(P_1*P_2) \longrightarrow \pi_0(Q_1)$\texttt{]}.

\msp{page 7}
We give ourselves a set $\mathcal{M}_0$, and
\[
\left\{
\begin{array}{rclccrcl}
\mathcal{M}_0\times\mathcal{M}_0 & \xrightarrow{\mathrm{Ch}} & \mathcal{H} & \text{denoted} & & (a,b) & \longmapsto & \mathrm{Ch}(a,b)\ \text{or}\ \mathrm{Hom\,ext}_{\mathcal{H}}(a,b)\\
\mathcal{M}_0\times\mathcal{M}_0\times\mathcal{M}_0 & \xrightarrow{\text{comp.}} & \mathrm{Fl}\,\mathcal{H} & & & (a,b,c) & \longmapsto & \bigl(\mathrm{Ch}(a,b)*\mathrm{Ch}(b,c)\longrightarrow\mathrm{Ch}(a,c)\bigr)
\end{array}
\right.
\]
\DERmark{ENV0205}

with two axioms

\begin{enumerate}
\item[a)] Associativity for the compositions, for $(a,b,c,d)\in\mathcal{M}_0^4$
\[
\begin{tikzcd}
\mathrm{Ch}(a,b)*\mathrm{Ch}(b,c)*\mathrm{Ch}(c,d) \arrow[r] \arrow[d] & \mathrm{Ch}(a,c)*\mathrm{Ch}(c,d) \arrow[d]\\
\mathrm{Ch}(a,b)*\mathrm{Ch}(b,d) \arrow[r] & \mathrm{Ch}(a,d).
\end{tikzcd}
\]
\DERmark{ENV0206}\DERmark{ENV0207}\DERmark{ENV0208}\DERmark{ENV0209}

\item[b)] Existence of two-sided unit objects
\[
1_a\in\pi_0(\mathrm{Ch}(a,a)) \qquad \text{for every } a\in\mathcal{M}_0
\]
\end{enumerate}

Then the $\pi_0(\mathrm{Ch}(a,b))$ and their composition define on $\mathcal{M}_0$ as set of objects a \textit{category} structure $\MOP$, having the $1_a$ as identities, and we find that
\[
(a,b)\longmapsto\mathrm{Ch}(a,b)
\]
can be interpreted in terms of a canonical \textit{bifunctor}
\[
\mathcal{M}^o_{0\,\Pi}\times\MOP \longrightarrow \mathcal{H},
\]
whose composite with
\[
\mathcal{H}\xrightarrow{\pi_0}\mathrm{Ens}
\]
is none other than the bifunctor $(a,b)\longrightarrow\mathrm{Hom}_{\MOP}(a,b)$,
\printedpage{6}
% head: A. Grothendieck : Dérivateurs, Ch. IX, Édition des Universités de Montpellier II et Paris VII
% foot: Version 12 mars 2002
\msp{page 8}
\[
\mathrm{Hom}_{\MOP}(a,b) \;\simeq\; \pi_0(\mathrm{Ch}(a,b)).
\]
This relation is the \textit{definition} of $\underline{\mathrm{Hom}}_{\MOP}$, but here the relation is understood as a % Reading resolved against PDF leaf 18: a single underline; see corrections.csv.
\textit{functorial} bijection in $a$, $b$.

We find in particular that if
\[
u : a'\longrightarrow a, \qquad v : b\longrightarrow b'
\]
are \textit{invertible} arrows in $\MOP$, then
\[
\mathrm{Ch}(a,b)\longrightarrow\mathrm{Ch}(a',b')
\]
corresponding in $\mathcal{H}$ is \texttt{[an]} \textit{isomorphism.}

We will apply all this mainly to the two situations deduced from a situation $(\mathcal{M},\Sigma)$, with
\[
\begin{array}{rlcllcl}
1^o) & \mathcal{H} & = & \mathrm{Hot}, & \mathrm{Ch}(a,b) & = & \mathrm{Ch}^{\Sigma}_{\infty}(\mathcal{M};a,b)\\
2^o) & \mathcal{H} & = & \mathrm{Ind\,Cat\,hot}, & \mathrm{Ch}(a,b) & = & \mathrm{CH}^{\Sigma}_{\infty}(\mathcal{M};a,b)\ \text{(\footnotemark[2])}.
\end{array}
\]
\footnotetext[2]{$*$ in Hot and Ind\,Cat\,hot is the cartesian product, nothing more.}
In both cases, we thus have a functor on $(\mathcal{M}\Sigma^{-1})^o\times(\mathcal{M}\Sigma^{-1})$, with values either in Hot, or in Ind\,Cat\,hot. (\textbf{NB} The two category structures on $\mathcal{M}_0$ deduced from the % [sic] des -> de, « déduites des ces deux foncteurs »
these two functors are canonically isomorphic.)
\[
\begin{tikzcd}[column sep=small]
 & \mathrm{Hot} & (a,b)\longmapsto\mathrm{Ch}^{\Sigma}_{\infty}(a,b)\\
(\mathcal{M}\Sigma^{-1})^o\times(\mathcal{M}\Sigma^{-1}) \arrow[ur] \arrow[dr] & & \\
 & \mathrm{Ind\,Cat\,hot} & (a,b)\longmapsto\mathrm{CH}^{\Sigma}_{\infty}(a,b)
\end{tikzcd}
\]

\msp{page 9}
As the canonical functors
\[
\begin{tikzcd}
\mathrm{Ind\,Cat\,hot} \arrow[r, "\mathrm{Ind}\,\varphi"] & \mathrm{Ind\,Hot} \arrow[r, "\mathrm{Ind}\,\psi"] & \mathrm{Ind\,Hot}_0^{\wedge} \arrow[r, "\varinjlim"] & \mathrm{Hot}_0^{\wedge},\\
\mathrm{Cat\,hot} \arrow[u, hook] \arrow[r, "\varphi"] & \mathrm{Hot} \arrow[u, hook] \arrow[r, "\psi"] & \mathrm{Hot}_0^{\wedge} \arrow[u, hook] \arrow[ur, "\mathrm{id}"'] &
\end{tikzcd}
\]
commuting with finite products, any theory of ``external Hom'' with values in one of these categories, defines one with values in each of those into which it maps.
\printedpage{7}
% head: A. Grothendieck : Dérivateurs, Ch. IX, Édition des Universités de Montpellier II et Paris VII
% foot: Version 12 mars 2002
Besides, these functors are not all fully faithful (this is in any case false for $\varphi$ and $\mathrm{Ind}\,\varphi$, and I do not know what the situation is for $\psi$ and $\mathrm{Ind}\,\psi$ -- the vertical arrows on the other hand are fully faithful (\footnote[3]{and also, $\mathrm{Ind\,Hot}_0^{\wedge}\xrightarrow{\varinjlim}\mathrm{Hot}_0^{\wedge}$ is \textit{not} fully faithful.}).) But they induce bijections for the $\pi_0$. So the category $\MOP$ deduced from a theory of external Hom with values in one of these categories, and the one % [sic] catégoires -> catégories, faute d'orthographe
deduced in another into which it maps, is the same up to canonical isomorphism. As far as this definition of $\MOP$ is concerned, nothing is changed then when we regard the $\mathrm{CH}_{\infty}$ as objects of Ind\,Hot, or only of $\mathrm{Ind\,Hot}_0^{\wedge}$, or even only of $\HotO$, or when we regard the $\mathrm{Ch}_{\infty}$ as objects
\msp{page 10}
in $\HotO$. But of course, the datum of the composition pairings loses precision each time (except perhaps when passing from Hot to $\HotO$, if $\psi$ is fully faithful \dots\,). When we take the images of the $\mathrm{CH}^{\Sigma}_{\infty}(a,b)$ and $\mathrm{Ch}^{\Sigma}_{\infty}(a,b)$ in $\HotO$ (the last category of the diagram), we find canonically isomorphic objects (functorially in $a$, $b$ in $\MOP\simeq\mathcal{M}\Sigma^{-1}$). When we regard them as being respectively in Ind\,Hot and in $\mathrm{Hot}\hookrightarrow\mathrm{Ind\,Hot}$, we find a functorial homomorphism
\[
\mathrm{CH}^{\Sigma}_{\infty}(a,b)\longrightarrow\mathrm{Ch}^{\Sigma}_{\infty}(a,b) \qquad\quad \text{in Ind\,Hot or in } \mathrm{Ind\,Hot}_0^{\wedge} \text{ as we choose.} \tag*{$(*)$}
\]
The $\varinjlim$ of $\mathrm{CH}^{\Sigma}_{\infty}(a,b)$ in Hot does not exist in general (unless the ind-object is essentially constant), but it exists in $\HotO$, and the canonical isomorphism
\[
\varinjlim_{\HotO}\mathrm{CH}^{\Sigma}_{\infty}(a,b)\xrightarrow{\ \sim\ }\mathrm{Ch}^{\Sigma}_{\infty}(a,b) \qquad\quad \text{in } \HotO \tag*{$(**)$}
\]
is the arrow deduced from $(*)$ in $\mathrm{Ind\,Hot}_0^{\wedge}$.

\msp{page 11}
In the two cases CH and Ch (which we are going to denote by the common symbol $C_{\infty}(a,b)$), the functoriality of this object $C_{\infty}(a,b)$ (even in the weakest sense, as an object of $\HotO$) with respect to $a$ and $b$ (even in the weakest sense, as a functor $\mathcal{M}^o\times\mathcal{M}\longrightarrow\HotO$ and not even $(\mathcal{M}\Sigma^{-1})^o\times(\mathcal{M}\Sigma^{-1})\longrightarrow\HotO$) is not at all evident from the direct definition in terms of the explicit direct system in Cat
\[
C_{\infty}(a,b) \;=\; (\mathrm{Ch}_{\tau_0^{2n}}(a,b))_{n\in\mathbf{N}},
\]
for an arrow $a'\longrightarrow a$ or $b\longrightarrow b'$ in $\mathcal{M}$, and even were it in $\Sigma$, in no way defines an arrow in Cat
\[
\begin{array}{rcl}
C_{\infty}(a,b) & \longrightarrow & C_{\infty}(a',b)\\
\text{or}\ \ C_{\infty}(a,b) & \longrightarrow & C_{\infty}(a,b').
\end{array}
\]
This difficulty recalls the one that \textsc{Quillen} had to overcome, in Homotopical Algebra, to face the difference between ``left homotopies'' and ``right homotopies''. I recall how
\printedpage{8}
% head: A. Grothendieck : Dérivateurs, Ch. IX, Édition des Universités de Montpellier II et Paris VII
% foot: Version 12 mars 2002
\msp{page 12}
we can nevertheless define such arrows, as resulting from a (much richer) arrow of ``composition'' of paths between $a'$ and $a$, and $a$ and $b$, or between $a$ and $b$, and $b$ and $b'$, that is in general then the pairing
\[
C_{\infty}(a,b)\times C_{\infty}(b,c)\longrightarrow C_{\infty}(a,c),
\]
whose direct definition is no less problematic, and can in any case \textit{not} be given as an arrow in Cat itself, or in Ind\,Cat. Interpreting e.g. $\mathrm{Ch}_{\infty}(a,b)$ in terms of \textit{infinite} paths (on both sides) from $a$ to $b$, in an obvious sense (but % [sic] cotés -> côtés, accent omis
``essentially finite'' on both sides), we can also define % [sic] cotés -> côtés, accent omis
\[
\mathrm{Ch}^{-}_{\infty}(a,b), \qquad \mathrm{Ch}^{+}_{\infty}(a,b)
\]
as the full subcategories formed respectively of paths which are totally degenerate in degrees $\geq 0$ (i.e. $a_0\longleftarrow a_1\longrightarrow a_2\longleftarrow a_3\cdots$ are all identities), which amounts to paths infinite only to the left, i.e. stopping at $a_0$, and symmetrically for $\mathrm{Ch}^{+}_{\infty}(a,b)$. Same story for the $\mathrm{CH}^{+}$, $\mathrm{CH}^{-}$.

\msp{page 13}
That said, we have this (returning to the ambivalent notion $C_{\infty}$):

\begin{enumerate}
\item[a)] the fully faithful inclusions
\[
C^{-}_{\infty}(a,b)\hookrightarrow C_{\infty}(a,b)\hookleftarrow C^{+}_{\infty}(a,b)
\]
are isomorphisms in Ind\,Cat\,hot resp. in Hot, and
\item[b)] We really do have a composition of paths
\[
C^{-}_{\infty}(a,b)\times C^{+}_{\infty}(b,c)\longrightarrow C_{\infty}(a,c).
\]
\end{enumerate}

The \textit{associativity} in the CH case moreover offers no difficulty, because of the possibility of reducing to the
\[
C_m(a,b)\times C_n(b,c)\times C_p(c,d)\rightrightarrows C_N(a,d)
\]
for $N$ large, and proving equality as arrows in Cat\,hot -- which will follow from the homotopy lemma % [sic] egalité -> égalité, accent omis
on intervals in Cat. In the case of $\mathrm{Ch}_{\infty}$ on the other hand I have not found a proof. (Only for lack of having tried seriously to get down to it? I don't much feel like it, seeing that it might be that $\mathrm{Hot}\longrightarrow\HotO$ is fully faithful, in which case the work becomes useless \dots\,)


\endgroup
% END UNIT Der-ChIX-678/Der-ChIX-678_p006-010.tex
% BEGIN UNIT Der-ChIX-678/Der-ChIX-678_p011-015.tex
\begingroup\makeatletter%
\NoAutoSpacing%
\newcommand{\printedpage}[1]{\marginpar{\footnotesize p.~#1}}%
\newcommand{\msp}[1]{\marginpar{\footnotesize [#1]}}%
\newcommand{\M}{\mathcal{M}}%
\newcommand{\Hot}{\mathrm{Hot}}%
\newcommand{\Ch}{\mathrm{Ch}}%
\newcommand{\uCh}{\underline{\mathrm{Ch}}}%
\newcommand{\Homext}{\mathrm{Hom\,ext}}%
\newcommand{\Lotimes}{\overset{\mathbf{L}}{\otimes}}%
\makeatother\let\printedpage\derpp\let\msp\dermsp%
\let\mspin\DERmspin%


\printedpage{9}
% head: A. Grothendieck : Dérivateurs, Ch. IX, Édition des Universités de Montpellier II et Paris VII
% foot: Version 12 mars 2002

\msp{page 14}

\section*{2\quad Comparison with Quillen}

A certain number of questions arise about this functor $(\M\Sigma^{-1})^o \times (\M\Sigma^{-1}) \longrightarrow \Hot$, and its relation to the classical constructions, in particular to the theory of \textsc{Quillen}. I will now write
\[
\begin{aligned}
\Ch^{\Sigma}_{\infty}(a,b) \;&=\; \Homext_{\Hot}(a,b) && \text{or simply}\\
&\phantom{\;=\;} \Homext(a,b).
\end{aligned}
\]

1$^o$). Let $f,g : a \rightrightarrows b$ be in $\M$, then (if $\M$ is a ``model category'') \textsc{Quillen} defines a set % [sic] Quilllen -> Quillen, three l's printed
\DERmark{ENV0210}

\[
\pi_1(a,b;f,g) \qquad \text{or} \qquad \pi_1(f,g)
\]
of homotopy classes from $f$ to $g$, in such a way that $\mathrm{Hom}_{\M}(a,b)$ becomes the set of objects of a groupoid, whose arrows are the elements of $\pi_1(f,g)$, for a natural notion of composition. I suspect that one has
\[
\pi_1(f,g) \;=\; \mathrm{Hom}_{\underbrace{\Pi(\uCh^{\Sigma}_{\infty}(a,b))}_{\substack{\text{fundamental}\\ \text{groupoid}\\ \text{of } \uCh^{\Sigma}_{\infty}(a,b)}}}(f,g), % [sic] fodamental -> fondamental, missing n
\]
when $f$ and $g$ are considered as ``infinite paths'' (by inserting each $f : a \longrightarrow b$ into
\[
\cdots \underbrace{a_{-2}}_{=a} \xleftarrow{\ \mathrm{id}\ } \underbrace{a_{-1}}_{=a} \xrightarrow{\ f\ } \underbrace{a_{0}}_{=b} \xleftarrow{\ \mathrm{id}\ } \underbrace{a_{1}}_{=b} \longrightarrow \cdots\ \ )
\]

\msp{page 15}

So the definition of the $\pi_1(f,g)$ (if that is indeed it) extends in a natural way to arbitrary infinite (essentially finite) paths.

2$^o$) Suppose that $\M$ admits an object both initial and final $0_{\M}$ ($\M$ ``pointed''). I say that then all the categories $\uCh^{\Sigma}_{\infty}(a,b)$ are pointed in a natural way: namely, by the infinite path from $a$ to $b$ defined (as above) starting from $f = 0_{a,b} : a \longrightarrow b$, the zero map from $a$ to $b$. One then feels like viewing the $\uCh^{\Sigma}_{\infty}(a,b)$ as defining elements in $\Hot_{\bullet}$ (``pointed Hot''). This is finer than saying that one gives oneself a neutral class
\[
e \;\longrightarrow\; \Ch(a,b) \qquad \text{in Hot,}
\]
for the functor from $\Hot_{\bullet}$ to Hot (or to $e\backslash\Hot$) is not faithful. The problem is to define the arrows
\[
\Ch(a,b)\times\Ch(b,c) \;\longrightarrow\; \Ch(a,c)
\]
in $\Hot_{\bullet}$, and not just in Hot, and with associativity (and the units) in $\Hot_{\bullet}$ as well. If one applies the definition with $\Ch^{-}(a,b)$ and $\Ch^{+}(b,c)$, the composite of the zero elements $a \xrightarrow{\ 0\ } b$ of $\Ch^{-}(a,b)$ and $b \xleftarrow{\ \mathrm{id}\ } b \xrightarrow{\ 0\ } c$ of $\Ch^{+}(b,c)$
\printedpage{10}
% head: A. Grothendieck : Dérivateurs, Ch. IX, Édition des Universités de Montpellier II et Paris VII
% foot: Version 12 mars 2002
\msp{page 16}
is the following element of $\Ch_{\infty}(a,c)$
\[
\begin{tikzcd}[row sep=huge]
 & a_{-2} & a_{-1} \arrow[l] \arrow[r] & a_{0} \arrow[d, equal] & a_{1} \arrow[l] \arrow[r] & a_{2} & \cdots \\
{}_{\mathrm{id}}\ \cdots & a & a \arrow[l, "\mathrm{id}"'] \arrow[r, "0"'] & b & b \arrow[l, "\mathrm{id}"'] \arrow[r, "0"'] & c & \cdots_{\mathrm{id}} \arrow[l, "\mathrm{id}"']
\end{tikzcd}
\]
this is \textit{not} the zero element of $\Ch_{\infty}(a,c)$ ! But one has a canonical arrow
\DERmark{ENV0211}

\[
e^{-}_{a,b}\circ e^{+}_{b,c} \;\;\longrightarrow\;\; e_{a,c} \qquad \text{in } \Ch_{\infty}(a,c)
\]
\DERmark{ENV0212}

which is the identity on all the vertices except those of indices 0 and 1, where its components are zero. There is a sorites to develop, to the effect that a pair
\DERmark{ENV0213}

\[
(f : K_{\bullet} \longrightarrow L_{\bullet}, c)
\]
of a functor between pointed categories, and of a path $c : f(e_{K_{\bullet}}) \longrightarrow e_{L_{\bullet}}$, defines in $\Hot_{\bullet}$ a canonical arrow
\DERmark{ENV0214}

\[
\mathrm{hot}_{\bullet}(K_{\bullet}) \;\longrightarrow\; \mathrm{hot}_{\bullet}(L_{\bullet})
\]
with transitivity if one has
\[
(f' : L_{\bullet} \longrightarrow M_{\bullet}, c' : f'(e_{L_{\bullet}}) \longrightarrow e_{M_{\bullet}}),
\]

\msp{page 17}

whence
\[
f'' = f'\circ f : K_{\bullet} \longrightarrow M_{\bullet}, \qquad \text{and}\ \ c'' : f''(e_{K_{\bullet}}) \longrightarrow e_{M_{\bullet}}, % [sic] closing parenthesis without opening parenthesis -> (f'' = f'\circ f : ..., c'' : ...), unbalanced parenthesis
\]
\DERmark{ENV0215}

defines thus
\[
\begin{tikzcd}
f'(f(e_{K_{\bullet}})) \arrow[r, "f'(c)"] \arrow[rr, bend right=12, "{c''\ \text{by definition}}"'] & f'(e_{L_{\bullet}}) \arrow[r, "c'"] & e_{M_{\bullet}}
\end{tikzcd}
\]

\noindent\llap{?\quad}So this does indeed give us pairings in $\Hot_{\bullet}$ (and also in $\mathrm{Cat\,hot}_{\bullet}$?). It is hard to believe that they could fail to be associative! Besides, $\Hot_{\bullet}$, with the product $\divideontimes$ with unit $S^0$ (\footnote[4]{No, one must take the product $\xi_{\bullet} \divideontimes \eta_{\bullet}$ ($\xi_{\bullet}\times\eta_{\bullet}$, with $\xi_{\bullet}\vee\eta_{\bullet}$ contracted to $e$).}), satisfies the earlier axiomatics for $\mathcal{H}$, and $\Hot_{\bullet} \longrightarrow \Hot$ is equipped with the lax monoidal comparison from this product (\footnote[5]{No, but one gets by all the same \ldots}), and with $\pi_0$. % Reading resolved against PDF leaf 22: \divideontimes, meaning smash; see corrections.csv.
\DERmark{ENV0216}\DERmark{ENV0217}


If all this works, one would thus have a theory of the $\Homext_{\bullet}$ with values in $\Hot_{\bullet}$; and probably also for the $\mathrm{CH}_{\infty}$ as objects of $\mathrm{Ind\,Hot}_{\bullet}$, and not only of $\mathrm{Ind\,Hot}$. But
\printedpage{11}
% head: A. Grothendieck : Dérivateurs, Ch. IX, Édition des Universités de Montpellier II et Paris VII
% foot: Version 12 mars 2002
in any case one cannot in general drop Cat hot, which remains important -- so one would have to have objects $\mathrm{CH}_{\infty}(a,b)$ in $\mathrm{Ind\,Cat\,hot}_{\bullet}$, to have a well-tailored setup.

The $\pi_0(C(a,b))$ will now be pointed sets, and this will be by the image of the zero map $a \xrightarrow{\ 0\ } b$, i.e. the composite
\[
a \longrightarrow 0 \longrightarrow b
\]

\msp{page 18}

which one must compose with the element deduced from the truncated path % [sic] l'elément -> l'élément, missing accent
\[
a \longrightarrow \underbrace{0}_{=a_0} \xleftarrow{\mathrm{id}} 0 \longrightarrow b \qquad \ldots
\]
\DERmark{ENV0218}


One would have to show that 0 is also a neutral object of $\M\Sigma^{-1}$, i.e. that
\[
\pi_0(a,0) = \pi_0(0,a) = \{0\} \qquad \text{for every } a \in \mathrm{Ob}\,\M
\]
and interpret the zero element in $\pi_0(a,b)$ as the zero map in $\M\Sigma^{-1}$. In general, one needs conditions which ensure (without $\M$ necessarily being pointed) that
\[
\M \xrightarrow{\ \gamma\ } \M\Sigma^{-1}
\]
commutes with finite sums, with finite products (so $\gamma(\varphi_{\M}) = \varphi_{\M\Sigma^{-1}}$, resp. $\gamma(e_{\M}) = e_{\M\Sigma^{-1}}$) (\footnote[6]{to do.}).

Suppose finally that $\M$ is \textit{additive,} and even $k$-linear ($k$ a fixed ring), i.e. one has in addition a ring homomorphism
\[
k \;\longrightarrow\; \mathrm{End}(\mathrm{id}_{\M}).
\]

One must assume some compatibility of $\Sigma$ with the additive structure, e.g. that the direct sum of two arrows in $\Sigma$ is in $\Sigma$, and probably something stronger, of the ``five lemma'' kind for $\Sigma$. (But one does not necessarily assume $\M$ to be abelian, so one would have trouble formulating a five lemma?) % [sic] dû mal -> du mal, accent on du

\msp{page 19}

But perhaps here a subtler axiom is needed, e.g. (if $\underbrace{\Sigma}_{\text{suspension}}$ and $\Omega$ exist (\footnote[7]{cf. below what this means.})) that the set of adjunction morphisms $\Sigma\Omega \longrightarrow \mathrm{id}$, $\mathrm{id} \longrightarrow \Omega\Sigma$ be \texttt{[made up of]} isomorphisms, which implies that $\Omega$ resp. $\Sigma$ is fully faithful. In any case, the natural expectation here is that $\Ch^{\Sigma}_{\infty}(a,b)$ be not only in $\Hot_{\bullet}$, but even in the category $D^{\leq 0}(k\text{-Mod})$, a full subcategory of the derived category $D^{-}(k\text{-Mod})$ consisting of the $K$ such that $H^{i}(K) = 0$ if $i > 0$. One can also denote it $D_{\geq 0}(k\text{-Mod})$, interpreting it in
\printedpage{12}
% head: A. Grothendieck : Dérivateurs, Ch. IX, Édition des Universités de Montpellier II et Paris VII
% foot: Version 12 mars 2002
terms of complexes of ``chains'' rather than of cochains. Here the pairings must correspond to the operation $* = \Lotimes_k$ in $D_{\geq 0}(k\text{-Mod})$, admitting $k$ as unit object, and
\[
\pi_0(K_{\bullet}) \;\simeq\; H_0(K_{\bullet}).
\]
So one must have
\[
\Homext(a,b) \Lotimes_k \Homext(b,c) \;\longrightarrow\; \Homext(a,c),
\]
where this time Hom ext denotes the object of $D_{\geq 0}(k\text{-Mod})$, instead of just an object of Hot, or of $\Hot_{\bullet}$. One uses the functors
\[
D_{\geq 0}(k\text{-Mod}) \longrightarrow \Hot_{\bullet} \longrightarrow \Hot
\]
\msp{page 20}
to make the link with the pointed theory in $\Hot_{\bullet}$, or \texttt{[the]} general theory in Hot. If one denotes by
\[
\gamma :\ D_{\geq 0}(k\text{-Mod}) \;\longrightarrow\; \Hot_{\bullet}
\]
the first of these functors, it does not commute with the products $*$ ($\Lotimes_k$ and $\times$ respectively), but one has, for $P_{\bullet}$, $Q_{\bullet}$ varying in $D_{\geq 0}(k\text{-Mod})$, a canonical functorial arrow % [sic] variable -> variables, agreement with P, Q
\[
\gamma(P_{\bullet} \Lotimes_k Q_{\bullet}) \xleftarrow{\ \mathit{can.}\ } \gamma(P_{\bullet})\times\gamma(Q_{\bullet})
\]
``$k$-bilinear'', which shows that giving pairings
\[
P_{\bullet} * Q_{\bullet} \xrightarrow{\ c_k\ } R_{\bullet},
\]
whence
\[
\gamma(P_{\bullet} * Q_{\bullet}) \xrightarrow{\ \gamma(c_k)\ } \gamma(R_{\bullet}), \tag*{$(*)$}
\]
is in either case finer than giving oneself only
\[
\gamma(P_{\bullet}) * \gamma(Q_{\bullet}) \xrightarrow{\ c_{\bullet}\ } \gamma(R_{\bullet}), \tag*{$(*')$}
\]
seeing that a datum $(*)$ will imply a datum $(*')$ by % [sic] donne -> donnée, missing letters
\[
\begin{tikzcd}
 & \gamma(P_{\bullet} * Q_{\bullet}) \arrow[dr, "\gamma(c)"] & \\
\gamma(P_{\bullet}) * \gamma(Q_{\bullet}) \arrow[ur, "\text{can.}"] \arrow[rr, "c_{\bullet}"] & & \gamma(R_{\bullet}).
\end{tikzcd}
\]
In view of this, one does indeed expect to have a \textit{strong} theory of the pairings
\printedpage{13}
% head: A. Grothendieck : Dérivateurs, Ch. IX, Édition des Universités de Montpellier II et Paris VII
% foot: Version 12 mars 2002
\msp{page 21}
\[
\Homext_k(a,b) \Lotimes_k \Homext_k(b,c) \;\longrightarrow\; \Homext_k(a,c),
\]
from which one will deduce the weaker pairings in $\Hot_{\bullet}$ (instead of $D_{\geq 0}(k\text{-Mod})$)
\[
\Homext_{\Hot_{\bullet}}(a,b) \times \Homext_{\Hot_{\bullet}}(b,c) \;\longrightarrow\; \Homext_{\Hot_{\bullet}}(a,c).
\]

3$^o$) Returning to the case $\M$ pointed and nothing more, \textsc{Quillen} defines, for $a$, $b$ in $\M$,
\[
\pi_1(a,b) \;=\; \pi_1(a,b\,;0_{a,b},0_{a,b})
\]
\DERmark{ENV0219}

as a bifunctor on $(\M\Sigma^{-1})^o\times(\M\Sigma^{-1})$ (with the notation of 1$^o$)). If what I put forward without proof in 1$^o$) is true, it must follow that one has a canonical isomorphism
\[
\pi_1(a,b) \;\simeq\; \pi_1(\Ch^{\Sigma}_{\infty}(a,b)),
\]
where $\Ch^{\Sigma}_{\infty}(a,b)$ is viewed as an object of $\Hot_{\bullet}$, so its $\pi_1$ is defined. One has in any case, without any reference to cofibrations and fibrations, a bifunctor with values in Groups, composed in
\[
(\M\Sigma^{-1})^o\times(\M\Sigma^{-1}) \xrightarrow{\ \ (a,b)\longmapsto \Ch^{\Sigma}_{\infty}(a,b)\ \ } \Hot_{\bullet} \xrightarrow{\ \pi_1\ } (\text{Groups}).
\]
If for fixed $a$, variable $b$,
\[
b \;\longmapsto\; \pi_1(a,b)
\]
is representable, one denotes the object representing it by $\Sigma(a)$ (\textit{suspension} -- do not confuse this $\Sigma$ with $\Sigma \subseteq \mathrm{Fl}(\M)$, which it is better now to denote by $W$, as in ``weak equivalence''). So

\msp{page 22}

\[
\mathrm{Hom}(\Sigma a,b) \;\simeq\; \pi_1(a,b),
\]
and $\Sigma a$ depends functorially on $a$, subject to existence. If it always exists, one has
\[
\Sigma :\ \M W^{-1} \;\longrightarrow\; \M W^{-1}.
\]
One defines likewise, subject to existence,
\[
\Omega :\ \M W^{-1} \;\longrightarrow\; \M W^{-1}
\]
by
\[
\mathrm{Hom}(a,\Omega b) \;\simeq\; \pi_1(a,b),
\]
and if both exist, one has the adjunction formula
\[
\mathrm{Hom}(\Sigma a,b) \;\simeq\; \mathrm{Hom}(a,\Omega b) \qquad (\simeq \pi_1\Homext_{\Hot_{\bullet}}(a,b)).
\]


\endgroup
% END UNIT Der-ChIX-678/Der-ChIX-678_p011-015.tex
% BEGIN UNIT Der-ChIX-678/Der-ChIX-678_p016-018.tex
\begingroup\makeatletter%
\newcommand{\printedpage}[1]{\marginpar{\footnotesize p.~#1}}%
\newcommand{\msp}[1]{\marginpar{\footnotesize [#1]}}%
\newcommand{\cercle}[1]{\tikz[baseline=(c.base)]\node[draw,circle,inner sep=1pt](c){#1};}%
\newcommand{\Homext}{\mathrm{Hom}\,\mathrm{ext}}%
\newcommand{\Chu}{\underline{\mathrm{Ch}}}%
\newcommand{\Hot}{\mathrm{Hot}}%
\makeatother\let\printedpage\derpp\let\msp\dermsp%
\let\mspin\DERmspin%


\printedpage{14}
% head: A. Grothendieck : Dérivateurs, Ch. IX, Édition des Universités de Montpellier II et Paris VII
% foot: Version 12 mars 2002

The existence of
\[
\begin{aligned}
\Omega :\ & \mathcal{M}W^{-1} \ \longrightarrow\ \mathcal{M}W^{-1}, \quad \text{or of}\\
\Sigma :\ & \mathcal{M}W^{-1} \ \longrightarrow\ \mathcal{M}W^{-1},
\end{aligned}
\]
or even of both at once, appears as an extremely strong hypothesis \cercle{a} on $(\mathcal{M}, W)$. This condition \textit{plus} the one \cercle{b} that the $\Chu_{\infty}(a,b)$ be homotopic to Thomason-Kan categories, or (which comes to the same thing, more or less?) that $Y \longmapsto \mathrm{Hom}_{\mathcal{M}(Y)W_Y^{-1}}(a_Y, b_Y)$ transform Hot-equivalences of Cat into bijections (even if only for fixed $Y$ \dots), seems the crucial test for knowing whether ``$(\mathcal{M}, \Sigma)$ \texttt{[rather} $(\mathcal{M}, W)$\texttt{]} lends itself to a homotopy theory ``in $\mathcal{M}$'' ''. I doubt that

\msp{page 23}

these conditions can be satisfied for finite $\mathcal{M}$, except in the trivial case where $\mathcal{M} = \emptyset$, where $\mathcal{M}\Sigma^{-1}$ \texttt{[rather} $\mathcal{M}W^{-1}$\texttt{]} is equivalent to the one-point category.

A third test \cercle{c} which seems important to me (but perhaps this follows from the rest), is the formulas in $\Hot_\bullet$
\[
\left\{
\begin{aligned}
\Homext_\bullet(\Sigma a, b) \ &\simeq\ \Homext_\bullet(a, \Omega b)\\
&\simeq\ \Omega\Homext_\bullet(a, b),
\end{aligned}
\right.
\tag*{$(*)$}
\]
where $\Homext_\bullet$ denotes the external Hom with values in $\Hot_\bullet$. One must view it separately (assuming the existence either of $\Omega$, or of $\Sigma$) as two formulas ``dual'' to one another,
\[
\left\{
\begin{aligned}
\Homext_\bullet(a, \Omega b) \ &\simeq\ \Omega\Homext_\bullet(a, b)\\
\Homext_\bullet(\Sigma a, b) \ &\simeq\ \Omega\Homext_\bullet(a, b),
\end{aligned}
\right.
\tag*{$(**)$}
\]
isomorphisms of \textit{groups} in $\Hot_\bullet$. (\textbf{NB} $\Omega b$ is a group object of $\mathcal{M}W^{-1}$, for it represents a contrafunctor with values in Gr, and likewise $\Sigma a$ is a cogroup of $\mathcal{M}W^{-1}$. The functor $y \longmapsto \Homext(a, y)$ must commute with finite products, and thus transform group objects into likewise. And likewise $x \longmapsto \Homext(x, b)$ must transform finite sums into finite products, hence cogroups into groups (\footnote[8]{The functor $\Omega$ must come from an endomorphism of the derivator $\mathbf{D} = \mathbf{D}_{\mathcal{M},W}$, and the latter must be left exact, and dually for $\Sigma$.}).)

I have convinced myself that the isomorphisms $(**)$ have a chance of existing only if one has the prerequisite condition \cercle{b}.

\par

\printedpage{15}
% head: A. Grothendieck : Dérivateurs, Ch. IX, Édition des Universités de Montpellier II et Paris VII
% foot: Version 12 mars 2002

\msp{page 24}

\section*{3\quad Attempt at constructing the functor $\Omega$}

There is probably little chance of proving anything about the functor
\[
\Omega :\ \mathcal{M}W^{-1} \longrightarrow \mathcal{M}W^{-1},
\]
without having in hand a more or less explicit \textit{construction} of it. I am going to try, taking my inspiration from \textsc{Quillen}, to work out general conditions on $(\mathcal{M}, W)$ (more general than those of \textsc{Quillen}) which ensure the existence of $\Omega$, and allow one to construct it. The geometric idea goes back to \textsc{Serre-Cartan}. The task is, for an object $S$ of $\mathcal{M}$, to construct an object $\mathrm{Ch}(S)$ \textit{in $\mathcal{M}$,} playing the role of a ``path space'', and which must fit into a factorization diagram for $\mathrm{diag}_S : S \longrightarrow S \times S$ (we assume that in $\mathcal{M}$ finite $\varprojlim$, hence finite products, exist)
\[
\begin{tikzcd}
S \arrow[d, "i_S"] \arrow[dd, bend right=40, "\mathrm{diag}_S"'] & \\
\mathrm{Ch}(S) \arrow[d, "p_S"] & \\
S \times S & e. \arrow[l]
\end{tikzcd}
\]
The essential property is that $p$ be a \textit{``fibration''} (in a suitable sense, to be worked out).

\msp{page 25}

As a subsidiary matter it may be useful (perhaps) that $i$ be a ``cofibration'' (\footnote[9]{and perhaps that $S \longrightarrow e$ also be a fibration?}). One defines $\Omega(S)$ as the fiber of $\mathrm{Ch}(S)$ over $S{\times}S$ (with respect to the unique morphism $e = e_{\mathcal{M}} \longrightarrow S\times S$, where $e_{\mathcal{M}}$ is the zero object of $\mathcal{M}$).
\[
\begin{tikzcd}[column sep=large]
\mathrm{Ch}(S) \arrow[d] & \Omega S \arrow[l] \arrow[d]\\
S \times S & e. \arrow[l]
\end{tikzcd}
\]

\par

\printedpage{16}
% head: A. Grothendieck : Dérivateurs, Ch. IX, Édition des Universités de Montpellier II et Paris VII
% foot: Version 12 mars 2002

This is an object of $\mathcal{M}$, which we shall regard mainly as an object of $\mathcal{M}W^{-1}$. We must make sure that we have a bijection functorial in $T$ in $\mathcal{M}W^{-1}$
\[
\underbrace{\pi_1\Chu^W_\infty(T, S)}_{\overset{\mathrm{def}}{=}\,\pi_1(T,S)} \ \simeq\ \underbrace{\mathrm{Hom}_{\mathcal{M}W^{-1}}(T, \Omega S)}_{\overset{\mathrm{def}}{=}\,\pi_0\Chu^W_\infty(T,\Omega S)}.
\]
\DERmark{ENV0220}\DERmark{ENV0221}

For this, one must compare
\[
\Chu^W_\infty(T, \Omega S) \qquad\quad \text{and} \qquad\quad \underbrace{\Omega_{\Hot_\bullet}(\Chu^W_\infty(T, S))}_{\overset{\substack{\text{``by definition''}\\ \text{for $\Omega$ of $\mathrm{Hot}_\bullet$}}}{=}\ \Chu_\infty(\mathrm{Ch}^W_\infty(T,S);0,0)} \qquad ,
\]
since it is a matter of constructing a bijection between their $\pi_0$.


\endgroup
% END UNIT Der-ChIX-678/Der-ChIX-678_p016-018.tex
\clearpage\setcounter{footnote}{0}\phantomsection\addcontentsline{toc}{section}{Chapter X}\markboth{Chapter X}{Chapter X}
% BEGIN UNIT Der-ChX-789/Der-ChX-789_p001-005.tex
\begingroup\makeatletter%
\newcommand{\printedpage}[1]{\marginpar{\footnotesize p.~#1}}%
\newcommand{\msp}[1]{\marginpar{\footnotesize [#1]}}%
\newcommand{\D}{\mathbf{D}}%
\newcommand{\Cat}{\mathrm{Cat}}%
\newcommand{\Fl}{\operatorname{Fl}}%
\newcommand{\Ob}{\operatorname{Ob}}%
\newcommand{\bDelta}{\boldsymbol{\Delta}}%
\makeatother\let\printedpage\derpp\let\msp\dermsp%
\let\mspin\DERmspin%


% ===================== page 1 (p001) =====================
\printedpage{}
% no page number printed (title page)
% head: (none)

\begin{center}
{\LARGE \textsf{THE DERIVATORS}}

\bigskip

{\large \textbf{ALEXANDRE GROTHENDIECK}}

\vspace{3em}

{\large \textbf{Chapter X}}

\medskip

{\large \textbf{Comparison of} Cat \textbf{with} $\bDelta^{\wedge}$}
\end{center}

\par

% ===================== page 2 (p002) =====================
\printedpage{}
% no page number printed (editorial note)
% head: (none)

\par\noindent\textbf{Note by the editors of the transcribed edition.}\par
This text was deciphered and transcribed in \LaTeXe{} by M. Künzer. It was edited by M. Künzer, J. Malgoire, and G. Maltsiniotis. The transcription is as faithful as possible to the manuscript. For the few obvious corrections, or rare comments by the editors, as well as for the original numbering of the manuscript pages, typewriter characters \texttt{[typewriter]} in square brackets are used. A question mark in square brackets means that one is not sure of the preceding word.

This edition is provisional. Remarks, comments and corrections are welcome. Send a message to:

\textbf{[e-mail address removed]}

\hfill G. Maltsiniotis\hspace*{3em}

\par

% ===================== page 3 (p003), printed 1 =====================
\printedpage{1}
% head: A. Grothendieck : Dérivateurs, Ch. X, Édition des Universités de Montpellier II et Paris VII
% foot: Version 9 avril 2002

\msp{page 1}

\section*{1\quad Recall of the notions of D-asphericity, and link with the ``basic localizer $W_{\mathbf{D}}$''}

For a derivator $\D$, the notions in question are the following (cf. I):

\begin{itemize}
\item[(i)] \textit{Object} of Cat that is $\D$-\textit{aspheric} (i.e. $Z \longrightarrow e$ is a $\D$-equivalence). NB I prefer to say ``$\D$-aspheric'' instead of ``$\D$-contractible'', which sometimes leads to confusion.

\item[\textcircled{ii}] $\D$-equivalence for an arrow $f : X \longrightarrow Y$ in Cat.

\item[(iii)] \textit{Functor} $f : X \longrightarrow Y$ that is $\D$-\textit{aspheric} (i.e. the $X/y$ are $\D$-aspheric) and $\D$-\textit{coaspheric} (the $y\backslash X$ are $\D$-aspheric). This is what I had called ``strict $\D$-equivalence'' (or ``costrict''), or homologically $\D$-cofinal functor (or cohomologically $\D$-cofinal), but the terminology ``aspheric'' seems to me the simplest and the best. (Fished out of PS, Pursuing Stacks.)

\item[(iv)] \textit{Functor} $f : X \longrightarrow Y$ that is $\D$-\textit{smooth,} $\D$-\textit{proper} (i.e. such that the fibers of the induced functors $X/x \longrightarrow Y/f(x)$, respectively $x\backslash X \longrightarrow f(x)\backslash Y$, are $\D$-aspheric).

\item[(v)] $\D$-\textit{fibration} $f : X \longrightarrow Y$ (self-dual notion), i.e. such that for $Y'' \xrightarrow{u} Y'$ over $Y$, if $Y''$, $Y'$ \texttt{[are]} Hot-aspheric, then $u_X : X'' \longrightarrow X'$ in
\[
\begin{tikzcd}
X \arrow[d] & X' \arrow[l] \arrow[d] & X'' \arrow[l] \arrow[d] \\
Y & Y' \arrow[l] & Y'' \arrow[l]
\end{tikzcd}
\]
is a $\D$-equivalence (\footnote{We shall also say that $f$ is $\D$-\textit{fibrant.} (\textbf{NB} There is no notion of ``$\D$-cofibrant'' to my knowledge.)}). The $\D$-\textit{fibration} is called \textit{trivial} (or \textit{aspheric}) if and only if its fibers are aspheric, or again if and only if $f$ is aspheric, or if and only if $f$ is coasperic. This also means that $f$ is a \textit{universal} $\D$-equivalence, % [sic] coaspérique -> coasphérique, letter h missing
\msp{page 2}
i.e. that for every $Y' \longrightarrow Y$, $X' = X \times_Y Y' \longrightarrow Y'$ is a $\D$-equivalence.

\item[(vi)] We also have weakened variants of the notion of $\D$-fibrations, namely the ``$\D$-\textit{left fibrations'',} or the $f$ that are ``\textit{cohomologically} $\D$-\textit{fibrant'',} and likewise the ``$\D$-\textit{right fibrations'',} i.e. the $f$ that are ``\textit{homologically} $\D$-\textit{fibrant''} (\footnote{New terminology: \textit{weak} cohomological (resp. homological) fibrations, cf. I, pages 111-126 for a corrected revision of the various notions of ``$\D$-fibrations''.}). The ``left''
\printedpage{2}
% head: A. Grothendieck : Dérivateurs, Ch. X, Édition des Universités de Montpellier II et Paris VII
% foot: Version 9 avril 2002
or cohomological notion for $f : X \longrightarrow Y$ means that the functors $X_y \hookrightarrow X/y$ ($y \in \Ob Y$) and $X/y \longrightarrow X/y'$ (for $(u : y \longrightarrow y') \in \Fl Y$) are $\D$-equivalences. And dually with the $X_y \hookrightarrow y\backslash X$, $y'\backslash X \longrightarrow y\backslash X$ for the ``right'' or ``homological'' notion. The $\D$-fibrations are none other than the universal left (cohomological) $\D$-fibrations or the universal right (homological) $\D$-fibrations.
\end{itemize}

Finally, this whole string of notions is deduced from the single notion of $\D$-equivalence (ii). If I denote by
\[
W_{\D} \subseteq \Fl(\Cat)
\]
the set of these $\D$-equivalences, one finds what in Pursuing Stacks I called a ``localizer'' (localizer) in Cat (\footnote{But beware, I assume here that the domain Diag of $\D$ is the whole of Cat!}). In Pursuing Stacks it was such a subset
\[
W \subseteq \Fl(\Cat),
\]
and not the datum of a derivator, that
\msp{page 3}
was taken as the fundamental starting datum for doing homotopical algebra of all kinds. I had moreover developed a string of notions of the same kind (i) to (vi), starting from this datum $W$. It could be that the statements for the notions I developed in I, in terms of a derivator, are valid more generally in terms of a localizer $W$, at least insofar as they still make sense. (So for the moment I set aside the strictly cohomological or homological statements, which appeal directly to $\D$.) Although it could be that under conditions to be made explicit on $W$, which imply that $(\Cat, W)$ do indeed define a derivator $\D_{\Cat,W} = \D$, one still manages to interpret $W$ as a $W_{\D}$, so that the cohomological and homological statements still make sense.

I feel like first looking in Pursuing Stacks at the axioms I have singled out for localizers, and checking whether they are all satisfied for
\msp{page 4}
$W_{\D}$. These conditions L 1 to L 5 are stated in Pursuing Stacks, §64 (page 164).

\textbf{L 1}\enspace Saturation of $W_D$. O.K.

This follows formally from the fact that we have a functor
\[
\Cat \xrightarrow{\varphi} C \qquad (= \Cat^o)
\]
such that for $u \in \Fl(\Cat)$, $u \in W_{\D} \iff \varphi(u)$ isomorphism. (Take for $\varphi(X)$ the subcategory of $\D(X)$ consisting of the $\D$-coefficients \textit{constant} on $X$).

\textbf{L 2}\enspace ``$\bDelta^1$ is $W$-aspheric over $e$'', i.e. for every $X$ in Cat, $X \times \bDelta^1 \longrightarrow X$ is in $W$.

This also means that $\bDelta^1 \longrightarrow e$ is a trivial $\D$-fibration, which is indeed the case.

\par

% ===================== page 5 (p005), printed 3 =====================
\printedpage{3}
% head: A. Grothendieck : Dérivateurs, Ch. X, Édition des Universités de Montpellier II et Paris VII
% foot: Version 9 avril 2002

\textbf{L 3}\enspace If $X \in \Ob\Cat$ has a final object, then $X \longrightarrow e$ is in $W$.

In fact, we require for derivators that this be true assuming only $X$ aspheric (in the absolute sense), and one even concludes from this that $X \longrightarrow e$ is $W$-aspheric, i.e. $X \times Y \xrightarrow{\mathrm{pr}_2} Y$ is in $W$ for every $Y$ in Cat (\footnote{In fact L 2 $\Longrightarrow$ L 3.}).

\textbf{L 4}\enspace If $f : X \longrightarrow Y$ is such that $X/y \longrightarrow Y/y$ are in $W$, then $f$ is in $W$.

This is OK.

\msp{page 5}
\textbf{L 5}\enspace If $X$, $Y$ are Cat-fibered over $S$, and $X \longrightarrow Y$ is an $S$-morphism that is \textit{cartesian,} inducing on the fibers $X_s \longrightarrow Y_s$ which are in $W$, then $f$ is in $W$.

For $W_{\D}$, this is proposition 6, page 71 of I.

These axioms L 1 to L 5 (equivalent to L 1, L 3$'$ ($\bDelta^1 \longrightarrow e$ is in $W$) and L 5) are in the nature of stability axioms, ensuring that $W$ is large enough. There is another condition which limits $W$ to not being too large:

\textbf{L a)}\enspace If $f \in W$, then $\pi_0(f)$ is bijective.

This is not satisfied for all the $W_{\D}$ – e.g. if the derivator $\D$ is \textit{zero,} then $W_{\D} = \Fl\Cat$! But for the derivators useful in practice, this always seems to be satisfied.

\textbf{L b), c)}\enspace The functor
\[
\Cat \longrightarrow \Cat\, W^{-1}
\]
commutes with finite sums and finite products (\footnote{This works if $(\Cat, W)$ comes from a Quillen triple (with Q-fibrations, Q-cofibrations, Q-weak equivalences).}).

To examine what this means in concrete terms on $W$ – I have not checked that it is O.K. for the $W_{\D}$.

\textbf{L d)}\enspace If $f : X \longrightarrow Y$ in $W$, then $f^o : X^o \longrightarrow Y^o$ as well.

This is corollary 2, page 75 of I.


\endgroup
% END UNIT Der-ChX-789/Der-ChX-789_p001-005.tex
\clearpage\setcounter{footnote}{0}\phantomsection\addcontentsline{toc}{section}{Chapter XI}\markboth{Chapter XI}{Chapter XI}
% BEGIN UNIT Der-ChXI-890/Der-ChXI-890_p001-005.tex
\begingroup\makeatletter%
\newcommand{\printedpage}[1]{\marginpar{\footnotesize p.~#1}}%
\newcommand{\msp}[1]{\marginpar{\footnotesize [#1]}}%
\newcommand{\Hot}{\mathrm{Hot}}%
\newcommand{\Cat}{\mathrm{Cat}}%
\newcommand{\Ob}{\operatorname{Ob}}%
\makeatother\let\printedpage\derpp\let\msp\dermsp%
\let\mspin\DERmspin%


% ===================== page 1 =====================
\printedpage{--}
% head: (none) -- title page, no page number printed

\begin{center}
\vspace*{5cm}
{\LARGE\textsf{THE DERIVATORS}}

\vspace{1.5em}
{\large\textbf{ALEXANDRE GROTHENDIECK}}

\vspace{3em}
{\large\textbf{Chapter XI}}

\vspace{0.5em}
{\large $\Hot$\textbf{-fibrations, proper functors and smooth functors etc.}\\
\textbf{(in} $\Cat$\textbf{)}}
\end{center}

\par

% ===================== page 2 =====================
\printedpage{--}
% head: (none) -- no page number printed

\vspace*{6cm}
\noindent \par\noindent\textbf{Note by the editors of the transcribed edition.}\par
This text was deciphered and transcribed in \LaTeXe\ by M. Künzer. It was edited by M. Künzer, J. Malgoire, and G. Maltsiniotis. The transcription is as faithful as possible to the manuscript. For the few obvious corrections, or rare comments by the editors, as well as for the original numbering of the manuscript pages, typewriter characters \texttt{[typewriter]} in brackets are used. A question mark in brackets means that one is not sure of the preceding word.

\noindent This edition is provisional. Remarks, comments and corrections are welcome. Send a message to:

\noindent\textbf{[e-mail address removed]}

\hfill G. Maltsiniotis\hspace*{2.5em}

\par

% ===================== page 3 =====================
\printedpage{1}
% head: A. Grothendieck : Dérivateurs, Ch. XI, Édition des Universités de Montpellier II et Paris VII
% foot: Version 28 mai 2002

\noindent\texttt{[There are no pages 1-126]}

\msp{page 127}

\section*{1\quad $\Hot$-fibrations in $\Cat$}

\noindent\textbf{Theorem 1.} \textit{Let $f : X \longrightarrow Y$ be in} $\Cat$\textit{. Equivalent conditions:}

\begin{itemize}[leftmargin=3em]
\item[\textit{a)}] \textit{For every $Y'' \xrightarrow{p} Y'$ over $Y$, whence $X'' = X \times_Y Y'' \xrightarrow{q} X' = X \times_Y Y'$, if $p$ is a} $\Hot$\textit{-equivalence, so is $q$.}

\item[\textit{a$'$)}] \textit{As in a), with $Y'' = e$, $Y'$ having a final object.}

\item[\textit{a$''$)}] \textit{As in a), with $Y'' = e$, $Y'$ having an initial object} (\footnote{keep a$''$) and b$''$), conditions dual to a$'$), b$'$) \texttt{[a$''$) and b$''$) were crossed out]}.})\textit{.}

\item[\textit{b$'$)}] \textit{For every $y \in \Ob Y$, $X_y \longrightarrow X/y$ is a} $\Hot$\textit{-equivalence, for every $y \longrightarrow y'$ in $Y$, $X/y \longrightarrow X/y'$ is a} $\Hot$\textit{-equivalence; finally these conditions still hold for $f' : X' \longrightarrow Y'$, for every base change $Y' \longrightarrow Y$ in} $\Cat$\textit{. (In fact, it suffices to require this for the first of the two conditions, concerning $X_y \longrightarrow X/y$.)}

\item[\textit{b$''$)}] \textit{Dual of b$'$), with $X_y \longrightarrow y\backslash X$ and $y'\backslash X \longrightarrow y\backslash X$.}

\item[\textit{c)}] \textit{For every locally constant sheaf of sets $F$ on $X$, $f_*(F)$ is locally constant and is computed ``fiber by fiber'', likewise for $R^1 f_*(F)$ if $F$ is a sheaf of groups, and for the $R^i f_*(F)$ if $F$ is an abelian sheaf. Moreover, we want these conditions to remain satisfied after every base change.} (\footnote{\textit{c$'$) (homological version). 1$^o$) in terms of direct images of locally constant groupoids, and 2$^o$) in terms of the $L_i f_!(L)$} \texttt{[rather $L_i f_!(L)$]}\textit{, for $L$ a locally constant abelian sheaf.}})
\end{itemize}
\DERmark{ENV0573}


\noindent\textsc{Proof}. We have the $\pm$ tautological implications
\[
\mathrm{a} \implies \mathrm{a}' \implies \mathrm{b}' \implies \mathrm{c},
\]
and it remains to prove that c $\implies$ a. Since a is self-dual, being equivalent to a$'$, b$'$, it is also equivalent to the conditions dual to these, a$''$ and b$''$ (cf. cor. 2, page 75). We are reduced to proving that if c) is satisfied, and if we have
\[
\begin{tikzcd}[row sep=large, column sep=large]
X \arrow[d, "f"'] & X' \arrow[l, "q"'] \arrow[d, "f'"] \\
Y & Y' \arrow[l, "p"]
\end{tikzcd}
\]
\printedpage{2}
% head: A. Grothendieck : Dérivateurs, Ch. XI, Édition des Universités de Montpellier II et Paris VII
% foot: Version 28 mai 2002
cartesian in $\Cat$, if $p$ is a $\Hot$-equivalence, so is $q$. We use the Artin-Mazur criterion.

% ===================== page 4 (printed 2) =====================

\noindent 1$^o$) If $F$ is a locally constant sheaf of sets on $X$, we must prove that $H^0(X,F) \xrightarrow{\sim} H^0(X', q^*F)$.

\msp{page 128}

\noindent But
\[
H^0(X,F) \simeq H^0(Y,\ \underbrace{f_*(F)}_{\substack{\text{locally constant}\\ \text{by hypothesis c)}}}\ )
\qquad
\underbrace{\xrightarrow{\qquad\sim\qquad}}_{\substack{\text{since } p\\ \text{Hot-equivalence}}}
\qquad
H^0(Y', p^* f_*(F)),
\]
but
\[
p^* f_*(F) \xrightarrow{\ \sim\ } f'_*(q^*(F)),
\]
since $f_*(F)$ and $f'_*(q^*F)$ are computed ``fiber by fiber'' by hypothesis. Thus
\[
H^0(X,F) \simeq H^0(Y', f'_*(q^*(F))) \simeq H^0(X', q^*(F)),
\]

\hfill q.e.d.

\noindent 2$^o$) If $F$ is a locally constant sheaf of groups on $X$, we must prove % [sic] faiscau -> faisceau, misspelling
\[
H^1(X,F) \xrightarrow{\ \sim\ } H^1(X', q^*F).
\]
The fact that it is injective already follows from 1$^o$), which implies that the functor
\[
\begin{array}{rcl}
F\text{-torsors on } X & \longrightarrow & q^*(F)\text{-torsors on } X' \\
P & \longmapsto & q^*P % [sic] f^*P -> q^*P, the functor is the inverse image by q from X to X'
\end{array}
\]
\DERmark{ENV0574}

is fully faithful, by the well-known argument using the $\mathcal{H}$om sheaves. We must see that this functor is essentially surjective.

\msp{page 129}

\noindent We have a homomorphism of exact sequences of pointed sets
\[
\begin{tikzcd}[column sep=large, row sep=large,
  every matrix/.append style={name=m},
  execute at end picture={
    \coordinate (lb) at ([xshift=0.4em]m-2-3.east);
    \coordinate (lt) at ([xshift=1.8em]m-2-3.east);
    \draw (lt |- m-1-3.north) -- (lb |- m-1-3.north) -- (lb |- m-2-3.south) -- (lt |- m-2-3.south);
    \coordinate (rb) at ([xshift=0.6em]m-2-4.east);
    \coordinate (rt) at ([xshift=-0.6em]m-2-4.east);
    \draw (rt |- m-1-3.north) -- (rb |- m-1-3.north) -- (rb |- m-2-3.south) -- (rt |- m-2-3.south);
  }]
H^1(Y, f_*F) \arrow[r, "u"] \arrow[d, "\alpha"] & H^1(X,F) \arrow[r, "v"] \arrow[d, "\beta"] & \overbrace{H^0(Y, R^1 f_*(F))}^{\ni\, \eta} \arrow[r, "w"] \arrow[d, "\gamma"] & H^2(Y, f_*F) \arrow[d, "\delta"] \\
H^1(Y', f'_*F') \arrow[r, "u'"] & \underbrace{H^1(X', F')}_{\ni\, \xi'} \arrow[r, "v'"] & \underbrace{H^0(Y', R^1 f'_*(F'))}_{\ni\, \eta'} \arrow[r, "w'"] & H^2(Y', f'_*F') % [sic] H^1(Y', f'_*F) -> H^1(Y', f'_*F'), the other entries of the lower row are in F'
\end{tikzcd}
\]
\DERmark{ENV0575}

(\footnote{where we have set $F' = q^*(F)$.}), where the sheaves $f_*F$, $R^1 f_*(F)$ are locally constant by hypothesis and are computed fiber by fiber, the sheaves $f'_*(F')$ and $R^1 f'_*(F')$ also, which implies that the latter are
\printedpage{3}
% head: A. Grothendieck : Dérivateurs, Ch. XI, Édition des Universités de Montpellier II et Paris VII
% foot: Version 28 mai 2002
canonically isomorphic (by the base change arrows) to the inverse images of their counterparts $f_*(F)$, $R^1 f_*(F)$. Since $p$ is a $\Hot$-equivalence, this implies that the arrows $\alpha$, $\gamma$ are \texttt{[some]} isomorphisms. Thus if $\xi' \in H^1(X', F')$, its image $\eta'$ in $H^0(Y', R^1 f'_*(F'))$ comes from a well-determined $\eta \in H^0(Y, R^1 f_*(F))$. We must see that this $\eta$ comes from a $\xi \in H^1(X,F)$. This (if I remember correctly) is possible if and only if a certain obstruction $w(\eta)$ in $H^2(Y, f_*(F))$ is in the neutral subset of $H^2$. \textbf{NB} The stack on $Y^\wedge$ of $F$-torsors on the $X^\wedge/U$ that lift $\eta$ is an $\mathcal{F} = f_*F$-gerbe $\mathcal{G}_\eta$, and we must see that it is neutral. But there is something better to do:

% ===================== page 5 (printed 3) =====================

\msp{page 130}

\noindent We note that this gerbe $\mathcal{G}$ is locally constant (since it corresponds to a locally constant band $f_*(G)$) and that $\mathcal{G}'$ is its inverse image. But since $p$ is a $\Hot$-equivalence, it follows that the functor
\[
\underline{\Gamma}(Y, \mathcal{G}) \longrightarrow \underline{\Gamma}(Y', \mathcal{G}')
\]
is an equivalence of categories, so $\xi'$ comes (up to isomorphism) from a unique $\xi$ of $\underline{\Gamma}(Y, \mathcal{G})$, q.e.d.

\noindent 3$^o$) Let $F$ be a locally constant abelian sheaf on $X$, we must see \texttt{[that]}
\[
H^i(X,F) \xrightarrow{\ \sim\ } H^i(X', F') \qquad (F' = q^*F), \tag*{$(*)$}
\]
for every $i \in \mathbf{N}$. But the two terms are the abutment of the Leray spectral sequences, with initial terms
\[
\begin{aligned}
E_2^{pq} &\simeq H^p(Y, R^q f_*(F)) \\
E^{\prime pq}_2 &\simeq H^p(Y', R^q f'_*(F')),
\end{aligned}
\]
and $(*)$ is associated with a homomorphism of spectral sequences, which for the initial terms is
\[
H^p(Y, R^q f_*(F)) \xrightarrow{\sim} H^p(Y', R^q f'_*(F')), \tag*{$(**)$} % [sic] R^q f'_*(F) -> R^q f'_*(F'), cf. the E'_2 term just above
\]
\DERmark{ENV0576}

associated with
\[
p^*(R^q f_*(F)) \longrightarrow R^q f'_*(F'). \tag*{$(*{*}*)$}
\]
But by hypothesis the $R^q f_*(F)$, $R^q f'_*(F')$ are locally constant and are computed fiber by fiber,

\msp{page 131}

\noindent whence it follows that the base change homomorphisms $(*{*}*)$ are \texttt{[some]} isomorphisms, and since we are dealing with locally constant sheaves and $p$ is a $\Hot$-equivalence, it follows that $(**)$ is \texttt{[an]} isomorphism, hence also $(*)$ on the abutments, q.e.d.

\noindent\textbf{Definition 1.} Under the equivalent conditions of theorem 1, we say that $f$ is a $\Hot$-\textit{fibration.}


\endgroup
% END UNIT Der-ChXI-890/Der-ChXI-890_p001-005.tex
% BEGIN UNIT Der-ChXI-890/Der-ChXI-890_p006-008.tex
\begingroup\makeatletter%
\derleqno%
\newcommand{\printedpage}[1]{\marginpar{\footnotesize p.~#1}}%
\newcommand{\msp}[1]{\marginpar{\footnotesize [#1]}}%
\newcommand{\Hot}{\mathrm{Hot}}%
\newcommand{\Ch}{\underline{\mathrm{Ch}}}%
\newcommand{\Ob}{\mathrm{Ob}\,}%
\newcommand{\Cat}{\mathrm{Cat}}%
\makeatother\let\printedpage\derpp\let\msp\dermsp%
\let\mspin\DERmspin%


\printedpage{4}
% head: A. Grothendieck : Dérivateurs, Ch. XI, Édition des Universités de Montpellier II et Paris VII
% foot: Version 28 mai 2002

\textbf{Proposition 1.} \textit{This notion is stable under base change (trivial), under composition (formal from a)). For $f$ to be a \textup{Hot}-fibration, it is necessary and sufficient that $f^{o}$ \textup{\texttt{[rather}} $f^{o}$\textup{\texttt{]}} be one.}
\DERmark{ENV0577}


\textbf{Proposition 2.} \textit{Let $f : X \longrightarrow S$ be a \textup{Hot}-fibration. Then $s \longmapsto \mathrm{hot}(X_s)$ is a local system on $S$, with values in \textup{Hot}. Thus the $X_s$ for $s$ in the same connected component of $S$ are \textup{Hot}-isomorphic to one another (but not canonically, if this component is not \textup{1}-connected).}

% (a vertical change bar is printed in the left margin beside Théorème 2)
\textbf{Theorem 2.} {\itshape Let
\begin{tikzcd}[ampersand replacement=\&, baseline=(\tikzcdmatrixname-1-1.base), column sep=small, row sep=small]
X' \arrow[rr, "f"] \arrow[dr, "p'"'] \& \& X \arrow[dl, "p"] \\
\& S \&
\end{tikzcd}
be a commutative diagram in \textup{Cat}, with $p$, $p'$ \textup{Hot}-fibrations, and let $S_0 \subseteq S$ be a subset of $S$ such that $S_0 \longrightarrow \pi_0(S)$ is surjective. For $f$ to be a \textup{Hot}-equivalence, it is necessary and sufficient that the $f_s : X'_s \longrightarrow X_s$ $(s \in S_0)$ be so.} % [sic] f_s : X_s -> Y_s -> f_s : X'_s -> X_s, the diagram has X' and X (cf. f_s : X'_s -> X_s p. 6)
\DERmark{ENV0578}


a) Sufficiency is proved as before, with the Artin-Mazur criterion.

\msp{page 132}

Let us note that the hypothesis on $S_0$ ensures that if $F$, $F'$ are two locally constant sheaves of sets on $S$, and $u : F \longrightarrow F'$ a homomorphism such that $u_s : F_s \xrightarrow{\ \sim\ } F'_s$ for $s \in S_0$, then $u$ is \texttt{[an]} isomorphism. % [sic] s \in S -> s \in S_0, otherwise the hypothesis on S_0 plays no role
\DERmark{ENV0579}


1$^{o}$) Let $F$ be a locally constant sheaf of sets on $X$, $F'$ its inverse image on $X'$, we must prove
\[
H^0(X,F) \ \xrightarrow{\ \sim\ }\ H^0(X',F').
\]
Now
\[
H^0(X,F) \simeq H^0(S, \underbrace{p_*(F)}_{\substack{\text{locally}\\ \text{constant on } S}}) \ \longrightarrow\ H^0(S, p'_*(F')) \simeq H^0(X',F'),
\] % [sic] H^0(Y, f_*(F)), H^0(Y, f'_*(F')) -> H^0(S, p_*(F)), H^0(S, p'_*(F')), the brace says "localement constant sur S" and (*) uses p_*, p'_*
\DERmark{ENV0580}\DERmark{ENV0581}

we must prove that
\[
p_*(F) \ \xrightarrow{\ \sim\ } p'_*(F') \tag*{$(*)$}
\]
is \texttt{[an]} isomorphism. Now by the hypothesis on $p$, $p'$, the two sheaves are locally constant and are computed fiber by fiber, from which it follows that $(*)$ is an isomorphism for the fibers at $s \in S_0$, hence an isomorphism everywhere. OK.

2$^{o}$) It already follows from this, if $F$ is a sheaf of \textit{groups,} that the functor
\[
(F\text{-torsors on } X) \ \longrightarrow\ (F'\text{-torsors on } X')
\]
\DERmark{ENV0582}

is fully faithful, hence $H^1(X,F) \longrightarrow H^1(X',F')$ \texttt{[is]} injective; it remains to prove that this functor is essentially surjective. Because of full faithfulness, this is \textit{local on $S$} (\textbf{NB} We see directly that the hypothesis made on the $s \in S_0$ implies the same hypothesis for every $s \in S$, using the fact that $p$, $p'$ are fibrations) (\footnote{The hypotheses made are indeed local on $S$ (at least if $S_0 = S$).}). Thus we are reduced to the case where $S = S/s_0$ has a final object $s_0$.

\printedpage{5}
% head: A. Grothendieck : Dérivateurs, Ch. XI, Édition des Universités de Montpellier II et Paris VII
% foot: Version 28 mai 2002

\msp{page 133}

We then have
\[
\begin{aligned}
H^1(X,F) &\ \simeq\ H^0(S, R^1p_*(F)) \\
H^1(X',F') &\ \simeq\ H^0(S, R^1p'_*(F')),
\end{aligned}
\]
and we see again that
\[
R^1p_*(F) \ \longrightarrow\ R^1p'_*(F')
\]
is \texttt{[an]} isomorphism (for it is computed fiber by fiber, and we use the hypothesis on the $f_s$), hence induces an isomorphism
\[
H^1(X,F) \ \longrightarrow\ H^1(X',F').
\]

3$^{o}$) If $F$ is even abelian, we must prove
\[
H^i(X,F) \ \xrightarrow{\ \sim\ }\ H^i(X',F')
\]
for every $i$, and this again follows from the Leray spectral sequence for $(p,F)$ and for $(p',F')$.

b) Conversely, suppose that $f$ is a Hot-equivalence, and let us prove that the $f_s : X'_s \longrightarrow X_s$ are Hot-equivalences. I am going to use the

% (a vertical change bar is printed in the left margin beside the Lemme)
\textbf{Lemma.} \textit{Let $S \in \Ob\Cat$, $S$ \textup{0}-connected. Then there exists an $\overline{S} \in \Ob\Cat$, $\overline{S}$ \textup{Hot}-aspheric, and $\overline{S} \longrightarrow S$ a \textup{Hot}-fibration.}

Let us see how this implies theorem 2 (converse part). We may suppose $S$ 0-connected, and we make the base change by $\overline{S} \longrightarrow S$

\msp{page 134}

as in the lemma:
\[
\begin{tikzcd}[column sep=small, row sep=small]
X' \arrow[rr, "f"] \arrow[dr, "p'"'] & & X \arrow[dl, "p"] & & & & \\
& S & & & \overline{X'} \arrow[rr, "\overline{f}"] \arrow[dr, "\overline{p'}"'] & & \overline{X} \arrow[dl, "\overline{p}"] \\
& & & & & \overline{S} \arrow[ullll, "\varphi"] &
\end{tikzcd}
\quad .
\]

Then $\overline{p'}$ and $\overline{p}$ are Hot-fibrations, and since $\varphi$ is a Hot-fibration, base change by $\varphi$ transforms the Hot-equivalence $f$ into a Hot-equivalence $\overline{f}$. Consider then $\overline{s} \in \Ob\overline{S}$, and the diagram
\[
\begin{tikzcd}
\overline{X'} \arrow[r, "\overline{f}"] & \overline{X} \\
\overline{X'}_{\overline{s}} \arrow[u, hook, "i'"] \arrow[r, "\overline{f}_{\overline{s}}"] & \overline{X}_{\overline{s}} \arrow[u, hook, "i"']
\end{tikzcd}
\qquad .
\]

\printedpage{6}
% head: A. Grothendieck : Dérivateurs, Ch. XI, Édition des Universités de Montpellier II et Paris VII
% foot: Version 28 mai 2002

Since $\{\overline{s}\} \longrightarrow \overline{S}$ is a Hot-equivalence, the arrow $i$ deduced from it by the base change $\overline{p} : \overline{X} \longrightarrow \overline{S}$ is one as well, since $\overline{p}$ is a Hot-fibration, and likewise for $i'$. Thus in the preceding commutative square, $\overline{f}$, $i$ and $i'$ are Hot-equivalences, hence so is $\overline{f}_{\overline{s}}$. But if $s = \varphi(\overline{s})$, $\overline{f}_{\overline{s}}$ is identified with $f_s : X'_s \longrightarrow X_s$. Besides, it is immediate that $\varphi$, being a Hot-fibration

\msp{page 135}

with nonempty source, and connected target, is surjective. (\textbf{NB} all the fibers of one and the same connected component of $S$ are Hot-equivalent, hence if one is nonempty, they all are.) This completes the proof, modulo the lemma.

It remains to prove it. Since $S \neq \emptyset$ \texttt{[rather} $S \neq \emptyset$\texttt{]}, I can choose an ``origin''
\DERmark{ENV0583}

\[
s_0 \in S,
\]
and look at the category
\[
\Ch_{\infty}(S; s_0, -)
\]
of infinite paths in $S$, with origin $s_0$, fiber of $\Ch_{\infty}(S)$ under the source functor
\[
\sigma : \Ch_{\infty}(S) \ \longrightarrow\ S.
\]
We know that the fibers of $\sigma$ are Hot-aspheric (VII), and the category constructed here is none other than (with the notations of VII, p.~152) than the category $\overline{S'}$ which enters into the canonical factorization of
\[
f : \{s_0\} = S' \ \longrightarrow\ S
\]
into
\[
\{s_0\} = S' \ \underbrace{\longrightarrow}_{\text{Hot-equivalence}}\ \overline{S'} \ \xrightarrow{\ \overline{f}\ }\ S,
\]
where here $\overline{f}$ is interpreted simply as the target functor (i.e.\ the functor induced on

\msp{page 136}

$\Ch_{\infty}(S; s_0, -)$ by the target functor
\[
\beta : \Ch_{\infty}(S) \ \longrightarrow\ S \ \ ).
\]
We proved in loc.\ cit.\ that $\overline{f}$ is a Hot-fibration (``proper'' and ``smooth''), but in a context which seemed to appeal tacitly to the theory of the derivator HOT. So to put the ``topological'' results of loc.\ cit.\ on a solid basis, we must develop the relevant notions of proper functor, smooth functor, without reference to a HOT theory. This is what I had in fact already done in 1983, in the days of Pursuing Stacks, but without a proper write-up on this subject. So I must now return to these notions, in a somewhat systematic way. Besides, I would have trouble constructing a single Hot-fibrant functor (except for the trivial examples, such as isomorphisms, or the products $Y \times Z \xrightarrow{\mathrm{pr}_1} Y$), if I did not have at my disposal a theory of smooth functors and proper functors (\footnote{I am exaggerating a little all the same -- after all it must be fairly obvious that the Thomason-Kan functors are Hot-fibrations (but there are still a few too few of them!).}). % [sic] dû mal -> du mal, "avoir du mal à" (partitive article, no accent)


\endgroup
% END UNIT Der-ChXI-890/Der-ChXI-890_p006-008.tex
\clearpage\setcounter{footnote}{0}\phantomsection\addcontentsline{toc}{section}{Chapter XII, 1–4}\markboth{Chapter XII, 1–4}{Chapter XII, 1–4}
% BEGIN UNIT Der-ChXIIa-012/Der-ChXIIa-012_p001-005.tex
\begingroup\makeatletter%
\newcommand{\printedpage}[1]{\marginpar{\footnotesize p.~#1}}%
\newcommand{\msp}[1]{\marginpar{\footnotesize [#1]}}%
\newcommand{\Cat}{\mathrm{Cat}}%
\newcommand{\Dl}{\boldsymbol{\Delta}}%
\newcommand{\fl}{\longrightarrow}%
\newcommand{\Ob}{\mathrm{Ob}\,}%
\newcommand{\Hot}{\mathrm{Hot}}%
\newcommand{\hot}{\mathrm{hot}}%
\newcommand{\uHom}{\underline{\mathrm{Hom}}}%
\newcommand{\id}{\mathrm{id}}%
\newcommand{\pr}{\mathrm{pr}}%
\makeatother\let\printedpage\derpp\let\msp\dermsp%
\let\mspin\DERmspin%


\printedpage{--}
% head: (none; no running head and no page number printed on this page)

\begin{center}
{\Large \textsf{THE DERIVATORS}}

\bigskip

\textbf{ALEXANDRE GROTHENDIECK}

\bigskip\bigskip

\textbf{Chapter XII}

\medskip

\textbf{Characterization of $W_\infty$. $W$-proper,\\
$W$-smooth functors etc. Amalgamated sums and\\
$W$-cocartesian squares in $\Cat$.}
\end{center}

\par

\printedpage{--}
% head: (none; no running head and no page number printed on this page)

\par\noindent\textbf{Note by the editors of the transcribed edition.}\par
This text was deciphered and transcribed in \LaTeXe{} by M. Künzer. It was edited by M. Künzer, J. Malgoire, and G. Maltsiniotis. The transcription is as faithful as possible to the manuscript. For the few obvious corrections, or rare comments by the editors, as well as for the original numbering of the manuscript pages, typewriter characters \texttt{[typewriter]} in brackets are used. A question mark in brackets means that one is not sure of the preceding word.

This edition is provisional. Remarks, comments and corrections are welcome. Send a message to:

\textbf{[e-mail address removed]}

\hfill G. Maltsiniotis\hspace*{2em}

\par

\printedpage{1}
% head: A. Grothendieck : Dérivateurs, Ch. XII, Édition des Universités de Montpellier II et Paris VII
% foot: Version 7 décembre 2002

\msp{page 1}

\hfill 19.11.90

\section*{I\quad Characterization of $W_\infty$.}

I would like to characterize $W_\infty$ as a subset of $\mathrm{Fl}(\Cat)$, by its properties as a localizer, and without appealing to knowledge of notions already using $W_\infty$, such as properness etc. I would like at least some properties of a localizer $W$, familiar for $W_\infty$, which are so strong that they imply $W \subseteq W_\infty$. Of course, in the proofs I allow myself to use everything I know about $W_\infty$.

Of course, at least a form of saturation will be needed, perhaps a very weak one, I presume the following will suffice (\footnote{(no, I had to strengthen it in 2 [?]).}).

\begin{itemize}[leftmargin=*]
\item[\textbf{W(1)}] (Saturation axiom.) a) The identities are in $W$, b) if $f$, $g$ \texttt{[are]} composable and two among $f$, $g$, $gf$ are in $W$, the third too, c) if $fg \in W$, $gf \in W$, then $f, g \in W$.
\end{itemize}

A particularly crucial property, which already distinguishes $W_\infty$ from the set of homotopisms (a notion too strict in many respects), is the axiom I called L4 bis:

\begin{itemize}[leftmargin=*]
\item[\textbf{W(2)}] (Localization axiom.) Let $f : X \fl Y$ be an arrow in $\Cat/S$, such that $\forall\ s \in \Ob S$, $f_s : X/s \fl Y/s$ is in $W$ (we shall say that $f$ is in $W$ ``locally over $S$\,''). Then $f \in W$.
\end{itemize}

\texttt{[The axioms W(1) etc. are also called} \textcircled{1} \texttt{etc. in the original.]}

In the case $Y = S$, this gives

\textbf{Corollary:} \textit{If $f : X \fl Y$ is $W$-aspheric, i.e. $X/y \fl Y/y$} \texttt{[rather $X/y \fl Y/y$]} \textit{in $W$ for every $y \in \Ob Y$, then $f \in W$} (\footnote{One can consider the validity of the corollary as a weakened variant W(2 bis) of W(2) above. In proposition 1, one uses only W(2 bis), a weak form of W(2).}).
\DERmark{ENV0636}


\textbf{Proposition 1.} \textit{If $W$ satisfies} W(1) \textit{and} W(2)\textit{, the following properties of $W$ are equivalent:}

\begin{itemize}[leftmargin=*]
\item[\textit{a)}] \textit{For every $X$ in $\Cat$, $X \times \Dl^1 \xrightarrow{\pr_1^X} X$ is in $W$ (i.e. $\Dl^1 \fl e$ is ``universally in $W$\,'').}

\item[\textit{b)}] \textit{If $X$ in $\Cat$ has a final object, then $X$ is $W$-aspheric, i.e. $X \fl e$ is in $W$.}

\item[\textit{c)}] \textit{Every homotopism is in $W$.}
\end{itemize}

\par

\printedpage{2}
% head: A. Grothendieck : Dérivateurs, Ch. XII, Édition des Universités de Montpellier II et Paris VII
% foot: Version 7 décembre 2002

\textsc{Proof.} b) $\Longrightarrow$ a) (\footnote{uses W(1) only in weak form, without the saturation condition c). NB Only b) $\Longrightarrow$ c) uses W(2 bis), the implications a) $\Longrightarrow$ c) $\Longrightarrow$ b) are valid under the sole hypothesis W(1).}), for to see that $\pr_1 \in W$, applying the corollary above, we must prove that $\forall\ x \in X$
\[ (X \times \Dl^1)/x \fl X/x \tag*{$(*)$} \]
is in $W$. Now
\[ (X \times \Dl^1)/x \simeq X/x \times \Dl^1, \]
\[
\begin{tikzcd}
\Dl^1 \arrow[d] & X \times \Dl^1 \arrow[l] \arrow[d] & X/x \times \Dl^1 \arrow[l] \arrow[d] \\
e & X \arrow[l] & X/x \arrow[l]
\end{tikzcd}
\]
\texttt{[the vertical arrow on the right is added]} and the two categories $X/x$, $X/x \times \Dl^1$ are $W$-aspheric since they have a final object, so $(*)$ is in $W$ by weak saturation W(1).

a) $\Longrightarrow$ c). I recall that
\[ f : X \fl Y \]
\msp{page 2}
is a homotopism if there exists
\[ g : Y \fl X \]
such that
\[ gf \sim \id_X, \quad fg \sim \id_Y, \tag*{$(*)$} \]
where the relation $(\sim)$ denotes homotopy in $\uHom(X,X)$ and $\uHom(Y,Y)$ respectively (defined by paths going from one element to the other). Let us prove the

\textbf{Lemma.} \textit{Suppose that $W$ satisfies} W(1) \textit{and the} a) \textit{of the proposition, i.e. $X \times \Dl^1 \fl X$ in $W$ for every $X$ in} $\Cat$. \textit{Then if $X$, $Y$} \texttt{[are]} \textit{in} $\Cat$ \textit{and $f, g$ in} $\uHom(X,Y)$ \textit{are homotopic, we have}
\[ \hot_W(f) = \hot_W(g). \]
\textit{Moreover, $f \in W$ if and only if $g \in W$.}

It suffices to prove this when there exists a homomorphism $f \fl g$. This can be interpreted as the existence of an arrow
\[ X \times \Dl^1 \xrightarrow{F} Y \]
whose restrictions to $X \times \{0\}$ and $X \times \{1\}$ respectively are $f$, $g$. To prove that $\hot_W(f) = \hot_W(g)$, it suffices to prove that
\[ i_0, i_1 : X \fl X \times \Dl^1 \qquad x \longmapsto (x,0),\ x \longmapsto (x,1) \]
\printedpage{3}
% head: A. Grothendieck : Dérivateurs, Ch. XII, Édition des Universités de Montpellier II et Paris VII
% foot: Version 7 décembre 2002
are equal in $\Hot_W$, since
\[ f = F \circ i_0, \quad g = F \circ i_1. \]
Now this follows from the fact that they are sections of $X \times \Dl^1$ over $X$, and that in $\Hot_W$ the projection $X \times \Dl^1 \xrightarrow{p} X$ is \texttt{[an]} isomorphism (hence $i_0 = i_1 = p^{-1}$ in $\Hot_W$). Moreover, the diagram
\[
\begin{tikzcd}
X \arrow[r, "i_0"] \arrow[dr, "f"'] & X \times \Dl^1 \arrow[d, "F"] \\
& Y,
\end{tikzcd}
\]
where $i_0 \in W$, shows that $f \in W$ if and only if $F \in W$, and similarly $g \in W$ if and only if $F \in W$, hence $f \in W \Longleftrightarrow g \in W$.

Let us return to $(*)$, we see that in $\Hot_W$ $g \circ f = \id_X$, $f \circ g = \id_Y$, so $f$ and $g$ are $\Hot_W$-isomorphisms inverse to one another. On the other hand, $gf \in W$, $fg \in W$, since $\id_X, \id_Y \in W$. So we win, if we have formulated W(1) strongly enough. We have thus proved a) $\Longrightarrow$ c).

c) $\Longrightarrow$ b) is trivial, since if $X$ has a final object, $X \fl e$ is a homotopism.

\msp{page 3}

We must of course in what follows suppose that $W$ satisfies the equivalent conditions of the proposition. We shall take e.g. the form

\begin{itemize}[leftmargin=*]
\item[\textbf{W(3)}] (Final element axiom.) If $X$ has a final object, then $X$ is $W$-aspheric.
\end{itemize}

\textbf{Corollary 1.} \textit{For $f : X \fl Y$ to be $W$-aspheric, it is necessary and sufficient that the $X/y$ be $W$-aspheric} ($y \in \Ob Y$).

For in order that
\begin{tikzcd}[baseline=(\tikzcdmatrixname-1-1.base), column sep=small, row sep=small]
X/y \arrow[rr] \arrow[dr] & & Y/y \arrow[dl] \\
& e &
\end{tikzcd}
be in $W$, since $Y/y$ is $W$-aspheric (final object), it is necessary and sufficient that $X/y$ be so as well.

\textbf{Corollary 2.} \textit{If $X$ in} $\Cat$ \textit{has an initial object, it is $W$-aspheric.} (Indeed, $X \fl e$ is a homotopism, and we apply prop. 1.) (\footnote{This is thus the statement dual to the axiom W(3) of the final object.})

I shall put down for the record

\textbf{Corollary 3.} \textit{Every homotopism in} $\Cat$ \textit{is in $W$. If $f, g : X \rightrightarrows Y$ in} $\Cat$ \textit{are homotopic, then they are equal in} $\Hot_W$\textit{, and $f$ is in $W$ if and only if $g$ is.}

\textbf{Proposition 2.} Let $f : X \fl Y$ be in $\Cat$. If $f \in W$, $f^\circ \in W$ \texttt{[rather $f^\circ \in W$].}
\DERmark{ENV0637}



\endgroup
% END UNIT Der-ChXIIa-012/Der-ChXIIa-012_p001-005.tex
% BEGIN UNIT Der-ChXIIa-012/Der-ChXIIa-012_p006-010.tex
\begingroup\makeatletter%
\newcommand{\printedpage}[1]{\marginpar{\footnotesize p.~#1}}%
\newcommand{\msp}[1]{\marginpar{\footnotesize [#1]}}%
\newcommand{\Ch}{\underline{\mathrm{Ch}}}%
\makeatother\let\printedpage\derpp\let\msp\dermsp%
\let\mspin\DERmspin%


\printedpage{4}
% head: A. Grothendieck : Dérivateurs, Ch. XII, Édition des Universités de Montpellier II et Paris VII
We use the proof of I, p. 75, with the help of the category called $\Phi(X)$ ($\mathrm{Ob}\,\Phi(X) = \mathrm{Fl}(X)$, but with a somewhat peculiar notion of arrow in $\Phi(X)$). We find that $f \in W \iff \Phi(f) \in W$, and likewise $f^\circ \in W \iff \Phi(f) \in W$.

NB. In the proof of proposition 2, as in that of proposition 1 (for the implication b) $\Longrightarrow$ a)), we used the localization property W(2) only in the weak form of the corollary, i.e. in the form $f$ $W$-aspheric $\Longrightarrow f \in W$.

\textbf{Corollary} (of prop. 2): \textit{Let $W \subseteq \mathrm{Fl}(\mathrm{Cat})$ satisfy the axiom} W(1) (saturation). \textit{Then the conjunction of the axioms} W(2), W(3) \textit{is equivalent to that of the dual axioms} W($2'$), W($3'$) (\footnote[5]{Likewise if in this statement we replace W(2) by W(2 bis).}).

It suffices to prove that 1,2,3 $\Longrightarrow$ $2'$,$3'$.

\msp{page 4}

For $3'$ \texttt{[rather} $3'$\texttt{]}, this is corollary 2 above. For $2'$ \texttt{[rather} $2'$\texttt{]}, this follows at once from proposition 2. (NB This statement would remain valid, if we would have taken W(2) in the weak form only, cf. the NB on the previous page.)
\DERmark{ENV0638}\DERmark{ENV0639}


I still have to put in an axiom of the form L5, L6 (cf. VI, p. 43-44). But the statement of these axioms uses the notion of Hot-fibered objects, or of perfect functors, i.e. calls on knowledge of $W_\infty$, which I would like to avoid (if I can). I must therefore introduce the notion of $W$-fibration, on the model of what I did in the context of derivators, in I (cf. p. 111-126). We introduce the notion of \textbf{D}-fibration on page 117 (prop. 5), by \texttt{[some]} equivalent conditions, some of which make sense in terms of the datum of $W$ alone. So let
\[ f : X \longrightarrow Y \]
be in Cat and consider the conditions

\begin{enumerate}
\item[a)] If $Y'' \longrightarrow Y'$ in $\mathrm{Cat}/Y$ is in $W$, then so is $X'' \longrightarrow X'$ (base change by $X \longrightarrow Y$ takes $W$ into $W$).
\item[a$'$)] As in a), but with $Y''$, $Y'$ $W$-aspheric.
\item[a$_1$)] As in a), in the particular case where $Y'$ has a final object and $Y''$ is the one-point category.
\item[a$_2$)] Dual of a$_1$), with an initial object of $Y'$.
\end{enumerate}

We thus have
\[
\begin{tikzcd}[row sep=large, column sep=1.4em]
 & & \mathrm{a} \arrow[d, Rightarrow] & \text{ultrastrong $W$-fibration} & \\
 & & \mathrm{a}' \arrow[dl, Rightarrow] \arrow[dr, Rightarrow] & \text{strong $W$-fibration} & \\
\text{left $W$-fibration} & \mathrm{a}_1 & & \mathrm{a}_2 & \text{right $W$-fibration}
\end{tikzcd}
\]
% foot: Version 7 décembre 2002
\printedpage{5}
% head: A. Grothendieck : Dérivateurs, Ch. XII, Édition des Universités de Montpellier II et Paris VII
and in the case $W = W_\infty$ all these conditions are equivalent. I have a mind \texttt{[to]} take the strongest of these conditions, but I must give myself some means of finding $W$-fibrations in so strong a sense. For example I would like to know that if $X$
\msp{page 5}
is in Cat, then $\Ch_\infty(X) \longrightarrow X \times X$ or $\Ch(X) \longrightarrow X \times X$ is a $W$-fibration. In both cases, and (say) for $W_\infty$, I prove it by proving that $p$ is both smooth, and proper. So we would have to go over the notions of smoothness and properness again in the context of the datum of a $W$.

If I remember correctly, property a$_1$) \texttt{[rather a}$_1$\texttt{)]} was something fairly easily checked, it sufficed to suppose that the
\DERmark{ENV0640}

\[ X_y \longrightarrow X/y \qquad \text{and} \qquad X/y \longrightarrow X/y' \]
(for $y \in \mathrm{Ob}\,Y$ resp. $u : y \longrightarrow y'$ in $\mathrm{Fl}\,Y$) are in $W$, and that this remains true under every base change - which amounts to condition a$_1$) on $f$. We could therefore formulate a condition on $W$ as follows:

\begin{enumerate}[leftmargin=*]
\item[\textbf{W(4)}] (Axiom of $W$-fibrations (\footnote[6]{See sufficient weakened variant, in the NB p. 10. See also [?] of variants pages 13, 14 below.}).) Let $f : X \longrightarrow Y$ in Cat satisfy the $W$-fibration condition a$_1$) above (p. 4), then it even satisfies a), i.e. base change by $f$ takes arrows in $W$ to arrows in $W$ (\footnote[7]{i.e. a left $W$-fibration is an ultrastrong fibration.}).
\end{enumerate}

Thanks to proposition 2, it follows that the dual condition a$_2$) on $f$ implies the strong $W$-fibration condition, so the axiom W(4), apparently asymmetric, also means that the four fibration conditions on $f$ are equivalent, it is therefore in fact symmetric.

\textbf{NB} I do not see, without additional axioms on $W$, how to pass from a$_1$ to a$'$ alone. The difficulty is to show, by means of a$_1$ alone, that if $Y$ is $W$-aspheric, the inclusion $X_y \longrightarrow X$ (where $y \in \mathrm{Ob}\,Y$) is in $W$. In the context of derivators \textbf{D}, this could be seen by \textbf{D}-cohomological means.

This very strong axiom suffices to imply that if $X$, $Y$ are $W$-fibered over $S$, and if $f : X \longrightarrow Y$ is an $S$-morphism which is in $W$, then $f_s : X_s \longrightarrow Y_s$ is in $W$ ($\forall\ s \in S$). Let us at least show the converse for now.

\msp{page 6}

\textbf{Proposition 3.} \textit{Let $X$, $Y$ be in} Cat$/S$\textit{, $f : X \longrightarrow Y$ an $S$-arrow. Suppose ($X$, $Y$ $W$-fibered) that $\forall\ s \in S$, $f_s : X_s \longrightarrow Y_s$ is in $W$, then $f \in W$} (\footnote[8]{\textit{Without} axiom W(4). This is the first time we need W(2) in its strong form, to the exclusion of W(2 bis).}).

We have to use the localization axiom, hence look at the $f/s$ induced by $f$
% foot: Version 7 décembre 2002
\printedpage{6}
% head: A. Grothendieck : Dérivateurs, Ch. XII, Édition des Universités de Montpellier II et Paris VII
\[
\begin{tikzcd}[row sep=large]
X/s \arrow[r, "f/s"'] & Y/s \\
X_s \arrow[u, "i_s^X"] \arrow[r, "f_s"] & Y_s \arrow[u, "i_s^Y"'] \quad .
\end{tikzcd}
\]

Condition a$_1$ suffices to imply that $i_s^X$ and $i_s^Y$ are in $W$, so $f/s$ is in $W$ if and only if $f_s$ is. We did not have to use axiom W(4), and it sufficed to suppose that $X$, $Y$ are ``$W$-cohomological'' fibered objects (or left $W$-fibered) over $S$, or dually, right $W$-fibered objects.

Moreover:

\textbf{Corollary} (\footnote[9]{Without axiom W(4).}). \textit{Let $S_0 \subseteq \mathrm{Ob}\,S$ be a subset of $S$ such that $S_0 \longrightarrow \pi_0(S)$ surjective. Then in proposition 3 it suffices to suppose $f_s \in W$ for $s \in S_0$.}

It suffices to see that if $u : s \longrightarrow s'$ is an arrow in $S$, then $f_s \in W \iff f_{s'} \in W$. We use the fact that $f_s \in W \iff f/s \in W$, and likewise for $f_{s'}$. So to see that $f/s \in W \iff f/s' \in W$. Now we have
\[
\begin{tikzcd}[row sep=large]
X/s \arrow[r, "f/s"] \arrow[d, "X/u"'] & Y/s \arrow[d, "Y/u"] \\
X/s' \arrow[r, "f/s'"] & Y/s' \quad ,
\end{tikzcd}
\]
where $X/u$ and $Y/u$ are in $W$, whence the result.

But let us return to the converse: if $f \in W$, then do we have $f_s \in W$? \texttt{[rather} $f_s \in W$\texttt{?]} There is an argument to this effect in VI, p. 23, which I am going to take up again. The idea is this. We consider the connected component $S_0$ of $S$ \texttt{[rather} $S_0$ \texttt{of} $S$\texttt{]} containing $s$, and we construct a
\DERmark{ENV0641}\DERmark{ENV0642}

\msp{page 7}
$W$-fibered object (in the \textit{strong sense} a))
\[ S' \longrightarrow S_0 \subseteq S \]
with $S'$ \textit{$W$-aspheric.} Since $X \longrightarrow Y$ over $S$ is in $W$, the same will therefore hold for $X' \longrightarrow Y'$ over $S'$. Since the fibers of $X'$, $Y'$ at $s'$ are identified with those of $X$, $Y$ at the image of $s'$ in $S$, we are reduced to the case where $S$ is $W$-aspheric. But then we have
\[
\begin{tikzcd}[row sep=large]
X \arrow[r, "f"] & Y \\
X_s \arrow[u, "i_s^X"] \arrow[r, "f_s"] & Y_s \arrow[u, "i_s^Y"'] \quad ,
\end{tikzcd}
\]
% foot: Version 7 décembre 2002
\printedpage{7}
% head: A. Grothendieck : Dérivateurs, Ch. XII, Édition des Universités de Montpellier II et Paris VII
where $i_s^X$, $i_s^Y$ are in $W$, since obtained from $\{s\} \longrightarrow S$ (which is in $W$, $S$ being $W$-aspheric) by the $W$-fibrant base changes $X \longrightarrow S$, $Y \longrightarrow S$. Hence $f \in W \implies f_s \in W$.

NB We used the fibration property for $X$, $Y$ only in the sense a$'$) (\footnote[10]{And even, in a$'$) we could require that $Y'' = e$, $Y'$ $W_\omega$-aspheric, where $W_\omega$ \texttt{[is]} defined on page 14, $W_\omega \subseteq W \cap W_\infty \ldots$}).

We must therefore prove, if we can, the

\textbf{Lemma} (?) (\footnote[11]{Is proved by means of axiom W(4) (+ some small later verifications).}). \textit{Let $S$ be a} $0$\textit{-connected object of} Cat\textit{. Then there exists a $W$-aspheric object $S'$ of} Cat\textit{, and a $W$-fibration (in the} strong \textit{sense of a), page 4) $S' \longrightarrow S$.}

The natural idea is to take, for fixed $s_0 \in \mathrm{Ob}\,S$,
\[ S' = \Ch(S; s_0, -) \qquad \text{or} \qquad S' = \Ch_\infty(S; s_0, -), \]
and the arrow
\[ \beta : c \longmapsto \text{target } c : S' \longrightarrow S. \]
We know that $S'$ is $W_\infty$-aspheric and that $\beta$ is $W_\infty$-fibrant (and even perfect). We must see that $S'$ is also $W$-aspheric, and $\beta$ $W$-fibrant. We must therefore go over the proof again in the case of $W_\infty$, and prove that it transforms to $W$, % [sic] aphérique -> asphérique, lettre s omise
\DERmark{ENV0643}

\msp{page 8}
given the axioms needed on $W$. I have a mind to look first at the choice $S' = \Ch(S; s_0, -)$, without passing to a direct limit which would force me to put a limit axiom on $W$. The delicate and essential homotopy properties of $\Ch(S)$ are established in VII, p. 157-169, namely that in the diagram
\[ S \xrightarrow{i_S} \Ch(S) \xrightarrow{p_S} S \times S, \]
$i_S \in W_\infty$ and $p_S$ is a $W_\infty$-fibration (in fact, \textit{proper} and \textit{smooth} for $W_\infty$) (\footnote[12]{In the statement of the ``essential properties'' of $\Ch(S)$, we can work with $W_\omega$, the smallest subset of Fl(Cat) satisfying axioms W(1), W(2 bis), W(3) (cf. page 14), so $W_\omega \subseteq W \cap W_\infty$.}). The fact that $i_S \in W_\infty$ \texttt{[rather} $i_S \in W_\infty$\texttt{]} implies that the source and target functors
\DERmark{ENV0644}

\[ \begin{tikzcd} \Ch(S) \arrow[r, shift left, "\sigma"] \arrow[r, shift right, "\beta"'] & S \end{tikzcd} \qquad \sigma = \mathrm{pr}_1 \circ p_S,\ \beta = \mathrm{pr}_2 \circ p_S \]
are in $W_\infty$ (since $\sigma i_S = \beta i_S = \mathrm{id}_S$). Since these are Hot-fibrations, it follows that the fibers $\Ch(S; s, -)$ resp. $\Ch(S; -, s)$ are Hot-aspheric. But even granting that $\sigma$, $\beta$ are in $W$ and are $W$-fibrations, I cannot conclude right away that the fibers are $W$-aspheric - there would be a vicious circle! So I need a direct proof.

Consider the functor
\[
\begin{tikzcd}[row sep=huge]
\Ch(S; s, -) \arrow[d, "p"'] \\
\underline{Ty} \quad ,
\end{tikzcd}
\]
% foot: Version 7 décembre 2002
\printedpage{8}
% head: A. Grothendieck : Dérivateurs, Ch. XII, Édition des Universités de Montpellier II et Paris VII
it is a Cat-fibrant functor, hence $W$-smooth (I will have to develop the notion of $W$-smoothness (\footnote[13]{$f : X \longrightarrow Y$ is $W$-smooth if the $X_y \longrightarrow y\backslash X$ are $W$-aspheric.}) carefully, I will assume for the moment the familiar formal properties, and will justify them later). Because of $W$-smoothness, to prove $p \in W$, it suffices to
\msp{page 9}
prove that the fibers of $p$ are $W$-aspheric (\footnote[14]{Here we use W(2 bis), not $W(2)$.}). Now the fibers are the $Z_\tau = \Ch_\tau(S; s, -)$, which are \textit{contractible} categories (i.e. $Z_\tau \longrightarrow e$ is a homotopism), hence $W$-aspheric by W(3) (cf. cor. 3, page 3).

On the other hand, if we grant that $\Ch(S) \xrightarrow{p_S} S{\times}S$ is $W$-smooth (and $W$-proper), hence $\Ch(S) \xrightarrow{\sigma_S} S$ also, having $W$-aspheric fibers as we have just seen, it follows that $\sigma_S \in W$, whence $S \xrightarrow{i_S} \Ch(S)$ also in $W$ (since $\sigma_S \circ i_S = \mathrm{id}_S$). So it remains only to prove that
\[ \Ch(S) \xrightarrow{p_S} S \times S \tag*{$(*)$} \]
is $W$-proper and $W$-smooth. It will follow that
\[ \Ch(S; s_0, -) \xrightarrow{\beta} S \]
is likewise $W$-proper and $W$-smooth, since obtained from $p_S$ by the base change
\[ S \longrightarrow S \times S \qquad s \longmapsto (s_0, s). \]

Next we must see that $W$-smooth and $W$-proper implies the $W$-fibration property \textit{in its strongest form} a) (whereas I foresee that a priori we will only have it in the weak forms a$_1$, a$_2$ - not even in the form a$'$). So it is here that the axiom W(4) of $W$-fibrations will be essential (\footnote[15]{in the weakened form W(4a$''$) (p. 13).})!

Here is the argument for the smoothness of $(*)$, in loc. cit. p. 161 ff. We must prove that for $c \in \mathrm{Ob}\,\Ch(S)$, and an arrow $(u,v) \in \mathrm{Fl}(S \times S)$ \texttt{[rather} $(u,v) \in$ \texttt{Fl(}$S \times S$\texttt{)]}, a certain category of ``liftings of $(u,v)$ with target $c$ '' \texttt{[rather ``liftings of} $(u,v)$ \texttt{with target} $c$ \texttt{'']}, $\mathcal{C} = \mathcal{C}(c,u,v)$, is $W$-aspheric (\footnote[16]{In fact, we prove that $\mathcal{C}(c,u,v)$ is \textit{contractible}, hence it is $W_\omega$-aspheric.}). For this, we compare it to a category $\mathcal{C}'$ with a final object, hence contractible and a fortiori $W$-aspheric, with the help of two functors
\DERmark{ENV0645}\DERmark{ENV0646}

\[ \begin{tikzcd} \mathcal{C} \arrow[r, shift left, "\varphi"] & \mathcal{C}', \arrow[l, shift left, "\psi"] \end{tikzcd} \]
and the homotopies $\varphi\psi \sim \mathrm{id}_{\mathcal{C}'}$ (no merit, $\mathcal{C}'$ being contractible)
\msp{page 10}
and $\psi\varphi \sim \mathrm{Id}_{\mathcal{C}}$ (we even find an arrow $\psi\varphi \longrightarrow \mathrm{id}_{\mathcal{C}}$). Thus $\varphi$ is a homotopism, hence a $W$-equivalence (cor. 3, p. 3), hence $\mathcal{C}$ is $W$-aspheric, $\mathcal{C}'$ being so. The properness of $p_S$ is equivalent to the smoothness of $p_{(S^\circ)}$, it is therefore proved at the same time.
% foot: Version 7 décembre 2002


\endgroup
% END UNIT Der-ChXIIa-012/Der-ChXIIa-012_p006-010.tex
% BEGIN UNIT Der-ChXIIa-012/Der-ChXIIa-012_p011-015.tex
\begingroup\makeatletter%
\NoAutoSpacing%
\newcommand{\printedpage}[1]{\marginpar{\footnotesize p.~#1}}%
\newcommand{\msp}[1]{\marginpar{\footnotesize [#1]}}%
\newcommand{\Ch}{\underline{\mathrm{Ch}}}%
\makeatother\let\printedpage\derpp\let\msp\dermsp%
\let\mspin\DERmspin%


\printedpage{9}
% head: A. Grothendieck : Dérivateurs, Ch. XII, Édition des Universités de Montpellier II et Paris VII

Thus, the lemma on page 7 is proved, using axiom W(4) which ensures that an arrow which is both $W$-smooth and $W$-proper satisfies the very strong condition a) of $W$-fibrations.

\textbf{NB} In fact, it would suffice, instead of axiom W(4) in its strong form, to assume that an $f : X \longrightarrow Y$ which is both ``homotopically smooth'' and ``homotopically proper'', i.e. the fibers of the $X/x \longrightarrow Y/y$ and of the $x\backslash X \longrightarrow y\backslash Y$ are contractible (hence $W$-aspheric for every $W$ satisfying W(1,2,3)), is a strong $W$-fibration. Let us state all the same what we have proved.

\textbf{Proposition 4.} \textit{If $W$ satisfies} W(1,2,3) (\footnote[17]{For W(2), it suffices in the form W(2 bis).}) (we always assume this in what follows) \textit{plus} W(4) (in a somewhat weakened sufficient form, cf. preceding remark), \textit{we have this: Let $f : X \longrightarrow Y$ be an arrow in} Cat$/S$\textit{, with $X$, $Y$ $W$-fibered over $S$. If $f \in W$, then $f_s : X_s \longrightarrow Y_s$ is in $W$ for every $s \in S$.}

I wonder whether I know enough to be able to prove
\[ W \subseteq W_\infty . \]
We must all the same eliminate joke cases, such as $W = \mathrm{Fl}(\mathrm{Cat})$, hence put in the axiom

\begin{itemize}[leftmargin=4em, labelwidth=3em, labelsep=0.6em]
\item[W(5)] (Axiom of connectedness (\footnote[18]{See on this subject the corollary on page 22.}).) If $f : X \longrightarrow Y$ is in $W$, $\pi_0(f) : \pi_0(X) \longrightarrow \pi_0(Y)$ \texttt{[is]} bijective.
\end{itemize}

\msp{page 11}I then factor $f$ as (\footnote[19]{Once again, in these arguments it pays to work with $W_\omega$.})
\[ X \xrightarrow{i} \bar{X} \xrightarrow{\bar{f}} Y, \]
with $\bar{f}$ a $W_\infty$-fibration and a $W$-fibration (\footnote[20]{in fact, $\bar{f}$ is both $W_\infty$-perfect (i.e. $W_\infty$-smooth and proper) and $W$-perfect (in fact, $W_\omega$-perfect).}), $i \in W \cap W_\infty$ - this works after all, % [sic] finallement -> finalement, spelling
by the construction in VII, p. 170-171 (of \textsc{Cartan-Serre}), starting from the properties of % [sic] propiétés -> propriétés, missing letter
$Y \longrightarrow \Ch(Y) \longrightarrow Y \times Y$ and nothing more. As $\bar{f} i = f \in W$ and $i \in W$, we deduce $\bar{f} \in W$. As $i \in W_\infty$, we conclude $f \in W_\infty \iff \bar{f} \in W_\infty$. Thus to prove that $f$ is in $W_\infty$, we are reduced to the case where $f$ is a $W$-fibration and a $W_\infty$-fibration. But then proposition 4 shows us that the fibers of $f$ are $W$-aspheric (\footnote[21]{we have used here only this property of the fibers, as a consequence of W(4). (We use W(4) in its weakest form W(4b$''$), cf. pages 13, 14).}), and it suffices to see that they are $W_\infty$-aspheric (proposition 3 (\footnote[22]{we use proposition 3 only in the case $Y = S$, for which W(2 bis) suffices.}) applied to $W_\infty$). Thus we are reduced to the case where $Y = e$, i.e. to proving that if $X$ in Cat is $W$-aspheric, it is $W_\infty$-aspheric.
% [sic] (footnote 21) consequence -> conséquence, missing accent

Let $x_0 \in \mathrm{Ob}\, X$, let us consider as before
\[ E = \Ch(X, x_0, -) \xrightarrow{p} X, \]
% foot: Version 7 décembre 2002
\printedpage{10}
% head: A. Grothendieck : Dérivateurs, Ch. XII, Édition des Universités de Montpellier II et Paris VII
we know that $E$ is $W \cap W_\infty$-aspheric, and that it is a fibration for $W$ and for $W_\infty$. As $X$ is $W$-aspheric (like $E$), $p \in W$, hence again by proposition 4 (\footnote[23]{in fact, condition W(4b$''$) on $W$.}), the fibers of $p$ are $W$-aspheric, in particular (by W(5)) 0-connected. We have therefore proved the

\msp{page 12}\textbf{Lemma.} (\footnote[24]{Uses W(1, 2 bis, 3, 4b$''$) on $W$.}) \textit{Let $(X, x_0)$ be in} Cat$_\bullet$\textit{. If $X$ is $W$-aspheric, then $\Omega(X, x_0) \overset{\mathrm{def}}{=} \Ch(X; x_0, x_0)$ is too} (assuming axioms W(1,2,3,4) on $W$)\textit{. Thus (by iteration) the $\Omega^n(X, x_0)$ are $W$-aspheric} ($\forall\, n \in \mathbf{N}$).

\textbf{Corollary.} \textit{If moreover $W$ satisfies the axiom} W(5) \textit{of connectedness, then $X$ is $W_\infty$-aspheric.}
\DERmark{ENV0647}


Indeed, ``we know'' that this means that all the $\Omega^n(X, x_0)$ are 0-connected. Now the connectedness axiom W(5) ensures that $W$-aspheric implies 0-connected.

But the ``we know'' is a tiny bit flimsy (\footnote[25]{This is proved quite simply in the corollary to ``theorem 1'', page 190.}). To justify it, it would be necessary, either % [sic] (footnote 25) theorème -> théorème, missing accent
that I make the link with classical 1/2-simplicial theory, and prove that the $\Omega^n$ which I construct are ``the right ones'', i.e. that their nerve does indeed give back (up to isomorphism in Hot) the classical $\Omega^n$ (??). In particular, that the $\pi_0(\Omega^n(X, x_0))$ is indeed ``everybody's'' $\pi_n(X, x_0)$. Or else that I prove directly that the 0-connectedness of the $\Omega^n(X, x_0)$ is equivalent to the $W_\infty$-asphericity of $X$, defined cohomologically in the manner of \textsc{Artin-Mazur}. So this is a piece of work that I shall absolutely have to do. But it is now no more than a simple piece of routine work, and I can state right now the

\textbf{Theorem 1.} (\footnote[26]{It even suffices to have W(4) in a weakened form, with ``fibrant'' functors in an ``absolute'' sense. Cf. the NB p. 10 - and above all, cf. pages 13, 14.},\footnote[27]{NB Instead of W(5), it suffices to assume $W \neq \mathrm{Fl}(\mathrm{Cat})$, cf. th. 2, p. 20.},\footnote[28]{The useful hypotheses on $W$ are in the end W(1), W(2 bis), W(3), W(4b$''$), and $W \neq \mathrm{Fl}(\mathrm{Cat})$.}) \textit{Let $W \subseteq$} Fl(Cat) \textit{satisfy the conditions} W(1,2,3,4,5) (axiom of saturation, of localization, of the final object, of fibration, of connectedness)\textit{. Then}
\[ W \subseteq W_\infty . \]

\textbf{Corollary.} \textit{We can characterize $W_\infty$ as the largest subset of} Fl(Cat) \textit{satisfying the}
\msp{page 13}\textit{conditions} W(1) \textit{to} W(5) \textit{above.}

Of course, one would like a more interesting axiom 6, which joined to the preceding ones ensures that $W = W_\infty$ (\footnote[29]{Cf. axiom W(6) on the next page.}). Or do these properties already ensure that $W = W_\infty$? To see this, perhaps I would have to know more about the structure of $W_\infty$-aspheric objects? If this were false, I am thinking of an axiom of ``homotopy descent'', which would imply that a $W_\infty$-aspheric object is also $W$-aspheric.

% foot: Version 7 décembre 2002
\printedpage{11}
% head: A. Grothendieck : Dérivateurs, Ch. XII, Édition des Universités de Montpellier II et Paris VII

I note right now that if some $W_i$ ($i \in I$) satisfy axioms W(1) to W(5), the same is true of $\bigcap_i W_i$ (\footnote[30]{This is true for any set of axioms in the nature of stability axioms.}). Thus there is a smallest set of arrows $W_\omega$ satisfying these conditions. If it is not equal to $W_\infty$, i.e. if axioms W(1) to W(5) do not already characterize $W_\infty$, it would be a really very interesting localizer.

I realize that I have used only a rather weak form of axiom W(4) - essentially, the validity of proposition 4. I am going to summarize different variants of the fibration axiom.

\textbf{W(4)}

\begin{itemize}[leftmargin=4em, labelwidth=2.5em, labelsep=0.6em, align=right]
\item[a$_1$)] If $f : X \longrightarrow Y$ in Cat is a left $W$-fibration, it is an ultrastrong $W$-fibration.

\item[a$_2$)] If $f : X \longrightarrow Y$ is a right $W$-fibration, it is an ultrastrong $W$-fibration.

\item[a$'$)] If $f : X \longrightarrow Y$ is $W$-perfect, i.e. both $W$-proper and $W$-smooth (\footnote[31]{which implies that it is a left $W$-fibration \textit{and} a right one.}), then $f$ is an ultrastrong $W$-fibration.

\item[a$''$)] If $f : X \longrightarrow Y$ is both homotopically proper and homotopically smooth (cf. NB page 10), $f$ is an ultrastrong $W$-fibration. (It is better to say: If $f$ is $W_\omega$-perfect, it is an ultrastrong $W$-fibration.)

\item[b)] Let $f : X \longrightarrow Y$ be an arrow in Cat$/S$, with $X$ and $Y$ $W$-fibered over $S$. If $f \in W$, then $\forall\, s \in S$, $f_s : X_s \longrightarrow Y_s$ is \texttt{[in]} $W$.

\item[b$'$)] Let $f : X \longrightarrow Y$ be an arrow in Cat$/S$ such that $X \longrightarrow S$ and $Y \longrightarrow S$ are $W_\omega$-perfect. If $f \in W$, then $\forall\,s \in S$, $f_s : X_s \longrightarrow Y_s$ is in $W$. \texttt{[Rather: Let $f : X \longrightarrow Y$ be an arrow in \textrm{Cat}$/S$ such that $X \longrightarrow S$ and $Y \longrightarrow S$ are $W_\omega$-perfect. If $f \in W$, then $\forall\, s \in S$, $f_s : X_s \longrightarrow Y_s$ is in $W$.]}

\item[b$''$)] \msp{page 14}Let $W_\omega \subseteq W \cap W_\infty$ be the smallest subset of Fl(Cat) satisfying conditions W(1,2,3) (\footnote[32]{It is better to take W(2 bis) instead of W(2), to find $W_\omega$ as small as possible!}) (axioms of saturation, of localization, of the final object). Let $f : X \longrightarrow Y$ be an arrow in Cat which is $W_\omega$-perfect. If $f \in W$, then the fibers of $f$ are $W$-aspheric \texttt{[rather $W$-aspheric]}.
\end{itemize}
\DERmark{ENV0648}\DERmark{ENV0649}


Among these properties of $W$, we have the tautological implications (except a$''$) $\Longrightarrow$ b), which is essentially proposition 4)

% foot: Version 7 décembre 2002
\printedpage{12}
% head: A. Grothendieck : Dérivateurs, Ch. XII, Édition des Universités de Montpellier II et Paris VII

\[
\begin{tikzcd}[row sep=huge, column sep=0pt]
\mathrm{a} \arrow[d, Rightarrow] & {}= (\mathrm{a}_1 \Longleftrightarrow \mathrm{a}_2) & \hspace{2.5cm} \\
\mathrm{a}' \arrow[d, Rightarrow] & & \smash{\left.\rule[-1.8cm]{0pt}{3.6cm}\right\}\ \text{\begin{tabular}{l}In fact,\\ all these\\ conditions are\\ equivalent\end{tabular}}} \\
\mathrm{a}'' \arrow[d, Rightarrow] & & \\
\mathrm{b} \arrow[d, Rightarrow] & & \\
\mathrm{b}' \arrow[d, Rightarrow] & & \\
\mathrm{b}'' & &
\end{tikzcd}
\]

(\footnote[33]{I return to the various variants of W(4) below, p. 53-56.}).

Well, in the lemma on page 12, it sufficed to assume that $W$ satisfies the weakest version b$''$) of the axiom of fibrations (seven versions!). Indeed, the fibrations $\bar{X} \longrightarrow Y$ and $E \longrightarrow X$ used in the proof of the lemma are $W_\omega$-perfect, and it sufficed to know that if they are in $W$, then the fibers are $W$-aspheric.

But I do not think that even the strongest version W(4a) of the axiom of fibrations, joined to the three preceding axioms and to W(5) (connectedness), implies $W = W_\infty$. It would at least be necessary to add to it the axiom

\begin{itemize}[leftmargin=4em, labelwidth=3em, labelsep=0.6em]
\item[\textbf{W(6)}] (Axiom of filtering limits.) Let $X = \varinjlim\limits_{i} X_i$ be a filtering direct limit in Cat. If the $X_i$ are $W$-aspheric, $X$ too.
\end{itemize}

This should be equivalent to the property, apparently stronger, in terms of arrows of $W$:

\textbf{Proposition 5.} \textit{Suppose we have} W(6) (in addition to W(1) to W(4) (\footnote[34]{W(4) in the weakest form W(4b$''$) .})),
\msp{page 15}\textit{and let}
\[ (f_i)_{i\in I} : (X_i)_{i\in I} \longrightarrow (Y_i)_{i\in I} \]
% foot: Version 7 décembre 2002
\printedpage{13}
% head: A. Grothendieck : Dérivateurs, Ch. XII, Édition des Universités de Montpellier II et Paris VII
\textit{a homomorphism of direct systems in} Cat\textit{, indexed by a filtering ordered set, let $X = \varinjlim X_i$, $Y = \varinjlim Y_i$, $f = \varinjlim f_i$, $f : X \longrightarrow Y$. If $f_i$ is in $W$ for every $i$, then $f \in W$.}

We factor each $f_i$ as
\[ X_i \xrightarrow{\alpha_i} \bar{X}_i \xrightarrow{\bar{f}_i} Y_i, \qquad \alpha_i \in W_\omega,\ \bar{f}_i \in \mathrm{Parf}\, W_\omega \]
(\footnote[35]{$W_\omega$ introduced in condition W(4b$''$), page 14.}) following the patented system. The construction being functorial, the $\bar{X}_i$ too form a direct system. Moreover, as the functor $S \longrightarrow \Ch(S)$ which gives rise to the canonical factorization of morphisms, commutes with filtering $\varinjlim$ (immediate), it follows that, if we put
\[ \bar{X} = \varinjlim_{i} \bar{X}_i \]
and look at the arrows
\[ X \xrightarrow{\alpha} \bar{X} \xrightarrow{\bar{f}} Y, \qquad \alpha = \varinjlim \alpha_i,\ \bar{f} = \varinjlim \bar{f}_i, \]
with $\bar{f}\alpha = {f}$, this factorization of $f$ is indeed its canonical factorization. Thus $\alpha \in W$, % [sic] $\bar{f}\alpha = \bar{f}$ -> $\bar{f}\alpha = f$, the composite should be f
\DERmark{ENV0650}

$\bar{f} \in \mathrm{Parf}\, W_\omega$. This having been seen, by virtue of prop. 3, $f \in W$, i.e. $\bar{f} \in W$, is true if the fibers of $\bar{f}$ are $W$-aspheric. But the fibers of $\bar{f}$ are interpreted as filtering direct limits of fibers of the $\bar{f}_i$, which are $W$-aspheric by condition W(4b$''$) on fibrations. Thus axiom W(6) gives what is needed.

\textbf{Conjecture} (bold? foolish?). \textit{If $W \subseteq$} Fl(Cat) \textit{satisfies the axioms} W(1) \textit{to} W(6) \textit{(with} W(4) \textit{in its weakest form} W(4b$''$) (\footnote[36]{and perhaps W(2) in the weaker form W(2 bis).})\textit{), then $W = W_\infty$.}

\msp{page 16}But it is quite possible that one must strengthen W(1) (saturation) a little, by requiring that every retract of an $f \in W$ be in $W$.

\textbf{Remark.} Let us consider, for $n \in \mathbf{N}$, the set $W_n \subseteq$ Fl(Cat) of ``$n$-homotopy equivalences'' in Cat:
\[ f \in W_n \iff \pi_0(f) \text{ bijective, and } \pi_i(f, x) : \pi_i(X, x) \longrightarrow \pi_i(Y, f(x)) \text{ bijective for } 0 < i \leq n. \]

I thought that this satisfied all that one can ask of a localizer, without having taken the trouble to write it down carefully. But this cannot satisfy W(1) through W(4b$''$), since we % [sic] Ma -> Mais, word truncated
do not have $W_n \subseteq W_\infty$!

W(1) is trivial from the definition (and we have without tears even the \textit{strong saturation} of $W_n$, % [sic] $W_i$ -> $W_n$ (also in the following lines), the remark concerns W_n
\DERmark{ENV0651}

i.e. if $f$ in Cat is \texttt{[an]} isomorphism in $\mathrm{Hot}_{W_n} = (\mathrm{Cat})W_n^{-1}$, then $f \in W_n$). I think I remember having checked W(2), by the same method as for $W_\infty$, using a cohomological criterion in the manner of \textsc{Artin-Mazur}, to express $f \in W_n$, and using the Leray spectral sequences for $X \longrightarrow S$ and $Y \longrightarrow S$. Condition W(3) is trivially satisfied since $W_n \supseteq W_\infty$ and it is for $W_\infty$. So it gets botched where I
\DERmark{ENV0652}\DERmark{ENV0653}\DERmark{ENV0654}\DERmark{ENV0655}

% foot: Version 7 décembre 2002


\endgroup
% END UNIT Der-ChXIIa-012/Der-ChXIIa-012_p011-015.tex
% BEGIN UNIT Der-ChXIIa-012/Der-ChXIIa-012_p016-020.tex
\begingroup\makeatletter%
\NoAutoSpacing%
\newcommand{\printedpage}[1]{\marginpar{\footnotesize p.~#1}}%
\newcommand{\msp}[1]{\marginpar{\footnotesize [#1]}}%
\newcommand{\Fl}{\mathrm{Fl}}%
\newcommand{\Cat}{\mathrm{Cat}}%
\newcommand{\card}{\operatorname{card}}%
\makeatother\let\printedpage\derpp\let\msp\dermsp%
\let\mspin\DERmspin%


\printedpage{14}
% head: A. Grothendieck : Dérivateurs, Ch. XII, Édition des Universités de Montpellier II et Paris VII
least expected: the fibration axiom. Let us take the weak form: Let $f : X \longrightarrow Y$ be a $W_\omega$-fibration and $f \in W_i$, are the fibers $X_s$ $W_i$-aspheric (\footnote[37]{$\omega > \infty$, $W_\omega \subset W_\infty \subset W_n$})? We may suppose $X$ and $Y$ 0-connected, let $x \in X_s$, the homotopy exact sequence gives
\[
\pi_{i+1}(X) \xrightarrow[\text{if } i \leq n-1]{\text{iso}} \pi_{i+1}(Y) \longrightarrow \pi_i(X_s) \longrightarrow \pi_i(X) \xrightarrow[\text{if } i \leq n]{\text{iso}} \pi_i(Y),
\]
so
\[
\pi_i(X_s) = 1 \qquad \text{if } 0 \leq i \leq n-1,
\]
but for $i = n$ we find\msp{page 17}
\[
\pi_n(X_s) = \begin{cases}\mathrm{Coker}\left(\pi_{n+1}(X) \longrightarrow \pi_{n+1}(Y)\right),&n\geq1,\\ \pi_1(Y)/\operatorname{im}\!\left(\pi_1(X)\longrightarrow\pi_1(Y)\right),&n=0\quad\text{(pointed set of cosets)}.\end{cases},
\]
\DERmark{ENV0656}

there is no reason for this to be zero, as it should if $X_s \longrightarrow e$ were in $W_n$.

It is true that sometimes one does strengthen the condition for $W_n$, let us say one defines $W'_n \subseteq W_n$ by
\[
f \in W'_n \iff \left\{
\begin{array}{ll}
\text{a)} & \pi_0(f) \text{ bijective,} \\[1ex]
\text{b)} & \pi_i(f,x) : \pi_i(X,x) \longrightarrow \pi_i(Y,f(x)) \text{ bijective for } 0 < i \leq n \text{ and } x \in X, \\[1ex]
\text{c)} & \pi_{n+1}(f,x) : \pi_{n+1}(X,x) \longrightarrow \pi_{n+1}(Y,f(x)) \ \textit{surjective,}
\end{array}
\right.
\]
that is to say also
\[
f \in W'_n \iff \left\{
\begin{array}{l}
\text{the ``homotopy fibers'' } Z \text{ of } X \longrightarrow Y \text{ (i.e. the fibers of } \bar{X} \xrightarrow{\bar{f}} Y\text{)} \\
\text{are } n\text{-connected, i.e. 0-connected and } \pi_i(Z_y,z) = 0 \text{ for } 1 \leq i \leq n.
\end{array}
\right.
\]
Then $X$ in Cat will be $W'_n$-aspheric if and only if it is $W_n$-aspheric, i.e. $n$-connected. So by form 2 of the definition of $W'_n$, axiom W(4b$''$) is satisfied. It must go wrong elsewhere! Perhaps saturation already.

\hspace*{3em}If $gf$ and $g$ are in $W'_i$, the same is not necessarily true of $f$,

for if $\bar{g}\bar{f}$ (and hence $\bar{g}$) is surjective, the same is not necessarily true for $\bar{f}$! So $W'_n$ does not deserve the name of ``localizer''. % [sic] par forcément -> pas forcément, misprint

On the other hand, $W_n$ seems excellent to me, if one does not ask the impossible of it - to satisfy a ``fibration axiom'' which does seem (more or less) to characterize $W_\infty$ alone among the ``good'' localizers. In any case the connectedness axiom W(5) is satisfied, and the filtering limits axiom W(6) as well. Moreover W(1) (and even strong saturation), W(2), W(3) are satisfied by every $W_{\mathbf{D}}$ associated with a derivator - which confirms their importance once again. The axioms W(4) (the ``draconian'' fibration axiom, $\pm$ special to $W_\infty$), W(5) (connectedness) and % [sic] connéxité -> connexité, misspelling (written connexité elsewhere)
\printedpage{15}
% head: A. Grothendieck : Dérivateurs, Ch. XII, Édition des Universités de Montpellier II et Paris VII
\msp{page 18}W(6) (limits) appear to me to be of a more special nature, I would not think of imposing any of them for the terminological use of ``basic localizer''.

\par
\printedpage{16}
% head: A. Grothendieck : Dérivateurs, Ch. XII, Édition des Universités de Montpellier II et Paris VII
\section*{II\quad Perplexities about the axioms and the terminology.}

In I and most of the time, in working with derivators, I have assumed \fbox{$\ W_{\mathbf{D}} \supseteq W_\infty\ $}, a hypothesis which seemed essential to me (wrongly (\footnote[38]{\textit{wrongly} it now turns out.}) or rightly). So we see that if $W_{\mathbf{D}}$ satisfies the fibration axiom W(4), even in its weakest form W(4b$''$), it is \textit{equal} to $W_\infty$. This does show that this axiom must be seen as practically a \textit{characterization} of $W_\infty$ (among the basic localizers which contain $W_\infty$ -- and these localizers seem to me the most important, and perhaps the only really useful ones).

The theory of \textbf{D}-fibrations in I, p. 111-126, was done in the spirit just mentioned, assuming $W_{\mathbf{D}} \supseteq W_\infty$. The notion of \textbf{D}-fibrations on which I had finally settled was consequently different from the notions considered here in a different spirit (page 4). It amounts to testing condition a) (almost always too strong, except for $W = W_\infty$), characterizing ultrastrong $W$-fibrations, by $Y'' \longrightarrow Y'$, where $Y''$, $Y'$ are both $W_\infty$-aspheric (which is more demanding than $W$-aspheric). It is only when $Y''$ and $Y'$ are $W_\infty$-aspheric (instead of $W$-aspheric), that one asks that $X'' \longrightarrow X'$ be in $W$. We thus find\msp{page 19} a less demanding notion than that of strong $W$-fibration. If we denote this condition here by b), we find (provided that $W_\infty \subseteq W$), with the notations of page 4 (\footnote[39]{Cf. the right notion of $W$-fibration in a context where one does not assume $W_\infty \subseteq W$, in th. 3 and def. 3, p. 51, 52.}) % [sic] arreté -> arrêté, missing circumflex
% In the print, a curved arrow starts at "b" in the diagram below and points up to "b)" in the line "par b), on trouve ..." above; it is drawn here to the empty node at the top left.
\[
\begin{tikzcd}[column sep=small, row sep=large]
{} & & & & \\
& & \mathrm{a} \arrow[d, Rightarrow] & & W\text{-ultrastrong fibration} \\
& & \mathrm{a}' \arrow[d, Rightarrow] & & W\text{-strong fibration} \\
& & \mathrm{b} \arrow[dl, Rightarrow] \arrow[dr, Rightarrow] \arrow[uuull, bend left=30] & & W\text{-fibration} \\
W\text{-left fibration} & \mathrm{a}_1 & & \mathrm{a}_2 & W\text{-right fibration.}
\end{tikzcd}
\]
a$_1$), a$_2$) were also called cohomological \textbf{D}-fibrations resp. homological \textbf{D}-fibrations (\footnote[40]{No, this is a still stronger notion, by dint of locally constant coefficients, cf. I, p. 121.}). When \textbf{D} satisfies Der 3a$'$ (I, p. 111) on locally constant \textbf{D}-coefficients on the $X$ that are $W_\infty$-aspheric, then the notions b), a$_1$), a$_2$) are equivalent. (I presume this is the case for all the \textbf{D} and all the $W$ that are really interesting.) On the other hand, I have the impression that the notion of strong $W$-fibration, or even ultrastrong, is hardly usable except in the case $W = W_\infty$, where it coincides with the preceding ones. That is to say that one can probably forget the terminology, and remember on the other hand that in the case $W = W_\infty$, the % [sic] cohomologique -> cohomologiques, agreement with "D-fibrations"
\printedpage{17}
% head: A. Grothendieck : Dérivateurs, Ch. XII, Édition des Universités de Montpellier II et Paris VII
$W$-fibrations have properties of extreme strength, expressed by the two properties singled out in W(4a) and W(4b).

\msp{page 20}In the theory of $\mathrm{HOT}_W$, influenced by the general reflection on derivators and axiom Der 3b, I have worked under the hypothesis $W \supseteq W_\infty$, but using a fibration axiom L5 $=$ W(4b) which, it turns out, ultimately implies $W = W_\infty$. I have the impression that so far (properties of $f_!$ for HOT) I have not used $W \supseteq W_\infty$, so it may be useful, since ``it costs no more'' (and forces one to be more attentive to what one uses), to drop $W \supseteq W_\infty$ (implying $W = W_\infty$), and on the other hand to put on $W$ axioms as strong (notably of fibration) as the properties of $W_\infty$, ``without depriving oneself of anything''. So I shall assume W(1), W(2), W(3), W(4a) (the strongest form of the fibration axiom, cf. pages 13, 14) - that is all that has been useful up to now. If one adds W(5) (connectedness axiom) to this it implies $W \subseteq W_\infty$. But for the moment I do not think I have a need for it.

But meditating on the question, I believe I have found the

\textbf{Theorem 2} (!) (\footnote[41]{Cf. for a stronger result XV 9.12, p. 133 bis. \textbf{NB} The result is more general, but not really stronger, since it does not give $W = \Fl(\Cat)$.}): \textit{Let} $W \subseteq \Fl(\Cat)$ \textit{satisfy} W(1,2,3,4b$''$)\textit{. If} $W$ \textit{does not satisfy} W(5)\textit{, then} $W = \Fl(\Cat)$.

\textsc{Proof}. Let $f : X \longrightarrow Y$ be in Cat, such that $\pi_0(f)$ is not bijective. Let $I$ be the discrete category having as set of objects the set $\pi_0(Y)$, so we have a commutative diagram
\[
\begin{tikzcd}[column sep=small]
X \arrow[rr, "f"] \arrow[dr, "p"'] & & Y \arrow[dl, "q"] \\
& I, &
\end{tikzcd}
\]
on the other hand a functor from a category into a discrete category is always $W_\omega$-perfect, so $p$, $q$ are $W$-fibrations. Since\msp{page 21} $f \in W$, it follows that $f_i \in W$ for all $i \in I$ (\footnote[42]{We use W(4) in a form stronger than W(4b$''$)!}).

If $\pi_0(f)$ were not surjective, taking $i$ which is not in the image, we would have $X_i = \emptyset \longrightarrow Y_i \neq \emptyset$, which would be in $W$. So let
\[
f : X = \emptyset \longrightarrow Y
\]
be in $W$. Then $\forall$ $Z$ in Cat, $f \times \mathrm{id}_Z : \underbrace{X \times Z}_{= \emptyset} \longrightarrow Y \times Z$ is in $W$ (since $Z$ is
\printedpage{18}
% head: A. Grothendieck : Dérivateurs, Ch. XII, Édition des Universités de Montpellier II et Paris VII
$W$-fibered over $e$), so $\emptyset \longrightarrow Y \times Z$ is in $W$. So too the projection
$\begin{tikzcd}[column sep=small, row sep=small, baseline=(\tikzcdmatrixname-1-1.base)]
Y \times Z \arrow[rr] & & Y. \\
& \emptyset \arrow[ul] \arrow[ur] &
\end{tikzcd}$
Since this projection is a $W$-fibration, it follows by W(4b$''$) that its fibers are $W$-aspheric. Since $Y \neq \emptyset$, there is such a fiber, it is $\simeq Z$. So $Z$ is $W$-aspheric. So we find that \textit{every} category $Z$ in Cat is $W$-aspheric. So every arrow $Z \longrightarrow Z'$ in Cat is in $W$, q.e.d.

Suppose now $\pi_0(f)$ surjective, but not injective. There would be a fiber of $\pi_0(X) \longrightarrow \pi_0(Y) = I$ which would be of cardinality $> 1$. Taking the corresponding $f_i$, we find $X_i \longrightarrow Y_i$ with $Y_i$ 0-connected and $\card \pi_0(X_i) \geq 2$. So we are reduced to the case
\[
f : X \longrightarrow Y \qquad\qquad \left\{
\begin{array}{l}
f \in W \\[1ex]
Y \text{ 0-connected} \\[1ex]
\card \pi_0(X) \geq 2.
\end{array}
\right.
\]
Let us consider the canonical factorization
\[
X \xrightarrow{i_f} \bar{X} \xrightarrow{\bar{f}} Y,
\]
with $i_f \in W$, so $\bar{f} \in W$. So by W(4b$''$) the fibers of $\bar{f}$ are $W$-aspheric. But we see at once, on decomposing $\bar{X}$ according to the connected components $X_i$ of $X$, that these fibers have a $\pi_0$ of cardinality $\geq \card \pi_0(X) \geq 2$. So we are reduced to the case where $Y = e$, i.e.
\[
X \ W\text{-aspheric}, \qquad \card \pi_0(X) \geq 2.
\]
Let us consider $X \times X$, which is also $W$-aspheric, and the diagonal map\msp{page 22}
\[
d : X \longrightarrow X \times X = Y,
\]
which is therefore in $W$. But $\pi_0(d)$ is identified with the diagonal map $\pi_0(X) \longrightarrow \pi_0(X) \times \pi_0(X)$, which is not surjective since $\card \pi_0(X) \geq 2$. So for the first part, $W = \Fl(\Cat)$.

So in more positive form:

\textbf{Corollary.} \textit{Let} $W \subseteq \Fl(\Cat)$ \textit{satisfy} W(1,2,3,4b$''$)\textit{, and be such that} $W \neq \Fl(\Cat)$. \textit{Then $W$ satisfies the connectedness axiom} W(5)\textit{, and consequently (th. 1, page 12)} $W \subseteq W_\infty$.


\endgroup
% END UNIT Der-ChXIIa-012/Der-ChXIIa-012_p016-020.tex
% BEGIN UNIT Der-ChXIIa-012/Der-ChXIIa-012_p021-025.tex
\begingroup\makeatletter%
\derleqno%
\newcommand{\printedpage}[1]{\marginpar{\footnotesize p.~#1}}%
\newcommand{\msp}[1]{\marginpar{\footnotesize [#1]}}%
\newcommand{\Ob}{\mathrm{Ob}\,}%
\makeatother\let\printedpage\derpp\let\msp\dermsp%
\let\mspin\DERmspin%


\printedpage{19}
% head: A. Grothendieck : Dérivateurs, Ch. XII, Édition des Universités de Montpellier II et Paris VII
% foot: Version 7 décembre 2002

\section*{III\quad Remarkable classes of functors related to $W$.}

Finally, I will call \textit{``basic localizer''} a subset $W$ of Fl(Cat) satisfying the following three conditions: % [sic] appelerai -> appellerai, spelling

\begin{enumerate}
\item[(i)] Saturation W(1) (cf. reminder p. 1).
\item[(ii)] Final element axiom W(3) (page 3).
\item[(iii)] The weakened variant W(2 bis) of the localization axiom W(2) given in the corollary on page 1 (formerly it was L4$'$ in Pursuing Stacks).
\end{enumerate}

When I need more from $W$, I will say so explicitly. These three axioms have sufficed for practically everything I have done with $W$ in Pursuing Stacks, the only exception being a single use (if I recall), somewhere, of a weakened variant of W(2) (for the \textit{cartesian} $S$-functors between \textit{fibered} $S$-categories - if on every fiber it is in $W$, it is itself in $W$).

We therefore suppose henceforth given a basic localizer
\[
W \subseteq \mathrm{Fl}(\mathrm{Cat}).
\]

\textbf{Proposition 1.} \textit{Let $f : X \longrightarrow Y$ be in} \textup{Cat}. \textit{Equivalent conditions:}

\begin{enumerate}
\item[\textit{a)}] \textit{For every local isomorphism $Y' \longrightarrow Y$, $f' : X' \longrightarrow Y'$ is $\in W$.}
\item[\textit{b)}] \textit{Same condition, in the particular case $Y' = Y/y \longrightarrow Y$, so $\forall\ y \in Y$, $X/y \simeq X \times_Y (Y/y)$ is $W$-aspheric.}
\end{enumerate}

\msp{page 23}\textbf{Definition 1.} Under these equivalent conditions a), b) of proposition 1, we say that $f$ is \textit{$W$-aspheric.} We say \texttt{[that]} $f$ is \textit{$W$-coaspheric} if $f^{\circ}$ is $W$-aspheric, i.e. if the $y\backslash X$ are $W$-aspheric (for $y \in \Ob Y$).

\textsc{Proof.} It suffices to prove b) $\Longrightarrow$ a). By definition of a local isomorphism, we see that condition b) is stable under base change $Y' \longrightarrow Y$, with $Y'$ a local isomorphism (since the $Y'/y' \xrightarrow{\sim} Y/y$, hence $X'/y' \xrightarrow{\sim} X/y$). It therefore suffices to see that if $f$ satisfies b), then $f \in W$. But this is axiom W(2 bis).

\begin{center}
\begin{tikzcd}[row sep=large]
X \arrow[d] & & & \\
Y & & Y' \arrow[ll] & \\
& Y/y \arrow[ul] & & Y'/y' \arrow[ll] \arrow[ul]
\end{tikzcd}
\end{center}

\printedpage{20}
% head: A. Grothendieck : Dérivateurs, Ch. XII, Édition des Universités de Montpellier II et Paris VII
% foot: Version 7 décembre 2002

\textbf{Remark.} If $f$ has a right adjoint $g$, then for every $y \in Y$, $X/y \simeq X/g(y)$, so $X/y$ has a final object and is therefore $W$-aspheric. Thus if $f$ has a right (left) adjoint, then $f$ is $W$-aspheric ($W$-coaspheric).

But even if $f$ has no right adjoint, there is still a generalized adjoint
\[
g : Y \longrightarrow X^{\wedge}
\]
and we have
\[
X/y \simeq X/g(y) \qquad \text{(for the fully faithful inclusion $X \hookrightarrow X^{\wedge}$),}
\]
so $\mathrm{Top}(X/y)$ is identified with the induced topos $\mathrm{Top}(X)/g(y)$. The condition $f$ $W$-aspheric therefore means that the sheaves $g(y)$ in $\mathrm{Top}(X)$ are
\msp{page 24}
$W$-aspheric. For many other equivalent formulations, I refer to Pursuing Stacks.

\textbf{Proposition 2.} \textit{Let $X \xrightarrow{f} Y \xrightarrow{g} Z$ be in} \textup{Cat}. \textit{If $f$ and $g$ are $W$-aspheric ($W$-coaspheric), likewise $gf$. If $gf$ and $f$ are $W$-aspheric ($W$-coaspheric), $g$ too.}

\textbf{Recall.} A morphism $p : S' \longrightarrow S$ in Cat is a \textit{local isomorphism} if $\forall\ s' \in \Ob S'$, the induced morphism
\[
p_{s'} : S'/s' \xrightarrow{\sim} S/s \qquad (s = p(s'))
\]
is an isomorphism: $p$ induces an isomorphism on the ``localized categories'' of $S'$ and $S$ at $s'$ and $s = p(s')$ respectively. It is immediate that this notion is stable under composition, that if $gf$ and $g$ are local isomorphisms, then $f$ is one too, that every isomorphism is a local isomorphism. Moreover, local isomorphisms are \textit{stable under base change,} which follows from the fact (trivial and important) that if we have a cartesian square in Cat

\begin{center}
\begin{tikzcd}[row sep=large]
X & X' \arrow[l, "q"'] \arrow[d, "f'"] \\
S \arrow[from=u, "f"'] & S', \arrow[l, "p"']
\end{tikzcd}
\end{center}

and if $x' \in \Ob X'$ has images $x$, $s'$, $s$ in $X$, $S'$, $S$, then the diagram

\begin{center}
\begin{tikzcd}[row sep=large]
X/x \arrow[d, "f/x"'] & X'/x' \arrow[l, "q/x'"'] \arrow[d, "f'/x'"] \\
S/s & S'/s' \arrow[l, "p/s'"']
\end{tikzcd}
\end{center}

is also cartesian. (So if $p/s'$ is \texttt{[an]} isomorphism, $q/x'$ too.)

\printedpage{21}
% head: A. Grothendieck : Dérivateurs, Ch. XII, Édition des Universités de Montpellier II et Paris VII
% foot: Version 7 décembre 2002

\msp{page 25}Let us return to proposition 2. The stability of the notion of aspheric morphism under composition follows at once from form a) of the condition. Likewise for the condition that if $gf$ and $f$ are $W$-aspheric, $g$ too - we must show that if $Z' \longrightarrow Z$ is \texttt{[a]} local isomorphism, then $g'$ is in $W$, now the hypothesis implies that $g'f'$ and $f'$ are.

\begin{center}
\begin{tikzcd}[row sep=large]
X \arrow[d, "f"'] & X' \arrow[l, "r"'] \arrow[d, "f'"] \\
Y \arrow[d, "g"'] & Y' \arrow[l, "q"'] \arrow[d, "g'"] \\
Z & Z' \arrow[l, "p"']
\end{tikzcd}
\end{center}

\textbf{Corollary 1.} \textit{The inverse image of an $f : X \longrightarrow Y$ which is $W$-aspheric (resp. $W$-coaspheric) under a local isomorphism (resp. under a local co-isomorphism) is $W$-aspheric ($W$-coaspheric).}

\begin{center}
\begin{tikzcd}[row sep=huge]
X \arrow[d] & X' \arrow[l] \arrow[d] & X'' \arrow[l] \arrow[d] \\
Y & Y' \arrow[l, "\text{local iso}"] & Y'' \arrow[l, "\text{local iso}"]
\end{tikzcd}
\end{center}

Follows formally from the definition by a) and from the fact that a composite of two local isomorphisms (resp. local co-isomorphisms) is likewise.

\textbf{NB} Of course, the notion of $W$-aspheric functor is not stable under arbitrary base change. For example, the inclusion $X = \{y_0\} \longrightarrow Y$ is $W$-aspheric if and only if the discrete categories $\mathrm{Hom}(y_0, y)$ are $W$-aspheric, i.e. (provided that we have W(5), the connectedness axiom) if and only if the $\mathrm{Hom}(y_0, y)$ have cardinality 1, i.e. $y_0$ \textit{initial object} of $Y$. But if $Y \neq \{y_0\}$, taking $y$ in $Y$, $y \neq y_0$, and setting $Y' = \{y\} \longrightarrow Y$, we have $X' = X \times_Y Y' = \emptyset$ and $\emptyset \longrightarrow e$ is not aspheric!

\msp{page 26}On the other hand, we expect the $W$-aspheric morphisms to be stable under \textit{$W$-smooth} base changes (which generalize local isomorphisms), and the $W$-coaspheric morphisms under $W$-proper base changes. Cf. below.

\textbf{Corollary 2.} \textit{Let $f : X \longrightarrow Y$ in} \textup{Cat} \textit{be an} \textup{equivalence} \textit{of categories. Then $f$ is $W$-aspheric and $W$-coaspheric.}

Indeed, $f$ has a left adjoint and a right adjoint, and we use the remark on page 23.

\textbf{Proposition 3.} \textit{Let $f : X \longrightarrow Y$ be in} \textup{Cat}. \textit{Equivalent conditions:}

\printedpage{22}
% head: A. Grothendieck : Dérivateurs, Ch. XII, Édition des Universités de Montpellier II et Paris VII
% foot: Version 7 décembre 2002

\begin{enumerate}
\item[\textit{a)}] \textit{For every $y \in \Ob Y$, the inclusion functor}
\[
X_y \longrightarrow X/y \tag*{$(*)$}
\]
\textit{is $W$-}co\textit{aspheric.}
\item[\textit{b)}] \textit{For every $x_0 \in X$, the induced morphism for the colocalized categories}
\[
x_0\backslash X \longrightarrow y_0\backslash Y \qquad \textit{(where $y_0 = f(x_0)$)}
\]
\textit{has $W$-aspheric fibers.}
\item[\textit{c)}] \textit{For every morphism $\boldsymbol{\Delta}^1 = Y' \longrightarrow Y$, let us consider $X' = X \times_Y Y'$ over $\boldsymbol{\Delta}^1$, the inclusion $X'_1 \longrightarrow X'$ of the closed fiber of $X'$ (fiber over the closed point of $Y'$) is $W$-coaspheric.}
\item[\textit{d)}] \textit{For every arrow $u_0 : y_0 \longrightarrow y_1$ in $Y$, and every $x_0$ in $X_{y_0}$, the category $X(x_0;u_0)$ formed of the arrows $u' : x_0 \longrightarrow x$ in $X$ with source $x_0$, and which lift $u_0$} \texttt{[rather} $u_0$\texttt{]}\textit{, is $W$-aspheric.}
\end{enumerate}
\DERmark{ENV0657}


\msp{page 27}\textsc{Proof.} To spell out condition a), we must spell out the category $\xi_0\backslash(X_y)$, for an arbitrary object $\xi_0 = (x_0, u_0)$ of $X/y$, $x_0 \in \Ob X$,
\[
u_0 : y_0 \longrightarrow y_1 = y \quad \text{in $Y$, with $y_0 = f(x_0)$.}
\]

This category $\xi_0\backslash X_{y_1}$ is then spelled out as the category of the $x_1$ in $X_{y_1}$, \textit{equipped} with an arrow from $\xi_0$ to the image $\xi_1 = (x_1, \mathrm{id}_{f(x_1)} = \mathrm{id}_{y_1})$ of $x_1$ in $X/y_1$ in $\xi_0$ \texttt{[``in} $\xi_0$\texttt{'' to be deleted]}, i.e. with an arrow $u$ from $x_0$ to $x_1$, such that $f(u)$ gives rise to the commutative triangle in $Y$, i.e. $f(u) = u_0$. % AUDIT resolved: regular CMTT12 editorial bracket; see corrections.csv, PDF leaf 32

\begin{center}
\begin{tikzcd}
x_1 \arrow[dd, dash] & & \\
& & x_0 \arrow[ull, "u"'] \arrow[dd, dash] \\
y_1 & y_1 \arrow[l, "\mathrm{id}"'] & \\
& & y_0 \arrow[ul, "u_0"'] \arrow[ull, "f(u)"]
\end{tikzcd}
\end{center}

So this is the category denoted $X(x_0;u_0)$ in d). We see at once that it is also isomorphic to the fiber of $x_0\backslash X \longrightarrow y_0\backslash Y$ at the element $u_0$ of $y_0\backslash Y$, referred to in b). Hence the equivalence of a) b) d). On the other hand, giving $u_0$ in $Y$ is equivalent to giving a functor $\boldsymbol{\Delta}^1 \longrightarrow Y$, as in c), and to express then that $\underbrace{X'_1}_{\simeq X_{y_1}} \longrightarrow X'$ is coaspheric, it suffices to express that the $x'\backslash X'_1$ ($x' \in \Ob X'$) are $W$-aspheric. Now this is trivial if $x' \in \Ob X'_1$ (initial object), so it suffices to verify it for $x'$ in $X'\backslash X'_1 = X'_0 \simeq X_{y_0}$. Then $x'$ is identified with an element $x_0$ of $X_{y_0}$, and we check at once that $x_0\backslash X'_1$ is isomorphic to the category $X(x_0;u_0)$ of d). This completes the proof.
\DERmark{ENV0658}


\printedpage{23}
% head: A. Grothendieck : Dérivateurs, Ch. XII, Édition des Universités de Montpellier II et Paris VII
% foot: Version 7 décembre 2002

\msp{page 28}\textbf{Definition 2.} Under the conditions of proposition 3, we say that $f$ is \textit{$W$-proper.} We say that $f$ is \textit{$W$-smooth} if $f^{\circ}$ is $W$-proper, i.e. if the inclusion morphisms $X_y \longrightarrow y\backslash X$ \texttt{[rather} $X_y \longrightarrow y\backslash X$\texttt{]} are $W$-aspheric etc. (conditions dual to those of proposition 3).

\textbf{Remark.} Condition a) of proposition 3 is satisfied in particular if the inclusion $(*)$ has a left adjoint (cf. remark on page 23). This means that the categories $X(x_0;u)$ have an initial object. This also means that the category $X$ over $Y$ is Cat-cofibered over $Y$ in the weak sense, i.e. that there exists, for $u_0 : y_0 \longrightarrow y_1$ an arrow in $Y$, and $x_0$ in $X_{y_0}$, an object $u_{0*}(x_0)$ in $X_{y_1}$ which represents the covariant functor $X_{y_1} \longrightarrow (\mathrm{Ens})$, $x_1 \longmapsto \mathrm{Hom}_u(x_0, x_1)$. Thus, if $X$ is Cat-cofibered in the weak sense over $Y$, and a fortiori if it is Cat-cofibered, then $f : X \longrightarrow Y$ is $W$-proper. Dually, if $f$ is weakly Cat-fibered, $f$ is \texttt{[}$W$\texttt{-]}smooth.
\DERmark{ENV0659}


\textbf{Proposition 4.} (\footnote[43]{Say also: If $X \xrightarrow{f} Y \xrightarrow{g} Z$, if $gf$ and $f$ $W$-smooth, then $g$ is smooth at the points of $\mathrm{Im}\, f$, ditto for proper.}) \textit{The notion of $W$-proper morphism is stable under composition, under base change. Local co-isomorphisms are proper (trivial from b)), a fortiori isomorphisms are proper, closed immersions are proper.}

\textit{Dual statement for smooth morphisms: stability under base change, under composition, local isomorphisms are smooth, in particular open immersions.}

\msp{page 29}Stability under base change is immediate from form c). Stability under composition \texttt{[is]} less trivial:

\begin{center}
\begin{tikzcd}[row sep=large]
{\scriptstyle x} &[-1.5em] X \arrow[d, "f"] & x\backslash X \arrow[l] \arrow[d] & \bar{X} \arrow[l] \arrow[d] \\
{\scriptstyle y} & Y \arrow[d, "g"] & y\backslash Y \arrow[l] \arrow[d] & \bar{Y} \arrow[l] \arrow[d] \\
{\scriptstyle z} & Z & z\backslash Z \arrow[l] & \{u\} \arrow[l]
\end{tikzcd}
\end{center}

In the diagram with cartesian squares opposite \texttt{[rather whose two squares on the right are cartesian]} (deduced from $x \in \Ob X$, whence $y = f(x)$, $z = g(y) = gf(x)$, and from $u \in \Ob z\backslash Z$), we must prove that $\bar{X}$ is $W$-aspheric. We know that $\bar{Y}$ is ($g$ proper) and that $\bar{X} \longrightarrow \bar{Y}$ has $W$-aspheric fibers. It therefore suffices to use proposition 5 below, and the immediate remark % AUDIT resolved: regular CMTT12 editorial bracket; see corrections.csv, PDF leaf 33

\textbf{Lemma.} \textit{If $f$ is proper, the same holds for $x\backslash f : x\backslash X \longrightarrow y\backslash Y$, for every $x$ in $X$, $y = f(x)$.}


\endgroup
% END UNIT Der-ChXIIa-012/Der-ChXIIa-012_p021-025.tex
% BEGIN UNIT Der-ChXIIa-012/Der-ChXIIa-012_p026-030.tex
\begingroup\makeatletter%
\newcommand{\printedpage}[1]{\marginpar{\footnotesize p.~#1}}%
\newcommand{\msp}[1]{\marginpar{\footnotesize [#1]}}%
\newcommand{\Cat}{\mathrm{Cat}}%
\newcommand{\Hot}{\mathrm{Hot}}%
\newcommand{\hot}{\mathrm{hot}}%
\newcommand{\Ob}{\mathrm{Ob}}%
\makeatother\let\printedpage\derpp\let\msp\dermsp%
\let\mspin\DERmspin%


\printedpage{24}
% head: A. Grothendieck : Dérivateurs, Ch. XII, Édition des Universités de Montpellier II et Paris VII
% foot: Version 7 décembre 2002
Indeed, if $u'$ is in $x\backslash X$ and $v'$ in $y\backslash Y$, and if $x'$, $y'$ are their images in $X$, $Y$ respectively, then
\[
u'\backslash(x\backslash X)\simeq x'\backslash X,\qquad v'\backslash(y\backslash Y)\simeq y'\backslash Y,
\]
and the arrow between the two is identified with the localization arrow
\[
x'\backslash X\longrightarrow y'\backslash Y.
\]

\textbf{Proposition 5.} \textit{Let $f : X\longrightarrow Y$ be in $\Cat$, suppose $f$ $W$-proper or $W$-smooth, and with $W$-aspheric fibers, then $f\in W$ (and even $f$ is} \textup{universally} \textit{in $W$, i.e. it is in $W$ and remains so under every base change).}

\leavevmode\msp{page 30}As the hypotheses are stable under base change, it suffices to prove $f\in W$. Suppose e.g. $X$ and $Y$ proper over $S$, let us prove that $f$ is $W$-\textit{aspheric}, i.e. that the $X/y$ ($y\in\Ob\,Y$) are $W$-aspheric. But $X_y\longrightarrow X/y$, being $W$-coaspheric by the properness hypothesis, is in $W$ (by W(2 bis)), and since $X_y$ is $W$-aspheric by hypothesis, $X/y$ too. But by W(2 bis), $f$ $W$-aspheric $\Longrightarrow f\in W$. The case $f$ smooth is dual.
\DERmark{ENV0660}


\textbf{Remark.} Thus, if $f$ is proper, if it has $W$-aspheric fibers it is ``universally $W$-aspheric'', i.e. is and remains $W$-aspheric under every base change. Conversely, if $f$ is universally in $W$ (not necessarily proper), then in particular its fibers are $W$-aspheric. But that of course doesn't imply that $f$ is proper, or smooth (\footnote[44]{But ``universally in $W$'' already implies ``universally $W$-aspheric'', since the notion of $W$-asphericity is defined by \texttt{[the]} proposition 1 a).}).

\textbf{Corollary.} \textit{If $W$ satisfies} \textup{W(2) (not just W(2 bis))}\textit{, then for an $S$-morphism $X\longrightarrow Y$ of categories which are both $W$-proper over $S$, or both $W$-smooth, if $f_s : X_s\longrightarrow Y_s$ is in $W$ for every $s\in S$, then $f\in W$.}

\leavevmode\msp{page 31}Suppose e.g. $f$ proper, let us consider, for $s\in\Ob\,S$
\[
\begin{tikzcd}[row sep=large, column sep=large]
X_s \arrow[r, "f_s"] \arrow[d, "i_s^X"'] & Y_s \arrow[d, "i_s^Y"] \\
X/s \arrow[r, "f/s"] & Y/s,
\end{tikzcd}
\]
as $X$, $Y$ \texttt{[are]} proper, the vertical arrows are in $W$ (much less than properness would suffice \dots), hence $f_s\in W\Longrightarrow f/s\in W$. Thus, ``$f$ is in $W$ locally over $S$'', which implies $f\in W$ by W(2).

\printedpage{25}
% head: A. Grothendieck : Dérivateurs, Ch. XII, Édition des Universités de Montpellier II et Paris VII
% foot: Version 7 décembre 2002
\textbf{Proposition 6.} \textit{Let}
\[
\begin{tikzcd}[row sep=large, column sep=large]
X & X' \arrow[l, "q"'] \arrow[d, "f'"] \\
S \arrow[from=u, "f"'] & S', \arrow[l, "p"]
\end{tikzcd}
\]
\textit{be cartesian in} $\Cat$. \textit{If $f$} \texttt{[is]} \textit{$W$-aspheric, $p$ $W$-smooth, then $f'$ is $W$-aspheric.}

We must prove \texttt{[that]} $X'\times_{S'}S'/s'$ is $W$-aspheric ($s'\in\Ob\,S'$). Let $s=p(s')$, we have
\[
\begin{tikzcd}[column sep=small]
S & & S' \arrow[ll, "p"'] & \\
& S/s \arrow[ul] & & S'/s', \arrow[ul] \arrow[ll, "p_0"]
\end{tikzcd}
\]
with $p_0$ $W$-smooth (lemma p.~29 \texttt{[rather 29]}) and with $W$-aspheric fibers, hence a universal $W$-equivalence (prop.~5). As the inverse image of $X$ over $S/s$ is $W$-aspheric (hypothesis for $f$), it \texttt{[follows]}
\DERmark{ENV0661}


\leavevmode\msp{page 32}that in the cartesian square opposite \texttt{[below]}
\[
(*)\qquad\qquad
\begin{tikzcd}[row sep=large, column sep=large]
X/s \arrow[d] & X'/s' \arrow[l, "W"'] \arrow[d] \\
S/s & S'/s' \arrow[l, "W"']
\end{tikzcd}
\]
the horizontal arrows are in $W$, and the first vertical arrow is so by the hypothesis on $f$, hence the second also, hence $X'/s'$ \texttt{[is]} $W$-aspheric.

\textbf{Corollary} (generalizes \texttt{[the]} prop.~6, obtained by taking $Y=S$). \textit{Let $p : S'\longrightarrow S$ be} \texttt{[a]} \textit{$W$-smooth morphism in} $\Cat$. \textit{Then the base change functor}
\[
\Cat/S\longrightarrow\Cat/S'
\]
\textit{is compatible with the localizers $W^{\mathrm{g}}_S$, $W^{\mathrm{g}}_{S'}$, i.e. transforms $f : X\longrightarrow Y$ which is in $W$} locally \textit{over $S$, into $f' : X'\longrightarrow Y'$ in $W$ locally over $S'$.}

This time, we look at the diagram similar to $(*)$
\[
\begin{tikzcd}[column sep=small]
X/s \arrow[dd] \arrow[dr, "f/s"] & & X'/s' \arrow[ll, "W"'] \arrow[dd] \arrow[dr, "f'/s'"] & \\
& Y/s \arrow[dl] & & Y'/s' \arrow[ll, "W"', crossing over] \arrow[dl] \\
S/s & & S'/s', \arrow[ll, "W"'] &
\end{tikzcd}
\]

\printedpage{26}
% head: A. Grothendieck : Dérivateurs, Ch. XII, Édition des Universités de Montpellier II et Paris VII
% foot: Version 7 décembre 2002
where the horizontal arrows are in $W$, hence $f/s\in W$ (hypothesis) implies $f'/s'\in W$.

\leavevmode\msp{page 33}Dually, if $p$ is $W$-proper, the inverse image functor is compatible with the localizers $W^{\mathrm{d}}_S$, $W^{\mathrm{d}}_{S'}$.

If $f : X\longrightarrow S$ is $W$-proper, we are going to define a canonical functor
\[
\begin{aligned}
\hot_*^{X/S,W} : S &\quad\longrightarrow\quad \Hot_W\overset{\mathrm{def}}{=}(\Cat)W^{-1}\\
\text{given on objects by}\qquad s &\quad\longmapsto\quad \hot_W(X_s).
\end{aligned}
\]
\DERmark{ENV0662}

We must define, for an arrow $u$ in $S$, $u : s\longrightarrow s'$, an arrow
\[
\hot_W(X_s)\longrightarrow\hot_W(X_{s'})
\]
in $\Hot_W$ - arrow which we can denote by $\hot_*^{X/S,W}(u)$. We define it by the diagram
\[
\begin{tikzcd}[row sep=large, column sep=large]
X_s \arrow[r, dotted] \arrow[d, "i_s^X"] & X_{s'} \arrow[d, "i_{s'}^X"] \\
X/s \arrow[r, "X/u"] & X/s',
\end{tikzcd}
\]
where the vertical arrows are $W$-coaspheric, hence in $W$. They are therefore isomorphisms in $\Hot_W$, whence the unique dotted arrow in $\Hot_W$, making the square commutative. From this definition it is immediate that
\[
u\longmapsto\hot_W(u)
\]
is compatible with composition, and units.

This functor is thus defined assuming

\leavevmode\msp{page 34}only that the
\[
X_y\longrightarrow X/y
\]
are in $W$, which is distinctly weaker than proper. If we suppose that this hypothesis is satisfied ``universally'' (i.e. remains satisfied after every base change $S'\longrightarrow S$) (\footnote[45]{We shall say that $f$ is $W$-\textit{preproper.}}), then we therefore have, for every $S'$ in $\Cat/S$, a functor
\[
\hot_*^{X'/S',W} : S'\longrightarrow\Hot_W,
\]
and we check at once that it is none other than the composite in

\printedpage{27}
% head: A. Grothendieck : Dérivateurs, Ch. XII, Édition des Universités de Montpellier II et Paris VII
% foot: Version 7 décembre 2002
\[
(*)\qquad\qquad
\begin{tikzcd}[column sep=small]
S' \arrow[rr, "p"] \arrow[dr, "\hot_*^{X'/S',W}"'] & & S \arrow[dl, "\hot_*^{X/S,W}"] \\
& \Hot_W &
\end{tikzcd}
\]
\DERmark{ENV0663}\DERmark{ENV0664}

\texttt{[rather} $\hot_*^{X'/S',W}$ \texttt{resp.} $\hot_*^{X/S,W}$\texttt{]}. We have dual observations if $X$ is smooth over $S$, more generally if the inclusions
\[
X_s\longrightarrow s\backslash X
\]
are in $W$. We then find a contrafunctor
\[
\left\{
\begin{aligned}
S^\circ &\quad\longrightarrow\quad \Hot_W\\
s &\quad\longmapsto\quad \hot_W(X_s)\\
u &\quad\longmapsto\quad \hot^*_{X/S,W}(u)\ \text{or}\ \hot^*_{X/S}(u)
\end{aligned}
\right.
\]

\leavevmode\msp{page 35}and a compatibility similar to $(*)$, if the hypothesis is satisfied \textit{universally} over $S$, e.g. $f$ $W$-smooth.

\textbf{Proposition 7.} \textit{Suppose $f : X\longrightarrow S$ both $W$-proper and $W$-smooth, more generally that the morphisms}
\[
X_s\longrightarrow X/s,\qquad X_s\longrightarrow s\backslash X\qquad(s\in S)
\]
\textit{are in $W$, and that this hypothesis remains satisfied after every base change. Then the corresponding functors} % [sic] verifiée -> vérifiée, accent missing
\[
\begin{aligned}
\lambda &\;=\;\hot_*^{X/S,W} : &&\quad S \quad\longrightarrow\quad \Hot_W\\
\lambda^\circ &\;=\;\hot^*_{X/S,W} : &&\quad S^\circ \quad\longrightarrow\quad \Hot_W
\end{aligned}
\]
\textit{transform every arrow $u$ of $S$ into} \texttt{[an]} \textit{isomorphism of} $\Hot_W$\textit{, and we have}
\[
\lambda^\circ(u)=\lambda(u)^{-1}\qquad\textit{in}\quad\Hot_W,
\]
\textit{i.e. $\lambda(u)$ and $\lambda^\circ(u)$ are inverses of each other.} % [sic] λ^c irc(u) -> λ°(u), command \circ misprinted as superscript c followed by "irc"
\DERmark{ENV0665}


\textsc{Proof}. The datum of $u : s_0\longrightarrow s_1$ is equivalent to that of
\[
S'=\boldsymbol{\Delta}^1\longrightarrow S,
\]
whence
\[
X'=X\times_S S'\qquad\text{over } S'=\boldsymbol{\Delta}^1,
\]
with the two fibers
\[
X'_0\simeq X_{s_0},\qquad X'_1\simeq X_{s_1}.
\]
The hypothesis ensuring the existence of $\lambda$ is expressed by the condition that
\[
X'_1\longrightarrow X'/1\simeq X'
\]
\printedpage{28}
% head: A. Grothendieck : Dérivateurs, Ch. XII, Édition des Universités de Montpellier II et Paris VII
% foot: Version 7 décembre 2002
\msp{page 36}is in $W$, so that the diagram
\[
X'_0\xrightarrow{i_0}X'\xleftarrow[W]{i_1}X'_1
\]
defines an arrow $\lambda(u) : X'_0\longrightarrow X'_1$ in $\Hot_W$, namely $\lambda(u)=i_1^{-1}i_0$. The hypothesis which ensures the existence of $\lambda^\circ$ is expressed by the symmetric condition $i_0\in W$, which makes it possible to define the arrow $\lambda^\circ(u) : X'_1\longrightarrow X'_0$ by $\lambda^\circ(u)=i_0^{-1}i_1$. When the two conditions are satisfied simultaneously, we therefore see that $\lambda(u)$, $\lambda^\circ(u)$ are inverses of each other.

\textbf{Remarks.} To define $\lambda$ ``universally'', it was not indispensable to consider the $X/s$ and the inclusions $X_s\longrightarrow X/s$, instead of that it sufficed to look at the base changes $S'=\boldsymbol{\Delta}^1\xrightarrow{u}S$, and to make the hypothesis that for every such base change, the inclusion $X'_1\xrightarrow{i_1}X'=X_{(\boldsymbol{\Delta}^1,u)}$ is in $W$. This hypothesis is automatically stable under arbitrary base change. (The $W$-properness of $f : X\longrightarrow S$ means that $X'_1\longrightarrow X'$ is not only in $W$, but even $W$-coaspheric.) Dually, for the ``universal'' definition of $\lambda^\circ$, it would suffice that for every base change $\boldsymbol{\Delta}^1\longrightarrow S$ as above, the inclusion $X'_0\xrightarrow{i_0}X'$ is in $W$.

\leavevmode\msp{page 37}The weakest reasonable notion of $W$-fibered, seems to me to be the one where we suppose both things at once. The hypotheses we have made up to now, excluding the strong form of $W$-properness or of $W$-smoothness, always amount to supposing that the \textit{base change by $f : X\longrightarrow S$} transforms certain arrows in $\Cat/S$ into \texttt{[an element of]} $W$ : To define $\lambda$, these arrows are
\[
\text{either } \{s\}\longrightarrow S/s,\qquad \text{or } \{1\}\longrightarrow\boldsymbol{\Delta}^1(\longrightarrow S),
\]
to define $\lambda^\circ$, they are
\[
\text{either } \{s\}\longrightarrow s\backslash S,\qquad \text{or } \{0\}\longrightarrow\boldsymbol{\Delta}^1(\longrightarrow S),
\]
the ``either'' depending on the definition we have adopted. In these examples, it is always a matter of the inclusion of the final object of a category $S'$, into the said category, or conversely of the initial object. This shows moreover (if we require universality), that the first hypothesis (ensuring $\lambda$) is \textit{exactly} that (case of the inclusion of a final object $s'$ - since for such an object, we have $S'/s'\simeq S'$), whereas the second hypothesis requires less, only for the case of $\boldsymbol{\Delta}^1$. And dually for $\lambda^\circ$ \texttt{[rather $\lambda^\circ$]}. But it should be noted that in practice, the most convenient way, and by far,
\DERmark{ENV0666}


\leavevmode\msp{page 38}to verify such properties, is either by the $W$-properness of $f$, or by $W$-smoothness - which reduces to verifying that certain well-specified categories are $W$-aspheric, without having to worry about base changes and about verifying that the arrows constructed in terms of arbitrary base changes are in $W$. And the use of $W$-properness and $W$-smoothness is all the more congenial, as these notions are stable under composition,


\endgroup
% END UNIT Der-ChXIIa-012/Der-ChXIIa-012_p026-030.tex
% BEGIN UNIT Der-ChXIIa-012/Der-ChXIIa-012_p031-035.tex
\begingroup\makeatletter%
\derleqno%
\newcommand{\printedpage}[1]{\marginpar{\footnotesize p.~#1}}%
\newcommand{\msp}[1]{\marginpar{\footnotesize [#1]}}%
\newcommand{\fl}{\longrightarrow}%
\newcommand{\uHom}{\underline{\mathrm{Hom}}}%
\newcommand{\uFib}{\underline{\mathrm{Fib}}}%
\newcommand{\Ob}{\operatorname{Ob}}%
\newcommand{\HOT}{\mathrm{HOT}}%
\newcommand{\Cat}{\mathrm{Cat}}%
\makeatother\let\printedpage\derpp\let\msp\dermsp%
\let\mspin\DERmspin%


\printedpage{29}
% head: A. Grothendieck : Dérivateurs, Ch. XII, Édition des Universités de Montpellier II et Paris VII
whereas all the derived notions, as well as the notions of fibrations to boot, have an unfortunate tendency \textit{not} to be stable under composition.

We could call pseudo-$W$-proper resp. quasi-$W$-proper (a little better than pseudo) what we have just brought out - I hear snickering! And if we want both at once - $\lambda$ and $\lambda^{\circ}$ associated as we said? I see the same principle of description with $S''\fl S'$ over $S$, of requiring that $X''\fl X'$ be transformed into \texttt{[elements of]} $W$. With the following choices which look nice

\begin{itemize}
\item[a$_1$)] \msp{page 39}$Y'$ has a final object, $Y''=e$. The definition thus comes down to the condition that the
\[ X_s\fl X/s, \qquad X/s'\fl X/s \]
(for $s\in\Ob S$ where $u : s'\fl s$ in $\mathrm{S}$) be in $W$, and that it stay that way after every base change $S'\fl S$.

\item[a$_2$)] $Y'$ has an initial object, $Y''=e$, this is the dual condition by dint of
\[ X_s\fl s\backslash X, \qquad s\backslash X\fl s'\backslash X. \]

\item[a$_0$)] $Y'=\boldsymbol{\Delta}^1$, $Y''=e$ (\footnote[46]{NB a$_1$), a$_2$) are denoted the same way on p. 4, but a$_0$) was forgotten.}).

\item[b)] $Y'$ is contractible, $Y''$ too (or if preferred, $Y'$ contractible, $Y''=e$)
\end{itemize}

thus the condition comes down to this: for every \textit{contractible} $Y'$ over $Y$ (notion independent of $W$, that is its main advantage) the inclusions of the fibers of $X'=X\times_Y Y'$ over $Y'$ into $X'$ are \texttt{[in]} $W$.

I feel that even for general $W$, b) remains in the realm of the verifiable, whereas we go soft if we go to the extreme of generality in this direction:

\begin{itemize}
\item[b$'$)] As b), but assuming only $Y'$ $W$-aspheric (\footnote[47]{This is essentially a$'$) of p. 4}).
\end{itemize}

Finally (condition a) of page 4)

\begin{itemize}
\item[c)] If $Y''\fl Y'$ is \texttt{[in]} $W$ without more, $X''\fl X'$ too.
\end{itemize}
% foot: Version 7 décembre 2002

\printedpage{30}
% head: A. Grothendieck : Dérivateurs, Ch. XII, Édition des Universités de Montpellier II et Paris VII
\msp{page 40}The only one of these notions that has the slightest desire to be stable under composition is the strongest c). On the other hand, the notion of a $W$-perfect arrow (cf. def. 3) is spontaneously stable under composition:

\textbf{Definition 3.} We say that $f : X\fl S$ is $W$-\textit{perfect} if it is both $W$-proper and $W$-smooth.

In the case $W=W_\infty$, all the variants on the preceding page, with the sole exception of a$_0$), are equivalent. It would be interesting to bring out minimal conditions on $W$ to ensure that this is so. (And the question arises whether in the case of $W_\infty$, a$_0$) suffices to imply the rest - I hardly expect it \dots)

In the general case, we have the diagram of implications (and even so I have limited myself)
\[
\begin{tikzcd}[column sep=small]
 & & \mathrm{c} \arrow[d, Rightarrow] & & \text{ultrastrong $W$-fibration} \\
 & & \mathrm{b}' \arrow[dd, Rightarrow] & & \text{$W$-strong $W$-fibration} \\
 & \text{$W$-perfect} \arrow[dr, Rightarrow] & & & \\
 & & \mathrm{b} \arrow[dl, Rightarrow] \arrow[dr, Rightarrow] & & \text{hot-strong $W$-fibration (\footnotemark[48])} \\
\text{left $W$-fibration} & \mathrm{a}_1 \arrow[dr, Rightarrow] & & \mathrm{a}_2 \arrow[dl, Rightarrow] & \text{right $W$-fibration} \\
 & & \mathrm{a}_0 & & \text{$W$-pseudofibration}
\end{tikzcd}
\]
\footnotetext[48]{It would seem that this is, in the end, the ``right notion'' of $W$-fibration, \textit{equivalent} to a$_1$ and a$_2$, which would be equivalent to each other. Cf. prop. 9, p. 48.}

\texttt{[``$W$-perfect $\Longrightarrow$ b'' circled]}.

\msp{page 41}\textbf{Proposition 8.} \textit{Let $f : X\fl S$ be in} Cat. \textit{The following conditions are equivalent.}

\begin{itemize}
\item[\textit{a)}] \textit{$f$ is a left $W$-fibration, i.e. condition} a$_1$ \textit{of page 39 is satisfied.}

\item[\textit{b)}] \textit{$f$ is $W$-preproper, and $\forall\; s_0\fl s_1$ in $S$, $X/s_0\fl X/s_1$ is in $W$.}

\item[\textit{c)}] \textit{$f$ is $W$-preproper, and a $W$-pseudofibration, i.e. condition a$_0$ of page 39 is satisfied.} % AUDIT resolved: CMTI12 text-italic a with CMR8 subscript; see corrections.csv, PDF leaf 40, ms41
\end{itemize}

(\footnote[49]{\textbf{Corollary.} \textit{When $f$ is} Cat-\textit{cofibrant, then these conditions are also equivalent to this: $\forall\; u : s_0\fl s_1$ in $S$, $u_* : X_{s_0}\fl X_{s_1}$ is in $W$.}}).
% foot: Version 7 décembre 2002

\printedpage{31}
% head: A. Grothendieck : Dérivateurs, Ch. XII, Édition des Universités de Montpellier II et Paris VII
\textbf{NB} These conditions are satisfied if $f$ is at the same time $W$-presmooth, cf. prop. 7.

\textsc{Proof}. We have a) $\Longrightarrow$ b) trivially (take $Y''=Y/s$, $Y''=Y/s'$ \texttt{[connection with the assertion? Cf. p. 4.]}).

b) $\Longleftrightarrow$ c) by the reasoning of proposition 7, for in terms of $X'=X_u=X\times_S(\boldsymbol{\Delta}^1,u)$, condition b) for a given $u : s_0\fl s_1$ is interpreted in the commutative diagram
\[
\begin{tikzcd}[column sep=small, row sep=small]
X'_0 \arrow[dr, "i'_0\in W"] \arrow[dd, "\wr"] & & X'_1 \arrow[dr, "i'_1\in W"] \arrow[dd, "\wr"] & \\
 & \underset{(\simeq X'_0)}{X'/0} \arrow[rr, "\varphi'", crossing over] \arrow[dd, "p_0"] & & X'/1\ (\simeq X') \arrow[dd, "p_1"] \\
X_{s_0} \arrow[dr, "i_0\in W"'] & & X_{s_1} \arrow[dr, "i_1\in W"] & \\
 & X/s_0 \arrow[rr, "\varphi"] & & X/s_1
\end{tikzcd}
\]

\texttt{[there are a few redundant $W$ in the original diagram]} where $i_0$, $i_1$, $i'_0$, $i'_1$ are in $W$ because of the hypothesis of $W$-preproperness (stable under base change).

Thus $p_0$, $p_1$ are \texttt{[in]} $W$, thus $\varphi\in W$ if and only if $\varphi'\in W$, now $\varphi'\in W$ is condition a$_0$.

\msp{page 42}Let us finally prove that b), c) imply a). Since c) (hence also b)) is stable under base change, it suffices to show this: If $S$ has a final object $e$, then for every $s\in\Ob S$, the inclusion $X_s\xrightarrow{\ i\ }X$ is in $W$. Now we have the commutative diagram
\[
\begin{tikzcd}
X_s \arrow[r, "i_s\in W"] \arrow[dr, "i"'] & X/s \arrow[d, "u\,\in\,W\text{ by b) }\texttt{[rather by b)]}"] \\
 & X/e=X,
\end{tikzcd}
\]
\DERmark{ENV0667}

which indeed shows $i=ui_s\in W$.

Dually, if $f$ is $W$-presmooth, then condition a$_0$ ($f$ a $W$-pseudofibration) is equivalent to a$_2$ ($f$ a right $W$-fibration), and also means that the
\[ s_1\backslash X\fl s_0\backslash X \]
are in $W$. But the notion of ``$W$-fibration'' that we land on (if we are not content with the paltry a$_0$!) is different in the two cases dual to each other. We would like in both cases to find ``hot-strong'' $W$-fibrations (condition b), at least by going all out (if necessary) - i.e. for $f$ $W$-perfect. So it comes down to this.

\textbf{Question:} \textit{Let $X\xrightarrow{\ f\ }S$, with $f$ $W$-perfect, $S$ contractible, $s\in\Ob S$. The inclusion $X_s\fl X$, is it in $W$?}
% foot: Version 7 décembre 2002

\printedpage{32}
% head: A. Grothendieck : Dérivateurs, Ch. XII, Édition des Universités de Montpellier II et Paris VII
\msp{page 43}We are going to assume only that $f$ is $W$-presmooth and a $W$-pseudofibration \texttt{[rather that $f$ is simply a right $W$-fibration, i.e that $f$ is $W$-presmooth and a $W$-pseudofibration]} (cf. prop. 8). We are going to reduce to the case where $X$ is even fibrant over $S$, by introducing $\Phi_S(X)$ \texttt{[fibered category defined by the functor $S^{\circ}\fl\Cat$, $s\longmapsto s\backslash X$]} and % [sic] i.e que -> i.e. que, missing period of the abbreviation
\[ X\xrightarrow{\ i_X\ }\Phi_S(X). \]
\DERmark{ENV0668}


The trouble is that I have done the theory of $\HOT_W$ up to now with the hypothesis (which turns out to be idiotic) that $W\supseteq W_\infty$ - I will therefore be obliged to review quickly under what more reasonable hypotheses on $W$ the familiar results are valid. For example I would like $i_X\in W$ (\footnote[50]{in fact, $i_X\in(W_\omega)^{\mathrm{d}}_S\subseteq W^{\mathrm{d}}_S\subseteq W_S$.}), so that we have
\[
\begin{tikzcd}[column sep=large, row sep=large]
X_s \arrow[r, "i_s\in W", "\text{since $X$ $W$-smooth}"'] \arrow[d, "j^X_s"'] & \Phi_S(X)_s=s\backslash X \arrow[d, "j_s^{\Phi_S(X)}"] \\
X \arrow[r, "i_X\in W"] & \Phi_S(X),
\end{tikzcd}
\]
we therefore see that to prove that the inclusion $j^X_s$ is \texttt{[in]} $W$, it suffices to prove the same thing for $\Phi_S(X)$, which is fibrant (not just $W$-smooth). Moreover, the pseudofibration hypothesis on $X$ surely passes to $\Phi_S(X)$ (I leave the two verifications for later (\footnote[51]{O.K.})). We are thus reduced to the case of
\[ X\in\Ob\uFib_0S, \qquad S\text{ contractible,} \]
to prove that the inclusion
\[ X_s\fl X \]
\msp{page 44}is in $W$. I am going to prove this:

\textbf{Lemma.} \textit{Let $S$ be in} Cat, \textit{$X\in\Ob\uFib_0S$, and let us consider, for an $S'$ in} Cat,
\[
\begin{aligned}
f &\quad\longmapsto\quad f^*(X/S)=X\times_S(S',f)\\
\uHom(S',S)^{\circ} &\quad\longrightarrow\quad \uFib_0S'.
\end{aligned}
\]
\textit{(This functor is defined thanks to the fact that $X$ is fibrant over $S$, and can therefore be interpreted, up to $S$-equivalence, in terms of a functor $S^{\circ}\fl$} Cat \textit{\dots\ (\footnote[52]{\textbf{NB} $\Phi_S(X)$ was even already split, defined by the functor $s\longmapsto s\backslash X$.})). The arrows of $\uHom(S',S)$ are transformed into arrows $\in W^{\mathrm{d}}_{S'}$.}
% foot: Version 7 décembre 2002

\printedpage{33}
% head: A. Grothendieck : Dérivateurs, Ch. XII, Édition des Universités de Montpellier II et Paris VII
Indeed, in terms of functors
\[ S'^{\circ}\fl\Cat, \]
if $X$ is defined by
\[ S^{\circ}\xrightarrow{\ \mathcal{X}\ }\Cat, \]
\DERmark{ENV0669}

then $f^*(X)$ is defined by $\mathcal{X}\circ f$, and the arrow
\DERmark{ENV0670}

\[ \mathcal{X}\circ g^{\mathrm{op}}\longrightarrow\mathcal{X}\circ f^{\mathrm{op}} \]
\DERmark{ENV0671}\DERmark{ENV0672}

defined by $u : f\fl g$ is none other than $\mathcal{X}*u$. Its value at $s'\in\Ob S'$ is obtained by taking
\DERmark{ENV0673}

\[ u(s') : f(s')\fl g(s'), \]
whence by $\mathcal{X}$
\DERmark{ENV0674}

\[ \mathcal{X}(u(s')) : \underbrace{\mathcal{X}(g(s'))}_{(\mathcal{X}\circ g)(s')}\fl\underbrace{\mathcal{X}(f(s'))}_{(\mathcal{X}\circ f)(s')} \]
\DERmark{ENV0675}\DERmark{ENV0676}\DERmark{ENV0677}\DERmark{ENV0678}\DERmark{ENV0679}

(\footnote[53]{to bring out a separate statement.}). But the hypothesis $X\in\uFib_0S$ means that the $\mathcal{X}(v)$, $v\in\mathrm{Fl}(S)$, are all in $W$. Thus
\DERmark{ENV0680}

\msp{page 45}$\mathcal{X}*u$ is in $W$ fiber by fiber, thus in $W^{\mathrm{d}}_{S'}$.

\textbf{Corollary 1.} \textit{The functor of the lemma, composed with}
\DERmark{ENV0681}

\[ \uFib_0S'\fl\HOT^{\mathrm{lc}}_W(S')\overset{\mathrm{def}}{=}(\uFib_0S')(W^{\mathrm{d}}_{S'})^{-1} \]
\DERmark{ENV0682}

\textit{factors through the fundamental groupoid of} $\uHom$ \textit{:}
\[
\begin{tikzcd}[row sep=huge]
\uHom(S',S)^{\circ} \arrow[r] \arrow[d] & \uFib_0S' \arrow[d] \\
\textstyle\prod(\uHom(S',S))^{\circ} \arrow[r] & \HOT^{\mathrm{lc}}_W(S').
\end{tikzcd}
\]
\textit{If $f$, $g$ are in the same connected component of $\uHom(S',S)$, then $f^*X$ and $g^*X$ are isomorphic (noncanonically, in general) in} $\HOT^{\mathrm{lc}}_W(S')$.

It is cleverer not to fix $X$ and to consider
\[ \uHom(S',S)\fl\uHom(\HOT^{\mathrm{lc}}_W(S),\HOT^{\mathrm{lc}}_W(S')) \tag*{$(*)$} \]
corresponding to
\[
\begin{aligned}
\uHom(S',S) &\quad\times\quad \HOT^{\mathrm{lc}}_W(S) &&\quad\longrightarrow\quad &&\HOT^{\mathrm{lc}}_W(S')\\
(f &\quad,\quad\ \xi) &&\quad\longmapsto\quad && f^*(\xi).
\end{aligned}
\]
The lemma will say that the arrows of $S'\fl S$ are transformed by $(*)$ into isomorphisms, whence
\[ \prod\uHom(S',S)\fl\uHom(\HOT^{\mathrm{lc}}_W(S),\HOT^{\mathrm{lc}}_W(S')), \]
% foot: Version 7 décembre 2002


\endgroup
% END UNIT Der-ChXIIa-012/Der-ChXIIa-012_p031-035.tex
% BEGIN UNIT Der-ChXIIa-012/Der-ChXIIa-012_p036-040.tex
\begingroup\makeatletter%
\newcommand{\printedpage}[1]{\marginpar{\footnotesize p.~#1}}%
\newcommand{\msp}[1]{\marginpar{\footnotesize [#1]}}%
\newcommand{\HOTlc}{\mathrm{HOT}^{\mathrm{lc}}_W}%
\newcommand{\Hot}{\mathrm{Hot}}%
\newcommand{\Hom}{\mathrm{Hom}}%
\newcommand{\hot}{\mathrm{hot}}%
\newcommand{\id}{\mathrm{id}}%
\newcommand{\Wd}{W^{\mathrm{d}}_S}%
\newcommand{\Wf}{W^{\mathrm{f}}_S}%
\makeatother\let\printedpage\derpp\let\msp\dermsp%
\let\mspin\DERmspin%


\printedpage{34}
% head: A. Grothendieck : Dérivateurs, Ch. XII, Édition des Universités de Montpellier II et Paris VII
% foot: Version 7 décembre 2002

\msp{page 46}
so if $f,g : S' \rightrightarrows S$ are two homotopic functors, then
\[ f^* \simeq g^* : \HOTlc(S) \rightrightarrows \HOTlc(S'). \]

The following follows at once

\textbf{Corollary 2.} \textit{If $f$ is a homotopism, $f^* : \HOTlc(S) \longrightarrow \HOTlc(S')$ is an equivalence of categories.}

Indeed, taking $g : S \longrightarrow S'$ such that
\[ fg \sim \id_S, \qquad gf \sim \id_{S'}, \]
we find that
\[ \begin{aligned}
g^*f^* \simeq (\id_S)^* &= \id_{\HOTlc(S)} \\
f^*g^* \simeq (\id_{S'})^* &= \id_{\HOTlc(S')},
\end{aligned} \]
so we find that (adjusting these functorial isomorphisms if need be) $f^*$ and $g^*$ are equivalences quasi-inverse to one another.

I have the impression (\footnote[54]{O.K.}) that all this should not use more than the standard axioms on $W$ as a basic localizer, at most W(2) instead of the weak form W(2 bis).

Let us return to $X$ smooth over contractible $S$. Corollary 2 implies that
\[ \begin{aligned}
\Hot_W &\longrightarrow \HOTlc(S) \\
\xi &\longmapsto p^*_S(\xi),
\end{aligned} \]
where $p_S : S \longrightarrow e$, is an equivalence. So in the category $\HOTlc(S)$, $X$ is isomorphic \msp{page 47} to a Cat-fibered object of the form $Z \times S \xrightarrow{\mathrm{pr}_2} S$, with $Z$ in Cat. Now for $X_1 = Z \times S \longrightarrow S$, it is immediate that $X_{1\,s} \longrightarrow X_1$ is in $W$, for a Cartesian product of arrows \texttt{[in]} $W$ (namely $\id_Z$ and $s \longrightarrow S$) is in $W$. (In fact, the inclusion is even a homotopism!) Let us now note that if $X_0$, $X_1$ are two objects in $\underline{\mathrm{Liss}}_0S$, such that there exists an $S$-arrow $X_0 \longrightarrow X_1$ which is in $\Wd$, then $X_{0\,s} \longrightarrow X_0$ is in $W$ if and only if $X_{1\,s} \longrightarrow X_1$ is. Indeed, we have
\[ \begin{tikzcd}
X_{0\,s} \arrow[r, "\in W"] \arrow[d, "i_s^{X_0}"'] & X_{1\,s} \arrow[d, "i_s^{X_1}"] & \\
X_0 \arrow[r, "\in W", ""{name=A, below}] & X_1 & \\
 & & \text{\footnotesize\begin{tabular}{@{}l@{}}here we probably use\\ W(2), not only W(2 bis),\end{tabular}}
\arrow[from=A, to=3-3, bend right]
\end{tikzcd} \]
\printedpage{35}
% head: A. Grothendieck : Dérivateurs, Ch. XII, Édition des Universités de Montpellier II et Paris VII
% foot: Version 7 décembre 2002
which proves the result.

It is true that the fact that $X = X_0$ and $X_1 = Z \times S$ are isomorphic in $\HOTlc(S)$ does not imply, by itself, the existence of a \msp{page 48} chain of isomorphisms coming from arrows in $\Wd$. But in fact, if in the statement of corollary 2 above we compare $X$ in $\underline{\mathrm{Fib}}_0S$ with $g^*f^*(X)$, thanks to a homotopy $\id_S \sim fg$, we do indeed find a sequence of functors $\varphi_i : S \longrightarrow S$, with $\varphi_0 = \id_S$, $\varphi_N = fg$, and $\varphi_i$ and $\varphi_{i+1}$ linked by an arrow in $\underline{\mathrm{Hom}}(S,S)$, so we have a sequence of objects $X_i = \varphi_i^*(X)$ of $\underline{\mathrm{Fib}}_0S$ (with $X_0 = X$, $X_N = g^*f^*X$), and $X_i$ and $X_{i+1}$ linked by an arrow \textit{in} $\underline{\mathrm{Fib}}_0S$ (not only in $\HOTlc(S)$), which arrow is in $\Wd = \Wf$. So we win, subject to various routine checks (\footnote[55]{O.K.}). So modulo these, we will have

\textbf{Proposition 9} (\footnote[56]{It's ok assuming that $W$ is \textit{strongly} saturated, i.e. $\hot_W(f)$ isomorphism $\Longrightarrow f \in W$.},\footnote[57]{So left $W$-fibration $\Longleftrightarrow$ right $W$-fibration $\Longleftrightarrow$ $W$-fibration.}). \textit{Suppose that $X$ in $\mathrm{Cat}/S$ is a left $W$-fibration (cf. prop. 8). Then it is even a $W$-fibration (in the sense of b), page 40, where it is called ``$W$-fibration \textup{hot}-strong''), in other words, if $S'' \longrightarrow S'$ is over $S$ with $S''$, $S'$ contractible, then $X'' \longrightarrow X'$ is in $W$ (\footnote[58]{we only find that $\hot_W(X'') \longrightarrow \hot_W(X')$ \texttt{[is an]} isomorphism.}). We probably have better: if $S'' \longrightarrow S'$ is a homotopism, then $X'' \longrightarrow X'$ is in $W$.}

\msp{page 49}
Let us prove (if we can) this nice complement. I point out that I had to make the hypothesis of $W$-presmoothness (or of $W$-preproperness) to be sure that the inverse image $X \times_S S'$, under $f' : S' \longrightarrow S$, really does coincide with $f^*$ on the $\HOTlc$ - otherwise I have no chance. % [sic] où -> ou, the conjunction "or" is meant
``We know'' (everything has to be looked at again!) that it's O.K. in the three cases $X$ $W$-smooth, $X$ $W$-proper, $X$ a $W$-fibration (but in what sense???). The check to be made (when working with the $\Wd$) is this:
\[ \begin{aligned}
&X \xrightarrow{i_X} \Phi_S(X) = Y \text{ is in } \Wd \text{ (without merit), and} \\
&\text{remains so under any base change.}
\end{aligned} \tag*{$(*)$} \]

But for $X, Y \in \mathrm{Ob}\,\mathrm{Cat}/S$ which are $W$-presmooth, $\Wd$-equivalence coincides with $\Wf$, which is preserved by base change, as is the hypothesis on $X$, $Y$. So we can rest easy.

We are reduced to proving that if $X$ is $W$-presmooth over $S$ and a $W$-pseudofibration, and if $S' \longrightarrow S$ is a homotopism, then $X' \longrightarrow X$ is in $W$. By the earlier argument using $X \longrightarrow \Phi_S(X)$, we reduce to the case where $X$ is split fibered.

\printedpage{36}
% head: A. Grothendieck : Dérivateurs, Ch. XII, Édition des Universités de Montpellier II et Paris VII
% foot: Version 7 décembre 2002

\msp{page 50}
\[ \begin{tikzcd}
X' \arrow[rr] \arrow[d] & & X \arrow[d] \\
S' \arrow[rr, "f"] \arrow[dr, "p'"'] & & S \arrow[dl, "p"] \\
 & e &
\end{tikzcd} \]

The arrow $X' = f^*(X) \longrightarrow X$ is interpreted in $\Hot_W$ as
\[ p'_!(f^*(\xi)) \longrightarrow p_!(\xi) \ , \qquad \text{where } \xi = \hot_{S,W}(X). \tag*{$(*)$} \]

I say that this arrow is \texttt{[an]} isomorphism, i.e. that for every $\eta \in \Hot_W$, the transposed arrow
\[ \Hom_W(p_!(\xi),\eta) \longrightarrow \Hom_W(p'_!(f^*(\xi)),\eta) \tag*{$(**)$} \]
is \texttt{[a]} bijection. Granting the adjunction formula, the arrow $(**)$ is interpreted as
\[ \Hom_W(\xi,\underbrace{p^*(\eta)}_{\eta_S}) \longrightarrow \Hom_W(f^*(\xi),\underbrace{p'^*(\eta)}_{f^*(\eta_S)}), \]
i.e.
\[ \Hom_W(\xi,\eta_S) \longrightarrow \Hom_W(f^*(\xi),f^*(\eta_S)), \]
that this is \texttt{[an]} isomorphism therefore follows from $f^*$ fully faithful on the $\HOTlc$. Now we have seen that $f^*$ is an equivalence (cor. 2, page 46).

\msp{page 51}
I am going to take up propositions 9 and 8 again, a bit badly put together.

\textbf{Theorem 3.} \textit{We assume that $W$ satisfies \textup{W(1,2,3)}. Let $f : X \longrightarrow S$. Consider the following conditions.}

\begin{itemize}[leftmargin=4em]
\item[\textit{a)}] \textit{$\forall\ S'' \longrightarrow S'$ over $S$, with $S''$, $S'$} contractible, \textit{$X'' \longrightarrow X'$ is in $W$.}

\item[\textit{a}$'$\textit{)}] \textit{$\forall\ S'$ over $S$, with $S'$ contractible, and $s' \in S'$, the inclusion} % AUDIT resolved: the label is a prime, not d prime; see corrections.csv, PDF leaf 46
\[ X_{s'} \hookrightarrow X' = X \times_S S' \]
\textit{is in $W$.}

\item[\textit{a}$_1$\textit{)}] \textit{As in a$'$), but assuming that $S'$ has a final object (instead of $S'$ contractible).} % AUDIT resolved: the reference is a prime; see corrections.csv, PDF leaf 46

\item[\textit{a}$_2$\textit{)}] \textit{As in a$'$), but assuming that $S'$ has an initial object.} % AUDIT resolved: the reference is a prime; see corrections.csv, PDF leaf 46

\printedpage{37}
% head: A. Grothendieck : Dérivateurs, Ch. XII, Édition des Universités de Montpellier II et Paris VII
% foot: Version 7 décembre 2002

\item[\textit{b}$_1$\textit{)}] \textit{$f$ is $W$-}preproper \textit{(i.e. a$'$) is satisfied when $s'$ is a final object of $S'$, with no other condition on $S'$ in $\mathrm{Cat}/S$), and moreover for every arrow $s_0 \longrightarrow s_1$ in $S$, the induced arrow $X/s_0 \longrightarrow X/s_1$ is in $W$ (i.e. a \textup{\texttt{[rather a]}}) is satisfied for $S/s_0 \longrightarrow S/s_1$ \textup{\texttt{[rather} $S/s_0 \longrightarrow S/s_1$\texttt{]}}).}

\item[\textit{b}$_2$\textit{)}] \textit{$f$ is $W$-}presmooth \textit{(i.e. a$'$) is satisfied when $s'$ is an initial object of $S'$, with no other condition on $S'$), and moreover the arrows $s_1\backslash X \longrightarrow s_0\backslash X$ \textup{\texttt{[rather} $s_1\backslash X \longrightarrow s_0\backslash X$\texttt{]}} (for $s_0 \longrightarrow s_1$ in $S$) are in $W$, (i.e. a \textup{\texttt{[rather a]}}) satisfied for $s_1\backslash S \longrightarrow s_0\backslash S$ \textup{\texttt{[rather} $s_1\backslash S \longrightarrow s_0\backslash S$\texttt{]}}).}

\item[\textit{c}$_1$\textit{)}] \textit{$f$ is $W$-preproper and a $W$-pseudofibration, i.e. condition a$'$) is satisfied in the following two cases: 1$^\circ$) $s'$ is a final object of $S'$, 2$^\circ$) $S' = \boldsymbol{\Delta}^1$.}

\item[\textit{c}$_2$\textit{)}] \textit{$f$ is $W$-presmooth and a $W$-pseudofibration, i.e. a$'$) in the two cases: 1$^\circ$) $s'$ is an initial object of $S'$, 2$^\circ$) $S' = \boldsymbol{\Delta}^1$.}
\end{itemize}
\DERmark{ENV0683}\DERmark{ENV0684}\DERmark{ENV0685}\DERmark{ENV0686}\DERmark{ENV0687}


\textit{All these 8 conditions are equivalent. They imply the following:}

\begin{itemize}[leftmargin=4em]
\item[\textit{d)}] \textit{Let $S'' \xrightarrow{u} S'$ be a homotopism in $\mathrm{Cat}/S$, and consider $u_X : X'' \longrightarrow X'$. Then $\hot_W(u_X)$ is \textup{\texttt{[an]}} isomorphism. (So if $W$ is strongly saturated, $u_X \in W$. So under this condition on $W$, d) is equivalent to the preceding conditions.)}
\end{itemize}

\msp{page 52}
\textbf{Definition 4.} We say that $f : X \longrightarrow Y$ in Cat is a \textit{$W$-fibration}, if it satisfies the equivalent conditions a) to c$_2$) of theorem 3.
\DERmark{ENV0688}


\textbf{Corollary 1.} \textit{If $f$ is $W$-proper or $W$-smooth, for $f$ to be a $W$-fibration, it is necessary and sufficient that $f$ be a $W$-pseudofibration, i.e. that for $S' = \boldsymbol{\Delta}^1 \longrightarrow S$, the inclusions of the two fibers}
\[ X'_0 \xrightarrow{i_0} X' \xleftarrow{i_1} X'_1 \]
\textit{of $X' = X \times_S S'$ be in $W$.}

\textbf{Corollary 2.} \textit{Suppose $f$ \textup{Cat}-fibrant. Then $f$ is a $W$-fibration if and only if for every arrow $u : s \longrightarrow s'$ in $S$, $u^* : X_{s'} \longrightarrow X_s$ is in $W$. Dual statement if $f$ is \textup{Cat}-cofibrant.} % [sic] pout -> pour, misprint

\textbf{Corollary 3.} \textit{Suppose $f$ $W$-preproper and $W$-presmooth (for example $f$ $W$-perfect). Then $f$ is a $W$-fibration.}

\textbf{Proposition 10.} \textit{Let $f : X \longrightarrow Y$ be an arrow in $\underline{\mathrm{Fib}}_WS$ (\footnote[59]{$X$ and $Y$ both preproper or both presmooth suffices.}) (full subcategory of $\mathrm{Cat}/S$ formed by the $W$-fibered objects over $S$). If $\forall\ s \in S$, $f_s : X_s \longrightarrow Y_s$ is in $W$, then $f \in W$.}

Let us consider
\printedpage{38}
% head: A. Grothendieck : Dérivateurs, Ch. XII, Édition des Universités de Montpellier II et Paris VII
% foot: Version 7 décembre 2002
\[ \begin{tikzcd}
X \arrow[rr, "f"] \arrow[dr, "i_X"'] & & Y \arrow[dr, "i_Y"] & \\
 & \Phi_S(X) \arrow[rr, "\Phi_S(f)"'] & & \Phi_S(Y),
\end{tikzcd} \]
where $i_X, i_Y \in W$, so $\Phi_S(f) \in W \Rightarrow f \in W$. Now the fact that $X$ \texttt{[is]} $W$-fibered over $S$ implies that $i_X \in \Wf$ (it even suffices that $X$ be presmooth \texttt{[rather presmooth]} over $S$), likewise for $Y$, $i_Y$. So the hypothesis $f \in \Wf$ implies $\Phi_S(f) \in \Wf$, so $\Phi_S(f) \in W$ by the \texttt{[the]} corollary to proposition 5 (page 30).
\DERmark{ENV0689}\DERmark{ENV0690}



\endgroup
% END UNIT Der-ChXIIa-012/Der-ChXIIa-012_p036-040.tex
% BEGIN UNIT Der-ChXIIa-012/Der-ChXIIa-012_p041-044.tex
\begingroup\makeatletter%
\NoAutoSpacing%
\newcommand{\printedpage}[1]{\marginpar{\footnotesize p.~#1}}%
\newcommand{\msp}[1]{\marginpar{\footnotesize [#1]}}%
\newcommand{\dashedbox}[1]{\tikz[baseline=(n.base)]\node[draw,dashed,inner sep=2pt](n){#1};}%
\newcommand{\Fib}{\mathrm{Fib}}%
\newcommand{\Parf}{\mathrm{Parf}}%
\newcommand{\HOT}{\mathrm{HOT}}%
\newcommand{\Hot}{\mathrm{Hot}}%
\newcommand{\hot}{\mathrm{hot}}%
\newcommand{\Hom}{\mathrm{Hom}}%
\makeatother\let\printedpage\derpp\let\msp\dermsp%
\let\mspin\DERmspin%


\printedpage{39}
% head: A. Grothendieck : Dérivateurs, Ch. XII, Édition des Universités de Montpellier II et Paris VII

\msp{page 53}

\section*{IV\quad On axiom W(4) and its variants.}

(\footnote[60]{Cf. also part VI (HOT$_W$), p. 186-196}). I have spent days trying to prove that all these variants are equivalent. I am going to give a list of those, among these variants, which in retrospect still seem to me to be of some use. (With notations different from those on pages 13-14.)

\begin{itemize}[leftmargin=5em,labelsep=0.5em]
\item[\fbox{W(4a)}] Every $W$-fibration (cf. th. 3, p. 51, def. 4, p. 52) is ultrastrong, which amounts to this: for
\[
\begin{tikzcd}
X & X' \arrow[l,"q"'] \arrow[d,"f'"] \\
Y \arrow[from=u,"f"'] & Y' \arrow[l,"p"']
\end{tikzcd}
\]
\DERmark{ENV0691}

cartesian in Cat, if $f \in \underline{\Fib}_W$, $p \in W$, then $q \in W$.

\item[W(4a$'$)] As in a), assuming $f \in \Parf_\omega$, i.e. $f$ $W_\omega$-perfect (hence a $W_\omega$-fibration), where $W_\omega \subseteq W$ is the smallest subset $\Sigma$ of Fl(Cat) satisfying W(1, 2 bis, 3). % [sic] petits -> petite, accord avec « partie »

\item[\fbox{W(4b)}] Consider a commutative diagram
$\begin{tikzcd}[ampersand replacement=\&, column sep=small, row sep=small, baseline=(X.base)]
|[alias=X]| X \arrow[rr,"f"] \arrow[dr,"p"'] \& \& Y \arrow[dl,"q"] \\
\& S \&
\end{tikzcd}$
in Cat, with $p, q \in \Fib_W$, $f \in W$.
Then $\forall s \in S$, $f_s : X_s \longrightarrow Y_s$ \texttt{[is in} $W$\texttt{]}.

\item[W(4b$'$)] As in b), assuming \texttt{[in addition]} $X$ and $Y$ split fibered over $S$, $f$ compatible with \texttt{[the]} splittings (hence a cartesian functor).

\item[W(4b$''$)] As in b), assuming $p$, $q$ $W$\textit{-perfect.}

\item[\fbox{W(4c)}] Let $f : X \longrightarrow Y$, with $f \in \Fib_W \cap W$, then the fibers of $f$ are $W$-aspheric.

\item[\dashedbox{W(4c$'$)}] As in c), assuming $f \in \Parf_\omega \cap W$. (This is the old W(4b$''$), stripped of the $W_\infty$-parasite.)
\end{itemize}

\msp{page 54}

I have circled the notions \texttt{[boxed here]} that have been really useful. In practice, one can bet that it is the strongest notion, W(4a), which is the only important one. Here are the tautological implications (with the exception of W(4a$'$) $\Longrightarrow$ W(4b), which is proved in the proof of proposition 4 (pages 6-10) (\footnote[61]{This proof uses a stronger notion of $W$-fibration than that of theorem 3, p. 51.})).

\printedpage{40}
% head: A. Grothendieck : Dérivateurs, Ch. XII, Édition des Universités de Montpellier II et Paris VII

\[
\begin{tikzcd}[column sep=small]
& \mathrm{a} \arrow[d,Rightarrow] & \\
& \mathrm{a}' \arrow[dd,Rightarrow,"{\text{\rule[0.6ex]{2cm}{0.4pt}\ \begin{tabular}[t]{@{}l@{}}doubtful implication,\\ one must weaken\\ b) a little, by\\ strengthening the\\ notion of\\ $W$-fibration\\ used\\ to state b)\end{tabular}}}"] & \\
& & \\
& \mathrm{b} \arrow[dl,Rightarrow] \arrow[dr,Rightarrow] \arrow[dd,Rightarrow] & \\
\mathrm{b}' & & \mathrm{b}'' \\
& \mathrm{c} \arrow[d,Rightarrow] & \\
& \mathrm{c}' &
\end{tikzcd}
\]

I have mainly sweated over showing that c$'$) $\Longrightarrow$ c), or c) $\Longrightarrow$ b), without getting there. The effort has allowed me to get rid of various false ideas, and to replace them with right ones, cf. below. There is all the same a positive result that stays afloat:

\textbf{Proposition 1.} \textit{Condition b) (and probably even b$'$) and b$''$)} (\footnote[62]{and assuming in the case of b$''$, that $W$ satisfies the condition on direct limits W(6 bis), cf. prop. 5 of the ``corollary'' of the, page 14.})\textit{, but I am not going to prove it here \dots) probably implies a) for $\bar{W}$, the strong saturation of $W$. In any case, if $W$ is strongly saturated, conditions a), a$'$), b) (and surely also b$'$) and b$''$)) are equivalent.}

\textsc{Proof} of b) $\Longrightarrow$ a), when $W$ is strongly saturated. In fact, assuming b) (and without even assuming $W$ strongly saturated), we are going to prove that if we have
\[
\begin{tikzcd}
X & X' \arrow[l,"q"'] \arrow[d,"f'"] \\
Y \arrow[from=u,"f"'] & Y' \arrow[l,"p"']
\end{tikzcd}
\]
cartesian in Cat, with
\msp{page 55}
$f$ a $W$-fibration, and $p \in W$, then $\underline{q \in \bar{W}}$, i.e. $\hot_W(q)$ \texttt{[is an]} isomorphism. For this, we interpret the arrow $q$ in $\Hot_W$ in the following way. Let
\[
\xi = \hot_{Y,W}(X/Y), \qquad \xi' = p^*(\xi),
\]
\printedpage{41}
% head: A. Grothendieck : Dérivateurs, Ch. XII, Édition des Universités de Montpellier II et Paris VII
which says that $q$ is identified with the canonical arrow
\[
H_\bullet^{\HOT_W}(Y',\xi') \longrightarrow H_\bullet^{\HOT_W}(Y,\xi) \qquad \text{in } \Hot_W = \HOT_W(e), \tag*{$(*)$}
\]
where we note moreover that $\xi \in \HOT_W(Y)$ is a coefficient \textit{locally constant} on $Y$.

Let us first establish this interpretation of $q$. Note that since $X \in \underline{\Fib}_W Y$, the inverse image $\xi' = p^*(\xi)$ is none other than $\hot_{Y',W}(X' = X \times_Y Y')$. On the other hand, if $g : Y \longrightarrow e$, $g' : Y' \longrightarrow e$ are the structural arrows, the two sides of $(*)$ are interpreted resp. as $g'_!(\xi')$ and $g_!(\xi)$ (the $g'_!$ and $g_!$ in the sense of the prederivator $\HOT_W$). Now the calculation of $g'_!$, $g_!$ gives us that the two sides are identified resp. with
\[
\hot_W(X'), \qquad \hot_W(X),
\]
and we admit that the arrow
\[
\hot_W(X') \longrightarrow \hot_W(X)
\]
is none other than $\hot_W(q)$. Besides, since $X$ is $W$-fibered over $Y$, practically by definition, $\xi$ is in $\HOT^{\mathrm{lc}}_W(Y)$. So it will suffice for us to prove the
\DERmark{ENV0692}


\msp{page 56}

\textbf{Corollary.} \textit{Suppose that $W$ satisfies W(4b), and let $p : S' \longrightarrow S$ be in $W$, $\xi \in \HOT^{\mathrm{lc}}_W(S)$, $\xi' = p^*(\xi)$. Then}
\[
H_\bullet(S',\xi') \xrightarrow{\ \sim\ } H_\bullet(S,\xi)
\]
\textit{is an isomorphism in} $\Hot_W$.

\textsc{Proof}. It suffices to prove that for every $\eta \in \mathrm{Ob}\,\Hot_W$ we have a bijection for the $\Hom(-,\eta)$. By the adjunction formula on $S$, $S'$, this map between the Hom in $\Hot_W$ is identified \texttt{[with]}
\[
\Hom_{\HOT_W(S)}(\xi,\eta_S) \longrightarrow \Hom_{\HOT_W(S')}(\xi',\eta_{S'}),
\]
where $\eta_S = g^*(\eta)$, $\eta_{S'} = g'^*(\eta)$ \texttt{[rather} $\eta_S = g^*(\eta)$\texttt{,} $\eta_{S'} = g'^*(\eta)$\texttt{, where} $g : S \longrightarrow e$\texttt{,} $g' : S' \longrightarrow e$ \texttt{are the structural arrows]}, hence $\eta_{S'} = p^*(\eta_S)$, and the arrow $(*)$ is the arrow on the Hom induced by the functor
\DERmark{ENV0693}

\[
p^*_{\mathrm{lc}} : \HOT^{\mathrm{lc}}_W(S) \longrightarrow \HOT^{\mathrm{lc}}_W(S').
\]
Now assuming axiom W(4b) (which we have not used so far), the said functor is an \textit{equivalence of categories} if $p \in W$ (cf. VI, p. 204).

\hfill q.e.d.

For fear of forgetting them, I am now going to make a list of the little questions I have elucidated along the way while scribbling.

\begin{itemize}[leftmargin=5em,labelsep=0.5em]
\item[1°)] If $X$, $Y$ are $W$-smooth over $S$, and $f : X \longrightarrow Y$ an $S$-morphism such that $f_s : X_s \longrightarrow Y_s$ is $W$-smooth for every $s \in \mathrm{Ob}\,S$, it does \textit{not} follow that $f$ is $W$-smooth. But this is nevertheless true when $X$ and $Y$ are split fibered over $S$ (which
\printedpage{42}
% head: A. Grothendieck : Dérivateurs, Ch. XII, Édition des Universités de Montpellier II et Paris VII
\msp{page 57}
already implies their $W$-smoothness over $S$), and $f$ compatible with the splittings and a fortiori cartesian. Dual statement for split cofibered categories and the strictly cartesian functors between them that are $W$-proper fiber by fiber: they are $W$-proper.

\item[2°)] I was interested, in a case of split fibered categories $X$, $Y$ over $S$ and a strictly cartesian functor $f : X \longrightarrow Y$, assumed $W$-perfect fiber by fiber, in concluding that $f$ should also be $W$-perfect. (It was a case where in addition one knows that $X$, $Y$ are also $W$-fibered over $S$, i.e. the base change functors $X_{s'} \longrightarrow X_s$, $Y_{s'} \longrightarrow Y_s$ are in $W$. But that was not much use \dots) The question that remained, in 1°), was that of the $W$\textit{-properness} of $f$. But this is hopeless (even for $Y = S$), by virtue of the brilliant % [sic] fonteurs -> foncteurs, lettre omise

\textbf{Proposition 2.} \textit{Suppose that} $S = 0$
$\begin{tikzcd}[ampersand replacement=\&, row sep=tiny, column sep=small, baseline=(Z.base)]
\& a \\
|[alias=Z]| \phantom{0} \arrow[ur] \arrow[dr] \& \\
\& b
\end{tikzcd}$
\textit{, and let $Z$ be a category over $S$. Then $\Phi(Z)$ is proper over $S$ if and only if $Z = \emptyset$.}

Let us then take in 2°) above $X = \Phi(Z)$, $Y = S$, $f : \Phi(Z) \longrightarrow S$ \texttt{[rather} $f : \Phi(Z) \longrightarrow S$\texttt{]} the structural morphism (so the $f_s$ are $W$-perfect), then $f$ is proper only if $X = \emptyset$. Yet, if we take $Z$ $W$-fibered over $S$, $\Phi(Z)$ is so too, and we have a counterexample to 2°).
\DERmark{ENV0694}

\end{itemize}

But we have all the same convinced ourselves of the following result:

\textbf{Proposition 3.} \textit{Let $X$, $Y$ be split fibered categories over $S$, $f : X \longrightarrow Y$ a strictly cartesian $S$-functor, such that}
\msp{page 58}
\textit{for every $s \in S$, $f_s : X_s \longrightarrow Y_s$ is} $W$-fibrant. \textit{Then,} if $W$ satisfies W(4a), $f$ \textit{is $W$-fibrant.}

(I was hoping to get that $f$ is $W$-fibrant ``for free'', to prove that W(4c) implies W(4b). But I had to assume W(4a) \dots)


\endgroup
% END UNIT Der-ChXIIa-012/Der-ChXIIa-012_p041-044.tex
\clearpage\setcounter{footnote}{0}\phantomsection\addcontentsline{toc}{section}{Chapter XII, 5}\markboth{Chapter XII, 5}{Chapter XII, 5}
% BEGIN UNIT Der-ChXIIb-012/Der-ChXIIb-012_p001-005.tex
\begingroup\makeatletter%
\newcommand{\printedpage}[1]{\marginpar{\footnotesize p.~#1}}%
\newcommand{\msp}[1]{\marginpar{\footnotesize [#1]}}%
\newcommand{\uHom}{\underline{\mathrm{Hom}}}%
\makeatother\let\printedpage\derpp\let\msp\dermsp%
\let\mspin\DERmspin%


\printedpage{43}
% head: A. Grothendieck : Dérivateurs, Ch. XII, Édition des Universités de Montpellier II et Paris VII ; foot: Version 7 décembre 2002

\msp{page 59}

\section*{V\quad Amalgamated sums and $W$-cocartesian squares in $\mathrm{Cat}$}

\subsection*{1\quad Generalities on $U \overset{i}{\hookrightarrow} X \overset{j}{\hookleftarrow} T$, and amalgamated sums.}

\subsubsection*{1.1\quad Gluing data between an open and a closed.}

Let $X$ be in Cat, $U$ an open in $X$, $T$ the closed complement. (\textbf{NB} The opens do indeed correspond to the opens of $\mathrm{Ob}\,X$, topologized by its natural preorder relation $x \leq y \iff \mathrm{Hom}_X(x,y) \neq \emptyset$, likewise for the closed ones. The topologies associated to preorder relations are those where every intersection of opens is still open, i.e. every union of closed sets is still closed. Thus the opens of such a topology are the closed sets of a dual topology, corresponding to the opposite preorder relation. The preorder relation associated to the topology is
\[
x \leq y \qquad\qquad \iff \qquad\qquad y \in \overline{\{x\}},
\]
\DERmark{ENV0714}

the topology associated to the preorder on $I$ is the one for which the opens are the unions of sets of the form $I/x$ (which form a basis of opens), or again the \textit{``sieves''} ($x \in U$, $y \leq x \Longrightarrow y \in U$), the closed sets are the unions of sets $x\backslash I$, that is to say the cosieves. For every subset $A$ of $I$, there is a smallest open containing $A$, which is also its smallest neighborhood, it is the union of the $I/x$ for $x \in A$. And dually there is a largest closed set contained in $A$, whose complement is the smallest open containing $B = I\backslash A$, hence $\bigcup_{x\in B} I/x$. The closure of $A$ is the union of the closures $x\backslash I$ of its points.)

\msp{page 60}

The situation $(X,U,T)$ is reconstructed, up to unique isomorphism, from the knowledge of

\begin{enumerate}[label={}]
\item[$1^\circ$)] The two categories $U$, $T$.
\item[$2^\circ$)] The bifunctor $\displaystyle\left\{\begin{aligned}
&H : U^\circ \times T \longrightarrow \mathrm{Ens}\\[1ex]
&H(u,t) = \mathrm{Hom}_X(u,t).
\end{aligned}\right.$
\end{enumerate}

Moreover, these data $1^\circ$) $2^\circ$) are subject to no restriction. We have, more precisely % [sic] en sont -> ne sont, negation required by the sense

\textbf{Proposition 1.} \textit{An isomorphism} \texttt{[rather an equivalence of categories]} \textit{between the category $\mathcal{T}$ whose objects are the triples $(X,U,T)$, with $X \in \mathrm{Ob}\,\mathrm{Cat}$, $U \in \mathrm{Ouv}(X)$, $T \in \mathrm{Ferm}(X)$, $U$ and $T$ complementary to one another, the morphisms from $(X,U,T)$ to $(X',U',T')$ being the functors $f : X \longrightarrow X'$ which induce $U \longrightarrow U'$, $T \longrightarrow T'$ (whence
\printedpage{44}
% head: A. Grothendieck : Dérivateurs, Ch. XII, Édition des Universités de Montpellier II et Paris VII ; foot: Version 7 décembre 2002
$U = f^{-1}(U')$, $T = f^{-1}(T')$), with the category $\mathcal{T}'$ of triples $(U,T,H)$, of two categories $U$, $T$ and of a bifunctor}
\[
H : U^\circ \times T \longrightarrow \mathrm{Ens}
\]
\hspace*{2em}\textit{(or again a presheaf $U^\circ \longrightarrow T^\vee$ on $U$, hence an object of} $\uHom(U^\circ,T^\vee)$\textit{)),} % [sic] préfaiscau -> préfaisceau, misspelling

\textit{a morphism from $(U,T,H)$ to $(U',T',H')$ being a triple}
\[
(f_U, f_T, u) \qquad\qquad
\left\{\begin{array}{l}
\left.\begin{array}{llll}
f_U : & U & \longrightarrow & U' \\[1ex]
f_T : & T & \longrightarrow & T'
\end{array}\right\}\text{functors}\\[2ex]
u : H \longrightarrow H' \circ (f_U^\circ \times f_T),\\
\text{i.e. } H(u,t) \longrightarrow H'(f_U(u), f_T(t))\\
\text{a functorial homomorphism}\\
\text{(arrow in } (U \times T^\circ)^\wedge \simeq \uHom(U^\circ,T^\vee)\text{)),}
\end{array}\right.
\]
\textit{composition of arrows obvious.}

If $U \hookrightarrow X \hookleftarrow T$ is a triple, and
\[
f_{U'} : U' \longrightarrow U,\ f_{T'} : T' \longrightarrow T
\]
arrows in Cat, there is a privileged way to ``glue'' $U'$ and $T'$, in such a way that the object $(X',U',T')$ maps into $(X,U,T)$ inducing

\msp{page 61}

$f_{U'}$, $f_{T'}$:
\[
\begin{tikzcd}[column sep=large, row sep=large]
U' \arrow[r, hook, "i'"] \arrow[d, "f_{U'}"'] & X' \arrow[d, "f"] & T' \arrow[l, hook', "j'"'] \arrow[d, "f_{T'}"] \\
U \arrow[r, hook, "i"] & X & T, \arrow[l, hook', "j"']
\end{tikzcd}
\]
it is obtained with the help of the bifunctors
\[
\left\{\begin{aligned}
&H'(u',t') \overset{\mathrm{def}}{=} H(f_{U'}(u'), f_{T'}(t'))\\
&\text{i.e. } H' = H \circ (f_{U'}^\circ \times f_{T'}).
\end{aligned}\right.
\]
\DERmark{ENV0715}

This triple, i.e. this object of $\mathcal{T}$, represents the following functor on $\mathcal{T}$ :
\[
\begin{aligned}(X'',U'',T'')\longmapsto{}&\mathrm{Hom}_{\mathcal{T}}((X'',U'',T''),(X,U,T))\\ &{}\times_{\substack{\mathrm{Hom}_{\mathrm{Cat}}(U'',U)\\{}\times \mathrm{Hom}_{\mathrm{Cat}}(T'',T)}}\bigl(\mathrm{Hom}_{\mathrm{Cat}}(U'',U') \times \mathrm{Hom}_{\mathrm{Cat}}(T'',T')\bigr)\end{aligned},
\]
\DERmark{ENV0716}

i.e. for every object $(X'',U'',T'')$ of $\mathcal{T}$ we have a canonical bijection, functorial in this object
\[
\begin{aligned}\mathrm{Hom}_{\mathcal{T}}((X'',U'',T''),(X',U',T')) &\simeq \mathrm{Hom}_{\mathcal{T}}((X'',U'',T''),(X,U,T))\\ &\quad{}\times_{\substack{\mathrm{Hom}_{\mathrm{Cat}}(U'',U)\\{}\times \mathrm{Hom}_{\mathrm{Cat}}(T'',T)}}\bigl(\mathrm{Hom}_{\mathrm{Cat}}(U'',U') \times \mathrm{Hom}_{\mathrm{Cat}}(T'',T')\bigr)\end{aligned}. \tag*{$(*)$}
\]
\DERmark{ENV0717}


\printedpage{45}
% head: A. Grothendieck : Dérivateurs, Ch. XII, Édition des Universités de Montpellier II et Paris VII ; foot: Version 7 décembre 2002

No doubt we even have
\[
\begin{aligned}\uHom_{\mathcal{T}}((X'',U'',T''),(X',U',T')) &\simeq \uHom_{\mathcal{T}}((X'',U'',T''),(X,U,T))\\ &\quad{}\times_{\substack{\uHom_{\mathrm{Cat}}(U'',U)\\{}\times \uHom_{\mathrm{Cat}}(T'',T)}}\bigl(\uHom_{\mathrm{Cat}}(U'',U') \times \uHom_{\mathrm{Cat}}(T'',T')\bigr)\end{aligned}, \tag*{$(**)$}
\]
\DERmark{ENV0718}

just as functorial, from which $(*)$ is deduced by passing to the sets of objects. We say that the triple $(X',U',T')$ is the triple \textit{induced} by the triple $(X,U,T)$, via the functors $f_{U'}$, $f_{T'}$.

\textbf{NB} If $f_{U'}$, $f_{T'}$ are fully faithful, then so is $f : X' \longrightarrow X$. We shall use this operation of privileged gluing mainly when $U' = U$, $f_{U'} = \mathrm{id}_U$, or conversely $T' = T$, $f_{T'} = \mathrm{id}_T$.

\subsubsection*{1.2\quad Relation with the gluing of topoi.}

The situation of a triple
\[
(\mathcal{X},\mathcal{U},\mathcal{T})
\]
of a topos $\mathcal{X}$, of an open $\mathcal{U}$ of $\mathcal{X}$, and of a

\msp{page 62}

complementary closed subtopos, is entirely similar (\textbf{NB} If $\mathcal{X} = \mathrm{Top}(X)$, with $X$ in Cat, then $\mathrm{Ouv}(\mathcal{X}) \simeq \mathrm{Ouv}(X)$, since the opens correspond to the subsheaves of the terminal sheaf on $\mathcal{X}$, i.e. of the terminal presheaf on $X$. This bijection is compatible with passage to complements.) Topos theory says that the pairs $(\mathcal{X},\mathcal{U},\mathcal{T})$ \texttt{[rather triples]} are reconstructed, up to an equivalence defined up to unique isomorphism, with the help of a triple
\[
\left\{\begin{aligned}
&\mathcal{U},\ \mathcal{T} \text{ two topoi}\\[1ex]
&\varphi : \mathrm{Fais}\,(\mathcal{U}) \longrightarrow \mathcal{T} \text{ a left exact functor (``accessible''),}
\end{aligned}\right.
\]
this functor $\varphi$ being obtained, when we start from
\[
\mathcal{U} \xrightarrow{\;i\;} \mathcal{X} \xleftarrow{\;j\;} \mathcal{T},
\]
as
\[
\varphi = j^*i_*.
\]
Moreover, the triple $(\mathcal{U},\mathcal{T},\varphi)$ can be chosen arbitrarily, with the sole restriction (practically harmless, since always satisfied for all the natural functors we encounter) that $\varphi$ be (besides being left exact) ``accessible''.

In the present case $\mathcal{X} = \mathrm{Top}(X)$, $\mathcal{U} = \mathrm{Top}(U)$, $\mathcal{T} = \mathrm{Top}(T)$, the functor $\varphi = j^*i_*$ is not only left exact, but it even commutes with arbitrary $\varprojlim$, since this is so (always) for $i_*$, and also (because the topos-theoretic $j$ is induced by the category-theoretic $j : T \longrightarrow X$) for $j^*$. In fact, we
\printedpage{46}
% head: A. Grothendieck : Dérivateurs, Ch. XII, Édition des Universités de Montpellier II et Paris VII ; foot: Version 7 décembre 2002
\msp{page 63}
find this :

\textbf{Proposition 2 :} \textit{Let us consider a situation of gluing two topoi $\mathcal{U}$, $\mathcal{T}$ by $\varphi : \mathrm{Fais}\,\mathcal{U} \longrightarrow \mathrm{Fais}\,\mathcal{T}$, a left exact functor. Let us suppose that $\mathcal{U} = \mathrm{Top}(U)$, $\mathcal{T} = \mathrm{Top}(T)$, with $U$ and $T$ in $\mathrm{Cat}$. For the glued topos $\mathcal{X}$ to be equivalent to a topos $\mathrm{Top}(X)$, with $X$ in $\mathrm{Cat}$, it is necessary and sufficient that the gluing functor $\varphi$ commute with arbitrary $\varprojlim$, or what amounts to the same thing, that it have a left adjoint :}
\[
\psi : \underbrace{\mathrm{Fais}\,\mathcal{T}}_{= T^\wedge} \longrightarrow \underbrace{\mathrm{Fais}\,\mathcal{U}}_{= U^\wedge}.
\]

We then expect that the gluing datum $\varphi$, or what amounts to the same thing the functor $\psi$ having a right adjoint, i.e. commuting with arbitrary $\varinjlim$, amounts to that of a gluing datum $H$
\[
H : U^\circ \times T \longrightarrow \mathrm{Ens}
\]
for $U$ and $T$. But we know well that the functors $\psi : T^\wedge \longrightarrow \mathcal{C}$ into a category $\mathcal{C}$ with arbitrary $\varinjlim$, which commute with the said $\varinjlim$, form a category equivalent (by the functor $\psi \longmapsto \psi|T(\subseteq T^\wedge)$) to that of functors
\[
\psi_0 : T \longrightarrow \mathcal{C} \quad\ (= \underbrace{U^\wedge}_{\uHom(U^\circ,\mathrm{Ens})} \text{ in this case}),
\]
\msp{page 64}
and giving such a functor $\psi_0$ does indeed amount to giving an $H$. Thus

\textbf{Corollary.} \textit{The category of topos-theoretic gluing data $\varphi : U^\wedge \longrightarrow T^\wedge$ between} $\mathrm{Top}(U)$, $\mathrm{Top}(T)$\textit{, such that the glued topos is a ``}paradisiacal\textit{ topos'', i.e. equivalent to a topos of the form $X^\wedge$ (with $X$ in} Cat\textit{), is equivalent to that of gluing data $H : U^\circ \times T \longrightarrow$} Ens \textit{between $U$ and $T$. The topos-theoretic datum $\varphi$ associated to $H$ is the functor right adjoint to $T^\wedge \longrightarrow U^\wedge$ deduced from $T \xrightarrow{H} U^\wedge$ by extension commuting with $\varinjlim$, or again it is given by}
\[
\varphi(\underbrace{F}_{\in U^\wedge})(t) \simeq \mathrm{Hom}_{U^\wedge}(H(t),F),
\]
\textit{where on the right-hand side $H$ is interpreted as a functor $T \longrightarrow U^\wedge$. The topos glued by $\varphi$ is canonically equivalent to} $\mathrm{Top}(X)$\textit{, where $X$ is the object of} Cat \textit{obtained by gluing with the help of $H$.}

From the point of view of sheaf theory, this gives the interpretation of $H$, as expressing the functor
\[
\varphi = j^*i_* : U^\wedge \longrightarrow T^\wedge
\]
by the preceding formula.

\printedpage{47}
% head: A. Grothendieck : Dérivateurs, Ch. XII, Édition des Universités de Montpellier II et Paris VII ; foot: Version 7 décembre 2002

\msp{page 65}

\subsubsection*{1.3\quad Arrows in $\mathrm{Cat}$, with source or target a glued object $(X,U,T)$.}

When $(X,U,T)$ is given by a gluing datum
\[
U,\ T,\ H : U^\circ \times T \longrightarrow \mathrm{Ens},
\]
the arrows
\[
f : X' \longrightarrow X
\]
in Cat are immediately described as follows : giving $f$ first defines
\[
U' = f^{-1}(U), \qquad\qquad T' = f^{-1}(T),
\]
a pair of an open of $X'$ and the complementary closed set, so that this datum is equivalent to that of

\begin{enumerate}[label={}]
\item[$1^\circ$)] a pair $(U',T')$ in $X'$
\item[$2^\circ$)] a morphism (in $\mathcal{T}$) from $(X',U',T')$ to $(X,U,T)$, which by 1.1 is equivalent to a triple
\[
\left\{\begin{array}{rccl}
f_{U'} : & U' & \longrightarrow & U\\[1ex]
f_{T'} : & T' & \longrightarrow & T\\[1ex]
u : & H_{U',T'} & \longrightarrow & H \circ (f_{U'}^\circ \times f_{T'}),
\end{array}\right.
\]
where we have put
\[
\left\{\begin{aligned}
&H_{U',T'} : U^{\prime\circ} \times T' \longrightarrow \mathrm{Ens}\\[1ex]
&H_{U',T'}(u',t') = \mathrm{Hom}_{X'}(u',t'),
\end{aligned}\right.
\]
i.e. $H_{U',T'}$ is the gluing datum which corresponds to the triple $(X',U',T')$ in $\mathcal{T}$. % [sic] conrrespond -> correspond, misspelling
\end{enumerate}

\msp{page 66}

We shall be more interested in the reverse situation, i.e. in the question of determining the functors
\[
\underbrace{X}_{= X_{U,T,H}} \quad \longrightarrow X'
\]
from $X$, obtained from $(U,T,H)$, to an arbitrary object of Cat. We find tautologically

\textbf{Proposition 3.} $\mathrm{Hom}_{\mathrm{Cat}}(X,X')$ \textit{is in canonical bijective correspondence, functorial in $X'$ in} Cat\textit{, with the set of triples}


\endgroup
% END UNIT Der-ChXIIb-012/Der-ChXIIb-012_p001-005.tex
% BEGIN UNIT Der-ChXIIb-012/Der-ChXIIb-012_p006-010.tex
\begingroup\makeatletter%
\newcommand{\printedpage}[1]{\marginpar{\footnotesize p.~#1}}%
\newcommand{\msp}[1]{\marginpar{\footnotesize [#1]}}%
\newcommand{\Hom}{\mathrm{Hom}}%
\newcommand{\uHom}{\underline{\mathrm{Hom}}}%
\newcommand{\Cat}{\mathrm{Cat}}%
\newcommand{\Ens}{\mathrm{Ens}}%
\newcommand{\id}{\mathrm{id}}%
\newcommand{\Ob}{\mathrm{Ob}\,}%
\makeatother\let\printedpage\derpp\let\msp\dermsp%
\let\mspin\DERmspin%


\printedpage{48}
% head: A. Grothendieck : Dérivateurs, Ch. XII, Édition des Universités de Montpellier II et Paris VII
% foot: Version 7 décembre 2002

\[
\left\{
\begin{aligned}
&\left.\begin{aligned} f_U &: U \longrightarrow X' \\ f_T &: T \longrightarrow X' \end{aligned}\right\}\text{arrows in Cat}\\
&u : H \longrightarrow H_{X'}(f_U^{\circ}\times f_T),\ \text{arrow in } \uHom(U^{\circ}\times T,\Ens),
\end{aligned}
\right.
\tag*{\texttt{[}$(*)$\texttt{]}}
\]
\textit{where for every $X'$ in} Cat\textit{, we have put}
\[
H_{X'} = ((x,y)\longmapsto \Hom_{X'}(x,y)) : X'^{\circ}\times X' \longrightarrow \Ens.
\]

Surely we will even have an isomorphism of categories
\[
\uHom(X,X') \simeq \text{category of triples } (f_U,f_T,u) \text{ as in } (*).
\]

\subsection*{1.4\quad Construction of amalgamated sums from an open immersion or closed immersion.} % [sic] où fermée -> ou fermée, conjonction « ou » et non adverbe « où »

For a construction of a sort of amalgamated sum, in terms of a triple $(X,U,T)$ $+\ \ldots$, \msp{page 67} which would be symmetric in $U$, $T$, I feel like looking first at the ``mixed'' situation, of a diagram in solid lines
\DERmark{ENV0719}

\[
\begin{tikzcd}
U \arrow[r, hook, "i"] \arrow[d, "f_U"'] & X \arrow[d, dotted, "f"] & T \arrow[l, hook', "j"'] \arrow[d, "f_T"] \\
U' \arrow[r, hook, dotted, "i'"] & X' & T', \arrow[l, hook', dotted, "j'"']
\end{tikzcd}
\tag*{$(*)$}
\]
which we would like to complete ``canonically'', if it can be done, to a morphism of triples (i.e. to an arrow in $\mathcal{T}$). The natural idea is that the $X'$ which fits into it would be the direct limit, in Cat, of the initial diagram. If $T'=T$, $f_T=\id_T$, this therefore means that $X'$ identifies with the amalgamated sum $X\amalg_U U'$, and conversely if $U'=U$, $f_U=\id_U$.

Let us denote by $X'$ the amalgamated sum \texttt{[rather the direct limit]}, and make explicit, for any $X''$ in Cat, the set of arrows $g$ from $X'$ to $X''$, as a functor in $X''$. It is the set of triples
\[
(g_X, g_{U'}, g_{T'}) \qquad
\left\{\begin{aligned} g_X &: X \longrightarrow X'' \\ g_{U'} &: U' \longrightarrow X'' \\ g_{T'} &: T' \longrightarrow X''\end{aligned}\right.
\qquad \text{with} \quad
\begin{aligned} g_X\circ i &= g_{U'} f_U \\ g_X \circ j &= g_{T'} f_T. \end{aligned}
\tag*{$(*)$}
\]
\printedpage{49}
% head: A. Grothendieck : Dérivateurs, Ch. XII, Édition des Universités de Montpellier II et Paris VII
% foot: Version 7 décembre 2002
\[
\begin{tikzcd}
U \arrow[r, hook, "i"] \arrow[d, "f_U"'] & X \arrow[dd, "g_X"'] & T \arrow[l, hook', "j"'] \arrow[d, "f_T"] \\
U' \arrow[dr, "g_{U'}"'] & & T' \arrow[dl, "g_{T'}"] \\
& X'' &
\end{tikzcd}
\]

Note that for such a triple, $g_X$ is partially known already in terms of $g_{U'}$, $g_{T'}$, by its restrictions $g_X\circ i$ and $g_X\circ j$ to $U$ and $T$ respectively, by the formulas
\[
g_U = g_{U'}f_U, \qquad g_T = g_{T'}f_T.
\]

\msp{page 68} Thus the datum of a triple $(*)$ is equivalent to that of a triple
\[
(g_{U'}, g_{T'}, g_X) \qquad
\left\{\begin{aligned}
&g_X \text{ an arrow } X\longrightarrow X'' \text{ corresponding to}\\
&g_U = g_{U'}\circ f_U \text{ and } g_T = g_{T'}\circ f_T \text{ already given (via } g_{U'}\text{, } g_{T'}\text{),}
\end{aligned}\right.
\]
so by virtue of 1.3. this corresponds to
\[
(g_{U'}, g_{T'}, u) \qquad
\left\{\begin{aligned}
g_{U'} &: U' \longrightarrow X''\\
g_{T'} &: T' \longrightarrow X''\\
u &: H \longrightarrow \underbrace{H_{X''}\left((g_{U'}\circ f_U)^{\circ}\times(g_{T'}\circ f_T)\right)}_{(u,t)\longmapsto \Hom_{X''}(g_{U'}f_U(u),g_{T'}f_T(t))}.
\end{aligned}\right.
\tag*{$(*)$}
\]
We are going to interpret the datum of $u$, in terms of a functor
\[
F^{*} = (H'\longmapsto H'\circ F) : \uHom(U'^{\circ}\times T',\Ens) \longrightarrow \uHom(U^{\circ}\times T,\Ens)
\tag*{$(**)$}
\]
where
\[
F : U^{\circ}\times T \longrightarrow U'^{\circ}\times T' \qquad \text{is given by } F = f_U^{\circ}\times f_T.
\]
We have an object $H''=H''_{g_{U'},g_{T'}}$ of the source category of $(**)$ determined in terms of the pair $(g_{U'},g_{T'})$, by
\[
H''(u',t') = \Hom_{X''}(g_{U'}(u'),g_{T'}(t')),
\]
and the target bifunctor of the arrow $u$ in $(*)$ is interpreted as
\[
F^{*}(H'') = H''\circ F : (u,t)\longmapsto \underbrace{H''(f_U(u),f_T(t))}_{=\Hom(g_{U'}f_U(u),g_{T'}f_T(t))} \quad .
\]
Thus, $u$ is interpreted as an arrow
\[
u : H \longrightarrow F^{*}(H'').
\]
Now the functor $F^{*}$ has a left adjoint $F_{!}$, i.e. we have
\[
\Hom(H,F^{*}(H'')) \simeq \Hom(F_{!}(H),H'').
\tag*{$(*{*}*)$}
\]

\printedpage{50}
% head: A. Grothendieck : Dérivateurs, Ch. XII, Édition des Universités de Montpellier II et Paris VII
% foot: Version 7 décembre 2002
\msp{page 69} with
\[
F_{!} : \uHom(U^{\circ}\times T,\Ens) \longrightarrow \uHom(U'^{\circ}\times T',\Ens).
\]
Put
\[
H' = F_{!}(H) : U'^{\circ}\times T' \longrightarrow \Ens,
\]
it is a bifunctor well determined by the sole datum of the initial diagram $(*)$ p. 67, or what amounts to the same thing, in terms of % [sic] déterminée -> déterminé, accord avec « bifoncteur » (masculin)
\[
(U,T,H,f_U : U\longrightarrow U', f_T : T\longrightarrow T').
\]

And we see that the datum of an arrow $X'\longrightarrow X''$, $X''$ a variable object of Cat, is equivalent to that of
\[
g_{U'} : U'\longrightarrow X'',\ g_{T'} : T'\longrightarrow X'',\ u : H'\longrightarrow H'' = H_{X''}(g_{U'}^{\circ}\times g_{T'}),
\]
that is to say, by virtue of 1.3., to that of an arrow
\[
X'_0 \longrightarrow X'',
\]
where $X'_0$ denotes the object of Cat obtained from $U'$, $T'$ by gluing by means of $H'$. This therefore shows that the direct limit sought is none other than $X'_0$. We have thus found the

\textbf{Theorem 1.} \textit{The direct limit $X'$ in} Cat \textit{of the diagram} $(*)$ \textit{page 67, defined in terms of a triple $(X,U,T)$ in $\mathcal{T}$ and arrows $f_U : U\longrightarrow U'$, $f_T : T\longrightarrow T'$ in} Cat\textit{, exists (without merit)} \msp{page 70} \textit{and the canonical arrows}
\[
U' \overset{i'}{\hookrightarrow} X' \overset{j'}{\hookleftarrow} T'
\]
\textit{are open and closed immersions complementary to one another. The gluing datum which defines $X'$ via $U'$, $T'$ is given by}
\[
H' = (\underbrace{f_U^{\circ}\times f_T}_{F})_{!}(H) : U'^{\circ}\times T' \longrightarrow \Ens,
\]
\textit{where $H$ is the gluing datum which defines $X$, i.e. $H : U^{\circ}\times T\longrightarrow$} Ens \textit{is the bifunctor $(u,t)\longmapsto \Hom_X(u,t)$.}

Thus we have, for $u'$ in $U'$ and $t'$ in $T'$
\[
\Hom_{X'}(u',t') = F_{!}(H)(u',t'),
\]
and it remains to make $F_{!}(H)$ explicit. Here $F=f_U^{\circ}\times f_T$, and $H$ is a covariant functor on $U^{\circ}\times T$, we can interpret it as a presheaf on $U\times T^{\circ}$, and hence write
\DERmark{ENV0720}

\[
\underbrace{F_{!}(H)(u',t')}_{\Hom_{X'}(u',t')} = \varinjlim_{(u',t')\backslash U\times T^{\circ}} \underbrace{\Hom_X(u,t)}_{H(u,t)} \ .
\tag*{$(*)$}
\]

\printedpage{51}
% head: A. Grothendieck : Dérivateurs, Ch. XII, Édition des Universités de Montpellier II et Paris VII
% foot: Version 7 décembre 2002
The category $(u',t')\backslash U\times T^{\circ}$ over which we take the $\varinjlim$ is made explicit as a product
\[
(u',t')\backslash U\times T^{\circ} \simeq (u'\backslash U)\times(T/t')^{\circ}.
\]
The maps
\[
\Hom_X(u,t)\longrightarrow \Hom_{X'}(u',t') \qquad (=F_{!}(H)(u',t'))
\]
\msp{page 71} deduced from this formula are none other than the following: $u$, $t$ are given as objects of $u'\backslash U$, $T/t'$, i.e. with arrows
\[
u'\xrightarrow{\alpha} f_U(u), \qquad f_T(t)\xrightarrow{\beta} t',
\]
whence
\[
\Hom_X(u,t) \xrightarrow{\text{via } f=f_X:X\to X'} \Hom_{X'}(\underbrace{f(u)}_{f_U(u)},\underbrace{f(t)}_{f_T(t)}) \xrightarrow{\text{via } \alpha,\ \beta} \Hom_{X'}(u',t').
\]

Let us note for the record that for every object $Z$ of Cat, and every presheaf $G$ on $Z$, we have the formula
\[
\varinjlim_{Z} G \simeq \pi_0(Z/G),
\]
where $Z/G$ is interpreted in terms of the canonical immersion $Z\longrightarrow Z^{\wedge}$ $(G\in \Ob Z^{\wedge})$. It is therefore a fibered category over $Z$ with discrete fibers, and we have in fact an isomorphism between the full subcategory of $\Cat/Z$ formed by such fibered categories with discrete fibers, with $Z^{\wedge}$. In this perfect dictionary between presheaves of sets and fibered categories with discrete fibers, the functor $\varinjlim$ on $Z^{\wedge}$ corresponds simply to the functor $\pi_0$ on the category of objects over $Z$ in Cat, itself induced by the functor (independent of $Z$)
\[
\Cat \xrightarrow{\pi_0} \Ens.
\]

\msp{page 72} We can therefore rewrite the fundamental formula $(*)$ p. 70, in terms of the functor $H$ defining the gluing $X$ of $U$ and of $T$, as
\[
\boxed{\underbrace{\Hom_{X'}(u',t')}_{H'(u',t')} \simeq \pi_0(\underbrace{(u',t')\backslash U\times T^{\circ}/H}_{((u'\backslash U)\times(T/t')^{\circ})/H}).}
\tag*{$(*)$}
\]

In what follows, we will probably only have to use the case where we have either $T'=T$, $f_T=\id_T$, or the reverse. We will work mainly in the first case, deducing what we will need for the dual case, by the principle of duality (automorphism $X\longmapsto X^{\circ}$ of the
\printedpage{52}
% head: A. Grothendieck : Dérivateurs, Ch. XII, Édition des Universités de Montpellier II et Paris VII
% foot: Version 7 décembre 2002
category Cat). In this case, in the fundamental formula $(*)$ the element $t'$ of $T'$ is written as an element $t$ of $T$, and the formula is rewritten as
\[
\boxed{\Hom_{X'}(u',t) \simeq \pi_0(\underbrace{u'\backslash U/t}_{\text{category of triples } (u,\alpha,\beta'),\ \left\{\begin{aligned} &u\in \Ob U\\ &\alpha : u\longrightarrow t\\ &\beta' : u'\longrightarrow f_U(u)\end{aligned}\right.}).}
\tag*{$(**)$}
\]

This is the formula that I had found directly in this case, without yet bothering to vary $T$ at the same time as $U$. I would like to make the link between the two formulas $(*)$ and $(**)$, by topological means, thus establish a canonical bijection, functorial etc.
\[
\pi_0(u'\backslash U/t) \xrightarrow{\sim} \pi_0(((u'\backslash U)\times(T/t)^{\circ})/H).
\]

\msp{page 73} We surely have a functor
\[
u'\backslash U/t \longrightarrow ((u'\backslash U)\times(T/t)^{\circ})/H
\]
which associates to a triple
\[
(u,\alpha,\beta') : u\in \Ob U,\ \alpha : u\longrightarrow t,\ \beta' : u'\longrightarrow f_U(u)
\]
the triple $(\bar u,\bar t,\xi)$ $(\xi\in H(\bar u,\bar t))$ on the right-hand side, with $\bar u$ in $u'\backslash U$, $\bar t$ in $T/t$, defined as follows:
\[
\bar u = (u,\beta'),\ \bar t = (t,\id_t), \underbrace{\xi}_{=\alpha} \in H(\bar u,\bar t) = \Hom(u,t).
\]
In short
\[
u'\backslash U/t = \text{fiber of } ((u'\backslash U)\times(T/t)^{\circ})/H \text{ at the initial object } \bar t \text{ of } (T/t)^{\circ}
\]
under the natural projection, composite arrow
\[
\begin{tikzcd}
((u'\backslash U)\times(T/t)^{\circ})/H \arrow[d] \\
(u'\backslash U)\times(T/t)^{\circ} \arrow[d, "\pi"] \\
(T/t)^{\circ}.
\end{tikzcd}
\]

We have to see that the inclusion of the fiber gives a bijection on the $\pi_0$, or what amounts to the same thing, that the inclusion of the fiber $u'\backslash U$ \textit{of} $\pi$ at the initial object $(t,\id_t)$ of $(T/t)^{\circ}$ gives a bijection on the $\varinjlim$ of the functor $H$. So the inclusion had better be aspheric, or at least co-cofinal. % [sic] co-cofinalle -> co-cofinale, orthographe


\endgroup
% END UNIT Der-ChXIIb-012/Der-ChXIIb-012_p006-010.tex
% BEGIN UNIT Der-ChXIIb-012/Der-ChXIIb-012_p011-015.tex
\begingroup\makeatletter%
\newcommand{\printedpage}[1]{\marginpar{\footnotesize p.~#1}}%
\newcommand{\msp}[1]{\marginpar{\footnotesize [#1]}}%
\newcommand{\Ob}{\mathrm{Ob}\,}%
\newcommand{\Hom}{\mathrm{Hom}}%
\makeatother\let\printedpage\derpp\let\msp\dermsp%
\let\mspin\DERmspin%


% ===================== page 11 =====================
\printedpage{53}
% head: A. Grothendieck : Dérivateurs, Ch. XII, Édition des Universités de Montpellier II et Paris VII

\msp{page 74}
But the inclusion of an initial object is always aspheric, and remains so under product with any object (i.e. $u'\backslash U$) of Cat (and even under any smooth base change \dots), whence the conclusion.

In what follows, we are mainly interested in the situation of amalgamated sums (where we change notation a little, now)
\[
\begin{tikzcd}
U \arrow[r, hook, "i"] \arrow[d, "p"'] & X \arrow[d, "q"'] & T \arrow[l, hook', dotted, "j"'] \arrow[d, dotted] \\
U' \arrow[r, hook, "i'"] & X' & T' \simeq T. \arrow[l, hook', dotted]
\end{tikzcd}
\tag*{$(*)$}
\]

The closed subcategory $T'$ complementary to $U'$ in $X'$ is then canonically isomorphic to the closed subcategory $T$ complementary to $U$ in $X$, the isomorphism being induced by $q$. Often we shall identify $T'$ and $T$, which is legitimate on viewing $X'$ as obtained from the gluing of $U'$ and $t$ \texttt{[rather }$T$\texttt{]}, by the bifunctor $U'^{\circ}\times T\longrightarrow \mathrm{Ens}$ made explicit in formula $(**)$ p.~72. We are mainly interested in isolating conditions which ensure the following two implications

\begin{itemize}
\item[a)] $i\in W\Longrightarrow i'\in W$ (and even $\Longleftrightarrow$).
\item[b)] $p\in W\Longrightarrow q\in W$ (and even $\Longleftrightarrow$).
\end{itemize}

\msp{page 75}
\textbf{Proposition 4.} \textit{Let $t\in\Ob T$. For $t\in\Ob\bar{U}$ to hold, it is necessary and sufficient that $q(t)\in\Ob\overline{U'}$. In other words, the ``boundary'' $\dot{T}=T\cap\bar{U}$ of $T$ in $X$ is ``the same'' as its boundary in $X'$.}

\textsc{Proof}. To say that $t\in\Ob\bar{U}$ also means
\[
U/t\neq\emptyset
\]
and implies, by continuity, that $q(t)\in\overline{q(U)}\subseteq\overline{U'}$. Conversely, if $q(t)\in\overline{U'}$, i.e. if there exists $u'\in U'$ such that
\[
\Hom(u',q(t))(=\pi_0(u'\backslash U/t))\neq\emptyset,
\]
it follows that $u'\backslash U/t\neq\emptyset$, a fortiori $U/t$ (which receives $u'\backslash U/t)\neq\emptyset$, q.e.d.

\textbf{Proposition 5.} \textit{Let}
\[
\begin{tikzcd}
Y \arrow[r, "i"] \arrow[d, "p"'] & X \arrow[d, "q"] \\
Y' \arrow[r, "i'"] & X'
\end{tikzcd}
\]
% foot: Version 7 décembre 2002
% ===================== page 12 =====================
\printedpage{54}
% head: A. Grothendieck : Dérivateurs, Ch. XII, Édition des Universités de Montpellier II et Paris VII
\textit{be a cocartesian square in \textup{Cat}. If $i$ is an open immersion (resp. closed), then $i'$ too, and the square is also cartesian. Ditto for $p$ and $q$. If $i$ and $p$ are open immersions (\footnote{\textbf{NB} We use \textit{``open immersion''}, \textit{``closed immersion''} in the \textit{strict sense}\,: it is a fully faithful \textit{monomorphism} whose image is a subcategory that is \textit{full,} either open or closed.}) (resp. \dots), then the situation is isomorphic to that of an object $X'$ of \textup{Cat}, of two open (resp. closed) subcategories $X$, $Y'$ of $X'$ having as union $X'$, as intersection} % [sic] de $X$ -> de $X'$, X and Y' are opens of X' (their union is X')
\DERmark{ENV0721}

\textit{$Y$. Conversely, in such a situation, the corresponding square is cocartesian, i.e. $X$} % [sic] situationm -> situation, (comma), typing error
\DERmark{ENV0722}

\texttt{[rather }$X'$\texttt{]} \textit{is identified with the amalgamated sum}
\[
X'\simeq X\amalg_{Y}Y'.
\]

Thus we also have proposition 15, page 158, and its corollaries.
% foot: Version 7 décembre 2002

% ===================== page 13 =====================
\printedpage{55}
% head: A. Grothendieck : Dérivateurs, Ch. XII, Édition des Universités de Montpellier II et Paris VII

\section*{2\quad Amalgamated sums and $W$-equivalences.}

\subsection*{2.1\quad Cocartesian square axiom and cube axiom (or ``Five Lemma'').}

We now suppose given a basic localizer
\[
W\subseteq\mathrm{Fl}(\mathrm{Cat}),
\]
so $W$ satisfies W(1, 2 bis, 3). We do not assume for the moment for the other standard axioms W(4) (fibration axiom), W(5) (connectedness axiom), W(6) (direct limit axiom). But we are now going to introduce two additional stability axioms, of a different nature, concerning amalgamated sums.

\begin{description}[leftmargin=5em, style=nextline]
\item[W(7 bis)] (Cocartesian square axiom.) Let $X$ be an object of Cat, $U$, $U'$ two open subcategories of $X$ such that $U\cup U'=X$, let us consider the corresponding cocartesian diagram (cf. prop. 5 in the margin of the page \texttt{[above]}):

\msp{page 76}
\[
\begin{tikzcd}[column sep=small]
& U\cap U' \arrow[dl, hook', "i'_0"'] \arrow[dr, hook, "i_0"] & \\
U \arrow[dr, hook, "i"'] & & U' \arrow[dl, hook', "i'"] \\
& X &
\end{tikzcd}
\tag*{$(*)$}
\]

Then
\[
i_0\in W\qquad\Longleftrightarrow\qquad i\in W
\]
(and symmetrically
\[
i'_0\in W\qquad\Longleftrightarrow\qquad i'\in W).
\]
\end{description}

\textbf{NB} By duality, this axiom is equivalent to the similar axiom, with \textit{closed subcategories} $T$, $T'$ of $X$ with union $X$.

\textbf{Corollary.} \textit{Assume \textup{W(7 bis)}. Consider the following conditions on the preceding square.}

\begin{itemize}
\item[\textit{a)}] \textit{$U\cap U'$ $W$-aspheric.}
\item[\textit{b)}] \textit{$U$, $U'$ $W$-aspheric.}
\item[\textit{c)}] \textit{$X$ $W$-aspheric.}
\end{itemize}

\textit{Then (a and b) implies c, and (b and c) implies a.}

\textbf{NB} Morally, (c and a) ought also to imply b, i.e. if $U\cap U'$ and $X$ are $W$-aspheric, the same is true of $U$ and $U'$. But it does not seem that the cocartesian square axiom allows one to prove this. On the other hand, the axiom which follows will imply the preceding corollary in its complete form.
% foot: Version 7 décembre 2002

% ===================== page 14 =====================
\printedpage{56}
% head: A. Grothendieck : Dérivateurs, Ch. XII, Édition des Universités de Montpellier II et Paris VII

\begin{description}[leftmargin=5em, style=sameline]
\item[W(7)] (``Cube axiom'', or ``Five Lemma''.) Let there be a commutative cubical diagram
\[
\begin{tikzcd}[row sep=small, column sep=small]
U_1\cap U'_1 \arrow[rr] \arrow[dr] \arrow[dd, "f_{U_1\cap U'_1}"'] & & U_1 \arrow[dr] \arrow[dd, "f_{U_1}"{swap,pos=.25}] & \\
& U'_1 \arrow[rr, crossing over] & & X_1 \arrow[dd, "f"] \\
U_2\cap U'_2 \arrow[rr] \arrow[dr] & & U_2 \arrow[dr] & \\
& U'_2 \arrow[rr] \arrow[from=uu, crossing over, "f_{U'_1}"{pos=.25}] & & X_2
\end{tikzcd}
\]
representing a morphism (vertical arrows) from a

\msp{page 77}
cocartesian square\quad $\vcenter{\hbox{\begin{tikzcd}
U_1\cap U'_1 \arrow[r] \arrow[d] & U_1 \arrow[d] \\
U'_1 \arrow[r] & X_1
\end{tikzcd}}}$\quad of the type considered in W(7 bis) ($X_1$ union of two open subcategories $U_1$, $U'_1$) into another (with subscripts 2). Consider the following conditions.

\begin{itemize}
\item[a)] $f_{U_1\cap U'_1}\in W$.
\item[b)] $f_{U_1}$ and $f_{U'_1}$ are \texttt{[in]} $W$.
\item[c)] $f\in W$.
\end{itemize}

Then any two of these conditions imply the third. In other words, when one of these conditions is satisfied, the other two are equivalent to each other.
\end{description}

\textbf{NB} This axiom is equivalent to the dual axiom, in terms of objects $X_1$, $X_2$ each represented as a union of two closed subcategories.

\textbf{Corollary.} \textit{Assume \textup{W(7)} satisfied. Then the corollary to \textup{W(7 bis)} (an axiom which we do not have to use here), in its completed form (cf. the NB following it) is valid: any two of the conditions a), b), c) imply the third.}

It suffices to look at the unique morphism from the diagram $(*)$, p.~76, into the trivial diagram\quad $\vcenter{\hbox{\begin{tikzcd}[column sep=small]
& e \arrow[dl] \arrow[dr] & \\
e \arrow[dr] & & e, \arrow[dl] \\
& e &
\end{tikzcd}}}$\quad corresponding to the case where $X$ is the final category $e$, and $U$, $U'$ equal to $X$. Then the corollary is nothing other than W(7) applied to this particular case.

This argument makes one feel that W(7) is probably stronger than W(7 bis). Prompted by this feeling, I notice with pleasure that it is well founded:
% foot: Version 7 décembre 2002

% ===================== page 15 =====================
\printedpage{57}
% head: A. Grothendieck : Dérivateurs, Ch. XII, Édition des Universités de Montpellier II et Paris VII

\msp{page 78}
\textbf{Proposition 6.} \textit{The axiom \textup{W(7)} (five lemma) implies \textup{W(7 bis)} (cocartesian square axiom).}

We look at the following cubical diagram
\[
\begin{tikzcd}[row sep=small, column sep=small]
U\cap U' \arrow[rr, "i_0"] \arrow[dr, "i'_0"] \arrow[dd, "f_{U\cap U'}=i_0"'] & & U' \arrow[dr, "i'"] \arrow[dd, "f_{U'}=\mathrm{id}"{swap,pos=.25}] & \\
& U \arrow[rr, crossing over, "i"] & & X \arrow[dd, "f=\mathrm{id}_X"] \\
U' \arrow[rr, "\mathrm{id}"] \arrow[dr, "i'"] & & U' \arrow[dr, "i'"] & \\
& X \arrow[rr, "\mathrm{id}"] \arrow[from=uu, crossing over, "f_U=i"{pos=.25}] & & X,
\end{tikzcd}
\]
from the given cocartesian square into the cocartesian square obtained from $X$ and its two open subcategories $X$, $U'$, with intersection $U'$. Here the two vertical arrows $f$ and $f_{U'}$ are isomorphisms, hence \texttt{[in]} $W$. With the notation of W(7), we therefore have c) ($f\in W$), hence a) and b) are equivalent, i.e. $i_0\in W$ and $i\in W$ are equivalent.

\hfill q.e.d.

\textbf{Theorem 2.} \textit{The basic localizer $W_\infty$ satisfies the axiom \textup{W(7)}.}

This is surely a consequence of the cohomological criterion of \textsc{Artin-Mazur}, taking into account the Mayer-Vietoris exact sequence for both of the two cocartesian squares, and the noncommutative variants of Mayer-Vietoris. I have not written out the verification in due form, and should probably repair this omission.

\msp{page 79}
\textbf{Remark.} The dual version (in the sense of \textsc{Quillen}) of the cocartesian square axiom W(7 bis), i.e. the corresponding axiom for cartesian squares, is not valid for $W_\infty$. The relevant statement would be the following: let
\[
\begin{tikzcd}
X & X' \arrow[l, "q"'] \arrow[d, "f'"] \\
S & S' \arrow[l, "p"']
\arrow[from=1-1, to=2-1, "f"']
\end{tikzcd}
\]
be a cartesian square, with $f$ a \textit{$W$-fibration.} Then $f\in W\Longleftrightarrow f'\in W$. If $W$ satisfies W(4c) (which is the case for $W_\infty$), then $f\in W$ does indeed imply that its fibers are $W$-aspheric, hence that it is a universal $W$-equivalence, i.e. $f'\in W\ \forall\, p:S'\longrightarrow S$. But for the reverse implication $f'\in W\Longrightarrow f\in W$, one must evidently assume that
\[
\pi_0(p):\pi_0(S')\longrightarrow\pi_0(S)
\]
% foot: Version 7 décembre 2002


\endgroup
% END UNIT Der-ChXIIb-012/Der-ChXIIb-012_p011-015.tex
% BEGIN UNIT Der-ChXIIb-012/Der-ChXIIb-012_p016-020.tex
\begingroup\makeatletter%
\newcommand{\printedpage}[1]{\marginpar{\footnotesize p.~#1}}%
\newcommand{\msp}[1]{\marginpar{\footnotesize [#1]}}%
\newcommand{\circled}[1]{\tikz[baseline=(c.base)]\node[draw,circle,inner sep=1pt](c){#1};}%
\newcommand{\Ob}{\operatorname{Ob}}%
\makeatother\let\printedpage\derpp\let\msp\dermsp%
\let\mspin\DERmspin%


\printedpage{58}
% head: A. Grothendieck : Dérivateurs, Ch. XII, Édition des Universités de Montpellier II et Paris VII
% foot: Version 7 décembre 2002

is \textit{surjective.} (E.g. if $S'=\emptyset$, we have $X'=\emptyset$, hence $f'\in W$, without that implying $f\in W$!) The fact that in the cocartesian case of W(7 bis), $i_0\in W\Longrightarrow i\in W$ must be regarded as ``going without saying'' - otherwise the square does not deserve, from the $W$-homotopic point of view, the name ``cocartesian''. On the other hand, the reverse implication $i\in W\Longrightarrow i_0\in W$ appears to me as a very particular feature of the derivator HOT$_W$ (when it has it), and notably % [sic] dérivatuer -> dérivateur, misprint
of HOT $=$ HOT$_{W_\infty}$. The beginning
\msp{page 80}
of the remark simply means that the dual derivator HOT$^\circ$ does not have this same property. I must refer in this connection to the two axioms for derivators corresponding to W(7 bis), W(7), which I will probably do shortly, so as to give a more precise meaning to the present comments.

\subsection*{2.2 Application to cocartesian squares deduced from an open or closed immersion.}

The axioms W(7), W(7 bis) (of the cube, and of the cocartesian square) concern cocartesian squares of a very particular nature. We are now going to see how to deduce from them statements of the same type for much more general cocartesian squares, including for squares which have nothing cocartesian about them (in Cat) - and which we will call % [sic] appelerons -> appellerons, misspelling
``homotopically cocartesian'' (relative to $W$), or simply $W$-cocartesian. (There will also be $W$-cartesian squares one day \dots) In the case where $W=W_\infty$, I will also call them Hot-cocartesian. % [sic] appelerai -> appellerai, misspelling

I will begin with the particular case of a cocartesian square associated with an open immersion (diagram ($*$), p. 74):
\[
(*)\qquad\qquad
\begin{tikzcd}
U \arrow[r, hook, "i"] \arrow[d, "p"'] & X \arrow[d, "q"'] & T \arrow[l, hook', dotted, "j"'] \arrow[d, dotted] \\
U' \arrow[r, hook, "i'"] & X' & T'\simeq T. \arrow[l, hook', dotted, "j'"']
\end{tikzcd}
\]
\msp{page 81}

The big trap here that we have to manage to exorcise is that it can happen that $i\in W$, without our having $i'\in W$, and also $p\in W$ without our having $q\in W$. I had convinced myself of this in 1983 while writing Pursuing Stacks (during the weeks of ``scratchwork''), and don't feel like looking for a counterexample here, unless it presents itself on its own (which it probably will).

I am going to give an important case where these troubles disappear:

\textbf{Proposition 6:} \textit{Consider the closures $\bar U$ and $\overline{U'}$ of $U$ and $U'$ in $X$, $X'$ respectively, and suppose that the inclusions}
\[
U\xrightarrow{i_0}\bar U,\qquad U'\xrightarrow{i'_0}\overline{U'}
\]
\printedpage{59}
% head: A. Grothendieck : Dérivateurs, Ch. XII, Édition des Universités de Montpellier II et Paris VII
% foot: Version 7 décembre 2002
\textit{are in $W$} (\footnote[64]{cf. th. 3 p. 107 for general conditions ensuring that this hypothesis is satisfied.}). \textit{Then}
\[
\left\{
\begin{array}{llll}
\circled{\textit{a}}\ i\in W & \iff & i'\in W & \textit{if we have }\mathrm{W}(7\text{ bis})\\
\circled{\textit{b}}\ p\in W & \iff & q\in W & \textit{if we have }\mathrm{W}(7).
\end{array}
\right.
\]

\textsc{Proof}. Let us first consider case \circled{a} of $i$, $i'$. Let
\[
k:\bar U\hookrightarrow X
\]
be the inclusion. The triangle
\begin{tikzcd}
U \arrow[r, "i_0"] \arrow[rr, bend right=8, "i"'] & \bar U \arrow[r, "k"] & X
\end{tikzcd}
, where $i_0\in W$, shows that
\[
i\in W\iff k\in W.
\]

On the other hand, $X$ is the union of the closed subcategories $\bar U$, $T$, so we have a cocartesian square in Cat
\[
(*)\qquad\qquad
\begin{tikzcd}
\bar U\cap T=\dot T \arrow[r, "k_0"] \arrow[d, "j_0"'] & T \arrow[d, "j"] \\
\bar U \arrow[r, "k"] & X.
\end{tikzcd}
\]

The axiom W(7 bis) (in its dual form (i.e. applied to the square formed by the opposite categories \dots)) implies
\msp{page 82}
that
\[
k\in W\iff k_0\in W,
\]
whence
\[
i\in W\iff (k_0:\dot T\longrightarrow T)\in W.
\]
Likewise we find, applying this to $(X',U')$ in place of $(X,U)$ % [sic] appliqunt -> appliquant, misprint
\[
i'\in W\iff (k'_0:\dot T'\longrightarrow T')\in W.
\]
Now we have $T\xrightarrow{\sim}T'$ (isomorphism), and by proposition 4 this isomorphism takes $\dot T$ to $\dot T'$. Thus
\[
k_0\in W\iff k'_0\in W\qquad\quad\text{(trivially)},
\]
whence a).

Let us prove b). The commutative square induced by ($*$), p. 80,
\printedpage{60}
% head: A. Grothendieck : Dérivateurs, Ch. XII, Édition des Universités de Montpellier II et Paris VII
% foot: Version 7 décembre 2002
\[
\begin{tikzcd}
U \arrow[r, hook, "i_0"] \arrow[d, "p"'] & \bar U \arrow[d, "\bar p"] \\
U' \arrow[r, hook, "i'_0"] & \overline{U'},
\end{tikzcd}
\]
where $i_0,i'_0\in W$ by hypothesis, shows that
\[
p\in W\iff\bar p\in W,
\]
so everything comes down to proving that
\[
\bar p\in W\overset{?}{\iff}q\in W.
\]
But $q$ induces a homomorphism from the square ($*$) deduced from $(X,\bar U,T)$ (with $X=\bar U\cup T$), into the similar square for $X'$
\[
\begin{tikzcd}
\overline{U'}\cap T'=\dot T' \arrow[r, "k'_0"] \arrow[d, "j'_0"'] & T' \arrow[d, "j'"] \\ % [sic] T -> T', the square for X' should have the corner T' (although T'\simeq T)
\overline{U'} \arrow[r, "k'"] & X',
\end{tikzcd}
\]
\DERmark{ENV0723}

the vertical arrows being denoted by $q_{\dot T}$, $q_T$, $\bar p$, $q$.
\msp{page 83}
We know that $q_{\dot T}$, $q_T$ are isomorphisms, hence in $W$, so the ``five lemma'' W(7) assures us that $\bar p\in W\iff q\in W$, q.e.d.

\textbf{Comment.} This proposition leads us to look for conditions on an open subcategory $U$ of an $X$, which ensure that the inclusion $i_0:U\longrightarrow\bar U$ is \texttt{[in]} $W$, and \textit{that this property is preserved by every cobase change $U\longrightarrow U'$.} Moreover, we seek conditions which are of a local nature on $\bar U$, more precisely on $\bar U\cap T=\dot U=\dot T$. We can motivate this concern (a posteriori) by saying that we would like the property $i_0\in W$ to be preserved, not only by an arbitrary cobase change $U\longrightarrow U'$, but also, in addition, by passing from $X$ to any full subcategory $X_1$ containing $U$, thus giving rise to a situation
\[
U\overset{i_Y}{\hookrightarrow}X_1,\qquad U\xrightarrow{(i_Y)_0}\bar U^{X_1}=\bar U\cap X_1
\]
\texttt{[}$Y=X_1$\texttt{?]}, similar to the starting situation. More generally still, we would like it to be preserved on passing from the triple $(X,U,T)$ to an ``induced'' triple $(X_1,U,T_1)$, by an arbitrary arrow $g:T_1\longrightarrow T$ with target $T$ in Cat, giving rise to the pair of arrows
\[
U\xrightarrow{\mathrm{id}}U,\ g:T_1\longrightarrow T
\]
(cf. pages 60, 61). The preceding case is the one where
\printedpage{61}
% head: A. Grothendieck : Dérivateurs, Ch. XII, Édition des Universités de Montpellier II et Paris VII
% foot: Version 7 décembre 2002
\msp{page 84}
$T_1$ is a full subcategory of $T$, $g$ the inclusion arrow. Let us note in this connection the sorital proposition % [sic] inclusion, Notons -> inclusion. Notons, comma before a capital letter
\DERmark{ENV0724}


\textbf{Proposition 7.} \textit{Equivalent conditions on the open subcategory $U$ of an object $X$ of} Cat (assuming W(7 bis)):

\begin{enumerate}[label={}, leftmargin=*]
\item[\textit{a)}] \textit{The inclusion $i_0:U\longrightarrow\bar U$ is $W$-aspheric, i.e. for every $t\in\dot U=\bar U\backslash U$, the category $U/t$ is $W$-aspheric.}
\item[\textit{b)}] \textit{The inclusion $i_0:U\longrightarrow\bar U$ is in $W$, and this condition remains satisfied, when we replace $(X,U)$ by $(X_1,U)$, where $(X_1,U)$ is deduced from $(X,U)$ by means of an arrow $T_1\longrightarrow T(=X\backslash U)$ in} Cat\textit{, as recalled above.}
\item[\textit{b$'$)}] \textit{As in b), in the case where $T_1$ is the one-point category $e$, and where moreover its image is in $\dot U$.}
\end{enumerate}

\textsc{Proof}. It is obvious that if in b), b$'$) we replace the condition ``$i_0\in W$'' by the stronger condition ``$i_0$ is $W$-aspheric'', the four \texttt{[rather three]} conditions a) to b$'$) are equivalent. It remains to prove that the condition that the $U/t$ be $W$-aspheric ($t\in\bar U$ - but this is trivial for $t\in U$, so it suffices to express it for $t\in\bar U\backslash U=\dot U$) is already implied by the weakest b$'$ and we guess that we have this:

\msp{page 85}

\textbf{Corollary.} (Assuming W(7 bis).) \textit{Let $X$ be in} Cat\textit{, $U$ an open subcategory of $X$ such that $X\backslash U=\overline{\{t\}}$ is the final category. Then the inclusion $U\xrightarrow{i}X$ is in $W$ if and only if it is $W$-aspheric, i.e. if and only if $U/t$ is $W$-aspheric (where $t$ is the unique element of $T$).}

We prove this corollary later (\footnote[65]{cf. cor. 4 of proposition 15, page 149.}). Proposition 7 will be of no use to us until then - it will have only a purely aesthetic role to play, to justify our deliberate intention of looking for conditions on $i:U\longrightarrow X$ (to ensure that $U\longrightarrow\bar U$ is in $W$ and remains so, when we replace $(X,U)$ by other pairs $(X_1,U_1)$ which are deduced from it by certain operations \dots) which
\msp{page 86}
are of a \textit{local nature} on $\bar U$ (in fact, at the points of $\dot U$), and which in any event imply the \textit{$W$-asphericity} of $U\longrightarrow\bar U$.

I had hoped that this asphericity condition (suggested by the notion ``of a Dwyer immersion'') would be stable under an arbitrary cobase change $U\longrightarrow U'$, which would at the same time have ensured that after such an operation, the open immersion $i':U'\longrightarrow X'$ obtained still satisfies the condition that $i'_0:U'\longrightarrow\overline{U'}$ is in $W$. We will see that this expectation is unfounded - the condition of $W$-asphericity is \textit{not} stable under cobase change, and it does \textit{not} even suffice to ensure that after a cobase change, $i'_0:U'\longrightarrow\overline{U'}$ is
\printedpage{62}
% head: A. Grothendieck : Dérivateurs, Ch. XII, Édition des Universités de Montpellier II et Paris VII
% foot: Version 7 décembre 2002
in $W$. A stronger notion will be needed, which is suggested by the case where $T=\{t\}$ is the final category and contained in $\bar U$ (the case to which we tried to reduce ourselves in proposition 7). In this case, as we said (preceding corollary), the condition that $i:U\longrightarrow\bar U=X$ is in $W$, is equivalent to the $W$-asphericity of $i$, i.e. to the $W$-asphericity of $U/t$, and likewise for $i':U'\longrightarrow\overline{U'}=X'$: it is in $W$ if and only if $U'/t'$ is $W$-aspheric, where
\msp{page 87}
$t'$ is the unique element of $T'(\simeq T)=X'\backslash U'$. So we need conditions on $U/t$ which ensure that not only $U/t$, but also all the similar categories $U'/t$ \texttt{[rather} $U'/t'$\texttt{]}, are $W$-aspheric. In the general case where $T$ is arbitrary, it is this condition that we are going to impose on the $U/t$ ($t\in\Ob T$), and it will turn out that this time, they are stable \texttt{[rather it is stable]} under every cobase change $U\xrightarrow{p}U'$, i.e. that the $U'/t'$ ($t'\in\Ob T'$) will then satisfy the same condition. It is this condition that we are now going to reformulate.

\subsection*{2.3 Totally $W$-aspheric objects of \textnormal{Cat}, and totally $W$-aspheric morphisms.}

Let us take up again the situation
\[
U\overset{i}{\hookrightarrow}X\overset{j}{\hookleftarrow}T,\qquad\qquad\text{with }T=\{t\}\text{, final category, }t\in\bar U\text{, i.e. }X=\bar U,
\]
and that of a cobase change $U\longrightarrow U'$, giving rise to a cocartesian square in Cat
\[
(*)\qquad\qquad
\begin{tikzcd}
U \arrow[r, hook, "i"] \arrow[d, "p"'] & X \arrow[d, "q"'] & T=\{t\} \arrow[l, hook', dotted, "j"'] \arrow[d, dotted, "\mathrm{iso}"] \\
U' \arrow[r, hook, "i'"] & X' & T'=\{t'\}. \arrow[l, hook', dotted, "j'"']
\end{tikzcd}
\]
The gluing datum $H$ between $U$ and $T$, $H:\underbrace{U^\circ\times T}_{\simeq U^\circ}\longrightarrow\mathrm{Ens}$, can be interpreted as a presheaf on $U$, $H\in\Ob U^\wedge:U^\circ\longrightarrow\mathrm{Ens}$,
\msp{page 88}
and we have
\[
U/t\simeq U/H,
\]
where $U/H$ is interpreted in terms of the canonical embedding
\[
U\hookrightarrow U^\wedge.
\]
Likewise, the gluing between $U'$ and $T'$ will be given by a presheaf
\[
H'\in\Ob U^{\prime\wedge}:U^{\prime\circ}\longrightarrow\mathrm{Ens},
\]


\endgroup
% END UNIT Der-ChXIIb-012/Der-ChXIIb-012_p016-020.tex
% BEGIN UNIT Der-ChXIIb-012/Der-ChXIIb-012_p021-025.tex
\begingroup\makeatletter%
\newcommand{\printedpage}[1]{\marginpar{\footnotesize p.~#1}}%
\newcommand{\msp}[1]{\marginpar{\footnotesize [#1]}}%
\newcommand{\cercle}[1]{\tikz[baseline=(c.base)]\node[draw,circle,inner sep=1pt](c){#1};}%
\newcommand{\defeq}{\overset{\mathrm{def}}{=}}%
\newcommand{\upi}{\underline{\pi}}%
\newcommand{\Cat}{\mathrm{Cat}}%
\newcommand{\Ob}{\mathrm{Ob}\,}%
\makeatother\let\printedpage\derpp\let\msp\dermsp%
\let\mspin\DERmspin%

\DERmark{ENV0725}


\printedpage{63}
% head: A. Grothendieck : Dérivateurs, Ch. XII, Édition des Universités de Montpellier II et Paris VII
and we will have
\[ U'/t' \simeq U'/H'. \]
\DERmark{ENV0726}

The relation between $H'$ and $H$ was made explicit in §1, as
\[ H' = p_!H, \qquad \text{where } p_! : U^{\wedge} \longrightarrow U'^{\wedge} \]
is the left adjoint of $p^* : U'^{\wedge} \longrightarrow U^{\wedge}$, given by
\[ p^*F' = F' \circ p. \]
We have
\[ p_!(H)(u') \simeq \pi_0(u'\backslash U/H). \]

\textbf{Definition 1.} \cercle{a} Let $U$ be an object of Cat, $H \in \Ob U^{\wedge}$ a presheaf on $U$. We say that $H$ is \textit{totally $W$-aspheric} if for every arrow $p : U \longrightarrow U'$ of Cat, $p_!(H)$ in $U'^{\wedge}$ is a $W$-aspheric object, i.e. $U'/p_!H$ is $W$-aspheric. \cercle{b} We say that an object $V$ of \msp{page 89} Cat is \textit{totally $W$-aspheric,} if the terminal functor $e_V$ on $V$ ($v \longmapsto \{\emptyset\}$) is totally $W$-aspheric, i.e. if for every arrow $f : V \longrightarrow V'$ in Cat, the presheaf
\[ \pi_0(f) \defeq f_!(e_V) \defeq (v' \longmapsto \pi_0(v'\backslash V)) \]
on $V'$ is $W$-aspheric, i.e. if the category
\[ \upi_0(f) \ \text{ or } \ \upi_0(V/V') \defeq V'/\pi_0(f) \]
is $W$-aspheric. \texttt{[This notion is different from the notion of total asphericity in Pursuing Stacks.]}

Let us note that the definition of the functors $p_!$ (or $f_!$) as adjoints, and their ``calculation'' in terms of the $\pi_0$ of certain categories, does not depend on the localizer $W$. Let us give a few details on this subject. We have an adjunction morphism in $V^{\wedge}$
\[ f^*f_!(F) \longleftarrow F, \]
hence
\[ V/F \longrightarrow V/f^*f_!(F) \simeq V \times_{V'} (V'/f_!(F)), \tag*{$(*)$} \]
the last isomorphism coming from the fact that we have commutativity, up to canonical isomorphism, in the diagram of functors
\[
\begin{tikzcd}[column sep=14em, row sep=large]
V'^{\wedge} \arrow[r, "f^*=(F'\longmapsto F'\circ f)"] \arrow[d, "F'\longmapsto V'/F'"'] & V^{\wedge} \arrow[d, "F\longmapsto V/F"] \\
\Cat/V' \arrow[r, "f^*_{\Cat}=(W'\longmapsto W'\times_{V'}V)"] & \Cat/V,
\end{tikzcd}
\]
\DERmark{ENV0727}

% foot: Version 7 décembre 2002
\printedpage{64}
% head: A. Grothendieck : Dérivateurs, Ch. XII, Édition des Universités de Montpellier II et Paris VII
\msp{page 90}
in other words, the operator $f^*$ on presheaves is interpreted, in geometric terms, as the operation $W' \longmapsto W' \times_{V'} V$ of base change by $V \longrightarrow V'$, on the categories over $V'$ which are fibered with discrete fibers (a property which is preserved by base change).

That said, giving an arrow $(*)$ is equivalent, by the universal property of base change, to giving an arrow (canonical in this case)
\[ f/F : V/F \longrightarrow V'/f_!(F) \]
making the square commutative
\[
\begin{tikzcd}[column sep=large, row sep=huge]
V/F \arrow[r, "\text{forgetting }F", "\substack{\text{or source}\\ \text{functor}}"'] \arrow[d, "f/F"'] & V \arrow[d, "f"] \\
V'/(F' = f_!(F)) \arrow[r, "\text{forgetting }F'"] & V'.
\end{tikzcd}
\tag*{$(*)$}
\]
In the case that interests us especially here, $F = e_V$, $V/F \longrightarrow V$ is an isomorphism, and we obtain a canonical factorization of functors \msp{page 91}
\[
\begin{tikzcd}[row sep=large]
 & V \arrow[dl, "\tilde{f}=f/e_V"'] \arrow[dd, "f"] \\
\begin{array}{r} \upi_0(V/V') \\ = V'/f_!(e_V) \end{array} \arrow[dr, "\text{can.}"'] & \\
 & V'.
\end{tikzcd}
\tag*{$(**)$}
\]
Let us make the $V'$-morphism $\tilde{f}$ explicit on the fibers,
\[ \tilde{f}_{v'} : V_{v'} \longrightarrow \upi_0(V/V')_{v'} \defeq \pi_0(v'\backslash V), \tag*{$(***)$} \]
where in the last term, we have identified the set $\pi_0(v'\backslash V)$ with the discrete category it defines. To make this arrow $(***)$ explicit, let us note that we have a canonical inclusion
\[ V_{v'} \longrightarrow v'\backslash V \]
and for every category $W$, a canonical (adjunction) morphism $W \longrightarrow \pi_0(W)$, hence a composite
\[ V_{v'} \longrightarrow v'\backslash V \longrightarrow \pi_0(v'\backslash V), \]
which is none other than $\tilde{f}_{v'}$ in $(***)$.

It should be easy to show that the factorization $(**)$ of $f$ is characterized, up to unique isomorphism, as a
% foot: Version 7 décembre 2002
\printedpage{65}
% head: A. Grothendieck : Dérivateurs, Ch. XII, Édition des Universités de Montpellier II et Paris VII
\msp{page 92}
factorization in $\Cat/V'$
\[
\begin{tikzcd}[column sep=small]
 & \tilde{V} \arrow[dr, "\pi"] & \\
V \arrow[ur, "\tilde{f}"] \arrow[rr, "f"'] & & V',
\end{tikzcd}
\]
by the following properties:
\begin{enumerate}[label={}, leftmargin=*]
\item[a)] $\pi$ is fibered with discrete fibers, i.e. a local isomorphism (i.e. $\tilde{V}$ is isomorphic in $\Cat/V'$ to a category $V'/F'$, with $F'$ in $V'^{\wedge}$). % [sic] V/F' -> V'/F', une catégorie au-dessus de V' définie par F' dans V'^∧ est de la forme V'/F'
\item[b)] $\tilde{f}_!(e_V) = e_{\tilde{V}}$, i.e. for every $\tilde{v} \in \Ob \tilde{V}$, the category $\tilde{v}\backslash V$ is 0-connected: the colocalizations of $V$ over $\tilde{V}$ are 0-connected. % [sic] pout -> pour, faute de frappe
\end{enumerate}
\DERmark{ENV0728}


Let us note here, in passing, that for \textit{every} morphism $\pi : W' \longrightarrow W$ in Cat which is fibered with discrete fibers, i.e. such that $W'$ is isomorphic in $\Cat/W$ to a category of the form $W/F$ ($F \in \Ob W^{\wedge}$), we have
\[ \pi_!(e_{W'}) = F, \qquad \text{i.e. } W' \underbrace{\ \simeq\ }_{\text{in } \Cat/W} W/\pi_!(e_{W'}). \tag*{$(*)$} \]
We must check that the functor
\[ w \longmapsto \pi_0(\underbrace{w\backslash(W/F)}_{\defeq w\backslash W/F}) \]
is canonically isomorphic to $F$. But the inclusion
\[ \underbrace{W'_w}_{= F(w) \text{ (discrete category)}} \qquad \longrightarrow w\backslash W' \]
\msp{page 93}
does indeed define as before
\[ F(w) = W'_w \longrightarrow \pi_0(w\backslash W'), \]
we must prove that this is bijective. Now, $W'$ being $W_\infty$-smooth over $W$, we know that the inclusion $W'_w \hookrightarrow w\backslash W'$ is in $W_\infty$, a fortiori induces an isomorphism on the $\pi_0$.

This established, if we have a factorization of $f : V \longrightarrow V'$ as $f = \pi\tilde{f}$ having the properties a) b) of the preceding page, we already see that we have
\[ f_!(e_V) \simeq \pi_!(\, \underbrace{\tilde{f}_!(e_V)}_{\simeq e_{\tilde{V}} \text{ by b)}} \,) \simeq \pi_!(e_{\tilde{V}}), \]
which already determines $\tilde{V}$ up to canonical isomorphism (thanks to a) as
\[ \tilde{V} \simeq V'/f_!(e_V) \defeq \upi_0(V/V'), \] % [sic] \underline{\pi}_0(V'/V) -> \underline{\pi}_0(V/V'), notation employée partout ailleurs (Définition 1, (**))
\DERmark{ENV0729}

% foot: Version 7 décembre 2002
\printedpage{66}
% head: A. Grothendieck : Dérivateurs, Ch. XII, Édition des Universités de Montpellier II et Paris VII
and it should not be very hard to convince oneself that $\tilde{f}$ is none other than the canonical arrow of page 91.

\textbf{Proposition 8.} \textit{For an object $H$ of $U^{\wedge}$ to be totally $W$-aspheric, it suffices that the object $U/H$ of} Cat \textit{be totally $W$-aspheric.}

\msp{page 94}
Let us indeed suppose this latter condition satisfied and prove that $H$ is totally $W$-aspheric, i.e. that for every $p : U \longrightarrow U'$, $U'/p_!(H)$ is $W$-aspheric.
\[
\begin{tikzcd}[row sep=large]
V \defeq U/H \arrow[r, "\alpha"] \arrow[dr, "f"'] & U \arrow[d, "p"] \\
 & U'
\end{tikzcd}
\]
Let $f$ be the composite functor in the diagram opposite \texttt{[above]}, we have $H \simeq \alpha_!(e_V)$ (by $(*)$ p. 92), hence
\[ p_!(H) \simeq p_!\alpha_!(e_V) \simeq (p\alpha)_!(e_V) \simeq f_!(e_V), \]
which is $W$-aspheric by the hypothesis on $V = U/H$.

Does the converse of \texttt{[the]} prop. 8 hold, i.e. if $H$ is totally $W$-aspheric, is the same true of $V = U/H$? The difficulty comes from the fact that this means that for \textit{every} arrow $f : V \longrightarrow V'$ in Cat with source $V$, we have $f_!(e_V)$ $W$-aspheric, whereas the hypothesis is equivalent to postulating it only for arrows that factor through $U$. Let us note that if we do not make this hypothesis on $f$, but suppose instead that $f$ fits into a commutative square \msp{page 95}
\[
\begin{tikzcd}[row sep=large]
V \arrow[r, "\alpha"] \arrow[d, "f"'] & U \arrow[d, "p"] \\
V' \arrow[r, "\alpha'"] & U',
\end{tikzcd}
\]
with $\alpha'$ also a local isomorphism, then we can conclude that $\upi_0(f) = \upi_0(V/V')$ is indeed aspheric. Indeed, we have a canonical isomorphism in $\Cat/U'$,
\[ \upi_0(\alpha'f) \simeq \upi_0(f), \qquad \text{i.e. } \ \upi_0(V/V') \simeq \upi_0(V/U'), \]
from the fact e.g. that the canonical factorization of $f$
% foot: Version 7 décembre 2002
\printedpage{67}
% head: A. Grothendieck : Dérivateurs, Ch. XII, Édition des Universités de Montpellier II et Paris VII
\[
\begin{tikzcd}[row sep=normal]
 & V \arrow[dl, "\tilde{f}"'] \arrow[dd, "f"] & \\
\upi_0(V/V') \arrow[dr, "\pi"'] & & \\
 & V' \arrow[dr, dotted, "\alpha'"'] & \\
 & & U'
\end{tikzcd}
\]
provides a factorization of the same type for $\alpha'f$, as
\[ \alpha'f = (\alpha'\pi) \circ \tilde{f}. \]
That said, we can apply the hypothesis to $p$ and we do indeed find the $W$-asphericity of $\upi_0(f) \simeq \upi_0(\alpha'f) = \upi_0(p\alpha)$.

\msp{page 96}
The natural idea, for a converse to \texttt{[the]} prop. 8, would be to take for $U'$ the amalgamated sum of $U$ and $V'$ over $V$, which will work provided we know that $\alpha'$ really is a local isomorphism just like $\alpha$. This is indeed the case e.g. if $V \longrightarrow U$ is an open immersion, i.e. $H$ a subpresheaf of the terminal presheaf. But I fear that this is not true in general - it looks a little too good to be true! This is indeed already false in the case
\[
\begin{tikzcd}[row sep=huge]
V = E_G \arrow[r, "\alpha\ \text{can}"] \arrow[d, "f"'] & B_G = U \\
V' = B_{G'}, &
\end{tikzcd}
\]
in the case e.g. where $G = \{\pm 1\}$, the group with two elements, and where the element $\alpha$ of $G'$, the image under $f$ of one of the two nonidentity arrows of $E_G$ is $\neq 1$ \texttt{[rather with square} $\neq 1$\texttt{]}. If $H'$ is the normal subgroup generated by $\alpha$ \texttt{[rather by} $\alpha^2$\texttt{]} in $G'$, the amalgamated sum is identified with $B_{G'/H'}$, but $B_{G'} \longrightarrow B_{G'/H'}$ is \textit{not} a local isomorphism (for $(B_{G'})/t_0 \longrightarrow (B_{G'/H'})/t'_0$ is not bijective on objects, being identified with $G' \overset{q}{\longrightarrow} G'/H' \ \ldots$).

But on the other hand the condition of total $W$-asphericity on an object $H$ of $U^{\wedge}$ already seems so demanding, that I can hardly see myself con\msp{page 97}structing an example where it is satisfied, without $U/H$ itself being totally $W$-aspheric.

The existence of totally $W$-aspheric objects in Cat is not so obvious a priori. That there are a good bunch of them is guaranteed by the
% foot: Version 7 décembre 2002


\endgroup
% END UNIT Der-ChXIIb-012/Der-ChXIIb-012_p021-025.tex
% BEGIN UNIT Der-ChXIIb-012/Der-ChXIIb-012_p026-030.tex
\begingroup\makeatletter%
\newcommand{\printedpage}[1]{\marginpar{\footnotesize p.~#1}}%
\newcommand{\msp}[1]{\marginpar{\footnotesize [#1]}}%
\newcommand{\upi}{\underline{\pi}}%
\newcommand{\uCh}{\underline{\mathrm{Ch}}}%
\newcommand{\Ob}{\operatorname{Ob}}%
\newcommand{\Hom}{\operatorname{Hom}}%
\newcommand{\bDelta}{\boldsymbol{\Delta}}%
\makeatother\let\printedpage\derpp\let\msp\dermsp%
\let\mspin\DERmspin%


\setcounter{footnote}{65}

\printedpage{68}
% head: A. Grothendieck : Dérivateurs, Ch. XII, Édition des Universités de Montpellier II et Paris VII

\noindent\textbf{Proposition 9.}

\begin{enumerate}[label=\textit{\alph*)}]
\item \textit{Every object $V$ of} Cat \textit{having a final object is totally $W$-aspheric - hence every presheaf of sets $H$} representable \textit{on a category $U$ (representability meaning precisely that $U/H$ has a final object) is totally $W$-aspheric.}

\item (\footnote{!! Probably \textit{false} as it stands, cf. correction on the next page.}) \textit{(Suppose} W(7 bis) \textit{is satisfied.) Let $V$ be in} Cat, \textit{$V'$, $V''$ two opens with union $V$. If $V'$, $V''$, $V'\cap V''$ are totally $W$-aspheric, the same holds for $V$.}
\end{enumerate}

\noindent\textsc{Proof} of a). Take $V$ with a final object and $f : V \longrightarrow V'$, whence the factorization
\[
\begin{tikzcd}[column sep=small, row sep=small]
 & \tilde{V} = \upi_0(V/V') \arrow[dr, "\pi"'] & \\
V \arrow[ur, "\tilde{f}"] \arrow[rr, "f"'] & & V'
\end{tikzcd}
\]
of $f$. One proves without difficulty (\footnote{Here is a simple proof: For every $V$ in Cat we have $e_V = \varinjlim_{v \text{ in } V} v$ (direct limit of the inclusion functor of $V$ into $V^\wedge$), for $\forall\, x$ in $V$
\[
\varinjlim_{v \text{ in } V} \Hom(x,v) \simeq \pi_0(x\backslash V) = e.
\]
Since $f_!$ commutes with $\varinjlim$ , we therefore have $f_!(e_V) = \varinjlim_{v \text{ in } V} f(v)$. If $v_0$ is a final object of $V$, this shows that $f_!(e_V) = f(v_0)$, q.e.d.}) (a matter of cofinality) the following:

\noindent\textbf{Corollary.} \textit{The image under $\tilde{f}$ of a final object of $V$ is a final object of $\tilde{V} = \upi_0(V/V')$.}

Since an object of Cat with a final object is $W$-aspheric by W(3), this establishes a).

Let us prove b). I feel like saying that if we have $f : V \longrightarrow W$ in Cat, then $\upi_0(V/W)$ is a union of two opens isomorphic to $\upi_0(V'/W)$, $\upi_0(V''/W)$, whose intersection is isomorphic to $\upi_0(V'\cap V''/W)$. This would mean neither more nor
\msp{page 98}
less than this: $f_!$ transforms the monomorphisms $V' \longrightarrow V$, $V'' \longrightarrow V$ (interpreted as arrows in $V^\wedge$, with target $e_V$) into monomorphisms, and the fiber product $V'\cap V'' = V' \times_V V''$ into the corresponding fiber product in $W^\wedge$. But the functor $f_!$ commutes with arbitrary $\varinjlim$ (being a left adjoint), and I see no a priori reason why it should commute (say) with fiber products (which would also imply that it transforms monomorphisms into likewise). If I remember correctly, this is ultrafalse in general. So I now very much doubt that statement b) is true, without an additional hypothesis. To ensure monomorphy, one thinks of requiring the existence of \textit{retractions} of $V$ onto $V'$, $V''$ and of $V'$ or $V''$ onto $V'\cap V''$.

I have just checked this, and give the
% foot: Version 7 décembre 2002
\printedpage{69}
% head: A. Grothendieck : Dérivateurs, Ch. XII, Édition des Universités de Montpellier II et Paris VII
\textbf{Corollary 1} (correction of \texttt{[the]} prop. 9, part b) (\footnote{NB There is a slightly more general statement, starting from a category $U$ and a diagram
$\begin{tikzcd}[ampersand replacement=\&, row sep=small]
V'\cap V'' \arrow[r] \arrow[d] \& V' \arrow[d] \\
V'' \arrow[r] \& V
\end{tikzcd}$
\textit{in} $U^\wedge$: if $V'$, $V''$, $V'\cap V''$ are totally $W$-aspheric, the same holds for $V$.}): \textit{The statement is valid if one assumes that there exists a} retraction \textit{$r'_0$ of $V''$ onto $V'\cap V''$, and similarly a retraction $r''_0$ of $V'$ onto $V'\cap V''$.} (\footnote{Corrected corollary also false - the retractions are not arrows of $V^\wedge$, and do not pass into $p_!$.}) % [resolved: see DER_XIIb-a_corrections.csv]
\[
\begin{tikzcd}
V'\cap V'' \arrow[r, "i''_0"'] \arrow[d, "i'_0"] & V' \arrow[d, "i'"] \arrow[l, dotted, bend right=50, "r''_0"'] \\
V'' \arrow[r, "i''"] \arrow[u, dotted, bend left=50, "r'_0"] & V \arrow[u, dotted, bend right=50, "r'"'] \arrow[l, dotted, bend left=50, "r''"]
\end{tikzcd}
\]

(\textbf{NB} We will see later that the existence of \textit{one} of these retractions suffices to ensure the conclusion.)

\textsc{Proof}. Since $V$ is the amalgamated sum of $V'$ and $V''$, there exists a retraction $r'$ of $V$ onto $V'$ (given by $\mathrm{id}_{V'}$ on $V'$, by $i''_0r'_0$ on $V''$) which induces $r'_0$ on $V''$, and similarly there exists
\msp{page 99}
a retraction $r''$ of $V$ onto $V''$, inducing $r''_0$ on $V'$. The existence of these retractions already ensures that for every functor $f : V \longrightarrow W$, the transform of the diagram
\[
\begin{tikzcd}[row sep=large]
V'\cap V'' \arrow[r] \arrow[d] & V' \arrow[d] \\
V'' \arrow[r] & V
\end{tikzcd}
\]
(interpreted as a diagram in $V^\wedge$) consists of monomorphisms (\footnote{\textbf{false,} already in the case where $V = \Psi$, $V'$, $V''$ its two canonical opens isomorphic to $\bDelta^1$, and $V'{\cap}V'' = e$.}) % [sic] intérprété -> interprété, misspelt accent
\DERmark{ENV0730}

\[
\begin{tikzcd}[row sep=large]
f_!(V'\cap V'') \arrow[r, hook] \arrow[d, hook] & f_!(V') \arrow[d, hook] \\
f_!(V'') \arrow[r, hook] & f_!(V),
\end{tikzcd}
\]
and moreover that $f_!(V') \cap f_!(V'')$ (which contains $f_!(V'\cap V'')$) is in fact \textit{equal to it.} It follows that we have a diagram both cartesian and cocartesian of open immersions in Cat
% foot: Version 7 décembre 2002
\printedpage{70}
% head: A. Grothendieck : Dérivateurs, Ch. XII, Édition des Universités de Montpellier II et Paris VII
\[
\begin{tikzcd}[row sep=large]
W/f_!(V'\cap V'') \arrow[r, hook] \arrow[d, hook] & W/f_!(V') \arrow[d, hook] \\
W/f_!(V'') \arrow[r, hook] & W/\Bigl(f_!(V) \overset{\mathrm{def}}{=} f_!(e_V)\Bigr),
\end{tikzcd}
\]
\DERmark{ENV0731}

which is also identified with
\[
\begin{tikzcd}[row sep=large]
\upi_0(V'\cap V''/W) \arrow[r, hook] \arrow[d, hook] & \upi_0(V'/W) \arrow[d, hook] \\
\upi_0(V''/W) \arrow[r, hook] & \underline{\pi_0}(V/W).
\end{tikzcd}
\]

\msp{page 100}
The hypothesis on $V'$, $V''$, $V'\cap V''$ implies that the three vertices of the square other than $\upi_0(V/W)$ are $W$-aspheric, hence by W(7 bis) the same holds for the latter, q.e.d.

\textbf{Corollary 2} (\footnote{More generally, if $I$ is a category \textit{ordered and finite,} having an initial object, then}). \textit{Assuming} W(7 bis)\textit{, the category} $\Psi = \begin{tikzcd}[ampersand replacement=\&, row sep=small, column sep=small, baseline=(B.base)]
 \& a \\
|[alias=B]| 0 \arrow[ur] \arrow[dr] \& \\
 \& b
\end{tikzcd}$ \textit{(or $0 \longleftarrow 1 \longrightarrow 2$) is totally $W$-aspheric} (\footnote{\textbf{False,} $\Psi$ is \textit{not} totally $W$-aspheric.}).

One applies corollary 1 to $\Psi$ and to its two open subcategories, $\Psi/a$, $\Psi/b$ (isomorphic to $\bDelta^1$), and to their intersection $\{0\} \simeq \bDelta^0 = e$. Since these three categories have a final object (unlike $\Psi$ itself, which has only an initial object \dots), they are totally $W$-aspheric. The corollary applies, the existence of the retractions $r'_0$, $r''_0$ being obvious here.

It is especially part b) (corrected) of proposition 9 (ms.97--100) which ensures the existence of totally $W$-aspheric categories other than those having a final object (\footnote{I now conjecture that these are the only ones! (\footnotemark)}\footnotetext{False, there are the filtering categories (if $W$ satisfies the axiom of limits).}). It is true that conversely one should still make sure that a $W$-aspheric object $X$ of Cat is not \textit{always} totally $W$-aspheric! The counterexample is not so obvious after all -
\DERmark{ENV0732}

\msp{page 101}
I spent hours before finding one and getting a little clarity! It suffices, in short, to exhibit an example of an arrow in Cat
\[
f : X \longrightarrow Y
\]
such that
% foot: Version 7 décembre 2002
\printedpage{71}
% head: A. Grothendieck : Dérivateurs, Ch. XII, Édition des Universités de Montpellier II et Paris VII
\begin{enumerate}[label=\alph*)]
\item $\underline{\pi_0}(f) \overset{\mathrm{def}}{=} \upi_0(X/Y) \xrightarrow{\ \sim\ } Y$, i.e. the $\pi_0(y\backslash X)$ \texttt{[rather the} $y\backslash X$\texttt{]} are 0-connected (for $y \in \Ob Y$),
\end{enumerate}
\DERmark{ENV0733}


and such that moreover

\begin{enumerate}[label=\alph*), start=2]
\item $X$ is $W$-aspheric, and $Y$ not.
\end{enumerate}

% a vertical bar is printed in the left margin beside the next paragraph (lines "La difficulté ..." to "... le foncteur")
The difficulty is the following, when one looks for ``stupid'' counterexamples: it is that hypothesis a) already implies, when $X$ is $W_\infty$-aspheric (or only 1-connected) that $Y$ is 1-connected: the functor
\[
Y' \longmapsto X' = X \times_Y Y'
\]
from the category of étale coverings $Y'$ of $Y$ to that of étale coverings $X'$ of $X$ is fully faithful (immediate verification), whence the stated result. So a counterexample must call upon the existence of a nonzero $\pi_i(Y,y)$, with $i \geq 2$. This implies
\msp{page 102}
that there is no counterexample with categories of the form $Y = Z/F$ (deduced from a functor $q : X \longrightarrow Z$ by taking $F = q_!(e_X) = \pi_0(X/Z)$), where $Z$ is a category such as $\bDelta^1$ (?), $\Phi$, $\Psi$, or $\begin{tikzcd}[ampersand replacement=\&, sep=small, baseline=(M.base)]
 \& {} \& \\
{} \arrow[ur] \arrow[dr] \& |[alias=M]| \phantom{x} \& {} \arrow[ul] \arrow[dl] \\
 \& {} \&
\end{tikzcd}$ (having the homotopy type of the circle) etc., for the connected components of such categories $Y$ are $K(\pi,1)$'s - the higher $\pi_i$ ($i \geq 2$) are zero.

This difficulty seen, the counterexample becomes obvious. I start from a 1-connected object \textit{whatsoever} $Y$ of $X$ \texttt{[rather of} Cat\texttt{]}, and I take $y_0 \in \Ob Y$,
\[
X = \uCh(Y; y_0, -) \xrightarrow{\ f\ } Y.
\]
We know that $X$ is $W_\infty$-aspheric, and that $f$ is $W_\infty$-perfect, hence the inclusions
\[
X_y \longrightarrow y\backslash X
\]
are $W_\infty$-equivalences. Hypothesis a) that the $y\backslash X$ be 0-connected is therefore equivalent to the same hypothesis on the fibers $X_y$, or
\msp{page 103}
again on a single fiber ($f$ being perfect, hence a $W_\infty$-fibration, and $Y$ 0-connected), namely
\[
X_{y_0} \simeq \Omega(Y, y_0).
\]
Its 0-connectedness is equivalent to the hypothesis of 1-connectedness of $Y$ (so far we have used only that of 0-connectedness). So for a counterexample, it suffices to take $Y$ 1-connected and not $\infty$-connected. The simplest case is the one where $Y$ has the homotopy type of the sphere $S^2$, realized e.g. by the ordered category shown opposite \texttt{[below]} with 6 vertices.
% foot: Version 7 décembre 2002

\printedpage{72}
% head: A. Grothendieck : Dérivateurs, Ch. XII, Édition des Universités de Montpellier II et Paris VII
\[
\begin{tikzcd}[row sep=small]
e_2^+ \bullet & \bullet\, e_2^- \\
e_1^+ \bullet \arrow[ur] \arrow[u] & \bullet\, e_1^- \arrow[ul] \arrow[u] \\
e_0^+ \bullet \arrow[ur] \arrow[u] & \bullet\, e_0^- \arrow[ul] \arrow[u]
\end{tikzcd}
\]
\DERmark{ENV0734}\DERmark{ENV0735}

\texttt{[moreover,} $e_0^+ \leq e_1^+ \leq e_2^+$\texttt{,} $e_0^- \leq e_1^- \leq e_2^-$\texttt{]} (\footnote{It is in short this example (up to topological homotopy) that I realize (without at first noticing it in the Appendix p. 10-12, but with finite categories).}). I summarize the result obtained.

\textbf{Proposition 10.} \textit{Let $Y$ be a small category which is} 1\textit{-connected ($y_0 \in \Ob Y$). Then} $\uCh(Y; y_0, -)$ \textit{is $W_\infty$-aspheric, but is not totally $W_\infty$-aspheric if $Y$ itself is not $\infty$-connected, i.e. $W_\infty$-aspheric.} (\footnote{But a more striking example is $\Psi = \begin{tikzcd}[ampersand replacement=\&, row sep=small, column sep=small, baseline=(A.base)]
 \& b \\
|[alias=A]| a \arrow[ur] \arrow[dr] \& \\
 \& c
\end{tikzcd}$ , which is $W_\infty$-aspheric, but not totally!}).

One may ask whether $Y$ $W_\infty$-aspheric implies that $\uCh(Y; y_0, -)$ is totally $W_\infty$-aspheric. I very much doubt it!

\msp{page 104}
Let us now return to the situation at the beginning of the section, diagram $(*)$, p.~87,
\[
\begin{tikzcd}
U \arrow[r, hook, "i"] \arrow[d, "p"'] & X \arrow[d, "q"'] & T = \{t\} \arrow[l, hook', dotted, "j"'] \arrow[d, dotted, "\mathrm{iso}"] \\
U' \arrow[r, hook, "i'"] & X' & T' = \{t'\}. \arrow[l, hook', dotted, "j'"']
\end{tikzcd}
\tag*{$(*)$}
\]
We can now give an example where $i$ is in $W$ (and even $W$-aspheric) without $i'$ being so:
\[
i \in W, \qquad i' \notin W. \tag*{$(**)$}
\]
One can even assume that $t$ is a final object of $X$ (which means that the object $H$ of $U^\wedge$ is isomorphic to $e_U$, the final object), and $t'$ a final object of $X'$ (i.e. $p_!(e_U) = e_{U'}$, i.e. the $u'\backslash U$ are 0-aspheric). In this case, since $X$, $X'$ are $W$-aspheric, the relations $(**)$ are equivalent to
\[
U \text{ is } W\text{-aspheric, and } U' \text{ not,} \tag*{$(***)$}
\]
which is precisely the situation we have just constructed (p. 101-103).

One would also like to find an example where in $(*)$ one has $p \in W$, $q \notin W$.
% foot: Version 7 décembre 2002


\endgroup
% END UNIT Der-ChXIIb-012/Der-ChXIIb-012_p026-030.tex
% BEGIN UNIT Der-ChXIIb-012/Der-ChXIIb-012_p031-035.tex
\begingroup\makeatletter%
\newcommand{\printedpage}[1]{\marginpar{\footnotesize p.~#1}}%
\newcommand{\msp}[1]{\marginpar{\footnotesize [#1]}}%
\newcommand{\hra}{\lhook\joinrel\longrightarrow}%
\newcommand{\Cat}{\mathrm{Cat}}%
\makeatother\let\printedpage\derpp\let\msp\dermsp%
\let\mspin\DERmspin%


\setcounter{footnote}{76}

\printedpage{73}
% head: A. Grothendieck : Dérivateurs, Ch. XII, Édition des Universités de Montpellier II et Paris VII
% foot: Version 7 décembre 2002

\msp{page 105}

I will again take $(X,U,\{t\})$ such that $t$ is a final object of $X$, but if we want $q\notin W$, we cannot suppose $t'$ a final element of $X'$. Instead of this, I am going to suppose

\begin{enumerate}[label=\alph*)]
\item $U$, $U'$ (and also $X$, having final object $t$) $W$-aspheric, hence $i\in W$ (and even $i$ $W$-aspheric), $p\in W$, but
\item $X'$ not $W$-aspheric, hence $i',q\notin W$.
\end{enumerate}

For this, let us note that
\[ U'/t'\simeq U'/p_!(e_U)\overset{\mathrm{def}}{=}\underline{\pi}_0(U/U'). \]
\DERmark{ENV0736}

The corollary to proposition 7 (cf. p. 85), not yet proved, tells us that $i'\in W$ (i.e. here, $X'$ $W$-aspheric) if and only if $i'$ is $W$-aspheric, i.e. $U'/t'$ is $W$-aspheric. We shall admit it here. Everything thus comes down to finding an arrow
\[ p:U\longrightarrow U' \]
in Cat with the following properties.

\begin{enumerate}[label=\alph*)]
\item $U$ and $U'$ are $W$-aspheric,
\item $\underline{\pi}_0(U/U')$ is not $W$-aspheric.
\end{enumerate}

We shall take up the preceding example again
\[ f:X\longrightarrow Y \]
(p. 101-103), taking
\[ U=X,\qquad \underline{\pi}_0(U/U')=Y, \]

\msp{page 106}

so it all comes down, starting from the given example $X\xrightarrow{f}Y$, to constructing an arrow
\[ \pi:Y\longrightarrow U' \tag*{$(*)$} \]
in Cat which satisfies the conditions

\begin{enumerate}[label=\alph*)]
\item $\pi$ is a local isomorphism,
\item $U'$ is $W$-aspheric.
\end{enumerate}

We shall then take
\[ p=\pi f. \]
But from the way $X$ has been constructed from a 1-connected and not $\infty$-connected $Y$, otherwise arbitrary, we are finally reduced to finding $(*)$ satisfying a), b), and moreover

\begin{enumerate}[label=\alph*), start=3]
\item \printedpage{74}
% head: A. Grothendieck : Dérivateurs, Ch. XII, Édition des Universités de Montpellier II et Paris VII
% foot: Version 7 décembre 2002
$Y$ is 1-connected, but not $\infty$-connected.
\end{enumerate}

(NB We now suppose $W=W_\infty$.) It amounts to the same thing to find

% printed with a vertical rule at the left of this block
\begin{quote}
$U'$ an \textit{aspheric} object of Cat, $F\in\operatorname{Ob}U'^\wedge$ such that $U'/F$ is 1-connected and not $\infty$-connected.
\end{quote}

One thinks for example of simply taking an open subcategory $U'_0$ of $U'$, which is 1-connected and not $\infty$-connected. Let us start with such a category $U'_0$ - for example a homotopy realization in Cat, of the sphere $S^2$. We immerse it as an open in an aspheric $U'$, by adding a final object to it. We have won.

\msp{page 107}

Let us make a statement in due form.

\textbf{Proposition 11.} \textit{One can find a situation} $(*)$ (p. 104 \texttt{[rather 104]}) \textit{with}
\DERmark{ENV0737}


\begin{enumerate}[label=\textit{\alph*)}]
\item \textit{$t$ a final object of $X$.}
\item \textit{$U$, $U'$ (and $X$) $W_\infty$-aspheric, hence $i,p\in W_\infty$ (and even $i$ $W_\infty$-aspheric).}
\item \textit{$X'$ not $W_\infty$-aspheric, hence $i',q\notin W_\infty$.}
\end{enumerate}

For this it suffices to start with a 1-connected object $Y$ of Cat, to adjoin to it a final object $u'_0$, which gives $U'$, to take $U=\underline{\mathrm{Ch}}(Y;y_0,-)$, $p:U\longrightarrow U'$ the composite $U\xrightarrow{\text{target}}Y\underset{\text{can.}}{\overset{\text{incl.}}{\hra}}U'$, and $X$ deduced from $U$ by adding a final object $t$ to it. We then have $U'/t'\simeq Y$, whence
\[ X'\ W_\infty\text{-aspheric (i.e. }i'\in W_\infty)\Longleftrightarrow Y\ W_\infty\text{-aspheric.} \]
We have the desired counterexample by taking $Y$ (1-connected but) not $\infty$-connected.

\subsection*{2.4\quad Application to open Dwyerian immersions.}

Let us give all the same the \textit{positive} statement we were after. % [sic] l'enoncé -> l'énoncé, missing accent

\textbf{Theorem 3.} \textit{Let $(X,U,T)$ be a triple $\in\mathcal{T}$ ($U$ open in $X$ in} Cat\textit{, $T=X\backslash U$), and for every arrow $p:U\longrightarrow U'$ an arrow of} Cat\textit{, the cocartesian square in} Cat

\[
\begin{tikzcd}
U \arrow[r, hook, "i"] \arrow[d, "p"'] & X \arrow[d, "q"'] & T \arrow[l, dotted, hook'] \arrow[d, dotted, "\text{iso}"] \\
U' \arrow[r, hook, "i'"] & X' & T'. \arrow[l, dotted, hook']
\end{tikzcd}
\]

\printedpage{75}
% head: A. Grothendieck : Dérivateurs, Ch. XII, Édition des Universités de Montpellier II et Paris VII
% foot: Version 7 décembre 2002

\msp{page 108}

\textit{Suppose that for every $t\in\dot{T}=\dot{U}$, the presheaf of sets}
\[ i^*(t)=(u\longmapsto\operatorname{Hom}(u,t))\in U^\wedge \]
\textit{is totally $W$-aspheric (def. 1, p. 88)} (\footnote{\textbf{NB} One says under these conditions that $i_0:U\longrightarrow\bar{U}$ is \textit{strictly $W$-aspheric} (cf. def. 2 next page).}) \textit{- e.g. that the category $U/t$ is totally $W$-aspheric (prop. 8, page 93). (NB the case where the $U/t$ have a final object, cf. prop. 9 a), i.e. the one where the inclusion $U\longrightarrow\bar{U}$ admits a right adjoint, is the one envisaged by Dwyer, that of open Dwyer immersions} (\footnote{I know of no other case where th. 3 applies - for lack of knowing of a totally $W$-aspheric category which does not have a final object, nor of an object $F$ of a category $U^\wedge$ which is totally $W$-aspheric without being representable, i.e. $U/F$ has a final object. Now, I know that \textit{filtering} categories are totally $W$-aspheric if $W$ satisfies the axiom of direct limits.})\textit{.) Then for every $p:U\longrightarrow U'$, the open immersion}
\[ i'_0:U'\longrightarrow\overline{U'} \]
\textit{is $W$-aspheric (and even, just like $i_0$, strictly $W$-aspheric), and a fortiori it is} \texttt{[in]} \textit{$W$. Consequently (prop. 6} \texttt{[p. 81]}\textit{) we have the following.}

\[
\begin{aligned}
&\text{\textcircled{a}}\qquad && i\in W && \Longleftrightarrow\quad && i'\in W\quad && \textit{assuming } \mathrm{W}(7\text{ bis}).\\
&\text{\textcircled{b}}\qquad && p\in W && \Longleftrightarrow\quad && q\in W\quad && \textit{assuming } \mathrm{W}(7).\ (\footnotemark)
\end{aligned}
\]
\footnotetext{Proof page 112.}

\textbf{Corollary 1.} \textit{Assume} W(7)\textit{. If $i:U\hra X$ is an open immersion such that $i_0:U\hra\bar{U}$ is totally} \texttt{[strictly suffices]} \textit{$W$-aspheric, then $i:U\longrightarrow X$ is a $W$-cofibration (``ultrastrong''), by which I mean this : for every arrow $U'\xrightarrow{f}U''$ in $U\backslash$}Cat\textit{,}

\msp{page 109}

\textit{such that $f\in W$, the arrow $g:X'\longrightarrow X''$ which is deduced from it by cobase change from the object $X$ of $U\backslash$}Cat\textit{, is also in $W$.}

\textsc{Proof of the corollary} : we apply theorem 3 to the open immersion $U'\overset{i'}{\hra}X'$ deduced from $U\overset{i}{\hra}X$ by the cobase change $U\xrightarrow{p}U'$, noting in addition that the hypothesis of theorem 3 on $i$ is preserved in passing from $i$ to $i'$.
\DERmark{ENV0738}


We must all the same introduce adequate terminology :

\textbf{Definition 2 :}

\begin{enumerate}[label=\alph*)]
\item An arrow $f:X\longrightarrow Y$ is said to be \textit{strictly $W$-aspheric} if for every $y\in\operatorname{Ob}Y$, $f^*(y)=(x\longmapsto\operatorname{Hom}_Y(f(x),y))$ is a totally $W$-aspheric object of $X^\wedge$ (cf. def. 1, page 88). It is said to be \textit{totally $W$-aspheric} if for every $y\in\operatorname{Ob}Y$, $X/y$ ($\simeq X/f^*(y)$) is a totally $W$-aspheric object of Cat.

\item \printedpage{76}
% head: A. Grothendieck : Dérivateurs, Ch. XII, Édition des Universités de Montpellier II et Paris VII
% foot: Version 7 décembre 2002
An open immersion $i:U\longrightarrow X$ is said to be \textit{$W$-Dwyerian} (resp. \textit{totally} \texttt{[$W$-]} \textit{Dwyerian}) if the induced immersion $i_0:U\longrightarrow\bar{U}$ is strictly (resp. totally) $W$-aspheric. A closed immersion $j:T\longrightarrow X$ is said to be \textit{$W$-Dwyerian} (resp. \textit{totally $W$-Dwyerian}) if and only if $j^\circ:T^\circ\longrightarrow X^\circ$ is an open $W$-Dwyerian immersion (resp totally $W$-Dwyerian).

\item An arrow $i:U\longrightarrow X$ of Cat (not necessarily an immersion)

\msp{page 110}

is called a \textit{$W$-cofibration,} if for every arrow $f:U'\longrightarrow U''$ in $U\backslash$Cat, the arrow $g:X'\longrightarrow X''$ which is deduced from it by cobase change by $X$ (object of $U\backslash$Cat) is in $W$ when $f$ is.
\end{enumerate}

\textbf{NB} This notion of $W$-cofibration is stable under composition and under cobase change, and the isomorphisms of Cat are $W$-cofibrations. If $f$ is a $W$-cofibration, $f^\circ$ is one and conversely.

(The corollary is stated by saying that an open or closed immersion which is $W$-Dwyerian is a $W$-cofibration.)

\begin{enumerate}[label=\alph*), start=4]
\item An arrow $i:U\longrightarrow X$ in Cat is called a \textit{trivial $W$-cofibration}, or a \textit{couniversal $W$-equivalence,} if for every $U'$ in $U\backslash$Cat, i.e. every arrow $p:U\longrightarrow U'$ in Cat, the arrow deduced from $i$,
\[ i':U'\longrightarrow X'=U'\amalg_U X\quad\texttt{[rather }U'\amalg_U X\texttt{]}\ \ , \]
\DERmark{ENV0739}

is in $W$. (NB It is immediate that this indeed implies that $i$ is a $W$-cofibration.)
\end{enumerate}

(This notion of trivial $W$-cofibration is also stable under composition and under cobase change, and the isomorphisms in Cat are such.)

\msp{page 111}

\textbf{Corollary 2.} (Assuming W(7).) \textit{Let $i:U\longrightarrow X$ be an open $W$-Dwyerian immersion. For $i$ to be a trivial $W$-cofibration, it is necessary and sufficient that $i\in W$, or what amounts to the same thing} (cf. page 82), \textit{that $\dot{T}\hra T$ be in $W$.}

This is obviously necessary, and sufficiency follows from theorem 3 a).

\textbf{Open questions.}

\begin{enumerate}[label=\alph*)]
\item I do not know whether an open immersion $i$ which is a $W$-cofibration, or even a trivial $W$-cofibration, is necessarily $W$-Dwyerian (\footnote{and perhaps even a Dwyer immersion?}). But by the corollary to proposition 7 (p. 85), I know that this is the case if we suppose $T$ isomorphic to the final category $e$. At least this follows : For not only $i:U\longrightarrow X$, but also all the open immersions $i_1:U\longrightarrow X_1$ which are deduced from it by the procedure of prop. 7 b) by means of arbitrary arrows $T_1\longrightarrow T$ with target $T$ in Cat, to be $W$-cofibrations, it is necessary and sufficient that $i$ be Dwyerian. % [sic] toute au moins -> tout au moins, agreement error

\item \printedpage{77}
% head: A. Grothendieck : Dérivateurs, Ch. XII, Édition des Universités de Montpellier II et Paris VII
% foot: Version 7 décembre 2002
Is there a good notion of a Dwyerian immersion, applicable to all immersions

\msp{page 112}

and including the case of open immersions and closed immersions?
\end{enumerate}

\textsc{Proof} of theorem 3 (p. 107). It suffices to prove that if $i$ is $W$-Dwyerian (i.e. $i_0$ strictly $W$-aspheric), then $i'$ is so as well - for it will follow that $i'_0$ is $W$-aspheric and a fortiori $\in W$. We must therefore verify that for every $t'\in\dot{U}'=\dot{T}$, the element $i'^*(t')=(u'\longmapsto\operatorname{Hom}_{X'}(u',t'))$ of $U'^\wedge$ is totally $W$-aspheric. Now $t'$ comes from a unique element $t$ of $\dot{T}$, and we know that we have % [sic] asphérique). alors -> asphérique), alors, period before lower-case word
\DERmark{ENV0740}\DERmark{ENV0741}\DERmark{ENV0742}


\[ i'^*(t')=p_!(i^*(t))\quad\texttt{[rather }i'^*(t')=p_!(i^*(t))\texttt{]},\qquad\text{where }\ p_!:U^\wedge\longrightarrow U'^\wedge. \]
\DERmark{ENV0743}


By hypothesis, $F=i^*(t)$ is a totally $W$-aspheric element of $U^\wedge$, which means that for every arrow $p:U\longrightarrow U'$ in Cat with source $U$, $p_!(F)$ is a $W$-aspheric element of $U'^\wedge$. In fact, we have the immediate lemma (by transitivity $(p'p)_!\simeq p'_!p_!$). % [sic] souce -> source, missing letter
\DERmark{ENV0744}


\textbf{Lemma.} \textit{Let $U$ be in} Cat\textit{, and $F$ in $U^\wedge$. If $F$ is totally $W$-aspheric (resp. if $U/F$ is totally $W$-aspheric), then for every arrow $p:U\longrightarrow U'$ in} Cat \textit{with source $U$, $F'=p_!(F)$ is a totally $W$-aspheric object of $U'^\wedge$ (resp.}

\msp{page 113}

\textit{$U'/F'$ is a totally $W$-aspheric object of} Cat\textit{.)}

\textbf{Corollary 3} (to theorem 3). \textit{If $i:U\longrightarrow X$ is an open or closed $W$-Dwyerian immersion (resp. totally $W$-Dwyerian), then for every arrow $p:U\longrightarrow U'$ in} Cat\textit{, $i':U'\longrightarrow X'=U'\amalg_U X$} \texttt{[rather $U'\amalg_U X$]} \textit{is so too.}
\DERmark{ENV0745}


(NB The non-resp. case is essentially theorem 3. It does not depend on any additional axiom on $W$, but is of interest only in view of the equivalences \textcircled{a} and \textcircled{b} (p. 108) which result from it, and which are valid assuming W(7 bis) resp. W(7) only.)

\textbf{Comment on motivation :} open or closed $W$-Dwyerian immersions, or totally $W$-Dwyerian ones, are stable under cobase change (by the preceding corollary), but do not at all seem to be so under composition. In this respect, ``Dwyer immersions'' (where one supposes the existence of the right adjoint to $i_0:U\longrightarrow\bar{U}$, in the case of an open immersion) behave better, being stable under composition. (\textbf{NB} It does not seem

\msp{page 114}

clear either that strictly $W$-aspheric arrows, or totally $W$-aspheric ones, are stable under composition - sorry!) My motivation for wearing myself out over the $W$-Dwyerian notion is twofold.

\begin{enumerate}[label=\alph*)]
\item I would like to have conditions at once comprehensible, and as general as possible, on an immersion (open e.g.), for it to be a $W$-cofibration, a visibly important property (just like the dual notion of $W$-fibration). % [sic] toute comme -> tout comme, agreement error
\end{enumerate}


\endgroup
% END UNIT Der-ChXIIb-012/Der-ChXIIb-012_p031-035.tex
% BEGIN UNIT Der-ChXIIb-012/Der-ChXIIb-012_p036-040.tex
\begingroup\makeatletter%
\NoAutoSpacing%
\newcommand{\printedpage}[1]{\marginpar{\footnotesize p.~#1}}%
\newcommand{\msp}[1]{\marginpar{\footnotesize [#1]}}%
\makeatother\let\printedpage\derpp\let\msp\dermsp%
\let\mspin\DERmspin%


\printedpage{78}
% head: A. Grothendieck : Dérivateurs, Ch. XII, Édition des Universités de Montpellier II et Paris VII
\begin{enumerate}
\item[b)] I would like, for later applications, to have at my disposal a class of open immersions having the following properties:
\begin{enumerate}
\item[1°)] They are $W$-cofibrations, and
\item[2°)] If $U \overset{i}{\hookrightarrow} X$ is in the class, and if $\varepsilon \xrightarrow{j} I$ is the ``source-target'' inclusion of the discrete category $\{0,1\}$ into the interval
\[
\Phi_n = 0 \longrightarrow 1 \longleftarrow 2 \longrightarrow \dots 2n-1 \longleftarrow 2n
\]
$(j(0)=0$, $j(1)=2n)$, $n\geq 2$, the inclusion $\alpha$ of $(U\times I)\cup(X\times\varepsilon)$ into $X\times I$ is in the class.
\end{enumerate}
Now this last stability property is \textit{not} satisfied by Dwyer immersions. (Take $i : \{0\} \longrightarrow \boldsymbol{\Delta}^1 = \{0,1\}$, the inclusion $\alpha$ is not a Dwyer immersion.)
\end{enumerate}

\msp{page 115}

\subsection*{2.5\quad Complement: criterion for total asphericity of a $W$-proper category over another.}

Let
\[
f : X \longrightarrow I
\]
be a proper category over an index category $I$. If $I$ has a final object $i_0$, then $X/i_0$ is identified with $X$, and consequently the inclusion
\[
X_{i_0} \longrightarrow X
\]
of the fiber of $X$ at $i_0$ is $W$-coaspheric \texttt{[rather }$W$\texttt{-coaspheric]}, and a fortiori is in $W$. Thus $X$ is aspheric if and only if $X_{i_0}$ is. We want to extend this result to a criterion for \textit{total} $W$-asphericity, more precisely to prove the
\DERmark{ENV0746}


\textbf{Proposition 11:} \textit{Let $f : X \longrightarrow I$ be $W$-proper, and let $i_0$ be a final object of $I$. If $X_{i_0}$ is totally $W$-aspheric} (def. 1, page 88), \textit{the same is true of $X$.} (Provided that $W$ satisfies the connectedness axiom W(5).)

The proof will be given below, after some preliminaries of independent interest. First of all, let us recall that if $X$ is $W$-proper over $I$, then it gives rise to a functor
\[
\begin{aligned}
I \quad &\longrightarrow \quad \mathrm{Hot}_W\\
i \quad &\longmapsto \quad \mathrm{hot}_W(X_i)(\simeq \mathrm{hot}_W(X/i));
\end{aligned}
\]
composing it with the functor
\[
\pi_0 : \mathrm{Hot}_W \longrightarrow \mathrm{Ens}
\]
(provided that the connectedness axiom W(5) (p. 10) is satisfied)
% foot: Version 7 décembre 2002
\printedpage{79}
% head: A. Grothendieck : Dérivateurs, Ch. XII, Édition des Universités de Montpellier II et Paris VII
\msp{page 116}
we find a functor
\[
i \longmapsto \underbrace{\pi_0(X_i)}_{\simeq\, \pi_0(X/i)} : I \longrightarrow \mathrm{Ens}.
\]
This being so, we have the sorital result

\textbf{Proposition 12:} \textit{Let $f : X \longrightarrow I$ be $W$-proper, and suppose that $W$ satisfies the connectedness axiom} W(5) (otherwise, $W$-properness will be of no use to us \dots) (\footnote[81]{It is better to say: Let $f$ be $W_0$-proper, where $W_0$ is the localizer of 0-connectedness \dots}). \textit{We then have a canonical bijection}
\[
\pi_0(X) \xleftarrow{\ \sim\ } \varinjlim_{I} \pi_0(X_i), \tag*{$(*)$}
\]
\textit{where the arrow is induced by the canonical maps $\pi_0(X_i) \longrightarrow \pi_0(X)$.}

\textsc{Proof.} Passing to $f^\circ : X^\circ \longrightarrow I^\circ$, we are reduced to proving the same result for $X$ \textit{$W$-smooth} over $I$, with the only difference that the functor
\[
i \longmapsto \pi_0(X_i)
\]
on $I$ will be contravariant, i.e. a presheaf, which offers the advantage of the intuition attached to the formalism of presheaves. Let
\[
g : I \longrightarrow e, \qquad \text{whence } gf = p : X \longrightarrow e,
\]
it is immediate, by the standard calculation
\[
p_!(F) = \pi_0(X/F) \qquad (F \in \mathrm{Ob}\, X^\wedge)
\]
(where we have identified $e^\wedge$ with (Ens)), that we have
\[
\pi_0(X) = p_!(e_X) = g_!(f_!(e_X)).
\]
Now the presheaf
\[
i \longmapsto \pi_0(X_i) \qquad (\simeq \pi_0(i\backslash X), \text{ since } X\ W\text{-smooth over } I)
\]
\msp{page 117}
is none other than $f_!(e_X)$. On the other hand, for any presheaf $G$ on $I$, we have
\[
g_!(G) \simeq \pi_0(I/G) \simeq \varinjlim_{I^\circ} G(i),
\]
\DERmark{ENV0747}

and applying this to $f_!(e_X)$, we find the formula $(*)$ of proposition 12.

Let us return to the situation with $f$ \textit{$W$-proper,} and consider a functor
\[
p : X \longrightarrow Y,
\]
we are interested in ``calculating''
\[
p_!(e_X) = (y \longmapsto \pi_0(y\backslash X)) \in \mathrm{Ob}\, Y^\wedge.
\]
Let us consider for this the following diagram
% foot: Version 7 décembre 2002
\printedpage{80}
% head: A. Grothendieck : Dérivateurs, Ch. XII, Édition des Universités de Montpellier II et Paris VII
\[
\begin{tikzcd}
y\backslash X_i \arrow[r] \arrow[d] & y\backslash X \arrow[r] \arrow[d] & y\backslash Y \arrow[d] \\
X_i \arrow[r] \arrow[d] & X \arrow[r, "p"] \arrow[d, "f"] & Y \\
\underset{\simeq\, e}{\{i\}} \arrow[r] & I\rlap{\quad ,} &
\end{tikzcd}
\]
where the three squares are cartesian. Since $y\backslash Y \longrightarrow Y$ is $W$-proper, the same is true of $y\backslash X \longrightarrow X$ by base change, hence of the composite $y\backslash X \longrightarrow X \longrightarrow I$ by transitivity. The diagram shows on the other hand that the fiber of $y\backslash X$ at $i \in \mathrm{Ob}\, I$ is identified with $y\backslash X_i$.

\msp{page 118}
Proposition 12 therefore gives us the formula
\[
\pi_0(y\backslash X) = \varinjlim_{I} \pi_0(y\backslash X_i),
\]
which formula cannot help being functorial in $y$, just as that of prop. 12 is in $\underline{\mathrm{Prop}}_W I$. Now for $y$ variable in $Y$, the presheaf on $Y$
\[
y \longmapsto \pi_0(y\backslash X_i)
\]
is none other than $p_{i\,!}(e_{X_i})$, where
\[
p_i : X_i \longrightarrow Y
\]
is the composite $X_i \longrightarrow X \xrightarrow{p} Y$ (induced by $p$). We have thus proved

\textbf{Corollary 1} (to prop. 12). \textit{Let $f : X \longrightarrow I$ be $W$-proper} (\footnote[82]{say $W_0$-proper \dots}), \textit{and $p : X \longrightarrow Y$ any functor. Let, for every $i \in \mathrm{Ob}\, I$, $p_i : X_i \longrightarrow Y$ be the functor induced on the fiber $X_i$ of $X$ at $i$. We then have a natural} functorial \textit{law}
\[
i \longmapsto p_{i\,!}(e_{X_i}) : I \longrightarrow Y^\wedge,
\]
\textit{and we have the formula}
\[
p_!(e_X) \xleftarrow{\ \sim\ } \varinjlim_{I} p_{i\,!}(e_{X_i}), \tag*{$(*)$}
\]
\textit{where the arrow is induced by the canonical morphisms}
\[
p_{i\,!}(e_{X_i}) \longrightarrow p_!(e_X)
\]
% foot: Version 7 décembre 2002
\printedpage{81}
% head: A. Grothendieck : Dérivateurs, Ch. XII, Édition des Universités de Montpellier II et Paris VII
\msp{page 119}
\textit{induced by the inclusions $X_i \longrightarrow X$.}

The functorial law of $p_{i\,!}(e_{X_i})$ in $i$ has been made explicit in the course of the proof.

Let us note that when $X$ over $I$ is cofibered and split, i.e. comes from a functor
\[
i \longmapsto X_i, \quad I \longrightarrow \mathrm{Cat},
\]
then the functorial law of $p_{i\,!}(e_{X_i})$ in $i$ is made explicit in a particularly pleasant way: to an arrow
\[
\alpha : i \longrightarrow j, \qquad \text{whence } \alpha_* : X_i \longrightarrow X_j,
\]
is associated
\[
p_!(\alpha) : p_{i\,!}(e_{X_i}) \longrightarrow p_{j\,!}(e_{X_j}),
\]
noting that we have a diagram
\[
\begin{tikzcd}
X_i \arrow[rr, "\alpha_*"] \arrow[dr, "p_i"', ""{name=A}] & & X_j \arrow[dl, "p_j", ""'{name=B}] \\
& Y &
\arrow[from=A, to=B, bend left, "\lambda_\alpha"] % [resolved: see DER_XIIb-a_corrections.csv]
\end{tikzcd}
\]
with a commutativity arrow
\[
\lambda_\alpha : p_i \longrightarrow p_j\alpha_*
\]
(not \texttt{[an]} isomorphism in general), whence
\[
p_{i\,!}(e_{X_i}) \xrightarrow{(\lambda_\alpha)_!} p_{j\,!}\underbrace{\alpha_{*\,!}(e_{X_i})}_{\longrightarrow\, e_{X_j}} \longrightarrow p_{j\,!}(e_{X_j}),
\]
no more doubt that's the one! (\footnote[83]{This arrow \texttt{[the arrow }$(\lambda_\alpha)_!$\texttt{]} should be made explicit with care!}).

\msp{page 120}
\textbf{Corollary 2:} \textit{Under the conditions of corollary 1, suppose that $I$ has a final object $i_0$. Then}
\[
(p_{i_0})_!(e_{X_{i_0}}) \longrightarrow p_!(e_X)
\]
\textit{is an isomorphism in $Y^\wedge$.}

Whence at once proposition 11: if $X_{i_0}$ is totally $W$-aspheric, then $X$ is too, since for every $p : X \longrightarrow Y$, $p_!(e_X) \simeq p_{i_0\,!}(e_{X_{i_0}})$, and this latter element of $Y^\wedge$ is $W$-aspheric % [sic] p_{i_0 !}(X_{i_0}) -> p_{i_0 !}(e_{X_{i_0}}), elsewhere the argument of p_{i!} is e_{X_i}
\DERmark{ENV0748}

by the hypothesis on $X_{i_0}$.

Let us note the special case:

\textbf{Corollary 3.} \textit{Let $Z$ be a totally $W$-aspheric object in} Cat. \textit{Then its product $Z\times I$ with any $I$ in} Cat \textit{having a final object is also totally $W$-aspheric.}
% foot: Version 7 décembre 2002
\printedpage{82}
% head: A. Grothendieck : Dérivateurs, Ch. XII, Édition des Universités de Montpellier II et Paris VII

\msp{page 121}

\section*{3\quad Amalgamated sums and integration\quad \begin{tikzcd}[ampersand replacement=\&, row sep=small, column sep=small, baseline=(\tikzcdmatrixname-2-1.base)]
\& X_1 \\
\int X_0 \arrow[ur] \arrow[dr] \& \\
\& X_2
\end{tikzcd}}

\subsection*{3.1\quad Canonical reduction to the case of a diagram of closed immersions}

Let us consider a diagram
\[
\begin{tikzcd}[row sep=small, column sep=small]
& X_1 \\
X_0 \arrow[ur, "\alpha_1"] \arrow[dr, "\alpha_2"'] & \\
& X_2
\end{tikzcd}
\tag*{$(*)$}
\]
in Cat. We can regard it as a contrafunctor
\[
\Phi^\circ \longrightarrow \mathrm{Cat},
\]
where
\[
\Phi = \left( \vcenter{\hbox{\begin{tikzcd}[row sep=small, column sep=small]
& 1 \arrow[dl] \\
0 & \\
& 2 \arrow[ul]
\end{tikzcd}}} \right),
\]
thus as defining a (split) fibered category $\mathcal{X}$ over $\Phi$. It is sometimes convenient to represent this category $\mathcal{X}$ by the ``integration sign'' % [sic] integration -> intégration, accent missing in the print
\DERmark{ENV0749}\DERmark{ENV0750}

\[
\begin{tikzcd}[row sep=small, column sep=small]
& X_1 \\
\int X_0 \arrow[ur] \arrow[dr] & \\
& X_2,
\end{tikzcd}
\]
and regard it as a sort of ``homotopy amalgamated sum'' of $X_1$ and $X_2$ under $X_0$, or as a ``homotopy direct limit'' of the original diagram $(*)$.

\msp{page 122}
By this one is to understand that (for a given $W \subseteq \mathrm{Fl}(\mathrm{Cat})$) the image $\mathrm{hot}_W(\mathcal{X})$ of $\mathcal{X}$ in $\mathrm{Hot}_W$ has a remarkable ``homotopy meaning'' in terms of $(*)$, more precisely, the diagram $(*)$ defines an object $\xi$ of $\mathrm{HOT}_W(\Phi)$ (which is a category of fractions of the category of all diagrams of type $(*)$ \dots), and if
\DERmark{ENV0751}\DERmark{ENV0752}

\[
\pi : \Phi \longrightarrow e
\]
denotes the absolute projection of $\Phi$, we have
\[
\pi_!(\xi) = \mathrm{hot}_W(\mathcal{X}),
\]
\DERmark{ENV0753}

% foot: Version 7 décembre 2002


\endgroup
% END UNIT Der-ChXIIb-012/Der-ChXIIb-012_p036-040.tex
% BEGIN UNIT Der-ChXIIb-012/Der-ChXIIb-012_p041-045.tex
\begingroup\makeatletter%
\newcommand{\printedpage}[1]{\marginpar{\footnotesize p.~#1}}%
\newcommand{\msp}[1]{\marginpar{\footnotesize [#1]}}%
\newcommand{\Hot}{\mathrm{Hot}}%
\newcommand{\HOT}{\mathrm{HOT}}%
\newcommand{\Cat}{\mathrm{Cat}}%
\newcommand{\cX}{\mathcal{X}}%
\makeatother\let\printedpage\derpp\let\msp\dermsp%
\let\mspin\DERmspin%


\printedpage{83}
% head: A. Grothendieck : Dérivateurs, Ch. XII, Édition des Universités de Montpellier II et Paris VII
by the known computation of the functor $\pi_!$, left adjoint to the functor $\pi^*$
\[
\HOT_W(e) \overset{\mathrm{def}}{=} \Hot_W \longrightarrow \HOT_W(\Phi)
\]
\DERmark{ENV0754}

(\footnote[84]{This assumes that we have W(2), not only W(2 bis).}). Now for every object $I$ in Cat, with absolute projection $p_I : I \longrightarrow e$, the functor
\[
p_{I\,!} : \HOT_W(I) \longrightarrow \HOT_W(e) = \Hot_W
\]
must be regarded as the ``homotopic'' equivalent of the direct limit functor $\varinjlim\limits_{I}$ on diagrams of type $I^{\circ}$ in the category $\Hot_W$ - which direct limits \textit{do not exist} in $\Hot_W$ (except
\msp{page 123}
for very special categories - the only ones I know are the discrete categories (which perhaps even have to be assumed finite as well)).

Let us note that we have a diagram of canonical functors
\[
\begin{tikzcd}[column sep=tiny, row sep=small, ampersand replacement=\&]
 \& X_0 \arrow[dl, "\alpha_1"'] \arrow[dr, "\alpha_2"] \& \\
X_1 \arrow[dr, "i_1"'] \& {\shortstack{\scriptsize not\\ \scriptsize commuta-\\ \scriptsize tive!}} \& X_2 \arrow[dl, "i_2"] \\
 \& \cX, \&
\end{tikzcd}
\tag*{$(*)$}
\]
which diagram has nothing commutative about it in Cat. \textit{It is however commutative in} $\Hot_W$ (whatever the basic localizer $W$), and even in Cat\,hot, for if we denote by
\[
i_0 : X_0 \longrightarrow \cX
\]
the canonical inclusion, we have canonical homomorphisms
\[
i_1\alpha_1 \longrightarrow i_0 \longleftarrow i_2\alpha_2,
\]
which shows that $i_1\alpha_1$ and $i_2\alpha_2$ are homotopic, hence define the same arrow in $\Hot_W$.

\msp{page 124}
Let us consider in $\Phi$ the two closed subsets isomorphic to $\boldsymbol{\Delta}^1$
\[
\begin{aligned}
F_1 &= 1\backslash\Phi &&= \{0,1\}\\
F_2 &= 2\backslash\Phi &&= \{0,2\},
\end{aligned}
\]
we have
\[
F_1\cap F_2 = 0\backslash\Phi = \{0\}, \qquad F_1\cup F_2 = \Phi.
\]

\printedpage{84}
% head: A. Grothendieck : Dérivateurs, Ch. XII, Édition des Universités de Montpellier II et Paris VII
Let
\[
\begin{aligned}
1\backslash\cX &= \cX\times_{\Phi} 1\backslash\Phi &&= \cX\times_{\Phi} F_1\\
2\backslash\cX &= \cX\times_{\Phi} 2\backslash\Phi &&= \cX\times_{\Phi} F_2
\end{aligned}
\]
be the inverse images of $F_1$, $F_2$ in $\cX$, we therefore have
\[
1\backslash\cX\cap 2\backslash\cX = X_0, \qquad 1\backslash\cX\cup 2\backslash\cX = \cX,
\]
so the diagram
\[
\begin{tikzcd}[column sep=small, row sep=small, ampersand replacement=\&]
 \& X_0 \arrow[dl, hook', "\alpha'_1"'] \arrow[dr, hook, "\alpha'_2"] \& \\
1\backslash\cX \arrow[dr, hook, "i'_1"'] \& \& 2\backslash\cX \arrow[dl, hook', "i'_2"] \\
 \& \cX, \&
\end{tikzcd}
\tag*{$(**)$}
\]
deduced from the similar diagram
\[
\begin{tikzcd}[column sep=small, row sep=small, ampersand replacement=\&]
 \& \{0\} \arrow[dl, hook'] \arrow[dr, hook] \& \\
F_1 \arrow[dr, hook] \& \& F_2 \arrow[dl, hook'] \\
 \& \Phi, \&
\end{tikzcd}
\]
is (like the latter) a square of closed immersions both cartesian and cocartesian. On the other hand, we have a canonical homomorphism
\msp{page 125}
from the diagram $(*)$ to $(**)$, which is the identity on the vertices $X_0$ and $\cX$, and which on $X_1$, $X_2$ is given by the canonical inclusions
\[
\underbrace{X_1}_{\text{\scriptsize $\cX_1$ fiber of $\cX$ at 1}} \longrightarrow 1\backslash\cX, \qquad
\underbrace{X_2}_{\text{\scriptsize $\cX_2$ fiber of $\cX$ at 2}} \longrightarrow 2\backslash\cX.
\]
These inclusions, as we know, are $W$-aspheric (they even have a right adjoint), a fortiori they are in $W$. Thus the noncommutative diagram $(*)$ is isomorphic, as a diagram in $\Hot_W$, to the commutative cartesian and cocartesian diagram $(**)$ formed of closed immersions. Since $\cX$ is indeed an amalgamated sum of its two closed subsets % [sic] fermées -> fermés, cf. « les deux fermés » p. 83
$1\backslash\cX$, $2\backslash\cX$ under their intersection $X_0$, this supports the intuition which sees $\cX$ as a ``homotopic amalgamated sum''. Let us give at once the following application:

\textbf{Proposition 13.} \textit{Let there be a diagram} $\begin{tikzcd}[column sep=small, row sep=tiny, ampersand replacement=\&, baseline=(\tikzcdmatrixname-2-1.base)]
 \& X_1 \\
X_0 \arrow[ur, "\alpha_1"] \arrow[dr, "\alpha_2"'] \& \\
 \& X_2
\end{tikzcd}$ \textit{in} Cat\textit{, whence the diagram} $(*)$\textit{, p. 123,}

\printedpage{85}
% head: A. Grothendieck : Dérivateurs, Ch. XII, Édition des Universités de Montpellier II et Paris VII
\[
\begin{tikzcd}[column sep=small, row sep=small, ampersand replacement=\&]
 \& X_0 \arrow[dl, "\alpha_1"'] \arrow[dr, "\alpha_2"] \& \\
X_1 \arrow[dr, "i_1"'] \& \& X_2 \arrow[dl, "i_2"] \\
 \& \cX, \&
\end{tikzcd}
\]
\msp{page 126}
\textit{commutative in} $\Hot_W$ \textit{(but not in} Cat\textit{, unless $X_0 = \emptyset$). Let us suppose that $W$ satisfies} W(7 bis) (cocartesian square axiom, cf. p. 75). \textit{Then we have}
\[
\left\{
\begin{aligned}
\alpha_1\in W &\iff i_2\in W\\
\alpha_2\in W &\iff i_1\in W. \text{ (\footnotemark[85])}
\end{aligned}
\right.
\]
\footnotetext[85]{\textbf{NB} For the implication $\Longrightarrow$ we need only the direct form of W(7 bis), i.e. (with the notations of XVI) of Loc(6$'$a).}
\textsc{Proof}. Using the preceding homomorphism from $(*)$ to $(**)$, we see that each of the arrows $\alpha_1$, $\alpha_2$, $i_1$, $i_2$ is in $W$, if and only if the corresponding arrow $\alpha'_1$, $\alpha'_2$, $i'_1$ or $i'_2$ is. Thus the above equivalences are equivalent to
\[
\begin{aligned}
\alpha'_1\in W &\iff i'_2\in W\\
\alpha'_2\in W &\iff i'_1\in W,
\end{aligned}
\]
which are true by virtue of W(7 bis) (dual formulation) applied to $\cX$ and to the two closed subsets $1\backslash\cX$, $2\backslash\cX$ of $\cX$.

\textbf{Corollary:} \textit{Let us consider} two \textit{diagrams of type} $\Phi^{\circ}$\textit{, and a homomorphism} $f = (f_0, f_1, f_2)$ \textit{from one to the other (cf. the diagram in solid lines below} \texttt{[below]}\textit{):}
\DERmark{ENV0755}

\[
\begin{tikzcd}[column sep=small, row sep=small, ampersand replacement=\&]
 \& X_1 \arrow[dd, "f_1"'] \arrow[drr, dotted, "i_1"] \& \& \\
X_0 \arrow[ur, "\alpha_1"] \arrow[drr, "\alpha_2"] \arrow[dd, "f_0"'] \& \& \& \cX \arrow[dd, dotted, "f_{\cX}"] \\
 \& X'_1 \arrow[drr, dotted, "i'_1"] \& X_2 \arrow[ur, dotted, "i_2"] \arrow[dd, "f_2"'] \& \\
X'_0 \arrow[ur, "\alpha'_1"] \arrow[drr, "\alpha'_2"] \& \& \& \cX', \\
 \& \& X'_2 \arrow[ur, dotted, "i'_2"] \&
\end{tikzcd}
\]
\textit{whence a homomorphism}
\msp{page 127}
\[
f_{\cX} = \int f : \quad \cX = \int
\begin{tikzcd}[column sep=small, row sep=tiny, ampersand replacement=\&, baseline=(\tikzcdmatrixname-2-1.base)]
 \& X_1 \\
X_0 \arrow[ur] \arrow[dr] \& \\
 \& X_2
\end{tikzcd}
\qquad\longrightarrow\qquad
\cX' = \int
\begin{tikzcd}[column sep=small, row sep=tiny, ampersand replacement=\&, baseline=(\tikzcdmatrixname-2-1.base)]
 \& X'_1 \\
X'_0 \arrow[ur] \arrow[dr] \& \\
 \& X'_2,
\end{tikzcd}
\]

\printedpage{86}
% head: A. Grothendieck : Dérivateurs, Ch. XII, Édition des Universités de Montpellier II et Paris VII
\textit{fitting into a cubical diagram (cf. the preceding diagram completed by dotted lines). Let us consider the following conditions.}

\begin{itemize}
\item[\textit{a)}] $f_0\in W$\textit{.}
\item[\textit{b)}] $f_1$ \textit{and} $f_2$ \textit{are} \texttt{[in]} $W$\textit{.}
\item[\textit{c)}] $f_{\cX} = \int(f_0, f_1, f_2)\in W$\textit{.}
\end{itemize}

\textit{Let us suppose the axiom} W(7) (``5 lemma'', cf. page 76) \textit{is satisfied. Then two among the conditions a), b), c) imply the third}
\[
\left\{
\begin{aligned}
a \mathbin{\&} b &\quad\Longrightarrow\quad c\\
b \mathbin{\&} c &\quad\Longrightarrow\quad a\\
c \mathbin{\&} a &\quad\Longrightarrow\quad b.
\end{aligned}
\right.
\]

Indeed, we have a homomorphism from the cubical diagram above to the similar diagram, where $X_1$, $X_2$, $X'_1$, $X'_2$ are replaced respectively by the closed subsets % [sic] pas -> par, misprint
\DERmark{ENV0756}

$1\backslash\cX$, $2\backslash\cX$ of $\cX$, and $1\backslash\cX'$, $2\backslash\cX'$ of $\cX'$, the two squares, upper and lower, being inclusions of closed subsets \texttt{[rather of closed subsets]} of $\cX$ resp. of $\cX'$. Thus, conditions a), b), c) are equivalent respectively to the similar conditions for the new cubical diagram
\DERmark{ENV0757}

\msp{page 128}
(a) and c) moreover not having changed, and b) giving
\begin{itemize}
\item[b$'$)] $f'_1 : 1\backslash\cX \longrightarrow 1\backslash\cX'$\ \ and\ \ $f'_2 : 2\backslash\cX \longrightarrow 2\backslash\cX'$\ \ are in $W$.
\end{itemize}
\DERmark{ENV0758}\DERmark{ENV0759}

\texttt{[rather} $f'_1 : 1\backslash\cX \longrightarrow 1\backslash\cX'$ \texttt{and} $f'_2 : 2\backslash\cX \longrightarrow 2\backslash\cX'$\texttt{]}). Then the conclusion of the corollary is equivalent to W(7) (dual formulation) in the case of the new cubical diagram, q.e.d.

\subsection*{3.2\quad Comparison of \ $\displaystyle\int
\begin{tikzcd}[column sep=small, row sep=tiny, ampersand replacement=\&, baseline=(\tikzcdmatrixname-2-1.base)]
 \& X_1 \\
X_0 \arrow[ur] \arrow[dr] \& \\
 \& X_2
\end{tikzcd}$ \ with an ordinary amalgamated sum.\\
$W$-homotopically cocartesian squares.}

The integrated category $\cX = \displaystyle\int
\begin{tikzcd}[column sep=small, row sep=tiny, ampersand replacement=\&, baseline=(\tikzcdmatrixname-2-1.base)]
 \& X_1 \\
X_0 \arrow[ur] \arrow[dr] \& \\
 \& X_2
\end{tikzcd}$ is a solution of the following universal problem, a 2-categorical variant of the one which defines the direct limit of diagrams. It is a matter, for $Y$ a variable object in Cat, of considering the systems of three functors $f_0$, $f_1$, $f_2$ and of the two functorial homomorphisms $\lambda_1$, $\lambda_2$:
\printedpage{87}
% head: A. Grothendieck : Dérivateurs, Ch. XII, Édition des Universités de Montpellier II et Paris VII
\[
\begin{tikzcd}[column sep=large, row sep=large, ampersand replacement=\&]
 \& X_0 \arrow[dl, "\alpha_1"'] \arrow[dr, "\alpha_2"] \arrow[dd, "f_0", ""{name=M, pos=0.62}] \& \\
X_1 \arrow[dr, "f_1"', ""{name=A, pos=0.45}] \& \& X_2 \arrow[dl, "f_2", ""{name=B, pos=0.45}] \\
 \& Y, \&
\arrow[from=A, to=M, bend left, shorten=3pt, "\lambda_1"]
\arrow[from=B, to=M, bend right, shorten=3pt, "\lambda_2"']
\end{tikzcd}
\tag*{$(*)$}
\]
\[
\begin{aligned}
\lambda_1 :\quad & f_1\alpha_1 \quad\longrightarrow\quad f_0\\
\lambda_2 :\quad & f_2\alpha_2 \quad\longrightarrow\quad f_0.
\end{aligned}
\]

\msp{page 129}
To give oneself an arrow of Cat
\[
\cX \longrightarrow Y
\]
``comes to the same thing'' as giving oneself a system $(f_0, f_1, f_2, \lambda_1, \lambda_2)$ as above, and this ``functorially in $Y$''. This is the special case of the general determination of the functors from the ``integrated'' fibered category of a contrafunctor
\[
I^{\circ} \longrightarrow \Cat, \qquad
\begin{aligned}
i &\quad\longmapsto\quad X_i\\
u &\quad\longmapsto\quad u^*
\end{aligned}
\]
to a category $Y$, in terms of a system of functors
\[
X_i \xrightarrow{f_i} Y,
\]
and of functorial homomorphisms
\[
f_i\circ u^* \xrightarrow{\lambda_u} f_j
\]
for $u : i \longrightarrow j$, whence $u^* : X_j \longrightarrow X_i$, with transitivity conditions for composite arrows, which we need not take into account in the present case, seeing that there are no composite arrows to look at in $I = \Phi$ (except for composites with an identity arrow, which we can dispense with looking at here \dots).

\msp{page 130}
If in particular we have a \textit{commutative} diagram in Cat


\endgroup
% END UNIT Der-ChXIIb-012/Der-ChXIIb-012_p041-045.tex
% BEGIN UNIT Der-ChXIIb-012/Der-ChXIIb-012_p046-050.tex
\begingroup\makeatletter%
\newcommand{\printedpage}[1]{\marginpar{\footnotesize p.~#1}}%
\newcommand{\msp}[1]{\marginpar{\footnotesize [#1]}}%
\newcommand{\sX}{\mathcal{X}}%
\newcommand{\tX}{\tilde{\mathcal{X}}}%
\makeatother\let\printedpage\derpp\let\msp\dermsp%
\let\mspin\DERmspin%

\DERmark{ENV0760}\DERmark{ENV0761}


\printedpage{88}
% head: A. Grothendieck : Dérivateurs, Ch. XII, Édition des Universités de Montpellier II et Paris VII
\[
Q \quad = \qquad
\begin{tikzcd}
 & X_0 \arrow[dl, "\alpha_1"'] \arrow[dr, "\alpha_2"] & \\
X_1 \arrow[dr, "f_1"'] & & X_2 \arrow[dl, "f_2"] \\
 & Y, &
\end{tikzcd}
\]
we can complete it by
\[
f_0 \overset{\mathrm{def}}{=} f_1\alpha_1 = f_2\alpha_2
\]
\DERmark{ENV0762}

into a diagram
\[
\begin{tikzcd}[row sep=large, column sep=large]
 & X_0 \arrow[dl, "\alpha_1"'] \arrow[dr, "\alpha_2"] \arrow[dd, "f_0", ""'{name=F0, pos=0.65}] & \\
X_1 \arrow[dr, "f_1"', ""{name=F1, pos=0.4}] & & X_2 \arrow[dl, "f_2", ""'{name=F2, pos=0.4}] \\
 & Y, &
\arrow[from=F1, to=F0, bend left, "\mathrm{id}", shorten <=2pt, shorten >=2pt]
\arrow[from=F2, to=F0, bend right, "\mathrm{id}"', shorten <=2pt, shorten >=2pt]
\end{tikzcd}
\]
which shows that we have a datum of the type $(*)$ above, thus defining an arrow
\[
\varphi_Q : \quad \sX = \int
\begin{tikzcd}[row sep=small, column sep=small]
 & X_1 \\
X_0 \arrow[ur] \arrow[dr] & \\
 & X_2
\end{tikzcd}
\qquad \longrightarrow \qquad Y
\]
which gives rise to a cubical diagram

\msp{page 131}
\[
\begin{tikzcd}[row sep=small, column sep=small]
 & X_1 \arrow[dd, "\mathrm{id}"'] \arrow[rrd, "i_1", ""'{name=I1}] & & \\
X_0 \arrow[ru, "\alpha_1"] \arrow[rrr, dotted, "i_0", ""'{name=I0, pos=0.8}] \arrow[rrd, "\alpha_2"] \arrow[dd, "\mathrm{id}"'] & & & \sX \arrow[dd, "\varphi"] \\
 & X_1 \arrow[rrd, "f_1"] & X_2 \arrow[ru, "i_2"', ""{name=I2}] \arrow[dd, "\mathrm{id}"'] & \\
X_0 \arrow[ru, "\alpha_1"] \arrow[rrr, dotted, "f_0"] \arrow[rrd, "\alpha_2"] & & & Y, \\
 & & X_2 \arrow[ru, "f_2"'] &
\arrow[from=I1, to=I0, "\lambda_1"', shorten <=1pt, shorten >=1pt]
\arrow[from=I2, to=I0, bend left, "\lambda_2", shorten <=1pt, shorten >=1pt] % Resolved from PDF leaf 11: a solid curve, unlike the genuinely dotted i_0 line.
\end{tikzcd}
\]
% foot: Version 7 décembre 2002

\printedpage{89}
% head: A. Grothendieck : Dérivateurs, Ch. XII, Édition des Universités de Montpellier II et Paris VII
commutative except for the top square, i.e. we have
\[
\varphi i_1 = f_1, \qquad \varphi i_2 = f_2, \qquad \text{and moreover} \quad \varphi i_0 = f_0,
\]
$\varphi$ moreover transforming the functorial homomorphisms
\[
\lambda_1 : i_1\alpha_1 \longrightarrow i_0, \qquad \lambda_2 : i_2\alpha_2 \longrightarrow i_0
\]
into the relevant identity functorial homomorphisms.

\setcounter{footnote}{85}
\textbf{Definition 3.} We say that the commutative diagram
\[
\begin{tikzcd}[row sep=small, column sep=small]
 & X_0 \arrow[dl, "\alpha_1"'] \arrow[dr, "\alpha_2"] & \\
X_1 \arrow[dr, "f_1"] & & X_2 \arrow[dl, "f_2"'] \\
 & Y &
\end{tikzcd}
\tag*{$(*)$}
\]
is \textit{$W$-homotopically cocartesian} (\footnote{or simply $W$-cocartesian, or that it defines $Y$ as a $W$-homotopic amalgamated sum of etc.}) if the canonical homomorphism
\[
\sX = \int
\begin{tikzcd}[row sep=small, column sep=small]
 & X_1 \\
X_0 \arrow[ur] \arrow[dr] & \\
 & X_2
\end{tikzcd}
\qquad \xrightarrow{\ \varphi\ } \qquad Y
\]
it defines is in $W$. (\footnote{Introduce $\varphi$ in a proposition in due form, and give its essential properties: strictly cartesian arrows are transformed into \textit{identity} arrows $+$ lemmas 1 and 2 of pp. 152, 154.})

\msp{page 132}
We then have the:

\textbf{Corollary 1.} \textit{Suppose} W(7 bis)\textit{. If the square} $(*)$ \textit{is $W$-homotopically cocartesian, then}
\[
\left\{
\begin{aligned}
\alpha_1 \in W &\quad\iff\quad f_2 \in W \\
\alpha_2 \in W &\quad\iff\quad f_1 \in W.
\end{aligned}
\right.
\]

This follows at once from the cubical diagram on the preceding page, where all the vertical arrows are in $W$, and from the similar statement concerning the top face of the cube, given in proposition 13, p. 125.

\textbf{Corollary 2.} \textit{Suppose} W(7)\textit{. Consider a cubical diagram}
% foot: Version 7 décembre 2002

\printedpage{90}
% head: A. Grothendieck : Dérivateurs, Ch. XII, Édition des Universités de Montpellier II et Paris VII
\[
\begin{tikzcd}[row sep=small, column sep=small]
 & X_1 \arrow[dd, "p_1"'] \arrow[rrd, "f_1"] & & \\
X_0 \arrow[ru, "\alpha_1"] \arrow[rrd, "\alpha_2"] \arrow[dd, "p_0"'] & & & Y \arrow[dd, "p"] \\
 & X'_1 \arrow[rrd, "f'_1"] & X_2 \arrow[ru, "f_2"] \arrow[dd, "p_2"'] & \\
X'_0 \arrow[ru, "\alpha'_1"] \arrow[rrd, "\alpha'_2"] & & & Y', \\
 & & X'_2 \arrow[ru, "f'_2"] &
\end{tikzcd}
\]
\textit{where the two faces, bottom and top, of the cube are $W$-homotopically cocartesian squares. Consider the conditions}

\msp{page 133}
\begin{enumerate}[label={}]
\item[\textit{a)}] $p_0 \in W$\textit{.}
\item[\textit{b)}] $p_1$ \textit{and} $p_2 \in W$\textit{.}
\item[\textit{c)}] $p \in W$\textit{.}
\end{enumerate}

\textit{Then two of these conditions imply the third.}

Follows at once from definition 3, and from the similar result (cor. page 126) for the similar cubical diagram, with $Y$, $Y'$ replaced by $\sX = \int
\begin{tikzcd}[row sep=small, column sep=small]
 & X_1 \\
X_0 \arrow[ur] \arrow[dr] & \\
 & X_2
\end{tikzcd}$
and $\sX' = \int
\begin{tikzcd}[row sep=small, column sep=small]
 & X'_1 \\
X'_0 \arrow[ur] \arrow[dr] & \\
 & X'_2
\end{tikzcd}$\;.

\section*{3.3\quad Proper integration versus smooth integration. Statement of equivalence between the two points of view. Comments on the case of arbitrary $W$-homotopic limits.}

Definition 3, and the preceding corollaries, are of interest only insofar as we have workable criteria, allowing us to assert that a commutative square really is homotopically cocartesian. This is what we are going to try to do shortly. There is also a perplexity, concerning the privileged role played, in the theory developed here, by closed immersions, to the exclusion of open immersions. We have in short reduced, from the
% foot: Version 7 décembre 2002
\printedpage{91}
% head: A. Grothendieck : Dérivateurs, Ch. XII, Édition des Universités de Montpellier II et Paris VII
\msp{page 134}
homotopical point of view, arbitrary diagrams $\begin{tikzcd}[row sep=small, column sep=small]
 & X_1 \\
X_0 \arrow[ur] \arrow[dr] & \\
 & X_2
\end{tikzcd}$ in Cat, to the case of a diagram of closed immersions, where the amalgamated sum takes on a familiar look. But why not just as well reduce it to the case of open immersions? We would then consider the initial diagram as defining a \textit{cofibered} category over
\[
\Phi^{\circ} = \Psi = \
\begin{tikzcd}[row sep=small, column sep=small]
 & 1 \\
0 \arrow[ur] \arrow[dr] & \\
 & 2,
\end{tikzcd}
\]
say $\tX$. We have in $\Psi$ the two opens
\[
\underbrace{\Psi/1}_{=\{0,1\}} = U_1, \qquad \underbrace{\Psi/2}_{=\{0,2\}} = U_2, \qquad U_1 \cap U_2 = \{0\},
\]
whence in $\tX$ two inverse image opens
\[
\tX/1, \qquad \tX/2, \qquad \text{with } \tX/1 \cap \tX/2 = X_0,
\]
whence a commutative diagram of open inclusions
\[
\begin{tikzcd}
 & X_0 \arrow[dl, "\alpha'_1"'] \arrow[dr, "\alpha'_2"] & \\
\tX/1 \arrow[dr, "i'_1"'] & & \tX/2 \arrow[dl, "i'_2"] \\
 & \tX, &
\end{tikzcd}
\tag*{$(*)$}
\]

\msp{page 135}
which receives the noncommutative diagram
\[
\begin{tikzcd}[row sep=large, column sep=large]
 & X_0 \arrow[dl, "\alpha_1"'] \arrow[dr, "\alpha_2"] \arrow[dd, dotted, "i_0", ""'{name=J0, pos=0.6}] & \\
X_1 \arrow[dr, "i_1"', ""{name=J1, pos=0.4}] & & X_2 \arrow[dl, "i_2", ""'{name=J2, pos=0.4}] \\
 & \tX, &
\arrow[from=J0, to=J1, dotted, bend right, "\tilde{\lambda}_1"', shorten <=2pt, shorten >=2pt]
\arrow[from=J0, to=J2, dotted, bend left, "\tilde{\lambda}_2", shorten <=2pt, shorten >=2pt]
\end{tikzcd}
\]
% foot: Version 7 décembre 2002
\printedpage{92}
% head: A. Grothendieck : Dérivateurs, Ch. XII, Édition des Universités de Montpellier II et Paris VII
the arrows
\[
X_1 \longrightarrow \tX/1, \qquad X_2 \longrightarrow \tX/2
\]
being in $W$ (and even $W$-coaspheric - they even have a left adjoint).

So we find, as before
\[
\left\{
\begin{aligned}
\alpha_1 \in W &\quad\iff\quad i_2 \in W \\
\alpha_2 \in W &\quad\iff\quad i_1 \in W,
\end{aligned}
\right.
\]
and likewise we can formulate a ``five lemma'' like the corollary on page 125 \texttt{[rather 126]}, but with $\tX$, $\tX'$ instead of $\sX$, $\sX'$.

Moreover, $\tX$ is the solution of a universal problem in Cat, in terms of data
\[
\begin{tikzcd}[row sep=large, column sep=large]
 & X_0 \arrow[dl, "\alpha_1"'] \arrow[dr, "\alpha_2"] \arrow[dd, "f_0", ""'{name=G0, pos=0.6}] & \\
X_1 \arrow[dr, "f_1"', ""{name=G1, pos=0.4}] & & X_2 \arrow[dl, "f_2", ""'{name=G2, pos=0.4}] \\
 & Y, &
\arrow[from=G0, to=G1, bend right, "\tilde{\lambda}_1"', shorten <=2pt, shorten >=2pt]
\arrow[from=G0, to=G2, bend left, "\tilde{\lambda}_2", shorten <=2pt, shorten >=2pt]
\end{tikzcd}
\]

\msp{page 136}
the only difference from the one formulated on page 128 being the opposite direction of the arrows $\lambda_1$, $\lambda_2$. For every commutative square
\[
\begin{tikzcd}
 & X_0 \arrow[dl, "\alpha_1"'] \arrow[dr, "\alpha_2"] & \\
X_1 \arrow[dr, "f_1"'] & & X_2 \arrow[dl, "f_2"] \\
 & Y, &
\end{tikzcd}
\tag*{$(*)$}
\]
we again find a canonical arrow
\[
\tilde{\varphi} : \tX \longrightarrow Y.
\]

It would be just as reasonable to say that the square $(*)$ is $W$-homotopically cocartesian, when $\tilde{\varphi}$ is \texttt{[in]} $W$, as to adopt definition 3, page 131 - we would find exactly the same corollaries. Would the two definitions be equivalent?? % [sic] seraient elles -> seraient-elles, hyphen of the inverted subject missing
% foot: Version 7 décembre 2002


\endgroup
% END UNIT Der-ChXIIb-012/Der-ChXIIb-012_p046-050.tex
% BEGIN UNIT Der-ChXIIb-012/Der-ChXIIb-012_p051-055.tex
\begingroup\makeatletter%
\newcommand{\printedpage}[1]{\marginpar{\footnotesize p.~#1}}%
\newcommand{\msp}[1]{\marginpar{\footnotesize [#1]}}%
\newcommand{\Cat}{\mathrm{Cat}}%
\newcommand{\Hom}{\mathrm{Hom}}%
\newcommand{\HOT}{\mathrm{HOT}}%
\newcommand{\Hot}{\mathrm{Hot}}%
\newcommand{\Fl}{\mathrm{Fl}}%
\newcommand{\sX}{\mathcal{X}}%
\newcommand{\tsX}{\tilde{\mathcal{X}}}%
\newcommand{\circled}[1]{\tikz[baseline=(char.base)]{\node[shape=circle,draw,inner sep=1pt] (char) {#1};}}%
\makeatother\let\printedpage\derpp\let\msp\dermsp%
\let\mspin\DERmspin%

\DERmark{ENV0763}\DERmark{ENV0764}


\printedpage{93}
% head: A. Grothendieck : Dérivateurs, Ch. XII, Édition des Universités de Montpellier II et Paris VII
% foot: Version 7 décembre 2002

I note that $\tsX$ being cofibered over $\Psi=\Phi^\circ$, $\tsX^\circ$ is fibered over $\Psi^\circ=\Phi$. More precisely, we will have
\[
\tsX^\circ = \int
\begin{tikzcd}[row sep=small, column sep=small]
 & (X_1)^\circ \\
(X_0)^\circ \arrow[ur, "\alpha_1^\circ"] \arrow[dr, "\alpha_2^\circ"'] & \\
 & (X_2)^\circ,
\end{tikzcd}
\]
\msp{page 137}
and it is thus a matter of comparing the two integrals
\[
\sX = \int
\begin{tikzcd}[row sep=small, column sep=small]
 & X_1 \\
X_0 \arrow[ur, "\alpha_1"] \arrow[dr, "\alpha_2"'] & \\
 & X_2
\end{tikzcd}
, \qquad
\tsX^\circ = \int
\begin{tikzcd}[row sep=small, column sep=small]
 & (X_1)^\circ \\
(X_0)^\circ \arrow[ur, "\alpha_1^\circ"] \arrow[dr, "\alpha_2^\circ"'] & \\
 & (X_2)^\circ
\end{tikzcd}
,
\]
in particular as objects of $\HOT_W(\Phi)$. For this, I take up again the construction of the functor
\[
X \longmapsto \Phi(X) : \Cat \longrightarrow \Cat
\]
from I, p. 72-75, but I prefer here to denote this functor by $\Sigma(X)$ (for $\Phi$ already denotes a small category $\begin{tikzcd}[row sep=tiny, column sep=small]
 & 1 \arrow[dl] \\
0 & \\
 & 2 \arrow[ul]
\end{tikzcd}$). We will use the homomorphisms functorial in $X$:
\[
\begin{tikzcd}[column sep=small]
 & \Sigma(X) \arrow[dl, "p_X"'] \arrow[dr, "q_X"] & \\
X & & X^\circ,
\end{tikzcd}
\]
which are $W$-aspheric, hence $\in W$. Moreover, $\Sigma(X^\circ)=\Sigma(X)$, $p_{X^\circ}=q_X$, $q_{X^\circ}=p_X$. Thus, if
\[
\begin{aligned}
X_* \quad &= \qquad
\begin{tikzcd}[ampersand replacement=\&, row sep=small, column sep=small]
 \& X_1 \\
X_0 \arrow[ur, "\alpha_1"] \arrow[dr, "\alpha_2"'] \& \\
 \& X_2
\end{tikzcd}
\qquad \in \quad \Hom(\Phi^\circ,\Cat)\\
X_*^\circ \quad &= \qquad
\begin{tikzcd}[ampersand replacement=\&, row sep=small, column sep=small]
 \& (X_1)^\circ \\
(X_0)^\circ \arrow[ur, "\alpha_1^\circ"] \arrow[dr, "\alpha_2^\circ"'] \& \\
 \& (X_2)^\circ
\end{tikzcd}
\qquad \in \quad \Hom(\Phi^\circ,\Cat),
\end{aligned}
\]
\msp{page 138}
we find
\[
\Sigma(X_*) =
\left(
\begin{tikzcd}[row sep=small, column sep=small]
 & \Sigma(X_1) \\
\Sigma(X_0) \arrow[ur, "\Sigma(\alpha_1)"] \arrow[dr, "\Sigma(\alpha_2)"'] & \\
 & \Sigma(X_2)
\end{tikzcd}
\right)
\]
and the arrows
\printedpage{94}
% head: A. Grothendieck : Dérivateurs, Ch. XII, Édition des Universités de Montpellier II et Paris VII
% foot: Version 7 décembre 2002
\[
\begin{tikzcd}[column sep=small]
 & \Sigma(X_*) \arrow[dl, "p_{X_*}"'] \arrow[dr, "q_{X_*}"] & \\
X_* & & X_*^\circ,
\end{tikzcd}
\]
in $\underline{\Hom}(\Phi^\circ,\Cat)$. These arrows are in $W$ term by term, so the homomorphisms in $\Cat$
\[
\begin{tikzcd}[column sep=small]
 & \int\Sigma(X_*) \arrow[dl, "\int p_{X_*}"'] \arrow[dr, "\int q_{X_*}"] & \\
\sX=\int X_* & & \int X_*^\circ=\tsX^\circ
\end{tikzcd}
\]
are in $W$.

Finally, the innocent commutative square
\[
Q \quad = \qquad
\begin{tikzcd}[row sep=small, column sep=small]
 & X_0 \arrow[dl] \arrow[dr] & \\
X_1 \arrow[dr] & & X_2 \arrow[dl] \\
 & Y &
\end{tikzcd}
\tag*{$(*)$}
\]
(and the corresponding diagram
\[
\widehat{Q} \quad = \qquad
\begin{tikzcd}[row sep=small, column sep=small]
 & X_0 \arrow[dl] \arrow[dr] & \\
X_1 \arrow[dr] & & X_2 \arrow[dl] \\
 & \sX=\int X_* \arrow[d, "\varphi_Q"] & \\
 & Y &
\end{tikzcd}
\qquad\qquad )
\tag*{$(**)$}
\]
gives rise, by transforming it
\msp{page 139}
by $X\longmapsto X^\circ$ and by $X\longmapsto\Sigma(X)$, and by taking into account the homomorphisms functorial in $X$ $\begin{tikzcd}[row sep=small, column sep=small]
 & \Sigma(X) \arrow[dl] \arrow[dr] & \\
X & & X^\circ
\end{tikzcd}$, to three diagrams of type $(**)$, i.e. modeled on the opposite category of
\[
\Gamma \quad = \qquad
\begin{tikzcd}[row sep=small, column sep=small]
 & 0 & \\
1 \arrow[ur, "\alpha_1"] & & 2 \arrow[ul, "\alpha_2"'] \\
 & 3 \arrow[ul, "i_1"] \arrow[ur, "i_2"'] & \\
 & 4 \arrow[u, "\varphi"'] &
\end{tikzcd}
\qquad
\text{\parbox[c]{0.35\textwidth}{\footnotesize (i.e. the free diamond (\footnotemark[88]), to which we have added an initial object)}} % [sic] auquelle -> à laquelle, grammatical form (catégorie is feminine)
\]
\DERmark{ENV0765}

\footnotetext[88]{false, for the square is not commutative}
\printedpage{95}
% head: A. Grothendieck : Dérivateurs, Ch. XII, Édition des Universités de Montpellier II et Paris VII
% foot: Version 7 décembre 2002
(\footnote[89]{We have $\alpha_1i_1\neq\alpha_2i_2$, but the relation $\alpha_1i_1\varphi=\alpha_2i_2\varphi$.}), namely
\[
\widehat{Q}, \qquad \widehat{\Sigma(Q)}, \qquad \widehat{Q^\circ},
\]
and two homomorphisms of diagrams (\texttt{[in]} $\underline{\Hom}(\Gamma^\circ,\Cat)$)
\[
\begin{tikzcd}[column sep=small]
 & \widehat{\Sigma(Q)} \arrow[dl] \arrow[dr] & \\
\widehat{Q} & & \widehat{Q^\circ}
\end{tikzcd}
\qquad
\text{(deduced from }
\begin{tikzcd}[column sep=small]
 & \Sigma(Q) \arrow[dl] \arrow[dr] & \\
Q & & Q^\circ
\end{tikzcd}
\text{),}
\tag*{$(***)$}
\]
which are in $W$ term by term (i.e. are in $W_\Gamma\subseteq\Fl(\underline{\Hom}(\Gamma^\circ,\Cat))$). Explicitly, at the vertices $i=0,1,2$, resp. $4$ of $\Gamma$, $(***)$ is none other than
\[
D_i =
\left(
\begin{tikzcd}[column sep=small]
 & \Sigma(X_i) \arrow[dl, "p_{X_i}"'] \arrow[dr, "q_{X_i}"] & \\
X_i & & X_i^\circ
\end{tikzcd}
\right)
\qquad \text{resp.} \qquad
\begin{tikzcd}[column sep=small]
 & \Sigma(Y) \arrow[dl, "p_Y"'] \arrow[dr, "q_Y"] & \\
Y & & Y^\circ,
\end{tikzcd}
\]
at vertex 3 it is
\[
\begin{tikzcd}[column sep=small]
 & \int\Sigma(X_*) \arrow[dl, "\int p_{X_*}"'] \arrow[dr, "\int q_{X_*}"] & \\
\sX=\int X_* & & \int X_*^\circ=\tsX^\circ,
\end{tikzcd}
\]
obtained by integration of the arrows
\[
\begin{tikzcd}[column sep=small]
 & \Sigma(X_*) \arrow[dl, "p_{X_*}"'] \arrow[dr, "q_{X_*}"] & \\
X_* & & X_*^\circ.
\end{tikzcd}
\]
\msp{page 140}
Looking at the restriction of the two morphisms of diagrams $(***)$ to the parts of the three diagrams of type $\Gamma^\circ$ over the tail $\{3,4\}$ of $\Gamma$ we find the following diagram in solid lines
\[
\begin{tikzcd}[column sep=3em, row sep=normal]
 & \int\Sigma(X_*) \arrow[dl, "W"'] \arrow[dr, "W"] \arrow[dd, "\varphi_{\Sigma(Q)}"'] & & \Sigma\tsX \arrow[dl, dotted, "W"'] \arrow[dr, dotted, "W"] \arrow[dd, dotted, "\Sigma\tilde{\varphi}_Q"] & \\
\sX \arrow[dd, "\varphi_Q"'] & & \tsX^\circ \arrow[dd, "\varphi_{Q^\circ}=(\tilde{\varphi}_Q)^\circ"'] & & \tsX \arrow[dd, dotted, "\tilde{\varphi}_Q"] \\
 & \Sigma(Y) \arrow[dl, "W"'] \arrow[dr, "W"] & & \Sigma(Y) \arrow[dl, dotted, "W"'] \arrow[dr, dotted, "W"] & \\
Y & & Y^\circ & & Y
\end{tikzcd}
\]
\printedpage{96}
% head: A. Grothendieck : Dérivateurs, Ch. XII, Édition des Universités de Montpellier II et Paris VII
% foot: Version 7 décembre 2002
which I have completed with a second panel in dotted lines. In this diagram of four adjoining squares, the horizontal sides of the squares represent arrows of $\Cat$ which are in $W$. We thus find: % [sic] qu j'ai -> que j'ai, missing letter
\[
\varphi_Q\in W \iff \varphi_{\Sigma Q}\in W \iff (\varphi_{Q^\circ}=\tilde{\varphi}_Q^\circ)\in W \iff \Sigma\tilde{\varphi}_Q\in W \iff \tilde{\varphi}_Q\in W.
\]
We also find that the five vertical arrows are isomorphic, as arrows of $\Hot_W$. To summarize:

\textbf{Proposition 14.} \textit{Let us consider a diagram (commutative square)}
\[
Q \quad = \qquad
\begin{tikzcd}[row sep=small, column sep=small]
 & X_0 \arrow[dl] \arrow[dr] & \\
X_1 \arrow[dr] & & X_2 \arrow[dl] \\
 & Y &
\end{tikzcd}
\]
\textit{in} $\Cat$\textit{, and the ``opposite'' diagram deduced by} $X\longmapsto X^\circ$
\[
Q^\circ \quad = \qquad
\begin{tikzcd}[row sep=small, column sep=small]
 & X_0^\circ \arrow[dl] \arrow[dr] & \\
X_1^\circ \arrow[dr] & & X_2^\circ \arrow[dl] \\
 & Y^\circ, &
\end{tikzcd}
\]
\msp{page 141}
\textit{and the corresponding arrows}
\[
\left\{
\begin{array}{rlclcl}
\varphi_Q : & \sX_Q & = &
\left(\displaystyle\int
\begin{tikzcd}[ampersand replacement=\&, row sep=small, column sep=small]
 \& X_1 \\
X_0 \arrow[ur] \arrow[dr] \& \\
 \& X_2
\end{tikzcd}
\right)
& \longrightarrow & Y \\[2ex]
\varphi_{Q^\circ} : & \sX_{Q^\circ} & = &
\left(\displaystyle\int
\begin{tikzcd}[ampersand replacement=\&, row sep=small, column sep=small]
 \& X_1^\circ \\
X_0^\circ \arrow[ur] \arrow[dr] \& \\
 \& X_2^\circ
\end{tikzcd}
\right)
& \longrightarrow & Y^\circ.
\end{array}
\right.
\]
\textit{Finally, let us also consider the canonical arrows $\tilde{\varphi}_Q$, $\tilde{\varphi}_{Q^\circ}$, relative to the proper integral $\int_{\mathrm{pr}}$} % [sic] relative -> relatives, agreement with « les flèches »
\printedpage{97}
% head: A. Grothendieck : Dérivateurs, Ch. XII, Édition des Universités de Montpellier II et Paris VII
% foot: Version 7 décembre 2002
\textit{(instead of the smooth integral $\int=\int_{\mathrm{lis}}$)}
\[
\left\{
\begin{array}{rlclcl}
\tilde{\varphi}_Q : & \tsX_Q & = &
\left(\displaystyle\int_{\mathrm{pr}}
\begin{tikzcd}[ampersand replacement=\&, row sep=small, column sep=small]
 \& X_1 \\
X_0 \arrow[ur] \arrow[dr] \& \\
 \& X_2
\end{tikzcd}
\right)
& \longrightarrow & Y \\[2ex]
\tilde{\varphi}_{Q^\circ} : & \tsX_{Q^\circ} & = &
\left(\displaystyle\int_{\mathrm{pr}}
\begin{tikzcd}[ampersand replacement=\&, row sep=small, column sep=small]
 \& X_1^\circ \\
X_0^\circ \arrow[ur] \arrow[dr] \& \\
 \& X_2^\circ
\end{tikzcd}
\right)
& \longrightarrow & Y^\circ.
\end{array}
\right.
\]
\textit{We then have the following:}

\begin{enumerate}[label={}, leftmargin=2.5em]\itshape
\item[\circled{\textit{a}}] $\varphi_{Q^\circ}\simeq(\tilde{\varphi}_Q)^\circ$, $\varphi_Q\simeq(\tilde{\varphi}_{Q^\circ})^\circ$ (from which it already follows that $\left\{\begin{array}{lcl}\varphi_Q\in W & \Longleftrightarrow & \tilde{\varphi}_{Q^\circ}\in W\\ \varphi_{Q^\circ}\in W & \Longleftrightarrow & \tilde{\varphi}_Q\in W\end{array}\right.$ (\footnote[90]{\upshape and that the arrows $\varphi_Q$ and $\tilde{\varphi}_{Q^\circ}$ on the one hand, $\varphi_{Q^\circ}$ and $\tilde{\varphi}_Q$ on the other hand, are isomorphic as arrows in $\Hot_W$.})).

\item[\circled{\textit{b}}] The relations
\[
\varphi_Q\in W, \qquad \varphi_{Q^\circ}\in W, \qquad \tilde{\varphi}_Q\in W, \qquad \tilde{\varphi}_{Q^\circ}\in W
\]
are equivalent, in other words: for the square to be $W$-homotopically cocartesian (smooth integral version), it is necessary and sufficient that the opposite square $Q^\circ$ be so, and these conditions are also equivalent to the corresponding conditions in terms of the proper integrals (instead of the smooth integrals).

\item[\circled{\textit{c}}] The four arrows $\varphi_Q$, $\varphi_{Q^\circ}$, $\tilde{\varphi}_Q$, $\tilde{\varphi}_{Q^\circ}$ of $\Cat$ as arrows in $\Hot_W$, are isomorphic in $\underline{\Fl}(\Hot_W)$. More precisely, as arrows of $\Cat$,
\msp{page 142}
they define isomorphic objects of $\HOT_W(\boldsymbol{\Delta}^1)$.

\item[\circled{\textit{d}}] More precisely still, the four diagrams in $\Cat$ of type $\Gamma^\circ$, where
\[
\Gamma = \;
\begin{tikzcd}[row sep=tiny, column sep=small]
 & 1 \arrow[dl, "\alpha_1"'] & & \\
0 & & 3 \arrow[ul, "i_1"'] \arrow[dl, "i_2"] & 4, \arrow[l, "\varphi"'] \\
 & 2 \arrow[ul, "\alpha_2"] & &
\end{tikzcd}
\]
$\alpha_1i_1\varphi=\alpha_2i_2\varphi$, into which the four arrows we have just considered fit, namely
\[
\widehat{Q} = 
\begin{tikzcd}[row sep=tiny, column sep=small]
 & X_1 \arrow[dr, "i_1"] & & \\
X_0 \arrow[ur, "\alpha_1"] \arrow[dr, "\alpha_2"'] \arrow[rr, dotted] & & \sX_Q \arrow[r, "\varphi_Q"] & Y, \\
 & X_2 \arrow[ur, "i_2"'] & &
\end{tikzcd}
\]
\end{enumerate}


\endgroup
% END UNIT Der-ChXIIb-012/Der-ChXIIb-012_p051-055.tex
% BEGIN UNIT Der-ChXIIb-012/Der-ChXIIb-012_p056-060.tex
\begingroup\makeatletter%
\newcommand{\printedpage}[1]{\marginpar{\footnotesize p.~#1}}%
\newcommand{\msp}[1]{\marginpar{\footnotesize [#1]}}%
\newcommand{\Cat}{\mathrm{Cat}}%
\newcommand{\uHom}{\underline{\mathrm{Hom}}}%
\newcommand{\HOT}{\mathrm{HOT}}%
\newcommand{\Hot}{\mathrm{Hot}}%
\newcommand{\tld}{^{\tilde{}}}%
\makeatother\let\printedpage\derpp\let\msp\dermsp%
\let\mspin\DERmspin%


\setcounter{footnote}{90}

\printedpage{98}
% head: A. Grothendieck : Dérivateurs, Ch. XII, Édition des Universités de Montpellier II et Paris VII
% foot: Version 7 décembre 2002
\[
\begin{tikzcd}[row sep=small, column sep=small]
 & X_1^\circ \arrow[dr,"i^{\prime}_{1}"] & & \\
\widehat{Q^\circ} = X_0^\circ \arrow[ur,"\alpha_1^\circ"] \arrow[dr,"\alpha_2^\circ"'] \arrow[rr,dotted] & & \mathcal{X}_{Q^\circ} \arrow[r] & Y^\circ, \\
 & X_2^\circ \arrow[ur,"i^{\prime}_{2}"'] & &
\end{tikzcd}
\]
\[
\begin{tikzcd}[row sep=small, column sep=small]
 & X_1 \arrow[dr] & & \\
\widehat{Q}\tld = X_0 \arrow[ur] \arrow[dr] \arrow[rr,dotted] & & \tilde{\mathcal{X}}_{Q} \arrow[r] & Y, \\
 & X_2 \arrow[ur] & &
\end{tikzcd}
\]
\[
\begin{tikzcd}[row sep=small, column sep=small]
 & X_1^\circ \arrow[dr] & & \\
\widehat{Q^\circ}\tld = X_0^\circ \arrow[ur] \arrow[dr] \arrow[rr,dotted] & & \tilde{\mathcal{X}}_{Q^\circ} \arrow[r] & Y^\circ, \\
 & X_2^\circ \arrow[ur] & &
\end{tikzcd}
\]
(\footnote{\textbf{NB} We should still complete these (noncommutative) diagrams with the dotted arrows \dots}) \textit{are isomorphic in} $\Hot_W(\Gamma)$\textit{. More precisely still, we have the isomorphisms in} $\uHom(\Gamma^\circ,\Cat)$
\[
(\widehat{Q}\tld)^\circ \simeq \widehat{Q^\circ},
\]
\textit{whence, applying it to $Q^\circ$ instead of $Q$,}
\[
(\widehat{Q^\circ}\tld)^\circ \simeq \widehat{Q},
\]
\textit{which implies (via the canonical isomorphisms valid in} $\HOT_W(S)$ \textit{for every object $S$ in} $\Cat$ \textit{and every} $X_* \in \mathrm{Ob}\,\uHom(S^\circ,\Cat)$\textit{,} $X_* \simeq X_*^\circ$ (\footnote{via the functor $\Sigma$ used above \dots})\textit{), that we have the canonical isomorphisms}
\DERmark{ENV0766}

\[
\left.\begin{aligned}
\widehat{Q} &\simeq \widehat{Q^\circ}\tld \\
\widehat{Q^\circ} &\simeq \widehat{Q}\tld
\end{aligned}\right\}
\qquad \text{\textit{in} } \HOT_W(\Gamma),
\]
\textit{but also directly}
\[
\widehat{Q} \simeq \widehat{Q^\circ}, \qquad \widehat{Q}\tld \simeq \widehat{Q^\circ}\tld \qquad \text{\textit{in} } \HOT_W(\Gamma)
\]
\msp{page 143}\textit{(and it remains to check that we thus have a} transitive \textit{system of canonical isomorphisms in} $\HOT_W(\Gamma)$\textit{, between the four objects considered, which must be essentially sorital \dots).}

\subsection*{Remarks.}

\begin{enumerate}
\item[1)] What we have just said in the case of diagrams in Cat of type $\Phi^\circ = \Psi$, seems still to make sense for diagrams of arbitrary type
\[
X_* : S^\circ \longrightarrow \Cat, \qquad s \longmapsto X_s
\]
\printedpage{99}
% head: A. Grothendieck : Dérivateurs, Ch. XII, Édition des Universités de Montpellier II et Paris VII
% foot: Version 7 décembre 2002
($S$ an object of Cat). We can integrate it in a smooth way, $\int X_* = \int_{\mathrm{lis}} X_*$, obtaining a fibered category (split) over $S$, or in a proper way, $\int_{\mathrm{pr}} X_*$, obtaining a split cofibered category over $S^\circ$. We can also consider $X_*^\circ$, the composite of $X_*$ with the automorphism $Y \longmapsto Y^\circ$ in Cat, and integrate it in either a proper or a smooth way. In all these cases, we find a diagram of type $\widehat{S}^\circ$, where $\widehat{S}$ is obtained from $S$ by adding an element which maps by a ``canonical'' (given) arrow to each of the others (\textit{without} commutation relations!) A
\msp{page 144}\textit{direct system} of morphisms in Cat
\[
f \qquad X_s \xrightarrow{f_s} Y
\]
then gives a morphism
\[
\int_{\mathrm{lis}} X_* \xrightarrow{\varphi_Q} Y,
\]
where $Q$ denotes the diagram, of type $\widehat{S}^\circ$, made up of the $X_s$, $Y$, the $u^* : X_t \longrightarrow X_s$ (for $u : s \longrightarrow t$ in $S$) and the $f_s$. Likewise we will have
\[
\int_{\mathrm{prop}} X_* \xrightarrow{\tilde{\varphi}_Q} Y,
\]
and if
\[
\Gamma = (\widehat{S} + \text{initial object}),
\]
we find diagrams $\widehat{Q}$, $\widehat{Q}\tld$ of type $\Gamma^\circ$ in Cat. Unless I am mistaken, the relations between $\widehat{Q}$, $\widehat{Q}\tld$, $\widehat{Q^\circ}$, $\widehat{Q^\circ}\tld$ are the same as those made explicit in proposition 14 above. I am not going to dwell on this with the requisite care, unless needed. We will say that a diagram $Q$ defines $Y$ as a $\varinjlim\limits_{I}$ that is $W$-homotopic \texttt{[where} $I = S^\circ$\texttt{]} if $\varphi_Q \in W$. It amounts to the same thing to say that $\varphi_{Q^\circ} \in W$, or that $\tilde{\varphi}_Q \in W$ or that $\tilde{\varphi}_{Q^\circ} \in W$ \dots
\msp{page 145}

\item[2)] The ``universal'' case, for
\begin{tikzcd}[baseline=(\tikzcdmatrixname-2-1.base), row sep=small, column sep=small]
 & X_1 \\
X_0 \arrow[ur] \arrow[dr] & \\
 & X_2
\end{tikzcd}
given, of a commutative square $Q = {}$%
\begin{tikzcd}[baseline=(\tikzcdmatrixname-2-1.base), row sep=tiny, column sep=small]
 & X_1 \arrow[dr] & \\
X_0 \arrow[ur] \arrow[dr] & & Y \\
 & X_2 \arrow[ur] &
\end{tikzcd}
in Cat, giving rise to
\[
\begin{tikzcd}[row sep=small, column sep=small]
 & & & X_1 & & & \\
\varphi_Q : &[1em] \left(\vrule height 6ex depth 6ex width 0pt\right. &[-1.5em] \int X_0 \arrow[ur] \arrow[dr] & & \left.\vrule height 6ex depth 6ex width 0pt\right) &[2em] \longrightarrow &[2em] Y, \\
 & & & X_2 & & &
\end{tikzcd}
\]
is that of $Y = X_1 \amalg_{X_0} X_2$, giving rise to % [sic] naissace -> naissance, letter n omitted
\printedpage{100}
% head: A. Grothendieck : Dérivateurs, Ch. XII, Édition des Universités de Montpellier II et Paris VII
% foot: Version 7 décembre 2002
\[
\begin{tikzcd}[row sep=small, column sep=small]
 & X_1 \arrow[dr] & \\
Q_0 = Q(X_*) = X_0 \arrow[ur] \arrow[dr] & & X_1 \amalg_{X_0} X_2. \\
 & X_2 \arrow[ur] &
\end{tikzcd}
\]

The general case reduces to this one, for we will have a commutative diagram
\[
\begin{tikzcd}[row sep=small, column sep=small]
 & & X_1 & & & \\
\left(\vrule height 6ex depth 6ex width 0pt\right. &[-1.5em] \int X_0 \arrow[ur] \arrow[dr] & & \left.\vrule height 6ex depth 6ex width 0pt\right) \arrow[rr,"{\varphi_{Q_0}=\varphi_{Q(X_*)}}"] \arrow[ddrr,"\varphi_Q"] &[4em] \hphantom{MMMM} &[4em] X_1 \amalg_{X_0} X_2 \arrow[dd,"\alpha_Q"] \\
 & & X_2 & & & \\
 & & & & & Y.
\end{tikzcd}
\]
\end{enumerate}

But there must be interesting cases where $\varphi_{Q_0}$ is \textit{not} in $W$, i.e. the Cat-cocartesian square
\begin{tikzcd}[baseline=(\tikzcdmatrixname-2-1.base), row sep=tiny, column sep=small]
 & X_1 \arrow[dr] & \\
X_0 \arrow[ur] \arrow[dr] & & X_1 \amalg_{X_0} X_2 \\
 & X_2 \arrow[ur] &
\end{tikzcd}
is not $W$-homotopically cocartesian, and where nevertheless we can find a commutative square $Q$
\msp{page 146}completing $X_*$, in such a way that this square is $W$-homotopically cocartesian \texttt{[rather cocartesian]}, i.e. that $\varphi_Q = \alpha_Q \varphi_{Q_0}$ is in $W$. This is all the more plausible since in Cat, amalgamated sums tend not to have good homotopy properties, even if one of the arrows $\alpha_1$, $\alpha_2$ is \texttt{[a]} monomorphism (or even an open immersion), contrary to what happens in $\boldsymbol{\Delta}^\wedge$, and more generally in categories of the form $Z^\wedge$ ($Z$ in Cat).

\subsection*{3.4\quad Criteria for homotopically cocartesian squares.}

The criteria I have obtained are all contained in the following statement:
\DERmark{ENV0767}


\textbf{Theorem 4} (\footnote{\textbf{NB} In this theorem, and corollaries 1, 2, 3, we need only the \textit{weak form} of W(7 bis), namely $i_0 \in W \Longrightarrow i \in W$ (cf. p. 79) (Loc(6$'$a) in the notation of XVI).},\footnote{\textbf{NB} All in all, the theorem (forward part) holds without any hypothesis W(7 bis) or weak W(7 bis) on $W$, and the same is true of its corollaries 1 to 3 (to the exclusion of corollaries 4, 5, which do use W(7 bis) in its full strength).}). \textit{Let us consider a commutative square diagram in} Cat
\[
Q \quad = \quad
\begin{tikzcd}
X_0 \arrow[r,"\alpha_1"] \arrow[d,"\alpha_2"'] \arrow[dr,dotted,"f_0"] & X_1 \arrow[d,"f_1"] \\
X_2 \arrow[r,"f_2"'] & Y,
\end{tikzcd}
\tag*{$(*)$}
\]
\printedpage{101}
% head: A. Grothendieck : Dérivateurs, Ch. XII, Édition des Universités de Montpellier II et Paris VII
% foot: Version 7 décembre 2002
\textit{and consider the canonical morphism}
\[
\begin{tikzcd}[row sep=small, column sep=small]
 & & X_1 & & & \\
\varphi_Q : \quad \mathcal{X} = \Bigl(\int &[-1.2em] X_0 \arrow[ur,"\alpha_1"] \arrow[dr,"\alpha_2"'] & & \Bigr) &[2em] \longrightarrow &[2em] Y, \\
 & & X_2 & & &
\end{tikzcd}
\]
\msp{page 147}\textit{the integral being the integral which is} smooth. \textit{(There is a dual statement concerning the proper integral of}\newline
\begin{tikzcd}[baseline=(\tikzcdmatrixname-2-1.base), row sep=small, column sep=small]
 & X_1 \\
X_0 \arrow[ur] \arrow[dr] & \\
 & X_2
\end{tikzcd}
\textit{.)}

\begin{enumerate}
\item[\textit{a)}] \textit{For $\varphi_Q$ to be $W$-smooth, it is necessary and sufficient that $f_1$, $f_2$ and $f_0 = f_1\alpha_1 = f_2\alpha_2$ be so.} (Assuming W(7 bis) (\footnote{For the forward part ($f_0$, $f_1$, $f_2$ smooth $\Longrightarrow \varphi_Q$ smooth), we need only W(7 bis) (\footnotemark), weak form. It is only for the converse (which has never yet been of use to me) that we need W(7 bis), in its full strength.}\footnotetext{We do not even need W(7 bis) in weak form, cf. annotation on page 156. We use a restrictive hypothesis on $W$ only in the converse ($\varphi_Q$ smooth $\Longrightarrow f_0$, $f_1$, $f_2$ smooth).})).

\item[\textit{b)}] \textit{The fiber of $\varphi_Q$ at an element $y$ of $Y$ is isomorphic to}
\[
\begin{tikzcd}[row sep=small, column sep=small]
 & & X_{1\,y} & \\
\int &[-1.2em] X_{0\,y} \arrow[ur] \arrow[dr] & &[2em] (\simeq \mathcal{X}_y). \\
 & & X_{2\,y} &
\end{tikzcd}
\tag*{$(**)$}
\]
\end{enumerate}

\textbf{Corollary 1.} (\footnote{Valid without any hypothesis on $W$ (other than Loc(1,2,3 bis)).}) \textit{Suppose $f_1$, $f_2$ and $f_1\alpha_1 = f_2\alpha_2$ $W$-smooth, and the integrated categories $(**)$ $W$-aspheric for every $y \in \mathrm{Ob}\,Y$. (NB This is still the smooth integral, but it does not really matter, for we have seen (cf. the upper part of the diagram p. 140, which implies that $\mathcal{X}$} \texttt{[is]} \textit{$W$-aspheric if and only if $\tilde{X}$ is) that one is $W$-aspheric if and only if the other is.) Then $\varphi_Q \in W$, i.e. the square $(*)$ is $W$-homotopically cocartesian} \texttt{[rather cocartesian]}\textit{. By duality, same conclusion if we suppose $f_1$, $f_2$, $f_1\alpha_1 = f_2\alpha_2$} proper \textit{instead of smooth. (For $Q$ $W$-homotopically cocartesian} \texttt{[rather cocartesian]}\textit{, if and only if $Q^\circ$ is, by} \texttt{[the]} \textit{prop. 14.)}
\DERmark{ENV0768}\DERmark{ENV0769}


\msp{page 148}Indeed, a morphism in Cat which is $W$-smooth and has $W$-aspheric fibers is in $W$ (and it is even a universal $W$-equivalence \dots).

\textbf{Corollary 2}. \textit{We assume that $W$ satisfies} W(7 bis) (cocartesian square axiom) (\footnote{unnecessary.}).
\printedpage{102}
% head: A. Grothendieck : Dérivateurs, Ch. XII, Édition des Universités de Montpellier II et Paris VII
% foot: Version 7 décembre 2002
\textit{Suppose the square $Q$ cocartesian in} Cat (\footnote{NB We do not really use that $Q$ is cocartesian, rather that $Q$ is \textit{cartesian,} and $f_2$ (hence also $\alpha_1$) an open immersion (resp. closed), and that for $y \in Y\backslash X_2$, $X_{1\,y}$ is $W$-aspheric. (In the case where $Q$ is cocartesian, $X_{1\,y}$ is $\simeq e$.)})\textit{, and that $\alpha_1$ is an open immersion (resp. closed) and $f_1$ smooth (resp. proper). Then $Q$ is $W$-homotopically cocartesian.} (\footnote{Here we need W(7 bis), forward form, i.e. Loc(6$'$a) (notation of XVI).})

It suffices to establish the non-resp'ed result, the other will follow from it by duality (since $Q$ is $W$-cocartesian if and only if $Q^\circ$ is). By hypothesis $f_1$ is $W$-smooth, on the other hand $\alpha_1$ and $f_2$ are open immersions, hence $W$-smooth, hence $f_1\alpha_1$ is $W$-smooth too. (Besides, $\alpha_2$ too, for the square is also cartesian, hence $\alpha_2$ obtained from $f_1$ (which is $W$-smooth by hypothesis) by base change $X_2 \longrightarrow Y$.) Thus $\varphi_Q$ is \texttt{[}$W$\texttt{-]}smooth, and it suffices to see that its fibers are $W$-aspheric. The calculation $(**)$ gives, for $y$ in $Y\backslash X_2$,
\[
\begin{tikzcd}[row sep=small, column sep=small]
 & & e & \\
\mathcal{X}_y \simeq \quad \int &[-1.2em] \emptyset \arrow[ur] \arrow[dr] & &[1em] \simeq e, \qquad \text{O.K.} \\
 & & \emptyset &
\end{tikzcd}
\]
And for $y$ in $X_2$
\[
\begin{tikzcd}[row sep=small, column sep=small]
 & & X_{1\,y} & \\
\mathcal{X}_y \simeq \int &[-1.2em] X_{0\,y} \arrow[ur,"\sim"] \arrow[dr] & &[2em] . \\
 & & e &
\end{tikzcd}
\]
\msp{page 149}Since $X_{0\,y} \longrightarrow X_{1\,y}$ is an isomorphism, hence is in $W$, it follows that $e \longrightarrow \mathcal{X}_y$ is too (prop. 13, p. 125), assuming W(7 bis), hence $\mathcal{X}_y$ is $W$-aspheric, O.K. (\footnote{We need only the forward form of W(7 bis), i.e. (notation of XVI) Loc(6$'$a).}).

\textbf{Corollary 3.} (Assuming W(7 bis) (\footnote{unnecessary.}).) \textit{Let $X$ be in} Cat\textit{, $U_1$, $U_2$ two open subcategories of $X$ such that $X = U_1 \cup U_2$, $U_0 = U_1 \cap U_2$ the open intersection. Then the square of open immersions}
\[
\begin{tikzcd}[row sep=large, column sep=large]
U_0 \arrow[r,hook] \arrow[d,hook] & U_1 \arrow[d,hook] \\
U_2 \arrow[r,hook] & X
\end{tikzcd}
\]
\textit{is $W$-cocartesian. (Dual statement for two closed subcategories $T_1$, $T_2$ such that $T_1 \cup T_2 = X$, with $T_0 = T_1 \cap T_2$.)}

We apply corollary 2. We know that the square is cocartesian in Cat (\footnote{We do not even need to know this, cf. marginal annotation to corollary 2.}), moreover $U_1 \lhook\joinrel\longrightarrow X$ is indeed $W$-smooth, being an open immersion.


\endgroup
% END UNIT Der-ChXIIb-012/Der-ChXIIb-012_p056-060.tex
% BEGIN UNIT Der-ChXIIb-012/Der-ChXIIb-012_p061-065.tex
\begingroup\makeatletter%
\newcommand{\printedpage}[1]{\marginpar{\footnotesize p.~#1}}%
\newcommand{\msp}[1]{\marginpar{\footnotesize [#1]}}%
\newcommand{\hto}[1][]{\overset{#1}{\lhook\joinrel\longrightarrow}}%
\newcommand{\circled}[1]{\tikz[baseline=(c.base)]\node[draw,circle,inner sep=1pt](c){#1};}%
\newcommand{\intdiag}[3]{\begin{tikzcd}[row sep=small,column sep=small,ampersand replacement=\&,baseline=(\tikzcdmatrixname-2-1.base)]%
 \& #2 \\%
 \displaystyle\int #1 \arrow[ur] \arrow[dr] \& \\%
 \& #3%
\end{tikzcd}}%
\makeatother\let\printedpage\derpp\let\msp\dermsp%
% Shared placement decision BUILD-MARGIN-EN-XIIb-152
\def\DERmargintarget{page 152}%
\def\msp#1{\def\DERmarginarg{#1}\ifx\DERmarginarg\DERmargintarget\marginpar{\raggedright\footnotesize\ttfamily\hspace*{20pt}[#1]}\else\dermsp{#1}\fi}%
\let\mspin\DERmspin%


\printedpage{103}
% head: A. Grothendieck : Dérivateurs, Ch. XII, Édition des Universités de Montpellier II et Paris VII

\noindent\textbf{Corollary 4.} (Assuming W(7 bis) (\footnote[104]{Here we need W(7 bis) in strong form, i.e. $i_0\in W \iff i\in W$ (cf. annotation th. 4).}).) \textit{Let $U \hto[i] X$ be an open immersion such that $T=X\backslash U$ is isomorphic to the one-point category $e$. For $i\in W$ to hold, it is necessary and sufficient that $i$ be $W$-aspheric, i.e. that $U/t$ (where $t$ is the unique object of $T$) be $W$-aspheric.}

It is of course enough to prove that
\[
i\in W \Longrightarrow U/t \ W\text{-aspheric.}
\]

\msp{page 150}
Consider the cartesian square
\[
\begin{tikzcd}[row sep=large,column sep=large]
U/t \arrow[r, hook, "i/t"] \arrow[d, "p"'] & X/t \arrow[d, "q"] \\
U \arrow[r, hook, "i"] & X,
\end{tikzcd}
\]
we are going to prove that it is $W$-cocartesian. Assuming W(7 bis), it will follow that $i\in W$ implies $i/t\in W$, and since $X/t$ is $W$-aspheric (having a final object), this means that $U/t$ is $W$-aspheric, which is what we want to prove.

It is enough, by virtue of corollary 2, to prove that the square is cocartesian in Cat, since $i/t$ is an open immersion ($i$ being one), and $X/t\longrightarrow X$ is a local isomorphism, hence $W$-smooth. We see that $(X/t)\backslash(U/t)\simeq e$ (one-point category), moreover the unique element of $(X/t)\backslash(U/t)$ is the final element $t_0 : t\xrightarrow{\mathrm{id}_t} t$ of $X/t$. Thus $H_{t_0}\simeq e_{U/t}$ (final presheaf), and $p_{!}(H_{t_0})\simeq H'_t \overset{\mathrm{def}}{=} (u\longmapsto \mathrm{Hom}_X(u,t))$, which characterizes the amalgamated sum in Cat
\DERmark{ENV0770}

\dots

Corollary 4 is none other than the unproved corollary of proposition 7 (p. 85).

Here is a variant of corollary 4.

\noindent\textbf{Corollary 5.} (Assuming W(7 bis) (\footnote[105]{strong form.}).) \textit{Let $i : U \hto X$ be an open immersion, let $t_0\in\mathrm{Ob}\,T$ be a \textup{closed} point of $T=X\backslash U$, such that $\{t_0\}$}
\msp{page 151}
\textit{is isomorphic to the one-point category, i.e. $\mathrm{Hom}_X(t_0,t_0)=\{\mathrm{id}\}$. Suppose that for every $t\in T\backslash\{t_0\}$, $U/t$ is $W$-aspheric. For $i\in W$ to hold, it is necessary and sufficient that $i$ be aspheric, i.e. that $U/t_0$ also be $W$-aspheric.}

We must prove that $i\in W$ implies that $U/t_0$ is $W$-aspheric. Consider the inclusion
\[
i_0 : U\longrightarrow U' \overset{\mathrm{def}}{=} X\backslash\{t_0\},
\]
\DERmark{ENV0771}

the assumption on the $U/t$, for $t\in T\backslash\{t_0\}$, means that $i_0$ is $W$-aspheric. Thus $i_0$ is in $W$, and since the composite $i : U\xrightarrow{i_0} U'\xrightarrow{i'} X$ is in $W$, this implies that
% foot: Version 7 décembre 2002
\printedpage{104}
% head: A. Grothendieck : Dérivateurs, Ch. XII, Édition des Universités de Montpellier II et Paris VII
$i' : U'\longrightarrow X$ is in $W$. We are now in the situation of corollary 4, which implies that $i'$ is $W$-aspheric. Thus $i=i'i_0$ is $W$-aspheric, as a composite of two $W$-aspheric morphisms, q.e.d.

\noindent\textsc{Proof of th. 4.}

\noindent\circled{a} Smoothness criterion for $\varphi_Q : \mathcal{X} = \intdiag{X_0}{X_1}{X_2} \longrightarrow Y$.

Since $\{1\}$, $\{2\}$ are opens of $\Phi = \begin{tikzcd}[row sep=small,column sep=small,baseline=(\tikzcdmatrixname-2-1.base)]
 & 1 \arrow[dl] \\
 0 & \\
 & 2 \arrow[ul]
\end{tikzcd}$, $i_1 : X_1\longrightarrow\mathcal{X}$ and $i_2 : X_2\longrightarrow\mathcal{X}$ are open immersions. Since $\varphi_Q\circ i_1 = f_1$,
\msp{page 152}
$\varphi_Q\circ i_2=f_2$, we already see that
\[
f_1 \text{ and } f_2 \text{ smooth} \iff \varphi_Q \text{ smooth on the open union of } i_1(X_1) \text{ and } i_2(X_2)\text{, i.e. on } \mathcal{X}\backslash i_0(X_0).
\]
It remains to express the smoothness of $\varphi_Q$ at the points $x_0$ of $i_0(X_0)$. Let $x_1$, $x_2$, $y$ be the images of $x_0$ in $X_1$, $X_2$, $Y$ respectively. We must express that the fibers of
\[
(\varphi_Q)/x_0 : \mathcal{X}/x_0\longrightarrow Y/y \tag*{$(*)$}
\]
are aspheric. For this we are going to make $\mathcal{X}/x_0$ and the arrow $(*)$ explicit in terms of the commutative square
\[
(**)\qquad\qquad
\begin{tikzcd}[row sep=large,column sep=large]
X_0/x_0 \arrow[r, "\alpha_1/x_0"] \arrow[d, "\alpha_2/x_0"'] & X_1/x_1 \arrow[d, "f_1/x_1"] \\
X_2/x_2 \arrow[r, "f_2/x_2"] & Y/y,
\end{tikzcd}
\]
induced by $Q$. I say

\noindent\textbf{Lemma 1.} \textit{Let there be a commutative square}
\[
Q \quad = \quad
\begin{tikzcd}[row sep=large,column sep=large]
X_0 \arrow[r, "\alpha_1"] \arrow[d, "\alpha_2"'] & X_1 \arrow[d, "f_1"] \\
X_2 \arrow[r, "f_2"] & Y
\end{tikzcd}
\]
\textit{in} Cat, \textit{$x_0\in\mathrm{Ob}\,X_0$, $x_1$, $x_2$, $y$ its images in $X_1$, $X_2$, $Y$ respectively, $Q/x_0$ the square $(**)$ above, induced by $Q$ on the localized categories $X_0/x_0$ etc., $\mathcal{X} = \intdiag{X_0}{X_1}{X_2}$. We then have a}
% foot: Version 7 décembre 2002
\printedpage{105}
% head: A. Grothendieck : Dérivateurs, Ch. XII, Édition des Universités de Montpellier II et Paris VII
\msp{page 153}
\textit{canonical isomorphism}
\[
\mathcal{X}/x_0 \quad\simeq\quad \intdiag{X_0/x_0}{X_1/x_1}{X_2/x_2} ,
\]
\textit{and the arrow $(\varphi_Q)/x_0$, the localized arrow of $\varphi_Q : \mathcal{X}\longrightarrow Y$ at $x_0$, is identified with the arrow $\varphi_{(Q/x_0)}$ :}
\[
\varphi_{(Q/x_0)} : \quad \intdiag{X_0/x_0}{X_1/x_1}{X_2/x_2} \qquad\longrightarrow\qquad Y/y.
\]

\noindent\textsc{Proof} of the Lemma. We know that
\[
\mathcal{X}/x_0\longrightarrow\mathcal{X}
\]
is Cat-fibered (split) (with discrete fibers), and since
\[
\mathcal{X}\longrightarrow\Phi
\]
is also Cat-fibered (split), it follows that the same is true of the composite arrow
\[
\mathcal{X}/x_0\longrightarrow\Phi.
\]
Thus to ``calculate'' $\mathcal{X}/x_0$, as a category over $\Phi$, we must make explicit its three fibers, and the two base change functors which connect them. We find at once that these fibers are canonically isomorphic to
\[
X_0/x_0, \qquad X_1/x_1, \qquad X_2/x_2,
\]
and we shall admit (!) that the base change functors
\msp{page 154}
\[
X_0/x_0\longrightarrow X_1/x_1, \qquad X_0/x_0\longrightarrow X_2/x_2
\]
(defined up to unique isomorphism) are none other than the localized arrows $\alpha_1/x_0$ and $\alpha_2/x_0$ of $\alpha_1 : X_0\longrightarrow X_1$ and $\alpha_2 : X_0\longrightarrow X_2$. This therefore shows that we have an isomorphism
\[
\mathcal{X}/x_0 \quad\simeq\quad \intdiag{X_0/x_0}{X_1/x_1}{X_2/x_2} ,
\]
% foot: Version 7 décembre 2002
\printedpage{106}
% head: A. Grothendieck : Dérivateurs, Ch. XII, Édition des Universités de Montpellier II et Paris VII
and we shall also admit that the arrow $(\varphi_Q)/x_0 : \mathcal{X}/x_0\longrightarrow Y/y$ is identified with $\varphi_{(Q/x_0)}$.

To express the $W$-smoothness of $\varphi_Q$ at $x_0\in\mathrm{Ob}\,X_0\subseteq\mathrm{Ob}\,\mathcal{X}$, we are thus reduced to expressing the $W$-asphericity of the fibers of
\[
\varphi_{(Q/x_0)} : \quad \intdiag{X_0/x_0}{X_1/x_1}{X_2/x_2} \qquad\longrightarrow\qquad Y/y.
\]
To make these fibers explicit, we need another auxiliary lemma

\noindent\textbf{Lemma 2} (\footnote[106]{NB This is b) of th. 4.}). \textit{Let there be a commutative square in} Cat
\[
Q \quad = \quad
\left(
\begin{tikzcd}[row sep=large,column sep=large]
X_0 \arrow[r, "\alpha_1"] \arrow[d, "\alpha_2"'] & X_1 \arrow[d, "f_1"] \\
X_2 \arrow[r, "f_2"] & Y
\end{tikzcd}
\right),
\]
\textit{whence}
\[
\varphi_Q : \quad \left( \intdiag{X_0}{X_1}{X_2} \right) \qquad\longrightarrow\qquad Y,
\]

\msp{page 155}
\textit{and let $y\in\mathrm{Ob}\,Y$. Then the fiber of $\varphi_Q$ at $y$}
\[
\mathcal{X}_y = \varphi_Q^{-1}(\ \underbrace{\{y\}}_{\text{one-point category}}\ )
\]
\textit{is given by (as a category over $\Phi$)}
\[
\mathcal{X}_y \quad\simeq\quad \intdiag{X_{0\,y}}{X_{1\,y}}{X_{2\,y}} .
\]

\noindent\textsc{Proof.} It is immediate that we have
\[
\mathcal{X}_y\cap\underbrace{\mathcal{X}_s}_{X_s}\simeq(X_s)_y \qquad\quad \text{for } s\in\mathrm{Ob}\,\Phi=\{0,1,2\},
\]
% foot: Version 7 décembre 2002
\printedpage{107}
% head: A. Grothendieck : Dérivateurs, Ch. XII, Édition des Universités de Montpellier II et Paris VII
and (chance does things well) the base change functors
\[
X_0\longrightarrow X_1, \qquad X_0\longrightarrow X_2
\]
send the (nonfull) subcategory $X_{0\,y}$ of $X_0$ into the (nonfull) subcategories $X_{1\,y}$ and $X_{2\,y}$ of $X_1$ and $X_2$ respectively.
\[
\begin{tikzcd}[row sep=large,column sep=large]
X_{s\,y} \arrow[r] \arrow[d] \arrow[dr, phantom, "\scriptstyle\text{cart.}"] & \mathcal{X}_y \arrow[r] \arrow[d] \arrow[dr, phantom, "\scriptstyle\text{cart.}"] & \{y\} \arrow[d] \\
X_s \arrow[r] \arrow[d] \arrow[dr, phantom, "\scriptstyle\text{cart.}"] & \mathcal{X} \arrow[r] \arrow[d] & Y \\
\{s\} \arrow[r] & \Phi &
\end{tikzcd}
\]
We thus have a fibered subcategory of $\mathcal{X}$, subject to checking that the strictly cartesian arrows
\[
\alpha_1(x_0)\longrightarrow x_0, \qquad \alpha_2(x_0)\longrightarrow x_0
\]
\DERmark{ENV0772}

in $\mathcal{X}$, for $x_0$ in $X_{0\,y}$, are in $\mathcal{X}_y$, i.e. are over the identity arrow of $y$. Now that is the essential property
\msp{page 156}
of the functor $\varphi_Q : \mathcal{X}\longrightarrow Y$, which we forgot to make explicit in 3.2. It really seems to work (like clockwork!)

Let us apply Lemma 2 to the square diagram $Q/x_0$ and to $\varphi_{(Q/x_0)}$ (page 154), we thus see that the fibers to be examined at the $u\in\mathrm{Ob}\,Y/y$ ($u : y'\longrightarrow y$) can be written as the integrals
\[
\mathcal{X}(x_0,u) \quad\overset{\mathrm{def}}{=}\quad \intdiag{(X_0/x_0)_u}{(X_1/x_1)_u}{(X_2/x_2)_u} .
\]
\DERmark{ENV0773}

Let us now use the fact that $f_1$, $f_2$ are $W$-smooth (expressing that $\varphi_Q$ is $W$-smooth in $\mathcal{X}\backslash X_0$), hence that $(X_1/x_1)_u$ and $(X_2/x_2)_u$ are $W$-aspheric. In the $W$-cocartesian square
\[
\begin{tikzcd}[row sep=large,column sep=small]
 & (X_1/x_1)_u \arrow[dr, "{i_1(x_0,u)}"] & \\
(X_0/x_0)_u \arrow[ur, "{\alpha_1(x_0,u)}"] \arrow[dr, "{\alpha_2(x_0,u)}"'] & & (\mathcal{X}/x_0)_u\simeq\mathcal{X}(x_0,u) \\
 & (X_2/x_2)_u \arrow[ur, "{i_2(x_0,u)}"'] &
\end{tikzcd}
\]
% foot: Version 7 décembre 2002


\endgroup
% END UNIT Der-ChXIIb-012/Der-ChXIIb-012_p061-065.tex
% BEGIN UNIT Der-ChXIIb-012/Der-ChXIIb-012_p066-070.tex
\begingroup\makeatletter%
\newcommand{\printedpage}[1]{\marginpar{\footnotesize p.~#1}}%
\newcommand{\msp}[1]{\marginpar{\footnotesize [#1]}}%
\newcommand{\Ob}{\operatorname{Ob}}%
\newcommand{\Fl}{\operatorname{Fl}}%
\newcommand{\Cat}{\mathrm{Cat}}%
\newcommand{\Fibdisc}{\underline{\mathrm{Fib\ disc}}}%
\newcommand{\cercle}[1]{\tikz[baseline=(c.base)]\node[draw,circle,inner sep=1pt](c){#1};}%
\makeatother\let\printedpage\derpp\let\msp\dermsp%
\let\mspin\DERmspin%


\printedpage{108}
% head: A. Grothendieck : Dérivateurs, Ch. XII, Édition des Universités de Montpellier II et Paris VII
(\footnote[107]{careful, this is not a commutative square. Say rather that $(\mathcal{X}/x_0)_u \simeq \int \begin{tikzcd}[ampersand replacement=\&, row sep=0.4em, column sep=0.4em, baseline=(\tikzcdmatrixname-2-1.base)] \& {} \\ {} \ar[ur]\ar[dr] \& \\ \& {} \end{tikzcd}$ is smooth over $\Phi$ with $W$-aspheric fibers (if $f_0$, $f_1$, $f_2$ are smooth), hence $W$-aspheric since $\Phi$ is.}), we must express that $(\mathcal{X}/x_0)_u$ is $W$-aspheric, knowing that $(X_1/x_1)_u$ and $(X_2/x_2)_u$ are. So this is equivalent to % Resolved from PDF leaf 31: the original diagram has exactly these two unlabelled arrows and no vertex labels.
\[
i_1(x_0,u) \in W \qquad (\text{or } i_2(x_0,u) \in W),
\]
hence also (assuming W(7 bis)) to
\msp{page 157}
\[
\alpha_1(x_0,u) \in W, \qquad \text{i.e. } (X_0/x_0)_u \ W\text{-aspheric.}
\]
To say that this is true for every $u$ in $Y/y$, means that $f_0 : X_0 \longrightarrow Y$ is smooth at $x_0$ \texttt{[rather at $x_0$]}.
\DERmark{ENV0774}


So we have proved a), in the more precise form:

\textbf{Corollary 6.} \textit{Let $x \in \Ob \mathcal{X}$ (where $\mathcal{X} = \int \begin{tikzcd}[ampersand replacement=\&, row sep=tiny, column sep=small, baseline=(\tikzcdmatrixname-2-1.base)] \& X_1 \\ X_0 \ar[ur]\ar[dr] \& \\ \& X_2 \end{tikzcd}$.) We want to express the $W$-smoothness of}
\[
\varphi_Q : \mathcal{X} \longrightarrow Y
\]
\textit{at $x$.}

\begin{enumerate}[label={}, leftmargin=2.5em]
\item[\textit{a)}] \textit{If $x \in \Ob X_i$ ($i \in \{1,2\}$), $\varphi_Q$ is $W$-smooth at $x$ if and only if $f_i : X_i \longrightarrow Y$ is.}
\item[\textit{b)}] \textit{Assume} \textup{W(7 bis)} (\footnote[108]{Unnecessary for the forward direction, assuming \texttt{[incomplete sentence]}.})\textit{. Let $x_0 \in \Ob X_0$, and suppose $f_1$, $f_2$ smooth at $x_1$, $x_2$. Then $\varphi_Q$ is $W$-smooth at $x_0$ if and only if $f_0 : X_0 \longrightarrow Y$ is.}
\end{enumerate}

\textit{So} (assuming W(7 bis)) \textit{in all three cases $x \in \Ob X_i$ ($i \in \{0,1,2\}$), under the corresponding hypotheses of a) and b), we find that $\varphi_Q$ is $W$-smooth at $x$ if and only if $f_i : X_i \longrightarrow Y$ is} \texttt{[not really for $i=0$]}.
\DERmark{ENV0775}


\cercle{b} As for b) of theorem 4, it is lemma 2 above. So everything is proved!

\msp{page 158}

\subsection*{3.5\quad Applications to amalgamated sums in a category $X^{\wedge}$ (with $X \in \Ob \Cat$).}

First a characterization of amalgamated sums involving an open immersion:

\textbf{Proposition 15.} \textit{Let}
% foot: Version 7 décembre 2002
\printedpage{109}
% head: A. Grothendieck : Dérivateurs, Ch. XII, Édition des Universités de Montpellier II et Paris VII
\[
Q : \qquad
\begin{tikzcd}
U \ar[r, "i"] \ar[d, "p"'] & X \ar[d, "q"] \\
U' \ar[r, "i'"] & X'
\end{tikzcd}
\]
\textit{be a commutative square in} $\Cat$\textit{. For this square to be cocartesian in} $\Cat$\textit{, with $i$ (hence $i'$) an open immersion, it is necessary and sufficient that the following conditions be satisfied:}

\begin{enumerate}[label={}, leftmargin=2.5em]
\item[\textit{a)}] \textit{$Q$ is cartesian.}
\item[\textit{b)}] \textit{$i'$ is an open immersion (hence $i$ too).}
\item[\textit{c)}] \textit{If $T' = X' \backslash U'$, and $T = X \times_{X'} T' = X \backslash U$, $T \longrightarrow T'$ is an isomorphism.}
\item[\textit{d)}] \textit{For every $t \in T$, if $t'$ is its image in $T'$, the functor induced by $p$}
\[
p/t : \underbrace{U/t}_{= V_t} \longrightarrow \underbrace{U'/t'}_{= V'_{t'}}
\]
\textit{satisfies $(p/t)_!(e_{U/t}) \simeq e_{U'/t'}$, or what amounts to the same, the ``colocal fibers'' of $p/t$, $v' \backslash V_t$, with $v'$ in $V'_{t'}$, are} $0$\textit{-connected.}
\end{enumerate}

\msp{page 159}
Immediate proof - this should rather go back to §1.

\textbf{Corollary 1.} \textit{Let $Q$ be a square that is} cartesian (cf. page 158), \textit{with $i'$ an} open immersion \textit{(hence $i$ too), and $q$} $W$-smooth \textit{(hence $p$ too). We assume that $W$ satisfies the axiom} W(5) \textit{of connectedness} (\footnote[109]{So it suffices $q$ $W_0$-smooth, where $W_0 = \{f \in \Fl \Cat \mid \pi_0(f) \text{ bijective}\}$.})\textit{. Let $T' = X' \backslash U'$, $T = X \backslash U$, suppose that $T \longrightarrow T'$ is an isomorphism. Then $Q$ is cocartesian.}
\DERmark{ENV0776}


It suffices to prove condition d) of prop. 15. Let us consider the square
\[
\begin{tikzcd}
U/t \ar[r, hook, "i/t"] \ar[d, "p/t"'] & X/t \ar[d, "q/t"] \\
U'/t' \ar[r, hook, "i'/t'"] & X'/t',
\end{tikzcd}
\tag*{$(*)$}
\]
I say that it is \textit{cartesian.} This follows from the sole assumption that $Q$ is cartesian, without any smoothness assumption on $q$ nor that $i'$ is an open immersion, and can be seen e.g. by looking at the cubical diagram:
% foot: Version 7 décembre 2002
\printedpage{110}
% head: A. Grothendieck : Dérivateurs, Ch. XII, Édition des Universités de Montpellier II et Paris VII
\[
\begin{tikzcd}[row sep=small, column sep=small]
& U/t \ar[rr] \ar[dl] \ar[dd] & & X/t \ar[dl] \ar[dd] \\
U'/t' \ar[rr, crossing over] \ar[dd] & & X'/t' & \\
& U \ar[rr] \ar[dl] & & X \ar[dl] \\
U' \ar[rr] & & X' \ar[from=uu, crossing over] &
\end{tikzcd}
\]
(\footnote[110]{NB Here $t$ is an arbitrary element of $\Ob X$, $t'$ its image in $X'$.}). In this diagram, the front and back face of the cube is cartesian, by definition of $U/t$ and of $U'/t'$
\msp{page 160}
respectively, and the base is by assumption. It follows that the top face of the cube is also a cartesian square.

This seen, since $q$ is smooth, the same is true of $q/t$ in the cartesian square $(*)$, hence also of $p/t$ by base change. On the other hand, the smoothness of $q$ implies that $q/t$ has $W$-aspheric fibers, hence by base change the same is true of $p/t$ (which is thus a universal $W$-equivalence, i.e. a trivial $W$-fibration). Moreover, because of the smoothness of $p/t$, for $v'$ in $U'/t'$, the inclusion % [sic] commme -> comme, repeated letter
\[
(U/t)_{v'} \longrightarrow v' \backslash (U/t)
\]
is in $W$, hence (by W(5)) defines a bijection on the $\pi_0$. Since $(U/t)_{v'}$ is $W$-aspheric by what precedes, it follows (by W(5)) that its $\pi_0$ has cardinality 1, hence the same is true of $v' \backslash (U/t)$, i.e. the latter is 0-connected. Now that is condition d) of proposition 15, which was to be established.

\msp{page 161}
\textbf{Corollary 2.} \textit{Let $Q$ be a cartesian square} (p. 158), \textit{with $i'$ (hence also $i$) an open immersion, and $q$ (hence also $p$) a} local isomorphism. \textit{Let $T = X \backslash U$, $T' = X' \backslash U'$ (hence $T \simeq X \times_{X'} T'$), and suppose $T \longrightarrow T'$ an isomorphism. Then $Q$ is cocartesian.}
\DERmark{ENV0777}


This is a special case of corollary 1, but with no explicit reference to a localizer $W$. We can take $W = W_{\infty}$, or $W = W_0$, to reduce to corollary 1. But it is more reasonable to say that in the cartesian square $(*)$, the assumption that $q$ is a local isomorphism implies that $q/t$ is an isomorphism, hence $p/t$ is too. It does indeed follow that the $v' \backslash (U/t)$ are 0-connected, this category being isomorphic to $v' \backslash (U'/t')$, hence has an initial object. % [sic] pendre -> prendre, missing letter

Now let $X$ be an object of Cat, and let us consider a cocartesian diagram in $X^{\wedge}$:
\[
\begin{tikzcd}
G \ar[r, "i"] \ar[d, "p"'] & F \ar[d, "q"] \\
G' \ar[r, "i'"] & F',
\end{tikzcd}
\tag*{$(*)$}
\]
% foot: Version 7 décembre 2002
\printedpage{111}
% head: A. Grothendieck : Dérivateurs, Ch. XII, Édition des Universités de Montpellier II et Paris VII
\msp{page 162}
where $i$ is a monomorphism. We know that the cocartesian character of the square is checked ``argument by argument'' by the cocartesian character of
\[
\begin{tikzcd}
G(x) \ar[r, "i(x)"] \ar[d, "p(x)"'] & F(x) \ar[d, "q(x)"] \\
G'(x) \ar[r, "i'(x)"] & F'(x),
\end{tikzcd}
\tag*{$(*)$}
\]
in Ens, $\forall\, x$ in $X$. Now such a cocartesian square in Ens, with $i(x)$ \texttt{[a]} monomorphism, is also cartesian. This implies that the square $Q$ in $X^{\wedge}$ is also cartesian.

Let us consider, on the other hand, the canonical functor
\[
\varphi_X : F \longmapsto X/F : X^{\wedge} \longrightarrow \mathrm{Cat}.
\]
It is better, moreover, to look at it as a functor
\[
\varphi_X : X^{\wedge} \longrightarrow \mathrm{Cat}/X,
\]
and I recall, for the record, that this functor is fully faithful, and that its essential image consists of the categories $X'$ over $X$ such that
\[
X' \longrightarrow X
\]
\msp{page 163}
is a local isomorphism, or what amounts to the same, that $X'$ is \textit{fibered} over $X$ with discrete fibers:
\[
\varphi_X : X^{\wedge} \xrightarrow{\ \sim\ } \underbrace{\Fibdisc\, X}_{\substack{\text{fibered categories}\\ \text{with discrete fibers}}} \ .
\]
Let us note that every arrow $X' \longrightarrow X''$ in $\Fibdisc\, X$ is a local isomorphism (tautological on $\begin{tikzcd}[ampersand replacement=\&, row sep=small, column sep=tiny, baseline=(\tikzcdmatrixname-1-1.base)] X'/x' \ar[rr] \ar[dr, "\text{iso.}"'] \& \& X''/x'' \ar[dl, "\text{iso.}"] \\ \& X/x \& \end{tikzcd}$), which shows that if we have a diagram
\[
\begin{tikzcd}[row sep=small, column sep=small]
X_1 \ar[dr] & & X_2 \ar[dl] \\
& X_0 &
\end{tikzcd}
\]
in $\Fibdisc\, X$, the fiber product $X_3 = X_1 \times_{X_0} X_2$ in Cat is such that $X_3 \longrightarrow X_2$ is a local isomorphism, hence also $X_3 \longrightarrow X$. That is to say $\Fibdisc(X)$ is stable under fiber products. This implies that fiber products exist in $\Fibdisc\, X$, and that they are the fiber products in Cat, or in Cat$/X$.
% foot: Version 7 décembre 2002

\printedpage{112}
% head: A. Grothendieck : Dérivateurs, Ch. XII, Édition des Universités de Montpellier II et Paris VII
\msp{page 164}
Likewise for finite products in $\Fibdisc\, X$, they are the finite products in Cat$/X$. In short

\textbf{Lemma 1.} \textit{The full subcategory} $\Fibdisc\, X$ \textit{of} $\Cat/X$ \textit{is stable under finite} $\varprojlim$\textit{, hence the} $\varprojlim$ \textit{exist in} $\Fibdisc\, X$ \textit{and the inclusion functor}
\[
\Fibdisc\, X \longrightarrow \Cat/X
\]
\textit{commutes with them, and in particular transforms cartesian squares of} $\Fibdisc\, X$ \textit{into cartesian squares of} $\Cat/X$\textit{, or what amounts to the same, of} $\Cat$\textit{. Moreover, by virtue of the equivalence}
\[
\varphi_X : X^{\wedge} \xrightarrow{\ \sim\ } \Fibdisc\, X \hookrightarrow \Cat/X,
\]
\textit{the functor $\varphi_X$ from $X^{\wedge}$ to} $\Cat/X$ \textit{is left exact, in particular it transforms cartesian squares of $X^{\wedge}$ into cartesian squares of} $\Cat/X$\textit{, or again of} $\Cat$\textit{.} % [sic] exacte -> exact, agreement with "le foncteur"

\msp{page 165}
Let us also note that being left exact, the functor $\varphi_X$ from $X^{\wedge}$ to Cat$/X$ transforms monomorphisms into monomorphisms. On the other hand:

\textbf{Lemma 2.} \textit{The monomorphisms which are local isomorphisms are exactly the open immersions. So in} $\Fibdisc\, X$\textit{, the monomorphisms are the open immersions, and the functor $\varphi_X : F \longmapsto X/F$ transforms monomorphisms into open immersions.}

Let us then return to the cocartesian diagram $(*)$ (p. 161), and to the corresponding diagram in Cat (and even in $\Fibdisc\, X$)
\[
\begin{tikzcd}
X/G \ar[r, hook, "X/i"] \ar[d, "X/p"'] & X/F \ar[d, "X/q"] & T \ar[l, hook', dotted] \ar[d, dotted] \\
X/G' \ar[r, hook, "X/i'"] & X/F' & T'. \ar[l, hook', dotted]
\end{tikzcd}
\]
Since $(*)$ is cartesian, the preceding square is cartesian, by what precedes. I want to prove that it is also cocartesian
\msp{page 166}
(in Cat$/X$, or in Cat, it amounts to the same), and for this I apply corollary 2 of proposition 15 (p. 161), taking into account that $X/q$ is a local isomorphism, and $X/i'$ (being a local isomorphism \texttt{[and a]} monomorphism) an open immersion. It suffices to prove that, $T'$ and $T$ being the closed subcategories complementary to the opens $X/G'$ of $X/F'$ and to $X/G$ of $X/F$, that $T \longrightarrow T'$ is \texttt{[an]} isomorphism.

{\small \texttt{[Now, $F$ and $G$, and $F'$ and $G'$, but not $p$ and $q$, exchange their roles.]}} % [sic] leur rôles -> leurs rôles, agreement

The verification is immediate: $T$ is the category of arrows $x \longrightarrow G$ in $X^{\wedge}$ (with source $x$ in $X$) which do not factor through $F$ \texttt{[rather through $F$]}, $T'$ that of arrows $x \longrightarrow G'$
\DERmark{ENV0778}

% foot: Version 7 décembre 2002


\endgroup
% END UNIT Der-ChXIIb-012/Der-ChXIIb-012_p066-070.tex
% BEGIN UNIT Der-ChXIIb-012/Der-ChXIIb-012_p071-072.tex
\begingroup\makeatletter%
\newcommand{\printedpage}[1]{\marginpar{\footnotesize p.~#1}}%
\newcommand{\msp}[1]{\marginpar{\footnotesize [#1]}}%
\makeatother\let\printedpage\derpp\let\msp\dermsp%
\let\mspin\DERmspin%


\printedpage{113}
% head: A. Grothendieck : Dérivateurs, Ch. XII, Édition des Universités de Montpellier II et Paris VII
which do not factor through $F'$, the functor $T \longrightarrow T'$ being the functor of composition of $x \longrightarrow G$ with $G \xrightarrow{q} G'$
\[
\begin{tikzcd}[row sep=large]
F \arrow[r, hook] \arrow[d] & G \arrow[d] \\
F' \arrow[r, hook] & G'
\end{tikzcd}
\]
(NB The composite does not factor through $F'$, for otherwise, because of the cartesian character of $Q$, $x \longrightarrow G$ would factor through $F$.) But for given $x$, the datum of $x \longrightarrow G$ not factoring through $F$ is that of an element \msp{page 167} of $G(x)\backslash F(x)$, and because of the cocartesian character of $Q$,
\[
(G(x)\backslash F(x)) \longrightarrow (G'(x)\backslash F'(x))
\]
is bijective, i.e. ``it amounts to the same thing'' to give $x \longrightarrow G$ not factoring through $F$, or $x \longrightarrow G'$ not factoring through $F'$.

Thus we have proved the

\textbf{Theorem 5.} \textit{Let $X$ be in} Cat\textit{, and}
\[
Q : \qquad
\begin{tikzcd}[row sep=large, baseline=(current bounding box.center)]
F \arrow[r, hook, "i"] \arrow[d, "p"'] & G \arrow[d, "q"] \\
F' \arrow[r, hook, "i'"] & G'
\end{tikzcd}
\]
\textit{a cocartesian square in $X^\wedge$, with $i$} \texttt{[a]} \textit{monomorphism. Then $i'$ is} \texttt{[a]} \textit{monomorphism, the square $Q$ is also cartesian, and the corresponding square in} Cat$/X$ \textit{or in} Cat
\[
\begin{tikzcd}[row sep=large]
X/F \arrow[r, hook, "X/i"] \arrow[d, "X/p"'] & X/G \arrow[d, "X/q"] \\
X/F' \arrow[r, hook, "X/i'"] & X/G'
\end{tikzcd}
\]
\msp{page 168} \textit{is both cartesian and cocartesian. Moreover, if $W$ is a basic localizer in} Cat\textit{, satisfying} W(7 bis) (in addition to W(1,2,3))\textit{, then the square $X/Q$ in} Cat \textit{is $W$-cocartesian.}

\printedpage{114}
% head: A. Grothendieck : Dérivateurs, Ch. XII, Édition des Universités de Montpellier II et Paris VII
This last statement follows from what preceded and from corollary 2 of theorem 4, taking into account that $X/q$ is a local isomorphism, hence is $W$-smooth.

\textbf{Corollary:} (\footnote[111]{\textbf{NB} If we only want $i \in W_X \Longrightarrow i' \in W_X$, $p \in W_X \Longrightarrow q \in W_X$, it suffices to have W(7 bis) \textit{weak}, i.e. $i_0 \in W \Longrightarrow i \in W$ in the statement of W(7 bis), p. 76 (Loc(6$'$a) with the notation of XVI).}) \textit{We have} (assuming W(7 bis)) \textit{the equivalences}
\[
\begin{aligned}
i \in W_X \quad &\Longleftrightarrow \quad i' \in W_X \\
p \in W_X \quad &\Longleftrightarrow \quad q \in W_X,
\end{aligned}
\]
\textit{where}
\[
W_X \subseteq \mathit{Fl}(X^\wedge)
\]
\textit{is defined by (cf. Pursuing Stacks \ldots)}
\[
u \in W_X \qquad\qquad \Longleftrightarrow \qquad\qquad X/u : X/F \longrightarrow X/G \quad \textit{is in } W.
\]

Thus, in a category $X^\wedge$, amalgamated sums involving a monomorphism behave perfectly from the homotopy point of view, more precisely for the localizers $W_X$, provided \msp{page 169} only that $W$ satisfies W(7 bis). In particular, the monomorphisms in $X^\wedge$ are ``$W_X$-cofibrations'', in a sense similar to that of definition 2 c) (pages 109, 110 \texttt{[rather 109, 110]}).
\DERmark{ENV0779}


\texttt{[page 170 does not exist.]}


\endgroup
% END UNIT Der-ChXIIb-012/Der-ChXIIb-012_p071-072.tex
\clearpage\setcounter{footnote}{0}\phantomsection\addcontentsline{toc}{section}{Chapter XII, 6}\markboth{Chapter XII, 6}{Chapter XII, 6}
% BEGIN UNIT Der-ChXIIc-30-03-03/Der-ChXIIc-30-03-03_p001-005.tex
\begingroup\makeatletter%
\newcommand{\printedpage}[1]{\marginpar{\footnotesize p.~#1}}%
\newcommand{\msp}[1]{\marginpar{\footnotesize [#1]}}%
\newcommand{\Ch}{\underline{\mathrm{Ch}}}%
\newcommand{\Hotp}[1]{\mathrm{Hot}_{\bullet\,#1}}%
\newcommand{\hotp}[1]{\mathrm{hot}_{\bullet\,#1}}%
\newcommand{\Ob}{\mathrm{Ob}\,}%
\newcommand{\HomEns}{\mathrm{Hom}_{\mathrm{Ens}}}%
\newcommand{\utpi}{\underset{\sim}{\pi}}%
\makeatother\let\printedpage\derpp\let\msp\dermsp%
\let\mspin\DERmspin%


\printedpage{115}
% head: A. Grothendieck : Dérivateurs, Ch. XII, Édition des Universités de Montpellier II et Paris VII
% foot: Version 30 mars 2003

\msp{page 171}

\section*{VI\quad Criterion for $W_\infty$-asphericity of the $\Omega^n(X,x)$.}

\subsection*{1\quad Cogitations on $W_n$.}

We are going to define, for $n \in \mathbf{N}$,
\[
W_n \subseteq \mathrm{Fl}(\mathrm{Cat})
\]
by the following conditions (where $f : X \longrightarrow Y$ is an arrow in Cat).

\noindent\begin{tabular}{@{}l c l >{\RaggedRight\arraybackslash}p{0.64\textwidth}@{}}
$f \in W_0$ & $\Longleftrightarrow$ & a) & $\pi_0(f)$ bijective.\\
 & $\Longleftrightarrow$ & b) & For every presheaf of sets $F$ \textit{constant} on $Y$, $H^0(Y,F) \longrightarrow H^0(X,f^*F)$ is bijective.\\
 & $\Longleftrightarrow$ & c) & The functor $F \longmapsto f^*F$ from presheaves of sets \textit{constant} on $Y$ to constant sheaves of sets on $X$ is fully faithful.
\end{tabular}

\texttt{[The terms sheaf or presheaf are used interchangeably here.]}

\textbf{NB} If $F_0$ is a set, $F_{0\,Y}$ the constant sheaf of sets defined by $F_0$ on $Y$, we have
\DERmark{ENV0780}

\[
\begin{aligned}
H^0(Y,F_{0\,Y}) &\simeq \HomEns(\pi_0(Y),F_0) && \text{(functorial in $Y$)}\\
\mathrm{Hom}_{Y^\wedge}(F_{0\,Y},G_{0\,Y}) &\simeq \HomEns(\pi_0(Y),\HomEns(F_0,G_0)) && \text{(functorial in $Y$),}
\end{aligned}
\]
which implies that a) $\Longrightarrow$ c) $\underset{\substack{\text{taking}\\ \mathrm{Hom}_{Y^\wedge}(e_Y,F)}}{\overset{\text{tautologically}}{\Longrightarrow}}$ b) $\Longrightarrow$ a).

\noindent\begin{tabular}{@{}l c l >{\RaggedRight\arraybackslash}p{0.64\textwidth}@{}}
$f \in W_1$ & $\Longleftrightarrow$ & a) & $\pi_0(f)$ bijective, and for every $x \in \Ob X$, $\pi_1(f,x) : \pi_1(X,x) \longrightarrow \pi_1(Y,f(x))$ is bijective. (NB It suffices to check this for \textit{one} $x$ in each connected component of $X$.)\\
 & $\Longleftrightarrow$ & a$'$) & $\Pi(X) \xrightarrow{\Pi f} \Pi(Y)$ is an equivalence of categories.\\
 & $\Longleftrightarrow$ & b) & $H^0(Y,F) \longrightarrow H^0(X,f^*F)$ bijective for every sheaf of sets \textit{locally constant} on $Y$, and $H^1(Y,F) \longrightarrow H^1(X,f^*F)$ bijective if in addition $F$ is a sheaf of groups.\\
 & $\Longleftrightarrow$ & b$'$) & As in b), but for $H^1$ we restrict ourselves to $F$ \textit{constant}, and require $H^1(f)$ \textit{surjective} only.\\
 & $\Longleftrightarrow$ & c) & The functor $F \longmapsto f^*F$ from the category of locally constant sheaves of sets on $Y$ to that of locally constant sheaves on $X$ is an equivalence of categories.
\end{tabular}

\par

\printedpage{116}
% head: A. Grothendieck : Dérivateurs, Ch. XII, Édition des Universités de Montpellier II et Paris VII
% foot: Version 30 mars 2003

\msp{page 172}

\textbf{NB} The group $\pi_1(X,x)$ is defined via the theory of the fundamental groupoid. We know that it is in canonical bijection with $\pi_0(\Omega(X,x))$, where $\Omega(X,x) = \Ch(X;x,x)$ (which allows us to have a monoidal composition law on $\Omega(X,x)$, giving the group law on $\pi_0$), or at one's choice $\Omega(X,x) = \Ch_\infty(X;x,x)$ (more convenient in certain situations, but not directly giving the composition law in $\pi_0(\Omega(X,x))$).

I am going to set, for $n \in \mathbf{N}$,
\[
\pi_n(X,x) = \pi_0(\Omega^n(X,x)),
\]
$\Omega^n$ denoting the $n$-th iterate of the functor $(X,x) \longmapsto \Omega(X,x)$ from the category $\mathrm{Cat}_\bullet$ of pointed objects of Cat to itself. I prefer here to work with the \underline{Ch} theory instead of $\Ch_\infty$. Sooner or later we will have to give a comparison statement \dots

\textbf{Theorem 1} (\footnote[112]{Probably \textit{false}. We should have a) $\Longrightarrow$ b), but b) $\Longrightarrow$ a) \texttt{[is]} probably false already for $n = 2$ (except in the case $X$, $Y$ 1-connected). Cf. corrected statement \texttt{[p. 178]}.}). \textit{Let $f : X \longrightarrow Y$ be in} Cat\textit{, and let $n \in \mathbf{N}$, $n \geq 2$, $S$ a subset of} $\Ob X$ \textit{such that the composite $S \longrightarrow \Ob X \longrightarrow \pi_0(X)$ is} \texttt{[an]} \textit{epimorphism. The following conditions are equivalent.}

\begin{itemize}
\item[\textit{a)}] \textit{$\pi_i(f,x)$ an isomorphism for every $x \in \Ob X$ and every $i \in [0,n]$.}

\item[\textit{a$'$)}] \textit{As in a), but with $x \in S$.}

\item[\textit{b)}] \textit{We have}
\[
\begin{array}{rcll}
H^0(Y,F) & \xrightarrow{\ \sim\ } & H^0(X,f^*F) & \textit{for $F$ a sheaf of sets}\\
 & & & \textit{locally constant on $Y$.}\\
H^1(Y,F) & \xrightarrow{\ \sim\ } & H^1(X,f^*F) & \textit{if in addition $F$ a sheaf of groups.}\\
H^i(Y,F) & \xrightarrow{\ \sim\ } & H^i(X,f^*F) & \textit{if in addition $F$ an abelian sheaf, $i \in [2,n]$.}
\end{array}
\]
\end{itemize}

\msp{page 173}

\textbf{NB} The equivalence of a) and a$'$) should be trivial, but for this we need to have defined the $\pi_n(X,x)$, for $x$ varying in $\Ob X$, as a local system on $X$, something I have not yet done. I am going to try to do it in a draft of air. First of all, the $\Omega(X,x)$ are the fibers of a perfect $W_\omega$-fibration (\footnote[113]{\textbf{NB} $W_\omega$ is the smallest localizer of Cat (satisfying W(1, 2 bis, 3)).}), fitting into the cartesian square
\[
\begin{tikzcd}
\Ch(X) \arrow[d, "{p_X=(\sigma_X,\beta_X)}"'] & \Omega(X,-) \arrow[l] \arrow[d] \\
X \times X & X. \arrow[l, "\mathrm{diag}_X"'] \arrow[ul, dotted, "i_X"']
\end{tikzcd}
\]

\par

\printedpage{117}
% head: A. Grothendieck : Dérivateurs, Ch. XII, Édition des Universités de Montpellier II et Paris VII
% foot: Version 30 mars 2003

We then know that we have a functor
\[
\begin{aligned}
\Pi X &\longrightarrow \mathrm{Hot}_{W_\omega}\\
x &\longmapsto \mathrm{hot}_{W_\omega}(\Omega(X,x)).
\end{aligned}
\tag*{$(*)$}
\]
Now we have, for $i \geq 1$,
\[
\pi_i(X,x) = \pi_{i-1}(\ \underbrace{\Omega(X,x)}_{\text{pointed object}}\ ).
\]
This gives the desired local system of the $\pi_i$, if only we had a functor a little more precise than $(*)$, namely
\[
\Pi X \longrightarrow \Hotp{W_\omega} \qquad \text{(\textit{pointed} $W_\omega$-homotopy category)}
\]
(\footnote[114]{But one must also check that the $\pi_n(X,x)$ ($n \geq 1$) are functors on $\Hotp{W_\infty}$, i.e. transform $W_\infty$-equivalences into likewise. This follows from W(4a) for $W_\infty$, which implies that $(X,x) \mapsto \Omega(X,x)$ factors through $\Hotp{W_\infty}$ (as a functor with values in $\Hotp{W_\infty}$).}). Now the arrow $i_X : X \longrightarrow \Ch X$ lifting $p_X$ gives a canonical section of $\Omega(X,-)$ over $X$ -
\msp{page 174}
precisely the one which assigns to every $x \in X$ the marked point $1_x$ of $\Omega(X;x,x)$ (``empty loop at $x$''). The question is thus this:

\begin{quote}
We have a $W$-fibration $\Omega \longrightarrow X$, with a given section $\sigma$. Does this define a canonical functor % [sic] donné -> donnée, accord avec « section »
\[
X \longrightarrow \Hotp{W},
\]
turning every arrow into \texttt{[an]} isomorphism?
\end{quote}

The natural idea is to take
\[
x \longmapsto \hotp{W}(\Omega/x)
\]
by pointing $\Omega/x$ at the marked point $\sigma(x)$ of $\Omega_x \hookrightarrow \Omega/x$. The trouble is that the transition morphisms $\Omega/x \longrightarrow \Omega/x'$ (for $x \xrightarrow{u} x'$ in $X$) are not compatible with the marked points. But we have an arrow (namely $\sigma(u)$) going
\msp{page 175}
from the image $\sigma(x)/x'$ of the marked point of $\Omega/x$ to the marked point $\sigma(x')/x'$ of $\Omega/x'$. This should suffice to define the arrow
\[
(\Omega/x,\sigma(x)/x) \longrightarrow (\Omega/x',\sigma(x')/x')
\]
as an arrow in $\Hotp{W}$, with transitivity for a composite $x \xrightarrow{u} x' \xrightarrow{v} x''$ in $X$. For the moment I will confine myself
\printedpage{118}
% head: A. Grothendieck : Dérivateurs, Ch. XII, Édition des Universités de Montpellier II et Paris VII
% foot: Version 30 mars 2003
\msp{page 176}
to an indication, leaving detailed verifications for later. If need be, I will use right now that for $i \in \mathbf{N}$, $i \geq 1$,
\[
\pi_i(f,x) : \pi_i(X,x) \longrightarrow \pi_i(Y,f(x))
\]
can be interpreted as coming from a homomorphism of local systems
\[
\utpi_i(X,-) \longrightarrow f^*(\utpi_i(Y,-))
\]
on $X$. Condition a) of Theorem 1 is interpreted by saying that this homomorphism is an \textit{isomorphism} for $i \in [1,n]$, and condition a$'$) then merely recalls that it suffices to test that this is so for \textit{one} point of each connected component of $X$.

The statement of Theorem 1 retains a meaning for $n = 1$, and in this case was ``recalled'' previously (and is $\pm$ tautological). This will allow us to proceed by induction on $n$, assuming that the theorem is proved for $n-1$. This means that we can suppose that conditions a) \textit{and} b) are satisfied by $f$ for $i \in [1,n-1]$, and confine ourselves to proving that, under these conditions, condition a) (or a$'$), for $i = n$, is equivalent to condition b) for $i = n$.

\msp{page 177}

So we simply have to prove this:

\textbf{Lemma} (\footnote[115]{Probably false.}). \textit{Suppose that $\pi_i(f,x)$} \texttt{[is an]} \textit{isomorphism for $i \in [0,n-1]$ and b) satisfied for $i \in [0,n-1]$, where $n \in \mathbf{N}$, $n \geq 2$, is fixed. Then the following conditions are equivalent:}

\begin{itemize}
\item[\textit{a)}] $\pi_n(f,x) : \pi_n(X,x) \longrightarrow \pi_n(Y,f(x))$ \texttt{[is an]} \textit{isomorphism for every $x \in \Ob X$.}

\item[\textit{b)}] $H^n(Y,F) \xrightarrow{\sim} H^n(X,f^*F)$ \texttt{[is an]} \textit{isomorphism for every locally constant abelian sheaf on $Y$.}
\end{itemize}

Let us note that $f$ induces an arrow in $\mathrm{Cat}_\bullet$,
\[
\Omega(f,x) : \Omega(X,x) \longrightarrow \Omega(Y,f(x)),
\]
whence
\[
\begin{tikzcd}[row sep=small]
\pi_i(\Omega(X,x)) \arrow[r] \arrow[d, dash, "\wr"'] & \pi_i(\Omega(Y,f(x))) \arrow[d, dash, "\wr"'] \\
\pi_{i+1}(X,x) \arrow[r] & \pi_{i+1}(Y,f(x)),
\end{tikzcd}
\]
identifying with $\pi_{i+1}(f,x)$:
\[
\pi_i(\Omega(f,x)) \simeq \pi_{i+1}(f,x).
\]


\printedpage{119}
% head: A. Grothendieck : Dérivateurs, Ch. XII, Édition des Universités de Montpellier II et Paris VII
% foot: Version 30 mars 2003

Hypothesis a) of the lemma implies that $\pi_i(\Omega(X,x)) \xrightarrow{\sim} \pi_i(\Omega(Y,y))$ for $i \in [0,n-1]$. Here we are dealing with the $\pi_i$ based at the unit points. We will probably have to prove that, because this is a morphism of monoids in Cat (satisfying certain conditions?) this remains true for the $\pi_i$ at \textit{every} point $\xi$ of $\Omega(X,x)$. By the induction hypothesis, we will thus know that $\Omega(X,x) \longrightarrow \Omega(Y,y)$ % [sic] \Omega(X,n) -> \Omega(X,x), le point de base est x
\DERmark{ENV0781}

\msp{page 178}
induces an isomorphism on cohomology with locally constant coefficients, up to and including dimension $n-1$.

Finally, I have $\pm$ convinced myself that the statement of Theorem 1 is false as it stands. Here is a statement that I hope is true, and that I am going to try to prove in a ``utilitarian'' spirit, using everything I know in homotopy and cohomology (and not just what I have developed ``from scratch'' in the modelizer Cat). % [sic] j'éspère -> j'espère, orthographe

\textbf{Theorem 1 corrected.} \textit{Let us consider, for a given integer $n \geq 2$, the following conditions on an arrow $f : X \longrightarrow Y$ in} Cat \textit{(or a continuous map of topological spaces, or a morphism in $\boldsymbol{\Delta}^\wedge$ etc. etc.):}

\begin{itemize}
\item[$a_n$)] \textit{If $Z$ is a ``homotopy fiber'' of $f$ (NB there is essentially one per connected component of $Y$), then $Z$ is $n$-connected, i.e.}
\msp{page 179}
\textit{(e.g.) $\pi_i(Z) = 0$ for $i \in [0,n]$.}

\item[$b_n$)] \textit{$\pi_0(f)$ is bijective, and $\forall\ x \in X$, $\pi_i(f,x) : \pi_i(X,x) \longrightarrow \pi_i(Y,f(x))$ is} \texttt{[an]} \textit{isomorphism for $i \in [1,n]$. (}\textbf{NB} \textit{It suffices to say this for} \textup{one} \textit{point $x$ of each connected component of $X$.)}

\item[$c_n$)] $H^i(Y,F) \longrightarrow H^i(X,f^*F)$ \texttt{[is an]} \textit{isomorphism for $0\leq i\leq n$ and $F$ a locally constant sheaf on $Y$ (with $F$ abelian if $i \geq 2$, $F$ a sheaf of groups if $i = 1$, $F$ a sheaf of sets if $i=0$).}
\end{itemize}
\DERmark{ENV0782}


\textit{We then have}
\[
a_n \ \Longrightarrow\ b_n \ \Longrightarrow \text{(\footnotemark[116])}\ c_n \ \Longrightarrow \text{(\footnotemark[117])}\ a_{n-1}
\]
\footnotetext[116]{if $X$, $Y$ 1-connected, this is really proved, otherwise, there are verifications to be made (which look sorital \dots).}%
\footnotetext[117]{\textbf{NB} Even if $X$, $Y$ 1-connected, it is not clear that $\mathrm{c}_n \Longrightarrow \mathrm{b}_n$, for one must show that $\pi_n(Z) \longrightarrow \pi_n(X)$ is zero, and one only finds that the composite with $\pi_n(X) \longrightarrow H_n(X)$ is zero. I presume that there are counterexamples.}%
\textit{(and I presume that each of these implications is strict).}

\textsc{Proof}. $\mathrm{a}_n \Longrightarrow \mathrm{b}_n$ is an immediate consequence of the homotopy exact sequence for a Hot-fibration $\bar{X} \xrightarrow{\bar{f}} Y$ which factors $X \xrightarrow{f} Y$ as $f = \bar{f} \circ i$, $i \in W_\infty$. In fact, $\mathrm{a}_n$ \textit{is equivalent} to $\mathrm{b}_n$ + the surjectivity of the $\pi_{n+1}(X,x) \longrightarrow \pi_{n+1}(Y,f(x))$, which of course does \textit{not} follow from $\mathrm{b}_n$, i.e. from the bijectivity of the $\pi_i(f,x)$ for $i \leq n$ (take some $K(\pi,n+1)$). Thus the implication is indeed strict.


\endgroup
% END UNIT Der-ChXIIc-30-03-03/Der-ChXIIc-30-03-03_p001-005.tex
% BEGIN UNIT Der-ChXIIc-30-03-03/Der-ChXIIc-30-03-03_p006-010.tex
\begingroup\makeatletter%
\NoAutoSpacing%
\newcommand{\printedpage}[1]{\marginpar{\footnotesize p.~#1}}%
\newcommand{\msp}[1]{\marginpar{\footnotesize [#1]}}%
\makeatother\let\printedpage\derpp\let\msp\dermsp%
\let\mspin\DERmspin%


\printedpage{120}
% head: A. Grothendieck : Dérivateurs, Ch. XII, Édition des Universités de Montpellier II et Paris VII

\msp{page 180}
$\mathrm{b}_n \Longrightarrow \mathrm{c}_n$. We can suppose that $f$ is a Hot-fibration (and even Hot-perfect if we feel like it \dots). We already know that for
\[
H^i(Y,F) \longrightarrow H^i(X,f^*F) \tag*{$(*)$}
\]
everything is O.K. if $i \in \{0,1\}$. So we can suppose $F$ abelian. If that gets us anywhere, we can suppose, by induction, (that the implication $\mathrm{b}_{n-1} \Longrightarrow \mathrm{c}_{n-1}$ is known, hence) that $(*)$ is \texttt{[an]} isomorphism for $i<n$. Remains to see it for $i=n$. We are reduced of course, since $\pi_0(f)$ is bijective, to the case where $X$, $Y$ are 0-connected, so there is essentially a single fiber $Z$. The homotopy exact sequence tells us that
\[
\pi_i(Z,z)=0 \qquad \text{for } i \le n-1 \ (Z \text{ is } (n-1)\text{-connected})
\]
and moreover
\[
\pi_n(Z,z) \longrightarrow \pi_n(X,z) \qquad \text{is zero}
\]
(since $\pi_n(X,z) \longrightarrow \pi_n(Y,y)$ \texttt{[is an]} isomorphism hence \texttt{[a]} monomorphism), so
\[
\pi_{n+1}(Y,y) \longrightarrow \pi_n(Z,z) \qquad \text{\texttt{[is an]} epimorphism.}
\]
As $Z$ is $(n-1)$-connected (hence also 1-connected), we have
\[
H^i(Z,G)=0\ , \qquad \text{if } 0<i<n\ , \tag*{$(*)$}
\]
for $G$ locally constant (hence constant) on $Z$.

\msp{page 181}
Let us write the Leray spectral sequence for $f : X \longrightarrow Y$ and $f^*F = F_X$:
\[
H^*(X,f^*F) \Longleftarrow E_2^{pq} = H^p(Y,H^q(Z,f^*F)),
\]
where ``$H^q(Z)$'' is to be viewed as a local system on $Y$. So by $(*)$ we have
\[
E_2^{pq}=0 \qquad \text{if } 0<q<n, \qquad
\begin{tikzpicture}[baseline=0pt, x=0.6cm, y=0.6cm]
\draw[->] (0,0) -- (0,4);
\draw[->] (0,0) -- (6,0);
\node[left] at (0,3) {$n$};
\foreach \i in {0,...,4} { \fill (\i,3) circle (2.5pt); \fill (\i,0) circle (2.5pt); }
\node at (1,1.5) {$0$};
\node at (3,1.5) {$0$};
\end{tikzpicture}
\qquad ,
\]
so we find the exact sequence
\[
0 \longrightarrow H^n(Y,F) \longrightarrow H^n(X,F_X) \xrightarrow{\ \alpha\ } H^0(Y,\ \underbrace{H^n(Z,\overbrace{F_Z}^{M_Z})}_{\begin{array}{l}\scriptstyle = \ \mathrm{Hom}(\pi_n(Z,z),M)\\ \scriptstyle \text{by \texttt{[the]} Hurewicz thm.}\end{array}}\ ) \longrightarrow H^{n+1}(Y,F).
\]

\printedpage{121}
% head: A. Grothendieck : Dérivateurs, Ch. XII, Édition des Universités de Montpellier II et Paris VII

Here I fix a point $x$ of $X$, $Z$ denotes the fiber of $f$ at $y=f(x)$, and $M$ that of $F$ at $y$, $M=F_y$. The group $\pi_1(Y)$ acts on it, for that matter $\pi \overset{\mathrm{def}}{=} \pi_1(X,x) \simeq \pi_1(Y,y)$ does too. The $H^0(Y,H^n(Z,F_Z))$ can be interpreted, unless I am mistaken, as
\DERmark{ENV0783}

\[
\mathrm{Hom}_\pi(\pi_n(Z,x),M),
\]
taking into account the action of $\pi=\pi_1(X,x)$ on $\pi_n(Z,x)$ as well. (To be checked.) We have to interpret the arrow
\[
H^n(X,F_X) \longrightarrow \mathrm{Hom}_\pi(\pi_n(Z,x),M),
\]
via the composite map \texttt{[where $z=x$, viewed as a point of $Z$]}
\[
\pi_n(Z,z) \longrightarrow \pi_n(X,x) \longrightarrow H_n(X,\mathbf{Z}). \tag*{$(*)$}
\]
If $F$ is \textit{constant} on $Y$ (e.g. if $Y$ \texttt{[is]} 1-connected), thus with value $M$, we have the canonical arrows
\[
H^n(X,F_X) = H^n(X,M) \longrightarrow \mathrm{Hom}(H_n(X,\mathbf{Z}),M) \longrightarrow \mathrm{Hom}(\pi_n(Z,z),M)\text{(\footnotemark[118])} \tag*{$(*')$}
\]
\footnotetext[118]{thanks to composition with $(*)$.}
\texttt{[read: }$\mathrm{Hom}_\pi$\texttt{]},

\msp{page 182}
and surely (apart from any hypothesis of type $\mathrm{b}_n$ or other) the direct canonical map ``of localization on $Y$'', followed by \dots,
\[
H^n(X,F_X) \longrightarrow H^0(Y,H^n(Z,F_Z)) \longrightarrow \mathrm{Hom}_\pi(\pi_n(Z),M) \tag*{$(**)$}
\]
is none other than the composite of the canonical maps above. In the present case where we suppose $\mathrm{b}_n$, the second arrow in $(**)$ is \texttt{[an]} isomorphism, so the composite ``is none other'' than the localization arrow itself. On the other hand $\pi_n(Z,z) \longrightarrow \pi_n(X,x)$ in $(*)$ is zero, so also $\pi_n(Z,z) \longrightarrow H_n(X,\mathbf{Z})$, so also the composite arrow $(*')$ deduced from it. Let us then return to the Leray exact sequence (previous page), we find that $H^n(Y,F) \xrightarrow{\ \sim\ } H^n(X,F_X)$.

How to recover the case where $F$ is \textit{not} constant? In the case $n=2$, this worked, by passing to the universal covers $\tilde{X}$, $\tilde{Y}$ of $X$, $Y$ (with respect to the chosen points $x$, $y$), and writing the spectral sequences
\[
\begin{tikzcd}[row sep=small, column sep=small]
H^*(Y,F) \arrow[d] & \Longleftarrow & E_2^{pq} \arrow[d] & = & H^p(\pi,H^q(\tilde{Y},F_{\tilde{Y}})) \\
H^*(X,F_X) & \Longleftarrow & E^{\prime pq}_2 & = & H^p(\pi,H^q(\tilde{X},F_{\tilde{X}})).
\end{tikzcd}
\]
The hypothesis $\mathrm{b}_n$ is preserved in passing from $f : X \longrightarrow Y$ to $\tilde{f} = \tilde{f}_x : \tilde{X} \longrightarrow \tilde{Y}$, so we know that
\printedpage{122}
% head: A. Grothendieck : Dérivateurs, Ch. XII, Édition des Universités de Montpellier II et Paris VII
\msp{page 183}
\[
E_2^{pq} \xrightarrow{\ \sim\ } E^{\prime pq}_2 \qquad \text{\texttt{[is an]} isomorphism if } q \le n. \tag*{$(*)$}
\]
In the case $n=2$, using that
\[
E_2^{p1} = E^{\prime p1}_2 = 0 \qquad \text{since } H^1(\tilde{X}) = H^1(\tilde{Y}) = 0, \qquad
\begin{tikzpicture}[baseline=0pt, x=0.5cm, y=0.5cm]
\draw[->] (0,0) -- (0,2.8);
\draw[->] (0,0) -- (5,0);
\node[above left] at (0,2.6) {$q$};
\node[left] at (0,2) {$2$};
\foreach \i in {0,...,3} { \fill (\i,2) circle (2.5pt); \fill (\i,0) circle (2.5pt); }
\node at (0,1) {$0$};
\node at (1,1) {$0$};
\node at (2,1) {$0$};
\end{tikzpicture}
\qquad ,
\]
we find two exact sequences
\[
\begin{tikzcd}
0 \arrow[r] & H^2(\pi,H^0(\tilde{Y})) \arrow[r] \arrow[d, "\text{iso}"] & H^2(Y) \arrow[r] \arrow[d] & H^0(\pi,H^2(\tilde{Y})) \arrow[r] \arrow[d, "\text{iso}"] & H^3(\pi,H^0(\tilde{Y})) \arrow[d, "\text{iso}"] \\
0 \arrow[r] & H^2(\pi,H^0(\tilde{X})) \arrow[r] & H^2(X) \arrow[r] & H^0(\pi,H^2(\tilde{X})) \arrow[r] & H^3(\pi,H^0(\tilde{X})),
\end{tikzcd}
\]
and a homomorphism from one to the other, which allows us to conclude by the 5 lemma that $H^2(Y,F) \longrightarrow H^2(X,F_X)$ is \texttt{[an]} isomorphism. But for arbitrary $n$, the relation $(*)$ on the $E_2$ terms (for $q \le n$) does not seem to me to imply the similar relation for the higher $E_n$, and from there for the $E_\infty$ terms, which would allow us to conclude.

So I return to the exact sequence p. 181, with $F$ a sheaf that is \textit{twisted.} I feel like finding a canonical map
\[
H^n(X,G) \longrightarrow \overbrace{\mathrm{Hom}_\pi(\pi_n(X,x),G_x)}^{\text{\scriptsize Homs \textit{additive} of course, i.e. of $\mathbf{Z}[\pi]$-modules}} \qquad ,
\]
where $G$ is any locally constant abelian sheaf on $X$, and $\pi = \pi_1(X,x)$. If I interpret
\msp{page 184}
$\pi_n(X,x)$ in terms of homotopy classes of pointed maps
\[
(S^n,x_0) \xrightarrow{\ \varphi\ } (X,x)
\]
I do have, for such a map $\varphi$,
\[
\varphi^* : H^n(X,G) \longrightarrow H^n(S^n,\ \overbrace{\varphi^*G}^{\simeq\,(G_x)_{S^n}}\ ) \simeq G_x
\]
(taking into account that $S^n$ is 1-connected, and that its $n$-cohomology with values in a constant sheaf $M$ is $M$.) Surely the map
\[
H^n(X,G) \longrightarrow G_x
\]
\printedpage{123}
% head: A. Grothendieck : Dérivateurs, Ch. XII, Édition des Universités de Montpellier II et Paris VII
thus defined depends only on the pointed homotopy class of $\varphi$, hence
\[
\pi_n(X,x) \times H^n(X,G) \longrightarrow G_x,
\]
and we must also prove a) that this is compatible with the actions of $\pi = \pi_1(X,x)$, and b) that it is additive in $\pi_n(X,x)$, to have the desired map
\[
H^n(X,G) \longrightarrow \mathrm{Hom}_\pi(\pi_n(X,x),G_x).
\]
If I admit this, it is only one step to admitting also that the arrow that really interests us (in the exact sequence p. 181)
\msp{page 185}
\[
H^n(X,F_X) \longrightarrow \mathrm{Hom}_\pi(\pi_n(Z,z),\underbrace{M}_{=\,F_y})
\]
is none other than the composite of the preceding arrow (previous page) for $G = F_X$, with
\[
\mathrm{Hom}_\pi(\pi_n(X,x),M) \longrightarrow \mathrm{Hom}_\pi(\pi_n(Z,x),M)
\]
deduced from
\[
\pi_n(Z,x) \longrightarrow \pi_n(X,x).
\]
Now this arrow is zero, by the hypothesis $\mathrm{b}_n$. So the arrow $\alpha$ in the exact sequence page 181 is zero, so $H^n(Y,F) \longrightarrow H^n(X,F_X)$ is \texttt{[an]} isomorphism, q.e.d. (but there is quite a bit of checking to do \dots).

$\mathrm{c}_n \Longrightarrow \mathrm{a}_{n-1}$. We can proceed by induction on $n$. (NB The statement makes sense even for $n=1$, where it is proved, and even $\mathrm{c}_1 \Longrightarrow \mathrm{b}_1$, and $\mathrm{b}_1 \Longrightarrow \mathrm{a}_0$, more generally $\mathrm{b}_n \Longrightarrow \mathrm{a}_{n-1}$, immediate from \texttt{[the]} homotopy exact sequence.) So we already know, in addition to the hypothesis $\mathrm{c}_n$, that we have $\mathrm{a}_{n-2}$, i.e. $Z$ is $(n-2)$-connected. We have to prove that it is even $(n-1)$-connected, i.e. $\pi_{n-1}(Z)=0$. The case $n=2$ requires % [sic] L’enoncé -> L’énoncé, missing accent on é
\msp{page 186}
perhaps special attention, since we do not even know yet that $Z$ is 1-connected. But we know that $\mathrm{c}_2$ implies $\mathrm{c}_1$, which is equivalent \texttt{[?]} to $\mathrm{b}_1$, so we have
\[
\pi_1(X,x) \longrightarrow \pi_1(Y,y) \qquad \text{isomorphism,}
\]
so
\[
\pi_1(Z,x) \longrightarrow \pi_1(X,x) \qquad \text{\texttt{[is]} \textit{zero,}}
\]
so
\[
\pi_2(Y,y) \longrightarrow \pi_1(Z,x) \qquad \text{\texttt{[is an]} epimorphism,}
\]
which implies at least that $\pi_1(Z,x)$ is \textit{abelian.} We must view it as a $\pi$-module, and prove its vanishing as such (using $\mathrm{c}_2$).

\printedpage{124}
% head: A. Grothendieck : Dérivateurs, Ch. XII, Édition des Universités de Montpellier II et Paris VII

Let us pass to the general case. We have for $F$ locally constant abelian on $Y$
\[
H^*(X,F_X) \Longleftarrow \underbrace{E_2^{pq} = H^p(Y,H^q(Z,F_Z))}_{\begin{array}{l}\scriptstyle = 0 \text{ if } 0<q<n-1\\ \scriptstyle \text{since } Z \text{ is } (n-2)\text{-connected}\end{array}}, \qquad
\begin{tikzpicture}[baseline=0pt, x=0.6cm, y=0.6cm]
\draw[->] (0,0) -- (0,5);
\draw[->] (0,0) -- (6,0);
\node[left] at (0,4) {$n-1$};
\foreach \i in {0,...,4} { \fill (\i,4) circle (2.5pt); \fill (\i,0) circle (2.5pt); }
\node at (1,2) {$0$};
\node at (3,2) {$0$};
\end{tikzpicture}
\qquad ,
\]
Hence
\[
0 \longrightarrow H^{n-1}(Y,F) \longrightarrow H^{n-1}(X,F_X) \longrightarrow \underbrace{H^0(Y,H^{n-1}(Z,F_Z))}_{\begin{array}{l}\scriptstyle \mathrm{Hom}_\pi(\pi_{n-1}(Z,x),F_x)\\ \scriptstyle \text{by Hurewicz}\end{array}} \longrightarrow H^n(Y,F) \longrightarrow H^n(X,F_X).
\]
The hypothesis $\mathrm{c}_n$ implies that we have
\msp{page 187}
\[
\mathrm{Hom}_\pi(\pi_{n-1}(Z,x),\underbrace{M}_{=\,F_x}) = 0
\]
and this \textit{whatever the $\pi$-module $M$,} since $\pi = \pi_1(X,x) \xrightarrow{\ \sim\ } \pi_1(Y,y)$. This implies (taking e.g. $M = \pi_{n-1}(Z,x)$!) that $\pi_{n-1}(Z,x)=0$. (But keeping in mind, in the case $n=2$, that we already know that $\pi_1(Z,x)$ is abelian).

So the theorem corrected on page 178 does look true. It is a bit sad, since in each of the implications % [sic] rectifie -> rectifié, past participle intended
\[
\mathrm{a}_n \ \Longrightarrow\ \mathrm{b}_n \ \Longrightarrow\ \mathrm{c}_n \ \Longrightarrow\ \mathrm{a}_{n-1}
\]
we lose a smidgen, and in the end we come out with the magnificent implication
\[
\mathrm{a}_n \ \Longrightarrow\ \mathrm{a}_{n-1}.
\]
In any case, it is $\pm$ clear that the ``right'' localizer $W_n$, the one that is to define the ``$n$-homotopy types'' as objects of $\mathrm{Hot}_{W_n}$, is defined by the condition $\mathrm{b}_n$.
\DERmark{ENV0784}


\section*{2\quad Criterion for $W_n$-asphericity.}

I am going to try to do better when $Y$ is reduced to a point. This time, I


\endgroup
% END UNIT Der-ChXIIc-30-03-03/Der-ChXIIc-30-03-03_p006-010.tex
% BEGIN UNIT Der-ChXIIc-30-03-03/Der-ChXIIc-30-03-03_p011-015.tex
\begingroup\makeatletter%
\newcommand{\printedpage}[1]{\marginpar{\footnotesize p.~#1}}%
\newcommand{\msp}[1]{\marginpar{\footnotesize [#1]}}%
\newcommand{\Ch}{\underline{\mathrm{Ch}}}%
\newcommand{\Fib}{\underline{\mathrm{Fib}}}%
\newcommand{\Fl}{\underline{\mathrm{Fl}}}%
\newcommand{\Hom}{\underline{\mathrm{Hom}}}%
\newcommand{\Cat}{\mathrm{Cat}}%
\makeatother\let\printedpage\derpp\let\msp\dermsp%
\let\mspin\DERmspin%


\printedpage{125}
% head: A. Grothendieck : Dérivateurs, Ch. XII, Édition des Universités de Montpellier II et Paris VII
% foot: Version 30 mars 2003

\msp{page 188}
want to be strict, and give (if I can) complete proofs in the framework in which I place myself at present, that of the modelizer Cat, and using only what I know how to prove there directly at present, without referring to semi-simplicial or topological theory.

\textbf{Theorem 2} (?). \textit{Let $X$ be in} Cat, \textit{$n \in \mathbf{N}$ an integer $\geq 1$. Then the following conditions are equivalent:}

\begin{itemize}
\item[\textit{$a_n$)}] \textit{$X$ is $0$-connected, and if $x$ is in $X$, $\pi_i(X,x) = 0$ for $i \in [1,n]$ (a condition which does not depend on the choice of $x$), i.e. the $\Omega^i(X,x)$ for $i \in [1,n]$ are $0$-connected.}
\item[\textit{$b_n$)}] \textit{For every group $G$, $G \longrightarrow H^0(X,G_X)$ is an isomorphism, and $H^1(X,G_X) = 0$. Moreover, if $G$ is abelian, we also have $H^i(X,G_X) = 0$, for $i \in [1,n]$.}
\end{itemize}

\textsc{Proof}. For $n = 1$ this equivalence is $\pm$ tautological and was recalled (in a slightly different form) at the beginning (p.~171). We proceed by induction. We may therefore assume that $n \geq 2$, that $X$ satisfies $\mathrm{a}_{n-1}$, $\mathrm{b}_{n-1}$, and prove that then $\mathrm{a}_n$, $\mathrm{b}_n$ are equivalent, i.e. $\pi_n(X,x) = 0$ is equivalent to $H^n(X,M) = 0$ for every abelian group $M$.
\DERmark{ENV0785}


\msp{page 189}
We use the $W_\infty$-fibration (perfect)
\[
\begin{tikzcd}[row sep=large]
E = \Ch(X;x,-) \arrow[d] & E\ W_\infty\text{-aspheric} \\
X &
\end{tikzcd}
\]
with fiber $Z$ at $x$ equal to $\Omega(X,x)$.

Suppose $\mathrm{a}_n$) for $X$, hence $a_{n-1}$) for $\Omega(X,x)$. By the induction hypothesis, we therefore have
\[
H^i(Z,M) = 0 \qquad \text{for } i \in [1,n-1],
\]
$M$ an abelian group. So in the Leray spectral sequence,
\[
H^*(E,M) \Longleftarrow E_2^{pq} = H^p(X,H^q(Z,M))
\]
we have
\[
E_2^{pq} = 0 \qquad \text{if } q \in [1,n-1], \qquad
\begin{tikzpicture}[baseline=0]
\draw[->] (0,0) -- (2.2,0);
\draw[->] (0,0) -- (0,2.1);
\foreach \x in {0,0.35,0.7,1.05,1.4} {
  \fill (\x,0) circle (1.5pt);
  \fill (\x,1.4) circle (1.5pt);
}
\node[left] at (-0.05,1.4) {$n$};
\node at (0.35,0.7) {$0$};
\node at (1.05,0.7) {$0$};
\end{tikzpicture}
\quad .
\]

\printedpage{126}
% head: A. Grothendieck : Dérivateurs, Ch. XII, Édition des Universités de Montpellier II et Paris VII
% foot: Version 30 mars 2003

We therefore have the exact sequence
\[
0 \longrightarrow H^n(X,M) \longrightarrow H^n(E,M)
\]
and since $E$ is $W_\infty$-aspheric, we have $H^n(E,M) = 0$, hence $H^n(X,M) = 0$.

Conversely, suppose $\mathrm{b}_n$), let us prove $\mathrm{a}_n$). We already know that we have $\mathrm{a}_{n-1}$), hence $\mathrm{a}_{n-2}$ for $\Omega(X,x)$, which implies
\[
H^i(Z,M) = 0 \qquad \text{for } i \in [1,n-2],
\]
hence
\[
E_2^{pq} = 0 \qquad \text{if } q \in [1,n-2].
\]
We therefore have the exact sequence

\msp{page 190}
\[
0 \longrightarrow \underbrace{H^{n-1}(X,M)}_{=\,0} \longrightarrow \underbrace{H^{n-1}(E,M)}_{=\,0} \longrightarrow H^0(X,H^{n-1}(Z,M)) \longrightarrow \underbrace{H^n(X,M)}_{=\,0} \longrightarrow H^n(E,M),
\]
hence
\[
H^0(X,\ \underbrace{H^{n-1}(Z,M)}_{\begin{array}{l}\text{local system}\\ \textit{constant}\text{ on } X,\\ \text{since } X \text{ is 1-connected}\end{array}}\ ) = 0,
\]
hence
\[
H^{n-1}(Z,M) = 0,
\]
which, by the induction hypothesis, implies $\pi_{n-1}(Z,z) = 0$, hence $\pi_n(X,x) = 0$, q.e.d.

\textbf{Corollary.} \textit{For $X$ to be $W_\infty$-aspheric, i.e. for us to have $G \xrightarrow{\ \sim\ } H^0(X,G)$, $H^1(X,G) = 0$} \texttt{[for]} \textit{every group} \texttt{[}$G$\texttt{]}\textit{, and $H^i(X,G) = 0$ if $i \geq 2$ for $G$ abelian, it is necessary and sufficient that} $1^\circ$) \textit{$X$ be} $0$\textit{-connected, and ($x$ being a point of $X$)} $2^\circ$) \textit{the $\Omega^i(X,x)$} \texttt{[be]} $0$\textit{-connected for $i \geq 1$.}

\msp{page 191}

\section*{3\quad On the philosophy of homotopy fibers.}

The criterion of asphericity (for $W_\infty$) that I have just obtained, thanks to the functor $\Ch(X)$ (or at one's choice, $\Ch_\infty(X)$) and its derivatives $\Omega(X,x)$ and the iterates, removes all doubt (if any could still remain) that the ``path spaces'' that I have constructed (quite laboriously!) are indeed ``the right ones''. I would now like to get as precise an idea as possible of their ``canonicity'', to what extent these constructions are independent of a choice (e.g. that between the theory $\Ch$, and $\Ch_\infty$). % [sic] aux foncteur -> au foncteur, number agreement

Everything comes down, in the end, to questions of factorization of an arbitrary arrow
\[
f : X \longrightarrow Y
\]
\printedpage{127}
% head: A. Grothendieck : Dérivateurs, Ch. XII, Édition des Universités de Montpellier II et Paris VII
% foot: Version 30 mars 2003
of Cat as
\[
X \xrightarrow{\ i_f\ } \bar{X} \xrightarrow{\ \bar{f}\ } Y,
\]
with
\[
i_f \in W, \qquad \bar{f} \in W\text{-fib}
\]
(the fact that $\bar{f}$ can even be chosen $W_\omega$-perfect, not seeming essential to me here).

This is indeed what is involved in constructing

\msp{page 192}
an object ``paths in $X$''
\[
\Ch(X),
\]
by factoring the diagonal map
\[
X \xrightarrow{\ \mathrm{diag}_X\ } X \times X,
\]
and also when we construct $\Omega(X,x)$ as the fiber at $x$ of a $W$-fibration
\[
E(X,x) \longrightarrow (X,x)
\]
which factors the inclusion
\[
\{x\} \longrightarrow X
\]
(thus $E(X,x)$ being $W$-aspheric).

We can be more or less demanding, and consider, either for a ``target'' $Y$ (or rather an $S$) fixed in Cat, a functorial factorization
\[
C : \Cat/S \longrightarrow \Fib_W S \hookrightarrow \Cat/S
\]
of the structural morphism $X \longrightarrow S$ of an object, as
\[
X \xrightarrow{\ i_X\ } C(X) \longrightarrow S,
\]
i.e. we give ourselves the functor $C$, \textit{plus}

\msp{page 193}
the functorial morphism
\[
i_C : \mathrm{id}_{\Cat/S} \longrightarrow C
\]
with
\[
(i_X : X \longrightarrow C(X)) \in W \qquad \forall\ X \text{ in } \Cat/S,
\]
or be still more demanding and vary the source and target of $f : X \longrightarrow Y (= S)$ at the same time, and ask for a functor
\[
\begin{aligned}
C : \Fl(\Cat) \quad &\longrightarrow \quad \Cat \\
f \quad &\longmapsto \quad C(f),
\end{aligned}
\]
\printedpage{128}
% head: A. Grothendieck : Dérivateurs, Ch. XII, Édition des Universités de Montpellier II et Paris VII
% foot: Version 30 mars 2003
with morphisms functorial in $f$
\[
\underbrace{s(f)}_{=\,X} \xrightarrow{\ i_f\ } \underbrace{C(f)}_{=\,\bar{X}} \xrightarrow{\ p_f\ } \underbrace{b(f)}_{=\,Y}
\]
($i$, $p$ arrows in $\Hom(\Fl(\Cat),\Cat)$), with
\[
i_f \in W, \qquad p_f \in W\text{-fib}
\]
for every arrow $f$ in Cat.

\msp{page 194}
In both points of view, in what sense do we have \textit{uniqueness} of such a construction?

Let us first take the first point of view, factorization of arrows with fixed target $S$, hence \texttt{[this is a]} question in $\Cat/S$. I forgot an important property, but it follows from what has been said: the functor
\[
\begin{tikzcd}[column sep=small]
X \longmapsto C(X) : & \Cat/S \arrow[rr] \arrow[dr, dotted] & & \Cat/S \\
& & \Fib_W S \arrow[ur, dotted, hook] &
\end{tikzcd}
\]
\textit{transforms} \texttt{[the]} \textit{arrows} \texttt{[of]} \textit{$W$ into likewise}, as we see from the commutative diagram
\[
\begin{tikzcd}[row sep=large, column sep=large]
X \arrow[r, "i_X"] \arrow[d, "f"'] & C(X) \arrow[d, "C(f)"] \\
Y \arrow[r, "i_Y"] & C(Y),
\end{tikzcd}
\]
where $i_X$, $i_Y$ are in $W$, hence
\[
f \in W \qquad \Longleftrightarrow \qquad C(f) \in W.
\]
Let
\[
(C',i')
\]
be a second functor, $\Cat/S \longrightarrow \Cat/S$,

\msp{page 195}
equipped with a morphism
\[
i' : \mathrm{id}_{\Cat/S} \longrightarrow C'
\]
which is in $W$ argument by argument, and such that $C'$ factors through $\Fib_W S$. We can compare $C$ and $C'$ by means of $C \circ C'$, for
\[
\begin{tikzcd}[column sep=small]
& C \circ C' & \\
C = C \circ \mathrm{id} \arrow[ur, "C\ast i'"] & & C' = \mathrm{id} \circ C', \arrow[ul, "i\ast C'"']
\end{tikzcd}
\]
\printedpage{129}
% head: A. Grothendieck : Dérivateurs, Ch. XII, Édition des Universités de Montpellier II et Paris VII
% foot: Version 30 mars 2003
i.e. we have, for every $X$ in $\Cat/S$,
\[
\begin{tikzcd}[column sep=small]
& C(C'(X)) & \\
C(X) \arrow[ur, "C(i'_X)"] & & C'(X), \arrow[ul, "i_{C'(X)}"']
\end{tikzcd}
\]
the two arrows being in $W$ (the one on the left because $C(W) \subseteq W$). These are arrows in $\Fib_W S$, and if we assume that $W$ satisfies W(4b), they induce $W$-equivalences in each fiber. This therefore already shows % [sic] le deux flèches -> les deux flèches, typo

\msp{page 196}
that up to isomorphism (perhaps noncanonical), the composite functor defined by $C$,
\[
\begin{tikzcd}
X \arrow[r, mapsto] & C(X) \arrow[r, mapsto] & (s \longmapsto \mathrm{hot}_W(C(X)_s)) \\
\Cat/S \arrow[r] \arrow[rr, bend right, "\varphi_C"'] & \Fib_W S \arrow[r] & \Hom(S^\circ,\mathrm{Hot}_W),
\end{tikzcd}
\]
does not depend on the particular choice of $C$, i.e. two such ``homotopy fibers'' functors (corresponding to two choices $C$, $C'$) are always isomorphic. But instead of using $C \circ C'$ to define such an isomorphism, we could have used $C' \circ C$, or some other procedure. We would like to understand whether there is genuine uniqueness, an intrinsic characterization of these isomorphisms $\varphi_C \xrightarrow{\ \sim\ } \varphi_{C'}$, independently of one construction or another. % [sic] aurai pu -> aurait pu, typo

I presume that $C \circ C'$ is itself also equipped with
\[
\mathrm{id}_{\Cat/S} \longrightarrow C \circ C',
\]

\msp{page 197}
and even a priori in two possible ways
\[
\begin{tikzcd}[column sep=large, row sep=large]
& C(C'(X)) & \\
C(X) \arrow[ur, "C(i'_X)"] & & C'(X) \arrow[ul, "i_{C'(X)}"'] \\
& X, \arrow[ul, "i_X"] \arrow[ur, "i'_X"'] \arrow[uu, dotted] &
\end{tikzcd}
\]
either by $C(i'_X) \circ i_X$, or by $i_{C'(X)} \circ i'_X$, fortunately they are equal, since for every arrow $X \xrightarrow{\ i'\ } X'$ in $\Cat/S$, the square % [sic] i_{C(X)} -> i_{C'(X)}, the diagram above has i_{C'(X)}
\DERmark{ENV0786}



\endgroup
% END UNIT Der-ChXIIc-30-03-03/Der-ChXIIc-30-03-03_p011-015.tex
% BEGIN UNIT Der-ChXIIc-30-03-03/Der-ChXIIc-30-03-03_p016-020.tex
\begingroup\makeatletter%
\newcommand{\printedpage}[1]{\marginpar{\footnotesize p.~#1}}%
\newcommand{\msp}[1]{\marginpar{\footnotesize [#1]}}%
\newcommand{\Hom}{\underline{\mathrm{Hom}}}%
\newcommand{\Fib}{\underline{\mathrm{Fib}}}%
\newcommand{\Cat}{\mathrm{Cat}}%
\newcommand{\HOT}{\mathrm{HOT}}%
\newcommand{\Hot}{\mathrm{Hot}}%
\newcommand{\hot}{\mathrm{hot}}%
\newcommand{\id}{\mathrm{id}}%
\makeatother\let\printedpage\derpp\let\msp\dermsp%
\let\mspin\DERmspin%


\printedpage{130}
% head: A. Grothendieck : Dérivateurs, Ch. XII, Édition des Universités de Montpellier II et Paris VII

\[
\begin{tikzcd}[column sep=small, row sep=small]
 & C(X') & \\
C(X) \arrow[ur, "C(i')"] & & X' \arrow[ul, "i_{X'}"'] \\
 & X \arrow[ul, "i_X"] \arrow[ur, "i'"'] &
\end{tikzcd}
\]

is commutative (by functoriality of $i_X : X \longrightarrow C(X)$), and we apply this to the case of $X' = C'(X)$, $i' = i'_X$.

This remark encourages us to make a leap into abstraction, and to introduce the category\msp{page 198}
\[
\mathcal{C} = \mathcal{C}_S = \begin{cases}
\text{category of pairs } (C,i), \text{ with} \\
C : \Cat/S \longrightarrow \Cat/S \text{ factoring through } \Fib_W S, \\
i : \id_{\Cat/S} \longrightarrow C \text{ such that } i_X \in W \ \forall\ X \text{ in } \Cat/S,
\end{cases}
\tag*{$(*)$}
\]
an arrow
\[
(C,i) \longrightarrow (C',i')
\]
being an arrow
\[
C \xrightarrow{u} C' \qquad \text{(in } \Hom(\Cat/S, \Cat/S))
\]
making the square \texttt{[read: the diagram]} commutative
\[
\begin{tikzcd}[column sep=small]
 & \id_{\Cat/S} \arrow[dl, "i"'] \arrow[dr, "i'"] & \\
C \arrow[rr, "u"] & & C'.
\end{tikzcd}
\]

In other words, $\mathcal{C}$ is defined as a full subcategory of % [sic] sous-catégiorie -> sous-catégorie, misspelled word
\[
\id_{\Cat/S}\backslash\Hom(\Cat/S, \Cat/S),
\]
namely the one whose objects $(C,i)$ satisfy the two conditions $(*)$ above (both involving the localizer $W$).

\msp{page 199}These conditions imply this

\begin{enumerate}
\item[a)] $\forall\ X$ in Cat/$S$, $(u_X : C(X) \longrightarrow C'(X)) \in W$, since
\[
\begin{tikzcd}[column sep=small]
 & X \arrow[dl, "i_X\in W"'] \arrow[dr, "i'_X\in W"] & \\
C(X) \arrow[rr, "u_X"] & & C'(X),
\end{tikzcd}
\]
\printedpage{131}
% head: A. Grothendieck : Dérivateurs, Ch. XII, Édition des Universités de Montpellier II et Paris VII
and

\item[b)] if $W$ satisfies W(4b), so that $u_X \in W^{\mathrm{f}}_S$, then this defines an \textit{isomorphism} between $\varphi_C(X)$ and $\varphi_{C'}(X)$ % [sic] ce définit -> ceci définit, missing word
\[
\begin{tikzcd}[row sep=small]
s \arrow[r, mapsto] & \hot_W(C(X)_s) \arrow[d, "\wr\, u_{X,s}"] \\
s \arrow[r, mapsto] & \hot_W(C'(X)_s).
\end{tikzcd}
\]
\end{enumerate}

Better still, in this case, the composite functor
\[
\mathcal{C} \times \Cat/S \longrightarrow \Fib_W S \longrightarrow \HOT^{\mathrm{lc}}_W(S)
\]
gives rise to
\[
\mathcal{C} \longrightarrow \Hom(\Cat/S, \HOT^{\mathrm{lc}}_W(S))
\]
which transforms every arrow of $\mathcal{C}$ into \texttt{[an]} \textit{isomorphism,} hence factors through the fundamental groupoid of $\mathcal{C}$,
\[
\Pi\mathcal{C} \longrightarrow \Hom(\Cat/S, \HOT^{\mathrm{lc}}_W(S)).
\]

\msp{page 200}As for the functor considered earlier
\[
\begin{aligned}
\mathcal{C} &\quad\longrightarrow\quad \underbrace{\Hom\,(\Cat/S, \Hom(S^\circ, \Hot_W))}_{\simeq\Hom((\Cat/S)\times S^\circ, \Hot_W)} \\
C &\quad\longmapsto\quad ((X,s) \longmapsto \hot_W(C(X)_s))\,,
\end{aligned}
\]
it is deduced from the preceding one by composing with
\[
\begin{tikzcd}[column sep=small]
\HOT^{\mathrm{lc}}_W(S) \arrow[rr] \arrow[dr, dotted] & & \Hom(S^\circ, \Hot_W), \\
 & \Hom(\Pi S^\circ, \Hot_W) \arrow[ur, dotted] &
\end{tikzcd}
\]
\DERmark{ENV0787}

more exactly, with the functor deduced from it
\[
\Hom(\Cat/S, \HOT^{\mathrm{lc}}_W(S)) \longrightarrow \Hom(\Cat/S, \Hom(S^\circ, \Hot_W)).
\]
\DERmark{ENV0788}


We see emerging here a plausible principle of canonicity of isomorphisms, as regards the functors
\[
C_W : \Cat/S \longrightarrow \HOT^{\mathrm{lc}}_W(S)
\]
defined by the $C \in \mathrm{Ob}\,\mathcal{C}$, by the relation
\[
\Pi\mathcal{C} \approx \ \text{one-point category, i.e. } \mathcal{C} \text{ is 1-connected.}
\]

Moreover, the arguments used so far are of so formal a nature, that the specific aspects of the Modelizer Cat
\printedpage{132}
% head: A. Grothendieck : Dérivateurs, Ch. XII, Édition des Universités de Montpellier II et Paris VII
\msp{page 201}seem to play practically no role in them, or a very unobtrusive role.

The square on page 197 shows above all that the category $\mathcal{C}$ is filtering - $(C,i)$, $(C',i')$ both map into $(C'',i'')$, with $C'' = C \circ C'$, and $i'' = (C * i') \circ i = (i * C') \circ i'$. There is even a law of composition in $\mathcal{C}$, by no means commutative or unital, that's for sure, but surely associative, and (if this composition is denoted by $*$) with functorial morphisms % [sic] (C,i') -> (C',i'), the second couple of the following formula is (C',i')
\DERmark{ENV0789}

\[
\xi \xrightarrow{\alpha_{\xi,\xi'}} \xi * \xi' \xleftarrow{\beta_{\xi,\xi'}} \xi'
\]
which ensure the upward filtering character of $\mathcal{C}$. These arrows seem somehow to replace the missing two-sided unit, and one expects them to have the same excellent properties as if one had such a unit, concerning the double composites $(\xi * \xi') * \xi'' \simeq \xi * (\xi' * \xi'')$,

\msp{page 202}
\[
(\xi * \xi') * \xi'' \xrightarrow{\sim} \xi * (\xi' * \xi'').
\]

There are two ways of mapping a binary product
\[
\xi * \xi', \qquad \xi' * \xi'', \qquad \xi * \xi''
\]
into the ternary product, and well, they must in each of the three cases, be equal. And likewise there are, this granted, again two ways of mapping $\xi$ into the ternary product, factoring either through $\xi * \xi'$, or through $\xi * \xi''$, and likewise for the other two functors $\xi'$, $\xi''$; and there again, these two ways must be equal, in each of these three cases - so $3 + 3 = 6$ equalities in all to write down and to check - we've seen worse!

The computational philosophy should be this: If we have an $n$-fold product
\[
\xi_1 * \xi_2 * \cdots * \xi_n
\]
(where I dispense with putting in parentheses \dots), for every partial $m$-fold product $(m \in [1, n-1])$\msp{page 203}
\[
\xi_{i_1} * \xi_{i_2} * \cdots * \xi_{i_m} \qquad\qquad (i_1 < i_2 < \cdots < i_m),
\]
using the associativity isomorphisms and the $\alpha_{\xi,\xi'}$, $\beta_{\xi,\xi'}$ left and right, there is nevertheless only one way of mapping it into the $n$-fold product. I really have the impression that this must be true. (Does this already imply that $\mathcal{C}$ is 1-connected?)

In fact, there is indeed a two-sided unit, but it is not in $\mathcal{C}$, but in the larger category
\[
\id_{\Cat/S}\backslash\Hom(\Cat/S, \Cat/S),
\]
which itself, like any category
\[
\id_{\mathcal{X}}\backslash\Hom(\mathcal{X}, \mathcal{X}),
\]
\printedpage{133}
% head: A. Grothendieck : Dérivateurs, Ch. XII, Édition des Universités de Montpellier II et Paris VII
is a monoidal category, i.e. with an associative \textit{and} unital product. The arrows $\alpha_{\xi,\xi'}$, $\beta_{\xi,\xi'}$ are the arrows that result from the presence of this unit (seeing that $\mathcal{C}$ is a full subcategory), and then it should become clear that we have the said compatibilities.

\msp{page 204}But even a monoidal category, is it 1-connected? For example $B_G$, for $G$ a group say, is indeed a monoidal category and even ``groupal'', a $G$-category (as strict as can be), and yet it is \textit{not} 1-connected!

Of course, we can short-circuit the question of the 1-connectedness of $\mathcal{C}$, by noting that after all we already have a priori a privileged object in
\[
\Hom(\Cat/S, \HOT^{\mathrm{lc}}_W(S)),
\]
i.e. a canonical functor
\[
\Cat/S \longrightarrow \HOT^{\mathrm{lc}}_W(S),
\]
namely the composite functor
\[
\Cat/S \underbrace{\longrightarrow}_{\substack{\text{localization}\\ \text{by } W_S}} (\Cat/S)W_S^{-1} \longrightarrow \underbrace{\HOT^{\mathrm{lc}}_W(S)}_{\overset{\mathrm{def}}{=} (\Fib_W S)(W^{\mathrm{f}}_S)^{-1}} \quad ,
\tag*{$(*)$}
\]
\DERmark{ENV0790}

where the second arrow is defined as the quasi-inverse of
\[
(\Fib_W S)(W^{\mathrm{f}}_S)^{-1} \longrightarrow (\Cat/S)W_S^{-1}
\tag*{$(**)$}
\]

\msp{page 205}deduced from the inclusion
\[
(\Fib_W S, W^{\mathrm{f}}_S) \hookrightarrow (\Cat/S, W_S).
\]

The fact that $(**)$ is an equivalence, is proved precisely by the fact of the existence of a functor $C : \Cat/S \longrightarrow \Cat/S$ and of an $i : \id_{\Cat/S} \longrightarrow C$, having the properties we know, i.e. from the fact that $\mathcal{C} \neq \emptyset$. Moreover, the datum of an object $C$ of $\mathcal{C}$ gives a convenient way of making the quasi-inverse of the said functor explicit. But this quasi-inverse has an existence that is canonical and independent of the choice of a $C$, i.e. of a mode of calculation. The fact that the composite functor $\varphi$ in $(*)$ can be ``lifted'' (up to isomorphism) to a functor % [sic] C \neq \emptyset -> \mathcal{C} \neq \emptyset, the non-emptiness concerns the category \mathcal{C}
\DERmark{ENV0791}\DERmark{ENV0792}

\[
C : \Cat/S \longrightarrow \Fib_W S
\]
between $\varphi(X) = \hot_{W,S}(X)$ and $\hot_{W,S}(C(X))$ and that the isomorphism can be made explicit, it too, by an $X$-functorial arrow
\[
i_X : X \longrightarrow C(X), \qquad i_X \in W,
\]
in Cat/$S$ itself (and not only in the localized $(\Cat/S)W_S^{-1}$),
\printedpage{134}
% head: A. Grothendieck : Dérivateurs, Ch. XII, Édition des Universités de Montpellier II et Paris VII
\msp{page 206}is a particular virtue of the \texttt{[illegible word]}, which should not for all that make us forget the intrinsic character of the arrows $(*)$, independently of any mode of calculation.

{\reversemarginpar\marginnote{$\|$}}% double vertical bar printed in the left margin beside the first lines of this paragraph
The key to the situation, and to the questions of canonicity, seems to me to be in the fact that $(**)$ is an equivalence of categories - and one must admit that, however familiar and almost banal it has finally become to me, this is a quite remarkable fact and by no means obvious a priori. (But probably one can also get it by the general \textit{fourbis} à la \textsc{Quillen-Thomason}, without fighting with a formalism of paths \dots) The ``canonical isomorphisms'' that we can establish between a ``theory $C$ '' and a ``theory $C'$ '' (via two different constructions, e.g., of categories of paths, or via the \textsc{Quillen-Thomason} approach, without paths) must always come back to those
\DERmark{ENV0793}


\msp{page 207}that we can define via the comparison of the ``theory $C$ '' on the one hand, the ``theory $C'$ '' on the other, with what we get directly from the composite arrow $(*)$ (p. 204). Finally, it seems to me that \textit{the existence} of a theory of the ``homotopy fibers'' (or ``$W$-homotopic'') of an arrow $X \longrightarrow S$ in Cat, results from the equivalence $(**)$ (p. 204), and this theory could probably to a certain extent be developed from there, without a choice of a functor $C$ or of a model category structure in the sense of \textsc{Quillen}. % [sic] qu'on peu définir -> qu'on peut définir, verb form

\section*{4\quad Axiomatization of ``$W$-homotopic fibers''.}

Yesterday's reflection makes me return to the axiomatics in ``categories of models'' % [sic] reflexion -> réflexion, missing accent
\[
(\mathcal{M}, W)
\]
(\footnote[119]{$W$ satisfies the axiom of weak saturation: isomorphisms \texttt{[are]} in $W$, and if two among $f$, $g$, $gf$ \texttt{[are]} in $W$, the third also.}) more general than Cat, in the spirit of VI (cf. theorem-scholie 2, p. 108-113). For the moment, we are going to suppose only $\mathcal{M}$ stable under fiber products (we do not need a final object). We introduce the set
\[
\mathrm{Fib}_W \subseteq \mathrm{Fl}(\mathcal{M})
\]

\msp{page 208}of \textit{$W$-fibrations} (``ultrastrong'' being understood!) by the condition

\begin{quote}
$f : X \longrightarrow Y$ is in $\mathrm{Fib}_W$ if and only if the base change functor by $X$,
\[
\begin{aligned}
Y' &\quad\longmapsto\quad X' = X \times_Y Y' \\
\mathcal{M}/Y &\quad\longrightarrow\quad \mathcal{M}/X
\end{aligned}
\]
\end{quote}


\endgroup
% END UNIT Der-ChXIIc-30-03-03/Der-ChXIIc-30-03-03_p016-020.tex
% BEGIN UNIT Der-ChXIIc-30-03-03/Der-ChXIIc-30-03-03_p021-025.tex
\begingroup\makeatletter%
\derleqno%
\NoAutoSpacing%
\newcommand{\printedpage}[1]{\marginpar{\footnotesize p.~#1}}%
\newcommand{\msp}[1]{\marginpar{\footnotesize [#1]}}%
\newcommand{\cM}{\mathcal{M}}%
\newcommand{\cF}{\mathcal{F}}%
\newcommand{\Ho}{\mathrm{Ho}}%
\newcommand{\ho}{\mathrm{ho}}%
\newcommand{\Fl}{\mathrm{Fl}}%
\newcommand{\Fib}{\underline{\mathrm{Fib}}}%
\newcommand{\Wu}{W^{\mathrm{u}}}%
\newcommand{\lc}{\mathrm{lc}}%
\newcommand{\circled}[1]{\tikz[baseline=(c.base)]\node[draw,circle,inner sep=1pt](c){#1};}%
\makeatother\let\printedpage\derpp\let\msp\dermsp%
\let\mspin\DERmspin%


\printedpage{135}
% head: A. Grothendieck : Dérivateurs, Ch. XII, Édition des Universités de Montpellier II et Paris VII

\begin{quote}
takes $W_Y$ into $W_X$, where $W_X \subseteq \Fl(\cM/X)$ \texttt{[is]} defined by
\DERmark{ENV0794}


$u \in W_X \iff \sigma(u) \in W$ \ ($\sigma : \cM/X \longrightarrow \cM$ the ``source'' functor, or ``forgetting $X$''). % [sic] \sigma(U) -> \sigma(u), the arrow is called u
\DERmark{ENV0795}

\end{quote}

(\footnote[120]{\textbf{NB} Give here the proposition: $\mathrm{Fib}_W$ \textit{stable under composition and under base change, and contains the isomorphisms, stable under direct factors.}}). Let us put, for every $S$ in $\cM$,
\[
\Ho^!_W(S) = (\cM/S)W_S^{-1},
\]
which plays the role of a category of relative $W$-homotopy types \textit{locally constant} over $S$. We would like to describe a contravariant law in $S$ for $\Ho^!_W(S)$, i.e. for
\[
f : S' \longrightarrow S
\]
an arbitrary arrow in $\cM$, define
\[
f^*_W : \Ho^!_W(S) \xrightarrow{?} \Ho^!_W(S'). \tag*{$(*)$}
\]
The difficulty is, of course, that the base change functor
\[
f^* : \cM/S \longrightarrow \cM/S'
\]
is \textit{not} in general compatible with the \msp{page 209} localizers $W_S$, $W_{S'}$. The notion of $W$-fibration thus serves in the first place to bring out the most obvious condition on an arrow $f : S' \longrightarrow S$ in $\cM$, for the base change functor to be compatible with the localizers $W_S$, $W_{S'}$ (\textit{that} is the definition), hence induces $(*)$ on the localized categories.

The \textit{``problem of homotopy fibers''} is precisely to define $(*)$ for \textit{every} arrow in $\cM$, not only for $W$-fibrations, in such a way that the $\Ho^!_W(S)$ form, for $S$ % [sic] tellle -> telle, letter doubled
variable, the fibers of a remarkable fibered category over $\cM$. When $\cM$ admits a final object, so that we can regard the arrows
\[
p : e \longrightarrow S \qquad (p \text{ as ``point''})
\]
in $\cM$ as ``points'' of $S$, the ``$W$-homotopy fiber'' of an arbitrary object $X/S$ of $\cM/S$ at the said point will by definition be the element of
\[
\Ho_W = \Ho_W(\cM) \simeq \Ho^!_W(e)
\]
defined by
\[
p^*(\ho^!_{W/S}(X/S)),
\]
\printedpage{136}
% head: A. Grothendieck : Dérivateurs, Ch. XII, Édition des Universités de Montpellier II et Paris VII
\msp{page 210}
where
\[
\ho^!_{W/S} : \cM/S \longrightarrow \Ho^!_W(S) \overset{\mathrm{def}}{=} (\cM/S)(W_S)^{-1}
\]
\DERmark{ENV0796}

is the canonical localization functor.

For every $S$ in $\cM$, let
\[
\Fib_W S \subseteq \cM/S
\]
be the full subcategory formed by the objects $X/S$ of $\cM/S$ such that $X \longrightarrow S$ is $W$-fibrant. We are going to define a variant
\[
\Ho^{\lc}_W(S)
\]
of $\Ho^!_W(S)$, as
\[
\Ho^{\lc}_W(S) \overset{\mathrm{def}}{=} (\Fib_W S)(\Wu_S)^{-1},
\]
\DERmark{ENV0797}

where

\begin{quote}
$\Wu_S \subseteq \Fl(\Fib_W S)$ is the localizer formed by the arrows $u : X \longrightarrow Y$ which are ``universally over $S$ in $W$'', i.e. such that for every arrow $S' \longrightarrow S$ of $\cM$ with target $S$, the arrow deduced by base change,
\[
u' : X' = X \times_S S' \longrightarrow Y' = Y \times_S S'
\]
is in $W_{S'}$.
\end{quote}

We have ``done what was needed'' so that $\Ho^{\lc}_W(S)$ is contravariant in $S$ for \textit{arbitrary} arrows $f : S' \longrightarrow S$ of $\cM$, they define
\[
f^{*\,\lc}_W : \Ho^{\lc}_W(S) \longrightarrow \Ho^{\lc}_W(S'), \tag*{$(**)$}
\]
\msp{page 211}
since the $W$-fibrations are stable under base change, and since consequently $f^* : \cM/S \longrightarrow \cM/S'$ induces
\[
f^*_{\Fib_W} : \Fib_W S \longrightarrow \Fib_W S'
\]
which, by definition of $\Wu_S$ and $\Wu_{S'}$, is compatible with the localizers
\[
f^*_{\Fib_W} : (\Fib_W S, \Wu_S) \longrightarrow (\Fib_W S', \Wu_{S'}),
\]
and thus passes to the categories of fractions as an arrow $(**)$.
\DERmark{ENV0798}


We may wonder why not be more generous, and define $\Ho^{\lc}_W(S)$ as $(\cM/S)(\Wu_S)^{-1}$, defining $\Wu_S$ in $\cM/S$ in the same way as just now in $\Fib_W S$ - which costs no more. A first answer is from experience the localizer $\Wu_S$ in the whole of $\cM/S$ turns out to be practically unusable, on the one hand because membership in $\Wu_S$ (for an arrow between objects of $\cM/S$ which are not very special, e.g. $W$-fibered over $S$) is unverifiable (except for isomorphisms and things like that \dots), and because on the other hand the category of fractions is prohibitively
\printedpage{137}
% head: A. Grothendieck : Dérivateurs, Ch. XII, Édition des Universités de Montpellier II et Paris VII
\msp{page 212}
big and can serve (it seems) no purpose. Another answer is that we want the category obtained to correspond well to the intuition of relative homotopy types \textit{locally constant} over $S$ - and we therefore want that for an arbitrary base change $f : S' \longrightarrow S$ by an

\begin{quote}
$S'$ in $\cM/S$, whence a composite functor
\[
\begin{tikzcd}[column sep=small]
\cF_S(\Wu_S)^{-1} \arrow[r, "f^{*\,\lc}_W"] \arrow[rr, dotted, bend right=15, "\varphi_{S'}"'] & \cF_{S'}(\Wu_{S'})^{-1} \arrow[r] & \Ho_W
\end{tikzcd}
\]

(\footnote[121]{$\cF$ here denotes a fibered subcategory of $\cM/\cM = \underline{\Fl}(\cM)$ over $\cM$, a possible substitute for the one formed by $\Fib_W S$ or by the whole of $\cM/S$.}) deduced from the functors of categories with localizers
\[
(\cF_S, \Wu_S) \xrightarrow{f^*} (\cF_{S'}, \Wu_{S'}) \xrightarrow{\text{forgetting } S'} (\cM, W)
\]
\end{quote}

(composite functor $\varphi_{S'}$ which can be viewed as the functor \textit{``homotopy fiber at the object $S'$ over $S$''}, this object being visualized as a ``point'' of $S$ ``with values in $S'$ '' following the general functorial yoga \dots), the functor
\[
\begin{aligned}
S' &\longmapsto \varphi_{S'}(\xi) \qquad \text{for } \xi \text{ in } \cF_S(\Wu_S)^{-1}\\
\cM/S &\longrightarrow \Ho_W
\end{aligned}
\]
take arrows in $W_S$ to isomorphisms: the ``$W$-homotopy fiber'' at $S'/S$ of an object of $\cF_S(\Wu_S)^{-1}$ (viewed as a $W$-homotopy type \textit{locally constant} over $S$) must be ``the same'', if $S'$ replaced by $S''$ with $S'' \longrightarrow S'$ given and in $W_S$. It is for this reason that it is reasonable to take $\cF_S$ ``small enough'' (namely $\Fib_W S$, which is \msp{page 213} exactly adjusted for this), rather than take the whole of $\cM/S$.

We can therefore say that the category $\Fib_W S$, and the localizer $\Wu_S$ have been constructed ad-hoc, as simply as possible, to give rise to a localized $\Ho^{\lc}_W(S)$ which corresponds to the intuition of locally constant relative $W$-homotopy types in the following two ways:

\begin{enumerate}[label={}]
\item[1°)] There is spontaneously a contravariant law in $S$, for arbitrary arrows $S' \longrightarrow S$ in $\cM$ (we have defined $\Wu_S$ just for this purpose accordingly).
\item[2°)] that the notion of ``$W$-homotopy fiber at $S'/S$'' of an object $\xi$ of $\Ho^{\lc}_W$ satisfy the preceding condition, i.e. an arrow $S'' \longrightarrow S'$ in $\cM/S$ which is in $W_S$ define an \textit{isomorphism} (in $\Ho_W$) of homotopy fibers
\[
\underbrace{\xi_{S''}}_{\varphi_{S''}(\xi)} \longrightarrow \underbrace{\xi_{S'}}_{\varphi_{S'}(\xi)}
\]
\DERmark{ENV0799}

(and we have chosen $\cF_S = \Fib_W(S)$ just for that purpose).
\end{enumerate}

\printedpage{138}
% head: A. Grothendieck : Dérivateurs, Ch. XII, Édition des Universités de Montpellier II et Paris VII

We thus have an excellent notion of $W$-homotopy fibers for $X$ $W$-fibered over $S$ - this time excellent for \textit{arbitrary} base changes - and there is no merit in that since we have defined everything on purpose for that! % [sic] quelconque -> quelconques, agreement with "des changements"

\msp{page 214}
To put it otherwise, we have defined \textit{two} notions of relative $W$-homotopy types over $S$, one so as to give such a homotopy type $\ho^!_{W,S}(X)$ for \textit{every} object of $\cM/S$, but there is a priori a good notion of base change only for the base changes $S' \longrightarrow S$ which are $W$-fibered, the other, which corresponds more directly to the intuition of a locally constant $W$-homotopy type over $S$, associates such a type $\ho^{\lc}_{W,S}(X)$ only to the $X$ of $\cM/S$ which are $W$-fibered over $S$, on the other hand this notion admits a good notion of inverse image and of ``$W$-homotopy fiber'' for arbitrary base changes.

A ``good'' theory of $W$-homotopy fibers must combine the two preceding notions into a single one, so that the advantages of the two add up, and the disadvantages cancel out! It is obtained with the help of a draconian axiom on the situation (whereas until now we have used nothing about $W \subseteq \Fl(\cM)$, except at most weak saturation, in addition to the existence of fiber products):

\msp{page 215}
\circled{a} (\textit{Homotopy fiber axiom.}) For every object $S$ in $\cM$, the canonical functor
\[
\Ho^{\lc}_W(S) \longrightarrow \Ho^!_W(S)
\]
deduced from the inclusion
\[
(\Fib_W(S), \Wu_S) \longrightarrow (\cM/S, W_S)
\]
is an \textit{equivalence of categories.}

This axiom implies that in $\Fib_W S$, $\Wu_S$ and the localizer induced by $W_S$ (which I still denote by $W_S$ by abuse of notation) have the same strong saturation. In fact, it is natural to supplement the preceding axiom with the following one.

\circled{b} (Adjunct axiom to the homotopy fiber axiom.) On $\Fib_W S$, the localizers $\Wu_S$ and $W_S$ coincide, i.e. we have this:

\begin{quote}
Let $X$, $Y$ be in $\Fib_W S$, $X \xrightarrow{u} Y$ an $S$-morphism such that $u \in W_S$, then $u \in \Wu_S$ (i.e. the morphism $u' : X' \longrightarrow Y'$ deduced by any base change $S' \longrightarrow S$ is still in $W$).
\end{quote}

\msp{page 216}
I presume that a) will rarely be satisfied without b) being so (and conversely for that matter), and that the proof of a) and b) will be done practically in a single breath.

If $\cM = \mathrm{Cat}$, $W$ a basic localizer, then the notion of $W$-fibration adapted here (read: ``ultrastrong'', in the specific terminology developed in Cat) is so strong, that it already practically implies b) (assuming W(4a$'$), it is true, cf. p. 53), and
\printedpage{139}
% head: A. Grothendieck : Dérivateurs, Ch. XII, Édition des Universités de Montpellier II et Paris VII
a) is hardly verifiable unless we have W(4a) (\footnote[122]{No, it seems to me that W(4b) should suffice \dots}) (loc. cit.). So here it is ``the fibration axiom'' on $W$, in its strong form W(4a), which is the keystone of a good theory of homotopy fibers.

Such a theory (returning to the general case) amounts to saying, in essence, that there are ``enough'' $W$-fibrations (``ultrastrong'') to be able to describe \textit{every} object of a $\cM/S$, up to $W_S$-equivalence, as a $W$-fibrant object over $S$. In the cases known to me, we have moreover this: we \msp{page 217} have for every $S$ in $\cM$ a subcategory $\cF_{0\,S}$ of $\Fib_W S$ (I am thinking for example of $\underline{\mathrm{Parf}}_W S$, or even $\underline{\mathrm{Parf}}_{W_\omega} S$, in the case where $\cM = \mathrm{Cat}$), such that for every full subcategory $\cF$ of $\cM/S$ such that
\[
\cF_{0\,S} \subseteq \cF \subseteq \cM/S
\]
(a fortiori, for a full $\cF$ such that $\Fib_W S \subseteq \cF \subseteq \cM/S$), the inclusions induce equivalences of categories
\[
\cF_{0\,S}(W_S)^{-1} \xrightarrow{\sim} \cF(W_S)^{-1} \xrightarrow{\sim} \underbrace{(\cM/S)(W_S)^{-1}}_{\overset{\mathrm{def}}{=} \mathrm{Ho}^!_W(S)}. % [sic] Hot -> Ho, elsewhere the notation is Ho^!_W(S)
\]
\DERmark{ENV0800}\DERmark{ENV0801}

It does not seem possible to fix the choice of $\cF_{0\,S}$ in ``the best possible'' way, once and for all in the axiomatic context adopted here. It is an additional element of structure in $\cM$ (besides the localizer $W$), on the same footing as the datum of the classes of $Q$-fibrations and $Q$-cofibrations, in the manner of \textsc{Quillen}. I do not know whether the axioms a), b) above (homotopy fiber axiom) already imply \msp{page 218} the variant of \circled{a}:

\circled{c} For every full subcategory $\cF$ of $\cM/S$ such that
\[
\Fib_W S \subseteq \cF \subseteq \cM/S
\]
the inclusions induce equivalences
\[
\underbrace{(\Fib_W S)(W_S)^{-1}}_{\begin{subarray}{l} = \ \Ho^{\lc}_W(S)\\ \text{assuming a)}\end{subarray}} \xrightarrow{\sim} \cF(W_S)^{-1} \xrightarrow{\sim} \underbrace{(\cM/S)(W_S)^{-1}}_{=\Ho^!_W(S)}.
\]

But I presume that in practice, when we can verify \circled{a}, we will be able to verify \circled{c} as well (\footnote[123]{Cf. page 240 ff. We should here develop a little more the formalism of $\Ho^{\lc}_W(S)$ and homotopy fibers, under the hypothesis \circled{a}\,\circled{b} \ above.}).


\endgroup
% END UNIT Der-ChXIIc-30-03-03/Der-ChXIIc-30-03-03_p021-025.tex
% BEGIN UNIT Der-ChXIIc-30-03-03/Der-ChXIIc-30-03-03_p026-030.tex
\begingroup\makeatletter%
\newcommand{\printedpage}[1]{\marginpar{\footnotesize p.~#1}}%
\newcommand{\msp}[1]{\marginpar{\footnotesize [#1]}}%
\newcommand{\circled}[1]{\tikz[baseline=(c.base)]\node[draw,circle,inner sep=1pt](c){#1};}%
\newcommand{\deff}{\overset{\mathrm{def}}{=}}%
\makeatother\let\printedpage\derpp\let\msp\dermsp%
\let\mspin\DERmspin%

\DERmark{ENV0802}


\printedpage{140}
% head: A. Grothendieck : Dérivateurs, Ch. XII, Édition des Universités de Montpellier II et Paris VII

\section*{5 Axiomatics of non-locally constant relative $W$-homotopy types $\mathrm{Ho}^{\mathcal{L}}_W(S)$.}

We have said nothing up to this point about possible relative $W$-homotopy types (over an object $S$ of $\mathcal{M}$) which would \textit{not} be ``locally constant''. Yet they were discussed in loc. cit. (VI, theorem-scholie 2, part II of the theorem pages 109-113.) I take the opportunity to come back to this. We must introduce, in each $\mathcal{M}/S$, a second localizer
\[
W^0_S \subseteq \mathrm{Fl}(\mathcal{M}/S), \qquad W^0_S \subseteq W_S
\]
(finer than $W_S$), which will make it possible to define
\msp{page 219}
\[
\mathrm{Ho}_W(S) \deff (\mathcal{M}/S)(W^0_S)^{-1}.
\]
That there is no \textit{canonical} choice which imposes itself here, deduced from the structure $(\mathcal{M},W)$ alone, is clear from the example of $(\mathrm{Cat},W)$, in which case we find \textit{two} symmetric choices, transformed into one another by the involution
\[
X \longmapsto X^\circ : (\mathrm{Cat},W) \xrightarrow{\sim} (\mathrm{Cat},W),
\]
namely $W^{\mathrm{g}}_S$ and $W^{\mathrm{d}}_S$. The choice of one or the other, say (to fix ideas) of $W^{\mathrm{g}}_S$, is associated with a notion of \textit{local isomorphism} (or of ``local co-isomorphism'' in the case of $W^{\mathrm{d}}_S$) which I would like to axiomatize. It does seem that for the main categories of models for Hot (and for its variants), we have in a fairly natural way such a notion of local isomorphism (without the encumbrance, most often, of an involutive automorphism providing a symmetric notion \dots).

We shall therefore suppose, as in loc. cit. p. 94-96, given a class
\[
\mathcal{L} \subseteq \mathrm{Fl}(\mathcal{M})
\]
of morphisms, playing the role of ``arrows of % [sic] jounant -> jouant, faute de frappe
\msp{page 220}
localization'' or ``of local isomorphisms'', satisfying the following conditions.

\begin{enumerate}
\item[$(*)$] $\mathcal{L}$ contains the isomorphisms, is stable under composition, and under base change.
\end{enumerate}

We put, for $S$ in $\mathcal{M}$,
\[
\mathcal{L}_S \subseteq \mathcal{M}/S \qquad \text{full subcategory of the } S'/S \text{ such that } S' \longrightarrow S \text{ \texttt{[is]} in } \mathcal{L}.
\]
We then define for every $S$ in $\mathcal{M}$
\[
\begin{aligned}
&W^{\mathcal{L}}_S \subseteq \mathrm{Fl}(\mathcal{M}/S)\\
&(X \xrightarrow{u} Y) \in W^{\mathcal{L}}_S \iff \forall\, S' \in \mathcal{L}_S,\ (u_{S'} : X \times_S S' \longrightarrow Y \times_S S') \text{ is in } W.
\end{aligned}
\]
% foot: Version 30 mars 2003

\printedpage{141}
% head: A. Grothendieck : Dérivateurs, Ch. XII, Édition des Universités de Montpellier II et Paris VII

This is therefore a variant of the localizer $W^{\mathrm{u}}_S$, where now we take as base changes only the ``$\mathcal{L}$-localizing'' arrows, i.e. $\in \mathcal{L}$. We thus have
\[
W^{\mathrm{u}}_S \subseteq W^{\mathcal{L}}_S \subseteq W_S
\]
as subsets of $\mathrm{Fl}(\mathcal{M}/S)$, and likewise in any full subcategory $\mathcal{F}$ of $\mathcal{M}/S$. In the case of $\underline{\mathrm{Fib}}_W S$, by virtue of axiom \circled{b} above ($W^{\mathrm{u}}_S = W_S$ in $\underline{\mathrm{Fib}}_W S$), the three localizers coincide.

We put
\[
\mathrm{Ho}^{\mathcal{L}}_W(S) \text{ (or } \mathrm{Ho}_W(S)\text{)} \ \deff (\mathcal{M}/S)(W^{\mathcal{L}}_S)^{-1},
\]
and we gave in loc. cit. (part II of the theorem-scholie) an interesting result on this invariant, for fixed $S$. My aim
\msp{page 221}
at present is to see what happens for variable $S$. Visibly, a notion of relative $W$-homotopy type over an arbitrary object $S$ of $\mathcal{M}$ is satisfactory, only if it allows inverse images by \textit{arbitrary} arrows
\[
f : S' \longrightarrow S
\]
in $\mathcal{M}$, which absolutely must define
\[
f^{*\mathcal{L}}_W \text{ or } f^*_W : \mathrm{Ho}^{\mathcal{L}}_W(S) \longrightarrow \mathrm{Ho}^{\mathcal{L}}_W(S').
\]
We are therefore faced with the same perplexity as just now - and it is \textit{not,} this time, the introduction of $\underline{\mathrm{Fib}}_W(S)$ that is going to save the situation, for
\[
(\underline{\mathrm{Fib}}_W S)(\underbrace{W^{\mathcal{L}}_S}_{\substack{= W_S \\ = W^{\mathrm{u}}_S}})^{-1} \longrightarrow (\mathcal{M}/S)(W^{\mathcal{L}}_S)^{-1}
\]
has no desire, this time, to be an equivalence of categories, but much rather an inclusion (fully faithful) of the relative $W$-homotopy types constant over $S$, into arbitrary $W$-homotopy types. An axiom is visibly needed to replace the ``homotopy fiber'' axiom.

\msp{page 222}

In the case of Cat, the class $\mathcal{L}$ of local isomorphisms (say) allows us to define, as above, the localizer $W^{\mathcal{L}}_S = W^{\mathrm{g}}_S$ in the $\mathrm{Cat}/S$, as consisting of the $u : X \longrightarrow Y$ over $S$ which are in $W$ and remain so under every base change $S' \longrightarrow S$ which is in $\mathcal{L}$. But there is reason to \textit{saturate} $\mathcal{L}$, by considering the class $\mathcal{L}^\wedge$ of arrows $f : S' \longrightarrow S$ of Cat such that the base change
\DERmark{ENV0803}

\[
f^* : \mathrm{Cat}/S \longrightarrow \mathrm{Cat}/S'
\]
sends $W^{\mathrm{g}}_S$ into $W^{\mathrm{g}}_{S'}$. This $\mathcal{L}^\wedge$ is tautologically stable under composition, and contains the isomorphisms (better, contains $\mathcal{L}$ = set of local isomorphisms). It is not
% foot: Version 30 mars 2003
\printedpage{142}
% head: A. Grothendieck : Dérivateurs, Ch. XII, Édition des Universités de Montpellier II et Paris VII
clear, from the definition given here, that it is stable under base change - we are going to strengthen the definition accordingly, by requiring, for $f : S' \longrightarrow S$ to be in $\mathcal{L}^\wedge$, that for every $T$ over $S$, let us consider $f_T : T' = S' \times_S T \longrightarrow T$, $f^*_T : \mathrm{Cat}/T \longrightarrow \mathrm{Cat}/T'$ be compatible with the localizers $W^{\mathrm{g}}_T$, $W^{\mathrm{g}}_{T'}$. Then $\mathcal{L}^\wedge$ is (just like $\mathcal{L}$) stable under composition and under base change, and moreover it contains $\mathcal{L}$. Moreover, for every $S$, we have $W^{\mathcal{L}}_S = W^{\mathcal{L}^\wedge}_S$. This little sorites could of course be carried out for any modelizer $(\mathcal{M},W)$. But the interesting thing here is that $\mathcal{L}^\wedge$ contains the set $\mathrm{Liss}_W$ of $W$-smooth arrows. (I do not know whether it is equal.) We could therefore have
\msp{page 223}
just as well started from the set of arrows $\mathrm{Liss}_W$, to define the localizers $W^{\mathrm{g}}_S$ as
\[
W^{\mathrm{g}}_S = W^{\mathrm{Liss}_W}_S \qquad \text{instead of } W^{\mathrm{g}}_S = W^{\mathrm{Isoloc}}_S,
\]
which is equivalent.

On the other hand, in working with the localized categories $\mathrm{HOT}^\circ_W(S) = (\mathrm{Cat}/S)(W^{\mathrm{g}}_S)^{-1}$, the role of $\underline{\mathrm{Fib}}_W S$ in the theory of the $\mathrm{HOT}^{\mathrm{lc}}_W(S)$ is played here by $\underline{\mathrm{Prop}}_W(S)$. We thus see that in the theory of non-locally constant relative $W$-homotopy types, $\mathrm{HOT}^\circ_W$ mode, the sets of arrows $\mathrm{Liss}_W$ and $\mathrm{Prop}_W$ play a role that is in a sense complementary: the arrows $f : S' \longrightarrow S$ of $\mathrm{Liss}_W$ are (practically) those which allow us directly to define
\[
f^{*\circ}_W : \mathrm{HOT}^\circ_W(S) \longrightarrow \mathrm{HOT}^\circ_W(S')
\]
as induced by
\[
f^* : (\mathrm{Cat}/S, W^{\mathrm{g}}_S) \longrightarrow (\mathrm{Cat}/S', W^{\mathrm{g}}_{S'}),
\]
whereas the arrows of $\mathrm{Prop}_W$ allow us to define for every $S$ in Cat the full subcategory
\DERmark{ENV0804}

\[
\underline{\mathrm{Prop}}_W S \subseteq \mathrm{Cat}/S,
\]
giving rise to an equivalence of categories
\[
\left\{
\begin{aligned}
&(\underline{\mathrm{Prop}}_W S)(W^{\mathrm{u}}_S)^{-1} \xrightarrow{\sim} \underbrace{(\mathrm{Cat}/S)(W^{\mathrm{g}}_S)^{-1}}_{= \mathrm{HOT}^\circ_W(S)} % ?? uncertain: exposant de HOT sous l'accolade -- 0 (zéro) ou \circ
\\
&\text{with}\\
&W^{\mathrm{u}}_S = W^{\mathrm{g}}_S \text{ on } \underline{\mathrm{Prop}}_W S
\end{aligned}
\right.
\]
\DERmark{ENV0805}

(which visibly is the announced substitute for ``the $W$-homotopy fiber axiom''). This
\msp{page 224}
equivalence, which allows us to describe $\mathrm{Hot}^\circ_W(S)$ in terms of $\underline{\mathrm{Prop}}_W S \subseteq \mathrm{Cat}/S$ and a localizer $W^{\mathrm{u}}_S$ \textit{stable under} \textbf{every} \textit{base change,} allows us on the other hand to define the expected contravariant law for $\mathrm{HOT}^\circ_W(S)$.
% foot: Version 30 mars 2003

\printedpage{143}
% head: A. Grothendieck : Dérivateurs, Ch. XII, Édition des Universités de Montpellier II et Paris VII

On the other hand, when we interchange the role of $\mathrm{Liss}_W$ and $\mathrm{Prop}_W$, we find the complementary theory of $\mathrm{HOT}_W(S)$,
\[
\left\{
\begin{aligned}
&\mathrm{HOT}_W(S) \deff (\mathrm{Cat}/S)(W^{\mathrm{d}}_S)^{-1} \xleftarrow{\sim} (\underline{\mathrm{Liss}}_W S)(W^{\mathrm{u}}_S)^{-1},\\
&\text{on } \underline{\mathrm{Liss}}_W S,\ W^{\mathrm{d}}_S = W^{\mathrm{u}}_S,
\end{aligned}
\right.
\]
the description of $\mathrm{HOT}_W(S)$ in terms of $\underline{\mathrm{Liss}}_W S$ and $W^{\mathrm{u}}_S$ allows us to define its contravariant character in $S$, for \textit{every} arrow $f : S' \longrightarrow S$ in Cat. Whereas if we want to be able to pass to the categories of fractions on $\mathrm{Cat}/S$ and $\mathrm{Cat}/S'$ in their entirety (via $W^{\mathrm{d}}_S$, $W^{\mathrm{d}}_{S'}$), we must restrict ourselves to $W$-proper base changes.

So the perplexity comes to me whether this complementary, or even this duality, between two sets of arrows, the ``smooth'' and the ``proper'', might not be a general fact which turned up at least in the categories of models close to the description of HOT or of a $\mathrm{Hot}_W$.

\msp{page 225}

It may well be that in $\boldsymbol{\Delta}^\wedge$ we have a perfect duality between these two kinds of notions (I shall have to examine the matter), perhaps even in any category $\mathcal{A}^\wedge$ (but in $\boldsymbol{\Delta}^\wedge$ there is an involutive automorphism $K \longmapsto K^\circ$, corresponding to the $X \longmapsto X^\circ$ in Cat \dots). And after all, even in the context of topological spaces, the notions of properness on the one hand, of local asphericity on the other, for an arrow $X \longrightarrow Y$, are well defined, and might well fit together as do the corresponding notions in Cat, and (who knows) give rise to two complementary versions, or ``dual'' to one another, of a notion of relative homotopy type.

Not being clear about this, I am a little undecided how to present an axiomatic description of the theory of the $\mathrm{Ho}^{\mathcal{L}}_W(S)$, for a general category of models $(\mathcal{M},W)$. Should I give myself from the outset
\[
\mathrm{Liss}, \mathrm{Prop} \subseteq \mathrm{Fl}(\mathcal{M}),
\]
stable under composition, under base change, containing the isomorphisms, and such that (perhaps)
\[
\mathrm{Liss} \cap \mathrm{Prop} \subseteq \mathrm{Fib}_W,
\]
\msp{page 226}
but above all they allow us to define, for $S$ in $\mathcal{M}$
\[
\left\{
\begin{array}{lcccl}
W^{\mathrm{g}}_S & \deff & W^{\mathrm{Liss}}_S & \subseteq & \mathrm{Fl}(\mathcal{M}/S)\\
W^{\mathrm{d}}_S & \deff & W^{\mathrm{Prop}}_S & \subseteq & \mathrm{Fl}(\mathcal{M}/S)
\end{array}
\right.
\]
\DERmark{ENV0806}\DERmark{ENV0807}

with the following properties:

\begin{enumerate}
\item[{\circled{a}}] If $f : S' \longrightarrow S$ is in Prop (resp. Liss), then $f^* : \mathcal{M}/S \longrightarrow \mathcal{M}/S'$ sends $W^{\mathrm{d}}_S$ into $W^{\mathrm{d}}_{S'}$ (resp. $W^{\mathrm{g}}_S$ into $W^{\mathrm{g}}_{S'}$),
\end{enumerate}
% foot: Version 30 mars 2003
\printedpage{144}
% head: A. Grothendieck : Dérivateurs, Ch. XII, Édition des Universités de Montpellier II et Paris VII
\noindent which is a pure tautology, and does not yet concern the relations of Prop with Liss, but on the other hand:

\begin{enumerate}
\item[{\circled{b}}]
\[
\left\{
\begin{aligned}
&(\underline{\mathrm{Liss}}S)(W^{\mathrm{u}}_S)^{-1} \xrightarrow{\sim} (\mathcal{M}/S)(W^{\mathrm{d}}_S)^{-1} \text{ \texttt{[is]} an equivalence}\\
&\text{and } W^{\mathrm{u}}_S = W^{\mathrm{d}}_S \text{ on } \underline{\mathrm{Liss}}S,
\end{aligned}
\right.
\]
and dually
\[
\left\{
\begin{aligned}
&(\underline{\mathrm{Prop}}S)(W^{\mathrm{u}}_S)^{-1} \xrightarrow{\sim} (\mathcal{M}/S)(W^{\mathrm{g}}_S)^{-1} \text{ \texttt{[is]} an equivalence}\\
&\text{and } W^{\mathrm{u}}_S = W^{\mathrm{g}}_S \text{ on } \underline{\mathrm{Prop}}S.
\end{aligned}
\right.
\]
\end{enumerate}

Moreover, there is reason to state \circled{c} two adjunction properties (dual to one another), for which I refer to VI, theorem-scholie 2, p. 108-113. Here the subsidiary question arises of simplifying the conditions of validity of the adjunction formula, somewhat disparate in loc. cit.

On the other hand, the question arises of developing the theory of the covariant laws $f_!$ and $f_*$ for $\mathrm{Ho}^{\mathrm{lc}}_W(S)$ and for $\mathrm{Ho}^{\mathrm{Prop}}_W(S)$, $\mathrm{Ho}^{\mathrm{Liss}}_W(S)$.

\msp{page 227}

I have good hope that the functor $f_!$ can be described, and its adjunction with $f^*$ established, at as little cost in this general context as in Cat. But I do not feel like dwelling on it at present.

\section*{6 $W$-homotopy cofiber in $\mathrm{Cat}$.}

On the other hand, I feel like examining here under what conditions on the basic localizer $W$, there exists in $(\mathrm{Cat},W)$ a theory dual to that of the $W$-homotopy fiber.

The key question is, of course, that of having a good grip on $W$-cofibrations. Now so far I know \textit{only} the open Dwyer immersions, the closed Dwyer immersions, and what follows from them by composition etc. The work of \textsc{Thomason} may motivate the hope that knowledge of those immersions, which seem here to play a role similar (dual) to that of $W$-perfect morphisms, will suffice to develop a good theory of ``$W$-homotopy cofibers''.

% trois traits verticaux imprimés dans la marge gauche de ce paragraphe
{\reversemarginpar\marginnote{\textbar\textbar\textbar}}I assume throughout this section that $W$ satisfies W(7 bis) (axiom of cocartesian squares), to have good properties of open Dwyer immersions (\footnote[124]{\begin{enumerate}
\item[a)] These are $W$-cofibrations (follows from b) and c)).
\item[b)] If
\end{enumerate}}), and more generally, for $W$-cofibrations. % ?? uncertain: coupure en fin de ligne "co-cartésiens" -- cocartésiens ou co-cartésiens
% foot: Version 30 mars 2003


\endgroup
% END UNIT Der-ChXIIc-30-03-03/Der-ChXIIc-30-03-03_p026-030.tex
% BEGIN UNIT Der-ChXIIc-30-03-03/Der-ChXIIc-30-03-03_p031-035.tex
\begingroup\makeatletter%
\newcommand{\printedpage}[1]{\marginpar{\footnotesize p.~#1}}%
\newcommand{\msp}[1]{\marginpar{\footnotesize [#1]}}%
\makeatother\let\printedpage\derpp\let\msp\dermsp%
\let\mspin\DERmspin%


\printedpage{145}
% head: A. Grothendieck : Dérivateurs, Ch. XII, Édition des Universités de Montpellier II et Paris VII
% foot: Version 30 mars 2003

\msp{page 228}
So the question arises of the dual factorization of an arrow $f : X \longrightarrow Y$ in Cat into
\[
\begin{tikzcd}[column sep=huge]
X \arrow[r, hook, "W\text{-cofibration}", "i_f"'] & \tilde{Y} \arrow[r, "\in W", "p_f"'] & Y.
\end{tikzcd}
\]

The results of \textsc{Thomason} ensure that this exists, with $i$ an open Dwyer immersion, and $p$ a $W$-fibration - at least in the case $W = W_\infty$. I am going to try to see it directly, by dually copying the construction of \textsc{Cartan-Serre}.

So it is first a matter of factoring
\[
X \amalg X \xrightarrow{\ \mathrm{codiag}_X\ } X,
\]
and for this I choose a fixed interval
\[
I \ = \ \Phi_n \ = \ (0 \longrightarrow 1 \longleftarrow 2 \longrightarrow \cdots \longrightarrow 2n-1 \longleftarrow 2n)
\]
with
\[
n \geq 2
\]
so that the inclusion
\[
\varepsilon = \text{discrete category } \{0,1\} \xrightarrow{\partial} I, \qquad \partial(0) = 0, \quad \partial(1) = 2n,
\]
is an open Dwyer immersion. This property is stable under Cat-fibrant base change, in particular on taking a product
\[
\underbrace{X \times \varepsilon}_{\simeq X \coprod X} \xrightarrow{i_X} X \times I,
\]
% continuation of a footnote begun on the preceding page of the edition (no footnote mark on this page)
\begingroup\renewcommand{\thefootnote}{}\footnotetext{%
\[
\begin{tikzcd}[ampersand replacement=\&]
U \arrow[r, hook, "i"] \arrow[d, "p"'] \& X \arrow[d, "q"] \\
U' \arrow[r, hook, "i'"] \& X'
\end{tikzcd}
\]
with $i$ an open Dwyer immersion, then
\[
\begin{aligned}
i \in W \quad &\Longleftrightarrow \quad i' \in W \\
p \in W \quad &\Longleftrightarrow \quad q \in W,
\end{aligned}
\]
so I need W(7 bis) \textit{and} even W(7) for the second implication.
\begin{enumerate}
\item[c)] Open Dwyer immersions \texttt{[are]} stable under cobase change and under composition, contain the isomorphisms.
\end{enumerate}}\endgroup
\printedpage{146}
% head: A. Grothendieck : Dérivateurs, Ch. XII, Édition des Universités de Montpellier II et Paris VII
% foot: Version 30 mars 2003
\msp{page 229}
on the other hand
\[
X \times I \xrightarrow{\mathrm{pr}_1} X
\]
is $W_\omega$-perfect, whence the desired factorization % [sic] factoriastion -> factorisation, misprint
\[
X \amalg X \xrightarrow{i_X} X \times I \ \xrightarrow{\ p_X = \mathrm{pr}_1^X\ } X
\]
- nothing could be simpler! How to use it to construct the factorization of an arbitrary arrow $f : X \longrightarrow Y$, by dualizing the constructions relating to path spaces? We define $Z(f)$ by the cocartesian square in Cat
\[
\begin{tikzcd}[column sep=5em, row sep=4em]
X \arrow[r, dotted, "\mathrm{inj}_1^X"'] & X \times I \arrow[r, "f'"] & Z(f) \arrow[d, dotted, bend left=45, "\,p_f"] \\
& X \arrow[u, hook, "\mathrm{inj}_0^X"'] \arrow[r, "f"] & Y, \arrow[u, hook, "W"']
\end{tikzcd}
\]
we have
\[
Z(f) \xrightarrow{\alpha} Y \qquad \text{induced by } \mathrm{pr}_1 : X \times I \longrightarrow X,
\]
and the desired factorization of $f$ is
\[
X \ \xrightarrow{\ i_f = f' \circ \mathrm{inj}_1^X\ } Z(f) \xrightarrow{p_f} Y.
\]

I say that
\[
\left\{
\begin{aligned}
& i_f \text{ is an open Dwyer immersion, hence a } W\text{-cofibration.} \\
& p_f \in W.
\end{aligned}
\right.
\]

\msp{page 230}
We see from the more complete diagram
\[
\begin{tikzcd}[column sep=6.5em, row sep=4.5em]
& X \times I \arrow[r, "f'"] \arrow[dd, dotted, bend left=45, "\mathrm{pr}_1"{pos=.60}] \arrow[dr, phantom, "\substack{\text{cocart.}\\ \text{and cart.}}"] & Z(f) \arrow[dd, dotted, bend left=60, "p_f"] \\
X \arrow[ur, dotted, "\mathrm{inj}_1^X"] \arrow[r, hook, "i_1^X"] & X \amalg X \arrow[u, hook, "i_X"'] \arrow[r, "f \coprod \mathrm{id}_X"] \arrow[dr, phantom, "\substack{\text{cocart.}\\ \text{and cart.}}"] & Y \amalg X \arrow[u, hook, "j"'] \\
& X \arrow[u, hook, "i_0^X"'] \arrow[r, "f"] & Y, \arrow[u, hook]
\end{tikzcd}
\]
\printedpage{147}
% head: A. Grothendieck : Dérivateurs, Ch. XII, Édition des Universités de Montpellier II et Paris VII
% foot: Version 30 mars 2003
\[
i_f = f' \circ \mathrm{inj}_1^X = f' i_X i_1^X = j \ \ \underbrace{(f \amalg \mathrm{id}_X) i_1^X}_{\substack{\text{canonical inclusion}\\ \text{of } X \text{ into}\\ Y \coprod X}} \ \ ,
\]
it is \texttt{[an]} open Dwyer immersion, as a composite of $X \longrightarrow Y \amalg X$ and $j : Y \amalg X \longrightarrow Z(f)$ which are such ($j$ obtained by cobase change from $i_X$, which is an open Dwyer immersion).
\DERmark{ENV0808}


On the other hand, $X \xrightarrow{\mathrm{inj}_0^X} X \times I$ is \texttt{[an]} open Dwyer immersion and in $W$, hence couniversally in $W$, hence $Y \longrightarrow Z(f)$ is also in $W$, hence also the retraction $p_f$.

On the other hand, the proposition dual to
\[
\underline{\mathrm{Fib}}_W S \cap W_S = W^{\mathrm{u}}_{S\,\mathrm{univ}}
\]
is valid for $W$-cofibrations:

\msp{page 231}
\textbf{Proposition 16} (\footnote[125]{Has nothing to do with Cat - a general thing when there are enough $W$-cofibrations.}). \textit{On the full subcategory $\underline{\mathrm{Cofib}}_W(S)$ of $S\backslash\mathrm{Cat}$, the localizer ${}_SW$ is equal to $({}_SW)^{\mathrm{univ}}$ (couniversal $W$-equivalences), i.e. we have this: Let}
\[
S \longrightarrow X, \qquad S \longrightarrow Y
\]
\textit{be $W$-cofibrations, and let $f : X \longrightarrow Y$ be a morphism of $S\backslash\mathrm{Cat}$ such that $f \in W$. Then for every cobase change $S \longrightarrow S'$, the corresponding arrow}
\[
f' : X' = X \amalg_S S' \longrightarrow Y' = Y \amalg_S S'
\]
\textit{is in $W$.}

\textsc{Proof}. We factor $S \longrightarrow S'$ into
\[
S \xrightarrow{i} S'_0 \xrightarrow{j} S',
\]
with $i$ an open Dwyer immersion (hence a $W$-cofibration), and $j \in W$. So we have the diagram with cocartesian squares
\[
\begin{tikzcd}
S \arrow[r] \arrow[d, hook, "i"'] & X \arrow[r, "f"] \arrow[d, hook] & Y \arrow[d, hook] \\
S'_0 \arrow[r] \arrow[d, "\in W"'] & X'_0 \arrow[r, "f'_0", "W"'] \arrow[d, "W"] & Y'_0 \arrow[d, "W"] \\
S' \arrow[r] & X' \arrow[r, "f'"] & Y'.
\end{tikzcd}
\]
Since $i$ is a $W$-cofibration and $f \in W$, we have $f'_0 \in W$. On the other hand, since $S \longrightarrow X$, hence $S'_0 \longrightarrow X'_0$ is
\printedpage{148}
% head: A. Grothendieck : Dérivateurs, Ch. XII, Édition des Universités de Montpellier II et Paris VII
% foot: Version 30 mars 2003
\msp{page 232}
a $W$-cofibration, $X'_0 \longrightarrow X'$ is in $W$, and for the same reason, $Y'_0 \longrightarrow Y'$ is too. So $f'$ is too. % [sic] Y'_0 \longrightarrow Y -> Y'_0 \longrightarrow Y', the diagram has Y'_0 \to Y'
\DERmark{ENV0809}


\textbf{NB} Here, this is a formal argument, valid (in dual form) in any $(\mathcal{M}, W)$ where we have a factorization statement for an arbitrary arrow $f : X \longrightarrow Y$ into $X \xrightarrow{i_f} \bar{X} \xrightarrow{\alpha} Y$, with $i_f \in W$ and $\alpha$ a $W$-fibration. Then, if $X$, $Y$ are $W$-fibered over $S$, every $S$-arrow from $X$ to $Y$ that is in $W$, is in $W^{\mathrm{u}}_S$. And dually, if every arrow $X \longrightarrow Y$ factors as $X \longrightarrow \tilde{Y} \longrightarrow Y$, with $X \longrightarrow \tilde{Y}$ a $W$-cofibration and $\tilde{Y} \longrightarrow Y$ in $W$.

A (small) worry comes to me here: have I at least proved that if
\[
S \longrightarrow X
\]
is a $W$-cofibration, then for every arrow $S \longrightarrow S'$ in Cat, the cocartesian square in Cat % [sic] S' \longrightarrow S -> S \longrightarrow S', the square below has S \to S'
\DERmark{ENV0810}

\[
\begin{tikzcd}
S \arrow[r] \arrow[d] & X \arrow[d] \\
S' \arrow[r] & X'
\end{tikzcd}
\]
is also $W$-cocartesian (\footnote[126]{in the old sense, special to Cat (def. 3, p. 131). To compare it with the general theory, cf. prop. 20, p. 262.})? Otherwise,
\msp{page 233}
what good would that square be?! To prove it, let us again factor $S \longrightarrow S'$ into
\[
\underset{\substack{\text{open}\\ \text{Dwyer}\\ \text{immersion}}}{S \hookrightarrow} S'_0 \underset{\in W}{\longrightarrow} S',
\]
whence the diagram of 2 cocartesian squares
\[
\begin{tikzcd}
S \arrow[r, "\mathrm{Cof}_W"] \arrow[d, hook, "D_W"'] & X \arrow[d, hook, "D_W"] \\
S'_0 \arrow[r, "\mathrm{Cof}_W"] \arrow[d, "\in W"'] & X'_0 \arrow[d, "\in W"] \\
S' \arrow[r, "\mathrm{Cof}_W"] & X',
\end{tikzcd}
\]
where $X'_0 \longrightarrow X'$ is $\in W$ by the hypothesis on $X$ under $S$. So we must verify:
\printedpage{149}
% head: A. Grothendieck : Dérivateurs, Ch. XII, Édition des Universités de Montpellier II et Paris VII
% foot: Version 30 mars 2003
\begin{enumerate}
\item[a)] A cocartesian square built on an arrow which is \texttt{[a]} Dwyer immersion, is $W$-cocartesian.
\item[b)] A commutative square in which two opposite sides are in $W$, is $W$-cocartesian.
\item[c)] The composite of two $W$-cocartesian squares is $W$-cocartesian.
\end{enumerate}

Let us provisionally admit b) and c), and prove a), for a cocartesian square
\[
\begin{tikzcd}
U \arrow[r, hook, "i"] \arrow[d, "p"'] & X \arrow[d, "q"] \\
U' \arrow[r, hook, "i'"] & X',
\end{tikzcd}
\]
where $i$ is an open Dwyer immersion. We proceed again as before, by factoring $p$ into $U \overset{j}{\hookrightarrow} U'_0 \xrightarrow{p'} U'$, where $j$ is \texttt{[an]} open Dwyer immersion, $p' \in W$. So we have a diagram of two adjoining cocartesian squares
\msp{page 234}
\[
\begin{tikzcd}
U \arrow[r, hook, "i \in D_W"] \arrow[d, hook, "D_W"'] & X \arrow[d, hook] \\
U'_0 \arrow[r, hook, "D_W"] \arrow[d, "p' \in W"'] & X'_0 \arrow[d, "q' \in W"] \\
U' \arrow[r, hook] & X',
\end{tikzcd}
\]
where the first square is made up of open immersions. We then know that it is $W$-cocartesian. In the second, we have $p' \in W$, hence $q' \in W$ since $i$ is a Dwyer immersion, hence a $W$-cofibration. So it is also $W$-cocartesian. We then conclude thanks to c) above. % [sic] immersions ouverts -> immersions ouvertes, agreement

\noindent Let us prove b), for
\begin{tikzcd}[baseline=(\tikzcdmatrixname-1-1.base)]
X_0 \arrow[r, "i_1"] \arrow[d, "i_2 \in W"'] & X_1, \arrow[d, "f_1 \in W"] \\
X_2 \arrow[r, "f_2"] & Y
\end{tikzcd}
by inserting it into
\begin{tikzcd}[baseline=(\tikzcdmatrixname-1-1.base)]
X_0 \arrow[r, "i_1"] \arrow[d, "i_2 \in W"'] & X_1, \arrow[d, "j_1"] \arrow[ddr, "f_1"] & \\
X_2 \arrow[r, "j_2"] \arrow[drr, "f_2"'] & \mathcal{X} \arrow[dr, "\varphi"] & \\
& & Y
\end{tikzcd}
where $\mathcal{X} = \displaystyle\int X_0$
\begin{tikzcd}[baseline=(\tikzcdmatrixname-2-1.base), row sep=small]
& X_1 \\
{} \arrow[ur] \arrow[dr] & \\
& X_2
\end{tikzcd}
.

Since $i_2 \in W$, we know that it follows that $j_1 \in W$, and since $f_1 \in W$, it follows that $\varphi \in W$, q.e.d.


\endgroup
% END UNIT Der-ChXIIc-30-03-03/Der-ChXIIc-30-03-03_p031-035.tex
% BEGIN UNIT Der-ChXIIc-30-03-03/Der-ChXIIc-30-03-03_p036-040.tex
\begingroup\makeatletter%
\newcommand{\printedpage}[1]{\marginpar{\footnotesize p.~#1}}%
\newcommand{\msp}[1]{\marginpar{\footnotesize [#1]}}%
\newcommand{\cX}{\mathcal{X}}%
\newcommand{\Ob}{\operatorname{Ob}}%
\makeatother\let\printedpage\derpp\let\msp\dermsp%
\let\mspin\DERmspin%

\DERmark{ENV0811}


\printedpage{150}
% head: A. Grothendieck : Dérivateurs, Ch. XII, Édition des Universités de Montpellier II et Paris VII
% foot: Version 30 mars 2003

It remains to prove c). I am going to deduce it from a

\textbf{Proposition 17.} \textit{Let there be a commutative square}
\[
Q \quad = \quad
\begin{tikzcd}[ampersand replacement=\&]
X_0 \arrow[r, "i_1"] \arrow[d, "i_2"'] \& X_1 \arrow[d, "f_1"] \\
X_2 \arrow[r, "f_2"] \& Y
\end{tikzcd}
\]
\textit{in} Cat\textit{, and let $\cX_Q$ be the fibered category over}
\[
\tilde{\Phi} \quad (\simeq \boldsymbol{\Delta}^1 \times \boldsymbol{\Delta}^1) \quad = \quad
\begin{tikzcd}[ampersand replacement=\&]
0 \& 1 \arrow[l] \\
2 \arrow[u] \& 3 \arrow[l] \arrow[u]
\end{tikzcd}
\]
\textit{which integrates it. Let $\cX$ be the inverse image of the closed subset}
\[
\begin{tikzcd}[ampersand replacement=\&]
0 \& 1 \arrow[l] \\
2 \arrow[u] \&
\end{tikzcd}
\]
\msp{page 235}
\textit{in $\cX_Q$, which integrates the diagram}
\[
\begin{tikzcd}[ampersand replacement=\&, row sep=small]
 \& X_1 \\
X_0 \arrow[ru] \arrow[rd] \& \\
 \& X_2
\end{tikzcd}
\]
\textit{. For $Q$ to be $W$-cocartesian, it is necessary and sufficient that the inclusion}
\[
i : \cX \longrightarrow \cX_Q
\]
\textit{be in $W$.}

Since $\cX_Q$ is smooth over $\tilde{\Phi}$, and 3 is an initial object of $\tilde{\Phi}$, it follows that the inclusion of the fiber of $\cX_Q$ at the vertex 3 of $\tilde{\Phi}$ is in $W$,
\[
Y \xrightarrow{i_Y} \cX_Q \qquad\qquad \in W.
\]

Consider then the diagram
\[
\begin{tikzcd}[ampersand replacement=\&, column sep=small]
\cX \arrow[rr, "\varphi_Q"] \arrow[rd, "i"', ""{name=direct,pos=.55,above right}] \& \& Y \arrow[ld, "i_Y\in W", ""{name=via,pos=.30,below left}] \\
 \& \cX_Q. \& \arrow[from=via,to=direct,dotted,shorten <=1pt,shorten >=1pt,"\alpha"']
\end{tikzcd}
\]

\printedpage{151}
% head: A. Grothendieck : Dérivateurs, Ch. XII, Édition des Universités de Montpellier II et Paris VII
% foot: Version 30 mars 2003

This diagram is not commutative, but there is a commutativity arrow
\[
\alpha : i_Y\varphi_Q \longrightarrow i.
\]

We know that this implies that
\[
i \in W \qquad\qquad \Longleftrightarrow \qquad\qquad (i_Y\varphi_Q) \in W
\]
(cf. lemma p. 2), and since $i_Y \in W$,
\[
(i_Y\varphi_Q) \in W \qquad\qquad \Longleftrightarrow \qquad\qquad \varphi_Q \in W,
\]
whence
\[
i \in W \qquad\qquad \Longleftrightarrow \qquad\qquad \varphi_Q \in W,
\]

\hfill q.e.d.

We can now prove

\msp{page 236}

\textbf{Proposition 18} (\footnote[127]{should be true in every category of models having enough cofibrations.}). (Assuming W(7 bis).) \textit{Consider a commutative diagram}
\[
\begin{tikzcd}[ampersand replacement=\&]
X_0 \arrow[r] \arrow[d] \arrow[rd, phantom, "1"] \& X_1 \arrow[r] \arrow[d] \arrow[rd, phantom, "2"] \& X_2 \arrow[d] \\
X'_0 \arrow[r] \& X'_1 \arrow[r] \& X'_2
\end{tikzcd}
\tag*{$(\ast)$}
\]
\textit{whence a composite square}
\[
\begin{tikzcd}[ampersand replacement=\&]
X_0 \arrow[r] \arrow[d] \arrow[rd, phantom, "3"] \& X_2 \arrow[d] \\
X'_0 \arrow[r] \& X'_2.
\end{tikzcd}
\]
\textit{Suppose that square} 1 \textit{is $W$-cocartesian. Then square} 2 \textit{is $W$-cocartesian if and only if square} 3 \textit{is.}

\textsc{Proof.}

\texttt{[Now, }$X_i$\texttt{ and }$X'_i$\texttt{ exchange their roles.]}

\printedpage{152}
% head: A. Grothendieck : Dérivateurs, Ch. XII, Édition des Universités de Montpellier II et Paris VII
% foot: Version 30 mars 2003

Let
\[
I \quad = \quad
\left(
\begin{tikzcd}[ampersand replacement=\&]
0 \arrow[d] \& 1 \arrow[l] \arrow[d] \& 2 \arrow[l] \arrow[d] \\
0' \& 1' \arrow[l] \& 2' \arrow[l]
\end{tikzcd}
\right)
\]
and let $\cX$ be the integral of the diagram $(\ast)$, a fibered category over $I$. We denote by $\cX_{010'}$ the full subcategory of $\cX$ which is the inverse image of the full subcategory % [sic] fibrée sur $\cX$ -> fibrée sur $I$, a category fibred over itself; the base is the category $I$ just defined
\DERmark{ENV0812}

\[
\underbrace{I_{010'}}_{=\{0,1,0'\}} \ \subseteq I,
\]
and likewise for every other subset $A$ of $\Ob I = \{0,1,2,0',1',2'\}$, giving a full subcategory $\cX_A$ of $\cX$. We put
\[
U \quad = \quad
\begin{tikzcd}[ampersand replacement=\&]
0 \arrow[d] \& 1 \arrow[l] \& 2 \arrow[l] \\
0' \& \&
\end{tikzcd}
, \qquad
V \quad = \quad
\begin{tikzcd}[ampersand replacement=\&]
0 \arrow[d] \& 1 \arrow[l] \arrow[d] \& 2 \arrow[l] \\
0' \& 1' \arrow[l] \&
\end{tikzcd}
,
\]
whence
\[
\cX_U = \cX_{0120'}, \qquad\qquad \cX_V = \cX_{0120'1'}.
\]
The arrows between the $\cX_A$ are the inclusions

\msp{page 237}

deduced from inclusions $A \hookrightarrow A'$ in $\mathfrak{P}(\Ob I)$. Thus, the $W$-cocartesian character of squares 1, 2, 3 is expressed respectively by the condition that the inclusions
\[
\begin{aligned}
\cX_{010'} \quad &\xrightarrow{\varphi_1} \quad \cX_{010'1'} \\
\cX_{121'} \quad &\xrightarrow{\varphi_2} \quad \cX_{121'2'} \\
\cX_{020'} \quad &\xrightarrow{\varphi_3} \quad \cX_{020'2'} % [sic] $\varphi_1$ -> $\varphi_3$, this inclusion is called $\varphi_3$ in the diagram that follows
\end{aligned}
\]
\DERmark{ENV0813}

be in $W$. Consider then the diagram of inclusion arrows
\printedpage{153}
% head: A. Grothendieck : Dérivateurs, Ch. XII, Édition des Universités de Montpellier II et Paris VII
% foot: Version 30 mars 2003
\[
\begin{tikzcd}[ampersand replacement=\&]
\cX_{020'} \arrow[r, "\varphi_3"] \arrow[d, "\alpha"'] \& \cX_{020'2'} \arrow[d, "\beta"] \\
\cX_U \arrow[r, "\psi_3"] \arrow[d, "\psi_1"'] \& \cX \arrow[d, "\mathrm{id}"] \\
\cX_V \arrow[r, "\psi_2"] \& \cX \\
\cX_{121'} \arrow[r, "\varphi_2"] \arrow[u, "\alpha'"] \& \cX_{121'2'}. \arrow[u, "\beta'"']
\end{tikzcd}
\]

I say that in this diagram we have
\[
\alpha, \beta, \alpha', \beta' \in W
\tag*{$(\ast)$}
\]
and
\[
\varphi_1 \in W \ \ \text{(i.e. square 1 $W$-cocartesian)} \qquad\qquad \Longleftrightarrow \qquad\qquad \psi_1 \in W.
\tag*{$(\ast\ast)$}
\]

The relations $(\ast)$ already imply that % [sic] déja -> déjà, missing grave accent
\[
\begin{aligned}
\varphi_2 \in W \quad &\Longleftrightarrow \quad \psi_2 \in W \\
\varphi_3 \in W \quad &\Longleftrightarrow \quad \psi_3 \in W,
\end{aligned}
\]
and under the $W$-cocartesian hypothesis $(\ast\ast)$, implying $\psi_1 \in W$,

\msp{page 238}

we find
\[
\psi_2 \in W \qquad\qquad \Longleftrightarrow \qquad\qquad \psi_3 \in W,
\]
whence
\[
\varphi_2 \in W \qquad\qquad \Longleftrightarrow \qquad\qquad \varphi_3 \in W,
\]
which is what was to be proved.

It remains to establish $(\ast)$ and $(\ast\ast)$. The relations $(\ast)$ come from the observation that if $i : A \longrightarrow A'$ is an inclusion of full subcategories of $I$, such that $i$ is $W$-\textit{aspheric} (i.e. $A/j$ $W$-aspheric for $j \in \Ob A'\backslash\Ob A$), then $\cX_A \longrightarrow \cX_{A'}$ is so too (smooth base change), hence is \texttt{[in]} $W$. We must verify $W$-asphericity for the four inclusions
\[
\begin{array}{rcll}
\{0,2,0'\} & \longrightarrow & U = \{0,1,2,0'\} & \qquad\qquad \boldsymbol{\Delta}^0 \\[1ex]
\{0,2,0',2'\} & \longrightarrow & I & \qquad\qquad \boldsymbol{\Delta}^0/\boldsymbol{\Delta}^1 \\[1ex]
\{1,2,1'\} & \longrightarrow & V = \{0,1,2,0',1'\} & \qquad\qquad \boldsymbol{\Delta}^1,\ \Delta^2 \\[1ex]
\{1,2,1',2'\} & \longrightarrow & I & \qquad\qquad \boldsymbol{\Delta}^1,\ \boldsymbol{\Delta}^1 \times \boldsymbol{\Delta}^1.
\end{array}
\]
\DERmark{ENV0814}


\printedpage{154}
% head: A. Grothendieck : Dérivateurs, Ch. XII, Édition des Universités de Montpellier II et Paris VII
% foot: Version 30 mars 2003

(NB I have marked in the margin the categories we obtain as $A/j$, up to isomorphism, in the four cases considered - they all have an initial object, hence are $W$-aspheric.)

It remains to prove $(\ast\ast)$. Let
\[
U' = \{0,1,0',1'\}, \qquad\qquad \text{hence } V = U \cup U',\ \ U \cap U' = \{0,1,0'\},
\]
where $U$, $U'$ are \textit{closed subsets} of $V$ which cover $V$. Thus we will likewise have
\[
\cX_V = \cX_U \cup \cX_{U'}, \qquad\qquad \cX_U \cap \cX_{U'} = \cX_{010'},
\]

\msp{page 239}

\[
\begin{tikzcd}[ampersand replacement=\&]
\cX_{U\cap U'} \arrow[r] \arrow[d, "\varphi_1"] \& \cX_U \arrow[d, "\psi_1"] \\
\cX_{U'} \arrow[r] \& \cX_V.
\end{tikzcd}
\]

Thus by W(7 bis) (dual case of closed immersions), we have
\[
\varphi_1 \in W \qquad\qquad \Longleftrightarrow \qquad\qquad \psi_1 \in W,
\]

\hfill q.e.d.

Let us note for the record the corollary which had motivated us.

\textbf{Corollary.} \textit{Assume} W(7 bis) (as throughout this section). \textit{Let there be a cocartesian square in} Cat
\[
\begin{tikzcd}[ampersand replacement=\&]
X_0 \arrow[r, "i_1"] \arrow[d, "i_2"'] \& X_1 \arrow[d] \\
X_2 \arrow[r] \& X_3,
\end{tikzcd}
\]
\textit{with $i_1$ (or $i_2$) a $W$-cofibration (e.g. a Dwyer immersion). Then this square is $W$-cocartesian.}

This digression (p. 232 ff.) is not necessary, for that matter, to formulate right now the conclusion of our cogitations on $W$-homotopic cofibers: \textit{There is such a theory in} Cat, provided only that $W$ satisfies W(7 bis). Indeed, everything follows from the fact that every arrow $f : X \longrightarrow Y$ in Cat factors as
\[
X \xrightarrow{i_f} \tilde{Y} \xrightarrow{p_f} Y,
\]
with $i_f$ a $W$-cofibration (and even an open Dwyer imm-


\endgroup
% END UNIT Der-ChXIIc-30-03-03/Der-ChXIIc-30-03-03_p036-040.tex
% BEGIN UNIT Der-ChXIIc-30-03-03/Der-ChXIIc-30-03-03_p041-045.tex
\begingroup\makeatletter%
\derleqno%
\newcommand{\printedpage}[1]{\marginpar{\footnotesize p.~#1}}%
\newcommand{\msp}[1]{\marginpar{\footnotesize [#1]}}%
\newcommand{\cercle}[1]{\tikz[baseline=(c.base)]\node[draw,circle,inner sep=1pt](c){#1};}%
\newcommand{\Fibu}{\underline{\mathrm{Fib}}}%
\newcommand{\Homu}{\underline{\mathrm{Hom}}}%
\newcommand{\Ho}{\mathrm{Ho}}%
\newcommand{\ho}{\mathrm{ho}}%
\newcommand{\lc}{\mathrm{lc}}%
\makeatother\let\printedpage\derpp\let\msp\dermsp%
\let\mspin\DERmspin%


\printedpage{155}
% head: A. Grothendieck : Dérivateurs, Ch. XII, Édition des Universités de Montpellier II et Paris VII
% foot: Version 30 mars 2003

\msp{page 240}
ersion), and $p_f \in W$. % beginning of a word continued from the preceding page (outside this range), as printed

\section*{7\quad Return to the axiomatics of the $W$-homotopic fiber.}

I come back to the reflection in n$^\circ$ 4. In the end, for
\[
(\mathcal{M}, W \subseteq \mathrm{Fl}(\mathcal{M}))
\]
given with $W$ weakly saturated, it appears that an excellent theory of $W$-homotopic fibers, satisfying the conditions \cercle{a}\cercle{b}\cercle{c} of n$^\circ$ 4 (pages 215-218), exists as soon as the following innocent-looking axiom is satisfied (in addition to the existence of the final object and fiber products (\footnote[128]{and even that! it seems to me that it suffices to assume that the $X \longrightarrow Y$ which are $W$-fibrant are ``squarable''.})).

\textbf{Left factorization:} (\footnote[129]{Under these conditions, $\mathcal{M}$ equipped with $(W, F_W)$ is a category with fibrations of K.S. \textsc{Brown-Anderson}, and ``everything is O.K.'', but one would have to check a) b) c) above (p. 215, 218), I am not sure I have done it.},\footnote[130]{\textbf{NB} The left factorization axiom is sufficient for the existence of a good theory of homotopy fibers (it seems) only in the case where every $X$ in $\mathcal{M}$ is $W$-fibrant, i.e. is in $\Fibu_W e$.}). Every arrow $f : X \longrightarrow Y$ of $\mathcal{M}$ factors as
\DERmark{ENV0815}

\[
X \xrightarrow{i_f} \bar{X} \xrightarrow{p_f} Y,
\]
with
\[
i_f \in W, \qquad p_f \in \mathrm{Fib}_W.
\]

The dual axiom gives rise to a good theory of $W$-homotopic cofibers.

\textbf{Right factorization:} Every arrow $f : X \longrightarrow Y$ in $\mathcal{M}$ factors as
\[
X \xrightarrow{j_f} \tilde{Y} \xrightarrow{q_f} Y,
\]
with
\[
q_f \in W, \qquad j_f \in \mathrm{Cofib}_W.
\]

\msp{page 241}
I would now like to develop somewhat the formalism of $W$-homotopic fibers, under the assumption (say) that $W$ satisfies the left factorization axiom.

Let us consider a cartesian diagram
\[
\begin{tikzcd}[ampersand replacement=\&]
X_1 \arrow[d, "p_1"'] \& Y \arrow[l, "q_2"'] \arrow[d, "q_1"] \\
S \& X_2 \arrow[l, "p_2"]
\end{tikzcd}
\]
\printedpage{156}
% head: A. Grothendieck : Dérivateurs, Ch. XII, Édition des Universités de Montpellier II et Paris VII
% foot: Version 30 mars 2003
in $\mathcal{M}$. If $p_1$ is in $\mathrm{Fib}_W$, then $Y$ over $X_2$ can be interpreted as the $W$-homotopic fiber of $X_1/S$ at $X_2/S$. More precisely, this is so for the class $\ho^{\lc}_{W,X_2}(Y)$,
\[
\ho^{\lc}_{W,X_2}(Y) \simeq p_2^*(\ho^{\lc}_{W,S}(X_1/S)). \tag*{$(*)$}
\]
But as an object over $X_1$, $Y$ can also be interpreted as a $W$-homotopic fiber of $X_2/S$ at $X_1/S$:
\[
\ho^{\lc}_{W,X_1}(Y/X_1) \simeq p_1^*(\ho^{\lc}_{W,S}(X_2/S)). \tag*{$(**)$}
\]
To tell the truth, one would have to prove it. For this, let us consider a factorization
\[
X_2 \xrightarrow{i_2} \bar{X}_2 \xrightarrow{\bar{p}_2} S \qquad i_2 \in W\ ,\quad \bar p_2 \in \mathrm{Fib}_W, % [sic] p_2 \in \mathrm{Fib}_W -> \bar{p}_2 \in \mathrm{Fib}_W, the factorisation just written is through \bar{p}_2
\]
\DERmark{ENV0816}

so that we find a diagram of cartesian squares
\[
\begin{tikzcd}[ampersand replacement=\&]
X_1 \arrow[d] \& \bar{Y} \arrow[l, "\bar{q}_2"'] \arrow[d] \& Y \arrow[l, "j_2"', "W"] \arrow[d] \\
S \& \bar{X}_2 \arrow[l, "\bar{p}_2"'] \& X_2. \arrow[l, "i_2"', "W"] % [sic] \bar{q}_2 (bottom left arrow) -> \bar{p}_2, the arrow \bar{X}_2 \to S is \bar{p}_2 of the factorisation above
\end{tikzcd}
\]
\DERmark{ENV0817}


\msp{page 242}
By definition of $p_1^*$, we have
\[
p_1^*(\ho^{\lc}_{W,S}(X_2/S)) \simeq \ho^{\lc}_{W,X_1}(\bar{Y}/X_1),
\]
on the other hand, since $X_1 \longrightarrow S$ is in $\mathrm{Fib}_W$, and $i_2 \in W$, we have $j_2 \in W$, hence
\[
\ho^{\lc}_{W,X_1}(Y/X_1) \xrightarrow{\sim} \ho^{\lc}_{W,X_1}(\bar{Y}/X_1),
\]
whence formula $(**)$.

Thus, the object $Y$ has two different homotopical interpretations, depending on whether one regards it as being over $X_1$, or over $X_2$. One would then like an interpretation which encompasses both at once. We are going to consider the diagram
\[
\begin{tikzcd}[ampersand replacement=\&]
X_1 \arrow[d] \& \\
S \& X_2 \arrow[l]
\end{tikzcd}
\]
of type $\Psi^\circ$, or
\[
\Psi \quad = \quad
\begin{tikzcd}[ampersand replacement=\&]
1 \& \\
0 \arrow[u] \arrow[r] \& 2
\end{tikzcd}
\qquad \text{or} \qquad
\begin{tikzcd}[ampersand replacement=\&, row sep=small, column sep=small]
 \& 1 \\
0 \arrow[ur] \arrow[dr] \& \\
 \& 2
\end{tikzcd}
\quad ,
\]
% the last "2" of the first diagram is printed overlapping the page number 156
\printedpage{157}
% head: A. Grothendieck : Dérivateurs, Ch. XII, Édition des Universités de Montpellier II et Paris VII
% foot: Version 30 mars 2003
from which $Y$ is obtained by fiber product, by completing this diagram to a diagram of type $\tilde{\Psi}^\circ$, where
\[
\tilde{\Psi} =
\begin{tikzcd}[ampersand replacement=\&, row sep=small, column sep=small]
 \& 1 \arrow[dr] \& \\
0 \arrow[ur] \arrow[dr] \& \& 3 \\
 \& 2 \arrow[ur] \&
\end{tikzcd}
\qquad (\simeq \boldsymbol{\Delta}^1 \times \boldsymbol{\Delta}^1).
\]

\msp{page 243}
For every category $I$, let
\[
\Ho_{\mathcal{M},W}(I) \text{ \ or \ } \Ho_W(I) \overset{\mathrm{def}}{=} \mathcal{M}(I)W(I)^{-1},
\]
\DERmark{ENV0818}

where
\[
\left\{
\begin{array}{rcl}
\mathcal{M}(I) & = & \Homu(I^\circ, \mathcal{M}) \\[1ex]
W(I) & \subseteq & \mathrm{Fl}(\mathcal{M}(I)) \text{ consists of the } u : F \longrightarrow G \text{ such that} \\[1ex]
 & & u_i : F(i) \longrightarrow G(i) \text{ in } W \text{ for every } i \text{ in } \mathrm{Ob}\, I.
\end{array}
\right.
\]

Thus, the diagram $\vcenter{\hbox{\begin{tikzcd}[ampersand replacement=\&, row sep=small, column sep=small]
 \& X_1 \arrow[dl] \\
S \& \\
 \& X_2 \arrow[ul]
\end{tikzcd}}}$ defines an object $\xi$ of $\mathcal{M}(\Psi)$, hence of $\Ho_W(\Psi)$, and the diagram $\vcenter{\hbox{\begin{tikzcd}[ampersand replacement=\&, row sep=small, column sep=small]
 \& X_1 \arrow[dl] \& \\
S \& \& Y \arrow[ul] \arrow[dl] \\
 \& X_2 \arrow[ul] \&
\end{tikzcd}}}$ can likewise be interpreted as an object $\tilde{\xi}$ of $\Ho_W(\tilde{\Psi})$, whose restriction to $\Psi$ is $\xi$.

This being laid down, I would like to define a functor
\[
i_* : \Ho_W(\Psi) \longrightarrow \Ho_W(\tilde{\Psi})
\]
(where
\[
i : \Psi \longrightarrow \tilde{\Psi}
\]
is the inclusion, which is an open immersion), which must correspond, in the case of an object $\vcenter{\hbox{\begin{tikzcd}[ampersand replacement=\&, row sep=small, column sep=small]
 \& X_1 \arrow[dl] \\
S \& \\
 \& X_2 \arrow[ul]
\end{tikzcd}}}$ of $\Ho_W(\Psi)$ with $p_1 \in \mathrm{Fib}_W$ \texttt{[}$:$ $X_1 \longrightarrow S$\texttt{]}, to the fiber product operation (\footnote[131]{\textbf{NB} This construction, which is done here rather laboriously, becomes immediate, and its main properties too, if one admits that the category with fibrations $(\mathcal{M}, W, \mathrm{Fib}_W)$ defines a \textit{derivator} on the domain of finite ordered sets, which can indeed be proved quite simply (with no more trouble than I am taking here).}).

\printedpage{158}
% head: A. Grothendieck : Dérivateurs, Ch. XII, Édition des Universités de Montpellier II et Paris VII
% foot: Version 30 mars 2003

\msp{page 244}
The construction of this functor $i_*$ should be essentially sorital, with what we have at our disposal. We shall first define, for $S$ \textit{fixed,} a canonical functor
\[
\Ho^{\lc}_W(S) \times \Ho^{\lc}_W(S) \longrightarrow \Ho_{W,S}(\Phi), \tag*{$(*)$}
\]
\DERmark{ENV0819}

where the second member is defined, by analogy with the $\Ho_W(I)$, by
\[
\Ho_{W,S}(I) \overset{\mathrm{def}}{=} ((\mathcal{M}/S)(I))(W_S(I))^{-1}
\]
\DERmark{ENV0820}

- i.e. it is the same definition, except that $(\mathcal{M}, W)$ is replaced by $(\mathcal{M}/S, W_S)$. As for $\Phi = \vcenter{\hbox{\begin{tikzcd}[ampersand replacement=\&, row sep=small, column sep=small]
1 \arrow[dr] \& \\
 \& 3 \\
2 \arrow[ur] \&
\end{tikzcd}}}$ , we regard it as a closed subcategory of $\tilde{\Psi}$, which is obtained from it by adjoining an initial object 0 (just as $\tilde{\Psi}$ is obtained from $\Psi$ by adjoining a final object 3). On the other hand, we have a tautological functor
\[
\Ho_{W,S}(\Phi) \longrightarrow \Ho_W(\tilde{\Psi}), \tag*{$(*')$}
\]
which on objects consists in completing a diagram $\vcenter{\hbox{\begin{tikzcd}[ampersand replacement=\&, row sep=small, column sep=small] % [sic] les objet -> les objets, plural agreement
X_1 \& \\
 \& Y \arrow[ul] \arrow[dl] \\
X_2 \&
\end{tikzcd}}}$ of $\mathcal{M}/S$ to the square diagram $\vcenter{\hbox{\begin{tikzcd}[ampersand replacement=\&, row sep=small, column sep=small]
 \& X_1 \arrow[dl] \& \\
S \& \& Y \arrow[ul] \arrow[dl] \\
 \& X_2 \arrow[ul] \&
\end{tikzcd}}}$ of $\mathcal{M}$.

\msp{page 245}
The composite of $(*)$ and $(*')$ will therefore give a functor
\[
\Ho^{\lc}_W(S) \times \Ho^{\lc}_W(S) \longrightarrow \Ho_W(\tilde{\Psi}), \tag*{$(**)$}
\]
already very close to the functor $i_*$ to be constructed. To construct $(*)$, we interpret (up to equivalence of categories) the first member as
\[
(\Fibu_W S \times \Fibu_W S)(W_S \times W_S)^{-1},
\]
so we have to construct a functor
\[
\Fibu_W S \times \Fibu_W S \longrightarrow \Ho_{W,S}(\Phi),
\]
which transforms arrows in $W_S \times W_S$ into isomorphisms. We shall take the functor
\[
(X_1, X_2) \ \longmapsto \
\left(
\begin{tikzcd}[ampersand replacement=\&]
X_1 \& X_1 \times_S X_2 \arrow[l] \arrow[d] \\
 \& X_2
\end{tikzcd}
\right),
\]
\printedpage{159}
% head: A. Grothendieck : Dérivateurs, Ch. XII, Édition des Universités de Montpellier II et Paris VII
% foot: Version 30 mars 2003
and we only have to check that if
\[
X_1 \xrightarrow{\alpha_1} X_1', \qquad X_2 \xrightarrow{\alpha_2} X_2'
\]
are two arrows in $\Fibu_W S$, whence
\[
X_1 \times_S X_2 \xrightarrow{\alpha = \alpha_1 \times_S \alpha_2} X_1' \times_S X_2',
\]
then
\[
\alpha_1 \in W,\ \alpha_2 \in W \qquad \Longrightarrow \qquad \alpha \in W,
\]
which is immediate from the factorization
\[
X_1 \times_S X_2 \quad \underbrace{\longrightarrow}_{\substack{\in W \\ \text{since } X_2 \in \Fibu_W S}} \quad X_1' \times_S X_2 \quad \underbrace{\longrightarrow}_{\substack{\in W \\ \text{since } X_1' \in \Fibu_W S}} \quad X_1' \times_S X_2'.
\]

\msp{page 246}
We have proved the

\textbf{Proposition 18.} \textit{Let $S$ be in $\mathcal{M}$. There exists, up to unique isomorphism, a functor and only one}
\[
\Ho^{\lc}_W(S) \times \Ho^{\lc}_W(S) \xrightarrow{\pi} \Ho_{W,S}(\Phi)
\]
\textit{giving rise to a square commutative up to isomorphism}
\[
\begin{tikzcd}[ampersand replacement=\&, column sep=large]
 \& {\scriptstyle (X_1,X_2)\ \longmapsto\ \begin{array}{@{}c@{}c@{}c@{}} \scriptstyle X_1 & & \\ & \nwarrow & \\ & & \scriptstyle X_1 \times_S X_2 \\ & \swarrow & \\ \scriptstyle X_2 & & \end{array}} \& \\[-3ex]
\Fibu_W S \times \Fibu_W S \arrow[rr] \arrow[d] \& \& (\mathcal{M}/S)(\Phi) \arrow[d] \\
\Ho^{\lc}_W(S) \times \Ho^{\lc}_W(S) \arrow[rr] \& \& \Ho_{W,S}(\Phi) \\[-4ex]
(= (\mathcal{M}/S)(W_S)^{-1} \times (\mathcal{M}/S)(W_S)^{-1}) \& \& (= \mathcal{M}/S(\Phi)(W_S(\Phi))^{-1}).
\end{tikzcd}
\]
\textit{For $X_1$, $X_2$ in $\mathcal{M}/S$, the diagram}
\[
\pi(X_1, X_2) \overset{\mathrm{def}}{=} \pi(\ho^{\lc}_{W,S}(X_1), \ho^{\lc}_{W,S}(X_2))
\]
\DERmark{ENV0821}

\textit{is ``computed'' as follows, up to canonical isomorphism: one chooses}
\[
X_1 \xrightarrow{\alpha_1} \bar{X}_1, \qquad X_2 \xrightarrow{\alpha_2} \bar{X}_2 \qquad \textit{arrows in } \mathcal{M}/S
\]
\textit{with}
\[
\alpha_1, \alpha_2 \in W_S, \qquad \bar{X}_1, \bar{X}_2 \in \mathrm{Ob}\, \Fibu_W S,
\]
\textit{which defines}
\[
\pi(\alpha_1, \alpha_2) : \pi(X_1, X_2) \xrightarrow{\sim} \pi(\bar{X}_1, \bar{X}_2) = \ho_{S,W,\Phi}
\left(
\begin{tikzcd}[ampersand replacement=\&]
\bar{X}_1 \& \bar{X}_1 \times_S \bar{X}_2 \arrow[l] \arrow[d] \\
 \& \bar{X}_2
\end{tikzcd}
\right).
\]


\endgroup
% END UNIT Der-ChXIIc-30-03-03/Der-ChXIIc-30-03-03_p041-045.tex
% BEGIN UNIT Der-ChXIIc-30-03-03/Der-ChXIIc-30-03-03_p046-050.tex
\begingroup\makeatletter%
\newcommand{\printedpage}[1]{\marginpar{\footnotesize p.~#1}}%
\newcommand{\msp}[1]{\marginpar{\footnotesize [#1]}}%
\newcommand{\Ho}{\mathrm{Ho}}%
\newcommand{\lc}{\mathrm{lc}}%
\newcommand{\uFib}{\underline{\mathrm{Fib}}}%
\makeatother\let\printedpage\derpp\let\msp\dermsp%
\let\mspin\DERmspin%


\setcounter{footnote}{131}

\printedpage{160}
% head: A. Grothendieck : Dérivateurs, Ch. XII, Édition des Universités de Montpellier II et Paris VII
% foot: Version 30 mars 2003

\msp{page 247}

\textbf{Corollary.} {\itshape Description of the functor
\[
(\xi_1,\xi_2)\longmapsto\xi_1\times_S\xi_2 : \Ho^{\lc}_W(S)\times\Ho^{\lc}_W(S)\longrightarrow\Ho^{\lc}_W(S).
\]
This functor really is a cartesian product in} $\Ho^{\lc}_W(S)$ (\footnote{but this is a special case of the cases of $\Ho_W$ (if $\mathcal{M}$ has binary products), replacing $\mathcal{M}$ by $\mathcal{M}/S$.}), {\itshape where finite products exist} (\footnote{Say also that the functors $f^* : \Ho^{\lc}_W(S)\longrightarrow\Ho^{\lc}_W(S')$ (for $f : S'\longrightarrow S$) commute with finite products.}).

To prove this last point, setting
\[
\mathcal{H}=\Ho^{\lc}_W(S),
\]
we must show that the functors
\[
\begin{tikzcd}[ampersand replacement=\&]
\mathcal{H} \arrow[r, shift left, "\delta"] \& \mathcal{H}\times\mathcal{H} \arrow[l, shift left, "\pi"]
\end{tikzcd}
\]
are adjoint, with
\[
\begin{array}{rclcrcl}
\delta & = & \mathrm{diag}_{\mathcal{H}} & : & \xi & \longmapsto & (\xi,\xi)\\
\pi & = & \text{``product''} & : & (\xi_1,\xi_2) & \longmapsto & \xi_1\times_S\xi_2.
\end{array}
\]
We are going to define the adjunction homomorphisms
\[
\begin{array}{rcllrcl}
\mathrm{id}_{\mathcal{H}} & \xrightarrow{\alpha} & \pi\delta & \text{i.e.} & \xi & \xrightarrow{\alpha_\xi} & \xi\times_S\xi\\
\mathrm{id}_{\mathcal{H}\times\mathcal{H}} & \xleftarrow{\beta} & \delta\pi & \text{i.e.} & (\xi_1,\xi_2) & \xleftarrow{\beta_{\xi_1,\xi_2}} & (\xi_1\times_S\xi_2,\xi_1\times_S\xi_2).
\end{array}
\]
For this, we describe $\mathcal{H}$ as
\[
\mathcal{H}=(\uFib_W S)(W_S)^{-1}
\]
and we note that $\delta$, $\pi$ come from adjoint functors at the level of $\mathcal{H}_0=\uFib_W S$
\[
\begin{tikzcd}[ampersand replacement=\&]
\mathcal{H}_0 \arrow[r, shift left, "\delta_0"] \& \mathcal{H}_0\times\mathcal{H}_0 \arrow[l, shift left, "\pi_0"]
\end{tikzcd},
\]
which give rise to $\alpha_0$, $\beta_0$
\[
\begin{array}{rcl}
\mathrm{id}_{\mathcal{H}_0} & \xrightarrow{\alpha_0} & \pi_0\delta_0\\
\mathrm{id}_{\mathcal{H}_0\times\mathcal{H}_0} & \xleftarrow{\beta_0} & \delta_0\pi_0.
\end{array}
\]

\printedpage{161}
% head: A. Grothendieck : Dérivateurs, Ch. XII, Édition des Universités de Montpellier II et Paris VII
% foot: Version 30 mars 2003

\msp{page 248}

We define $\alpha$, $\beta$ by passage to the categories of fractions for $\alpha_0$, $\beta_0$. The 2 compatibilities, being satisfied for $\alpha_0$, $\beta_0$, are also satisfied for $\alpha$, $\beta$. Q.e.d.

I single out here the

\textbf{Lemma.} {\itshape Let \begin{tikzcd}[ampersand replacement=\&]
\mathcal{X} \arrow[r, shift left, "f"] \& \mathcal{Y} \arrow[l, shift left, "g"]
\end{tikzcd} be a pair of adjoint functors, and let $\Sigma$, $\Sigma'$ be localizers in $\mathcal{X}$, $\mathcal{Y}$ such that $f(\Sigma)\subseteq\Sigma'$, $g(\Sigma')\subseteq\Sigma$, so that we have a pair of functors
\[
\begin{tikzcd}[ampersand replacement=\&]
\mathcal{X}\Sigma^{-1} \arrow[r, shift left, "\bar{f}"] \& \mathcal{Y}\Sigma'^{-1} \arrow[l, shift left, "\bar{g}"]
\end{tikzcd}.
\]
Then $\bar{f}$, $\bar{g}$ are adjoint, the adjunction arrows being deduced from those for $f$, $g$.}

\textbf{Corollary.} {\itshape Let $(\mathcal{M},W)$ be a category of models, suppose that in $\mathcal{M}$ binary products (resp. finite products) exist, and that for $u_1 : X_1\longrightarrow X'_1$ and $u_2 : X_2\longrightarrow X'_2$, if $u_1,u_2\in W$, then $u_1\times u_2 : X_1\times X_2\longrightarrow X'_1\times X'_2$ is in $W$. Then in $\mathcal{M}W^{-1}$ binary products (resp. finite products) exist, and the localization functor $\mathcal{M}\longrightarrow\mathcal{M}W^{-1}$ commutes with them.}

We apply the lemma to
\[
\begin{tikzcd}[ampersand replacement=\&]
\mathcal{M} \arrow[r, shift left, "\delta"] \& \mathcal{M}\times\mathcal{M} \arrow[l, shift left, "\pi"]
\end{tikzcd}
\qquad \text{and } W,\ W\times W,
\]
and for the final object,
\[
\begin{tikzcd}[ampersand replacement=\&]
\mathcal{M} \arrow[r, shift left, "e"] \& e \arrow[l, shift left, "\bullet"]
\end{tikzcd}
\qquad\qquad \text{and } W,\ W_0=\{\mathrm{id}_e\}.
\]

\msp{page 249}

I am now going to construct the functor
\[
i_* : \Ho_W(\Psi)\longrightarrow\Ho_W(\tilde{\Psi}).
\]
I must first prove the

\textbf{Lemma 1.} {\itshape Let $\mathcal{F}$ be the full subcategory of $\mathcal{M}(\Psi)$, category of diagrams \begin{tikzcd}[ampersand replacement=\&, row sep=small, column sep=small]
 \& X_1 \arrow[dl, "p_1"'] \\
S \& \\
 \& X_2 \arrow[ul, "p_2"]
\end{tikzcd} in $S$ \texttt{[rather in }$\mathcal{M}$\texttt{, over }$S$\texttt{]}, consisting of those for which $p_1$, $p_2$ are in Fib$_W$. Let $W_{\mathcal{F}}$ be the localizer induced on $\mathcal{F}$ by $W(\Psi)$ (termwise $W$-equivalences). Then
\[
\mathcal{F}W_{\mathcal{F}}^{-1}\xrightarrow{\sim}\mathcal{M}(\Psi)W(\Psi)^{-1}\ (\overset{\mathrm{def}}{=}\Ho_W(\Psi)).
\]}
\DERmark{ENV0822}


\printedpage{162}
% head: A. Grothendieck : Dérivateurs, Ch. XII, Édition des Universités de Montpellier II et Paris VII
% foot: Version 30 mars 2003

We are going to prove this lemma below. Assuming it, to construct $i_*$, it suffices to construct
\[
\mathcal{F}\longrightarrow\Ho_W(\tilde{\Psi})
\]
which maps $W_{\mathcal{F}}$ to isomorphisms. I shall take, of course,
\[
\begin{gathered}
\begin{tikzcd}[ampersand replacement=\&, row sep=small, column sep=small]
 \& X_1 \arrow[dl] \& \& \& \& X_1 \arrow[dl] \& \\
S \& \& \longmapsto \& \& S \& \& X_1\times_S X_2 \arrow[ul] \arrow[dl] \\
 \& X_2 \arrow[ul] \& \& \& \& X_2 \arrow[ul] \&
\end{tikzcd}\\
\mathcal{F}\quad\longrightarrow\quad\Ho_W(\tilde\Psi).
\end{gathered}
\] % [sic] \Ho_W(\Psi) -> \Ho_W(\tilde{\Psi}), the image is a diagram of type \tilde{\Psi} and the functor to be constructed is \mathcal{F}\longrightarrow\Ho_W(\tilde{\Psi})
\DERmark{ENV0823}

Let us check that it transforms arrows of $W_{\mathcal{F}}$ into isomorphisms.

\msp{page 250}

\[
\begin{tikzcd}[ampersand replacement=\&, row sep=small, column sep=small]
 \& X_1\times_S X_2 \arrow[dl] \arrow[dd] \& \& X'_1\times_{S'}X'_2 \arrow[ll, "f_3"'] \arrow[dl] \arrow[dd] \\
X_1 \arrow[dd] \& \& X'_1 \arrow[ll, crossing over, "f_1"'] \& \\
 \& X_2 \arrow[dl] \& \& X'_2 \arrow[ll, "f_2"'] \arrow[dl] \\
S \& \& S' \arrow[ll, "f_0"'] \arrow[from=uu, crossing over] \&
\end{tikzcd}
\]

We must prove that if $f_0,f_1,f_2\in W$, then $f_3\in W$. But let us consider the inverse image of the cartesian square $Q$ under $f_0 : S'\longrightarrow S$, whence a cartesian square $Q_{S'}$ and a factorization
\[
\begin{tikzcd}[ampersand replacement=\&]
Q \& Q_{S'} \arrow[l, "\tilde{f}"'] \& Q' \arrow[l, "g"'] \arrow[ll, bend left, "f"]
\end{tikzcd}
\]
of $Q'\longrightarrow Q$. Since $f_0\in W$, and since $X_1$, $X_2$, $X_1\times_S X_2$ are in $\uFib_W S$, the horizontal arrows realizing $Q_{S'}\longrightarrow Q$ are all in $W$. This, plus the fact that $f_1,f_2\in W$, imply that in the morphism $g$, the arrows $g_1$, $g_2$ are in $W$. It follows that the arrow
\[
g_3 : X'_1\times_{S'}X'_2\longrightarrow(X_1\times_S X_2)\times_S S'
\]
is as well (theory of the product in $\uFib_W S'$ and $\Ho^{\lc}_W(S')$), hence $f_3=\tilde{f}_3g_3\in W$, q.e.d.
\DERmark{ENV0824}


\msp{page 251}

Let us note that we have a canonical isomorphism
\[
\alpha : i^*i_*\xrightarrow{\sim}\mathrm{id}_{\Ho_W(\Psi)}. \tag*{$(*)$}
\]
\DERmark{ENV0825}

I feel like proving that the functors
\[
\begin{tikzcd}[ampersand replacement=\&]
\Ho_W(\Psi) \arrow[r, shift right, "i_*"'] \& \Ho_W(\tilde{\Psi}) \arrow[l, shift right, "i^*"']
\end{tikzcd}
\]
\DERmark{ENV0826}

\printedpage{163}
% head: A. Grothendieck : Dérivateurs, Ch. XII, Édition des Universités de Montpellier II et Paris VII
% foot: Version 30 mars 2003
are adjoint, and that $(*)$ is one of the adjunction arrows. We must define the other adjunction arrow
\[
\beta : i_*i^*\longleftarrow\mathrm{id}_{\Ho_W(\tilde\Psi)}. \tag*{$(**)$}
\]
\DERmark{ENV0827}

But for this, I must use the

\textbf{Lemma 2.} Let $\mathcal{G}$ be the full subcategory of $\mathcal{M}(\tilde{\Psi})$, consisting of the \begin{tikzcd}[ampersand replacement=\&]
X_1 \arrow[d, "p_1"'] \& X_3 \arrow[l, "q_2"'] \arrow[d, "q_1"] \\
S \& X_2 \arrow[l, "p_2"]
\end{tikzcd} such that $p_1,p_2\in\mathrm{Fib}_W$ and let $W_{\mathcal{G}}$ be the localizer induced by $W(\tilde{\Psi})$ (termwise $W$-equivalences). Then the inclusion $\mathcal{G}\longrightarrow\mathcal{M}(\tilde{\Psi})$ induces an equivalence
\DERmark{ENV0828}\DERmark{ENV0829}\DERmark{ENV0830}

\[
\mathcal{G}W_{\mathcal{G}}^{-1}\xrightarrow{\sim}\mathcal{M}(\tilde{\Psi})W(\tilde{\Psi})^{-1}\overset{\mathrm{def}}{=}\Ho_W(\tilde{\Psi}).
\]
\DERmark{ENV0831}\DERmark{ENV0832}\DERmark{ENV0833}\DERmark{ENV0834}


I shall for the moment assume this

\msp{page 252}

lemma, just like lemma 1 (p. 249). I am going to consider $i^*$, $i_*$ as arrows between the categories $\mathcal{F}W_{\mathcal{F}}^{-1}$, $\mathcal{G}W_{\mathcal{G}}^{-1}$. They are induced by the arrows % [sic] categories -> catégories, missing accent
\[
\begin{tikzcd}[ampersand replacement=\&]
\mathcal{F} \arrow[r, shift right, "i_{0\,*}"'] \& \mathcal{G} \arrow[l, shift right, "i_0^*"']
\end{tikzcd}
\]
between $\mathcal{F}$ and $\mathcal{G}$ themselves, % [sic] eux mêmes -> eux-mêmes, missing hyphen
\[
\begin{tikzcd}[ampersand replacement=\&]
 \& X_1 \arrow[dd] \& X_3 \arrow[l] \arrow[dd] \& \& X_1 \arrow[dd] \& \\
i_0^* :\quad \mathcal{G}\quad\longrightarrow\quad\mathcal{F} : \& \& \& \longmapsto \& \& \\
 \& S \& X_2 \arrow[l] \& \& S \& X_2 \arrow[l]
\end{tikzcd}
\]
\[
\begin{tikzcd}[ampersand replacement=\&]
 \& X_1 \arrow[dd] \& \& \& X_1 \arrow[dd] \& X_1\times_S X_2 \arrow[l] \arrow[dd] \\
i_{0\,*} :\quad \mathcal{F}\quad\longrightarrow\quad\mathcal{G} : \& \& \& \longmapsto \& \& \\
 \& S \& X_2 \arrow[l] \& \& S \& X_2. \arrow[l]
\end{tikzcd}
\]
These two functors are adjoint to one another, and compatible with the localizers $W_{\mathcal{F}}$, $W_{\mathcal{G}}$. Thus the functors deduced from $i_0^*$, $i_{0\,*}$ by passage to the categories of fractions are also adjoint (cf. lemma p. 248).

Thus, modulo the proof of lemmas 1 and 2, we have proved this:

\printedpage{164}
% head: A. Grothendieck : Dérivateurs, Ch. XII, Édition des Universités de Montpellier II et Paris VII
% foot: Version 30 mars 2003

\msp{page 253}

\textbf{Theorem.} {\itshape Let $(\mathcal{M},W)$ be a category of models ($W\subseteq\mathit{Fl}\,\mathcal{M}$) satisfying the left factorization axiom (p. 240), thus giving rise to a good theory of $W$-homotopic fibers. Let
\[
\begin{tikzcd}[ampersand replacement=\&, row sep=small, column sep=small]
 \& 1 \& \& \& \& 1 \arrow[dr] \& \\
\Psi=0 \arrow[ur] \arrow[dr] \& \& {} \arrow[r, "i"] \& {} \& \tilde{\Psi}=0 \arrow[ur] \arrow[dr] \& \& 3 \\
 \& 2 \& \& \& \& 2 \arrow[ur] \&
\end{tikzcd}
\]
(open immersions in} $\mathrm{Cat}${\itshape ), and let us consider the restriction functor
\[
i^* : \Ho_W(\tilde{\Psi})\longrightarrow\Ho_W(\Psi).
\]
This functor admits a right adjoint
\[
i_* : \Ho_W(\Psi)\longrightarrow\Ho_W(\tilde{\Psi}).
\]
This functor $i_*$ is ``computed'' as follows: its value at \begin{tikzcd}[ampersand replacement=\&]
X_1 \arrow[d, "p_1"'] \& \\
S \& X_2 \arrow[l, "p_2"]
\end{tikzcd} is canonically isomorphic (in} $\Ho_W(\tilde{\Psi})${\itshape ) to \begin{tikzcd}[ampersand replacement=\&]
X_1 \arrow[d] \& X_1\times_S X_2 \arrow[l] \arrow[d] \\
S \& X_2 \arrow[l]
\end{tikzcd}, if we suppose that $p_1$ or $p_2$ is a $W$-fibration. In the general case it is computed by factoring $p_1$ as
\[
X_1\xrightarrow{i_1}\bar{X}_1\xrightarrow{p'_1}S,\qquad i_1\in W,\ p'_1\in\mathit{Fib}_W,
\]
and by taking \begin{tikzcd}[ampersand replacement=\&]
\bar{X}_1 \arrow[d] \& \bar{X}_1\times_S X_2 \arrow[l] \arrow[d] \\
S \& X_2 \arrow[l]
\end{tikzcd}, or by

\msp{page 254}

proceeding symmetrically, or by simultaneously factoring $p_1$ and $p_2$, and taking % [sic] simultanement -> simultanément, missing accent
\[
\begin{tikzcd}[ampersand replacement=\&]
\bar{X}_1 \arrow[d] \& \bar{X}_1\times_S\bar{X}_2 \arrow[l] \arrow[d] \\
S \& \bar{X}_2. \arrow[l]
\end{tikzcd}
\]}


\endgroup
% END UNIT Der-ChXIIc-30-03-03/Der-ChXIIc-30-03-03_p046-050.tex
% BEGIN UNIT Der-ChXIIc-30-03-03/Der-ChXIIc-30-03-03_p051-055.tex
\begingroup\makeatletter%
\NoAutoSpacing%
\tikzcdset{ampersand replacement=\&}%
\newcommand{\printedpage}[1]{\marginpar{\footnotesize p.~#1}}%
\newcommand{\msp}[1]{\marginpar{\footnotesize [#1]}}%
\newcommand{\Ho}{\mathrm{Ho}}%
\newcommand{\Fib}{\mathrm{Fib}}%
\makeatother\let\printedpage\derpp\let\msp\dermsp%
\let\mspin\DERmspin%


\printedpage{165}
% head: A. Grothendieck : Dérivateurs, Ch. XII, Édition des Universités de Montpellier II et Paris VII
% foot: Version 30 mars 2003

\textit{The functorial law is obtained with the help of lemma 1 p. 249, which allows us to restrict ourselves to the arguments $\begin{tikzcd}[ampersand replacement=\&] X_1 \arrow[d, "p_1"'] \& \\ S \& X_2 \arrow[l, "p_2"] \end{tikzcd}$ such that $p_1,p_2\in \mathit{Fib}_W$, so that we find a functor}
\[
\begin{tikzcd}[ampersand replacement=\&]
X_1 \arrow[d] \& \\
S \& X_2 \arrow[l]
\end{tikzcd}
\qquad\longmapsto\qquad
\begin{tikzcd}[ampersand replacement=\&]
X_1 \arrow[d] \& X_1\times_S X_2 \arrow[l] \arrow[d] \\
S \& X_2, \arrow[l]
\end{tikzcd}
\]
\textit{compatible with the localizers, and passing to the categories of fractions} $\Ho_W(\Psi)$ \textit{and} $\Ho_W(\tilde{\Psi})$ \textit{respectively.}

\textbf{NB} By proceeding in a similar, simpler way (with no need to resort to lemma 1), we find the existence of a \textit{left adjoint} $i_!$ to $i^* : \Ho_W(\tilde\Psi)\longrightarrow\Ho_W({\Psi})$, subject to the existence of an initial object $\varphi$ in $\mathcal{M}$. This functor $i_!$ is given by % [sic] Cet foncteur -> Ce foncteur, "foncteur" begins with a consonant
\DERmark{ENV0835}


\msp{page 255}
\[
i_!\left(
\begin{tikzcd}[ampersand replacement=\&]
X_1 \arrow[d] \& \\
S \& X_2 \arrow[l]
\end{tikzcd}
\right)
\qquad = \qquad
\left(
\begin{tikzcd}[ampersand replacement=\&]
X_1 \arrow[d] \& \varphi \arrow[l] \arrow[d] \\
S \& X_2 \arrow[l]
\end{tikzcd}
\right),
\]
the functorial law being evident.

This is not yet the end of the sorites! We must still interpret the composite functor
\[
\Ho_W(\Psi)\xrightarrow{\ i_*\ }\Ho_W(\tilde{\Psi})\xrightarrow{\ \varepsilon_3^*\ }\Ho_W(e)=\Ho_W,
\]
where
\[
\varepsilon_3 : e\xrightarrow{\ \sim\ }\{3\}
\qquad\longrightarrow\qquad
\begin{tikzcd}[ampersand replacement=\&, column sep=small, row sep=small]
\& 1 \arrow[dl] \& \\
3 \& \& 0 \arrow[ul] \arrow[dl] \\
\& 2 \arrow[ul] \&
\end{tikzcd}
\]
is the inclusion of vertex 3 in $\tilde{\Psi}$. If $\Ho_W(\Psi)$ is interpreted, in accordance with lemma 1, as $\mathcal{F}W_{\mathcal{F}}^{-1}$, this composite functor is none other than
\[
\begin{tikzcd}[ampersand replacement=\&]
X_1 \arrow[d] \& \\
S \& X_2 \arrow[l]
\end{tikzcd}
\qquad\longmapsto\qquad X_1\times_S X_2
\]
\printedpage{166}
% head: A. Grothendieck : Dérivateurs, Ch. XII, Édition des Universités de Montpellier II et Paris VII
% foot: Version 30 mars 2003
(for $\left(\begin{tikzcd}[ampersand replacement=\&] X_1 \arrow[d, "p_1"'] \& \\ S \& X_2 \arrow[l, "p_2"] \end{tikzcd}\right)\in \mathrm{Ob}\,\mathcal{F}$, i.e. $p_1,p_2\in\Fib_W$.)

We want to interpret this functor as a $W$-homotopic $\varprojlim$, i.e. viewed as a ``$W$-homotopic fiber product''. To be precise, we have this:

\msp{page 256}
\textbf{Corollary 1.} \textit{Consider $p_\Psi : \Psi\longrightarrow e$, and the functor}
\[
p_\Psi^* : \Ho_W(e)=\Ho_W\longrightarrow\Ho_W(\Psi),
\]
\textit{deduced by passage to the categories of fractions from the functor}
\[
\begin{array}{ccc}
\mathcal{M} & \longrightarrow & \mathcal{M}(\Psi) \\
X & \longrightarrow & \left(
\begin{tikzcd}[ampersand replacement=\&]
X \arrow[d] \& \\
X \& X \arrow[l]
\end{tikzcd}
\right).
\end{array}
\]
\textit{This functor admits the preceding functor $\varepsilon_3^*i_*$ as a right adjoint, i.e.}
\[
p_{\Psi\,*}=\varepsilon_3^*i_* : \Ho_W(\Psi)\longrightarrow\Ho_W.
\]

\textsc{Proof}. We can regard $p_\Psi$ as the composite % [sic] p_\psi -> p_\Psi, the subscript is the capital \Psi everywhere else
\DERmark{ENV0836}

\[
\Psi\xrightarrow{\ i\ }\tilde{\Psi}\xrightarrow{\ p_{\tilde{\Psi}}\ }e,
\]
and we therefore have
\[
p_\Psi^*\simeq i^*(p_{\tilde{\Psi}})^*.
\]
Thus the existence and the determination of a right adjoint for $p_\Psi^*$ will follow from the existence and the determination of right adjoints $i_*$, $(p_{\tilde{\Psi}})_*$ of $i^*$, $(p_{\tilde{\Psi}})^*$, via
\[
(p_\Psi)_*\simeq(p_{\tilde{\Psi}})_*i_*.
\]

\msp{page 257}
Thus corollary 1 follows from the formula
\[
p_{\tilde{\Psi}\,*}=\varepsilon_3^*,
\]
which itself is a special case of the

\printedpage{167}
% head: A. Grothendieck : Dérivateurs, Ch. XII, Édition des Universités de Montpellier II et Paris VII
% foot: Version 30 mars 2003

\textbf{Lemma.} \textit{Let $I$ be in} Cat\textit{, with final object $e$, consider the inclusion}
\[
\varepsilon : e\longrightarrow I,\qquad \textit{and } p : I\longrightarrow e.
\]
\textit{Let $(\mathcal{M},W)$ be a category of models, and consider}
\[
\Ho_W(I)\xleftarrow{\ p^*\ }\Ho_W.
\]
\textit{This functor admits a right adjoint} \texttt{[rather right]} \textit{$p_*$, which is isomorphic to $\varepsilon^*$.}
\DERmark{ENV0837}


\textsc{Proof}. Consider the pair of adjoint functors
\[
\begin{tikzcd}[ampersand replacement=\&, column sep=large]
\mathcal{M}(I) \arrow[r, shift right, "p_*=\varprojlim\limits_I"'] \& \mathcal{M} \arrow[l, shift right, "p^*"']
\end{tikzcd}\ ;
\]
as $I$ admits the final object $e$, the functor $p_*$ is isomorphic to $\varepsilon^*$. These functors $p^*$, $\varepsilon^*$ are compatible with the localizers $W(I)$, $W$, so we conclude by the lemma on page 248.

It only remains to prove lemmas 1 and 2. I will do it tomorrow.

\msp{page 258}
Today, I still feel like giving the sorites of $W$-cartesian squares in a

\textbf{Proposition 19 :} \textit{Let $\begin{tikzcd}[ampersand replacement=\&] X_1 \arrow[d, "p_1"'] \& X_3 \arrow[l, "f_1"'] \arrow[d, "f_2"] \\ S \& X_2 \arrow[l, "p_2"] \end{tikzcd}$ be a commutative square in $\mathcal{M}$. Consider factorizations of $p_1$, $p_2$ as}
\[
X_1\xrightarrow{\ i_1\ }\bar{X}_1\xrightarrow{\ \bar{p}_1\ }S,\qquad X_2\xrightarrow{\ i_2\ }\bar{X}_2\xrightarrow{\ \bar{p}_2\ }S,
\]
\textit{$i_1,i_2\in W$, $\bar{p}_1,\bar{p}_2\in F_W$} \texttt{[$=$ Fib$_W$]}\textit{, and consider the arrows induced by $X_3\longrightarrow X_1\times_S X_2$ (deduced from $Q$) via $\begin{tikzcd}[ampersand replacement=\&, row sep=small] \& \bar{X}_1\times_S X_2 \\ X_1\times_S X_2 \arrow[ur] \arrow[r] \arrow[dr] \& \bar{X}_1\times_S\bar{X}_2 \\ \& X_1\times_S\bar{X}_2 \end{tikzcd}$, so that we find a commutative diagram}
\[
\begin{tikzcd}[ampersand replacement=\&, row sep=small]
\& \bar{X}_1\times_S X_2 \arrow[d, "j_2"'] \\
X_3 \arrow[ur, "\pi_1"] \arrow[r, "\pi_0"] \arrow[dr, "\pi_2"] \& \bar{X}_1\times_S\bar{X}_2 \\
\& X_1\times_S\bar{X}_2, \arrow[u, "j_1"]
\end{tikzcd}
\]
\textit{where $j_1=i_1\times_S\mathrm{id}_{\bar{X}_2}$ and $j_2=\mathrm{id}_{\bar{X}_1}\times_S i_2$ are in $W$ (since $i_1$ and $i_2$ are, and $\bar{X}_1$ and $\bar{X}_2$ are in $\underline{\mathrm{Fib}}_W S$). We therefore have}
\DERmark{ENV0838}

\[
(\pi_1\in W)\qquad\iff\qquad(\pi_2\in W)\qquad\qquad(\iff\ \pi_0\in W),
\]
\textit{which proves that the condition $\pi_1\in W$, which seems to depend on the choice of the factorization through $i_1$, does not depend on it, and likewise for $\pi_2\in W$, and for $\pi_0\in W$ (which seems to depend on both $i_1$, and $i_2$).}

\printedpage{168}
% head: A. Grothendieck : Dérivateurs, Ch. XII, Édition des Universités de Montpellier II et Paris VII
% foot: Version 30 mars 2003

\msp{page 259}
\textbf{Definition.} \textit{When these conditions are satisfied, we say that the square $Q$ is $W$-cartesian.}

\textbf{Corollary 1.} \textit{If $Q$ is $W$-cartesian, then the image of $Q$ in} $\Ho_W(\tilde{\Psi})$ \textit{is in the essential image of}
\[
i_* : \Ho_W(\Psi)\longrightarrow\Ho_W(\tilde{\Psi}),
\]
\textit{i.e. the canonical arrow in} $\Ho_W(\tilde{\Psi})$
\[
Q\longrightarrow i_*i^*Q
\]
\textit{is} \texttt{[an]} \textit{isomorphism. The converse is true if $W$ is strongly saturated.}

\textbf{Corollary 2.} \textit{Consider a commutative diagram in $\mathcal{M}$}
\[
\begin{tikzcd}[ampersand replacement=\&]
X_0' \arrow[d] \& X_1' \arrow[l] \arrow[d] \arrow[dl, phantom, "\textit{1}"] \& X_2' \arrow[l] \arrow[d] \arrow[dl, phantom, "\textit{2}"] \\
X_0 \& X_1 \arrow[l] \& X_2, \arrow[l]
\end{tikzcd}
\tag*{(D)}
\]
\textit{whence two commutative squares 1, 2, and a composite square}
\[
\begin{tikzcd}[ampersand replacement=\&]
X_0' \arrow[d] \& X_2' \arrow[l] \arrow[d] \arrow[dl, phantom, "\textit{3}"] \\
X_0 \& X_2. \arrow[l]
\end{tikzcd}
\]
\textit{Suppose square 1 is $W$-cartesian. Then square 2 is $W$-cartesian if and only if square 3 is.}

\textsc{Proof}. Consider a factorization of $\begin{tikzcd}[ampersand replacement=\&, row sep=small] X_0' \arrow[d] \\ X_0 \end{tikzcd}$ as
\[
\begin{tikzcd}[ampersand replacement=\&, column sep=small, row sep=small]
\bar{X}_0' \arrow[dr] \& X_0' \arrow[l] \arrow[d] \\
\& X_0
\end{tikzcd}
\]
with $(X_0'\longrightarrow\bar{X}_0')\in W$, and $(\bar{X}_0'\longrightarrow X_0)\in\Fib_W$. We can then replace $X_0'$ by $\bar{X}_0'$ in (D), obtaining a diagram $(\bar{D})$, and we see at once
\printedpage{169}
% head: A. Grothendieck : Dérivateurs, Ch. XII, Édition des Universités de Montpellier II et Paris VII
% foot: Version 30 mars 2003
\msp{page 260}
that each of the squares 1(D), 2(D), 3(D) is $W$-cartesian if and only if the corresponding square $1(\bar{D})$, $2(\bar{D})$ or $3(\bar{D})$ is. This shows that we can assume $X_0'\longrightarrow X_0$ to be in $\Fib_W$. Consider then the commutative diagram deduced from (D)
\DERmark{ENV0839}

\[
\begin{tikzcd}[ampersand replacement=\&, column sep=small]
X_0' \arrow[d] \& \tilde{X}_1'\overset{\mathrm{def}}{=}X_0'\times_{X_0}X_1 \arrow[l] \arrow[d] \& X_1' \arrow[l, "\alpha"'] \& X_2' \arrow[l] \arrow[d] \\
X_0 \& X_1 \arrow[l] \& \& X_2, \arrow[ll]
\end{tikzcd}
\]
\DERmark{ENV0840}

where $\alpha\in W$ by virtue of the $W$-cartesian hypothesis on square 1. We see again that square 2 is $W$-cartesian if and only if the corresponding square with $X_1'$ replaced by $\tilde{X}_1'$ is, as for squares 1, 3, they do not change under this modification. This shows that we can assume square 1 \textit{cartesian.} Consider then the diagram
\[
\begin{tikzcd}[ampersand replacement=\&]
X_0' \arrow[d, "\mathrm{Fib}_W"'] \& X_1' \arrow[l] \arrow[d, "\mathrm{Fib}_W"] \arrow[dl, phantom, "\text{cart.}"] \& \tilde{X}_2'\overset{\mathrm{def}}{=}X_1'\times_{X_1}X_2 \arrow[l] \arrow[d, "\mathrm{Fib}_W"] \arrow[dl, phantom, "\text{cart.}"] \& X_2' \arrow[l, "\beta"'] \\
X_0 \& X_1 \arrow[l] \& X_2. \arrow[l]
\end{tikzcd}
\]
\DERmark{ENV0841}

Then proposition 19 and the definition show that square 2 is $W$-cartesian if and only if $\beta\in W$, and likewise for square 3. Which completes the proof. % [sic] seulenment -> seulement, misspelling

\textbf{Corollary 2.} \textit{Let $\begin{tikzcd}[ampersand replacement=\&] X_1 \arrow[d, "p_1"'] \& X_3 \arrow[l, "f_2"'] \arrow[d, "f_1"] \\ S \& X_2 \arrow[l, "p_2"] \end{tikzcd}$ be a $W$-cartesian square. Then}
\[
\begin{aligned}
p_1\in W\quad&\Longrightarrow\quad f_1\in W\\
p_2\in W\quad&\Longrightarrow\quad f_2\in W,
\end{aligned}
\]
\textit{more precisely, if $p_i\in W$, then $Q$ is $W$-cartesian if and only if $f_i\in W$.}

\msp{page 261}
Suppose $p_1\in W$. Then in the factorization of $p_1$ as $X_1\longrightarrow\bar{X}_1\longrightarrow S$, we can take $\bar{X}_1=S$. Thus the square is $W$-cartesian if and only if
\[
(X_3\longrightarrow\bar{X}_1\times_S X_2=X_2)\in W,
\]
i.e. $f_1\in W$, q.e.d.

Of course, we cannot in general assert the reverse implications: $f_i\in W$ by no means implies $p_i\in W$.


\endgroup
% END UNIT Der-ChXIIc-30-03-03/Der-ChXIIc-30-03-03_p051-055.tex
% BEGIN UNIT Der-ChXIIc-30-03-03/Der-ChXIIc-30-03-03_p056-060.tex
\begingroup\makeatletter%
\NoAutoSpacing%
\newcommand{\printedpage}[1]{\marginpar{\footnotesize p.~#1}}%
\newcommand{\msp}[1]{\marginpar{\footnotesize [#1]}}%
\newcommand{\Ho}{\mathrm{Ho}}%
\newcommand{\Cof}{\mathrm{Cof}}%
\makeatother\let\printedpage\derpp\let\msp\dermsp%
\let\mspin\DERmspin%


\printedpage{170}
% head: A. Grothendieck : Dérivateurs, Ch. XII, Édition des Universités de Montpellier II et Paris VII

All these results, and the notion of a $W$-cartesian square, dualize (using the appropriate axiom of right factorization, cf. page 240) into a theory of $W$-homotopic amalgamated sums, and of $W$-cocartesian squares. This time we work with the category
\[
\Phi \;=\; \left(
\begin{tikzcd}[ampersand replacement=\&, row sep=small, column sep=small]
 \& 1 \arrow[dl] \\
0 \& \\
 \& 2 \arrow[ul]
\end{tikzcd}
\right)
\qquad\quad \overset{j}{\hookrightarrow} \qquad\quad
\tilde{\Phi} \;=\; \left(
\begin{tikzcd}[ampersand replacement=\&, row sep=small, column sep=small]
 \& 1 \arrow[dl] \& \\
0 \& \& 3 \arrow[ul] \arrow[dl] \\
 \& 2 \arrow[ul] \&
\end{tikzcd}
\right),
\]
$j$ being a closed immersion. We have a restriction functor
\[
j^* : \Ho_W(\tilde{\Phi}) \longrightarrow \Ho_W(\Phi),
\]
for which we construct a left adjoint
\[
j_! : \Ho_W(\Phi) \longrightarrow \Ho_W(\tilde{\Phi}).
\]
And we find a natural notion of a
\msp{page 262}
$W$-cocartesian square, which seems irreproachable from every point of view, and of perfect simplicity. % [sic] tout points -> tous points, agreement of "tout"

This leads us to look again at the notion in question in Cat, where it looked very clever, deep and a little obscure, with integrals and all that. This is the moment to realize that ``that's all it is''. So to prove the

\textbf{Proposition 20.} \textit{Suppose} $\mathcal{M} = \mathrm{Cat}$\textit{, $W$ a basic localizer satisfying W(7 bis). Consider a commutative square in} Cat
\[
Q = \left(
\begin{tikzcd}[ampersand replacement=\&]
S \arrow[r] \arrow[d] \& X_1 \arrow[d] \\
X_2 \arrow[r] \& X_3
\end{tikzcd}
\right).
\]
\textit{For it to be $W$-cocartesian} \texttt{[rather} $W$\texttt{-cocartesian]} \textit{in the sense of the theory developed here (for general categories of models), it is necessary and sufficient that}
\DERmark{ENV0842}

\[
\varphi_Q : \quad \mathcal{X}_Q = \left(
\begin{tikzcd}[ampersand replacement=\&, row sep=small, column sep=small]
  \& X_1 \\
\int S \arrow[ur] \arrow[dr] \& \\
 \& X_2
\end{tikzcd}
\right) \qquad\quad \longrightarrow \qquad\quad X_3
\]
\DERmark{ENV0843}


\texttt{[rather}
\[
\varphi_Q : \quad \mathcal{X}_Q = \left(
\begin{tikzcd}[ampersand replacement=\&, row sep=small, column sep=small]
 \& X_1 \\
\int S \arrow[ur] \arrow[dr] \& \\
 \& X_2
\end{tikzcd}
\right) \qquad\quad \longrightarrow \qquad\quad X_3\text{\texttt{]}}
\]
\printedpage{171}
% head: A. Grothendieck : Dérivateurs, Ch. XII, Édition des Universités de Montpellier II et Paris VII
\textit{be in $W$.} (NB The interest of this criterion is that it is ``intrinsic'' to $Q$, not calling on a factorization of $\beta_1$ or of $\beta_2$ \texttt{[where} $\beta_i : S \longrightarrow X_i$\texttt{]} into $S \xrightarrow{\;\in \Cof_W\;} \tilde{X}_i \xrightarrow{\in W} X_i$. Cf. interesting application in section V 3.4., p. 146 ff., th. 4 and corollaries.)

\msp{page 263}
\textsc{Proof.} Consider a factorization of $S \longrightarrow X_1$ as
\[
S \xrightarrow{i_1} \tilde{X}_1 \xrightarrow{p_1} X_1, % [sic] X -> X_1, the factorisation is of S -> X_1
\]
\DERmark{ENV0844}

with $i_1$ an open Dwyer immersion, $p_1 \in W$. So we have a commutative square
\[
\tilde{Q} = \left(
\begin{tikzcd}[ampersand replacement=\&]
S \arrow[r, hook, "\mathrm{Dw}"] \arrow[d] \& \tilde{X}_1 \arrow[d] \\
X_2 \arrow[r] \& X_3
\end{tikzcd}
\right)
\]
and a homomorphism $\tilde{Q} \xrightarrow{\alpha} Q$ which is the identity on the three vertices other than $\tilde{X}_1$, and which is $p_1$ at $\tilde{X}_1$. We deduce from this a commutative diagram
\[
\begin{tikzcd}[ampersand replacement=\&, row sep=small, column sep=small]
\mathcal{X}_Q \arrow[rr, leftarrow, "\mathcal{X}_\alpha = \lambda"] \arrow[dr, "\varphi_Q"'] \& \& \mathcal{X}_{\tilde{Q}} \arrow[dl, "\varphi_{\tilde{Q}}"] \\
 \& X_3, \&
\end{tikzcd}
\]
\DERmark{ENV0845}

and we have $\lambda \in W$ (for it is a $\Phi$-arrow of fibered categories over $\Phi$, which induces a $W$-equivalence on each of three fibers.) Thus $\varphi_Q \in W \iff \varphi_{\tilde{Q}} \in W$, on the other hand it is clear that $Q$ is $W$-cocartesian if and only if $\tilde{Q}$ is. So we may assume that $S \hookrightarrow X_1$ is \texttt{[an]} open (Dwyer) immersion (\footnote[134]{Here we use W(3) (strong form), i.e. Loc(3 bis) in the notation of XVI. But we can do without it, cf. below.}). Applying this
\msp{page 264}
with the roles of $X_1$ and $X_2$ interchanged, we see that we may moreover assume that $S \longrightarrow X_2$ is also \texttt{[an]} open Dwyer immersion. Consider then
\[
\begin{tikzcd}[ampersand replacement=\&]
S \arrow[r, hook] \arrow[d, hook] \arrow[dr, phantom, "Q'"] \& X_1 \arrow[d, hook] \& \\
X_2 \arrow[r, hook] \& X_1 \amalg_S X_2 \arrow[r, "\pi"] \& X_3 \\
 \& \mathcal{X}_Q \arrow[u, "\varphi_{Q'}"] \arrow[ur, "\varphi_Q"'] \&
\end{tikzcd}
\]
\printedpage{172}
% head: A. Grothendieck : Dérivateurs, Ch. XII, Édition des Universités de Montpellier II et Paris VII
We have $\varphi_Q = \pi\varphi_{Q'}$, on the other hand we know that $\varphi_{Q'} \in W$ (cor. 3, p. 149). Thus
\[
\pi \in W \qquad\quad \iff \qquad\quad \varphi_Q \in W,
\]
where $\pi \in W$ means that $Q$ is $W$-cocartesian, q.e.d.

\msp{page 265}

\section*{8\quad Appendix : Return to the axioms W(7), W(7 bis) (pages 75, 76).}

(\footnote[135]{Provisional reflection, superseded by XVI.}). We have seen (prop. 6, p. 78) that W(7 bis) implies W(7). Let us prove a partial converse. % [sic] Reflexion -> Réflexion, missing accent; supplanté -> supplantée, agreement (in footnote 135)

\textbf{Proposition.} \textit{Suppose} W(7 bis) (cocartesian square axiom) \textit{satisfied, and let us place ourselves under the conditions of the statement of} W(7), \textit{and more generally, in the case of a cubical diagram}
\[
\begin{tikzcd}[ampersand replacement=\&, row sep=small, column sep=small]
 \& X_2 \arrow[rr] \arrow[dd, "f_2" near end] \& \& X_3 \arrow[dd, "f_3"] \\
X_0 \arrow[ur] \arrow[rr, crossing over] \arrow[dd, "f_0"'] \& \& X_1 \arrow[ur] \& \\
 \& Y_2 \arrow[rr] \& \& Y_3 \\
Y_0 \arrow[ur] \arrow[rr] \& \& Y_1 \arrow[ur] \arrow[from=uu, crossing over, "f_1"' near end] \&
\end{tikzcd}
\]
\textit{where the two squares $X$ and $Y$ (upper and lower face of the cube) are $W$-cocartesian. Consider the following conditions on this diagram}
\begin{enumerate}[label={}]
\item[\textit{a)}] \textit{$f_0 \in W$.}
\item[\textit{b)}] \textit{$f_1$ and $f_2 \in W$.}
\item[\textit{c)}] \textit{$f_3 \in W$.}
\end{enumerate}
\textit{Then we have the implications}
\[
a \text{ \textit{\&} } b \Longrightarrow c, \qquad\quad b \text{ \textit{\&} } c \Longrightarrow a.
\]
\textit{Moreover, if $W$ is stable under direct factors (e.g. $W$ strongly saturated), then we also have}
\[
a \text{ \textit{\&} } c \Longrightarrow b,
\]
\textit{thus the conjunction} \texttt{[of]} \textit{any two among the conditions a) b) c) implies the third.}

\textsc{Proof.} Using the functorial factorization
\printedpage{173}
% head: A. Grothendieck : Dérivateurs, Ch. XII, Édition des Universités de Montpellier II et Paris VII
\msp{page 266}
of an arbitrary arrow (here the $f_i$) as $pi$, with $i$ an open Dwyer immersion and $p \in W$ (p. 229), we are reduced to the case where the $f_i$ are open Dwyer immersions. But for any arrow $g : S \longrightarrow T$ which is in $\Cof_W$, let $Z(g)$ be its mapping cone % [sic] remené -> ramené, spelling
\[
\begin{tikzcd}[ampersand replacement=\&, row sep=small]
S \arrow[rr, "g"] \arrow[dr] \& \& T \arrow[dr, "\alpha(g)"] \& \\
 \& e \arrow[rr, "\bar{g}"] \& \& Z(g),
\end{tikzcd}
\]
and
\[
\alpha(g) : T = \mathrm{but}(g) \longrightarrow Z(g)
\]
the canonical arrow (\footnote[136]{This is a bit quick, for it is not clear that $X_1' \amalg_{X_0} X_2' \longrightarrow X_1 \amalg_{X_0} X_2$ is \texttt{[in]} $W$, without already using one of the three ingredients of W(7) (namely $f_0, f_1, f_2 \in W \implies f_3 \in W$). It seems then that our argument works only with an additional hypothesis of this sort \dots But there is \textit{Brown's lemma,} cf. XVI \S3, notably p. 20.}). This is functorial in $g$, so that the given cubical diagram extends to a double cube
\[
\begin{tikzcd}[ampersand replacement=\&, row sep=small, column sep=small]
 \& X_2 \arrow[rr] \arrow[dd, "f_2" near end] \& \& X_3 \arrow[dd, "f_3"] \\
X_0 \arrow[ur] \arrow[rr, crossing over] \arrow[dd, "f_0"'] \& \& X_1 \arrow[ur] \& \\
 \& Y_2 \arrow[rr] \arrow[dd, "\alpha_2" near end] \& \& Y_3 \arrow[dd, "\alpha_3"] \\
Y_0 \arrow[ur] \arrow[rr, crossing over] \arrow[dd, "\alpha_0"'] \& \& Y_1 \arrow[ur] \arrow[from=uu, crossing over, "f_1"' near end] \& \\
 \& Z_2 \arrow[rr] \& \& Z_3 \\
Z_0 \arrow[ur] \arrow[rr] \& \& Z_1 \arrow[ur] \arrow[from=uu, crossing over, "\alpha_1"' near end] \&
\end{tikzcd}
\]
with
\[
Z_i = Z(f_i),\ \alpha_i = \alpha(f_i) \qquad\quad \text{for } i \in \{0,1,2,3\}.
\]
Note that W(7 bis) implies that if $g \in \Cof_W$,
\msp{page 267}
we have
\[
g \in W \qquad\quad \iff \qquad\quad \bar{g} \in W, \qquad\quad \text{i.e. } Z(g)\ W\text{-aspheric},
\]
in particular
\[
f_i \in W \qquad\quad \iff \qquad\quad Z_i\ W\text{-aspheric}.
\]
Thus, the statement to be proved about the initial cube reduces to a statement about the square
\printedpage{174}
% head: A. Grothendieck : Dérivateurs, Ch. XII, Édition des Universités de Montpellier II et Paris VII
\[
\begin{tikzcd}[ampersand replacement=\&, row sep=small, column sep=small]
 \& Z_1 \arrow[dr] \& \\
Z_0 \arrow[ur] \arrow[dr] \& \& Z_3, \\
 \& Z_2 \arrow[ur] \&
\end{tikzcd}
\]
namely that the three conditions
\begin{enumerate}[label={}]
\item[a$'$)] $Z_0$ $W$-aspheric.
\item[b$'$)] $Z_1$ and $Z_2$ $W$-aspheric.
\item[c$'$)] $Z_3$ $W$-aspheric.
\end{enumerate}
giving rise to the implications
\[
\mathrm{a}' \text{ \& } \mathrm{b}' \Longrightarrow \mathrm{c}', \qquad\quad \mathrm{b}' \text{ \& } \mathrm{c}' \Longrightarrow \mathrm{a}',
\]
and that if $W$ is stable under direct factors, then
\[
\mathrm{a}' \text{ \& } \mathrm{c}' \Longrightarrow \mathrm{b}'.
\]

This is more or less the corollary on page 76, supplemented on page 77 (corollary), except that there the arrows are assumed to be open immersions, and the square cartesian (hence also cocartesian) in Cat, i.e. $Z_3 = Z_1 \cup Z_2$, $Z_0 = Z_1 \cap Z_2$. But
\msp{page 268}
what matters for the validity of this statement is (as one might suspect) that $Z_*$ be $W$-cocartesian. So we must first prove this.

\textbf{Lemma.} \textit{If the cubical diagram p. 265 is such that the faces $X$, $Y$ are $W$-cocartesian, and the $f_i$ ($i \in \{0,1,2,3\}$) $W$-cofibrations, then the square $Z$ of the cofibers of the $f_i$ is itself also $W$-cocartesian.}

This lemma follows from the theory of the ``unfolding of cubical diagrams'', relative to a category of models $(\mathcal{M}, W)$ admitting a theory of $W$-homotopic cofibers, which I really shall have to write! (But the existence for $(\mathrm{Cat}, W)$ of such a theory of $W$-homotopic cofibers, and hence of that of unfoldings of diagrams in $\mathcal{M}$ of type $(\boldsymbol{\Delta}^1)^I$, $I$ a finite index set, depends only on W(7 bis) - we do not have to call on W(7).)

\msp{page 269}
This ``unfolding'' is an infinite diagram all of whose infinite sequences, in the three directions of the coordinate axes, are ``suspension sequences''. I reproduce the useful part of the unfolding. % [sic] deploiement -> déploiement, missing accent


\endgroup
% END UNIT Der-ChXIIc-30-03-03/Der-ChXIIc-30-03-03_p056-060.tex
% BEGIN UNIT Der-ChXIIc-30-03-03/Der-ChXIIc-30-03-03_p061-064.tex
\begingroup\makeatletter%
\newcommand{\printedpage}[1]{\marginpar{\footnotesize p.~#1}}%
\newcommand{\msp}[1]{\marginpar{\footnotesize [#1]}}%
\newcommand{\Asph}{\mathrm{Asph}}%
\newcommand{\Cof}{\mathrm{Cof}}%
\makeatother\let\printedpage\derpp\let\msp\dermsp%
\let\mspin\DERmspin%


\printedpage{175}
% head: A. Grothendieck : Dérivateurs, Ch. XII, Édition des Universités de Montpellier II et Paris VII
% foot: Version 30 mars 2003

% dangling arrows (from X_2, X_3, Y_2, Y_3, Z_2, Z_3) as printed
\[
\begin{tikzcd}[ampersand replacement=\&, column sep=small, row sep=small]
 \& \& {} \& \& {} \& \& U \arrow[dd, "\mu"] \\
 \& X_2 \arrow[rr] \arrow[dd] \arrow[ur, shorten >=1ex] \& \& X_3 \arrow[rr] \arrow[dd] \arrow[ur, shorten >=1ex] \& \& T \arrow[ur] \arrow[dd] \& \\
X_0 \arrow[rr, crossing over] \arrow[ur] \arrow[dd] \& \& X_1 \arrow[rr, crossing over] \arrow[ur] \arrow[dd, crossing over] \& \& S \arrow[ur, "\lambda"] \arrow[dd, crossing over] \& \& U' \arrow[dd] \\
 \& Y_2 \arrow[rr] \arrow[dd] \arrow[ur, shorten >=1.5ex] \& \& Y_3 \arrow[rr] \arrow[dd] \arrow[ur, shorten >=1.5ex] \& \& T' \arrow[ur] \arrow[dd] \& \\
Y_0 \arrow[rr, crossing over] \arrow[ur] \arrow[dd] \& \& Y_1 \arrow[rr, crossing over] \arrow[ur] \arrow[dd, crossing over] \& \& S' \arrow[ur, "\lambda'"] \arrow[dd, crossing over] \& \& U'' \\
 \& Z_2 \arrow[rr] \arrow[ur, shorten >=1.5ex] \& \& Z_3 \arrow[rr] \arrow[ur, shorten >=1.5ex] \& \& T'' \arrow[ur] \& \\
Z_0 \arrow[rr] \arrow[ur] \& \& Z_1 \arrow[rr] \arrow[ur] \& \& S'' \arrow[ur, "\lambda''"] \& \&
\end{tikzcd}
\]

The condition that $X$ be $W$-cocartesian is expressed by $\lambda \in W$, or again by $U \in \Asph_W$ (for $U$ identifies with the homotopy cofiber $Z(\lambda)$ of $\lambda$ \dots\,), likewise $Y$ $W$-cocartesian is expressed by $U' \in \Asph_W$, and $Z$ $W$-cocartesian by $U'' \in \Asph_W$. Now $U'' \sim Z(\mu)$, where $\mu : U \longrightarrow U'$ is in $W$ if $U$, $U'$ are $W$-aspheric, hence $U'' = Z(\mu)$ is then $W$-aspheric. This proves the lemma.

This shows that, assuming W(7 bis) (in fact, we have used only the ``universal'' part of W(7 bis)), namely that in
\begin{tikzcd}[ampersand replacement=\&, column sep=small, row sep=small, baseline=(\tikzcdmatrixname-2-2.base)]
T_0 \arrow[rr] \arrow[dr, "i"'] \& \& T_2 \arrow[dr, "i'"] \& \\
 \& T_1 \arrow[rr] \& \& T_3
\end{tikzcd}
a cocartesian diagram of open immersions, $i \in W \Longrightarrow i' \in W$ (instead of $\Longleftrightarrow$),

\msp{page 270}

the validity of W(7) reduces to the case where the square $Y$ is the trivial square
\begin{tikzcd}[ampersand replacement=\&, column sep=small, row sep=small, baseline=(\tikzcdmatrixname-2-1.base)]
 \& e \arrow[dr] \& \\
e \arrow[ur] \arrow[dr] \& \& e, \\
 \& e \arrow[ur] \&
\end{tikzcd}
$e$ the final object of Cat.

We can now conclude: We have (assuming the forward part of W(7 bis))
\[
\mathrm{a}' \ \& \ \mathrm{b}' \Longrightarrow \mathrm{c}',
\]
for if $Z_0$, $Z_1$ are $\in \Asph_W$, then $i_1 \in W$, hence $j_1 \in W$, and if in addition $Z_2 \in \Asph_W$, we therefore have $Z_3 \in \Asph_W$.
\[
\begin{tikzcd}[ampersand replacement=\&, column sep=small, row sep=small]
 \& Z_1 \arrow[dr, "j_2"] \& \\
Z_0 \arrow[ur, "i_1"] \arrow[dr, "i_2"'] \& \& Z_3 \\
 \& Z_2 \arrow[ur, "j_1"'] \&
\end{tikzcd}
\]

If we also have the reverse (``unusual'') part of W(7 bis), then
\[
\mathrm{b}' \ \& \ \mathrm{c}' \Longrightarrow \mathrm{a}',
\]
for if $Z_1$ and $Z_3$ are $W$-aspheric, then $j_2 \in W$, hence $i_2 \in W$, and if in addition $Z_2$ is $W$-aspheric, we conclude $Z_0$ $W$-aspheric.

\printedpage{176}
% head: A. Grothendieck : Dérivateurs, Ch. XII, Édition des Universités de Montpellier II et Paris VII
% foot: Version 30 mars 2003

Finally, suppose that $W$ is \textit{stable under direct factors,} and W(7 bis) satisfied in forward form. I say that then
\[
\mathrm{c}' \ \& \ \mathrm{a}' \Longrightarrow \mathrm{b}',
\]
i.e.
\[
Z_0 \text{ and } Z_3 \ \in \Asph_W \qquad\quad \Longrightarrow \qquad\quad Z_1 \text{ and } Z_2 \ \in \Asph_W.
\]
Factoring $i_1$ and $i_2$ as
\[
\begin{aligned}
Z_0 &\xrightarrow{i'_1} Z'_1 \xrightarrow{\alpha_1} Z_1, \\
Z_0 &\xrightarrow{i'_2} Z'_2 \xrightarrow{\alpha_2} Z_2,
\end{aligned}
\]
with $i'_1$, $i'_2$ open Dwyer immersions, and $\alpha_1, \alpha_2 \in W$, we are reduced to the case where

\msp{page 271}

$i_1$, $i_2$ are \textit{open Dwyer immersions,} and where the square is cocartesian (which implies that it is $W$-cocartesian). Then making the cobase change by $Z_0 \xrightarrow{p_0} Z'_0 = e$, we find a homomorphism of cocartesian squares
\[
\begin{tikzcd}[ampersand replacement=\&, column sep=small, row sep=small]
 \& Z_1 \arrow[rr] \arrow[dd, "p_1" near end] \& \& Z_3 \arrow[dd, "p_3"] \\
Z_0 \arrow[ur] \arrow[rr, crossing over] \arrow[dd, "p_0"'] \& \& Z_2 \arrow[ur] \& \\
 \& Z'_1 \arrow[rr] \& \& Z'_3, \\
e = Z'_0 \arrow[ur] \arrow[rr] \& \& Z'_2 \arrow[ur] \arrow[from=uu, crossing over, "p_2"' near end] \&
\end{tikzcd}
\]
where all the faces of the cube are $W$-cocartesian, and since $p_0 \in W$ ($Z_0$ being $W$-aspheric), the other vertical arrows $p_1$, $p_2$, $p_3$ are \texttt{[in]} $W$, hence the square $Z'_*$ also satisfies $Z'_0$ and $Z'_3$ \texttt{[}$W$\texttt{-]}aspheric (and in addition $Z'_0 = e$), and we must prove $Z'_1$, $Z'_2$ $W$-aspheric (whence likewise for $Z_1$, $Z_2$). Thus, we are reduced to the particular case where $Z_0 = e$.

But in this case there exists a retraction (trivial!) of $Z_1$ onto $Z_0$, hence also of $Z_3$ onto $Z_2$. Thus, $Z_2$ is a direct

\msp{page 272}

factor of a $W$-aspheric object, which implies, by the hypothesis on $W$, that $Z_2$ is $W$-aspheric. By symmetry, $Z_1$ is too. We have won!

\textbf{NB} In fact, these arguments are general, and give this:

\setcounter{footnote}{136}
\textbf{Theorem.} \textit{Let $(\mathcal{M}, W)$ be a model category, admitting a theory of the $W$-homotopy cofiber \textup{(}\footnote{But we only know how to verify the existence of such a theory in $(\mathrm{Cat}, W)$ if we accept the weak form of W(7), i.e. $f_0, f_1, f_2 \in W \implies f_3 \in W$, or some special case of it. But \textit{Brown's lemma} saves the situation, cf. XVI \S3.}\textup{)} (i.e. we have a factorization of every arrow $f : X \longrightarrow Y$ as $pi$, with $i \in \Cof_W$, $p \in W$.) We therefore have a notion of $W$-cocartesian square in $\mathcal{M}$ \textup{(}\footnote{to check that the factorization suffices, even if not all objects are fibrant over $e$. But we can then replace $\mathcal{M}$ by $\mathcal{M}_f$, and develop this formalism of homotopy fibers accordingly \dots}\textup{)}. Let
\printedpage{177}
% head: A. Grothendieck : Dérivateurs, Ch. XII, Édition des Universités de Montpellier II et Paris VII
% foot: Version 30 mars 2003
then $f : X \longrightarrow Y$ be a homomorphism from a $W$-cocartesian square $X$ to another $Y$, $f = (f_0, f_1, f_2, f_3)$.}
\DERmark{ENV0846}

\[
\begin{tikzcd}[ampersand replacement=\&, column sep=small, row sep=small]
 \& X_2 \arrow[rr] \arrow[dd, "f_2" near end] \& \& X_3 \arrow[dd, "f_3"] \\
X_0 \arrow[ur] \arrow[rr, crossing over] \arrow[dd, "f_0"'] \& \& X_1 \arrow[ur] \& \\
 \& Y_2 \arrow[rr] \& \& Y_3 \\
Y_0 \arrow[ur] \arrow[rr] \& \& Y_1 \arrow[ur] \arrow[from=uu, crossing over, "f_1"' near end] \&
\end{tikzcd}
\]

\textit{Let us consider the conditions on $f$}

\begin{enumerate}
\item[\textit{a)}] \textit{$f_0 \in W$.}
\item[\textit{b)}] \textit{$f_1$ and $f_2 \in W$.}
\item[\textit{c)}] \textit{$f_3 \in W$.}
\end{enumerate}

\textit{Then}

\begin{enumerate}
\item[\textit{1°)}] \textit{a \& b $\Longrightarrow$ c.}

\msp{page 273}

\item[\textit{2°)}] \textit{Suppose that $\mathcal{M}$ has a final object $e$. Then the implication}
\[
\textit{b \& c} \Longrightarrow \textit{a}
\]
\textit{for}~every \textit{homomorphism $f : X \longrightarrow Y$ of $W$-cocartesian squares is equivalent to the following axiom on $W$:}

Axiom of conservative $W$-cocartesian squares (\footnote{cf. argument page 78 which shows that the implication b \& c $\Longrightarrow$ a entails the conservative axiom.}): \textit{If}
\[
\begin{tikzcd}[ampersand replacement=\&, column sep=small, row sep=small]
 \& Z_1 \arrow[dr, "j_2"] \& \\
Z_0 \arrow[ur, "i_1"] \arrow[dr, "i_2"'] \& \& Z_3 \\
 \& Z_2 \arrow[ur, "j_1"'] \&
\end{tikzcd}
\]
\textit{is $W$-cocartesian, then $j_1 \in W \Longrightarrow i_1 \in W$ (hence in fact $i_1 \in W \Longleftrightarrow j_1 \in W$ and $i_2 \in W \Longleftrightarrow j_2 \in W$).}

\textit{This axiom is equivalent to the conjunction of the following two axioms.}

\textbf{Cons 1.} \textit{If we have a cocartesian square}
\printedpage{178}
% head: A. Grothendieck : Dérivateurs, Ch. XII, Édition des Universités de Montpellier II et Paris VII
% foot: Version 30 mars 2003
\[
\begin{tikzcd}[ampersand replacement=\&, column sep=small, row sep=small]
S \arrow[rr, "f"] \arrow[dr] \& \& T \arrow[dr] \& \\
 \& e \arrow[rr] \& \& U,
\end{tikzcd}
\]
\DERmark{ENV0847}

\textit{with $f \in \Cof_W$, then $U \in \Asph_W \Longrightarrow f \in W$ (hence $U \in \Asph_W \Longleftrightarrow f \in W$).}

\textbf{Cons 2.} \textit{If we have a cocartesian square,}
\DERmark{ENV0848}

\[
\begin{tikzcd}[ampersand replacement=\&, column sep=small, row sep=small]
 \& Z_1 \arrow[dr] \& \\
Z_0 \arrow[ur, "i_1"] \arrow[dr, "i_2"'] \& \& Z_3, \\
 \& Z_2 \arrow[ur] \&
\end{tikzcd}
\]
\textit{with $i_1, i_2 \in \Cof_W$, then: If $Z_1$, $Z_2$, $Z_3$ are $W$-aspheric, $Z_0$ is too.}

\msp{page 274}

\item[\textit{3°)}] \textit{Suppose that every direct factor of a $W$-aspheric object is $W$-aspheric, which is the case in particular if $W$ itself is stable under direct factors, a fortiori if $W$ is strongly saturated. Suppose in addition that $\mathcal{M}$ has} \texttt{[a]} \textit{final object. Then we have (under the general conditions at the beginning)}
\[
\textit{c \& a} \Longrightarrow \textit{b,}
\]
\textit{i.e. if $f_0$, $f_3$ are} \texttt{[in]} \textit{$W$, likewise $f_1$, $f_2$.}
\end{enumerate}

This last part can be unfolded as follows:

\textbf{Corollary 1.} \textit{Under the additional hypothesis} \textup{3\,°)} \textit{on $W$, the validity of c \& a $\Longrightarrow$ b (for every arrow $f$ between $W$-cocartesian squares) is equivalent to each of the following two other axioms on $W$ (which are therefore equivalent):}

\begin{enumerate}
\item[\textit{1)}] \textit{If we have a $W$-cocartesian square}
\[
\begin{tikzcd}[ampersand replacement=\&, column sep=small, row sep=small]
 \& Z_1 \arrow[dr] \& \\
Z_0 \arrow[ur] \arrow[dr] \& \& Z_3, \\
 \& Z_2 \arrow[ur] \&
\end{tikzcd}
\]
\textit{then $Z_0$ and $Z_3$ $W$-aspheric $\Longrightarrow$ $Z_1$ and $Z_2$ $W$-aspheric.}
\item[\textit{1$'$)}] \textit{As in 1), with $Z_0 = e$ (in the case where we suppose that $\mathcal{M}$ admits} \texttt{[a]} \textit{final object $e$).}
\end{enumerate}

\msp{page 275}

\textbf{Corollary 2.} \textit{Suppose that $W$ satisfies the axiom of conservative $W$-cocartesian squares, that every direct factor of a $W$-aspheric object is $W$-aspheric (e.g. $W$ stable under direct factors), finally that $\mathcal{M}$ admits a final object. Then for every homomorphism $f : X \longrightarrow Y$ of $W$-cocartesian squares, the conjunction of two of the conditions a) b) c) on p. 272 implies the third.} % [sic] pout -> pour, typographical error


\endgroup
% END UNIT Der-ChXIIc-30-03-03/Der-ChXIIc-30-03-03_p061-064.tex
\clearpage\setcounter{footnote}{0}\phantomsection\addcontentsline{toc}{section}{Chapter XII, appendix}\markboth{Chapter XII, appendix}{Chapter XII, appendix}
% BEGIN UNIT Der-ChXIIapp-19-06-03/Der-ChXIIapp-19-06-03_p001-005.tex
\begingroup\makeatletter%
\newcommand{\printedpage}[1]{\marginpar{\footnotesize p.~#1}}%
\newcommand{\msp}[1]{\marginpar{\footnotesize [#1]}}%
\newcommand{\Hom}{\mathrm{Hom}}%
\newcommand{\Fl}{\mathrm{Fl}}%
\newcommand{\Ob}{\mathrm{Ob}}%
\newcommand{\Cat}{\mathrm{Cat}}%
\newcommand{\Ens}{\mathrm{Ens}}%
\newcommand{\pr}{\mathrm{pr}}%
\newcommand{\Ker}{\mathrm{Ker}}%
\newcommand{\Coker}{\mathrm{Coker}}%
\newcommand{\mon}{\mathrm{mon}}%
\newcommand{\deq}{\mathrel{\overset{\mathrm{def}}{=}}}%
\makeatother\let\printedpage\derpp\let\msp\dermsp%
\let\mspin\DERmspin%

\DERmark{ENV0695}


\printedpage{}
% no page number printed (title page)
% head: (none)

\begin{center}
{\LARGE\textsf{THE DERIVATORS}}

\bigskip

\textbf{ALEXANDRE GROTHENDIECK}

\bigskip\bigskip

{\large\textbf{Appendix to Chapter XII}}
\end{center}

\par

\printedpage{1}
% head: A. Grothendieck : Dérivateurs, Ch. XII, Édition des Universités de Montpellier II et Paris VII
% foot: Version 20 juin 2003

\msp{page 1}
\textbf{Proposition} (?) (\footnote{Cf. pages 11 ff. for the key counterexamples.}). \textit{Let $X$ be in} $\Cat$\textit{, with $b$, $c$ in $X$ such that}

\begin{itemize}
\item[\textit{a)}] \textit{$b\times c$ is in $X$.}
\item[\textit{b)}] \textit{$\forall$ $x$ in $X$, $b(x)$ ($\deq$} $\Hom_X(x,b)$\textit{) and $c(x)$ are $\neq\emptyset$, i.e. $X/b \longrightarrow X$ and $X/c \longrightarrow X$ are surjective on objects.}
\end{itemize}

\textit{Then $X$ is $\mathbf{W}_\infty$-aspheric.} (\footnote{False, cf. below. $X$ is 1-connected, but not necessarily 2-connected.})

The fact that $b$ is a greatest element of $\Ob X$ already ensures 0-connectedness.

For 1-connectedness, let us prove that in the fundamental groupoid $\Pi(X)$, $a \deq b\times c$ becomes a final element.
\[
\begin{tikzcd}
b\times c \arrow[r] \arrow[d] & b \\
c &
\end{tikzcd}
\]

For this, let us first prove that $\forall$ $x$ in $X$, all the elements of $a(x) = b(x)\times c(x)$ are equal in $\Fl\Pi(X)$. Indeed, since $b\times c \longrightarrow b$ is \texttt{[an]} isomorphism, hence \texttt{[a]} monomorphism in $\Pi(X)$, this already shows that if $\beta,\beta'\in b(x)$, $\gamma,\gamma'\in c(x)$, whence $(\beta,\gamma)$, $(\beta,\gamma')$, $(\beta',\gamma')\in b(x)\times c(x)$, we have $(\beta,\gamma) = (\beta,\gamma')$ in $\Pi(X)$ (since $\pr_1(\beta,\gamma) = \pr_1(\beta,\gamma')$, and likewise (using that $\pr_2$ is \texttt{[an]} isomorphism, hence \texttt{[a]} monomorphism) $(\beta,\gamma') = (\beta',\gamma')$). Thus all the $(\beta,\gamma)\in b(x)\times c(x) = a(x)$ are equal in $\Fl\Pi(X)$. To finish proving that $a$ is a final object of $\Pi(X)$, it suffices to see that every arrow \textit{in} $\Pi(X)$ from an $x\in\Ob X$ to $a$, comes from an arrow $x\longrightarrow a$ in $X$. But an arrow in $\Pi(X)$ is obtained by a path in $X$
\[
\begin{tikzcd}
x = x_0 \arrow[r, "u_1"] \arrow[dr, "\alpha_0"] & x_1 \arrow[d, "\alpha_1"'] & x_2 \arrow[l, "u_2"'] \arrow[r] \arrow[dl, "\alpha_2"'] & {} & \cdots & {} & x_{2N} = a \arrow[l, "u_{2N}"'] \\
& a. & & & & &
\end{tikzcd}
\]

Let us take an arrow $\alpha_0 : x\longrightarrow a$ in $X$, I say that the % [sic] a_0 -> a, the arrows α_i go to a (no a_0 is defined)
\DERmark{ENV0696}

\msp{page 2}
homotopy class of the path $c$ (i.e. its image in $\Fl(\Pi(X))$) is equal to that of $\alpha_0$. % [sic] u_0 -> \alpha_0 (?), no u_0 is defined; also the parenthesis opened at "(i.e." is not closed
\DERmark{ENV0697}\DERmark{ENV0698}

To see this, let us choose for every $i\in[0,2N]$ an arrow $\alpha_i : x_i\longrightarrow a$. (It exists thanks to hypothesis b)). I say that for each $i\in[0,2N]$, the triangle
\printedpage{2}
% head: A. Grothendieck : Dérivateurs, Ch. XII, Édition des Universités de Montpellier II et Paris VII
% foot: Version 20 juin 2003
\[
\begin{tikzcd}[column sep=small]
x_0 \arrow[rr, "c_i"] \arrow[dr, "\alpha_0"] & & x_i \arrow[dl, "\alpha_i"'] \\
& a &
\end{tikzcd}
\]
in $\Pi(X)$ is commutative, where $c_i$ denotes the homotopy class of the path from $x_0$ to $x_i$ defined by the $u_j$ ($j\in[1,i]$) \texttt{[rather} ($i\in[1,2N]$)\texttt{]}. We do this by induction on $i$, the assertion is trivial if $i = 0$ ($c_0$ represents the identity map). The diagram % [sic] identitique -> identique, non-existent word
\DERmark{ENV0699}

\[
\begin{tikzcd}[column sep=small]
x_0 \arrow[rr, "c_i"] \arrow[dr, "\alpha_0"] \arrow[rrrr, bend left=32, "c_{i+1}"] & & x_i \arrow[dl, "\alpha_i"'] \arrow[rr, "u_{i+1}"] & & x_{i+1} \arrow[dlll, "\alpha_{i+1}"'] \\ % verified: solid curved arrow; bend increased for label clearance
& a & & &
\end{tikzcd}
\]
shows that it suffices to prove that for $i\in[0,2N-1]$ the triangle
\[
\begin{tikzcd}[column sep=small]
x_i \arrow[rr, "u_{i+1}"] \arrow[dr, "\alpha_i"'] & & x_{i+1} \arrow[dl, "\alpha_{i+1}"] \\
& a &
\end{tikzcd}
\]
is commutative in $\Pi(X)$. (\textbf{NB} $u_{i+1}$ is in the direct or retrograde direction, according as % [sic] directe -> direct, agreement with "sens"
$i$ is even or odd.) But this follows from the fact that the image of $\Hom_X(x_i,a)$ (resp. $\Hom_X(x_{i+1},a)$) % [sic] paire ou impaire -> pair ou impair, agreement with "i"
\msp{page 3}
in $\Fl\Pi(X)$ is reduced to a single element.

Thus we find that $X$ \textit{is} 1-\textit{connected.} So there is hope! It remains to prove
\[
H^i(X,M) = 0, \qquad i\geq 2, \quad M \text{ constant coefficient.}
\]

I have looked at the case when we have this: $b = c$, $\Ob X = \{b, b\times b\}$, $\Hom_X(b\times b,b)$ \texttt{[rather} $\Hom(b\times b, b)$\texttt{]} $= \{$ disjoint union of the sets $u\circ\pr_1$ and $u\circ\pr_2$, with $u\in\Hom_X(b,b) = G$ an arbitrary monoid $\}$. Then $X^\circ$ is fibered, via $I\longmapsto b^I$, over the full subcategory $S$ of $\Ens$ consisting of $\{0\}$ and $\{0,1\}$, and thus $X$ is cofibered (in particular proper) over $S^\circ \deq T$. (NB $S^\circ = T$ has an initial object, but not a final object. Its cohomology with values in an arbitrary presheaf is not always trivial, cf. below.)
\DERmark{ENV0700}


We find \texttt{[``We find'' crossed out]}
\[
H^n(X,M) \Longleftarrow H^p(T, \underbrace{R^q f_*(M_X)}_{= \mathcal{H}^q}), \qquad \text{where } f : X \longrightarrow S^\circ
\]
\texttt{[}$n = p+q$\texttt{]}. I admit here (provisionally) this: % [sic] J'admet -> J'admets, verb ending
\[
\begin{aligned}
H^p(T,F) &= 0 \ , &&\text{for } p\geq 2, \text{ any } F\in T^\wedge \ ,\\
H^0(T,F) &= \Ker\delta \ , &&\\
H^1(T,F) &= \Coker\delta \ , && % [sic] H^2(T,F) -> H^1(T,F), the first line gives H^p(T,F) = 0 for p >= 2
\end{aligned}
\]
\DERmark{ENV0701}

\printedpage{3}
% head: A. Grothendieck : Dérivateurs, Ch. XII, Édition des Universités de Montpellier II et Paris VII
% foot: Version 20 juin 2003
where $\delta : F_0 \longrightarrow F_1$ is defined as the difference $\delta = p_1^* - p_2^*$, where $p_1$, $p_2$ are the two projections of $\{0,1\}\simeq\{0\}^2$ in $T = S^\circ$. Moreover we have
\[
(\mathcal{H}^q)_0 = H^q(G,M), \qquad (\mathcal{H}^q)_1 = H^q(G\times G, M),
\]
and
\[
p_1^*,\ p_2^*\ :\ H^q(G,M) \rightrightarrows H^q(G{\times}G, M) \qquad \text{are the transposes of } \pr_1,\ \pr_2 : G\times G\longrightarrow G.
\]

\msp{page 4}
We readily check that
\[
H^0(T,\mathcal{H}^q) = 0 \qquad \text{if } q\geq 1,
\]
and there remains
\[
H^n(X,M) \simeq H^1(T,\mathcal{H}^{n-1}) = \Coker(\delta : H^{n-1}(G) \longrightarrow H^{n-1}(G\times G)). % [sic] H^1(S, ...) -> H^1(T, ...), the cohomology is taken on T = S°
\]
\DERmark{ENV0702}

This gives, for $n = 2$,
\[
\begin{aligned}
H^2(X,M) &\simeq \Coker\Big(\delta : \Hom_\mon(G,M) \longrightarrow \underbrace{\mathit{Hom}_\mon(G\times G, M)}_{\Hom_\mon(G,M)\times\Hom_\mon(G,M)} \ \Big)\\
&\simeq \Coker\Big(u\longmapsto (u,-u) : H \longrightarrow H\times H\Big),
\end{aligned}
\]
where $H$ is the abelian group $\Hom_\mon(G,M)$.

Whence
\[
H^2(X,M) \ \simeq\ H = \Hom_\mon(G,M).
\]
Thus
\begin{quote}
$X$ is 2-connected if and only if the enveloping \textit{abelian} group of the monoid $G$ is zero.
\end{quote}

Thus, if $G$ is an abelian group $\neq 0$, $X$ \textit{is not} 2-connected.

\textbf{Corollary.} $\Psi = {}$\begin{tikzcd}[ampersand replacement=\&, row sep=tiny, column sep=small, baseline=(A.base)]
\& b_0 \\
|[alias=A]| a_0 \arrow[ur] \arrow[dr] \& \\
\& c_0
\end{tikzcd}
\textit{is} not \textit{totally} $\mathbf{W}_\infty$\textit{-aspheric; i.e. if we have a diagram} \begin{tikzcd}[ampersand replacement=\&, row sep=tiny, column sep=small, baseline=(B.base)]
\& b \\
|[alias=B]| a \arrow[ur] \arrow[dr] \& \\
\& c
\end{tikzcd}
\textit{in a category} $X\in\Ob\Cat$\textit{, such that its amalgamated sum in $X^\wedge$ is equal to $e_X$, it does not follow that $X$ is $\mathbf{W}_\infty$-aspheric.} (But perhaps will be willing to be 1-connected?)

\printedpage{4}
% head: A. Grothendieck : Dérivateurs, Ch. XII, Édition des Universités de Montpellier II et Paris VII
% foot: Version 20 juin 2003

\msp{page 5}
\textbf{Question.} \textit{Are there other totally} $\mathbf{W}_\infty$\textit{-aspheric objects of} $\Cat$\textit{, than the categories having a final object?} I now suspect not -- hence that one cannot improve substantially on the already known notion of an open Dwyer immersion.

In the diagram
\[
\begin{tikzcd}
U \arrow[r, hook, "i"] \arrow[d, "p"'] \arrow[dr, phantom, "\substack{\text{cocart.}\\ \text{and cart.}}"] & X \arrow[d, "q"'] & T \arrow[l, hook', dotted, "j"'] \arrow[d, dotted, "\text{iso}"] \\
U' \arrow[r, hook, "i'"] & X' & T', \arrow[l, hook', dotted, "j'"']
\end{tikzcd}
\tag*{$(*)$}
\]
let us take
\[
\begin{aligned}
U = \Psi, \quad & T = \{t\} = e \text{ final category,}\\
& t \text{ final object of } X,\\
& \text{i.e. } H_t = (u\longmapsto \Hom_X(u,t)) = e_U \in \Ob U^\wedge \ \mbox{(\footnotemark)}.
\end{aligned}
\]
\footnotetext{i.e. $X$ obtained from $U = \Psi$ by adding to it a final object $t$ -- thus $X\simeq\boldsymbol{\Delta}^1\times\boldsymbol{\Delta}^1$,
\[
\begin{tikzcd}[ampersand replacement=\&,
execute at end picture={
\draw[rounded corners] ([xshift=-10pt,yshift=10pt]A.center) -- ([xshift=-10pt,yshift=-10pt]C.center) -- ([xshift=10pt,yshift=-10pt]D.center) -- ([xshift=10pt,yshift=3pt]D.center) -- ([xshift=3pt,yshift=3pt]C.center) -- ([xshift=3pt,yshift=10pt]A.center) -- cycle;
\node (P) at ([xshift=-35pt,yshift=-22pt]C.center) {$\Psi$};
\draw[-stealth] (P) -- ([xshift=-11pt,yshift=-4pt]C.center);
}]
|[alias=A]| {} \arrow[r] \& {} \\
|[alias=C]| {} \arrow[u] \arrow[r] \& |[alias=D]| {} \arrow[u]
\end{tikzcd}
\ .
\]}
Let us take for $U'$ the category spoken of in the hypothetical (and false) ``proposition'' p. 1 with $b = c$, more precisely a category with two objects $b$ and $b\times b$, a group (or even a monoid) $G$ of automorphisms \texttt{[(or even of endomorphisms)]} of $b$, and the other arrows those which follow from these, requiring that $b\times b$ really be the cartesian square of $b$. We find a cofibered category (hence proper) over the category $T$ opposite to that of finite sets of the form $\{0\}\cup\{0,1\}$ (two objects) with fibers $B_G$ and $B_{G\times G}$ respectively ($B_\pi =$ category with a single object, and endomorphism monoid $\pi$). The functor
\[
U = \Psi \longrightarrow U'
\]
is the one given by the diagram

\msp{page 6}
\[
\begin{tikzcd}[row sep=small]
& b \\
b\times b \arrow[ur, "\pr_1"] \arrow[dr, "\pr_2"'] & \\
& b
\end{tikzcd}
\]


\endgroup
% END UNIT Der-ChXIIapp-19-06-03/Der-ChXIIapp-19-06-03_p001-005.tex
% BEGIN UNIT Der-ChXIIapp-19-06-03/Der-ChXIIapp-19-06-03_p006-010.tex
\begingroup\makeatletter%
\tikzcdset{every diagram/.append style={baseline={([yshift=-.5ex]current bounding box.center)}}}%
\newcommand{\printedpage}[1]{\marginpar{\footnotesize p.~#1}}%
\newcommand{\msp}[1]{\marginpar{\footnotesize [#1]}}%
\newcommand{\Hom}{\mathrm{Hom}}%
\newcommand{\id}{\mathrm{id}}%
\newcommand{\Winf}{\mathbf{W}_{\infty}}%
\makeatother\let\printedpage\derpp\let\msp\dermsp%
\let\mspin\DERmspin%


\printedpage{5}
% head: A. Grothendieck : Dérivateurs, Ch. XII, Édition des Universités de Montpellier II et Paris VII
% foot: Version 20 juin 2003

in $U'$, and hypotheses a), b) of page 1 ensure that the amalgamated sum $\varinjlim$ of this diagram in $U'^{\wedge}$ is $e_{U'}$, i.e. that
\[ p_!(e_U) = e_{U'}. \]
It follows that in the amalgamated sum, $X'$ considered in $(*)$, the unique element $t'$ of $T'$ is a final object of $X'$. We have seen that $U'$ is 1-connected, but not 2-connected if the enveloping abelian group of $G$ (i.e. $G/$(commutator subgroup) in the case where $G$ is a group) is not reduced to 1. The most economical example is obtained by taking $G = \{\pm 1\}$ the group with two elements. This is thus a case where in $(*)$ we have
\[ i \in W \ \ \text{(and even $i$ aspheric),} \qquad \text{but } i' \notin W, \]
and which has the advantage over the more sophisticated example of page 103 (prop. 10) of involving only finite categories -- and even very small ones!

Let us now proceed as in proposition 11, denoting by $U''$ the category obtained from $U'$ by adding to it a final
\msp{page 7}
object, and let
\[ p' : U' \longrightarrow U'' \]
be the inclusion, so that
\[
\begin{aligned}
p'_!(e_{U'}) \quad &= \quad (p'p)_!(e_U) = U' \quad && \text{as an object in } U''^{\wedge}\\
& && \text{(subpresheaf of the final presheaf } e_{U''}\text{).}
\end{aligned}
\]
Consider
\[
\begin{tikzcd}
U \arrow[r, hook, "i"] \arrow[d, "{p''=p'p}"'] \arrow[dr, phantom, "\begin{smallmatrix}\text{cocart.}\\ \text{and cart.}\end{smallmatrix}"] & X \arrow[d, "q''"'] & T \arrow[l, hook', dotted, "j"'] \arrow[d, dotted, "\text{iso}"] \\
U'' \arrow[r, hook, "i''"] & X'' & T''. \arrow[l, hook', dotted, "j''"']
\end{tikzcd}
\]
This time $U$, $X$ \textit{and} $U''$ are $\Winf$-aspheric (and even all three contractible), hence $i$, $p''$ are in $W$, yet $X''$ is not $\Winf$-aspheric (hence $i'', q'' \notin \Winf$), for in order for it to be so, i.e. for $i'' \in \Winf$, it would be necessary (prop. 7, cor. page 85 (\footnote[4]{not yet proved!})) that $i''$ be $\Winf$-aspheric, i.e. that $U''/t''$ be $\Winf$-aspheric. But by construction $U''/t'' \simeq U'$, which is \textit{not} $\Winf$-aspheric.

Here is an even simpler example of a functor $f : \Psi \longrightarrow X$ such that $f_!(e_{\Psi})$ is not $\Winf$-aspheric, i.e. of a diagram
\[
\begin{tikzcd}
& b \\
a \arrow[ur, "\beta_0"] \arrow[dr, "\gamma_0"'] & \\
& c
\end{tikzcd}
\]
\printedpage{6}
% head: A. Grothendieck : Dérivateurs, Ch. XII, Édition des Universités de Montpellier II et Paris VII
% foot: Version 20 juin 2003
\msp{page 8}
in a category $V$, such that the direct limit $L = b \amalg_a c$ \textit{in} $V^{\wedge}$ is a non-aspheric object of $V^{\wedge}$, i.e. $V/L$ is non-aspheric. I know that, for such a counterexample, $a \longrightarrow b$ and $a \longrightarrow c$ must not be monomorphisms. The idea is to take a universal example where neither of these arrows is \texttt{[a]} monomorphism, by taking
\msp{page 9}
\[
V = \left(
\begin{tikzcd}
& & b \\
0 \arrow[r, shift left, "\lambda"] \arrow[r, shift right, "\mu"'] & a \arrow[ur, "\beta"] \arrow[dr, "\gamma"'] & \\
& & c
\end{tikzcd}
\right)
\]
with
\[ \beta\lambda = \beta\mu \qquad \text{and} \qquad \gamma\lambda = \gamma\mu \]
the only two nonidentity arrows, besides the four arrows marked $\beta$, $\gamma$, $\lambda$, $\mu$.

We find that
\[ b \amalg_a c \xrightarrow{\ \sim\ } e_V \qquad \text{(amalgamated sum in } V^{\wedge}\text{).} \] % [sic] e_X -> e_V, the category considered here is V
\DERmark{ENV0703}


Yet $V$ is not aspheric, nor even 1-connected: it is homotopic to the open subcategory $V_0 = (0 \overset{\lambda}{\underset{\mu}{\rightrightarrows}} a)$, for we have a retraction $V \xrightarrow{\ r\ } V_0$ (sending $b$ and $c$ to $a$), and an arrow $ir \longrightarrow \id_V$ (where $i : V_0\hookrightarrow V$ is the inclusion). Thus $V$ has the same homotopy type \textit{as} $V_0$, which is a circle. So it is in our interest to take this $V$, where $\tilde{V}$ obtained from $V$ by adding a final element $e$, % [sic] i : X_0 \hookrightarrow X -> i : V_0 \hookrightarrow V, the inclusion of V_0 in V is meant
\DERmark{ENV0704}

\[
\tilde{V} = \left(
\begin{tikzcd}
& & b \arrow[dr] & \\
0 \arrow[r, shift left] \arrow[r, shift right] & a \arrow[ur] \arrow[dr] & & e \\
& & c \arrow[ur] &
\end{tikzcd}
\right),
\]
as $U'$, for the counterexamples on pages 6 and 7.

\msp{page 10}
\subsection*{2.6. Counterexamples for total asphericity.}

I have finally convinced myself that
\[
\Psi = \left\{
\begin{tikzcd}
& b_0 \\
a_0 \arrow[ur] \arrow[dr] & \\
& c_0
\end{tikzcd}
\right\}
\]
\printedpage{7}
% head: A. Grothendieck : Dérivateurs, Ch. XII, Édition des Universités de Montpellier II et Paris VII
% foot: Version 20 juin 2003
is \textit{not} totally aspheric (for the localizer $\Winf$ -- in the present number, we shall assume $W = \Winf$), or what amounts to the same thing, that we can have a diagram of type $\Psi$
\[
\begin{tikzcd}
& b \\
a \arrow[ur, "\beta"] \arrow[dr, "\gamma"'] & \\
& c
\end{tikzcd}
\tag*{$(*)$}
\]
in a category $Z \in \mathrm{Ob}\,\mathrm{Cat}$, whose direct limit $b \amalg_a c$ in $Z^{\wedge}$ is the final presheaf $e_Z$, i.e. such that we have
\[ b(z) \amalg_{a(z)} c(z) \xrightarrow{\ \sim\ } \text{one-point set} \]
for every object $z$ of $Z$ (where $a(z) = \Hom_Z(z,a)$ \texttt{[rather} $\Hom_Z(z,a)$\texttt{]}) without $Z$ being aspheric for all that, nor even 1-connected (\footnote[5]{false -- it is necessarily 1-connected (\footnotemark[6]), but not 2-connected.}).\footnotetext[6]{cf. comment XII, p. 101.} The simplest example I have found, is the ``universal'' example of a diagram $(*)$ where neither one nor the other of the two arrows $a \longrightarrow b$, $a \longrightarrow c$ is \texttt{[a]} monomorphism (for if it is, then we can prove that $Z$ is indeed aspheric, as we shall see later on). It is the category
\DERmark{ENV0705}

\[
Z = \left\{
\begin{tikzcd}
& & b_0 \\
0_0 \arrow[r, shift left, "\lambda_0"] \arrow[r, shift right, "\mu_0"'] & a_0 \arrow[ur, "\beta_0"] \arrow[dr, "\gamma_0"'] & \\
& & c_0
\end{tikzcd}
\right\}
\]
\msp{page 11}
obtained from $\Psi$ by adding to it an object $0_0$, connected to $a_0$ by two distinct arrows $\lambda_0$, $\mu_0$, such that
\[ \beta_0\lambda_0 = \beta_0\mu_0 \qquad \text{and} \qquad \gamma_0\lambda_0 = \gamma_0\mu_0, \]
these two composite arrows being, together with the four arrows $\beta_0$, $\gamma_0$, $\lambda_0$, $\mu_0$ marked on the diagram, the only (six) nonidentity arrows of $Z$. $\Psi$ immersed in $Z$ as a closed subcategory, with complement the one-point category reduced to $0_0$. (NB The closed immersion $\Psi \longrightarrow Z$ is not a Dwyer immersion, i.e. (as $Z$ is the smallest open subcategory of $Z$ containing $\Psi$) is not coaspheric, for the category $0\backslash\Psi$, cofibered over $\Psi$ with discrete fibers, of cardinality 2 at $a_0$ and 1 at $b_0$ and $c_0$, is the category
\[
\begin{tikzcd}
& \tilde{b}_0 \\
a'_0 \arrow[ur] \arrow[ddr] & \\
a''_0 \arrow[uur, crossing over] \arrow[dr] & \\
& \tilde{c}_0
\end{tikzcd}
\]
\printedpage{8}
% head: A. Grothendieck : Dérivateurs, Ch. XII, Édition des Universités de Montpellier II et Paris VII
% foot: Version 20 juin 2003
i.e.
\[
\begin{tikzcd}
& {} & \\
{} \arrow[ur] \arrow[dr] & & {} \arrow[ul] \arrow[dl] \\
& {} &
\end{tikzcd}
\quad ,
\]
which has the homotopy type of the circle.) Nor the open immersion $\{0_0\} \longrightarrow Z$, for $\{0_0\}/a_0$ is the discrete category of cardinality 2.

The inclusion $i : Z_0 \longrightarrow Z$ of
\[ Z_0 \ = \ (\ 0_0 \overset{\lambda_0}{\underset{\mu_0}{\rightrightarrows}} a_0\ ) \]
into $Z$ is a homotopism (\footnote[7]{false, there is no retraction at all! In fact, $Z$ is the \textit{suspension} of the $\rightrightarrows$, which has the homotopy type (topological) of the circle, and it therefore has the homotopy type of the 2-sphere.}) (for there exists an obvious retraction $r$ of $Z$ onto $Z_0$ (\footnotemark[7]), and an arrow $ir \longrightarrow \id_Z$), hence $Z$ is homotopic to $Z_0$, which has the homotopy type of the circle, hence is not 1-connected.

\msp{page 12}
On the other hand, the verification of
\[ b_0 \amalg_{a_0} c_0 \ = \ e_Z \qquad \text{in } Z^{\wedge} \]
is immediate.

We are going to use this arrow
\[ \Psi \longrightarrow Z \]
to construct interesting (and heartrending) examples of amalgamated sums
\[
\begin{tikzcd}
U \arrow[r, hook, "i"] \arrow[d, "p"'] & X \arrow[d, "q"'] \\
U' \arrow[r, hook, "i'"] & X',
\end{tikzcd}
\]
with $i$ an open immersion, such that we even have
\[ i :\ U \longrightarrow \overline{U} \qquad \text{is aspheric, hence in } W, \]
whereas nevertheless we shall have
\[ i' :\ U' \longrightarrow \overline{U'} \qquad \notin W, \]
and what is better (or worse),
\[ p \in W \qquad \text{but} \qquad q \notin W \]
\printedpage{9}
% head: A. Grothendieck : Dérivateurs, Ch. XII, Édition des Universités de Montpellier II et Paris VII
% foot: Version 20 juin 2003
(and possibly, at the same time, $i \in W$ but $i' \notin W$). Moreover (unlike the constructions of propositions 10, 11 (p. 101-107)), the categories considered may all be finite. I am particularly interested in the case where $X$ is a ``square''
\[ X = \Phi_n \times \Phi_m \qquad (n,m \geq 2), \]
and where $U$ is its ``boundary''
\[ U = \dot{X} \overset{\mathrm{def}}{=} (\dot{\Phi}_n \times \Phi_m) \cup (\Phi_n \times \dot{\Phi}_m). \]
\DERmark{ENV0706}


\msp{page 13}
\[
\begin{tikzcd}
0 \arrow[r] \arrow[d] & a \arrow[d] & a' \arrow[l] \arrow[r] \arrow[d] & \cdots \\
b \arrow[r] & t_1 & \cdots \arrow[l] & \\
b' \arrow[u] \arrow[r] \arrow[d] & {} \arrow[u] & \cdots \arrow[l] \arrow[u] & \\
\cdots & & &
\end{tikzcd}
\]

\textbf{Proposition.} \textit{Let $n,m \geq 2$, and consider the inclusion $U = \partial X \overset{i}{\hookrightarrow} X$ of the boundary of the ``square'' $X = \Phi_n \times \Phi_m$ into $X$. It is a $W$-aspheric open immersion, to be precise, the $U/t$ ($t \in \dot{U} = \bar{U}\backslash U$) are isomorphic to $e$ (one-point category) or to $\Phi =$ \begin{tikzcd}[ampersand replacement=\&, row sep=small, column sep=small] \& {} \arrow[dl] \\ {} \& \\ \& {} \arrow[ul] \end{tikzcd} (category with a final element), except for the four elements $t_i$ of the four ``corners'' of $\dot{U}$, where it is the category}
\[
\tilde{\Psi} = \left(
\begin{tikzcd}
& & a' \arrow[dl] \\
& a & \\
0 \arrow[ur] \arrow[dr] & & \\
& b & \\
& & b' \arrow[ul]
\end{tikzcd}
\right)
\]
\textit{having no final element. Thus $U \xrightarrow{\ i\ } X$ is not an open Dwyer immersion (but $U \longrightarrow X\backslash\{t_1,t_2,t_3,t_4\}$ is one \ldots). This being so, I say that $i : U \longrightarrow X$ is not a $W$-cofibration.}


\endgroup
% END UNIT Der-ChXIIapp-19-06-03/Der-ChXIIapp-19-06-03_p006-010.tex
% BEGIN UNIT Der-ChXIIapp-19-06-03/Der-ChXIIapp-19-06-03_p011-014.tex
\begingroup\makeatletter%
\newcommand{\printedpage}[1]{\marginpar{\footnotesize p.~#1}}%
\newcommand{\msp}[1]{\marginpar{\footnotesize [#1]}}%
\newcommand{\dotto}{\mathrel{\tikz[baseline=-0.5ex]\draw[dotted,-latex] (0,0)--(1.6em,0);}}%
\makeatother\let\printedpage\derpp\let\msp\dermsp%
\let\mspin\DERmspin%


\printedpage{10}
% head: A. Grothendieck : Dérivateurs, Ch. XII, Édition des Universités de Montpellier II et Paris VII
% foot: Version 20 juin 2003
\textsc{Proof}. Let
\[ V \hookrightarrow U \]
\DERmark{ENV0707}

be the inclusion into $U = \partial X$ of the corner-category \texttt{[}$V =$\texttt{]} $\vcenter{\hbox{\begin{tikzcd}[row sep=small, column sep=small]
0 \arrow[r] \arrow[d] & a \\
b &
\end{tikzcd}}}$ (cf. figure), isomorphic to $\Psi$. One sees at once that this is a closed Dwyer immersion (the smallest open of $U$ containing $\Psi$ is just the category $\tilde{\Psi} = U/t_1$ from earlier), hence a $W$-cofibration. If $U \longrightarrow X$ were a $W$-cofibration, the same would hold for the composite $V \longrightarrow X$. We are going to see that this is not the case.

Let $\overline{V} = \vcenter{\hbox{\begin{tikzcd}[row sep=small, column sep=small]
0 \arrow[r] \arrow[d] & a \arrow[d] \\
b \arrow[r] & t_1
\end{tikzcd}}}$ be the closure of $V$ in $X$,
\msp{page 14}and consider the inclusion
\[ \overline{V} \longrightarrow X, \]
I say that \texttt{[it]} is a closed Dwyer immersion, i.e. that for every $s \in W\backslash\overline{V}$, where $W$ is the smallest open of $X$ containing $\overline{V}$, i.e.
\[
W = \quad
\begin{tikzcd}[row sep=small, column sep=small]
0 \arrow[r] \arrow[d] & a \arrow[d] & a' \arrow[l] \arrow[d] \\
b \arrow[r] & {} & a'' \arrow[l] \\
b' \arrow[u] \arrow[r] & b'' \arrow[u] & c, \arrow[l] \arrow[u]
\end{tikzcd}
\]
the category $s\backslash\overline{V}$ has an initial object. (The categories $s\backslash\overline{V}$ that one obtains, from the five possible choices for $s$, are $\boldsymbol{\Delta}^1$ (if $s = a'$ or $b'$) or $e$ (for $s = a''$, $b''$ or $c$)). Thus $\overline{V} \longrightarrow X$ is a $W$-cofibration.

This being established, one knows that there exists, for [?] the inclusion (open immersion)
\[ V \hookrightarrow \overline{V} \qquad \text{i.e. } \Psi \hookrightarrow \boldsymbol{\Delta}^1 \times \boldsymbol{\Delta}^1, \]
a functor $V \longrightarrow V'$, giving rise to a cocartesian square
\[
\begin{tikzcd}
V \arrow[r, hook, "i"] \arrow[d, "p"'] & \overline{V} \arrow[d, "\overline{p}"] \\
V' \arrow[r, hook, "i'"] & \overline{V'}
\end{tikzcd}
\tag*{$(*)$}
\]
with
\[ i, p \in W \ \text{(for $V$, $V'$, $\overline{V}$ $W$-aspheric), but } i', \overline{p} \notin W \ \text{(i.e. $\overline{V'}$ \textit{not} $W$-aspheric).} \]

I recall that one first takes
\printedpage{11}
% head: A. Grothendieck : Dérivateurs, Ch. XII, Édition des Universités de Montpellier II et Paris VII
% foot: Version 20 juin 2003
\[
V'_0 = \left(
\begin{tikzcd}[row sep=small]
 & & a \\
0' \arrow[r, shift left, "\lambda"] \arrow[r, shift right, "\mu"'] & 0 \arrow[ur, "\alpha"] \arrow[dr, "\beta"'] & \\
 & & b
\end{tikzcd}
\right)
\qquad \text{with} \quad
\left\{
\begin{aligned}
\alpha\lambda &= \alpha\mu \\
\beta\lambda &= \beta\mu
\end{aligned}
\right.
\]
and the ``identical'' closed immersion of $V$ into $V'_0$ \texttt{[rather} $V'_0$\texttt{]}
\DERmark{ENV0708}

\[ p_0 : V \longrightarrow V'_0 \]
\msp{page 15}(which is such that
\[ p_{0!}(e_V) = e_{V'_0} \ \ ), \]
then for $V'$ one takes the category obtained by adding a final object to $V'_0$
\[
V' = \left(
\begin{tikzcd}[row sep=small, column sep=small]
 & & a \arrow[dr, "\rho"] & \\
0' \arrow[r, shift left, "\lambda"] \arrow[r, shift right, "\mu"'] & 0 \arrow[ur, "\alpha"] \arrow[dr, "\beta"'] & & e \\
 & & b \arrow[ur, "\sigma"'] &
\end{tikzcd}
\right)
\qquad \text{with} \quad
\left\{
\begin{aligned}
\alpha\lambda &= \alpha\mu \\
\beta\lambda &= \beta\mu \\
\rho\alpha &= \sigma\beta
\end{aligned}
\right.
\]
and for $p$ the composite functor
\[
\begin{tikzcd}
V \arrow[r, "p_0"] \arrow[rr, bend right, "p"'] & V'_0 \arrow[r, "p_1"] & V',
\end{tikzcd}
\]
where $p_1$ is the inclusion (open). One then finds the square $(*)$, with $\overline{V'}$ not $W$-aspheric, or what amounts to the same thing, $i' \notin W$. For since $\overline{V'}\backslash V' \simeq \overline{V}\backslash V$ is the one-point category, $i' \in W$ is equivalent to $V'/t$ $W$-aspheric (\footnote[8]{assuming (W 7 bis).}) (cf. cor. 4, p. 149 in XII). Now by construction $V'/t \simeq V'_0$, which is not $W$-aspheric (in the case, say, where $W \subseteq W_1$, for $\pi_1(V'_0)=0$ and $\pi_2(V'_0) \simeq \mathbf{Z}$).
\DERmark{ENV0709}


\msp{page 16}Let us then consider the diagram of cocartesian squares
\[
\begin{tikzcd}
V \arrow[r, hook, "i"] \arrow[d, "p"'] & \overline{V} \arrow[r, hook] \arrow[d, "\overline{p}"] & X \arrow[d, "q"] \\
V' \arrow[r, hook, "i'"] & \overline{V'} \arrow[r, hook] & X'
\end{tikzcd}
\qquad .
\]

\printedpage{12}
% head: A. Grothendieck : Dérivateurs, Ch. XII, Édition des Universités de Montpellier II et Paris VII
% foot: Version 20 juin 2003
Since $\overline{V} \hookrightarrow X$ is \texttt{[a]} closed Dwyer immersion, the square on the right is $W$-cocartesian, hence $(\overline{p} \in W) \iff (q \in W)$, hence (since $\overline{p} \notin W$) we have
\[ p \in W, \qquad \text{but} \qquad q \notin W, \]
which indeed shows that $V \longrightarrow X$ is not a $W$-cofibration. Factoring $V \longrightarrow X$ through $U$ (instead of $\overline{V}$), one finds
\[
\begin{tikzcd}
V \arrow[r] \arrow[d, "p\in W"'] & U \arrow[r] \arrow[d, "q_U\in W"] & X \arrow[d, "q\notin W"] \\
V' \arrow[r] & U' \arrow[r] & X'
\end{tikzcd}
\qquad ,
\]
where $q_U \in W$ since $p \in W$ and $V \longrightarrow U$ is \texttt{[a]} Dwyer immersion. This makes explicit that $U \longrightarrow X$ is not a $W$-cofibration, since
\[ q_U \in W, \qquad \text{but} \qquad q \notin W, \]
\hfill q.e.d.

\msp{page 17}\textbf{Gluing in} Cat \textbf{of an open and a complementary closed.}

\texttt{[Suppose given]}
\[ X_0 \underbrace{\hookrightarrow}_{\text{open}} X \ \underbrace{\hookleftarrow}_{\substack{\text{closed}\\ \text{comple-}\\ \text{mentary}}} X_1 \]
in Cat, \texttt{[and let]}
\[ \tilde{X}_{01} \subseteq \mathrm{Fl}(X) \]
be \texttt{[the]} full subcategory of arrows with source in $X_0$, target in $X_1$.

\textbf{Question.} \textit{This diagram}
\[
\begin{tikzcd}[column sep=small]
 & & \tilde{X}_{01} \arrow[dll, "p"' near start, "\text{fibrant, fibers the $x_0\backslash X_1$}"' near end] \arrow[drr, "q" near start, "\text{cofibrant, fibers the $X_0/x_1$}" near end] & & \\
X_0 \arrow[ddrr, "i"'] & & & & X_1 \arrow[ddll, "j"] \\
 & {} \arrow[rr, bend left] & & {} & \\
 & & X, & &
\end{tikzcd}
\tag*{$(*)$}
\]
\printedpage{13}
% head: A. Grothendieck : Dérivateurs, Ch. XII, Édition des Universités de Montpellier II et Paris VII
% foot: Version 20 juin 2003
\textit{is it $W$-cocartesian?}

OK if $i$ $W$-aspheric or $j$ $W$-coaspheric, i.e. $X \longrightarrow \boldsymbol{\Delta}^1$ smooth or proper -- this is a case where one has either $X_0 \xrightarrow{W} X$, or $X_1 \xrightarrow{W} X$. But even assuming $i$ or $j$ in $W$, and moreover $X$ aspheric (hence $X_0$ and $X$ aspheric, or $X_1$ and $X$ aspheric), and $W = \mathbf{W}_\infty$, I have not been able to prove that the diagram is cocartesian, i.e. $q \in W$ resp. $p \in W$ (\footnote[9]{I assume that $W$ satisfies Loc (6).}). (Nor if $X_0 \simeq e$, $X$, $X_1$ $W$-aspheric $\Longrightarrow$ $e\backslash X_1$ $W$-aspheric, i.e. $X_1 \longrightarrow X$ $W$-coaspheric $\ldots$\,) But I have at least proved, in all cases, that % [sic] j (upright) -> $j$, letter denoting a functor printed in roman type
\DERmark{ENV0710}\DERmark{ENV0711}

\[
\begin{tikzcd}[row sep=small]
 & X_0 & & \\
\mathcal{X} \overset{\mathrm{def}}{=} \displaystyle\int \tilde{X}_{01} \arrow[ur] \arrow[dr] & & {} \arrow[r, "\varphi"] & X \\
 & X_1 & &
\end{tikzcd}
\tag*{$(**)$} % verified: CMSY10 glyph, \mathcal{X}; corrections.csv, PDF leaf 14
\]
\DERmark{ENV0712}

is in $\mathbf{W}_1$.

Something \textit{false} (unless $i$ or $j$ in $W = \mathbf{W}_\infty$) if one replaces $\tilde{X}_{01}$ by $\tilde{X}'_{01}$, formed of the triples $(x_0, x_1, c)$, with $x_0 \in \mathrm{Ob}\,X_0$, $x_1 \in \mathrm{Ob}\,X_1$, $c : x_0 \dotto x_1$ a \textit{path} from $x_0$ to $x_1$ in $X$! We have
\msp{page 18}therefore a ``topological'' situation entirely peculiar to the category of models Cat, which does not seem to have an analog in categories such as Esp or $\boldsymbol{\Delta}^\wedge$, where however the notions of open and complementary closed (with a grain of salt in $\boldsymbol{\Delta}^\wedge$) have a meaning.

The quasi-trivial categorical argument which establishes that $(**)$ is in $\mathbf{W}_1$, i.e. induces an equivalence of categories between $\mathrm{Rev}(X)$ and
\[ \mathrm{Rev}(\mathcal{X}) \ \simeq \ \mathrm{Rev}(X_0) \overset{2}{\times}_{\mathrm{Rev}(\tilde{X}_{01})} \mathrm{Rev}(X_1), \]
seems to me to have to extend to higher dimensions, by dint of $n$-gerbes on $X$ -- but the foundations for the theory of $n$-gerbes are lacking. (And I have not had the courage to spell out the argument even for $n = 2$ (where however \textsc{Giraud} provides what is needed) (\footnote[10]{But I have nevertheless made a sketch of a verification for $n = 2$, which convinces me that we do indeed have $\varphi \in \mathbf{W}_2$. And I have hardly any doubt that we even have $\varphi \in \mathbf{W}_\infty$.}).) So I conjecture that $\varphi$ is at the very least in $\mathbf{W}_\infty$, i.e. that the diagram $(*)$ is $W$-cocartesian if $W \supseteq \mathbf{W}_\infty$ (and satisfies Loc (6 a$'$)). % [sic] je n'ai pas guère (footnote 10) -> je n'ai guère, double negation pas + guère
\DERmark{ENV0713}



\endgroup
% END UNIT Der-ChXIIapp-19-06-03/Der-ChXIIapp-19-06-03_p011-014.tex
\clearpage\setcounter{footnote}{0}\phantomsection\addcontentsline{toc}{section}{Chapter XIII}\markboth{Chapter XIII}{Chapter XIII}
% BEGIN UNIT Der-ChXIII-25-01-04/Der-ChXIII-25-01-04_p001-005.tex
\begingroup\makeatletter%
\newcommand{\printedpage}[1]{\marginpar{\footnotesize p.~#1}}%
\newcommand{\msp}[1]{\marginpar{\footnotesize [#1]}}%
\newcommand{\Fl}{\mathrm{Fl}}%
\newcommand{\M}{\mathcal{M}}%
\makeatother\let\printedpage\derpp\let\msp\dermsp%
\let\mspin\DERmspin%


\printedpage{--} % no page number printed on this page
% head: (none)

\begin{center}
\vspace*{5cm}
{\LARGE\sffamily THE DERIVATORS}

\vspace{1.5cm}
{\large\textbf{ALEXANDRE GROTHENDIECK}}

\vspace{2.5cm}
{\large\textbf{Chapter XIII}}

\vspace{0.5cm}
{\large\textbf{Categories of models (1)}}
\end{center}

\par

\printedpage{--} % no page number printed on this page
% head: (none)

\vspace*{6cm}
\noindent \par\noindent\textbf{Note by the editors of the transcribed edition.}\par
This text has been deciphered and transcribed in \LaTeXe{} by M. Künzer. It has been edited by M. Künzer, J. Malgoire, and G. Maltsiniotis. The transcription is as faithful as possible to the manuscript. For the few obvious corrections, or rare comments by the editors, as well as for the original numbering of the manuscript pages, typewriter characters \texttt{[typewriter]} in brackets are used. A question mark in brackets means that one is not sure of the preceding word. The indexes, the list of axioms, the glossary, and the bibliographical references at the end of the chapter have been compiled by the editors. The page numbers in the indexes are those of the original.

\noindent This edition is provisional. Remarks, comments and corrections are welcome. Send a message to:

\noindent\textbf{[e-mail address removed]}

\hfill G. Maltsiniotis\hspace*{2cm}

\par

\printedpage{1}
% head: A. Grothendieck : Dérivateurs, Ch. XIII, Édition des Universités de Montpellier II et Paris VII
% foot: Version 25 janvier 2004

\msp{page 1}

\section*{1\quad Quillen couples $(\Phi, \Psi)$}

I want first to give a weakened variant of the notion of a Quillen triple (Quillen's ``model category''). This will not be an entirely self-dual notion -- thus there will be \textit{two} weakened variants, dual to one another.

But even before that, a sorites on the LLP (left lifting property) and RLP (right lifting property) (\footnote[1]{give formal definition, with reference to \textsc{Quillen}.}). Let $\M$ be any category (most often, $\M$ will be large). We consider subsets
\[
\Phi,\ \Psi\ \subset\ \Fl(\M)
\]
and we define maps
\[
\boxed{\begin{array}{ccc}
\Phi & \longmapsto & \Phi_* \\[1ex]
\Psi & \longmapsto & \Psi^*
\end{array}}
\qquad
\Fl(\M) \longrightarrow \Fl(\M)
\]
\texttt{[rather} $\mathfrak{P}(\Fl(\M) \longrightarrow \mathfrak{P}(\Fl(\M))$\texttt{,}\newline\texttt{where for every set} $E$\texttt{,} $\mathfrak{P}(E)$ \texttt{denotes the set of subsets of} $E$\,\texttt{]} by % [sic] \mathfrak{P}(\Fl(\M) -> \mathfrak{P}(\Fl(\M)), a closing parenthesis is missing in the first term
\[
\begin{aligned}
\Phi_* \ &= \ \{\text{set of } p \in \Fl(\M) \text{ that have the RLP with respect to } \Phi\} \\
\Psi^* \ &= \ \{\text{set of } i \in \Fl(\M) \text{ that have the LLP with respect to } \Psi\},
\end{aligned}
\]
and then
\[
\boxed{\begin{array}{ccc}
\Phi & \longmapsto & \widetilde{\Phi} \\[1ex]
\Psi & \longmapsto & \overline{\Psi}
\end{array}}
\qquad
\Fl(\M) \longrightarrow \Fl(\M)
\]
\texttt{[rather} $\mathfrak{P}(\Fl(\M) \longrightarrow \mathfrak{P}(\Fl(\M))$\texttt{]} by % [sic] \mathfrak{P}(\Fl(\M) -> \mathfrak{P}(\Fl(\M)), a closing parenthesis is missing in the first term
\[
\begin{aligned}
\widetilde{\Phi} \ &= \ (\Phi_*)^* \\
\overline{\Psi} \ &= \ (\Psi^*)_* \ .
\end{aligned}
\]

\msp{page 2}

I summarize the most immediate properties of these operations in a

\noindent\textbf{Proposition 1.}

\begin{enumerate}[label={}]
\item[a)] $\Phi \longmapsto \Phi_*$, $\Psi \longmapsto \Psi^*$ \textit{are decreasing maps from} $\Fl(\M)$ \texttt{[rather} $\mathfrak{P}(\Fl(\M))$\texttt{]} \textit{to itself, for the order relation} $\subset$\textit{, and} $\Phi \longmapsto \widetilde{\Phi}$\textit{,} $\Psi \longmapsto \overline{\Psi}$ \textit{increasing maps.}

\printedpage{2}
% head: A. Grothendieck : Dérivateurs, Ch. XIII, Édition des Universités de Montpellier II et Paris VII
% foot: Version 25 janvier 2004

\item[b)] \textit{We have}
\[
\Phi\ \subset\ \widetilde{\Phi},
\]
\textit{and for $\Phi = \widetilde{\Phi}$ to hold, it is necessary and sufficient that $\Phi$ be of the form $\Psi^*$. We then say that $\Phi$ is} \textup{left closed} (\footnote[2]{Give formal definition.})\textit{. The set of subsets of} $\Fl(\M)$ \textit{that are left closed is stable under arbitrary intersections} (\footnote[3]{more precisely, $\bigcap_i(\Psi_i^*) = (\bigcup_i \Psi_i)^*$}), \textit{and $\widetilde{\Phi}$ is the smallest closed subset of} $\Fl(\M)$ \textit{that contains $\Phi$ and the largest subset $\Phi'$ of} $\Fl(\M)$ \textit{such that $\Phi'_* = \Phi_*$ (or such that $\Phi'_* \supset \Phi_*$).}

\textit{Symmetric statements for $\Psi$: we have}
\[
\Psi\ \subset\ \overline{\Psi},
\]
\textit{we have $\Psi = \overline{\Psi}$ if and only if $\Psi$} \texttt{[is]} \textit{of the form $\Phi_*$, we then say that $\Psi$ is} \textup{right closed.} \textit{The set of subsets of} $\Fl(\M)$ \textit{that are right closed is stable under arbitrary intersections} (\footnote[4]{more precisely, we have $\bigcap_i(\Phi_{i\,*}) = (\bigcup_i \Phi_i)_*$}), \textit{and $\overline{\Psi}$ is the smallest right closed subset containing $\Psi$, and also the largest subset $\Psi'$ of} $\Fl(\M)$ \textit{such that $\Psi'^* = \Psi^*$ (or such that $\Psi'^* \supset \Psi^*$).}

\item[] \hspace*{-2.5em}\msp{page 3}

\item[c)] \textit{The maps}
\[
\Phi \ \longmapsto\ \Phi_*, \qquad\quad \Psi\ \longmapsto\ \Psi^*
\]
\textit{establish mutually inverse bijections between the set of left closed subsets $\Phi$ of} $\Fl(\M)$\textit{, and the set of right closed subsets $\Psi$. This is an anti-isomorphism of ordered sets. We then have}
\[
(\Phi \cap \Phi')_* \ = \quad
\begin{array}{l}
\textit{right closed subset of } \Fl(\M) \ \textit{``generated by } \Phi_*,\ \Phi'_*\textit{'',} \\
\text{i.e. \texttt{[the]} \textit{smallest right closed subset containing} } \Phi_*, \\
\Phi'_*, \text{ i.e. } \overline{\Phi_* \cup \Phi'_*}
\end{array}
\]
\texttt{[for} $\Phi$\texttt{,} $\Phi'$ \texttt{left closed subsets]}\textit{, more generally}
\[
(\bigcap_i \Phi_i)_* \ = \ \overline{\bigcup_i (\Phi_{i\,*})}
\]
\textit{for every family $(\Phi_i)_{i \in I}$ of left closed subsets of} $\Fl(\M)$\textit{. Dual formulas} \texttt{[for right closed subsets]}
\[
\left\{
\begin{aligned}
(\Psi \cap \Psi')^* \ &= \ \widetilde{\Psi^* \cup \Psi'^*} \\
(\bigcap_j \Psi_j)^* \ &= \ \widetilde{\bigcup_j (\Psi_j^*)}\ .
\end{aligned}
\right.
\]

\item[d)] \textit{If} $\Phi, \Psi \subset \Fl(\M)$\textit{, then we have}
\[
\Phi \subset \Psi^* \iff \Psi \subset \Phi_*, \qquad\quad \textit{we then write } \Phi \longleftrightarrow \Psi,
\]
\textit{this means that $\Phi$ has the LLP for $\Psi$, or that $\Psi$ has the RLP for $\Phi$.}
\end{enumerate}

\printedpage{3}
% head: A. Grothendieck : Dérivateurs, Ch. XIII, Édition des Universités de Montpellier II et Paris VII
% foot: Version 25 janvier 2004

\msp{page 4}

\noindent\textbf{Proposition 2.} \textit{Let $\Phi$ be a left closed subset of} $\Fl(\M)$\textit{. Then $\Phi$ satisfies the conditions} a) b) c) \textit{below:}

\begin{enumerate}[label={}]
\item[a)] \textit{$\Phi$ contains the isomorphisms, and is stable under composition, and under cobase change (}i.e. \textit{if we have}
$\vcenter{\hbox{\begin{tikzcd}
A \arrow[r, "i"] \arrow[d] & B \arrow[d] \\
A' \arrow[r, "i'"] & B'
\end{tikzcd}}}$
\textit{a cocartesian square in $\M$, with $i \in \Phi$, we have $i' \in \Phi$).}

\item[b)] \textit{$\Phi$ stable under direct factors (as a subset of the set of objects of} $\underline{\Fl}(\M)$\textit{).}

\item[c)] \textit{$\Phi$ ``stable under ordinal direct limits'', in the following sense: Let $I$ be a well-ordered set, $(A_i)_{i \in I}$ a direct system in $\M$, such that the $u_{ji} : A_i \longrightarrow A_j$ are in $\Phi$, and $A_{i_0} \xleftarrow{\ \sim\ } \varinjlim\limits_{i<i_0} A_i$ if $i_0 \in I$ is a limit element. Suppose that $A = \varinjlim\limits_{I} A_i$ exists. Then $\forall\ i_0 \in I$, the canonical arrow $A_{i_0} \longrightarrow A$ is in $\Phi$.}
\end{enumerate}

There is a dual statement for a right closed subset $\Psi$; it contains the isomorphisms, is stable under composition and under base change, under direct factors, and finally stable under ordinal inverse limits.

\msp{page 5}

We will have a partial converse in accessible categories. It is for the statement of this converse that I have set apart condition \textit{b}) of stability under direct factors, instead of lumping it together with \textit{a}).

\noindent\textbf{Definition 1.}

\begin{enumerate}[label={}]
\item[\textit{a})] (For the record.) RLP, LLP, the relation $\Phi \longleftrightarrow \Psi$ between subsets of $\Fl(\M)$ (nonsymmetric relation!).

\item[\textit{b})] (For the record.) Subset of $\Fl(\M)$ left closed, right closed.

\item[\textit{c})] Let $\Phi \longleftrightarrow \Psi$. We say that the pair $(\Phi, \Psi)$ has the \textit{Quillen factorization property}, if \textit{every} arrow $f$ of $\M$ factors as
\[
f = pi, \qquad\quad i \in \Phi,\ p \in \Psi.
\]
We say that $(\Phi, \Psi)$ is a \textit{Quillen couple} if it is a pair associated by \texttt{[the]} properties LLP and RLP ($\Phi \longleftrightarrow \Psi$), and \texttt{[if]} they \texttt{[the subsets} $\Phi$ \texttt{and} $\Psi$\texttt{]} are closed (\textit{i.e.} $\Phi = \Psi^*$, $\Psi = \Phi_*$), and if they have the Quillen factorization property.
\end{enumerate}

I state right now the essential technical result, ensuring the existence of many Quillen couples:


\endgroup
% END UNIT Der-ChXIII-25-01-04/Der-ChXIII-25-01-04_p001-005.tex
% BEGIN UNIT Der-ChXIII-25-01-04/Der-ChXIII-25-01-04_p006-010.tex
\begingroup\makeatletter%
\newcommand{\printedpage}[1]{\marginpar{\footnotesize p.~#1}}%
\newcommand{\msp}[1]{\marginpar{\footnotesize [#1]}}%
\newcommand{\cM}{\mathcal{M}}%
\newcommand{\Fl}{\mathrm{Fl}}%
\newcommand{\Cof}{\mathrm{Cof}}%
\newcommand{\Fib}{\mathrm{Fib}}%
\newcommand{\Ob}{\mathrm{Ob}}%
\makeatother\let\printedpage\derpp\let\msp\dermsp%
\let\mspin\DERmspin%


\setcounter{footnote}{4}

\printedpage{4}
% head: A. Grothendieck : Dérivateurs, Ch. XIII, Édition des Universités de Montpellier II et Paris VII
% foot: Version 25 janvier 2004
\msp{page 6}\textbf{Theorem 1.} \textit{Suppose that $\cM$ is} accessible, \textit{and stable under (small) direct limits. (There should suffice stability under amalgamated sums, and under filtering $\varinjlim$, or even ordinal \dots).} \textit{Let $\Phi_0$ be a} small \textit{subset of} $\Fl(\cM)$\textit{, $\Psi=(\Phi_0)_*$, and $\Phi$ a subset of} $\Fl(\cM)$ \textit{such that}
\[
\Phi_0\ \subset\ \Phi\ \subset\Psi^*\ (=\widetilde{\Phi}_0).
\]
\textit{Then:}

\begin{itemize}[leftmargin=2.5em]
\item[a)] \textit{For $\Phi$ to be left closed,} i.e. $\Phi=\widetilde{\Phi}_0$\textit{, it is necessary and sufficient that it satisfy conditions} a) b) c) \textit{of proposition 2. When this is so, the couple $(\Phi,\Psi)$ is a Quillen couple.}

\item[b)] \textit{Suppose only that $\Phi$ satisfies conditions} a), c) \textit{of proposition 2 (excluding condition} b) \textit{on direct factors) – thus we cannot assert that $\Phi$ is left closed. Nevertheless, $(\Phi,\underbrace{\Psi}_{=\,\Phi_*})$ satisfies the Quillen factorization property.}
\end{itemize}

\textbf{Remark.} In the proof, the existence of $\Phi_0\subset\Phi$, which plays the role of a small set of ``generators'' (on the left) for $\Phi$, plays an essential role. It would be worth developing the notion of \textit{accessible subsets}
\msp{page 7}
of $\Fl(\cM)$, which would imply the existence of a small set $\Phi_0\subset\Phi$ such that $\Phi\subset\widetilde{\Phi}_0$, and which would have the property that
\[
\begin{aligned}
&\Phi\ \text{accessible}\quad\Longrightarrow\quad\Phi_*\ \text{accessible}\\
&\Psi\ \text{accessible}\quad\Longrightarrow\quad\Psi^*\ \text{accessible.}
\end{aligned}
\]
Any couple $(\Phi,\Psi)$ of associated closed subsets ($\Phi=\Psi^*$, $\Psi=\Phi_*$) and accessible, would thus be a Quillen couple.

\section*{2\quad Weak Quillen categories, and construction thereof (Conjectural theorem-scholie)}

Throughout what follows, we give ourselves a \textit{category of models}
\[
\left\{
\begin{array}{ll}
\mbox{$(\cM,W)$, where $W\subset\Fl(\cM)$} & \\[3ex]
\mbox{\textit{$W$ weakly saturated} (\footnotemark) :} &
\begin{array}[c]{@{}l@{\quad}p{6.5cm}@{}}
a) & $W$ contains the isomorphisms.\\[1ex]
b) & If two among $f$, $g$, $gf$ \texttt{[are]} in $W$, the third also.
\end{array}
\end{array}
\right.
\]
\footnotetext{This is axiom M 5 of \textsc{Quillen}.}

\printedpage{5}
% head: A. Grothendieck : Dérivateurs, Ch. XIII, Édition des Universités de Montpellier II et Paris VII
% foot: Version 25 janvier 2004
We want to describe additional data a little weaker than those of a Quillen model category. We suppose given subsets
\[
C,\ F,\ TC,\ TF\ \subset\ \Fl(\cM),
\]
with the following properties. As a preliminary axiom, I lay down

\begin{itemize}[leftmargin=3em,labelsep=0.6em]
\item[\textbf{M 0}] Stability of $\cM$ under finite $\varprojlim$ and finite $\varinjlim$ \texttt{[rather $\varinjlim$]}.

\item[\textbf{M 1}]\msp{page 8}Lifting-extension property
\[
C\longleftrightarrow TF,\qquad TC\longleftrightarrow F.
\]
\DERmark{ENV0584}


\item[\textbf{M 2}] Factorization property for $(C,TF)$, $(TC,F)$.

\item[\textbf{M 3}] Stability of $F$ under composition and under base change, of $C$ under composition and cobase change, $F$ and $C$ contain the isomorphisms.

\item[\textbf{M 4}] Stability of $TF$ under composition and base change, of $TC$ under composition and cobase change, and $TF$ and $TC$ contain the isomorphisms.

\item[\textbf{M}$'$] $TC\subset C\cap W$, $TF\subset F\cap W$.

\item[\textbf{M}$''$] $C\subset\ \underbrace{\Cof_W}_{\text{$W$-cofibrations}}\ ,\ F\subset\ \underbrace{\Fib_W}_{\text{$W$-fibrations}}\ .$
\end{itemize}

\texttt{[A morphism $p:X\longrightarrow Y$ of $\cM$ is a $W$-fibration if for every diagram of cartesian squares of the form:}
\[
\begin{tikzcd}
X \ar[d,"p"'] & X' \ar[l] \ar[d] & X'' \ar[l,"t"'] \ar[d] \\
Y & Y' \ar[l] & Y'' \ar[l,"s"] \ ,
\end{tikzcd}
\]
\texttt{if $s\in W$, then $t\in W$. The notion of $W$-cofibration is the dual notion.]}

In the case of \textsc{Quillen} (Q-\textit{categories of models}), M$'$ is strengthened to
\[
TC=C\cap W,\qquad TF=F\cap W,
\]
and there is no axiom M$''$. Axiom M$''$ is equivalent to the property of ``propriety'', corresponding to Q-categories of models that are \textit{proper.}

\textbf{Proposition 3.} \textit{Suppose conditions} M 0 \textit{(for the record),} M 1, M 2, M$'$, M$''$ \textit{are satisfied. Then replacing $C$, $TC$ by the saturations for the stability properties of containing the isomorphisms, of being stable under composition and under cobase change, and $F$, $TF$ by those of containing the isomorphisms and of being stable under composition and under base change, we find $C'$, $TC'$, $F'$, $TF'$ satisfying}
\printedpage{6}
% head: A. Grothendieck : Dérivateurs, Ch. XIII, Édition des Universités de Montpellier II et Paris VII
% foot: Version 25 janvier 2004
\msp{page 9}\textit{all the axioms} M 1 to M 4, M$'$, M$''$. % [sic] $M''$ (italic M) -> M$''$ (roman M), the axiom is printed with roman M everywhere else

\textsc{Proof}. Proposition 2 assures us that M 1 continues to be satisfied, and for M 2 this is trivial. M 3 and M 4 will hold by construction. On the other hand, M$''$ remains valid, since $\Cof_W$ has the stabilities required for $C'$, and $\Fib_W$ those required for $F'$. It remains to prove M$'$, \textit{i.e.}
\DERmark{ENV0585}

\[
TC'\ \subset\ C'\cap W,\qquad TF'\subset F'\cap W,
\]
and for this it suffices to see that $C'\cap W$ resp. $F'\cap W$ has the stabilities required for $TC'$ resp. $TF'$. This follows from the

\textbf{Lemma.} \textit{Let $C$ be a subset of} $\Cof_W$\textit{. If $C$ contains the isomorphisms, the same is true of $C\cap W$. If moreover $C$ is stable under composition resp. under cobase change, the same is true} \texttt{[of]} \textit{$C\cap W$. Dual statement for a subset $F$ of} $\Fib_W$\textit{.}

This comes from the fact that
\[
\boxed{\ \Cof_W\cap W=W^{\mathrm{univ}}\ (\mbox{\footnotemark})\ }
\stackrel{\mathrm{def}}{=}
\begin{array}{l}
\text{set of arrows $A\longrightarrow B$ such that for every}\\
A\longrightarrow A',\ A'\longrightarrow B'=A'\vee_A B\ \ \text{is in $W$.}
\end{array}
\]
\DERmark{ENV0586}

\footnotetext{arrows \textit{couniversally} in $W$, or $W$-\textit{equivalences that are couniversal}.}
is itself stable under composition and under cobase change, and of course we have, if $C\subset\Cof_W$,
\[
C\cap W\ =\ C\cap(\Cof_W\cap W)\ =\ C\cap W^{\mathrm{univ}}.
\]

\msp{page 10}So it remains to prove the

\textbf{Corollary 1.} \textit{Suppose there exist subsets $C$, $F$, $TC$, $TF$ $\subset$} $\Fl(\cM)$ \textit{satisfying} M 2, M$'$ (\footnote{We do not need all of M$'$, only $TC\subset W$, $TF\subset W$.}), M$''$. \textit{Then}
\[
\begin{aligned}
\Cof_W\cap W\ &=\ W^{\mathrm{univ}}\\
\Fib_W\cap W\ &=\ W_{\mathrm{univ}}\ \ (\mbox{\footnotemark}).
\end{aligned}
\]
\footnotetext{Arrows $f:X\longrightarrow Y$ \textit{universally} in $W$, \textit{i.e.} such that for every $Y'\longrightarrow Y$, $f':X'=X\times_Y Y'\longrightarrow Y'$ is in $W$.}

To see this, let us show e.g. the second (the first will follow from it by duality).
\[
\begin{tikzcd}
X \ar[d,"f"'] & X' \ar[l] \ar[d,"f'"] \\
Y & Y' \ar[l]
\end{tikzcd}
\qquad\qquad\qquad
f\in\Fib_W,\ f\in W\Longrightarrow f'\in W.
\]

\printedpage{7}
% head: A. Grothendieck : Dérivateurs, Ch. XIII, Édition des Universités de Montpellier II et Paris VII
% foot: Version 25 janvier 2004
We have better:

\textbf{Proposition 4.} \textit{Suppose that every arrow of $\cM$ factors as $pi$, with $i\in W$ and $p\in$} $\Fib_W$\textit{. Then, if $X$, $Y$ are in} $\underline{\Fib}_W S\subset\cM/S$ \texttt{[full subcategory of $\cM/S$ consisting of the $r:Z\longrightarrow S$, where $r\in\Fib_W$]}\textit{, then every $f:X\longrightarrow Y$ in $\cM/S$ which is in} $W_S$ \texttt{[consisting of the arrows over $S$ which are in $W$]} \textit{is in $W^{\mathrm{u}}_S$,} i.e. \textit{it remains in $W$ after any base change $S'\longrightarrow S$.}
\[
\begin{tikzcd}[column sep=small]
X \ar[rr,"f"] \ar[dr,"p"'] & & Y \ar[dl,"q"] \\
& S &
\end{tikzcd}
\qquad\qquad\qquad
p,q\in\Fib_W
\]

\textit{Dual statement for the arrows of} $S\backslash\underline{\Cof}_W\subset S\backslash\cM$ \texttt{[full subcategory of $S\backslash\cM$ consisting of the $i:S\longrightarrow Z$, where $i\in\Cof_W$]}\textit{, such an arrow which is in $W$ is $S$-couniversally in $W$, if every arrow of $\cM$ factors as $qj$, with $j\in$} $\Cof_W$ \textit{and $q\in W$.}

\msp{page 11}\textsc{Proof}.
\[
\begin{tikzcd}[column sep=small]
X \ar[dr,"f\in W"] \ar[dd] & & \overline{X'} \ar[ll] \ar[dr,"\overline{f'}\in W"] \ar[dd] & & X' \ar[ll,"i_X"'] \ar[dr,"f'"] \ar[dd] & \\
& Y \ar[dl] & & \overline{Y'} \ar[ll,crossing over] \ar[dl] & & Y' \ar[ll,"i_Y"',crossing over] \ar[dl] \\
S & & \overline{S'} \ar[ll,"F"] & & S' \ar[ll,"i"',"TC"] &
\end{tikzcd}
\]
We factor $S'\longrightarrow S$ as $S'\overset{i}{\longrightarrow}\overline{S'}\overset{p}{\longrightarrow}S$, with $i\in W$, $p\in\Fib_W$, by M 2. Thus $f\in W\Longrightarrow\overline{f'}\in W$. Since $X,Y\in\Ob\,\underline{\Fib}_W S$, we have $\overline{X'},\overline{Y'}\in\Ob\,\underline{\Fib}_W\overline{S'}$, and since $i:S'\longrightarrow\overline{S'}$ is in $W$, hence also $i_X$, $i_Y$, we have $f'\in W$ by M 5 (weak saturation of $W$).

\textbf{Comment on Proposition 3.} We can be even more demanding, by also requiring in M 3, M 4 stability under direct factors, and (when only M 1, M 2, M$'$, M$''$ % [sic] M2 -> M 2, printed without the space used everywhere else
\DERmark{ENV0587}

are satisfied) by saturating $C$, $TC$, $F$, $TF$ also for this additional stability. Proposition 2 again ensures that M 1 continues to hold (and M 2, M 3, M 4 trivially), we see that this is still so for M$'$, M$''$ provided we strengthen the weak saturation axiom M 5 on $W$ by
\[
\textit{stability of $W$ under direct factors,}
\]
which indeed implies that $\Cof_W$, $\Fib_W$ are stable under direct factors, as are $W^{\mathrm{univ}}$ and $W_{\mathrm{univ}}$.

\msp{page 12}It does not seem that axioms M 0 to M 4, plus M$'$ and M$''$ suffice to develop the exact sequences of suspension and cosuspension. A tiny little bit more is needed:

\printedpage{8}
% head: A. Grothendieck : Dérivateurs, Ch. XIII, Édition des Universités de Montpellier II et Paris VII
% foot: Version 25 janvier 2004
\textbf{Definition 2.} I call a \textit{left pre-Q-category of models,} the data of $\cM$ and of $W$, $C$, $F$, $TC$, $TF$ $\subset$ $\Fl(\cM)$, satisfying axioms M 0 to M 5, M$'$, M$''$, plus the strengthened form of the first inclusion in M$'$
\[
\text{M}'_0\ \text{(g)}\qquad\qquad\qquad TC=C\cap W\qquad\quad(\text{not only } TC\subset C\cap W).
\]
Likewise, we define the \textit{right pre-Q-category of models,} by the condition
\[
\text{M}'_0\ \text{(d)}\qquad\qquad\qquad\qquad\qquad\qquad TF=F\cap W.
\]

NB When the two conditions are satisfied at once, then the data of $TC$, $TF$ already follow from those of $W$, $C$, $F$, and we simply have a ``proper'' Q-category of models (in the sense of \textsc{Quillen}).

\textbf{Scholie.} For me, the introduction of $C$, $TC$, $F$, $TF$ is purely auxiliary – what matters to me are the properties of $W$. Thus, the existence of $C$, $TC$, $F$, $TF$ satisfying the weak axioms M 0 to M 4, M$'$, M$''$, since it gives the factorizations
\msp{page 13}
spoken of in \texttt{[the]} proposition 4, entails the remarkable property of $W$ stated in that proposition, and in the corollary which precedes it. Likewise, these factorization properties (\footnote{in a slightly strengthened form it is true – a \textit{functorial} factorization for $f$ is needed \dots}) ensure the existence of a good theory of $W$-homotopic fibers and cofibers without any need for reference in this theory to $C$, $TC$, $F$, $TF$ (whose introduction gives at most some small supplements \dots). But I do not see how, from the factorization hypotheses alone, to prove the crucial exactness property for $W$-cartesian, and $W$-cocartesian squares, necessary to be able to develop the \textit{exact} sequence of cosuspension and of suspension. It is to get this exactness that the auxiliary, and transitory, introduction of the paraphernalia $(C,TC,F,TF)$ is useful, the additional axiom on the left $TC=C\cap W$ ensuring (unless I am mistaken) the exactness of $W$-cartesian squares, and the axiom on the right $TF=F\cap W$ that of $W$-cocartesian squares. (Something to be checked with more care than I have done so far.) Once this is obtained, and possibly likewise for the $(\cM(I),W(I))$ for every diagram-category $I$, we can
\DERmark{ENV0588}

\msp{page 14}
totally forget these auxiliary sets – I feel like forgetting them in any case.

My hope is that for a given $W$, even if we do not manage to find an associated Quillen triple structure $(W,C,F)$ (e.g. because there might not be one), we can find both quintuples $(W,C,F,TC,TF)$ which are left pre-Quillen, and quintuples (\textit{other} quintuples!) which are right pre-Quillen.

Here is the kind of statement in proper form, for the moment conjectural or even entirely heuristic, that I have in mind:

\textbf{Conjectural theorem-scholie:} \textit{We suppose that $\cM$ is stable under ``arbitrary'' $\varinjlim$, and under finite $\varprojlim$, and that $W$ is stable under filtering $\varinjlim$ (in addition to weak saturation).}


\endgroup
% END UNIT Der-ChXIII-25-01-04/Der-ChXIII-25-01-04_p006-010.tex
% BEGIN UNIT Der-ChXIII-25-01-04/Der-ChXIII-25-01-04_p011-015.tex
\begingroup\makeatletter%
\newcommand{\printedpage}[1]{\marginpar{\footnotesize p.~#1}}%
\newcommand{\msp}[1]{\marginpar{\footnotesize [#1]}}%
\newcommand{\circled}[1]{\tikz[baseline=(c.base)]\node[draw,circle,inner sep=1pt](c){#1};}%
\newcommand{\Cof}{\mathrm{Cof}}%
\newcommand{\Fib}{\mathrm{Fib}}%
\newcommand{\Fl}{\mathrm{Fl}}%
\newcommand{\univ}{\mathrm{univ}}%
\makeatother\let\printedpage\derpp\let\msp\dermsp%
\let\mspin\DERmspin%


\printedpage{9}
% head: A. Grothendieck : Dérivateurs, Ch. XIII, Édition des Universités de Montpellier II et Paris VII

\begin{itemize}[leftmargin=2.5em]
\item[\circled{1}] \textit{Then} $\Cof_W$ \textit{and} $W^{\univ}$ \textit{have the stability properties} a) \textit{and} c) \textit{of proposition} 2 (p.~4), \textit{and also the stability property} b) \textit{under direct factors, provided that} $W$ \textit{itself has it.}

\item[\circled{2}] \textit{Suppose} $\mathcal{M}$ accessible, \textit{so that we have theorem} 1. \textit{Equivalent conditions on} $W$ \textit{:}

\msp{page 15}

\begin{itemize}[leftmargin=2.5em]
\item[a)] \textit{There is a small subset} $C_0$ \textit{of} $\Cof_W$\textit{, and} $TC_0$ \textit{of} $W^{\univ}$\textit{, such that}
\[ TF \overset{\mathrm{def}}{=} C_{0*} \subset W, \qquad (TC_0)_* \overset{\mathrm{def}}{=} F \subset \Fib_W. \]
\DERmark{ENV0589}\DERmark{ENV0590}


\item[a$'$)] \textit{There is a quadruple} $(C,TC,F,TF)$ \textit{of subsets of} $\Fl(\mathcal{M})$\textit{, satisfying conditions} M~1 \textit{to} M~4, M$'$, M$''$ \textit{on page} 6 (\footnote[10]{with $C$, $TC$ left closed, $F$, $TF$ right closed.}).

\item[a$''$)] \textit{As in} a$'$)\textit{, assuming} M$'$ \textit{in strengthened form with} M$'_0$ (g), i.e. \textit{we have a left pre-Q-category of models structure (}$TC = C \cap W$\textit{).}

\item[b)] \textit{We have}
\[ (\Cof_W)_* \ \subset\ W, \qquad (W^{\univ})_* \ \subset\ \Fib_W \]
\textit{(thus} $(\Cof_W)_* \subset \Fib_W \cap W$ \textit{since} $(\Cof_W)_* \subset (W^{\univ})_*$ (\footnote[11]{$W^{\univ} \subset \Cof_W$.})\textit{).}

\item[b$'$)] \textit{We have}
\[ (\Cof_W)_* \ \subset\ W_{\univ}, \qquad (W^{\univ})_* \ \subset\ \Fib_W. \]
\end{itemize}

\item[\circled{3}] \textit{For there to be a quadruple which, with} $W$\textit{, defines a right pre-Q-category of models (}$TF = F \cap W$\textit{)} (\footnote[12]{We are still assuming $\mathcal{M}$ accessible in 3° and 4°, with in addition $F$, $TF$ (and $C$, $TC$) right closed (resp. left closed).})\textit{, it is necessary and sufficient that we have} b$'$) \textit{above, as well as}
\[ (W^{\univ})_* \cap W_{\univ} \ \ \text{\itshape is right closed.} \]

\msp{page 16}

\item[\circled{4}] \textit{The preliminary assumptions on} $(\mathcal{M}, W)$\textit{, including accessibility, as well as the equivalent conditions in}\circled{2}\textit{, are stable under passage from} $(\mathcal{M}, W)$ \textit{to} $(\mathcal{M}(I), W(I))$\textit{. But what about the condition in}\circled{3} \textit{? Is it a consequence of the conditions in} \circled{2} \textit{(}i.e. \textit{is it equivalent to them) ?} % [sic] leurs est-elle -> lui est-elle, grammar
\end{itemize}

I am going to sketch the proof of the theorem-scholie, to pin down the technical work that remains to be done.

1°) Is routine verification.

2°) Here are the $\pm$ tautological implications between the conditions under consideration : % [sic] les implication -> les implications, agreement
\DERmark{ENV0591}

\[ a'' \Longrightarrow a' \overset{?}{\Longrightarrow} a \Longrightarrow b \Longleftarrow b', \tag*{$(*)$} \]
% foot: Version 25 janvier 2004
\printedpage{10}
% head: A. Grothendieck : Dérivateurs, Ch. XIII, Édition des Universités de Montpellier II et Paris VII
but I see that $a' \Longrightarrow a$ is not at all tautological, and perhaps even false after all~! (\footnote[13]{That is why I modified it, strengthening $a'$) by \texttt{[the condition of note 10]}.}). We would have to assume that in \texttt{[the]} quadruple $(C,F,TC,TF)$, $F$ and $TF$ are right closed, and more precisely, that $F = (TC)_*$, $TF = C_*$, which will be the case e.g. if we have a ``closed'' Q-category of models structure (``closed model category''), in the sense of \textsc{Quillen}. The other three implications in $(*)$ really are tautological.

$a \Longrightarrow a'$. We can in $a$) assume
\[ TC_0 \ \subset\ C_0 \]
(replacing $C_0$ by $C_0 \cup TC_0$ if necessary), then, thanks to
\msp{page 17}
1°) ($W^{\univ}$ and $\Cof_W$ are left closed), setting
\[ \widetilde{TC_0} = TC, \qquad \widetilde{C_0} = C, \qquad \text{\texttt{[rather}}\ \ \widetilde{TC_0} = TC, \qquad \widetilde{C_0} = C, \text{\texttt{]}} \]
\DERmark{ENV0592}

we will have $TC \subset W^{\univ}$, $C \subset \Cof_W$. Passing from $TC_0$ to $TC$, from $C_0$ to $C$, we have lost only the fact that these subsets of $\Fl(\mathcal{M})$ are small. But by theorem 1, setting
\[ F \ =\ (TC)_*, \qquad TF \ =\ C_*, \]
the pairs of closed subsets (left -- right)
\[ (C,\, TF), \qquad (TC,\, F) \]
satisfy the Quillen factorization condition. We will then have properties M~1 to M~4, plus the inclusions $TC \ \subset\ C \cap W$ (since $TC \ \subset\ W^{\univ} \ \subset\ W$, and $TC \ \subset\ C$ by construction), and $C \subset \Cof_W$ of M$'$ and M$''$. It remains to check $TF \subset F \cap W$, $F \subset \Fib_W$, but $F$ $(= (TC)_*)$ $\ \subset \Fib_W$ by assumption, and $TF$ $(= C_*)$ $\ \subset W$ too~; on the other hand $TF \subset F$ follows from $TC \subset C$.

$a' \Longrightarrow a''$. We thus want, in addition to conditions M~1 to M~4 and M$'$, M$''$, the condition
\[ TC \ =\ C \cap W\ (= C \cap W^{\univ}). \]
The idea is to replace the pair
\msp{page 18}
\[ C,\ TC \qquad (C \subset \Cof_W,\ TC \subset C \cap W), \]
giving rise to
\[ F = (TC)_*, \qquad TF = C_*, \]
by
\[ C,\ C \cap W\ (\supset TC), \]
% foot: Version 25 janvier 2004
\printedpage{11}
% head: A. Grothendieck : Dérivateurs, Ch. XIII, Édition des Universités de Montpellier II et Paris VII
which does not change the pair $(C, TF)$, but does change the pair $(TC, F)$, replaced by
\[ \underbrace{(TC)'}_{\supset TC} \ =\ \underbrace{C \cap W}_{= C\cap W^{\univ}}, \qquad F' \ =\ \underbrace{(TC')_*}_{\subset F} \]
(we have replaced $F$ by a smaller $F'$) (\footnote[14]{NB $TC'$ will still be left closed as an intersection of two left closed subsets.}). The only problem is whether $(TC', F')$ still satisfies the Quillen factorization condition. These are closed subsets (left for $TC'$, right for $F'$) ``dual'' to one another,
\[ TC' = F'^*, \qquad F' = TC'_* \ , \]
but since we do not know whether $TC'$ is of the form $\widetilde{TC'_0}$, with $TC'_0$ small, we are not sure that the conclusion of theorem 1 applies. But if we assume in $a'$) that $C$ is an \textit{accessible} subset of $\Fl(\mathcal{M})$ (\textit{i.e.}, already being left closed, that $C = \widetilde{C_0}$, $C_0$ small),
\msp{page 19}
then I presume that we will be able to prove that the same is true of $C \cap W$ ($W$ itself being accessible, thanks to the condition of stability under small filtering $\varinjlim$), so theorem 1 really does apply to $TC' = C \cap W$. But if we keep $a'$) as it stands, then we would probably have to prove that $a'$) \textit{implies} $a'$) in the stronger form, with $C$ \textit{accessible.} But since we have $a' \Longrightarrow a$, and we have just established that $a \Longrightarrow a'$ in this stronger form, we win (modulo the conjectural fact assumed about the accessible \texttt{[rather accessible]} subsets of $\Fl(\mathcal{M})$).
\DERmark{ENV0593}


It thus remains to prove
\[ b \ \Longrightarrow\ a, \qquad b \ \Longrightarrow\ b', \]
and it suffices to prove the first implication, since we know that $a$) $\Longrightarrow$ $a'$), and that $a'$) implies that
\[ \Fib_W \cap W \ =\ W_{\univ} \]
(cor.~1, page 10), so $(\Cof_W)_* \subset W$ and $(W^{\univ})_* \subset \Fib_W$ (\textit{i.e.} $b$)) imply, since $(\Cof_W)_* \subset (W^{\univ})_*$ (because $W^{\univ} \subset \Cof_W$), $(\Cof_W)_* \subset W \cap \Fib_W = W_{\univ}$, which is the part of $b'$) to be established (the other $(W^{\univ})_* \subset \Fib_W$ being tautologically contained in $b$).

\msp{page 20}

Let us now note that $b$) is nothing other than the weakened variant of $a$), when we omit specifying that $C_0$, $TC_0$ must be small. So the only problem, for $b \Longrightarrow a$, is to see this :

\textbf{Conjectural lemma} (\footnote[15]{This lemma looks a bit false, but it should be possible to fix that \dots}).

\begin{itemize}[leftmargin=2.5em]
\item[\circled{1}] \textit{If we have a left closed subset} $\Phi$ \textit{of} $\Fl(\mathcal{M})$\textit{, such that} $\Phi_* \subset \Fib_W$ \textit{(resp.} $\Phi_* \subset W$\textit{), then there is a subset} of small size $\Phi_0 \subset \Phi$\textit{, such that} $\Phi_{0*} \subset \Fib_W$ \textit{(resp.} $(\Phi_0)_* \subset W$\textit{).}

\item[\circled{2}] \textit{If} $C$\textit{,} $C'$ \textit{are two left closed subsets of} $\Fl(\mathcal{M})$\textit{, if} $C$ \textit{and} $C'$ \textit{are accessible, the same is true of} $C \cap C'$\textit{.}
\end{itemize}
\DERmark{ENV0594}

% foot: Version 25 janvier 2004

\printedpage{12}
% head: A. Grothendieck : Dérivateurs, Ch. XIII, Édition des Universités de Montpellier II et Paris VII
I believe this really is true under the assumptions made on $\mathcal{M}$, $W$ -- which must imply that in a suitable sense, $W$ is an accessible subset of the set $\Fl(\mathcal{M})$ of arrows of the accessible category $\mathcal{M}$. It should follow from this that a subset of $\mathcal{M}$ constructed explicitly in terms of $W$, such as $\Fib_W$, is itself accessible, and the lemma should come out of that.

Thus, 2° is proved, modulo this accessibility lemma, which I will have to prove $\pm$ at the same time as theorem 1 (which I have only stated here).

\msp{page 21}

Let us prove 3°. First suppose there is a ``closed'' quadruple $(C, TC, F, TF)$, fitting into a ``closed Quillen pre-quintuple'' which is a right Q-category of models, \textit{i.e.}
\[ TF \ =\ F \cap W \ =\ F \cap W_{\univ}. \]
Let us then prove \texttt{[that]}
\[ \underbrace{(W^{\univ})_* \cap W_{\univ}}_{A_W} \qquad \text{is right closed.} \tag*{$(*)$} \]
Indeed, we have
\[ TC \ \subset\ W^{\univ}, \]
hence
\[ (W^{\univ})_* \ \subset\ (TC)_* \ =\ F, \]
so
\[ A = (W^{\univ})_* \cap W_{\univ} \ \subset\ F \cap W_{\univ} \ \overset{\text{(by hyp.)}}{=}\ TF, \]
now $TF$ is right closed and contained in $W_{\univ}$, hence
\[ \overline{A} \ \subset\ TF \ \subset\ W_{\univ}, \]
and since $(W^{\univ})_*$ is right closed, we also have
\[ \overline{A} \ \subset\ (W^{\univ})_*, \]
hence $\overline{A} \subset W_{\univ} \cap (W^{\univ})_* = A$, so $A$ is right closed, q.e.d.

Let us prove the converse, \textit{i.e.} suppose the conditions of 2° are satisfied, \textit{plus} condition $(*)$ above. We must construct the quadruple.

\msp{page 22}

Let us first note, to get our bearings, that the relations $b'$) of 2°), and the transposed relations (taking into account that $\Cof_W$ and $W^{\univ}$ are left closed, by 1°), namely
\[ \left\{ \begin{array}{lllllll}
(\Cof_W)_* & \subset & W_{\univ} & , & (W^{\univ})_* & \subset & \Fib_W \\
(W_{\univ})^* & \subset & \Cof_W & , & (\Fib_W)^* & \subset & W^{\univ},
\end{array} \right. \]
giving (\footnote[16]{taking into account $W^{\univ} \subset \Cof_W$, $W_{\univ} \subset \Fib_W$ and the transposed inclusions.}) the two square diagrams of inclusions
% foot: Version 25 janvier 2004
\printedpage{13}
% head: A. Grothendieck : Dérivateurs, Ch. XIII, Édition des Universités de Montpellier II et Paris VII
\[
\begin{tikzcd}[column sep=small, row sep=small]
 & W^{\univ} \arrow[dr, hook] & \\
(\Fib_W)^* \arrow[ur, hook] \arrow[dr, hook] & & \Cof_W \\
 & (W_{\univ})^* \arrow[ur, hook] &
\end{tikzcd}
\qquad\qquad (\text{\footnotemark[17]})
\tag*{(I)}
\]
\footnotetext[17]{\textbf{NB} We have these squares of inclusions, whenever we have a Quillen system $(W, C, F, TC, TF)$ that is \textit{closed.}}
and
\[
\begin{tikzcd}[column sep=small, row sep=small]
 & (W^{\univ})_* \arrow[dl, hook'] & \\
\Fib_W & & (\Cof_W)_*\ , \arrow[ul, hook'] \arrow[dl, hook'] \\
 & W_{\univ} \arrow[ul, hook', "\alpha"'] &
\end{tikzcd}
\tag*{(II)}
\]
of which the first is the transpose (under $\Phi \longmapsto \Phi^*$) of the second, and where the 8 sets of arrows considered, except for $\Fib_W$ and $W_{\univ}$ (linked by $\alpha$ in square II) are closed (the four sets in I left closed, and the two sets in II not involved in $\alpha$, right closed).

\msp{page 23}

The ``closed Quillen pre-quintuples'', associated with $W$, correspond to the choice of left closed subsets $TC$ and $C$, fitting into diagram (I) as shown in I$'$ below, and give rise by transposition to $F$ and $TF$ as in II$'$ below,
\[
\begin{tikzcd}[column sep=small, row sep=small]
 & & W^{\univ} \arrow[dr, hook] & \\
 & TC \arrow[ur, hook] \arrow[dr, hook] & & \Cof_W \\
(\Fib_W)^* \arrow[ur, hook] \arrow[dr, hook] & & C \arrow[ur, hook] & \\
 & (W_{\univ})^* \arrow[ur, hook] & &
\end{tikzcd}
\tag*{(I$'$)}
\]
\[
\begin{tikzcd}[column sep=small, row sep=small]
 & & (W^{\univ})_* \arrow[dl, hook'] & \\
 & F \arrow[dl, hook'] & & (\Cof_W)_* \arrow[ul, hook'] \arrow[dl, hook'] \\
\Fib_W & & TF \arrow[ul, hook'] \arrow[dl, hook'] & \\
 & W_{\univ}\ , \arrow[ul, hook'] & &
\end{tikzcd}
\tag*{(II$'$)}
\]
the ``catch'' thus being to choose $TC$ and $C$ large enough that we have
\[ \underbrace{(TC)_*}_{= F} \ \subset\ \Fib_W, \qquad \underbrace{C_*}_{= TF} \ \subset\ W_{\univ} \qquad \text{--} \]
% foot: Version 25 janvier 2004


\endgroup
% END UNIT Der-ChXIII-25-01-04/Der-ChXIII-25-01-04_p011-015.tex
% BEGIN UNIT Der-ChXIII-25-01-04/Der-ChXIII-25-01-04_p016-020.tex
\begingroup\makeatletter%
\newcommand{\printedpage}[1]{\marginpar{\footnotesize p.~#1}}%
\newcommand{\msp}[1]{\marginpar{\footnotesize [#1]}}%
\newcommand{\univ}{\mathrm{univ}}%
\newcommand{\Cof}{\mathrm{Cof}}%
\newcommand{\Fib}{\mathrm{Fib}}%
\newcommand{\cercle}[1]{\mbox{\tikz[baseline=(c.base)]\node[draw,circle,inner sep=1pt](c){#1};}}%
\newcommand{\clarrow}{\mathrel{\leftarrow\mkern-4mu{+}\mkern-4mu\rightarrow}}%
\makeatother\let\printedpage\derpp\let\msp\dermsp%
\let\mspin\DERmspin%


\printedpage{14}
% head: A. Grothendieck : Dérivateurs, Ch. XIII, Édition des Universités de Montpellier II et Paris VII
% foot: Version 25 janvier 2004

\msp{page 24}
\noindent and the inclusions stipulated in 2$^\circ$, $b'$) ensure that we can indeed find $TC$ and $C$, if need be taking them as large as possible, $TC = W^{\univ}$, $C = \Cof_W$ -- in which case, moreover, we will have
\[
TC \;=\; C \cap W^{\univ},
\]
\textit{i.e.} we have a left pre-Q-category of models. But now, on the other hand, the point is to ensure that
\[
TF \;=\; F \cap W_{\univ},
\]
which implies, as we have just seen, (noting $(W^{\univ})_* \subset F$) that
\[
\overline{(W^{\univ})_* \cap W_{\univ}} \;\subset\; \overline{F \cap W_{\univ}} = \overline{TF} = TF \subset W_{\univ}. \tag*{$(*)$}
\]
Conversely, if $(*)$ is satisfied, let us take
\[
F = (W^{\univ})_*
\]
as small as possible, \textit{i.e.} $C = W^{\univ}$ \texttt{[rather} $TC = W^{\mathtt{univ}}$\texttt{]} as large as possible. We must therefore take
\[
TF \;=\; F \cap W_{\univ} \;=\; (W^{\univ})_* \cap W_{\univ},
\]
and this set is indeed right closed by hypothesis.

\textbf{Corollary.} \textit{Suppose condition} b$'$) \textit{of the theorem-scholie is satisfied. The following conditions} a), b) \textit{on $W$ are equivalent (and imply the condition of} 3$^\circ$ \textit{of the theorem-scholie).}

\msp{page 25}
\begin{itemize}[leftmargin=3em]
\item[\cercle{a}] $(W^{\univ})_* \cap \overline{W_{\univ}} = (\Cof_W)_*$.

\item[\cercle{b}] $\widetilde{W^{\univ}\cup(W_{\univ})^*} = \Cof_W$, i.e. $\Cof_W$ \textit{is the left closed subset generated by} $W^{\univ}$ \textit{and} $(W_{\univ})^*$\textit{,} % doubt resolved: printed tilde ends before the final star; preserve the edition reading, see corrections.csv (PDF leaf 16).
\end{itemize}
\DERmark{ENV0595}


\noindent\textit{and these conditions imply}

\begin{itemize}[leftmargin=3em]
\item[\cercle{c}] \textit{that there exists a structure of a} closed Quillen triple $(W, C, F)$ \textit{including} $W$ \textit{(thus with} $TC = W \cap C$, $TF = W \cap F$\textit{) (which will be closed in Quillen's sense if and only if} $W$ \texttt{[is]} \textit{stable under direct factors -- for $C$ and $F$ are already spontaneously stable).}
\end{itemize}

\noindent Indeed, it suffices then to take
\[
TC \;=\; W^{\univ}, \qquad\quad C \;=\; \Cof_W .
\]
But I doubt that the existence of a closed Quillen triple $(W, C, F)$ suffices to imply conditions \textit{a}), \textit{b}) above, or even the weaker condition
\[
(W^{\univ})_* \cap W_{\univ} \;=\; (\Cof_W)_*
\]
which means \textit{exactly} that $C = \Cof_W$, $TC = W^{\univ}$, and their transposes $TF = (\Cof_W)_*$, $F = (W^{\univ})_*$ form a closed Quillen category.

\par
\printedpage{15}
% head: A. Grothendieck : Dérivateurs, Ch. XIII, Édition des Universités de Montpellier II et Paris VII
% foot: Version 25 janvier 2004

\msp{page 26}
\noindent Conditions on $(\mathcal{M}, W)$ with
\begin{itemize}[leftmargin=3em]
\item[\textit{a})] $\mathcal{M}$ stable under finite $\varprojlim$ and (small) $\varinjlim$,
\item[\textit{b})] $W$ stable under filtering $\varinjlim$,
\item[\textit{c})] $\mathcal{M}$ accessible.
\end{itemize}

\[
\mathrm{I}\ \left\{
\begin{array}{@{}l@{\quad}p{11cm}@{}}
\cercle{a} & existence of a weak Quillen system $(W, C, TC, F, TF)$ (closed if one insists) \\[1ex]
\Updownarrow & \\[1ex]
\cercle{b} & existence of a left exact weak Quillen system ($TC = C \cap W$). \\[1ex]
\Updownarrow & \\[1ex]
\cercle{c} & $\underbrace{(\Cof_W)_* \subset W,\ (W^{\univ})_* \subset \Fib_W}_{(*)}$ ($\Longleftrightarrow$ $(\Cof_W)_* \subset W_{\univ}$, $(W^{\univ})_* \subset \Fib_W$).
\end{array}\right.
\]

\noindent$\Uparrow$
\[
\mathrm{II}\ \left\{
\begin{array}{@{}l@{\quad}p{11cm}@{}}
\cercle{a} & existence of a right exact weak Quillen system ($TF = F \cap W$) \\[1ex]
\Updownarrow & \\[1ex]
\cercle{d} & We have $(*)$, and moreover $(W^{\univ})_* \cap W_{\univ}$ right closed. \\[1ex]
\Updownarrow & \\[1ex]
\cercle{b} & existence of a closed Quillen system (closed model category), $TC = C \cap W$, $TF = F \cap W$, $F = (TC)_*$, $TF = C_*$, $TC = F^*$, $C = (TF)^*$. \\[1ex]
\Updownarrow & \\[1ex]
\cercle{c} & Putting $TC = W^{\univ}$, $F = (TC)_*$, $TF = F \cap W_{\univ}$, $C = (TF)^*$, we find together with $W$ a closed Quillen system. % [sic] $C = (TF)_*$ -> $C = (TF)^*$, cf. the line b) above where $C = (TF)^*$
\end{array}\right.
\]
\DERmark{ENV0596}


\noindent$\Uparrow$
\[
\mathrm{III}\ \left\{
\begin{array}{@{}p{11.5cm}@{}}
$C = \Cof_W$ and $TC = W^{\univ}$ are part (together with $W$) of a closed Quillen system (thus $F = (W^{\univ})_*$, $TF = (\Cof_W)_*$, $W = TF \circ TC$). \\[1ex]
$\Updownarrow$ \\[1ex]
\texttt{[We have]} $(*)$, and moreover $(W^{\univ})_* \cap W_{\univ} = (\Cof_W)_*$.
\end{array}\right.
\]

\noindent$\Uparrow$
\[
\mathrm{IV}\ \left\{
\begin{array}{@{}p{11.5cm}@{}}
We have $(*)$ and $\Cof_W = \widetilde{(W_{\univ})^* \cup W^{\univ}}$\ \ . \\[1ex]
$\Updownarrow$ \\[1ex]
We have $(*)$ and $(W^{\univ})_* \cap \overline{W_{\univ}} = (\Cof_W)_*$.
\end{array}\right.
\]

\par
\printedpage{16}
% head: A. Grothendieck : Dérivateurs, Ch. XIII, Édition des Universités de Montpellier II et Paris VII
% foot: Version 25 janvier 2004

\noindent Case $\mathcal{M} = \mathrm{Cat}$, $W = \mathbf{W}_\infty$. We know that II is satisfied, by \textsc{Thomason}. Thus we have
\[
(\Cof_W)_* \subset W_{\univ}, \qquad\quad (W^{\univ})_* \subset \Fib_W, \qquad\quad (W^{\univ})_* \cap W_{\univ} \ \text{ right closed.}
\]
I have been able to prove the first two inclusions directly (\textit{i.e.} I), but not that $(W^{\univ})_* \cap W_{\univ}$ is right closed. (I should at least try!) (\footnote[18]{It would suffice to know that $\overline{W_{\univ}} \cap \Fib_W = W_{\univ}$, which would imply that III and IV are equivalent (\footnotemark[19]).}\footnotetext[19]{But this is false, e.g. if $X$ \texttt{[is a]} nonempty discrete category, $X \longrightarrow e$ is in $\overline{W_{\univ}} \cap \Fib_W$.}). We do not have IV, \textit{i.e.} we have
\[
(W^{\univ})_* \cap \overline{W_{\univ}} \;\neq\; (\Cof_W)_*
\]
for
\[
(W^{\univ})_* \cap \overline{W_{\univ}} \not\subset W_{\univ}
\]
(contains the $X \longrightarrow e$ with $X$ nonempty and discrete), so we have
\DERmark{ENV0597}

\[
(W^{\univ})_* \cap \overline{W_{\univ}} \;\;\underset{\neq}{\supset}\;\; \underbrace{(W^{\univ})_* \cap W_{\univ}}_{\text{is \textit{not} the orthogonal of } W^{\univ} \cup (W_{\univ})^*} \quad . % doubt resolved: printed strict-superset sign confirmed; see corrections.csv (PDF leaf 18).
\]

\msp{page 27}

\section*{3\quad Quillen category (second version)}

We are still given a pair
\[
(\mathcal{M}, W), \qquad\quad W \subset \mathrm{Fl}(\mathcal{M}), \qquad\quad W \text{ weakly saturated,}
\]
$\mathcal{M}$ satisfies M 0 (stability under finite $\varprojlim$ and finite $\varinjlim$).

\noindent\textbf{Definition 1 :} A \textit{weak Quillen structure} associated to the category of models $(\mathcal{M}, W)$ is the datum of four subsets $(C, TC, F, TF)$ of $\mathrm{Fl}(\mathcal{M})$, satisfying the following conditions.

\begin{itemize}[leftmargin=3em,labelsep=0.5em,align=left]
\item[\textbf{M 1}] $C \longleftrightarrow TF$, $TC \longleftrightarrow F$ (\textit{i.e.} $TF \subset C_*$, $F \subset (TC)_*$).
\item[\textbf{M 2}] Quillen factorization property for $(C, TF)$ and $(TC, F)$.
\item[\textbf{M 3}] Stability of $C$ under composition and under cobase change, of $F$ under composition and base change -- $C$ and $F$ contain the isomorphisms.
\item[\textbf{M 4}] Similar stability for $TC$ and $TF$.
\item[$\mathbf{M}'$] $TC \subset C \cap W$, $TF \subset F \cap W$.
\item[$\mathbf{M}''$] $C \subset \Cof_W$, $F \subset \Fib_W$.
\end{itemize}

\noindent We say that the Quillen structure is \textit{exact} if it satisfies the additional axiom

\par
\printedpage{17}
% head: A. Grothendieck : Dérivateurs, Ch. XIII, Édition des Universités de Montpellier II et Paris VII
% foot: Version 25 janvier 2004

\noindent\textbf{M 2 bis}\qquad $W = TF \circ TC$.

\noindent We say that it is \textit{left exact} if it satisfies the additional axiom

\noindent\textbf{M$'_0$ (g)}\qquad $TC = C \cap W$ (instead of only $TC \subset C \cap W$),

\noindent\textit{right exact} if it satisfies

\noindent\textbf{M$'_0$ (d)}\qquad $TF = F \cap W$ (instead of only $TF \subset F \cap W$) (\footnote[20]{\textbf{NB} If $C$, $F$, $W$ \texttt{[are]} stable under direct factors (thus also $C \cap W$ and $F \cap W$), then left exactness $TC = C \cap W$ \textit{is equivalent} to right exactness $TF = F \cap W$, and implies that $C$, $TC$, $F$, $TF$ are \textit{closed} (on the left for $C$, $TC$, on the right for $F$, $TF$).}).

\msp{page 28}
\noindent\textbf{Remark.}
\begin{itemize}[leftmargin=3em]
\item[a)] A left exact or right exact Quillen structure is exact.
\item[b)] If we have a datum $(C, TC, F, TF)$ satisfying M 1, M 2, M$'$, $M''$, when we saturate $C$, $TC$ on the left and $F$, $TF$ on the right, for the stabilities specified in M 3, M 4, we find a weak Quillen structure. (All the axioms are obvious, except M$'$, for which I refer to p. 9, 10, in particular prop. 4.) If we had M 2 bis, this will still hold (a fortiori) after saturation, \textit{i.e.} we have an exact Quillen structure.
\end{itemize}
\DERmark{ENV0598}\DERmark{ENV0599}


\noindent\textbf{Definition 1 bis.} The weak Quillen structure is called a \textit{Quillen structure} (without qualification) if it is left and right exact, or (as one says) ``biexact'', then $(C, F)$ already determines $(C, TC, F, TF)$. (That corresponds exactly to the ``model categories'' of \textsc{Quillen} \texttt{[proper]}.) We say that the weak Quillen structure is \textit{closed} if we have % [sic] determine -> détermine, missing accent
\[
C \clarrow TF\ , \qquad\quad TC \clarrow F\ ,
\]
\textit{i.e.} $TF = C_*$, $C = (TF)^*$ and $F = (TC)_*$, $TC = F^*$ (\footnote[21]{This also means that $C$, $TF$, $TC$, $F$ are stable under direct factors. We say that the weak Quillen structure is \textit{left closed} (resp. \textit{right closed}) if $C$, $TC$ (resp. $F$, $TF$) are left closed (resp. right closed).}). A weak Quillen structure that is both biexact (\textit{i.e.} not weak at all!) and closed is called a \textit{perfect} Quillen structure. (That corresponds exactly to the ``closed model category'' \texttt{[proper]} of \textsc{Quillen}.)

\msp{page 29}
\noindent\textbf{Remark.} For a weak Quillen structure to be perfect, it suffices that it be biexact and that we have
\[
C = (TF)^*, \qquad\quad F = (TC)_*
\]
(which already implies that $TC = F^*$, $TF = C_*$, \textit{i.e.} that $TC$ is left closed and $TF$ right closed). This is proved in Quillen. It also implies that $W$ is strongly saturated -- it's the land of plenty!

\par
\printedpage{18}
% head: A. Grothendieck : Dérivateurs, Ch. XIII, Édition des Universités de Montpellier II et Paris VII
% foot: Version 25 janvier 2004

\noindent\textbf{Scholie :} The pair $TC \longleftrightarrow F$, satisfying the Quillen factorization condition, and satisfying $TC \subset W$, $F \subset \Fib_W$, already allows one to develop a satisfactory theory of $W$-homotopic fibers in $\mathcal{M}$ -- but without an exactness property for $W$-cartesian squares. Dually, the datum of $C \longleftrightarrow TF$, with factorization, and $TF \subset W$, $C \subset \mathrm{Cofib}_W$, gives a good theory of $W$-homotopic cofibers, but without an exactness property for $W$-cocartesian squares. The weak Quillen structure brings together these two data dual to one another, with in addition the additional link between these two data, by $TC \subset C$, $TF \subset F$ (stipulated in $M'$). This datum already ensures that % [sic] donnés -> données, feminine noun as in "ces deux données duales"
\[
\underbrace{\mathcal{M}_{\mathrm{cf}}}_{\substack{\text{full subcategory}\\ \text{of $\mathcal{M}$ consisting of}\\ \text{objects fibrant and}\\ \text{cofibrant at once}}} W_{\mathrm{cf}}^{-1} \;\xrightarrow{\ \sim\ }\; \mathcal{M}W^{-1} = \mathrm{Ho}_W
\]
\msp{page 30}
is an equivalence, but not that the localization functor
\[
\mathcal{M}_{\mathrm{cf}} \longrightarrow \mathcal{M}_{\mathrm{cf}}W_{\mathrm{cf}}^{-1}
\]
is a simple passage to the quotient on arrows, \textit{i.e.} that if $X$, $Y$ are in $\mathcal{M}_{\mathrm{cf}}$, then
\[
\mathrm{Hom}_{\mathcal{M}}(X, Y) \;\longrightarrow\; \underbrace{\mathrm{Hom}_W(X, Y)}_{\overset{\mathrm{def}}{=}\,\mathrm{Hom}_{\mathrm{Ho}_W}(X, Y)} \tag*{$(*)$}
\]
\DERmark{ENV0600}

is surjective. This is ensured by the exactness axiom M 2 bis,
\[
W \;=\; TF \circ TC.
\]
With this axiom alone, we have no nice characterization of the equivalence relation in $\mathrm{Hom}_{\mathcal{M}}(X, Y)$ defined by the arrow $(*)$ (for $X$, $Y$ in $\mathcal{M}_{\mathrm{cf}}$).

If we have left exactness,
\[
TC = C \cap W \qquad\quad (\Longleftrightarrow\ TC = C \cap W^{\univ}),
\]
then this equivalence relation is the one generated by the relation of \textit{right Q-homotopy} (\footnote[22]{It would seem that the relation of right Q-homotopy is an equivalence relation, as soon as $Y$ \texttt{[is]} in $\mathcal{M}_{\mathrm{f}}$ and without assuming the weak Quillen structure to be right or left exact.}, \footnote[23]{\textbf{NB} According to \textsc{Baues}, it would seem that on $\mathcal{M}_{\mathrm{cf}}$, this relation of right homotopy is equivalent to the relation of left homotopy, and that the latter does not depend on the cylinder chosen \ldots}) (implied by left Q-homotopy), and the quotient category thus defined
\[
\pi_{\mathrm{d}}\mathcal{M}_{\mathrm{cf}} \;=\; \pi\mathcal{M}_{\mathrm{cf}} \;\simeq\; \mathrm{Ho}_W,
\]
maps into $\pi_{\mathrm{d}}\mathcal{M}_{\mathrm{c}}$,
\[
\pi\mathcal{M}_{\mathrm{cf}} \longrightarrow \pi_{\mathrm{d}}\mathcal{M}_{\mathrm{c}},
\]
and this functor can be interpreted (unless I am mistaken)


\endgroup
% END UNIT Der-ChXIII-25-01-04/Der-ChXIII-25-01-04_p016-020.tex
% BEGIN UNIT Der-ChXIII-25-01-04/Der-ChXIII-25-01-04_p021-025.tex
\begingroup\makeatletter%
\derleqno%
\newcommand{\printedpage}[1]{\marginpar{\footnotesize p.~#1}}%
\newcommand{\msp}[1]{\marginpar{\footnotesize [#1]}}%
\newcommand{\cM}{\mathcal{M}}%
\newcommand{\circled}[1]{\tikz[baseline=(c.base)]\node[draw,circle,inner sep=1pt](c){#1};}%
\makeatother\let\printedpage\derpp\let\msp\dermsp%
\let\mspin\DERmspin%


\setcounter{footnote}{23}

\printedpage{19}
% head: A. Grothendieck : Dérivateurs, Ch. XIII, Édition des Universités de Montpellier II et Paris VII
% foot: Version 25 janvier 2004
\msp{page 31}
as a \textit{right adjoint} (fully faithful) of the localization functor
\[
\pi_{\mathrm{d}}\cM_{\mathrm{c}} \longrightarrow (\pi_{\mathrm{d}}\cM_{\mathrm{c}})W_{\mathrm{c}}^{-1} \;\simeq\; \mathrm{Ho}_W,
\]
and $\mathrm{Ho}_W$ is obtained from $\pi_{\mathrm{d}}\cM_{\mathrm{c}}$ by a \textit{calculus of left fractions.} Dually, if we have right exactness, \textit{i.e.} $TF = F\cap W$ ($\Longleftrightarrow$ $TF = F\cap W_{\mathrm{univ}}$), then the relevant equivalence relation on $\mathrm{Hom}_{\cM}(X,Y)$, for $X$, $Y$ in $\cM_{\mathrm{cf}}$, is left Q-homotopy (implied by right Q-homotopy), and this time the natural arrow from $\pi_{\mathrm{g}}\cM_{\mathrm{cf}} = \pi\cM_{\mathrm{cf}} \simeq \mathrm{Ho}_W$ to $\pi_{\mathrm{d}}\cM_{\mathrm{f}}$ \texttt{[rather }$\pi_{\mathtt{g}}\cM_{\mathtt{f}}$\texttt{]},
\[
\underbrace{\pi\cM_{\mathrm{cf}}}_{\simeq\,\mathrm{Ho}_W} \longrightarrow \pi_{\mathrm{d}}\cM_{\mathrm{f}} \qquad \texttt{[rather }\pi_{\mathtt{g}}\cM_{\mathtt{f}}\texttt{]}
\]
is a \textit{left adjoint} (fully faithful) of the localization functor
\[
\pi_{\mathrm{d}}\cM_{\mathrm{f}} \longrightarrow \cM_{\mathrm{f}}W_{\mathrm{f}}^{-1} \;\simeq\; \mathrm{Ho}_W \qquad \texttt{[rather }\pi_{\mathtt{g}}\cM_f \longrightarrow \cM_{\mathtt{f}}W_{\mathtt{f}}^{-1} \;\simeq\; \mathrm{Ho}_W\texttt{]}, % [sic] \cM_f (italic f, twice) -> \cM_{\mathrm{f}}, elsewhere the index f is upright
\]
\DERmark{ENV0601}

and the localized category is made explicit by a calculus of right fractions.

In the biexact case (\textit{i.e.} that of genuine Quillen structures), the two Q-homotopy relations in $\cM_{\mathrm{cf}}$ (more generally in $\mathrm{Hom}(X,Y)$, if $X$ in $\cM_{\mathrm{c}}$ and $Y$ in $\cM_{\mathrm{f}}$)
\msp{page 32}
coincide, and even become particularly simple. (For if $f, g : X \rightrightarrows Y$, with $X$ cofibrant and $Y$ fibrant, then $f\sim g$ can be tested on \textit{any} ``cylinder object'' for $X$, or on \textit{any} ``path object'' for $Y$.)

In the case where the weak Quillen structure is left exact, \textit{i.e.} where the relevant equivalence relation in $\mathrm{Fl}\,\cM_{\mathrm{cf}}$ is right Q-homotopy, then the exactness axiom is satisfied for $W$-cartesian squares, but I doubt that the same holds for $W$-cocartesian squares (\footnote{It seems nevertheless that this is the case, according to \textsc{Baues}, cf. corresponding marginal annotation page 30.}). (Unless we make an additional hypothesis, e.g. right exactness of the Quillen structure.) If we have only exactness (but neither left nor right), it seems that we cannot infer from it the exactness either of $W$-cartesian squares or of $W$-cocartesian squares, hence not that of cosuspension sequences (loop sequences), nor of suspension sequences. We might therefore think that the exactness condition, all by itself, is not of much
\msp{page 33}
use for the homotopy formalism in the categories of fractions $\mathrm{Ho}_W$ etc. But we do nevertheless have this: let $\Psi = {}$\begin{tikzcd}[row sep=tiny, column sep=small, baseline=(\tikzcdmatrixname-2-1.base)]
 & 1 \\
0 \arrow[ru] \arrow[rd] & \\
 & 2
\end{tikzcd}, and suppose that $(\cM,W)$ \textit{and} $(\cM(\Psi),W(\Psi))$ \texttt{[where for every small category }$I$\texttt{, }$\cM(I)$\texttt{ denotes the category of presheaves on }$I$\texttt{ with values in }$\cM$\texttt{, and }$W(I)$\texttt{ is the subset of }\texttt{Fl}$(\cM(I))$\texttt{ consisting of morphisms}
\printedpage{20}
% head: A. Grothendieck : Dérivateurs, Ch. XIII, Édition des Universités de Montpellier II et Paris VII
% foot: Version 25 janvier 2004
\texttt{of presheaves that are in }$W$\texttt{ ``argument by argument'']} admit an exact weak Quillen structure. Then in $\cM$ the $W$-cartesian squares are $W$-exact, so we have in $\cM$ long exact cosuspension sequences (exact loop sequences) (\footnote{it seems that I botched it here, and I doubt that the result mentioned is correct.}). (The hypothesis on $(\cM(\Psi),W(\Psi))$ serves only to ensure that every arrow of $\mathrm{Ho}_{W(\Psi)} = \cM(\Psi)W(\Psi)^{-1}$ is isomorphic to an arrow coming from an arrow of $\cM(\Psi)$ itself. But this is an apparently very delicate property, and I know of no other way to obtain it than via the existence of an exact weak Quillen structure.) Dually, if $(\cM,W)$ and $(\cM(\Phi),W(\Phi))$ admit exact weak Quillen structures, then the $W$-cocartesian squares in $\cM$ are $W$-exact, and we find (if $\cM$ is pointed) long exact suspension sequences. % [sic] W-exactes -> W-exacts, agreement with "carrés" (cf. "W-exacts" below)

\msp{page 34}
This shows that if, for a given category of models $(\cM,W)$, we are interested in ``doing homotopical algebra'' for all the pairs $(\cM(I),W(I))$ at once, with $I$ in Cat, it is enough (if we want to have a good theory of $W$-homotopic fibers and cofibers in these categories, including the property of $W$-exactness for $W$-cartesian squares and for $W$-cocartesian squares) to have for each $(\cM(I),W(I))$ an exact weak Quillen structure, nothing more.

This therefore shows us the interest of developing criteria, for a $(\cM,W)$, ensuring that this is indeed so. On this subject, I note that I know how to check the exactness of a Quillen structure \textit{only if} I have left exactness, or right exactness. This means that I must in fact bring out conditions on $(\cM,W)$ ensuring (say) that the $(\cM(I),W(I))$ admit left exact weak Quillen structures.

\msp{page 35}
The weak Quillen structures for the $(\cM(I),W(I))$ will be deduced from a weak Quillen structure for $(\cM,W)$, in the following way: they will all be \textit{``right closed'',} \textit{i.e.}
\[
F(I) \;=\; (TC(I))_*, \qquad TF(I) \;=\; (C(I))_*, \tag*{(1)}
\]
in particular (for $I=e$)
\[
F \;=\; (TC)_*, \qquad TF \;=\; C_*,
\]
which already implies (since $C\subset \mathrm{Cof}_W$, $TC\subset W^{\mathrm{univ}}$)
\[
(\mathrm{Cof}_W)_* \;\subset\; W, \qquad (W^{\mathrm{univ}})_* \;\subset\; \mathrm{Fib}_W \tag*{(2)}
\]
(whence $(\mathrm{Cof}_W)_*\subset W_{\mathrm{univ}}$). By virtue of (1), to define the weak Quillen structures for the $(\cM(I),W(I))$, it suffices to define $C(I)$, $TC(I)$, and even, since we want left exactness, \textit{i.e.}
\[
TC(I) \;=\; C(I)\cap W(I) \tag*{(3)}
\]
\printedpage{21}
% head: A. Grothendieck : Dérivateurs, Ch. XIII, Édition des Universités de Montpellier II et Paris VII
% foot: Version 25 janvier 2004
(in particular
\[
TC \;=\; C\cap W \quad ), \tag*{(4)}
\]
it remains to define $C(I)$. We shall define $C(I)$ ``argument by argument'', \textit{i.e.} for $f\in\mathrm{Fl}(\cM(I))$
\DERmark{ENV0602}

\[
f\in C(I) \qquad\Longleftrightarrow\qquad \alpha_s^*(f)\in C \ \text{ for every } s\in\mathrm{Ob}\,I, \tag*{(5)}
\]
where
\[
\alpha_s : e\longrightarrow I \qquad \texttt{[is]}\ \text{defined by } \alpha_s(e)=s.
\]
It follows moreover at once, by (3),
\[
f\in TC(I) \qquad\Longleftrightarrow\qquad \alpha_s^*(f)\in TC \ \text{ for every } s\in\mathrm{Ob}\,I. \tag*{(6)}
\]

\msp{page 36}
We therefore suppose that $(\cM,W)$ admits a left exact and right closed weak Quillen structure $(C,TC,F,TF)$, defined in terms of
\[
C\;\subset\;\mathrm{Fl}(\cM),
\]
by
\[
TC \;=\; C\cap W, \qquad F \;=\; (TC)_*, \qquad TF = C_*.
\]
For this to really give a weak Quillen structure (necessarily left exact and right closed), we simply have to check the following conditions.

\begin{itemize}[leftmargin=3em]
\item[\circled{i}] Factorization for $(C, TF=C_*)$ and $(TC=C\cap W, F=(TC)_*)$ – we shall come back to this point.
\item[\circled{ii}] Stability of $C$ under composition and cobase change, and $C$ contains the isomorphisms (from which the similar stability of $TC = C\cap W$ will follow, taking into account that $C\subset\mathrm{Cof}_W$) (NB The appropriate stabilities for $F$, $TF$ are trivial by definition $F=(TC)_*$, $TF=C_*$.)
\item[\circled{iii}] $C\subset\mathrm{Cof}_W$, $\underbrace{C_*}_{=\,TF}\subset W$, $(C\cap W)_*\subset\mathrm{Fib}_W$.
\end{itemize}

To check, similarly, that $(C(I),TC(I),F(I),TF(I))$ is a weak Quillen structure for $W(I)$, we must check the conditions similar to \circled{i}\,\circled{ii}\,\circled{iii} . Now (\textit{ii}) for $C(I)$ follows at once from condition (\textit{ii}) for $C$. It remains to examine (\textit{iii}) and (\textit{i}).

\printedpage{22}
% head: A. Grothendieck : Dérivateurs, Ch. XIII, Édition des Universités de Montpellier II et Paris VII
% foot: Version 25 janvier 2004
\msp{page 37}
The inclusions of (\textit{iii}).

The inclusion
\[
C\;\subset\;\mathrm{Cof}_W
\]
implies that every $f\in C(I)$ is in $\mathrm{Cof}_W$ argument by argument. This implies that $f$ is in $\mathrm{Cof}_{W(I)}$, as we see at once from the definition – I make it explicit in the

\textbf{Lemma 1:} \textit{Let $f : X_\bullet\longrightarrow Y_\bullet$ be an arrow of $\cM(I)$. For $f$ to be in} $\mathrm{Cof}_{W(I)}$ (resp. \textit{in} $\mathrm{Fib}_{W(I)}$), \textit{it} suffices \textit{that for every $s$ in} $\mathrm{Ob}\,I$, \textit{$\alpha_s^*(f)$ be in} $\mathrm{Cof}_W$ (resp. \textit{in} $\mathrm{Fib}_W$)\textit{. Likewise (for the record) for a condition for $f\in W(I)$ – there it is necessary and sufficient that $\alpha_s^*(f)\in W\ \forall\, s\in\mathrm{Ob}\,I$.}

We shall apply this same lemma to establish the inclusions
\[
C(I)_* \;\subset\; W(I), \qquad (C(I)\cap W(I))_* \;\subset\; \mathrm{Fib}_{W(I)}.
\]
So it suffices to see that
\[
\begin{aligned}
\alpha_s^*(C(I)_*) \;&\subset\; C_* &&\qquad (\subset\; W)\\
\alpha_s^*((\underbrace{C(I)\cap W(I)}_{(C\cap W)(I)})_*) \;&\subset\; (C\cap W)_* &&\qquad (\subset\; \mathrm{Fib}_W),
\end{aligned}
\]
and more generally
\[
\alpha_s^*(\Phi(I)_*) \;\subset\; \Phi_* \tag*{(7)}
\]
for any ( ?) subset $\Phi$ of $\mathrm{Fl}(\cM)$, thus defining $\Phi(I)\subset\mathrm{Fl}(\cM(I))$ term by term.

\msp{page 38}
Making this explicit, we find the condition
\[
\alpha_{s\,!}(\Phi)\subset\Phi(I) \qquad \text{for every } s\in\mathrm{Ob}\,I,
\]
for it to work, \textit{i.e.} we want that
\[
i\in\Phi \qquad\Longrightarrow\qquad \alpha_{s\,!}(i)\in\Phi(I), \quad \textit{i.e.}\ \ \alpha_t^*(\alpha_{s\,!}(i))\in\Phi\ \ \forall\, t\in\mathrm{Ob}\,I.
\]
Now the functor $\alpha_t^*\alpha_{s\,!}$ is made explicit as
\[
A \longmapsto \underbrace{\mathrm{Hom}_I(t,s)\times A}_{\substack{\text{sum of } \mathrm{Hom}_I(t,s)\\ \text{copies of } A \text{ in } \cM}},
\]
so we must suppose $\Phi$ \textit{stable under direct sums} (of the type considered $\mathrm{Hom}_I(t,s)\times A$). This is indeed the case for $\Phi = C$ and $\Phi = C\cap W$ and for direct sums that are \textit{finite,} thanks to (\textit{ii}). We thus find the

\printedpage{23}
% head: A. Grothendieck : Dérivateurs, Ch. XIII, Édition des Universités de Montpellier II et Paris VII
% foot: Version 25 janvier 2004
\textbf{Lemma 2.}

\begin{itemize}[leftmargin=3em]
\item[a)] \textit{Let $I$ in} Cat \textit{be such that direct sums of type} $\mathrm{Hom}_I(t,s)$ \textit{exist in $\cM$. (E.g. we suppose that $I$ is} finite, \textit{more generally that the} $\mathrm{Hom}_I(t,s)$ \textit{are finite sets; or that ``small'' direct sums exist in $\cM$.) Let $\Phi$ be a subset of} $\mathrm{Fl}(\cM)$\textit{, stable under this type of direct sums, and let} $\Phi(I)\subset\mathrm{Fl}(\cM(I))$ \textit{be the corresponding subset. Then for every $s\in\mathrm{Ob}\,I$, we have}

\msp{page 39}
\[
\alpha_s^*(\Phi(I)_*) \;\subset\; \Phi_*.
\]
\item[b)] \textit{Suppose that the} $\mathrm{Hom}_I(t,s)$ \textit{are finite (e.g. $I$ finite). Then the hypothesis of} a) \textit{applies for $\Phi=C$ and for $\Phi=C{\cap}W$, so condition} (iii) \textit{p.} 36 \textit{is satisfied for the structure $(W(I),C(I),TC(I),F(I),TF(I))$ in $\cM(I)$. Same conclusion without a finiteness hypothesis on the} $\mathrm{Hom}_I(t,s)$\textit{, when we suppose $\cM$ stable under ``small'' direct sums, and $C$ and $W$ also stable under such direct sums.}
\end{itemize}

\textbf{NB} Given the scholie above, it is perfectly reasonable to restrict ourselves if need be to categories $I$ that are \textit{finite} – given that the category of these $I$ is stable under the products $I\longmapsto I\times\Psi$, $I\longmapsto I\times\Phi$.

It remains, to be happy, to check the factorization condition \circled{i} for $(C,TF)$, $(TC,F)$, but also and above all for the $(C(I),TF(I),TC(I),F(I))$. This means that even starting from a closed Quillen structure (land of Cockaigne)
\msp{page 40}
does not get us very far all by itself, for lack of knowing that the factorization axiom is preserved, in passing from $(\cM,W,\ldots)$ to $(\cM(I),W(I),\ldots)$.

It is here that the factorization theorem of no.~1 should prove providential. But to be able to apply it to the pairs $(C(I),C(I)_*)$, $(TC(I),TC(I)_*)$, we need (in addition to the accessibility of $\cM$, which is a quite innocuous condition) very strong stability conditions for $C(I)$, $TC(I)$ (\footnote{Plus the \textit{small} systems of generators for $C$, $TC$.}). To tell the truth, there remains the stability of $C$ and of $C\cap W$ under ordinal direct limits (which will imply the same stability for the $C(I)$, $(C{\cap}W)(I) = C(I){\cap}W(I)$).

Truth to tell, I see little chance of finding a weak Quillen structure $(C,TC,F,TF)$ associated with $W$, and satisfying all the stated conditions, including that of stability of $C$ and of $C\cap W$ under ordinal $\varinjlim$, without supposing that $W$ itself is already stable under (small) filtering direct limits. % [sic] associé -> associée, agreement with "structure"

\msp{page 41}
We can think of using a \textit{functorial} factorization of the arrows in $\cM$ as
\[
\begin{aligned}
f = pi \quad \text{with}\quad & i\in TC, && p\in F\\
\text{or}\quad & i\in C, && p\in TF,
\end{aligned}
\]


\endgroup
% END UNIT Der-ChXIII-25-01-04/Der-ChXIII-25-01-04_p021-025.tex
% BEGIN UNIT Der-ChXIII-25-01-04/Der-ChXIII-25-01-04_p026-030.tex
\begingroup\makeatletter%
\newcommand{\printedpage}[1]{\marginpar{\footnotesize p.~#1}}%
\newcommand{\msp}[1]{\marginpar{\footnotesize [#1]}}%
\newcommand{\limind}{\mathop{\underrightarrow{\mathrm{lim}}}}%
\newcommand{\limproj}{\mathop{\underleftarrow{\mathrm{lim}}}}%
\newcommand{\M}{\mathcal{M}}%
\newcommand{\Fl}{\mathrm{Fl}}%
\newcommand{\Ob}{\mathrm{Ob}}%
\newcommand{\Cof}{\mathrm{Cof}}%
\newcommand{\Fib}{\mathrm{Fib}}%
\newcommand{\univ}{\mathrm{univ}}%
\newcommand{\HO}{\mathrm{HO}}%
\makeatother\let\printedpage\derpp\let\msp\dermsp%
\let\mspin\DERmspin%


\printedpage{24}
% head: A. Grothendieck : Dérivateurs, Ch. XIII, Édition des Universités de Montpellier II et Paris VII
to deduce a similar factorization for the arrows $f$ of $\M(I)$, as
\[
\begin{aligned}
f = pi \quad \text{with}\quad & i \in TC(I), && p \in F(I)\\
\text{or}\quad & i \in C(I), && p \in TF(I),
\end{aligned}
\]
where $C(I)$, $TC(I)$, $F(I)$, $TF(I)$ are here defined via $C$, $TC$, $F$, $TF$ ``argument by argument''. But the trouble is that by proceeding in this naive way, we will no longer necessarily have $C(I) \longleftrightarrow TF(I)$, $TC(I) \longleftrightarrow F(I)$. We have to shrink $TF(I)$ and $F(I)$ considerably, to have an RLP property with respect to $C(I)$ resp. $TC(I)$.

So I can hardly see any other way out, in this general context, than by assuming the stability of $W$ under filtering $\limind$ (which is after all not very nasty).

\noindent\msp{page 42}

Willy-nilly, we come back to the situation considered in paragraph 2. We are going to restate the essentials of what was sensed then:

\textbf{Theorem 2} (still a little conjectural). {\itshape Let $(\M,W)$ be a category of models, with $W$ (as always) weakly saturated, but in addition we assume $\M$ stable under small $\limind$ and under finite $\limproj$, and $W$ stable under filtering $\limind$, finally $\M$ accessible. Consider the following conditions:

\begin{enumerate}
\item[\textup{a)}] There exists a weak Quillen structure for $W$, left exact, right closed, $(C,TC,F,TF)$.
\item[\textup{a$'$)}] As in \textup{a)}, but with in addition $C$, $TC$ left closed (\textup{i.e.} we have a closed weak Quillen structure), and $C$, $TC$ of the form $\widetilde{C_0}$, $\widetilde{TC_0}$, where $C_0$ and $TC_0$ are small.
\item[\textup{a$''$)}] There exist $C_0 \subset \Cof_W$, $TC_0 \subset W^{\univ}$, small subsets, such that $C_{0*} \subset W$, $(TC_0)_* \subset \Fib_W$.
\item[\textup{b)}] We have
\[
(\Cof_W)_* \;\subset\; W_{\univ}, \qquad\quad (W^{\univ})_* \;\subset\; \Fib_W
\]
(whence it follows, $\Cof_W$ and $W^{\univ}$ being left closed, that
\[
(W_{\univ})^* \subset \Cof_W, \qquad\quad (\Fib_W)^* \subset W^{\univ} \qquad\quad \text{\itshape ).}
\]
\end{enumerate}

\noindent\msp{page 43}

These conditions imply that for every $I$ in \textup{Cat}, there exists a left exact and closed weak Quillen structure for $W(I)$ in $\M(I)$. Such a structure is defined, in terms of a similar structure $C$, $TC = C \cap W$, $F = (TC)_*$, $TF = C_*$ for $W$, by}
\[
C(I) \;=\; \Bigl\{ f \in \Fl(\M(I)) \Bigm| \alpha_s^*(f) \in C \ \ \forall\, s \in \Ob I \Bigr\},
\]

\printedpage{25}
% head: A. Grothendieck : Dérivateurs, Ch. XIII, Édition des Universités de Montpellier II et Paris VII
\noindent\textit{then}
\[
\begin{aligned}
TC(I) \;&=\; C(I)\cap W(I) \;=\; \Bigl\{ f \in \Fl(\M(I)) \Bigm| \alpha_s^*(f) \in TC \ \ \forall\, s \in \Ob I \Bigr\}\\
F(I) \;&=\; (TC(I))_*\\
TF(I) \;&=\; C(I)_*\,.
\end{aligned}
\]

\textbf{Corollary-Scholie.} \textit{These conditions imply that for every $I$ in} Cat, \textit{$(\M(I),W(I))$ gives rise to a good theory of $W(I)$-homotopic fibers and cofibers, with $W(I)$-cartesian $W(I)$-exact squares, and $W(I)$-cocartesian $W(I)$-exact squares – hence with long} exact \textit{sequences of cosuspension and of suspension, in the case where $\M$, hence also $\M(I)$, is a category with a zero object.}

\noindent\msp{page 44}

I will take care of the proof of the preceding theorem, as well as of the factorization theorem of paragraph 1, when I have my notes on accessible categories and functors and cardinal filtrations. In the meantime, I already feel like making a little more explicit the formalism one arrives at by means of the hypotheses made above. (I am pretty sure at least that condition $a''$) suffices to imply all the others. And it is indeed in the form $a''$) that these conditions are going to be verified in practice, I presume, notably in the case where $\M = $ Cat.)

\section*{4\quad The variances $u^*$, $u_*$, $u_{!}$ for $(\M(I),W(I))$ (I variable)}

I assume the conditions of the preceding paragraph are satisfied for $(\M,W)$ (cf. theorem p. 42), and choose a closed and left exact weak Quillen structure $(C,TC,F,TF)$ for $W$, whence for every $I$ in Cat

\noindent\msp{page 45}

\noindent a closed and left exact weak Quillen structure $(C(I),TC(I),F(I),TF(I))$ for $W(I)$, in $\M(I)$, described on page 43. Now let
\[
u \;:\; I' \longrightarrow I
\]
be an arrow in Cat, whence an arrow
\[
u^* \;:\; (\M(I),W(I)) \;\longrightarrow\; (\M(I'),W(I')).
\]
We know, since $\M$ is stable under arbitrary $\limind$ and also (being moreover accessible) under arbitrary $\limproj$, that $u^*$ admits left and right adjoint functors,
\[
u_{!},\ u_* \;:\; \M(I') \;\longrightarrow\; \M(I),
\]
\printedpage{26}
% head: A. Grothendieck : Dérivateurs, Ch. XIII, Édition des Universités de Montpellier II et Paris VII
and I propose here to summarize the most immediate properties of $u^*$ and of $u_{!}$, $u_*$ with respect to the weak Quillen structures under consideration.

\textbf{Proposition 1:} \textit{With the preceding notations, we have the relations:}
\[
u^*(C(I)) \subset C(I'), \qquad u^*(TC(I)) \subset TC(I'). \tag*{(1)}
\]
\[
u_*(F(I')) \subset F(I), \qquad u_*(TF(I')) \subset TF(I). \tag*{(2)}
\]
\[
\begin{aligned}
u^*(F(I)) \subset F(I') &\iff u_{!}(TC(I')) \subset TC(I)\\
u^*(TF(I)) \subset TF(I') &\iff u_{!}(C(I')) \subset C(I).
\end{aligned} \tag*{(3)}
\]

On the other hand, I do not know anything positive to say about the inclusions
\[
\begin{aligned}
u_*(C(I')) \;&\subset\; C(I), \;\; u_*(TC(I')) \;\subset\; TC(I)\\
u_{!}(F(I')) \;&\subset\; F(I), \;\; u_{!}(TF(I')) \;\subset\; TF(I).
\end{aligned}
\]
Except that for a $u$ that is at all general ($u$ not an isomorphism, say), $u_*$ has no desire

\noindent\msp{page 46}

\noindent to be compatible with the structures $C$, $TC$ (a property of homotopical right exactness, morally, since $u_*$ is left exact and in general not at all right exact), nor $u_{!}$ with the structures $F$, $TF$ (dual comment to be given, $u_{!}$ is \textit{right} exact, and does not feel like being ``homotopically left exact''). % [sic] exacte -> exact, accord avec « foncteur » / « exact à gauche »

\textsc{Proof} of (1), (2), (3). Relations (1) are obvious by definition. Relations (2), (3) are scarcely less so, in view of the following general sorital property, whose place would be in paragraph 1:

\textbf{Proposition 2.} \textit{Let}
\[
\begin{tikzcd}
C' \arrow[r, shift left, "u_{!}"] & C \arrow[l, shift left, "u^*"]
\end{tikzcd}
\]
\textit{be a pair of adjoint functors $(u_{!},u^*)$, and let $\Phi' \subset \Fl(C')$, $\Psi \subset \Fl(C)$. Then}
\[
(u_{!}(\Phi') \longleftrightarrow \Psi) \qquad\quad \iff \qquad\quad (\Phi' \longleftrightarrow u^*(\Psi)).
\]

\textbf{Corollary.} (This is the same statement, with only the notation changed.) \textit{Suppose that $u^* : C \longrightarrow C'$ admits a} right \textit{adjoint $u_*$ (not necessarily a left adjoint)}
\[
\begin{tikzcd}
C' \arrow[r, shift right, "u_*"'] & C, \arrow[l, shift right, "u^*"']
\end{tikzcd}
\]

\printedpage{27}
% head: A. Grothendieck : Dérivateurs, Ch. XIII, Édition des Universités de Montpellier II et Paris VII
\noindent\textit{let $\Phi \subset \Fl(C)$, $\Psi' \subset \Fl(C')$, then}
\[
(u^*(\Phi) \longleftrightarrow \Psi') \qquad\quad \iff \qquad\quad (\Phi \longleftrightarrow u_*(\Psi')).
\]

This is proposition 2 for \begin{tikzcd} C'_0 \arrow[r, shift left, "u_{0\,!}"] & C_0 \arrow[l, shift left, "u_0^*"] \end{tikzcd}, with $C'_0 = C$, $C_0 = C'$, $(u_{0\,!} : C \longrightarrow C') = u^*$, and $(u_0^* : C' \longrightarrow C) = u_*$.

\noindent\msp{page 47}

For the inclusions (2), we use the corollary of proposition 2 for the pairs $\Phi = TC(I)$, $\Psi' = F(I')$, resp. $\Phi = C(I)$, $\Psi' = TF(I')$ \texttt{[rather} $\Psi' = TF(I')$\texttt{]}, and apply (1) which gives $u^*(\Phi) \longleftrightarrow \Psi'$, whence $\Phi \longleftrightarrow u_*(\Psi')$, which is nothing other than (2). For the equivalences (3), we use proposition 2, putting $\Psi = F(I)$, $\Phi' = TC(I')$ resp. $\Psi = TF(I)$, $\Phi' = C(I')$.
\DERmark{ENV0603}


\textbf{Corollary 1} (of proposition 1): \textit{Suppose $u : I' \longrightarrow I$ is such that for every $i \in \Ob I$, the category $i\backslash I'$ is discrete} (\footnote[27]{More generally, it suffices that every connected component $Z_\alpha$ of $i\backslash I'$ have an initial object $e_\alpha$, so that $\limind\limits_{i\backslash I'} X'(i') \simeq \coprod_{\alpha\in\pi_0(i\backslash I')} X'(e_\alpha)$. This is the case in particular if $I' \longrightarrow I$ is a local isomorphism, e.g. an open immersion.}) \textit{(e.g. $I'$ discrete). Then we have}
\[
\left\{
\begin{aligned}
u_{!}(C(I')) \;&\subset\; C(I), \;\; u_{!}(TC(I')) \;\subset\; TC(I)\\
u^*(F(I)) \;&\subset\; F(I'), \;\; u^*(TF(I)) \;\subset\; TF(I').
\end{aligned}
\right.
\]

Indeed, the relations on the second line are equivalent (because of (3)) to those on the first. We have to prove that if $f$ in $\Fl\M(I')$ is in $C(I')$ resp. $TC(I')$, $u_{!}(f)$ is in $C(I)$ resp. $TC(I)$, which also means that $\alpha_s^*(u_{!}(f))$ is in $C$ resp. $TC$, for every $s \in \Ob I$. Now the functor
\[
\alpha_s^* u_{!} \;:\; \M(I') \;\longrightarrow\; \M
\]
is given explicitly by the well-known formula
\[
\alpha_s^*(u_{!}(X')) \;\overset{\mathrm{def}}{=}\; u_{!}(X')(s) \;\simeq\; \limind\limits_{i'\in s\backslash I'} X'(i').
\]
\DERmark{ENV0604}


Since

\noindent\msp{page 48}

\noindent $s\backslash I'$ is discrete, we therefore have the isomorphism functorial in $X'$ in $\M(I')$
\[
\alpha_s^*(u_{!}(X')) \;\simeq \coprod_{(s',\, x \,:\, s \,\longrightarrow\, u(s')) \text{ in } \Ob\,(s\backslash I')} X'(s'),
\]
so if $f \,:\, X' \longrightarrow Y'$, $\alpha_s^*(u_{!}(f))$ is made explicit as a direct sum in $\underline{\mathrm{Fl}}(\M)$ of arrows of the form
\[
f_{s'} \;:\; X'(s') \;\longrightarrow\; Y'(s'),
\]
\printedpage{28}
% head: A. Grothendieck : Dérivateurs, Ch. XIII, Édition des Universités de Montpellier II et Paris VII
which are in $C$ resp. $TC$ since $f$ is in $C(I')$ resp. $TC(I')$. As $C$, $TC$ are stable under arbitrary direct sums ($C$, $TC$ being left closed), we win.

\textbf{Corollary 2.} \textit{Let $I$ be in} Cat, \textit{and $f \in \Fl(\M(I))$. For $f$ to be in $F(I)$} (resp. \textit{in $TF(I)$\,), it} is necessary \textit{that $\forall\, s \in \Ob I$, $f_s : X(s) \longrightarrow Y(s)$ be so.} (But this condition does not look sufficient at all!) \textit{In particular, let $X \in \Ob\M(I)$, then $X \in \Ob\M(I)_f \Longrightarrow X(s) \in \Ob\M_f \ \ \forall\, s \in \Ob I$.}

We apply corollary 1 to $\alpha_s : e \longrightarrow I$.

\textbf{Scholie.} Proposition 1 (relation (1)) shows that the functor $u^*$ transforms $W$-cocartesian squares into likewise, long exact sequences of suspension into likewise – this is ``the homotopical right exactness'' of $u^*$. We expect the same to hold for cartesian squares and long exact sequences

\noindent\msp{page 49}

\noindent of cosuspension, but to prove it some work will be needed, which we will do in what follows.

Relations (2) show that the functor $u_*$ transforms $W$-cartesian squares into likewise, hence long exact sequences of cosuspension into likewise. I do \textit{not} expect the same to hold in general for $W$-cocartesian squares – \textit{i.e.} that $u_*$ be ``homotopically right exact'', and not only left exact. But I should at least do some examples (with $I'$, $I$ finite ordered categories).

As for $u_{!}$, the proposition says nothing very tangible about it. Nevertheless we expect $u_{!}$ to transform $W$-cocartesian squares into likewise – that is the least one can expect! Or perhaps I am mistaken, in confusing the functors $u_*$, $u_{!}$ considered here, with the similar functors
\[
\begin{tikzcd}
\HO_W(I') \arrow[r, shift left, "u_{!}^W"] \arrow[r, shift right, "u_*^W"'] & \HO_W(I),
\end{tikzcd}
\]
left resp. right adjoints to
\[
\begin{tikzcd}
\HO_W(I') & \HO_W(I), \arrow[l, "u_W^*"']
\end{tikzcd}
\]
induced by
\[
u^* \;:\; (\M(I),W(I)) \;\longrightarrow\; (\M(I'),W(I')).
\]
But in general $u_*$, $u_{!}$ are not compatible

\noindent\msp{page 50}

\noindent with the localizers $W(I')$, $W(I)$, so they must not be confused with the functors $u_{!}^W$, $u_*^W$. So perhaps I was hasty, in what I said about $u_*$ and $u_{!}$. I will be obliged to study the situation very closely, when I try to construct the functors $u_*^W$ and $u_{!}^W$. For the moment, I still have no assurance that the conditions considered for $(\M,W)$ suffice to ensure the existence of $u_*^W$, $u_{!}^W$ – not even for $I$, $I'$ finite ordered sets!

\textbf{Corollary 1.} \textit{Let $X \in \M(I)_{\mathrm{f}}$, then for every $\alpha : s \longrightarrow t$ in $I$ which is} \texttt{[a]} monomorphism, \textit{$X(\alpha) : X(t) \longrightarrow X(s)$ is in $F$.}

% foot: Version 25 janvier 2004


\endgroup
% END UNIT Der-ChXIII-25-01-04/Der-ChXIII-25-01-04_p026-030.tex
% BEGIN UNIT Der-ChXIII-25-01-04/Der-ChXIII-25-01-04_p031-035.tex
\begingroup\makeatletter%
\newcommand{\printedpage}[1]{\marginpar{\footnotesize p.~#1}}%
\newcommand{\msp}[1]{\marginpar{\footnotesize [#1]}}%
\newcommand{\circled}[1]{\tikz[baseline=(c.base)]\node[draw,circle,inner sep=1pt](c){#1};}%
\makeatother\let\printedpage\derpp\let\msp\dermsp%
\let\mspin\DERmspin%


\printedpage{29}
% head: A. Grothendieck : Dérivateurs, Ch. XIII, Édition des Universités de Montpellier II et Paris VII
(We apply \texttt{[the]} corollary 1 (p. 47) to $\tilde{\alpha} : \boldsymbol{\Delta}_1 \longrightarrow I$.)

\textbf{Proposition 3.} (\footnote[28]{Compare with the lemma p. 63 and its proof (more exactly, the corollary of the lemma).}) \textit{Let $X$ be in $\mathcal{M}(\boldsymbol{\Delta}_1)$, $X = (X_0 \xleftarrow{\alpha} X_1)$. Then $X$ is in $\mathcal{M}(\boldsymbol{\Delta}_1)_{\mathrm{f}}$ if and only if $X_0$, $X_1$} \texttt{[are]} \textit{in $\mathcal{M}_{\mathrm{f}}$ and $\alpha \in F$.}

\textbf{Corollary 2.} \textit{Let $I' \overset{u}{\hookrightarrow} I$ be} \texttt{[an]} open immersion. \textit{Then for $X \in \operatorname{Ob}\mathcal{M}(I)_{\mathrm{f}}$, $\varprojlim\limits_{I} X(i) \longrightarrow \varprojlim\limits_{I'} X(i)$ is in $F$.}

\section*{5\quad Quillenizable categories of models}

I return to the general case of a category of models $(\mathcal{M},W)$.

% marginal vertical bar printed beside the following Définition 1
\textbf{Definition 1.} We say that $(\mathcal{M},W)$ is \textit{Quillenizable} if there exists an \textit{exact} weak Quillen structure $(C,TC,F,TF)$ for $\mathcal{M}$.

That is a self-dual condition, but one which I know how to verify in practice only by means of either left or right exactness -- thus by a stronger condition which destroys the symmetry. It is to preserve it as far as possible that I introduce the preceding definition, as well as the

\textbf{Definition 2.} Let, in addition to $(\mathcal{M},W)$, Diag $\subset$ Cat be a full subcategory of Cat. We say that
\msp{page 51}$W$ is \textit{Quillenizable on} Diag, if for every $I$ in Diag, $(\mathcal{M}(I),W(I))$ is Quillenizable.

In practice, I will use this definition in the case where Diag $=$ Cat, or else, that it satisfy at least the stability properties useful for the domain of a derivator: stable under finite $\varprojlim$, under finite sums -- possibly under finite $\varinjlim$, etc (\footnote[29]{I will assume that Diag contains at least the categories coming from finite ordered sets.}). % [sic] eventuellement -> éventuellement, missing accent

\textbf{Example 1:} That of paragraphs 3, 4, where we isolated conditions on $(\mathcal{M},W)$, with $\mathcal{M}$ a ``big'' category and $\mathcal{M}$ accessible, for $W$ to be Quillenizable on the whole of Cat. This is the case which seems to me at present the most interesting.

\textbf{Example 2:} Let $\mathcal{A}$ be an abelian category, and $\mathcal{M} = K^{\bullet}(\mathcal A)$ the category of cochain complexes with values in $\mathcal{A}$, $W = W_{\mathcal{A}}$ the set of quasi-isomorphisms, hence % [sic] K^\bullet(A) -> K^\bullet(\mathcal{A}), plain italic A where calligraphic A is used everywhere else
\DERmark{ENV0605}

\[ \mathcal{M}W^{-1} \;=\; K^{\bullet}(\mathcal{A})W_{\mathcal{A}}^{-1} \;\simeq\; D^{\bullet}(\mathcal{A}), \qquad \text{derived category of $\mathcal{A}$.} \]
We also define $\mathcal{M}^{+}$, $\mathcal{M}^{-}$, $\mathcal{M}^{\mathrm{b}}$, with the induced localizers $W_{\mathcal{A}}^{+}$, $W_{\mathcal{A}}^{-}$, $W_{\mathcal{A}}^{\mathrm{b}}$, and the categories of fractions $D^{+}(\mathcal{A})$, $D^{-}(\mathcal{A})$, $D^{\mathrm{b}}(\mathcal{A})$. Or
\msp{page 52}again, for $I$ in Cat,
\[ \mathcal{M}(I) \;=\; K^{\bullet}(\mathcal{A}(I)) \]
% foot: Version 25 janvier 2004
\printedpage{30}
% head: A. Grothendieck : Dérivateurs, Ch. XIII, Édition des Universités de Montpellier II et Paris VII
\[ W(I) \;=\; \underbrace{W_{\mathcal{A}(I)}}_{\text{quasi-isomorphisms}} \;\subset\; \mathrm{Fl}(K^{\bullet}(\mathcal{A}(I))), \]
and likewise for $\mathcal{M}^{\alpha}(I)$, $W^{\alpha}(I)$ $(\alpha \in \{+,-,\mathrm{b}\})$, at least \textit{if $I$ is finite.} There are also the subcategories $K^{\geq 0}(\mathcal{A})$ and $K^{\leq 0}(\mathcal{A})$ of cochain resp. chain complexes in the strict sense (with $K^{i} = 0$ if $i < 0$, resp. if $i > 0$ \texttt{[rather if $i < 0$, resp. if $i > 0$]}). We will have for arbitrary $I$
\DERmark{ENV0606}

\[ \left\{ \begin{aligned} K^{\geq 0}(\mathcal{A})(I) &\;\simeq\; K^{\geq 0}(\mathcal{A}(I)) \\ W_{\mathcal{A}}^{\geq 0}(I) &\;=\; W_{\mathcal{A}(I)}^{\geq 0}, \end{aligned} \right. \]
and likewise for $K^{\leq 0}$, $W^{\leq 0}$.

According to \textsc{Quillen} page 1.2, $(K^{+}(\mathcal{A}),W_{\mathcal{A}}^{+})$ is Quillenizable if $\mathcal{A}$ has enough injectives, $(K^{-}(\mathcal{A}),W_{\mathcal{A}}^{-})$ is so if $\mathcal{A}$ has enough projectives. He says nothing for $(K^{\bullet}(\mathcal{A}),W_{\mathcal{A}})$ without more -- yet one does cohomological algebra with $D^{\bullet}(\mathcal{A})$ without more, and even without making an assumption of the existence of injectives or projectives. I have somewhat forgotten what \textsc{Verdier} did on this subject -- I should all the same look at his thesis (!), or else the exposition of \textsc{Hartshorne}.

\msp{page 53}At any rate, we have here examples of a quite different type from those considered in 1°. Thus, there is no need to assume that $\mathcal{M}$ is stable under small $\varinjlim$ -- $\mathcal{M}$ can even be a small category, it is so if and only if $\mathcal{A}$ itself is. If $\mathcal{A}$ is stable under (small) filtering $\varinjlim$, and is accessible (e.g. has a small generating family, and filtering $\varinjlim$ are exact), then $K^{\bullet}(\mathcal{A})$ etc. (\footnote[30]{!!}) have the same properties, and $(K^{+}(\mathcal{A}),W_{\mathcal{A}}^{+})$ is Quillenizable on the whole of Cat (\footnote[31]{careful, $K^{+}(\mathcal{A})$ is \textit{not} stable under infinite filtering $\varinjlim$ -- so one cannot apply the carpet of paragraphs 2, 3 directly to $K^{+}$.}). (This is the case where we know that $\mathcal{A}$ has enough injectives, and cohomological algebra is done mostly by dint of injective resolutions \dots) In this case we know that the % [sic] elle même -> elle-même, missing hyphen
\[ u_{*} \;:\; D^{+}(I',\mathcal{A}) \longrightarrow D^{+}(I,\mathcal{A}) \]
exist as right adjoints of the
\[ u^{*} \;:\; D^{+}(I,\mathcal{A}) \longrightarrow D^{+}(I',\mathcal{A}), \]
but I am no longer too sure whether we know how to define $u_{*}$ also on all of $D(I',\mathcal{A})$, and I doubt that $u_{!}$ exists without an assumption on $u$ (of the sort: finite cohomological dimension). Probably we must always have
\[ u_{!} \;:\; D^{-}(I',\mathcal{A}) \longrightarrow D^{-}(I,\mathcal{A}), \]
\msp{page 54}left adjoint to
\[ u^{*} \;:\; D^{-}(I,\mathcal{A}) \longrightarrow D^{-}(I',\mathcal{A}). \]
% foot: Version 25 janvier 2004

\printedpage{31}
% head: A. Grothendieck : Dérivateurs, Ch. XIII, Édition des Universités de Montpellier II et Paris VII
But I no longer really know whether we should expect to be able to define $u_{!}$ on all of $D(I',\mathcal{A})$, apart from assumptions like finite cohomological dimension. The sad truth is that I have $\pm$ forgotten the theory of derived functors! (I ought to ask \textsc{Jean Malgoire} again to send me canonical texts.)

\textbf{Scholie.} If $(\mathcal{M},W)$ is Quillenizable on Diag $\subset$ Cat (and even without its being Quillenizable), it defines a prederivator $\mathbf{D}_{\mathcal{M},W}$ on Diag. We will sometimes set
\[ \mathrm{HO}_{\mathcal{M},W}(I) \;\text{ or }\; \mathrm{HO}_{W}(I) \;=\; \mathbf{D}_{\mathcal{M},W}(I). \]
The Quillenizable hypothesis implies, on each of these homotopy categories, the (relatively weak) ``triangulated'' structure, coming from the \textit{exact} suspension and cosuspension sequences (\footnote[32]{In the case where $\mathcal{M}$ is pointed by \texttt{[a]} neutral object. Otherwise, there is the structure associated to cartesian and cocartesian squares. It would be quite interesting to axiomatize such a species of structure.}). (The Quillenizable character ensures the exactness of these sequences \dots) The functors $u^{*} : \mathrm{HO}_{W}(I) \longrightarrow \mathrm{HO}_{W}(I')$
\msp{page 55}are ``homotopically exact'' in the sense of these ``triangulated'' (noncommutative) structures, \textit{i.e.} transform exact suspension resp. cosuspension sequences into likewise. % [sic] des ces -> de ces, grammar

Which axioms of derivators are satisfied under the Quillenizable hypothesis alone? \textbf{Der} 1 \texttt{[see chapter I, p. 43]} is so without any hypothesis at all on the category of models $(\mathcal{M},W)$:
\[ \mathrm{HO}_{W}(\coprod_{\alpha} I_{\alpha}) \;\simeq\; \prod_{\alpha} \mathrm{HO}_{W}(I_{\alpha}) \]
for a \textit{finite} family of objects $I_{\alpha}$ in Diag. The localization axiom \textbf{Der} 2 \texttt{[see chapter I, p. 43]} is so under the hypothesis that $W$ be \textit{strongly saturated,} and Quillenizable. We use the fact that every arrow of $\mathrm{HO}_{W}(I)$ comes, up to isomorphism, from an arrow in $\mathcal{M}(I)$. (That was one of the points which remained pending in VI, in the theory of $\mathrm{HOT}_{W}$, $W \subset \mathrm{Fl}(\mathrm{Cat})$ a basic localizer \dots)

The axiom \textbf{Der} 3, new version \texttt{[see chapter I, pp. 111-112]}, involves a basic localizer $\mathbf{W} \subset \mathrm{Fl}(\mathrm{Cat})$ (not necessarily $\mathbf{W}_{\infty}$) (\footnote[33]{It is natural here to assume that $\mathbf{W}$ satisfies all the ``good axioms'' on basic localizers W (1, 2, 3, 4, 5, 6, 7) \texttt{[see chapter XII]}.}). The strongest form would be this: if
\[ u \;:\; I' \longrightarrow I \]
\msp{page 56}is in $\mathbf{W}$, we want the functor $u^{*}$ to induce an \textit{equivalence}
\[ u_{\mathrm{lc}}^{*} \;:\; \mathrm{HO}_{W}^{\mathrm{lc}}(I) \xrightarrow{\;\sim\;} \mathrm{HO}_{W}^{\mathrm{lc}}(I'). \]
But I doubt that we are equipped to verify so strong an axiom, without having at least a theory of the operations $u_{!}$ or $u_{*}$. One could on the off chance posit this axiom by taking $\mathbf{W} = \mathbf{W}_{\alpha} \overset{\mathrm{def}}{=}$ the smallest basic localizer satisfying W (1 to 7) (the largest of
\DERmark{ENV0607}

% foot: Version 25 janvier 2004
\printedpage{32}
% head: A. Grothendieck : Dérivateurs, Ch. XIII, Édition des Universités de Montpellier II et Paris VII
these localizers being $\mathbf{W}_{\infty}$). For the moment, the way to verify this axiom (in terms of the datum $(\mathcal{M},W)$) is for me a $\pm$ total mystery. (But I know how to verify it for $\mathrm{HOT}_{W}$, when $W \subset \mathrm{Fl}(\mathrm{Cat})$ satisfies W (1, 2, 3, 4, 5, \dots).)

I am going to look into it, once I have a reasonable command of the formalism of $u_{*}$, $u_{!}$ relative to the $u : I' \longrightarrow I$ in Diag. The Quillenizable hypothesis alone seems too coarse to ensure the existence of $u_{!}$ or of $u_{*}$ in interesting cases. % [sic] u_\star -> u_*, star symbol where u_* is used elsewhere
\DERmark{ENV0608}


\noindent\msp{page 57}

\section*{6\quad Quillenizable category of models, and functors $u_{!}$, $u_{*}$}

\textbf{Definition 1:}

\begin{itemize}[leftmargin=2.5em]
\item[\circled{a}] We say that $(\mathcal{M},W)$ is \textit{uniformly Quillenizable on the left} on Diag, if we can find a weak Quillen structure $(C,TC,F,TF)$ for $\mathcal{M}$ which is \textit{closed,} and which moreover has the following property: for $I$ in Cat, let $C(I)$, $TC(I) \subset \mathrm{Fl}(\mathcal{M}(I))$ be defined by
\[ \begin{aligned} &(u : X \longrightarrow Y) \in C(I) \text{ (resp. } TC(I)) \text{ if and only if } \ \forall\, s \in \operatorname{Ob} I, \\ &u_{s} : X_{s} \longrightarrow Y_{s} \text{ is in } C \text{ (resp. } TC). \end{aligned} \]
\DERmark{ENV0609}

Let moreover
\[ F(I) \;=\; TC(I)_{*}, \qquad TF(I) \;=\; C(I)_{*}. \]
The condition in view is that for every $I$ \textit{in} Diag, $(C(I),TC(I),F(I),TF(I))$ is a \textit{closed exact} weak Quillen structure for $W(I)$ (\footnote[34]{$(C,TC,F,TF)$ is then called a \textit{uniform left Quillenization} of $(\mathcal{M},W)$ on Diag.}).

\item[\circled{b}] We say that $(\mathcal{M},W)$ is \textit{Quillenizable on Diag uniformly and exactly} \texttt{[on the]} \textit{left}, if we have in addition $TC = C \cap W$, \textit{i.e.} if $(C,TC,F,TF)$ is left exact, hence defined by $C$ alone, via
\[ TC = C \cap W, \qquad F = (TC)_{*}, \qquad TF = C_{*} \qquad \text{(\footnotemark[35])} \]
\footnotetext[35]{We say that $(C,\ldots)$, or simply $C$, defines a \textit{uniform and left exact Quillenization} of $(\mathcal{M},W)$ on Diag.}
\end{itemize}

We saw in paragraph 3 that if $\mathcal{M}$ is stable under small $\varinjlim$ and accessible, and $W$ stable under small filtering $\varinjlim$, finally if
% foot: Version 25 janvier 2004
\printedpage{33}
% head: A. Grothendieck : Dérivateurs, Ch. XIII, Édition des Universités de Montpellier II et Paris VII
\msp{page 58}$W$ admits a weak Quillen structure $(C,TC,F,TF)$ which is closed and left exact -- then the conditions of definition 1 are satisfied.

On the other hand, let us note that proposition 1 page 45 and its corollaries, including proposition 3 (page 50) and its corollaries, are valid provided that the Quillenizations of the $(\mathcal{M}(I),W(I))$ under consideration are defined ``uniformly on the left'' by a uniform left Quillenization (not necessarily left exact) of $(\mathcal{M},W)$ on Diag.

Uniform left Quillenization will serve us to establish the existence of homotopy direct images $u_{*}$, for certain $u$. Uniform right Quillenization, dually, will have as its main consequence the existence of homotopy $u_{!}$ for certain $u : I' \longrightarrow I$ in Diag.

I am going to summarize what I have been able to establish so far in this direction, in the

\msp{page 59}\textbf{Theorem-scholie:} \textit{Let $(C,TC,F,TF)$ be a uniform left Quillenization of $W$ on} Diag $\subset$ Cat.

\begin{itemize}[leftmargin=3em]
\item[1)] \textit{Let $u : I' \longrightarrow I$ be an arrow in} Diag. \textit{We assume} Diag \textit{matched to $\mathcal{M}$ in such a way that for $u$, an arrow in} Diag, \textit{the $\varprojlim$ in $\mathcal{M}$ (of type $I'/i$) required to ensure the existence of}
\[ u_{*} \;:\; \mathcal{M}(I') \longrightarrow \mathcal{M}(I) \]
\textit{exist in $\mathcal{M}$. (This is O.K. if $I$, $I'$} \texttt{[are]} \textit{finite, $\mathcal{M}$ being assumed stable under finite $\varprojlim$ in any case.) We know} (prop. 1, page 45) \textit{that $u_{*}$ induces}
\DERmark{ENV0610}

\[ u_{*}^{\mathrm{f}} \;:\; \mathcal{M}(I')_{\mathrm{f}} \longrightarrow \mathcal{M}(I)_{\mathrm{f}}\,. \tag*{$(*)$} \]
\textit{We say that $u$ is} admissible (on the left) for the given uniform left Quillenization of $W$ on Diag, \textit{if $u_{*}^{\mathrm{f}}$ is compatible with the localizers $W(I')_{\mathrm{f}}$, $W(I)_{\mathrm{f}}$,} i.e. \textit{if $f' : X' \longrightarrow Y'$ in $\mathcal{M}(I')$, $X'$, $Y'$ fibrant, $f' \in W(I')$ implies $u_{*}(f') \in W(I)$. We say that $Z$ in} Diag \textit{is} admissible (for the given uniform left Quillenization of $W$ on Diag \dots) \textit{if the functor $Z \longrightarrow e$ is admissible. This being so:}
\begin{itemize}[leftmargin=2.5em]
\item[a)] \textit{For $u$ to be admissible, it suffices that the $I'/s$ for $s \in \operatorname{Ob} I$, be so. NB there is a converse, when $I$ is ordered, hence the $I'/s \longrightarrow I'$ open immersions.}
\item[b)] \msp{page 60}\textit{Every finite ordered set is admissible. More generally (barring error) every ordered set $I$ such that the topological space associated to $I^{\circ}$ is noetherian (}i.e. \textit{every increasing sequence of open sets of $I^{\circ}$ is stationary, or again every increasing sequence of closed sets of $I$ is stationary), and if moreover every irreducible closed subset of $I^{\circ}$ admits a generic point,} i.e. \textit{$I^{\circ}$ is (not only a noetherian space, but also) a topological space which is} sober. \textit{(For $I$, this means that every upward filtering open subset of $I$ admits a greatest element.)} % [sic] fermées -> fermés, agreement with masculine (ouverts / fermés)
\end{itemize}
\end{itemize}
% foot: Version 25 janvier 2004


\endgroup
% END UNIT Der-ChXIII-25-01-04/Der-ChXIII-25-01-04_p031-035.tex
% BEGIN UNIT Der-ChXIII-25-01-04/Der-ChXIII-25-01-04_p036-040.tex
\begingroup\makeatletter%
\newcommand{\printedpage}[1]{\marginpar{\footnotesize p.~#1}}%
\newcommand{\msp}[1]{\marginpar{\footnotesize [#1]}}%
\newcommand{\cM}{\mathcal{M}}%
\newcommand{\HO}{\mathrm{HO}}%
\newcommand{\Hom}{\mathrm{Hom}}%
\newcommand{\Ob}{\mathrm{Ob}}%
\newcommand{\rf}{\mathrm{f}}%
\newcommand{\rc}{\mathrm{c}}%
\makeatother\let\printedpage\derpp\let\msp\dermsp%
\let\mspin\DERmspin%


\printedpage{34}
% head: A. Grothendieck : Dérivateurs, Ch. XIII, Édition des Universités de Montpellier II et Paris VII

\begin{itemize}[leftmargin=3em]
\item[c)] \textit{In particular, if $I'$, $I$ are finite ordered sets, then every $u : I' \longrightarrow I$ is admissible. The same holds, more generally, if $I$, $I'$ are ordered sets and $I'^{\circ}$ is a noetherian sober space.}
\end{itemize}

\begin{itemize}[leftmargin=2em]
\item[2)] \textit{Suppose $u : I' \longrightarrow I$ admissible. Consider}
\[
u_*^W \ : \ \underbrace{\cM(I')_{\rf}W(I')_{\rf}^{-1}}_{\simeq\, \HO_W(I')} \ \longrightarrow \ \underbrace{\cM(I)_{\rf}W(I)_{\rf}^{-1}}_{\HO_W(I)},
\]
\textit{induced by $u_*$. This functor is right adjoint to $u_W^* : \HO_W(I) \longrightarrow \HO_W(I')$.}

\msp{page 61}
\textit{Moreover, the canonical homomorphism (for $s \in \Ob I$, $X' \in \Ob \HO_W(I')$)}
\[
u_*^W(X')(s) \ \simeq \ p_{s\,*}(X'|I'/s) \ \overset{\mathrm{def}}{=} \ \underset{I'/s}{\underleftarrow{\mathrm{holim}}}{}_W(X'|I'/s)
\]
\DERmark{ENV0611}

\textit{(where $p_s : I'/s \longrightarrow e$) is an isomorphism, if we suppose that the $I'/s$ are admissible (we shall then say that $u$ is} strictly \textit{admissible).}
\end{itemize}

\textbf{Corollary.} \textit{On the domain formed by finite ordered sets, the prederivator} $\HO_W$ \textit{satisfies (in addition to axioms} \textbf{Der}~1, \textbf{Der}~2\textit{) axioms} \textbf{Der}~4 \textit{and} \textbf{Der}~5 \textit{(existence of the direct images $u_*$, right adjoints of the functors $u^*$, and standard calculation of the fibers of a direct image).}

\textsc{Proof} of the theorem-scholie.

$1^\circ$) \textit{a}) is immediate. For \textit{b}), we proceed by induction on the cardinality $n$ of the finite ordered set $I$, the cases $n \in \{0,1\}$ being trivial. Suppose then $n \geq 2$, and the theorem proved for the $n' < n$. I distinguish two cases: if $I$ has a greatest element $a$, then, denoting by $\alpha : e \longrightarrow I$ the functor with value $a$, by $p : I \longrightarrow e$ the projection, we know that we have
\[
p_* \ \simeq \ \alpha^*, \qquad \textit{i.e.} \ \varprojlim_{I} X_i \simeq X_a,
\]
whence the fact that on the whole of $\cM(I)$
\msp{page 62}
(not only on $\cM(I)_{\rf}$), $p_*$ is compatible with the localizers $W(I)$, $W$.

If $I$ has no greatest element, then the set $M$ of its maximal elements has cardinality $\geq 2$, we can therefore write it $M = M_1 \amalg M_2$, and we will have % [sic] des ses -> de ses, typographical slip
\[
I \ = \ I_1 \cup I_2, \qquad \text{where } I_1 = I/M_1,\ I_2 = I/M_2
\]
are two open subsets of $I$, different from $I$, hence of cardinalities $n_1$, $n_2 \ < \ n$. Consider also $I_0 = I_1 \cap I_2$, it is an open subset of cardinality $< n$. The diagram
\[
\begin{tikzcd}[row sep=small]
 & I_1 \arrow[dr, dotted] & \\
I_0 \arrow[ur] \arrow[dr] & & I = I_1 \cup I_2 \\
 & I_2 \arrow[ur, dotted] &
\end{tikzcd}
\]
% foot: Version 25 janvier 2004
\printedpage{35}
% head: A. Grothendieck : Dérivateurs, Ch. XIII, Édition des Universités de Montpellier II et Paris VII
corresponds to an arrow
\[
\begin{tikzcd}[row sep=small, column sep=small]
 & & & 1 & \\
u \ : \ I \arrow[r] & \Psi = 0 \arrow[urr, start anchor=east, end anchor=south west] \arrow[drr, start anchor=east, end anchor=north west] & & & , \\
 & & & 2 &
\end{tikzcd}
\]
with
\[
I_1 \ = \ I/1, \qquad I_2 \ = \ I/2, \qquad I_0 \ = \ I/0.
\]
The induction hypothesis and \textit{a}) imply that $u$ is admissible. As a composite of admissible functors is admissible (tautologically, via $(vu)_* \simeq v_*u_*$), we are finally reduced to the
\msp{page 63}

\textbf{Lemma.} \textit{The category} \begin{tikzcd}[row sep=small, column sep=small, baseline=(\tikzcdmatrixname-2-1.base)]
 & 1 \\
\Psi = 0 \arrow[ur] \arrow[dr] & \\
 & 2
\end{tikzcd} \quad \textit{is admissible.}

This will follow from the

\textbf{Corollary} (\footnote[36]{\textbf{NB} There is an obvious variant, characterizing the $X$ such that $X \longrightarrow e$ is in $TF(\Psi)$ as those for which $X_0$, $X_1$, $X_2$ are fibrant and \textit{trivial,} and $\alpha_1$, $\alpha_2 \in TF$.}). \textit{Let} \begin{tikzcd}[row sep=small, column sep=small, baseline=(\tikzcdmatrixname-2-1.base)]
 & X_1 \arrow[dl, "\alpha_1"'] \\
X_\bullet = X_0 & \\
 & X_2 \arrow[ul, "\alpha_2"]
\end{tikzcd} \quad \textit{be in $\cM(\Psi)$. For $X_\bullet$ to be fibrant, it is necessary and sufficient that we have}

\begin{itemize}[leftmargin=3em]
\item[a)] \textit{$X_0$, $X_1$, $X_2$ are fibrant, and}
\item[b)] \textit{$\alpha_1$, $\alpha_2 \in F$.}
\end{itemize}

The necessity of \textit{a}) is a special case of corollary 2, page 48, that of \textit{b}) is a special case of corollary 1 page 50. For sufficiency (which we shall not need for the theorem), we know how to show that in the commutative diagram in solid lines
\[
% Audit resolved: PDF leaf 7 confirms interrupted vertical sides A_0 -> X_0 and A_2 -> X_2, obscured by the displaced B labels; the printed geometry is retained. See corrections.csv.
\begin{tikzcd}
A_1 \arrow[r] \arrow[dr, "TC"] \arrow[dd] & A_0 \arrow[dr, "TC"] \arrow[d, dash] & A_2 \arrow[l] \arrow[dr, "TC"] \arrow[d, dash] & \\
 & B_1 \arrow[r] \arrow[dl, dotted, "h_1"'] \arrow[d] & B_0 \arrow[dl, dotted, "h_0"'] \arrow[d] & B_2 \arrow[l] \arrow[dl, dotted, "h_2"'] \\
X_1 \arrow[r, "\alpha_1"'] & X_0 & X_2 \ , \arrow[l, "\alpha_2"] &
\end{tikzcd}
\]
with three arrows marked $TC$, we can find $h_1$, $h_0$, $h_2$ leaving the diagram commutative (three triangles and two squares to keep commutative). For this, we begin by choosing $h_0$ so as to make the middle vertical triangle commutative -- this is possible since $X_0$ is fibrant, $A_0 \longrightarrow B_0$ in $TC$.

% foot: Version 25 janvier 2004
\printedpage{36}
% head: A. Grothendieck : Dérivateurs, Ch. XIII, Édition des Universités de Montpellier II et Paris VII
\msp{page 64}

Next, we see separately for $h_1$, $h_2$, that we can find them with each satisfying the two commutativity constraints (1 vertical triangle and one tilted square). For $h_1$ e.g. existence is seen from the auxiliary square diagram (solid lines)
\[
\begin{tikzcd}
A_1 \arrow[r, "TC"] \arrow[dd] & B_1 \arrow[dd] \arrow[ddl, dotted, "h_1"'] \arrow[dr, dotted, "\alpha_1^B"] & \\
 & & B_0 \arrow[dl, dotted, "h_0"] \\
X_1 \arrow[r, "F", "\alpha_1^X"'] & X_0 \ \ . &
\end{tikzcd}
\]

\textbf{NB} This is the proof of proposition 3, page 50.

The corollary (the ``necessary'' part) having been seen, the lemma follows from it, for we know that if we have a homomorphism $f$ of cartesian squares $X$, $Y$
\[
\begin{tikzcd}[row sep=small, column sep=small]
 & X_1 \arrow[dl, "p_1^X"'] \arrow[dd, "f_1"'] & & X_3 \arrow[ll] \arrow[dl] \arrow[dd, "f_3"] \\
X_0 \arrow[dd, "f_0"'] & & X_2 \arrow[ll, crossing over, "p_2^X"'] & \\
 & Y_1 \arrow[dl, "p_1^Y"'] & & Y_3 \arrow[ll] \arrow[dl] \\
Y_0 & & Y_2 \ \ , \arrow[ll, "p_2^Y"'] \arrow[from=uu, crossing over, "f_2"] &
\end{tikzcd}
\]
with
\[
p_1^X,\ p_2^X,\ p_1^Y,\ p_2^Y \ \ \text{fibrant \texttt{[belonging to} $F$\texttt{]} (even $p_1^X$ and $p_1^Y$ fibrant is enough),}
\]
and
\[
f_0,\ f_1,\ f_2 \ \in \ W,
\]
then $f_3 \in W$. (This was seen in the theory of $W$-homotopic fiber products \texttt{[see chapter XII, p. 253]}, and the direct verification
\msp{page 65}
is immediate.) Of the relevant initial hypothesis that \begin{tikzcd}[row sep=small, column sep=small, baseline=(\tikzcdmatrixname-2-1.base)]
 & X_1 \arrow[dl] \\
X_0 & \\
 & X_2 \arrow[ul]
\end{tikzcd} and \begin{tikzcd}[row sep=small, column sep=small, baseline=(\tikzcdmatrixname-2-1.base)]
 & Y_1 \arrow[dl] \\
Y_0 & \\
 & Y_2 \arrow[ul]
\end{tikzcd}
are fibrant, we only need to use a little bit, e.g. that $p_1^X$ and $p_1^Y$ are fibrant, which ensures that the two squares $X$, $Y$ are $W$-homotopically cartesian, which is enough then to imply that $f_0$, $f_1$, $f_2 \in W \Longrightarrow f_3 \in W$.

I am going to try to generalize this reasoning, to handle the case of an ordered set $I$ which may be infinite, but such that $I^\circ$ is noetherian and sober, hence every increasing sequence of closed subsets of $I$, and every decreasing sequence of open subsets, is stationary, and moreover every upward filtering open subset has a greatest element.

% foot: Version 25 janvier 2004
\printedpage{37}
% head: A. Grothendieck : Dérivateurs, Ch. XIII, Édition des Universités de Montpellier II et Paris VII

% margin: a vertical bar is printed in the left margin beside the lines "prouver ceci : (*) Si I est tel que ... alors I aussi."
We want to prove that $I$ is admissible, and by ``noetherian induction'' we are reduced to proving this: $(*)$ If $I$ is such that for every open subset $I' \subset I$, with $I' \neq I$, $I'$ is admissible, then $I$ too.

$(*)$ granted, let us prove that $I = I_0$ is admissible.

\msp{page 66}
If it were not, by $(*)$ there would exist $I_1 \subset I_0$, $I_1 \neq I_0$, with $I_1$ not admissible. Likewise, there would be $I_2 \subset I_1$, $I_2 \neq I_1$, not admissible etc. -- we would have constructed a strictly decreasing sequence of open subsets $I_n$ of $I$, contrary to the noetherian character of $I^\circ$.

Let us prove $(*)$, and again distinguish two cases.

\textit{a}) $I^\circ$ is irreducible, hence has a generic point, \textit{i.e.} a least element, \textit{i.e.} $I$ has a greatest element. This case was already treated earlier -- every category with a final object is admissible.

\textit{b}) $I^\circ$ is reducible, hence a union of two closed subsets $I_1^\circ$, $I_2^\circ$ distinct from $I^\circ$, so $I$ is a union of two open subsets $I_1$, $I_2$ distinct from $I$. Then by the hypothesis in $(*)$, $I_1$, $I_2$ and $I_0 = I_1 \cap I_2$ are admissible. We then finish as before, being reduced to the admissibility of \begin{tikzcd}[row sep=small, column sep=small, baseline=(\tikzcdmatrixname-2-1.base)]
 & 1 & \\
\Psi = 0 \arrow[ur] \arrow[dr] & & ! \\
 & 2 &
\end{tikzcd}

\msp{page 67}
Let us prove part \textit{c}) of the theorem-scholie, part $1^\circ$). This follows from \textit{a}), once we have seen that for $s \in I$, $I'/s$ satisfies the same hypotheses as $I'$. Now $I$ being ordered, $I/s$ is an open subcategory of $I$, hence $I'/s$ an open subcategory of $I'$, with $(I'/s)^\circ$ a closed ordered subset of $I'^{\circ}$. The latter being noetherian and sober, the same holds for every closed subset, whence the conclusion.

Let us prove part $2^\circ$) of the theorem-scholie. Consider
\[
\begin{tikzcd}[row sep=large]
\cM(I')_{\rc} \arrow[r, hook, dotted] & \cM(I') & \cM(I')_{\rf} \arrow[l, hook', dotted] \arrow[d, "u_*^{\rf}"'] \\
\cM(I)_{\rc} \arrow[r, hook, dotted] \arrow[u, "u_{\rc}^*"] & \cM(I) \arrow[u, "u^*"] & \cM(I)_{\rf} \ , \arrow[l, hook', dotted]
\end{tikzcd}
\]
where $u^*$ and a fortiori $u_{\rc}^*$ are compatible with $W(I)$, $W(I')$, and by hypothesis $u_*$ is compatible with $W(I')_{\rf}$, $W(I)_{\rf}$. Passing to the categories of fractions, we find

% foot: Version 25 janvier 2004
\printedpage{38}
% head: A. Grothendieck : Dérivateurs, Ch. XIII, Édition des Universités de Montpellier II et Paris VII
\msp{page 68}
\[
\begin{tikzcd}[row sep=large]
\cM(I')_{\rc}W(I')_{\rc}^{-1} \arrow[r, dotted, "\sim"] & \HO_W(I') & \cM(I')_{\rf}W(I')_{\rf}^{-1} \arrow[l, dotted, "\sim"'] \arrow[d, "u_*^W"'] \\
\cM(I)_{\rc}W(I)_{\rc}^{-1} \arrow[r, dotted, "\sim"] \arrow[u, "u_W^*"] & \HO_W(I) & \cM(I)_{\rf}W(I)_{\rf}^{-1} \arrow[l, dotted, "\sim"']
\end{tikzcd}
\]
\DERmark{ENV0612}\DERmark{ENV0613}

\texttt{[rather} $\cM(I')_{\rf}W(I')_{\rf}^{-1}$ \texttt{(resp.} $\cM(I)_{\rf}W(I)_{\rf}^{-1}$\texttt{) instead of} $\cM(I')_{\rf}$ \texttt{(resp.} $\cM(I)_{\rf}$\texttt{)]} and we want to show that $u_W^*$ and $u_*^W$ are adjoint, more precisely find a bifunctorial isomorphism, for $F$ \textit{in} $\cM(I)_{\rc}$ and $G'$ \textit{in} $\cM(I')_{\rf}$
\[
\Hom_W(u^*F, G') \xrightarrow{\ \sim\ } \Hom_W(F, u_*G') \tag*{$(*)$}
\]
(where $\Hom_W$ denotes the Hom in $\HO_W(I')$ resp. $\HO_W(I)$). As $F$ is in $\cM(I)_{\rc}$, $u^*F$ is in $\cM(I')_{\rc}$, and as $G'$ is in $\cM(I')_{\rf}$, $u_*G'$ is in $\cM(I)_{\rf}$. It follows that the two sides of $(*)$ are quotients of the two sides in the canonical bijection
\[
\Hom_{\cM(I')}(u^*F, G') \xrightarrow[\ \sim\ ]{\alpha} \Hom_{\cM(I)}(F, u_*G'). \tag*{$(**)$}
\]
We simply have to show that the latter transforms the relevant equivalence relation $R'$ on the first side,
\msp{page 69}
into the one $R$ on the second.

I say that in any case we have
\[
R \ \leq \ \alpha(R')
\]
\textit{without} any admissibility hypothesis on $u$, \textit{i.e.} that the inverse bijection of $(**)$ gives by passage to quotients a homomorphism % [sic] invese -> inverse, missing letter
\[
\Hom_W(u^*F, G') \longleftarrow \Hom_W(F, u_*G'), \tag*{$(***)$}
\]
necessarily surjective, and moreover functorial in $F$ \texttt{[,} $G'$ \texttt{crossed out]}, variables in $\HO_W(I)$ and $\HO_W(I')$ respectively. Indeed, we define the arrow $(***)$ \texttt{[rather} $(***)$\texttt{]} as the composite
\DERmark{ENV0614}

\[
\begin{tikzcd}[column sep=large]
\Hom_W(F, u_*G') \arrow[r, "\underbrace{\hspace{2.6cm}}_{\text{\footnotesize\begin{tabular}{@{}c@{}}since $u^*$ is\\ compatible with\\ $W(I)$, $W(I')$ on\\ the whole of $\cM(I)$,\\ and we interpret\\ the two sides\\ as Homs\\ in $\cM(I)W(I)^{-1}$\\ and $\cM(I')W(I')^{-1}$.\end{tabular}}}"'] & \Hom_W(u^*F, u^*u_*G') \arrow[dr, "\text{\footnotesize\begin{tabular}{@{}l@{}}composition\\ with the adjunction\\ homomorphism $u^*u_*G' \longrightarrow G'$\end{tabular}}"] & \\
 & & \Hom_W(u^*F, G').
\end{tikzcd}
\]
% foot: Version 25 janvier 2004


\endgroup
% END UNIT Der-ChXIII-25-01-04/Der-ChXIII-25-01-04_p036-040.tex
% BEGIN UNIT Der-ChXIII-25-01-04/Der-ChXIII-25-01-04_p041-045.tex
\begingroup\makeatletter%
\newcommand{\printedpage}[1]{\marginpar{\footnotesize p.~#1}}%
\newcommand{\msp}[1]{\marginpar{\footnotesize [#1]}}%
\newcommand{\M}{\mathcal{M}}%
\newcommand{\Hom}{\mathrm{Hom}}%
\newcommand{\HO}{\mathrm{HO}}%
\newcommand{\Ob}{\mathrm{Ob}\,}%
\newcommand{\Fl}{\mathrm{Fl}\,}%
\newcommand{\Filt}{\mathrm{Filt}}%
\newcommand{\card}{\mathrm{card}\,}%
\newcommand{\rf}{\mathrm{f}}%
\makeatother\let\printedpage\derpp\let\msp\dermsp%
\let\mspin\DERmspin%


\printedpage{39}
% head: A. Grothendieck : Dérivateurs, Ch. XIII, Édition des Universités de Montpellier II et Paris VII
% foot: Version 25 janvier 2004

This therefore implies that the functor on $\HO_W(I)$ $F \longmapsto \Hom_{\HO_{W}(I')}(u^*F,G')$ is, if not representable, at least a term-by-term quotient of the functor represented by $u_*G'$ (\footnote[37]{But it is not clear, without an additional hypothesis on the situation, that this representable functor can be chosen in functorial dependence on $G'$, viewed in $\HO_W(I')$.}).

\msp{page 70}

But suppose now that $u$ is admissible. Then I am going to define the arrow in the original direction $(*)$ (p. 68) as follows: I choose an arrow
\[
i' \ : \ u^*F \longrightarrow F' \qquad \text{with } i' \in TC \subset C \cap W,\ F' \in \Ob \M(I')_{\rf}
\]
(so in fact $F' \in \Ob \M(I')_{\mathrm{cf}}$, since $u^*F$ is cofibrant), and by adjunction the arrow $i'$ corresponds to
\[
i \ : \ F \longrightarrow u_*F'.
\]
I consider the diagram (commutative!)
\[
\begin{tikzcd}[column sep=large, row sep=large]
\Hom(F',G') \arrow[r, "\alpha"] \arrow[d, "{\scriptstyle \lambda:f'\,\longmapsto\, f'\circ i'}"] & \Hom(u_*F',u_*G') \arrow[d, "{\scriptstyle \mu:f\,\longmapsto\, f\circ i}"] \\
\Hom(u^*F,G') \arrow[r, "\text{adj.}"] & \Hom(F,u_*G')\ .
\end{tikzcd}
\]

We note that we still have $F'$ cofibrant (since $u^*F$ is and $i' \in TC \subset C$), so $\Hom_W(F',G')$ is a quotient of $\Hom(F',G')$, and $\lambda$ induces an isomorphism
\[
\Hom_W(F',G') \xrightarrow{\ \sim\ } \Hom_W(u^*F,G') % [sic] u^*F' -> u^*F, the target of lambda is Hom(u^*F,G')
\]
\DERmark{ENV0615}

(since $i' \in TC \subset W$). Since $F'$, $G' \in \Ob \M(I')_{\rf}$ and $u_*^{\rf}$ is compatible with $W(I')_{\rf}$, $W(I)_{\rf}$, it follows that the horizontal arrow $\alpha$
\msp{page 71}
is compatible with the equivalence relations defined by
\[
\Hom(F',G') \longrightarrow \Hom_{W(I')}(F',G') \quad \text{resp.} \ \ \Hom(u_*F',u_*G') \longrightarrow \Hom_{W(I)}(u_*F',u_*G')\ .
\]
And the same holds for the arrow $\mu : f \longmapsto f \circ i$, trivially. So the adjunction arrow $\Hom(u^*F,G') \longrightarrow \Hom(F,u_*G')$ is compatible with the relations $R$, $R'$, which finishes proving that $(**)$, page 68, induces a bijection $(*)$, page 68.

This therefore shows

\begin{enumerate}
\item[1°)] that for a fixed object $G'$ in $\HO_W(I')$ (which we can always suppose, up to isomorphism, to come from $\M(I')_{\rf}$), the functor
\[
F \ \longmapsto\ \Hom_{\HO_{W}(I')}(u^*_W F,G')
\]
is representable -- \textit{i.e.} the functor $u^*_W$ admits a right adjoint, and
\DERmark{ENV0616}

\item[2°)] that this right adjoint is computed as $u_*(G')$ (when the argument $G'$ is taken fibrant) -- since the homomorphism $({*}{*}{*})$, page 69, is \textit{functorial in $F$}. (A posteriori it follows that it is also functorial in $G'$ \dots)
\end{enumerate}
\DERmark{ENV0617}


\printedpage{40}
% head: A. Grothendieck : Dérivateurs, Ch. XIII, Édition des Universités de Montpellier II et Paris VII
% foot: Version 25 janvier 2004

\msp{page 72}

The last statement in 2° of the theorem-scholie follows from the computation of $u_*^W$, taking into account that $I'/s \longrightarrow I'$ is a local isomorphism, so the inverse image functor sends
\[
\M(I')_{\rf} \longrightarrow \M(I'/s)_{\rf},
\]
so if $G'$ is in $\M(I')_{\rf}$,
\[
u_*(G')(s) \ \simeq\ p_{s\,*}(G'|I'/s)
\]
is a computation of $p^W_{s\,*}(G'|I'/s)$.

This completes the proof of the theorem-scholie on direct images.

\textbf{Comments.}

1) This theorem seems to me far from satisfactory, it is a makeshift. I had hoped, with the hypotheses of paragraph 3, that we would have the existence of $u_*^W$ for \textit{every} arrow $u : I' \longrightarrow I$ in Cat. I now doubt that this is the case, but have not yet succeeded in making a counterexample. The most important test cases for the question of admissibility of a category $I$ seem to me the following:

\begin{enumerate}
\item[\textit{a})] $I$ an infinite discrete category (say countable). The question here comes down to this: let $f_i : X_i \longrightarrow Y_i$ be in $W$, with $X_i$, $Y_i$ fibrant, \textit{i.e.} such that $X_i\longrightarrow e$, $Y_i\longrightarrow e$ are in $F$ \texttt{[\textit{i.e.}} $X_i \longrightarrow e$\texttt{,} $Y_i \longrightarrow e$ \texttt{in} $F$\texttt{]}. Then is $f = \prod f_i : \prod X_i \longrightarrow \prod Y_i$ in $W$?

\msp{page 73}

\item[\textit{b})] $I$ the category % [sic] I (upright roman) -> $I$, italic everywhere else for the category
\[
0 \longrightarrow 1 \longrightarrow 2 \longrightarrow 3 \longrightarrow \cdots,
\]
\DERmark{ENV0618}\DERmark{ENV0619}

\textit{i.e.} the ordered category associated with the ordered set $\mathbf{N}$, for the usual order relation $m \leq n$. Thus $\M(I)$ consists of the \textit{inverse} systems in $\M$, indexed by $\mathbf{N}$. One easily proves (proceeding as in the corollary to the lemma p. 63) that $\M(I)_{\rf}$ consists of the inverse systems
\[
X_0 \longleftarrow X_1 \longleftarrow X_2 \longleftarrow \cdots \longleftarrow X_n \xleftarrow{\ p_n\ } X_{n+1} \longleftarrow \cdots
\]
such that the $X_i$ are in $\M_{\rf}$, and the $p_n$ in $F$. The question is therefore: if $X_\bullet$, $Y_\bullet$ are fibrant and $f_\bullet : X_\bullet \longrightarrow Y_\bullet$ is such that $f_i \in W$ for $i \in \mathbf{N}$, then is $\varprojlim f_i : X_\infty \longrightarrow Y_\infty$ in $W$?

If the answer is affirmative for $\mathbf{N}$ (for a $W$ and a given Quillenization), it % [sic] reponse -> réponse, accent missing
is also so for every countable discrete category, as one sees at once. But I do not see how to deduce even \textit{a}) from the axioms, say for a closed Quillen triple (closed model category) such that $\M$ is moreover stable under filtering $\varinjlim$, that these are exact, and that there is
\printedpage{41}
% head: A. Grothendieck : Dérivateurs, Ch. XIII, Édition des Universités de Montpellier II et Paris VII
% foot: Version 25 janvier 2004
\msp{page 74}
a small set of generators in $\M$ (so $\M$ is accessible), that $W$ is stable under filtering $\varinjlim$. I should ask \textsc{Thomason} whether he knows a counterexample. % [sic] connait -> connaît, circumflex missing

\item[\textit{c})] The two preceding examples were ordered categories. As a typical example of a non-ordered category, where I do not see what happens, I point out $I = B_G$, $G$ a group (so a single object, with an automorphism group $G$). Also the case of $\boldsymbol{\Delta}[1] = \{\boldsymbol{\Delta}_0, \boldsymbol{\Delta}_1\}$ (full subcategory of $\boldsymbol{\Delta}$) is not clear to me.
\end{enumerate}

2) One would like to be able to treat the existence of $u_!^W$ dually, for want of better. Under the starting hypothesis of paragraph 4 on $(\M,W)$, this therefore leads to a change of tactics. This time, we shall have to assume the existence of a closed Quillenization that is \textit{right exact} for $W$ (which is more demanding, it seems, than such a left exact structure, in the case under consideration), so
\[
TF \ = \ F \cap W,
\]
\DERmark{ENV0620}

\msp{page 75}
and moreover also $C$, $TC$ such that we have
\[
C \ = \ \widetilde{C_0}, \qquad TC \ = \ \widetilde{TC_0}
\]
with $C_0$, $TC_0$ small (\footnote[38]{One then says that $C$, $TC$ (as left closed subsets \texttt{[rather on the left]}) are ``accessible''.}). We shall now set
\DERmark{ENV0621}

\[
\begin{aligned}
F(I) \ &= \ \{f \in \Fl \M(I) \ \mid \ f_s \in F \quad\ \ \forall\, s \in \Ob I\}\\
TF(I) \ &= \ \{f \in \Fl \M(I) \ \mid \ f_s \in TF \ \ \forall\, s \in \Ob I\},
\end{aligned}
\]
then
\[
\begin{aligned}
C(I) \ &= \ TF(I)^*\\
TC(I) \ &= \ F(I)^*\ .
\end{aligned}
\]
The hope (which I shall have to check later) is that in this way we obtain a right-uniform Quillenization of $W$, over the whole of Cat. The only problem is the validity of the factorization axioms for $f \in \Fl \M(I)$ as $f = pi$, with $i \in TC$, $p \in F$, and also with $i \in C$, $p \in TF$. % [sic] f \in F, f \in TF -> p \in F, p \in TF, the factorisation is f = pi
\DERmark{ENV0622}

I have good hope that this is true -- this reduces essentially to seeing that $(C(I),TF(I))$ and $(TC(I),F(I))$, say $(\Phi(I),\Psi(I))$, satisfy, in addition to $\Phi(I) = \Psi(I)^*$, $\Psi(I) = \Phi(I)_*$ (\textit{i.e.} $\Psi(I)$ is right closed), and that $\Phi(I)$ is accessible (just like $\Phi = \Phi(e)$). If this works, we shall thus have the existence
\printedpage{42}
% head: A. Grothendieck : Dérivateurs, Ch. XIII, Édition des Universités de Montpellier II et Paris VII
% foot: Version 25 janvier 2004
\msp{page 76}
of $u_!$ for $u$ a functor between finite ordered categories, and under slightly more general conditions. (E.g. $u : I' \longrightarrow I$ with $I$ ordered, $I'$ ordered and with underlying space noetherian and sober.) Perhaps this time, the hypothesis that $W$ is stable under filtering $\varinjlim$ will be of use after all! Thus, the case of discrete categories, and of left filtering categories such as $\mathbf{N}^\circ$, will no longer pose any problems because of this. So there is hope! Every ordered set which is a filtering $\varinjlim$ of \textit{closed} noetherian and sober ordered subsets, is then $W$-admissible on the right. But what about categories of type $\Gamma = B_G$ e.g. ? ? ?

\msp{page 77}

\section*{7\quad Cardinal filtrations and accessible subsets of a big category}

I refer to the theory of accessible functors and categories in SGA 4, I, §9 (p. 138-178). I give myself a big category (always a $\mathfrak{U}$-category) $E$ with a \textit{cardinal filtration} (increasing)
\[
\Bigl(\Filt^{\pi}(E)\Bigr)_{\pi \geq \pi_0}, \qquad \pi_0 \text{ a fixed infinite cardinal,}
\]
by the strictly full subcategories $\Filt^{\pi}(E)$. We have \texttt{[SGA 4, I, 9.12]}

\begin{enumerate}
\item[\textit{a})] $\Filt^{\pi}(E)$ essentially small, \textit{i.e.} equivalent to a small category (\textit{i.e.} to a category belonging to $\mathfrak{U}$).
\item[\textit{b})] $E$ satisfies $L_{\pi_0}$ (\textit{i.e.} $E$ is stable under $\varinjlim$ of type $I$, when $I$ is an ordered set large relative to $\pi_0$), and $\Filt^{\pi}(E)$ is stable under filtering direct limits indexed by $I$ large relative to $\pi_0$, and such that $\card I \leq \pi$.
\item[\textit{c})] For every $\pi \geq \pi_0$, and for every $X \in \Ob E$, we have
\[
X \ \simeq\ \varinjlim_{i \in I} X_i,
\]
where the $X_i$ \texttt{[are]} in $\Filt^{\pi}(E)$, and where $I$ is an ordered set large relative to $\pi$. Moreover, if $X \in \Ob \Filt^{\pi'}(E)$, with $\pi' > \pi$, we can take $I$ such that $\card I \leq \pi'^{\pi}$ (so if $\pi' = 2^{c'}$ with $c' \geq \pi$, whence $(\pi')^{\pi} = (2^{c'})^{\pi} = 2^{c'\pi} = 2^{c'} = \pi'$, we shall have $\card I \leq \pi'$). (\textbf{NB} If we assume the generalized continuum hypothesis, we shall have $\pi'^{\pi} = \pi'$ for $\pi<\operatorname{cf}(\pi')$, hence $\card I \leq \pi'$.) (\footnote[39]{NB \textit{b}) and \textit{c}) imply that $E = \bigcup_{\pi \geq \pi_0} \Filt^{\pi}(E)$.})
\DERmark{ENV0623}

\end{enumerate}

\printedpage{43}
% head: A. Grothendieck : Dérivateurs, Ch. XIII, Édition des Universités de Montpellier II et Paris VII
% foot: Version 25 janvier 2004

\msp{page 78 crossed out}

\msp{page 79}

Let us consider the representation in \textit{c}),
\[
X \ = \ \varinjlim_{I} X_i \qquad\qquad
\left\{
\begin{aligned}
& X_i \in \Filt^{\pi},\\
& I \text{ large relative to } \pi.
\end{aligned}
\right.
\]
Let
\[
\tilde{I} \ = \ \{I' \subset I \mid \card I' \leq \pi,\ \ I' \text{ large relative to } \pi_0\},
\]
endowed with the inclusion order relation. For
\[
I' \ \in\ \tilde{I}, \qquad \text{let}\ \ X_{I'} = \varinjlim_{I'} X_i\ ,
\]
this limit exists by $L_{\pi_0}$ (hypothesis \textit{b}), it is in $\Filt^{\pi}$ by \textit{b}), so we have a direct system
\[
(X_{I'})_{I' \in \tilde{I}}, \qquad X_{I'} \in \Filt^{\pi},
\]
with
\[
\card \tilde{I} \ \leq\ (\card I)^{\pi},
\]
and if $X \in \Filt^{\pi'}$ $(\pi' > \pi)$, taking
\[
\card I \ \leq\ \pi'^{\pi},
\]
we shall have $\card \tilde{I} \leq (\pi'^{\pi})^{\pi} = \pi'^{\pi^2} = \pi'^{\pi}$, hence
\[
\card \tilde{I} \ \leq\ \pi'^{\pi}.
\]
I say that $\tilde{I}$ is itself also \textit{large relative to $\pi$,} just like $I$. Indeed,
\msp{page 80}
let
\[
(I_\alpha)_{\alpha \in A}, \qquad I_\alpha \in \tilde{I}, \qquad \card A \ \leq\ \pi,
\]
let us prove that there exists $I' \in \tilde{I}$ which is an upper bound for the $I_\alpha$, \textit{i.e.} which contains $I'_0 = \bigcup_{\alpha \in A} I_\alpha$. We shall have $\card I'_0 \leq \pi^2 = \pi$ (since $\card A \leq \pi$ and $\card I_\alpha \leq \pi$ for every $\alpha$), and it suffices to see this:

\texttt{[In fact, the assertion is obvious: since} $\card I'_0 \leq \pi$\texttt{, and since} $I$ \texttt{is large relative to} $\pi$\texttt{,} $I'_0$ \texttt{admits an upper bound} $i'_0$ \texttt{in} $I$\texttt{, and one can take} $I' = I'_0 \cup \{i'_0\}$\texttt{.]}

\textbf{Lemma 1:} \textit{Let $\pi_0$, $\pi$ be two cardinals, with $\pi \geq \pi_0$, $I$ an ordered set large relative to $\pi$, and $I'_0$ a subset of $I$ of cardinality $\leq \pi$. Then there exists a subset $I'$ of $I$ satisfying the conditions:} (\footnote[40]{Provided one assumes $\pi > \pi_0$, and moreover $\pi^{\pi_0} = \pi$ (e.g. $\pi = 2^c$ with $c \geq \pi_0$).})

\begin{enumerate}
\item[a)] $I' \supset I'_0$.
\item[b)] $\card I' \leq \pi$.
\end{enumerate}


\endgroup
% END UNIT Der-ChXIII-25-01-04/Der-ChXIII-25-01-04_p041-045.tex
% BEGIN UNIT Der-ChXIII-25-01-04/Der-ChXIII-25-01-04_p046-050.tex
\begingroup\makeatletter%
\newcommand{\printedpage}[1]{\marginpar{\footnotesize p.~#1}}%
\newcommand{\msp}[1]{\marginpar{\footnotesize [#1]}}%
\newcommand{\card}{\operatorname{card}}%
\newcommand{\Filt}{\mathrm{Filt}}%
\newcommand{\Ob}{\operatorname{Ob}}%
\newcommand{\circled}[1]{\tikz[baseline=(c.base)]\node[draw,circle,inner sep=1pt](c){#1};}%
\makeatother\let\printedpage\derpp\let\msp\dermsp%
\let\mspin\DERmspin%


\printedpage{44}
% head: A. Grothendieck : Dérivateurs, Ch. XIII, Édition des Universités de Montpellier II et Paris VII

\begin{itemize}
\item[c)] \textit{$I'$ large relative to $\pi_0$.}
\end{itemize}

\texttt{[The lemma as stated is obvious for the same reasons as above (even without the hypothesis of the footnote). The proof that follows in fact proves a finer result:}

\texttt{\textsc{Lemma} 1$'$: \textit{Let $\pi_0$, $\pi$ be two cardinals, with $\pi_0<\pi$ and $\pi^{\pi_0}=\pi$, $I$ an ordered set large relative to $\pi_0$, and $I'_0$ a subset of $I$ of cardinality $\le\pi$. Then there exists a subset $I'$ of $I$ satisfying the conditions:}}

\begin{itemize}
\item[\texttt{a)}] $I'\supset I'_0$\texttt{.}
\item[\texttt{b)}] $\card I'\le\pi$\texttt{.}
\item[\texttt{c)}] \texttt{\textit{$I'$ large relative to $\pi_0$. ]}}
\end{itemize}

We are going to try to construct $I'$ by transfinite induction, defining the $I'_\alpha\subset I$ as follows:
\[
\begin{aligned}
I'_0 &= I'_0\;, \\
I'_\alpha &= \textstyle\bigcup_{\beta<\alpha} I'_\beta \quad \text{if $\alpha$ is a limit ordinal,} \\
I'_{\alpha+1} &= \left\{\parbox{0.6\textwidth}{derived from $I'_\alpha$ as follows: for every subset $J$ of $I'_\alpha$ such that $\card J\le\pi_0$, which is not bounded above in $I'_\alpha$, we choose an upper bound $i_J$ in $I$ (\texttt{[which]} exists, since $I$ \texttt{[is]} large relative to $\pi$, hence to $\pi_0$) and $I'_{\alpha+1}$ is obtained by adjoining these $i_J$ to $I'_\alpha$.}\right.
\end{aligned}
\]

\msp{page 81}
Let then $\alpha_0$ be the first ordinal of cardinality $>\pi_0$, and consider
\[
I' \;=\; I'_{\alpha_0} \;=\; \bigcup\limits_{\substack{\alpha<\alpha_0\\ \textit{i.e.}\ \card\alpha\le\pi_0}} I'_\alpha .
\]
Then $I'$ contains $I'_0$, it is large relative to $\pi_0$, but what is its cardinality? Let us prove by transfinite induction that
\[
\card I'_\alpha \;\le\pi \qquad\quad \text{if } \alpha<\alpha_0,\ \ \textit{i.e.}\ \ \card\alpha\le\pi_0 .
\]
This is O.K. for $\alpha=0$ by the hypothesis on $I'_0$. If it is true for $I'_\alpha$, is it true for $I'_{\alpha+1}$? We adjoin to $I'_\alpha$ a set of elements $i_J$ whose cardinality is $\le\card\mathfrak{P}_{\le\pi_0}(I'_\alpha)$, hence of cardinality $\le\card(I'_\alpha)^{\pi_0}\le\pi^{\pi_0}$, and if we suppose
\[
\pi^{\pi_0} \;=\; \pi,
\]
we win. Finally suppose that $\alpha$ is a limit ordinal (of cardinality $\le\pi_0$). So we have
\[
I'_\alpha \;=\; \bigcup_{\beta<\alpha} I'_\beta ,
\]
so
\printedpage{45}
% head: A. Grothendieck : Dérivateurs, Ch. XIII, Édition des Universités de Montpellier II et Paris VII
\msp{page 82}
\[
\card I'_\alpha \;\le\; \underbrace{(\card\alpha)\pi}_{\le\pi\pi=\pi^2=\pi} \qquad\quad \text{since the } \card I'_\beta\le\pi,
\]
we win again.

Finally we find, by the same argument,
\[
\card I'_{\alpha_0} \;\le\; \underbrace{(\card\alpha_0)\pi\le\pi^2}_{\substack{\text{if }\card\alpha_0\le\pi,\\ \textit{i.e.}\text{ if }\pi_0<\pi}} = \pi,
\]
so we win again.

Thus, we have found a second direct system
\[
(X_{\tilde\imath})_{\tilde\imath\in\tilde I}
\]
of the same type as the initial system, having as direct limit $X$ (that it has the same direct limit is immediate). But what have we gained? This: in $\tilde I$ every subset of cardinality $\le\pi$ and large relative to $\pi_0$ is not only \textit{bounded above} in $\tilde I$, but even admits a least upper bound there. This follows from % margin: double vertical rule beside the four lines "inductive est immédiat" ... "supérieure. Cela résulte du"

\msp{page 83}
\textbf{Lemma 2:} \textit{Let $I$ be large relative to $\pi>\pi_0$, with $\pi^{\pi_0}=\pi$} \texttt{[it suffices that $\pi\ge\pi_0$]} \textit{(cf. lemma 1). Let $(I'_\alpha)_{\alpha\in A}$ be a set of subsets of $I$ which are in $\tilde I$,} i.e.

\begin{itemize}
\item[a)] \textit{of cardinality $\le\pi$, and}
\item[b)] \textit{large relative to $\pi_0$,}
\end{itemize}

\textit{and suppose that} $\card A\le\pi$ \textit{and that $A$, ordered by inclusion, is large relative to $\pi_0$. Then $I'=\bigcup_{\alpha\in A} I'_\alpha$ is still in $\tilde I$,} i.e. \textit{$I'$ is of cardinality $\le\pi$ (since of cardinality $\le\pi^2=\pi$), and moreover, large relative to $\pi_0$.} (Immediate verification.)

Thus, we have proved this:

% margin: single vertical rule beside Proposition 1, from here to the end of the page
\textbf{Proposition 1.} \textit{Let $E$ be} \texttt{[a]} \textit{big category, with cardinal filtration} $(\Filt^\pi)_{\pi\ge\pi_0}$. \textit{Let $\pi$ be a cardinal with}
\[
\pi>\pi_0, \qquad\quad \pi^{\pi_0}=\pi \qquad\quad \textit{(e.g. $\pi=2^c$, with $c\ge\pi_0$)} \quad \texttt{[}\pi\ge\pi_0\ \texttt{suffices].}
\]
\textit{Then for every $X$ in $E$, we can find a direct system $(X_i)_{i\in I}$, with limit $X$, with}

\begin{itemize}
\item[a)] \[ \left\{ \begin{aligned} & X_i\in\Filt^\pi \ \ \textit{for every $i\in I$,} \\ & \ I \ \textit{large relative to $\pi$,} \end{aligned} \right. \]
\end{itemize}

\textit{satisfying the following additional conditions:}

\begin{itemize}
\item[b)] \textit{Every subset $J$ of $I$ of cardinality $\le\pi$, and large relative to $\pi_0$, admits a least upper bound $i_J$ in $I$.}
\item[c)] \textit{We have, under these conditions}
\end{itemize}

\printedpage{46}
% head: A. Grothendieck : Dérivateurs, Ch. XIII, Édition des Universités de Montpellier II et Paris VII
\msp{page 84}
\[
X_{i_J} \xleftarrow{\ \sim\ } \varinjlim_{i\in J} X_i .
\]
\textit{Moreover, if $X\in\Filt^{\pi'}$ with $\pi'>\pi$, we can take $I$ of cardinality $\le\pi'^{\pi}$ (hence of cardinality $\le\pi'$ if $\pi'^{\pi}=\pi'$, e.g. if $\pi'=2^{c'}$, with $c'\ge\pi$).}

\textbf{Definition.} We call a direct system $(X_i)_{i\in I}$, satisfying conditions \textit{a}), \textit{b}), \textit{c}) above, a \textit{well-adapted direct system} (to the given cardinal filtration, for the given cardinals $\pi_0$, $\pi$), and the corresponding representation
\[
X \xleftarrow{\ \sim\ } \varinjlim X_i
\]
a \textit{well-adapted representation} of $X$ in terms of objects of $\Filt^\pi$.

\textbf{NB} When $E$ is stable under small filtering $\varinjlim$, the situation simplifies considerably. We can then take
\[
\tilde I \;=\; \{I'\subset I \mid \card I'\le\pi,\ I' \text{ filtering}\}
\]
(without burdening ourselves with the condition $I'$ large relative to $\pi_0$) and again define the direct system $(X_{I'})_{I'\in\tilde I}$ by
\[
X_{I'} \;=\; \varinjlim_{i\in I'} X_i \;.
\]
We again have
\msp{page 85}
\[
\card\tilde I \;\le\; (\card I)^\pi \;\le\; \pi'^{\pi} \qquad\quad \text{if }\ \card I\le\pi'^{\pi}
\]
(in particular, if $X\in\Filt^{\pi'}$, and we take $I$ as indicated in condition $c$ of cardinal filtrations). Moreover, this time, the fact that $\tilde I$ is large relative to $\pi$ is trivial, we even have better, apart from $I$ being filtering:

\begin{itemize}
\item[$(*)$] \textit{Every} subset $(I'_\alpha)_{\alpha\in A}$ of $\tilde I$ which is filtering and of cardinality $\le\pi$ admits a least upper bound $\tilde\imath_A$,
\end{itemize}

namely the union (which is indeed filtering). Moreover, the direct system $(X_{\tilde\imath})_{\tilde\imath\in\tilde I}$ has the stronger property that
\[
X_{\tilde\imath_A} \xleftarrow{\ \sim\ } \varinjlim_{\tilde\imath\in A} X_{\tilde\imath} \;. % [sic] \varprojlim (arrow under lim pointing left) -> \varinjlim, inductive limit of the inductive system X_ĩ over the filtrant part A
\]
\DERmark{ENV0624}

Thus we have

\textbf{Proposition 1 bis.} \textit{Suppose $E$ is stable under small filtering $\varinjlim$, and} $\Filt^\pi$ \textit{stable under these, where $\pi\ge\pi_0$ is given. Then for every $X$ in $E$, we can find a direct system $(X_i)_{i\in I}$ with direct limit $X$, with}

\begin{itemize}
\item[a)] $X_i\in\Filt^\pi$ \textit{for every $i\in I$, $I$ an ordered set large relative to $\pi$.}
\item[b)] \textit{Every filtering subset $J$ of $I$ of cardinality $\le\pi$ admits a least upper bound $i_J$ in $I$.}
\item[c)] \textit{Under these conditions}
\end{itemize}

\printedpage{47}
% head: A. Grothendieck : Dérivateurs, Ch. XIII, Édition des Universités de Montpellier II et Paris VII
\msp{page 86}
\[
X_{i_J} \;=\; \varinjlim_{i\in J} X_i \;.
\]
Such direct systems will be called \textit{``perfectly adapted''} (to the given cardinal filtration, for the given cardinal $\pi$ – $\pi_0$ no longer figures in it, except as the beginning of the cardinal filtration and nothing more).

Since in practice, we will have to handle mostly $E$ stable under small filtering $\varinjlim$, and cardinal filtrations stable accordingly (strengthening of condition \textit{b}) of cardinal filtrations), it is this case, and the notion of \textit{perfectly adapted representation} of $X$ as $\varinjlim\limits_{I} X_i$ ($X_i\in\Filt^\pi$), that I will put forward. The ``general'' case, where we suppose only \texttt{[that]} $E$ satisfies $L_{\pi_0}$, is more convoluted and is handled less prettily (\footnote[41]{We will call a cardinal filtration \textit{strict,} if the $\underline{\Filt}^\pi$ are stable under filtering $\varinjlim$ of cardinality $\le\pi$ (and these small filtering limits exist in $E$).}).

\textbf{Definition 2.} Let $E$ be equipped with a strict cardinal filtration $(\underline{\Filt}^\pi)_{\pi\ge\pi_0}$. Let $M$ be a subset of $\Ob E$, $\underline{M}$ the full subcategory of $E$ which it defines, and let $\pi$ be a cardinal $\ge\pi_0$. We say that $M$ is \textit{$\pi$-accessible,} or that $\underline{M}$ is \textit{$\pi$-accessible}, % [sic] où que -> ou que, disjunction "ou" not the relative "où"
\msp{page 87}
if the following conditions are satisfied. For every direct system
\[
(X_i)_{i\in I}
\]
perfectly adapted to the cardinal filtration for $\pi$ (hence $I$ a filtering ordered set such that every filtering subset $J$ of cardinality $\le\pi$ has a least upper bound $i_J$, and $X_{i_J}\xleftarrow{\ \sim\ }\varinjlim\limits_{j\in J} X_j$); let $I_M$ be the set of $i\in I$ such that $X_i\in M$. With this understood, we want this: for $X$ ($=\varinjlim X_i$) to be in $M$, it is necessary and sufficient that $I_M$ be cofinal in $I$ (so that $X$ appears as a filtering $\varinjlim$, with index set $I_M$ large relative to $\pi$, of elements of $M\cap\underline{\Filt}^\pi\overset{\mathrm{def}}{=}M^\pi$).
\DERmark{ENV0625}


\msp{page 88}
I feel like taking up again the definition of the $\pi$-accessibility of $M$, for a given cardinal filtration, by way of equivalent conditions. For every direct system $(X_i)_{i\in I}$ in $E$ \texttt{[rather in $E$]}, I keep the notation
\DERmark{ENV0626}

\[
I_M \;=\; \{i\in I \mid X_i\in M\}.
\]
I set, for every cardinal $\pi\ge\pi_0$,
\[
M^\pi \;=\; M\cap\Filt^\pi .
\]

\begin{itemize}
\item[\circled{a$_\pi$}] $M$ is the set of objects of the form
\[
\varinjlim_{i\in I} X_i \;,
\]
where the $X_i$ are in $M^\pi$, and $I$ is large relative to $\pi$.
\end{itemize}

\printedpage{48}
% head: A. Grothendieck : Dérivateurs, Ch. XIII, Édition des Universités de Montpellier II et Paris VII

\begin{itemize}
\item[\circled{b$_\pi$}] Let $(X_i)_{i\in I}$ be a direct system in $E$ \texttt{[rather in $E$]} perfectly adapted to the cardinal filtration for $\pi$, \textit{i.e.} $I$ filtering and every filtering subset $J$ of $I$ of cardinality $\le\pi$ admits a least upper bound $i_J$, and $X_{i_J}\xleftarrow{\ \sim\ }\varinjlim\limits_{J} X_i$, and finally $X_i\in\Filt^\pi$ for every $i\in I$. Let $X=\varinjlim X_i$. In order to have $X\in M$, it is necessary and sufficient that $I_M$ be cofinal in $I$.
\item[\circled{c$_\pi$}] $\underline{M}^\pi=\underline{M}\cap\underline{\Filt}^\pi$ is a subcategory of $\underline{M}$ generating by strict epimorphisms, $\underline{M}$ is stable under $\varinjlim$ large relative to $\pi$, finally for every $X\in M$, $\underline{M}^\pi/X$ is a filtering category.
\end{itemize}
\DERmark{ENV0627}


\msp{page 89}
I presume that these conditions are equivalent.

$b_\pi\Longrightarrow a_\pi$. We suppose \textit{b}) holds, we want to prove \textit{a}). First that every $X\in M$ is of the stated form $\varinjlim\limits_{I} X_i$, with the $X_i\in M^\pi$ and $I$ large relative to $\pi$. By \texttt{[the]} proposition 1, we know that we have a perfect adapted presentation for $\pi$, $X=\varinjlim\limits_{i\in I} X_i$. By \textit{b}) ``it is necessary'', the subset $I_M$ of $I$ is cofinal in $I$, so since $I$ is large relative to $\pi$, so is $I_M$, and we will have (setting $J=I_M$)
\[
X \;=\; \varinjlim_{i\in J} X_i, \qquad\quad X_i\in M^\pi,\ \ J \text{ large relative to } \pi,
\]
O.K. Conversely, let us prove that if $X$ can be put in the form stated in \textit{a}), it is in $M$. This will follow from

\textbf{Lemma 1:} \textit{If condition} b) (\textit{or} a)) \textit{is satisfied, $M$ is stable under direct limits large relative to $\pi$.}

We must prove that if
\msp{page 90}
\[
X \;=\; \varinjlim_{i\in I} X_i \;,
\]
with $I$ large relative to $\pi$, the $X_i$ in $M$, then $X\in M$. Condition \textit{b}) ``it is sufficient'' ensures this when we suppose in addition the $X_i$ in $\Filt^\pi$ (hence in $M^\pi$), so that we will have $I_M=I$, hence $I_M$ cofinal, hence $X\in M$ (\footnote[42]{Not quite – we assume more there, about the direct system. We cannot avoid, even in the case $X_i\in M^\pi$ $\forall\, i\in I$, making the argument of the following pages.}). If we do not suppose the $X_i$ in $\Filt^\pi$, we must work a little more.

Go straight to the part of the proof below, page 96, which proves \textit{b}) $\Longrightarrow$ \textit{a}). The rest of the proof will prove the

\textbf{Lemma 2.} \textit{If every element of Filt$^\pi$ is $\pi$-accessible, then} b$_\pi$) \textit{(and even} a$_\pi$\textit{) implies that $M$ is stable under $\varinjlim$ large relative to $\pi$.}


\endgroup
% END UNIT Der-ChXIII-25-01-04/Der-ChXIII-25-01-04_p046-050.tex
% BEGIN UNIT Der-ChXIII-25-01-04/Der-ChXIII-25-01-04_p051-055.tex
\begingroup\makeatletter%
\newcommand{\printedpage}[1]{\marginpar{\footnotesize p.~#1}}%
\newcommand{\msp}[1]{\marginpar{\footnotesize [#1]}}%
\newcommand{\barlongleftrightarrow}{\mathrel{\text{\ooalign{$\longleftrightarrow$\cr\hidewidth$\mid$\hidewidth\cr}}}}%
\newcommand{\idx}[2]{\noindent #1\dotfill #2\par}%
\makeatother\let\printedpage\derpp\let\msp\dermsp%
\let\mspin\DERmspin%


\printedpage{49}
% head: A. Grothendieck : Dérivateurs, Ch. XIII, Édition des Universités de Montpellier II et Paris VII
% foot: Version 25 janvier 2004

\msp{page 91}Let
\[ \mathcal{C} \;=\; \underline{M}^\pi/\mathcal{X} \]
be the following category. Its objects are the triples
\[ (Z,i,\alpha) \quad \text{such that} \quad Z\in M^\pi,\ i\in I,\ \alpha : Z\longrightarrow X_i \ , \]
the arrows
\[ (Z,i,\alpha) \;\longrightarrow\; (Z',i',\alpha') \]
are the pairs $(u,\lambda)$ with
\[ u : Z\longrightarrow Z', \qquad \lambda : i\longrightarrow i', \qquad (\textit{i.e.}\ \text{we have } i\le i' \text{ and } \lambda \text{ is the unique arrow from } i \text{ to } i'), \]
such that the diagram
\[
\begin{tikzcd}[row sep=large, column sep=large]
Z \arrow[r,"\alpha"] \arrow[d,"u"'] & X_i \arrow[d,"\text{transition } \lambda_{i',i}"] \\
Z' \arrow[r,"\alpha'"] & X_{i'}
\end{tikzcd}
\]
commutes. We can therefore define $\mathcal{C}$ as the fiber product in the cartesian square
\[
\begin{tikzcd}[row sep=large, column sep=large]
 & \mathcal{C} \arrow[r,"\text{cofibered}"] \arrow[d] \arrow[dl,dotted,"\mathcal{Z}"'] & I \arrow[d,"\mathcal{X} : i \longmapsto X_i"] \\ % [sic] X \longmapsto X_i -> i \longmapsto X_i, the functor \mathcal{X} goes from I, so the variable should be i
\underline{M}^\pi & \underline{M}^\pi/E \arrow[l,dotted,"\text{source}"'] \arrow[r,"\text{target}","\text{cofibered}"'] & E.
\end{tikzcd}
\]
\DERmark{ENV0628}

The fiber $\mathcal{C}_i$ of $\mathcal{C}$ at $i\in I$ is $M^\pi/X_i$ \texttt{[rather} $M^\pi/X_i$\texttt{]}. We have, as in \textit{b}) $\Longrightarrow$ first half of \textit{a}), that $\underline{M}^\pi$ is a
\DERmark{ENV0629}

\msp{page 92}system of strict generators of $\underline{M}$ \texttt{[this is not clear without further hypotheses on} $\underline{M}$\texttt{]}, and $X_i\in\operatorname{Ob}\underline{M}$ by hypothesis,
\[ X_i \;=\; \varinjlim_{Z\in\operatorname{Ob}\underline{M}^\pi/X_i} Z \;=\; \varinjlim_{\mathcal{C}_i}\mathcal{Z}|\mathcal{C}_i, \]
so the given direct system $(X_i)_{i\in I}$ on $I$ with values in $E$, is deduced from the direct system (not ordered in general) on $\mathcal{C}$, with values in $\underline{M}^\pi$, by ``integrating over the fibers'' of $\mathcal{C}\longrightarrow I$. Besides, since $\underline{M}^\pi/E\longrightarrow E$ is cofibered, hence proper, the same is true of $\mathcal{C}\longrightarrow I$, so the $\varinjlim$ of $\mathcal{Z}$ over the fibers $\mathcal{C}_i$ ($=$ $\underline{M}^\pi/X_i$) do form a covariant functor $I\longrightarrow E$, whose direct limit is that of $\mathcal{Z}$ (\footnote[43]{\textbf{NB} If we admit the validity of $a$), then it already follows that $\varinjlim\limits_{i\in I} X_i \simeq \varinjlim\limits_{\alpha\in C} Z_\alpha$ is in $M$. (Stability of $M$ under $\varinjlim$ large relative to $\pi$).}).
\DERmark{ENV0630}


\printedpage{50}
% head: A. Grothendieck : Dérivateurs, Ch. XIII, Édition des Universités de Montpellier II et Paris VII
% foot: Version 25 janvier 2004

\msp{page 93}Let, for every $i\in I$, $\tilde{\mathcal{C}}_i$ be the ordered set
\[ \tilde{\mathcal{C}}_i = \{\text{subcategories } \mathcal{C}'_i \text{ of } \mathcal{C}_i = \underline{M}^\pi/X_i \mid \operatorname{card}\mathcal{Z}(\mathcal{C}'_i)\le\pi,\ \ \mathcal{C}'_i \text{ filtering}\}. \]
I admit for the moment the

\textbf{Lemma} (\footnote[44]{Uses the tacit hypothesis that the elements of $\mathrm{Filt}^\pi$ are $\pi$-accessible.}). \textit{For every $X\in M$, the category $\underline{M}^\pi/X$ is filtering. (Under hypothesis} b), \textit{or that every object of $M$ is of the form stated in} a).\textit{)}

It follows, since $X_i\in M$, that $\mathcal{C}_i$ is filtering, so the union of the $\mathcal{C}'_i\in\tilde{\mathcal{C}}_i$ is $\mathcal{C}_i$; the ordered set $\tilde{\mathcal{C}}_i$ is filtering, and every filtering subset $(\mathcal{C}'_{i_\alpha})_{\alpha\in A}$ of cardinality $\le\pi$ of $\tilde{\mathcal{C}}_i$ has a supremum $\mathcal{C}'_{i_A}$ there, namely the union of the $\mathcal{C}'_{i_\alpha}$. $\tilde{\mathcal{C}}_i$ is therefore a $\pi$-adapted ordered set.

We can regard the ordered sets $\tilde{\mathcal{C}}_i$, for varying $i\in I$, as the fibers of an increasing map of ordered sets
\[ p \;:\; \tilde{\mathcal{C}} \;\longrightarrow\; I, \]
\DERmark{ENV0631}

by writing
\[ \left\{ \begin{array}{l}
(\mathcal{C}'_i\subset\mathcal{C}_i) \;\le\; (\mathcal{C}'_j\subset\mathcal{C}_j) \\
\text{if and only if } i\le j, \\
\text{and the canonical transition functor } \mathcal{C}_i\longrightarrow\mathcal{C}_j \text{ maps } \mathcal{C}'_i \text{ into } \mathcal{C}'_j.
\end{array}\right. \]
We thus see that $\tilde{\mathcal{C}}$ too is filtering and large relative to $\pi$.

\msp{page 94}We define a direct system on $\tilde{\mathcal{C}}$, with values in $\underline{\mathrm{Filt}}^\pi$, by setting
\[ Z_{\mathcal{C}'_i} \;=\; \varinjlim_{\alpha\,\in\,\mathcal{C}'_i\ (\subset\,\mathcal{C}_i\,=\,\underline{M}^\pi/X_i)} Z_\alpha \]
(where an $\alpha\in\mathcal{C}_i$ consists, I recall, of $Z_\alpha\xrightarrow{u_\alpha}X_i$, with $Z_\alpha\in M^\pi\subset\underline{\mathrm{Filt}}^\pi$). These $Z_{\mathcal{C}'_i}$ are no longer necessarily in $M^\pi$, but only in $\mathrm{Filt}^\pi$. (Then what was the use of defining the $\mathcal{C}_i$ themselves as $\underline{M}^\pi/X_i$, instead of $\underline{\mathrm{Filt}}^\pi/X_i$?)

The direct system
\[ (Z_{\mathcal{C}'_i})_{\mathcal{C}'_i\in\tilde{\mathcal{C}}_i} \]
is perfectly well adapted to the cardinal filtration for the cardinal $\pi$, and it is immediate that its $\varinjlim$ is isomorphic to $\varinjlim\limits_{\alpha\in\operatorname{Ob}\mathcal{C}_i=\underline{M}^\pi/X_i} Z_\alpha = X_i\in M$ . Thus by \textit{b}), the set $\tilde{\mathcal{C}}_M$ of the $\mathcal{C}'_i\in\tilde{\mathcal{C}}_i$ such that $Z_{\mathcal{C}'_i}\in M$ (so $Z_{\mathcal{C}'_i}\in M^\pi$) is cofinal in $\tilde{\mathcal{C}}_i$. It follows therefore that the set $\tilde{\mathcal{C}}_M$ is also cofinal in $\tilde{\mathcal{C}}$.

\printedpage{51}
% head: A. Grothendieck : Dérivateurs, Ch. XIII, Édition des Universités de Montpellier II et Paris VII
% foot: Version 25 janvier 2004

\msp{page 95}We also check immediately that
\[ \varinjlim_{\alpha\in\tilde{\mathcal{C}}} Z_\alpha \;=\; \varinjlim X_i \;=\; X, \]
so, $\tilde{\mathcal{C}}_M$ being cofinal,
\[ \varinjlim_{\alpha\in\tilde{\mathcal{C}}_M} Z_\alpha \;=\; X. \]
All in all, we have managed to represent $X$ in the form
\[ X \;=\; \varinjlim_{\alpha\in A} X_\alpha, \]
where this time the $X_\alpha$ are \textit{in} $M^\pi$ (not just in $M$), and $A$ is large relative to $\pi$. But we do not know that every filtering subset $B$ of cardinality $\le\pi$ of $A$ has a supremum $\alpha_B$ in $A$, and still less that $X_{\alpha_B} = \varinjlim\limits_{\alpha\in B} X_\alpha$ (\footnote[45]{If we assume $(a_\pi)$, this argument proves at least that $M$ is stable under $\varinjlim$ large relative to $\pi$.}). So we cannot yet apply \textit{b}) directly. If it were only a matter of proving the second half of \textit{a}), where we already assume the $X_i\in M^\pi$ and $I$ large relative to $\pi$, we have gained nothing at all by our labors!

\msp{page 96}But let us return to this case
\[ X \;=\; \varinjlim X_i \;\left|\; \begin{array}{l} X_i\in M^\pi \\[1.5ex] I \text{ large relative to } \pi, \end{array}\right. \]
and introduce again
\[ \tilde{I} \;=\; \{I'\subset I \mid I' \text{ filtering, of cardinality } \le\pi\}, \]
so this time $\tilde{I}$ is indeed adapted to $\pi$ (filtering, and every filtering subset of cardinality $\le\pi$ has a supremum), and if we set
\[ X_{\tilde{\imath}} \;=\; \varinjlim_{i\in I'_{\tilde{\imath}}} X_i, \]
we see that the direct system
\[ (X_{\tilde{\imath}})_{\tilde{\imath}\in\tilde{I}} \]
is perfectly adapted to $\pi$ and the cardinal filtration. Moreover, I say that the set $\tilde{I}_M$ is \textit{cofinal} in $\tilde{I}$. Indeed, $\tilde{I}_M$ contains the subset $\tilde{I}_0$ consisting of the $I'\in\tilde{I}$ which have a greatest element -- and $\tilde{I}_0$ is indeed cofinal! Now, we can apply
\printedpage{52}
% head: A. Grothendieck : Dérivateurs, Ch. XIII, Édition des Universités de Montpellier II et Paris VII
% foot: Version 25 janvier 2004
\msp{page 97}the ``it suffices'' in \textit{b}), which concludes that $X = \varinjlim X_i\in M$, q.e.d. % [sic] conclue -> conclut, verb form (indicative of conclure)

We have proved $b\Longrightarrow a$, without using the hypothesis that every object of $\mathrm{Filt}^\pi$ is $\pi$-accessible. But we have used this hypothesis to prove that $b$ (or $a$) implies that $M$ is stable under $\varinjlim$ large relative to $\pi$ (lemma 2, p. 90).

$a\overset{?}{\Longrightarrow} b$. The ``it suffices'' in \textit{b}) is an immediate consequence of \textit{a}). It remains to prove that if, in the $\pi$-adapted direct system $(X_i)_{i\in I}$ with limit $X$, we have $X\in M$, then $I_M$ is cofinal in $I$. But by hypothesis \textit{a}), we will have $X = \varinjlim\limits_J Y_j$, with $J$ large relative to $\pi$, the $Y_j$ in $M^\pi$. Let us admit that $\mathrm{Filt}^\pi\subset E_\pi$, \textit{i.e.} every element of $\mathrm{Filt}^\pi$, in particular the $X_i$, are $\pi$-accessible.

I think I will not manage to prove $a_\pi\Longrightarrow b_\pi$, but only this: If $\pi'$ is a cardinal $>\pi$ (say the successor of $\pi$ in the ordinal sequence of cardinals), then $a_\pi\Longrightarrow b_{\pi'}$.
\DERmark{ENV0632}


\msp{page 98}This will follow from

\textbf{Lemma 3.} \textit{Let $X\simeq\varinjlim\limits_I X_i$, $X\simeq\varinjlim\limits_J Y_j$ be two representations of $X$ in the form of filtering $\varinjlim$, with $I$ large relative to $\pi'$, $J$ large relative to $\pi$, with $\pi$, $\pi'$ two cardinals such that $\pi'>\pi$. We assume $(X_i)_{i\in I}$ $\pi'$-adapted,} i.e. \textit{for every filtering subset $I'$ of $I$ of cardinality $\le\pi'$, $I'$ has a supremum $i_{I'}$, and $X_{i_{I'}} = \varinjlim\limits_{i\in I'} X_i$. Then for every $i\in I$, we can find a subset $I'$ of $I$, containing $i$,} large relative to $\pi$, \textit{of cardinality $\le\pi'$, and a subset $J'$ of $J$, finally an isomorphism $I'\xrightarrow[\sim]{\lambda}J'$ and a strictly increasing map $\mu : J'\longrightarrow I'$ and homomorphisms} \texttt{[end of chapter]}

\par
\printedpage{53}
% head: A. Grothendieck : Dérivateurs, Ch. XIII, Édition des Universités de Montpellier II et Paris VII
% foot: Version 25 janvier 2004

\section*{Index of notations}

\idx{LLP, RLP}{1}
\idx{$\Phi_*$, $\Psi^*$}{1}
\idx{$\widetilde{\Phi}$, $\overline{\Psi}$}{1}
\idx{$\Phi \longleftrightarrow \Psi$}{3}
\idx{$\mathrm{Cof}_W$, $\mathrm{Fib}_W$}{8}
\idx{$W^{\mathrm{univ}}$}{9}
\idx{$W_{\mathrm{univ}}$}{10}
\idx{$\underline{\mathrm{Fib}}_W S\subset\mathcal{M}/S$}{10}
\idx{$W_S$}{10}
\idx{$W_S^{\mathrm{u}}$}{10}
\idx{$S\backslash\underline{\mathrm{Cof}}_W$}{10}
\idx{$\Phi \barlongleftrightarrow \Psi$}{28}
\idx{$\mathcal{M}_{\mathrm{cf}}$, $W_{\mathrm{cf}}$}{29}
\idx{$\mathrm{Ho}_W$}{29}
\idx{$\mathrm{Hom}_W$}{30}
\idx{$\pi_{\mathrm{d}}$, $\pi$}{30}
\idx{$\mathcal{M}_{\mathrm{c}}$}{30}
\idx{$W_{\mathrm{c}}$}{31}
\idx{$\pi_{\mathrm{g}}$}{31}
\idx{$\mathcal{M}_{\mathrm{f}}$}{31}
\idx{$f\sim g$}{32}
\idx{$\Psi,\ \Phi,\ \mathcal{M}(\Psi),\ W(\Psi),\ \mathcal{M}(\Phi),\ W(\Phi)$}{33}
\idx{$\mathcal{M}(I)$, $W(I)$}{34}
\idx{$F(I),\ TF(I),\ C(I),\ TC(I)$}{35}
\idx{$\alpha_s$}{35}
\idx{$u^*$, $u_*$, $u_!$}{45}
\idx{$u^*_W$, $u^W_*$, $u^W_!$}{49}
\idx{$\mathrm{HO}_W(I)$}{49}
\idx{Diag}{50}
\idx{$\mathrm{HO}_{\mathcal{M},W}(I)$, $\mathbf{D}_{\mathcal{M},W}(I)$}{54}
\idx{$\mathrm{HO}^{\mathrm{lc}}_W(I)$, $u^*_{\mathrm{lc}}$}{56}
\idx{$u^{\mathrm{f}}_*$}{59}
\idx{$u^*_{\mathrm{c}}$}{67}
\idx{$\mathfrak{U}$}{77}
\idx{$\mathrm{Filt}^\pi(E)$}{77}


\endgroup
% END UNIT Der-ChXIII-25-01-04/Der-ChXIII-25-01-04_p051-055.tex
% BEGIN UNIT Der-ChXIII-25-01-04/Der-ChXIII-25-01-04_p056-060.tex
\begingroup\makeatletter%
\newcommand{\printedpage}[1]{\marginpar{\footnotesize p.~#1}}%
\newcommand{\msp}[1]{\marginpar{\footnotesize [#1]}}%
\newcommand{\entry}[2]{\noindent #1\dotfill\ #2\par}%
\newcommand{\Fl}{\mathrm{Fl}}%
\makeatother\let\printedpage\derpp\let\msp\dermsp%
\let\mspin\DERmspin%


\printedpage{54}
% head: A. Grothendieck : Dérivateurs, Ch. XIII, Édition des Universités de Montpellier II et Paris VII
% foot: Version 25 janvier 2004

\entry{$L_\pi$}{77}
\entry{$\tilde{I}$}{79}
\entry{$\underline{M}$}{86}
\entry{$I_M$}{87}
\entry{$M^\pi$}{87}

\par

\printedpage{55}
% head: A. Grothendieck : Dérivateurs, Ch. XIII, Édition des Universités de Montpellier II et Paris VII
% foot: Version 25 janvier 2004

\section*{Terminological index}

\entry{accessible (left closed subset)}{75}
\entry{admissible (on the left)}{59}
\entry{category of models}{7}
\entry{closed Quillen category}{25}
\entry{left closed, right closed}{2}
\entry{Quillen couple}{5}
\entry{weakly saturated}{7}
\entry{cardinal filtration}{77}
\entry{strict cardinal filtration}{86}
\entry{arrow couniversally in $W$}{9}
\entry{arrow universally in $W$}{10}
\entry{large relative to $\pi$}{77}
\entry{path object}{32}
\entry{cylinder object}{32}
\entry{$\pi$-accessible (subset of a category equipped with a strict cardinal filtration)}{86}
\entry{left, right pre-Q-category of models}{12}
\entry{left, right pre-Quillen}{14}
\entry{closed Quillen pre-quintuple}{21}
\entry{proper}{8}
\entry{Quillen factorization property}{5}
\entry{Q-category of models}{8}
\entry{right Q-category of models}{21}
\entry{``closed'' Q-category of models}{16}
\entry{left, right Q-homotopy}{30}
\entry{Quillenizable}{50}
\entry{Quillenizable on Diag}{51}
\entry{Quillenizable on Diag uniformly and left exactly}{57}
\entry{uniformly left Quillenizable on Diag}{57}
\entry{uniform right Quillenization on Diag}{58}
\entry{uniform left Quillenization on Diag}{57}
\entry{uniform and left exact Quillenization on Diag}{57}
\entry{well-adapted representation}{84}
\entry{perfectly adapted representation}{86}
\entry{strictly admissible}{61}
\entry{Quillen structure}{28}

\par

\printedpage{56}
% head: A. Grothendieck : Dérivateurs, Ch. XIII, Édition des Universités de Montpellier II et Paris VII
% foot: Version 25 janvier 2004

\entry{exact Quillen structure}{27}
\entry{left exact, right exact Quillen structure}{27}
\entry{weak Quillen structure}{27}
\entry{biexact weak Quillen structure}{28}
\entry{closed, left closed, right closed weak Quillen structure}{28}
\entry{perfect Quillen structure}{28}
\entry{closed Quillen system}{26}
\entry{weak Quillen system}{26}
\entry{left exact, right exact weak Quillen system}{26}
\entry{well-adapted direct system}{84}
\entry{perfectly adapted direct system}{86}
\entry{closed Quillen triple}{25}
\entry{Quillen triple}{14}
\entry{$W$-cofibration}{8}
\entry{couniversal $W$-equivalence}{9}
\entry{$W$-fibration}{8}

\par

\printedpage{57}
% head: A. Grothendieck : Dérivateurs, Ch. XIII, Édition des Universités de Montpellier II et Paris VII
% foot: Version 25 janvier 2004

\section*{List of axioms}

\noindent $\mathcal{M}$ category

\noindent $W \ \subset\ \Fl(\mathcal{M})$ , \qquad $C$, $F$, $TC$, $TF \ \subset\ \Fl(\mathcal{M})$

\begin{description}[leftmargin=4.5em,labelwidth=4em]
\item[M 0] The category $\mathcal{M}$ admits finite inverse limits and finite direct limits.

\item[M 1] The arrows belonging to $C$ have the left lifting property relative to those belonging to $TF$, and the arrows belonging to $F$ have the right lifting property relative to those belonging to $TC$.

\item[M 2] Every arrow of $\mathcal{M}$ decomposes into an arrow of $C$ followed by an arrow of $TF$, and into an arrow of $TC$ followed by an arrow of $F$.

\item[M 3] The subset $F$ of $\Fl(\mathcal{M})$ is stable under composition and base change, $C$ is stable under composition and cobase change, $F$ and $C$ contain the isomorphisms.

\item[M 4] The subset $TF$ of $\Fl(\mathcal{M})$ is stable under composition and base change, $TC$ is stable under composition and cobase change, $TF$ and $TC$ contain the isomorphisms.

\item[M 5] If in a commutative triangle of $\mathcal{M}$ two of the three arrows are in $W$, the same holds for the third ; $W$ contains the isomorphisms.

\item[M$'$] $TC \ \subset\ C \cap W$, $TF \ \subset\ F \cap W$.

\item[M$''$] The arrows belonging to $C$ are $W$-cofibrations, and the arrows belonging to $F$ are $W$-fibrations.

\item[M$'_{\mathbf{0}}$(g)] $TC = C \cap W$.

\item[M$'_{\mathbf{0}}$(d)] $TF = F \cap W$.

\item[M 2\,bis] $W = TF \circ TC$.
\end{description}

\par

\printedpage{58}
% head: A. Grothendieck : Dérivateurs, Ch. XIII, Édition des Universités de Montpellier II et Paris VII
% foot: Version 25 janvier 2004

\setcounter{footnote}{0}
\renewcommand{\thefootnote}{\fnsymbol{footnote}}
\section*{Glossary}

\subsection*{Section 1}

\begin{itemize}
\item $\Phi \ \subset\ \Fl(\mathcal{M})$

$\Phi$ is \textit{left closed} : there exists \ $\Psi \ \subset\ \Fl(\mathcal{M})$ such that $\Phi$ consists of the arrows having the left lifting property relative to the arrows of $\Psi$.

\item $\Psi \ \subset\ \Fl(\mathcal{M})$

$\Psi$ is \textit{right closed} : there exists \ $\Phi \ \subset\ \Fl(\mathcal{M})$ such that $\Psi$ consists of the arrows having the right lifting property relative to the arrows of $\Phi$.

\item $\Phi$, $\Psi \ \subset\ \Fl(\mathcal{M})$

$(\Phi, \Psi)$ has the \textit{Quillen factorization property} :
\begin{itemize}
\item[a)] the arrows of $\Phi$ have the left lifting property relative to the arrows of $\Psi$ ;
\item[b)] every arrow $f$ of $\mathcal{M}$ factors as
\[
f = pi \quad , \qquad i \in \Phi \ , \quad p \in \Psi \ .
\]
\end{itemize}

\item $(\Phi, \Psi)$ \textit{Quillen couple} :
\begin{itemize}
\item[a)] $(\Phi, \Psi)$ has the Quillen factorization property ;
\item[b)] $\Phi$ is left closed and $\Psi$ right closed.
\end{itemize}
\end{itemize}

\subsection*{Section 2}

\begin{itemize}
\item $W \ \subset\ \Fl(\mathcal{M})$

\textit{$W$ weakly saturated} :
\begin{itemize}
\item[a)] $W$ contains the isomorphisms ;
\item[b)] if two among the three arrows of a commutative triangle are in $W$, the third too.
\end{itemize}

\item $W \ \subset\ \Fl(\mathcal{M})$

A morphism $p : X \longrightarrow Y$ of $\mathcal{M}$ is a \textit{$W$-fibration} ($p \in \mathrm{Fib}_W$) if for every diagram of cartesian squares
\end{itemize}


\endgroup
% END UNIT Der-ChXIII-25-01-04/Der-ChXIII-25-01-04_p056-060.tex
% BEGIN UNIT Der-ChXIII-25-01-04/Der-ChXIII-25-01-04_p061-065.tex
\begingroup\makeatletter%
\newcommand{\printedpage}[1]{\marginpar{\footnotesize p.~#1}}%
\newcommand{\msp}[1]{\marginpar{\footnotesize [#1]}}%
\newcommand{\Fl}{\mathrm{Fl}}%
\newcommand{\Filt}{\mathrm{Filt}}%
\newcommand{\Mo}{\mathcal{M}}%
\newcommand{\Mzg}{$\mathbf{M}'_{\mathbf{0}}\mathbf{(g)}$}%
\newcommand{\Mzd}{$\mathbf{M}'_{\mathbf{0}}\mathbf{(d)}$}%
\renewcommand{\thefootnote}{\fnsymbol{footnote}}%
\makeatother\let\printedpage\derpp\let\msp\dermsp%
\let\mspin\DERmspin%
\edef\derfnsaved{\the\value{footnote}}\setcounter{footnote}{0}%


\printedpage{59}
% head: A. Grothendieck : Dérivateurs, Ch. XIII, Édition des Universités de Montpellier II et Paris VII

\[
\begin{tikzcd}
X \arrow[d, "p"'] & X' \arrow[l] \arrow[d] & X'' \arrow[l, "t"'] \arrow[d] \\
Y & Y' \arrow[l] & Y'' \arrow[l, "s"]
\end{tikzcd}
\quad ,
\]

if $s \in W$, then $t \in W$.

\begin{itemize}
\item[\textbullet] $W \ \subset\ \Fl(\Mo)$

A morphism $i : A \longrightarrow B$ of $\Mo$ is a $W$-\textit{cofibration} ($i \in \mathrm{Cof}_W$) if for every diagram of cocartesian squares
\[
\begin{tikzcd}
A'' \arrow[d] & A' \arrow[l, "s"'] \arrow[d] & A \arrow[l] \arrow[d, "i"] \\
B'' & B' \arrow[l, "t"] & B \arrow[l]
\end{tikzcd}
\quad ,
\]
if $s \in W$, then $t \in W$.

\item[\textbullet] $W \ \subset\ \Fl(\Mo)$

A morphism $p : X \longrightarrow Y$ of $\Mo$ is \textit{universally in} $W$, or a $W$-\textit{universal equivalence} ($p \in W_{\mathrm{univ}}$) if for every cartesian square
\[
\begin{tikzcd}
X \arrow[d, "p"'] & X' \arrow[l] \arrow[d, "p'"] \\
Y & Y' \arrow[l]
\end{tikzcd}
\quad ,
\]
the arrow $p'$ belongs to $W$.

\item[\textbullet] $W \ \subset\ \Fl(\Mo)$

A morphism $i : A \longrightarrow B$ of $\Mo$ is \textit{couniversally in} $W$, or a $W$-\textit{couniversal equivalence} ($i \in W^{\mathrm{univ}}$) if for every cocartesian square
\[
\begin{tikzcd}
A' \arrow[d, "i'"'] & A \arrow[l] \arrow[d, "i"] \\
B' & B \arrow[l]
\end{tikzcd}
\quad ,
\]
the arrow $i'$ belongs to $W$.
\end{itemize}
% foot: Version 25 janvier 2004

\printedpage{60}
% head: A. Grothendieck : Dérivateurs, Ch. XIII, Édition des Universités de Montpellier II et Paris VII

\begin{itemize}
\item[\textbullet] \textit{Left pre-Q-category of models} (resp. \textit{right}):

$\Mo$ and $W$, $C$, $F$, $TC$, $TF$ $\subset$ $\Fl(\Mo)$ satisfying axioms \textbf{M~0} to \textbf{M~5}, $\mathbf{M}'$, $\mathbf{M}''$ and \Mzg{} (resp. \Mzd).

\item[\textbullet] \textit{Q-category of models}, \textit{Quillen triple}:

model category in the sense of Quillen.

\item[\textbullet] \textit{``Closed'' Q-category of models}, \textit{closed Quillen triple}, \textit{closed Quillen category}, \textit{closed Quillen system}:

closed model category in the sense of Quillen.

\item[\textbullet] \textit{Left pre-Quillen} (resp. \textit{right}):

left pre-Q-category of models (resp. right).

\item[\textbullet] \textit{Closed Quillen pre-quintuple}:

$W$, $C$, $F$, $TC$, $TF$ $\subset$ $\Fl(\Mo)$ satisfying axioms \textbf{M~0} to \textbf{M~5}, $\mathbf{M}'$, $\mathbf{M}''$, $C$, $TC$ being left closed and $F$, $TF$ right closed.

\item[\textbullet] \textit{Weak Quillen system}:

$W$, $C$, $F$, $TC$, $TF$ $\subset$ $\Fl(\Mo)$ satisfying axioms \textbf{M~0} to \textbf{M~5}, $\mathbf{M}'$, $\mathbf{M}''$.

\item[\textbullet] \textit{Left exact weak Quillen system} (resp. \textit{right}):

$W$, $C$, $F$, $TC$, $TF$ $\subset$ $\Fl(\Mo)$ satisfying axioms \textbf{M~0} to \textbf{M~5}, $\mathbf{M}'$, $\mathbf{M}''$ and \Mzg{} (resp. \Mzd).
\end{itemize}

\section*{Section 3}

\begin{itemize}
\item[\textbullet] $(\Mo,W)$ \textit{category of models}:

$\Mo$ a category satisfying axiom \textbf{M~0};

$W$ $\subset$ $\Fl(\Mo)$ weakly saturated.

\item[\textbullet] \textit{Weak Quillen structure associated with the category of models} $(\Mo,W)$:

quadruple $(C,F,TC,TF)$ of subsets of $\Fl(\Mo)$ satisfying axioms \textbf{M~1} to \textbf{M~5}, $\mathbf{M}'$, $\mathbf{M}''$.

\item[\textbullet] \textit{Exact weak Quillen structure}:

weak Quillen structure satisfying in addition axiom \textbf{M~2\,bis}.

\item[\textbullet] \textit{Left exact weak Quillen structure} (resp. \textit{right}):

weak Quillen structure satisfying in addition axiom \Mzg{} (resp. \Mzd).
\end{itemize}
% foot: Version 25 janvier 2004

\printedpage{61}
% head: A. Grothendieck : Dérivateurs, Ch. XIII, Édition des Universités de Montpellier II et Paris VII

\begin{itemize}
\item[\textbullet] \textit{Left closed weak Quillen structure} (resp. \textit{right}):

weak Quillen structure $(C,F,TC,TF)$ such that $C$, $TC$ (resp. $F$, $TF$) are left closed (resp. right closed).

\item[\textbullet] \textit{Closed weak Quillen structure}:

weak Quillen structure both left and right closed.

\item[\textbullet] \textit{Quillen structure}:

``biexact'' weak Quillen structure:

weak Quillen structure both left and right exact

$=$ proper Quillen model category.

\item[\textbullet] \textit{Perfect Quillen structure}:

``biexact'' and closed weak Quillen structure

$=$ proper closed Quillen model category.
\end{itemize}

\section*{Section 5}

\begin{itemize}
\item[\textbullet] Category of models $(\Mo,W)$ \textit{Quillenizable}:

there exists an \textit{exact} weak Quillen structure $(C,F,TC,TF)$ associated with $(\Mo,W)$.

\item[\textbullet] $(\Mo,W)$ \textit{Quillenizable on} Diag $\subset$ Cat:

for every $I$ in Diag, the category of models $(\Mo(I),\, W(I))$ is Quillenizable.
\end{itemize}

\section*{Section 6}

\begin{itemize}
\item[\textbullet] Category of models $(\Mo,W)$ \textit{uniformly left Quillenizable on} Diag:

there exists a \textit{closed} weak Quillen structure $(C,F,TC,TF)$ for $(\Mo,\, W)$ such that for every $I$ in Diag, if one defines $C(I)$ and $TC(I)$ ``argument by argument'' (\footnote{A morphism of functors $u : X \longrightarrow Y$ from $I^{\circ}$ to $\Mo$ being in $C(I)$ (resp. $TC(I)$) if and only if for every object $i$ of $I$, $u_i$ is in $C$ (resp. $TC$).}), and $F(I)$ and $TF(I)$ by the right lifting property with respect to $TC(I)$ and $C(I)$ respectively (\footnote{A morphism of $\Mo(I)$ being in $F(I)$ (resp. $TF(I)$) if and only if it has the right lifting property with respect to the arrows belonging to $TC(I)$ (resp. $C(I)$).}), the quadruple $(C(I),TC(I),F(I),TF(I))$ is a \textit{closed exact} weak Quillen structure for $(\Mo(I),\, W(I))$. One then says that $(C,F,TC,TF)$ is a \textit{uniform left Quillenization} of $(\Mo,\, W)$ on Diag).
\end{itemize}
% foot: Version 25 janvier 2004

\printedpage{62}
% head: A. Grothendieck : Dérivateurs, Ch. XIII, Édition des Universités de Montpellier II et Paris VII

\begin{itemize}
\item[\textbullet] Category of models $(\Mo,W)$ \textit{Quillenizable on} Diag \textit{uniformly and exactly on the left}:

$(\Mo,W)$ admits a uniform left Quillenization on Diag $(C,F,TC,TF)$ such that $TC = C \cap W$ (in other words, such that $(C,F,TC,TF)$ is a \textit{left exact} weak Quillen structure). One then says that $(C,F,TC,TF)$ is a \textit{uniform and left exact Quillenization} of $(\Mo,\, W)$ on Diag. % [sic] strucure -> structure, letter t omitted
\DERmark{ENV0633}

\end{itemize}

\section*{Section 7}

\begin{itemize}
\item[\textbullet] $I$ a filtering ordered set, \quad $\pi$ a cardinal

$I$ \textit{large relative to} $\pi$:

every subset of $I$ of cardinality less than or equal to $\pi$ admits an upper bound in $I$.

\item[\textbullet] $E$ a category

A \textit{cardinal filtration} of $E$ is an increasing filtration
\[
\Big(\Filt^{\pi}(E)\Big)_{\pi\geq\pi_0} \qquad , \qquad\qquad \pi_0 \text{ a fixed infinite cardinal,}
\]
of $E$ by strictly full subcategories $\Filt^{\pi}(E)$, indexed by the cardinals $\pi \geq \pi_0$ belonging to a given universe, and satisfying the following conditions:
\begin{itemize}
\item[\textit{a})] for every $\pi \geq \pi_0$, $\Filt^{\pi}(E)$ is equivalent to a small category;
\item[\textit{b})] the category $E$ admits filtering direct limits indexed by ordered sets $I$ large relative to $\pi_0$, and for every $\pi \geq \pi_0$, $\Filt^{\pi}(E)$ is stable under the said limits if $\mathrm{card}\, I \leq \pi$;
\item[\textit{c})] For every $\pi \geq \pi_0$, and for every $X \in \mathrm{Ob}\, E$, there exists an isomorphism
\[
X \ \simeq\ \varinjlim_{i\in I} X_i \ ,
\]
where $(X_i)_{i\in I}$ is a filtering direct system in $E$, indexed by an ordered set $I$ large relative to $\pi$, and where the $X_i$ are objects of $\Filt^{\pi}(E)$. Moreover, if $X \in \mathrm{Ob}\,\Filt^{\pi'}(E)$, $\pi' \geq \pi$, one can take $I$ such that $\mathrm{card}\, I \leq \pi'^{\pi}$.
\end{itemize}

\item[\textbullet] $E$ a category, \quad $(\Filt^{\pi}(E))_{\pi\geq\pi_0}$ a cardinal filtration of $E$, \quad $\pi \geq \pi_0$ a cardinal

A \textit{well-adapted direct system} (to the given cardinal filtration, for the given cardinals $\pi_0$, $\pi$) is a filtering direct system $(X_i)_{i\in I}$ satisfying the following conditions:
\begin{itemize}
\item[\textit{a})] $I$ is large relative to $\pi$, and for every $i \in I$, $X_i \in \Filt^{\pi}(E)$;
\item[\textit{b})] Every subset $J$ of $I$ of cardinality $\leq \pi$, and large relative to $\pi_0$, admits a least upper bound $i_J$, and we have
\[
X_{i_J} \ \xleftarrow{\ \sim\ }\ \varinjlim_{i\in J} X_i \ .
\]
% foot: Version 25 janvier 2004
\end{itemize}

\printedpage{63}
% head: A. Grothendieck : Dérivateurs, Ch. XIII, Édition des Universités de Montpellier II et Paris VII

If $X$ is an object of $E$, a \textit{well-adapted representation} of $X$, in terms of objects of $\Filt^{\pi}(E)$, is the datum of a well-adapted direct system $(X_i)_{i\in I}$, to the given cardinal filtration, for the cardinals $\pi_0$, $\pi$, and of an isomorphism
\[
X \ \xleftarrow{\ \sim\ }\ \varinjlim_{i\in I} X_i \ .
\]

\item[\textbullet] $E$ a category admitting small filtering direct limits

A \textit{strict cardinal filtration} of $E$ is a cardinal filtration $(\Filt^{\pi}(E))_{\pi\geq\pi_0}$ such that for every cardinal $\pi \geq \pi_0$, $\Filt^{\pi}(E)$ is stable under filtering direct limits of cardinality $\leq\pi$.
\DERmark{ENV0634}


\item[\textbullet] $E$ a category admitting small filtering direct limits

$(\Filt^{\pi}(E))_{\pi\geq\pi_0}$ a strict cardinal filtration of $E$, \quad $\pi \geq \pi_0$ a cardinal

A \textit{perfectly adapted direct system} (to the given cardinal filtration, for the cardinal $\pi$) is a filtering direct system $(X_i)_{i\in I}$ satisfying the following conditions:
\begin{itemize}
\item[\textit{a})] $I$ is large relative to $\pi$, and for every $i \in I$, $X_i \in \Filt^{\pi}(E)$;
\item[\textit{b})] Every filtering subset $J$ of $I$ of cardinality $\leq \pi$, admits a least upper bound $i_J$, and we have
\[
X_{i_J} \ \xleftarrow{\ \sim\ }\ \varinjlim_{i\in J} X_i \ .
\]
\end{itemize}
If $X$ is an object of $E$, a \textit{perfectly adapted representation} of $X$, in terms of objects of $\Filt^{\pi}(E)$, is the datum of a perfectly adapted direct system $(X_i)_{i\in I}$, to the given cardinal filtration, for the cardinal $\pi$, and of an isomorphism
\DERmark{ENV0635}

\[
X \ \xleftarrow{\ \sim\ }\ \varinjlim_{i\in I} X_i \ .
\]
\end{itemize}
% foot: Version 25 janvier 2004


\setcounter{footnote}{\derfnsaved}%
\endgroup
% END UNIT Der-ChXIII-25-01-04/Der-ChXIII-25-01-04_p061-065.tex
% BEGIN UNIT Der-ChXIII-25-01-04/Der-ChXIII-25-01-04_p066-066.tex
\begingroup\makeatletter%
\newcommand{\printedpage}[1]{\marginpar{\footnotesize p.~#1}}%
\newcommand{\msp}[1]{\marginpar{\footnotesize [#1]}}%
\makeatother\let\printedpage\derpp\let\msp\dermsp%
\let\mspin\DERmspin%


\printedpage{64}
% head: A. Grothendieck : Dérivateurs, Ch. XIII, Édition des Universités de Montpellier II et Paris VII
% foot: Version 25 janvier 2004

\section*{References}

\noindent H. J. Baues, ``Algebraic Homotopy'', Cambridge studies in advanced mathematics, vol. 15 (1989).

\noindent R. Hartshorne, ``Residues and Duality'', Lecture Notes in Mathematics, vol. 20, Springer-Verlag (1966).

\noindent D. G. Quillen, ``Homotopical Algebra'', Lecture Notes in Mathematics, vol. 43, Springer-Verlag (1967).

\noindent R. W. Thomason, \textit{Cat as a Closed Model Category}, Cahiers de Topologie et Géométrie Différentielle, Vol. XXI-3, pp. 305-324 (1980).

\noindent J.-L. Verdier, ``Des catégories dérivées des catégories abéliennes'', Astérisque, vol. 239 (1996).

\noindent SGA 4, I, ``Théorie des Topos et Cohomologie Étale des Schémas'', Séminaire de Géométrie Algébrique du Bois-Marie, directed by M. Artin, A. Grothendieck, J.-L. Verdier, Lecture Notes in Mathematics, vol. 269, Springer-Verlag (1972).


\endgroup
% END UNIT Der-ChXIII-25-01-04/Der-ChXIII-25-01-04_p066-066.tex
\clearpage\setcounter{footnote}{0}\phantomsection\addcontentsline{toc}{section}{Chapter XIV}\markboth{Chapter XIV}{Chapter XIV}
% BEGIN UNIT Der-ChXIV-07-10-10/Der-ChXIV-07-10-10_p001-005.tex
\begingroup\makeatletter%
\newcommand{\printedpage}[1]{\marginpar{\footnotesize p.~#1}}%
\newcommand{\msp}[1]{\marginpar{\footnotesize [#1]}}%
\makeatletter%
\renewcommand\@makefnmark{\hbox{\textsuperscript{\normalfont(\@thefnmark)}}}%
\renewcommand\@makefntext[1]{\parindent 1em\noindent\hb@xt@1.8em{\hss\textsuperscript{\normalfont\@thefnmark}}#1}%
\makeatother%
\makeatother\let\printedpage\derpp\let\msp\dermsp%
\let\mspin\DERmspin%


% ===================== PDF page 1 =====================
\printedpage{}% no page number printed on this page
% head: (none)

\begin{center}
\vspace*{6cm}

{\LARGE \textsf{THE DERIVATORS}}

\vspace{1cm}

{\large \textbf{ALEXANDRE GROTHENDIECK}}

\vspace{2cm}

{\large \textbf{Chapter XIV}}

\vspace{0.3cm}

{\large \textbf{h-cartesian and h-cocartesian squares}}
\end{center}

\par

% ===================== PDF page 2 =====================
\printedpage{}% no page number printed on this page
% head: (none)

\vspace*{6cm}

\noindent \par\noindent\textbf{Note by the editors of the transcribed edition.}\par
This text has been deciphered and transcribed in \LaTeXe{} by M. Künzer. It has been edited by M. Künzer, J. Malgoire, and G. Maltsiniotis. The transcription is as faithful as possible to the manuscript. For the few obvious corrections, or rare comments by the editors, as well as for the original numbering of the manuscript pages, typewriter characters \texttt{[typewriter]} in square brackets are used. A question mark in square brackets means that one is not sure of the preceding word.

\noindent This chapter is unfinished. In the first section, Grothendieck introduces a notion of left (resp. right) ``quadrangulated category''. The text breaks off abruptly at the beginning of the second section, devoted to examples, when Grothendieck realizes, probably, that the main example he has in mind (the homotopy category of a category of models equipped with homotopically cartesian (resp. cocartesian) squares) does \textit{not} satisfy the axiomatics of quadrangulated categories. This is briefly explained in an afterword by the editors.

\noindent This edition is provisional. Remarks, comments and corrections are welcome. Send a message to:

\noindent \textbf{[e-mail address removed]}

\hfill G. Maltsiniotis\hspace*{2cm}

\par

% ===================== PDF page 3 =====================
\printedpage{1}
% head: A. Grothendieck : Dérivateurs, Ch. XIV, Édition des Universités de Montpellier II et Paris VII
% foot: Version 7 octobre 2010

\msp{page 1}

\section*{1\quad Axioms for quadrangulated structures (an aborted attempt)}

Let $(\mathcal{M},W)$ be a category of models admitting a theory of the homotopy fiber. Then we have in $\mathrm{Ho}_W=\mathcal{M}W^{-1}$ a notion of \textit{$W$-cartesian square} (or $W$-homotopically cartesian), by which we mean a square diagram
\[
\begin{tikzcd}
X_1 \arrow[d] & X_3 \arrow[l] \arrow[d] \\
X_0 & X_2 \arrow[l]
\end{tikzcd}
\]
in $\mathrm{Ho}_W$ coming, up to isomorphism, from a $W$-cartesian square in $\mathcal{M}$ – thus isomorphic, in $\underline{\mathrm{Hom}}(\boldsymbol{\Delta}_1\times\boldsymbol{\Delta}_1,\mathrm{Ho}_W)$, to a \textit{cartesian} square \texttt{[in $\mathcal{M}$]} % [sic] a un carré -> à un carré, preposition printed without its grave accent
\[
\begin{tikzcd}
Y_1 \arrow[d, "p_1"'] & Y_3\ (\simeq Y_1\times_{Y_0}Y_2) \arrow[l] \arrow[d] & \\
Y_0 & Y_2 \arrow[l, "p_2"] & ,
\end{tikzcd}
\]
with $p_1,\ p_2\in\mathrm{Fib}_W$. (In fact, it suffices to have $p_1$ or $p_2$ in $\mathrm{Fib}_W$ to ensure the $W$-cartesian character of the cartesian square under consideration.) I would like to axiomatize this situation, mainly to make a list of the essential useful properties of this notion, and of the dual notion of $W$-cocartesian square (or $W$-homotopically cocartesian), in the case where $(\mathcal{M},W)$ admits a theory of the homotopy cofiber. The species of structure, namely that of ``quadrangulated categories'' (left, right, two-sided) is a variant, in a \msp{page 2} nonadditive case, of the notion of triangulated category isolated by \textsc{Verdier}. Like the latter, this notion is in fact too weak – it is \textit{not} ``the right one'' – but it is of provisional practical interest. The ``right'' notion (if there is one apart from that of derivator) should work with a finer notion of square and of $W$-cartesian square, working systematically in $\mathrm{HO}_W(\boldsymbol{\Delta}_1\times\boldsymbol{\Delta}_1)=\mathcal{M}(\boldsymbol{\Delta}_1\times\boldsymbol{\Delta}_1)W(\boldsymbol{\Delta}_1\times\boldsymbol{\Delta}_1)^{-1}$ (where the ``real'' $W$-cartesian squares live!), and similar categories, such as $\mathrm{HO}_W(\Psi)$, $\mathrm{HO}_W(\Phi)$. But to include them in an axiomatics making no mention of a category of models $(\mathcal{M},W)$, manifestly amounts to putting one's finger in the gears of the notion of derivator, and that is not my purpose for the moment.

\noindent\textbf{Definition 1.} I call a \textit{``category with h-cartesian squares'',} or a \textit{``left quadrangulated category'',} a category $\mathcal{H}$, equipped with a subset Cart of the set $\mathrm{Hom}(\boldsymbol{\Delta}_1\times\boldsymbol{\Delta}_1,\mathcal{H})$
\printedpage{2}
% head: A. Grothendieck : Dérivateurs, Ch. XIV, Édition des Universités de Montpellier II et Paris VII
% foot: Version 7 octobre 2010
of square diagrams \texttt{[by square we mean commutative square]}
\[
\begin{tikzcd}
X_1 \arrow[d] & X_3 \arrow[l] \arrow[d] \\
X_0 & X_2 \arrow[l]
\end{tikzcd}
\]
in $\mathcal{H}$, squares called \textit{``h-cartesian''} (h as in ``homotopically''), this \msp{page 3} datum being subject to the following axioms, Cart (1) to Cart (\ldots).

\noindent Since I am not sure I know all the essential properties in advance, the number of these axioms remains open for the moment.

\noindent\textbf{Cart (1) (Symmetry).} If a square
\[
\begin{tikzcd}
X_1 \arrow[d] & X_3 \arrow[l] \arrow[d] \\
X_0 & X_2 \arrow[l]
\end{tikzcd}
\]
is h-cartesian, likewise the symmetric square (obtained by interchanging the roles of $X_1$ and $X_2$).

\noindent\textbf{Cart (2) (Trivial squares).} Let
\[
\begin{tikzcd}
X_1 \arrow[d, "p_1"'] & X_3 \arrow[l, "q_2"'] \arrow[d, "q_1"] \\
X_0 & X_2 \arrow[l, "p_2"]
\end{tikzcd}
\]
be a square in $\mathcal{H}$, with $p_1$ an isomorphism. Then the square is h-cartesian if and only if $q_1$ is an isomorphism.

\noindent\textbf{Cart (3) (Composite squares).} Let there be a commutative diagram
\[
\begin{tikzcd}
X'_0 \arrow[d] \arrow[dr, phantom, "\mathrm{I}"] & X'_1 \arrow[l] \arrow[d] \arrow[dr, phantom, "\mathrm{II}"] & X'_2 \arrow[l] \arrow[d] & \\
X_0 & X_1 \arrow[l] & X_2 \arrow[l] & ,
\end{tikzcd}
\]
whence the squares I, II and III:
\[
\begin{tikzcd}
X'_0 \arrow[d] & X'_2 \arrow[l] \arrow[d] \\
X_0 & X_2 \arrow[l]
\end{tikzcd}
\]
\printedpage{3}
% head: A. Grothendieck : Dérivateurs, Ch. XIV, Édition des Universités de Montpellier II et Paris VII
% foot: Version 7 octobre 2010
(which we call the \textit{composite square} (``horizontal'') of the squares I and II). Suppose I h-cartesian. Then II is h-cartesian if and only if III is \texttt{[see editors' afterword]}.

\noindent\textbf{Corollary.} \textit{Let there be a commutative cubical diagram, which we can regard as a homomorphism from one square} (\textit{upper face of the cube}) \textit{to another} (\textit{lower face}):
\[
\begin{tikzcd}
 & {} \arrow[dl] \arrow[dd] & & {} \arrow[ll] \arrow[dl] \arrow[dd] \\
{} \arrow[dd] & & {} \arrow[ll, crossing over] & \\
 & {} \arrow[dl] & & {} \arrow[ll] \arrow[dl] \\
{} & & {} \arrow[ll] \arrow[from=uu, crossing over] & .
\end{tikzcd}
\]
\textit{If all the faces of the cube}\footnote{In fact, the lower face, and the two parallel lateral faces (e.g. the front one and the back one) suffice.} \textit{except the upper face are} \msp{page 4} \textit{h-cartesian squares (namely the lower face and the four lateral faces), the same holds for the upper face.}\footnote{\textbf{Corollary} (of Cart (3) and Cart (2)): \textit{every square isomorphic to an h-cartesian square is h-cartesian.}}

\noindent These axioms are in the nature of stability axioms, giving us criteria asserting that certain squares are h-cartesian, if such and such conditions are fulfilled (e.g. that certain other squares are already h-cartesian). I have the impression of having exhausted here the relevant properties of this type (\textit{i.e.} that the others should be deducible from them). At present, I am dealing with properties of existence and uniqueness of h-cartesian squares. % [sic] carrés h-cartésien -> carrés h-cartésiens, missing plural agreement

\noindent\textbf{Cart (4) (Fiber product axiom).} Every diagram
\[
\begin{tikzcd}
X_1 \arrow[d] & \\
X_0 & X_2 \arrow[l]
\end{tikzcd}
\]
in $\mathcal{H}$ can be completed to an h-cartesian square
\[
\begin{tikzcd}
X_1 \arrow[d] & X_3 \arrow[l] \arrow[d] & \\
X_0 & X_2 \arrow[l] & .
\end{tikzcd}
\]
The following axiom (\textit{cf.} its corollary) will ensure among other things that the h-cartesian square in question is \textit{unique up to isomorphism} – but one should note that it is in general not unique ``up to unique isomorphism''.


\endgroup
% END UNIT Der-ChXIV-07-10-10/Der-ChXIV-07-10-10_p001-005.tex
% BEGIN UNIT Der-ChXIV-07-10-10/Der-ChXIV-07-10-10_p006-010.tex
\begingroup\makeatletter%
\newcommand{\printedpage}[1]{\marginpar{\footnotesize p.~#1}}%
\newcommand{\msp}[1]{\marginpar{\footnotesize [#1]}}%
\newcommand{\uHom}{\underline{\mathrm{Hom}}}%
\newcommand{\bDelta}{\boldsymbol{\Delta}}%
\newcommand{\ho}{\mathrm{ho}}%
\makeatother\let\printedpage\derpp\let\msp\dermsp%
\let\mspin\DERmspin%


\printedpage{4}
% head: A. Grothendieck : Dérivateurs, Ch. XIV, Édition des Universités de Montpellier II et Paris VII

\textbf{Cart (5) (Exactness axiom)} Let us consider two h-cartesian squares
\[
X = \left(\begin{tikzcd}
X_1 \ar[d] & X_3 \ar[l] \ar[d] \\
X_0 & X_2 \ar[l]
\end{tikzcd}\right)
\qquad \text{and} \qquad
Y = \left(\begin{tikzcd}
Y_1 \ar[d] & Y_3 \ar[l] \ar[d] \\
Y_0 & Y_2 \ar[l]
\end{tikzcd}\right)
\]
and a homomorphism $p = (p_0, p_1, p_2)$ between the diagrams of type
\[
\Phi \ = \
\begin{tikzcd}
1 \ar[d] & \\
0 & 2 \ar[l]
\end{tikzcd}
\]
associated with them (\textit{cf.} the solid-line diagram below):

\msp{page 5}
\[
\begin{tikzcd}[row sep=small, column sep=small]
& X_1 \ar[dl] \ar[dd, "p_1" near start] & & X_3 \ar[ll] \ar[dl] \ar[dd, dotted, "p_3"] \\
X_0 \ar[dd, "p_0"'] & & X_2 \ar[ll, crossing over] & \\
& Y_1 \ar[dl] & & Y_3 \ar[ll] \ar[dl] \\
Y_0 & & Y_2 \ar[ll] \ar[from=uu, crossing over, "p_2"' near end] & .
\end{tikzcd}
\]

Then:
\begin{itemize}[leftmargin=2.5em]
\item[a)] There exists a homomorphism $\bar{p} : X \longrightarrow Y$ inducing $p$, \textit{i.e.} there exists an arrow $p_3$ so as to complete the commutative cubical diagram (thus respecting the commutativity constraint for the two squares adjacent to the edge $p_3$).
\item[b)] If moreover $p_0$, $p_1$, $p_2$ are isomorphisms, the same is true of $p_3$ \texttt{[see editors' afterword]}.
\end{itemize}

Taking the case where $p_0$, $p_1$, $p_2$ are identities, we find:

\textbf{Corollary 1.} \textit{Under the conditions of} Cart (4)\textit{, if we have} two \textit{h-cartesian squares $X$, $X'$ which complete}
\[
\begin{tikzcd}[row sep=small]
& X_1 \ar[dl] & \\
X_0 & & , \\
& X_2 \ar[ul] &
\end{tikzcd}
\]
\textit{there exists a homomorphism from $X$ to $X'$ which induces the identity on the vertices with indices} 0, 1, 2\textit{, and such a homomorphism is always an isomorphism.}

But this isomorphism is not necessarily unique, in other words, an automorphism of an h-cartesian square can induce the identity on the vertices $X_0$, $X_1$, $X_2$, without being the identity.
% foot: Version 7 octobre 2010

\printedpage{5}
% head: A. Grothendieck : Dérivateurs, Ch. XIV, Édition des Universités de Montpellier II et Paris VII

\msp{page 6}
\textbf{Corollary 2.} \textit{Let us consider a commutative square}
\[
Y = \left(\begin{tikzcd}
Y_1 \ar[d] & Y_3 \ar[l] \ar[d] \\
Y_0 & Y_2 \ar[l]
\end{tikzcd}\right) \quad ,
\]
\textit{and an object $T$ of $\mathcal{H}$, whence (putting} $\mathrm{Hom}(T,X) = X(T)$\textit{) a commutative square in} (Ens)
\[
\begin{tikzcd}
Y_1(T) \ar[d] & Y_3(T) \ar[l] \ar[d] \\
Y_0(T) & Y_2(T) \ar[l]
\end{tikzcd} \quad ,
\]
\textit{and hence}
\[
Y_3(T) \longrightarrow Y_1(T) \times_{Y_0(T)} Y_2(T) \quad . \tag*{$(*)$}
\]
\textit{If $Y$ is h-cartesian, this map} $(*)$ \textit{is} surjective \textit{whatever $T$ may be.}

We apply Cart (5) to $Y$ and to the constant square $X$ with value $T$
\[
X = \quad
\begin{tikzcd}
T \ar[d, "\mathrm{id}"'] & T \ar[l, "\mathrm{id}"'] \ar[d, "\mathrm{id}"] \\
T & T \ar[l, "\mathrm{id}"']
\end{tikzcd} \quad ,
\]
noting that giving $p = (p_0, p_1, p_2)$ from
\[
\begin{tikzcd}
T \ar[d] & \\
T & T \ar[l]
\end{tikzcd}
\qquad\qquad \text{to} \qquad\qquad
\begin{tikzcd}
Y_1 \ar[d] & \\
Y_0 & Y_2 \ar[l]
\end{tikzcd}
\]
amounts to giving an element of the right-hand side of $(*)$,
\[
\begin{tikzcd}
Y_1 \ar[d] & T \ar[l, "p_1"'] \ar[dl, "p_0"] \ar[d, "p_2"] \\
Y_0 & Y_2 \ar[l]
\end{tikzcd} \quad ,
\]
and that the elements of $Y_3(T)$ which map (by $(*)$) to this element are none other than the $p_3$ whose existence axiom Cart (5) asserts.

\textbf{N.B.} We shall note that the surjective map $(*)$ is in general \textit{not} bijective, \textit{i.e.} the h-cartesian square is in general not cartesian.
% foot: Version 7 octobre 2010

\printedpage{6}
% head: A. Grothendieck : Dérivateurs, Ch. XIV, Édition des Universités de Montpellier II et Paris VII

\msp{page 7}
\textbf{N.B.} There is a dual notion of \textit{category with h-cocartesian squares,} or \textit{right quadrangulated category.} Most often, the two structures are present simultaneously. We should bring out properties which relate the two dual structures on one and the same category. Thus, unless I am mistaken, in the case $\mathcal{H}$ additive, the two kinds of squares, the h-cartesian and the h-cocartesian ones, should \textit{coincide.} (I shall have to come back to this later.) In the case $\mathcal{H}$ ``pointed'', there is the adjunction of the functors $\Sigma$, $\Omega$ linked, one to the right quadrangulated structure, the other to the left quadrangulated structure. And there are, when there are internal or external \underline{Hom}, possible properties of $\uHom(X,Y)$ with respect to h-cocartesian quadrangles in $X$, h-cartesian in $Y$ -- but this is hardly a link between the two structures \ldots

\section*{2\quad Examples}

A first trivial example, is the one where $\mathcal{H}$ is stable under fiber products, and where the h-cartesian squares are defined as the cartesian squares. Axioms Cart (1) to Cart (5) are then well-known, $\pm$ tautological properties of the notion of a cartesian square. And
\msp{page 8}
in a sense, the general case must reduce to this banal special case.

The only type of other case that I know of, is the one where $\mathcal{H} = \mathcal{M}W^{-1}$, where $W$ is a localizer admitting a theory of the $W$-homotopy fiber (\textit{cf.} chapter XII). In this case, as we have said, we take as h-cartesian squares in $\mathcal{H}$ those which are isomorphic, in $\uHom(\bDelta_1 \times \bDelta_1, \mathcal{H})$, to a $W$-cartesian square in $\mathcal{M}$. We still have to verify axioms Cart (1) to Cart (5). I think that I have little chance of getting there except by means of a

\textbf{Proposition 1.} \textit{Under hypotheses on $(\mathcal{M}, W)$ which will be brought out below, satisfied in any case if $(\mathcal{M}, W)$ is Quillenizable, on the category of finite ordered sets} (\textit{cf.} XIII, p. 51)\textit{, we have this: let}
\[
X = \left(\begin{tikzcd}
X_1 \ar[d] & X_3 \ar[l] \ar[d] \\
X_0 & X_2 \ar[l]
\end{tikzcd}\right)
\]
\textit{be a square in $\mathcal{M}$, and let} $\ho_W(X)$ \textit{be its image in} $\uHom(\bDelta_1 \times \bDelta_1, \mathcal{H})$\textit{. For} $\ho_W(X)$ \textit{to be h-cartesian, it is} necessary (\textit{and it is sufficient, of course}) \textit{that $X$ be $W$-cartesian} \texttt{[see editors' afterword]}\textit{.}
% foot: Version 7 octobre 2010

\printedpage{7}
% head: A. Grothendieck : Dérivateurs, Ch. XIV, Édition des Universités de Montpellier II et Paris VII

\section*{Editors' afterword}

The proposition above is false. We recall that if $\mathcal{M}$ is a Quillen closed model category with class of weak equivalences $W$, a commutative square
\[
X = \left(\begin{tikzcd}
X_1 \ar[d] & X_3 \ar[l] \ar[d] \\
X_0 & X_2 \ar[l]
\end{tikzcd}\right)
\]
in $\mathcal{M}$ is $W$-cartesian if and only if there exists a \textit{cartesian} square
\[
Y = \left(\begin{tikzcd}
Y_1 \ar[d] & Y_3 \ar[l] \ar[d] \\
Y_0 & Y_2 \ar[l]
\end{tikzcd}\right)
\]
with $Y_1 \longrightarrow Y_0$ a fibration, and $Y_0$ and $Y_2$ fibrant (this last condition being superfluous if the model category $\mathcal{M}$ is right proper), and a morphism of squares $s : X \longrightarrow Y$, $s = (s_0, s_1, s_2, s_3)$, \textit{i.e.} a commutative cube
\[
\begin{tikzcd}[row sep=small, column sep=small]
& X_1 \ar[dl] \ar[dd, "s_1" near start] & & X_3 \ar[ll] \ar[dl] \ar[dd, "s_3"] \\
X_0 \ar[dd, "s_0"'] & & X_2 \ar[ll, crossing over] & \\
& Y_1 \ar[dl] & & Y_3 \ar[ll] \ar[dl] \\
Y_0 & & Y_2 \ar[ll] \ar[from=uu, crossing over, "s_2"' near end] & ,
\end{tikzcd}
\]
such that $s$ is a weak equivalence argument by argument, \textit{i.e.} $s_0, s_1, s_2, s_3 \in W$. It is equivalent to require that for \textit{every} commutative diagram with weak equivalences as vertical arrows
\[
\begin{tikzcd}
X_1 \ar[r] \ar[d] & X_0 \ar[d] & X_2 \ar[l] \ar[d] \\
Y_1 \ar[r] & Y_0 & Y_2 \ar[l]
\end{tikzcd}
\]
such that $Y_1 \longrightarrow Y_0$ is a fibration, and $Y_0$, $Y_2$ fibrant objects, the canonical arrow $X_3 \longrightarrow Y_1 \times_{Y_0} Y_2$ be a weak equivalence.

Let us now suppose that the model category $\mathcal{M}$ is pointed, in other words that it admits an object $*$ both initial and final, and let $X$ be a fibrant object of $\mathcal{M}$. Let us choose a decomposition of the arrow $* \longrightarrow X$, into a weak equivalence followed by a fibration
\[
* \longrightarrow C \longrightarrow X \quad ,
\]
and form the cartesian square
% foot: Version 7 octobre 2010
\printedpage{8}
% head: A. Grothendieck : Dérivateurs, Ch. XIV, Édition des Universités de Montpellier II et Paris VII
\[
\begin{tikzcd}
C \ar[d] & \Omega X \ar[l] \ar[d] \\
X & * \ar[l]
\end{tikzcd} \quad ,
\]
where by definition $\Omega X = C \times_X *$ is the \textit{loop space} of $X$. This square is $W$-cartesian. Moreover, it is isomorphic in $\uHom(\bDelta_1 \times \bDelta_1,\ \mathcal{M}W^{-1})$ to the trivial square
\[
\begin{tikzcd}
* \ar[d] & \Omega X \ar[l] \ar[d] \\
X & * \ar[l]
\end{tikzcd} \quad .
\]

Nevertheless, the latter is \textit{not} in general $W$-cartesian. Indeed, if it were, the unique arrow $s_3 : \Omega X \longrightarrow \Omega X$ in $\mathcal{M}$ making the cube commute
\[
\begin{tikzcd}[row sep=small, column sep=small]
& * \ar[dl] \ar[dd] & & \Omega X \ar[ll] \ar[dl] \ar[dd, dotted, "s_3"] \\
X \ar[dd, "\parallel"'] & & * \ar[ll, crossing over] & \\ % Audit resolved: vertical identity mark on X; see corrections.csv, PDF leaf 10.
& C \ar[dl] & & \Omega X \ar[ll] \ar[dl] \\
X & & * \ar[ll] \ar[from=uu, crossing over] &
\end{tikzcd}
\]
would be a weak equivalence. Now, since the arrow $s_3$ factors through $* \times_X * \simeq *$, this would imply that $\Omega X \to *$ is a weak equivalence. This condition is very restrictive. If, for example, $\mathcal{M}$ is the model category of pointed topological spaces, and $X$ a pointed CW-complex, it is equivalent to the contractibility of the connected component of $X$ containing the base point. Similarly, if $\mathcal{M}$ is the model category of complexes of a Grothendieck abelian category and $X$ a complex, it means that $X$ is acyclic.

Analogous considerations show that the $h$-cartesian squares of $\mathcal{M}W^{-1}$, defined in section 2 as being the commutative squares of $\mathcal{M}W^{-1}$ which are isomorphic in $\uHom(\bDelta_1 \times \bDelta_1,\ \mathcal{M}W^{-1})$ to a $W$-cartesian square, do not satisfy all the conditions of Cart (3) and Cart (5) of section 1. Indeed, let us continue to suppose that $\mathcal{M}$ is pointed, and keep the above notation. Let us consider the commutative diagram in $\mathcal{M}W^{-1}$
\[
\begin{tikzcd}
* \ar[d] & \Omega X \ar[l] \ar[d] & \Omega X \ar[l, "*_{\Omega X}"'] \ar[d] \\
X & * \ar[l] & * \ar[l]
\end{tikzcd} \quad ,
\]
where $*_{\Omega X}$ denotes the unique endomorphism of $\Omega X$ factoring through $*$. It follows from what precedes that the left-hand square and the composite square are $h$-cartesian. On the other hand, the right-hand square is not in general. Indeed, if this square were isomorphic in $\uHom(\bDelta_1 \times \bDelta_1,\ \mathcal{M}W^{-1})$ to a $W$-cartesian square, the arrow $*_{\Omega X}$ would be an isomorphism
% foot: Version 7 octobre 2010


\endgroup
% END UNIT Der-ChXIV-07-10-10/Der-ChXIV-07-10-10_p006-010.tex
% BEGIN UNIT Der-ChXIV-07-10-10/Der-ChXIV-07-10-10_p011-011.tex
\begingroup\makeatletter%
\newcommand{\printedpage}[1]{\marginpar{\footnotesize p.~#1}}%
\newcommand{\msp}[1]{\marginpar{\footnotesize [#1]}}%
\newcommand{\cM}{\mathcal{M}}%
\makeatother\let\printedpage\derpp\let\msp\dermsp%
\let\mspin\DERmspin%


\printedpage{9}
% head: A. Grothendieck : Dérivateurs, Ch. XIV, Édition des Universités de Montpellier II et Paris VII
% foot: Version 7 octobre 2010
of $\cM W^{-1}$, hence a weak equivalence of $\cM$, which would imply the same very restrictive conditions as above. Thus one of the implications of axiom Cart 3 is not satisfied.

Likewise the commutative cube in $\cM W^{-1}$
\[
\begin{tikzcd}[row sep=small, column sep=small]
 & * \arrow[dl] \arrow[dd] & & \Omega X \arrow[ll] \arrow[dl] \arrow[dd, "*_{\Omega X}"] \\
X \arrow[dd, "\shortparallel"'] & & * \arrow[ll, crossing over] & \\
 & * \arrow[dl] & & \Omega X \arrow[ll] \arrow[dl] \\
X & & * \arrow[ll] \arrow[from=uu, crossing over] & \mathllap{,}
\end{tikzcd}
\]
whose horizontal faces are $h$-cartesian, shows that condition $(b)$ of axiom Cart (5) is not satisfied in general.


\endgroup
% END UNIT Der-ChXIV-07-10-10/Der-ChXIV-07-10-10_p011-011.tex
\clearpage\setcounter{footnote}{0}\phantomsection\addcontentsline{toc}{section}{Chapter XV}\markboth{Chapter XV}{Chapter XV}
% BEGIN UNIT Der-ChXV-05-11-10/Der-ChXV-05-11-10_p001-005.tex
\begingroup\makeatletter%
\derleqno%
\newcommand{\printedpage}[1]{\marginpar{\footnotesize p.~#1}}%
\newcommand{\msp}[1]{\marginpar{\footnotesize [#1]}}%
\newcommand{\M}{\mathcal{M}}%
\newcommand{\Fl}{\mathrm{Fl}}%
\newcommand{\uFl}{\underline{\mathrm{Fl}}}%
\newcommand{\uNerf}{\underline{\mathrm{Nerf}}}%
\makeatother\let\printedpage\derpp\let\msp\dermsp%
\let\mspin\DERmspin%


% ===================== PDF page 1 =====================
\printedpage{---}
% no page number printed on this page
% head: (none)

\begin{center}
{\Large\textsf{THE DERIVATORS}}

\bigskip

\textbf{ALEXANDRE GROTHENDIECK}

\vspace{2cm}

\textbf{Chapter XV}

\medskip

\textbf{Factorization theorems}

\textbf{(Models (2))}
\end{center}

\par

% ===================== PDF page 2 =====================
\printedpage{---}
% no page number printed on this page
% head: (none)

\par\noindent\textbf{Note by the editors of the transcribed edition.}\par
This text was deciphered and transcribed in \LaTeXe{} by M. Künzer. It was edited by M. Künzer, J. Malgoire, and G. Maltsiniotis. The transcription is as faithful as possible to the manuscript. Only a few obvious minor corrections have been made. For the rare additions or comments by the editors, as well as for the original numbering of the manuscript pages, typewriter characters \texttt{[typewriter]} in brackets are used.

This edition is provisional. Remarks, comments and corrections are welcome. Send a message to:

\textbf{[e-mail address removed]}

\hfill G. Maltsiniotis\hspace*{2cm}

\par

% ===================== PDF page 3 =====================
\printedpage{1}
% head: A. Grothendieck : Dérivateurs, Ch. XV, Édition des Universités de Montpellier II et Paris VII
% foot: Version 5 novembre 2010

\msp{page 1}

\section*{1\quad Fundamental construction: the transfinite factorization sequence of an arrow $\boldsymbol{f}$.}

Let $\M$ be a $\mathfrak{U}$-category, stable under ``small'' (\textit{i.e.} $\mathfrak{U}$-small) filtering $\varinjlim$. Let
\[
C_0 \;\subset\; \Fl(\M)
\tag*{(1.1)}
\]
be a \textit{small} set of arrows, we propose to construct for every arrow $f$ of $\M$, a factorization
\[
f \;=\; p\,i \,, \qquad i \;\in\; \widetilde{C_0} = ((C_0)_*)^* \;, \;\; p \;\in\; (C_0)_* \;.
\tag*{(1.2)}
\]
For this, we are going to construct an ordinal sequence of functors
\[
\left\{
\begin{array}{rcl}
f & \longmapsto & X_\alpha(f) \;=\; X_\alpha(C_0,f) \\
X_\alpha : \uFl(\M) & \longrightarrow & \M
\end{array}
\right.
\tag*{(1.3)}
\]
and functorial homomorphisms
\[
i_\alpha(f) : \mathrm{source}\, f \;\longrightarrow\; X_\alpha(f) \;,
\tag*{(1.4)}
\]
\[
p_\alpha(f) : X_\alpha(f) \;\longrightarrow\; \mathrm{but}\, f \;,
\tag*{(1.5)}
\]
so that we shall have the diagram (where $X = \mathrm{source}\, f$, $Y = \mathrm{but}\, f$)
\[
\begin{tikzcd}[column sep=small]
X \arrow[rr, "f"] \arrow[dr, "i_\alpha(f)"'] & & Y \\
& X_\alpha(f) \arrow[ur, "p_\alpha(f)"'] &
\end{tikzcd}
\qquad ,
\tag*{(1.6)}
\]
this diagram being commutative
\[
f \;=\; p_\alpha(f) i_\alpha(f) \;.
\tag*{(1.7)}
\]
We can also say that we have a section of the canonical functor of ``composition'', transpose of $\boldsymbol{\Delta}_1 \longrightarrow \boldsymbol{\Delta}_2$, given by $0 \longmapsto 0$, $1 \longmapsto 2$:
\[
\left\{
\begin{array}{l}
\begin{tikzcd}[ampersand replacement=\&]
\uNerf_2(\M) \arrow[r, shift left, "\gamma"] \& \uNerf_1(\M) \arrow[l, dotted, shift left, "\Phi_\alpha"]
\end{tikzcd}
\;=\; \uFl(\M) \;, \\[1ex]
\\[1ex]
\\
\end{array}
\right.
\tag*{(1.8)}
\]
$\Phi$ as in ``factorization'',
\printedpage{2}
% head: A. Grothendieck : Dérivateurs, Ch. XV, Édition des Universités de Montpellier II et Paris VII
% foot: Version 5 novembre 2010
\msp{page 2}
where
\[
\Phi_\alpha(f) \;=\; \{\; X \xrightarrow{i_\alpha(f)} X_\alpha(f) \xrightarrow{p_\alpha(f)} Y \;\}
\tag*{(1.9)}
\]
with $p_\alpha(f)i_\alpha(f) = f$. The functoriality of 1.4 and 1.5, or that of $\Phi_\alpha$, \textit{i.e.} of 1.9, is expressed by the fact that every commutative square
\[
\begin{tikzcd}[row sep=large]
X \arrow[r, "f"] \arrow[d, "u_X"'] & Y \arrow[d, "u_Y"] \\
X' \arrow[r, "{f'}"] & Y'
\end{tikzcd}
\tag*{(1.10)}
\]
in $\M$, interpreted as an arrow $u = (u_X, u_Y)$ from $f$ to $f'$ in $\uFl(\M)$, gives rise to the commutative double square
\[
\begin{tikzcd}[row sep=large]
X \arrow[r, "i_\alpha(f)"] \arrow[d, "u_X"] & X_\alpha(f) \arrow[r, "p_\alpha(f)"] \arrow[d, "X_\alpha(u)"] & Y \arrow[d, "u_Y"] \\
X' \arrow[r, "{i_\alpha(f')}"] & X_\alpha(f') \arrow[r, "{p_\alpha(f')}"] & Y' \;.
\end{tikzcd}
\tag*{(1.11)}
\]
\textbf{N.B.} When we want to make explicit the dependence of these constructions on $C_0$, we shall write
\[
X_\alpha(C_0,f) \;, \qquad i_\alpha(C_0,f) \;, \qquad p_\alpha(C_0,f) \;, \qquad \Phi_\alpha(C_0,f) \;,
\]
or
\[
X_\alpha(C_0,-) \;\;\text{or}\;\; X_\alpha^{C_0} \;, \qquad i_\alpha(C_0,-) \;\;\text{or}\;\; i_\alpha^{C_0} \qquad \textit{etc.} \;,
\]
instead of $X_\alpha(f)$, $i_\alpha(f)$, \dots

Here are the essential properties of the ordinal direct system of the $X_\alpha(f)$ or $\Phi_\alpha(f)$.

\msp{page 3}

\textbf{1.} $X_0(f) = X$, $i_0(f) = \mathrm{id}_X$, $p_0(f) = f$ (\textit{i.e.} the factorization $\Phi_0(f)$ of $f$ is the trivial factorization $X \xrightarrow{\mathrm{id}_X} X \xrightarrow{f} Y$).

\textbf{2.}
\[
\begin{aligned}
X_{\alpha+1}(f) &\;=\; X_1(p_\alpha(f)) \\
i_{\alpha+1}(f) &\;=\; i_1(p_\alpha(f)) \circ i_\alpha(f) \\
p_{\alpha+1}(f) &\;=\; p_1(p_\alpha(f))
\end{aligned}
\]
\[
\begin{tikzcd}
\underbrace{X}_{= \,\mathrm{source}\, f} \arrow[r, "i_\alpha(f)"] \arrow[rrr, bend left=20, "f"] & X_\alpha(f) \arrow[rr, "p_\alpha(f)"] \arrow[dr, bend right=40, "i_1(p_\alpha(f))"'] & & Y \\
& & \underbrace{X_1(p_\alpha(f))}_{= X_{\alpha+1}(f)} \arrow[ur, bend right=40] &
\end{tikzcd}
\]

\par

% ===================== PDF page 5 =====================
\printedpage{3}
% head: A. Grothendieck : Dérivateurs, Ch. XV, Édition des Universités de Montpellier II et Paris VII
% foot: Version 5 novembre 2010

or better
\[
\begin{tikzcd}[row sep=large]
X \arrow[rrr, "f"] \arrow[d, "\mathrm{id}"] & & & Y \arrow[d, "\mathrm{id}"] \\
X \arrow[r, "i_\alpha(f)"] \arrow[d, "\mathrm{id}"] & X_\alpha(f) \arrow[rr, "p_\alpha(f)"] \arrow[d, "\mathrm{id}"] & & Y \arrow[d, "\mathrm{id}"] \\
X \arrow[r] & X_\alpha(f) \arrow[r, "i_1(p_\alpha(f))"] & \underbrace{X_1(p_\alpha(f))}_{= X_{\alpha+1}(f)} \arrow[r, "p_1(p_\alpha(f))"] & Y \;.
\end{tikzcd}
\tag*{(1.12)}
\]
The transition morphism $X_\alpha(f) \longrightarrow X_{\alpha+1}(f)$ is
\[
X_\alpha(f) \;\; \xrightarrow{\;\;i_1(p_\alpha(f))\;\;} \;\; X_1(p_\alpha(f)) \;=\; X_{\alpha+1}(f) \;.
\]

\textbf{3.} If $\alpha$ is a limit ordinal, we have
\[
\left\{
\begin{array}{rcl}
X_\alpha(f) & = & \varinjlim\limits_{\alpha'<\alpha} X_{\alpha'}(f). \\[3ex]
i_\alpha(f) & = & \parbox[t]{8cm}{canonical morphism from $X = X_0(f)$ into the direct limit.} \\[5ex]
p_\alpha(f) & = & \text{inverse limit of the } p_{\alpha'}(f) : X_{\alpha'}(f) \longrightarrow Y.
\end{array}
\right.
\]
(We can regard the direct limit defining $X_\alpha(f)$ as a limit in $X\backslash\M/Y\dots$)

The transition morphisms
\[
X_{\alpha'}(f) \;\longrightarrow\; X_\alpha(f) \qquad (\; \alpha' \;<\; \alpha \;)
\]
are the canonical morphisms relative to the direct system of the $X_{\alpha'}(f)$.

\msp{page 4}

We can therefore say that the ordinal sequence of functorial factorizations $\Phi_\alpha(f)$ of variable $f$, is deduced from the functorial factorization $\Phi_1(f)$ of $f$ as
\[
\begin{tikzcd}
X \arrow[r, "i_1(f)"] \arrow[rr, bend left=20, "f"] & X_1(f) \arrow[r, "p_1(f)"] & Y
\end{tikzcd}
\]
by ``transfinite iteration on the right'', \textit{i.e.} by regarding the operations $\Phi_\alpha$ as functors
% footnote mark (1) printed beside the following display; \footnotemark/\footnotetext used because \footnote does not work inside math
\[
\left\{
\begin{array}{rcl}
p_\alpha :\; \uFl(\M) & \longrightarrow & \uFl(\M) \\
f & \longmapsto & p_\alpha(f) \;,
\end{array}
\right.
\qquad\qquad \text{\footnotemark}
\]
\footnotetext{\textbf{N.B.} $p_\alpha$ respects the category structures \textit{over} $\M$, via the target functors: $\uFl(\M) \longrightarrow \M$, \textit{i.e.} $\uFl(\M)$ is regarded as an object of $\mathrm{CAT}/\M$.}
equipped with functorial homomorphisms (over $\M$)
\[
\mathrm{id}_{\uFl(\M)} \;\longrightarrow\; p_\alpha \;.
\]


\endgroup
% END UNIT Der-ChXV-05-11-10/Der-ChXV-05-11-10_p001-005.tex
% BEGIN UNIT Der-ChXV-05-11-10/Der-ChXV-05-11-10_p006-010.tex
\begingroup\makeatletter%
\newcommand{\printedpage}[1]{\marginpar{\footnotesize p.~#1}}%
\newcommand{\msp}[1]{\marginpar{\footnotesize [#1]}}%
\newcommand{\cM}{\mathcal{M}}%
\newcommand{\Flu}{\underline{\mathrm{Fl}}}%
\newcommand{\Fl}{\mathrm{Fl}}%
\newcommand{\fU}{\mathfrak{U}}%
\newcommand{\fV}{\mathfrak{V}}%
\makeatother\let\printedpage\derpp\let\msp\dermsp%
\let\mspin\DERmspin%


\printedpage{4}
% head: A. Grothendieck : Dérivateurs, Ch. XV, Édition des Universités de Montpellier II et Paris VII
% foot: Version 5 novembre 2010

Thus, if we consider
\[
\left\{
\begin{array}{rl}
\Flu(\cM) & \text{as an object of CAT}/\cM \text{ by the functor } f \longmapsto \text{target}\, f\\[1ex]
\text{where} & \text{CAT is the category of categories \textit{elements of} } \fV\text{\textit{,} a universe}\\
 & \text{such that } \fU \in \fV\text{, whence Cat } \overset{\text{(def)}}{=} \mathrm{Cat}^{\fU},\ \cM,\ \Flu(\cM) \in \text{CAT} ,
\end{array}
\right.
\]
we have an endofunctor $p_1$ of $\Flu(\cM)$ over $\cM$, and a functorial homomorphism
\[
i_1 \ : \ \mathrm{id}_{\Flu(\cM)/\cM} \longrightarrow p_1 \ , \qquad \textit{i.e.} \quad
\begin{tikzcd}[column sep=small]
X \arrow[rr, "f"] \arrow[dr, "i_1(f)"'] & & Y \\
 & \underbrace{X_1(f)}_{=\,\text{source}\, p_1(f)} \arrow[ur, "p_1(f)"'] &
\end{tikzcd} \ ,
\tag*{(1.13)}
\]
and the description of the ordinal sequence of endofunctors $p_\alpha$ of $\Flu\,\cM/\cM$ follows from this datum alone, taking into account that small filtering limits exist in the fibers of $\Flu\,\cM \xrightarrow{\text{target}} \cM$.

\msp{page 5}

\[
\begin{tikzcd}[column sep=normal]
X \arrow[r, "i_1(f)"] \arrow[dr, "f"'] & X_1(f) \arrow[r, "i_1(p_1(f))"] \arrow[d, "p_1(f)"] & X_2(f) \arrow[r] \arrow[dl, "p_2(f)"] & \cdots \arrow[r] & X_\alpha(f) \arrow[r, "i_1(p_\alpha(f))"] \arrow[dlll, "p_\alpha(f)"] & X_{\alpha+1}(f) \arrow[r] \arrow[dllll, "p_{\alpha+1}(f) = p_1(p_\alpha(f))"] & \cdots \\
 & Y & & & & &
\end{tikzcd}
\tag*{(1.14)}
\]

So, to finish defining the ordinal sequence of the $X_\alpha(f)$ (equipped with the $i_\alpha(f)$, $p_\alpha(f)$), or of the $p_\alpha(f)$ (equipped with the $i_\alpha(f)$), it remains to describe $X_1(f)$, $i_1(f)$, $p_1(f)$.

The essential property of this triple that we shall need is the following:

\begin{itemize}[leftmargin=5em, labelwidth=4em, labelsep=1em, align=left]
\item[(1.15)] \textbf{Fact}$(C_0)$ (factorization property with respect to $C_0$).

For every commutative diagram
\[
D : \qquad
\begin{tikzcd}
A \arrow[r] \arrow[d, "i"'] & X \arrow[d, "f"] \\
B \arrow[r] & Y
\end{tikzcd}
\qquad \text{with } i \in C_0 \ ,
\]
whence a solid-arrow diagram
\[
\begin{tikzcd}
A \arrow[r] \arrow[d, "i"'] & X \arrow[d, "i_D"] \arrow[dr, dotted, "i_1(f)"] \arrow[dd, dotted, bend right=60, "f"'] & \\
B \arrow[r] & X_D \overset{\mathrm{def}}{=} X \amalg_A B \arrow[r, dotted, "u_D"] \arrow[d, "p_D"] & X_1(f) \arrow[dl, dotted, "p_1(f)"] \\
 & Y \ , &
\end{tikzcd}
\]
\DERmark{ENV0849}

\printedpage{5}
% head: A. Grothendieck : Dérivateurs, Ch. XV, Édition des Universités de Montpellier II et Paris VII
% foot: Version 5 novembre 2010
there exists a homomorphism from the factorization $f = p_D i_D$ of $f$ to the factorization $f = p_1(f)i_1(f)$.
\end{itemize}

To formulate this condition, I have used the existence of the amalgamated sums $X_D$, \textit{i.e.} the condition on $C_0$:

\msp{page 6}

\begin{itemize}[leftmargin=5em, labelwidth=4em, labelsep=1em, align=left]
\item[(1.16)] The arrows $i \in C_0$ are cosquarable, \textit{i.e.} for a diagram
\[
\begin{tikzcd}
A \arrow[r, "u"] \arrow[d, "i"'] & X \\
B &
\end{tikzcd}
\]
with $i \in C_0$, the amalgamated sum exists.
\end{itemize}

But we can also formulate 1.15 as follows, independently of the hypothesis (1.16):

\begin{itemize}[leftmargin=5em, labelwidth=4em, labelsep=1em, align=left]
\item[(1.17)] For every commutative square
\[
\begin{tikzcd}
A \arrow[r, "u"] \arrow[d, "i"'] & X \arrow[d, "f"] \\
B \arrow[r, "v"] & Y
\end{tikzcd}
\qquad ,
\]
with $i \in C_0$, there exists a solid-arrow commutative square
\[
\begin{tikzcd}
A \arrow[r, "u"] \arrow[d, "i"'] & X \arrow[d, "i_1(f)"] \\
B \arrow[r, "v_1"] \arrow[dr, dotted, "v"] & X_1(f) \arrow[d, dotted, "p_1(f)"] \\
 & Y
\end{tikzcd}
\]
such that $v = p_1(f)v_1$.
\end{itemize}

This has as a consequence the

\textbf{Corollary 1} (of 1.17). \textit{Let} $f \in \Fl(\cM)$\textit{, and} $\alpha$ \textit{be an ordinal} $(\in \fU)$. \textit{For every commutative square}
\[
\begin{tikzcd}
A \arrow[r, "u_\alpha"] \arrow[d, "i"'] & X_\alpha(f) \arrow[d, "p_\alpha(f)"] \\
B \arrow[r, "v"] & Y
\end{tikzcd}
\]
\printedpage{6}
% head: A. Grothendieck : Dérivateurs, Ch. XV, Édition des Universités de Montpellier II et Paris VII
% foot: Version 5 novembre 2010
\textit{there exists a commutative square}
\[
\begin{tikzcd}
A \arrow[r, "u_\alpha"] \arrow[d, "i"'] & X_\alpha(f) \arrow[d, "\text{trans.}"] \arrow[dd, dotted, bend left=60, "p_\alpha(f)"] \\
B \arrow[r, "v_1"] \arrow[dr, dotted, "v"] & X_{\alpha+1}(f) \arrow[d, dotted, "p_{\alpha+1}(f)"] \\
 & Y
\end{tikzcd}
\]
\textit{such that}
\[
v \ = \ p_{\alpha+1}(f)\ v_1 \ .
\]

\msp{page 7}

\textbf{Corollary 2.} \textit{Let $\pi$ be a cardinal} $(\in \fU)$ \textit{such that for every arrow $i : A \longrightarrow B$ in $C_0$, the source $A$ is $\pi$-accessible. Let $\alpha$ be a nonzero limit ordinal of cofinality strictly greater than $\pi$} (i.e. \textit{such that the ordered set $I_\alpha$ of ordinals $\alpha' < \alpha$ is large relative to $\pi$}) -- \textit{e.g. $\alpha$ the first ordinal whose cardinality is $> \pi$. Then for every arrow $f \in \Fl(\cM)$, we have $p_\alpha(f) \in (C_0)_*$.}
\DERmark{ENV0850}


\textsc{Proof}. This means that for every commutative square
\[
\begin{tikzcd}
A \arrow[r, "u"] \arrow[d, "i"'] & X_\alpha(f) \arrow[d, "p_\alpha(f)"] \\
B \arrow[r, "v"] \arrow[ur, dotted, "w"] & Y \ ,
\end{tikzcd}
\tag*{$(*)$}
\]
with $i \in C_0$, the dotted arrow $w$ (with $wi = u$, $p_\alpha(f)w = v$) exists. Now
\[
X_\alpha(f) = \varinjlim_{\alpha' \in I_\alpha} X_{\alpha'}(f) \ ,
\]
so $A$ being $\pi$-accessible and $I_\alpha$ large relative to $\pi$, $u$ comes from an arrow $u_{\alpha'} : A \longrightarrow X_{\alpha'}(f)$, and (if necessary replacing $\alpha'$ by $\alpha'' > \alpha'$ and $u_{\alpha'}$ by the composite $u_{\alpha''} : A \longrightarrow X_{\alpha'}(f) \longrightarrow X_{\alpha''}(f)$) we may suppose $p_{\alpha'}(f)u_{\alpha'} = vi$, \textit{i.e.} that $(*)$ comes from a commutative square
\[
\begin{tikzcd}
A \arrow[r, "u_{\alpha'}"] \arrow[d, "i"'] \arrow[rr, dotted, bend left=30, "u"] & X_{\alpha'}(f) \arrow[r, dotted, "\text{trans.}"'] \arrow[d, "p_{\alpha'}(f)"'] & X_\alpha(f) \arrow[dl, dotted, "p_\alpha(f)"] \\
B \arrow[r, "v"] & Y \ , &
\end{tikzcd}
\]

\printedpage{7}
% head: A. Grothendieck : Dérivateurs, Ch. XV, Édition des Universités de Montpellier II et Paris VII
% foot: Version 5 novembre 2010
and applying corollary 1, we find
\[
\begin{tikzcd}
A \arrow[r, "u_{\alpha'}"] \arrow[d, "i"'] & X_{\alpha'}(f) \arrow[d, "\text{trans.}"] & \\
B \arrow[r, "v_1"] \arrow[dr, dotted, "v"] & X_{\alpha'+1}(f) \arrow[r, dotted, "\text{trans.}"] \arrow[d, "p_{\alpha^{\prime}+1}(f)"] & X_\alpha(f) \arrow[dl, dotted, "p_\alpha(f)"] \\ % [sic] p_{\alpha+1}(f) -> p_{\alpha'+1}(f), the arrow starts at X_{\alpha'+1}(f); ?? uncertain: subscript of the label, tightly set -- \alpha+1 or \alpha'+1
 & Y &
\end{tikzcd}
\]
\DERmark{ENV0851}

with $v = p_{\alpha^{\prime}+1}(f)v_1$, and we shall take % [sic] p_{\alpha+1}(f)v_1 -> p_{\alpha'+1}(f)v_1, corollaire 1 is applied to \alpha' (arrow from X_{\alpha'+1}(f))
\DERmark{ENV0852}

\[
w \ = \ \mathrm{trans}_{\alpha,\alpha'+1} \circ v_1 \ .
\]

\msp{page 8}

But we want a factorization of $f$, an arbitrary arrow in $\cM$, as
\[
f \ = \ p\ i \ , \qquad \text{with} \ \ p \in (C_0)_* \ \ \textit{and} \ \ i \in \widetilde{C_0} \ ,
\]
so we must still state a property, ensuring that the $i_\alpha(f)$ are in $\widetilde{C_0}$. We shall even have something more precise. Let us introduce

\begin{itemize}[leftmargin=10em, labelwidth=9em, labelsep=1em, align=left]
\item[(1.18)\hfill $C'_0 \ = $] smallest subset of $\Fl(\cM)$ containing $C_0$, stable under cobase change (whenever it exists).
\item[\hfill $C''_0 \ = $] set of arrows which are either in $C'_0$, or amalgamated sums of a possibly empty family of such arrows $A \xrightarrow{i_\alpha} B_\alpha$ with the same source $A$, $i_\alpha \in C'_0$; the arrow $A \longrightarrow B = \displaystyle\coprod_{\alpha}{}_{(A)} B_\alpha$ is in $C''_0$.
\end{itemize}
\DERmark{ENV0853}


(Condition 1.16 implies that the elements of $C'_0$ are cosquarable, and this suffices, together with the condition of stability of $\cM$ under filtering $\varinjlim$, to ensure that the infinite amalgamated sums envisaged exist in $\cM$.)

\textbf{N.B.} $C''_0$ is \textit{equal} to the smallest set of arrows containing $C_0$ and stable under cobase change and under small amalgamated sums. \texttt{[If small sums exist in $\cM$, then $C''_0$ is also the set of arrows obtained by cobase change from sums of arrows of $C_0$.]}

\setcounter{footnote}{1}% the edition numbers these footnotes 2 and 3
\begin{itemize}[leftmargin=10em, labelwidth=9em, labelsep=1em, align=left]
\item[(1.19)\hfill $C'''_0 \ = $] smallest subset of $\Fl(\cM)$ containing $C''_0$, and moreover stable under strict ordinal direct limits, \textit{i.e.} for an ordinal direct system $(A_\alpha)_{\alpha<\alpha_0}$, such that for $\alpha$ a limit ordinal we have $A_\alpha \xleftarrow{\ \sim\ } \varinjlim\limits_{\alpha' \in I_\alpha} A_{\alpha'}$, if the transition morphisms $i_{\beta,\alpha} : A_\alpha \longrightarrow A_\beta$ are in $C'''_0$, so too the canonical morphisms
\[
A_\alpha \xrightarrow{\ i_\alpha\ } \varinjlim_{\beta \in I_{\alpha_0}} A_\beta \ . \quad \footnotemark
\]
\end{itemize}
\footnotetext{$C_0 \subset C'_0 \subset C''_0 \subset C'''_0$}% footnote 2: mark printed in the display, text printed at the foot of p. 8

\printedpage{8}
% head: A. Grothendieck : Dérivateurs, Ch. XV, Édition des Universités de Montpellier II et Paris VII
% foot: Version 5 novembre 2010
\msp{page 9}

Proposition 2, p. 4 of XIII ensures that we have
\[
C'''_0 \ \subset \ \widetilde{C_0} \ .
\]
This being so, here is the second essential property of the factorization $f = p_1(f)i_1(f)$ (the first being 1.17):

\begin{itemize}[leftmargin=5em, labelwidth=4em, labelsep=1em, align=left]
\item[(1.20)] \qquad\qquad For every arrow $f \in \Fl(\cM)$, we have $i_1(f) \in C''_0$.
\end{itemize}

It follows at once, by ordinal induction:

\begin{itemize}[leftmargin=5em, labelwidth=4em, labelsep=1em, align=left]
\item[(1.21)] \qquad For every arrow $f \in \Fl(\cM)$ and every ordinal $\alpha \in \fU$, we have $i_\alpha(f) \in C'''_0$, a fortiori $i_\alpha(f) \in \widetilde{C_0}$.
\end{itemize}

So if $\alpha$ is large relative to $\pi$ (\textit{cf.} cor. 2, page 7), the factorization $f = p_\alpha(f)i_\alpha(f)$ satisfies
\[
\left\{
\begin{array}{l}
f = p_\alpha(f)i_\alpha(f)\text{, with } i_\alpha(f) \in C'''_0 \subset \widetilde{C_0} = \Phi \text{ \texttt{[not to be}}\\
\text{\texttt{confused with }}\Phi_\alpha\text{\texttt{]}, } p_\alpha(f) \in C_{0*} = \Psi \text{, when } \alpha \text{ is}\\
\text{large relative to } \pi\text{, the sources of the } i \in C_0 \ \pi\text{-accessible.}
\end{array}
\right.
\tag*{(1.22)}
\]

It remains to construct the functor $f \longmapsto X_1(f)$, with $i_1(f)$, $p_1(f)$, satisfying conditions (1.17) and (1.20), under the hypothesis (1.16) that the $i \in C_0$ are cosquarable. We then consider, for $f \in \Fl(\cM)$ fixed, the set $\Delta(f)$ (or $\Delta(C_0, f)$) of \textit{all} the diagrams $D$ of the form (1.15), with $i \in C_0$, and we put
\[
X_1(f) \ = \ \coprod_{D \in \Delta(f)}{}_{(X)} X_D
\tag*{(1.23)}
\]
(amalgamated sum under $X$, \textit{i.e.} sum indexed
\msp{page 10}
by $\Delta(f)$ in $X\backslash\cM$), the morphism $X \xrightarrow{i_1(f)} X_1(f)$ being the costructural morphism, and the morphism $X_1(f) \xrightarrow{p_1(f)} Y$ being the one with the $p_D$ as components. In short, we proceed ``as economically as possible'' (and as naturally as possible) to ensure the validity of (1.17), in the equivalent form (1.15). On the other hand, by definition of $C'_0$, $C''_0$, the $i_D$ are in $C'_0$, and $i_1(f)$ is in $C''_0$, which is (1.20)\footnote{One should all the same make explicit the functorial dependence of $X_1(f)$ on $f$\,!!!}.

To sum up:

\textbf{Theorem 1.} \textit{Let $\cM$ be a $\fU$-category stable under filtering $\fU$-direct limits, $C_0 \subset \Fl(\cM)$ a small set of arrows of $\cM$. We assume the arrows $i \in C_0$ to be cosquarable. This ensures the definition of an endofunctor $f \longmapsto p_1(f)$ of $\Flu(\cM)$, and of a functorial homomorphism $i_1 : \mathrm{id}_{\Flu(\cM)} \longrightarrow p_1$,} i.e. \textit{$i_1(f) : f \longrightarrow p_1(f)$, inducing the identity on the target objects of $f$ and of $p_1(f)$,}
\[
\begin{tikzcd}[column sep=small]
X \arrow[rr, "i_1(f)"] \arrow[dr, "f"'] & & X_1(f) \arrow[dl, "p_1(f)"] \\
 & Y \ , &
\end{tikzcd}
\]


\endgroup
% END UNIT Der-ChXV-05-11-10/Der-ChXV-05-11-10_p006-010.tex
% BEGIN UNIT Der-ChXV-05-11-10/Der-ChXV-05-11-10_p011-015.tex
\begingroup\makeatletter%
\newcommand{\printedpage}[1]{\marginpar{\footnotesize p.~#1}}%
\newcommand{\msp}[1]{\marginpar{\footnotesize [#1]}}%
\newcommand{\uFl}{\underline{\mathrm{Fl}}}%
\newcommand{\Fl}{\mathrm{Fl}}%
\newcommand{\Filt}{\mathrm{Filt}}%
\newcommand{\cM}{\mathcal{M}}%
\newcommand{\card}{\operatorname{card}}%
\newcommand{\circled}[1]{\tikz[baseline=(c.base)]\node[draw,circle,inner sep=0.5pt](c){#1};}%
\newcommand{\pfn}[2]{\textsuperscript{(}\footnote[#1]{#2}\textsuperscript{)}}%
\makeatother\let\printedpage\derpp\let\msp\dermsp%
\let\mspin\DERmspin%


\printedpage{9}
% head: A. Grothendieck : Dérivateurs, Ch. XV, Édition des Universités de Montpellier II et Paris VII
% foot: Version 5 novembre 2010

{\itshape whence by iteration an ordinal direct system of endofunctors $f\longmapsto p_\alpha(f)$ of $\uFl(\cM)$, and of homomorphisms $i_\alpha : \mathrm{id}_{\uFl(\cM)} \longrightarrow p_\alpha$, \textup{i.e.}}
\[
\begin{tikzcd}[column sep=small]
X \arrow[rr, "i_\alpha(f)"] \arrow[dr, "f"'] & & X_\alpha(f) \arrow[dl, "p_\alpha(f)"] \\
& Y\,, &
\end{tikzcd}
\]
{\itshape forming an ordinal direct system of factorizations of $f$ functorial in $f$.}

\msp{page 11}

{\itshape This being so, we have:}
\begin{enumerate}[label={}, leftmargin=*]
\item[a)] {\itshape For every ordinal $\alpha$, we have $i_\alpha(f)\in\Phi$.}
\item[b)] {\itshape Suppose that the sources of the $i\in C_0$ are accessible and then let $\pi$ be a cardinal in $\mathfrak{U}$ such that they are $\pi$-accessible. Then for every nonzero limit ordinal $\alpha$ of cofinality strictly greater than $\pi$, we have $p_\alpha(f)\in C_{0*}$, hence}
\[
f \;=\; p_\alpha(f)\, i_\alpha(f) \qquad\quad \textit{with} \qquad\quad i_\alpha(f)\in\Phi\,, \qquad\quad p_\alpha(f)\in\Psi\,.
\]
\DERmark{ENV0854}

\end{enumerate}

We have a slightly more precise result, as follows:

\textbf{Corollary 1.} \textit{Let $\Phi_0$ be a subset of} $\Fl(\cM)$ \textit{such that}
\[
C_0''' \;\subset\; \Phi_0 \;\subset\; \Phi \qquad\quad (\textit{hence}\ \ C_0\subset\Phi_0\subset\Phi)\,,
\tag*{(1.24)}
\]
\textit{where $C_0'''\subset\Phi$ is defined in} 1.19, \textit{p.} 8. \textit{For example, if $\Phi_0$ is a subset of} $\Fl(\cM)$ \textit{satisfying}
\[
C_0 \;\subset\; \Phi_0 \;\subset\; \Phi\,,
\]
\textit{and stable under cobase change, under small amalgamated sums} (\textit{cf.} (1.18), p.~8), \textit{under strict ordinal $\varinjlim$} (\textit{cf.} (1.19), p.~8), \textit{it satisfies} (1.24). (\textbf{N.B.} \textit{If $\cM$ admits simple amalgamated sums, such $\Phi_0$ exist trivially, it suffices to take $\Phi_0=\Phi$.})

\msp{page 12}

\textit{This being so, for every arrow $f$ in $\cM$, and every ordinal $\alpha$, we have}
\[
i_\alpha(f) \;\in\; \Phi_0\,.
\]

\textbf{Corollary 2.} \textit{Suppose that the sources of the arrows $i\in C_0$ are accessible, and let $\Phi_0$ be as in Corollary} 1 (e.g. $\Phi_0=C_0'''$). \textit{Then the pair $(\Phi_0,\Psi)$ is a Quillen precouple,} i.e.
\begin{enumerate}[label={}, leftmargin=*]
\item[a)] $\Phi_0\longleftrightarrow\Psi$, i.e. $\Psi\subset\Phi_{0*}$ (\textit{or what amounts to the same thing, $\Phi_0\subset\Psi^*$}), \textit{and}
\item[b)] \textit{we have the Quillen factorization property.}
\end{enumerate}
\textit{This therefore implies that $\Phi$ $(=\widetilde{C_0}=\widetilde{\Phi_0})$ consists of the arrows which are direct factors of arrows of $\Phi_0$. Thus if $\Phi_0$ is stable under direct factors, we have $\Phi_0=\Phi$,} i.e. \textit{$\Phi_0$ is left closed.}

\printedpage{10}
% head: A. Grothendieck : Dérivateurs, Ch. XV, Édition des Universités de Montpellier II et Paris VII
% foot: Version 5 novembre 2010

\section*{2\quad Variation of the factorizations $f=p_\alpha i_\alpha$ with the generating subset $C_0$}

I say that the ordinal direct system of functors $p_\alpha=p_\alpha^{C_0} : \uFl(\cM)\longrightarrow\uFl(\cM)$,
\msp{page 13}
equipped with the arrows
\[
i_\alpha^{C_0} \;:\; \mathrm{id}_{\uFl(\cM)} \longrightarrow p_\alpha^{C_0}\,,
\]
depends functorially on $C_0$, when one regards $C_0$ as an element of the ordered set of \textit{small} subsets of $\Fl(\cM)$, consisting of cosquarable arrows (thus of the ordered set
\[
\mathfrak{P}_{\mathfrak{U}\text{-small}}(\Fl_{\text{cosquar.}}(\cM)) \quad ).
\]
If we have an inclusion
\[
C_0 \;\hookrightarrow\; C_1\,,
\tag*{(2.1)}
\]
we deduce from it for every ordinal $\alpha$ a functorial homomorphism
\[
u_\alpha \;=\; u_\alpha^{C_1,C_0} \;:\; X_\alpha^{C_0} \;\longrightarrow\; X_\alpha^{C_1}\,,
\tag*{(2.2)}
\]
making commutative the squares deduced from $f : X\longrightarrow Y$ in $\cM$,
\[
\begin{tikzcd}
X \arrow[r, "i_\alpha^{C_0}(f)"] \arrow[d, "\mathrm{id}"'] & X_\alpha^{C_0}(f) \arrow[r, "p_\alpha^{C_0}(f)"] \arrow[d, "u_\alpha(f)"] & Y \arrow[d, "\mathrm{id}"] \\
X \arrow[r, "i_\alpha^{C_1}(f)"] & X_\alpha^{C_1}(f) \arrow[r, "p_\alpha^{C_1}(f)"] & Y\,,
\end{tikzcd}
\tag*{(2.3)}
\]
and compatible with the transition morphisms for variable $\alpha$,
\[
\begin{tikzcd}
X_\alpha^{C_0} \arrow[r, "{\mathrm{trans}_{\beta,\alpha}^{C_0}}"] \arrow[d, "u_\alpha"'] & X_\beta^{C_0} \arrow[d, "u_\beta"] \\
X_\alpha^{C_1} \arrow[r, "{\mathrm{trans}_{\beta,\alpha}^{C_1}}"] & X_\beta^{C_1}\,.
\end{tikzcd}
\tag*{(2.4)}
\]

\msp{page 14}

To define the $u_\alpha$, satisfying the preceding conditions, it suffices to define them for $\alpha=1$ ``and to iterate'':
\[
\begin{tikzcd}[column sep=small]
& X_1^{C_0}(f) \arrow[dd, "u_1(f)"] \arrow[dr, "p_1^{C_0}(f)"] & \\
X \arrow[ur, "i_1^{C_0}(f)"] \arrow[dr, "i_1^{C_1}(f)"'] & & Y \\
& X_1^{C_1}(f) \arrow[ur, "p_1^{C_1}(f)"'] & {}.
\end{tikzcd}
\tag*{(2.5)}
\]
Now $X_1^{C_0}(f)$ is defined as a direct sum in $X\backslash\cM/Y$, relative to an index set $\Delta(C_0,f)$ \textit{contained} in $\Delta(C_1,f)$ -- and the family of objects whose direct sum is taken being induced from a (non-small) family whose index set would be $\Delta(\Fl_{\text{cosquar.}}(\cM),f)$, namely consisting of \textit{all} the square diagrams
\printedpage{11}
% head: A. Grothendieck : Dérivateurs, Ch. XV, Édition des Universités de Montpellier II et Paris VII
% foot: Version 5 novembre 2010
\[
D : \qquad
\begin{tikzcd}
A \arrow[r, "u"] \arrow[d, "i"'] & X \arrow[d, "f"] \\
B \arrow[r, "v"] & Y
\end{tikzcd}
\qquad\qquad \text{with $i$ cosquarable and nothing more ,}
\]
the definition of the $X_D$, $i_D$, $p_D$ associated to $D$ depending only on $D$, and not on the set $C_0$ or $C_1$ of cosquarable arrows which contains $i$. This said, the arrow $u_1(f)$
\msp{page 15}
in (2.5) is none other than the canonical inclusion of a partial amalgamated sum, indexed by $\Delta(C_0,f)$, into the amalgamated sum corresponding to the larger index set. The functoriality of this arrow $u_1(f)$, for variable $f$, ``is immediate'', it follows from the commutativity of the square, associated to a morphism $v=(v_X,v_Y)$ of arrows
\[
\begin{tikzcd}
X \arrow[r, "f"] \arrow[d, "v_X"'] & Y \arrow[d, "v_Y"] \\
X' \arrow[r, "f'"] & Y'\,,
\end{tikzcd}
\tag*{(2.6)}
\]
namely the diagram of maps between index sets
\[
\begin{tikzcd}
\Delta(C_0,f) \arrow[r, hook, "u_1(f)"] \arrow[d, "{\begin{array}{r}\text{functorial dependence}\\ \text{on $f$ of $\Delta(C_0,f)$}\end{array}}"'] & \Delta(C_1,f) \arrow[d, "{\begin{array}{l}\text{functorial dependence}\\ \text{on $f$ of $\Delta(C_1,f)$}\end{array}}"] \\
\Delta(C_0,f') \arrow[r, hook, "u_1(f')"] & \Delta(C_1,f')
\end{tikzcd}
\tag*{(2.7)}
\]
(commutativity which amounts to saying that $\Delta(C_0,f)$ may be regarded as a functor in the pair $(C_0,f)$ as an object of the product category $\underline{\mathfrak{P}_{\text{sm.}}(\Fl_{\text{cosquar.}}(\cM))}\times\uFl(\cM)$ \ldots). The commutative square (2.7) is the ``indices'' part of a commutative square between families of objects $\mathcal{X}(C_0,f)$, $\mathcal{X}(C_0,f')$, $\mathcal{X}(C_1,f)$, $\mathcal{X}(C_1,f')$,
\DERmark{ENV0855}\DERmark{ENV0856}\DERmark{ENV0857}\DERmark{ENV0858}

\msp{page 16}
indexed by these sets, namely the $X_D^{C_0}$ and the $X_D^{C_1}$ relative to the diagrams $D$ involving, either the arrow $f$, or the arrow $f'$ \ldots

Now the canonical construction of the ordinal direct system $p_\alpha$, with the $i_\alpha : \mathrm{id}_{\uFl(\cM)}\longrightarrow p_\alpha$, is functorial with respect to the initial datum $p_1$, $i_1$ -- I do not spell out the detail of the functoriality, although I should \ldots

\section*{3\quad Upper bounds for cardinals}

I now suppose that $\cM$ is equipped with a \textit{cardinal filtration}
\[
(\Filt^\pi)_{\pi\geq\pi_0} \qquad\quad (\pi_0,\ \pi \text{ cardinals in } \mathfrak{U})
\tag*{(3.1)}
\]
(SGA 4 I 9.12, p.~149), the $\Filt^\pi$ stable under simple amalgamated sums (whenever they exist), hence also under infinite amalgamated sums indexed by sets
\printedpage{12}
% head: A. Grothendieck : Dérivateurs, Ch. XV, Édition des Universités de Montpellier II et Paris VII
% foot: Version 5 novembre 2010
of indices of cardinality $\leq\pi$. We further suppose that every object of $\Filt^{\pi_0}$ is accessible, and we denote by \circled{$\pi_1$} the smallest cardinal such that $\pi_1\geq\pi_0$ and that the objects of $\Filt^{\pi_0}$ are $\pi_1$-accessible. Denoting, for every cardinal, by $\cM_\pi$ the subset of $\mathrm{Ob}\,\cM$ consisting of the $\pi$-accessible objects, we will then have for every $\pi\geq\pi_1$ (\textit{loc. cit.} 9.16)
\[
\cM_\pi \;\subset\; \Filt^\pi \;\subset\; \cM_{(\pi^{\pi_0})} \qquad\quad (\text{if } \pi\geq\pi_1)\,,
\tag*{(3.2)}
\]

\msp{page 17}

thus if $\pi^{\pi_0}=\pi$, e.g. if $\pi=2^c$ with $c\geq\pi_0$, we will have
\[
\Filt^\pi \;=\; \cM_\pi\,.
\tag*{(3.2 bis)}
\]
(If we admit the generalized continuum hypothesis, we have $\pi^{\pi_0}=\pi$ for every $\pi>\pi_0$, hence $\Filt^\pi=\cM_\pi$ if $\pi\geq\pi_1$ and $\pi>\pi_0$.)

I put, for every cardinal $\pi\geq\pi_0$,
\[
\Fl^\pi(\cM) \;=\; \{\ f\in\Fl(\cM) \mid \text{source and target of $f$ are in } \Filt^\pi\ \}
\tag*{(3.3)}
\]
-- which probably gives a cardinal filtration on $\uFl(\cM)$ (but that is not our concern \ldots). My aim now, it is to bring out conditions on $C_0\subset\Fl_{\text{cosquar.}}(\cM)$ (small subset of $\Fl(\cM)$) and on the cardinal $\pi$, and the ordinal $\alpha$, so that we have
\[
\begin{aligned}
f\in\Fl^\pi(\cM) \;\;\overset{?}{\Longrightarrow}\;\; & X_\alpha(f)\in\Filt^\pi \\
& (\textit{i.e.}\ p_\alpha(f)\in\Fl^\pi(\cM),\ \text{or again } i_\alpha(f)\in\Fl^\pi(\cM))\,.
\end{aligned}
\tag*{(3.4)}
\]
More generally, I propose to give an estimate of the ``cardinality'' of $X_\alpha(f)$, for fixed $C_0$, as a function of $C_0$, of $f$, of $\alpha$, more exactly, as a function of the three cardinals $\pi(C_0)$, $\pi(f)$, $\pi(\alpha)$ defined as
\[
\left\{
\begin{array}{rcl}
\pi(C_0) & = & \text{smallest cardinal such that } C_0\subset\Fl^\pi(\cM) \\[1ex]
\pi(f) & = & \pi(\{f\}) \;=\; \text{smallest cardinal such that } f\in\Fl^\pi(\cM) \\[1ex]
\pi(\alpha) & = & \card\alpha\,.
\end{array}
\right.
\tag*{(3.5)}
\]
Thus I am looking for a cardinal
\msp{page 18}
function of cardinal variables $\varphi$, such that we have
\[
X_\alpha(f) \;\in\; \Filt^\pi \qquad\quad \text{if} \qquad\quad \pi \;=\; \varphi(\underbrace{\pi(C_0)}_{r},\underbrace{\pi(f)}_{s},\underbrace{\pi(\alpha)}_{t})\,.
\tag*{(3.6)}
\]
The crucial estimate will surely concern the case $\alpha=1$. Let us note that $X_1^{C_0}(f)=X_1(f)$ is defined as an amalgamated sum (in general infinite) of objects $X_D$, for a set $\Delta(C_0,f)=\Delta(f)$ of diagrams, namely the commutative diagrams
\[
\begin{tikzcd}
A \arrow[r] \arrow[d, "i"'] & X \arrow[d, "f"] \\
B \arrow[r] & Y
\end{tikzcd}
\]
\printedpage{13}
% head: A. Grothendieck : Dérivateurs, Ch. XV, Édition des Universités de Montpellier II et Paris VII
% foot: Version 5 novembre 2010
with $i\in C_0$. For a fixed $i$, and a fixed $f$, the cardinality of the set of diagrams of this type is bounded above by that of $\mathrm{Hom}(A,X)\times\mathrm{Hom}(B,Y)$. We have $A,B\in\Filt^r$, $X,Y\in\Filt^s$, so we need an upper bound for
\[
\card\mathrm{Hom}_{\cM}(A,X)\,, \qquad\quad A\;\in\;\Filt^r\,, \qquad\quad X\;\in\;\Filt^s\,.
\]

\msp{page 19}

Here are the essential upper bounds in this direction:

\textbf{Proposition 3.1.} \textit{Let $\pi_0$ be the initial cardinal of the cardinal filtration} $(\Filt^\pi)_{\pi\geq\pi_0}$ \textit{under consideration, $\pi_1$ the smallest cardinal $\geq\pi_0$ such that the objects of} $\Filt^{\pi_0}$ \textit{are $\pi_1$-accessible, $C_1=$} card $\underbrace{\widetilde{\uFl^{\pi_1}(\cM)}}_{=\uFl(\cM^{\pi_1})}$\pfn{4}{\textit{cf.} below the notation $\widetilde{\mathcal{C}}$.}. \textit{For every cardinal $\pi\geq\pi_0$, I denote by $\cM^\pi$ the strictly full subcategory of $\cM$ whose set of objects is} $\Filt^\pi$. \textit{Finally, if $\mathcal{C}$ is a category, I denote by $\widetilde{\mathcal{C}}$ the set of isomorphism classes of objects of $\mathcal{C}$.}\pfn{5}{Let us point out that we always have $\card\widetilde{\uFl(\mathcal{C})}\geq\card\widetilde{\mathcal{C}}$ since $\widetilde{\mathcal{C}}\hookrightarrow\widetilde{\uFl(\mathcal{C})}$ via $\mathrm{cl}(X)\longmapsto\mathrm{cl}(\mathrm{id}_X)$.}
% ?? uncertain: straight underlines under "le plus petit" and under "pleine" -- underlined / not underlined (a wavy tilde and an underbrace stand close by)

$1^\circ$) \textit{If $r\geq\sup(\pi_1,C_1)$, we have}
\[
\card\,(\underbrace{\widetilde{\uFl^r(\cM)}}_{=\uFl(\cM^r)}) \;\leq\; 2^{(r^{\pi_0})} \quad\ (=2^r,\ \textit{if } r=2^c,\ \textit{with } c\geq\pi_0)\,.
\tag*{(3.7)}
\]

$2^\circ$) \textit{If $r,s\geq\sup(\pi_1,C_1)$, and if $A\in\Filt^r$, $X\in\Filt^s$, we have}
% (3.8) and (3.9) share one brace on the right
\[
\begin{array}{l}
\text{(3.8)}\\[3ex]
\text{(3.9)}
\end{array}
\qquad\qquad
\begin{array}{r@{\;\;\leq\;\;}l}
\card\mathrm{Hom}_{\cM}(A,X) & s^{(r^{\pi_0})}\\[3ex]
\card\widetilde{\uFl(\cM^r/X)} & s^{(r^{\pi_0})}
\end{array}
\qquad\quad
\left\{
\begin{array}{lll}
= & s^r & \textit{if } r=2^c,\ \textit{with } c\geq\pi_0\,,\\[3ex]
= & s & \textit{if } s=2^{c'},\ \textit{with } c'\geq r^{\pi_0}\,.
\end{array}
\right.
\]

\textbf{N.B.} If we admit the generalized continuum hypothesis, which implies that $\pi^{\pi'}=\pi$ if $\pi>\pi'$ ($\pi$ an infinite cardinal, $\pi'\geq1$), then taking $r$ such that $r>\pi_0$, the preceding formulas give (with moreover $r\geq\sup(\pi_1,C_1)$):
\[
\card\widetilde{\uFl^r(\cM)} \;\;\leq\;\; 2^r\,,
\tag*{(3.7 bis)}
\]
\[
\card\mathrm{Hom}_{\cM}(A,X) \;\leq\; s^r \qquad\quad (\text{if } A\in\Filt^r,\ X\in\Filt^s)\,,
\tag*{(3.8 bis)}
\]
\[
\card\widetilde{\uFl(\cM^r/X)} \;\leq\; s^r \qquad\quad (\text{if } X\in\Filt^s)\,.
\tag*{(3.9 bis)}
\]
Without the continuum hypothesis, these formulas are also valid for $r$ of the form $2^c$, with $c\geq\pi_0$. It is these last formulas which look the most natural. For (3.8 bis) and % [sic] qui on l'air -> qui ont l'air, verb form
(3.9 bis), I point out that the right-hand side $s^r$ is the cardinality of the set $\mathrm{Hom}_{\mathrm{Ens}}(A,X)$, where $A$, $X$ are two sets such that $\card A=r$, $\card X=s$. We therefore cannot hope to improve (3.8 bis)\,! For the


\endgroup
% END UNIT Der-ChXV-05-11-10/Der-ChXV-05-11-10_p011-015.tex
% BEGIN UNIT Der-ChXV-05-11-10/Der-ChXV-05-11-10_p016-020.tex
\begingroup\makeatletter%
\newcommand{\printedpage}[1]{\marginpar{\footnotesize p.~#1}}%
\newcommand{\msp}[1]{\marginpar{\footnotesize [#1]}}%
\newcommand{\M}{\mathcal{M}}%
\newcommand{\Filt}{\mathrm{Filt}}%
\newcommand{\Fl}{\mathrm{Fl}}%
\newcommand{\uFl}{\underline{\mathrm{Fl}}}%
\newcommand{\uFilt}{\underline{\mathrm{Filt}}}%
\newcommand{\card}{\operatorname{card}}%
\newcommand{\Hom}{\operatorname{Hom}}%
\newcommand{\Ob}{\operatorname{Ob}}%
\newcommand{\wavy}[1]{\widetilde{#1}}%
\makeatother\let\printedpage\derpp\let\msp\dermsp%
\let\mspin\DERmspin%


\printedpage{14}
% head: A. Grothendieck : Dérivateurs, Ch. XV, Édition des Universités de Montpellier II et Paris VII
% foot: Version 5 novembre 2010
\msp{page 20}
first formula (3.7 bis), let us note that we have equality when $\M$ is the category of pairs $(A,B)$, $A$ a set and $B$ a subset of $A$, with the obvious arrows, and we define $\M^\pi$ (for $\pi\geq\pi_0$) as the strictly full subcategory consisting of the $(A,B)$ such that $\card A\leq\pi$.
\DERmark{ENV0859}


\textsc{Proof of} 3.1. Let us first prove formula (3.8),
\[
\card\Hom(A,X)\ \leq\ s^{(r^{\pi_0})}\qquad\quad
\left\{\begin{array}{l}
\text{if } A\in\Filt^r,\ X\in\Filt^s,\\[1.5ex]
r\geq\pi_1,\ s\geq\sup(\pi_1,C_1).
\end{array}\right.
\]

(\textbf{N.B.} I shall not use $r\geq C_1$, postulated in 3.1 $1^\circ$.) I write \texttt{[``I write'' crossed out]}
\[
\begin{aligned}
A\ &=\ \varinjlim_I A_i\ ,\quad \card I\ \leq\ r^{\pi_0}\ ,\quad A_i\ \in\ \Filt^{\pi_0}\ ,\\
X\ &=\ \varinjlim_J X_j\ ,\quad \card J\ \leq\ s^{\pi_1}\ ,\quad X_j\ \in\ \Filt^{\pi_1}\ ,\qquad\qquad J\text{ large relative to }\pi_1\ .
\end{aligned}
\]
We shall have
\[
\Hom(A,X)\ \simeq\ \varprojlim_{i\in I}\ \underbrace{\Hom(A_i,\,\varinjlim_J X_j)}_{\overset{(\text{\footnotemark[6]})}{\simeq}\ \varinjlim\limits_{j\in J}\ \underbrace{\Hom(A_i,X_j)}_{\text{of cardinality }\leq C_1}}\ ,
\]
\footnotetext[6]{since the $A\in\Filt^{\pi_0}$ are $\pi_1$-accessible.}
thus for every $i\in I$
\[
\card\Hom(A_i,X)\ \leq\ \card(J)\times C_1\ \leq\ s^{\pi_1}\times C_1\ =\ s^{\pi_1}\ ,
\]
\DERmark{ENV0860}

since $s\geq C_1$, hence $s^{\pi_1}\geq C_1$, whence
\[
\card(\,\varprojlim_{i\in I}\Hom(A_i,X))\ \leq\ (s^{\pi_1})^{\card I}\ \leq\ \underbrace{(s^{\pi_1})^{(r^{\pi_0})}}_{s^{\pi_1 r^{\pi_0}}}\ . % [sic] j\in I -> i\in I, the limit is over the index i of A_i
\]
\DERmark{ENV0861}

Since $r\geq\pi_1$, we have $r^{\pi_0}\geq\pi_1$, hence $r^{\pi_0}\pi_1=r^{\pi_0}$, whence
\[
\card\Hom(A,X)\ \leq\ s^{(r^{\pi_0})}\ ,
\]
this is indeed formula (3.8). To derive from it an estimate for $\M^r/X$ (with $X\in\Filt^s$), \msp{page 21} I shall first bound above the cardinality of the isomorphism classes of objects of $\M^r/X$, \textit{i.e.} of the arrows up to isomorphism (in $\M^r$)
\[
A\ \longrightarrow\ X\ ,\qquad\quad X\text{ fixed in }\Filt^s,\ A\text{ variable in }\M^r\ .
\]
We shall have
\[
A\ =\ \varinjlim_{i\in I} A_i\qquad\quad\text{with } A_i\in\Filt^{\pi_0},\ I\text{ ordered set of cardinality}\leq r^{\pi_0}\ .
\]

\printedpage{15}
% head: A. Grothendieck : Dérivateurs, Ch. XV, Édition des Universités de Montpellier II et Paris VII
% foot: Version 5 novembre 2010
For fixed $I$, the cardinality of the set of direct systems indexed by $I$, in $\M^{\pi_0}/X$, up to isomorphism, can be bounded above as follows: for every $i\in I$, one must choose $A_i$ in $\M^{\pi_0}$ (but only its class in $\wavy{\M^{\pi_0}}$ matters), so this is bounded above by card $\wavy{\M^{\pi_0}}\ \leq C_0$~\textsuperscript{(}\footnote[7]{Here we have put $C_0=\card\wavy{\uFl^{\pi_0}}$ , hence $C_0\leq C_1$.}\textsuperscript{)}, and for all the $i$, by $\card I\times C_0\leq r^{\pi_0}\times C_0$, and if $r\geq C_0$, hence $r^{\pi_0}\geq C_0$, we find an upper bound on the possible choices of the system of \textit{objects} $A_i$ in $\M^{\pi_0}$, up to isomorphism, given by $r^{\pi_0}$. If in addition we specify the transition morphisms $A_i\longrightarrow A_j$ for $i<j$, this gives for each comparable pair, an element of $I\times I$, a choice bounded above by $C_0$, hence again an upper bound of $\ \underbrace{\card(I\times I)}_{=(\card I)^2=\card I}\cdot C_0\leq r^{\pi_0}\cdot C_0=r^{\pi_0}$. Thus the set of direct systems in question, up to isomorphism, for fixed $I$, has cardinality $\leq r^{\pi_0}\cdot r^{\pi_0}=r^{\pi_0}$. Finally, the \textit{homomorphisms} of a given direct system into $X$ are bounded above by $\card I\times s^{(r^{\pi_0})}$ (by (3.8), already proved), hence by
\DERmark{ENV0862}

\[
r^{\pi_0}\cdot s^{(r^{\pi_0})}\ =\ s^{(r^{\pi_0})}\qquad\quad(\text{since } r^{\pi_0}\leq s^{(r^{\pi_0})})\ .
\]

Altogether, for fixed $I$, the set of direct systems up to isomorphism, indexed by $I$, with values in $\M^r/X$, is \msp{page 22} of cardinality bounded above by $\quad\underbrace{r^{\pi_0}}_{\text{\scriptsize\begin{tabular}{@{}c@{}}number of\\ direct systems\\ in $\M^{\pi_0}$, up to\\ isomorphism\end{tabular}}}\quad\cdot\quad\underbrace{s^{(r^{\pi_0})}}_{\text{\scriptsize\begin{tabular}{@{}c@{}}number of\\ arrows from\\ such a direct system\\ into $X$\end{tabular}}}\quad=s^{(r^{\pi_0})}$.
\DERmark{ENV0863}


But we must now let $I$ vary, up to isomorphism, in the category of ordered sets (filtering \textit{etc.}) of cardinality $\leq\lambda=r^{\pi_0}$. We see at once that this set of isomorphism classes (which can be interpreted as a quotient set of a certain set of subsets of a fixed set $I_0$ of cardinality $\lambda$) has cardinality $\leq 2^\lambda$. We thus finally find
\[
\card\wavy{\M^r/X}\ \ \leq\ s^{(r^{\pi_0})}\cdot 2^{(r^{\pi_0})}\ =\ s^{(r^{\pi_0})}\ ,
\]
on condition that
\[
r\ \geq\ \sup(\pi_1,C_0)\ ,\qquad\quad s\ \geq\ \sup(\pi_1,C_1)\ .
\]
We must prove the similar inequality for card $\wavy{\uFl(\M^r/X)}$ \ (3.9). Now, since for $A$, $B$ in $\M^r$, we have, if $r\geq\sup(\pi_1,C_1)$,
\DERmark{ENV0864}

\[
\card\Hom_{\M}(A,B)\ \leq\ r^{(r^{\pi_0})}
\]
by virtue of (3.8) where we put $s=r$, we find that
\[
\begin{aligned}
\card\wavy{\uFl(\M^r/X)}\ &\leq\ \Bigl(\card\wavy{\M^r/X}\,\Bigr)^2\cdot r^{(r^{\pi_0})}\\
&\leq\ s^{(r^{\pi_0})}\cdot r^{(r^{\pi_0})}\\
&=\ s^{(r^{\pi_0})}\qquad\text{if } s\geq r\ .
\end{aligned}
\]
(We shall look below at the cases $s\leq r$~\textsuperscript{(}\footnote[8]{\textit{cf.} cor. 2, page 23.}\textsuperscript{)}). This is indeed formula (3.9).

\printedpage{16}
% head: A. Grothendieck : Dérivateurs, Ch. XV, Édition des Universités de Montpellier II et Paris VII
% foot: Version 5 novembre 2010
\msp{page 23}
Let us turn to the category $\M^r$ itself. The argument on page 21 shows that the set of isomorphism classes of objects of $\M^r$ indexed by a fixed $I$, of cardinality $\leq r^{\pi_0}$, as $A=\varinjlim\limits_i A_i$ with $A_i\in\Filt^{\pi_0}$, has cardinality bounded above by $r^{\pi_0}$ (if $r\geq C_0$). If we let $I$ vary, we find an upper bound of $r^{\pi_0}\cdot 2^{(r^{\pi_0})}=2^{(r^{\pi_0})}$. We conclude from this as before that the set $\ \wavy{\Fl\,\M^r}$ has cardinality bounded above by
\[
(2^{(r^{\pi_0})})^2\times r^{(r^{\pi_0})}\qquad\quad(\text{if } r\geq\sup(\pi_1,C_1))\ ,
\]
hence by $r^{(r^{\pi_0})}$. I say that we have
\[
r^{(r^{\pi_0})}\ =\ 2^{(r^{\pi_0})}\ ,
\]
indeed, we have $\geq$, and to prove $\leq$, we note that $r<2^r$, hence $r^{(r^{\pi_0})}\leq(2^r)^{(r^{\pi_0})}=2^{r\cdot r^{\pi_0}}=2^{(r^{\pi_0})}$, O.K.! This thus establishes 3.7. This argument in fact proves the

\textbf{Lemma 3.2.} \textit{Let $r$ be a cardinal, $r'$ a cardinal $\geq r$. Then we have}
\[
r^{r'}\ =\ 2^{r'}\qquad\quad(=s^{r'}\ \textit{for all}\ 2\leq s\leq r)\ .
\]

\textbf{Corollary 1.} \textit{Let $r\geq\sup(\pi_1,C_1)$, and $A,X\in\Filt^r$. Then}
\DERmark{ENV0865}

\[
\card\Hom_{\M}(A,X)\ \leq\ 2^{(r^{\pi_0})}\qquad\quad(=2^r\ \textit{if}\ r=2^c,\ \textit{with}\ c\geq\pi_0)\ . \tag*{(3.10)}
\]

We apply (3.8), which gives the upper bound $r^{(r^{\pi_0})}$, and lemma 3.2.

We shall note that we find the same upper bound for the cardinality of $\wavy{\uFl(\M^r)}$ as a whole (3.7)! Likewise, (3.9) gives:

\textbf{Corollary 2.} \textit{Let $r\geq\sup(\pi_1,C_1)$, and $X\in\Ob\M^r$. Then}
\[
\card\wavy{\uFl(\M^r/X)}\ \leq\ 2^{(r^{\pi_0})}\qquad\quad(=2^r\ \textit{if}\ r=2^c,\ \textit{with}\ c\geq\pi_0)\ . \tag*{(3.11)}
\]
so this is formula (3.9) for $s\leq r$, left pending.

\msp{page 24}
Let us return to the cardinal bound for the $X_\alpha(f)$, where $f : X\longrightarrow Y$ is in $\Fl^s(\M)$, \textit{i.e.} $X,Y\in\Filt^s$. We assume moreover $C_0\subset\Fl^r(\M)$,
\[
\begin{aligned}
&C_0\ \subset\ \Fl^r(\M)\ =\ \Fl(\M^r)\ ,\qquad\quad(f:X\longrightarrow Y)\in\Fl^s(\M)\ ,\ \ \textit{i.e.}\ X,Y\in\Filt^s\ ,\\
&\text{with } s\geq r\ .
\end{aligned} \tag*{(3.10)}
\]

\texttt{[Twice (3.10).]} There is a stupid difficulty, which is that $\M^r$ is equivalent to a small category, but is not itself small. Even if $C_0$ contained only a single isomorphism class of arrows of $\M^r$, its cardinality could be (for fixed $r$) arbitrarily large. Thus, in what follows, I must assume the condition of ``cardinal non-plethoricity'' for $C_0$:
\[
\parbox{0.72\linewidth}{$\forall\,i\in C_0$, denoting by $C_0(i)$ the set of $i'\in C_0$ such that $i\simeq i'$ in $\Fl(\M)$, I require that $\card C_0(i)\leq 2^{(r^{\pi_0})}$ (which is also the value of the upper bound for the cardinality of $\wavy{\uFl^r(\M)}=\wavy{\uFl(\M^r)}$\,).} \tag*{(3.11)}
\]

\printedpage{17}
% head: A. Grothendieck : Dérivateurs, Ch. XV, Édition des Universités de Montpellier II et Paris VII
% foot: Version 5 novembre 2010
\texttt{[Twice (3.11).]} This thus implies that we have $\card C_0\leq(2^{(r^{\pi_0})})^2=2^{(r^{\pi_0})}$,
\[
\card C_0\ \leq\ 2^{(r^{\pi_0})}\ , \tag*{(3.12)}
\]
which is equivalent to (3.11) (given $C_0\subset\Fl^r(\M)$, postulated in (3.10)).

It follows by (3.8) that we have
\[
\card\Delta(C_0,f)\ \leq\ \underbrace{\card(C_0)}_{\leq 2^{(r^{\pi_0})}}\cdot s^{(r^{\pi_0})}\cdot s^{(r^{\pi_0})}\ \leq\ s^{(r^{\pi_0})}\ .
\]

On the other hand, the $X_D$ whose amalgamated sum under $X$ we take, to obtain $X_1(f)$, are in $\M^s$, if we make the following assumption. \msp{page 25}
\[
\left|\ \parbox{0.72\linewidth}{The $\M^\pi=\uFilt^\pi$ are stable under simple amalgamated sums, and under filtering $\varinjlim$ of cardinality $\leq\pi$ (not necessarily large relative to $\pi_0$) ,}\right. \tag*{(3.13)}
\]
taking into account that $r\leq s$, hence $\M^r\subset\M^s$, hence $A$, $B$, $X$ in $X_D=B\amalg_A X$ are in $\M^s$. The \textit{finite} partial amalgamated sums of the $X_D$ are therefore also in $\M^s$, by induction on the cardinality of the index set. The infinite amalgamated sum thus appears as a filtering $\varinjlim$ of objects of $\M^s$, the index set being the set $\mathfrak{P}^*_{\mathfrak{f}}\Delta(C_0,f)$ of nonempty finite subsets of $\Delta(C_0,f)$. This set has cardinality bounded above by $s'$, if $s'$ is an \textit{infinite} cardinal such that $\card\Delta(C_0,f)\leq s'$. We can therefore take $s'=s^{(r^{\pi_0})}$. We have therefore proved, taking into account (3.13) applied to $\M^{s'}$, the % ?? uncertain: symbol in X_D = B \amalg_A X -- \amalg or a doubled \sqcup; subscript of \mathfrak{P}^* -- fraktur f or roman f

\textbf{Lemma 3.3.} \textit{We assume} \textup{(3.10), i.e.} $C_0\subset\Fl^r(\M)$, $f\in\Fl^s(\M)$, \textit{with} $s\geq r\geq\sup(\pi_1,C_1)$, \textit{and moreover that $C_0$ satisfies the non-plethoricity condition} \textup{(3.11)} \textup{(}\textit{which we can always assume, at the cost of reducing $C_0$ without changing $\widetilde{C_0}=\Phi$ and $C_{0*}=\Psi$}\textup{)}\textit{, and that the cardinal filtration satisfies the stability condition (harmless!)} \textup{(3.13)}\textit{. Put}
\[
s'\ =\ s^{(r^{\pi_0})}\ . \tag*{(3.14)}
\]
\textit{Then $X_1(C_0,f)$ is in $\M^{s'}$. More generally and more prettily, for fixed $C_0$ and $r$} \textup{(}\textit{with $C_0\subset\Fl^r(\M)$ and satisfying the non-plethoricity condition} \textup{(3.11)} \textit{relative to $r$}\textup{)}\textit{, for every cardinal} \msp{page 26} \textit{of the form}
\[
s'\ =\ s^t\ ,\qquad\quad\textit{with}\ t\geq r^{\pi_0},\ s\geq 2 \tag*{(3.15)}
\]
\textit{(whence $s'\geq 2^{(r^{\pi_0})}>r^{\pi_0}\geq r$), $\uFl^{s'}(\M)$ is stable under the functor $p_1^{C_0}:\uFl(\M)\longrightarrow\uFl(\M)$.}

Indeed, if we apply the first part of the lemma to the pair $(r,s')$, we find that $X_1(f)$ is in $\M^{s''}$, with $s''=s'^{(r^{\pi_0})}$. But
\[
s''\ =\ s'^{(r^{\pi_0})}\ =\ (s^t)^{(r^{\pi_0})}\ =\ s^{tr^{\pi_0}}\ \underset{\substack{\vert\\ \text{\scriptsize since $t\geq r^{\pi_0}$,}\\ \text{\scriptsize hence $tr^{\pi_0}=t$}}}{=}\ s^t\ =\ s'\ ,
\]
\printedpage{18}
% head: A. Grothendieck : Dérivateurs, Ch. XV, Édition des Universités de Montpellier II et Paris VII
% foot: Version 5 novembre 2010
so we see that $f\in\Fl^{s'}$ implies $X_1(f)\in\M^{s'}$, or what amounts to the same thing, $p_1(f)\in\Fl^{s'}$.

\textbf{Corollary.} \textit{Let $C_0$, $r$, $s'=s^t$ $(t\geq r^{\pi_0},\ s\geq 2)$ be as above, and let $\alpha$ be an ordinal number such that}
\[
\card\alpha\ \leq\ s'\ =\ s^t\ .
\]
\textit{Then $\uFl^{s'}(\M)$ is stable under the functor $p_\alpha=p_\alpha^{C_0}$.}

We see this at once by transfinite induction, taking into account lemma 3.3, and the stability assumption (3.13).

To summarize:

\textbf{Theorem 2.} \textit{Let $\M$ be a category stable under small filtering $\varinjlim$, and endowed with a cardinal filtration} $(\Filt^\pi)_{\pi\geq\pi_0}$\textit{, such that the} $\Filt^\pi$ \textit{are stable under filtering $\varinjlim$ of cardinality $\leq\pi$ (not necessarily large relative to $\pi_0$), and under the simple amalgamated sums that exist in $\M$. We assume that the elements of} $\Filt^{\pi_0}$ \textit{are accessible, and we denote by $\pi_1$ the smallest cardinal $\geq\pi_0$, such that the objects of} $\Filt^{\pi_0}$ \textit{are $\pi_1$-accessible. We put} \msp{page 27}
\[
c_1\ =\ \card\Fl(K_1),\qquad K_1\subset\M^{\pi_1}\ \text{a chosen small full skeleton}
\]
\DERmark{ENV0866}

\textit{(where for every $\pi\geq\pi_0$, $\M^\pi$ denotes the strictly full subcategory of $\M$ defined by} $\Filt^\pi$\textit{, and} $\Fl^\pi(\M)=\Fl(\M^\pi)$\textit{).}

\textit{Let $C_0$ be a small subset of} $\Fl(\M)$\textit{, consisting of cosquarable arrows, and put}
\[
\Phi\ =\ \widetilde{C_0}\ =\ \Psi^*\ ,\qquad\quad\textit{where}\ \ \Psi=C_{0*}=\Phi_*\ .
\]
\textit{Let $r$ be a cardinal} $\geq\sup(\pi_1,c_1)$\textit{, and such that we have $C_0\subset\Fl^r(\M)$. We assume that $C_0$ is irredundant (}i.e.\textit{ two isomorphic objects of $C_0$ are equal), or more generally, that $C_0$ satisfies the very harmless non-plethoricity condition} \textup{(3.11)} \textit{relative to $r$. We consider the transfinite sequence of functors $p_\alpha=p_\alpha^{C_0}$ on $\uFl(\M)$} \textup{(}\textit{cf.} \textup{th. 1)}\textit{, giving rise to the functorial factorizations of an $f\in\uFl(\M)$, $f:X\longrightarrow Y$, as}
\[
X\ \xrightarrow{i_\alpha(f)}\ X_\alpha(f)\ \xrightarrow{p_\alpha(f)}\ Y\ . \tag*{(3.16)}
\]

\textit{We then have the following:}

$1^\circ$) \textit{We have (for the record) $i_\alpha(f)\in\Phi$ for every $\alpha$, and moreover,} if $\alpha$ is large relative to $r^{\pi_0}$, \textit{we also have $p_\alpha(f)\in\Psi$. If $r'$ is a cardinal $>r^{\pi_0}$ (e.g. the successor cardinal of $r^{\pi_0}$), we can take for $\alpha$ the first ordinal whose cardinality is equal to $r'$.}

$2^\circ$) \textit{For every cardinal $s$ of the form}
\[
s\ =\ c^t\qquad\quad\textit{with}\ \ t\geq r^{\pi_0},\ c\geq 2
\]
\textit{(thus $s=2^t$ if we take $c\leq t$,} cf. \textit{lemma 3.2), $\uFl^s(\M)$ is stable under the $p_\alpha$ such that} $\card\alpha\leq s$\textit{. If we take e.g. for $\alpha$ the first ordinal whose cardinality is $>r^{\pi_0}$, then $\uFl^s(\M)$ is stable under $p_\alpha$, which by $1^\circ$) thus realizes a functorial factorization} \textup{(3.16)} \textit{with $i_\alpha(f)\in\Phi$, $p_\alpha(f)\in\Psi$.}


\endgroup
% END UNIT Der-ChXV-05-11-10/Der-ChXV-05-11-10_p016-020.tex
% BEGIN UNIT Der-ChXV-05-11-10/Der-ChXV-05-11-10_p021-025.tex
\begingroup\makeatletter%
\newcommand{\printedpage}[1]{\marginpar{\footnotesize p.~#1}}%
\newcommand{\msp}[1]{\marginpar{\footnotesize [#1]}}%
\newcommand{\Fl}{\mathrm{Fl}}%
\newcommand{\id}{\mathrm{id}}%
\newcommand{\qpre}{\mathrel{\leftarrow\!\!\!\!\mapsto}}%
\makeatother\let\printedpage\derpp\let\msp\dermsp%
\let\mspin\DERmspin%


\printedpage{19}
% head: A. Grothendieck : Dérivateurs, Ch. XV, Édition des Universités de Montpellier II et Paris VII
% foot: Version 5 novembre 2010

\msp{page 28}
\textsc{Proof}. Everything has been proved, if we recall that the objects of $\mathcal{M}^r$ are $r^{\pi_0}$-accessible (\textit{cf.} (3.2)).

\section*{4\quad Complements on Quillen structures}

\textbf{Definition 4.1.} Let $\mathcal{M}$ be a category. By a \textit{Quillen precouple} in $\mathcal{M}$ we mean a pair $(\Phi,\Psi)$ of subsets of $\Fl(\mathcal{M})$, satisfying the following two conditions.
\begin{itemize}[leftmargin=2.5em]
\item[\textit{a})] We have $\Phi\longleftrightarrow\Psi$, \textit{i.e.} $\Psi\subset\Phi_*$, or what amounts to the same, $\Phi\subset\Psi^*$ -- or again: for every commutative square in $\mathcal{M}$
\[
\begin{tikzcd}[row sep=large, column sep=large]
A \arrow[r, "u"] \arrow[d, "i"'] & X \arrow[d, "p"] \\
B \arrow[ur, dotted, "w" description] \arrow[r, "v"] & Y \;,
\end{tikzcd}
\tag*{(4.1.1)}
\]
with $i\in\Phi$, $p\in\Psi$, there exists $w$ such that $wi=u$, $pw=v$.
\item[\textit{b})] Every $f\in\Fl(\mathcal{M})$ factors as
\[
f \;=\; p\,i\;, \qquad \text{with } i\in\Phi,\, p\in\Psi\ .
\]
\end{itemize}
We write
\[
\Phi \;\qpre\; \Psi
\tag*{(4.1.2)}
\] % ?? uncertain: the printed relation sign -- a double arrow with a small vertical bar in the middle (rendered by \qpre)
to indicate that $(\Phi,\Psi)$ is a Quillen precouple.

We say that the Quillen precouple is \textit{closed}, if $\Phi$ is left closed, $\Psi$ right closed. We shall see below (cor. 2, page 31) that this implies (hence is equivalent to) % [sic] clos à droit -> clos à droite, gender agreement with « à droite »
\[
\Phi \;=\; \Psi^*\;, \qquad \Psi \;=\; \Phi_*\ .
\tag*{(4.1.3)}
\]

\msp{page 29}
A closed Quillen precouple is also called a Quillen couple.

\textbf{Proposition 4.2.} \textit{Let} $(\Phi,\Psi)$ \textit{be a Quillen precouple in} $\mathcal{M}$\textit{, and let} $j\in\Fl(\mathcal{M})$\footnote[9]{\textbf{N.B.} We use only that $\Phi\longleftrightarrow\Psi$, and that $j$ can be put in the form $pi$, with $i\in\Phi$, $p\in\Psi$.}. % footnote mark printed as "(9)"
\textit{We then have the equivalent conditions on} $j$ \textit{:}
\begin{itemize}[leftmargin=2.5em]
\item[a)] $j\in\widetilde{\Phi}$.
\item[b)] $j\in\Psi^*$.
\item[c)] $j$ \textit{is a direct factor of an} $i\in\Phi$.
\item[c$'$)] $j$ \textit{is a direct factor of an} $i\in\Phi$ \textit{having the same source as} $j$\textit{, more precisely, there exists a diagram}
\[
\begin{tikzcd}[row sep=large, column sep=large]
A \arrow[r, "j"] \arrow[d, shift right=1ex, "\id_A"'] & B \arrow[d, shift right=1ex, "q"'] \\
A \arrow[u, dotted, shift right=1ex, "\id_A"'] \arrow[r, "i"] & B' \;, \arrow[u, dotted, shift right=1ex, "p"']
\end{tikzcd}
\tag*{(4.2.1)}
\]
\printedpage{20}
% head: A. Grothendieck : Dérivateurs, Ch. XV, Édition des Universités de Montpellier II et Paris VII
% foot: Version 5 novembre 2010
\textit{with} $i\in\Phi$\textit{, and}
\[
qj \;=\; i\;, \qquad pi \;=\; j\;, \qquad pq \;=\; \id_B\;,
\tag*{(4.2.2)}
\]
\textit{so that} $j$ \textit{is a direct factor of} $i\in\Phi$\textit{, via} $(\id_A,q)$\textit{, admitting the retraction} $(\id_A,p)$.
\item[d)] \textit{For every commutative solid-arrow diagram}
\[
\begin{tikzcd}[column sep=small]
A \arrow[rr, "j"] \arrow[dr, "i"'] & & B \arrow[dl, dotted, bend left=40, "q"] \\
& B' \arrow[ur, "p"'] &
\end{tikzcd}
\;,
\tag*{(4.2.3)}
\]
i.e. \textit{every factorization of} $j$ \textit{as} $pi$\textit{, if we have} $p\in\Psi$\textit{, there exists a section} $q$ \textit{of} $B'$ \textit{over} $B$ \textit{in} $A\backslash\mathcal{M}$\textit{,} i.e. \textit{an arrow} $q:B\longrightarrow B'$ \textit{such that}
\[
pq \;=\; \id_B \qquad \textit{and} \qquad qj \;=\; i\ .
\]
\end{itemize}

\textsc{Proof} of the proposition. It proceeds by the following sequence of implications, all $(\pm)$ tautological:
\[
(c') \;\Longrightarrow\; (c) \;\Longrightarrow\; (a) \;\Longrightarrow\; (b) \;\Longrightarrow\; (d) \;\Longrightarrow\; (c')\;,
\]
\msp{page 30}
of which the first four use only the hypothesis $\Phi\longleftrightarrow\Psi$ (from $(a)$ of definition 4.1), and only the last, $(d)\Longrightarrow(c')$, calls on the Quillen factorization hypothesis ($(b)$ in 4.1); and even then it suffices to assume it for $j$ only!

$(c')\Longrightarrow(c)$ pure tautology, $(c)\Longrightarrow(a)$ since we know that $\widetilde{\Phi}$ is stable under direct factors (and contains $\Phi$), $(a)\Longrightarrow(b)$ since $\widetilde{\Phi}\subset\Psi^*$ ($\widetilde{\Phi}$ being by definition the smallest left closed subset of $\Fl(\mathcal{M})$ containing $\Phi$, and $\Psi^*$ being left closed), $(b)\Longrightarrow(d)$ since we write the given diagram (4.2.3) (solid arrows) as % [sic] (4.2.2) -> (4.2.3), the diagram given in condition d) is (4.2.3)
\DERmark{ENV0867}

\[
\begin{tikzcd}[row sep=large, column sep=large]
A \arrow[r, "u"] \arrow[d, "j\in\Psi^*"'] & X \arrow[d, "p\in\Psi"] \\
B \arrow[ur, dotted, "q" description] \arrow[r, "v"] & Y \;,
\end{tikzcd}
\tag*{$(*)$}
\]
putting $Y=B$, $v=\id_B$ in it, and rebaptizing $X$ as $B'$, and $u$ as $i$: the hypothesis $(b)$, $j\in\Psi^*$, means that in a solid-arrow diagram $(*)$, there always exists $q$ such that $qj=u$, $pq=v$, and expressing this in the particular case where $v=\id_B$, we find $(d)$.

It remains to prove $(d)\Longrightarrow(c')$, this time using the factorization hypothesis on $j$, \textit{i.e.} that we can write $j$ in the form $pi$:
\[
\begin{tikzcd}[column sep=small]
A \arrow[rr, "j"] \arrow[dr, "i"'] & & B \arrow[dl, dotted, bend left=40, "q"] \\
& B' \arrow[ur, "p"'] &
\end{tikzcd}
\;,
\]
with $i\in\Phi$, $p\in\Psi$. By virtue of $(d)$, since $p\in\Psi$, there thus exists $q:B\longrightarrow B'$ satisfying
\printedpage{21}
% head: A. Grothendieck : Dérivateurs, Ch. XV, Édition des Universités de Montpellier II et Paris VII
% foot: Version 5 novembre 2010
\msp{page 31}
\[
pq \;=\; \id_B\;, \qquad qj \;=\; i\;,
\]
which proves $(c')$, since $i\in\Phi$.

\textbf{Corollary 1.} \textit{Let} $(\Phi,\Psi)$ \textit{be a Quillen precouple. Then we have}
\[
\begin{aligned}
\widetilde{\Phi} \;\;=\;\; \Psi^* \;\;=\;\; &\textit{set of objects of } \underline{\Fl}(\mathcal{M}) \textit{ which are direct fac-}\\
&\textit{tors of elements of } \Phi\ ,
\end{aligned}
\tag*{(4.3)}
\] % printed line break of the display falls inside "fac-teurs"; kept as printed
\textit{and dually}
\[
\begin{aligned}
\overline{\Psi} \;\;=\;\; \Phi_* \;\;=\;\; &\textit{set of objects of } \underline{\Fl}(\mathcal{M}) \textit{ which are direct fac-}\\
&\textit{tors of elements of } \Psi\ .
\end{aligned}
\tag*{(4.4)}
\] % printed line break of the display falls inside "fac-teurs"; kept as printed

\textbf{Corollary 2.} \textit{The following conditions on the Quillen precouple} $(\Phi,\Psi)$ \textit{are equivalent.}
\begin{itemize}[leftmargin=2.5em]
\item[a)] \textit{The precouple is} closed, i.e. $\Phi=\widetilde{\Phi}$\textit{,} $\Psi=\overline{\Psi}$\textit{.}
\item[b)] $\Phi=\Psi^*$\textit{,} $\Psi=\Phi_*$\textit{.}
\item[c)] $\Phi$ \textit{and} $\Psi$ \textit{are stable under direct factors.}
\end{itemize}
Now let
\[
W \;\subset\; \Fl(\mathcal{M})
\tag*{(4.5)}
\]
be a localizer, satisfying the usual conditions
\[
\left\{
\begin{aligned}
&\textit{a})\quad W \text{ contains the isomorphisms.}\\[1ex]
&\textit{b})\quad \text{If } X\xrightarrow{\;f\;}Y\xrightarrow{\;g\;}Z \text{ in } \mathcal{M}\text{, if two of } f,\ g,\ gf \text{ are}\\
&\phantom{\textit{b})}\quad \text{in } W\text{, then the third too.}
\end{aligned}
\right.
\tag*{(4.6)}
\]
($W$ mildly saturated). Suppose moreover given four other subsets
\[
C\ ,\ TC\ , \qquad F\ ,\ TF
\tag*{(4.7)}
\]
of $\Fl(\mathcal{M})$, satisfying the following conditions.
\[
\left\{
\begin{aligned}
&\textit{a})\quad (C,TF) \text{ and } (TC,F) \text{ are Quillen precouples.}\\[1ex]
&\textit{b})\quad TC\subset C\cap W,\ TF\subset F\cap W.
\end{aligned}
\right.
\tag*{(4.8)}
\]

\msp{page 32}
These properties are spelled out as follows:
\begin{itemize}[leftmargin=2.5em]
\item[\textbf{M\,1}] Lifting-extension properties, $C\longleftrightarrow TF$ and $TC\longleftrightarrow F$.
\item[\textbf{M\,2}] Factorization property for every arrow $f$ of $\mathcal{M}$ as $f=pi$, with, at our choice, $i\in C$, $p\in TF$, or $i\in TC$, $p\in F$.
\item[\textbf{M$'$}] $TC\subset C\cap W$, $TF\subset F\cap W$.
\end{itemize}

\printedpage{22}
% head: A. Grothendieck : Dérivateurs, Ch. XV, Édition des Universités de Montpellier II et Paris VII
% foot: Version 5 novembre 2010
This is therefore part of the axioms of a weak Quillen structure for $W$ (XIII, p. 27, def. 1) -- the stability axioms M\,3, M\,4, and the ``properness'' (propriety) axioms M$''$ are missing, namely $F\subset\mathrm{Fib}_W$, $C\subset\mathrm{Cof}_W$.

\textbf{Proposition 4.3.} \textit{Let} $\mathcal{M}$ \textit{be a category,} $(C,TF)$ \textit{and} $(TC,F)$ \textit{two Quillen precouples in} $\mathcal{M}$\textit{, satisfying}
\[
TC \;\subset\; C\;, \qquad TF \;\subset\; F\;,
\tag*{(4.3.1)}
\]
\textit{and let} $W$ \textit{be a subset of} $\Fl(\mathcal{M})$ \textit{of which for the moment we assume only} 4.6, (b)\footnote[10]{We use only that if $X\xrightarrow{\;f\;}Y\xrightarrow{\;g\;}Z$, and if $gf\in W$, then $f\in W\Longleftrightarrow g\in W$. We do not need $W$ to contain the isomorphisms, nor to be stable under composition.}\textit{, and} % footnote mark printed as "(10)"
\[
W \;\supset\; TC\ ,\ TF\ ,
\tag*{(4.3.2)}
\]
\textit{so we have}
\[
TC \;\subset\; C\cap W\ , \qquad TF \;\subset\; F\cap W\ .
\tag*{(4.3.3)}
\]
\textit{Let us consider the following extension-lifting conditions} (a), (a$'$)\textit{, and the factorization conditions} (b), (b$'$), (c), (d) (\textit{recalling that} $\widetilde{TC}=F^*$\textit{,} $\overline{TF}=C_*$\textit{, by virtue of proposition} 4.2 (\textit{cor.} 1) \textit{applied to the Quillen precouples} $(TC,F)$\textit{,} $(C,TF)$):
\begin{itemize}[leftmargin=2.5em]
\item[a)] $C\cap W\subset\widetilde{TC}$\quad (\textit{hence} $TC\subset C\cap W\subset\widetilde{TC}$).
\item[a$'$)] $F\cap W\subset\overline{TF}$\quad (\textit{hence} $TF\subset F\cap W\subset\overline{TF}$).
\item[b)] $W\subset TF\circ\widetilde{TC}$.
\item[b$'$)] $W\subset\overline{TF}\circ TC$.
\item[c)] $W\subset\overline{TF}\circ\widetilde{TC}$.
\item[d)] $W\subset TF{\circ}TC$ (\textit{hence, by} (4.3.2)\textit{,} $W=TF{\circ}TC$ \textit{if} $W$ \textit{is stable under composition}).\footnote[11]{\textbf{N.B.} Condition $(a)$ also means that $(C\cap W,F)$ is still a Quillen precouple, resp. $(b)$ means that $(C,F\cap W)$ is one -- and then this pair of Quillen precouples this time satisfies $(d)$, \textit{i.e.} $W=TF'\circ TC'$.} % footnote mark printed as "(11)"
\end{itemize}

\msp{page 33}
\textit{Then all the conditions except the last are equivalent} (\textit{and they are therefore implied by the latter}). \textit{Thus when} $TC$ \textit{or} $TF$ \textit{are stable under direct factors} (\textit{e.g. if} $TC$ \textit{is left closed, or} $TF$ \textit{right closed}), \textit{all the conditions under consideration are equivalent} (\textit{since by corollary} 2 \textit{to proposition} 4.2\textit{, we then have} $TC=\widetilde{TC}$ \textit{resp.} $TF=\overline{TF}$\textit{, hence} (b$'$) \textit{resp.} (b) \textit{is then equivalent to} (d)).

\textsc{Proof}. It proceeds by the diagram of implications
\[
\begin{tikzcd}[column sep=small, row sep=small]
& (d) \arrow[dl, Rightarrow] \arrow[dr, Rightarrow] & \\
(b) \arrow[dr, Rightarrow] & \mathrm{Q} & (b') \arrow[dl, Rightarrow] \\
& (c) \arrow[dl, Rightarrow] \arrow[dr, Rightarrow] & \\
(a) \arrow[uu, Rightarrow] & & (a') \arrow[uu, Rightarrow]
\end{tikzcd}
\;,
\]
\printedpage{23}
% head: A. Grothendieck : Dérivateurs, Ch. XV, Édition des Universités de Montpellier II et Paris VII
% foot: Version 5 novembre 2010
where the four implications of the upper square Q are tautological. It thus suffices to establish the implications
\[
(c) \;\Longrightarrow\; (a) \;\Longrightarrow\; (b)\;,
\]
the implications $(c)\Longrightarrow(a')\Longrightarrow(b')$ follow from these by duality.

Suppose then $(c)$, \textit{i.e.}
\[
W \;\subset\; \overline{TF}\circ\widetilde{TC}\;,
\tag*{$(*)$}
\]
and let us prove $(a)$, \textit{i.e.} that
\[
C\cap W \;\subset\; \widetilde{TC}\ .
\]
But if $j\in C\cap W$, to prove that $j\in\widetilde{TC}$, let us write $j=pi$, $i\in\widetilde{TC}$, $p\in\overline{TF}$, which is possible by the hypothesis $(*)$ and $j\in W$.

\msp{page 34}
\[
\begin{tikzcd}[column sep=small]
A \arrow[rr, "j\in C"] \arrow[dr, "i\in\widetilde{TC}"'] & & B \arrow[dl, dotted, bend left=60, "q"] \\
& B' \arrow[ur, "p\in\overline{TF}"'] &
\end{tikzcd}
\;.
\]
Let us now use the hypothesis $j\in C$ and $C\longleftrightarrow\overline{TF}$, we thus find (prop. 4.2, $(a)\Longrightarrow(d)$) an arrow $q:B\longrightarrow B'$ such that $pq=\id_B$, $qj=i$, so $j$ appears as a direct factor of $i$. Since $i\in\widetilde{TC}$ and $\widetilde{TC}$ is stable under direct factors, we conclude $j\in\widetilde{TC}$, \quad \textit{q.e.d.}
\DERmark{ENV0868}


$(a)\Longrightarrow(b)$. We thus assume
\[
C\cap W \;\subset\; \widetilde{TC}\;,
\tag*{$(*)$}
\]
and we want to prove $(b)$, \textit{i.e.}
\[
W \;\subset\; TF\circ\widetilde{TC}\ .
\]
Let then $f\in W$, and let us factor $f$ as $pi$, with $i\in C$, $p\in TF$. (O.K. because of the Quillen precouple $(C,TF)$). Since $TF\subset W$, we have $p\in W$, and since $pi\in W$, we conclude $i\in W$ (mild saturation of $W$ -- we do have to use it after all!). We thus have $i\in C{\cap}W$, hence $i\in\widetilde{TC}$ by $(*)$, whence $f\in TF\circ\widetilde{TC}$, \quad \textit{q.e.d.}

\textbf{Corollary 1.} \textit{Suppose} $TC$ \textit{stable under direct factors. Then all the conditions of proposition} 4.3 \textit{are equivalent, and are equivalent to the relation}
\[
C\cap W \;=\; TC\ .
\tag*{(4.3.4)}
\]
\textit{Dually, if} $TF$ \textit{is stable under direct factors, the conditions of proposition} 4.3 \textit{are all equivalent and are equivalent to}
\[
F\cap W \;=\; TF\ .
\tag*{(4.3.4$'$)}
\]


\endgroup
% END UNIT Der-ChXV-05-11-10/Der-ChXV-05-11-10_p021-025.tex
% BEGIN UNIT Der-ChXV-05-11-10/Der-ChXV-05-11-10_p026-030.tex
\begingroup\makeatletter%
\newcommand{\printedpage}[1]{\marginpar{\footnotesize p.~#1}}%
\newcommand{\msp}[1]{\marginpar{\footnotesize [#1]}}%
\newcommand{\Fl}{\mathrm{Fl}}%
\newcommand{\Cof}{\mathrm{Cof}}%
\newcommand{\deq}{\overset{\mathrm{def}}{=}}%
\newcommand{\univ}{\mathrm{univ}}%
\makeatother\let\printedpage\derpp\let\msp\dermsp%
\let\mspin\DERmspin%

\DERmark{ENV0869}


\printedpage{24}
% head: A. Grothendieck : Dérivateurs, Ch. XV, Édition des Universités de Montpellier II et Paris VII
% foot: Version 5 novembre 2010

\msp{page 35}
Indeed, by corollary 2 of proposition 4.2, the hypothesis on $TC$ implies that $TC = \widetilde{TC}$, so since $TC \subset C{\cap}W$, the hypothesis $C{\cap}W \subset \widetilde{TC}$ $(= TC)$ is equivalent to $C{\cap}W = TC$.

\textbf{Corollary 2.} \textit{Suppose $C$ and $W$ stable under direct factors. Then the equivalent conditions of proposition} 4.3 \textit{are equivalent to}
\[
C\cap W \;=\; \widetilde{TC}\ . \tag*{(4.3.5)}
\]
\textit{Dually, if $F$, $W$ are stable under direct factors, the equivalent conditions of proposition} 4.3 \textit{are equivalent to}
\[
F\cap W \;=\; \overline{TF}\ . \tag*{(4.3.5$'$)}
\]

Indeed, if $C$ and $W$ are stable under direct factors, the same is true of $C\cap W$, and hence if $C\cap W\subset\widetilde{TC}$, since $C\cap W$ contains $TC$, we conclude $\widetilde{C\cap W}=\widetilde{TC}$. But since $C\cap W\longleftrightarrow F$, $(C\cap W,F)$ is a Quillen precouple, and $C\cap W$ is stable under direct factors, $C\cap W$ is left closed by corollary 2 of proposition 4.2, \textit{i.e.} $C\cap W=\widetilde{C\cap W}$, whence $C\cap W=\widetilde{TC}$, \quad \textit{q.e.d.}
\DERmark{ENV0870}


\textbf{Corollary 3.} \textit{Suppose the equivalent conditions} (a) \textit{to} (c) \textit{of} 4.3 \textit{satisfied, and moreover} $W$ stable under direct factors\footnote[12]{In fact, it suffices to suppose that we have $W\supset\widetilde{TC}$, $\overline{TF}$, \textit{i.e.} that every direct factor of an element of $W$ which is either in $C$, or in $F$, is still in $W$. Then $(\widetilde{C},\widetilde{TC},\bar{F},\overline{TF})$ define a closed Quillen structure, and we know that this implies that in fact, $W$ too is stable under direct factors.}. \textit{Then the preliminary conditions} (4.3.1) \textit{and} (4.3.2) \textit{remain satisfied when one replaces $(C,TF)$ and $(TC,F)$ by their closures $(C',TF')=(\widetilde{C},\overline{TF})$ and $(TC',F')=(\widetilde{TC},\bar{F})$, and the equivalent conditions of proposition} 4.3 \textit{are still satisfied for these couples, in the strong forms}
\[
\left\{
\begin{aligned}
C'\cap W &\;=\; TC' && (=\widetilde{TC'}\,)\\
F'\cap W &\;=\; TF' && (=\overline{TF'}\,)\\
W &\;=\; TF'\circ TC'\ .
\end{aligned}
\right.
\tag*{(4.3.6)}
\]

\msp{page 36}
\textit{The triple $(W,C',F')$ is a closed Quillen triple $($corresponding to a ``closed model category'' of Quillen, at least if one assumes $\mathcal{M}$ stable under finite $\varinjlim$ and finite $\varprojlim$, and that $W$ is mildly saturated$)$.}\footnote[13]{This gives the\par
\textbf{Corollary 4.} \textit{Let $W$ be a mildly saturated localizer, and $(C,TF)$, $(TC,F)$ two closed Quillen couples such that $TC\subset C\cap W$, $TF\subset F\cap W$. Then the following conditions are equivalent.}
\begin{enumerate}[leftmargin=*]
\item[a)] $TC=C\cap W$.
\item[b)] $TF=F\cap W$.
\item[c)] $W=TF\circ TC$.
\item[d)] $(W,C,F)$ \textit{is a closed Quillen triple.}
\end{enumerate}}
\DERmark{ENV0871}


\textsc{Proof.} The relations (4.3.1) imply
\[
\widetilde{TC}\;\subset\;\widetilde{C}\qquad\text{and}\qquad\overline{TF}\;\subset\;\bar{F}\ ,
\]
\printedpage{25}
% head: A. Grothendieck : Dérivateurs, Ch. XV, Édition des Universités de Montpellier II et Paris VII
% foot: Version 5 novembre 2010
\textit{i.e.}
\[
TC'\;\subset\;C'\ ,\qquad TF'\;\subset\;F'\ ,
\]
\textit{i.e.} (4.3.1) for the new Quillen couples. The relation 4.3.2, $W\supset TC,TF$, implies
\[
W\;\supset\;\widetilde{TC}\ ,\ \overline{TF}\ ,\qquad \textit{i.e.}\qquad W\;\supset\;TC'\ ,\ TF'\ ,
\]
by virtue of corollary 2 to proposition 4.2 and the fact that $W$ is stable under direct factors; so we have (4.3.2) for the new couples. Since by hypothesis we have
\[
W\;=\;\overline{TF}\circ\widetilde{TC}\ ,
\]
this reads
\[
W\;=\;TF'\circ TC'\ ,
\]
so the new couples satisfy the same conditions, which this time, the couples being closed, take the strong form (4.3.6).

Let us prove that the triple $(W,C',F')$ satisfies Quillen's axioms for a closed Quillen triple. We do have, by (4.3.6),
\[
\left\{
\begin{aligned}
C' &\;=\;(TF')^*\;=\;(F'\cap W)^*\ , &\qquad F' &\;=\;(TC')_*\;=\;(C'\cap W)_*\ ,\\
W &\;=\;TF'\circ TC'\;=\;C'_*\circ F'^*\ ,
\end{aligned}
\right.
\]
which corresponds to the definition of Quillen's \textit{closed} model categories: Quillen's axioms M\,1, M\,2, M\,3, M\,4 are trivially satisfied (M\,3 and M\,4 because
\msp{page 37}
$C$ and $F$, $TC$ and $TF$ are left closed resp. right closed). Axiom M\,5 is the mild saturation of $W$.

\textbf{4.4.} Here is the idea for constructing more or less ``all'' the closed Quillen structures $(C,TC,F,TF)$ associated to a given localizer $W\subset\Fl(\mathcal M)$. We suppose that $\mathcal{M}$ satisfies the conditions of the preceding paragraph, concerning the existence of filtering $\varinjlim$, of a cardinal filtration, accessibility in $\mathcal{M}$ (that of the elements of $\mathrm{Filt}^{\pi_0}$). We start from a small subset
\DERmark{ENV0872}

\[
C_0\;\subset\;\Fl(\mathcal{M})
\]
such that
\[
C_{0*}\;\subset\;W\ , \tag*{(4.4.1)}
\]
and we set
\[
\left\{
\begin{aligned}
C &\;=\;\widetilde{C_0}\ , &\qquad TC &\;=\;C\cap W\ ,\\
TF &\;=\;C_*\;=\;C_{0*}\ , &\qquad F &\;=\;(TC)_*\ .
\end{aligned}
\right.
\tag*{(4.4.2)}
\]
There is only one problem, since we already know that $(C,TF)$ is a Quillen couple by the factorization theorem in §1 (p. 10): we need the factorization for the couple $(TC,F\;=\;(TC)_*)$. At that point, the preliminary conditions (4.3.1), (4.3.2) of proposition 4.3 are satisfied:
\[
TC\;\subset C\ ,\qquad TF\;\subset F\ ,\qquad W\;\supset\;TC\ ,\ TF\ ,
\]
as well as condition $(a)$, $C\cap W\subset\widetilde{TC}$, since $C\cap W=TC$. So we find, by the equivalence of conditions $(a)$, $(a')$, $(c)$, and since $TC$ and $TF$ are closed,
\[
TF\;=\;F\cap W\ ,\qquad W\;=\;TF\circ TC\ ,
\]
\printedpage{26}
% head: A. Grothendieck : Dérivateurs, Ch. XV, Édition des Universités de Montpellier II et Paris VII
% foot: Version 5 novembre 2010
\msp{page 38}
on the other hand the hypothesis on $W$ implies that $W$ is stable under direct factors, $C$ is too, being left closed, so $TC=C{\cap}W$ also, so by proposition 4.2, corollary 2, we have $TC\;=\;\widetilde{TC}$, $TC$ is left closed. So we win, by the last assertion of corollary 3, proposition 4.3 (p. 35).
\DERmark{ENV0873}


In fact, independently of the particular conditions on $\mathcal{M}$, $W$, we have just proved this:

\textbf{Lemma 4.4.1} (corollary of proposition 4.3.). \textit{Let $\mathcal{M}$ be a category, $W$ a localizer in $\mathcal{M}$ mildly saturated and stable under direct factors. Then the closed Quillen structures $(C,F)$ associated to $W$ are obtained as follows: one takes a subset $C_0$ of} $\Fl(\mathcal{M})$ \textit{such that}
\[
TF\;\deq\;C_{0*}\;\subset\;W\ ,
\]
\textit{one sets}
\[
C\;=\;\widetilde{C_0}\;=\;TF^*\ ,\qquad TC\;=\;C\cap W\ ,\qquad F\;=\;(TC)_*\ ,
\]
\textit{and one must manage to get $(C,TF)$ and $(TC,F)$ to satisfy Quillen's factorization condition,} i.e. \textit{be Quillen precouples.} (In fact, $C$, $TF$ and $F$ being closed by construction, and $TC=C\cap W$ being stable under direct factors ($C$ and $W$ being so), they will even be Quillen \textit{couples}.)

\msp{page 39}
The additional hypotheses made earlier on $\mathcal{M}$ served to ensure that if I take $C_0$ small, then the factorization works for $(C,TF)$ automatically. There remains the problem for $TC$, $F$.

Let us note that if it works, then $TC=\widetilde{TC}=F^*$, so $TC$ is left closed, and \textit{a fortiori} satisfies the corresponding stability conditions:
\[
\left\{
\begin{tabular}{rp{0.6\textwidth}}
\textit{a}) & Stable under cobase change, under composition, under direct factors (these last two are ensured automatically, since $C$, $W$ are stable).\\[1ex]
\textit{b}) & Stable under strict ordinal $\varinjlim$.
\end{tabular}
\right.
\tag*{(4.4.3)}
\]

The stability hypothesis on $W$ serves (among other things) to guarantee stability $(b)$ for $TC$ (in fact, for this alone it would suffice that $W$ be stable under strict ordinal $\varinjlim$). There remains the stability of $TC=C{\cap}W$ under cobase change, which is problematic since $W$ is in general \textit{not} stable under cobase change. It is to ensure this stability that I find myself practically compelled here to impose
\[
C_0\;\subset\;\Cof_W\ , \tag*{(4.4.4)}
\]
I say that with the stability hypothesis made on $W$, this implies
\[
C\;\deq\;\widetilde{C_0}\subset\Cof_W\ . \tag*{(4.4.5)}
\]
\printedpage{27}
% head: A. Grothendieck : Dérivateurs, Ch. XV, Édition des Universités de Montpellier II et Paris VII
% foot: Version 5 novembre 2010
\msp{page 40}
Indeed, we have:

\textbf{Proposition 4.4.2.} \textit{Let $\mathcal{M}$ be stable under small filtering direct $\varinjlim$, and $W\subset$} $\Fl(\mathcal{M})$ \textit{a localizer $($mildly saturated$)$} stable under the said limits. \textit{Then} $\Cof_W$ \textit{satisfies the stability conditions} (4.4.3) (``\textit{left Quillen stability}''), \textit{and is also stable under arbitrary filtering $\varinjlim$.}

Indeed, condition $(a)$ is satisfied without any additional condition on $W$. On the other hand, it is immediate that this condition with respect to filtering $\varinjlim$ implies that $\Cof_W$ \textit{is stable under filtering direct limits} in $\underline{\Fl}(\mathcal{M})$.

\textbf{Corollary 4.4.3.} \textit{Under the conditions of proposition} 4.4.2, \textit{and assuming that every arrow factors as a $W$-cofibration followed by a $W$-equivalence,} $\Cof_W\cap W=W^{\univ}$ $($\textit{set of co-universal $W$-equivalences in $\mathcal{M}$}$)$ \textit{is stable under filtering direct limits, and hence ``left stable'' in the sense of} (4.4.3).
\DERmark{ENV0874}


\textbf{Corollary 4.4.4.} \textit{Under the conditions of proposition} 4.4.2, \textit{suppose moreover that simple amalgamated sums exist in $\mathcal{M}$, and let $C_0$ be a small subset of} $\Cof_W$ \textit{whose elements are accessible, and consider}
\[
C\;=\;\widetilde{C_0}\ ,\qquad TC\;\deq\;C\cap W\ .
\]
\textit{Then we have}
\[
C\;\subset\;\Cof_W\ ,\qquad TC\;=\;C\cap W^{\univ}\;\subset\;W^{\univ}\ ,
\]
\DERmark{ENV0875}

\msp{page 41}
\textit{and $C$ and $TC$ are left $Q$-stable.}

Indeed, by theorem 1 of paragraph 1 (p. 10) and its corollary, $\widetilde{C_0}$ is the smallest subset of $\Fl(\mathcal{M})$ containing $C_0$ and left Q-stable. Since $\Cof_W$ is left Q-stable and contains $C_0$, we therefore have $C=\widetilde{C_0}\subset\Cof_W$. So $TC\deq C{\cap}W\subset\Cof_W{\cap}W\deq W^{\univ}$, so $TC=C\cap W^{\univ}$. Since $W^{\univ}$ is itself also left $Q$-stable (preceding corollary), the same is true of $C\cap W$, \quad \textit{q.e.d.}

\textbf{Corollary 4.4.5.} \textit{Suppose that every object of $\mathcal{M}$ is accessible. Under the conditions of the preceding corollary, setting}
\[
TF\;=\;C_{0*}\;=\;C_*\ ,\qquad F\;=\;(TC)_*\ ,
\]
\textit{the couple $(C,TF)$ is a Quillen couple. For the same to hold for the couple $(TC,F)$, it suffices that there exist a} small \textit{subset $TC_0$ of $C\cap W$, such that $\widetilde{TC_0}=TC$ $(\deq C\cap W)$.}

The first assertion is, as we saw earlier, theorem 1 of paragraph 1. The second follows from this same theorem. We can be more precise in this direction:

\textbf{Proposition 4.4.6.} \textit{Let $\mathcal{M}$ be a category stable under small filtering $\varinjlim$, under amalgamated sums, and all of whose objects are accessible, $W$ a localizer}
\msp{page 42}
\textit{mildly saturated in $\mathcal{M}$, stable under filtering $\varinjlim$, $C_0$ a} small \textit{subset of} $\Fl(\mathcal{M})$, \textit{such that}
\[
C_0\;\subset\;\Cof_W\ ,\qquad \underbrace{C_{0*}}_{(=TF)}\;\subset\;W\ ,
\]
\printedpage{28}
% head: A. Grothendieck : Dérivateurs, Ch. XV, Édition des Universités de Montpellier II et Paris VII
% foot: Version 5 novembre 2010
\textit{let us set}
\[
C\;=\;\widetilde{C_0}\qquad TC\;=\;C\cap W\ ,\qquad TF\;=\;C_*\ ,\qquad F\;=\;(TC)_*\ ,
\]
\textit{and let $TC_0$ be a small subset of $TC=C\cap W$. Put $F_0=(TC_0)_*$. We then have}
\DERmark{ENV0876}

\[
\left\{
\begin{aligned}
C &\;\subset\;\Cof_W\ , &\qquad TC &\;=\;C\cap W^{\univ}\quad(\textit{for the record})\\
\widetilde{TC_0} &\;\subset\;TC\ , & F_0 &\supset F\quad(\textit{the latter trivial})\ ,
\end{aligned}
\right.
\tag*{(4.4.6)}
\]
\textit{and the following conditions are equivalent:}
\begin{enumerate}[leftmargin=*]
\item[a)] $\widetilde{TC_0}=TC\quad(\textit{we have}\subset)$\,;
\item[a$'$)] $F_0=F\quad(\textit{we have}\supset)$\,;
\item[b)] $TF=F_0\cap W\quad(\textit{we have}\subset)$\footnote[14]{\textbf{N.B.} We see easily that $\boxed{TF=F\cap W,}$ but that does not suffice to have $TF=F_0\cap W$ for a suitable $TC_0$.}\,;
\item[c)] $W=TF\circ\widetilde{TC_0}\quad(\textit{we have}\supset)$\footnote[15]{\textbf{N.B.} We know that $\boxed{W=TF\circ TC,}$ which is in no way sufficient to imply $(c)$\,!}\,;
\end{enumerate}
\textit{and imply that $(W,C,F)$ is a closed Quillen triple.}

\textsc{Proof.} The inclusion $\widetilde{TC_0}\;\subset\;TC$ follows from the left Q-stability of $TC$ (cor. 4.4.4) and from the fact that $TC_0\subset TC$ is small (one then applies theorem 1 just as for $C\subset\Cof_W\ldots$). The equivalence of $(a)$, $(b)$, $(c)$ is none other than proposition 4.3 applied to $W$ and the two Quillen couples $(C,TF)$, $(\widetilde{TC_0},F_0)$ (taking into account theorem 1 of paragraph 1). That we then have a closed Quillen triple was seen in corollary 3 of proposition 4.3 (p. 35). It remains to prove the equivalence of $(a)$ and $(a')$. We have $(a)\Longrightarrow(a')$, since $(TC_0)_*=(\widetilde{TC_0})_*=F_0$. Conversely, $(a')$ implies $\widetilde{TC_0}=\widetilde{TC}$, hence the factorization theorem for $(TC,F)$ by virtue of the factorization theorem, corollary 1 (p. 11) applied to $TC_0\subset TC\subset\widetilde{TC_0}$ and to $F_0=F=(TC_0)_*$.

\msp{page 43}
Thus, the only problem, once $C_0$ (small) has been found satisfying $C_0\subset\Cof_W$, $C_{0*}\subset W$, is to find a \textit{small} subset $TC_0\subset TC$ which ``generates $TC\;=\;C\cap W$ '', \textit{i.e.} which satisfies the equivalent conditions $(a)$, $(a')$, $(b)$, $(c)$ of the preceding proposition. I have spent days trying to prove the existence of $TC_0$, assuming $\mathcal{M}$ equipped with a cardinal filtration, and $W$ an accessible subset of $\mathcal{M}$, without having succeeded. Of course, all that really matters to us is that $(TC,F)$ satisfy the factorization condition. But I see no way of proving it, except by ``generating'' $TC$ on the left by a \textit{small} subset $TC_0$.

\section*{5\quad Preliminary technical results on accessible subsets of $\mathrm{Ob}\,\mathcal{M}$, stable under filtering direct limits}

Yet I think I can arrive at something usable, by working systematically with $C$'s ($\subset\Cof_W$, left Q-closed, left Q-generated by a \textit{small} subset $C_0$ of $C$) which are stable under all filtering direct limits.


\endgroup
% END UNIT Der-ChXV-05-11-10/Der-ChXV-05-11-10_p026-030.tex
% BEGIN UNIT Der-ChXV-05-11-10/Der-ChXV-05-11-10_p031-035.tex
\begingroup\makeatletter%
\newcommand{\printedpage}[1]{\marginpar{\footnotesize p.~#1}}%
\newcommand{\msp}[1]{\marginpar{\footnotesize [#1]}}%
\newcommand{\cM}{\mathcal{M}}%
\newcommand{\cC}{\mathcal{C}}%
\newcommand{\Filt}{\mathrm{Filt}}%
\newcommand{\uFilt}{\underline{\mathrm{Filt}}}%
\newcommand{\Ss}{\mathrm{Ss}}%
\newcommand{\Ob}{\mathrm{Ob}\,}%
\newcommand{\Img}{\mathrm{Im}\,}%
\newcommand{\Fl}{\mathrm{Fl}}%
\newcommand{\uFl}{\underline{\mathrm{Fl}}}%
\newcommand{\card}{\mathrm{card}\,}%
\newcommand{\isoto}{\xrightarrow{\ \sim\ }}%
\newcommand{\hookto}{\lhook\joinrel\longrightarrow}%
\newcommand{\circled}[1]{\tikz[baseline=(c.base)]\node[draw,circle,inner sep=1pt](c){#1};}%
\makeatletter%
\renewcommand\@makefnmark{\hbox{\textsuperscript{\normalfont(\@thefnmark)}}}%
\renewcommand\@makefntext[1]{\parindent 1em\noindent\hbox to 1.8em{\hss\textsuperscript{\normalfont\@thefnmark}}#1}%
\makeatother%
\makeatother\let\printedpage\derpp\let\msp\dermsp%
\let\mspin\DERmspin%


\printedpage{29}
% head: A. Grothendieck : Dérivateurs, Ch. XV, Édition des Universités de Montpellier II et Paris VII
% foot: Version 5 novembre 2010

I again assume the following conditions on $\cM$:
\msp{page 44}

\begin{enumerate}[nosep]
\item[\textit{a})] $\cM$ stable under small direct limits (and not only filtering ones).
\item[\textit{b})] $\cM$ endowed with a cardinal filtration $(\Filt^{\pi})_{\pi\geq\pi_0}$.
\item[\textit{c})] The objects of $\cM$ are accessible (\textbf{N.B.} It suffices to assume this for the objects of $\Filt^{\pi_0}$).
\end{enumerate}

Condition ($b$) implies that $\cM$ contains a small generating family by strict epimorphisms (e.g. a small subcategory $\cC$ of $\uFilt^{\pi_0} = \cM^{\pi_0}$ such that the inclusion $\cC \longrightarrow \uFilt^{\pi_0}$ is an equivalence of categories). By SGA 4 I 8.12.7, it follows that $\cM$ is stable under small inverse limits. There follows the

\textbf{Lemma 5.1} (under conditions ($a$), ($b$)). \textit{For every object $X$ of $\cM$, the set} $\Ss(X)$ \textit{of subobjects of $X$, endowed with the natural order relation ``of inclusion'', admits} inf \textit{and} sup \textit{of arbitrary size.}

inf, since these are expressed as inverse limits in $\cM$. Now the existence of inf implies that of sup, since $\sup_J = \inf_{J'}$, where $J'$ is the set of upper bounds of $J$ in $\Ss(X)$.

\textbf{Corollary 5.2.} \textit{For every arrow $f : Z \longrightarrow X$ in $\cM$, there exists, among the subobjects $X'$ of $X$ through which $f$ factors, a smallest one} % [sic] lesquelles -> lesquels, "sous-objets" is masculine
\msp{page 45}
\textit{(which we will denote} $\Img f$\textit{). The morphism $Z \longrightarrow \Img f$ is an epimorphism.}

This last assertion comes from this one:

\textbf{Lemma 5.3.} \textit{Let $\cM$ be a category in which kernels of double arrows exist. Let $f : Z \longrightarrow Y$ be an arrow of $\cM$. Then $f$ is an epimorphism if and only if $f$ does not factor through a subobject that is} strict, $Y_0$, \textit{of $Y$ different from $Y$.} % [sic] g -> f, the arrow has been named f
\DERmark{ENV0877}


We are going to make an additional (practically innocuous) hypothesis on the $\Filt^{\pi}$:

\begin{enumerate}[nosep]
\item[\textit{d})] $\Filt^{\pi}$ is stable under direct limits indexed by index categories $I$ such that $\card \Fl(I) \leq \pi$\footnote[16]{This condition will be called ($d_0$) below.}. (It suffices to assume stability under initial object, simple amalgamated sums, and filtering $\varinjlim$ indexed by an ordered set $I$ of cardinality $\leq \pi$.) Moreover, if $f : Z \longrightarrow X$ is an arrow of $\cM$ such that $Z$ or $X$ is in $\Filt^{\pi}$, then $\Img f \in \Filt^{\pi}$.
\end{enumerate}

This last condition on $\Img f$ splits into two:
\msp{page 46}

\begin{enumerate}[nosep]
\item[$d_1$)] Every subobject of an object of $\Filt^{\pi}$ is in $\Filt^{\pi}$.
\item[$d_2$)] If $Z$ is in $\Filt^{\pi}$, then every quotient $Z'$ of $Z$ such that $\Img p = Z'$ (where $p : Z \longrightarrow Z'$ is the canonical morphism) is in $\Filt^{\pi}$.\footnote[17]{It would seem that we will not need ($d_1$), only ($d_2$). But in fact, ($d_1$) must be a consequence of ($d_2$), and ($d_2$) is a consequence of the first part of ($d$) (stability under finite $\varinjlim$) if every epimorphism in $\cM$ is strict, hence effective (if the $\Filt^{\pi}$ are stable under finite $\varprojlim$ as well \ldots).}
\end{enumerate}

\printedpage{30}
% head: A. Grothendieck : Dérivateurs, Ch. XV, Édition des Universités de Montpellier II et Paris VII
% foot: Version 5 novembre 2010

Of course, in practice it will always be true that $\Filt^{\pi}$ is stable under passage to subobjects and quotients of an object (not only the particular quotients at issue in ($d_2$)). % [sic] closing parenthesis of "(pas seulement" missing, parentheses unbalanced as printed
\DERmark{ENV0878}

Besides, when every bimorphism (\textit{i.e.} monomorphism \textit{and} epimorphism) of $\cM$ is an isomorphism, then ($d_2$) is indeed equivalent to stability of $\Filt^{\pi}$ under quotients of objects.

\textbf{Proposition 5.3.} \texttt{[Twice 5.3.]} \textit{Suppose conditions} (a), (b), (d) \textit{are satisfied for $\cM$. Let $\pi$ be a cardinal $\geq \pi_0$, and $X$ an object of $\cM$. Let} $\Ss_{\pi}(X)$ \textit{be the set of subobjects of $X$ that are in} $\Filt^{\pi}$\textit{. This set, ordered by the order relation induced by} $\Ss(X)$\textit{, is large relative to $\pi$ (and in particular filtering). If, moreover, filtering direct limits of monomorphisms in $\cM$ are monomorphisms}\footnote[18]{This is an \textit{exactness condition} on $\cM$, which I will denote as:\par
\noindent\circled{e}\quad Monomorphisms in $\cM$ are stable under filtering $\varinjlim$ in $\uFl(\cM)$.} \textit{(e.g. if filtering direct limits are left exact), or if the objects of} $\Filt^{\pi}$ \textit{are $\pi$-accessible, we have $\varinjlim\limits_{\alpha\in\Ss_{\pi}(X)} Z_{\alpha} \isoto X$}\footnote[19]{Say also that $\Ss_{\pi}(X)$ is $\pi$-adapted, and the direct system $(Z_{\alpha})_{\alpha}$ as well.}.

\msp{page 47}
\textsc{Proof}. Let $(Z_{\alpha})_{\alpha\in A}$ be a family of subobjects of $X$ such that $Z_{\alpha} \in \Filt^{\pi}$, and $\card A \leq \pi$. Let $Z = \coprod_{\alpha\in A} Z_{\alpha}$, we have $Z \in \Filt^{\pi}$ by ($d_0$). The smallest subobject of $X$ that is an upper bound for the $Z_{\alpha}$ is none other than $\Img (Z \longrightarrow X)$, which is in $\Filt^{\pi}$ by ($d_2$). Let us prove that $\varinjlim\limits_{\alpha\in\Ss_{\pi}(X)} Z_{\alpha} \isoto X$ is an isomorphism, when we know that it is a monomorphism. Now
\DERmark{ENV0879}\DERmark{ENV0880}

\[
X \;=\; \varinjlim_{I} X_i \ ,
\]
where $(X_i)_{i\in I}$ is a filtering direct system, $I$ an ordered set large relative to $\pi$, the $X_i$ in $\Filt^{\pi}$. For every $i \in I$, let $X'_i$ be the image of $X_i$ in $X$. We therefore have $X'_i \in \Filt^{\pi}$ by ($d_2$). Using the fact that the $X_i \longrightarrow X'_i$ are epimorphisms, we see at once that $X = \varinjlim\limits_{I} X'_i$. Consider the map
\[
\begin{aligned}
I \quad &\longrightarrow \quad \Ss_{\pi}(X)\\
i \quad &\longmapsto \quad X'_i \ .
\end{aligned}
\]
This is an increasing map, let $J$ be its image, a filtering subset, large relative to $\pi$, of $\Ss_{\pi}(X)$, such that
\[
\varinjlim_{\alpha\in J} Z_{\alpha} \ \isoto \ X \ .
\]
\msp{page 48}
We therefore have
\[
\underbrace{\varinjlim_{\alpha\in J} Z_{\alpha}}_{\simeq X} \ \longrightarrow \ \varinjlim_{\alpha\in\Ss_{\pi}X} Z_{\alpha} \ \xrightarrow{\ \varphi\ } \ X \ ,
\]
the composite being the identity. If we know that $\varphi$ is a monomorphism, it follows that it is an isomorphism (a monomorphism having a section). Now $\varphi$ will be a monomorphism if every filtering $\varinjlim$ of monomorphisms is a monomorphism, and then we win.

\printedpage{31}
% head: A. Grothendieck : Dérivateurs, Ch. XV, Édition des Universités de Montpellier II et Paris VII
% foot: Version 5 novembre 2010

If we assume that the elements of $\Filt^{\pi}$ are $\pi$-accessible, we win as well, for we find:

\textbf{Corollary 5.4.} \textit{If the objects of} $\Filt^{\pi}$ \textit{are $\pi$-accessible (e.g. if $\pi \geq \pi_1$ and $\pi = \pi^{\pi_0}$, e.g. $\pi \geq \pi_1$ and $\pi = 2^{c}$, with $c \geq \pi_0$,} cf. (3.2) \textit{p. 16), then for every filtering subset $J$ of} $\Ss(X)$ \textit{large relative to $\pi$, such that $\varinjlim\limits_{\alpha\in J} Z_{\alpha} \isoto X$, $J$ is cofinal in} $\Ss_{\pi}(X)$ \textit{(hence $\varinjlim\limits_{\alpha\in J} Z_{\alpha} \isoto \varinjlim\limits_{\alpha\in\Ss(X)} Z_{\alpha}$).}

This is immediate.

\msp{page 49}
\textbf{Scholie 5.5.} When we make on $\cM$ the fairly innocuous exactness hypothesis that filtering $\varinjlim$ of monomorphisms are monomorphisms, \textit{or} restrict ourselves to cardinals $\pi$ such that the objects of $\Filt^{\pi}$ are $\pi$-accessible, we can \textit{canonize} the representation of an object $X$ as $\varinjlim\limits_{i\in I} X_i$, with $X_i \in \Filt^{\pi}$ and $I$ large relative to $\pi$, by taking $I = \Ss_{\pi}(X)$, and the corresponding direct system. We also have the following pleasant additional property: every filtering subset $J$ of $I$, of cardinality $\leq \pi$, admits a supremum $i_J$, and (if we assume that $\cM$ satisfies condition ($e$), page 46, note 18), we will have
\[
X_{i_J} \;=\; \varinjlim_{i\in J} X_i \ ;
\]
in other words, we have a representation of $X$ by a direct system that is (perfectly) \textit{$\pi$-adapted} (\textit{cf.} XIII, p. 85, 86).

\msp{page 50}
Now let
\[
M \;\subset\; \Ob\cM
\]
be a subset of $\Ob\cM$, stable under the isomorphism relation. If $M$ is $\pi$-accessible (according to condition $\mathrm{b}_{\pi}$ of XIII, p. 88), then putting
\[
M^{\pi} \;=\; M \cap \Filt^{\pi} \ ,
\]
$M^{\pi}$ determines $M$ by the following condition on $X \in \Ob\cM$.
\[
\begin{aligned}
&X \in M \iff \text{the subset } \Ss^{M}_{\pi}(X) \text{ of } \Ss_{\pi}(X) \text{ consisting of the } X_{\alpha} \in \Ss_{\pi}(X)\\
&\text{such that } X_{\alpha} \in M \ (\iff X_{\alpha} \in M^{\pi}) \text{ is cofinal in } \Ss_{\pi}(X).
\end{aligned}
\tag*{(5.1)}
\]
Note that when $X \in \Filt^{\pi}$, this criterion is trivially valid, since then $\Ss_{\pi}(X)$ admits $X$ as greatest element, and the cofinal subsets are those containing this final object. Thus (5.1) is interesting only for the $X \in \Ob\cM \backslash \Filt^{\pi}$. Note that for \textit{every} subset $M^{\pi}$ of $\Filt^{\pi}$, we can define a subset $M$ of $\Ob\cM$ by condition (5.1), and we will trivially have
\[
M^{\pi} \;=\; M \cap \Filt^{\pi} \ . \tag*{(5.2)}
\]
But probably, although $M$ satisfies (5.1), it is not for all that necessarily $\pi$-accessible. Indeed, for it to be so, it must be stable under $\varinjlim$
\printedpage{32}
% head: A. Grothendieck : Dérivateurs, Ch. XV, Édition des Universités de Montpellier II et Paris VII
% foot: Version 5 novembre 2010
\msp{page 51}
large relative to $\pi$ of elements of $M^{\pi}$, which implies that every direct factor of an element of $M^{\pi}$ is in $M$, hence also in $M^{\pi}$, according to the stability condition ($d$) on $\Filt^{\pi}$. This is a particular condition on the subset $M^{\pi}$ of $\Filt^{\pi}$. So, we will have to have a condition on $\cM$ stronger than the validity of (5.1), namely:
\[
\parbox{0.72\textwidth}{For every $X \in \Ob\cM$ and for every subset $J$ of $\Ss_{\pi}(X)$, such that
\begin{enumerate}[nosep]
\item[\textit{a})] $J$ is cofinal in $\Ss_{\pi}(X)$ and
\item[\textit{b})] If $J' \subset J$ is a filtering subset of $J$, of cardinality $\leq \pi$, then its supremum $\alpha_{J'}$ in $\Ss_{\pi}(X)$ belongs to $J$,
\end{enumerate}
we have $X \in M$ if and only if $J^{M} \overset{\mathrm{def}}{=} J \cap \Ss^{M}_{\pi}(X)$ is cofinal in $J$ (or equivalently, in $\Ss_{\pi}(X)$)\,\footnotemark[20].}
\tag*{(5.3)}
\]
\DERmark{ENV0881}

\footnotetext[20]{Instead of $J$ cofinal in $\Ss_{\pi}(X)$, it even suffices to have $\varinjlim\limits_{\alpha\in J} X_{\alpha} \isoto X$ -- the corresponding condition on $M$ thus becomes more demanding. But if the objects of $\Filt^{\pi}$ are $\pi$-accessible, this condition is in fact equivalent to the weaker condition.}

If we already admit the ``it suffices'' of (5.1), then it suffices in (5.3) to assume $X \in M$ and require that $J^{M}$ then be cofinal in $J$.

But note the
\msp{page 52}

\textbf{Proposition 5.6.}\footnote[21]{False as it stands, true if we assume that $M^{\pi}$ is stable under filtering $\varinjlim$.} \textit{Suppose $\pi$ chosen in such a way that the objects of} $\Filt^{\pi}$ \textit{are $\pi$-accessible (e.g. $\pi = 2^{c}$, with $c \geq \pi_1$,} cf. (3.2), \textit{p. 16). Then, for a subset $M$ of} $\Ob\cM$\textit{, $M$ is $\pi$-accessible if and only if it satisfies criterion} (5.3).

We have seen that this is necessary (tautologically, in view of proposition 5.3 and scholie 5.5, bottom of page 49). Let us prove that this is sufficient, by first establishing the

\textbf{Lemma 5.7.} \textit{Under the conditions of} 5.6 \textit{on $\pi$, and condition} (e) \textit{on $\cM$} (p. 46), \textit{for every $\pi$-adapted representation $X = \varinjlim\limits_{i\in I} X_i$ of an object $X$ of $\cM$, the subset $I'$ of those $i \in I$ such that $X_i \longrightarrow X$ is a monomorphism, is cofinal in $I$. Moreover, $I'$ is stable under suprema in $I$ of filtering subsets $J$ of $I'$ of cardinality $\leq \pi$.}

Let us show how this lemma 5.7 implies the ``it suffices'' of proposition 5.6. So let $X = \varinjlim\limits_{i\in I} X_i$ be a $\pi$-adapted representation of $X$ ($X_i \in \Filt^{\pi}$, $I$ large relative to $\pi$, \textit{etc.}). We must prove that, if 5.3 is valid, then $X \in M$ if and only if the set $I_M$ is cofinal in $I$.

\msp{page 53}
Assume $X \in M$. Consider the subset $I' \subset I$ described in the lemma. We have a map $I' \longrightarrow \Ss_{\pi}(X)$, namely $i \longmapsto X'_i = \Img (X_i \hookto X)$, and we have
\[
X \simeq \varinjlim_{i\in I} X_i \ \simeq \ \varinjlim_{i\in I'} X_i \ \simeq \ \varinjlim_{i\in I'} X'_i \ \simeq \ \varinjlim_{\alpha\in J} Z_{\alpha} \ ,
\]
where $J$ is the image of $I'$ in $\Ss_{\pi}(X)$. Since the objects of $\Filt^{\pi}$ are $\pi$-accessible, we conclude that $J$ is cofinal in $\Ss_{\pi}(X)$. By 5.3, it follows that the set $J_M$ of those $\alpha \in J$
\printedpage{33}
% head: A. Grothendieck : Dérivateurs, Ch. XV, Édition des Universités de Montpellier II et Paris VII
% foot: Version 5 novembre 2010
such that $Z_{\alpha} \in M$ is cofinal in $J$, hence $\neq \emptyset$. Thus $I'_M$, which is the inverse image of $J_M$ in $I'$, is in any case nonempty. Applying the same result while replacing $I'$ by the subset $I'_{\geq i'_0}$, we see that $I'_M$ is cofinal in $I'$, hence also in $I$, since $I'$ is cofinal in $I$; thus $I_M \supset I'_M$ is cofinal in $I$, \quad \textit{q.e.d.}

Conversely, suppose that $I_M$ is cofinal in $I$, and prove $X \in M$. The difficulty here comes from the fact that the representation $X = \varinjlim\limits_{i\in I_M} X_i$ is no longer adapted to $\pi$ -- $I_M$ is not necessarily stable in $I$ under suprema (in $I$) of filtering subsets in $I_M$, $I'$ of cardinality $\leq \pi$. In short, we must
\msp{page 54}
prove that $M$ is stable under $\varinjlim$ large relative to $\pi$ without more, and I now doubt that this follows from just the hypotheses made here, which only concern direct systems where the transition morphisms are monomorphisms.

Here, I think, is a counterexample. Let us take
\[
M \quad = \quad \parbox[t]{0.6\textwidth}{set of objects of $\cM$ isomorphic to a fixed object $X_0 \in \Filt^{\pi}$, such that $X_0$ has no endomorphism which is a monomorphism other than the identity.}
\tag*{(5.4)}
\]
Then condition 5.3 is satisfied. For:
\begin{enumerate}[nosep]
\item[\textit{a})] If $X \in M$, \textit{i.e.} $X \simeq X_0$, then $X \in \Filt^{\pi}$, hence $\Ss_{\pi}(X)$ admits $X$ as final object, and $J$, being cofinal, contains this element, hence $J_M$ (which also contains it) is cofinal.
\end{enumerate}
And:
\begin{enumerate}[nosep]
\item[\textit{b})] If $J_M$ is cofinal, hypothesis (5.4) implies that the set of transition morphisms between the $Z_{\alpha}$ ($\alpha \in J_M$) consists of isomorphisms, hence the direct limit $X$ of the $Z_{\alpha}$ in question is isomorphic to each of them, hence $X \simeq X_0$, \textit{i.e.} $X \in M$.
\end{enumerate}
Yet, $M$ is $\pi$-accessible only if it is stable under $\varinjlim$ large relative to $\pi$, which implies that it is stable under \textit{direct factors.} But hypothesis
\msp{page 55}
(5.4) in no way rules out that $X_0$ admits a projector $p$ ($p^{2} = p$) that is nontrivial. This will have an image $X'$ in $\cM$, which will be a direct factor of $X_0$ not isomorphic to $X_0$, \textit{i.e.} not in $M$.

These difficulties disappear if we assume that $M^{\pi}$ is stable under filtering $\varinjlim$ indexed by ordered sets $I$ of cardinality $\leq \pi$. Indeed, in this case, in the preceding proof, the subset $I_M$ is stable under suprema (in $I$) of subsets $J$ of $I_M$ that are filtering and of cardinality $\leq \pi$. Thus the direct system $(X_i)_{i\in I_M}$ is still $\pi$-adapted, and we can apply lemma 5.7 to it, to find a cofinal subset $I'_M \subset I_M$, such that the $X_i \longrightarrow X$ for $i \in I'_M$ are monomorphisms. The image $J$ of $I'_M$ in $\Ss_{\pi}(X)$ is then a subset of $\Ss_{\pi}(X)$ such that $\varinjlim\limits_{\alpha\in J} Z_{\alpha} = X$, hence (by the hypothesis made on $\pi$) cofinal in $\Ss_{\pi}(X)$ (5.4), with $Z_{\alpha} \in M^{\pi}$ for every $\alpha \in J$. Thus by 5.3 we have $X \in M$, \textit{q.e.d.}

\msp{page 56}
It remains to prove lemma 5.7. With the hypotheses and notation of the lemma, so let $i_0 \in I$, we must find $i' \in I'$ such that $i' \geq i_0$. Consider the image $X'_i$ of $X_i$ in $X$, and


\endgroup
% END UNIT Der-ChXV-05-11-10/Der-ChXV-05-11-10_p031-035.tex
% BEGIN UNIT Der-ChXV-05-11-10/Der-ChXV-05-11-10_p036-040.tex
\begingroup\makeatletter%
\newcommand{\printedpage}[1]{\marginpar{\footnotesize p.~#1}}%
\newcommand{\msp}[1]{\marginpar{\footnotesize [#1]}}%
\newcommand{\Filt}{\mathrm{Filt}}%
\newcommand{\Ob}{\mathrm{Ob}}%
\newcommand{\Ss}{\mathrm{Ss}}%
\newcommand{\M}{\mathcal{M}}%
\newcommand{\Ci}{\mathcal{C}}%
\makeatother\let\printedpage\derpp\let\msp\dermsp%
\let\mspin\DERmspin%


\printedpage{34}
% head: A. Grothendieck : Dérivateurs, Ch. XV, Édition des Universités de Montpellier II et Paris VII
% foot: Version 5 novembre 2010
the inclusion
\[
\underbrace{X'_{i_0}}_{\in\, \Filt^\pi} \;\longrightarrow\; X \;\simeq\; \varinjlim X_i \;.
\]
Since $X'_{i_0} \in \Filt^\pi$, hence is $\pi$-accessible, there is an $i_1 > i_0$ such that $X'_{i_0} \longrightarrow X$ factors through
\[
X'_{i_0} \;\longrightarrow\; X_{i_1} \;.
\]
Applying this same result to $X_{i_1}$, $X'_{i_1}$ \textit{etc.}, we find an infinite sequence of morphisms
\[
\begin{tikzcd}[column sep=small]
X_{i_0} \arrow[dr] & & X_{i_1} \arrow[dr] & & X_{i_2} \arrow[dr] & \\
& X'_{i_0} \arrow[ur] & & X'_{i_1} \arrow[ur] & & X'_{i_2} \quad \cdots
\end{tikzcd}
\]
such that the morphisms composed two at a time are, on the upper line, the transitions
\[
X_{i_m} \longrightarrow X_{i_n} \qquad (\; i_m \;<\; i_n \;) \;,
\]
and on the lower line, the inclusions
\[
X'_{i_m} \hookrightarrow X'_{i_n} \qquad (\; \hookrightarrow X \;) \;.
\]
Let $i_\omega = \sup_n i_n$ in $I$, hence
\[
X_{i_\omega} \;=\; \varinjlim_{n} X_{i_n} \;\simeq\; \varinjlim_{n} X'_{i_n} \;\simeq\; X'_{i_\omega} \;,
\]
so $X_{i_\omega} \longrightarrow X$ is a monomorphism.

\msp{page 57}

The last assertion of 5.7 follows at once from the fact that a filtering direct limit of monomorphisms in $\M$ is a monomorphism (condition $(e)$, p.~46).

We can make proposition 5.6 obtained in this way more precise:

\textbf{Proposition 5.7}. \texttt{[Twice 5.7.]} \textit{Let $\M$ be a category satisfying conditions} (a), (b), (c), (d) (p.~44, 45) and (e) (p.~46). \textit{Let $M$ be a subset of} $\Ob\M$\textit{, stable under isomorphism, $\pi$ a cardinal such that the objects of} $\Filt^\pi$ \textit{are $\pi$-accessible. The following conditions are equivalent:} % [sic] Soit $\M$ (calligraphique) -> Soit $M$, la partie de Ob M est notée M partout ailleurs
\DERmark{ENV0882}


\begin{enumerate}[label={}]
\item[a)] \textit{$M$ is $\pi$-accessible, and stable under filtering $\varinjlim$.}
\item[b)] \textit{$M$ is $\pi$-accessible, and $M \cap \Filt^\pi$ is stable under filtering $\varinjlim$ of cardinality $\leq \pi$.}
\item[c)] \textit{There is a subset $M^\pi$ of} $\Filt^\pi$\textit{, stable under filtering $\varinjlim$ of cardinality $\leq \pi$ and such that for every $X \in \Ob\M$, putting $J=\Ss_\pi(X)$, we have $X \in M$ if and only if the set $J_M$ of the $Z_\alpha \in J$ such that $Z_\alpha \in M^\pi$ is cofinal in $J$.}
\end{enumerate}
\DERmark{ENV0883}


(\textit{Moreover, under condition} (c), \textit{we have}
\printedpage{35}
% head: A. Grothendieck : Dérivateurs, Ch. XV, Édition des Universités de Montpellier II et Paris VII
% foot: Version 5 novembre 2010
\msp{page 58}
$M^\pi = M \cap \Filt^\pi$.)
\begin{enumerate}[label={}]
\item[c$'$)] \textit{There is a subset $M_0$ of} $\Filt^\pi$ \textit{such that $M$ is the smallest subset of $\M$ containing $M_0$ and stable under filtering direct limits.} (\textit{Moreover, we can then take $M_0$} small.)\footnote[22]{\textbf{N.B.} We have
\[
\begin{tikzcd}[ampersand replacement=\&]
(a) \arrow[r, Leftrightarrow] \arrow[d, Rightarrow] \& (b) \arrow[d, Rightarrow] \\
(c') \arrow[r, Rightarrow, "?"] \& (c)
\end{tikzcd}
\quad ,
\]
without conditions on $\pi$ (except $\pi \geq \pi_0$).}.
\DERmark{ENV0884}

\end{enumerate}

\textbf{Corollary 5.8.} \textit{The maps}
\[
M \;\longmapsto\; M^\pi = M \cap \Filt^\pi \;,
\]
\[
M^\pi \quad\longmapsto\quad \parbox[t]{0.5\textwidth}{\textit{the set of $X \in \Ob\M$ such that the set $\Ss_\pi^{M^\pi}(X)$ of the $Z_\alpha \in \Ss_\pi(X)$ which are in $M^\pi$ is cofinal} (= \textit{smallest subset of $\Ob\M$ containing $M^\pi$ and stable under filtering $\varinjlim$})}
\]
\DERmark{ENV0885}

\textit{establish bijections, inverse to one another, between the set of $M \subset \Ob\M$ satisfying the equivalent conditions of proposition} 5.7\textit{, and the set of subsets $M^\pi$ of} $\Filt^\pi$ \textit{stable under filtering $\varinjlim$ of cardinality $\leq \pi$.}

\textsc{Proof of} 5.7. The implication $(b) \Longrightarrow (c)$ is tautological, in view of the definition of the $\pi$-accessibility of $M$, and of proposition 5.3. And $(a) \Longrightarrow (b)$ comes from the stability of $\Filt^\pi$ under filtering $\varinjlim$ of cardinality $\leq \pi$ (condition $(d)$).

It therefore remains to prove $(c) \Longrightarrow (a)$. (It is here that we must use the special hypothesis on $\pi$.) Once this is seen, we shall know the equivalence of $(a)$, $(b)$, $(c)$, and on the other hand, it is immediate that $(a) \Longrightarrow (c')$, by
\msp{page 59}
taking $M_0 = M \cap \Filt^\pi$. On the other hand, $(c') \Longrightarrow (c)$ (without hypotheses on $\pi\,$?), since we can suppose that the subset $M_0$ of $\Filt^\pi$ is stable under filtering $\varinjlim$ of cardinality $\leq \pi$. But then, knowing the equivalence of $(a)$, $(b)$, $(c)$, whence the corollary at once, it follows that the envelope stable under filtering $\varinjlim$ of $M_0$ in $\M$ is none other than the accessible subset deduced from $M_0$ that is in question in $(c)$, and we win. % [sic] $(b)$ $(c)$ -> $(b)$, $(c)$, virgule manquante
\DERmark{ENV0886}


So it all comes down to proving that, if all the objects of $\Filt^\pi$ are $\pi$-accessible, then $(c)$ implies $(a)$. Let us first show that $(c)$ implies the membership criterion (5.3) (p.~51), so by proposition 5.6 ``revised and corrected'' in note 21, it already follows that $(c) \Longrightarrow (b)$, and it will remain to establish $(b) \Longrightarrow (a)$, \textit{i.e.} that $(b)$ implies the stability of $M$ under arbitrary filtering $\varinjlim$ (not only large relative to $\pi$).
\DERmark{ENV0887}


That $(c)$ implies (5.3) will follow from the following lemma, applied to $I = \Ss_\pi(X)$:

\textbf{Lemma 5.8.} \texttt{[Twice 5.8.]} \textit{Let $I$ be an ordered set that is} $\pi$-adapted, i.e. \textit{filtering and such that every filtering subset $K$ of $I$, of cardinality $\leq \pi$, has a supremum $i_K$ in $I$. I say that a subset $J$ of $I$ is} $\pi$-cofinal, \textit{if it is cofinal, and if for every filtering subset $K$ of $J$ of cardinality $\leq \pi$, its supremum $i_K$ in $I$ is}
\printedpage{36}
% head: A. Grothendieck : Dérivateurs, Ch. XV, Édition des Universités de Montpellier II et Paris VII
% foot: Version 5 novembre 2010
\msp{page 60}
\textit{in $J$.} (\textit{It follows, then, that $J$ too is $\pi$-adapted, for the induced order structure.}) \textit{This said, the intersection of two $\pi$-cofinal subsets of $I$ is again $\pi$-cofinal.}

It is enough to see that it is cofinal, for its stability under suprema in $I$ of filtering subsets of cardinality $\leq \pi$ is tautological (if it is true for some $J_\alpha$, it is true for $\bigcap_\alpha J_\alpha$). Let $i_0 \in I$, and let us construct an element of $J \cap J'$ that dominates $i_0$. We construct step by step a double sequence
\[
\begin{tikzcd}
i_1 \arrow[dr] & & i_2 \arrow[dr] & & \cdots \\
& i'_1 \arrow[ur] & & i'_2 \arrow[ur] &
\end{tikzcd}
\]
with the $i_\alpha$ in $J$, the $i'_\alpha$ in $J'$, and $i_1 \geq i_0$. We thus have
\[
\sup i_\alpha \;=\; \sup i'_\alpha \;,
\]
and if $i$ is the common value of the two sides, we have $i \in J \cap J'$.

This seen, to prove the validity of (5.3), we note that the subset $J' = \Ss_\pi^{(M^\pi)}(X)$ of $I = \Ss_\pi(X)$, consisting of the $Z_\alpha \subset X$ such that $Z_\alpha \in M^\pi$, is stable under suprema of filtering subsets of cardinality $\leq \pi$, since $M^\pi$ is stable
\msp{page 61}
under $\varinjlim$ of cardinality $\leq \pi$. So if $J$ is a $\pi$-cofinal subset of $\Ss_\pi(X)$, $J \cap J'$ is cofinal in $J$, \textit{i.e.} in $I$, if and only if $J'$ is cofinal in $I$ (\textit{i.e.} if and only if $X \in M$, by condition $(c)$). For $J \cap J'$ cofinal $\Longrightarrow$ $J'$ cofinal tautologically, and $J'$ cofinal implies $J'$ $\pi$-cofinal, hence $J \cap J'$ $\pi$-cofinal by lemma 5.8, hence cofinal.

It remains only to prove that $(b) \Longrightarrow (a)$, \textit{i.e.} the stability of $M$ under arbitrary filtering direct limits, assuming $(b)$. Here again, this should not use the special hypothesis on $\pi$, nor even the exactness hypothesis $(e)$ on $\M$, but be a general story of accessibility:

\textbf{Proposition 5.9.} \textit{Let $\M$ be a category satisfying} (a), (b), (c), (d$_0$) (pages 44, 45), \textit{$\pi \geq \pi_0$ a cardinal, $M$ a $\pi$-accessible subset of} $\Ob\M$, i.e. \textit{satisfying condition $b_\pi$ of} XIII, p.~88. \textit{Then $M$ is stable under filtering $\varinjlim$ if and only if $M^\pi \overset{\mathrm{def}}{=} M \cap \Filt^\pi$ is stable under filtering $\varinjlim$ of cardinality $\leq \pi$.}
\DERmark{ENV0888}


I don't feel like spelling out the proof here, having the painful feeling of forever redoing the same proof (here and in XIII). I shall have to come back to this in a more settled situation.

\msp{page 62}

I am going to sketch the proof all the same. Let $X = \varinjlim\limits_{I} X_i$, $I$ filtering, the $X_i$ in $M$, we must prove that $X \in M$ (assuming the stability hypothesis on $M^\pi = M \cap \Filt^\pi$).
\printedpage{37}
% head: A. Grothendieck : Dérivateurs, Ch. XV, Édition des Universités de Montpellier II et Paris VII
% foot: Version 5 novembre 2010
\[
\begin{tikzcd}
& \Ci \arrow[r] \arrow[d] \arrow[dl, dotted, "\mathfrak{z}"'] \arrow[dr, phantom, "\text{\footnotesize cart.}"] & I \arrow[d, "\varphi"] \\
\underline{M}^\pi & \underline{M}^\pi/\underline{M} \arrow[l, dotted, "\text{\footnotesize ``source''}"] \arrow[r, "\text{\footnotesize ``target''}"'] & \underline{M}
\end{tikzcd}
\] % resolved: the PDF glyph is fraktur z; see corrections.csv, leaf 4.

$\Ci$ is a cofibered category over $I$, whose fibers are
\[
\Ci_i \;\simeq M^\pi/X_i \;,
\]
now since $X_i \in M$, we see, because of the $\pi$-accessibility condition, that $M^\pi/X_i$ is \textit{filtering.} We are going to single this out as

\textbf{Lemma 5.10.} \textit{Let $M$ be a $\pi$-accessible subset of $\M$. Then for every $Y \in M$, $M^\pi/Y$ is a category that is} filtering, \textit{and $\Ss_\pi^{M^\pi}(Y)$ is cofinal in it.}
\DERmark{ENV0889}


\msp{page 63}

Let us consider the canonical functor
\[
I \;=\; \Ss_\pi^{M^\pi}(Y) \xrightarrow{\;\varphi\;} \underline{M}^\pi/Y \;=\; I' \;,
\]
\DERmark{ENV0890}

where the source category $I$ is filtering. By SGA 4 I 8.1.3, this functor is cofinal if and only if it satisfies two conditions F\,1, F\,2, and then $I'$ is also filtering. Condition F\,1 is the usual condition that every $i' \in \Ob I'$ is ``dominated'' by a $\varphi(i)$, \textit{i.e.} every $Z \longrightarrow Y$, with $Z \in M^\pi$, is ``dominated'' by a \textit{monomorphism} $Z_\alpha \hookrightarrow Y$, with $Z_\alpha \in M^\pi$, which follows from the fact that it is dominated by $Z_\alpha \hookrightarrow Y$ with $Z_\alpha \in \Filt^\pi$, \textit{i.e.} $Z_\alpha \in \Ss_\pi(Y)$, and that among these $Z_\alpha$ those which are in $M^\pi$ are cofinal (since $Y \in M$). The other condition\footnote[23]{Since $I$ is filtering.} is that a double arrow $i' \overset{f'}{\underset{g'}{\rightrightarrows}} \varphi(i)$ can
\msp{page 64}
be equalized by an arrow $\varphi(h)$, with $h : i \longrightarrow j$. But in fact, for an $i' \in \Ob I'$, \textit{i.e.} a given $Z \longrightarrow Y$ ($Z \in M^\pi$), and an $i \in \Ob I$, \textit{i.e.} a given $Z_\alpha \hookrightarrow Y$ ($Z_\alpha \in M^\pi$), there is at most \textit{one} arrow in $I' = \underline{M}^\pi/Y$ from $Z \longrightarrow Y$ to $Z_\alpha \hookrightarrow Y$. So condition F\,2 is trivially satisfied (we take $h = \mathrm{id}_i$).

\textbf{Corollary 5.11.} \textit{If $(X_i)_{i \in I}$ is a filtering direct system, with the $X_i$ in $M^\pi$, then the category $\Ci = \underline{M}^\pi/\underline{M} \times_{\underline{M}} (I,\varphi)$, constructed on page} 62, \textit{is filtering, and the canonical functor} (\textit{cofibrant}) \textit{$\Ci \longrightarrow I$ is cofinal.}
\DERmark{ENV0891}


That $\Ci$ is filtering and $\Ci \longrightarrow I$ cofinal follows formally from the fact that $I$ is, that the fibers of $\Ci$ are, and that $\Ci$ is cofibered over $I$. (Deserves a lemma, for future reference.)

\textbf{N.B.} We can replace $\Ci$ by $\Ci_0$ whose fibers are the
\[
\Ss_\pi^{M^\pi}(X_i) \;\subset\; \underline{M}^\pi \;; \qquad \text{cofinal in } \Ci_i = \underline{M}^\pi/X_i \;,
\]
\msp{page 65}
we find a subcategory of $\Ci$ cofinal in $\Ci$, and which this time is ordered just like $I$.

\printedpage{38}
% head: A. Grothendieck : Dérivateurs, Ch. XV, Édition des Universités de Montpellier II et Paris VII
% foot: Version 5 novembre 2010
We have
\[
X_i \;=\; \varinjlim_{Z_\alpha \in \Ci_i} Z_\alpha \;,
\]
hence
\[
X \;=\; \varinjlim_{\alpha \in \Ci} Z_\alpha \;,
\]
altogether we have arrived at representing $X$ as a filtering $\varinjlim$ of elements $Z_\alpha$ which are (not only in $M$, but even) in $M^\pi$.

Let us change the notation, taking
\[
X \;=\; \varinjlim_{i \in I} X_i \;, \qquad I \text{ filtering ordered set, } X_i \in M^\pi \;\; \forall\, i \in I \;,
\]
let $I'$ be the set of \textit{filtering} subsets $J \subset I$ such that $\mathrm{card}\, J \leq \pi$, $I'$ ordered by inclusion. Then $I'$ is large relative to $\pi$, and we have a direct system
\[
(X_J)_{J \in I'} \;, \qquad X_J \;=\; \varinjlim_{i \in J} X_i \;.
\]
By the stability hypothesis on $M^\pi$ (which has not yet been used), the $X_J$ are in $M^\pi$. We therefore have
\[
X \;=\; \varinjlim_{J \in I'} X_J \;, \qquad \text{the } X_J \text{ in } M^\pi\text{, } I' \text{ large relative to } \pi \;,
\]
so $X \in M$ since $M$ is $\pi$-admissible, \quad \textit{q.e.d.}

\msp{page 66}

\textbf{Definition 5.12.}
\begin{enumerate}[label={}]
\item[1)] Let $\M$ be a category satisfying conditions $(a)$, $(b)$, $(c)$, $(d)$ (p.~44-46). Let $\pi \geq \pi_0$ be a cardinal. A subset $M$ of $\Ob\M$ is said to be \textit{L-$\pi$-accessible,} if it is $\pi$-accessible and stable under small filtering $\varinjlim$\footnote[24]{We then say that $M$ is \textit{L-stable.}}, or, what amounts to the same thing (5.9), if it is $\pi$-accessible and if $M^\pi = M \cap \Filt^\pi$ is stable under filtering $\varinjlim$ of cardinality $\leq \pi$. (\textbf{N.B.} If the objects of $\Filt^\pi$ are $\pi$-accessible, 5.8 tells us that this also means that it suffices for this that there is a subset $M_0$ of $\Filt^\pi$ such that $M$ is the L-envelope (envelope stable under filtering $\varinjlim$) of $M_0$ -- a condition which is in any case always necessary, by taking $M_0 = M^\pi$.)
\item[2)] Let $M$ be a subset of $\Ob\M$. We say that it is \textit{L-accessible,} if it satisfies one of the following four equivalent conditions:
\begin{enumerate}[label={}]
\item[\textit{a})] There is a cardinal $\pi$ such that $M$ is L-$\pi$-accessible.
\item[\textit{a}$'$)] $M$ is accessible and L-stable.
\item[\textit{b})] There is a small subset $M_0$ of $\Ob\M$, such that $M$ is the L-envelope of $M_0$\footnote[25]{This condition does not depend on the chosen cardinal filtration of $\M$.}.
\item[\textit{b}$'$)] There is a subset $M_0$ of $\M$ contained in a subset $\Filt^\pi$ (or, what amounts to the same thing, such that the full subcategory $\underline{M}_0$ of $\M$ is essentially small, \textit{i.e.} equivalent to a small category, such that $M$ is the \textit{L}-envelope of $M_0$).
\end{enumerate}
\end{enumerate}


\endgroup
% END UNIT Der-ChXV-05-11-10/Der-ChXV-05-11-10_p036-040.tex
% BEGIN UNIT Der-ChXV-05-11-10/Der-ChXV-05-11-10_p041-045.tex
\begingroup\makeatletter%
\newcommand{\printedpage}[1]{\marginpar{\footnotesize p.~#1}}%
\newcommand{\msp}[1]{\marginpar{\footnotesize [#1]}}%
\newcommand{\Filt}{\mathrm{Filt}}%
\newcommand{\Fl}{\mathrm{Fl}}%
\newcommand{\Ob}{\mathrm{Ob}}%
\newcommand{\card}{\operatorname{card}}%
\newcommand{\uHom}{\underline{\mathrm{Hom}}}%
\newcommand{\uFl}{\underline{\mathrm{Fl}}}%
\newcommand{\into}{\lhook\joinrel\longrightarrow}%
\makeatother\let\printedpage\derpp\let\msp\dermsp%
\let\mspin\DERmspin%


\printedpage{39}
% head: A. Grothendieck : Dérivateurs, Ch. XV, Édition des Universités de Montpellier II et Paris VII
% foot: Version 5 novembre 2010

\msp{page 67}
The equivalence of $(a)$ and $(a')$, or of $(b)$ and $(b')$ is more or less tautological. The equivalence of $(a)$ and $(b)$ follows at once from the N.B. included in $1^{\circ}$, taking into account that the cardinals $\pi$ such that the objects of $\Filt^{\pi}$ are $\pi$-accessible are cofinal in the set of cardinals in $\mathfrak{U}$. (\textit{Cf.} 3.2, page 16).

\textbf{Proposition 5.13.} \textit{Every finite union of L-$\pi$-accessible subsets of} $\Ob\mathcal{M}$ \textit{is L-$\pi$-accessible. The intersection of a family $(M_{\alpha})_{\alpha\in A}$ of L-$\pi$-accessible subsets, with} $\card A \leq \pi$\textit{, is L-$\pi$-accessible.}

\textsc{Proof}. The first assertion comes from the fact that a \textit{finite} union of L-stable subsets of $M$ is L-stable, and that we have likewise for the notion of $\pi$-accessibility. The two things are moreover proved practically in the same way, using the

\textbf{Lemma 5.14.} \textit{Let $I$ be a filtering ordered set, $(I_{\alpha})_{\alpha\in A}$ a} finite \textit{nonempty family of subsets of $I$, such that $\bigcup I_{\alpha} = I$. Then there exists $\alpha \in A$ such that $I_{\alpha}$ is cofinal.}

\msp{page 68}
\textsc{Proof} by contradiction. If for no $\alpha$, $I_{\alpha}$ were cofinal, we would conclude from this that for every $\alpha$ there exists an $i_{\alpha}$ such that $I/i_{\alpha} \cap I_{\alpha} = \emptyset$. Let $i$ be an upper bound for the $i_{\alpha}$, we will thus have $I/i \cap I_{\alpha} = \emptyset \ \ \forall\, \alpha \in A$, contrary to the hypothesis that $\bigcup I_{\alpha} = I$.

The case of an \textit{intersection} of L-$\pi$-accessible subsets: it is of course L-stable, so it remains to prove that it is $\pi$-accessible, \textit{i.e.} satisfies condition $\mathrm{b}_{\pi}$ of \textit{loc. cit.}

\msp{page 69}
Let then
\[
X \;=\; \varinjlim_{i\in I} X_{i} \; , \qquad\quad X_{i} \in \Filt^{\pi} \; ,
\]
be a $\pi$-adapted representation of $X$, with $I$ a $\pi$-adapted ordered set. Thus we have
\[
\begin{aligned}
X \in M_{\alpha} \quad &\Longleftrightarrow \quad \text{the set } I_{M_{\alpha}} \text{ of } i \in I \text{ such that } X_{i} \in M_{\alpha}\\
&\phantom{\Longleftrightarrow} \quad \text{is cofinal in } I.
\end{aligned}
\tag*{$(*)$}
\]
Moreover, since $M_{\alpha}^{\pi} = M_{\alpha} \cap \Filt^{\pi}$ is stable under filtering $\varinjlim$ of cardinality $\leq \pi$, we see at once that the subset $I_{M_{\alpha}}$ of $I$ is stable under suprema in $I$ of filtering subsets $J$ of $I_{M_{\alpha}}$ of cardinality $\leq \pi$. Thus $I_{M_{\alpha}}$ cofinal is equivalent to $I_{M_{\alpha}}$ $\pi$-\textit{cofinal} (\textit{cf.} 5.8, page 59). Let $M = \bigcap_{\alpha\in A} M_{\alpha}$, I want to prove that
\[
X \in M \quad \Longleftrightarrow \quad \text{the set } I_{M} \ (= \textstyle\bigcap I_{M_{\alpha}}) \text{ is cofinal in } I.
\tag*{$(**)$}
\]
Now $(*)$ implies that $X \in M$ $(= \bigcap M_{\alpha})$ is equivalent to $I_{M_{\alpha}}$ $\pi$-cofinal for every $\alpha$. So all in all we want to prove
\[
I_{M_{\alpha}} \text{ cofinal in } I \text{ for every } \alpha \quad\ \Longleftrightarrow \quad\ \textstyle\bigcap I_{M_{\alpha}} \text{ is cofinal in } I.
\tag*{$(***)$}
\]

\msp{page 70}
We of course have $\Longleftarrow$, so it remains to prove $\Longrightarrow$. If the index set $A$ is finite, this follows from Lemma 5.8, taking into account that the cofinality of the $I_{M_{\alpha}}$ implies their $\pi$-cofinality. We must therefore simply improve 5.8, by the

\printedpage{40}
% head: A. Grothendieck : Dérivateurs, Ch. XV, Édition des Universités de Montpellier II et Paris VII
% foot: Version 5 novembre 2010

\textbf{Lemma 5.15.} \textit{Let $\pi$ be a cardinal, $I$ a $\pi$-adapted filtering ordered set, and let $(I_{\alpha})_{\alpha\in A}$ be a family of subsets $I_{\alpha}$ that are $\pi$-cofinal} (cf. 5.8), \textit{such that} $\card A \leq \pi$. \textit{Then $\bigcap I_{\alpha}$ is $\pi$-cofinal.}

It suffices to prove that it is cofinal, for the property of stability under suprema of filtering subsets of cardinality $\leq \pi$ is obvious. Replacing $I$ by $I/i_{0}$, and the $I_{\alpha}$ by their traces on $I/i_{0}$ if need be, the question is simply to prove that $\bigcap_{\alpha} I_{\alpha} \neq \emptyset$. We are going to construct an element $s$ of $\bigcap I_{\alpha}$ by transfinite induction. Let us take a well-order on $A$, and prove by transfinite induction
\msp{page 71}
that for every $\alpha \in A$, the set
\[
J_{\alpha} \;=\; \bigcap_{\beta<\alpha} I_{\beta}
\]
is cofinal. If we know that $J_{\alpha}$ is cofinal, the same holds for $J_{\alpha+1} = J_{\alpha} \cap I_{\alpha}$, by 5.8. So we are reduced to proving that $J_{\alpha}$ is cofinal if $\alpha$ is a limit ordinal, thus to proving Lemma 5.15 for a \textit{decreasing} transfinite sequence of $\pi$-cofinal subsets $J_{\beta}$ $(\beta < \alpha)$ when $\card \alpha \leq \pi$, such that $J_{\beta} = \bigcap_{\beta'<\beta} J_{\beta'}$ when $\beta$ is a limit. To prove that the intersection is $\pi$-cofinal -- and we are reduced to proving that it is nonempty. We are going to construct by transfinite induction an \textit{increasing} sequence of elements
\[
i_{\beta} \;\in\; J_{\beta} \qquad\quad (\, \beta \;\leq\; \alpha \,) \; .
\]
If $i_{\beta}$ is constructed, since $J_{\beta+1}$ is cofinal, there exists $i_{\beta+1} \geq i_{\beta}$ in $J_{\beta+1}$, O.K. If $\beta$ is a limit ordinal, we take
\[
i_{\beta} \;=\; \sup_{\beta'<\beta} i_{\beta'} \; ,
\]
$i_{\beta}$ is in each $J_{\beta_{0}}$ for $\beta_{0} < \beta$,
\DERmark{ENV0892}

\msp{page 72}
since the $i_{\beta'}$ for $\beta_{0} \leq \beta' < \beta$ are all in $J_{\beta_{0}} \supset J_{\beta'}$, hence also their limit, hence
\[
i_{\beta} \;\in\; \bigcap_{\beta_{0}<\beta} J_{\beta_{0}} \;=\; J_{\beta} \; .
\]
So transfinite induction works, and thus gives us an element $s_{\alpha}$ of $J_{\alpha}$\,!

I am also going to prove in due form the assertion for $\bigcup_{\alpha\in A} M_{\alpha}$, when $A$ is finite. Setting this time $M = \bigcup_{\alpha} M_{\alpha}$, we must prove, using the relations $(*)$ (p. 69), that we have
\DERmark{ENV0893}

\[
\begin{aligned}
&X \in M &&\Longleftrightarrow \quad &&\text{the set } \ I_{M} \ \ (= \ \textstyle\bigcup_{\alpha} I_{M_{\alpha}})\\
&(\textit{i.e.} \ \exists \ \alpha \text{ with } X \in M_{\alpha}) && &&\text{is cofinal in } I \ ,
\end{aligned}
\tag*{$(**)$}
\]
and using $(*)$ the first term of this equivalence to be established means that there exists $\alpha$ with $I_{M_{\alpha}}$ cofinal in $I$. Of course, if one of the $I_{M_{\alpha}}$ is cofinal in $I$, a fortiori their union. But conversely, if the union is cofinal, it is filtering and a union of the subsets $I_{M_{\alpha}}$, one of which must therefore be cofinal in it by Lemma 5.14, hence also in $I$, \textit{q.e.d.}
\DERmark{ENV0894}


\msp{page 73}
\textbf{Corollary 5.16} (of Proposition 5.13). \textit{Every finite union of L-accessible subsets of $\mathcal{M}$ is L-accessible. Every intersection $\bigcap M_{\alpha}$ of a small family of L-accessible subsets of $\mathcal{M}$, is L-accessible.}

\printedpage{41}
% head: A. Grothendieck : Dérivateurs, Ch. XV, Édition des Universités de Montpellier II et Paris VII
% foot: Version 5 novembre 2010

\section*{6\quad Abortive attempt at a generation criterion for L-accessible subsets of $\Fl(\mathcal{M})$.}

Let $\mathcal{M}$ be a category equipped with a cardinal filtration $(\Filt^{\pi})_{\pi\geq\pi_{0}}$, and let $S$ be an object of Cat, let us consider
\[
\mathcal{M}(S) \;=\; \uHom(S,\mathcal{M}) \; ,
\]
and equip $\mathcal{M}(S)$ with the filtration by the
\[
\mathcal{M}^{\pi}(S) \;=\; \uHom(S,\mathcal{M}^{\pi} = \underline{\Filt}^{\pi}(\mathcal{M})) \; .
\]
I say that this is a cardinal filtration of $\mathcal{M}(S)$. We must check conditions $(a)$, $(b)$, $(c)$ of SGA 4 I 9.12. Conditions $(a)$ ($\Filt^{\pi}$ essentially small) and $(b)$ (condition $\mathrm{L}_{\pi_{0}}$, and stability of $\Filt^{\pi}$ under $\varinjlim\limits_{I}$ , $I$ filtering large relative to $\pi_{0}$ and of cardinality $\leq \pi$) pass at once from $\mathcal{M}$ to $\mathcal{M}(S)$. Condition $(c)$ is more delicate (representation
\msp{page 74}
of an object as a filtering $\varinjlim$ large relative to $\pi$ of objects of $\Filt^{\pi}$). I am not going to redo the work, done in a much more general and more delicate case in SGA 4 I §9, namely in Corollary 9.24 of Theorem 9.22 (p. 177). The only restriction is that there the objects of $\mathcal{M}$ are assumed accessible, and denoting by $\pi_{1}$ the smallest cardinal such that $\pi_{1} \geq \pi_{0}$ and that the objects of $\Filt^{\pi_{0}}$ are $\pi_{1}$-accessible, we restrict ourselves for the cardinal filtration of $\uHom(S,\mathcal{M})$ to the cardinals \fbox{$\pi \geq c = 2^{c_{0}}$,} where
\[
c_{0} \;=\; \sup(\pi_{1}, \card \Fl(S)) \; .
\]
This is a reasonable and innocuous restriction.

For the moment, I am interested above all in the case where $S = \boldsymbol{\Delta}_{1}$, hence in
\[
\uHom(\boldsymbol{\Delta}_{1},\mathcal{M}) \;=\; \uFl(\mathcal{M}) \; ,
\]
in which case we have $c_{0} = \pi_{1}$, so it suffices to restrict ourselves to the cardinals $\pi \geq 2^{\pi_{1}}$ to have a good cardinal filtration on $\uFl(\mathcal{M})$.

It is thus the set of the two conditions $(b)$, $(c)$ on $\mathcal{M}$ (p. 44) that is stable under passage to the $\uHom(S,\mathcal{M}) = \mathcal{M}(S)$. The same holds for condition $(a)$ (stability under small direct limits), and also for condition\msp{page 75}
$(d)$ (p. 45). Likewise the exactness condition $(e)$ (p. 46) on filtering $\varinjlim$ (they must transform monomorphisms into monomorphisms).

To be on the safe side, I will make on $\mathcal{M}$ the assumptions $(a)$ to $(e)$ inclusive (although $(e)$ seems a little artificial to me, we should be able to eliminate it, with a second wind\,!). These conditions thus have the desirable property of stability under passage to the $\mathcal{M}(S)$, and in particular under passage to $\Fl(\mathcal{M})$. Increasing $\pi_{0}$ if need be, we can therefore assume that the $\Fl^{\pi}(\mathcal{M}) = \mathrm{Hom}(\boldsymbol{\Delta}_{1},\mathcal{M}^{\pi})$ $(\pi \geq \pi_{0})$ also form a cardinal filtration of $\uFl(\mathcal{M})$.

We thus have at our disposal the notions of $\pi$-accessible, accessible, L-$\pi$-accessible, L-accessible subsets of $\Fl(\mathcal{M}) \ = \ \Ob\uFl(\mathcal{M})$. These notions are stable under finite unions and intersections, and the absolute notions (without $\pi$) stable under ``arbitrary'' intersections (\textit{i.e.} for \textit{small} families of subsets). This will be something particularly\printedpage{42}%
% head: A. Grothendieck : Dérivateurs, Ch. XV, Édition des Universités de Montpellier II et Paris VII
% foot: Version 5 novembre 2010
\msp{page 76}
crucial, the case of an intersection $C \cap W$\,!

The main result, for which I practiced the scales of the preceding section, is the following:

\setcounter{footnote}{25}
\textbf{Theorem 6.1}\footnote{Botched proof, statement perhaps \textit{false}\,!!}. \textit{Let $\mathcal{M}$ be a category satisfying the stated conditions:}
\begin{enumerate}[label={}]
\item[a)] \textit{Stability under small} $\varinjlim$ \textit{(not necessarily filtering).}
\item[b)] $\mathcal{M}$ \textit{equipped with a cardinal filtration} $(\Filt^{\pi})_{\pi\geq\pi_{0}}$.
\item[c)] \textit{The objects of $\mathcal{M}$ are accessible.}
\item[d)] \textit{The} $\Filt^{\pi}$ \textit{are stable under passage to subobjects and under} $\varinjlim$ \textit{indexed by categories $I$ such that} $\card \Fl(I) \leq \pi$\textit{, and if $Z \in \Filt^{\pi}$, every quotient $Z'$ of $Z$ (such that $Z \longrightarrow Z'$ does not factor through a subobject $Z'_{0}$ of $Z'$, different from $Z'$) is in} $\Filt^{\pi}$.
\item[e)] \textit{Filtering direct limits in $\mathcal{M}$ transform monomorphisms into monomorphisms.}
\end{enumerate}
\DERmark{ENV0895}

\textit{Let $C$ be an L-accessible subset of} $\Fl(\mathcal{M}) = \Ob\uFl(\mathcal{M})$, i.e. (5.12) \textit{a subset stable under} $\varinjlim$ \textit{that are filtering, and accessible. Suppose moreover that $C$ is made up of monomorphisms}\footnote{In the end, we get rid of this assumption.}\textit{, and that $C$ contains the} isomorphisms. \textit{Then there exists a small subset $C_{0} \subset C$, such that $C \subset \widetilde{C_{0}}$. More precisely, if $\pi$ is a cardinal such that $C$ is $\pi$-accessible, and if we set $C^{\pi} = C \cap \Fl^{\pi}(\mathcal{M})$, then we have}
\msp{page 77}
\[
C \;\subset\; \widetilde{C^{\pi}} \;=\; \widetilde{C_{0}} \; ,
\tag*{(6.1.1)}
\]
\textit{where $C_{0}$ is any small subset of $C^{\pi}$ such that every object of $C^{\pi}$ is isomorphic to an object of $C_{0}$.}

\textbf{Corollary 6.2.} \textit{We have $C_{*} = C_{0\,*}$, $\widetilde{C} = \widetilde{C_{0}}$.}

\textbf{Corollary 6.3.} \textit{Equivalent conditions on the L-accessible subset $C$ of} $\Fl(\mathcal{M})_{\mathrm{mono}}$\textit{, containing the isomorphisms:}
\begin{enumerate}[label={}]
\item[a)] \textit{$C = \widetilde{C}$.}
\item[b)] \textit{$C$ is stable under cobase change and under composition.}
\item[b$'$)] \textit{$C$ is stable under cobase change, and under simple amalgamated sums,} i.e. \textit{if $A \longrightarrow B_{1}$, $A \longrightarrow B_{2}$ are in $C$, then $A \longrightarrow B_{1} \amalg_{A} B_{2}$ as well.}
\item[c)] \textit{$C^{\pi}$ is stable in} $\Fl^{\pi}(\mathcal{M}) = \Fl(\mathcal{M}^{\pi})$ \textit{under cobase change and under composition.}
\item[c$'$)] \textit{$C^{\pi}$ is stable in} $\Fl(\mathcal{M}^{\pi})$ \textit{under cobase change and under amalgamated sums} (cf. b$'$).
\end{enumerate}
\textit{When these conditions are satisfied, then $(C,C_{*})$ is a Quillen couple, which falls under Theorem} 1 \textit{of section} 1 (p. 10).

The last assertion follows from Theorem 6.1 and from the corollary to Theorem 1 of section 1. We have
\printedpage{43}%
% head: A. Grothendieck : Dérivateurs, Ch. XV, Édition des Universités de Montpellier II et Paris VII
% foot: Version 5 novembre 2010
\msp{page 78}
on the other hand tautologically
\[
\begin{tikzcd}
(a) \arrow[r, Rightarrow] & (b) \arrow[r, Rightarrow] \arrow[d, Rightarrow] & (b') \arrow[d, Rightarrow] \\
 & (c) \arrow[r, Rightarrow] & (c') \ ,
\end{tikzcd}
\]
and it remains to prove $(c') \Longrightarrow (a)$, which we will do below.

Let us then prove 6.1. In (6.1.1) the equality $\widetilde{C_{0}} = \widetilde{C^{\pi}}$ is tautological. So everything comes down to proving
\[
C \;\subset\; \widetilde{C^{\pi}} \; .
\]
One difficulty here is that we do \textit{not} know that $\widetilde{C^{\pi}}$ is stable under filtering $\varinjlim$ -- we know it only for $C$. Let then
\[
A \;\into\; B
\]
be an arrow in $C$, let us prove that it is in $\widetilde{C^{\pi}}$. For this, we are going to construct, by transfinite induction, a \textit{strictly} increasing transfinite sequence of subobjects $A_{\alpha}$ of $B$ containing $A = A_{0}$,
\[
\underbrace{A}_{=\,A_{0}} \;\; \overset{i_{\alpha}}{\into} \;\; A_{\alpha} \;\; \overset{i_{\beta,\alpha}}{\into} \;\; A_{\beta} \xrightarrow{\ j_{\beta}\ } B \; ,
\]
such that we have
\[
i_{\alpha} \;\in\; \widetilde{C^{\pi}} \; , \qquad\quad \text{more generally the } i_{\beta,\alpha} \;\in\; \widetilde{C^{\pi}} \; , \qquad\quad j_{\alpha} \;\in\; C \; ,
\]
and that we have
\[
A_{\alpha} \;=\; \varinjlim_{\beta<\alpha} A_{\beta} \qquad\quad \text{if } \alpha \text{ limit ordinal} \; .
\]

\msp{page 79}
This already specifies the recursive definition for $\alpha = \alpha_{0}$ a limit ordinal. Let us note that the $A_{\alpha_{0}}$ constructed in terms of the $A_{\beta}$ $(\beta < \alpha_{0})$ does satisfy the conditions that
\begin{enumerate}[label={}]
\item[\textit{a})] the $i_{\alpha_{0},\beta} : A_{\beta} \longrightarrow A_{\alpha_{0}}$ are in $\widetilde{C^{\pi}}$, for $\widetilde{C^{\pi}}$ is stable under ``strict ordinal direct limits'' (\textit{cf.} 1.19, p. 8),
\item[\textit{b})] $j_{\alpha_{0}} \, : \, A_{\alpha_{0}} \longrightarrow B$ is in $C$, for it is the filtering direct limit of the system of the $A_{\beta} \longrightarrow B$ $(\beta < \alpha_{0})$ which are in $C$, and $C$ is stable under these limits.
\end{enumerate}
It remains then to specify how we pass from $A_{\alpha}$ to $A_{\alpha+1}$. This comes down to proving the

\textbf{Lemma 6.4.} \textit{Let $f : A \into B$ be in $C$, such that $f$ is not an isomorphism. Then there exists a subobject $A_{1}$ of $B$,} distinct from $A$, \textit{such that $A \longrightarrow A_{1}$ is in $\widetilde{C^{\pi}}$, and $A_{1} \longrightarrow B$ in $C$.}

If this lemma is admitted, this means that the transfinite construction continues as long as we find $A_{\alpha}$ such that $A_{\alpha} \neq B$, \textit{i.e.} $A_{\alpha} \into B$ is not an isomorphism. But since the $A_{\alpha}$ are all different, if $\mu$ is the cardinality of the set of subobjects of $B$ (we know that it is a
\DERmark{ENV0896}



\endgroup
% END UNIT Der-ChXV-05-11-10/Der-ChXV-05-11-10_p041-045.tex
% BEGIN UNIT Der-ChXV-05-11-10/Der-ChXV-05-11-10_p046-050.tex
\begingroup\makeatletter%
\newcommand{\printedpage}[1]{\marginpar{\footnotesize p.~#1}}%
\newcommand{\msp}[1]{\marginpar{\footnotesize [#1]}}%
\newcommand{\M}{\mathcal{M}}%
\newcommand{\cercle}[1]{\tikz[baseline=(X.base)]\node[draw,circle,inner sep=0.6pt](X){#1};}%
\newcommand{\Filt}{\mathrm{Filt}}%
\newcommand{\Mon}{\mathrm{Mon}}%
\newcommand{\Ss}{\mathrm{Ss}}%
\newcommand{\Fl}{\mathrm{Fl}}%
\newcommand{\Ob}{\mathrm{Ob}}%
\newcommand{\Cof}{\mathrm{Cof}}%
\makeatother\let\printedpage\derpp\let\msp\dermsp%
\let\mspin\DERmspin%


\printedpage{44}
% head: A. Grothendieck : Dérivateurs, Ch. XV, Édition des Universités de Montpellier II et Paris VII

\msp{page 80}
small set, because of the hypotheses made on $\M$), the ordinals $\alpha$ for which $A_\alpha$ exists satisfy $\mathrm{card}\,\alpha \leq \mu$. So there must exist an $\alpha$ such that $A_\alpha = B$. But then $i_\alpha : A \longrightarrow A_\alpha$ is equal to $f$, hence $f \in \widetilde{C^{\pi}}$, \quad \textit{q.e.d.}
\DERmark{ENV0897}


It remains then to prove 6.4. For this, let $I$ be the set of subobjects of $B$ which are in $\Filt^{\pi}$, ordered by inclusion. (If $B$ is in $\Filt^{\pi}$, $A$ too, so $f$ is already in $C^{\pi} \subset \widetilde{C^{\pi}}$, and there is nothing to prove.) It is a filtering set large relative to $\pi$ (and even $\pi$-\textit{adapted}), whose $\varinjlim$ is $B$ (5.3, p. 46). Consider the
\[
A_i \;=\; B_i \cap A \;,
\]
they are in $\Filt^{\pi}$, as subobjects of objects of $\Filt^{\pi}$ (condition $(d)$ on the cardinal filtration). Thus, the arrow $f$ appears as a filtering direct limit of the arrows
\[
f_i \;:\; A_i \longrightarrow B_i \;.
\]
It is trivial that the system of the $A_i$ is cofinal in $\Ss_{\pi}(A)$, hence $A = \varinjlim A_i$, whence $f = \varinjlim f_i$.

\msp{page 81}
On the other hand, I say that the direct system of the $f_i$ in $\underline{\Fl}(\M)$ is \textit{well adapted to $\pi$.} We already know that $I$ is well adapted to $\pi$, the sup in $I$ of a filtering subset $J$ of cardinal $\leq \pi$ being the sup in the ordered set $\Ss(B)$ (which sup is in $\Ss_{\pi}(B)$). It is essential, for our argument, to know that if $i_J = \sup J$, then not only $B_{i_J} = \varinjlim\limits_{i\in J} B_i$, but \marginpar{$\|$}% marginal double vertical bar printed in the left margin at this line
also $A_{i_J} = \varinjlim\limits_{i\in J} A_i$, so that the direct system of the $f_i$ is $\pi$-adapted in $\underline{\Fl}(\M)$. It seems that we therefore need an exactness condition on $\M$ close to $(e)$, namely:

\begin{itemize}
\item[\cercle{e$'$}] The filtering $\varinjlim$ functors commute with base change.
\end{itemize}

(It suffices here in the very particular case of a base change by a monomorphism $A \longrightarrow B$, and of a filtering direct limit of subobjects of $B$ with inclusion morphisms as transitions.) Let us note that \cercle{e}, \cercle{e$'$} generally arise as a consequence of the condition % [sic] Il le suffit -> Il suffit, ungrammatical as printed
\DERmark{ENV0898}


\begin{itemize}
\item[(e$_0$)] The filtering $\varinjlim$ in $\M$ are left exact.
\end{itemize}
\DERmark{ENV0899}


\msp{page 82}
With this laid down, we thus see, $C$ being $\pi$-accessible and $I$ large relative to $\pi$, that the set $I_C$ of $i \in I$ such that $f_i \in C$ is \textit{cofinal} in $I$. But since $f_i \in\Fl^\pi(\mathcal M)$, $f_i \in C$ is equivalent to $f_i \in C^{\pi}$. Since $A \longrightarrow B$ is not an isomorphism, there exists an $i_0$ such that $B_{i_0} \not\subset A$, then an $i \geq i_0$ such that $f_i \in C^{\pi}$, and also $B_i \not\subset A$ (since $B_i \supset B_{i_0}$). Thus we have proved this:
\DERmark{ENV0900}


\begin{itemize}
\item[($*$)] Under the hypotheses of (6.4) and with the hypothesis \cercle{e$'$} on $\M$, there exists a subobject $B'$ of $B$ satisfying the following conditions.
\begin{itemize}
\item[\textit{a})] $B' \in \Filt^{\pi}$.
\item[\textit{b})] $B' \not\subset A$.
\item[\textit{c})] Setting $A' = B' \cap A$, the inclusion
\[
i' \;:\; A' \;\longrightarrow\; B'
\]
is in $C^{\pi}$.
\end{itemize}

\printedpage{45}
% head: A. Grothendieck : Dérivateurs, Ch. XV, Édition des Universités de Montpellier II et Paris VII

Moreover, the set $J \subset \Ss_{\pi}(B)$ formed by the $B'$ in question is cofinal in $\Ss_{\pi}(B)$, and even $\pi$-cofinal.
\DERmark{ENV0901}

\end{itemize}

We still have to prove this last point -- that $J$ (which we know is cofinal) is even $\pi$-cofinal, \textit{i.e.} stable under sup (in $\Ss_{\pi}(B)$) of filtering subsets $J' \subset J$ such that $\mathrm{card}\, J' \leq \pi$. But this follows at once from the fact that $C$ being stable under

\msp{page 83}
filtering direct limits, $C^{\pi} = C \cap \Filt^{\pi}$ is stable under filtering $\varinjlim$ of cardinal $\leq \pi$, by virtue of condition $(d)$ on $\Filt^{\pi}$. Consider now the cocartesian square
\[
\begin{tikzcd}
A' \arrow[r, hook, "i'"] \arrow[d, hook] & B' \arrow[d] & & \\
A \arrow[r, "i"] & C' = A \amalg_{A'} B' \arrow[r, "j"] & B \;.
\end{tikzcd}
\]
Since $i' \in C^{\pi}$, we have $i \in \widetilde{C^{\pi}}$ (for $\widetilde{C^{\pi}}$ stable under cobase change). On the other hand, I need the fact that $j$ is in $C$, and I do not even know that it is a monomorphism! So we have to go back to the indexed notation
\[
\begin{tikzcd}
A_i \arrow[r, hook, "\alpha_i"] \arrow[d, hook] & B_i \arrow[d] & \\
A \arrow[r, hook, "\beta_i"] & C_i \arrow[r, "\gamma_i"] & B \;.
\end{tikzcd}
\]
Since the formation of amalgamated sums commutes with filtering direct limits, we see that $\varinjlim C_i \simeq A \amalg_{\varinjlim A_i} \varinjlim B_i = A \amalg_A B = B$, hence the direct limit of the $\gamma_i$ is an isomorphism (essentially $\mathrm{id}_B$). On the other hand, the same argument shows that the direct system of the $C_i$ $(i \in J)$ in $\M$ is $\pi$-adapted, hence also that of the $\gamma_i$\footnote[28]{\textit{Except} that we do not know that the $C_i$ are in $\Filt^{\pi}$ -- in fact, they will not be in general, since we do not assume that $A, B \in \Filt^{\pi}$.}
\msp{page 84}
in $\underline{\Fl}(\M)$. Since $\varinjlim \gamma_i$ is an isomorphism, it is in $C$. Since $C$ is $\pi$-accessible, it follows that the set $J'$ formed by the $i \in J$ such that $\gamma_i \in C$ is cofinal as well\footnote[29]{\textit{Botched, cf.} preceding annotation.}. So we find $i$ with $\alpha_i \in C^{\pi}$, $\gamma_i \in C$, which establishes lemma 6.4, and completes the proof of theorem 6.1.

\section*{7\quad Construction of closed Quillen triples with $\boldsymbol{C} = \mathbf{Mon} \subset \mathbf{Cof}_{\boldsymbol{W}}$}

Finally, all the beautiful developments of paragraphs 5, 6 do not give me the desired ``sweeping result''. Perhaps it is false -- perhaps the property of generation
\printedpage{46}
% head: A. Grothendieck : Dérivateurs, Ch. XV, Édition des Universités de Montpellier II et Paris VII
``$^{\sim}$'' by a small subset (for accessible, or L-accessible sets) is satisfied much more rarely, than I would have imagined. To do things properly, I would have to work through to the % resolved: the quoted raised wave denotes Q-closure; see corrections.csv, leaf 13.
\msp{page 85}
end a few examples, e.g. that of $\Cof_W$, and of $W^{\mathrm{univ}}$ in Cat, for a more or less general basic localizer, and especially in the case $W = \mathbf{W}_{\infty}$. Surely $\Cof_W$ and $W^{\mathrm{univ}}$ are accessible (I ought to check it nevertheless), and if the hypothetical ``theorem'' of paragraph 6 is false for them, at least I will know for sure. But there still remains the question: if $C = \widetilde{C}$ is stable under filtering $\varinjlim$, and of the form $\widetilde{C_0}$, $C_0$ small, and if $W$ is L-accessible as well, then is $C \cap W$ also of the form $\widetilde{TC_0}$, with $TC_0$ small? I have just tried once again to prove it, assuming moreover $C \subset$ mono, and see no way of getting there. So a counterexample would be needed to dispel these perplexities! And the best would be in Cat, while we're at it.

There nevertheless remains afloat a tangible positive result, which I will now try to pin down properly and prove:

\msp{page 86}
\textbf{Theorem 7.1.} \textit{Let $\M$ be a category stable under small filtering $\varinjlim$, $M_0 \subset \Ob\,\M$ a subset of} $\Ob\,\M$ \textit{generating by strict monomorphisms, more precisely we assume this:}

\begin{quote}
\itshape
For every monomorphism $i : A \hookrightarrow B$ in $\M$ which is not an isomorphism, there exists a monomorphism $B' \overset{j}{\hookrightarrow} B$ having the following properties:
\begin{itemize}
\item[\textup{a)}] $B' \in M_0$.
\item[\textup{b)}] $j$ does not factor through $i$.
\item[\textup{c)}] $A' = A \times_B B'$ exists in $\M$, as does $C = A \amalg_{A'} B'$, $A'$ is in $M_0$, and in the canonical diagram
\[
\begin{tikzcd}
A' \arrow[r, hook, "i'"] \arrow[d, hook] & B' \arrow[d, hook] \arrow[dr, hook, dotted, "j"] & \\
A \arrow[r, hook, "\alpha"] \arrow[rr, hook, dotted, bend right=20, "i"'] & C \arrow[r, "\beta"] & B
\end{tikzcd}
\tag*{(7.1.1)}
\]
the arrow $\beta$ is a monomorphism.
\end{itemize}
\end{quote}

(\textbf{N.B.} These conditions are satisfied if we assume, in addition to the condition that $M_0$ be generating for monomorphisms, that $M_0$ contains, along with every object $B'$, and for every object $A'$ which is either a quotient, or a subobject of $B'$, an object $A'_0$ isomorphic to $A'$; and if moreover $\M$ satisfies the following exactness conditions:
\begin{itemize}
\item[1$^{\circ}$)] Every arrow $f$ of $\M$ factors as $ip$, with $i$ a monomorphism, $p$ an epimorphism, and:
\item[2$^{\circ}$)] If $A_1$, $A_2$ are two subobjects of an object $B$ of $\M$, then the intersection $A_0 = A_1 \times_B A_2$ exists, as does the
\printedpage{47}
% head: A. Grothendieck : Dérivateurs, Ch. XV, Édition des Universités de Montpellier II et Paris VII
\msp{page 87}
amalgamated sum $C = A_1 \amalg_{A_0} A_2$, and the canonical morphism $C \longrightarrow B$ is a monomorphism.)
\end{itemize}

\textit{Assume moreover that $\M$ satisfies condition} \cercle{e}, p. 76, i.e. \textit{that every filtering $\varinjlim$ of monomorphisms is a monomorphism.}

1$^{\circ}$) \textit{Let} $\Mon \subset \Fl(\M)$ \textit{be the set of monomorphisms, and} $\Mon_0$ \textit{the subset of} $\Mon$ \textit{formed by the $i \in$} $\Mon$ \textit{such that source and target belong to $M_0$. We then have}
\[
\Mon \;\subset\; \widetilde{\Mon_0} \quad , \qquad \text{\textit{whence}} \quad
\left\{
\begin{aligned}
\widetilde{\Mon} \quad &= \quad \widetilde{\Mon_0} \\
(\Mon)_{*} \quad &= \quad (\Mon_0)_{*} \;.\text{\footnotemark[30]}
\end{aligned}
\right.
\tag*{(7.1.2)}
\]
\footnotetext[30]{\textbf{N.B.} Here we have a criterion for injectivity of objects in the categories $\M/Y$ -- it suffices to test the injectivity condition by monomorphisms $i$ which are in $\Mon_0$.}

2$^{\circ}$) \textit{If $M_0$ is} small \textit{and if the objects of $M_0$ are accessible} (e.g. $\M$ satisfies \cercle{c}, p. 76), \textit{and if} $\Mon$ \textit{is stable under cobase change, then there exists a functorial factorization of every arrow $f$ of $\M$ as}
\[
f \;=\; p\, i \;,
\]
\textit{with $i \in$} $\Mon$, \textit{$p \in$} $(\Mon)_{*}$. (\textbf{N.B.} An arrow $p : X \longrightarrow Y$ of $\M$ is in $(\Mon)_{*}$ if and only if it makes $X$ an \textit{injective} object of $\M/Y$ -- thus the factorization assertion ensures that in the induced categories $\M/Y$ ``there are enough injectives''.) \textit{Moreover, we have}
\[
\Mon \;=\; \widetilde{\Mon_0} \quad .
\tag*{(7.1.3)}
\]

3$^{\circ}$) \textit{Let $W$ be a subset of} $\Fl(\M)$ \textit{which is} L-accessible (i.e. \textit{accessible, and stable under filtering direct limits}). (\textit{We now assume that $\M$ satisfies the conditions} (a), (b), (c), (d), (e) \textit{on page 76, whence a good theory of L-accessible subsets of} $\Ob\,(\M)$ \textit{or of} $\Fl(\M)$.) \textit{We assume, more precisely, that $W$ satisfies the conditions:}

\msp{page 88}
\[
\left\{
\begin{array}{@{\quad}l@{\quad}p{0.6\textwidth}}
\text{a)} & \textit{$W$ is L-accessible} (\textit{i.e.} stable under filtering $\varinjlim$, and accessible). \\[1.5ex]
\text{b)} & \textit{Let $\alpha$, $\beta$ be composable, with $\beta\alpha$ and $\alpha$ in $W$, then $\beta \in W$.} (\textbf{N.B.} This condition will be satisfied if $W$ is a mildly saturated localizer, and it is to this situation that we are going to apply our statement.) \\[1.5ex]
\text{c)} & $\Mon \cap W$ \textit{is stable under cobase change.}
\end{array}
\right.
\tag*{(7.1.4)}
\]
\textit{and that $\M$ satisfies the additional conditions}
\[
\left\{
\begin{array}{@{\quad}l@{\quad}p{0.6\textwidth}}
\text{e$'$)} & \textit{The filtering $\varinjlim$ commute with base change.} \\[1.5ex]
\text{f)} & \textit{The exactness condition} (7.2.1) \textit{below.}
\end{array}
\right.
\tag*{(7.1.5)}
\]
\textit{Let $\pi$ be a cardinal such that $W$ is $\pi$-accessible, and $M_0 \subset$} $\Filt^{\pi}$ (\textit{hence $M_0$ essentially small}). \textit{Then we have}
\[
W \cap \Mon \;\subset\; \widetilde{(W \cap \Mon)^{\pi}}
\tag*{(7.1.6)}
\]
\printedpage{48}
% head: A. Grothendieck : Dérivateurs, Ch. XV, Édition des Universités de Montpellier II et Paris VII
(\textit{where for every subset $C$ of} $\Fl(\M)$, \textit{we set} $C^{\pi} = C \cap \Fl^{\pi}(\M)$), \textit{hence}
\[
\widetilde{W \cap \Mon} \;=\; \widetilde{(W \cap \Mon)^{\pi}} \;,
\tag*{(7.1.7)}
\]
\textit{hence $\widetilde{W \cap \Mon}$ is Q-generated by a} small \textit{subset of} $W \cap \Mon$.

\textbf{Corollary 7.2.} \textit{Let $\M$ satisfy} (a), (b), (c), (d), (e) \textit{and the additional conditions:}
\[
\left\{
\begin{array}{@{\quad}l@{\quad}p{0.6\textwidth}}
\text{e$'$)} & \textit{The filtering $\varinjlim$ commute with base change.} \\[1.5ex]
\text{f)} & \textit{If $A_1$, $A_2$ are two subobjects of an object $B$ of $\M$, and $A_0 = A_1 \cap A_2$, then $A_1 \amalg_{A_0} A_2 \hookrightarrow B$ is a monomorphism}\,\footnotemark[31].
\end{array}
\right.
\tag*{(7.2.1)}
\]
\footnotetext[31]{Another way of saying it: if $C = \sup_{\Ss(B)}(A_1, A_2)$, the cartesian square
\[
\begin{tikzcd}[ampersand replacement=\&]
A_0 \arrow[r] \arrow[d] \& A_1 \arrow[d] \\
A_2 \arrow[r] \& C
\end{tikzcd}
\]
is also cocartesian.}

\textit{Let $W$ be a subset of} $\Fl(\M)$ \textit{satisfying the following conditions.}
\[
\left\{
\begin{array}{@{\quad}l@{\quad}p{0.6\textwidth}}
\text{a)} & \textit{$W$ is a mildly saturated localizer.} \\[1.5ex]
\text{b)} & \textit{$W$ is L-accessible,} i.e. \textit{it is stable under filtering $\varinjlim$, and is accessible.} \\[1.5ex]
\text{c)} & $\Mon \cap W$ \textit{is stable under cobase change,} i.e. \textit{if we have a diagram which is} cocartesian
\begin{center}
\begin{tikzcd}[ampersand replacement=\&]
A \arrow[r, hook, "i"] \arrow[d, "\alpha"'] \arrow[dr, phantom, "\scriptstyle\text{cocart.}"] \& B \arrow[d, "\beta"] \\
A' \arrow[r, "i'"] \& B'
\end{tikzcd}
\end{center}
\textit{with $i$ a monomorphism belonging to $W$, then $i'$ is a monomorphism belonging to $W$.}\,\footnotemark[32]
\end{array}
\right.
\tag*{(7.2.3)}
\]
\footnotetext[32]{\textit{Cf.} below for a fourth condition $(d)$ on $W$, which ensures that $W$ fits into a Quillen triple.}

\msp{page 89}
\textit{Under these conditions, we have}
\[
\underbrace{\Mon}_{C} \cap W \;=\; \widetilde{W \cap \Mon} \quad , \qquad \text{i.e. } \Mon \cap W \text{ \textit{is Q-closed on the left}} \;,
\tag*{(7.2.4)}
\]


\endgroup
% END UNIT Der-ChXV-05-11-10/Der-ChXV-05-11-10_p046-050.tex
% BEGIN UNIT Der-ChXV-05-11-10/Der-ChXV-05-11-10_p051-055.tex
\begingroup\makeatletter%
\newcommand{\printedpage}[1]{\marginpar{\footnotesize p.~#1}}%
\newcommand{\msp}[1]{\marginpar{\footnotesize [#1]}}%
\newcommand{\Mon}{\mathrm{Mon}}%
\newcommand{\Ob}{\mathrm{Ob}}%
\newcommand{\Fl}{\mathrm{Fl}}%
\newcommand{\Ss}{\mathrm{Ss}}%
\newcommand{\card}{\mathrm{card}}%
\newcommand{\cercle}[1]{\tikz[baseline=(c.base)]\node[draw,circle,inner sep=1pt](c){#1};}%
\makeatother\let\printedpage\derpp\let\msp\dermsp%
\let\mspin\DERmspin%


\printedpage{49}
% head: A. Grothendieck : Dérivateurs, Ch. XV, Édition des Universités de Montpellier II et Paris VII
% foot: Version 5 novembre 2010

\textit{and}
\[
\text{\textit{there exists a small subset} } TC_0 \subset \Mon, \text{ \textit{such that} } \Mon \cap W = \widetilde{TC_0}\ . \tag*{(7.2.5)} % [sic] (7.2.5) -> numéro probablement autre (peut-être (7.2.3)), le même numéro (7.2.5) est repris plus bas pour une autre formule et précède (7.2.4)
\]

\textit{(More precisely, if} $\pi \geq \pi_0$ \textit{is a cardinal such that} $W$ \textit{is} $\pi$\textit{-accessible, we can take} $TC_0 = (\Mon \cap W)^{\pi}$ \textit{(or a} small \textit{subset of} $(\Mon \cap W)^{\pi}$ \textit{having the same isomorphism classes).) Thus} (th.~1 of §1) \textit{the couple} $(TC = \Mon \cap W,\ F = (TC)_*)$ \textit{is a closed Quillen couple, and there is even a} functorial \textit{factorization of every arrow} $f$ \textit{as}
\[
f \ = \ p\, i \ ; \qquad i \ \in \ TC \ , \qquad p \ \in \ F \ . \tag*{(7.2.4)}
\]
\textit{If we set}
\[
C \ = \ \Mon \ , \qquad TC \ = \ \Mon \cap W \ , \qquad F \ = \ (TC)_* \ , \qquad TF \ = \ C_* \ , \tag*{(7.2.5)}
\]
\textit{the triple} $(W, C, F)$ \textit{is a closed Quillen triple, if and only if we have the following additional conditions.}
\[
\left\{
\begin{array}{lll}
\text{(7.2.6)} & \cercle{g}\ \text{\textit{(on $\mathcal{M}$ only)}} & \parbox[t]{0.36\textwidth}{Mon \textit{stable under cobase change.}} \\[4ex]
\text{(7.2.7)} & \underbrace{TF}_{\overset{\mathrm{def}}{=}\,(\Mon)_*} \ \subset \ W & \parbox[t]{0.36\textwidth}{i.e. \textit{if} d) \textit{every} $f : X \longrightarrow Y$ \textit{which makes} $X$ \textit{an injective object of} $\mathcal{M}/Y$\textit{, is in} $W$ \textit{(hence also in} $W_{\mathrm{univ}}$\textit{).}}
\end{array}
\right.
\]
\DERmark{ENV0902}


\msp{page 90}

\textsc{Proof of} 7.1. We can give a somewhat more delicate statement, by supposing given from the outset, not a subset $M_0$ of $\Ob \mathcal{M}$, but a subset
\[
\Mon_0 \ \subset \ \Mon \ , \tag*{(7.1)}
\]
satisfying the following condition

\begin{quote}
For every monomorphism $i : A \longrightarrow B$ in $\mathcal{M}$ which is not an isomorphism, there exists $i' \in \Mon_0$, $i' : A' \longrightarrow B'$, and an arrow $i' \longrightarrow i$ in $\underline{\Fl}(\mathcal{M})$,
\DERmark{ENV0903}

\[
\begin{tikzcd}[row sep=large, column sep=large]
A' \arrow[r, "{i'}"] \arrow[d] & B' \arrow[d] \\
A \arrow[r, "i"] & B \ ,
\end{tikzcd}
\tag*{(7.2)}
\]
such that the corresponding arrow
\[
C \ = \ A \amalg_{A'} B' \ \longrightarrow \ B
\]
is a monomorphism, and $B' \longrightarrow B$ does not factor through $A$.
\end{quote}

That is what was aimed at by condition (7.1.1) on $M_0$, which implies (7.2) on $\Mon_0$ defined in terms of $M_0$ as said in 1°.

\printedpage{50}
% head: A. Grothendieck : Dérivateurs, Ch. XV, Édition des Universités de Montpellier II et Paris VII
% foot: Version 5 novembre 2010

In what follows, we suppose given a subset $\Mon_0 \subset \Fl(\mathcal{M})$ satisfying (7.2). In practice, it will often be deduced from a subset $M_0$ of $\Ob \mathcal{M}$, generating for monomorphisms, as made explicit in 1°, but we do not need the fact that $i' \longrightarrow i$ is a monomorphism (\textit{i.e.} (7.1.1) is stronger than what we need). E.g. if $B' \in M_0$ and $B' \longrightarrow B$ is given, not factoring through $A$, we shall be able to take $A' = B' \times_B A$ and $i' : A' \longrightarrow B'$ (to ensure $i' \in \Mon_0$, we shall therefore have to suppose $M_0$ stable under passage to subobjects), and chances are

\msp{page 91}

that $C \longrightarrow B$ will then indeed be a monomorphism \ldots

1°) We must prove
\[
\Mon \ \subset \ \widetilde{\Mon_0}
\]
assuming (7.2) (and the stability of $\mathcal{M}$ under filtering $\varinjlim$). This will follow directly from the

\textbf{Lemma 7.3.} \textit{Let $\mathcal{M}$ be a category stable under filtering direct limits, and $C_0$, $C$ subsets of} $\Fl(\mathcal{M})$\textit{. We shall suppose}
\begin{enumerate}[label={}]
\item[\cercle{a}] \textit{$C$ made up of monomorphisms, and}
\item[\cercle{b}] \textit{that the following condition (generalizing} (7.2) \textit{above) is satisfied.}
\end{enumerate}

\begin{quote}
\textit{For every arrow $i : A \hookrightarrow B$ in $C$ such that $i$ is not an isomorphism, there exists $i' : A' \longrightarrow B'$ in $\widetilde{C_0}$, and an arrow $i' \longrightarrow i$ in} $\underline{\Fl}(\mathcal{M})$\textit{, such that in the diagram}
\[
\begin{tikzcd}[row sep=large, column sep=large]
A' \arrow[r, "{i'}"] \arrow[d, hook] & B' \arrow[d] \arrow[rd, dotted] & \\
A \arrow[r, "\alpha"] \arrow[rr, dotted, bend right, "i"'] & D \ = \ A \amalg_{A'} B' \arrow[r, "\beta"] & B \ ,
\end{tikzcd}
\tag*{(7.3.1)}
\]
\textit{the arrow $\beta$ is in $C$, and $B' \longrightarrow B$ does not factor through $A$ (or, what amounts to the same, the subobject of $B$ defined by $D$ and $\beta$ is not equal to the one defined by $A$ and $i$).}
\end{quote}

\textit{We suppose moreover}
\begin{enumerate}[label={}]
\item[\cercle{c}] \textit{$C$ stable under filtering} $\varinjlim$\textit{, and}
\item[\cercle{d}] \textit{that for every $B \in \Ob \mathcal{M}$, the set of subobjects of $B$ is small.} (A most innocuous condition! This is a condition \textit{on} $\mathcal{M}$, not involving $C$).
\end{enumerate}
\textit{Then we have}
\[
C \ \subset \ \widetilde{C_0} \ , \qquad \text{\textit{hence} } \widetilde{C} \ \subset \ \widetilde{C_0} \ , \tag*{(7.3.2)}
\]
\textit{hence, if $C_0 \subset C$, we conclude}

\msp{page 92}
\[
\widetilde{C} \ = \ \widetilde{C_0} \ . \tag*{(7.3.3)}
\]

\printedpage{51}
% head: A. Grothendieck : Dérivateurs, Ch. XV, Édition des Universités de Montpellier II et Paris VII
% foot: Version 5 novembre 2010

(This lemma immediately implies the first assertion of 7.1, 1°.) Let us prove this lemma, \textit{i.e.} that every $i \in C$ is in $\widetilde{C_0}$, by constructing by transfinite induction a strictly increasing sequence of subobjects of $B$,
\[
A_\alpha \ \overset{j_\alpha}{\hookrightarrow} \ B \ , \qquad \text{with the } j_\alpha \in C \ ,
\]
and with:

\cercle{1°} $A_0 = A$. We shall prove that $A \xrightarrow{\ i_\alpha \,=\, i_{\alpha,0}\ } A_\alpha$ in $\widetilde{C_0}$, and there will be a final ordinal $\alpha$ such that $A_\alpha = B$\,; we shall therefore conclude from this that $i = i_\alpha$ is in $\widetilde{C_0}$.

If $i$ is an isomorphism, we know that $i \in \widetilde{C_0}$, there is nothing to prove (the transfinite construction stops at stage 0). The two steps of the construction, after the initial step 1°, are:

\cercle{2°} Let $A_\alpha$ be constructed. If $A_\alpha = B$, the construction is stopped. If $A_\alpha \xrightarrow{j_\alpha} B$ is not an isomorphism, we apply (7.3.1) to $j_\alpha \in C$, to find a $D \overset{j}{\hookrightarrow} B$, thus $j \in C$ and $D \supsetneqq A_\alpha$. Moreover,
\[
i_{\alpha+1,\alpha} \ : \ A_\alpha \ \longrightarrow \ A_{\alpha+1} \ \overset{\mathrm{def}}{=} \ D \qquad \text{is in } \widetilde{C_0} \ , \tag*{(7.3.4)}
\]
\DERmark{ENV0904}

since obtained from $i' \in \widetilde{C_0}$ by cobase change.

\cercle{3°} If $\alpha$ is a limit ordinal, we set
\[
A_\alpha \ = \ \varinjlim_{\beta < \alpha} A_\beta \ , \tag*{(7.3.5)}
\]

\msp{page 93}

and the canonical arrow
\[
A_\alpha \ \xrightarrow{j_\alpha} \ B
\]
is in $C$, since the $j_\beta$ are and $C$ is stable under filtering direct limits.

We have thus completed the description of the transfinite systems. For every $\alpha$, we have
\[
\card \alpha \ \leq \ \card \Ss(B) \ ,
\]
since the map $\beta \longmapsto A_\beta$ from the set $I_\alpha = \{\beta \mid \beta < \alpha\}$ to $\Ss(B)$ is \textit{strictly} increasing. Thus there exists an ordinal $\alpha$ such that $\card \alpha \leq \card \Ss(B)$, such that $A_\alpha = B$. But then $i \in \widetilde{C_0}$, as a transfinite direct limit (or composite) of arrows belonging to $\widetilde{C_0}$. (The condition $\forall\ \beta$, $i_{\beta+1,\beta} \in \widetilde{C_0}$ suffices to imply that all the transition morphisms $i_{\beta',\beta}$ are in $\widetilde{C_0}$.)
\DERmark{ENV0905}


Let us return to 7.1, 2°. It remains to prove that if Mon is stable under cobase change, then we have
\[
\Mon \ = \ \widetilde{\Mon_0} \quad .
\]

But since Mon is already stable under composition, and under arbitrary filtering direct limits (hence also under direct factors), it follows by corollary 2 of theorem 1, §1 (p. 11), that $\Mon = \widetilde{\Mon_0}$\ \ . The factorization statement is theorem 1 of §1.

\printedpage{52}
% head: A. Grothendieck : Dérivateurs, Ch. XV, Édition des Universités de Montpellier II et Paris VII
% foot: Version 5 novembre 2010

\msp{page 94}

Let us prove 7.1, 3°, \textit{i.e.} formula (7.1.6),
\[
W \cap \Mon \ \subset \ \widetilde{(W \cap \Mon)^{\pi}} \ .
\]

For this, we again apply lemma 7.3 to
\[
C_0 \ = \ (W \cap \Mon)^{\pi} \ , \qquad C \ = \ W \cap \Mon \ .
\]

The conditions (\textit{a}), (\textit{b}), (\textit{c}), (\textit{d}) of the lemma are satisfied. This is clear for all, except for (\textit{b}), which calls for a non-tautological verification. Let then
\[
f \ : \ A \hookrightarrow B \ , \qquad f \in \Mon \cap W,\ f \text{ not an isomorphism} \ .
\]

We regard $B$ as a direct limit\footnote[33]{\textit{Of a $\pi$-adapted direct system.}}
\[
B \ = \ \varinjlim_{I} B_i \ , \qquad I \ = \ \Ss_{\pi}(B) \ .
\]
For every $i$, we set
\[
A_i \ = \ A \cap B_i \ = \ A \times_B B_i \ ,
\]
and
\[
f_i \ : \ A_i \ \hookrightarrow \ B_i \ ,
\]
hence
\[
f \ = \ \varinjlim_{i \in I} f_i \ .
\]
We need to know that the direct system of arrows $f_i$ is $\pi$-adapted, which leads us to postulate the following additional exactness condition on $\mathcal{M}$:

\msp{page 95}
\[
e')\quad \text{the filtering } \varinjlim \text{ in } \mathcal{M} \text{ commute with base change}\ , \tag*{(7.4)}
\]

it even suffices to suppose it in the very special case we have considered here.

Since $W$ is $\pi$-accessible, and $(f_i)_{i \in I}$ is $\pi$-adapted, since $\varinjlim f_i \in W$, it follows that the set $I_W$ of $i \in I$ such that $f_i \in W$ is cofinal. Since $A \overset{f}{\hookrightarrow} B$ is not an isomorphism, there is therefore an $i \in I_W$ such that $B_i \hookrightarrow B$ does not factor through $A$. We then have the diagram
\[
\begin{tikzcd}[row sep=large, column sep=large]
A_i \arrow[r, hook, "f_i"] \arrow[d, hook] \arrow[rd, phantom, "\text{\footnotesize cocart.}"] & B_i \arrow[d, hook] \arrow[rd, dotted, hook, "\text{\footnotesize incl.}"] & \\
A \arrow[r, hook, "\alpha_i"] \arrow[rr, dotted, hook, bend right, "f"'] & C_i \arrow[r, "\beta_i"] & B \ ,
\end{tikzcd}
\]

\printedpage{53}
% head: A. Grothendieck : Dérivateurs, Ch. XV, Édition des Universités de Montpellier II et Paris VII
% foot: Version 5 novembre 2010
with $f_i \in (C \cap W)^{\pi} = C_0$, hence $\alpha_i \in \widetilde{C_0}$. We need to know that $\alpha_i \in W$, hence an additional condition on $W$:

\begin{quote}
If we have a cocartesian square
\[
\begin{tikzcd}[row sep=large, column sep=large]
A \arrow[r, hook, "i"] \arrow[d] & B \arrow[d] \\
A' \arrow[r, "{i'}"] & B' \ ,
\end{tikzcd}
\tag*{(7.5)}
\]
with $i \in \Mon \cap W$, then $i' \in W$\footnote[34]{\textbf{N.B.} This condition is satisfied.}.
\end{quote}

(I have strengthened (7.1.4) accordingly).

\msp{page 96}

Since by hypothesis $f = \beta_i \alpha_i$ is in $W$, and $\alpha_i \in W$, it follows by (7.1.4), (\textit{b}) that $\beta_i \in W$. We shall have to know moreover that $\beta_i$ is a monomorphism, hence $\beta_i \in \Mon \cap W = C$, which will finish establishing the missing condition (\textit{b}) of 7.3. This is the exactness condition (7.2.1), which I evidently have to add to the hypotheses in 7.1, 3°. Having done that, lemma 7.3 applies and establishes 7.1, 3°.
\DERmark{ENV0906}


\textsc{Proof of} 7.2. Let us set
\[
C \ = \ \Mon \ , \qquad TC \ = \ C \cap W \ , \qquad TC_0 \ = \ TC^{\pi} \ ,
\]
hence by 7.1, 3°, we have
\[
\underbrace{TC_0}_{\substack{\text{essentially}\\ \text{small}}} \ \subset \ TC \ \subset \ \widetilde{TC_0} \ = \ \widetilde{TC} \ ,
\]
and by theorem 1, §1 (p. 10-12), cor. 2, to check that $TC = \widetilde{TC_0}$ (\textit{i.e.} $TC$ is Q-closed on the left), it amounts to the same to see that it is Q-stable on the left, \textit{i.e.} stable under composition, under cobase change, under strict ordinal $\varinjlim$, under direct factors. Now the last two are implied by the fact that $TC = C \cap W$ is stable under filtering $\varinjlim$, seeing that $W$ is, and $C$ also by condition \cercle{e} (p. 76) on $\mathcal{M}$. On the other hand, $C = \Mon$ and $W$ being stable under composition, likewise for $C \cap W$. So there remains cobase

\msp{page 97}

change, and this is ensured by 7.2.3, (\textit{c}) (which we already needed in 7.1, 3°, to establish $C \cap W \subset \widetilde{TC_0}$). (\textbf{N.B.} the hypotheses of 7.2 are practically those of 7.1, 3°, except that we suppose that $W$ is a localizer (understood: mildly saturated), to be precise we have used, up to now, as an additional hypothesis, only the stability of $W$ under composition.)

It remains to establish that $(W, C = \Mon, F = (TC)_*)$ is a closed Quillen triple, if and only if the conditions (7.2.6) and (7.2.7), \textit{i.e.}
\[
\left\{
\begin{aligned}
& C = \Mon \ \ \text{stable under cobase change,} \\
& TF = C_* = (\Mon)_* \subset W,
\end{aligned}
\right.
\]


\endgroup
% END UNIT Der-ChXV-05-11-10/Der-ChXV-05-11-10_p051-055.tex
% BEGIN UNIT Der-ChXV-05-11-10/Der-ChXV-05-11-10_p056-060.tex
\begingroup\makeatletter%
\newcommand{\printedpage}[1]{\marginpar{\footnotesize p.~#1}}%
\newcommand{\msp}[1]{\marginpar{\footnotesize [#1]}}%
\newcommand{\Mon}{\mathrm{Mon}}%
\newcommand{\Fl}{\mathrm{Fl}}%
\newcommand{\Cof}{\mathrm{Cof}}%
\newcommand{\Fib}{\mathrm{Fib}}%
\newcommand{\Cat}{\mathrm{Cat}}%
\newcommand{\univ}{\mathrm{univ}}%
\newcommand{\chap}{^{\wedge}}%
\makeatother\let\printedpage\derpp\let\msp\dermsp%
\let\mspin\DERmspin%


\printedpage{54}
% head: A. Grothendieck : Dérivateurs, Ch. XV, Édition des Universités de Montpellier II et Paris VII
% foot: Version 5 novembre 2010
are satisfied. Necessity is obvious, by definition of closed Quillen triples, and for sufficiency, we note that it follows from theorem 7.1, 2$^\circ$, and from the part of 7.2 already proved, that $(C,TF)$, $(TC,F)$ are closed Quillen couples, and we have moreover
\[
TC \;=\; C\cap W\ , \qquad TF \;\subset\; F\cap W
\]
(since $TF\subset W$ by (7.2.7), and $TF\subset F$ since $TC\subset C$). We know
\msp{page 98}
(paragraph 4.3, corollary 4, page 36) that it follows that $(W,C,F)$ is a closed Quillen triple. This completes the proof of 7.1 and 7.2.

Let us sum up the essentials:

\textbf{Corollary 7.4.}\footnote[35]{\textbf{N.B.} It is ultimately this corollary that is the main result of the present section.} \textit{Let $\mathcal{M}$ be a category satisfying the conditions} (a), (b), (c), (d), (e) (p.~76) \textit{and} (e$'$), (f) ((7.2.1), p.~88), \textit{finally} (g) (7.2.6), i.e. \textit{stability of the set} Mon \textit{of monomorphisms of $\mathcal{M}$ under cobase change. Let $W\subset \Fl(\mathcal{M})$ be a localizer in $\mathcal{M}$, which is} stable under filtering $\varinjlim$ \textit{and moreover} (\textit{a subtler condition, but in practice innocuous}) accessible. \textit{We put}
\[
C \;=\; \Mon\ , \qquad TC \;=\; C\cap W\ , \qquad TF \;=\; C_*\ , \qquad F \;=\; (TC)_*\ .
\]
\textit{Then:}
\begin{itemize}[leftmargin=*]
\item[1$^\circ$)] \textit{$(C,TF)$ is a Quillen couple} (\textit{a statement independent of the datum $W$}), \textit{and even a Quillen couple with functorial factorization.}
\item[2$^\circ$)] \textit{For $(TC,F)$ to be a Quillen couple, it is necessary and sufficient that $TC$ be stable under cobase change.} (\textit{And then there is functorial factorization.})\footnote[36]{\textbf{N.B.} In cases 1$^\circ$ and 2$^\circ$, $C$, $TC=C\cap W$ are Q-generated on the left by \textit{small} subsets.}.
\item[3$^\circ$)] \textit{For $(W,C,F)$ to be a closed Quillen triple, it is necessary and sufficient that the following conditions be satisfied.}
\[
\left\{
\begin{aligned}
&\text{a)}\quad TC \overset{\mathrm{def}}{=} W\cap \Mon\ \textit{stable under cobase change.}\\
&\text{b)}\quad TF \overset{\mathrm{def}}{=} (\Mon)_* = \textit{maps of injective type} \subset W.
\end{aligned}
\right.
\tag*{(7.4.1)}
\]
\DERmark{ENV0907}\DERmark{ENV0908}

\item[4$^\circ$)] \msp{page 99}\textit{For $(W,C,F)$ to be a closed and left proper Quillen triple} (i.e. $C\subset \Cof_W$), \textit{it is necessary and sufficient that we have condition} (7.4.1), (b) (i.e. \textit{$W$ contains the morphisms of injective type}), \textit{and moreover}
\[
\begin{aligned}
&C\subset \Cof_W,\ \textit{or again: for every cocartesian square}\\
&\qquad\qquad\begin{tikzcd}[ampersand replacement=\&, row sep=large, column sep=large]
A \arrow[r, "i"] \arrow[d, "\alpha"'] \& B \arrow[d, "\beta"] \\
A' \arrow[r, "i'"] \& B'\ ,
\end{tikzcd}\\
&\textit{with $i$ a monomorphism, then } \alpha\in W \Longrightarrow \beta\in W.
\end{aligned}
\tag*{(7.4.2)}
\]

\item[5$^\circ$)] \printedpage{55}
% head: A. Grothendieck : Dérivateurs, Ch. XV, Édition des Universités de Montpellier II et Paris VII
% foot: Version 5 novembre 2010
\textit{For $(W,C,F)$ to be a closed Quillen triple and} proper, \textit{it is necessary and sufficient that the two conditions of 4$^\circ$ be satisfied, and also that we have:}
\[
\begin{aligned}
&F\subset \Fib_W,\ \textit{or again: for every cartesian square}\\
&\qquad\qquad\begin{tikzcd}[ampersand replacement=\&, row sep=large, column sep=large]
X \arrow[r, "q"] \arrow[d, "f"'] \& X' \arrow[d, "f'"] \\
Y \arrow[r, "p"'] \& Y'\ ,
\end{tikzcd}\\
&\textit{if } f'\in F \ (\text{i.e. } f'\in (C\cap W)_*)\ \textit{and } p\in W,\ \textit{we have } q\in W.
\end{aligned}
\tag*{(7.4.3)}
\]
\DERmark{ENV0909}

\end{itemize}

\textsc{Proof of} 7.4. Parts 1$^\circ$, 2$^\circ$, 3$^\circ$ are just a summary of what has been established. For 4$^\circ$, we have to see that the stated conditions imply (7.4.1), $(a)$, now we have $C\subset \Cof_W$, hence $C\cap W\subset \Cof_W\cap W = W^{\univ}$
\msp{page 100}
hence $C\cap W = C\cap W^{\univ}$, which is stable under cobase change, since $C$ and $W^{\univ}$ are. As for 5$^\circ$, it is a tautology, taking 4$^\circ$ into account.

\section*{8\quad Application to the categories of models $(A\chap,W_A)$}

Let $A$ be in Cat, and let us consider in $A\chap$ the localizer $W_A$, inverse image under
\[
\left\{
\begin{aligned}
i_A : A\chap \quad &\longrightarrow \quad \Cat\\
F \quad &\longmapsto \quad A/F
\end{aligned}
\right.
\tag*{(8.1)}
\]
of a given basic localizer $W$, satisfying the conditions
\[
\mathrm{W}\,(\underbrace{1}_{\text{sat.}},\underbrace{2}_{\text{loc.}},\underbrace{3}_{\text{final obj.}}\,,\underbrace{5}_{\text{conn. }\pi_0}\,,\underbrace{6}_{\varinjlim},\underbrace{7\text{ bis}}_{\substack{\text{coc. squares}\\ \text{or ``M.V.''}}}\,,\underbrace{8}_{\text{acc.}}),\quad
\text{\parbox[t]{4.3cm}{where W\,(8) is the accessibility of $W$ (thus $W$ is L-accessible if it also satisfies W\,(6)).}}
\tag*{(8.2)}
\]

To be able however to apply the developments on accessible subsets of a big category in §5, we must verify the conditions $(a)$, $(b)$, $(c)$, $(d)$, $(e)$ of §5, both in Cat, and (especially) in $A\chap$.

a) Stability under small filtering $\varinjlim$, O.K.

b) Cardinal filtration. Well known for $A\chap = \underline{\mathrm{Hom}}(A^{\mathrm{op}},\mathrm{Ens})$, starting from that of Ens by the cardinals (\textit{cf.} SGA 4 I, 9.24, p.~177). For Cat, we define a cardinal filtration
\DERmark{ENV0910}

\[
(\Cat^{\pi})_{\pi\geq\pi_0}\ , \qquad \text{where } \pi_0 = \text{cardinal of countable sets}\ ,
\tag*{(8.3)}
\]
\printedpage{56}
% head: A. Grothendieck : Dérivateurs, Ch. XV, Édition des Universités de Montpellier II et Paris VII
% foot: Version 5 novembre 2010
\msp{page 101}
by the condition, for $A$ in Cat,
\[
A\in \mathrm{Filt}^{\pi}(\Cat) \;=\; \Cat^{\pi} \qquad\Longleftrightarrow\qquad \operatorname{card}\Fl(A) \;\leq\; \pi\ .
\tag*{(8.4)}
\]

We must verify the conditions $(a)$, $(b)$, $(c)$ of \textit{loc. cit.} 9.12. Condition $(a)$ is obvious, $(b)$ is contained in the stronger condition of stability of $\Cat^{\pi}$ under $\varinjlim\limits_{I}$, for $I$ a category such that $\operatorname{card}\Fl(I)\leq\pi$ (this is condition $(d_0)$ of these notes, p.~45-46). Let us note that $\Cat^{\pi}$ is also stable under finite $\varprojlim$ (and if $\pi$ is of the form $2^c$, it is even stable under inverse limits $\varprojlim\limits_{I}$, with $\operatorname{card}\Fl(I)\leq c\ \ldots$). Condition $(c)$ is easily verified by the method of §5, noting that for $X$ in Cat, the set
\[
\mathrm{Ss}_{\pi}(X)
\]
of subcategories of $X$ which are in $\Cat^{\pi}$ is upward filtering and large relative to $\pi$, and that $X$ is its direct limit. On the other hand, it is also clear that if $X\in\Cat^{\pi'}$ with $\pi'>\pi$, then
\[
\operatorname{card}\mathrm{Ss}_{\pi}(X) \;\leq\; \pi'^{\pi}\ ,
\]
which finishes establishing condition $(c)$.

Moreover, condition $(d)$ on p.~45 is immediately verified for the $\Cat^{\pi}$, and also for the $A^{\wedge\pi}$ (for $\pi\geq\operatorname{card}\Fl(A)$).

Also that monomorphisms are stable under filtering $\varinjlim$ (condition $(e)$, page 46).

\msp{page 102}
Let us pass to the additional conditions $(e')$, $(f)$, $(g)$ required in corollary 7.4. Condition $(e')$ (commutation of filtering $\varinjlim$ with base change) is satisfied in $A\chap$ and in Cat. Condition $(f)$ says that for two subobjects $A_1$, $A_2$ of an object $B$, if $A_0=A_1\cap A_2$, then
\[
A_1\amalg_{A_0}A_2 \longrightarrow B
\]
is a monomorphism. I doubt that this is true in Cat, and it does not matter, for it is to $A\chap$, not to Cat, that we shall want to apply 7.4 -- and it is obviously true in $A\chap$, and in any topos, being so in (Ens). Finally, stability of monomorphisms under cobase change is surely false in Cat, but true in $A\chap$ -- so it is O.K.

Since $W$ is stable under filtering $\varinjlim$ (W\,(6)), and $i_A$ commutes with them, $W_A$ is also stable. Similarly, $W_A$ is L-accessible, by the

\textbf{Lemma 8.1.} \textit{Let $\mathcal{M}$, $\mathcal{M}'$ be two categories satisfying} (a), (b), (c), (d), (e), \textit{and $\varphi:\mathcal{M}\longrightarrow\mathcal{M}'$ a functor commuting with filtering $\varinjlim$. Then for any L-accessible subset $M'$ of $\mathcal{M}'$, $\varphi^{-1}(M')$ is L-accessible.}

\msp{page 103}
\textit{To be precise, let $\pi$ be a cardinal such that $M'$ is $\pi$-accessible and $\varphi$ maps $\mathcal{M}^{\pi}$ into $\mathcal{M}'^{\pi}$ (for example, for this last condition, it suffices $\pi=2^c$ for $c$ large enough,} cf. loc. cit. 9.14, p.~153). \textit{Then $M=\varphi^{-1}(M')$ is $\pi$-accessible.}

Immediate verification.

We therefore have the

\printedpage{57}
% head: A. Grothendieck : Dérivateurs, Ch. XV, Édition des Universités de Montpellier II et Paris VII
% foot: Version 5 novembre 2010
\textbf{Lemma 8.2.} \textit{Let $A$ be in} Cat. \textit{Then $(A\chap,W_A)$ satisfies the preliminary conditions of corollary} 7.4, \textit{as well as the following conditions.}
\[
\underbrace{\Mon(A\chap)}_{\substack{\textit{set of}\\ \textit{monomorphisms}\\ \textit{in } A\chap}} \quad\subset\ \Cof_{W_A}\ ,
\tag*{(8.2.1)}
\]
\[
\Mon\cap W_A \qquad \textit{stable under cobase change}\ ,
\tag*{(8.2.2)}
\]
\textit{in other words} (\textit{taking into account that} Mon \textit{is stable under cobase change}), \textit{for a cocartesian square}
\[
\begin{tikzcd}[row sep=large, column sep=large]
A \arrow[r, hook, "i"] \arrow[dr, "\alpha"'] & B \arrow[dr, "\beta"] & \\
& A' \arrow[r, hook, "i'"] & B'\ ,
\end{tikzcd}
\tag*{(8.2.3)}
\]
\textit{with $i$ a monomorphism, if $\alpha\in W_A$ (resp. $i\in W_A$), we have $\beta\in W_A$ (resp. $i'\in W_A$). In fact, we have better:}
\[
i\ \in\ W_A \ \Longleftrightarrow\ i'\ \in\ W_A\ , \qquad \alpha\ \in\ W_A \ \Longleftrightarrow\ \beta\ \in\ W_A\ .
\tag*{(8.2.4)}
\]

\msp{page 104}
It remains to verify (8.2.4), which was seen in XII \ldots

We therefore conclude from 7.4 1$^\circ$, 2$^\circ$, 3$^\circ$:

\textbf{Corollary 8.3.} \textit{Let us put}
\[
\left\{
\begin{aligned}
C \;&=\; \Mon(A\chap)\ , & TC \;&=\; C\cap W_A\ ,\\
F \;&=\; (TC)_*\ , & TF \;&=\; C_*\ .
\end{aligned}
\right.
\tag*{(8.3.1)}
\]
\DERmark{ENV0911}

\textit{Then $(C,TF)$ and $(TC,F)$ are Quillen couples with functorial factorization.} (\textbf{N.B.} \textit{$C$ and $TC$ being moreover Q-generated on the left by small subsets of} $\Fl(A\chap)$.) \textit{For $(W_A,C,F)$ to be a closed Quillen triple, it is necessary and sufficient that we have}
\[
\underbrace{TF}_{\overset{\mathrm{def}}{=}\,(\Mon)_*} \quad\subset\ W_A\ ,
\tag*{(8.3.2)}
\]
\DERmark{ENV0912}

i.e. \textit{that every arrow $f:X\longrightarrow Y$ of injective type in $A\chap$} (i.e. \textit{which makes $X$ an injective object of $A\chap/Y$}) \textit{be in $W_A$.}

\textit{When this condition is satisfied, the Quillen triple $(W_A,C,F)$ is left proper, by} 8.2.1.

It is not clear that it is necessarily proper-plain-and-simple, \textit{i.e.} that we have
\[
F\ \subset \Fib_W \qquad (?)\ .
\tag*{(8.3.3)}
\]

We will come back to this later. For the moment, the essential thing is to examine the condition (8.3.2) on $A$.

\printedpage{58}
% head: A. Grothendieck : Dérivateurs, Ch. XV, Édition des Universités de Montpellier II et Paris VII
% foot: Version 5 novembre 2010
\msp{page 105}
\textbf{Theorem 8.4.} \textit{Let $W$ be a basic localizer in} Cat \textit{satisfying} $\mathrm{W}\,(1,2,3,6,7,8)$, \textit{notations of} XII, (i.e. \textit{all the axioms considered so far, except perhaps the one} W\,(4) \textit{on fibrations} \texttt{[and the one W\,(5) on connectedness]}), \textit{let us suppose moreover $W$ stable under direct factors} (\textit{strengthening} W\,(1\,bis) \textit{of the axiom of mild saturation} W\,(1)).\footnote[37]{\textbf{N.B.} For the validity of 1$^\circ$, 3$^\circ$ we do not need W\,(6) ( $\varinjlim$ ) or W\,(8) (accessibility), nor even W\,(7) (cocartesian square axiom). The latter is already necessary to ensure part $(a)$ of corollary 8.5. But for part 2$^\circ$ of theorem 8.4, we need \textit{all} the axioms of the basic localizer $W$, except only W\,(4) \texttt{[and W\,(5)]}.} \textit{Let $A$ be an object of} Cat, \textit{and $W_A=i_A^{-1}(W)$ the localizer of $A\chap$ deduced from $W$. With the notations of} 8.3:
\begin{itemize}[leftmargin=*]
\item[1$^\circ$)] \textit{The following conditions are equivalent.}
\begin{itemize}[leftmargin=*]
\item[a)] \textit{We have} (8.3.2),
\[
TF\ \subset\ W_A\ ,
\]
i.e. \textit{every arrow of injective type in $A\chap$ is in $W_A$.}
\item[b)] \textit{Let}
\[
L \;=\; L_A
\tag*{(8.4.1)}
\]
\textit{be the Lawvere element in $A\chap$, hence the presheaf}
\[
a \quad\longmapsto\quad
\begin{aligned}
&\textit{set of subpresheaves of $a$}\\
&(\text{i.e. }\textit{set of opens of $A/a$},\\
&\text{i.e. }\textit{of sieves of $A/a$})\ ,
\end{aligned}
\]
\textit{and let us consider the arrow $L\longrightarrow e_A$ (final object of $A\chap$). We want this arrow to be $W_A$-aspheric,} i.e. \textit{to be universally in $W_A$ -- or what comes to the same thing, that for every $a\in A$, $L_A\times a$ $(=L_{A/a})$ be a $W_A$-aspheric object}
\msp{page 106}
\textit{of $A\chap$ (or, what comes to the same thing, of $(A/a)\chap$).}
\item[c)] \textit{There exists an object $I$ of $A\chap$ such that $I\longrightarrow e_A$ is $W_A$-aspheric, with two sections $\delta_0$, $\delta_1$ of $I$ over $e_A$, with} $\mathrm{Ker}(\delta_0,\delta_1)=\emptyset_A$ \textit{(empty presheaf).}
\item[d)] \textit{The localized categories $A/a=B$ are} test categories \textit{for $W$} (cf. \textit{Pursuing Stacks}, \textit{§65, especially pages 172-173), which amounts to saying that we have this, for each of the categories $B=A/a$:}
\[
\begin{aligned}
&\textit{The functors}\\
&\left\{
\begin{aligned}
&i_B:\ \ B\chap \ \longrightarrow\ \Cat\ \ (\ F\ \longmapsto\ B/F\ )\\
&j_B:\ \ \Cat\ \longrightarrow\ B\chap\ \ (\ C\ \longmapsto\ (b\longmapsto \mathrm{Hom}(B/b,C))\ )
\end{aligned}
\right.
\end{aligned}
\tag*{(8.4.2)}
\]
(\textit{where $j_B$ is the right adjoint of $i_B$}) \textit{are both ``model preserving''}, i.e.
\[
\left\{
\begin{aligned}
W_B \;&=\; i_B^{-1}(W) \qquad \textit{(by definition of $W_B$, in fact)}\\
W \;&=\; j_B^{-1}(W_B)\ ,
\end{aligned}
\right.
\tag*{(8.4.3)}
\]
\end{itemize}
\end{itemize}


\endgroup
% END UNIT Der-ChXV-05-11-10/Der-ChXV-05-11-10_p056-060.tex
% BEGIN UNIT Der-ChXV-05-11-10/Der-ChXV-05-11-10_p061-065.tex
\begingroup\makeatletter%
\newcommand{\printedpage}[1]{\marginpar{\footnotesize p.~#1}}%
\newcommand{\msp}[1]{\marginpar{\footnotesize [#1]}}%
\newcommand{\Simp}{\boldsymbol{\Delta}}%
\newcommand{\Cub}{\square}%
\makeatother\let\printedpage\derpp\let\msp\dermsp%
\let\mspin\DERmspin%


\printedpage{59}
% head: A. Grothendieck : Dérivateurs, Ch. XV, Édition des Universités de Montpellier II et Paris VII

\begin{itemize}[label={},leftmargin=4.5em]
\item \textit{and the localized functors}
\[
\begin{tikzcd}
\mathrm{Hot}_{B,W} \arrow[r, shift left, "\overline{i_B}"] & \mathrm{Hot}_W \arrow[l, shift left, "\overline{j_B}"]
\end{tikzcd}
\]
\textit{are equivalences of categories. Moreover, these equivalences of categories are quasi-inverse to one another, the adjunction morphisms being deduced from the adjunction morphisms for the pair of adjoint functors $(i_B, j_B)$, these adjunction morphisms being moreover in $W$ resp. in $W_B$.}\footnote[38]{\textbf{N.B.} Condition $(d)$, seemingly very strong, is equivalent to the seemingly more innocuous condition:\par
$d'$) the functor $j_A : \mathrm{Cat} \longrightarrow A^\wedge$ transforms elements of $W$ into $W_A$-aspheric arrows;\par\noindent
and it even suffices to require it for $C \longrightarrow e$ in Cat, where $C$ is an object of Cat having a final object.\par\noindent
\texttt{[Condition ($d'$) above is too strong and is not equivalent to condition ($d$); on the other hand the weaker condition considered next is indeed equivalent to ($d$)].}}
\end{itemize}

\msp{page 107}

\begin{enumerate}[leftmargin=2.5em]
\item[$2^\circ$)] \textit{If the equivalent conditions of $1^\circ$ are satisfied, then the triple}
\[
(W_A\,,\ C = \mathrm{Mon}_{A^\wedge}\,,\ F = (C \cap W_A)_*)
\]
\textit{is a (closed) left proper Quillen triple.}
\item[$3^\circ$)] \textit{If the equivalent conditions of $1^\circ$ are satisfied, and if moreover $A$ is $W$-aspheric, then (and only then) $A$ itself is a} test category \textit{for $W$, which amounts to saying that the conditions on the category $B$ in} (d) \textit{are satisfied when $B$ is} either \textit{$A$, or a localized category $A/a$.}
\end{enumerate}

\textbf{Corollary 8.5.} \textit{The equivalent conditions of $1^\circ$ and of $3^\circ$ are satisfied in each of the following cases (and this for any $W$, satisfying the stated hypotheses).}
\begin{enumerate}[leftmargin=2.5em]
\item[a)] $A = \Simp$, $\Simp^{\sim}$, $\Cub$, $\Cub^{\sim}$ \textit{(categories of standard combinatorial simplices, ordered or not, and of standard combinatorial cubes, ordered or not \ldots)}\footnote[39]{$\Cub^{\sim}$ is essentially the ``universal'' category $A$, stable under finite products, and containing a ``separating interval'' $(I,\delta_0,\delta_1)$ as in $(b)$.}.
\item[b)] \textit{$A$ is stable under finite products}\footnote[40]{In fact, it suffices that $I \longrightarrow e$ be squarable, \textit{i.e.} that $A$ have a final object, and that the products $I \times a$ ($a \in \mathrm{Ob}\,A$) exist in $A$.}\textit{, and contains an object $I$ such that $I$ has two sections $\delta_0$, $\delta_1$ over the final object, with $\delta_0$, $\delta_1$ disjoint,} i.e. $\mathrm{Ker}(\delta_0,\delta_1) = \emptyset_{A^\wedge}$ \textit{(empty subpresheaf of the final presheaf).}
\end{enumerate}

\msp{page 108}

Case $(b)$ of the corollary falls under condition $(c)$ in $1^\circ$. (Moreover, $A$ having a final object, is $W$-aspheric, so we are under the conditions of $3^\circ$.) The cases $\Simp$, $\Simp^{\sim}$, $\Cub$ (and $\Cub^{\sim}$, which in fact falls under case $(b)$) are treated in Pursuing Stacks \S34 (p. 47-50) \texttt{[see also \S36]}.

\textsc{Proof of theorem} 8.4. The essential part, technically, is of course $1^\circ$\footnote[41]{Part $2^\circ$ follows from it, thanks to corollary 8.3.}. We have the $\pm$ tautological implications
\[
(a) \ \Longrightarrow\ (b) \ \Longrightarrow\ (c)\ ,
\]
where $(a) \Longrightarrow (b)$ comes from the fact that $L$ is obviously an injective object of $A^\wedge$. (The Lawvere object of any topos is an injective object.) It remains to establish
\[
(c) \ \Longrightarrow\ (d) \ \Longrightarrow\ (a)\ .
\]
% foot: Version 5 novembre 2010

\printedpage{60}
% head: A. Grothendieck : Dérivateurs, Ch. XV, Édition des Universités de Montpellier II et Paris VII

That we have $(c) \Longrightarrow (d)$ (and $(c) \Longrightarrow (b)$) is established in Pursuing Stacks, see the summary of the $W$-test notions in Pursuing Stacks, \S65, p. 168-179; the relevant statement here is proposition 3, page 173.

It remains to establish $(d) \Longrightarrow (a)$ – but now we're on easy street! I note here the following lemma, whose place would rather be at the end of \S4 of the generalities on Quillen triples:

\msp{page 109}

\textbf{Lemma 8.6.} \textit{Let $\mathcal{M}$ be a category, $(C,TF)$ a Quillen precouple in $\mathcal{M}$, and $W \subset \mathrm{Fl}(\mathcal{M})$ a subset of} $\mathrm{Fl}(\mathcal{M})$ \textit{stable under direct factors. Equivalent conditions:}
\begin{enumerate}[leftmargin=2.5em]
\item[a)] \textit{$TF \subset W$.}
\item[b)] \textit{Every arrow $f$ in $\mathcal{M}$ factors as $f = pi$, with $i \in C$, $p \in W$.}
\item[b$'$)] \textit{As in} (b), \textit{with $f$ in $TF$.}
\end{enumerate}

\textsc{Proof}. The implications $(a) \Longrightarrow (b) \Longrightarrow (b')$ are tautological and do not appeal to the stability hypothesis on $W$. But for a given $f \in TF$, to see that it is in $W$, if $W$ is stable under direct factors, it suffices for that particular $f$ to factor as $pi$, with $i \in C$, $p \in W$. Indeed, we see at once that then $f$ is a direct factor of $p$ (since $i \in C$), and since $p \in W$, we have $f \in W$, \quad \textit{q.e.d.} % [sic] $i \in W$ -> $f \in W$, the conclusion to be proved concerns f
\DERmark{ENV0913}


Let us return to 8.4, proof of $(d) \Longrightarrow (a)$.

\texttt{[There follows a crossed-out passage, at the bottom of page 109 and at the top of page 110.]}

\msp{page 110}

In the end, my plan of proof changes, for although $(d)$ looks very, very strong as a condition, as it stands I do not see, in the end, how to deduce $(a)$ from it. But in fact we can consider that the equivalence of $(b)$, $(c)$, $(d)$ is put here for the record, it is established in \textit{loc. cit.}, so I only have to concern myself with the implication $(c) \Longrightarrow (a)$, which will follow
\msp{page 111}
from lemma 8.6, which reduces us to proving that every arrow $f : X \longrightarrow Y$ in $A^\wedge$ factors as $pi$, with $i$ a monomorphism and $p \in W_A$. We use the mapping cone argument \texttt{[rather mapping cylinder]}, starting from
\[
\begin{tikzcd}[column sep=small]
e_A \amalg e_A \arrow[rr, hook, "{(\delta_0,\delta_1)}", "\text{mono}"'] \arrow[dr] & & I \arrow[dl, "(W_A)_{\mathrm{univ}}"] \\
& e_A\ , &
\end{tikzcd}
\]
whence for every $X$ in $A^\wedge$, taking the product with $X$,
\[
\begin{tikzcd}[column sep=small]
X \amalg X \arrow[rr, hook, "{(\delta_0^X,\delta_1^X)}", "\text{mono}"'] \arrow[dr, "{(\mathrm{id},\mathrm{id})}"] & & X \times I \arrow[dl, "\mathrm{pr}_1^X \in W_A"] \\
& X\ , &
\end{tikzcd}
\]
% foot: Version 5 novembre 2010
\printedpage{61}
% head: A. Grothendieck : Dérivateurs, Ch. XV, Édition des Universités de Montpellier II et Paris VII
and we construct the mapping cone \texttt{[rather mapping cylinder]} of $f$ by the cocartesian diagram
\[
\begin{tikzcd}[column sep=4em,row sep=4em]
 & X \arrow[r, "f"] \arrow[d, hook, shift right, "\delta_0^X"'] \arrow[dr, phantom, "\text{cocart.}"] & Y \arrow[d, hook, shift right, "\beta"'] \\
X \arrow[r, dotted, "\delta_1^X"] \arrow[rr, dotted, bend right, "i_f"'] & X \times I \arrow[r, "f'"] \arrow[u, dotted, shift right, "\mathrm{pr}_1^X"'] & Z(f)\ , \arrow[u, dotted, shift right, "p_f"']
\end{tikzcd}
\qquad
\begin{array}{c}
p_f \text{ deduced from} \\
\mathrm{pr}_1^X : X \times I \to X
\end{array}
\tag*{(8.6.1)}
\]
the desired factorization of $f$ is
\[
f \ = \ p_f\, i_f\ , \qquad \text{where}\quad i_f = f'\delta_1^X\ .
\]
Here, since $\mathrm{pr}_1^X \in W_A$, we conclude that
\msp{page 112}
$\delta_0^X$, $\delta_1^X$ (which are sections of it) are in $W_A$, since $\delta_0^X$ is a monomorphism, we thus have $\delta_0^X \in TC$, hence ($TC$ being stable under cobase change) $\beta \in W_A$, hence finally $p_f$ (which is a retraction of $Z(f)$ onto $Y$) is also in $W_A$. It remains to see that $i_f$ is a monomorphism, which is seen from the more complete diagram (\textit{cf.} XII, p. 230) that I have in front of me: $i_f$ appears as a composite of two monomorphisms $X \longrightarrow Y \amalg X$ and $Y \amalg X \longrightarrow Z(f)$ \ldots{} Theorem 8.4 is proved.

I will all the same isolate the everlasting formal argument (which I have used mostly in dual form) in a

\textbf{Lemma 8.7.} \textit{Let $\mathcal{M}$ be a category, $W$ a mildly saturated localizer in $\mathcal{M}$, $C$ a subset stable under composition and under cobase change. We propose, for $X$ fixed in $\mathcal{M}$, to find a functorial factorization of every arrow $f : X \longrightarrow Y$} (i.e. \textit{of the costructural morphism of a $Y$ in $X\backslash\mathcal{M}$) as $p_f i_f$, with $i_f \in C$, $p_f \in W$. We assume for this that we have}
\msp{page 113}
\textit{a factorization of this type for $X \amalg X \xrightarrow{(\,\mathrm{id}\,,\mathrm{id}\,)} X$,} i.e. \textit{an object $I_X$ and a factorization}
\[
\begin{tikzcd}[column sep=6em]
X \amalg X \arrow[r, hook, "{i_X = (\delta_0^X,\delta_1^X)}"] \arrow[rr, dotted, bend right, "{(\mathrm{id}_X,\mathrm{id}_X)}"'] & I_X \arrow[r, "p_X"] & X\ ,
\end{tikzcd}
\]
\textit{with $i_X \in C$, $p_X \in W$. We then associate to every $f$ the mapping cone} \texttt{[rather cylinder]} \textit{$Z(f)$ as in diagram} (8.6.1) \textit{(where one must read $I_X$ instead of $X \times I$), and we then have $p_f \in W$, $i_f \in C$ as desired. But one must assume, to know that $i_f \in C$, that $X \in \mathcal{M}_C$ ($C$-cofibrant objects),} i.e. \textit{that for every $Y$, $X \longrightarrow Y \amalg X$ is in $C$. (Moreover, we have tacitly assumed $\mathcal{M}$ stable under finite $\varinjlim$.)}
% foot: Version 5 novembre 2010

\printedpage{62}
% head: A. Grothendieck : Dérivateurs, Ch. XV, Édition des Universités de Montpellier II et Paris VII

\textit{The complete diagram is}
\[
\left\{
\begin{array}{c}
\begin{tikzcd}[ampersand replacement=\&,column sep=4em,row sep=3.5em]
 \& X \arrow[r, "f"] \arrow[d, "\mathrm{inj}_0^X"'] \arrow[dr, phantom, "\text{cocart.}"] \& Y \arrow[d] \\
X \arrow[r, "\mathrm{inj}_1^X"] \arrow[dr, dotted, "\delta_1^X"'] \& X \amalg X \arrow[r, "f''"] \arrow[d, "i_X"'] \arrow[dr, phantom, "\text{cocart.}"] \& Y \amalg X \arrow[d, "j"] \\
 \& I_X \arrow[r, "f'"'] \arrow[uu, dotted, bend right=45, "p_X"{pos=.12,swap}] \& Z(f) \arrow[uu, dotted, bend right=45, "p_f"{pos=.25,swap}]
\end{tikzcd}
\\[1ex]
f = p_f i_f, \quad \textit{where} \quad i_f = f' \circ \delta_1^X = j \circ (\underbrace{f'' \circ \mathrm{inj}_1^X}_{\lambda})\ .
\end{array}
\right.
\tag*{(8.7.1)}
\]

\msp{page 114}

We have $i_X \in C$ by hypothesis, hence $j \in C$ (stability of $C$ under cobase change), we have $\lambda : X \longrightarrow Y \amalg X$ in $C$ since $X$ in $\mathcal{M}_C$, hence the composite $i_f = j \circ \lambda$ is in $C$ (stability of $C$ under composition) – we win!

\textbf{Comments.} For the equivalence of $(a)$, $(b)$, $(c)$ in theorem 8.4, $1^\circ$, we thus did not have to call on the developments of Pursuing Stacks, which serve only to establish the equivalence of the three conditions in question with the ``modelizing'' condition $(d)$. To summarize:
\begin{enumerate}[leftmargin=2.5em]
\item[$1^\circ$)] The implications
\[
(a) \ \Longrightarrow\ (\,(b) \ \Longleftrightarrow\ (c) \ \Longleftrightarrow\ (d)\,)
\]
are valid under just the conditions
\[
\mathrm{W}\,(1,\ 2\,\mathrm{bis},\ 3)\ ,
\]
and for the equivalence of $(a)$ with the other conditions, it suffices to have
\[
\mathrm{W}\,(1\,\mathrm{bis},\ 2\,\mathrm{bis},\ 3)
\]
(where W\,(1\,bis) strengthens W\,(1) by the direct factor condition, whereas W\,(2\,bis) weakens W\,(2), requiring only that $W$-aspheric morphisms be in $W$). The conclusions of 8.4, $3^\circ$, concerning the case where $A$ is $W$-aspheric, are also valid under the conditions W\,(1, 2\,bis, 3) alone.
\end{enumerate}

\msp{page 115}

\begin{enumerate}[leftmargin=2.5em]
\item[$2^\circ$)] Corollary 8.5, referring to conditions $(b)$, $(c)$, $(d)$ of $1^\circ$ (to the exclusion of $(a)$, which assumes $W$ stable under direct factors), is valid for a $W$ satisfying
\[
\mathrm{W}\,(1,\ 2\,\mathrm{bis},\ 3,\ \underbrace{7\,\mathrm{bis}}_{\substack{\text{square}\\ \text{cocartesian}}}\,)\ .
\]
% foot: Version 5 novembre 2010
\printedpage{63}
% head: A. Grothendieck : Dérivateurs, Ch. XV, Édition des Universités de Montpellier II et Paris VII
\item[$3^\circ$)] For the validity of 8.4, $2^\circ$, \textit{i.e.} to have a Quillen triple $(W_A,C,F)$, we need (almost) the whole range of conditions on $W$,
\[
\mathrm{W}\,(1\,\mathrm{bis},\ 2\,\mathrm{bis},\ 3,\ \underbrace{\ 6\ }_{\text{limits}},\ \underbrace{7\,\mathrm{bis}}_{\substack{\text{square}\\ \text{cocartesian}}},\ \underbrace{\ 8\ }_{\text{accessibility}})\ ^{\footnotemark[42]}\ ,
\]
\footnotetext[42]{\textbf{N.B.} We should renumber the axioms W\,(\ldots).}
to the exclusion of W\,(4), and being content for W\,(2) with the weakened form W\,(2\,bis) and no more. (Even the connectedness axiom W\,(5) was not used, but I had initially included it, since there is little chance that we will ever have to work with $W$'s that do not satisfy it! It might moreover be, in the presence of the other axioms, or only of W\,(1\,bis, 2\,bis, 3, 7\,bis), that if W\,(5) is \textit{not} satisfied, we have $W = \mathrm{Fl}\,\mathrm{Cat}$ in its entirety, hence $W_A = \mathrm{Fl}\,A^\wedge$!)
\end{enumerate}

\msp{page 116}

\section*{9\quad Revision of the notations, for the axioms of a basic localizer\footnote[43]{\textit{Cf.} XVII \texttt{[rather XVI]} for a definitive version (?).}}

\begin{description}[style=sameline,leftmargin=6.5em]
\item[Loc\,(1)] (Saturation).
\begin{enumerate}[leftmargin=2.5em]
\item[$a$)] Isomorphisms are in $W$.
\item[$b$)] If $f$, $g$ are composable, and two among $f$, $g$, $gf$ are in $W$, the third as well.
\item[$c$)] If $X \underset{g}{\overset{f}{\rightleftarrows}} Y$, if $gf$ and $fg$ are in $W$, then $f$ and $g$ as well.
\end{enumerate}
\end{description}
\noindent Strengthened variant:
\begin{description}[style=sameline,leftmargin=6.5em]
\item[Loc\,(1\,bis)] Loc\,(1) plus
\begin{enumerate}[leftmargin=2.5em]
\item[$d$)] $W$ is stable under direct factors.
\end{enumerate}
\item[Loc\,(2)] (Final object axiom). If $X$ in Cat has a final object, then $X$ is $W$-aspheric, \textit{i.e.} $X \longrightarrow e$ is in $W$.
\item[Loc\,(3)] (Localization axiom, weak form). Let $f : X \longrightarrow Y$ be such that for every $y \in Y$, $X/y$ is $W$-aspheric (\textit{i.e.} the morphism $f/y : X/y \longrightarrow Y/y$ deduced from $f$ by base change $Y/y \longrightarrow Y$ is in $W$). Then $f \in W$.
\end{description}
\noindent Strengthened variant of Loc\,(3):
\begin{description}[style=sameline,leftmargin=6.5em]
\item[Loc\,(3\,bis)] (Localization axiom, strong form). Let $f : X \longrightarrow Y$ be a morphism in $\mathrm{Cat}/S$, such that $\forall\ s \in S$, $f/s : X/s \longrightarrow Y/s$ (deduced from $f$ by base change $S/s \longrightarrow S$) is in $W$. Then $f \in W$. (\textbf{N.B.} Unless I am mistaken, it suffices to require it if $X$, $Y$ are split Cat-fibered over $S$, and $f$ $S$-cartesian compatible with the splitting – assuming Loc\,(1, 2)). % [sic] il suffit le l'exiger -> il suffit de l'exiger, "le" for "de"
\end{description}
% foot: Version 5 novembre 2010


\endgroup
% END UNIT Der-ChXV-05-11-10/Der-ChXV-05-11-10_p061-065.tex
% BEGIN UNIT Der-ChXV-05-11-10/Der-ChXV-05-11-10_p066-070.tex
\begingroup\makeatletter%
\newcommand{\printedpage}[1]{\marginpar{\footnotesize p.~#1}}%
\newcommand{\msp}[1]{\marginpar{\footnotesize [#1]}}%
\newcommand{\cercle}[1]{\tikz[baseline=(c.base)]\node[draw,circle,inner sep=0.8pt](c){#1};}%
\newcommand{\Cat}{\mathrm{Cat}}%
\newcommand{\id}{\mathrm{id}}%
\newcommand{\inj}{\mathrm{inj}}%
\newcommand{\Fib}{\mathrm{Fib}}%
\newcommand{\Fl}{\mathrm{Fl}}%
\makeatother\let\printedpage\derpp\let\msp\dermsp%
\let\mspin\DERmspin%


\printedpage{64}
% head: A. Grothendieck : Dérivateurs, Ch. XV, Édition des Universités de Montpellier II et Paris VII
% foot: Version 5 novembre 2010

\msp{page 117}

\begin{description}[labelwidth=5.5em,leftmargin=6.5em,style=sameline]
\item[Loc\,(4)] (Cocartesian square axiom). Let $X$ be in Cat, $X_1$, $X_2$ two opens with union $X$, $X_0 = X_1\cap X_2$, whence a cocartesian diagram in Cat,
\[
\begin{tikzcd}
X_0 \arrow[r, "i_1"] \arrow[d, "i_2"'] & X_1 \arrow[d, "j_2"] \\
X_2 \arrow[r, "j_1"] & X\ .
\end{tikzcd}
\]
We then have $i_1\in W \iff j_1\in W$ (and symmetrically $i_2\in W \iff j_2\in W$).

\item[Loc\,(6)]\footnote[44]{There will be a more ``basic'' Loc\,(5) below.} (Axiom of $W$-fibrations). There are variants \cercle{a}, (a$'$), \cercle{b}, (b$'$), (b$''$), (c), (c$'$), \textit{cf.} XII, p. 53-59. The strongest form \cercle{a} is this:

Let $f : X\longrightarrow Y$ in Cat be a $W$-fibration in the old sense, \textit{i.e.} such that for every $Y''\longrightarrow Y'$ in $\Cat/Y$, with $Y'$, $Y''$ contractible, the arrow $X''\longrightarrow X'$ obtained by the base change $X\longrightarrow Y$ is in $W$. Then $f$ is in $\Fib_W$, \textit{i.e.} the base change functor by $f$, $\Cat/Y\longrightarrow\Cat/X$, transforms arrows in $W$ into arrows in $W$.

\item[Loc\,(7)] (Direct limit axiom). $W$ is stable under filtering direct limits.
\end{description}

Strengthened form of Loc\,(7):

\begin{description}[labelwidth=5.5em,leftmargin=6.5em,style=sameline]
\item[Loc\,(7\,bis)] (L-accessibility). $W$ is stable under filtering $\varinjlim$ \textit{and} accessible, \textit{i.e.} there exists a small subset $W_\pi$ of $\Fl(\Cat)$, such that $W$ is the smallest subset of $\Fl(\Cat)$ stable under filtering $\varinjlim$, containing $W_\pi$.
\end{description}

\msp{page 118}

\begin{description}[labelwidth=5.5em,leftmargin=6.5em,style=sameline]
\item[Loc\,(8)]\footnote[45]{I have changed the numbering, to take account of a Loc\,(5) below, which is again a stability axiom \texttt{[originally the axioms Loc\,(6) to Loc\,(8) were denoted Loc\,(5) to Loc\,(7)]}.} (Connectedness axiom). $f\in W\Longrightarrow \pi_0(f)$ bijective.
\end{description}

\textbf{N.B.} All the axioms, except the last (and a smidgen of the next to last) are stability axioms – they all say that $W$ is ``large'': it contains certain arrows, and if it contains such and such arrows, it contains such and such others. Loc\,(8), and to a certain extent Loc\,(7\,bis), are ``limiting'' axioms on the size of $W$: Loc\,(8) tells us that $W\subset\mathbf{W}_0$ (= set % [sic] Loc (7) -> Loc (8), apparently: the inclusion in W_0 is the axiome de connexité Loc (8)
\DERmark{ENV0914}

of $f\in\Fl\,\Cat$ such that $\pi_0(f)$ is bijective), Loc\,(7\,bis) tells us of the existence of a cardinal such that $W$ is equal to the L-envelope of $W^\pi$ (hence $W\subset\text{L-env}(W^\pi)$).

I have tried to arrange the axioms in order of practical importance. I ended up adding the connectedness axiom at the end, for in the theoretical developments on general localizers, % [sic] théorique -> théoriques, agreement with "développements"
it practically never comes in. But I recall (XII, p. 20, th. 2) that if $W$ satisfies Loc\,(1, 2, 3, 6), and if Loc\,(8) is \textit{not} satisfied, then $W=\Fl(\Cat)$\,; thus in the presence of Loc\,(1, 2, 3, 6), Loc\,(8) \textit{is equivalent} to $W\neq\Fl(\Cat)$ in its entirety.

\msp{page 119}

The axioms that were used for 8.4, 1°, 3° and 8.5, $(b)$ are
\[
\mathrm{Loc}\,(1\,\mathrm{bis},\ 2,\ 3)\ .
\]

\printedpage{65}
% head: A. Grothendieck : Dérivateurs, Ch. XV, Édition des Universités de Montpellier II et Paris VII
% foot: Version 5 novembre 2010

Those which are necessary for 8.5, $(a)$ are
\[
\mathrm{Loc}\,(1\,\mathrm{bis},\ 2,\ 3,\ 4)\ .
\]
Finally, those necessary for 8.4, 2° are
\[
\mathrm{Loc}\,(1\,\mathrm{bis},\ 2,\ 3,\ 4,\ 7\,\mathrm{bis})\ ,
\]
and they imply the \textit{strong saturation of $W$} (\textit{cf.} next paragraph).

I must repair an omission, having somewhat treated with contempt (!) the questions of connected components, up to now.

\textbf{Proposition 9.1.} \textit{Suppose that $W$ satisfies} Loc\,(1, 2, 3, 4) \textit{and no more. Let $f_i : X_i\longrightarrow Y_i$ $(i\in I)$ be a} finite \textit{family of arrows of} Cat. \textit{For $f=\coprod f_i : \coprod X_i\longrightarrow\coprod Y_i$ to be in $W$, it suffices that $\forall\ i\in I$, we have $f_i\in W$. If $W$ also satisfies} Loc\,(7) (\textit{direct limit axiom}), \textit{then this conclusion remains valid without a finiteness condition on $I$.}

\textsc{Proof}. The case of finite $I$ is handled by induction on $n=\operatorname{card}I$, we are immediately reduced to the case where $n=2$\footnote[46]{The infinite case reduces to this under the hypothesis Loc\,(7), by passage to the limit.}. We first see the

\msp{page 120}

\textbf{Corollary 9.2.} \textit{Under the hypothesis} Loc\,(1, 4), \textit{let $f : X\longrightarrow Y$ be an arrow of} Cat, \textit{and $Z$ an object of} Cat. \textit{Then $f\in W$ if and only if $f\amalg\id_Z : X\amalg Z\longrightarrow Y\amalg Z$ is in $W$.}

Immediate from the cocartesian square, where $\inj_0^X$ is an open Dwyer immersion,
\[
\begin{tikzcd}[column sep=large, row sep=large]
X \arrow[r, "f"] \arrow[d, "\inj_0^X"'] & Y \arrow[d, "\inj_0^Y"] \\
X\amalg Z \arrow[r, "f'=f\amalg\id_Z"'] & Y\amalg Z\ ,
\end{tikzcd}
\]
which shows that $f\in W\iff f'\in W$.

By iterating this result, we find, if
\[
f_1\ :\ X_1\longrightarrow Y_1\ ,\qquad f_2\ :\ X_2\longrightarrow Y_2
\]
are given, that
\[
f_1\ \in\ W\ ,\ f_2\ \in\ W\quad\Longrightarrow\quad f_1\amalg f_2\ \in\ W\ ,
\]
by looking at the diagram of two cocartesian squares:
\[
\begin{tikzcd}[column sep=large, row sep=large]
X_1 \arrow[r, "f_1"] \arrow[d] & Y_1 \arrow[d] \\ % [sic] f -> f_1, the arrow X_1 -> Y_1 is f_1 (cf. f'_1 = f_1 ∐ id_{X_2})
X_1\amalg X_2 \arrow[r, "f'_1=f_1\amalg\id_{X_2}"'] \arrow[d, "f'_2=\id_{X_1}\amalg f_2"'] & Y_1\amalg X_2 \arrow[d, "g_2=\id_{Y_1}\amalg f_2"] \\
X_1\amalg Y_2 \arrow[r, "g_1=f_1\amalg\id_{Y_2}"'] & Y_1\amalg Y_2\ .
\end{tikzcd}
\]
\DERmark{ENV0915}


\printedpage{66}
% head: A. Grothendieck : Dérivateurs, Ch. XV, Édition des Universités de Montpellier II et Paris VII
% foot: Version 5 novembre 2010

If $f_1$ and $f_2$ are in $W$, the same is true of $f'_1$, $f'_2$, $g_1$, $g_2$, by 9.2, hence also of $f_1\amalg f_2=g_2f'_1$.

\msp{page 121}

But I see no way of proving the converse, that $f_1\amalg f_2\in W\Longrightarrow f_1,f_2\in W$. Thanks to what has been proved, we are reduced to the case where $f_1$, $f_2$ are open Dwyer immersions, but it does not seem that this gets us a hair's breadth further\,! In short, we need to know that the canonical maps $Y_i\hookrightarrow\coprod Y_i$ are in $\Fib_W$. Since $Y_i$ is a direct summand, we are reduced to the case of a canonical immersion $Y_1\hookrightarrow Y_1\amalg Y_2=Y$: if $f : X\longrightarrow Y$ is in $W$, is the same true of $X_1=X\times_Y Y_1\longrightarrow Y_1$\,? This is a property which, of course, would be a consequence of the old W\,(4\,a) (XII, p. 53), or just W\,(4\,a$'$), hence of the current Loc\,(6\,a$'$). Even the form W\,(4\,c$'$) $=$ Loc\,(6\,c$'$), the weakest of all, suffices, as we see by factoring $X\longrightarrow Y$ as $X\xrightarrow{\ i\ }X'\xrightarrow{\ p\ }Y$, with $i\in W$ (more precisely, $i_1\ :\ X_1\longrightarrow X'_1$ and $i_2\ :\ X_2\longrightarrow X'_2$ in $W$\,!), and $p\ \in\ \mathrm{Parf}_W$. By Loc\,(6\,c$'$), $f\in W$, hence $p\in W$ implies that the fibers of $p$ are $W$-aspheric, which implies that $p$ is in $W_{\mathrm{univ}}$, hence $p_1 : X'_1\longrightarrow Y_1$ is in $W$.
\DERmark{ENV0916}


\msp{page 122}

So let us note for conscience's sake:

\textbf{Corollary 9.3.} \textit{If $W$ satisfies} Loc\,(1, 2, 3, 6\,c$'$) (\textit{where} Loc\,(6\,c$'$) \textit{is the weakest of all the variants of the fibration axiom}), \textit{then for any family $(f_i : X_i\longrightarrow Y_i)_{i\in I}$ of arrows in} Cat, \textit{$f=\coprod f_i\in W\Longrightarrow$ the $f_i$ are in $W$ $\forall\ i\in I$}\footnote[47]{In fact, I do not know how to prove the implication in the reverse direction, assuming only Loc\,(1, 2, 3, 6\,c$'$).}. \textit{And if $i : Y_0\hookrightarrow Y$ is an arrow which is both an open immersion and a closed immersion} (i.e. \textit{an isomorphism with a direct summand of $Y$}), \textit{then $i$ is in} $\Fib_W$, i.e. \textit{if $f : X\longrightarrow X'$ is an arrow in} $\Cat/Y$, \textit{then $f\in W$ implies $(f_0 : X_0\longrightarrow X'_0)\in W$, where $X_0=X|Y_0$, $X'_0=X'|Y_0$.} % [sic] f_0 : X_0 -> Y_0 -> f_0 : X_0 -> X'_0, apparently: f is a flèche X -> X' and X'_0 is defined but otherwise unused
\DERmark{ENV0917}


But it is rather scandalous to have to use such a strong axiom, to conclude from it such an innocuous result\,! I note this:

\textbf{Proposition 9.4.} \textit{Let $W\subset\Fl(\Cat)$, satisfying} Loc\,(1) \textit{and no more. The following conditions are equivalent} (\textit{and implied by} Loc\,(1, 2, 3, 4, 6\,c$'$)).
\begin{enumerate}[label={}, leftmargin=2.5em]
\item[a)] \textit{If $f_1 : X_1\longrightarrow Y_1$, $f_2 : X_2\longrightarrow Y_2$ are two arrows of} Cat, \textit{then}
\[
f_1\amalg f_2\in W\iff f_1\in W,\ f_2\in W\,.
\]
\item[b)] \textit{The following two conditions on $W$ are satisfied:}

\msp{page 123}

\begin{enumerate}[label={}, leftmargin=2.5em]
\item[1°)] \textit{If $f\ :\ X\longrightarrow Y$ is in $W$, then it is so couniversally under $\emptyset$,} i.e. \textit{for every $Z$ in} Cat, \textit{$f\amalg\id_Z : X\amalg Z\longrightarrow Y\amalg Z$ is in $W$.} (\textbf{N.B.} \textit{This is equivalent to: every finite sum of arrows $f_i\in W$ is in $W$ – and likewise even for infinite sums, if} Loc\,(7), \textit{the limit axiom, is satisfied.})
\item[2°)] \textit{Let $Y_0\hookrightarrow Y$ be an inclusion of a direct summand of an object of} Cat. \textit{Then if $X\xrightarrow{\ f\ }Y$ is in $W$, so is $X_0=X|Y_0\longrightarrow Y_0$.} (\textbf{N.B.} \textit{This is equivalent to the condition that the inclusion of a direct summand is always a $W$-fibration.})
\end{enumerate}

\textit{Moreover, if $\pi_0(Y)$ is finite, or if $W$ satisfies the limit axiom} Loc\,(7), \textit{it suffices to require 2° for $Y_0$ a connected component of $Y$.}

\printedpage{67}
% head: A. Grothendieck : Dérivateurs, Ch. XV, Édition des Universités de Montpellier II et Paris VII
% foot: Version 5 novembre 2010

\item[c)] \textit{For every} finite \textit{family $(f_i)_{i\in I}$ of arrows $f_i : X_i\longrightarrow Y_i$ in} Cat,
\[
\coprod f_i\in W\iff f_i\ \in\ W\ \ \textit{for all}\ i\in I.
\]
\end{enumerate}
\textit{Moreover, if $W$ satisfies} Loc\,(7), \textit{then it is unnecessary in} (c) \textit{to restrict ourselves to finite $I$.} (\textit{And in any case, the equivalent conditions} (a), (b), (c) \textit{imply for} every \textit{family of arrows $(f_i)_{i\in I}$, $\coprod f_i\in W\Longrightarrow f_i\in W$ $\forall\ i\in I$.})

\msp{page 124}

According to 9.3, these conditions are satisfied assuming Loc\,(1, 2, 3, 4, 6\,c$'$), but Loc\,(6\,c$'$) is a steam hammer for a fly\,! One would rather like to add an axiom to the axioms Loc\,(1) to Loc\,(4) and Loc\,(6) to Loc\,(8), which together with Loc\,(8) would be the most innocuous axiom of all.

\begin{description}[labelwidth=5.5em,leftmargin=6.5em,style=sameline]
\item[Loc\,(5)] (Decomposition axiom, weak form). If $(f_i)_{i\in I}$ is a finite family of arrows of Cat, then $f=\coprod f_i$ is in $W$ if and only if the $f_i$ are. (These are the equivalent conditions of 9.4.)\footnote[48]{It may also be convenient to introduce the still weaker condition:
\begin{description}[nosep,labelwidth=4.5em,leftmargin=5em,style=sameline]
\item[Loc\,(5\,a)] If $f_1$, $f_2$ are in $W$, then $f_1\amalg f_2\in W$.\\
This is also condition $(b)$, 1° of proposition 9.4.
\item[Loc\,(5\,b)] If $f_1\amalg f_2$ is in $W$, then $f_1$ and $f_2$ too.\\
$\iff$ If $Y_0\subset Y$ is a direct summand, then $Y_0\hookrightarrow Y$ is in $\Fib_W$.
\end{description}}

\item[Loc\,(5\,bis)] (Decomposition axiom, strong form). Like Loc\,(5), but without assuming $I$ finite.\footnote[49]{\texttt{[Let us rewrite:]}
\begin{description}[nosep,labelwidth=4.5em,leftmargin=5em,style=sameline]
\item[Loc\,(5)] \textit{a}) If $f_1,f_2\in W$, then $f_1\amalg f_2\in W$.\\
\textit{b}) If $f_1\amalg f_2\in W$, then $f_1,f_2\in W$\\
($\iff b'$) every direct summand $Y_0\hookrightarrow Y$ is in $\Fib_W$\,).
\item[Loc\,(5\,bis)] Like Loc\,(5), but $(a)$ replaced by ($a$\,bis) a family (not necessarily finite).
\end{description}}
\end{description}

Thus, assuming the limit axiom Loc\,(7), the two preceding axioms are equivalent. Moreover, let us note this:

\textbf{Corollary 9.5} (of proposition 9.4). \textit{The axiom} Loc\,(5\,bis) \textit{is equivalent to each of the following properties:}
\begin{enumerate}[label={}, leftmargin=3em]
\item[\cercle{a}] \textit{For every $Y$ of} Cat, \textit{decomposed into its connected components $Y_i$ $(i\in I=\pi_0(Y))$, and for every arrow $f\ :\ X\longrightarrow Y$ in} Cat, \textit{we have $f\ \in\ W$ if and only if the $f_i : X_i=X|Y_i\longrightarrow Y_i$ are in $W$.}
\item[\cercle{b}] \textit{For $Y$, $Y_i$ as in} \cercle{a}, \textit{for every}

\msp{page 125}

\textit{arrow $f\ :\ X\longrightarrow X'$ in $\Cat/Y$, we have $f\ \in\ W$ if and only if $\forall\ i\ \in\ I$, the morphism $f_i : X_i\longrightarrow X'_i$ over $Y_i$ is in $W$.}
\end{enumerate}

\textbf{Corollary 9.6.} \textit{If the axiom} Loc\,(5) \textit{is satisfied, then the connectedness axiom} Loc\,(8) \textit{is equivalent to this: for every arrow $f\ :\ X\longrightarrow Y$ in $W$, with $Y$ $0$-connected, $X$ is $0$-connected.}

\printedpage{68}
% head: A. Grothendieck : Dérivateurs, Ch. XV, Édition des Universités de Montpellier II et Paris VII
% foot: Version 5 novembre 2010

N.B. The axioms Loc\,(8) (connectedness) and Loc\,(5\,bis) (of decomposition) seem to be of a similar nature, both being linked to the decomposition of the objects of Cat into connected components. (Nevertheless, I think they are independent.) I presume that we will never really have to work with $W$ which do not satisfy these two axioms\,! (In addition to Loc\,(1, 2, 3), which for their part constitute the ``vital minimum''.) But apparently there is no need for them in the general theoretical developments so far. (But the connectedness axiom Loc\,(8) could play a role, when one speaks of the reciprocal determination of the structures of contractibility, and of $W$-asphericity, \textit{cf.} Pursuing Stacks.) % [sic] "(Mais l'axiome ... Pursuing Stacks." -> closing parenthesis missing, the parenthesis opened at "(Mais" is never closed
\DERmark{ENV0918}


\msp{page 126}

Let us note this:

\textbf{Proposition 9.7.}\footnote[50]{\textit{Cf.} stronger version in 9.9 below (rid of the hypothesis Loc\,(5).)} \textit{Suppose that $W$ satisfies} Loc\,(1, 2, 3, 5), i.e. \textit{it is a basic localizer satisfying the decomposition condition} (\textit{weak form}). \textit{If the axiom} Loc\,(8) \textit{of connectedness is} not \textit{satisfied, then the subcategory} $\mathrm{Hot}^*_W$ \textit{of the $X\ \in$} Cat \textit{which are nonempty is equivalent to the one-point category} (\textit{hence} $\mathrm{Hot}_W$ \textit{is equivalent either to the one-point category, if $\emptyset\longrightarrow e$ is an isomorphism in} $\mathrm{Hot}_W$, \textit{or to $\boldsymbol{\Delta}_1$ otherwise}\footnote[51]{The first case is also the one where $\exists\ f\in W$ with $\pi_0(f)$ \textit{not surjective,} or again $\emptyset\longrightarrow X$ in $W$, with $X$ nonempty.}.)

This is a variant of th. 2 of XII (p. 20), where we assumed Loc\,(6\,c$'$), stronger than Loc\,(5) and even than Loc\,(5\,bis), but with the stronger conclusion that $W=\Fl(\Cat)$.

\textsc{Proof}. Let
\[
f\ :\ X\longrightarrow Y\ ,\qquad f\ \in\ W\ ,
\]
be such that $\pi_0(f)$ is not \textit{surjective.} If $Y_0$ is a connected component which is not in $\mathrm{Im}\pi_0(f)$, then $X_0=X|Y_0$ is empty, and by Loc\,(5), $\emptyset\longrightarrow Y_0$ is in $W$. Thus $\forall\ Z$ in Cat, % [sic] Loc (7) -> Loc (5), apparently: the step X_0 -> Y_0 in W uses the decomposition hypothesis Loc (5) of 9.7, not Loc (7)
\DERmark{ENV0919}

$\underbrace{\emptyset\times Z}_{=\,\emptyset}\longrightarrow Y_0\times Z$ is in $W$, and for every arrow $g : Z\longrightarrow Z'$, $\id_{Y_0}\times g : Y_0\times Z\longrightarrow Y_0\times Z'$ is in $W$. Taking $g$ such that $\pi_0(g)$ is not injective, e.g. $e\amalg e\longrightarrow e$, we find an arrow $\varphi=\id_{Y_0}\times g$ in $W$ such that $\pi_0(\varphi)$ is not injective.

\msp{page 127}

Let then $\varphi\ :\ I\longrightarrow J$ be in $W$, with $\pi_0(\varphi)$ not injective. Proceeding as before using Loc\,(5), we are reduced to the case where in addition $J$ is $0$-connected. (\textbf{N.B.} In both cases, we could also have invoked 9.6.) Let us choose two elements $i_0$, $i_1$ in $I$ in two distinct connected components, let $j_0$, $j_1$ be their images in $J$. We are going to treat $(I,i_0,i_1)$ and $(J,j_0,j_1)$ as ``homotopy intervals''. The fact that $i_0$, $i_1$ belong to distinct components of $I$ implies that two arrows $f_0,f_1 : X\rightrightarrows Y$ are ``homotopic with respect to $(I,i_0,i_1)$'', \textit{i.e.} there exists $h\ :\ I\times X\longrightarrow Y$ in Cat such that $f_0\ =\ h\partial_0^X$, $f_1=h\partial_1^X$, where
\DERmark{ENV0920}

\[
\partial_0^X\ ,\ \partial_1^X\ :\ X\ \longrightarrow\ I\times X
\]
are given by $x\longmapsto(i_\alpha,x)$ $(\alpha\in\{0,1\})$. To find $h$, we decompose $I$ as
\[
I\ =\ I_0\amalg I_1\ ,\qquad\quad \text{with}\quad i_0\in I_0,\ i_1\in I_1\ ,
\]
whence $I\times X\ \simeq\ (I_0\times X)\amalg(I_1\times X)$, and we define $h$ by its restrictions $h_0$, $h_1$ to the two pieces $I_0\times X$, $I_1\times X$, namely
\[
h_0\ =\ f_0\circ\mathrm{pr}_2^{I_0}\ ,\qquad\quad h_1\ =\ f_1\circ\mathrm{pr}_2^{I_1}\ ,
\]


\endgroup
% END UNIT Der-ChXV-05-11-10/Der-ChXV-05-11-10_p066-070.tex
% BEGIN UNIT Der-ChXV-05-11-10/Der-ChXV-05-11-10_p071-075.tex
\begingroup\makeatletter%
\newcommand{\printedpage}[1]{\marginpar{\footnotesize p.~#1}}%
\newcommand{\msp}[1]{\marginpar{\footnotesize [#1]}}%
\newcommand{\marginbar}{\marginpar{$\Vert$}}%
\newcommand{\Hot}{\mathrm{Hot}}%
\newcommand{\Hom}{\mathrm{Hom}}%
\newcommand{\Cat}{\mathrm{Cat}}%
\newcommand{\hot}{\mathrm{hot}}%
\newcommand{\Ob}{\mathrm{Ob}}%
\newcommand{\card}{\operatorname{card}}%
\newcommand{\pr}{\mathrm{pr}}%
\newcommand{\id}{\mathrm{id}}%
\makeatletter%
\renewcommand\@makefnmark{\hbox{\@textsuperscript{\normalfont(\@thefnmark)}}}%
\renewcommand\@makefntext[1]{\parindent 1em\noindent\hbox{\@textsuperscript{\normalfont\@thefnmark}}\,#1}%
\makeatother%
\makeatother\let\printedpage\derpp\let\msp\dermsp%
\let\mspin\DERmspin%


\printedpage{69}
% head: A. Grothendieck : Dérivateurs, Ch. XV, Édition des Universités de Montpellier II et Paris VII
% foot: Version 5 novembre 2010
where
\[
\pr_2^{I_\alpha}\ :\ I_\alpha\times X \longrightarrow X\ .
\]

\msp{page 128}Let us now consider the diagram (solid lines)
\[
\begin{tikzcd}[column sep=small]
& X\amalg X \arrow[dl, "{(\partial_0^X,\partial_1^X)}"'] \arrow[dr, "{(\partial_0^{\prime X},\partial_1^{\prime X})}"] & \\
I\times X \arrow[rr, "\varphi\times\id_X\,=\,\varphi_X"] \arrow[dr, "h"'] & & J\times X \arrow[dl, dotted, "h'"] \\
& Y\ , &
\end{tikzcd}
\]
where $\partial_0^{\prime X}$, $\partial_1^{\prime X}$ are defined as $\partial_0^X$, $\partial_1^X$, but using $j_0,j_1\in\Ob J$ instead of $i_0,i_1\in\Ob I$. If there existed in Cat an arrow $h'$ (dotted) such that $h'\circ\varphi_X=h$, it would follow that $h'$ is a homotopy of $f_0$ and $f_1$ along $(J,j_0,j_1)$. Since $J$ is 0-connected, it would follow at once that $f_0$, $f_1$ are homotopic in the usual sense in Cat, hence that $f_0$ and $f_1$ have the same image in $\Hot_W$. But if one cannot find an arrow $h'$ in Cat, at least one can find it in $\Hot_W$, since $\varphi_X\in W$, hence an isomorphism in $\Hot_W$ (and $h'$ in $\Hom_W(J\times X,Y)$ will even be uniquely determined). Thus in $\Hot_W$, we have
\[
f_0\ =\ h'\partial_0^{\prime X}\ ,\qquad\quad f_1\ =\ h'\partial_1^{\prime X}\ .
\]

\msp{page 129}To prove that
\[
f_0\ =\ f_1\qquad\text{in}\quad\Hom_W(X,Y)\ ,
\]
it suffices to see that
\[
\partial_0^{\prime X}\ =\ \partial_1^{\prime X}\qquad\text{in}\quad\Hom_W(X,J\times X)\ .
\]
Now in fact $\partial_0^{\prime X}$ and $\partial_1^{\prime X}$ are homotopic, for $e\underset{j_1}{\overset{j_0}{\rightrightarrows}}J$ are homotopic ($j_0$ and $j_1$ belonging to the same connected component of $J$), hence also their product with $\id_X$.

One sees therefore that for any two objects $X$, $Y$ of Cat, all the arrows from $X$ to $Y$ are equal in Hot. The argument given even proves that all the arrows \textit{in} $\Hot_W$ from $X$ to $Y$ are equal. This means that the category $\Hot_W^*$ (where for any two objects $X$, $Y$, $\Hom_{\Cat}(X,Y)$, and a fortiori $\Hom_W(X,Y)$ is $\neq\emptyset$) is equivalent to the one-point category. As for $\Hot_W$ itself, whose set of objects is $\{\emptyset\}\amalg\Ob\Hot_W^*$, with $\emptyset$ an initial object, it will therefore be equivalent to the one-point category, or else to $\boldsymbol{\Delta}_1$, according as $\emptyset\longrightarrow e$ is an isomorphism in $\Hot_W$, or not.

\msp{page 130}\textbf{Comment 9.8.} If the axioms Loc\,(1, 2, 3, 5)~\footnote[52]{In fact, this comment holds even in the absence of the assumption Loc\,(5), \textit{cf.} 9.9 below.} are satisfied, then the connectedness axiom Loc\,(8) \textit{is equivalent} to each of the following:

\printedpage{70}
% head: A. Grothendieck : Dérivateurs, Ch. XV, Édition des Universités de Montpellier II et Paris VII
% foot: Version 5 novembre 2010
\begin{enumerate}[leftmargin=*]
\item[\textit{a})] $\Hot_W$ is equivalent neither to $e$, nor to $\boldsymbol{\Delta}_1$.
\item[\textit{b})] $\Hot_W$ has infinitely many isomorphism classes.
\item[\textit{c})] $f_0:X_0\longrightarrow Y_0$ (arrow given in advance in Cat, such that $\pi_0(f_0)$ is not an isomorphism, and $X_0\neq\emptyset$) is \textit{not} an isomorphism in $\Hot_W$.
\item[\textit{d})] $X_0$ and $Y_0$ are not isomorphic in $\Hot_W$ (where $X_0$, $Y_0$ are two objects given in advance in Cat, with $X_0$, $Y_0$ nonempty and $\card\pi_0(X_0)\neq\card\pi_0(Y_0)$).
\end{enumerate}

% marginal mark: a double vertical rule in the left margin beside the first two lines of this paragraph
\marginbar One can therefore say that the connectedness axiom plays the role of a \textit{nontriviality axiom} for the localized category $\Hot_W$ of Cat. It will be satisfied automatically for all the basic localizers $W$ that we will be led to construct or to introduce. But we will practically never have to invoke it as an assumption, with a view to certain conclusions.

I finally realize that in proposition 9.7 the assumption Loc\,(5) is unnecessary:

\textbf{Proposition 9.9.}~\footnote[53]{\textit{Cf.} better statement 9.12 (p. 133 bis).} \textit{Let $W$ be a basic localizer in} Cat (i.e. \textit{we have} Loc\,(1, 2, 3) \textit{and nothing more}), \textit{plus the condition that the} sum \textit{of $f_1,f_2\in W$ is in $W$}~\footnote[54]{This condition Loc\,(5\,a) holds in particular if $W$ also satisfies one of the two axioms Loc\,(4) (\textit{cf.} 9.1), or Loc\,(6\,c$'$) (\textit{cf.} 9.3).}. \textit{If the connectedness axiom is not satisfied, then} $\Hot_W$ \textit{is equivalent to the one-point category} (\textit{if $\emptyset\longrightarrow e$ is an isomorphism in} $\Hot_W$~\footnote[55]{This also means that there exists $Z\neq\emptyset$ in Cat, such that $\emptyset\longrightarrow Z$ is an isomorphism in $\Hot_W$.}) \textit{or to $\boldsymbol{\Delta}_1$} (\textit{in the opposite case}), \textit{and} $\Hot_W^*$ ($\overset{\mathrm{def}}{=}\Hot_W\backslash\{\emptyset\}$) \textit{is equivalent to the one-point category.}
\DERmark{ENV0921}


\msp{page 131}The condition on sums ensures that in $\Hot_W$ the sum of two objects exists, and that $\Cat\longrightarrow\Hot_W$ commutes with sums of two objects. This said, let
\[
f\ :\ X\longrightarrow Y \tag*{$(*)$} % the label (*) is printed at the left margin
\]
be an arrow such that $f\in W$, but $\pi_0(f)$ not surjective. Let us first suppose $X\neq\emptyset$ (the least obvious case). If $Y_0$ is a nonempty summand of $Y$ such that $f$ factors through $Y_1=Y\backslash Y_0$, then $f$ is the sum of
\[
f_1\ :\ X\ =\ X_1\longrightarrow Y_1\qquad\quad\text{and}\qquad\quad f_0\ :\ X_0\ =\ \emptyset\longrightarrow Y_0
\]
in Cat, hence also in $\Hot_W$. Thus we have a sum $f\ =\ f_0\amalg f_1$ in $\Hot_W$ which is an isomorphism, hence for every $Z$ in $\Hot_W$,
\[
\Hom_W(Y_0,Z)\times\Hom_W(Y_1,Z)\ \xrightarrow{\ \sim\ }\ \underbrace{\Hom_W(X_0,Z)}_{=\,e}\times\Hom_W(X_1,Z)
\]
(\textbf{N.B.} $\emptyset$ is an initial object of $\Hot_W$, being so in Cat), \textit{i.e.} the composite
\[
Z(Y_0)\times Z(Y_1)\ \xrightarrow{\pr_2}\ Z(Y_1)\ \xrightarrow{Z(f_1)}\ Z(X_1)
\]
is bijective. Since $X\ =\ X_1\neq\emptyset$, hence also $Y_1\neq\emptyset$, and $Y_0\neq\emptyset$, the sets $Z(Y_0)$, $Z(Y_1)$, $Z(X_1)$ are $\neq\emptyset$ if $Z\neq\emptyset_{\Cat}$. The composite being bijective, the first arrow $\pr_2$ is injective, which implies
\printedpage{71}
% head: A. Grothendieck : Dérivateurs, Ch. XV, Édition des Universités de Montpellier II et Paris VII
% foot: Version 5 novembre 2010
\msp{page 132}that $Z(Y_0)$ is reduced to a point (the two factors $Z(Y_0)$, $Z(Y_1)$ being $\neq\emptyset$). This means that $Y_0$ is an initial object in $\Hot_W^*$. It follows in particular that $Y_0\amalg Y_0$ is also an initial object, and the canonical map $Y_0\amalg Y_0\longrightarrow Y_0$ is therefore an isomorphism in $\Hot_W$, whereas the induced map on the $\pi_0$ is not injective. (It is $\pi_0(Y_0)\amalg\pi_0(Y_0)\longrightarrow\ \pi_0(Y_0)$, with $\pi_0(Y_0)$ nonempty.) This case will be treated below.

If in $(*)$ $X$ is empty, $Y$ nonempty (since $\pi_0(f)$ not surjective), then $Y$ is an initial object of $\Hot_W$, and one sees again that $Y\amalg Y\longrightarrow Y$ is an isomorphism in $\Hot_W$, although not inducing an isomorphism of the $\pi_0$. Proposition 9.9 will therefore follow from the

\textbf{Lemma 9.10.} \textit{Let $W$ be a basic localizer and nothing more} (\textit{axioms} Loc\,(1, 2, 3)). \textit{Suppose there exists $f:I\longrightarrow J$ in} Cat \textit{such that} $\hot_W(f)$ \textit{is an isomorphism, but $\pi_0(f)$} is not injective~\footnote[56]{Or else such that $I$ is 0-connected, and $J$ not 0-connected (\textit{i.e.} $\card\pi_0(I)\ =\ 1$, $\card\pi_0(J)\ \geq\ 2$). This second case is treated by the same method essentially.}. \textit{Then} $\Hot_W^*$ \textit{is equivalent to the one-point category $e$, and} $\Hot_W$ \textit{is equivalent to $e$} (\textit{if $\emptyset\longrightarrow e$ is an isomorphism in} $\Hot_W$) \textit{or to $\boldsymbol{\Delta}_1$} (\textit{in the opposite case}).

\msp{page 133}The proof is the one given for 9.7: one proves that if $X$, $Y$ are two objects of Cat, \textit{nonempty}, then two arrows $f_0,f_1:X\rightrightarrows Y$ are equal in $\Hot_W$, by a homotopy argument, choosing two points $i_0$, $i_1$ of $I$ which are in two distinct connected components, and such that their images $j_0$, $j_1$ in $J$ are in the same component.

\textbf{Proposition 9.11.} \textit{Let $W$ be a basic localizer} (i.e. \textit{we have} Loc\,(1, 2, 3)). \textit{Then the following conditions are equivalent.}
\begin{enumerate}[leftmargin=*]
\item[a)] \textit{We have} Loc\,(5\,bis) (\textit{decomposition axiom, strong form}) \textit{and} Loc\,(8) (\textit{connectedness axiom}).
\item[b)] \textit{We have} Loc\,(5\,bis), \textit{and} $\Hot_W$ \textit{is not equivalent to $e$ or to $\boldsymbol{\Delta}_1$} (cf. \textit{comment} 9.8).
\item[c)] \textit{For every $f:X\longrightarrow Y$, we have $f\in W$ if and only if $\pi_0(f)$ is an isomorphism, and for every connected component $X_i$ of $X$, if $Y_i$ denotes the corresponding connected component of $Y$, the induced arrow $f_i:X_i\longrightarrow Y_i$ is in $W$.}
\end{enumerate}

The reason why I do not isolate this property as one of the axioms Loc\,(\ldots) is that it is \textit{not} a stability axiom, but an axiom of a mixed nature.

\msp{page 133 bis}\textbf{Proposition 9.12.} \textit{Let $W$ be a basic localizer} (\textit{nothing more}), i.e. \textit{satisfying} Loc\,(1, 2, 3). \textit{The following conditions are equivalent.}
\begin{enumerate}[leftmargin=*]
\item[a)] \textit{$W$ satisfies the connectedness axiom} Loc\,(8), i.e. $f\in W\Longrightarrow\pi_0(f)$ \textit{bijective.}
\item[b)] \textit{If $f,g:X\rightrightarrows Y$ in} Cat \textit{are such that} $\hot_W(f)=\hot_W(g)$, \textit{then} $\pi_0(f)=\pi_0(g)$.
\item[c)] \textit{There exist $X$, $Y$ in} Cat \textit{such that the map} $\hot_W^{X,Y}:\Hom_{\Cat}(X,Y)\longrightarrow\Hom_W(X,Y)$ \textit{is not constant.}
\item[c$'$)] \textit{The subcategory} $\Hot_W^*$ \textit{of} $\Hot_W$, \textit{consisting of the nonempty objects, is not equivalent to the one-point category.}
\item[c$''$)] \textit{The category} $\Hot_W$ \textit{is equivalent neither to the category $e$ nor to the category $\boldsymbol{\Delta}_1$.}
\item[d)] \textit{The set of isomorphism classes of objects of} $\Hot_W$ \textit{is infinite.}
\end{enumerate}

\printedpage{72}
% head: A. Grothendieck : Dérivateurs, Ch. XV, Édition des Universités de Montpellier II et Paris VII
% foot: Version 5 novembre 2010
\textsc{Proof}. We have $(a)\Longrightarrow(b)$, since $(a)$ is equivalent to saying that the functor $\pi_0:\Cat\longrightarrow\mathrm{Ens}$ factors through $\Hot_W$. Moreover the corresponding functor
\[
\Hot_W\ \longrightarrow\ (\mathrm{Ens})
\]
is essentially surjective -- there is a right inverse up to isomorphism, $I\longmapsto\hot_W(I)$. This proves that $(a)$ implies $(d)$, and we have $(d)\Longrightarrow(c')$, $(d)\Longrightarrow(c'')$ tautologically. A more or less tautological argument shows moreover that if $\Hot_W^*$ is equivalent to $e$, then $\Hot_W$ is equivalent to $e$ or to $\boldsymbol{\Delta}_1$ (for $\emptyset$ is an initial object of $\Hot_W$) -- hence $(c'')\Longrightarrow(c')$. Thus we have more or less tautologically
\DERmark{ENV0922}

\[
\begin{tikzcd}
(\mathrm{a}) \arrow[r, Rightarrow] \arrow[d, Rightarrow] & (\mathrm{d}) \arrow[r, Rightarrow] & (\mathrm{c}'') \arrow[r, Rightarrow] & (\mathrm{c}') \\
(\mathrm{b}) \arrow[rr, Rightarrow] & & (\mathrm{c}) \arrow[ru, Rightarrow] &
\end{tikzcd}
\]
where the implications $(b)\Longrightarrow(c)\Longrightarrow(c')$ are also trivial \texttt{[it is rather the implication (c$'$) $\Longrightarrow$ (c) that is trivial~; the implication (c) $\Longrightarrow$ (c$'$) is proved after the statement of corollary 9.13, before the proof of the latter]}. Thus
\msp{page 133 ter}everything comes down to proving that $(c)\Longrightarrow(a)$, or again:

\textbf{Corollary 9.13.} \textit{If the connectedness axiom is not satisfied, then if $X$, $Y$ are two objects of} Cat, \textit{all the arrows $X\longrightarrow Y$ in} Cat \textit{are equal in} $\Hot_W$, i.e. \textit{the map} $\Hom_{\Cat}(X,Y)\longrightarrow\Hom_W(X,Y)$ \textit{is constant.}

This therefore implies that all the objects of Cat$^*$, \textit{i.e.} all the nonempty objects of Cat, are isomorphic in $\Hot_W$, and one finds a canonical isomorphism between them, with transitivity. Let then $\mathcal{H}$ be the quotient category of Cat, having the same objects as Cat, and having exactly one arrow between two nonempty objects of Cat and exactly one arrow from $\emptyset$ to every other object of Cat (and no arrow from a nonempty object of Cat to $\emptyset$). It is therefore a category equivalent to $\boldsymbol{\Delta}_1$. We have seen that $\Cat\longrightarrow\Hot_W$ factors as
\[
\Cat\ \xrightarrow{\text{can.}}\ \mathcal{H}\ \longrightarrow\ \Hot_W\ .
\]
If $\Cat\longrightarrow\mathcal{H}$ transforms every arrow belonging to $W$ into an isomorphism, then it is immediate that $\mathcal{H}$ is a solution to the universal problem defining $\Hot_W$, hence that $\mathcal{H}\longrightarrow\Hot_W$ is an isomorphism. This also means that $W$ does not contain an arrow $\emptyset\longrightarrow X$ with $X$ nonempty (in which case $W$ contains all arrows). Thus in this case $\Hot_W\simeq\boldsymbol{\Delta}_1$. In the opposite case, let $\mathcal{H}'$ be the category obtained from $\mathcal{H}$ by inverting the arrows $\emptyset\longrightarrow X$, \textit{i.e.} the 1-connected groupoid on the set of objects $\Ob\Cat$. It is no longer a quotient category of Cat, but we still have a factorization
\[
\Cat\ \xrightarrow{\text{epi}}\ \mathcal{H}\ \ \xrightarrow{\text{loc.}}\ \mathcal{H}'\ \longrightarrow\ \Hot_W\ ,
\]
\printedpage{73}
% head: A. Grothendieck : Dérivateurs, Ch. XV, Édition des Universités de Montpellier II et Paris VII
% foot: Version 5 novembre 2010
\msp{page 133 quater}and this time it is $\mathcal{H}'$ that is a solution of the universal problem ($\Cat\longrightarrow\mathcal{H}'$ transforming arrows of $W$ into isomorphisms). Thus in this case, $\mathcal{H}'\longrightarrow\Hot_W$ is an isomorphism, and $\Hot_W$ is equivalent to $e$. In both cases, $\Hot_W^*$ is equivalent to $e$.

Thus, proposition 9.12 follows from corollary 9.13, which it is now a matter of proving.

Let $f:I\longrightarrow J$ in Cat be such that $\pi_0(f)$ is not injective. We have seen in the proof of 9.7 how one deduces from this (by choosing two elements $i_0$, $i_1$ in two distinct connected components of $I$, such that $j_0=f(i_0)$ and $j_1=f(i_1)$ are in the same connected component of $J$) that for $X$, $Y$ in Cat, two arrows $f_0,f_1:X\rightrightarrows Y$ are always equal in $\Hot_W$ (\textit{cf.} lemma 9.10).

Let now
\[
f\ :\ I\longrightarrow J
\]
in Cat be such that $\pi_0(f)$ is not surjective, hence $J=J_0\amalg J_1$, with $f$ factoring through $J_0$, and $J_1\neq\emptyset$. Let us then consider the discrete category with two objects,
\[
\varepsilon\ =\ \{0,1\}\ ,
\]
and two arrows
\[
g_0,\ g_1\ :\ J\ \rightrightarrows\ \varepsilon\ ,
\]
$g_0$ being constant with value 0, and $g_1$ being equal to 0 on $J_0$, to 1 on $J_1$. We have
\[
g_0f\ =\ g_1f\ ,
\]
hence, since $f\in W$,
\[
g_0\ =\ g_1\qquad\quad\text{in}\quad\Hot_W\ .
\]
One concludes from this, by taking a point $j_1$ in $J_1$, and the corresponding inclusion $\alpha_{j_1}:e\longrightarrow J$, that
\msp{page 133 quinquies}
\[
g_0\alpha_{j_1}\ =\ g_1\alpha_{j_1}\ ,
\]
\textit{i.e.} the two sections of $\varepsilon$ over $e$ are equal in $\Hot_W$. Since the functor $\Cat\longrightarrow\Hot_W$ commutes with finite products \texttt{[this property is not established under the sole conditions Loc\,(1,2,3)~; it is proved in Pursuing Stacks, section 64, on replacing Loc\,(3) by the stronger axiom Loc\,(3\,bis)]}, it follows that for $X$ in Cat, the two arrows
\[
X\ \underset{\partial_1^X}{\overset{\partial_0^X}{\rightrightarrows}}\ X\times\varepsilon\ =\ X\amalg X
\]
are equal in $\Hot_W$. If then
\[
f_0,\ f_1\ :\ X\ \rightrightarrows\ Y
\]
are two arrows in Cat, they are of the form
\[
f_0\ =\ h\partial_0^X\ ,\qquad\quad f_1\ =\ h\partial_1^X\ ,
\]
with
\[
h\ =\ (f_0,f_1)\ :\ X\amalg X\ \longrightarrow\ Y\ , % [sic] X\amalg X \longrightarrow X -> X\amalg X \longrightarrow Y, h\partial_i^X = f_i : X \to Y requires target Y
\]
\DERmark{ENV0923}

and since $\partial_0^X=\partial_1^X$ in $\Hot_W$, one concludes $f_0=f_1$ in $\Hot_W$, \quad \textit{q.e.d.}


\endgroup
% END UNIT Der-ChXV-05-11-10/Der-ChXV-05-11-10_p071-075.tex
% BEGIN UNIT Der-ChXV-05-11-10/Der-ChXV-05-11-10_p076-077.tex
\begingroup\makeatletter%
\newcommand{\printedpage}[1]{\marginpar{\footnotesize p.~#1}}%
\newcommand{\msp}[1]{\marginpar{\footnotesize [#1]}}%
\newcommand{\HOT}{\mathrm{HOT}}%
\newcommand{\hot}{\mathrm{hot}}%
\makeatother\let\printedpage\derpp\let\msp\dermsp%
\let\mspin\DERmspin%


\printedpage{74}
% head: A. Grothendieck : Dérivateurs, Ch. XV, Édition des Universités de Montpellier II et Paris VII
% foot: Version 5 novembre 2010

\msp{page 134}

\section*{10\quad Application to the derivator $\HOT_W$}

\textbf{Theorem 10.1.} \textit{Let $W$ be a localizer in} Cat \textit{satisfying the conditions} Loc\,(1, 2, 3, 4, 7\,bis) (\textit{axioms of strong-form saturation} (\textit{direct factors}), \textit{of the final object, of weak-form localization, of the cocartesian square, of L-accessibility}~\textsuperscript{(57)}\footnotetext[57]{\textbf{N.B.} Loc\,(1\,bis), \textit{i.e.} $W = W^{\natural}$, is already implied by the condition that $W$ be stable under filtering $\varinjlim$.}). \textit{Let} $\HOT_W$ \textit{be the prederivator on} (Cat) \textit{associated to} $(\mathrm{Cat}, W)$ \textit{in the usual way,}
\[
I \longmapsto \mathrm{Cat}(I)\,W(I)^{-1} \;=\; \underline{\mathrm{Hom}}(I^{\circ}, \mathrm{Cat})\,W(I)^{-1}\,. \tag*{(10.1.1)}
\]
\textit{One then has the following.}

1\,$^{\circ}$) \textit{For every arrow $u : I \longrightarrow J$ in} Cat, \textit{the functor $u_!$} (\textit{left adjoint to $u^*$}) \textit{exists} (\textbf{N.B.} For this one does not need Loc\,(4) or Loc\,(7\,bis), it suffices that $W$ be a basic localizer \texttt{[one nevertheless needs the strong-form localization axiom Loc\,(3\,bis)]}), \textit{and it is computed in the standard way} (\textit{axioms} Der~4\,bis \textit{and} Der~5\,bis). \textit{The functor $u_*$, right adjoint to $u^*$, exists}~\textsuperscript{(58)}\footnotetext[58]{\textit{Always\,! Cf.} XIX, p. 29.} [\textit{if $u$ is of the form $v \times \mathrm{id}_S$, where $v : I_0 \longrightarrow J_0$ is an arrow between finite ordered categories, and $S$ an arbitrary object of} Cat], \textit{and this functor $u_*$ is computed in the standard way} (\textit{axioms} Der~4 \textit{and} Der~5). \textit{Thus one has a theory of $W$-cartesian squares and of $W$-cocartesian squares in the categories} $\HOT_W(S)$, \textit{and of long sequences}
\msp{page 135}
\textit{of suspension and of cosuspension in the categories} $\HOT^{\bullet}_W(S)$. (\textit{More precisely, these are theories relative to the induced derivators} $\HOT_{W,S}$, $\HOT^{\bullet}_{W,S}$ \textit{on $S$,} $I \longmapsto \HOT_W(S \times I)$ \textit{resp.} $I \longmapsto \HOT^{\bullet}_W(S \times I)$ \ldots). \textit{Moreover, the cocartesian squares and the cartesian squares relative to the induced derivators} $\HOT_{W,S}$ \textit{are} exact, i.e. \textit{the long suspension and cosuspension sequences are exact} (\textit{axioms of right and left exactness} Der~6\,bis \textit{and} Der~6).

2\,$^{\circ}$) \textit{The prederivator} $\HOT_W$ \textit{also satisfies the axiom} Der~1 \textit{of finite decomposition}~\textsuperscript{(59)}\footnotetext[59]{Without any axiom on $W$\,! \texttt{[One only needs the identities to be in} $W$\texttt{.]}} (\textit{for} $I = \coprod_{\alpha} I_{\alpha}$), \textit{and the axiom} Der~2 \textit{of localization. Moreover, $W$ is strongly saturated,} i.e. $f \in W \iff \hot_W(f)$ \textit{isomorphism.}

3\,$^{\circ}$) \textit{Concerning the axiom} Der~3 \textit{and its variants} (\textit{cf.} I, pages 111-112, 125), \textit{one has this:}

a)~\textsuperscript{(60)}\footnotetext[60]{Here one needs only Loc\,(1\,ter, 2, 3) (Loc\,(1\,ter) is strong saturation, \textit{i.e.} $f \in W \iff \hot_W(f)$ isomorphism).}. \textit{Let $f : S' \longrightarrow S$ be an arrow in} Cat. \textit{For it to be a} $\HOT_W$\textit{-equivalence,} i.e. \textit{for the functor} $f^* : \HOT_W(S) \longrightarrow \HOT_W(S')$ \textit{to induce an} equivalence \textit{between the full subcategories formed by the objects} constant \textit{on $S$, $S'$, it is necessary and sufficient that $f$ be in $W$. The corresponding statement in the induced derivators} $\HOT_{W,I}$ \textit{is also valid, for $I \neq \emptyset$. Thus a fortiori,} $\HOT_W$ \textit{and} $\HOT_{W,I}$ \textit{satisfy}
\msp{page 136}
\printedpage{75}
% head: A. Grothendieck : Dérivateurs, Ch. XV, Édition des Universités de Montpellier II et Paris VII
% foot: Version 5 novembre 2010
\textit{the final object axiom,} i.e. \textit{if $S$ in} Cat \textit{has a final object, then $S \longrightarrow e$ is a} $\HOT_W$\textit{-equivalence, and even a} $\HOT_{W,I}$\textit{-equivalence for every $I$ in} Cat.

b) \textit{Suppose that $W$ satisfies} Loc\,(1\,ter, 2, 3, 6\,b) (\textit{thus the axiom} Loc\,(6\,b) \textit{of $W$-fibrations, variant} b). \textit{Then for every $f : S' \longrightarrow S$ in $W$,} $\HOT^{\mathrm{lc}}_W(S) \longrightarrow \HOT^{\mathrm{lc}}_W(S')$ \textit{is an} equivalence of categories (\textit{the converse being clear by} (a)). \textit{Likewise for} % [sic] clair -> claire, accord avec « réciproque »
\DERmark{ENV0924}\DERmark{ENV0925}

\[
\HOT^{\mathrm{lc}}_{W,I}(S) \longrightarrow \HOT^{\mathrm{lc}}_{W,I}(S') \qquad (\text{i.e. \textit{for}}\ \ \HOT^{\mathrm{lc}/I}_W(S \times I) \longrightarrow \HOT^{\mathrm{lc}/I}_W(S' \times I)\ ).
\]
\textit{In other words, for the derivator} $\HOT_W$ \textit{and for all the} $\HOT_{W,I}$ ($\mathbf{D}$ \textit{being one of the these derivators})~\textsuperscript{(61)}\footnotetext[61]{To be checked carefully in the case $\mathbf{D} = \HOT_{W,I}$.}, \textit{every $\mathbf{D}$-equivalence is a $\mathbf{D}$-equivalence which is} 2\textit{-complete,} i.e. \textit{for $f : S' \longrightarrow S$, for $f^* : \mathbf{D}(S) \longrightarrow \mathbf{D}(S')$ to induce an equivalence between the subcategories $\mathbf{D}^{\mathrm{lc}}(S) \longrightarrow \mathbf{D}^{\mathrm{lc}}(S')$, it already suffices that it induce a fully faithful functor on the full subcategories formed by the constant objects, and this is equivalent} (\textit{if $I \neq \emptyset$}) \textit{to the condition that $f \in W$.} % [sic] des ces dérivateurs -> de ces dérivateurs, faute de frappe

\msp{page 137}

\textbf{Comment.} In this summary statement of what I know how to prove at present for $\HOT_W$, I have somewhat mixed the various sources – what is put ``for the record'' from VI (and which does not depend on Loc\,(4, 7 or 7\,bis)), \textit{i.e.} what depends essentially on the construction of $u_!$ given in VI (thus what concerns $u_!$ in 1$^{\circ}$, \textit{minus} right exactness, and all that concerns 3$^{\circ}$, $(a)$, $(b)$).

The proof of 10.1 will follow from the developments of the following section \texttt{[which does not exist]}, taking into account the main result of section 8 (th. 8.4), and of section 7 (cor. 7.4), which are the ``hard'' ingredients of Theorem 10.1.


\endgroup
% END UNIT Der-ChXV-05-11-10/Der-ChXV-05-11-10_p076-077.tex
\clearpage\setcounter{footnote}{0}\phantomsection\addcontentsline{toc}{section}{Chapter XVI}\markboth{Chapter XVI}{Chapter XVI}
% BEGIN UNIT der-ch16/der-ch16_p001-005.tex
\begingroup\makeatletter%
\newcommand{\printedpage}[1]{\marginpar{\footnotesize p.~#1}}%
\newcommand{\msp}[1]{\marginpar{\footnotesize [#1]}}%
\newcommand{\circled}[1]{\tikz[baseline=(c.base)]\node[draw,circle,inner sep=1pt](c){#1};}%
\newcommand{\Cat}{\mathrm{Cat}}%
\newcommand{\Ob}{\mathrm{Ob}\,}%
\newcommand{\hot}{\mathrm{hot}}%
\makeatother\let\printedpage\derpp\let\msp\dermsp%
\let\mspin\DERmspin%


% ===================== page 1 (PDF) =====================
\printedpage{[unnumbered]}
% head: (none printed)

\begin{center}
{\sffamily\LARGE THE DERIVATORS}

\bigskip

{\large\textbf{ALEXANDRE GROTHENDIECK}}

\vspace{3em}

{\large\textbf{Chapter XVI}}

\medskip

{\large\textbf{The basic localizers}}
\end{center}

\par

% ===================== page 2 (PDF) =====================
\printedpage{[unnumbered]}
% head: (none printed)

\par\noindent\textbf{Note by the editors of the transcribed edition.}\par
This text was deciphered and transcribed in \LaTeXe{} by M. Künzer. It was edited by M. Künzer, J. Malgoire, and G. Maltsiniotis. The transcription is as faithful as possible to the manuscript. Only a few obvious minor corrections have been made. For the rare additions or comments by the editors, as well as for the original numbering of the manuscript pages, typewriter characters \texttt{[typewriter]} in brackets are used.

This edition is provisional. Remarks, comments and corrections are welcome. Send a message to:

\textbf{[e-mail address removed]}

\hfill G. Maltsiniotis\hspace*{2cm}

\par

% ===================== page 3 (PDF) =====================
\printedpage{1}
% head: A. Grothendieck : Dérivateurs, Ch. XVI, Édition des Universités de Montpellier II et Paris VII
% foot: Version 18 mars 2019

\msp{page 1}

\section*{1\quad Axioms Loc(1) to Loc(3) (and variants): basic localizers}

I feel like taking up again the exposition of the axioms useful for a basic localizer
\[
W \subset \mathrm{Fl}(\Cat),
\]
guiding myself this time more by the geometric meaning of the axioms, than by the $\pm$ technical aliases one encounters (or does not encounter) in checking them. I am therefore revising in this spirit the presentation of XV §9.

\begin{description}
\item[Loc (1)] (\textit{saturation}).
\begin{itemize}
\item[\textit{a})] Iso $\subset W$.
\item[\textit{b})] If two among $f$, $g$, $gf$ are in $W$, the third too.
\item[\textit{c})] If $gf$ and $fg$ are in $W$, $f$ and $g$ too.
\end{itemize}
\end{description}

The conditions (\textit{a}) and (\textit{b}) alone, I am going to call them (for a general $W \subset \mathrm{Fl}\,\mathcal{M}$) \textit{weak saturation} of the localizer. The conditions (\textit{a}), (\textit{b}), (\textit{c}), \textit{mild saturation} (``mild saturation''). Other variants:

\begin{description}
\item[Loc (1 bis)] The conditions (\textit{a}), (\textit{b}), (\textit{c}) above, plus
\begin{itemize}
\item[\textit{d})] $W$ stable under direct factors.
\end{itemize}
\end{description}

This is of course implied by:

\begin{description}
\item[Loc (1 ter)] $W$ is \textit{strongly saturated, i.e.} if $f \in \Cat$ is such that $\hot_W(f)$ is an isomorphism, then $f \in W$.
\item[Loc (2)] (\textit{Final object}). If $X$ in Cat has a final object, then $X$ is $W$-aspheric, \textit{i.e.} $X \longrightarrow e$ is in $W$.
\end{description}

Sometimes, this axiom is a consequence of the following weakened variant.

\begin{description}
\item[Loc (2$_0$)] (Segment axiom). $\boldsymbol{\Delta}^1$ is $W$-aspheric.
\end{description}

\msp{page 2}

\begin{description}
\item[Loc (3)] (\textit{localization,} weak form). Let $f : X \longrightarrow S$ be an arrow in Cat which is $W$\textit{-aspheric, i.e.} such that $\forall\ s \in \Ob S$, the localized category $X/s$ over $S$ at $s$, $X/s = X \times_S S/s$, is $W$-aspheric. Then $f \in W$.\footnote[1]{In a final version, one should interchange the notations for Loc (3) and for Loc (3 bis), since the latter appears to me to be the more important of the two - it is satisfied for all the localizers that have been used (to my knowledge), and is valid for every localizer of the form $W_{\mathbf{D}}$, where $\mathbf{D}$ is a derivator (or just a prederivator satisfying Der (4,5) \texttt{[Der (2) is also needed; the numbering of the axioms is that of chapter I, p. 68]}).} % mark printed as superscript (1); [sic] utilisées -> utilisés, agreement with "localiseurs"
\end{description}

\textbf{Recall 1.} Loc (1, 2, 3) imply that we have
\[
W = W^\circ \qquad \textit{i.e. } f \in W \iff f^\circ \in W,
\]
\printedpage{2}
% head: A. Grothendieck : Dérivateurs, Ch. XVI, Édition des Universités de Montpellier II et Paris VII
% foot: Version 18 mars 2019
where we denote by $X \longmapsto X^\circ$, $f \longmapsto f^\circ$ the functor of passage to the opposite category, functor
\[
X \longmapsto X^\circ : \Cat \xrightarrow{\ \sim\ } \Cat
\]
which is an \textit{automorphism} of Cat (thus commutes with all $\varinjlim$ and $\varprojlim$ ). It follows that the set of axioms Loc (1, 2, 3) is equivalent to Loc (1, 2$^{\mathrm{opp}}$, 3$^{\mathrm{opp}}$), where Loc (2$^{\mathrm{opp}}$) is the ``dual'' axiom of Loc (2), namely that every object of Cat having an initial object is $W$-aspheric, and likewise Loc (3$^{\mathrm{opp}}$) the axiom that every $f : X \longrightarrow S$ which is $W$\textit{-coaspheric, i.e.} such that the colocalized categories $s\backslash X$ of $X$ over $S$ are $W$-aspheric, is in $W$.

\textbf{Definition 1.1.} By a \textit{basic localizer} we mean a subset $W$ of Fl(Cat) satisfying the axioms Loc (1, 2, 3) (\textit{i.e.} mild saturation, final object axiom, localization axiom).

\msp{page 3}

\textbf{Recall 1.2.} If $W$ is a basic localizer, then every \textit{homotopism} $f : X \longrightarrow Y$ of Cat is in $W$ (it is essentially for this reason that we imposed (\textit{c}) in mild saturation Loc (1)). If $f, g : X \rightrightarrows Y$ are homotopic (\textit{i.e.} belong to the same connected component of $\underline{\mathrm{Hom}}(X, Y)$), then $\hot_W(f) = \hot_W(g)$ and $f \in W \iff g \in W$.

Here is the strong form of Loc (3).

\begin{description}
\item[Loc (3 bis)] (\textit{Localization,} strong form). Let $S$ be in Cat and $f : X \longrightarrow Y$ a morphism in $\Cat/S$ such that ``$f$ is a $W$-equivalence locally over $S$'', \textit{i.e.} for every $s \in \Ob S$, the induced morphism
\[
f/s : X/s \longrightarrow Y/s
\]
for the localized categories of $X$, $Y$ over $S$ at $s$, is in $W$. Then $f \in W$.
\end{description}

The axiom Loc (3) is nothing but the special case of Loc (3 bis) corresponding to the case where $Y = S$.

\textbf{Proposition 1.3.} \textit{Let $W \subset \mathrm{Fl}(\Cat)$. Then we have the following equivalences}
\DERmark{ENV0926}

\[
\begin{aligned}
\text{Loc (1, 2, 3)} &\iff \text{Loc (1, $2_0$, $3'$)}\\
\text{Loc (1, 2, 3 bis)} &\iff \text{Loc (1, $2_0$, $3'$ bis),}
\end{aligned}
\]
\textit{where the axioms} Loc ($3'$) \textit{and} Loc ($3'$ bis) \textit{are formulated as follows:}

\msp{page 4}

\begin{description}
\item[Loc (3$'$)] \circled{\textit{a}} (Axiom of Cat-fibrations, weak direct form). Let $S$ be in Cat, $X$ a split fibered category over $S$, such that the fibers $X_s$ ($s \in \Ob S$) are $W$-aspheric. Then $X \longrightarrow S$ is in $W$.

\circled{\textit{b}} Let $X$ be in $\Cat/S$, then $i_X : X \longrightarrow \Phi_S(X)$ is in $W$\footnote[2]{(\textit{b}) concerns the \textit{comparison} of an arbitrary $X$ over $S$, with the split fibered category $\Phi_S(X)$ over $S$.}. % mark printed as superscript (2)
\end{description}

\par

% ===================== page 5 (PDF) =====================
\printedpage{3}
% head: A. Grothendieck : Dérivateurs, Ch. XVI, Édition des Universités de Montpellier II et Paris VII
% foot: Version 18 mars 2019

\begin{description}
\item[Loc (3$'$ bis)] \circled{\textit{a}} (Axiom of Cat-fibrations, strong direct form). Let $S$ be in Cat, $X$, $Y$ two split fibered categories over $S$, $f : X \longrightarrow Y$ a cartesian $S$-functor compatible with the splittings, such that $\forall\ s \in \Ob S$, the induced functor on the fibers
\[
f_s : X_s \longrightarrow Y_s
\]
is in $W$. Then $f \in W$.

\circled{\textit{b}} As in Loc ($3'$) (\textit{b}).
\end{description}

\textbf{N.B.} In fact, assuming Loc (1, 2), the statements Loc ($3'$), Loc ($3'$ bis) imply the similar statements without splittings, and without even assuming $f$ cartesian in Loc ($3'$ bis), \textit{cf.} proof below.

\textsc{Proof.} \textit{a}) Loc (1, 2, 3) $\Longrightarrow$ Loc (1, $2_0$, $3'$). It suffices to prove Loc ($3'$). Let us prove Loc ($3'$ $a$). For this, under the conditions of Loc ($3'$ $a$), let us prove that $X \longrightarrow S$ is $W$-coaspheric, \textit{i.e.} that the categories
\[
s\backslash X = X \times_S s\backslash S
\]
are $W$-aspheric. It will indeed follow from this that $f \in W$ by virtue of Loc (3) applied to $f^\circ$ and of $W = W^\circ$. But
\msp{page 5}
let us consider the canonical inclusion
\[
i_s : X_s \lhook\joinrel\longrightarrow s\backslash X,
\]
it follows from the fact that $p : X \longrightarrow S$ is Cat-fibrant that this functor is $W$-aspheric (it even has a right adjoint, \textit{i.e.} the categories $X_s/u$, where $u \in \Ob(s\backslash X)$, have a final object). Thus by Loc (3) we have $i_s \in W$, and since $X_s$ is $W$-aspheric by hypothesis, likewise $s\backslash X$, q.e.d.

Let us prove Loc ($3'$ $b$): We note that $i_X : X \longrightarrow \Phi_S(X)$ is $W$\textit{-coaspheric} (it even has a left adjoint), hence $i_X^\circ$ is $W$-aspheric, hence in $W$, hence $i_X\in W$ since $W = W^\circ$. % [sic] i_X^\circ \in W^\circ -> i_X \in W (or i_X \in W^\circ), apparently intended: from i_X^\circ \in W one gets i_X \in W^\circ = W
\DERmark{ENV0927}


$a$ bis) Loc (1, 2, 3 bis) $\Longrightarrow$ Loc (1, $2_0$, $3'$ bis). This time we have to prove only Loc ($3'$ bis). For (\textit{b}), this is already done. Remains to prove (\textit{a}). If $f$ is as in the statement Loc ($3'$ bis $a$), let us prove that for every $s$ in $S$, the induced morphism
\[
s\backslash f : s\backslash X \longrightarrow s\backslash Y
\]
is in $W$. Since $W = W^\circ$, we will therefore have $(s\backslash f)^\circ \in W$, and applying Loc (3 bis) to $f^\circ$, we find $f^\circ \in W$, hence $f \in W$, and we will be done. But $s\backslash f$ fits into the commutative square
\msp{page 6}
\[
\begin{tikzcd}
X_s \arrow[r, "f_s"] \arrow[d, "i_s^X"'] & Y_s \arrow[d, "i_s^Y"] \\
s\backslash X \arrow[r, "s\backslash f"] & s\backslash Y,
\end{tikzcd}
\]


\endgroup
% END UNIT der-ch16/der-ch16_p001-005.tex
% BEGIN UNIT der-ch16/der-ch16_p006-010.tex
\begingroup\makeatletter%
\newcommand{\printedpage}[1]{\marginpar{\footnotesize p.~#1}}%
\newcommand{\msp}[1]{\marginpar{\footnotesize [#1]}}%
\newcommand{\Loc}{\mathrm{Loc}}%
\makeatletter%
\renewcommand\@makefnmark{\hbox{\textsuperscript{(\@thefnmark)}}}%
\renewcommand\@makefntext[1]{\noindent\textsuperscript{\@thefnmark}\,#1}%
\makeatother%
\makeatother\let\printedpage\derpp\let\msp\dermsp%
\let\mspin\DERmspin%


\setcounter{footnote}{2}

\printedpage{4}
% head: A. Grothendieck : Dérivateurs, Ch. XVI, Édition des Universités de Montpellier II et Paris VII
where, as we saw in $(a)$, $i_s^X$ and $i_s^Y$ are in $W$. By hypothesis we have $f_s \in W$, hence $s\backslash f \in W$, q.e.d.

\textit{b}) Loc (1, $2_0$, $3'$) $\Longrightarrow$ Loc (1, 2, 3). Let us first prove Loc (2). Applying Loc ($3'$ \textit{a}) to the projection
\[
\boldsymbol{\Delta}^1 \times X \longrightarrow X \qquad \text{(where $X \in \operatorname{Ob}\mathrm{Cat}$),}
\]
which is split fibered with fibers isomorphic to $\boldsymbol{\Delta}^1$, hence $W$-aspheric by Loc ($2_0$), we see that this projection is in $W$. This allows us to apply the standard homotopy arguments (\textit{cf.} XII, p. 1 \ldots) to conclude that a homotopism is in $W$, hence a contractible category is $W$-aspheric, a fortiori a category with a final object (or with an initial object) is $W$-aspheric, \textit{i.e.} we have Loc (2).

Let us prove Loc (3). Since Loc (1, 2, 3) $\Longleftrightarrow$ Loc (1, $2^\circ$, $3^\circ$), it suffices to prove Loc ($3^\circ$), since we have proved Loc ($2^\circ$). Let
\[
f : X \longrightarrow S \qquad \text{$W$-\textit{co}aspheric,}
\]
let us prove $f \in W$. Let us introduce the \textit{fibered} category over $S$ whose fibers are the $s\backslash X$, namely $\Phi_S(X)$, and consider the canonical inclusion
\[
i_X : X \longrightarrow \Phi_S(X).
\]
\msp{page 7}By Loc ($3'$ \textit{b}) it is in $W$. On the other hand, the fibers $s\backslash X$ of $\Phi_S(X)$ are $W$-aspheric by hypothesis, hence by Loc ($3'$ \textit{a}) $\Phi_S(X) \longrightarrow S$ is in $W$. Since $i_X \in W$, it follows that $(f : X \longrightarrow S) \in W$, q.e.d.

\textit{b} bis) Loc (1, $2_0$, $3'$ bis) $\Longrightarrow$ Loc (1, 2, 3 bis). We have already proved Loc (1, 2, 3) (since Loc ($3'$ bis) contains Loc ($3'$) as a special case). Let us therefore prove Loc (3 bis) assuming all that already. It will again suffice to prove it in dual form, since
\DERmark{ENV0928}

\[
\text{Loc (1, 2, 3 bis)} \iff \text{Loc(1, $2^\circ$, 3 bis$^\circ$),}
\]
\msp{page 8}and Loc ($2^\circ$) is known from Loc (1, 2, 3) $\Longleftrightarrow$ Loc (1, $2^\circ$, $3^\circ$). It thus remains to prove that if
\[
f : X \longrightarrow Y \qquad \text{in } \mathrm{Cat}/S
\]
is such that
\[
s\backslash f : s\backslash X \longrightarrow s\backslash Y
\]
is in $W$ for every $s \in \operatorname{Ob} S$, then $f \in W$. Let us then consider the commutative square
\[
\begin{tikzcd}
X \arrow[r, "i_X"] \arrow[d, "f"'] & \Phi_S(X) \arrow[d, "\Phi_S(f)"] \\
Y \arrow[r, "i_Y"] & \Phi_S(Y),
\end{tikzcd}
\]
where $i_X, i_Y \in W$ by Loc ($3'$ bis \textit{b}). The hypothesis on $f$ means that $\Phi_S(f)$ induces $W$-equivalences on the fibers, hence by Loc ($3'$ bis) $\Phi_S(f) \in W$, hence $f \in W$ by Loc (1), q.e.d.
% foot: Version 18 mars 2019

\printedpage{5}
% head: A. Grothendieck : Dérivateurs, Ch. XVI, Édition des Universités de Montpellier II et Paris VII
\section*{2\quad Axioms Loc (4) and Loc (5) (of direct sums, of connectedness)\protect\footnote{We are dealing here with axioms that one can regard as being of a more or less ``trivial'' nature. In all cases known to me, they are indeed trivially satisfied.}}

\begin{description}[leftmargin=7em, labelwidth=6em]
\item[Loc (4)] (Axiom of \textit{direct sums}). For a family $(f_i)_{i\in I}$ of arrows in Cat, if $f = \coprod f_i$, we have $f \in W$ if and only if $f_i \in W$ for every $i \in I$.
\end{description}

This axiom splits into two, into the ``if'' and the ``only if'', which I am going to formulate in the most economical form:

\textbf{Proposition 2.1.} \textit{The axiom} Loc (4) \textit{of sums is equivalent to the conjunction of the following two axioms:}

\begin{description}[leftmargin=7em, labelwidth=6em]
\item[Loc (4 a)] \msp{page 9}Let $(f_i)_{i\in I}$ be a family of arrows in Cat, if the $f_i \in W$, then $f = \coprod f_i \in W$.
\item[Loc (4 b)] Let $S$ be in Cat, and $S_0 \hookrightarrow S$ a direct summand (subcategory both open and closed). Let $X$ be in Cat$/S$. If $X \longrightarrow S$ is in $W$, then so is $X_0 = X|S_0 \longrightarrow S_0$. (In other words, the inclusion $S_0 \hookrightarrow S$ is a ``$W$-fibration'', \textit{i.e.} is in $\mathrm{Fib}_W$.)
\end{description}

\textbf{Corollary 2.2.} \textit{Consider the variant of} Loc (4 a) \textit{obtained by restricting to finite families.} (\textit{It is equivalent to} Loc (4 a) \textit{when $W$ is stable under} $\varinjlim$ \textit{that are filtering,} cf. \textit{axiom} Loc (8) \textit{below}). \textit{This variant is equivalent to}
\DERmark{ENV0929}


\begin{description}[leftmargin=7.5em, labelwidth=6.5em]
\item[Loc (4 a$_0$)] Every object $X$ of Cat is ``$W$-cofibrant'', \textit{i.e.} $\emptyset \longrightarrow X$ is in $\mathrm{Cof}_W$, \textit{i.e.} for every $f : Y \longrightarrow Y'$ which is in $W$, the corresponding morphism $\mathrm{id}_X \amalg f : X \amalg Y \longrightarrow X \amalg Y'$ is in $W$.
\end{description}

\textbf{Reminder 2.3.} In any case, if $W$ is a basic localizer (which we shall assume from now on, unless otherwise stated), the \textit{finite} analogue of Loc (4 a) for products is true, \textit{i.e.} if we have a \textit{finite} family of arrows $(f_i)_{i\in I}$ in Cat, if the $f_i$ are in $W$, then $f = \prod f_i \in W$. \texttt{[For this, one must assume that $W$ satisfies the localization axiom, \textit{strong} form, Loc 3 bis.]} The analogue for an infinite family is satisfied for practically no basic localizer.

\msp{page 10}\textbf{Proposition 2.4.} \textit{Let $W$ be a basic localizer. Then finite products exist in} $\mathrm{Hot}_W$\textit{, and the functor} $\mathrm{hot}_W : \mathrm{Cat} \longrightarrow \mathrm{Hot}_W \overset{\mathrm{def}}{=} (\mathrm{Cat})W^{-1}$ \textit{commutes with them. If $W$ satisfies} Loc (4 a) (resp. Loc (4 a$_0$)), \textit{then arbitrary sums} (resp. \textit{finite}) \textit{exist in} $\mathrm{Hot}_W$\textit{, and the functor} $\mathrm{hot}_W$ \textit{commutes with them. Moreover, in this case sums are distributive with respect to products (since this is so in} Cat).

The proof follows simply from the general result:
\DERmark{ENV0930}


\textbf{Lemma 2.5.} \textit{Let $(\mathcal{M}, W)$ be a category of models ($W$ any subset of} $\mathrm{Fl}\,\mathcal{M}$). \textit{Suppose that products} (resp. \textit{finite products,} resp. \textit{binary products}) \textit{exist in $\mathcal{M}$, and that moreover for every family $(f_i)_{i\in I}$ of arrows of $\mathcal{M}$} (resp. \textit{every finite family,}
% foot: Version 18 mars 2019
\printedpage{6}
% head: A. Grothendieck : Dérivateurs, Ch. XVI, Édition des Universités de Montpellier II et Paris VII
resp. \textit{binary family}) \textit{$f_i \in W$ for every $i \in I$ implies $f = \prod f_i \in W$. Then products of the type envisaged exist in} $\mathrm{Ho}_W = \mathcal{M}W^{-1}$\textit{, and the functor} $\mathrm{ho}_W : \mathcal{M} \longrightarrow \mathcal{M}W^{-1}$ \textit{commutes with them. Finally, if $\mathcal{M}$ admits a final object $e$, then $e$ is also a final object of} $\mathrm{Ho}_W$.

\texttt{[It is not clear that the assertion is true in this generality for \textit{infinite} products; it is true if $W$ is the class of weak equivalences of a \textit{Quillen model category.}]}

\msp{page 11}We deduce from this the dual statements concerning direct sums and initial objects.

The proof is a special case of the following very useful result, applied to the pair of adjoint functors
\[
\begin{tikzcd}
\mathcal{M} \arrow[r, shift left, "\mathrm{diag}"] & \mathcal{M}^I \arrow[l, shift left, "\text{\scriptsize\shortstack{products\\ in $\mathcal{M}$\\ indexed by $I$}}"]
\end{tikzcd}
\qquad\qquad :
\]

\textbf{Lemma 2.6}\footnote{\textit{cf.} XII, lemma on page 248.}. \textit{Let $(\mathcal{M}, W)$, $(\mathcal{M}', W')$ be two categories of models, and}
\[
\begin{tikzcd}
\mathcal{M} \arrow[r, shift left, "\varphi"] & \mathcal{M}' \arrow[l, shift left, "\psi"]
\end{tikzcd}
\]
\textit{a pair of adjoint functors compatible with the localizers,} i.e. \textit{such that}
\[
\varphi(W) \subset W', \qquad \psi(W') \subset W.
\]
\textit{Then the localized functors}
\[
\begin{tikzcd}
\mathcal{M}W^{-1} \arrow[r, shift left, "\bar{\varphi}"] & \mathcal{M}'W'^{-1} \arrow[l, shift left, "\bar{\psi}"]
\end{tikzcd}
\]
\textit{are also adjoint, the adjunction arrows for $\bar{\varphi}\bar{\psi}$ and $\bar{\psi}\bar{\varphi}$ being deduced from those for $\varphi\psi$ and $\psi\varphi$ by composition with the localization functors $\mathcal{M} \longrightarrow \mathcal{M}W^{-1}$, $\mathcal{M}' \longrightarrow \mathcal{M}'W'^{-1}$.}

\begin{description}[leftmargin=7em, labelwidth=6em]
\item[Loc (5)] \msp{page 12}(Axiom of \textit{connectedness} or of \textit{non-triviality}).

Let $f \in W$, then $\pi_0(f)$ is bijective.
\end{description}

This property of $W$ is in constant use in practice, when we work with $W$ instead of restricting ourselves to constructing $W$'s and playing around with the axioms and their stability under passage to $\underline{\mathrm{Hom}}(I, M)$'s (where here $\mathcal{M} = \mathrm{Cat}$). On the other hand, we can also regard it as an axiom of \textit{non-triviality} for $\mathrm{Hot}_W$. I recall in this connection:

\textbf{Proposition 2.7.} \textit{Let $W$ be a basic localizer. Equivalent conditions:}
\begin{itemize}[leftmargin=3em]
\item[a)] $W$ \textit{satisfies the connectedness axiom} Loc (5).
\item[b)] \textit{The functor} $\pi_0 : \mathrm{Cat} \longrightarrow \mathrm{Ens}$ \textit{factors through a functor}
% foot: Version 18 mars 2019
\printedpage{7}
% head: A. Grothendieck : Dérivateurs, Ch. XVI, Édition des Universités de Montpellier II et Paris VII
\[
\begin{tikzcd}
\mathrm{Cat} \arrow[r, dotted] \arrow[rr, dotted, bend right, "\pi_0"'] & \mathrm{Hot}_W \arrow[r, "\pi_0"] & \mathrm{Ens}.
\end{tikzcd}
\]
\item[c)] \textit{The category} $\mathrm{Hot}^*_W = \mathrm{Hot}_W\backslash\{\emptyset\}$ \textit{is not equivalent to the one-point category.}
\item[c$'$)] \textit{The category} $\mathrm{Hot}_W$ \textit{is equivalent neither to $e$ nor to $\boldsymbol{\Delta}^1$.}\footnote{\textbf{NB} If $W$ is strongly saturated (condition Loc (1 ter)), then the conditions opposite (\textit{i.e.} Loc (5)) are equivalent to this: % [sic] "(condition Loc (1 ter)," -> "(condition Loc (1 ter)),", parenthèse ouvrante non refermée
\begin{itemize}[leftmargin=2.5em]
\item[d)] \textit{$W$ is equal neither to Fl}(Cat) \textit{in its entirety, nor to Fl}(Cat$^*$) $\cup$ $\{\mathrm{id}_\emptyset\}$ (i.e. \textit{Fl}(Cat)$\backslash W$ \textit{is neither empty, nor equal to the set of arrows $\emptyset \longrightarrow X$ with $X \neq \emptyset$}).
\end{itemize}
If $W$ satisfies Loc (7 D$_{W_\omega}$), then $W$ satisfies Loc (5) if and only if $W \neq \mathrm{Fl}(\mathrm{Cat})$ (\textit{cf.} XII, th. 2, p. 20).}
\DERmark{ENV0931}

\end{itemize}

For the proof, and for other equivalent forms of the connectedness axiom, \msp{page 13}I refer to XV 9.12 (p. 133 bis). For the name ``connectedness axiom'', I can offer the following explanation 2.8 and 2.9.

\textbf{Corollary 2.8.} \textit{If the connectedness axiom} Loc (5) \textit{is satisfied, then the $W$-aspheric objects are} 0\textit{-connected} (\textit{tautology}), \textit{and the converse is true if $W$ is strongly saturated.} (\textit{Or also, if $W$ satisfies the axiom of fibrations} Loc (7) \textit{below, since by} XII, th. 2, p. 20, \textit{in this case,} Loc (5) \textit{is equivalent to $W \neq \mathrm{Fl}(\mathcal{M})$ \ldots}).

For if the connectedness axiom is not satisfied, then for every nonempty $X$ in Cat, $X \longrightarrow e$ is an isomorphism, hence ($W$ being strongly saturated) $X \longrightarrow e$ is in $W$, \textit{i.e.} $X$ is $W$-aspheric, hence there are $W$-aspheric objects that are not 0-connected. Let us also note: % [sic] "non 0-connexe" -> "non 0-connexes", accord avec "des objets"

\textbf{Corollary 2.9.} \textit{Suppose that $W$ satisfies the axiom} Loc (4) \textit{of direct sums} (\textit{or only} Loc (4 b)). \textit{For $W$ to satisfy the connectedness axiom, it is necessary and sufficient that for every arrow $f : X \longrightarrow Y$ with $Y$} 0\textit{-connected, if $f$ is in $W$, then $X$ is} 0\textit{-connected.} (\textit{This is necessary of course without assuming} Loc (4 b) \textit{or anything at all.})

\msp{page 14}\texttt{[The numbers 2.7, 2.8 and 2.9 appear twice.]}

\textbf{2.7.} Let us now consider the composite functor
\[
\mathrm{Ens} \xrightarrow{I \longmapsto I_{\mathrm{Cat}}} \mathrm{Cat} \xrightarrow{\quad \mathrm{hot}_W \quad} \mathrm{Hot}_W, \qquad I \longmapsto I_{\mathrm{Hot}_W} \text{ or } I_W, \tag*{(2.7.1)}
\]
where the first functor is the one that assigns to the set $I$ the discrete category $I_{\mathrm{Cat}}$ having $I$ as its set of objects. (Most often, we shall identify $I_{\mathrm{Cat}}$ with $I$ in the notation, just as for an ordered set $J$, one identifies $J$ and the ``ordered category'' % [sic] "catégorie ordonné" -> "catégorie ordonnée", accord
defined by $J$. Here $I$ can be regarded as equipped with the discrete order.) We have, moreover
\[
I_{\mathrm{Cat}} \simeq I \times e \ \ (\overset{\mathrm{def}}{=} \text{sum of $I$ copies of $e$}), \tag*{(2.7.2)}
\]
\DERmark{ENV0932}

so when Loc (4 a) is satisfied, so $\mathrm{hot}_W$ commutes with direct sums, we shall have a canonical isomorphism
\[
\begin{aligned}
I_W \simeq I \times e_W \quad &\text{(where $e_W$ is the final object of $\mathrm{Hot}_W$,} \\
&\text{\textit{i.e.} $e$ regarded as an object of $\mathrm{Hot}_W$).}
\end{aligned} \tag*{(2.7.3)}
\]
% foot: Version 18 mars 2019

\printedpage{8}
% head: A. Grothendieck : Dérivateurs, Ch. XVI, Édition des Universités de Montpellier II et Paris VII
\textbf{Proposition 2.8.} \textit{Suppose the connectedness axiom} Loc (5) \textit{is satisfied, whence a functor}
\[
\pi_0 : \mathrm{Hot}_W \longrightarrow \mathrm{Ens}. \tag*{(2.8.1)}
\]
\textit{This functor is} left adjoint \textit{to the functor $I \longmapsto I_W$} (2.7.1), i.e. \textit{we have a bijection functorial in $(X, I)$ in} $\mathrm{Hot}_W \times \mathrm{Ens}$:
\msp{page 15}
\[
\mathrm{Hom}_{\mathrm{Ens}}(\pi_0(X), I) \simeq \mathrm{Hom}_W(X, I_W). \tag*{(2.8.2)}
\]
\textit{The functor $\pi_0$ is also} right adjoint \textit{to the} same \textit{functor $I \longmapsto I_W$,} i.e. \textit{we have a functorial bijection} % [sic] "un bijection" -> "une bijection", accord
\[
\mathrm{Hom}_{\mathrm{Ens}}(I, \pi_0(X)) \simeq \mathrm{Hom}_W(I_W, X). \tag*{(2.8.3)}
\]

Taking $I = e$ (one-point set), we find the:

\textbf{Corollary 2.9.} \textit{We have the isomorphism functorial in $X$ in} $\mathrm{Hot}_W$
\[
\boxed{\mathrm{Hom}_W(e_W, X) \simeq \pi_0(X).} \tag*{(2.8.4)}
\]

\textbf{Comment 2.10.} The adjunction formula (2.8.2) is, moreover, familiar, whenever we have a category $\mathcal{H}$ (such as Cat, or $\mathrm{Hot}_W$ \textit{etc.}) % [sic] "où" -> "ou", conjonction de coordination
\DERmark{ENV0933}

equipped with a functor $\pi_0$, $\pi_0(X)$ corresponding to the intuition of a decomposition of an object into ``connected components''. (I am not going to spell out here the general conditions under which one defines such a functor.) On the other hand, the adjunction in the reverse direction of these same two functors is a phenomenon quite special to categories \msp{page 16}of the type $\mathrm{Hot}_W$. If we take into account the description of $I_W$ as $I \times e_W$, the unusual adjunction formula (2.8.3) is equivalent, of course, to formula (2.8.4) - it is deduced from it by ``raising to the power $I$''. We can reverse the steam, and in a general category $\mathcal{H}$, \textit{define} a functor $\pi_0 : \mathcal{H} \longrightarrow \mathrm{Ens}$ by formula (2.8.4), % [sic] "par le formule" -> "par la formule", accord
\[
\pi_0(X) \overset{\mathrm{def}}{=} \mathrm{Hom}(e_{\mathcal{H}}, X) \qquad (= \Gamma_{\mathcal{H}}(X), \text{ with a familiar notation}),
\]
\DERmark{ENV0934}

where $e_{\mathcal{H}}$ is a final object of $\mathcal{H}$. We shall therefore have, putting $I_{\mathcal{H}} = I \times e_{\mathcal{H}}$, the formula
\[
\mathrm{Hom}(I_{\mathcal{H}}, X) \xrightarrow{\ \sim\ } \mathrm{Hom}(I, \pi_0(X))
\]
(\textit{i.e.} formula (2.8.3)), saying that we have a pair $(I_{\mathcal{H}}, \pi_0)$ of adjoint functors. Suppose on the other hand that for every set $I$, the canonical adjunction map
\[
I \longrightarrow \pi_0(I_{\mathcal{H}}) = \mathrm{Hom}_{\mathcal{H}}(e_{\mathcal{H}}, e_{\mathcal{H}} \times I)
\]
is \textit{bijective}\footnote{that tends to be ultrafalse if $\mathcal{H}$ is additive. On the other hand, true if $\mathcal{H}$ is a category with sums distributive with respect to products (or even universal), whose final object is \textit{connected} \ldots}, \textit{i.e.} that $I \longmapsto I_{\mathcal{H}}$ is \textit{fully faithful.} We then define, for every $I$, a canonical map
% foot: Version 18 mars 2019


\endgroup
% END UNIT der-ch16/der-ch16_p006-010.tex
% BEGIN UNIT der-ch16/der-ch16_p011-015.tex
\begingroup\makeatletter%
\newcommand{\printedpage}[1]{\marginpar{\footnotesize p.~#1}}%
\newcommand{\msp}[1]{\marginpar{\footnotesize [#1]}}%
\newcommand{\Hom}{\operatorname{Hom}}%
\newcommand{\Ob}{\operatorname{Ob}}%
\newcommand{\Fl}{\operatorname{Fl}}%
\newcommand{\Hot}{\mathrm{Hot}}%
\newcommand{\Cat}{\mathrm{Cat}}%
\newcommand{\Ens}{\mathrm{Ens}}%
\newcommand{\id}{\mathrm{id}}%
\makeatletter%
\renewcommand\@makefnmark{\hbox{\@textsuperscript{\normalfont(\@thefnmark)}}}%
\renewcommand\@makefntext[1]{\parindent 1em\noindent\hb@xt@1.8em{\hss\@textsuperscript{\normalfont\@thefnmark}}#1}%
\makeatother%
\makeatother\let\printedpage\derpp\let\msp\dermsp%
\let\mspin\DERmspin%


\printedpage{9}
% head: A. Grothendieck : Dérivateurs, Ch. XVI, Édition des Universités de Montpellier II et Paris VII
% foot: Version 18 mars 2019

\[
\begin{tikzcd}[column sep=small]
\Hom(X,I_{\mathcal{H}}) \arrow[rr] \arrow[dr,dotted,"\pi_0^{X,I_{\mathcal{H}}}"'] & & \Hom(\pi_0(X),I) \arrow[dl,dotted,"\sim"'] \\ % ?? uncertain: label of the dotted right arrow, a tilde printed along the arrow -- \sim
& \Hom(\pi_0(X),\pi_0(I_{\mathcal{H}})). &
\end{tikzcd}
\]

The ``bi-adjunction'' of the functors $\pi_0$ and $I\longmapsto I_{\mathcal{H}}$ amounts to saying that this arrow is bijective,
\msp{17}
or again, that the arrow $\pi_0^{X,I_{\mathcal{H}}}$ :
\[
\Hom(X,I_{\mathcal{H}})\longrightarrow\Hom(\Gamma_{\mathcal{H}}(X),\Gamma_{\mathcal{H}}(I_{\mathcal{H}})\simeq I)
\]
is bijective.

\textbf{Corollary 2.11.} \textit{The functor} $I\longmapsto I_W$ \textit{from} Ens \textit{to} $\Hot_W$ \textit{is fully faithful,} i.e. \textit{the adjunction homomorphism}
\[
\pi_0(I_W)\longrightarrow I
\]
\textit{is bijective.} (\textit{Or again,} % [sic] Où encore -> Ou encore, the accent makes no sense here
\DERmark{ENV0935}

\textit{$\pi_0$ is a localization functor.}) % [sic] le localisation -> de localisation, wrong word
\DERmark{ENV0936}


Indeed, this bijectivity is obvious when one works in Cat, and follows from this at once in $\Hot_W$.

\textsc{Proof of} 2.8. The adjunction formula (2.8.2) follows from the corresponding formula for Cat, obvious, and from lemma 2.6 on adjoint functors between categories of models. The adjunction in the other direction reduces to formula (2.8.4). We have canonical arrows
\[
\begin{tikzcd}
\Hom_W(e_W,X) \arrow[r,shift left,"\alpha"] & \pi_0(X), \arrow[l,shift left,"\beta"]
\end{tikzcd}
\]
where $\alpha$ is $\pi_0^{e_W,X}$, given that $\pi_0(e_W)=\pi_0(e_{\Cat})=e_{\Ens}$, and $\Hom(e_{\Ens},\pi_0(X))\simeq\pi_0(X)$. The arrow $\beta$ is the composite arrow
\[
\Ob X\xrightarrow{\ \sim\ }\Hom_{\Cat}(e_{\Cat},X)\xrightarrow{\mathrm{hot}_W}\Hom_W(e_W,X)
\]
\msp{18}
noting that this arrow passes to the quotient as $\pi_0(X)\longrightarrow\Hom_W(e_W,X)$, since two homotopic maps from $e_{\Cat}$ to $X$ are equal in $\Hot_W$. I say that
\[
\beta\alpha=\id\ ,\qquad\qquad \alpha\beta=\id\ .
\]
The relation $\alpha\beta=\id$ amounts to saying that the arrow $i_x : e_{\Cat}\longrightarrow X$ defined by $x\in\Ob X$, gives $\pi_0(i_x) : e_{\Ens}\longrightarrow\pi_0(X)$ whose value is the class of $x$ in $\pi_0(X)$, this is tautological. To prove that $\beta\alpha=\id$, it then suffices to prove that $\beta$ is surjective, \textit{i.e.} the

\textbf{Corollary 2.12.} \textit{For every $X$ in} Cat, \textit{the map}
\[
\underbrace{\Hom_{\Cat}(e_{\Cat},X)}_{\simeq\Ob X}\longrightarrow\Hom_W(e_W,X)
\]
\printedpage{10}
% head: A. Grothendieck : Dérivateurs, Ch. XVI, Édition des Universités de Montpellier II et Paris VII
% foot: Version 18 mars 2019
\textit{is} surjective (\textit{assuming the connectedness axiom} Loc (5)).

\textsc{Proof of} 2.12. If $X=\emptyset$, the assertion amounts to saying that $\Hom_W(e_W,\emptyset)=\emptyset$, but more generally, we have the

\textbf{Corollary 2.13.} (Assuming Loc (5).) \textit{Let $X$ be nonempty in} Cat, \textit{then} $\Hom_W(X,\emptyset)=\emptyset$. (\textit{Thus $\emptyset$ is an initial object that is} strict \textit{in} $\Hot_W$.)

Indeed, if we had $X\longrightarrow\emptyset$ in $\Hot_W$,
\msp{page 13'}\texttt{[there is a shift in the numbering]}
we would deduce $\pi_0(X)\longrightarrow\emptyset$, or $\pi_0(X)=\emptyset$, absurd\,!

There remains the case where $X$ is nonempty. This follows from the general fact :

\textbf{Lemma 2.14.} \textit{Let $\mathcal{M}$ be a category having an initial object that is} strict (\textit{if} $\Hom_{\mathcal{M}}(X,\emptyset_{\mathcal{M}})\neq\emptyset$, \textit{we have} $X=\emptyset_{\mathcal{M}}$), \textit{a final object $e$, and such that for every $X\neq\emptyset_{\mathcal{M}}$, $X\longrightarrow e$ has a section,} i.e. $\Hom(e,X)\neq\emptyset$. \textit{Let} $W\subset\Fl\mathcal{M}$
\begin{itemize}[leftmargin=3em,labelwidth=2.5em,align=right]
\item[1$^{\circ}$)] \textit{not containing any arrows} $\emptyset_{\mathcal{M}}\longrightarrow X$ \textit{with} $X\neq\emptyset_{\mathcal{M}}$, \textit{and}
\item[2\,$^{\circ}$)] \textit{such that for every diagram}
\begin{tikzcd}[baseline=(\tikzcdmatrixname-1-1.base)]
X_1 & X_2 \arrow[l,"u_2"'] \\
e \arrow[u,"u_1"] \arrow[ur,dotted,"s"'] &
\end{tikzcd}
\textit{with} $u_2\in W$ $\exists$ $s$ in $\mathcal{M}$ \textit{such that} $u_2s=u_1$

\textit{in} $\mathcal{M}W^{-1}$.
\end{itemize}
\textit{Then}
\begin{itemize}[leftmargin=3em,labelwidth=2.5em,align=right]
\item[a)] $\emptyset_{\mathcal{M}}$ \textit{is an initial object that is} strict \textit{in} $\mathcal{M}W^{-1}$.
\item[b)] \textit{For every} $X\neq\emptyset_{\mathcal{M}}$, $\Hom_{\mathcal{M}}(e,X)\longrightarrow\Hom_W(e,X)$ \textit{is surjective.}
\end{itemize}

\textsc{Proof}. That $\emptyset_{\mathcal{M}}$ is an initial object of $\mathcal{M}W^{-1}$ has been seen (last assertion of lemma 2.5, dual form), let us prove that it is strict, \textit{i.e.} that if $\Hom_W(X,\emptyset_{\mathcal{M}})\neq\emptyset$, we have $X=\emptyset_{\mathcal{M}}$. Indeed, if we had $u : X\longrightarrow\emptyset_{\mathcal{M}}$ in $\mathcal{M}W^{-1}$, $u$ would be represented by
\[
X=X_0\xleftarrow{\ W\ }X_1\longrightarrow X_2\xleftarrow{\ W\ }\cdots X_{2n-2}\longrightarrow X_{2n-1}\xleftarrow{\ W\ }X_{2n}=\emptyset_{\mathcal{M}}.
\]
We proceed by induction on the half-length of this chain to see that $X=\emptyset_{\mathcal{M}}$. Indeed, the hypothesis on $W$ implies that $X_{2n-1}=\emptyset_{\mathcal{M}}$, then the one that $\emptyset_{\mathcal{M}}$ is strict implies
$X_{2n-2}=\emptyset_{\mathcal{M}}$, whence the conclusion. % [sic] X_{2n-1} -> X_{2n-2}, the strictness of the initial object gives X_{2n-2}
\DERmark{ENV0937}


\msp{page 14'}
Let us prove the surjectivity of $\Hom_{\mathcal{M}}(e,X)\longrightarrow\Hom_W(e,X)$ for $X\neq\emptyset_{\mathcal{M}}$, by showing that every reduced path of length 2n+1 is ``homotopic'' to a reduced path of length $2n-1$ :
\[
e=X_0\longrightarrow X_1\xleftarrow{\ W\ }X_2\longrightarrow X_3\longleftarrow\ \cdots\ \longleftarrow X_{2n}\longrightarrow X_{2n+1}=X.
\]
Indeed, we can replace the section
\[
\begin{tikzcd}
e=X_0 \arrow[r,"u_1"'] \arrow[rr,dotted,bend right,"s"'] & X_1 & X_2 \arrow[l,"W"',"u_2"] \arrow[r,"u_3"'] & X_3 & X_4 \arrow[l,"W"']
\end{tikzcd}
\]
\printedpage{11}
% head: A. Grothendieck : Dérivateurs, Ch. XVI, Édition des Universités de Montpellier II et Paris VII
% foot: Version 18 mars 2019
by
\[
\begin{tikzcd}
e \arrow[r,"u_3s"'] & X_3 & \cdots, \arrow[l,"W"']
\end{tikzcd}
\]
where $s$ is a section of $X_2$ over $e$ such that
\[
u_2s=u_1\qquad\quad\textit{in}\qquad\quad\mathcal{M}W^{-1}\ .
\]
This therefore proves by induction on $n$ that the path is homotopic to a path of length 1, $e\longrightarrow X$, q.e.d.

We must therefore, in the present case, check 1$^{\circ}$) (this follows from Loc (5)), and 2$^{\circ}$)
\DERmark{ENV0938}

\[
\begin{tikzcd}
X_1 & X_2 \arrow[l,"u_2\in W"'] & \\
e \arrow[u,"u_1"] \arrow[ur,dotted,"s"'] & & .
\end{tikzcd}
\]
But since $u_2\in W$, we have $\pi_0(u_2)$ an isomorphism, hence
\msp{page 15'}
surjective, hence $\exists$ $s$ such that $u_2s$ is homotopic to $u_1$ in Cat, which implies $u_2s=u_1$ in $\Hot_W$, q.e.d.

\textbf{Final comment.} The connectedness axiom Loc (5) differs from all the other axioms Loc (1) to Loc (8) (but with the exception of Loc (8 bis) on accessibility) in that it is not a stability axiom. The other axioms Loc (1) to Loc (8) (and also Loc (8 bis) on accessibility and even the later axioms % [sic] ultérieures -> ultérieurs, agreement with axiomes
of exactness~\footnote[7]{there will in fact be no need for them, thanks to the providential theory of K. S. \textsc{Brown}\,!}) are satisfied if $W=\Fl\mathcal{M}$ - the connectedness axiom is therefore the only one that stands in the way of the \textit{triviality} of the homotopy theory defined by $W$. This is to say that it cannot be a consequence of any of the others. On the other hand, the axiom Loc (4) of direct sums is not independent of the others. Let us point out in this connection (reminder) :
\[
\left\{
\begin{array}{rl}
\text{Loc (4 a$_0$)} & \text{is a consequence of the cocartesian square axiom Loc (6 a) (direct form).}\\[1ex]
\text{Loc (4 a)} & \text{is a consequence of the preceding Loc (6 a)}\\
 & \text{+ the axiom of filtering direct limits Loc (8).}\\[3ex]
\text{Loc (4 b)} & \text{is a consequence of the fibration axiom Loc (7). Thus}\\[1ex]
\text{Loc (4)} & \text{is a consequence of Loc (6 a, 7, 8).}
\end{array}
\right.
\]

\printedpage{12}
% head: A. Grothendieck : Dérivateurs, Ch. XVI, Édition des Universités de Montpellier II et Paris VII
% foot: Version 18 mars 2019
\msp{page 16'}

\section*{3\quad Axiom Loc (6) of the cocartesian square (or of ``Mayer-Vietoris'', or of $W$-cofibrations)~\footnote[8]{It is the name ``\textit{axiom of $W$-cofibrations}'' that strikes me at present as the best, -- it is this name that seems to me best to capture the conceptual and not technical meaning of the axiom.}}

I see several variants, and am not yet totally clear in my mind about the mutual implications. I begin with the more complicated variant, but which we will need in any case :

\begin{itemize}[leftmargin=5em,labelwidth=4em,align=left]
\item[\textbf{Loc 6}] (Axiom of $W$-cofibrations, or 5-lemma, or axiom of homomorphisms of cocartesian squares). Let us consider in Cat a commutative cubical diagram
\[
\begin{tikzcd}[row sep=small,column sep=small]
X_0 \arrow[rr] \arrow[dr] \arrow[dd,"f_0"'] & & X_2 \arrow[dr] \arrow[dd,"f_2" near start] & \\
& X_1 \arrow[rr,crossing over] & & X_3 \arrow[dd,"f_3"] \\
Y_0 \arrow[rr] \arrow[dr] & & Y_2 \arrow[dr] & \\
& Y_1 \arrow[rr] \arrow[from=uu,crossing over,"f_1"' near start] & & Y_3,
\end{tikzcd}
\]
\end{itemize}
a homomorphism $f=(f_0,f_1,f_2,f_3)$ from a square diagram $X$ to another $Y$. We suppose that $X$ comes from an object $X_3$ of Cat from two opens $X_1$, $X_2$ of $X_3$ \textit{covering} $X_3$, and from their intersection $X_0=X_1\cap X_2$, and likewise for $Y$. (Let us call these diagrams in Cat MV diagrams (``\textit{open}'') - ``\textit{Mayer-Vietoris}''.) We then have the following~\footnote[9]{It is good to introduce properties of the diagram
\begin{itemize}[leftmargin=3em,labelwidth=2.5em,align=right]
\item[1$^{\circ}$)] $f_0\in W$.
\item[2$^{\circ}$)] $f_1$ and $f_2$ in $W$.
\item[3$^{\circ}$)] $f_3$ in $W$.
\end{itemize}
and to formulate parts a) b) c) of Loc (6) as follows
\begin{itemize}[leftmargin=3em,labelwidth=2.5em,align=right]
\item[a)] 1$^{\circ}$ + 2$^{\circ}$ $\Longrightarrow$ 3$^{\circ}$.
\item[b)] 2$^{\circ}$ + 3$^{\circ}$ $\Longrightarrow$ 1$^{\circ}$.
\item[c)] 1$^{\circ}$ + 3$^{\circ}$ $\Longrightarrow$ 2$^{\circ}$.
\end{itemize}} :
\begin{itemize}[leftmargin=3em,labelwidth=2.5em,align=right]
\item[a)] $f_0,f_1,f_2\in W\Longrightarrow f_3\in W$~\footnote[10]{N.B. it seems that in Loc (6 a), the statement suffices on assuming $f_0$, $f_2$ to be \textit{isomorphisms,} \textit{cf.} p. 22.}.
\item[b)] $f_1,f_2,f_3\in W\Longrightarrow f_0\in W$.
\item[c)] $f_0,f_3\in W\Longrightarrow f_1,f_2\in W$~\footnote[11]{One should not include c in Loc 6 \textit{cf.} p. 26.}.
\end{itemize}
\noindent\footnote[12]{N.B. We will see below (3.3) that Loc 6 c already implies Loc (6 a, b), \textit{i.e.} all of Loc (6). And if $W$ satisfies Loc (1 bis), \textit{i.e.} stable under direct factors, then Loc (6 a, b) $\Longrightarrow$ Loc (6 c).}.

\printedpage{13}
% head: A. Grothendieck : Dérivateurs, Ch. XVI, Édition des Universités de Montpellier II et Paris VII
% foot: Version 18 mars 2019
\msp{page 17'}
\noindent\footnote[13]{I am going to denote by Loc (6) the axiom Loc (6 a, b), without including c in it.}.

\textbf{N.B.} I distinguish these conditions as Loc (6 a), Loc (6 b), Loc (6 c). I presume that we will never have to make use of Loc (6 b), without also assuming Loc (6 a), but the reverse is surely not the case. Thus, $W_0=\{f\in\Fl\Cat\mid\pi_0(f)\text{ bijective}\}$ satisfies a), but not b). I presume (without having checked) that it is the same for the $W_n$, $n\in\mathbf{N}$, of ``\textit{$n$-homotopy equivalence}''.

It could be that Loc (6 a) and Loc (6 a, b) are equivalent to the following weaker, and simpler, form. (If so, it is in the simplest form, of course, that they should be stated.)

\begin{itemize}[leftmargin=5em,labelwidth=4em,align=left]
\item[\textbf{Loc 6$'$}] (Axiom of cocartesian squares, or of Mayer-Vietoris). Let
\[
\begin{tikzcd}[row sep=small,column sep=small]
X_0 \arrow[rr,"i_2"] \arrow[dr,"i_1"] & & X_2 \arrow[dr,"j_1"] & \\
& X_1 \arrow[rr,"j_2"] & & X_3
\end{tikzcd}
\]
be an MV diagram in Cat. Then
\begin{itemize}[leftmargin=3em,labelwidth=2.5em,align=right]
\item[a)] If $i_1\in W$, $j_1\in W$.
\item[b)] If $j_1\in W$, $i_1\in W$.
\end{itemize}
\end{itemize}

\msp{page 18'}
\textbf{Proposition 3.1.}
\[
\begin{aligned}
\text{Loc (6 a)}\quad &\Longrightarrow\quad \text{Loc (6$'$ a),}\\
\text{Loc (6 b)}\quad &\Longrightarrow\quad \text{Loc (6$'$ b),}\\
\text{Loc (6 c)}\quad &\Longrightarrow\quad \text{Loc (6$'$ a, b)}\ \footnotemark[14].
\end{aligned}
\]
\footnotetext[14]{N.B. we even have Loc (6 c) $\Longrightarrow$ $\underbrace{\text{Loc 6 (a, b)}}_{\text{Loc (6)}}$.}

This is seen on the cubical diagram
\[
\text{(Q)}\qquad\qquad
\begin{tikzcd}[row sep=small,column sep=small]
X_0 \arrow[rr,"i_2"] \arrow[dr,"i_1"] & & X_2 \arrow[dr,"j_1"] \arrow[from=dd,"f_2=\id"' near start] & \\
& X_1 \arrow[rr,"j_2",crossing over] & & X_3 \\
X_0 \arrow[uu,"f_0=\id"] \arrow[rr,"i_2"] \arrow[dr,"\id"] & & X_2 \arrow[dr,"\id"] & \\
& X_0 \arrow[rr,"i_2"] \arrow[uu,crossing over,"f_1=i_1" near end] & & X_2, \arrow[uu,"f_3=j_1"']
\end{tikzcd}
\]
where the $f_0$, $f_1$, $f_2$, $f_3$ of the statement of Loc (6) become $\id_{X_0}$, $i_1$, $\id_{X_2}$, $j_1$. Loc (6 a) therefore gives $i_1\in W\Longrightarrow j_1\in W$, \textit{i.e.} Loc (6$'$ a), and Loc (6 c) $j_1\in W\Longrightarrow i_1\in W$, \textit{i.e.} Loc (6$'$ b). For Loc (6 b) $\Longrightarrow$ Loc (6$'$ b), we use (\textit{cf.} XII, p. 78) the cube


\endgroup
% END UNIT der-ch16/der-ch16_p011-015.tex
% BEGIN UNIT der-ch16/der-ch16_p016-020.tex
\begingroup\makeatletter%
\newcommand{\printedpage}[1]{\marginpar{\footnotesize p.~#1}}%
\newcommand{\msp}[1]{\marginpar{\footnotesize [#1]}}%
\newcommand{\Cof}{\mathrm{Cof}}%
\newcommand{\Fl}{\mathrm{Fl}}%
\newcommand{\Mono}{\mathrm{Mono}}%
\newcommand{\univ}{\mathrm{univ}}%
\newcommand{\id}{\mathrm{id}}%
\makeatletter%
\renewcommand\@makefnmark{\hbox{\textsuperscript{\normalfont(\@thefnmark)}}}%
\renewcommand\@makefntext[1]{\parindent 1em\noindent\hbox{\textsuperscript{\normalfont\@thefnmark}}\,#1}%
\makeatother%
\makeatother\let\printedpage\derpp\let\msp\dermsp%
\let\mspin\DERmspin%


\noindent\printedpage{14}% head: A. Grothendieck : Dérivateurs, Ch. XVI, Édition des Universités de Montpellier II et Paris VII
% foot: Version 18 mars 2019
\[
\begin{tikzcd}[ampersand replacement=\&, row sep=small, column sep=small]
X_0 \arrow[rr, "i_2"] \arrow[dr, "i_1"] \arrow[dd, "i_1"'] \& \& X_2 \arrow[dr, "j_1"] \arrow[dd, "j_1" near start] \& \\
\& X_1 \arrow[rr, "j_2" near start, crossing over] \& \& X_3 \arrow[dd, "\id"] \\
X_1 \arrow[rr, "j_2" near end] \arrow[dr, "\id"'] \& \& X_3 \arrow[dr, "\id"'] \& \\
\& X_1 \arrow[rr, "j_2"] \arrow[from=2-2, "\id"' near start, crossing over] \& \& X_3,
\end{tikzcd}
\tag*{(Q$'$)} % checked against PDF leaf 16: id labels the front vertical X_1 to X_1, not the horizontal X_1 to X_3.
\]
where $f_0 = i_1$, $f_1 = \id_{X_1}$, $f_2 = j_1$, $f_3 = \id_{X_3}$, hence % [sic] $f_1 = j_1$ -> $f_2 = j_1$, $f_1$ is already given as $\id_{X_1}$ and $f_2$ is missing
\DERmark{ENV0939}

\[
j_1 \in W \Longrightarrow i_1 \in W, \ \ \textit{i.e.}\ \text{Loc (6 b)} \Longrightarrow \text{Loc (6$'$ b)},
\]
and if we have Loc (6 c), we see that $i_1 \in W \Longrightarrow j_1 \in W$, \textit{i.e.} Loc (6$'$ a).

\textbf{Comment 3.2.} Loc (6 a) appears to me as the ``vital minimum'' for establishing the existence of a $W$-homotopic theory of cofibers in Cat. For this, one must indeed have at hand a serviceable criterion ensuring that an arrow $i : A \longrightarrow B$ in Cat is in $\Cof_W$.

\msp{page 19}We find, when we have Loc (6 a) (in dual form, \textit{i.e.} for the \textit{closed ones,} it's all the same), that the \textit{open Dwyer immersions} are in $\Cof_W$, which suffices to have \textit{a good theory of $W$-homotopic cofibers, in which MV squares are $W$-cocartesian} (something we absolutely insist on). % [sic] a laquelle -> à laquelle, missing accent

As for Loc (6 b), if I have Loc (6$'$ b) + Loc (6 a) instead, I should manage to prove Loc (6). In XII Appendix (pages 265 ff.) I prove, it seems, that (with the present notation) Loc (6$'$) implies Loc (6), and even Loc (6 c) if $W$ is stable under direct factors (condition Loc (1 bis)). But the proof is incomplete, since I tacitly assume that the Dwyer immersions are in $\Cof_W$, something I had only established using axiom W (7), in fact using Loc (6 a) (which is part of W (7)).

\msp{page 20}But here the idea comes to me of the lemma of K. S. \textsc{Brown}, cited by \textsc{Baues}, which I am going to give again here:

\textbf{Lemma 3.2.1.} (K. S. \textsc{Brown}) \textit{Let $(\mathcal{M}, W)$ be a localizer} ($W$ \textit{mildly saturated}), \textit{equipped with a set $C \subset \Fl\,\mathcal{M}$ satisfying the conditions:}
\begin{enumerate}
\item[1)] \textit{$C$ stable under composition and under cobase change} (N.B. \textit{one has to assume the existence of the amalgamated sums}
$\begin{tikzcd}[ampersand replacement=\&, baseline=(\tikzcdmatrixname-1-1.base), row sep=small, column sep=small]
X_0 \arrow[rr, "i_1"] \arrow[dr, "i_2"'] \& \& X_1 \arrow[dr] \& \\
\& X_2 \arrow[rr] \& \& X_3
\end{tikzcd}$
\textit{only if $i_1 \in C$}), \textit{contains the isomorphisms.}
\item[2)] \textit{$\mathcal{M}$ contains an initial object, $\emptyset_{\mathcal{M}}$, and for every object $X$ of $\mathcal{M}$, $\emptyset_{\mathcal{M}} \longrightarrow X$ is in $C$.} (N.B. (1) \textit{and} (2) \textit{imply that $\mathcal{M}$ is stable under finite sums.})
\item[3)] \textit{Every arrow $f$ of $\mathcal{M}$ factors as $pi$, with $i \in C$, $p \in W$.}
\item[4)] \textit{$C \cap W$ is stable under cobase change,} i.e. \textit{if}
\[
\begin{tikzcd}[ampersand replacement=\&, row sep=small, column sep=small]
X_0 \arrow[rr, "i_1"] \arrow[dr, "i_2"'] \& \& X_1 \arrow[dr] \& \\
\& X_2 \arrow[rr, "j_1"'] \& \& X_3
\end{tikzcd}
\tag*{$(*)$}
\]
\printedpage{15}% head: A. Grothendieck : Dérivateurs, Ch. XVI, Édition des Universités de Montpellier II et Paris VII
% foot: Version 18 mars 2019
\textit{is cocartesian and $i_1 \in C \cap W$, then $j_1 \in C \cap W$.} (N.B. \textit{$j_1 \in C$ is already contained in} (1), \textit{it remains to require $j_1 \in W$.})
\end{enumerate}
\textit{Under these conditions, we have $C \subset \Cof_W$.}

One can also state the lemma as follows

\textbf{Corollary 3.2.2.} \textit{Suppose that $C$ satisfies} (1), (2), (3). \textit{Then}
\msp{page 21}\textit{the following conditions on $C$ are equivalent.}
\begin{enumerate}
\item[a)] \textit{Condition} (4) \textit{above,} i.e. \textit{stability of $C \cap W$ under cobase change.}
\item[a$'$)] \textit{$C \cap W \subset W^{\univ}$.}
\item[b)] \textit{$C \subset \Cof_W$.}
\item[b$'$)] \textit{For every cocartesian square}
\[
\begin{tikzcd}[ampersand replacement=\&, row sep=small, column sep=small]
X_0 \arrow[rr, "i"] \arrow[dr, "p"'] \arrow[to=2-4, phantom, "\textit{coc.}" description] \& \& X_1 \arrow[dr, "q"] \& \\
\& X'_0 \arrow[rr, "j"] \& \& X'_1
\end{tikzcd}
\]
\textit{with $i \in C$, $p \in W$, we have $q \in W$.}
\end{enumerate}
\textsc{Proof}. We have
\[
\begin{aligned}
&\text{b} \iff \text{b}'\\
&\text{a} \iff \text{a}'
\end{aligned}
\]
since $C$ is stable under cobase change. We have b $\Longrightarrow$ a$'$ since $\Cof_W \cap W = W^{\univ}$ thanks to hypothesis 3°). We have a $\Longrightarrow$ b$'$ by Brown's lemma.

We are going to apply this to the case $\mathcal{M} = \mathrm{Cat}$, $C =$ open Dwyer immersions. Conditions 1), 2) do not depend on $W$, they are satisfied, and 3) as well as soon as $W$ is a basic localizer. As for 4), it is a consequence of the old W (7 bis), \textit{i.e.} Loc (6$'$) (we
\msp{page 22}apparently need Loc (6$'$ a \textit{and} b).) We conclude $C \subset \Cof_W$ - it's wonderful! (But I don't know how to prove the lemma \dots)

This legitimizes the proof given in XII \textit{loc. cit.}, and we get the

\textbf{Proposition 3.3.} \textit{The condition} Loc (6$'$) (\textit{cocartesian square axiom, strong form} (a) \textit{and} (b)) is equivalent \textit{to} Loc (6 a, b) (\textit{axiom of homomorphisms of cocartesian squares}) \textit{and is} implied \textit{by} Loc (6 c). \textit{If $W$ satisfies} Loc (1 bis), i.e. \textit{is stable under direct factors, then} Loc (6$'$) is equivalent \textit{to} Loc (6 c) \textit{and to} Loc (6 a b c).

\textbf{Comment 3.4.} Thus, we get a theory of $W$-homotopic cofibers \textit{either} under hypothesis Loc (6$'$) (equivalent to the more complicated Loc (6)), or under the weaker (and more complicated) hypothesis Loc (6 a).

% left margin: a vertical line is printed alongside the following paragraph
In the first case, it is a theory of a very special nature, since the $W$-cocartesian squares are ``reversible'' (\textit{i.e.} we have not only $i_1 \in W \Longrightarrow j_1 \in W$, but also the reverse implication).

There remains then the question whether Loc (6$'$ a) suffices for a theory of $W$-homotopic cofibers, or what surely amounts to the same thing, whether it already implies Loc (6 a) (hence is equivalent to it). This hangs entirely on the \textit{question}:

\printedpage{16}% head: A. Grothendieck : Dérivateurs, Ch. XVI, Édition des Universités de Montpellier II et Paris VII
% foot: Version 18 mars 2019
\msp{page 23}% left margin: a question mark followed by a double vertical line alongside the first three lines of the following paragraph
\marginpar{\footnotesize ?\,$\Vert$}does the axiom Loc (6$'$ a) (cocartesian square axiom, direct form) imply that the set of open Dwyer immersions which are in $W$ ($W$-trivial Dwyer immersions) is stable under cobase change? In fact, we see that for this it would suffice to verify the axiom of the cube (or of homomorphisms of cocartesian squares), direct form Loc (6 a), in the case where $f_0$ and $f_2$ are identities, hence to prove that then $f_1 \in W \Longrightarrow f_3 \in W$ (which would directly give $\mathrm{Dw} \subset \Cof_W$). But I see no way of proving it \dots

Suppose now that condition Loc (6 a) is satisfied, so that we have a ``good theory of $W$-homotopic cofibers'', hence also a theory of $W$-cocartesian squares (\textit{cf.} XII, p. 258-262). There are two questions that arise, concerning the validity, with these weaker hypotheses, of the developments of XII:
\begin{enumerate}
\item[1°)] Given a commutative square
\[
\begin{tikzcd}[ampersand replacement=\&, row sep=small, column sep=small]
X_0 \arrow[rr] \arrow[dr] \& \& X_2 \arrow[dr] \& \\
\& X_1 \arrow[rr] \& \& X_3
\end{tikzcd}
\]
in Cat, is it true that this square is $W$-cocartesian if and only if the arrow
\[
\text{canonical}\ \displaystyle\int
\begin{tikzcd}[ampersand replacement=\&, row sep=small, column sep=small]
\& X_1 \\
X_0 \arrow[ur] \arrow[dr] \& \\
\& X_2
\end{tikzcd}
\quad \longrightarrow \ X_3 \text{ is in } W \text{ ? }\text{\footnotemark[15]}
\]
\footnotetext[15]{Better, in all cases, when we have a diagram
$\begin{tikzcd}[ampersand replacement=\&, baseline=(\tikzcdmatrixname-2-1.base), row sep=tiny, column sep=small]
\& X_1 \\
X_0 \arrow[ur] \arrow[dr] \& \\
\& X_2
\end{tikzcd}$, its integral $\displaystyle\int$ is a $W$-homotopic amalgamated sum of said diagram (assuming Loc (6 a) and no more, and even only Loc (6$'$ a)).}
\end{enumerate}
\msp{page 24}
\begin{enumerate}
\item[2°)] Is it still true that if $A$ is in Cat, and $u : F \lhook\joinrel\longrightarrow G$ a homomorphism in $A^{\wedge}$, then $u \in \Cof_{W_A}$, and if $u \in W$, $u \in W_A^{\univ}$, hence do we have
\[
\Mono_{A^{\wedge}} \subset \Cof_{W_A}, \qquad \Mono_{A^{\wedge}} \cap W_A \subset W_A^{\univ} \text{ ? }\text{\footnotemark[16]}
\]
\footnotetext[16]{N.B. Loc (6$'$ a) suffices in place of ( Loc (6 a)).}
\end{enumerate}
But 1°) is proved in XII, prop. 20, p. 262 - I have just checked that the proof is valid under the present hypotheses \footnote[17]{But it uses in addition Loc (3 bis). But with another simpler proof that I have just added (p 264), we can do without it (and we get a more complete view of the situation).}, provided that I check that corollary 3, p. 149 in XII is valid, \textit{i.e.} % [sic] (footnote 17) vison -> vision, missing letter
\begin{enumerate}
\item[1° bis)] If
$\begin{tikzcd}[ampersand replacement=\&, baseline=(\tikzcdmatrixname-1-1.base)]
X_0 \arrow[r] \arrow[d] \& X_1 \arrow[d] \\
X_2 \arrow[r] \& X
\end{tikzcd}$
is an MV square in Cat, then
$\displaystyle\int
\begin{tikzcd}[ampersand replacement=\&, row sep=small, column sep=small]
\& X_1 \\
X_0 \arrow[ur] \arrow[dr] \& \\
\& X_2
\end{tikzcd}
\quad \longrightarrow \quad X$ is in $W$ (hence, assuming 1°), the square is $W$-cocartesian).
\end{enumerate}

\printedpage{17}% head: A. Grothendieck : Dérivateurs, Ch. XVI, Édition des Universités de Montpellier II et Paris VII
% foot: Version 18 mars 2019
To check this, I am led to check that theorem 4 in loc. cit. (p. 146) is still valid.

I have just done so, and have annotated theorem 4 and its corollaries accordingly (some results (cor. 4 and cor. 5) using W (7 bis) in its full strength, \textit{i.e.} Loc (6$'$), \textit{i.e.} Loc (6 a \textit{and} b). But corollary 3 of theorem 4 is O.K.

I have also just checked the validity of 2°), by checking the proof of the corollary of theorem 5 (XII, p. 168) and annotating it accordingly.

\msp{page 25}3.5. Let us sum up: we have
\[
\text{Loc (6 c)} \ \Longrightarrow \ \text{Loc (6)} \ \iff \ \text{Loc (6$'$)}
\tag*{(3.5.1)}
\]
and
\[
\text{If } \ W = W^{\natural}\,, \ \text{ we have } \ \text{Loc (6 c)} \ \iff \ \text{Loc (6)} \ \iff \ \text{Loc (6$'$)}.
\tag*{(3.5.2)}
\]
On the other hand
\[
\text{Loc (6 a)} \ \Longrightarrow \ \text{Loc (6$'$ a)}\,,
\tag*{(3.5.3)}
\]
and I do not know whether the reverse implication
\[
\text{Loc (6$'$ a)} \ \overset{?}{\Longrightarrow} \ \text{Loc (6 a)}
\]
is valid \footnote[18]{This is the case in the presence of Loc (8, 9).}. To establish Loc (6 a), it suffices in any case to do so when $f_0$ and $f_2$ are isomorphisms.

The axiom Loc (6 a) seems to be the ``vital minimum'' for $W$ to give rise to a ``good theory of $W$-homotopic cofibers'', hence also of $W$-cocartesian squares. More precisely, I feel like proving this:

\textbf{Proposition 3.6.} \textit{The axiom} Loc (6 a) \textit{is equivalent to the conjunction of the following two properties.}

\printedpage{18}% head: A. Grothendieck : Dérivateurs, Ch. XVI, Édition des Universités de Montpellier II et Paris VII
% foot: Version 18 mars 2019
\begin{itemize}[leftmargin=7em, labelwidth=6em]
\item[\textbf{Loc (6 A)}] Cat has ``enough $W$-cofibrations'',
\msp{page 26}\textit{i.e.} every arrow $f$ of Cat factors as $f = pi$, with $i \in \Cof_W$ and $p \in W$. (This also means that Cat equipped with $(W, \Cof_W)$ is a category with cofibrations of K. S. \textsc{Brown} - \textsc{Anderson}, and all the results of \textsc{Brown-Quillen} theory are applicable. In particular, we have a good theory of $W$-cocartesian squares.)
\item[\textbf{Loc (6 B)}] If $X$ in Cat is the union of two opens $U$, $V$, with intersection $S$, then the square
\[
\begin{tikzcd}[ampersand replacement=\&, row sep=large, column sep=large]
S \arrow[r] \arrow[d] \& U \arrow[d] \\
V \arrow[r] \& X
\end{tikzcd}
\]
is $W$-cocartesian \footnote[19]{In fact, Loc (6 B) is a consequence of Loc (6 A), which therefore is \textit{equivalent} to Loc (6 a). (\textit{Cf.} sketch of proof in 4.4.15, pages 66-68.)}.
\end{itemize}
On the other hand, condition Loc (6 A) can be replaced by the condition (apparently weaker)
\begin{itemize}[leftmargin=7em, labelwidth=6em]
\item[\textbf{Loc (6 A$'$)}] Every Dwyer morphism is in $\Cof_W$.
\end{itemize}
\textbf{Corollary 3.7.} \textit{The condition} Loc (6), i.e. Loc (6 a, b) (\textit{or again the equivalent condition} Loc (6$'$)) \textit{is equivalent to the conjunction of conditions} Loc (6 A) (\textit{or} Loc (6 A$'$)) \textit{and} Loc (6 B) \textit{above}, \textit{and of the condition}
\DERmark{ENV0940}

\begin{itemize}[leftmargin=7em, labelwidth=6em]
\item[\textbf{Loc (6 C)}] For every $W$-cocartesian square
\[
\begin{tikzcd}[ampersand replacement=\&, row sep=large]
X_0 \arrow[r, "i"] \arrow[d] \& X_1 \arrow[d] \\
X_2 \arrow[r, "j"] \& X
\end{tikzcd}
\]
in Cat, $j \in W \Longrightarrow i \in W$.
\end{itemize}
\msp{page 27}On the other hand, I suspect that Loc (6 c) is hardly ever satisfied unless $W = W^{\natural}$ (the most common case of course), in which case Loc (6 c) \textit{is equivalent} to Loc (6). There is probably no reason to include it in Loc (6), any more than $W = W^{\natural}$ in Loc (1).


\endgroup
% END UNIT der-ch16/der-ch16_p016-020.tex
% BEGIN UNIT der-ch16/der-ch16_p021-025.tex
\begingroup\makeatletter%
\newcommand{\printedpage}[1]{\marginpar{\footnotesize p.~#1}}%
\newcommand{\msp}[1]{\marginpar{\footnotesize [#1]}}%
\newcommand{\Fib}{\mathrm{Fib}}%
\newcommand{\Fibf}{\mathrm{Fibf}}%
\newcommand{\Parf}{\mathrm{Parf}}%
\newcommand{\Catfib}{\mathrm{Catfib}}%
\newcommand{\Fibsc}{\mathrm{Fibsc}}%
\newcommand{\Cat}{\mathrm{Cat}}%
\newcommand{\HOT}{\mathrm{HOT}}%
\makeatother\let\printedpage\derpp\let\msp\dermsp%
\let\mspin\DERmspin%


\printedpage{19}
% head: A. Grothendieck : Dérivateurs, Ch. XVI, Édition des Universités de Montpellier II et Paris VII

\section*{4\quad Axiom Loc (7) of $W$-fibrations}

4.1. Recall the notion of a \textit{weak $W$-fibration} (XII p. 51, 52): the strongest version of the 8 equivalent conditions for $f : X \longrightarrow S$ to be a weak fibration is this.

\begin{itemize}[leftmargin=2cm]
\item[\textbf{(Ff)}] $\forall\ S'' \longrightarrow S'$ in $\Cat/S$, with $S'$, $S''$ \textit{contractible,} the corresponding morphism $X'' = X \times_S S'' \longrightarrow X' = X \times_S S'$ is in $W$.
\end{itemize}

When $W$ is strongly saturated, then this implies the condition, apparently stronger:

\begin{itemize}[leftmargin=2cm]
\item[\textbf{(Ff$'$)}] $\forall\ S'' \longrightarrow S'$ in $\Cat/S$, such that $S'' \longrightarrow S'$ is a homotopism, $X'' \longrightarrow X'$ is in $W$.
\end{itemize}

I denote by $\Fibf_W$ the set of weak fibrations. I summarize in a diagram of inclusions the main classes of arrows in Cat that we have encountered, akin to classes of ``fibrations'' (for $W$, in a suitable sense):

\msp{page 28}

\[
\begin{tikzcd}
 & & \Fib_W \arrow[d] & & \\
\Parf_\omega \arrow[r] & \Parf_W \arrow[r] & \Fibf_W & \underbrace{\Catfib^W}_{\mathcal{F}_W} \arrow[l] & \underbrace{\Catfib^\omega}_{\mathcal{F}_\omega} \arrow[l] \ ,
\end{tikzcd}
\]
where $\Fib_W$ denotes the $W$-fibrations ($f : S' \longrightarrow S$ is in $\Fib_W$ if and only if the base change functor $X \longmapsto X' = X \times_S S'$ from $\Cat/S$ to $\Cat/S'$ is compatible with the localizers $W_S$ and $W_{S'}$), Catfib the Cat-fibrations, $\Catfib^W$ (resp. $\Catfib^\omega$) the Cat-fibrations such that the base change functors are in $W$ (resp. are in $W_\omega$, \textit{cf.} def. below) \footnote[20]{N.B. $\Catfib_0^W = \Catfib \cap \Fibf_W$, $\Catfib_0^\omega = \Catfib \cap \Fibf_{W_\omega}$}, $\Parf_W$ the $W$-perfect morphisms, \textit{i.e.} $W$-smooth and $W$-proper, finally $\Parf_\omega$ the $W_\omega$-perfect morphisms, where $W_\omega$ is the smallest localizer satisfying Loc (1, 2, 3). Among the kinds of ``fibrations'' considered, it is only this last one, together with $\Catfib^\omega$, that is independent of $W$.

As such, it seems to me to play a role similar, dual, to that of Dwyer morphisms, and I presume that one can make them play a similar role. In any \msp{page 29} case, we know that every arrow $f$ of Cat factors as $pi$, with $i \in W_\omega$ and $p \in \Parf_\omega$, hence $i \in W$ for every $W$ satisfying Loc (1, 2, 3). We also know that $\Catfib^\omega$ satisfies the strong factorization statement, $f$ factors as $f = pi$, $i \in W_\omega$, $p \in \Catfib_0^\omega$. \footnote[21]{Thus the factorization property holds for all these classes, except perhaps $\Fib_W$.}

Let us also note that $\Parf_\omega$ and $\Parf_W$, as well as $\Fib_W$, are stable not only under base change (this is the case for the 5 classes considered, which moreover contain all isomorphisms), but also under composition. (Moreover, $\Fib_W$ is stable under direct factors, and I presume that the same holds for $\mathrm{Liss}_W$ and for $\mathrm{Prop}_W$, hence for $\Parf_W$, in particular for $\Parf_{W_\omega}$, but I have not taken the trouble to check it.) On the other hand, a priori $\Fibf_W$ is not in general stable under composition, and $\Catfib^W$, $\Catfib^\omega$ have no reason
% foot: Version 18 mars 2019
\printedpage{20}
% head: A. Grothendieck : Dérivateurs, Ch. XVI, Édition des Universités de Montpellier II et Paris VII
either to be stable under composition \footnote[22]{this should be made explicit by examples.}.

4.2. Finally, it seems more ``illuminating'' to introduce, instead of the well-determined basic localizer $W_\omega$, a basic localizer % [sic] déterminée -> déterminé, accord avec « localiseur » (masculin)
\[
W_0 \subset \mathrm{Fl}(\Cat),
\]
which will play the role of a sort of ``reference localizer'' for formulating certain properties of $W$. We will always assume that
\[
W \supset W_0.
\]

\msp{page 30}

Depending on the case, $W_0$ may be $W_\omega$, or some other basic localizer defined as the ``smallest'' satisfying given stability axioms. Or it may be $\mathbf{W}_\infty$ (the usual weak equivalences), when one is interested in localizers that contain $\mathbf{W}_\infty$ (the only ones, apparently, with which one has worked up to now in homotopy). It may also be $W$, and the statements obtained by specializing the results obtained to the case where $W = W_0$ are by no means tautological. The diagram on page 28 is then replaced by the following diagram of inclusions, more illuminating for our purpose:
\[
\begin{tikzcd}
 & \Fib_W \arrow[d] & \\
\Parf_W \arrow[r] & \Fibf_W & \mathcal{F}_W \arrow[l] \\
\Parf_{W_0} \arrow[r] \arrow[u] & \Fibf_{W_0} \arrow[u] & \mathcal{F}_{W_0} \arrow[l] \arrow[u] \\
 & \Fib_{W_0}, \arrow[u] &
\end{tikzcd}
\tag*{(4.2.1)}
\]
where for every basic localizer $W$, I put for short
\[
\mathcal{F}_W = \Catfib^W,
\]
which thus also defines $\mathcal{F}_{W_0}$. The class $\Fib_W$ \msp{page 31} should be regarded as the focus of our attention; most of the properties of $W$ that we are going to review here, will be expressed by an inclusion of the form
\[
\mathcal{F} \subset \Fib_W,
\]
where $\mathcal{F}$ is one of the seven other classes of diagram (4.2.1); with the exception at most of the case $\mathcal{F} = \Fib_{W_0}$, \textit{i.e.} of the inclusion $\Fib_{W_0} \subset \Fib_W$. As a property of $W$, this inclusion seems very weak and for that reason not of much use, except of course in the case where $\Fib_{W_0} \longrightarrow \Fibf_{W_0}$ is an equality, so that the inclusion considered is equivalent to $\Fibf_{W_0} \subset \Fib_W$ (and implies that $\Parf_{W_0}$ and $\mathcal{F}_{W_0}$ are also contained in $\Fib_W$). This on the other hand, as we shall see, is a fairly strong and very useful property of $W$, implying in particular the fact that $(\Cat, W, \Fib_W)$ is a category with fibrations.
% foot: Version 18 mars 2019

\printedpage{21}
% head: A. Grothendieck : Dérivateurs, Ch. XVI, Édition des Universités de Montpellier II et Paris VII

\msp{page 32}

For purposes of later reference, I am going to formulate now the various main aspects of \textit{the $W$-fibration axiom} Loc (7). The formulation which I had arrived at (\textit{cf.} I, VI, XII) used the notion which I now call ``\textit{weak} $W$-fibrations'' (the former $W$-fibration, now dethroned!), and in its strongest form was written as the inclusion % [sic] W-fibrations faible -> W-fibration faible, accord du nombre
\[
\Fibf_W \subset \Fib_W.
\]
This notion of weak $W$-fibration, just like that of $W$-perfect morphism, however useful they may be, nevertheless seems to me a little technical, and I prefer to give a formulation that calls only on more familiar notions, more elementary in appearance. This is why I am going to replace $\Fibf_W$ by the subclass $\mathcal{F}_W = \Catfib^W$, which (we shall see) in fact comes to the same thing.

\msp{page 33}

This class $\mathcal{F}_W$ is extremely natural, and even imposes itself in the yoga of derivators. Indeed, in the construction of $\HOT_W(S)$ for $S$ in Cat, one passes to a category of fractions of the category
\[
\Cat(S) = \underline{\mathrm{Hom}}(S^\circ, \Cat),
\]
\DERmark{ENV0941}

by a localizer $W(S)$ (which we had generally denoted by $W_S^{\mathrm{g}}$ or $W_S^{\mathrm{u}}$ \ldots). This category $\Cat(S)$ is on the other hand isomorphic to the category
\[
\Fibsc\, S
\]
of fibered and split categories over $S$, with as arrows the $S$-functors which respect the splitting (and in particular cartesian) (\textit{cf.} VI, and especially the diagram p. 175). We have seen (\textit{loc. cit.}) that the localized category $\HOT_W(S)$ is equivalent to the localized category of the very closely related category
\[
\Fib\, S \qquad \text{(which I would now prefer to denote by $\Catfib\, S$ or $\Catfib_S$),}
\]
formed by the \textit{fibered} categories over $S$, with as \msp{page 34} arrows only the cartesian functors. (One should note however that the canonical functor $\Fibsc\, S \longrightarrow \Fib\, S$ is faithful but not fully, nevertheless it induces an equivalence on the localized categories, defining in the obvious way, by $W$-equivalence fiber by fiber, \textit{i.e.} as $W_S^{\mathrm{u}}$, the relevant localizer on $\Fib\, S$.) Thus for this reason of principle, and moreover throughout the development of the theory of $\HOT_W$, the category $\Fib\, S$ and thereby, the fibrant functors, have played an essential role. The class of arrows which we denote here by $\mathcal{F}_W$ or $\Catfib^W$ (one could also denote it by $\Fib^W$, were it not for the risk of confusion with $\Fib_W$) is then interpreted very naturally starting from there. An arrow $f : X \longrightarrow S$ is in $\mathcal{F}_W$ if and only if it is a) fibrant, and thereby defines in a particularly direct way an object of $\HOT_W(S)$, since up to $S$-equivalence, $X/S$ comes from an object of $\Fibsc\, S$ or again of $\Cat(S)$, \msp{page 35} and b) the object in $\HOT_W(S)$ defined by $f$ is in $\HOT_W^{\mathrm{lc}}(S)$, \textit{i.e.} in the essential image of the strictly full subcategory $\Fibsc_0^W S$ of $\Fibsc\, S$ (under the canonical functor
% foot: Version 18 mars 2019
\printedpage{22}
% head: A. Grothendieck : Dérivateurs, Ch. XVI, Édition des Universités de Montpellier II et Paris VII
$\Fibsc \longrightarrow \HOT_W(S)$), or what comes to the same thing (at least if $W$ is strongly saturated) that the functor % [sic] saturée -> saturé, accord avec W (« W fortement saturé » ailleurs)
\[
\begin{aligned}
S^{\circ} \quad &\longrightarrow \quad \mathrm{Hot}_W \\
s \quad &\longmapsto \quad \mathrm{hot}_W(X_s)
\end{aligned}
\]
be a ``locally constant'' functor. Thus, we shall regard the $f \in \mathcal{F}_W$ as the most economical and most natural arrows in Cat, which define locally constant relative $W$-homotopy types.

Let us also note that all the classes of arrows (4.2.1) are stable under base change (and they contain the isomorphisms). The categories $\Fib_W$, $\Parf_W$ and their $W_0$ variants are also stable under composition, on the other hand, this is not the case for $\Fibf_W$ and $\mathcal{F}_W$ and their $W_0$ variants, except under special axioms on $W$ or on $W_0$. It follows at once \msp{page 36} that if $\mathcal{F}$ is one of these classes, the condition % [sic] « sur W_0). » -> parenthèse fermante sans parenthèse ouvrante correspondante
\DERmark{ENV0942}

\[
\mathcal{F} \subset \Fib_W
\]
is equivalent to this one:

\begin{itemize}[label={},leftmargin=3.5cm]
\item For every cartesian diagram
\[
\begin{tikzcd}
X \arrow[d, "f"'] & X' \arrow[l, "q"'] \arrow[d, "f'"] \\
S & S', \arrow[l, "p"]
\end{tikzcd}
\tag*{(4.2.2)}
\]
with $f \in \mathcal{F}$ and $p \in W$, we have $q \in W$.

(N.B. We also have $f' \in \mathcal{F}$, by the stability recalled above.)
\end{itemize}

4.3. Here then (\textit{cf.} below) is the relevant axiom, Loc (7) with its three aspects a, b, c \footnote[23]{I finally added a fourth F to it (for ``fibration'').}. We have seen (and shall see again) that we have the implications
\[
\text{a} \qquad \Longrightarrow \qquad \text{b} \qquad \Longrightarrow \qquad \text{c} \qquad (\Longrightarrow \qquad \text{F}).
\]
When we refer to Loc (7) without further specification, it will be understood that this is the strongest form Loc (7 a), which thus implies the other two. No less is needed to be completely at ease. But I recall that the weakened variant Loc (7b) is sufficient to develop the theory of the functors $f_!$ for the categories $\HOT_W(S)$ \footnote[24]{Here my memory is playing a trick on me - for this one has no need at all of the already very strong axiom Loc (7b), but only of Loc (3 bis), which does seem to be satisfied for all ``reasonable'' basic localizers. % [sic] qu'il soit vérifiée -> qu'il soit vérifié, accord avec « il »
} (\textit{cf.} VI).

\msp{page 37}

This is moreover a surprising thing, since via the general yoga of derivators associated with categories of models, one would rather expect that to have left adjoints $f_!$ of the inverse image functors $f^*$, one must make hypotheses on $W$ of type \textit{$W$-cofibration}
% foot: Version 18 mars 2019
\printedpage{23}
% head: A. Grothendieck : Dérivateurs, Ch. XVI, Édition des Universités de Montpellier II et Paris VII
(essentially, Loc (6 a)). It seems unlikely to me that Loc (7 b), or even Loc (7), implies Loc (6 a), or even its weaker variant Loc (6$'$ a) (axiom of the cocartesian square, or of Mayer-Vietoris, weak form). If on the other hand these axioms are not \textit{satisfied,} whereas Loc (7) is, there is a chance that this will give an example where there is a formalism of the $f_!$, without the left exactness axiom being satisfied.

As for the form Loc (7 c), weaker still, it is nevertheless sufficient (as we shall see below) to ensure that Cat has enough $W$-fibrations \msp{page 38} (a result which is by no means obvious, and could only become apparent to me after I had learned of the substantial lemma of K. S. \textsc{Brown}, in his beautiful work where he introduces categories with fibrations and with cofibrations).

We will also give below some very simple cases, where the condition Loc (7 b) already implies Loc (7) ($W$ strongly saturated, \textit{i.e.} Loc (1 ter), and satisfying Loc (3 bis)), or even where Loc (7 c) implies Loc (7) (the case where $W$ satisfies the direct limit axiom Loc (8)). I point out in passing that I know of no basic localizer satisfying the axiom Loc (8) (which already implies $W = W^{\natural}$, \textit{i.e.} Loc (1 bis)), without $W$ being strongly saturated (\textit{i.e.} even satisfying Loc (1 ter)).
\DERmark{ENV0943}


\begin{itemize}[label={},leftmargin=0.6cm]
\item \textbf{Loc (7)}
\begin{itemize}[leftmargin=1.2cm]
\item[a)] We have $\mathcal{F}_W \subset \Fib_W$, \textit{i.e.} for every cartesian diagram (4.2.2) in Cat, with $f \in \mathcal{F}_W$ (\textit{i.e.} $f$ Cat-fibrant and weakly $W$-fibrant) and $p \in W$, we have $q \in W$.
\item[b)] For every commutative diagram
\[
\begin{tikzcd}[column sep=small]
X \arrow[rr, "f"] \arrow[dr, "p"'] & & Y \arrow[dl, "q"] \\
 & S &
\end{tikzcd}
\]
\msp{page 39}
with $p, q \in \mathcal{F}_W$ and $f \in W$, we have $f \in W_S^{\mathrm{u}}$ (or, what comes to the same thing since $X$, $Y$ are $W$-smooth over $S$, $f_s \in W$ for every $s \in \mathrm{Ob}\, S$) \footnote[25]{\textbf{N.B.} The proof of $B(\mathcal{F}_W) \Longrightarrow B(\Fibf_W)$ in fact shows that it suffices to postulate b) when $f$ is cartesian, and even when $X$, $Y$ are split over $S$ and $f$ compatible with the splitting.}.
\item[c)] We have
\[
\mathcal{F}_W \cap W \subset W^{\mathrm{u}},
\]
in other words, for every arrow $f : X \longrightarrow S$ in Cat such that $f \in \mathcal{F}_W$ (one could call these arrows Cat-$W$-fibrations, it being understood that they are also Cat-fibrations, but not necessarily $W$-fibrations), if $f \in W$, then $f \in W_S^{\mathrm{u}}$ (or, what comes to the same thing, the fibers $X_s$ of $f$ ($s \in \mathrm{Ob}\, S$) are $W$-aspheric).
\end{itemize}
\end{itemize}
% foot: Version 18 mars 2019


\endgroup
% END UNIT der-ch16/der-ch16_p021-025.tex
% BEGIN UNIT der-ch16/der-ch16_p026-030.tex
\begingroup\makeatletter%
\newcommand{\printedpage}[1]{\marginpar{\footnotesize p.~#1}}%
\newcommand{\msp}[1]{\marginpar{\footnotesize [#1]}}%
\newcommand{\cF}{\mathcal{F}}%
\newcommand{\cM}{\mathcal{M}}%
\newcommand{\Cat}{\mathrm{Cat}}%
\newcommand{\Fib}{\mathrm{Fib}}%
\newcommand{\Fibf}{\mathrm{Fibf}}%
\newcommand{\Parf}{\mathrm{Parf}}%
\newcommand{\Fl}{\mathrm{Fl}}%
\newcommand{\Comp}{\mathrm{Comp}}%
\newcommand{\Catfib}{\mathrm{Catfib}}%
\newcommand{\dottedimplies}{\mathrel{\tikz[baseline=-0.6ex]{\draw[dotted, double, double distance=1.5pt, -{Implies}] (0,0) -- (0.9,0);}}}%
\makeatother\let\printedpage\derpp\let\msp\dermsp%
\let\mspin\DERmspin%


\printedpage{24}
% head: A. Grothendieck : Dérivateurs, Ch. XVI, Édition des Universités de Montpellier II et Paris VII
% foot: Version 18 mars 2019

\begin{itemize}[leftmargin=4em]
\item[F)] \footnote[26]{added after the fact.} Cat has sufficiently many $W$-fibrations, \textit{i.e.} Cat equipped with $(W, \Fib_W)$ is a category with fibrations (proper, and even a Brown one, \textit{i.e.} $X \longrightarrow e$ is in $\Fib_W$ for every $X$ in Cat). % [sic] (propre, et même de Brown, -> parenthesis never closed, closing ")" apparently missing
\end{itemize}
\DERmark{ENV0944}\DERmark{ENV0945}


It is clear that c is the special case of b obtained by taking $Y = S$, $q = \mathrm{id}_S$, so b $\Longrightarrow$ c. The implication a $\Longrightarrow$ b is less evident, it is established first in XII, and we will return to it below (4.4.2 below). The rest of the present paragraph is devoted to studying various variants of the preceding conditions a) b) c) and the implications between them.

\msp{page 40}

4.4. Let, more generally,
\[
\cF \subset \Fl(\Cat)
\]
be a set of arrows in Cat, thinking more particularly of the case where $\cF$ is one of the seven sets in diagram (4.2.1) (p. 30) other than $\Fib_W$. We can formulate three properties for $\cF$ similar to those formulated in Loc (7 a b c) in the case of $\cF_W$, let us denote these properties
\[
\mathrm{A}(\cF)\,,\ \mathrm{B}(\cF)\,,\ \mathrm{C}(\cF)\,,\ \text{\footnotemark[27]}
\]
\footnotetext[27]{\textbf{N.B.} These conditions make sense if $\cF \subset \Fl(\cM)$, if $(\cM, W)$ is an arbitrary category of models.}
thus
\[
\left\{
\begin{array}{lcl}
\mathrm{A}(\cF) & : & \cF \subset \Fib_W. \\[1ex]
\mathrm{B}(\cF) & : & \Fl(\cF_S) \cap W_S \subset W_S^{\mathrm{u}} \\
 & & \text{(where $\cF_S$ denotes the full subcategory of $\Cat/S$} \\
 & & \text{consisting of the $f : X \longrightarrow S$ with $f \in \cF$).} \\[1ex]
\mathrm{C}(\cF) & : & \cF \cap W \subset W^{\mathrm{u}}.
\end{array}
\right.
\tag*{(4.4.1)}
\]
It is clear that we have $\mathrm{B}(\cF) \Longrightarrow \mathrm{C}(\cF)$, and I am going to give a condition on $\cF$ which implies that we also have $\mathrm{A}(\cF) \Longrightarrow \mathrm{B}(\cF)$.

\msp{page 41}

We will apply it in the following cases: % [sic] (Nous -> opening parenthesis never closed in this paragraph, closing ")" apparently missing
\DERmark{ENV0946}

\[
\cF \;=\; \cF_W,\ \Fibf_W,\ \Parf_W,\ \cF_{W_0},\ \Fibf_{W_0}.
\tag*{(4.4.2)}
\]
I point out moreover that we will see below (cor. 4.4.10) that
\[
\mathrm{A}(\cF_W) \iff \mathrm{A}(\Fibf_W), \qquad \mathrm{A}(\cF_{W_0}) \iff \mathrm{A}(\Fibf_{W_0}),
\tag*{(4.4.3)}
\]
and likewise for the variants B, C, so that for the implication we want to establish, there is no need to distinguish the cases of $\cF_W$ and $\Fibf_W$, of $\cF_{W_0}$ and $\Fibf_{W_0}$. Moreover it will suffice to look at the sets $\cF_{W_0}$, $\Fibf_{W_0}$, $\Parf_{W_0}$, since the case of the sets relative to $W$ is a special case of these (taking $W_0 = W$). \footnote[28]{This will thus make a priori $3 \times 6 = 18$ different conditions, but because of (4.4.3) there are (up to equivalence) only $4 \times 3 = 12$, and in fact two among these four triples of conditions reduce to a single one (for $\cF = \Parf_W$ and $\cF = \Parf_{W_0}$), so that in all only \textit{eight} different conditions remain, which we summarize in the summary diagram (4.4.7.1) on page 48. There are even only \textit{four,} $\mathrm{A}_{W_0}$, $\mathrm{B}_{W_0}$, $\mathrm{C}_{W_0}$, $\mathrm{D}_{W_0}$, as functions of $W_0$, the other four A, B, C, D are obtained by specializing to the case where $W_0 = W$.}
\DERmark{ENV0947}\DERmark{ENV0948}


\printedpage{25}
% head: A. Grothendieck : Dérivateurs, Ch. XVI, Édition des Universités de Montpellier II et Paris VII
% foot: Version 18 mars 2019

\textbf{Lemma 4.4.1.} \footnote[29]{Holds in an arbitrary category of models.} \textit{Suppose that $\cF$ satisfies the following condition: every arrow $f$ of} Cat \textit{factors as $f = pi$, with $i \in W$, $p \in \cF$. Then the condition} $\mathrm{A}(\cF)$ \textit{on $\cF$ implies that in} Cat \textit{there are ``enough $W$-fibrations'', and consequently} (\textit{by a formal argument well} % [sic] f \in \cF -> p \in \cF, the factorisation is f = pi and lemme 4.4.3 3) restates this condition with p \in \cF
\DERmark{ENV0949}

\msp{page 42}
\textit{known}), \textit{that we have} $\mathrm{B}(\Fib_W)$, \textit{and a fortiori} (\textit{since by hypothesis} $\cF \subset \Fib_W$) $\mathrm{B}(\cF)$.

\textbf{N.B.} Indeed it is clear that if we have $\cF \subset \cF'$, then
\[
\mathrm{A}(\cF') \Longrightarrow \mathrm{A}(\cF) \qquad \text{and likewise for B, C.}
\tag*{$(*)$}
\]
Now the preliminary condition on $\cF$ is indeed satisfied in the case of a set of one of the forms $\cF \;=\; \cF_W$, $\Fibf_W$, $\Parf_W$. It suffices to see this in the case of the subsets $\cF_W$, $\Parf_W$ of $\Fibf_W$, \textit{i.e.} that every arrow $f$ of Cat factors as
\[
f = pi,
\]
with $i \in W$, and $p$ being, as desired, in $\Parf_W$, or in $\cF_W$. Now in the case of $\Parf_W$, \marginpar{$\|$}this is what is established in VII (Categories of paths (2)) thanks to the formalism of the $\underline{\mathrm{Ch}}(S) \longrightarrow S \times S$. To find a factorization with $f = qj$, $j \in W$, $q \in \cF_W$, it suffices to see that $p$ factors in this way as $p = q_0 j_0$, we will have % [sic] établie -> établi, agreement with "ce qui"
\[
f = \underbrace{q_0}_{\in \cF_W} \underbrace{(j_0 i)}_{\in W},
\]
\DERmark{ENV0950}\DERmark{ENV0951}\DERmark{ENV0952}

so we can assume $f \in \Parf_W$. But if $X$ is $W$-perfect over $S$, hence a
\msp{page 43}
weak $W$-fibered object over $S$, then in the factorization
\begin{center}
\begin{tikzcd}[column sep=small]
X \arrow[rr, "i"] \arrow[dr, "f"'] & & \Phi_S(X) \arrow[dl, "p"] \\
 & S &
\end{tikzcd}
\end{center}
with $i \in W$ and $p \in \Catfib$ (\textit{cf.} VI), we know that $p$ is also a weak $W$-fibration, \textit{i.e.} $p \in \cF_W$, whence the conclusion. Thus we find the

\textbf{Corollary 4.4.2.} \textit{Suppose that $\cF$ is one of the six sets} (4.4.2) (i.e. \textit{one of the sets in diagram} (4.2.1), \textit{excluding $\mathit{Fib}_W$ and $\mathit{Fib}_{W_0}$}). \textit{Then we have}
\[
\mathrm{A}(\cF) \;\Longrightarrow\; \mathrm{B}(\cF) \;\Longrightarrow\; \mathrm{C}(\cF)\,.
\]

\textbf{Lemma 4.4.3.} \textit{Let $(\cM, W)$ be a category of models, and let $\cF \subset \Fl(\cM)$ satisfy the following conditions}
\begin{itemize}[leftmargin=3em]
\item[1)] \textit{$\cF$ stable under base change, composition, contains the isomorphisms.}
\item[2)] \textit{$\cM$ has a final object, and for every $X \in \mathrm{Ob}\,\cM$, $X \longrightarrow e$ is in $\cF$.}
\item[3)] \textit{Every arrow $f$ of $\cM$ factors as $f = pi$, with $i \in W$, $p \in \cF$} (\textit{this is the condition of lemma} 4.4.1).
\end{itemize}

\printedpage{26}
% head: A. Grothendieck : Dérivateurs, Ch. XVI, Édition des Universités de Montpellier II et Paris VII
% foot: Version 18 mars 2019
\noindent\textit{Then the conditions} $\mathrm{A}(\cF)$, $\mathrm{B}(\cF)$, $\mathrm{C}(\cF)$ \textit{on $\cF$ are equivalent.}

Indeed, by 4.4.1, it remains to prove the implication
\[
\mathrm{C}(\cF) \;\Longrightarrow\; \mathrm{A}(\cF),
\]

\msp{page 44}

\textit{i.e.} this: if $\cF$ also satisfies the condition
\begin{itemize}[leftmargin=3em]
\item[4)] $\cF \cap W \subset W^{\mathrm{u}}$ (or what amounts to the same thing, since $\cF$ is stable under base change, that $\cF \cap W$ is stable under base change),
\end{itemize}
then we conclude from this
\[
\cF \subset \Fib_W.
\]
Now this is none other than the lemma of K. S. \textsc{Brown}, in the form dual to the one given above (3.2.1, page 20).

Among the six sets of arrows considered in (4.4.2), the hypotheses 1) 2) 3) are satisfied for $\Parf_W$ and $\Parf_{W_0}$, to the exclusion of the other four. (The others satisfy all the conditions, with the sole exception of stability under composition \footnote[30]{\textbf{N.B.} For $\Fib_{W_0}$, it is the conditions 1) 2) that hold, and on the other hand, condition 3) that is a problem. Let us denote it by $\mathrm{F}(W_0)$. If it is satisfied (\textit{i.e.} if Cat has enough $W_0$-fibrations), then we also have equivalence of the conditions $\mathrm{A}(\Fib_{W_0})$, $\mathrm{B}(\Fib_{W_0})$, $\mathrm{C}(\Fib_{W_0})$.}). We thus find:

\textbf{Corollary 4.4.4.} \textit{The conditions} $\mathrm{A}(\Parf_W)$, $\mathrm{B}(\Parf_W)$, $\mathrm{C}(\Parf_W)$ \textit{are equivalent} (we denote this condition on $W$ by D). \textit{Likewise the conditions} $\mathrm{A}(\Parf_{W_0})$, $\mathrm{B}(\Parf_{W_0})$, $\mathrm{C}(\Parf_{W_0})$ \textit{are equivalent.}

\msp{page 45}

(We will denote the condition obtained on the pair $W \supset W_0$ of basic localizers by $D_{W_0}$.)

Thus we have (since $\Parf_{W_0} \subset \Parf_W$)
\[
\mathrm{D} = \mathrm{D}_W \;\Longrightarrow\; \mathrm{D}_{W_0}\,,
\tag*{(4.4.4)}
\]
and if $W'_0 \subset W_0$, we will have
\[
\mathrm{D}_{W_0} \;\Longrightarrow\; \mathrm{D}_{W'_0}\,.
\tag*{(4.4.5)}
\]

\textbf{Corollary 4.4.5.} \textit{Let us denote by} A, B, C \textit{respectively the conditions} Loc (7 a), Loc (7 b), Loc (7 c). \textit{Then we have the implications}
\[
\mathrm{A} \;\Longrightarrow\; \mathrm{B} \;\Longrightarrow\; \mathrm{C}\,.
\tag*{(4.4.5.1)}
\]
\textit{Moreover, if we denote} (\textit{for a set $\cF$ of arrows of} Cat) \textit{by} $\Comp(\cF)$ \textit{the condition on $\cF$ that $\cF$ be stable under composition, we have the equivalences}
\[
\mathrm{A} \iff (\mathrm{C} + \Comp(\cF_W)) \iff (\mathrm{C} + \Comp(\Fibf_W)).
\tag*{(4.4.5.2)}
\]

\printedpage{27}
% head: A. Grothendieck : Dérivateurs, Ch. XVI, Édition des Universités de Montpellier II et Paris VII
% foot: Version 18 mars 2019

\textsc{Proof}. The first assertion is contained in corollary 4.4.2, applied to the case $\cF = \cF_W$. On the other hand, if $\cF_W$ satisfies $\Comp(\cF_W)$, it satisfies conditions 1) 2) 3) of Brown's lemma 4.4.3, so combined with C, it implies A. It remains, for the first equivalence, to prove that
\[
A \;\Longrightarrow\; \Comp(\cF_W).
\]

\msp{page 46}

But by definition we have
\[
\cF_W = \Catfib \cap \Fibf_W
\tag*{$(*)$}
\]
and we know that Catfib is stable under composition. So $\cF_W$ is too, provided that the composite $gf$ of two arrows in $\cF_W$ is in $\Fibf_W$. Now by hypothesis A, which means $\cF_W \subset \Fib_W$, we have $f, g \in \Fib_W$, hence $gf \in \Fib_W$, since $\Fib_W$ is stable under composition, a fortiori $gf \in \Fibf_W$.

On the other hand, the preceding formula $(*)$ shows that $\Comp(\Fibf_W) \Longrightarrow \Comp(\cF_W)$, so we get for free the implication $\Longleftarrow$ in the second equivalence (4.4.5.2). The implication in the opposite direction then amounts to
\[
\mathrm{A} \;\Longrightarrow\; \Comp(\Fibf_W),
\]

\msp{page 47}

which will be evident thanks to the result (given below)
\[
\mathrm{A} = \mathrm{A}(\cF_W) \iff \mathrm{A}(\Fibf_W) \qquad \text{and likewise for B, C,}
\]
which thus means that
\[
\cF_W \subset \Fib_W \qquad \text{implies} \qquad \Fibf_W \subset \Fib_W \qquad (\text{so } \Fibf_W = \Fib_W),
\]
now $\Fib_W$ being stable under composition, the same is true of $\Fibf_W$, q.e.d.

\textbf{Corollary 4.4.6.} \textit{Suppose that $\cF_{W_0}$ is stable under composition, which is the case in particular} (\textit{by} 4.4.5.2 \textit{applied to $W_0$ instead of $W$}) \textit{if $W_0$ satisfies} Loc (7). \textit{Then the conditions} $\mathrm{A}(\cF_{W_0})$, $\mathrm{B}(\cF_{W_0})$, $\mathrm{C}(\cF_{W_0})$ (\textit{which we also denote simply by} $\mathrm{A}_{W_0}$, $\mathrm{B}_{W_0}$, $\mathrm{C}_{W_0}$) \textit{are equivalent.} (\textbf{\textit{N.B.}} \textit{If we do not make the preceding hypothesis on $W_0$, we will only have} $\mathrm{A}_{W_0} \Longrightarrow \mathrm{B}_{W_0} \Longrightarrow \mathrm{C}_{W_0}$, \textit{by} 4.4.2.)

4.4.7. We have now gone through all the conditions of ``fibration'' type on $W$, or on the pair $(W, W_0)$, which have seemed useful to me. I am going to lay them out here in a (more or less) summary diagram: % [sic] parues -> paru, agreement of the past participle

\msp{page 48}

\printedpage{28}
% head: A. Grothendieck : Dérivateurs, Ch. XVI, Édition des Universités de Montpellier II et Paris VII
% foot: Version 18 mars 2019

\begin{center}
\begin{tikzcd}[column sep=2.3em, row sep=small]
 & \overbrace{\mathrm{A}_{W_0}}^{\cF_{W_0},\,\Fibf_{W_0}} \arrow[dd, Rightarrow] & & \\
\overbrace{\mathrm{A}}^{\cF_W,\,\Fibf_W} \arrow[ur, Rightarrow] \arrow[dd, Rightarrow] & & & \\
 & \mathrm{B}_{W_0} \arrow[dd, Rightarrow] & & \\
\mathrm{B} \arrow[ur, Rightarrow] \arrow[dd, Rightarrow] & & & \\
 & \mathrm{C}_{W_0} \arrow[dr, Rightarrow, "\bullet"] \arrow[uuuu, Rightarrow, dotted, bend right=40, "\text{\scriptsize if } W_0\ \text{\scriptsize (7)}"'] & & \\
\mathrm{C} \arrow[ur, Rightarrow] \arrow[dr, Rightarrow, "\bullet"] & & \overbrace{\mathrm{D}_{W_0}}^{\Parf_{W_0}} \arrow[r, Rightarrow, "\bullet"] & \mathrm{F} : (\Cat, W, \Fib_W) \text{ a category with fibrations.} \\
 & \overbrace{\mathrm{D}}^{\Parf_W} \arrow[ur, Rightarrow] & &
\end{tikzcd}
\end{center}

\begin{center}
\begin{tikzcd}[row sep=large]
 & \mathrm{A} \arrow[r, Rightarrow] \arrow[d, Rightarrow] & \mathrm{A}_{W_0} \arrow[d, Rightarrow] & \\
\cF_W, \Fibf_W \smash{\left\{\vrule height 2.6cm depth 2.6cm width 0pt\right.} & \mathrm{B} \arrow[r, Rightarrow] \arrow[d, Rightarrow] & \mathrm{B}_{W_0} \arrow[d, Rightarrow] & \smash{\left.\vrule height 2.6cm depth 2.6cm width 0pt\right\}} \quad \text{\scriptsize\begin{tabular}{@{}l@{}}revised version\\ of the diagram\end{tabular}} \\
 & \mathrm{C} \arrow[r, Rightarrow] \arrow[d, Rightarrow] & \mathrm{C}_{W_0} \arrow[d, Rightarrow] & \\
 & \underbrace{\mathrm{D}}_{\Parf_W} \arrow[r, Rightarrow] & \underbrace{\mathrm{D}_{W_0}}_{\Parf_{W_0}} \arrow[r, Rightarrow] & \mathrm{F},
\end{tikzcd}
\end{center}

diagram which I am going to comment on. All the solid-line implications in the diagram have already been obtained or are tautological, except for the three $\overset{\bullet}{\Longrightarrow}$ to which I will return. The dotted implication $\mathrm{C}_{W_0} \dottedimplies \mathrm{A}_{W_0}$ is a conditional implication, subject to the condition marked on the diagram as $W_0$ (7), which means: $W_0$ satisfies Loc (7). This implication, \textit{i.e.} the equivalence of $\mathrm{A}_{W_0}$, $\mathrm{B}_{W_0}$, $\mathrm{C}_{W_0}$ under the condition $W_0$ (7), was seen in 4.4.6. I have marked, above the two groups (A, B, C) and ($\mathrm{A}_{W_0}$, $\mathrm{B}_{W_0}$, $\mathrm{C}_{W_0}$) as well as above D and D$'$ (which each correspond to such a group of conditions $\mathrm{A}(\cF)$, $\mathrm{B}(\cF)$, $\mathrm{C}(\cF)$ all equivalent, with $\cF \;=\; \Parf_W$ or $\cF = \Parf_{W_0}$) the
\msp{page 49}
sets of arrows which give rise to them respectively: it will be noted that there are two for the two triples named first, A, B, C being definable via $\cF_W$ or $\Fibf_W$ indifferently, and likewise $\mathrm{A}_{W_0}$, $\mathrm{B}_{W_0}$, $\mathrm{C}_{W_0}$ via $\cF_{W_0}$ or $\Fibf_{W_0}$, as already stated (4.4.3) and as we will see below (cor. 4.4.10, pages 54-55). On the other hand, the reduced triples D and $\mathrm{D}_{W_0}$ are each defined by just one $\cF$, namely by $\Parf_W$ and % [sic] pour les pour les -> pour les, words repeated; [sic] Fib_W / Fib_{W_0} -> Fibf_W / Fibf_{W_0}, the diagram and (4.4.3) give Fibf
\DERmark{ENV0953}\DERmark{ENV0954}\DERmark{ENV0955}



\endgroup
% END UNIT der-ch16/der-ch16_p026-030.tex
% BEGIN UNIT der-ch16/der-ch16_p031-035.tex
\begingroup\makeatletter%
\newcommand{\printedpage}[1]{\marginpar{\footnotesize p.~#1}}%
\newcommand{\msp}[1]{\marginpar{\footnotesize [#1]}}%
\newcommand{\Fib}{\mathrm{Fib}}%
\newcommand{\Fibf}{\mathrm{Fibf}}%
\newcommand{\Parf}{\mathrm{Parf}}%
\newcommand{\Fl}{\mathrm{Fl}}%
\newcommand{\inc}{\mathrm{inc}}%
\newcommand{\Ob}{\mathrm{Ob}}%
\newcommand{\dotimp}{\overset{\bullet}{\Longrightarrow}}%
\newcommand{\dottedimpl}{\mathrel{\begin{tikzcd}[ampersand replacement=\&, column sep=small, baseline=-0.5ex]{}\arrow[r,Rightarrow,dotted]\&{}\end{tikzcd}}}%
\newcommand{\circled}[1]{\tikz[baseline=(c.base)]\node[draw,circle,inner sep=1pt](c){#1};}%
\makeatother\let\printedpage\derpp\let\msp\dermsp%
\let\mspin\DERmspin%


\setcounter{footnote}{30}

\printedpage{29}
% head: A. Grothendieck : Dérivateurs, Ch. XVI, Édition des Universités de Montpellier II et Paris VII
% foot: Version 18 mars 2019
$\Parf_{W_0}$ respectively. (In all, this brings in six distinct sets $\mathcal{F}$, namely the six sets of diagram (4.2.1) other than $\Fib_W$ and $\Fib_{W_0}$.)

The non-tautological implication $\mathrm{C} \dotimp \mathrm{D}$ is a special case of the implication $\mathrm{C}_{W_0} \dotimp \mathrm{D}_{W_0}$, which is contained in what has just been announced, namely $\mathrm{C}_{W_0}(\mathcal{F}_{W_0}) \Longrightarrow \mathrm{C}_{W_0}(\Fibf_{W_0})$, taking into account that $\Parf_{W_0} \subset \Fibf_{W_0}$ (\textit{cf.} $(*)$, p. 42). Taking
\DERmark{ENV0956}

\msp{page 50}
$W_0 = W$, we thus find the implication $\mathrm{A} \Longrightarrow \mathrm{D}$, which is however contained in the other implications of the diagram, as a composite in the chain of implications
\[
\mathrm{A} \Longrightarrow \mathrm{B} \Longrightarrow \mathrm{C} \Longrightarrow \mathrm{D} \Longrightarrow \mathrm{D}_{W_0} \dotimp \mathrm{F}.
\]

It remains, in these general comments, to say a word about this last implication
\[
\mathrm{D}_{W_0} \dotimp \mathrm{F},
\]
\textit{i.e.} about the fact that the weakest of all the conditions A to $\mathrm{D}_{W_0}$ considered here already implies that Cat has enough $W$-fibrations, \textit{i.e.} is a category with fibrations for $(W, \Fib_W)$. This is clear when we take $\mathrm{D}_{W_0}$ in the form $\mathrm{A}(\Parf_{W_0})$, \textit{i.e.}
\DERmark{ENV0957}

\[
\Parf_{W_0} \subset \Fib_W.
\]
\DERmark{ENV0958}

As every arrow $f$ of Cat factors as $f = pi$, with
\[
i \in W_0 \subset W \quad\text{and}\quad p \in \Parf_{W_0} \subset \Fib_W\, , % [sic] f \in \Parf_{W_0} -> p \in \Parf_{W_0}, the factor of f = pi lying in Parf is p
\]
\DERmark{ENV0959}

we find the factorization with $i \in W$, $p \in \Fib_W$. % [sic] f \in \Fib_W -> p \in \Fib_W, the factorisation f = pi requires p \in \Fib_W
\DERmark{ENV0960}


\msp{page 51}
4.4.8. To finish establishing the implications of the diagram p. 48, it only remains to establish the equivalences
\[
\left\{
\begin{aligned}
&\mathrm{A}(\mathcal{F}_{W_0}) &&\Longleftrightarrow &&\mathrm{A}(\Fibf_{W_0})\\
&\mathrm{B}(\mathcal{F}_{W_0}) &&\Longleftrightarrow &&\mathrm{B}(\Fibf_{W_0})\\
&\mathrm{C}(\mathcal{F}_{W_0}) &&\Longleftrightarrow &&\mathrm{C}(\Fibf_{W_0}),
\end{aligned}
\right.
\tag*{(4.4.8.1)}
\]
which will give the similar equivalences for $W$, on taking $W_0 = W$.

The first equivalence (4.4.8.1) will be a consequence of the following ($\pm$ tautological) lemma:

\textbf{Lemma 4.4.9.} \textit{Let} $(\mathcal{M}, W)$ \textit{be a category of models,} $\mathcal{F}, \mathcal{G} \subset \Fl(\mathcal{M})$, \textit{such that every} $f \in \mathcal{F}$, $f : X \longrightarrow S$, \textit{can be written} $f = pi$, \textit{with} $i \in W^u_S$, $p \in \mathcal{G}$. \textit{Then} $\mathrm{A}(\mathcal{G}) \Longrightarrow \mathrm{A}(\mathcal{F})$, i.e.
\[
\mathcal{G} \subset \Fib_W \;\Longrightarrow\; \mathcal{F} \subset \Fib_W,
\tag*{(4.4.9.1)}
\]
\textit{and likewise} $\mathrm{C}(\mathcal{G}) \Longrightarrow \mathrm{C}(\mathcal{F})$, i.e.
\[
\mathcal{G} \cap W \subset W^u \;\Longrightarrow\; \mathcal{F} \cap W \subset W^u.
\tag*{(4.4.9.2)}
\]

\textsc{Proof}. A) Indeed, if $\mathcal{G} \subset \Fib_W$, the factorization $f = pi$ of $f$ is such that $p \in \Fib_W$. Now a composite $pi$, with $p \in \Fib_W$ and $i \in W^{\mathrm{u}}_S$, is in $\Fib_W$, as one verifies at once on the diagram of cartesian squares
\printedpage{30}
% head: A. Grothendieck : Dérivateurs, Ch. XVI, Édition des Universités de Montpellier II et Paris VII
% foot: Version 18 mars 2019
\[
\begin{tikzcd}
X \arrow[d, "W^{\mathrm{u}}_S"'] & X' \arrow[l] \arrow[d, "(W)"'] & X'' \arrow[l] \arrow[d, "(W)"'] \\
\overline{X} \arrow[d, "\Fib_W"'] & \overline{X'} \arrow[l] \arrow[d] & \overline{X''} \arrow[l, "(W)"'] \arrow[d] \\
S & S' \arrow[l] & S'' \arrow[l, "W"']
\end{tikzcd}
\]
on which one reads that $X'' \longrightarrow X'$ is in $W$, q.e.d.

\msp{page 52}
This lemma applies in particular to the case where $\mathcal{F} = \Fibf_{W_0}$, $\mathcal{G} = \mathcal{F}_{W_0}$, for $f$ factors as
\[
\begin{tikzcd}
X \arrow[dr, "i"] \arrow[dd, "f"'] & \\
& \Phi_S(X) \arrow[dl, "p"] \\
S, &
\end{tikzcd}
\]
and we know that $i \in W^{\mathrm{u}}_{0\,S} \subset W^{\mathrm{u}}_S$, and $p \in \mathcal{F}_{W_0} \subset \mathcal{F}_W$, which are the desired conditions.

C) Similar proof: Let $f \in \mathcal{F} \cap W$, factoring as $f = gi$, with $i \in W^{\mathrm{u}}_S$, $g \in \mathcal{G}$, we will then have $g \in W$, hence $g \in \mathcal{G} \cap W \subset W^{\mathrm{u}}$ (by hypothesis $\mathrm{C}(\mathcal{G})$), a fortiori we have $g \in W^{\mathrm{u}}_S$, hence $f \in W^{\mathrm{u}}_S$ as a composite of two arrows of $W^{\mathrm{u}}_S$, and as the target of $f$ is $S$ itself, this means $f \in W^{\mathrm{u}}$, q.e.d.
\DERmark{ENV0961}

\[
\begin{tikzcd}
X \arrow[dr, "i\in W^{\mathrm{u}}_S"] \arrow[dd, "\mathcal{F}\cap W"'] & \\
& \bar{X} \arrow[dl, "g"] \\
S &
\end{tikzcd}
\]

I do not see a way of also proving
\[
\mathrm{B}(\mathcal{G}) \qquad\quad \Longrightarrow \qquad\quad \mathrm{B}(\mathcal{F}),
\]
with only the hypotheses of the preceding lemma on the pair $\mathcal{F}$, $\mathcal{G}$. It would be necessary that for every $S$ in $\mathcal{M}$, and every arrow
\msp{page 53}
$u : X \longrightarrow Y$ in $\mathcal{M}/S$, we could insert the diagram obtained in $\mathcal{M}$
\[
\begin{tikzcd}[column sep=small]
X \arrow[rr, "u"] \arrow[dr, "f"'] & & Y \arrow[dl, "g"] \\
& S &
\end{tikzcd}
\]
\printedpage{31}
% head: A. Grothendieck : Dérivateurs, Ch. XVI, Édition des Universités de Montpellier II et Paris VII
% foot: Version 18 mars 2019
into a commutative diagram below
\[
\begin{tikzcd}[column sep=small]
X \arrow[rr, "u"] \arrow[d, "i"'] & & Y \arrow[d, "j"] \\
\bar{X} \arrow[rr, "\bar{u}"] \arrow[dr, "p"'] & & \bar{Y} \arrow[dl, "q"] \\
& S &
\end{tikzcd}
\tag*{(4.4.9.3)}
\]
with
\[
f = pi, \qquad\quad g = qj, \qquad\quad i, j \in W^{\mathrm{u}}_S, \qquad\quad p, q \in \mathcal{G}.
\]
This will in any case be so if for every $S$, and a variable object $X$ of $\mathcal{M}/S$, we can find a factorization \textit{functorial in} $X$ of its structural morphism $f_X$ as $p_X i_X$, with $i_X \in W^{\mathrm{u}}_S$ and $p_X \in \mathcal{G}$, \textit{i.e.} if $i$, $p$ are obtained from a functor
\[
\varphi : \mathcal{F}_S \longrightarrow \mathcal{M}/S \ \text{\footnotemark},
\]
\footnotetext{$\mathcal{F}_S$ is the full subcategory of $\mathcal{M}/S$ consisting of the $X$ such that $f_X \in \mathcal{F}$.}
and a morphism of functors
\[
i : \underbrace{\inc_{\mathcal{F}_S}}_{\substack{\text{inclusion}\\ \text{functor}}} \longrightarrow \varphi.
\]
Assuming the existence of a diagram having the stated properties, for every arrow $u$ in $\mathcal{F}_S$, or only every arrow $u$ which is also in $W_S$, we can now prove that $u \in W_S \Longrightarrow u \in W^{\mathrm{u}}_S$, provided that the similar property is true for $\mathcal{G}$ in place of $\mathcal{F}$.

\msp{page 54}
Indeed, the hypothesis on $u \in W$ implies $\bar{u} \in W$, and the hypothesis on $\mathcal{G}$ then implies $\bar{u} \in W^{\mathrm{u}}_S$. As $W^{\mathrm{u}}_S$ is a localizer in $\mathcal{M}/S$, we do conclude $u \in W^{\mathrm{u}}_S$, q.e.d.

Thus we have obtained

\textbf{Corollary 4.4.10.} \textit{Suppose, under the conditions of lemma} 4.4.9, \textit{that every triangular commutative diagram} $\begin{tikzcd}[column sep=small, row sep=small, ampersand replacement=\&]
X \arrow[rr, "u"] \arrow[dr, "f"'] \& \& Y \arrow[dl, "g"] \\
\& S \&
\end{tikzcd}$ \textit{can be inserted into a diagram of type} (4.4.9.3) \textit{as above} (\textit{with the explicit relations on the arrows occurring in it}). \textit{Then}
\[
\mathrm{B}(\mathcal{G}) \;\Longrightarrow\; \mathrm{B}(\mathcal{F}).
\]
\textit{Thus if moreover} $\mathcal{G} \subset \mathcal{F}$, \textit{so that} $\mathrm{A}(\mathcal{F}) \Longrightarrow \mathrm{A}(\mathcal{G})$ \textit{tautologically, and likewise for} B \textit{and for} C, \textit{we find}
\[
\begin{aligned}
\mathrm{A}(\mathcal{F}) &\;\Longleftrightarrow\; \mathrm{A}(\mathcal{G})\\
\mathrm{B}(\mathcal{F}) &\;\Longleftrightarrow\; \mathrm{B}(\mathcal{G})\\
\mathrm{C}(\mathcal{F}) &\;\Longleftrightarrow\; \mathrm{C}(\mathcal{G}).
\end{aligned}
\]

We can apply this corollary in the case that concerns us here, with
\[
\mathcal{G} \overset{\mathrm{def}}{=} \mathcal{F}_{W_0} \subset \mathcal{F} \overset{\mathrm{def}}{=} \Fibf_{W_0},
\]
\DERmark{ENV0962}\DERmark{ENV0963}

\printedpage{32}
% head: A. Grothendieck : Dérivateurs, Ch. XVI, Édition des Universités de Montpellier II et Paris VII
% foot: Version 18 mars 2019
\msp{page 55}
seeing that in this case we do have a \textit{functorial} factorization of an $f \in \Fibf_{W_0}$, $f : X \longrightarrow S$, as $f = pi$, $i \in W^{\mathrm{u}}_{0\,S} \subset W^{\mathrm{u}}_S$, $p \in \mathcal{F}_{W_0} = \mathcal{G}$, namely the one recalled on page 43:
\[
\begin{tikzcd}[column sep=small]
X \arrow[rr, "i_X"] \arrow[dr, "f=f_X"'] & & \bar{X} = \Phi_S(X) \arrow[dl, "p_X"] \\
& S. &
\end{tikzcd}
\]

This thus proves the announced equivalences (4.4.8.1) and finishes establishing the validity of the implications of the diagram on page 48.

These implications, with the sole exception of the dotted implication $\mathrm{C}_{W_0} \dottedimpl \mathrm{A}_{W_0}$ (which assumes that we have $W_0$ (7)), make no appeal to any particular hypothesis on the basic localizers $W$ and $W_0$ (apart from the inclusion $W \supset W_0$). But we can, using lemma 4.4.10, give remarkable conditional implications, under the sole condition that $W$, or $W_0$, satisfy the innocuous \textit{axiom}
\msp{page 56}
\textit{of limits} Loc (8), and even only in the weak form of stability of localizers under $\varinjlim$ of \textit{sequences}. Let us first give them in technical form, before unfolding the consequences in a theorem-scholie:

\textbf{Corollary 4.4.11.} \textit{Suppose that} $W_0$ \textit{satisfies axiom} Loc (8), \textit{or only} Loc (8 bis) (\textit{stability under} $\varinjlim$ \textit{of sequences}). \textit{Then} $\mathrm{D}_{W_0} \Longrightarrow \mathrm{A}_{W_0}$ (\textit{hence the conditions} $\mathrm{A}_{W_0}$, $\mathrm{B}_{W_0}$, $\mathrm{C}_{W_0}$, $\mathrm{D}_{W_0}$ \textit{are equivalent}). \textit{In particular} (\textit{taking} $W_0 = W$), \textit{if} $W$ \textit{satisfies} Loc (8 bis), \textit{then we have} $\mathrm{D}_W \Longrightarrow \mathrm{A}_W$ (\textit{hence the conditions} $\mathrm{A}_W$, $\mathrm{B}_W$, $\mathrm{C}_W$, $\mathrm{D}_W$ \textit{are equivalent}).

Indeed, the implication to be proved can also be written
\[
\mathrm{A}(\Parf_{W_0}) \;\Longleftrightarrow\; \mathrm{A}(\Fibf_{W_0}).
\]

The non-tautological implication $\Longrightarrow$ is covered by lemma 4.4.9. We must check that assuming Loc (8 bis) for $W_0$, every $f \in \Fibf_{W_0}$, $f : X \longrightarrow S$, factors as $f = pi$, with $i \in W^{\mathrm{u}}_{0\,S}$, $p \in \Parf_{W_0}$. Now we know\footnote{and this is an entirely unexpected find, a little like the lemma of K. S. \textsc{Brown} \ldots} (VI, p. 193 ff.) % [sic] f \in \Fib_{W_0} -> f \in \Fibf_{W_0}, the implication to be proved concerns A(\Fibf_{W_0}) and the next line has \Fibf_{W_0}
\DERmark{ENV0964}\DERmark{ENV0965}\DERmark{ENV0966}

\msp{page 57}
that axiom Loc (8 bis) implies that $f \in \Fibf_{W_0}$ factors as
\[
\begin{tikzcd}[column sep=small]
X \arrow[rr, "i_\infty"] \arrow[dr, "f"'] & & \bar{X} = \Phi_\infty(X) \arrow[dl, "p_\infty"] \\
& S, &
\end{tikzcd}
\qquad
\text{\parbox[t]{0.33\textwidth}{\footnotesize (category of paths infinite on one side, relativized over $S$, \textit{i.e.} of pairs $(x,c)$, with $x \in \Ob X$ and $c : f(x) \longrightarrow s$ an infinite path in $S$, with source $f(x)$.)}}
\]
with $p_\infty \in \Parf_{W_0}$ (as we want), and $i_\infty \in W^{\mathrm{u}}_{0\,S}$, and as $W^{\mathrm{u}}_{0\,S} \subset W^{\mathrm{u}}_S$, we do have $i_\infty \in W^{\mathrm{u}}_S$.

\textbf{Corollary 4.4.12.} \textit{Suppose that} $W$ and $W_0$ \textit{satisfy} Loc (8 bis). \textit{Then we have} $\mathrm{D}_{W_0} \Longrightarrow \mathrm{A}$ (\textit{hence the eight conditions} A, B, C, D, $\mathrm{A}_{W_0}$, $\mathrm{B}_{W_0}$, $\mathrm{C}_{W_0}$, $\mathrm{D}_{W_0}$ \textit{are equivalent}).

\printedpage{33}
% head: A. Grothendieck : Dérivateurs, Ch. XVI, Édition des Universités de Montpellier II et Paris VII
% foot: Version 18 mars 2019
This means indeed that
\[
\mathrm{A}(\Parf_{W_0}) \;\Longrightarrow\; \mathrm{A}(\Fibf_W),
\]
and this is again covered by lemma 4.4.9. Indeed, the above factorization of an arrow $f$ does not depend on the choice of a localizer, and $p_\infty$ is $W$-perfect for \textit{every} localizer satisfying Loc (8 bis), hence in particular for $W_0$. And the arrow $i_\infty$ is in $W^{\mathrm{u}}_S$, since the arrow $f$ under consideration is in $\Fibf_W$.
\DERmark{ENV0967}


\msp{page 58}
I am going to recapitulate the results obtained in this section in a theorem, which will constitute the main result of the present paragraph on the axioms of $W$-fibrations.

\textbf{Theorem 4.4.13.} \textit{Let} $W$ \textit{be a basic localizer in} Cat. \textit{For every localizer} $W_0$ \textit{contained in} $W$, \textit{we have made explicit a set of four conditions on the pair} $W_0$, $W$, \textit{which for} $W$ \textit{fixed, depend on} $W_0$ \textit{alone, denoted}
\[
\mathrm{A}_{W_0},\ \mathrm{B}_{W_0},\ \mathrm{C}_{W_0},\ \mathrm{D}_{W_0},
\tag*{(4.4.13.1)}
\]
\textit{and defined respectively as}
\[
\mathrm{A}(\mathcal{F}_{W_0}),\ \mathrm{B}(\mathcal{F}_{W_0}),\ \mathrm{C}(\mathcal{F}_{W_0}), \qquad\quad \textit{and}
\]
\[
\mathrm{A}(\Parf_{W_0}) \Bigl(\Longleftrightarrow\; \mathrm{B}(\Parf_{W_0}) \;\Longleftrightarrow\; \mathrm{C}(\Parf_{W_0})\Bigr),
\]
\textit{with the notations} $\mathrm{A}(\cdot)$, $\mathrm{B}(\cdot)$, $\mathrm{C}(\cdot)$ \textit{introduced in} (4.4.1). \textit{Taking} $W_0 = W$, \textit{we find the conditions on} $W$
\[
\mathrm{A},\ \mathrm{B},\ \mathrm{C},\ \mathrm{D}
\tag*{(4.4.13.2)}
\]
\textit{with} $\mathrm{A} = \mathrm{A}_W$, $\mathrm{B} = \mathrm{B}_W$, $\mathrm{C} = \mathrm{C}_W$, $\mathrm{D} = \mathrm{D}_W$, \textit{of which the first three are none other than} Loc (7 a), Loc (7 b), Loc (7 c) (cf. 4.3, \textit{page} 38).

\circled{1} \textit{For these conditions we have the following implications.}

\begin{enumerate}[leftmargin=2.5em]
\item[a)] \textit{We have the implications}
\[
\mathrm{A}_{W_0} \;\Longrightarrow\; \mathrm{B}_{W_0} \;\Longrightarrow\; \mathrm{C}_{W_0} \;\Longrightarrow\; \mathrm{D}_{W_0},
\tag*{(4.4.13.3)}
\]
(\textit{whence in particular}

\msp{page 59}
\[
\mathrm{A}_W \;\Longrightarrow\; \mathrm{B}_W \;\Longrightarrow\; \mathrm{C}_W \;\Longrightarrow\; \mathrm{D}_W\ ).
\tag*{(4.4.13.4)}
\]
\item[b)] \textit{If} $W'_0 \subset W_0$ \textit{is a basic localizer contained in} $W_0$, \textit{we have} (\textit{tautologically})
\[
\mathrm{A}_{W_0} \;\Longrightarrow\; \mathrm{A}_{W'_0}, \qquad\quad \mathrm{B}_{W_0} \;\Longrightarrow\; \mathrm{B}_{W'_0}, \qquad\quad \mathrm{C}_{W_0} \;\Longrightarrow\; \mathrm{C}_{W'_0}.
\]
\textit{We thus find, by} (a) \textit{and} (b) (\textit{taking there for the pair} $(W, W_0)$ \textit{as} $(W_0, W'_0)$),
\end{enumerate}


\endgroup
% END UNIT der-ch16/der-ch16_p031-035.tex
% BEGIN UNIT der-ch16/der-ch16_p036-040.tex
\begingroup\makeatletter%
\newcommand{\printedpage}[1]{\marginpar{\footnotesize p.~#1}}%
\newcommand{\msp}[1]{\marginpar{\footnotesize [#1]}}%
\newcommand{\cercle}[1]{\tikz[baseline=(c.base)]\node[draw,circle,inner sep=1pt](c){#1};}%
\newcommand{\fourche}{\begin{tikzcd}[ampersand replacement=\&, row sep=0.5em, column sep=1em, baseline=(\tikzcdmatrixname-2-1.base)]%
 \& X_1 \\%
 X_0 \arrow[ur] \arrow[dr] \& \\%
 \& X_2%
\end{tikzcd}}%
\newcommand{\carre}{\begin{tikzcd}[ampersand replacement=\&, row sep=0.5em, column sep=1em, baseline=(\tikzcdmatrixname-2-1.base)]%
 \& X_1 \arrow[dr] \& \\%
 X_0 \arrow[ur] \arrow[dr] \& \& X \\%
 \& X_2 \arrow[ur] \&%
\end{tikzcd}}%
\makeatletter%
\renewcommand{\@makefnmark}{\hbox{\textsuperscript{\normalfont(\@thefnmark)}}}%
\renewcommand{\@makefntext}[1]{\parindent 1em\noindent\hbox to 1.8em{\hss\textsuperscript{\normalfont\@thefnmark}}\,#1}%
\makeatother%
\makeatother\let\printedpage\derpp\let\msp\dermsp%
\let\mspin\DERmspin%


\begin{enumerate}[leftmargin=2.5em]
\item[] \printedpage{34}%
% head: A. Grothendieck : Dérivateurs, Ch. XVI, Édition des Universités de Montpellier II et Paris VII
\textit{the rectangle of implications}
\[
\begin{tikzcd}[row sep=large, column sep=large, execute at end picture={
\draw[decorate,decoration={brace,amplitude=6pt}] ([xshift=-14pt]A.west |- C.south) -- ([xshift=-14pt]A.north west) node[midway,left=8pt]{$\mathcal{F}_W,\mathrm{Fibf}_W$};
\draw[decorate,decoration={brace,amplitude=6pt}] ([xshift=14pt]AW.north east) -- ([xshift=14pt]AW.east |- CW.south) node[midway,right=8pt]{$\mathcal{F}_{W_0},\mathrm{Fibf}_{W_0}$};
}]
|[alias=A]| \mathrm{A} \arrow[r,Rightarrow] \arrow[d,Rightarrow] & |[alias=AW]| \mathrm{A}_{W_0} \arrow[d,Rightarrow] \\
\mathrm{B} \arrow[r,Rightarrow] \arrow[d,Rightarrow] & \mathrm{B}_{W_0} \arrow[d,Rightarrow] \\
|[alias=C]| \mathrm{C} \arrow[r,Rightarrow] \arrow[d,Rightarrow] & |[alias=CW]| \mathrm{C}_{W_0} \arrow[d,Rightarrow] \\
\underbrace{\mathrm{D}}_{\mathrm{Parf}_W} \arrow[r,Rightarrow] & \underbrace{\mathrm{D}_{W_0}}_{\mathrm{Parf}_{W_0}}\,,
\end{tikzcd}
\tag*{(4.4.13.5)}
\]
\textit{where we have recalled, next to the braces, what the set, or the sets $\mathcal{F}\subset\mathrm{Fl}(\mathrm{Cat})$ are which define the conditions considered (as one of the conditions $\mathrm{A}(\mathcal{F})$, $\mathrm{B}(\mathcal{F})$, $\mathrm{C}(\mathcal{F})$ or $\mathrm{D}(\mathcal{F})=\mathrm{A}(\mathcal{F})$ in the case where the three conditions are equivalent). The weakest of the eight conditions considered is $\mathrm{D}_{W_0}$, the strongest} A. (\textit{And for $W_0$ variable, we find the weakest condition of all by taking $D_{W_\omega}$, where $W_\omega$ is the smallest of the basic localizers in} Cat.)

\item[c)] \textit{We have moreover, for the weakest condition $\mathrm{D}_{W_0}$, the implication}
\[
\mathrm{D}_{W_0} \implies \mathrm{F},
\]
\msp{page 60}%
\textit{where} $\mathrm{F}=\text{Loc (7 F)}$ \textit{is the condition that} Cat \textit{has enough $W$-fibrations,} i.e. \textit{that it is a category of fibrations for} $(W,\mathrm{Fib}_W)$.
\end{enumerate}

\noindent\cercle{2}\enspace{\itshape If $W$ and $W_0$ satisfy \textup{Loc (8 bis)} (axiom of countable limits), then \textup{$\mathrm{D}_{W_0}\Longrightarrow\mathrm{A}$, i.e.} the 8 conditions considered are equivalent. Thus if $W$ satisfies \textup{Loc (8 bis)}, considering the smallest basic localizer $W_L$ satisfying \textup{Loc (8 bis)} (so that $W\supset W_L$), we see that the eight conditions
\[
\mathrm{A},\ \mathrm{B},\ \mathrm{C},\ \mathrm{D},\ \mathrm{A}_{W_L},\ \mathrm{B}_{W_L},\ \mathrm{C}_{W_L},\ \mathrm{D}_{W_L}
\]
are equivalent~\footnote[33]{N.B. The weakest form, $\mathrm{D}_{W_{\mathrm{L}}}$, of these conditions, is the relation \fbox{$\mathrm{Parf}_{W_{\mathrm{L}}}\cap W\subset W^{\mathrm{u}}$.}}. In particular, under the same (quite innocuous) hypothesis on $W$, the conditions \textup{Loc (7) $=$ Loc (7 a), Loc (7 b)} and \textup{Loc (7 c)} are equivalent, and they are also equivalent to the condition \textup{$\mathrm{D}=\mathrm{A}(\mathrm{Parf}_W)$, i.e.} to
\[
\mathrm{Parf}_W\subset\mathrm{Fib}_W \qquad (\textit{special case of } \text{Loc (7 a)}),
\]
and to its two seemingly weaker variants \textup{$\mathrm{B}(\mathrm{Parf}_W)$, $\mathrm{C}(\mathrm{Parf}_W)$} (special case of \textup{Loc (7 b) and of Loc (7 c)} respectively).}

% foot: Version 18 mars 2019

\noindent\cercle{3}\enspace{\itshape Suppose that $W_0$ satisfies \textup{Loc (8 bis)} (without any particular hypothesis on $W$). Then
\printedpage{35}%
% head: A. Grothendieck : Dérivateurs, Ch. XVI, Édition des Universités de Montpellier II et Paris VII
\msp{page 61}%
the four conditions $\mathrm{A}_{W_0}$, $\mathrm{B}_{W_0}$, $\mathrm{C}_{W_0}$, $\mathrm{D}_{W_0}$ are equivalent, and the diagram of implications of \textup{1°)} (p. \textup{59}) therefore reduces to the linear diagram
\[
\mathrm{A} \implies \mathrm{B} \implies \mathrm{C} \implies \mathrm{D} \implies \mathrm{D}_{W_0}\ (\Longleftrightarrow \mathrm{A}_{W_0}\ldots),
\]
whose first term is \textup{Loc (7 a) $=$ Loc (7)}, and the last is equivalent to the inclusion
\[
\mathcal{F}_{W_0}\subset\mathrm{Fib}_W,
\]
special case of \textup{$\mathrm{A}=$ Loc (7 a)}, or again to the inclusion (apparently stronger)
\[
\mathrm{Fibf}_{W_0}\subset\mathrm{Fib}_W,
\]
(special case of the equivalent form
\[
\mathrm{Fibf}_W\subset\mathrm{Fib}_W,
\]
of the axiom} Loc (7 a)).
\DERmark{ENV0968}


\noindent\cercle{4}\enspace\textit{Finally, we have the following more technical conditional implications.}
\begin{enumerate}[leftmargin=2.5em]
\item[a)] \textit{If $W_0$ satisfies Loc (7), then the conditions $\mathrm{A}_{W_0}$, $\mathrm{B}_{W_0}$, $\mathrm{C}_{W_0}$ are equivalent, so diagram (4.4.13.5) reduces to the diagram}
\[
\begin{tikzcd}[row sep=large, column sep=large]
\mathrm{A} \arrow[d,Rightarrow] & \\
\mathrm{B} \arrow[d,Rightarrow] & \\
\mathrm{C} \arrow[r,Rightarrow] \arrow[d,Rightarrow] & \mathrm{C}_{W_0} \arrow[d,Rightarrow] \\
\mathrm{D} \arrow[r,Rightarrow] & \mathrm{D}_{W_0}\ .
\end{tikzcd}
\]
\DERmark{ENV0969}


\item[b)] \textit{If $W$ is strongly saturated} (\textit{axiom} Loc (1 ter), \textit{which}
\msp{page 62}%
\textit{is satisfied in the great majority of cases, and is even imperative for constructing a prederivator satisfying the localization axiom} Der 2) \textit{and satisfies} Loc (3 bis) (\textit{strong localization axiom,} cf. §1), \textit{then we have} $\mathrm{B}\Longrightarrow\mathrm{A}$, \textit{hence}
\[
\mathrm{A} \iff \mathrm{B},
\]
i.e. Loc (7 a) $=$ Loc (7) \textit{is equivalent to} Loc (7 b).
\end{enumerate}

\noindent\textsc{Proof}. Everything stated has been seen in the preceding pages, with the sole exception of 4°, b. But this implication $\mathrm{B}\Longrightarrow\mathrm{A}$ for $W$ strongly saturated is proved in XII (prop. 1, p. 54). I recall that the proof makes essential use of the % [sic] prouvé -> prouvée, accord avec « implication »
% foot: Version 18 mars 2019
\printedpage{36}%
% head: A. Grothendieck : Dérivateurs, Ch. XVI, Édition des Universités de Montpellier II et Paris VII
formalism of the prederivator $\mathrm{HOT}_W$ (constructed and studied mainly in VI), and the theory of the functors $f_!$, and in particular the result (valid if $W$ satisfies Loc (3 bis, 7 b)) that for an arrow $f:S'\longrightarrow S$ in $W$, the corresponding arrow
\[
f^*_{\mathrm{lc}}:\mathrm{HOT}^{\mathrm{lc}}_W(S)\longrightarrow\mathrm{HOT}^{\mathrm{lc}}_W(S')
\]
is an equivalence of categories. So the implication $\mathrm{B}\Longrightarrow\mathrm{A}$, too, is a result which is not at all obvious.

\noindent\msp{page 63}%
4.4.14. I summarize the implications obtained so far, including the conditional implications, in the following summary diagram:
\[
\begin{tikzpicture}[baseline=(C.base),font=\small]
\node[draw,circle,inner sep=1pt] (A) at (0,0) {$\mathrm A$};
\node (AW) at (2,0) {$\mathrm A_{W_0}$};
\node (B) at (0,-1.2) {$\mathrm B$};
\node (BW) at (2,-1.2) {$\mathrm B_{W_0}$};
\node (C) at (0,-2.4) {$\mathrm C$};
\node (CW) at (2,-2.4) {$\mathrm C_{W_0}$};
\node (D) at (0,-3.6) {$\mathrm D$};
\node (DW) at (2,-3.6) {$\mathrm D_{W_0}$};
\node (F) at (7.8,-3.6) {$\mathrm F$};
\node[draw,circle,inner sep=1pt] (Ab) at (2,-4.8) {$\mathrm A$};
\node[anchor=west] at (2.25,-4.8) {$\,,$};
\foreach \p/\q in {A/AW,B/BW,C/CW,D/DW,A/B,B/C,C/D,AW/BW,BW/CW,CW/DW,DW/F}
  \draw[double,double distance=1pt,-{Implies}] (\p) -- (\q);
\draw[dotted,->] (B.west) to[out=145,in=215] node[left,font=\scriptsize] {$W$ (1 ter, 3 bis)} (A.west);
\draw[dotted,->] (CW.east) to[out=70,in=-70] node[right,font=\scriptsize] {$W_0$ (7)} (AW.east);
\draw[dotted,->] (DW.east) to[out=60,in=-40,looseness=1.45] node[right,font=\scriptsize] {$W_0$ (8 bis)} (AW.east);
\draw[dotted,->] (D.south east) -- node[below left,font=\scriptsize] {$W$ (8 bis)} (Ab.north west);
\draw[dotted,->] (DW) -- node[right,font=\scriptsize] {$W$ (8 bis) $+$ $W_0$ (8 bis)} (Ab);
\node[anchor=south west,align=left,font=\scriptsize] at (3.4,-3.45) {N.B. we shall see from the example of $\mathbf{W}_0$\\
that this implication is \textit{strict,}\\
even under the hypothesis\\
Loc (1 ter, 2, 3 bis, 4, 5, 6 a, 8, 9).};
\end{tikzpicture}
\]
\DERmark{ENV0970}

where I have circled the strongest and most important condition of all, $\mathrm{A}=$ Loc (7), which appears at two opposite ends of the diagram.

I would like to make one more comment on the ingredients of the results proved here, and in particular of the equivalence of the eight conditions A to $\mathrm{D}_{W_0}$, when $W$ and $W_0$ satisfy Loc (8 bis):

\msp{page 64}%
\begin{enumerate}[leftmargin=2.5em]
\item[1)] The theory of \textit{categories of paths} \underline{Ch}$(X)$ (VII), relying on that of \textit{smooth, proper, perfect morphisms.}
\item[2)] The lemma of K. S. \textsc{Brown} (3.2.1, page 20).
\item[3)] For the implication $\mathrm{B}\implies\mathrm{A}$, \textit{i.e.} Loc (7 b) $\implies$ Loc (7) (under the hypothesis $W$ (1 ter, 3 bis)), I use the theory of the prederivator $\mathrm{Hot}_W$ (VI), and more particularly that of the functors $f_!$ (defined under the sole hypothesis $W$ (3 bis), strengthening the automatic condition $W$ (3) (forming part of the definition of a basic localizer)).
\end{enumerate}
% foot: Version 18 mars 2019

\leavevmode\printedpage{37}%
% head: A. Grothendieck : Dérivateurs, Ch. XVI, Édition des Universités de Montpellier II et Paris VII
It is clear that I would not have managed to see clearly, without the fundamental work of K. S. \textsc{Brown}. Yet there still remain open questions. Thus, if we do not assume that $W$ satisfies Loc (8 bis), but if need be Loc (3 ter) (and I have not yet had occasion to work with a localizer which does not satisfy it \ldots), is it nevertheless true that we still have the implication
\[
\mathrm{D} \qquad \Longrightarrow \qquad \mathrm{A},
\]
so that the conditions A, B, C, D are equivalent (and in particular Loc (7 c) already implies Loc (7) in all its force)? It is clear in all this that
\msp{page 65}%
I have not applied myself enough to developing examples and counterexamples. Not to mention that it seems to me that all the basic localizers on Cat which we encounter ``in practice'' already satisfy Loc (8) and a fortiori Loc (8 bis), so for them we do indeed have the expected implication. The same remark applies to the hypothetical implication
\[
\mathrm{D}_{W_0} \overset{?}{\implies} \mathrm{A},
\]
under hypothesis Loc (8 bis) on $W_0$ and on $W$. An obvious candidate for a counterexample already to the first question would be $W_\omega=$ smallest basic localizer on Cat, of which it is almost obvious that it does \textit{not} satisfy Loc (8) nor even Loc (8 bis). But there is no reason why it should satisfy F either, so that (if it does not satisfy it, hence not D either) the implication $\mathrm{D}\Longrightarrow\mathrm{A}$ is indeed valid. This brings us to the question, a more substantial one, whether we might not have the implication % [sic] un qu'il est -> dont il est (?), construction défectueuse
\DERmark{ENV0971}

\msp{page 66}%
\[
\mathrm{F} \overset{?}{\implies} \mathrm{A}\ \text{\footnotemark[34]},
\]
\footnotetext[34]{This is false, \textit{cf.} below, the implication Loc (7 $\mathrm{D}_{W_\omega}$) $\Longrightarrow$ Loc (7 F) is \textit{strict.}}%
at least under suitable fairly innocuous conditions on $W$, of the sort Loc (1 ter, 3 bis, 8 bis).

This is probably the most substantial and the most interesting question which is raised by the diagram on page 63, and which no longer calls for the introduction of an auxiliary basic localizer $W_0$. In view of Brown's lemma, it does not seem to me absolutely ruled out that such a miracle should indeed take place!

4.4.15 (Digression on the axiom Loc (6) of cofibrations.) This draws my attention to the in some sense dual question: If Cat has enough $W$-\textit{co}fibrations, \textit{i.e.} if equipped with $(W,\mathrm{Cof}_W)$ it is a category with cofibrations, does this already imply the axiom Loc (6 a) (which had been isolated, essentially, to ensure that this is indeed so). Proposition 3.6. (p. 25) tells us that it amounts to the same thing to prove that, if $W$ satisfies Cof, then the open M. V. squares are
\msp{page 67}%
$W$-cocartesian. But it seems to me now that we can prove this, under the sole additional condition Loc (3 bis)~\footnote[35]{and perhaps, in addition, Loc (1 ter), \textit{i.e.} $W$ strongly saturated.}, which ensures the existence of the functors $f_!$, and % [sic] cocartésien -> cocartésiens, accord avec « les carrés » ; [sic] (note 35) saturée -> saturé, accord avec « W » (localiseur)
% foot: Version 18 mars 2019
\printedpage{38}%
% head: A. Grothendieck : Dérivateurs, Ch. XVI, Édition des Universités de Montpellier II et Paris VII
in particular the construction of Holim of a diagram $\fourche$ \quad as $\int\fourche$ \quad .
So to say that an open $\mathcal{M}$ diagram $\carre$ is $W$-cocartesian, must amount to saying that the canonical arrow
\[
\int\fourche \qquad \longrightarrow \qquad X
\]
is in $W$; at least if $W$ is assumed strongly saturated. Now this has indeed been seen in XII 3.4 (theorem 4 and its \textit{corollary} 3, p. 149). In loc. cit. I state the conclusions under the hypothesis which I called W (7 bis) (the present Loc (6)), but I have just checked that for its part \textit{in the forward direction,} hence also for its corollaries 1 to 3, we do not need
\msp{page 68}%
this hypothesis (as I indicate in marginal annotations).

4.16. The result I have just proved (without carefully checking the proof) seems to me as delicate, as it is satisfying. This makes the hesitant ``conjecture'' of 4.4.14, namely that
\[
\mathrm{F} \overset{?}{\Longrightarrow} \mathrm{A} \qquad \text{hence} \qquad \mathrm{F}\Longleftrightarrow\mathrm{A}
\]
assuming that $W$ satisfies Loc (1 ter, 3 bis, 8 bis), decidedly more inviting. For one begins to get used to the most astounding miracles \ldots

Yet I have just checked that this conjecture is false, using the example of
\[
\mathbf{W}_0=\{f\in\mathrm{Fl}(\mathrm{Cat})\mid\pi_0(f)\ \mathrm{iso}\}\,,
\]
which satisfies F, but not A, nor even $\mathrm{D}_{W_\omega}$. Thus the implication
\[
\mathrm{D}_{W_\omega} \implies \mathrm{F}
\]
and a fortiori the implications
\[
\left\{
\begin{aligned}
&\text{Loc (7 a)} \implies \mathrm{F},\ \text{Loc (7 b)} \implies \mathrm{F},\ \text{Loc (7 c)} \implies \mathrm{F}\\
&\mathrm{D}_{W_0} \implies \mathrm{F}
\end{aligned}
\right.
\]
are \textit{strict,} even when restricting ourselves to basic localizers $W$ satisfying (like $\mathbf{W}_0$)
\[
\text{Loc (1 ter, 2, 3 bis, 4, 5, 6 a, 8, 9)}\ \text{\footnotemark[36]}.
\]
\footnotetext[36]{\textbf{N.B. $\mathbf{W}_0$} does \textit{not} satisfy Loc (6 b), nor Loc (7 $\mathrm{D}_{W_\omega}$) (\textit{cf.} below), on the other hand Loc (7 F).}
% foot: Version 18 mars 2019


\endgroup
% END UNIT der-ch16/der-ch16_p036-040.tex
% BEGIN UNIT der-ch16/der-ch16_p041-045.tex
\begingroup\makeatletter%
\newcommand{\printedpage}[1]{\marginpar{\footnotesize p.~#1}}%
\newcommand{\msp}[1]{\marginpar{\footnotesize [#1]}}%
\newcommand{\uFl}{\underline{\mathrm{Fl}}}%
\newcommand{\uCat}{\underline{\mathrm{Cat}}}%
\newcommand{\Parf}{\mathrm{Parf}}%
\makeatother\let\printedpage\derpp\let\msp\dermsp%
\let\mspin\DERmspin%


\printedpage{39}
% head: A. Grothendieck : Dérivateurs, Ch. XVI, Édition des Universités de Montpellier II et Paris VII

\msp{page 69}
That $\mathbf{W}_0$ satisfies this set of axioms is checked by AQT \texttt{[âne qui trotte: at a donkey's trot, i.e. by routine verification]}. (\textbf{N.B.} That it satisfies Loc (1 ter, 2, 3 bis) is already a consequence of the fact that $\mathbf{W}_0$ is of the form $W_{\mathbf{D}}$, where $\mathbf{D}$ ($=\mathbf{D}_{\mathbf{W}_0}$ in this case) is a prederivator satisfying Der (3, 4, 5), which is indeed the case for $\mathbf{D}_{\mathbf{W}_0}$, which is a genuine \textit{derivator.}) We must check that every $f$ in Cat factors as $f = pi$, with $i \in \mathbf{W}_0$, $p$ a $\mathbf{W}_0$-fibration. Now if $f : X \longrightarrow Y$, $I = \pi_0(X)$, we take $\bar{X} = \coprod_{\alpha\in I} Y_{f(\alpha)}$, whence a factorization $f = pi$, where $\pi_0(i)$ is an isomorphism, \textit{i.e.} $i \in \mathbf{W}_0$, and where $p : \bar{X} \longrightarrow Y$ is, over the connected component $Y_j$ of $Y$, given by the product $\bar{X}|Y_j \simeq Y_j \times \underbrace{\pi_0(f)^{-1}(j)}_{\text{discrete category}}$ . The morphism $p$ is indeed a $\mathbf{W}_0$-fibration, for thanks to Loc (4), it suffices to check this on each $Y_j$, and there it is deduced from an arrow $I_j \longrightarrow e$ (a $\mathbf{W}_0$-fibration) by base change $Y_j \longrightarrow e$.

On the other hand, that $\mathbf{W}_0$ does \textit{not} even satisfy $\mathrm{D}_{W_\omega}$ has already been seen (I, p. 87),
\msp{page 70}
by realizing the map of the $n$-covering $f : S^1 \longrightarrow S^1$ in Cat by a $W_\omega$-perfect morphism (in fact, by a locally trivial Cat-fibration), a very easy thing\footnote[37]{Another way of seeing it: we \textit{know} (XII) that $\mathrm{W}$ (5, 7 $\mathrm{D}_W$) $\Longrightarrow W \subset \mathbf{W}_\infty$, now we have $\mathbf{W}_\infty \subset \mathbf{W}_0$, and the inclusion is \textit{strict\,!}}.

I think that by now the whole set of different variants of the axiom of $W$-fibrations is more or less understood, except that it is still not clear, without the axiom Loc (8 bis), whether the implications in the diagram
\[
\begin{tikzcd}[row sep=large, column sep=large]
\mathrm{A} \arrow[r, Rightarrow] \arrow[d, Rightarrow] & \mathrm{A}_{W_0} \arrow[d, Rightarrow] \\
\mathrm{B} \arrow[r, Rightarrow] \arrow[d, Rightarrow] & \mathrm{B}_{W_0} \arrow[d, Rightarrow] \\
\mathrm{C} \arrow[r, Rightarrow] \arrow[d, Rightarrow] & \mathrm{C}_{W_0} \arrow[d, Rightarrow] \\
\mathrm{D} \arrow[r, Rightarrow] & \mathrm{D}_{W_0}
\end{tikzcd}
\]
are strict or not. But it seems to me that all the interesting basic localizers encountered so far satisfy Loc (8), and even Loc (1 ter, 2, 3 bis, 4, 5, 6 a, 8, 9) \marginpar{?}(perhaps even Loc (7 F) - I would have to test on the $\mathbf{W}_n$), so the unresolved perplexity seems to me of no great consequence. % [sic] résolu -> résolue, accord avec « la perplexité » ; « ? » imprimé dans la marge gauche à la hauteur de cette ligne

\msp{page 71}

\printedpage{40}
% head: A. Grothendieck : Dérivateurs, Ch. XVI, Édition des Universités de Montpellier II et Paris VII

\section*{5\quad Limits axiom Loc (8), and the accessibility axiom Loc (9)}

It is the axiom:

\textbf{Loc (8)}\quad $W$ is stable under filtering $\varinjlim$.

A useful weakened variant is the following:

\textbf{Loc (8 bis)}\quad $W$ is stable under countable $\varinjlim$.

Explicitly, if we have two filtering direct systems in Cat
\[
\begin{aligned}
\underline{X} &= (X_i)_{i\in I}\\
\underline{Y} &= (Y_i)_{i\in I},
\end{aligned}
\]
and an arrow
\[
\underline{f} = (f_i)_{i\in I} : \underline{X} \longrightarrow \underline{Y},
\]
given by $f_i : X_i \longrightarrow Y_i$, putting $X = \varinjlim X_i$, $Y = \varinjlim Y_i$, $f = \varinjlim f_i : X \longrightarrow Y$, Loc (8) tells us that if the $f_i \in W$, then $f \in W$. Here $I$ is
\msp{page 72}
an arbitrary (small) filtering category, and after replacing it by a cofinal category if necessary, we may suppose $I$ an ordered set. The axiom Loc (8 bis) corresponds to the case where $I = \mathbf{N}$.

Note that thanks to the construction of the category of paths, we have a factorization, functorial in $f$, of every arrow $f : X \longrightarrow Y$ of Cat as $f = pi$:
\[
X \xrightarrow{\quad i = i_f\quad} \bar{X}(f) = \bar{X} \xrightarrow{\quad p = p_f\quad} Y,
\]
with $i \in W_\omega$, $p \in \Parf_{W_\omega}$, a fortiori $i \in W$, $p \in \Parf_W$. This functor $f \longmapsto (\bar{X}(f), p_f)$ commutes with direct limits. Thus, if we apply this to each of the $f_\alpha : X_\alpha \longrightarrow Y_\alpha$, factoring it as
\[
X_\alpha \xrightarrow{i_\alpha} \bar{X}_\alpha \xrightarrow{p_\alpha} Y_\alpha,
\]
we find a direct system of diagrams of type $\boldsymbol{\Delta}^2$ (factoring the $f_i$), and on passing to the $\varinjlim$ we find
\[
X \xrightarrow{i} \bar{X} \xrightarrow{p} Y,
\]
which is none other than the canonical factorization
\msp{page 73}
of $f = \varinjlim f_i$ as $pi$, with $i \in W_\omega$, $p \in \Parf_{W_\omega}$ (N.B. Whereas neither $W_\omega$, nor $\Parf_{W_\omega}$ are stable under filtering $\varinjlim$, for sure\,!). We therefore have $i \in W$, and the condition $f \in W$, \textit{i.e.} $pi \in W$, is therefore equivalent to $p \in W$. This reduces us, for checking the axiom Loc (8) (resp. Loc (8 bis)), to the case where the $f_i$ are in $\Parf_{W_\omega}$.

Suppose now that $W$ satisfies $\mathrm{D}_{W_\omega}$ (the weakest of all the variants considered here of the axiom of $W$-fibrations Loc (7), with the sole exception of F). Then we have
\[
p \in W \qquad\quad \Longleftrightarrow \qquad\quad \text{the fibers of $p$ are $W$-aspheric,}
\]
\printedpage{41}
% head: A. Grothendieck : Dérivateurs, Ch. XVI, Édition des Universités de Montpellier II et Paris VII
the implication $\Longrightarrow$ resulting from the hypothesis $\mathrm{D}_{W_\omega}$ $\Parf_{W_\omega} \cap W \subset W^{\mathrm{u}}$, and $\Longleftarrow$ resulting from the sorites of $W$-smooth morphisms \textit{etc.}, via Loc (3). Likewise, the hypothesis that the $f_i \in W$ is equivalent to the one that the fibers of the $p_i$ are $W$-aspheric. On the other hand, it is immediate that
\msp{page 74}
every fiber of $p$ is isomorphic to the $\varinjlim\limits_{I}$ of the corresponding fibers of the $p_\alpha$. In this way, we find:

\textbf{Proposition 5.1.} \textit{Suppose that $W$ satisfies the axiom of $W$-fibrations} Loc (7 $\mathrm{D}_{W_\omega}$) (\textit{which is therefore the case if it satisfies one of the stronger axioms} Loc (7 a \textit{or} b \textit{or} c)). \textit{For $W$ to satisfy} Loc (8) (resp. \textit{to satisfy} Loc (8 bis)), \textit{it is necessary and sufficient that the set of objects in} Cat \textit{which are $W$-aspheric be stable under filtering $\varinjlim$} (resp. \textit{under $\varinjlim$ of countable sequences}).

5.2. Note that the category Cat is a paratopos, a fortiori it is \textit{accessible} (\textit{cf.} XVIII, §8), and the same therefore holds for $\uFl(\mathrm{Cat})$. Thus we have in $\mathrm{Fl}(\mathrm{Cat}) = \mathrm{Ob}\,\uFl(\mathrm{Cat})$ the notion of an \textit{accessible subset.} We shall therefore say that a subset $W$ of $\mathrm{Fl}\,\mathrm{Cat}$ (more generally of a category $\mathrm{Fl}(\mathcal{M})$, where $\mathcal{M}$ is an accessible category) is \textit{accessible,} if it is so with respect to the ambient category
\msp{page 75}
$\uFl\,\mathrm{Cat}$. This being so, the axiom Loc (9) is the following.

\textbf{Loc (9)}\quad $W$ is an \textit{accessible} subset of $\mathrm{Fl}(\mathrm{Cat}) = \mathrm{Ob}\,\uFl(\mathrm{Cat})$.

5.3. Let us recall that this somewhat technical notion simplifies considerably, if we suppose that $W$ satisfies the limits axiom Loc (8). Then $W$ is accessible if and only if there exists a \textit{small} subset $W_0$ of Fl(Cat), such that $W$ is the ``L$_1$-envelope'' of $W_0$, \textit{i.e.} the smallest subset of Fl(Cat) containing $W_0$ and L$_1$-stable (\textit{i.e.} stable under filtering $\varinjlim$). If this is so, it is clear that we can take $W_0$ of the form
\[
W_\pi = W \cap \mathrm{Fl}_\pi(\mathrm{Cat}),
\]
where $\pi$ is a suitable infinite cardinal, and where $\mathrm{Fl}_\pi(\mathrm{Cat}) \overset{\mathrm{def}}{=} \mathrm{Fl}(\uCat_\pi)$, with $\uCat_\pi$ denoting the strictly full subcategory of Cat consisting of the objects $X$ such that $\mathrm{card}\,X \leq \pi$ (where by definition, $\mathrm{card}\,X = \mathrm{card}\,\mathrm{Fl}\,X$).
\DERmark{ENV0972}


\msp{page 76}
Note that the filtration we have just described on Cat (indexed by the small infinite cardinals) is none other than its canonical filtration, \textit{i.e.} $\uCat_\pi$ is the full subcategory of Cat consisting of the $\pi$-accessible objects of Cat, and likewise in $\uFl(\mathrm{Cat})$. We thus find, thanks to the developments in XVIII, §13, the following characterization, which I recall here for reference (restricting myself for simplicity to the case where $W$ satisfies Loc (8)):

\textbf{Proposition 5.3.1.}\footnote[38]{\textit{cf. loc. cit.} 13.36.} \textit{Let $W$ be a subset of} Fl(Cat) \textit{stable under filtering $\varinjlim$, and let, for every infinite cardinal $\pi$, $W_\pi = W \cap \mathrm{Fl}_\pi(\mathrm{Cat})$, and $\underline{W}$, $\underline{W}_\pi$ be the corresponding full subcategories of} $\uFl(\mathrm{Cat})$. \textit{We say that $W$ is} $\pi$-accessible \textit{if $W$ is the L-envelope of $W_\pi$ (thus $W$ is accessible if and only if there exists an infinite cardinal $\pi$ such that $W$ is $\pi$-accessible). This being so, the following conditions are equivalent:}

\printedpage{42}
% head: A. Grothendieck : Dérivateurs, Ch. XVI, Édition des Universités de Montpellier II et Paris VII

\begin{enumerate}[label={}]
\item[a)] \textit{$W$ is $\pi$-accessible.}
\item[b)] \textit{For every direct system $(f_i)_{i\in I}$,}
\msp{page 77}
\textit{$L_1$-$\pi$-adapted in} $\uFl\,\mathrm{Cat}$ (cf. \textit{below}), \textit{with direct limit $f$, we have $f \in W$ if and only if the full subcategory $I_W$ of the $i \in \mathit{Ob}I$ such that $f_i \in W$ is cofinal in $I$.}
\item[b$'$)] (\textit{When we have chosen, for every $f \in$} Fl(Cat)\textit{, an $L_1$-$\pi$-adapted representation of $f$ as a $\varinjlim$}): \textit{If $f \in W$, then the full subcategory $I_W$ defined in} b) \textit{is cofinal in $I$.}
\end{enumerate}

Let us recall here the notion of an L$_1$-$\pi$-adapted direct system. The conditions are:
\begin{enumerate}[label={}]
\item[1°)] $I$ is large relative to $\pi$, and stable under filtering $\varinjlim$ of cardinal $\leq \pi$.
\item[2°)] The functor $i \longmapsto x_i : I \longrightarrow \uFl(\mathrm{Cat})$ commutes with the $\varinjlim$ in question.
\item[3°)] The $x_i$ are in $\uFl(\mathrm{Cat})_\pi \overset{\mathrm{def}}{=} \uFl_\pi(\mathrm{Cat})$.
\end{enumerate}
\DERmark{ENV0973}

Note that in the case of a ``parabelian paratopos'' such as Cat or $\uFl\,\mathrm{Cat}$, there is a canonical way, $\pi$ being given, to represent an object $f$ as the $\varinjlim$ of an L$_1$-$\pi$-adapted direct system. Here,
\[
f : X \longrightarrow Y
\]
being given, we consider the set $I$ of
\msp{page 78}
pairs $(X_i \subset X, Y_i \subset Y)$ of subcategories of $X$ and $Y$, such that we have $X_i, Y_i \in \mathrm{Ob}\,\mathrm{Cat}_\pi$ (\textit{i.e.} $\mathrm{card}\,(X_i), \mathrm{card}\,(Y_i) \leq \pi$) and that $f(X_i) \subset Y_i$. It is immediate that $I$ is filtering, for the obvious order relation
\[
(X_i, Y_i) \subset (X_j, Y_j) \qquad\quad \Longleftrightarrow \qquad\quad X_i \subset X_j \ \text{ and } \ Y_i \subset Y_j, % [sic] X_i \subset Y_i et X_j \subset Y_j -> X_i \subset X_j et Y_i \subset Y_j, la relation d'ordre compare les couples (X_i,Y_i) et (X_j,Y_j)
\]
\DERmark{ENV0974}\DERmark{ENV0975}\DERmark{ENV0976}\DERmark{ENV0977}

and even large relative to $\pi$, that $I$ is stable under filtering $\varinjlim$ of cardinal $\leq \pi$, and that the functor
\[
i \longmapsto (f_i : X_i \longrightarrow Y_i) : I \longrightarrow \uFl\,\mathrm{Cat}
\]
commutes with the said $\varinjlim$ , finally the $f_i$ are in $\uFl(\mathrm{Cat})_\pi$ by construction. Thus we have an L$_1$-$\pi$-adapted direct system, and it is clear that its direct limit is canonically isomorphic to $f$. This having been seen:

\textbf{Corollary 5.3.2.} \textit{The condition that $W$ be $\pi$-accessible therefore means that for every $f : X \longrightarrow Y \in W$, and for every subcategory $X_i$ of $X$ and $Y_i$ of $Y$, of cardinal $\leq \pi$ (with $f(X_i) \subset Y_i$, but in fact this is unnecessary, after enlarging $Y_i$ if necessary), there exist subcategories} % [sic] (avec f(X_i) \subset Y_i), mais ... agrandir Y_i) -> (avec f(X_i) \subset Y_i, mais ... agrandir Y_i), parenthèse fermante surnuméraire
\DERmark{ENV0978}

\msp{page 79}
\textit{$X_j \supset X_i$ and $Y_j \supset Y_i$ of cardinal $\leq \pi$, such that $f$ induces $f_j : X_j \longrightarrow Y_j$ and that we have $f_j \in W$.}

The notion of accessibility of $W$ seems to me as well understood as that of limits, I no longer see any mystery in it. On the other hand, it seems to me that all the basic localizers used to date satisfy it (just as they do the axiom Loc (8)). Here, I shall confine myself to establishing:

\printedpage{43}
% head: A. Grothendieck : Dérivateurs, Ch. XVI, Édition des Universités de Montpellier II et Paris VII

\textbf{Theorem 5.4.}\footnote[39]{Well deserves the name of \textit{theorem,} especially if we include in it the whole list of Loc axioms mentioned opposite\,! This theorem is not at all obvious, starting from the cohomological definition of $\mathbf{W}_\infty$ - it is a result (``well known'') which seems deep to me.} \textit{The basic localizer $\mathbf{W}_\infty$ (``usual'' weak equivalences) satisfies the axiom} Loc (9) (\textit{in fact, it satisfies} all \textit{the axioms seen so far, namely} Loc (1 ter, 2, 3 bis, 4, 5, 6, 7, 8, 9)). \textit{More precisely, $\mathbf{W}_\infty$ is $\aleph_0$-accessible (where $\aleph_0$, or $\pi_0$, is the first infinite cardinal).} % [sic] les axiomes vu -> les axiomes vus, accord du participe

I do not return here to the verification of the other axioms, and confine myself to proving Loc (9), accessibility, using what may be useful to me. I
\msp{page 80}
confine myself to indicating the proof (rather technical) in its broad outlines. For the moment, $W$ is again an arbitrary basic localizer, except that it satisfies Loc (8).

\textbf{Lemma 5.5.}\footnote[40]{We suppose that $W$ satisfies Loc (8).}
\begin{enumerate}[label={}]
\item[a)] \textit{For the criterion} 5.3.1 b) \textit{of $\pi$-accessibility of $W$ to be satisfied, it is necessary and sufficient that it be so when the $f_i$ are $W_\omega$-perfect.}
\item[b)] \textit{Suppose that $\mathit{Parf}_{W_\omega} \cap W \subset W^u$,} i.e. \textit{that $W$ satisfies} Loc (7 $\mathrm{D}_{W_\omega}$) (\textit{the weakest of the variants of the axiom of $W$-fibrations, except only for} Loc (7 F)). \textit{Suppose moreover (to be on the safe side) that $W$ satisfies} \texttt{[unfinished sentence]}. \texttt{[?Then]} \textit{$W$ is $\pi$-accessible if and only if the set} $\mathrm{Asph}_W \subset \mathrm{Ob}\,\mathrm{Cat}$ \textit{of $W$-aspheric objects of} Cat \textit{is $\pi$-accessible. Or again, that the subset} $\mathrm{Asph}_W \subset \mathrm{Ob}\,\mathrm{Cat}_{\textit{0-connected}}$ \textit{be $\pi$-accessible, where} $\mathrm{Cat}_{\textit{0-connected}}$ \textit{denotes the full subcategory of} Cat \textit{consisting of} 0-\textit{connected objects.}
\end{enumerate}
Essentially, the proof is the same as that of 5.1, a little more technical. The essential tool is still the functorial factorization $f = pi$ of an arrow of Cat, with $i \in W_\omega$, $p \in \Parf_{W_\omega}$, and the property of
\msp{page 81}
commutation with filtering $\varinjlim$ of the said factorization.

Let us now return to the case of $W = \mathbf{W}_\infty$. We must prove that $\mathrm{Asph}_{\mathbf{W}_\infty} \subset \mathrm{Ob}\,\mathrm{Cat}_{\text{0-connected}}$ is $\pi_0$-accessible. Thus to prove the criterion corresponding to 5.3.1 b), placing ourselves now in $\mathrm{Cat}_{\text{0-connected}}$ instead of $\mathrm{Fl}\,\mathrm{Cat}$. We therefore have
\[
X = \varinjlim_{i\in I} X_i,
\]
an L$_1$-$\pi$-admissible representation of $X$ $\mathbf{W}_\infty$-aspheric, to prove that for every $i$ in $I$, $\exists j$ in $I$ and $i \longrightarrow j$, with $X_j$ $\mathbf{W}_\infty$-aspheric.

I admit here the AQT \texttt{[âne qui trotte: at a donkey's trot, i.e. by routine verification]} lemma:

\textbf{Lemma 5.6.} \textit{Let $\mathbf{W}_0 = \{f \in \mathrm{Fl}\,\mathrm{Cat} \mid \pi_0(f) \ \ \textit{isomorphisme}\}$. Then $\mathbf{W}_0$ is $\pi_0$-accessible (and $L_1$-stable). A fortiori, the set of} 0-\textit{connected objects of} Cat \textit{is $\pi_0$-accessible.}

This reduces us, after replacing $I$ by a cofinal category if necessary, to supposing that the $X_i$ are 0-connected. (But I forgot that this reduction has already been made\,! So one must view 5.6 as an ingredient


\endgroup
% END UNIT der-ch16/der-ch16_p041-045.tex
% BEGIN UNIT der-ch16/der-ch16_p046-050.tex
\begingroup\makeatletter%
\newcommand{\printedpage}[1]{\marginpar{\footnotesize p.~#1}}%
\newcommand{\msp}[1]{\marginpar{\footnotesize [#1]}}%
\makeatletter%
\renewcommand\@makefnmark{\hbox{\textsuperscript{(\@thefnmark)}}}%
\renewcommand\@makefntext[1]{\parindent 1em\noindent\hbox to 1.8em{\hss\textsuperscript{\@thefnmark}\,}#1}%
\makeatother%
\makeatother\let\printedpage\derpp\let\msp\dermsp%
\let\mspin\DERmspin%


\printedpage{44}
% head: A. Grothendieck : Dérivateurs, Ch. XVI, Édition des Universités de Montpellier II et Paris VII
% foot: Version 18 mars 2019

\msp{page 82}
of the proof of 5.5 b).)

Let $x$ be an object (point) of $X$, represented by a point $x_{i_0}$ of $X_{i_0}$; after restricting to the cofinal category $i_0\backslash I$, let $x_i$ be its image in $X_i$. Let us consider (via my standard construction of categories of paths) the $\Omega^1(X,x)$, $\Omega^1(X_i,x_i)$. We find
\DERmark{ENV0979}

\[ \Omega^1(X,x) \simeq \varinjlim_{I} \Omega^1(X_i,x_i), \]
and one checks that this is again an $\mathrm{L}_1$-$\pi_0$-adapted direct system. Iterating, we find for $n \in \mathbf{N}^*$
\[ \Omega^n(X,x) \simeq \varinjlim_{I} \Omega^n(X_i,x_i), \]
$\mathrm{L}_1$-$\pi_0$-adapted direct limit. Using 5.6, we conclude from this that for every $n$, the subset (let us suppose $I$ an ordered set, what the devil, that's enough!)
\[ I_n = \{ i \in I \mid \Omega^n(X_i,x_i) \ \text{ is 0-connected} \} \]
is cofinal in $I$, moreover it is stable under filtering $\varinjlim$ of cardinality $\leq \pi_0$. Thus % checked against PDF leaf 46: correct as printed; Theorem 5.4 explicitly uses pi_0 for the first infinite cardinal.
\[ I_\infty \overset{\mathrm{def}}{=} \bigcap I_n \]
\DERmark{ENV0980}

is also cofinal in $I$ (\textit{cf.} XVIII, 13.10.2, p. 302).

\msp{page 83}
Now $i \in I_\infty$ means that all the $\Omega^n(X_i,x_i)$ are 0-connected, which also means that $X_i$ is $\mathbf{W}_\infty$-aspheric (since it is already 0-connected) (\textit{cf.} XII, cor. p. 190). This completes the proof of theorem 5.4.

\section*{6\quad The condition Loc (5, 7 D$_{W_\omega}$) implies $W \subset \mathbf{W}_\infty$}

I am going to give here again, for the record, the proof of theorem 1 of XII:

\textbf{Theorem 6.1.} \textit{Let $W$ be a basic localizer satisfying} D$_{W_\omega}$ \textit{(where $W_\omega$ is the smallest of all basic localizers). Then we have $W = \mathrm{Fl}\,\mathrm{Cat}$,} or \textit{$W \subset \mathbf{W}_\infty$ (according as $W$ does not satisfy the connectedness axiom} Loc (5), \textit{or does satisfy it.})

I will give the proof assuming that $W$ satisfies Loc (5). For the proof of the fact that in the contrary case, we have $W = \mathrm{Fl}(\mathrm{Cat})$, I refer to XII, th. 2, p. 20. (See also the very close result, more general and weaker, recalled in prop. 2.7, page 12.)

\textbf{Lemma 6.2.} \textit{Let $W$, $W'$ be two basic localizers, $W$ satisfying} Loc (7 D$_{W_\omega}$). \textit{For $W \subset W'$ to hold, it is necessary and sufficient that} $\mathrm{Asph}_W \subset \mathrm{Asph}_{W'}$.

\msp{page 84}
It suffices of course to prove ``it suffices'', so let us suppose that every $W$-aspheric object is $W'$-aspheric, and let us prove that $W \subset W'$. So let $f : X \longrightarrow Y$ be in $W$, let us prove $f \in W'$. For this, I use the canonical factorization $f = pi$,
\[ X \xrightarrow{\ i \in W_\omega\ } \bar{X} \xrightarrow{\ p \in \mathrm{Parf}_{W_\omega}\ } Y. \]
\printedpage{45}
% head: A. Grothendieck : Dérivateurs, Ch. XVI, Édition des Universités de Montpellier II et Paris VII
% foot: Version 18 mars 2019
We therefore also have $i \in W$, and since $f = pi \in W$, we have $p \in W$, hence $p \in \mathrm{Parf}_{W_\omega} \cap W$. Now D$_{W_\omega}$ means that this set is contained in $W^{\mathrm{u}}$, so we have $p \in W^{\mathrm{u}}$, and in particular the fibers of $p$ are $W$-aspheric. They are therefore $W'$-aspheric, and since $p$ is $W_\omega$-perfect and a fortiori $W'$-perfect, it follows that $p \in W'$. We also have $i \in W_\omega \subset W'$, hence $i \in W'$, hence also $f = pi \in W'$, q.e.d.

Let us then return to the proof of 6.1. By the lemma, it suffices to prove that if $X$ in Cat is $W$-aspheric, it is $\mathbf{W}_\infty$-aspheric. But by the hypothesis Loc (5) we already know that $X$ is 0-connected, and choosing a point $x$ in $X$, its being $\mathbf{W}_\infty$-aspheric is expressed by the condition
\msp{page 85}
that the $\Omega^n(X,x)$ (for $n \geq 1$) are 0-connected (\textit{cf.} XII, cor. p. 190 already cited). For this, because of Loc (5) it suffices to prove that they are $W$-aspheric, and by induction it suffices to prove that $\Omega^1(X,x)$ is. Now by definition of $\Omega^1(X,x)$, this is the fiber at $x$ of an arrow in Cat
\[ E = \underline{\mathrm{Ch}}(X,x,-) \xrightarrow{\ p \in \mathrm{Parf}_{W_\omega}\ } X, \]
where the source $E$ is $W_\omega$-aspheric, a fortiori $W$-aspheric. Since $X$ is as well, we therefore have $p \in W$, hence $p \in \mathrm{Parf}_{W_\omega} \cap W \subset W^{\mathrm{u}}$, hence $p$ has $W$-aspheric fibers, and in particular $\Omega^1(X,x)$ is $W$-aspheric, q.e.d.

\textbf{Corollary 6.3.} (Characterization of $\mathbf{W}_\infty$.) \textit{$\mathbf{W}_\infty$ is the largest basic localizer satisfying} Loc (5, 7 D$_{W_\omega}$) \textit{- hence also the largest basic localizer satisfying} Loc (5, 7).

I think that we will find below an ``internal'' characterization of $\mathbf{W}_\infty$, via a ``descent axiom'' (very strong) Loc (10),
\msp{page 86}
as the unique basic localizer satisfying Loc (5, 7 D$_{W_\omega}$, 10) (or again Loc (5, 7, 10)).

6.4. \textit{Comment on the axioms for $W$-fibrations of the type} Loc (7 \ldots), \textit{and in particular on} Loc (7 D$_{W_\omega}$). \textit{Conjectural relation with the} ``simple homotopy types'' \textit{of} \textsc{Whitehead}. % [sic] du types -> du type (or des types), grammatical agreement
\DERmark{ENV0981}


a) It seems to me that \textit{all} the basic localizers with which we have tacitly worked up to now (the $\mathbf{W}_n$ ($n \in \mathbf{N}^*$), homological and cohomological equivalence, rational equivalence \ldots) correspond to localizers $W \supset \mathbf{W}_\infty$. By theorem 6.1, it therefore follows that \textit{none} of them, apart from $\mathbf{W}_\infty$ itself, satisfies the very weakest version D$_{W_\omega}$ of axiom Loc (7 \ldots), with the sole exception of Loc (7 F). Thus, the localizers $W$ which satisfy it are with the exception of $\mathbf{W}_\infty$ itself, of a quite new type - localizers which are strictly contained in $\mathbf{W}_\infty$.
\DERmark{ENV0982}


b) This raises, of course, the question of the existence of basic localizers contained in $\mathbf{W}_\infty$ and
\msp{page 87}
distinct from $\mathbf{W}_\infty$. I have not even proved that $W_\omega \neq \mathbf{W}_\infty$, although I have not the slightest doubt that this is indeed the case. \texttt{[Cisinski has proved that} $W_\omega = \mathbf{W}_\infty$\texttt{.]} A fortiori, I do not know whether the smallest basic localizer satisfying Loc (7 D$_{W_\omega}$) is different from $\mathbf{W}_\infty$. I would be quite particularly interested in the smallest localizer $\tilde{\mathbf{W}}_\infty \subset \mathbf{W}_\infty$ satisfying all the stability axioms Loc (1 ter, 2, 3 bis, 4, 6, 7, 8) (\textbf{N.B.} it satisfies
\printedpage{46}
% head: A. Grothendieck : Dérivateurs, Ch. XVI, Édition des Universités de Montpellier II et Paris VII
% foot: Version 18 mars 2019
Loc (5) necessarily, since $\tilde{\mathbf{W}}_\infty \subset \mathbf{W}_\infty \subset \mathbf{W}_0$, moreover Loc (7) is equivalent, in the presence of Loc (8), to the very weak form Loc (7 D$_{W_\omega}$): $\mathrm{Parf}_{W_\omega} \cap W \subset W^{\mathrm{u}}$. Finally, one shows that Loc (4) is a consequence of the other axioms, \textit{cf.} below.) I now conjecture that $\tilde{\mathbf{W}}_\infty$ (which gives rise to a homotopy formalism just as complete as the everlasting ``weak equivalence'') is indeed different from $\mathbf{W}_\infty$ \texttt{[no by Cisinski]}. And I suspect
\msp{page 88}
that the category $\mathrm{Hot}_{\tilde{\mathbf{W}}_\infty}$ deserves to be called the category of \textit{``simple homotopy types''} (in the style of \textsc{Whitehead})\footnote[41]{Botched (?), \textit{cf.} comment on page 92.}. A first decisive test is to check that the principal cohomological invariant of \textsc{Whitehead} of a simple homotopy type, namely an element in a $K(\mathbf{Z}[\pi_1])$ (if I remember correctly\footnote[42]{no, it is a finer invariant \ldots}) - or something close to it, can already be defined in terms of the objects in $\mathrm{Hot}_{\bullet\,\tilde{\mathbf{W}}_\infty}$ (the pointed Hot). And secondly, if one associates to a finite ``simplicial mock-up'' (describing combinatorially a finite triangulation) the ordered set of its simplices, viewed as an object of Cat, and the corresponding object in $\mathrm{Hot}_{\tilde{\mathbf{W}}_\infty}$, the latter does not change\footnote[43]{This is obvious even for $W_\omega$, since the homotopy type \textit{in the strict sense does not change!} - No, there is confusion here between topological homotopy and combinatorial homotopy, \textit{cf.} p. 92 \ldots}, up to canonical isomorphism, when one performs on this ``polyhedron'' under consideration an ele\msp{page 89}mentary operation in the style of \textsc{Whitehead}. This moreover suggests trying to define such elementary operations in the whole of Cat, or at least in the subcategory of (small) ordered sets, in the hope of arriving (perhaps) at a more direct construction of the localizer relevant to simple homotopy type.

c) I end with a comment on the meaning of axiom Loc (7 D$_{W_\omega}$), as opposed to the weaker axiom Loc (7 F) which only ensures the existence of a ``structure with fibrations'' $(W, \mathrm{Fib}_W)$ on Cat. I also point out in passing that Loc (7 F), just like Loc (5) and like Loc (9), and unlike all the other axioms reviewed up to now, is \textit{not} a stability axiom: it asserts, not that certain arrows (under suitable hypotheses) are in $W$, but the existence of a factorization. Thus it is
\msp{page 90}
not at all clear that the intersection of two basic localizers satisfying Loc (7 F) satisfies it.

Loc (7 D$_{W_\omega}$) is exactly what is needed to ensure that the canonical factorization of an arrow $f : X \longrightarrow Y$ of Cat as $f = pi$, with $i \in W_\omega$, $p \in \mathrm{Parf}_{W_\omega}$, is also a factorization of the type stipulated by the factorization axiom Loc (7 F): we have (of course) $i \in W$, but also $p \in \mathrm{Fib}_W$ (since D$_{W_\omega}$ is equivalent to its strong form $\mathrm{Parf}_{W_\omega} \subset \mathrm{Fib}_W$). This also means that the canonical factorization of $f$ can be used to ``compute'' its homotopy fibers, not only at the points of $Y$, but also with respect to a base change morphism $Y' \longrightarrow Y$.

In the absence of this axiom,
\msp{page 91}
\printedpage{47}
% head: A. Grothendieck : Dérivateurs, Ch. XVI, Édition des Universités de Montpellier II et Paris VII
% foot: Version 18 mars 2019
assuming only Loc (7 F), we know that there exists a theory of $W$-homotopy fibers, but we have a priori no practical means of ``computing'' the said fibers, for lack of a manageable criterion sufficient to ensure $p \in \mathrm{Fib}_W$. I therefore have the impression that axiom Loc (7 F) is going to be of meager use. (It is true that the example of $\mathbf{W}_0$ ought to make me cautious! I would have to look at the case of the $\mathbf{W}_n$, whether they satisfy Loc (7 F) \ldots). To my mind, axiom Loc (7 D$_{W_\omega}$) is the exact ``bare minimum'' for being at ease in handling $W$-homotopy fibers. It is remarkable that this implies that one must have $W \subset \mathbf{W}_\infty$, and that this at the same time excludes \textit{all} the variants considered to date of the notion of homotopy type!

\msp{page 92}
6.5. Finally, the relation glimpsed with the ``simple homotopy types'' of \textsc{Whitehead} vanishes\footnote[44]{But perhaps not after all - the question remains open.}: indeed, the latter are strictly finer than ordinary homotopy types (\textit{strict}), and Whitehead's invariant is \textit{not} an invariant of ordinary homotopy type. It is true that when one considers the simplicial skeleton of a triangulated polyhedron as an object of Cat, ordinary (continuous) homotopy between two such polyhedra is interpreted, in Cat, not (it seems to me) by the existence of a homotopism between the corresponding objects of Cat, but by their being isomorphic in $\mathrm{Hot}_{\mathbf{W}_\infty} \overset{\mathrm{def}}{=} \mathrm{Cat}(\mathbf{W}_\infty)^{-1}$. But the method of simplicial approximations of \textsc{Brouwer} shows that if the polyhedra $X$ and $Y$ are homotopic, then there exist simplicial subdivisions $X'$ and $Y'$ of $X$ and of $Y$ which are ``combinatorially homotopic'' (in a suitable sense). So I presume that the simplicial mock-ups of $X'$ and $Y'$ are then homotopic in Cat. It would follow that % [sic] plus fin -> plus fins, agreement with "ces derniers"
\DERmark{ENV0983}

\msp{page 93}
$X$ and $Y$, I mean their mock-ups, are indeed homotopic in Cat, if it is true that a simplicial subdivision (barycentric subdivision would suffice, iterating \ldots) $X'$ of $X$ has a mock-up in Cat which is homotopic to it. The question raised here seems interesting to me, clearly one must find the answer to it. It already arises for any ordered set $I$, which we will suppose finite to simplify, taking the set (ordered by inclusion) $I' = \mathrm{sd}(I)$ of totally ordered subsets of $I$. We know that the topological realizations of $I$ and of $\mathrm{sd}(I)$ are canonically isomorphic, but are $I$ and $\mathrm{sd}(I)$ homotopic in Cat? We do have a canonical increasing map % [sic] ensemble ordonnée -> ensemble ordonné, agreement
\[ f : \mathrm{sd}(I) \longrightarrow I, \]
which associates to every totally ordered subset of $I$ its greatest element. But is there an increasing map in the opposite direction, which is a homotopism inverse to the preceding one up to homotopy? I have just checked in the case

\msp{page 94} missing

\msp{page 95}
of the boundary of the 2-simplex

\begin{center}
\begin{tikzpicture}[baseline=(b.base)]
\node (b) at (0.95,0.3) {};
\draw (0,0) -- (1.9,0) -- (0.95,0.95) -- cycle;
\fill (0,0) circle (2.2pt);
\fill (1.9,0) circle (2.2pt);
\fill (0.95,0.95) circle (2.2pt);
\end{tikzpicture}
\qquad\qquad\qquad or \qquad\qquad\qquad
\begin{tikzpicture}[baseline=(c.base)]
\node (c) at (0,-0.1) {};
\draw (0,0) circle (0.48);
\fill (90:0.48) circle (2.2pt);
\fill (210:0.48) circle (2.2pt);
\fill (330:0.48) circle (2.2pt);
\end{tikzpicture}
\qquad\quad ,
\end{center}

barycentric subdivision
\printedpage{48}
% head: A. Grothendieck : Dérivateurs, Ch. XVI, Édition des Universités de Montpellier II et Paris VII
% foot: Version 18 mars 2019

\begin{center}
\begin{tikzpicture}[baseline=(b.base)]
\node (b) at (0.95,0.3) {};
\draw (0,0) -- (0.95,0.95) -- (1.9,0) -- cycle;
\fill (0,0) circle (2.2pt);
\fill (0.95,0) circle (2.2pt);
\fill (1.9,0) circle (2.2pt);
\fill (0.475,0.475) circle (2.2pt);
\fill (1.425,0.475) circle (2.2pt);
\fill (0.95,0.95) circle (2.2pt);
\end{tikzpicture}
\DERmark{ENV0984}

\qquad\qquad\qquad or \qquad\qquad\qquad
\begin{tikzpicture}[baseline=(c.base)]
\node (c) at (0,-0.1) {};
\draw (0,0) circle (0.48);
\fill (30:0.48) circle (2.2pt);
\fill (90:0.48) circle (2.2pt);
\fill (150:0.48) circle (2.2pt);
\fill (210:0.48) circle (2.2pt);
\fill (270:0.48) circle (2.2pt);
\fill (330:0.48) circle (2.2pt);
\end{tikzpicture}
\end{center}

(hexagon), that for every increasing map $g : I \longrightarrow \mathrm{sd}(I)$, its topological realization is homotopic to a constant map, so $g$ cannot be a homotopism in Cat (for this would imply that its topological realization is a homotopism).

This therefore shows that the notion of homotopism in Cat is \textit{strictly finer} than that for the topological realizations, which once again restores substance to the hope that Whitehead's invariant could be defined as a homotopy invariant on Cat, or at least on a suitable full subcategory, such as that of ordered sets (possibly restricted to be finite), or of simplicial mock-ups. And even, that the set of Wh-homotopy types (simple homotopy types) can be
\msp{page 96}
canonically embedded in a suitable $\mathrm{Hot}_W$, where $W$ would be a basic localizer strictly finer than $\mathbf{W}_\infty$, e.g. (who knows?) $\tilde{\mathbf{W}}_\infty$. The first piece of work one would like to do is to construct such a $W$ \textit{ad hoc}, so that Whitehead's invariant can be written as a functor on $\mathrm{Hot}_{\bullet W}$. To begin with, one would already have to make sure that this invariant can be viewed as an invariant of homotopy type, \textit{i.e.} that it is interpreted as a functor on $\mathrm{Hot}_{\bullet}\mathrm{Cat} \simeq \mathrm{Cat}(\,\underbrace{\text{homotop}}_{\substack{\text{set of}\\ \text{homotopisms}}}\,)^{-1}$. This work will perhaps also suggest the definition of ``higher Whitehead invariants'' (involving the $\pi_i$, not only the $\pi_1$), which could perhaps serve (with the one of dimension 1) to describe the hypothetical localizer $W$ (``of Whitehead'').

\msp{page 97}

\section*{7\quad Comments on the axioms}

7.1. The main reason for being of a basic localizer $W$ is to define a theory of homotopy, and more precisely a prederivator $\mathrm{HOT}_W$, of which one hopes that it will even be a derivator, or that it will at least satisfy some of the axioms of derivators. The axioms Loc (1) to Loc (9) are for the most part oriented toward this objective, except that Loc (1, 2, 3) constitutes an absolute ``bare minimum'', to be able to begin working with $W$. The main questions which arise seem to me to be the following: % [sic] d'homotopie’ -> d'homotopie (or “théorie d'homotopie”), closing quote without opening quote
\DERmark{ENV0985}

\begin{enumerate}[label={}, leftmargin=2.5em]
\item[1)] When is Cat equipped with $(W, \mathrm{Cof}_W)$, resp. with $(W, \mathrm{Fib}_W)$, a category with cofibrations resp. with fibrations, so that the results of K. S. \textsc{Brown} are applicable?
\item[2)] Under what conditions, on $W$ and on $f$, do $f_!$ and $f_*$ exist? And when do we have, in addition, exactness of cocartesian resp. cartesian squares?

\msp{page 98}
\item[3)] More generally, in terms of the axioms imposed on $W$, which are the axioms Der (\ldots) that are satisfied for the prederivator $\mathrm{HOT}_W$? (Which could also depend on the choice of the domain Diag $\subset$ Cat. The prederivator is defined a priori on the whole of Cat). In particular, % [sic] quelles sont les axiomes -> quels sont les axiomes, agreement; [sic] que sont vérifiés -> qui sont vérifiés, relative pronoun
\end{enumerate}


\endgroup
% END UNIT der-ch16/der-ch16_p046-050.tex
% BEGIN UNIT der-ch16/der-ch16_p051-055.tex
\begingroup\makeatletter%
\newcommand{\printedpage}[1]{\marginpar{\footnotesize p.~#1}}%
\newcommand{\msp}[1]{\marginpar{\footnotesize [#1]}}%
\makeatletter%
\renewcommand{\@makefnmark}{\hbox{\normalfont\@textsuperscript{(\@thefnmark)}}}%
\renewcommand{\@makefntext}[1]{\noindent\hbox{\normalfont\@textsuperscript{\@thefnmark}}\,#1}%
\makeatother%
\makeatother\let\printedpage\derpp\let\msp\dermsp%
\let\mspin\DERmspin%


\printedpage{49}
% head: A. Grothendieck : Dérivateurs, Ch. XVI, Édition des Universités de Montpellier II et Paris VII
% foot: Version 18 mars 2019

\begin{enumerate}[label=4)]
\item when is $\mathrm{HOT}_W$ really and truly a derivator on all of Cat, \textit{i.e.} when does it satisfy all the axioms?
\end{enumerate}

A question which seems to me less important (in particular than 1°) is the following:

\begin{enumerate}[label=1 bis)]
\item When does $W$ fit into a closed Quillen triple $(W,C,F)$ (preferably proper).
\end{enumerate}

I shall go over these questions one by one.

7.2. We have seen (but without my having made a careful check of the converse) that $(W,\mathrm{Cof}_W)$ defines a structure with cofibrations if the axiom \textit{Loc (6 a)} is satisfied, and that this condition is also necessary if Loc (3 bis) is satisfied\footnote[45]{Check all the same whether Loc (3 bis) is not already used in the forward direction! \textbf{O.K.}}.

\msp{page 99}
Then the open Dwyer immersions (or, at one's choice, the closed Dwyer immersions) are in $\mathrm{Cof}_W$, and the standard mapping cone construction defines, for every arrow $f$ in Cat, a \textit{functorial} factorization
\[
\begin{aligned}
f = pi \quad & i \text{ open Dwyer immersion (hence in } \mathrm{Cof}_W\text{),}\\
& p \in W.
\end{aligned}
\]
This, plus the construction of the functor $f_!$, and in particular of $W$-amalgamated sums
\begin{tikzcd}[row sep=tiny, column sep=small, baseline=(\tikzcdmatrixname-2-1.base)]
 & X_1 \\
\displaystyle\int X_0 \arrow[ur] \arrow[dr] & \\
 & X_2
\end{tikzcd}, given in VI (theory of $\mathrm{HOT}_W$) and XII, gives us perfect mastery (it seems to me) over the ``left homotopy structure'' defined by $W$, under the sole axiom Loc (6 a) (Cube axiom, direct form). I recall that the functors $f_!$ are constructed in VI under the sole hypothesis \textit{Loc (3 bis)} ($f$ any arrow of Cat), but unless I am mistaken, the construction
\msp{page 100}
and the $W$-homotopic interpretation of
\begin{tikzcd}[row sep=tiny, column sep=small, baseline=(\tikzcdmatrixname-2-1.base)]
 & X_1 \\
\displaystyle\int X_0 \arrow[ur] \arrow[dr] & \\
 & X_2
\end{tikzcd}
given in XII do not even depend on this hypothesis\footnote[46]{\textit{cf.} XII, page 264} - but I cannot guarantee it. It seems to me that XII (and in particular the long §5 on amalgamated sums) is written under the hypothesis that W (1, 2, 3), \textit{i.e.} Loc (1, 2, 3bis), are satisfied.

7.3. It is a rather extraordinary thing that the functors $f_!$ for $\mathrm{HOT}_W$ exist, under the sole hypothesis \textit{Loc (3 bis)} (localization axiom, strong form) - and they are ``computed'' in a magically simple way! I am pretty sure that under this hypothesis alone, \textit{the exactness axiom for cocartesian squares} is \textit{not} satisfied\footnote[47]{An example would be needed!}. I see no way of checking it, except via \textsc{Brown}'s theory, under the hypothesis \textit{Loc (6 a)} (\textit{cf.} 7.2).

But under the sole condition Loc (3 bis) the standard calculation of the fibers of an $f_!(\xi)$
\printedpage{50}
% head: A. Grothendieck : Dérivateurs, Ch. XVI, Édition des Universités de Montpellier II et Paris VII
% foot: Version 18 mars 2019
\msp{page 101}
is valid, in other words the prederivator $\mathrm{HOT}_W$ satisfies Der (4$^*$, 5$^*$) - and moreover Der (6$^*$) (exactness ``on the left'') if $W$ satisfies Loc (6 a). % [sic] d'autre termes -> d'autres termes, missing plural s

7.4. Let us come to the question whether Cat equipped with $(W,\mathrm{Fib}_W)$ is a category with fibrations. This is the axiom \textit{Loc (7 F),} but I know scarcely any way of checking it (apart from the very, very special case $W = \mathbf{W}_0$) other than by means of at least the axiom \textit{Loc (7 $D_{W_\omega}$),} or one or another of the stronger variants of Loc (7). Loc (7 $\mathrm{D}_{W_\omega}$) has been discussed at length above, \textit{cf.} in particular 6.4 (pages 89-91). In my view, this axiom ensures a mastery over the right $W$-homotopic structure ($W$-homotopic fibers, exact cosuspension sequences) comparable to that of Loc (6 a) for the
\msp{page 102}
left structure. In particular, we have the \textit{functorial} factorization of every arrow $f$ of Cat as
\[
f = pi, \qquad i \in W_\omega \subset W,\ p \in \mathrm{Parf}_{W_\omega} \subset \mathrm{Fib}_W .
\]
But \textit{the axiom Loc (7 $D_{W_\omega}$) seems to me to be of a considerably stronger nature than Loc (6 a),} which (unless I am mistaken) is satisfied for all the localizers that have been considered so far. (But I admit not having looked at the case of the $\mathbf{W}_n$, $n \in \mathbf{N}$, except $\mathbf{W}_0$.) On the other hand, the axiom Loc (7 $\mathrm{D}_{W_\omega}$) implies $W \subset \mathbf{W}_\infty$ (theorem 6.1), a truly draconian relation (\textit{cf.} comment 6.4 a).

Besides, the axiom Loc (7 $\mathrm{D}_{W_\omega}$) alone, as far as I know, does not imply the existence of $f_*$ for \textit{every} arrow in Cat, but only for
\msp{page 103}
arrows of the form $f \times \mathrm{id}_X : S \times X \longrightarrow T \times X$, where $f : S \longrightarrow T$ is a morphism between finite ordered sets. Or more generally, between ordered sets $S$, $T$ such that $S^\circ$ is artinian, \textit{i.e.} $\mathrm{S}$ noetherian, and $\mathrm{Esp}(S^\circ)$ is noetherian (this last hypothesis being unnecessary under the axiom Loc (8) on limits, \textit{cf.} XVII, th. 3.1, p. 50). So this is not the result dreamed of.

7.5. The only case I know where $f_*$ exists for \textit{every} arrow in Cat, is the case where $W$ satisfies
\[
\textit{Loc (6 a, 8, 9)}\ \ \text{\footnotemark[49]\,\footnotemark[50]}, \tag*{$(*)$} % [sic] marks printed (48) (49) -> the two footnotes printed below are numbered 49 and 50, numbering mismatch between marks and footnotes
\]
\DERmark{ENV0986}

\footnotetext[49]{Check whether the condition Loc (3 bis) is not also needed in addition to Loc (6 a, 8, 9). Check made, it is O.K. (without Loc (3 bis)).}
\footnotetext[50]{\textbf{N.B.} On further reflection, it appears that even Loc (6$'$ a, 8, 9) suffices.}
that is to say the axiom of weak $W$-cofibrations, already mentioned at length and finally those of limits and accessibility, \textit{i.e.} $W$ is $\mathrm{L}_1$-accessible in \underline{Fl}(Cat). The relevant key result is given in XV 8.4 (p. 105-107), which must be combined with XV 7.4 (p. 98-99) to make
\msp{page 104}
good use of it, and to deduce theorem 10.1 (XV, p. 134-135). Better than the existence of the $f_*$, the surprisingly weak set of axioms I have just spelled out suffices to ensure that $\mathrm{HOT}_W$ is a \textit{derivator} on all of Cat, in the full sense of the term, \textit{i.e.} that it satisfies Der (1, 2, 3, 4, 5, 6, 4$^*$, 5$^*$, 6$^*$). This is obtained by replacing (thanks to XV 8.4) Cat by
\printedpage{51}
% head: A. Grothendieck : Dérivateurs, Ch. XVI, Édition des Universités de Montpellier II et Paris VII
% foot: Version 18 mars 2019
$A^\wedge$, where $A$ is a suitable test category, $W$ by $W_A \subset \mathrm{Fl}\, A^\wedge$, and proving that $W_A$ fits into a closed Quillen triple $(W,C,F)$, with $C = \mathrm{Mono}(A^\wedge)$, which determines $F$ as $(C \cap W_A)_*$. We then apply XV 7.4 to this situation (results summarized in XIX, th. 2.3 (p. 19), 2.6 (p. 25), 2.9 (p. 28). We recover the existence of the $f_!$ (with a method of calculation quite different from the one in Cat!), and in addition we find the existence of $f_*$ (plus the exactness in cor. [?] 5
\msp{page 105}
of $W$-cartesian squares) - which for me, for a good while, was \textit{the} big problem for the theory of $\mathrm{HOT}_W$!).

I find it staggering that one arrives at such a strong result with such ridiculously weak axioms. More particularly, that one has all the desirable properties of the right homotopy structure, without having to make any hypothesis of the $W$-fibration type on $W$! These axioms do \textit{not} imply the axiom Loc (7) (equivalent to Loc (7 $\mathrm{D}_{W_\omega}$) in the presence of Loc (8)), as we have seen from the example of $\mathbf{W}_0$. Do they at least imply Loc (7 F), \textit{i.e.} the existence of enough $W$-fibrations? I would not venture any conjecture on this subject! But one must bear in mind that in the absence of Loc (7) (and even
\msp{page 106}
assuming Loc (7 F)), while knowing that the right homotopy structure associated with $W$ is perfect in every respect, one has a priori no mastery over it, in particular for computing $W$-homotopic fibers and constructing $W$-cartesian squares. Even on reducing all these constructions to $A^\wedge$ (taking, if one insists, $A = \boldsymbol{\Delta}$, \textit{i.e.} reducing to the semi-simplicial context), though one knows that there is a \textit{functorial} factorization $f = pi$ with $i \in \mathrm{TC} \subset W_A$, $p \in F \subset \mathrm{Fib}_{W_A}$, I see no concrete way of making it explicit. Unless the consideration of infinite paths, which (in the absence of Loc (7)) turns out to be inoperative in Cat saves the day in $\boldsymbol{\Delta}^\wedge$? ?

\msp{page 107}
A question which ought to be easy to clear up, but on which I do not want to linger just now.

Let us note in passing that \textsc{Quillen}'s strong saturation result, for closed Quillen triples, shows that the set of axioms $(*)$ implies the axiom Loc (1 ter), \textit{i.e.} that $W$ is \textit{strongly saturated.}

Let us also note that unlike what happens for the $f_!$, I know how to prove the existence of $f_*$, and in particular of cosuspension sequences, only under conditions (either Loc (7 $\mathrm{D}_{W_\omega}$), or Loc (6a, 8, 9)) which already ensure that the $W$-cartesian squares are exact, and consequently the cosuspension sequences exact.

I think that in the preceding pages, I have said all that I know how to say that is relevant, concerning questions 1, 2, 4 raised in 7.1. These answers se\msp{page 108}em satisfactory to me - more satisfactory than I would have dared hope during the first two months of my cogitations! There remain questions 3 and 1 bis, which I shall now examine.

\printedpage{52}
% head: A. Grothendieck : Dérivateurs, Ch. XVI, Édition des Universités de Montpellier II et Paris VII
% foot: Version 18 mars 2019
7.6. On the different derivator axioms: under what conditions on $W$ is one of these axioms satisfied for $\mathrm{HOT}_W$\footnote[51]{For axioms Der (1) to Der (5$^*$), \textit{cf.} I, p. 67-68. The axioms Der (6, 6$^*$) are the exactness axioms, whose necessity became apparent to me only later.}? As for axioms Der (4, 5, 6, 4$^*$, 5$^*$, 6$^*$), I think I have said all that I know how to say about them, and have perfectly manageable criteria for validity, in the preceding pages. Let us summarize nevertheless: % [sic] ce validité -> de validité, misprint
\[
\left\{
\begin{array}{lcl}
\text{structure with cofibrations} & : & \text{Loc (6) (+ Loc (3 bis) ?)} \\[1ex]
\text{structure with fibrations} & : & \text{Loc (7 }\mathrm{D}_{W_\omega}\text{)} \\[1ex]
\text{Der (4}^*\text{, 5}^*\text{)} & : & \text{Loc (3 bis)} \\[1ex]
\text{Der (4}^*\text{, 5}^*\text{, 6}^*\text{)} & : & \text{Loc (3 bis, 6a) or Loc (6}'\text{ a, 8, 9)} \\[1ex]
\text{Der (4, 5, 6) (on Cat) (and all the rest !)} & : & \text{Loc (6}'\text{ a, 8, 9)} \\[1ex]
\text{Der (4, 5, 6) (on Diag $=$ finite ordered sets)} & : & \text{Loc (7 }\mathrm{D}_{W_\omega}\text{) (or only Loc (7 F)).}
\end{array}
\right.
\]

There remain the axioms Der (1) (decomposition), Der (2) (localization), Der (3) (unit object), which I shall now go over. One checks
\msp{page 109}
quasi-tautologically that the axiom Der 1 is satisfied with no condition on $W$ (not even Loc (2, 3)!) - in fact, it is satisfied for every prederivator $\mathbf{D}_{(\mathcal{M},W)}$ defined by a category of models.

The localization axiom Der (2), on the other hand, is delicate in nature. It is valid, more generally, for a category of models $(\mathcal{M},W)$ (defining a prederivator $\mathbf{D} = \mathbf{D}_{(\mathcal{M},W)}$) and an object $S$ of Cat, when $(\mathcal{M}(S),W(S))$ is associated with a structure with fibrations or with cofibrations, which by K. S. \textsc{Brown} implies that every arrow of $\mathrm{HO}_W(S) = \mathcal{M}(S)W(S)^{-1}$ is ``realized'' by an arrow of $\mathcal{M}(S)$ - and then the result is obvious. The condition I have just stated is satisfied when $(\mathcal{M},W)$ itself is associated with a structure with cofibrations or fibrations,
\msp{page 110}
and if moreover we assume that we are in one of the following two cases: % [sic] l'un de deux cas -> l'un des deux cas, apparent misprint
\begin{enumerate}[label=\alph*)]
\item $S$ is a finite ordered set (which leads us to restrict ourselves, for want of anything better, to Diag $=$ subcategory of Cat formed by the finite ordered sets).
\item The factorization stipulated in the factorization axiom $f = pi$ (with $i \in \mathrm{Cof}$, $p \in W$ (cofibration case) or $i \in W$, $p \in \mathrm{Fib}$ (fibration case)) can be chosen functorial in $f$.
\end{enumerate}
Now, in the case of Cat, in the cases where one knows how to prove that $W$ is associated with a structure with cofibrations (axiom \textit{Loc (6 a)}) or with fibrations (axiom \textit{Loc (7 $D_{W_\omega}$)}), one does indeed find a functorial factorization. Thus each of these axioms implies Der (2). % [sic] où à fibrations -> ou à fibrations, accent misprint

\msp{page 111}
There remains the contractibility axiom: if $S$ in Cat is contractible, is it $\mathrm{HOT}_W$-aspheric, \textit{i.e.} is the functor
\[
\mathrm{HOT}_W(e) = \mathrm{Hot}_W \longrightarrow \mathrm{HOT}_W(S)
\]
\printedpage{53}
% head: A. Grothendieck : Dérivateurs, Ch. XVI, Édition des Universités de Montpellier II et Paris VII
% foot: Version 18 mars 2019
fully faithful? When one has the theory of $f_!$ ``Cat version'' (thus under the sole axiom \textit{Loc (3 bis)}), one finds the following (noted in VI, p. 202)\footnote[52]{under Loc (3 bis).}:
\begin{quote}
Let $f : S \longrightarrow T$ be an arrow in Cat. For it to be a $\mathrm{HOT}_W$-equivalence, \textit{i.e.} for $f^* : \mathrm{HOT}_W(T) \longrightarrow \mathrm{HOT}_W(S)$ to induce a \textit{fully faithful} functor on the subcategory of constant objects of $\mathrm{HOT}_W(T)$, it is necessary and sufficient that $\mathrm{hot}_W(f) : \mathrm{hot}_W(S) \longrightarrow \mathrm{hot}_W(T)$ be an isomorphism. Thus it \textit{suffices} that $f \in W$, and this condition is also necessary if $W$ is strongly saturated (axiom Loc (1 ter)).
\end{quote}
To ease my conscience, I spell out the proof here. The existence of the functors
\msp{page 112}
$f_!$, in particular of the homology objects $H_\bullet^{\mathrm{HOT}_W}(S,\xi)$ (or simply $H_\bullet^W(S,\xi)$) for $\xi \in \mathrm{HOT}_W(e) = \mathrm{Hot}_W$, makes it possible to interpret homologically the condition of $\mathrm{HOT}_W$-equivalence, as meaning that for every $\xi = \mathrm{hot}_W(X)$ in $\mathrm{Hot}_W$, the canonical arrow in $\mathrm{Hot}_W$ % [sic] (où simplement -> (ou simplement, accent misprint
\[
H_\bullet^W(S,\xi) \longrightarrow H_\bullet^W(T,\xi) \tag*{$(*)$}
\]
is an isomorphism. Now the calculation of the $f_!$ (in the case of $p : S \longrightarrow e$ and of $T \longrightarrow e$) shows us that the two sides of $(*)$ are none other (up to isomorphism) than $\mathrm{hot}_W(X \times S)$ and $\mathrm{hot}_W(X \times T)$, the arrow being induced by
\[
\mathrm{id}_X \times f : X \times S \longrightarrow X \times T. \tag*{$(*)$}
\]
Thus the condition under consideration means that this arrow is an isomorphism in $\mathrm{Hot}_W$ for every $X$ in Cat. Taking $X = e$, one sees that this implies that $\mathrm{hot}_W(f)$ is an isomorphism, and this suffices, since the
\msp{page 113}
functor $\mathrm{hot}_W : \mathrm{Cat} \longrightarrow \mathrm{Hot}_W$ commutes with finite products (Loc (1, 2, 3) suffices) \texttt{[rather Loc (1, 2, 3 bis)]} and the arrow $(*)$ in $\mathrm{Hot}_W$ identifies with
\[
\mathrm{id}_{\mathrm{hot}_W(X)} \times \mathrm{hot}_W(f)\,.
\]
This therefore shows that if $S$ is a $W$-aspheric object, \textit{i.e.} such that $S \longrightarrow e$ is in $W$ (and a fortiori, if $S$ is contractible), then $S \longrightarrow e$ is $\mathrm{HOT}_W$-aspheric.

To summarize, we get the table
\[
\left\{
\begin{array}{lcl}
\text{Der (1)} & : & \text{with no hypothesis on } W \\[1ex]
\text{Der (2)} & : & \text{Loc (6 a) \textit{or} Loc (7 }\mathrm{D}_{W_\omega}\text{)} \\[1ex]
\text{Der (3)} & : & \text{Loc (3 bis) \textit{or} Loc (6}'\text{ a, 8, 9)\ \ \footnotemark[53].}
\end{array}
\right.
\]
\footnotetext[53]{\textbf{N.B.} We shall see below that Loc (6 a, 8, 9) already implies Loc (3 bis).}

7.7. It only remains to justify that Loc (6$'$ a, 8, 9) implies Der (3). But then, as we have seen, $\mathrm{HOT}_W$ can be defined by a very special Quillen triple. But even for a merely special Quillen triple $(\mathcal{M},W)$, the prederivator $\mathbf{D}$ which it defines satisfies (in addition to Der (4$^*$, 5$^*$, 6$^*$)) the conditions Der (1) (trivially), Der (2)


\endgroup
% END UNIT der-ch16/der-ch16_p051-055.tex
% BEGIN UNIT der-ch16/der-ch16_p056-059.tex
\begingroup\makeatletter%
\newcommand{\printedpage}[1]{\marginpar{\footnotesize p.~#1}}%
\newcommand{\msp}[1]{\marginpar{\footnotesize [#1]}}%
\newcommand{\circled}[1]{\tikz[baseline=(c.base)]\node[draw,circle,inner sep=1pt](c){#1};}%
\makeatother\let\printedpage\derpp\let\msp\dermsp%
\let\mspin\DERmspin%


\printedpage{54}
% head: A. Grothendieck : Dérivateurs, Ch. XVI, Édition des Universités de Montpellier II et Paris VII
\msp{page 114}
(\textit{cf.} p. 110, case b)) and \textit{in addition} Der (3). (Thus in the case of a \textit{very} special Quillen triple, the prederivator obtained is indeed a derivator.) I will spell out the proof here.

I will denote by $(\mathcal{M},\Sigma)$ the given category of models ($\Sigma\subset\operatorname{Fl}\mathcal{M}$), by $\mathbf{D}=\mathbf{D}_{\mathcal{M},\Sigma}$ the corresponding prederivator, and by
\DERmark{ENV0987}

\[ W=W_{\mathbf{D}} \] % [sic] W_D -> W_{\mathbf{D}}, index printed as italic D, elsewhere bold D
\DERmark{ENV0988}\DERmark{ENV0989}

the associated localizer in Cat, consisting of the ``$\mathbf{D}$-equivalences''. We see at once that $W$ is strongly saturated (true for any $W_{\mathbf{D}}$, $\mathbf{D}$ a prederivator), \textit{i.e.} satisfies Loc (1 ter), a fortiori Loc (1), on the other hand it satisfies Loc (3 bis$^*$), as a consequence of the fact that $\mathbf{D}$ satisfies Der (4$^*$, 5$^*$). A fortiori it satisfies Loc (3$^*$), and to see that every $X$ contractible in Cat is $W$-aspheric, \textit{i.e.} $X\longrightarrow e$ in $W$, we know (\textit{cf.} §1)
\DERmark{ENV0990}

\msp{page 115}
that it suffices to prove that this is so if $X$ has an initial object (axiom Loc (2$^*$) dual to axiom Loc (2)). At that point, we will know that $W$ is a basic localizer, and moreover it satisfies Loc (3$^*$ bis), hence also Loc (3 bis) (strong form of the localization axiom).

Thus \textit{we must prove} that if $S$ has an initial object, then $S$ is $W$-aspheric, \textit{i.e.} $S\longrightarrow e$ is in $W$, \textit{i.e.} $p^*:\mathbf{D}(e)\longrightarrow\mathbf{D}(S)$ is fully faithful. It amounts to the same thing to say that the adjunction arrow
\DERmark{ENV0991}

\[ p_!p^*\xi\longrightarrow\xi \]
is an isomorphism, for every $\xi$ in $\mathbf{D}(e)$. Now the computation of $p_!(\eta)$ (for $\eta\in\operatorname{Ob}\mathbf{D}(S)$, \textit{i.e.} $\eta\in\operatorname{Ob}\mathcal{M}(S)$) is done as follows: we choose an arrow
\[ \eta'\xrightarrow{\Sigma_S}\eta \]
with $\eta'$ ``cofibrant'' (for the relevant cofibration structure on $\mathcal{M}(S)$), and we
\msp{page 116}
set \footnote[54]{operation $p_!$ at the level of $\mathcal{M}(S)$.}
\[ p_!(\operatorname{ho}_{\Sigma,S}(\eta))=\operatorname{ho}_\Sigma(p_!(\eta')). \]
\DERmark{ENV0992}\DERmark{ENV0993}

When $\eta=p^*(\xi)$, the relevant adjunction arrow in $\operatorname{Ho}_\Sigma=\mathbf{D}(e)$ is obtained by composition
\[ p_!(\eta')\longrightarrow p_!(p^*(\xi))\xrightarrow{\text{adj.}}\xi, \tag*{$(*)$} \]
where here $p_!$ is still taken in the ordinary sense ($\mathcal{M}(S)\longrightarrow\mathcal{M}(e)=\mathcal{M}$). Now since $S$ has an initial object 0, we have
\[ \begin{aligned}
p_!(\eta') &\ \simeq\ \eta'(0)\\
p_!(p^*(\xi)) &\ \simeq\ p^*(\xi)(0)\ \simeq\ \xi,
\end{aligned} \]
and the adjunction arrow in $(*)$ is the identity, the first arrow is the arrow $\eta'(0)\longrightarrow\eta(0)$ induced by $\eta'\longrightarrow\eta$, it is in $\Sigma$, hence $(*)$ is in $\Sigma$, hence it is an isomorphism in $\operatorname{Ho}_\Sigma=\mathbf{D}(e)$, q.e.d.

The same argument works, moreover, if $\mathcal{M}$ is only a category with cofibrations or with fibrations, but then restricting oneself to the case where $S$ is a finite ordered set.
% foot: Version 18 mars 2019

% Original hard page break suppressed for the reflowed reader.
\printedpage{55}
% head: A. Grothendieck : Dérivateurs, Ch. XVI, Édition des Universités de Montpellier II et Paris VII
\msp{page 117}
I will set out separately what we have just obtained:

\textbf{Proposition 7.7.1.} \textit{Let} $(\mathcal{M},W)$ \textit{be a category of models, associated to a closed Quillen triple which is} special $(W,C,F)$\textit{. Then the prederivator} $\mathbf{D}=\mathbf{D}_{\mathcal{M},W}$ \textit{on} Cat \textit{satisfies the axioms} Der (1, 2, 3, 4$^*$, 5$^*$, 6$^*$)\textit{. If moreover the Quillen triple is} very \textit{special} (\textit{thus in particular} $C=\operatorname{Mono}(\mathcal{M})$), \textit{then it satisfies in addition} Der (4, 5, 6), i.e. \textit{it is a derivator in the full sense of the term. On the other hand, the localizer} $W=W_{\mathbf{D}}$ \textit{in} Cat \textit{associated to} $\mathbf{D}$ (\textit{in the case of a special Quillen triple and no more}) \textit{is a} basic localizer, \textit{satisfying in addition} Loc (1 ter, 3 bis).

\textbf{N.B.} What is said here for the case of a very special Quillen triple (which does not call upon the theory of basic localizers
\msp{page 118}
in Cat) shows once again that if $W$ a basic localizer in Cat satisfies Loc (6$'$ a, 8, 9), then $\mathrm{HOT}_W$ is a genuine derivator. We saw moreover in 7.6 (page 111) that we then have $W=W_{\mathrm{HOT}_W}$ \footnote[55]{But this was under the hypothesis Loc (3 bis)! Beware of the vicious circle.}, hence it follows:

\textbf{Corollary 7.7.2:} \textit{We have the implication}
\[ \text{Loc } (6'\text{ a, 8, 9})\ \Longrightarrow\ \text{Loc(1 ter, 3 bis)}\ \text{\footnotemark[56]}. \]
\footnotetext[56]{Loc (3 bis) not proved, probably false.}

Let us also note, for the record:

\textbf{Corollary 7.7.3.} \textit{Let} $(\mathcal{M},W)$ \textit{be a category of models, associated to a category with cofibrations. Then on the domain} Diag \textit{consisting of finite ordered sets, the prederivator} $\mathbf{D}=\mathbf{D}_{\mathcal{M},W}$ \textit{satisfies the axioms} Der (1, 2, 3, 4$^*$, 5$^*$, 6$^*$), \textit{and dually. Thus if} $W$ \textit{is associated both to a cofibration structure and to a fibration structure,} $\mathbf{D}$ \textit{is a derivator on} Diag.

\msp{page 119}
\texttt{[The following six lines seem out of context.]}

\textit{then the base change arrow}
\[ p^*f_!\longrightarrow f'_!q^* \]
\textit{between functors} $\mathrm{HOT}_W(X)\longrightarrow\mathrm{HOT}_W(S')$ \textit{is an isomorphism.}

This is practically trivial, considering, in the case $f$ $W$-smooth, that if $Z$ over $X$ is $W$-smooth over $X$, it is so over $S$; and in the case $p$ $W$-proper, using the last statement (``omission'') in \circled{6}, p. 199.

7.8. I finish by putting in the spotlight what, at present, seems to me the main substance of all these ``reminders'' and as it were the quintessence of the theory of basic localizers
% margin: three vertical rules printed in the left margin, beside the footnotes and down to the foot of the page
% foot: Version 18 mars 2019
\printedpage{56}
% head: A. Grothendieck : Dérivateurs, Ch. XVI, Édition des Universités de Montpellier II et Paris VII
developed up to now:
\[ \left\{ \begin{array}{rcl}
\text{Loc (3 bis)} & \Longrightarrow & \text{Der (4$^*$, 5$^*$)}\\[1ex]
\text{Loc (3 bis, 6 a)} & \Longrightarrow & \text{Der (1, 2, 3, 4$^*$, 5$^*$, 6$^*$)}\\[1ex]
\text{Loc (6$'$ a, 8, 9)} & \Longrightarrow & \text{all the axioms of derivators (+ Loc (1 ter)).}
\end{array} \right. \]

I consider the axioms Loc (3 bis), Loc (8), Loc (9) as ``harmless'' axioms - all the localizers one defines by ordinary cohomological or homotopical means, satisfy them. There remains the axiom Loc (6 a) or Loc (6$'$ a) of $W$-cofibrations (weak form),
\msp{page 120}
which visibly plays a crucial role in the theory of $\mathrm{HOT}_W$. On the other hand, the axioms Loc (4) (decomposition), Loc (5) (connectedness) and Loc (7) and its variants (and of course also Loc (6 b)) do not enter in at all! This is particularly impressive, as regards the axiom of $W$-fibrations Loc (7) and its variants (all equivalent to Loc (7) in the presence of Loc (8), with the sole exception of Loc (7 F)). As we saw, Loc (7) is not at all implied by Loc (6 a, 8, 9) ($\mathbf{W}_0$ does not satisfy it), at most it could be that Loc (7 F) is. The same is true of the connectedness axiom Loc (5), as one sees from the example of $W=\mathrm{Fl}(\mathrm{Cat})$, which satisfies all the axioms seen so far, with the sole exception
\msp{page 121}
of Loc (5). As for the question of the independence of the axiom Loc (4) from Loc (6 a, 8, 9), we are going to examine this kind of question in the following paragraph.

The set of axioms Loc (6$'$ a, 8, 9) now seems to me so important, that I feel like giving it a name, and calling a \textit{special} (basic) \textit{localizer} in Cat \footnote[57]{I hesitate whether I should not add the axiom Loc (3 bis), which ensures that $W=W_{\mathbf{D}}$, $\mathbf{D}=\mathrm{HOT}_W$.}, a basic localizer which satisfies them. These are thus also the ones for which we know that $\mathrm{HOT}_W$ is a genuine \textit{derivator.} If one cannot dismiss as ``pathological'' a structure as rich and delicate as a derivator, it follows at once that one cannot dismiss special localizers in Cat either (even if one affects to look down on the
\msp{page 122}
general basic localizers!). Indeed, $W$ is recovered uniquely with the help of $\mathbf{D}=\mathrm{HOT}_W$, or of any derivator equivalent to it, as the set of $\mathbf{D}$-equivalences in Cat \footnote[58]{But this is proved only assuming Loc (3 bis).}: it is the basic localizer (corresponding to a notion of ``universal'' equivalence, \textit{i.e.} in the category Cat of all the elementary ``spaces'') canonically associated to the derivator. It is therefore by no means a sufficiently fine fixed reference localizer, such as $\mathbf{W}_\infty$, that is fit to play a role!

This brings us again to a question of the kind already raised: if $W\subset\mathbf{W}_\infty$ is the smallest localizer which satisfies Loc (6$'$ a, 8) (or Loc (3 bis, 6$'$ a, 8)), can we say a) that $W$ is even \textit{special}, \textit{i.e.}
\msp{page 123}
satisfies Loc (9) (accessibility) as well?, and moreover b) do we have $W\neq\mathbf{W}_\infty$? I am convinced that this is the case (there is less merit for that $W$ in being different from $\mathbf{W}_\infty$,
% foot: Version 18 mars 2019
\printedpage{57}
% head: A. Grothendieck : Dérivateurs, Ch. XVI, Édition des Universités de Montpellier II et Paris VII
than for the $\tilde{\mathbf{W}}_\infty$ considered in 6.4 (p. 87)).

7.9. Omission: question 1 bis (p. 98), on the existence of a closed Quillen structure on Cat, corresponding to $W$. The only result I know in this direction is that of \textsc{Thomason}, explicitly constructing such a Quillen structure (moreover proper), in the case where $W=\mathbf{W}_\infty$. I have not really looked at the proof of \textsc{Thomason}, and do not know whether it can be extended to more general basic localizers, for example to those satisfying Loc (6$'$ a, 8, 9). If this were so, it would moreover prove that these localizers
\msp{page 124}
have enough cofibrations and enough fibrations, hence that they satisfy Loc (6 a) and Loc (7 F). (As for Loc (6 a), it seems that one can prove it directly very simply, \textit{cf.} the following paragraph. But here again, one would have to look at some examples (homological or cohomological equivalences, $\mathbf{W}_n$, \textit{i.e.} the homotopy $n$-equivalences), and see to begin with whether Loc (7 F) is satisfied.)

It does not seem to me that the theory of \textsc{Thomason} has been up to now \texttt{[illegible words]} to be very useful, technically speaking. Psychologically, I think it has been important, for giving confidence that Cat really is an excellent category of models for Hot. I was really delighted and stimulated to learn of this theorem. But the notion of a Dwyer map (which I learned in 1983 in the article of \textsc{Thomason}) has been, technically, much more important for me: it was exactly what was missing from my panoply of ``homotopical'' notions and intuitions in Cat \dots
% foot: Version 18 mars 2019


\endgroup
% END UNIT der-ch16/der-ch16_p056-059.tex
\clearpage\setcounter{footnote}{0}\phantomsection\addcontentsline{toc}{section}{Chapter XVII}\markboth{Chapter XVII}{Chapter XVII}
% BEGIN UNIT der-ch17/der-ch17_p001-005.tex
\begingroup\makeatletter%
\newcommand{\printedpage}[1]{\marginpar{\footnotesize p.~#1}}%
\newcommand{\msp}[1]{\marginpar{\footnotesize [#1]}}%
\newcommand{\Fl}{\mathrm{Fl}}%
\makeatother\let\printedpage\derpp\let\msp\dermsp%
\let\mspin\DERmspin%


% ===================== PDF page 1 =====================
\printedpage{--}
% head: (none)
% no page number printed on this page

\begin{center}
\vspace*{5cm}

{\Large\textsf{THE DERIVATORS}}

\bigskip

\textbf{ALEXANDRE GROTHENDIECK}

\vspace{2cm}

\textbf{Chapter XVII}

\medskip

\textbf{Categories with fibrations and cofibrations}\\
\textbf{(Models (3))}
\end{center}

\par

% ===================== PDF page 2 =====================
\printedpage{--}
% head: (none)
% no page number printed on this page

\vspace*{5cm}

\par\noindent\textbf{Note by the editors of the transcribed edition.}\par
This text was deciphered and transcribed in \LaTeXe{} by M. Künzer. It was edited by M. Künzer, J. Malgoire, and G. Maltsiniotis. The transcription is as faithful as possible to the manuscript. Only a few obvious minor corrections have been made. For the rare additions or comments by the editors, as well as for the original numbering of the manuscript pages, typewriter characters \texttt{[typewriter]} in brackets are used.

This edition is provisional. Remarks, comments and corrections are welcome. Send a message to:

\textbf{[e-mail address removed]}

\hspace*{\fill}G. Maltsiniotis\hspace*{2cm}

\par

% ===================== PDF page 3 =====================
\printedpage{1}
% head: A. Grothendieck : Dérivateurs, Ch. XVII, Édition des Universités de Montpellier II et Paris VII
% foot: Version 9 février 2019

\msp{page 1}

\section*{1\quad Preliminaries\footnote{\textbf{N.B.} These preliminaries will be completely taken up again and superseded in section 2. Not to be read!}}

\textbf{Proposition 1.} \textit{Let $\mathcal{M}$ be a category stable under arbitrary $\varinjlim$, $C \subset \Fl(\mathcal{M})$ containing the isomorphisms, stable under cobase change and under composition} of transfinite length. \textit{Let, for every small ordered set $I$, $C_I \subset \Fl(\mathcal{M}(I^{\circ}))$ be the set of arrows $u_\bullet : X_\bullet \longrightarrow Y_\bullet$ such that $\forall\ i_0 \in I$, in the diagram} % ?? uncertain: superscript of I -- \circ or degree sign

\[
\begin{tikzcd}[column sep=large]
X'_{i_0} = \varinjlim\limits_{i<i_0} X_i \arrow[r] \arrow[d, "\alpha_{i_0} = \varinjlim\limits_{i<i_0} u_i"'] \arrow[dr, phantom, "\text{cocart.}"] & X_{i_0} \arrow[d, "\beta_{i_0}"'] \arrow[dd, dotted, shift left=3ex, "u_{i_0}"] & \\
Y'_{i_0} = \varinjlim\limits_{i<i_0} Y_i \arrow[r] & \tilde{Y}_{i_0} \arrow[d, "\tilde{u}_{i_0}"'] & \text{\footnotesize (defined as an amalgamated sum)} \\
& Y_{i_0}\ , &
\end{tikzcd}
\tag*{$(*)$}
\]

\textit{the arrow $\tilde{u}_{i_0}$ is in $C$\footnote{Condition stable under passage from $I$ to an \textit{open} ordered subset $I'$ of $I$.}. I say that if $I$ is} well ordered, \textit{then $u_\bullet \in C_I$ implies}
\begin{itemize}[leftmargin=2.5em]
\item[a)] \textit{that $\varinjlim u_i : \varinjlim X_i \longrightarrow \varinjlim Y_i$ is in $C$, and}
\item[b)] $\forall\ i \in I$, $(u_i : X_i \longrightarrow Y_i) \in C$\footnote{Say also that we have
\begin{itemize}[leftmargin=2em]
\item[b$'$)] \textit{If $i' \leq i$, then in the diagram similar to $(*)$,}
\end{itemize}
\[
\begin{tikzcd}[ampersand replacement=\&]
X'_{i'} \arrow[r] \arrow[d] \arrow[dr, phantom, "\text{cocart.}"] \& X'_i \arrow[d] \\
Y'_{i'} \arrow[r] \& \tilde{Y}'_i \arrow[d, "u_{i',i}"] \\
\& Y'_i\ ,
\end{tikzcd}
\]
\textit{we have} $u_{i',i} \in C$ (we recover $(*)$ by taking $i' = i_0$, $i = i_0 + 1$). Case b) is obtained by taking $i' = 0$, $i$ arbitrary, whence $(X'_i \longrightarrow Y'_i) \in C$. Applying it to $i+1$ instead of $i$, we find $(u_i : X_i \longrightarrow Y_i) \in C$.}.
\end{itemize}

. % a lone full stop printed on its own line

\textbf{Corollary.} \textit{Let $W \subset \operatorname{Fl}(\mathcal M)$ be a mildly saturated localizer, such that $TC \overset{\mathrm{def}}{=} C \cap W$ is stable under cobase change, and under composition} of transfinite length. \textit{Let $u_\bullet : X_\bullet \longrightarrow Y_\bullet$ be in $C_I$. Then $u_\bullet \in W_{I^{\circ}}$ (i.e. $\forall\ i \in I$, we have $u_i \in W$) if and only if $u_\bullet \in (TC)_I$, i.e. if and only if for every $i \in I$, the arrow $\tilde{u}_i$ in $(*)$ is in $TC$, or what}
\DERmark{ENV0994}\DERmark{ENV0995}

\msp{page 2 does not exist}
\printedpage{2}
% head: A. Grothendieck : Dérivateurs, Ch. XVII, Édition des Universités de Montpellier II et Paris VII
% foot: Version 9 février 2019
\msp{page 3}
\textit{amounts to the same thing, in $W$ (since by hypothesis it is already in $C$). If this condition is satisfied, then $\varinjlim u_i : \varinjlim X_i \longrightarrow \varinjlim Y_i$ is in $W$.}

% ===================== PDF page 4 (continued) =====================

\textsc{Proof} of the corollary. Suppose $u_\bullet \in (TC)_I$, then by the proposition, $\varinjlim\limits_{i\in I} u_i$ is in $TC$, or what amounts to the same thing, in $W$. On the other hand, applying this same result to the direct systems induced on $I_{<i_0}$, we find that in $(*)$, the arrow $\alpha_{i_0}$ is in $TC$, hence also $\beta_{i_0}$ by cobase change, hence also $u_{i_0} = \tilde{u}_{i_0}\beta_{i_0}$. So it remains to prove that conversely, if we assume the $u_{i_0} \in W$, then the $\tilde{u}_{i_0}$ are themselves also in $W$ (i.e. in $TC = C \cap W$), or what amounts to the same thing, that the $\beta_{i_0}$ are in $W$, or again in $TC$. But since $TC$ is stable under cobase change, it suffices for this to prove that $\alpha_{i_0} \in TC$, or again $\alpha_{i_0} \in W$. So we are reduced to proving that
\DERmark{ENV0996}

\begin{itemize}[leftmargin=2.5em]
\item[c)] $u \in W_{I^{\circ}} \cap C_I \Longrightarrow (\,\varinjlim\limits_{i<i_0} u_i) \in W$ for every $i_0 \in I$,
\end{itemize}
so, changing notation, to proving \texttt{[incomplete sentence]}. We will prove a) b) c) together by transfinite induction on the ordinal number $\alpha$ such that $I \simeq I_\alpha = \{\text{set of ordinals } i < \alpha\}$.

\msp{page 4}

1$^\circ$) Case $\alpha = 0$, i.e. $I = \emptyset$: it is trivial.

2$^\circ$) If it's ok for $I_\alpha$, also for $I_{\alpha+1} = \{\beta \mid \beta \leq \alpha\}$. Here, $I_{\alpha+1}$ having a greatest element $\alpha$, a) is contained in b) (for $i = \alpha$, greatest element of $I_{\alpha+1}$.) To prove b), it remains to prove that $u_\alpha \in C$. We look at the diagram $(*)$ for $i_0 = \alpha$. By the induction hypothesis (form a), applied to $I_\alpha$, we have $\alpha_{i_0} \in C$, hence by cobase change $\beta_{i_0} \in C$, hence by composition $u_{i_0} = \tilde{u}_{i_0}\beta_{i_0} \in C$. To prove c), it remains to prove that $\varinjlim\limits_{i<\alpha} u_i \in W$. But by the induction hypothesis applied to $I_\alpha = \{i \mid i < \alpha\}$, we already know that $(u_i)_{i<\alpha} \in (TC)_{I_\alpha}$, hence by the proved part of the corollary or of the proposition, $\varinjlim\limits_{i<\alpha} u_i \in TC$.

3$^\circ$) If $\alpha$ is a limit ordinal, and if a) b) c) hold for $I_\beta$ with $\beta < \alpha$, they also hold for $I_\alpha$. This time, b) is contained in the induction
\msp{page 5}
hypothesis, we must prove a), i.e. that
\[
u_\alpha \overset{\mathrm{def}}{=} \varinjlim_{i<\alpha} u_i \ :\ \underbrace{X_\alpha}_{\overset{\mathrm{def}}{=} \varinjlim\limits_{i<\alpha} X_i} \longrightarrow \underbrace{Y_\alpha}_{\overset{\mathrm{def}}{=} \varinjlim\limits_{i<\alpha} Y_i}
\]
\DERmark{ENV0997}\DERmark{ENV0998}\DERmark{ENV0999}

is in $C$. As for c), it is already contained in the induction hypothesis. So it finally remains to prove
\[
u_\alpha \ \in \ C
\]
when we assume that the proposition is proved for the ordinals $\beta < \alpha$.

We are going to construct a third ordinal system on $I_\alpha$, by introducing for every $0 \leq i < \alpha$, i.e. $i \in I_\alpha$, $Z_i$ by the cocartesian square
\printedpage{3}
% head: A. Grothendieck : Dérivateurs, Ch. XVII, Édition des Universités de Montpellier II et Paris VII
% foot: Version 9 février 2019
\[
\begin{tikzcd}
X_i \arrow[r] \arrow[d] \arrow[dr, phantom, "\text{cocart.}"] & X_\alpha \arrow[d] \\
Y_i \arrow[r] & Z_i\ .
\end{tikzcd}
\]

The transition arrows between the $Z_i$ are obvious, and we will have $\varinjlim\limits_{i} Z_i \xrightarrow{\ \sim\ } Y_\alpha$. Thus, $u_\alpha : X_\alpha (= Z_{-1}) \longrightarrow Y_\alpha$ is decomposed into a transfinite composite of morphisms
\msp{page 6}
\[
Z_i \xrightarrow{v_{i+1}} Z_{i+1} \qquad (-1 \leq i < \alpha)
\tag*{$(*)$}
\]
and
\[
\varinjlim_{i<i_0} Z_i \xrightarrow{v_{i_0}} Z_{i_0} \qquad i_0 \text{ limit ordinal } < \alpha.
\tag*{$(**)$}
\]

(\textbf{N.B.} The case where $i_0$ is \textit{not} a limit ordinal is none other than $(*)$ for $i = i_0 - 1$.) Everything comes down to proving that these arrows are in $C$ (hence their transfinite composite too). % [sic] $\beta_0 - 1$ -> $i_0 - 1$, the case discussed is that of $i_0$
\DERmark{ENV1000}


Case $(*)$.
\[
\begin{tikzcd}
X_i \arrow[r] \arrow[d, "u_i"'] \arrow[dr, phantom, "\text{cocart.}"] & X_{i+1} \arrow[r] \arrow[d, "\beta_{i+1}"'] \arrow[dr, phantom, "\text{cocart.}"] & X_\alpha \arrow[d] \\
Y_i \arrow[r] & \tilde{Y}_{i+1} \arrow[r] \arrow[d, "\tilde{u}_{i+1}"'] \arrow[dr, phantom, "\text{cocart.}"] & Z_i \arrow[d, "v_{i+1}"] \\
& Y_{i+1} \arrow[r] & Z_{i+1}\ .
\end{tikzcd}
\tag*{\footnotemark}
\]
\footnotetext{To give oneself confidence -- one could treat the general case right away.} % [sic] toute de suite -> tout de suite, French idiom

Since $\tilde{u}_{i+1} \in C$, $v_{i+1} \in C$ by cobase change, OK. We must however treat separately the case $i = -1$, i.e. the arrow $X_\alpha \longrightarrow Z_0$, which is included in
\DERmark{ENV1001}

\[
\begin{tikzcd}
X_0 \arrow[r] \arrow[d, "u_0"'] \arrow[dr, phantom, "\text{cocart.}"] & X_\alpha = Z_{-1} \arrow[d, "v_0"] \\
Y_0 \arrow[r] & Z_0\ ,
\end{tikzcd}
\]
where $u_0 \in C$, hence $v_0 \in C$, OK.

Case $(**)$. If $i_0 < \alpha$ is a limit ordinal, let us inspect $(**)$. We have by passage to the limit
\[
Z'_{i_0} \overset{\mathrm{def}}{=} \varinjlim_{i<i_0} Z_i \ \simeq\ Y'_{i_0} \amalg_{X'_{i_0}} X_\alpha
\]
\DERmark{ENV1002}

(cf. diagram $(*)$ for the definition of $X'_{i_0}$, $Y'_{i_0}$), fitting into the following diagram grafted onto $(*)$


\endgroup
% END UNIT der-ch17/der-ch17_p001-005.tex
% BEGIN UNIT der-ch17/der-ch17_p006-010.tex
\begingroup\makeatletter%
\newcommand{\printedpage}[1]{\marginpar{\footnotesize p.~#1}}%
\newcommand{\msp}[1]{\marginpar{\footnotesize [#1]}}%
\newcommand{\defeq}{\overset{\mathrm{def}}{=}}%
\newcommand{\cocart}{\textup{cocart.}}%
\makeatother\let\printedpage\derpp\let\msp\dermsp%
\let\mspin\DERmspin%

\DERmark{ENV1003}


\setcounter{footnote}{4}

\printedpage{4}
% head: A. Grothendieck : Dérivateurs, Ch. XVII, Édition des Universités de Montpellier II et Paris VII
% foot: Version 9 février 2019

\msp{page 7}

\[
\begin{tikzcd}
X'_{i_0} \arrow[r] \arrow[d, "\alpha_{i_0}"'] \arrow[dr, phantom, "\cocart"] & X_{i_0} \arrow[r] \arrow[d, "\beta_{i_0}"'] \arrow[dr, phantom, "\cocart"] & X_{\alpha} \arrow[d] \\
Y'_{i_0} \arrow[r] & \tilde{Y}_{i_0} \arrow[r] \arrow[d, "\tilde{u}_{i_0}"'] \arrow[dr, phantom, "\cocart"] & Z'_{i_0} \arrow[d, "v_{i_0}"] \\
& Y_{i_0} \arrow[r] & Z_{i_0}\ .
\end{tikzcd}
\]

We have $\tilde{u}_{i_0} \in C$ by hypothesis, hence $v_{i_0} \in C$, which completes the proof.

\textbf{Proposition 2.} \textit{$(\mathcal{M}, C)$ and where applicable $W$, as in proposition 1 and its corollary, $I$ a locally finite ordered set, i.e. such that $\forall\ i_0 \in I$, $I/i_0$ is finite. Let}
\[
u_{\bullet}\ : X_{\bullet} \longrightarrow Y_{\bullet}
\]
\textit{an arrow in\footnote{in fact, the hypothesis that $X_{\bullet}$, $Y_{\bullet}$ in $\mathcal{M}(I^{\circ})_{\mathrm{c}}$ is unnecessary, both in proposition 2 and in its two corollaries.} $\mathcal{M}(I^{\circ})_{\mathrm{c}} =$ full subcategory of direct systems on $I$, with values in $\mathcal{M}$, such that $\underbrace{\emptyset_{\mathcal{M}(I^{\circ})}}_{\substack{\textup{initial}\\ \textup{object}\\ \textup{of } \mathcal{M}(I^{\circ})}} \longrightarrow Z_{\bullet}$ is in $C_I$ (defined in prop. 1), i.e. such that $\forall\ i_0 \in I$, the arrow}
\[
Z'_{i_0} \ = \ \varinjlim_{i < i_0} Z_i \ \longrightarrow \ Z_{i_0}
\]
\textit{is in $C$, and suppose $u_{\bullet} \in C_I$. Then}
\begin{enumerate}
\item[a)] $\varinjlim\limits_{i \in I} u_{\bullet} : X_I \defeq \varinjlim\limits_{i \in I} X_i \longrightarrow Y_I \defeq \varinjlim\limits_{i \in I} Y_i$ \textit{is in $C$,}
\msp{page 8}\\
\textit{and}
\item[b)] \textit{$\forall\ i_0 \in I$, we have $\alpha_{i_0}, \beta_{i_0} \in C$ in diagram $(*)$ p. 1.}
\end{enumerate}

\textbf{Corollary.} \textit{Moreover (assuming $W$ given as in prop. 1, cor. 1), the following conditions on a $u_{\bullet}$ between $C$-cofibrant $X_{\bullet}$, $Y_{\bullet}$ are equivalent.}
\begin{enumerate}
\item[$1^{\circ}$)] \textit{$u_{\bullet} \in (TC)_I$, i.e. the $\tilde{u}_{i_0} \in TC$ ($\forall\ i_0 \in I$).}
\item[$2^{\circ}$)] \textit{$u_{\bullet} \in C_I \cap W_{I^{\circ}}$, i.e. $\forall\ i_0 \in I$, $\tilde{u}_{i_0} \in C$, $u_{i_0} \in W$.} % [sic] T_I -> C_I, the explanation that follows gives \tilde u_{i_0} \in C, i.e. the condition C_I
\item[$3^{\circ}$)] \textit{For every $i_0 \in I$, we have}
\[
\alpha_{i_0},\ \beta_{i_0},\ \tilde{u}_{i_0},\ u_{i_0} \ \in \ TC
\]
\DERmark{ENV1004}

\textit{(i.e. $\alpha_{i_0}, \tilde{u}_{i_0} \in TC$, which implies $\beta_{i_0} \in TC$, then $u_{i_0} \in TC$ by cobase change and by composition).}
\end{enumerate}
\textit{And these conditions imply}
\begin{enumerate}
\item[$4^{\circ}$)] \textit{that $u_I : X_I \longrightarrow Y_I$ is in $TC$.}
\end{enumerate}

Let us see how or to what extent the corollary follows from the proposition. We have of course % [sic] sur -> sûr, accent missing (printed "bien sûr" elsewhere)
\printedpage{5}
% head: A. Grothendieck : Dérivateurs, Ch. XVII, Édition des Universités de Montpellier II et Paris VII
% foot: Version 9 février 2019
\[
\begin{tikzcd}[column sep=small]
& 3^{\circ} \arrow[dl, Rightarrow] \arrow[dr, Rightarrow] & \\
1^{\circ} & & 2^{\circ}
\end{tikzcd}
\]
tautologically. Besides, the proposition applied to $TC$ instead of $C$ shows that we have $1^{\circ} \Longrightarrow 3^{\circ}$ \& $4^{\circ}$, for by applying it to the sets $I_{<i_0}$ we find that $\alpha_{i_0} \in TC\ \forall\ i_0 \in I$. So proposition 2 implies % [sic] X -> C, the proposition is applied to TC instead of C
\DERmark{ENV1005}

\[
(1^{\circ} \Longleftrightarrow 3^{\circ}) \Longrightarrow 2^{\circ}\ (\&\ 4^{\circ}).
\]

\msp{page 9}

So to prove $2^{\circ} \Longrightarrow 3^{\circ}$, it would remain to prove that $2^{\circ}$ suffices to imply $4^{\circ}$ (and this for any locally finite ordered set $I$) -- for applying this to the $I_{<i_0}$, we find that $2^{\circ}$ implies $\alpha_{i_0} \in TC = C \cap W$ (not only $\alpha_{i_0} \in C$, according to proposition 2), hence $\beta_{i_0} \in TC \subset W$, hence $\tilde{u}_{i_0} \in W$ by saturation, hence $\tilde{u}_{i_0} \in TC$, whence $3^{\circ}$.

Note on the other hand that b) in proposition 2 is a consequence of a), applied to the restriction of $u_{\bullet}$ to $I_{<i_0}$. All in all, there are therefore only the two things to prove
\[
\left\{
\begin{aligned}
&\text{a)} & u_{\bullet} \in C_I &\quad\Longrightarrow\quad u_I \in C \\
&\text{b)} & u_{\bullet} \in C_I \cap W &\quad\Longrightarrow\quad u_I \in C \cap W\ \ (\defeq TC) \\
& & &\qquad\qquad\ \ (\text{and } u_{\bullet} \in (TC)_I)
\end{aligned}
\right.
\tag*{$(*)$}
\]

We shall prove this by reduction to the special case of proposition 1, already proved.

Besides, I also feel like proving the

\textbf{Corollary 2.} \textit{$\mathcal{M}$, $C$ (and possibly $W$). The hypotheses on $X_{\bullet}, Y_{\bullet} \in \mathcal{M}(I^{\circ})_{\mathrm{c}}$ and on $u_{\bullet} \in C_I$, finally on $I$ locally finite, are those of proposition 2. For every subset $J$ of $I$, let}
\[
u_J \defeq \varinjlim_{i \in J} u_i \ :\ X_J \defeq \varinjlim_{i \in J} X_i \ \longrightarrow\ Y_J \defeq \varinjlim_{i \in J} Y_i\ .
\]
\textit{Let $I' \subset I''$ be two subsets} open \textit{in $I$, and consider the diagram}

\msp{page 10}

\[
\begin{tikzcd}
X_{I'} \arrow[r, "\tau_X"] \arrow[d, "\alpha = u_{I'}"'] \arrow[dr, phantom, "\cocart"] & X_{I''} \arrow[d, "\beta"'] \arrow[dd, bend left=30, "u_{I''}"] & \\
Y_{I'} \arrow[r, "\tilde{\tau}_Y"] \arrow[dr, dotted, "\tau_Y"'] & \tilde{Y}_{I''} \arrow[d, "\tilde{u}"'] & \text{\footnotesize (defined as an amalgamated sum)} \\
& Y_{I''}\ . &
\end{tikzcd}
\tag*{$(*)$}
\]

\printedpage{6}
% head: A. Grothendieck : Dérivateurs, Ch. XVII, Édition des Universités de Montpellier II et Paris VII
% foot: Version 9 février 2019

\textit{Then we have}
\[
\left\{
\begin{aligned}
&\textit{$u_{I'}, u_{I''} \in C$ (hence also $\beta \in C$) (i.e. for every $J$ open in $I$,}\\
&\textit{$u_J : X_J \longrightarrow Y_J$ is in $C$)\,\footnotemark, and $\tau_X, \tau_Y \in C$ (hence also $\tilde{\tau}_Y \in C$)}\\
&\textit{(i.e. for $I' \subset I''$ an inclusion of open subsets, and $X_{\bullet}$ $C$-cofibrant in $\mathcal{M}(I^{\circ})$,}\\
&\textit{we have $X_{I'} \longrightarrow X_{I''}$ in $C$)\,\footnotemark, and finally, even}\\
&\qquad\qquad\qquad\qquad\boxed{\tilde{u} \in C.}
\end{aligned}
\right.
\]
\addtocounter{footnote}{-1}\footnotetext{This is of course contained in prop. 2, applied to $u_{\bullet}|_J$.}\stepcounter{footnote}\footnotetext{It is for this assertion $\tau_X \in C$ (whence $\tilde{\tau}_Y, \tau_Y \in C$, provided $\tilde{u} \in C$) that we need the hypothesis $X_{\bullet} \in \mathcal{M}(I^{\circ})_{\mathrm{c}}$, $\tilde{u} \in C$ true without this hypothesis.}

\textbf{N.B.} The case of proposition 2 in the form $(*)$ a), page 9, (i.e. also the fact that the $u_{I'}$, $u_{I''}$, \dots\ are in $C$) is obtained by taking $I'' = I$, $I' = \emptyset$, the diagram becomes
\[
\begin{tikzcd}
\emptyset \arrow[r] \arrow[d, "\mathrm{id}"'] \arrow[dr, phantom, "\cocart"] & X_I \arrow[d, "\mathrm{id}"] \\
\emptyset \arrow[r] & X_I \arrow[d, "\tilde{u} = u_I"] \\
& Y_I\ ,
\end{tikzcd}
\]
and the conclusion is therefore $u_I \in C$\footnote{The case $(Y_{I'} \xrightarrow{\tau_Y} Y_{I''}) \in C$ is obtained by taking $X_{\bullet} = \emptyset_{\bullet}$, we shall have
\[
\begin{tikzcd}[ampersand replacement=\&]
\emptyset \arrow[r, "\tau_X", "\sim"'] \arrow[d] \arrow[dr, phantom, "\cocart"] \& \emptyset \arrow[d] \\
Y_{I'} \arrow[r, "\mathrm{id}"] \arrow[dr, "\tau_Y"'] \& \tilde{Y}_{I''} = Y_{I'} \arrow[d, "\tilde{u} = \tau_Y"] \\
\& Y_{I''}\ .
\end{tikzcd}
\]}. The \textit{definition} of $C_I$, i.e. the condition $u \in C_I$, is recovered by requiring $\tilde{u} \in C$ in $(*)$ above, in the case $I'' = I/i_0$, $I' = I_{<i_0}$. It is remarkable
\msp{page 11}
that this definition of $C_I$ gives a conclusion as strong as $\tilde{u} \in C$ for \textit{every} pair $I' \subset I''$ of open subsets of $I$.

It is in the end in this form that we are going to prove proposition 2 -- it is \textit{this} now that seems to me \textit{the} right statement in this direction. With the supplement, of course, that if moreover $u_{\bullet} \in W_{I^{\circ}}$, i.e. the $u_i$ are in $W$, then $u_{\bullet}$ is even in $(TC)_I$, or again
\printedpage{7}
% head: A. Grothendieck : Dérivateurs, Ch. XVII, Édition des Universités de Montpellier II et Paris VII
% foot: Version 9 février 2019
simply that the $u_{I'}$ \textit{are in} $TC$ (whence it will therefore follow that $\tilde{u}$ in $(*)$, p. 10, is also in $TC$).

We are going to proceed by transfinite construction. We may of course suppose that $I'' = I$. We are nevertheless going to prove first the special case where $I' = \emptyset$, i.e. prove at least that $u_I \in C$. We construct by transfinite induction a strictly increasing sequence of open subsets
\[
I_{\alpha} \ \subset \ I,
\]
ending at
\[
I_{\alpha_0} \ = \ I,
\]
by setting
\[
\left\{
\begin{aligned}
I_0 \ &= \ \emptyset\,, \\
I_{\alpha} \ &= \ \textstyle\bigcup_{\beta < \alpha} I_{\beta} && \text{if $\alpha$ a limit ordinal}\,, \\
I_{\alpha+1} \ &= \ I_{\alpha} \cup (I/s_{\alpha}) && \text{if $I_{\alpha} \neq I$, and choosing $s_{\alpha} \in I \backslash I_{\alpha}$}\,.
\end{aligned}
\right.
\]

\msp{page 12}

(We end up having $I_{\alpha_0} = I$, of course.) We set
\[
\begin{aligned}
X_{\alpha} \ &= \ X_{I_{\alpha}} \ = \ \varinjlim_{i \in I_{\alpha}} X_i \\
Y_{\alpha} \ &= \ Y_{I_{\alpha}} \ = \ \varinjlim_{i \in I_{\alpha}} Y_i \\
u_{\alpha} \ &= \ u_{I_{\alpha}} \ \ :\ \ X_{\alpha} \longrightarrow Y_{\alpha}\ ,
\end{aligned}
\]
homomorphism of transfinite direct systems. We shall have
\[
X_{\alpha_0} \ = \ X_I, \qquad Y_{\alpha_0} \ = \ Y_I, \qquad u_{\alpha_0} \ = \ u_I,
\]
and we must therefore prove that $u_{\alpha_0} \in C$. For this, we apply proposition 1, so it suffices to prove that $(u_{\alpha}) \in C_{I_{\alpha_0+1}}$ (where $I_{\alpha_0+1} = \{\text{ordinals } i \mid i < \alpha_0 + 1, \text{ i.e. } i \leq \alpha_0\}$). So we must prove that the arrows $\tilde{u}_{i_0}$ in $(*)$ p. 1 (where this time the $i$, $i_0$ are ordinals, $i_0 \leq \alpha_0$) are in $C$. I am going to distinguish the two cases where $\alpha = i_0$ is a limit ordinal, which is trivial for then $X_{i_0} \xleftarrow{\ \sim\ } \varinjlim\limits_{i < i_0} X_i$, $Y_{i_0} \xleftarrow{\ \sim\ } \varinjlim\limits_{i < i_0} Y_i$, and $\tilde{u}_{i_0} = \mathrm{id}$, and the case $\alpha$ not a limit, i.e.
\[
\begin{tikzcd}
X_{I_{\alpha-1}} \arrow[r] \arrow[d, "u_{\alpha-1}"'] & X_{I_{\alpha}} \arrow[d, "\beta_{\alpha}"'] \arrow[dd, bend left=30, "u_{\alpha}"] \\
Y_{I_{\alpha-1}} \arrow[r] & \tilde{Y}_{I_{\alpha}} \arrow[d, "\tilde{u}_{\alpha}"'] \\
& Y_{I_{\alpha}}\ ,
\end{tikzcd}
\]
\printedpage{8}
% head: A. Grothendieck : Dérivateurs, Ch. XVII, Édition des Universités de Montpellier II et Paris VII
% foot: Version 9 février 2019
\msp{page 13}
to prove that $\tilde{u}_{\alpha} \in C$. When we suppose in addition $u_\bullet\in W_{I^{\circ}}$, i.e. the $u_i \in W$, we must in addition prove that $\tilde{u}_{\alpha} \in TC = C \cap W$. Proposition 2 and its corollary 1 will follow from this. To prove $\tilde{u}_{\alpha} \in C$ resp. $\tilde{u}_{\alpha} \in TC$, we may suppose $I_{\alpha} = I$ (changing notation if need be). We therefore set
\DERmark{ENV1006}

\[
\begin{aligned}
I_{\alpha} \ &= \ I, & I_{\alpha-1} \ &= \ I' \\
I/s_{\alpha-1} \ &= \ I_0, & I_0 \cap I' \ &= \ I'_0\ ,
\end{aligned}
\]
so that we have a commutative diagram of open inclusions
\[
\begin{tikzcd}
I_0 \cap I' = I'_0 \arrow[r, hook] \arrow[d, hook] & I_0 \arrow[d, hook] \\
I' \arrow[r, hook] & I = I_0 \cup I'.
\end{tikzcd}
\]
We shall use only the fact that $I_0$, hence $I'_0$, is finite, i.e. that in the pair $(I' \subset I)$, $I$ can be written in the form $I = I' \cup I_0$, with $I_0$ an open subset that is \textit{finite,} to conclude from this that in the diagram
\[
\begin{tikzcd}
X_{I'} \arrow[r, "\tau^X_{I',I}"] \arrow[d, "u_{I'}"'] \arrow[dr, phantom, "\cocart"] & X_I \arrow[d, "v_{I',I}"'] \arrow[dd, bend left=30, "u_I"] \\
Y_{I'} \arrow[r] \arrow[dr, "\tau^Y_{I',I}"'] & \tilde{Y}_I \arrow[d, "\substack{\tilde{u} =\\ u_{I',I}}"'] \\
& Y_I
\end{tikzcd}
\]
we have $\tilde{u} \in C$ resp. $\tilde{u} \in TC = C \cap W$.

\msp{page 14}

We need general notation, on the one hand for an inclusion $J \subset K$ of open subsets of $I$, giving rise to the diagram
\[
\begin{tikzcd}
X_J \arrow[r, "\tau^X_{J,K}"] \arrow[d, "u_J"'] \arrow[dr, phantom, "\cocart"] & X_K \arrow[d, "\tilde{u}_{J,K}"'] \arrow[dd, dotted, bend left=30, "u_K"] & \\
Y_J \arrow[r] \arrow[dr, dotted, "\tau^Y_{J,K}"'] & Y_{J,K} \arrow[d, "u_{J,K}"'] & \text{\footnotesize (formerly $\tilde{Y}_K$)} \\
& Y_K\ , &
\end{tikzcd}
\tag*{(1)}
\]
and for inclusions of three nested open subsets $J \subset K \subset L$ of $I$,


\endgroup
% END UNIT der-ch17/der-ch17_p006-010.tex
% BEGIN UNIT der-ch17/der-ch17_p011-015.tex
\begingroup\makeatletter%
\newcommand{\printedpage}[1]{\marginpar{\footnotesize p.~#1}}%
\newcommand{\msp}[1]{\marginpar{\footnotesize [#1]}}%
\newcommand{\cocart}{\text{\scriptsize cocart.}}%
\makeatletter%
  \renewcommand\@makefnmark{\hbox{\textsuperscript{\normalfont(\@thefnmark)}}}%
  \renewcommand\@makefntext[1]{\parindent 1em\noindent\hbox to 1.8em{\hss\textsuperscript{\normalfont\@thefnmark}}#1}%
\makeatother%
\makeatother\let\printedpage\derpp\let\msp\dermsp%
\let\mspin\DERmspin%


\printedpage{9}
% head: A. Grothendieck : Dérivateurs, Ch. XVII, Édition des Universités de Montpellier II et Paris VII
% foot: Version 9 février 2019
\[
\begin{tikzcd}
X_J \arrow[r] \arrow[d, "u_J"'] \arrow[dr, phantom, "\cocart"] & X_K \arrow[r] \arrow[d, "\tilde{u}_{J,K}"'] \arrow[dr, phantom, "\cocart"] & X_L \arrow[d, "\tilde{u}_{J,L}"'] \arrow[dd, bend left=40, "\tilde{u}_{K,L}"] \\
Y_J \arrow[r] & Y_{J,K} \arrow[r] \arrow[d, "u_{J,K}"'] \arrow[dr, phantom, "\cocart"] & Y_{J,L} \arrow[d, "u^L_{J,K}"'] \arrow[dd, bend left=40, "u_{J,L}"] \\
 & Y_K \arrow[r] & Y_{K,L} \arrow[d, "u_{K,L}"'] \\
 & & Y_L\ ,
\end{tikzcd}
\tag*{(2)}
\]
giving rise to a new arrow $u^L_{J,K}$, deduced from $u_{J,K}$ by the cobase change \msp{page 15} $Y_{J,K} \longrightarrow Y_{J,L}$. We shall use this diagram in the two cases where $(J \subset K \subset L)$ is one of the two triples
\[
\left\{
\begin{aligned}
& I'_0 \hookrightarrow I_0 \hookrightarrow I \\
& I'_0 \hookrightarrow I' \hookrightarrow I,
\end{aligned}
\right.
\]
arising from the cocartesian square (in Cat) (MV square)
\[
\begin{tikzcd}
I'_0 \arrow[r, hook] \arrow[d, hook] & I_0 \arrow[d, hook] \\
I' \arrow[r, hook] & I\ .
\end{tikzcd}
\tag*{(3)}
\]

We must insert the two corresponding diagrams of type $(*)$ \texttt{[on page 10]} into one and the same diagram. We shall note in this connection that the squares in $\mathcal{M}$ deduced from (3) by applying to it $J \longmapsto X_J$ and $J \longmapsto Y_J$
\[
\begin{matrix} (4_X) \\[12ex] (4_Y) \end{matrix}
\qquad\qquad
\begin{tikzcd}[column sep=small, row sep=small]
X_{I'_0} \arrow[rr] \arrow[dr] \arrow[dd, dotted, "u_{I'_0}"] & & X_{I_0} \arrow[dr] \arrow[dd, dotted, "u_{I_0}"] & \\
 & X_{I'} \arrow[rr] \arrow[dd, dotted, "u_{I'}"'] & & X_I \arrow[dd, dotted, "u_I"'] \\
Y_{I'_0} \arrow[rr] \arrow[dr] & & Y_{I_0} \arrow[dr] & \\
 & Y_{I'} \arrow[rr] & & Y_I
\end{tikzcd}
\qquad\qquad
\left.\vrule height 9ex depth 9ex width 0pt\right\}\ (4)
\]
are cocartesian in $\mathcal{M}$. (Transitivity of the $\varinjlim$ \ldots). % [sic] cocartésien -> cocartésiens, agreement with « les carrés »

\msp{page 16}
We have a homomorphism $(u_{I'_0}, u_{I_0}, u_{I'}, u_I)$ deduced from $u$ from the square $(4_X)$ into the square $(4_Y)$, % [sic] les carré -> le carré, singular noun with plural article
and this homomorphism too fits into a diagram, including the two avatars of (2) already mentioned. Here is the relevant diagram: % [sic] signalé -> signalés, agreement with « les deux avatars »
\printedpage{10}
% head: A. Grothendieck : Dérivateurs, Ch. XVII, Édition des Universités de Montpellier II et Paris VII
% foot: Version 9 février 2019
% the footnote mark (9) stands in the printed label of the display; \footnotemark/\footnotetext used because \footnote cannot be set inside a display tag
\[
\begin{tikzcd}[column sep=small, row sep=small]
 & X_{I_0} \arrow[rrr] \arrow[dd] & & & X_I \arrow[dd] \\
X_{I'_0} \arrow[ur] \arrow[rr] \arrow[dd] & & X_{I'} \arrow[urr] \arrow[dd] & & \\
 & Y_{I'_0,I_0} \arrow[rrr] \arrow[dd, "u_{I'_0,I_0}"] & & & Y_{I'_0,I} \arrow[dd, "u^I_{I'_0,I_0}"] \arrow[dddl] \\
Y_{I'_0} \arrow[ur] \arrow[rr] \arrow[dr, dotted] \arrow[dddrr, dotted] & & Y_{I'_0,I'} \arrow[urr] \arrow[ddd] & & \\
 & Y_{I_0} \arrow[rrr] \arrow[dddrrr, dotted] & & & Y_{I_0,I} \arrow[ddd, "u_{I_0,I}"] \\
 & & & Y_{I',I} \arrow[ddr, "u_{I',I}"] & \\
 & & Y_{I'} \arrow[ur] \arrow[drr, dotted] & & \\
 & & & & Y_I\ .
\end{tikzcd}
\tag*{(5~\footnotemark[9])}
\]
\footnotetext[9]{unfathomable diagram, to be replaced by diagram (6) below.}
\DERmark{ENV1007}


The upper cube is the cube in $\mathcal{M}$ with 6 cocartesian faces deduced from the three arrows from $X_{I'_0}$ to $X_{I_0}$, $X_{I'}$, $Y_{I'_0}$. % [sic] faces cocartésiens -> faces cocartésiennes, gender agreement
We shall need to know that the \textit{base} of the cube is a cocartesian square. The front part of the diagram is the avatar of (2) corresponding to the inclusions $I'_0 \subset I' \subset I$, the back part the avatar for the inclusion $I'_0 \subset I_0 \subset I$. So the two other vertical squares, below the upper cube, are also cocartesian, as is the slanting dotted square, which is none other than the base $(4_Y)$ of diagram (4). From these features of diagram (5) it follows easily (although \msp{page 17} I put in a good day of hit-or-miss scribbling, before finally realizing it \ldots) that the diamond-shaped square in solid lines
\[
\begin{tikzcd}[column sep=small, row sep=small]
 & Y_{I'_0,I} \arrow[dl] \arrow[dr] & \\
Y_{I',I} \arrow[dr] & & Y_{I_0,I} \arrow[dl] \\
 & Y_I &
\end{tikzcd}
\]
is also cocartesian. We see it trivially on the unfolded diagram
\printedpage{11}
% head: A. Grothendieck : Dérivateurs, Ch. XVII, Édition des Universités de Montpellier II et Paris VII
% foot: Version 9 février 2019
\[
\begin{tikzcd}[column sep=tiny, row sep=small]
 & & Y_{I'_0} \arrow[dl] \arrow[dr] & & \\
 & Y_{I'_0,I'} \arrow[dl] \arrow[dr] & \cocart & Y_{I'_0,I_0} \arrow[dl] \arrow[dr, "u_{I'_0,I_0}"] & \\
Y_{I'} \arrow[dr] & \cocart & Y_{I'_0,I} \arrow[dl] \arrow[dr, "u^I_{I'_0,I_0}"] & \cocart & Y_{I_0} \arrow[dl] \\
 & Y_{I',I} \arrow[dr, "u_{I',I}"] & \cocart & Y_{I_0,I} \arrow[dl] & \\
 & & {?\simeq Y_I \simeq Y_{I'} \amalg_{Y_{I'_0}} Y_{I_0}\ ,} & &
\end{tikzcd} % ?? uncertain: coproduct sign -- \amalg or \sqcup
\tag*{(6)}
\]
which is finally the relevant diagram (diagram (5) was of no use, it is a wrong end of things that I took hold of). We see on diagram (6) that to prove $u_{I',I} \in C$ resp. $TC$, it \msp{page 18} suffices to prove $u_{I'_0,I_0} \in C$ resp. $TC$. But this is an assertion concerning the inclusion of ordered sets $I'_0 \subset I_0$ that are \textit{finite.} We are thus reduced to proving the following lemma, where $I_0$ becomes $I$:

\textbf{Lemma.} \textit{Let $I$ be a finite ordered set, $u_\bullet : X_\bullet \longrightarrow Y_\bullet$ an arrow in $\mathcal{M}(I^\circ)$, where $\mathcal{M}$ is a category equipped with a set $C \subset \mathrm{Fl}(\mathcal{M})$ of arrows such that $C$ is stable under cobase change and under composition, and contains the isomorphisms. We suppose $u_\bullet \in C_I$, i.e. that for every $i \in I$, the arrow $u_{I_{<i},I/i}$ is in $C$. Then for every open subset $I'$ of $I$, the arrow $u_{I',I}$ is in $C$}~\footnote[10]{\textbf{N.B.} Of course, $u_{J,K} \in C$ if $J \subset K$ are two open subsets of $I$, as we see by applying the lemma to $K$ instead of $I$.}\textit{.}

\textbf{Corollary.} \textit{If moreover $W \ \subset\ \mathrm{Fl}(\mathcal{M})$ is a mildly saturated localizer, such that $TC \overset{\mathrm{def}}{=} W \cap C$ is stable under cobase change (and is so in any case under composition, and contains the isomorphisms), then we have}
\DERmark{ENV1008}

\[
C_I \cap W_{I^\circ} \ = \ (TC)_I,
\]
\textit{thus if $u_\bullet \in C_I$ (cf. above) and the $u_i \in W$, then $u_\bullet \in (TC)_I$, hence applying what is said for $C$ to $TC$, we have $u_{I',I} \in TC$ for every} \msp{page 19} \textit{open subset $I'$ of $I$.}

We shall prove the lemma and its corollary at the same time, proceeding by induction on $\operatorname{card} I = n$. If $n = 0$, i.e. $I = \emptyset$, the statements are trivial. By transitivity, using diagrams of type (2), we are reduced to the case where $I = I' \cup \{i_0\}$, where $i_0 \in I \backslash I'$. We again set $I_0 = I/i_0$, $I'_0 = I_0 \cap I'$, and the reasoning we have just carried out shows us that it suffices to prove that $u_{I'_0,I_0} \in C$ resp. $\in TC$. If $I_0 \neq I$, this follows from the induction hypothesis. % [sic] I'_0 \neq I -> I_0 \neq I, the alternative « Sinon, i_0 est un objet final de I » requires I_0 = I
\DERmark{ENV1009}

Otherwise, $i_0$ is a final object of $I$, and $I' = I_{<i_0}$, and then $u_{I',I} \in C$ is included in the definition of $C_I$. Remains then to prove that $u_{I',I} \in TC$ if the $u_i \in W$. We have the diagram
\printedpage{12}
% head: A. Grothendieck : Dérivateurs, Ch. XVII, Édition des Universités de Montpellier II et Paris VII
% foot: Version 9 février 2019
\[
\begin{tikzcd}
X_{I'} \arrow[r] \arrow[d, "u_{I'}"'] \arrow[dr, phantom, "\cocart"] & X_I \arrow[d, "\tilde{u}_{I',I}"'] \arrow[dd, dotted, bend left=40, "u_I"] \\
Y_{I'} \arrow[r] & Y_{I',I} \arrow[d, "u_{I',I}"'] \\
 & Y_I\ ,
\end{tikzcd}
\]
\msp{page 20} and the induction hypothesis implies $u_{I'} \in TC$, hence $\tilde{u}_{I',I} \in TC$ by cobase change. On the other hand, the arrow $u_I : X_I \longrightarrow Y_I$ identifies with $X_{i_0} \longrightarrow Y_{i_0}$ ($i_0$ being a final element of $I$), so it is in $W$. Thus by saturation, $u_{I',I} \in W$, hence $u_{I',I} \in TC$, q.e.d.

Thus, proposition 2 is proved, as well as its corollary 1 -- and we see that if in proposition 2 and its corollary, we suppose $I$ \textit{finite,} then it is unnecessary in the hypotheses for $C$, resp. $TC = C \cap W$, to suppose stability under composition that is \textit{transfinite,} ordinary composition suffices. But it remains to prove corollary 2, which I have just proved in the case $I$ \textit{finite,} without having to appeal to a ``transfinite'' hypothesis. We must now prove that it is still OK in the locally finite case, supposing $C$ stable under transfinite composition. (There is no longer any need to return to the case of $W$, for we now already know that $C_I \cap W_{I^\circ} = (TC)_I$, by corollary 1 already proved.)

\msp{page 21}
The proof is essentially the same as for the special case already treated $I' = \emptyset$, $I'' = I$, and we could just as well have written up the proof in this setting, with the help of a strictly increasing transfinite sequence of open subsets in $I$, $I_\alpha$, with \fbox{$I_0 = I',$} $I_{\alpha_0} = I''$ $(= I)$, and
\[
\begin{aligned}
I_\alpha \ &= \ I_{\alpha-1} \cup I/s_{\alpha-1} && \text{if $\alpha$ nonlimit, $I_{\alpha-1} \neq I$, choosing $s_{\alpha-1} \in I \backslash I_{\alpha-1}$,} \\
I_\alpha \ &= \ \textstyle\bigcup_{\beta<\alpha} I_\beta && \text{if $\alpha$ limit,}
\end{aligned}
\]
and again setting
\[
X_\alpha \ = \ X_{I_\alpha}, \qquad Y_\alpha \ = \ Y_{I_\alpha}, \qquad u_\alpha = u_{I_\alpha} \ : \ X_\alpha \longrightarrow Y_\alpha.
\]
Let $\alpha_0$ be such that $I_{\alpha_0} = I$ (terminal index of the sequence), we must prove that in
\[
\begin{tikzcd}
X_0 \arrow[r] \arrow[d] & X_{\alpha_0} \arrow[d] \\
Y_0 \arrow[r] & Y_{0,\alpha_0} \arrow[d, "u_{0,\alpha_0}"] \\
 & Y_{\alpha_0}\ ,
\end{tikzcd}
\]
\printedpage{13}
% head: A. Grothendieck : Dérivateurs, Ch. XVII, Édition des Universités de Montpellier II et Paris VII
% foot: Version 9 février 2019
\msp{page 22} we have
\[
u_{0,\alpha_0} \ \in \ C.
\]

Taking into account the marginal supplement to proposition 1 (which is proved in the same way \ldots), we are reduced to proving that the homomorphism of the ordinal systems under consideration is in $C_{I_{\alpha_0+1}}$ (where $I_{\alpha_0+1}$ is the set of ordinals $\leq \alpha_0$). Now, that this is so is proved exactly as earlier, except that we must first make a verification of $(*)$ p.~1 ($\tilde{u}_\bullet \in C$) for $i_0 = 0$. Then diagram $(*)$ p.~1 is
\[
\begin{tikzcd}
\emptyset \arrow[r] \arrow[d, "\mathrm{iso}"] & X_0 = X_{I'} \arrow[d, "\mathrm{iso}"', "\mathrm{id}"] \\
\emptyset \arrow[r] & X_0 = X_{I'} \arrow[d, "u_0 = u_{I'}"] \\
 & Y_0 = Y_{I'}\ ,
\end{tikzcd}
\]
and we must therefore verify $u_0 = u_{I'} \in C$. Now this is contained in proposition 2 itself, already proved, which we must therefore use. % [sic] prouvé -> prouvée, agreement with « la proposition 2 »
The rest is proved as before. Thus, corollary 2 of proposition 2 is also proved.

\textbf{Comments:} There is an uneasiness, for proposition 1 and proposition 2 (as well as their corollaries) have essentially the same conclusion, with hypotheses on $I$ which are \msp{page 23} different: well-ordered in one case, locally finite in another. The natural idea is to give a common statement, which covers both at once, with $I$ artinian, i.e. every decreasing sequence is stationary. What is needed is a hypothesis on $I$ on the one hand, on $u_\bullet : X_\bullet \longrightarrow Y_\bullet$ in $\mathcal{M}(I^\circ)$ on the other hand, which ensures this:

\begin{enumerate}
\item[$(*)$] $\forall\ i_0 \in I$, and every open subset $J'$ of $J = I/i_0$, the arrow $u_{J',J}$ in the diagram
\[
\begin{tikzcd}
X_{J'} \arrow[r] \arrow[d, "u_{J'}"'] \arrow[dr, phantom, "\cocart"] & X_J \arrow[d, "\tilde{u}_{J',J}"'] \arrow[dd, bend left=40, "u_J"] \\
Y_{J'} \arrow[r] & Y_{J',J} \arrow[d, "u_{J',J}"'] \\
 & Y_J
\end{tikzcd}
\]
is in $C$.
\end{enumerate}

\textbf{Question:} is it correct that it suffices, for the validity of $(*)$, to suppose % [sic] exacte -> exact, agreement with « il »
\begin{enumerate}
\item[a)] $u_\bullet \in C_I$, i.e. $u_{J',J} \in C$ when $J = I/i_0$, $J' = J \backslash \{i_0\}$, and
\item[b)] \textit{$I$ locally finite?}
\end{enumerate}


\endgroup
% END UNIT der-ch17/der-ch17_p011-015.tex
% BEGIN UNIT der-ch17/der-ch17_p016-020.tex
\begingroup\makeatletter%
\newcommand{\printedpage}[1]{\marginpar{\footnotesize p.~#1}}%
\newcommand{\msp}[1]{\marginpar{\footnotesize [#1]}}%
\newcommand{\Fl}{\mathrm{Fl}}%
\newcommand{\M}{\mathcal{M}}%
\makeatother\let\printedpage\derpp\let\msp\dermsp%
\let\mspin\DERmspin%


\printedpage{14}
% head: A. Grothendieck : Dérivateurs, Ch. XVII, Édition des Universités de Montpellier II et Paris VII

\msp{page 24}

\textbf{Lemma.} \textit{Let $P \subset I$ be the set of $i_0 \in I$ such that for every open subset $K$ of $J = I/i_0$, the arrow $u_{K,J}$ is in $C$. This said, let $i_0 \in I$ be such that every $i < i_0$ is in $P$. Then $i_0 \in P$.} (We assume $u_\bullet \in C_I$, but $I$ arbitrary.)

Indeed, let $K \subset J = I/i_0$ be an open subset. If $K = J$, it is trivial that $u_{K,J} \in C$ (it is an isomorphism). If $K \neq J$, we have % [sic] Si K = J (second occurrence) -> Si K \neq J, the first case already treats K = J
\DERmark{ENV1010}

\[
K \ \subset\ J' \overset{\mathrm{def}}{=} J\backslash\{i_0\} \ \subset\ J,
\]
\DERmark{ENV1011}

and to prove that $u_{K,J} \in C$, it suffices to prove that $u_{K,J'}, u_{J',J} \in C$. Now $u_{J',J} \in C$ by the hypothesis $u_\bullet \in C_I$. And $u_{K,J'} \in C$ by what has been proved for the case of $J'$, taking into account that by the hypothesis on $i_0$, every $i \in J'$ (i.e. every $i < i_0$) is in $P$. % [sic] par ce que a été prouvé -> par ce qui a été prouvé, grammar

\textbf{Corollary.} \textit{If $I$ is artinian, $P = I$.}

For any subset $P$ of $I$ having the property just stated ($I_{<i_0} \subset P \Longrightarrow i_0 \in P$) is equal to $I$. If indeed there were $x_0 \notin P$, there would exist $x_1 < x_0$ with $x_1 \notin P$, and by induction we would construct a strictly decreasing sequence in $I$, absurd. % [sic] avec x_0 \notin P -> avec x_1 \notin P, the new element must lie outside P
\DERmark{ENV1012}\DERmark{ENV1013}\DERmark{ENV1014}


\msp{page 25}

\section*{2\quad Description of $C_I \subset \Fl(\M(I^\circ))$, for $I$ an ordered set}

Let $\M$ be a category, stable under (small) $\varinjlim$ , $C \subset \Fl(\M)$ (``cofibrations''). We assume
\[
\left\{
\begin{array}{lll}
\text{a)} & C & \text{contains the isomorphisms and is stable}\\
 & & \text{under cobase change.}\\[1ex]
\text{b)} & C & \text{stable under transfinite composition.}
\end{array}
\right.
\tag*{(2.1)}
\]

Let $I$ be a (small) ordered set, and consider
\[
\M(I^\circ) \ =\ \underline{\mathrm{Hom}}(I,\M),
\]
we are going to define a subset
\[
C_I \ \subset\ \Fl(\M(I^\circ))
\]
as follows, for $u_\bullet : X_\bullet \longrightarrow Y_\bullet$ an arrow in $\M(I^\circ)$ :
\[
u_\bullet \in C_I \quad\Longleftrightarrow\quad
\begin{array}{l}
\forall\ i_0 \in I \text{ and for every open subset } J \subset I/i_0 = K,\\
\text{the arrow } u_{J,K} \text{ is in } C.
\end{array}
\tag*{(2.2)}
\]

Here is the definition of $u_{J,K}$. For every open subset $J$ of $I$, we set
\[
\left\{
\begin{array}{rclrcl}
X_J & = & \varinjlim\limits_{i\in J} X_i, & Y_J & = & \varinjlim\limits_{i\in J} Y_i,\\[2ex]
u_J : X_J & \longrightarrow & Y_J, & u_J & = & \varinjlim\limits_{i\in J} u_i,
\end{array}
\right.
\tag*{(2.3)}
\]
and we consider, for $K \subset J$, the diagram
% foot: Version 9 février 2019

\printedpage{15}
% head: A. Grothendieck : Dérivateurs, Ch. XVII, Édition des Universités de Montpellier II et Paris VII

\msp{page 26}

\[
\begin{tikzcd}
X_K \arrow[r, "\tau^X_{K,J}"] \arrow[d, "u_K"'] \arrow[dr, phantom, "\text{\scriptsize cocart.}"] & X_J \arrow[d, "\tilde{u}_{K,J}"'] \arrow[dd, bend left=40, "u_J"] \\
Y_K \arrow[r, "\tau^Y_{K,J}"] & Y_{K,J} \rlap{$\ = Y_K \amalg_{X_K} X_J$} \arrow[d, "u_{K,J}"'] \\
 & Y_J \ .
\end{tikzcd}
\tag*{(2.4)}
\]
\DERmark{ENV1015}


\textbf{Theorem 2.1.} \textit{Let $\M$ be a $\mathfrak{U}$-category stable under small $\varinjlim$ , $C \subset \Fl(\M)$, satisfying the stability conditions} \textup{(2.1)} \textit{(under cobase change and under transfinite composition), $I$ a small ordered set, $u_\bullet : X_\bullet \longrightarrow Y_\bullet$ an arrow in $\M(I^\circ)$.}
\begin{itemize}[leftmargin=3em]
\item[$1^\circ$)] \textit{We assume $u_\bullet \in C_I$ (cf.} \textup{(2.2)}\textit{). Let $K \subset J$ be two open subsets of $I$. Then $u_{K,J}$ (defined in} \textup{(2.4)}\textit{) is in $C$.}
\item[$2^\circ$)] \textit{Suppose $I$ is artinian. In order to have $u_\bullet \in C_I$, i.e. for $u_{K,J}$ to be in $C$ for every pair $(K \subset J)$ of open subsets such that $J$ has a greatest element, i.e. is of the form $I/i_0$ (i.e. also, by} $1^\circ$\textit{), for} \textup{every} \textit{pair of open subsets $K \subset J$), (it is necessary and) it suffices that this be so in the particular case where $J = I/i_0$, $J' = I'/i_0 \overset{\mathrm{def}}{=} I/i_0 - \{i_0\} = \{i \in I \mid i < i_0\}$.}

\msp{page 27}
\item[$3^\circ$)] \textit{Suppose again $I$ is artinian}\footnote[11]{I do not know whether this hypothesis is necessary, cf. comment p. 39.}\textit{. Let $W \subset \Fl(\M)$ be a mildly saturated localizer, such that $TC \overset{\mathrm{def}}{=} C \cap W$ also satisfies the stability conditions} \textup{(2.1)} \textit{(under cobase change, under transfinite composition). The following conditions on an arrow $u_\bullet : X_\bullet \longrightarrow Y_\bullet$ in $\M(I^\circ)$ are equivalent.}
\begin{itemize}[leftmargin=2em]
\item[a)] \textit{$u_\bullet \in (TC)_I$, i.e. for every $i_0 \in I$ and $K$ an open subset of $J = I/i_0$, the arrow $u_{K,J}$ is in $TC$.}
\item[b)] \textit{$u_\bullet \in C_I \cap W_{I^\circ}$, i.e. for every $i_0 \in I$, we have $u_{i_0} \in W$, and for every open subset $K$ of $I/i_0 = J$, we have $u_{K,J} \in C$.}
\end{itemize}
\DERmark{ENV1016}\DERmark{ENV1017}\DERmark{ENV1018}\DERmark{ENV1019}

\end{itemize}

Let us give the main corollaries right away.

\textbf{Corollary 2.2.} \textit{Suppose $(u_\bullet : X_\bullet \longrightarrow Y_\bullet) \in C_I$. Then}
\begin{itemize}[leftmargin=3em]
\item[a)] \textit{$\forall\ i \in I$, $u_i : X_i \longrightarrow Y_i$ is in $C$.}
\item[b)] \textit{The arrow}
\[
u_I \ :\ X_I \longrightarrow Y_I
\tag*{(2.2.1)}
\]
\textit{is in $C$. More generally, for every open subset $J$ of $I$, the arrow}
\[
u_J \ :\ X_J \longrightarrow Y_J
\tag*{(2.2.2)}
\]
\textit{is in $C$.}
\end{itemize}
% foot: Version 9 février 2019

\printedpage{16}
% head: A. Grothendieck : Dérivateurs, Ch. XVII, Édition des Universités de Montpellier II et Paris VII

\msp{page 28}

For a), we apply $1^\circ$) (or simply the \textit{definition} of $C_I$) in the case of the inclusion of open subsets $\emptyset \hookrightarrow I/i$, since we have
\[
u_{\emptyset,I/i} \ \simeq\ u_i \ .
\]
For b), we apply $1^\circ$) to the inclusion $\emptyset \longrightarrow I$, since
\[
u_{\emptyset,I} \ \simeq\ u_I \ ,
\]
and more generally
\[
u_{\emptyset,J} \ \simeq\ u_J \ .
\]

\textbf{Corollary 2.3.} \textit{Suppose $u_\bullet \in C_I \cap W_{I^\circ}$. Then for every inclusion of open subsets $K \subset J$ of $I$, we have $u_{K,J} \in TC$, and in particular (taking $K = \emptyset$), we have $u_J \in TC$. Finally, for every $i \in I$, we have $u_i \in TC$.}

\textbf{Corollary 2.4.} \textit{Let $X_\bullet$ in $\M(I^\circ)_{\mathrm{c}}$ be a $C$-cofibrant direct system, i.e. such that $\emptyset_I \longrightarrow X_\bullet$ is in $C_I$, i.e. such that for every $i_0 \in I$ and every open subset $K$ of $J = I/i_0$, the transition arrow $\tau_{K,J} : X_K = \varinjlim\limits_{i\in K} X_i \longrightarrow X_J \simeq X_{i_0}$ is in $C$. Then for} \textup{every} \textit{pair of open subsets $K \subset J$, the arrow}
\[
\tau_{K,J} \ :\ \varinjlim_{i\in K} X_i \longrightarrow \varinjlim_{i\in J} X_i
\]
\textit{is in $C$}\footnote[12]{and conversely, characterizes the $C$-cofibrant objects.}\textit{. In particular (taking $K = I/k$, $J = I/j$ for $k \leq j$), if $k \leq j$ are in $I$, the transition morphism}
\[
\tau_{k,j} \ :\ X_k \longrightarrow X_j
\]
\msp{page 29}
\textit{is in $C$. Finally, if $W \subset \Fl(\M)$ is as in} \textup{(2.1.\ $3^\circ$)}\textit{, and if $\emptyset_I \longrightarrow X$ is in $W_{I^\circ}$, i.e. the $X_i$ are $W$-coaspheric, then $X_\bullet$ is in $\M(I^\circ)_{TC}$ (where $TC = C \cap W$), and for every open subset $J$ of $I$, $X_J = \varinjlim\limits_{i\in J} X_i$ is in $\M_{TC}$, hence $X_J$ is $W$-coaspheric, and for every pair of open subsets $K \subset J$, the transition arrow $\tau_{K,J}$ is in $TC$.}

Let us point out finally the following complement to theorem 2.1, which will follow from its proof :

\textbf{Corollary 2.5.} \textit{Suppose that the topological space associated with $I$ (whose open subsets are the open subsets of $I$, and the $I/i$ form a basis of open subsets) is} \textup{noetherian}\textit{, i.e. every increasing sequence of open subsets is stationary. Then all the conclusions in theorem} \textup{2.1} \textit{and its corollaries remain valid, assuming only the stability of $C$ or of $TC = C \cap W$ under simple composition (instead of transfinite composition), in addition to the other two stabilities.}
% foot: Version 9 février 2019

\printedpage{17}
% head: A. Grothendieck : Dérivateurs, Ch. XVII, Édition des Universités de Montpellier II et Paris VII

\msp{page 30}

\textsc{Proof} of theorem 2.1 and of corollary 2.5.

2.1. $1^\circ$) We may assume $J = I$, $K \subset I$. We construct a strictly increasing transfinite sequence of open subsets $I_\alpha$, with
\[
\left\{
\setlength{\arraycolsep}{4pt}
\begin{array}{rcll}
I_0 & = & K, & \\[1ex]
I_{\alpha+1} & = & I_\alpha \cup I/s_\alpha & \text{where } s_\alpha \in I\backslash I_\alpha \text{ (if } I_\alpha \neq I\text{, otherwise, the transfinite construction stops),}\\[1ex]
I_\alpha & = & \bigcup_{\beta<\alpha} I_\beta & \text{if } \alpha \text{ a limit ordinal.}
\end{array}
\right.
\]

There exists an $\alpha_0$ such that we have $I_{\alpha_0} = I$. We then see as in paragraph 1
\msp{page 31}
that for every $\alpha$, we have
\[
u_{I_\alpha,I_{\alpha+1}} \ \in\ C,
\tag*{$(*)$}
\]
for it is deduced from $u_{(I/s_\alpha)\cap I_\alpha,I/s_\alpha}$ by cobase change. Cf. diagram (6), p. 17, which ought to be extracted as a lemma, stripped of the technical peculiarities of the situation, i.e. a lemma on homomorphisms of cocartesian squares :
\DERmark{ENV1020}


\textbf{Lemma 2.6.} \textit{Let there be a homomorphism of cocartesian squares $X_\bullet$, $Y_\bullet$, in a category $\M$ with amalgamated sums (\textbf{N.B.} it suffices that the amalgamated sums occurring above exist) :}
\[
\begin{tikzcd}[row sep=small, column sep=small]
X_0 \arrow[rr] \arrow[dr] \arrow[dd, "f_0"'] & & X_1 \arrow[dr] \arrow[dd, "f_1" near end] & \\
 & X'_0 \arrow[rr, crossing over] & & X'_1 \arrow[dd, "f'_1"] \\
Y_0 \arrow[rr] \arrow[dr] & & Y_1 \arrow[dr] & \\
 & Y'_0 \arrow[rr] \arrow[from=uu, crossing over, "f'_0" near start] & & Y'_1 \ .
\end{tikzcd}
\tag*{(2.6.1)}
\]

\textit{Let $Z_0$, $Z_1$ be the two lateral squares and consider, for every commutative square}
$Z = \begin{tikzcd}[ampersand replacement=\&, baseline=(\tikzcdmatrixname.center)]
X \arrow[r] \arrow[d] \& X' \arrow[d] \\
Y \arrow[r] \& Y'
\end{tikzcd}$ \textit{, the arrow $u_Z$ of the diagram}
\[
\begin{tikzcd}
X \arrow[r] \arrow[d] \arrow[dr, phantom, "\text{\scriptsize cocart.}"] & X' \arrow[d] & \\
Y \arrow[r] & \tilde{Z} \arrow[r, "u_Z"'] & Y' \ .
\end{tikzcd}
\]
% foot: Version 9 février 2019

\printedpage{18}
% head: A. Grothendieck : Dérivateurs, Ch. XVII, Édition des Universités de Montpellier II et Paris VII

\msp{page 32}

\textit{Then the arrow $u_{Z_1}$ is deduced from $u_{Z_0}$ by cobase change, more precisely, we have a cocartesian square, where} $T = \begin{tikzcd}[ampersand replacement=\&, baseline=(\tikzcdmatrixname.center)]
X_0 \arrow[r] \arrow[d] \& X_1 \arrow[d] \\
Y_0 \arrow[r] \& Y_1
\end{tikzcd}$ \textit{and} $T' = \begin{tikzcd}[ampersand replacement=\&, baseline=(\tikzcdmatrixname.center)]
X'_0 \arrow[r] \arrow[d] \& X'_1 \arrow[d] \\
Y'_0 \arrow[r] \& Y'_1
\end{tikzcd}$ \textit{denote the back and front faces of the cube} \textup{(2.6.1) :}
\[
\begin{tikzcd}
\tilde{Z}_0 \arrow[r, "u_{Z_0}"] \arrow[d] \arrow[dr, phantom, "\text{\scriptsize cocart.}"] & Y'_0 \arrow[d] \\
\tilde{Z}_1 \arrow[r, "u_{Z_1}"] & Y'_1 \ ,
\end{tikzcd}
\tag*{$(*)$}
\]
\textit{which we see on the base of the diagram}
\[
\begin{tikzcd}[row sep=small, column sep=small]
X_0 \arrow[rr] \arrow[dr] \arrow[dd, "f_0"'] & & X_1 \arrow[dr] \arrow[dd, "f_1" near end] & & \\
 & X'_0 \arrow[rr, crossing over] & & X'_1 \arrow[dd] \arrow[dddr, "f'_1"] & \\
Y_0 \arrow[rr] \arrow[dr] \arrow[dr, phantom, shift left=2ex, "\text{\tiny cocart.}"] & & Y_1 \arrow[dr] & & \\
 & \tilde{Z}'_0 \arrow[rr] \arrow[dr, "u_{Z_0}"'] \arrow[from=uu, crossing over] & & \tilde{Z}'_1 \arrow[dr, "u_{Z_1}"'] & \\
 & & Y'_0 \arrow[rr] \arrow[from=uuul, crossing over, "f'_0" near end] \arrow[ur, phantom, "\text{\tiny cocart.}"] & & Y'_1 \ ,
\end{tikzcd} % ?? uncertain: arrow to which the label f'_0 belongs and exact positions of the two "cocart." labels -- f'_0 printed beside the long diagonal X'_0 -> Y'_0 near Z̃'_0 -> Z̃'_1
\tag*{$(**)$}
\]
\textit{in which the cube is a totally cocartesian cube, deduced from the cocartesian square $X$ by ``exponentiating'' it by the arrow $f_0 : X_0 \longrightarrow Y_0$} \texttt{[rather, by the arrow} $(f_0,f_1)$ \texttt{of} $\underline{\Fl}(\M)$\texttt{]}\textit{.}

The base of the cube, and the square
\msp{page 33}
composed of this one and of $(*)$ (base of the polyhedral diagram $(**)$) being cocartesian, the same holds for $(*)$.
\DERmark{ENV1021}


\textbf{Lemma 2.7.} \textit{Let $(I_\alpha)_{\alpha\leq\alpha_0}$ be an increasing ordinal sequence of open subsets of $I$. Then $u_{I_0,I_{\alpha_0}} : Y_{I_0,I_{\alpha_0}} \longrightarrow Y_{I_{\alpha_0}}$ is a transfinite composite, indexed by the set of ordinals $0 \leq \alpha < \alpha_0$, of a sequence of arrows $v_\alpha$, where each $v_\alpha$ is deduced from $u_{I_\alpha,I_{\alpha+1}}$ by cobase change.}
\DERmark{ENV1022}\DERmark{ENV1023}


To see this, we introduce the transfinite sequence of objects $Z_\alpha = Y_{I_\alpha,I_{\alpha_0}}$, defined by the diagram
% foot: Version 9 février 2019


\endgroup
% END UNIT der-ch17/der-ch17_p016-020.tex
% BEGIN UNIT der-ch17/der-ch17_p021-025.tex
\begingroup\makeatletter%
\newcommand{\printedpage}[1]{\marginpar{\footnotesize p.~#1}}%
\newcommand{\msp}[1]{\marginpar{\footnotesize [#1]}}%
\newcommand{\M}{\mathcal{M}}%
\newcommand{\Mc}{\mathcal{M}_{\mathrm{c}}}%
\makeatother\let\printedpage\derpp\let\msp\dermsp%
\let\mspin\DERmspin%


\printedpage{19}
% head: A. Grothendieck : Dérivateurs, Ch. XVII, Édition des Universités de Montpellier II et Paris VII

\[
\begin{tikzcd}
X_{I_\alpha} \arrow[r] \arrow[d] & X_{I_{\alpha_0}} \arrow[d] \\
Y_{I_\alpha} \arrow[r] & Y_{I_\alpha,I_{\alpha_0}} \rlap{${}= Z_\alpha$} \arrow[d] \\
& Y_{I_{\alpha_0}} \rlap{\ ,}
\end{tikzcd}
\]
we have $Z_0 = Y_{I_0,I_{\alpha_0}}$, $Z_{\alpha_0} = Y_{I_{\alpha_0},I_{\alpha_0}} \simeq Y_{I_{\alpha_0}}$, $Z_\alpha = \varinjlim\limits_{\beta<\alpha} Z_\beta$ if $\alpha$ \msp{34} a limit ordinal. So it remains to determine the $v_{\alpha,\alpha+1} : Z_\alpha \longrightarrow Z_{\alpha+1}$ (for $\alpha < \alpha_0$), which we see from the diagram
\DERmark{ENV1024}\DERmark{ENV1025}\DERmark{ENV1026}\DERmark{ENV1027}

\[
\begin{tikzcd}
X_{I_\alpha} \arrow[r] \arrow[d] & X_{I_{\alpha+1}} \arrow[r] \arrow[d] & X_{I_{\alpha_0}} \arrow[d] \\
Y_{I_\alpha} \arrow[r] & Y_{I_\alpha,I_{\alpha+1}} \arrow[r] \arrow[d, "u_{I_\alpha,I_{\alpha+1}}"] & Y_{I_\alpha,I_{\alpha_0}} \rlap{${}= Z_\alpha$} \arrow[d, "v_{\alpha,\alpha+1}"] \\
& Y_{I_{\alpha+1}} \arrow[r] & Y_{I_{\alpha+1},I_{\alpha_0}} \rlap{${}= Z_{\alpha+1}\ .$}
\end{tikzcd}
\]
\DERmark{ENV1028}


Let us resume the proof of 2.1. We are going to prove 1° by the two preceding lemmas. Let us note that when the space underlying $I$ is noetherian, then the transfinite construction is necessarily finite, so we do not need stability of $C$ under infinite transfinite composites.

\msp{35}

Let us prove 2°. We therefore assume that for every $i_0 \in I$, we have $u_{I_{<i_0},I/i_0} \in C$, we want to deduce from this that $u_\bullet \in C_I$, i.e. $u_{J,I/i_0} \in C$ for \textit{every} open set.
\DERmark{ENV1029}


\noindent\textbf{Lemma 2.8.} (Under the artinian hypothesis on $I$.) \textit{Let $P \subset I$ be the set of $i_0 \in I$ such that $u_\bullet|I/i_0$ is in $C_{I/i_0}$, i.e. that for every $i \leq i_0$, and every open set $K$ of $I/i$, we have $u_{K,I/i} \in C$. Let $i_0 \in I$ be such that $I_{<i_0} \subset P$. Then $i_0 \in P$.}
\DERmark{ENV1030}


\noindent\textsc{Proof}. It remains to prove that if $K \subset I/i_0 = J$ is an open set, then $u_{K,J} \in C$. We may assume $K \neq I/i_0$, i.e. $K \subset I_{<i_0} \overset{\mathrm{def}}{=} J'$. But then by the initial hypothesis on $u_\bullet$, we have $u_{J',J} \in C$, and on the other hand, by the hypothesis on $i_0$, $u_\bullet|J'$ is in $C_{J'}$. So by theorem 2.1. 1°, we have $u_{K,J'} \in C$. So by transitivity $u_{K,J} \in C$, q.e.d.
\DERmark{ENV1031}


\noindent\textbf{Corollary 2.9.} \textit{If $I$ is artinian, then $P = I$, i.e. $u_\bullet \in C_I$.}

This therefore proves 2.1. 2°. Let us note that if $\mathrm{Esp}(I)$ is noetherian, the same is true of every open set (such as $J'$ above), so \msp{36} for the conclusion of the lemma, no need to assume $C$ stable under infinite composition.
\DERmark{ENV1032}


\printedpage{20}
% head: A. Grothendieck : Dérivateurs, Ch. XVII, Édition des Universités de Montpellier II et Paris VII

Let us prove 2.1. 3°. We have a) $\Longrightarrow$ b), for by 2.1. 1°, $u_\bullet \in (TC)_I$ implies that for $i \in I$, $u_i : X_i \longrightarrow Y_i$ is in $TC$ (cf. cor. 2.2 a). So it remains to prove b) $\Longrightarrow$ a), i.e. if $u_\bullet \in C _I\cap W_{I^\circ}$, then $u_\bullet \in (TC)_I$, i.e. the $u_{K,I/i_0}$ are in $TC$. We already know that $u_{K,I/i_0} \in C$ by hypothesis $u_\bullet \in C_I$, so we must prove that it is in $W$. We are reduced to proving that $u_K \in W$, i.e. $u_K \in TC$, for by the diagram (where $J = I/i_0$)
\DERmark{ENV1033}\DERmark{ENV1034}

\[
\begin{tikzcd}
X_K \arrow[r] \arrow[d, "u_K"'] \arrow[dr, phantom, "\text{cocart.}"] & X_J \arrow[d, "\tilde{u}_{K,J}"'] \arrow[dd, bend left=50, "u_J"] \\
Y_K \arrow[r] & Y_{K,J} \arrow[d, "u_{K,J}"'] \\
& Y_J \rlap{\ ,}
\end{tikzcd}
\]
$u_K \in TC$ implies $\tilde{u}_{K,J} \in TC$, and since $u_J \simeq u_{i_0} \in W$ by hypothesis, it indeed follows that $u_{K,J} \in W$ by saturation. But we are going to go about it differently:

\msp{37}

\noindent\textbf{Lemma 2.9.} \textit{$I$ an arbitrary ordered set, $C$, $W$ as in th.} 2.1, 1° \textit{and} 3°, \textit{$u_\bullet : X_\bullet \longrightarrow Y_\bullet$ in $C_I \cap W_{I^\circ}$. Let $P$ be the set of $i_0 \in I$ such that $u_\bullet|I/i_0$ is in $(TC)_{I/i_0}$. If $i_0 \in I$ is such that $I_{<i_0} \subset P$, then $i_0 \in P$.}

Indeed, it remains to prove that for every open set $K$ of $J = I/i_0$, $K \neq J$, we have $u_{K,J} \in TC$. But for this, it suffices to see that $u_K, u_J \in W$ (as we had just seen). Now $u_J \simeq u_{i_0} \in W$ by hypothesis. On the other hand $u_\bullet|K \in (TC_K)$ since $K \subset I_{<i_0} \subset P$, so by theorem 2.1, 1°, we conclude that $u_K \in TC$.

If now we assume $I$ artinian, it follows from 2.9 that $P = I$, so $u_\bullet \in (TC)_I$, q.e.d. % [sic] 3.9 -> 2.9, the lemma just stated is Lemme 2.9
\DERmark{ENV1035}


This completes the proof of theorem 2.1, and also of corollary 2.5, seeing that for the proof of lemma 2.9 as well, the infinite hypotheses on $C$, $C \cap W$ are unnecessary if $\mathrm{Top}(I)$ is noetherian, hence its open sets are also noetherian.

\noindent\textbf{Remarks.}
\begin{enumerate}[leftmargin=2.5em]
\item[1)] In 2.1. 2°, the hypothesis $I$ artinian is not superfluous, as we see by taking $I = \mathbf{N}^{\mathrm{opp}} \simeq \mathbf{Z}^{\leq 0}$, \msp{38} and for an arrow % [sic] n'est par superflue -> n'est pas superflue, negation "pas" intended
\[ \emptyset_{\M(I)} \longrightarrow Y_\bullet\ . \]

To say that this arrow is in $C_I$ means that the transition morphisms $Y_i \xrightarrow{u_i} Y_{i+1}$ are in $C$, \textit{and} that the $Y_i$ are $C$-cofibrant, i.e. $\emptyset_{\M} \longrightarrow Y_i$ is in $C$, for every $i$. But the condition considered in 2.1. 2° is only that the $u_i$ are in $C$. So by taking for $Y_\bullet$ a constant direct system with value $Y$, with $Y$ not $C$-cofibrant, i.e. $\emptyset_{\M} \longrightarrow Y$ not in $C$, we have a counterexample.
\DERmark{ENV1036}


With the same $I$, we see that $u_\bullet : X_\bullet \longrightarrow Y_\bullet$ is in $C_I$ if and only if
\DERmark{ENV1037}

\begin{enumerate}[leftmargin=2.5em]
\item[a)] the $u_i \in C$, and
\item[b)] the $u_{i,i+1} : Y_{i,i+1} \longrightarrow Y_{i+1}$ are in $C$.
\end{enumerate}

\printedpage{21}
% head: A. Grothendieck : Dérivateurs, Ch. XVII, Édition des Universités de Montpellier II et Paris VII

If moreover the $u_i$ are in $W$, hence in $TC$, we see as above that $u_{i,i+1} \in TC$ for every $i$, hence $u_\bullet \in (TC)_I$. So here we still have
\[ (TC)_I \ = \ C_I \cap W_{I^\circ}\ , \]
although $I$ is not artinian. I do not know whether the artinian hypothesis in 2.1. 3° is superfluous or not.
\item[2)] Let $u_\bullet \in \mathrm{Fl}(\M(I^\circ))$\footnote[13]{$I^\circ = I^{\mathrm{opp}}$.}, so $u_\bullet$ is an object of $\underline{\mathrm{Fl}}(\M(I^\circ)) \simeq \M(I^\circ \times \boldsymbol{\Delta}^1) = \M((I \times (\boldsymbol{\Delta}^1)^{\mathrm{opp}})^{\mathrm{opp}})$. We check that for this object to be $C$-cofibrant, it is necessary and sufficient that
\begin{enumerate}[leftmargin=2.5em]
\item[1°)] $X_\bullet$ be $C$-cofibrant, and
\item[2°)] $u_\bullet \in C_I$ (which implies that $Y_\bullet$ is also cofibrant).
\end{enumerate}
\end{enumerate}
\msp{39}

\section*{3\quad Categories of models with cofibrations}

Let $\M$ be a category equipped with two sets of arrows
\[ W,\ C\ \subset\ \mathrm{Fl}(\M)\ . \]
We assume that $\M$ contains an initial object, and:
\begin{description}[leftmargin=3em, labelwidth=2.5em]
\item[Cof$_0$] $W$ is mildly saturated, i.e. contains the isomorphisms and if $f$, $g$ are composable, and if two among $f$, $g$, $gf$ are in $W$, the third too.
\item[Cof$_1$] $C$ contains the isomorphisms, is stable under composition and under cobase change.
\item[Cof$_2$] $TC \overset{\mathrm{def}}{=} C \cap W$ is stable under cobase change.
\item[Cof$_3$] Every arrow $f$ in $\M$ is written
\[ f \ = \ pi \]
\DERmark{ENV1038}

with $i \in C$, $p \in W$.
\end{description}
These are the axioms of K. S. \textsc{Brown} (dual form), who moreover assumes that the objects of $\M$ are $C$-cofibrant, i.e. $\emptyset_{\M} \longrightarrow X$ is in $C$. I will not make this assumption. The subcategory $\Mc$ of $\M$ formed of the \msp{40} $C$-cofibrant objects satisfies the axioms of \textsc{Brown}, moreover one proves that
\[ \mathrm{Ho}_W(\Mc) \xrightarrow{\ \sim\ } \mathrm{Ho}_W(\M) \]
is an equivalence. Thus the results of \textsc{Brown} concerning $\Mc$, equipped with the sets of arrows induced by $W$, $C$, are applicable to $\mathrm{Ho}_W(\M)$ and to its avatars (such as $\mathrm{Ho}_W(S\backslash\M)$, $\mathrm{Ho}_W(\M/T)$ and $\mathrm{Ho}_W(S\backslash\underset{u_0}{\M}/T)$, \ where $u_0 : S \longrightarrow T$ is given). It turns out that for the most part, the standard homotopy results (in particular exact sequence of cofibrations relating to the homotopy operations ``on the right'' (or should one say on the left? yet these are questions of right exactness, not left exactness)) are valid. I will have to return to this in greater detail later on. In particular, there is a good theory of $W$-cocartesian squares (and \textit{in} $\Mc$\textit{,} there are enough $W$-cofibrations\footnote[14]{More precisely, by K. S. \textsc{Brown}, \textit{in} $\Mc$\textit{,} the objects of $C$ are $W_{\mathrm{c}}$-cofibrations.}, \ldots),
\printedpage{22}
% head: A. Grothendieck : Dérivateurs, Ch. XVII, Édition des Universités de Montpellier II et Paris VII
\msp{41}
and moreover, these are \textit{exact.} The category $\mathrm{Ho}_W(\M) \simeq \mathrm{Ho}_W(\Mc)$ ``is computed'' by a calculus of left fractions in $\pi\Mc$, where $\pi\Mc$ is the category obtained from $\Mc$ by replacing the $\mathrm{Hom}_{\Mc}(X,Y) = \mathrm{Hom}_{\M}(X,Y)$ in it by the quotient by a suitable homotopy relation:
\begin{description}[leftmargin=9em, labelwidth=8em, font=\normalfont]
\item[$f,g : X \rightrightarrows Y$] are QB-homotopic (in the sense of \textit{Quillen-Brown}) if there exists an $X' \xrightarrow{\alpha} X$ in $W$ (and we may assume $X' \in \mathrm{Ob}\,\Mc$ as well) such that $f\alpha \overset{1}{\sim} g\alpha$ in the sense of \textsc{Quillen}, i.e. there exists a cylinder object
\begin{tikzcd}[baseline=(current bounding box.center)]
I_{X'} \arrow[r, "\in W"] & X' \\
X' \vee X' \arrow[u, "\in C"] \arrow[ur, dotted] &
\end{tikzcd}
for $X'$, such that $(f\alpha, g\alpha)\ :\ X' \vee X' \longrightarrow Y$ factors through an $h' : I_{X'} \longrightarrow Y$.
\end{description}

It is from this criterion that the (left) exactness of $W$-cocartesian squares follows.

I want to see how a cofibration structure on $\M$ extends to the categories $\M(I^\circ) = \underline{\mathrm{Hom}}(I, \M)$. We have in $\M(I^\circ)$ \msp{42} the set $W_{I^\circ}$, and $C_I$ (it is probably better to denote it also by $C_{I^\circ}$) defined in the preceding paragraph. It is a matter of checking the axioms. The axiom $\mathrm{Cof}_0$ is immediate, and well known for $I$ an arbitrary category. % [sic] \underline{Hom}(I, M) -> \underline{Hom}(I, \mathcal{M}), italic M printed where the script \mathcal{M} is used everywhere else
\DERmark{ENV1039}


$\mathrm{Cof}_1$ : $C_{I^\circ}$ contains the isomorphisms, that is clear, because $C$ contains them. Stable under composition $X_\bullet \xrightarrow{u} Y_\bullet \xrightarrow{v} Z_\bullet$, let us look at $K \subset J$, put $w = v \circ u$,
\[
\begin{tikzcd}[execute at end picture={\draw[decorate,decoration={brace,mirror,amplitude=6pt}] ([yshift=-2pt]B.south west) -- ([yshift=-2pt]E.south east) node[midway,below=7pt]{$\scriptstyle w_{K,J}$};}]
X_K \arrow[r] \arrow[d] & X_J \arrow[d] & & \\
Y_K \arrow[r] \arrow[d] & Y_{K,J} \arrow[r, "u_{K,J}"] \arrow[d] & Y_J \arrow[d] & \\
Z_K \arrow[r] & |[alias=B]| \bullet \arrow[r, "u'_{K,J}"] & Z_{K,J} \arrow[r, "v_{K,J}"] & |[alias=E]| Z_J \rlap{\ ,}
\end{tikzcd}
\]
\vspace{1ex}

$w_{K,J}$ is the composite of $v_{K,J}$, which is in $C$, and $u'_{K,J}$, which is too since obtained from $u_{K,J}$ by cobase change.

(\textbf{N.B.} We have already made the argument with a single $u_\bullet$, and \textit{three} open sets $J \subset K \subset L$. It \msp{43} would be appropriate to introduce the notion of a $C$-square
\begin{tikzcd}[baseline=(current bounding box.center)]
X_0 \arrow[r] \arrow[d] & X_1 \arrow[d] \\
X_2 \arrow[r] & X_3
\end{tikzcd}, as a square such that
\printedpage{23}
% head: A. Grothendieck : Dérivateurs, Ch. XVII, Édition des Universités de Montpellier II et Paris VII
$X_1 \amalg_{X_0} X_2 \longrightarrow X_3$ is in $C$ (\textbf{N.B.} notion symmetric for the symmetry functors of squares), and say that the composite of two $C$-squares is a $C$-square.)

Cobase change
\[
\begin{tikzcd}
X_\bullet \arrow[r, "u"] \arrow[d] & Y_\bullet \arrow[d] \\
X'_\bullet \arrow[r, "u'"] & Y'_\bullet \rlap{\ .}
\end{tikzcd}
\]

Let us consider $K \subset J$, we have
\[
\begin{tikzcd}[row sep=large, column sep=large]
& X_K \arrow[r] \arrow[d] \arrow[ddl] \arrow[dr, phantom, "\text{\scriptsize cocart.}"] & X_J \arrow[d] \arrow[dr] & \\
& Y_K \arrow[r] \arrow[ddl] \arrow[dl, phantom, "\text{\scriptsize cocart.}"] & Y_{K,J} \arrow[r, "u_{K,J}"] \arrow[ddl] \arrow[d, phantom, "\text{\scriptsize cocartesian}"] & Y_J \arrow[ddl] \\
X'_K \arrow[r, crossing over] \arrow[d] \arrow[dr, phantom, "\text{\scriptsize cocart.}"] & X'_J \arrow[d] \arrow[dr, crossing over] \arrow[from=1-3, crossing over] & {} & \\
Y'_K \arrow[r] & Y'_{K,J} \arrow[r, "u'_{K,J}"] & Y'_J \rlap{\ ,} &
\end{tikzcd}
\]
where the three lateral squares marked ``cocart.'' are cocartesian, as is the square in an inclined plane. It follows, first that the fourth lateral square of the cube is itself also cocartesian, \msp{44} then that the square having two sides $u_{K,J}$, $u'_{K,J}$ is cocartesian. (\textbf{N.B.} compare the similar diagram p. 32 -- there would be a common lemma to extract, distinguishing two cases \ldots)

$\mathrm{Cof}_2$. To prove that $C_{I^\circ} \cap W_{I^\circ}$ is stable under cobase change, I see no way of doing it, if I do not already know that we have
\[ C_{I^\circ} \cap W_{I^\circ} \ = \ (TC)_{I^\circ}\ , \tag*{$(*)$} \]
and then we apply the preceding result (where $W$ did not enter) to $TC$ instead of $C$. But to have $(*)$, I must assume that we are in one of the following two cases (so as to be able to apply theorem 2.1 and its corollary 2.5).
\begin{enumerate}[leftmargin=2.5em]
\item[a)] $I$ finite, or more generally $I$ artinian \textit{and} $\mathrm{Top}(I)$ is noetherian.
\item[b)] $I$ artinian, and moreover $C$ and $TC$ are stable under transfinite composition.
\end{enumerate}
$\mathrm{Cof}_3$. The factorization axiom seems to me a priori the most delicate\footnote[15]{The argument that follows applies to other factorization situations, for example for a given Quillen precouple $(\Phi, \Psi)$ in $\M$.}. But in the end this comes out without too much trouble. Let there be given % [sic] facorisation -> factorisation, missing letter
\[ f \ : \ X_\bullet \longrightarrow Y_\bullet \]


\endgroup
% END UNIT der-ch17/der-ch17_p021-025.tex
% BEGIN UNIT der-ch17/der-ch17_p026-030.tex
\begingroup\makeatletter%
\newcommand{\printedpage}[1]{\marginpar{\footnotesize p.~#1}}%
\newcommand{\msp}[1]{\marginpar{\footnotesize [#1]}}%
\makeatother\let\printedpage\derpp\let\msp\dermsp%
\let\mspin\DERmspin%


\printedpage{24}
% head: A. Grothendieck : Dérivateurs, Ch. XVII, Édition des Universités de Montpellier II et Paris VII

\msp{45}in $\mathcal{M}(I^\circ)$, let us consider the quadruples
\[
(J,\ Z,\ i,\ p)\ ,
\]
\DERmark{ENV1040}

where $J\subset I$ is an open set, $Z$ an object of $\mathcal{M}(J^\circ)$, and $i$, $p$ are arrows
\[
\left\{
\begin{aligned}
& X|J \xrightarrow{\ i\ } Z \xrightarrow{\ p\ } Y|J \\
& \text{with } i\in C_{J^\circ},\ p\in W_{J^\circ}.
\end{aligned}
\right.
\]
\DERmark{ENV1041}

We order the quadruples by extension. If we have a totally ordered set, of $(J_\alpha,Z_\alpha,i_\alpha,p_\alpha)$, putting $J=\bigcup J_\alpha$, we find unique $Z$, $i$, $p$ on $J$ which extend % [sic] \bigcup I_\alpha -> \bigcup J_\alpha, the union is of the open sets J_\alpha
\DERmark{ENV1042}

the given $Z_\alpha$, $i_\alpha$, $p_\alpha$, moreover we will have $i\in C_{J^\circ}$, $p\in W_{J^\circ}$, since membership in $C_{J^\circ}$ resp. $W_{J^\circ}$ is a local question on $J$. So by Zorn, we are reduced to proving that if $(J,Z,i,p)$ is maximal, we have $J=I$ – whence the desired factorization $f=pi$. So to prove that if % [sic] d'ou -> d'où, missing accent
$J\neq I$, the quadruple is not maximal. Let us consider the nonempty closed set $I\backslash J$. $I$ is assumed artinian, there is a \textit{minimal element} $i_0$ of $I\backslash J$. Let us put
\msp{46}
\[
J' \;=\; J\cup\{i_0\} \;=\; J\cup I/i_0,
\]
and let us show that we can extend the triple $Z$, $i$, $p$ to $Z'$, $i'$, $p'$ on $J'$. Putting $K'=I/i_0$, % [sic] j' -> i', the triplet is Z, i, p
\DERmark{ENV1043}

$K=J\cap K'=I_{<i_0}$, we thus have % [sic] J\cap K -> J\cap K', as in the diagram below
\DERmark{ENV1044}

\[
\begin{tikzcd}
K=J\cap K' \arrow[r, hook] \arrow[d, hook] & K'=I/i_0 \arrow[d, hook] \\
J \arrow[r, hook] & J'=J\cup K'.
\end{tikzcd}
\]

We then see that the extensions of $(Z,i,p)$ from $J$ to $J'$ correspond to those of $(Z|K,i|K,p|K)$ to $K'$. For this, we note that for an arrow $\varphi$ between objects of $\mathcal{M}(J'^\circ)$, this arrow is in $C_{J'^\circ}$ resp. $W_{J'^\circ}$, if and only if its restrictions to $J$ and to $K'$ are.
\DERmark{ENV1045}


We are thus reduced to the standard case of an ordered set with final object $e_0$, an ordered set which we will denote by $I$, and the ordered set $I'$ of objects $i<e_0$, i.e. $I'=I\backslash\{e_0\}$. % [sic] I\backslash\{i_0\} -> I\backslash\{e_0\}, I' consists of the i < e_0
\DERmark{ENV1046}

We have $f:X_\bullet\longrightarrow Y_\bullet$, an arrow in $\mathcal{M}(I^\circ)$, a factorization of $f|I'$ as $p'\circ i'$ with $i'\in C_{I'^\circ}$, $p'\in W_{I'^\circ}$, % [sic] W_{I^\circ} -> W_{I'^\circ}, factorisation of f|I' on I'
\DERmark{ENV1047}

\msp{47}
and it is a matter of extending it to a similar factorization of $f$. We note that giving an object $X_\bullet$ of $\mathcal{M}(I^\circ)$ is equivalent to giving an object $X'_\bullet$ of $\mathcal{M}(I'^\circ)$ (the restriction of $X$ to $I'$) and an arrow
\[
X'_{I'} \xrightarrow{\ \alpha_X\ } X_{e_0}=X_0
\]
in $\mathcal{M}$, where $X'_{I'}=\varinjlim\limits_{i'\in I'}X'_{i'}$. Thus giving the arrow $f:X_\bullet\longrightarrow Y_\bullet$ amounts to giving $f'=f|I':X'_\bullet\longrightarrow Y'_\bullet$, and an arrow $f_0:X_0\longrightarrow Y_0$ making the square commutative
\printedpage{25}
% head: A. Grothendieck : Dérivateurs, Ch. XVII, Édition des Universités de Montpellier II et Paris VII
\[
\begin{tikzcd}
X'_{I'} \arrow[r, "\alpha_X"] \arrow[d, "f'_{I'}"'] & X_0 \arrow[d, "f_0"] \\
Y'_{I'} \arrow[r, "\alpha_Y"] & Y_0\ .
\end{tikzcd}
\]

Likewise, finding a factorization of $f$ as $pi$, which extends a given factorization $f'=p'i'$ on $I'$, amounts to giving a commutative diagram
\DERmark{ENV1048}

\[
\begin{tikzcd}[column sep=small, row sep=small]
& X'_{I'} \arrow[rr, "\alpha_X"] \arrow[dl, "i'_{I'}"'] \arrow[dd, "f'_{I'}"] & & X_0 \arrow[dl, "i_0"'] \arrow[dd, "f_0"] \\
Z_{I'} \arrow[rr, "\alpha_Z"'] \arrow[dr, "p'_{I'}"'] & & Z_0 \arrow[dr, "p_0"'] & \\
& Y'_{I'} \arrow[rr, "\alpha_Y"'] & & Y_0\ ,
\end{tikzcd}
\tag*{$(*)$}
\]
\msp{48 blank}
\msp{49}
where the object $Z_0$, and the three arrows $\alpha_Z$, $i_0$, $p_0$ are to be determined. If we have $i'\in C_{I'^\circ}$, $p'\in W_{I'^\circ}$, % [sic] p -> p', the given factorisation is f'=p'i'
\DERmark{ENV1049}

then we conclude by theorem 2.1. $1^\circ$ (cor. 2.2) that
\[
i'_{I'}\ \in\ C\ ,
\]
and if we want $i\in C_{I^\circ}$, $p\in W_{I^\circ}$, it only remains to require
\[
p_0\in W, \qquad (i'_{I',I}:\underbrace{Z_{I'}\amalg_{X_{I'}}X_0}_{Z'_0}\longrightarrow Z_0)\in C,
\]
by theorem 2.1. $2^\circ$. We will then construct the diagram $(*)$, from the similar diagram below
\[
\begin{tikzcd}
& X_{I'} \arrow[rr, "\alpha_X"] \arrow[dl, "i'_{I'}\in C"'] \arrow[dd, "\text{cocart.}" {pos=0.25}, "f'_{I'}" {pos=0.75}] & & X_0 \arrow[dl, "i'_0\in C"'] \arrow[dd, "f_0" {pos=0.25}] \arrow[dr, dotted, "i_0=j_0i'_0"] & \\
Z_{I'} \arrow[rr, "\alpha'_Z"] \arrow[dr, "p'_{I'}"'] & & Z'_0=Z_{I',I} \arrow[rr, dotted, "j_0\in C"] \arrow[dr, "p'_0"'] & & Z_0 \arrow[dl, dotted, "p_0\in W"] \\
& Y_{I'} \arrow[rr, "\alpha_Y"] & & Y_0 &
\end{tikzcd}
\]
(with $Z_0$ of $(*)$ replaced by $Z'_0=Z_{I',I}$) by factoring the canonical arrow $p'_0:Z'_0\longrightarrow Y_0$ as
\[
p'_0 \;=\; p_0j_0\ , \qquad \text{with } p_0\in W,\ j_0\in C\ ,
\]
and putting
\[
i_0 \;=\; j_0i'_0\ ,
\]
which completes the construction.

We have thus proved the

\printedpage{26}
% head: A. Grothendieck : Dérivateurs, Ch. XVII, Édition des Universités de Montpellier II et Paris VII

\msp{50}\textbf{Theorem 3.1.} \textit{Let $(\mathcal{M},W,C)$ be a category with cofibrations (cf. page \textup{39}), $\mathcal{M}$ stable under finite $\varinjlim$. For any ordered set $I$, let us consider the corresponding tuple $(\mathcal{M}(I^\circ),W_{I^\circ},C_{I^\circ})$, where $C_{I^\circ}$ is defined in paragraph \textup{2} (cf. on this point theorem \textup{2.1}). Suppose that $I$ is an \textup{artinian} ordered set (e.g. $I$ locally finite, i.e. the $I/i$ finite). Suppose moreover that we are in one of the following two cases:} % [sic] odonné -> ordonné, missing letter
\begin{itemize}[leftmargin=2.5em]
\item[a)] \textit{The space \textup{Esp}$(I)$ is noetherian, i.e. every increasing sequence of open sets of $I$ is stationary (e.g. $I$ finite).}
\item[b)] \textit{$C$ and $TC=C\cap W$ are stable, not only under ordinary composition, but also under transfinite composition. In this case, we also assume $\mathcal{M}$ stable under small filtering $\varinjlim$ (hence under arbitrary small $\varinjlim$).}
\end{itemize}
\textit{Under these conditions, $(\mathcal{M}(I^\circ),W_{I^\circ},C_{I^\circ})$ is also a category with cofibrations.}

\textbf{Proposition 3.2.} \textit{Let $u:I\longrightarrow I'$ be an increasing map between ordered sets, whence a pair of adjoint functors}
\[
\begin{tikzcd}
\mathcal{M}(I^\circ) \arrow[r, shift left, "u_{!}"] & \mathcal{M}(I'^\circ) \arrow[l, shift left, "u^*"]
\end{tikzcd}\ .
\tag*{(3.2.1)}
\]
\msp{51}
\begin{itemize}[leftmargin=2.5em]
\item[a)] \textit{The functor $u^*$ is compatible with the localizers $W_{I'^\circ}$, $W_{I^\circ}$. If $I$ is discrete, or $u$ étale, then $u^*$ maps $C_{I'^\circ}$ into $C_{I^\circ}$. The same is still true when $u$ is locally cofiltering, i.e. if for every $i'\in\operatorname{Ob}I'$, the category $i'\backslash I$ is locally cofiltering, i.e. its connected components are cofiltering (condition $\mathrm{PS}_{\mathrm{g}}$ \textup{1} on $i'\backslash I$)} \textsuperscript{(}\footnote[16]{no need to speak of ``weakly'' locally cofiltering, since these are ordered sets where there is no need to distinguish between the strong and weak notions.}\textsuperscript{)}.
\item[b)] \textit{In all cases, $u_!$ transforms $C_{I^\circ}$ into $C_{I'^\circ}$, $TC_{I^\circ}$ into $TC_{I'^\circ}$, finally $\mathcal{M}(I^\circ)_{\mathrm{c}}$ into $\mathcal{M}(I'^\circ)_{\mathrm{c}}$.}
\end{itemize}
\DERmark{ENV1050}


\textsc{Proof}. Let us prove a), in the most awkward case, where we assume only that $u$ is locally weakly cofiltering. We apply II 7.4 (p. 92), which tells us that $\forall i\in\operatorname{Ob}I$,
\msp{52}
considering the morphism induced by $u$
\[
I/i \longrightarrow I'/i' \qquad (\text{where } i'=u(i)),
\]
this is \textit{(globally)} weakly cofiltering, and that for every open set $U$ of $I/i$, denoting by $U'$ the open set generated by its image, the induced morphism $v:U\longrightarrow U'$ is still weakly cofiltering, and a fortiori \textit{cofinal.} Let then $f':X'\longrightarrow Y'$ be an arrow in $C_{I'^\circ}$, let us prove that
\[
\underbrace{u^*(f')}_{f}\ :\ \underbrace{u^*(X')}_{X}\ \longrightarrow\ \underbrace{u^*(Y')}_{Y}
\]
is in $C_{I^\circ}$, i.e. that for every $i\in I$, and $U$ an open set of $I/i$, in the canonical diagram
\printedpage{27}
% head: A. Grothendieck : Dérivateurs, Ch. XVII, Édition des Universités de Montpellier II et Paris VII
\[
\begin{tikzcd}
X_U \arrow[r] \arrow[d] \arrow[dr, phantom, "\text{cocart.}"] & X(i) \arrow[d] \\
Y_U \arrow[r] & Y_{U,i} \arrow[d, "v_{U,i}"] \\
& Y(i)\ ,
\end{tikzcd}
\tag*{$(*)$}
\]
\msp{53}
the arrow $v_{U,i}$ is in $C$. The diagram $(*)$ maps into the similar diagram for $f':X'\longrightarrow Y'$, $i'=u(i)$, and the open set $U'\subset I'/i'$ generated by $u(U)$:
\[
\begin{tikzcd}[column sep=small, row sep=small]
X_U \arrow[rr] \arrow[dr, "f_U"] \arrow[dd, "c_U^{X'}"'] & & X(i) \arrow[dr] \arrow[dd, "\text{cocart.}"' {pos=0.2}, "\mathrm{id}"' {pos=0.75}] & & \\
& Y_U \arrow[rr] \arrow[dd, "c_U^{Y'}"' {pos=0.3}] & & Y_{U,i} \arrow[dr, "v_{U,i}"] \arrow[dd, "c^{f'}_{U,i}"] & \\
X'_{U'} \arrow[rr, "\text{cocart.}"' {pos=0.6}] \arrow[dr, "f'_{U'}"'] & & X'(i') \arrow[dr] & & Y(i) \arrow[dd, "\mathrm{id}"] \\
& Y'_{U'} \arrow[rr] & & Y'_{U',i'} \arrow[dr, "v_{U',i'}"'] & \\
& & & & Y'(i')\ .
\end{tikzcd}
\tag*{$(**)$}
\]

By cofinality, the two vertical arrows $c_U^{X'}$ and $c_U^{Y'}$ are isomorphisms, in addition to $X(i)\xrightarrow{\ id\ }X'(i')$, so the fourth arrow of the cube, $c^{f'}_{U,i}$, is an isomorphism. Since $v_{U',i'}$ is in $C$, $v_{U,i}$ too, we win.

Let us prove b). So let
\[
f\ :\ X\longrightarrow Y
\]
be an arrow in $C_I$, and let
\[
f' \;=\; u_!(f)\ :\ u_!(X)=X' \longrightarrow u_!(Y)=Y',
\]
let us prove that it is in $C_{I'^\circ}$, hence that
\msp{54}
for two open sets $U'\subset V'$ of $I'$ (where we can assume that $V'=I'/i'$), in the base-diagram of $(**)$, $v_{U',V'}$ is in $C$. Let $U=u^{-1}(U')$, $V=u^{-1}(V')$, and let us consider the % ?? uncertain: hyphen at line end in "diagramme-base" -- kept as a compound hyphen, or "diagrammebase"
diagram similar to $(**)$, defined by the two inclusions of open sets $U'\subset V'$, $U\subset V$. This time, since the formation of $u_!$ commutes with restrictions to open sets, by transitivity of $\varinjlim$ , the three vertical arrows $c^X_{U'}$, $c^Y_{U'}$, $c^X_{V'}$ are isomorphisms, hence so is $c^f_{U',V'}$. In the same % [sic] vertivales -> verticales, misprint
way, the arrow $Y_V\longrightarrow Y'_{V'}$ (in place of $Y(i)\longrightarrow Y'(i')$) is an isomorphism. Since $v_{U,V}\in C$, the same is true of $v_{U',V'}$, which proves that $f'\in C_{I'^\circ}$. We proceed in the same way for $TC_I$. Finally, the assertion that $u_!$ maps $\mathcal{M}(I^\circ)_{\mathrm{c}}$ into $\mathcal{M}(I'^\circ)_{\mathrm{c}}$ follows from the definitions and from $u_!(C_{I^\circ})\subset C_{I'^\circ}$.

\textbf{Corollary 3.3} (of prop. 3.2 and th. 3.1). \textit{Let $u:I\longrightarrow I'$ be as in \textup{3.2}. Suppose moreover $I$ artinian, and}
\printedpage{28}
% head: A. Grothendieck : Dérivateurs, Ch. XVII, Édition des Universités de Montpellier II et Paris VII
\msp{55}
\textit{that we are for $\mathcal{M}$, $W$, $C$, $I$ under one of the conditions \textup{a), b)} of theorem \textup{3.1}. Then $u_!$ induces an arrow}
\[
u_{!\,\mathrm{cof}}\ :\ \mathcal{M}(I^\circ)_{\mathrm{c}} \longrightarrow \mathcal{M}(I'^\circ)_{\mathrm{c}}\ \subset\ \mathcal{M}(I'^\circ)
\]
\textit{which transforms arrows of $W(I^\circ)$ into arrows of $W(I'^\circ)$ (i.e. compatible with the localizers). It thus induces a functor on the localized categories, whence a}
\[
u_!^W\ :\ \mathrm{Ho}_W(I^\circ) \longrightarrow \mathrm{Ho}_W(I'^\circ) \qquad \mbox{\textsuperscript{(}\footnotemark[17]\textsuperscript{)}}.
\tag*{(3.3.1)}
\]
\footnotetext[17]{\textbf{N.B.} This arrow is defined, even without assuming $I'$ ordered, since $u_{!\,\mathrm{cof}}$ still transforms arrows of $W(I^\circ)$ into arrows of $W(I'^\circ)$. Cf. comment p. 86.}

\textsc{Proof}. By theorem 3.1, we know that $\mathcal{M}(I^\circ)$, $W(I^\circ)$, $C_{I^\circ}$ is a structure of a category of models with cofibrations. We also know that $TC_{I^\circ}=C_{I^\circ}\cap W(I^\circ)$, by th. 2.1. $3^\circ$ (p. 27). Thus by 3.2, $u_!$ maps $C_{I^\circ}\cap W(I^\circ)$ into $TC_{I'^\circ}=C_{I'^\circ}\cap W(I'^\circ)$, and a fortiori, into arrows in $W(I'^\circ)$. Moreover, by the theory of categories with cofibrations, we find that every arrow $f$ between cofibrant objects is a composite $f=pi$, where $i\in C_{I^\circ}$, and where $p$ is in $W(I^\circ)$, and moreover a \textit{left inverse} of an arrow in $C_{I^\circ}\cap W(I^\circ)$, i.e. of an arrow of $(TC)_{I^\circ}$. When
\DERmark{ENV1051}

\msp{56}
$f\in W(I^\circ)$, then $i\in W(I^\circ)$, hence $i\in W(I^\circ)\cap C_{I^\circ}=(TC)_{I^\circ}$. Altogether, every arrow $f\in W(I^\circ)_{\mathrm{c}}$ is a composite $pi$ in $\mathcal{M}(I^\circ)_{\mathrm{c}}$, where $i\in(TC)_{I^\circ}$, and $p$ is a left inverse of a $j\in(TC)_{I^\circ}$. Since $u_!$ transforms $(TC)_{I^\circ}$ into $(TC)_{I'^\circ}$, it transforms $f$ into $f'=p'i'$, with $i'\in(TC)_{I'^\circ}\subset W_{I'^\circ}$, and $p'$ a left inverse of a $j'=u_!(j)$ which is also in $W_{I'^\circ}$. But then $p'\in W(I'^\circ)$, hence also $f'=p'i'\in W_{I'^\circ}$, which proves that $u_!(W(I^\circ)_{\mathrm{c}})\subset W(I'^\circ)$. The last assertion comes from the following lemma, applied to $(\mathcal{M}(I^\circ),W(I^\circ),C_{I^\circ})$ in place of $(\mathcal{M},W,C)$:
\DERmark{ENV1052}


\textbf{Lemma 3.4.} \textit{Let $(\mathcal{M},W,C)$ be a category with cofibrations. Then the inclusion $\mathcal{M}_{\mathrm{c}}\subset\mathcal{M}$ induces an equivalence}
\[
\mathcal{M}_{\mathrm{c}}W_{\mathrm{c}}^{-1} \xrightarrow{\ \sim\ } \mathcal{M}W^{-1}=\mathrm{Ho}_W\ .
\]

We will give the proof (by AQT \texttt{[âne qui trotte: at a donkey's trot, i.e. by routine verification]}) later.

This completes the proof of 3.3.

\textbf{3.5.} The question arises whether the functor $u_!^W$ of (3.3.1) is indeed a left adjoint
\msp{57}
to the functor in the opposite direction $u^*_W$, induced by $u^*$. Let us note that this is the case if $u^*$ transforms $\mathcal{M}(I'^\circ)_{\mathrm{c}}$ into $\mathcal{M}(I^\circ)_{\mathrm{c}}$, and when we assume that $I'$ is also artinian and also satisfies one of the hypotheses a) b) of theorem 3.1. For then $(\mathcal{M}(I^\circ),W(I^\circ),C_{I^\circ})$ is a category with cofibrations by theorem 3.1, we can apply the preceding lemma to it, so $u^*_W$ and $u_!^W$ are realized from a pair of functors
\[
\begin{tikzcd}
\mathcal{M}(I^\circ)_{\mathrm{c}} \arrow[r, shift left, "u_{!}^{\mathrm{c}}"] & \mathcal{M}(I'^\circ)_{\mathrm{c}} \arrow[l, shift left, "u_{\mathrm{c}}^*"]
\end{tikzcd}\ ,
\]


\endgroup
% END UNIT der-ch17/der-ch17_p026-030.tex
% BEGIN UNIT der-ch17/der-ch17_p031-035.tex
\begingroup\makeatletter%
\newcommand{\printedpage}[1]{\marginpar{\footnotesize p.~#1}}%
\newcommand{\msp}[1]{\marginpar{\footnotesize [#1]}}%
\newcommand{\Esp}{\operatorname{Esp}}%
\newcommand{\Fl}{\operatorname{Fl}}%
\newcommand{\Ob}{\operatorname{Ob}}%
\newcommand{\Hom}{\operatorname{Hom}}%
\newcommand{\Ho}{\operatorname{Ho}}%
\newcommand{\id}{\mathrm{id}}%
\newcommand{\pr}{\mathrm{pr}}%
\newcommand{\M}{\mathcal{M}}%
\makeatother\let\printedpage\derpp\let\msp\dermsp%
\let\mspin\DERmspin%


\printedpage{29}
% head: A. Grothendieck : Dérivateurs, Ch. XVII, Édition des Universités de Montpellier II et Paris VII
% foot: Version 9 février 2019
which are visibly adjoint to one another (since induced by a pair of adjoint functors % [sic] adjoint -> adjoints, accord avec « lesquels »
$u_!$, $u^*$), and moreover compatible with the localizers. So we can apply the adjoint functor lemma for the localized categories. To sum up:

\textbf{Corollary 3.5.} \textit{Suppose that $I$, $I'$ are artinian, and that we are in one of the following two cases:}
\begin{enumerate}
\item[a)] \textit{The spaces} $\Esp(I)$, $\Esp(I')$ \textit{are noetherian (e.g. $I$, $I'$ finite).}
\msp{page 58}
\item[b)] \textit{$C$ and $TC = C\cap W$ are stable under transfinite composition.}
\end{enumerate}
\textit{Suppose moreover that $u^*$ maps $\M(I'^\circ)_{\mathrm{c}}$ into $\M(I^\circ)_{\mathrm{c}}$, which is the case in particular if $u$ is locally cofiltering, i.e. the localized ordered sets $i'\backslash I$ are locally cofiltering. Then the functor $u_!^W$ constructed in (3.3.1) is left adjoint to $u_W^*$.}

\textbf{Remark 3.6.} When we have a pair of adjoint functors
\[
\begin{tikzcd}
\mathcal{C} \arrow[r, shift left, "u_!"] & \mathcal{C}' \arrow[l, shift left, "u^*"]
\end{tikzcd}
\ ,
\]
with the localizers $W$, $W'$ of $\mathcal{C}$, $\mathcal{C}'$, $u^*$ compatible with the localizers, and if moreover we have a full subcategory $\mathcal{C}_0$ of $\mathcal{C}$ such that
\begin{enumerate}
\item[a)] $u_!|\mathcal{C}_0$ transforms arrows of $W$ into arrows of $W'$, and
\item[b)] the inclusion $\mathcal{C}_0 \hookrightarrow \mathcal{C}$ induces an \textit{equivalence} $\mathcal{C}_0W_0^{-1} \xrightarrow{\sim} \mathcal{C}W^{-1}$ (where $W_0 = W\cap\Fl(\mathcal{C}_0)$),
\end{enumerate}
we find two canonical functors
\[
\begin{tikzcd}
\mathcal{C}W^{-1} \arrow[r, shift left, "u_!^W"] & \mathcal{C}'W'^{-1} \arrow[l, shift left, "u_W^*"]
\end{tikzcd}
\ ,
\]
where $u_!^W$ is defined only up to canonical isomorphism.

\msp{page 59}
Contrary to what \textsc{Anderson} claims, it is not true in general that these two % [sic] par -> pas, faute de frappe
functors are adjoint, without additional hypotheses. Here is a counterexample: we take $\mathcal{C} = \mathcal{C}'$, $u_! = u^* = \id_{\mathcal{C}}$, $W' = \operatorname{Is}(\mathcal{C})$, $W = \Fl(\mathcal{C})$, so $u_W^*$ is the canonical functor
\[
\underbrace{\Pi\,\mathcal{C}}_{\text{fundamental groupoid}} \qquad \longleftarrow \mathcal{C}\ .
\]
Suppose that $\mathcal{C}$ is 1-connected, i.e. $\Pi\,\mathcal{C}$ is equivalent to $e$. Let $e\in\Ob\mathcal{C}$ be such that $\operatorname{End}_{\mathcal{C}}(e)$ is the trivial group, and take $\mathcal{C}_0 = \{e\}$. Then the hypothesis on $\mathcal{C}_0$ is satisfied. But if $u_!^W$ were left adjoint to $u_W^*$, $e$ would be an initial object of $\mathcal{C}$. For example if $\mathcal{C}$ has a final object $e$ (which implies that $\mathcal{C}$ is 1-connected and $e$ rigid), this final object would have to be initial if \textsc{Anderson} were right with ``it is clear that \ldots''.

\msp{page 60}
Nevertheless:

\textbf{Theorem 3.7.} \textit{Let $(\M,W,C)$ be a category with cofibrations stable under finite $\varinjlim$, and $u : I \longrightarrow I'$ an increasing map between artinian ordered sets. We suppose that we are under one of the following hypotheses.}
\begin{enumerate}
\item[a)] $\Esp(I)$, $\Esp(I')$ \textit{are noetherian, e.g. $I$, $I'$ finite.}

\printedpage{30}
% head: A. Grothendieck : Dérivateurs, Ch. XVII, Édition des Universités de Montpellier II et Paris VII
% foot: Version 9 février 2019
\item[b)] \textit{$C$ and $TC \overset{\mathrm{def}}{=} C\cap W$ are stable under transfinite composition, and $\M$ stable under small filtering $\varinjlim$.}
\end{enumerate}
\DERmark{ENV1053}

\textit{Then the functor} $u_!^W : \Ho_W(I^\circ) \longrightarrow \Ho_W(I'^\circ)$ \textit{constructed in} (3.3.1) \textit{is left adjoint to the functor $u_W^*$ in the opposite direction.}

\textsc{Proof.} We use the well-known diagram
\[
\begin{tikzcd}[row sep=large, column sep=large]
 & I \arrow[d, "\partial_1"'] \arrow[dd, bend left=60, "w"] \\
I \arrow[r, shift left, "\partial_0\ (\text{open})"] \arrow[d, "u"'] \arrow[dr, phantom, "\text{cocart.}"] & I\times\boldsymbol{\Delta}^1 \arrow[l, dotted, shift left, "\pr_1"] \arrow[d] \\
I' \arrow[r, shift left, "\text{open imm.}"] & C(u) \arrow[l, dotted, shift left, "v"]
\end{tikzcd}
\]
to factor $u$ as
\[
u \;=\; vw\ ,
\]
where $w$ is a \textit{closed immersion,} and $v$ is obtained from $\pr_1 : I\times\boldsymbol{\Delta}^1 \longrightarrow I$ by push-out. It is a retraction of the mapping cone C(u) onto the open $I'$, whose closed complement is identified with $I$ via $w$.

\msp{page 61}
The gluing functor between $I'$, $I$
\[
H \;:\; I'^\circ\times I \;\longrightarrow\; \mathrm{Ens}
\]
is given by
\[
\underbrace{H(i',i)}_{=\,\Hom_{C(u)}(i',i)} \;=\; \Hom_{I'}(i',u(i))\ . \tag*{$(*)$}
\]
The retraction $v$ from $C(u)$ to $I'$ maps $i$ to $u(i)$, and an arrow $i'\longrightarrow i$ corresponding to $i'\longrightarrow u(i)$ in $I'$ via $(*)$, is transformed by $v$ into that same arrow $i'\longrightarrow u(i)$. Let us calculate, % [sic] transformé -> transformée, accord avec « une flèche »
for a given $i'_0$, the category $i'_0\backslash\Bigl(C(u),\ v : C(u)\longrightarrow I'\Bigr)$, of pairs of an object $\xi$ of $C(u)$ % [sic] des couple -> des couples, accord du pluriel
and an arrow $i'_0\longrightarrow v(\xi)$ in $I'$. We see that $i'_0\backslash C(u)$ admits the open $i'_0\backslash I'$, with the closed complement $i'_0\backslash I$,
\[
\begin{tikzcd}[column sep=large, row sep=huge]
I' \arrow[r, hook, "\text{open}"'] & C(u) \arrow[d] & I \arrow[l, hook', "\begin{array}{l}\text{closed}\\ \text{complement}\\ \text{of } I'\end{array}"'] \\
 & I'\ ,
\end{tikzcd}
\]
these two categories being glued by the bifunctor
\[
\begin{aligned}
(i'_0\backslash I')^\circ \times (i'_0\backslash I) &\quad\longrightarrow\quad \mathrm{Ens}\\
\Bigl((\alpha : i'_0\longrightarrow i'), (\beta : i'_0\longrightarrow u(i), i)\Bigr) &\quad\longmapsto\quad \Hom_{i'_0\backslash I'}(i',u(i))
\end{aligned}
\]
\printedpage{31}
% head: A. Grothendieck : Dérivateurs, Ch. XVII, Édition des Universités de Montpellier II et Paris VII
% foot: Version 9 février 2019
induced by $(*)$\,; i.e. the arrows from $(\alpha : i'_0\longrightarrow i') \in \Ob\,(i'_0\backslash I')$ to $(\beta : i'_0\longrightarrow u(i)) \in i'_0\backslash I$ are the arrows $i' \xrightarrow{\gamma} u(i)$ such that the triangle
\[
\begin{tikzcd}
 & i'_0 \arrow[dl, "\alpha"'] \arrow[dr, "\beta"] & \\
i' \arrow[rr, "\gamma"] & & u(i)
\end{tikzcd}
\]
commutes.

\msp{page 62}
If $i' = i'_0$, $\alpha = \id_{i'_0}$, then this commutativity amounts to $\gamma = \beta$, which means that there is one and only one arrow from $\alpha$ to $(i,\beta)$. This therefore shows that $(i'_0,\id_{i'_0})$ is an initial object of $i'_0\backslash C(u)$. We have therefore proved the lemma:

\textbf{Lemma 3.8.} \textit{Let $u : I\longrightarrow I'$ be an arrow in} (Cat)\textit{, let us factor it as}
\[
I \xrightarrow{\;w\;} C(u) \xrightarrow{\;v\;} I'\ ,
\]
\textit{where $C(u)$ is the mapping cone defined on page} 60\textit{, $w$ a closed immersion with complement $I'$, $v$ the canonical retraction of $C(u)$ onto $I'$. This being so, for every $i'_0\in\Ob I'$, the category $i'_0\backslash(C(u),v)$ admits an initial object – hence $v : C(u)\longrightarrow I'$ is (as thoroughly as can be) cofiltering.}

We can now prove 3.7. First of all, the hypotheses made on $I$, $I'$ carry over to $C(u)$ (artinian hypothesis, and possibly noetherian hypothesis on $\Esp(C(u))$). So
\msp{page 63}
we can apply the preceding results, in particular 3.5, to the arrows $v$, $w$. The preceding lemma indeed shows that $v$ satisfies the conditions of 3.5, so the functor $v_!^W$ constructed in 3.3.1 is left adjoint to $v_W^*$. Consider on the other hand the functor $w_!^W$ constructed in 3.3.1, I say that the same conclusion applies. This comes from the fact that $w$ is a closed immersion. Indeed, in this case it is immediate that
\[
w_! \;:\; \M(I^\circ) \;\longrightarrow\; \M(I'^\circ), \qquad\quad \text{where } J = I' = C(u), % [sic] I = C(u) -> I' = C(u), w a pour but le mapping-cone C(u)
\]
\DERmark{ENV1054}\DERmark{ENV1055}

is compatible with the localizers, since we have
\[
\left\{
\begin{array}{lcll}
w_!(X)(j) & \simeq & X(j) & \text{if } j\in I,\\[1ex]
w_!(X)(j) & \simeq & \emptyset_{\M} & \text{if } j\in J\backslash I,
\end{array}
\right.
\]
by the calculation of the fibers
\[
w_!(X)(j) \;\simeq\; \varinjlim_{i\in I/j} X(i)\ .
\]
So the definition of $w_!^W$ given in 3.3.1 amounts in fact to passing directly to the localized categories for $w_!$, without having to restrict ourselves to $\M(I^\circ)_{\mathrm{c}}$, and the adjoint functor lemma tells us that
\printedpage{32}
% head: A. Grothendieck : Dérivateurs, Ch. XVII, Édition des Universités de Montpellier II et Paris VII
% foot: Version 9 février 2019
\msp{page 64}
this functor is left adjoint to $w_W^*$.

Since we obviously have
\[
u_!^W \;\simeq\; v_!^W\circ w_!^W\ , \qquad\quad u_W^* \;\simeq\; w_W^*\circ v_W^*\ ,
\]
we conclude from this that $u_!^W$ is indeed left adjoint to $u_W^*$, q.e.d.

\textbf{Corollary 3.9.} \textit{Let $(\M,W,C)$ be a category with cofibrations. Then the prederivator $\mathbf{D}_W$, restricted to the category} $\mathrm{Diag}_0$ \textit{consisting of the categories $I$ that are} ordered, \textit{noetherian and such that} $\Esp(I^\circ)$ \textit{is noetherian, satisfies the axioms} Der 1 (decomposition), Der 2 (localization), Der 3 (final object), Der 4 bis (existence of $u_!$), Der 5 bis (mode of calculation of the $u_!$), Der 6 (exactness of $W$-cocartesian squares). \textit{If $C$, $TC = C\cap W$ are stable under transfinite composition, we have the same result on the category} Diag \textit{consisting of the noetherian ordered categories $I$, nothing more.}
\DERmark{ENV1056}\DERmark{ENV1057}\DERmark{ENV1058}


\msp{page 65}
\textbf{3.10.} Thus, to have a derivator in the full sense of the term, either on $\mathrm{Diag}_0$, or on Diag, all that is missing is Der 4 (existence of the $u_*$), Der 5 (calculation of the $u_*$) and Der 6 (exactness of $W$-cartesian squares). We are going to obtain them in principle with the aid of a \textit{category with fibrations} structure $(\M,F,W)$ associated to $(\M,W)$, i.e. the kind of structure dual to the one introduced % [sic] associé -> associée, accord avec « structure »
on page 39. These two structures do not have to be related to each other in any way other than by the fact that they correspond to the same $W$. They are considerably easier to construct than Quillen structures $(W,C,F)$, where $C$, $F$ must be related in an extremely strict way, with the properties $C \longleftrightarrow \underbrace{TF}_{=\,F\cap W}\ ,\ \underbrace{TC}_{=\,C\cap W} \longleftrightarrow F$ (of which there is no mention in the axiomatics of K. S. \textsc{Brown}), and factorization for $(C,TF)$ and $(TC,F)$.

\msp{page 66}
But returning to the construction of the $u_!^W$, there remains the problem of constructing them for an arrow $u : I\longrightarrow I'$ \textit{whatsoever} in Cat. % [sic] u : I' -> I -> u : I -> I', u_!^W va de Ho_W(I°) dans Ho_W(I'°)
\DERmark{ENV1059}


To tell the truth, the construction of the functor
\[
u_!^W \;:\; \Ho_W(I^\circ) \;\longrightarrow\; \Ho_W(I'^\circ) \tag*{(3.3.1)}
\]
is valid without supposing that $I'$ is an ordered category – it can be any category. This comes from the fact that the functor induced by $u_!$,
\[
u_{!\,\mathrm{cof}} \;:\; \M(I^\circ)_{\mathrm{c}} \;\longrightarrow\; \M(I'^\circ)
\]
is still compatible with the localizers, as we see by the calculation of the fibers
\[
u_!(X)(i') \;\simeq\; \varinjlim_{i\in I/i'} X(i)\ . % [sic] u_!(X)(i) -> u_!(X)(i'), la fibre est prise en i'
\]
\DERmark{ENV1060}

Thus, we are reduced to the cases where $I' = e$, at the cost of replacing $I$ by $I/i'$, and noting that % [sic] aux cas -> au cas, singulier
$(I/i'\longrightarrow I$ being étale) the inverse image of $X$ on $I/i'$ is cofibrant (cf. 3.2.a). % [sic] cofibrant -> cofibrante, accord avec « l'image inverse »

\msp{page 67}
But there arises again the perplexity whether the functor (3.3.1) thus constructed is indeed left adjoint to $u_W^*$. We can think of using once again the same factorization as $u = vw$, by
\printedpage{33}
% head: A. Grothendieck : Dérivateurs, Ch. XVII, Édition des Universités de Montpellier II et Paris VII
% foot: Version 9 février 2019
using the mapping cone. The difficulty, however, is that now $J = C(u)$ is no longer an ordered category, so we do not have a ready-made definition of $C_{J^\circ}\subset\Fl(\M(J^\circ))$, giving rise to a subcategory $\M(J^\circ)_{\mathrm{c}}$ etc. In fact, it would seem that we are practically paralyzed, so long as we do not manage to construct on $\M(J^\circ)$ also a category with cofibrations structure, behaving well with respect to the functors $v_!$. The % [sic] au foncteurs -> aux foncteurs, accord du pluriel
``proof'' by \textsc{Anderson} of the existence of $u_!$ for \textit{every} arrow in Cat, is a total flop (even for $I\longrightarrow e$, $I$ a finite ordered set – but there I was able to set things right), for he reduces to the case where $I$ is ordered and locally finite, and he
\msp{page 68}
assumes without further ado that \textit{in this case} $u_!^W$ is left adjoint to $u_W^*$, and \textit{moreover,} that if $u$ is $\mathbf{W}_0$-connected, hence $u^* : \M(I'^\circ)\longrightarrow\M(I^\circ)$ fully faithful, that the same is true of $u_W^*$ obtained by passing to the localized categories. Now even if we grant that $u_!^W$ is left adjoint to
$u_W^*$, for lack of having an adjunction morphism close at hand ``down below'' (in $\Ho_W(I'^\circ)$), % [sic] u_*^W -> u_W^*, notation employée partout ailleurs
\DERmark{ENV1061}

obtained from the adjunction morphism in $\M(I^\circ)$, it is not at all clear that $u_W^*$ is fully faithful, i.e. $u_!^W$ a localization functor. Counterexample 3.6 clearly shows that, even if we have adjunction (i.e. if $e$ is an initial object of $\mathcal{C} = \mathcal{C}'$), the functor $u^* = \id_{\mathcal{C}}$ may well be fully faithful (and even an equivalence of categories), the localized functor has no reason to be so: this would prove that every object of Cat having an initial object % [sic] pouverait -> prouverait, faute de frappe
is equivalent to $e$\,!

\msp{page 69}
\textbf{3.11.} Let us say that we have a \textit{strict category with cofibrations,} if in the factorization axiom $\mathrm{Cof}_3$ (cf. p. 39), the factorization can be taken functorially in $f$. It amounts to $\pm$ the same thing to say that the factorization axiom remains valid, when we think of a category $\M(I)$ for $I$ in Cat, defining $W(I)$ and $C(I)$ there (not to be confused with $C_I$) fiber by fiber. This amounts to the same thing at least if $\M$ is equivalent to a small category. Otherwise, we must assume functoriality for every \textit{small} full subcategory of $\underline{\mathrm{Fl}}(\M)$. It is better to define the strict category with cofibrations by this stability property.

We then see that the prederivator $\mathbf{D} = \mathbf{D}_{\M}$ on Cat is a derivator \textit{relative to} $\mathcal{D} = \mathrm{Diag}_0$ defined above (noetherian ordered categories $I$, such that $\Esp(I^\circ)$ is noetherian), or even relative to $\mathcal{D} = \mathrm{Diag}$ (noetherian ordered categories $I$, nothing more) when $C$ and $TC$ are stable under transfinite composition. By this I mean that $u_!$ exists when $u$ is an arrow of $\mathcal{D}$,
\DERmark{ENV1062}\DERmark{ENV1063}\DERmark{ENV1064}

\msp{page 70}
and that this remains valid when we pass to an induced derivator $\mathbf{D}_Z$ ($Z$ arbitrary in Cat)\,; in other words, that $v_!$ exists when $v$ is of the form $\id_Z\times u$, where $u : I\longrightarrow I'$ is an % [sic] d'autre termes -> d'autres termes, accord du pluriel
arrow of $\mathcal{D}$, and that the calculation is done \textit{over} $I'$ according to the rule Der 5 bis. This therefore means that the base change arrow for the $f_!$ is an isomorphism in the case % [sic] un isomorphismes -> un isomorphisme, accord du singulier
that follows:
\[
\begin{tikzcd}[row sep=huge, column sep=huge]
I\times Z \arrow[d, "u\times\id_Z"] & (i'\backslash I)\times Z \arrow[l] \arrow[d, "\pr_2"] \\
I'\times Z & \{i'\}\times Z \arrow[l]
\end{tikzcd}
\ ,
\]
where $i'$ is an arbitrary object of $I'$.


\endgroup
% END UNIT der-ch17/der-ch17_p031-035.tex
% BEGIN UNIT der-ch17/der-ch17_p036-040.tex
\begingroup\makeatletter%
\newcommand{\printedpage}[1]{\marginpar{\footnotesize p.~#1}}%
\newcommand{\msp}[1]{\marginpar{\footnotesize [#1]}}%
\newcommand{\circled}[1]{\tikz[baseline=(c.base)]\node[draw,circle,inner sep=1pt](c){#1};}%
\newcommand{\defeq}{\overset{\mathrm{def}}{=}}%
\newcommand{\postilde}{\tilde{}}%
\newcommand{\Fl}{\mathrm{Fl}}%
\newcommand{\Cof}{\mathrm{Cof}}%
\makeatother\let\printedpage\derpp\let\msp\dermsp%
\let\mspin\DERmspin%

\DERmark{ENV1065}


\printedpage{34}
% head: A. Grothendieck : Dérivateurs, Ch. XVII, Édition des Universités de Montpellier II et Paris VII
% foot: Version 9 février 2019
The \textit{strict} hypothesis for a category with cofibrations seems fairly innocuous to me, in most cases known to me, it is satisfied. But even with this hypothesis, I see no way to prove the existence of $u_!$ for \textit{every} arrow of Cat, and this is of course a very awkward situation. One would very much like to isolate conditions that ensure the validity of Der 4 bis and of Der 5 bis (existence of $u_!$ for
\msp{page 71}
\textit{every} arrow of Cat, and way of computing the fibers \dots).

\section*{4\quad Special categories with cofibrations}

\textbf{4.1.} I now suppose that
\begin{itemize}[leftmargin=3em]
\item[\circled{a}] $\mathcal{M}$ is stable under filtering $\varinjlim$, and
\item[\circled{b}] that its objects are accessible,
\end{itemize}
so that we are under the conditions for applying the factorization theorem (XV, \S1). We suppose:
\begin{itemize}[leftmargin=3em]
\item[\circled{c}] given $W, C \subset \Fl(\mathcal{M})$, satisfying the conditions $\Cof_1$, $\Cof_2$ of page 39. For $\Cof_3$, cf. below.
\item[\circled{d}] $C$, $TC \defeq C \cap W$ \textit{stable under transfinite composition} (this is the case if the structure comes from a closed Quillen structure $(W, C, F)$).
\end{itemize}
\[
TF \defeq C_* \;\subset\; \Fl(\mathcal{M})\ , \tag*{(4.1)}
\]
it follows from the factorization theorem that every arrow $f$ of $\mathcal{M}$ factors as
\[
f \;=\; pi\ , \qquad p \in TF,\ i \in C\ ,
\]
\textit{provided} we make the hypothesis
\begin{itemize}[leftmargin=3em]
\item[\circled{e}] \mbox{}
\[
\exists\, C_0 \subset C,\ \textit{small} \text{ subset of } C, \text{ such that } C \subset \tilde{C}_0, \text{ i.e. } TF = (C_0)_*. \tag*{(4.3)}
\]
\end{itemize}
Under these conditions, if we suppose
\begin{itemize}[leftmargin=3em]
\item[\circled{f}] \mbox{}
\[
TF \subset W\ , \tag*{(4.4)}
\]
\end{itemize}
\msp{page 72}
then the factorization axiom $\Cof_3$ is satisfied. Conversely, if $\Cof_3$ is satisfied, and if $W$ is stable under direct factors, then we have 4.4. Thus if $W$ is stable under direct factors, $\Cof_3$ is equivalent to (4.4).

\textbf{4.2.} It seems to me that these conditions, in practice, will hardly ever be fulfilled, unless $W$ (which is given in the very first place, satisfying $\Cof_0$ \textsuperscript{(}\footnote[18]{i.e. weak saturation.}\textsuperscript{)}), is stable under $\varinjlim$ that are
\DERmark{ENV1066}

\printedpage{35}
% head: A. Grothendieck : Dérivateurs, Ch. XVII, Édition des Universités de Montpellier II et Paris VII
% foot: Version 9 février 2019
filtering. At that point, to construct the cofibration structure, we start from $C_0$ a small subset of $\Fl(\mathcal{M})$, taken so that we have
\[
(C_0)_* \;\subset\; W\ , \tag*{(4.5)}
\]
we take for $C$ at choice the saturation $C'_0$ of $C_0$ under cobase change and transfinite composition, or $C = \tilde{C}_0 =$ set of direct factors of $C'_0$. The conditions $\Cof_1$ (stabilities of $C$), $\Cof_3$ (factorization) are satisfied automatically, and it remains to take care of $\Cof_2$, which comes down to the stability of $C{\cap}W$ under cobase change (stability under composition and that it contain the isomorphisms, being immediate). This is a delicate condition, and the only standard way I see to ensure it, is by supposing
\[
C_0 \;\subset\; \Cof_W\ ; \tag*{(4.6)}
\]
\msp{page 73}
the stability of $W$ under filtering $\varinjlim$ (hence also under direct factors) then ensures that we also have
\[
C \;\subset\; \Cof_W\ , \tag*{(4.7)}
\]
and consequently
\[
\left\{
\begin{aligned}
& C \cap W \subset \Cof_W \cap W = W^{\mathrm{univ}}, \\
& \text{hence }\ TC \defeq C \cap W = C \cap W^{\mathrm{univ}}\ ,
\end{aligned}
\right. \tag*{(4.8)}
\]
which ensures the stability of $C \cap W$ under cobase change. The stability of $TC$ under transfinite composition follows from that of $C$, and from the stability of $W$ under transfinite composition.

\textbf{4.3.} I return to the general conditions\circled{a} to\circled{f} of 4.1, without necessarily supposing $W$ stable under filtering $\varinjlim$. I now want to extend the cofibration structure to the $\mathcal{M}(I)$, $I$ arbitrary in Cat. For this we do not (I think) have to use the developments of sections 2, 3, but we will make the connection with these later on.

Let us denote by $I_0$ the discrete category having the same objects as $I$, whence a canonical functor
\[
v = v_I \ :\ I_0 \longrightarrow I\ . \tag*{(4.9)}
\]
The idea is to define the structure of $\mathcal{M}(I)$ via the evident structure of $\mathcal{M}(I_0) \simeq \mathcal{M}^{I_0}$, by playing on the pair of adjoint functors
\[
\begin{tikzcd}
\mathcal{M}(I_0) \arrow[r, shift left, "\varphi = v_!"] & \mathcal{M}(I) \arrow[l, shift left, "\psi = v^*"]
\end{tikzcd}
\ .
\]
\msp{page 74}
We are going to present things in a slightly more general way.

\textbf{Definition 4.4.} I call a \textit{special} cofibration structure a structure $(\mathcal{M}, W, C)$ satisfying conditions a) to f) of 4.1, which I reformulate:
\begin{itemize}[leftmargin=3em]
\item[\circled{a}] $\mathcal{M}$ stable under small direct limits, and the objects of $\mathcal{M}$ are accessible.
\printedpage{36}
% head: A. Grothendieck : Dérivateurs, Ch. XVII, Édition des Universités de Montpellier II et Paris VII
% foot: Version 9 février 2019
\item[\circled{b}] $W$ is weakly saturated (i.e. $\Cof_0$ of p. 39), $C$ stable under cobase change and under composition, $TC \defeq C{\cap}W$ stable under cobase change (and under composition automatically).
\item[\circled{b$'$}] $C$ and $TC$ are stable even under composition that is \textit{transfinite.}
\item[\circled{c}] $\exists$ small subset $C_0 \subset C$ such that $C \subset \tilde{C}_0$ ($\Longrightarrow \tilde{C} = \tilde{C}_0$).
\item[\circled{d}] $TF \defeq C_* \subset W$ (which implies, together with\circled{a} to\circled{c}, the factorization axiom $\Cof_3$ of page 39).
\end{itemize}

\textbf{4.5.} We are going to give a general recipe for extending such a structure, when we have a pair of adjoint functors
\[
\begin{tikzcd}
\mathcal{M} \arrow[r, shift left, "v_!"] & \mathcal{M}' \arrow[l, shift left, "v^*"]
\end{tikzcd}
\ ,
\]
$\mathcal{M}$ being equipped with a special cofibration structure $(W, C)$, and $\mathcal{M}'$ satisfying the preliminary condition\circled{a}. We will then define the subsets $W'$, $C'$ of $\Fl(\mathcal{M}')$
\msp{page 75}
by
\[
\left\{
\begin{aligned}
W' \;&=\; (v^*)^{-1}(W) \\
C' \;&=\; \text{closure stable under cobase change} \\
&\phantom{\;=\;}\ \ \text{and transfinite composition of } v_!(C),
\end{aligned}
\right. \tag*{(4.5.1)}
\]
thus
\[
v_!(C) \;\subset\; C' \;\subset\; (v_!(C))\postilde \;=\; \tilde{C}'\ . \tag*{(4.5.2)}
\]
I say that $(W', C')$, and at choice $(W', \tilde{C}')$, is a special cofibration structure on $\mathcal{M}'$. We are going to check the various conditions of 4.4. Condition a) is already assured.

b) \textsuperscript{(}\footnote[19]{\textbf{N.B.} It is better to start by proving c).}\textsuperscript{)}. The weak saturation of $W'$ follows from that of $W$, via (4.5.1). The stability of $C'$ under cobase change and under composition (even transfinite) is assured by definition of $C'$, and likewise for $\tilde{C}'$ (being Q-closed on the left). Stability of $TC' \defeq C'{\cap}W'$ under cobase change? It suffices to see that if we have a cocartesian square
\[
\begin{tikzcd}
X'_0 \arrow[r, "i'"] \arrow[d] \arrow[dr, phantom, "\text{\scriptsize cocart.}"] & X'_1 \arrow[d] \\
X'_2 \arrow[r, "j'"] & X'_3
\end{tikzcd}
\]
in $\mathcal{M}'$, with $i'$ in $TC' = C' \cap W'$, then $j' \in W'$ (for we already know $j' \in C'$), i.e.
\marginpar{$||$}% marginal sign || printed in the left margin of this line
$v^*(j') \in W$. We are going to suppose that $v^*$ \textit{commutes with amalgamated sums,} and we are thus reduced to proving that $v^*(i') \in TC$. Thus % [sic] $v^*(j') \in W'$ -> $v^*(j') \in W$, since $W' = (v^*)^{-1}(W)$
\DERmark{ENV1067}

\printedpage{37}
% head: A. Grothendieck : Dérivateurs, Ch. XVII, Édition des Universités de Montpellier II et Paris VII
% foot: Version 9 février 2019
\msp{page 76}
we must prove the

\textbf{Lemma 4.5.2.} \textit{Under the conditions} (4.5.3) \textit{and} (4.5.4) \textit{below, we have}
\[
v^*(C') \;\subset\; C\ , \qquad v^*(TC') \;\subset\; TC\ ,
\]
\textit{and if $C = \tilde{C}$, we have $v^*(\tilde{C}') \subset C$, $v^*(T\tilde{C}') \subset TC$.}

For $v^*(C') \subset C$, we must first prove
\[
v^*(v_!(C)) \;\subset\; C\ , \tag*{(4.5.3)}
\]
and I see no way -- we have to take it as a hypothesis. Moreover, that $v^*$ commutes with arbitrary $\varinjlim$ (cobase change already admitted -- now we need filtering direct limits)
\[
v^* \text{ commutes with small } \varinjlim \text{ of any kind} \tag*{(4.5.4)}
\]
to conclude from (4.5.3) that $v^*(C' =$ closure of $v_!(C)$ under cobase change and transfinite composition) is also in $C$. Passing to $v^*(TC') \subset TC$, by what precedes, $v^*(C') \subset C$, it suffices to prove $v^*(TC') \subset W$, but $v^*(TC') \subset v^*(W') \subset W$ by definition of $W'$.

Let us suppose $C = \tilde{C}$, i.e. $C$ stable under direct factors, and let us prove $v^*(\tilde{C}') \subset C$, $v^*(T\tilde{C}') \subset TC$. By c) below, $\tilde{C}'$ is formed of the direct factors
\msp{page 77}
of elements of $C'$, hence $v^*(C') \subset C$ implies $v^*(\tilde{C}') \subset C$ since $C$ stable under direct factors.

Let us finally prove $v^*(T\tilde{C}') \subset TC \defeq C \cap W$. As $T\tilde{C}' = \tilde{C}' \cap W'$, this follows from $v^*(\tilde{C}') \subset C$, and from $v^*(W') \subset W$.

Let us prove b$'$). The stability of $C'$ and of $\tilde{C}'$ under transfinite composition is assured, it remains to prove the stability of $TC' = C' \cap W'$ and of $T\tilde{C}' = \tilde{C}' \cap W'$. It suffices to see that the transfinite composite $f$ in question is in $W'$, hence that $v^*(f) \in W$, but this follows from the fact that $v^*$ commutes with transfinite composition by (4.5.4), and that $v^*$ takes $TC'$ and $T\tilde{C}'$ into $TC$.

\textit{Let us prove} c) for $C'$ (hence also for $\tilde{C}'$). Let $C_0$ be a small subset of $C$ such that $C \subset \tilde{C}_0$, i.e. $C$ is contained in the closure of $C_0$ under the three operations of cobase change, of transfinite composition, of direct factors. We then see that $v_!(C)$ is in the closure of $v_!(C_0)$ under the operations, hence also $C'$, and we can then take $C'_0 = v_!(C_0)$.

\textit{Let us finally prove} d), it is the same assertion for $C'$ and for $\tilde{C}'$, since
\[
C'_* \;=\; \tilde{C}'_*\ (\ \defeq\ TF'\ )\ ,
\]
\msp{page 78}
we must therefore prove
\[
TF' \defeq C'_* \;\subset\; W'\ .
\]
This will follow from the
\printedpage{38}
% head: A. Grothendieck : Dérivateurs, Ch. XVII, Édition des Universités de Montpellier II et Paris VII
% foot: Version 9 février 2019

\textbf{Lemma 4.5.3.} \textit{We have $TF' = (v^*)^{-1}(TF)$.}

This is formal from the definition of $TF'$ as
\[
TF' \;=\; (v_!(C))_*\ ,
\]
of $TF$ as
\[
TF \;=\; C_*\ ,
\]
and from the adjunction of $v_!$, $v^*$. As $TF \subset W$, we conclude
\[
TF' \;=\; (v^*)^{-1}(TF) \;\subset\; (v^*)^{-1}(W) = W'\ ,
\]
q.e.d.

Thus we have proved this:

\textbf{Proposition 4.6.} \textit{Let us consider a pair of adjoint functors}
\[
\begin{tikzcd}
\mathcal{M} \arrow[r, shift left, "v_!"] & \mathcal{M}' \arrow[l, shift left, "v^*"]
\end{tikzcd}
\ ,
\]
\textit{where $\mathcal{M}$ is equipped with a cofibration structure that is} special $(W, C)$\textit{. We further suppose the following conditions satisfied.}
\[
\begin{aligned}
&1^\circ)\quad \mathcal{M}' \textit{ stable under small } \varinjlim\ \textit{, its elements are accessible.} \\
&2^\circ)\quad v^* \textit{ commutes with small } \varinjlim\ . \\
&3^\circ)\quad v^*(v_!(C)) \subset C.
\end{aligned} \tag*{(4.6.1)}
\]
\msp{page 79}
\textit{Let us consider the subsets $W'$, $C'$ of} $\Fl(\mathcal{M}')$ \textit{defined by}
\[
\left\{
\begin{aligned}
W' \;&=\; (v^*)^{-1}(W), \\
C' \;&=\; \textit{closure of } v_!(C) \textit{ under cobase change and under composition} \\
&\phantom{\;=\;}\ \textit{that is transfinite.}
\end{aligned}
\right. \tag*{(4.6.2)}
\]
\textit{Then $(\mathcal{M}', W', C')$ is a special cofibration structure. If moreover $C = \tilde{C}$, i.e. $C$ is stable under direct factors, then $(\mathcal{M}', W', \tilde{C}')$ is also a special cofibration structure} \textsuperscript{(}\footnote[20]{\textbf{N.B.} This is a general fact for special categories with cofibrations, and should be made explicit below (cf. 5.1). The hypothesis on $C$ is unnecessary here (or rather, it suffices instead that $W$ be stable under direct factors \textit{and} under transfinite composition, which is not a more cheering hypothesis \dots).}\textsuperscript{)}.

\textbf{4.7.} I would like, with a view to statements of type 3.7 (establishing the existence of functors of type $u_!^W$), to have a hold on $TC' = C'{\cap}W'$, and on $T\tilde{C}' = \tilde{C}'{\cap}W'$, to be precise, I would like to have the relation
\[
v_!(TC)\ \subset TC' \subset v_!(TC)\postilde\ , \tag*{(4.7.1)}
\]
or even the slightly stronger relation
\[
\underbrace{v_!(TC)\ \subset TC'}_{} \ \overset{\text{tautol.}}{\subset}\ \underbrace{T\tilde{C}' \;\subset\; (v_!(TC))\postilde}_{}\ , \tag*{(4.7.2)}
\]


\endgroup
% END UNIT der-ch17/der-ch17_p036-040.tex
% BEGIN UNIT der-ch17/der-ch17_p041-045.tex
\begingroup\makeatletter%
\newcommand{\printedpage}[1]{\marginpar{\footnotesize p.~#1}}%
\newcommand{\msp}[1]{\marginpar{\footnotesize [#1]}}%
\newcommand{\defeq}{\overset{\mathrm{def}}{=}}%
\newcommand{\fnm}[1]{\textsuperscript{(}\footnotemark[#1]\textsuperscript{)}}%
\newcommand{\cl}{^{\tilde{}}}% tilde printed as exponent after a symbol or a parenthesis
\makeatother\let\printedpage\derpp\let\msp\dermsp%
\let\mspin\DERmspin%

\DERmark{ENV1068}


\printedpage{39}
% head: A. Grothendieck : Dérivateurs, Ch. XVII, Édition des Universités de Montpellier II et Paris VII
which decomposes into two
\[
\left\{
\begin{array}{rcl}
v_{!}(TC) & \subset & TC' \\
T\tilde{C}' & \subset & (v_{!}(TC))\cl.
\end{array}
\right.
\tag*{(4.7.3)}
\]
The less delicate of the two is the \msp{page 80} first, $v_{!}(TC) \subset TC'$, which amounts, by virtue of the definitions, implying
\[
v_{!}(TC) \subset v_{!}(C) \subset C' \, ,
\]
to
\[
v_{!}(TC) \;\subset\; W' \; ,
\]
i.e. by definition of $W'$, to
\[
v^{*}v_{!}(TC) \;\subset\; TC \; ,
\tag*{(4.7.4)}
\]
which is a hypothesis of type (4.6.1) $3^{\circ}$), to be added to the hypotheses (4.6.1).

The hypothetical relation
\[
\underbrace{T\tilde{C}'}_{=\,\tilde{C}'\cap W'} \;\subset\; (v_{!}(TC))\cl, \quad \text{or only} \quad TC' \subset (v_{!}(TC))\cl,
\tag*{$(*)$}
\]
appears to me to be of a distinctly more delicate nature. Since $T\tilde{C}' \supset TC' \supset v_{!}(TC)$, it is equivalent to the \textit{equality} of $(T\tilde{C}')_{*}$ resp. $(TC')_{*}$ and $(v_{!}(TC))_{*} = (v^{*})^{-1}(TC)_{*}$. Let us put
\[
\left\{
\begin{array}{lcl}
(TC)_{*} & \defeq & F \\
(TC')_{*} & \defeq & F' \;\supset\; (T\tilde{C}')_{*} \;=\; F'_{\sim} \; .
\end{array}
\right.
\tag*{(4.7.5)}
\]
We therefore have, by virtue of the inclusion
\[
v_{!}(TC) \;\subset\; TC' \; ( \;\subset\; T\tilde{C}' \; ) \; ,
\]
the relation
\[
F'_{\sim} \;\subset\; F' \;\subset\; (v^{*})^{-1}(F) \; ,
\tag*{(4.7.6)}
\]
and the question is therefore whether we have equality $F' = (v^{*})^{-1}(F)$, or even $F'_{\sim} = (v^{*})^{-1}(F)$ (whence $F'_{\sim} = F'$).

\msp{page 81}

In fact, we show fairly easily that we have $F'_{\sim} = F'$, i.e. $T\tilde{C}' \defeq \tilde{C}' \cap W' \subset (C' \cap W')\cl$ -- it is a general story about special categories with cofibrations \textsuperscript{(}\footnote[21]{cf. prop. 5.1 below.}\textsuperscript{)}. But I hardly see how to prove (under suitable hypotheses) that $F' = (v^{*})^{-1}(F)$. I shall have to come back to this if need be \textsuperscript{(}\footnote[22]{see however prop. 5.8 below.}\textsuperscript{)}. For the moment, I shall confine myself to reaping the fruit of 4.6 in a particular case:

\textbf{Proposition 4.8.} \textit{Let} $(\mathcal{M}, W, C)$ \textit{be a special category with cofibrations} (definition 4.4, page 74)\textit{. Let $I$ be an object of} Cat\textit{, $I_0$ the associated discrete category, let us consider $v : I_0 \longrightarrow I$ and} % [sic] page 94 -> page 74, apparently: the same déf. 4.4 is cited as "page 74" on printed p. 42
\DERmark{ENV1069}

% foot: Version 9 février 2019
\printedpage{40}
% head: A. Grothendieck : Dérivateurs, Ch. XVII, Édition des Universités de Montpellier II et Paris VII
\[
\begin{tikzcd}
\mathcal{M}(I_0) \arrow[r, shift left, "v_{!}"] & \mathcal{M}(I) \arrow[l, shift left, "v^{*}"]
\end{tikzcd} \; ,
\]
\textit{finally let us endow $\mathcal{M}(I_0) = \mathcal{M}^{I_0}$ with the structure of a category with cofibrations,} product \textit{of those of the factors, i.e. let us define $W(I_0)$, $C(I_0)$ ``fiber by fiber''. (This coincides with the general definition given in paragraph} 2\textit{, for $I_0$ endowed with the discrete order structure, and makes $\mathcal{M}(I_0)$ a special category with cofibrations.) Then the couple $(v_{!}, v^{*})$ satisfies} \msp{page 82} \textit{the conditions of proposition} 4.6\textit{, as well as the condition supplementary to} (4.6.1)\textit{,}
\[
4^{\circ}) \qquad v^{*}(v_{!}(TC(I_0)))\subset TC(I_0) \; .
\tag*{(4.8.1)}
\]
\DERmark{ENV1070}

\textit{Thus we find a structure of a special category with cofibrations on $\mathcal{M}(I)$, by putting}
\[
W(I) \;=\; (v^{*})^{-1}(W(I_0))
\tag*{(4.8.2)}
\]
\textit{(this is in fact equivalent to the usual definition of $W(I)$) and}
\[
C(I) \quad = \quad
\begin{array}[t]{l}
\textit{closure of } v_{!}(C(I_0)) \textit{ under the operations of cobase}\\
\textit{change and transfinite composition.}
\end{array}
\tag*{(4.8.3)}
\]
\textit{We have moreover under these conditions:}
\[
\underbrace{TF(I)}_{\defeq\,(C(I))_{*}} \;=\; (v^{*})^{-1}( \; \underbrace{TF(I_0)}_{\substack{\defeq\,(C_{I_0})_{*},\\ \textit{is also defined}\\ \textit{fiber by fiber}}} \; ) \; ,
\tag*{(4.8.8)}
\]
\[
\left\{
\begin{array}{rcl}
v_{!}( \, \overbrace{(TC)(I_0)}^{C(I_0)\cap W(I_0)} \, ) & \subset & TC(I) \qquad\qquad\qquad \text{\fnm{23}} \\
F(I) \defeq (TC(I))_{*} & \subset & (v^{*})^{-1}(F(I_0)).
\end{array}
\right.
\tag*{(4.8.9)}
\]
\footnotetext[23]{These inclusions should give $TC(I)=(v_!(TC(I_0)))\cl$ and $F(I)=(v^*)^{-1}(F(I_0))$, but I do not know how to prove it without making stronger hypotheses, of the type \textit{special Quillen category.}}
\DERmark{ENV1071}


\textsc{Proof}. We already know that $\mathcal{M}(I)$ is stable under small $\varinjlim$ , and that every element is accessible, as follows at once from the appropriate properties of $\mathcal{M}$. It is also clear that $v^{*}$ commutes with small $\varinjlim$ . It remains to prove (4.6.1. $3^{\circ}$), and $4^{\circ}$ in (4.8.1), i.e. \msp{page 83} that $v^{*}v_{!} : \mathcal{M}(I_0) \longrightarrow \mathcal{M}(I_0)$ maps $C$ into $C$, $TC$ into $TC$. Now we have
\[
v_{!}(F_0)(i) \;=\; \varinjlim_{(\alpha,j)\in i\backslash I_0} F_0(j) \qquad\quad \text{\fnm{24}}
\]
\footnotetext[24]{where $j \in \mathrm{Ob}\, I_0 = \mathrm{Ob}\, I$, and $\alpha : i \longrightarrow v(j) = j$.}
and $i\backslash I_0$ is the discrete category over $I_0$, whose fiber at $j \in \mathrm{Ob}\, I_0 \;=\; \mathrm{Ob}\, I$ is $\mathrm{Hom}_I(i,j)$:
\[
i\backslash I_0 \;\;=\;\; \coprod_{j\in \mathrm{Ob}\, I} \text{(discrete category associated to $\mathrm{Hom}_I(i,j)$)} \; ,
\]
% foot: Version 9 février 2019
\printedpage{41}
% head: A. Grothendieck : Dérivateurs, Ch. XVII, Édition des Universités de Montpellier II et Paris VII
thus
\[
v_{!}(F_0)(i) \;=\; \coprod_{j\in \mathrm{Ob}\, I = \mathrm{Ob}\, I_0} \underbrace{F_0(j)^{\mathrm{Hom}_I(i,j)}}_{\substack{\text{direct sum}\\ \text{of } \mathrm{Hom}_I(i,j)\\ \text{copies of } F_0(j)}} \, ,
\]
so we find
\[
v^{*}v_{!}F_0( \, \underbrace{\; i \;}_{\in \mathrm{Ob}\, I_0 \,=\, \mathrm{Ob}\, I} \, ) \;=\; v_{!}(F_0)(i) \;=\; \text{\tikz[baseline=0pt]\draw[->] (0,0) .. controls (0.4,0) .. (0.5,0.45);} \quad .
\]
% the curved arrow points up to the underbraced expression F_0(j)^{Hom_I(i,j)} of the preceding display
Since $C$ and $TC$ are stable under arbitrary direct sums (being so under finite direct sums, because of the finite stability conditions on $C$ and $TC$, and also under infinite sums because of the axiom of transfinite stability \dots\textsuperscript{(}\footnote[25]{deserves a general lemma \dots}\textsuperscript{)}), it does indeed follow that $v^{*}v_{!}$ maps $C(I_0)$ and $TC(I_0)$ into themselves. % [sic] somme infinies -> sommes infinies, agreement with the plural adjective

Thus, the first assertion of 4.8 follows from 4.6. Moreover, (4.8.8) and (4.8.9) were seen in 4.7, under the general conditions of 4.6 and assuming in addition the $4^{\circ}$ of (4.7.4).

\msp{page 84}

\section*{5\quad Relation of special categories with cofibrations to special Quillen categories}

Let $(\mathcal{M}, W, C)$ be a special category with cofibrations. I introduce
\[
TF \;=\; C_{*}, \qquad\quad TC \;=\; C \cap W, \qquad\quad F \;=\; (TC)_{*} \; ,
\tag*{(5.1)}
\]
so we have
\[
\begin{tikzcd}[row sep=large, column sep=large]
TC \arrow[r, hook] \arrow[d, leftrightarrow] & C \arrow[d, leftrightarrow, "\shortmid"{marking, sloped}] \\
F & TF \arrow[l, hook']
\end{tikzcd} \; ,
\tag*{(5.2)}
\]
$(C, TF)$ being a right closed Quillen precouple. We then know (XV 4.2 cor. 1, page 31) that $\tilde{C}$ consists of the direct factors of elements of $C$. Let us put
\[
T\tilde{C} \;=\; \tilde{C} \cap W \; ,
\tag*{(5.3)}
\]
whence
\[
TC \;\subset\; T\tilde{C} \; .
\tag*{(5.4)}
\]

\textbf{Proposition 5.1} \textit{Let $(\mathcal{M}, W, C)$ be a special category with cofibrations. We have $T\tilde{C} \subset (TC)\cl$ (and even $T\tilde{C} \subset (TC)^{\natural}$ \textsuperscript{(}\footnote[26]{\textbf{N.B.} For a subset $A$ of $\mathrm{Fl}(\mathcal{M})$, $A^{\natural}$ denotes the set of direct factors of elements of $A$.}\textsuperscript{)}), or what amounts to the same thing by virtue of (5.4), $TC$ and $T\tilde{C}$ have the same left $Q$-closure:}
\[
 (T\tilde{C})\cl = (TC)\cl \qquad\quad \textit{(and even } T\tilde{C} \subset (TC)^{\natural}\textit{),}
\tag*{(5.5)}
\]
\DERmark{ENV1072}

% foot: Version 9 février 2019
\printedpage{42}
% head: A. Grothendieck : Dérivateurs, Ch. XVII, Édition des Universités de Montpellier II et Paris VII
\textit{whence}
\[
(\tilde{C})_{*} = C_{*} \, (= TF), \qquad\quad (T\tilde{C})_{*} = (TC)_{*} \, (= F) \qquad\quad \text{\fnm{27}}.
\tag*{(5.6)}
\]
\footnotetext[27]{These relations do not depend on the hypothesis $W$ stable under direct factors.}
\textit{Moreover, if $W$ stable under direct factors (or if merely every direct factor of an arrow of $TC$ is in $W$), and if moreover $W$ is stable under transfinite composition, then $(\mathcal{M}, W, \tilde{C})$ is a special category with cofibrations. Moreover, $T\tilde{C}$ is left $Q$-saturated, hence} (by 5.5)
\[
T\tilde{C} \;=\; (TC)^{\natural} \; .
\tag*{(5.7)}
\]

\msp{page 85}

\textsc{Proof}. Let us prove
\[
T\tilde{C} \;\subset\; (TC)\cl \, ,
\]
and more precisely, that every element of $T\tilde{C}$ is a direct factor of an element of $TC$:
\[
T\tilde{C} \subset (TC)^{\natural} \; .
\]
\DERmark{ENV1073}

Let then $f \in T\tilde{C} \defeq \tilde{C}\cap W$, let us consider their factorization of $f$ as $pi$, with $i \in C$, $p \in TF$. Since $TF \subset W$, we have $p \in W$, hence $i \in W$, hence $i \in C \cap W \defeq TC$. Since $f \in \tilde{C}$, $p \in TF$, there exists $q : Y \longrightarrow X'$ such that $pq = \mathrm{id}_Y$, $qf = i$,
\DERmark{ENV1074}\DERmark{ENV1075}

\[
\begin{tikzcd}[column sep=small, row sep=small]
X \arrow[rr, "f\in \tilde{C}\cap W"] \arrow[dr, "i\in C"'] & & Y \arrow[dl, dotted, shift right, "q"'] \\
& X' \arrow[ur, shift right, "p\in TF"'] &
\end{tikzcd}
\]
\DERmark{ENV1076}\DERmark{ENV1077}

so $f$ appears as a direct factor of $i \in TC$, q.e.d.

Let us prove that $(\mathcal{M}, W, \tilde{C})$ is a special category with cofibrations (def. 4.4, page 74). Condition a) on $\mathcal{M}$ settled, condition b), the part on $W$, too. The stability of $\tilde{C}$ under cobase change and under composition OK, also (for b$'$) the stability under transfinite composition. Stability of $T\tilde{C} = \tilde{C} \cap W$ under cobase change: it remains to prove that if we have a cocartesian square
\[
\begin{tikzcd}[row sep=large, column sep=large]
X \arrow[r, "i"] \arrow[d] \arrow[dr, phantom, "\text{cocart.}"] & Y \arrow[d] \\
X' \arrow[r, "i'"] & Y'
\end{tikzcd} \; ,
\]
if $i \in \tilde{C}\cap W$, then $i' \in W$ (for we already know that $i' \in \tilde{C}$). But the preceding proof shows that it is a direct factor \textit{in} $X\backslash\mathcal{M}$ of an arrow in $TC$, so $i'$ is a direct factor of an arrow of $TC \subset W$. So if $W$ is stable under direct factors, $i' \in W$.
% foot: Version 9 février 2019

\printedpage{43}
% head: A. Grothendieck : Dérivateurs, Ch. XVII, Édition des Universités de Montpellier II et Paris VII
\msp{page 86}

This finishes proving b). For b$'$), it remains to see that $T\tilde{C}$ is stable under transfinite composition. I see no other way of seeing it, than by assuming that $W$ itself is stable under transfinite composition (which seems scarcely less demanding than the stability of $W$ under $\varinjlim$ \texttt{[filtering]}). This finishes proving b, b$'$). The condition c) on $C$ implies the same condition on $\tilde{C}$. Finally, since $\tilde{C}_{*} = C_{*}$, the condition d) is the same for $C$ and $\tilde{C}$. This finishes proving that $(\mathcal{M}, W, \tilde{C})$ is a special category with cofibrations. (\textbf{N.B.} we needed two additional hypotheses on $W$, stability under direct factors and under transfinite composition, to establish that $T\tilde{C}$ is stable under cobase change and under transfinite composition). The last assertion of 5.1 and the relation (5.7) which results from it, amounts to saying that $T\tilde{C} = \tilde{C} \cap W$ is stable under direct factors, which is obvious by the same property of $\tilde{C}$, and the hypothesis made on $W$ of stability under direct factors (or only that the direct factors of elements of $T\tilde{C}$ are in $W$, hence in $T\tilde{C}$). % [sic] Le dernière -> La dernière, gender agreement with "assertion"

\msp{page 87}

\textbf{Proposition 5.2.} \textit{Let $(\mathcal{M}, W, C)$ be a special category with cofibrations. Then with the notations of} (5.1)\textit{, we have}
\[
W \;=\; TF \circ TC,
\tag*{(5.8)}
\]
\[
TF \;=\; F \cap W.
\tag*{(5.9)}
\]

\textsc{Proof}. Let $f \in W$, let us write $f = pi$, with $i \in C$, $p \in TF$, hence $p \in W$, hence $i \in W$, hence $i \in C \cap W = TC$, whence (5.8).

Let us prove (5.9). We have $TF \subset F \cap W$, let us prove conversely that an $f \in F \cap W$ is in $TF$. Let us write $f = pi$, with $p \in TF$, $i \in TC$.
\[
\begin{tikzcd}[column sep=small, row sep=small]
X \arrow[rr, "f\in F\cap W"] \arrow[dr, "i\in TC"'] & & Y \\
& X' \arrow[ur, "p\in TF"'] &
\end{tikzcd}
\]
Since $f \in F$, $i \in TC$, and $TC \longleftrightarrow F$ by definition of $F$, $\exists\, r : X' \longrightarrow X$ such that $ri = \mathrm{id}_X$, $fr = p$. So $f$ is a direct factor of $p \in TF$, and since $TF = C_{*}$ is stable under direct factors, we conclude that $f \in TF$.

\textbf{Corollary 5.3.} \textit{Let $(\mathcal{M}, W, C)$ be a special category with cofibrations. For $(W, C, F)$, $(F \defeq (TC)_{*})$ to be a Quillen triple in $\mathcal{M}$, it is necessary and sufficient that the couple $(TC, F)$ satisfy the factorization axiom (i.e. that it be a Quillen precouple). For this Quillen structure to be closed, it is necessary and sufficient that $C$ be stable under direct factors.} % [sic] soit clos -> soit close, agreement with "cette structure"

\msp{page 88}

\textsc{Proof}. The necessity of the factorization axiom for $(TC, F)$, in order to have a Quillen triple, is tautological. For sufficiency, already having the two Quillen precouples $(C, TF)$, $(TC, F)$ with
% foot: Version 9 février 2019


\endgroup
% END UNIT der-ch17/der-ch17_p041-045.tex
% BEGIN UNIT der-ch17/der-ch17_p046-050.tex
\begingroup\makeatletter%
\newcommand{\printedpage}[1]{\marginpar{\footnotesize p.~#1}}%
\newcommand{\msp}[1]{\marginpar{\footnotesize [#1]}}%
\newcommand{\cl}{^{\tilde{}}}% tilde printed as exponent after a symbol or a parenthesis
\newcommand{\defeq}{\overset{\mathrm{def}}{=}}%
\newcommand{\encircled}[1]{\tikz[baseline=(X.base)]\node[draw,circle,inner sep=1pt](X){#1};}%
\newcommand{\encoval}[1]{\tikz[baseline=(X.base)]\node[draw,rounded corners=4pt,inner sep=2pt](X){#1};}%
\makeatother\let\printedpage\derpp\let\msp\dermsp%
\let\mspin\DERmspin%

\DERmark{ENV1078}


\printedpage{44}
% head: A. Grothendieck : Dérivateurs, Ch. XVII, Édition des Universités de Montpellier II et Paris VII
\[
\begin{tikzcd}
TC \arrow[r, hook] \arrow[d, leftrightarrow] & C \arrow[d, leftrightarrow] \\
F & TF \arrow[l, hook'] \ ,
\end{tikzcd}
\tag*{(5.2)}
\]
and $TC = C \cap W$, the fact that $(W, C, F)$ is a Quillen triple is corollary 4 p. 36 of XV: The conditions $TC = C \cap W$, $TF = F \cap W$, $W = TF \circ TC$ are equivalent, and equivalent to the fact that $(W, C, F)$ is a Quillen triple.
\DERmark{ENV1079}


By Quillen's definition of triples that are \textit{closed,} by the properties
\[
\underbrace{W = TF \circ TC}_{\substack{\text{true for all}\\ \text{the Quillen}\\ \text{triples}}}, \qquad C = (TF)^*, \qquad \underbrace{F = (TC)_*}_{\substack{\text{true in}\\ \text{the present case}}}\ , % [sic] le triples -> les triples, agreement with "tous"
\]
it remains to express that $C = (TF)^*$, which is equivalent to $C$ stable under direct factors (cf. loc. cit. cor. 1, p. 31), q.e.d.

\textbf{Remark 5.4.} The statements 5.2 and 5.3, and the relation $T\tilde{C} \subset (TC)^\natural$ in proposition 5.1, do not call upon the whole hypothesis that $(\mathcal{M}, W, C)$ be a special category with cofibrations, but only that we have a category with cofibrations, with $C \mathrel{\leftarrow\!\!\!+\!\!\!\rightarrow} \underbrace{TF}_{\defeq C_*} \subset W$, % ?? uncertain: arrow symbol between C and TF -- double-headed arrow crossed by a small bar/plus sign
\msp{page 89}
i.e. a category with $W$, $C$ satisfying Cof$_0$, Cof$_1$, Cof$_2$ of page 39, plus $TF \defeq C_* \subset W$, plus the factorization axiom for $(C, TF)$ (which implies Cof$_3$).

\textbf{Corollary 5.5.} \textit{Let $(\mathcal{M}, W, C)$ be a special category with cofibrations, such that there exists a} small \textit{subset $TC_0 \subset TC = C \cap W$, with}
\[
TC \ \subset \ TC_0\cl\ ,
\tag*{(5.5.1)}
\]
\textit{or what amounts to the same thing, $(TC)\cl = (TC_0)\cl$, i.e.}
\[
\underbrace{F}_{\defeq (TC)_*} \ = \ (TC_0)_*\ .
\tag*{(5.5.2)}
\]
\textit{Then $(W, C, F)$ is a Quillen triple. This triple is closed if and only if $C$ is stable under direct factors.}

Indeed, the factorization theorem in XV §1 implies that under the hypothesis made on $TC_0 \subset TC$, the pair $(TC, F)$ satisfies the factorization condition. Indeed, by conditions b, b$'$ of definition 4.4, $TC$ has the desired stability properties (contains the isomorphisms, is stable under cobase change and under transfinite composition).

\msp{page 90}
\textbf{Definition 5.6.} The Quillen triples deduced as in 5.5 from a special category with cofibrations are called \textit{special.} Thus a Quillen triple $(W, C, F)$ in $\mathcal{M}$ is special if it satisfies conditions a) to d) in definition (4.4), concerning $(W, C)$ exclusively, plus the following two conditions.
% foot: Version 9 février 2019

\printedpage{45}
% head: A. Grothendieck : Dérivateurs, Ch. XVII, Édition des Universités de Montpellier II et Paris VII
\begin{itemize}[leftmargin=3em]
\item[\encircled{c$'$}] $\exists$ small subset $TC_0$ of $TC$, such that $TC \subset (TC_0)\cl$.
\item[\encircled{e}] $F = (TC)_*$.
\end{itemize}

\textbf{Corollary 5.7.} \textit{Let $(W, C, F)$ be a Quillen triple that is} closed \textit{in $\mathcal{M}$. For this triple to be special, it is necessary and sufficient that it satisfy the following conditions.}
\begin{itemize}[leftmargin=3em]
\item[\encircled{a}] \textit{$\mathcal{M}$ stable under small $\varinjlim$\,, and its elements are accessible.}
\item[\encoval{c, c$'$}] \textit{There exists a small subset $C_0$ of $C$, $TC_0$ of $TC$, such that we have}
\[
C \ = \ C_0\cl\ , \qquad TC \ = \ (TC_0)\cl\ .
\tag*{(5.7.1)}
\]
\end{itemize}
Indeed, all the other conditions a) to e), including c$'$) of 5.6, are automatically satisfied.

\msp{page 91}
\textbf{Proposition 5.8.} \textit{Let us consider a pair of adjoint functors}
\[
\begin{tikzcd}
\mathcal{M} \arrow[r, shift left, "v_!"] & \mathcal{M}' \arrow[l, shift left, "v^*"] \ ,
\end{tikzcd}
\tag*{(5.8.1)}
\]
\textit{where $\mathcal{M}$ is equipped with a special closed Quillen triple $(W, C, F = (TC)_*)$, with $TC = C \cap W$} (cf. definition 5.6). \textit{We suppose the following conditions satisfied} (compare prop. 4.6, relations (4.6.1), and (4.7.4)):
\[
\left\{
\begin{array}{ll}
\text{a)} & \mathcal{M}' \ \textit{stable under small}\ \varinjlim\ \textit{, its elements are accessible.}\\[1.5ex]
\text{b)} & v^* \ \textit{commutes with small}\ \varinjlim\ .\\[1.5ex]
\text{c)} & v^*(v_!(C)) \subset C,\ v^*(v_!(TC)) \subset TC.
\end{array}
\right.
\tag*{(5.8.2)}
\]
\textit{Let us define $(W', C', F')$ in $\mathcal{M}'$} (compare prop. 4.6, notably 4.6.2) \textit{by}
\[
\left\{
\begin{array}{lcl}
W' & = & (v^*)^{-1}(W)\\[1ex]
C' & = & v_!(C)\cl\\[1ex]
F' & = & (TC')_*\ , \qquad \textit{where}\ \ TC' \ \defeq \ C' \cap W'\ .
\end{array}
\right.
\tag*{(5.8.3)}
\]
\textit{Then $(W', C', F')$ is a closed special Quillen triple in $\mathcal{M}'$, and moreover we have}
\[
F' \ = \ (v^*)^{-1}(F)\ , \qquad TF' \ = \ (v^*)^{-1}(TF)\ ,
\tag*{(5.8.4)}
\]
\[
TC' \ = \ (v_!(TC))\cl\ .
\tag*{(5.8.5)}
\]

\msp{page 92}
\textsc{Proof}. By 4.6, we already know that $(\mathcal{M}', W', C')$ is a special category with cofibrations. (\textbf{N.B.} What is denoted $C'$ here, was denoted $\tilde{C}'$ in proposition 4.6. We must use the last assertion of 4.6, where we assume $C$ stable under direct factors, OK since $(W, C, F)$ assumed closed.) To prove that $(W', C', F' \defeq (TC')_*)$ is a Quillen triple
\DERmark{ENV1080}

% foot: Version 9 février 2019
\printedpage{46}
% head: A. Grothendieck : Dérivateurs, Ch. XVII, Édition des Universités de Montpellier II et Paris VII
that is special, it remains, by 5.3 and definition 5.6, to prove that there exists a small subset $TC'_0$ of $TC'$, such that
\[
(\ TC'_0 \ \subset\ )\ TC' \ \subset \ (TC'_0)\cl\ .
\tag*{(5.8.6)}
\]
(This Quillen triple will moreover be closed, since $C'$ is Q-closed on the left by definition.) So let $TC_0$ be a small subset of $TC$ such that
\[
TC \ \subset \ (TC_0)\cl\ ,
\tag*{(5.8.7)}
\]
which exists by the hypothesis on $(W, C, F)$. We shall take
\[
TC'_0 \ \defeq \ v_!(TC_0)\ .
\tag*{(5.8.8)}
\]
We saw in 4.7 that we have
\[
v_!(TC) \ \subset \ TC'
\tag*{(5.8.9)}
\]
(follows from the hypothesis $v^*(v_!(TC)) \subset TC$ in (5.8.2. c)), a fortiori
\[
TC'_0 \ = \ v_!(TC_0) \ \subset \ TC'\ ,
\]
\msp{page 93}
it remains to prove
\[
TC' \ \overset{?}{\subset} \ (TC'_0)\cl\ ,
\tag*{(5.8.10)}
\]
which is equivalent to $(TC')\cl \subset (TC'_0)\cl$, or again to
\[
(TC'_0)_* \ \overset{?}{\subset} \ (TC')_* \ \defeq \ F'\ ,
\tag*{(5.8.10 bis)}
\]
whereas we have a priori the reverse inclusion (deduced from $TC'_0 \subset TC'$)
\[
F' \ \subset \ (TC'_0)_*\ .
\tag*{(5.8.11)}
\]
Now by adjunction, we have
\[
\begin{aligned}
(\ \underbrace{TC'_0}_{v_!(TC_0)}\ )_* \ &= \ (v_!(TC_0))_*\\
&= \ (v^*)^{-1}(\ \underbrace{(TC_0)_*}_{= (TC)_* \defeq F}\ )\\
&= \ (v^*)^{-1}(F)\ ,
\end{aligned}
\]
so (5.8.11) is written
\[
F' \ \subset \ (v^*)^{-1}(F)\ , \qquad \text{i.e.}\ \ v^*(F') \ \subset \ F
\tag*{(5.8.12)}
\]
(compare (4.7.6)), and we must show that this inclusion is an equality (on which I had gotten stuck p. 80, 81 \ldots), which is none other than the first of the two relations (5.8.4). But it is in the initial form (5.8.10) (transposed from (5.8.4)) that we are going to prove the thing, finally. Note that we have
% foot: Version 9 février 2019
\printedpage{47}
% head: A. Grothendieck : Dérivateurs, Ch. XVII, Édition des Universités de Montpellier II et Paris VII
\msp{page 94}
\[
(\underbrace{v_!(TC_0)}_{TC'_0})\cl \ = \ (v_!(TC))\cl\ ,
\]
since $TC$ is the closure of $TC_0$ under the operations of cobase change, of transfinite composition, and of direct factors, and $v_!$ commutes with these operations. So we must simply show
\[
TC' \ \subset \ (v_!(TC))\cl
\tag*{(5.8.12 bis)}
\]
(where there is no longer any question of $TC_0$). But the factorization theorem applies to the pair
\[
\Big( (v_!(TC))\cl,\ (v_!(TC))_* = (v^*)^{-1}(F) \Big)\ ,
\]
since $(v_!(TC))\cl = (TC'_0)\cl$, $TC'_0$ small. So if we have $f' \in TC' = C' \cap W'$, let us factor $f'$ as $p'i'$, with
\[
i' \ \in \ (v_!(TC))\cl\ , \qquad p' \ \in \ (v^*)^{-1}(F)\ .
\tag*{(5.8.13)}
\]
\[
\begin{tikzcd}[column sep=small, row sep=small]
X' \arrow[rr, "f' \in C' \cap W'"] \arrow[dr, "i'"'] & & Y' \arrow[dl, dotted, shift right, "q'"'] \\
& X'_1 \arrow[ur, shift right, "p'"'] &
\end{tikzcd}
\tag*{$(*)$}
\]
Note that we have
\msp{page 95}
\[
(v_!(TC))\cl \ \subset \ TC' \ (\ \defeq \ C' \cap W'\ )
\tag*{(5.8.14)}
\]
for $v_!(TC) \subset TC' \subset C'$, hence $(v_!(TC))\cl \subset \tilde{C}' = C'$, and it remains to prove that $v_!(TC)\cl \subset W'$, i.e. $v^*((v_!(TC))\cl) \subset W$. Now I say that we even have
\[
v^*((v_!(TC))\cl) \ \subset \ TC\ ,
\]
i.e.
\[
v_!(TC)\cl \ \subset \ (v^*)^{-1}(TC)\ .
\tag*{(5.8.15)}
\]
This follows from the fact that $TC$ is Q-closed, hence Q-saturated on the left, which implies, since $v^*$ commutes with small $\varinjlim$\,, that $(v^*)^{-1}(TC)$ is so as well. On the other hand, it contains $v_!(TC)$, hence also the left Q-saturation of $v_!(TC)$. But the latter is \textit{equal} to the Q-closure $(v_!(TC))\cl$, by the factorization theorem, seeing that $TC'_0 \subset v_!(TC) \subset (TC'_0)\cl$, with $TC'_0$ small. % [sic] (v^*)^{-1}(T) -> (v^*)^{-1}(TC), the set in question is TC
\DERmark{ENV1081}


Thus, (5.8.13) and (5.8.14) imply $i' \in TC' = C' \cap W'$, hence $i' \in W'$, whence $p' \in W'$, i.e. $v^*(p') \in W$, and since by hypothesis
% foot: Version 9 février 2019
\printedpage{48}
% head: A. Grothendieck : Dérivateurs, Ch. XVII, Édition des Universités de Montpellier II et Paris VII
\msp{page 96}
(5.8.13) we have $v^*(p') \in F$, we have $v^*(p') \in TF$, i.e.
\[
p' \ \in \ (v^*)^{-1}(TF) \ = \ TF' \ (\ \defeq \ C'_*\ )\ , % [sic] C' -> C'_*, cf. (5.8.16) TF' ( = C'_* )
\]
\DERmark{ENV1082}

the equality $(v^*)^{-1}(TF) = TF'$, i.e.
\[
(v^*)^{-1}(C_*) \ = \ (v_!(C))_*\ ,
\]
resulting from the adjunction of $v_!$, $v^*$. Altogether, we therefore have
\[
p' \ \in \ TF' \ (\ = \ C'_*\ )\ ,
\tag*{(5.8.16)}
\]
and since $f' \in C'$, diagram $(*)$, page 94, shows us that there exists
\DERmark{ENV1083}

\[
q' \ : \ Y' \longrightarrow X'_1\ ,
\]
with $p'q' = \mathrm{id}$, $q'f' = i'$. So $f'$ is a direct factor of $i' \in v_!(TC)\cl$, so we have $f' \in v_!(TC)\cl$, whence the relation
\[
TC' \ \subset \ v_!(TC)\cl\ ,
\]
which finishes proving that $(W', C', F')$ is a special and closed Quillen triple.

The two relations (5.8.4) were proved along the way (the second being essentially trivial, and the first representing the main difficulty in the proof). Remains to establish (5.8.5). But we already know (5.8.9) that
\[
v_!(TC) \ \subset \ TC'\ ,
\]
\msp{page 97}
whence
\[
v_!(TC)\cl \ \subset \ TC',
\]
since $TC'$ is Q-closed on the left (or because that was proved in (5.8.14)), we have already proved the inclusion in the reverse direction (5.8.12 bis), an assertion equivalent to (5.8.4) (and which represented the main difficulty).

\textbf{Remark 5.9.} If we do not assume that the structure on $\mathcal{M}$ is closed, then the proof given of 5.8 shows that we nevertheless find on $\mathcal{M}'$ a special Quillen structure, by defining $C'$ this time as in 4.6. But I doubt that we will ever have to work with special Quillen triples which are not also closed. In any case, if $W$ is stable under small filtering $\varinjlim$\footnote[28]{which implies stability under direct factors, and also implies the same property for $W' \defeq (v^*)^{-1}(W)$.} (and I doubt that one finds special Quillen triples outside this case), then if $(W, C, F)$ is a special Quillen triple, the same is true of $(W, \tilde{C} = C^\natural, F)$, and I presume that we have every interest in replacing the initial triple by this one. % footnote mark printed in the text as superscript (28)
% foot: Version 9 février 2019


\endgroup
% END UNIT der-ch17/der-ch17_p046-050.tex
% BEGIN UNIT der-ch17/der-ch17_p051-055.tex
\begingroup\makeatletter%
\newcommand{\printedpage}[1]{\marginpar{\footnotesize p.~#1}}%
\newcommand{\msp}[1]{\marginpar{\footnotesize [#1]}}%
\newcommand{\Dl}{\boldsymbol{\Delta}}%
\newcommand{\tl}{\tilde{}}%
\newcommand{\Ob}{\mathrm{Ob}}%
\newcommand{\Fl}{\mathrm{Fl}}%
\makeatother\let\printedpage\derpp\let\msp\dermsp%
\let\mspin\DERmspin%


\printedpage{49}
% head: A. Grothendieck : Dérivateurs, Ch. XVII, Édition des Universités de Montpellier II et Paris VII
% foot: Version 9 février 2019

\msp{page 98}

\section*{6\quad Application to the construction of derivators (associated with special Quillen triples)}

\textbf{Theorem 6.1.}\footnote[29]{This is the main result in this part XVII of the notes.} \textit{Let $(W,C,F)$ be a special closed Quillen triple on the category $\mathcal{M}$} (cf. def. 5.6, p. 90, and cor. 5.7). \textit{For every $I$ in} Cat\textit{, let us consider the category}
\[
\mathcal{M}(I) \;=\; \underline{\mathrm{Hom}}(I^{\circ},\mathcal{M}) \;, \tag*{(6.1.1)}
\]
\textit{and let us consider}
\[
W(I),\ F(I) \;\subset\; \Fl(\mathcal{M}(I)) \;,
\]
\textit{defined by}
\[
\begin{aligned}
&\text{\textit{$u : X \longrightarrow Y$ in $\mathcal{M}(I)$ is in $W(I)$ (resp. $F(I)$) if and only}}\\
&\text{\textit{if $\forall\, i \in \Ob I$, $u(i) : X(i) \longrightarrow Y(i)$ is in $W$ (resp. $F$).}}
\end{aligned}
\tag*{(6.1.2)}
\]

\begin{itemize}[leftmargin=2.5em]
\item[a)] \textit{There exists a closed Quillen structure (necessarily unique) in $\mathcal{M}(I)$, having $W(I)$ and $F(I)$ as the sets of quasi-isomorphisms and of fibrations respectively. We thus also have}
\[
\left\{
\begin{aligned}
TF(I) \;&=\; F(I)\cap W(I) \;=\; \{u\in \Fl(\mathcal{M}(I)) \mid u(i)\in TF \ \ \forall\, i\in \Ob I\},\\
C(I) \;&=\; TF(I)^{*}.
\end{aligned}
\right.
\tag*{(6.1.3)}
\]
\textit{This Quillen structure is} special.

\msp{page 99}
\item[b)] \textit{Let $I_0$ be the discrete category having the same set of objects as $I$, and}
\[
v_I \ \ \text{\textit{or}}\ \ v \;:\; I_0 \longrightarrow I
\tag*{(6.1.4)}
\]
\textit{the canonical functor. Let us denote by}
\[
C(I) \;=\; (TF(I))^{*}, \qquad TC(I) \;=\; C(I)\cap W(I) \;=\; F(I)^{*}
\tag*{(6.1.5)}
\]
\textit{the sets of cofibrations resp. of trivial cofibrations in $\mathcal{M}(I)$, and let us use the same notations $C(I_0)$, $TC(I_0)$ for $I_0$. We then have}
\[
\begin{aligned}
&\text{\textit{$u : X_0 \longrightarrow Y_0$ in $\mathcal{M}(I_0)$ is in $C_0$ (resp. $TC_0$) if and only if}}\\
&\text{\textit{for every $i$ in $\Ob I$, $u(i) : X_0(i) \longrightarrow Y_0(i)$ is in $C$ (resp. $TC$).}}
\end{aligned}
\tag*{(6.1.6)}
\]
\textit{Finally we have}
\[
\left\{
\begin{aligned}
C(I) \;&=\; (v_!(C(I_0)))\tl\\
TC(I) \;&=\; (v_!(TC(I_0)))\tl\,.
\end{aligned}
\right.
\tag*{(6.1.7)} % ?? uncertain: shape of the small superscript mark after the closing parenthesis -- tilde (\tilde{}) or breve (\breve{})
\]

\item[c)] \textit{Let}
\[
u \;:\; I \longrightarrow I'
\]
\textit{be an arrow in} Cat\textit{, let us consider the couple of adjoint functors}
\[
\begin{tikzcd}
\mathcal{M}(I) \arrow[r, shift left, "u_!"] & \mathcal{M}(I') \arrow[l, shift left, "u^{*}"]
\end{tikzcd}
\;.
\tag*{(6.1.8)}
\]

\printedpage{50}
% head: A. Grothendieck : Dérivateurs, Ch. XVII, Édition des Universités de Montpellier II et Paris VII
% foot: Version 9 février 2019

\msp{page 100}
\textit{This couple satisfies Quillen's conditions for couples of adjoint functors between Quillen categories} (Hom. Alg. Chp. I, §4, th. 3, p. 4.5), \textit{so that they define by passage to the localized categories a couple of adjoint functors}
\[
\begin{tikzcd}
\mathrm{Ho}_W(I) \arrow[r, shift left, "u_!^{W}"] & \mathrm{Ho}_W(I') \arrow[l, shift left, "u^{*}_{W}"]
\end{tikzcd}
\;.
\tag*{(6.1.9)}
\]
\textit{We have, more precisely, the relations}
\[
\left\{
\begin{aligned}
u_!(C(I)) \;&\subset\; C(I')\\
u_!(TC(I)) \;&\subset\; TC(I')\;,
\end{aligned}
\right.
\tag*{(6.1.10)}
\]
\textit{whence}
\[
(u_!(\mathcal{M}(I)_{\mathrm{c}})) \;\subset\; \mathcal{M}(I')_{\mathrm{c}}\;,
\tag*{(6.1.11)}
\]
\textit{and moreover}
\[
\begin{aligned}
&\text{\textit{for every arrow $f$ in $\mathcal{M}(I)_{\mathrm{c}}$, such that $f \in W(I)$, we have}}\\
&\text{\textit{$u_!(f) \in W(I')$, i.e. $(u_!)_{\mathrm{c}} : \mathcal{M}(I)_{\mathrm{c}} \longrightarrow \mathcal{M}(I')_{\mathrm{c}}$ is compatible}}\\
&\text{\textit{with the localizers.}}
\end{aligned}
\tag*{(6.1.12)}
\]
\textit{On the other hand we have}
\[
\left\{
\begin{aligned}
u^{*}(F(I')) \;&\subset\; F(I)\\
u^{*}(TF(I')) \;&\subset\; TF(I)\\
u^{*}(W(I')) \;&\subset\; W(I)\;,
\end{aligned}
\right.
\tag*{(6.1.13)}
\]
\textit{hence}
\[
\left\{
\begin{aligned}
&u^{*}(\mathcal{M}(I')_{\mathrm{f}}) \;\subset\; \mathcal{M}(I)_{\mathrm{f}}\\
&u^{*}_{\mathrm{f}} : \mathcal{M}(I')_{\mathrm{f}} \longrightarrow \mathcal{M}(I)_{\mathrm{f}} \text{ \textit{is compatible with the localizers.}}
\end{aligned}
\right.
\tag*{(6.1.14)}
\]

\msp{page 101}
\item[d)] \textit{Suppose that $I$ is an ordered category, and let us consider the set $C_I \subset \Fl(\mathcal{M}(I))$ defined in §2}\footnote[30]{In paragraph 2, this was denoted $C_{I^{\circ}}$ instead of $C_I$.} (cf. page 25, cf. also th. 2.1, p. 26, and its corollaries). \textit{We then have}
\[
\left\{
\begin{aligned}
C_I \;&\supset\; C(I)\\
TC_I \;&\supset\; TC(I)\;,
\end{aligned}
\right.
\tag*{(6.1.15)}
\]
\textit{and we have equality when $I$ is an ordered set that is} Noetherian.
\DERmark{ENV1084}

\end{itemize}

\textsc{Proof}. Assertion a) is a special case of 5.8, applied to the couple of adjoint functors
\[
\begin{tikzcd}
\mathcal{M}(I_0) \arrow[r, shift left, "v_!"] & \mathcal{M}(I) \arrow[l, shift left, "v^{*}"]
\end{tikzcd}
\;.
\]
The conditions (5.8.2) are satisfied, this was seen in (4.8). Proposition 5.8 then tells us of the existence of a closed Quillen triple structure on $\mathcal{M}(I)$, whose set of quasi-isomorphisms $W(I)$ and set of fibrations $F(I)$, are given by the first formula (5.8.3) and by the first formula (5.8.4), which is equivalent to the definitions (6.1.2). (\textbf{N.B.} The fact that $\mathcal{M}(I_0)$, with the obvious structure $W(I_0)$, $C(I_0)$, $F(I_0)$ (product), is
\printedpage{51}
% head: A. Grothendieck : Dérivateurs, Ch. XVII, Édition des Universités de Montpellier II et Paris VII
% foot: Version 9 février 2019
\msp{page 102}
indeed a special closed Quillen category, is $\pm$ tautological.)

\textit{Let us prove} b). The relation (6.1.6) is tautological. The relations (6.1.7) are contained in 5.8 (cf. second relation in (5.8.3) for $C' = v_!(C)\tl$, i.e. $C(I) = (v_!(C(I_0)))\tl$, and (5.8.5) for $TC(I) = (v_!(TC(I_0)))\tl$).

\textit{Let us prove} c). The first assertion is, by \textsc{Quillen}, a special case of the relations (6.1.10) to (6.1.14), it is these that have to be established. The relations (6.1.13) and consequently (6.1.14) are tautological by virtue of the definitions, so it is a matter of establishing \textit{(6.1.10)} and (6.1.12).

For the formulas (6.1.10), I use the commutative diagram
\[
\begin{tikzcd}[row sep=large, column sep=large]
I \arrow[d, "u"'] & I_0 \arrow[l, "v"'] \arrow[d, "u_0"] \\
I' & I'_0\;, \arrow[l, "v'"']
\end{tikzcd}
\]
and I note that the functors $u_!^{\mathcal{M}}$, $u_{0\,!}^{\mathcal{M}}$, commuting with small direct limits,
\msp{page 103}
behave well with respect to the operation
\[
\Phi \longmapsto \Phi\tl\,,
\]
when $\Phi$ is a subset of $\Fl(\mathcal{M}(I))$ resp. of $\Fl(\mathcal{M}(I_0))$ which is Q-accessible, i.e. such that there exists a \textit{small} subset $\Phi_0 \subset \Phi$ with $\Phi \subset \tilde{\Phi}_0$, more precisely, we have
\[
u_!(\Phi\tl) \;\subset\; (u_!(\Phi))\tl\,.
\]
This noted, we have by b)
\[
C(I) \;=\; (v_!(C(I_0)))\tl\,,
\]
whence
\[
u_!(C(I)) \;\subset\; (\underbrace{u_!(v_!(C(I_0)))}_{(uv)_!(C(I_0))})\tl\,,
\]
and the right-hand side can also be written
\[
(v'_!u_{0\,!}(C(I_0)))\tl \;\overset{?}{\subset}\; (v'_!(C(I'_0)))^{\sim} \;=\; C(I')\;. % [sic] (v'_!(I'_0)) -> (v'_!(C(I'_0))), the C( ) around I'_0 apparently omitted
\]
\DERmark{ENV1085}

Thus everything comes down to showing that
\[
u_{0\,!}(C(I_0)) \;\subset\; C(I'_0)\;,
\tag*{(6.1.16)}
\]
which follows at once from the computation of $u_{0\,!}$ by means of direct sums, and from the fact that $C$ is stable under arbitrary direct sums. This thus proves $u_!(C(I)) \subset C(I')$, and we prove in the same way $u_!(TC(I)) \subset TC(I')$, which reduces in the same way to the special case
\[
u_{0\,!}(TC(I_0)) \;\subset\; TC(I'_0)\;.
\tag*{(6.1.17)}
\]

\printedpage{52}
% head: A. Grothendieck : Dérivateurs, Ch. XVII, Édition des Universités de Montpellier II et Paris VII
% foot: Version 9 février 2019

\msp{page 104}

It remains to prove (6.1.12), and we proceed as for 3.3 (p. 54), using the fact that an arrow $f$ in $\mathcal{M}(I)_{\mathrm{c}}$ which is in $W(I)$, can be written in $\mathcal{M}(I)_{\mathrm{c}}$ in the form $f = pi$, where $i \in TC(I)$, and where $p$ is a left inverse (retraction) of an arrow $j \in TC(I)$. We will therefore have $u_!(f) = u_!(p)u_!(i)$, and $u_!(i) \in TC(I')$ by (6.1.10), and $u_!(p)$ is a left inverse of $u_!(j) \in TC(I')$. We then have $u_!(i), u_!(j), u_!(p) \in W(I')$, whence $u_!(f) \in W(I')$, q.e.d.

Let us finally prove d). In view of the definition of $C_I$, to establish $C(I) \subset C_I$, it suffices to see this:

\textbf{Corollary 6.2.}\footnote[31]{\textbf{N.B.} This result generalizes, under the general conditions in 6.2.5.} \textit{Let $I$ be arbitrary in} Cat\textit{, let $U \subset V$ be two closed sets in $I$ (i.e. open sets in $I^{\circ}$), and let $f : X_{\bullet} \longrightarrow Y_{\bullet}$ be an arrow in $C(I)$ (resp. in $TC(I)$). Let us also put}
\msp{page 105}
\[
X_U \;=\; \varinjlim_{i\in U} X(i), \qquad X_V \;=\; \varinjlim_{i\in V} X(i)\;, % [sic] second index i\in U -> i\in V, X_V should be the limit over V
\]
\DERmark{ENV1086}

\textit{and likewise for $Y$, and let us consider the diagram}
\[
\begin{tikzcd}[row sep=large, column sep=large]
X_U \arrow[r, "\text{trans.}"] \arrow[d, "u_U"'] \arrow[dr, phantom, "\text{cocart.}"] & X_V \arrow[d] \arrow[dd, dotted, bend left=60, "u_V"] \\
Y_U \arrow[r] \arrow[dr, dotted, "\text{trans.}"'] & Y_{U,V} \arrow[d, "u_{U,V}"'] \\
 & Y_V\;.
\end{tikzcd}
\tag*{(6.2.1)}
\]
\textit{Then the arrow $u_{U,V}$ is in $C$ (resp. $TC$).}

\textsc{Proof} of 6.2. Let us admit that the restriction of an arrow of $C(I)$ resp. $TC(I)$ to a closed set $V$ of $I$ is in $C(V)$ resp. $TC(V)$ (cf. \dots{} below). This allows us to assume $V = I$. Giving $U$, which is an open set of $I^{\circ}$, is equivalent to giving an arrow $v : I^{\circ} \longrightarrow \Dl^1$ (taking the fiber at 0), or again to giving an arrow $u : I \longrightarrow (\Dl^1)^{\circ} \simeq \Dl^1$ (taking the fiber at 1). Now % [sic] le fibre -> la fibre, gender (cf. "la fibre en 0" just before)
\[
u \;:\; I \longrightarrow \Dl^1\;, \qquad U \;=\; I_1 \overset{\mathrm{def}}{=} u^{-1}(\{1\})\;.
\tag*{(6.2.2)}
\]
\DERmark{ENV1087}


\msp{page 106}
We then have
\[
X_U \;=\; u_!(X)(1), \quad X_I \;=\; u_!(X)(0)\;,
\]
and the transition morphism $X_U \longrightarrow X_I$ is deduced from the arrow $0 \xrightarrow{\ \partial\ } 1$ in $\Dl^1$, and likewise for $Y$. Let
\[
\xi \;=\; u_!(X), \qquad \eta \;=\; u_!(Y)\;,
\]
\printedpage{53}
% head: A. Grothendieck : Dérivateurs, Ch. XVII, Édition des Universités de Montpellier II et Paris VII
% foot: Version 9 février 2019
so the diagram (6.2.1) is identified with the diagram deduced from
\[
\varphi \;:\; \xi \longrightarrow \eta\;, \qquad \varphi \;=\; u_!(f)\;,
\]
as
\[
\begin{tikzcd}[row sep=large, column sep=large]
\xi(1) \arrow[r, "\xi(\partial)"] \arrow[d, "\varphi(1)"'] \arrow[dr, phantom, "\text{cocart.}"] & \xi(0) \arrow[d] \arrow[dd, dotted, bend left=60, "\varphi(0)"] \\
\eta(1) \arrow[r] \arrow[dr, dotted, "\eta(\partial)"'] & \tilde{\eta} \arrow[d, "\tilde{\varphi}"'] \\
 & \eta(0) \qquad .
\end{tikzcd}
\tag*{(6.2.3)}
\]
\DERmark{ENV1088}

Now by 6.1.10 applied to $u : I \longrightarrow \Dl^1$, the relation $f \in C(I)$ resp. $f \in TC(I)$ implies $\varphi \in C(\Dl^1)$ resp. $\varphi \in TC(\Dl^1)$. Thus to establish 6.1, we are reduced to the special case where $I = \Dl^1$, $U = \{1\}$, $V = I$, hence to proving the

\textbf{Lemma 6.2.4.} \textit{Let $\varphi : \xi \longrightarrow \eta$ be an arrow in $\mathcal{M}(\Dl^1)$. If $\varphi \in C(\Dl^1)$ (resp.}
\msp{page 107}
\textit{$\varphi \in TC(\Dl^1)$), then the arrow $\tilde{\varphi}$ in the diagram} (6.2.3) \textit{is in $C$ (resp. $TC$).}

More generally, and more precisely:

\textbf{Lemma 6.2.5.} \textit{Let $\mathcal{M}$ be a category stable under amalgamated sums, and $\Phi$ a subset Q-closed on the left} (we will take $\Phi = C$, and $\Phi = TC$, in the case of lemma 6.2.4). \textit{Let $\Psi = \Phi_{*}$, and let us consider, for every $I$ in} Cat\textit{, $\Phi(I)$, $\Psi\langle I\rangle$ in $\mathcal{M}(I)$, defined by:}
\[
\begin{aligned}
&\text{\textit{$u : X \longrightarrow Y$ in $\mathcal{M}(I)$ is in $\Psi\langle I\rangle$ if and only if $\forall\, i \in \Ob I$,}}\\
&\text{\textit{$u(i) : X(i) \longrightarrow Y(i)$ is in $\Psi$,}}
\end{aligned}
\tag*{(6.2.5.1)}
\]
\[
\Phi(I) \;=\; \Psi\langle I\rangle^{*}.
\tag*{(6.2.5.2)}
\]

Let us take the case $I = \Dl^1 = \{0 \xrightarrow{\ \partial\ } 1\}$, and let $u : X \longrightarrow Y$ be an arrow in $\mathcal{M}(\Dl^1)$, corresponding to a commutative square
\[
\begin{tikzcd}[row sep=large, column sep=large]
X(1) \arrow[r, "X(\partial)"] \arrow[d, "u(1)"'] & X(0) \arrow[d, "u(0)"] \\
Y(1) \arrow[r, "Y(\partial)"] & Y(0)\;.
\end{tikzcd}
\tag*{(6.2.5.3)}
\]

Let us consider the corresponding diagram


\endgroup
% END UNIT der-ch17/der-ch17_p051-055.tex
% BEGIN UNIT der-ch17/der-ch17_p056-060.tex
\begingroup\makeatletter%
\newcommand{\printedpage}[1]{\marginpar{\footnotesize p.~#1}}%
\newcommand{\msp}[1]{\marginpar{\footnotesize [#1]}}%
\makeatother\let\printedpage\derpp\let\msp\dermsp%
\let\mspin\DERmspin%


\noindent\printedpage{54}\msp{page 108}
% head: A. Grothendieck : Dérivateurs, Ch. XVII, Édition des Universités de Montpellier II et Paris VII
\[
\begin{tikzcd}
X(1) \arrow[r, "X(\partial)"] \arrow[d, "u(1)"'] \arrow[dr, phantom, "\text{cocart.}"] & X(0) \arrow[d] \arrow[dd, dotted, bend left=40, "u(0)"] \\
Y(1) \arrow[r] \arrow[dr, dotted, "Y(\partial)"'] & \tilde{Y} \arrow[d, "\tilde{u}"'] \\
 & Y(0)
\end{tikzcd}
\qquad .
\tag*{(6.2.5.4)}
\]

For us to have $u \in \Phi(\boldsymbol{\Delta}^1)$, it is necessary and sufficient that we have
\[
u(0),\ u(1),\ \tilde{u} \ \in\ \Phi\ .
\tag*{(6.2.5.5)}
\]

I have already checked this (in dual form), but can no longer find the place (perhaps I did not take notes). I don't feel like looking for the proof here.

To finish proving 6.1. d), i.e. the equality $C(I) = C_I$ for $I$ noetherian, I feel like proving something still more general, as follows: % [sic] plus générale -> plus général, agreement with « quelque chose »
\DERmark{ENV1089}


\textbf{Corollary 6.3.} \textit{Let $(\mathcal{M}, \Phi, \Psi)$ and $I \longmapsto \Phi(I) \subset \mathrm{Fl}(\mathcal{M}(I))$ be as in \textup{6.2.5.} Let $I$ be an ordered set, and let us also consider $\Phi_I$ defined (in the dual situation) on page \textup{25 (2.2)} (but with dual notation \dots). I say that we have}
\[
\Phi_I \ \supset\ \Phi(I)\ ,
\tag*{(6.3.1)}
\]
\msp{page 109}\textit{and we have equality if $I$ is noetherian, i.e. $I^{\circ}$ is artinian.}

\textsc{Proof}. The proof of $\Phi(I) \subset \Phi_I$ is the one we have just given. It remains to prove that if $I$ is noetherian, we have $\Phi_I \subset \Phi(I)$. Given $i : A \longrightarrow B$ in $\Phi_I$, we must prove that $i \in \Phi(I)$, % [sic] u -> i, the arrow in question is i : A -> B
\DERmark{ENV1090}

i.e. that for a commutative diagram in $\mathcal{M}(I)$,
\[
\begin{tikzcd}
A \arrow[r, "\alpha"] \arrow[d, "i"'] & X \arrow[d, "p"] \\
B \arrow[r, "\beta"'] \arrow[ur, dotted, "h"] & Y
\end{tikzcd}
\tag*{(6.3.2)}
\]
with $p \in \Psi\langle I\rangle$, i.e. $p(i) : X(i) \longrightarrow Y(i)$ in $\Psi$ for every $i$, $\exists\ h : B \longrightarrow X$ leaving the diagram commutative. We are going to return to the dual notation, working in $\mathcal{M}(I^{\circ})$, with $I$ artinian. We consider the set $\mathcal{H}$ of pairs % [sic] des couple -> des couples, plural
$(U, h_U)$ of an open set $U$ of $I$ and an $h_U : B|U \longrightarrow X|U$, making the restriction to $U$ of (6.3.2) commutative. % [sic] commutatif -> commutative, agreement with « la restriction »
We order $\mathcal{H}$ by extension, we note that every totally ordered subset $(U_\alpha, h_{U_\alpha})$ in $\mathcal{H}$ has a least upper bound, namely $U = \bigcup_\alpha U_\alpha$, $h_U$ defined by the $h_{U_\alpha}$ on the $U_\alpha$.
% foot: Version 9 février 2019

\printedpage{55}\msp{page 110}
% head: A. Grothendieck : Dérivateurs, Ch. XVII, Édition des Universités de Montpellier II et Paris VII
(\textbf{N.B.} The set $\mathcal{H}$ is not empty, for $(\emptyset, h_{\emptyset})$ is in $\mathcal{H}$ -- but this is contained in the preceding statement that every totally ordered subset of $\mathcal{H}$ has a least upper bound in it \dots). By Zorn, there exists a maximal element $(U, h_U)$. To prove that $U = I$, and we will be done. Or again that if $(U, h_U)$ such that $U \neq I$, it is \textit{not} maximal. Let $F = I\backslash U \neq \emptyset$ be the complementary closed set. By the artinian hypothesis on $I$, $F$ has a minimal element $i_0$. Let
\[
U' \ =\ U \cup \overbrace{I/i_0}^{V} \ =\ U \cup \{i_0\}\ ,
\]
we are going to extend $h_U$ to $h_{U'}$ on $U'$, we will be done. For this, let $V_0 = V \cap U = V\backslash\{i_0\}$, we are reduced to extending the restriction $h_{V_0}$ of $h_U$ to $V_0$, from $V_0$ to $V$. All told, we may suppose $I = I/i_0$ equipped with a final element $i_0$, $U = I\backslash\{i_0\}$, and show that a $(U, h_U) \in \mathcal{H}$ extends to $I$. Let $A'$, $B'$, $X'$, $Y'$ be the restrictions of $A$, $B$, $X$, $Y$ to $U$ (which we may also denote by $I'$). The datum of $A$ % [sic] La donné -> La donnée, missing feminine ending
is equivalent to that of $A'$, that of $A_0 = A(i_0) \in \mathrm{Ob}\,\mathcal{M}$, and of an arrow from
\[
A'_{I'} \ \overset{\mathrm{def}}{=}\ \varinjlim\limits_{i' \in I'} A'(i') \ \xrightarrow{\ \lambda_A\ }\ A_0\ (\ \simeq A_I\ )\ ,
\]
\DERmark{ENV1091}

\msp{page 111}and likewise for $B$, $X$, $Y$. More precisely, the category $\mathcal{M}(I^{\circ})$ can be interpreted as that of triples $X = (X', X_0, \lambda_X)$, where $X' \in \mathcal{M}(I'^{\circ})$, $X_0 \in \mathcal{M}$, and $\lambda_X : X'_{I'} \longrightarrow X_0$. In these terms, the datum of the diagram (6.3.2) is equivalent to that of a similar square diagram
\[
\begin{tikzcd}
A' \arrow[r, "\alpha'"] \arrow[d, "i'"'] & X' \arrow[d, "p'"] \\
B' \arrow[r, "\beta'"'] \arrow[ur, dotted, "h'"] & Y'
\end{tikzcd}
\tag*{(6.3.3)}
\]
\textit{in} $\mathcal{M}(I'^{\circ})$, whence a diagram in $\mathcal{M}$
\[
\begin{tikzcd}
A'_{I'} \arrow[r, "\alpha'_{I'}"] \arrow[d, "i'_{I'}"'] & X'_{I'} \arrow[d, "p'_{I'}"] \\
B'_{I'} \arrow[r, "\beta'_{I'}"'] \arrow[ur, dotted, "h'_{I'}"] & Y'_{I'}\ ,
\end{tikzcd}
\tag*{(6.3.4)}
\]
of another square diagram \textit{in} $\mathcal{M}$
\[
\begin{tikzcd}
A_0 \arrow[r, "\alpha_0"] \arrow[d, "i_0"'] & X_0 \arrow[d, "p_0"] \\
B_0 \arrow[r, "\beta_0"'] & Y_0\ ,
\end{tikzcd}
\tag*{(6.3.5)}
\]
and finally of a homomorphism $\lambda = (\lambda_A, \lambda_B, \lambda_X, \lambda_Y)$ from the square (6.3.4) to the square (6.3.5):
% foot: Version 9 février 2019
\printedpage{56}\msp{page 112}
% head: A. Grothendieck : Dérivateurs, Ch. XVII, Édition des Universités de Montpellier II et Paris VII
\[
\begin{tikzcd}
A'_{I'} \arrow[rr, "\alpha'_{I'}"] \arrow[dr, "i'_{I'}"'] \arrow[dd, "\lambda_A"'] & & X'_{I'} \arrow[dr, "p'_{I'}"] \arrow[dd, "\lambda_X" near start] & \\
 & B'_{I'} \arrow[ur, dotted, "h'_{I'}"] \arrow[rr, crossing over, "\beta'_{I'}" near end] & & Y'_{I'} \arrow[dd, "\lambda_Y"] \\
A_0 \arrow[rr, "\alpha_0" near start] \arrow[dr, "i_0"'] & & X_0 \arrow[dr, "p_0"] & \\
 & B_0 \arrow[from=2-2, crossing over, "\lambda_B"' near start] \arrow[ur, dotted, "h_0"'] \arrow[rr, "\beta_0"] & & Y_0
\end{tikzcd}
\qquad\qquad \text{\footnotemark[32]}
\tag*{(6.3.6)}
\]
\footnotetext[32]{\textbf{N.B.} The condition $p \in \Psi\langle I\rangle$ is expressed by $p' \in \Psi\langle I'\rangle$ and \fbox{$p_0 \in \Psi.$}}

On the other hand, the question of extending $(I', h') \in \mathcal{H}$ to $(I, h)$, amounts to that of finding $h_0 : B_0 \longrightarrow X_0$ in $\mathcal{M}$, so that together with $h'_{I'}$ this again makes a commutative diagram of (6.3.6), which amounts to saying that on the one hand the two triangles formed in the base square are commutative:
\[
p_0h_0 \ =\ \beta_0, \qquad\quad h_0i_0 \ =\ \alpha_0\ ,
\tag*{(6.3.7)}
\]
and on the other hand the diagonal square defined by $h'_{I'}$, $h_0$ is commutative (which expresses that we have an arrow $h$ from $B$ to $X$):
\[
\lambda_X h'_{I'} \ =\ h_0\lambda_B\ . % [sic] h'_{X'} -> h'_{I'}, the arrow is h'_{I'} everywhere else (cf. (6.3.16))
\tag*{(6.3.8)}
\]
\DERmark{ENV1092}


\msp{page 113}But the diagram (6.3.6), where the densely dotted arrow $h'_{I'}$ is already in hand, and $h_0$ in spaced dots remains to be found, gives rise to the diagram, in which is brought out the arrow $\tilde{i} \in \Phi$ which expresses the hypothesis $i \in \Phi_I$, % [sic] C(i) -> C(I), cf. i' \in C(I')
in view of $i' \in \Phi_{I'}$:
\DERmark{ENV1093}\DERmark{ENV1094}\DERmark{ENV1095}\DERmark{ENV1096}

\[
\begin{tikzcd}
A'_{I'} \arrow[rr, "\alpha'_{I'}"] \arrow[dr, "i'_{I'}"'] \arrow[dd, "\lambda_A"'] & & X'_{I'} \arrow[dr, "p'_{I'}"] \arrow[dd, "\lambda_X" near start] & \\
 & B'_{I'} \arrow[ur, dotted, "h'_{I'}"] \arrow[rr, crossing over, "\beta'_{I'}" near end] & & Y'_{I'} \arrow[dd, "\lambda_Y"] \\
A_0 \arrow[ur, phantom, "\text{cocart.}"] \arrow[rr, "\alpha_0" near start] \arrow[dr, "r"'] \arrow[ddr, dotted, "i_0"'] & & X_0 \arrow[dr, "p_0\in\Psi"] & \\
 & \tilde{B} \arrow[from=2-2, crossing over, "s"' near start] \arrow[ur, dotted, "\tilde{h}_0"] \arrow[rr, dotted, "\tilde{\beta}_0"] \arrow[d, "\tilde{i}\in\Phi"'] & & Y_0 \\
 & B_0\ , \arrow[uur, dotted, "h_0"'] \arrow[urr, "\beta_0"'] & &
\end{tikzcd}
\tag*{(6.3.9)}
\]
% foot: Version 9 février 2019
\printedpage{57}
% head: A. Grothendieck : Dérivateurs, Ch. XVII, Édition des Universités de Montpellier II et Paris VII
where the densely dotted arrow $\tilde{h}_0$ from $\tilde{B} \overset{\mathrm{def}}{=} A_0 \amalg_{A'_{I'}} B'_{I'}$ to $X_0$, is defined by the commutativity conditions
\DERmark{ENV1097}

\[
\tilde{h}_0r \ =\ \alpha_0, \qquad\quad \tilde{h}_0s \ =\ \lambda_X h'_{I'}\ ,
\tag*{(6.3.10)}
\]
which is legitimate, seeing that
\[
\begin{aligned}
\alpha_0\lambda_A \ &=\ (\lambda_X h'_{I'})i'_{I'} \\
&=\ \lambda_X(h'_{I'}i'_{I'}) \\
&=\ \lambda_X\alpha'_{I'}.
\end{aligned}
\]

\msp{page 114}As for the arrow $\tilde{\beta}_0 : \tilde{B} \longrightarrow Y_0$, it is given by the commutativity condition
\[
\tilde{\beta}_0r \ =\ p_0\alpha_0, \qquad\quad \tilde{\beta}_0s \ =\ \lambda_Y\beta'_{I'}\ ,
\tag*{(6.3.11)}
\]
which is again legitimate, since
\[
(p_0\alpha_0)\lambda_A \ =\ (\lambda_Y\beta'_{I'})i'_{I'}
\]
by virtue of the commutativity of the cube (6.3.6). To establish the full commutativity of the diagram (6.3.9) (solid lines, and densely dotted lines for $h'_{I'}$, $\tilde{h}_0$ and $\tilde{\beta}_0$ -- excluding $h_0$ not yet constructed), we must still verify on the one hand the relations % [sic] les relation -> les relations, plural
\[
\left\{
\begin{aligned}
\tilde{h}_0r \ &=\ \alpha_0, \quad p_0\tilde{h}_0 \ =\ \tilde{\beta}_0, \\
\lambda_X h'_{I'} \ &=\ \tilde{h}_0s
\end{aligned}
\right.
\qquad\qquad \text{\footnotemark[33]},
\tag*{(6.3.12)}
\]
\footnotetext[33]{\textbf{N.B.} the first and the third of these relations are already known, cf. (6.3.10).}
establishing the commutativity of the upper cube with the diagonal square (diagram of type (6.3.6), the dotted arrow $h_0$ included in it), on the other hand the relations
\[
\left\{
\begin{aligned}
\tilde{i}r \ &=\ i_0 \qquad \text{(and $\tilde{i}s = \lambda_B$ for the record)} \\
\beta_0\tilde{i} \ &=\ \tilde{\beta}_0\ ,
\end{aligned}
\right.
\tag*{(6.3.13)}
\]
where the first line is already known by the definition of $\tilde{B}$ and of $\tilde{i}$. The second
\msp{page 115}reduces to the two relations
\[
\begin{aligned}
(\beta_0\tilde{i})r \ &=\ \beta_0(\tilde{i}r) \\
&=\ \beta_0i_0 \\
&=\ p_0\alpha_0 \\
&=\ \tilde{\beta}_0r\ ,
\end{aligned}
\]
OK by commutativity of the base square of (6.3.5),
\[
\begin{aligned}
(\beta_0\tilde{i})s \ &=\ \beta_0(\tilde{i}s) \\
&=\ \beta_0\lambda_B \\
&=\ \lambda_Y\beta'_{I'} \\
&=\ \tilde{\beta}_0s\ ,
\end{aligned}
\]
% foot: Version 9 février 2019
\printedpage{58}
% head: A. Grothendieck : Dérivateurs, Ch. XVII, Édition des Universités de Montpellier II et Paris VII
OK by commutativity of the front square in (6.3.6).

So it remains to verify the second relation (6.3.12),
\[
p_0\tilde{h}_0 \ =\ \tilde{\beta}_0\ ,
\]
which reduces to
\[
\begin{aligned}
(p_0\tilde{h}_0)r \ &=\ p_0(\tilde{h}_0r) \\
&=\ p_0\alpha_0 \\
&=\ \tilde{\beta}_0r\ ,
\end{aligned}
\]
OK by (6.3.11),
\[
\begin{aligned}
(p_0\tilde{h}_0)s \ &=\ p_0(\tilde{h}_0s) \\
&=\ p_0(\lambda_X h'_{I'}) \\
&=\ (p_0\lambda_X)h'_{I'} \\
&=\ (\lambda_Y p'_{I'})h'_{I'} \\
&=\ \lambda_Y(p'_{I'}h'_{I'}) \\
&=\ \lambda_Y\beta'_{I'} \\
&=\ \tilde{\beta}_0s\ ,
\end{aligned}
\]
OK.

This being settled, I say that the constraints (6.3.7), (6.3.8) on $h_0$ amount to the commutativity of the two triangles in the square % [sic] les carré -> le carré, singular (one square (6.3.14))
\[
\begin{tikzcd}
\tilde{B} \arrow[r, "\tilde{h}_0"] \arrow[d, "\tilde{i}\in\Phi"'] & X_0 \arrow[d, "p_0\in\Psi"] \\
B_0 \arrow[r, "\beta_0"'] \arrow[ur, dotted, "h_0"] & Y_0\ ,
\end{tikzcd}
\tag*{(6.3.14)}
\]
\msp{page 116}i.e. to the two relations
\[
h_0\tilde{i} \ =\ \tilde{h}_0\ , \qquad\quad p_0h_0 \ =\ \beta_0\ .
\tag*{(6.3.15)}
\]

The second of these relations is none other than the first relation (6.3.7), and it remains to see that the first relation (6.3.15) is equivalent to the two remaining relations in (6.3.7), (6.3.8) taken together:
\[
h_0i_0 \ =\ \alpha_0\ , \qquad\quad \lambda_X h'_{I'} \ =\ h_0\underbrace{\lambda_B}_{\tilde{i}s}
\tag*{(6.3.16)}
\]
% foot: Version 9 février 2019


\endgroup
% END UNIT der-ch17/der-ch17_p056-060.tex
% BEGIN UNIT der-ch17/der-ch17_p061-062.tex
\begingroup\makeatletter%
\newcommand{\printedpage}[1]{\marginpar{\footnotesize p.~#1}}%
\newcommand{\msp}[1]{\marginpar{\footnotesize [#1]}}%
\newcommand{\cercle}[1]{\tikz[baseline=(c.base)]\node[draw,circle,inner sep=1pt](c){#1};}%
\makeatother\let\printedpage\derpp\let\msp\dermsp%
\let\mspin\DERmspin%


\printedpage{59}
% head: A. Grothendieck : Dérivateurs, Ch. XVII, Édition des Universités de Montpellier II et Paris VII
% foot: Version 9 février 2019

Now the relation in question $h_0\tilde{i} = \tilde{h}_0$ is equivalent to the set of the two relations
\[
\begin{aligned}
(h_0\tilde{i})r &= h_0(\tilde{i}r)\\
&= h_0i_0\\
&\overset{?}{=} \alpha_0\\
&= \tilde{h}_0r,
\end{aligned}
\]
equivalent to $h_0i_0 = \alpha_0$,
\[
\begin{aligned}
(h_0\tilde{i})s &= h_0(\tilde{i}s)\\
&= h_0\lambda_B\\
&\overset{?}{=} \lambda_Xh'_{I'}\\
&= \tilde{h}_0s,
\end{aligned}
\]
equivalent to $h_0\lambda_B = \lambda_Xh'_{I'}$, that's OK.

This being so, in the square (6.3.14), the hypotheses $i \in\Phi_I$, $p \in \Psi\langle I\rangle$ imply $\tilde{i}_0 \in \Phi$, $p_0 \in \Psi$. Thus the desired arrow $h_0$ exists, q.e.d. % [sic] p_o -> p_0, index printed as letter o, elsewhere index 0 (h_0, i_0)
\DERmark{ENV1098}\DERmark{ENV1099}


This completes the proof of corollary 6.3, and thereby also, that of the last part, d), of 6.1. We are happy\,!

\msp{page 117}

\textbf{Proposition 6.4.}\footnote[34]{Subject to the existence of the relevant $\varinjlim$ and $\varprojlim$ in $\mathcal{M}$. We may assume for simplicity $\mathcal{M}$ stable under $\varinjlim$ and $\varprojlim$ .} % mark printed as superscript (34)
\textit{Let $\mathcal{M}$ be a category, and $\Phi \longleftrightarrow \Psi$ a Q-closed couple in Q-duality ($\Phi = \Psi^*$, $\Psi = \Phi_*$). Let $u : I \longrightarrow I'$ be an arrow in} Cat\textit{, whence a sequence of adjoint functors}
\[
\begin{tikzcd}[column sep=large]
\mathcal{M}(I) \arrow[r, shift left=3, "u_!"] \arrow[r, shift right=3, "u_*"'] & \mathcal{M}(I') \arrow[l, "u^*"']
\end{tikzcd}\ .
\]
\textit{We define $\Phi(I)$, $\Psi\langle I\rangle$, $\Phi(I')$, $\Psi\langle I'\rangle$ as in} 6.2.5. \textit{We then have:}

\cercle{a}
\[
\left\{
\begin{aligned}
u_!(\Phi(I)) &\ \subset\ \Phi(I')\ ,\\
u^*(\Psi\langle I'\rangle) &\ \subset\ \Psi\langle I\rangle\ .
\end{aligned}
\right.
\tag*{(6.4.1)}
\]

\cercle{b} \textit{We have the equivalence}
\[
u^*(\Phi(I')) \subset \Phi(I) \qquad\Longleftrightarrow\qquad u_*(\Psi\langle I\rangle) \subset \Psi\langle I'\rangle\ .
\tag*{(6.4.2)}
\]

\cercle{c} \textit{The equivalent relations} (6.4.2) \textit{hold in each of the following three cases:}
\begin{enumerate}[label={}, leftmargin=2.5em]
\item 1°) \textit{$u$ is coétale.}
\item 2°) \textit{$I$ is discrete.}
\item 3°) \textit{For every $i' \in \operatorname{Ob} I'$, the connected components of $I/i'$ admit a final object.}
\end{enumerate}

\par
\printedpage{60}
% head: A. Grothendieck : Dérivateurs, Ch. XVII, Édition des Universités de Montpellier II et Paris VII
% foot: Version 9 février 2019

\textbf{Corollary 6.5.} \textit{Let $(W,C,F)$ be a special closed Quillen triple in $\mathcal{M}$, and let us define, for every $I$ in} Cat\textit{, $C(I)$, $TC(I)$ and $F(I)$, $TF(I)$ as in theorem} 6.1. a)\textit{. We then have}

\msp{page 118}

\textit{the equivalences}
\[
\Bigl(u^*(C(I'))\ \subset\ C(I)\Bigr) \quad\Longleftrightarrow\quad \Bigl(u_*(TF)\ \subset\ TF'\Bigr)\ ,
\tag*{(6.5.1)}
\]
\[
\Bigl(u^*(TC(I'))\ \subset\ TC(I)\Bigr) \quad\Longleftrightarrow\quad \Bigl(u_*(F)\ \subset\ F'\Bigr)\ .
\tag*{(6.5.2)}
\]
\textit{The equivalent conditions} (6.5.1) \textit{as well as} (6.5.2) \textit{are satisfied when we are in one of the cases} 1°) 2°) 3°) \textit{of} 6.4. c).

\textsc{Proof} of 6.4. The assertions \cercle{a} , \cercle{b} are $\pm$ tautological, by adjunction, compare XIII, prop. 2, and cor., page 46. To prove c), it suffices to treat case 3°), for 1° $\Longrightarrow$ 3°, 2° $\Longrightarrow$ 3°. But the computation of
\[
u_*(X)(i')\ \simeq\ \varprojlim_{i\in I/i'} X(i)
\]
shows that if the connected components $(I/i')_\alpha$ of $I/i'$ each have a final object $e_\alpha$, then
\[
u_*(X)(i')\ \simeq\ \prod_\alpha X(e_\alpha)\ .
\]
Now $\Psi$, being right closed, is stable under arbitrary products, so $X \in \Psi\langle I\rangle$ implies that the $u_*(X)(i')$ are in $\Psi$, so $u_*(X)$ is in $\Psi\langle I'\rangle$, q.e.d.
\DERmark{ENV1100}



\endgroup
% END UNIT der-ch17/der-ch17_p061-062.tex
\clearpage\setcounter{footnote}{0}\phantomsection\addcontentsline{toc}{section}{Chapter XVIII}\markboth{Chapter XVIII}{Chapter XVIII}
% BEGIN UNIT der-ch18/der-ch18_p001-005.tex
\begingroup\makeatletter%
\newcommand{\printedpage}[1]{\marginpar{\footnotesize p.~#1}}%
\newcommand{\msp}[1]{\marginpar{\footnotesize [#1]}}%
\newcommand{\uHom}{\underline{\mathrm{Hom}}}%
\newcommand{\Ens}{\mathrm{Ens}}%
\newcommand{\Fais}{\mathrm{Fais}}%
\newcommand{\pstop}{\mathrm{pstop}}%
\makeatother\let\printedpage\derpp\let\msp\dermsp%
\let\mspin\DERmspin%


\printedpage{} % no page number printed
% head: (none)

\begin{center}
{\Large\textsf{THE DERIVATORS}}

\bigskip

\textbf{ALEXANDRE GROTHENDIECK}

\bigskip\bigskip

\textbf{Chapter XVIII}

\medskip

\textbf{Accessible categories and sets}
\end{center}

\par

\printedpage{} % no page number printed
% head: (none)

\par\noindent\textbf{Note by the editors of the transcribed edition.}\par
This text was deciphered and transcribed in \LaTeXe{} by M. Künzer. It was edited by M. Künzer, J. Malgoire, and G. Maltsiniotis. The transcription is as faithful as possible to the manuscript. Only a few obvious minor corrections have been made. For the rare additions or comments by the editors, as well as for the original numbering of the manuscript pages, typewriter characters \texttt{[typewriter]} in brackets are used. A few footnotes added by the editors are indicated by ``~Ed.~N.~''. The page numbering of this chapter begins at page 50, because originally this chapter was the continuation of chapter XIX which ends in the middle of page 50.

This edition is provisional. Remarks, comments and corrections are welcome. Send a message to:

\textbf{[e-mail address removed]}

\hfill G. Maltsiniotis\hspace*{2cm}

\par

\printedpage{1}
% head: A. Grothendieck : Dérivateurs, Ch. XVIII, Édition des Universités de Montpellier II et Paris VII
% foot: Version 17 décembre 2025

\msp{page 50}

\section*{4\quad Paratopos and the operation $\mathcal{M}\boxtimes\mathcal{N}$}

\subsection*{4.1\quad Pseudo-topos. Problem of the product.}

I call a \textit{pseudo-$\mathfrak{U}$-topos category}, or simply a \textit{pseudo-topos}, a $\mathfrak{U}$-category $\mathcal{M}$ stable under $\mathfrak{U}$-direct limits, and which admits a \textit{small} subcategory $\mathcal{M}_0$ that is \textit{strictly generating} (i.e. generating by strict epimorphisms). Then $\mathcal{M}$ also admits $\mathfrak{U}$-inverse limits (SGA 4 I 8.12.7). (But the example of (Gr), cf. loc. cit. 8.12.9, shows that $\mathcal{M}^\circ$ is not necessarily, for all that, a pseudo-$\mathfrak{U}$-topos, for lack of $\mathcal{M}$ admitting a small subcategory that is strictly \textit{co}generating. But theorem 2.6 \texttt{[of chapter XIX]} gives conditions on $\mathcal{M}$ (mainly of the exactness type) % [sic] par forcément -> pas forcément, misprint; also the parenthesis opened at "(Mais" is never closed
\msp{page 51}
for $\mathcal{M}$ to admit enough injectives -- which will then imply the existence of a small strictly cogenerating subcategory, hence that $\mathcal{M}^\circ$ is also a pseudo-topos (\footnote{This is the case in particular if $\mathcal{M}$ is a topos.})).
\DERmark{ENV1101}


I visualize pseudo-topos categories as being the category of ``sheaves'' on an object $\mathcal{X}$, visualized as a sort of ``space'' which I distinguish from the category $\mathcal{M}=\Fais(\mathcal{X})$ of sheaves on $\mathcal{X}$. It is $\mathcal{X}$ that will be for this the ``pseudo-topos''. Pseudo-topoi form a 2-category, the variance in $\mathcal{X}$ being opposite to that in $\Fais(\mathcal{X})=\mathcal{M}$. To be precise, if $\mathcal{X}$ corresponds to $\mathcal{M}$, $\mathcal{Y}$ to $\mathcal{N}$, a morphism $f$ from $\mathcal{X}$ to $\mathcal{Y}$ will be given by a functor
\[
f^* \;:\; \mathcal{N} \longrightarrow \mathcal{M}
\tag*{(4.1.1)}
\]
in the opposite direction, commuting with arbitrary direct $\varinjlim$ (\footnote{Say here that the objects of $\mathcal{M}$ are visualized as generalized ``opens'' of $\mathcal{X}$. We have $\mathcal{M}\xrightarrow{\sim}\uHom^!(\mathcal{M}^\circ,\Ens)$, i.e. $\mathcal{M}$ is indeed interpreted as the ``sheaves'' on this category $\mathcal{M}$ of opens \ldots}). (\textbf{N.B.} if $\mathcal{X}$, $\mathcal{Y}$ are even topoi, i.e. $\mathcal{M}$, $\mathcal{N}$ topos categories, and if we want $f$ to be a morphism of topoi, we will require in addition that $f^*$ be in addition \textit{left exact}. We will say that a \textit{morphism of pseudo-topoi is exact} if in addition $f^*$ is left exact.)

\msp{page 52}
By virtue of the dual of loc. cit. 8.12.7, this is equivalent to the fact that $f^*$ admits a right adjoint $f_*$,
\[
f_* \;:\; \mathcal{M} \longrightarrow \mathcal{N}\;,
\tag*{(4.1.2)}
\]
which therefore commutes with small $\varprojlim$\,. The datum of $f_*$ satisfying this property therefore determines $f^*$ up to unique isomorphism. But in general, it is not enough to give oneself such an $f_*$, to be able to assert that it corresponds to an $f^*$, i.e. that it admits a left adjoint\,; for this, we would have to know that the functors $\mathcal{M}\longrightarrow\Ens$ which commute with small $\varprojlim$ are representable, which will be the case if $\mathcal{M}$ admits a small subcategory that is \textit{co}generating, i.e. if $\mathcal{M}^\circ$ is also a pseudo-topos category. % [sic] qu'il admette un adjoint à droite -> adjoint à gauche, f^* is the left adjoint of f_*
\DERmark{ENV1102}


\par

\printedpage{2}
% head: A. Grothendieck : Dérivateurs, Ch. XVIII, Édition des Universités de Montpellier II et Paris VII
% foot: Version 17 décembre 2025

It is probably wise to proceed as in the case of topoi, and to call a \textit{morphism of pseudo-topoi} $f:\mathcal{X}\longrightarrow\mathcal{Y}$ a quadruple
\[
(f^*,f_*,\alpha,\beta)
\tag*{(4.1.3)}
\]
\msp{page 53}
of two adjoint functors $(f^*,f_*)$, and of the adjunction morphisms. When $f^*$ is an equivalence, which is equivalent to $f_*$ an equivalence, or $\alpha$ and $\beta$ isomorphisms, we will say that we have an \textit{equivalence of pseudo-topoi}, and for all practical purposes, we will then be able to regard $\mathcal{X}$ and $\mathcal{Y}$ as essentially identical objects.

For two pseudo-topoi $\mathcal{X}$, $\mathcal{Y}$, I feel like introducing (as in the case of topoi) a product pseudo-topos
\[
\mathcal{X}\times\mathcal{Y}\;,
\tag*{(4.1.4)}
\]
which would have the properties of a 2-product in the 2-category of pseudo-topoi (\footnote{It turns out in the end that what we construct is \textit{not} a 2-product in the 2-category under consideration, so it is better to denote it by a sign such as $\mathcal{X}\boxtimes\mathcal{Y}$. But if $\mathcal{X}$, $\mathcal{Y}$ are two topoi, $\mathcal{X}\boxtimes\mathcal{Y}$ is a 2-product in the 2-category of topoi.}). I should specify that by definition
\[
\begin{aligned}
\uHom_{\pstop}(\mathcal{X},\mathcal{Y}) \;&\simeq\; \uHom_!(\mathcal{F}(\mathcal{Y}),\mathcal{F}(\mathcal{X}))^\circ\\
&\simeq\; \uHom^{!'}(\mathcal{F}(\mathcal{X}),\mathcal{F}(\mathcal{Y})),
\end{aligned}
\tag*{(4.1.5)}
\]
where the subscript ! (resp. the superscript !) indicates the full subcategory in the category $\uHom$ under consideration consisting of the functors which commute with small $\varinjlim$ (resp. with small
\msp{page 54}
$\varprojlim$\,), and the accent in the $\uHom^{!'}$ in the last member indicates that we require in addition that the functor in question have a left adjoint.

So let
\[
\mathcal{M} \;=\; \mathcal{F}(\mathcal{X})\;, \qquad \mathcal{N} \;=\; \mathcal{F}(\mathcal{Y})\;,
\tag*{(4.1.6)}
\]
and let $\mathcal{Z}$ be a third pseudo-topos, described by
\[
\mathcal{P} \;=\; \mathcal{F}(\mathcal{Z})\;.
\tag*{(4.1.7)}
\]
We want to make explicit, up to equivalence, the product category $\uHom_{\pstop}(\mathcal{Z},\mathcal{X})\times\uHom_{\pstop}(\mathcal{Z},\mathcal{Y})$, i.e.
\[
\uHom_!(\mathcal{M},\mathcal{P})^\circ \;\times\; \uHom_!(\mathcal{N},\mathcal{P})^\circ.
\tag*{(4.1.8)}
\] % [sic] \mathcal{Q} -> \mathcal{P}, \mathcal{Q} is not defined; \mathcal{Z} is described by \mathcal{P} (4.1.7)
\DERmark{ENV1103}

I am going to introduce for this a category
\printedpage{3}
% head: A. Grothendieck : Dérivateurs, Ch. XVIII, Édition des Universités de Montpellier II et Paris VII
% foot: Version 17 décembre 2025
\msp{page 55}
\[
\mathcal{M}\boxtimes\mathcal{N}
\tag*{(4.1.9)}
\]
by the formula
\[
\mathcal{M}\boxtimes\mathcal{N} \;\simeq\; \uHom^{!!}(\mathcal{M}^\circ\times\mathcal{N}^\circ,\Ens)\;,
\tag*{(4.1.10)}
\]
where the sign $\uHom^{!!}$ denotes the full subcategory of the category $\uHom$, consisting of the bi-contrafunctors % doubt resolved: retain the diplomatic bi-contrafoncteurs reading; see corrections.csv, PDF leaf 5
\[
F \;:\; \mathcal{M}^\circ\times\mathcal{N}^\circ \longrightarrow \Ens\;,
\tag*{(4.1.10.1)}
\]
such that for every $x$ in $\mathcal{M}$, the functor
\[
y \longmapsto F(x,y) \;:\; \mathcal{N}^\circ \longrightarrow \Ens
\tag*{(4.1.10.2)}
\]
is in $\uHom^!(\mathcal{N}^\circ,\Ens)$, i.e. transforms small $\varinjlim$ into $\varprojlim$\,, and that, symmetrically, for every $y$ in $\mathcal{N}$, the functor
\[
x \longmapsto F(x,y) \;:\; \mathcal{M}^\circ \longrightarrow \Ens
\tag*{(4.1.10.3)}
\]
is in $\uHom^!(\mathcal{M}^\circ,\Ens)$. But this also means that these functors are \textit{representable}\,; so we also find
\[
\mathcal{M}\boxtimes\mathcal{N} \xleftarrow{\sim} \uHom^!(\mathcal{N}^\circ,\mathcal{M})
\tag*{(4.1.11)}
\]
and
\[
\mathcal{M}\boxtimes\mathcal{N} \xleftarrow{\sim} \uHom^!(\mathcal{M}^\circ,\mathcal{N})
\tag*{(4.1.12)}
\]
\msp{page 56}
using (in each of these formulas) on both sides the equivalences of categories
\[
\left\{
\begin{aligned}
&\mathcal{M} &&\xrightarrow{\sim} &&\uHom^!(\mathcal{M}^\circ,\Ens)\\
&\mathcal{N} &&\xrightarrow{\sim} &&\uHom^!(\mathcal{N}^\circ,\Ens)\;.
\end{aligned}
\right.
\tag*{(4.1.13)}
\]
We can also say that the category $\mathcal{M}\boxtimes\mathcal{N}$ is equivalent to that of \textit{pairs} of functors
\[
\left\{
\begin{aligned}
&\varphi &&: &&\mathcal{M}^\circ &&\longrightarrow &&\mathcal{N}\\
&\psi &&: &&\mathcal{N}^\circ &&\longrightarrow &&\mathcal{M}\;,
\end{aligned}
\right.
\tag*{(4.1.14)}
\]
equipped with two adjunction morphisms
\[
\left\{
\begin{aligned}
&\psi\varphi^\circ &&\xleftarrow{\alpha} &&\mathrm{id}_{\mathcal{M}}\\
&\varphi\psi^\circ &&\xleftarrow{\beta} &&\mathrm{id}_{\mathcal{N}}
\end{aligned}
\right.
\tag*{(4.1.15)}
\]
satisfying the two usual compatibilities expressing the adjunction of $\varphi^\circ$ and of $\psi$ (or, equivalently, of $\psi^\circ$ and of $\varphi$). The functor
\[
(\varphi,\psi) \longmapsto \varphi \qquad (\text{resp.}\;\; (\varphi,\psi) \longmapsto \psi)
\]
gives an equivalence of the category of these pairs with $\uHom^!(\mathcal{M}^\circ,\mathcal{N})$ (resp. with $\uHom^!(\mathcal{N}^\circ,\mathcal{M})$), so that we have a commutative square of canonical equivalences of functors:


\endgroup
% END UNIT der-ch18/der-ch18_p001-005.tex
% BEGIN UNIT der-ch18/der-ch18_p006-010.tex
\begingroup\makeatletter%
\newcommand{\printedpage}[1]{\marginpar{\footnotesize p.~#1}}%
\newcommand{\msp}[1]{\marginpar{\footnotesize [#1]}}%
\newcommand{\cM}{\mathcal{M}}%
\newcommand{\cN}{\mathcal{N}}%
\newcommand{\cX}{\mathcal{X}}%
\newcommand{\cY}{\mathcal{Y}}%
\newcommand{\cF}{\mathcal{F}}%
\newcommand{\uHom}{\underline{\mathrm{Hom}}}%
\newcommand{\Hom}{\mathrm{Hom}}%
\newcommand{\Ens}{\mathrm{Ens}}%
\newcommand{\Top}{\mathrm{Top}}%
\newcommand{\Ob}{\mathrm{Ob}\,}%
\newcommand{\defeq}{\overset{\mathrm{def}}{=}}%
\newcommand{\bt}{\boxtimes}%
\newcommand{\tbt}{\mathbin{\widetilde{\boxtimes}}}%
\newcommand{\linj}{\lhook\joinrel\longrightarrow}%
\makeatother\let\printedpage\derpp\let\msp\dermsp%
\let\mspin\DERmspin%

\DERmark{ENV1104}


\printedpage{4}
% head: A. Grothendieck : Dérivateurs, Ch. XVIII, Édition des Universités de Montpellier II et Paris VII

\msp{page 57}

\[
\begin{tikzcd}[column sep=small]
 & \cM \tbt \cN \arrow[dl, "\sim" sloped] \arrow[dr, "\sim" sloped] & \\
\uHom^{!}(\cM^{\circ},\cN) \arrow[dr, "\sim" sloped] & & \uHom^{!}(\cN^{\circ},\cM) \arrow[dl, "\sim" sloped] \\
 & \underbrace{\cM \bt \cN}_{\defeq \uHom^{!!}(\cM^{\circ}\times\cN^{\circ},\Ens)} &
\end{tikzcd}
\tag*{(4.1.16)}
\]

\texttt{[where $\cM \tbt \cN$ is the category of pairs defined on the preceding page.]}

I say that $\cM \bt \cN$ \textit{is a pseudo-topos category}. It is clear, \texttt{[in]} addition, that it admits small $\varprojlim$ , which are computed argument by argument in $\uHom(\cM^{\circ}\times\cN^{\circ},\Ens)$, or at one's choice in $\uHom^{!}(\cM^{\circ},\cN)$, or $\uHom^{!}(\cN^{\circ},\cM)$. How to construct small $\varinjlim$ ? There is no question of invoking an argument using cogenerators. Besides, one must also find a small strictly generating subcategory. Evidently, there is something nontrivial to prove!

\section*{4.2\quad Muddles over $\cM \bt \cN$}

Let us return to the general situation of a pseudo-topos $\cX$, described by the pseudo-topos category $\cM$. If $\cN$ is any category, we define the category

\msp{page 58}
\[
\mathrm{Fais}(\cX,\cN) \qquad \text{or} \qquad \cF(\cX,\cN)
\]
of sheaves on $\cX$, with values in $\cN$, by
\[
\cF(\cX,\cN) \ \defeq\ \uHom^{!}(\cM^{\circ},\cN) \ .
\tag*{(4.2.1)}
\]
We thus have a fully faithful functor
\[
\cF(\cX,\cN) \ \defeq\ \uHom^{!}(\cM^{\circ},\cN) \ \linj\ \Hom^{!}(\cN^{\circ},\overline{\cM}) \ ,
\tag*{(4.2.2)}
\]
where $\overline{\cM}$ denotes the subcategory of $\cM^{\wedge}$ (\textbf{N.B.} $\cM^{\wedge}$ is not in general a $\mathcal{U}$-category never mind that) formed by the \textit{representable} functors $\cM^{\circ}\longrightarrow\Ens$ (which is a $\mathfrak{U}$-category, equivalent to $\cM$ by the canonical inclusion $\cM \linj \overline{\cM}$). Thus we have up to unique isomorphism a factorization of (4.2.2) as
\[
\cF(\cX,\cN) \ \overset{\text{fully faithful}}{\lhook\joinrel\relbar\joinrel\relbar\joinrel\longrightarrow}\ \uHom^{!}(\cN^{\circ},\cM) \ ,
\tag*{(4.2.3)}
\]
the essential image consisting of the functors
\[
\psi \ :\ \cN^{\circ} \longrightarrow \cM
\]
such that for every $x$ in $\cM$,
\[
y \longmapsto \Hom_{\mathcal M}(x,\psi(y)) \ :\ \cN^{\circ} \longrightarrow \Ens
\]
\DERmark{ENV1105}

is representable, i.e. of functors which
% foot: Version 17 décembre 2025
\printedpage{5}
% head: A. Grothendieck : Dérivateurs, Ch. XVIII, Édition des Universités de Montpellier II et Paris VII
\msp{page 59}
admit a \textit{left adjoint}. When $\cN$ itself is a pseudo-topos category, this condition is automatic (it follows from the hypothesis that $\psi$ commutes with small $\varprojlim$ ). Thus in this case, (4.2.3) is an equivalence of categories.

Now let $S$ ($S$ as in ``site'') be a \textit{small} subcategory of $\cM$, generating by strict epimorphisms. We know that this means that the composite functor
\[
\cM \xrightarrow{\ \sim\ } \uHom^{!}(\cM^{\circ},\Ens) \xrightarrow{\text{restriction to } S} S^{\wedge} \ (\ \defeq \uHom(S^{\circ},\Ens)\ )
\tag*{(4.2.4)}
\]
is fully faithful, and this implies more generally, that for any category $\cN$, the restriction functor
\[
\cF(\cX,\cN) \ \defeq\ \uHom^{!}(\cM^{\circ},\cN) \ \overset{\text{fully faithful}}{\lhook\joinrel\relbar\joinrel\relbar\joinrel\longrightarrow}\ \uHom(S^{\circ},\cN)
\tag*{(4.2.5)}
\]
is fully faithful. Indeed, the second member, via the fully faithful embedding
\[
\cN \ \overset{\text{fully faithful}}{\lhook\joinrel\relbar\joinrel\relbar\joinrel\longrightarrow}\ \uHom^{!}(\cN^{\circ},\Ens) \ \linj\ \cN^{\wedge} \ ,
\]
embeds into
\[
\uHom(S^{\circ},\uHom^{!}(\cN^{\circ},\Ens)) \ \simeq\ \Hom^{!}(\cN^{\circ},S^{\wedge}) \ ,
\]
and the full faithfulness of (4.2.5) is therefore equivalent to that of the composite functor
\[
\cF(\cX,\cN) \ \defeq\ \uHom^{!}(\cM^{\circ},\cN) \ \longrightarrow\ \uHom^{!}(\cN^{\circ},S^{\wedge}) \ \simeq\ \uHom(S^{\circ},\uHom^{!}(\cN^{\circ},\Ens)) \ .
\tag*{(4.2.6)}
\]

\msp{page 60}
Now we have commutativity in the square
\[
\begin{tikzcd}[column sep=large, row sep=large]
\cF(\cX,\cN) \defeq \uHom^{!}(\cM^{\circ},\cN) \arrow[r, "\text{\scriptsize(4.2.5)}"] \arrow[d, hook, "\text{\scriptsize(4.2.3)}"', "\text{\scriptsize fully faithful}"] & \uHom(S^{\circ},\cN) \arrow[d, hook, "\text{\scriptsize fully faithful}"] \\
\uHom^{!}(\cN^{\circ},\cM) \arrow[r, hook, "\text{\scriptsize(4.2.4)}", "\text{\scriptsize fully faithful}"'] & \underbrace{\uHom^{!}(\cN^{\circ},S^{\wedge})}_{\simeq\, \uHom(S^{\circ},\uHom^{!}(\cN^{\circ},\Ens))}
\end{tikzcd}
\tag*{(4.2.7)}
\]
(where the vertical arrows are equivalences, if $\cN$ itself is a pseudo-topos category). Let us note on the other hand that we have an equivalence
\[
\uHom(S^{\circ},\cN) \ \xleftarrow{\ \sim\ }\ \uHom^{!}(S^{\wedge\circ},\cN) \ \defeq\ \cF(\Top(S),\cN) \ ,
\tag*{(4.2.8)}
\]
where $\Top(S)$ is the pseudo-topos (in fact, a topos) defined by the pseudo-topos category $S^{\wedge}$. For this, let us note that we have
\[
\uHom^{!}(S^{\wedge\circ},\cN) \ \xrightarrow{\ \sim\ }\ \uHom_{!}(S^{\wedge},\cN^{\circ})^{\circ} \ ,
\]
and on the other hand the restriction functor to $S \linj S^{\wedge}$ gives
\[
\uHom_{!}(S^{\wedge},\cN^{\circ}) \ \xrightarrow{\ \sim\ }\ \uHom(S,\cN^{\circ}) \ ,
\]
whence
% foot: Version 17 décembre 2025
\printedpage{6}
% head: A. Grothendieck : Dérivateurs, Ch. XVIII, Édition des Universités de Montpellier II et Paris VII
\msp{page 61}
\[
\uHom^{!}(S^{\wedge\circ},\cN) \ \xrightarrow{\ \sim\ }\ (\uHom(S,\cN^{\circ}))^{\circ} \ \xrightarrow{\ \sim\ }\ \uHom(S^{\circ},\cN) \ ,
\]
\DERmark{ENV1106}

which is none other than (4.2.8). Thus:

\textbf{Proposition 4.2.1.} \textit{Let $\cM=\cF(\cX)$ be a pseudo-topos category, equipped with a strictly generating full subcategory $S$, $\cN$ any $\mathfrak{U}$-category, then the following ``restriction to $S$'' functor is fully faithful.}
\[
\cF(\cX,\cN) \ \defeq\ \uHom^{!}(\cM^{\circ},\cN) \ \longrightarrow\ \uHom(S^{\circ},\cN) \ \xleftarrow{\ \sim\ }\ \underbrace{\uHom^{!}(S^{\wedge\circ},\cN)}_{\defeq\, \cF(\Top(S),\cN)} \ ,
\tag*{(4.2.1.1)}
\]
\DERmark{ENV1107}

\textit{and it fits into the commutative square (4.2.7) of fully faithful functors.}

\textbf{Corollary 4.2.2.} \textit{Suppose that $\cN=\cF(\cY)$ is also a pseudo-topos category, equipped with a strictly generating full subcategory $T$. Then}
\[
\cF(\cX,\cN) \ \simeq\ \cF(\cY,\cM) \ \simeq\ \underbrace{\uHom^{!!}(\cM^{\circ}\times\cN^{\circ},\Ens)}_{=\,\cM\bt\cN} \ ,
\]
\textit{and the restriction functor}
\[
\begin{aligned}
\cM\bt\cN \ \defeq\ \uHom^{!!}(\cM^{\circ}\times\cN^{\circ},\Ens) \quad &\longrightarrow \quad \uHom(S^{\circ}\times T^{\circ},\Ens)\\
&\simeq \quad (S\times T)^{\wedge}\\
&\defeq \quad \cF(\Top(S\times T))
\end{aligned}
\tag*{(4.2.2.1)}
\]
\textit{is fully faithful.}

The first assertion was seen in 4.1. For the full faithfulness of (4.2.2.1), we note that this functor is interpreted, up to equivalence of categories, as the composite of the fully faithful functor (4.2.1.1), and of the functor
\msp{page 62}
analogous for $\cN$ (hence it too fully faithful)
\[
\cF(\cY,S^{\wedge}) \ \defeq\ \uHom^{!}(\cN^{\circ},S^{\wedge}) \ \xrightarrow{\ \text{restr. to } T\ }\ \underbrace{\uHom(T^{\circ},S^{\wedge})}_{\uHom(S^{\circ}\times T^{\circ},\Ens)} \ ,
\tag*{(4.2.2.2)} % [sic] \uHom^{!}(\cN^{\wedge},S^{\wedge}) -> \uHom^{!}(\cN^{\circ},S^{\wedge}), cf. (4.2.2.3) where the same F(Y,S^) is Hom!(N°,S^)
\]
\DERmark{ENV1108}

taking into account the symmetry equivalence
\[
\cF(\cY,S^{\wedge}) \ \defeq\ \uHom^{!}(\cN^{\circ},S^{\wedge}) \ \simeq\ \uHom^{!}(S^{\wedge\circ},\cN) \ .
\tag*{(4.2.2.3)}
\]

All in all, when $\cM$, $\cN$ are both pseudo-topos categories, equipped with strictly generating subcategories $S$ resp. $T$, we have a canonical fully faithful embedding
\[
\cM\bt\cN \ \overset{\rho}{\linj}\ (S\times T)^{\wedge} \ ,
\tag*{(4.2.9)}
\]
% foot: Version 17 décembre 2025
\printedpage{7}
% head: A. Grothendieck : Dérivateurs, Ch. XVIII, Édition des Universités de Montpellier II et Paris VII
which we subsequently interpret as $(f_{S}^{\cM}\times f_{T}^{\cN})_{*}$, where
\[
\left\{
\begin{aligned}
f_{S}^{\cM} \ &:\ \cX \ \longrightarrow\ \Top S\\
f_{T}^{\cN} \ &:\ \cY \ \longrightarrow\ \Top T
\end{aligned}
\right.
\tag*{(4.2.10)}
\]
are the morphisms of pseudo-topoi, given by the fully faithful embeddings (direct image functors)
\[
\left\{
\begin{aligned}
(f_{S}^{\cM})_{*} \ &:\ \cM \ \longrightarrow\ S^{\wedge}\\
(f_{T}^{\cN})_{*} \ &:\ \cN \ \longrightarrow\ T^{\wedge} \ ,
\end{aligned}
\right.
\tag*{(4.2.11)}
\]
which do indeed admit left adjoints
\msp{page 63}
\[
\left\{
\begin{aligned}
(f_{S}^{\cM})^{*} \ &:\ S^{\wedge} \ \longrightarrow\ \cM\\
(f_{T}^{\cN})^{*} \ &:\ T^{\wedge} \ \longrightarrow\ \cN \ ,
\end{aligned}
\right.
\tag*{(4.2.12)}
\]
namely the canonical extensions (by commutation with small $\varinjlim$ ) of the inclusion functors $S \linj \cM$, $T \linj \cN$.

The essential question, now, is evidently to construct a left adjoint for the functor (4.2.9). It amounts to the same thing, for $(s,t)$ in $S\times T$, to define an object
\[
s\bt t \ \in\ \Ob\cM\bt\cN \ ,
\tag*{(4.2.13)}
\]
such that we have, for $F$ varying in $\cM\bt\cN$, an isomorphism functorial in $F$,
\[
\underbrace{\rho(F)(s,t)}_{\defeq\, F(s,t)} \ \simeq\ \Hom(s\bt t,F) \ .
\]
Thus we must prove this:

\textbf{Lemma 4.2.3} (\footnote[4]{Not proved; perhaps false!}). \textit{Let $\cM$, $\cN$ be two pseudo-topos categories, and $s\in\Ob\cM$, $t\in\Ob\cN$, and let us consider on $\cM\bt\cN\defeq\uHom^{!!}(\cM^{\circ}\times\cN^{\circ},\Ens)$ the functor $F\longmapsto F(s,t)$. This functor is} representable \textit{(by an object which I denote by $s\bt t$).}

The natural idea would be to define $s\bt t$ by
\[
(s\bt t)(x,y) \ =\ s(x)\times t(y) \ =\ \Hom_{\cM}(x,s)\times\Hom_{\cN}(y,t) \ ,
\]

\msp{page 64}
which does define a functor
\[
\cM^{\circ}\times\cN^{\circ} \longrightarrow \Ens \ .
\]
Unfortunately, for one of the arguments fixed, say for $y$ fixed, this functor is \textit{not} compatible with inverse limits, e.g. it does not transform the final object of $\cM^{\circ}$ corresponding to the initial object of $\cM$ into a final object of Ens, but into $t(y)=\Hom_{\cN}(y,t)$. % [sic] "de $\cM$)" -> closing parenthesis without opening one, unbalanced as printed
\DERmark{ENV1109}

% foot: Version 17 décembre 2025
\printedpage{8}
% head: A. Grothendieck : Dérivateurs, Ch. XVIII, Édition des Universités de Montpellier II et Paris VII
Nevertheless, I think that $s\bt t$ must indeed be definable as the product of two objects of $\cM\bt\cN$, depending functorially on $s$ resp. on $t$ (\footnote[5]{also false, cf. below.}). We are thus led right now, whether we like it or not, to try to define canonical functors, commuting with small $\varinjlim$ ,
\[
\cM \longrightarrow \cM\bt\cN \longleftarrow \cN
\tag*{(4.2.13)} % [sic] (4.2.13) -> (4.2.14), the label (4.2.13) is already used for s⊠t ∈ Ob M⊠N
\]
(which will be the inverse image functors for the canonical projections
\begin{tikzcd}[row sep=small, column sep=small, baseline=(\tikzcdmatrixname-2-1.base)]
 & \cX \\
\cX\times\cY \arrow[ur] \arrow[dr] & \\
 & \cY
\end{tikzcd}
) (\footnote[6]{The case where $\cM$, $\cN$ are topos-categories, shows us that these functors should be described as $x\longmapsto x\bt e_{\cN}$, $y\longmapsto e_{\cM}\bt y$. But it is not clear that we have outside the case of topoi, and even granting the existence of the $x\bt y$, the formula (valid for a topos) $x\bt y\simeq(x\bt e_{\cN})\times(e_{\cM}\bt y)$.}).

\section*{4.3\quad The operation $\cM\bt\cN$ and the accessibility axiom}

I spent the day yesterday, trying to establish the existence of the $x\bt y$ in $\cM\bt\cN$, without succeeding – at least not without making an additional accessibility hypothesis,
\msp{page 65}
about which it is becoming obvious now that it has to be made in any case. I am going all the same to state the positive of my perplexities.

\textbf{Proposition-perplexity 4.3.1.} \textit{Let $\cM$, $\cN$ be two pseudo-topos categories, whence $\cM\bt\cN$ as in 4.1. Let $S\subset\cM$, $T\subset\cN$ be two small strictly generating subcategories, and let us consider the corresponding fully faithful inclusion:}
\[
\begin{tikzcd}[column sep=large]
\underbrace{\cM\bt\cN}_{\defeq\, \uHom^{!!}(\cM^{\circ}\times\cN^{\circ},\Ens)} \arrow[r, hook, "\alpha"'] & S^{\wedge}\bt T^{\wedge} \ \simeq\ (S\times T)^{\wedge}. \arrow[l, dotted, bend right, "\varphi\,?"']
\end{tikzcd}
\tag*{(4.3.1.1)}
\]
\textit{Let us consider the following conditions:}
\begin{enumerate}[label={}]
\item[\textit{a)}] \textit{$\cM\bt\cN$ is a pseudo-topos.}
\item[\textit{b)}] \textit{For $x\in\Ob\cM$, $y\in\Ob\cN$, $x\bt y\in\Ob\cM\bt\cN$ exists, i.e. the functor (commuting with small $\varprojlim$ )}
\[
F \ \longmapsto\ F(x,y) \ , \qquad \cM\bt\cN \longrightarrow \Ens
\tag*{(4.3.1.2)}
\]
\textit{is representable.}
\item[\textit{b$'$)}] \textit{As in b), with $x\in\Ob S$, $y\in\Ob T$.}
\item[\textit{c)}] \textit{The functor $\alpha$ (which commutes with small $\varprojlim$ ) admits a left adjoint}
\[
\varphi \ :\ (S\times T)^{\wedge} \longrightarrow \cM\bt\cN \ .
\]
\end{enumerate}
\textit{We then have the implications}
\[
b \ \Longrightarrow\ b' \ \Longleftrightarrow\ c \ \Longrightarrow\ a \ .
\tag*{(4.3.1.3)}
\]
% foot: Version 17 décembre 2025


\endgroup
% END UNIT der-ch18/der-ch18_p006-010.tex
% BEGIN UNIT der-ch18/der-ch18_p011-015.tex
\begingroup\makeatletter%
\newcommand{\printedpage}[1]{\marginpar{\footnotesize p.~#1}}%
\newcommand{\msp}[1]{\marginpar{\footnotesize [#1]}}%
\newcommand{\Hom}{\operatorname{Hom}}%
\newcommand{\Ob}{\operatorname{Ob}}%
\newcommand{\pr}{\operatorname{pr}}%
\newcommand{\card}{\operatorname{card}}%
\newcommand{\Fl}{\operatorname{Fl}}%
\newcommand{\Filt}{\operatorname{Filt}}%
\newcommand{\Ind}{\operatorname{Ind}}%
\makeatother\let\printedpage\derpp\let\msp\dermsp%
\let\mspin\DERmspin%


\setcounter{footnote}{6}

\printedpage{9}
% head: A. Grothendieck : Dérivateurs, Ch. XVIII, Édition des Universités de Montpellier II et Paris VII

\noindent\textit{Moreover, if c is satisfied, we have this: for $s\in S$, $t\in T$, let us consider}
\[
(s,t)\ \in\ (S\times T)\ \subset\ (S\times T)^\wedge\ \simeq\ S^\wedge\boxtimes T^\wedge\ ,
\]
\textit{which is also written out as}
\[
\left\{
\begin{aligned}
&\rlap{$\displaystyle (s,t)\ =\ \pr_1^*(s)\times\pr_2^*(t)\ \simeq\ s\overset{S^\wedge,T^\wedge}{\boxtimes}t\in\Ob\,(S\times T)^\wedge\ ,$}\\
&\textit{i.e.}\ (s,t)(s',t')\ =\ s(s')\times t(t') &&=\ \Hom_S(s',s)\times\Hom_T(t',t)\\
& &&=\ \Hom_{S\times T}((s',t'),(s,t))\ ,
\end{aligned}
\right.
\tag*{(4.3.1.4)}
\]
\[
\begin{tikzcd}
S\times T \arrow[r, "\mathrm{pr}_2"] \arrow[d, "\mathrm{pr}_1"] & T\\
S\ ,
\end{tikzcd}
\]
\msp{page 66}\textit{we have}
% note: the label (4.3.1.4) is printed for this display as well as for the preceding one
\[
s\overset{\mathcal M,\mathcal N}{\boxtimes}t\ \simeq\ \varphi(s,t)\ ,
\tag*{(4.3.1.4)}
\]
\textit{and $\varphi(S\times T)=\{s\boxtimes t\mid s\in S,\ t\in T\}$ is a strictly generating subset of $\mathcal M\boxtimes\mathcal N$ (}\footnote{\textbf{N.B.} Of course we have $(s,t)=s\overset{S^\wedge,T^\wedge}{\boxtimes}t$.}\textit{).}

\noindent\textbf{N.B.} I know of no example where the conditions under consideration are not satisfied, but nevertheless now doubt that even the weakest always is, i.e. that $\mathcal M\boxtimes\mathcal N$ is a pseudo-topos as soon as $\mathcal M$ and $\mathcal N$ are.

\noindent\textsc{Proof of} 4.3.1. The implication b $\Longrightarrow$ b$'$ is tautological. The implication b$'$ $\Longleftrightarrow$ c follows from the

% marginal vertical bar printed beside Lemme 4.3.2 (from "Lemme 4.3.2." to "soit représentable.")
\noindent\textbf{Lemma 4.3.2.} \textit{Let $\mathcal E$, $X$ be $\mathfrak U$-categories with $X$ small, $\mathcal E$ stable under small $\varinjlim$ ,}
\DERmark{ENV1110}

\[
\alpha\ :\ \mathcal E\longrightarrow X^\wedge
\]
\textit{a functor. For $\alpha$ to admit a left adjoint $\varphi$, it is necessary and sufficient that for every $x\in\Ob X$, the functor}
\[
\left\{
\begin{array}{ccl}
\xi & \longmapsto & \Hom_{X^\wedge}(x,\alpha(\xi))\ =\ \alpha(\xi)(x)\\
\mathcal E & \longrightarrow & \mathrm{Ens}
\end{array}
\right.
\]
\DERmark{ENV1111}

\textit{be representable.}

\noindent This is obviously necessary. Sufficiency follows from the fact that $X\subset X^\wedge$ is strict%
% foot: Version 17 décembre 2025
\printedpage{10}%
% head: A. Grothendieck : Dérivateurs, Ch. XVIII, Édition des Universités de Montpellier II et Paris VII
\msp{page 67}%
ly generating in $X^\wedge$, i.e. that every object $F$ of $X^\wedge$ is a direct limit $F=\varinjlim\limits_I F_i$, the $F_i$ in $X$ (we take $I=X/F$), whence
\[
\begin{aligned}
\Hom(F,\alpha(\xi))\ &\simeq\ \varprojlim_I\ \underbrace{\Hom(F_i,\alpha(\xi))}_{\Hom(\rho(F_i),\xi)}\\
&\simeq\ \Hom(\varinjlim_I\,\rho(F_i),\xi)\ ,
\end{aligned}
\]
\DERmark{ENV1112}

so the functor $\xi\longmapsto\Hom(F,\alpha(\xi))$ is represented by $\rho(F)\overset{\mathrm{def}}{=}\varinjlim\limits_I\rho(F_i)$.
\DERmark{ENV1113}


\noindent Let us prove c $\Longrightarrow$ a. For the existence of small $\varinjlim$ in $\mathcal M\boxtimes\mathcal N$, and the existence of a strictly generating subset, we invoke the general lemma:

% marginal vertical bar printed beside Lemme 4.3.3 (from "Lemme 4.3.3." to "comme")
\noindent\textbf{Lemma 4.3.3.} \textit{Let $\alpha:\mathcal E\hookrightarrow\mathcal F$ be a fully faithful functor admitting a left adjoint $\rho$.}
\begin{enumerate}[label={}]
\item[\textcircled{a}] \textit{If the $\varinjlim$ of a given type $I$ exist in $\mathcal F$, they exist in $\mathcal E$, and are computed as}
\[
\varinjlim_I{}_{\mathcal E}\,\xi_i\ =\ \rho(\varinjlim_I{}_{\mathcal F}\,\alpha(\xi_i))\ .
\tag*{(4.3.3.1)}
\]
\item[\textcircled{b}] \textit{If $U$ is a small strictly generating subset of $\mathcal F$, $\rho(U)$ is a small strictly generating subset in $\mathcal E$.}
\end{enumerate}
\DERmark{ENV1114}


\noindent Indeed, if $\xi$ denotes the right-hand side \texttt{[of equality 4.3.3.1]}, and $\eta$ a variable object of $\mathcal E$, we will have
\[
\begin{aligned}
\Hom_{\mathcal E}(\xi,\eta)\ &\overset{\text{adj.}}{\simeq}\ \Hom_{\mathcal F}(\varinjlim_I\alpha(\xi_i),\alpha(\eta))\\
&\simeq\ \varprojlim_I\quad\underbrace{\Hom_{\mathcal F}(\alpha(\xi_i),\alpha(\eta))}_{\simeq\,\Hom_{\mathcal E}(\xi_i,\eta)\ \ \text{since $\alpha$ fully faithful}}\quad .
\end{aligned}
\]
\msp{page 68}
\[
\simeq\ \ \varprojlim_I\Hom_{\mathcal E}(\xi_i,\eta)\ ,
\]
which proves that $\xi$ is a direct limit of the $\xi_i$ in $\mathcal E$ (\footnote{To prove b), we note that the full faithfulness of $\alpha$ implies that $\rho$ is essentially surjective; as $\rho$ commutes with small $\varinjlim$ and every object of $\mathcal F$ is of the form $\varinjlim\limits_I x_i$, the $x_i$ in $U$, we will have $\rho(x)=\varinjlim\limits_I\rho(x_i)$, the $\rho(x_i)$ in $\rho(U)$, we win.}).
\DERmark{ENV1115}


\noindent This proves that (assuming c)) $\mathcal M\boxtimes\mathcal N$ is stable under small $\varinjlim$ , so it remains to prove that $\mathcal M\boxtimes\mathcal N$ admits a small strictly generating family. This will follow from b) in the lemma, since $S\times T$ strictly generates $(S\times T)^\wedge$. This proves at the same time the last assertion of the proposition, subject to proving formula (4.3.1.4). We must therefore prove an isomorphism functorial in $F\in\Ob\mathcal M\boxtimes\mathcal N$, % [sic] (4.1.3.4) -> (4.3.1.4), the formula s ⊠ t ≃ φ(s,t) to be proved is labelled (4.3.1.4) on p. 9
\DERmark{ENV1116}

\[
\Hom(\rho(s,t),F)\ \simeq\ F(s,t)\ ,
\]
but by adjunction we have
\[
\Hom(\rho(s,t),F)\ \simeq\ \Hom((s,t),\alpha(F))\ =\ \alpha(F)(s,t)\ =\ F(s,t)\ ,
\]
OK.

% foot: Version 17 décembre 2025
\printedpage{11}
% head: A. Grothendieck : Dérivateurs, Ch. XVIII, Édition des Universités de Montpellier II et Paris VII

\noindent\msp{page 69}\textbf{Proposition 4.3.4.} \textit{Let $\mathcal M$ be a $\mathfrak U$-category. The following conditions are equivalent.}
\begin{enumerate}[label={}]
\item[\textit{a)}] \textit{$\mathcal M$ is a $\mathfrak U$-pseudo-topos category, i.e. it is stable under $\mathfrak U$-direct limits, and admits a $\mathfrak U$-small strictly generating subcategory $S$.}
\item[\textit{b)}] \textit{There exist a $\mathfrak U$-small category $S$, and a functor $i:\mathcal M\hookrightarrow S^\wedge$ that is} fully faithful, \textit{having a left adjoint.}
\end{enumerate}

\noindent\textsc{Proof}. We have a) $\Longrightarrow$ b), for it suffices to take for $S$ in b) the one specified in a), the hypothesis on $S$ implies that the canonical functor
\[
\mathcal M\ \longrightarrow\ S^\wedge\ ,\qquad\quad \text{composite of } \mathcal M\hookrightarrow\ \mathcal M^\wedge\ \xrightarrow{\ \text{restriction}\ }S^\wedge
\]
is fully faithful. That it admits a left adjoint follows from lemma 4.3.2. The implication b) $\Longrightarrow$ a) follows from lemma 4.3.3.

\noindent\textbf{Corollary 4.3.5.} \textit{Let $\mathcal M$ be a $\mathfrak U$-category. For $\mathcal M$ to be a $\mathfrak U$-pseudo-topos that is} accessible \textit{(i.e. all of whose objects are accessible), it is necessary and sufficient that there exist $S$ in} Cat \textit{and a functor fully faithful and} accessible $\mathcal M\overset{\alpha}{\hookrightarrow}S^\wedge$\textit{, admitting a left adjoint $\rho$.}

\noindent\msp{page 70}\textsc{Proof}. It suffices to note that if $S$, $\alpha$ are as in 4.3.4 b), then $\alpha$ is accessible if and only if the $\rho(S)$ are made up of accessible objects in $\mathcal M$; as this subset is strictly generating, we know that this implies that \textit{every} object of $\mathcal M$ is accessible. (Cf. SGA 4 I 9.9.)

\noindent I am going to try to give a more complete statement:

\noindent\textbf{Theorem 4.3.6.} \textit{Let $\mathcal M$ be a $\mathfrak U$-category. Equivalent conditions:}
\begin{enumerate}[label={}]
\item[\textit{a)}] \textit{$\mathcal M$ stable under small $\varinjlim$ , and it admits a small full subcategory $S$, strictly generating, and consisting of accessible objects of $\mathcal M$.}
\item[\textit{a$'$)}] \textit{$\mathcal M$ is a weakly accessible $\mathfrak U$-pseudo-topos category, i.e. satisfies a), with in addition} every \textit{element of $\mathcal M$ (not only those of $S$) accessible in $\mathcal M$.}
\item[\textit{b)}] \textit{There exist a small category $S$, and a functor $\alpha:\mathcal M\hookrightarrow S^\wedge$ that is} fully faithful and accessible, \textit{such that $\alpha$ admits a left adjoint.}
\item[\textit{c)}] \msp{page 71}\textit{$\mathcal M$ is stable under small $\varinjlim$ , is accessible, and admits a cardinal filtration} $(\Filt^\pi\mathcal M)_{\pi\geqslant\pi_0}$.
\item[\textit{d)}] \textit{There exist a cardinal $\pi_1$, a small category $\mathcal M_{\pi_1}$, stable under direct limits of cardinality $\leqslant\pi_1$, such that we have an equivalence of categories}
\[
\mathcal M\ \xleftarrow{\ \sim\ }\ \Ind_{\pi_1}\mathcal M_{\pi_1}\ ,
\tag*{(4.3.6.1)}
\]
\textit{where the right-hand side denotes the full subcategory of} $\Ind(\mathcal M_{\pi_1})$ \textit{consisting of the ind-objects $(X_i)_{i\in I}$ of $\mathcal M_{\pi_1}$ such that $I$ is large relative to $\pi_1$.}
\end{enumerate}

\noindent\textsc{Proof of} 4.3.6 (\footnote{For a complete proof (the one given here is botched from a certain point on), cf. §4.6 (p. 139-148).}). We have already seen a$'$ $\Longrightarrow$ a $\Longrightarrow$ b $\Longrightarrow$ a$'$, i.e. the equivalence of a, a$'$, b, let PTA be this condition. We are going to prove the implications
\[
\mathrm{PTA}\ \Longrightarrow\ \mathrm{c}\ \Longrightarrow\ \mathrm{d}\ \Longrightarrow\ \mathrm{PTA}\ ,
\]
% foot: Version 17 décembre 2025
\printedpage{12}%
% head: A. Grothendieck : Dérivateurs, Ch. XVIII, Édition des Universités de Montpellier II et Paris VII
which will establish the theorem.

\noindent\textit{Let us prove PTA $\Longrightarrow$ c.} Let $S$ be a small strictly generating full subcategory. Thus, \msp{page 72}for every $x$ in $\mathcal M$ we have
\[
x\ =\ \varinjlim_{s\in S/x}s
\tag*{(4.3.6.2)}
\]
(cf. SGA 4 I 7.2 (i)). Let $\pi_0=\card\Fl(S)$, and for every cardinal $\pi\geqslant\pi_0$, let
\[
\begin{array}{lcl}
\Filt^\pi\mathcal M & = & \text{full subcategory of $\mathcal M$ consisting of objects $x$ such that}\\
 & & \card\,(\Fl(S/x))\leqslant\pi\ .
\end{array}
\tag*{(4.3.6.3)}
\]
I want to prove that we thus get a cardinal filtration on $\mathcal M$, i.e. that conditions a) b) c) of SGA 4 I 9.12 are satisfied. This is proved in loc. cit. 9.13 under slightly different hypotheses, the condition of accessibility of $\mathcal M$ being replaced by the one (much less pleasant!) that the strict epimorphic families of arrows in $\mathcal M$ are strict universal epimorphic families. This condition is used only in the proof of condition b), namely the stability of $\Filt^\pi\mathcal M$ under $\varinjlim\limits_I$ of cardinality $\leqslant\pi$ (without, moreover, having to assume that $I$ is filtering, nor a fortiori large relative to $\pi_0$). It must moreover be said that the definition of $\Filt^\pi\mathcal M$ given in loc. cit. is not quite the same as (4.3.6.3), \msp{page 73}but it is
% the printed line break "car-/dinal" inside the display is kept as printed
\[
\begin{array}{lcl}
\Filt'^{\pi}\mathcal M & = & \text{full subcategory of $\mathcal M$ consisting of objects $x$ such that there}\\
 & & \text{exists a strict epimorphic family $x_i\xrightarrow{\alpha_i}x$ of car-}\\
 & & \text{dinality $\leqslant\pi$, with $x_i\in S\ \forall\ i\in I$.}
\end{array}
\tag*{(4.3.6.3$'$)}
\]
Thus we have
\[
\Filt^\pi\mathcal M\ \subset\ \Filt'^{\pi}\mathcal M\ ,
\]
but in loc. cit. it is proved that for $x$ in $\Filt'^{\pi}\mathcal M$, we have
\[
\card\Fl(S/x)\ \leqslant\ \pi^{\pi_0}\ ,
\]
whence
\[
\Filt^\pi\mathcal M\ \subset\ \Filt'^{\pi}\mathcal M\ \subset\ \Filt^{\pi^{\pi_0}}\mathcal M\ ,
\]
so if $\pi=\pi^{\pi_0}$, e.g. $\pi=2^c$ with $c\geqslant\pi_0$, then we get
\[
\Filt^\pi\mathcal M\ =\ \Filt'^{\pi}\mathcal M\ .
\]

\noindent So the two filtrations look $\pm$ equivalent. But the proof given of the stability of $\Filt'^{\pi}$ really seems to make essential use of the universality condition on strict epimorphic families (hence effective epimorphic, with the hypotheses made implying the existence of small $\varprojlim$ ). So I must, willy-nilly, do the work over again! % [sic] on l'air -> ont l'air, verb avoir (3rd person plural) intended

% foot: Version 17 décembre 2025
\printedpage{13}
% head: A. Grothendieck : Dérivateurs, Ch. XVIII, Édition des Universités de Montpellier II et Paris VII

\noindent\msp{page 74}Condition a): $\Filt^\pi$ equivalent to a small category. This is immediate, since by (4.3.6.2), $\Filt^\pi$ consists of $\varinjlim\limits_I$ in $\mathcal M$ of objects of $S$ (a small category), with $I$ such that $\card\Fl I\leqslant\pi$.
\DERmark{ENV1117}


\noindent Condition b): $\Filt^\pi$ stable under $\varinjlim\limits_I$ , when $\card\Fl I\leqslant\pi$ (\footnote{when $I$ filtering and large relative to $\pi_0$ -- I prove it only in this case, alas!}). So let $(x_i)_{i\in I}$, where
\DERmark{ENV1118}

\[
\begin{array}{ccc}
I & \longrightarrow & \mathcal M\\[1ex]
i & \longmapsto & x_i
\end{array}
\]
a direct system in $\mathcal M$, with $x_i$ in $\Filt^\pi$. Let us consider the fiber product
\[
\begin{tikzcd}
I \arrow[d, "\varphi"'] & S/I = J \arrow[l, "\text{cofibered}"'] \arrow[d] \arrow[dr, dotted, "\psi"] & & \\
\mathcal M & S/\mathcal M \arrow[l, "\text{target}"', "\text{Cat-cofibered}"] \arrow[r, "\text{source}"] & S \arrow[r, hook] & \mathcal M\ ,
\end{tikzcd}
\tag*{(4.3.6.4)}
\]
where $J$ is Cat-cofibered over $I$, its fiber at $i\in\Ob I$ being isomorphic to $S/x_i$, so $\card\Fl J_i\leqslant\pi$. Since by hypothesis $\card\Fl I\leqslant\pi$, we conclude that $\card\Fl J\leqslant\pi^3=\pi$. (\textbf{N.B.} The set of arrows $v:j_0\longrightarrow j_1$ of $J$ over an arrow $u:i_0\longrightarrow i_1$ of $I$ is ``contained'' in the product set $\Ob J_{i_0}\times\Fl J_{i_1}$, by the map
\[
v\longmapsto(\operatorname{source}v,v':u_*(j_0)\longrightarrow j_1)\ ,
\]
hence of cardinality bounded above by $\pi^2$.) We \msp{page 75}then have
\[
x\ \overset{\mathrm{def}}{=}\ \varinjlim_I x_i\ =\ \varinjlim_J\ \underbrace{\xi_j}_{=\,\psi(j)}\qquad\quad\mbox{(\footnotemark)}\ ,
\]
\DERmark{ENV1119}

\footnotetext{We use the distributivity property of $\varinjlim$ , relative to a Cat-cofibrant functor
\[
\begin{tikzcd}[ampersand replacement=\&]
J \arrow[r] \arrow[dr] \& I\\
\& \mathcal M\ ,
\end{tikzcd}
\]
made explicit above.}%
where this time the limit is a limit of objects of $S$, taken in $\mathcal M$. It follows that $x\in\Filt'^{\pi}$, this is not quite $x\in\Filt^\pi$. (That implies $x\in\Filt^{\pi^{\pi_0}}$, so $x\in\Filt^\pi$ if $\pi=\pi^{\pi_0}$, e.g. $\pi=2^c$ with $c\geqslant\pi_0$.) But this is the moment to use the fact that the objects of $S$ are accessible, and we will take $\pi_0$ such that the objects of $S$ are $\pi_0$-accessible, and this time we will take $I$ filtering and large relative to $\pi_0$ all the same (to restrict ourselves to what is required in condition b) of loc. cit.). We must prove that $\card\,(\Fl(S/x))\leqslant\pi$ (which means $x\in\Filt^\pi$). But since the objects of $S$ are $\pi_0$-accessible and $I$ large relative to $\pi_0$, we see at once that
\[
S/x\ \simeq\ \varinjlim_i S/x_i\ ,
\]
% foot: Version 17 décembre 2025


\endgroup
% END UNIT der-ch18/der-ch18_p011-015.tex
% BEGIN UNIT der-ch18/der-ch18_p016-020.tex
\begingroup\makeatletter%
\newcommand{\printedpage}[1]{\marginpar{\footnotesize p.~#1}}%
\newcommand{\msp}[1]{\marginpar{\footnotesize [#1]}}%
\newcommand{\uHom}{\underline{\mathrm{Hom}}}%
\newcommand{\Ind}{\mathrm{Ind}}%
\newcommand{\Hom}{\mathrm{Hom}}%
\newcommand{\card}{\mathrm{card}}%
\newcommand{\Fl}{\mathrm{Fl}}%
\newcommand{\Ob}{\mathrm{Ob}}%
\newcommand{\Filt}{\mathrm{Filt}}%
\newcommand{\hookxrightarrow}[2][]{\lhook\joinrel\xrightarrow[#1]{#2}}%
\makeatother\let\printedpage\derpp\let\msp\dermsp%
\let\mspin\DERmspin%


\setcounter{footnote}{11}

\printedpage{14}
% head: A. Grothendieck : Dérivateurs, Ch. XVIII, Édition des Universités de Montpellier II et Paris VII

whence
\[
\Fl\,S/x \;\simeq\; \varinjlim_{i} \Fl\,S/x_i \;,
\]
hence
\[
\card\Fl(S/x) \;\leqslant\; (\underbrace{\card I}_{\leqslant\pi})\underbrace{(\sup_{i}\card\Fl(S/x_i))}_{\leqslant\pi} \;\leqslant \pi^2 = \pi \;,
\]

\hfill q.e.d.

\msp{page 76}
% margin: hand-drawn bracket mark in the left margin beside the next two lines
Let us prove c) (\footnote{\textit{Botched proof}, only works if $\pi^{\pi_0} = \pi$\,!!}), i.e. that for every $\pi' > \pi \geqslant \pi_0$ and every $x \in \Ob\Filt^{\pi'}$, one can write
\[
x \;=\; \varinjlim_{I} x_i \qquad I \text{ large relative to } \pi\text{, the } x_i \text{ in } \Filt^{\pi}\text{, } \card I \leqslant \pi'^{\pi}. \tag*{(4.3.6.5)}
\]
% margin: hand-drawn bracket mark in the left margin beside the next three lines
The proof goes as in loc. cit. (p. 151-152), slightly changing the definition of $I$, as denoting the set of full subcategories $i$ of $S/x$ which are filtering, of cardinality $\leqslant \pi$ \textit{and moreover large relative to} $\pi_0$ (\footnote{It is not certain that $I$ is large relative to $\pi$, unless one assumes $\pi^{\pi_0} = \pi$\,!}). (We may suppose $S/x$ filtering by enlarging $S$ so that $S$ is stable in $\mathcal{M}$ under small finite $\varinjlim$, hence the categories $S/x$ too, which implies that they are filtering. It is \texttt{[incomplete sentence]})

We have thus proved PTA $\Longrightarrow$ c).

\textit{Let us prove c) $\Longrightarrow$ d).} Let
\[
\pi_1 \;=\; 2^c \qquad \text{with} \qquad c \;\geqslant\; \pi_0 \;.
\]
We then know that
\[
\Filt^{\pi_1}\mathcal{M} \;=\; \mathcal{M}_{\pi_1} \;\overset{\mathrm{def}}{=}\; \text{full subcategory of $\mathcal{M}$ of the $\pi_1$-accessible objects of $\mathcal{M}$}
\]
\DERmark{ENV1120}

(loc. cit. 9.16), which implies in particular that $\mathcal{M}_{\pi_1}$ is stable under $\varinjlim$ of cardinality $\leqslant \pi_1$, by virtue of loc. cit. 9.9 (page 145). Moreover,
\msp{page 77}
by loc. cit. 9.18 b) (p. 156), we have
\[
\Ind_{\pi_1}(\mathcal{M}_{\pi_1}) \xrightarrow{\;\sim\;} \mathcal{M} \;,
\]
which proves d).

\textit{Let us prove d) $\Longrightarrow$ PTA.} By loc. cit. 9.18 a) (p. 156), we see that the objects of $\mathcal{M}_{\pi_1}$ are $\pi_1$-accessible in $\Ind_{\pi_1}(\mathcal{M}_{\pi_1})$, on the other hand it is clear that $\mathcal{M}_{\pi_1}$ is a strictly generating full subcategory of $\Ind_{\pi_1}(\mathcal{M}_{\pi_1})$, so it remains to see that $\Ind_{\pi_1}(\mathcal{M}_{\pi_1})$ is stable under small $\varinjlim$ , i.e. the

\textbf{Lemma 4.3.7.} \textit{Let $\mathcal{M}_\pi$ be a $\mathfrak{U}$-category stable under direct limits of cardinality $\leqslant \pi$, where $\pi$ is a given infinite cardinal. Then} $\Ind_\pi\mathcal{M}_\pi$ \textit{is stable under small direct limits. Moreover, $\mathcal{M} = \Ind_\pi\mathcal{M}_\pi$, equipped with $\mathcal{M}_\pi \longrightarrow \mathcal{M}$, is a solution of a $2$-universal problem, namely
\printedpage{15}
% head: A. Grothendieck : Dérivateurs, Ch. XVIII, Édition des Universités de Montpellier II et Paris VII
that for every category $\mathcal{N}$ stable under small direct limits, the restriction functor induces an equivalence}
\[
\uHom_{!}(\mathcal{M},\mathcal{N}) \xrightarrow{\;\sim\;} \uHom_{!\,\pi}(\mathcal{M}_\pi,\mathcal{N}) \;, \tag*{(4.3.7.1)}
\]
\textit{where the sign ! (resp. $!\,\pi$) denotes the full subcategory of the category} $\uHom$\textit{, formed by the functors which commute with small} $\varinjlim$ \textit{(resp. with direct limits of cardinality $\leqslant \pi$).}

\msp{page 78}
Let us first prove that the $\varinjlim$ of cardinality $\leqslant \pi$ \textit{in} $\mathcal{M}_\pi$ are also $\varinjlim$ \textit{in} $\mathcal{M} = \Ind_\pi\mathcal{M}_\pi$ (\footnote{But \textit{beware}, $\mathcal{M}_\pi \longrightarrow \Ind(\mathcal{M}_\pi)$ does \textit{not} commute with these $\varinjlim$ !}). We thus have
\[
x \xleftarrow{\;\sim\;} \varinjlim_{I} x_i \qquad \textit{in } \mathcal{M}_\pi \;,
\]
i.e.
\[
\Hom(x,y) \xrightarrow{\;\sim\;} \varprojlim_{I} \Hom(x_i,y) \qquad \text{if } y \in \Ob\mathcal{M}_\pi \;, \tag*{$(*)$}
\]
let us prove that this remains true for $y$ in $\mathcal{M}$, i.e. $y$ of the form
\[
y \;=\; \varinjlim_{J} y_j \;, \qquad J \text{ large relative to } \pi\text{, the } y_j \text{ in } \mathcal{M}_\pi.
\]
We have indeed
\[
\begin{aligned}
\Hom(x,y) \quad &\simeq \quad \varinjlim_{J} \Hom(x,y_j)\\
&\overset{(*)}{\simeq} \quad \varinjlim_{J}\, \varprojlim_{I} \Hom(x_i,y_j)\\
&\overset{\text{SGA 4 I 9.8}}{\simeq} \quad \varprojlim_{I}\, \varinjlim_{J} \Hom(x_i,y_j)\\
&\simeq \quad \varprojlim_{I} \Hom(x_i,y) \;,
\end{aligned}
\]
altogether
\[
\Hom(x,y) \;\simeq\; \varprojlim_{I} \Hom(x_i,y),
\]
q.e.d.

To prove that the small $\varinjlim$ are representable in $\mathcal{M} = \Ind_\pi(\mathcal{M}_\pi)$, as they are in $\Ind(\mathcal{M}_\pi)$ (SGA 4 I 8.9.5 b, we use the hypothesis that the \textit{finite} $\varinjlim$ are representable in $\mathcal{M}_\pi$), it remains to see that $\Ind_\pi(\mathcal{M}_\pi)$ is stable in $\mathcal{M}_\pi$ under these limits.

\msp{page 79}

\section*{4.4\quad Study of $\Ind_\pi(C)$ (for $C$ a small category stable under $\varinjlim$ of cardinality $\leqslant \pi$)}

\textbf{4.4.1.} In the end, the proof of Lemma 4.3.7 is not so obvious as I anticipated, and I must devote some technical preliminaries to it, which I am going to collect here. Throughout what follows, $C$ represents a $\mathfrak{U}$-category equivalent to a small category.

\printedpage{16}
% head: A. Grothendieck : Dérivateurs, Ch. XVIII, Édition des Universités de Montpellier II et Paris VII

\textbf{Proposition 4.4.2.} \textit{Let $F \in C^{\wedge}$ be an ind-representable functor, i.e. such that $C/F$ is filtering. Let, on the other hand, $i \longmapsto x_i$ be a functor $I \longrightarrow C$, where $I$ is a small filtering category, and let}
\[
\alpha \;:\; x \;=\; \varinjlim_{I}{}_{C^{\wedge}} x_i \longrightarrow F \qquad \text{(\footnotemark)} \tag*{(4.4.2.1)}
\]
\footnotetext{Giving $\alpha$ is also equivalent to giving a homomorphism of ind-objects $(x_i)_{i\in I} \longrightarrow F$, $F$ being considered as a (generalized) ind-object indexed by $C/F$.}
\textit{an inverse system of arrows $x_i \longrightarrow F$ (}\textbf{N.B.} \textit{the direct limit $x$ is taken in $C^{\wedge}$, and is computed argument by argument), which is equivalent to a factorization $\varphi_\alpha : I \longrightarrow C/F$ of $I \longrightarrow C$ as}
\[
I \xrightarrow{\;\varphi_\alpha\;} C/F \hookxrightarrow[\psi]{\textit{can.}} C \;.
\]
\textit{That said, for $\alpha$ (4.4.2.1) to be an isomorphism in} $\Ind'(C) \subset C^{\wedge}$ \textit{(full subcategory of ind-representable functors on $C$), it is necessary and sufficient that $\varphi_\alpha$ be} cofinal\textit{, or what amounts to the same thing, $I$ being filtering, that it satisfy the two usual conditions:}
\begin{itemize}[leftmargin=2.5em]
\item[F\,1)] $\forall\, \xi \in \Ob C/F$, $\exists\, i \in \Ob I$ \textit{and an arrow} $\xi \longrightarrow \varphi_\alpha(i)$.
\item[F\,2)] $\forall\, \xi \in \Ob C/F$, \textit{and} $i \in \Ob I$, \textit{and every double arrow}

\msp{page 80}
\[
\xi \rightrightarrows \varphi_\alpha(i) \;,
\]
\textit{there exists an arrow $i \xrightarrow{\;u\;} j$ in $I$ such that $\varphi_\alpha(u)$ equalizes the double arrow.}
\end{itemize}
\textsc{Proof}. If $\varphi_\alpha$ is cofinal, we conclude from this by definition that $\varinjlim_I \psi\varphi_\alpha \longrightarrow \varinjlim_{C/F}\psi$ is an isomorphism, now this arrow is none other than $\alpha$. Conversely, suppose that this arrow is an isomorphism, hence $C/F \xleftarrow{\;\sim\;} C/\varinjlim x_i$, and it remains to see that the canonical functor
\DERmark{ENV1121}

\[
I \;\longrightarrow\; C/\varinjlim{}_{C^{\wedge}} x_i
\]
is cofinal, i.e. satisfies F\,1) and F\,2). The verification is $\pm$ tautological, using the fact that for every $y$ in $C$, we have
\[
\Hom(y,\, \varinjlim{}_{C^{\wedge}} x_i) \;\xleftarrow{\;\sim\;}\; \varinjlim{}_{\mathrm{Ens}}\Hom(y,x_i),
\]
and the ``calculation'' of filtering direct limits in Ens.

% margin: vertical bar in the left margin beside Corollaire 4.4.3
\textbf{Corollary 4.4.3.} \textit{Let $\pi$ be an infinite ordinal. For $F$ to be in} $\Ind'_\pi(C)$ \textit{(essential image of} $\Ind_\pi(C)$ \textit{in $C^{\wedge}$), it is necessary and sufficient that there exist an ordered category $I$ large relative to $\pi$, and a cofinal functor $I \longrightarrow C/F$.}

We must therefore give conditions, on a category $J$ $(= C/F)$, for there to exist an ordered category $I$ large relative to $\pi$, and a cofinal functor $I \longrightarrow J$.

\msp{page 81}
Such a category will be called \textit{large relative to} $\pi$. (It is immediate that in the case where $J$ is itself ordered, this terminology is consistent with that already adopted for ordered sets, i.e. that if there exists an increasing cofinal map $\varphi : I \longrightarrow J$, with $I$ large relative to $\pi$, then $J$ is itself large relative to $\pi$.) Such a category is
\printedpage{17}
% head: A. Grothendieck : Dérivateurs, Ch. XVIII, Édition des Universités de Montpellier II et Paris VII
automatically filtering (SGA 4 I 8.1.3 b)). Let us also note that if we have a cofinal functor $J \longrightarrow J'$, if $J$ is large relative to $\pi$, $J'$ is too (by transitivity of cofinality).

\textbf{Proposition 4.4.4.} \textit{Let $J$ be a $\mathfrak{U}$-category, $\pi$ an infinite cardinal. Equivalent conditions:}
\begin{itemize}[leftmargin=4em]
\item[\textit{a)}] \textit{$J$ is large relative to $\pi$, i.e. there exists an ordered set large relative to $\pi$ and a cofinal functor $\varphi : I \longrightarrow J$.}
\item[\textit{b)}] \textit{$J$ satisfies the following two conditions:}
\begin{itemize}[leftmargin=4em]
\item[$\mathrm{PF}_{\pi}\,1$] $\mathrm{Ord}\,J$ \textit{is large relative to $\pi$, i.e. for every family $(y_\alpha)_{\alpha\in A}$ of objects of $J$, with} $\card A \leqslant \pi$\textit{, there exists an upper bound $y$ in $J$ (i.e. a $y$ and arrows $y_\alpha \xrightarrow{\;v_\alpha\;} y$).}
\item[$\mathrm{PF}_{\pi}\,2$] \textit{For $y_0$, $y_1$ in $J$ and every set $(u_\alpha)_{\alpha\in A}$ of arrows $y_0 \xrightarrow{\;u_\alpha\;} y_1$ with} $\card A \leqslant \pi$\textit{, there exists an arrow $y_1 \xrightarrow{\;u'\;} y_2$ which equalizes them, i.e. such that $u'u_\alpha = u'u_\beta \;\forall\, \alpha,\beta \in A$. (It suffices to require this when} $\card A \geqslant 2$.\textit{)}
\end{itemize}
\item[\textit{c)}] \textit{For every subset $A$ of} $\Fl J$\textit{, with} $\card A \leqslant \pi$\textit{, exists a subcategory $J'$ of $J$, with} $\card\Fl J' \leqslant \pi$\textit{, such that $A \subset \Fl J'$ and $J'$ has a final object.}
\item[\textit{d)}] \textit{Let $I$ be the set of subcategories $J'$ of $J$ such that} $\card J' \leqslant \pi$ \textit{and having a final object}

\msp{page 82}
\textit{and let us order $I$ by inclusion. Then $I$ is large relative to $\pi$, moreover the canonical functor $I \longrightarrow J$ $(J' \longmapsto e_{J'} =$ final object of $J')$ is cofinal} (\footnote{Introduce $I(J)$ before the statement of the proposition. \textbf{N.B.} The cofinality condition $I(J) \longrightarrow J$ is not automatic, as we see by taking $J = B_G$, $G = \{\pm 1\}$.},\footnote{\textbf{N.B.} The cofinality hypothesis is superfluous if $J$ is rigid, i.e. every endomorphism of an object of $J$ is the identity.},\footnote{\ Ed. note.\ \ In this condition (d), the ``\textit{moreover the functor \dots\ is cofinal}'' is added by Grothendieck \textit{a posteriori}, which explains his vain attempt to establish it in the implication d $\Longrightarrow$ a below, an implication which becomes obvious after this addition. On the other hand, conditions (c) and (d) are \textit{not} implied by the equivalent conditions (a) and (b). For example, if $J$ is the category having a single object and a single nonidentity morphism $p$ such that $p^2 = p$, then $J$ satisfies conditions (a) and (b), but \textit{not} conditions (c) and (d).}).
\end{itemize}
(\footnote{Other conditions:
\begin{itemize}[leftmargin=3em]
\item[e)] For every set $A$ with $\card A \leqslant \pi$, $\mathrm{diag}_J : J \longrightarrow J^A$ is cofinal.
\item[e$'$)] (If $A_0$ is a given set, of cardinality $= \pi$.) $J \longrightarrow J^{A_0}$ is cofinal.
\end{itemize}
See also 4.9.6.9.1 and 4.9.6.9.2, pages 237, 238.})

\textsc{Proof}.
\begin{itemize}[leftmargin=7em,labelwidth=6em,align=left]
\item[a $\;\Longrightarrow\;$ b] immediate verification, using the criterion F 1, F 2 for cofinality for $I$ filtering.
\item[b $\;\Longrightarrow\;$ c] by AQT \texttt{[âne qui trotte: at a donkey's trot, i.e. by routine verification]}.
\item[c $\;\Longrightarrow\;$ d] by AQT \texttt{[âne qui trotte: at a donkey's trot, i.e. by routine verification]}.
\end{itemize}
% foot: Version 17 décembre 2025

\printedpage{18}
% head: A. Grothendieck : Dérivateurs, Ch. XVIII, Édition des Universités de Montpellier II et Paris VII
\begin{itemize}[leftmargin=7em,labelwidth=6em,align=left]
\item[d $\;\Longrightarrow\;$ a] Let us consider the ordered set large relative to $\pi$ constructed in the statement of d), and the \textit{functor}
\[
I \;\longrightarrow\; J
\]
given by $J' \longmapsto e_{J'}$ $=$ final object of $J'$. (Verification that it is a functor immediate.) This functor is surjective, hence satisfies the cofinality condition F 1. Condition F 2 also follows from d) by looking at the three objects $J'$, $j \xrightarrow{\;u_1\;} e_{J'}$, $j \xrightarrow{\;u_2\;} e_{J'}$ of $I$,
\[
\begin{tikzcd}[column sep=small, row sep=small]
J', & e_{J'} \\
j \arrow[ur, shift left, "u_1"] \arrow[ur, shift right, "u_2"'] &
\end{tikzcd}
\]
and by taking a $J''$ which is an upper bound for all three, then $e_{J'} \longrightarrow e_{J''}$ equalizes $u_1$, $u_2$ (\footnote{A bit brief, the case where $e_{J'} = j$ escapes this argument, and requires that one \texttt{[incomplete sentence]}}).
\end{itemize}

% margin: vertical bar in the left margin beside Corollaire 4.4.5
\textbf{Corollary 4.4.5.} \textit{Let $F \in \Ob C^{\wedge}$. For $F$ to belong to} $\Ind'_\pi(C)$\textit{, it is necessary and sufficient that $C/F$ satisfy the set of four equivalent conditions of 4.4.4.}

\textbf{Corollary 4.4.6.} \textit{Let $J$ be a category stable under} $\varinjlim$ \textit{of cardinality $\leqslant \pi$. Then $J$ is large relative to $\pi$.}

Indeed, conditions $\mathrm{PF}_{\pi}\,1$, $\mathrm{PF}_{\pi}\,2$ are visibly satisfied.

\msp{page 83}
\textbf{4.4.6.} I give myself a functor % [sic] 4.4.6. -> 4.4.7. (?), numéro déjà employé par le Corollaire 4.4.6 qui précède
\[
J \xrightarrow{\;\varphi\;} \Ind'_\pi C \;, \tag*{(4.4.6.1)}
\]
$J$ a small category, and we look for a $\varinjlim$ in $\Ind'_\pi C$, postulating if need be the existence of $\varinjlim$ of cardinality $\leqslant \pi$ in $C$, and the hypothesis $\card\Fl J \leqslant \pi$. The most natural procedure for reducing the question of the existence of $\varinjlim$ in $\Ind'_\pi C$, to the similar question in $C$, is made explicit in SGA 4 I 8.8.3. We consider the fully faithful inclusion
\[
C \;\hookxrightarrow{\;\alpha\;}\; \Ind'_\pi C \;, \tag*{(4.4.6.2)}
\]
which commutes with $\varinjlim$ of cardinality $\leqslant \pi$ (those which are representable in $C$, cf. p. 78 and reference to SGA 4 I 9.8), whence
\[
\uHom(J,C) \;\hookxrightarrow{\;\alpha\;}\; \uHom(J,\Ind'_\pi C) \tag*{(4.4.6.3)}
\]
also fully faithful, and commuting with $\varinjlim$ of cardinality $\leqslant \pi$ (representable in the first member). This arrow fits into the commutative square relative to the datum of a $j \in \Ob J$, and of the inclusion functors
\[
e \xrightarrow{\;\varepsilon_j\;} J \;, \qquad \varepsilon_j(e) \;=\; j \;; \tag*{(4.4.6.4)} % [sic] e_j(e) -> \varepsilon_j(e), la flèche est notée \varepsilon_j
\]
\DERmark{ENV1122}

% foot: Version 17 décembre 2025


\endgroup
% END UNIT der-ch18/der-ch18_p016-020.tex
% BEGIN UNIT der-ch18/der-ch18_p021-025.tex
\begingroup\makeatletter%
\newcommand{\printedpage}[1]{\marginpar{\footnotesize p.~#1}}%
\newcommand{\msp}[1]{\marginpar{\footnotesize [#1]}}%
\newcommand{\uHom}{\underline{\mathrm{Hom}}}%
\newcommand{\Hom}{\mathrm{Hom}}%
\newcommand{\Ind}{\mathrm{Ind}}%
\newcommand{\Ob}{\operatorname{Ob}}%
\newcommand{\Fl}{\operatorname{Fl}}%
\newcommand{\card}{\operatorname{card}}%
\makeatother\let\printedpage\derpp\let\msp\dermsp%
\let\mspin\DERmspin%


\printedpage{19}
% head: A. Grothendieck : Dérivateurs, Ch. XVIII, Édition des Universités de Montpellier II et Paris VII
% foot: Version 17 décembre 2025

\msp{page 84}
\[
\begin{tikzcd}
\uHom(J,C) \arrow[r, "\alpha^J"] \arrow[d, "\varepsilon_j^*"] & \uHom(J,\Ind'_\pi C) \arrow[d, "{\varepsilon_j^{*\prime}}"] & \ni \varphi \arrow[d, mapsto] \\
C \arrow[r, "\alpha"] & \Ind'_\pi(C) & \ni \varphi(j)\ ,
\end{tikzcd}
\qquad (\text{\footnotemark[21]})
\tag*{(4.4.6.2)}
\]
\footnotetext[21]{The functors $\varepsilon_j^*$, $\varepsilon_j^{*\prime}$ are the evaluation functors at $j$.}
which induces a functor
\[
\underbrace{\uHom(J,C)/\varphi}_{\overset{\mathrm{def}}{=}\, I} \xrightarrow{\varepsilon_j^\varphi} C/\varphi(j)\ . \tag*{(4.4.6.3)} % [sic] \uHom(J/C) -> \uHom(J,C), everywhere else written Hom(J,C)
\]
\DERmark{ENV1123}\DERmark{ENV1124}

If $C$ is stable under $\varinjlim$ relative to a set $K$ of the type of conditions $\leqslant \pi$ (e.g. finite $\varinjlim$), the same holds for $\uHom(J,C)$, and also for $\uHom(J,C)/\varphi$, and for $C/\varphi(j)$, by virtue of the fact that $\alpha^J$ and $\alpha$ commute with the $\varinjlim$ of the type envisaged, moreover (4.4.6.3) commutes with the limits in question. Let us therefore suppose $C$ stable at least under finite $\varinjlim$, so the source and target categories in (4.4.6.3) are filtering. (The fact that the functor commutes with finite $\varinjlim$ does not seem useful in the present context.) This will serve us to have the convenient cofinality criterion SGA 4 I 8.1.3 b) in terms of F\,1, F\,2.

\msp{page 85}
Let us recall that the cofinality of the functors (4.4.6.3), when it holds ( ?), implies the fact that $\varphi$ is a direct limit of the composite functor
\[
\begin{tikzcd}[row sep=small]
I \arrow[r, "\psi_J"] \arrow[d, equal] & \uHom(J,C) \arrow[r, "\alpha^J"] \arrow[dd, dotted, "\varepsilon_j^*"] & \uHom(J,\Ind'_\pi(C)) \arrow[dd, dotted, "{\varepsilon_j^{*\prime}}"] \\
\uHom(J,C)/\varphi \arrow[d, "\varepsilon_j^\varphi"'] & & \\
C/\varphi(j) \arrow[r, dotted, "\psi_j"] & C \arrow[r, dotted, "\alpha"] & \Ind'_\pi(C)\ ,
\end{tikzcd}
\tag*{(4.4.6.4)}
\]
along the filtering category $I$, more precisely that $\varphi$ is the $\varinjlim$ ``argument by argument in $j$'', which means that for every $j \in \Ob J$, $\varphi(j) = \varepsilon_j^{*\prime}(\varphi)$ is a direct limit of $\varepsilon_j^{*\prime}(\alpha^J\psi) = (\alpha\psi_j)\varepsilon_j^\varphi$. (Now we do indeed know that $\varphi(j)$ is the limit of $\alpha\psi$, since $\varphi(j) \in \Ob C^\wedge$ is ind-representable, hence also that of $(\alpha\psi_j)\varepsilon_j^\varphi$ if $\varepsilon_j^\varphi$ is cofinal.)
\DERmark{ENV1125}\DERmark{ENV1126}\DERmark{ENV1127}


Thus, admitting the cofinality of (4.4.6.3) and interpreting $\psi_J$ in (4.4.6.4)

\msp{page 86}
\[
\psi_J \ :\ I \longrightarrow \uHom(J,C)
\]
as a functor
\[
I \times J \xrightarrow{\psi} C\ , % [sic] \Psi -> \psi, the functor I x J -> C is called \psi in the next formula and in (4.4.7.1)
\]
\DERmark{ENV1128}

we find the ``indicial representation'' of $\varphi$
\[
\varphi(j) \ =\ \underbrace{\varinjlim_{i\in I} \psi(i,j)}_{\substack{\text{direct limit}\\ \text{taken in } \Ind'_\pi(C)\\ \text{of objects of } C \hookrightarrow \Ind'_\pi(C)}} \qquad ,
\]

\printedpage{20}
% head: A. Grothendieck : Dérivateurs, Ch. XVIII, Édition des Universités de Montpellier II et Paris VII
% foot: Version 17 décembre 2025
whence
\[
\begin{aligned}
\underbrace{\varinjlim_{j\in J}{}_{\,\Ind'_\pi C}\varphi(j)}_{\text{subject to existence}} \quad &= \quad \varinjlim_{j\in J}\, \varinjlim_{i\in I} \psi(i,j)\\
&= \quad \varinjlim_{i\in I}\, \varinjlim_{j\in J} \psi(i,j)\ ,
\end{aligned}
\]
all direct limits being taken in $\Ind'_\pi(C)$. But if the $\varinjlim$ of type $J$ are representable in $C$, and $\card J \leqslant \pi$, we know that the $\varinjlim\limits_{j\in J} \psi(i,j)$ \textit{in} $C$ is also their $\varinjlim$ in $\Ind'_\pi C$, and we therefore find that $\varinjlim\limits_{j\in J} \varphi(j)$ exists in $\Ind'_\pi(C)$ if and only if $\varinjlim\limits_{i\in I} \Psi(i)$ exists in $\Ind'_\pi(C)$, where for $i \in I$, we set
\[
\Psi(i) \ =\ \varinjlim_{j\in J} \psi(i,j) \qquad \textit{in } C.
\]

\msp{page 87}
Let us summarize what we have obtained:

\textbf{Lemma 4.4.7.} \textit{Let $C$ be a small category, $\pi$ an infinite cardinal, $J$ a category such that $\card \Fl J \leqslant \pi$, $\varphi : J \longrightarrow \Ind'_\pi(C)$ a functor. Suppose the $\varinjlim$ of type $J$ to exist in $C$, and that the functors $\varepsilon_j^\varphi$ ($j \in \Ob J$) of (4.4.6.3) are \textup{cofinal} (which is made explicit in a simple way especially if source and target are filtering, something established if $C$ is stable under finite $\varinjlim$). Then $\varphi$ has a ``canonical indicial representation''}
\[
\psi \ :\ I \times J \longrightarrow C \tag*{(4.4.7.1)}
\]
\textit{with}
\[
I \ =\ \uHom(J,C)/\varphi\ , \tag*{(4.4.7.2)}
\]
\textit{i.e. we have the isomorphism functorial in $j \in \Ob J$}
\[
\varphi(j) \ \simeq\ \varinjlim_{I} \psi(i,j) \qquad (\text{\footnotemark[22]})\ . \tag*{(4.4.7.3)}
\]
\footnotetext[22]{direct limit in $\Ind'_\pi$ of objects of $C$.}
\textit{Let us set for every $i \in \Ob I$}
\[
\Psi(i) \ =\ \varinjlim_{J} \psi(i,j) \qquad (\text{\footnotemark[23]})\ . \tag*{(4.4.7.4)}
\]
\footnotetext[23]{direct limit taken in $C$, but valid also in $\Ind'_\pi C$.}
\textit{Then $\varinjlim\limits_{j\in J} \varphi(j)$ (in $\Ind'_\pi C$) exists if and only if}
\msp{page 88}
\textit{the $\varinjlim\limits_{i\in I} \Psi(i)$ exists (in $\Ind'_\pi C$ as well), and the two are equal.}

\textbf{Corollary 4.4.8.} \textit{Suppose that $C$ is stable under $\varinjlim$ of cardinality $\leqslant \pi$ (not only under those of type $J$). Then $\uHom(J,C)$ and $I = \uHom(J,C)/\varphi$ have the same property, $I$ is large relative to $\pi$, and $\varinjlim\limits_{j\in J} \varphi(j) \ =\ \varinjlim\limits_{i\in I} \Psi(i)$ exists in $\Ind'_\pi(C)$, subject to cofinality of the functors (4.4.6.3).}

\printedpage{21}
% head: A. Grothendieck : Dérivateurs, Ch. XVIII, Édition des Universités de Montpellier II et Paris VII
% foot: Version 17 décembre 2025
Indeed it suffices to verify this

\textbf{Lemma 4.4.8.1.} \textit{Let $I$ be a category large relative to $\pi$, and $\Psi : I \longrightarrow C$ a functor. Then $\varinjlim\limits_{i} \Psi(i)$ exists in $\Ind'_\pi(C)$, and it is the direct limit in $C^\wedge$, or again ``it is'' the ind-object large relative to $\pi$ represented by $\Psi$.} (In other words (\footnote[24]{This is an abuse, this stability is a distinctly stronger statement! I prove it on page 99.}), $\Ind'_\pi(C) \subset C^\wedge$ is a strictly full subcategory \textit{stable} under $\varinjlim$ large relative to $\pi$.)

\msp{page 89}
This lemma is a tautology: if $F$ denotes the direct limit in $C^\wedge$, it is clear by definition (taking a cofinal functor $I_0 \longrightarrow I$, with $I_0$ \textit{ordered} and large relative to $\pi$, so that the $\varinjlim$ can be taken along $I_0$) that $F$ is in $\Ind'_\pi(C)$, and then it is a fortiori a $\varinjlim$ in $\Ind'_\pi C$.

To verify the existence of the $\varinjlim$ of cardinality $\leqslant \pi$ in $\Ind'_\pi C$, in the case where $C$ itself is stable under the said limits, it therefore suffices, thanks to Corollary 4.4.8, to prove the cofinality of the functors (4.4.6.3). This will follow from the

\textbf{Lemma 4.4.9.} \textit{Let $C$ be a small category, $\pi$ an infinite cardinal, $J$ a small category such that for all $j, j' \in \Ob J$, we have $\card \Hom(j,j') \leqslant \pi$. Under these conditions, for every $j \in \Ob J$, the functor (4.4.6.3) \textup{admits a left adjoint,} hence a fortiori it is cofinal (since the categories colocalized over $C/\varphi(j)$ have an initial object, a fortiori they are $0$-connected).}

\msp{page 90}
\textsc{Proof of} 4.4.9. Let us return to the diagram (4.4.6.2), from which we deduce (4.4.6.3), we are going to rewrite it in a simplified way:
\DERmark{ENV1129}

\[
\begin{tikzcd}
B \arrow[r, "\beta"] \arrow[d, shift left, "\varepsilon"] & B' \arrow[d, shift left, "\varepsilon'"] & \ni y'_0 \arrow[d, mapsto] \\
C \arrow[r, "\gamma"] \arrow[u, dotted, shift left, "\lambda"] & C' \arrow[u, dotted, shift left, "\lambda'"] & \ni z'_0\ .
\end{tikzcd}
\tag*{(4.4.9.1)}
\]
I say that $\varepsilon$, $\varepsilon'$ have \textit{left adjoints} $\lambda$ resp. $\lambda'$, and that the canonical ``base change'' homomorphism
\[
\lambda'\gamma \longrightarrow \beta\lambda \tag*{(4.4.9.2)}
\]
is an \textit{isomorphism.} Indeed, the existence of the adjoints of $\varepsilon$, $\varepsilon'$, namely of the functors $\varepsilon_{j\,!}$ and $\varepsilon'_{j\,!}$ relative to the inclusion
\[
\varepsilon_j \ :\ e \longrightarrow J
\]
and to the functor categories $\uHom(J,C)$ resp. $\uHom(J,\Ind'_\pi C)$, resulting from the existence of sums of type $\Hom(j,j')$ ($j' \in \Ob J$) in $C$ (already established), \textit{and in} $\Ind'_\pi C$\textit{ ;} and the fact that the functor $C \longrightarrow \Ind'_\pi C$ commutes
\msp{page 91}
with the said sums (something also already established, since $\card\, (\Hom(j,j')) \leqslant \pi$) then implies that (4.4.9.2) is an isomorphism. Subject to the pending verification (the existence % [sic] (4.4.4.2) -> (4.4.9.2), the isomorphism in question is the change-of-base arrow (4.4.9.2)
\DERmark{ENV1130}

\printedpage{22}
% head: A. Grothendieck : Dérivateurs, Ch. XVIII, Édition des Universités de Montpellier II et Paris VII
% foot: Version 17 décembre 2025
of sums of cardinality $\leqslant \pi$ in $\Ind'_\pi(C)$), Lemma 4.4.9 will follow from the more general result:

\textbf{Lemma 4.4.10.} \textit{Consider a commutative square (4.4.9.1) (solid lines) of categories, where $\varepsilon$ admits a left adjoint $\lambda$, $\varepsilon'$ a left adjoint $\lambda'$, and the base change arrow (4.4.9.2) is an isomorphism. Then for every $y'_0$ in $B'$, setting $\varepsilon'(y'_0) = z'_0$, the functor induced by $(\varepsilon, \varepsilon')$,}
\[
\overline{B} \ =\ B/y'_0 \xrightarrow{\bar{\varepsilon}} \overline{C} = C/z'_0\ ,
\]
\textit{admits a left adjoint $\overline{\lambda}$, given by}
\[
\begin{tikzcd}[column sep=small]
\overline{\lambda}(\gamma(z) \xrightarrow{v} z'_0) \ \ = & \Bigl(\lambda'\gamma(z) \simeq \beta(\lambda(z)) \arrow[r, "{\lambda'(v)}"] \arrow[ddr, "{\substack{\overline{\lambda}(v) \,=\, \text{composite}\\ \text{adj} \circ \lambda'(v)}}"'] & \lambda'(z'_0) \arrow[d, equal] \\
 & & \lambda'\varepsilon'(y'_0) \arrow[d, "\text{adj.}"] \\
 & & y'_0\Bigr)\ \ .
\end{tikzcd}
\]

\msp{page 92}
The proof is carried out by a tedious AQT \texttt{[âne qui trotte: at a donkey's trot, i.e. by routine verification]} (I confess not to have pushed to the last detail \ldots).

To finish the proof of 4.4.9, there remains to prove one last ( ?)

\textbf{Lemma 4.4.11.} \textit{Let $C$ be a small category, stable under sums of cardinality $\leqslant \pi$. Then $\Ind'_\pi C$ is stable under sums of cardinality $\leqslant \pi$ (and, for the record, the functor $\alpha : C \longrightarrow \Ind'_\pi C$ commutes with them).}

Applying Corollary 4.4.8 to Lemma 4.4.7, we are reduced to proving the cofinality of (4.4.6.3) when $J$ is discrete, plus the fact that $\uHom(J,C)/\varphi$ is large relative to $\pi$, in this case. But giving $\varphi$ amounts to giving a family $(\varphi_j)_{j\in J}$ of objects of $\Ind'_\pi(C)$, and the functor (4.4.6.3) is made explicit as the projection
\[
\prod_{j\in J} (C/\varphi_j) \longrightarrow C/\varphi_j
\]
from the total product onto one of the factors. So we are reduced to this:

\msp{page 93}
\textbf{Lemma 4.4.12.} \textit{Let $(C_j)_{j\in J}$ be a family of small categories large relative to $\pi$, indexed by a small set $J$. Then $\prod C_j$ is large relative to $\pi$, and the projections (for $j_0 \in J$) $\prod C_j \longrightarrow C_{j_0}$ are cofinal.}

Let us first prove cofinality, which is the most important for us. This is obvious if $J = \{j_0\}$, otherwise we have $J = J' \cup \{j_0\}$, and $\prod_{j\in J} C_j \simeq C' \times C_{j_0}$, where $C' = \prod_{j\in J'} C_j$, so we must see that $C' \times C_{j_0} \longrightarrow C_{j_0}$ is cofinal, i.e. that every category $\xi\backslash(C' \times C_{j_0})$ (for $\xi \in \Ob C_{j_0}$) is 0-connected, now it is isomorphic to $C' \times (\xi\backslash C_{j_0})$, and $\xi\backslash C_{j_0}$ being 0-connected, it remains to see that $C'$ is 0-connected. We should be careful in this regard that an \textit{infinite} product of 0-connected categories is not necessarily 0-connected, so cofinality is not completely tautological! But if we prove that the category $\prod_{j\in J} C_j$ is
\printedpage{23}
% head: A. Grothendieck : Dérivateurs, Ch. XVIII, Édition des Universités de Montpellier II et Paris VII
% foot: Version 17 décembre 2025
\msp{page 94}
large relative to $\pi$, hence filtering, this implies that it is 0-connected, and we can apply this same result to $C' = \prod_{j\in J'} C_j$, and we win.

So we are reduced to proving that $\prod C_j$ is large relative to $\pi$. This can be seen either from the criterion $\mathrm{PF}_\pi\,1$, $\mathrm{PF}_\pi\,2$ of 4.4.4 b, or by reducing to the case where the $C_j$ are ordered, but that looks longer. Let us note that we prove in the same way that a (small) product of filtering categories is again a filtering category.

\textbf{Remarks.}
\begin{enumerate}[label={}, leftmargin=*]
\item[1.] This argument in no way shows (even if $C$ is stable under $\varinjlim$ of cardinality $\leqslant \pi$) that for every $\varphi \in \uHom(J,\Ind(C))$, the functor similar to (4.4.6.3)
\[
\uHom(J,C)/\varphi \longrightarrow C/\varphi(j)
\]
is cofinal, and I doubt that it always is. The difficulty comes from the fact that $C \longrightarrow \Ind\, C$ does not commute with sums of cardinality $\leqslant \pi$, only with finite sums,
\msp{page 95}
so Lemma 4.4.10 cannot be invoked.
\item[2.] We have just proved, via 4.4.8 and 4.4.9 (whose proof is completed with 4.4.12), that if $C$ is a small category stable under $\varinjlim$ of cardinality $\leqslant \pi$, then $\Ind_\pi C$ is also stable under these limits. Moreover, the direct systems of this type $J$ in $\Ind_\pi C$ admit an indicial representation
\[
\psi \ :\ I \times J \longrightarrow C\ ,
\]
with $I$ large relative to $\pi$. And we can always suppose, at the cost of taking $I' \longrightarrow I$ cofinal with $I'$ an ordered set large relative to $\pi$, that $I$ is moreover an ordered set.
\end{enumerate}

That is the means then, for the proof of the

\textbf{Theorem 4.4.13.} \textit{Let $C$ be a category equivalent to a small category, and stable under $\varinjlim$ of cardinality $\leqslant \pi$, where $\pi$ is a given infinite cardinal. Then:}
\msp{page 96}

\begin{enumerate}[label={}, leftmargin=*]
\item[\textit{a)}] \textit{(For the record.) The canonical functor}
\[
C \longrightarrow \Ind_\pi(C) \xrightarrow{\ \sim\ } \Ind'_\pi(C)
\]
\textit{commutes with $\varinjlim$ of cardinality $\leqslant \pi$.}
\item[\textit{b)}] \textit{(For the record.) The direct systems of cardinality $\leqslant \pi$ in $\Ind_\pi(C)$ (i.e. the functors $\varphi : J \longrightarrow \Ind_\pi(C)$, with $\card \Fl(J) \leqslant \pi$) admit an indicial representation}
\[
\psi \ :\ I \times J \longrightarrow C\ ,
\]
\textit{with $I$ an ordered set large relative to $\pi$, and admit a direct limit in $\Ind_\pi C$ which are computed by}
\[
\underbrace{\varinjlim_{j\in J} \varphi(j)}_{\text{in } \Ind_\pi(C)} \ \simeq\ \text{``}\,\varinjlim_{i\in I}\,\text{''}\Psi(i)
\]
\end{enumerate}


\endgroup
% END UNIT der-ch18/der-ch18_p021-025.tex
% BEGIN UNIT der-ch18/der-ch18_p026-030.tex
\begingroup\makeatletter%
\newcommand{\printedpage}[1]{\marginpar{\footnotesize p.~#1}}%
\newcommand{\msp}[1]{\marginpar{\footnotesize [#1]}}%
\newcommand{\Ind}{\operatorname{Ind}}%
\newcommand{\uHom}{\underline{\mathrm{Hom}}}%
\newcommand{\Ob}{\operatorname{Ob}}%
\newcommand{\card}{\operatorname{card}}%
\newcommand{\Fl}{\mathrm{Fl}}%
\newcommand{\PF}{\mathrm{PF}}%
\makeatother\let\printedpage\derpp\let\msp\dermsp%
\let\mspin\DERmspin%


\printedpage{24}
% head: A. Grothendieck : Dérivateurs, Ch. XVIII, Édition des Universités de Montpellier II et Paris VII
% foot: Version 17 décembre 2025

\textit{with}
\[
\Psi(i) \;=\; \varinjlim_{j\in J}{}_{C}\,\psi(i,j)\;.
\]
\textit{c)} $\Ind_\pi C$ \textit{is stable under small direct limits} (\footnote{also add, in d), corollary 4.4.16, page 101.}).

It remains to establish this last point. Since we already have finite $\varinjlim$ by b, it remains to prove the existence of small filtering $\varinjlim$, or, if preferred, of infinite sums, which will be nicer in fact. But
\msp{page 97}
we can just as well treat the general case of a $\varinjlim$ of arbitrary type $J$ ($J$ small). For this, let
\[
I \;=\; I(J) \;=\; \begin{aligned}[t]
&\text{set of subcategories } J' \text{ of } J \text{ such that}\\
&\card\Fl(J')\leqslant\pi,
\end{aligned}
\]
let us order $I$ by inclusion. It is obvious that $I\neq\emptyset$ (for $\emptyset\in I$), and even that $I$ \textit{is large relative to} $\pi$, for if $(J'_\alpha)_{\alpha\in A}$ with $\card(A)\leqslant\pi$ is a family of elements of $I$, then the subcategory of $J$ generated by the $J'_\alpha$ is still in $I$ (immediate verification), and is an upper bound for the $J'_\alpha$. In fact, the consideration of $I(J)$ associated with $J$ will allow us to prove this, which generalizes 4.4.13 c), taking b) into account:

\textbf{Corollary 4.4.14} (\footnote{Compare also 4.11.3, page 270.}). \textit{Let $\mathcal{M}$ be a category, $\pi$ an infinite cardinal such that}
\begin{enumerate}
\item[a)] \textit{$\mathcal{M}$ is}
\msp{page 98}
\textit{stable under $\varinjlim$ of cardinality $\leqslant\pi$, and}
\item[b)] \textit{$\mathcal{M}$ is stable under $\varinjlim$ over ordered sets $I$ large relative to $\pi$.}
\end{enumerate}
\textit{Then $\mathcal{M}$ is stable under small $\varinjlim$ .}

(\footnote{\textbf{Corollary 4.4.14.1.} $\mathcal{M}$ being as in 4.4.14, and $f:\mathcal{M}\longrightarrow\mathcal{E}$ being a functor which commutes with $\varinjlim$ of cardinality $\leqslant\pi$, and with small $\varinjlim$ large relative to $\pi$, then $f$ commutes with small $\varinjlim$ .}).

Indeed, if $\varphi:J\longrightarrow\mathcal{M}$ is given, we construct from it
\[
\psi\;:\;I\;=\;I(J)\;\longrightarrow\;\mathcal{M}
\]
by
\[
\psi(J')\;=\;\varinjlim_{J'}\varphi(J')\;,
\]
which exists by the hypothesis on $\mathcal{M}$. It is then immediate that $\varinjlim\limits_{J}\varphi$ exists if and only if $\varinjlim\limits_{I}\psi$ exists, now this existence is assured by the fact that $I$ is large relative to $\pi$, and the hypothesis on $\mathcal{M}$.

To finish proving theorem 4.4.13, we are thus reduced to the particular case of direct limits $\varinjlim\limits_{J}\varphi(j)$, when $\varphi:J\longrightarrow\Ind'_\pi C$ is given, with $J$ an
\printedpage{25}
% head: A. Grothendieck : Dérivateurs, Ch. XVIII, Édition des Universités de Montpellier II et Paris VII
% foot: Version 17 décembre 2025
\msp{page 99}
ordered set large relative to $\pi$. But this will follow from the more precise result:

\textbf{Corollary 4.4.15} (\footnote{compare with 4.7.3.1, which ought to be made explicit here.}). \textit{Let $C$ be a small category. Then the strictly full subcategory} $\Ind'_\pi C$ \textit{of $C^\wedge$ is stable in $C^\wedge$ under $\varinjlim$ large relative to $\pi$, in other words, these $\varinjlim$} exist in $\Ind'_\pi(C)$, \textit{and the inclusion functor commutes with them.} % doubt resolved: existent is upright in the PDF; see corrections.csv, PDF leaf 27

\textsc{Proof}. Let $I$ be an ordered set, large relative to $\pi$, and $i\longmapsto F_i$ a functor $I\longrightarrow\Ind'_\pi(C)$, $F\;=\;\varinjlim\limits_{I}F_i$ in $C^\wedge$, let us prove that $F\in\Ob\Ind'_\pi(C)$, i.e. (4.4.3) that $C/F$ is large relative to $\pi$. I use criterion b) of 4.4.4, we have to check $\PF_\pi\,1$, $\PF_\pi\,2$. Let
\[
x_\alpha\xrightarrow{\;u_\alpha\;}F\;=\;\varinjlim F_i
\]
be objects of $C/F$ indexed by $A$ with $\card A\leqslant\pi$. Each $u_\alpha$ comes from $x_\alpha\longrightarrow F_{i_\alpha}$, and taking a common upper bound for the $i_\alpha$ ($I$ large relative to $\pi$), we find an $i\in I$ and $v_\alpha:x_\alpha\longrightarrow F_i$. Since $F_i$ is in $\Ind'_\pi(C)$, i.e. $C/F_i$ large relative to $\pi$, considering the $v_\alpha$ as objects of $C/F_i$, they are bounded above by an object $x\longrightarrow F_i$ of $C/F_i$. We thus have
\[
\begin{tikzcd}
x_\alpha \arrow[r] \arrow[d] & F_{i_\alpha} \arrow[d] \arrow[dr, dotted] & \\
x \arrow[r] & F_i \arrow[r] & F\;,
\end{tikzcd}
\]
which shows that the $x_\alpha\longrightarrow F$ are bounded above by $x\longrightarrow F$,
\msp{page 100}
which proves $\PF_\pi\,1$. Let us prove $\PF_\pi\,2$, i.e. let $x\xrightarrow{\;u\;}F$, $y\xrightarrow{\;v\;}F$ be two objects of $C/F$, and $(f_\alpha)_{\alpha\in A}$ a family of cardinality $\leqslant\pi$ of arrows from $u$ to $v$. We know that $u$, $v$ come from $u_i:x\longrightarrow F_i$, $v_i:y\longrightarrow F_i$ for $i$ large enough.
\[
\begin{tikzcd}[column sep=small]
x \arrow[rr, "f_\alpha"] \arrow[dr, "u"'] & & y \arrow[dl, "v"] \\
& F &
\end{tikzcd}
\]
\[
\begin{tikzcd}[column sep=small]
x \arrow[rr, "f_\alpha"] \arrow[dr, "u_i"'] & & y \arrow[dl, "v_i"] \\
& F_i &
\end{tikzcd}
\]

For every $\alpha$, we can choose $i_\alpha\geqslant i$ such that $v_if_\alpha$ and $u_i$ become equal on composing with $F_i\longrightarrow F_{i_\alpha}$, and since $\card A\leqslant\pi$ and $I$ large relative to $\pi$, there exists an upper bound $i'$ for the $i_\alpha$. All told, we can assume (replacing $i$ by $i'$) that $v_if_\alpha=u_i$ for every $\alpha$. Then $u_i$, $v_i$ being considered as objects of $C/F_i$, the $f_\alpha$ are arrows of $C/F_i$ from $u_i$ to $v_i$. Since $C/F_i$ is large relative to $\pi$ (by the hypothesis $F_i\in\Ob\Ind'_\pi(C)$), we find that there exists an equalizer $(y,v_i)\xrightarrow{\;g\;}(z,w_i)$.
\printedpage{26}
% head: A. Grothendieck : Dérivateurs, Ch. XVIII, Édition des Universités de Montpellier II et Paris VII
% foot: Version 17 décembre 2025
\[
\begin{tikzcd}
x \arrow[r, shift left=2, "(f_\alpha)"] \arrow[r] \arrow[r, shift right=2] \arrow[dr, "u_i"'] & y \arrow[r, "g"] \arrow[d, "v_i"'] & z \arrow[dl, "w_i"] \\
& F_i &
\end{tikzcd}
\]

Then the image of $g$ in $C/F$ (via $F_i\longrightarrow F$) equalizes the arrows $f_\alpha$, as arrows of $C/F$, which proves $\PF_\pi\,2$ and completes the proof of 4.4.15, and thereby also that of theorem 4.4.13, and of lemma 4.3. This also completes the proof of the first assertion of lemma
\DERmark{ENV1131}

\msp{page 101}
4.3.7 (p. 77), which completes the proof of theorem 4.3.6 (p. 70-71). It remains to establish the universal property stated in 4.3.7, which I am going to recall:

\textbf{Corollary 4.4.16} (\footnote{Could be incorporated as part d) in theorem 4.4.13, for it is an important complement, and not so obvious after all.}). \textit{Let $C$ be a category equivalent to a small category, stable under $\varinjlim$ of cardinality $\leqslant\pi$, where $\pi$ is a given infinite cardinal, and let us consider the canonical inclusion functor}
\[
C\xrightarrow{\;\alpha\;}\Ind_\pi(C)\;. \tag*{(4.4.16.1)}
\]
\textit{This functor is} 2\textit{-universal for functors $f:C\longrightarrow\mathcal{E}$, where $\mathcal{E}$ is a $\mathcal{U}$-category stable under small $\varinjlim$ , and $f$ a functor commuting with these direct limits. In other words, we have}
\begin{enumerate}
\item[\textit{1$^\circ$)}] \textit{$\alpha$ has the preceding properties, i.e.} $\Ind_\pi(C)$ \textit{is stable under $\varinjlim$ , and the functor $\alpha$ commutes with $\varinjlim$ of cardinality $\leqslant\pi$.}
\item[\textit{2$^\circ$)}] \textit{For every $\mathfrak{U}$-category $\mathcal{E}$ stable under small direct limits, the functor}
\[
\begin{aligned}
g\;&\longmapsto\;g\circ\alpha\\
\uHom_{!}(\Ind_\pi C,\mathcal{E})\;&\longrightarrow\;\uHom_{!\,\pi}(C,\mathcal{E})
\end{aligned}
\tag*{(4.4.16.2)}
\]
\textit{is an equivalence of categories, where the subscript $!$ (resp. $!\,\pi$) indicates the strictly full subcategory of the category} $\uHom$
\msp{page 102}
\textit{under consideration, formed by the functors which commute with small $\varinjlim$ (resp. with those of cardinality $\leqslant\pi$).}
\end{enumerate}

We can thus also say that $\Ind_\pi(C)$ plays the role of an ``envelope'' stable under small $\varinjlim$ of the category $C$, when the $\varinjlim$ of cardinality $\leqslant\pi$, which already exist in $C$, are respected by the embedding of $C$ in its envelope. It is a by no means tautological result that it suffices for this to add to $C$ some $\varinjlim$ \textit{large relative to} $\pi$ of objects of $C$, and that more precisely, the envelope category is realized by $\Ind_\pi C$ (or, if preferred and less plethorically, by $\Ind'_\pi C\subset C^\wedge$). % [sic] les plongement -> le plongement, number agreement

\textsc{Proof}. 1°) is contained in 4.4.13 a) and c). Let us therefore prove 2°). It is obvious that $\alpha$ is a fully faithful functor whose essential image strictly generates
\printedpage{27}
% head: A. Grothendieck : Dérivateurs, Ch. XVIII, Édition des Universités de Montpellier II et Paris VII
% foot: Version 17 décembre 2025
$\mathcal{M}=\Ind_\pi C$, whence it follows that the functor (4.4.16.2) is fully faithful, being induced by the fully faithful functor
\[
\alpha^*_{\mathcal{E}}\;:\;\uHom_{!}(\mathcal{M},\mathcal{E})\;\longrightarrow\;\uHom(C,\mathcal{E})
\]
\DERmark{ENV1132}

\msp{page 103}
on the full subcategories under consideration. It thus remains to show that (4.4.16.2) is essentially surjective. So let $f:C\longrightarrow\mathcal{E}$ be in $\uHom_{!\,\pi}(C,\mathcal{E})$, let us consider its canonical extension % [sic] les sous-catégorie pleines -> les sous-catégories pleines, number agreement
\[
\left\{
\begin{aligned}
&Lf\;:\;\Ind(C)\;\longrightarrow\;\mathcal{E}\\
&Lf((X_i)_{i\in I})\;=\;\varinjlim_{i\in I}f(X_i)
\end{aligned}
\right.
\tag*{(4.4.16.3)}
\]
(cf. SGA 4 I 8.7.3), and let
\[
\psi(f)\;=\;Lf|\Ind_\pi(C)\;. \tag*{(4.4.16.4)}
\]
Let us recall that $Lf$ commutes with small filtering $\varinjlim$, but since the inclusion functor
\[
\Ind_\pi C\;\longrightarrow\;\Ind(C) \tag*{(4.4.16.5)}
\]
does not commute with these $\varinjlim$ , it is not clear a priori that $\psi(f)$ commutes with the said limits, nor a fortiori that it commutes with small $\varinjlim$ without more. But we have seen that (4.4.16.5) commutes with $\varinjlim$ \textit{large relative to} $\pi$, which is contained in corollary 4.4.15 (since the \textit{filtering} direct limits in $\Ind(C)$ are the limits \textit{in} $C^\wedge$): to say that $\Ind'_\pi(C)\longrightarrow\Ind'(C)$ commutes with a certain type of filtering $\varinjlim$, amounts to saying that this is so for the composite with $\Ind'_\pi(C)\hookrightarrow C^\wedge$, which is established by 4.4.15. From this it follows that $\psi(f)$ \textit{commutes with} $\varinjlim$ \textit{large relative to} $\pi$. To establish that $g=\psi(f)$ commutes
\msp{page 104}
with small $\varinjlim$ , it thus remains to see that it commutes with $\varinjlim$ of cardinality $\leqslant\pi$, by virtue of 4.4.14.1 (margin page 98). To establish this, we use theorem 4.4.13 b, representing a functor $\varphi:J\longrightarrow\Ind_\pi(C)$ (where $\card\Fl J\leqslant\pi$) in indexed form by
\[
\Phi\;:\;I\times J\;\longrightarrow\;C\;, \tag*{(4.4.16.6)}
\]
$I$ being an ordered set large relative to $\pi$. We thus have
\[
\varphi(j)\;=\;\varinjlim_{i\in I}{}_{\Ind_\pi C}\,\Phi(i,j)\;,
\]
and since $g=\psi(f)$ commutes with $\varinjlim$ of type $I$ by what we have seen, we will have
\[
g(\varphi(j))\;=\;\varinjlim_{i\in I}{}_{(\mathcal{E})}\;g(\Phi(i,j))\;,
\]
whence
\[
\varinjlim_{j\in J}{}_{(\mathcal{E})}\,g(\varphi(j))\;=\;\varinjlim_{j\in J}\;\varinjlim_{i\in I}\;\underbrace{g(\Phi(i,j))}_{=\,f(\Phi(i,j))}\qquad\text{(limits in $\mathcal{E}$)}\;. \tag*{$(*)$}
\]
\printedpage{28}
% head: A. Grothendieck : Dérivateurs, Ch. XVIII, Édition des Universités de Montpellier II et Paris VII
% foot: Version 17 décembre 2025

On the other hand, we have
\[
\begin{aligned}
\varinjlim_{j\in J}{}_{\Ind_\pi}\,\varphi(j)\;&\simeq\;\varinjlim_{j\in J}\;\varinjlim_{i\in I}\;\Phi(i,j)\\
&\simeq\;\varinjlim_{i\in I}\;\underbrace{\varinjlim_{j\in J}\;\Phi(i,j)}_{\substack{\text{can be computed}\\ \text{in $C$, say}\\ \Psi(i)\,\in\,\Ob C}}\;,
\end{aligned}
\]
thus
\[
\varinjlim_{j\in J}{}_{\Ind_\pi}\,\varphi(j)\simeq\varinjlim_{i\in I}\;\Psi(i)\;, \tag*{$(**)$}
\]
\msp{page 105}
whence by applying $g$, which commutes with $\varinjlim$ of type $I$, and now recalling that $g|C\simeq f$,
\[
g(\,\varinjlim_{J}{}_{\Ind_\pi C}\,\varphi(j))\;\simeq\;\varinjlim_{i\in I}\;f(\Psi(i))\;. \tag*{$(*{*}*)$}
\]
We have on the other hand
\[
\Psi(i)\;\overset{\mathrm{def}}{=}\;\varinjlim_{j\in J}\;\Phi(i,j)\;,\qquad\qquad \varinjlim\;\text{ in } C\;,
\]
\DERmark{ENV1133}

and since $f$ commutes with $\varinjlim$ of type $J$ (for $\card\Fl J\leqslant\pi$), we conclude from this
\[
f(\Psi(i))\;\simeq\;\varinjlim_{j\in J}\;f(\Phi(i,j))\;,
\]
thus by $(*{*}*)$
\[
g(\,\varinjlim_{j\in J}\;\varphi(j))\;\simeq\;\varinjlim_{i\in I}\;\varinjlim_{j\in J}\;f(\Phi(i,j))\;. \tag*{$(*{*}{*}*)$}
\]
By composing with $(*)$ and using again the commutativity of $\varinjlim\limits_{I}$ with $\varinjlim\limits_{J}$ , we thus find
\[
g(\,\varinjlim\;\varphi(j))\;\overset{\sim}{\longleftarrow}\;\varinjlim\;g(\varphi(j))\;,
\]
which proves that $g$ commutes with direct limits of type $J$, and finishes proving that $g$ is \texttt{[in]} $\uHom_{!}(\Ind_\pi(C),\mathcal{E})$. Since $g|C\simeq f$, we conclude that the functor (4.4.16.2) is essentially surjective, which finishes proving 4.4.16. Moreover, we have constructed an explicit quasi-inverse functor, $\psi$,
\msp{page 106}
fitting into the commutative diagram
\[
\begin{tikzcd}[column sep=large, row sep=large]
\uHom_{!\,\pi}(C,\mathcal{E}) \arrow[r, "\psi", "(\sim)"'] \arrow[d, hook, "\text{fully faithful incl.}"] & \uHom_{!}(\Ind_\pi(C),\mathcal{E}) \arrow[r] \arrow[d, hook, dotted, "\text{fully faithful}"] & \uHom(\Ind_\pi(C),\mathcal{E}) \\
\uHom(C,\mathcal{E}) \arrow[r, "L", "\substack{\sim\\ \text{(cf. SGA 4 I 8.7.3)}}"'] & \uHom_{!\,\mathrm{filt}}(\Ind(C),\mathcal{E}) \arrow[ur, "R\ \text{(restriction)}"'] &
\end{tikzcd}
\qquad.
\tag*{(4.4.16.6)} % [sic] (4.4.16.6) -> (4.4.16.7)?, the label (4.4.16.6) is already used for Phi on the preceding page
\]
\DERmark{ENV1134}



\endgroup
% END UNIT der-ch18/der-ch18_p026-030.tex
% BEGIN UNIT der-ch18/der-ch18_p031-035.tex
\begingroup\makeatletter%
\newcommand{\printedpage}[1]{\marginpar{\footnotesize p.~#1}}%
\newcommand{\msp}[1]{\marginpar{\footnotesize [#1]}}%
\newcommand{\card}{\operatorname{card}}%
\newcommand{\Fl}{\operatorname{Fl}}%
\newcommand{\Ob}{\operatorname{Ob}}%
\newcommand{\Ind}{\operatorname{Ind}}%
\newcommand{\Filt}{\operatorname{Filt}}%
\newcommand{\Hom}{\operatorname{Hom}}%
\newcommand{\Aut}{\operatorname{Aut}}%
\newcommand{\uHom}{\underline{\mathrm{Hom}}}%
\newcommand{\uFilt}{\underline{\mathrm{Filt}}}%
\newcommand{\circled}[1]{\tikz[baseline=(c.base)]\node[draw,circle,inner sep=1pt](c){#1};}%
\makeatother\let\printedpage\derpp\let\msp\dermsp%
\let\mspin\DERmspin%


\setcounter{footnote}{29}

\printedpage{29}
% head: A. Grothendieck : Dérivateurs, Ch. XVIII, Édition des Universités de Montpellier II et Paris VII
% foot: Version 17 décembre 2025

\section*{4.5\quad Canonical cardinal filtration of a $\mathfrak{U}$-paratopos category}

\textbf{Definition 4.5.1.} By a \textit{$\mathfrak{U}$-paratopos category} we mean a category $\mathcal{M}$ which is a $\mathfrak{U}$-pseudo-topos and accessible, i.e. which satisfies the equivalent conditions of theorem 4.3.6 (p. 70). We say that it is a \textit{$\pi$-$\mathfrak{U}$-paratopos category}, if it is $\pi$-accessible, i.e. if it admits a small generating subcategory whose objects are $\pi$-accessible, or what amounts to the same thing, if % [sic] soient π-accessible -> soient π-accessibles, accord du pluriel
$\mathcal{M}$ is equivalent to a category of the form $\Ind'_{\pi}(\mathcal{M}_{\pi})$, where $\mathcal{M}_{\pi}$ is a small category (or: a category equivalent to a small category), and stable under $\varinjlim$ of cardinality $\leqslant \pi$. (And we can take for $\mathcal{M}_{\pi}$ the full subcategory of $\mathcal{M}$ consisting of the objects
\msp{page 107}
which are $\pi$-accessible).

In what follows, we are going, more generally, to consider a category $C$ equivalent to a small category, stable under $\varinjlim$ of cardinality $\leqslant \pi_0$ (\footnote{There is no need to assume this right away.}), and the category
\[
\mathcal{M} \ \overset{\mathrm{def}}{=}\ \Ind'_{\pi_0}C \ \subset\ C^{\wedge}\ ,
\tag*{(4.5.1)}
\]
\DERmark{ENV1135}

where $\pi_0$ is a given infinite cardinal. I propose to introduce on $\mathcal{M}$ a canonical cardinal filtration:

\textbf{Proposition 4.5.2} (\footnote{Is proved (and is true) only if $\pi^{\pi_0} = \pi$, see correction 4.5.6, page 116.}). \textit{Let $F \in \Ob\mathcal{M} \subset \Ob C^{\wedge}$, and let $\pi$ be a cardinal $\geqslant \pi_1 =$} $\sup(\pi_0, \card\Fl C_0)$\textit{, where $C_0$ is a small full subcategory of $C$ such that $C_0 \overset{\sim}{\hookrightarrow} C$. The following conditions are equivalent:}
\begin{enumerate}[leftmargin=3em]
\item[\textit{a)}] $\card\Fl(C_0/F) \leqslant \pi$.
\item[\textit{a$'$)}] $\card\Ob\,(C_0/F) \leqslant \pi$.
\item[\textit{a$''$)}] \textit{For every $x$ in $C_0$,} $\card F(x) \leqslant \pi$.
\item[\textit{b)}] \textit{There exists a small filtering category $I$ (not necessarily ordered),} large relative to $\pi_0$\textit{, with} $\card I \leqslant \pi$\textit{, and a direct system $(x_i)_{i\in I}$ in $C$, such that we have}
\[
F \ \simeq\ \varinjlim_{i\in I} x_i\ ,
\tag*{(4.5.2.1)}
\]
\textit{where the $\varinjlim$ is taken in $\mathcal{M}$, or what amounts to the same thing (by virtue of 4.4.15, $I$ being large relative to $\pi_0$) in $C^{\wedge}$, i.e. $F$ is isomorphic to the ind-representable functor defined by the ind-object $\underline{x} = (x_i)_{i\in I}$ of $C$.}
\item[\textit{c)}] \textit{$F$ is $\pi$-accessible in $\mathcal{M}$.}
\item[\textit{c$'$)}] \textit{$F$ is $\pi$-accessible in $C^{\wedge}$.}
\end{enumerate}
\msp{page 108}

\textsc{Proof}. We have the $\pm$ tautological implications
\[
\begin{tikzcd}[row sep=small]
 & \mathrm{a}' \arrow[r, Rightarrow] & \mathrm{a}'' & \\
\mathrm{a} \arrow[ur, Rightarrow] \arrow[dr, Rightarrow] & & & \\
 & \mathrm{b} \arrow[r, Rightarrow] & \mathrm{c}' \arrow[r, Rightarrow] & \mathrm{c}\ ,
\end{tikzcd}
\tag*{$(*)$}
\]
\printedpage{30}
% head: A. Grothendieck : Dérivateurs, Ch. XVIII, Édition des Universités de Montpellier II et Paris VII
% foot: Version 17 décembre 2025
where a$'$ $\Longrightarrow$ a$''$ comes from the canonical inclusion $F(x) = \Hom_{\mathcal{M}}(x,F) \hookrightarrow \Ob C/F$, a $\Longrightarrow$ b is obtained by taking $I = C_0/F$ and noting that $I$ is stable under direct limits of cardinality $\leqslant \pi_0$ ($C$ being so), hence is large relative to $\pi_0$ (\footnote{We do not need this stability, in any case $C_0/F$ is large relative to $\pi_0$, by virtue of 4.4.3, p. 80.}). We have b $\Longrightarrow$ c$'$ by virtue of the fact that the strictly full subcategory $C^{\wedge}_{\pi}$ of $C^{\wedge}$ defined by
\[
\left\{
\begin{array}{lcl}
C^{\wedge}_{\pi} & = & \text{strictly full subcategory of } C^{\wedge} \text{ consisting of the}\\
 & & F \in C^{\wedge} \text{ which are } \pi\text{-accessible in }C^{\wedge}.
\end{array}
\right.
\tag*{(4.5.2.2)}
\]
\DERmark{ENV1136}

is stable under $\varinjlim$ of cardinality $\leqslant \pi$ (SGA 4 I 9.9). Finally c$'$ $\Longrightarrow$ c comes from the fact that the inclusion functor $\mathcal{M} \hookrightarrow C^{\wedge}$ is fully faithful and $\pi_0$-accessible, i.e. commutes with direct limits large relative to $\pi_0$, and a fortiori with those large relative to $\pi$. By the implications of the diagram above, it therefore suffices to prove that we have
% marginal mark (handwritten squiggle) in the left margin of the edition beside the following display
\[
\mathrm{c} \overset{(\mbox{\footnotemark[33]})}{\Longrightarrow} a \qquad\quad \text{and} \qquad\quad \mathrm{a}'' \ \Longrightarrow\ \mathrm{a}\ .
\]
\footnotetext[33]{is proved (and is true) only if $\pi^{\pi_0} = \pi$.}
\setcounter{footnote}{33}

\msp{page 109}

\textit{Let us prove} a$''$ $\Longrightarrow$ a. Let us consider the map
\[
\Fl C/F \ \longrightarrow\ \Fl C
\tag*{$(*)$}
\]
induced by the étale functor $C' = C/F \longrightarrow C$, then the set of arrows above an % [sic] induit -> induite, accord avec « l'application »
arrow $\alpha : x \longrightarrow y$ of $C$ is in one-to-one correspondence with $\Ob C'_y = F(y)$, since such an arrow $\alpha'$ is known when one knows its target, which is an arbitrary object of $F(y)$ ($C' \longrightarrow C$ being étale, i.e. fibrant with discrete fibers). Thus the fiber of $(*)$ at $\alpha$ is of % [sic] le fibre -> la fibre, genre
cardinality $= \card F(y) \leqslant \pi$, so we have
\[
\card\Fl C/F \ \leqslant\ \pi \underbrace{\card\Fl C}_{\leqslant \pi_1} \ \leqslant\ \pi \underbrace{\ \pi_1\ }_{\leqslant \pi} \ \leqslant\ \pi^2 = \pi\ ,
\]
which proves a.

Let us finally prove c $\Longrightarrow$ a. For this, we are going to use the first inclusion $\mathcal{M}_{\pi} \subset \Filt^{\pi}\mathcal{M}$ of SGA 4 I 9.16, which applies as soon as we prove the

% marginal mark (handwritten squiggle) in the left margin of the edition beside Corollaire 4.5.3
\textbf{Corollary 4.5.3} (\footnote{Not proved, condition c) of cardinal filtrations perhaps not satisfied. Probably \textit{false} in general.}). \textit{Under the conditions of 4.5.2, let, for every cardinal $\pi \geqslant \pi_1$,} $\Filt^{\pi}(\mathcal{M})$ \textit{be the strictly full subcategory of $\mathcal{M}$ consisting of the $F \in \Ob\mathcal{M}$ which satisfy condition a) of 4.5.2. Then the family of} $(\Filt^{\pi})_{\pi\geqslant\pi_1}$ \textit{is a}
\msp{page 110}
\textit{cardinal filtration of $\mathcal{M}$.}

Condition a) of loc. cit. 9.12 (the $\uFilt^{\pi}$ are equivalent to small categories) is obvious. Let us prove b), i.e. the stability of $\uFilt^{\pi}$ under filtering $\varinjlim$ $F = \varinjlim F_i$ \textit{large relative to} $\pi_0$, of cardinality $\leqslant \pi$. Since $\mathcal{M} \hookrightarrow C^{\wedge}$ commutes with these limits, for $x$ in $C$
\[
F(x) \ =\ \varinjlim_{i} F_i(x)\ ,
\]
\printedpage{31}
% head: A. Grothendieck : Dérivateurs, Ch. XVIII, Édition des Universités de Montpellier II et Paris VII
% foot: Version 17 décembre 2025
whence $\card F(x) \leqslant \card\Ob I \cdot \pi \leqslant \pi \cdot \pi = \pi$, which (by the equivalence a $\Longleftrightarrow$ a$''$) implies $F \in \Ob\Filt^{\pi}$. Let us prove condition c) of loc. cit. For this, we will need the fact, more precise than b), which we have just established:

\textbf{Corollary 4.5.4} (\footnote{Proved only if $\pi^{\pi_0} = \pi$\,! Probably \textit{false} in general.}). \textit{For every $\pi \geqslant \pi_1$,} $\Filt^{\pi}$ \textit{is stable under $\varinjlim$ (in $\mathcal{M}$) of cardinality $\leqslant \pi$ (not necessarily filtering, a fortiori not necessarily large relative to $\pi_0$).}

It suffices, by virtue of what is already proved of lemma 4.4.14 (\footnote{We need a variant of lemma 4.4.14, for the $\varinjlim$ of cardinality $\leqslant \pi'$, which is proved in the same way.}), to show that $\Filt^{\pi}$ % [sic] de du lemme -> du lemme, mot répété
\DERmark{ENV1137}

is stable under $\varinjlim$ \textit{of cardinality} $\leqslant \pi_0$. If $\varphi : J \longrightarrow \Filt^{\pi}$ is a functor with $\card J \leqslant \pi_0$, we put it in indexed form
\[
\psi\ :\ I \times J \ \longrightarrow\ C\ ,
\]
thanks to 4.4.13 b, with $I$ large relative to $\pi_0$.
\DERmark{ENV1138}


\msp{page 111}

I recall that we could take
\[
I \ =\ \uHom(J,C)/\varphi\ ,
\]
where $\uHom(J,C) \hookrightarrow \uHom(J, \mathcal{M} = \Ind'_{\pi_0}C)$, at the cost of not requiring $I$ to be ordered. We would like to prove that $\card\Fl I \leqslant \pi$. The trouble is that we only arrive at the upper bound
\DERmark{ENV1139}

\[
\card\Fl I \ \leqslant\ \pi^{\pi_0}\ .
\]
(Even if $J$ is discrete, we do no better, since then
\[
I \ \simeq\ \prod_{j\in J} C/\varphi(j)\ , \qquad \text{whence} \quad \Fl I = \prod_{j\in J} \Fl C/\varphi(j)\ ,
\]
and if $\card\Fl C/\varphi(j) = \pi$ for every $j$, we find $\card\Fl I = \pi^{\pi_0}$ -- there is no way to do better\,!)

So I cannot manage to prove 4.5.4 directly, which would have come in handy for proving condition c) of loc. cit. But to tell the truth, we have done here the work of the proof of theorem 4.3.6, part b $\Longrightarrow$ c -- we have already proved there that the filtration we have just constructed (by condition a) of 4.5.2) is a cardinal filtration. Condition c) of SGA 4 I 9.12 is proved on page 76. But I am going to redo
\msp{page 112}
the proof more seriously.

So let $F \in \Ob\mathcal{M}$, consider
\[
J \ =\ C/F\ ,
\]
it is a category stable under $\varinjlim$ of cardinality $\leqslant \pi_0$, hence large relative to $\pi_0$, and we have
\[
F \ \simeq\ \varinjlim_{j\in J} x_j \qquad\quad (\,\varinjlim \ \text{in } \mathcal{M})
\]
\printedpage{32}
% head: A. Grothendieck : Dérivateurs, Ch. XVIII, Édition des Universités de Montpellier II et Paris VII
% foot: Version 17 décembre 2025
denoting by $j \longmapsto x_j$ the canonical functor $C/F \longrightarrow C$. Then let $I$ be the set of subcategories $J'$ of $J$ which are \textit{of cardinality $\leqslant \pi$ and large relative to $\pi_0$, $I$} ordered by inclusion, and for every $J' \in I$, let us put
\[
X_{J'} \ =\ \varinjlim_{j\in J'} x_j\ .
\]
As $J'$ is large relative to $\pi_0$, the direct limit in $\mathcal{M}$ is also a limit in $C^{\wedge}$, whence it follows at once, since $\card\Ob J' \leqslant \card\Fl J' \leqslant \pi$, that $X_{J'}$ is in $\Filt^{\pi}$. Then
\[
J' \ \longmapsto\ X_{J'} \ \in\ \Ob\Filt^{\pi}
\]
is a functor
\[
I \ \longrightarrow\ \Filt^{\pi}\ ,
\]
whose direct limit is that of the $x_j$, hence isomorphic to $F$. On the other hand, it is immediate that $I$ is large relative to $\pi$, i.e. that if $(J'_{\alpha})_{\alpha\in A}$ is
\msp{page 113}
a family of elements of $I$, it is bounded above by a $J' \in I$. This follows from the

% marginal mark (handwritten squiggle) in the left margin of the edition beside Lemme 4.5.5
\textbf{Lemma 4.5.5} (\footnote{Is proved \textit{only} if $\pi^{\pi_0} = \pi$\,! Probably false in general.}). \textit{Let $\pi$, $\pi_0$, with $\pi > \pi_0$, be infinite cardinals, and let $J$ be a category large relative to $\pi_0$. Then for every subset $\Phi$ of} $\Fl J$\textit{, such that} $\card\Phi \leqslant \pi$\textit{, there exists a subcategory $J'$ of $J$} large relative to $\pi_0$\textit{, such that} $\card J' \leqslant \pi$ (\footnote{This is OK if we require only $\card J' \leqslant \pi^{\pi_0}$.})\textit{, and} $\Fl J' \supset \Phi$.

This lemma being admitted, we have therefore represented an arbitrary object $F$ of $\mathcal{M}$ as a $\varinjlim\limits_{i\in I} X_i$, with $I$ an ordered set large relative to $\pi$, the $X_i$ in $\Filt^{\pi}$. It remains to prove that if $\pi' > \pi$ and $F \in \Filt^{\pi'}$, then we can take $I$ such that $\card I \leqslant \pi'^{\pi}$. But $F \in \Filt^{\pi'}$ means that $J = C/F$ satisfies $\card\Fl J \leqslant \pi'$. Then the set of subsets of cardinality $\leqslant \pi$ of $\Fl J$ has a cardinality bounded above by $\pi'^{\pi}$, a fortiori $\card I \leqslant \pi'^{\pi}$, which finishes proving that $(\Filt^{\pi})_{\pi\geqslant\pi_1}$ is a cardinal filtration on $\mathcal{M}$.

\msp{page 114}

To prove the lemma, we begin by saturating $\Phi$ under finite composition, we find a subcategory $J'_0$ of $J$ such that $\card\Fl J'_0 \leqslant \pi$ \ldots{} But I notice that to prove the lemma, I again need $\pi^{\pi_0} = \pi$ -- an envelope of $J'$ large relative to $\pi_0$ is going to have a cardinality bounded above (at best) by $\pi^{\pi_0}$.
\DERmark{ENV1140}


Let us take e.g.
\[
J \ =\ \mathfrak{P}(S) \qquad\quad \text{(ordered set)}
\]
with $\card S = \pi > \pi_0$, so $J$ is large relative to $\pi_0$. Let us consider the map
\[
S \ \longrightarrow\ J = \mathfrak{P}(S)\ , \qquad\quad s \longmapsto \{s\}\ ,
\]
and let us take for $\Phi \subset \Fl J$ the set of $\mathrm{id}_{\{s\}}$. Thus we are looking for a subset $J'$ of $J$, containing $S$, and large relative to $\pi_0$ for the order induced by $J$. We see at once that the smallest subset $J'$ that does the job is the set $\mathfrak{P}_{\pi_0}(S)$ of subsets of $S$ of cardinality $\leqslant \pi_0$. But we find
\[
\card J' \cdot \pi_0! \ =\ \pi^{\pi_0}\ ,
\]
\printedpage{33}
% head: A. Grothendieck : Dérivateurs, Ch. XVIII, Édition des Universités de Montpellier II et Paris VII
% foot: Version 17 décembre 2025
\msp{page 115}
where
\[
\begin{aligned}
\pi_0! \ &=\ \card\Aut\,(E_0) \qquad\quad (E_0 \text{ a set of cardinality } \pi_0)\\
&=\ \pi_0^{\pi_0}\\
&=\ 2^{\pi_0}\ ,
\end{aligned}
\]
so
\[
\card J' \cdot 2^{\pi_0} \ =\ \pi^{\pi_0}\ .
\tag*{$(*)$}
\]
On the other hand, it is immediate that $\card J' \geqslant \card\mathfrak{P}(E_0) = 2^{\pi_0}$ (where $E_0$ is a subset of $S$ of cardinality $\pi_0$), so the left-hand side of $(*)$ is equal to $\card J'$, so we find
\[
\card J' \ =\ \pi^{\pi_0}\ ,
\]
sorry\,!

Finally, I have convinced myself that proposition 4.5.2, and the preceding corollaries and lemmas by which I tried to establish it, is true only for those cardinals $\pi > \pi_0$ such that
\[
\pi^{\pi_0} \ =\ \pi\ ,
\]
which, assuming the generalized continuum hypothesis, is equivalent to
\[
\pi^{\pi_0} \ <\ 2^{\pi}\ (\ =\ \pi^{\pi}\ )
\]
(or again to $\pi^{\pi_0} \neq 2^{\pi}$ $(= \pi^{\pi})$). I am going to give
\msp{page 116}
here a corrected statement.

\textbf{Proposition 4.5.6.}
\begin{enumerate}[leftmargin=3em]
\item[\circled{1\textsuperscript{o}}] \textit{Under the general hypotheses of 4.5.2 on $C$, $\pi_0$, $\pi_1$, $\pi$, the conditions a, a$'$, a$''$, b, c$'$ of this proposition are equivalent (provided $\pi \geqslant \pi_1$)} (\footnote{\textbf{N.B.} For this equivalence, the hypothesis that $C$ be stable under $\varinjlim$ of cardinality $\leqslant \pi$ is unnecessary, and also for c$'$) to imply c), and to be equivalent to it if $\pi = \pi^{\pi_0} \geqslant \pi_1$, cf. 4.7.1, p. 161.})\textit{, and imply condition c ($F$ is $\pi$-accessible). When $\pi^{\pi_0} = \pi$ (in particular when $\pi = 2^{c}$, with $c \geqslant \pi_0$), then all the conditions a to c$'$ of 4.5.2 are equivalent.}
\item[\circled{2\textsuperscript{o}}] \textit{More generally and more precisely, for $\pi \geqslant \pi_1$ given, let us consider the conditions on $\pi$:}
\begin{enumerate}[leftmargin=2em]
\item[\textit{1)}] $\pi^{\pi_0} = \pi$.
\item[\textit{2)}] $\Filt^{\pi}(\mathcal{M})$ \textit{(defined in 4.5.3) is stable under $\varinjlim$ of cardinality $\leqslant \pi$.}
\item[\textit{3)}] \textit{$\forall$ $F$ in $\mathcal{M}$, we have $F \simeq \varinjlim\limits_{I} X_i$, with $I$ an ordered set large relative to $\pi$, the $X_i$ in} $\Filt^{\pi}(\mathcal{M})$\textit{. Moreover, if $\pi' > \pi$ and $F \in$} $\Filt^{\pi'}(\mathcal{M})$\textit{, we can take $I$ of cardinality $\leqslant \pi'^{\pi}$.}
\item[\textit{4)}] \textit{As in 3), without the bound on cardinalities, and assuming only $I$ a category large relative to $\pi$ (not necessarily an ordered category).} % [sic] catégorie ordonné -> catégorie ordonnée, accord
\item[\textit{5)}] \textit{Condition c) of 4.5.2 implies the other conditions (they are equivalent, by 1°), in other words}
\[
\mathcal{M}_{\pi} \ \subset\ \Filt^{\pi} \qquad\quad \textit{(so } \mathcal{M}_{\pi} = \Filt^{\pi}\textit{, since we know that } \Filt^{\pi} \subset \mathcal{M}_{\pi}\textit{)}\ ,
\]
\end{enumerate}
\end{enumerate}


\endgroup
% END UNIT der-ch18/der-ch18_p031-035.tex
% BEGIN UNIT der-ch18/der-ch18_p036-040.tex
\begingroup\makeatletter%
\newcommand{\printedpage}[1]{\marginpar{\footnotesize p.~#1}}%
\newcommand{\msp}[1]{\marginpar{\footnotesize [#1]}}%
\newcommand{\Filt}{\operatorname{Filt}}%
\newcommand{\Ob}{\operatorname{Ob}}%
\newcommand{\card}{\operatorname{card}}%
\newcommand{\Fl}{\operatorname{Fl}}%
\newcommand{\Hom}{\operatorname{Hom}}%
\newcommand{\uHom}{\underline{\operatorname{Hom}}}%
\newcommand{\Ens}{\mathrm{Ens}}%
\newcommand{\Ind}{\operatorname{Ind}}%
\newcommand{\circled}[1]{\tikz[baseline=(c.base)]\node[draw,circle,inner sep=1pt](c){#1};}%
\makeatother\let\printedpage\derpp\let\msp\dermsp%
\let\mspin\DERmspin%


\setcounter{footnote}{39}

\printedpage{34}
% head: A. Grothendieck : Dérivateurs, Ch. XVIII, Édition des Universités de Montpellier II et Paris VII
% foot: Version 17 décembre 2025

\msp{page 117}

\begin{quote}
\textit{where $\mathcal{M}_\pi$ denotes the strictly full subcategory of $\mathcal{M}$ formed of the $\pi$-accessible objects.}
\end{quote}

\begin{quote}
\textit{This being so, we have the implications}
\[
\mathrm{1)} \ \Longrightarrow\ \mathrm{2)} \ \Longrightarrow\ \mathrm{3)} \ \Longrightarrow\ \mathrm{4)} \ \Longrightarrow\ \mathrm{5)} \ ,
\]
\textit{and in the particular case where $\mathcal{M} = $} (Ens), \textit{$C$ = full subcategory formed of the sets of cardinal $\leqslant \pi_0$, we also have} 5) $\Longrightarrow$ 1), \textit{i.e. all these conditions are equivalent.}
\end{quote}

\textsc{Proof.}

1) $\Longrightarrow$ 2). Cf. proof of 4.5.4.

2) $\Longrightarrow$ 3). The category $J = \Filt^\pi/F$ is stable under $\varinjlim$ of cardinal $\leqslant \pi$, $\Filt^\pi$ being so, hence $J$ is large relative to $\pi$, moreover $F = \varinjlim\limits_{j\in J} X_j$, where $j \longmapsto X_j$ is the source functor $\Filt^\pi/F \longrightarrow \Filt^\pi$. Whence the desired representation of $F$, replacing $J$ if need be by a cofinal ordered set $I$, e.g. by taking for $I$ the set of subcategories $J$ of $I$ such that $\card \Fl J \leqslant \pi$ \texttt{[this second part of the sentence ``e.g. ...'' calls for correction]}. The bound $\card I \leqslant \pi'^{\pi}$, for $F \in \Ob \Filt^{\pi'}$, is obtained by AQT \texttt{[âne qui trotte: at a donkey's trot, i.e. by routine verification]}.

The implication 3) $\Longrightarrow$ 4) is tautological.

4) $\Longrightarrow$ 5) (\footnote{compare lemma 4.5.14, page 133.}), for if $F \in \Ob \mathcal{M}_\pi$, we find by 4) that $F$ is a direct factor of an object of $\Filt^\pi$, and $\Filt^\pi$ being stable
\DERmark{ENV1141}

\msp{page 118}
under direct factors (argument by argument), whence $F \in \Ob \Filt^\pi$.
\DERmark{ENV1142}


Finally, if $\mathcal{M} = \Ens$, $C = \Ens(\pi_0)$. We note that
\[
\Filt^\pi \ = \ \{ \ X \in \Ob \Ens \mid \card\,(X)^{\pi_0} \ \leqslant \ \pi \ \} \ .
\]
On the other hand, we easily find
\[
\begin{aligned}
(\Ens)_\pi \quad &\overset{\mathrm{def}}{=} \quad \text{full subcategory of (Ens) formed} \\
&\phantom{\overset{\mathrm{def}}{=}} \quad \text{of the $\pi$-accessible objects} \\
&= \quad \{ X \in \Ob \Ens \mid \card X \leqslant \pi \} \ .
\end{aligned}
\]
\DERmark{ENV1143}\DERmark{ENV1144}

Thus we have $(\Ens)_\pi \subset \Filt^\pi$ if and only if $\pi^{\pi_0} \leqslant \pi$, i.e. $\pi^{\pi_0} = \pi$, q.e.d. (\footnote{Cf. proof of part 1°) of 4.5.6, below (bottom of page and following page).}).

\textbf{Corollary 4.5.7.} \textit{For every cardinal $\pi \geqslant \pi_1$, we have}
\[
\Filt^\pi \ \subset \ \mathcal{M}_\pi \ \subset \ \Filt^{\pi^{\pi_0}} \ . \tag*{(4.5.7.1)}
\]

The inclusion $\Filt^\pi \subset \mathcal{M}_\pi$ is the implication a $\Longrightarrow$ c of part 1°) of 4.5.6. To prove $\mathcal{M}_\pi \subset \Filt^{\pi^{\pi_0}}$, we put $\pi' = \pi^{\pi_0}$ and note that $\pi'^{\pi_0} = \pi^{\pi_0 \cdot \pi_0} = \pi^{\pi_0} = \pi'$. Thus we can
\printedpage{35}
% head: A. Grothendieck : Dérivateurs, Ch. XVIII, Édition des Universités de Montpellier II et Paris VII
% foot: Version 17 décembre 2025
apply 1° to $\pi'$ instead of $\pi$, and we find that $\mathcal{M}_{\pi'} = \Filt^{\pi'}$. Since $\mathcal{M}_\pi \subset \mathcal{M}_{\pi'}$, we conclude $\mathcal{M}_\pi \subset \Filt^{\pi'}$, q.e.d.

I forgot to prove 1°) in 4.5.6. In view of the diagram of implications $(\ast)$, p. 108, and the implication $\mathrm{a}'' \Longrightarrow \mathrm{a}$ already proved (p. 109),
\msp{page 119}
it remains only to prove
\[
\mathrm{c}' \ \Longrightarrow \ \mathrm{a}'' \ ,
\]
i.e. that $F$ $\pi$-accessible \textit{in} $C^\wedge$ implies $\card F(x) \leqslant \pi \ \forall \ x \in C$. This will follow from the somewhat more general result:

\textbf{Lemma 4.5.8.}
\begin{itemize}[leftmargin=3em]
\item[\textit{1°)}] \textit{In the category $\mathcal{E} = $} (Ens), \textit{the $\pi$-accessible objects ($\pi$ an infinite cardinal) are the sets $X$ such that} $\card X \leqslant \pi$.
\item[\textit{2°)}] (\footnote{Cf. generalization 4.14.18, p. 391 b) and c).}). \textit{Let $C$ be a category equivalent to a small category, such that $\forall \ x, y \in \Ob C$, we have} $\card \Hom(x,y) \leqslant \pi$. \textit{Let $\mathcal{E}$ be a $\mathfrak{U}$-category stable under $\varinjlim$ large relative to $\pi$, and weakly $\pi$-accessible (i.e. admitting a generating subcategory formed of $\pi$-accessible objects), and stable under products of cardinal $\leqslant \pi$. Then for an object $F$ of $\uHom(C,\mathcal{E})$ to be $\pi$-accessible, it is} necessary \textit{that $\forall \ x \in \Ob C$, $F(x)$ be a $\pi$-accessible object of $\mathcal{E}$, and this condition is also sufficient if} $\card \operatorname{\text{red}} \Fl(C) \leqslant \pi$ (\footnote{even \textit{without} any condition of stability of $\mathcal{E}$ under products of cardinal $\leqslant \pi$.}).
\end{itemize}

\textsc{Proof.} 1°) $\card X \leqslant \pi \Longrightarrow X$ $\pi$-accessible in $\mathcal{E} = (\Ens)$, for $X = \coprod_{x \in X} \{x\}$, and we apply SGA 4 I 9.9. Conversely, suppose $X$ $\pi$-accessible, and write $X \simeq \varinjlim\limits_{i \in I} X_i$, where $I$ is the set of subsets of $X$ which are of cardinal $\leqslant \pi$. Since $I$ is large relative to $\pi$, the identity map $X \xrightarrow{\ \mathrm{id}\ } X$ factors through a set $X_i$, hence $X_i = X$, hence $\card X = \card X_i \leqslant \pi$, q.e.d.

\msp{page 120}

2°) Consider, for $x \in \Ob C$, the functor
\[
e \xrightarrow{\ \varepsilon_x\ } C \ ,
\]
whence a functor
\[
\varepsilon_x^* \ : \ \uHom(C,\mathcal{E}) \longrightarrow \mathcal{E} \ , \qquad F \longmapsto F(x) \ ,
\]
admitting a right adjoint
\[
\varepsilon_{x\,*} \ : \ \mathcal{E} \ \longrightarrow \ \uHom(C,\mathcal{E}) \ ,
\]
given by
\[
\varepsilon_{x\,*}(E)(y) \ = \ E^{\Hom(y,x)}
\]
\DERmark{ENV1145}

(subject to the existence of products of cardinal $\leqslant \pi$). This functor $\varepsilon_{x\,*}$ is $\pi$-accessible by virtue of loc. cit. 9.8.

This being so, let $F$ be a $\pi$-accessible object of $\uHom(C,\mathcal{E})$, we see that $\varepsilon_x^*(F)$ is $\pi$-accessible, $\varepsilon_x^*$ having a $\pi$-accessible right adjoint:
\[
E \ \longmapsto \ \Hom_{\mathcal{E}}(\underbrace{F(x)}_{\varepsilon_x^*(F)}, E) \ \simeq \ \Hom_{\uHom(C,\mathcal{E})}(F, \varepsilon_{x\,*}(E))
\]
\printedpage{36}
% head: A. Grothendieck : Dérivateurs, Ch. XVIII, Édition des Universités de Montpellier II et Paris VII
% foot: Version 17 décembre 2025
is the composite of the $\pi$-accessible functor $\varepsilon_{x\,*}$, and the $\pi$-accessible covariant functor represented by $F$ in $\uHom(C,\mathcal{E})$.

Conversely, if the $F(x)$ are $\pi$-accessible, let us prove that $F$ is so, provided $\card \Fl(C) \leqslant \pi$. Done by AQT \texttt{[âne qui trotte: at a donkey's trot, i.e. by routine verification]}, I am not going to spell it out here.

\msp{page 121}

\textbf{Proposition 4.5.9.} \textit{Let $C$ be a category equivalent to a small category, and let $\mathcal{M} = $} $\Ind'_{\pi_0}(C)$, \textit{$\pi_0$ a given infinite cardinal. Assume $C$ stable under $\varinjlim$ of cardinal $\leqslant \pi_0$. For every cardinal $\pi \geqslant \pi_0$, let $\mathcal{M}_\pi \supset C$ be the full subcategory of $\mathcal{M}$ formed of the $\pi$-accessible objects of $\mathcal{M}$. Let, on the other hand,} $\pi_1 = \sup(\pi_0, \card \Fl C)$. \textit{Then $(\mathcal{M}_\pi)_{\pi \geqslant \pi_1}$ is a cardinal filtration of $\mathcal{M}$} (\footnote{\textbf{N.B.} All the conditions for cardinal filtrations are satisfied with the exception, when $\pi \geqslant \pi_0$, of that of the bound on cardinals $\card I \leqslant \pi'^{\pi}$ condition c).})\textit{, and moreover, for $\pi \geqslant \pi_0$, $\mathcal{M}_\pi$ is stable in $\mathcal{M}$ under $\varinjlim$ of cardinal $\leqslant \pi$.}

\textsc{Proof.} The stability of $\mathcal{M}_\pi$ is a reminder, cf. SGA 4 I 9.9. Since we have
\[
\mathcal{M}_\pi \ \subset \ \Filt^{\pi^{\pi_0}}
\]
for $\pi \geqslant \pi_1$ (cf. 4.5.7), we therefore see that the $\mathcal{M}_\pi$ are equivalent to small categories (condition a) for cardinal filtrations). Condition b) for cardinal filtrations is contained in the stability property of the $\mathcal{M}_\pi$ which we have just recalled. Finally let us prove condition c). So let $\pi \geqslant \pi_0$, and let $x$ be in $\mathcal{M}$, consider the category
\[
I \ = \ \mathcal{M}_\pi/x \ ,
\]
it is stable under $\varinjlim$ of cardinal $\leqslant \pi$, by the similar stability of $\mathcal{M}_\pi$, hence $I$ is large relative to $\pi$ and equivalent to a small category, on the other hand, we have
\msp{page 122}
\[
x \ = \ \varinjlim_{i \in I} x_i \ ,
\]
the $x_i$ in $\mathcal{M}_\pi$ ( $\varinjlim$ in $\mathcal{M}$ of the functors $I = \mathcal{M}_\pi/x \longrightarrow \mathcal{M}_\pi$). Replacing $I$ if need be by a cofinal ordered category, this proves the qualitative part of c. Suppose finally $x \in \Ob \mathcal{M}_{\pi'}$, where $\pi' > \pi$, let us show that we then have
\[
\card I_0, \ \card I' \ \leqslant \ \pi'^{\pi} \qquad \text{(\footnotemark, \footnotemark)} \ , \tag*{(4.5.9.1)}
\]
\addtocounter{footnote}{-1}%
\footnotetext{where $I_0$, $I'_0$ are skeletal subcategories of $I$, $I'$ such that $I_0 \hookrightarrow I$, $I'_0 \hookrightarrow I'$ are equivalences \texttt{[Grothendieck has crossed out the subscript 0 of $I'$ in 4.5.9.1]}.}%
\stepcounter{footnote}%
\footnotetext{We only prove it if $\pi \geqslant \pi_1$.}%
which will finish proving that we have a cardinal filtration. We need a bound on the cardinals.

\textbf{Lemma 4.5.10} (\footnote{Give a common statement with 4.5.11, and first prove 4.5.11.}). \textit{Let $\pi$ be a cardinal $\geqslant \pi_1$. Then we have} % [sic] énonce -> énoncé, missing accent (in footnote 47)
\[
\card \Fl((\Filt^\pi)_0) \ \leqslant \ 2^\pi \ . \tag*{(4.5.10.1)}
\]

Indeed we have, denoting by $C^\wedge(\pi)$ the strictly full subcategory of $C^\wedge$ formed of the $F : C^\circ \longrightarrow \Ens$ with $\card F(x) \leqslant \pi$ for every $x$ in $C$ (i.e. $F$ factors through $\Ens(\pi)$), and by $C^\wedge(\pi)_0$ an
\printedpage{37}
% head: A. Grothendieck : Dérivateurs, Ch. XVIII, Édition des Universités de Montpellier II et Paris VII
% foot: Version 17 décembre 2025
\msp{page 123}
exhaustive skeletal subcategory,
\[
\card C^\wedge(\pi)_0 \ \leqslant \ 2^\pi \ , \tag*{(4.5.10.2)}
\]
whence (4.5.10.1), since for $\pi \geqslant \pi_1$, $\Filt^\pi$ is a full subcategory of $C^\wedge(\pi)$. To establish (4.5.10.2), we note that the set of isomorphism classes of $\Ens(\pi)$ is identified with the set of cardinals $\leqslant \pi$, which set is of cardinal $\leqslant \pi$. Replacing $\Ens(\pi)$ by $\Ens(\pi)_0$, $C$ by $C_0$, we must bound $\card \Fl(\uHom(C_0, \Ens(\pi)_0))$. Now $\Ob \uHom(C_0, \Ens(\pi)_0)$ is contained in the set of maps from $\Fl C_0$ to $\Fl(\Ens(\pi)_0)$, the first is bounded by $\pi_1$, the second by $\pi^\pi \times \pi \times \pi = \pi^{2\pi} = \pi^\pi = 2^\pi$ (\textbf{N.B.} $\pi < 2^\pi$, whence $\pi^\pi \leqslant (2^\pi)^\pi = 2^{\pi^2} = 2^\pi$), hence $\Ob \uHom$ bounded by $(2^\pi)^{\pi_1} = 2^{\pi\pi_1} = 2^\pi$. On the other hand, if $F$, $G$ are in $\uHom(C_0, \Ens(\pi))$, $\Hom(F,G)$ is bounded by $(\pi^\pi)^{\pi_1} = \pi^\pi = 2^\pi$, hence $\Fl \uHom$ is bounded by $2^\pi \times 2^\pi \times 2^\pi$, q.e.d.
\DERmark{ENV1146}\DERmark{ENV1147}\DERmark{ENV1148}\DERmark{ENV1149}\DERmark{ENV1150}\DERmark{ENV1151}\DERmark{ENV1152}\DERmark{ENV1153}\DERmark{ENV1154}


\textbf{Corollary 4.5.11.} \textit{Let $F \in \Ob \Filt^\pi$, $F' \in \Ob \Filt^{\pi'}$, with $\pi' \geqslant \pi \geqslant \pi_1$. Then}
\[
\card \Hom(F, F') \ \leqslant \ \pi'^{\pi} \ . \tag*{(4.5.11.1)}
\]

\msp{page 124}

It is essentially the same proof, slightly simpler. (We should have started there.)

\textbf{Corollary 4.5.12.} \textit{If $\pi' \geqslant \pi \geqslant \pi_1$, we have for $x \in \Ob \mathcal{M}_{\pi'}$}
\[
\card\,(\mathcal{M}_\pi/x)_0 \ \leqslant \ \pi'^{(\pi^{\pi_0})} \ . \tag*{(4.5.12)}
\]

Indeed,
\[
\card \Ob\,(\mathcal{M}_\pi/x)_0 \ \overset{\text{by 4.5.11}}{\leqslant} \ \underbrace{\underbrace{\card \Ob\,(\mathcal{M}_\pi)_0}_{\leqslant 2^{(\pi^{\pi_0})} \text{by 4.5.10}} \times (\pi'^{\pi_0})^{\pi^{\pi_0}}}_{= \pi'^{(\pi^{\pi_0})}} \ ,
\]
i.e.
\[
\card \Ob\,(\mathcal{M}_\pi/x)_0 \ \leqslant \ \pi'^{(\pi^{\pi_0})} \ ,
\]
whence
\[
\card \Fl(\mathcal{M}_\pi/x)_0 \ \leqslant \ \pi'^{(\pi^{\pi_0})} \times \pi'^{(\pi^{\pi_0})} \times (\pi'^{\pi_0})^{(\pi^{\pi_0})} \ \leqslant \ \pi'^{(\pi^{\pi_0})} \ .
\]
This therefore gives (returning to 4.5.9.1)
\[
\card I_0 \ \leqslant \ \pi'^{(\pi^{\pi_0})} \ .
\]

On the other hand, $I'$ can be defined in terms of $I_0$ as the set of subcategories of $I_0$ which are of cardinal $\leqslant \pi$, hence their number is bounded by $(\card I_0)^\pi \ \leqslant \ (\pi'^{\pi^{\pi_0}})^\pi = \pi'^{\pi\pi^{\pi_0}} = \pi'^{(\pi^{\pi_0})}$.

But this is not quite the required bound, since
\[
\pi'^{\pi} \leqslant \pi'^{(\pi^{\pi_0})} \qquad \text{(\footnotemark)},
\]
\footnotetext{\textbf{N.B.} We have equality if we have $\pi' = 2^{c'}$ with $c' \geqslant \pi^{\pi_0}$, or $\pi = 2^c$ with $c \geqslant \pi_0$.}%
and since it is possible, a priori, that the inequality
\printedpage{38}
% head: A. Grothendieck : Dérivateurs, Ch. XVIII, Édition des Universités de Montpellier II et Paris VII
% foot: Version 17 décembre 2025
\msp{page 125}
is strict! Apparently, the ``bounds''
\[
\mathcal{M}_\pi \ \subset \ \Filt^{\pi^{\pi_0}} \ , \qquad \mathcal{M}_{\pi'} \ \subset \ \Filt^{\pi'^{\pi_0}}
\]
are too crude. So it is a matter of getting a better grasp of the categories $\mathcal{M}_\pi$, for $\pi \geqslant \pi_0$ (\footnote{But we can use 4.5.10, 4.5.11 to conclude that for $x \in \Ob \mathcal{M}_{\pi'} \subset \Ob \Filt^{\pi'^{\pi_0}}$, we have $\card \Fl(C_0/x) \leqslant (\pi'^{\pi_0})^{\pi_0} = \pi'^{\pi_0}$, we will therefore have $x = \varinjlim\limits_{J} x_j$ ($x_j \in \Ob C$, $\varinjlim$ filtering, $\card \Fl J \leqslant \pi'^{\pi_0}$), and introducing the set $I$ of subcategories $J' \subset J$ which are of cardinal $\leqslant \pi$, and putting $x_{J'} = \varinjlim\limits_{j \in J'} x_j$, we find $x = \varinjlim\limits_{J' \in I} x_{J'}$, where this time $I$ is an \textit{ordered set} large relative to $\pi$, the $x_{J'} \in \mathcal{M}_\pi$, and $\card I \leqslant (\card J)^\pi \leqslant (\pi'^{\pi_0})^\pi = \pi'^{\pi\pi_0} = \pi'^{\pi}$, this is the right bound!}). I am going to venture a

\textbf{Proposition 4.5.13.} \textit{The hypotheses and notations are those of 4.5.9. Let $\pi$ be a cardinal $\geqslant \pi_0$, and let $x \in \Ob \mathcal{M}$. Consider the conditions:}
\begin{itemize}[leftmargin=3em]
\item[\textit{a})] \textit{There exists a category $I$, with} $\card \Fl I \leqslant \pi$, \textit{and a functor $i \longmapsto x_i : I \longrightarrow C$, such that $x$ is isomorphic to $\varinjlim x_i$.}
\item[\textit{a}$'$)] \textit{As in a), with $I$ filtering.}
\item[\textit{a}$''$)] \textit{As in a), with $I$ ordered and filtering.}
\item[\textit{b})] \textit{There exists a} cofinal \textit{functor $I \longrightarrow C/x$, with} $\card \Fl I \leqslant \pi$.
\item[\textit{b}$'$)] \textit{The filtering category $C/x$ is ``essentially of cardinal $\leqslant \pi$'', i.e. there exists an ordered category} \texttt{[no, cf. b$''$)]} filtering $I$, $\card \Fl I \leqslant \pi$, \textit{and a cofinal functor $I \longrightarrow C/x$.}
\item[\textit{b}$''$)] \textit{As in b$'$), with $I$ moreover ordered.}
\item[\textit{c})] \textit{There exists a strict epimorphic family of arrows $(u_\alpha)_{\alpha \in A}$, $x_\alpha \longrightarrow x$, with $x_\alpha\in\Ob C$ for every $\alpha\in A$, indexed by a set $A$ such that} $\card A \leqslant \pi$.
\item[\textit{d})] \textit{We have $x \in \Ob \mathcal{M}_\pi$, i.e. $x$ is $\pi$-accessible in $\mathcal{M}$.}
\end{itemize}
\DERmark{ENV1155}


\msp{page 126}

\textit{This being so:}
\begin{itemize}[leftmargin=3em]
\item[\circled{1\textsuperscript{\textit{o}}}] \textit{The first six conditions $a$, $a'$, $a''$, $b$, $b'$, $b''$ are equivalent, say $P(\pi)$, and imply the last two,}
\[
\begin{tikzcd}[column sep=small]
 & P(\pi) \arrow[dl, Rightarrow] \arrow[dr, Rightarrow] & \\
c & & d \ .
\end{tikzcd}
\]
\textit{If $\pi \geqslant \pi_1$, then $P(\pi)$ is equivalent to $c$.}
\item[\circled{2\textsuperscript{\textit{o}}}] \textit{Let $\Filt'^{\pi} \subset \mathcal{M}$ be the strictly full subcategory of $\mathcal{M}$ formed of the objects which satisfy condition $P(\pi)$ above, and $\Filt^\pi$ the strictly full subcategory defined in 4.5.3 (p. 109). We then have}
\[
\Filt^\pi \ \subset \ \Filt'^{\pi} \ \subset \ \mathcal{M}_\pi \ , \tag*{(4.5.13.1)}
\]
\textit{and for there to be equality of the three terms, it suffices that we have $\pi^{\pi_0} = \pi$. (Cf. prop. 4.5.6, page 116). Moreover, $\Filt'^{\pi}$ is stable in $\mathcal{M}$ under} sums \textit{of cardinal $\leqslant \pi$.}
\item[\circled{3\textsuperscript{\textit{o}}}] \textit{The following conditions on $C$, $\pi_0$, $\pi$ are equivalent.}
\end{itemize}


\endgroup
% END UNIT der-ch18/der-ch18_p036-040.tex
% BEGIN UNIT der-ch18/der-ch18_p041-045.tex
\begingroup\makeatletter%
\newcommand{\printedpage}[1]{\marginpar{\footnotesize p.~#1}}%
\newcommand{\msp}[1]{\marginpar{\footnotesize [#1]}}%
\newcommand{\Filt}{\operatorname{Filt}}%
\newcommand{\card}{\operatorname{card}}%
\newcommand{\Fl}{\operatorname{Fl}}%
\newcommand{\Ob}{\operatorname{Ob}}%
\newcommand{\Ind}{\operatorname{Ind}}%
\makeatother\let\printedpage\derpp\let\msp\dermsp%
\let\mspin\DERmspin%


\setcounter{footnote}{49}

\printedpage{39}
% head: A. Grothendieck : Dérivateurs, Ch. XVIII, Édition des Universités de Montpellier II et Paris VII

\begin{itemize}
\item[\textit{1)}] $\Filt'^{\pi}$ \textit{is stable in $\mathcal{M}$ under} $\varinjlim$ \textit{of cardinality $\leqslant \pi$.}
\item[\textit{1$'$)}] $\Filt'^{\pi}$ \textit{is stable under cokernels of double arrows.}
\item[\textit{2)}] \textit{Every object $x$ of $\mathcal{M}$ is of the form $\varinjlim\limits_{I} x_i$, where $x_i$ in} $\Filt'^{\pi}$\textit{, $I$ large relative to $\pi$.}
\item[\textit{3)}] \textit{Condition d) of 1°) implies $P(\pi)$ (hence is equivalent to it), i.e. we have}
\[ \mathcal{M}_{\pi} \ = \ \Filt'^{\pi}. \]
\end{itemize}

\msp{page 127}
\textit{These conditions are satisfied in the following three cases:}
\begin{itemize}
\item[\textit{A)}] \textit{$\pi^{\pi_0} = \pi$ (as a reminder, this is contained in 2°).}
\item[\textit{B)}] \textit{The strict epimorphic families in $\mathcal{M}$ are strict universal epimorphic families.}
\item[\textit{C)}] \textit{Filtering direct limits in $\mathcal{M}$ commute with base change (which is the case in particular if the filtering} $\varinjlim$ \textit{are left exact)} (\footnote{I suppose that this condition amounts to the same condition on $C$, for the filtering $\varinjlim$ of cardinality $\leqslant \pi$.}).
\end{itemize}

\textbf{Remark.} The filtration of $\mathcal{M}$ by the $\Filt'^{\pi}$ is none other than the one considered in SGA 4 I 9.13, where one makes in addition assumption B) on $\mathcal{M}$, which afterwards was no longer so much to my taste. At any rate, according to (4.5.13.1), this filtration comes closer to the ``good'' filtration by the $\mathcal{M}_{\pi}$, than the one I had still been inclined not so long ago to prefer to it, the one by the $\Filt^{\pi}$. It has moreover over the latter the advantage, in the case B) considered, of being a cardinal filtration (which is not necessarily the case for $\Filt^{\pi}$, at least if there exist infinite cardinals $\pi \geqslant \pi_0$ with $\pi^{\pi_0} \neq \pi$). I now suspect that in general (and if there exist infinite cardinals $\pi > \pi_0$ such that $\pi^{\pi_0} \neq \pi$), outside conditions A) B) C), the filtration by the $(\Filt'^{\pi})$ is \textit{not} a
\msp{page 128}
cardinal filtration. In any case, 4.5.6 and the present proposition, 3°, show that if either $(\Filt^{\pi})_{\pi\geqslant\pi_2}$, or $(\Filt'^{\pi})_{\pi\geqslant\pi_2}$ is a cardinal filtration ($\pi_2$ being a given cardinal $\geqslant \pi_1$), then it coincides with the filtration by the $\mathcal{M}_{\pi}$. So apparently, in all the cases where I know how to construct a cardinal filtration, on a paratopos at least (i.e. on a category $\mathcal{M}$ stable under small $\varinjlim$ , and all of whose objects are accessible -- the existence of a cardinal filtration implying moreover that of a small strictly generating subcategory \dots), this is none other than the one by the $\mathcal{M}_{\pi}$ -- something I had not noticed when writing SGA 4 I 9. % [sic] m'étais par aperçu -> m'étais pas aperçu, "par" for the negation "pas"

\textsc{Proof of 4.5.13.}

\textit{Let us prove 1°.} We have the tautological implications
\[
\begin{tikzcd}[row sep=large]
\mathrm{b}'' \arrow[r, Rightarrow] \arrow[d, Rightarrow] & \mathrm{b}' \arrow[r, Rightarrow] \arrow[d, Rightarrow] & \mathrm{b} \arrow[d, Rightarrow] \\
\mathrm{a}'' \arrow[r, Rightarrow] & \mathrm{a}' \arrow[r, Rightarrow] & \mathrm{a} \text{ ,}
\end{tikzcd}
\qquad\qquad \text{(\footnotemark) ,}
\]
\footnotetext{taking into account that $C$ is strictly generating in $\mathcal{M}$, hence $x \simeq \varinjlim\limits_{i\in C/x} x_i$.}
and to prove the equivalence of these six conditions, it suffices to prove
\[ \mathrm{a} \ \Longrightarrow \ \mathrm{b}'' \ . \]
% foot: Version 17 décembre 2025

\printedpage{40}
% head: A. Grothendieck : Dérivateurs, Ch. XVIII, Édition des Universités de Montpellier II et Paris VII

\msp{page 129}
If we have
\[ x \ \simeq \ \varinjlim_{I} x_i \ , \qquad\quad \card \Fl I \ \leqslant \ \pi \ , \]
let $I'$ be the set of finite \textit{subdiagrams} $J \subset I$ of $I$, ordered by inclusion, so $I'$ is a filtering ordered set, and from the fact that $C$ is stable under finite $\varinjlim$, we find a functor
\[ J \longmapsto x_J \ = \ \varinjlim_{i\in J} x_i \ : \ I' \longrightarrow C \ , \]
and we will have
\[ x \ \simeq \ \varinjlim_{J\in I'} x_J \ . \]
On the other hand, we have
\[
\begin{aligned}
\card I' \ &= \ \card \mathfrak{P}_f( \ \underbrace{\Fl(I)}_{\text{of card $\leqslant \pi$}} \ ) \\
&\leqslant \ \pi + \pi^2 + \ \cdots \ + \pi^n + \ \dots \\
&\leqslant \ \pi\aleph_0 \ = \ \pi \ ,
\end{aligned}
\]
q.e.d. To establish 1°, it remains to establish
\[
\begin{tikzcd}[column sep=small]
& P(\pi) \arrow[dl, Rightarrow] \arrow[dr, Rightarrow] & \\
\mathrm{c} & & \mathrm{d} \text{ ,}
\end{tikzcd}
\]
and the reverse implication $\mathrm{c} \Longrightarrow P(\pi)$ when $\pi \geqslant \pi_1$. But $P(\pi) \Longrightarrow \mathrm{c}$ is a tautology, and $P(\pi) \Longrightarrow \mathrm{d}$ since $\mathcal{M}_{\pi}$ is stable under $\varinjlim$ of cardinality $\leqslant \pi$. It remains to prove that $\mathrm{c} \Longrightarrow P(\pi)$ if $\pi \geqslant \pi_1$.

\msp{page 130}
Suppose then that we have a strict epimorphic family
\[ x_{\alpha} \longrightarrow x \qquad\quad (\alpha \in A,\ \card A \leqslant \pi) \ , \]
this therefore means that $x$ is the direct limit of a direct system indexed by the category $B$ having as set of objects $A \amalg (A\times A)$, the nonidentity arrows being represented by the canonical arrows
\[
\begin{tikzcd}[column sep=small]
& x_{\alpha} \times_x x_{\beta} \arrow[dl, "p_{\alpha,\beta}"'] \arrow[dr, "q_{\alpha,\beta}"] & \\
x_{\alpha} & & x_{\beta} \text{ .}
\end{tikzcd}
\]

(\textbf{N.B.} $B$ is not an ordered category, for for the object $(\alpha,\alpha)$, represented by $x_{\alpha}\times_x x_{\alpha}$, % [sic] n'est par une -> n'est pas une, "par" for the negation "pas"
there are \textit{two} arrows into the object $\alpha$, represented by $x_{\alpha}$. But we obviously have $\card \Fl(B) \leqslant \pi$.) If we knew that $C$ is stable under binary products and under passage to subobjects (here the $x_{\alpha} \times_x x_{\beta} \ \subset \ x_{\alpha} \times x_{\beta}$), we would have proved condition a). Failing that, let us again use the fact that $C$ is generating by epimorphisms (even strict ones) and choose for every pair $(\alpha,\beta)$ an epimorphic set of arrows:
\[ \underbrace{x_{\alpha,\beta,\gamma}}_{\in \Ob C_0} \xrightarrow{u_{\alpha,\beta,\gamma}} x_{\alpha} \times_x x_{\beta} \ , \]
% foot: Version 17 décembre 2025

\printedpage{41}
% head: A. Grothendieck : Dérivateurs, Ch. XVIII, Édition des Universités de Montpellier II et Paris VII
where $C_0$ is a reduced exhaustive subcategory of $C$. The arrow $u_{\alpha,\beta,\gamma}$ is known when we know the composites
\msp{page 131}
\[
\begin{tikzcd}
& x_{\alpha} \\
x_{\alpha,\beta,\gamma} \arrow[ur, "v_{\alpha,\beta,\gamma}"] \arrow[dr, "w_{\alpha,\beta,\gamma}"'] & \\
& x_{\beta} \text{ ,}
\end{tikzcd}
\]
so for $x_{\alpha,\beta,\gamma} \in \Ob C_0$ fixed, the possible choices for $u_{\alpha,\beta,\gamma}$ have cardinality bounded by $\pi_1\pi_1 = \pi_1$, moreover we have $\card \Ob C_0 \leqslant \pi_1$, so the total number of arrows $v_{\alpha,\beta,\gamma}$, $w_{\alpha,\beta,\gamma}$ (for $\alpha$, $\beta$ fixed) is bounded by $\pi_1 \cdot \pi_1 = \pi_1$. So the total number of arrows $u_{\alpha,\beta,\gamma}$, for $\alpha$, $\beta$, $\gamma$ variable, is bounded by $\pi \cdot \pi_1 = \pi$ (if $\pi \geqslant \pi_1$). So $x$ is represented by a category $I$ having as objects the disjoint union of $A$ and the set of triple indices $\alpha,\beta,\gamma$, and for each such index, exactly two nonidentity arrows $v_{\alpha,\beta,\gamma}$, $w_{\alpha,\beta,\gamma}$ with this object as source, and going respectively into $x_{\alpha}$, $x_{\beta}$. We thus have $\card \Fl I \leqslant \pi$, and we have a functor $I \longrightarrow C$, $\alpha \longmapsto x_{\alpha}$, $(\alpha,\beta,\gamma) \longmapsto x_{\alpha,\beta,\gamma}$, such that $x \ = \ \varinjlim\limits_{i\in I} x_i$. This completes the proof of 1°).

\msp{page 132}
\textit{Let us prove 2°.} The first inclusion (4.5.13.1) is tautological, the second merely expresses $P(\pi) \Longrightarrow \mathrm{d}$. The fact that $\pi^{\pi_0} = \pi$ implies equality of the extreme terms is a reminder of 4.5.6. It remains to prove that $\Filt'^{\pi}$ is stable in $\mathcal{M}$ under sums of cardinality $\leqslant \pi$. But if $(x_{\alpha})_{\alpha\in A}$ is a family of objects of $\Filt'^{\pi}$, with $\card A \leqslant \pi$, and if $x = \coprod_{\alpha\in A} x_{\alpha}$, taking for every $x_{\alpha}$ a functor $\varphi_{\alpha} : I_{\alpha} \longrightarrow C$, with $\card \Fl I_{\alpha} \leqslant \pi$, such that $x_{\alpha}$ is the $\varinjlim$ in $\mathcal{M}$ of this functor, we find that $x$ is the $\varinjlim$ of the functor $\varphi : I \longrightarrow C$ defined by the $\varphi_{\alpha}$, where $I = \coprod_{\alpha} I_{\alpha}$ (sum in Cat). It is obvious that $\card \Fl(I) \leqslant \pi\cdot\pi = \pi$, which completes the proof of 2°.

\textit{Let us prove 3°.} We have the tautological implications $1 \Longrightarrow 1'$, and $1' \Longrightarrow 1$, since $\Filt'^{\pi}$ is stable under sums of cardinality $\leqslant \pi$. Thus $1 \Longleftrightarrow 1'$, and implies that for every $x$ in $\mathcal{M}$,
\msp{page 133}
$I \ = \ \Filt'^{\pi}/x$ is stable under $\varinjlim$ of cardinality $\leqslant \pi$, hence is large relative to $\pi$. Since $x \simeq \varinjlim\limits_{i\in I} x_i$, this proves that $1 \Longrightarrow 2$. The implication $2 \Longrightarrow 3$ is standard, and follows from $\Filt'^{\pi} \subset \mathcal{M}_{\pi}$, I am going to spell out in a

\textbf{Lemma 4.5.14.} \textit{Let $\mathcal{M}$ be a category stable under small} $\varinjlim$ \textit{, $\pi$ an infinite cardinal, $\mathcal{M}_0$ a strictly full subcategory of $\mathcal{M}$, stable under direct factors. Consider the conditions}
\DERmark{ENV1156}

\begin{itemize}
\item[\textit{a)}] \textit{$\mathcal{M}_0$ is a strictly generating subcategory, and stable in $\mathcal{M}$ under} $\varinjlim$ \textit{of cardinality $\leqslant \pi$.}
\item[\textit{b)}] \textit{For every object of $\mathcal{M}$, we have an isomorphism}
\[ x \ \simeq \ \varinjlim_{i\in I} x_i \ , \]
\textit{the $x_i$ in $\mathcal{M}_0$, $I$ large relative to $\pi$.}
% foot: Version 17 décembre 2025

\printedpage{42}
% head: A. Grothendieck : Dérivateurs, Ch. XVIII, Édition des Universités de Montpellier II et Paris VII
\item[\textit{b$'$)}] \textit{As in b), with $I$ a category that is} ordered.
\item[\textit{c)}] \textit{$\mathcal{M}_{\pi} \subset \mathcal{M}_0$ (where $\mathcal{M}_{\pi}$ is the strictly full subcategory of $\mathcal{M}$ consisting of the $\pi$-accessible objects).}
\end{itemize}
\textit{We thus have the implications}
\[ a \ \Longrightarrow \ b \ \Longleftrightarrow \ b' \ \Longrightarrow \ c \ , \]
\textit{and if we assume $\mathcal{M}_0 \subset \mathcal{M}_{\pi}$, and $\mathcal{M}_{\pi}$ a strictly generating subcategory of $\mathcal{M}$, then the four conditions $a$, $b$, $b'$, $c$ are equivalent, and equivalent to $\mathcal{M}_0 = \mathcal{M}_{\pi}$} (\footnote{cf. proof of proposition 4.5.6, $4 \Longrightarrow 5$, page 117.}).

\msp{page 134}
Finally (returning to the proof of 4.5.13, 3°), the implication $3) \Longrightarrow 1)$ comes from the fact that $\mathcal{M}_{\pi}$ is stable under $\varinjlim$ of cardinality $\leqslant \pi$.

It remains to prove that these conditions, i.e. the stability of $\Filt'^{\pi}$ under cokernels of double arrows, is satisfied in cases A) B) C), case A) being moreover already known. Case B) follows immediately from the transitivity of strict universal epimorphic families, and this is moreover also contained explicitly (stability of $\Filt'^{\pi}$ under $\varinjlim$ of cardinality $\leqslant \pi$) in SGA 4 I 9.13 (if not in the statement, at least it is pointed out in the proof). When writing SGA 4 I 9, I therefore knew that we have $\mathcal{M}_{\pi} \subset \Filt'^{\pi}$ for every $\pi$, but the inclusion, more obvious however in the case at hand, $\Filt'^{\pi} \subset \mathcal{M}_{\pi}$ (when we assume the objects of $C$ $\pi_1$-accessible) had apparently escaped me\,?

It remains to deal with case C). $\Big($I have the impression that it is a little more general than case B), i.e. that condition B) on $\mathcal{M}$ implies C), but not conversely. Thus, condition C) is verified as soon as $\mathcal{M}$ corresponds to a species of algebraic structure defined in
\msp{page 135}
terms of $\varprojlim$ \textit{which are finite}, i.e. when $\mathcal{M}$ is of the form $\Ind(C)$, $C$ stable under $\varinjlim$ \textit{which are finite}, whereas I doubt that condition B) is always satisfied for these categories. For example in the category Cat (which is a paratopos of this species), the filtering $\varinjlim$ are left exact, but I do not believe that the strict epimorphic families are universal. % [sic] épimorphiques stricts -> épimorphiques strictes, agreement with "familles"
Let us take e.g. two group homomorphisms $G_0 \xrightarrow{u_0} G$, $G_1 \xrightarrow{u_1} G$ such that $G_0 * G_1 \longrightarrow G$ is surjective, then this two-element family is strictly epimorphic in Gr and in Cat (\footnote{!, this is \textit{false} in general, and also false in the particular case below for nontrivial $H$.}), but is not necessarily universally strictly epimorphic, even if e.g. $G_0 = G_1 = H$, $G = H \times H$, and $u_0, u_1$ are injections $H \rightrightarrows H \times H$, given by $u_0(h) = (h,1)$, and $u_1(h) = (1,h)$. If I take the diagonal map $H \longrightarrow H \times H$, the inverse images of $u_0$, $u_1$ are the same object of $\mathcal{M}/H$, namely $e \longrightarrow H$, and this family is strictly epimorphic only if $H$ is the trivial group.$\Big)$
\DERmark{ENV1157}


Case C) will follow from the more precise result:

\textbf{Corollary 4.5.15.} \textit{Under the general conditions of 4.5.13, suppose that in $\mathcal{M}$ the} $\varinjlim$
\msp{page 136}
that are filtering \textit{commute with base change (e.g. the filtering} $\varinjlim$ \textit{are exact). Then for every $\pi \geqslant \pi_1$,} $\Filt'^{\pi}$ \textit{consists of the $x \in \Ob \mathcal{M}$ satisfying the following condition.}
% foot: Version 17 décembre 2025

\printedpage{43}
% head: A. Grothendieck : Dérivateurs, Ch. XVIII, Édition des Universités de Montpellier II et Paris VII
\begin{itemize}
\item[\textit{c$'$)}] \textit{There exists a family of arrows $(u_{\alpha})_{\alpha\in A}$, $u_{\alpha} : x_{\alpha} \longrightarrow x$, in $\mathcal{M}$, which is} universally strictly epimorphic \textit{(i.e. covering for the canonical topology of $\mathcal{M}$), with $x_{\alpha} \in \Ob C$ $\forall\ \alpha \in A$, and} $\card A \leqslant \pi$.
\end{itemize}
\textsc{Proof} : we must prove that this condition is equivalent to condition c) of 4.5.13 1°). Since we trivially have $\mathrm{c}' \Longrightarrow \mathrm{c}$, it remains to show that $\mathrm{c} \Longrightarrow \mathrm{c}'$. But we saw in 4.5.13 1°) that c implies a$''$, i.e. that we have an isomorphism
\[ x \ = \ \varinjlim_{i\in I} x_i \ , \]
the $x_i$ in $C$, and $I$ a \textit{filtering ordered} set, with $\card I \leqslant \pi$. It remains to see that if $x$ in $\mathcal{M}$ is a \textit{filtering} direct limit of objects $x_i$, then the set of arrows $x_i \longrightarrow x$ is strictly epimorphic \textit{universally.} Now this is immediate, for for $x' \longrightarrow x$, we have $x' = \varinjlim\limits_{i} x'_i$ (by assumption C), where $x'_i = x' \times_x x_i$, so the family $x'_i \longrightarrow x'$ is strictly epimorphic, q.e.d.

Let us now prove that in case C, the condition of 3°, 1 is satisfied, i.e. that $\Filt'^{\pi}$ is stable under $\varinjlim$ of cardinality $\leqslant \pi$. This
\msp{page 137}
follows at once from the

\textbf{Corollary 4.5.16.} (of 4.5.15). \textit{Under the conditions of 4.5.15 (i.e. condition C) of 4.5.13, 3°), let $x$ be an object of $\mathcal{M}$, and suppose that there exists a family} strictly epimorphic \underline{universally} $(x_{\alpha} \longrightarrow x)_{\alpha\in A}$ \textit{of arrows of $\mathcal{M}$, with $x_{\alpha} \in$} $\Ob \Filt'^{\pi}$ \textit{for every $\alpha \in A$, and} $\card A \leqslant \pi$. \textit{Then $x \in$} $\Ob \Filt'^{\pi}$.
\DERmark{ENV1158}


Indeed, the assumption on $x_{\alpha}$ is equivalent, by virtue of 4.5.15, to the existence of a strict universal epimorphic family of arrows $(x_{\alpha\beta} \longrightarrow x_{\alpha})_{\beta\in B_{\alpha}}$, with $\card B_{\alpha} \leqslant \pi$ and the $x_{\alpha\beta}$ in $C$. Let $B = \coprod_{\alpha\in A} B_{\alpha}$, we then have $\card B \leqslant \pi^2 = \pi$, and on the other hand, the family of composite arrows $x_{\alpha\beta} \longrightarrow x_{\alpha} \longrightarrow x$ indexed by $B$ is universally strictly epimorphic (by transitivity of this type of family), hence $x \in \Ob \Filt'^{\pi}$, q.e.d.

With all that, we have still not given a characterization of $\mathcal{M}_{\pi}$ itself, which is better than $\Filt^{\pi}$ and $\Filt'^{\pi}$ from the categorical point of view, for only the $\mathcal{M}_{\pi}$ define a cardinal filtration, and moreover they have the desirable stability properties in this
\msp{page 138}
context, for $\varinjlim$ of cardinality $\leqslant \pi$. The only nontautological criterion I manage to extract in the general case (for every $\pi \geqslant \pi_0$, and without any additional exactness assumption on $\mathcal{M}$) is the following, much less ``explicit'' or ``computational'' than the very concrete criteria for $\Filt^{\pi}$ and for $\Filt'^{\pi}$:

\textbf{Proposition 4.5.17.} \textit{Let $\pi$ be a cardinal $\geqslant \pi_0$. Then $\mathcal{M}_{\pi}$ is equal to the smallest strictly full subcategory of $\mathcal{M}$, containing $C$, and stable under} $\varinjlim$ \textit{of cardinality $\leqslant \pi$.} (It is important to note that it \textit{is not enough} to say: stable under $\varinjlim\limits_{I}$ , with \textit{$I$ large relative to} $\pi_0$, and $\card I \leqslant \pi$. We then find $\Filt^{\pi}$, which risks being a strictly smaller category than $\mathcal{M}_{\pi}$. It is also the category of the $\varinjlim\limits_{I} x_i$, with $x_i \in C$ $\forall\ i \in \Ob I$, and $\card \Fl I \leqslant \pi$ -- since that category is $\Filt'^{\pi}$.)
\DERmark{ENV1159}


\textsc{Proof}. Let $\mathcal{M}'$ be the category of the statement, we thus have $\mathcal{M}' \subset \mathcal{M}_{\pi}$, since $\mathcal{M}_{\pi}$ contains $C$ and is stable under the $\varinjlim$ of cardinality $\leqslant \pi$. Conversely, lemma 4.5.14, $\mathrm{a} \Longrightarrow \mathrm{c}$.
% foot: Version 17 décembre 2025


\endgroup
% END UNIT der-ch18/der-ch18_p041-045.tex
% BEGIN UNIT der-ch18/der-ch18_p046-050.tex
\begingroup\makeatletter%
\newcommand{\printedpage}[1]{\marginpar{\footnotesize p.~#1}}%
\newcommand{\msp}[1]{\marginpar{\footnotesize [#1]}}%
\newcommand{\cardred}{\mathop{\mathrm{card}\,\textup{red}}\nolimits}%
\newcommand{\Ind}{\mathrm{Ind}}%
\newcommand{\circled}[1]{\tikz[baseline=(c.base)]\node[draw,circle,inner sep=1pt](c){#1};}%
\makeatother\let\printedpage\derpp\let\msp\dermsp%
\let\mspin\DERmspin%


\printedpage{44}
% head: A. Grothendieck : Dérivateurs, Ch. XVIII, Édition des Universités de Montpellier II et Paris VII
% foot: Version 17 décembre 2025

\msp{page 139}

\section*{4.6\quad Characterization of paratopos categories (proof of 4.3.6)}

It is time to return to the proof of Theorem 4.3.6 (p. 70-71), the one I had given being incomplete (and erroneous – I asserted that $(\mathrm{Filt}^{\pi})_{\pi\geqslant\pi_0}$ is a cardinal filtration, which is false if one assumes only $\pi^{\pi_0}=\pi$ for every $\pi\geqslant\pi_1$).

We had seen the equivalence of conditions a, a$'$, b, namely PTA. I am going to finish the proof by the logical scheme
\[
\text{PTA} \ \Longrightarrow\ \text{d)} \ \Longrightarrow\ \text{c)} \ \Longrightarrow\ \text{PTA}\ . \tag*{$(*)$}
\]

Let us prove PTA $\Longrightarrow$ d), i.e. that if $\mathcal{M}$ is an accessible pseudo-topos, then there exists an infinite cardinal $\pi_0$, and a small category $C_0$ stable under $\varinjlim$ of cardinality $\leqslant\pi_0$, and an equivalence of categories % [sic] alors existe -> alors il existe, missing subject "il"
\DERmark{ENV1160}

\[
\mathcal{M} \xrightarrow{\ \sim\ } \Ind'_{\pi_0}(C_0)\ . \tag*{(4.6.1)}
\]

More precisely, we have this:

\textbf{Proposition 4.6.1.} \textit{Let $\mathcal{M}$ be a paratopos, i.e. a $\mathfrak{U}$-category $\mathcal{M}$ stable under small $\varinjlim$ , admitting a full generating subcategory $S$ equivalent to a small category, consisting of accessible objects of $\mathcal{M}$. Let $\pi_0$ be a cardinal such that the objects of $S$ are $\pi_0$-accessible (i.e. $S\subset\mathcal{M}_{\pi_0}$), and let $C$ be the smallest strictly full subcategory of $\mathcal{M}$, containing $S$ and stable under $\varinjlim$ of cardinality $\leqslant\pi_0$. Then} % [sic] équivalent -> équivalente, agreement with the feminine "sous-catégorie"
\begin{itemize}[leftmargin=2.5em]
\item[\textit{1°)}] \textit{$C$ is equivalent to a small category.}

\msp{page 140}

\item[\textit{2°)}] \textit{$C=\mathcal{M}_{\pi_0}$, strictly full subcategory of $\mathcal{M}$ consisting of the $\pi_0$-accessible objects.}
\item[\textit{3°)}] \textit{The canonical functor}
\[
\Ind'_{\pi_0}(C) \xrightarrow[\sim]{\texttt{[}\varphi\texttt{]}} \mathcal{M} % ?? uncertain: font of the brackets around \varphi above the arrow -- typewriter (editorial) brackets or ordinary math brackets
\]
\textit{(defined up to unique isomorphism by the condition that its composite with the inclusion $C\hookrightarrow\Ind'_{\pi_0}(C)$ be isomorphic to the inclusion $C\longrightarrow\mathcal{M}$, and that $\varphi$ commute with small $\varinjlim$ , or again with filtering $\varinjlim$ large relative to $\pi_0$ – cf. 4.4.16, page 101) is an equivalence of categories} (\footnote[54]{thus we also have $\Ind'_{\pi_0}(C_0)\xrightarrow{\ \sim\ }\mathcal{M}$, where $C_0$ is a \textit{small} category equivalent to $C$.}).
\end{itemize}

\textsc{Proof}. 1°) We are going to construct $C$ by transfinite induction, setting
\[
\begin{aligned}
C_0 &\quad=\quad && \text{strictly full subcategory of $\mathcal{M}$ generated by $S$.}\\
C_{\alpha+1} &\quad=\quad && \text{strictly full subcategory of $\mathcal{M}$ generated by the $\varinjlim\limits_{I}$ in}\\
& && \text{$\mathcal{M}$ of objects of $C_\alpha$, with $I$ a category such that $\operatorname{card}\operatorname{Fl}I\leqslant\pi_0$.}\\
C_\alpha &\quad=\quad && \textstyle\bigcup_{\alpha'<\alpha}C_{\alpha'}\ \text{ if $\alpha$ is a limit ordinal.}
\end{aligned}
\]

This uniquely defines the $C_\alpha$ (for every small ordinal $\alpha$). By the definition of $C$, we see by transfinite induction that the $C_\alpha$ are contained in $C$. Let $\alpha_0$
\DERmark{ENV1161}

\printedpage{45}
% head: A. Grothendieck : Dérivateurs, Ch. XVIII, Édition des Universités de Montpellier II et Paris VII
% foot: Version 17 décembre 2025
\msp{page 141}
% ?? uncertain: a hand-drawn mark (shaped roughly like a "2" or a bracket) stands in the left margin beside lines 3--5 of this page -- marginal annotation reproduced in the image, not typeset text
be the first ordinal such that $\operatorname{card}\alpha_0>\pi_0$, so the set $I_{\alpha_0}$ of ordinals $\alpha<\alpha_0$ is large relative to $\pi_0$. I say that we have $C_{\alpha_0}=C$. For this, it suffices to see that $C_{\alpha_0}$ is stable under $\varinjlim$ of cardinality $\leqslant\pi_0$. But if $I$ is a category such that $\operatorname{card}\operatorname{Fl}I\leqslant\pi_0$, and $\varphi : I\longrightarrow C_{\alpha_0}=\bigcup_{\alpha<\alpha_0}C_\alpha$ a functor, then the definition of $\alpha_0$ implies that $\varphi$ factors as $\varphi_\alpha : I\longrightarrow C_\alpha$, with $\alpha<\alpha_0$ (\footnote[55]{A bit short. One must invoke the fact that $I_{\alpha_0}$ is large relative to $\pi_0$.}). Thus $\varinjlim\varphi\in\operatorname{Ob}C_{\alpha+1}\subset\operatorname{Ob}C_{\alpha_0}$, which shows the required stability. % [sic] de cardinal \leqslant\pi -> de cardinal \leqslant\pi_0, the stability to be proved concerns cardinal \leqslant\pi_0
\DERmark{ENV1162}\DERmark{ENV1163}


To see that $C$ is equivalent to a small category, it suffices to see this for the $C_\alpha$, proceeding by transfinite induction. This follows from the following lemma, immediate:

\textbf{Lemma 4.6.1.1.}
\begin{itemize}[leftmargin=2.5em]
\item[\textit{1°)}] \textit{The direct limit $C$ of a filtering direct system $(C_i)_{i\in I}$ of categories $C_i$ equivalent to small categories is equivalent to a small category. If $\cardred(C_i)$ $\leqslant\pi$ for every $i$, and $\operatorname{card}I\leqslant\pi$, then $\cardred(C)\leqslant\pi$ ($\pi$ being an} % [sic] système inductive -> système inductif, agreement with the masculine "système"

\msp{page 142}

\textit{infinite cardinal). Here, for any category $\mathcal{A}$, I put $\operatorname{card}\mathcal{A}=\operatorname{card}\operatorname{Fl}\mathcal{A}$, $\cardred(\mathcal{A})$ $=\operatorname{card}(\mathcal{A}_0)$, where $\mathcal{A}_0$ is a reduced category equivalent to $\mathcal{A}$.}
\item[\textit{2°)}] \textit{Let $C_0\subset\mathcal{M}$ be a full subcategory of a category $\mathcal{M}$, $\pi_0$ an infinite cardinal, and let $C'_0$ be the strictly full subcategory of $\mathcal{M}$ generated by the direct limits in $\mathcal{M}$ (representable in $\mathcal{M}$) of functors $I\longrightarrow C_0$, with $I$ a category such that $\operatorname{card}I\leqslant\pi_0$. If $C_0$ is equivalent to a small category, the same holds for $C'_0$. More precisely, if $\cardred(C_0)\leqslant\pi_1$, $\pi_1$ an infinite cardinal, then $\cardred(C'_0)\leqslant$ $\pi_1^{\pi_0}$} (\footnote[56]{This bound is not proved and probably false if one does not make somewhat stricter assumptions.}). % [sic] n'est par prouvée -> n'est pas prouvée, misprint in footnote 56
\end{itemize}

\textbf{Corollary 4.6.1.2.} \textit{$\mathcal{M}$ and $C_0$ being as in 2°) above, let $C$ be the smallest strictly full subcategory of $\mathcal{M}$ containing $C_0$, and stable under $\varinjlim$ of cardinality $\leqslant\pi_0$ (those which are representable in $\mathcal{M}$). If $C_0$ is equivalent to a small category, the same holds for $C$. More precisely, if $\cardred(C_0)\leqslant\pi_1$, where $\pi_1$ is an infinite cardinal, then}

\msp{page 143}

\textit{we have $\cardred(C)\leqslant\pi_1^{\pi_0}$} (\footnote[57]{bound not proved, probably false, cf. annotation above.}).

3°) is proved by the preceding transfinite proof, we prove by transfinite induction that
\[
\cardred C_\alpha \ \leqslant\ \pi_1^{\pi_0}
\]
for every $\alpha$. This is OK for $\alpha=0$ by hypothesis (since $\pi_1\leqslant\pi_1^{\pi_0}$). If it is true for $C_\alpha$, then 2°) ensures that
\[
\cardred C_{\alpha+1} \ \leqslant\ (\pi_1^{\pi_0})^{\pi_0} \ =\ \pi_1^{\pi_0\cdot\pi_0} \ =\ \pi_1^{\pi_0}\ , % [sic] C_\alpha -> C_{\alpha+1}, this is the successor step of the induction
\]
\DERmark{ENV1164}

finally, if $\alpha$ is a limit ordinal $<\alpha_0$, we will have $C_\alpha=\bigcup_{\alpha'<\alpha}C_{\alpha'}$, hence
\[
\cardred C_\alpha\leqslant\ \pi_0\cdot\pi_1^{\pi_0} \ =\ \pi_1^{\pi_0} % [sic] \longrightarrow -> \leqslant, an inequality is intended
\]
\DERmark{ENV1165}

(since $\pi_1^{\pi_0}\geqslant\pi_0$), by virtue of 1° above, and taking into account the fact that $\operatorname{card}I_\alpha\leqslant\pi_0$. Let $\pi'_0$ be the successor cardinal of $\pi_0$, then the same bound gives us
\[
\cardred C_{\alpha_0} \ \leqslant\ \pi'_0\pi_1^{\pi_0}\ .
\]
\printedpage{46}
% head: A. Grothendieck : Dérivateurs, Ch. XVIII, Édition des Universités de Montpellier II et Paris VII
% foot: Version 17 décembre 2025
But we have $\pi_1^{\pi_0}\geqslant2^{\pi_0}>\pi_0$, hence $\pi_1^{\pi_0}\geqslant\pi'_0$, hence $\pi_1^{\pi_0}\geqslant\pi'_0$, hence $\pi'_0\pi_1^{\pi_0}=\pi_1^{\pi_0}$, whence
\[
\cardred C_{\alpha_0} \ \leqslant\ \pi_1^{\pi_0}\ ,
\]
which is the announced bound.

\msp{page 144}

Let us prove 4.6.1 2°, i.e. $C=\mathcal{M}_{\pi_0}$. We have $C\subset\mathcal{M}_{\pi_0}$, since by hypothesis $S\subset\mathcal{M}_{\pi_0}$, and we know that $\mathcal{M}_{\pi_0}$ is stable under $\varinjlim$ of cardinality $\leqslant\pi_0$. On the other hand, we have $\mathcal{M}_{\pi_0}\subset C$, by virtue of 4.5.14 a $\Longrightarrow$ c applied to $C$, which is strictly generating in $\mathcal{M}$, $S$ being so, and $C\supset S$, and stable under $\varinjlim$ of cardinality $\leqslant\pi_0$ by construction.
\DERmark{ENV1166}


\textit{Let us finally prove 3°}. This is equivalent to the following:

\textbf{Lemma 4.6.1.3.} \textit{Let $\mathcal{M}$ be a $\mathfrak{U}$-category stable under} small \textit{$\varinjlim$ , $\pi_0$ a cardinal, $\mathcal{M}_0$ a subcategory of $\mathcal{M}$} strictly generating, \textit{stable under $\varinjlim$ of cardinality $\leqslant\pi_0$, and consisting of $\pi_0$-accessible objects. (By virtue of 4.6.1 2° already proved, applied to $S=\mathcal{M}_0$, if $\mathcal{M}_0$ is strictly full, we therefore have $\mathcal{M}_0=\mathcal{M}_{\pi_0}$.) Then the canonical functor}
\[
\Ind'_{\pi_0}(\mathcal{M}_0) \ \longrightarrow\ \mathcal{M}
\]
\textit{is an equivalence of categories.}

The fact that this functor is fully faithful follows at once from the hypothesis that the objects of $\mathcal{M}_0$ are $\pi_0$-accessible: % [sic] facteur -> foncteur, misprint
\DERmark{ENV1167}

\[
\begin{aligned}
\operatorname{Hom}_{\mathcal{M}}(\varinjlim_{i\in I}\underbrace{X_i}_{\in\mathcal{M}_0},\ \varinjlim_{j\in J}\underbrace{Y_j}_{\in\mathcal{M}_0}\ ) &\ \simeq\ \varprojlim_{i\in I}\ \underbrace{\operatorname{Hom}_{\mathcal{M}}(X_i,\ \varinjlim_{j\in J}Y_j)}_{\varinjlim\limits_{j\in J}\operatorname{Hom}(X_i,Y_j)}\\
&\ \simeq\ \underline{\operatorname{Hom}}_{\Ind_{\pi_0}(\mathcal{M}_0)}((X_i),(Y_j))\ .
\end{aligned}
\]
\DERmark{ENV1168}


\msp{page 145}

It remains to see that this functor is essentially surjective. Since $\mathcal{M}_0$ is strictly generating, it follows that every $x\in\mathcal{M}$ is written
\[
x \ =\ \varinjlim_{i\in I}x_i\ ,
\]
where $I=\mathcal{M}_0/x$, and the $x_i\in\mathcal{M}_0$. The hypothesis that $\mathcal{M}_0$ is stable under $\varinjlim$ of cardinality $\leqslant\pi_0$ implies the same property for $I=\mathcal{M}_0/x$, so $I$ is large relative to $\pi_0$. Thus $(x_i)_{i\in I}$ defines an object of $\Ind_{\pi_0}\mathcal{M}_0$, whose image in $\mathcal{M}$ is isomorphic to $x$, q.e.d. (Compare SGA 4 I 9.18.)

\textbf{Corollary 4.6.2.} \textit{Under the conditions of 4.6.1, let us take for $\pi_0$} the smallest \textit{cardinal such that the objects of $S$ are $\pi_0$-accessible. For every $\pi\geqslant\pi_0$, let $\mathcal{M}_\pi$ be the strictly full subcategory of $\mathcal{M}$ consisting of the $\pi$-accessible objects, which is also (by 4.6.1 2°) the closure of $C$ in $\mathcal{M}$ stable under $\varinjlim$ of cardinality $\leqslant\pi$.}
\begin{itemize}[leftmargin=2.5em]
\item[\textit{1°)}] \textit{Let $\pi_1=\cardred S$. We then have}
\[
\cardred\mathcal{M}_\pi \ \leqslant\ \pi_1^{\pi} \qquad (=2^{\pi}\ \textit{if}\ \pi\geqslant\pi_1\geqslant2). \tag*{(4.6.2.1)}
\]
\DERmark{ENV1169}

\printedpage{47}
% head: A. Grothendieck : Dérivateurs, Ch. XVIII, Édition des Universités de Montpellier II et Paris VII
% foot: Version 17 décembre 2025
\msp{page 146}
\item[\textit{2°)}] \textit{$(\mathcal{M}_\pi)_{\pi\geqslant\pi_1}$ is a cardinal filtration of $\mathcal{M}$, where $\pi_1=\sup(\pi_0,\operatorname{card}S)$.}
\end{itemize}

\textsc{Proof}. The bound (4.6.2.1) is contained in 4.6.1.2, where $\pi_0$ has become $\pi$. If $\pi\geqslant\pi_1$, we have
\[
\pi_1^{\pi} \ =\ 2^{\pi} \tag*{$(*)$}
\]
for we have $\pi_1\geqslant2$, whence $\pi_1^{\pi}\geqslant2^{\pi}$, on the other hand, we have $\pi_1\leqslant2^{\pi_1}$, whence $\pi_1^{\pi}\leqslant(2^{\pi_1})^{\pi}=2^{\pi_1\pi}=2^{\pi}$, since $\pi_1\pi=\pi$ for $\pi\geqslant\pi_1$; whence $(*)$.

Assertion 2° is included for the record; by virtue of $\mathcal{M}\simeq\Ind'_{\pi_0}\mathcal{M}_{\pi_0}$, we may assume that we have $\mathcal{M}=\Ind'_{\pi_0}(\mathcal{M}_{\pi_0})$, and we can then apply 4.5.9. One would moreover like to say that $(\mathcal{M}_\pi)_{\pi\geqslant\pi_0}$ is a cardinal filtration, without restricting oneself to $\pi\geqslant\pi_1$. This is true except for the bound on cardinals in condition c) of the definition of cardinal filtrations (SGA 4 I 9.12):
\[
\underbrace{x}_{\in\operatorname{Ob}\mathcal{M}_{\pi'}} \ =\ \varinjlim_{i\in I}\underbrace{x_i}_{\in\operatorname{Ob}\mathcal{M}_{\pi}}\ , \qquad\qquad \text{$I$ large relative to $\pi$ and $\operatorname{card}I\leqslant\pi'^{\pi}$}
\]
for $\pi'\geqslant\pi$. One must assume, for this bound, that $\pi\geqslant\pi_1$.

\msp{page 147}

We can now finish the proof of 4.3.6, by proving the last two implications $(*)$ (p. 139). We have d $\Longrightarrow$ c, i.e. if $\mathcal{M}\simeq\Ind'_{\pi_0}(C)$, $C$ stable under direct limits of cardinality $\leqslant\pi_0$, and $C$ equivalent to a small category, then $\mathcal{M}$ is stable under small $\varinjlim$ (4.4.13, p. 95-96), its elements are accessible, and it admits a cardinal filtration (cf. 4.5.9). But one must check all the same that the elements of $C$ are $\pi_0$-accessible. This is contained in SGA 4 I 9.18, which asserts that the full faithfulness of the functor $\Ind_{\pi_0}(C)\longrightarrow\mathcal{M}$ \textit{is equivalent} to the fact that the elements of $C$ are $\pi_0$-accessible in $\mathcal{M}$. The proof in loc. cit. was not clear, the ``it suffices'' was seen bottom of page 144. The ``it is necessary'' does not seem clear to me, with the generality given in loc. cit. – and now I see that it is even false, as one sees by taking for $C$ the unit category $e$, so $\Ind'_{\pi_0}C\simeq e$, and the assertion affirming that if $e$ is fully embedded in $\mathcal{M}$ (thus by a \textit{rigid} object of $\mathcal{M}$, i.e. an object $a$ such that $\operatorname{End}a=\{\mathrm{id}_a\}$), then $a$ is $\aleph_0$-accessible. This is false in the category of rings, where $\mathbf{R}$ is a rigid object not $\aleph_0$-accessible!

\msp{page 148}

In the case we are concerned with, $\mathcal{M}=\Ind'_{\pi_0}C$, one must clearly use the explicit calculation of the $\varinjlim$ large relative to $\pi_0$ in $\Ind'_{\pi_0}C\subset C^{\wedge}$, they are indeed calculated ``argument by argument'', cf. 4.4.15.

This completes the proof of d $\Longrightarrow$ c in $(*)$, page 139. Finally the implication c $\Longrightarrow$ PTA is tautological, the existence of a cardinal filtration implying that of a small strictly generating subcategory. This completes the proof of 4.3.6.
\DERmark{ENV1170}


I am going to give again a slightly more precise formulation of 4.3.6., specifying the cardinals, and adding to it an additional condition, inspired by SGA 4 I 9.11 (revised and corrected, cf. an annotation in the copy \texttt{[of Grothendieck]} of SGA 4 I (\footnote[58]{Ed.~note.\ \ This annotation says ``cf. simpler result above, if $C$ is \textit{strictly} generating'' and
% footnote 58 continues at the foot of printed page 48
the result in question, also in a marginal annotation in Grothendieck's copy, is:\par
\textbf{Corollary.} \textit{Let $F$ admit} \texttt{[a]} \textit{subcategory generating by strict epimorphisms, consisting of accessible objects, then} every \textit{object of $F$ is accessible (use 7.2. (i) b)).}})).

\printedpage{48}
% head: A. Grothendieck : Dérivateurs, Ch. XVIII, Édition des Universités de Montpellier II et Paris VII
% foot: Version 17 décembre 2025

\textbf{Theorem 4.6.3.} \textit{Let $\mathcal{M}$ be a $\mathfrak{U}$-category.}
\begin{itemize}[leftmargin=2.5em]
\item[\circled{1\textsuperscript{\textit{o}}}] \textit{The following conditions are equivalent for an infinite cardinal $\pi_0$:}
\begin{itemize}[leftmargin=3em]
\item[$a_{\pi_0}$\textit{)}] \textit{$\mathcal{M}$ is a $\mathfrak{U}$-$\pi_0$-paratopos category, i.e. $\mathcal{M}$ is stable under small $\varinjlim$ , and contains a small full} \texttt{[strictly]} \textit{generating subcategory $S$, consisting of $\pi_0$-accessible objects of $\mathcal{M}$.}

\msp{page 149}

\item[$b_{\pi_0}$\textit{)}] \textit{There exist a small category $S$, and a fully faithful and $\pi_0$-accessible functor $\mathcal{M}\xrightarrow{\ i\ }S^{\wedge}$, admitting a left adjoint.}
\item[$c_{\pi_0}$\textit{)}] \textit{There exists a small category $C$, stable under $\varinjlim$ of cardinality $\leqslant\pi_0$, such that $\mathcal{M}\simeq\Ind'_{\pi_0}(C)$.}
\item[$c'_{\pi_0}$\textit{)}] \textit{As c), with $C$ only equivalent to a small category.}
\end{itemize}
\item[\textit{2°)}] \textit{These conditions imply that $\mathcal{M}$ is stable under small $\varinjlim$ and small $\varprojlim$ , and the condition:}
\begin{itemize}[leftmargin=3em]
\item[$e_{\pi_0}$\textit{)}] \textit{For every category $I$ large relative to $\pi_0$, the functor}
\[
L \ =\ \varinjlim_{I}\ :\ \underline{\operatorname{Hom}}(I,\mathcal{M}) \ \longrightarrow\ \mathcal{M}
\]
\textit{is left exact (hence exact).}
\end{itemize}
\textit{This condition implies in particular the following:}
\begin{itemize}[leftmargin=3em]
\item[$\alpha(\pi_0)$] \textit{On the category of double arrows of $\mathcal{M}$, the functor} Ker \textit{(with values in $\mathcal{M}$) is $\pi_0$-accessible.}
\item[$\beta(\pi_0)$] \textit{For every arrow $y\longrightarrow x$ in $\mathcal{M}$, the base change functor $\mathcal{M}/x\longrightarrow\mathcal{M}/y$ is $\pi_0$-accessible.}
\end{itemize}
\DERmark{ENV1171}

(\footnotemark[59]). % the text of footnote 59 is not printed on pages 46--50 of the PDF

\msp{page 150}

\item[\textit{3°)}] \textit{The following conditions are equivalent.}
\begin{itemize}[leftmargin=3em]
\item[\textit{a)}] \textit{$\mathcal{M}$ is a $\mathfrak{U}$-paratopos category, i.e. is stable under small $\varinjlim$ , and contains a small full subcategory $S$ generating by strict epimorphisms, and consisting of accessible objects.} % [sic] épimorphismes strictes -> épimorphismes stricts, agreement with the masculine "épimorphismes"
\item[\textit{a$'$)}] \textit{As a), but requiring that $\mathcal{M}$ be weakly accessible, i.e. that} every \textit{object of $\mathcal{M}$ be accessible.}
\item[\textit{b)}] \textit{There exist a small category $S$ and a fully faithful and accessible functor $\mathcal{M}\hookrightarrow S^{\wedge}$ admitting a left adjoint.}
\item[\circled{c}] \textit{There exist an infinite cardinal $\pi_0$, a small category $S$ admitting direct limits of categories $I$ with $\operatorname{card}\operatorname{Fl}I\leqslant\pi_0$, and an equivalence $\mathcal{M}\simeq$ $\Ind'_{\pi_0}(S)$.}
\item[\textit{c$'$)}] \textit{As c, assuming only $S$ equivalent to a small category.}
\item[\textit{d)}] \textit{$\mathcal{M}$ is stable under small $\varinjlim$ , is weakly accessible, i.e. all its objects are accessible, and admits a cardinal filtration.}
\end{itemize}
\DERmark{ENV1172}

\end{itemize}


\endgroup
% END UNIT der-ch18/der-ch18_p046-050.tex
% BEGIN UNIT der-ch18/der-ch18_p051-055.tex
\begingroup\makeatletter%
\newcommand{\printedpage}[1]{\marginpar{\footnotesize p.~#1}}%
\newcommand{\msp}[1]{\marginpar{\footnotesize [#1]}}%
\newcommand{\circled}[1]{\tikz[baseline=(c.base)]\node[draw,circle,inner sep=1pt](c){#1};}%
\newcommand{\Mc}{\mathcal{M}}%
\newcommand{\U}{\mathfrak{U}}%
\makeatother\let\printedpage\derpp\let\msp\dermsp%
\let\mspin\DERmspin%


\printedpage{49}
% head: A. Grothendieck : Dérivateurs, Ch. XVIII, Édition des Universités de Montpellier II et Paris VII
% foot: Version 17 décembre 2025
% The mark of footnote 59 is on the preceding page of the edition; its text is printed at the foot of this page.
\footnotetext[59]{\textbf{Question.} Condition $e_{\pi_0}$ (or even, the pair $\alpha(\pi_0) + \beta(\pi_0)$ plus the existence of a small strictly generating subcategory (and the stability of $\Mc$ under small $\varinjlim$ ), does it imply that $\Mc$ is a $\pi_0$-paratopos ? I presume not.)}

\begin{itemize}[leftmargin=4em]
\item[\circled{e}] (\footnote[60]{This is the prettiest intrinsic characterization of paratopoi in the ``economy'' mode, the other, the structural statement, is \circled{c} .}). \textit{$\Mc$ is stable under small $\varinjlim$ , contains a small strictly generating full subcategory, and there exists an infinite cardinal $\pi_0$ such that condition $\mathrm{e}_{\pi_0}$ of 2°) is satisfied.}
\item[\circled{e$'$}] \textit{As e, but requiring, instead of $\mathrm{e}_{\pi_0}$, the weaker condition conjunction of $\alpha(\pi_0)$ and $\beta(\pi_0)$ in 2°.}
\end{itemize}

\msp{page 151}

\textsc{Proof}. Part 1° has already been proved. The first assertion in 2°), namely the stability of $\Mc$ under small $\varprojlim$ , is well known (follows from stability under small $\varinjlim$ , and from the existence of a small generating family, cf. SGA 4 I 8.12.7, dual version of the marginal annotation). The validity of condition $\mathrm{e}_{\pi_0}$ follows from the fact that it holds in $C^{\wedge}$ (where $C$ is the category mentioned in 1°, $\mathrm{c}_{\pi_0}$), and from the fact that the $\varprojlim$ in $\Mc$, and the $\varinjlim$ large relative to $\pi_0$ in $\Mc$, are computed in $C^{\wedge}$. We moreover find in the same way : % [sic] assertions -> assertion, singular subject "La première"
\DERmark{ENV1173}


\textbf{Corollary 4.6.4.} \textit{If $\Mc$ is a $\U$-$\pi_0$-paratopos category, then for every cardinal $\pi \geqslant \pi_0$, and every category $J$ such that} $\mathrm{card}\, J \leqslant \pi$\textit{, the functor} % [sic] tout catégorie -> toute catégorie, gender agreement
\[
\underline{\mathrm{Hom}}(J,\Mc) \xrightarrow{\ \varprojlim\limits_{J}\ } \Mc
\]
\textit{is $\pi$-accessible, or what amounts to the same thing, for every small category $I$ large relative to $\pi$, the functor}
\[
\underline{\mathrm{Hom}}(I,\mathcal M) \xrightarrow{\ \varinjlim\limits_{I}\ } \Mc % [sic] M -> \mathcal{M}, the first display has \mathcal{M}
\]
\DERmark{ENV1174}

\textit{commutes with the $\varprojlim$ of cardinality $\leqslant \pi$.}

\msp{page 152}

On the other hand, the fact that $\mathrm{e}_{\pi_0}$ implies $\alpha(\pi_0)$ and $\beta(\pi_0)$ is immediate.

Part 3°), as regards the equivalence of conditions a) to d), is contained in 1°, and the equivalence with d) has also been seen (the only non-tautological thing being that condition c), namely $\Mc \simeq \mathrm{Ind}_{\pi_0}(C)$ with $C$ small and stable under $\varinjlim$ of cardinality $\leqslant \pi_0$, implies the existence of a cardinal filtration, which was seen in 4.5.9, and reaffirmed in 4.6.2 2° above). On the other hand, that these conditions a) to d) imply \circled{e} , and that \circled{e} implies e$'$, was seen in 2°. So it remains to see that the conditions of e$'$ imply a), and for this, it remains to establish that the elements of $\Mc$ are accessible. But this is what is established in SGA 4 I 9.11, except that the statement of condition c), second case, is botched, and should be replaced by condition $\beta(\pi_0)$ (condition $\alpha(\pi_0)$ appearing in condition a) of loc. cit.).

\msp{page 153}

\textbf{Remark 4.6.5.} The developments of sections 4.4, 4.5, 4.6 considerably sharpen and strengthen, in many respects, the theory of cardinal filtrations of SGA 4 I 9. This is so, at least, in the case of $\U$-categories $\Mc$ stable under small $\varinjlim$ . In loc. cit. one works with slightly more general $\U$-categories, those satisfying
\printedpage{50}
% head: A. Grothendieck : Dérivateurs, Ch. XVIII, Édition des Universités de Montpellier II et Paris VII
% foot: Version 17 décembre 2025
the \textit{condition} $L$ : there exists an infinite cardinal $\pi_0$ such that $\Mc$ is stable under small $\varinjlim$ \textit{large relative to} $\pi_0$. I presume however that one should be able to improve in a similar way the theory of \textit{accessible} cardinal filtrations, taking inspiration from corollary 9.18 of loc. cit., establishing an equivalence of categories
\[
\Mc \;\simeq\; \mathrm{Ind}_{\pi_0}(C)\;,
\]
for a cardinal $\pi_0$ and a suitable small category $C$. Here, one can no longer assume $C$ stable under this or that type of $\varinjlim$ . Nevertheless, 4.4.15 tells us that without any further hypothesis on $C$, $\Mc \overset{\mathrm{def}}{=} \mathrm{Ind}_{\pi_0}(C)$ does satisfy $\mathrm{L}_{\pi_0}$, i.e. is stable (and even stable in $C^{\wedge}$)
\DERmark{ENV1175}


\msp{page 154}

under small $\varinjlim$ large relative to $\pi_0$.

This seen, I feel like calling \textit{accessible} a category $\Mc$ which admits a cardinal filtration, and which is moreover weakly accessible \texttt{[i.e.]} every object of which is accessible. And to venture a % [sic] filtration cardinal -> filtration cardinale, gender agreement

\textbf{Conjecture 4.6.6} (\footnote[61]{All this will be made considerably more precise in the theory which will be developed in sections 4.7, 4.16.}). \textit{Let $\Mc$ be a $\U$-category. Equivalent conditions :} % [sic] (in footnote 61) développé -> développée, agreement with "la théorie"
\begin{itemize}[leftmargin=4em]
\item[\textit{a)}] \textit{$\Mc$ satisfies condition $L$, and admits a small strictly generating category consisting of accessible objects.}
% printed: a vertical double arrow $\Uparrow$ at the left, between a) and b)
\item[\llap{$\Uparrow$\ }\textit{b)}] \textit{There exists a small category $C$, and an infinite cardinal $\pi_0$, such that $\Mc$ is equivalent to the category $\mathrm{Ind}'_{\pi_0}(C)$.}
% printed: a vertical double arrow $\Uparrow$ at the left, between b) and c)
\item[\llap{$\Uparrow$\ }\textit{c)}] \textit{$\Mc$ is ``accessible'', i.e. is weakly accessible (every object accessible) and admits a cardinal filtration.}
\end{itemize}
We have b $\Longrightarrow$ a $\Longleftarrow$ c tautologically, and c $\Longrightarrow$ b was seen in loc. cit. 9.18, so we have the sequence
\[
\mathrm{c} \;\Longrightarrow\; \mathrm{b} \;\Longrightarrow\; \mathrm{a}\;,
\]
so it would remain to establish a $\Longrightarrow$ c.

\msp{page 155}

The work accomplished (and the sweat shed) in 4.5 will moreover save us, when the time comes, from quite natural errors in the description of a cardinal filtration, starting from a strictly generating subcategory $S$ : it is neither $\mathrm{Filt}^{\pi}$, nor $\mathrm{Filt}'^{\pi}$, but something subtler -- perhaps one must take the category $\Mc_{\pi}$ again, or else saturate $S$ under the direct limits of cardinality $\leqslant \pi$ which are willing to be representable in $\Mc$ : we then find a $\mathrm{Filt}''^{\pi} \subset \Mc_{\pi}$, and if these $\mathrm{Filt}''^{\pi}$ do indeed give a cardinal filtration, we must necessarily have $\Mc_{\pi} \subset \mathrm{Filt}''^{\pi}$, hence finally $\mathrm{Filt}''^{\pi} = \Mc_{\pi}$. So I expect that in any event, the cardinal filtration (if it exists) will be given by the $\Mc_{\pi}$. Condition a) of loc. cit. 9.12 ($\Mc_{\pi}$ equivalent to a small category), and the bound on cardinals in loc. cit. c), should not cause any difficulties, and condition b) is OK. The difficulty

\msp{page 156}

comes from the fact that $\Mc_{\pi}$ not even being stable under sums and cokernels of double arrows (when these are not represented in $\Mc$), the category $\Mc_{\pi}/x$ (for $x$ in $\Mc$) has no reason to be filtering, as far as I can see, nor a fortiori
\printedpage{51}
% head: A. Grothendieck : Dérivateurs, Ch. XVIII, Édition des Universités de Montpellier II et Paris VII
% foot: Version 17 décembre 2025
to be large relative to $\pi$. So it is not too clear how then to represent $x$ as $\varinjlim\limits_{I} x_i$, the $x_i$ in $\Mc_{\pi}$ and $I$ large relative to $\pi$.

It must indeed be said that in SGA 4 I 9, I give no \textit{example} of a cardinal filtration on a category where arbitrary small $\varinjlim$ would not be representable (cf. loc. cit. 9.13). Perhaps after all it is against nature to want them beyond ? ? Perhaps one must wait until forced to it, to try to develop a theory outside the framework of paratopoi.

\msp{page 157}

\textbf{4.6.7.} We are now in a position to reformulate the litanies of somewhat technical conditions a) to f) of theorem 2.6 (p. 25 in XVIII) concerning the ``special on the left'' character of the set $\mathrm{Mon}(\Mc) \subset \mathrm{Fl}(\Mc)$ :

\textbf{Proposition 4.6.7} (\footnote[62]{Perhaps one could call these categories ``\textit{parabelian paratopos}'' categories ? ?}). \textit{The preliminary conditions on the $\U$-category $\Mc$, ensuring the validity of theorem XVIII 2.6, in a very slightly modified form (namely by modifying in 2.6 d) the condition of stability of} $\mathrm{Filt}^{\pi}$ \textit{under quotients, into an exactness property d) below), are equivalent to the set of the following six conditions a) to f).} % [sic] modifié -> modifiée, agreement with "forme"; équivaut -> équivalent, plural subject "Les conditions"
\DERmark{ENV1176}

\[
\left\{
\begin{array}{ll}
\textit{a)} & \Mc\ \textit{is stable under small}\ \varinjlim\ . \\[1.5ex]
\textit{b)} & \Mc\ \textit{admits a small strictly generating subcategory.} \\[1.5ex]
\textit{c)} & \textit{Filtering direct limits in}\ \Mc\ \textit{are exact.}
\end{array}
\right.
\]
\textit{These conditions imply that $\Mc$ is a} $\U$-paratopos category \textit{(cf. 4.6.3, 3°), and in particular that the objects of $\Mc$ are} accessible, \textit{and that $\Mc$ admits a cardinal filtration $(\Mc_{\pi})_{\pi \geqslant \pi_0}$ such that for every $\pi \geqslant \pi_0$, $\Mc_{\pi}$ is stable under $\varinjlim$ of cardinality $\leqslant \pi$ (in fact, we will have $\Mc_{\pi} =$ full subcategory of the $\pi$-accessible objects of $\Mc$). Also, they imply the stability of $\Mc$ under small $\varprojlim$ . On the other hand, the following three conditions are stability conditions for monomorphisms.}
\begin{itemize}[leftmargin=4em]
\item[\textit{d)}] \textit{Let $f : X \longrightarrow Y$ be an arrow in $\Mc$ and $R = X \times_Y X \rightrightarrows X$, then $X/R \longrightarrow Y$ is a monomorphism -- in other words, $f$ factors as a composite $ip$, with $i$ monomorphic and $p$ an effective epimorphism.}
\end{itemize}

\msp{page 158}

\begin{itemize}[leftmargin=4em]
\item[\textit{e)}] \textit{If $A_1$, $A_2$ are two subobjects of an object $B$ of $\Mc$, and $A_0 = A_1 \cap A_2 \overset{\mathrm{def}}{=} A_1 \times_B A_2$, then $A_1 \amalg_{A_0} A_2 \longrightarrow B$ is a monomorphism. (Note the close kinship of conditions d) and e).)}
\item[\textit{f)}] \textit{The set} $\mathrm{Mon}(\Mc)$ \textit{of monomorphisms in $\Mc$ is stable under cobase change.}
\end{itemize}
\DERmark{ENV1177}

In fact, it is not clear that these conditions, and in particular d), suffice to entail the stability of the $\Mc_{\pi}$ under quotients (only under effective quotients), as postulated in 2.6 d). But in the statement and the proof given in XV cor. 7.4 (p. 98-99), we were content with a somewhat weaker form of this stability property, restricting ourselves to the quotients $Z'$ of $Z \in \mathrm{Ob}\,\Mc_{\pi}$ \textit{such} that there is no subobject $Z'_0$ of $Z'$, distinct from $Z'$, and such that $Z \longrightarrow Z'$ factors through $Z'_0$ ; we therefore required that $Z' \in \mathrm{Ob}\,\Mc_{\pi}$. But if condition d) is satisfied, it follows at once that a quotient $Z'$ of $Z$ having this additional property is necessarily effective, from which it does follow that $Z' \in \mathrm{Ob}\,\Mc_{\pi}$.

\printedpage{52}
% head: A. Grothendieck : Dérivateurs, Ch. XVIII, Édition des Universités de Montpellier II et Paris VII
% foot: Version 17 décembre 2025
\msp{page 159}

\textbf{Corollary 4.6.8.} \textit{Conditions d), e), f) of 4.6.7 are satisfied if $\Mc$ is an abelian category. So the set of conditions a) to f) is satisfied in each of the following two cases :}
\begin{itemize}[leftmargin=4em]
\item[\textit{1)}] \textit{$\Mc$ is an abelian category where filtering direct limits are representable and exact, and admitting a small strictly generating category} (\footnote[63]{This is the case studied in ``Tohoku''.}).
\item[\textit{2)}] \textit{$\Mc$ is a topos.}
\end{itemize}
\textbf{Remark 4.6.9.} I am not sure I know other examples where the conditions of 4.6.7 are verified, than those noted in 4.6.8 above (\footnote[64]{or which follow trivially from them (products etc.)}). One will note that case 1°) mentioned there is that of Tohoku, where I establish the existence of ``enough injective objects'' in such an abelian category. In Tohoku, I confine myself to assuming the existence of a small subcategory, \textit{generating for monomorphisms}, a condition apparently weaker than the one stated in 4.6.8 1°). But in an abelian category, the two notions (and their various variants in SGA 4 I 7.1) are equivalent by virtue of SGA 4 I 7.3 (i). % [sic] par sûr -> pas sûr, negation "ne ... pas"

\msp{page 160}

On the other hand, it seems that conditions e) f) are practically never verified for the categories defined by the species of algebraic structure ``definable by finite products'' and nothing more, such as groups, rings (of all kinds), etc. But the category of \textit{commutative} groups, being abelian, is an exception. (One can moreover interpret it as $\Mc \boxtimes \Mc$, where $\Mc = (\mathrm{Gr})$, and one can then wonder whether there are some generalizations to paratopos categories of this form, $\Mc$ being a paratopos category \ldots) % [sic] tel les groupes -> tels les groupes, agreement; quelque généralisations -> quelques généralisations, plural

\msp{page 161}

\section*{4.7\quad $\mathrm{Ind}_{\pi}C$ ($C$ arbitrary) and generalized cardinal filtrations.}

I feel the need to launch into developments on this subject, to arrive at a better understanding of the notion of the \textit{accessible part} of a paratopos, and of the stability theorem for paratopoi under certain key operations on them.

In what follows, $C$ denotes a category equivalent to a small category, $C_0$ a skeletal exhaustive subcategory, $\pi_0$ an infinite cardinal. We set
\[
\Mc \;=\; \mathrm{Ind}_{\pi_0}C \;\subset\; C^{\wedge} \qquad \text{(\footnotemark[65])}\;. \tag*{(4.7.1)}
\]
\footnotetext[65]{\textbf{N.B.} I now denote by $\mathrm{Ind}C$, $\mathrm{Ind}_{\pi}C$ what was previously denoted $\mathrm{Ind}'C$, $\mathrm{Ind}'_{\pi}C$.}
We propose to introduce on $\Mc$ various $\pm$ ``cardinal'' filtrations, by corresponding properties of an object $F$ of $\Mc \subset C^{\wedge}$, which I am now going to introduce, taking inspiration from 4.5.2 (p. 107) and 4.5.6 (p. 116).

\textbf{Proposition 4.7.1.} \textit{Set} $\pi_1 = \sup(\pi_0, \mathrm{card}\, C_0)$. \textit{Let $\pi \geqslant \pi_0$ be a cardinal, and let us consider the following conditions (depending on $\pi$) on an object $F$ of $\Mc$ :}
\begin{itemize}[leftmargin=4em]
\item[\textit{a)}] $\mathrm{card}\, C_0/F \leqslant \pi$ \textit{(i.e.} $\mathrm{card}\,\mathrm{Fl}\, C_0/F \leqslant \pi$\textit{).}
\item[\textit{a$'$)}] $\mathrm{card}\,\mathrm{Ob}\, C_0/F \leqslant \pi$.
\item[\textit{a$''$)}] $\mathrm{card}\, F(x) \leqslant \pi$ \textit{for every $x \in$} $\mathrm{Ob}\, C$ \textit{(}\textbf{N.B.} \textit{it suffices to require this for $x$ in $C_0$).}
\end{itemize}

\printedpage{53}
% head: A. Grothendieck : Dérivateurs, Ch. XVIII, Édition des Universités de Montpellier II et Paris VII
% foot: Version 17 décembre 2025
\msp{page 162}

\begin{itemize}[leftmargin=4em]
\item[\textit{b)}] \textit{We have $F = \varinjlim\limits_{i \in I} x_i$, with the $x_i \in$} $\mathrm{Ob}\, C$\textit{, $I$ a category (not necessarily ordered) large relative to $\pi_0$, and such that} $\mathrm{card}\, I \leqslant \pi$. \textit{(}\textbf{N.B.} \textit{by virtue of 4.4.2 (p.79), for given $I$, this is equivalent to the existence of a cofinal functor $I \longrightarrow C/F$, or again $I \longrightarrow C_0/F$.)} (\footnote[66]{\textbf{N.B.} We shall denote by $\mathrm{Filt}^{\pi} \subset \Mc$ the strictly full subcategory of $\Mc$ consisting of the $F$ satisfying this condition.}).
\item[\textit{b$'$)}] \textit{We have $F \in$} $\mathrm{Ob}\,\mathrm{Filt}'^{\pi}$\textit{, where} $\mathrm{Filt}'^{\pi}$ \textit{is the smallest strictly full subcategory of $\Mc$ containing $C$, and stable under $\varinjlim$ large relative to $\pi_0$ and of cardinality $\leqslant \pi$.}
\item[\textit{c)}] \textit{We have $F \in$} $\mathrm{Ob}\,\Mc_{\pi}$\textit{, i.e. $F$ is a $\pi$-accessible object of $\Mc$.}
\item[\textit{c$'$)}] \textit{We have $F \in$} $\mathrm{Ob}\, C^{\wedge}_{\pi}$\textit{, i.e. $F$ is a $\pi$-accessible object of $C^{\wedge}$.}
\end{itemize}
\textit{This being set, we have the following diagram of implications.}
\[
\begin{tikzcd}[column sep=large, row sep=large]
\mathrm{a} \arrow[r, Rightarrow] \arrow[dr, Rightarrow] & \mathrm{a}' \arrow[r, Rightarrow] & \mathrm{a}'' \arrow[dr, Rightarrow] & & \\
& \mathrm{b} \arrow[r, Rightarrow] & \mathrm{b}' \arrow[r, Rightarrow] & \mathrm{c}' \arrow[dr, Rightarrow] \arrow[ulll, dotted, bend right=50, "\text{if } \pi \geqslant \pi_1"'] & \\
& & & & \mathrm{c}\;, \arrow[ull, dotted, bend left=30, "\text{if } \pi = \pi^{\pi_0}"']
\end{tikzcd}
\]
\textit{where the two dotted implications, c$'$ $\Longrightarrow$ a and c $\Longrightarrow$ b, are conditional subject to the condition $\pi \geqslant \pi_1$, resp. $\pi = \pi^{\pi_0}$.}

\msp{page 163}

\textbf{Corollary 4.7.2.}
\begin{itemize}[leftmargin=4em]
\item[\textit{1°)}] \textit{If $\pi \geqslant \pi_1$, the six conditions a, a$'$, a$''$, b, b$'$, c$'$ are equivalent, and the condition in question (which can be written $F \in$} $\mathrm{Ob}\,\mathrm{Filt}^{\pi}$ \textit{or $F \in$} $\mathrm{Ob}\,\mathrm{Filt}'^{\pi}$ \textit{indifferently) implies c :} % [sic] où -> ou, alternative "ou ... indifféremment"
\[
\mathrm{Filt}^{\pi} \;=\; \mathrm{Filt}'^{\pi} \;\subset\; \Mc_{\pi} \qquad \textit{if } \pi \geqslant \pi_1 \tag*{(4.7.2.1)}
\]
\DERmark{ENV1178}

\textit{(whereas in general, for $\pi \geqslant \pi_0$, we have}
\[
\mathrm{Filt}^{\pi} \;\subset\; \mathrm{Filt}'^{\pi}\ (\,\subset C^{\wedge}_{\pi} \cap \Mc\,) \;\subset\; \Mc_{\pi} \qquad \textit{if } \pi \geqslant \pi_0\textit{).} \tag*{(4.7.2.2)}
\]
\item[\textit{2°)}] \textit{When $\pi = \pi^{\pi_0}$, e.g. when $\pi$ is of the form $2^{c}$, with $c \geqslant \pi_0$, then conditions b, b$'$, c, c$'$ are equivalent, i.e. we have}
\[
\mathrm{Filt}^{\pi} \;=\; \mathrm{Filt}'^{\pi} \;=\; (C^{\wedge})_{\pi} \cap \Mc \;=\; \Mc_{\pi}\;. \tag*{(4.7.2.3)} % [sic] (4.7.3.2) -> (4.7.2.3), the neighbouring labels are 4.7.2.1, 4.7.2.2, 4.7.2.4
\]
\DERmark{ENV1179}

\item[\textit{3°)}] \textit{When $\pi \geqslant \pi_1$} and \textit{$\pi = \pi^{\pi_0}$, then the seven conditions a), a$'$), a$''$), b) , b$'$), c) , c$'$) are equivalent, so in particular}
\[
\mathrm{Filt}^{\pi} \;=\; \mathrm{Filt}'^{\pi} \;=\; C^{\wedge}_{\pi} \cap \Mc \;=\; \Mc_{\pi} \qquad (\pi \geqslant \pi_1 \textit{ and } \pi = \pi^{\pi_0}\textit{).} \tag*{(4.7.2.4)}
\]
\end{itemize}


\endgroup
% END UNIT der-ch18/der-ch18_p051-055.tex
% BEGIN UNIT der-ch18/der-ch18_p056-060.tex
\begingroup\makeatletter%
\newcommand{\printedpage}[1]{\marginpar{\footnotesize p.~#1}}%
\newcommand{\msp}[1]{\marginpar{\footnotesize [#1]}}%
\newcommand{\Filt}{\mathrm{Filt}}%
\newcommand{\card}{\operatorname{card}}%
\newcommand{\Ob}{\operatorname{Ob}}%
\newcommand{\Fl}{\operatorname{Fl}}%
\newcommand{\Ind}{\mathrm{Ind}}%
\newcommand{\cardred}{\text{card red}\,}%
\makeatother\let\printedpage\derpp\let\msp\dermsp%
\let\mspin\DERmspin%


\printedpage{54}
% head: A. Grothendieck : Dérivateurs, Ch. XVIII, Édition des Universités de Montpellier II et Paris VII

\textsc{Proof} of 4.7.1. The implications a $\Longrightarrow$ a$'$ $\Longrightarrow$ a$''$ are tautological, as is b $\Longrightarrow$ b$'$. The implication b$'$ $\Longrightarrow$ c$'$ comes from the stability of $C^{\wedge}_{\pi}$ under $\varinjlim$ of cardinality $\leqslant \pi$, and from the fact that the $\varinjlim$ of the type under consideration (in particular, large relative to $\pi_0$) \textit{in} $\mathcal{M}$ are the same as in $C^{\wedge}$, i.e. are computed ``argument by argument'' (cf. 4.4.15, p. 99). This, together with the fact that the inclusion $\mathcal{M} \hookrightarrow C^{\wedge}$ is fully faithful, also implies c$'$ $\Longrightarrow$ c.

\msp{page 164}
The implication a$''$ $\Longrightarrow$ c$'$ is seen as in 4.5.6 (p. 119) by applying lemma 4.5.8 1° and 2°, to the case of $C^{\wedge} = \underline{\mathrm{Hom}}(C^{\circ}, \mathcal{E})$, where $\mathcal{E} =$ (Ens). The reverse conditional implication c$'$ $\Longrightarrow$ a$''$, for $\pi \geqslant \pi_1$, is seen as in loc. cit., by applying this lemma 4.5.8. So it still remains to establish the three implications
\[
\left\{
\begin{array}{ccll}
\text{a} & \Longrightarrow & \text{b} & \\[4pt]
\text{a}'' & \Longrightarrow & \text{a} & \text{if } \pi \geqslant \pi_1 \\[4pt]
\text{c} & \Longrightarrow & \text{b} & \text{if } \pi = \pi^{\pi_0}.
\end{array}
\right. \tag*{$(*)$}
\]
But a $\Longrightarrow$ b is tautological, on taking $I = C_0/F$, since we know that $C_0/F$ is large relative to $\pi_0$ and that $F = \varinjlim\limits_{i \in I} x_i$, practically by definition of $\mathcal{M} = \underbrace{\Ind_{\pi_0}(C)}_{\simeq \Ind_{\pi_0} C_0} \subset C^{\wedge}$. So it remains to establish the two conditional implications of $(*)$. The implication a$''$ $\Longrightarrow$ a for $\pi \geqslant \pi_1$ is established as in 4.5.2 (cf. page 109), it is essentially sorital. It remains to establish
\DERmark{ENV1180}

\[
\text{c} \;\Longrightarrow\; \text{b} \qquad \text{if} \;\; \pi \;=\; \pi^{\pi_0} \;. \tag*{$(*)$}
\]
We will need the auxiliary result:

\msp{page 165}
\textbf{Lemma 4.7.3.} \textit{Let $\pi \geqslant \pi_0$ be such that $\pi^{\pi_0} = \pi$ (thus in fact $\pi > \pi_0$). For every object $F$ of $\mathcal{M}$, we have}
\[
F \;=\; \varinjlim_{i \in I} x_i \qquad \text{(\footnotemark[67])}\;. \tag*{(4.7.3.1)}
\]
\footnotetext[67]{the $x_i$ in $\Filt^{\pi}$, $I$ an ordered set large relative to $\pi$.}
\textit{If in addition $\pi' > \pi$ and $F \in \Ob \Filt^{\pi'}$, then we can take $I$ such that}
\[
\card I \;\leqslant\; \pi'^{\pi} \;. \tag*{(4.7.3.2)}
\]

Indeed, we have
\[
F \;=\; \varinjlim_{j \in J} x_j \qquad \text{with } J = C_0/F \qquad \text{(\footnotemark[68])}\;,
\]
\footnotetext[68]{where $J$ can also be replaced by a category cofinal in [?] $C_0/F$ large relative to $\pi_0$.}
with $J$ large relative to $\pi_0$. Then let
\[
I \;\;=\;\; \parbox[t]{0.62\textwidth}{set of subcategories $J'$ of $J$ \textit{large relative to} $\pi_0$, and such that $\card J' \leqslant \pi$.} \tag*{(4.7.3.3)}
\]

\printedpage{55}
% head: A. Grothendieck : Dérivateurs, Ch. XVIII, Édition des Universités de Montpellier II et Paris VII

I say that $I$ (ordered by inclusion) \textit{is large relative to} $\pi$ (\footnote[69]{cf. proof 4.7.5, page 168-171.}). This follows at once from lemma 4.5.5, p. 113, (stated at first, erroneously, assuming only $\pi > \pi_0$ and not $\pi^{\pi_0} = \pi$). Let us then put, for every $J' \in I$,
\[
x_{J'} \;=\; \varinjlim_{j \in J'} x_j \;,
\]
we thus have
\[
x_{J'} \;\in\; \Ob \Filt^{\pi}
\]
by definition of $I$ and of $\Filt^{\pi}$. On the other hand,
\[
J' \;\longmapsto\; x_{J'} \;:\; I \;\longrightarrow\; \Filt^{\pi}
\]
\msp{page 166}
is a functor, hence
\[
\varinjlim_{J' \in I} x_{J'} \;\simeq\; \varinjlim_{j \in J} x_j \;\simeq\; F \;, % [sic] x -> F, the colimit is the object F of (4.7.3.1)
\]
\DERmark{ENV1181}

whence the first assertion of 4.7.3. There remains the bound on cardinals for $F \in \Ob \Filt^{\pi'}$. In this case, replacing $C_0/F$ if need be by a cofinal category large relative to $\pi_0$, $\card J \leqslant \pi'$. We then see at once that $\card I \leqslant \pi'^{\pi}$, which is the announced bound.
\DERmark{ENV1182}\DERmark{ENV1183}


We can now prove the implication $(*)$, page 164, i.e.
\[
\mathcal{M}_{\pi} \;\subset\; \Filt^{\pi} \qquad (\text{for } \pi = \pi^{\pi_0}).
\]
So let $F$ be in $\mathcal{M}_{\pi}$, i.e. $\pi$-accessible, let us write $F$ in the form (4.7.3.1), with $I$ large relative to $\pi$. By definition of $\pi$-accessibility, applied to $F \xrightarrow{\mathrm{id}} F \simeq \varinjlim F_i$, we see that $F$ is a direct factor of an $F_i$, hence of an object of $\Filt^{\pi}$. Not knowing that $\Filt^{\pi}$ is stable under direct factors, we cannot conclude from this that $F$ is in $\Filt^{\pi}$ (so I cannot prove c $\Longrightarrow$ b), but only that $F$ is in $\Filt'^{\pi}$, which, for its part, is stable under $\varinjlim$ large relative to $\pi_0$, of given cardinality $\pi$ $> \pi_0$, hence also \msp{page 167} stable under direct factors. This completes the proof of 4.7.1, and at the same time we get:
\DERmark{ENV1184}


% marginal vertical bar beside Corollaire 4.7.4 and its statement
\textbf{Corollary 4.7.4.} \textit{If $\pi \geqslant \pi_0$ is such that $\pi^{\pi_0} = \pi$, then we have}
\[
\Filt'^{\pi} \;=\; (\Filt^{\pi})^{\natural} \;, \tag*{(4.7.4.1)}
\]
\textit{i.e.} $\Filt'^{\pi}$ \textit{consists of the objects $F$ of $\mathcal{M}$ which are direct factors of an object of} $\Filt^{\pi}$.

Indeed, since $\Filt^{\pi} \subset \Filt'^{\pi}$, and since $(\Filt^{\pi})^{\natural} \subset \Filt'^{\pi}$, and since $\Filt'^{\pi}$ is stable under direct factors, we have $(\Filt^{\pi})^{\natural} \subset \Filt'^{\pi}$ for every $\pi \geqslant \pi_0$. But conversely, since $\Filt'^{\pi} \subset \mathcal{M}_{\pi}$, the preceding argument shows that for $\pi = \pi^{\pi_0}$, we have $\Filt'^{\pi} \subset \mathcal{M}_{\pi} \subset (\Filt^{\pi})^{\natural}$, which proves 4.7.4.
\DERmark{ENV1185}


\textbf{Remark 4.7.4.1} \texttt{[twice 4.7.4.1]}. Of course, by 4.7.2, 1°, as soon as $\pi \geqslant \pi_1$, we even have $\Filt^{\pi} = \Filt'^{\pi}$ (without having to assume $\pi = \pi^{\pi_0}$), hence \textit{if} $\pi > \pi_0$ (which will be the case if $\pi_1 > \pi_0$, i.e. $\pi_1 \neq \pi_0$). $\Filt'^{\pi}$ being stable under direct factors, $\Filt^{\pi}$ is so too, (4.7.4.1) is satisfied trivially. But if $\pi = \pi_1 = \pi_0$, we have $\Filt^{\pi} = \Filt^{\pi_0} = \Filt'^{\pi_0} =$ subcategory of functors \textit{representable} in $C$ (equivalent to $C$), and the latter is equal to $C^{\natural}$ in $\mathcal{M}$ only if $C$ is Karoubian, cf. \ldots{} below. % [sic] la cas -> le cas, gender of "cas"

\printedpage{56}
% head: A. Grothendieck : Dérivateurs, Ch. XVIII, Édition des Universités de Montpellier II et Paris VII

\msp{page 168}
\textbf{4.7.5.} For the proof of 4.7.1 to be complete, I would still have to give a proof of 4.5.5, p. 113, which is the key technical lemma for establishing the

\textbf{Lemma 4.7.5.1} (used in the proof of 4.7.3). \textit{Let $\pi_0$, $\pi$ be two infinite cardinals, with $\pi \geqslant \pi_0$ and $\pi^{\pi_0} = \pi$ (hence $\pi > \pi_0$). Let $J$ be a category large relative to $\pi_0$. Then the set $I$, ordered by inclusion, of subcategories $J'$ of $J$ which are large relative to $\pi_0$, and such that} $\card J' \leqslant \pi$, \textit{is large relative to $\pi$. (Moreover, if} $\card J \leqslant \pi'$, \textit{we have} $\card I \leqslant \pi'^{\pi}$\textit{.)}

So we must show the

\textbf{Corollary 4.7.5.2} (\footnote[70]{Compare with the proof of 4.6.1.2.}) (this is lemma 4.5.5 not yet proved). \textit{Under the conditions of 4.7.5, let $K \subset \mathrm{Fl} J$ be such that} $\card K \leqslant \pi$. \textit{Then there exists $J' \in I$ such that $K \subset J'$.}

We know, $\pi$ being infinite, that the category $K'$ generated by $K$ is also of cardinality $\leqslant \pi$, so we can suppose that $K$ ``is'' a subcategory of $J$. We are going to construct a $J' \in I$ containing $K$, by constructing by transfinite recursion an increasing sequence $K_{\alpha}$ of subcategories of $J$, with $\card K_{\alpha} \leqslant \pi$, indexed by the ordinals $\alpha$ such that \fbox{$\card \alpha \leqslant \pi$}:

\msp{page 169}
\begin{enumerate}
\item[1°)] $K_0 = K$.
\item[2°)] $K_{\alpha} = \bigcup_{\beta < \alpha} K_{\beta}$ ($= \varinjlim K_{\beta}$) if $\alpha$ is a limit ordinal (hence $\card K_{\alpha} \leqslant \pi$, since the $\card K_{\beta} \leqslant \pi$).
\item[3°)] $K_{\alpha+1}$ obtained from $K_{\alpha}$ in the following way. If $K_{\alpha}$ is large relative to $\pi_0$, we take $K_{\alpha+1} = K_{\alpha}$ (hence, in view of 2°, we will have $K_{\alpha'} = K_{\alpha}$ for $\alpha' \geqslant \alpha$ -- the transfinite construction stops \ldots). Otherwise, $K_{\alpha+1}$ is obtained as follows:
\begin{enumerate}
\item[a)] for every subset $A$ of cardinality $\leqslant \pi_0$ of $\Ob K_{\alpha}$, we choose an upper bound $x_A$ in $J$ (we feel like taking it in $K_{\alpha}$, if there is an upper bound in $K_{\alpha}$)\,;
\item[b)] for every pair $x$, $y$ of objects of $K_{\alpha}$, and every subset of cardinality $\leqslant \pi_0$ of $\mathrm{Hom}_{K_{\alpha}}(x,y)$, we choose an arrow $u_B : y \longrightarrow z_B$ in $J$, which equalizes the $v_{\beta} \in B$ (and we take this arrow in $K_{\alpha}$, when possible)\,;
\item[c)] we take for $K_{\alpha+1}$ the smallest subcategory containing the $\mathrm{id}_{x_A}$ and the $u_B$.
\end{enumerate}
\end{enumerate}

We must check that
\[
\card K_{\alpha+1} \;\leqslant\; \pi \;.
\]
The number of possible choices for $A \in \mathfrak{P}_{\leqslant \pi_0}(\Ob K_\alpha)$ \msp{page 170} is $\leqslant \pi^{\pi_0} = \pi$, the number of possible choices for $B \in \mathfrak{P}_{\pi_0}(\mathrm{Fl}\, K_\alpha)$ is also $\leqslant \pi^{\pi_0} = \pi$, so the number of arrows of the form $\mathrm{id}_{x_A}$ or $u_B$ is $\leqslant \pi + \pi = \pi$. But then the subcategory $K_{\alpha+1}$ of $J$ generated by this set of arrows is of cardinality $\leqslant \pi$. Then let % [sic] \leqslant 1 -> \leqslant \pi, the statement to be verified is card K_{\alpha+1} \leqslant \pi
\DERmark{ENV1186}\DERmark{ENV1187}\DERmark{ENV1188}

\printedpage{57}
% head: A. Grothendieck : Dérivateurs, Ch. XVIII, Édition des Universités de Montpellier II et Paris VII
$\alpha_0$ be the least ordinal such that
\[
\card \alpha_0 \;>\; \pi_0 \;,
\]
we thus have
\[
\card \alpha_0 \;=\; \text{successor cardinal of } \pi_0 \leqslant \pi \;,
\]
since $\pi^{\pi_0} = \pi$, whence $\pi > \pi_0$ (for $\pi^{\pi_0} = \pi_0^{\pi_0} \geqslant 2^{\pi_0} > \pi_0$). Thus, $\alpha_0$ is among the admissible ordinals $\alpha$ (i.e. such that $\card \alpha \leqslant \pi$). I say that $K_{\alpha_0}$ \textit{is large relative to} $\pi_0$ (which will prove 4.7.5.2, on taking $J' = K_{\alpha_0}$). This follows immediately from the fact that
\[
K_{\alpha_0} \;=\; \bigcup_{\alpha < \alpha_0} K_{\alpha} \qquad (\text{\textit{increasing} union}),
\]
% double marginal vertical bar beside the following two lines
and that the set of $\alpha < \alpha_0$ is \textit{large relative to} $\pi_0$ (\footnote[71]{spell this out in a

\textbf{Lemma.} \textit{Let $\pi_0$ be a cardinal, $\alpha_0 = \alpha_{\pi'_0}$ ($\pi'_0$ the successor of $\pi_0$) the least ordinal such that} $\card \alpha_0 > \pi_0$, \textit{then $I_{\alpha_0} = \{\alpha \mid \alpha < \alpha_0\}$ is large relative to $\pi_0$.}}). So every set of objects or arrows of $K_{\alpha_0}$, of cardinality $\leqslant \pi_0$, is contained in some $K_{\alpha}$ with $\alpha < \alpha_0$.

\msp{page 171}
This implies, if it is a set of objects, that there is a common upper bound in $K_{\alpha+1}$, a fortiori in $K_{\alpha_0}$. If it is a subset of a $\mathrm{Hom}_{K_{\alpha}}(x,y)$, that there is an equalizer $y \longrightarrow z$ in $K_{\alpha+1}$, a fortiori in $K_{\alpha_0}$. Thus the criterion $\mathrm{PF}_{\pi_0}$ 1,2 of 4.4.4 (page 89) is satisfied by $K_{\alpha_0}$, which completes the proof of 4.7.5.2.
\DERmark{ENV1189}


\textbf{Corollary 4.7.6} (of 4.7.1). \textit{For every cardinal $\pi \geqslant \pi_0$, we have}
\[
\underbrace{\Filt'^{\pi}}_{\overset{\textit{if } \pi \geqslant \pi_1}{=}\ \Filt^{\pi}} \;\subset\; \mathcal{M}_{\pi} \;\subset\; \underbrace{\Filt'^{\pi^{\pi_0}}}_{\overset{\textit{if } \pi^{\pi_0} \geqslant \pi_1}{=}\ \Filt^{\pi^{\pi_0}}} \qquad (= (\Filt^{\pi^{\pi_0}})^{\natural} \textit{ by cor. 4.7.4})\,, \tag*{(4.7.6.1)} % ?? uncertain: tiny condition under the second brace -- "si \pi^{\pi_0} \geqslant \pi_1" or "si \pi^{\pi}_0 \geqslant \pi_1"
\]
\DERmark{ENV1190}

\textit{therefore if $\pi \geqslant \pi_1$, we also have}
\[
\Filt^{\pi} \;\subset\; \mathcal{M}_{\pi} \;\subset\; \Filt^{\pi^{\pi_0}} \qquad \textit{(if $\pi^{\pi_0} \geqslant \pi_1$).} \tag*{(4.7.6.2)}
\]

Indeed, $\Filt'^{\pi} \subset \mathcal{M}_{\pi}$ is none other than the implication b$'$ $\Longrightarrow$ c of 4.7.1, on the other hand, we have $\pi \leqslant \pi^{\pi_0}$, whence
\[
\mathcal{M}_{\pi} \;\subset\; \mathcal{M}_{\pi^{\pi_0}} \;=\; \Filt'^{\pi^{\pi_0}} \;,
\]
where the last equality comes from $\mathcal{M}_{\pi'} = \Filt'^{\pi'}$ if $\pi'^{\pi_0} = \pi'$, i.e. from (4.7.2.3) (page 163), putting $\pi' = \pi^{\pi_0}$, taking into account that $\pi'^{\pi_0} = (\pi^{\pi_0})^{\pi_0} = \pi^{\pi_0 \cdot \pi_0} = \pi^{\pi_0} = \pi'$.
\DERmark{ENV1191}


\msp{page 172}
\textbf{Proposition 4.7.7.} \textit{Consider $C$, $C_0$, $\mathcal{M} = \Ind_{\pi_0}(C)$, $\pi_1 = \sup(\pi_0, \card \mathrm{Fl}\, C_0)$ as before. For every cardinal $\pi \geqslant \pi_0$, let $\mathcal{F}^{\pi}$ be the strictly full subcategory} $\Filt'^{\pi}$ \textit{(resp. $(C^{\wedge})_{\pi} \cap \mathcal{M}$, resp. $\mathcal{M}_{\pi}$) in $\mathcal{M}$. Then the increasing filtration $(\mathcal{F}^{\pi})_{\pi \geqslant \pi_0}$ of $\mathcal{M}$ by the subcategories $\mathcal{F}^{\pi}$ satisfies the following conditions.}
\begin{enumerate}
\item[\textit{a)}] \textit{For every $\pi \geqslant \pi_0$, $\mathcal{F}^{\pi}$ is equivalent to a small category.}
\item[\textit{b)}] \textit{$\mathcal{M}$ satisfies $\mathrm{L}_{\pi_0}$ (i.e. is stable under filtering $\varinjlim$ large relative to $\pi_0$), and for every $\pi \geqslant \pi_0$, $\mathcal{F}^{\pi}$ is stable under filtering limits large relative to $\pi_0$ and of cardinality $\leqslant \pi$ [?].}
\printedpage{58}
% head: A. Grothendieck : Dérivateurs, Ch. XVIII, Édition des Universités de Montpellier II et Paris VII
\item[\textit{c)}] \textit{If $\pi \geqslant \pi_0$ is} such that \fbox{$\pi^{\pi_0} = \pi$,} \textit{then for every $F$ in $\mathcal{M}$, we have an isomorphism}
\[
F \;=\; \varinjlim_{i \in I} x_i \;,
\]
\DERmark{ENV1192}

\textit{the $x_i$ in $\mathcal{F}^{\pi}$, and $I$ an ordered category large relative to $\pi$. If in addition $\pi' > \pi$ and $F \in\Ob\mathcal F^{\pi'}$, we can take $I$ such that} % [sic] F^{\pi'} (italic F) -> \mathcal{F}^{\pi'}, the filtration is denoted \mathcal{F}^{\pi}
\DERmark{ENV1193}\DERmark{ENV1194}

\[
\card I \;\leqslant\; \pi'^{\pi} \qquad \underline{\textit{if}}\ \pi'^{\pi} \geqslant \pi_1. \tag*{(4.7.7.1)}
\]
\item[\textit{d)}] \textit{If $\pi \geqslant \pi_0$, we have $\mathcal{F}^{\pi} \subset \mathcal{M}_{\pi}$, i.e. the objects of $\mathcal{F}^{\pi}$ are $\pi$-accessible.}
\end{enumerate}
\textbf{N.B.} These are exactly, apart from d), the conditions a), b), c) for cardinal filtrations of SGA 4 I 9.7.12, \msp{page 173} except only that in c), we restrict ourselves here to the ordinals $\pi$ such that $\pi^{\pi_0} = \pi$ (\footnote[72]{and \textit{in addition}, in the bound (4.7.7.1), to the cardinals $\pi' > \pi$ such that we also have $\pi'^{\pi} \geqslant \pi_1$ (which is automatic if we have $\pi \geqslant \pi_1$, which is the case, given that $\pi = \pi^{\pi_0} \geqslant 2^{\pi_0}$, if we have $2^{\pi_0} \geqslant \pi_1$.)}). The additional condition d) is an ``accessibility condition'', cf. below.

\textsc{Proof of} 4.7.7. a) Since we have
\[
\Filt'^{\pi} \;\subset\; (C^{\wedge})_{\pi} \cap \mathcal{M} \;\subset\; \mathcal{M}_{\pi} \;,
\]
we are reduced to proving that $\mathcal{M}_{\pi}$ is equivalent to a small category. We are reduced to the case $\pi \geqslant \pi_1$, $\pi = \pi^{\pi_0}$, in which case we have $\mathcal{M}_{\pi} = \Filt^{\pi}$. In fact, we find:

\textbf{Corollary 4.7.7.1} \texttt{[twice 4.7.7.1]} (\footnote[73]{Compare with 4.7.19.}). \textit{If $\pi \geqslant \pi_1$, we have}
\[
\left\{
\begin{array}{rcll}
\cardred \Filt^{\pi} & \leqslant & 2^{\pi} & \\[4pt]
\cardred \Filt'^{\pi} & \leqslant & \cardred \mathcal{M}_{\pi} \;\leqslant\; 2^{(\pi^{\pi_0})} & \qquad \textit{($= 2^{\pi}$ if $\pi^{\pi_0} = \pi$).}
\end{array}
\right.
\]

The second bound follows from the first and from
\[
\mathcal{F}^{\pi} \;\subset\; \mathcal{M}_{\pi} \;\subset\; \Filt^{\pi^{\pi_0}} \qquad (\text{for } \pi^{\pi_0} \geqslant \pi_1) \qquad \text{(\footnotemark[74]).}
\]
\footnotetext[74]{Cf. (4.7.6.2).}
For the first, we use condition a$''$) of 4.7.1 (which characterizes $\Filt^{\pi}$ for $\pi \geqslant \pi_1$). I pass over the details of the verification.

\msp{page 174}
Let us prove 4.7.7 b). The condition $\mathrm{L}_{\pi_0}$ is known, there remains the stability of $\mathcal{F}^{\pi}$. In the case of $\Filt'^{\pi}$, it is established by definition. In the cases $C^{\wedge}_{\pi} \cap \mathcal{M}$ and $\mathcal{M}_{\pi}$, it follows from the stability of the notion of $\pi$-accessibility under $\varinjlim$ of cardinality $\leqslant \pi$, either in $C^{\wedge}$ or in $\mathcal{M}$ (taking into account that $C^{\wedge}$ and $\mathcal{M}$ have a generating family $C$ formed of $\pi_0$-accessible objects, hence $\pi$-accessible).
\DERmark{ENV1195}


Let us prove c): this is essentially lemma 4.7.3, taking into account that we have
\[
\Filt^{\pi} \;\subset\; \mathcal{F}^{\pi} \;,
\]


\endgroup
% END UNIT der-ch18/der-ch18_p056-060.tex
% BEGIN UNIT der-ch18/der-ch18_p061-065.tex
\begingroup\makeatletter%
\newcommand{\printedpage}[1]{\marginpar{\footnotesize p.~#1}}%
\newcommand{\msp}[1]{\marginpar{\footnotesize [#1]}}%
\newcommand{\Ob}{\operatorname{Ob}}%
\newcommand{\Filt}{\mathrm{Filt}}%
\newcommand{\Ind}{\operatorname{Ind}}%
\newcommand{\card}{\operatorname{card}}%
\newcommand{\cardred}{\operatorname{card\,red}}%
\makeatother\let\printedpage\derpp\let\msp\dermsp%
\let\mspin\DERmspin%


\noindent\printedpage{59}% head: A. Grothendieck : Dérivateurs, Ch. XVIII, Édition des Universités de Montpellier II et Paris VII
% foot: Version 17 décembre 2025
except that one must take care with the cardinal bound, $\card I \leqslant \pi'^{\pi}$. But if $F \in \Ob\mathcal{F}^{\pi'}$, since by (4.7.6.1)
\[
\mathcal{F}^{\pi'} \;\subset\; \underbrace{\mathcal{M}_{\pi'}}_{\subset\mathcal{M}_{\pi'^{\pi}}\subset\Filt'^{\pi'^{\pi}}} \;\subset\; \Filt'^{\pi'^{\pi_0}} ,
\]
hence a fortiori
\[
\mathcal{F}^{\pi'} \;\subset\; \Filt'^{\pi'^{\pi}} \underbrace{\;=\;}_{\text{if } \pi'^{\pi}\geqslant\pi_1} \Filt^{\pi'^{\pi}} , % [sic] first \Filt^{\pi'^{\pi}} -> \Filt'^{\pi'^{\pi}}, otherwise the equality under the brace would be trivial (cf. the preceding display)
\]
\DERmark{ENV1196}

whence, by 4.7.3, with $\pi'$ replaced by $\pi'^{\pi}=\pi''$, the possibility of taking $I$ with
\[
\card I \;\leqslant\; \pi''^{\pi} \;=\; (\pi'^{\pi})^{\pi} \;=\; \pi'^{\pi\cdot\pi} \;=\; \pi'^{\pi} ,
\]
which completes the proof.

\noindent\msp{page 175}\textbf{Remark 4.7.7.2.} I was obliged, in the statement of condition c), to bring in a second cardinal $\pi_1$, and this for lack of knowing whether, for $\pi''$ such that $\pi''=\pi''^{\pi_0}$, $\Filt^{\pi''}$ is stable under direct factors. (To do things properly, one would have to either prove that it is stable, hence that we have in 4.7.1 c $\Longrightarrow$ b, i.e. $\mathcal{M}_\pi\subset\Filt^{\pi}$ (and not only c $\Longrightarrow$ b$'$, i.e. $\mathcal{M}_\pi\subset\Filt'^{\pi}$) if $\pi=\pi^{\pi_0}$, or else make a counterexample!) This somewhat strange ingredient, we shall alas be obliged to drag it along in the definition of generalized cardinal filtrations, if we do not want to forbid ourselves to regard in the present case $(\mathcal{F}^{\pi})_{\pi\geqslant\pi_0}$ as such a filtration, and to allow ourselves only to regard $(\mathcal{F}^{\pi})_{\pi\geqslant\pi_1}$ as such, which would seem ridiculous to me. In practice (even if theoretically the introduction of $\pi_1$ in the description of a generalized cardinal filtration were indeed necessary) we shall hardly have cases where it is necessary to make mention of $\pi_1$ in condition c. Indeed, $\pi\geqslant\pi_0$ and $\pi^{\pi_0}=\pi$ imply $\pi\geqslant 2^{\pi_0}$, and then $\pi'^{\pi}\geqslant 2^{(2^{\pi_0})}$. If therefore $\pi_1\leqslant 2^{(2^{\pi_0})}$, i.e. $\card C_0\leqslant 2^{(2^{\pi_0})}$, then the condition $\pi'^{\pi}\geqslant\pi_1$ in c) is automatically satisfied, and there is therefore no need to mention it. Now it is very rare to have to deal with small categories $C_0$, whose cardinality is $>$ the power % [sic] faut -> faute, "faute de savoir" intended; [sic] \pi^{\pi'}\geqslant\pi_1 -> \pi'^{\pi}\geqslant\pi_1, the condition is written \pi'^{\pi}\geqslant\pi_1 elsewhere (cf. 4.7.8)
\DERmark{ENV1197}

\msp{page 176}
of the continuum $c$, and a fortiori, such that $\card C_0 > 2^{c}$. If on the other hand $\card C_0\leqslant 2^{c}$, since $2^{c}\leqslant 2^{(2^{\pi_0})}$ (since $\pi_0\geqslant\aleph_0$), we shall indeed have $\pi_1\leqslant 2^{(2^{\pi_0})}$, and we can rest easy.

\textbf{Definition 4.7.8.} Let $\mathcal{M}$ be a $\mathfrak{U}$-category. We call a $\mathfrak{U}$-\textit{generalized cardinal filtration} (or simply generalized cardinal filtration, when $\mathfrak{U}$ is understood) the datum of an infinite cardinal $\pi_0$, and of a filtration $(\mathcal{F}^{\pi})_{\pi\geqslant\pi_0}$ of $\mathcal{M}$ by strictly full subcategories $\mathcal{F}^{\pi}$, indexed by the $\mathfrak{U}$-cardinals $\geqslant\pi_0$, and satisfying conditions a) b) c) stated in 4.7.7, where in c) we put
\[
\pi_1 \;=\; \sup(\pi_0,\cardred\mathcal{F}^{\pi_0}) \, . \tag*{(4.7.8.1)}
\]
(Thus if $\cardred\mathcal{F}^{\pi_0}\leqslant 2^{(2^{\pi_0})}$, there is no need to take account of the condition $\pi'^{\pi}\geqslant\pi_1$, cf. 4.7.7.1.) We say that the generalized cardinal filtration (fcg) is \textit{accessible}, if it satisfies in addition condition d). We say that it is \textit{canonical} if we even have
\[
\mathcal{F}^{\pi} \;=\; \mathcal{M}_{\pi} \tag*{(4.7.8.2)}
\]
for every $\pi\geqslant\pi_0$.

\textbf{Proposition 4.7.9.} \textit{Let $(\mathcal{F}^{\pi})_{\pi\geqslant\pi_0}$ be an fcg on the $\mathfrak{U}$-category $\mathcal{M}$, and $\pi>\pi_0$ a cardinal such that $\pi^{\pi_0}=\pi$. Then we have}
\printedpage{60}% head: A. Grothendieck : Dérivateurs, Ch. XVIII, Édition des Universités de Montpellier II et Paris VII
% foot: Version 17 décembre 2025
\msp{page 177}
\[
\left\{
\begin{aligned}
&\mathcal{M}_\pi\subset\mathcal{F}^{\pi}\quad\text{\textit{(if $\pi=\pi^{\pi_0}$),}}\\
&\text{\textit{hence} }\mathcal{M}_\pi=\mathcal{F}^{\pi}\text{ \textit{if moreover the fcg is} accessible.}
\end{aligned}
\right.
\tag*{(4.7.9.1)}
\]

\textsc{Proof} (\footnote[75]{\texttt{[In what follows, $\mathcal{F}^{\pi}$ has sometimes been written $\mathcal{F}_{\pi}$, often corrected to $\mathcal{F}^{\pi}$ by the author. We write $\mathcal{F}^{\pi}$ throughout.]}}). Applying condition c) for fcg to an $F\in\Ob\mathcal{M}_\pi$, we find that $F$ is a direct factor of an $x_i\in\Ob\mathcal{F}^{\pi}$. Since $\mathcal{F}^{\pi}$ is stable under $\varinjlim$ large relative to $\pi_0$ and of cardinality $\leqslant\pi$, and since \textit{in addition} $\pi>\pi_0$, we conclude that:

\textbf{Lemma 4.7.9.1} \texttt{[twice 4.7.9.1]}. \textit{If $\pi>\pi_0$ (\footnote[76]{But $\mathcal{F}^{\pi_0}$ is not necessarily stable under direct factors, as we see by taking $\mathcal{M}=\Ind_{\pi_0}(C)$, $C$ a small non-Karoubian category, and the cardinal filtration for the $\Filt'^{\pi}$ -- we note that $\Filt'^{\pi_0}=\Filt^{\pi_0}\simeq e$, since a category large relative to $\pi_0$ and of cardinality $\leqslant\pi_0$ is cofinal to $e$.}), $\mathcal{F}^{\pi}$ is stable under direct factors.}

Whence it follows that $F\in\mathcal{F}^{\pi}$, whence (4.7.9.1).

\textbf{Corollary 4.7.10.} \textit{Suppose the fcg accessible. For every cardinal $\pi\geqslant\pi_0$, we have}
\[
\mathcal{F}^{\pi} \;\subset\; \mathcal{M}_\pi \;\subset\; \mathcal{F}^{\pi^{\pi_0}} \; ( \;=\; \mathcal{M}_{\pi^{\pi_0}}) \, . \tag*{(4.7.10.1)}
\]

Indeed, $\mathcal{F}^{\pi}\subset\mathcal{M}_\pi$ by the definition of accessible fcg, and $\mathcal{M}_\pi\subseteq\mathcal{M}_{\pi^{\pi_0}}$ (since $\pi\leqslant\pi^{\pi_0}$), now $\mathcal{M}_{\pi^{\pi_0}}=\mathcal{F}^{(\pi^{\pi_0})}$ by virtue of (4.7.9.1) applied to $\pi'=\pi^{\pi_0}$ instead of $\pi$ (since $\pi'^{\pi_0}=\pi'$), whence (4.7.10.1). We shall note that this inclusion generalizes (4.7.9.1), it gives $\mathcal{F}^{\pi}=\mathcal{M}_\pi$ if $\pi=\pi^{\pi_0}$.
\DERmark{ENV1198}


\textbf{4.7.11.} It seems to me that the only useful generalized cardinal filtrations are those which are \textit{accessible.} In what follows, I am going to confine myself to this case. Since in that case $\mathcal{F}^{\pi}\subset\mathcal{M}_\pi$, we conclude from this
\msp{page 178}
that the canonical functor
\[
\Ind_\pi(\mathcal{F}^{\pi}) \;\hookrightarrow\; \mathcal{M} \tag*{(4.7.11.1)}
\]
extending the inclusion $\mathcal{F}^{\pi}\subset\mathcal{M}$, so as to commute with $\varinjlim$ large relative to $\pi$, is \textit{fully faithful.} Moreover, when $\pi^{\pi_0}=\pi$, condition c) for fcg tells us among other things that this functor is essentially surjective, hence that it is an \textit{equivalence of categories}:
\[
\Ind_\pi(\mathcal{F}^{\pi}) \xrightarrow{\ \sim\ } \mathcal{M} \qquad (\text{if } \pi=\pi^{\pi_0}). \tag*{(4.7.11.2)}
\]

\textbf{Proposition 4.7.12.} \textit{Let $(\mathcal{F}^{\pi})_{\pi\geqslant\pi_0}$ be an accessible fcg (fcga) on the $\mathfrak{U}$-category $\mathcal{M}$, and let $\pi$ be a cardinal $\geqslant\pi_0$. The following conditions are equivalent.}
\begin{enumerate}[label=\textit{\alph*)}]
\item \textit{The fully faithful functor (4.7.11.1) is essentially surjective, i.e. is an equivalence of categories.}
\item \textit{For every $F\in\Ob\mathcal{M}$, $\mathcal{F}^{\pi}/F$ is a filtering category and large relative to $\pi$.}
\end{enumerate}

\noindent\printedpage{61}% head: A. Grothendieck : Dérivateurs, Ch. XVIII, Édition des Universités de Montpellier II et Paris VII
% foot: Version 17 décembre 2025
\textit{When these conditions are satisfied, and if $\pi>\pi_0$, we have}
\[
\mathcal{F}^{\pi} \;=\; \mathcal{M}_\pi \, . \tag*{(4.7.12.1)}
\]

\textsc{Proof}. The equivalence of a) and b) follows at once from 4.4.2 and from the definition of categories large relative to $\pi$ (p.~81). That these
\msp{page 179}
conditions imply $\mathcal{F}^{\pi}=\mathcal{M}_\pi$, i.e. $\mathcal{M}_\pi\subset\mathcal{F}^{\pi}$ in the case where $\pi>\pi_0$, is seen as for (4.7.9.1) by using Lemma 4.7.9.1.

\textbf{Definition 4.7.13.}
\begin{enumerate}[label={$\arabic*^{\circ}$)}]
\item A cardinal $\pi\geqslant\pi_0$ is called a \textit{useful cardinal} for the fcga $(\mathcal{F}^{\pi})_{\pi\geqslant\pi_0}$, if it satisfies the equivalent conditions a) and b) of 4.7.12, and if in addition in the case where $\pi=\pi_0$, we have $\mathcal{F}^{\pi}=\mathcal{M}_\pi$. (Thus if $\pi$ is a useful cardinal, we always have $\mathcal{F}^{\pi}=\mathcal{M}_\pi$.) In the contrary case, we say that $\pi$ is a \textit{superfluous cardinal} (for the given fcga). % [sic] le fcga donnée -> la fcga donnée, gender agreement (feminine)
\item A cardinal $\pi$ is called \textit{regular} for the given fcg, if it is useful, and if in addition the cardinal bound stated in (4.7.7.1) is valid.
\end{enumerate}

\textbf{4.7.14} Thus, all cardinals $\pi\geqslant\pi_0$ such that $\pi^{\pi_0}=\pi$ are useful, and even regular for the given \texttt{[generalized]} cardinal filtration. Let us note that if all cardinals $\pi\geqslant\pi_0$ are useful, then we have $\mathcal{F}^{\pi}=\mathcal{M}_\pi$ for every $\pi\geqslant\pi_0$, hence the fcga is canonical. Let us also note that in all cases, if $(\mathcal{F}^{\pi})_{\pi\geqslant\pi_0}$ is an fcga, then $(\mathcal{M}_\pi)_{\pi\geqslant\pi_0}$ is one too.

\noindent\msp{page 180}Indeed, condition a) for fcg (cf. 4.7.7) follows from the same condition for $(\mathcal{F}^{\pi})$, and from (4.7.10.1), $\mathcal{M}_\pi\subset\mathcal{F}^{\pi^{\pi_0}}$. Condition b) is a general property of $\pi$-accessibility, taking into account that $(\mathcal{F}^{\pi_0})_0$ is a small generating subcategory consisting of $\pi_0$-accessible objects of $\mathcal{M}$. Condition d) is satisfied tautologically, and it remains to check condition c). But we have seen that for $\pi=\pi^{\pi_0}$, we have $\mathcal{M}_\pi=\mathcal{F}^{\pi}$ (4.7.9.1), hence condition c) for $(\mathcal{M}_\pi)$ follows from the same condition for $(\mathcal{F}^{\pi})$. Thus, we find:

\textbf{Proposition 4.7.15.} \textit{If $(\mathcal{F}^{\pi})_{\pi\geqslant\pi_0}$ is an fcga on the $\mathfrak{U}$-category $\mathcal{M}$, so is $(\mathcal{M}_\pi)_{\pi\geqslant\pi_0}$. Moreover, every useful (resp. regular) cardinal for the first filtration, is so for the second (called the} associated canonical \textit{fcg of the given} \texttt{[generalized]} \textit{cardinal filtration).}

The last assertion follows from the fact that if $\pi$ is useful for $(\mathcal{F}^{\pi})$, then we have $\mathcal{F}^{\pi}=\mathcal{M}_\pi$.

\textbf{4.7.16.} It seems to me that one can say that the canonical fcg associated with a given fcga
\msp{page 181}
is always ``better'' than the latter (if the latter is not canonical, i.e. if the two are not equal), because of the following two facts:
\begin{enumerate}[label=\alph*)]
\item precisely because it is \textit{canonical}, i.e. depends on no other choice than that of the cardinal $\pi_0$ alone, and
\item because there are a priori more useful cardinals and regular cardinals for the canonical fcg.
\end{enumerate}
So all in all, it seems to me that the \texttt{[generalized]} cardinal filtrations the only useful ones are those which are canonical.

\noindent\printedpage{62}% head: A. Grothendieck : Dérivateurs, Ch. XVIII, Édition des Universités de Montpellier II et Paris VII
% foot: Version 17 décembre 2025
Let us note that a cardinal filtration (in the sense of SGA 4 I 9, which I retain here) which is in addition \textit{accessible} (i.e. $\mathcal{F}^{\pi}\subset\mathcal{M}_\pi$ for every $\pi\geqslant\pi_0$), is nothing other than a canonical fcga, such that \textit{all} cardinals are regular (a fortiori useful) for the said filtration. I presume (without however having constructed an example) that the canonical \texttt{[generalized]} cardinal filtrations of categories of the form $\mathcal{M}=\Ind_{\pi_0}(C)$, $C$ a small \textit{Karoubian} category (so that $\pi_0$ is a useful cardinal) are not necessarily fc (not generalized, i.e. in the SGA sense), i.e. I presume that it happens that there are cardinals $\pi\geqslant\pi_0$ not regular, and even superfluous.

\noindent\msp{page 182}I must finally study what happens when we change the initial cardinal $\pi_0$ of an fcga.

\textbf{Proposition 4.7.17.} \textit{Let $(\mathcal{F}^{\pi})_{\pi\geqslant\pi_0}$ be an fcga on the $\mathfrak{U}$-category $\mathcal{M}$, and let $\pi'_0$ be a cardinal $\geqslant\pi_0$. Then $(\mathcal{F}^{\pi})_{\pi\geqslant\pi'_0}$ is an fcga on $\mathcal{M}$, which is canonical if the one we start with is canonical. For a cardinal $\pi\geqslant\pi'_0$, it is useful for $(\mathcal{F}^{\pi})_{\pi\geqslant\pi'_0}$ if and only if it is so for $(\mathcal{F}^{\pi})_{\pi\geqslant\pi_0}$. For it to be regular for $(\mathcal{F}^{\pi})_{\pi\geqslant\pi'_0}$, it} suffices \textit{that it be so for $(\mathcal{F}^{\pi})_{\pi\geqslant\pi_0}$, and this is also necessary when we assume in addition that $2^{\pi}\geqslant\pi'_1\overset{\mathrm{def}}{=}\sup(\pi'_0,\cardred\mathcal{F}^{\pi'_0})$, a fortiori when $\pi\geqslant\pi'_1$.}
\DERmark{ENV1199}


\textsc{Proof}. Conditions a) b) are trivially satisfied by the second filtration, being satisfied by the first, as is d) if it is by the first. The same is true of condition c) if $\pi=\pi^{\pi'_0}$, which indeed implies $\pi^{\pi_0}=\pi$, since in general
\[
\pi \;\leqslant\; \pi^{\pi_0} \;\leqslant\; \pi^{\pi'_0} \, .
\]
It remains to examine the cardinal bound
\[
\card I \;\leqslant\; \pi'^{\pi} \qquad \text{if } F\in\Ob\mathcal{F}^{\pi'} \ \underline{\text{and}}\ \pi'^{\pi}\geqslant\pi'_1 \, , % [sic] \longrightarrow -> \leqslant, a majoration (inequality) of cardinals is meant
\]
\DERmark{ENV1200}

where this time $\pi'_1=\sup(\pi'_0,\cardred\mathcal{F}^{\pi'_0})$. In any event, we have $\pi'_1\geqslant\pi_1$, hence
\msp{page 183}
$\pi'^{\pi}\geqslant\pi'_1$ implies $\pi'^{\pi}\geqslant\pi_1$, hence we can apply condition c) for fcg for $(\mathcal{F}^{\pi})_{\pi\geqslant\pi_0}$, to conclude that it is still verified. So in short, \textit{all} the conditions for fcg (resp. fcga) are trivially verified for the \textit{restricted} fcg of the first. Likewise, the fact that a cardinal $\pi'\geqslant\pi'_0$ is useful is trivially independent of which of the two filtrations we consider, it depends only on the datum of the cardinal $\pi'$ and of the subcategory $\mathcal{F}^{\pi'}$ of $\mathcal{M}$. There remains the cardinal bound, in the case where we are interested in the \textit{regularity} of $\pi'$ for the fcga considered. The bound required is the same, but the precondition for the validity of this bound, namely
\[
\pi'^{\pi} \;\geqslant\; \pi'_1 \qquad \text{resp.} \qquad \pi'^{\pi} \;\geqslant\; \pi_1 \, ,
\]
is not the same. Since $\pi'_1\geqslant\pi_1$, we see that the precondition is less restrictive in the case of the first filtration, than of the second, whence the ``it suffices'' in the last assertion. But when $2^{\pi}\geqslant\pi'_1$, even the more restrictive one is automatic\msp{page 184}ally verified, since $\pi'^{\pi}\geqslant 2^{\pi}\geqslant\pi'_1$, and then the regularity of $\pi$ for the second filtration implies the regularity for the first, q.e.d. % [sic] le première -> la première, gender agreement (feminine)

\noindent\printedpage{63}% head: A. Grothendieck : Dérivateurs, Ch. XVIII, Édition des Universités de Montpellier II et Paris VII
% foot: Version 17 décembre 2025
\textbf{Corollary 4.7.18.} \textit{Suppose $\pi_0$ regular for the fcga (\footnote[77]{It suffices for this that $\pi_0$ be useful, i.e. $\mathcal{M}\xleftarrow{\ \sim\ }\Ind_{\pi_0}(\mathcal{F}^{\pi_0}=\mathcal{M}_{\pi_0})$, this implies regular, cf. 4.8.5.}). If we restrict the fcga $(\mathcal{F}^{\pi})_{\pi\geqslant\pi_0}$ to the cardinals $\pi\geqslant\pi'_0$, where $\pi'_0\geqslant\pi_1$, then we find an fcga where the cardinal bound (for $\pi=\pi^{\pi_0}$) is valid without the precondition $\pi'^{\pi}\geqslant\pi'_1$; more precisely, we then have}
\[
\pi'_1 \;\leqslant\; 2^{({\pi'_0}^{\pi_0})} \qquad \text{\textit{(a fortiori $\pi'_1\leqslant 2^{(2^{\pi'_0})}$)},} \tag*{(4.7.18.1)}
\]
\textit{hence the precondition $\pi'^{\pi}\geqslant\pi'_1$ is automatically verified if $\pi\geqslant\pi_0$, $\pi=\pi^{\pi'_0}$, $\pi'>\pi$ (cf. remark 4.7.7.2).}

This is therefore equivalent to the following cardinal bound (applied to the case of $\pi=\pi'_0$).

\textbf{Lemma 4.7.19.} \textit{Suppose $\pi_0$ regular for the fcga $(\mathcal{F}^{\pi})_{\pi\geqslant\pi_0}$. Then for every cardinal $\pi\geqslant\pi_1$ ($\overset{\mathrm{def}}{=}\sup(\pi_0,\cardred\mathcal{F}^{\pi_0})$), we have}
\DERmark{ENV1201}

\[
\cardred\mathcal{F}^{\pi} \;\leqslant\; 2^{(\pi^{\pi_0})} \, , \tag*{(4.7.19.1)}
\]
\textit{and if $\pi'\geqslant\pi_1$ and $F\in\Ob\mathcal{F}^{\pi'}$, we have}
\[
\cardred\mathcal{F}^{\pi}/F \;\leqslant\; \pi'^{(\pi^{\pi_0})} \qquad \text{(\footnotemark[78])} \, . \tag*{(4.7.19.2)}
\]
\footnotetext[78]{\textbf{N.B.} I do not know whether we can replace in these bounds $\pi^{\pi_0}$ by $\pi_0$.}

\noindent\msp{page 185}\textsc{Proof of the lemma}. Since $\pi_0$ is regular, putting $C=\mathcal{F}^{\pi_0}=\mathcal{M}_{\pi_0}$, we find $\mathcal{M}\simeq\Ind_{\pi_0}(C)$, hence we can assume $\mathcal{M}=\Ind_{\pi_0}(C)$, and apply the results proved for categories of this form. Since $\mathcal{F}^{\pi}\subset\mathcal{M}_\pi$, $\mathcal{F}^{\pi'}\subset\mathcal{M}_{\pi'}$, the bounds (4.7.19.1) and (4.7.19.2) are contained in % [sic] (4.7.9.1) et (4.7.9.2) -> (4.7.19.1) et (4.7.19.2), the majorations of Lemme 4.7.19 are meant
\DERmark{ENV1202}

\[
\left\{
\begin{array}{rcll}
\cardred\mathcal{M}_\pi & \leqslant & 2^{(\pi^{\pi_0})} & \\
\cardred\mathcal{M}_\pi/F & \leqslant & \pi'^{(\pi^{\pi_0})} & \text{if } F\in\Ob\mathcal{M}_{\pi'},\ \text{with}\\
 & & & \pi'\geqslant\pi\geqslant\pi_1.
\end{array}
\right\}
\ \text{for } \pi\geqslant\pi_1 \qquad \text{(\footnotemark[79])}. \tag*{$(*)$}
\]
\footnotetext[79]{Compare with bounds $\cardred\Filt^{\pi}\leqslant 2^{\pi}$, $\cardred\Filt^{\pi}/F\leqslant\pi'^{\pi}$, if $F\in\Ob\Filt^{\pi'}$, with $\pi'\geqslant\pi\geqslant\pi_1$, bounds more pleasant than those of (4.7.19). We should all the same make sure on an example whether the nicer bounds might not be valid in the general case.}
The first of these bounds is nothing other than the second bound in 4.7.7.1. Let us note that we do indeed have
\[
2^{(\pi^{\pi_0})} \;\leqslant\; 2^{(2^{\pi})} \, ,
\]
which indeed follows from
\[
\pi^{\pi_0} \;\leqslant\; 2^{\pi} \, ,
\]
which we prove by noting that $\pi<2^{\pi}$, whence $\pi^{\pi_0}\leqslant(2^{\pi})^{\pi_0}=2^{\pi\cdot\pi_0}=2^{\pi}$ since $\pi\geqslant\pi_0$. It remains to establish the second bound
\[
\cardred\mathcal{M}_\pi/F \;\leqslant\; \pi'^{(\pi^{\pi_0})} \qquad \text{if } F\in\Ob\mathcal{M}_{\pi'},\ \text{with } \pi'\geqslant\pi\geqslant\pi_1 \, . \tag*{$(*)$}
\]
(We do not need it for 4.7.18, but it will probably be useful elsewhere.) We use $\mathcal{M}_{\pi'}\subset\Filt^{\pi'^{\pi_0}}$ if $\pi'\geqslant\pi_1$ (4.7.6.2), hence if $F\in\Ob\mathcal{M}_{\pi'}$, we have


\endgroup
% END UNIT der-ch18/der-ch18_p061-065.tex
% BEGIN UNIT der-ch18/der-ch18_p066-070.tex
\begingroup\makeatletter%
\newcommand{\printedpage}[1]{\marginpar{\footnotesize p.~#1}}%
\newcommand{\msp}[1]{\marginpar{\footnotesize [#1]}}%
\newcommand{\Ob}{\operatorname{Ob}}%
\newcommand{\Fl}{\operatorname{Fl}}%
\newcommand{\card}{\operatorname{card}}%
\newcommand{\Hom}{\operatorname{Hom}}%
\newcommand{\Filt}{\operatorname{Filt}}%
\newcommand{\Ind}{\operatorname{Ind}}%
\newcommand{\circled}[1]{\tikz[baseline=(c.base)]\node[draw,circle,inner sep=0.8pt](c){#1};}%
\makeatother\let\printedpage\derpp\let\msp\dermsp%
\let\mspin\DERmspin%


\printedpage{64}
% head: A. Grothendieck : Dérivateurs, Ch. XVIII, Édition des Universités de Montpellier II et Paris VII

\msp{page 186}

\[ \card F(x) \leqslant \pi'^{\pi_0} \qquad \forall\, x \in \Ob C , \] % [sic] \longrightarrow -> \leqslant, a cardinal inequality is expected (cf. the next display)
\DERmark{ENV1203}


thus if $G \in \Ob \mathcal{M}_\pi \subset \Ob \Filt^{\pi^{\pi_0}}$, whence
\[ \card G(x) \leqslant \pi^{\pi_0} \qquad \forall\, x \in \Ob C , \]
we find
\[ \card \Hom(G,F) \leqslant \underbrace{\underbrace{\card(\Ob C)}_{\leqslant \pi_1}\cdot(\pi'^{\pi_0})^{(\pi^{\pi_0})}}_{\pi'^{\pi_0\cdot\pi^{\pi_0}} = \pi'^{(\pi^{\pi_0})}} , \]
whence
\[ \card \Ob \mathcal{M}_\pi/F \leqslant \underbrace{\underbrace{\card(\Ob \mathcal{M}_\pi)}_{\leqslant 2^{(\pi^{\pi_0})}}\cdot\pi'^{(\pi^{\pi_0})}}_{\gamma = \pi'^{\pi^{\pi_0}}} , \]
and finally
\[ \card \Fl \mathcal{M}_\pi/F \leqslant \underbrace{\underbrace{\gamma\cdot\gamma}_{\gamma = \pi'^{(\pi^{\pi_0})}}\cdot\underbrace{\card(\Fl \mathcal{M}_\pi)}_{\leqslant 2^{(\pi^{\pi_0})}}}_{\pi'^{(\pi^{\pi_0})}} , \]
which is the announced upper bound (the same as the one proved in 4.5.12, the proof being essentially the same).

\msp{page 187}

\textbf{4.7.20.} A canonical fcg is called \textit{maximal} (or \textit{strictly canonical}) if it is not deduced from a strictly more extensive canonical fcg $(\mathcal{F}^\pi)_{\pi\geqslant\pi_{00}}$ (thus with $\pi_{00} < \pi_0$). So if $\mathcal{M}$ admits an fcga, then it even admits a strictly canonical fcg, and the latter is moreover unique. We will come back to this in the next section. % [sic] où strictement -> ou strictement, alternative is meant; [sic] n'est par -> n'est pas, negation
\DERmark{ENV1204}


\section*{4.8\quad $\mathfrak{U}$-accessible categories}

\textbf{Proposition 4.8.1.} \textit{Let $\mathcal{M}$ be a $\mathfrak{U}$-category, $\pi_0$ an infinite cardinal. Consider the following conditions:}
\begin{itemize}[label={}]
\item[\textit{a)}] \textit{$\mathcal{M}$ satisfies $L_{\pi_0}$ (i.e. is stable under small $\varinjlim$ large relative to $\pi_0$), and contains a small full subcategory $S$} strictly generating, \textit{consisting of $\pi_0$-accessible objects.}
\item[\textit{b)}] \textit{There exists a small category $C$, and an equivalence of categories}
\[ \mathcal{M} \simeq \Ind_{\pi_0}(C) . \]
\item[\textit{c)}] \textit{$\mathcal{M}$ admits an accessible generalized cardinal filtration $(\mathcal{F}^\pi)_{\pi\geqslant\pi_0}$, having $\pi_0$ as initial cardinal, and such that $\pi_0$ is a useful cardinal for the said filtration.}
\end{itemize}
% foot: Version 17 décembre 2025

\printedpage{65}
% head: A. Grothendieck : Dérivateurs, Ch. XVIII, Édition des Universités de Montpellier II et Paris VII

\msp{page 188}
\begin{itemize}[label={}]
\item[\textit{c$'$)}] \textit{As in c), with $\pi_0$ a regular cardinal (and not only useful) for the accessible generalized cardinal filtration in question, and this filtration being moreover assumed} canonical.
\end{itemize}
\textit{We then have the diagram of implications}
\[ c' \Longleftrightarrow c \Longleftrightarrow b \Longrightarrow a . \]

\textbf{4.8.1.1.} I must here make the terminology precise, as to what I mean by a (full) subcategory $C$ \textit{strictly generating} in $\mathcal{M}$. For this, I refer to SGA 4 I 7.2, (i), to the equivalent conditions b) and c) of loc. cit.:
\begin{itemize}[label={}]
\item[\circled{b}] $\forall\, F \in \Ob \mathcal{M}$, the canonical arrow
\[ \varinjlim_{i\in C/F} x_i \longrightarrow F \]
is an isomorphism.
\item[\circled{c}] The canonical functor
\[ \begin{array}{ccc}
\mathcal{M} & \longrightarrow & C^{\wedge} \\[1ex]
F & \longmapsto & (\underbrace{x}_{\in \Ob C} \longmapsto \underbrace{\Hom_{\mathcal{M}}(x,F)}_{\in \Ob (\mathrm{Ens})})
\end{array} \]
is fully faithful.
\end{itemize}
These conditions imply:
\begin{itemize}[label={}]
\item[a)] $C$ is generating in $\mathcal{M}$ by strict epimorphisms, i.e. for every $F$ in $\mathcal{M}$, the family of arrows $x \longrightarrow F$ with source in $C$, with target $F$, is a strict epimorphic family.
\end{itemize}
And when $\mathcal{M}$ is stable under fiber products, this condition a) is equivalent to the preceding ones. In loc. cit. I assert

\msp{page 189}

this equivalence without a restrictive hypothesis on $\mathcal{M}$ – but in a marginal annotation I assert that this assertion is \textit{false}, outside the hypothesis of the existence of fiber products in $\mathcal{M}$. Yet I have just checked the proof given in loc. cit. of the implication a $\Longrightarrow$ b, and it is correct. But perhaps the annotation results from a confusion between condition a), and the following condition, weaker a priori, and which resembles it like a brother:
\begin{itemize}[label={}]
\item[a$'$)] There exists a family of arrows $u_i : x_i \longrightarrow F$, with source in $C$, and with target $F$, which is a strict epimorphic family (cf. SGA 4 I 10.3).
\end{itemize}
One feels like saying that when one \textit{enlarges} a strict epimorphic family of arrows, by adding other arrows with the same target, it must remain a strict epimorphic family (``a fortiori'') – which gives the impression that a$'$) must imply a), hence be equivalent to it. But if one tries to prove this hypothetical property % [sic] agrandisse -> agrandit, verb form
% foot: Version 17 décembre 2025
\printedpage{66}
% head: A. Grothendieck : Dérivateurs, Ch. XVIII, Édition des Universités de Montpellier II et Paris VII
\msp{page 190}
of stability for strict epimorphic families, one sees that one only gets there if one assumes that the initial family is moreover a \textit{universal epimorphic family} (if fiber products exist in $\mathcal{M}$), or (subject to the preceding proviso) if for every $x' \longrightarrow F$, the sieve of $\mathcal{M}/F'$ inverse image of the sieve of $\mathcal{M}/F$ generated by the $x_i \longrightarrow F$ is an ``epimorphic sieve''. This does indeed suggest that there must be cases where condition a$'$) is satisfied, but not condition a) (equivalent to b and to c).

But we are then confronted with the very unpleasant situation: if $C$ is a strictly generating subcategory of $\mathcal{M}$, it does not thereby follow that a full subcategory $C'$ of $\mathcal{M}$ such that $C' \supset C$ is also strictly generating\,! Yet I have always used previously, without thinking twice about it, that if $C$ is strictly generating, the same is true of $C' \supset C$. So I shall have to be very attentive now, % [sic] qui si -> que si, conjunction
\msp{page 191}
and check whether this $\pm$ tacit error renders certain proofs inoperative, or even, makes certain statements false. Also, I would need to give an example where $C$ is strictly generating, but $C' \supset C$ is not.

Finally, on reflection, these feared troubles disappear. I am going to state a result in due form:

\textbf{Proposition 4.8.2.} \textit{Let $C$, $\mathcal{M}$ be two $\mathfrak{U}$-categories, $i : C \longrightarrow \mathcal{M}$ a functor, whence a functor}
\[ \left\{\begin{aligned}
&\varphi : \mathcal{M} \longrightarrow C^{\wedge}\\
&\varphi(X) = (x \longmapsto \Hom_{\mathcal{M}}(i(x),X))
\end{aligned}\right. \tag*{(4.8.2.1)} \]
\textit{(composite functor $\mathcal{M} \xrightarrow{\textit{can.}} \mathcal{M}^{\wedge} \underbrace{\xrightarrow{i^*}}_{\substack{\textit{i.e. composition}\\ \textit{with } i}} C^{\wedge}$). For every $X$ in $\mathcal{M}$, we consider the category $C/X \simeq C/\varphi(X)$, and the canonical injection} [\,?] \textit{(in $\mathcal{M}^{\vee}$)}
\[ \varinjlim_{x\in C/X} i(x) \longrightarrow X . \tag*{(4.8.2.2)} \]
\begin{itemize}[label={}]
\item[\textit{1°)}] \textit{Consider the following conditions.}
\begin{itemize}[label={}]
\item[\circled{a}] \textit{For every $X \in \Ob \mathcal{M}$, the family of arrows $i(x) \longrightarrow X$, indexed by $\Ob C/X$, is a strict epimorphic family.}

\msp{page 192}
\item[\circled{b}] \textit{For every $X \in \Ob \mathcal{M}$, the canonical arrow 4.8.2.2. is an isomorphism, i.e. makes $X$ a $\varinjlim$ of the canonical functors $C/X \longrightarrow \mathcal{M}$.}
\item[\circled{c}] \textit{The functor $\varphi$ (4.8.2.1) is fully faithful.}
\end{itemize}
\textit{We then have the implications}
\[ \mathrm{b} \Longleftrightarrow \mathrm{c} \Longrightarrow \mathrm{a} , \]
\textit{and if we assume $i$ fully faithful, then the three conditions} a, b, c \textit{are equivalent.} (\textbf{N.B.} We say that the functor $i : C \longrightarrow \mathcal{M}$ is \textit{strictly generating}, if the equivalent conditions b, c are satisfied. If $C$ is a subcategory of $\mathcal{M}$ (not
% foot: Version 17 décembre 2025
\printedpage{67}
% head: A. Grothendieck : Dérivateurs, Ch. XVIII, Édition des Universités de Montpellier II et Paris VII
necessarily full), we say that $C$ is \textit{strictly generating} if the inclusion functor is strictly generating, i.e. satisfies the equivalent conditions b), c). For this, it is therefore necessary that condition a) be satisfied, i.e. that for every $X$ in $\mathcal{M}$, the family of arrows $x \longrightarrow X$ with source in $C$, with target $X$, be a strict epimorphic family, and this condition is also sufficient when $C$ is a full subcategory of $\mathcal{M}$.)
\item[\textit{2°)}] (\footnote[80]{\texttt{[2° is on page 193.]}}). \textit{Suppose that $i$ is strictly generating} (\footnote[81]{It is even enough that it satisfy the weaker condition a) of 1°).}). \textit{Then for every object $X$ of $\mathcal{M}$, the family of arrows $i(x) \longrightarrow X$, indexed by $\Ob C/X$, is not only a strict epimorphic family (condition a) in 1°), but moreover, every family of arrows with target $X$, including the preceding family, is also a strict epimorphic family.}
\item[\textit{3°)}] \textit{Consider a commutative diagram of functors of $\mathfrak{U}$-categories:}

\msp{page 193}
\[ \begin{tikzcd}[column sep=small]
C \arrow[rr, "\alpha"] \arrow[dr, "i"'] & & C' \arrow[dl, "i'"] \\
& E &
\end{tikzcd} \]
\textit{If the functor $i$ is strictly generating} (\footnote[82]{In fact, it is enough that it satisfy the weaker condition a) of 1°).}), \textit{and if moreover $i'$ is fully faithful, then $i'$ is strictly generating.}
\DERmark{ENV1205}

\end{itemize}

\textbf{Corollary 4.8.3.} \textit{Let $\mathcal{M}$ be a $\mathfrak{U}$-category, $C$ a strictly generating subcategory, $C'$ a} full \textit{subcategory of $\mathcal{M}$ which contains $C$, then $C'$ is strictly generating.}

This is indeed a special case of 4.8.2 3°.

\msp{page 194}

\textsc{Proof of 4.8.2}. The proof of 1°) is that of SGA 4 I 7.2 (i). (Where one assumes that $C$ is a full subcategory of $\mathcal{M}$, but the proof given in fact proves the statement given here.) Let us prove 2°). Let us therefore add to the family of
\[ i(x) \longrightarrow X \qquad (x \in \Ob C/X) \]
a family
\[ X_\alpha \xrightarrow{p_\alpha} X \qquad (\alpha \in A) \]
of new arrows, and prove that the total family thus obtained is still a strict epimorphic family. We must therefore prove that if $Y$ is an object of $\mathcal{M}$, and if we have a family of arrows $u_x : i(x) \longrightarrow Y$, $u_\alpha : X_\alpha \longrightarrow Y$, satisfying the usual compatibility condition for pairs of arrows:
\[ (v_x,v_y) \qquad\qquad \begin{tikzcd}
& i(x) \arrow[dr, dotted, "p_x"] & \\
Z \arrow[ur, "v_x"] \arrow[dr, "v_y"'] & & X\ , \\
& i(y) \arrow[ur, dotted, "p_y"'] &
\end{tikzcd} \]
or
% foot: Version 17 décembre 2025
\printedpage{68}
% head: A. Grothendieck : Dérivateurs, Ch. XVIII, Édition des Universités de Montpellier II et Paris VII
\[ (v_x,v_\alpha) \qquad\qquad \begin{tikzcd}
& i(x) \arrow[dr, dotted, "p_x"] & \\
Z \arrow[ur, "v_x"] \arrow[dr, "v_\alpha"'] & & X\ , \\
& X_\alpha \arrow[ur, dotted, "p_\alpha"'] &
\end{tikzcd} \]
\msp{page 195}
or
\[ (v_\alpha,v_\beta) \qquad\qquad \begin{tikzcd}
& X_\alpha \arrow[dr, dotted] & \\
Z \arrow[ur, "v_\alpha"] \arrow[dr, "v_\beta"'] & & X\ , \\
& X_\beta \arrow[ur, dotted] &
\end{tikzcd} \]
namely the relations
\[ \left\{\begin{aligned}
u_xv_x &= u_yv_y\\
u_xv_x &= u_\alpha v_\alpha\\
u_\alpha v_\alpha &= u_\beta v_\beta ,
\end{aligned}\right. \tag*{($*$)} \]
then there exists an arrow
\[ u : X \longrightarrow Y \]
(necessarily unique, by property a) of strict epimorphicity of the initial family indexed by $C/X$) such that we have
\DERmark{ENV1206}

\[ \left\{\begin{aligned}
up_x &= u_x \quad \text{for every } x \in \Ob C/X\\ % [sic] v_x -> u_x, the arrow u p_x goes to Y (cf. up_\alpha = u_\alpha below)
up_\alpha &= u_\alpha \quad \text{for every } \alpha \in A. % [sic] v_\alpha -> u_\alpha, cf. the relation up_\alpha = u_\alpha below
\end{aligned}\right. \tag*{($*$)} \]
\DERmark{ENV1207}\DERmark{ENV1208}

Now by property a), there exists a unique $u : X \longrightarrow Y$, in such a way that the first of the relations ($*$) is satisfied, for every $x \in \Ob C/X$. So it is only a matter of proving that for this $u$, we have for every $\alpha$ the relation

\msp{page 196}
\[ up_\alpha = u_\alpha . \]
Now, applying property a) to $X_\alpha$ instead of $X$, we see that the family of arrows % [sic] veut -> voit, sense of the sentence
\DERmark{ENV1209}

\[ i(y) \xrightarrow{q_y} X_\alpha \]
(indexed by $C/X_\alpha$) is a strict epimorphic family, a fortiori epimorphic. Thus to prove $up_\alpha = u_\alpha$, it is enough to prove that we have
\[ (up_\alpha)q_y = u_\alpha q_y \tag*{($**$)} \]
for every $y$ in $C/X_\alpha$. But this is a consequence of compatibility ($*$) (second formula), applied to the following case,
% foot: Version 17 décembre 2025


\endgroup
% END UNIT der-ch18/der-ch18_p066-070.tex
% BEGIN UNIT der-ch18/der-ch18_p071-075.tex
\begingroup\makeatletter%
\newcommand{\printedpage}[1]{\marginpar{\footnotesize p.~#1}}%
\newcommand{\msp}[1]{\marginpar{\footnotesize [#1]}}%
\newcommand{\Ob}{\operatorname{Ob}}%
\newcommand{\Ind}{\operatorname{Ind}}%
\newcommand{\Filt}{\operatorname{Filt}}%
\newcommand{\card}{\operatorname{card}}%
\newcommand{\End}{\operatorname{End}}%
\newcommand{\Fl}{\operatorname{Fl}}%
\newcommand{\Kar}{\operatorname{Kar}}%
\makeatother\let\printedpage\derpp\let\msp\dermsp%
\let\mspin\DERmspin%


\printedpage{69}
% head: A. Grothendieck : Dérivateurs, Ch. XVIII, Édition des Universités de Montpellier II et Paris VII

\[
\begin{tikzcd}[column sep=small]
 & i(y) \arrow[dr, dotted, "p_\alpha q_y=p_x"] & \\
Z=i(y) \arrow[ur, "\mathrm{id}"] \arrow[dr, "q_y"'] & & X \ , \\
 & X_\alpha \arrow[ur, dotted, "p_\alpha"'] &
\end{tikzcd}
\]
we therefore have, setting $x=(y,p_\alpha q_y)\in\Ob C/X$,
\[
\underbrace{u_x\circ\mathrm{id}_{i(y)}}_{u_x} \ = \ u_\alpha q_y \ ,
\]
so to prove $(**)$, it remains to verify
\[
u_x \ = \ \underbrace{(up_\alpha)q_y}_{u(p_\alpha q_y)\,=\,up_x} \ ,
\]
\msp{page 197}
which is indeed verified by definition of $u_x$. This finishes proving $2^\circ$.

Let us prove $3^\circ$. Since $i'$ is fully faithful, by $1^\circ$, to prove that $i'$ is strictly generating, it suffices to establish that for every object $X$ of $\mathcal{M}$, the family of arrows $i'(x')\longrightarrow X$, indexed by $\Ob C'/X$, is a strict epimorphic family. Now this family contains the family of $\underbrace{i'(\alpha(x))}_{=\,i(x)}\longrightarrow X$, indexed by $\Ob C/X$. Thus it follows from the hypothesis on $i$, and from $2^\circ$, that this family is a strict epimorphic family. We've won! Whew \dots

\textsc{Proof of 4.8.1.} We have
\[
\mathrm{c}' \ \Longrightarrow \ \mathrm{c} \ \Longrightarrow \ \mathrm{b} \ \Longrightarrow \ \mathrm{a} \, ,
\]
the first implication tautological, the second resulting from the fact that $\pi_0$ a useful cardinal for the fcga under consideration, implies that
\DERmark{ENV1210}

\[
\mathcal{M} \ \xleftarrow{\ \sim\ } \ \Ind_{\pi_0}(\mathcal{F}^{\pi_0}) \ ,
\]
$\mathcal{F}^{\pi_0}$ equivalent to a small category, whence b).

\msp{page 198}
Finally b) $\Longrightarrow$ a) follows from the fact that $\Ind_{\pi_0}(C)$ is stable under $\varinjlim$ large relative to $\pi_0$ (4.4.15), and that $C$ is a strictly generating subcategory of $\Ind_{\pi_0}(C)$, made up of $\pi_0$-accessible objects. We need to spell out the

\textbf{Lemma 4.8.4.} \textit{Let $C$ be a $\mathfrak{U}$-category, $\pi_0$ an infinite cardinal. Then $C$ is strictly generating in $\mathcal{M}=\Ind C$ and in $\mathcal{M}=\Ind_{\pi_0}C$.}

\msp{page 199}
Indeed, the canonical functor $\mathcal{M}\longrightarrow C^\wedge$ is fully faithful, and we apply criterion c) of 4.8.2.

To finish proving 4.8.1, it therefore remains to establish the implication
\[
\mathrm{a} \ \Longrightarrow \ \mathrm{c}' \, .
\]

\printedpage{70}
% head: A. Grothendieck : Dérivateurs, Ch. XVIII, Édition des Universités de Montpellier II et Paris VII

Let us suppose then that $\mathcal{M}$ satisfies $\mathrm{L}_{\pi_0}$, and admits a small strictly generating subcategory $S$, made up of $\pi_0$-accessible objects, hence a canonical functor
\[
\Ind_{\pi_0}S \ \longrightarrow \ \mathcal{M} \ ,
\]
extending the inclusion $S\hookrightarrow\mathcal{M}$, so as to commute with the $\varinjlim$ large relative to $\pi_0$, and because $S\subset\mathcal{M}_{\pi_0}$, this functor is \textit{fully faithful.} But how to prove that it is essentially surjective? This will not be true, moreover, without a further hypothesis on $S$. The natural idea is to saturate $S$ under the $\varinjlim$
\msp{page 200}
of cardinality $\leqslant\pi_0$ representable in $\mathcal{M}$, and in particular under direct factors, but even then (and even admitting that $\mathcal{M}_{\pi_0}$ itself is equivalent to a small category, and taking $C=\mathcal{M}_{\pi_0}$), it is nevertheless not clear that the fully faithful functor $\Ind_{\pi_0}C\longrightarrow\mathcal{M}$ is essentially surjective, which would also amount to saying that for every $F$ in $\mathcal{M}$, $C/F$ is large relative to $\pi_0$. It is in fact not even clear that $C/F$ is filtering.

So in the end, I doubt that we have a) $\Longrightarrow$ c$'$) or just a) $\Longrightarrow$ b), the situation is visibly not well understood.

But let us prove all the same
\[
\mathrm{b} \ \Longrightarrow \ \mathrm{c}' \, .
\]
For this, we may assume $\mathcal{M}=\Ind_{\pi_0}C$, $C$ small. We then know that the filtration of $\mathcal{M}$ by the $(\mathcal{M}_\pi)_{\pi\geqslant\pi_0}$ is an fcga (prop. 4.7.7 and definition 4.7.8), and even an fcg that is \textit{canonical.} Moreover, I say that the cardinal $\pi_0$ is \textit{useful} and even regular for this filtration (which will finish proving c$'$). We have a slightly more precise result:

\msp{page 201}
\textbf{Corollary 4.8.5.} \textit{Let $C$ be a category equivalent to a small category, $\pi_0$ an infinite cardinal, $\mathcal{M}=\Ind_{\pi_0}(C)\subset C^\wedge$. Then:}
\begin{enumerate}
\item[\textit{1$^\circ$)}] $\pi_0$ \textit{is a cardinal} useful \textit{and even regular for the canonical fcg} $(\mathcal{M}_\pi)_{\pi\geqslant\pi_0}$ \textit{of} $\mathcal{M}$.
\item[\textit{2$^\circ$)}] $\mathcal{M}_{\pi_0}=C^\natural$ \textit{(full subcategory of $\mathcal{M}=\Ind_{\pi_0}C$ made up of the direct factors in $\mathcal{M}$ of objects of $C$), and $\mathcal{M}_{\pi_0}$ is a Karoubi envelope of $C$.}
\end{enumerate}
(\footnote{\textbf{Corollary 4.8.5.1.} \textit{If $C$ is Karoubian, then $\mathcal{M}_{\pi_0}\subset C^\wedge$ is the category made up of the functors} representable \textit{in} $C\overset{\sim}{\hookrightarrow}\mathcal{M}_{\pi_0}$.}, \footnote{\textbf{4.8.5.2.} Spell this out too: Let $u : C\longrightarrow C'$ be a functor, with $C$, $C'$ equivalent to small categories. Equivalent conditions:
\begin{enumerate}
\item[a)] $u$ induces $\Kar C\xrightarrow{\ \sim\ }\Kar C'$.
\item[a$'$)] $u$ fully faithful, and every object of $C'$ is isomorphic to a direct factor of an object of the essential image.
\item[b)] $\Ind(u) : \Ind(C)\longrightarrow\Ind(C')$ is an equivalence.
\item[c)] $\Ind_{\pi_0}(u) : \Ind_{\pi_0}(C)\longrightarrow\Ind_{\pi_0}(C')$ is an equivalence ($\pi_0$ a given infinite cardinal).
\item[d)] $u_! : C^\wedge\longrightarrow C'^\wedge$ is an equivalence.
\item[d$'$)] $u_* : C'^\wedge\longrightarrow C^\wedge$ is an equivalence \texttt{[rather} $u_* : C^\wedge\longrightarrow C'^\wedge$\texttt{]}.
\item[d$''$)] $u^* : C'^\wedge\longrightarrow C^\wedge$ is an equivalence.
\end{enumerate}}).

\printedpage{71}
% head: A. Grothendieck : Dérivateurs, Ch. XVIII, Édition des Universités de Montpellier II et Paris VII

\textsc{Proof.}

$1^\circ$) We already know that
\[
\Ind_{\pi_0}\mathcal{M}_{\pi_0} \ \longrightarrow \ \mathcal{M} \tag*{$(*)$}
\]
is fully faithful, and to see that it is essentially surjective, we note that we have
\[
\mathcal{M} \ \simeq \ \Ind_{\pi_0}C \ \hookrightarrow \ \Ind_{\pi_0}\mathcal{M}_{\pi_0} \ \longrightarrow \ \mathcal{M} \ ,
\]
the composite being isomorphic to the identity functor, is essentially surjective, hence so is $(*)$. This therefore proves that $\pi_0$ is a useful cardinal for the filtration under consideration, let us prove that this cardinal is even regular, hence that for
\[
F \ \in \ \Ob\mathcal{M}_\pi \ , \qquad \text{with } \pi>\pi_0,\ \pi^{\pi_0}\geqslant\pi_1
\]
we have
\[
F \ = \ \varinjlim_I x_i \ , \qquad x_i\in\Ob\mathcal{M}_{\pi_0},\ I \text{ large relative to } \pi_0,\ \card I\leqslant\pi^{\pi_0}.
\]
\msp{page 202}
For this, we use (4.7.6.2), page 171,
\[
\mathcal{M}_\pi \ \subset \ \Filt^{\pi^{\pi_0}} \qquad (\text{if } \pi\geqslant\pi_0,\ \pi^{\pi_0}\geqslant\pi_1),
\]
and the definition of $\Filt^{\pi'}$ (for $\pi'\geqslant\pi_0$) shows us that this is the conclusion sought.

Let us then prove $2^\circ$). Since $\mathcal{M}_{\pi_0}$ is stable under the finite direct $\varinjlim$ which are representable in $\mathcal{M}$ (SGA 4 I 9.9), it is stable under direct factors, and since it contains $C$, we therefore have $C^\natural\subset\mathcal{M}_{\pi_0}$. The reverse inclusion $\mathcal{M}_{\pi_0}\subset C^\natural$ is seen in the standard way, by writing an object $F$ of $\mathcal{M}_{\pi_0}$ as $\varinjlim\limits_I x_i$, the $x_i$ in $C$, $I$ large relative to $\pi_0$, then $F$ is a direct factor of one of the $x_i$.

The fact that $C^\natural$ is a \textit{Karoubi envelope} of $C$ in $\mathcal{M}\subset C^\wedge$ comes from the fact that $\mathcal{M}$ is stable under direct factors, and that $C$ is a full subcategory of $\mathcal{M}$, cf. SGA 4 I (\,?\,\dots).

\textbf{4.8.6. Open questions.} We have thus proved 4.8.1 revised and corrected. There remains the question:
\begin{enumerate}
\item[1)] Give an example where 4.8.1 a) holds, but not b) etc. (\footnote{It suffices to find a nonregular cardinal, for the canonical fcg of an $\mathcal{M}=\Ind_{\pi_0}C$.}).
\item[2)] Bring out serviceable additional conditions which, joined to a), imply b), hence c, c$'$.
\end{enumerate}

\msp{page 203}
Let us also note for the record the following questions.
\begin{enumerate}
\item[3)] Let $\mathcal{M}=\Ind_{\pi_0}(C)$, and $\pi>\pi_0$. Then is $\Filt^\pi$ stable under direct factors? Is this true at least if $\pi=\pi^{\pi_0}$, or what amounts to the same thing (cf. 4.7.3), do we have in this case $\Filt^\pi=\Filt'^{\pi}$? If so, it would be unnecessary in the description 4.7.7 of cardinal filtrations, in the bound on cardinals (4.7.7.1), $\card I\leqslant\pi'^{\pi}$, to put in the restrictive hypothesis $\pi'^{\pi}\geqslant\pi_1$.
\end{enumerate}

To have an affirmative answer to this question, for a given $\pi$, it would suffice (and probably be necessary) to have an affirmative answer to this one: Let $\mathcal{D}$, $\mathcal{D}_0$ be two categories large relative to $\pi_0$, $\mathcal{D}_0$ a direct factor of $\mathcal{D}$, $\mathcal{D}$ admitting a cofinal category $I$ large relative to $\pi_0$, of cardinality $\leqslant\pi$. Then is the same true of $\mathcal{D}_0$? If $p$ is a retraction of $\mathcal{D}$ onto $\mathcal{D}_0$, the difficulty is that I do not know whether $p$ is a cofinal functor. In fact, I see that
\printedpage{72}
% head: A. Grothendieck : Dérivateurs, Ch. XVIII, Édition des Universités de Montpellier II et Paris VII
\msp{page 204}
this is not always the case, even if $\mathcal{D}$ is finite. Let us take e.g. for $\mathcal{D}_0$ a category having final object $e$, and an arrow $x\longrightarrow e$ which is not an isomorphism, and admitting a section (e.g. $\mathcal{D}_0$ the full subcategory of $\boldsymbol{\Delta}$ formed by $\boldsymbol{\Delta}_0$, $\boldsymbol{\Delta}_1$: \begin{tikzcd}[cramped] \boldsymbol{\Delta}_0 \arrow[r, shift left, "\partial_0"] \arrow[r, shift right, "\partial_1"'] & \boldsymbol{\Delta}_1 \arrow[l, bend left=60, "\kappa"] \end{tikzcd} ). Let $\mathcal{D}$ be obtained from $\mathcal{D}_0$ by adding a final object $\tilde{e}$. Then $\mathcal{D}_0$, $\mathcal{D}$ are filtering, and we get a retraction of $\mathcal{D}$ onto $\mathcal{D}_0$, sending $\tilde{e}$ to $x$. This functor is not cofinal, for a cofinal functor must send $\tilde{e}$ to $e$.

More generally, if $\mathcal{D}_0$ is a category such that there exists an arrow
\[
\alpha : e_{\mathcal{D}_0^\wedge} \ \longrightarrow \ x_0 \tag*{$(*)$}
\]
in $\mathcal{D}_0^\wedge$, with $x_0\in\Ob\mathcal{D}_0$, then we see that $\mathcal{D}_0$ is large relative to $\pi_0$ whatever the infinite cardinal $\pi_0$, and that on the other hand, again denoting by $\mathcal{D}$ the category obtained from $\mathcal{D}_0$ by adding a final object to it, there is a retraction of $\mathcal{D}$ onto $\mathcal{D}_0$ which sends $e_{\mathcal{D}}$ to $x_0$. This retraction is cofinal only if $x_0$ is a final object of $\mathcal{D}_0$, i.e. if and only if $(*)$ is an isomorphism. Let us suppose this is not so. But there is a cofinal functor from the subcategory $I_0$ of $\mathcal{D}_0$, such that $\Ob I_0=\{x_0\}$, $\Fl I_0=\End_{I_0}(x_0)=\{p,\mathrm{id}_{x_0}\}$, where $p$ is the composite $x_0\xrightarrow{\text{can.}}e_{\mathcal{D}_0^\wedge}\xrightarrow{\alpha}x_0$. So no counterexample this way!

\msp{page 205}
\begin{enumerate}
\item[4)] Let us consider $\mathcal{M}=\Ind_{\pi_0}(C)$, $C$ small, and let us consider the canonical filtration $(\mathcal{M}_\pi)_{\pi\geqslant\pi_0}$. Give examples where there are non-useful cardinals, or better, where for every cardinal $\pi$, exists a cardinal $\pi'>\pi$ which is non-useful. (This will give the example asked for in $1^\circ$.) On the other hand, give examples of useful cardinals which are not regular -- if possible, as before, such cardinals as large as we want.
\end{enumerate}

\textbf{Definition 4.8.7.} Let $\mathcal{M}$ be a $\mathfrak{U}$-category, $\pi_0$ an infinite cardinal. We say that $\mathcal{M}$ is $\pi_0$-\textit{accessible} if $\mathcal{M}$ is equivalent to a category $\Ind_{\pi_0}(C)$, $C$ a small category, i.e. if it satisfies the equivalent conditions b, c, c$'$ of 4.8.1 (implying condition a), probably strictly weaker). We say that $\mathcal{M}$ is accessible if there exists a cardinal such that $\mathcal{M}$ is $\pi_0$-accessible (\textbf{N.B.} if $\pi_0$ is such a cardinal, then every cardinal $\pi>\pi_0$ such that $\pi^{\pi_0}=\pi$, in particular the cardinals $\pi'^{\pi_0}$, with $\pi'\geqslant 2$, have the same property), or what
\msp{page 206}
amounts to the same thing by 4.8.1, if $\mathcal{M}$ admits an fcga, or again a canonical fcg.

\textbf{Comment 4.8.8.} This definition is in contradiction with the terminology introduced in SGA 4 I 9.4 (p. 114); we shall henceforth designate by the names ``weakly $\pi_0$-accessible category'', ``weakly accessible category'' the notions introduced in 9.4 -- i.e. the categories admitting a full generating subcategory (not even strictly generating) whose objects are $\pi_0$-accessible (resp. accessible).

Let us note that if $\mathcal{M}$ is $\pi_0$-accessible (resp. accessible), the condition of 4.8.1 tells us that $\mathcal{M}$ is weakly $\pi_0$-accessible (resp. weakly accessible), and even that there exists a small full \textit{strictly} generating subcategory, which is made up of $\pi_0$-accessible objects (resp. accessible). I very much doubt that this last necessary condition is also sufficient for
\printedpage{73}
% head: A. Grothendieck : Dérivateurs, Ch. XVIII, Édition des Universités de Montpellier II et Paris VII
$\mathcal{M}$ to be $\pi_0$-accessible (resp. accessible). But by 4.6.3.3 (p. 150), this is so if we assume $\mathcal{M}$ stable under small $\varinjlim$ , or what amounts to the same thing, assuming $\mathrm{L}_{\pi_0}$, stable under direct limits of cardinality $\leqslant\pi_0$.
\DERmark{ENV1211}


\msp{page 207}
\textbf{Proposition 4.8.9.} \textit{Let $\mathcal{M}$ be an accessible $\mathfrak{U}$-category (4.8.7), $\pi_0$ a cardinal such that $\mathcal{M}$ satisfies $L_{\pi_0}$. Equivalent conditions:}
\begin{enumerate}
\item[\textit{a)}] \textit{$\mathcal{M}$ is a paratopos category (cf. 4.6.3, $3^\circ$, p. 150).}
\item[\textit{b)}] \textit{$\mathcal{M}$ is stable under small direct limits.}
\item[\textit{b$'$)}] \textit{$\mathcal{M}$ is stable under direct limits of cardinality $\leqslant\pi_0$.}
\item[\textit{b$''$)}] \textit{(If $\mathcal{M}$ is equivalent to a categories} $\Ind_{\pi_0}C$\textit{, with $C$ a small category} that is Karoubian.\textit{) The category $C$ is stable (in itself) under direct limits of cardinality $\leqslant\pi_0$.} % [sic] une catégories -> une catégorie, accord du nombre
\item[\textit{c)}] \textit{$\mathcal{M}$ is stable under small inverse limits.}
\item[\textit{c$'$)}] \textit{(Under the hypothesis $\mathcal{M}\simeq\Ind_{\pi_0}(C)$, and when we give ourselves a cardinal $\pi$ such that $\pi=\pi^{\pi_0}\geqslant\pi_1=\sup(\pi_0,\card\textup{red}\,C)$.) $\mathcal{M}_\pi$ is stable under finite $\varprojlim$, and $\mathcal{M}$ is stable under small infinite products (or again, under filtering $\varprojlim$).}
\item[\textit{c$''$)}] \textit{There exists an infinite cardinal $\pi$, and a small category $\mathcal{D}$ stable under finite $\varprojlim$, such that $\mathcal{M}\simeq\Ind_\pi(\mathcal{D})$, moreover $\mathcal{M}$ stable under infinite products (or again under filtering $\varprojlim$).}
\end{enumerate}
(\footnote{\textbf{N.B.} The equivalence of c$'$, c$''$ with the other conditions is not established (here, but only in 4.11 below.)}, \footnote{\textbf{N.B.} The criteria c$'$, c$''$ seem to me more technical and less useful \texttt{[c$'$, c$''$ above are put in parentheses]}.}).

\textsc{Proof.} The following implications are tautological, or already $\pm$ known:
\[
\begin{tikzcd}[column sep=large, row sep=large]
\mathrm{a} \arrow[r, Rightarrow] & \mathrm{b} \arrow[r, Rightarrow, "\text{cf. 4.4.14}"] \arrow[dl, Rightarrow] & \mathrm{b}' \arrow[r, Rightarrow, "\bullet"] & \mathrm{b}'' \\
\mathrm{c}\ ( \arrow[r, Rightarrow, "\bullet"] & \mathrm{c}' \arrow[r, Rightarrow] & \mathrm{c}'')\ , &
\end{tikzcd}
\]
\msp{page 208}
where only the dot-marked implications b$'\Longrightarrow$ b$''$ and c $\Longrightarrow$ c$'$ call for a proof. For b$'\Longrightarrow$ b$''$, we note that the hypothesis $C$ Karoubian implies that $C\simeq\mathcal{M}_{\pi_0}$ (4.8.5, $2^\circ$), now we know that $\mathcal{M}_{\pi_0}$ is stable in $\mathcal{M}$ under direct limits of cardinality $\leqslant\pi_0$ (and these are representable in $\mathcal{M}$ by b$'$), hence $\mathcal{M}_{\pi_0}$ admits $\varinjlim$ of the type under consideration, hence $C$ too.

c $\Longrightarrow$ c$'$. Let us note that if $\mathcal{M}=\Ind_{\pi_0}(C)$, then the $\varprojlim$ representable in $\mathcal{M}$ are computed in $C^\wedge$. It follows that if $\pi\geqslant\pi_1$, then $\Filt^\pi$ is stable under finite $\varprojlim$ ($\Filt^\pi$ then being given by condition a$''$ of 4.7.1), on the other hand, if in addition $\pi^{\pi_0}=\pi$, $\pi$ is a useful cardinal, hence $\mathcal{M}\simeq\Ind_\pi(\mathcal{D})$, where $\mathcal{D}=\Filt^\pi=\mathcal{M}_\pi$. It now suffices to see that $\Ind_\pi(\mathcal{D})$ is then stable under finite $\varprojlim$, to conclude that the same holds for $\mathcal{M}$, hence that $\mathcal{M}$ is stable under arbitrary inverse limits, if it is under small products or under small filtering $\varprojlim$, whence c$'$. It therefore remains to prove the following lemma. % [sic] petites produites -> petits produits, accord


\endgroup
% END UNIT der-ch18/der-ch18_p071-075.tex
% BEGIN UNIT der-ch18/der-ch18_p076-080.tex
\begingroup\makeatletter%
\newcommand{\printedpage}[1]{\marginpar{\footnotesize p.~#1}}%
\newcommand{\msp}[1]{\marginpar{\footnotesize [#1]}}%
\newcommand{\Ind}{\operatorname{Ind}}%
\newcommand{\Filt}{\operatorname{Filt}}%
\newcommand{\card}{\operatorname{card}}%
\newcommand{\red}{\operatorname{\text{red}}}%
\newcommand{\Ob}{\operatorname{Ob}}%
\newcommand{\Hom}{\operatorname{Hom}}%
\newcommand{\uHom}{\underline{\operatorname{Hom}}}%
\newcommand{\End}{\operatorname{End}}%
\newcommand{\Ens}{\mathrm{Ens}}%
\newcommand{\cM}{\mathcal{M}}%
\newcommand{\cD}{\mathcal{D}}%
\newcommand{\defeq}{\overset{\mathrm{def}}{=}}%
\makeatother\let\printedpage\derpp\let\msp\dermsp%
\let\mspin\DERmspin%

\DERmark{ENV1212}


\printedpage{74}
% head: A. Grothendieck : Dérivateurs, Ch. XVIII, Édition des Universités de Montpellier II et Paris VII

\textbf{Lemma 4.8.10.} \textit{Let $C$ be a small category, $\pi_0$ an infinite cardinal, $\pi_1 = \sup(\pi_0, \card\red C)$. Suppose $\Ind_{\pi_0}C$ is stable under finite $\varprojlim$ (resp. and of a certain type). Then for every $\pi \geqslant \pi_1$, $\Filt^{\pi}$ is stable under the said $\varprojlim$ . (Thus also $\cM_{\pi} = \Filt^{\pi}$, if $\pi = \pi^{\pi_0}$.) More} % [sic] \Ind_{\pi_0} stable -> \Ind_{\pi_0}C stable, the argument C appears to be omitted (cf. the following sentence)
\DERmark{ENV1213}

\msp{page 209}
\textit{generally, let $J$ be a small category such that $\Ind_{\pi_0}C$ is stable under $\varprojlim$ of type $J$, and let $c = \card\red J$, suppose that $\pi^{\max(1,c)} = \pi$ (e.g. $c$ finite). Then $\Filt^{\pi}$ is stable in $\cM$ under $\varprojlim$ of type $J$. (Thus the same holds for $\cM_{\pi} = \Filt^{\pi}$, if in addition $\pi^{\pi_0} = \pi$.)}
\DERmark{ENV1214}


We use criterion a$''$ of 4.7.1, namely $\card F(x) \leqslant \pi$ for every $x \in C$, for membership in $\Filt^{\pi}$. We may suppose $J$ reduced, hence $\card\Ob J \leqslant c$. Since the inclusion functor $\Ind_{\pi_0}(C) \longrightarrow C^{\wedge}$ commutes with arbitrary $\varprojlim$, the computation of $\varprojlim\limits_{j\in J} x_j$ in $C^{\wedge}$ fiber by fiber shows that in $C^{\wedge}$, we have
\[
X \ \defeq\ \varprojlim_{j\in J} x_j \ \hookrightarrow\ \underbrace{\prod_{j\in\Ob J} x_j}_{\text{product in } C^{\wedge}} \quad ,
\] % [sic] \varinjlim_{j\in J} x_j -> \varprojlim_{j\in J} x_j, the text concerns the projective limit (cf. the preceding sentence and (4.8.10.1))
\DERmark{ENV1215}\DERmark{ENV1216}

whence
\[
\card X(x) \ \leqslant\ \prod_{j\in\Ob J} \underbrace{\card x_j(x)}_{\leqslant\pi} \ \leqslant\ \pi^{\max(1,c)}\ ,
\]
\DERmark{ENV1217}

whence (with no other hypothesis on $\pi$ than $\pi \geqslant \pi_1$)
\[
\varprojlim_{j\in J} x_j \ \in\ \Ob\Filt^{\pi^{\max(1,c)}}\ , \tag*{(4.8.10.1)}
\]
\DERmark{ENV1218}

whence the desired conclusion if $\pi^{\max(1,c)} = \pi$, q.e.d.
\DERmark{ENV1219}


\msp{page 210}

We are going to prove the additional implications
\[
\text{b}'' \ \Longrightarrow\ \text{b}\ , \qquad \text{c} \ \Longrightarrow\ \text{b}\ , \qquad \text{c}'' \ \overset{?}{\Longrightarrow}\ \text{c}\ ,
\]
which completes the proof.

The implication b$''$ $\Longrightarrow$ a is contained in 4.6.3.

For the implication c $\Longrightarrow$ b, we note that the stability of $\cM$ under small $\varinjlim$ is equivalent to the assertion that every small inverse limit of representable functors $\cM \xrightarrow{\ \xi_i\ } \Ens$ is representable. Now the objects of $\cM$ being accessible ($\cM$ being accessible), the representable functors $\xi_i$ are accessible, hence so is their $\varprojlim$ (SGA 4 I 9.7). On the other hand, the functor $\xi$ is left exact, as a $\varprojlim$ of representable hence left exact functors. Thus the fact that $\xi$ is representable follows from the % [sic] exactes -> exacts, agreement with « foncteurs »

\textbf{Lemma 4.8.11.} \textit{Let $\cM$ be an accessible category, stable under finite $\varprojlim$, $\cM \longrightarrow$} Ens \textit{a functor. For $\xi$ to be prorepresentable, it is necessary and sufficient that it be left exact and accessible. Thus (SGA 4 I 8.12.8.1) if $\cM$ is stable under small $\varprojlim$ , $\xi$ is representable if and only if it commutes with the said $\varprojlim$ , and is accessible.}
\DERmark{ENV1220}


% foot: Version 17 décembre 2025

\printedpage{75}
% head: A. Grothendieck : Dérivateurs, Ch. XVIII, Édition des Universités de Montpellier II et Paris VII

\msp{page 211}

The ``it is necessary'' is included as a reminder, we have just seen it as a consequence of the fact that every object of $\cM$ is accessible. The ``it is sufficient'' is essentially SGA 4 I 8.13.2.b), except that in loc. cit. one supposes $\cM$ equipped with an fc that is \textit{strict.} But the proof given in loc. cit. also shows for an fcga, choosing in loc. cit. (p. 137) $\pi$ such that $\xi$ is $\pi$-accessible, and moreover $\pi$ is useful for the generalized cardinal filtration under consideration.

It remains to prove c$''$ $\overset{?}{\Longrightarrow}$ c, subject to this really being true. This would be a consequence of the following lemma (where $\cD$, $\pi$ are replaced by $C$, $\pi_0$).

\textbf{Lemma 4.8.12} (?) (\footnote[88]{Lemma not proved here -- cf. proof 4.11.7 (unfinished).}). \textit{Let $C$ be a small category, $\pi_0$ an infinite cardinal. If $C$ is stable under finite $\varprojlim$, then $\Ind_{\pi_0}C$ as well.} % [sic] C_0 -> C, the statement then speaks only of C
\DERmark{ENV1221}


The analogous assertion for $\Ind(C)$, instead of $\Ind_{\pi}C$, is established in SGA 4 I 8.9.1 c). But it is not clear without more that $\Ind_{\pi}C$ is stable in $\Ind(C)$ under finite $\varprojlim$, in the case under consideration. More precisely, if $F = \varprojlim\limits_{j\in J} F_j$, the $F_j$ in $\Ind_{\pi}C$,
\msp{page 212}
$F$ in $\Ind(C) \subset C^{\wedge}$, we know that $C/F$ is filtering and we have to see that it is filtering large relative to $\pi_0$. We are reduced to the three cases where $J = \emptyset$, $J = e \amalg e$, and $J = \vcenter{\hbox{\begin{tikzcd}[ampersand replacement=\&, column sep=small, row sep=small]
\bullet \arrow[dr] \& \& \bullet \arrow[dl] \\
\& \bullet \&
\end{tikzcd}}}$. The first two cases are essentially trivial (for $e$, and a product of two categories large relative to $\pi_0$, are large relative to $\pi_0$), the third case on the other hand does not look trivial a priori. This is visibly related to the question of indicial representations of functors $J \longrightarrow \Ind_{\pi_0}(C)$, for $\card\red J \leqslant \pi_0$ -- i.e. the validity of 4.4.13 b) outside the hypothesis that $C$ is stable under $\varinjlim$ of cardinality $\leqslant \pi_0$. We will have to come back to this very soon, for this question now appears to me to be a key question, for the stability properties of the notion of accessible category (without any hypothesis of existence of $\varinjlim$ or $\varprojlim$ ). % [sic] qu'il est filtrante -> qu'il est filtrant, agreement with « il » (C/F)

% marginal handwritten mark (curved stroke) in the left margin beside Proposition 4.8.13
\textbf{Proposition 4.8.13} (\footnote[89]{\textit{Not proved.} But OK if $\pi_0 > \card\red C$, cf. 4.15.10, p. 462.}). \textit{Let $C$ be a category equivalent to a small category, and let $\pi_0 \geqslant \card\red C$ be an infinite cardinal. Then $\cM = \Ind_{\pi_0}C = \cM_{\pi_0}$ (i.e. all the objects of $\Ind_{\pi_0}C$ are $\pi_0$-accessible), and $\cM$ is a Karoubi envelope of $C$.}

\msp{page 213}

Let us point out at once the following corollaries.

% marginal handwritten mark (curved stroke) in the left margin beside Corollaire 4.8.14
\textbf{Corollary 4.8.14} (\footnote[90]{\textit{Not proved.} Yes after all, cf. 4.15.10, and 4.15.11, pages 462-463.}). \textit{Let $\cM$ be a $\mathfrak{U}$-category. Equivalent conditions:}
\begin{enumerate}[leftmargin=*]
\item[\textit{a)}] \textit{$\cM$ is accessible, and is equivalent to a small category.}
\item[\textit{a$'$)}] \textit{$\cM$ is accessible, and there exists an infinite cardinal $\pi_0$ such that all the objects of $\cM$ are $\pi_0$-accessible (i.e. $\cM = \cM_{\pi_0}$).}
\item[\textit{b)}] \textit{$\cM$ is equivalent to a small category, and is Karoubian, i.e. stable under direct factors} (\footnote[91]{\textbf{N.B.} This implies that $\cM$ and $\cM^{\circ}$ are accessible, and I conjecture that the converse is true as well.}). % [sic] sont accessible -> sont accessibles, plural agreement (footnote 91)
\end{enumerate}

The equivalence of a) and a$'$) is clear, since $\cM$ admits a canonical generalized cardinal filtration $(\cM_{\pi})_{\pi\geqslant\pi_0}$, consisting of categories equivalent to small categories (condition
% foot: Version 17 décembre 2025
\printedpage{76}
% head: A. Grothendieck : Dérivateurs, Ch. XVIII, Édition des Universités de Montpellier II et Paris VII
a) for fcg), and such that every such subcategory is contained in one of the $\cM_{\pi}$ (follows from condition b) for fcg). Since $\cM_{\pi_0}$ is equivalent to a small category, and is Karoubian (being stable under filtering $\varinjlim$ large relative to $\pi_0$), it is also clear that a$'$) $\Longrightarrow$ b). It remains to prove b) $\Longrightarrow$ a$'$), now this follows at once from 4.8.13 applied to $\cM$, which moreover shows that for $\pi_0$ one can take any infinite cardinal such that $\pi_0 \geqslant \card\red\cM$.

\msp{page 214}

\textbf{4.8.15.} \textsc{Proof of} 4.8.13. We are going to prove:

\textbf{Lemma 4.8.15.1.} \textit{Under the conditions of 4.8.13, let $F \in \Ob\cM$, $x \in \Ob C$, then}
\[
\card F(x) \ \leqslant\ \pi_0\ . \tag*{$(*)$}
\]

This therefore means, by 4.7.1, a$''$ $\Longrightarrow$ b for $\pi \geqslant \pi_1$, taking into account that here $\pi_1 = \pi_0$, $\pi = \pi_0$, that we have $F \in \Filt^{\pi_0}$, whence a fortiori $F \in \cM_{\pi_0}$, which implies $\cM = \cM_{\pi_0}$, i.e. proposition 4.8.13, taking 4.8.5 into account. So it all comes down to proving $(*)$.
\DERmark{ENV1222}\DERmark{ENV1223}


Let us write
\[
F \ = \ \varinjlim_{i\in I} x_i\ , \qquad\quad \text{the } x_i \in \Ob C,\ I \text{ ordered set large relative to } \pi_0,
\]
whence
\[
F(x) \ = \ \varinjlim_{i\in I} \Hom_C(x, x_i)\ .
\]
We may suppose
\[
\card C \ \leqslant \ \pi_0\ ,
\]
\DERmark{ENV1224}

replacing $C$ by an equivalent reduced category if need be.

For every $\xi \in \Ob C$, let
\[
I_{\xi} \ = \ \{ i \in I \mid x_i = \xi \}\ ,
\]
then
\[
I \ = \ \bigcup_{\xi\in\Ob C} I_{\xi}\ ,
\]
\msp{page 215}
a union of a family of subsets of $I$ of cardinality $\card\Ob C \leqslant \pi_0$. Since $I$ is large relative to $\pi_0$, we see that this implies that one of the $I_{\xi}$ is cofinal in $I$. Replacing $I$ by $I_{\xi}$ if need be, we may suppose $I_{\xi} = I$, i.e. that all the $x_i$ are equal to the same element $\xi$. Replacing $C$ by the full subcategory $C_0 = \{\xi\}$ if need be, we may suppose $C$ reduced to a single element. (For if we prove that $F$, as an object of $\Ind_{\pi_0}(C_0)$, is a direct factor of an object of $C_0$ (i.e. of $\xi$), the same will hold for $F$ as an object of $\Ind C$). Thus we may suppose
\DERmark{ENV1225}

\[
C \ = \ \{\xi\}\ ,
\]
thus $C$ given by the monoid
\[
M \ = \ \End_C(\xi)\ , \qquad\quad \card M \ \leqslant\pi_0\ ,
\]
\DERmark{ENV1226}

and giving the direct system $(x_i)_{i\in I}$ amounts to giving a system of elements
\[
\alpha_{ji} \ \in\ M \qquad\quad (i, j \in I,\ i \leqslant j)
\]
% foot: Version 17 décembre 2025
\printedpage{77}
% head: A. Grothendieck : Dérivateurs, Ch. XVIII, Édition des Universités de Montpellier II et Paris VII
satisfying
\[
\alpha_{ii} \ = \ 1\ , \qquad\quad \alpha_{kj}\alpha_{ji} \ = \ \alpha_{ki} \ \ \text{if } i \leqslant j \leqslant k.
\]
We have
\[
F(\xi) \ = \ \varinjlim_{i\in I} M_i \qquad\quad
\left\{
\begin{array}{l}
\text{where } M_i = M \ \forall\, i \in I, \text{ and } M_i \longrightarrow M_j \\
\text{is given by} \\[2ex]
\qquad\qquad u \longmapsto \alpha_{ji}u\ . \\[4ex]
\end{array}
\right.
\]

\msp{page 216}

I have been able to reduce, by similar arguments (\footnote[92]{For these reductions, cf. 4.15.8 (p. 453-463). The ultimate reduction is found in 4.15.8.9 (p. 461), and its corollary 4.15.10.}), to the case where the monoid $M = \End(\xi)$ is a monoid of monomorphisms, i.e.
\[
uv \ = \ uv' \ \ \Longrightarrow\ \ v \ = \ v' \qquad\quad \text{if } u, v, v' \in M\ ,
\] % [sic] u, v, v' \in \cM -> u, v, v' \in M, the elements belong to the monoid M
\DERmark{ENV1227}

which implies that the transition arrows in the direct system of the $M_i = M$ are monomorphisms. It follows that one can have $\card F(\xi) \leqslant \pi_0$ ($I$ being large relative to $\pi_0$) only if $\alpha_{ji_0}$ is an isomorphism for $j \geqslant i_0$, and $i_0$ large enough, i.e. if the induced direct system for $j \geqslant i_0$ is constant, a fortiori that the direct system under consideration is essentially constant. But I have just spent nearly four hours trying to prove this point, without succeeding! This seems more delicate to prove than I had expected, but I still believe that it really is true.

\msp{page 217}

\subsection*{4.9\quad The indicial representation theorem for $\Ind_{\pi}(C)$}

\textbf{4.9.1} (\footnote[93]{\texttt{[The numbers 4.9.1 to 4.9.6 are used twice, 4.9.3 three times.]}}). Let $C$ be a category equivalent to a small category, $\pi_0$ an infinite cardinal,
\[
\cM \ = \ \Ind_{\pi_0}(C) \ \subset\ C^{\wedge}\ , \tag*{(4.9.1)}
\]
$J$ a category equivalent to a small category. Consider the fully faithful inclusion $C \hookrightarrow \cM$, and the corresponding inclusion of $\mathfrak{U}$-categories
\[
\uHom(J, C) \ \hookrightarrow\ \uHom(J, \cM) \qquad\quad \text{(fully faithful).} \tag*{(4.9.2)}
\]
We know, since $C \subset \cM_{\pi_0}$, that the objects of $\uHom(J, C)$ give $\pi_0$-\textit{accessible} objects of $\uHom(J, \cM)$, at least if we suppose
\[
\card\red J \ \leqslant\ \pi_0\ , \tag*{(4.9.3)}
\]
which we shall do henceforth. (Cf. 4.5.8, 2°, p. 119.) As the right-hand side of (4.9.2) satisfies $\mathrm{L}_{\pi_0}$ ($\cM$ satisfies it), it follows that this functor (4.9.2) extends to
\[
\Ind_{\pi_0}(\uHom(J, C)) \ \hookrightarrow\ \uHom(J,\ \underbrace{\cM}_{=\,\Ind_{\pi_0}(C)}\ )\ , \tag*{(4.9.3)}
\]
% foot: Version 17 décembre 2025
\printedpage{78}
% head: A. Grothendieck : Dérivateurs, Ch. XVIII, Édition des Universités de Montpellier II et Paris VII
and this functor is \textit{fully faithful} thanks to the accessibility result we have just stated. I am interested in establishing cases where this functor is essentially surjective, i.e. is an \textit{equivalence} of categories. Let us note that for
\[
\varphi \ :\ J \longrightarrow \cM
\]
to belong to the essential image, it is necessary and sufficient that there exist a functor
\[
\Phi \ :\ I \ \longrightarrow\ \uHom(J, C)\ ,\ I \text{ large relative to } \pi_0 \text{ (and we may suppose } I \text{ ordered)}\ , \tag*{(4.9.4)}
\]

\msp{page 218}

such that we have
\[
\varphi \ \xleftarrow{\ \sim\ }\ \varinjlim_{i\in I} \Phi(i)\ , \tag*{(4.9.5)}
\]
or what amounts to the same thing, such that we have
\[
\varphi(j) \ \xleftarrow{\ \sim\ }\ \varinjlim_{i\in I} \Phi(i, j)\ , \qquad\quad \text{for } j \in J\ , \tag*{(4.9.6)}
\]
where we have put
\[
\Phi(i, j) \ = \ \Phi(i)(j)\ ,
\]
thus interpreting $\Phi$ as a functor
\[
\Phi \ :\ I \times J \ \longrightarrow\ C\ ; \tag*{(4.9.7)}
\]
the functorial law of $\varphi(j)$ given by (4.9.6) being evident. In other words, the $\varphi$ in the essential image of (4.9.3) are those which admit an \textit{indicial representation} by a functor $\Phi$ in the form (4.9.4), or in the form (4.9.7), with \textit{in addition} $I$ large relative to $\pi_0$. I recall in this connection the % [sic] a ce sujet -> à ce sujet, missing accent
\DERmark{ENV1228}


\textbf{Proposition 4.9.2.} \textit{Under the preceding conditions on $C$, $\pi_0$, $J$ (in particular} $\card\red J \leqslant \pi_0$\textit{), if $\varphi : J \longrightarrow \cM \defeq \Ind_{\pi_0}(C)$, for $\varphi$ to belong to the essential image of the fully faithful functor (4.9.3), it is necessary and sufficient that it satisfy the following two conditions:}

\msp{page 219}

\begin{enumerate}[leftmargin=*]
\item[\textit{a)}] $\uHom(J, C)/\varphi$ \textit{is large relative to $\pi_0$.}
\item[\textit{b)}] \textit{For every $j \in J$, the functor}
\[
\uHom(J, C)/\varphi \ \longrightarrow\ C/\varphi(j) \qquad\quad
\left\{
\parbox{0.35\textwidth}{\textit{functor between filtering categories large relative to $\pi_0$, by virtue of a),}}
\right. \tag*{(4.9.2.1)}
\]
\textit{induced by the functor}
\[
\left\{
\begin{array}{rcl}
\uHom(J, C) & \longrightarrow & C \\
F & \longmapsto & F(j)\ ,
\end{array}
\right. \tag*{(4.9.2.2)}
\]
\textit{is} cofinal.
\end{enumerate}
% foot: Version 17 décembre 2025


\endgroup
% END UNIT der-ch18/der-ch18_p076-080.tex
% BEGIN UNIT der-ch18/der-ch18_p081-085.tex
\begingroup\makeatletter%
\newcommand{\printedpage}[1]{\marginpar{\footnotesize p.~#1}}%
\newcommand{\msp}[1]{\marginpar{\footnotesize [#1]}}%
\newcommand{\uHom}{\underline{\mathrm{Hom}}}%
\newcommand{\Ind}{\mathrm{Ind}}%
\newcommand{\Ob}{\operatorname{Ob}}%
\newcommand{\Fl}{\operatorname{Fl}}%
\newcommand{\card}{\operatorname{card}}%
\newcommand{\cardred}{\mathop{\textup{card\,red}}\nolimits}%
\newcommand{\cM}{\mathcal{M}}%
\newcommand{\rbox}[1]{\tikz[baseline=(X.base)]\node[draw,rounded corners=4pt,inner sep=2pt](X){\textup{#1}};}%
\makeatother\let\printedpage\derpp\let\msp\dermsp%
\let\mspin\DERmspin%


\printedpage{79}
% head: A. Grothendieck : Dérivateurs, Ch. XVIII, Édition des Universités de Montpellier II et Paris VII

The main result in the present section is the following:

\textbf{Theorem 4.9.3.} \textit{Under the general conditions of 4.9.1 for $C$, $\pi$, $J$ (in particular} $\cardred J \leqslant \pi$\textit{), the fully faithful functor (4.9.3) is an equivalence of categories in each of the following two cases:}
\DERmark{ENV1229}

\begin{itemize}[label={}]
{\itshape
\item[a)] $C$ is stable under countable filtering $\varinjlim$.
\item[b)] There exists a cardinal $\mu < \pi$, and a cardinal $\mu' > \mu$ \textup{(}\textbf{\upshape N.B.} We may suppose that $\mu'$ is the successor of $\mu$, hence $\mu' \leqslant \pi$), such that $C$ is stable under direct limits large relative to $\mu$, and of cardinality $\mu'$.
}
\DERmark{ENV1230}

\end{itemize}
(\footnote[94]{Common case: there exists a limit ordinal $\alpha$, such that 1) $\card \alpha \leqslant \pi$. 2) The direct limits \texttt{[of]} type $I_\alpha = \{\beta \mid \beta < \alpha\}$ are representable in $C$.}).
\DERmark{ENV1231}


\textbf{Remark 4.9.4.}
\begin{itemize}[label={}]
\item[a)] This theorem, case a), contains theorem 4.4.13, part b) (p. 95-96), where one supposes that $C$ is stable under direct limits of cardinality $\leqslant \pi_0$, the hypothesis a) of 4.9.3 being considerably weaker.
\item[b)] \msp{page 220}I presume that in the absence of a restrictive hypothesis of type a) or b) above, the fully faithful functor (4.9.3) is not always essentially surjective, even in the particular case where $J$ is a finite category, e.g. $J = B_G$, $G = \{\pm 1\}$. For a counterexample in the very closely related case where we take $\cM = \Ind(C)$ (instead of $\Ind_{\pi_0}(C)$), and where we consider the functor (fully faithful if $J$ is a category equivalent to a \textit{finite} category) % [sic] équivalent -> équivalente, agreement with « une catégorie »
\[
\Ind(\uHom(J,C)) \longrightarrow \uHom(J,\Ind(C)),
\]
I refer to exercise SGA 4 I 8.8.8.
\end{itemize}
The main application of theorem 4.9.3 consists in the

\textbf{Corollary 4.9.5.} \textit{Let $\cM$ be an} accessible \textit{category (4.8.7), $J$ a category equivalent to a small category. Then the category $\cM(J^\circ) = \uHom(J,\cM)$ is accessible. More precisely, let $\pi_0$ be a} useful \textit{cardinal for $\cM$ (i.e. such that $\Ind_{\pi_0}(\cM_{\pi_0}) \xrightarrow{\sim} \cM$, i.e. a cardinal such that $\cM$ is $\pi_0$-accessible), let $\pi$ be an infinite cardinal such that $\pi^{\pi_0} = \pi \geqslant \cardred J$} (\footnote[95]{We thus have $\pi = \pi^{\pi_0} \geqslant 2^{\pi_0} > \pi_0$.}), \textit{i.e. $\pi = c^{\pi_0}$, where $c$ is any cardinal $c \geqslant \sup(\pi_0, \cardred J)$, then $\uHom(J,\cM)$ is $\pi$-accessible, more precisely we have an equivalence of categories} % doubt resolved against PDF; see corrections.csv
\[
\uHom(J,\cM) \xleftarrow{\ \sim\ } \Ind_\pi(\uHom(J,\cM_\pi)) \ . \tag*{(4.9.5.1)}
\]

\msp{page 221}We know indeed, with the assumptions made on $\pi_0$, $\pi$, that $\pi$ is useful for $\cM$, i.e. that
\[
\cM \xleftarrow{\ \sim\ } \Ind_\pi(\cM_\pi) \ ,
\]
on the other hand, $\cM_\pi$ is stable under $\varinjlim$ large relative to $\pi_0$ and of cardinality $\leqslant \pi$. We may therefore apply 4.9.3 b), taking $C = \cM_\pi$, $\pi$ and $\pi_0$ as here (hence $\pi_0 < 2^{\pi_0} \leqslant \pi$), and $\pi'_0 = \pi$.

\printedpage{80}
% head: A. Grothendieck : Dérivateurs, Ch. XVIII, Édition des Universités de Montpellier II et Paris VII

\textbf{4.9.6} (\footnote[96]{Preliminary results (for the proof of 4.9.3).}). The rest of the present section 4.9 is devoted to the proof of 4.9.3, i.e. to verifying, for $J \longrightarrow \cM$, conditions a), b) of 4.9.2. In all that follows, we therefore fix a
\[
\varphi \ : \ J \longrightarrow \cM \ = \ \Ind_\pi(C) \ . \tag*{(4.9.6.1)}
\]
We may suppose $J$ reduced, hence small and $\card J \leqslant \pi_0$. We shall consider the functor deduced from the functor (4.9.2.1)
\[
\uHom(J,C)/\varphi \longrightarrow \prod_{j \in \Ob J} C/\varphi(j) \ , \tag*{(4.9.6.2)}
\]
from which the functors in question are deduced by composing with the projection functors
\[
\prod C/\varphi(j) \longrightarrow C/\varphi(j_0) \ . \tag*{(4.9.6.3)}
\]
Since the factor categories $C/\varphi(j)$ are filtering (and even large relative to $\pi$), their product is also filtering, and even large relative to $\pi$ (immediate), and likewise
\msp{page 222}
for the partial products. It follows at once that the projections (4.9.6.3) are cofinal, by the

\textbf{Lemma 4.9.6.1.} \textit{Let $I$, $I'$ be two categories with $I'$ \textup{0}-connected (e.g. filtering), then the projection functor $pr_1 : I \times I' \longrightarrow I$ is cofinal.}

Indeed, for $i \in \Ob I$ \texttt{[rather} $i \in \Ob I$\texttt{]}, we have
\DERmark{ENV1232}

\[
i\backslash(I \times I') \ \simeq \ (i\backslash I) \times I' \ ,
\]
now $i\backslash I$ and $I'$ being 0-connected, the same is true of their product.

\textbf{Lemma 4.9.6.2} (\footnote[97]{Beware, it is not true that an infinite product of 0-connected categories is 0-connected!}). \textit{Let $(I_\alpha)_{\alpha \in A}$ be a family of} filtering \textit{categories (resp. large relative to $\pi$). Then $I_A \stackrel{\mathrm{def}}{=} \prod_{\alpha \in A} I_\alpha$ is filtering (resp. large relative to $\pi$), and for every subset $B \subset A$, the canonical projection $I_A \longrightarrow I_B$ is cofinal.} % [sic] cofinal -> cofinale, agreement with « la projection »
\DERmark{ENV1233}


(\footnote[98]{Cf. the elegant proof below, 4.9.6.6.}). The stability of the property ``filtering'' or ``large relative to $\pi$'' by AQT \texttt{[âne qui trotte: at a donkey's trot, i.e. by routine verification]}. We have, if $B' = A\backslash B$, $I_A \simeq I_B \times I_{B'} \longrightarrow I_B$, and we apply lemma 4.9.6.1 to conclude that this functor is cofinal, $I_{B'}$ being filtering, hence 0-connected.

\msp{page 223}Since the functor (4.9.6.3) is cofinal, to establish that (4.9.2.1) is so, it suffices to prove that the functor (4.9.6.2) is so, which we shall do. Let us note moreover that the said functor is \textit{faithful}, for it fits into the commutative square
\[
\begin{tikzcd}[column sep=large, row sep=large]
\uHom(J,C)/\varphi \arrow[r] \arrow[d] & \prod C/\varphi(j) \arrow[d] \\
\uHom(J,C) \arrow[r] & C^{\Ob J}\ ,
\end{tikzcd}
\]
\printedpage{81}
% head: A. Grothendieck : Dérivateurs, Ch. XVIII, Édition des Universités de Montpellier II et Paris VII
where the vertical arrows and the second horizontal arrow are faithful. Ultimately, to establish theorem 4.9.3, we shall study the \textit{faithful} functor (4.9.6.2)
\[
\begin{aligned}
&f \ : \ \uHom(J,C)/\varphi \ \longrightarrow \ \textstyle\prod_{j \in J_0} C/\varphi(j) \\ % [sic] \Phi -> \varphi, the functor is written \varphi everywhere else
&\text{(where we have put } J_0 = \Ob J,\ J_1 = \Fl J) \ ,
\end{aligned}
\tag*{$(*)$}
\]
\DERmark{ENV1234}

and we shall prove:
\begin{itemize}[label={}]
\item[\textcircled{a}] This faithful functor $(*)$ is \textit{cofinal.}
\item[\textcircled{b}] The source category $\uHom(J,C)/\varphi$ is large relative to $\pi_0$.
\end{itemize}
For this we shall bring out an auxiliary lemma on the cofinality of \textit{faithful} functors:

\msp{page 224}\textbf{Lemma 4.9.6.3} (\footnote[99]{This lemma is part of the general sorites on cofinality and filtering categories.}). \textit{Let $f : I \longrightarrow I'$ be a functor, and let us consider the conditions:}
{\itshape
\begin{itemize}[label={}]
\item[a)] $I$ filtering, $f$ cofinal, i.e. (in the presence of $I$ filtering) $f$ satisfies the familiar conditions \rbox{F 1}, \rbox{F 2}.
\item[b)] $I'$ is filtering, and the $i'\backslash I$ (for $i' \in \Ob I'$) are filtering.
\item[b$'$)] $I'$ \textup{0}-connected, satisfies PF (2) \textup{(}\footnote[100]{e.g. PF 2 is satisfied if $I'$ is \textit{ordered}, without having to suppose in advance $I'$ filtering, only 0-connected.}\textup{)} (equalization of double arrows), and the categories $i'\backslash I$ are filtering for every $i'\in\Ob I'$.
\item[b$''$)] $I'$ is filtering, the $i'\backslash I$ pseudo-filtering, i.e. satisfy condition PF $1'$.
\item[c)] $I'$ filtering, and $f$ satisfies F 1, as well as the following condition.
\begin{itemize}[label={}]
\item[\rbox{F 3}] For every pair $(i_0,i_1)$ of objects of $I$ and every arrow $f(i_0) \xrightarrow{u'} f(i_1)$, there exist an object $i_2$ of $I$, and an arrow $v' : f(i_1) \longrightarrow f(i_2)$, such that $v'$ and $v'u'$ can be lifted to arrows $i_1 \xrightarrow{v} i_2$ and $i_0 \xrightarrow{w} i_2$ (i.e. there exist $i_2$ and $v : i_1 \longrightarrow i_2$, $w : i_0 \longrightarrow i_2$ such that we have $f(v)u' = f(w)$).
\end{itemize}
\end{itemize}
\[
\begin{tikzcd}
i_0 \arrow[rr, bend left=20, "w"] & i_1 \arrow[r, "v"] & i_2 % doubt resolved against PDF; see corrections.csv
\end{tikzcd}
\]
\[
\begin{tikzcd}
f(i_0) \arrow[r, "u'"] \arrow[rr, dotted, bend right=20, "f(w)"'] & f(i_1) \arrow[r, "f(v) = v'"] & f(i_2)
\end{tikzcd}
\]
\begin{itemize}[label={}]
\item[1)] We have the implications
\[
a \iff b \iff b' \implies b'' \iff c \ ,
\]
and if $f$ is \textup{faithful} (e.g. $I$ ordered), the five conditions a), b), b$'$), b$''$), c) are equivalent.
\end{itemize}
}
\DERmark{ENV1235}


\printedpage{82}
% head: A. Grothendieck : Dérivateurs, Ch. XVIII, Édition des Universités de Montpellier II et Paris VII
{\itshape
\begin{itemize}[label={}]
\item[2)] \msp{page 225}Let $\pi$ be an infinite cardinal, and let us consider the conditions:
\begin{itemize}[label={}]
\item[$a_\pi$)] $I$ large relative to $\pi$, $f$ cofinal, i.e. F 1, F 2.
\item[$b_\pi$)] $I'$ large relative to $\pi$, the $i'\backslash I$ ($i' \in \Ob I'$) also.
\item[$b_\pi{}''$)] $I'$ large relative to $\pi$, the $i'\backslash I$ ($i' \in \Ob I'$) are $\pi$-pseudo-filtering, i.e. satisfy PF $1_\pi{}'$. % [sic] PS 1_π' -> PF 1_π', the conditions are elsewhere called PF
\item[$c_\pi$)] $I'$ large relative to $\pi$, and $f$ satisfies F 1 and
\begin{itemize}[label={}]
\item[\rbox{F $3_\pi$}] \textup{(}\footnote[101]{\textbf{N.B.} The condition F $3_\pi$ is tautologically satisfied if $f$ is \textit{fully faithful.}}\textup{)}. For every family $(i_\alpha)_{\alpha \in A}$ of objects of $I$, with \textup{card} $A \leqslant \pi$, and $i \in \Ob I$, and a family of arrows $u'_\alpha : f(i_\alpha) \longrightarrow f(i)$, exists $j \in \Ob I$ and an arrow $v' : f(i) \longrightarrow f(j)$, such that $v'$ \textup{and} the $v'u'_\alpha$ come from arrows of $I$, i.e. there exist arrows $v : i \longrightarrow j$, $w_\alpha : i_\alpha \longrightarrow j$ such that $f(v)u'_\alpha = f(w_\alpha)$ $\forall\ \alpha \in A$.
\end{itemize}
\end{itemize}
\[
\begin{tikzcd}[row sep=huge]
i_\alpha \arrow[d, dash] \arrow[rr, dotted, bend left=20, "w_\alpha"] & i \arrow[r, dotted, "v"] \arrow[d, dash] & j \arrow[d, dash] \\
f(i_\alpha) \arrow[r, "u'_\alpha"] & f(i) \arrow[r, "v' = f(v)"] & f(j)
\end{tikzcd}
\]
This being laid down, we have
\[
a_\pi \iff b_\pi \implies c_\pi \implies b_\pi{}'' \ ,
\]
and if $f$ is faithful, or $I$ ordered, the four conditions $a_\pi$, $b_\pi$, $b_\pi{}''$, $c_\pi$ are equivalent.
\end{itemize}
}
\DERmark{ENV1236}

(\footnote[102]{Proof of the proposition given in 4.9.6.7, page 228.}).
\DERmark{ENV1237}


\textbf{4.9.6.4. N.B.} In the case of interest to us, where $f$ is the functor $(*)$ (p. 223), we know that $I'$ is large relative to $\pi$. We would like to establish that $I$ is so and that $f$ is cofinal, i.e. condition $\mathrm{a}_\pi$. We shall prove it by taking $\mathrm{c}_\pi$ (which implies $\mathrm{a}_\pi$ since $f$ is faithful). It therefore remains to
\DERmark{ENV1238}

\msp{page 226}
prove, for the functor $f$ in $(*)$, p. 223, the following two conditions:
\begin{itemize}[label={}]
\item[F 1)] For every system of arrows in $\cM$
\[
X(j) \xrightarrow{p(j)} \varphi(j) \ , \qquad j \in J_0 = \Ob J, \text{ the } X(j) \text{ in } \Ob C \ ,
\]
there exists a functor $F \ : \ J \longrightarrow C$ and $F \xrightarrow{q} \varphi$, finally a system of arrows $X(j) \xrightarrow{u(j)} F(j)$ in the $C/\varphi(j)$, i.e. making the diagrams commutative % [sic] commutatif -> commutatifs, agreement with « les diagrammes »
\[
\begin{tikzcd}[column sep=small]
X(j) \arrow[rr, "u(j)"] \arrow[dr, "p(j)"'] & & F(j) \arrow[dl, "q(j)"] \\
& \varphi(j)\ . &
\end{tikzcd}
\]
\item[F $3_\pi$)] For every system $(F_\alpha)_{\alpha \in A}$ ($\card A \leqslant \pi$) of elements in $\uHom(J,C)$, and arrows % [sic] \uHom(C,J) -> \uHom(J,C), the F_α are functors J -> C
\[
F_\alpha \xrightarrow{p_\alpha} \varphi \qquad \text{in } \uHom(J,\cM),
\]
\DERmark{ENV1239}

and for $G$ in $\uHom(J,C)$ and $G \xrightarrow{q} \varphi$ in $\uHom(J,\cM)$, given a system of arrows $u_\alpha(j)$,
\printedpage{83}
% head: A. Grothendieck : Dérivateurs, Ch. XVIII, Édition des Universités de Montpellier II et Paris VII
\[
\begin{tikzcd}[column sep=small]
F_\alpha(j) \arrow[rr, "u_\alpha(j)", "\text{comm.}"'] \arrow[dr, "p_\alpha(j)"'] & & G(j) \arrow[dl, "q(j)"] \\
& \varphi(j)\ , &
\end{tikzcd}
\qquad u_\alpha(j) \text{ arrow in } C/\varphi(j) \ ,
\]
there exists $H$ in $\uHom(J,C)$ and $H \xrightarrow{q} \varphi$ in $\uHom(J,\cM)$ and an arrow $G \xrightarrow{v} H$ in $\uHom(J,C)/\varphi$, such that for every $\alpha \in A$, the system of composite arrows
\DERmark{ENV1240}

\msp{page 227}
\[
\begin{tikzcd}[row sep=huge]
F_\alpha(j) \arrow[r, "u_\alpha(j)", "\text{comm.}"'] \arrow[dr, dotted, "p_\alpha(j)"'] & G(j) \arrow[r, "v(j)", "\text{comm.}"'] \arrow[d, dotted, "q(j)"] & H(j) \arrow[dl, dotted, "q(j)"] \\
& \varphi(j) &
\end{tikzcd}
\]
(for $j \in J_0 = \Ob J$) comes from a \textit{natural transformation} $w_\alpha : F_\alpha \longrightarrow H$. % [sic] J_0 \in \Ob J -> J_0 = \Ob J, J_0 was defined as Ob J
\DERmark{ENV1241}\DERmark{ENV1242}

\end{itemize}
To prove 4.9.3, we are thus reduced to proving the two preceding propositions F 1), F $3_\pi$), under one of the two conditions a) or b) of 4.9.3.

\textbf{Corollary 4.9.6.5} (of 4.9.6.3). \textit{Let $(f_\alpha : I_\alpha \longrightarrow I'_\alpha)$ be a family of cofinal functors, the source categories being filtering. Then the product functor $f = \prod f_\alpha : I = \prod_\alpha I_\alpha \longrightarrow I' = \prod_\alpha I'_\alpha$ is cofinal.}

Indeed, we must verify that the categories $i'\backslash I$, for $i' \in \Ob I'$, are 0-connected. Now if $i' = (i'_\alpha)_{\alpha \in A}$, we have
\[
i'\backslash I \ \simeq \ \prod_\alpha (i'_\alpha\backslash I_\alpha) \ ,
\]
and by 4.9.6.3, 1°) a $\Longrightarrow$ b, the $i'_\alpha\backslash I_\alpha$ are filtering, hence their product is filtering
\DERmark{ENV1243}

\msp{page 228}
(cf. 4.9.6.2), hence 0-connected, q.e.d. % [sic] 4.4.6.2 -> 4.9.6.2, the lemma on products of filtrantes is 4.9.6.2
\DERmark{ENV1244}


\textbf{Corollary 4.9.6.6.} \textit{Let $\pi$ be an infinite cardinal. Every product of categories large relative to $\pi$ is large relative to $\pi$.} % [sic] 4.4.6.6 -> 4.9.6.6, numbering of the section and footnote 98
\DERmark{ENV1245}


Let $I = \prod_{\alpha \in A} I_\alpha$, let $B$ be a set of cardinality $\leqslant \pi$, we must prove that $I \xrightarrow{\mathrm{diag}_I} I^B$ is cofinal. Now this functor is identified with the product of the similar functors $I_\alpha \longrightarrow I_\alpha^B$, which are cofinal, the source categories being moreover filtering. Thus we may conclude by corollary 4.9.6.5.

\textbf{4.9.6.7.} \textsc{Proof of} 4.9.6.3. (Part 1°.)

a $\Longrightarrow$ b. We already know that a) implies that $I'$ is filtering (SGA 4 I 8.1.3 b). Remains to prove that for every $i' \in \Ob I'$, $i'\backslash I$ is filtering (\footnote[103]{Short-circuit: we already know that $i'\backslash I$ is 0-connected, on the other hand it is \textit{locally} filtering, since coétale over $I$ which is filtering, hence it is filtering.}). Two conditions PF $1'$ and PF 2 to verify (in addition to the fact that $i'\backslash I \neq \emptyset$, which is known by cofinality of $f$, implying that $i'\backslash I$ is 0-connected).
\begin{itemize}[label={}]
\item[1)] Let $(i,u')$, $(j,v')$ be two elements of $i'\backslash I$, we must see that they are bounded above by a third element. Now $I$ being filtering, $i$, $j$ are bounded above by $k$. We would therefore be done if we knew that $f(u_1)u' = f(v_1)v'$, say $w'$, for then $(k,w')$ would do.
\end{itemize}


\endgroup
% END UNIT der-ch18/der-ch18_p081-085.tex
% BEGIN UNIT der-ch18/der-ch18_p086-090.tex
\begingroup\makeatletter%
\newcommand{\printedpage}[1]{\marginpar{\footnotesize p.~#1}}%
\newcommand{\msp}[1]{\marginpar{\footnotesize [#1]}}%
\newcommand{\bs}{\backslash}%
\newcommand{\encadre}[1]{\tikz[baseline=(X.base)]\node[draw,rounded corners,inner sep=2pt](X){#1};}%
\makeatother\let\printedpage\derpp\let\msp\dermsp%
\let\mspin\DERmspin%


\printedpage{84}
% head: A. Grothendieck : Dérivateurs, Ch. XVIII, Édition des Universités de Montpellier II et Paris VII

\[
\begin{tikzcd}
{} & i \arrow[drr, "u_1"] \arrow[dd, dash] & {} & {} \\
{} & {} & j \arrow[r, "v_1"] & k \arrow[dd, dash] \\
i' \arrow[r, "u'"] \arrow[drr, "v'"'] & f(i) \arrow[drr] & {} & {} \\
{} & {} & f(j) \arrow[r] \arrow[uu, dash, crossing over] & f(k)
\end{tikzcd}
\]

\begin{itemize}[label={}]
\item \msp{page 229}But we apply F 2 for $f$ to the double arrow in question $i' \rightrightarrows f(k)$, we find $w : k \longrightarrow \bar{k}$,
\[
\begin{tikzcd}
{} & k \arrow[r, "w"] \arrow[dd, dash] & \bar{k} \arrow[dd, dash] \\
{} & {} & {} \\
i' \arrow[r, shift left] \arrow[r, shift right] & f(k) \arrow[r] & f(\bar{k})
\end{tikzcd}
\]
such that
\[
\underbrace{f(w)(f(u_1)u')}_{= f(wu_1)u'} = \underbrace{f(w)(f(v_1)v')}_{f(wv_1)v'} \ ,
\]
and replacing $k$ by $\bar{k}$, and $u_1$, $v_1$ by $wu_1$, $wv_1$, we find what we want.
\end{itemize}

\begin{enumerate}[label=\arabic*), start=2]
\item PF (2) for $i'\bs I$, i.e. we give ourselves a diagram opposite \texttt{[below]},
\[
\begin{tikzcd}
{} & i \arrow[r, shift left, "u_1"] \arrow[r, shift right, "u_2"'] \arrow[dd, dash] & j \arrow[r, dotted, "w"] \arrow[dd, dash] & j' \\
{} & {} & {} & {} \\
i' \arrow[r, "u'"] & f(i) \arrow[r, shift left, "f(u_1)"] \arrow[r, shift right, "f(u_2)"'] & f(j) \arrow[r, dotted, "f(w)"] & f(j') \ ,
\end{tikzcd}
\]
with $f(u_1)u' = f(u_2)u'$, say $v'$ (double arrow in $i'\bs I$ from $(i,u')$ to $(j,v')$), we want to equalize it on the right. Now $I$ being filtering, there exists $w : j \longrightarrow j'$ such that $wu_1 = wu_2$, now this arrow defines an arrow $(j,v') \longrightarrow (j', f(w)v')$, which equalizes the given double arrow. % [sic] f(w')v' -> f(w)v', the arrow just defined is w : j --> j'
\DERmark{ENV1246}

\end{enumerate}

b $\Longrightarrow$ b$'$ is clear.

b) $\Longrightarrow$ a). Since the $i'\bs I$ are filtering, hence 0-connected, $f$ is cofinal, it remains to prove that $I$ is filtering, hence to check PF 0 ($I$ 0-connected), PF 1, PF 2 for $I$. $I$ is 0-connected, since $I'$ is, and $f$ cofinal, i.e. $\mathbf{W}_0$-coaspheric, whence PF 0.

\msp{page 230}Let us prove PF 1, i.e. that if
\begin{tikzcd}[row sep=tiny, column sep=small, baseline=(\tikzcdmatrixname-2-1.base)]
{} & i_1 \\
i_0 \arrow[ur] \arrow[dr] & {} \\
{} & i_2
\end{tikzcd}
is given, we can complete it to a commutative square.
\[
\begin{tikzcd}[row sep=small, column sep=small]
{} & i_1 \arrow[dr, dotted, "v_1"] & {} \\
i_0 \arrow[ur, "u_1"] \arrow[dr, "u_2"'] & {} & i_3 \\
{} & i_2 \arrow[ur, dotted, "v_2"'] & {}
\end{tikzcd}
\]
% foot: Version 17 décembre 2025

\printedpage{85}
% head: A. Grothendieck : Dérivateurs, Ch. XVIII, Édition des Universités de Montpellier II et Paris VII

\[
\begin{tikzcd}[row sep=small, column sep=small]
{} & f(i_1) \arrow[dr, "f(v_1)"] & {} \\
f(i_0) \arrow[ur, "f(u_1)"] \arrow[dr, "f(u_2)"'] & {} & f(i_3) \\
{} & f(i_2) \arrow[ur, "f(v_2)"'] & {}
\end{tikzcd}
\]

But let us consider the objects $\xi_1 = (i_1, f(u_1))$, $\xi_2 = (i_2, f(u_2))$ of $f(i_0)\bs I$, and using PF 1$'$ for this category, we find $\xi_3 = (i_3, w')$ in the said category, and arrows from $\xi_1$, $\xi_2$ to $\xi_3$, i.e. $v_1 : i_1 \longrightarrow i_3$, \texttt{[}$v_2 : i_2 \longrightarrow i_3$\texttt{,]} such that $\underbrace{f(v_1)f(u_1)}_{f(v_1u_1)} = w'$, $\underbrace{f(v_2)f(u_2)}_{f(v_2u_2)} = w'$, hence
\[
f(v_1u_1) \ = \ f(v_2u_2) \ .
\]
Thus, we have a double arrow $i_0 \overset{v_1u_1}{\underset{v_2u_2}{\rightrightarrows}} i_3$, equalized by $f$. If we prove PF 2 for $I$, we find $i_3 \xrightarrow{\alpha} \bar{i}_3$ such that $\alpha(v_1u_1) = \alpha(v_2u_2)$, i.e. $(\alpha v_1)u_1 = (\alpha v_2)u_2$, and replacing $i_3$ by $\bar{i}_3$, $v_1$, $v_2$ by $\alpha v_1$, $\alpha v_2$ if necessary, we have won.

It remains to prove PF 2 for $I$, i.e. to equalize a double arrow
\[
i_0 \ \overset{u_1}{\underset{u_2}{\rightrightarrows}} \ i_1 \ .
\]
\DERmark{ENV1247}


\msp{page 231}Now $I'$ satisfying PF 2, we can equalize the double arrow $f(u_1)$, $f(u_2)$ by an arrow $f(i_1) \xrightarrow{v'} i'$ in $I'$, and with $i'\bs I$ nonempty, we can suppose that $i'$ is of the form $f(i_2)$.
\[
\begin{tikzcd}
i_0 \arrow[r, shift left, "u_1"] \arrow[r, shift right, "u_2"'] \arrow[dd, dash] & i_1 \arrow[r] & i_2 \\
{} & {} & {} \\
f(i_0) \arrow[r, shift left, "f(u_1)"] \arrow[r, shift right, "f(u_2)"'] & f(i_1) \arrow[r, "v'"] & f(i_2) % doubt resolved against PDF; see corrections.csv
\end{tikzcd}
\]
Let us admit for a moment condition F 3 of c), which we are going to apply to $(i_1, i_2, v')$. There thus exist $i_2 \xrightarrow{v} \bar{i}_2$, $w : i_1 \longrightarrow \bar{i}_2$ such that we have $f(v)v' = f(w)$. It follows that $f(w)$ equalizes $f(u_1)$ and $f(u_2)$, i.e. that we have $f(wu_1) = f(wu_2)$, i.e. the initial double arrow is replaced by $(\bar{u}_1 = wu_1, \bar{u}_2 = wu_2)$, such that $f(\bar{u}_1) = f(\bar{u}_2)$. In short, we are reduced to the case of a double arrow $(u_1, u_2)$ such that $f(u_1) = f(u_2)$, say $v'$. But then let $\xi_1 = (i_1, v')$, $\xi_0 = (i_0, \mathrm{id}_{f(i_0)})$, we thus have a double arrow $(u_1, u_2) : \xi_0 \rightrightarrows \xi_1$ in $f(i_0)\bs I$. This category % [sic] \xi_1 = (i_0, v') -> \xi_1 = (i_1, v'), u_1, u_2 go from i_0 to i_1
\DERmark{ENV1248}\DERmark{ENV1249}

\msp{page 232}is filtering, so this double arrow is equalizable on the right, which implies at once that the double arrow $(v_1, v_2)$ \textit{in $I$} is also equalizable on the right.

So it remains to prove that b$'$) implies condition F 3. But let us consider the two objects $\xi_0 = (i_0, \mathrm{id}_{f(i_0)})$, $\xi_1 = (i_1, u')$ of $f(i_0)\bs I$, and use the filtering condition PF 1 for this category, so we find an object $\xi_2 = (i_2, v')$, and two arrows
\begin{tikzcd}[row sep=tiny, column sep=small, baseline=(\tikzcdmatrixname-2-2.base)]
\xi_0 \arrow[dr, "w"] & {} \\
{} & \xi_2 \\
\xi_1 \arrow[ur, "v"'] & {}
\end{tikzcd}
, i.e. $i_0 \xrightarrow{w} i_2$
% foot: Version 17 décembre 2025
\printedpage{86}
% head: A. Grothendieck : Dérivateurs, Ch. XVIII, Édition des Universités de Montpellier II et Paris VII
and $i_1 \xrightarrow{v} i_2$ such that we have $f(w)\mathrm{id}_{f(i_0)} = f(v)u' = v'$, whence $f(w) = f(v)u'$, which is condition F 3.
\DERmark{ENV1250}

\[
\begin{tikzcd}
i_0 \arrow[rr, dotted, bend left, "w"] \arrow[dd, dash] & i_1 \arrow[r, dotted, "v"] \arrow[dd, dash] & i_2 \arrow[dd, dash] \\
{} & {} & {} \\
f(i_0) \arrow[r, "u'"] \arrow[rr, dotted, bend right, "v'"'] & f(i_1) & f(i_2)
\end{tikzcd}
\]

Thus we have established
\[
\text{a} \ \Longleftrightarrow \ \text{b} \ \Longleftrightarrow \ \text{b}' \ ,
\]
let us now prove

a $\Longrightarrow$ c. We already know that a) implies $I'$ filtering, and it remains to establish F 3. Now $I$ being filtering, there exist arrows
\begin{tikzcd}[row sep=tiny, column sep=small, baseline=(\tikzcdmatrixname-2-2.base)]
i_0 \arrow[dr, "w"] & {} \\
{} & i_2 \\
i_1 \arrow[ur, "v"'] & {}
\end{tikzcd}
. We
\msp{page 233}do not necessarily have $f(w) = f(v)u'$, we have here a double arrow $f(i_0) \rightrightarrows f(i_2)$ in $I'$. But by condition F 2 for $f$, there exists $i_2 \xrightarrow{\alpha} \bar{i}_2$ in $I$ such that $f(\alpha)$ equalizes the said double arrow, hence $f(\underbrace{\alpha w}_{\bar{w}}) = f(\underbrace{\alpha v}_{\bar{v}})u'$, and replacing $i_2$ by $\bar{i}_2$, and $w$, $v$ by $\bar{w} = \alpha w$, $\bar{v} = \alpha v$ if necessary, we win. % [sic] f(i_0) = f(i_1) -> f(i_0) \rightrightarrows f(i_2), the double arrow f(w), f(v)u' goes from f(i_0) to f(i_2)

It remains to prove (for part $1^\circ$) that if $f$ is \textit{faithful}, then c) $\Longrightarrow$ b). It thus remains to prove that c) implies that for $i'$ in $I'$, $i'\bs I$ is filtering. We already know by F 1 that $i'\bs I$ is nonempty, it remains to prove conditions PF 1$'$, PF 2 for $i'\bs I$. With $I'$ filtering, we can complete
\DERmark{ENV1251}

\[
\begin{tikzcd}[row sep=small, column sep=small]
{} & f(i_1) \arrow[dr, dotted] & {} \\
i' \arrow[ur] \arrow[dr] & {} & j' \ , \\
{} & f(i_2) \arrow[ur, dotted] & {}
\end{tikzcd}
\]
and by F 1 we can suppose $j'$ of the form $f(i_3)$.
% foot: Version 17 décembre 2025

\printedpage{87}
% head: A. Grothendieck : Dérivateurs, Ch. XVIII, Édition des Universités de Montpellier II et Paris VII

\[
\begin{tikzcd}[column sep=small]
{} & i_1 \arrow[drr, dotted, "v_1\,?"] \arrow[dd, dash] & {} & {} & {} \\
{} & {} & {} & i_3 \arrow[r] \arrow[dd, dash] & \bar{i}_3 \arrow[dd, dash] \\
{} & f(i_1) \arrow[drr, "v_1'"] & i_2 \arrow[ur, dotted, "v_2\,?"'] \arrow[dd, dash, crossing over] & {} & {} \\
i' \arrow[ur, "u_1'"] \arrow[dr, "u_2'"'] & {} & {} & f(i_3) \arrow[r] & f(\bar{i}_3) \\
{} & {} & f(i_2) \arrow[ur, "v_2'"'] & {} & {}
\end{tikzcd}
\]
Applying F 3 to $i_1$, $i_3$, $v_1'$, we can suppose (replacing $i_3$ by $\bar{i}_3$ if necessary)
\msp{page 234}that $v_1'$ comes from a $v_1 : i_1 \longrightarrow i_3$. Applying F 3 again to $i_2$, $i_3$, $v_2'$, we can suppose that $v_2'$ also comes from a $v_2 : i_2 \longrightarrow i_3$. Then $\xi_3 = (i_3, v_1'u_1' = v_2'u_2')$ is a common upper bound of $\xi_1 = (i_1, u_1')$, $\xi_2 = (i_2, u_2')$.

Let us finally prove PF 2 for $i'\bs I$, so we have a diagram opposite \texttt{[below]} with $f(u_1)u' = f(u_2)u'$ (double arrow in $i'\bs I$), we want to equalize in $i'\bs I$.
\[
\begin{tikzcd}
{} & i_0 \arrow[r, shift left, "u_1"] \arrow[r, shift right, "u_2"'] & i_1 \arrow[rr, bend left, "w"] & i_2 \arrow[r] & \bar{i}_2 \\
{} & {} & {} & {} & {} \\
i' \arrow[r, "u'"] & f(i_0) \arrow[r, shift left, "f(u_1)"] \arrow[r, shift right, "f(u_2)"'] & f(i_1) \arrow[r, dotted, "v'"] & f(i_2) \arrow[r] & f(\bar{i}_2)
\end{tikzcd}
\]
But $I'$ being filtering, we can equalize $(f(u_1), f(u_2))$ by $f(i_1) \longrightarrow j'$, and because of F 1 we can suppose $j' = f(i_2)$. Applying F 3 once more to $(i_1, i_2, v')$, we find $i_1 \longrightarrow \bar{i}_2$ such that $f(w)$ equalizes $f(u_1)$, $f(u_2)$:
\[
f(wu_1) \ = \ f(wu_2) \ .
\]
But $f$ being \textit{faithful}, this implies $wu_1 = wu_2$, and we have thus equalized the double arrow of $i'\bs I$ by $w : (i_1, -) \longrightarrow (\bar{i}_2, -)$.

\subsection*{4.9.6.8. N.B.}

\begin{enumerate}[label=\arabic*$^\circ$)]
\item If we do not suppose $f$ to be faithful, we see that c) nevertheless implies that the $i'\bs I$ satisfy PF 1$'$ (they are therefore pseudo-filtering and 0-connected). So $f$ is cofinal.

\msp{page 235}We see, by the same argument, that $I$ satisfies PF 1$'$, i.e. is pseudo-filtering. Thus, if we introduce the condition:
\begin{itemize}[label=\encadre{b$''$}]
\item $I'$ is filtering, the $i'\bs I$ ($i' \in \mathrm{Ob}\, I'$) are pseudo-filtering (i.e. satisfy PF 1$'$) (hence $f$ cofinal).
\end{itemize}
\DERmark{ENV1252}

We find the implications
\[
(\text{a} \ \Longleftrightarrow \ \text{b} \ \Longleftrightarrow \ \text{b}') \ \ \Longrightarrow \ \ (\text{b}'' \ \Longleftrightarrow \ \text{c}) \ ,
\]
and the equivalent conditions c), b$''$) imply that $I$ is also pseudo-filtering, and that for every double arrow $\xi_0 \overset{u_1}{\underset{u_2}{\rightrightarrows}} \xi_1$ in an $i'\bs I$ or in $I$, there exists an arrow $\xi_1 \xrightarrow{v} \xi_2$ in the same category, such that $vu_1$ and $vu_2 \ : \ \xi_0 \rightrightarrows \xi_2$ have \textit{the same image}
% foot: Version 17 décembre 2025
\printedpage{88}
% head: A. Grothendieck : Dérivateurs, Ch. XVIII, Édition des Universités de Montpellier II et Paris VII
\textit{in $I'$}. So if $f$ is faithful, then these conditions c, b$''$ \texttt{[rather b$''$]} imply that $I$ and the $i'\bs I$ also satisfy the double arrow condition PF 2, hence are filtering, so in this case the five conditions a) b) b$'$) c) b$''$) \texttt{[rather b$''$)]} are equivalent. But also when $I$ is ordered, we know [\,?] that the double arrow condition in $I$ is automatically satisfied.
\DERmark{ENV1253}\DERmark{ENV1254}


\item \msp{page 236}Let us also note that in condition b$'$), we cannot drop the hypothesis that $I'$ satisfies PF 2 (in addition to $I'$ 0-connected and the $i'\bs I$ filtering). To see this, let us take $I' = I$, $f = \mathrm{id}$, so hypothesis a ($\Longleftrightarrow$ b $\Longleftrightarrow$ b$'$) means that $I$ is filtering, $f$ being of course cofinal. But weakened b$'$ means that $I$ is 0-connected and the $i\bs I$ are filtering, and we know that this does \textit{not} imply that $I$ is filtering, even if we suppose $I$ pseudo-filtering. A simple example is given by $I$ defined by
\[
\mathrm{Ob}\, I \ = \ \mathbf{N} \ , \qquad \mathrm{Hom}_I(i,j) \ = \ \left\{
\begin{array}{rcccc}
\emptyset & \text{if} & i & > & j \\[1ex]
\{\mathrm{id}_i\} & \text{if} & i & = & j \\[1ex]
\{\pm 1\} & \text{if} & i & < & j \ ,
\end{array}
\right.
\]
the composition of arrows being carried out by multiplication of the signs $\varepsilon_{k,j}\varepsilon'_{j,i} = (\varepsilon\varepsilon')_{k,i}$ (putting $\mathrm{id}_i = (+1)_{i,i}$). The categories $i\bs I$ are isomorphic to one another and are \textit{ordered} and pseudo-filtering, hence filtering, on the other hand, $I$ is pseudo-filtering, yet it is not filtering, for it does not satisfy PF 2.

\item On the other hand, if $f$ is not faithful, then c) (equivalent to $b''$) does not imply a), as we see by taking $I' = e$ ; then b$''$) is equivalent to $I$ pseudo-filtering, which does not imply that $I$ is filtering.
\end{enumerate}

\subsection*{\msp{page 237}\textbf{4.9.6.9.} \textsc{Proof of 4.9.6.3, part} $2^\circ$.}

This is essentially ``the same'' proof as for part $1^\circ$).

a$_\pi$) $\Longrightarrow$ b$_\pi$). We can short-circuit the direct proof, by looking, for a set $A$ of cardinality $\leqslant \pi$, at the diagram
\[
\begin{tikzcd}
I \arrow[r, "\delta_A"] \arrow[d, "f"'] & I^A \arrow[d, "f^A"] \\
I' \arrow[r, "\delta'_A"] & I'^A
\end{tikzcd}
\]
\DERmark{ENV1255}

\texttt{[rather $\delta'_A$]}. Since $f$ is cofinal, $I$ filtering, $f^A$ is also cofinal (4.9.6.5), on the other hand, $I$ large relative to $\pi$ means that $\delta_A$ is cofinal ($\forall$ $A$ with $\mathrm{card}\, A \leqslant \pi$). So $f^A\delta_A = \delta'_Af$ is cofinal, hence $\delta'_Af$ and $f$ are, hence $\delta'_A$ also, which means that $I'$ is large relative to $\pi$.
\DERmark{ENV1256}\DERmark{ENV1257}\DERmark{ENV1258}\DERmark{ENV1259}


It remains to prove that for $i'$ in $I'$, $i'\bs I$ is large relative to $\pi$. This follows from the fact that $i'\bs I \xrightarrow{p} I$ is coétale, that $i'\bs I$ is 0-connected, and from the
\DERmark{ENV1260}


\textbf{Lemma 4.9.6.9.1.} \textit{Let $I$ be a filtering category (resp. large relative to $\pi$). Then for every $J$ over $I$ coétale over $I$ and $0$-connected, $J$ is filtering (resp. large relative to $\pi$).}
% foot: Version 17 décembre 2025


\endgroup
% END UNIT der-ch18/der-ch18_p086-090.tex
% BEGIN UNIT der-ch18/der-ch18_p091-095.tex
\begingroup\makeatletter%
\newcommand{\printedpage}[1]{\marginpar{\footnotesize p.~#1}}%
\newcommand{\msp}[1]{\marginpar{\footnotesize [#1]}}%
\newcommand{\uGamma}{\underline{\Gamma}}%
\newcommand{\Ind}{\operatorname{Ind}}%
\newcommand{\Ob}{\operatorname{Ob}}%
\newcommand{\card}{\operatorname{card}}%
\newcommand{\trois}[1]{\mathrel{\overset{#1}{\vcenter{\offinterlineskip\hbox{$\longrightarrow$}\vskip0.5pt\hbox{$\longrightarrow$}\vskip0.5pt\hbox{$\longrightarrow$}}}}}%
\makeatother\let\printedpage\derpp\let\msp\dermsp%
\let\mspin\DERmspin%


\printedpage{89}
% head: A. Grothendieck : Dérivateurs, Ch. XVIII, Édition des Universités de Montpellier II et Paris VII
% foot: Version 17 décembre 2025

The filtering case is already known (cf. II \dots). For the ``large relative to $\pi$'' case, we are going to use the

\textbf{Corollary 4.9.6.9.2.} \textit{Let $J$ be a category. For $J$ to be large relative to $\pi$, it}
\msp{page 238}
\textit{is necessary and sufficient that it be $0$-connected, and satisfy PF$_\pi$ 2 (cf. page 81), as well as the following variant of PF$_\pi$ 1$'$.}
\begin{itemize}[leftmargin=4em]
\item[\textit{PF$_\pi$ 1)}] \textit{For every object $x_0$ of $J$, $x_0\backslash J$ satisfies PF$_\pi$ 1$'$, i.e. for every family of arrows $x_0 \xrightarrow{u_\alpha} y_\alpha$ ($\alpha \in A$, $\card A \leqslant \pi$), there exists an arrow $x_0 \xrightarrow{v} z$ and arrows $v_\alpha : y_\alpha \longrightarrow z$, with $v_\alpha u_\alpha = v \ \forall\, \alpha \in A$.}
\end{itemize}

\textsc{Proof}. If $J$ is large relative to $\pi$, then PF$_\pi$ 1) follows at once from the conjunction of PF$_\pi$ 1$'$ and PF$_\pi$ 2, hence the ``only if''. Conversely, suppose that $J$ satisfies PF$_\pi$ 1, PF$_\pi$ 2, and is 0-connected, we must prove that it satisfies PF$_\pi$ 1$'$. But we already know that the conditions in question, which imply PF 1, PF 2, 0-connected, imply that $J$ is filtering. Thus if we have a family $(x_\alpha)_{\alpha\in A}$, $\card A \leqslant \pi$, choosing an $x_0 \in \Ob J$, we can for every $\alpha$ find $y_\alpha$ dominating $x_0$ and $x_\alpha$, hence arrows
\begin{tikzcd}[row sep=0.4em, column sep=1em, baseline=(\tikzcdmatrixname-2-2.base)]
x_\alpha \arrow[dr] & \\
& y_\alpha \\
x_0 \arrow[ur] &
\end{tikzcd}, we can apply PF$_\pi$ 1 to them, hence arrows $y_\alpha \longrightarrow z$, and by composing, arrows $x_\alpha \longrightarrow z$, q.e.d. % [sic] leurs appliquer -> leur appliquer, grammar

Let us now prove the lemma. It suffices (by the preceding corollary) to prove that $J$ satisfies PF$_\pi$ 1 and PF$_\pi$ 2. Now PF$_\pi$ 1 is immediate, since $\forall\, x_0$ in $J$, $x_0\backslash J \xrightarrow{\ \sim\ } y_0\backslash I$, where $y_0$ is
\msp{page 239}
the image of $x_0$ in $I$. As for PF$_\pi$ 2, if we have $x_0 \trois{u_\alpha} x_1$ in $J$, hence $g(x_0) \trois{g(u_\alpha)} g(x_1)$, we equalize in $I$ by $g(x_1) \xrightarrow{w} y$, but since $g$ is coétale, $\exists\, x_1 \xrightarrow{v} x_2$ in $J$ such that $g(v) = w$, so we have $g(vu_\alpha) = g(vu_\beta)$ for $\alpha, \beta \in A$, and since $g$ is faithful, we conclude from this that $vu_\alpha = vu_\beta$, so $v$ equalizes $(u_\alpha)_{\alpha\in A}$, q.e.d.

Let us prove b$_\pi$) $\Longrightarrow$ a$_\pi$): If $I'$ and the $i'\backslash I$ are large relative to $\pi$, then $I$ is too (and $f$ is cofinal, the $i'\backslash I$ being 0-connected, since filtering). Since we already know that $I$ is 0-connected ($I'$ being so and $f$ cofinal), it remains, by 4.9.6.9.2, to verify PF$_\pi$ 1 and PF$_\pi$ 2.

\msp{page 240}
The proof given in 1° for the similar filtering case (implication b$'$ $\Longrightarrow$ a, pages 231-232), applies mutatis mutandis. (Property F 3 has already been established in part 1°), as a consequence of b, and can therefore be used freely.)

Thus we have proved a$_\pi$ $\Longleftrightarrow$ b$_\pi$. Let us prove a$_\pi$ $\Longrightarrow$ c$_\pi$. We already know that a$_\pi$ implies that $I'$ is large relative to $\pi$, and of course condition F 1, it remains to establish the validity of condition F 3$_\pi$. By F 3, $\forall\, \alpha \in A$ we can find $j_\alpha$ in $I$ and $w^0_\alpha : i_\alpha \longrightarrow j_\alpha$, $v^0_\alpha : i \longrightarrow j_\alpha$ such that % [sic] v^0_\alpha : i \longrightarrow j -> v^0_\alpha : i \longrightarrow j_\alpha, the following diagram has v^0_\alpha : i \to j_\alpha
\DERmark{ENV1261}

\[
f(v^0_\alpha)u'_\alpha \ = \ f(w^0_\alpha) \ .
\]
\DERmark{ENV1262}


\printedpage{90}
% head: A. Grothendieck : Dérivateurs, Ch. XVIII, Édition des Universités de Montpellier II et Paris VII
% foot: Version 17 décembre 2025

\[
\begin{tikzcd}[column sep=large, row sep=huge]
i_\alpha \arrow[rr, bend left=20, "w^0_\alpha"] \arrow[d, dash] & i \arrow[r, "v^0_\alpha"] \arrow[d, dash] & j_\alpha \arrow[r, "\lambda_\alpha"] \arrow[d, dash] & j \arrow[d, dash] \\
f(i_\alpha) \arrow[r, "u'_\alpha"] & f(i) \arrow[r, "f(v^0_\alpha)"] & f(j_\alpha) \arrow[r] & f(j)
\end{tikzcd}
\]
\DERmark{ENV1263}


Since $I$ satisfies PF 1$_\pi$ (cf. 4.9.6.9.2), there exist $j$ in $I$ and $\lambda_\alpha : j_\alpha \longrightarrow j$, such that $\lambda_\alpha v^0_\alpha = \lambda$ is independent of $\alpha$, then we have for every $\alpha \in A$ % [sic] h -> j, the arrows \lambda_\alpha go to j
\DERmark{ENV1264}

\[
f(\lambda)u'_\alpha \ \Bigl( \ = \ f(\lambda_\alpha)\underbrace{f(v^0_\alpha)u'_\alpha}_{f(w^0_\alpha)} \ \Bigr) \ = \ f(\underbrace{\lambda_\alpha w^0_\alpha}_{w_\alpha}) \ ,
\]
\DERmark{ENV1265}

so $\lambda : i \longrightarrow j$ and the $w_\alpha = \lambda_\alpha w^0_\alpha : i_\alpha \longrightarrow j$ satisfy the desired conditions.

c$_\pi$ $\Longrightarrow$ b$''_\pi$. We must prove that if c$_\pi$ is satisfied, then the $i'\backslash I$ satisfy PF 1$'_\pi$.

\msp{page 241}
Let then $\xi_\alpha = (i_\alpha, u'_\alpha)$ ($\alpha \in A$, $\card A \leqslant \pi$) be a family of elements of $i'\backslash I$, since $I'$ satisfies PF 1$_\pi$ (cf. 4.9.6.9.2), there exist $j'$ in $I'$ and arrows $u''_\alpha : f(i_\alpha) \longrightarrow j'$, such that $u''_\alpha u'_\alpha = u'$ independent of $\alpha$. Since $f$ satisfies F 1, we can assume that $j'$ is of the form $f(j)$. Applying F 3$_\pi$ to the $i_\alpha$, $j$, $u''_\alpha$, we find $v : j \longrightarrow k$ and $w_\alpha : i_\alpha \longrightarrow k$, such that
\[
f(v)u''_\alpha \ = \ f(w_\alpha) \qquad \forall\, \alpha \in A \ ,
\]
hence
\[
f(w_\alpha)u'_\alpha \ = \ f(v)\,\underbrace{u''_\alpha u'_\alpha}_{u'} \ = \ f(v)u' \qquad \text{independent of $\alpha$,}
\]
say $w'$, so the $w_\alpha$ define an arrow from the $\xi_\alpha$ to $\xi = (k, w')$, q.e.d.

b$''_\pi$ $\Longrightarrow$ b$_\pi$ if $f$ is faithful. We are reduced to proving that if b$''_\pi$ is satisfied, $f$ faithful, then the categories $i'\backslash I$, which we already know satisfy PF 1$'_\pi$, also satisfy PF 2$_\pi$. Let then $\xi = (i, u')$, $\eta = (j, v')$ be two objects of $i'\backslash I$, and $(u_\alpha)_{\alpha\in A}$ a family of arrows $\xi \longrightarrow \eta$ (with $\card A \leqslant \pi$), hence arrows in $I$ such that $f(u_\alpha)u' = v' \ \forall\, \alpha \in A$.
\DERmark{ENV1266}

\[
\begin{tikzcd}[column sep=large]
i \arrow[r, shift left=2, "u_\alpha"] \arrow[r] \arrow[r, shift right=2] & j \arrow[r, dotted, "v"] & k
\end{tikzcd}
\]
\[
\begin{tikzcd}[column sep=large]
i' \arrow[r, "u'"] \arrow[rr, bend right=15, "v'"'] & f(i) \arrow[r, shift left=2] \arrow[r] \arrow[r, shift right=2] & f(j) \arrow[r] & f(k)
\end{tikzcd}
\]

Since $I'$ is large relative to $\pi$, we can equalize by $f(j) \longrightarrow k'$, and since $k'\backslash I \neq \emptyset$, we can assume $k'$ of the form $f(k)$. Now
\msp{page 242}
by the implication b$''$ $\Longrightarrow$ a $\Longrightarrow$ c of part 1°) (cf. page 235) and since b$''_\pi$ $\Longrightarrow$ b$''$, we already know that $f$ satisfies F 3. Applying it to $j$, $k$ and $f(j) \longrightarrow f(k)$, we can
\printedpage{91}
% head: A. Grothendieck : Dérivateurs, Ch. XVIII, Édition des Universités de Montpellier II et Paris VII
% foot: Version 17 décembre 2025
assume that $f(j) \longrightarrow f(k)$ comes from $v : j \longrightarrow k$. We shall therefore have $f(v)f(u_\alpha) = f(v)f(u_\beta)$ $\forall\, \alpha, \beta \in A$, so $f(vu_\alpha) = f(vu_\beta)$. Since $f$ is faithful, this implies that $vu_\alpha = vu_\beta$, i.e. $v$ equalizes the $u_\alpha$. Besides, $v$ can be regarded as an arrow \textit{in} $i'\backslash I$, from $(j, v')$ to $(k, f(v)v')$. Q.e.d.

\textbf{4.9.6.10.} When $I$ is \textit{ordered}, then the four conditions a$_\pi$, b$_\pi$, c$_\pi$, $b''_\pi$ are equivalent, since the $i'\backslash I$ are ordered, and it is clear that then b$''_\pi$ $\Longleftrightarrow$ b$_\pi$. % [sic] claire -> clair, grammar

\msp{page 243}

\section*{4.10\quad The indicial representation theorem: proof}

We are first going to reformulate theorem 4.9.3 in a more general form, in terms of \textit{cofibered} categories. Let
\[
\mathcal{D} \longrightarrow J \tag*{(4.10.1)}
\]
be a cofibrant functor, where $J$ is equivalent to a small category, and likewise for the fiber categories $\mathcal{D}_j$ ($j \in \Ob J$). We give ourselves an infinite cardinal $\pi$. Let
\[
\mathcal{M}_j \ = \ \Ind_\pi(\mathcal{D}_j) \qquad (j \in \Ob J) \ , \tag*{(4.10.2)}
\]
then the $\mathcal{M}_j$ are the fibers of a cofibered category $\mathcal{N}$ over $J$, which I also denote by $\Ind_\pi(\mathcal{D}/J)$, so we put
\[
\mathcal{N} \ = \ \Ind_\pi(\mathcal{D}/J) \ . \tag*{(4.10.3)}
\]
We thus have a fully faithful inclusion of cofibered categories over $J$
\[
\mathcal{D} \hookrightarrow \mathcal{N} \ ,
\]
hence an inclusion
\[
\uGamma(\mathcal{D}/J) \longrightarrow \uGamma(\mathcal{N}/J) \ . \tag*{(4.10.4)}
\]

\msp{page 244}
In the case where $\mathcal{D}$ is of the form $C \times J$, then $\mathcal{N}$ is of the form $\mathcal{M} \times J$, with $\mathcal{M} = \Ind_\pi(C)$, and the inclusion (4.10.4) is none other (up to isomorphism) than
\DERmark{ENV1267}

\[
\underline{\operatorname{Hom}}(J, C) \longrightarrow \underline{\operatorname{Hom}}(J, \mathcal{M}) \ .
\]
Returning to the general case, we also see that:
\begin{enumerate}[label={}, leftmargin=3em]
\item[1°)] The second member of (4.10.4) is stable under $\varinjlim$ large relative to $\pi$, which are computed fiber by fiber.
\item[2°)] The first member of (4.10.4) consists of $\pi$-accessible objects of the second.
\end{enumerate}
\DERmark{ENV1268}\DERmark{ENV1269}

Hence a \textit{fully} faithful functor
\[
\Ind_\pi(\uGamma(\mathcal{D}/J)) \hookrightarrow \uGamma(\mathcal{N}/J) \ . \tag*{(4.10.4)}
\]
\setcounter{footnote}{103}%
(\footnote{\texttt{[The numbers 4.10.1 to 4.10.4 are used twice.]}}).

\printedpage{92}
% head: A. Grothendieck : Dérivateurs, Ch. XVIII, Édition des Universités de Montpellier II et Paris VII
% foot: Version 17 décembre 2025

\textbf{Theorem 4.10.1.} \textit{Suppose that we have}
\[
\mathrm{card}\,\textup{red}\, J \ \leqslant \ \pi \ , \tag*{(4.10.1.1)}
\]
\textit{and that there exists a limit ordinal $\alpha$, such that:}
\begin{enumerate}[label={}, leftmargin=3em]
\item[\textit{1)}] $\card \alpha \leqslant \pi$.
\item[\textit{2)}] \textit{The direct limits of}
\msp{page 245}
\textit{type $I_\alpha = \{\beta \mid \beta < \alpha\}$ are representable in the categories $\mathcal{D}_j$ ($j \in \Ob J$), and for every arrow $\alpha : j \longrightarrow j'$ in $J$ \texttt{[too many $\alpha$]}, $\alpha_* : \mathcal{D}_j \longrightarrow \mathcal{D}_{j'}$ commutes with the said limits.}
\DERmark{ENV1270}

\end{enumerate}
\textit{Then the canonical fully faithful functor (4.10.4) is an equivalence of categories.}

This will be contained in the slightly more precise result:

\textbf{Corollary 4.10.2.} \textit{The hypotheses being those of 4.10.1, with in addition}
\[
\card J \ \leqslant \ \pi \tag*{(4.10.2.1)}
\]
\textit{(which is equivalent to (4.10.1.1) if $J$ is reduced, a case to which we can of course reduce in 4.10.1). Let $\varphi \in \Ob \uGamma(\mathcal{N}/J)$, and let us consider the canonical} faithful \textit{functor:}
\[
\uGamma(\mathcal{D}/J)/\varphi \longrightarrow \prod_{j\in J} \mathcal{D}_j/\varphi(j) \ . \tag*{(4.10.2.2)}
\]
\DERmark{ENV1271}

\textit{Then the source category $\uGamma(\mathcal{D}/J)/\varphi$ is large relative to $\pi$, and the functor in question is cofinal.}

\msp{page 246}
From this, the theorem follows in an essentially sorital way, just as in the preceding special case 4.9.3 (cf. page 223) (\footnote{A bit short\,! The argument ought to be spelled out with care – there are two or three pages' worth of it\,!}).

To establish the corollary, since the target category of (4.10.2.2) is large relative to $\pi$, as a product of such categories (cf. 4.9.6.6, page 228), and the functor faithful, it suffices by virtue of 4.9.6.3, 2°, to verify conditions F 1 and F 3$_\pi$ (page 225) for the functor in question.

We are going to interpret the functor (4.10.2.2) in terms of a cofibered category
\[
\mathcal{D}_\varphi \longrightarrow J \ , \tag*{(4.10.5)}
\]
whose fibers are the $\mathcal{D}_j/\varphi(j)$. We note that if $\alpha : j \longrightarrow j'$ is an arrow of $J$, then we have
\[
\left\{
\begin{aligned}
&\alpha_* : \mathcal{D}_j \ \longrightarrow \ \mathcal{D}_{j'} \\
&\text{and} \ \ \alpha_*(\varphi(j)) \ \longrightarrow \ \varphi(j') \ ,
\end{aligned}
\right.
\]
\DERmark{ENV1272}


\printedpage{93}
% head: A. Grothendieck : Dérivateurs, Ch. XVIII, Édition des Universités de Montpellier II et Paris VII
% foot: Version 17 décembre 2025

\msp{page 247}
hence a composite functor
\[
\mathcal{D}_j/\varphi(j) \ \longrightarrow \ \mathcal{D}_{j'}/\alpha_*(\varphi(j)) \ \longrightarrow \ \mathcal{D}_{j'}/\varphi(j') \ ,
\]
\DERmark{ENV1273}

with transitivity up to canonical isomorphism etc. We note that we have a canonical isomorphism
\[
\uGamma(\mathcal{D}_\varphi/J) \ \simeq \ \uGamma(\mathcal{D}/J)/\varphi \ ,
\]
so that the faithful canonical functor (4.10.2.2) is identified with the functor
\[
\uGamma(\mathcal{D}_\varphi/J) \longrightarrow \prod_{j\in\Ob J} (\mathcal{D}_\varphi)_j \ . \tag*{(4.10.6)}
\]

Corollary 4.10.2 will then be a consequence of the following more general result:

\textbf{Corollary 4.10.3.} \textit{Let $\mathcal{E} \longrightarrow J$ be a cofibrant functor, and $\pi$ an infinite cardinal. We assume:}
\begin{enumerate}[label={}, leftmargin=3em]
\item[\textit{a)}] $\card J \leqslant \pi$.
\item[\textit{b)}] \textit{The fiber categories $\mathcal{E}_j$ are equivalent to small categories, and large relative to $\pi$.}
\item[\textit{c)}] \textit{There exists a limit ordinal $\alpha$ such that:}
\begin{enumerate}[label={}, leftmargin=3em]
\item[\textit{1°)}] $\card \alpha \leqslant \pi$.
\item[\textit{2°)}] \msp{page 248}\textit{The direct limits of type $I_\alpha$ are representable in the categories $\mathcal{E}_j$, and the cobase change functors $\lambda_* : \mathcal{E}_j \longrightarrow \mathcal{E}_{j'}$ commute with them. (This is the case e.g. if the} countable \textit{filtering direct limits are representable in the $\mathcal{E}_j$, and if the $\alpha_*$ commute with them\,; or if there exists a cardinal $\pi_0 < \pi$ and a cardinal $\pi'_0 > \pi_0$, such that the direct limits large relative to $\pi_0$ and of cardinality $\leqslant \pi'_0$ are representable in the $\mathcal{E}_j$, and the $\alpha_*$ commute with them.)}
\end{enumerate}
\DERmark{ENV1274}\DERmark{ENV1275}

\end{enumerate}
\textit{Under these conditions, the faithful canonical functor}
\[
\uGamma(\mathcal{E}/J) \longrightarrow \prod_{j\in\Ob J} \mathcal{E}_j \tag*{(4.10.3.1)}
\]
\textit{is cofinal, and the source category is large relative to $\pi$.}

It is $\pm$ immediate that the hypotheses of 4.10.1 imply that those of 4.10.3 are satisfied for $\mathcal{E} = \mathcal{D}_\varphi \longrightarrow J$ (4.10.5). For
\msp{page 249}
this, we use the

% marginal vertical bar printed in the left margin beside the following statement
\textbf{Lemma 4.10.4.} \textit{Let $C$ be a category, and $\pi$ an infinite cardinal. Then the inclusion functor $C \longrightarrow \mathcal{M} \overset{\mathrm{def}}{=} \Ind_\pi C$ commutes with $\varinjlim$ of cardinality $\leqslant \pi$.}
\DERmark{ENV1276}


Indeed, let $I$ be such that $\card I \leqslant \pi$, and $(x_i)_{i\in I}$ a direct system in $C$, having a limit $x$ there. On the other hand, let $X = \varinjlim\limits_{I'} \xi_{i'}$ be an object of $\Ind_\pi C$, with $I'$ large relative to $\pi$. We have


\endgroup
% END UNIT der-ch18/der-ch18_p091-095.tex
% BEGIN UNIT der-ch18/der-ch18_p096-100.tex
\begingroup\makeatletter%
\newcommand{\printedpage}[1]{\marginpar{\footnotesize p.~#1}}%
\newcommand{\msp}[1]{\marginpar{\footnotesize [#1]}}%
\newcommand{\uGamma}{\underline{\Gamma}}%
\newcommand{\Hom}{\operatorname{Hom}}%
\newcommand{\Ob}{\mathrm{Ob}\,}%
\newcommand{\card}{\mathrm{card}\,}%
\makeatother\let\printedpage\derpp\let\msp\dermsp%
\let\mspin\DERmspin%


\printedpage{94}
% head: A. Grothendieck : Dérivateurs, Ch. XVIII, Édition des Universités de Montpellier II et Paris VII
thus
\[
\begin{aligned}
\Hom_{\mathcal{M}}(x,X) &\simeq \varinjlim_{I'} \Hom(x,\xi_{i'})\\
&\simeq \varinjlim_{I'} \varprojlim_{I} \Hom(x_i,\xi_{i'})\\
&\overset{\text{SGA 4 I 9.8}}{\simeq} \varprojlim_{I} \underbrace{\varinjlim_{I'} \Hom(x_i,\xi_{i'})}_{\simeq \Hom(x_i,X)}\\
&\simeq \varprojlim_{I} \Hom(x_i,X)\ ,
\end{aligned}
\]
which proves that $x$ is also a $\varinjlim$ of the $x_i$ in the whole of $\mathcal{M}$.

\msp{250}
Thus, we are reduced to proving 4.10.3. Applying 4.9.6.3, 2°, we are brought back to proving, for the functor under consideration, conditions F 1 and F $3_\pi$ (page 225).
\DERmark{ENV1277}


In the arguments that follow, we are going to use conditions a), b) of 4.10.3, but not, for the moment, condition c.

\textbf{Lemma 4.10.5.} \textit{Under the conditions of 4.10.3 a) b) (excluding c), for (4.10.3.1) to be cofinal and $\Gamma(\mathcal{E}/J)$ large relative to $\pi$, it is necessary and sufficient that the functor under consideration satisfy condition F 1, recalled below, and its variant F$'$ 1 stated below:}
\begin{itemize}[leftmargin=3em]
\item[F\,1)] {\itshape For every system $(X(j))_{j\in\Ob J}$ of elements $X(j)\in\Ob\mathcal{E}_j$, exists a \textup{section} $G$ of $\mathcal{E}$ over $J$, and a system of arrows
\[
X(j) \longrightarrow G(j) \qquad (j\in ObJ)\ . \tag*{\textup{(4.10.5.1)}}
\]}
\item[F$'$1)] {\itshape Under the conditions of the preceding \textup{F\,1}, suppose in addition a section $F$ of $\mathcal{E}$ over $J$ \msp{251} and a system of arrows
\[
F(j) \longrightarrow X(j) \qquad (j\in\Ob J)\ ,
\]
then $G$ and the system of arrows (4.10.5.1) can be chosen in such a way that the composite arrows
\[
F(j) \longrightarrow X(j) \longrightarrow G(j)
\]
define a homomorphism of sections
\[
F \longrightarrow G
\]
(in the category $\uGamma(\mathcal{E}/J)$).}
\end{itemize}
(\footnote[106]{\textbf{N.B.} Conditions F 1 and F$\,'$\,1 are equivalent to saying that the functor under consideration satisfies conditions F1 and F3, seeing also (by virtue of 4.9.6.3, 1°) that $\uGamma(\mathcal{E}/J)$ is \textit{filtering} without more, and the faithful functor (4.10.3.1) cofinal. This already implies that $\uGamma(\mathcal{E}/J)$ is large relative to $\pi$.}).

\textsc{Proof}. We have to see that F 1 \& F$'$ 1 is equivalent to F 1 \& F $3_\pi$. The implication $\Longleftarrow$ is $\pm$ immediate, it is $\Longrightarrow$ that we shall need, so
\[
\text{F 1 \& F}'\text{ 1} \quad\Longrightarrow\quad \text{F } 3_\pi\ .
\]
% foot: Version 17 décembre 2025

\printedpage{95}
% head: A. Grothendieck : Dérivateurs, Ch. XVIII, Édition des Universités de Montpellier II et Paris VII
For this, let us restate F $3_\pi$ in the present context. We are given a family $(F_\alpha)_{\alpha\in A}$, $\card A\leqslant\pi$, of objects of $\uGamma(\mathcal{E}/J)$, an object $F$ of the same category, and for every $j$ in $\Ob J$, a system of arrows
\[
F_\alpha(j) \xrightarrow{u_\alpha(j)} F(j)\ ,
\]
but without supposing that for $\alpha$ fixed, $j$ variable

\msp{252}
this defines a homomorphism of sections $F_\alpha\longrightarrow F$. We seek a section $G$ of $\mathcal{E}$ over $J$, and a homomorphism of sections $F\longrightarrow G$, such that for every $\alpha$, the system of composites
\[
F_\alpha(j) \longrightarrow F(j) \longrightarrow G(j)
\]
defines a homomorphism of sections $F_\alpha\longrightarrow G$. So let $\lambda : j\longrightarrow j'$ be an arrow of $J$, let us consider the squares (\textit{noncommutative}) % [sic] F_\alpha -> G_\alpha -> F_\alpha -> G, the section sought is G (cf. p. 95 below: F_\alpha -> G)
\DERmark{ENV1278}

\[
\begin{tikzcd}
F_\alpha(j) \arrow[r,"u_\alpha(j)"] \arrow[d,"F_\alpha(\lambda)"'] \arrow[dr,phantom,"!"] & F(j) \arrow[d,"F(\lambda)"] \\
F_\alpha(j') \arrow[r,"u_\alpha(j')"] & F(j')
\end{tikzcd}
\tag*{$(*)$}
\]
(indexed by $\alpha\in A$). For fixed $j'$ in $J$, the number of these squares is bounded above by $\card\Ob(J/j')\times\card A\leqslant\pi\cdot\pi=\pi$. It follows, since $\mathcal{E}_{j'}$ is large relative to $\pi$, that there exists an object $X(j')$ of $\mathcal{E}_{j'}$, and an arrow $v(j') : F(j')\longrightarrow X(j')$ in such a way that for every $\alpha\in A$, and every $\lambda\in\Ob J/j'$, the square deduced from $(*)$

\msp{253}
by composing with $F(j')\longrightarrow X(j')$ is commutative. We make this choice for every object $j'$ of $J$. Let us then apply F$'$ 1 to $F$ and to the system of arrows $F(j)\xrightarrow{v(j)}X(j)$, we find $G$ and a system of arrows
\[
w(j)\ :\ X(j) \longrightarrow G(j)
\]
in such a way that the composite arrows
\[
\begin{tikzcd}[column sep=small]
w(j)v(j) : & F(j) \arrow[rr] \arrow[dr,dotted] & & G(j) \\
& & X(j) \arrow[ur,dotted] &
\end{tikzcd}
\]
define a homomorphism of sections. I say that this homomorphism has the desired property for F $3_\pi$, i.e. that for every $\alpha\in A$, the system of composites
\[
F_\alpha(j) \longrightarrow F(j) \longrightarrow G(j)
\]
(for $j\in\Ob J$ variable) defines a homomorphism of sections $F_\alpha\longrightarrow G$. Indeed, for this it suffices to consider, for an arrow $\lambda$ of $J$, the
% foot: Version 17 décembre 2025
\printedpage{96}
% head: A. Grothendieck : Dérivateurs, Ch. XVIII, Édition des Universités de Montpellier II et Paris VII
\msp{254}
diagram
\[
\begin{tikzcd}
F_\alpha(j) \arrow[r,"u_\alpha(j)"] \arrow[d,"F_\alpha(\lambda)"'] \arrow[dr,phantom,"!"] & F(j) \arrow[r,"v(j)"] \arrow[d,"F(\lambda)"'] & X(j) \arrow[r,"w(j)"] \arrow[d,phantom,"\text{comm.}"] & G(j) \arrow[d,"G(\lambda)"] \\
F_\alpha(j') \arrow[r,"u_\alpha(j')"] & F(j') \arrow[r,"v(j')"] & X(j') \arrow[r,"w(j')"] & G(j')\ .
\end{tikzcd}
\]
Lemma 4.10.5 being proved, everything thus comes down to proving condition F 1, and its complement F$'$ 1, which will be proved in the same breath.

\textbf{4.10.6.} So we are given the system of objects
\[
X(j)\ \in\ \Ob\mathcal{E}_j \qquad (j\in\Ob J)\ , \tag*{(4.10.6.1)}
\]
and as the case may be, an object $F$ of $\uGamma(\mathcal{E}/J)$, and a system of arrows
\[
F(j) \xrightarrow{u_j} X(j)\ , \tag*{(4.10.6.2)} % [sic] (4.6.10.2) -> (4.10.6.2), numbering of 4.10.6; cited as (4.10.6.2) below
\]
\DERmark{ENV1279}

and we propose to construct $G$ in $\uGamma(\mathcal{E}/J)$, and a system of arrows
\[
X(j) \xrightarrow{w(j)} G(j)\ , \tag*{(4.10.6.3)}
\]
in such a way that (if we have the additional datum of the $u(j)$) the composite arrows
\[
F(j) \xrightarrow{u(j)} X(j) \xrightarrow{w(j)} G(j) % [sic] v(j) -> w(j), the arrow X(j) -> G(j) is w(j) in (4.10.6.3)
\]
\DERmark{ENV1280}


\msp{255}
are functorial in $j$.

We are going to make a transfinite construction on the ordinals $\alpha$ such that
\[
\card\alpha\ \leqslant\ \pi
\]
of a direct system
\[
(X_\alpha(j))_\alpha
\]
of objects of $\prod_{j\in\Ob J}\mathcal{E}_j$. We are first going to ignore the possible datum (4.10.6.2). We shall put
\[
X_0(j)\ =\ X(j)\ , \tag*{(4.10.6.4)}
\]
and we are going to give the prescription for passing from $X_\alpha$ to $X_{\alpha+1}$. The construction will also include, for every arrow $\lambda : j'\longrightarrow j$ of $J$ and every $\alpha$, the datum of
\[
X_\alpha(j') \xrightarrow{t_\alpha(\lambda)} X_{\alpha+1}(j)\ , \qquad \text{arrow (in $\mathcal{E}$) over $\lambda$,} \tag*{(4.10.6.5)} % [sic] X_\alpha(j) -> X_\alpha(j'), the source of an arrow over \lambda : j' -> j (cf. p. 97: X_\alpha(j') -> X_{\alpha+1}(j))
\]
\DERmark{ENV1281}

and these arrows, and the transition arrows $X_\alpha(j)\xrightarrow{\tau_\alpha(j)}X_{\alpha+1}(j)$, must
% foot: Version 17 décembre 2025
\printedpage{97}
% head: A. Grothendieck : Dérivateurs, Ch. XVIII, Édition des Universités de Montpellier II et Paris VII
\msp{256}
satisfy the following three conditions:
\begin{itemize}[leftmargin=3em]
\item[(i)] If $\lambda$ is an identity arrow $\mathrm{id}_j : j\longrightarrow j$, then \fbox{$t_\alpha(\lambda)=\tau_\alpha(j)$.} % [sic] id_J : J -> J -> id_j : j -> j, the identity arrow of an object j of J
\item[(ii)] In the diagram
\[
\begin{tikzcd}
X_\alpha(j') \arrow[r,"\tau_\alpha(j')"] \arrow[dr,"t_\alpha(\lambda)"] & X_{\alpha+1}(j') \arrow[r] \arrow[dr,"t_{\alpha+1}(\lambda)"] \arrow[d,phantom,"\bullet"] & X_{\alpha+2}(j') \arrow[r] \arrow[dr] & X_{\alpha+3}(j') \arrow[r] \arrow[dr,dotted] & {} \\
X_\alpha(j) \arrow[r] & X_{\alpha+1}(j) \arrow[r,"\tau_{\alpha+1}(j)"] & X_{\alpha+2}(j) \arrow[r,"\tau_{\alpha+2}(j)"] & X_{\alpha+3}(j) \arrow[r] & {,}
\end{tikzcd}
\] % [sic] bottom row printed with j' (X_\alpha(j'), X_{\alpha+1}(j'), \tau_{\alpha+1}(j'), X_{\alpha+2}(j'), \tau_{\alpha+2}(j'), X_{\alpha+3}(j')) -> j, the arrows t(\lambda) go to the objects at j (cf. the text below: \tau_{\alpha+1}(j), \tau_{\alpha+2}(j))
\DERmark{ENV1282}\DERmark{ENV1283}

$\tau_{\alpha+2}(j)$ ``makes the square $\bullet$ commutative'', i.e. equalizes $t_{\alpha+1}(\lambda)\tau_\alpha(j')$ and $\tau_{\alpha+1}(j)t_\alpha(\lambda)$:
\[
\boxed{\tau_{\alpha+2}(j)t_{\alpha+1}(\lambda)\tau_\alpha(j')\ =\ \tau_{\alpha+2}(j)\tau_{\alpha+1}(j)t_\alpha(\lambda)\ .}
\]
\item[(iii)] If we have two consecutive arrows $\lambda : j''\longrightarrow j'$, $\mu : j'\longrightarrow j$, we thus have the diagram
\[
\begin{tikzcd}
X_\alpha(j'') \arrow[r] \arrow[dr,"t_\alpha(\lambda)"] \arrow[ddr,"t_\alpha(\mu\lambda)"] & X_{\alpha+1}(j'') \arrow[r] \arrow[dr] & X_{\alpha+2}(j'') \arrow[r] \arrow[dr] & X_{\alpha+3}(j'') \\
X_\alpha(j') \arrow[r] & X_{\alpha+1}(j') \arrow[r] \arrow[dr,"t_{\alpha+1}(\mu)"] & X_{\alpha+2}(j') \arrow[r] \arrow[dr] & X_{\alpha+3}(j') \\
X_\alpha(j) \arrow[r] & X_{\alpha+1}(j) \arrow[r,"\tau_{\alpha+1}(j)"] & X_{\alpha+2}(j) \arrow[r,"\tau_{\alpha+2}(j)"] & X_{\alpha+3}(j)\ ,
\end{tikzcd}
\]
\[
\boxed{\tau_{\alpha+2}(j)t_{\alpha+1}(\mu)t_\alpha(\lambda)\ =\ \tau_{\alpha+2}(j)\tau_{\alpha+1}(j)t_\alpha(\mu\lambda)}
\]
\end{itemize}

\msp{257}
Thus, to define the $X_{\alpha+1}(j)$, we have to take account of all these constraints. So:
\begin{itemize}[leftmargin=3em]
\item[1°)] For fixed $j$, $X_{\alpha+1}(j)$ must be the target of the arrows
\[
X_\alpha(j') \xrightarrow{t_\alpha(\lambda)} X_{\alpha+1}(j)\ ,
\]
for the $\lambda\ :\ j'\longrightarrow j$ in $\Ob J/j$ -- including for $\lambda=\mathrm{id}_j$, in which case $t_\alpha(\lambda)\ :\ X_\alpha(j)\longrightarrow X_{\alpha+1}(j)$ will play the role of the transition morphism $\tau_\alpha(j)$.
\item[2°)] If $\alpha$ is of the form $\alpha'+2$, i.e. if it is neither a limit ordinal, nor the successor of a limit ordinal, then we require in addition, for the transition morphisms $\tau_\alpha(j)$, that they equalize the pairs of arrows provided for in (ii) and (iii) above.
\end{itemize}
Since the set of these data in 1°), or of these constraints in 2°),
\msp{258}
has cardinality $\leqslant\pi$, and the category $\mathcal{E}_j$ is large relative to $\pi$, we see that we can indeed find $X_{\alpha+1}(j)$ and $t_\alpha(\lambda)\ :\ X_\alpha(j')\longrightarrow X_{\alpha+1}(j)$, so as to satisfy (if $\alpha=\alpha'{+}2$) the constraints recalled in 2°. This completes the description of the construction of $X_{\alpha+1}(\bullet)$, in terms of $X_\beta(\bullet)$ $(\beta\leqslant\alpha)$ and the $t_\beta(\lambda)$ $(\beta<\alpha)$.

We must finally define $X_\alpha(\bullet)$ for $\alpha$ a limit ordinal. This time, we ask it only to receive the $X_\beta(\bullet)$ for $\beta<\alpha$. Since the set $I_\alpha$ of these $\beta$ has cardinality $\leqslant\pi$, and the $\mathcal{E}_j$ are large relative to $\pi$, we can indeed find the $X_\alpha(j)$.
% foot: Version 17 décembre 2025

\printedpage{98}
% head: A. Grothendieck : Dérivateurs, Ch. XVIII, Édition des Universités de Montpellier II et Paris VII
\msp{259}
We have thus proved the essential part of

\textbf{Lemma 4.10.7.} \textit{Under the conditions of 4.10.3 a) b), and for a given system $X(\bullet)\in\Ob\prod_{j\in\Ob J}\mathcal{E}_j$, we can find a direct system $(X_\alpha(j))_\alpha$ in $\prod\mathcal{E}_j$, indexed by the ordinals $\alpha$ such that} $\card\alpha\leqslant\pi$, \textit{and for every arrow $\lambda : j'\longrightarrow j$ in $J$ arrows $t_\alpha(\lambda) : X_\alpha(j')\longrightarrow X_{\alpha+1}(j)$, in such a way that we have the following conditions.} % [sic] (X_\alpha(i))_\alpha -> (X_\alpha(j))_\alpha, the index is j \in Ob J; [sic] X_\alpha(j') -> X_\alpha(j) -> X_\alpha(j') -> X_{\alpha+1}(j), cf. (4.10.6.5) and p. 97
\DERmark{ENV1284}\DERmark{ENV1285}

\begin{itemize}[leftmargin=3em]
\item[\textit{1°)}] \textit{$X_0(\bullet)=X(\bullet)$.}
\item[\textit{2°)}] \textit{The conditions (i) (ii) (iii) (p. 256).}
\item[\textit{3°)}] \textit{If we are given $F$ in \textup{Ob}$\,\uGamma(\mathcal{E}/J)$, and arrows $u(j) : F(j)\longrightarrow X(j)=X_0(j)$, we may suppose that for every arrow $\lambda : j'\longrightarrow j$ in $J$, we have}
\[
\begin{tikzcd}
F(j') \arrow[r,"u(j')"] \arrow[d,"F(\lambda)"'] & X_0(j') \arrow[dr,"t_0(\lambda)"] & \\
F(j) \arrow[r,"u(j)"] & X_0(j) \arrow[r,"\tau_0(j)"] & X_1(j)
\end{tikzcd}
\qquad\qquad \text{commutative.}
\tag*{$(*)$}
\]
\end{itemize}

\msp{260}
It only remains to prove this last complement 3°. This is a constraint only on the construction of $X_1(\bullet)$ and of the $t_0(\lambda)$. So suppose that we have constructed $X_1(\bullet)$ without taking this constraint into account. For fixed $j$, we must therefore ``make commutative'' all these squares, i.e. all these double arrows $(t_0(\lambda)u(j'),\tau_0(j)u(j)F(\lambda))$ with target $X_1(j)$. Since the set of these double arrows has cardinality $\card\Ob(J/j)\leqslant\pi$, we can find $X_1(j)\longrightarrow X'_1(j)$ in $\mathcal{E}_j$, equalizing all these double arrows, so at the cost of replacing $X_1(j)$ by $X'_1(j)$, we win.
\DERmark{ENV1286}


\msp{261}
It is finally time to appeal to hypothesis c) of 4.10.3. So let $\alpha$ be a limit ordinal, as stipulated in this condition. Since $\card\alpha\leqslant\pi$, a fortiori $\card\beta\leqslant\pi$ if $\beta<\alpha$, we thus have a direct system
\[
(X_\beta(\bullet))_{\beta\in I_\alpha}
\]
in $\prod_{j\in J}\mathcal{E}_j$. Since the $\varinjlim$ of type $I_\alpha$ exist in $\mathcal{E}_j$, we can therefore consider
\[
G(j)\ =\ \varinjlim_{\beta\in I_\alpha} X_\beta(j)\ \in\ \Ob\mathcal{E}_j\ .
\]
If $\lambda : j'\longrightarrow j$ is an arrow of $J$, we deduce from it a $\lambda$-arrow
\[
G(\lambda)\ :\ G(j')=\varinjlim X_\beta(j') \longrightarrow G(j)\simeq\varinjlim X_{\beta+1}(j)
\]
\DERmark{ENV1287}

by $G(\lambda)=\varinjlim t_\beta(\lambda)$. We must
\msp{262}
take note that the $t_\beta(\lambda)$ do not define a homomorphism of direct systems, but by condition (ii), p. 256, they define a homomorphism of ind-objects, so
\DERmark{ENV1288}

% foot: Version 17 décembre 2025


\endgroup
% END UNIT der-ch18/der-ch18_p096-100.tex
% BEGIN UNIT der-ch18/der-ch18_p101-105.tex
\begingroup\makeatletter%
\newcommand{\printedpage}[1]{\marginpar{\footnotesize p.~#1}}%
\newcommand{\msp}[1]{\marginpar{\footnotesize [#1]}}%
\newcommand{\cercle}[1]{\tikz[baseline=(X.base)]\node[draw,circle,inner sep=1pt](X){\upshape #1};}%
\newcommand{\defeq}{\overset{\mathrm{def}}{=}}%
\newcommand{\cardred}{\operatorname{card}\,\text{\upshape red}\,}%
\newcommand{\Hom}{\underline{\operatorname{Hom}}}%
\newcommand{\uGamma}{\underline{\Gamma}}%
\newcommand{\Ind}{\operatorname{Ind}}%
\newcommand{\into}{\lhook\joinrel\longrightarrow}%
\makeatother\let\printedpage\derpp\let\msp\dermsp%
\let\mspin\DERmspin%

\DERmark{ENV1289}


\printedpage{99}
% head: A. Grothendieck : Dérivateurs, Ch. XVIII, Édition des Universités de Montpellier II et Paris VII
% foot: Version 17 décembre 2025

a homomorphism on the $\varinjlim$ (\footnote{Careful, it is here that we tacitly use the fact that the functor $\lambda_* : \mathcal{E}_{j'} \longrightarrow \mathcal{E}_j$ commutes with $\varinjlim$ of type $I_\alpha$, which implies that the inclusions $\mathcal{E}_j \into \mathcal{E}$ commute with this type of $\varinjlim$ .}). Condition (i) ensures that $\lambda \longmapsto G(\lambda)$ transforms identities into identities, and condition (iii) that we have $G(\mu\lambda) = G(\mu)G(\lambda)$. Thus $G$ appears as a section of $\mathcal{E}$ over $J$. On the other hand, we have arrows
\[
X(j) \defeq X_0(j) \longrightarrow G(j) = \varinjlim_{\beta\in I_\alpha} X_\beta(j) \ ,
\]
which proves F (1). On the other hand, condition 3°) of the lemma proves the validity of F$'$ (1). This completes the proof of 4.10.3, and thereby of 4.10.1 and of 4.9.3.

\msp{page 263}

\textbf{Corollary 4.10.8} (of theorem 4.10.1). \textit{Let $\mathcal{N} \longrightarrow J$ be a cofibrant functor, such that the fibers are accessible categories (4.8), the cobase change functors $\lambda_*$ : $\mathcal{N}_{j'} \longrightarrow \mathcal{N}_j$ accessible, finally $J$ equivalent to a small category. Then the category $\mathcal{M} = \uGamma(\mathcal{N}/J)$ is accessible. More precisely, let $\pi_0$ be an infinite cardinal such that the categories $\mathcal{N}_j$ are $\pi_0$-accessible, and that the functors $\lambda_*$ are $\pi_0$-accessible, and let $\pi$ be an infinite cardinal such that we have:}
\begin{enumerate}[leftmargin=*]
\item[\cercle{a}] $\pi^{\pi_0} = \pi$.
\item[\cercle{b}] $\forall\ \lambda : j' \longrightarrow j$ \textit{in} $J$, $\lambda_* : \mathcal{N}_{j'} \longrightarrow \mathcal{N}_j$ \textit{maps} $(\mathcal{N}_{j'})_{\pi_0}$ \textit{into} $(\mathcal{N}_j)_{\pi_0}$.
\item[\cercle{c}] $\cardred J \leqslant \pi$.
\end{enumerate}
\DERmark{ENV1290}

\textit{Then $\mathcal{M}$ is $\pi$-accessible, and more precisely, we have}
\[
\mathcal{M} \ \simeq \ \Ind_\pi(\mathcal{M}_\pi) \ ,
\]
\textit{where $\mathcal{M}_\pi$ is the full subcategory of $\mathcal{M}$ consisting of the sections $F$ of $\mathcal{N}$ over $J$ such that we have $F(j) \in (\mathcal{N}_j)_\pi \ \forall\ j \in \operatorname{Ob} J$,} \msp{page 264} \textit{$(\mathcal{N}_j)_\pi$ denoting the full subcategory of $\mathcal{N}_j$ consisting of the $\pi$-accessible objects of $\mathcal{N}_j$.}
\DERmark{ENV1291}\DERmark{ENV1292}\DERmark{ENV1293}


Indeed, the hypothesis at the beginning implies the existence of $\pi_0$ and of $\pi$ satisfying the stated conditions. Conditions a) and b) on $\pi$, plus the fact that the functors $\lambda_*$ are $\pi_0$-accessible, imply that the strictly full subcategory $\mathcal{N}_\pi$ of $\mathcal{N}$ whose fibers are the $(\mathcal{N}_j)_\pi$, is a cofibered subcategory, and that moreover $\mathcal{N}$ is $J$-equivalent to $\Ind_\pi(\mathcal{N}_\pi/J)$. On the other hand, the fibers of $\mathcal{N}_\pi/J$ are stable under $\varinjlim$ large relative to $\pi_0$ and of cardinality $\leqslant \pi$, so the ``ordinal'' condition of 4.10.1 is satisfied (taking $\mathcal{D} = \mathcal{N}_\pi$ there). Since moreover $\operatorname{card} J \leqslant \pi$, we conclude by 4.10.1 that $\mathcal{M} \simeq \Ind_\pi(\mathcal{M}_\pi)$.

\msp{page 265}

\textbf{Corollary 4.10.9.} \textit{Let $\mathcal{N} \longrightarrow J$ be a ``cofibration in paratopoi'' (\footnote{ambiguous term, cf. the double explanation.}), i.e. a cofibrant functor such that the fibers are paratopos categories, and the cobase change functors are inverse image functors (resp. direct image functors) for morphisms of paratopoi, i.e. they commute with (small) $\varinjlim$ (resp. with small $\varprojlim$ ). Assume furthermore $J$ equivalent to a small category. Then $\uGamma(\mathcal{N}/J)$ is a paratopos category.}

Indeed, it is an accessible category by virtue of 4.10.8, on the other hand the hypothesis on the $\lambda_*$ implies at once that $\uGamma(\mathcal{N}/J)$ is stable under small $\varprojlim$ (resp. under small $\varprojlim$
\printedpage{100}
% head: A. Grothendieck : Dérivateurs, Ch. XVIII, Édition des Universités de Montpellier II et Paris VII
% foot: Version 17 décembre 2025
\texttt{[rather} $\varinjlim$ \texttt{?]}). Thus $\uGamma(\mathcal{N}/J)$ is a paratopos category by virtue of 4.8.9 b $\Longrightarrow$ a (resp. c $\Longrightarrow$ a).

\textbf{Corollary 4.10.10.} \textit{Let $\mathcal{N} \longrightarrow J$ be a category that is} fibered \textit{(not cofibered!), such that the fibers are paratopos categories and the functors} \msp{page 266} \textit{of base change are direct image functors for morphisms of paratopoi (i.e. accessible functors, commuting with small $\varprojlim$ ). Assume furthermore $J$ equivalent to a small category. Then $\uGamma(\mathcal{N}/J)$ is a paratopos category.}

Indeed, the hypothesis implies that for every arrow $\lambda : j' \longrightarrow j$ in $J$, the functor $\lambda^*$ : $\mathcal{N}_j \longrightarrow \mathcal{N}_{j'}$ admits a left adjoint $\lambda_! : \mathcal{N}_{j'} \longrightarrow \mathcal{N}_j$, whence it follows that $\mathcal{N} \longrightarrow J$ is also cofibrant, the cobase change functors commuting with small $\varinjlim$ . Thus we can apply 4.10.9, which gives the desired result.
\DERmark{ENV1294}\DERmark{ENV1295}


\textbf{Question 4.10.11.} \textit{Let $\mathcal{N} \longrightarrow J$ be a functor that is} fibrant, \textit{such that the $\mathcal{N}_j$ are accessible categories, and the base change functors $\lambda^*$ accessible, $J$ being equivalent to a small category. I} \msp{page 267} \textit{conjecture that under these conditions, $\uGamma(\mathcal{N}/J) \defeq \Hom_J(J,\mathcal{N})$ is accessible, and that the same is true of $\Hom\,\underline{\operatorname{cart}}_J(J,\mathcal{N})$.}
\DERmark{ENV1296}


This is what is proved in SGA 4 I 9.24 under slightly stronger hypotheses (with the help of the old notion too strong for my taste, of cardinal filtrations). I would be surprised if the (somewhat longish) proof of loc. cit. did not apply to the general case. But we would like to have a statement that covers this one and 4.10.8, in the case of a functor $\mathcal{N} \longrightarrow J$ that would be assumed neither fibrant nor cofibrant, but \ldots(???).

\section*{4.11\quad Application to the existence of limits in a category of the form $\Ind_{\pi_0}(C)$.}

\textbf{Proposition 4.11.1.} \textit{Let} \msp{page 268} \textit{$C$ be a category equivalent to a small category, $\pi$ an infinite cardinal, put $\mathcal{M} = \Ind_\pi(C)$. Let $\Phi$ be a set of categories of reduced cardinality $\leqslant \pi$. We propose to examine whether $\mathcal{M}$ is ``stable under $\varinjlim$ of type $\Phi$'' (i.e. stable under $\varinjlim$ of type $I$, whatever $I \in \Phi$ may be). It suffices for this that $C$ is stable under $\varinjlim$ of type $\Phi$, and that moreover there exists a limit ordinal $\alpha$ with $\operatorname{card}\alpha \leqslant \pi$ and such that $C$ is stable under $\varinjlim$ of type $I_\alpha$ (e.g. such that $I_\alpha \in \Phi$). Conversely, if $\mathcal{M}$ is stable under $\varinjlim$ of type $\Phi$, and if we assume moreover $C$ Karoubian (i.e. stable under direct factors), then $C$ is stable under $\varinjlim$ of type $\Phi$. In either case, the inclusion functor $C \into \mathcal{M}$ commutes with the $\varinjlim$ under consideration.}

(\footnote{Make explicit here the corollaries of types 4.13.20, 4.13.21 (p. 341).}).

\textsc{Proof}. a) Sufficient condition. We apply 4.9.3 to conclude that
\printedpage{101}
% head: A. Grothendieck : Dérivateurs, Ch. XVIII, Édition des Universités de Montpellier II et Paris VII
% foot: Version 17 décembre 2025
\msp{page 269}
for every category $J$ of reduced cardinality $\leqslant \pi$ (in particular for $J \in \Phi$), every functor $\varphi : J \longrightarrow \mathcal{M} = \Ind_\pi(C)$ admits an indicial representation $\Phi : J \times I \longrightarrow C$, with $I$ large relative to $\pi$. This, plus the stability of $C$ under $\varinjlim$ of type $J$, suffices to establish the existence of the $\varinjlim$ of $\varphi$.

The converse comes from the fact that if $C$ is Karoubian, then its essential image in $\mathcal{M}$ is none other than $\mathcal{M}_\pi$, which is stable in $\mathcal{M}$ under any $\varinjlim$ of cardinality $\leqslant \pi$ representable in $\mathcal{M}$. If $\mathcal{M}$ is stable under $\varinjlim$ of type $\Phi$, the same is true of $\mathcal{M}_\pi$, hence of $C$, q.e.d.

\textbf{Corollary 4.11.2.} \textit{For $\mathcal{M}$ to be stable under filtering $\varinjlim$ (resp. under small $\varinjlim$ ), it suffices that $C$ be so under filtering $\varinjlim$ (resp. under $\varinjlim$ )} of cardinality $\leqslant \pi$, \textit{and this is also necessary if $C$ is Karoubian.}

Suppose $C$ stable under filtering $\varinjlim$ \msp{page 270} (resp. under $\varinjlim$ of cardinality $\leqslant \pi$), then by 4.11.1 $\mathcal{M}$ is stable under this same type of direct limits. On the other hand, it is stable under $\varinjlim$ large relative to $\pi$. It thus suffices to prove the
\DERmark{ENV1297}


\textbf{Lemma 4.11.3.} \textit{Let $\mathcal{M}$ be a category, $\pi$ a cardinal. Suppose that $\mathcal{M}$ is stable under filtering $\varinjlim$ of cardinality $\leqslant \pi$ (resp. under $\varinjlim$ of cardinality $\leqslant \pi$, resp. under sums of cardinality $\leqslant \pi$),} and \textit{under $\varinjlim$ large relative to $\pi$. Then $\mathcal{M}$ is stable under small filtering $\varinjlim$ (resp. under small $\varinjlim$ , resp. under small sums).} % [sic] des cardinal -> de cardinal, agreement (cf. "de cardinal" elsewhere in the statement)

This lemma is proved like 4.4.14, page 97.

There remains the converse in 4.11.2, which follows from the similar assertion in 4.11.1.

\textbf{Corollary 4.11.4.} \textit{When is $\mathcal{M} = \Ind_\pi C$ stable under sums? If there exists an ordinal $\alpha$ such that $\operatorname{card}\alpha \leqslant \pi$, and such that $C$ is stable under $\varinjlim$ of type $I_\alpha$, then it suffices that $C$ be stable under sums of cardinality $\leqslant \pi$. This condition is also necessary if $C$ is Karoubian.}

(\footnote{For other variants of 4.11.1, cf. 4.13.19 (p. 340) and corollaries.}).

Is proved like 4.11.2, using 4.11.3.

\msp{page 271}

\textbf{Proposition 4.11.5} (\footnote{?, proof not finished.}). \textit{Let $\Phi$ be a set of categories $I$ equivalent to finite categories, and ``rigid'', i.e. for every $i \in \operatorname{Ob} I$, $\Hom_I(i,i) = \{\mathrm{id}_i\}$. Let $C$ be a category equivalent to a small category, $\pi$ an infinite cardinal, $\mathcal{M} = \Ind_\pi C$. For $\mathcal{M}$ to be stable under $\varinjlim$ (resp. under $\varprojlim$ ) of type $\Phi$, it} suffices \textit{that $C$ be so. In the non-resp. case, if $C$ is Karoubian, this condition is also necessary.}

This last assertion is contained in the similar assertion in 4.11.1. As for the first assertion, this follows in the usual way from the

\textbf{Lemma 4.11.6} (\footnote{?, proof not spelled out is perhaps false!}). \textit{Let $J$ be a category equivalent to a finite category, $J$ rigid, $C$ a category equivalent to a small category, $\pi$ an infinite cardinal. Then every functor}
\printedpage{102}
% head: A. Grothendieck : Dérivateurs, Ch. XVIII, Édition des Universités de Montpellier II et Paris VII
% foot: Version 17 décembre 2025
\textit{$J \longrightarrow \mathcal{M} \defeq \Ind_\pi C$ admits an indicial representation $J \times I \longrightarrow C$, with $I$ large relative to $\pi$, in other words, the} \msp{page 272} \textit{canonical fully faithful functor}
\[
\Ind_\pi(\Hom(J,C)) \ \longrightarrow \ \Hom(J,\mathcal{M})
\]
\textit{deduced from the inclusion $C \into \mathcal{M}$, hence $\Hom(J,C) \into \Hom(J,\mathcal{M})$ (cf. 4.9), is an equivalence of categories.} (\textbf{N.B.} I do not know whether this is so without a rigidity hypothesis on $J$.)

\textsc{Proof}. This is proved like the similar result for $\Ind(C)$ instead of $\Ind_\pi C$, cf. SGA 4 I 8.8.5.

\msp{page 273}

% a vertical rule is printed in the left margin beside Corollaire 4.11.7
\textbf{Corollary 4.11.7} (\footnote{subject to the validity of 4.11.5.}). \textit{For $\mathcal{M} = \Ind_\pi C$ to be stable under $\varinjlim$ (resp. under $\varprojlim$ ) that are} finite, \textit{it} suffices \textit{that $C$ be so. And this is also necessary in the non-resp. case, when $C$ is assumed Karoubian.}

\section*{4.12\quad Accessible functors between accessible categories}

\textbf{Proposition 4.12.1} (\footnote{Would go back to 4.8.}). \textit{Let $\mathcal{M}$, $\mathcal{N}$ be two accessible categories (4.8), and $f : \mathcal{M} \longrightarrow \mathcal{N}$ an accessible functor. Then there exist infinite cardinals $\pi_0$, $\pi_1$ such that $\pi_1 \geqslant \pi_0$ and such that for every cardinal $\pi \geqslant \pi_1$, we have}
\[
f(\mathcal{M}_\pi) \ \subset \ \mathcal{N}_{\pi^{\pi_0}} \ , \tag*{(4.12.1.1)} % [sic] (4.12.1) -> (4.12.1.1), the formula is cited below as (4.12.1.1) and the next one is numbered (4.12.1.2)
\]
\DERmark{ENV1298}

\textit{and in particular}
\[
f(\mathcal{M}_\pi) \ \subset \ \mathcal{N}_\pi \qquad \textit{if } \pi = \pi^{\pi_0}. \tag*{(4.12.1.2)}
\]
\textit{More precisely, we can take for $\pi_0$ a useful cardinal for $\mathcal{M}$} \msp{page 274} \textit{and such that $f$ is $\pi_0$-accessible (e.g. the smallest cardinal of this type). Then let $\pi_1$ be a cardinal such that we have $f(\mathcal{M}_{\pi_0}) \subset \mathcal{N}_{\pi_1}$, $\pi_1 \geqslant \cardred \mathcal{M}_{\pi_0}$, and finally such that $\mathcal{N}_{\pi_1}$ is} \texttt{[a]} \textit{strictly generating subcategory of $\mathcal{N}$. The pair $(\pi_0,\pi_1)$ satisfies the stated conditions.}

\textsc{Proof}. Indeed we will have
\[
\mathcal{M} \ \xleftarrow{\ \sim\ } \Ind_{\pi_0}(\mathcal{M}_{\pi_0})
\]
\DERmark{ENV1299}

and $f$, commuting with $\varinjlim$ large relative to $\pi_0$, will be known (up to unique isomorphism) by its restriction to $\mathcal{M}_{\pi_0}$, which factors through $\mathcal{N}_{\pi_1}$,
\[
f|\mathcal{M}_{\pi_0} : \mathcal{M}_{\pi_0} \ \longrightarrow \ \mathcal{N}_{\pi_1} \ .
\]
To prove (4.12.1.1), putting $\pi' = \pi^{\pi_0} \geqslant \pi$, and noting that $\pi'^{\pi_0} = \pi'$ \texttt{[rather} $\pi'^{\pi_0} = \pi'$\texttt{]},
\DERmark{ENV1300}

\[
f(\mathcal{M}_\pi) \ \subset \ f(\mathcal{M}_{\pi'}) \ \overset{?}{\subset} \ \mathcal{N}_{\pi'} \ ,
\]
\printedpage{103}
% head: A. Grothendieck : Dérivateurs, Ch. XVIII, Édition des Universités de Montpellier II et Paris VII
% foot: Version 17 décembre 2025
we see that it suffices to prove (4.12.1.2) (for $\pi'$). So we can assume $\pi = \pi^{\pi_0}$ and prove (4.12.1.2). But since $\pi_1 \geqslant \cardred \mathcal{M}_{\pi_0}$, we know that if $\pi$ is such that $\pi \geqslant \pi_1$, $\pi^{\pi_0} = \pi$, then $\mathcal{M}_\pi$ consists of the $\varinjlim\limits_{I} x_i$, the $x_i \in \mathcal{M}_{\pi_0}$, and $I$ large relative to $\pi_0$, and $\operatorname{card} I \leqslant \pi$.
\DERmark{ENV1301}


\msp{page 275}

We then have ($f$ being $\pi_0$-accessible)
\[
f(x) \ = \ \varinjlim_{I} f(x_i) \ ,
\]
the $f(x_i)$ are in $\mathcal{N}_{\pi_1}$, hence in $\mathcal{N}_\pi$, now $\mathcal{N}_\pi$ stable under $\varinjlim$ of cardinality $\leqslant \pi$ (since $\mathcal{N}$ is generated by $\mathcal{N}_{\pi_1}$, consisting of objects that are $\pi_1$-accessible and a fortiori $\pi$-accessible, cf. SGA 4 I 9.9). Hence $f(x) \in \mathcal{N}_\pi$, q.e.d.

\textbf{Corollary 4.12.2.} \textit{Assume that $\pi_0$ is also useful for $\mathcal{N}$, and that $\pi_1 \geqslant \cardred \mathcal{N}_{\pi_0}$. Let $\pi$ be a cardinal $\geqslant \pi_1$, and such that $\pi^{\pi_0} = \pi$. Then $f$ is identified, up to equivalence of categories and isomorphism of functors, with the functor}
\[
\Ind(f_\pi) \ : \ \Ind_\pi(\mathcal{M}_\pi) \ \longrightarrow \ \Ind_\pi(\mathcal{N}_\pi) \ ,
\]
\textit{where $f_\pi : \mathcal{M}_\pi \longrightarrow \mathcal{N}_\pi$ is the functor induced by $f$.}

\textbf{Proposition 4.12.3} (\footnote{goes back to 4.7.}). \textit{Let $C$ be a category equivalent to a small category, $\mathcal{N}$ a $\mathfrak{U}$-category, $\pi_0$ an infinite cardinal.} \msp{page 276} \textit{We assume that $\mathcal{N}$ satisfies $L_{\pi_0}$ (stability under $\varinjlim$ large relative to $\pi_0$). Let $\mathcal{M} = \Ind_{\pi_0} C$, $\Hom_{!\,\pi_0}(\mathcal{M},\mathcal{N})$ the strictly full subcategory of $\Hom(\mathcal{M},\mathcal{N})$ consisting of the $f$ : $\mathcal{M} \longrightarrow \mathcal{N}$ that are $\pi_0$-accessible. Then the restriction functor}
\[
\Hom_{!\,\pi_0}(\mathcal{M},\mathcal{N}) \ \longrightarrow \ \Hom(C,\mathcal{N}) \tag*{(4.12.3.1)}
\]
\textit{is an equivalence of categories. We find a quasi-inverse functor as the composite}
\[
\Hom(C,\mathcal{N}) \ \longrightarrow \ \Hom(\underbrace{\Ind_{\pi_0}(C)}_{\mathcal{M}},\Ind_{\pi_0}(\mathcal{N})) \ \longrightarrow \ \Hom(\mathcal{M},\mathcal{N}) \ , \tag*{(4.12.3.2)}
\]
\textit{where the second arrow is the composition functor $F \longmapsto L \circ F$ with the direct limit functor $L : \Ind_{\pi_0}\mathcal{N} \longrightarrow \mathcal{N}$. The composite functor (4.12.3.2) has image contained in $\Hom_{!\,\pi_0}(\mathcal{M},\mathcal{N})$ and induces an equivalence between this last category and $\Hom(C,\mathcal{N})$.} % [sic] (4.11.3.2) -> (4.12.3.2), the composed functor is the display just labelled (4.12.3.2)
\DERmark{ENV1302}


\msp{page 277}

\textit{In other words, we can say that the canonical inclusion functor}
\[
C \ \longrightarrow \ \Ind_{\pi_0} C \ = \ \mathcal{M}
\]
\textit{is} 2-universal \textit{for functors from $C$ into categories $\mathcal{N}$ satisfying $L_{\pi_0}$, the latter being regarded as $0$-objects of a $2$-category where the $1$-objects are the functors $\mathcal{N} \longrightarrow \mathcal{N}'$ that are $\pi_0$-accessible, i.e. commute with $\varinjlim$ large relative to $\pi_0$.}

This statement is the analogue, for arbitrary small $C$, of 4.4.16 (page 101), valid for $C$ stable under $\varinjlim$ of cardinality $\leqslant \pi_0$. There are of course other variants, in relation to commuting with other types of $\varinjlim$ ( filtering $\varinjlim$, sums, finite direct limits etc.) considered in the preceding section 4.11.

The proof of 4.12.3 is sorital.


\endgroup
% END UNIT der-ch18/der-ch18_p101-105.tex
% BEGIN UNIT der-ch18/der-ch18_p106-110.tex
\begingroup\makeatletter%
\newcommand{\printedpage}[1]{\marginpar{\footnotesize p.~#1}}%
\newcommand{\msp}[1]{\marginpar{\footnotesize [#1]}}%
\newcommand{\M}{\mathcal{M}}%
\newcommand{\N}{\mathcal{N}}%
\newcommand{\cH}{\mathcal{H}}%
\newcommand{\Hom}{\underline{\mathrm{Hom}}}%
\newcommand{\Ind}{\operatorname{Ind}}%
\newcommand{\Ob}{\operatorname{Ob}}%
\newcommand{\Fl}{\operatorname{Fl}}%
\newcommand{\card}{\operatorname{card}}%
\newcommand{\red}{\mathop{\textup{red}}\nolimits}%
\newcommand{\uS}{\underline{S}}%
\makeatother\let\printedpage\derpp\let\msp\dermsp%
\let\mspin\DERmspin%


\printedpage{104}
% head: A. Grothendieck : Dérivateurs, Ch. XVIII, Édition des Universités de Montpellier II et Paris VII
% foot: Version 17 décembre 2025
\msp{page 278}
\textbf{Corollary 4.12.4.} \textit{Let $\M$ and $\N$ be two accessible categories, and let $\pi_0$ be an infinite cardinal such that $\M$ is $\pi_0$-accessible, i.e. such that $\pi_0$ is a useful cardinal for $\M$. Let $\pi_1 = \sup(\pi_0,\allowbreak \card\red\M_{\pi_0})$. Then for every cardinal $\pi \geqslant \pi_1$ useful for $\M$ (and in particular, for every $\pi \geqslant \pi_1$ such that $\pi^{\pi_0} = \pi$), the category of functors $\M \longrightarrow \N$ that are $\pi$-accessible is itself accessible.}

Indeed, we will have $\M \simeq \Ind_\pi\M_\pi$, whence by 4.12.3
\[
\Hom_{!\,\pi}(\M,\N) \xrightarrow{\sim} \Hom(\M_\pi,\N) ,
\]
now $\Hom(\M_\pi,\N)$ is accessible by 4.9.5 (p. 220). We can moreover make this point more precise. Suppose $\pi_0$ chosen so that $\pi_0$ is useful not only for $\M$, but for $\N$ as well. Recall that we have $\card\Fl\M_\pi \leqslant 2^\pi$ (cf. 4.7.7, page 173).

\msp{page 279}
Let us therefore put
\[
\pi' = 2^\pi ,
\]
we therefore have $\pi'^{\pi_0} = 2^{\pi\pi_0} = 2^\pi = \pi'$, and $\card\red\M_\pi \leqslant \pi'$. Thus 4.9.5 gives us that
\[
\Hom(\M_\pi,\N) \simeq \Ind_{\pi'}(\Hom(\M_\pi,\N_{\pi'})) , % [sic] \M_{\pi'} -> \N_{\pi'}, the second argument should lie in N, as in Corollaire 4.12.5
\]
\DERmark{ENV1303}

whence:

\textbf{Corollary 4.12.5.} \textit{Under the conditions of 4.12.4, if $\pi_0$ is useful for $\M$} and \textit{for $\N$, then for $\pi \geqslant \pi_1$, $\Hom_{!\,\pi}(\M,\N)$ is $\pi'$-accessible, where $\pi' = 2^\pi$, more precisely, we have an equivalence of categories}
\[
\Hom_{!\,\pi}(\M,\N) \simeq \Ind_{\pi'}(\Hom(\M_\pi,\N_{\pi'})) .
\]

\textbf{Remark 4.12.6.} If $\M$, $\N$ are accessible, the category $\Hom_{\mathrm{acc}}(\M,\N)$ of accessible functors is not always accessible. If $\N$ satisfies $\mathrm{L}_{\pi_0}$, this category satisfies $\mathrm{L}_{\pi_0}$, % [sic] par toujours -> pas toujours, negation "n'est pas" intended
and the $\varinjlim$ large relative to $\pi_0$ are computed argument by argument.

\msp{page 280}
If we could find a subcategory (equivalent to a small category) $\cH_\pi$ of $\cH = \Hom_{\mathrm{acc}}(\M,\N)$, for a cardinal $\pi \geqslant \pi_0$ such that every object $F$ of $\cH$ could be written $\varinjlim\limits_{I} F_i$, the $F_i$ in $\cH_\pi$, $I$ large relative to $\pi$, then if the objects of $\cH_\pi$ are $\pi'$-accessible functors from $\M$ to $\N$, the same would hold for every $F$ in $\cH$. Thus there would exist an infinite cardinal $\pi'$ such that every accessible functor from $\M$ to $\N$ is $\pi'$-accessible. But unless I am mistaken, if e.g. $\M = \N = \mathrm{Ens}$, there is no such cardinal. I do not want to dwell here on this question, and make a counterexample.

\msp{page 281}

\section*{4.13\quad Accessible sets of objects of an accessible category}

\textbf{Definition 4.13.1.} Let $\M$ be an accessible category, $S$ a subset of $\Ob\M$, $\uS$ the full subcategory of $\M$ defined by $S$. We say that $S$ is an \textit{accessible subset} of $\Ob\M$, or again that $\uS$ is a \textit{subcategory ``accessible in $\M$''}, if
\printedpage{105}
% head: A. Grothendieck : Dérivateurs, Ch. XVIII, Édition des Universités de Montpellier II et Paris VII
% foot: Version 17 décembre 2025
\begin{enumerate}
\item[a)] $\uS$ is strictly full, i.e. for $x$, $y$ two isomorphic objects of $\M$, $x \in S$ implies $y \in S$,
\item[b)] $\uS$ is accessible,
\item[c)] the inclusion functor $\uS \longrightarrow \M$ is accessible.
\end{enumerate}
Let $\pi_0$ be a useful cardinal for $\M$, i.e. such that $\M \simeq \Ind_{\pi_0}(\M_{\pi_0})$, let $\pi_1 = \sup(\pi_0, \card\red\M_{\pi_0})$. On the other hand let $\pi'_0$ be a cardinal such that

\msp{page 282}
$\uS$ (assumed accessible in $\M$) is $\pi'_0$-\textit{accessible}, i.e. $\pi'_0$-useful for $\uS$, and moreover such % [sic] \pi_0-accessible -> \pi'_0-accessible, the following "i.e. \pi'_0-utile" refers to \pi'_0
\DERmark{ENV1304}

that $f : \uS \hookrightarrow \M$ is $\pi'_0$-accessible. Let $\pi'_1 \geqslant \pi'_0$ be a cardinal such that $\uS_{\pi'_0} \subset \M_{\pi'_1}$, $\pi'_1 \geqslant \card\red\uS_{\pi'_0}$. It follows from 4.12.1 that for every cardinal
\[
\pi \geqslant \pi'_1 , \qquad \text{such that } \pi^{\pi'_0} = \pi, \tag*{$(*)$}
\]
we must have
\[
\uS_\pi \subset \M_\pi , \qquad \text{hence } \uS_\pi \subset \M_\pi \cap \uS,
\]
moreover, since $\uS \hookrightarrow \M$ is $\pi'_0$-accessible and a fortiori $\pi$-accessible, and fully faithful, it follows at once that $\M_\pi \cap S \subset S_\pi$, whence
\[
S_\pi = \M_\pi \cap S . % [sic] \M \cap S -> \M_\pi \cap S, the preceding lines concern \M_\pi \cap S
\]
\DERmark{ENV1305}


\msp{page 283}
Moreover, condition $(*)$ implies that $\pi$ is useful for $\uS$, hence
\[
\uS \xleftarrow{\sim} \Ind_\pi(\uS_\pi) .
\]
$\Big($ Moreover, \textit{if} $\pi'_0 \geqslant \pi_0$, the relation $\pi^{\pi'_0} = \pi$ implies a fortiori $\pi^{\pi_0} = \pi$, and if $\pi'_1 \geqslant \pi_1$, we see that $(*)$ implies that $\pi$ is also useful for $\M$, i.e.
\[
\M \xleftarrow{\sim} \Ind_\pi(\M_\pi) .
\]
On the other hand, the inclusion functor $\uS \hookrightarrow \M$ is obtained essentially from $\uS_\pi \hookrightarrow \M_\pi$ by passing to the categories $\Ind_\pi$. (\footnote[116]{\texttt{[The content of this parenthesis is]} ``optional''.}). $\Big)$

This proves among other things the ``only if'' in the following proposition:

\textbf{Proposition 4.13.2.} \textit{Let $\M$ be an accessible category, $\pi_0$ an infinite cardinal useful for $\M$, $S$ a subset of} $\Ob\M$, \textit{$\uS$ the corresponding full subcategory of $\M$.}

\msp{page 284}
\begin{enumerate}
\item[\textit{a)}] \textit{For $S$ to be accessible, it is necessary and sufficient that there exist an infinite cardinal $\pi$ and a subset $S(\pi)$ of $\M_\pi$, such that $S$ consists of the objects of $\M$ that are isomorphic to an object of the form $\varinjlim\limits_{i\in I} x_i$, with the $x_i$ in $S(\pi)$ and $I$ large relative to $\pi$. In other words, it is necessary and sufficient that there exist an infinite cardinal $\pi$ and a full subcategory % [sic] \pi_1 -> \pi, the rest of the statement uses \pi
$\uS(\pi)$ of $\M_\pi$, such that $\uS$ is the essential image of the canonical fully faithful functor}
\DERmark{ENV1306}

\[
\Ind_\pi(\uS(\pi)) \longrightarrow \M \qquad \bigl(\xleftarrow{\sim} \Ind_\pi(\M_\pi)\ \text{(\footnotemark[117])}\bigr) , \tag*{(4.13.2.1)} % [sic] \Ind_\pi(\M) -> \Ind_\pi(\M_\pi), footnote 117 refers to \pi utile pour \M, i.e. \M \simeq \Ind_\pi(\M_\pi)
\]
\DERmark{ENV1307}

\footnotetext[117]{if $\pi$ useful for $\M$.}
\textit{extending the canonical inclusion $\uS(\pi) \hookrightarrow \M$ (cf. 4.12.3).}
\item[\textit{b)}] \printedpage{106}
% head: A. Grothendieck : Dérivateurs, Ch. XVIII, Édition des Universités de Montpellier II et Paris VII
% foot: Version 17 décembre 2025
\textit{If $S$ is accessible, we can find cardinals $\pi'_0$, $\pi'_1$ such that $\pi'_1 \geqslant \pi'_0 \geqslant \pi_0$ (where $\pi_0$ is a given cardinal useful for $\M$)}
\msp{page 285}
\textit{and such that every cardinal $\pi \geqslant \pi'_1$ with $\pi^{\pi'_0} = \pi$ satisfies the conditions of a), taking $S(\pi) = S \cap \Ob\M_\pi$, i.e. $\uS(\pi) = \uS \cap \M_\pi$. Under these conditions, $\pi$ is useful for $\uS$} and \textit{for $\M$, $S(\pi)$ is a} strictly \textit{full subcategory of $\M_\pi$, stable in $\M_\pi$ under direct factors (or what amounts to the same thing, $\M_\pi$ being Karoubian, a Karoubian subcategory of $\M_\pi$).}
\DERmark{ENV1308}

\item[\textit{c)}] \textit{We can take $\pi'_0$ to be the smallest cardinal $\geqslant \pi_0$ useful for $\uS$ and such that $\uS \hookrightarrow \M$ is $\pi'_0$-accessible, and $\pi'_1 = \sup(\pi'_0, \card\red\uS_{\pi'_0})$.}
\end{enumerate}

Part b) as well as c) has already been proved, which implies the ``only if'' in part a). It remains to prove the ``if'' in a). But the hypothesis implies that
\[
\uS \xleftarrow{\sim} \Ind_\pi(S_\pi) ,
\]
which implies that $\uS$ is accessible and even $\pi$-accessible, and it also implies that $\uS$ is a \textit{strictly} full subcategory of $\M$. It remains to see

\msp{page 286}
\texttt{[that]} the inclusion functor is accessible, in fact it is $\pi$-accessible, since (4.13.2.1) is (cf. 4.12.3).

\textbf{Definition 4.13.3.} Let $\M$, $S \subset \Ob\M$, $\uS \subset \M$ be as in 4.12.1, $\pi$ an infinite cardinal. We say that $S$ is a $\pi$-\textit{accessible} subset of $\Ob\M$, or that $\uS$ is a $\pi$-accessible subcategory in $\M$, if $S$, $\pi$ satisfy the conditions of 4.13.2 a), i.e. $\uS$ is the essential image % [sic] 4.13.3 a) -> 4.13.2 a), a self-reference of the definition, the conditions a) are those of Proposition 4.13.2
\DERmark{ENV1309}

of $\Ind_\pi(S(\pi))$, where $S(\pi)$ is a full subcategory of $\M_\pi$.

Note that, at the cost of enlarging $S(\pi)$ to $S'(\pi) = {\natural}$-closure of $S(\pi)$ in $\M_\pi$, we can suppose that $S(\pi)$ is a \textit{strictly full and Karoubian} subcategory

\msp{page 287}
of $\M_\pi$. We then know that, $S(\pi)$ being Karoubian, $S(\pi) \longrightarrow (\underbrace{\Ind_\pi(S(\pi))}_{})_\pi$ is an equivalence of categories. Since $\Ind_\pi(S(\pi)) \simeq \uS$ and since $S(\pi)$ is strictly full in $\M_\pi$, hence also in $\M$ and a fortiori in $\uS$, that we have
\[
\uS(\pi) = \uS_\pi , \qquad \text{subcategory of } \uS \text{ of the } \pi\text{-accessible objects of } S. \tag*{(4.13.1)}
\]
I say that we will also have
\[
\uS(\pi) = \uS \cap \M_\pi . \tag*{(4.13.2)}
\]
As c is clear, it remains to see that $\uS{\cap}\M_\pi \subset \uS(\pi)$, now it is immediate that $\uS{\cap}\M_\pi \subset S_\pi$, % [sic] c -> an inclusion sign (\supset), a letter c printed where the reverse inclusion seems meant
and we conclude by (4.13.1). We thus find:

\textbf{Proposition 4.13.4.} \textit{Let $\M$ be an accessible category, $\pi$ an infinite cardinal. Let $\Sigma_\pi$ be the set of}

\msp{page 288}
\textit{strictly full subcategories $\uS(\pi)$ of $\M_\pi$ that are Karoubian, or what amounts to the same thing, stable in $\M_\pi$ under direct factors. Let $\Sigma'_\pi$ be the set of subcategories of}
\printedpage{107}
% head: A. Grothendieck : Dérivateurs, Ch. XVIII, Édition des Universités de Montpellier II et Paris VII
% foot: Version 17 décembre 2025
\textit{$\M$ that are $\pi$-accessible in $\M$. Consider the maps}
\[
\begin{tikzcd}
\Sigma_\pi \arrow[r, shift left, "\alpha_\pi"] & \Sigma'_\pi \arrow[l, shift left, "\beta_\pi"]
\end{tikzcd}
\tag*{(4.13.4.1)}
\]
\textit{defined as follows: If $\uS(\pi) \in \Sigma_\pi$, $\alpha_\pi(\uS(\pi))$ is the essential image of the fully faithful functor} $\Ind_\pi(\uS(\pi)) \longrightarrow \M$. \textit{If $\uS \in \Sigma'_\pi$, $\beta_\pi(\uS) = \uS \cap \M_\pi$. The maps $\alpha_\pi$, $\beta_\pi$ are bijections, inverse to one another.}
\DERmark{ENV1310}


We have not had to suppose that $\pi$ is useful for $\M$, nor that $\pi$ is bounded below by a cardinal useful for $\M$, i.e. that we have $\pi \geqslant \pi_0$, where $\pi_0$ is the smallest cardinal useful for $\M$.

\msp{page 289}
\textbf{4.13.5.} Let $\uS$ be a \textit{strictly full} subcategory of $\M$ (accessible category), $\pi'_0$ an infinite cardinal. To say that $\uS$ is $\pi'_0$-accessible in $\M$ means that the following conditions are satisfied.
\begin{enumerate}
\item[1)] $\pi'_0$ is useful for $\uS$, i.e. $\uS \xleftarrow{\sim} \Ind_{\pi'_0}(\uS_{\pi'_0})$.
\item[2)] $\uS \hookrightarrow \M$ is $\pi'_0$-accessible.
\item[3)] $\uS_{\pi'_0} \hookrightarrow \M_{\pi'_0}$ (which implies that $S_{\pi'_0}$ is equivalent to a small category).
\end{enumerate}
Suppose these conditions satisfied, and let $\pi$ be a cardinal $> \pi'_0$. Is $\uS$ $\pi$-accessible in $\M$, i.e. are the similar conditions satisfied for $\pi$? This is obvious for 2°), but may be false for 1°, 2° \texttt{[rather 3°\,?]}. But let us call on 4.13.2 c). Put
\[
\pi'_1 = \sup(\pi'_0, \card\red S_{\pi'_0}, \underbrace{\pi_0}_{\substack{\text{smallest infinite}\\ \text{cardinal useful}\\ \text{for } \M}} ) .
\]

\msp{page 290}
Then for $\pi \geqslant \pi'_1$ such that $\pi^{\pi'_0} = \pi$, $\uS$ is $\pi$-accessible in $\M$ (4.13.2 b and c). Suppose e.g.
\[
\pi'_0 \geqslant \pi_0
\]
(in practice this will always be satisfied) and moreover
\[
\pi'_0 \geqslant \card\red\M_{\pi_0} ,
\]
i.e. suppose
\[
\pi'_0 \geqslant \sup(\pi_0, \card\red\M_{\pi_0}) = \pi_1 .
\]
Then
\[
\pi'_1 = \sup(\pi'_0, \card\red S_{\pi'_0}) \leqslant \sup(\pi'_0, \card\red\M_{\pi'_0}) \leqslant 2^{({\pi'_0}^{\pi_0})}
\]
(4.7.7.1, p. 177). So we find this:

\textbf{Corollary 4.13.6.} \textit{Let $\pi'_0$ be a cardinal $\geqslant \pi_1 = \sup(\pi_0, \card\red\M_{\pi_0})$ (\footnote[118]{\textbf{N.B.} $\pi_0$ is the smallest cardinal useful for $\M$.}), and let $\uS$ be a $\pi'_0$-accessible subcategory in $\M$. Then for every cardinal $\pi$ such that $\pi \geqslant 2^{({\pi'_0}^{\pi_0})}$, and such that we have $\pi^{\pi'_0} = \pi$ (i.e. $\pi$ of the form $c^{\pi'_0}$, with $c$ a cardinal $\geqslant 2^{({\pi'_0}^{\pi_0})}$), $\uS$ is $\pi$-accessible in $\M$. (Moreover, these cardinals are regular for $\uS$} and \textit{for $\M$, hence $\uS \simeq \Ind_\pi(\uS_\pi)$, $\M \simeq \Ind_\pi(\M_\pi)$).}

(\footnote[119]{Continue here with 4.13.23 etc.}, \footnote[120]{Put here lemma 4.16.7.11 noted in the margin p. 292.}).

\printedpage{108}
% head: A. Grothendieck : Dérivateurs, Ch. XVIII, Édition des Universités de Montpellier II et Paris VII
% foot: Version 17 décembre 2025
\msp{page 291}
\textbf{Proposition 4.13.7} (\footnote[121]{Move it after 4.13.23.}). \textit{Let $\M$ be an accessible category, $(S_\alpha)_{\alpha\in A}$ a family of subsets of} $\Ob\M$, \textit{indexed by a small set $A$, $(\uS_\alpha)_{\alpha\in A}$ the corresponding family of accessible subcategories in $\M$. Let $S = \bigcap_{\alpha\in A} S_\alpha$, $\uS = \bigcap_{\alpha\in A} \uS_\alpha$, so that $\uS$ is the full subcategory of $\M$ defined by $S \subset$} $\Ob\M$. \textit{This being so, if the $S_\alpha$ are accessible, so is $S$ – or again, if the $\uS_\alpha$ are accessible in $\M$, so is $\uS$.}

This follows from the more precise result:

\textbf{Corollary 4.13.8} (\footnote[122]{Cf. more reasonable and more correct reformulation page 302, 302 bis.}). \textit{With the notation of 4.13.7 and 4.13.6, let, for every $\alpha \in A$, $\pi'_\alpha$ be a cardinal $\geqslant \pi_1$ such that $S_\alpha$ is $\pi'_\alpha$-accessible. Let}
\[
\begin{aligned}
\pi'_0 &= \sup\nolimits_{\alpha\in A} \pi'_\alpha \\
\pi''_0 &= \sup(\card A, \sup\nolimits_\alpha 2^{({\pi'_\alpha}^{\pi_0})})
\end{aligned}
\]
\textit{and let $\pi \geqslant \pi''_0$ be a cardinal such that $\pi^{\pi'_0} = \pi$, i.e. of the form $c^{\pi'_0}$, with $c \geqslant \pi'_0$. Then $S$ is $\pi$-accessible (i.e. $\uS$ is $\pi$-accessible in $\M$). In particular, $S$ is $2^{\pi''_0}$-accessible.} (\footnote[123]{This looks a bit outlandish! We must moreover suppose that $\pi$ is not a limit cardinal!}).

\msp{page 292}
\textsc{Proof}. Let $\alpha \in A$. Since $\pi'_\alpha \leqslant \pi'_0$ and $\pi^{\pi'_0} = \pi$, we have a fortiori
\[
\pi^{\pi'_\alpha} = \pi , \qquad \text{moreover } \pi \geqslant \pi''_0 \geqslant 2^{({\pi'_\alpha}^{\pi_0})} , % doubt resolved against PDF; see corrections.csv
\]
\DERmark{ENV1311}

therefore by 4.13.6, $\uS_\alpha$ is $\pi$-accessible in $\M$. On the other hand, since $\uS_\alpha$ is $\pi'_\alpha$-accessible, it is stable in $\M$ under $\varinjlim$ large relative to $\pi'_\alpha$, a fortiori under $\varinjlim$ large relative to $\pi'_0$ ($\geqslant \pi'_\alpha$). Moreover, $\pi^{\pi'_0} = \pi$ implies $\pi \geqslant 2^{\pi'_0} > \pi'_0$. Thus 4.13.8 follows from the
\DERmark{ENV1312}


\textbf{Corollary 4.13.9.} \textit{Let $\pi > \pi_0$ be a cardinal such that $\uS_\alpha$ is $\pi$-accessible in $\M$ for every $\alpha \in A$, and such that} $\card A \leqslant \pi$. \textit{Suppose moreover that there exists a cardinal $\pi'_0$ such that $\pi_0 \leqslant \pi'_0 < \pi$, such that the $\uS_\alpha$ are stable in $\M$ under $\varinjlim$ large relative to $\pi'_0$. Then $\uS = \bigcap_{\alpha\in A} \uS_\alpha$ is $\pi$-accessible in $\M$.}

(\footnote[124]{This would be the moment to formulate here Lemma 4.16.7.11 (p. 488), which should even be moved up with the first generalities on $\pi$-accessible subsets to as early as before 4.13.7 (immediately after 4.13.6).}).

Put, for every $\alpha \in A$,
\[
\begin{aligned}
\uS_\alpha(\pi) &\overset{\mathrm{def}}{=} \uS_\alpha \cap \M_\pi \\
\uS(\pi) &\overset{\mathrm{def}}{=} \uS \cap \M_\pi = \bigcap \uS_\alpha(\pi) .
\end{aligned}
\]
\DERmark{ENV1313}\DERmark{ENV1314}

Each $\uS_\alpha(\pi)$ is stable in $\M$ under $\varinjlim$ large relative to $\pi'_0$ and of cardinal $\leqslant \pi$ (since this is so for $\uS_\alpha$ and $\M_\pi$) (so the same is true of their intersection

\msp{page 293}
$\uS(\pi)$. A fortiori, $\uS(\pi)$ is a strictly full and Karoubian subcategory of $\M_\pi$.)


\endgroup
% END UNIT der-ch18/der-ch18_p106-110.tex
% BEGIN UNIT der-ch18/der-ch18_p111-115.tex
\begingroup\makeatletter%
\newcommand{\printedpage}[1]{\marginpar{\footnotesize p.~#1}}%
\newcommand{\msp}[1]{\marginpar{\footnotesize [#1]}}%
\newcommand{\cercle}[1]{\tikz[baseline=(c.base)]\node[draw,circle,inner sep=1pt](c){#1};}%
\newcommand{\incl}{\lhook\joinrel\longrightarrow}%
\newcommand{\card}{\operatorname{card}}%
\newcommand{\Ob}{\operatorname{Ob}}%
\newcommand{\Ind}{\operatorname{Ind}}%
\makeatother\let\printedpage\derpp\let\msp\dermsp%
\let\mspin\DERmspin%


\printedpage{109}
% head: A. Grothendieck : Dérivateurs, Ch. XVIII, Édition des Universités de Montpellier II et Paris VII

On the other hand, each $\underline{S}_\alpha$ being stable in $\mathcal{M}$ under $\varinjlim$ large relative to $\pi'_0$, the same holds for their intersection $\underline{S}$. A fortiori, the latter is stable under $\varinjlim$ large relative to $\pi$, so the canonical fully faithful functor $\Ind_\pi(\underline{S}(\pi))\longrightarrow\mathcal{M}$ factors through $S$ as
\[ \Ind_\pi\underline{S}(\pi) \incl \underline{S} \; , \]
we must prove that this functor is essentially surjective. So let $x\in S=\Ob\underline{S}$, it all comes down to proving that
\[ \varinjlim_{i\in\underline{S}(\pi)/x} x_i \xrightarrow{\ \sim\ } x \; , \tag*{$(*)$} \]
and that $\underline{S}(\pi)/x$ is large relative to $\pi$.

Note that, since $\pi'_0\geqslant\pi_0$, and $\pi^{\pi'_0}=\pi$, we will have $\pi^{\pi_0}=\pi$, so $\pi$ is a useful cardinal for $\mathcal{M}$, so $I=\mathcal{M}_\pi/x$ is large relative to $\pi$, and we have % [sic] grande devant $x$ -> grande devant $\pi$, apparently intended (one speaks of being large before a cardinal)
\DERmark{ENV1315}

\[ \varinjlim_{i\in I=\mathcal{M}_\pi/x} x_i \xrightarrow{\ \sim\ } x \; . \tag*{$(**)$} \]
Thus, if we admit that $\underline{S}(\pi)/x$ is large relative to $\pi$,
\msp{page 294}
the fact that $(*)$ is an isomorphism is equivalent to saying that the canonical functor
\[ J \overset{\mathrm{def}}{=} \underline{S}(\pi)/x \incl \mathcal{M}_\pi/x \; = \; I \] % [sic] $\mathcal{M}_\pi = I$ -> $\mathcal{M}_\pi/x = I$, since $I=\mathcal{M}_\pi/x$ above and in the display for $J_\alpha$ below
\DERmark{ENV1316}\DERmark{ENV1317}

is \textit{cofinal} (cf. 4.4.2, p. 79). But we see at once that this functor is an inclusion functor of a strictly full subcategory (stable under $\varinjlim$ large relative to $\pi'_0$, of cardinality $\leqslant\pi$).

Note now that the arrow $(**)$ similar to $(*)$, with $\underline{S}(\pi)$ replaced by $\underline{S}_\alpha(\pi)$, is indeed an isomorphism, since $x\in\Ob\underline{S}\subset\Ob\underline{S}_\alpha$, and $\underline{S}_\alpha$ is $\pi$-accessible. This is interpreted as the fact that the inclusion functor
\[ J_\alpha \overset{\mathrm{def}}{=} \underline{S}_\alpha(\pi)/x \incl \mathcal{M}_\pi/x \; = \; I \]
\DERmark{ENV1318}

is cofinal. This is again an inclusion of a strictly full subcategory $J_\alpha$ of $I$, stable under $\varinjlim$ large relative to $\pi'_0$ and of cardinality $\leqslant\pi$, and we have
\[ J \; = \; \bigcap J_\alpha \; . \]

Thus the desired condition will follow from the
\msp{page 295}

% marginal hand-drawn mark (manuscript) in the left margin beside the following lemma
\textbf{Lemma 4.13.10} (\footnote[125]{We must assume that $\pi$ is not a limit cardinal, to be able to conclude that $I$ is stable under $\varinjlim$ of type $I_{\alpha_\pi}$, and the $J_r$ stable in $I$ under these limits.},\footnote[126]{Cf. more general form of the lemma in 4.13.10.2, p. 302.}). \textit{Let $\pi'_0$ and $\pi>\pi'_0$ be infinite cardinals, $I$ a category equivalent to a small category, $(J_\alpha)_{\alpha\in A}$ a family of subcategories $J_\alpha$ of $I$. We assume:}
\begin{itemize}
\item[\cercle{a}] \textit{$I$ is large relative to $\pi$ (}\footnote[127]{It suffices that $I$ be large relative to every $\pi'<\pi$.}\textit{).} % footnote 127 printed at the foot of p. 110
\item[\cercle{b}] \textit{$I$ is stable under $\varinjlim$ of cardinality $\leqslant\pi$, and large relative to $\pi'_0$ (}\footnote[128]{Here stability under $\varinjlim$ of type $I_{\alpha_\pi}$ suffices, cf. page 302 [?].}\textit{).} % footnote 128 printed at the foot of p. 110
\end{itemize}
% foot: Version 17 décembre 2025

\printedpage{110}
% head: A. Grothendieck : Dérivateurs, Ch. XVIII, Édition des Universités de Montpellier II et Paris VII
\begin{itemize}
\item[\cercle{c}] \textit{$\forall\ r\in A$, $J_r$ is stable in $I$ under the direct limits envisaged in b) (in particular, $J_r$ is a} strictly full \textit{subcategory of $I$).} % [sic] envisagés -> envisagées, agreement with « limites »
\item[\cercle{d}] \textit{$\forall\, r\in A$, $J_r$ is a} cofinal \textit{subcategory of $I$ (}large relative to $\pi$\textit{).}
\item[\cercle{e}] $\card A\leqslant\pi$.
\end{itemize}
\textit{Let $J=\bigcap_{r\in A}J_r$. Then $J$ is still cofinal in $I$, and large relative to $\pi$ (and moreover stable in $I$ under $\varinjlim$ of cardinality $\leqslant\pi$ and large relative to $\pi'_0$).}

In other words, if $I$ is large relative to $\pi$, the intersection $J$ of a family $(J_\alpha)_{\alpha\in A}$ of strictly full subcategories cofinal and large relative to $\pi$, is still
\msp{page 296}
cofinal and large relative to $\pi$, \textit{provided that} $\card A\leqslant\pi$, and that we postulate a stability property of $I$ and of the $J_\alpha$ in $I$, for direct limits of cardinality $\leqslant\pi$ and large relative to a suitable given cardinal $\pi'_0<\pi$.

\textbf{4.13.10.1.} Note moreover that if $f : J\longrightarrow I$ is a cofinal \textit{fully faithful} functor, with $I$ large relative to $\pi$, then $J$ is automatically large relative to $\pi$. This follows e.g. from 4.9.6.3 2° (p. 225) by the criterion $\mathrm{c}_\pi$ (since for $f$ fully faithful, condition F $3_\pi$ is tautologically satisfied). We can also see it thus: Let $B$ be a set with $\card B\leqslant\pi$, we must prove that $J\longrightarrow J^B$ is cofinal, now we have
\[
\begin{tikzcd}
J \arrow[r, "\text{cof.\,?}"] \arrow[d, "\text{cof.}"'] & J^B \arrow[d, "\text{cof.}"] \\
I \arrow[r, "\text{cof.}"] & I^B \; ,
\end{tikzcd}
\]
where three of the four arrows of the square are cofinal, $J^B\longrightarrow I^B$ being so as a product of cofinal functors between \textit{filtering} categories (cf. 4.9.6.5, page 227) -- indeed we already know that $J$ is filtering.

\msp{page 297}
Thus the composite
\[ J \longrightarrow J^B \xrightarrow{\text{cof.}} I^B \]
is cofinal. But since the second functor $J^B\longrightarrow I^B$ is fully faithful, it follows at once that $J\longrightarrow J^B$ is cofinal (for $\xi\backslash J$ for $\xi\in\Ob J^B$ is the same, up to isomorphism, whether we consider $\xi$ as an object of $J^B$, or of $I^B$).

Thus in the statement of condition d) of the lemma, it is unnecessary to specify that $J_\alpha$ is large relative to $\pi$, and likewise in the conclusion about $J$.

Since $J\incl I$ is fully faithful, to see that this functor is cofinal, it suffices to see that it satisfies F 1, i.e. that for every $i\in\Ob I$, there exists $j\in J$ with $i\longrightarrow j$. We are going to construct $j$ as the limit of an ordinal direct system of objects $i_\alpha$ of $I$, indexed by the ordinals $\alpha$ such that $\card\alpha<\pi$. Denote by $\alpha_\pi$ the smallest ordinal such that % [sic] le plus petite ordinal -> le plus petit ordinal, agreement with « ordinal »
% foot: Version 17 décembre 2025
\printedpage{111}
% head: A. Grothendieck : Dérivateurs, Ch. XVIII, Édition des Universités de Montpellier II et Paris VII
\msp{page 298}
% marginal hand-drawn mark (manuscript) in the left margin beside the following two lines
$\card\alpha_\pi=\pi$. The ordered set $I_{\alpha_\pi}$ of ordinals $\alpha<\alpha_\pi$ being large relative to $\pi'_0$ (\footnote[129]{false if $\pi$ is a limit cardinal.}) and of cardinality $\pi$, we will therefore be able to take the limit $\varinjlim\limits_{\alpha\in I_{\alpha_\pi}} i_\alpha$ in $I$. % [sic] $I_{\alpha_0}$ -> $I_{\alpha_\pi}$, the index set just defined is $I_{\alpha_\pi}$
\DERmark{ENV1319}


Construction of $(i_\alpha)_{\alpha\in I_{\alpha_\pi}}$. Let % [sic] $I_{\alpha_0}$ -> $I_{\alpha_\pi}$, the index set just defined is $I_{\alpha_\pi}$
\DERmark{ENV1320}

\[ \alpha_\pi \; = \; \text{smallest ordinal } \alpha \text{ such that } \card\alpha=\pi, \tag*{(1)} \]
so this is the ordinal characterized by
\[ \left\{
\begin{aligned}
&\card\alpha_\pi \; = \; \pi \\
&\alpha<\alpha_\pi \;\; \Longrightarrow \;\; \card\alpha \; < \; \pi \; ,
\end{aligned}
\right. \tag*{(2)} \]
% marginal hand-drawn mark (manuscript) in the left margin beside the following paragraph
or again by the fact that the ordered set $I_{\alpha_\pi}$ of ordinals $\alpha<\alpha_\pi$ is \textit{perfectly well ordered}. Thus, $I_{\alpha_\pi}$ \textit{is large relative to $\pi'$ for every cardinal $\pi'<\pi$} (\footnote[130]{False if $\pi$ is a limit cardinal.}). In conditions b) c) of 4.13.10, it will suffice to assume $I$ stable under $\varinjlim$ of type $I_{\alpha_\pi}$, and the $J_r$ ($r\in A$) stable in $I$ under this type of limits. % [sic] 4.3.10 -> 4.13.10, the conditions b) c) are those of Lemme 4.13.10
\DERmark{ENV1321}


\textbf{Lemma 4.13.10.1.} \textit{Let $\pi$ be an infinite cardinal, $A$ a nonempty set such that} $\card A\leqslant\pi$. \textit{Then}
\msp{page 299}
\textit{there exists a map}
\[ \varphi \; : \; I_{\alpha_\pi} \longrightarrow A \]
\textit{such that for every $r\in A$, $\varphi^{-1}(\{r\})$ is a cofinal subset of $I_{\alpha_\pi}$, or what is more precise,} $\card\varphi^{-1}(\{r\})=\pi$. % [sic] ce que est -> ce qui est, relative pronoun subject

Indeed, it suffices to consider a set $A'$ of cardinality $\pi$ (e.g. $A'=I_{\alpha_\pi}$), the map
\[ \psi \; = \; \mathrm{pr}_1 \; : \; A\times A'\longrightarrow A \; , \]
and to give oneself a bijection
\[ A\times A' \xrightarrow{\ \sim\ } I_{\alpha_\pi} \]
which exists since
\[ \card\,(A\times A') \; = \; \card A\cdot\card A' \; = \; \card A' \; = \; \pi \; = \; \card I_{\alpha_\pi} \; , \]
since $0\neq\card A\leqslant\card A'=\pi$ and $\pi$ infinite. We assume $\varphi$ chosen as in 4.13.10.1, and we will construct the direct system $(i_\alpha)_{\alpha\in I_{\alpha_\pi}}$ in such a way that we have
\[ i_\alpha \; \in \; J_{\varphi(\alpha)} \; . \]

Setting
\[ j \; = \; \varinjlim_{\alpha\in I_{\alpha_\pi}} i_\alpha \; , \]
% foot: Version 17 décembre 2025
\printedpage{112}
% head: A. Grothendieck : Dérivateurs, Ch. XVIII, Édition des Universités de Montpellier II et Paris VII
\msp{page 300}
for every $r$, since
\[ I_r \overset{\mathrm{def}}{=} \varphi^{-1}(r) \; \subset \; I_{\alpha_\pi} \]
\DERmark{ENV1322}

is cofinal in $I_{\alpha_\pi}$, we will therefore have
\[ j \; = \; \varinjlim_{\alpha\in I_r} i_\alpha \; , \]
and by virtue of $(*)$, the $i_\alpha$ for $\alpha\in I_r$ are in $J_r$, so by passage to the limit, using hypothesis c), we find that $j\in J_r$, so
\[ j \; \in \; \bigcap_{r\in A} J_r \; = \; J \; , \]
and since we will choose
\[ i_0 \; = \; i \; , \tag*{$(**)$} \]
we will indeed have found $i\longrightarrow j$ with $j\in J$. So it suffices to construct the direct system $(i_\alpha)$ satisfying $(*)$ and $(**)$. % [sic] avec $i\in J$ -> avec $j\in J$, it is $j$ that was shown to lie in $J$
\DERmark{ENV1323}


If the direct system has been constructed up to the stage $i_\alpha\in\Ob J_{\varphi(\alpha)}$, we continue with
\[ i_\alpha \longrightarrow i_{\alpha+1} \; \in \; \Ob J_{\varphi(\alpha+1)} \; , \]
which is possible since $J_{\varphi(\alpha+1)}$ is
\msp{page 301}
cofinal. So it remains to construct $i_\alpha$ for $\alpha$ a limit ordinal. Now the set of $i_\beta$ for $\beta<\alpha$ is of cardinality $<\pi$, and since by a) $I$ is large relative to $\pi$ (it suffices that it be so relative to every cardinal $<\pi$), there exists an $i\in\Ob I$ which dominates all the $i_\beta$. Let us choose
\[ i_\beta \xrightarrow{v_\beta} i \; . \]

We will not have a priori commutativity in
\[
\begin{tikzcd}
i_{\beta'} \arrow[r, "\tau_{\beta,\beta'}"] \arrow[rr, bend right=15, "v_{\beta'}"'] & i_\beta \arrow[r, "v_\beta"] & i \; ,
\end{tikzcd}
\] % doubt resolved against PDF; see corrections.csv
but since the set of double arrows $i_{\beta'}\mathrel{\overset{v_\beta\tau_{\beta,\beta'}}{\underset{v_{\beta'}}{\rightrightarrows}}} i$ is of cardinality $<\pi$, $I$ being large relative to $\pi$, there exists $i\longrightarrow i'$ which equalizes them all. Since $J_{\varphi(\alpha)}$ is cofinal in $I$, there exists $i''\in J_{\varphi(\alpha)}$, and $i'\longrightarrow i''$. We take $i_\alpha=i''$, with the composite homomorphisms
\DERmark{ENV1324}

\[
\begin{tikzcd}
i_\beta \arrow[r, "v_\beta"] \arrow[rrr, dotted, bend right=12, "\tau_{\alpha,\beta}"'] & i \arrow[r] & i' \arrow[r] & i''
\end{tikzcd}
\] % doubt resolved against PDF; see corrections.csv
% foot: Version 17 décembre 2025
\printedpage{113}
% head: A. Grothendieck : Dérivateurs, Ch. XVIII, Édition des Universités de Montpellier II et Paris VII
\msp{page 302}
as transition morphisms. This completes the proof of 4.13.10 and thereby of 4.13.9, 4.13.8.

\textbf{Remark 4.13.10.2.} Here is what we really had to use, to establish the conclusion of 4.13.10, in terms of $I$, $\pi$, $(J_r)_{r\in A}$ (and without introducing any $\pi'_0$).
\begin{itemize}
\item[\cercle{a$'$}] $I$ is large relative to every cardinal $\pi'<\pi$.
\item[\cercle{b$'$}] $I$ is stable under $\varinjlim$ of type $I_{\alpha_\pi}$.
\item[\cercle{c$'$}] $\forall\, r\in A$, $J_r$ is stable in $I$ under $\varinjlim$ of type $I_{\alpha_\pi}$.
\item[\cercle{d$'$}] $\forall\, r\in A$, $J_r$ is cofinal in $I$.
\item[\cercle{e$'$}] $\card A\leqslant\pi$.
\end{itemize}
(\textbf{N.B.} d$'$) and e$'$) are none other than d) and e) of 4.13.10, whereas a$'$), b$'$), c$'$) are weakened versions of a), b), c) respectively.) % [sic] 4.3.10 -> 4.13.10, the conditions d) e) are those of Lemme 4.13.10
\DERmark{ENV1325}


\textbf{4.13.8 (reformulation).} \textit{Let, for $r\in A$, $\pi''_r\geqslant\pi'_r\geqslant\pi_0$ (= given useful cardinal for $\mathcal{M}$) be as specified in 4.13.2 b) (and c), let $\pi'=\sup_{r\in A}\pi'_r$, $\pi''=\sup_{r\in A}\pi''_r$, so that for every cardinal $\pi$ with $\pi^{\pi'}=\pi\geqslant\pi''$, $\underline{S}_r$ is $\pi$-accessible in $\mathcal{M}$ (and $\pi$ useful for $\underline{S}_r$} and \textit{for $\mathcal{M}$). Let $c$ be a} non-limit \textit{cardinal such that $c\geqslant$} $\card A$ \textit{and $c>\pi'$ (e.g. the successor of} $\sup(\card A,\pi')$\textit{). Let us take $\pi$ such that}
\DERmark{ENV1326}

\[ \pi^{\pi'} \; = \; \pi \; \geqslant \; \sup(\pi'',c) \; . \]
\textit{Then $\underline{S}=\bigcap_{r\in A}\underline{S}_r$ is $\pi$-accessible.}

Indeed, if $\alpha_c$ is the first ordinal of cardinality $c$, since $c$ is not a limit, $I_{\alpha_c}$ is large relative to every cardinal $c'<c$, in particular relative to $\pi'$. Now for every $x\in\Ob\mathcal{M}$, $I=\mathcal{M}_\pi/x$ and the $J_r=\underline{S}_r(\pi)/x$ are stable under $\varinjlim$ large relative to $\pi'$ and of cardinality $\leqslant\pi$,
\DERmark{ENV1327}

\msp{page 302 bis}
and the $J_r$ stable in $I$ for these limits, so this is in particular the case for $\varinjlim$ of type $I_{\alpha_c}$. We can therefore apply result 4.13.10.2 to conclude that if the $J_r$ are cofinal in $I$ (which is large relative to $\pi$, $\pi$ being useful for $\mathcal{M}$), then $J=\underline{S}(\pi)/x$ is so as well. This then implies, as we have seen, that $\underline{S}$ is $\pi$-accessible in $\mathcal{M}$.
\DERmark{ENV1328}


\textbf{4.13.10.3.} We can reformulate the preceding reformulation, using the terminology introduced below (4.13.23):

\textit{Let $c$ be a cardinal, $\alpha_c$ the smallest ordinal $\alpha$ such that} $\card\alpha=c$, \textit{$\pi$ a cardinal $\geqslant c$. We assume}
\begin{itemize}
\item[$(*)$] \textit{$\pi$ is useful for $\mathcal{M}$, and $\mathcal{M}$ is stable under $\varinjlim$ of type $I_{\alpha_c}$. (This second condition is the case e.g. if $c$ is} non-limit \textit{and $c>\pi_0$, $\pi_0$ being a cardinal such that $\mathcal{M}$ satisfies \textit{2}$_{\pi_0}$, e.g. a} useful \textit{cardinal for $\mathcal{M}$ (}\footnote[131]{For then $I_{\alpha_c}$ is large relative to $\pi_0$.}\textit{)).}
\end{itemize}
\textit{Then the set of subcategories $\underline{S}$ of $\mathcal{M}$ which are $\alpha_c$-$\pi$-accessible in $\mathcal{M}$ (i.e. which are $\pi$-accessible and stable in $\mathcal{M}$ under $\varinjlim$ of type $I_{\alpha_c}$) is stable under intersection of cardinality $\leqslant c$. More generally, if $\Phi$ is a set of categories such that} % [sic] accessible -> accessibles, agreement with « sous-catégories »
\begin{itemize}
\item[\textit{1°)}] \textit{$I\in\Phi\Longrightarrow$} $\card I\leqslant\pi$ \textit{and}
\item[\textit{2°)}] \textit{$I_{\alpha_c}\in\Phi$,}
\end{itemize}
\textit{then the set of subcategories of $\mathcal{M}$ which are $\Phi$-$\pi$-accessible in $\mathcal{M}$, is stable under intersection of cardinality $\leqslant c$.}
% foot: Version 17 décembre 2025


\endgroup
% END UNIT der-ch18/der-ch18_p111-115.tex
% BEGIN UNIT der-ch18/der-ch18_p116-120.tex
\begingroup\makeatletter%
\newcommand{\printedpage}[1]{\marginpar{\footnotesize p.~#1}}%
\newcommand{\msp}[1]{\marginpar{\footnotesize [#1]}}%
\newcommand{\ilim}{\varinjlim\limits}%
\newcommand{\uS}{\underline{S}}%
\newcommand{\Ob}{\operatorname{Ob}}%
\newcommand{\circlabel}[1]{\tikz[baseline=(c.base)]\node[draw,circle,inner sep=1pt](c){#1};}%
\makeatother\let\printedpage\derpp\let\msp\dermsp%
\let\mspin\DERmspin%


\printedpage{114}
% head: A. Grothendieck : Dérivateurs, Ch. XVIII, Édition des Universités de Montpellier II et Paris VII
% foot: Version 17 décembre 2025

\msp{page 303}\noindent\textbf{Definition 4.13.11} (\footnote[132]{Recast with 4.13.23 etc.}). Let $\mathcal{M}$ be an accessible category, $\pi_0$ a useful cardinal for $\mathcal{M}$, $\pi > \pi_0$ a cardinal. A category $I$ is said to be $\pi_0$-$\pi$-\textit{adapted} if
\begin{enumerate}
\item[a)] $I$ is equivalent to a small category,
\item[b)] $I$ is large relative to $\pi$, and
\item[c)] $I$ is stable under $\varinjlim$ of cardinality $\leqslant \pi$, and large relative to $\pi_0$ (\footnote[133]{\textbf{N.B.} If $\pi = \pi_0$, condition c) is equivalent to saying that $I$ is stable under direct factors.}).
\end{enumerate}
\noindent A direct system $(x_i)_{i\in I}$ of type $I$ in $\mathcal{M}$ is said to be $\pi_0$-$\pi$-\textit{adapted}, if $I$ is $\pi_0$-$\pi$-adapted, i.e. satisfies conditions a) b) c), and if moreover we have
\begin{enumerate}
\item[d)] the functor $i \longmapsto x_i : I \longrightarrow \mathcal{M}$ commutes with the $\varinjlim$ (large relative to $\pi_0$ and of cardinality $\leqslant \pi$) considered in c), and
\item[e)] the $x_i$ are in $\mathcal{M}_\pi$.
\end{enumerate}

\noindent\textbf{Proposition 4.13.12.} \textit{Let $\mathcal{M}$ be a $\pi_0$-accessible category, and $\pi > \pi_0$ a useful cardinal for $\mathcal{M}$. Then every $x \in \Ob\mathcal{M}$ admits a representation as $x \xleftarrow{\ \sim\ } \ilim_{j\in J} x_j$, where the direct system $(x_j)_{j\in J}$ in $\mathcal{M}$ is $\pi_0$-$\pi$-adapted. More precisely, one can take $J = \mathcal{M}_\pi/x$.}

\msp{page 304}\noindent Indeed, the verification of properties a) to d) for the canonical direct system, with $\varinjlim$ isomorphic to $x$, is immediate.

The preceding proposition does not say whether one can take an ordered set for $J$. In this case, condition c) means that for every subset $J'$ of $J$ large relative to $\pi_0$ and of cardinality $\leqslant \pi$, $J$ admits a least upper bound $i_{J'}$ in $J$, and condition d) that we then have $x_{i_{J'}} \xleftarrow{\ \sim\ } \ilim_{i\in J'} x_i$.
\DERmark{ENV1329}\DERmark{ENV1330}


The natural idea, to construct a $\pi_0$-$\pi$-adapted direct system $(x_j)_{j\in J}$, with $J$ ordered, starting from the one we have canonically, with $J = \mathcal{M}_\pi/x$, is to take for $I$ the set, ordered by inclusion, of subcategories $J'$ of $J_0$ (small full subcategory of $J$ equivalent to $J$) which are of cardinality $\leqslant \pi$ and large relative to $\pi_0$. (Compare (4.7.3.3), page 165). This ordered set is indeed (small and) \textit{large relative to} $\pi$, if we assume that $\pi^{\pi_0} = \pi$ (cf. proof 4.7.5, p. 168-171).

\msp{page 305}\noindent\textbf{Lemma 4.13.13.} \textit{Let $\pi_0$, $\pi > \pi_0$ be infinite cardinals such that $\pi^{\pi_0} = \pi$, and $J$ a small category, large relative to $\pi$. Let $I$ be the set of subcategories $J'$ of $J$ which are large relative to $\pi_0$, and of cardinality $\leqslant \pi$. Then $I$ is $\pi_0$-$\pi$-adapted.}

It remains to see that for every ordered set $A$ large relative to $\pi_0$ and of cardinality $\leqslant \pi$, if $r \longmapsto J_r$ is an increasing family of elements of $I$ indexed by $A$, then
\[
J \ \overset{\mathrm{def}}{=} \ \bigcup_{r\in A} J_r \ = \ \varinjlim_{r\in A} J_r
\]
\DERmark{ENV1331}

is in $I$, i.e. is a subcategory of $J$ of cardinality $\leqslant \pi$ (which is clear) and \textit{large relative to} $\pi_0$. This will follow from the
\DERmark{ENV1332}


\printedpage{115}
% head: A. Grothendieck : Dérivateurs, Ch. XVIII, Édition des Universités de Montpellier II et Paris VII
% foot: Version 17 décembre 2025

\noindent\textbf{Lemma 4.13.13.1.} \textit{Let $(J_r)_{r\in A}$ be a direct system of categories large relative to $\pi_0$, indexed by an ordered set $A$ large relative to $\pi_0$. Then $J = \ilim_{r\in A} J_r$ is large relative to $\pi_0$.}

\msp{page 306}\noindent We must see that if $B$ is a set of cardinality $\leqslant \pi_0$, then the diagonal functor
\[
\delta \ : \ J \longrightarrow J^B \tag*{$(*)$}
\]
is cofinal. Now since $A$ is large relative to $\pi_0$, $\operatorname{card} B \leqslant \pi_0$, products indexed by $B$ commute with the $\varinjlim$ of type $A$ (SGA 4 I 9.8 applied to Cat and to $\pi_0$). Thus the arrow $(*)$ is interpreted as the direct limit of the similar arrows
\[
\delta_r \ : \ J_r \longrightarrow J_r^B \ , \tag*{$(**)$}
\]
which are cofinal, hence $\delta$ is cofinal as a direct limit of cofinal functors.

\noindent\textbf{Lemma 4.13.13.2.} \textit{Under the hypotheses of 4.13.13, let $\mathcal{M}$ be a category satisfying $\mathrm{L}_{\pi_0}$, and $(x_j)_{j\in J}$ (where $J \longrightarrow \mathcal{M}$) a direct system in $\mathcal{M}$, indexed by $J$. For every $J' \in I$, put}
\[
x_{J'} \ = \ \varinjlim_{j\in J'} x_j
\]
\msp{page 307}\textit{(which exists in $\mathcal{M}$, since $\mathcal{M}$ satisfies $\mathrm{L}_{\pi_0}$). Then the $(x_{J'})_{J'\in I}$ form a direct system in $\mathcal{M}$, and $I \longrightarrow \mathcal{M}$ commutes with direct limits large relative to $\pi_0$ and of cardinality $\leqslant \pi$. If $\mathcal{M}$ is $\pi_0$-accessible, and the $x_j$ ($j \in \Ob J$) are in $\mathcal{M}_\pi$, then $(x_{J'})_{J'\in I}$ is a $\pi_0$-$\pi$-adapted direct system in $\mathcal{M}$.}

All this is essentially trivial, and gives:

\noindent\textbf{Corollary 4.13.14} (of 4.13.12). \textit{Under the conditions of 4.13.12, if $\pi^{\pi_0} = \pi$, one can for every $x \in \Ob\mathcal{M}$ find a representation $x \simeq \ilim_{i\in I} x_i$, where $(x_i)_{i\in I}$ is a $\pi_0$-$\pi$-adapted direct system with $I$} an ordered set. \textit{Moreover, if $\pi' > \pi$ and if $x \in \mathcal{M}_{\pi'}$, one can take $I$ such that we have}
\[
\operatorname{card} I \ \leqslant \ {\pi'}^{\pi} \ .
\]

This last addendum is immediate,\msp{page 308} via the usual estimates.

\noindent\textbf{Theorem 4.13.15} (\footnote[134]{Cf. more general and more symmetric version of theorem 4.13.15 in 4.13.22.}). \textit{Let $\mathcal{M}$ be a $\pi_0$-accessible category, and $\pi > \pi'_0 \geqslant \pi_0$ such that $\pi^{\pi_0} = \pi$} (\footnote[135]{One can restrict oneself to assuming $\pi$ useful for $\mathcal{M}$ (instead of $\pi^{\pi_0} = \pi$), then one finds that a, a$'$, b, c are equivalent and imply c$'$, but are perhaps not equivalent to c$'$.})\textit{. Let $\uS$ be a strictly full subcategory of $\mathcal{M}$, and let $\uS(\pi) = \uS \cap \mathcal{M}_\pi$. The following conditions are equivalent.}
\begin{enumerate}
\item[\circlabel{a}] \textit{$\uS$ is $\pi$-accessible in $\mathcal{M}$ (4.13.3) and $\uS$ is stable in $\mathcal{M}$ under direct limits large relative to $\pi'_0$.}
\end{enumerate}

\printedpage{116}
% head: A. Grothendieck : Dérivateurs, Ch. XVIII, Édition des Universités de Montpellier II et Paris VII
% foot: Version 17 décembre 2025

\begin{enumerate}
\item[a$'$)] \textit{$\uS$ is $\pi$-accessible in $\mathcal{M}$, and $\uS(\pi)$ is stable in $\mathcal{M}_\pi$ under direct limits large relative to $\pi'_0$ and of cardinality $\leqslant \pi$} (\footnote[136]{give a name to this last condition, say $\mathrm{L}(\pi'_0, \pi)$ (for $\uS(\pi)$ in $\mathcal{M}_\pi$).})\textit{.}
\item[b)] \textit{For every object $x$ of $\mathcal{M}$, we have this: the strictly full subcategory $\uS(\pi)/x$ of $J = \mathcal{M}_\pi/x$ is cofinal in $\mathcal{M}_\pi/x$ if and only if $x \in \Ob\uS$. Moreover, $\uS(\pi)$ is stable in $\mathcal{M}_\pi$ under $\varinjlim$ large relative to $\pi'_0$ and of cardinality $\leqslant \pi$.}
\item[\circlabel{c}] \textit{For every object $x$ of $\mathcal{M}$ and for every representation $x \simeq \ilim_{i\in I} x_i$, where $(x_i)_{i\in I}$ is a $\pi'_0$-$\pi$-adapted direct system, (4.13.11), we have this: the strictly full subcategory $I_{\uS}$ of $I$ formed by the $i \in \Ob I$ such that $x_i \in \Ob\uS$, is cofinal in $I$ if and only if $x \in \Ob\uS$. Moreover, $\underline{S}(\pi)$ is stable in $\mathcal{M}_\pi$ under $\varinjlim$ large relative to $\pi'_0$ and of cardinality $\leqslant \pi$.} % [sic] S(\pi) -> \underline{S}(\pi), printed without the underline used everywhere else for this subcategory
\end{enumerate}
\DERmark{ENV1333}

\msp{page 309}
\begin{enumerate}
\item[c$'$)] \textit{As in d)} \texttt{[rather c)]}\textit{, with $I$ an ordered set.}
\end{enumerate}
\noindent\textbf{Corollary 4.13.16.} \textit{The hypotheses on $\mathcal{M}$, $\pi_0$, $\uS$ being those of 4.13.15, for there to exist cardinals $\pi > \pi'_0 \geqslant \pi_0$ satisfying the conditions $\pi^{\pi'_0} = \pi$ (a fortiori $\pi^{\pi_0} = \pi$) and the equivalent conditions a) to d$'$)} \texttt{[rather c$'$)]} \textit{of 4.13.15, it (is necessary and it) suffices that $\uS$ be accessible in $\mathcal{M}$. More precisely, if $\uS$ is accessible in $\mathcal{M}$, let $\pi'_0 \geqslant \pi_0$ be a cardinal such that $\uS$ is $\pi'_0$-accessible in $\mathcal{M}$, and let $\pi$ be a cardinal $> \pi'_0$ such that $\pi^{\pi'_0} = \pi$. Then $(\uS, \pi'_0, \pi)$ satisfies the equivalent conditions of 4.13.15.}

Indeed, the ``it is necessary'' is obvious from condition a) of 4.13.15. For the ``it suffices'' we note that $\uS$ $\pi'_0$-accessible in $\mathcal{M}$ implies that $\uS$ is stable in $\mathcal{M}$ under $\varinjlim$ large relative to $\pi'_0$. Thus to verify a), it suffices to verify that $\uS$ is $\pi$-accessible in $\mathcal{M}$. Now since $\uS$ is\msp{page 310} $\pi'_0$-accessible and $\pi > \pi'_0$ satisfies $\pi^{\pi'_0} = \pi$, we conclude from this that $\uS$ is $\pi$-accessible, i.e. $\uS \xleftarrow{\ \sim\ } \mathrm{Ind}_\pi(\uS_\pi)$. Thus by 4.13.3, it remains to see that $\uS_\pi \subset \mathcal{M}_\pi$, or what amounts to the same thing, that $\uS_\pi \subset \uS(\pi) \overset{\mathrm{def}}{=} \uS \cap \mathcal{M}_\pi$. (In fact, we shall even have equality, for it is clear that $\uS(\pi) \subset \uS_\pi \ldots$). Now $\uS(\pi)$ is a subcategory of $\uS$ stable under $\varinjlim$ large relative to $\pi'_0$ and of cardinality $\leqslant \pi$ (since $\uS$ and $\mathcal{M}_\pi$ are stable under these limits \textit{in} $\mathcal{M}$, and these limits in $\uS$ and in $\mathcal{M}$ are likewise), and $\uS(\pi)$ contains $\underline{S}_{\pi'_0}$. Thus $\uS(\pi)$ contains $\uS_\pi$, by 4.7.1 (page 161), implication c) $\Longrightarrow$ b$'$), applied to $\uS$, $\pi'_0$, $\pi$ instead of $\mathcal{M}$, $\pi_0$, $\pi$. This therefore proves corollary 4.13.16. % [sic] S_{\pi'_0} -> \underline{S}_{\pi'_0}, printed without the underline
\DERmark{ENV1334}\DERmark{ENV1335}\DERmark{ENV1336}


\textsc{Proof of} 4.13.15. Here are the $\pm$ tautological implications,
\[
\begin{tikzcd}[row sep=small]
 & & & \mathrm{c}' \\
\mathrm{a} \arrow[r, Rightarrow] & \mathrm{a}' \arrow[r, Rightarrow, "\bullet"] & \mathrm{c} \arrow[ur, Rightarrow] \arrow[dr, Rightarrow] & \\
 & & & \mathrm{b}\ ,
\end{tikzcd}
\tag*{$(*)$}
\]
\msp{page 311}which I am going to spell out a little. We have a $\Longrightarrow$ a$'$ since if a is satisfied, $\uS$ \textit{and} $\mathcal{M}_\pi$ are stable \textit{in} $\mathcal{M}$ under $\varinjlim$ large relative to $\pi'_0$ and of cardinality $\leqslant \pi$, hence also their intersection $\uS(\pi)$,
\printedpage{117}
% head: A. Grothendieck : Dérivateurs, Ch. XVIII, Édition des Universités de Montpellier II et Paris VII
% foot: Version 17 décembre 2025
which is therefore also stable in $\mathcal{M}_\pi$ under the said limits. The implication c $\Longrightarrow$ c$'$ is a tautology, and c $\Longrightarrow$ b as well, in view of 4.13.12, which tells us that the canonical direct system $(x_i)_{i\in I}$, with $I = \mathcal{M}_\pi/x$, with limit $x$ is $\pi_0$-$\pi$-adapted. (\textbf{N.B.} Since $\pi^{\pi_0} = \pi$ and $\pi'_0 \geqslant \pi_0$, we also have $\pi^{\pi_0} = \pi$, hence $\pi$ is useful for $\mathcal{M}$, so that 4.13.12 applies.)
\DERmark{ENV1337}


It remains to prove the ``$\pm$ tautology'' a$'$ $\Longrightarrow$ c. So let us take a $\pi'_0$-$\pi$-adapted representation
\DERmark{ENV1338}

\[
x \ \simeq \ \varinjlim_{i\in I} x_i \ ,
\]
let us prove that
\[
x \ \in \ \Ob\uS \qquad\quad \Longleftrightarrow \qquad\quad I_{\uS} \text{ cofinal in } I.
\]
The implication $\Longleftarrow$ comes from the fact that, $\uS$ being $\pi$-accessible in $\mathcal{M}$, is stable in $\mathcal{M}$ under $\varinjlim$ large relative to $\pi$ (\footnote[137]{We use here, for $\Longleftarrow$, of the hypothesis $(x_i)_{i\in I}$ $\pi_0$-$\pi$-adapted, only the sole fact that $I$ is large relative to $\pi$.}). It remains to prove\msp{page 312} that if $x \in \Ob\uS$, then $I_{\uS}$ is cofinal in $I$. So let us consider the diagram of categories
\[
\begin{tikzcd}
I \arrow[r, "\varphi"] & J = \mathcal{M}_\pi/x \\
I' = I_{\uS} \arrow[u, hook] \arrow[r, "\varphi'"] & J' = \uS(\pi)/x \ , \arrow[u, hook, "\alpha"']
\end{tikzcd}
\]
this is a cartesian square ($I' = I_{\uS}$ is the inverse image under $\varphi : I \longrightarrow J$ of the strictly full subcategory $J'$ of $J$). We have the following properties
\begin{enumerate}
\item[1)] $I$ is large relative to $\pi$, and $\varphi$ cofinal ($\Longrightarrow$ $J$ large relative to $\pi$).
\item[2)] $I$ stable under $\varinjlim$ large relative to $\pi'_0$ and of cardinality $\leqslant \pi$, and $\varphi$ commutes with these $\varinjlim$ .
\item[3)] $J'$ is stable under $\varinjlim$ of the preceding type and $\alpha$ commutes with them (\footnote[138]{This is where the stability hypothesis made in a$'$) comes in, for by a$'$), $\uS$ is $\pi$-accessible in $\mathcal{M}$.}).
\item[4)] $J'$ cofinal in $J$.
\end{enumerate}
\DERmark{ENV1339}

\noindent We have just concluded that $I' = \varphi^{-1}(J')$\msp{page 313} is cofinal in $I$. As it is a full subcategory (and even strictly full), it suffices to prove property F 1, i.e. that for every $i \in \Ob I$, exists $i' \in \Ob I'$ and an arrow $i \longrightarrow i'$.

We are going to construct $i'$ again by transfinite construction, on the ordinal interval $I_{\alpha_\pi}$, by constructing a direct system
\[
(i_\alpha)_{\alpha\in I_{\alpha_\pi}} \ , \qquad\quad i_\alpha \in \Ob I \ ,
\]
of objects of $I$, and a direct system $(j_\alpha)_{\alpha\in I_{\alpha_\pi}}$, $j_\alpha \in \Ob J'$, transition arrows denoted $\tau_{\beta,\alpha} : i_\alpha \longrightarrow i_\beta$, $\tau'_{\beta,\alpha} : j_\alpha \longrightarrow j_\beta$, finally for every $\alpha$ in $I_{\alpha_\pi}$ arrows
\[
\begin{tikzcd}
\varphi(i_\alpha) \arrow[r, "u_\alpha"] \arrow[rr, dotted, bend right, "\varphi(\tau_{\alpha+1,\alpha})"'] & j_\alpha \arrow[r, "v_\alpha"] \arrow[rr, dotted, bend left, "\tau'_{\alpha+1,\alpha}"] & \varphi(i_{\alpha+1})\ ( \arrow[r, dotted, "u_{\alpha+1}"] & j_{\alpha+1}\ )
\end{tikzcd}
\tag*{$(*)$}
\]
\printedpage{118}
% head: A. Grothendieck : Dérivateurs, Ch. XVIII, Édition des Universités de Montpellier II et Paris VII
% foot: Version 17 décembre 2025
in $J$, such that
\[
\left\{
\begin{aligned}
v_\alpha u_\alpha \ &= \ \varphi(\tau_{\alpha+1,\alpha}) \\
u_{\alpha+1} v_\alpha \ &= \ \tau'_{\alpha+1,\alpha}
\end{aligned}
\right.
\qquad\qquad \text{(\footnotemark[139])} \ .
\tag*{$(**)$}
\]
\footnotetext[139]{cf. also $(****)$, next page.}

\msp{page 314}\noindent Let us show how the construction of such a system of objects and arrows with
\[
i_0 \ = \ i \tag*{$(***)$}
\]
will imply the existence of $i' \in \Ob I'$ and of $i \longrightarrow i'$. Indeed, since $I_{\alpha_\pi}$ is large relative to $\pi'_0$ (it is large relative to every cardinal $< \pi$), there exists
\[
i' \ \overset{\mathrm{def}}{=} \ \varinjlim_{\alpha\in I_{\alpha_\pi}} i_\alpha % [sic] i\in I_{\alpha_\pi} -> \alpha\in I_{\alpha_\pi}, the index of the limit is \alpha
\]
\DERmark{ENV1340}\DERmark{ENV1341}

in $I$, and we have a transition arrow
\[
i \ = \ i_0 \longrightarrow i' \ ,
\]
so it remains to prove that $i' \in \Ob I'$, i.e. $\varphi(i') \in \Ob J'$. Now $\varphi$ commuting with the $\varinjlim$ of type $I_{\alpha_\pi}$, we shall have
\[
\varphi(i') \ = \ \varinjlim \varphi(i_\alpha) \ .
\]
\DERmark{ENV1342}

I still need compatibilities
\[
\begin{tikzcd}
\varphi(i_\alpha) \arrow[r, "u_\alpha"] \arrow[d, "\varphi(\tau_{\beta,\alpha})"'] & j_\alpha \arrow[d, "\tau'_{\beta,\alpha}"] & {} \arrow[d, phantom, "\text{and}" description] & j_\alpha \arrow[r, "v_\alpha"] \arrow[d, "\tau'_{\beta,\alpha}"'] & \varphi(i_{\alpha+1}) \arrow[d, "\varphi(\tau_{\beta+1,\alpha+1})"] \\
\varphi(i_\beta) \arrow[r, "u_\beta"] & j_\beta & {} & j_\beta \arrow[r, "v_\beta"] & \varphi(i_{\beta+1})
\end{tikzcd}
\tag*{$(****)$} % [sic] j_\alpha (lower right corner of the first square) -> j_\beta, target of \tau'_{\beta,\alpha} and of u_\beta
\]
\DERmark{ENV1343}

forgotten earlier.

\msp{page 315}\noindent The compatibilities $(****)$ imply that we have homomorphisms
\[
\begin{tikzcd}
\varphi(i') \arrow[r, shift left, "u"] & j \arrow[l, shift left, "v"]
\end{tikzcd}
\qquad\qquad \text{where } j \overset{\mathrm{def}}{=} \varinjlim_{\alpha\in I_{\alpha_\pi}} j_\alpha \quad \text{(limit \textit{in} $J'$, hence in $J$)} \ ,
\]
\DERmark{ENV1344}\DERmark{ENV1345}

and the compatibilities $(**)$ imply that
\[
vu \ = \ \mathrm{id}_{\varphi(i')} \ , \qquad\quad uv = \mathrm{id}_j \ , % [sic] \mathrm{id}_J -> \mathrm{id}_j, uv is an endomorphism of j
\]
\DERmark{ENV1346}\DERmark{ENV1347}

hence that $u$, $v$ are isomorphisms inverse to each other. Since $J'$ is a \textit{strictly} full subcategory of $J$, it does indeed follow that $\varphi(i') \in J'$, i.e. $i' \in I'$, q.e.d.
\DERmark{ENV1348}\DERmark{ENV1349}


So everything comes down to constructing the
\[
\underbrace{i_\alpha}_{\in \Ob I} \ , \ \tau_{\beta,\alpha} \ , \qquad\quad \underbrace{j_\alpha}_{\in \Ob J'} \ , \ \tau'_{\beta,\alpha} \ , \qquad\quad u_\alpha \ , \ v_\alpha \ ,
\]


\endgroup
% END UNIT der-ch18/der-ch18_p116-120.tex
% BEGIN UNIT der-ch18/der-ch18_p121-125.tex
\begingroup\makeatletter%
\newcommand{\printedpage}[1]{\marginpar{\footnotesize p.~#1}}%
\newcommand{\msp}[1]{\marginpar{\footnotesize [#1]}}%
\newcommand{\circled}[1]{\tikz[baseline=(c.base)]\node[draw,circle,inner sep=1pt](c){#1};}%
\newcommand{\ovalbox}[1]{\tikz[baseline=(c.base)]\node[draw,rounded corners=6pt,inner sep=2pt](c){#1};}%
\newcommand{\Ob}{\operatorname{Ob}}%
\newcommand{\card}{\operatorname{card}}%
\newsavebox{\dctauAB}%
\newsavebox{\dctauABone}%
\newsavebox{\dctauAAzero}%
\makeatother\let\printedpage\derpp\let\msp\dermsp%
\let\mspin\DERmspin%


\sbox{\dctauAB}{\circled{$\scriptstyle\tau_{\alpha,\beta}$}}% boxes for the circled labels of ms. 318 (p. 120)
\sbox{\dctauABone}{\circled{$\scriptstyle\tau_{\alpha,\beta+1}$}}
\sbox{\dctauAAzero}{\circled{$\scriptstyle\tau_{\alpha,\alpha_0}$}}
\printedpage{119}
% head: A. Grothendieck : Dérivateurs, Ch. XVIII, Édition des Universités de Montpellier II et Paris VII
satisfying the constraints $(**)$, $(***)$, $(****)$. We are going to make the construction in three steps.

\circled{1} Start:
\[
i_0 \;=\; i\ .
\]
Since $J'$ is cofinal in $J$, exists $j_0\in\Ob J'$ and
\[
\varphi(i_0)\xrightarrow{u_0} j_0\ .
\]

\msp{page 316}

\circled{2} Suppose that the $i_\beta$, $j_\beta$ have been constructed for $\beta<\alpha$, as well as the $\tau_{\beta',\beta}:i_\beta\longrightarrow i_{\beta'}$, $\tau'_{\beta',\beta}:j_\beta\longrightarrow j_{\beta'}$ for $\beta,\beta'<\alpha$, finally the $u_\beta:\varphi(i_\beta)\longrightarrow j_\beta$, $v_\beta:j_\beta\longrightarrow\varphi(i_{\beta+1})$ for $\beta<\alpha$ resp. for $\beta+1<\alpha$, so as to satisfy the constraints $(**)$, $(****)$, i.e. commutativities in
\[
\left\{
\begin{array}{l}
\left.
\begin{array}{ll}
\circled{1} &
\begin{tikzcd}[ampersand replacement=\&]
\varphi(i_\beta) \arrow[r, "u_\beta"] \arrow[rr, bend right=15, "\varphi(\tau_{\beta+1,\beta})"'] \& j_\beta \arrow[r, "v_\beta"] \& \varphi(i_{\beta+1})
\end{tikzcd}
\\[5ex]
\circled{2} &
\begin{tikzcd}[ampersand replacement=\&]
j_\beta \arrow[r, "v_\beta"] \arrow[rr, dotted, bend right=15, "\tau'_{\beta+1,\beta}"'] \& \varphi(i_{\beta+1}) \arrow[r, "u_{\beta+1}"] \& j_{\beta+1}
\end{tikzcd}
\end{array}
\right\}\quad \textit{if}\ \beta+1<\alpha
\\[8ex]
\begin{array}{ll}
\circled{3} &
\begin{tikzcd}[ampersand replacement=\&]
\varphi(i_\beta) \arrow[r, "u_\beta"] \arrow[d, "\varphi(\tau_{\beta',\beta})"'] \& j_\beta \arrow[d, "\tau'_{\beta',\beta}"] \\
\varphi(i_{\beta'}) \arrow[r, "u_{\beta'}"] \& j_{\beta'}
\end{tikzcd}
\qquad\qquad \text{if } \beta<\beta'<\alpha
\end{array}
\\[10ex]
\begin{array}{ll}
\circled{4} &
\begin{tikzcd}[ampersand replacement=\&]
j_\beta \arrow[r, "v_\beta"] \arrow[d, "\tau'_{\beta',\beta}"'] \& \varphi(i_{\beta+1}) \arrow[d, "\varphi(\tau_{\beta'+1,\beta+1})"] \\
j_{\beta'} \arrow[r, "v_{\beta'}"] \& \varphi(i_{\beta'+1})
\end{tikzcd}
\qquad\qquad \text{if } \beta<\beta' \text{ and } \beta'+1<\alpha\ ,
\end{array}
\end{array}
\right.
\]

\msp{page 317}

we must construct
% foot: Version 17 décembre 2025

\printedpage{120}
% head: A. Grothendieck : Dérivateurs, Ch. XVIII, Édition des Universités de Montpellier II et Paris VII
\begin{itemize}[leftmargin=*]
\item[1)] $i_\alpha\in\Ob I$,\quad and a homomorphism $(i_\beta)_{\beta\in I_\alpha}\longrightarrow i_\alpha$, i.e. $\tau_{\alpha,\beta}$ : $i_\beta\longrightarrow i_\alpha$ $(\beta\in I_\alpha)$, compatible with the translations $\tau_{\beta',\beta}$ for $\beta<\beta'<\alpha$.
\item[2)] $j_\alpha\in\Ob J'$,\quad and a homomorphism $(j_\beta)_{\beta\in I_\alpha}\longrightarrow j_\alpha$, i.e. $\tau'_{\alpha,\beta}$ : $j_\beta\longrightarrow j_\alpha$ $(\beta\in I_\alpha)$, compatible with the translations $\tau'_{\beta',\beta}$ for $\beta<\beta'<\alpha$.
\item[3)] An arrow
\[
u_\alpha:\varphi(i_\alpha)\longrightarrow j_\alpha\ .
\]
\item[4)] If $\alpha$ is not a limit ordinal, i.e. $\alpha=\alpha_0+1$, an arrow
\[
v_{\alpha_0}:j_{\alpha_0}\longrightarrow\varphi(i_{\alpha_0+1})=\varphi(i_\alpha)\ ;
\]
\end{itemize}

and this so as to satisfy the constraints \circled{1}\,\circled{2}\,\circled{3}\,\circled{4} when we assume $\beta+1\leqslant\alpha$ (instead of $\beta+1<\alpha$) in \circled{1} and \circled{2}, and $\beta'\leqslant\alpha$ (instead of $\beta'<\alpha$) in \circled{3}, and $\beta'+1\leqslant\alpha$ (instead of $\beta'+1<\alpha$) in \circled{4}. So for the constraints $1^\circ$, $2^\circ$, $4^\circ$, we need take them into account \textit{only} if $\alpha$ is not a limit ordinal, so $\alpha=\alpha_0+1$, and setting there % [sic] il n'y a en tenir -> il n'y a à en tenir, the preposition « à » appears to be missing
\DERmark{ENV1350}

\[
\beta' \;=\; \alpha_0\ .
\]
As for the constraint $3^\circ$, we must take it
\msp{page 318}
into account for $\beta'=\alpha$. So in all, we must require
\[
(3_\alpha)(\beta)\qquad\qquad
\begin{tikzcd}
\varphi(i_\beta) \arrow[r, "u_\beta"] \arrow[d, "\varphi(\usebox{\dctauAB})"'] & j_\beta \arrow[d, "\tau'_{\alpha,\beta}" {draw, circle, inner sep=1pt}] \\
\varphi(i_\alpha) \arrow[r, "u_\alpha" {draw, circle, inner sep=1pt}] & j_\alpha % [sic] j_\beta -> j_\alpha, the target of \tau'_{\alpha,\beta} and u_\alpha is j_\alpha
\end{tikzcd}
\qquad \text{commutative for } \beta<\alpha\ (\footnotemark[140]),
\]
\DERmark{ENV1351}

\footnotetext[140]{\textbf{N.B.} It suffices to require it for $\beta=\alpha_0$, in the case $\alpha=\alpha_0+1$, cf. below, p. 323.}
(by restriction on $\alpha$), and if $\alpha=\alpha_0+1$ is a non-limit ordinal, we must require in addition
\[
(4_\alpha)(\beta)\qquad\qquad
\begin{tikzcd}
j_\beta \arrow[r, "v_\beta"] \arrow[d, "\tau'_{\alpha_0,\beta}"'] & \varphi(i_{\beta+1}) \arrow[d, "\varphi(\usebox{\dctauABone})"] \\
j_{\alpha_0} \arrow[r, "v_{\alpha_0}" {draw, circle, inner sep=1pt}] & \varphi(i_{\alpha_0+1})=\varphi(i_\alpha)
\end{tikzcd}
\qquad \beta<\alpha_0
\]
\[
(1_\alpha)\qquad\qquad
\begin{tikzcd}
\varphi(i_{\alpha_0}) \arrow[r, "u_{\alpha_0}"] \arrow[rr, bend right=15, "\varphi(\usebox{\dctauAAzero})"'] & j_{\alpha_0} \arrow[r, "v_{\alpha_0}" {draw, circle, inner sep=1pt}] & \varphi(i_{\alpha_0+1})=\varphi(i_\alpha)
\end{tikzcd}
\]
% foot: Version 17 décembre 2025

\printedpage{121}
% head: A. Grothendieck : Dérivateurs, Ch. XVIII, Édition des Universités de Montpellier II et Paris VII
\[
(2_\alpha)\qquad\qquad
\begin{tikzcd}
j_{\alpha_0} \arrow[r, "v_{\alpha_0}" {draw, circle, inner sep=1pt}] \arrow[rr, bend right=15, "\tau'_{\alpha,\alpha_0}"' {draw, circle, inner sep=1pt}] & \varphi(i_{\alpha_0+1})=\varphi(i_\alpha) \arrow[r, "u_\alpha" {draw, circle, inner sep=1pt}] & j_\alpha\ .
\end{tikzcd}
\]

In these compatibilities, I have circled the arrows $\tau_{\alpha,\beta}$, $\tau'_{\alpha,\beta}$ $(\beta<\alpha)$, $\Bigl(\ \tau_{\alpha,\beta+1}\ (\beta+1<\alpha)\ \Bigr)$, $v_{\alpha_0}$, $\Bigl(\ \tau_{\alpha,\alpha_0}\ \Bigr)$, $u_\alpha$, $\Bigl(\ \tau'_{\alpha,\alpha_0}\ \Bigr)$, which were introduced in the data 1) 2) 3) 4), page 317. % [sic] donnés -> données, the noun « données » is meant

\msp{page 319}

We are going to distinguish then the two cases.

\textit{2 A)} $\alpha=\alpha_0+1$.

To construct $i_\alpha$ and the $\tau_{\alpha,\beta}:i_\beta\longrightarrow i_\alpha$, it suffices to construct $i_\alpha$ and
\[
\circled{$\tau_{\alpha,\alpha_0}$}\ :\ i_{\alpha_0}\ \longrightarrow\ i_\alpha=i_{\alpha_0+1}\ .
\]
Similarly for $j_\alpha$ and the $\tau'_{\alpha,\beta}$, it suffices to construct
\[
\circled{$\tau'_{\alpha,\alpha_0}$}\ :\ j_{\alpha_0}\ \longrightarrow\ j_\alpha\ .
\]
We must moreover give ourselves
\[
j_{\alpha_0}\ \xrightarrow{\circled{$v_{\alpha_0}$}}\ \varphi(\underbrace{\,i_\alpha\,}_{=i_{\alpha_0+1}})\ \xrightarrow{\circled{$u_\alpha$}} j_\alpha\ ,
\]
so in all two objects $i_\alpha$, $j_\alpha$ (in $I$ and in $J'$ respectively) and four arrows $\tau_{\alpha,\alpha_0}$, $\tau'_{\alpha,\alpha_0}$, $v_{\alpha_0}$, $u_\alpha$. We need four compatibilities (page 318), where the last two $3_\alpha(\beta)$, $4_\alpha(\beta)$ depend on a parameter $\beta<\alpha$, i.e. $\beta\leqslant\alpha_0$, or on $\beta<\alpha_0$.

\msp{page 320 struck out}

\msp{page 321}

We are first going to deal with the first two compatibilities $1_\alpha$, $2_\alpha$, in terms of the datum of
\[
i_{\alpha_0}\ \in\ \Ob I\ ,\qquad j_{\alpha_0}\ \in\ \Ob J'\ ,\qquad u_{\alpha_0}\ :\ \varphi(i_{\alpha_0})\ \longrightarrow\ j_{\alpha_0}\ .
\]
The two desired compatibilities for the two objects $i_\alpha$, $j_\alpha$ and the four arrows $\tau_{\alpha,\alpha_0}$, $\tau'_{\alpha,\alpha_0}$, $v_{\alpha_0}$, $u_\alpha$ are indicated in the diagram
\[
(*)\qquad\qquad
\begin{tikzcd}
\varphi(i_{\alpha_0}) \arrow[r, "u_{\alpha_0}"] \arrow[rr, bend right=15, "\varphi(\tau_{\alpha,\alpha_0})"'] & j_{\alpha_0} \arrow[r, "v_{\alpha_0}\ ?"] \arrow[rr, bend left=20, "\tau'_{\alpha,\alpha_0}"] & \varphi(\underbrace{\,i_\alpha\,}_{?}) \arrow[r, "u_\alpha\ ?"] & j_\alpha\ .
\end{tikzcd}
\]

We note that since $\varphi$ is cofinal, $\exists\ i_\alpha=i\in\Ob I$ and $j_{\alpha_0}\xrightarrow{v_{\alpha_0}}\varphi(i)$. But we want the composite
\[
w\ :\ \varphi(i_{\alpha_0})\ \xrightarrow{u_{\alpha_0}}\ j_{\alpha_0}\ \xrightarrow{v_{\alpha_0}}\ \varphi(i)
\]
to be of the form $\varphi(\tau)$, where $\tau=\tau_{\alpha,\alpha_0}:i_{\alpha_0}\longrightarrow i_\alpha$ is an arrow in $I$. Since $\varphi$ is a cofinal functor of filtering categories, we know
% foot: Version 17 décembre 2025
\printedpage{122}
% head: A. Grothendieck : Dérivateurs, Ch. XVIII, Édition des Universités de Montpellier II et Paris VII
\msp{page 322}
that we can find $i\xrightarrow{\rho}i'$ in $I$, such that the composite arrow $\varphi(i_{\alpha_0})\xrightarrow{w}\varphi(i)\xrightarrow{\varphi(\rho)}\varphi(i')$ is of the form $\varphi(\tau)$. So at the cost of replacing $i$ by $i'$, and $v_{\alpha_0}$ by the composite $v'_{\alpha_0}$ : $j_{\alpha_0}\xrightarrow{v_{\alpha_0}}\varphi(i)\xrightarrow{\varphi(\rho)}\varphi(i')$, we can suppose that $v_{\alpha_0}u_{\alpha_0}$ is of the form $\varphi(\tau)$. This thus gives us in principle the object $i=i_\alpha$, and the arrows $v_{\alpha_0}:j_{\alpha_0}\longrightarrow\varphi(i)$ and $\tau_{\alpha,\alpha_0}:i_{\alpha_0}\longrightarrow i_\alpha$, so as to satisfy the first compatibility in $(*)$. It remains to construct $j_\alpha$, $\tau'_{\alpha,\alpha_0}$, $u_\alpha$ so as to satisfy the second compatibility $(*)$. Now $J'$ being cofinal in $J$, there exists $j_\alpha$ in $J'$ and an arrow $u_\alpha:\varphi(i_\alpha)\longrightarrow j_\alpha$. On the other hand, since $J'\longrightarrow J$ is fully faithful, the composite arrow $j_{\alpha_0}\xrightarrow{v_{\alpha_0}}\varphi(i_\alpha)\xrightarrow{u_\alpha}j_\alpha$ in $J$, between objects of $J'$, comes from a unique arrow $\tau'_{\alpha,\alpha_0}:j_{\alpha_0}\longrightarrow j_\alpha$ in $J'$.

\msp{page 323}

Thus, we have constructed the two objects and the four desired arrows, so as to satisfy the constraints $1_\alpha$, $2_\alpha$ of page 318. Are the constraints $3_\alpha(\beta)$, $4_\alpha(\beta)$ then satisfied automatically?
\[
\begin{tikzcd}
\varphi(i_\beta) \arrow[r, "u_\beta"] \arrow[d] & j_\beta \arrow[d] \\
\varphi(i_{\alpha_0}) \arrow[r, "u_{\alpha_0}"] \arrow[d, "\varphi(\tau_{\alpha,\alpha_0})"'] & j_{\alpha_0} \arrow[d, "\tau'_{\alpha,\alpha_0}"] \\
\varphi(i_\alpha) \arrow[r, "u_\alpha"] & j_\alpha
\end{tikzcd}
\qquad\qquad \beta<\alpha, \text{ i.e. } \beta\leqslant\alpha_0\ .
\]
The upper square is commutative by the induction hypothesis, for the total (composite) square to be commutative, it suffices that the lower square be so. We thus have a double arrow $(\tau'_{\alpha,\alpha_0}u_{\alpha_0},u_\alpha\varphi(\tau_{\alpha,\alpha_0})):\varphi(i_{\alpha_0})\rightrightarrows j_\alpha$, it may not be trivial, but we know, $J'$ being filtering, that we can equalize it by an arrow $j_\alpha\longrightarrow\bar{j}_\alpha$ in $J'$. So at the cost of replacing $j_\alpha$ by % [sic] \varphi(i_\alpha) -> \varphi(i_{\alpha_0}), both composites of the double arrow start at \varphi(i_{\alpha_0})
\DERmark{ENV1352}

\msp{page 324}
$\bar{j}_\alpha$, $u_\alpha$ by the composite $\bar{u}_\alpha:\varphi(i_\alpha)\longrightarrow j_\alpha\longrightarrow\bar{j}_\alpha$ and $\tau'_{\alpha,\alpha_0}$ by the composite $\bar{\tau}'_{\alpha,\alpha_0}:j_{\alpha_0}\longrightarrow j_\alpha\longrightarrow\bar{j}_\alpha$ (without touching $i_\alpha$, $v_{\alpha_0}$, $\tau_{\alpha,\alpha_0}$), we can obtain commutativity, whence $3_\alpha(\beta)$. % [sic] \varphi(i_0) -> \varphi(i_\alpha), u_\alpha has source \varphi(i_\alpha); [sic] \tau_{\alpha_0,\alpha} -> \tau_{\alpha,\alpha_0}, the arrow is written \tau_{\alpha,\alpha_0} everywhere else
\DERmark{ENV1353}\DERmark{ENV1354}


There remains the constraint $4_\alpha(\beta)$. For this, let us consider the composite diagram
\[
\begin{tikzcd}
j_\beta \arrow[r, "v_\beta"] \arrow[d, "\tau'_{\alpha_0,\beta}"'] & \varphi(i_{\beta+1}) \arrow[d, "\varphi(\tau_{\alpha,\beta+1})"] \\
j_{\alpha_0} \arrow[r, "v_{\alpha_0}"] & \varphi(i_{\alpha_0+1})=\varphi(i_\alpha)\ .
\end{tikzcd}
\]

If $\alpha_0$ is not a limit ordinal, but of the form $\alpha_1+1$, we can again reduce to the case where $\beta=\alpha_1$, proceeding as above. But we can proceed in all cases, without such an artifice (\,!), by considering the double arrows
\[
(*)\qquad\qquad
\begin{tikzcd}
j_\beta \arrow[r, shift left, "\varphi(\tau_{\alpha,\beta+1})v_\beta"] \arrow[r, shift right, "v_{\alpha_0}\tau'_{\alpha_0,\beta}"'] & \varphi(i_\alpha)
\end{tikzcd}
\]
% foot: Version 17 décembre 2025
\printedpage{123}
% head: A. Grothendieck : Dérivateurs, Ch. XVIII, Édition des Universités de Montpellier II et Paris VII
deduced from the preceding square. Since
\msp{page 325}
the set of these double arrows is of cardinality $\card\alpha<\pi$, and $J$ is large relative to $\pi$ (it would suffice that it be so for the cardinals $\pi'<\pi$), we can find an arrow $\varphi(i_\alpha)\xrightarrow{u}j$ in $J$ which equalizes all these double arrows. Since $\varphi$ is cofinal, we can suppose that $j$ is of the form $\varphi(i)$. Since $\varphi$ is a cofinal functor of filtering categories, there then exists $i\xrightarrow{\rho}i'$ in $I$, such that the composite $\varphi(i_\alpha)\xrightarrow{u}\varphi(i)\xrightarrow{\varphi(\rho)}\varphi(i')$ is of the form $\varphi(\sigma)$, for $\sigma:i_\alpha\longrightarrow i'$. In all, we have found $\sigma:i_\alpha\longrightarrow i'$, such that $\varphi(\sigma)$ $(=\varphi(\rho)u):\varphi(i_\alpha)\longrightarrow\varphi(i')$ equalizes all the double arrows considered in $(*)$, p. 324. So it is necessary to replace $i_\alpha$ by $\bar{i}_\alpha=i'$, $v_{\alpha_0}$ by the composite $j_{\alpha_0}\longrightarrow\varphi(i_\alpha)\xrightarrow{\varphi(\sigma)}\varphi(\bar{i}_\alpha)$, $\tau_{\alpha,\alpha_0}$
\msp{page 326}
by the composite $\bar{\tau}_{\alpha,\alpha_0}:i_{\alpha_0}\xrightarrow{\tau_{\alpha,\alpha_0}}i_\alpha\xrightarrow{\sigma}\bar{i}_\alpha$.

To sum up, here is the recipe for constructing, first \ovalbox{$i_\alpha$, $v_{\alpha_0}$, $\tau_{\alpha,\alpha_0}$,}and \textit{then} \ovalbox{$j_\alpha$, $u_\alpha$, $\tau'_{\alpha,\alpha_0}$}.

We proceed first as on page 321, 322 to construct the first triple, so as to satisfy the first compatibility $(*)$ (p. 321), i.e. $1_\alpha$. We then ``correct'' this choice of $i_\alpha$, $v_{\alpha_0}$, $\tau_{\alpha,\alpha_0}$, by replacing $i_\alpha$ by $\bar{i}_\alpha$ for $\sigma:i_\alpha\longrightarrow\bar{i}_\alpha$, $v_{\alpha_0}$ by $\bar{v}_{\alpha_0}=\varphi(\sigma)v_{\alpha_0}$, and $\tau_{\alpha,\alpha_0}$ by $\bar{\tau}_{\alpha,\alpha_0}=\sigma\tau_{\alpha,\alpha_0}$, so as to obtain the compatibilities $4_\alpha(\beta)$ for $\beta<\alpha_0$, as we have just indicated. We will no longer touch these choices of $i_\alpha$, $v_{\alpha_0}$, $\tau_{\alpha,\alpha_0}$, in the rest of the construction. We construct $(j_\alpha,u_\alpha,\tau'_{\alpha,\alpha_0})$ so as to satisfy the second compatibility $(*)$ (p. 321), as
\msp{page 327}
indicated on page 322. Then, we correct this choice, by replacing $j_\alpha$ by $\bar{j}_\alpha$, with $\rho:j_\alpha\longrightarrow\bar{j}_\alpha$, and $u_\alpha$ by $\bar{u}_\alpha=\rho u_\alpha$, $\tau'_{\alpha,\alpha_0}$ by $\bar{\tau}'_{\alpha,\alpha_0}=\rho\tau'_{\alpha,\alpha_0}$, so as to satisfy the constraints $3_\alpha(\beta)$ of page 318.

This completes the construction in the case where $\alpha=\alpha_0+1$.

\textit{Case 2 B: $\alpha$ is a limit ordinal.}

So we must construct the data 1) 2) 3) of page 317 (to the exclusion of datum 4), so as to satisfy the sole constraint $3_\alpha(\beta)$.

We begin to construct $i_\alpha$ and the $\tau_{\alpha,\beta}$ in $I$. This is possible, for $I_\alpha=\{\beta\mid\beta<\alpha\}$ is of cardinality $\pi_\alpha<\pi$ – it therefore suffices to suppose $I$ large relative to the cardinals $\pi'<\pi$.
\DERmark{ENV1355}


We then likewise construct $j_\alpha$ in $J'$ and the $\tau'_{\alpha,\beta}$ $(\beta<\alpha)$, using the % [sic] J_\alpha -> j_\alpha, the object of J' is meant
\DERmark{ENV1356}\DERmark{ENV1357}

\msp{page 328}
fact that $J'$ is large relative to $\card\alpha=\card I_\alpha$. We must moreover construct an arrow
\[
u_\alpha\ :\ \varphi(i_\alpha)\ \longrightarrow\ j_\alpha\ .
\]
This is possible, for $J'$ being cofinal in $J$, there exists $j'$ in $J'$ and $\varphi(i_\alpha)\longrightarrow j'$, and $J'$ being filtering, there exists $\bar{j}_\alpha$ and arrows % [sic] \varphi(i_\alpha)\longrightarrow J' -> \varphi(i_\alpha)\longrightarrow j', the arrow goes to the object j'
\DERmark{ENV1358}

$\begin{tikzcd}[ampersand replacement=\&, baseline=(\tikzcdmatrixname-1-1.base)]
j_\alpha \arrow[rr] \& \& \bar{j}_\alpha \\
\& j' \arrow[ur] \&
\end{tikzcd}$, so at the cost of replacing $j_\alpha$ by $\bar{j}_\alpha$, the $\tau'_{\alpha,\beta}$ by $\bar{\tau}'_{\alpha,\beta}$, we do indeed find an arrow $u_\alpha:\varphi(i_\alpha)\longrightarrow j_\alpha$. We must then deal with the % [sic] target j_\alpha in the small diagram -> \bar{j}_\alpha, the two arrows go to the object \bar{j}_\alpha just introduced
\DERmark{ENV1359}

% foot: Version 17 décembre 2025


\endgroup
% END UNIT der-ch18/der-ch18_p121-125.tex
% BEGIN UNIT der-ch18/der-ch18_p126-130.tex
\begingroup\makeatletter%
\newcommand{\printedpage}[1]{\marginpar{\footnotesize p.~#1}}%
\newcommand{\msp}[1]{\marginpar{\footnotesize [#1]}}%
\DeclareRobustCommand{\card}{\operatorname{card}}%
\DeclareRobustCommand{\Ob}{\operatorname{Ob}}%
\DeclareRobustCommand{\Ind}{\operatorname{Ind}}%
\makeatother\let\printedpage\derpp\let\msp\dermsp%
\let\mspin\DERmspin%


\printedpage{124}
% head: A. Grothendieck : Dérivateurs, Ch. XVIII, Édition des Universités de Montpellier II et Paris VII

compatibilities $3_\alpha(\beta)$, hence double arrows
\[
\begin{tikzcd}
\varphi(i_\beta) \arrow[r, shift left, "\tau'_{\alpha,\beta}u_\beta"] \arrow[r, shift right, "u_\alpha\varphi(\tau_{\alpha,\beta})"'] & j_\alpha \; .
\end{tikzcd}
\]
As the set of these double arrows has cardinality $\card\alpha < \pi$, and $J'$ is large relative to the said cardinal, we
\msp{page 329}
thus find an arrow
\[
\rho : j_\alpha \longrightarrow j
\]
in $J$ which equalizes them all. As $J'$ is cofinal, we can suppose $j$ in $J'$, say $\bar{j}_\alpha$. Thus, at the cost of replacing $j_\alpha$ by $\bar{j}_\alpha$, and the $\tau'_{\alpha,\beta}$ by $\bar{\tau}'_{\alpha,\beta} = \rho\tau'_{\alpha,\beta}$, finally $u_\alpha : \varphi(i_\alpha) \longrightarrow j_\alpha$ by $\bar{u}_\alpha = \rho u_\alpha : \varphi(i_\alpha) \longrightarrow \bar{j}_\alpha$, we satisfy all the desired constraints.

I summarize the result (``$\pm$ tautological''!) proved on pages 313-329.

% a vertical change bar is printed in the left margin alongside the statement from "Lemme-clef" to "suivantes :"
\textbf{Key lemma 4.13.17.} \textit{Let us consider a cartesian square in} Cat
\[
\begin{tikzcd}
I \arrow[r, "\varphi"] & J \\
I' \arrow[u, hook] \arrow[r] & J' \arrow[u, hook, "\psi"'] \; .
\end{tikzcd}
\]

\textit{We make the following assumptions:}

\msp{page 330}

\begin{enumerate}[label={}]
\item[\textit{1)}] \textit{$\varphi$ and $\psi$ are cofinal, $\psi$ is the inclusion of a strictly full subcategory $J'$ of $J$, $I$ is filtering (hence $J$ and $J'$ as well).}
\end{enumerate}

\textit{We seek conditions for the inclusion of the strictly full subcategory $I' \longrightarrow I$ to be cofinal (in which case we see at once that $I'$ is filtering, and $I' \longrightarrow J'$ is also cofinal). For this, we suppose that there exists a \textup{limit ordinal} $\alpha_0$ having the following additional properties:}

\begin{enumerate}[label={}]
\item[\textit{2)}] \textit{$I$ and $J'$ are stable under $\varinjlim$ of type $I_{\alpha_0}$ $(= \{\beta \mid \beta \ \textit{ordinal}\ < \alpha_0\})$, and the functors $\varphi$, $\psi$ commute with these limits.}
\item[\textit{3)}] \textit{For every $\alpha \in I_{\alpha_0}$, $I$ is large relative to $\card\alpha$. (Thus, $\varphi$ being cofinal, $J$ is also large relative to $\card\alpha$, and likewise, $\psi$ being fully faithful, and cofinal, $J'$ is large relative to $\card\alpha$.)}
\end{enumerate}

\textit{Under these conditions, $I'$ is cofinal in $I$ (hence also large relative to $\card\alpha$, for every $\alpha \in I_{\alpha_0}$).}

The proof is the one given above. We did not have to use that the ordinal $\alpha_0 = \alpha_\pi$ was of this particular kind, only that it was a limit ordinal.
\DERmark{ENV1360}


\printedpage{125}
% head: A. Grothendieck : Dérivateurs, Ch. XVIII, Édition des Universités de Montpellier II et Paris VII

\msp{page 331}

This lemma implies, as we saw on page 312, the implication a$'$ $\Longrightarrow$ c in the diagram of implications $(*)$ on page 310, which is thus now proved.
\DERmark{ENV1361}


To finish proving theorem 4.13.15, it remains to prove the implications
\[
\begin{tikzcd}[row sep=small]
\mathrm{c}' \arrow[dr, Rightarrow] & & \\
& \mathrm{a}' \arrow[r, Rightarrow] & \mathrm{a} \; . \\
\mathrm{b} \arrow[ur, Rightarrow] & &
\end{tikzcd}
\]

Let us therefore prove (b or c$'$) $\Longrightarrow$ a$'$. It remains to prove that $\underline{S}$ is $\pi$-accessible in $\mathcal{M}$, i.e. that the fully faithful functor
\[
\Ind_\pi(\underline{S}(\pi)) \longrightarrow \mathcal{M} \tag*{$(*)$}
\]
has essential image $\underline{S}$. Thus we must prove that if $x \in \Ob\mathcal{M}$, then
\[
x \in \Ob\underline{S} \quad \overset{(?)}{\Longleftrightarrow} \quad
\begin{minipage}[c]{0.55\textwidth}
$x$ is in the essential image of $\Ind_\pi(\underline{S}(\pi))$, i.e. $x$ is isomorphic to $\varinjlim\limits_{i\in I} x_i$, the $x_i$ in $\Ob\underline{S}(\pi)$, $I$ large relative to $\pi$.
\end{minipage}
\tag*{$(*)$}
\]

\msp{page 332}

$\Bigg($ Let us first suppose that $\underline{S}$ satisfies b) (\footnote[141]{Clumsy, we can treat the two cases b) and c) at the same time, cf. bottom of page 333.}). Thus we know that
\[
x \in \Ob\underline{S} \qquad \Longleftrightarrow \qquad \underline{S}(\pi)/x \text{ is cofinal in } \mathcal{M}_\pi/x \; . \tag*{$(**)$}
\]
\DERmark{ENV1362}


We want to deduce the equivalence $(*)$ from the equivalence $(**)$.

Let us prove $\Longrightarrow$ in $(*)$, i.e. let us suppose $x \in \Ob\underline{S}$, and prove that $x$ is in the essential image of $\Ind_\pi\underline{S}(\pi)$. By virtue of $(**)$, $\underline{S}(\pi)/x$ is a subcategory (strictly full) cofinal in $\mathcal{M}_\pi/x$, thus since $\mathcal{M}_\pi/x$ is large relative to $\pi$, $\underline{S}(\pi)/x$ is as well. As we know that
\[
x \ = \ \varinjlim_{i\in\mathcal{M}_\pi/x} x_i
\]
($\pi$ being a useful cardinal for $\mathcal{M}$), it follows by cofinality that we also have
\[
x \ = \ \varinjlim_{i\in\underline{S}(\pi)/x} x_i \; ,
\]
and since $\underline{S}(\pi)/x$ is large relative to $\pi$, this implies that $x$ belongs to the essential image of $\Ind_\pi(\underline{S}(\pi))$.

Let us prove $\Longleftarrow$ in $(*)$, i.e. let us suppose that $x$ is in the essential image of $\Ind_\pi(\underline{S}(\pi))$,
\printedpage{126}
% head: A. Grothendieck : Dérivateurs, Ch. XVIII, Édition des Universités de Montpellier II et Paris VII
\msp{page 333}
\[
x \ = \ \varinjlim_{i\in I} x_i \; , \qquad\quad \text{the } x_i \text{ in } \underline{S}(\pi),\ I \text{ large relative to } \pi.
\]

Then the composite functor $i \longmapsto x_i$,
\[
I \ \longrightarrow \ \underline{S}(\pi)/x \ \hookrightarrow \ \mathcal{M}_\pi/x
\]
is cofinal (cf. 4.4.2, page 79), hence $\underline{S}(\pi)/x \longrightarrow \mathcal{M}_\pi/x$ satisfies condition F 1 (since the composite satisfies it), and since $\underline{S}(\pi)/x \longrightarrow \mathcal{M}_\pi/x$ is fully faithful, this already implies that this functor is cofinal. Thus by $(**)$ we conclude from this $x \in \Ob\underline{S}$, q.e.d.

Let us now suppose that $S$ satisfies c$'$ (instead of b), and let us again prove the equivalence $(*)$. For this, let us write
\[
\square \qquad x \ = \ \varinjlim_{i\in I} x_i \qquad\quad
\left\{
\begin{minipage}[c]{0.5\textwidth}
the $x_i \in \Ob\mathcal{M}_\pi$, $I$ \textit{ordered set}, $(x_i)_{i\in I}$ is $\pi_0$-$\pi$-adapted (\footnotemark[142]).
\end{minipage}
\right.
\]
\footnotetext[142]{\textbf{N.B.} That exists, thanks to 4.13.14.}

which is possible since $\pi$ is useful for $\mathcal{M}$ $\Bigg)$. Let us define $I_{\underline{S}}$ as in the statement of c), hence as the inverse image of $\underline{S}(\pi)/x \hookrightarrow \mathcal{M}_\pi/x$ under $I \longrightarrow \mathcal{M}_\pi/x$. By hypothesis d$'$) \texttt{[rather c$'$)]}, we have
\DERmark{ENV1363}

\[
x \in \Ob\underline{S} \qquad \Longleftrightarrow \qquad I_{\underline{S}} \text{ is cofinal in } I \; . \tag*{$(**')$}
\]

Using $(**')$, let us again prove $(*)$ (page 331). In fact, we shall use only that we have
\msp{page 334}
the isomorphism $\square$, with $I$ large relative to $\pi$, and that we have the equivalence $(**')$, without having to use that $I$ is ordered. Thus the proof also holds in case c).

Let us prove $\Longrightarrow$ in $(*)$, i.e. let us suppose $x \in \Ob\underline{S}$, let us prove that $x$ is in the essential image of $\Ind_\pi(\underline{S}(\pi))$. By virtue of $(**')$, $I_{\underline{S}}$ is cofinal in $I$, and since $I_{\underline{S}} \longrightarrow I$ is a fully faithful inclusion, it follows that $I_{\underline{S}}$ is also large relative to $\pi$, on the other hand we shall have by cofinality
\[
x \ = \ \varinjlim_{i\in I_{\underline{S}}} x_i \; ,
\]
where this time the $x_i$ are in $\underline{S}(\pi) = \underline{S}\cap\mathcal{M}_\pi$, whence it follows that $x$ is in the essential image of $\Ind_\pi(\underline{S}(\pi))$.

Let us prove $\Longleftarrow$ in $(*)$, i.e. let us suppose that we have
\[
x \ = \ \varinjlim_{j\in J} x_j \; , \qquad\quad \text{the } x_j \text{ in } \underline{S}(\pi),\ J \text{ large relative to } \pi \; .
\]

\printedpage{127}
% head: A. Grothendieck : Dérivateurs, Ch. XVIII, Édition des Universités de Montpellier II et Paris VII

\msp{page 335}

Invoking 4.4.2 once again, we find that $J \longrightarrow \mathcal{M}_\pi/x$ is cofinal, and since this functor factors through the fully faithful inclusion functor $\underline{S}(\pi)/x \hookrightarrow \mathcal{M}_\pi/x$, it follows as before that the latter is cofinal. We shall admit for a moment that this implies that the inverse image $I_{\underline{S}}$ of this strictly full subcategory, under the likewise cofinal functor $I \longrightarrow \mathcal{M}_\pi/x$, is cofinal in $I$. This being so, the equivalence $(**')$ (p. 333) implies $x \in \Ob\underline{S}$, which is what we wanted.

Thus it remains to prove that if
\[
J' \ = \ \underline{S}(\pi)/x \hookrightarrow \mathcal{M}_\pi/x \ = \ J
\]
is cofinal, then $I' = I_{\underline{S}} \longrightarrow I$ is as well. It seems that we really do have to invoke again the stability conditions $\mathrm{L}(\pi_0,\pi)$ made in d$'$) \texttt{[rather c$'$)]} which had not been necessary for case c) and
\msp{page 336}
the (a bit) nasty key lemma 4.13.17, where we shall take (as in the proof written
before the statement of this lemma) % [sic] cet lemme -> ce lemme, demonstrative adjective "cet" before a consonant
$\alpha = \alpha_\pi \overset{\mathrm{def}}{=}$ smallest ordinal $\alpha$ such that $\card\alpha = \pi$. Thus $I_\alpha$ has cardinality $\pi$, and is large relative to every cardinal $\pi' < \pi$, and in particular relative to $\pi'_0$. Assumption $2^\circ$) of 4.13.17 is satisfied for $I$, by definition of a $\pi_0$-$\pi$-adapted direct system, and for $J'$, thanks to assumption $\mathrm{L}(\pi_0,\pi)$ on $\underline{S}(\pi)$. Likewise, assumption $3^\circ$) ($I$ large relative to every cardinal $< \pi$) is satisfied by definition of $\pi_0$-$\pi$-adapted systems -- $I$ is even large relative to $\pi$. Assumption $1^\circ$) has already been seen.
\DERmark{ENV1364}


We have thus proved (by a circular route)
\[
\begin{tikzcd}[row sep=small]
& & & \mathrm{c}' & \\
\mathrm{a} \arrow[r, Rightarrow] & \Bigg(\mathrm{a}' \arrow[r, Leftrightarrow] & \mathrm{c} \arrow[ur, Leftrightarrow] \arrow[dr, Rightarrow] & & \Bigg)\, , \\ % Verified against PDF leaf 9: the oblique equivalence has both arrowheads.
& & & \mathrm{b} &
\end{tikzcd}
\]

\texttt{[and also b $\Longrightarrow$ c,]} having to invoke the nasty key lemma for the \textit{two} implications
\[
\mathrm{a}' \ \Longrightarrow \ \mathrm{c} \qquad\quad \text{and} \qquad\quad \mathrm{c}' \ \Longrightarrow \ \mathrm{a}'
\]
(hence for two implications \textit{in opposite directions}, a strange thing \dots).

\msp{page 337}

It only remains still to prove the implication
\[
\mathrm{a}' \ \Longrightarrow \ \mathrm{a} \; ,
\]
which is essentially sorital.

This completes the proof of the accessibility criterion 4.13.15 for a full subcategory $\underline{S}$ of $\mathcal{M}$.

We are going to give a somewhat decanted variant of 4.13.15, which (with slightly weaker assumptions) captures the essential part of its substance. Statement 4.13.15 is an almost immediate consequence of it:

\printedpage{128}
% head: A. Grothendieck : Dérivateurs, Ch. XVIII, Édition des Universités de Montpellier II et Paris VII

\msp{page 338}

\textbf{Corollary 4.13.18} (of 4.13.15). \textit{Let $\mathcal{M}$ be an accessible category, $\pi$ a useful cardinal for $\mathcal{M}$, $\alpha_0$ a \textup{limit ordinal} such that we have $\card\alpha_0 \leqslant \pi$. Let $I_{\alpha_0} = \{\alpha \mid \alpha < \alpha_0\}$, and suppose that $\mathcal{M}$ is stable under $\varinjlim$ of type $I_{\alpha_0}$.}

\textit{Let then $\underline{S}$ be a strictly full subcategory of $\mathcal{M}$, $\underline{S}(\pi) = \underline{S}\cap\mathcal{M}_\pi$. Let us denote by $L(\alpha_0)$ the condition: $\underline{S}(\pi)$ is stable in $\mathcal{M}_\pi$ (or what amounts to the same thing, in $\mathcal{M}$) under $\varinjlim$ of type $I_{\alpha_0}$. Then the following conditions are equivalent.}

\begin{enumerate}[label={}]
\item[\textit{a)}] \textit{$\underline{S}$ is $\pi$-accessible in $\mathcal{M}$, and is stable in $\mathcal{M}$ under $\varinjlim$ of type $I_{\alpha_0}$.}
\item[\textit{a$'$)}] \textit{$\underline{S}$ is $\pi$-accessible and $\underline{S}(\pi)$ satisfies $L(\alpha_0)$.}
\item[\textit{b)}] \textit{$\underline{S}(\pi)$ satisfies $L(\alpha_0)$, and for every $x \in \Ob\mathcal{M}$, we have this: $x \in \Ob\underline{S}$ if and only if the strictly full subcategory $\underline{S}(\pi)/x$ of $\mathcal{M}(\pi)/x$ is cofinal in $\mathcal{M}(\pi)/x$} \texttt{[read $\mathcal{M}_\pi/x$]}.
\item[\textit{c)}] \textit{$\underline{S}(\pi)$ satisfies $L(\alpha_0)$, moreover for every $x \in \Ob\mathcal{M}$, and every representation}
\[
x \ \simeq \ \varinjlim_{i\in I} x_i
\]
\msp{page 339}
\textit{which is ``$\alpha_0$-$\pi$-adapted'', i.e. such that}
\begin{enumerate}[label={}]
\item[\textit{$1^\circ$)}] \textit{$I$ equivalent to a small category, and large relative to $\pi$,}
\item[\textit{$2^\circ$)}] \textit{the $\varinjlim$ of type $I_{\alpha_0}$ exist in $I$,}
\item[\textit{$3^\circ$)}] \textit{the functor $I \longrightarrow \mathcal{M}$, $i \longmapsto x_i$, commutes with the $\varinjlim$ of type $I_{\alpha_0}$, and finally}
\item[\textit{$4^\circ$)}] \textit{the $x_i \in \Ob\mathcal{M}_\pi$}
\end{enumerate}
\textit{(compare 4.13.11), we have this (where we denote by $I_{\underline{S}}$ the strictly full subcategory of $I$ formed by the $i$ such that $x_i \in \Ob\underline{S}$): we have $x \in \Ob\underline{S}$ if and only if $I_{\underline{S}}$ is cofinal in $I$.}
\item[\textit{c$'$)}] \textit{(\textup{If} $I_{\alpha_0}$ is large relative to $\pi_0$, and $\pi^{\pi_0} = \pi$, to be able to invoke 4.13.14, cf. p. 335.) As in c), but assuming $I$ ordered.}
\end{enumerate}
(\footnote[143]{\textbf{N.B.} If we do not suppose $I_{\alpha_0}$ large relative to $\pi_0$ and % [sic] suppose par -> suppose pas, negation "ne ... pas"
$\pi^{\pi_0} = \pi$, then all the conditions from a) to c) (a, a$'$, b, c) are equivalent and \textit{imply} c, but are perhaps not equivalent to it.}).

\textsc{Proof}. It is carried out following the same logical diagram as that of 4.13.15,
\[
\begin{tikzcd}[row sep=small]
& & & \mathrm{c}' \arrow[dr, Rightarrow, "\textbf{!}"] & \nearrow\ \begin{tabular}[b]{@{}r@{}}\footnotesize here, we must suppose\\ \footnotesize $I_{\alpha_0}$ large relative to $\pi_0$\\ \footnotesize and $\pi^{\pi_0} = \pi$.\end{tabular} & \\
\mathrm{a} \arrow[r, Rightarrow] & \mathrm{a}' \arrow[r, Rightarrow, "\textbf{!}"] & \mathrm{c} \arrow[ur, Rightarrow] \arrow[dr, Rightarrow] & & \mathrm{a}' \arrow[r, Rightarrow] & \mathrm{a} \; , \\
& & & \mathrm{b} \arrow[ur, Rightarrow] & &
\end{tikzcd}
\]
where we have marked with the sign\,! the two implications which are ``nasty'', in the sense that they call on the key lemma 4.13.17 (pages 329-330), which is proved by a somewhat twisted transfinite construction. The proof of the seven implications of the diagram is essentially
the one given on the preceding pages 311-337, % [sic] celle donné -> celle donnée, agreement with feminine "celle"


\endgroup
% END UNIT der-ch18/der-ch18_p126-130.tex
% BEGIN UNIT der-ch18/der-ch18_p131-135.tex
\begingroup\makeatletter%
\newcommand{\printedpage}[1]{\marginpar{\footnotesize p.~#1}}%
\newcommand{\msp}[1]{\marginpar{\footnotesize [#1]}}%
\newcommand{\Ind}{\operatorname{Ind}}%
\newcommand{\Ob}{\operatorname{Ob}}%
\newcommand{\Ord}{\operatorname{Ord}}%
\newcommand{\card}{\operatorname{card}}%
\newcommand{\cardred}{\mathop{\textup{card\,red}}\nolimits}%
\newcommand{\cercle}[1]{\tikz[baseline=(c.base)]\node[draw,circle,inner sep=1pt](c){\textup{#1}};}%
\makeatother\let\printedpage\derpp\let\msp\dermsp%
\let\mspin\DERmspin%


\setcounter{footnote}{143}

\printedpage{129}
% head: A. Grothendieck : Dérivateurs, Ch. XVIII, Édition des Universités de Montpellier II et Paris VII
% foot: Version 17 décembre 2025

\msp{page 340}
with the exception of a$'$ $\Longrightarrow$ a, which is a little simpler in the present case. We are going to bring it out in the form of a somewhat more general lemma. (\textbf{N.B.} The proof of a$'$ $\Longrightarrow$ a in 4.13.15 had been omitted as ``sorital''\,!, cf. page 337\,\dots)

\textbf{Lemma 4.13.19.} \textit{Let $\mathcal{M}$ be an accessible category, and $\underline{S}$ an accessible subcategory in $\mathcal{M}$, $\pi$ a useful cardinal for $\mathcal{M}$ and such that $\underline{S}$ is $\pi$-accessible, i.e. $\underline{S}$ is the essential image of $\Ind_\pi(\underline{S}(\pi))$ under the fully faithful canonical functor $\Ind_\pi(\underline{S}(\pi))\longrightarrow\mathcal{M}$, where $\underline{S}(\pi)=\underline{S}\cap\mathcal{M}_\pi$. We assume that there exists a limit ordinal $\alpha$, such that $\card\alpha\leqslant\pi$ and such that $\mathcal{M}_\pi$ and $\underline{S}(\pi)$ $(=\underline{S}_\pi)$ are stable under $\varinjlim$ of type $I_\alpha$, and that the inclusion functor $\underline{S}(\pi)\longrightarrow\mathcal{M}_\pi$ commutes with them. Let $\Phi$ be a set of categories such that $\forall\ I\in\Phi$, we have $\cardred I\leqslant\pi$, and suppose that $\mathcal{M}_\pi$ is stable under $\varinjlim$ of type $\Phi$, and that $\underline{S}(\pi)$ is stable in $\mathcal{M}_\pi$ under the said limits (or what amounts to the same thing, that it is stable in itself under these $\varinjlim$\,, and that the inclusion functor commutes with them). Then $\mathcal{M}$ and $\underline{S}$ are stable under $\varinjlim$ of type $\Phi$, and the inclusion functor $\underline{S}\longrightarrow\mathcal{M}$ commutes with them, i.e. $\mathcal{M}$ is stable under these limits, and $\underline{S}$ stable in $\mathcal{M}$ under these same limits.} % [sic] « soient stables par lim de type I_alpha » printed twice -> once, repetition
\DERmark{ENV1365}\DERmark{ENV1366}


\msp{page 341}
\textsc{Proof}. We apply 4.11.1 (p. 267-268) separately to $\mathcal{M}\simeq\Ind_\pi\mathcal{M}_\pi$ and to $\underline{S}\simeq\Ind_\pi\underline{S}(\pi)$, to conclude that $\mathcal{M}$ and $\underline{S}$ are separately stable under $\varinjlim$ of type $\Phi$. It remains to see that the functor $\underline{S}\longrightarrow\mathcal{M}$ commutes with them. This is seen as in the proof of 4.11.1, more generally we get this:

\textbf{Lemma 4.13.20} (cor. to 4.11.1). \textit{Let $f_0 : C\longrightarrow C'$ be a functor between categories equivalent to small categories, $\pi$ an infinite cardinal, $f:\mathcal{M}\stackrel{\mathrm{def}}{=}\Ind_\pi C\longrightarrow\mathcal{M}'\stackrel{\mathrm{def}}{=}\Ind_\pi C'$ the canonical extension $\Ind_\pi(f_0)$ of $f_0$. Let $\Phi$ and $\alpha$ be as in 4.11.1, and suppose that $I_\alpha\in\Phi$ and that $f_0$ commutes with $\varinjlim$ of type $\Phi$. Then $f$ commutes with these same direct limits.} (\footnote{It is better to take functors $f_0 : C\longrightarrow\mathcal{N}$, and $\mathcal{N}$ [?] that $f_0$ commutes with $\varinjlim$ of type $\Phi$, so the same is true of its extension $\Ind_\pi(C,\mathcal{N})\longrightarrow\mathcal{N}$. (\textbf{N.B.} We assume that $\mathcal{N}$ satisfies $\mathrm{L}_\pi$.)}).
\DERmark{ENV1367}\DERmark{ENV1368}


\textbf{Corollary 4.13.21.} \textit{Under the conditions of 4.13.20, suppose that $\Phi$ contains the categories of cardinality $\leqslant\pi$ (resp. the filtering categories of cardinality $\leqslant\pi$, resp. the filtering categories large relative to $\pi_0$ of cardinality $\leqslant\pi$, resp. the discrete categories of cardinality $\leqslant\pi$). Then $f$ commutes with small $\varinjlim$ (resp. with small filtering $\varinjlim$, resp. with small filtering $\varinjlim$ large relative to $\pi_0$, resp. with small sums). (Converse clear and ``left to the reader''.)}

Seen as in 4.4.14 and 4.4.14.1 (p. 97, 98).

\msp{page 342}
Here at last, it seems to me, is the definitive form (\,?) of 4.13.15, 4.13.18.

\textbf{Theorem 4.13.22} (\footnote{cf. variant 4.13.36, page 358.}). \textit{Let $\mathcal{M}$ be an accessible category, $\pi$ an infinite cardinal, useful for $\mathcal{M}$, and $\Phi$ a set of categories satisfying the two conditions:}
\begin{itemize}[leftmargin=2.5em]
\item[\textit{1°)}] $J\in\Phi$ \textit{implies} $\cardred J\leqslant\pi$.
\item[\textit{2°)}] \textit{There exists a} limit ordinal $\alpha_0$ \textit{with} $I_{\alpha_0}\in\Phi$.
\end{itemize}
\textit{We assume $\mathcal{M}$ stable under $\varinjlim$ of type $\Phi$. Let $\underline{S}$ be a} strictly full \textit{subcategory of $\mathcal{M}$, and $\underline{S}(\pi)=\underline{S}\cap\mathcal{M}_\pi$. Let us denote by $L(\Phi)$ the condition on $\underline{S}(\pi)$: $\underline{S}(\pi)$ is stable in} % [sic] L(\varphi) -> L(\Phi), the condition is later written L(\Phi)
\DERmark{ENV1369}

\printedpage{130}
% head: A. Grothendieck : Dérivateurs, Ch. XVIII, Édition des Universités de Montpellier II et Paris VII
% foot: Version 17 décembre 2025
\textit{$\mathcal{M}$ (or what amounts to the same thing, in $\mathcal{M}_\pi$) under $\varinjlim$ of type $\Phi$. This being said, the following conditions on $\underline{S}$ are equivalent:} % [sic] où ce qui revient -> ou ce qui revient, conjunction "ou"
\DERmark{ENV1370}

\begin{itemize}[leftmargin=2.5em]
\item[\cercle{a}] \textit{$\underline{S}$ is $\pi$-accessible in $\mathcal{M}$, and stable in $\mathcal{M}$ under $\varinjlim$ of type $\Phi$.}
\item[a$'$)] \textit{$\underline{S}$ is $\pi$-accessible in $\mathcal{M}$, and $\underline{S}(\pi)$ satisfies $L(\Phi)$.}
\end{itemize}
\msp{page 343}
\begin{itemize}[leftmargin=2.5em]
\item[b)] \textit{$\underline{S}(\pi)$ satisfies $L(\Phi)$, and for every $x\in\Ob\mathcal{M}$, we have this: $x\in\Ob\underline{S}$ if and only if the strictly full subcategory $\underline{S}(\pi)/x$ of $\mathcal{M}(\pi)/x$ is cofinal in $\mathcal{M}(\pi)/x$} \texttt{[read} $\mathcal{M}_\pi/x$\texttt{]}.
\item[\cercle{c}] \textit{$\underline{S}(\pi)$ satisfies $L(\Phi)$, and for every $x\in\Ob\mathcal{M}$ and every representation $x\simeq\varinjlim\limits_{i\in I}x_i$ which is $\Phi$-$\pi$-adapted (cf. definition below), we have this: $x\in\Ob\underline{S}$ if and only if $I_{\underline{S}}$ is cofinal in $I$, where $I_{\underline{S}}$ is the strictly full subcategory of $I$ formed by the $i\in\Ob I$ such that $x_i\in\Ob\underline{S}$ (or again, $x_i\in\Ob\underline{S}(\pi)$).} \textbf{N.B.} \textit{We say that a direct system $(x_i)_{i\in I}$ in $\mathcal{M}$ is $\Phi$-$\pi$-adapted if}
\begin{itemize}[leftmargin=2em]
\item[\textit{1)}] \textit{$I$ equivalent to a small category, and $I$ large relative to $\pi$,}
\item[\textit{2)}] \textit{the $\varinjlim$ of type $\Phi$ exist in $I$,}
\item[\textit{3)}] \textit{the functor $i\longmapsto x_i$ from $I$ to $\mathcal{M}$ commutes with $\varinjlim$ of type $\Phi$, and finally}
\item[\textit{4)}] \textit{the $x_i\in\Ob\mathcal{M}_\pi$.}
\end{itemize}
\item[c$'$)] \textit{(If $\pi^{\pi_0}=\pi$, where $\pi_0$ is a cardinal $\leqslant\pi$ useful for $\mathcal{M}$, and if $J\in\Phi$ implies that $J$ is large relative to $\pi_0$.) As in c), with $I$} ordered.
\end{itemize}
(\footnote{\textbf{N.B.} When the restrictive conditions opposite \texttt{[parenthesis in c$'$)]} are not satisfied, then a, a$'$, b, c are still equivalent, and imply c$'$, but are perhaps not equivalent to c$'$).},\footnote{\textbf{N.B.} One should restate the auxiliary conditions [?] (If \dots), cf. p. 359-360.}).

\textbf{Definition 4.13.23} (\footnote{\textbf{N.B.} It would be worthwhile to move this definition up (\textit{without} b), even before 4.13.15.}). Let $\mathcal{M}$ be an accessible category, $\pi$ a cardinal, $\Phi$ a set of categories, satisfying the two conditions:
\begin{itemize}[leftmargin=2.5em]
\item[a)] $J\in\Phi\Longrightarrow\cardred J\leqslant\pi$.
\item[b)] There exists a limit ordinal $\alpha_0$, such that $I_{\alpha_0}\in\Phi$ (\footnote{If it is to give a definition, one can dispense with condition b).}).
\end{itemize}
\msp{page 344}
Let $S$ be a subset of $\Ob\mathcal{M}$, $\underline{S}$ the full subcategory it defines. We say that $S$ is a \textit{$\Phi$-$\pi$-accessible subset of} $\Ob\mathcal{M}$, or that $\underline{S}$ is a \textit{$\Phi$-$\pi$-accessible subcategory in} $\mathcal{M}$, if the following conditions are satisfied.
\begin{itemize}[leftmargin=2.5em]
\item[1°)] $\pi$ is useful for $\mathcal{M}$, $\mathcal{M}$ is stable under $\varinjlim$ of type $\Phi$ (\footnote{Preliminary condition on $\mathcal{M}$.}).
\item[2°)] $\underline{S}$ is a $\pi$-accessible subcategory in $\mathcal{M}$.
\item[3°)] $\underline{S}$ is stable in $\mathcal{M}$ under $\varinjlim$ of type $\Phi$.
\end{itemize}
(\footnote{\textbf{N.B.} When $\Phi=\{I_{\alpha_0}\}$, we also say that $\underline{S}$ is $\alpha_0$-$\pi$-accessible in $\mathcal{M}$. If $\Phi$ consists of the categories of cardinality $\leqslant\pi$ and large relative to $\pi_0$, we say that $\underline{S}$ is $\pi_0$-$\pi$-accessible in $\mathcal{M}$.}).

\textbf{N.B.} When the preliminary condition 1°) \textit{on} $\mathcal{M}$ is satisfied, the validity of conditions 2°) 3°) is therefore equivalent to the condition that $\underline{S}$ be \textit{strictly} full in $\mathcal{M}$, and that in
\printedpage{131}
% head: A. Grothendieck : Dérivateurs, Ch. XVIII, Édition des Universités de Montpellier II et Paris VII
% foot: Version 17 décembre 2025
addition $\underline{S}$ satisfy condition a) of theorem 4.13.22. On the other hand, this condition is equivalent to each of conditions a$'$, b, c of this same theorem, and also to the last condition c$'$, provided that we have $\pi^{\pi_0}=\pi$ (for a cardinal $\pi_0<\pi$ useful for $\mathcal{M}$) and that $J\in\Phi$ implies that $J$ is large relative to $\pi_0$.

Conceptually, it is the definition given in 4.13.23, i.e. condition \cercle{a} of 4.13.22, which seems the clearest, or its immediate variant a$'$. But from the practical point of view, of working with accessible subsets, it is by property \cercle{c} which seems the most convenient. The main weight of theorem 4.13.22 is indeed
\msp{page 345}
the implication a $\Longrightarrow$ c (via a $\Longrightarrow$ a$'$ $\overset{!}{\Longrightarrow}$ c). So one must above all see this theorem as giving a remarkable criterion for membership in $\Ob\underline{S}$ of an object $x$ of $\mathcal{M}$ which is presented by a $\Phi$-$\pi$-adapted direct system:
\[ x\ \in\ \Ob\underline{S}\qquad\quad\Longleftrightarrow\qquad\quad I_{\underline{S}}\ \text{cofinal}\ , \tag*{$(*)$} \]
and in this equivalence, it is the implication $\Longrightarrow$ which is the non-sorital part, and remarkable.

\textbf{Proposition 4.13.24.} \textit{Let $\mathcal{M}$, $\pi$, $\Phi$ be as in 4.13.23} (\footnote{\textbf{N.B.} The condition $I_{\alpha_0}\in\Phi$ is not needed.})\textit{, we assume that condition 1° of 4.13.23 is satisfied ($\pi$ useful for $\mathcal{M}$, $\mathcal{M}$ stable under $\varinjlim$ of type $\Phi$). Let ${\Sigma'}^{\Phi}_{\pi}$ be the set of subcategories $\underline{S}$ of $\mathcal{M}$ which are $\Phi$-$\pi$-accessible in $\mathcal{M}$, and $\Sigma^{\Phi}_{\pi}$ the set of subcategories $\underline{S}(\pi)$ of $\mathcal{M}_\pi$ which are strictly full, and stable in $\mathcal{M}_\pi$ under $\varinjlim$ of type $\Phi$. (These subcategories are necessarily Karoubian, since stable under $\varinjlim$ of type $I_{\alpha_0}$, hence also under direct factors.) Then the bijections (4.13.4.1) (page 288) induce mutually inverse bijections}
\[
\begin{tikzcd}
\Sigma^{\Phi}_{\pi} \arrow[r, shift left] & {\Sigma'}^{\Phi}_{\pi} \arrow[l, shift left]
\end{tikzcd}\ .
\]

\msp{page 346}
% marge : trait manuscrit (reproduit en image) en regard de la Proposition 4.13.25
\textbf{Proposition 4.13.25.} \textit{Let $\mathcal{M}$, $\pi$, $\Phi$ be as in 4.13.23, with $I_{\alpha_\pi}\in\Phi$ ($\alpha_\pi$ being the smallest ordinal of cardinality $\pi$} \texttt{[rather} $\pi$\texttt{]}\textit{). Then the set ${\Sigma'}^{\Phi}_{\pi}$ of subcategories of $\mathcal{M}$ $\Phi$-$\pi$-accessible in $\mathcal{M}$ is stable under intersection of families of cardinality $\leqslant\pi$} (\footnote{If $\pi$ is not a limit cardinal. Otherwise, one must say: families of cardinality $<\pi$. Cf. 4.13.10.3, where I give a slightly more general corrected statement.}).
\DERmark{ENV1371}


(\textbf{N.B.} The bijection $\Sigma^{\Phi}_{\pi}\xrightarrow{\ \sim\ }{\Sigma'}^{\Phi}_{\pi}$ of 4.13.24 respects the natural order relations of inclusion, hence transforms intersections into intersections.)

Seen as in 4.13.7 (which is contained in the present statement), using lemma 4.13.10.

\textbf{4.13.26.} Let $\mathcal{M}$, $\pi$, $\Phi$ be as in 4.13.23, and let $\underline{S}$ be a \textit{strictly full} subcategory of $\mathcal{M}$, $f:\underline{S}\longrightarrow\mathcal{M}$ the inclusion functor. The condition that $\underline{S}$ be $\Phi$-$\pi$-accessible in $\mathcal{M}$ breaks up into:
\begin{itemize}[leftmargin=2.5em]
\item[1)] $\pi$ is useful for $\mathcal{M}$ \textit{and} for $\underline{S}$.
\item[2)] $\mathcal{M}$ and $\underline{S}$ are stable under $\varinjlim$ of type $\Phi$.
\item[3)] $f$ commutes with the said limits.
\item[4)] $f$ is $\pi$-accessible.
\item[5)] $f(\underline{S}_\pi)\subset\mathcal{M}_\pi$.
\end{itemize}

\printedpage{132}
% head: A. Grothendieck : Dérivateurs, Ch. XVIII, Édition des Universités de Montpellier II et Paris VII
% foot: Version 17 décembre 2025
\msp{page 347}
Let then $\pi'>\pi$ be another cardinal, when can we be sure that if $\underline{S}$ is $\Phi$-$\pi$-accessible in $\mathcal{M}$, it is also $\Phi$-$\pi'$-accessible\,? Conditions 2) 3) 4) are visibly stable under passage from $\pi$ to $\pi'$. Conditions 1) and 5) remain. We know that if $\pi'^{\pi}=\pi'$, then condition 1) remains valid for $\pi'$. More generally, if $\pi_0$ is a cardinal $\leqslant\pi$ such that $\pi_0$ is useful for $\mathcal{M}$ and for $\underline{S}$, then if $\pi'>\pi_0$ is such that $\pi'^{\pi_0}=\pi'$, then $\pi'$ is also useful for $\mathcal{M}$ and for $\underline{S}$. If we admit the GCH (generalized continuum hypothesis), this is so for every $\pi'>\pi_0$, hence for every $\pi'>\pi$.

Condition 5) remains, for which an answer is given by 4.12.1 (p. 273). If $\pi_0\leqslant\pi$ is chosen so that $f$ is $\pi_0$-accessible (e.g. $\pi_0=\pi$), and if
\[
\begin{aligned}
\pi'_1 &= \cardred\underline{S}_{\pi_0} \qquad\quad (\,\leqslant\cardred\underline{S}(\pi)\,)\\
\pi''_1 &= \sup(\pi'_1,\pi)\ ,
\end{aligned}
\]
we see that if $\pi'\geqslant\pi''_1$ (i.e. $\pi'\geqslant\pi'_1,\pi$) is such that $\pi'^{\pi_0}=\pi'$, then we do indeed have
\[ f(\underline{S}_{\pi'})\ \subset\ \mathcal{M}_{\pi'}\ . \]

\msp{page 348}
If then in addition $\pi'>\pi\geqslant\pi_0$ (as assumed at the start), and if we assume the GCH, then the condition $\pi'^{\pi_0}=\pi'$ is automatic, and only the condition $\pi'\geqslant\pi''_1$ remains, i.e. in the end $\pi'\geqslant\cardred\underline{S}_{\pi_0}$. To sum up:

\textbf{Proposition 4.13.27.} \textit{Under the conditions of 4.13.23 (without condition b) on $\Phi$), let $\underline{S}$ be a subcategory of $\mathcal{M}$ which is $\Phi$-$\pi$-accessible. Let $\pi_0\leqslant\pi$ be a cardinal which is useful for $\underline{S}$ and for $\mathcal{M}$, and such that $f:\underline{S}\longrightarrow\mathcal{M}$ is $\pi_0$-accessible, i.e. such that $\underline{S}$ is stable in $\mathcal{M}$ under $\varinjlim$ large relative to $\pi_0$ (one can e.g. take $\pi_0=\pi$). Then for every cardinal $\pi'>\pi$ such that}
\DERmark{ENV1372}

\begin{itemize}[leftmargin=2.5em]
\item[\textit{a)}] \textit{$\pi'^{\pi_0}=\pi'$ and}
\item[\textit{b)}] \textit{$\pi'\geqslant\cardred\underline{S}_{\pi_0}$,}
\end{itemize}
\textit{$\underline{S}$ is $\Phi$-$\pi'$-accessible in $\mathcal{M}$. Thus if $\pi'>\pi$ is such that $\pi'^{\pi}=\pi'$, $\pi'>\cardred\underline{S}_\pi$, then $\underline{S}$ is $\Phi$-$\pi'$-accessible. (It suffices e.g. to take $\pi'=2^c$, with $c\geqslant\sup(\pi,\cardred\underline{S}(\pi))$.) If we admit the GCH, it therefore suffices to have $\pi'>\sup(\pi,\cardred\underline{S}(\pi))$.}

\msp{page 349}
\textbf{Proposition 4.13.28} (\footnote{To be moved further down, after 4.13.31, which is the first important application of 4.13.22, whereas 4.13.28 is proved without having to use 4.13.22 or its variants.}). \textit{Let $\mathcal{M}$, $\pi$, $\Phi$ be as in 4.13.23. Suppose in addition that $J\in\Phi$ implies that $J$ is} filtering. \textit{Then the set $\tilde{\Sigma}^{\Phi}_{\pi}$ of $\Phi$-$\pi$-accessible subsets of $\mathcal{M}$ is stable under finite unions. More generally, let $\lambda$ be a cardinal $<\pi$ such that $\forall\ J\in\Phi$, $J$ is large relative to $\lambda$. Then the set $\tilde{\Sigma}^{\Phi}_{\pi}$ is stable under unions $(S_r)_{r\in A}$, indexed by sets $A$ of cardinality $\leqslant\lambda$.}

\textbf{Corollary 4.13.29.} \textit{The set of accessible subsets of} $\Ob\mathcal{M}$ \textit{is stable under ``small'' unions -- i.e. if $(S_r)_{r\in A}$ is a small family of accessible subsets of} $\Ob\mathcal{M}$, \textit{then $\bigcup_{r\in A}S_r$ is also accessible.}

The proof of 4.13.28 rests on the following lemma (I do not spell out the proof): % [sic] par la démonstration -> pas la démonstration, sense of the sentence
\DERmark{ENV1373}


\printedpage{133}
% head: A. Grothendieck : Dérivateurs, Ch. XVIII, Édition des Universités de Montpellier II et Paris VII
% foot: Version 17 décembre 2025
\msp{page 350}
\textbf{Lemma 4.13.30.} \textit{Let $I$ be a category large relative to $\pi$, and $(J_r)_{r\in A}$ a family of full subcategories such that $\card A\leqslant\pi$, such that $\bigcup_{r\in A}\Ob J_r$ is cofinal in $\Ord(I)$, i.e. for every $i\in\Ob I$ $\exists\ j\in\bigcup_{r\in A}\Ob J_r$, and an arrow $i\longrightarrow j$. Then there exists $r\in A$ such that $J_r$ is cofinal in $I$.}

\textbf{Theorem 4.13.31.} \textit{Let $\mathcal{M}$, $\mathcal{N}$ be accessible categories, $f:\mathcal{M}\longrightarrow\mathcal{N}$ an accessible functor, $T\subset\Ob\mathcal{N}$ a subset of $\mathcal{N}$, $\underline{T}\subset\mathcal{N}$ the corresponding strictly full subcategory of $\mathcal{N}$, $S=f^{-1}(T)$, $\underline{S}=f^{-1}(\underline{T})$ the associated full subcategory of $\mathcal{M}$. If $T$ is accessible, then $S$ is accessible\,; in other words, if $\underline{T}$ is accessible in $\mathcal{N}$, $\underline{S}$ is accessible in $\mathcal{M}$.}

This follows from the more precise statement:

\textbf{Corollary 4.13.32.} \textit{Let $\pi_0$ and $\pi>\pi_0$ be cardinals having the following properties:}

\msp{page 351}
\begin{itemize}[leftmargin=2.5em]
\item[\textit{a)}] \textit{$\underline{T}$ is $\pi_0$-$\pi$-accessible in $\mathcal{N}$ (4.13.23).}
\item[\textit{b)}] \textit{$\pi_0$ is useful for $\mathcal{M}$ (it is for $\mathcal{N}$, because of a)).}
\item[\textit{c)}] \textit{$f(\mathcal{M}_\pi)\subset\mathcal{N}_\pi$.}
\item[\textit{d)}] \textit{$\mathcal{M}$ satisfies $L_{\pi_0}$, and $f$ is $\pi_0$-accessible.}
\end{itemize}
\textit{(The existence of such a pair $(\pi_0,\pi)$ follows at once from 4.13.16 and from 4.12.1.) Then $\underline{S}$ is $\pi_0$-$\pi$-accessible in $\mathcal{M}$.}
\DERmark{ENV1374}


Since $f$ commutes with $\varinjlim$ large relative to $\pi_0$ by virtue of d), and $\underline{T}$ is stable under these limits by a) (cf. 4.13.23, 3°), the same is therefore true of $\underline{S}$, i.e. 3° of definition 4.13.23 is satisfied. On the other hand, condition 1° ($\pi$ useful for $\mathcal{M}$, $\mathcal{M}$ stable under $\varinjlim$ large relative to $\pi_0$ and of cardinality $\leqslant\pi$) is verified by b) and d). Thus it remains to verify condition 2°, i.e. that $\underline{S}$ is $\pi$-accessible. Let us put
\[
\begin{aligned}
\underline{T}(\pi) &= \underline{T}\cap\mathcal{N}_\pi\\
\underline{S}(\pi) &= \underline{S}\cap\mathcal{M}_\pi\ \ \underbrace{\ \subset\ }_{\text{thanks to c)}}\ \ f^{-1}(\underline{T}(\pi))\ .
\end{aligned}
\]
Since $\underline{S}$ is stable in $\mathcal{M}$ under $\varinjlim$ large relative to $\pi_0$ (which exist in $\mathcal{M}$ by d), and a fortiori under $\varinjlim$ large relative to $\pi$, we have a fully faithful canonical functor
\[ \Ind_\pi\underline{S}(\pi)\ \longrightarrow\ \underline{S}\ . \]

\msp{page 352}
Besides, $\underline{T}$ being strictly full in $\mathcal{N}$, $\underline{S}$ is strictly full in $\mathcal{M}$, and it remains to verify that the preceding functor is essentially surjective. Let then $x\in\Ob\underline{S}$. We will have, thanks to b),
\[ x\ =\ \varinjlim_{I}x_i\ ,\qquad\quad x_i\in\mathcal{M}_\pi,\ I=\mathcal{M}_\pi/x\ . \]
Besides, the preceding direct system is $\pi_0$-$\pi$-adapted, i.e. (p. 339):
\begin{itemize}[leftmargin=2.5em]
\item[1)] $I$ equivalent to a small category, $I$ large relative to $\pi$ (follows from the fact that $\pi$ is useful for $\mathcal{M}$, cf.\,\cercle{b}\,).
\end{itemize}


\endgroup
% END UNIT der-ch18/der-ch18_p131-135.tex
% BEGIN UNIT der-ch18/der-ch18_p136-140.tex
\begingroup\makeatletter%
\newcommand{\printedpage}[1]{\marginpar{\footnotesize p.~#1}}%
\newcommand{\msp}[1]{\marginpar{\footnotesize [#1]}}%
\newcommand{\Ob}{\operatorname{Ob}}%
\newcommand{\Fl}{\operatorname{Fl}}%
\newcommand{\card}{\operatorname{card}}%
\newcommand{\Ind}{\operatorname{Ind}}%
\newcommand{\Hom}{\operatorname{Hom}}%
\newcommand{\id}{\mathrm{id}}%
\newcommand{\limind}{\varinjlim}%
\newcommand{\cercle}[1]{\tikz[baseline=(c.base)]\node[draw,circle,inner sep=1pt](c){#1};}%
\newcommand{\defeq}{\overset{\mathrm{def}}{=}}%
\makeatother\let\printedpage\derpp\let\msp\dermsp%
\let\mspin\DERmspin%

\DERmark{ENV1375}


\printedpage{134}
% head: A. Grothendieck : Dérivateurs, Ch. XVIII, Édition des Universités de Montpellier II et Paris VII
% foot: Version 17 décembre 2025
\begin{itemize}[leftmargin=2.5em]
\item[2)] The direct limits large relative to $\pi_0$ and of cardinality $\leqslant \pi$ exist in $I$ (follows from $\mathrm{L}_{\pi_0}$ on $\mathcal{M}$ (cf.\,\cercle{d})). % [sic] est de cardinal -> et de cardinal, coordination ("grandes devant π0 et de cardinal ≤ π")
\item[3)] The functor $\xi : I \longrightarrow \mathcal{M}$ commutes with them.
\item[4)] The $x_i$ are in $\mathcal{M}_\pi$.
\end{itemize}
\DERmark{ENV1376}

Since $f$ commutes with the $\limind$ large relative to $\pi_0$ (cf.\,\cercle{d}), a fortiori with those large relative to $\pi$, we therefore have
\DERmark{ENV1377}

\[ f(x) \;\simeq\; \varinjlim_{I} f(x_i) \ , \]
I say that the direct system $(f(x_i))_{i\in I}$ in $\mathcal{N}$ is $\pi_0$-$\pi$-adapted. Conditions 1), 2) are \textit{ne varietur} (concern only $I$), the functor $I \longrightarrow \mathcal{N}$ is the composite $I \xrightarrow{\xi} \mathcal{M} \xrightarrow{f} \mathcal{N}$, hence commutes with the $\limind$ large relative to $\pi_0$ and of cardinality $\leqslant \pi$ (\footnote[155]{this commutation should have a name, condition $\mathrm{L}(\pi_0,\pi)$ on a functor.}), since this is so for $\xi$ and $f$, finally the $f(x_i)$ are in $\mathcal{N}_\pi$ since by \cercle{c} $f(\mathcal{M}_\pi) \subset \mathcal{N}_\pi$.

\msp{page 353}

Moreover $f(x) \in \Ob \underline{T}$ (since $x \in \Ob \underline{S}$, where $\underline{S} \defeq f^{-1}(\underline{T})$). Thus we can apply 4.13.22 a $\Longrightarrow$ c to the situation $(\mathcal{N}, \pi_0, \pi, \underline{T})$ and to the direct system $(f(x_i))_{i\in I}$, to conclude that the subcategory $I'$ of $I$ consisting of the $i$ such that $f(x_i) \in \Ob \underline{T}$ \texttt{[rather Ob}~$\underline{T}$\texttt{]}, i.e. such that $x_i \in \underline{S}$, hence $x_i \in \underline{S}(\pi)$, is \textit{cofinal} in $I$. Thus we have
\DERmark{ENV1378}\DERmark{ENV1379}

\[ x \;\simeq\; \varinjlim_{i\in I'} x_i \ , \]
with this time the $x_i$ in $\underline{S}(\pi)$, and $I'$ large relative to $\pi$ (since cofinal and full in $I$, which is large relative to $\pi$), which proves that $x$ is in the essential image of $\Ind_\pi(\underline{S}(\pi))$, q.e.d.

\textbf{Question 4.13.33.} We have just seen that if $f : \mathcal{M} \longrightarrow \mathcal{N}$ is an accessible functor between accessible categories, then the inverse image of an accessible subset of $\Ob \mathcal{N}$ is an accessible subset of $\Ob \mathcal{M}$. I conjecture that the similar statement, not for direct images, but for \textit{essential images}, is
\msp{page 354}
true, i.e. that for every subcategory $\underline{S}$ of $\mathcal{M}$ accessible in $\mathcal{M}$, its essential image in $\mathcal{N}$ is accessible in $\mathcal{N}$ (\footnote[156]{False, even if we take the $\natural$-closure of the essential image (``$\natural$-essential image'').}). In this question, we can assume that $\underline{S} = \mathcal{M}$, and the question comes down to this: does every accessible functor $f : \mathcal{M} \longrightarrow \mathcal{N}$ between accessible categories factor into an (accessible) \textit{essentially surjective} functor, and an accessible fully faithful functor $\mathcal{M} \longrightarrow \underline{T} \longrightarrow \mathcal{N}$ (where $\underline{T}$ is itself accessible) (\footnote[157]{We must at least restrict ourselves to essentially $\natural$-surjective – but even that is too optimistic.})~? If so, this would seem to me perhaps the deepest result (perhaps the only truly deep one) in category theory. % [sic] vraie -> vrai, agreement with « l'énoncé similaire »

Here is a question that seems to me very close. Let $\mathcal{M}$ be an accessible category, $A$ a small full subcategory. Suppose that $\mathcal{M}$ satisfies $\mathrm{L}_{\pi_0}$ (e.g. $\pi_0$ useful for $\mathcal{M}$), and let $\underline{S}$ be the smallest strictly full subcategory of $\mathcal{M}$ containing $A$, and stable under limits large relative to $\pi_0$. I conjecture that $\underline{S}$ is accessible in $\mathcal{M}$. % [sic] sous-catégories -> sous-catégorie, singular after « une »
\DERmark{ENV1380}


\msp{page 355}

Likewise, if $\Phi$ is a small set of small categories, such that $\mathcal{M}$ is stable under $\limind$ of type $\Phi$, we can take the envelope $\underline{S}$ of $A$ for the $\limind$ large relative to $\pi_0$ \textit{and} the
\printedpage{135}
% head: A. Grothendieck : Dérivateurs, Ch. XVIII, Édition des Universités de Montpellier II et Paris VII
% foot: Version 17 décembre 2025
$\limind$ of type $\Phi$, and I again conjecture that $\underline{S}$ is accessible. In this direction, we have at least the following very partial result.

\textbf{Proposition 4.13.34.} \textit{Let $\mathcal{M}$ be an accessible category stable under small filtering $\limind$, and let $\underline{S}$ be a strictly full subcategory of $\mathcal{M}$. Equivalent conditions:}
\begin{itemize}[leftmargin=2.5em]
\item[\textit{a)}] \textit{$\underline{S}$ is accessible in $\mathcal{M}$, and stable in $\mathcal{M}$ under small filtering $\limind$.}
\item[\textit{b)}] \textit{There exists a small full subcategory $A$ of $\mathcal{M}$, such that $\underline{S}$ is the smallest strictly full subcategory of $\mathcal{M}$ containing $A$ and stable under small filtering $\limind$.}
\end{itemize}

\msp{page 356}

We have a) $\Longrightarrow$ b), for if $\underline{S}$ is $\pi$-accessible in $\mathcal{M}$, it is immediate (without any other stability condition on $\underline{S}$) that $\underline{S}$ is the smallest strictly full subcategory of $\mathcal{M}$ containing $\underline{S}(\pi) \defeq \underline{S}\cap\mathcal{M}_\pi$ and stable under $\limind$ large relative to $\pi$. If then we already know that $\underline{S}$ is even stable under \textit{all} small filtering $\limind$, it follows that $\underline{S}$ is the smallest subcategory of $\mathcal{M}$ having this stronger stability property, and containing $A$.

Let us prove b $\Longrightarrow$ a. Let $\pi$ be a useful cardinal for $\mathcal{M}$, and such that $A \subset \mathcal{M}_\pi$. Let $B$ be the smallest strictly full subcategory of $\mathcal{M}$ containing $A$ and stable under filtering $\limind$ of cardinality $\leqslant \pi$. Thus we have $B \hookrightarrow \mathcal{M}_\pi$ and $B \subset \underline{S}$, and $\underline{S}$ is also the L-envelope of $B$. Besides $B$ is stable under
\msp{page 357}
direct factors. Let us then consider the essential image $\underline{S}'$ of $\Ind_\pi(B)$ in $\mathcal{M}$. We have $\underline{S}' \subset \underline{S}$, and everything comes down to proving that $\underline{S}' = \underline{S}$, from which it will indeed follow that $\underline{S}$ is $\pi$-accessible in $\mathcal{M}$, a fortiori accessible in $\mathcal{M}$, whence a). To prove that $\underline{S}' = \underline{S}$, by definition of $\underline{S}$ as the L-envelope of $A$, it suffices to see that $\underline{S}'$ is stable under small filtering $\limind$. Now this follows from the fact that $B \subset \mathcal{M}_\pi$ is stable in $\mathcal{M}_\pi$ under filtering $\limind$ of cardinality $\leqslant \pi$, and from 4.13.19, which implies at least, since $B = \underline{S}' \cap \mathcal{M}_\pi$ (by virtue of 4.13.4) that $\underline{S}'$ is stable under filtering $\limind$ \textit{of cardinality} $\leqslant \pi$. That it is even stable under arbitrary small filtering $\limind$ follows from this, and from the fact that it is so under $\limind$ large relative to $\pi$. (Cf. 4.13.21 \dots) % [sic] filtrante quelconques -> filtrantes quelconques, plural agreement

\msp{page 358}

\textbf{Definition 4.13.35.} Let $\mathcal{M}$ be an accessible category, stable under small filtering $\limind$, $S$ a subset of $\Ob \mathcal{M}$, $\underline{S}$ the corresponding full subcategory in $\mathcal{M}$. We say that $S$ is \textit{L-accessible}, or that $\underline{S}$ is L-accessible in $\mathcal{M}$, if $\underline{S}$ is accessible in $\mathcal{M}$, and in addition stable under small filtering $\limind$ (\footnote[158]{\textbf{N.B.} The set of L-accessible subsets of $\Ob \mathcal{M}$ is stable under small intersections, and under \textit{finite} unions. The inverse image of an L-accessible subset under a functor (\footnotemark[159]) $f : \mathcal{M} \longrightarrow \mathcal{N}$ of accessible categories is L-accessible.}).\footnotetext[159]{an \textit{$L_0$-functor}: which commutes with small filtering $\limind$.}

\textbf{Proposition 4.13.36.} \textit{Let $\mathcal{M}$, $S$, $\underline{S}$ be as in 4.13.35, and let $\pi$ be a useful cardinal for $\mathcal{M}$. Suppose $\underline{S}$ stable in $\mathcal{M}$ under small filtering $\limind$, or only $\underline{S}(\pi) \defeq \underline{S} \cap \mathcal{M}_\pi$ stable under small filtering $\limind$ of cardinality $\leqslant \pi$. Equivalent conditions:}
\begin{itemize}[leftmargin=2.5em]
\item[\textit{a)}] \textit{$\underline{S}$ is $\pi$-accessible. (It will automatically be stable under small filtering $\limind$, hence L-accessible.)}
\item[\textit{b)}] \textit{For every $x \in \Ob \mathcal{M}$, we have $x \in \Ob \underline{S}$ if and only if $\underline{S}(\pi)/x$ is cofinal in $\mathcal{M}_\pi/x$.}
\printedpage{136}
% head: A. Grothendieck : Dérivateurs, Ch. XVIII, Édition des Universités de Montpellier II et Paris VII
% foot: Version 17 décembre 2025
\item[\textit{c)}] \textit{For every $x \in \Ob \mathcal{M}$, and every representation $x \simeq \varinjlim\limits_{i\in I} x_i$}
\msp{page 359}
\textit{which is L-$\pi$-adapted (cf. page 343, c), we have this: $x \in \Ob S$ if and only if $I_{\underline{S}}$ is cofinal in $I$.}
\item[\textit{c}$'$\textit{)}] \textit{As in c), with $I$ ordered.}
\item[\textit{c}$''$\textit{)}] \textit{(If we have chosen, for every $x \in \Ob \mathcal{M}$,} one \textit{L-$\pi$-adapted representation $x \simeq \varinjlim\limits_{I} x_i$): $x \in \Ob \underline{S} \Longleftrightarrow$ the set $I_{\underline{S}}$ is cofinal in $I$.}
\end{itemize}
\DERmark{ENV1381}

\textsc{Proof}. This is essentially the statement of 4.13.22 in the special case where $\Phi =$ set of \textit{filtering} categories of cardinality $\leqslant \pi$, but where, to simplify, we have assumed beforehand that $\underline{S}(\pi)$ satisfies $\mathrm{L}(\Phi)$. The equivalence of a) b) c) is contained in 4.13.22, as is c $\Longrightarrow$ c$'$. But for c$'$ $\Longrightarrow$ a, the auxiliary hypothesis $\pi^{\pi_0} = \pi$ and that the $J \in \Phi$ will be large relative to $\pi_0$, is not satisfied here, and we need an independent proof, modeled on that of pages 333-336. For this it suffices, at the place \texttt{[of]} $\square$, page 333, to invoke the % [sic] un démonstration -> une démonstration, gender agreement

\textbf{Lemma 4.13.37.} \textit{Let $\mathcal{M}$ be a category}
\msp{page 360}
\textit{that is accessible and stable under filtering $\limind$, $\pi$ a useful cardinal for $\mathcal{M}$. Then for every $x \in \Ob \mathcal{M}$, we have a representation}
\[ x \;=\; \varinjlim_{i\in I} x_i \]
\textit{L-$\pi$-adapted (i.e. with}
\begin{itemize}[leftmargin=2.5em]
\item[\textit{1)}] \textit{$I$ equivalent to a small category, and $I$ large relative to $\pi$,}
\item[\textit{2)}] \textit{$I$ stable under filtering $\limind$ of cardinality $\leqslant \pi$,}
\item[\textit{3)}] \textit{$I \longrightarrow \mathcal{M}$ commutes with these $\limind$ ,}
\item[\textit{4)}] \textit{$x_i \in \Ob \mathcal{M}_\pi \ \forall\, i \in I$),} % [sic] d) -> 4), the items are numbered 1), 2), 3)
\end{itemize}
\DERmark{ENV1382}

\textit{with in addition $I$} ordered.

Indeed, we start from
\[ x \;=\; \varinjlim_{j\in J} x_j \ , \qquad\quad J \;\simeq\; \underbrace{(\mathcal{M}_\pi/x)_0}_{\substack{\text{reduced subcategory}\\ \text{equivalent to } \mathcal{M}_\pi/x}} \qquad . \]
We define $I$ as the set of \textit{filtering} subcategories $J'$ of $J$, of cardinality \texttt{[``cardinality'' crossed out]} $\leqslant \pi$, $I$ ordered by inclusion, and we set, for $J' \in I$,
\[ x_{J'} \;=\; \varinjlim_{j\in J'} x_j \ . \]
It is immediate that we thus obtain a direct system in $\mathcal{M}$, indexed by $I$.

\msp{page 361}

On the other hand, it is clear that $I$ is small, I say that $I$ is stable under filtering $\limind$ of cardinality $\leqslant \pi$, and that the functor $J' \longmapsto x_{J'} : I \longrightarrow \mathcal{M}$ commutes with them, finally that the $x_{J'}$ are in $\mathcal{M}_\pi$. It remains to see that $I$ is large relative to $\pi$, i.e. that if we have $(J'_r)_{r\in A}$, with $\card A \leqslant \pi$, the $J'_r$ in $I$, then $\exists\, J' \in I$ which dominates the $J'_r$. Let $K$ be the subcategory of $J$ generated by the $J'_r$, we easily see that $\card K \leqslant \pi$. Everything \texttt{[``everything'' crossed out]} comes down to proving:

\printedpage{137}
% head: A. Grothendieck : Dérivateurs, Ch. XVIII, Édition des Universités de Montpellier II et Paris VII
% foot: Version 17 décembre 2025
\textbf{Lemma 4.13.38} (\footnote[160]{\textbf{N.B.} For the proof of 4.13.37, it suffices to prove 4.13.38 in the case where $J$ is ordered, in which case the contortions at the end are unnecessary.}). \textit{Let $J$ be a category large relative to $\pi$, $K$ a subcategory of $J$ of cardinality $\leqslant \pi$. Then there exists a subcategory $K'$ that is} filtering \textit{in $J$, containing $K$, of cardinality $\leqslant \pi$. If we assume that we are not, in the case where $J$ contains a maximal object, and that $K$ contains all the maximal objects of $J$,}
\msp{page 362}
\textit{we can even take $K'$ such that $\Ob K' = \Ob K \cup \{e\}$, where $e$ is a} final object \textit{of $K'$, and $K$ is a full subcategory of $K'$. (Thus if $K' \neq K$, $K'$ is isomorphic to the right cone constructed on $K$, ``obtained from $K$ by adding a final object to it''.)}

\textsc{Proof}. If $K$ has a final object, we take $K' = K$. Otherwise, since $J$ is large relative to $\card \Ob K$, there exists a $j \in \Ob J$ which dominates all the objects of $K$. Let us choose, for every $k \in \Ob K$, an arrow
\[ \alpha_k \;:\; k \longrightarrow j \ . \]
For every arrow
\[ u \;:\; k' \longrightarrow k \]
in $K$, let us consider the double arrow
\[ (\alpha_k \circ u, \alpha_{k'}) \;:\; k' \rightrightarrows j \ . \]
The set of these double arrows, with target $j$, is of cardinality $\leqslant \card \Fl K \leqslant \pi$, so $J$ being large relative to $\pi$, there exists an arrow $\rho : j \longrightarrow e$ in $J$ which equalizes all these double arrows. Let $K' \subset J$, having as objects $\Ob K \cup \{e\}$, and having as arrows, in addition to
\msp{page 363}
those of $K$, the arrows $\rho\alpha_k : k \longrightarrow e$ (for $k \in \Ob K$), and $\id_e$. \textit{If} $e \notin \Ob K$, we see that this is indeed a subcategory, admitting $e$ as final object and $K$ as full subcategory. If we had $e \in \Ob K$, and if there exists an object $e'$ of $J$ which dominates $e$ and which is not in $K$, we immediately reduce to the preceding case, replacing $e$ by $e'$.

$\Biggl($ (\footnote[161]{The construction which follows is contained in the more general and more precise statement 4.13.39.}). Thus there remains the case where there exists an object $e$ in $K$ which is a greatest element in $K$, and such that every object of $J$ dominating $e$ is in $K$. This implies, $J$ being filtering, that $e$ is a greatest object of $J$. (For if $x \in J$, $\exists\, x'$ a common upper bound for $x$ and $e$, hence $x'$ is in $K$, hence dominated by $e$, which therefore dominates $x$.) We see as before that there exists an arrow $p : e \longrightarrow j$ in $J$, which equalizes all the double arrows $k \rightrightarrows e$ with target $e$ \textit{in}~$K$. Composing with $j \longrightarrow e$ if necessary, we can assume $j = e$. Thus we have an endomorphism $p$ of $e$ in $J$, such that $pu = pv$ for every double arrow $k \overset{u}{\underset{v}{\rightrightarrows}} e$ in $K$, with target $e$.

\msp{page 364}

For $k \in \Ob K$, let
\[ f_k \;:\; k \longrightarrow e \]
be the common value of the $pu$, for $u \in \Hom_K(k,e)$. We see at once that if we have $k' \xrightarrow{v} k$ in $K$, then
\[ f_{k'} \;=\; f_k \circ v \ . \]
\printedpage{138}
% head: A. Grothendieck : Dérivateurs, Ch. XVIII, Édition des Universités de Montpellier II et Paris VII
% foot: Version 17 décembre 2025
In particular, taking $k' = e$, we find
\[ f_k \circ v \;=\; f_e \,(=\ p\ ) \qquad\quad \text{if } v : e \longrightarrow k \ . \]
Let then $K'$ be the subcategory of $J$ having the same objects as $K$, and as arrows, in addition to those of $K$, the composite arrows $k \xrightarrow{u} e \xrightarrow{p^r} e \xrightarrow{v} k'$, with $u$, $v$ in $K$, and $r \geqslant 1$, i.e. the composite arrows $k \xrightarrow{f_k} e \xrightarrow{p^{r'}} e \xrightarrow{v} k'$, with $k, k' \in \Ob K$, $v \in \Fl K$, $r' \geqslant 0$. We must see that it is stable under composition (to indeed have a subcategory). This is clear for $vu$ if $u$ or $v$ are in $\Fl K$. If \texttt{[neither]} $u$ nor $v$ is, we must take a composite
\[
\begin{tikzcd}[column sep=large]
k \arrow[r, "u_1"] & e \arrow[r, "p^r", "r\geqslant 1"'] & e \arrow[r, "u_2"] \arrow[rr, bend right=20] & k' \arrow[r, "v_1"] & e \arrow[r, "p^s", "s\geqslant 1"'] & e \arrow[r, "v_2"] & k'' \ ,
\end{tikzcd}
\]
with $u_1$, $u_2$, $v_1$, $v_2$ in $\Fl K$, $r, s \geqslant 1$. Now we have
\[ p(v_1u_2) \;=\; p \ , \qquad\quad \text{whence} \qquad\quad p^s(v_1u_2) \;=\; p^{s-1}(pv_1u_2) \;=\; p^{s-1}p \;=\; p^s \ , \]
so the composite is
\[ k \xrightarrow{\;u_1\;} e \xrightarrow{\;p^{r+s}\;} e \xrightarrow{\;v_2\;} k'' \ , \]
which
\msp{page 365}
is therefore indeed in $\Fl K'$.

Thus, $K'$ is a subcategory of $J$, but is it filtering? Condition PS 1$'$ is obvious, since $e$ is a greatest element of $J$, there remains condition PS 2 of the double arrow. We easily see that this condition is satisfied in $K'$, if and only if it is satisfied for every double arrow
\[ e \overset{p^r}{\underset{p^s}{\rightrightarrows}} e \qquad\quad \text{with } r, s \geqslant 0 \ . \]
(Indeed, we immediately reduce to the double arrows with target $e$, and note that every arrow of $K'$ with target $e$ can be written as a composite $k \xrightarrow{u} e \xrightarrow{p^r} e \xrightarrow{v} e$, with $r \geqslant 0$ and $u$, $v$ arrows in $K$. Composing with $e \xrightarrow{p} e$, we then find, since $pv = p$, $k \xrightarrow{u} e \xrightarrow{p^{r+1}} e$, that is $k \xrightarrow{f_k} e \xrightarrow{p^r} e$.) But we see at once that the possibility of equalizing $(p^r, p^s)$ in $K'$ is equivalent to the existence of $t \geqslant 0$ such that $p^{r+t} = p^{s+t}$. Thus $K'$ is filtering if and only if $r \longmapsto p^r$ is a stationary sequence (take $r = 0$, $s = 1$, whence $p^r = p^{r+1}$ for $r$ large). This is
\DERmark{ENV1383}\DERmark{ENV1384}

\msp{page 366}
not necessarily the case.

Let us therefore set $K = K_0$, $K' = K_1$, and construct by induction an increasing sequence of subcategories $K_i$ ($i \in \mathbf{N}$), with $\Ob K_i = \Ob K$, and (for $i \geqslant 1$) $K_i$ being obtained from $K_{i-1}$ by adding to it an arrow $p_i : e \longrightarrow e$, as above. Thus every double arrow in a $K_i$ can be equalized in $K_{i+1}$. Let us take $K_\infty = \limind K_i$, $K_\infty$ is a \textit{filtering} subcategory of $J$, containing $K$ and of cardinality $\leqslant \pi$, q.e.d. $\Biggr)$
\DERmark{ENV1385}


We can make this construction more precise as follows:


\endgroup
% END UNIT der-ch18/der-ch18_p136-140.tex
% BEGIN UNIT der-ch18/der-ch18_p141-145.tex
\begingroup\makeatletter%
\newcommand{\printedpage}[1]{\marginpar{\footnotesize p.~#1}}%
\newcommand{\msp}[1]{\marginpar{\footnotesize [#1]}}%
\newcommand{\uHom}{\underline{\mathrm{Hom}}}%
\newcommand{\Fl}{\operatorname{Fl}}%
\newcommand{\Ob}{\operatorname{Ob}}%
\newcommand{\End}{\operatorname{End}}%
\newcommand{\Ind}{\operatorname{Ind}}%
\newcommand{\id}{\mathrm{id}}%
\newcommand{\cardred}{\mathop{\textup{card\,red}}}%
\newcommand{\card}{\operatorname{card}}%
\newcommand{\edbr}[1]{\textup{\texttt{#1}}}%
\makeatother\let\printedpage\derpp\let\msp\dermsp%
\let\mspin\DERmspin%


\printedpage{139}
% head: A. Grothendieck : Dérivateurs, Ch. XVIII, Édition des Universités de Montpellier II et Paris VII
% foot: Version 17 décembre 2025

\noindent\textbf{Corollary 4.13.39.} {\itshape Under the conditions of 4.13.38, suppose that $K$ contains an object $e$ maximal in $J$. Let $\Sigma$ be the set of endomorphisms $p$ of $e$ in $J$, such that for every double arrow $k \overset{u}{\underset{v}{\rightrightarrows}} e$ \textup{in} $K$, we have $pu=pv$. Then $\Sigma$ is stable under composition, and there exists a nonempty subset $\Sigma_\infty$ \edbr{[}$=\Sigma_0$\edbr{]} of $\Sigma$, such that for $p,p'\in\Sigma_0$, exists
\msp{page 367}
$q\in\Sigma_0$ such that $qp=qp'$ (i.e. such that the subcategory $J_0$ of $J$ whose set of objects is $\{e\}$, and set of arrows $\{\id_e\}\cup\Sigma_\infty$, is filtering). If $\Sigma_0$ is chosen thus, let $\mathcal{F}$ be the union of $\Fl K$ and of the set of composite arrows
\[
k \xrightarrow{\ u\ } e \xrightarrow{\ p\ } e \xrightarrow{\ v\ } k'\ , \tag*{$(*)$}
\]
with $u,v\in\Fl K$, $p\in\Sigma_\infty$. Then $\mathcal{F}$ is the set of arrows of a subcategory $K'$ of $J$, having the same objects as $K$, and admitting $e$ as greatest element. Moreover, $K'$ contains $K$ and is filtering.}
\DERmark{ENV1386}\DERmark{ENV1387}


\noindent\textsc{Proof}. That $\Sigma$ is stable under composition is clear. To construct $\Sigma_\infty$ having the stated properties, we note that, $J$ being large relative to $\pi$, $\Sigma$ is nonempty (\footnote[162]{cf. argument below for construction of $p_i$.}). We construct by induction a sequence of elements
\[
p_0\ ,\ p_1\ ,\ \ldots\ ,\ p_n\ ,\ \ldots
\]
of $\Sigma$, taking for $p_0$ any element of $\Sigma$, and for $i\geqslant 1$, taking $p_i\in\Sigma$ such that for every double arrow $(p,p')$ of the submonoid $M_{i-1}$ of $\End_J(e)$ generated by $(p_j)_{0\leqslant j\leqslant i-1}$, we have $p_ip=p_ip'$. Such a $p_i$ exists, for taking the set of double arrows of the preceding form $(p,p')$, with $p,p'\in M_{i-1}$, and of those of the form $(u,v) : k \rightrightarrows e$, with $u,v\in\Fl K$, the set of
\DERmark{ENV1388}

\msp{page 368}
all these double arrows has cardinality $\leqslant\pi$ ($M_{i-1}$ being countable, hence of cardinality $\leqslant\pi$), hence there exists an arrow $e\longrightarrow j$ in $J$ which equalizes them all ($J$ being large relative to $\pi$). Composing if necessary with $j\longrightarrow e$ ($e$ being a maximal object of $J$), we find the desired $p_i : e\longrightarrow e$.

Let then $\Sigma_0$ be the multiplicative subset of $\Sigma$ generated by the $p_i$ ($i\in\mathbf{N}$), direct limit of the similar sets $\Sigma_i$ generated by the $p_j$ for $j\in[0,i]$. If $p$, $p'$ are in $\Sigma_\infty$, they are in one and the same $\Sigma_i$, hence we will have $p_{i+1}p=p_{i+1}p'$, which proves that $\Sigma_\infty$ has the required filtering property. % [sic] il sont -> ils sont, agreement of the verb with the plural subject

Let us now prove that $\mathcal{F}$ is the set of arrows of a subcategory $K'$ of $J$, having $\Ob K$ as set of objects. It remains to establish stability under composition $vu$. But if $u$ and $v$ are in $\Fl K$, or if one of the two is in $\Fl K$, the other not (hence of the form $(*)$), this stability is evident. Remains to see the case where neither \edbr{[}$u$\edbr{]} nor $v$ is in $\Fl K$, hence both are of the form $(*)$:
\[
k \xrightarrow{\ u\ } e \xrightarrow{\ p\ } e \xrightarrow{\ v\ } k' \xrightarrow{\ u'\ } e \xrightarrow{\ p'\ } e \xrightarrow{\ v'\ } k''\ ,
\]
with $u,v,u',v'\in\Fl K$, $p,p'\in\Sigma_\infty$. Since $(u'v,\id_e)$ is a double arrow of
\printedpage{140}
% head: A. Grothendieck : Dérivateurs, Ch. XVIII, Édition des Universités de Montpellier II et Paris VII
% foot: Version 17 décembre 2025
\msp{page 369}
$K$ with target $e$, we have
\[
p'u'v \;=\; p'\id_e \;=\; p'\ ,
\]
hence the composite in question is also written
\[
k \xrightarrow{\ u\ } e \xrightarrow{\ p'p\ } e \xrightarrow{\ v'\ } k''\ ,
\]
where $p'p\in\Sigma_\infty$ since $\Sigma_\infty$ is stable under composition, hence it is in $\mathcal{F}$, q.e.d.

It is clear that $K'$ contains $K$, and having the same objects as $K$ which admits $e$ as greatest element, that $e$ is a greatest element of $K'$. Remains to prove that $K'$ is filtering. Since $K'$ has a greatest element, it remains to see that every double arrow $k\rightrightarrows l$ in $K'$ can be equalized in $K'$. Composing if necessary with $l\longrightarrow e$, we can suppose $l=e$. Now we see that every arrow $k\longrightarrow e$ in $K'$ is of the form
\[
k \xrightarrow{\ u\ } e \xrightarrow{\ p\ } e \xrightarrow{\ v\ } e\ ,
\]
with $u,v\in\Fl K$ and $p\in\Sigma_\infty\cup\{\id_e\}$. Composing if necessary with $p_0\in\Sigma_\infty$ fixed, this arrow becomes (since $p_0v=p_0\id_e=p_0$)
\[
k \xrightarrow{\ u\ } e \xrightarrow{\ p_0p\ } e\ , \qquad\quad \text{i.e. } p_0pu\ .
\]
For an arrow $v'p'u' : k\longrightarrow e$, we find $p_0p'u'$. Let then $q\in\Sigma_0$ be such that $q(p_0p)=q(p_0p')$, say $r$ ($q$ exists since $p_0p,p_0p'\in\Sigma_\infty$). Thus $(q(p_0pu),q(p_0p'u'))=(ru,ru')$, and we have $ru=ru'$ since $r\in\Sigma_\infty$. All told, $qp_0$ equalizes the two arrows $vpu$ and $v'p'u'$, q.e.d.

\msp{page 370}

\section*{4.14\quad Accessibility criteria for big categories}

\noindent\textbf{Theorem 4.14.1.} {\itshape Let $\mathcal{C}$ be a category equivalent to a small category, $\Phi$, $\Psi$ two small sets of categories, equivalent to small categories. For every category $I$, we denote by $I_{\bullet}$ (resp. ${}_{\bullet}I$) the right cone (resp. left) on $I$, obtained from $I$ by adding a final (resp. initial) object $e=e_I$. Let $\Lambda_1$ (resp. $\Lambda_2$) be a set of pairs $(I,\varphi)$ (resp. $(I,\psi)$), where $I\in\Phi$ (resp. $I\in\Psi$), and where
\[
\varphi\ :\ I _{\bullet}\longrightarrow \mathcal{C} \qquad\quad \text{(resp. }\psi\ :\ {}_{\bullet}I \longrightarrow \mathcal{C}\ \ \text{)} \tag*{\textup{(4.14.1.1)}}
\]
is a functor. For every category $\mathcal{N}$ stable (\footnote[163]{This stability property is not necessary for the definition of $\uHom^{\Lambda_2}_{\Lambda_1}(\mathcal{C},\mathcal{N})$, and it is almost certain that the theorem is valid without making this assumption on $\mathcal{N}$. Cf. 4.14.16, where we see that it is in any case unnecessary for $\Psi$.}), under $\varinjlim$ of type $\Phi$ and under $\varprojlim$ of type $\Psi$ we denote by
\[
\uHom^{\Lambda_2}_{\Lambda_1}(\mathcal{C},\mathcal{N})\ \subset\ \uHom(\mathcal{C},\mathcal{N}) \tag*{\textup{(4.14.1.2)}}
\]
the strictly full subcategory of $\uHom(\mathcal{C},\mathcal{N})$ formed of the functors $f:\mathcal{C}\longrightarrow\mathcal{N}$ having the following property:
\begin{enumerate}[leftmargin=2.5em]
\item[\textit{a)}] $\forall\ (I,\varphi)\in\Lambda_1$, the composite functor
\[
I_{\bullet} \xrightarrow{\ \varphi\ } \mathcal{C} \xrightarrow{\ f\ } \mathcal{N}
\]
is $I$-left exact, i.e. makes $(f\varphi)(e_I)$ a $\varinjlim\limits_{I}$ of $f\varphi|I$.
\printedpage{141}
% head: A. Grothendieck : Dérivateurs, Ch. XVIII, Édition des Universités de Montpellier II et Paris VII
% foot: Version 17 décembre 2025
\msp{page 371}
\item[\textit{b)}] For every $(I,\psi)\in\Lambda_2$, the composite functor
\[
{}_{\bullet}I \xrightarrow{\ \psi\ } \mathcal{C} \xrightarrow{\ f\ } \mathcal{N}
\]
is $I$-right exact, i.e. makes $(f\psi)(e_I)$ an inverse limit of $(f\psi)|I$.
\end{enumerate}
\noindent With these notations: if $\mathcal{N}$ is an accessible category, the same is true of $\uHom^{\Lambda_2}_{\Lambda_1}(\mathcal{C},\mathcal{N})$.}
\DERmark{ENV1389}\DERmark{ENV1390}


\noindent\textbf{Comment.} We will use this theorem chiefly when we have $\Lambda_2=\emptyset$ or $\Lambda_1=\emptyset$, in which case we write $\uHom_{\Lambda_1}(\mathcal{C},\mathcal{N})$, resp. $\uHom^{\Lambda_2}(\mathcal{C},\mathcal{N})$, for the category $\uHom^{\Lambda_2}_{\Lambda_1}(\mathcal{C},\mathcal{N})$. We will apply the theorem chiefly when $\mathcal{C}$ itself is stable under $\varinjlim$ of type $\Phi$ and under $\varprojlim$ of type $\Psi$, and when $\Lambda_1$ (resp. $\Lambda_2$) is already formed of functors $\varphi$ (resp. $\psi$) which are $I$-left exact (resp. right exact). Thus $\uHom^{\Lambda_2}_{\Lambda_1}(\mathcal{C},\mathcal{N})$ is then formed of the $f:\mathcal{C}\longrightarrow\mathcal{N}$ which commute with certain $\varinjlim$ of type $\Phi$, and with certain $\varprojlim$ of type $\Psi$. In particular, we find:

\msp{page 372}

\noindent\textbf{Corollary 4.14.2.} {\itshape Let $\mathcal{C}$ be a category equivalent to a small category, $\mathcal{N}$ a $\mathfrak{U}$-category, $\Phi$ and $\Psi$ two small sets of categories $I$, equivalent to small categories. We suppose $\mathcal{C}$ and $\mathcal{N}$ stable under $\varinjlim$ of type $\Phi$, under $\varprojlim$ of type $\Psi$. Let $\uHom^{\Psi}_{\Phi}(\mathcal{C},\mathcal{N})$ be the strictly full subcategory of $\uHom(\mathcal{C},\mathcal{N})$ formed of the $f:\mathcal{C}\longrightarrow\mathcal{N}$ which commute with $\varinjlim$ of type $\Phi$ and with $\varprojlim$ of type $\Psi$. Then, if $\mathcal{N}$ is accessible, the same is true of $\uHom^{\Psi}_{\Phi}(\mathcal{C},\mathcal{N})$.}

Here are some remarkable special cases of this corollary:

\noindent\textbf{Corollary 4.14.3.} {\itshape Let $\mathcal{M}$, $\mathcal{N}$ be two accessible categories, and $\pi_0$ an infinite cardinal, such that $\mathcal{M}$ satisfies $L_{\pi_0}$. Let $\uHom_{\pi_0}(\mathcal{M},\mathcal{N})$ be the strictly full subcategory of $\uHom(\mathcal{M},\mathcal{N})$ formed of the $\pi_0$-accessible functors. Then $\uHom_{\pi_0}(\mathcal{M},\mathcal{N})$ is accessible.}

\msp{page 373}

(The case where $\pi_0$ is useful for $\mathcal{M}$ has already been treated.)

\noindent\textsc{Proof}. In the general case, let $\pi>\pi_0$ be a useful cardinal for $\mathcal{M}$, so that we have
\[
\mathcal{M} \xleftarrow{\ \sim\ } \Ind_\pi(\mathcal{C}) \qquad\quad \text{with } \mathcal{C}=\mathcal{M}_\pi\ .
\]
We can suppose $\mathcal{M}=\Ind_\pi(\mathcal{C})$. The fact that $\mathcal{M}$ satisfies $\mathrm{L}_{\pi_0}$ implies that $\mathcal{C}$ is stable under $\varinjlim$ large relative to $\pi_0$, and of cardinality $\leqslant\pi$. Let $\Phi$ be the set of categories of this nature. $\Phi$ is not small, but there exists a small $\Phi_0\subset\Phi$, such that every $I\in\Phi$ is isomorphic to an $I_0\in\Phi_0$, which will suffice for the validity of 4.14.1 and 4.14.2 (where here we will put $\Psi=\emptyset$). We can consider $\uHom_{!\,\pi_0}(\mathcal{M},\mathcal{N})$ as a strictly full subcategory of $\uHom_{!\,\pi}(\mathcal{M},\mathcal{N})\xrightarrow{\ \sim\ }\uHom(\mathcal{C},\mathcal{N})$ (cf. 4.12.3, p. 275). For a functor $f:\mathcal{M}\longrightarrow\mathcal{N}$ which is $\pi$-accessible, to say that it is even $\pi_0$-accessible amounts to saying that its restriction to $\mathcal{C}$, i.e. its image in $\uHom(\mathcal{C},\mathcal{N})$, commutes with $\varinjlim$ of type $\Phi$. (Cf. 4.13.20, page 341 (\footnote[164]{The reference is not quite adequate. We should gather in a separate paragraph all the statements of this type, some of which use the indicial representation theorem.}), where we take $\alpha=\alpha_\pi$, noting that $J_{\alpha_\pi}$ large relative to $\pi_0$ since $\pi>\pi_0$).

\msp{page 374}

Thus we find
\[
\uHom_{!\,\pi_0}(\mathcal{M},\mathcal{N}) \xrightarrow{\ \sim\ } \uHom_{\Phi}(\mathcal{C},\mathcal{N})\ ,
\]
\printedpage{142}
% head: A. Grothendieck : Dérivateurs, Ch. XVIII, Édition des Universités de Montpellier II et Paris VII
% foot: Version 17 décembre 2025
and corollary 4.14.2 gives the desired conclusion.

\noindent\textbf{Corollary 4.14.4.} {\itshape Let $\mathcal{M}$, $\mathcal{N}$ be two accessible categories, $\Phi$ and $\Psi$ two sets of categories as in 4.14.2, we suppose that $\mathcal{M}$ and $\mathcal{N}$ satisfy $L_{\pi_0}$, and \edbr{[are]} stable under $\varinjlim$ of type $\Phi$ and under $\varprojlim$ of type $\Psi$. Let $\uHom^{\Psi}_{\Phi,\pi_0}(\mathcal{M},\mathcal{N})$ be the strictly full subcategory of $\uHom(\mathcal{M},\mathcal{N})$ formed of the $\pi_0$-accessible functors which commute with $\varinjlim$ of type $\Phi$, and with $\varprojlim$ of type $\Psi$. Then $\uHom^{\Psi}_{\Phi,\pi_0}(\mathcal{M},\mathcal{N})$ is accessible.}

\noindent\textsc{Proof}. This time, we choose $\pi>\pi_0$ such that we have $\cardred I\leqslant\pi$ for every $I$ in $\Phi$ or in $\Psi$. The category under consideration is then interpreted (putting again $\mathcal{C}=\mathcal{M}_\pi$) as
\msp{page 375}
\[
\uHom^{\Psi}_{\Phi,\pi_0}(\mathcal{M},\mathcal{N}) \xrightarrow{\ \sim\ } \uHom^{\Psi}_{\Phi'}(\mathcal{C},\mathcal{N})\ ,
\]
where $\Phi'$ is formed of the categories which are either in $\Phi$, or which are large relative to $\pi_0$ and of cardinality $\leqslant\pi$. It therefore suffices to apply 4.14.2 once more.

\noindent\textbf{Corollary 4.14.5.} {\itshape Let us consider the following three sets:
\begin{enumerate}[leftmargin=2.5em]
\item[\textit{a)}] $\Phi_1$ is formed of all categories equivalent to small categories.
\item[\textit{b)}] $\Phi_2$ as a), but restricting to filtering categories.
\item[\textit{c)}] $\Phi_3$ formed of the discrete categories, and of those which are large relative to $\pi_0$ ($\pi_0$ a given infinite cardinal).
\end{enumerate}
\noindent Let $\Phi$ be one of the sets of categories $\Phi_1$, $\Phi_2$, $\Phi_3$, $\Phi_2\cup\Phi_3$ or the union of one of these and a small set $\Phi_0$ of essentially small categories (\footnote[165]{We could give here a statement which includes 4.14.4 by adding a
\begin{enumerate}[leftmargin=2.5em]
\item[d)] $\Phi_4$ formed of the categories large relative to $\pi_0$,
\end{enumerate}
\noindent and take $\Phi$ of the form of $\Phi_1$, $\Phi_2$, $\Phi_4$, $\Phi_3$, $\Phi_3\cup\Phi_2$, $\Phi_3\cup\Phi_4$, or the union of these and a $\Phi_0$.}). Let $\Psi$ be a small set of categories. Let $\mathcal{M}$, $\mathcal{N}$ be accessible categories, stable under $\varinjlim$ of type $\Phi$, and under $\varprojlim$ of type $\Psi$. Then $\uHom^{\Psi}_{\Phi}(\mathcal{M},\mathcal{N})$ is accessible.}

\msp{page 376}

\noindent\textsc{Proof}. Let $\pi'_0\geqslant\pi_0$ be a useful cardinal for $\mathcal{M}$, $\pi$ a cardinal having the following properties.
\begin{enumerate}[leftmargin=2.5em]
\item[a)] $\pi\geqslant\sup(\pi'_0,\cardred\mathcal{M}_{\pi'_0})$, $\pi>\pi_0$.
\item[b)] $\pi^{\pi'_0}=\pi$.
\item[c)] $I\in\Phi_0\cup\Psi\Longrightarrow\cardred I\leqslant\pi$.
\end{enumerate}
\noindent Let $\mathcal{C}=\mathcal{M}_\pi$, whence $\mathcal{M}\simeq\Ind_\pi(\mathcal{C})$. In all the cases under consideration, $\Phi$ contains the set of essentially small categories which are large relative to $\pi_0$, a fortiori those which are large relative to $\pi$, hence the category $\uHom^{\Psi}_{\Phi}$ is interpreted as a strictly full subcategory of
\[
\uHom(\mathcal{C},\mathcal{N})\ .
\]
If on the other hand $\Phi_\pi$ denotes the set of $I\in\Phi$ such that $\card I\leqslant\pi$, we find
\[
\uHom^{\Psi}_{\Phi}(\mathcal{M},\mathcal{N})\ \simeq\ \uHom^{\Psi}_{\Phi_\pi}(\mathcal{C},\mathcal{N})\ ,
\]
and we apply 4.14.2.

\printedpage{143}
% head: A. Grothendieck : Dérivateurs, Ch. XVIII, Édition des Universités de Montpellier II et Paris VII
% foot: Version 17 décembre 2025
\msp{page 377}

In particular, we find:

\noindent\textbf{Corollary 4.14.6.} {\itshape Let $\mathcal{M}$, $\mathcal{N}$ be two accessible categories, suppose $\mathcal{M}$ stable under small $\varinjlim$ (i.e. $\mathcal{M}$ is a paratopos category). Then the category $\uHom_{!}(\mathcal{M},\mathcal{N})$ of functors $f:\mathcal{M}\longrightarrow\mathcal{N}$ which commute with small $\varinjlim$ is accessible. If $\mathcal{N}$ is itself also a paratopos, $\uHom_{!}(\mathcal{M},\mathcal{N})$ is a paratopos.}

\noindent\textsc{Proof}. The first assertion is the special case of 4.14.5 corresponding to the case where $\Phi=\Phi_1$, $\Psi=\emptyset$. The second assertion follows from it, since $\uHom_{!}(\mathcal{M},\mathcal{N})$ is visibly stable under small $\varinjlim$ , hence it is a paratopos being accessible (4.8.9, page 207).
\DERmark{ENV1391}


\noindent\textbf{Corollary 4.14.7.} {\itshape Let $\Phi$, $\Psi$, $\mathcal{M}$ be as in 4.14.5. Then for every accessible category $\mathcal{N}$, the category
\msp{page 378}
\[
\uHom^{\Phi}_{\Psi}(\mathcal{M}^{\circ},\mathcal{N})\ \simeq\ (\uHom^{\Psi}_{\Phi}(\mathcal{M},\mathcal{N}^{\circ}))^{\circ}
\]
is accessible.}

\noindent\textsc{Proof}. We proceed as for 4.14.5, to make explicit
\[
\uHom^{\Psi}_{\Phi}(\mathcal{M},\mathcal{N}^{\circ}) \xrightarrow{\ \sim\ } \uHom^{\Psi}_{\Phi_\pi}(\mathcal{C},\mathcal{N}^{\circ})\ ,
\]
whence
\[
\uHom^{\Phi}_{\Psi}(\mathcal{M}^{\circ},\mathcal{N}) \xrightarrow{\ \sim\ } \uHom^{\Phi_\pi}_{\Psi}(\mathcal{C}^{\circ},\mathcal{N})\ ,
\]
and we apply 4.14.2 to $\mathcal{C}^{\circ}$, $\Psi$, $\Phi_\pi$.

\noindent\textbf{Corollary 4.14.8.} {\itshape Let $\mathcal{M}$ and $\mathcal{N}$ be two paratopos categories. Then
\[
\mathcal{M}\boxtimes\mathcal{N}\ \simeq\ \uHom^{!}(\mathcal{M}^{\circ},\mathcal{N}) \qquad\quad \text{(cf. 4.2 and 4.3)}
\]
is a paratopos.}

Indeed, by virtue of 4.14.7 (with $\Phi=\Phi_1$ [?], $\Psi=\emptyset$), it is an accessible category.
\DERmark{ENV1392}


\msp{page 379}

On the other hand, it is immediate that it is stable under small $\varprojlim$ , $\mathcal{N}$ being so. Thus we conclude again by 4.8.9 (page 207) that it is a paratopos category.

\noindent\textbf{Corollary 4.14.9.} {\itshape Let $\mathcal{M}$, $\mathcal{N}$ be two paratopos categories, and $\uHom^{!}(\mathcal{M},\mathcal{N})$ the category of functors $\mathcal{M}\longrightarrow\mathcal{N}$ which commute with small $\varprojlim$ . We then have an equivalence of categories
\[
\begin{aligned}
\uHom^{!}(\mathcal{M},\mathcal{N}) &\ \xrightarrow{\ \sim\ }\ \uHom_{!}(\mathcal{N},\mathcal{M})^{\circ}\\
f &\ \longmapsto\ \text{left adjoint of } f\ ,
\end{aligned}
\]
\noindent\textup{\textbf{!}}\ \textup{\textbf{!}} and consequently the category \textup{\textbf{opposite}} to $\uHom^{!}(\mathcal{M},\mathcal{N})$ is accessible.} % Verified against PDF leaf 25: a nonverbal marginal curved mark; see the correction register.

This follows from the

\noindent\textbf{Lemma 4.14.10} (for the record). {\itshape Let $\mathcal{M}$ be \edbr{[a]} pseudo-topos category, $f:\mathcal{M}\longrightarrow\mathcal{N}$ a functor. For $f$ to admit a left adjoint (resp. a right adjoint), it is necessary and sufficient that $f$ commute with small $\varprojlim$ (resp. with small $\varinjlim$ ).}


\endgroup
% END UNIT der-ch18/der-ch18_p141-145.tex
% BEGIN UNIT der-ch18/der-ch18_p146-150.tex
\begingroup\makeatletter%
\newcommand{\printedpage}[1]{\marginpar{\footnotesize p.~#1}}%
\newcommand{\msp}[1]{\marginpar{\footnotesize [#1]}}%
\newcommand{\uHom}{\underline{\mathrm{Hom}}}%
\newcommand{\uFl}{\underline{\mathrm{Fl}}}%
\newcommand{\uIs}{\underline{\mathrm{Is}}}%
\newcommand{\uMon}{\underline{\mathrm{Mon}}}%
\newcommand{\uEp}{\underline{\textup{Ép}}}%
\newcommand{\uGamma}{\underline{\Gamma}}%
\newcommand{\defeq}{\overset{\mathrm{def}}{=}}%
\newcommand{\cardred}{\mathop{\textup{card\,red}}}%
\newcommand{\Ob}{\operatorname{Ob}}%
\newcommand{\Ind}{\operatorname{Ind}}%
\newcommand{\Ens}{\mathrm{Ens}}%
\newcommand{\circled}[1]{\tikz[baseline=(c.base)]\node[circle,draw,inner sep=1pt](c){\upshape #1};}%
\makeatother\let\printedpage\derpp\let\msp\dermsp%
\let\mspin\DERmspin%

\DERmark{ENV1393}


\printedpage{144}
% head: A. Grothendieck : Dérivateurs, Ch. XVIII, Édition des Universités de Montpellier II et Paris VII
% foot: Version 17 décembre 2025

\msp{page 380}
This is equivalent to the special case:

\textbf{Corollary 4.14.11.} \textit{Let $\mathcal{M}$ be a paratopos, and $f$ a covariant (resp. contravariant) functor from $\mathcal{M}$ to \textup{Ens}. For $f$ to be representable, it is necessary and sufficient that $f$ commute with $\varprojlim$ .}

This has been seen -- the case of contrafunctors being the most obvious (cf. 4.1 -- it suffices that $\mathcal{M}$ be a paratopos, and even that $\mathcal{M}$ admit a small generating full subcategory, cf. SGA 4 I 8.12.7 (for the dual statement). In the case of a covariant functor, the difficulty comes from the fact that $\mathcal{M}$ does not necessarily admit a small cogenerating subcategory. One must use the existence of a (generalized) cardinal filtration, cf. loc. cit. 8.13.2.) % [sic] Dans les cas -> Dans le cas, singular "un foncteur covariant"
\DERmark{ENV1394}\DERmark{ENV1395}


After all these fine corollaries, we are ripe for tackling the proof of 4.14.1!

\msp{page 381}
\textsc{Proof of} 4.14.1. We know that $\mathcal{M} = \uHom(\mathcal{C},\mathcal{N})$ is accessible (4.9.5, p. 220), and $\uHom^{\Lambda_2}_{\Lambda_1}(\mathcal{C},\mathcal{N})$ is a strictly full subcategory of this category. It suffices to prove that it is an accessible subcategory in $\mathcal{M}$ (4.13.1, p. 281). For every $\xi = (I,\varphi) \in \Lambda_1$, we have a functor
\[
L_\xi \ : \ \mathcal{M} \ = \ \uHom(\mathcal{C},\mathcal{N}) \longrightarrow \uFl(\mathcal{N}) \ ,
\]
given by
\[
L_\xi(f) \ = \ \text{canonical arrow } \varinjlim_{i\in I} f\varphi(i) \longrightarrow f\varphi(e_I) \ .
\]
I say that the functor $L_\xi$ is \textit{accessible.} Indeed, let $\pi_0$ be a cardinal such that $\mathcal{N}$ satisfies $\mathrm{L}_{\pi_0}$, the same therefore holds for $\mathcal{M} = \uHom(\mathcal{C},\mathcal{N})$ and for $\uFl(\mathcal{N}) \defeq \uHom(\boldsymbol{\Delta}_1,\mathcal{N})$, the limits in question are computed argument by argument. I say that $L_\xi$ is $\pi_0$-accessible, i.e. commutes with $\varinjlim$ large relative to $\pi_0$. It suffices to see that this is so for the functors composed with the functors
\[
\begin{tikzcd}
\uFl(\mathcal{N}) \arrow[r, shift left, "\sigma"] \arrow[r, shift right, "\rho"'] & \mathcal{N}
\end{tikzcd}
\]
``source'' and ``target'', i.e. the functors
\[
f \longmapsto \varinjlim_{i\in I} f\varphi(i) \qquad \text{and} \qquad f \longmapsto f\varphi(e_I) \ .
\]

\msp{page 382}
Now the first functor, $f \longmapsto \varinjlim\limits_{i\in I} f\varphi(i)$, is the composite of
\[
f \longmapsto f\circ\varphi_I \ : \ \uHom(\mathcal{C},\mathcal{N}) \longrightarrow \uHom(I,\mathcal{N})
\]
(which commutes with the $\varinjlim$ that are representable in $\mathcal{N}$, and in particular with the $\varinjlim$ large relative to $\pi_0$) and of the functor
\[
\varinjlim_{I} \ : \ \uHom(I,\mathcal N) \longrightarrow \mathcal{N} \ , % [sic] \uHom(I,N) -> \uHom(I,\mathcal{N}), italic N printed instead of the script N used everywhere else
\]
\DERmark{ENV1396}

which also commutes with direct limits representable in $\mathcal{N}$, and in particular with those large relative to $\pi_0$. As for the functor $f \longmapsto (f\circ\varphi)(e_I) = f(\varphi(e_I))$, it is clear that it too commutes with the $\varinjlim$ representable in $\mathcal{N}$, so it is $\pi_0$-accessible.

\printedpage{145}
% head: A. Grothendieck : Dérivateurs, Ch. XVIII, Édition des Universités de Montpellier II et Paris VII
% foot: Version 17 décembre 2025

This being so, let us note that
\[
\uHom_{\Lambda_1}(\mathcal{C},\mathcal{N}) \ = \ \bigcap_{\xi\in\Lambda_1} \uHom_{\{\xi\}}(\mathcal{C},\mathcal{N}) \ ,
\]
and that $\uHom_\xi(\mathcal{C},\mathcal{N})$ is the inverse image, under the preceding accessible functor $L_\xi$, of the strictly full subcategory $\uIs(\mathcal{N})$ of $\uFl(\mathcal{N})$ consisting of the \textit{isomorphisms} in $\mathcal{N}$. If we prove that the latter is accessible in $\uFl(\mathcal{N})$, it follows that
\msp{page 383}
$\uHom_\xi(\mathcal{C},\mathcal{N})$ is accessible in $\mathcal{M}$ (4.13.31, page 350). Thus, $\Lambda_1$ being small, it follows that $\uHom_{\Lambda_1}(\mathcal{C},\mathcal{N})$ is accessible in $\mathcal{M}$ (4.13.7, page 291).
\DERmark{ENV1397}\DERmark{ENV1398}


We prove in the same way that $\uHom^{\Lambda_2}(\mathcal{C},\mathcal{N})$ is accessible in $\mathcal{M} = \uHom(\mathcal{C},\mathcal{N})$. Here, we must consider the functors % [sic] \mathrm{Hom}(\mathcal{C},\mathcal{N}) -> \uHom(\mathcal{C},\mathcal{N}), printed without underline, elsewhere underlined
\DERmark{ENV1399}

\[
L'_\xi \ : \ \mathcal{M} = \uHom(\mathcal{C},\mathcal{N}) \longrightarrow \uFl(\mathcal{N})
\]
defined by
\[
L'_\xi \ = \ \text{canonical arrow } f\varphi(e'_I) \longrightarrow \varprojlim_{i\in I} f\varphi(i) \ .
\]
To prove that this functor is accessible, we proceed as before, taking however $\pi_0$ large enough that we have
\[
\left\{
\begin{array}{l}
I \in \Psi \ \Longrightarrow \ \cardred I \leqslant \pi_0. \\[1ex]
\mathcal{N}_{\pi_0} \text{ is generating in } \mathcal{N}.
\end{array}
\right.
\]
I say that then $L'_\xi$ is $\pi_0$-accessible: indeed, proceeding as before for $L_\xi$, we are reduced to seeing that
\[
\varprojlim_I: \uHom(I,\mathcal{N}) \longrightarrow \mathcal{N}
\]
\DERmark{ENV1400}

is $\pi_0$-accessible. Now this is contained in SGA 4 I 9.8.
\DERmark{ENV1401}


\msp{page 384}
Thus, we have proved that the two strictly full subcategories
\[
\uHom_{\Lambda_1}(\mathcal{C},\mathcal{N}) \ , \qquad \uHom^{\Lambda_2}(\mathcal{C},\mathcal{N})
\]
of $\mathcal{M} \defeq \uHom(\mathcal{C},\mathcal{N})$ are accessible in $\mathcal{M}$. The same therefore holds for their intersection (4.13.7)
\[
\uHom^{\Lambda_2}_{\Lambda_1}(\mathcal{C},\mathcal{N}) \ = \ \uHom_{\Lambda_1}(\mathcal{C},\mathcal{N}) \cap \uHom^{\Lambda_2}(\mathcal{C},\mathcal{N}) \ ,
\]
q.e.d.

It remains to prove the

\textbf{Lemma 4.14.12.} \textit{Let $\mathcal{N}$ be an accessible category. Then the category $\uFl(\mathcal{N}) \defeq \uHom(\boldsymbol{\Delta}_1,\mathcal{N})$ is accessible (4.9.5), and the strictly full subcategory $\uIs(\mathcal{N})$ consisting of the invertible arrows in $\mathcal{N}$ is accessible in $\uFl(\mathcal{N})$. (The same holds for the strictly full subcategory $\uMon(\mathcal{N})$ consisting of the monomorphisms in $\mathcal{N}$.)}

\textsc{Proof}. Let $\mathcal{C}$ be a small strictly generating subcategory of $\mathcal{N}$, and consider the small family of functors
\[
h_x \ = \ (y \longmapsto \mathrm{Hom}(x,y)) \ : \ \mathcal{N} \longrightarrow \Ens \ ,
\]
\printedpage{146}
% head: A. Grothendieck : Dérivateurs, Ch. XVIII, Édition des Universités de Montpellier II et Paris VII
% foot: Version 17 décembre 2025
\msp{page 385}
whence functors
\[
\uFl(h_x) \ : \ \uFl(\mathcal{N}) \longrightarrow \uFl(\Ens) \ . % [sic] h_\xi -> h_x, the functors defined just before are h_x
\]
\DERmark{ENV1402}

For an arrow $u : y \longrightarrow z$ in $\mathcal{N}$ to be an isomorphism (resp. a monomorphism), it is necessary and sufficient that for every $x \in \Ob\mathcal{C}$, $h_x(u)$ be so. Thus, $\uIs(\mathcal{N})$ (resp. $\uMon(\mathcal{N})$) appears as the intersection of the inverse images, under the functors $\uFl(h_x)$ above, of the subcategory $\uIs(\Ens)$ (resp. $\uMon(\Ens)$) of $\uFl(\Ens)$. This reduces us to the case where $\mathcal{N} = (\Ens)$. But if $\mathcal{C}$ denotes the strictly full subcategory of Ens consisting of the \textit{finite} sets, these form a strictly generating family consisting of objects of \textit{finite presentation}, and since Ens is stable under small $\varinjlim$ , hence under filtering $\varinjlim$, it follows that we have % [sic] h_\xi -> h_x, the functors defined just before are h_x
\DERmark{ENV1403}

\[
\Ind(\mathcal{C}) \longrightarrow (\Ens) \ , \tag*{$(*)$}
\]
whence at once (SGA 4 I 8.8.5, p. 104)

\msp{page 386}
\[
\Ind(\uFl(\mathcal{C})) \xrightarrow{\ \sim\ } \uFl(\Ens) \ . \tag*{$(*)$}
\]
So the desired accessibility result follows from the following two lemmas:

\textbf{Lemma 4.14.13.} \textit{The preceding equivalence $(*)$ induces equivalences}
\[
\begin{aligned}
\Ind(\uIs(\mathcal{C})) \quad &\xrightarrow{\ \sim\ } \quad \uIs(\Ens) \\
\Ind(\uMon(\mathcal{C})) \quad &\xrightarrow{\ \sim\ } \quad \uMon(\Ens)
\end{aligned}
\]
\textit{(where $\mathcal{C} = (\mathrm{Ensf})$ is the full subcategory of \textup{(Ens)} consisting of the \textup{finite} sets.)}

\textbf{Lemma 4.14.14} (\footnote[166]{\textbf{N.B.} Up to equivalence, one finds the accessible categories $\mathcal{M}$ stable under small filtering $\varinjlim$, and such that $\mathcal{M}_{\mathrm{pf}}$ (subcategory of objects of \textit{finite} presentation) is strictly generating.}). \textit{Let $\mathcal{C}$ be a category equivalent to a small category. Then:}
\begin{itemize}[nosep]
\item[\textit{a)}] \textit{$\mathcal{M} \defeq \Ind(\mathcal{C})$ is an accessible category.}
\item[\textit{b)}] \textit{For every full subcategory $\mathcal{D}$ in $\mathcal{C}$, the essential image of $\Ind\mathcal{D}$ in $\mathcal{M} = \Ind(\mathcal{C})$ is an L-accessible subcategory in $\mathcal{M}$.}
\end{itemize}
\textsc{Proof} of the two lemmas immediate and ``left to the reader''. The b) in 4.14.14 comes from the fact that the functor $\Ind(\mathcal{D}) \longrightarrow \Ind(\mathcal{C})$ coming from a functor $\mathcal{D} \longrightarrow \mathcal{C}$ (here an
\msp{page 387}
inclusion functor) is an L-functor (i.e. commutes with small filtering $\varinjlim$), and a fortiori is accessible (and even $\aleph_0$-accessible).

\textbf{4.14.15.} One cannot help wondering whether for $\mathcal{N}$ accessible, the strictly full subcategory $\uEp(\mathcal{N})$ in $\uFl(\mathcal{N})$, consisting of the epimorphisms in $\mathcal{N}$, is accessible in $\uFl(\mathcal N)$. It is immediate that if $\mathcal{N}$ satisfies $\mathrm{L}_{\pi_0}$, then $\uEp(\mathcal{N})$ is stable in $\uFl(\mathcal{N})$ under the $\varinjlim$ of type $\mathrm{L}_{\pi_0}$, a fortiori under direct factors. It would therefore have a hard time not being accessible in $\uFl(\mathcal{N})$. Let us note all the same: % [sic] dû mal -> du mal, "avoir du mal à"
\DERmark{ENV1404}


\textbf{Corollary 4.14.15.1} (of 4.14.12). \textit{Let $\mathcal{N}$ be an accessible category stable under amalgamated sums $y \amalg_x z$} (\footnote[167]{I am almost sure that this hypothesis is superfluous. Compare page 390.}). \textit{Then $\uEp(\mathcal{N})$ is an accessible subcategory in $\uFl(\mathcal{N})$.}
\DERmark{ENV1405}


\printedpage{147}
% head: A. Grothendieck : Dérivateurs, Ch. XVIII, Édition des Universités de Montpellier II et Paris VII
% foot: Version 17 décembre 2025

\msp{page 388}
Let us indeed consider the functor
\[
\begin{aligned}
\gamma \ : \ \uFl(\mathcal{N}) \quad &\longrightarrow \quad \uFl(\mathcal{N}) \\
(x \longrightarrow y) \quad &\longmapsto \quad (y \amalg_x y \xrightarrow{\ \ \text{codiag}\ \ } y) \ .
\end{aligned}
\]
This functor commutes with the types of $\varinjlim$ that are representable in $\mathcal{N}$, in particular, it is accessible. On the other hand, $\uEp(\mathcal{N})$ is the inverse image under $\gamma$ of the strictly full subcategory $\uIs(\mathcal{N})$ of $\uFl(\mathcal{N})$. So the corollary follows from the fact that by 4.14.12 the latter is accessible in $\uFl(\mathcal{N})$.
\DERmark{ENV1406}


\textbf{4.14.16.} I presume that 4.14.1 remains valid without assuming $\mathcal{N}$ stable under the $\varinjlim$ of type $\Phi$ and under the $\varprojlim$ of type $\Psi$. If we have an essentially small category $I$, the question is whether the strictly full subcategory $\uHom_{\square}(I_\bullet,\mathcal{N})$ of $\uHom(I_\bullet,\mathcal{N})$, consisting of the functors satisfying
\msp{page 389}
condition a) of 4.14.1, is accessible in $\uHom(I_\bullet,\mathcal{N})$, and likewise for $\uHom^{\square}({}_\bullet I,\mathcal{N})$. For this second case, it looks OK, for if $\mathcal{C}$ is a small strictly generating subcategory of $\mathcal{N}$, and if for every $x \in \Ob$ \texttt{[}$\mathcal{C}$\texttt{]} we again consider the functor
\[
h_x \ : \ \mathcal{N} \longrightarrow \Ens \ ,
\]
whence
\[
\bar{h}_x \ : \ \uHom({}_\bullet I,\mathcal{N}) \longrightarrow \uHom({}_\bullet I,\Ens) \ ,
\]
we see that $\uHom^{\square}({}_\bullet I,\mathcal{N})$ is the intersection of the inverse images of the subcategories $\uHom^{\square}({}_\bullet I,\Ens)$ of $\uHom({}_\bullet I,\Ens)$. So the latter being accessible in $\uHom({}_\bullet I,\Ens)$, by 4.14.1 ((Ens) being stable under small $\varprojlim$ ), it does indeed follow that $\uHom^{\square}({}_\bullet I,\mathcal{N})$ is so in $\uHom({}_\bullet I,\mathcal{N})$, q.e.d.

The case of $\uHom_{\square}(I_\bullet,\mathcal{N})$ in $\uHom(I_\bullet,\mathcal{N})$ seems more delicate, in the case where
\msp{page 390}
\[
I = 
\begin{tikzcd}[sep=small]
 & {} \\
\bullet \arrow[ur] \arrow[dr] & \\
 & {}
\end{tikzcd}
\quad , \qquad \text{thus} \qquad
I_\bullet = 
\begin{tikzcd}[sep=small]
 & {} \arrow[dr] & \\
\bullet \arrow[ur] \arrow[dr] & & {} \\
 & {} \arrow[ur] &
\end{tikzcd}
\quad ,
\]
\DERmark{ENV1407}

if the envisaged statement is true, we will conclude from it that $\uEp(\mathcal{N})$ is accessible in $\uFl(\mathcal{N})$, for every accessible category, for it is the inverse image of $\uHom_{\square}(I_\bullet,\mathcal{N})$ under the canonical (accessible) functor
\[
\begin{aligned}
\uFl(\mathcal{N}) \quad &\longrightarrow \quad \uHom(I_\bullet,\mathcal{N}) \\
(x \xrightarrow{\ u\ } y) \quad &\longmapsto \quad
\left(
\begin{tikzcd}[sep=small, ampersand replacement=\&]
 \& y \arrow[dr, "\mathrm{id}_y"] \& \\
x \arrow[ur, "u"] \arrow[dr, "u"'] \& \& y \\
 \& y \arrow[ur, "\mathrm{id}_y"'] \&
\end{tikzcd}
\right) \ .
\end{aligned}
\]

\printedpage{148}
% head: A. Grothendieck : Dérivateurs, Ch. XVIII, Édition des Universités de Montpellier II et Paris VII
% foot: Version 17 décembre 2025

This is to say that the envisaged generalization of 4.14.1 is not idle, but is probably useful.

\msp{page 391}
\textbf{Theorem 4.14.17} (\footnote[168]{Cf. generalization in \texttt{[the]} theorem page 422, th. 4.14.26.}). \textit{Let $p : \mathcal{E} \longrightarrow \mathcal{C}$ be a cofibrant or fibrant functor, with $\mathcal{C}$ essentially small (i.e. equivalent to a small category), and such that the fiber categories are accessible, and the cobase change (cofibrant case) or base change (fibrant case) functors associated with the arrows $u$ of $\mathcal{C}$ are accessible. Then the category $\uGamma(\mathcal{E}/\mathcal{C})$ of sections of $\mathcal{E}$ over $\mathcal{C}$ is accessible, and for every $s \in \Ob\mathcal{C}$, the ``evaluation at $s$'' functor $\varepsilon_s : \uGamma(\mathcal{E}/\mathcal{C}) \longrightarrow \mathcal{E}_s$ is accessible.}

\textbf{Corollary 4.14.18.} \textit{Let $\pi_0$ be an infinite cardinal such that the fiber categories satisfy the condition $L_{\pi_0}$, and that the cobase change (resp. base change) functors between them likewise satisfy $L_{\pi_0}$. (Such a $\pi_0$ exists.)}
\begin{itemize}[nosep]
\item[\circled{a}] \textit{Then $\mathcal{M} = \uGamma(\mathcal{E}/\mathcal{C})$ satisfies the condition $\underline{\mathrm{L}}_{\pi_0}$, and the functors $\varepsilon_s$ too, i.e. the $\varinjlim$ large relative to $\pi_0$ in $\uGamma(\mathcal{E}/\mathcal{C})$ are computed argument by argument.}
\item[\circled{b}] \textit{Let $\pi \geqslant \sup(\pi_0,\cardred\mathcal{C})$ be a cardinal. For an object $f$ of $\mathcal{M}$ to belong to $\mathcal{M}_\pi$, it \textup{suffices} that $\forall\ s \in \Ob\mathcal C$, $f(s) \defeq \varepsilon_s(x)$ \textup{\texttt{[rather} $\varepsilon_s(f)$\texttt{]}} be in $\mathcal{E}_{s\,\pi}$. (Compare 4.5.8, page 119.)} % [sic] \Ob C -> \Ob\mathcal{C}, italic C printed instead of script C
\item[\circled{c}] \textit{Suppose $\pi_0$ chosen so as} % [sic] pi_0 -> \pi_0, the letters "pi" printed instead of the Greek letter
\msp{page 392}

\textit{to satisfy the following additional conditions: $\pi_0$ is useful for $\mathcal{M} = \uGamma(\mathcal{E}/\mathcal{C})$, and for each of the fiber categories $\mathcal{E}_s$ ($s \in \Ob\mathcal{C}$) (such a $\pi_0$ exists). Let}
\DERmark{ENV1408}\DERmark{ENV1409}

\[
\pi_1 \ = \ \sup(\pi_0, \cardred\mathcal{C}, \cardred\mathcal{M}_{\pi_0}) \ .
\]
\textit{Then for every cardinal $\pi \geqslant \pi_1$ and such that $\pi^{\pi_0} = \pi$, we have}
\[
\mathcal{M}_\pi \ = \ \left\{
\begin{array}{l}
\textit{strictly full subcategory of $\mathcal{M}$ consisting of the $f \in$} \\
\textit{$\Ob\mathcal{M}$ such that $f(s) \in \Ob\mathcal{E}_{s\,\pi}$ $\forall\ s \in S$ \textup{\texttt{[rather}} $s \in$} \\
\textit{$\Ob C$\textup{\texttt{]}}.} % [sic] \Ob C -> \Ob\mathcal{C}, italic C printed instead of script C
\end{array}
\right.
\]
\textit{If moreover $\pi$ is chosen so that we have}
\[
\pi \ \geqslant \ \sup_{s\in\Ob\mathcal{C}} \cardred\mathcal{E}_{s\,\pi_0} \ ,
\]
\textit{then the strictly full subcategory $\mathcal{E}(\pi)$ of $\mathcal{E}$ such that $\Ob\mathcal{E}(\pi) = \bigcup_{s\in\Ob\mathcal{C}} \Ob\mathcal{E}_{s\,\pi}$ is a $\mathcal{C}$-cofibered (resp. $\mathcal{C}$-fibered) subcategory of $\mathcal{E}$, and we have}
\[
\uGamma(\mathcal{E}/\mathcal{C})_\pi \xleftarrow{\ \sim\ } \uGamma(\mathcal{E}(\pi)/\mathcal{C}), % [sic] C -> \mathcal{C}, italic C printed instead of script C
\]
\DERmark{ENV1410}

\textit{thus ($\pi$ being useful for $\mathcal{M}=\uGamma(\mathcal{E}/\mathcal{C})$) we have an equivalence of categories}
\DERmark{ENV1411}

\[
\uGamma(\mathcal{E}/\mathcal{C}) \xleftarrow{\ \sim\ } \Ind_\pi(\uGamma(\mathcal{E}(\pi)/\mathcal{C})) \ .
\]
\end{itemize}


\endgroup
% END UNIT der-ch18/der-ch18_p146-150.tex
% BEGIN UNIT der-ch18/der-ch18_p151-155.tex
\begingroup\makeatletter%
\newcommand{\printedpage}[1]{\marginpar{\footnotesize p.~#1}}%
\newcommand{\msp}[1]{\marginpar{\footnotesize [#1]}}%
\newcommand{\uGamma}{\underline{\Gamma}}%
\newcommand{\uFl}{\underline{\mathrm{Fl}}}%
\newcommand{\uIs}{\underline{\mathrm{Is}}}%
\newcommand{\uHom}{\underline{\mathrm{Hom}}}%
\newcommand{\Ob}{\mathrm{Ob}\,}%
\newcommand{\Fl}{\mathrm{Fl}\,}%
\newcommand{\id}{\mathrm{id}}%
\newcommand{\pr}{\mathrm{pr}}%
\newcommand{\diag}{\mathrm{diag}\,}%
\newcommand{\cardred}{\operatorname{card\,red}}%
\newcommand{\cocart}{\mathrm{cocart}}%
\newcommand{\cart}{\mathrm{cart}}%
\newcommand{\cc}{\mathrm{c}}%
\newcommand{\circled}[1]{\tikz[baseline=(c.base)]\node[draw,circle,inner sep=0.8pt](c){\upshape #1};}%
\makeatother\let\printedpage\derpp\let\msp\dermsp%
\let\mspin\DERmspin%


\printedpage{149}
% head: A. Grothendieck : Dérivateurs, Ch. XVIII, Édition des Universités de Montpellier II et Paris VII

\msp{page 393}
All these assertions are easily proved in the order in which they are given, using 4.14.17 and the general results already proved. Of course, we can begin by replacing $\mathcal{C}$ by an equivalent reduced subcategory $\mathcal{C}_0$ and $\mathcal{E}$ by $\mathcal{E}_0 = \mathcal{E}|\mathcal{C}$ \texttt{[rather} $\mathcal{E}|\mathcal{C}_0$\texttt{]}, so as to reduce (in b) and c)) to the case where $\operatorname{card}\mathcal{C} \leqslant \pi$.

\textbf{Corollary 4.14.19.} \textit{Under the conditions of 4.14.17, the subcategory $\uGamma_{\cocart}(\mathcal{E}/\mathcal{C})$ of $\mathcal{M} = \uGamma(\mathcal{E}/\mathcal{C})$ (case $p$ cofibrant), resp. the subcategory $\uGamma_{\cart}(\mathcal{E}/\mathcal{C})$ (case $p$ fibrant) is accessible in $\mathcal{M}$, hence accessible and the inclusion functor} \texttt{[is]} \textit{accessible.}

(\textbf{N.B.} These two categories are strictly full in $\uGamma(\mathcal{E}/\mathcal{C})$, so the assertion that they are accessible in it \textit{is equivalent} to the last assertion.)

\textsc{Proof of} 4.14.19. We may assume $\mathcal{C}$ reduced, hence small. Let us consider the case where $p$ is cofibrant. For every arrow $\lambda : s \longrightarrow t$ in $\mathcal{C}$, let us consider the functor
\[
\varphi_\lambda \ : \ \mathcal{M} \longrightarrow \uFl(\mathcal{E}_t)
\]
defined by
\msp{page 394}
\[
\varphi_\lambda(f) \ = \ (\lambda_*(f(s)) \xrightarrow{f(\lambda)^\natural} f(t)) \ .
\]
It is an accessible functor, since its composite with the source and target functors are the functors $\lambda_*\varepsilon_s$ and $\varepsilon_t$, hence accessible, $\varepsilon_s$, $\varepsilon_t$ and $\lambda_*$ being accessible. Let us consider the strictly full subcategory
\[
\uIs(\mathcal{E}_t) \ \subset \ \uFl(\mathcal{E}_t) \ ,
\]
we know that it is accessible in $\uFl(\mathcal{E}_t)$ (4.14.12). Thus its inverse image under $\varphi_\lambda$ is accessible in $\mathcal{M}$, and consequently so is
\[
\uGamma_{\cocart}(\mathcal{E}/\mathcal{C}) \ \overset{\mathrm{def}}{=} \ \bigcap_{\lambda\in\Fl\mathcal{C}} \varphi_\lambda^{-1}(\uIs(\mathcal{E}_t))
\]
\DERmark{ENV1412}\DERmark{ENV1413}\DERmark{ENV1414}

% [sic] \uGamma_{\cart} -> \uGamma_{\cocart}, the case treated is p cofibrant
% [sic] \lambda\in\Ob\mathcal{C} -> \lambda\in\Fl\mathcal{C}, \lambda runs over arrows of \mathcal{C}
as the intersection of a \textit{small} family of subcategories accessible in $\mathcal{M}$, q.e.d. The case where $p$ is fibrant is treated in a similar way, by introducing the functors
\[
\psi_\lambda \ : \ \mathcal{M} = \uGamma(\mathcal{E}/\mathcal{C}) \longrightarrow \uFl(\mathcal{E}_s)
\]
given by
\[
\psi_\lambda(f) \ = \ (f(s) \xrightarrow{f(\lambda)^\flat} \lambda^* f(t)) \ .
\]

\msp{page 395}
\textbf{Corollary 4.14.20.} \textit{Let $\pi_0$ be an infinite cardinal chosen as in 4.14.18 a).}
\begin{itemize}[leftmargin=2.5em]
\item[\circled{a}] \textit{Then $\mathcal{M}_{\cc} = \uGamma_{\cocart}(\mathcal{E}/\mathcal{C})$ (resp. $\mathcal{M}_{\cc} = \uGamma_{\cart}$) is $L_{\pi_0}$-stable in $\mathcal{M} = \uGamma(\mathcal{E}/\mathcal{C})$, or what amounts to the same thing (as it is strictly full), it satisfies $L_{\pi_0}$, and the inclusion functor satisfies $L_{\pi_0}$. The $\varinjlim$ large relative to $\pi_0$ in this category are computed argument by argument.}
\item[\circled{b}] \textit{For every cardinal $\pi \geqslant \pi_0$, every object of $\mathcal{M}_{\cc}$ which is $\pi$-accessible in $\mathcal{M}$ is so (tautologically) in $\mathcal{M}_{\cc}$. If $\pi \geqslant \sup(\pi_0, \cardred\mathcal{C})$, then (by 4.14.18 b) every object $f$ of $\mathcal{M}_{\cc}$ such that $f(s) \in \Ob\mathcal{E}_{s\,\pi}$ $\forall\ s \in S$} \texttt{[rather} $s \in \Ob C$\texttt{]} \textit{is in $\mathcal{M}_\pi$} \textup{and} \textit{in $\mathcal{M}_{\cc\,\pi}$.}
% foot: Version 17 décembre 2025

\printedpage{150}
% head: A. Grothendieck : Dérivateurs, Ch. XVIII, Édition des Universités de Montpellier II et Paris VII

\item[\circled{c}] \textit{Suppose $\pi_0$ chosen as in 4.14.18 c, and in addition that $\mathcal{M}_{\cc}$ is $\pi_0$-accessible in $\mathcal{M}$. (Such a $\pi_0$ exists.) We thus have, for $\pi \geqslant \pi_1 = \sup(\pi_0, \cardred\mathcal{C}, \cardred\mathcal{M}_{\pi_0})$ such that $\pi^{\pi_0} = \pi$,}
\[
(\mathcal{M}_{\cc})_\pi \ = \ \mathcal{M}_\pi{\cap}\mathcal{M}_{\cc} \ = \quad
\begin{array}{l}
\textit{Strictly full subcategory of $\mathcal{M}_{\cc}$ consisting of the $f$}\\
\textit{such that $f(s) \in \Ob\mathcal{E}_{s\,\pi}$ $\forall\ s \in \Ob\mathcal{C}$.}
\end{array}
\]
\DERmark{ENV1415}


\msp{page 396}
\textit{If in addition $\pi$ is chosen so that we have $\pi \geqslant \sup_{s\in\Ob\mathcal{C}} \cardred\mathcal{E}_{s\,\pi_0}$, then}
\[
\left\{
\begin{array}{rrcl}
 & \mathcal{M}_{\cc\,\pi} & \xleftarrow{\ \sim\ } & \uGamma_{\cocart}(\mathcal{E}(\pi)/\mathcal{C})\\[1ex]
\textit{resp.} & \mathcal{M}_{\cc\,\pi} & \xleftarrow{\ \sim\ } & \uGamma_{\cart}(\mathcal{E}(\pi)/\mathcal{C}) \ ,
\end{array}
\right.
\]
\textit{hence ($\pi$ being useful for $\mathcal{M}_{\cc}$) an equivalence of categories}
\[
\begin{array}{rrcl}
 & \uGamma_{\cocart}(\mathcal{E}/\mathcal{C}) & \xleftarrow{\ \sim\ } & \mathrm{Ind}_\pi(\uGamma_{\cocart}(\mathcal{E}(\pi)/\mathcal{C}))\\
\textit{resp.}\quad & \uGamma_{\cart}(\mathcal{E}/\mathcal{C}) & \xleftarrow{\ \sim\ } & \mathrm{Ind}_\pi(\uGamma_{\cart}(\mathcal{E}(\pi)/\mathcal{C})).
\end{array}
\]
\end{itemize}

\textsc{Proof}. Everything is clear in the order in which it is given using the two preceding corollaries.

\textbf{Corollary 4.14.21.} \textit{Let us consider a diagram of functors}
\[
\begin{tikzcd}[column sep=small, row sep=small]
\mathcal{E}_1 \arrow[dr, "p_1"] & & \mathcal{E}_2 \arrow[dl, "p_2"'] \\
& \mathcal{E}_0 \ , &
\end{tikzcd}
\tag*{$(*)$}
\]
\textit{the categories and the functors being accessible.}

\msp{page 397}
\begin{itemize}[leftmargin=2.5em]
\item[\textit{a)}] \textit{The 2-fiber product}
\[
\mathcal{M}_0 \ = \ \mathcal{E}_1 \overset{2}{\times}_{\mathcal{E}_0} \mathcal{E}_2
\]
\textit{is accessible, and the canonical functors $\mathcal{M} \longrightarrow \mathcal{E}_1$, $\mathcal{M} \longrightarrow \mathcal{E}_2$ are accessible.} % [sic] \mathcal{M} \longrightarrow \mathcal{E}_1, \mathcal{M} \longrightarrow \mathcal{E}_2 -> \mathcal{M}_0 \longrightarrow \mathcal{E}_1, \mathcal{M}_0 \longrightarrow \mathcal{E}_2, assertion a) concerns \mathcal{M}_0
\item[\textit{b)}] (\footnote[169]{This is contained in 4.14.22.1 above (p. 406), applied to $\Pi = \mathcal{E}_1{\times}\mathcal{E}_2$, $\Pi' = \mathcal{E}_0$, and $F, G : \Pi \rightrightarrows \Pi'$ given by $F = p_1\circ\pr_1$ and $G = p_2\circ\pr_2$. On the other hand, a) and c) follow easily from b).}). \textit{The category $\mathcal{M}$ of triples $(x_1,x_2,u)$ with $x_1 \in \Ob\mathcal{E}_1$, $x_2 \in \Ob\mathcal{E}_2$, and}
\[
u \ : \ p_1(x_1) \ \longrightarrow \ p_2(x_2) \ ,
\]
\textit{is accessible, and the canonical functors $\mathcal{M} \longrightarrow \mathcal{E}_1$, $\mathcal{M} \longrightarrow \mathcal{E}_2$ are accessible.}
\item[\textit{c)}] \textit{The strictly full subcategory $\mathcal{M}_0$ of $\mathcal{M}$ is accessible in $\mathcal{M}$.}
\end{itemize}

\noindent\textsc{Proof}. Let $\mathcal{C}$ be the category
\[
\begin{tikzcd}[column sep=small]
1 \arrow[dr] & & 2 \arrow[dl] \\
& 0 \ , &
\end{tikzcd}
\]
and $\mathcal{E}$ the cofibered category over $\mathcal{C}$ defined by the diagram $(*)$. We are under the conditions of 4.14.17. The category $\Gamma = \uGamma(\mathcal{E}/\mathcal{C})$ is interpreted
% foot: Version 17 décembre 2025
\printedpage{151}
% head: A. Grothendieck : Dérivateurs, Ch. XVIII, Édition des Universités de Montpellier II et Paris VII
\msp{page 398}
as the category of quintuples $(x_0,x_1,x_2,u_1,u_2)$, with $x_i \in \Ob\mathcal{E}_i$ for $i \in \{0,1,2\}$, and $u_1, u_2 \in \Fl\mathcal{E}_0$,
\DERmark{ENV1416}

\[
\begin{tikzcd}[column sep=small, row sep=small]
p_1(x_1) \arrow[dr, "u_1"] & & p_2(x_2) \arrow[dl, "u_2"'] \\
& x_0 \ . &
\end{tikzcd}
\tag*{$(**)$}
\]
By 4.14.18, this category $\uGamma$ is accessible. The category $\mathcal{M}$ is equivalent to the strictly full subcategory $\Gamma_1$ of $\Gamma$, consisting of the quintuples for which we have $u_2 \in \uIs\mathcal{E}_0$. We have indeed functors quasi-inverse to each other % [sic] \Gamma_0 -> \Gamma_1, the subcategory is called \Gamma_1 in what follows
\DERmark{ENV1417}

\[
\begin{tikzcd}
\mathcal{M} \arrow[r, shift left, "\alpha"] & \Gamma_1 \arrow[l, shift left, "\beta"]
\end{tikzcd}
\]
given by
\[
\begin{aligned}
\alpha(x_1,x_2,u) \ &= \ (p_2(x_2),x_1,x_2,u,\id_{p_2(x_2)})\\
\beta(x_0,x_1,x_2,u_1,u_2) \ &= \ (x_1,x_2,u_2^{-1}u_1) \ ,
\end{aligned}
\]
we then have
\[
\begin{aligned}
\beta\alpha \quad &= \quad \id_{\mathcal{M}}\\
\alpha\beta \quad &\xrightarrow{\ \sim\ } \quad \id_{\Gamma_1} \ .
\end{aligned}
\]
On the other hand, $\Gamma_1$ is obviously accessible in $\Gamma$, as the inverse image of $\uIs(\mathcal{E}_0)$ (which % [sic] accessible dans \mathcal{M} -> accessible dans \Gamma, \Gamma_1 is a subcategory of \Gamma
\DERmark{ENV1418}

\msp{page 399}
is accessible in $\uFl(\mathcal{E}_0)$ under the accessible functor
\[
\uGamma \ \longrightarrow \ \uFl(\mathcal{E}_0)
\]
given by $(x_0, x_1, x_2, u_1, u_2) \longmapsto u_2$. This proves that $\mathcal{M}$, which is equivalent to $\Gamma_1$, is also accessible, and the accessibility of the canonical functor $\mathcal{M}_0 \longrightarrow \mathcal{E}_i$ ($i \in \{1,2\}$) follows from that of $\Gamma \longrightarrow \mathcal{E}_i$, implying that of the composite $\Gamma_1 \hookrightarrow \Gamma \longrightarrow \mathcal{E}_i$. This therefore proves b). To prove a), it suffices to see that $\mathcal{M}_0$ is accessible in $\mathcal{M}$, now it is again the inverse image of $\uIs(\mathcal{E}_0) \subset \uFl(\mathcal{E}_0)$ under the canonical accessible functor % [sic] (x_1, x_1, x_2, u_1, u_2) -> (x_0, x_1, x_2, u_1, u_2), the quintuplets begin with x_0
\DERmark{ENV1419}\DERmark{ENV1420}

\[
\mathcal{M} \ \longrightarrow \ \uFl(\mathcal{E}_0) \ ,
\]
so $\mathcal{M}_0$ is indeed accessible in $\mathcal{M}$, as the inverse image of a subcategory accessible in $\uFl(\mathcal{E}_0)$. We can moreover also interpret $\mathcal{M_0}$, up to equivalence of categories, as % [sic] donc \mathcal{M} est bien accessible -> donc \mathcal{M}_0 est bien accessible, \mathcal{M}_0 is the subcategory in question
\DERmark{ENV1421}\DERmark{ENV1422}\DERmark{ENV1423}

% [sic] \uFl(\mathcal{E}) -> \uFl(\mathcal{E}_0), as in the preceding display
% [sic] interpréter \mathcal{M} -> interpréter \mathcal{M}_0, as in the following display
\msp{page 400}
\[
\mathcal{M}_0 \ \simeq \ \uGamma_{\cocart}(\mathcal{E}/\mathcal{C}) \ ,
\]
the canonical functors $\mathcal{M}_0 \longrightarrow \mathcal{E}_i$ ($i \in \{1,2\}$) being interpreted as evaluation functors, which gives a) directly, as a special case of 4.14.19. % [sic] i \in \{0,1\} -> i \in \{1,2\}, the canonical functors of a) go to \mathcal{E}_1, \mathcal{E}_2
\DERmark{ENV1424}

% foot: Version 17 décembre 2025

\printedpage{152}
% head: A. Grothendieck : Dérivateurs, Ch. XVIII, Édition des Universités de Montpellier II et Paris VII

\textbf{4.14.22.} \textsc{Proof of} 4.14.17. The case where $p$ is cofibrant has already been proved (th. 4.10.8, p. 263). We are going to see that the case where $p$ is fibrant reduces to the preceding theorem (in a very special case), using to the full the general results on accessible parts of an accessible category. Let us note that assertions a) and b) of corollary 4.14.18 are proved in AQT \texttt{[âne qui trotte: at a donkey's trot, i.e. by routine verification]} fashion without having to appeal to 4.14.17. Part a) already implies that the evaluation functors $\varepsilon_s : \mathcal{M} \longrightarrow \mathcal{E}_s$ are accessible. It therefore remains to prove
\msp{page 401}
that $\mathcal{M} = \uGamma(\mathcal{E}/\mathcal{C})$ is accessible, in the case where $p$ is \textit{fibrant}. We shall moreover assume that $\mathcal{C}$ is small (replacing $\mathcal{C}$ by an equivalent small full subcategory if need be).

Let
\[
\begin{aligned}
\Pi \ &= \ \prod\nolimits_{s\in\Ob\mathcal{C}} \mathcal{E}_s\\
\Pi' \ &= \ \prod\nolimits_{\lambda\in\Fl\mathcal{C}} \mathcal{E}_{s(\lambda)} \ ,
\end{aligned}
\]
where $s(\lambda)$ denotes the \textit{source} $s$ of the arrow $\lambda : s \longrightarrow t$ in $\mathcal{C}$. $\Pi$, $\Pi'$ are accessible categories, as (small) products of accessible categories. (Trivial special case of 4.10.8, for a discrete index category.) We are going to define two functors
\[
\Pi \ \overset{F}{\underset{G}{\rightrightarrows}} \ \Pi'
\]
by defining for every $\lambda : s_0 \longrightarrow t_0$ in $\Fl\mathcal{C}$ two functors
\msp{page 402}
\[
\left\{
\begin{array}{lcll}
\Pi & \overset{F_\lambda}{\underset{G_\lambda}{\rightrightarrows}} & \mathcal{E}_{s(\lambda)=s_0} & \\
F_\lambda & = & \pr_{s_0} & \text{projection of } \Pi = \prod\nolimits_s \mathcal{E}_s \text{ onto the factor } \mathcal{E}_{s_0}\\
G_\lambda & = & \lambda^*\pr_{t_0} & \text{composite } \Pi \xrightarrow{\pr_{t_0}} \mathcal{E}_{t_0} \xrightarrow{\lambda^*} \mathcal{E}_{s_0}\\
 & & & \text{\texttt{[rather} } \lambda^*\pr_{t_0} \text{ \texttt{resp.} } \Pi \xrightarrow{\pr_{t_0}} \mathcal{E}_{t_0} \xrightarrow{\lambda^*} \mathcal{E}_{s_0}\text{\texttt{]}}
\end{array}
\right.
\qquad \bigl(\mbox{\footnotemark[170]}, \mbox{\footnotemark[171]}\bigr) \ .
\]
\DERmark{ENV1425}\DERmark{ENV1426}\DERmark{ENV1427}\DERmark{ENV1428}\DERmark{ENV1429}\DERmark{ENV1430}

\footnotetext[170]{It is immediate that the $F_\lambda$, $G_\lambda$, hence $F$ and $G$, are \textit{accessible} functors.}
\footnotetext[171]{\textbf{N.B.} We assume chosen for every $\lambda \in \Fl\mathcal{C}$ an inverse image functor $\lambda^* : \mathcal{E}_t \longrightarrow \mathcal{E}_s$. We assume that if $\lambda$ is an identity arrow, $\id_s$, \texttt{[we]} choose $\lambda^* = \id_{\mathcal{E}_s}$.}

Let $f \in \mathcal{M} = \uGamma(\mathcal{E}/\mathcal{C})$, then for every $\lambda \in \Fl\mathcal{C}$, $f$ defines an arrow
\[
\alpha_\lambda(f) \ : \quad \underbrace{f(s)}_{= F_\lambda(\varepsilon(f))} \quad \xrightarrow{f(\lambda)^\flat} \quad \underbrace{\lambda^*(f(t))}_{= G_\lambda(\varepsilon(f))} \ ,
\]
\textit{functorial} in $f$. We have a functor
\[
\varepsilon \ = \ (\varepsilon_s)_{s\in\Ob\mathcal{C}} \ : \ \mathcal{M} \ \longrightarrow \ \Pi \ ,
\]
and for every $\lambda \in \Fl\mathcal{C}$ a functorial homomorphism
\[
\alpha_\lambda \ : \ F_\lambda \circ \varepsilon \ \longrightarrow \ G_\lambda \circ \varepsilon \ ,
\]
whence a functorial homomorphism
\[
\alpha \ = \ (\alpha_\lambda)_{\lambda\in\Fl\mathcal{C}} \ : \ F \circ \varepsilon \ \longrightarrow \ G \circ \varepsilon \ .
\]
Let then $\mathcal{M}'$ be the following category:
% foot: Version 17 décembre 2025
\printedpage{153}
% head: A. Grothendieck : Dérivateurs, Ch. XVIII, Édition des Universités de Montpellier II et Paris VII
\msp{page 403}
\[
\mathcal{M}' \quad = \quad \parbox[t]{0.5\textwidth}{category of pairs $(X,u)$, with $X$ in $\Ob\Pi$, and $u : F(X) \longrightarrow G(X)$ an arrow in $\Pi'$.}
\]
The functorial homomorphism $\alpha : F \circ \varepsilon \longrightarrow G \circ \varepsilon$ defines a functor
\[
\mathcal{M} \ \longrightarrow \ \mathcal{M}' \ ,
\]
and it is immediate that this functor is fully faithful. In fact, the category $\mathcal{M}'$ is isomorphic to the category consisting of the pairs $(f_0,f_1)$ of two maps
\[
\begin{array}{lcccccc}
f_0 : & \Ob\mathcal{C} & \longrightarrow & \Ob\mathcal{E} & s & \longmapsto & f_0(s)\\[1ex]
f_1 : & \Fl\mathcal{C} & \longrightarrow & \Fl\mathcal{E} & \lambda & \longmapsto & f_1(\lambda)
\end{array}
\]
in such a way that $f_1$ and $f_0$ are compatible with the source and target maps, i.e. that they define a homomorphism of \textit{diagrams} $f : \diag\mathcal{C} \longrightarrow \diag\mathcal{E}$, and that this homomorphism is a right inverse of $\mathrm{diag}(p) : \diag\mathcal{E} \longrightarrow \diag\mathcal{C}$.

\msp{page 404}
The preceding functor $\mathcal{M} \hookrightarrow \mathcal{M}'$ is interpreted as the functor which associates to every section $f$ of $\mathcal{E}$ over $\mathcal{C}$, the corresponding morphism of diagrams. It is therefore the inclusion of a \textit{strictly full} subcategory of $\mathcal{M}'$. Thus to prove that $\mathcal{M}$ is accessible, it suffices to prove: % [sic] toute sections -> toute section, agreement of toute
\DERmark{ENV1431}

\begin{itemize}[leftmargin=2.5em]
\item[a)] $\mathcal{M}'$ is accessible.
\item[b)] $\mathcal{M}$ is accessible in $\mathcal{M}'$.
\end{itemize}
Let us first prove b) (admitting a). The subcategory $\mathcal{M}$ of $\mathcal{M}'$ is defined by two types of conditions, on an object $f = (f_0,f_1)$ of $\mathcal{M}'$:
\begin{itemize}[leftmargin=2.5em]
\item[1°)] $\forall\ s \in \Ob\mathcal{C}$, $f_1(\id_s) = \id_{f_0(s)}$.
\item[2°)] If $s \xrightarrow{\lambda} t \xrightarrow{\mu} u$ are two composable arrows in $\mathcal{C}$, then
\[
f_1(\mu\lambda) \ = \ f_1(\mu)f_1(\lambda) \ .
\]
\end{itemize}
These conditions are made explicit the one and \texttt{[the other by]} the equality of a double arrow in $\mathcal{E}_s$,
\msp{page 405}
namely
\[
f_0(s) \overset{f_1(\id_s)}{\underset{\id_{f_0(s)}}{\rightrightarrows}} f_0(s)
\]
in case 1°), and
\[
f_0(s) \overset{f_1(\mu\lambda)^\flat}{\underset{(f_1(\mu)f_1(\lambda))^\flat}{\rightrightarrows}} (\mu\lambda)^*(f_0(u))
\]
in case 2°). In either case, we have a functor
\[
\mathcal{M}' \ \xrightarrow{\gamma_s \text{ resp. } \gamma_{\lambda,\mu}} \ \uHom(I,\mathcal{E}_s)
\]
% foot: Version 17 décembre 2025


\endgroup
% END UNIT der-ch18/der-ch18_p151-155.tex
% BEGIN UNIT der-ch18/der-ch18_p156-160.tex
\begingroup\makeatletter%
\newcommand{\printedpage}[1]{\marginpar{\footnotesize p.~#1}}%
\newcommand{\msp}[1]{\marginpar{\footnotesize [#1]}}%
\newcommand{\uHom}{\underline{\mathrm{Hom}}}%
\newcommand{\uFl}{\underline{\mathrm{Fl}}}%
\newcommand{\uIs}{\underline{\mathrm{Is}}}%
\newcommand{\uGamma}{\underline{\Gamma}}%
\newcommand{\bDelta}{\boldsymbol{\Delta}}%
\newcommand{\Ens}{\mathrm{Ens}}%
\newcommand{\Ob}{\mathrm{Ob}\,}%
\newcommand{\Fl}{\mathrm{Fl}}%
\newcommand{\Hom}{\mathrm{Hom}}%
\makeatother\let\printedpage\derpp\let\msp\dermsp%
\let\mspin\DERmspin%


\printedpage{154}
% head: A. Grothendieck : Dérivateurs, Ch. XVIII, Édition des Universités de Montpellier II et Paris VII
in the category of double arrows of $\mathcal{E}_s$ ($I$ denoting the double-arrow category of type $0\rightrightarrows 1$), a functor which is (trivially) accessible. Let then $I\xrightarrow{\;i\;}\bDelta_1$ be the canonical functor, and
\[
\underbrace{\uHom(\bDelta_1,\mathcal{E}_s)}_{\simeq\,\uFl(\mathcal{E}_s)}\ \hookrightarrow\ \uHom(I,\mathcal{E}_s)
\]
the functor $u\longmapsto u\circ i$, which is a strictly fully faithful functor, having as its image the subcategory of double arrows $(u,v)$ such that $u=v$. As the inclusion functor is accessible and $\uHom(\bDelta_1,\mathcal{E}_s)$ and $\uHom(I,\mathcal{E}_s)$ accessible
\msp{page 406}
(a special case of 4.14.1, or also of 4.10.8!), we see that the strictly full subcategory of $\uHom(I,\mathcal{E}_s)$ of double arrows equalized in $\mathcal{E}_s$, is accessible in $\uHom(I,\mathcal{E}_s)$. So, its inverse image in $\mathcal{M}'$ is accessible. So the intersection of these inverse images for $s\in\Ob\mathcal{C}$ and $(\lambda,\mu)\in\Fl_2(\mathcal{C})$ varying, is so as well ($\mathcal{C}$ being small).

So everything comes down to proving that $\mathcal{M}'$ is accessible, which is contained in the

% vertical bar printed in the left margin alongside Lemme 4.14.22.1
\textbf{Lemma 4.14.22.1.} \textit{Let $\Pi$, $\Pi'$ be two accessible categories, $F,G:\Pi\rightrightarrows\Pi'$ two accessible functors, and $\mathcal{M}'$ the category of pairs $(X,u)$, with $X\in\Ob\Pi$ and $u:F(X)\longrightarrow G(X)$ an arrow of $\Pi'$. Then $\mathcal{M}'$ is accessible, and the canonical functor $\mathcal{M}'\longrightarrow\Pi$ is accessible.} % [sic] (X,u), X \in Ob Pi ... F(x) -> the same letter (X or x) throughout, capital X and lowercase x used for the same object
\DERmark{ENV1432}


\textsc{Proof of the Lemma.} The datum of $\Pi$, $\Pi'$, $F$, $G$ defines a cofibered category $\mathcal{F}$ over the category $I=(0\rightrightarrows 1)$
\msp{page 407}
above. The category $\uGamma(\mathcal{F}/I)$ is interpreted, up to isomorphism, as the category of triples $(x,u,v)$ with $x\in\Ob\Pi$, $x'\in\Ob\Pi'$, and $u$ and $v$ arrows
\[
\begin{tikzcd}[column sep=small]
F(x) \arrow[dr, "u"'] & & G(x) \arrow[dl, "v"] \\
& x' &
\end{tikzcd}
\]
in $\Pi'$. The category $\mathcal{M}'$ is equivalent to the strictly full subcategory of $\uGamma(\mathcal{F}/I)$, consisting of triples $(x,u,v)$ such that $v$ is an isomorphism, i.e. to the subcategory of $\uGamma(\mathcal{F}/I)$, the inverse image of $\uIs(\Pi')$ under the functor $(x,u,v)\longmapsto v$,
\[
\uGamma(\mathcal{F}/I)\ \longrightarrow\ \uFl(\Pi')\ .
\]
Now by 4.10.8, $\uGamma(\mathcal{F}/I)$ is accessible, and the preceding functor is of course accessible. Finally, by 4.14.12,
\msp{page 408}
$\uIs(\Pi')$ is accessible in the accessible category $\uFl(\Pi')$. So its inverse image $\mathcal{M}'$ in $\uGamma(\mathcal{F}/I)$ is accessible in $\uGamma(\mathcal{F}/I)$, hence accessible. So the same holds for the equivalent category, q.e.d.

Thus, thanks to the general results on accessible parts of accessible categories, theorem 4.14.17, in the case where $p$ is fibrant, is reduced to the cofibrant case, in the very special case where the index category is $I=(0\rightrightarrows 1)$. One could moreover make a similar reduction in the case where $p$ is cofibrant, by a proof practically % [sic] aux cas cofibrant -> au cas cofibrant, singular noun
\DERmark{ENV1433}

% foot: Version 17 décembre 2025
\printedpage{155}
% head: A. Grothendieck : Dérivateurs, Ch. XVIII, Édition des Universités de Montpellier II et Paris VII
identical. In other words, theorem 4.10.8 reduces to the case where the index category is $I=(0\rightrightarrows 1)$. But I suspect that the proof of the theorem in this special case, is not essentially simpler than in the general case.

\msp{page 409}

\textbf{4.14.23.} One may ask about finding a statement ensuring the accessibility of a category $\uGamma(\mathcal{E}/\mathcal{C})$, which covers the two cases considered in 4.14.17. Here is what I can propose. I am going to make the following assumption on the functors $p:\mathcal{E}\longrightarrow\mathcal{C}$ (in addition to that of $\mathcal{C}$ being essentially small): % [sic] les foncteur -> le foncteur, singular functor p

\begin{itemize}[leftmargin=*]
\item[H)] For every $\lambda\in\Fl\mathcal{C}$, the category $\mathcal{E}_\lambda=\mathcal{E}\times_{\mathcal{C}}(\bDelta_1,i_\lambda)$, where $i_\lambda:\bDelta_1\longrightarrow\mathcal{C}$ is the functor defined by $\lambda$, is accessible, and the canonical functor $\mathcal{E}_\lambda\xrightarrow{\;p_\lambda\;}\bDelta_1$ is accessible.
\end{itemize}

Note that if $\lambda:s\longrightarrow t$ in $\mathcal{C}$ is such that $\mathcal{E}_\lambda$ and $p_\lambda$ are accessible, then $\mathcal{E}_s$ and $\mathcal{E}_t$ are accessible, since isomorphic to the inverse image categories under $p_\lambda:\mathcal{E}_\lambda\longrightarrow\bDelta_1$ of the subcategories $\{0\}$ resp. $\{1\}$ of $\Delta_1$, which are accessible in the accessible category $\bDelta_1$, so % [sic] sont accessible -> sont accessibles, plural agreement
\msp{page 410}
$\mathcal{E}_{\lambda_0}$ and $\mathcal{E}_{\lambda_1}$ are accessible in $\mathcal{E}_\lambda$, hence are accessible, hence also $\mathcal{E}_s$ and $\mathcal{E}_t$. Thus, assumption H) implies that the fiber categories are accessible. Note moreover:

\textbf{Lemma 4.14.23.1.} \textit{Suppose that $p:\mathcal{E}\longrightarrow\mathcal{C}$ is cofibrant or fibrant, and let $(\lambda:s\longrightarrow t)\in\Fl(\mathcal{C})$. For $\mathcal{E}_\lambda$ to be accessible and $\mathcal{E}_\lambda\xrightarrow{\;p_\lambda\;}\bDelta_1$ accessible, it is necessary and sufficient that $\mathcal{E}_s$, $\mathcal{E}_t$ be accessible and $\lambda_*:\mathcal{E}_s\longrightarrow\mathcal{E}_t$ resp. $\lambda^*:\mathcal{E}_t\longrightarrow\mathcal{E}_s$ be an accessible functor.}

\textsc{Proof.} We may suppose that $\mathcal{C}=\bDelta_1$, and that $\mathcal{E}$ is the cofibered (resp. fibered) category over $\bDelta_1$ defined by a functor $\varphi:\mathcal{E}_0\longrightarrow\mathcal{E}_1$ (resp. $\psi:\mathcal{E}_1\longrightarrow\mathcal{E}_0$), and to show that $\mathcal{E}$ and $p:\mathcal{E}\longrightarrow\bDelta_1$ are accessible if and only if $\mathcal{E}_0$ and $\mathcal{E}_1$ and $\varphi$ (resp. $\psi$) are.

\msp{page 411}

Let us prove the ``only if''. We have already seen that if $\mathcal{E}$, $p$ are accessible, $\mathcal{E}_0$ and $\mathcal{E}_1$ are accessible in $\mathcal{E}$, hence accessible, and the inclusion functors % [sic] dans E_s -> dans E, here C = Delta_1 and the subcategories lie in E
\DERmark{ENV1434}

\[
\mathcal{E}_0\ \hookrightarrow\ \mathcal{E}\ \hookleftarrow\ \mathcal{E}_1
\]
accessible. To say that $p$ is cofibrant (resp. fibrant) means that the inclusion functor $\mathcal{E}_1\hookrightarrow\mathcal{E}$ admits a left adjoint $\bar{\varphi}$ (resp. that the functor $\mathcal{E}_0\hookrightarrow\mathcal{E}$ admits a right adjoint $\tilde{\psi}$). Then $\varphi$ (resp. $\psi$) is the composite $\mathcal{E}_0\longrightarrow\mathcal{E}\xrightarrow{\;\bar{\varphi}\;}\mathcal{E}_1$ (resp. the composite $\mathcal{E}_1\hookrightarrow\mathcal{E}\xrightarrow{\;\tilde{\psi}\;}\mathcal{E}_0$). So to prove that $\varphi$ (resp. $\psi$) is accessible, it suffices to prove that $\bar{\varphi}$ (resp. $\tilde{\psi}$) is. This is a special case of the % [sic] E -> E_1 (target of psi-tilde composite) -> E_0, psi goes from E_1 to E_0
\DERmark{ENV1435}\DERmark{ENV1436}\DERmark{ENV1437}


\textbf{Lemma 4.14.23.2.} \textit{Let $\mathcal{E}$, $\mathcal{E}'$ be two accessible categories} (\footnote[172]{\textbf{N.B.} It suffices that $\mathcal{E}$ admit a small generating subcategory, and that the objects of $\mathcal{E}'$ be accessible.}), \textit{and $\mathcal{E}\underset{G}{\overset{F}{\rightleftarrows}}\mathcal{E}'$ a pair of adjoint functors between $\mathcal{E}$ and $\mathcal{E}'$. Then $F$ and $G$ are accessible.}

\msp{page 412}

Indeed, $F$ is accessible, since it commutes with arbitrary small $\varinjlim$, a fortiori it satisfies $\mathrm{L}_{\pi_0}$ (where $\pi_0$ is such that $\mathcal{E}$, $\mathcal{E}'$ satisfy $\mathrm{L}_{\pi_0}$). Let us prove that $G$ is accessible. If $S$ is a
% foot: Version 17 décembre 2025
\printedpage{156}
% head: A. Grothendieck : Dérivateurs, Ch. XVIII, Édition des Universités de Montpellier II et Paris VII
small generating set of elements in $\mathcal{E}$, so that the small family of functors
\[
h_\xi\ =\ (x\longmapsto\Hom_{\mathcal{E}}(\xi,x))\ :\ \mathcal{E}\ \longrightarrow\ \Ens
\]
is conservative, to prove that $G$ is accessible, it suffices to prove that the $h_\xi\circ G$ ($\xi\in S\subset\Ob\mathcal{E}$) are. Now
\[
h_\xi\circ G(x')\ =\ \Hom(\xi,G(x'))\ \simeq\ \Hom_{\mathcal{E}'}(F(\xi),x')\ , % [sic] L_xi -> h_xi, the functors were named h_xi
\]
\DERmark{ENV1438}

and since $F(\xi)$ is an accessible object of $\mathcal{E}'$, the functor
\[
x'\ \longmapsto\ (h_\xi\circ G)(x')\ \simeq\ \Hom_{\mathcal{E}'}(F(\xi),x') % [sic] L_xi -> h_xi, the functors were named h_xi
\]
\DERmark{ENV1439}

is indeed accessible, q.e.d.

This proves the ``only if'' in 4.14.23.1. Let us prove the ``if'', i.e. suppose $\mathcal{E}_s$, $\mathcal{E}_t$ accessible, and $\varphi$ (resp. $\psi$) accessible.

\msp{page 413}

In both cases, $\mathcal{E}_0$ is an open of $\mathcal{E}$, $\mathcal{E}_1$ is the complementary closed, and $\mathcal{E}$ is the category ``glued'' with the help of the bifunctor
\[
H\ :\ \mathcal{E}_0^{\circ}\times\mathcal{E}_1\ \longrightarrow\ \Ens
\]
given by
\[
\left\{
\begin{aligned}
H(x_0,x_1)\ &=\ \Hom(\varphi(x_0),x_1)\\
&\hspace{-2em}\text{resp.}\\
H(x_0,x_1)\ &=\ \Hom(x_0,\psi(x_1))\ ,
\end{aligned}
\right.
\]
according as we are in the case $p$ cofibrant or $p$ fibrant. We are going to use the

\textbf{Proposition 4.14.24} (\footnote[173]{cf. more general statements and nice proof, 4.15.1 and 4.15.4 (p. 425, 432).}). \textit{Let $\mathcal{E}$ be a category, $\mathcal{E}_0$ an open subcategory, $\mathcal{E}_1$ the complementary closed subcategory. Equivalent conditions:}
\begin{itemize}[leftmargin=*]
\item[\textit{a)}] \textit{$\mathcal{E}$ is accessible, and $\mathcal{E}\longrightarrow\bDelta_1$ (with fibers $\mathcal{E}_0$, $\mathcal{E}_1$) is accessible.}
\item[\textit{b)}] \textit{$\mathcal{E}$ is accessible, and $\mathcal{E}_0$, $\mathcal{E}_1$ are accessible in $\mathcal{E}$.}
\item[\textit{c)}] \textit{$\mathcal{E}_0$, $\mathcal{E}_1$ are accessible, and for every $x_0\in\Ob\mathcal{E}_0$, the functor}
\[
x_1\ \longmapsto\ \uHom_{\mathcal{E}}(x_0,x_1)\ :\ \mathcal{E}_1\ \longrightarrow\ \Ens
\]
\msp{page 414}
\textit{is accessible, and there exists $\pi$ such that for every $x_1\in\Ob\mathcal{E}_1$, the functor}
\[
x_0\ \longmapsto\ \uHom_{\mathcal{E}}(x_0,x_1)^{\circ}\ :\ \mathcal{E}_0\ \longrightarrow\ (\Ens)^{\circ}
\]
\textit{is $\pi$-accessible.}
\end{itemize}
% foot: Version 17 décembre 2025

\printedpage{157}
% head: A. Grothendieck : Dérivateurs, Ch. XVIII, Édition des Universités de Montpellier II et Paris VII
Let us admit this lemma \texttt{[rather this proposition]}, let us finish the proof of 4.14.23.1. We suppose $\mathcal{E}_0$, $\mathcal{E}_1$ and $\varphi$ (resp. $\psi$) accessible, we want to prove condition a) of 4.14.24. For this, it suffices to prove condition c). Now we have, in the case $p$ cofibrant, % [sic] 4.14.23.3 -> 4.14.24, the conditions a), c) are those of Proposition 4.14.24
\DERmark{ENV1440}

\[
\Hom_{\mathcal{E}}(x_0,x_1)\ \simeq\ \Hom_{\mathcal{E}_1}(\varphi(x_0),x_1)\ ,
\]
this is indeed accessible in $x_1$, since $\mathcal{E}_1$ is accessible, and $\pi$-accessible in $x_0$ if $\varphi$ is, as a composite of $\varphi$ with $\mathcal{E}_1\longrightarrow\Ens^{\circ}$, $\Hom(-,x_1)$, which is. In the case $p$ fibrant, we have
\DERmark{ENV1441}

\[
\Hom_{\mathcal{E}}(x_0,x_1)\ \simeq\ \Hom_{\mathcal{E}_0}(x_0,\psi(x_1))\ ,
\]
this is accessible in $x_1$, as the composite of the accessible functor $\psi:\mathcal{E}_1\longrightarrow\mathcal{E}_0$, and the likewise accessible functor $\Hom(x_0,-):\mathcal{E}_0\longrightarrow\Ens$, which is $\pi$-accessible in $x_0$ tautologically.

To finish the proof of 4.14.23.1, it therefore remains to prove proposition 4.14.24. We have already noted that we have a $\Longrightarrow$ b, and b $\Longrightarrow$ c clear, since the functor considered $\mathcal{E}_1\longrightarrow\Ens$ is the composite $\mathcal{E}_1\overset{\text{incl.}}{\hookrightarrow}\mathcal{E}\xrightarrow{\;\Hom_{\mathcal{E}}(x_0,-)\;}\Ens$, is accessible as a composite of two
\DERmark{ENV1442}

\msp{page 415}
accessible functors. So it remains to prove c $\Longrightarrow$ a. Let then $\pi_0$ be a useful cardinal for $\mathcal{E}_0$, and such that the functors
\[
x_0\ \longmapsto\ H(x_0,x_1)\ :\ \mathcal{E}_0\longrightarrow\Ens^{\circ}
\]
are $\pi_0$-accessible, let $\pi$ be a cardinal $\geqslant\sup(\pi_0,\mathrm{card}\,\text{red}\,\mathcal{E}_{0\,\pi_0})$ such that $\pi^{\pi_0}=\pi$ (hence useful for $\mathcal{E}_0$), useful for $\mathcal{E}_1$, and such that the functors
\[
x_1\ \longmapsto\ H(x_0,x_1)\ :\ \mathcal{E}_1\longrightarrow\Ens
\]
are $\pi$-accessible, for $x_0\in\Ob\mathcal{E}_{0\,\pi_0}$. Let $\mathcal{E}(\pi)$ be the strictly full subcategory of $\mathcal{E}$, such that
\[
\Ob\mathcal{E}(\pi)\ =\ \Ob\mathcal{E}_{0\,\pi}\cup\Ob\mathcal{E}_{1\,\pi}\ .
\]
I say that
\begin{itemize}[leftmargin=*]
\item[1°)] $\mathcal{E}$ satisfies $\mathrm{L}_\pi$, and $\mathcal{E}_0$ and $\mathcal{E}_1$ are stable in $\mathcal{E}$ under $\mathrm{L}_\pi$, the functors $\mathcal{E}_0\longrightarrow\mathcal{E}$, $\mathcal{E}_1\longrightarrow\mathcal{E}$, $\mathcal{E}\longrightarrow\bDelta_1$ satisfy $\mathrm{L}_\pi$ (\footnote[174]{it suffices that $\mathcal{E}_0$ and $\mathcal{E}_1$ satisfy $\mathrm{L}_\pi$, and that
\[
x_0\ \longmapsto\ H(x_0,x_1)\ :\ \mathcal{E}_0\longrightarrow\Ens^{\circ}
\]
satisfy $\mathrm{L}_\pi$ $\forall\ x\in\Ob\mathcal{E}$.}).
\item[2°)] $\mathcal{E}(\pi)=\mathcal{E}_\pi$ (whence a fully faithful functor
\[
\mathrm{Ind}_\pi(\mathcal{E}(\pi))\ \longrightarrow\ \mathcal{E}\ \ ).
\]
\item[3°)] The said functor is essentially surjective,
\msp{page 416}
hence an equivalence of categories.
\end{itemize}
Note that 2°) and 3°) imply that $\mathcal{E}$ is accessible, and 1°) that $\mathcal{E}\longrightarrow\bDelta_1$ is accessible, whence condition a) of the lemma we are to establish. So everything comes down to proving 1°, 2°, 3° above.
% foot: Version 17 décembre 2025

\printedpage{158}
% head: A. Grothendieck : Dérivateurs, Ch. XVIII, Édition des Universités de Montpellier II et Paris VII
1°) Since $\mathcal{E}_1$ is closed in $\mathcal{E}$, it is clear that $\mathcal{E}_1$ is stable under all $\varinjlim$ representable in $\mathcal{E}$, and more precisely, that for a functor $I\longrightarrow\mathcal{E}_1$, its direct limit in $\mathcal{E}_1$ exists if and only if it exists in $\mathcal{E}$, and the two are isomorphic. Suppose $I$ large relative to $\pi$, and let $\varphi:I\longrightarrow\mathcal{E}$, $i\longmapsto x_i$. Let us distinguish two cases.
\DERmark{ENV1443}


a) $\varphi$ takes its values in $\mathcal{E}_0$. Let $x_0$ be its direct limit \textit{in} $\mathcal{E}_0$. I say that this is also a direct limit in $\mathcal{E}$. We must prove that for $x_1\in\Ob\mathcal{E}_1$, the arrow
\[
\Hom_{\mathcal{E}}(x_0,x_1)\ \longrightarrow\ \varprojlim_{i\in I}\Hom_{\mathcal{E}}(x_i,x_1)
\]
is isomorphic. Now this is the case
\msp{page 417}
by the assumption on $\pi$,
\[
\xi\ \longmapsto H(\xi,x_1)\ :\ \mathcal{E}_0\ \longrightarrow\ (\Ens)^{\circ}
\]
satisfies $\mathrm{L}_\pi$ for every $x_1$ in $\Ob\mathcal{E}_1$.

b) $\varphi$ does not factor through $\mathcal{E}_0$, i.e. $\exists\ i_0\in\Ob I$ with $\varphi(i_0)\in\Ob\mathcal{E}_1$. Then we have $\varphi(i)\in\Ob\mathcal{E}_1$ for every $i\in\Ob I[i_0]$, where $I[i_0]$ denotes the strictly full subcategory of $I$ consisting of the $i$ such that $i\geqslant i_0$. This category is cofinal in $I$, hence large relative to $\pi$ like $I$, and for $\varinjlim\limits_{I}x_i$ to exist, it is necessary and sufficient that $\varinjlim\limits_{I[i_0]}x_i$ exist, which is the case by the cases a) already treated. % [sic] les cas a) déjà traité -> le cas a) déjà traité, singular agreement

This therefore proves at the same time that $\mathcal{E}$ satisfies $\mathrm{L}_\pi$, and that $\mathcal{E}_0$ and $\mathcal{E}_1$ are $\mathrm{L}_\pi$-stable. The proof shows at the same time that the functor $p:\mathcal{E}\longrightarrow\bDelta_1$ satisfies $\mathrm{L}_\pi$.

2°) From this it follows at once that $\mathcal{E}_\pi\subset\mathcal{E}(\pi)$. Let us prove that we also have
\[
\mathcal{E}(\pi)\ \subset\ \mathcal{E}_\pi\ ,
\]
and let us again distinguish two cases, for
\msp{page 418}
$x\in\Ob \mathcal{E}(\pi)$, according as $x$ is in $\mathcal{E}_{0\,\pi}$ or $\mathcal{E}_{1\,\pi}$. If it is in $\mathcal{E}_{1\,\pi}$, to prove that % [sic] E(pi) -> \mathcal{E}(pi), script letter elsewhere
\DERmark{ENV1444}\DERmark{ENV1445}

\[
\Hom_{\mathcal{E}}(x,\underbrace{\varinjlim y_i}_{y})\ \xleftarrow{\ \sim\ }\ \varinjlim_{i}\Hom_{\mathcal{E}}(x,y_i)
\tag*{$(*)$}
\]
(for a direct system $(y_i)_{i\in I}$ in $\mathcal{E}$ large relative to $\pi$), if the $y_i$ are in $\mathcal{E}_0$, hence $y\in\mathcal{E}_0$, then both sides are empty, and this is OK. Otherwise, we may suppose (replacing $I$ by an $I[i_0]$ if need be) that the $y_i$ are in $\mathcal{E}_1$, and then this is OK too, since $y$ is also the $\varinjlim$ in $\mathcal{E}_1$, by 1°). % [sic] les x_i -> les y_i, the inductive system is (y_i)
\DERmark{ENV1446}


If $x$ is in $\mathcal{E}_{0\,\pi}$, then $(*)$ is clear if the $y_i$ are in $\mathcal{E}_0$. Otherwise, we may again suppose that the $y_i$ are in $\mathcal{E}_1$, and $(*)$ reads % [sic] y_i sont dans E_0 (second occurrence) -> E_1, the case "sinon" and the following "les y_i dans E_1"
\DERmark{ENV1447}

\[
H(x,\varinjlim y_i)\ \xleftarrow{\ \sim\ }\ \varinjlim H(x,y_i)\ ,
\]
for $x\in\Ob\mathcal{E}_{0\,\pi}$, the $y_i$ in $\mathcal{E}_1$, $I$ large relative to $\pi$. In other words, we must prove that the functor
\[
x_1\ \longmapsto\ H(x_0,x_1)\ :\ \mathcal{E}_1\ \longrightarrow\ \Ens
\]
% foot: Version 17 décembre 2025


\endgroup
% END UNIT der-ch18/der-ch18_p156-160.tex
% BEGIN UNIT der-ch18/der-ch18_p161-165.tex
\begingroup\makeatletter%
\derleqno%
\newcommand{\printedpage}[1]{\marginpar{\footnotesize p.~#1}}%
\newcommand{\msp}[1]{\marginpar{\footnotesize [#1]}}%
\newcommand{\cE}{\mathcal{E}}%
\newcommand{\cC}{\mathcal{C}}%
\newcommand{\cM}{\mathcal{M}}%
\newcommand{\Ob}{\operatorname{Ob}}%
\newcommand{\Fl}{\operatorname{Fl}}%
\newcommand{\Hom}{\operatorname{Hom}}%
\newcommand{\End}{\operatorname{End}}%
\newcommand{\Ind}{\operatorname{Ind}}%
\newcommand{\Ens}{\mathrm{Ens}}%
\newcommand{\pr}{\mathrm{pr}}%
\newcommand{\id}{\mathrm{id}}%
\newcommand{\card}{\operatorname{card}}%
\newcommand{\cardred}{\operatorname{card}\,\text{\upshape red}}%
\newcommand{\defeq}{\overset{\mathrm{def}}{=}}%
\makeatother\let\printedpage\derpp\let\msp\dermsp%
\let\mspin\DERmspin%

\DERmark{ENV1448}


\setcounter{footnote}{174}

%=====================================================================
\printedpage{159}
% head: A. Grothendieck : Dérivateurs, Ch. XVIII, Édition des Universités de Montpellier II et Paris VII

\msp{page 419}
is $\pi$-accessible, when $x_0$ is in $\cE_{0\,\pi}$. Now by the hypothesis on $\pi$, it is $\pi$-accessible if $x_0$ is in $\cE_{0\,\pi_0}$. But if $x_0$ is in $\cE_{0\,\pi}$, we will have % [sic] x -> x_0, the formula that follows concerns x_0
\DERmark{ENV1449}

\[
x_0 \;=\; \varinjlim_{i\in I} x_{0\,i} \qquad
\left\{
\begin{array}{l}
I\ \ \text{large}\ \ \text{relative to}\ \ \pi_0,\\
\text{the}\ \ x_{0\,i}\ \ \text{in}\ \ \cE_{0\,\pi_0},\\
\card I \leqslant \pi
\end{array}
\right.
\]
(for $\pi_0$ is useful for $\cE_0$ by hypothesis, and $\pi$ satisfies $\pi^{\pi_0} = \pi \geqslant \cardred \cE_{0\,\pi_0}$). As $x_0 \longmapsto H(x_0,x_1) : \cE_0 \longrightarrow \Ens^{\circ}$ satisfies $\mathrm{L}_{\pi_0}$, for every $x_1$ in $\cE_1$, we will therefore have % [sic] x_0 \longrightarrow H(x_0,x_1) -> x_0 \longmapsto H(x_0,x_1), elsewhere such functors are written with \longmapsto
\DERmark{ENV1450}\DERmark{ENV1451}

\[
H(x_0) \;\simeq\; \varprojlim_{i\in I} H(x_{0\,i})
\]
(where for $x_0$ in $\cE_0$, $H(x_0)$ denotes the functor $x_1 \longmapsto H(x_0,x_1) : \cE_1 \longrightarrow (\Ens)$). As % [sic] H(x) -> H(x_0), the functor depends on x_0
\DERmark{ENV1452}

the functors $H(x_{0\,i})$ are $\pi$-accessible, and $\card I \leqslant \pi$, it follows that their $\varprojlim\limits_{I}$ $H(x_0)$ is $\pi$-accessible (SGA 4 I 9.6 (ii)). This therefore completes the proof of 2°).

3°). To prove that $\Ind_\pi(\cE(\pi)) \longrightarrow \cE$ is essentially surjective, it suffices to see that every $x \in \Ob\cE_0$ and every $x \in \Ob\cE_1$
\msp{page 420}
is in the essential image. Now this follows from the hypothesis that $\pi$ is useful for $\cE_0$ and $\cE_1$, and from 1°) which implies that $\cE_0 \longrightarrow \cE$ and $\cE_1 \longrightarrow \cE$ satisfy $\mathrm{L}_\pi$.

This completes the proof of proposition 4.14.24, and likewise the following result, more precise than the implication c) $\Longrightarrow$ a).

\textbf{Corollary 4.14.25.} \textit{Let $\cE$ be a category, $\cE_0$ an open subcategory, $\cE_1$ the complementary closed subcategory,}
\[
H \;=\; ((x_0,x_1) \longmapsto \Hom_{\cE}(x_0,x_1)) \;:\; \cE_0^\circ \times \cE_1 \;\longrightarrow\; \Ens
\]
\textit{the gluing functor.}
\begin{enumerate}
\item[\textit{1°)}] \textit{Let $\pi$ be an infinite cardinal. In order that $\cE$ and $\cE_0$ satisfy $L_\pi$, and the inclusion functor $\cE_0 \hookrightarrow \cE$ likewise, it is necessary and sufficient that $\cE_0$ and $\cE_1$ satisfy $L_\pi$, and that for every $x_1$ in $\cE_1$, the functor $x_0 \longmapsto H(x_0,x_1) : \cE_0 \longrightarrow \Ens^\circ$ satisfy $L_\pi$. Then $\cE_1$ and $\cE_1 \hookrightarrow \cE$ also satisfy $L_\pi$,}
\msp{page 421}
\textit{and we have}
\DERmark{ENV1453}

\[
\cE_\pi \;\subset\; \cE(\pi) \;\defeq\;
\begin{array}{l}
\text{\itshape strictly\ \ full\ \ subcategory\ \ of\ \ $\cE$\ \ such\ \ that}\\
\Ob\cE(\pi) = \Ob\cE_{0\,\pi} \cup \Ob\cE_{1\,\pi}.
\end{array}
\]
\item[\textit{2°)}] \textit{Suppose that condition c) of 4.14.24 is satisfied, and let $\pi_0$ be an infinite cardinal such that $\pi_0$ is useful for $\cE_0$, and that the functors $x_0 \longmapsto H(x_0,x_1) : \cE_0 \longrightarrow \Ens^\circ$ defined by the $x_1 \in \Ob\cE_1$ are $\pi_0$-accessible. Let $\pi$ be a cardinal such that $\pi^{\pi_0} = \pi \geqslant \sup(\pi_0, \cardred\cE_{0\,\pi_0})$, and such that for every $x_0 \in \Ob\cE_{0\,\pi_0}$, the functor $x_1 \longmapsto H(x_0,x_1) : \cE_1 \longrightarrow \Ens$ is $\pi$-accessible. Then we have}
\[
\cE_\pi \;=\; \cE(\pi)\ .
\]
\DERmark{ENV1454}

% foot: Version 17 décembre 2025

%=====================================================================
\printedpage{160}
% head: A. Grothendieck : Dérivateurs, Ch. XVIII, Édition des Universités de Montpellier II et Paris VII

\item[\textit{3°)}] \textit{If $\pi_0$ and $\pi$ are as in 2°), and if moreover $\pi$ is useful for $\cE_1$, then $\pi$ is useful for $\cE$, hence (by virtue of 2°)}
\[
\cE \;\xleftarrow{\ \sim\ }\; \Ind_\pi(\cE(\pi))\ .
\]
\end{enumerate}
We can now state the sought-for generalization of 4.14.17:

\msp{page 422}
\textbf{Theorem 4.14.26} (generalization of 4.14.17). \textit{Let $p : \cE \longrightarrow \cC$ be a functor, with $\cC$ small. We assume that $p$ satisfies hypothesis H) of page 409, which is equivalent (using 4.14.24) to the following:}
\DERmark{ENV1455}

\begin{enumerate}
\item[\textit{H$'$)}] \textit{$\forall\ s \in \Ob\cC$, the fiber category $\cE_s$ is accessible, and for every $\lambda : s \longrightarrow t$ in $\Fl\cC$, the functor}
\[
\left\{
\begin{array}{lcl}
H_\lambda : \cE_s^\circ \times \cE_t & \longrightarrow & \Ens\\[1ex]
\qquad (x_0,x_1) & \longmapsto & \Hom_\lambda(x_0,x_1)
\end{array}
\right.
\]
\textit{satisfies the following two conditions:}
\begin{enumerate}
\item[\textit{a)}] \textit{$\forall\ x_0 \in \Ob\cE_s^\circ$, $x_1 \longmapsto H_\lambda(x_0,x_1) : \cE_t \longrightarrow \Ens$ is accessible.}
\item[\textit{b)}] \textit{There exists an infinite cardinal $\pi_\lambda$, such that for every $x_1 \in \Ob\cE_t$, the functor $x_0 \longmapsto H_\lambda(x_0,x_1) : \cE_s \longrightarrow (\Ens)^\circ$ is $\pi_\lambda$-accessible.}
\end{enumerate}
\end{enumerate}
\textit{This being so, the category $\underline{\Gamma}(\cE/\cC)$ of sections of $\cE$ over $\cC$ is accessible, and the functors $\varepsilon_s : \underline{\Gamma}(\cE/\cC) \longrightarrow \cE_s$ are accessible. We can be more precise in a way similar to 4.14.18 (statement in due form left to the reader).}

\msp{page 423}
\textsc{Proof} of 4.14.26. It is carried out exactly like that of 4.14.17 (p. 400 ff.), but is conceptually simpler. We assume $\cC$ small, and introduce the categories
\[
\Pi \;=\; \prod_{s\,\in\,\Ob\cC} \cE_s\ , \qquad \Pi' \;=\; \prod_{\lambda\,\in\,\Fl\cC} \cE_\lambda\ ,
\]
where
\[
\cE_\lambda \;=\; \cE \times_{\cC} (\boldsymbol{\Delta}_1, i_\lambda)\ ,
\]
and the two canonical functors
\[
\Pi \;\underset{G}{\overset{F}{\rightrightarrows}}\; \Pi'
\]
with components
\[
\begin{aligned}
F_\lambda \;&=\; i_\lambda \circ \pr_{s(\lambda)}\\
G_\lambda \;&=\; j_\lambda \circ \pr_{b(\lambda)}\ ,
\end{aligned}
\]
where for $\lambda : s \longrightarrow t$, we set
\[
s(\lambda) \;=\; s\ , \qquad b(\lambda) \;=\; t\ ,
\]
and $i_\lambda$, $j_\lambda$ are the canonical inclusions
% In the print, the vertical arrow p_lambda and Delta_1 are enclosed together in a pair of large parentheses.
\[
\begin{tikzcd}
\cE_s \arrow[r,"i_\lambda"] \arrow[d,dash,"\wr"] & \cE_\lambda \arrow[dd,"p_\lambda"] & \cE_t\ . \arrow[l,"j_\lambda"'] \arrow[d,dash,"\wr"] \\
(\cE_\lambda)_0 & & (\cE_\lambda)_1 \\
& \boldsymbol{\Delta}_1 &
\end{tikzcd}
\]
% foot: Version 17 décembre 2025

%=====================================================================
\printedpage{161}
% head: A. Grothendieck : Dérivateurs, Ch. XVIII, Édition des Universités de Montpellier II et Paris VII

\msp{page 424}
We denote by $\cM'$ the category
\[
\cM' \;=\;
\begin{array}{l}
\text{category of pairs } (X,u)\text{, with}\\
X \in \Ob\Pi,\ u : F(X) \longrightarrow G(X),
\end{array}
\] % [sic] F(x), G(x) -> F(X), G(X), the object is named X
\DERmark{ENV1456}

which is constructible by 4.14.22.1 is accessible (for $\Pi$ and $\Pi'$ and the functors $F$ and $G$ are, by hypothesis H) made on $p$). We can interpret it as the category of sections % [sic] \pi' -> \Pi', the category introduced is \Pi'
\DERmark{ENV1457}

of $\mathrm{Diag}(\cE)$ over $\mathrm{Diag}(\cC)$. We have a strictly fully faithful functor
\[
\cM \;\defeq\; \underline{\Gamma}(\cE/\cC) \;\hookrightarrow\; \cM'\ ,
\]
which lets us interpret $\cM$ as a strictly full subcategory of $\cM'$. We then see as in the case of 4.14.17 that this subcategory is accessible in $\cM'$ -- this is essentially sorital. It follows that $\cM$ is accessible and that the functors $\cM \longrightarrow \cE_s$ are also, as composites of
\[
\cM \;\overset{\text{\upshape acc.}}{\hookrightarrow}\; \cM' \;\overset{\text{\upshape acc.}}{\longrightarrow}\; \Pi \;\xrightarrow{\ \ \pr_s\ \ \text{\upshape acc.}\ \ }\; \cE_s\ ,
\]
\DERmark{ENV1458}\DERmark{ENV1459}


\hfill q.e.d.

\msp{page 425}

\subsection*{4.15\quad Accessible category over another}

\textbf{Proposition 4.15.1.} \textit{Let $f : \cE \longrightarrow \cC$ be a functor with $\cC$ small. We assume in addition that condition H (p. 409) is satisfied, which we make explicit here in the equivalent form H$'$:}
\begin{enumerate}
\item[\textit{H 1)}] \textit{For every $s \in \Ob\cC$, $\cE_s$ is accessible.}
\item[\textit{H 2)}] \textit{For every $\lambda \in \Fl\cC$, $\lambda : s \longrightarrow t$, there exists an infinite cardinal $\pi_\lambda$ such that for every $x_1 \in \Ob\cE_t$, the functor}
\[
x_0 \;\longmapsto\; \Hom_\lambda(x_0,x_1) \;:\; \cE_s \;\longrightarrow\; (\Ens)^\circ
\]
\textit{is $\pi_\lambda$-accessible ($\cE_s$ satisfying in addition $L_{\pi_\lambda}$).}
\item[\textit{H 3)}] \textit{For every $\lambda \in \Fl\cC$, $\lambda : s \longrightarrow t$, and for every $x_0 \in \Ob\cE_s$, the functor}
\[
x_1 \;\longmapsto\; \Hom_\lambda(x_0,x_1) \;:\; \cE_t \;\longrightarrow\; (\Ens)
\]
\textit{is accessible.}
\end{enumerate}
\textit{Then the inclusion functors $i_s : \cE_s \longrightarrow \cE$ are accessible. Moreover, if $\cC$ ``has no projectors'' (i.e. if every endomorphism $p$ of an object $s$ of $\cC$ such that $p^2 = p$, is the identity), then $\cE$ is accessible, and $f : \cE \longrightarrow \cC$ is accessible.}

This follows from the more precise result:
% foot: Version 17 décembre 2025

%=====================================================================
\printedpage{162}
% head: A. Grothendieck : Dérivateurs, Ch. XVIII, Édition des Universités de Montpellier II et Paris VII

\msp{page 426}
\textbf{Corollary 4.15.2} (\footnote{For a) and b) we use only H 1 and H 2, to the exclusion of H 3.}). \textit{The hypotheses on $f : \cE \longrightarrow \cC$ are those of 4.15.1 ($\cC$ small, and $f$ satisfies $H' =$ H 1) 2) 3)).}
\begin{enumerate}
\item[\textit{a)}] \textit{Let $s \in \Ob\cC$ and let $\pi_1$ be an infinite cardinal such that for every $\lambda : s \longrightarrow t$ (arrow of $\cC$ with source $s$), we have $\pi_1 \geqslant \pi_\lambda$ (where $\pi_\lambda$ is defined in the statement of H 2). Then the functor $i_s : \cE_s \longrightarrow \cE$ is $\pi_1$-accessible.} % [sic] source t -> source s, lambda : s -> t has source s (cf. p. 163 "de source s")
\item[\textit{b)}] \textit{We assume $\cC$ has no projectors. Let $\pi_1$ be a cardinal $\geqslant \sup(\pi_\lambda, \card\cC)$. Then $\cE$ satisfies $L_{\pi_1}$, and $f : \cE \longrightarrow \cC$ is $\pi_1$-accessible.}
\item[\textit{c)}] \textit{Suppose again that $\cC$ has no projectors (e.g. a ``rigid'' category, e.g. a preordered category). Let $\pi_0$, $\pi_1$ be cardinals such that:}
\begin{enumerate}
\item[\textcircled{\footnotesize $1^o$}] $\forall\ s \in \Ob\cC$, $\pi_0$ is useful for $\cE_s$ and $\pi_0\geqslant\sup_{\lambda\in\Fl\cC}\pi_\lambda$ \textit{(such a $\pi_0$ exists thanks to H 1 and to the fact that $\cC$, hence $\Ob\cC$ small).}
\item[\textcircled{\footnotesize $2^o$}] $\pi_1 \geqslant \sup(\pi_0, \card\cC, \sup_{\lambda\in\Fl\cC}\pi_\lambda, \sup_{s\in\Ob\cC} \cardred \cE_{s\,\pi_0})$.
\item[\textcircled{\footnotesize $3^o$}] \textit{$\forall\lambda \in \Fl\cC$, and $x_0 \in \Ob\cE_{s,\pi_0}$, the functor $H_\lambda(x_0)$ of H 3, $x_1 \longmapsto \Hom_\lambda(x_0,x_1) : \cE_t \longrightarrow \Ens$ is $\pi_1$-accessible.}
\end{enumerate}
\DERmark{ENV1460}\DERmark{ENV1461}\DERmark{ENV1462}

\textit{Let $\pi$ be a cardinal such that $\pi^{\pi_0} = \pi \geqslant \pi_1$. Then by b), $\cE$ satisfies $L_{\pi_1}$, hence $L_\pi$, and $f : \cE \longrightarrow \cC$ is $\pi_1$-accessible. Moreover,}
\[
i_s(\cE_{s\,\pi}) \;\subset\; \cE_\pi\ ,
\]
\textit{or again}
\[
\cE(\pi) \;\subset\; \cE_\pi\ , \tag*{$(*)$}
\]
\textit{where $\cE(\pi)$ is the full subcategory of $\cE$ such that}
\[
\Ob\cE(\pi) \;=\; \bigcup_{s\in\Ob\cC} \Ob\,(\cE_{s\,\pi})\ .
\]
\textit{As $\cE$ satisfies $L_{\pi_1}$ by b), hence a fortiori $L_\pi$, we therefore have by $(*)$ a functor} fully faithful
\DERmark{ENV1463}

\[
\Ind_\pi(\cE(\pi)) \;\xrightarrow{\ \sim\ }\; \cE\ ,
\]
\DERmark{ENV1464}

\textit{and this functor is an equivalence of categories (hence $\cE$ is accessible).}
\item[\textit{d)}] \textit{In fact, under the hypotheses of c) on $\cC$, $\pi_0$, $\pi_1$, $\pi$, we have $\cE(\pi) = \cE_\pi$ (not only $\cE(\pi) \subset \cE_\pi$).}
\end{enumerate}

\msp{page 427}
\textsc{Proof of 4.15.1 and 4.15.2.} It is clear that 4.15.1 is contained in 4.15.2 a) and c). We will therefore prove 4.15.2.

a) Let $(x_i)_{i\in I}$ be a direct system in $\cE_s$, with $I$ large relative to $\pi_1$, and $x$ its direct limit in $\cE_s$ (which exists, for $\cE_s$ satisfies $\mathrm{L}_{\pi_\lambda}$, hence $\mathrm{L}_{\pi_1}$). We must prove that it is a limit in $\cE$, hence that for every $y \in \Ob\cE$, the canonical arrow
\DERmark{ENV1465}\DERmark{ENV1466}

\[
\Hom_{\cE}(x,y) \;\longrightarrow\; \varprojlim_{I} \Hom_{\cE}(x_i,y) \tag*{$(*)$}
\]
% foot: Version 17 décembre 2025

%=====================================================================
\printedpage{163}
% head: A. Grothendieck : Dérivateurs, Ch. XVIII, Édition des Universités de Montpellier II et Paris VII
is an isomorphism. Now $\Hom_{\cE}(x,y)$ and the $\Hom_{\cE}(x_i,y)$ are objects of $\Ens/\Hom_{\cC}(s,t)$, where $s = f(x)$, $t = f(y)$, and the transition arrows are arrows in $\Ens/\Hom_{\cC}(s,t)$. This shows that the bijectivity of $(*)$ comes down to that of
\[
\Hom_\lambda(x,y) \;\longrightarrow\; \varprojlim_{I} \Hom_\lambda(x_i,y) \tag*{$(*)$}
\]
for every $\lambda \in \Hom_{\cC}(f(x),f(y))$ \texttt{[rather} $\Hom_{\cC}(f(x),f(y))$\texttt{]}. For $\lambda$ fixed, and $y$ fixed in $\Ob\cE_t$, the bijectivity in $(*)$, for every direct system large relative to $\pi_1$ in $\cE_s$, with limit $x$, means that the functor
\DERmark{ENV1467}\DERmark{ENV1468}\DERmark{ENV1469}\DERmark{ENV1470}

\[
x \;\longmapsto\; \Hom_\lambda(x,y) \;:\; \cE_s \;\longrightarrow\; (\Ens)^\circ
\]
is $\pi_1$-accessible. Now this is indeed the case, by the hypothesis on $\pi_1$, whatever the arrow $\lambda : s \longrightarrow t$ with source $s$, and whatever $y \in \Ob\cE_t$.
\DERmark{ENV1471}


\msp{page 428}
Let $(x_i)_{i\in I}$ be a direct system large relative to $\pi_1$, and $(s_i)_{i\in I}$ the image of this system in $\cC$. We will prove below the

% In the left margin, beside the statement of Lemme 4.15.3 and beside the following paragraph "Mais attention ...", the edition prints a wavy marginal sign (one for each).
\textbf{Lemma 4.15.3} (\footnote{Cf. proof page 452.},\footnote{subject to the validity of 4.8.13 (But this is now established, in a sufficient form, in 4.15.10).}). \textit{Let $\cC$ be a category, $\pi$ a cardinal, such that} $\card\cC \leqslant \pi$\textit{, and that $\cC$ has no projectors. Then for every direct system $(s_i)_{i\in I}$ large relative to $\pi$ in $\cC$, there exists a cofinal subset $J$ in $I$ such that the induced system $(s_i)_{i\in J}$ is} constant, \textit{i.e. such that the $s_i$ are equal to an $s \in \cC$, and that the transition morphisms are all equal to} $\id_s$.

But watch out, we rely in proving 4.15.3 on theorem 4.8.13 (page 212-216), which I still have not proved. In the case where $\cC$ is rigid (i.e. $\End_{\cC}(s) = \{\id_s\}\ \forall\ s \in \Ob\cC$), the proof is however immediate, cf. below.

\msp{page 429}
To prove that $x = \varinjlim\limits_{I} x_i$ exists in $\cE$, and commutes with $f$, we can replace $I$ by a cofinal subcategory. This brings us back, thanks to the lemma, to the case where $(s_i)_{i\in I}$ is constant, hence to the case where the direct system comes from a direct system in a fiber category $\cE_s$. But by a), we know that its limit $x$ in $\cE_s$ is a limit in $\cE$, whence the existence of the latter, and commutation with $f$, since $\varinjlim s_i = s$ ( $\varinjlim$ of a constant filtering direct system).

c) We must prove that if $x \in \Ob\cE_s$ is $\pi$-accessible in $\cE_s$, it is so in $\cE$, i.e. that for every direct system $(y_i)_{i\in I}$ large relative to $\pi$, the canonical arrow
\[
\varinjlim_{i} \Hom(x,y_i) \;\longrightarrow\; \Hom(x,y) \tag*{$(**)$}
\]
is bijective. Now as by 2°) $\pi \geqslant \pi_1 \geqslant \sup(\card\cC, \sup_\lambda \pi_\lambda)$, it follows again from the lemma % [sic] card S -> card \mathcal{C}, no S is defined; the condition 2° bounds card \mathcal{C}
\DERmark{ENV1472}

that, possibly replacing $I$ by a cofinal subcategory, we can assume that the direct system $(y_i)_{i\in I}$ comes from a system in a fiber category $\cE_t$, and that $y$ is its $\varinjlim$ in $\cE_t$. But then the sets $\Hom(x,y)$ and $\Hom(x,y_i)$, and the transition morphisms between them, are again in $\Ens/\Hom(s,t)$, and we are reduced
% foot: Version 17 décembre 2025


\endgroup
% END UNIT der-ch18/der-ch18_p161-165.tex
% BEGIN UNIT der-ch18/der-ch18_p166-170.tex
\begingroup\makeatletter%
\newcommand{\printedpage}[1]{\marginpar{\footnotesize p.~#1}}%
\newcommand{\msp}[1]{\marginpar{\footnotesize [#1]}}%
\newcommand{\cE}{\mathcal{E}}%
\newcommand{\cC}{\mathcal{C}}%
\newcommand{\Ob}{\operatorname{Ob}}%
\newcommand{\Fl}{\operatorname{Fl}}%
\newcommand{\Hom}{\operatorname{Hom}}%
\newcommand{\Ind}{\operatorname{Ind}}%
\newcommand{\card}{\operatorname{card}}%
\newcommand{\red}{\operatorname{\text{red}}}%
\newcommand{\id}{\mathrm{id}}%
\newcommand{\Ens}{\mathrm{Ens}}%
\newcommand{\circled}[1]{\tikz[baseline=(c.base)]\node[draw,circle,inner sep=1pt](c){#1};}%
\makeatother\let\printedpage\derpp\let\msp\dermsp%
\let\mspin\DERmspin%


\printedpage{164}
% head: A. Grothendieck : Dérivateurs, Ch. XVIII, Édition des Universités de Montpellier II et Paris VII

\msp{page 430}to prove that for every $\lambda : s \longrightarrow t$, the arrow
\[
\varinjlim_{i} \Hom_{\lambda}(x, y_i) \ \longrightarrow \ \Hom_{\lambda}(x, y)
\]
is bijective. For fixed $x$ and $\lambda$, and varying direct systems in $\cE_t$, large relative to $\pi$, this means that the functor
\[
H_{\lambda}(x) \ = \ (y \longmapsto \Hom_{\lambda}(x, y)) \ : \ \cE_t \ \longrightarrow \ \Ens
\]
is $\pi$-accessible, and this for any
\[
x \ \in \ \Ob \cE_{s\,\pi} \ .
\]
Now this is the case if $x \in \Ob \cE_{s\,\pi_0}$, by virtue of hypothesis $3^{\circ}$) (then $H_{\lambda}(x_0)$ is $\pi_1$-accessible, % [sic] $s \in \Ob \cE_{s\,\pi_0}$ -> $x_0 \in \Ob \cE_{s\,\pi_0}$, the condition concerns an object of $\cE_{s\,\pi_0}$
\DERmark{ENV1473}

and a fortiori $\pi$-accessible). We have to pass from there to the case where $x \in \Ob \cE_{s\,\pi}$. Now since $\pi_0$ is useful for $\cE_s$ by $1^{\circ}$, we have
\[
x \ = \ \varinjlim_{j \in J} x_j \ , \qquad x_j \in \Ob \cE_{s\,\pi_0},\ J \text{ large relative to } \pi_0 \ ,
\]
and since by $2^{\circ}$, $\pi^{\pi_0} = \pi \geqslant \pi_1 \geqslant \sup(\pi_0, \card \red \cE_{s\,\pi_0})$, we conclude from this that we can take $J$ so that
\[
\card J \ \leqslant \ \pi \ .
\]
On the other hand, by virtue of $\pi \geqslant \pi_1 \ \geqslant \ \pi_{\lambda}$, the functor
\[
\left\{
\begin{array}{lcl}
x_0 & \longmapsto & H_{\lambda}(x_0) \\[1ex]
\cE_s & \longrightarrow & (\underline{\Hom}(\cE_t, \Ens))^{\circ}
\end{array}
\right.
\]
\DERmark{ENV1474}

is $\pi_0$-accessible, so we have
\[
H_{\lambda}(x) \ \simeq \ \varprojlim_{j \in J} H_{\lambda}(x_j) \ .
\]
\DERmark{ENV1475}


\msp{page 431}Since $x_j \ \in \ \Ob \cE_{s\,\pi_0}$, $H_\lambda(x_j)$ is a $\pi_1$-accessible functor, and since $\card J \leqslant \pi$, the $\varprojlim\limits_{j \in J} H_{\lambda}(x_j)$ is also $\pi$-accessible (cf. SGA 4 I 9.8).

Thus we have proved $\cE_{s\,\pi} \subset \cE_\pi$ for every $s \in \Ob \cC$, i.e. the inclusion
\DERmark{ENV1476}\DERmark{ENV1477}

\[
\cE(\pi) \ \subset \ \cE_{\pi} \ ,
\]
whence the fully faithful functor
\[
\Ind_{\pi}(\cE(\pi)) \longrightarrow \cE \ , \tag*{$(*)$}
\]
since $\cE$ satisfies $\mathrm{L}_{\pi}$. It remains to prove that it is essentially surjective. But if $x \in \Ob \cE$, $s = f(x)$, we have $x \in \Ob \cE_s$, and since $\pi_0$ is useful for $\cE_s$ by $1^{\circ}$), and $\pi^{\pi_0} = \pi > \pi_0$,
$\pi$ is also useful for $\cE_s$, so we have \textit{in} $\cE_s$ % [sic] utile pour $\pi$ -> utile pour $\cE_s$, usefulness is a property relative to the category $\cE_s$
\DERmark{ENV1478}

\[
x \ = \ \varinjlim_{I} x_i \ , \qquad \text{the } x_i \in \cE_{s\,\pi} \subset \cE(\pi),\ I \text{ large relative to } \pi \ ,
\]
\printedpage{165}
% head: A. Grothendieck : Dérivateurs, Ch. XVIII, Édition des Universités de Montpellier II et Paris VII
and by a) (since $\pi \geqslant \pi_1 \geqslant \sup_{\lambda \in \Fl \cC} \pi_{\lambda}$) $x$ is also a limit of the $x_i$ in $\cE$, so it is in the essential image of $(*)$.

d) It remains to prove that we have $\cE_{\pi} \ \subset \ \cE(\pi)$. But we know (by 4.8.~\dots), since $\cE \ \simeq \ \Ind_{\pi}(\cE(\pi))$, that $\cE_{\pi} \ = \ \cE(\pi)^{\natural}$, i.e. that for $x \ \in \ \Ob \cE_{\pi}$, $\exists \ y \ \in \ \cE(\pi)$ and a projector $p$ in $y$, such that $x$ is isomorphic to the image of $p$ in $\cE$. Now $f(p)$ is a projector in $f(y) = s$, so by the hypothesis on $\cC$, $f(p) = \id_s$, i.e. $p$ is a projector in $\cE_s$. Since therefore $x$ is in $\cE_s$, and
\DERmark{ENV1479}


\msp{page 432}even in $\cE_{s\,\pi}$, since $\cE_{s\,\pi}$ is stable under direct factors in $\cE_s$. This proves that $x \in \Ob \cE(\pi)$, q.e.d.

We can present 4.15.1 in a more congenial form.

\textbf{Corollary 4.15.4.} \textit{Let $f : \cE \longrightarrow \cC$ be a functor of $\mathfrak{U}$-categories. The following conditions are equivalent:}
\begin{enumerate}
\item[\textit{a)}] \textit{For every functor $\cC' \longrightarrow \cC$, with $\cC'$ small and without projectors, $\cE' \overset{\mathrm{def}}{=} \cE \times_{\cC} \cC'$ is accessible, and $f' : \cE' \longrightarrow \cC'$ is accessible.}
\item[\textit{b)}] \textit{As in a), but with $\cC' = \boldsymbol{\Delta}_1$, in other words, this is condition $H$ p.~409: for every $\lambda \in \Fl \cC$, $\cE_{\lambda} = \cE \times_{\cC} (\boldsymbol{\Delta}_1, i_{\lambda})$ is accessible, and $\cE_{\lambda} \longrightarrow \boldsymbol{\Delta}_1$ is accessible.}
\item[\textit{c)}] \textit{The condition $H' = H\ \mathit{1} + H\ \mathit{2} + H\ \mathit{3}$, made explicit in the statement of 4.15.1.}
\end{enumerate}
\DERmark{ENV1480}

\textsc{Proof}. It is done in a circle,
\[
\mathrm{a} \ \Longrightarrow \ \mathrm{b} \ \Longrightarrow \ \mathrm{c} \ \Longrightarrow \ \mathrm{a} \ ,
\]
the implication a $\Longrightarrow$ b being tautological, b $\Longrightarrow$ c being almost so (cf. 4.14.24, a $\Longrightarrow$ b $\Longrightarrow$ c). In fact, conditions b and c are conditions on the categories $\cE_\lambda \xrightarrow{f_{\lambda}} \boldsymbol{\Delta}_1$ for $\lambda \in \Fl \cC$, hence stable under base change. To prove
\DERmark{ENV1481}


\msp{page 433}c $\Longrightarrow$ a, we can therefore assume that $\cC' = \cC$, $\cC' \longrightarrow \cC$ the identity functor, and prove that if c) is satisfied and if $\cC$ is small and without projectors, then $\cE$ is accessible and $f$ accessible. Now this is what 4.15.1 asserts.

\textbf{Remark 4.15.5.} If $f : \cE \longrightarrow \cC$ is a fully faithful functor, it is immediate that condition H is satisfied (for $\cE_{\lambda} \longrightarrow \boldsymbol{\Delta}_1$ will be fully faithful, so $\cE_{\lambda}$ is equivalent
to $e$ or to $\boldsymbol{\Delta}_1$, hence accessible, and every functor from $\cE_{\lambda}$ to any category is % [sic] où -> ou, disjunction "e ou Δ1"
accessible \dots). But if we assume $\cC$ small \textit{without} assuming $\cC$ without projectors, it does not follow that $\cE$ is accessible, or what amounts to the same thing (being essentially small), that $\cE$ is Karoubian. (The most trivial example is the one where $\cE = \cC$, $f = \id_{\cC}$, $\cE$ is Karoubian only if $\cC$ is!) But even assuming $\cC$ a small Karoubian category, it does not necessarily follow that $\cE$ is. Indeed, \textit{every} small category $\cE$ (Karoubian or not) embeds as a full subcategory in its Karoubian envelope $\cC$ -- so it suffices to take $\cE$ non-Karoubian. The example

\msp{page 434}that is most economical is
\[
\Ob \cE \ = \ \{0\} \ , \qquad \Fl(\cE) \ = \ \{\id_0, p\} \ , \qquad \text{with } p^2 = p \neq \id_0 \ ,
\]
\printedpage{166}
% head: A. Grothendieck : Dérivateurs, Ch. XVIII, Édition des Universités de Montpellier II et Paris VII
the Karoubi envelope of $\cE$ is given by
\[
P \ = \ \left\{
\begin{tikzcd}
1 \arrow[r, shift right, "\sigma"'] & 0 \arrow[l, shift right, "\rho"'] \arrow[loop right, "p"]
\end{tikzcd}
\qquad
\begin{aligned}
\rho\sigma \ &= \ \id_1 \\
\sigma\rho \ &= \ p \ \neq \ \id_0
\end{aligned}
\quad \Longrightarrow \ p^2 = p
\right\} .
\]
This raises the question of finding a statement similar to 4.15.1, giving simple serviceable conditions, stronger than $\mathrm{H}' = \mathrm{H}$ (1+2+3), on $f : \cE \longrightarrow \cC$, for $\cE$ to be accessible and $f$ accessible, assuming $\cC$ \textit{Karoubian}. Or again, serviceable necessary and
sufficient conditions (stronger than 4.15.4) on $f : \cE \longrightarrow \cC$, so that every \textit{small} % [sic] suffisants -> suffisantes, agreement with "conditions"
\textit{Karoubian} category $\cC'$, and every functor $\cC' \longrightarrow \cC$, the category $\cE' = \cE \times_{\cC} \cC'$ be accessible, and the functor $\cE' \longrightarrow \cC'$ accessible. A natural idea is that it would suffice for this conclusion to be true when $\cC'$ is one of the two categories $\boldsymbol{\Delta}_1$, $P$ (where $P$ is the category with two objects, with three nonidentity arrows, defined above). But I doubt that this is so (\footnote[178]{That should work on the other hand if we take the categories $0 \longrightarrow 1$,
\[
\begin{array}{c}
2 \longrightarrow \underbrace{\begin{tikzcd}[ampersand replacement=\&]
1 \arrow[r, shift right, "\sigma"'] \& 0 \arrow[l, shift right, "\rho"'] \arrow[loop right, "p"]
\end{tikzcd}}_{P}
\quad \text{and} \quad
2 \longleftarrow \underbrace{\begin{tikzcd}[ampersand replacement=\&]
1 \arrow[r, shift right, "\sigma"'] \& 0 \arrow[l, shift right, "\rho"'] \arrow[loop right, "p"]
\end{tikzcd}}_{P}
\qquad (\text{\footnotemark[179]}) \ , \\
\nwarrow \ \begin{array}{c} \text{3 objects and} \\ \text{5 nonidentity} \\ \text{arrows} \end{array} \ \nearrow
\end{array}
\]
thanks to 4.15.7, page 440.}\footnotetext[179]{these categories represent respectively the functors $\mathrm{Cat} \longrightarrow \Ens$ assigning to a category $\cC$ the diagrams in $\cC$ consisting of $(x, y, p, u)$, with $p$ a projector on $x$ \texttt{[rather in $\cC$]}, and $u : y \longrightarrow x^{\natural}$ (i.e. $u : y \longrightarrow x$ with $u = pu$), resp. $u : x^{\natural} \longrightarrow y$ (i.e. $u : x \longrightarrow y$ with $u = up$).}).
\DERmark{ENV1482}


We will all the same try a somewhat very special statement in this direction:

% Verified on PDF leaf8: a non-textual wavy marginal mark beside items a), b); no mathematical token omitted.
\msp{page 435}\textbf{Proposition 4.15.6} (\footnote[180]{\textit{False?}}). \textit{Let $f : \cE \longrightarrow \cC$ be a functor of $\mathfrak{U}$-categories. We assume:}
\begin{enumerate}
\item[\textit{a)}] \textit{$f$ is fibrant or cofibrant} (\footnote[181]{The fibrant case is not proved, perhaps \textit{false} -- and likewise for the cofibrant case!}).
\item[\textit{b)}] \textit{The fiber categories $\cE_s$ ($s \in \Ob \cC$) are accessible, and for every arrow $\lambda : s \longrightarrow t$ in $\cC$, the base change functor $\lambda^* : \cE_t \longrightarrow \cE_s$ (resp. cobase change $\lambda_* : \cE_s \longrightarrow \cE_t$) is accessible.}
\item[\textit{c)}] \textit{$\cC$ is essentially small and Karoubian.}
\end{enumerate}
\textit{Then $\cE$ is accessible, and $f$ is accessible.}

\textsc{Proof}. We have already seen that conditions a) and b) imply $\mathrm{H}'$, i.e. H (1+2+3). Besides, if $\cC'$ is a category and $\cC' \xrightarrow{\ \sim\ } \cC$ an equivalence, so then $\cE' \overset{\mathrm{def}}{=} \cE \times_{\cC} \cC' \longrightarrow \cE$ is an equivalence, so the desired conclusion for $\cE$, $f$ is equivalent to the same conclusion for $\cE'$, $f' : \cE' \longrightarrow \cC'$. So we can assume that $\cC$ is small (and even reduced, if one insists). By virtue of 4.15.2 a), we thus see that the inclusion functors $\cE_s \longrightarrow \cE$ are accessible (and even $\pi_1$-accessible) as soon as $\pi_1 \geqslant \sup_{\lambda \in \Fl \cC} \pi_{\lambda}$. We also see
\DERmark{ENV1483}


\msp{page 436}by AQT \texttt{[âne qui trotte: at a donkey's trot, i.e. by routine verification]} that we can find cardinals $\pi_0$, $\pi_1$ with the following properties:

\printedpage{167}
% head: A. Grothendieck : Dérivateurs, Ch. XVIII, Édition des Universités de Montpellier II et Paris VII
\begin{enumerate}
\item[\circled{1\textsuperscript{o}}] $\pi_0$ is useful for $\cE_s$, $\forall \ s \in \Ob \cC$.
\item[$2^{\circ}$)] $\pi_1 \geqslant \sup(\pi_0, \card \cC, \sup_{\lambda \in \Fl \cC} \pi_{\lambda}, \sup_{s \in \Ob \cC} \card \red \cE_{s\,\pi_0})$.
\item[\circled{3\textsuperscript{o}}] $\forall \ (\lambda : s \longrightarrow t) \in \Fl \cC$ and $x_0 \in \Ob \cE_{s\,\pi_0}$, the functor
\[
H_{\lambda}(x_0) \ : \ (x_1 \longmapsto \Hom_{\lambda}(x_0, x_1)) \ : \ \cE_t \ \longrightarrow \ (\Ens)
\]
\DERmark{ENV1484}

is $\pi_1$-accessible. (\textbf{N.B.} If $f$ is cofibrant, this means that $\forall \ \lambda \in \Fl \cC$, $\lambda : s \longrightarrow t$,
\fbox{$\lambda_*(\cE_{s\,\pi_0}) \ \subset \ \cE_{t\,\pi_1}$}. If $f$ is fibrant, it suffices to require that $\forall \ \lambda \in \Fl \cC$, $\lambda^* : \cE_t \longrightarrow \cE_s$ % [sic] $\cE_{s\,\pi_1}$ -> $\cE_{t\,\pi_1}$, $\lambda_*$ goes to $\cE_t$, cf. the later display "(contenu dans 3°)"
\DERmark{ENV1485}

be $\pi_1$-accessible.) But in this case, we also want \fbox{$\lambda^*(\cE_{t\,\pi_0}) \ \subset \ \cE_{s\,\pi_1}$}. % [sic] $\lambda^*(\cE_{s\,\pi_0})$ -> $\lambda^*(\cE_{t\,\pi_0})$, $\lambda^*$ is defined on $\cE_t$, cf. the later display "(contenu dans 3°)"
\DERmark{ENV1486}

\item[\circled{4\textsuperscript{o}}] The functors $\lambda^*$ (case $f$ fibrant) resp. $\lambda_*$ (case $f$ cofibrant) are $\pi_1$-accessible.
\end{enumerate}
Let $\pi$ then be a cardinal such that
\[
\pi \ = \ \pi^{\pi_0} \ \geqslant \ \pi_1 \ , \tag*{$(*)$}
\]
and let $\cE(\pi)$ be the (strictly) full subcategory of $\cE$ defined by
\[
\Ob \cE(\pi) \ = \ \bigcup_{s \, \in \, \Ob \cC} \Ob \cE_{s\,\pi} \ .
\]
Thanks to $(*)$, and to $1^{\circ}$, $4^{\circ}$ and to
\[
\left\{
\begin{array}{l}
\pi_1 \ \geqslant \ \sup_s(\card \red \cE_{s\,\pi_0}) \qquad \text{(contained in $2^{\circ}$)} \\[1ex]
\lambda^*(\cE_{t\,\pi_0}) \ \subset \ \cE_{s\,\pi_1} \qquad \text{resp.} \qquad \lambda_*(\cE_{s\,\pi_0}) \ \subset \ \cE_{t\,\pi_1} \qquad \text{(contained in $3^{\circ}$)} \ ,
\end{array}
\right.
\]
we see that $\cE(\pi)$ is a

\msp{page 437}fibered subcategory, resp. a cofibered subcategory, over $\cC$ (the functor $\cE(\pi) \longrightarrow \cE$ thus being cartesian resp. cocartesian). Up to $\cC$-equivalence, $\cE$ is then reconstructed from $\cE(\pi)$, by passing to the $\Ind_{\pi}(\cE_{s\,\pi})$, which form a fibered resp. cofibered system over $\cC$ via the functors $\Ind_{\pi}(\lambda^*_{\pi})$ resp. $\Ind_{\pi}(\lambda^{\pi}_*)$ \dots

Let us prove that $\cE$ satisfies $\mathrm{L}_{\pi}$ and that $f : \cE \longrightarrow \cC$ is $\pi$-accessible (\footnote[182]{Here we use only $\pi \geqslant \card \red \cC$.}). Since $\cC$ is Karoubian, and $\pi \geqslant \pi_1 \geqslant \card \cC$, $\cC$ satisfies $\mathrm{L}_{\pi}$ and the ind-objects of $\cC$ large relative to $\pi$ are essentially constant (4.8.13). Let $I$ then be large relative to $\pi$, and $I \longrightarrow \cE$, $i \longmapsto x_i$, be a functor, we have to construct a $\varinjlim$ in $\cE$, using the already known $\varinjlim$ $s$ of the $s_i = f(x_i)$. We will take an $x$ in $\cE_s$. Let us first take the case where $f$ is cofibrant. We see that the $\beta_{i\,*}(x_i)$ form a direct system in $\cE_s$, since for $i \leqslant j$, we have
\DERmark{ENV1487}\DERmark{ENV1488}\DERmark{ENV1489}

\[
\beta_{i\,*}(x_i) \ = \ \beta_{j\,*} \underbrace{\beta_{ji\,*}(x_i)}_{\longrightarrow \, x_j} \ \longrightarrow \ \beta_{j\,*}(x_j) \ . % [sic] $\beta_{i\,*}(x_j)$ -> $\beta_{i\,*}(x_i)$, the object transported from $s_i$ is $x_i$
\]
\DERmark{ENV1490}

\[
\begin{tikzcd}[row sep=small]
x_i \arrow[rr, "\alpha_i"] & & x \, ? \\
s_i \arrow[rr, "\beta_i"] \arrow[dr, "\beta_{ji}"'] & & s \\
& s_j \arrow[ur, "\beta_j"'] &
\end{tikzcd}
\]
We set

\printedpage{168}
% head: A. Grothendieck : Dérivateurs, Ch. XVIII, Édition des Universités de Montpellier II et Paris VII
\msp{page 438}
\[
x \ = \ \varinjlim_{I} \beta_{i\,*}(x_i) \ , \qquad \varinjlim \ \textit{in } \cE_s \ .
\]
We know that this is a $\varinjlim$ in $\cE$, hence that for every $y$ in $\cE$, we have
\[
\Hom(x, y) \xrightarrow{\ \sim\ } \varprojlim_{I} \Hom(\beta_{i\,*}(x_i), y) \ ,
\]
or what amounts to the same thing, for every arrow $\lambda : s \longrightarrow t$, and $y$ in $\Ob \cE_t$, we have
\[
\underbrace{\Hom_{\lambda}(x, y)}_{= \, \Hom_{\cE_t}(\lambda_* x, y)} \ \simeq \ \varprojlim \Hom_{\cE_t}(\underbrace{\lambda_* \beta_{i\,*}(x_i)}_{(\lambda \beta_i)_*(x_i)}, y) \ .
\]
\DERmark{ENV1491}

This also follows from the fact that $\lambda_*$ is $\pi$-accessible (being $\pi_1$-accessible) \dots{} We easily check that this is a $\varinjlim$ of the $x_i$ in $\cE$. And of course, this limit commutes with $f$ by construction.

But in the case where $f$ is fibrant, I do not see how to go about constructing a limit in $\cE_s$. It is true that we can assume that all the $s_i$ are equal to $s$, but the transitions $\beta_{ji}$ are not the identity, nor even isomorphisms. We would have to take $\cC = P$ to try to understand -- I am not going to insist (\footnote[183]{I have checked that if $f$ is fibrant, then $\cE$ is \textit{Karoubian}. So there is hope \dots}).

\msp{page 439}In the case where $f$ is fibrant, let us admit (as an additional hypothesis if need be) that $\cE$ satisfies $L$ and that $f$ is accessible, and take $\pi$ large enough for $\cE$, $f$ to satisfy $\mathrm{L}_{\pi}$. We can then finish the proof of 4.15.6. We have to prove
\[
\cE(\pi) \ \subset \ \cE_{\pi} \ ,
\]
hence that if $x \in \Ob \cE_{s\,\pi}$, and $(y_i)_{i \in I}$ is a direct system in $\cE$ large relative to $\pi$, with limit $y$, we have
\[
\varinjlim \Hom_{\cE}(x, y_i) \xrightarrow{\ \sim\ } \Hom_{\cE}(x, y) \ .
\]
So we have to prove that every arrow $x \xrightarrow{u} y$ factors, in an essentially unique way $u_i$, through a $y_i$. But the arrow $s \xrightarrow{\ \lambda \, = \, f(u)\ } t$ in $\cC$ factors, in an essentially unique way, through some $\lambda_i \, : \, s \longrightarrow t_i$. So we are reduced to finding a \textit{$\lambda_i$-morphism} $x \longrightarrow y_i$ for large $i$, which factors. In the fibrant case, this means that we are looking for $x \longrightarrow \lambda^*_i(y_i)$ which factors $x \longrightarrow \lambda^*(y)$, now the $\lambda^*_i(y_i)$ form a direct system in $\cE_s$, and its limit \textit{in} $\cE_s$ is $\lambda^*(y)$ (\textit{to be checked}). Since $x$ is in $\cE_{s\,\pi}$, the given hom\texttt{[omorphism]} $x \longrightarrow \lambda^*(y)$ factors in

\msp{page 440}an essentially unique way through a $\lambda^*_i(y_i)$.

This time, however, it is the case where $f$ is cofibrant that poses a problem. In fact, I cannot work it out.

Thus, we do find a functor
\[
\Ind_{\pi}(\cE(\pi)) \ \longrightarrow \ \cE \ ,
\]


\endgroup
% END UNIT der-ch18/der-ch18_p166-170.tex
% BEGIN UNIT der-ch18/der-ch18_p171-175.tex
\begingroup\makeatletter%
\newcommand{\printedpage}[1]{\marginpar{\footnotesize p.~#1}}%
\newcommand{\msp}[1]{\marginpar{\footnotesize [#1]}}%
\newcommand{\ux}{\underline{x}}%
\newcommand{\cC}{\mathcal{C}}%
\newcommand{\Ob}{\operatorname{Ob}}%
\newcommand{\card}{\operatorname{card}}%
\newcommand{\id}{\mathrm{id}}%
\newcommand{\Hom}{\operatorname{Hom}}%
\newcommand{\Ind}{\operatorname{Ind}}%
\newcommand{\Img}{\operatorname{Im}}%
\newcommand{\Ker}{\operatorname{Ker}}%
\newcommand{\defeq}{\overset{\mathrm{def}}{=}}%
\makeatother\let\printedpage\derpp\let\msp\dermsp%
\let\mspin\DERmspin%

\DERmark{ENV1492}


\setcounter{footnote}{183}

\printedpage{169}
% head: A. Grothendieck : Dérivateurs, Ch. XVIII, Édition des Universités de Montpellier II et Paris VII
which is essentially surjective, but for lack of knowing that $\mathcal{E}(\pi) \subset \mathcal{E}_\pi$ in the cofibrant case, it is not clear that this functor is an equivalence of categories.

\textbf{4.15.7.} I must postpone the proof of lemma 4.15.3 (p. 428), assuming 4.8.13. (Not yet proved -- I am going to come back below to the proof of 4.8.13.) We are going to use the

\textbf{Lemma 4.15.7.1.} \textit{Let $\cC$ be a category, $\ux = (x_i)_{i\in I}$ an ind-object of $\cC$. For $\ux$ to be essentially constant, i.e. isomorphic, in} $\Ind(\cC)$\textit{, to the constant ind-object defined by an object}

\msp{page 441}
\textit{of $\cC$, it is necessary and sufficient that there exist a cofinal functor $\varphi : J \longrightarrow I$, with $J$ filtering, in such a way that the reindexed system $(x_j)_{j\in J}$, with $x_j = x_{\varphi(j)}$, have the following property: for every $j \in J$, one can find a projector $p_j$ in $x_j$, in such a way that:}
\begin{itemize}
\item[\textit{1°)}] \textit{For every arrow $\lambda : j \longrightarrow j'$ in $J$, we have commutativity in}
\[
\begin{tikzcd}
x_j \arrow[r, "x_\lambda"] \arrow[d, "p_j"] & x_{j'} \arrow[d, "p_{j'}"] \\
x_j \arrow[r, "x_\lambda"] & x_{j'}
\end{tikzcd}
\]
\textit{(where the $x_\lambda$ are the transition arrows), i.e. $p_{j'}x_\lambda = x_\lambda p_j$.}
\item[\textit{2°)}] \textit{For every $j \in \Ob J$, there exists an arrow $\lambda : j \longrightarrow j'$ in $J$ such that $p_{j'}x_\lambda = x_\lambda p_j = x_\lambda$.}
\item[\textit{3°)}] \textit{For every $j \in \Ob J$, the projector $p_j$ in $x_j$ admits an image $x_j^\natural$.} % [sic] un image -> une image, gender agreement
\item[\textit{4°)}] \textit{For every arrow $\lambda : j \longrightarrow j'$ in $J$, the arrow $x_\lambda^\natural : x_j^\natural \longrightarrow x_{j'}^\natural$ induced by $x_\lambda$ (thanks to 1°) is an isomorphism.}
\end{itemize}
\textit{Moreover, we can suppose $J$ of the form $i_0\backslash I$, the cofinal functor $J \longrightarrow I$ being the canonical functor $i_0\backslash I \longrightarrow I$ ($i_0$ being a suitable fixed object of $I$)} \texttt{[incomplete sentence]}
\DERmark{ENV1493}


\msp{page 442}
Let us note at once two corollaries.

\textbf{Corollary 4.15.7.2.} \textit{Suppose that $\cC$ is without projectors. Then $\ux = (x_i)_{i\in I}$ is an essentially constant ind-object if and only if there exists a cofinal functor, $J \longrightarrow I$, with $J$ filtering (\footnote{and we can suppose $J = i_0\backslash I \xrightarrow{\text{can.}} I$}), in such a way that the reindexed direct system $(x_j)_{j\in J}$ is isomorphic} (in the category $\underline{\mathrm{Hom}}(J,\cC)$) \textit{to a} constant \textit{direct system (i.e. defined by a constant functor $J \longrightarrow \cC$: all the $x_j$ equal to one and the same $x$, and the transition morphisms identities).}

Sufficiency is clear, for necessity we apply 4.15.7.1, which tells us that we can reindex so as to satisfy conditions 1°) to 4°) (for suitable projectors $p_j$ in the $x_j$). But the hypothesis on $\cC$ implies that $p_j = \id_{x_j}\ \forall\, j \in \Ob J$, and conditions 1°, 2°, 3° are satisfied tautologically for such a choice, condition 4° means that the transition morphisms are isomorphisms. But it is immediate that for a given direct system $(x_j)_{j\in J}$, with $J$ 1-\textit{connected} (e.g. $J$ filtering), the condition that the $x_\lambda$ are isomorphisms % [sic] x_j -> x_\lambda, the transition morphisms are meant
\DERmark{ENV1494}

is necessary and sufficient for this system to be isomorphic to a constant direct system, whence the conclusion.
% foot: Version 17 décembre 2025

\printedpage{170}
% head: A. Grothendieck : Dérivateurs, Ch. XVIII, Édition des Universités de Montpellier II et Paris VII
\msp{page 443}
\textbf{Corollary 4.15.7.3.} \textit{Suppose that $I$ is large relative to $\pi_0$, where $\pi_0$ is an infinite cardinal such that $\pi_0 \geqslant \card \cC$. Then $(x_i)_{i\in I}$ is essentially constant if and only if we can find a cofinal functor $J \longrightarrow I$ with $J$ large relative to $\pi_0$ (and ordered if we like) (\footnote{The proof (slightly rearranged) shows that we can take for $J$ a subcategory (not full) of a category $i_1\backslash I$, for a suitable object $i_1$ of $I$. Thus we have $\card J \leqslant \pi$ if $\card I \leqslant \pi$, $\pi$ being a given infinite cardinal.}), in such a way that the reindexed system $(x_j)_{j\in J}$ satisfies the following conditions:}
\begin{itemize}
\item[\textit{1°)}] \textit{The map $j \longmapsto x_j$ from $\Ob J$ to $\Ob\cC$ is constant, i.e. there exists an $x_0 \in \Ob\cC$ such that $x_j = x_0\ \forall\, j \in \Ob J$.}
\item[\textit{2°)}] \textit{There exists a projector $p$ in $x_0$ such that for $j < j'$ in $J$, we have $\tau_{j',j} = p$.}
\item[\textit{3°)}] \textit{The projector $p$ admits an image in $\cC$.}
\end{itemize}
\textsc{Proof}. These conditions obviously imply those of 4.15.7.1, taking $p_j = p\ \forall\, j \in \Ob J$, at least if $J$ has no greatest element (which causes a difficulty for 2°). But if $J$ has a final object, the same is true of $I$, and it is clear then that every direct system indexed by $I$ is essentially constant. So, in these cases, conditions 1°) 2°) 3°) above imply that $\ux$ is essentially constant. Converse\msp{page 444}ly, suppose that $\ux$ is essentially constant, and let us prove that we can reindex as indicated. The case where $I$ has a final object is trivial, we shall take $J = e$, $p = \id$. So suppose that $I$ has no final object. Since $I$ large relative to $\pi_0$ and $\pi_0 \geqslant \card\Ob\cC$, we see that there exists an $x_0 \in \Ob\cC$ such that the full subcategory $I' = I_{x_0}$ of $I$, consisting of the $i$ such that $x_i = x_0$, is cofinal in $I$. Since $(x_i)_{i\in I}$ is essentially constant, the same is true of $(x_{i'})_{i'\in I'}$, and we can apply criterion 4.15.7.1 to it, with \marginpar{$\Vert$}% double vertical bar printed in the left margin beside the two following lines
\DERmark{ENV1495}

$J \longrightarrow I'$ cofinal. The proof will show moreover that if $I'$ is large relative to $\pi_0$ (an infinite cardinal with no other conditions), we can suppose $J$ large relative to $\pi_0$ (\footnote{take this into account in the proof of 4.15.7.1.}). Moreover, it is clear that we can suppose $J$ ordered. Since $\pi_0 \geqslant \card\Hom(x_0,x_0)$, we see again that there exists a projector $p$ in $x_0$ such that the subset $J_1 = J'$ of $J$, consisting of the $j$ such that $p_j = p$, is cofinal. Being cofinal, it too is large relative to $\pi_0$. So we can suppose that all the $p_j$ are equal to one and the same $p$. Conditions 1°-4° become:
\begin{itemize}
\item[1°)] $p\tau_{j'j} = \tau_{j'j}p$ for $j \leqslant j'$.
\item[2°)] For every $j$, exists $j' \geqslant j$ with
\[
p\tau_{j'j} \;=\; \tau_{j'j}p \;=\; \tau_{j' j}\ . % [sic] \tau_{jj'} -> \tau_{j'j}, index order used everywhere else
\]
\DERmark{ENV1496}

\item[3°)] The image $x_0^\natural$ of the projector $p$ in $x_0$ exists.
\item[4°)] The endomorphism $\tau^\natural_{j'j}$ of $x_0^\natural$ is an isomorphism, for $j \leqslant j'$.
\end{itemize}
\msp{page 445}
Now we have the

\textbf{Lemma 4.15.7.4.} \textit{Let $\cC$ be a category, $x_0$ an object of $\cC$, $p$ a projector in $x_0$, $\tau$ an endomorphism of $x_0$ satisfying the relations}
\[
\tau p \;=\; p\tau \;=\; \tau\ , \tag*{$(*)$}
\]
% foot: Version 17 décembre 2025
\printedpage{171}
% head: A. Grothendieck : Dérivateurs, Ch. XVIII, Édition des Universités de Montpellier II et Paris VII
\textit{and such that the automorphism of $x_0^\natural \defeq$} $\Img p$ \textit{(defined e.g. as an object of $\cC^\wedge$, $x_0^\natural =$} $\Ker(\id_{x_0},p)$) % [sic] "(défini ... Ker(id_{x_0},p)" -> a closing parenthesis is missing, the parenthesis opened before "défini" is never closed
\DERmark{ENV1497}

\textit{be the identity. Then we have $\tau = p$.}

\textsc{Proof of} 4.15.7.4. We are reduced at once, by the standard procedure, to the case where $\cC = \mathrm{Ens}$, where the verification is immediate: the endomorphisms $\tau$ of the set $E$ $(= x)$ which satisfies $(*)$ correspond one-to-one to the endomorphisms $u$ of $E^\natural = pE$, via
\DERmark{ENV1498}

\[
\begin{cases}
u \;\longmapsto\; \sigma u \rho \\
\tau \;\longmapsto\; \tau|E^\natural
\end{cases}
\qquad
\text{(where }
\begin{tikzcd}[column sep=small]
E^\natural \arrow[r, hook, shift left, "\sigma"] & E \arrow[l, shift left, "\rho"]
\end{tikzcd}
\text{ are the canonical arrows)}\ ,
\]
so $\tau|E^\natural = \id \Longrightarrow \tau = \sigma\rho = p$.
\DERmark{ENV1499}


Let us choose a $j_0 \in J$, and consider the homomorphisms
\[
\tau^\natural_{j,j_0} \;:\; x_0^\natural \longrightarrow x_0^\natural
\]
induced by the $\tau_{j,j_0}$ with $j \geqslant j_0$. Here again, since the set of possible isomorphisms is of

\msp{page 446}
cardinality $\leqslant \pi_0$, and $j_0\backslash J$ is large relative to $\pi_0$, we see that there exists $J' \subset j_0\backslash J$, cofinal in $j_0\backslash J$, hence in $J$, such that for $j \in J'$, the $\tau^\natural_{j,j_0}$ are equal to one and the same $\tau : x_0^\natural \xrightarrow{\sim} x_0^\natural$. So if $j', j'' \in J'$, with $j' \leqslant j''$, we have
\DERmark{ENV1500}\DERmark{ENV1501}

\[
\underbrace{\tau^\natural_{j''j_0}}_{\tau} \;=\; \tau^\natural_{j''j'}\,\underbrace{\tau^\natural_{j'j_0}}_{\tau}\ ,
\]
\DERmark{ENV1502}\DERmark{ENV1503}\DERmark{ENV1504}

whence
\[
\tau^\natural_{j''j'} \;=\; \id_{x_0^\natural}\ ,
\]
\DERmark{ENV1505}

since the arrow $\tau$ is an isomorphism. So the direct system reindexed by $J'$, condition 4° is replaced by the stricter condition \fbox{4° bis) $\tau^\natural_{j',j} = \id_{x_0^\natural}$.} But condition 2° then becomes, thanks to 4.15.7.4,
\[
\boxed{\text{2° bis) For every } j\text{, exists } j' \geqslant j \text{ with } \tau_{j',j} = p}\ .
\]
So, we can suppose 1°, 2° bis, 3°, 4° bis (we again denote by $J$, instead of $J'$, the indexing ordered set large relative to $\pi_0$).

Let $K$ be the ordered set defined by

\msp{page 447}
\[
\begin{cases}
\Ob K \;=\; \Ob J \\
\text{if } j,j' \in \Ob J\text{, we have } j \overset{K}{\leqslant} j' \text{ if and only if } j \overset{\mathrm{J}}{=} j'\text{, or } j \overset{J}{<} j' \text{ and } \tau_{j',j} = p\,.
\end{cases}
\]
It is immediate that this is an \textit{order relation} on $J$ (since $p^2 = p$), finer than the initial order relation. The direct system induced on $K$ satisfies conditions 1°, 2°, 3° of 4.15.7.3. On the other hand, $K \longrightarrow J$ is cofinal, hence so is the composite $K \longrightarrow J \longrightarrow I$. It remains to prove that $K$ is large relative to $\pi_0$, hence that every subset $K' \subset K$ such that $\card K' \leqslant \pi_0$ has an upper bound in $K$. It has an upper bound \textit{in} $J$, say $j$. By hypothesis 2 bis, there exists $j' \geqslant j$ such that $\tau_{j'j} = p$, and we can suppose $j' > j$. (Indeed, if $\tau_{j'j} = p$, we see that
% foot: Version 17 décembre 2025
\printedpage{172}
% head: A. Grothendieck : Dérivateurs, Ch. XVIII, Édition des Universités de Montpellier II et Paris VII
\msp{page 448}
the initial form 2°) of condition 2° bis that we have $\tau_{j''j} = p$ for every $j'' \geqslant j'$, so if $j$ is not a final object of $J$, % [sic] par -> pas, misprint
we can suppose $j' > j$. But $J$ has no final object, for otherwise $I$ would have one, a case we had eliminated. So there exists $j'$ such that $j' \overset{K}{>} j$.) Now we have this (which we apply to the case $i \in K'$):
\[
i \leqslant j \ \text{ and } \ j \overset{K}{<} j' \qquad \Longrightarrow \qquad i \overset{K}{<} j'\ .
\]
\[
\underbrace{x_i}_{= x_0} \xrightarrow{\tau_{j,i}} x_j \xrightarrow{\tau_{j',j}=p} x_{j'}\ ,
\]
\[
\tau_{j'i} \;=\; \tau_{j'j}\tau_{ji} = \underbrace{p\tau_{ji}}_{\overset{\text{1°}}{=}\,\tau_{ji}p} \;\overset{?}{=}\; p\ ,
\]
we have indeed
\[
\left.
\begin{aligned}
\tau_{j',i} &\;=\; p\tau_{j',i} \\
\tau_{j',i} &\;=\; \tau_{j',i}p % [sic] - -> =, the second relation of (*) is meant
\end{aligned}
\right\}
\quad \overset{\text{by 4.15.7.4}}{\Longrightarrow} \quad % Doubt resolved against PDF leaf 14: no underline; the two strokes belong to the implication arrow.
\tau_{j',i} \;=\; p\ .
\]
\DERmark{ENV1506}

This completes the proof of corollary 4.15.7.3.

\textbf{Corollary 4.15.7.5} (of 4.15.7.3). \textit{Let $(x_i)_{i\in I}$ be an essentially constant direct system large relative to $\pi_0 \geqslant$} $\card\cC$\textit{, in $\cC$} without projectors. \textit{Then there exists $J \longrightarrow I$ cofinal, with $J$ ordered large relative to $\pi_0$, and such that the direct system induced on $J$ is} constant.

\textsc{Let us prove lemma} 4.15.7.1. We must make explicit the existence of two arrows in $\Ind(\cC)$,
\[
z \underset{\beta}{\overset{\alpha}{\rightleftarrows}} \text{"}\varinjlim_{i\in I}\text{"} x_i \;=\; \ux\ ,
\]
with $\beta\alpha = \id_z$, $\alpha\beta = \id_{\ux}$. The datum

\msp{page 449}
of $\alpha$ is equivalent to that of an element of $\varinjlim\limits_{i\in I} \Hom(z,x_i)$, and is therefore given by an % [sic] x -> \alpha, parallel with "la donnée de \beta" below
\DERmark{ENV1507}

\[
\alpha_{i_0} \;:\; z \longrightarrow x_{i_0}\ ,
\]
the datum of $\beta$ is equivalent to that of an element of $\varprojlim\limits_{i\in I} \Hom(x_i,z)$, i.e. of arrows
\[
\beta_i \;:\; x_i \longrightarrow z\ ,
\]
with the usual compatibility condition. At the cost of replacing $I$ by $I'$, with $I' \longrightarrow I$ cofinal, we can suppose $I$ ordered (and large relative to $\pi_0$, if the initial $I$ was), and the compatibility condition on the $\beta_i$ reads
\[
\beta_j\tau_{ji}=\beta_i \qquad \text{if } i \leqslant j \text{ in } I,
\]
\DERmark{ENV1508}

where $\tau_{ji} : x_i \longrightarrow x_j$ denotes the transition arrow. We must make explicit, in terms of these data, the conditions
\[
\beta\alpha=\mathrm{id}_z,\quad\alpha\beta \;=\; \id_{\ux}\ .
\]
\DERmark{ENV1509}

% foot: Version 17 décembre 2025

\printedpage{173}
% head: A. Grothendieck : Dérivateurs, Ch. XVIII, Édition des Universités de Montpellier II et Paris VII
\[
\begin{tikzcd}
z \arrow[r, "\alpha_{i_0}"] \arrow[dr] & x_{i_0} \arrow[r, "\tau_{i,i_0}"] \arrow[d, "\beta_{i_0}"] & x_i \arrow[dl, "\beta_i"] \\
 & z &
\end{tikzcd}
\]
The relation $\beta\alpha = \id_z$ is equivalent to
\DERmark{ENV1510}

\[
\beta_{i_0}\alpha_{i_0} \;=\; \id_z\ , \tag*{$(*_{i_0})$}
\]
it implies moreover, on setting for $i \geqslant i_0$,
\[
\alpha_i \;\defeq\; \tau_{i,i_0}\alpha_{i_0} \;:\; z \longrightarrow x_i\ ,
\]
that we have for $i \geqslant i_0$
\[
\beta_i\alpha_i \;=\; \id_z \qquad \forall\, i \geqslant i_0\ . \tag*{$(*_i)$}
\]

\msp{page 450}
Let us then set
\[
p_i \;=\; \alpha_i\beta_i \;:\; x_i \longrightarrow x_i\ ,
\]
it is clear that we have
\[
\begin{cases}
p_i^2 \;=\; p_i \qquad \text{($p_i$ a projector)} \\
\alpha_i \;:\; z \longrightarrow x_i \qquad \text{induces an isomorphism} \\
\qquad z \xrightarrow{\sim} x_i^\natural \;\defeq\; \Img p_i \;=\; \Ker(\id_{x_i},p_i)\ .
\end{cases}
\]
The desired cofinal ordered set $J$ is $i_0\backslash I$, it is filtering, and is large relative to $\pi_0$ if $I$ is. To tell the truth, the reduction to the case $I$ ordered is unnecessary, we can take $J = i_0\backslash I$ directly, it is immediate that $J$ is filtering ($I$ being so), and $i_0\backslash I \longrightarrow I$ cofinal, finally that if $I$ is large relative to an infinite cardinal $\pi_0$, the same is true of $i_0\backslash I$.

Let us express that $\alpha\beta = \id_{\ux}$; now we see that $\alpha\beta$ is defined by the endomorphism $(p_j)_{j\in J}$ of the direct system $(x_j)_{j\in J}$ -- the fact that this really is an endomorphism of this system, i.e. that we have relation 1°) of 4.15.7.1,
\DERmark{ENV1511}\DERmark{ENV1512}

\[
p_{j'}x_\lambda \;=\; x_\lambda p_j \qquad \text{if } \lambda : j \longrightarrow j' \text{ in } J, % [sic] p_\lambda -> p_j, cf. relation 1°) of 4.15.7.1
\]
\DERmark{ENV1513}

is immediate, it reads
\[
\alpha_{j'}\underbrace{\beta_{j'}x_\lambda}_{\beta_j} \;=\; \underbrace{x_\lambda\alpha_j}_{\alpha_{j'}}\beta_j\ , \qquad \text{now both sides are equal to } \alpha_{j'}\beta_j.
\]
% foot: Version 17 décembre 2025


\endgroup
% END UNIT der-ch18/der-ch18_p171-175.tex
% BEGIN UNIT der-ch18/der-ch18_p176-180.tex
\begingroup\makeatletter%
\newcommand{\printedpage}[1]{\marginpar{\footnotesize p.~#1}}%
\newcommand{\msp}[1]{\marginpar{\footnotesize [#1]}}%
\newcommand{\card}{\operatorname{card}}%
\makeatother\let\printedpage\derpp\let\msp\dermsp%
\let\mspin\DERmspin%


\printedpage{174}
% head: A. Grothendieck : Dérivateurs, Ch. XVIII, Édition des Universités de Montpellier II et Paris VII

\msp{page 451}
Saying that it is the identity, is expressed by condition 2°) on $(x_j,p_j)_{j\in J}$. Condition 3° is clear by construction, $p_j$ admitting an image isomorphic to $z$ via $\alpha_j$. On the other hand, one sees \texttt{[right]} away that for an arrow $\lambda : j \longrightarrow j'$ of $J$, one has commutativity in % [sic] est clair -> est claire, agreement with "condition"
\[
\begin{tikzcd}
x_j \arrow[r, "x_\lambda"] & x_{j'} \\
x_j^\natural \arrow[u, hook] \arrow[r, "x_\lambda"] & x_{j'}^\natural \arrow[u, hook] \\
z \arrow[u, "\wr", "\overline{\alpha_j}"'] \arrow[r, "\mathrm{id}_z", "\sim"'] & z \ , \arrow[u, "\wr", "\overline{\alpha_{j'}}"']
\end{tikzcd}
\]
where $\overline{\alpha_j}$, $\overline{\alpha_{j'}}$ are the isomorphisms defined by $\alpha_j$, $\alpha_{j'}$ respectively. This proves that $x_\lambda^\natural$ is an isomorphism.

This proves that if $(x_i)_{i\in I}$ is essentially constant, there exists $J\longrightarrow I$ as indicated. The converse is also true, for one can suppose $J = I$, since $(x_i)_{i\in I}$ and $(x_j)_{j\in J}$ define isomorphic ind-objects. Then the direct system of the $x_j^\natural$ has as transition morphisms isomorphisms, and as $J$ is filtering hence 1-connected, one easily deduces from this that it is isomorphic
\msp{page 452}
to a constant direct system, with value $z$. Choosing such an isomorphism
\[
z \xrightarrow[\sim]{\overline{\alpha_i}} x_i^\natural \ ,
\]
one finds $z \xrightarrow{\alpha} \underline{x}$, and the system of the $p_i$ defines some
\[
\begin{tikzcd}
x_i \arrow[r, "\overline{p_i}"] \arrow[rr, bend right=12, "\beta_i \overset{\mathrm{def}}{=} \overline{\alpha_i}^{-1}\overline{p_i}"'] & x_i^\natural \arrow[r, "\overline{\alpha_i}^{-1}"] & z \ ,
\end{tikzcd}
\] % [sic] \overline{\alpha_j} (index j) -> \overline{\alpha_i}, the arrow goes from x_i^\natural, cf. z -> x_i^\natural above
\DERmark{ENV1514}\DERmark{ENV1515}\DERmark{ENV1516}

and one easily sees that one will have
\[
\beta\alpha \ = \ id_z \ , \qquad \alpha\beta \ = \ \mathrm{id}_{\underline{x}} \ ,
\]

\hfill q.e.d.

\textsc{I can now prove} 4.15.3. It is indeed essentially the conjunction of 4.15.7.5 (p. 448), and of 4.8.13, which tells us essentially that if $\mathcal{C}$ is \textit{Karoubian} (a condition satisfied if $\mathcal{C}$ is without projectors), then every direct system in $\mathcal{C}$, large relative to $\pi_0$, where $\pi_0 \geqslant \card\,\text{red}\,\mathcal{C}$, is essentially constant,

\hfill q.e.d.

\printedpage{175}
% head: A. Grothendieck : Dérivateurs, Ch. XVIII, Édition des Universités de Montpellier II et Paris VII

\msp{page 453}
\textbf{4.15.8.} I want to return here to the unfinished proof of 4.8.13 (p. 212-216). I take up the proof where I had left it (p. 215). One must prove that
\[
F(\xi) \ = \ \varinjlim_{i\in I} M_i
\]
has cardinality $\leqslant \pi_0$. I announced that one can reduce to the case where $M$ is a monoid of monomorphisms, i.e. that $uv = uv' \Longrightarrow v = v'$ (for $u,v,v' \in M$). I am going to make this argument explicit.

For every $i\in I$, let
\[
N_i \ = \ \left\{ u \in M \ \middle| \ \begin{array}{l}
\text{the set } I_{i,u} \text{ of } j\in I \text{ such that}\\
j \geqslant i,\ \alpha_{ji} = u \text{ is cofinal}\\
\text{in } I\text{, i.e. for every } j_0\in I,\\
\text{there exists } j \in I_{i,u}\text{, with } j \geqslant j_0
\end{array} \right\} .
\tag*{(4.15.8.1)}
\]
Since $I$, hence $i\backslash I$ is large relative to $\pi_0$, and $\card M \leqslant \pi_0$, one sees at once that
\[
N_i \ \neq \ \emptyset \ .
\tag*{(4.15.8.2)}
\]
$\Big($ If on the other hand $i' \geqslant i$, one finds
\[
N_i \ = \ N_{i'}\alpha_{i'i} \ .
\tag*{(4.15.8.3)}
\]
The inclusion $N_{i'}\alpha_{i'i} \subset N_i$ is $\pm$ tautological, let us prove $N_i \subset N_{i'}\alpha_{i'i}$. Let $u \in N_i$, the set $J = I_{i,u} \cap (i'\backslash I)$ of $j \geqslant i'$ such that $\alpha_{ji} = u$ is cofinal in $I$, hence large relative to $\pi_0$.
\[
\begin{tikzcd}
i \arrow[r] \arrow[rr, bend right=12, "u = \alpha_{ji}"'] & i' \arrow[r, "\alpha_{ji'}"] & j
\end{tikzcd}
\]
Thus there exists a $v \in M$ such that the set $J'$ of $j\in J$ such that $\alpha_{ji'} = v$
\msp{page 454}
is cofinal in $J$, hence in $I$. Since this set $J'$ is contained in $I_{i',v}$, this implies that $I_{i',v}$ is cofinal in $I$, i.e. $v \in N_{i'}$. On the other hand one has, for $j \in J'$,
\[
\underbrace{\alpha_{ji'}}_{v}\alpha_{i'i} \ = \ \underbrace{\alpha_{ji}}_{u} \ ,
\]
whence $u = v\alpha_{i'i}$ with $v\in N_{i'}$, whence $u \in N_{i'}\alpha_{i'i}$, q.e.d. $\Big)$ % [sic] N_i\alpha_{i'i} -> N_{i'}\alpha_{i'i}, the inclusion to be proved is N_i \subset N_{i'}\alpha_{i'i}
\DERmark{ENV1517}


For $i \in I$, let
\[
\alpha_i \ : \ \underbrace{\mathrm{Hom}(\xi, \underbrace{x_i}_{=\,\xi})}_{=\,M} \ \longrightarrow \ \underbrace{\mathrm{Hom}(\xi,F) \ = \ F(\xi)}_{\simeq\ \varinjlim\limits_{i} \mathrm{Hom}(\xi,x_i)}
\]
\printedpage{176}
% head: A. Grothendieck : Dérivateurs, Ch. XVIII, Édition des Universités de Montpellier II et Paris VII
be the canonical morphism. One finds this:
\begin{center}
\makebox[0.20\linewidth][l]{(4.15.8.4)}%
\begin{minipage}[c]{0.78\linewidth}
If $u \in N_i$, then the equivalence relation $R_i$ on $M$ defined by $\alpha_i$ ($v \overset{R_i}{\sim} v' \Longleftrightarrow \alpha_i(v) = \alpha_i(v')$) is equal to that defined by the map $v \longmapsto uv$ from $M$ to $M$, i.e.
\[
\alpha_i(v) \ = \ \alpha_i(v') \qquad \Longleftrightarrow \qquad uv \ = \ uv' \ .
\]
\end{minipage}
\end{center}

This is immediate, since $I_{i,u}$ being cofinal in $I$, one has $\alpha_i(v) = \alpha_i(v')$ if and only if $\exists\ j \in I_{i,u}$ with $\alpha_{ji}v = \alpha_{ji}v'$, now $\alpha_{ji} = u$. % [sic] C'est immédiate -> C'est immédiat, agreement

For $u \in M$, let us denote by $R_u$ the equivalence relation on $M$ defined by $v \longmapsto uv$. The set of equivalence relations on \texttt{[}$M$\texttt{]}
\msp{page 455}
obtained in this way has cardinality $\leqslant \card M \leqslant \pi_0$, hence the set of equivalence relations $R_i$ ($i \in I$) also has cardinality $\leqslant \pi_0$. This shows, $I$ being large relative to $\pi_0$, that there exists an equivalence relation of the form $R_u$ ($u \in M$), and a cofinal set $I'$ in $I$, such that $\forall\ i \in I'$ one has $R_i = R_u$. This shows that, at the cost of replacing $I$ by the cofinal ordered set $I'$, that one can suppose that one has
\DERmark{ENV1518}

\[
R_i \ = \ R_u \qquad \forall\, i \in I \ .
\tag*{(4.15.8.5)}
\]
Let then
\[
E \ = \ M/R_u \ = \ M/R_i \qquad (\forall\ i \in I) \ ,
\]
one therefore sees that the $\alpha_{ji} \in M$ (for $i \leqslant j$) are such that the map
\[
\tau_{\alpha_{ji}} \ : \ v \ \longmapsto \ \alpha_{ji}v \ : \ M \ \longrightarrow \ M
\]
is compatible with the equivalence relation $R_u$, hence passes to the quotient as a map
\[
\overline{\alpha_{ji}} \ : \ E \ \longrightarrow \ E \ ,
\]
\DERmark{ENV1519}

which is moreover \textit{injective.}

Let $M'$ be the submonoid of $M$ formed of the $u \in M$ which have the preceding stability property, as well as the injectivity property we have just stated, and $\overline{M'}$ the image monoid of $M'$ in the monoid
\[
\mathrm{Mon}(E) \ \subset \ \mathrm{End}_{\mathrm{Ens}}(E)
\]
\msp{page 456}
of injective maps from $E$ to itself. One therefore sees that the given direct system, defined by the $\alpha_{ji}$, takes its values in the submonoid $M'$ of $M$ (i.e. of the corresponding subcategory of the category $B_M$), and that the direct limit $F(\xi)$ which interests us (to prove $\card F(\xi) \leqslant \pi_0$) is isomorphic to the direct limit of the direct system of a system of sets all equal to $E$, the transition morphisms being the $\overline{\alpha_{ji}} \in \overline{M'} \subset \mathrm{Mon}(E)$. Moreover, one will have $\card \overline{M'} \leqslant \card M \leqslant \pi_0$. The desired result $\card F(\xi) \leqslant \pi_0$ will then follow from the

\textbf{Key lemma 4.15.8.1.} \textit{Let $M$ be a monoid of monomorphisms (i.e. such that $uv = uv'$ $\Longrightarrow v = v'$, for $u,v,v' \in M$), and let $(\alpha_{ji})_{i\leqslant j}$ be a direct system in $B_M$, indexed by an ordered set $I$ large relative to $\pi_0$, where $\pi_0$ is a cardinal}
\printedpage{177}
% head: A. Grothendieck : Dérivateurs, Ch. XVIII, Édition des Universités de Montpellier II et Paris VII
\msp{page 457}
\textit{which is infinite and such that} $\card M \leqslant \pi_0$. \textit{Then there exists an $i_0 \in I$ such that for every $i \geqslant i_0$, $\alpha_{i,i_0}$ is an isomorphism.}
\DERmark{ENV1520}


If one applies this lemma to the case of $\overline{M'}$ and of the $\overline{\alpha_{ji}}$, one therefore finds that there exists an $i_0 \in I$ such that the $\overline{\alpha_{ji_0}}$ are isomorphisms. This implies that $F(\xi) \simeq \varinjlim\limits_{i} \underbrace{E_i}_{E}$ is isomorphic to $E_{i_0}$, hence has cardinality equal to $\card E_{i_0} = \card E \leqslant \pi_0$, which will complete the proof.

Moreover, statement 4.8.13 implies 4.15.8.1, for in the direct system of the $(M_i = M, \alpha_{ji})$, with $\varinjlim$ equal to $F(\xi)$, if there did \textit{not} exist an $i_0$ as stipulated in 4.15.8.1, then one would easily conclude from this that one would have $\card F(\xi) > \pi_0$, contrary to 4.8.13. We are going to bring this out in a lemma:

\msp{page 458}
\textbf{Lemma 4.15.8.2.} \textit{Let $(E_i)_{i\in I}$ be a direct system of sets, indexed by an ordered set $I$ large relative to $\pi_0$. We suppose that there exists $i_0 \in I$ such that}
\begin{enumerate}[nosep]
\item[\textit{a)}] \textit{$i \geqslant i_0$ implies} $\card E_i = \pi_0$, \textit{and}
\item[\textit{b)}] \textit{$j \geqslant i \geqslant i_0$ implies that $u_{ji} : E_i \longrightarrow E_j$ is a monomorphism. Suppose moreover}
\item[\textit{c)}] \textit{$\forall\ i \geqslant i_0$ $\exists\ j > i$ such that $u_{ji}$ is} not \textit{an isomorphism (i.e. is not surjective). Let $E_\infty = \varinjlim\limits_{i\in I} E_i$. Then} $\card E_\infty > \pi_0$.
\end{enumerate}
This lemma therefore implies the \textit{equivalence} of 4.8.13 and the ``key lemma'' 4.15.8.1. To prove 4.15.8.2, we are going to prove something a little more precise:

\textbf{Corollary 4.15.8.3.} \textit{Under the conditions of 4.15.8.2, let $\alpha_0$ be the first ordinal such that} $\card \alpha_0 > \pi_0$, \textit{and let $I_{\alpha_0}$ be the set of ordinals $\beta < \alpha$ (an ordered set which is large relative to $\pi_0$). There then exists a strictly increasing map}
\[
I_{\alpha_0} \ \longrightarrow \ I \ , \qquad \alpha \ \longmapsto \ i_\alpha \ , \qquad \textit{mapping } 0 \textit{ to } i_0 \ ,
\]
\msp{page 459}
\textit{such that for every $\alpha \in I_{\alpha_0}$, one has}
\[
\underbrace{u_{i_{\alpha+1},i_\alpha}}_{\substack{\textit{it is a } \text{monomorphism,}\\ \textit{by b, since } i_{\alpha+1} > i_\alpha \geqslant i_0}} \qquad\qquad \text{\textit{is not an isomorphism} (\footnotemark[187]).}
\]
\footnotetext[187]{\textit{better}: $\forall\ \alpha \in I_{\alpha_0}$ $\varinjlim\limits_{\beta\in I_\alpha} E_{i_\beta} \longrightarrow E_{i_\alpha}$ is not surjective.}
\textit{One then has}
\[
E'_\infty \ \overset{\mathrm{def}}{=} \ \underbrace{\varinjlim_{\alpha\in I_{\alpha_0}} E_{i_\alpha}}_{} \ \hookrightarrow \ \varinjlim_{i\in I} E_i \ \overset{\mathrm{def}}{=} E_\infty
\]
\DERmark{ENV1521}\DERmark{ENV1522}

\textit{and}
\[
\card E'_\infty \ \geqslant \ \underbrace{\card I_{\alpha_0}}_{=\ \textit{successor of } \pi_0} \ > \ \pi_0 \ ,
\] % [sic] E'_\alpha -> E'_\infty, the following line has \card E'_\infty
\DERmark{ENV1523}

\textit{whence} $\card E_\infty \geqslant \card E'_\infty > \pi_0$.

\textsc{Proof}. The construction of the $i_\alpha$ by transfinite induction is immediate, using c) and the fact that $I$ is large relative to $\pi_0$, and that for every $\alpha \in I_{\alpha_0}$, the set $I_\alpha$
\printedpage{178}
% head: A. Grothendieck : Dérivateurs, Ch. XVIII, Édition des Universités de Montpellier II et Paris VII
of $\beta < \alpha$ has cardinality $\leqslant \pi_0$. The fact that $E'_\infty \longrightarrow E_\infty$ is injective follows from b). The fact that
\[
\card E'_\infty \ \geqslant \ \card I_{\alpha_0}
\]
follows from the fact that $E'_\infty$ is an increasing union of subsets identified with $E_{i_\alpha}$, and for every
\msp{page 460}
$\alpha$ one can find a
\[
\xi_\alpha \ \in \ E_{i_\alpha}\backslash \bigcup_{\beta<\alpha_0} E_{i_\beta} \ .
\]
\DERmark{ENV1524}

The map
\[
\begin{aligned}
\alpha \ &\longmapsto \ \xi_\alpha\\
I_{\alpha_0} \ &\longrightarrow \ E'_\infty
\end{aligned}
\]
is injective, hence $\card E'_\infty \geqslant \card I_{\alpha_0}$.

% vertical bar printed in the left margin beside the whole of Corollaire 4.15.8.4
\textbf{Corollary 4.15.8.4.} \textit{For given $M$ and $\pi_0$, for the key lemma 4.15.8.1 to be valid, it is necessary and sufficient that it be so in the case where $I = I_{\alpha_0}$, $\alpha_0$ being the first ordinal such that} $\card \alpha_0 > \pi_0$. \textit{It amounts to the same thing that one cannot find a direct system $(\alpha_{ji})_{j\geqslant i}$ indexed by $I_{\alpha_0}$, with values in $M$, such that for $i < j$, one has $\alpha_{ji}$ noninvertible in $M$.} % [sic] \card \alpha_0 > \alpha_0 -> \card \alpha_0 > \pi_0, cf. Corollaire 4.15.8.3
\DERmark{ENV1525}


\msp{page 461}
I am going to summarize what has been obtained in a summary statement:

\textbf{Theorem 4.15.9.} \textit{Let $\mathcal{C}$ be a small category, $\pi$ a cardinal such that $\pi \geqslant$} $\card\mathcal{C}$. \textit{The following conditions are equivalent (and I conjecture that they are always satisfied):}
\begin{enumerate}[nosep]
\item[\textit{a)}] \textit{$\mathcal{M} \overset{\mathrm{def}}{=} \mathrm{Ind}_\pi\mathcal{C}$ is reduced to $\mathcal{M}_\pi$.}
\item[\textit{a$'$)}] \textit{(If $\mathcal{C}$ is Karoubian.) Every direct system $\underline{x} = (x_i)_{i\in I}$ in $\mathcal{C}$, large relative to $\pi$, is essentially constant.}
\item[\textit{a$''$)}] \textit{For every direct system $(x_i)_{i\in I}$ in $\mathcal{C}$, large relative to $\pi$, there exists a cofinal functor $J \longrightarrow I$, with $J$ an ordered set large relative to $\pi$, and such that the reindexed system $(x_j)_{j\in J}$ satisfies the following conditions:}
\begin{enumerate}[nosep]
\item[\textit{1)}] \textit{All the $x_j$ are equal to one and the same element $x$ of $\mathcal{C}$.}
\item[\textit{2)}] \textit{There exists a projector $p$ in $x$, such that one has $\alpha_{j'j} = p$ for $j < j'$ ($j,j' \in J$), where $\alpha_{j'j} : x_j \longrightarrow x_{j'}$ is the transition arrow.}
\end{enumerate}
\DERmark{ENV1526}

\item[\textit{b)}] \textit{For every $x \in \mathrm{Ob}\,\mathcal{C}$, and for every quotient monoid $M'$ of $M = \mathrm{End}_{\mathcal{C}}(x)$, such that $M'$ is a monoid of monomorphisms (i.e. $uv = uv' \Longrightarrow v = v'$ if $u,v,v' \in M'$), and for every direct system large relative to $\pi$ with values in $B_{M'}$ (defined by a system of elements $(\alpha_{ji})_{i\leqslant j,\ i,j\in I}$ of an ordered set $I$ large relative to $\pi$, with $\alpha_{kj}\alpha_{ji} = \alpha_{ki}$ for $i \leqslant j \leqslant k$), there exists $i_0 \in I$ such}
\msp{page 462}
\textit{that $\alpha_{ii_0}$ is invertible for $i \geqslant i_0$, i.e. $\alpha_{ji}$ is invertible for $j \geqslant i \geqslant i_0$.}
\DERmark{ENV1527}

\item[\textit{b$'$)}] \textit{$M'$ being a monoid of monomorphisms as in b) (hence} $\card M' \leqslant \card\mathcal{C} \leqslant$ \textit{$\pi$), let moreover $\alpha_0$ be the smallest ordinal such that} $\card\alpha_0 > \pi$, \textit{and let $I = I_{\alpha_0}$ be the ordered set of ordinals $\alpha < \alpha_0$ (which set is large relative to $\pi$). There exists} no \textit{direct system in $B_{M'}$, such that for $i < j$ in $I$, $\alpha_{ji}$ is a noninvertible element in $M'$ (hence $u \longmapsto \alpha_{ji}u : M' \longrightarrow M'$ is injective but not surjective).}
\end{enumerate}


\endgroup
% END UNIT der-ch18/der-ch18_p176-180.tex
% BEGIN UNIT der-ch18/der-ch18_p181-185.tex
\begingroup\makeatletter%
\newcommand{\printedpage}[1]{\marginpar{\footnotesize p.~#1}}%
\newcommand{\msp}[1]{\marginpar{\footnotesize [#1]}}%
\newcommand{\card}{\operatorname{card}}%
\newcommand{\redop}{\mathop{\textnormal{red}}\nolimits}%
\newcommand{\Hom}{\operatorname{Hom}}%
\newcommand{\Ob}{\operatorname{Ob}}%
\newcommand{\Ind}{\operatorname{Ind}}%
\newcommand{\Ens}{\mathrm{Ens}}%
\newcommand{\defeq}{\overset{\mathrm{def}}{=}}%
\newcommand{\into}{\lhook\joinrel\longrightarrow}%
\makeatother\let\printedpage\derpp\let\msp\dermsp%
\let\mspin\DERmspin%

\DERmark{ENV1528}


\printedpage{179}
% head: A. Grothendieck : Dérivateurs, Ch. XVIII, Édition des Universités de Montpellier II et Paris VII
% foot: Version 17 décembre 2025

Let us point out at once the

\textbf{Corollary 4.15.10.} \textit{The preceding equivalent conditions are satisfied, if $\pi > \card\mathcal{C}$.}

Indeed, let $\alpha_1$ be the smallest ordinal such that $\card\alpha_1 > \card\mathcal{C}$. Thus $\pi_1 \defeq \card\alpha_1$ is the successor of $\pi_0 \defeq \card\mathcal{C}$, and we have $\pi \geqslant \pi_1$, hence (if $\alpha_0$ is as specified in b$'$ above) $\card\alpha_0 > \pi_1$, hence $\alpha_1 < \alpha_0$. Let (with the notation of b$'$), for $\alpha \in I$ \texttt{[rather} $\alpha \in I_{\alpha_0}$\texttt{]}, i.e.\ $\alpha < \alpha_0$,
\[
M_\alpha \;=\; \varinjlim_{i<\alpha} \underbrace{M_i}_{=\,M}\ ,
\]
we thus see, since the $\alpha_{ji}$ are injective and not surjective, that
\[
\card M_\alpha \;\geqslant\; \card\alpha\ ,
\]
\msp{463}
and that for $\alpha < \beta$, we have
\[
M_\alpha \;\into\; M_\beta \qquad \text{injective}\ .
\]
Let us apply this to the case where $\alpha = \alpha_1$, thus we find
\[
\card M_{\alpha_1} \;\geqslant\; \card\alpha_1 \;=\; \pi_1\ ,
\]
on the other hand, taking $\beta = \alpha_1 + 1$, whence $M_\beta \simeq M$, we find an injection
\[
M_{\alpha_1} \;\into\; M\ ,
\]
whence
\[
\pi_1 \;\leqslant\operatorname{card}M_{\alpha_1}\leqslant\operatorname{card}M\leqslant\; \pi_0\ ,
\]
\DERmark{ENV1529}

a contradiction, which establishes b$'$.

\textbf{Corollary 4.15.11.} \textit{Let $\mathcal{C}$ be an essentially small category, $\pi$ a cardinal such that $\pi > \card\redop\mathcal{C}$. Let $\mathcal{M} \defeq \Ind_\pi\mathcal{C}$, then $\mathcal{M} = \mathcal{M}_\pi$, and it is a Karoubi envelope of $\mathcal{C}$.}

We may assume $\mathcal{C}$ reduced, hence $\card\mathcal{C} < \pi$, and we apply the preceding corollary.

Thus, the disputed result 4.8.13 (p.~212) is finally proved in a slightly weaker form, with $\pi > \card\redop\mathcal{C}$ instead of $\pi \geqslant \card\redop\mathcal{C}$. By the same token, the applications of 4.8.13 we have made so far (in particular lemma 4.15.3).

\msp{464}

\subsection*{4.16\quad Examples and questions}

\textbf{4.16.0.} I had thought that every ``reasonable'' functor (or more precisely, every ``logically accessible'' functor) between accessible categories ought to be accessible. This is unfortunately not true, since a composite of two Hom functors
\[
\Ens \longrightarrow \Ens\ , \qquad\quad E \longmapsto \underbrace{\Hom(\Hom(E,I),J)}_{J^{(I^E)}}
\]
\printedpage{180}
% head: A. Grothendieck : Dérivateurs, Ch. XVIII, Édition des Universités de Montpellier II et Paris VII
% foot: Version 17 décembre 2025
is not accessible, if $I$, $J$ are both of cardinality $\geqslant 2$. Taking e.g.\ $I = J = \{0,1\}$, we find the functor
\[
E \longmapsto \mathfrak{P}(\mathfrak{P}(E))\ ,
\]
which is not accessible.

\textbf{4.16.1.} Second disillusionment: I had thought that for every accessible category $\mathcal{M}$, every % [sic] toute catégories -> toute catégorie, singular required after "toute"
strictly full subcategory $\mathcal{N}$ stable under direct factors and ``reasonable'' (e.g. ``logically accessible'') ought to be accessible in $\mathcal{M}$. This is not so, even in the particular case of the inverse image of the accessible subcategory in Ens
\msp{465}
\[
\underbrace{\Ens^*}_{=\,\{E\,\in\Ens\,\mid\, E\neq\emptyset\}} \quad\subset\ \Ens
\]
under an accessible \textit{representable} contrafunctor
\[
\begin{aligned}
\mathcal{M} \quad &\longrightarrow \quad (\Ens)^\circ \\
x \quad &\longmapsto \quad \Hom_{\mathcal{M}}(x,z)\ ,
\end{aligned}
\]
i.e.\ the open (hence strictly full) subcategory $\mathcal{M}_z$ of $\mathcal{M}$ consisting of the objects $x$ such that $\Hom_{\mathcal{M}}(x,z) \neq \emptyset$, is not always accessible in $\mathcal{M}$. Here, I have not arrived at % [sic] n'est par toujours -> n'est pas toujours, misprint "par" for "pas"
as elementary an example (cf.\ appendix) -- I took $\mathcal{M} = (k\text{-comm.\ alg.})$, $k$ an infinite field and $z$ the initial object $k$ of $\mathcal{M}$ -- I found that $\mathcal{M}_z$ is not $\mathrm{L}_\pi$-stable in $\mathcal{M}$, whatever the cardinal $\pi$ % [sic] quelque soit -> quel que soit, standard spelling of the locution
-- but this example used that $\mathfrak{U} = \mathfrak{U}_1$, the smallest universe containing $\mathbf{N}$. I do not know whether for every other universe, one can find a field $k$ of cardinality $\pi \in \mathfrak{U}$ large enough, so that every $\pi$-ultrafilter on a set $E \in \mathfrak{U}$ is trivial. Perhaps
\msp{466}
such an example, where not the construction of the pair $(\mathcal{M},\mathcal{N})$, but the proof of the fact that $\mathcal{N}$ is not accessible in $\mathcal{M}$, appeals to particular hypotheses on the ambient universe, should not be considered as (altogether!) ``logically accessible''. To be really convincing, one would need an example of a more elementary nature.

% margin: a vertical bar is printed in the left margin beside the three lines from "L_pi-stable dans M pour pi assez grand" to "logiquement accessible) serait accessible dans M"
There remains the possibility that there is a metamathematical theorem ensuring that every logically accessible pair $(\mathcal{M},\mathcal{N})$, with $\mathcal{M}$ accessible, and $\mathcal{N}$ a strictly full subcategory $\mathrm{L}_\pi$-stable in $\mathcal{M}$ for $\pi$ large enough, is accessible in $\mathcal{M}$ (\footnote[188]{This is false, even for a closed subcategory, cf.\ 4.16.7.10 below.}). This would apply in particular to the case of a \textit{closed} subcategory: every closed subcategory (that is logically accessible) would be accessible in $\mathcal{M}$. The simplest test example would be that of the closed complement $\mathcal{M}'_z$ of the open subcategory $\mathcal{M}_z$ above (\footnote[189]{\textbf{N.B.} it is accessible in $\mathcal{M}$ if and only if $\mathcal{M}_z$ is! % [sic] elles est -> elle est, agreement with singular verb
}), i.e.\ the category of the $x \in \Ob\mathcal{M}$ such that $\Hom_{\mathcal{M}}(x,z) = \emptyset$, or again
\msp{467}
the inverse image of the accessible closed subcategory
\[
e \;=\; \{\emptyset\} \;\subset\; (\Ens)
\]
reduced to $\{\emptyset\}$, under the representable functor $x \longmapsto \Hom(x,z) : \mathcal{M} \longrightarrow \Ens^{\circ}$. It could thus be that such a category is always accessible in $\mathcal{M}$. But I have been unable to
\DERmark{ENV1530}

\printedpage{181}
% head: A. Grothendieck : Dérivateurs, Ch. XVIII, Édition des Universités de Montpellier II et Paris VII
% foot: Version 17 décembre 2025
prove it, even in the particular case above, where $\mathcal{M} = (k\text{-comm.\ alg.})$, $k$ being an infinite commutative field. The difficulty is (in all cases) the following: to find, for fixed $z$, a cardinal $\pi > \pi_0$ such that for every $x \in \Ob\mathcal{M}$, and every system (or \textit{one} system) $\pi_0$-$\pi$-adapted $(x_i)_{i\in I}$ with direct limit $x$ ($\pi_0$ a given cardinal, useful for $\mathcal{M}$), \textit{if} $\Hom_{\mathcal{M}}(x,z)$ ($\simeq \varprojlim \Hom_{\mathcal{M}}(x_i,z)$) is empty, then $\exists\ i_0 \in I$ such that $\Hom_{\mathcal{M}}(x_{i_0},z)$ is already empty. In the example $\mathcal{M} = (k\text{-comm.\ alg.})$, one must therefore prove that one can find a cardinal $\pi$
\DERmark{ENV1531}

such that for every $k$-algebra $A$ with $\Hom_{\mathcal{M}}(A,k) = \emptyset$, there exists a sub-$k$-algebra $A_i$ of % [sic] pour tout k-algèbre -> pour toute k-algèbre, "algèbre" is feminine
cardinality $\leqslant \pi$, such that $\Hom_k(A_i,k) = \emptyset$ (\footnote[190]{This is false, cf.\ marginal annotation p.~466.}). I tend to believe that this must be true (and even without
\msp{468}
assuming that $k$ is a field, and also taking the condition $\Hom_{\mathcal{M}}(A,Z) = \emptyset$, where $Z$ is an arbitrary object of $\mathcal{M}$). This sort of question seems surprisingly delicate \dots % [sic] étonnement -> étonnamment, adverb required

\textbf{4.16.2.} It is possible, more generally, that the following question has an affirmative answer: Let
\[
f \;:\; \mathcal{M} \longrightarrow \mathcal{N}^\circ
\]
be an accessible functor, with $\mathcal{M}$, $\mathcal{N}$ accessible categories, and let $\underline{T}$ be a subcategory of $\mathcal{N}$ accessible in $\mathcal{N}$, and such that there exists a cardinal $\pi$ having the following properties: $\mathcal{N}^\circ$ satisfies $\mathrm{L}_\pi$ (i.e.\ $\mathcal{N}$ stable under $\varprojlim$ large relative to $\pi$), and $\underline{T}$ is stable in $\mathcal{N}$ under these same limits. Then is $\underline{S} \defeq f^{-1}(\underline{T})$ accessible in $\mathcal{M}$ (\footnote[191]{False, with $\mathcal{N} = \Ens$ and a representable functor $\mathcal{M} \longrightarrow \Ens^\circ$, taking $\underline T = \{\emptyset\}$, cf.\ marginal annotations p.~466, 467.})? This generalizes the case $\mathcal{N} = (\Ens)$, $\underline{T} = \{\emptyset\}$, $f$ a functor represented by $z \in \Ob\mathcal{M}$, considered just now. With the same type of example, but taking $\underline T = (\Ens)^*$
\DERmark{ENV1532}\DERmark{ENV1533}\DERmark{ENV1534}

\msp{469}
instead of $\underline T = \{\emptyset\}$, we see that the additional stability condition on $\underline{T}$ is not
\DERmark{ENV1535}

superfluous. (\textbf{N.B.} It serves to ensure that $\underline{S}$ is $\mathrm{L}_\pi$-stable in $\mathcal{M}$ for large $\pi$.) % [sic] Elle sont -> Elles sont, agreement with plural verb
\DERmark{ENV1536}


\textbf{4.16.3.} Related question: I had hoped that for an accessible functor between accessible categories % [sic] Questions apparentée -> Question apparentée, agreement
\DERmark{ENV1537}

\[
f \;:\; \mathcal{N} \longrightarrow \mathcal{M}\ ,
\]
the $\natural$-essential image, consisting of the $y \in \Ob\mathcal{M}$ which are isomorphic to a direct factor of an object of the form $f(x)$, with $x \in \Ob\mathcal{N}$, would be accessible in $\mathcal{M}$. The example above (not at all elementary!) of the subcategory $\mathcal{M}_z$ of $\mathcal{M}$ (consisting of the $x \in \Ob\mathcal{M}$ such that $\Hom_{\mathcal{M}}(x,z) \neq \emptyset$) shows that this is not so. Indeed, $\mathcal{M}_z$ is the essential image, and also the $\natural$-essential image, of the functor
\DERmark{ENV1538}

\[
\mathcal{M}/z \longrightarrow \mathcal{M}\ ,
\]
now $\mathcal{M}/z$ is accessible, and the preceding functor accessible, if $\mathcal{M}$ is accessible (verification by AQT \texttt{[âne qui trotte: at a donkey's trot, i.e. by routine verification]}, cf.\ below, more general statement).

\msp{470}

We get the impression that the reason why we have such counterexamples, is that we have $\natural$-essential images which are not $\mathrm{L}_\pi$-stable in $\mathcal{M}$ for any cardinal $\pi \in \mathfrak{U}$. This leads to modifying the conjecture as follows: let $\pi$ be a cardinal such that $\mathcal{M}$ satisfies $\mathrm{L}_\pi$,
\printedpage{182}
% head: A. Grothendieck : Dérivateurs, Ch. XVIII, Édition des Universités de Montpellier II et Paris VII
% foot: Version 17 décembre 2025
and let $\underline{S}(\pi)$ be the smallest strictly full subcategory of $\mathcal{M}$, $\mathrm{L}_\pi$-stable in $\mathcal{M}$, and containing the essential image of $f$. (One could call $\underline{S}(\pi)$ the $\mathrm{L}_\pi$-essential image.) Then $\underline{S}(\pi)$ is accessible in $\mathcal{M}$. Moreover, the (decreasing) family of the $\underline{S}(\pi)$, for varying $\pi$, is stationary, i.e.\ there exists a cardinal $\pi$ such that $\pi' \geqslant \pi$ implies $\underline{S}(\pi') = \underline{S}(\pi)$. It is then clear that $\underline{S}(\pi)$ is the smallest subcategory of $\mathcal{M}$ that is accessible in $\mathcal{M}$ and contains the essential image of $f$. The same question arises for a functor $g^\circ : \mathcal{N}^\circ \longrightarrow \mathcal{M}$, where $g : \mathcal{N} \longrightarrow \mathcal{M}^\circ$ is an accessible functor, $\mathcal{N}$ and $\mathcal{M}$ accessible. I conjecture that the answer is again affirmative.

\msp{471}

% margin: a hand-drawn curved stroke is printed in the left margin beside the last three lines of the following paragraph ("pleine M de C^ qui soit ..." to "sans pourtant être accessible.")
Let $\mathcal{M}$ be a category satisfying condition L, and admitting a small strictly generating full subcategory $\mathcal{C}$, consisting of accessible objects. We thus have a fully faithful and accessible functor
\[
\mathcal{M} \;\into\; \mathcal{C}^\wedge\ ,
\]
$\mathcal{M}$ $\mathrm{L}_\pi$-stable in $\mathcal{C}^\wedge$ for suitable $\pi$. If $\overline{\mathcal{M}}$ is the essential image, so that $\mathcal{M} \simeq \overline{\mathcal{M}}$, it is thus a strictly full subcategory $\mathrm{L}_\pi$-stable in $\mathcal{C}^\wedge$. So we expect it to be accessible in $\mathcal{C}^\wedge$ (if the situation described is ``logically accessible''), or what amounts to the same, $\mathcal{M}$ to be accessible. (Compare 4.8.1, p.~187.) But it should not be difficult (\footnote[192]{OK by taking a \textit{closed} $\underline{S}$ not accessible in $\mathcal{M}$ of an accessible category $\mathcal{M}$, which we embed in a $\mathcal{C}^\wedge$. But these $\underline{S}$ do \textit{not} have a small strictly generating family! % [sic] ces S n'a -> ce S n'a, plural determiner with singular verb
}), by making a number $\pi_{\mathfrak{U}}$ of choices (where $\pi_{\mathfrak{U}} = \card\mathfrak{U}$), to make examples all the same (``logically inaccessible'') of a strictly full subcategory $\mathcal{M}$ of $\mathcal{C}^\wedge$ that is $\mathrm{L}_\pi$-stable, but not accessible in $\mathcal{C}^\wedge$. This $\mathcal{M}$ thus admits a small strictly generating subcategory consisting of accessible objects (\footnote[193]{This is not clear, for this the inclusion functor would have to be a left adjoint!}) and satisfies $\mathrm{L}_\pi$, without however being accessible.

\msp{472}

\textbf{4.16.4} (\footnote[194]{move down after 4.16.5.}). Let $\mathcal{M}$ be a category such that $\mathcal{M}$ and $\mathcal{M}^\circ$ are accessible. I conjecture that then $\mathcal{M}$ is essentially small. More generally, suppose $\mathcal{M}$ accessible, and that $\mathcal{M}$ admits a small subcategory which is \textit{strictly co}generating. Then is $\mathcal{M}$ essentially small? More generally, if $\mathcal{M}$ admits a small strictly generating subcategory, and also a small strictly cogenerating subcategory, i.e.\ if $\mathcal{M}$ and $\mathcal{M}^\circ$ are both equivalent to full subcategories of a $\mathcal{C}^\wedge$, with $\mathcal{C}$ small, then is $\mathcal{M}$ essentially small?

\textbf{4.16.5.} Let $I$ be a small category, $I_\bullet$ the right cone on $I$ (obtained by adding a final object $e$), $\mathcal{M}$ an accessible category,
\[
\underline{\Hom}_{\square}(I_\bullet,\mathcal{M}) \;\into\; \underline{\Hom}(I_\bullet,\mathcal{M})
\]
the strictly full subcategory of the (accessible) category $\underline{\Hom}(I_\bullet,\mathcal{M})$, consisting of the functors $f$ which make $f(e)$ the direct limit in $\mathcal{M}$ of the $f(i)$ ($i \in \Ob I$). This subcategory is $\mathrm{L}_\pi$-stable for large $\pi$, so there is reason
\msp{473}
to conjecture that it is accessible in $\underline{\Hom}(I_\bullet,\mathcal{M})$. $\Big($ I recall that we saw in 4.14 that this is OK if the $\varinjlim$ of type $I$ are representable in $\mathcal{M}$. And that an affirmative answer, in the case where
\printedpage{183}
% head: A. Grothendieck : Dérivateurs, Ch. XVIII, Édition des Universités de Montpellier II et Paris VII
% foot: Version 17 décembre 2025
\[
\begin{tikzcd}[row sep=small, column sep=small]
 & 1 \\
I = 0 \arrow[ur] \arrow[dr] & \\
 & 2
\end{tikzcd}
\quad ,\qquad\quad \text{thus} \qquad\quad
\begin{tikzcd}[row sep=small, column sep=small]
 & 1 \arrow[dr] & \\
I_\bullet = 0 \arrow[ur] \arrow[dr] & & 3 \\
 & 2 \arrow[ur] &
\end{tikzcd}
\quad ,
\]
\DERmark{ENV1539}

implies that the subcategory $\underline{\text{Ép}}(\mathcal{M})$ in $\underline{\mathrm{Fl}}(\mathcal{M})$ is accessible in $\underline{\mathrm{Fl}}(\mathcal{M})$.$\Big)$

\textbf{N.B.} If $\mathcal{M}$ admits a small \textit{co}generating subcategory, then the question falls under conjecture 4.16.2 concerning an accessible functor $\mathcal{M} \longrightarrow \mathcal{N}^\circ$, here $\underline{\Hom}(I_\bullet,\mathcal{M}) \longrightarrow \underline{\Hom}(I^\circ \times \mathcal{C},\Ens)^\circ \dots$

\textbf{4.16.6.} There remains the question of the relative accessibility of a category $\mathcal{E}$ over a small category $\mathcal{C}$. I now think, thanks to 4.15.9, and its corollary 4.15.10, that the following statement must ``fall out'' by AQT \texttt{[âne qui trotte: at a donkey's trot, i.e. by routine verification]}:

\textbf{Proposition 4.16.6.1.} \textit{Let $f : \mathcal{E} \longrightarrow \mathcal{C}$ be a functor, with $\mathcal{C}$ a small Karoubian category. We assume that $f$ satisfies condition $H' = H$ (1+2+3) of 4.15.2 (page 415), where H 2, H 3 suffice, for every arrow $\lambda$} other than an identity \textit{of $\mathcal{C}$. We assume moreover that for a category $\mathcal{C}'$ of one of the two types we are going to describe, and a functor $\mathcal{C}' \longrightarrow \mathcal{C}$, $\mathcal{E}' = \mathcal{E} \times_{\mathcal{C}} \mathcal{C}'$ is accessible, and $\mathcal{E}' \longrightarrow \mathcal{C}'$ accessible.}
\DERmark{ENV1540}


\msp{474}

$\Big($ \textit{We take for $\mathcal{C}'$ the category}
\[
\mathcal{Q} = \left\{ \quad
\begin{tikzcd}[column sep=2.2em]
\xi_0 = \xi^\natural \arrow[r, shift left, "\sigma"] \arrow[dr, "u_0"'] & \xi \arrow[l, shift left, "\rho"] \arrow[d, "u = u_0\rho"] \arrow[loop right, "p"] \\
 & \eta
\end{tikzcd}
\qquad
\begin{array}{rclll}
\rho\sigma & = & \mathrm{id}_{\xi_0} & & \\[2ex]
\sigma\rho & = & p\ , & \text{\textit{whence}} & \left\{\begin{array}{rcl} p^2 & = & p \\ p\sigma & = & \sigma \\ \rho p & = & \rho \end{array}\right. \\[4ex]
u & = & u_0\rho\ , & \text{\textit{whence}} & \left\{\begin{array}{rcl} up & = & u \\ u\sigma & = & u_0 \end{array}\right.
\end{array}
\right.
\]
\textit{(three objects $\xi$, $\xi_0$, $\eta$ and five nonidentity arrows $\sigma$, $\rho$, $p$, $u_0$, $u$, linked by the relations above), or the category (isomorphic to $Q^\circ$)} % Doubt resolved against PDF leaf 25: Q is ordinary mathematical italic, retained.
\[
\mathcal{Q}' = \left\{ \quad
\begin{tikzcd}
\xi_0 \arrow[r, shift left, "\sigma"] & \xi \arrow[l, shift left, "\rho"] \arrow[loop right, "p"] \\
 & \eta \arrow[ul, "u_0"] \arrow[u, "u"]
\end{tikzcd}
\qquad
\begin{array}{lcl}
\left.\begin{array}{rcl} \rho\sigma & = & \mathrm{id}_{\xi_0} \\ \sigma\rho & = & p \end{array}\right\} & \Longrightarrow & \left\{\begin{array}{rcl} p^2 & = & p \\ p\sigma & = & \sigma \\ \rho p & = & \rho \end{array}\right. \\[4ex]
\begin{array}{rcl} u & = & \sigma u_0 \end{array} & \Longrightarrow & \left\{\begin{array}{rcl} pu & = & u \\ \rho u & = & u_0 \end{array}\right.
\end{array}
\right.
\]
\textit{We restrict ourselves to functors $\mathcal{Q} \longrightarrow \mathcal{C}$, $\mathcal{Q}' \longrightarrow \mathcal{C}$ such that $p$ is transformed into an arrow} other than an identity \textit{of $\mathcal{C}$ (\footnote[195]{I am not sure if this is reasonable -- to be checked with care!}) (thus the condition relative to $\mathcal{Q}$, $\mathcal{Q}'$ becomes empty, if $\mathcal{C}$ has no projectors).} $\Big)$

\textit{Under these conditions, $\mathcal{E}$ is accessible, and $f : \mathcal{E} \longrightarrow \mathcal{C}$ is accessible.}


\endgroup
% END UNIT der-ch18/der-ch18_p181-185.tex
% BEGIN UNIT der-ch18/der-ch18_p186-190.tex
\begingroup\makeatletter%
\newcommand{\printedpage}[1]{\marginpar{\footnotesize p.~#1}}%
\newcommand{\msp}[1]{\marginpar{\footnotesize [#1]}}%
\newcommand{\Ob}{\operatorname{Ob}}%
\newcommand{\Hom}{\operatorname{Hom}}%
\newcommand{\Ind}{\operatorname{Ind}}%
\newcommand{\card}{\operatorname{card}}%
\newcommand{\cM}{\mathcal{M}}%
\newcommand{\cE}{\mathcal{E}}%
\newcommand{\cC}{\mathcal{C}}%
\newcommand{\cQ}{\mathcal{Q}}%
\newcommand{\uS}{\underline{S}}%
\newcommand{\uT}{\underline{T}}%
\newcommand{\U}{\mathfrak{U}}%
\makeatother\let\printedpage\derpp\let\msp\dermsp%
\let\mspin\DERmspin%


\setcounter{footnote}{195}

\printedpage{184}
% head: A. Grothendieck : Dérivateurs, Ch. XVIII, Édition des Universités de Montpellier II et Paris VII
% foot: Version 17 décembre 2025

\msp{page 475}
We then find the

\begingroup\itshape
\noindent\textbf{\upshape Corollary 4.16.6.2.} Let $f : \cE \longrightarrow \cC$ be a functor. Equivalent conditions:
\begin{enumerate}[label={}]
\item[a)] For every category $\cC'$ that is \textup{small} and \textup{Karoubian}, and every functor $\cC' \longrightarrow \cC$, $\cE' = \cE \times_{\cC} \cC'$ is accessible, and the functor $\cE' \longrightarrow \cC'$ is accessible.
\item[b)] As in a), but $\cC'$ being one of the three categories $\Delta_1$, $\cQ$, $\cQ'$ (where $\cQ$, $\cQ'$ are defined in 4.16.6.1 above).
\item[b$'$)] \textup{(}\footnote{?, to be checked carefully!}\textup{)}. (If $\cC$ is without projectors.) As in b), but restricting to the case $\cC' = \boldsymbol{\Delta}_1$ (to the exclusion of $\cQ$, $\cQ'$).
\end{enumerate}
\DERmark{ENV1541}

\endgroup

\begingroup\itshape
\noindent\textbf{\upshape Corollary 4.16.6.3.} Suppose $f : \cE \longrightarrow \cC$ is fibrant or cofibrant. Then the preceding equivalent conditions imply that:
\begin{enumerate}[label={}]
\item[1$^\circ$)] The $\cE_s$ are accessible ($s \in \Ob \cC$).
\item[2$^\circ$)] For every arrow $\lambda : s \longrightarrow t$ in $\cC$, the functor $\lambda^* : \cE_t \longrightarrow \cE_s$ (resp. $\cE_s \longrightarrow \cE_t$) is accessible.
\end{enumerate}
\noindent On the other hand, the collection of these conditions 1$^\circ$, 2$^\circ$) is equivalent to the validity of b) above, in the case $\cC' = \boldsymbol{\Delta}_1$, or again to the validity
\msp{page 476}
of a) above, when $\cC'$ is \textup{without projectors}.
\endgroup

\textbf{Question 4.16.6.4.} \textit{Do the preceding conditions 1$^\circ$ and 2$^\circ$ imply those of 4.16.6.2? In other words, if the base category $\cC$ is equal to $\cQ$ or to $\cQ'$, then do these conditions imply that $\cE$ is accessible, and $f : \cE \longrightarrow \cC$ accessible?}

It seems stupid not to know the answer!

\textbf{4.16.6.5.} I feel like calling a category \textit{accessible over $\cC$}, a category $\cE$ over $\cC$ satisfying the equivalent conditions of 4.16.6.2 -- at least in the case where $\cC$ is assumed small. For this terminology to be really satisfactory, we would have to know how to prove that these conditions imply this: for every accessible category $\cM$ and every accessible functor $\cM \longrightarrow \cC$, the category $\cE \times_{\cC} \cM$ is accessible, and the functor $\cE \times_{\cC} \cM \longrightarrow \cM$ is accessible. % [sic] satisfaisant -> satisfaisante, accord avec « terminologie »
\DERmark{ENV1542}


\msp{page 476bis}
The idea is that one would call (under a smallness hypothesis) \textit{accessible category $\cE$} over another $\cC$, a category having the property we have just stated -- but now assuming in addition that $\cC$ satisfies $\mathrm{L}_\pi$, I imagine, for every accessible functor $\cM \longrightarrow \cC$, with $\cM$ accessible, $\cE \times_{\cC} \cM$ is accessible and the canonical functor $\cE \times_{\cC} \cM \longrightarrow \cE$ is accessible.

\msp{page 477}
\subsection*{4.16.7. Return to the closed subcategories of an accessible category (4.16.1).}

\textbf{Proposition 4.16.7.1.} \textit{Let $\cM$ be an accessible category, $(\uS_i)_{i\in I}$ a family of closed subcategories accessible in $\cM$. We do not assume that \textup{card}\,$I$ is small. Let $\uS$ be the strictly full subcategory of $\cM$ such that $\Ob \uS = \bigcup_{i\in I} \Ob \uS_i$. Then for $\uS$ % [sic] On ne suppose que -> On ne suppose pas que, « pas » apparemment omis
\printedpage{185}
% head: A. Grothendieck : Dérivateurs, Ch. XVIII, Édition des Universités de Montpellier II et Paris VII
% foot: Version 17 décembre 2025
to be accessible in $\cM$, it is necessary and sufficient that there exist a \textup{small} subset $J \subset I$, such that $\Ob \uS = \bigcup_{i\in J} \Ob\underline S_i$. (In particular, if $I$ is ordered, and $i \longmapsto \uS_i$ increasing and large relative to every small cardinal, it is necessary and sufficient that there exist an $i \in I$ such that $\uS = \uS_i$, i.e. such that $j \geqslant i$ implies $\uS_j = \uS_i$.)}
\DERmark{ENV1543}


This is obviously sufficient, by virtue of 4.13\dots{} Let us prove conversely that if $\uS$ is accessible in $\cM$, then the preceding condition is satisfied. Let $\pi$ be a cardinal $\geqslant \pi_0$ such that $\uS$ is $\pi$-accessible in $\cM$, hence
\[
\uS \ \xleftarrow{\ \sim\ }\ \Ind_\pi \uS(\pi) \ , \qquad\qquad \text{where } \uS(\pi) = \uS \cap \cM_\pi.
\]

\msp{page 478}
Let $\uS(\pi)_0$ be a skeletal subcategory of $\uS(\pi)$, such that $\uS(\pi)_0 \hookrightarrow \uS(\pi)$ is an equivalence, so that $\uS(\pi)_0$ is small. For every $x \in \Ob \uS(\pi)_0$, let $i_x \in I$ be such that $x \in \Ob \uS_{i_x}$. Let $J$ be the set of all the $i_x$, which is a small subset of $I$, and let $\uS_J$ be the strictly full subcategory of $\cM$ with set of objects $\bigcup_{j\in J} \Ob \uS_j$. Thus $\uS(\pi)_0$ is contained in $\uS_J$, hence also $\uS(\pi)$, since $\uS_J$ is strictly full: % [sic] que est -> qui est, pronom relatif sujet
\[
\uS(\pi) \ \hookrightarrow\ \uS_J \ .
\]
Moreover $\uS_J$ is accessible in $\cM$ (loc. cit.), and is closed in $\cM$ as a union of closed subcategories, \textit{a fortiori it satisfies $L_\pi$}. Thus, since $\uS(\pi) \subset \uS_J$, the essentially surjective functor $\Ind_\pi(\uS(\pi)) \longrightarrow \uS$ factors through $\uS_J$. Hence $\uS_J \hookrightarrow \uS$ is essentially surjective, and since $\uS_J$ is strictly full, we conclude that $\uS_J = \uS$, q.e.d.

\msp{page 479}
\textbf{Proposition 4.16.7.2.} \textit{Let $\cM$ be an accessible category, $x$ an object of $\cM$, and ${}_x\cM$ the closed subcategory of $\cM$ generated by $x$, consisting of the $y \in \Ob \cM$ such that $\Hom(x,y) \neq \emptyset$, i.e. the essential image of $x\backslash\cM \longrightarrow \cM$. Then ${}_x\cM$ is a subcategory of $\cM$ accessible in $\cM$.}

Indeed, it is the inverse image of the subcategory $\mathrm{Ens}^*$ accessible in $\mathrm{Ens}$ under the representable (hence accessible) functor
\[
h_x \ :\ y \longmapsto \Hom(x,y) \ :\ \cM \ \longrightarrow\ \mathrm{Ens} \ ,
\]
hence is accessible in $\cM$ by virtue of 4.13\dots

\textbf{Corollary 4.16.7.3.} \textit{Let $\cM$ be an accessible category, $\uS$ a closed subcategory. For $\uS$ to be accessible, it is necessary and sufficient that there exist a \textup{small} subset $J$ of $\uS$, such that $\Ob \uS = \bigcup_{x\in J} \Ob\,(\,{}_x\cM)$.}

Indeed, we have
\[
\Ob \uS \ =\ \bigcup_{x\in\Ob\uS} \Ob\,(\,{}_x\cM) \ ,
\]
i.e. $\uS$ appears as the union of
\msp{page 480}
the family (not necessarily small) of closed subcategories ${}_x\cM$ of $\cM$ (indexed by $x \in \Ob \uS$), accessible in $\cM$ by virtue of 4.16.7.2. Thus our assertion is a particular case of 4.16.7.1.

\printedpage{186}
% head: A. Grothendieck : Dérivateurs, Ch. XVIII, Édition des Universités de Montpellier II et Paris VII
% foot: Version 17 décembre 2025

One can also, for every small subset $J$ of $\Ob \uS$, consider the closed subcategory $\uS_J \subset \uS$ of $\cM$ generated by $J$. The set $I$ of the subsets $J$ in question, ordered by inclusion, is large relative to every cardinal $\pi \in \U$, and $\uS$ is an increasing union of the $\uS_J$. On the other hand, the proposition tells us that $\uS$ is accessible in $\cM$ if and only if there exists a $J \in I$, such that $\uS_J = \uS$, i.e. such that $\uS_J = \uS_{J'}$ for $J \subset J'$.

\begingroup\itshape
\noindent\textbf{\upshape Proposition 4.16.7.4.} Let $\cM$ be an accessible category. The following conditions are equivalent:
\begin{enumerate}[label={}]
\item[a)] There exists a subcategory $\uS$ of $\cM$ closed in $\cM$, and not accessible in $\cM$.
\item[b)] Let $I$ be the set of ordinals $\alpha$ such that $\card \alpha \in \U$, i.e. of the
\msp{page 481}
ordinals $\alpha$ which are in $\U$. Thus $I = I_{\alpha_0}$, where $\alpha_0$ is the smallest ordinal which does not belong to $\U$, and $I$ is large relative to every small cardinal $\pi$, i.e. relative to every cardinal in $\U$ \textup{(}\footnote{introduce $I = I_{\U}$ before the statement of the proposition.}\textup{)}. The condition under consideration is the following:

There exists a transfinite sequence % [sic] transfini -> transfinie, accord avec « suite »
\[
\underline{x} \ =\ (x_\alpha)_{\alpha\in I}
\]
of elements of $\Ob \cM$, such that we have
\[
\Hom(x_\beta, x_\alpha) \ =\ \emptyset \qquad\qquad \textit{if } \beta < \alpha.
\]
To be precise, for such a transfinite sequence, let $\uS$ or ${}_{\underline{x}}\cM$ be the closed subcategory of $\cM$ generated by the family of the $x_\alpha$, i.e. such that $\Ob \uS = \bigcup_{\alpha\in I} \Ob\ {}_{x_\alpha}\cM = \bigcup_{\alpha\in I} x_\alpha\backslash\Ob \cM$. Then ${}_{\underline{x}}\cM$ is a closed subcategory not accessible in $\cM$, and every closed subcategory not accessible in $\cM$ is obtained in this way.
\end{enumerate}
\endgroup

\textsc{Proof}. Let us consider a sequence $\underline{x} = (x_\alpha)_{\alpha\in I}$ as above. If $J$ is a small subset of $I$, there exists an $\alpha \in I$ which is an upper bound for the $\beta \in J$, hence $\Hom(x_\beta, x_\alpha) = \emptyset$, for $\beta \in J$,
\msp{page 482}
which means that $x_\alpha$ does not belong to any of the $\uS_\beta = {}_{x_\beta}\cM$, for $\beta \in J$. Thus $\uS$ is not the union of a small subfamily of the family of the ${}_{x_\alpha}\cM$, hence by virtue of 4.16.7.1, $\uS$ is not accessible in $\cM$. % [sic] n'est par réunion -> n'est pas réunion ; n'est par accessible -> n'est pas accessible, « par » pour « pas » (2 fois)

Conversely, let $\uS$ be such a subcategory in $\cM$. Let us consider a well-order on $\Ob \uS$ (objects $s_\alpha$), and construct the sequence of the $x_\alpha$ ($\geqslant s_\alpha$) by transfinite induction (\footnote{\textbf{N.B.} We have $\card \uS \notin \U$, otherwise, $\uS$ would be small, hence accessible in $\cM$, contrary to the hypothesis.}).

1) We choose $x_0 =$ first element of $\Ob \uS$ (exists, since $\Ob \uS \neq \emptyset$, i.e. $\uS \neq \emptyset$, since otherwise, $\uS$ would be accessible in $\cM$).

2) If the $x_\beta$ ($\beta < \alpha$) have been constructed, we take for $x_\alpha$ the first element of $\uS$ bounding the $x_\beta$ and which does not belong to the closed subcategory of $\cM$ generated by the $x_\beta$ ($\beta < \alpha$). (This closed subcategory is $\subset \uS$, and $\neq \uS$, since $\uS$ is not accessible in $\cM$.) Thus we do have
\[
\Hom(x_\beta, x_\alpha) \ =\ \emptyset \qquad\qquad \text{if } \beta < \alpha.
\]

\printedpage{187}
% head: A. Grothendieck : Dérivateurs, Ch. XVIII, Édition des Universités de Montpellier II et Paris VII
% foot: Version 17 décembre 2025

\msp{page 483}
Thus, the transfinite induction continues over all of $I$. We see by transfinite induction that we have, for every $\alpha \in I$,
\[
s_\alpha \ \leqslant\ x_\alpha
\]
(since $x_\alpha > x_\beta \geqslant s_\beta$ for $\beta < \alpha$, hence $x_\alpha > s_\beta$ for every $\beta < \alpha$, which is equivalent to $x_\alpha \geqslant s_\alpha$). Moreover, we also find, by transfinite induction, that
\DERmark{ENV1544}

\[
\beta \ <\ \alpha \qquad\quad \Longrightarrow \qquad\quad
\text{\parbox[c]{5.5cm}{$s_\beta$ belongs to the closed subcategory $\uS_\alpha$ of $\cM$ generated by the $x_{\beta'}$ with $\beta' < \alpha$.}}
\quad ,
\tag*{$(*)$}
\]
since by construction $x_\alpha$ is the \textit{first} of the elements of $\uS$ which is \textit{not} in the said subcategory $\uS_\alpha$ of $\cM$. This shows that every $s_\beta$ belongs to the closed subcategory ${}_{\underline{x}}\cM$ of $\cM$ generated by the family of \textit{all} the $x_\alpha$, so we find that $\uS \subset {}_{\underline{x}}\cM \subset \uS$, hence $\uS = {}_{\underline{x}}\cM$, q.e.d.
\DERmark{ENV1545}


\begingroup\itshape
\noindent\textbf{\upshape Proposition 4.16.7.5.} For a given universe $\U$ (containing $\mathbf{N}$ as an element), the following conditions are equivalent:
\DERmark{ENV1546}

\begin{enumerate}[label={}]
\item[a)] There exists a $\U$-accessible category $\cM$ having a closed subcategory not accessible % [sic] 1) -> a), la démonstration renvoie à a) et b)
\msp{page 484}
in $\cM$.
\DERmark{ENV1547}

\item[b)] There exists a category $\cC \in \U$, and a transfinite sequence of elements of
\[
\cC^\wedge \ \overset{\mathrm{def}}{=}\ \underline{\mathrm{Hom}}(\cC^{\circ}, \mathrm{Ens}_{\U}) \ ,
\]
\DERmark{ENV1548}\DERmark{ENV1549}

$\underline{x} = (x_\alpha)_{\alpha\in I}$, indexed by the set of all ordinals $\alpha \in \U$, and such that
\[
\Hom(x_\alpha, x_\beta) \ =\ \emptyset \qquad\qquad \textit{if } \alpha < \beta.
\]
\end{enumerate}
\endgroup

\textsc{Proof}. We have b) $\Longrightarrow$ a) by taking $\cM = \cC^\wedge$, by virtue of the preceding proposition. We see likewise a) $\Longrightarrow$ b), by embedding $\cM$ in a category $\cC^\wedge$ ($\cC$ being a small strictly generating subcategory of $\cM$).

I have not been able (\footnote{But yes, in the end, cf. 4.16.7.10 below.}) to construct a sequence $(x_\alpha)_{\alpha\in I}$ as required in this proposition, nor to prove that such a sequence does not exist, even in a particular case such as $\cM = (k\text{-comm. alg.})$, $k$ a commutative ring. I have not
\msp{page 485}
even proved that one cannot find a transfinite sequence of \textit{fields} $k_\alpha$, having the stated properties.

A fortiori, I have not been able to clear up the question whether an $\underline{\mathrm{L}}$-stable subcategory of an accessible category $\cM$ is always accessible in $\cM$ (\footnote{This is false, cf. below.}). It could be that there is no such case. It could also be that there are some, even with a closed subcategory (which is stable under every type of $\varinjlim$ representable in $\cM$, except for the $\varinjlim$ over an empty index category $\ldots$).

\printedpage{188}
% head: A. Grothendieck : Dérivateurs, Ch. XVIII, Édition des Universités de Montpellier II et Paris VII
% foot: Version 17 décembre 2025

For \textit{open} subcategories, I recall that I constructed (in the appendix) an example of such a subcategory, not accessible in $\cM$. Concerning open subcategories, I point out here the statement (almost tautological):

\msp{page 486}
\textbf{Proposition 4.16.7.6} (\footnote{This proposition and its corollary are contained in 4.16.7.8 below, there is no need to make a separate statement of them.}). \textit{Let $\cM$ be an accessible category, $\uS$ an \textup{open} subcategory of $\cM$. For $\uS$ to be accessible in $\cM$, it is necessary and sufficient that it be $\underline{L}$-stable in $\cM$, i.e. that there exist a cardinal $\pi_0$ such that $\cM$ satisfies $L_{\pi_0}$, and that $\uS$ be stable in $\cM$ under $\varinjlim$ of type $L_{\pi_0}$. More precisely, if this condition is satisfied, then for every cardinal $\pi$ such that $\pi \geqslant \pi_0$ and that $\pi$ is useful for $\cM$, $\uS$ is $\pi$-accessible in $\cM$. Or again: if $\pi$ is a useful cardinal for $\cM$, then $\uS$ is $\pi$-accessible in $\cM$ if and only if it is $L_\pi$-stable in $\cM$.}

This last assertion follows at once from the following corollary, itself tautological.

\textbf{Corollary 4.16.7.7.} \textit{Let $\pi$ be a cardinal as indicated above, let $\uS(\pi) = \uS \cap \cM_\pi$, and let $x \in \Ob \cM$ and
\[
x \ \simeq\ \varinjlim_{i\in I} x_i \qquad\qquad
\begin{cases}
x_i \in \Ob \cM_\pi \\[1ex]
I \ \textit{large relative to}\ \pi
\end{cases}
\]
be a representation of $x$ as the essential image of an object of $\Ind_\pi \cM_\pi$. Then $x \in \Ob \uS$ $\Longleftrightarrow$ $x_i \in \Ob \uS(\pi)$ $\forall\, i \in I$.}

\msp{page 487}
I have finally cleared up the question of the accessibility of open subcategories and closed subcategories of an accessible category:

\begingroup\itshape
\noindent\textbf{\upshape Proposition 4.16.7.8.} Let $\cM$ be an accessible category, $\uS$ a closed subcategory of $\cM$, $\uT$ the open subcategory complementary to $\uS$, $\pi$ an infinite cardinal useful for $\cM$. The following conditions are equivalent:
\begin{enumerate}[label={}]
\item[a)] $\uS$ is $\pi$-accessible in $\cM$.
\item[a$'$)] $\uS$ is the closed subcategory generated by $\uS(\pi) \overset{\mathrm{def}}{=} \uS \cap \cM_\pi$, i.e. for $x \in \Ob \cM$, we have $x \in \Ob \uS$ if and only if $\exists\, \xi \in \Ob \uS(\pi)$ such that $\Hom_{\cM}(\xi, x) \neq \emptyset$.
\item[b)] $\uT$ is $\pi$-accessible in $\cM$.
\item[b$'$)] $\uT$ is $\underline{L}_\pi$-stable in $\cM$.
\end{enumerate}
\DERmark{ENV1550}

\endgroup

\noindent Let us note at once the

\begingroup\itshape
\noindent\textbf{\upshape Corollary 4.16.7.9.} For $\cM$, $\uS$, $\uT$ as above, equivalent conditions:
\begin{enumerate}[label={}]
\item[a)] $\uS$ is accessible in $\cM$.
\item[a$'$)] There exists a \textup{small} subset $A$ of $\Ob \cM$, such that $\uS$ is the closed subcategory ${}_A\cM$ generated by $A$, i.e. $\uS = \bigcup_{\xi\in A} {}_\xi\cM$, where ${}_\xi\cM$ is the closed subcategory of $\cM$ consisting of the $x \in \Ob \cM$ such that $\Hom(\xi, x) \neq \emptyset$.
\item[b)] \msp{page 488}$\uT$ is accessible in $\cM$.
\end{enumerate}
\endgroup


\endgroup
% END UNIT der-ch18/der-ch18_p186-190.tex
% BEGIN UNIT der-ch18/der-ch18_p191-195.tex
\begingroup\makeatletter%
\newcommand{\printedpage}[1]{\marginpar{\footnotesize p.~#1}}%
\newcommand{\msp}[1]{\marginpar{\footnotesize [#1]}}%
\newcommand{\cM}{\mathcal{M}}%
\newcommand{\cN}{\mathcal{N}}%
\newcommand{\cF}{\mathcal{F}}%
\newcommand{\cA}{\mathcal{A}}%
\newcommand{\Ob}{\operatorname{Ob}}%
\newcommand{\Hom}{\operatorname{Hom}}%
\newcommand{\Ind}{\operatorname{Ind}}%
\newcommand{\Ker}{\operatorname{Ker}}%
\newcommand{\Max}{\operatorname{Max}}%
\newcommand{\Ann}{\operatorname{Ann}}%
\newcommand{\defeq}{\overset{\mathrm{def}}{=}}%
\newcommand{\incl}{\lhook\joinrel\longrightarrow}%
\makeatother\let\printedpage\derpp\let\msp\dermsp%
\let\mspin\DERmspin%

\DERmark{ENV1551}


\printedpage{189}
% head: A. Grothendieck : Dérivateurs, Ch. XVIII, Édition des Universités de Montpellier II et Paris VII

\begin{enumerate}
\item[\textit{b$'$)}] \textit{$\underline{T}$ is $\underline{L}$-stable in $\cM$.}
\end{enumerate}

\noindent\textbf{Corollary 4.16.7.10.} \textit{Let $\cM$ be an accessible category. The following conditions are equivalent} (\footnote[202]{Example where these conditions are satisfied: $\cM = (k\text{-comm. alg.})$, $k$ an infinite commutative field, $\mathfrak{U}$ being the smallest universe such that $\mathbf{N}\in\mathfrak{U}$.}).
\begin{enumerate}
\item[\textit{a)}] \textit{There exists a closed subcategory of $\cM$, not accessible in $\cM$.}
\item[\textit{b)}] \textit{There exists an open subcategory in $\cM$, not accessible in $\cM$.}
\item[\textit{c)}] \textit{There exists a transfinite sequence}
\[
\underline{x} \ = \ (x_\alpha)_{\alpha\in I}
\]
\textit{of objects of $\cM$, indexed by the set $I$ of ordinals of the reference universe $\mathfrak{U}$, such that we have}
\[
\Hom(x_\alpha,x_\beta) \ = \ \emptyset \qquad \text{\textit{for} } \alpha<\beta .
\]
\end{enumerate}

\noindent Indeed, the equivalence of a) and b) is clear from the preceding corollary, and that of a) and c) is none other than 4.16.7.4.

\noindent\textsc{Proof of} 4.16.7.8. We shall have to use the

\noindent\textbf{Lemma 4.16.7.11.} \textit{Let $\cM$ be an accessible category, $\pi$ a useful cardinal for $\cM$, $\underline{S}$ a full subcategory of $\cM$, $\underline{S}(\pi)\defeq\underline{S}\cap\cM_\pi$, whence a fully faithful functor}
\[
\varphi \ : \ \Ind_\pi(\underline{S}(\pi)) \longrightarrow \cM \ .
\]

\msp{page 489}
\begin{enumerate}
\item[\textit{a)}] \textit{Let $x\in\Ob\cM$. For $x$ to belong to the essential image of the functor $\varphi$, it is necessary and sufficient that the full subcategory $\underline{S}(\pi)/x$ of the category $\cM_\pi/x$ (which is filtering and even large relative to $\pi$) be cofinal (hence also large relative to $\pi$).} % [sic] \psi -> \varphi, the functor of the lemma is called \varphi in the display above
\item[\textit{b)}] \textit{For $\underline{S}$ to be $\pi$-accessible in $\cM$, it is necessary and sufficient that for every $x\in\Ob\cM$, we have}
\[
x \ \in \ \Ob\underline{S} \qquad\iff\qquad \underline{S}(\pi)/x \text{ \textit{is cofinal in} } \cM_\pi/x
\]
\DERmark{ENV1552}

\textit{(or again, that $\underline{S}$ satisfy $L_\pi$, and that for every $x$ in $\cM$, we have}
\[
x \ \in\Ob\underline{S} \qquad\Longrightarrow\qquad \underline{S}(\pi)/x \text{ \textit{cofinal in} } \cM_\pi/x\text{\textit{.)}}
\]
\end{enumerate}

\noindent\textsc{Proof}. Part b) is a tautological consequence of a), which I shall prove. To say that $x$ is in the essential image of $\varphi$, means that
\DERmark{ENV1553}

\[
x \ \simeq \ \varinjlim_{i\in I} x_i \ , \tag*{$(*)$}
\]
with $I$ large relative to $\pi$ and the $x_i$ in $\underline{S}(\pi)$. We know that $(*)$ implies ($I$ being filtering, and $\cM$ being equivalent to $\Ind_\pi(\cM_\pi)$ since $\pi$ useful for $\cM$) that $I\longrightarrow\cM_\pi/x$ is cofinal (4.4.2, page 79). Now this functor factors through the fully faithful inclusion functor $\underline{S}(\pi)/x\incl\cM_\pi/x$, which is cofinal if and only if it satisfies condition F 1. It is immediate that this condition is satisfied when it is satisfied for the composite functor in question, which is indeed the case
\DERmark{ENV1554}

\printedpage{190}
% head: A. Grothendieck : Dérivateurs, Ch. XVIII, Édition des Universités de Montpellier II et Paris VII
\msp{page 490}
by hypothesis.

\noindent Conversely, suppose that $\underline{S}(\pi)/x$ is cofinal in $\cM_\pi/x$. Since we have
\[
x \xleftarrow{\ \sim\ } \varinjlim_{\cM_\pi/x} x_i \ ,
\]
$\pi$ being useful for $\cM$, we will also have by cofinality
\[
x \xleftarrow{\ \sim\ } \varinjlim_{\underline{S}(\pi)/x} x_i \ , % [sic] \underline{S}(\pi)/\pi -> \underline{S}(\pi)/x, the cofinal subcategory is \underline{S}(\pi)/x
\]
\DERmark{ENV1555}

now $\cM_\pi/x$ being large relative to $\pi$, the same holds for the cofinal full subcategory $\underline{S}(\pi)/x$. Thus $x$ is in the essential image of $\Ind_\pi\underline{S}(\pi)$, q.e.d.
\DERmark{ENV1556}


\noindent\textbf{Lemma 4.16.7.12.} \textit{Let $I$ be a filtering category, $J$ a closed subcategory, $K$ the complementary open subcategory. We then have the equivalences}
\[
J \ \neq \ \emptyset \iff J \text{ \textit{cofinal} } \iff K \text{ \textit{not cofinal} } \iff K \ \neq \ I \ .
\]

\noindent\textsc{Proof} immediate.

\noindent I return to the proof of 4.16.7.8. We apply Lemma 4.16.7.11 to $\cM$, $\pi$, and to the two
\msp{page 491}
subcategories $\underline{S}$, $\underline{T}$ of $\cM$, noting that for $x\in\Ob\cM$, $\underline{S}(\pi)/x$ is a closed subcategory of $\mathcal M_\pi/x$, whose open complement is $\underline{T}(\pi)/x$, so that we can apply Lemma 4.16.7.12. We find:
\DERmark{ENV1557}

\[
\begin{aligned}
\underline{S}\ \pi\text{-accessible} \quad &\overset{4.16.7.11}{\iff} \quad \forall\ x\ \in\ \Ob\cM,\ \text{we have } (x\ \in\ \Ob\underline{S}\iff\underline{S}(\pi)/x \text{ cofinal in}\\
&\qquad\quad \cM_\pi/x)\\
&\iff \quad \forall\, x\in\Ob\cM,\ \text{we have } \Bigl( \qquad \underbrace{x\notin\Ob\underline{S}}_{\text{(i.e. $x\in\Ob\underline{T}$, by 4.16.7.12)}}\\
&\qquad\qquad\qquad \iff \qquad \underbrace{\underline{S}(\pi)/x \text{ not cofinal in } \cM_\pi/x}_{\text{(i.e. $\underline{T}(\pi)/x$ cofinal in $\cM_\pi/x$, by 4.16.7.12)}} \ \Bigr)\\
&\overset{4.16.7.11}{\iff} \quad \underline{T} \text{ is } \pi\text{-accessible.}
\end{aligned}
\]
\DERmark{ENV1558}

We thus have the equivalence of a) and b). On the other hand, $\underline{S}(\pi)/x$ cofinal in $\cM_\pi/x$ means simply $\underline{S}(\pi)/x\neq\emptyset$ (4.16.7.12), so $\underline{S}$ $\pi$-accessible means that $\underline{S}$ consists of the $x\in\Ob\cM$ such that $\underline{S}(\pi)/x\neq\emptyset$, which also means that $\underline{S}$ is the closed subcategory in $\cM$ generated by $\underline{S}(\pi)$, whence a $\iff$ a$'$. On the other hand, the equivalence b $\iff$ b$'$ is none other than 4.16.7.6 (last formulation). % [sic] 4.6.7.6 -> 4.16.7.6, the surrounding references are all to 4.16.7.x
\DERmark{ENV1559}\DERmark{ENV1560}


\msp{page 492}

\section*{4.17\quad Return to paratopoi, and to the operations $\cM\boxtimes\cN$, $\underline{\operatorname{Hom}}_{!}(\cM,\cN)$}

\noindent\texttt{[This paragraph does not exist.]}

\printedpage{191}
% head: A. Grothendieck : Dérivateurs, Ch. XVIII, Édition des Universités de Montpellier II et Paris VII

\msp{page 1}

\section*{Appendix}

\subsection*{Ultraproducts of fields, and accessible categories}

\noindent Let $(k_\xi)_{\xi\in E}$ be a family of fields, indexed by a set $E$. $\Bigl($ Let $\Phi$ be a nontrivial filter in $E$, i.e. a subset $\Phi\subset\mathfrak{P}(E)$ such that:
\begin{enumerate}
\item[a)] $\emptyset\notin\Phi$, i.e. $A\in\Phi\Longrightarrow A\neq\emptyset$.
\item[b)] $A,B\in\Phi\Longrightarrow A\cap B\in\Phi$.
\item[c)] If $A\in\Phi$, $A\subset B$, then $B\in\Phi$.
\item[d)] $\Phi\neq\emptyset$ ($\iff E\in\Phi$).
\end{enumerate}
(This amounts to saying that $\Phi$ is a closed subset of the ordered set $\mathfrak{P}(E)$, \textit{left filtering}, and does not contain the least element of $\mathfrak{P}(E)$, i.e. not equal to all of $\mathfrak{P}(E)$.) $\Bigr)$ For every subset $J$ of $E$, put
\[
k_J \ = \ \prod_{\xi\in J} k_\xi \ ,
\]
we thus have an inverse system
\[
(k_J)_{J\in\mathfrak{P}(E)} \ , \qquad \text{i.e. } J\subset J' \text{ gives } k_{J'}\longrightarrow k_J \ ,
\]
and the system
\[
(k_J)_{J\in\Phi}
\]
is a direct system of rings, indexed by the filtering ordered set $\Phi^o$. Put % [sic] d'annaux -> d'anneaux, misspelling
\[
k_\Phi \ = \ \varinjlim_{J\in\Phi} k_J \ .
\]
If $\Psi$ is a filter finer than $\Phi$, i.e.
\msp{page 2}
such that $\Psi\supset\Phi$, we find a canonical homomorphism
\[
k_\Phi \longrightarrow k_\Psi \ ,
\]
so if $\cF(E)=\cF$ is the set of filters on $E$, ordered by inclusion, we find a direct system
\[
\Phi \longmapsto k_\Phi \ : \ \cF \longrightarrow \Ann \ .
\]
We have not yet used that the $k_\xi$ are fields, in fact it would suffice that they be objects of a category $\cA$ stable under small products and under small filtering $\varinjlim$. Take now $\cA=\Ann$ (category of unital rings, not necessarily commutative), and suppose that the $k_\xi$ are fields. % [sic] produis -> produits, misspelling

\printedpage{192}
% head: A. Grothendieck : Dérivateurs, Ch. XVIII, Édition des Universités de Montpellier II et Paris VII

\noindent\textbf{Proposition 1.} \textit{Suppose that the $k_\xi$ ($\xi\in E$) are} fields (\footnote[203]{NB This condition is necessary if one wants $k_\Phi$ to be a field for every ultrafilter.}). \textit{For $k_\Phi$ to be a field, it is necessary and sufficient that $\Phi$ be an ultrafilter, i.e. that for every decomposition $E=A\amalg B$, we have $A\in\Phi$ or $B\in\Phi$.} (NB If $A\in\Phi$, we have $B\notin\Phi$, since otherwise, we would have $A\cap B\in\Phi$, hence $\emptyset\in\Phi$, absurd.)

\noindent Necessity: Let, for every subset $A$ of $E$,
\[
f_A\in k_E \qquad \text{be defined by} \qquad f_A(\xi) \ = \ \left\{
\begin{array}{lll}
1_{k_\xi} & \text{if} & \xi\in A\\[1ex]
0_{k_\xi} & \text{if} & \xi\in E\backslash A \ .
\end{array}\right.
\]

\msp{page 3}
\noindent Let
\[
f_{A,\Phi} \ \in k_\Phi
\]
be the image of $f_A$ under the canonical homomorphism $k_E\longrightarrow k_\Phi$. If $E=A\amalg B$, we will have
\[
f_A \ + \ f_B \ = \ 1 \ , \qquad f_A\, f_B \ = \ 0 \ ,
\]
whence
\[
f_{A,\Phi} \ + \ f_{B,\Phi} \ = \ 1 \ , \qquad f_{A,\Phi}\, f_{B,\Phi} \ = \ 0 \ .
\]
If $k_\Phi$ is a field, this implies that we have $f_{A,\Phi}=0$ or $f_{B,\Phi}=0$. But if e.g. $f_{A,\Phi}=0$, we see that this means that there exists a $J\in\Phi$ such that $f_A|J=0$, which also means that $J\subset B$. Likewise, if $f_{B,\Phi}=0$, we will have $J\subset A$. Thus $\Phi$ is an ultrafilter.

\noindent Conversely, suppose that $\Phi$ is an ultrafilter. Let us prove that $k_\Phi$ is a field, i.e. that for $f\in k_\Phi$, $f\neq0\Longrightarrow f$ invertible. Note that $f$ comes from an element of a $k_J$ ($J\in\Phi$), hence is the canonical image of a $g\in k_E$. To say that $f=0$ means that $\exists\ J\in\Phi$ such that $g|J=0$. The opposite $f\neq0$ therefore means that for every $J\in\Phi$, we have $g|J\neq0$.
\DERmark{ENV1561}\DERmark{ENV1562}


\msp{page 4}
\noindent Let then $A$ be the set of $\xi\in E$ such that $g(\xi)\neq0$, $B=E\backslash A$ the set of $\xi\in E$ such that $g(\xi)=0$. We thus have $B\notin\Phi$ since $f\neq0$, and since $\Phi$ is an ultrafilter, we therefore have $A\in\Phi$. But let $g'\in k_E$ be defined by $g'(\xi)=g(\xi)^{-1}$ if $\xi\in A$, $g'(\xi)=1$ if $\xi\in B$, and let $f'$ be the image of $g'$ in $k_\Phi$. We then see that we have % [sic] B = E\A ... tels que g(\xi) \neq \emptyset -> g(\xi) = 0, the complement of A is the set where g vanishes
\DERmark{ENV1563}\DERmark{ENV1564}\DERmark{ENV1565}

\[
f\ f' \ = \ f' \  \ f \ = \ 1 \ ,
\]
\DERmark{ENV1566}

which indeed shows that $k_\Phi$ is a field.

\noindent Suppose now that all the $k_\xi$ ($\xi\in E$) are equal to one and the same field $k$. Then the $k_J$ form an inverse system of $k$-algebras, the $k_\Phi$ a direct system of $k$-algebras. If $k$ is commutative, the same holds for the $k_J$, $k_\Phi$, and in particular for $\Phi$ an ultrafilter, $k_\Phi$ is an \textit{extension} of the field $k$.

\noindent Let us return to the general case, but the $k_\xi$ commutative. Let $\cF_0\subset\cF$ be the set of ultrafilters on $E$. Note that for every $\Phi\in\cF$, $k_E\longrightarrow k_\Phi$ is surjective. Thus if $\Phi\in\cF_0$, $k_\Phi$ identifies with a field quotient of the ring $k_E$.
\DERmark{ENV1567}


\printedpage{193}
% head: A. Grothendieck : Dérivateurs, Ch. XVIII, Édition des Universités de Montpellier II et Paris VII

\msp{page 5}
\noindent For every ring $\cA$, let $\Max(\cA)$ be the maximal spectrum, i.e. the set of maximal ideals of $\cA$. We thus find a canonical map
\[
\left\{
\begin{array}{lll}
\cF_0 & \longrightarrow & \Max(k_E)\\[1ex]
\ \Phi & \longmapsto & \Ker(k_E\longrightarrow k_\Phi) \ \defeq \ \mathfrak{m}_\Phi \ .
\end{array}\right. \tag*{$(*)$}
\]

\noindent\textbf{Proposition 2.} \textit{The preceding map is} bijective.

\noindent Let us first prove that it is injective, this follows at once from the

\noindent\textbf{Corollary 1} (\footnote[204]{\textbf{NB} For arbitrary $\Phi\in\cF$, $\mathfrak{m}_\Phi\defeq\Ker(\chi_\Phi : k_E\longrightarrow k_\Phi)$ determines $\Phi$, since
\[
A \ \in \ \Phi \qquad\iff\qquad \chi_\Phi(f_A) \ = \ 1
\]
(but one must assume the $K_\xi\neq0$, i.e. $0_\xi\neq1_\xi\ \ldots$).}). \textit{Let, for $\Phi\in\cF_0$,}
\[
\chi_\Phi \ : \ k_E \longrightarrow k_\Phi
\]
\textit{be the canonical map. Then for $A\in\mathfrak{P}(E)$, we have}
\[
A \ \notin \ \Phi \qquad\iff\qquad \chi_\Phi(f_A) \ = \ 0 \ , \quad \text{\textit{i.e.} } f_A\in\mathfrak{m}_\Phi
\]
\textit{(i.e. $A\in\Phi\iff\chi_\Phi(f_A)\neq0$, i.e. $f_A\notin\mathfrak{m}_\Phi$)} (\footnote[205]{\textbf{NB} We thus have
\[
(\chi_\Phi(f_A) \ \neq \ 0) \qquad\iff\qquad (\chi_\Phi(f_A) \ = \ 1) \ .
\]}).

\noindent Indeed, $\chi_\Phi(f_A)=0$ means that there exists $B\in\Phi$ such that $f_A|B=0$, i.e. such that $B\subset E\backslash A$; but this is equivalent to $E\backslash A\in\Phi$, hence to $A\notin\Phi$ since $\Phi$ is an ultrafilter, q.e.d.

\msp{page 6}
\noindent Remains to prove that the map $(*)$ is surjective. So let $\mathfrak{m}$ be a maximal ideal of $k_E$, and let
\[
\Phi \ \defeq \ \{A\in\mathfrak{P}(E)\mid\chi(f_A) \ = \ 1\} \ ,
\]
where $\chi : k_E\longrightarrow k_E/\mathfrak{m}=\kappa(\mathfrak{m})$ is the canonical map. I say that $\Phi$ is an ultrafilter. Indeed:
\begin{enumerate}
\item[a)] $\emptyset\notin\Phi$ since $f_\emptyset=0$, whence $\chi(f_\emptyset)=0\neq1$ in $\kappa(\mathfrak m)=k_E/\mathfrak m$.
\item[b)] If $A,B\in\Phi$, we have $A\cap B\in\Phi$, since $f_{A\cap B}=f_A\cdot f_B$, whence $\chi(f_{A\cap B})=\chi(f_A)\chi(f_B)=1$.
\item[c)] If $A\subset B$ and $A\in\Phi$, then $B\in\Phi$. Indeed, we have $f_A=f_Af_B$, whence $\chi(f_A)=\chi(f_A)\chi(f_B)$, and since $\chi(f_A)=1$, we find $\chi(f_B)=1$.
\item[d)] We have $E\in\Phi$, i.e. $\chi(f_E)=1$, now $f_E=1$, so $\chi(f_E)=1$.
\end{enumerate}
\DERmark{ENV1568}

Thus, $\Phi$ is a filter -- we have used only the fact that $k_E/\mathfrak{m}$ is a ring $\neq0$ (to conclude that $\emptyset\notin\Phi$). Let us prove that if $k_E/\mathfrak{m}$ is an integral domain, $\Phi$ is an \textit{ultrafilter}, i.e.:
\DERmark{ENV1569}\DERmark{ENV1570}

\begin{enumerate}
\item[e)] If $A\notin\Phi$, then $B\defeq E\backslash A\in\Phi$.
\end{enumerate}
Indeed, we have $f_A+f_B=1$, $f_Af_B=0$, whence


\endgroup
% END UNIT der-ch18/der-ch18_p191-195.tex
% BEGIN UNIT der-ch18/der-ch18_p196-200.tex
\begingroup\makeatletter%
\newcommand{\printedpage}[1]{\marginpar{\footnotesize p.~#1}}%
\newcommand{\msp}[1]{\marginpar{\footnotesize [#1]}}%
\newcommand{\limind}{\mathop{\underrightarrow{\mathrm{lim}}}}%
\newcommand{\limproj}{\mathop{\underleftarrow{\mathrm{lim}}}}%
\makeatother\let\printedpage\derpp\let\msp\dermsp%
\let\mspin\DERmspin%


\printedpage{194}
% head: A. Grothendieck : Dérivateurs, Ch. XVIII, Édition des Universités de Montpellier II et Paris VII

\noindent\hspace*{1em}\textcircled{a}\quad $\chi(f_A)+\chi(f_B)=1$, and\\
\hspace*{1em}\textcircled{b}\quad $\chi(f_A)\chi(f_B)=0$.\\
$\chi(f_A)\neq 1$, we conclude from \textcircled{a} $\chi(f_B)\neq 0$, whence by \textcircled{b} $\chi(f_A)=0$, whence from \textcircled{a} $\chi(f_B)=1$, whence $B\in\Phi$.

\msp{page 7}
Let us finally prove that $\mathfrak{m}=\mathfrak{m}_\Phi$. Since $\mathfrak{m}_\Phi$ is a maximal ideal, it suffices to prove
\[
\mathfrak{m}_\Phi\ \subset\ \mathfrak{m}\ .
\]
Now if $A\in\mathfrak{P}(E)$ is such that $A\notin\Phi$, or $f_A\in\mathfrak{m}_\Phi$ (cor.~1), we also have $f_A\in\mathfrak{m}$. Thus it suffices to prove the

\textbf{Lemma.} \textit{Let $\Phi\in\mathcal{F}_0$. Then $\mathfrak{m}_\Phi$ is the ideal generated by the $f_A$ for $A\in\mathfrak{P}(E)\backslash\Phi$. More precisely, the $f_A$ are in $\mathfrak{m}_\Phi$, and every $f\in\mathfrak{m}_\Phi$ is of the form $ff_A$, with $A$ as above.}

Indeed, $f\in\mathfrak{m}_\Phi$ means that there exists $B\in\Phi$ such that $f|B=0$. But this also means that $f=ff_A$, where $A=E\backslash B$ \texttt{[rather} $A=E\backslash B$\texttt{]}, whence $A\notin\Phi$.
\DERmark{ENV1571}


This proves proposition 2, and at the same time:

\textbf{Corollary 2.} \textit{Every quotient ring of $k_E=\prod_{\xi\in E}k_\xi$ which is an integral domain is a field.}

Let us now note that if $\Phi\in\mathcal{F}$ is any filter on $E$, then $k_\Phi$ is a quotient ring of $k_E$, hence
\[
\operatorname{Max}(k_\Phi)\ \hookrightarrow\ \underbrace{\operatorname{Max}(k_E)}_{\simeq\,\mathcal{F}_0}\ . \tag*{$(*)$}
\]
This being so, we find:

% a vertical bar is printed in the left margin beside the statement of Corollaire 3
\textbf{Corollary 3} (\footnote[206]{Works for any system of rings $(k_\xi)_{\xi\in E}$ \textit{which are nonzero}, and $\Phi,\Psi\in\mathcal{F}$, $\Psi$ finer than $\Phi$, i.e. $\Phi\subset\Psi$ if and only if $\mathfrak{m}_\Phi\subset\mathfrak{m}_\Psi$ (no need for $\Psi$ to be an ultrafilter).}). \textit{The preceding inclusion identifies} $\operatorname{Max}(k_\Phi)$ \textit{with the set of ultrafilters $\Psi$ on $E$ which are} finer \textit{than $\Phi$.}

Indeed, if $\Psi$ is such an ultrafilter, we
\msp{page 8}
have the commutative diagram
\[
\begin{tikzcd}
k_E \arrow[r, "\chi_\Phi"] \arrow[rr, bend right=15, "\chi_\Psi"'] & k_\Phi \arrow[r] & k_\Psi
\end{tikzcd}\ , \tag*{$(**)$}
\]
which shows that $\mathfrak{m}_\Psi$ \texttt{[rather} $\mathfrak{m}_\Psi$\texttt{]} is in the image of $(*)$. Conversely, suppose that $\Psi\in\mathcal{F}_0$ comes from $\operatorname{Max}(k_\Phi)$, i.e. that we have $\mathfrak{m}_\Phi\subset\mathfrak{m}_\Psi$ (where $\mathfrak{m}_\Phi$ is the kernel of $\chi_\Phi$), or again that we have a factorization $(**)$ of $\chi_\Psi$. Let us prove that $\Psi$ is finer than $\Phi$, i.e. that $\Psi\supset\Phi$. So let $A\in\Phi$, let us prove that we have $A\in\Psi$. But $A\in\Phi$ implies \texttt{[the]} fact that $\chi_\Phi(f_A)=1$, hence its image in $k_\Psi$ is equal to 1, hence we have $A\in\Psi$ by corollary 1.
\DERmark{ENV1572}


\printedpage{195}
% head: A. Grothendieck : Dérivateurs, Ch. XVIII, Édition des Universités de Montpellier II et Paris VII

\msp{page 9}
Let us now suppose that all the $k_\xi$ are equal \texttt{[to]} one and the same ring \fbox{$\ k\neq 0\ $}\,. Thus $k_J$, $k_\Phi$ are ``rings under $k$'' $k\longrightarrow k_J\longrightarrow k_\Phi$, moreover $k\longrightarrow k_J$ (if $J\neq\emptyset$) and $k\longrightarrow k_\Phi$ are injective. We ask the question when $k\longrightarrow k_\Phi$ is \textit{bijective}, i.e. surjective. Since $k_\Phi\simeq k_E/\mathfrak{m}_\Phi$, this therefore means that $k_E=k+\mathfrak{m}_\Phi$, i.e. that every $f\in k_E$ can be written
\[
f\ =\ c\ +\ g\qquad
\begin{cases}
c\ \in\ k\ , & \text{i.e. } c(\xi)=\underbrace{c_0}_{\in\,k}\ \forall\,\xi\in E\ ,\\[2ex]
g\ \in\ \mathfrak{m}_\Phi & \text{i.e. } \exists\,A\in\Phi \text{ with } g|A=0\ .
\end{cases}
\]
But for $f\in k_E$, $c\in k$ and $A\in\Phi$ fixed, the existence of $g\in k_E$ such that $g|A=0$ and $f=c+g$ is equivalent to $f|A=c_A$ (constant function on $A$, with value $c$), i.e. $A\subset f^{-1}(\{c\})$. For $f$, $c$ fixed, the existence of $A\in\Phi$ satisfying the relation is equivalent to $f^{-1}(\{c\})\in\Phi$. And for $f$ fixed, the existence of $c$ is equivalent to that of a fiber $E_c$ of $f:E\longrightarrow k$, which is in $\Phi$. Thus:

\textbf{Proposition 3.} \textit{Let $\Phi\in\mathcal{F}$. For $k\longrightarrow k_\Phi$ to be surjective, it is necessary and sufficient that for every family of subsets} essentially disjoint \textit{$(E_\xi)_{\xi\in k}$ of $E$ indexed by $k$, with union $E$, there exist a $\xi\in k$ such that $E_\xi\in\Phi$. A fortiori, it} is necessary \textit{that $\Phi$ be an ultrafilter (since} $\operatorname{card}k\geqslant 2$ \textit{by hypothesis).}

\msp{page 10}
Indeed, giving such a family $(E_\xi)_{\xi\in k}$ is equivalent to giving a map $E\longrightarrow k$, i.e. an element of $k_E$. % [sic] $f_E$ -> $k_E$, an application $E\to k$ is an element of $k_E$
\DERmark{ENV1573}


\textbf{Scholie.} The property we have just made explicit for a filter $\Phi$ on $E$, depends on $k$ only through $\pi=\operatorname{card}k$. We can also state it by saying that for every \textit{partition} $(E_i)_{i\in I}$ of $E$ into a set $I$ of subsets such that $\operatorname{card}I\leqslant\pi$, there is one (and only one) of the $E_i$ which is in $\Phi$. We then say that $\Phi$ is a $\pi$-\textit{ultrafilter.}

If $\pi\geqslant 2$ is finite, this property is equivalent to saying that $\Phi$ is an ultrafilter. Besides, whatever $\pi$, the ultrafilters $\Phi$ \textit{which are trivial} (i.e. such that there exists $\xi\in E$ with $\{\xi\}\in\Phi$, or again such that $\Phi=\{A\in\mathfrak{P}(E)\mid\xi\in A\}$) are $\pi$-ultrafilters. If $\pi$ is infinite, it is not clear that there exist nontrivial $\pi$-ultrafilters. We know that such exist on $E$ \textit{only} if $\operatorname{card}E\notin\mathfrak{U}_1$, where $\mathfrak{U}_1$ is the smallest universe such that $\mathbf{N}\in\mathfrak{U}_1$, i.e. such
\msp{page 11}
that $\mathfrak{U}$ contains as an element an infinite set.

If $\mathfrak{U}=\mathfrak{U}_1$, and if $k$ is an infinite field, e.g. $k=\mathbf{Q}$, then for every set $E\in\mathfrak{U}$ and for every \textit{nontrivial} ultrafilter $\Phi$ on $E$, $k_\Phi$ is a \textit{nontrivial} extension of $k$. But if $\mathfrak{U}=\mathfrak{U}_\alpha$ is an arbitrary universe, and $\pi$ the first cardinal in $\mathfrak{U}$ which does not belong to a strictly smaller universe $\mathfrak{U}'\subset\mathfrak{U}$, I presume that for every $E\in\mathfrak{U}$, every $\pi$-ultrafilter on $E$ is trivial. % [sic] $k_E$ est une extension -> $k_\Phi$ est une extension, the extension of $k$ attached to the ultrafilter $\Phi$ is $k_\Phi$
\DERmark{ENV1574}


\printedpage{196}
% head: A. Grothendieck : Dérivateurs, Ch. XVIII, Édition des Universités de Montpellier II et Paris VII

\subsection*{Application to an example of a nonaccessible subcategory.}

Let $\mathcal{M}$ be a $\mathfrak{U}$-accessible category, and $a\in\operatorname{Ob}\mathcal{M}$, whence a functor
\[
\begin{aligned}
f^o\ :\ \mathcal{M}^o\ &\longrightarrow\ \mathrm{Ens}\\
x\ &\longmapsto\ \operatorname{Hom}_{\mathcal{M}}(x,a)\ \overset{\mathrm{def}}{=}\ a(x)\ ,
\end{aligned}
\]
\DERmark{ENV1575}

a functor commuting with arbitrary $\limproj$, hence such that $f:\mathcal{M}\longrightarrow(\mathrm{Ens})^o$ commutes with arbitrary $\limind$, and in particular $f$ is accessible. Let us consider the strictly full subcategory
\msp{page 12}
\[
\underbrace{\mathrm{Ens}^*}_{=\,\mathrm{Ens}\backslash\{\emptyset\}}\ \subset\ \mathrm{Ens}
\]
of nonempty sets, we know that it is accessible, but it is not stable under $\limproj$ of type $\mathrm{L}_\pi$, however large we take $\pi$. % [sic] n'est par stable -> n'est pas stable, typo

Let us consider
\[
\underbrace{\underline{S}\ =\ \underline{S}_a}_{\overset{\mathrm{def}}{=}\,f^{-1}(\mathrm{Ens}^*)}\ \subset\ \mathcal{M}\ ,
\]
\DERmark{ENV1576}

strictly full subcategory of $\mathcal{M}$ consisting of the $x\in\operatorname{Ob}\mathcal{M}$ such that
\[
a(x)\ \neq\ \emptyset\ ,\qquad\quad \text{i.e. } x\leqslant a \text{ for the canonical preorder of } \mathcal{M}\ .
\]
It is therefore an \textit{open} subcategory of $\mathcal{M}$. Its closed complement $\underline{S}'=\underline{S}'_a$ is the closed subcategory which is the inverse image, under $f$, of the strictly full subcategory of Ens reduced to the single object $\{\emptyset\}$. This subcategory of Ens is the standard final category. It is stable in Ens under arbitrary $\limind$, and under arbitrary filtering $\limproj$. Thus, $\underline{S}'$ is stable in $\mathcal{M}$ under arbitrary filtering $\limind$. Thus, $\underline{S}'$ is stable in $\mathcal{M}$ under arbitrary filtering $\limind$, in particular, it is $\mathrm{L}_\pi$-stable for $\pi$ sufficiently large. Let us make this explicit:

\msp{page 13}
\[
\begin{array}{cl}
\underbrace{\underline{S}'\ =\ \underline{S}'_a}_{}\ \subset\ \mathcal{M} & \\
\Big\|\ \text{\scriptsize def} & \text{closed subcategory of } \mathcal{M}\\
f^{-1}(\{\emptyset\})\ = & \text{consisting of the } x\ \in\ \operatorname{Ob}\mathcal{M} \text{ such}\\
 & \text{that } a(x)=\emptyset.
\end{array}
\]

Since $\underline{S}'$, being closed in $\mathcal{M}$, is stable under nonempty $\limind$ and in particular $\mathrm{L}_\pi$-stable for $\pi$ large, there is more reason to expect $\underline{S}'$ to be accessible in $\mathcal{M}$ than for its complement $\underline{S}$, for which it is not clear that it has this stability property, but only stability under nonempty $\limproj$, and in particular under direct factors.

In fact, we shall see that it can happen that $\underline{S}$ and $\underline{S}'$ are both nonaccessible in $\mathcal{M}$:

\textbf{Proposition 4.} \textit{Let $\mathfrak{U}$ be a universe, containing an infinite cardinal $\pi_0$ such that for every $E\in\mathfrak{U}$, every $\pi_0$-ultrafilter on $E$ is trivial. Let then $k\in\mathfrak{U}$ be a field of cardinality $\pi'_0\geqslant\pi_0$, and let $\mathcal{M}$ be the category}
\printedpage{197}
% head: A. Grothendieck : Dérivateurs, Ch. XVIII, Édition des Universités de Montpellier II et Paris VII
\msp{page 14}
\textit{of commutative $k$-algebras. Thus $\mathcal{M}$ is a $\mathfrak{U}$-paratopos, admitting $k$ as initial object. Let $\underline{S}'$ be the strictly full subcategory, stable under small nonempty $\limind$, consisting of the objects $A$ of $\mathcal{M}$ such that $\operatorname{Hom}_{\mathcal{M}}(A,k)=\emptyset$, and $\underline{S}$ the complementary strictly full subcategory, consisting of the objects $A$ such that $\operatorname{Hom}_{\mathcal{M}}(A,k)\neq\emptyset$. Then $\underline{S}'$ is a subcategory} closed \textit{in $\mathcal{M}$, and $\underline{S}$ the complementary open subcategory, hence $\underline{S}$ (just as $\underline{S}'$) is stable under direct factors. This being so:} % [sic] second $\operatorname{Hom}_{\mathcal{M}}(A,k)=\emptyset$ -> $\operatorname{Hom}_{\mathcal{M}}(A,k)\neq\emptyset$, $\underline{S}$ is the complement of $\underline{S}'$
\DERmark{ENV1577}

\begin{enumerate}[label={}]
\item[\textit{a)}] \textit{There is no cardinal $\pi$ such that $\underline{S}$ is $L_\pi$-stable in $\mathcal{M}$, a fortiori $\underline{S}$ is not accessible in $\mathcal{M}$.}
\item[\textit{b)}] (\footnote[207]{?, not proved.}). \textit{Although $\underline{S}'$ is $L_0$-stable in $\mathcal{M}$ (i.e. stable under filtering $\limind$, and even, as we have said, under every nonempty $\limind$), $\underline{S}'$ is nevertheless not accessible in $\mathcal{M}$.}
\end{enumerate}
\noindent\textsc{Proof}. Assertion a) will follow from the

\textbf{Corollary 1.} \textit{Under the conditions of proposition 4, let $E\in\mathfrak{U}$ be an infinite set, $\Phi_0$ the filter of complements of finite subsets of $E$. Then $\operatorname{Hom}_{\mathcal{M}}(k_{\Phi_0},k)=\emptyset$, a fortiori, for every filter $\Phi$ on $E$ finer than $\Phi_0$, we have $\operatorname{Hom}_{\mathcal{M}}(k_\Phi,k)=\emptyset$.}

\msp{page 15}
Here we have put, of course, for every filter $\Phi$ on $E$,
\[
k_\Phi\ =\ \limind_{J\in\Phi}k_J\ .
\]
We know that for every commutative $k$-algebra $\mathcal{A}$, $\operatorname{Hom}(\mathcal{A},k)$ is in 1-1 correspondence with the subset of $\operatorname{Max}(\mathcal{A})$ consisting of the maximal ideals $\mathfrak{m}$ of $\mathcal{A}$ such that $k\longrightarrow\mathcal{A}/\mathfrak{m}$ is an isomorphism. If $\mathcal{A}$ is of the form $k_\Phi$, $\Phi$ a filter on $E$, this set is in one-to-one correspondence with the set of $\pi'_0$-ultrafilters $\Psi$ on $E$ finer than $\Phi$, where $\pi'_0=\operatorname{card}(k)$ (proposition 3, and cor.~3 of prop.~2, page 7). Since $\pi'_0\geqslant\pi_0$, such an ultrafilter is trivial by hypothesis. Thus we find
\[
\operatorname{Hom}_{\mathcal{M}}(k_\Phi,k)\ \simeq\ \text{set of } \xi\in E \text{ such that } \xi\in\textstyle\bigcap_{A\in\Phi}A\ . \tag*{$(*)$}
\]
\DERmark{ENV1578}

But saying that $\Phi$ is finer than $\Phi_0$ also means that for every $\xi\in E$, we have $E\backslash\{\xi\}\in\Phi$, i.e. $\exists\,A\in\Phi$ with $A\subset E\backslash\{\xi\}$, i.e. $\xi\notin\bigcap_{A\in\Phi}A$. Hence we then have $\operatorname{Hom}_{\mathcal{M}}(k_\Phi,k)=\emptyset$.
\DERmark{ENV1579}


The conclusion of corollary 1 is expressed by the relation
\[
k_\Phi\ \notin\ \operatorname{Ob}\underline{S}\ ,\qquad\quad \text{i.e. } k_\Phi\in\operatorname{Ob}\underline{S}'\ ,
\]
on the other hand, we have by definition
\[
k_\Phi\ =\ \limind_{J\in\Phi}k_J\ ,
\]
\msp{page 16}
and it is clear (since these $J\in\Phi$ are $\neq\emptyset$), that the $k_J$ in question are in $\underline{S}$. Thus, for a filter $\Phi$ as in corollary 1, we find a direct system of objects of $\underline{S}$, indexed by $\Phi^o$, whose direct limit is not in $\underline{S}$. Thus a) of proposition 4 follows by noting that for every cardinal $\pi\in\mathfrak{U}$, we can find a set $E\in\mathfrak{U}$, and a filter $\Phi$ on $E$ finer than $\Phi_0$, such that the ordered set $\Phi^o$ is large relative to $\pi$. To find such a $\Phi$, it suffices to take $E$ of cardinality $\pi'>\pi$, and for $\Phi$ the set
\printedpage{198}
% head: A. Grothendieck : Dérivateurs, Ch. XVIII, Édition des Universités de Montpellier II et Paris VII
of complements of subsets of $E$ of cardinality $\leqslant\pi$, so that $\Phi^o$ is isomorphic to the ordered set of subsets of $E$ of cardinality $\leqslant\pi$, which is indeed large relative to $\pi$. This proves a) of proposition 4.

Let us now prove b). Let us note the

\textbf{Lemma.} \textit{Let $\mathcal{M}$ be an accessible category stable under $\limind$ of type $I_\alpha$, $\alpha$ a given limit ordinal, $\underline{S}'$ a closed subcategory, $\pi$ a useful cardinal}
\msp{page 17}
\textit{for $\mathcal{M}$, such that} $\operatorname{card}\alpha\leqslant\pi$\textit{. For every $x\in\operatorname{Ob}\underline{S}'$, let us choose an $\omega$-$\pi$-adapted representation ($\omega$ being the first infinite ordinal)}
\[
x\ \simeq\ \limind_{i\in I_x}x_i
\]
\textit{(thus $I_x$ large relative to $\pi$, stable under $\limind$ of type $I_\alpha$, $i\longmapsto x_i:I\longrightarrow\mathcal{M}$ commutes with the said limits, finally the $x_i\in\operatorname{Ob}\mathcal{M}_\pi$). For $\underline{S}'$ to be $\pi$-accessible in $\mathcal{M}$, it is necessary and sufficient that for every $x\in\operatorname{Ob}\underline{S}'$, there exist an $i\in I_x$ such that $x_i\in\operatorname{Ob}\underline{S}'$.}

Indeed, note that $\underline{S}'$, being closed, is stable in $\mathcal{M}$ under every nonempty $\limind$ representable in $\mathcal{M}$, and in particular under $\limind$ of type $I_\alpha$. So we can apply 4.13.22 (p.~342) to $\mathcal{M}$, $\underline{S}'$, $\pi$ and to $\Phi=\{I_\omega\}$. Then a) $\Longrightarrow$ c) implies that if $\underline{S}'$ is $\pi$-accessible in $\mathcal{M}$, then for every $x\in\operatorname{Ob}\underline{S}'$, the set $I_{x,S'}$ of the $i\in I_x$ such that $x_i\in\operatorname{Ob}\underline{S}'$ is cofinal in $I_x$. This implies that it is nonempty. (And conversely too, $\underline{S}'$ being a closed subcategory in $\mathcal{M}$, since if $x_i\in\underline{S}'$, then for $j\in\operatorname{Ob}I_x$ with $i\leqslant j$, we have $x_j\in\underline{S}'$, i.e. $I_{x,\underline{S}'}$ contains the essential image of $i\backslash I_x$, which is cofinal in $I$ \dots) This proves the ``only if''. Let us prove the ``if''. So let us prove that $\underline{S}'$ is equal to the essential image of $\operatorname{Ind}_\pi(\underline{S}'(\pi))$. It contains it, since $\underline{S}'$ is stable under $\limind$ of type $\mathrm{L}_\pi$. % Doubt resolved against PDF leaf 40: the first S-prime subscript is not underlined, retained. % [sic] $x_i\in\operatorname{Ob}\underline{S}$ -> $x_i\in\operatorname{Ob}\underline{S}'$, the set $I_{x,\underline{S}'}$ concerns the $x_i$ in $\underline{S}'$
\DERmark{ENV1580}


\msp{page 18}
So it remains to prove that every $x\in\operatorname{Ob}\underline{S}'$ is in the essential image. Now the hypothesis assures us that $x\simeq\limind\limits_{i\in I_x}x_i$ is isomorphic to $\limind\limits_{i\in I_{x,\underline{S}'}}x_i$ (for $I_{x,\underline{S}'}$ is cofinal in $I_x$), where this time the $x_i$ are in $\underline{S}'$. As $I_x$ large relative to $\pi$, likewise the cofinal full subcategory $I_{x,\underline{S}'}$, so we win.

\textbf{Corollary.} \textit{Under the preceding conditions, if $\underline{S}'$ is $\pi$-accessible, the same is true of the complementary open subcategory $\underline{S}$,} provided \textit{that $\underline{S}$ is stable in $\mathcal{M}$ under $\limind$ of type $L_\pi$.}

We must prove that $\underline{S}$ is equal to the essential image of $\operatorname{Ind}_\pi\underline{S}(\pi)$. The ``provided'' hypothesis ensures that this essential image is $\subset\underline{S}$. Conversely, every [\,?] $x\in\operatorname{Ob}\underline{S}$ is in this essential image. We will have
\[
x\ \simeq\ \limind_{i\in I}x_i\ ,\qquad\quad x_i\in\operatorname{Ob}\mathcal{M}_\pi,\ I=\mathcal{M}_\pi/x,\ \text{large relative to } \pi\ ,
\]
and since $x\in\operatorname{Ob}\underline{S}$ and $\underline{S}$ open, the $x_i$ are in $\underline{S}$, hence in $\underline{S}(\pi)$, OK.

In the case which concerns us, where $\mathcal{M}=(k\text{-alg.})$, as $\underline{S}$ is precisely \textit{not} stable in $\mathcal{M}$ under $\mathrm{L}_\pi$ (whatever $\pi$) (that is what made us conclude that
\msp{page 19}
$\underline{S}$ is not accessible in $\mathcal{M}$), we therefore cannot infer from this inaccessibility of $\underline{S}$, that of $\underline{S}'$. Let us note the


\endgroup
% END UNIT der-ch18/der-ch18_p196-200.tex
% BEGIN UNIT der-ch18/der-ch18_p201-201.tex
\begingroup\makeatletter%
\newcommand{\printedpage}[1]{\marginpar{\footnotesize p.~#1}}%
\newcommand{\msp}[1]{\marginpar{\footnotesize [#1]}}%
\newcommand{\Hom}{\operatorname{Hom}}%
\newcommand{\card}{\operatorname{card}}%
\makeatother\let\printedpage\derpp\let\msp\dermsp%
\let\mspin\DERmspin%


\printedpage{199}
% head: A. Grothendieck : Dérivateurs, Ch. XVIII, Édition des Universités de Montpellier II et Paris VII
% foot: Version 17 décembre 2025

\noindent\textbf{Corollary.} \textit{Let $\pi$ be an infinite cardinal. For the closed subcategory $\underline{S}'$ of $\mathcal{M}$ (consisting of commutative $k$-algebras $A$ such that $\Hom_{\mathcal{M}}(A,k)=\emptyset$) to be $\pi$-accessible, it is necessary and sufficient that for every $k$-algebra $A$ (of cardinality ${}>\pi$), such that $\Hom_{\mathcal{M}}(A,k)=\emptyset$, there exist a subalgebra $A_i$} of cardinality $\leqslant\pi$, \textit{such that $\Hom_{\mathcal{M}}(A_i,k)=\emptyset$.} % [sic] tout k-algèbre -> toute k-algèbre, « algèbre » est féminin

\noindent(We apply the lemma, taking for $\alpha$ the first infinite ordinal.)

\noindent\textbf{Question.} \textit{Does this hold, simply taking $\pi\geqslant\card k$~?} The question of the existence of a cardinal $\pi$, having the property indicated above, seems delicate to me, and I do not know the answer. If one wants to construct a counterexample for every $\pi$, the preceding developments suggest taking for $A$ an infinite product of nontrivial extensions of $k$, or a ring $k_{\Phi}$ as in the proof of proposition 4 a).


\endgroup
% END UNIT der-ch18/der-ch18_p201-201.tex
\clearpage\setcounter{footnote}{0}\phantomsection\addcontentsline{toc}{section}{Chapter XIX}\markboth{Chapter XIX}{Chapter XIX}
% BEGIN UNIT der-ch19/der-ch19_p001-005.tex
\begingroup\makeatletter%
\newcommand{\printedpage}[1]{\marginpar{\footnotesize p.~#1}}%
\newcommand{\msp}[1]{\marginpar{\footnotesize [#1]}}%
\newcommand{\Qo}{\mathrel{\leftharpoondown\joinrel\relbar\joinrel\rightharpoonup}}   % lower half-head left, upper half-head right
\newcommand{\Qorev}{\mathrel{\leftharpoonup\joinrel\relbar\joinrel\rightharpoondown}} % mirrored form of (1.2)
\newcommand{\Qog}{\mathrel{\leftarrow\joinrel\relbar\joinrel\rightharpoonup}}         % full head left
\newcommand{\Qod}{\mathrel{\leftharpoondown\joinrel\relbar\joinrel\rightarrow}}       % full head right
\newcommand{\Qogd}{\mathrel{\leftarrow\joinrel\relbar\joinrel\rightarrow}}            % full heads both sides
\newcommand{\Qx}[2]{\mathrel{#1\joinrel\mapstochar#2}}%
\newcommand{\Qxx}[2]{\mathrel{#1\joinrel\mapstochar\mkern2mu\mapstochar#2}}%
\newcommand{\tld}{\tilde{}}   % left Q-closure mark (small tilde in exponent)
\newcommand{\rcl}{\bar{}}       % right Q-closure mark (horizontal bar in exponent)
\makeatother\let\printedpage\derpp\let\msp\dermsp%
\let\mspin\DERmspin%


\printedpage{(unnumbered)}
% head: (none; title page)

\begin{center}
{\Large\textsf{THE DERIVATORS}}

\bigskip

\textbf{ALEXANDRE GROTHENDIECK}

\bigskip\bigskip

\textbf{Chapter XIX}

\medskip

\textbf{Models (4)}
\end{center}

\par

\printedpage{(unnumbered)}
% head: (none; editorial note)

\par\noindent\textbf{Note by the editors of the transcribed edition.}\par
This text was deciphered and transcribed in \LaTeXe\ by M. Künzer. It was edited by M. Künzer, J. Malgoire, and G. Maltsiniotis. The transcription is as faithful as possible to the manuscript. Only a few obvious minor corrections have been made. For the rare additions or comments of the editors, as well as for the original numbering of the manuscript pages, typewriter characters \texttt{[typewriter]} in square brackets are used.

This edition is provisional. Remarks, comments and corrections are welcome. Send a message to:

\textbf{[e-mail address removed]}

\hfill G. Maltsiniotis\hspace*{2cm}

\par

\printedpage{1}
% head: A. Grothendieck : Dérivateurs, Ch. XIX, Édition des Universités de Montpellier II et Paris VII
% foot: Version 11 février 2019

\msp{page 1}

\section*{1\quad Reminders on Quillen couples (special Quillen couples)}

I am going to take up again the notations and terminology for Quillen's LLP \texttt{[left lifting property]} and RLP \texttt{[right lifting property]} relations. I write (for a category $\mathcal{M}$, and $\Phi, \Psi \subset \mathrm{Fl}(\mathcal{M})$)
\[
\Phi \Qo \Psi \qquad \iff \qquad (\Phi \subset \Psi^{*}) \qquad \iff \qquad (\Psi \subset \Phi_{*})\ ,
\tag*{(1.1)}
\]
i.e. $\Phi$ has the LLP with respect to $\Psi$, or again $\Psi$ the RLP with respect to $\Phi$. I also say that $(\Phi, \Psi)$ is a \textit{couple} which is Q-\textit{orthogonal.} This notion is not symmetric, but I sometimes allow myself (exceptionally) to write
\[
\Psi \Qorev \Phi \qquad \text{as a synonym for} \qquad \Phi \Qo \Psi\ .
\tag*{(1.2)}
\]
$\Phi$ is Q-\textit{orthogonal on the left} to $\Psi$, $\Psi$ is Q-\textit{orthogonal on the right} to $\Phi$, and Q-orthogonality on the left (resp. on the right) is denoted by the half-arrowhead pointing downward (resp. upward). I write
\[
\left\{
\begin{aligned}
&\Phi \quad \Qog \quad \Psi \quad \iff \quad (\Phi = \Psi^{*})\\
&\Phi \quad \Qod \quad \Psi \quad \iff \quad (\Psi = \Phi_{*})\\
&\Phi \quad \Qogd \quad \Psi \quad \iff \quad (\Phi = \Psi^{*} \ \text{ and } \ \Psi = \Phi_{*})\ .
\end{aligned}
\right.
\tag*{(1.3)}
\]
These three cases of strengthened Q-orthogonality thus respectively imply $\Phi = \Phi^{\tld}$
\DERmark{ENV1581}

resp. $\Psi = \Psi^{\rcl}$ resp. $\Phi = \Phi^{\tld}$ and $\Psi = \Psi^{\rcl}$. We say that $\Phi$ is Q-\textit{left closed} if $\Phi = \Phi^{\tld}$, that % reading resolved: DER_XIX_corrections.csv
$\Psi$ is Q-\textit{right closed} if $\Psi = \Psi^{\rcl}$, that the Q-orthogonal couple $(\Phi, \Psi)$ is Q-\textit{left closed} (resp. \textit{right}) if $\Phi$ (resp. $\Psi$) is left Q-closed (resp. right). We say that it is Q-\textit{closed} if the two conditions % reading resolved: DER_XIX_corrections.csv

\msp{page 2}

are satisfied simultaneously. When no confusion is to be feared, I shall sometimes say simply that $\Phi$ is Q-closed, or that $\Psi$ is Q-closed, without specifying whether it is on the left or % [sic] où -> ou, conjunction "or" intended
on the right. But for a Q-orthogonal couple $(\Phi, \Psi)$ (where I prefer to write it $\Phi \Qo \Psi$), Q-closed % [sic] ou -> où, relative adverb intended
always means left and right Q-closed.

When a Q-orthogonal couple $\Phi \Qo \Psi$ has Quillen's factorization property, I indicate it by the vertical bar crossing the horizontal, thus
\[
\Phi \ \Qx{\leftharpoondown}{\rightharpoonup} \ \Psi\ ,
\tag*{(1.4)}
\]
and we define in the same way the meaning of the graphisms
\[
\left\{
\begin{aligned}
&\Phi \quad \Qx{\leftharpoondown}{\rightharpoonup} \quad \Psi\\
&\Phi \quad \Qx{\leftarrow}{\rightharpoonup} \quad \Psi\\
&\Phi \quad \Qx{\leftharpoondown}{\rightarrow} \quad \Psi\\
&\Phi \quad \Qx{\leftarrow}{\rightarrow} \quad \Psi\ .
\end{aligned}
\right.
\tag*{(1.5)} % [sic] (1.3) -> (1.5), number already used above and (1.5) is missing
\]
\DERmark{ENV1582}


\par

\printedpage{2}
% head: A. Grothendieck : Dérivateurs, Ch. XIX, Édition des Universités de Montpellier II et Paris VII
% foot: Version 11 février 2019

A Q-\textit{orthogonal} Q-\textit{factorizing} couple $\Phi \Qx{\leftharpoondown}{\rightharpoonup} \Psi$ is called a \textit{Quillen precouple.} If it is moreover Q-closed, i.e. if we have $\Phi \Qx{\leftarrow}{\rightarrow} \Psi$, we say that it is a \textit{Quillen couple.}

$\Phi$ is said to be Q-\textit{left saturated} if it contains the isomorphisms and is stable under cobase change, under (small) transfinite composition, and under direct factors. Dually $\Psi$ is said to be Q-\textit{right saturated} if it

\msp{page 3}

contains the isomorphisms, is stable under base change, under transfinite composition, and under direct factors. If $\Phi$ is left Q-saturated, it is stable under (small) direct sums, i.e. if the $f_i : A_i \longrightarrow B_i$ ($i \in I$, $I$ small) are in $\Phi$, $f = \coprod f_i : \coprod A_i \longrightarrow \coprod B_i$ also, and likewise for (small) finite or infinite amalgamated sums (when we work in an $A\backslash\mathcal{M}$ instead of $\mathcal{M}$). We have the dual statement for $\Psi$, when $\Psi$ is right Q-saturated.

In fact, there are weaker notions of Q-saturation than the one I have just given: Q-\textit{simple saturation} (on the left: $\Phi$ contains the isomorphisms, stable under binary composition, under cobase change), Q-\textit{transfinite saturation} (we require in addition stability under transfinite composition), finally Q-\textit{total saturation,} or Q-saturation for short (when we require in addition stability under direct factors). We have the implications
\[
\left\{
\begin{aligned}
&\text{left Q-closed}\\
&\Downarrow\\
&\text{Q-saturated (total)}\\
&\Downarrow\\
&\text{transfinitely Q-saturated}\\
&\Downarrow\\
&\text{simply Q-saturated.}
\end{aligned}
\right.
\tag*{(1.6)}
\]
For $\Phi \subset \mathrm{Fl}(\mathcal{M})$, we have its (left) Q-closure $\Phi^{\tld}$ (which is the smallest left Q-closed subset of $\mathrm{Fl}(\mathcal{M})$ containing $\Phi$) and likewise its (total) left Q-saturation, its transfinite Q-saturation, its simple Q-saturation

\msp{page 4}

with the tautological inclusions
\[
\Phi \ \subset \
\underbrace{\Phi'}_{\substack{\text{Q-saturation}\\ \text{simple on the left}}}
\ \subset \
\underbrace{\Phi''}_{\substack{\text{Q-saturation}\\ \text{transfinite on the left}}}
\ \subset \
\underbrace{\Phi'''}_{\substack{\text{Q-saturation}\\ \text{total on the left}}}
\ \subset \
\underbrace{\Phi^{\tld}}_{\substack{\text{Q-closure}\\ \text{on the left}}}
\ ,
\tag*{(1.7)}
\]
and dually for $\Psi$. Between $\Phi$ and $\Phi'$ there is also a $\Phi'_0$,
\[
\Phi \ \subset \ \Phi'_0 \ \subset \ \Phi'\ ,
\tag*{(1.8)}
\]
which is the Q-\textit{saturation under cobase change} (also containing, by definition, all isomorphisms). To construct all these saturations, we proceed as follows:

\par

\printedpage{3}
% head: A. Grothendieck : Dérivateurs, Ch. XIX, Édition des Universités de Montpellier II et Paris VII
% foot: Version 11 février 2019

\begin{enumerate}[leftmargin=*]
\item[1)] $\Phi'_0 \quad = \quad$ set of arrows $i : A \longrightarrow B$ fitting into a cocartesian square
\[
\begin{tikzcd}
A_0 \arrow[r, "i_0"] \arrow[d] & B_0 \arrow[d]\\
A \arrow[r, "i"] & B\ ,
\end{tikzcd}
\]
with $i_0 \in \Phi$ or $i_0$ an identity (in which case we find for $i$ the isomorphisms).

\item[2)] $\Phi' \quad = \quad$ set of multiple composite arrows
\[
i_n \, i_{n-1} \ \cdots \ i_1\ ,
\]
with $i_\alpha \in \Phi'_0$.

\item[3)] $\Phi'' \quad = \quad$ set of transfinite composites of arrows $\in \Phi'_0$.
\end{enumerate}

We feel like putting
\[
\Phi''' \ = \ \Phi^{\prime\prime\natural} \ \overset{\mathrm{def}}{=} \ \text{set of direct factors of elements of } \Phi'',
\]
\DERmark{ENV1583}

and it is a fact that $\Phi^{\prime\prime\natural} \subset \Phi'''$ trivially, but it is not clear in general that we have % [sic] par -> pas, typo
\msp{page 5}
equality, for lack of being able to show that $\Phi^{\prime\prime\natural}$ is stable, not only under direct factors, but also under cobase change and under transfinite composition. However, let us recall:

\textbf{Theorem 1.1.} \textit{Suppose $\mathcal{M}$ stable under small $\varinjlim$ , and that its elements are accessible. Suppose moreover that $\Phi$ is small, more generally, that there exists a small subset $\Phi_0 \subset \Phi$ such that $\Phi \subset \Phi_0{}^{\tld}$, hence}
\[
\Psi \ \overset{\mathrm{def}}{=} \ \Phi_{*} \ = \ (\Phi_0)_{*} \qquad \text{\footnotemark}.
\tag*{(1.1.1)} % footnote mark printed as superscript (1)
\]
\DERmark{ENV1584}

\footnotetext{\textbf{N.B.} We thus have $\Phi \Qod \Psi$, a fortiori $\Phi_0 \Qod \Psi$.}
\textit{Under these conditions, if $\Phi''_0$ denotes the transfinite left \textup{Q}-saturation of $\Phi_0$, then we have $\Phi''_0 \Qx{\leftharpoondown}{\rightarrow} \Psi$, i.e. every arrow $f$ of $\mathcal{M}$ factors as $f = pi$, with $i \in \Phi''_0$, $p \in \Psi$\,\footnote{There is even \textit{functorial} factorization in $f$, a property which deserves a notation, e.g. $\Phi''_0 \Qxx{\leftharpoondown}{\rightarrow} \Psi$.}. A fortiori, if $\Phi''$ denotes the transfinite left \textup{Q}-saturation of $\Phi$, we have} % footnote mark printed as superscript (2)
\[
\Phi'' \ \Qx{\leftharpoondown}{\rightarrow} \ \Psi\ .
\]

\bigskip

\textbf{Corollary 1.2.} \textit{Under the conditions of the theorem, we have}
\[
\Phi''' \ \text{(total left Q-saturation of $\Phi$)} \ = \
\underbrace{\Phi^{\prime\prime\natural}}_{\substack{\text{set of}\\ \text{direct factors of elements}\\ \text{of $\Phi''$}}}
\ = \ \Phi^{\tld}\ .
\tag*{(1.2.1)}
\]
\textit{More generally (in appearance only!), if $\Phi_1$ is a subset of \textup{$\mathrm{Fl}(\mathcal{M})$} which is transfinitely left \textup{Q}-saturated, and such that}
\[
\Phi_0 \ \subset \ \Phi_1 \ \subset \ \Phi^{\tld} \ = \ \Phi_0{}^{\tld}\ ,
\]


\endgroup
% END UNIT der-ch19/der-ch19_p001-005.tex
% BEGIN UNIT der-ch19/der-ch19_p006-010.tex
\begingroup\makeatletter%
\derleqno%
\newcommand{\printedpage}[1]{\marginpar{\footnotesize p.~#1}}%
\newcommand{\msp}[1]{\marginpar{\footnotesize [#1]}}%
\newcommand{\Qorth}{\mathrel{\leftharpoondown\mkern-6mu\relbar\mkern-6mu\rightharpoonup}}%
\newcommand{\Qorthx}{\mathrel{\leftharpoondown\mkern-9mu{+}\mkern-9mu\rightharpoonup}}%
\newcommand{\Qorthxfull}{\mathrel{\leftharpoondown\mkern-9mu{+}\mkern-9mu\rightarrow}}%
\newcommand{\Qcouple}{\mathrel{\leftarrow\mkern-9mu{+}\mkern-9mu\rightarrow}}%
\newcommand{\clos}{\longleftrightarrow}%
\newcommand{\fl}{\longrightarrow}%
\makeatother\let\printedpage\derpp\let\msp\dermsp%
\let\mspin\DERmspin%


\printedpage{4}
% head: A. Grothendieck : Dérivateurs, Ch. XIX, Édition des Universités de Montpellier II et Paris VII
% foot: Version 11 février 2019

\msp{page 6}
\textit{then}
\begin{enumerate}[label={}]
\item[a)] \textit{one has $\Phi_1 \Qorthxfull \Psi$ and} % reading resolved: DER_XIX_corrections.csv
\item[b)] $\Phi_1^{\sim} = \Phi^{\sim} = \Phi_1^{\natural}$.
\end{enumerate}
This corollary follows from the theorem by virtue of XV prop. 4.2, of which I recall here the essentials:

\textbf{Proposition 1.3.} \textit{Let $\Phi \Qorthx \Psi$. Then $\Phi^{\sim} = \Phi^{\natural} = \Psi^{*}$, and dually one has $\Psi^{-} = \Psi^{\natural} = \Phi_{*}$.}

Let us also note, in this line of thought:

\textbf{Corollary 1.4.} \textit{Let $\Phi \Qorth \Psi$, and let $W$ be an arbitrary subset of $\mathrm{Fl}(\mathcal{M})$. In order to have $\Phi \subset W^{\natural}$, it suffices that every arrow of $\mathcal{M}$ factor as $f = pi$, with $i \in W$, $p \in \Psi$. (And this condition is obviously also necessary if $\Phi \Qorthx \Psi$, and $W$ stable under direct factors.)}

One says that a subset $\Phi_0$ of $\Phi$ is Q-\textit{generating on the left} for $\Phi$ if it satisfies the condition
\[
\Phi \ \subset \Phi_0^{\sim}\ , \tag*{(1.9)}
\]
and one says that $\Phi$ is a \textit{special} subset if $\Phi$ admits a \textit{small} left Q-generating subset. If $\Phi_0 \subset \Phi$ is small, and if $\mathcal{M}$ satisfies the conditions of theorem 1.1, then $\Phi_0$ is a Q-generating subset (on the left) if and only if its total Q-saturation (on the left) $\Phi_0'''$ contains $\Phi$ (a condition apparently stronger than $\Phi \subset \Phi_0^{\sim}$, since in general one has $\Phi_0''' \subset \Phi_0^{\sim}$). This therefore also means, via 1.2, that every \msp{page 7} element of $\Phi$ is a direct factor of a transfinite composite of arrows obtained from
\[
\Phi_0 \cup (\text{Identities of } \mathcal{M})
\]
by cobase change.

I recall also the $\pm$ tautological sorites of adjoint functors:

\textbf{Proposition 1.5.} \textit{Let}
\[
\begin{tikzcd}
\mathcal{M}' \arrow[r, shift right=1, "v^{*}"'] & \mathcal{M} \arrow[l, shift right=1, "v_{!}"']
\end{tikzcd}
\tag*{(1.5.1)}
\]
\textit{be a pair of adjoint functors, let}
\[
\Phi \ \subset\ \mathrm{Fl}(\mathcal{M})\ , \qquad \Psi' \ \subset\ \mathrm{Fl}(\mathcal{M}')\ ,
\]
\textit{then one has}
\[
(v_{!}(\Phi) \Qorth \Psi') \qquad \Longleftrightarrow \qquad (\Phi \Qorth v^{*}(\Psi'))\ . \tag*{(1.5.2)}
\]

\textbf{Corollary 1.6.} \textit{For every subset $\Phi$ of} $\mathrm{Fl}(\mathcal{M})$,
\[
(v_{!}(\Phi))_{*} \ =\ (v^{*})^{-1}(\Phi_{*})\ . \tag*{(1.6.1)}
\]

\printedpage{5}
% head: A. Grothendieck : Dérivateurs, Ch. XIX, Édition des Universités de Montpellier II et Paris VII
% foot: Version 11 février 2019

Indeed, applying 1.5 to $\Phi$ and to $\{f'\}$, where $f' \in \mathrm{Fl}(\mathcal{M}')$, one finds
\[
f' \in (v_{!}(\Phi))_{*} \qquad \Longleftrightarrow \qquad v^{*}(f') \in \Phi_{*}\ ,
\]
which is equivalent to (1.6.1). Dually:

\textbf{Corollary 1.7.} \textit{For every subset $\Psi'$ of $\mathrm{Fl}(\mathcal{M}')$, one has}
\DERmark{ENV1585}\DERmark{ENV1586}

\[
(v^{*}(\Psi'))^{*} \ =\ (v_{!})^{-1}(\Psi'^{*}). \tag*{(1.7.1)}
\]

Indeed, if $f \in \mathrm{Fl}(\mathcal{M})$, one has by (1.5.2)
\[
f \in (v^{*}(\Psi'))^{*} \qquad \Longleftrightarrow \qquad v_{!}(f) \in\Psi'^{*}\ . % [sic] \Psi^{*} -> \Psi'^{*}, cf. (1.7.1)
\]
\DERmark{ENV1587}

\textsuperscript{(}\footnote[3]{give here 1.9 page 9.}\textsuperscript{)}.

\msp{page 8}
A Q-orthogonal couple $\Phi \Qorth \Psi$ is said to be \textit{special} if it satisfies the following conditions.
\[
\left\{
\begin{array}{ll}
\text{a)} & \parbox[t]{0.7\textwidth}{$\mathcal{M}$ is stable under small $\varinjlim$ , the objects of $\mathcal{M}$ are accessible. (One then says that $\mathcal{M}$ is ``special''.)} \\[3ex]
\text{b)} & \text{The couple is closed: } \Phi \clos \Psi. \\[1ex]
\text{c)} & \Phi \text{ is \textit{special.}}
\end{array}
\right.
\tag*{(1.10)}
\]
This therefore implies, by 1.1, that $\Phi \Qcouple \Psi$ is a Quillen couple. If $\mathcal{M}$ satisfies a), the special Q-orthogonal couples in $\mathcal{M}$, or again the special Quillen couples, thus correspond one-to-one to the subsets $\Phi$ of $\mathrm{Fl}(\mathcal{M})$ which are \textit{special} and left closed (or what amounts to the same thing, left Q-saturated), by associating to such a $\Phi$ the couple $\Phi \Qcouple \Phi_{*}$.

This being so:

\marginpar{$\Vert$}% double vertical bar printed in the left margin beside Corollaire 1.8
\textbf{Corollary 1.8.} \textit{Suppose that $\mathcal{M}$, $\mathcal{M}'$ satisfy condition} a) \textit{of} (1.10), \textit{and let $(\Phi, \Psi)$ be a special Quillen couple in $\mathcal{M}$. Then $\Bigl(\Phi' = v_{!}(\Phi)^{\sim}, \Psi' = (v^{*})^{-1}(\Psi)\Bigr)$ is a special Quillen couple in $\mathcal{M}'$.}

Indeed, by corollary 1.6 one has
\[
\Phi'_{*} \ =\ (v_{!}(\Phi))_{*} \ =\ (v^{*})^{-1}(\Psi) \ =\ \Psi'\ ,
\]
hence $\Phi' \clos \Psi'$. It remains to prove that $\Phi'$ is special. Let $\Phi_0$ be a small subset of $\Phi$ such that $\Phi \subset \Phi_0^{\sim}$, put $\Phi'_0 \ =\ v_{!}(\Phi_0) \subset \Phi'$, \msp{page 9} it suffices to see that one has
\[
v_{!}(\Phi) \ \subset\ v_{!}(\Phi_0)^{\sim}\ , \tag*{($*$)}
\]
from which it will follow that one has $\Phi' \overset{\mathrm{def}}{=} (v_{!}(\Phi))^{\sim} \subset v_{!}(\Phi_0)^{\sim}$. But one has
\DERmark{ENV1588}

\[
v_{!}(\Phi_0)^{\sim} \ =\ (v_{!}(\Phi_0)_{*})^{*}\ ,
\]
\printedpage{6}
% head: A. Grothendieck : Dérivateurs, Ch. XIX, Édition des Universités de Montpellier II et Paris VII
% foot: Version 11 février 2019
and
\[
v_{!}(\Phi_0)_{*} \ =\ (v^{*})^{-1}(\underbrace{\Phi_{0\,*}}_{\Psi}) \ =\ (v^{*})^{-1}(\Psi) \ \overset{\mathrm{def}}{=}\ \Psi'\ ,
\]
\DERmark{ENV1589}

hence
\[
v_{!}(\Phi_0)^{\sim} \ =\ \Psi'^{*} \ =\ \Phi' \ =\ v_{!}(\Phi)^{\sim}\ ,
\]
whence ($*$). This argument does not appeal to the ``special'' hypothesis on $\mathcal{M}$, and can be made explicit in the

\textbf{Corollary 1.9.} \textsuperscript{(}\footnote[4]{should come before 1.8.}\textsuperscript{)} \textit{Under the general conditions of} 1.5 \textit{($v_{!} : \mathcal{M} \fl \mathcal{M}'$ being a functor admitting a right adjoint), let $\Phi_0$, $\Phi$ be two subsets of} $\mathrm{Fl}(\mathcal{M})$ \textit{such that $\Phi \subset \Phi_0^{\sim}$ (or what amounts to the same thing, $\Phi^{\sim} \subset \Phi_0^{\sim}$). Then $v_{!}(\Phi) \subset v_{!}(\Phi_0)^{\sim}$, i.e. $v_{!}(\Phi)^{\sim} \subset v_{!}(\Phi_0)^{\sim}$.}

Indeed, this relation is also equivalent to $v_{!}(\Phi_0)_{*} \ \subset\ v_{!}(\Phi)_{*}$, that is (by 1.6) $(v^{*})^{-1}(\Phi_{0\,*}) \ \subset\ (v^{*})^{-1}(\Phi_{*})$, and is implied by $\Phi_{0\,*} \ \subset\ \Phi_{*}$. The latter on the other hand is equivalent to the hypothesis $\Phi^{\sim} \subset \Phi_0^{\sim}$, q.e.d.

\msp{page 10}
If $\Psi$ is a subset of $\mathrm{Fl}(\mathcal{M})$, one obtains from it for every $I$ in Cat a subset
\[
\left\{
\begin{array}{l}
\Psi(I) \subset \mathrm{Fl}(\mathcal{M}(I)), \quad \text{with}\ \ \mathcal{M}(I) = \underline{\mathrm{Hom}}(I^{\circ}, \mathcal{M})\ , \\[1ex]
f \in \Psi(I) \iff f(i) \in \Psi \ \ \forall\, i \in \mathrm{Ob}\, I\ .
\end{array}
\right.
\tag*{(1.11)}
\]
If one introduces the discrete category $I_0$ underlying $\mathrm{Ob}\, I$, and the canonical functor
\[
v_I = v \ :\ I_0 \fl I\ , \tag*{(1.12)}
\]
one finds a pair of adjoint functors
\[
\begin{tikzcd}[row sep=small]
\mathcal{M}(I) \arrow[r, shift right=1, "v^{*}"'] & \mathcal{M}(I_0) \arrow[l, shift right=1, "v_{!}"'] \\
( & ) \arrow[l, dotted, "v_{*}"]
\end{tikzcd}
\ \simeq\ \mathcal{M}^{I_0}\ ,
\tag*{(1.13)}
\]
and one has
\[
\Psi(I) \ =\ (v^{*})^{-1}(\Psi(I_0))\ .
\]
Suppose now that $\Phi$ contains the identities and that one has
\DERmark{ENV1590}

\[
\left\{
\begin{array}{rcl}
\Psi & = & \Phi_{*} \\[1ex]
\text{whence}\ \ \Psi(I_0) & = & \Phi(I_0)_{*}
\end{array}
\right.
\tag*{(1.14)}
\]
(immediate verification), one then concludes by 1.6
\[
\Psi(I) \ =\ (v_{!}(\Phi(I_0)))_{*} \ =\ (\Phi\langle I\rangle)_{*}\ , \tag*{(1.15)}
\]
with
\[
\Phi\langle I\rangle \ \overset{\mathrm{def}}{=}\ v_{!}(\Phi(I_0))^{\sim} \ =\ \Psi(I)^{*}\ . \tag*{(1.16)}
\]
\DERmark{ENV1591}

If one has a functor
\[
u \ :\ I \fl I'\ , \tag*{(1.17)}
\]
whence a triple of adjoint functors
\printedpage{7}
% head: A. Grothendieck : Dérivateurs, Ch. XIX, Édition des Universités de Montpellier II et Paris VII
% foot: Version 11 février 2019
\msp{page 11}
\[
\left\{
\begin{tikzcd}[column sep=large]
\mathcal{M}(I) \arrow[r, shift left=3, "u_{!}"] \arrow[r, shift right=3, "u_{*}"] & \mathcal{M}(I') \arrow[l, "u^{*}"']
\end{tikzcd}
\right.
,
\tag*{(1.18)}
\]
one will have in a $\pm$ tautological way
\[
\left\{
\begin{array}{rcl}
u_{!}(\Phi\langle I\rangle) & \subset & \Phi\langle I'\rangle \\
u^{*}(\Psi(I')) & \subset & \Psi(I)\ .
\end{array}
\right.
\tag*{(1.19)}
\]
Corollary 1.8 gives:

\textbf{Proposition 1.10.} \textit{Suppose that $\Phi \Qcouple \Psi$ is a} special \textit{Quillen couple in $\mathcal{M}$, i.e. of the form $(\Phi, \Phi_{*})$, with $\Phi$ a special left Q-closed subset in $\mathcal{M}$ ($\mathcal{M}$ satisfying condition} (1.10. a))\textit{. Then for every $I$ in} Cat, \textit{$(\Phi\langle I\rangle, \Psi(I))$ is a special couple in $\mathcal{M}(I)$, in particular it satisfies the factorization theorem.}

Let us consider the dual construction, which leads us to introduce, starting from a couple $\Phi \clos \Psi$ in $\mathcal{M}$, the couple
\[
\Phi(I) \ \clos\ \Psi[I]\ , \tag*{(1.20)}
\]
where $\Phi(I)$ is defined by the general formula (1.11) (where $\Psi$ is replaced by $\Phi$), and where $\Psi[I]$ is defined either as $\Phi(I)_{*}$, or as
\[
\Psi[I] \ =\ v_{*}(\Psi(I_0))^{-} \ =\ \Phi(I)_{*}\ . \tag*{(1.21)}
\]

\msp{page 12}
One then has, for $u : I \fl I'$
\[
\left\{
\begin{array}{rcl}
u^{*}(\Phi(I')) & \subset & \Phi(I) \\
u_{*}(\Psi[I]) & \subset & \Psi[I']\ .
\end{array}
\right.
\tag*{(1.22)}
\]
There still arises the question whether the couple $(\Phi(I), \Psi[I])$ satisfies the factorization theorem, under suitable hypotheses on $(\Phi, \Psi)$. I can only venture the

\textbf{Question-conjecture 1.11.} \textit{Suppose that $(\Phi, \Psi)$ is a} special \textit{Quillen couple, is the same true of $(\Phi(I), \Psi[I])$ for every $I$ in} (Cat)\textit{? In other words, if $\mathcal{M}$ is special, condition} (1.10. a))\textit{, and if $\Phi$ is a} special \textit{left Q-closed subset of $\mathcal{M}$, i.e. admitting a small Q-generating subset, is it true that for every $I$ in} Cat, \textit{$\Phi(I)$ is also special? I conjecture that yes, making if need be additional hypotheses on $\mathcal{M}$ (existence of a cardinal filtration?). E.g. if $\Phi_0$ is Q-generating in $\Phi$, and small, then is $\Phi_0(I)$ Q-generating in $\Phi(I)$?}

\textbf{Proposition 1.12} \textsuperscript{(}\footnote[5]{cf. XVII 6.4 (p. 117), but where the notations for $\Phi\langle I\rangle$, $\Psi(I)$ are reversed.}\textsuperscript{)}. \textit{Let $\Phi \clos \Psi$ in $\mathcal{M}$, let $u : I \fl I'$ be an arrow in} Cat, \textit{whence}
\[
\left\{
\begin{tikzcd}[column sep=large]
\mathcal{M}(I) \arrow[r, shift left=3, "u_{!}"] \arrow[r, shift right=3, "u_{*}"] & \mathcal{M}(I') \arrow[l, "u^{*}"']
\end{tikzcd}
\right.
.
\tag*{(1.12.1)}
\]
\printedpage{8}
% head: A. Grothendieck : Dérivateurs, Ch. XIX, Édition des Universités de Montpellier II et Paris VII
% foot: Version 11 février 2019
\textit{One then has} (1.19), \textit{i.e.}
\[
\left\{
\begin{array}{rcl}
u_{!}(\Phi\langle I\rangle) & \subset & \Phi\langle I'\rangle \\
u^{*}(\Psi(I')) & \subset & \Psi(I)\ ,
\end{array}
\right.
\tag*{(1.13.1)}
\]
\textit{and moreover the equivalence}
\[
u^{*}(\Phi\langle I'\rangle)\subset\Phi\langle I\rangle \qquad \Longleftrightarrow \qquad u_{*}(\Psi(I)) \subset \Psi(I')\ . \tag*{(1.13.2)}
\]
\DERmark{ENV1592}

\textit{These equivalent relations are satisfied in each of the three cases}~\textsuperscript{(}\footnote[6]{\textbf{N.B.} One has $1^{\circ} \Longrightarrow 3^{\circ}$, $2^{\circ} \Longrightarrow 3^{\circ}$.}\textsuperscript{)} \textit{:}
\begin{enumerate}[label={}]
\item[1°)] \textit{$u$ is coétale.}
\item[2°)] \textit{$I$ is discrete.}
\item[3°)] $\forall\, i' \in \mathrm{Ob}\, I'$\textit{, every connected component of $I/i'$ admits a final object.}
\end{enumerate}
\msp{page 13}

\section*{2\quad Reminders on closed Quillen triples and special Quillen triples}

Let $(\mathcal{M}, W)$ be a \textit{category of models,} i.e. $\mathcal{M}$ equipped with a localizer $W \subset \mathrm{Fl}(\mathcal{M})$, i.e. a subset satisfying
\[
\parbox{0.75\textwidth}{$W$ contains the isomorphisms, and if $f$, $g$ are composable arrows, if two among $f$, $g$, $gf$ are in $W$, the third also \textit{(weak saturation).}}
\tag*{(2.1)}
\]
I recall that the datum of a \textit{closed Quillen triple} $(W, C, F)$\textit{,} or of a closed Quillen structure, on $\mathcal{M}$, having $W$ as set of quasi-isomorphisms, amounts to the datum of two Quillen couples
\[
C \Qcouple TF\ , \qquad TC \ \Qcouple\ F\ , \tag*{(2.2)}
\]
satisfying
\[
TC \ \subset\ C \cap W\ , \qquad TF \ \subset\ F \cap W\ , \tag*{(2.3)}
\]
plus one of the three equivalent conditions:
\[
\left\{
\begin{array}{lrcll}
\text{a)} & TC & = & C \cap W & \qquad (\text{i.e. } C \cap W \subset TC) \\[1ex]
\text{b)} & TF & = & F \cap W & \qquad (\text{i.e. } F \cap W \subset TF) \\[1ex]
\text{c)} & W & = & TF \circ TC\ . &
\end{array}
\right.
\tag*{(2.4)}
\]
Let us also recall that this implies that $W$ \textit{is strongly saturated,} and a fortiori that $W = W^{\natural}$, i.e. $W$ stable under direct factors.

Let us also recall that by virtue of (1.4), one has:


\endgroup
% END UNIT der-ch19/der-ch19_p006-010.tex
% BEGIN UNIT der-ch19/der-ch19_p011-015.tex
\begingroup\makeatletter%
\newcommand{\printedpage}[1]{\marginpar{\footnotesize p.~#1}}%
\newcommand{\msp}[1]{\marginpar{\footnotesize [#1]}}%
\newcommand{\M}{\mathcal{M}}%
\newcommand{\Fl}{\mathrm{Fl}}%
\newcommand{\Ho}{\mathrm{Ho}}%
\newcommand{\cc}{\mathrm{c}}%
\newcommand{\ddarr}{\mathrel{\twoheadleftarrow\joinrel\twoheadrightarrow}}%
\newcommand{\harp}{\mathrel{\leftharpoondown\joinrel\relbar\joinrel\rightharpoonup}}%
\makeatletter%
\renewcommand\@makefnmark{\hbox{\@textsuperscript{\normalfont(\@thefnmark)}}}%
\renewcommand\@makefntext[1]{\noindent\hbox{\@textsuperscript{\normalfont\@thefnmark}}#1}%
\newcommand{\Qcouple}{\mathrel{\leftarrow\mkern-9mu{+}\mkern-9mu\rightarrow}}%
\makeatother%
\makeatother\let\printedpage\derpp\let\msp\dermsp%
\let\mspin\DERmspin%


\printedpage{9}
% head: A. Grothendieck : Dérivateurs, Ch. XIX, Édition des Universités de Montpellier II et Paris VII
% foot: Version 11 février 2019

\msp{page 14}

\begin{itemize}[leftmargin=4em]
\item[(2.5)] Given the two couples (2.2), with $TC \subset C$ (equivalent to $TF \subset F$), for the first relation (2.3) to hold, i.e. $TC \subset W$, it is necessary, and sufficient if $W = W^{\natural}$, that every arrow $f$ of $\M$ (or only every arrow of $TC$) factor as $f = pi$, with $i \in W$, $p \in F$. Dually, for $TF \subset W$ to hold, it is necessary and sufficient if $W = W^{\natural}$, that every arrow $f$ of $\M$ (or only every arrow of $TF$) factor as $f = pi$, with $i \in C$, $p \in W$. % [sic] i \in W, f \in C -> i \in C, p \in W (?), dual of the preceding factorisation with p \in F
\end{itemize}
\DERmark{ENV1593}\DERmark{ENV1594}\DERmark{ENV1595}


We say that the Quillen structure $(W,C,F)$ is \textit{special} if the two Quillen couples (2.2) are special, i.e. if the following conditions are satisfied.
\DERmark{ENV1596}


\begin{itemize}[leftmargin=4em]
\item[(2.6)]
\begin{itemize}
\item[a)] $\M$ is stable under small $\varinjlim$, and its elements are accessible ($\M$ ``special'').
\item[b)] There exist \textit{small} subsets $C_0 \subset C$, $TC_0 \subset TC$ which are Q-generating on the left, i.e. such that we have
\[
\left\{
\begin{array}{rclcrclcrcl}
C & \subset & C_0^{\sim}\ , & \text{i.e.} & C & = & C_0^{\sim}\ , & \text{i.e.} & TF & = & C_{0*}\ ;\\
TC & \subset & (TC_0)^{\sim}\ , & \text{i.e.} & TC & = & (TC_0)^{\sim}\ , & \text{i.e.} & F & = & (TC_0)_*\ .
\end{array}\right. % reading resolved: DER_XIX_corrections.csv
\]
\end{itemize}
\end{itemize}

Thus, the datum of such a \textit{special} Quillen triple in $\M$, when $\M$ satisfies (2.6. a), amounts to the datum of the two subsets
\msp{page 15}
\[
TC \ \subset\ C
\tag*{(2.7)}
\]
of $\Fl(\M)$, satisfying the conditions:
\[
\left\{
\begin{array}{ll}
C,\ TC\ \text{``special''} &
\left\{\begin{array}{ll}
\text{a)} & \text{$C$, $TC$ are Q-closed on the left.}\\[1ex]
\text{b)} & \parbox[t]{0.45\textwidth}{$C$, $TC$ are Q-generated on the left by \textit{small} subsets $C_0$, $TC_0$ of $C$, $TC$ respectively.}
\end{array}\right.\\[5ex]
 & \quad\ \begin{array}{ll}
\text{c)} & TC \subset W,\ TF \overset{\mathrm{def}}{=} C_* \subset W.\\[1ex]
\text{d)} & \parbox[t]{0.45\textwidth}{We have one of the three equivalent additional conditions of (2.4), i.e.:}\\[5ex]
 & \begin{array}{lrcllrcl}
\alpha) & TC & = & C \cap W\ , & \text{i.e.} & C \cap W & \subset & TC\ .\\
\beta) & TF & = & F \cap W\ , & \text{i.e.} & F \cap W & \subset & TF\ .\\
\gamma) & W & = & TF \circ TC\ . & & & &
\end{array}
\end{array}
\end{array}\right.
\tag*{(2.8)}
\]
\DERmark{ENV1597}


\printedpage{10}
% head: A. Grothendieck : Dérivateurs, Ch. XIX, Édition des Universités de Montpellier II et Paris VII
% foot: Version 11 février 2019

\textbf{Proposition 2.1.} \textit{Let}
\[
\begin{tikzcd}
\M' \arrow[r, shift right, "v^*"'] & \M \arrow[l, shift right, "v_!"']
\end{tikzcd}
\tag*{(2.1.1)} % [sic] (2.11) -> (2.1.1), numbering sequence (2.1.2), (2.1.3), ...
\]
\DERmark{ENV1598}

\textit{be a pair of adjoint functors, $\M$ being equipped with a closed Quillen structure} $(W,C,F)$. \textit{Put}
\[
\left\{
\begin{array}{rcl}
W' & = & (v^*)^{-1}(W)\\
F' & = & (v^*)^{-1}(F)\\
TF' & = & F' \cap W' \ = \ (v^*)^{-1}(TF)\\
C' & = & v_!(C)^{\sim}\\
TC' & = & v_!(TC)^{\sim}\ .
\end{array}\right.
\tag*{(2.1.2)}
\]

\begin{itemize}
\item[1°)] \textit{We then have}
\[
\left\{
\begin{array}{rclcrcl}
C' & \longleftrightarrow & TF'\ , & & TC' & \longleftrightarrow & F'\ ,\\
TC' & \subset & C'\ , & & TF' & \subset & F'\ .
\end{array}\right.
\tag*{(2.1.3)}
\]

\textit{For the data} (2.1.2) \textit{in $\M'$ to correspond there to a Quillen triple $(W',C',F')$, it is necessary and sufficient that the following conditions be satisfied:}
\msp{page 16}
\[
\left\{
\begin{array}{ll}
\text{a)} & \parbox[t]{0.7\textwidth}{\textit{Factorization for $C' \longleftrightarrow TF'$ and $TC' \longleftrightarrow F'$, i.e. the relations $C' \Qcouple TF'$ and $TC' \Qcouple F'$,}}\\[4ex]
\text{b)} & TC' \subset W',
\end{array}\right. % reading resolved: DER_XIX_corrections.csv
\tag*{(2.1.4)}
\]
\textit{or again, in place of} b)\textit{, the equivalent condition} (in the presence of a)):
\[
\left\{
\begin{array}{ll}
\text{b$'$)} & \parbox[t]{0.7\textwidth}{\textit{Every arrow $f'$ in $\M'$ (not only $f' \in TC'$) factors as $f' = p'i'$, with $i' \in W'$, $p' \in F'$.}}
\end{array}\right.
\tag*{(2.1.5)}
\]
\item[2°)] \textit{Suppose} (2.1.4. a) \textit{is satisfied (factorization). Then, for $(W',C',F')$ to be a closed Quillen triple, it} suffices \textit{that we have}
\[
\text{b$''$)}\qquad \boxed{\ v^*v_!(TC)\ \subset\ TC\ }\ ,
\tag*{(2.1.6)}
\]
\textit{and that in addition we be in one of the following three cases:}
\begin{itemize}
\item[A)] \textit{$v^*$ admits a right adjoint $v_*$.}
\item[B)] \textit{$v^*$ commutes with small} $\varinjlim$\textit{, $(W,C,F)$ is special, and $\M'$ is special, i.e. satisfies condition} (2.6. a).
\item[C)] \textit{$v^*$ is fully faithful, and $v^*v_!(F) \subset F$.}
\end{itemize}
\end{itemize}
Let us note at once the

\textbf{Corollary 2.2.} \textit{Suppose that $(W,C,F)$ is a special Quillen triple, that $\M'$ is special, that $v^*$ commutes with small} $\varinjlim$\textit{, finally that we have} (2.1.6)\textit{, i.e. $v^*v_!(TC) \subset TC$. Then $(W',C',F')$ is a special Quillen triple.}

\printedpage{11}
% head: A. Grothendieck : Dérivateurs, Ch. XIX, Édition des Universités de Montpellier II et Paris VII
% foot: Version 11 février 2019

\msp{page 17}
Indeed, by 2.1. 2°, case B, to see that $(W',C',F')$ is a closed Quillen triple, one must check (2.1.4. a). But this follows from 1.8, applied to the two special Quillen couples $(C,TF)$, $(TC,F)$ in $\M$: the couples $(C',TF')$ and $(TC',F')$ are special Quillen couples, which proves at the same time that the triple $(W',C',F')$ is special, q.e.d. % [sic] (C',TC') -> (C',TF'), image of the couple (C,TF)
\DERmark{ENV1599}\DERmark{ENV1600}


Let us prove 2.1. The first two relations (2.1.3) are contained in 1.6, the relations $TC' \subset C'$, $TF' \subset F'$ are tautological. The criterion (2.1.4) for $(W',C',F')$ to be a closed Quillen triple, is contained in the reminder on page 13, taking into account that we already know that $TF' = F' \cap W'$ (relation (2.4. b)). That the condition (2.1.4. b), in the presence of (2.1.4. a), is equivalent to the condition (2.1.5. b$'$), is contained in the first assertion (2.5), noting that we have $W' = W'^{\natural}$ (which follows from the definition of $W'$ and from $W = W^{\natural}$).

It remains to prove 2°), i.e. that condition b) is satisfied assuming a) and b$''$) (2.1.6) in each of the three cases A), B), C).

\textit{Case A.} By definition of $W' = (v^*)^{-1}(W)$, to prove that $TC' \subset W'$, it suffices to prove
\[
v^*(TC')\ \subset\ W\ ,  % [sic] W' -> W, v^*(TC') lies in M and W' = (v^*)^{-1}(W)
\]
\DERmark{ENV1601}

and a fortiori, it suffices to prove
\[
v^*(TC')\ \subset\ TC\ ,\qquad\quad \text{i.e. } v^*(TC') \harp F\ .
\tag*{(2.1.7)} % reading resolved: DER_XIX_corrections.csv
\]
Now by 1.5, applied to the pair of
\msp{page 18}
adjoint functors $(v^*,v_*)$, this relation is equivalent to
\[
TC' \harp v_*(F)\ .
\]
Since $TC' \overset{\mathrm{def}}{=} v_!(TC)^{\sim}$, this is equivalent to
\DERmark{ENV1602}

\[
v_!(TC) \harp v_*(F)\ ,
\]
or again (applying 1.5 again to $(v_!(TC),F)$ instead of $(TC',F)$), to
\[
v^*v_!(TC) \harp F\ ,\qquad\quad \text{i.e. } v^*v_!(TC) \subset \underbrace{TC}_{=\,F^*}\ ,
\]
now this is hypothesis b$''$).

\textit{Case B.} It again suffices to prove 2.1.7. Now, since $v_!(TC)$ admits a small Q-generating set on the left, $v_!(TC)^{\sim}$ is the closure of $v_!(TC)$ under the operations of cobase change, transfinite composition, direct factors. Since $v^*$ commutes with all these operations (as it commutes with small $\varinjlim$), and since $TC$ in $\Fl(\M)$ is stable under these operations, we thus see that to show that $v^*$ maps $v_!(TC)^{\sim}$ into $TC$, it suffices to see that it maps $v_!(TC)$ into $TC$, which is again hypothesis b$''$).

\textit{Case C.} We are going to check b) in the equivalent form b$'$). So let $f' \in \Fl(\M')$, $f' : X' \longrightarrow Y'$,
\printedpage{12}
% head: A. Grothendieck : Dérivateurs, Ch. XIX, Édition des Universités de Montpellier II et Paris VII
% foot: Version 11 février 2019
\msp{page 19}
put
\[
f\ =\ v^*(f'),
\]
and factor $f$ as $pi$, with $i \in TC$, $p \in F$, whence
\[
v_!v^*(f') = v_!(f) = v_!(p)v_!(i)\ .
\]
But $v^*$ being fully faithful, $v_!v^* \xrightarrow{\ \sim\ } \mathrm{id}$, so we find a factorization of $f'$ as $v_!(p)v_!(i)$. It therefore suffices to see that
\[
v_!(p)\ \in\ F'\ ,\qquad\quad v_!(i)\ \in\ W'\ ,
\]
which is equivalent to
\[
v^*v_!(p) \in F\ ,\qquad\quad v^*v_!(i) \in W\ ,
\]
and follows from the hypotheses that $v^*v_!(F) \subset F$ and $v^*v_!(TC) \subset TC \subset W$, q.e.d.

\textbf{N.B.} I have noted cases A) and C) for conscience's sake, with no precise application in view. The most useful case will fall under case B), and deserves to be spelled out again in a

\textbf{Theorem 2.3}\footnote[7]{cf. scholie 2.10, which should precede theorem 2.3.}. \textit{Let $(W,C,F)$ be a special Quillen triple in $\M$. For every $I$ in} Cat\textit{, let us consider the triple $(W(I),C\langle I\rangle,F(I))$ in $\M(I) = \underline{\mathrm{Hom}}(I^{\circ},\M)$, where the notations are those of page 10, so $W(I)$, $F(I)$ and $TF(I)$ are defined fiber by fiber, and $C\langle I\rangle = (TF(I))^* = (v_{I\,!}(C(I_0)))^{\sim}$}\footnote[8]{where $v_I : I_0 \longrightarrow I$ ($I_0$ discrete category having the same objects as I).}\textit{.}
\DERmark{ENV1603}

\begin{itemize}
\item[1°)] \textit{For every $I$ in} Cat\textit{, $(W(I),C\langle I\rangle,F(I))$ is}
\msp{page 20}
\textit{a special Quillen triple in $\M(I)$.}
\item[2°)] \textit{If $u : I \longrightarrow I'$ is an arrow in} Cat\textit{, whence}
\[
\begin{tikzcd}
\M(I) \arrow[r, shift left, "u_!"] & \M(I'), \arrow[l, shift left, "u^*"]
\end{tikzcd}
\tag*{(2.3.1)}
\]
\textit{then}
\[
\left\{
\begin{array}{l}
\left.\begin{array}{rcl}
u^*(W(I')) & \subset & W(I)\\
u^*(F(I')) & \subset & F(I)\\
u^*(TF(I')) & \subset & TF(I)
\end{array}\right\}\textit{tautology}\\
\begin{array}{rcl}
u_!(C\langle I\rangle) & \subset & C\langle I'\rangle\\
u_!(TC\langle I\rangle) & \subset & TC\langle I'\rangle\ ,
\end{array}
\end{array}\right.
\tag*{(2.3.2)}
\]
\textit{finally, $u_!$ induces a functor}
\[
\M(I)_{\cc} \longrightarrow \M(I')_{\cc} \qquad\quad \text{\footnotemark[9]}
\tag*{(2.3.3)}
\]
\footnotetext[9]{Subcategories of cofibrant objects.}
compatible with the localizers $W(I)_{\cc} = W(I)|\M(I)_{\cc}$ \textit{and} $W(I')_{\cc} = W(I')|\M(I')_{\cc}$\textit{, hence induces a functor}
\[
u_!^W\ :\ \Ho_W(I) \longrightarrow \Ho_W(I')\ .
\tag*{(2.3.4)}
\]
\printedpage{13}
% head: A. Grothendieck : Dérivateurs, Ch. XIX, Édition des Universités de Montpellier II et Paris VII
% foot: Version 11 février 2019
\textit{This functor is a left adjoint to the functor $u^*_W$,}
\[
u^*_W\ :\ \Ho_W(I') \longrightarrow \Ho_W(I)\ ,
\tag*{(2.3.5)}
\]
\textit{induced by $u^*$} (2.3.1). % [sic] (3.3.1) -> (2.3.1), the adjunction displayed in (2.3.1)
\DERmark{ENV1604}

\end{itemize}
\textsc{Proof of 2.3.}

\textit{Let us prove} 1°). We consider as in 1.10 the functor
\[
v = v_I\ :\ I_0 \longrightarrow I
\]
and the pair of adjoint functors
\[
\begin{tikzcd}
\M(I) \arrow[r, shift right, "v^*"'] & \M(I_0)\ , \arrow[l, shift right, "v_!"']
\end{tikzcd}
\tag*{(2.3.6)}
\]
to which we can apply 2.2. It remains
\msp{page 21}
only to check
\[
v^*v_!(TC)\ \subset\ TC\ ,
\tag*{(2.3.7)}
\]
now we have, for $X = (X_i)_{i \in I_0} \in \M(I_0)$,
\[
\begin{aligned}
v_!(X)(i) &\ \simeq\ \textstyle\coprod_{j \in I_0} v_{j\,!}(X_j)(i)\\
&\ \simeq\ \textstyle\coprod_{j \in I_0}\ \ \underbrace{X_j^{\mathrm{Hom}(i,j)}}_{\substack{\text{direct sum}\\ \text{of }\mathrm{Hom}(i,j)\text{ copies}}}\qquad ,
\end{aligned}
\]
and relation (2.3.7) follows from the fact that $TC$, being Q-closed on the left, is stable under small direct sums.

\textit{Let us prove} 2°). Formulas (2.3.2) are essentially tautological, the last two follow from the third and the second, by 1.5 applied to the pair of adjoint functors $(u_!,u^*)$. The relation $u_!(C\langle I\rangle) \subset C\langle I'\rangle$, and the fact that $u_!$ transforms initial object into initial object, imply (2.3.3). That $u_!$ restricted to $\M(I)_{\cc}$ is compatible with the localizers comes from $u_!(TC\langle I\rangle) \subset TC\langle I'\rangle$, and from the fact that every arrow $f$ in $\M(I)_{\cc}$ which is in $W(I)$, decomposes as $f = pi$, with $i \in TC$ and $p$ a left inverse of a $j \in TC\langle I\rangle$. This seen, together with the first two formulas (2.3.2) shows that we are under the conditions
\msp{page 22}
of Quillen's adjunction theorem, whence the adjunction of $u_!^W$ and $u^*_W$, q.e.d.

\textbf{Corollary 2.4.} \textit{Let $W$ be a localizer in $\M$, such that there exists a special Quillen triple $(W,C,F)$ including $W$. Then the prederivator} $\Ho_W$ \textit{on} Cat \textit{satisfies the axioms Der 1 (decomposition), Der 2 (localization), Der 3 (final object), Der 4 bis (existence of the $u_!^W$), Der 5 bis (computation of the fibers of the $u_!^W$), Der 6 bis (left exactness, i.e. left exactness of $W$-cocartesian squares in the} $\Ho_W(I)$\textit{), Der 6 (right exactness of $W$-cartesian squares).} % [sic] Der 4 bis -> Dér 4 bis, accent missing as printed

We have just proved Der 4 bis, and Der 5 bis is practically contained in the manner of obtaining of $u_!$, taking into account proposition 1.12, which implies that for $X$ in $\M(I)_{\cc}$, and $i'$ in $I'$, the inverse image of $X$ in $\M(i'\backslash I)$ is in $\M(i'\backslash I)_{\cc}$ \dots


\endgroup
% END UNIT der-ch19/der-ch19_p011-015.tex
% BEGIN UNIT der-ch19/der-ch19_p016-020.tex
\begingroup\makeatletter%
\derleqno%
\newcommand{\printedpage}[1]{\marginpar{\footnotesize p.~#1}}%
\newcommand{\msp}[1]{\marginpar{\footnotesize [#1]}}%
\newcommand{\deq}{\overset{\mathrm{def}}{=}}%
\newcommand{\Ho}{\mathrm{Ho}}%
\newcommand{\Mono}{\mathrm{Mono}}%
\newcommand{\Fl}{\mathrm{Fl}}%
\newcommand{\Filt}{\mathrm{Filt}}%
\newcommand{\Cof}{\mathrm{Cof}}%
\newcommand{\qpair}{\mathrel{\leftharpoondown\joinrel\relbar\joinrel\rightharpoonup}}%
\makeatletter%
\renewcommand\@makefnmark{\hbox{\@textsuperscript{\normalfont(\@thefnmark)}}}%
\renewcommand\@makefntext[1]{\parindent 1em\noindent\hbox{\@textsuperscript{\normalfont\@thefnmark}}#1}%
\newcommand{\Qcouple}{\mathrel{\leftarrow\mkern-9mu{+}\mkern-9mu\rightarrow}}%
\makeatother%
\makeatother\let\printedpage\derpp\let\msp\dermsp%
\let\mspin\DERmspin%

\DERmark{ENV1605}


\printedpage{14}
% head: A. Grothendieck : Dérivateurs, Ch. XIX, Édition des Universités de Montpellier II et Paris VII
% foot: Version 11 février 2019

\msp{page 23}
The decomposition axiom is satisfied for every prederivator of the form $\mathbf{D}_{\mathcal{M},W}$. The left exactness axiom Der 6 bis is valid, provided that the $(\mathcal{M}(I),W(I))$ fit into a structure with cofibrations, which is the case. Likewise for the localization axiom Der 2, provided that in addition $W$ is strongly saturated, which is also the case. Finally, the final object axiom says that if $I$ has a final object, then
\DERmark{ENV1606}

\[ \Ho_W(e) \longrightarrow \Ho_W(I) \]
is fully faithful, or again that the adjunction morphism
\[ u_!u^*(\xi) \longrightarrow \xi \]
is an isomorphism. But by virtue of 1.12, case 3°, in this case $u^*$ maps $\mathcal{M}(e)_{\mathrm{c}} = \mathcal{M}_{\mathrm{c}}$ into $\mathcal{M}(I)_{\mathrm{c}}$, so the computation of the pair of adjoint functors $u_!^W$, $u_W^*$ can be carried out using the restrictions of $u_!$, $u^*$ to $\mathcal{M}(I)_{\mathrm{c}}$ and $\mathcal{M}_{\mathrm{c}}$ respectively. Since the adjunction morphism $u_!^{\mathrm{c}}u_{\mathrm{c}}^* \longrightarrow \mathrm{id}$ is an isomorphism, the same therefore holds for $u_!^W u_W^* \longrightarrow \mathrm{id}$, which is deduced from it by localizing, q.e.d.

\textbf{2.5.} To finish constructing $\Ho_W$ as a derivator on all of Cat, one must \msp{page 24} therefore establish Der 4, Der 5: existence of the $u_*^W$, and computation of the fibers of $u_*^W$. As for Der 6, it is obtained automatically, from the fact that each $(\mathcal{M}(I),W(I))$ fits into a fibration structure $(W(I),F(I))$. (Omission in the proof of 2.4.) The natural idea is now to use the triple structure
\[ (W(I),C(I),\ \underbrace{F[I]}_{\deq TC(I)_*}\ )\ , \tag*{(2.5.1)} \]
but the difficulty is to prove the factorization for a couple
\[ C(I) \longleftrightarrow TF[I]\ , \qquad \underbrace{TC(I)}_{\deq C(I)\cap W(I)} \quad \longleftrightarrow F[I]\ , \tag*{(2.5.2)} \]
for want of being able to show, under the general conditions of 2.3 alone (special Quillen triple) that $C(I)$ and $TC(I)$ are also special. (Cf. question 1.11.) But with suitable conditions on $\mathcal{M}$, we shall manage to get by if $C = \mathit{mono}$. The relevant statement, a little technical, is found in XV cor. 7.4 (p. 98-99), which I am going to recall here.

\printedpage{15}
% head: A. Grothendieck : Dérivateurs, Ch. XIX, Édition des Universités de Montpellier II et Paris VII
% foot: Version 11 février 2019

\msp{page 25}
\textbf{Theorem 2.6.} \textit{Let $\mathcal{M}$ be a category satisfying the following conditions.}
\[ \left\{ \begin{minipage}{0.72\textwidth}\itshape
\begin{itemize}[leftmargin=2.5em]
\item[\upshape a)] $\mathcal{M}$ is ``special'', i.e. stable under small $\varinjlim$ , and its elements are accessible.
\item[\upshape b)] Filtering direct limits in $\mathcal{M}$ are left exact (and in $\mathcal{M}$ finite $\varprojlim$ are representable – follows from \textup{a} and \textup{c}).
\item[\upshape c)] $\mathcal{M}$ has a cardinal filtration $(\Filt^\pi)_{\pi\geq\pi_0}$.
\item[\upshape d)] The $\Filt^\pi$ are stable under $\varinjlim$ indexed by categories $I$ such that $\operatorname{card}\Fl(I) \leq \pi$, and under quotients ($Z \in \Filt^\pi$ and $Z \longrightarrow Z'$ an epimorphism implies $Z' \in \Filt^\pi$).
\item[\upshape e)] If $A_1$, $A_2$ are two subobjects of an object $B$ of $\mathcal{M}$, and $A_0 = A_1 \cap A_2$, then $A_1 \amalg_{A_0} A_2 \longrightarrow B$ is a monomorphism.
\item[\upshape f)] If we have a cocartesian square
\begin{center}
\begin{tikzcd}
A_0 \arrow[r, "i"] \arrow[d] & A_1 \arrow[d] \\
A_2 \arrow[r, "i'"] & A_3
\end{tikzcd}
\end{center}
with $i$ a monomorphism, then $i'$ is a monomorphism.
\end{itemize}
\end{minipage}\right. \tag*{(2.6.1)} \]
\DERmark{ENV1607}\DERmark{ENV1608}

$^{(10,11,12,13)}$\footnotetext[10]{\textbf{N.B.} a + c also means that $\mathcal{M}$ is a ``$\mathfrak{U}$-paratopos category'', and implies the existence of a cardinal filtration satisfying d) (except the assertion concerning the quotients $Z'$ of $Z$). Thus we are dealing with a $\mathfrak{U}$-paratopos satisfying in addition b) e) f). (Is it then a topos\,? \textit{No,} for e), f) are O.K. in \textit{abelian categories\,!})}\footnotetext[11]{\textbf{N.B.} Condition d) is probably unnecessary, at least if epimorphisms are effective, for d) will be satisfied in any case for $\pi = 2^c$, with $c \geq \pi_0$.}\footnotetext[12]{\textbf{N.B.} Condition e) is \textit{not} satisfied for the category (Gr), and the theorem fails for this category, which does \textit{not} have ``enough injectives'', cf. SGA 4 I 8.12.9. On the other hand, e) and f) are satisfied by \textit{abelian} categories.}\footnotetext[13]{Cf. XVIII 4.6.7 (p. 157) for a more congenial formulation of the set of conditions a) to f) for $\mathcal{M}$.}. \textit{(We shall then say that $\mathcal{M}$ is} ``very special''\textit{). Let moreover}
\[ W \ \subset \Fl(\mathcal{M}) \tag*{(2.6.2)} \]
\textit{be a localizer in $\mathcal{M}$, which satisfies the conditions}
\[ \left\{ \begin{minipage}{0.72\textwidth}\itshape
\begin{itemize}[leftmargin=2.5em]
\item[\upshape a)] $W$ is stable under filtering $\varinjlim$ $^{\textup{(14)}}$.
\item[\upshape b)] $W$ is moreover $\mathrm{L}_1$-accessible, i.e. there exists a \textup{small} subset $W_0 \subset W$ such that $W$ is the $\mathrm{L}_1$-envelope of $W_0$, i.e. such that $W$ is the smallest subset of $\Fl(\mathcal{M})$ containing $W_0$, and stable under filtering $\varinjlim$. % [sic] la L_1-enveloppe de W -> la L_1-enveloppe de W_0, W_0 being the small generating part just introduced
\end{itemize}
\end{minipage}\right. \tag*{(2.6.3)} \]
\DERmark{ENV1609}\DERmark{ENV1610}\DERmark{ENV1611}

\footnotetext[14]{\textbf{N.B.} a) already implies that $W = W^\natural$.}

\printedpage{16}
% head: A. Grothendieck : Dérivateurs, Ch. XIX, Édition des Universités de Montpellier II et Paris VII
% foot: Version 11 février 2019

\msp{page 26}
\textit{Let us then put}
\[ C \ = \ \Mono(\mathcal{M}), \qquad TC \ = \ C\cap W\ , \qquad TF \ = \ C_*\ , \qquad F = (TC)_*\ . \tag*{(2.6.4)} \]
\textit{Under these conditions:}
\begin{itemize}[leftmargin=2.5em]
\item[1°)] \textit{$C \Qcouple TF$, more precisely $C \qpair TF$ is a special Quillen couple.} % reading resolved: DER_XIX_corrections.csv
\item[2°)] \textit{For $TC \qpair F$ to be a Quillen couple, it is necessary and sufficient that $TC$ be stable under cobase change, and then it is a special Quillen couple.} % reading resolved: DER_XIX_corrections.csv
\item[3°)] \textit{For $(W,C,F)$ to be a Quillen triple (which will necessarily be special by \textup{1°, 2°}, and in particular closed), it is necessary and sufficient that the following conditions be satisfied.}
\[ \left\{ \begin{array}{ll}
\text{a)} & TC\ (\deq W\cap\Mono)\ \text{\itshape is stable under cobase change.}\\[1ex]
\text{b)} & TF\ (\deq (\Mono)_*) \ = \ \text{\itshape ``maps of injective type''}\ \subset W.
\end{array}\right. \tag*{(2.6.5)~$^{(15)}$} \]
\footnotetext[15]{\textbf{N.B.} These conditions (2.6.5) are equivalent to saying that $(W,C)$ is a structure of a category with cofibrations on $\mathcal{M}$.}
\textit{And condition} b) \textit{is also equivalent to:}
\[ \begin{minipage}{0.7\textwidth}
\begin{itemize}[leftmargin=2.5em]
\item[b$'$)] \textit{Every arrow $f$ of $\mathcal{M}$ (or only of $TF$) factors as $f = pi$, with $i$ a monomorphism, $p \in W$.}
\end{itemize}
\end{minipage} \tag*{(2.6.6)} \]
\end{itemize}

\textbf{Corollary 2.7.} \textit{$\mathcal{M}$ and $W$ being as in} (2.6.1), (2.6.3), \textit{for $(W,C,F)$ to be a Quillen triple that is} left proper \textit{(i.e. $C \subset \Cof_W$), it is necessary and sufficient that $W$ satisfy condition} (2.6.5. b) \textit{(i.e. $W$ contains the arrows of injective type), and moreover}
\msp{page 27}

\begin{itemize}[leftmargin=7em]
\item[] $C \subset \Cof_W$, \textit{i.e. the condition:}
\item[a$'$)] \textit{For every cocartesian diagram}
\[ \begin{tikzcd}
A_0 \arrow[r, "i"] \arrow[d, "f"'] & A_1 \arrow[d, "f'"'] \\
A_2 \arrow[r, "i'"] & A_3
\end{tikzcd} \tag*{(2.7.7)} \]
\textit{with $i$ a monomorphism, if $f \in W$, then $f' \in W$. (Stability of $W$ under cobase change along monomorphisms.)}
\end{itemize}

\textit{Under these conditions, this Quillen triple is special, and a fortiori closed.}

This corollary follows from 2.6. 3°), since $C \subset \Cof_W$ implies $TC \deq C\cap W = C\cap W^{\mathrm{univ}}$ (for $\Cof_W \cap W \ = \ W^{\mathrm{univ}}$), hence that $TC$ is stable under cobase change, whence condition (2.6.5. a).

\printedpage{17}
% head: A. Grothendieck : Dérivateurs, Ch. XIX, Édition des Universités de Montpellier II et Paris VII
% foot: Version 11 février 2019

\textbf{2.8.} We say that a \textit{Quillen structure} $(W,C,F)$ on $\mathcal{M}$ is \textit{very special,} if it is obtained as in theorem 2.6, i.e. if we have the conditions:
\DERmark{ENV1612}

\[ \left\{ \begin{minipage}{0.72\textwidth}
\begin{itemize}[leftmargin=2.5em]
\item[a)] $\mathcal{M}$ is very special, i.e. satisfies conditions (2.6.1).
\item[b)] $W$ is stable under small filtering $\varinjlim$, and is accessible (conditions (2.6.3)).
\item[c)] $C = \Mono$.
\end{itemize}
\end{minipage}\right. \tag*{(2.8.1)} \]
\DERmark{ENV1613}

Thus for a localizer $W$ in $\mathcal{M}$ to correspond to a very special Quillen structure \msp{page 28} in $\mathcal{M}$ (which will necessarily be unique), it is necessary and sufficient that $\mathcal{M}$ and $W$ satisfy conditions (2.8.1. a)\,b)), and that moreover we have:
\[ \left\{ \begin{minipage}{0.65\textwidth}
\begin{itemize}[leftmargin=2.5em]
\item[a)] $TC \deq \Mono\cap W$ is stable under cobase change.
\item[b)] $TF \deq (\Mono)_* \subset W$, i.e. $W$ contains the arrows of injective type in $\mathcal{M}$.
\end{itemize}
\end{minipage}\right. \tag*{(2.8.2)} \]
\footnote[16]{We also say that the localizer $W$ is very special, if it defines a very special Quillen structure.}.

We now find the

\textbf{Theorem 2.9.} \textit{Let} $(W, C = \Mono, F = (TC)_* = (W\cap\Mono)_*)$ \textit{be a very special Quillen triple in a category $\mathcal{M}$} (2.8). \textit{Then:}
\begin{itemize}[leftmargin=2.5em]
\item[1°)] \textit{For every $I$ in} Cat, $(W(I),C(I),F[I])$ \textit{is a very special Quillen triple in $\mathcal{M}(I)$, and in particular we have}
\[ C(I) \ = \ \Mono(\mathcal{M}(I))\ , \tag*{(2.9.1)} \]
\textit{and moreover}
\[ F[I]\cap W(I) \ = \ (TF)[I]\ . \tag*{(2.9.2)} \]
\item[2°)] \textit{For every arrow $u : I \longrightarrow I'$ in} Cat, \textit{whence the pair of adjoint functors}
\[ \begin{tikzcd}
\mathcal{M}(I) \arrow[r, shift right, "u_*"'] & \mathcal{M}(I') \arrow[l, shift right, "u^*"']
\end{tikzcd}\ , \tag*{(2.9.3)} \]
\textit{we have}
\[ \left\{ \begin{array}{rcl}
u^*(W(I')) & \subset & W(I)\\
u^*(C(I')) & \subset & C(I)\\
u^*(TC(I')) & \subset & TC(I)\\
u_*(F[I]) & \subset & F[I']\\
u_*(TF[I]) & \subset & TF[I']\ ,
\end{array}\right. \tag*{(2.9.4)} \]
\textit{finally $u_*$ induces a functor}
\printedpage{18}
% head: A. Grothendieck : Dérivateurs, Ch. XIX, Édition des Universités de Montpellier II et Paris VII
% foot: Version 11 février 2019
\msp{page 29}
\[ u_*^{\mathrm{f}} \ :\ \mathcal{M}(I)_{\mathrm{f}} \longrightarrow \mathcal{M}(I')_{\mathrm{f}} \qquad\quad ^{(17)} \tag*{(2.9.5)} \]
\footnotetext[17]{full subcategories of fibrant objects.}
compatible with the localizers $W(I)_{\mathrm{f}}$ \textit{and} $W(I')_{\mathrm{f}}$ \textit{induced by $W(I)$, $W(I')$ respectively, hence induces a functor}
\[ u_*^W \ :\ \Ho_W(I) \longrightarrow \mathrm{Ho}_W(I')\ . \tag*{(2.9.6)} \] % [sic] italic Ho_W(I') -> upright Ho_W(I'), inconsistent with the upright Ho_W elsewhere
\DERmark{ENV1614}

\textit{This functor is a} right adjoint \textit{to the functor}
\[ u_W^* \ :\ \mathrm{Ho}_W(I') \longrightarrow \Ho_W(I) \tag*{(2.9.7)} \] % [sic] italic Ho_W(I') -> upright Ho_W(I'), inconsistent with the upright Ho_W elsewhere
\DERmark{ENV1615}

\textit{induced by $u^*$} (2.9.3).
\end{itemize}
\textsc{Proof of 2.9.} To prove 1°), we must first check that the preliminary conditions (2.8.1. a)\,b)) on $\mathcal{M}$, $W$ are stable under passage from $\mathcal{M}$, $W$ to $\mathcal{M}(I)$, $W(I)$, which is done by AQT \texttt{[âne qui trotte: at a donkey's trot, i.e. by routine verification]}. Besides, one checks right away that
\[ C(I) \ = \ \Mono(\mathcal{M}(I))\ , \]
and theorem 2.6 (1°, 2°) then implies that
\[ \begin{array}{ccccc}
C(I) & \qpair & C(I)_* & \deq & TF[I]\\
\underbrace{TC(I)}_{\deq C(I)\cap W(I)} & \qpair & (TC(I))_* & \deq & F[I]
\end{array} \] % reading resolved: DER_XIX_corrections.csv
are special Quillen couples. (Noting for the second, that $TC(I)$ is indeed \msp{page 30} stable under cobase change, $TC$ being so.) To check that we have a very special Quillen triple, it remains by virtue of 2.8 to check that
\[ TF[I] \ \subset\ W(I)\ , \]
which is written
\[ v^*(TF[I]) \ \subset\ W(I_0)\ , \]
where
\[ v \ :\ I_0 \longrightarrow I \]
is defined as in (1.12) (p. 10). In fact, we shall even have
\[ v^*(TF[I]) \  \subset\ \underbrace{TF(I_0)}_{= F(I_0)\cap W(I_0)}\quad , \]
\DERmark{ENV1616}

and the proof is dual to that of
\[ v^*v_!(TC(I_0)) \ \subset\ TC(I_0) \]
(cf. (2.3.7), p. 21), it does not depend on questions of factorization in $\mathcal{M}(I)$, and comes here from the stability of $TF$ under small products. This proves 1°.

The proof of 2° is now the dual of that of (2.3. 2°). Since this makes us repeat the same proof twice, this shows that we have not ``set up the sorites'' quite carefully enough. So by way of making up for this, I add a


\endgroup
% END UNIT der-ch19/der-ch19_p016-020.tex
% BEGIN UNIT der-ch19/der-ch19_p021-025.tex
\begingroup\makeatletter%
\newcommand{\printedpage}[1]{\marginpar{\footnotesize p.~#1}}%
\newcommand{\msp}[1]{\marginpar{\footnotesize [#1]}}%
\newcommand{\M}{\mathcal{M}}%
\newcommand{\Cat}{\mathrm{Cat}}%
\newcommand{\Ho}{\mathrm{Ho}}%
\newcommand{\Fl}{\mathrm{Fl}}%
\newcommand{\Cof}{\mathrm{Cof}}%
\newcommand{\Diag}{\mathrm{Diag}}%
\newcommand{\id}{\mathrm{id}}%
\newcommand{\deq}{\overset{\mathrm{def}}{=}}%
\newcommand{\Qs}{\tilde{\ }}%
\newcommand{\fact}{\mathrel{\leftarrow\mkern-11mu\twoheadrightarrow}}%
\newcommand{\harp}{\mathrel{\leftharpoondown\joinrel\relbar\joinrel\rightharpoonup}}%
\makeatother\let\printedpage\derpp\let\msp\dermsp%
\let\mspin\DERmspin%

\DERmark{ENV1617}


\printedpage{19}
% head: A. Grothendieck : Dérivateurs, Ch. XIX, Édition des Universités de Montpellier II et Paris VII
% foot: Version 11 février 2019

\msp{page 31}
\textbf{Scholie 2.10.} \textit{Let $(W,C,F)$ be a closed Quillen triple in $\M$. For $I$ in \textup{Cat}, for $(W(I),C\langle I\rangle,F(I))$ to be a closed Quillen triple for $\M(I)$, it is necessary and sufficient that the couples $C\langle I\rangle\longleftrightarrow TF(I)$ and $TC\langle I\rangle\longleftrightarrow F(I)$ be factorizing. Suppose that this condition is satisfied for $I$ and for $I'$, and let $u : I\longrightarrow I'$ be an arrow from $I$ to $I'$. Then the assertions of \textup{2.3. 2°} hold, in particular we find a functor $u^{W}_{!} : \Ho_W(I)\longrightarrow\Ho_W(I')$, left adjoint to $u^{*}_{W}$\footnote[18]{Say also that the computation of the fibers of $u^{W}_{!}(X)$, provided for in Der 5 bis, is OK.}. These results dualize at once (by $\M\longmapsto\M^{\circ}$, $I\longmapsto I^{\circ}$) for the triples $(W(I),C(I),F[I])$, and for the construction of $u^{W}_{*}$, right adjoint to $u^{*}_{W}$. In particular, for the preceding triple to be a Quillen triple (necessarily closed), it is necessary and sufficient that the two couples $C(I)\longleftrightarrow TF[I]$ and $TC(I)\longleftrightarrow F[I]$ be factorizing; and when this condition is satisfied for $I$ and for $I'$ (for a given arrow $u : I\longrightarrow I'$ in \textup{Cat}), then $u^{W}_{*}$ exists and is constructed dually to the construction of $u^{W}_{!}$ made explicit in \textup{2.3. 2°}.}

\msp{page 32}
% a vertical bar is printed in the left margin alongside Corollaire 2.11
\textbf{Corollary 2.11.} \textit{The prederivator $\mathbf{D}_{\M,W}$ defined for a \textup{very special} localizer $W$ in $\M$ is a derivator in the full sense of the term, i.e. it satisfies all the axioms Der 1 to Der 6 (including Der 4 bis and Der 5 bis and Der 6 bis) of derivators.}

\textbf{N.B.} To my own surprise, I have not been able to find a case of a localizer $W$ which defines a derivator on all of Cat (with existence of $u^{W}_{!}$, $u^{W}_{*}$), \textit{except} for $W$ a very special localizer, hence fitting into a very special Quillen triple -- or when one can reduce to this in the standard way. (For example in the case where $\M=\Cat$, $W$ a basic localizer satisfying Loc 6 a) and Loc 8, 8 bis)\footnote[19]{axiom of cocartesian squares (direct version), axiom of direct limits, accessibility.} -- the monomorphisms in $W$ flatly refuse to be stable under cobase change!) I have absolutely no idea whether the spectacular result ``proved'' by \textsc{Anderson} is true or false. It looks almost too good to be true \ldots{}Besides, I know of no special Quillen structure which is not very special -- sorry! % [sic] stable -> stables, agreement with « les monomorphismes »

\msp{page 33}

\section*{3\quad Case of categories with cofibrations, of categories with fibrations}

I recall the axioms of a \textit{cofibration structure} $(W,C)$ on a category $\M$.

\begin{itemize}[leftmargin=4em,labelwidth=3.5em,align=left]
\item[Cof$_0$ :] $W$ is a localizer.
\item[Cof$_1$ :] $C$ contains the isomorphisms, is stable under cobase change, under composition.
\item[Cof$_2$ :] $C\cap W$ is stable under cobase change (and also under composition, and contains the isomorphisms, by virtue of Cof$_0$, Cof$_1$).
\item[Cof$_3$ :] Every arrow $f$ in $\M$ factors as $f=pi$, with $i\in C$, $p\in W$.
\end{itemize}

(Axioms of K. S. \textsc{Brown}, slightly generalized by \textsc{Anderson}). There is an optional ``axiom
\printedpage{20}
% head: A. Grothendieck : Dérivateurs, Ch. XIX, Édition des Universités de Montpellier II et Paris VII
% foot: Version 11 février 2019
of continuity'' or \textit{``transfinite axiom''}, which is the following:

\begin{itemize}[leftmargin=4em,labelwidth=3.5em,align=left]
\item[Cof$_4$ :] $\M$ stable under small filtering $\varinjlim$, and $C$ and $TC = C\cap W$ are stable under transfinite composition.
\end{itemize}

Thus, Cof$_1$, Cof$_2$, Cof$_4$ can be summarized by saying that $C$ and $TC$ are transfinitely Q-saturated on the left.

The cofibration structure is called \textit{special} if it satisfies Cof$_4$, if $\M$ is special, and if $C$ admits a \textit{small} subset $C_0$ which is Q-generating on the left, i.e. such that
\[
C \ \subset\ C_0\Qs\,, \qquad \text{i.e.}\quad \tilde{C}=C_0\Qs. \tag*{(3.1)}
\]

\msp{page 34}
Then, by the factorization theorem 1.1, we have
\[
C \ \fact\ TF \ \deq\ C_* \ ,
\]
and the factorization axiom Cof$_3$ implies (cf. 1.4) $TF\subset W^{\natural}$. More generally, we have:

\textbf{Proposition 3.1.} \textit{Suppose that $\M$ is special (stable under small $\varinjlim$ , and accessible elements), that $(W,C)$ satisfies Cof$_0$, Cof$_1$, and that $C$ is stable under transfinite composition (which, together with Cof$_1$, is stated by saying that $C$ is transfinitely Q-saturated on the left), finally that $C$ admits a small subset $C_0$ Q-generating on the left, so that $C\fact TF\deq C_*$. Then for the factorization axiom Cof$_3$ to be satisfied, it suffices that we have}
\[
TF \ \subset\ W \ , \tag*{(3.1.1)}
\]
\textit{and it is necessary to have $TF\subset W^{\natural}$, hence if $W=W^{\natural}$, the axiom Cof$_3$ is equivalent to} (3.1.1).

The cofibration structure is called \textit{semi-closed} if $C$ is Q-closed on the left, i.e. if we have $C=C\Qs$ (which implies that $C$ is stable under transfinite composition), \textit{closed} if in addition $TC=(TC)\Qs$ (which implies the transfinite axiom Cof$_4$).

\msp{page 35}
I do not return here to the theory of K. S. \textsc{Brown} on categories with cofibrations. It is, to my knowledge, the most economical set of axioms known, for ensuring on the one hand a good grip on $\Ho_W$ (through a corresponding notion of homotopy of maps, and a calculus of fractions in $\Pi\M_{\mathrm{c}}$), and on the other hand the existence of the \textit{exact} sequence of cofibrations, hence a theory of $W$-homotopic cofibers. (\textbf{N.B.} When the structure is \textit{proper,} i.e. $C\subset\Cof_W$, then the existence of this structure amounts simply to saying that ``$W$ admits enough $W$-cofibrations'', i.e. that every $f\in\Fl(\M)$ factors as $pi$, with $i\in\Cof_W$, $p\in W$; and we can therefore take $C=\Cof_W$. In the case where $W$ is stable under small filtering $\varinjlim$, $C=\Cof_W$ and $TC=W^{\mathrm{univ}}$ satisfy in addition the transfinite axiom Cof$_4$.) Finally, K. S. \textsc{Brown} shows that every arrow $i : X\longrightarrow Y$ in $C$, \textit{with $X$, $Y$ cofibrant} (i.e. $\emptyset_{\M}\longrightarrow X$, $\emptyset_{\M}\longrightarrow Y$ in $C$) is in $\Cof_W$ (``Brown's Lemma'').

I am interested here in the question of extending the cofibration structure to a category of the form $\M(I)$, and of using this extension to construct
\msp{page 36}
the functors $u^{W}_{!}$. Let us note already that the theory of K. S. \textsc{Brown}, with the supplements of \textsc{Anderson}\footnote[20]{revised and arranged by me.}, gives essentially this:

\printedpage{21}
% head: A. Grothendieck : Dérivateurs, Ch. XIX, Édition des Universités de Montpellier II et Paris VII
% foot: Version 11 février 2019

\textbf{Proposition 3.2.} \textit{Let $W\subset\Fl(\M)$ be a localizer associated with a cofibration structure $(C,W)$ in $\M$.}
\begin{itemize}[leftmargin=3em,labelwidth=2.5em,align=right]
\item[1°)] \textit{On the category \textup{$\Diag_0$} of \textup{finite} ordered sets, this prederivator satisfies the axioms Der 1 (decomposition), Der 2 (localization), Der 3 (final object), Der 4 bis (existence of $u^{W}_{!}$), Der 5 bis (computation of the fibers of $u^{W}_{!}$), Der 6 bis (left exactness).}
\item[2°)] \textit{This is still true on the larger subcategory \textup{$\Diag_1$} in \textup{Cat}, consisting of \textup{noetherian} ordered sets $I$ such that the topological space associated with $I^{\circ}$ is noetherian (every increasing sequence of objects of $I$, or of open sets of $I^{\circ}$, is stationary).}
\item[3°)] \textit{Finally, if the transfinite axiom Cof$_4$ is satisfied, i.e. if $C$ and $TC$ are stable under transfinite composition, then the conclusion of \textup{1°} remains valid on the full subcategory \textup{$\Diag_2$} of \textup{Cat} consisting of ordered sets which are merely noetherian.}
\end{itemize}

\textbf{Corollary 3.3.} \textit{Suppose that we have Cof$_3$ in the stronger form of a factorization
\msp{page 37}
$f=pi$ \textup{functorial} in $f$, which is the case in particular if the cofibration structure is special. Then for every $I$ in \textup{Cat}, $\M(I)$ equipped with $(W(I),C(I))$ is a category with cofibrations, which satisfies the transfinite axiom Cof$_4$ if $(W,C)$ satisfies it. Thus we can apply to it the results of \textup{3.2}, which imply in particular that the axioms Der 1, 2, 3 ( ?) hold on all of \textup{Cat}, and that $u^{W}_{!}$ exists when $u$ is of the form $\id_I\times v$, with $v : J\longrightarrow K$ an arrow in \textup{$\Diag_1$} (resp. in \textup{$\Diag_2$}, if the transfinite axiom is satisfied).}

\textbf{3.4.} The key to the theory is in the extension of the cofibration structure $(W,C)$ from $\M$ to a category $\M(I)$, where $I$ is an ordered set, as $(W(I),C_I)$. The main facts are stated in XVII, in the theorems p. 26-27, p. 50, p. 60. It is when $I$ is \textit{noetherian,} and we assume either that $\mathrm{Esp}(I^{\circ})$ is also noetherian (e.g. $I$ finite), or that $(W,C)$ satisfies the transfinite axiom Cof$_4$, that we prove that $(W(I),C_I)$ is also a cofibration structure on $\M(I)$.

\msp{page 38}
An important fact for the construction of the functors $u^{W}_{!}$, for $u : I\longrightarrow I'$ an arrow between noetherian ordered categories, is the relation
\[
C_I\cap W(I) \ = \ (TC)_I \ , \tag*{(3.4.1)}
\]
from which it follows that we have, for $u$ as above,
\[
u_!(C_I\cap W(I)) \ \subset\ C_{I'}\cap W(I') \ , \tag*{(3.4.2)}
\]
an essential ingredient of the construction of the functor $u^{W}_{!}$. To extend this construction to the case of a $\pm$ arbitrary arrow in Cat, there arises therefore the problem of extending the definition of $C_I$, and the formulas (3.4.1) and (3.4.2), to the case of $\pm$ arbitrary categories in Cat. Now we have proved:

\textbf{Proposition 3.5.} \textit{If $\Phi=\Psi^*$, with $\Psi=\Phi_*$, is a left \textup{Q-closed} subset of \textup{$\Fl(\M)$}, and $I$ a noetherian ordered set, $\Phi_I$ the extension of $\Phi$ to a subset of \textup{$\Fl(\M(I))$} defined by \textup{\textsc{Anderson}}, then we have}
\DERmark{ENV1618}

\[
\Phi_I \ = \ \Psi(I)^* \ = \ v_!(\Phi(I_0))\Qs \ , \tag*{(3.5.1)}
\]
\textit{where $v : I_0\longrightarrow I$ is the canonical functor from the discrete category defined by \textup{$\mathrm{Ob}\,I$}.}

\printedpage{22}
% head: A. Grothendieck : Dérivateurs, Ch. XIX, Édition des Universités de Montpellier II et Paris VII
% foot: Version 11 février 2019

This is (with different notations) lemma 6.2.5 (p. 107) of XVII, whose proof I do not repeat here.

This gives the key to a definition of $\Phi_I$ as $\Phi\langle I\rangle$ (with the notations of paragraphs 1, 2),
\msp{page 39}
for arbitrary $I$ in Cat,
\[
\Phi_I \ \deq\ \Phi\langle I\rangle \ \deq\ \Psi(I)^* \ = \ v_!(\Phi(I_0))\Qs \ , \tag*{(3.5.2)}
\]
where $\Psi=\Phi_*$, and where we assume $\Phi$ Q-closed on the left, i.e. $\Phi=\Psi^*$. In the case where we assume only that $\Phi$ is transfinitely Q-saturated, without more, we would like to set
\[
\Phi_I \ \deq\ \text{transfinite Q-saturation of } v_!(\Phi(I_0)) \ , \tag*{(3.5.3)}
\]
and I presume that for $I$ ordered noetherian, this coincides with \textsc{Anderson}'s definition. But I will restrict myself in what follows to the case where $C=\Phi$ is Q-closed on the left.

Let us then return to the case of a cofibration structure $(W,C)$ in $\M$, assuming $C$ and $TC=C\cap W$ Q-closed on the left. Let us define for every $I$ in Cat $C_I=C\langle I\rangle$ as in (3.5.2). We then have, for $u : I\longrightarrow I'$ an arrow in Cat,
\[
u_!(C\langle I\rangle) \ \subset\ C\langle I'\rangle \ , \tag*{(3.5.4)}
\]
but we have little reason here to hope to have (3.4.2) as well, without having the formula
\[
C\langle I\rangle\cap W(I) \ \overset{?}{=}\ (TC)\langle I\rangle \qquad\qquad \text{\footnote[21]{\textbf{N.B.} We have the inclusion $\supset$, since $v^*v_!TC(I_0)\subset TC(I_0)$.}}. \tag*{(3.5.5)}
\]
We have proved that this formula is valid if $I$ is ordered noetherian (XVII th. 2.1, p. 26-27). This
\msp{page 40}
brings us to the

\textbf{Question 3.6.} \textit{Let $\M$ be a category stable under small limits, and let $(W,C)$ be a cofibration structure in $\M$ such that $C$ and $TC=C\cap W$ are Q-closed on the left. Let $I$ be an object of \textup{Cat}, do we then have}
\[
C\langle I\rangle\cap W(I) \ = \ (TC)\langle I\rangle \text{ ?} \tag*{(3.6.1)}
\]

We saw in the preceding paragraph that this is so if $(W,C)$ fits into a closed Quillen triple $(W,C,F=(TC)_*)$, and when we assume in addition that the couples
\[
C\langle I\rangle \ \longleftrightarrow\ TF(I) \ , \qquad TC\langle I\rangle \ \longleftrightarrow\ F(I)
\]
are factorizing, in which case $(W(I),C\langle I\rangle,F(I))$ is a closed Quillen structure in $\M(I)$. Besides, $(W,C,F)$ is a Quillen structure (necessarily closed) if and only if the couples $C\longleftrightarrow \underbrace{TF}_{\deq C_*}$, $TC\longleftrightarrow \underbrace{\ F\ }_{\deq (TC)_*}$ themselves are factorizing, and if % [sic] factorisant -> factorisants, agreement with « les couples »
\printedpage{23}
% head: A. Grothendieck : Dérivateurs, Ch. XIX, Édition des Universités de Montpellier II et Paris VII
% foot: Version 11 février 2019
in addition $TF\subset W$ (which is automatic if $W=W^{\natural}$, for we will have in any case, by Cof$_3$, $TF\subset W^{\natural}$).
\DERmark{ENV1619}


But can we get by with less, to be able to conclude (3.6.1)? Would it not suffice to know that $TC\langle I\rangle\longleftrightarrow F(I)$ is factorizing (which has little chance of being true for an $I$ $\pm$ general, if we do not already have $TC\fact F$, factorization in $\M$ itself). This is indeed the case:

\msp{page 41}
\textbf{Proposition 3.7.} \textit{Under the general conditions of \textup{(3.6)} on $\M$, $W$, $C$, $I$, suppose that the couple $TC\langle I\rangle\longleftrightarrow F(I)$ (where $F\deq (TC)_*$) is factorizing, i.e. $TC\langle I\rangle\fact F(I)$. Then the formula \textup{(3.6.1)} is true.}

Indeed, if $f\in C\langle I\rangle\cap W(I)$, we factor $f$ as $pi$, with $i\in TC\langle I\rangle$, $p\in F(I)$. Since $f,i\in W(I)$, we will have $p\in W(I)$, i.e. $p\in TF(I)$, and because of $\underbrace{C\langle I\rangle}_{\ni\, f}\harp\underbrace{TF(I)}_{\ni\, p}$, there exists $q : Y\longrightarrow X'$ such that $pq=\id_Y$, $qf=i$, % reading resolved: DER_XIX_corrections.csv; % [sic] p : Y -> q : Y, the arrow Y -> X' is q (see the diagram and pq = id_Y)
\DERmark{ENV1620}

\[
\begin{tikzcd}[column sep=small, row sep=large]
X \arrow[rr, "f\in C\langle I\rangle\cap W(I)"] \arrow[dr, "i\in TC\langle I\rangle"'] & & Y \arrow[dl, dotted, shift right=1ex, "q"'] \\
& X' \arrow[ur, "p\in F(I)"'] \ , &
\end{tikzcd}
\]
so $f$ is a direct factor of $i\in TC\langle I\rangle$. Since $TC\langle I\rangle$ is stable under direct factors, we win.

\textbf{Corollary 3.8.} \textit{Let $(W,C,F)$ be a closed Quillen structure in $\M$, and suppose that $TC=C\cap W$ is a Q-special subset of $\M$, i.e. the couple $TC\fact F$ a special couple (i.e. $\M$ is \textup{special}, and there exists a small subset $TC_0$ Q-generating on the left for $TC$). Then for every $I$ in \textup{Cat}, we have}
\[
C\langle I\rangle\cap W(I) \ = \ TC\langle I\rangle \ . \tag*{(3.8.1)}
\]

\msp{page 42}
Apparently, we have not had to assume that $C$ too is a Q-special subset of $\Fl(\M)$. But as we need the factorization for $C\longleftrightarrow TF$ (to have at least a cofibration structure on $\M$), we have hardly any way of establishing it without assuming $C$ special. Not to mention that to be able to make (3.8.1) explicit for the construction of $u^{W}_{!}$, we need to know that $(W(I),C\langle I\rangle)$ does indeed define a cofibration structure on $\M(I)$, which can hardly be done (it seems to me) without having the factorization for $C\langle I\rangle\longleftrightarrow TF(I)$ as well; something I do not manage to establish for a general $I$ in Cat, except by means of the hypothesis of specialness on $C$ as well.

\textbf{3.9.} To get to the bottom of it, we would need a few concrete and illuminating examples, of the following type:

a) \texttt{[There is no b).]} Closed Quillen structure $(W,C,F)$ on $\M$, such that for the associated prederivator, $u^{W}_{!}$ does not always exist, or $u^{W}_{*}$ does not always exist. I am thinking in particular of the case of an arrow $u : I\longrightarrow I'$ with $I'=e$ (presuming that if in this case, $u^{W}_{!}$ resp. $u^{W}_{*}$ exists, then it exists in all cases). Such an example would be even more % [sic] n'existe par toujours -> n'existe pas toujours, typing error
\DERmark{ENV1621}



\endgroup
% END UNIT der-ch19/der-ch19_p021-025.tex
% BEGIN UNIT der-ch19/der-ch19_p026-028.tex
\begingroup\makeatletter%
\newcommand{\printedpage}[1]{\marginpar{\footnotesize p.~#1}}%
\newcommand{\msp}[1]{\marginpar{\footnotesize [#1]}}%
\newcommand{\limind}{\varinjlim}%
\newcommand{\limproj}{\varprojlim}%
\newcommand{\Fl}{\mathrm{Fl}}%
\newcommand{\Is}{\mathrm{Is}}%
\newcommand{\Ho}{\mathrm{Ho}}%
\newcommand{\qpair}{\mathrel{\leftarrow\mkern-4mu\mapstochar\mkern-2mu\rightarrow}}%
\newcommand{\Qcouple}{\mathrel{\leftarrow\mkern-9mu{+}\mkern-9mu\rightarrow}}%
\makeatother\let\printedpage\derpp\let\msp\dermsp%
\let\mspin\DERmspin%


\printedpage{24}
% head: A. Grothendieck : Dérivateurs, Ch. XIX, Édition des Universités de Montpellier II et Paris VII
% foot: Version 11 février 2019
\msp{page 43}
instructive, of course, if $\mathcal{M}$ is stable under small $\limind$ , if its objects are accessible etc. and if $W$ is L-accessible (i.e. stable under small filtering $\limind$, and moreover accessible \dots). The first candidate for $I$ that I think of would be a $B_G$ ($G$ a group, e.g. $G=\pm1$), it is indeed the example at the antipodes of ordered sets, which, for their part, do not seem to cause too much trouble (assuming the noetherian hypothesis at least).

The following examples are not all that convincing, alas!

\textbf{Example 3.10.} Let $\mathcal{M}$ be any category, then the couples $(\Fl(\mathcal{M}),\Is(\mathcal{M}))$ and $(\Is(\mathcal{M}),\Fl(\mathcal{M}))$ are both Quillen couples,
\[
\Fl(\mathcal{M}) \Qcouple \Is(\mathcal{M}) \ , \qquad \Is(\mathcal{M}) \Qcouple \Fl(\mathcal{M}) \ . \tag*{(3.10.1)} % reading resolved: DER_XIX_corrections.csv
\]
It follows that if we set
\[
W \ = \ \Is(\mathcal{M}) \ , \qquad C \ = \ F \ = \ \Fl(\mathcal{M}) \ , \tag*{(3.10.2)}
\]
$(W,C,F)$ is a closed Quillen triple \textsuperscript{(}\footnote[22]{at least if we assume that $\mathcal{M}$ is stable under $\limind$ and $\limproj$ that are \textit{finite.}}\textsuperscript{)}. Moreover, for every $I$ in Cat, we have $W(I)=\Is(\mathcal{M}(I))$, hence
\[
\Ho_W(I) \ = \ \mathcal{M}(I) \ . \tag*{(3.10.3)}
\]
So the question of the existence of the functors $u_!^W$, $u_*^W$ reduces to that of the usual functors $u_!$, $u_*$, left and right adjoints to $u^* : \mathcal{M}(I') \longrightarrow \mathcal{M}(I)$.

\msp{page 44}
It is well known that the existence of these functors for a given $u$, calls for the existence of certain $\limind$ resp. $\limproj$ in $\mathcal{M}$; e.g. if $I'=e$, it is necessary and sufficient that the $\limind$ (resp. $\limproj$ ) of type $I^\circ$ be representable in $\mathcal{M}$. But with the hypotheses made on $\mathcal{M}$, there is no problem if $u : I \longrightarrow I'$ is a functor between categories that are \textit{finite.} We are going to give an example where this is no longer so.

\textbf{Example 3.11.} Let $k$ be a commutative ring, and $\mathcal{M}=K^{\mathrm{b}}(k\text{-Ab})$ the category of bounded cochain complexes of $k$-modules, $W$ the set of quasi-isomorphisms, hence
\[
\Ho_W \ = \ D^{\mathrm{b}}(k\text{-Ab}) \ ,
\]
the bounded derived category of the abelian category of $k$-modules. We know, using the existence of enough projectives in $k$-Ab, that this structure is associated with a closed Quillen structure $(W,C,F)$, described by \textsc{Quillen}
\msp{page 45}
in Homotopical Algebra -- except that he takes $D^-(k)$ instead of $D^{\mathrm{b}}(k)$. Perhaps, to verify the axioms, we must assume that $k$ has finite cohomological dimension.

Let us take $I=B_G$, where $G$ is a group (e.g. $G=\{\pm1\}$), and $u : I \longrightarrow e$. Then $\mathcal{M}(I)$ is the category of bounded complexes of $k[G]$-modules. We see that $u_!^W$ resp. $u_*^W$ exists if and only if for every such bounded complex, or simply for a plain $k[G]$-module $M$, the complex $L\Gamma^G(M)$ resp. $R\Gamma_G(M)$ belongs to $D^{\mathrm{b}}(k)$, or what amounts to the same,
\printedpage{25}
% head: A. Grothendieck : Dérivateurs, Ch. XIX, Édition des Universités de Montpellier II et Paris VII
% foot: Version 11 février 2019
if and only if we have
\[
\begin{aligned}
H_i(G,M) \ &= \ 0 \quad \text{for } i \gg 0 \ ,\\
\text{resp.}\quad H^i(G,M) \ &= \ 0 \quad \text{for } i \gg 0 \ ,
\end{aligned}
\]
for every $k[G]$-module $M$. This is practically never the case; when we take e.g. $k=\mathbf{Z}$ and $G$ a finite group, this is the case only if $G=e$. The most economical example seems to me to be $G=\pm1$, $k=\mathbf{F}_2$.

\msp{page 46}
So here $u_!^W$ and $u_*^W$ do not exist, for $I \longrightarrow e$ with $I$ a finite category. But also $\mathcal{M}$ is not stable under small $\limind$ .

\textbf{Example 3.12.} Let $\mathcal{M}$ be a category equipped with a strongly saturated localizer $W$, such that $\mathcal{M} \longrightarrow \underbrace{\mathcal{M}W^{-1}}_{=\mathcal{H}}$ admits an adjoint -- a left one, say. So we have a situation of two adjoint functors % [sic] muni -> munie, agreement with « une catégorie »
\[
\begin{tikzcd}[column sep=huge]
\mathcal{H} \arrow[r, shift left, "v_!\ \text{fully faithful}"] & \mathcal{M} \arrow[l, shift left, "v^*\ \text{localization}"]
\end{tikzcd}
\ , \tag*{(3.12.1)}
\]
and $W=(v^*)^{-1}(\mathrm{Iso}(\mathcal{H}))$. We find that for every $I$ in Cat, we have
\[
\mathrm{Ho}_W(I) \ \simeq \ \mathcal{H}(I) \ \overset{\mathrm{def}}{=} \ \mathrm{Hom}(I^\circ,\mathcal{H}). \tag*{(3.12.2)} % [sic] HO_W -> Ho_W, notation Ho_W used in (3.10.3) and Exemple 3.11
\]
\DERmark{ENV1622}\DERmark{ENV1623}

So we can think of making an example where in $\mathcal{M}$ the small $\limind$ and $\limproj$ exist, but not in $\mathcal{H}$, so that the $u_!^W$ or $u_*^W$ do not always exist. But this seems difficult to construct, seeing that in the situation envisaged, we easily see that the small $\limproj$ exist in $\mathcal{H}$. So if in $\mathcal{H}$ % [sic] envisagé -> envisagée, agreement with « la situation »
\msp{page 47}
there exists a small cogenerating family (and for this, it is enough that such a family exist in $\mathcal{M}$), there will also be small $\limind$ , by virtue of SGA 4 I 8.12.7. Likewise, they exist if $\mathcal{H}$ is a full subcategory of $\mathcal{M}$ \textit{stable under small} $\limind$ , and it seems practically impossible to make an example (3.12.1) without $\mathcal{H}$ having this stability property. Not to mention that there is even the following perplexity:

\marginpar{! \textbar}\textbf{Question 3.13.} \textit{Is there a $\mathfrak{U}$-category (\,$\mathfrak{U}$ a given universe) in which the small $\limind$ are representable, without the same being true of the small $\limproj$ ?} \textbf{N.B.} Such a category does not admit a small generating family, a fortiori it is not equivalent to a small category.

\msp{page 48}
\textbf{Question 3.14.} \textsc{Anderson} claims that for the prederivator defined by a category with cofibrations satisfying the transfinite axiom Cof$_4$, the functors $u_!^W$ exist. His proof is totally botched, nevertheless I cannot manage to make a counterexample. Could the result he announces be true? I do not even know of an example of a localizer $W$, in a category $\mathcal{M}$ \textit{with small $\limind$ and small $\limproj$ ,} such that $u_!^W$ or $u_*^W$ do not always exist -- but I am convinced that there are examples (even with $\mathcal{M}$ a small category?). I am less convinced of this if we suppose $W$ stable under filtering $\limind$, and if need be, moreover, that the elements of $\mathcal{M}$ are accessible (plus possibly other innocuous categorical properties of $\mathcal{M}$, such as the existence of a cardinal
\printedpage{26}
% head: A. Grothendieck : Dérivateurs, Ch. XIX, Édition des Universités de Montpellier II et Paris VII
% foot: Version 11 février 2019
\msp{page 49}
filtration etc.)

\textbf{Question 3.15.} This brings me back to the particularly elusive and irritating question, of the existence of small Q-generating sets, in a subset $\Phi$ of $\Fl(\mathcal{M})$, when we suppose (say) that $\mathcal{M}$ is stable under small $\limproj$ , that its elements are accessible and that it admits a cardinal filtration. Thus we have the question: does an L-accessible subset of $\Fl(\mathcal{M})$, which we assume if need be to be Q-saturated on the left (or what amounts to the same, containing the isomorphisms, and stable under cobase change, under composition and under direct factors) admit a \textit{small subset} $\Phi_0$ that is Q-generating on the left? % [sic] Q-saturé -> Q-saturée, agreement with « une partie »

\marginpar{! \textbar}On the other hand, I have not at present constructed a single example of a couple $\Phi \longleftrightarrow \Psi$, in a category $\mathcal{M}$, that is not factorizing! The case where $\mathcal{M}$ is special, was of course the most interesting. Clearly, in all
\msp{page 50}
this it is examples that are lacking. The key to an understanding does seem to lie in the questions of factorization between sets $\Phi \longleftrightarrow \Psi$. One ought to make a systematic study of this, but I shrink from the prospect of a digression that would perhaps take me yet more weeks! % [sic] Le clef -> La clef, « clef » is feminine


\endgroup
% END UNIT der-ch19/der-ch19_p026-028.tex
\clearpage\section*{Apparatus of corrections}\addcontentsline{toc}{section}{Apparatus of corrections}\markboth{Apparatus of corrections}{Apparatus of corrections}
\subsection*{Chapter I}
\DERapparatusitem{ENV0001}{preamble;\allowbreak{}\allowbreak{} local\allowbreak{} macro\allowbreak{} definition}{\textbackslash{}\allowbreak{}text\{\allowbreak{}def\}\allowbreak{}}{\textbackslash{}\allowbreak{}mathrm\{\allowbreak{}def\}\allowbreak{}}{TYPESETTING:\allowbreak{}7;\allowbreak{}\allowbreak{} NOTATION\_\allowbreak{}DER:\allowbreak{}17}
\DERapparatusitem{ENV0002}{11;\allowbreak{}\allowbreak{} ms13,\allowbreak{}\allowbreak{} prose\allowbreak{} parenthesis\allowbreak{} around\allowbreak{} the\allowbreak{} simplicial-set\allowbreak{} example}{}{)\allowbreak{}}{DER\_\allowbreak{}I-a:\allowbreak{}8}
\DERapparatusitem{ENV0003}{18;\allowbreak{}\allowbreak{} Edition18;\allowbreak{}\allowbreak{} ms26,\allowbreak{}\allowbreak{} diagonal\allowbreak{} and\allowbreak{} products}{admits\allowbreak{} a\allowbreak{} left\allowbreak{} adjoint}{admits\allowbreak{} a\allowbreak{} right\allowbreak{} adjoint}{DER\_\allowbreak{}I-a:\allowbreak{}14}
\DERapparatusitem{ENV0004}{18;\allowbreak{}\allowbreak{} running\allowbreak{} text\allowbreak{} or\allowbreak{} formula}{\textbackslash{}\allowbreak{}text\{\allowbreak{}def\}\allowbreak{}}{\textbackslash{}\allowbreak{}mathrm\{\allowbreak{}def\}\allowbreak{}}{TYPESETTING:\allowbreak{}7;\allowbreak{}\allowbreak{} NOTATION\_\allowbreak{}DER:\allowbreak{}17}
\DERapparatusitem{ENV0005}{21;\allowbreak{}\allowbreak{} Edition21;\allowbreak{}\allowbreak{} ms31,\allowbreak{}\allowbreak{} abstract\allowbreak{} mate\allowbreak{} construction}{the\allowbreak{} left\allowbreak{} adjoint\allowbreak{} functors}{the\allowbreak{} right\allowbreak{} adjoint\allowbreak{} functors}{DER\_\allowbreak{}I-a:\allowbreak{}15}
\DERapparatusitem{ENV0006}{23;\allowbreak{}\allowbreak{} running\allowbreak{} text\allowbreak{} or\allowbreak{} formula}{\textbackslash{}\allowbreak{}text\{\allowbreak{}def\}\allowbreak{}}{\textbackslash{}\allowbreak{}mathrm\{\allowbreak{}def\}\allowbreak{}}{TYPESETTING:\allowbreak{}7;\allowbreak{}\allowbreak{} NOTATION\_\allowbreak{}DER:\allowbreak{}17}
\DERapparatusitem{ENV0007}{25;\allowbreak{}\allowbreak{} running\allowbreak{} text\allowbreak{} or\allowbreak{} formula}{\textbackslash{}\allowbreak{}text\{\allowbreak{}def\}\allowbreak{}}{\textbackslash{}\allowbreak{}mathrm\{\allowbreak{}def\}\allowbreak{}}{TYPESETTING:\allowbreak{}7;\allowbreak{}\allowbreak{} NOTATION\_\allowbreak{}DER:\allowbreak{}17}
\DERapparatusitem{ENV0008}{25;\allowbreak{}\allowbreak{} running\allowbreak{} text\allowbreak{} or\allowbreak{} formula}{\textbackslash{}\allowbreak{}text\{\allowbreak{}def\}\allowbreak{}}{\textbackslash{}\allowbreak{}mathrm\{\allowbreak{}def\}\allowbreak{}}{TYPESETTING:\allowbreak{}7;\allowbreak{}\allowbreak{} NOTATION\_\allowbreak{}DER:\allowbreak{}17}
\DERapparatusitem{ENV0009}{26;\allowbreak{}\allowbreak{} running\allowbreak{} text\allowbreak{} or\allowbreak{} formula}{\textbackslash{}\allowbreak{}text\{\allowbreak{}def\}\allowbreak{}}{\textbackslash{}\allowbreak{}mathrm\{\allowbreak{}def\}\allowbreak{}}{TYPESETTING:\allowbreak{}7;\allowbreak{}\allowbreak{} NOTATION\_\allowbreak{}DER:\allowbreak{}17}
\DERapparatusitem{ENV0010}{27;\allowbreak{}\allowbreak{} edition\allowbreak{} p27;\allowbreak{}\allowbreak{} ms49;\allowbreak{}\allowbreak{} description\allowbreak{} before\allowbreak{} square\allowbreak{} (\allowbreak{}*)\allowbreak{}}{'}{}{DER\_\allowbreak{}I-a:\allowbreak{}23}
\DERapparatusitem{ENV0011}{30;\allowbreak{}\allowbreak{} edition\allowbreak{} p30;\allowbreak{}\allowbreak{} ms54;\allowbreak{}\allowbreak{} first\allowbreak{} displayed\allowbreak{} dual\allowbreak{} fibre\allowbreak{} inclusion}{Y}{X}{DER\_\allowbreak{}I-a:\allowbreak{}25}
\DERapparatusitem{ENV0012}{32;\allowbreak{}\allowbreak{} edition\allowbreak{} p32;\allowbreak{}\allowbreak{} ms57;\allowbreak{}\allowbreak{} right\allowbreak{} strict\allowbreak{} fibre,\allowbreak{}\allowbreak{} and\allowbreak{} following\allowbreak{} cohomology\allowbreak{} square}{}{'}{DER\_\allowbreak{}I-a:\allowbreak{}26}
\DERapparatusitem{ENV0013}{32;\allowbreak{}\allowbreak{} edition\allowbreak{} p32;\allowbreak{}\allowbreak{} ms57;\allowbreak{}\allowbreak{} right\allowbreak{} strict\allowbreak{} fibre,\allowbreak{}\allowbreak{} and\allowbreak{} following\allowbreak{} cohomology\allowbreak{} square}{}{'}{DER\_\allowbreak{}I-a:\allowbreak{}26}
\DERapparatusitem{ENV0014}{33;\allowbreak{}\allowbreak{} edition\allowbreak{} p33;\allowbreak{}\allowbreak{} ms58;\allowbreak{}\allowbreak{} right\allowbreak{} direct\allowbreak{} image\allowbreak{} converse}{any\allowbreak{} right\allowbreak{} exact\allowbreak{} derivator}{any\allowbreak{} left\allowbreak{} exact\allowbreak{} derivator}{DER\_\allowbreak{}I-a:\allowbreak{}29}
\DERapparatusitem{ENV0015}{33;\allowbreak{}\allowbreak{} edition\allowbreak{} p33;\allowbreak{}\allowbreak{} ms58;\allowbreak{}\allowbreak{} end\allowbreak{} of\allowbreak{} homological\allowbreak{} paragraph}{a\allowbreak{} left\allowbreak{} exact\allowbreak{} derivator}{a\allowbreak{} right\allowbreak{} exact\allowbreak{} derivator}{DER\_\allowbreak{}I-a:\allowbreak{}30}
\DERapparatusitem{ENV0016}{34;\allowbreak{}\allowbreak{} printed34;\allowbreak{}\allowbreak{} ms60;\allowbreak{}\allowbreak{} product-square\allowbreak{} lifting\allowbreak{} category}{\_\allowbreak{}0,\allowbreak{}\allowbreak{} v'\_\allowbreak{}0)\allowbreak{}\textbackslash{}\allowbreak{}backslash(\allowbreak{}J/\allowbreak{}j}{'\_\allowbreak{}0\textbackslash{}\allowbreak{}backslash\allowbreak{} J)\allowbreak{}/\allowbreak{}(\allowbreak{}j\_\allowbreak{}0,\allowbreak{}\allowbreak{} v}{DER\_\allowbreak{}I-a:\allowbreak{}41}
\DERapparatusitem{ENV0017}{34;\allowbreak{}\allowbreak{} printed34;\allowbreak{}\allowbreak{} ms60;\allowbreak{}\allowbreak{} product-square\allowbreak{} lifting\allowbreak{} category}{which\allowbreak{} is\allowbreak{} contractible\allowbreak{} having\allowbreak{} an\allowbreak{} initial\allowbreak{} object}{which\allowbreak{} is\allowbreak{} contractible,\allowbreak{}\allowbreak{} being\allowbreak{} a\allowbreak{} product\allowbreak{} of\allowbreak{} a\allowbreak{} category\allowbreak{} having\allowbreak{} an\allowbreak{} initial\allowbreak{} object\allowbreak{} and\allowbreak{} a\allowbreak{} category\allowbreak{} having\allowbreak{} a\allowbreak{} final\allowbreak{} object}{DER\_\allowbreak{}I-a:\allowbreak{}41}
\DERapparatusitem{ENV0018}{35;\allowbreak{}\allowbreak{} running\allowbreak{} text\allowbreak{} or\allowbreak{} formula}{\textbackslash{}\allowbreak{}text\{\allowbreak{}def\}\allowbreak{}}{\textbackslash{}\allowbreak{}mathrm\{\allowbreak{}def\}\allowbreak{}}{TYPESETTING:\allowbreak{}7;\allowbreak{}\allowbreak{} NOTATION\_\allowbreak{}DER:\allowbreak{}17}
\DERapparatusitem{ENV0019}{37;\allowbreak{}\allowbreak{} printed37;\allowbreak{}\allowbreak{} ms64;\allowbreak{}\allowbreak{} Proposition2prime\allowbreak{} opening}{\$\allowbreak{}\textbackslash{}\allowbreak{}DD\$\allowbreak{}\allowbreak{} is\allowbreak{} a\allowbreak{} left\allowbreak{} derivator}{\$\allowbreak{}\textbackslash{}\allowbreak{}DD\$\allowbreak{}\allowbreak{} is\allowbreak{} a\allowbreak{} right\allowbreak{} derivator}{DER\_\allowbreak{}I-a:\allowbreak{}42}
\DERapparatusitem{ENV0020}{37;\allowbreak{}\allowbreak{} edition\allowbreak{} p37;\allowbreak{}\allowbreak{} ms64;\allowbreak{}\allowbreak{} Proposition2-prime\allowbreak{} condition\allowbreak{} c}{Y}{X}{DER\_\allowbreak{}I-a:\allowbreak{}37}
\DERapparatusitem{ENV0021}{39;\allowbreak{}\allowbreak{} printed39;\allowbreak{}\allowbreak{} ms67;\allowbreak{}\allowbreak{} Theorem1(\allowbreak{}a)\allowbreak{}}{(\allowbreak{}prop.\allowbreak{}\allowbreak{} 2)\allowbreak{}}{(\allowbreak{}prop.\allowbreak{}\allowbreak{} 2\$\allowbreak{}'\$\allowbreak{})\allowbreak{}}{DER\_\allowbreak{}I-b:\allowbreak{}1}
\DERapparatusitem{ENV0022}{41;\allowbreak{}\allowbreak{} printed41;\allowbreak{}\allowbreak{} ms69;\allowbreak{}\allowbreak{} Theorem2(\allowbreak{}dprime)\allowbreak{}}{so\allowbreak{} \$\allowbreak{}Y'\$\allowbreak{}\allowbreak{} is\allowbreak{} a\allowbreak{} category\allowbreak{} with\allowbreak{} final\allowbreak{} object\allowbreak{} \$\allowbreak{}e\$\allowbreak{}}{where\allowbreak{} \$\allowbreak{}e\textbackslash{}\allowbreak{}longrightarrow\allowbreak{} Y'\$\allowbreak{}\allowbreak{} selects\allowbreak{} a\allowbreak{} final\allowbreak{} object\allowbreak{} of\allowbreak{} \$\allowbreak{}Y'\$\allowbreak{}}{DER\_\allowbreak{}I-b:\allowbreak{}5}
\DERapparatusitem{ENV0023}{41;\allowbreak{}\allowbreak{} printed41;\allowbreak{}\allowbreak{} ms69;\allowbreak{}\allowbreak{} Theorem2\allowbreak{} proof}{Y}{X}{DER\_\allowbreak{}I-b:\allowbreak{}3}
\DERapparatusitem{ENV0024}{41;\allowbreak{}\allowbreak{} printed41;\allowbreak{}\allowbreak{} ms69;\allowbreak{}\allowbreak{} Theorem2\allowbreak{} proof}{Y}{X}{DER\_\allowbreak{}I-b:\allowbreak{}4}
\DERapparatusitem{ENV0025}{41;\allowbreak{}\allowbreak{} printed41;\allowbreak{}\allowbreak{} ms70;\allowbreak{}\allowbreak{} conclusion\allowbreak{} of\allowbreak{} cprimeprime\allowbreak{} implies\allowbreak{} a}{this\allowbreak{} implies\allowbreak{} that\allowbreak{} \$\allowbreak{}f\$\allowbreak{}\allowbreak{} is\allowbreak{} a\allowbreak{} co-strict}{this\allowbreak{} implies\allowbreak{} that\allowbreak{} \$\allowbreak{}q\$\allowbreak{}\allowbreak{} is\allowbreak{} a\allowbreak{} co-strict}{DER\_\allowbreak{}I-b:\allowbreak{}6}
\DERapparatusitem{ENV0026}{41;\allowbreak{}\allowbreak{} printed41;\allowbreak{}\allowbreak{} ms70;\allowbreak{}\allowbreak{} Corollary1}{j}{f}{DER\_\allowbreak{}I-b:\allowbreak{}7}
\DERapparatusitem{ENV0027}{41;\allowbreak{}\allowbreak{} printed41;\allowbreak{}\allowbreak{} ms70;\allowbreak{}\allowbreak{} Definition4}{Under\allowbreak{} the\allowbreak{} equivalent\allowbreak{} conditions\allowbreak{} of\allowbreak{} theorem\allowbreak{} 1}{Under\allowbreak{} the\allowbreak{} equivalent\allowbreak{} conditions\allowbreak{} of\allowbreak{} theorem\allowbreak{} 2}{DER\_\allowbreak{}I-b:\allowbreak{}8}
\DERapparatusitem{ENV0028}{42;\allowbreak{}\allowbreak{} printed42;\allowbreak{}\allowbreak{} ms71;\allowbreak{}\allowbreak{} footnote32}{Y}{S}{DER\_\allowbreak{}I-b:\allowbreak{}9}
\DERapparatusitem{ENV0029}{42;\allowbreak{}\allowbreak{} printed42;\allowbreak{}\allowbreak{} ms71;\allowbreak{}\allowbreak{} footnote32}{If\allowbreak{} \$\allowbreak{}f\$\allowbreak{}\allowbreak{} \textbackslash{}\allowbreak{}texttt\{\allowbreak{}[\allowbreak{}is]\allowbreak{}\}\allowbreak{}\allowbreak{} proper}{If\allowbreak{} \$\allowbreak{}p\$\allowbreak{}\allowbreak{} and\allowbreak{} \$\allowbreak{}q\$\allowbreak{}\allowbreak{} are\allowbreak{} \textbackslash{}\allowbreak{}texttt\{\allowbreak{}[\allowbreak{}is]\allowbreak{}\}\allowbreak{}\allowbreak{} proper}{DER\_\allowbreak{}I-b:\allowbreak{}10}
\DERapparatusitem{ENV0030}{43;\allowbreak{}\allowbreak{} printed43;\allowbreak{}\allowbreak{} ms72;\allowbreak{}\allowbreak{} q\_\allowbreak{}X\allowbreak{} label}{\$\allowbreak{}\textbackslash{}\allowbreak{}Cat\$\allowbreak{}-fibrant\textbackslash{}\allowbreak{}\textbackslash{}\allowbreak{}(\allowbreak{}hence\allowbreak{} smooth)\allowbreak{}}{\$\allowbreak{}\textbackslash{}\allowbreak{}Cat\$\allowbreak{}-cofibrant\textbackslash{}\allowbreak{}\textbackslash{}\allowbreak{}(\allowbreak{}hence\allowbreak{} proper)\allowbreak{}}{DER\_\allowbreak{}I-b:\allowbreak{}12}
\DERapparatusitem{ENV0031}{44;\allowbreak{}\allowbreak{} printed44;\allowbreak{}\allowbreak{} ms73;\allowbreak{}\allowbreak{} W-properness\allowbreak{} note}{\$\allowbreak{}p\_\allowbreak{}X\$\allowbreak{}\allowbreak{} is\allowbreak{} \$\allowbreak{}W\$\allowbreak{}-smooth}{\$\allowbreak{}p\_\allowbreak{}X\$\allowbreak{}\allowbreak{} is\allowbreak{} \$\allowbreak{}W\$\allowbreak{}-proper}{DER\_\allowbreak{}I-b:\allowbreak{}14}
\DERapparatusitem{ENV0032}{44;\allowbreak{}\allowbreak{} printed44;\allowbreak{}\allowbreak{} ms74;\allowbreak{}\allowbreak{} q\_\allowbreak{}Z\allowbreak{} strict-equivalence\allowbreak{} test}{z\_\allowbreak{}0\textbackslash{}\allowbreak{}bs\textbackslash{}\allowbreak{}Phi(\allowbreak{}Z)\allowbreak{}}{\textbackslash{}\allowbreak{}Phi(\allowbreak{}Z)\allowbreak{}/\allowbreak{}z\_\allowbreak{}0\textasciicircum{}\allowbreak{}\{\allowbreak{}\textbackslash{}\allowbreak{}mathrm\{\allowbreak{}op\}\allowbreak{}\}\allowbreak{}}{DER\_\allowbreak{}I-b:\allowbreak{}16}
\DERapparatusitem{ENV0033}{44;\allowbreak{}\allowbreak{} printed44;\allowbreak{}\allowbreak{} ms74;\allowbreak{}\allowbreak{} q\_\allowbreak{}Z\allowbreak{} strict-equivalence\allowbreak{} test}{z\_\allowbreak{}0\textbackslash{}\allowbreak{}bs\textbackslash{}\allowbreak{}Phi(\allowbreak{}Z)\allowbreak{}}{\textbackslash{}\allowbreak{}Phi(\allowbreak{}Z)\allowbreak{}/\allowbreak{}z\_\allowbreak{}0\textasciicircum{}\allowbreak{}\{\allowbreak{}\textbackslash{}\allowbreak{}mathrm\{\allowbreak{}op\}\allowbreak{}\}\allowbreak{}}{DER\_\allowbreak{}I-b:\allowbreak{}16}
\DERapparatusitem{ENV0034}{47;\allowbreak{}\allowbreak{} printed47;\allowbreak{}\allowbreak{} ms78;\allowbreak{}\allowbreak{} uninteresting\allowbreak{} empty-fibre\allowbreak{} cases}{\$\allowbreak{}Y'\_\allowbreak{}0\allowbreak{} =\allowbreak{}\allowbreak{} \textbackslash{}\allowbreak{}emptyset\$\allowbreak{},\allowbreak{}\allowbreak{} where\allowbreak{} \$\allowbreak{}Y'\_\allowbreak{}1\allowbreak{} =\allowbreak{}\allowbreak{} \textbackslash{}\allowbreak{}emptyset\$\allowbreak{}}{\$\allowbreak{}Y'\_\allowbreak{}0\allowbreak{} =\allowbreak{}\allowbreak{} \textbackslash{}\allowbreak{}emptyset\$\allowbreak{},\allowbreak{}\allowbreak{} or\allowbreak{} \$\allowbreak{}Y'\_\allowbreak{}1\allowbreak{} =\allowbreak{}\allowbreak{} \textbackslash{}\allowbreak{}emptyset\$\allowbreak{}}{DER\_\allowbreak{}I-b:\allowbreak{}17}
\DERapparatusitem{ENV0035}{50;\allowbreak{}\allowbreak{} ms81,\allowbreak{}\allowbreak{} first\allowbreak{} cohomology\allowbreak{} display}{,\allowbreak{}M)\allowbreak{}\allowbreak{} \textbackslash{}\allowbreak{}longrightarrow\allowbreak{} H\textasciicircum{}\allowbreak{}*\_\allowbreak{}\{\allowbreak{}\textbackslash{}\allowbreak{}DD\}\allowbreak{}(\allowbreak{}X'}{',\allowbreak{}M)\allowbreak{}\allowbreak{} \textbackslash{}\allowbreak{}longrightarrow\allowbreak{} H\textasciicircum{}\allowbreak{}*\_\allowbreak{}\{\allowbreak{}\textbackslash{}\allowbreak{}DD\}\allowbreak{}(\allowbreak{}X}{DER\_\allowbreak{}I-b:\allowbreak{}19}
\DERapparatusitem{ENV0036}{50;\allowbreak{}\allowbreak{} ms81,\allowbreak{}\allowbreak{} second\allowbreak{} cohomology\allowbreak{} display}{\_\allowbreak{}*(\allowbreak{}M\_\allowbreak{}X)\allowbreak{})\allowbreak{}\allowbreak{} \textbackslash{}\allowbreak{}longrightarrow\allowbreak{} H\textasciicircum{}\allowbreak{}*\_\allowbreak{}\{\allowbreak{}\textbackslash{}\allowbreak{}DD\}\allowbreak{}(\allowbreak{}Y,\allowbreak{}\allowbreak{} p'\_\allowbreak{}*(\allowbreak{}M\_\allowbreak{}\{\allowbreak{}X'\}\allowbreak{}}{'\_\allowbreak{}*(\allowbreak{}M\_\allowbreak{}\{\allowbreak{}X'\}\allowbreak{})\allowbreak{})\allowbreak{}\allowbreak{} \textbackslash{}\allowbreak{}longrightarrow\allowbreak{} H\textasciicircum{}\allowbreak{}*\_\allowbreak{}\{\allowbreak{}\textbackslash{}\allowbreak{}DD\}\allowbreak{}(\allowbreak{}Y,\allowbreak{}p\_\allowbreak{}*(\allowbreak{}M\_\allowbreak{}X}{DER\_\allowbreak{}I-b:\allowbreak{}20}
\DERapparatusitem{ENV0037}{50;\allowbreak{}\allowbreak{} ms81,\allowbreak{}\allowbreak{} direct-image\allowbreak{} display}{\_\allowbreak{}*(\allowbreak{}M\_\allowbreak{}X)\allowbreak{}\allowbreak{} \textbackslash{}\allowbreak{}longrightarrow\allowbreak{} p'\_\allowbreak{}*(\allowbreak{}M\_\allowbreak{}\{\allowbreak{}X'\}\allowbreak{}}{'\_\allowbreak{}*(\allowbreak{}M\_\allowbreak{}\{\allowbreak{}X'\}\allowbreak{})\allowbreak{}\allowbreak{} \textbackslash{}\allowbreak{}longrightarrow\allowbreak{} p\_\allowbreak{}*(\allowbreak{}M\_\allowbreak{}X}{DER\_\allowbreak{}I-b:\allowbreak{}21}
\DERapparatusitem{ENV0038}{50;\allowbreak{}\allowbreak{} ms82,\allowbreak{}\allowbreak{} special\allowbreak{} coefficient\allowbreak{} after\allowbreak{} Question}{p\_\allowbreak{}*(\allowbreak{}M\_\allowbreak{}Y)\allowbreak{}}{p\_\allowbreak{}*(\allowbreak{}M\_\allowbreak{}\{\allowbreak{}Y'\}\allowbreak{})\allowbreak{}}{DER\_\allowbreak{}I-b:\allowbreak{}22}
\DERapparatusitem{ENV0039}{51;\allowbreak{}\allowbreak{} ms83,\allowbreak{}\allowbreak{} opfibration\allowbreak{} corollary}{u}{y}{DER\_\allowbreak{}I-b:\allowbreak{}24}
\DERapparatusitem{ENV0040}{51;\allowbreak{}\allowbreak{} ms84,\allowbreak{}\allowbreak{} adjunction\allowbreak{} counit}{}{F}{DER\_\allowbreak{}I-b:\allowbreak{}27}
\DERapparatusitem{ENV0041}{52;\allowbreak{}\allowbreak{} ms84,\allowbreak{}\allowbreak{} Proposition6(\allowbreak{}a)\allowbreak{}}{Y}{X'}{DER\_\allowbreak{}I-b:\allowbreak{}28}
\DERapparatusitem{ENV0042}{52;\allowbreak{}\allowbreak{} ms84,\allowbreak{}\allowbreak{} footnote40}{But\allowbreak{} OK\allowbreak{} if\allowbreak{} \$\allowbreak{}S\$\allowbreak{}\allowbreak{} is\allowbreak{} \$\allowbreak{}\textbackslash{}\allowbreak{}DD\$\allowbreak{}-contractible.\allowbreak{}}{But\allowbreak{} OK\allowbreak{} if\allowbreak{} \$\allowbreak{}S\$\allowbreak{}\allowbreak{} is\allowbreak{} completely\allowbreak{} \$\allowbreak{}\textbackslash{}\allowbreak{}DD\$\allowbreak{}-aspheric.\allowbreak{}}{DER\_\allowbreak{}I-b:\allowbreak{}29}
\DERapparatusitem{ENV0043}{52;\allowbreak{}\allowbreak{} ms86,\allowbreak{}\allowbreak{} opening\allowbreak{} sentence}{Y}{X}{DER\_\allowbreak{}I-b:\allowbreak{}30}
\DERapparatusitem{ENV0044}{53;\allowbreak{}\allowbreak{} ms87,\allowbreak{}\allowbreak{} monodromy\allowbreak{} paragraph}{given\allowbreak{} up\allowbreak{} to\allowbreak{} unique\allowbreak{} isomorphism}{given\allowbreak{} up\allowbreak{} to\allowbreak{} isomorphism}{DER\_\allowbreak{}I-b:\allowbreak{}33}
\DERapparatusitem{ENV0045}{57;\allowbreak{}\allowbreak{} ms93,\allowbreak{}\allowbreak{} condition(\allowbreak{}b)\allowbreak{}}{\textasciicircum{}\allowbreak{}G}{}{DER\_\allowbreak{}I-b:\allowbreak{}42}
\DERapparatusitem{ENV0046}{58;\allowbreak{}\allowbreak{} ms95,\allowbreak{}\allowbreak{} quotient-category\allowbreak{} relations}{-}{+}{DER\_\allowbreak{}I-b:\allowbreak{}46}
\DERapparatusitem{ENV0047}{58;\allowbreak{}\allowbreak{} ms95,\allowbreak{}\allowbreak{} quotient-category\allowbreak{} relations}{+\allowbreak{} p\_\allowbreak{}0\textasciicircum{}\allowbreak{}-\allowbreak{} =\allowbreak{}\allowbreak{} p\_\allowbreak{}1\textasciicircum{}\allowbreak{}+\allowbreak{} p\_\allowbreak{}0\textasciicircum{}\allowbreak{}+}{-\allowbreak{} p\_\allowbreak{}0\textasciicircum{}\allowbreak{}+\allowbreak{} =\allowbreak{}\allowbreak{} p\_\allowbreak{}1\textasciicircum{}\allowbreak{}+\allowbreak{} p\_\allowbreak{}0\textasciicircum{}\allowbreak{}-}{DER\_\allowbreak{}I-b:\allowbreak{}46}
\DERapparatusitem{ENV0048}{60;\allowbreak{}\allowbreak{} ms98,\allowbreak{}\allowbreak{} first\allowbreak{} paragraph\allowbreak{} closing\allowbreak{} parenthesis}{by\allowbreak{} virtue\allowbreak{} of\allowbreak{} \textbackslash{}\allowbreak{}textbf\{\allowbreak{}Der\}\allowbreak{}\allowbreak{} 3\allowbreak{} (\allowbreak{}\$\allowbreak{}a'\$\allowbreak{})\allowbreak{}.\allowbreak{}}{by\allowbreak{} virtue\allowbreak{} of\allowbreak{} \textbackslash{}\allowbreak{}textbf\{\allowbreak{}Der\}\allowbreak{}\allowbreak{} 3\allowbreak{} (\allowbreak{}\$\allowbreak{}a'\$\allowbreak{})\allowbreak{})\allowbreak{}.\allowbreak{}}{DER\_\allowbreak{}I-b:\allowbreak{}49}
\DERapparatusitem{ENV0049}{60;\allowbreak{}\allowbreak{} ms99,\allowbreak{}\allowbreak{} definition\allowbreak{} of1-complete}{We\allowbreak{} say\allowbreak{} that\allowbreak{} \$\allowbreak{}f\$\allowbreak{}\allowbreak{} and\allowbreak{} a\allowbreak{} 1-}{We\allowbreak{} say\allowbreak{} that\allowbreak{} \$\allowbreak{}f\$\allowbreak{}\allowbreak{} is\allowbreak{} a\allowbreak{} 1-}{DER\_\allowbreak{}I-b:\allowbreak{}51}
\DERapparatusitem{ENV0050}{61;\allowbreak{}\allowbreak{} ms100,\allowbreak{}\allowbreak{} comparison\allowbreak{} with\allowbreak{} Proposition2}{a\allowbreak{} left\allowbreak{} \$\allowbreak{}\textbackslash{}\allowbreak{}DD\$\allowbreak{}\allowbreak{} fibration}{a\allowbreak{} left\allowbreak{} \$\allowbreak{}\textbackslash{}\allowbreak{}DD\$\allowbreak{}-fibration}{DER\_\allowbreak{}I-b:\allowbreak{}53}
\DERapparatusitem{ENV0051}{62;\allowbreak{}\allowbreak{} Printed62;\allowbreak{}\allowbreak{} ms102;\allowbreak{}\allowbreak{} Example}{If\allowbreak{} \$\allowbreak{}f\$\allowbreak{}\allowbreak{} is\allowbreak{} \textbackslash{}\allowbreak{}Hot-proper}{If\allowbreak{} \$\allowbreak{}f\$\allowbreak{}\allowbreak{} is\allowbreak{} a\allowbreak{} \textbackslash{}\allowbreak{}Hot-fibration\allowbreak{} and\allowbreak{} is\allowbreak{} \textbackslash{}\allowbreak{}Hot-proper}{DER\_\allowbreak{}I-b:\allowbreak{}54}
\DERapparatusitem{ENV0052}{62;\allowbreak{}\allowbreak{} Printed62;\allowbreak{}\allowbreak{} ms102;\allowbreak{}\allowbreak{} footnote54}{prop.\allowbreak{}\allowbreak{} 3\$\allowbreak{}'\$\allowbreak{}.\allowbreak{}}{prop.\allowbreak{}\allowbreak{} 4.\allowbreak{}}{DER\_\allowbreak{}I-b:\allowbreak{}55}
\DERapparatusitem{ENV0053}{63;\allowbreak{}\allowbreak{} Printed63;\allowbreak{}\allowbreak{} ms106;\allowbreak{}\allowbreak{} final\allowbreak{} sentence}{then\allowbreak{} \$\allowbreak{}X''\allowbreak{} \textbackslash{}\allowbreak{}to\allowbreak{} X'\$\allowbreak{}\allowbreak{} too.\allowbreak{}}{then\allowbreak{} \$\allowbreak{}X''\allowbreak{} \textbackslash{}\allowbreak{}to\allowbreak{} X'\$\allowbreak{}\allowbreak{} too)\allowbreak{}.\allowbreak{}}{DER\_\allowbreak{}I-b:\allowbreak{}57}
\DERapparatusitem{ENV0054}{65;\allowbreak{}\allowbreak{} Printed65;\allowbreak{}\allowbreak{} ms109;\allowbreak{}\allowbreak{} lemma\allowbreak{} final\allowbreak{} formula}{q'}{f}{DER\_\allowbreak{}I-b:\allowbreak{}60}
\DERapparatusitem{ENV0055}{65;\allowbreak{}\allowbreak{} Printed65;\allowbreak{}\allowbreak{} ms110;\allowbreak{}\allowbreak{} paragraph\allowbreak{} after\allowbreak{} Yoneda\allowbreak{} display}{''}{}{DER\_\allowbreak{}I-b:\allowbreak{}61}
\DERapparatusitem{ENV0056}{66;\allowbreak{}\allowbreak{} Printed66;\allowbreak{}\allowbreak{} ms111;\allowbreak{}\allowbreak{} Der3\allowbreak{} b\allowbreak{} parenthesis}{\textbackslash{}\allowbreak{}DD(\allowbreak{}Z)\allowbreak{}\_\allowbreak{}\{\allowbreak{}\textbackslash{}\allowbreak{}text\{\allowbreak{}loc.\allowbreak{}\allowbreak{} const.\allowbreak{}\}\allowbreak{}\}\allowbreak{}}{\textbackslash{}\allowbreak{}DD(\allowbreak{}Z)\allowbreak{}\_\allowbreak{}\{\allowbreak{}\textbackslash{}\allowbreak{}text\{\allowbreak{}const.\allowbreak{}\}\allowbreak{}\}\allowbreak{}}{DER\_\allowbreak{}I-b:\allowbreak{}64}
\DERapparatusitem{ENV0057}{69;\allowbreak{}\allowbreak{} Printed69;\allowbreak{}\allowbreak{} ms117;\allowbreak{}\allowbreak{} Proposition5(\allowbreak{}a2)\allowbreak{}}{has\allowbreak{} a\allowbreak{} cofinal\allowbreak{} object}{has\allowbreak{} an\allowbreak{} initial\allowbreak{} object}{DER\_\allowbreak{}I-b:\allowbreak{}65}
\DERapparatusitem{ENV0058}{70;\allowbreak{}\allowbreak{} Printed70;\allowbreak{}\allowbreak{} ms117;\allowbreak{}\allowbreak{} Definition3}{of\allowbreak{} proposition\allowbreak{} 4,\allowbreak{}}{of\allowbreak{} proposition\allowbreak{} 5,\allowbreak{}}{DER\_\allowbreak{}I-b:\allowbreak{}66}
\DERapparatusitem{ENV0059}{71;\allowbreak{}\allowbreak{} Printed71;\allowbreak{}\allowbreak{} ms119;\allowbreak{}\allowbreak{} following\allowbreak{} paragraph}{}{'}{DER\_\allowbreak{}I-b:\allowbreak{}69}
\DERapparatusitem{ENV0060}{72;\allowbreak{}\allowbreak{} Printed72;\allowbreak{}\allowbreak{} ms122;\allowbreak{}\allowbreak{} last\allowbreak{} sentence}{the\allowbreak{} notion\allowbreak{} of\allowbreak{} \$\allowbreak{}\textbackslash{}\allowbreak{}D\$\allowbreak{}-equivalence\allowbreak{} already\allowbreak{} implies}{the\allowbreak{} notion\allowbreak{} of\allowbreak{} cohomological\allowbreak{} \$\allowbreak{}\textbackslash{}\allowbreak{}D\$\allowbreak{}-equivalence\allowbreak{} already\allowbreak{} implies}{DER\_\allowbreak{}I-b:\allowbreak{}72}
\subsection*{Chapter II}
\DERapparatusitem{ENV0061}{1;\allowbreak{}\allowbreak{} Printed\allowbreak{} p.\allowbreak{}1;\allowbreak{}\allowbreak{} ms.\allowbreak{}1;\allowbreak{}\allowbreak{} first\allowbreak{} proposition}{\textbackslash{}\allowbreak{}alpha(\allowbreak{}x')\allowbreak{}\allowbreak{} \textbackslash{}\allowbreak{}longrightarrow\allowbreak{} x}{x\_\allowbreak{}0\allowbreak{} \textbackslash{}\allowbreak{}longrightarrow\allowbreak{} \textbackslash{}\allowbreak{}alpha(\allowbreak{}x')\allowbreak{}}{DER\_\allowbreak{}II-a:\allowbreak{}1}
\DERapparatusitem{ENV0062}{1;\allowbreak{}\allowbreak{} Printed\allowbreak{} pp.\allowbreak{}1–2;\allowbreak{}\allowbreak{} ms.\allowbreak{}1–2;\allowbreak{}\allowbreak{} sentence\allowbreak{} after\allowbreak{} unit}{which\allowbreak{} also\allowbreak{} means\allowbreak{} that\allowbreak{} for\allowbreak{} every\allowbreak{} arrow}{which\allowbreak{} is\allowbreak{} ensured\allowbreak{} when\allowbreak{} for\allowbreak{} every\allowbreak{} arrow}{DER\_\allowbreak{}II-a:\allowbreak{}2}
\DERapparatusitem{ENV0063}{2;\allowbreak{}\allowbreak{} Printed\allowbreak{} p.\allowbreak{}2;\allowbreak{}\allowbreak{} ms.\allowbreak{}3;\allowbreak{}\allowbreak{} triple\allowbreak{} defining\allowbreak{} the\allowbreak{} comma\allowbreak{} category}{g}{f'}{DER\_\allowbreak{}II-a:\allowbreak{}3}
\DERapparatusitem{ENV0064}{4;\allowbreak{}\allowbreak{} Printed\allowbreak{} p.\allowbreak{}4;\allowbreak{}\allowbreak{} ms.\allowbreak{}6;\allowbreak{}\allowbreak{} transported\allowbreak{} section\allowbreak{} formula}{\textbackslash{}\allowbreak{}xi}{\textbackslash{}\allowbreak{}xi'}{DER\_\allowbreak{}II-a:\allowbreak{}6}
\DERapparatusitem{ENV0065}{4;\allowbreak{}\allowbreak{} Printed\allowbreak{} p.\allowbreak{}4;\allowbreak{}\allowbreak{} ms.\allowbreak{}6;\allowbreak{}\allowbreak{} transported\allowbreak{} section\allowbreak{} formula}{\textbackslash{}\allowbreak{}xi}{\textbackslash{}\allowbreak{}xi'}{DER\_\allowbreak{}II-a:\allowbreak{}6}
\DERapparatusitem{ENV0066}{4;\allowbreak{}\allowbreak{} Printed\allowbreak{} p.\allowbreak{}4;\allowbreak{}\allowbreak{} ms.\allowbreak{}6;\allowbreak{}\allowbreak{} transported\allowbreak{} section\allowbreak{} formula}{\textbackslash{}\allowbreak{}xi}{\textbackslash{}\allowbreak{}xi'}{DER\_\allowbreak{}II-a:\allowbreak{}6}
\DERapparatusitem{ENV0067}{4;\allowbreak{}\allowbreak{} Printed\allowbreak{} p.\allowbreak{}4;\allowbreak{}\allowbreak{} ms.\allowbreak{}6;\allowbreak{}\allowbreak{} transported\allowbreak{} section\allowbreak{} formula}{\textbackslash{}\allowbreak{}xi}{\textbackslash{}\allowbreak{}xi'}{DER\_\allowbreak{}II-a:\allowbreak{}6}
\DERapparatusitem{ENV0068}{4;\allowbreak{}\allowbreak{} running\allowbreak{} text\allowbreak{} or\allowbreak{} formula}{\textbackslash{}\allowbreak{}text\{\allowbreak{}def\}\allowbreak{}}{\textbackslash{}\allowbreak{}mathrm\{\allowbreak{}def\}\allowbreak{}}{TYPESETTING:\allowbreak{}7;\allowbreak{}\allowbreak{} NOTATION\_\allowbreak{}DER:\allowbreak{}17}
\DERapparatusitem{ENV0069}{4;\allowbreak{}\allowbreak{} Printed\allowbreak{} p.\allowbreak{}4;\allowbreak{}\allowbreak{} ms.\allowbreak{}6;\allowbreak{}\allowbreak{} completed\allowbreak{} cone}{\textbackslash{}\allowbreak{}simeq}{\textbackslash{}\allowbreak{}in}{DER\_\allowbreak{}II-a:\allowbreak{}7}
\DERapparatusitem{ENV0070}{4;\allowbreak{}\allowbreak{} Printed\allowbreak{} p.\allowbreak{}4;\allowbreak{}\allowbreak{} ms.\allowbreak{}6;\allowbreak{}\allowbreak{} paragraph\allowbreak{} after\allowbreak{} cqfd}{\_\allowbreak{}0}{}{DER\_\allowbreak{}II-a:\allowbreak{}8}
\DERapparatusitem{ENV0071}{4;\allowbreak{}\allowbreak{} Printed\allowbreak{} p.\allowbreak{}4;\allowbreak{}\allowbreak{} ms.\allowbreak{}6;\allowbreak{}\allowbreak{} map\allowbreak{} to\allowbreak{} the\allowbreak{} constant\allowbreak{} coefficient}{\textbackslash{}\allowbreak{}xi\textasciicircum{}\allowbreak{}*(\allowbreak{}u'}{u\textasciicircum{}\allowbreak{}*(\allowbreak{}\textbackslash{}\allowbreak{}xi'(\allowbreak{}x')\allowbreak{}}{DER\_\allowbreak{}II-a:\allowbreak{}9}
\DERapparatusitem{ENV0072}{4;\allowbreak{}\allowbreak{} Printed\allowbreak{} p.\allowbreak{}4;\allowbreak{}\allowbreak{} ms.\allowbreak{}6;\allowbreak{}\allowbreak{} constant\allowbreak{} presheaf\allowbreak{} section\allowbreak{} prose}{}{\_\allowbreak{}0}{DER\_\allowbreak{}II-a:\allowbreak{}10}
\DERapparatusitem{ENV0073}{4;\allowbreak{}\allowbreak{} Printed\allowbreak{} p.\allowbreak{}4;\allowbreak{}\allowbreak{} ms.\allowbreak{}6;\allowbreak{}\allowbreak{} parenthesis\allowbreak{} after\allowbreak{} c\_\allowbreak{}x0}{}{'}{DER\_\allowbreak{}II-a:\allowbreak{}11}
\DERapparatusitem{ENV0074}{4;\allowbreak{}\allowbreak{} Printed\allowbreak{} p.\allowbreak{}4;\allowbreak{}\allowbreak{} ms.\allowbreak{}6;\allowbreak{}\allowbreak{} boxed\allowbreak{} arrow\allowbreak{} v}{'}{}{DER\_\allowbreak{}II-a:\allowbreak{}12}
\DERapparatusitem{ENV0075}{4;\allowbreak{}\allowbreak{} Printed\allowbreak{} p.\allowbreak{}4;\allowbreak{}\allowbreak{} ms.\allowbreak{}6;\allowbreak{}\allowbreak{} final\allowbreak{} phrase}{on\allowbreak{} \$\allowbreak{}x\_\allowbreak{}0\textbackslash{}\allowbreak{}backslash\allowbreak{} C'\$\allowbreak{}\allowbreak{} resp.\allowbreak{}\allowbreak{} \$\allowbreak{}x\_\allowbreak{}1\textbackslash{}\allowbreak{}backslash\allowbreak{} C'\$\allowbreak{}}{on\allowbreak{} \$\allowbreak{}x\_\allowbreak{}1\textbackslash{}\allowbreak{}backslash\allowbreak{} C'\$\allowbreak{}\allowbreak{} resp.\allowbreak{}\allowbreak{} \$\allowbreak{}x\_\allowbreak{}0\textbackslash{}\allowbreak{}backslash\allowbreak{} C'\$\allowbreak{}}{DER\_\allowbreak{}II-a:\allowbreak{}13}
\DERapparatusitem{ENV0076}{5;\allowbreak{}\allowbreak{} Printed\allowbreak{} p.\allowbreak{}5;\allowbreak{}\allowbreak{} ms.\allowbreak{}7;\allowbreak{}\allowbreak{} upper-left\allowbreak{} arrows\allowbreak{} of\allowbreak{} large\allowbreak{} diagram}{1}{0}{DER\_\allowbreak{}II-a:\allowbreak{}14}
\DERapparatusitem{ENV0077}{5;\allowbreak{}\allowbreak{} Printed\allowbreak{} p.\allowbreak{}5;\allowbreak{}\allowbreak{} ms.\allowbreak{}7;\allowbreak{}\allowbreak{} upper-left\allowbreak{} arrows\allowbreak{} of\allowbreak{} large\allowbreak{} diagram}{0}{1}{DER\_\allowbreak{}II-a:\allowbreak{}14}
\DERapparatusitem{ENV0078}{6;\allowbreak{}\allowbreak{} running\allowbreak{} text\allowbreak{} or\allowbreak{} formula}{\textbackslash{}\allowbreak{}text\{\allowbreak{}def\}\allowbreak{}}{\textbackslash{}\allowbreak{}mathrm\{\allowbreak{}def\}\allowbreak{}}{TYPESETTING:\allowbreak{}7;\allowbreak{}\allowbreak{} NOTATION\_\allowbreak{}DER:\allowbreak{}17}
\DERapparatusitem{ENV0079}{6;\allowbreak{}\allowbreak{} Printed\allowbreak{} p.\allowbreak{}6;\allowbreak{}\allowbreak{} ms.\allowbreak{}9;\allowbreak{}\allowbreak{} unit\allowbreak{} internal-Hom\allowbreak{} display}{\textbackslash{}\allowbreak{}Homint(\allowbreak{}\textbackslash{}\allowbreak{}varepsilon\_\allowbreak{}\{\allowbreak{}C'\}\allowbreak{},\allowbreak{}}{\textbackslash{}\allowbreak{}Homint\_\allowbreak{}\{\allowbreak{}C'\}\allowbreak{}(\allowbreak{}\textbackslash{}\allowbreak{}varepsilon\_\allowbreak{}\{\allowbreak{}C'\}\allowbreak{},\allowbreak{}\textbackslash{}\allowbreak{}alpha\textasciicircum{}\allowbreak{}*}{DER\_\allowbreak{}II-a:\allowbreak{}15}
\DERapparatusitem{ENV0080}{6;\allowbreak{}\allowbreak{} Printed\allowbreak{} p.\allowbreak{}6;\allowbreak{}\allowbreak{} ms.\allowbreak{}10;\allowbreak{}\allowbreak{} projection\allowbreak{} display}{G\allowbreak{} \textbackslash{}\allowbreak{}otimes\allowbreak{} \textbackslash{}\allowbreak{}alpha\_\allowbreak{}!G'\allowbreak{} \textbackslash{}\allowbreak{}xrightarrow\{\allowbreak{}\textbackslash{}\allowbreak{}\allowbreak{} \textbackslash{}\allowbreak{}sim\textbackslash{}\allowbreak{}\allowbreak{} \}\allowbreak{}\allowbreak{} \textbackslash{}\allowbreak{}alpha\_\allowbreak{}!(\allowbreak{}\textbackslash{}\allowbreak{}alpha\textasciicircum{}\allowbreak{}*G\allowbreak{} \textbackslash{}\allowbreak{}otimes\allowbreak{} G')\allowbreak{}}{\textbackslash{}\allowbreak{}alpha\_\allowbreak{}!(\allowbreak{}\textbackslash{}\allowbreak{}alpha\textasciicircum{}\allowbreak{}*G\allowbreak{} \textbackslash{}\allowbreak{}otimes\allowbreak{} G')\allowbreak{}\allowbreak{} \textbackslash{}\allowbreak{}xrightarrow\{\allowbreak{}\textbackslash{}\allowbreak{}\allowbreak{} \textbackslash{}\allowbreak{}sim\textbackslash{}\allowbreak{}\allowbreak{} \}\allowbreak{}\allowbreak{} G\allowbreak{} \textbackslash{}\allowbreak{}otimes\allowbreak{} \textbackslash{}\allowbreak{}alpha\_\allowbreak{}!G'}{DER\_\allowbreak{}II-a:\allowbreak{}16}
\DERapparatusitem{ENV0081}{7;\allowbreak{}\allowbreak{} Printed\allowbreak{} p.\allowbreak{}7;\allowbreak{}\allowbreak{} ms.\allowbreak{}11;\allowbreak{}\allowbreak{} topological-space\allowbreak{} paragraph}{\textbackslash{}\allowbreak{}mathcal}{}{DER\_\allowbreak{}II-a:\allowbreak{}17}
\DERapparatusitem{ENV0082}{11;\allowbreak{}\allowbreak{} Introduction\allowbreak{} to\allowbreak{} Proposition2.\allowbreak{}3,\allowbreak{}\allowbreak{} edition\allowbreak{} p.\allowbreak{}11}{\textbackslash{}\allowbreak{}mathrm\{\allowbreak{}Fais\}\allowbreak{}(\allowbreak{}\textbackslash{}\allowbreak{}sX)\allowbreak{}\allowbreak{} =\allowbreak{}\allowbreak{} S\textasciicircum{}\allowbreak{}\{\allowbreak{}\textbackslash{}\allowbreak{}wedge\}\allowbreak{}}{\textbackslash{}\allowbreak{}mathrm\{\allowbreak{}Fais\}\allowbreak{}(\allowbreak{}\textbackslash{}\allowbreak{}sX)\allowbreak{}\allowbreak{} \textbackslash{}\allowbreak{}simeq\allowbreak{} S\textasciicircum{}\allowbreak{}\{\allowbreak{}\textbackslash{}\allowbreak{}sim\}\allowbreak{}}{DER\_\allowbreak{}II-a:\allowbreak{}45}
\DERapparatusitem{ENV0083}{11;\allowbreak{}\allowbreak{} Proposition\allowbreak{} 2.\allowbreak{}3(\allowbreak{}b)\allowbreak{},\allowbreak{}\allowbreak{} edition\allowbreak{} p.\allowbreak{}11}{U}{V}{DER\_\allowbreak{}II-a:\allowbreak{}23}
\DERapparatusitem{ENV0084}{13;\allowbreak{}\allowbreak{} Corollary\allowbreak{} 2.\allowbreak{}7,\allowbreak{}\allowbreak{} final\allowbreak{} quantified\allowbreak{} reformulation,\allowbreak{}\allowbreak{} edition\allowbreak{} p.\allowbreak{}13}{or\allowbreak{} again\allowbreak{} that\allowbreak{} for\allowbreak{} two\allowbreak{} objects\allowbreak{} \$\allowbreak{}a\$\allowbreak{},\allowbreak{}\allowbreak{} \$\allowbreak{}b\$\allowbreak{}\allowbreak{} in\allowbreak{} \$\allowbreak{}X\$\allowbreak{}}{or\allowbreak{} again\allowbreak{} that\allowbreak{} \$\allowbreak{}X\$\allowbreak{}\allowbreak{} be\allowbreak{} nonempty\allowbreak{} and\allowbreak{} that\allowbreak{} for\allowbreak{} two\allowbreak{} objects\allowbreak{} \$\allowbreak{}a\$\allowbreak{},\allowbreak{}\allowbreak{} \$\allowbreak{}b\$\allowbreak{}\allowbreak{} in\allowbreak{} \$\allowbreak{}X\$\allowbreak{}}{DER\_\allowbreak{}II-a:\allowbreak{}26}
\DERapparatusitem{ENV0085}{13;\allowbreak{}\allowbreak{} Corollary\allowbreak{} 2.\allowbreak{}8(\allowbreak{}b)\allowbreak{},\allowbreak{}\allowbreak{} edition\allowbreak{} p.\allowbreak{}13}{of\allowbreak{} \$\allowbreak{}0\$\allowbreak{}-connected\allowbreak{} objects\allowbreak{} \$\allowbreak{}U\$\allowbreak{},\allowbreak{}\allowbreak{} \$\allowbreak{}V\$\allowbreak{}}{of\allowbreak{} \$\allowbreak{}0\$\allowbreak{}-connected\allowbreak{} objects\allowbreak{} \$\allowbreak{}U\$\allowbreak{},\allowbreak{}\allowbreak{} \$\allowbreak{}V,\allowbreak{}W\$\allowbreak{}}{DER\_\allowbreak{}II-a:\allowbreak{}27}
\DERapparatusitem{ENV0086}{15;\allowbreak{}\allowbreak{} Proof\allowbreak{} after\allowbreak{} equalizer\allowbreak{} square,\allowbreak{}\allowbreak{} manuscript\allowbreak{} p.\allowbreak{}24\allowbreak{} /\allowbreak{}\allowbreak{} edition\allowbreak{} p.\allowbreak{}15}{c)\allowbreak{}\allowbreak{} holds}{b)\allowbreak{}\allowbreak{} holds}{DER\_\allowbreak{}II-a:\allowbreak{}28}
\DERapparatusitem{ENV0087}{15;\allowbreak{}\allowbreak{} Same\allowbreak{} paragraph,\allowbreak{}\allowbreak{} manuscript\allowbreak{} p.\allowbreak{}24}{\textbackslash{}\allowbreak{}X\textbackslash{}\allowbreak{}times\textbackslash{}\allowbreak{}X}{\textbackslash{}\allowbreak{}X}{DER\_\allowbreak{}II-a:\allowbreak{}29}
\DERapparatusitem{ENV0088}{16;\allowbreak{}\allowbreak{} Continuation\allowbreak{} of\allowbreak{} proof,\allowbreak{}\allowbreak{} manuscript\allowbreak{} p.\allowbreak{}25\allowbreak{} /\allowbreak{}\allowbreak{} edition\allowbreak{} p.\allowbreak{}16}{proof\allowbreak{} of\allowbreak{} d)\allowbreak{}\allowbreak{} \$\allowbreak{}\textbackslash{}\allowbreak{}Longrightarrow\$\allowbreak{}\allowbreak{} c)\allowbreak{}}{proof\allowbreak{} of\allowbreak{} d)\allowbreak{}\allowbreak{} \$\allowbreak{}\textbackslash{}\allowbreak{}Longrightarrow\$\allowbreak{}\allowbreak{} a)\allowbreak{}}{DER\_\allowbreak{}II-a:\allowbreak{}30}
\DERapparatusitem{ENV0089}{18;\allowbreak{}\allowbreak{} Footnote36,\allowbreak{}\allowbreak{} edition\allowbreak{} p.\allowbreak{}18}{if\allowbreak{} \$\allowbreak{}C\$\allowbreak{}\allowbreak{} is\allowbreak{} a\allowbreak{} groupoid}{if\allowbreak{} \$\allowbreak{}G\$\allowbreak{}\allowbreak{} is\allowbreak{} a\allowbreak{} group}{DER\_\allowbreak{}II-a:\allowbreak{}35}
\DERapparatusitem{ENV0090}{19;\allowbreak{}\allowbreak{} Proof\allowbreak{} of\allowbreak{} Corollary\allowbreak{} 2.\allowbreak{}14,\allowbreak{}\allowbreak{} edition\allowbreak{} p.\allowbreak{}19}{S}{T}{DER\_\allowbreak{}II-a:\allowbreak{}38}
\DERapparatusitem{ENV0091}{19;\allowbreak{}\allowbreak{} Cohomological\allowbreak{} vanishing\allowbreak{} statement,\allowbreak{}\allowbreak{} manuscript\allowbreak{} p.\allowbreak{}31}{}{\textbackslash{}\allowbreak{}mathcal\allowbreak{} }{DER\_\allowbreak{}II-a:\allowbreak{}39}
\DERapparatusitem{ENV0092}{19;\allowbreak{}\allowbreak{} Induction\allowbreak{} hypothesis,\allowbreak{}\allowbreak{} manuscript\allowbreak{} p.\allowbreak{}31}{every\allowbreak{} topos\allowbreak{} that\allowbreak{} is\allowbreak{} \textbackslash{}\allowbreak{}textit\{\allowbreak{}totally\}\allowbreak{}}{every\allowbreak{} topos\allowbreak{} that\allowbreak{} is\allowbreak{} \textbackslash{}\allowbreak{}textit\{\allowbreak{}locally\allowbreak{} totally\}\allowbreak{}}{DER\_\allowbreak{}II-a:\allowbreak{}40}
\DERapparatusitem{ENV0093}{19;\allowbreak{}\allowbreak{} Underbrace\allowbreak{} below\allowbreak{} H\textasciicircum{}\allowbreak{}i(\allowbreak{}U\textasciicircum{}\allowbreak{}*,\allowbreak{}F)\allowbreak{},\allowbreak{}\allowbreak{} manuscript\allowbreak{} p.\allowbreak{}31}{\textbackslash{}\allowbreak{}text\{\allowbreak{}for\allowbreak{} \}\allowbreak{}\allowbreak{} i\textbackslash{}\allowbreak{}\textbackslash{}\allowbreak{}\allowbreak{} \textbackslash{}\allowbreak{}text\{\allowbreak{}variable\}\allowbreak{}}{\textbackslash{}\allowbreak{}text\{\allowbreak{}for\allowbreak{} variable\allowbreak{} \textbackslash{}\allowbreak{}v\{\allowbreak{}C\}\allowbreak{}ech\allowbreak{} degree\}\allowbreak{}\textbackslash{}\allowbreak{}\textbackslash{}\allowbreak{}\allowbreak{} \textbackslash{}\allowbreak{}text\{\allowbreak{}at\allowbreak{} fixed\allowbreak{} \}\allowbreak{}i}{DER\_\allowbreak{}II-a:\allowbreak{}41}
\DERapparatusitem{ENV0094}{21;\allowbreak{}\allowbreak{} Corollary\allowbreak{} 2.\allowbreak{}17,\allowbreak{}\allowbreak{} edition\allowbreak{} p.\allowbreak{}21}{e\_\allowbreak{}Y}{Y\textasciicircum{}\allowbreak{}\textbackslash{}\allowbreak{}wedge}{DER\_\allowbreak{}II-a:\allowbreak{}42}
\DERapparatusitem{ENV0095}{21;\allowbreak{}\allowbreak{} Canonical\allowbreak{} factorization\allowbreak{} in\allowbreak{} Corollary\allowbreak{} 2.\allowbreak{}17}{p}{i}{DER\_\allowbreak{}II-a:\allowbreak{}43}
\DERapparatusitem{ENV0096}{21;\allowbreak{}\allowbreak{} After\allowbreak{} proof\allowbreak{} of\allowbreak{} Corollary\allowbreak{} 2.\allowbreak{}17,\allowbreak{}\allowbreak{} edition\allowbreak{} p.\allowbreak{}21}{examples\allowbreak{} therefore\allowbreak{} of\allowbreak{} categories}{examples\allowbreak{} of\allowbreak{} categories}{DER\_\allowbreak{}II-a:\allowbreak{}44}
\DERapparatusitem{ENV0097}{21;\allowbreak{}\allowbreak{} Proof\allowbreak{} of\allowbreak{} Corollary2.\allowbreak{}18,\allowbreak{}\allowbreak{} edition\allowbreak{} p.\allowbreak{}21}{the\allowbreak{} final\allowbreak{} object\allowbreak{} \$\allowbreak{}F\$\allowbreak{}\allowbreak{} of\allowbreak{} \$\allowbreak{}X/\allowbreak{}F\$\allowbreak{}}{the\allowbreak{} final\allowbreak{} object\allowbreak{} of\allowbreak{} \$\allowbreak{}(\allowbreak{}X/\allowbreak{}F)\allowbreak{}\textasciicircum{}\allowbreak{}\textbackslash{}\allowbreak{}wedge\$\allowbreak{},\allowbreak{}\allowbreak{} corresponding\allowbreak{} to\allowbreak{} \$\allowbreak{}F\$\allowbreak{}\allowbreak{} in\allowbreak{} \$\allowbreak{}X\textasciicircum{}\allowbreak{}\textbackslash{}\allowbreak{}wedge/\allowbreak{}F\$\allowbreak{},\allowbreak{}}{DER\_\allowbreak{}II-a:\allowbreak{}46}
\DERapparatusitem{ENV0098}{21;\allowbreak{}\allowbreak{} printed21;\allowbreak{}\allowbreak{} ms36;\allowbreak{}\allowbreak{} Question,\allowbreak{}\allowbreak{} equivalent\allowbreak{} cofinal\allowbreak{} formulation}{Y}{X}{DER\_\allowbreak{}II-a:\allowbreak{}47}
\DERapparatusitem{ENV0099}{22;\allowbreak{}\allowbreak{} printed22;\allowbreak{}\allowbreak{} ms36;\allowbreak{}\allowbreak{} proof2.\allowbreak{}19}{p}{d}{DER\_\allowbreak{}II-a:\allowbreak{}48}
\DERapparatusitem{ENV0100}{preamble;\allowbreak{}\allowbreak{} local\allowbreak{} macro\allowbreak{} definition}{\textbackslash{}\allowbreak{}text\{\allowbreak{}def\}\allowbreak{}}{\textbackslash{}\allowbreak{}mathrm\{\allowbreak{}def\}\allowbreak{}}{TYPESETTING:\allowbreak{}7;\allowbreak{}\allowbreak{} NOTATION\_\allowbreak{}DER:\allowbreak{}17}
\DERapparatusitem{ENV0101}{24;\allowbreak{}\allowbreak{} printed24;\allowbreak{}\allowbreak{} ms37;\allowbreak{}\allowbreak{} original\allowbreak{} reading\allowbreak{} in\allowbreak{} (\allowbreak{}3.\allowbreak{}2)\allowbreak{}}{\textbackslash{}\allowbreak{}uHom(\allowbreak{}I,\allowbreak{}(\allowbreak{}\textbackslash{}\allowbreak{}Ens)\allowbreak{})\allowbreak{}}{\textbackslash{}\allowbreak{}uHom(\allowbreak{}I\textasciicircum{}\allowbreak{}o,\allowbreak{}(\allowbreak{}\textbackslash{}\allowbreak{}Ens)\allowbreak{})\allowbreak{}}{DER\_\allowbreak{}II-a:\allowbreak{}50}
\DERapparatusitem{ENV0102}{25;\allowbreak{}\allowbreak{} printed25;\allowbreak{}\allowbreak{} ms38;\allowbreak{}\allowbreak{} properness\allowbreak{} row\allowbreak{} (\allowbreak{}3.\allowbreak{}8)\allowbreak{}}{(\allowbreak{}\$\allowbreak{}i'\allowbreak{} =\allowbreak{}\allowbreak{} u(\allowbreak{}i)\allowbreak{}\$\allowbreak{})\allowbreak{}}{(\allowbreak{}\$\allowbreak{}i\allowbreak{} =\allowbreak{}\allowbreak{} u(\allowbreak{}i')\allowbreak{}\$\allowbreak{})\allowbreak{}}{DER\_\allowbreak{}II-a:\allowbreak{}52}
\DERapparatusitem{ENV0103}{25;\allowbreak{}\allowbreak{} printed25;\allowbreak{}\allowbreak{} ms39;\allowbreak{}\allowbreak{} universal-equivalence\allowbreak{} row\allowbreak{} (\allowbreak{}3.\allowbreak{}8)\allowbreak{}}{\textbackslash{}\allowbreak{}text\{\allowbreak{}the\allowbreak{} \$\allowbreak{}W\_\allowbreak{}0\$\allowbreak{}-fibrations\allowbreak{} (\allowbreak{}or\allowbreak{} weak\allowbreak{} \$\allowbreak{}W\_\allowbreak{}0\$\allowbreak{}-fibrations\}\allowbreak{}}{\textbackslash{}\allowbreak{}text\{\allowbreak{}the\allowbreak{} \$\allowbreak{}W\_\allowbreak{}0\$\allowbreak{}-fibrations\allowbreak{} (\allowbreak{}or\allowbreak{} weak\allowbreak{} \$\allowbreak{}W\_\allowbreak{}0\$\allowbreak{}-fibrations)\allowbreak{}\}\allowbreak{}}{DER\_\allowbreak{}II-a:\allowbreak{}53}
\DERapparatusitem{ENV0104}{27;\allowbreak{}\allowbreak{} printed27;\allowbreak{}\allowbreak{} ms41;\allowbreak{}\allowbreak{} paragraph\allowbreak{} before3.\allowbreak{}1}{u(\allowbreak{}\textbackslash{}\allowbreak{}alpha)\allowbreak{}}{\textbackslash{}\allowbreak{}alpha}{DER\_\allowbreak{}II-a:\allowbreak{}56}
\DERapparatusitem{ENV0105}{27;\allowbreak{}\allowbreak{} printed27;\allowbreak{}\allowbreak{} ms41;\allowbreak{}\allowbreak{} Proposition3.\allowbreak{}1(\allowbreak{}a)\allowbreak{}}{I}{u}{DER\_\allowbreak{}II-a:\allowbreak{}57}
\DERapparatusitem{ENV0106}{27;\allowbreak{}\allowbreak{} printed27;\allowbreak{}\allowbreak{} ms42;\allowbreak{}\allowbreak{} paragraph\allowbreak{} after3.\allowbreak{}1}{The\allowbreak{} greatest\allowbreak{} difference\allowbreak{} between\allowbreak{} the\allowbreak{} notion\allowbreak{} of\allowbreak{} \$\allowbreak{}W\_\allowbreak{}0\$\allowbreak{}-fibration,\allowbreak{}\allowbreak{} and\allowbreak{} that\allowbreak{} of\allowbreak{} weak\allowbreak{} \$\allowbreak{}W\_\allowbreak{}0\$\allowbreak{}-fibration,\allowbreak{}\allowbreak{} seems\allowbreak{} to\allowbreak{} me\allowbreak{} that\allowbreak{} \$\allowbreak{}W\_\allowbreak{}0\$\allowbreak{}-perfect\allowbreak{} morphisms\allowbreak{} (\allowbreak{}i.\allowbreak{}e.\allowbreak{}\allowbreak{} \$\allowbreak{}W\_\allowbreak{}0\$\allowbreak{}-smooth\allowbreak{} and\allowbreak{} \$\allowbreak{}W\_\allowbreak{}0\$\allowbreak{}-proper)\allowbreak{},\allowbreak{}\allowbreak{} and\allowbreak{} a\allowbreak{} fortiori\allowbreak{} the\allowbreak{} \$\allowbreak{}W\$\allowbreak{}-perfect\allowbreak{} ones\allowbreak{} for\allowbreak{} a\allowbreak{} basic\allowbreak{} localizer\allowbreak{} \$\allowbreak{}W\allowbreak{} \textbackslash{}\allowbreak{}subseteq\allowbreak{} W\_\allowbreak{}0\$\allowbreak{},\allowbreak{}\allowbreak{} are\allowbreak{} necessarily\allowbreak{} weak\allowbreak{} \$\allowbreak{}W\$\allowbreak{}-fibrations,\allowbreak{}\allowbreak{} but\allowbreak{} (\allowbreak{}\textbackslash{}\allowbreak{}textit\{\allowbreak{}even\allowbreak{} if\allowbreak{} \$\allowbreak{}W\allowbreak{} =\allowbreak{}\allowbreak{} W\_\allowbreak{}\textbackslash{}\allowbreak{}omega\$\allowbreak{},\allowbreak{}\}\allowbreak{}\allowbreak{} the\allowbreak{} finest\allowbreak{} of\allowbreak{} all\allowbreak{} basic\allowbreak{} localizers)\allowbreak{}\allowbreak{} are\allowbreak{} not\allowbreak{} \$\allowbreak{}W\$\allowbreak{}-fibrations.\allowbreak{}}{The\allowbreak{} greatest\allowbreak{} difference\allowbreak{} between\allowbreak{} the\allowbreak{} notion\allowbreak{} of\allowbreak{} \$\allowbreak{}W\_\allowbreak{}0\$\allowbreak{}-fibration,\allowbreak{}\allowbreak{} and\allowbreak{} that\allowbreak{} of\allowbreak{} weak\allowbreak{} \$\allowbreak{}W\_\allowbreak{}0\$\allowbreak{}-fibration,\allowbreak{}\allowbreak{} seems\allowbreak{} to\allowbreak{} me\allowbreak{} that\allowbreak{} \$\allowbreak{}W\_\allowbreak{}0\$\allowbreak{}-perfect\allowbreak{} morphisms\allowbreak{} (\allowbreak{}i.\allowbreak{}e.\allowbreak{}\allowbreak{} \$\allowbreak{}W\_\allowbreak{}0\$\allowbreak{}-smooth\allowbreak{} and\allowbreak{} \$\allowbreak{}W\_\allowbreak{}0\$\allowbreak{}-proper)\allowbreak{}\allowbreak{} are\allowbreak{} necessarily\allowbreak{} weak\allowbreak{} \$\allowbreak{}W\_\allowbreak{}0\$\allowbreak{}-fibrations,\allowbreak{}\allowbreak{} but\allowbreak{} are\allowbreak{} not\allowbreak{} necessarily\allowbreak{} \$\allowbreak{}W\_\allowbreak{}0\$\allowbreak{}-fibrations.\allowbreak{}\allowbreak{} The\allowbreak{} \$\allowbreak{}W\$\allowbreak{}-perfect\allowbreak{} morphisms,\allowbreak{}\allowbreak{} for\allowbreak{} a\allowbreak{} basic\allowbreak{} localizer\allowbreak{} \$\allowbreak{}W\allowbreak{} \textbackslash{}\allowbreak{}subseteq\allowbreak{} W\_\allowbreak{}0\$\allowbreak{},\allowbreak{}\allowbreak{} are\allowbreak{} weak\allowbreak{} \$\allowbreak{}W\$\allowbreak{}-fibrations.\allowbreak{}}{DER\_\allowbreak{}II-a:\allowbreak{}58}
\DERapparatusitem{ENV0107}{28;\allowbreak{}\allowbreak{} printed28;\allowbreak{}\allowbreak{} ms44;\allowbreak{}\allowbreak{} footnote53}{U}{\textbackslash{}\allowbreak{}mathcal\{\allowbreak{}X\}\allowbreak{}}{DER\_\allowbreak{}II-a:\allowbreak{}60}
\DERapparatusitem{ENV0108}{28;\allowbreak{}\allowbreak{} printed28;\allowbreak{}\allowbreak{} ms44;\allowbreak{}\allowbreak{} footnote53}{f}{u}{DER\_\allowbreak{}II-a:\allowbreak{}61}
\DERapparatusitem{ENV0109}{28;\allowbreak{}\allowbreak{} printed28;\allowbreak{}\allowbreak{} ms44;\allowbreak{}\allowbreak{} footnote54,\allowbreak{}\allowbreak{} condition\allowbreak{} a\textasciicircum{}\allowbreak{}IV}{'\allowbreak{} \textbackslash{}\allowbreak{}longrightarrow\allowbreak{} I'}{\textbackslash{}\allowbreak{}longrightarrow\allowbreak{} I}{DER\_\allowbreak{}II-a:\allowbreak{}62}
\DERapparatusitem{ENV0110}{28;\allowbreak{}\allowbreak{} printed28;\allowbreak{}\allowbreak{} ms44;\allowbreak{}\allowbreak{} (\allowbreak{}3.\allowbreak{}2.\allowbreak{}3)\allowbreak{}\allowbreak{} right\allowbreak{} underbrace}{}{u\textasciicircum{}\allowbreak{}*}{DER\_\allowbreak{}II-a:\allowbreak{}63}
\DERapparatusitem{ENV0111}{preamble;\allowbreak{}\allowbreak{} local\allowbreak{} macro\allowbreak{} definition}{\textbackslash{}\allowbreak{}text\{\allowbreak{}def\}\allowbreak{}}{\textbackslash{}\allowbreak{}mathrm\{\allowbreak{}def\}\allowbreak{}}{TYPESETTING:\allowbreak{}7;\allowbreak{}\allowbreak{} NOTATION\_\allowbreak{}DER:\allowbreak{}17}
\DERapparatusitem{ENV0112}{29;\allowbreak{}\allowbreak{} ms45-46,\allowbreak{}\allowbreak{} 3.\allowbreak{}4(\allowbreak{}a)\allowbreak{},\allowbreak{}\allowbreak{} i.\allowbreak{}e.\allowbreak{}\allowbreak{} preceding\allowbreak{} F0,\allowbreak{}F1}{be\allowbreak{} nonempty,\allowbreak{}\allowbreak{} i.\allowbreak{}e.\allowbreak{}\allowbreak{} that\allowbreak{} we\allowbreak{} have}{be\allowbreak{} nonempty\allowbreak{} (\allowbreak{}condition\allowbreak{} F\allowbreak{} 1)\allowbreak{};\allowbreak{}\allowbreak{} F\allowbreak{} 0\allowbreak{} is\allowbreak{} also\allowbreak{} required.\allowbreak{}\allowbreak{} Thus\allowbreak{} we\allowbreak{} must\allowbreak{} have}{DER\_\allowbreak{}II-b:\allowbreak{}2}
\DERapparatusitem{ENV0113}{29;\allowbreak{}\allowbreak{} printed29;\allowbreak{}\allowbreak{} ms46;\allowbreak{}\allowbreak{} F0}{f}{u}{DER\_\allowbreak{}II-a:\allowbreak{}65;\allowbreak{}\allowbreak{} DER\_\allowbreak{}II-b:\allowbreak{}1}
\DERapparatusitem{ENV0114}{29;\allowbreak{}\allowbreak{} printed29;\allowbreak{}\allowbreak{} ms46;\allowbreak{}\allowbreak{} footnote57}{If\allowbreak{} \$\allowbreak{}I'\$\allowbreak{}\allowbreak{} is\allowbreak{} locally\allowbreak{} weakly\allowbreak{} filtering\allowbreak{} (\allowbreak{}=\allowbreak{}\allowbreak{} prefiltering)\allowbreak{}}{If\allowbreak{} \$\allowbreak{}I'\$\allowbreak{}\allowbreak{} satisfies\allowbreak{} right\allowbreak{} PS\allowbreak{} 1}{DER\_\allowbreak{}II-a:\allowbreak{}66}
\DERapparatusitem{ENV0115}{30;\allowbreak{}\allowbreak{} printed30;\allowbreak{}\allowbreak{} ms46;\allowbreak{}\allowbreak{} Corollary3.\allowbreak{}4(\allowbreak{}c)\allowbreak{},\allowbreak{}\allowbreak{} final\allowbreak{} assertion}{'}{}{DER\_\allowbreak{}II-a:\allowbreak{}67;\allowbreak{}\allowbreak{} DER\_\allowbreak{}II-b:\allowbreak{}3}
\DERapparatusitem{ENV0116}{30;\allowbreak{}\allowbreak{} printed30;\allowbreak{}\allowbreak{} ms46;\allowbreak{}\allowbreak{} Corollary3.\allowbreak{}4(\allowbreak{}d)\allowbreak{},\allowbreak{}\allowbreak{} hypothesis}{prefiltering\allowbreak{} (\allowbreak{}=\allowbreak{}\allowbreak{} locally\allowbreak{} weakly\allowbreak{} filtering,\allowbreak{}\allowbreak{} i.\allowbreak{}e.\allowbreak{}\allowbreak{} PS\allowbreak{} 1)\allowbreak{}}{right\allowbreak{} prefiltering\allowbreak{} (\allowbreak{}i.\allowbreak{}e.\allowbreak{}\allowbreak{} satisfying\allowbreak{} right\allowbreak{} PS\allowbreak{} 1)\allowbreak{}}{DER\_\allowbreak{}II-a:\allowbreak{}68}
\DERapparatusitem{ENV0117}{31;\allowbreak{}\allowbreak{} printed31;\allowbreak{}\allowbreak{} ms48;\allowbreak{}\allowbreak{} final\allowbreak{} PS2\allowbreak{} argument}{\textbackslash{}\allowbreak{}longrightarrow}{\textbackslash{}\allowbreak{}xrightarrow\{\allowbreak{}\textbackslash{}\allowbreak{}gamma\}\allowbreak{}}{DER\_\allowbreak{}II-a:\allowbreak{}69}
\DERapparatusitem{ENV0118}{33;\allowbreak{}\allowbreak{} printed33;\allowbreak{}\allowbreak{} ms51;\allowbreak{}\allowbreak{} slice\allowbreak{} restriction\allowbreak{} notation}{(\allowbreak{}u/\allowbreak{}i)\allowbreak{}}{\textbackslash{}\allowbreak{}alpha\_\allowbreak{}i}{DER\_\allowbreak{}II-a:\allowbreak{}70;\allowbreak{}\allowbreak{} DER\_\allowbreak{}II-b:\allowbreak{}4}
\DERapparatusitem{ENV0119}{33;\allowbreak{}\allowbreak{} printed33;\allowbreak{}\allowbreak{} ms52;\allowbreak{}\allowbreak{} Proposition4.\allowbreak{}1(\allowbreak{}a\allowbreak{} double-prime)\allowbreak{}}{}{F}{DER\_\allowbreak{}II-a:\allowbreak{}72;\allowbreak{}\allowbreak{} DER\_\allowbreak{}II-b:\allowbreak{}6}
\DERapparatusitem{ENV0120}{33;\allowbreak{}\allowbreak{} printed33;\allowbreak{}\allowbreak{} ms52;\allowbreak{}\allowbreak{} Proposition4.\allowbreak{}1(\allowbreak{}d)\allowbreak{}}{I}{i}{DER\_\allowbreak{}II-a:\allowbreak{}73;\allowbreak{}\allowbreak{} DER\_\allowbreak{}II-b:\allowbreak{}7}
\DERapparatusitem{ENV0121}{33;\allowbreak{}\allowbreak{} printed33;\allowbreak{}\allowbreak{} ms52;\allowbreak{}\allowbreak{} Proposition4.\allowbreak{}1(\allowbreak{}d\allowbreak{} prime)\allowbreak{}}{\$\allowbreak{}i\textbackslash{}\allowbreak{}bs\allowbreak{} I'\allowbreak{} \textbackslash{}\allowbreak{}longrightarrow\allowbreak{} i\textbackslash{}\allowbreak{}bs\allowbreak{} I\$\allowbreak{}\allowbreak{} is\allowbreak{} cofinal}{\$\allowbreak{}i\textbackslash{}\allowbreak{}bs\allowbreak{} I'\allowbreak{} \textbackslash{}\allowbreak{}longrightarrow\allowbreak{} i\textbackslash{}\allowbreak{}bs\allowbreak{} I\$\allowbreak{}\allowbreak{} is\allowbreak{} co-cofinal}{DER\_\allowbreak{}II-a:\allowbreak{}74;\allowbreak{}\allowbreak{} DER\_\allowbreak{}II-b:\allowbreak{}8}
\DERapparatusitem{ENV0122}{34;\allowbreak{}\allowbreak{} ms54,\allowbreak{}\allowbreak{} fibre\allowbreak{} of\allowbreak{} (\allowbreak{}4.\allowbreak{}10)\allowbreak{}}{u}{\textbackslash{}\allowbreak{}alpha}{DER\_\allowbreak{}II-b:\allowbreak{}11}
\DERapparatusitem{ENV0123}{38;\allowbreak{}\allowbreak{} ms61,\allowbreak{}\allowbreak{} definition\allowbreak{} of\allowbreak{} alpha\_\allowbreak{}i}{F}{I}{DER\_\allowbreak{}II-b:\allowbreak{}15}
\DERapparatusitem{ENV0124}{38;\allowbreak{}\allowbreak{} ms62,\allowbreak{}\allowbreak{} first\allowbreak{} Hom\allowbreak{} in\allowbreak{} (\allowbreak{}**)\allowbreak{}}{\textbackslash{}\allowbreak{}Hom(\allowbreak{}I/\allowbreak{}i,\allowbreak{}\allowbreak{} F\_\allowbreak{}i)\allowbreak{}}{\textbackslash{}\allowbreak{}Hom(\allowbreak{}\textbackslash{}\allowbreak{}xi\_\allowbreak{}\{\allowbreak{}I/\allowbreak{}i\}\allowbreak{},\allowbreak{}\allowbreak{} F\_\allowbreak{}i)\allowbreak{}}{DER\_\allowbreak{}II-b:\allowbreak{}16}
\DERapparatusitem{ENV0125}{38;\allowbreak{}\allowbreak{} ms62,\allowbreak{}\allowbreak{} constant\allowbreak{} coefficient\allowbreak{} in\allowbreak{} (\allowbreak{}**)\allowbreak{}}{\textbackslash{}\allowbreak{}xi'\_\allowbreak{}\{\allowbreak{}I'/\allowbreak{}i\}\allowbreak{}}{\textbackslash{}\allowbreak{}xi\_\allowbreak{}\{\allowbreak{}I'/\allowbreak{}i\}\allowbreak{}}{DER\_\allowbreak{}II-b:\allowbreak{}17}
\DERapparatusitem{ENV0126}{38;\allowbreak{}\allowbreak{} ms62,\allowbreak{}\allowbreak{} constant\allowbreak{} coefficient\allowbreak{} in\allowbreak{} (\allowbreak{}**)\allowbreak{}}{\textbackslash{}\allowbreak{}xi'\_\allowbreak{}\{\allowbreak{}I'/\allowbreak{}i\}\allowbreak{}}{\textbackslash{}\allowbreak{}xi\_\allowbreak{}\{\allowbreak{}I'/\allowbreak{}i\}\allowbreak{}}{DER\_\allowbreak{}II-b:\allowbreak{}17}
\DERapparatusitem{ENV0127}{38;\allowbreak{}\allowbreak{} ms62,\allowbreak{}\allowbreak{} counit\allowbreak{} (\allowbreak{}***)\allowbreak{}}{(\allowbreak{}u/\allowbreak{}i)\allowbreak{}\_\allowbreak{}!(\allowbreak{}u/\allowbreak{}i)\allowbreak{}\textasciicircum{}\allowbreak{}*(\allowbreak{}\textbackslash{}\allowbreak{}xi\_\allowbreak{}\{\allowbreak{}I'/\allowbreak{}i\}\allowbreak{})\allowbreak{}\allowbreak{} \textbackslash{}\allowbreak{}xrightarrow\{\allowbreak{}\textbackslash{}\allowbreak{}text\{\allowbreak{}adj.\allowbreak{}\}\allowbreak{}\}\allowbreak{}\allowbreak{} \textbackslash{}\allowbreak{}xi\_\allowbreak{}\{\allowbreak{}I'/\allowbreak{}i\}\allowbreak{}}{(\allowbreak{}u/\allowbreak{}i)\allowbreak{}\_\allowbreak{}!(\allowbreak{}u/\allowbreak{}i)\allowbreak{}\textasciicircum{}\allowbreak{}*(\allowbreak{}\textbackslash{}\allowbreak{}xi\_\allowbreak{}\{\allowbreak{}I/\allowbreak{}i\}\allowbreak{})\allowbreak{}\allowbreak{} \textbackslash{}\allowbreak{}xrightarrow\{\allowbreak{}\textbackslash{}\allowbreak{}text\{\allowbreak{}adj.\allowbreak{}\}\allowbreak{}\}\allowbreak{}\allowbreak{} \textbackslash{}\allowbreak{}xi\_\allowbreak{}\{\allowbreak{}I/\allowbreak{}i\}\allowbreak{}}{DER\_\allowbreak{}II-b:\allowbreak{}18}
\DERapparatusitem{ENV0128}{39;\allowbreak{}\allowbreak{} ms64,\allowbreak{}\allowbreak{} second\allowbreak{} restriction\allowbreak{} map\allowbreak{} of\allowbreak{} beta\_\allowbreak{}*G}{\textbackslash{}\allowbreak{}lambda}{\textbackslash{}\allowbreak{}mu}{DER\_\allowbreak{}II-b:\allowbreak{}19}
\DERapparatusitem{ENV0129}{39;\allowbreak{}\allowbreak{} ms63bis,\allowbreak{}\allowbreak{} paragraph\allowbreak{} after\allowbreak{} definition4.\allowbreak{}6}{one\allowbreak{} of\allowbreak{} the\allowbreak{} weaker\allowbreak{} conditions}{one\allowbreak{} of\allowbreak{} the\allowbreak{} stronger\allowbreak{} conditions}{DER\_\allowbreak{}II-b:\allowbreak{}21}
\DERapparatusitem{ENV0130}{40;\allowbreak{}\allowbreak{} ms64bis,\allowbreak{}\allowbreak{} question\allowbreak{} after\allowbreak{} proof4.\allowbreak{}7}{without\allowbreak{} being\allowbreak{} \$\allowbreak{}\textbackslash{}\allowbreak{}D\$\allowbreak{}-aspheric}{without\allowbreak{} being\allowbreak{} \$\allowbreak{}\textbackslash{}\allowbreak{}D\$\allowbreak{}-biaspheric}{DER\_\allowbreak{}II-b:\allowbreak{}22}
\DERapparatusitem{ENV0131}{40;\allowbreak{}\allowbreak{} ms65,\allowbreak{}\allowbreak{} footnote60}{except\allowbreak{} for\allowbreak{} the\allowbreak{} notion\allowbreak{} of\allowbreak{} \$\allowbreak{}\textbackslash{}\allowbreak{}D\$\allowbreak{}-connectedness}{including\allowbreak{} the\allowbreak{} notion\allowbreak{} of\allowbreak{} \$\allowbreak{}\textbackslash{}\allowbreak{}D\$\allowbreak{}-connectedness}{DER\_\allowbreak{}II-b:\allowbreak{}23}
\DERapparatusitem{ENV0132}{41;\allowbreak{}\allowbreak{} ms66,\allowbreak{}\allowbreak{} (\allowbreak{}4.\allowbreak{}10)\allowbreak{},\allowbreak{}\allowbreak{} first\allowbreak{} pullback\allowbreak{} target}{X}{Y}{DER\_\allowbreak{}II-b:\allowbreak{}24}
\DERapparatusitem{ENV0133}{41;\allowbreak{}\allowbreak{} running\allowbreak{} text\allowbreak{} or\allowbreak{} formula}{\textbackslash{}\allowbreak{}text\{\allowbreak{}def\}\allowbreak{}}{\textbackslash{}\allowbreak{}mathrm\{\allowbreak{}def\}\allowbreak{}}{TYPESETTING:\allowbreak{}7;\allowbreak{}\allowbreak{} NOTATION\_\allowbreak{}DER:\allowbreak{}17}
\DERapparatusitem{ENV0134}{42;\allowbreak{}\allowbreak{} running\allowbreak{} text\allowbreak{} or\allowbreak{} formula}{\textbackslash{}\allowbreak{}text\{\allowbreak{}def\}\allowbreak{}}{\textbackslash{}\allowbreak{}mathrm\{\allowbreak{}def\}\allowbreak{}}{TYPESETTING:\allowbreak{}7;\allowbreak{}\allowbreak{} NOTATION\_\allowbreak{}DER:\allowbreak{}17}
\DERapparatusitem{ENV0135}{43;\allowbreak{}\allowbreak{} running\allowbreak{} text\allowbreak{} or\allowbreak{} formula}{\textbackslash{}\allowbreak{}text\{\allowbreak{}def\}\allowbreak{}}{\textbackslash{}\allowbreak{}mathrm\{\allowbreak{}def\}\allowbreak{}}{TYPESETTING:\allowbreak{}7;\allowbreak{}\allowbreak{} NOTATION\_\allowbreak{}DER:\allowbreak{}17}
\DERapparatusitem{ENV0136}{43;\allowbreak{}\allowbreak{} ms69,\allowbreak{}\allowbreak{} definition\allowbreak{} of\allowbreak{} A\textasciicircum{}\allowbreak{}cup}{\textbackslash{}\allowbreak{}mathbf}{\textbackslash{}\allowbreak{}mathcal}{DER\_\allowbreak{}II-b:\allowbreak{}26}
\DERapparatusitem{ENV0137}{43;\allowbreak{}\allowbreak{} running\allowbreak{} text\allowbreak{} or\allowbreak{} formula}{\textbackslash{}\allowbreak{}text\{\allowbreak{}def\}\allowbreak{}}{\textbackslash{}\allowbreak{}mathrm\{\allowbreak{}def\}\allowbreak{}}{TYPESETTING:\allowbreak{}7;\allowbreak{}\allowbreak{} NOTATION\_\allowbreak{}DER:\allowbreak{}17}
\DERapparatusitem{ENV0138}{43;\allowbreak{}\allowbreak{} ms69,\allowbreak{}\allowbreak{} (\allowbreak{}D\_\allowbreak{}0)\allowbreak{},\allowbreak{}\allowbreak{} q\allowbreak{} in\allowbreak{} Q}{q\allowbreak{} \textbackslash{}\allowbreak{}in\allowbreak{} Q\allowbreak{} \textbackslash{}\allowbreak{}mid}{q\allowbreak{} \textbackslash{}\allowbreak{}in\allowbreak{} \textbackslash{}\allowbreak{}mathcal\{\allowbreak{}Q\}\allowbreak{}\allowbreak{} \textbackslash{}\allowbreak{}mid}{DER\_\allowbreak{}II-b:\allowbreak{}27}
\DERapparatusitem{ENV0139}{45;\allowbreak{}\allowbreak{} ms72,\allowbreak{}\allowbreak{} first\allowbreak{} paragraph}{subset\allowbreak{} of\allowbreak{} \$\allowbreak{}Q\$\allowbreak{}\allowbreak{} containing\allowbreak{} \$\allowbreak{}B\$\allowbreak{}}{subset\allowbreak{} of\allowbreak{} \$\allowbreak{}\textbackslash{}\allowbreak{}mathcal\allowbreak{} Q\$\allowbreak{}\allowbreak{} containing\allowbreak{} \$\allowbreak{}B\$\allowbreak{}}{DER\_\allowbreak{}II-b:\allowbreak{}31}
\DERapparatusitem{ENV0140}{45;\allowbreak{}\allowbreak{} running\allowbreak{} text\allowbreak{} or\allowbreak{} formula}{\textbackslash{}\allowbreak{}text\{\allowbreak{}def\}\allowbreak{}}{\textbackslash{}\allowbreak{}mathrm\{\allowbreak{}def\}\allowbreak{}}{TYPESETTING:\allowbreak{}7;\allowbreak{}\allowbreak{} NOTATION\_\allowbreak{}DER:\allowbreak{}17}
\DERapparatusitem{ENV0141}{45;\allowbreak{}\allowbreak{} running\allowbreak{} text\allowbreak{} or\allowbreak{} formula}{\textbackslash{}\allowbreak{}text\{\allowbreak{}def\}\allowbreak{}}{\textbackslash{}\allowbreak{}mathrm\{\allowbreak{}def\}\allowbreak{}}{TYPESETTING:\allowbreak{}7;\allowbreak{}\allowbreak{} NOTATION\_\allowbreak{}DER:\allowbreak{}17}
\DERapparatusitem{ENV0142}{47;\allowbreak{}\allowbreak{} ms75,\allowbreak{}\allowbreak{} question\allowbreak{} before\allowbreak{} ms76}{\$\allowbreak{}\textbackslash{}\allowbreak{}alpha\_\allowbreak{}0\allowbreak{} :\allowbreak{}\allowbreak{} \textbackslash{}\allowbreak{}cP\allowbreak{} \textbackslash{}\allowbreak{}longrightarrow\allowbreak{} \textbackslash{}\allowbreak{}cQ\$\allowbreak{},\allowbreak{}\allowbreak{} or\allowbreak{} \$\allowbreak{}\textbackslash{}\allowbreak{}beta\_\allowbreak{}0\allowbreak{} :\allowbreak{}\allowbreak{} \textbackslash{}\allowbreak{}cQ\allowbreak{} \textbackslash{}\allowbreak{}longrightarrow\allowbreak{} \textbackslash{}\allowbreak{}cP\$\allowbreak{}}{\$\allowbreak{}\textbackslash{}\allowbreak{}beta\_\allowbreak{}0\allowbreak{} :\allowbreak{}\allowbreak{} \textbackslash{}\allowbreak{}cP\allowbreak{} \textbackslash{}\allowbreak{}longrightarrow\allowbreak{} \textbackslash{}\allowbreak{}cQ\$\allowbreak{},\allowbreak{}\allowbreak{} or\allowbreak{} \$\allowbreak{}\textbackslash{}\allowbreak{}alpha\_\allowbreak{}0\allowbreak{} :\allowbreak{}\allowbreak{} \textbackslash{}\allowbreak{}cQ\allowbreak{} \textbackslash{}\allowbreak{}longrightarrow\allowbreak{} \textbackslash{}\allowbreak{}cP\$\allowbreak{}}{DER\_\allowbreak{}II-b:\allowbreak{}34}
\DERapparatusitem{ENV0143}{47;\allowbreak{}\allowbreak{} ms77,\allowbreak{}\allowbreak{} (\allowbreak{}5.\allowbreak{}20)\allowbreak{},\allowbreak{}\allowbreak{} second\allowbreak{} equality}{A\textasciicircum{}\allowbreak{}\{\allowbreak{}\textbackslash{}\allowbreak{}natural\}\allowbreak{}\allowbreak{} \&\allowbreak{}\allowbreak{} =\allowbreak{}\allowbreak{} \&\allowbreak{}\allowbreak{} \textbackslash{}\allowbreak{}sup\textbackslash{}\allowbreak{}nolimits\_\allowbreak{}\{\allowbreak{}p\allowbreak{} \textbackslash{}\allowbreak{}in\allowbreak{} A\}\allowbreak{}\allowbreak{} p\textasciicircum{}\allowbreak{}\{\allowbreak{}\textbackslash{}\allowbreak{}natural\}\allowbreak{},\allowbreak{}}{A\textasciicircum{}\allowbreak{}\{\allowbreak{}\textbackslash{}\allowbreak{}natural\}\allowbreak{}\allowbreak{} \&\allowbreak{}\allowbreak{} =\allowbreak{}\allowbreak{} \&\allowbreak{}\allowbreak{} \textbackslash{}\allowbreak{}\{\allowbreak{}p\textbackslash{}\allowbreak{}in\textbackslash{}\allowbreak{}cP\textbackslash{}\allowbreak{}mid\allowbreak{} p\textbackslash{}\allowbreak{}leq\allowbreak{} (\allowbreak{}\textbackslash{}\allowbreak{}sup\textbackslash{}\allowbreak{}nolimits\_\allowbreak{}\{\allowbreak{}a\textbackslash{}\allowbreak{}in\allowbreak{} A\}\allowbreak{}a)\allowbreak{}\textasciicircum{}\allowbreak{}\{\allowbreak{}\textbackslash{}\allowbreak{}natural\}\allowbreak{}\textbackslash{}\allowbreak{}\}\allowbreak{}\allowbreak{} =\allowbreak{}\allowbreak{} \textbackslash{}\allowbreak{}\{\allowbreak{}p\textbackslash{}\allowbreak{}in\textbackslash{}\allowbreak{}cP\textbackslash{}\allowbreak{}mid\allowbreak{} p\textbackslash{}\allowbreak{}leq\allowbreak{} (\allowbreak{}\textbackslash{}\allowbreak{}sup\textbackslash{}\allowbreak{}nolimits\_\allowbreak{}\{\allowbreak{}a\textbackslash{}\allowbreak{}in\allowbreak{} A\}\allowbreak{}a\textasciicircum{}\allowbreak{}\{\allowbreak{}\textbackslash{}\allowbreak{}natural\}\allowbreak{})\allowbreak{}\textasciicircum{}\allowbreak{}\{\allowbreak{}\textbackslash{}\allowbreak{}natural\}\allowbreak{}\textbackslash{}\allowbreak{}\}\allowbreak{},\allowbreak{}}{DER\_\allowbreak{}II-b:\allowbreak{}35}
\DERapparatusitem{ENV0144}{48;\allowbreak{}\allowbreak{} ms79,\allowbreak{}\allowbreak{} absolute\allowbreak{} bipolar\allowbreak{} closed\allowbreak{} subsets}{the\allowbreak{} convex\allowbreak{} circled}{the\allowbreak{} nonempty\allowbreak{} convex\allowbreak{} circled}{DER\_\allowbreak{}II-b:\allowbreak{}37}
\DERapparatusitem{ENV0145}{49;\allowbreak{}\allowbreak{} ms79,\allowbreak{}\allowbreak{} Galois\allowbreak{} example,\allowbreak{}\allowbreak{} closed\allowbreak{} group\allowbreak{} subsets}{the\allowbreak{} subgroups\allowbreak{} of\allowbreak{} \$\allowbreak{}G\$\allowbreak{}.\allowbreak{}}{the\allowbreak{} subgroups\allowbreak{} of\allowbreak{} \$\allowbreak{}G\$\allowbreak{}\allowbreak{} closed\allowbreak{} for\allowbreak{} the\allowbreak{} Krull\allowbreak{} topology.\allowbreak{}}{DER\_\allowbreak{}II-b:\allowbreak{}38}
\DERapparatusitem{ENV0146}{49;\allowbreak{}\allowbreak{} ms80,\allowbreak{}\allowbreak{} lifting\allowbreak{} property\allowbreak{} parenthesis}{with\allowbreak{} respect\allowbreak{} to\allowbreak{} \$\allowbreak{}\textbackslash{}\allowbreak{}Psi\$\allowbreak{},\allowbreak{}\allowbreak{} \$\allowbreak{}\textbackslash{}\allowbreak{}Phi\$\allowbreak{}\allowbreak{} the\allowbreak{} RLP}{with\allowbreak{} respect\allowbreak{} to\allowbreak{} \$\allowbreak{}\textbackslash{}\allowbreak{}Phi\$\allowbreak{},\allowbreak{}\allowbreak{} \$\allowbreak{}\textbackslash{}\allowbreak{}Phi\$\allowbreak{}\allowbreak{} the\allowbreak{} RLP}{DER\_\allowbreak{}II-b:\allowbreak{}39}
\DERapparatusitem{ENV0147}{50;\allowbreak{}\allowbreak{} ms82,\allowbreak{}\allowbreak{} inverse\allowbreak{} image\allowbreak{} of\allowbreak{} classes\allowbreak{} by\allowbreak{} phi}{\textbackslash{}\allowbreak{}Phi\textbackslash{}\allowbreak{}times\textbackslash{}\allowbreak{}Psi}{\textbackslash{}\allowbreak{}Psi\textbackslash{}\allowbreak{}times\textbackslash{}\allowbreak{}Phi}{DER\_\allowbreak{}II-b:\allowbreak{}40}
\DERapparatusitem{ENV0148}{50;\allowbreak{}\allowbreak{} ms82,\allowbreak{}\allowbreak{} condition\allowbreak{} (\allowbreak{}**)\allowbreak{},\allowbreak{}\allowbreak{} dotted\allowbreak{} arrow}{\textbackslash{}\allowbreak{}varphi}{f}{DER\_\allowbreak{}II-b:\allowbreak{}41}
\DERapparatusitem{ENV0149}{51;\allowbreak{}\allowbreak{} ms83,\allowbreak{}\allowbreak{} (\allowbreak{}6.\allowbreak{}3)\allowbreak{},\allowbreak{}\allowbreak{} explanatory\allowbreak{} parenthesis}{\$\allowbreak{}A\$\allowbreak{}\allowbreak{} is\allowbreak{} stable\allowbreak{} under\allowbreak{} base\allowbreak{} change\allowbreak{} by\allowbreak{} \$\allowbreak{}p\textbackslash{}\allowbreak{}in\allowbreak{} L\$\allowbreak{}}{the\allowbreak{} base\allowbreak{} changes\allowbreak{} of\allowbreak{} arrows\allowbreak{} of\allowbreak{} \$\allowbreak{}A\$\allowbreak{}\allowbreak{} by\allowbreak{} \$\allowbreak{}p\textbackslash{}\allowbreak{}in\allowbreak{} L\$\allowbreak{}\allowbreak{} belong\allowbreak{} to\allowbreak{} \$\allowbreak{}W\$\allowbreak{}}{DER\_\allowbreak{}II-b:\allowbreak{}43}
\DERapparatusitem{ENV0150}{51;\allowbreak{}\allowbreak{} ms84,\allowbreak{}\allowbreak{} lower\allowbreak{} set\allowbreak{} of\allowbreak{} transposed\allowbreak{} relation}{Q}{P}{DER\_\allowbreak{}II-b:\allowbreak{}46}
\DERapparatusitem{ENV0151}{52;\allowbreak{}\allowbreak{} ms84,\allowbreak{}\allowbreak{} pullback\allowbreak{} on\allowbreak{} slices\allowbreak{} in\allowbreak{} general\allowbreak{} M}{\textbackslash{}\allowbreak{}mathrm\{\allowbreak{}Cat\}\allowbreak{}}{\textbackslash{}\allowbreak{}mathcal\allowbreak{} M}{DER\_\allowbreak{}II-b:\allowbreak{}47}
\DERapparatusitem{ENV0152}{52;\allowbreak{}\allowbreak{} ms84,\allowbreak{}\allowbreak{} pullback\allowbreak{} on\allowbreak{} slices\allowbreak{} in\allowbreak{} general\allowbreak{} M}{\textbackslash{}\allowbreak{}mathrm\{\allowbreak{}Cat\}\allowbreak{}}{\textbackslash{}\allowbreak{}mathcal\allowbreak{} M}{DER\_\allowbreak{}II-b:\allowbreak{}47}
\DERapparatusitem{ENV0153}{52;\allowbreak{}\allowbreak{} ms85,\allowbreak{}\allowbreak{} transporter\allowbreak{} definition\allowbreak{} slice\allowbreak{} categories}{\textbackslash{}\allowbreak{}mathrm\{\allowbreak{}Cat\}\allowbreak{}}{\textbackslash{}\allowbreak{}mathcal\allowbreak{} M}{DER\_\allowbreak{}II-b:\allowbreak{}48}
\DERapparatusitem{ENV0154}{52;\allowbreak{}\allowbreak{} ms85,\allowbreak{}\allowbreak{} transporter\allowbreak{} definition\allowbreak{} slice\allowbreak{} categories}{\textbackslash{}\allowbreak{}mathrm\{\allowbreak{}Cat\}\allowbreak{}}{\textbackslash{}\allowbreak{}mathcal\allowbreak{} M}{DER\_\allowbreak{}II-b:\allowbreak{}48}
\DERapparatusitem{ENV0155}{54;\allowbreak{}\allowbreak{} ms86,\allowbreak{}\allowbreak{} (\allowbreak{}6.\allowbreak{}8)\allowbreak{},\allowbreak{}\allowbreak{} transporter\allowbreak{} notation}{}{\_\allowbreak{}W}{DER\_\allowbreak{}II-b:\allowbreak{}49}
\DERapparatusitem{ENV0156}{54;\allowbreak{}\allowbreak{} running\allowbreak{} text\allowbreak{} or\allowbreak{} formula}{\textbackslash{}\allowbreak{}text\{\allowbreak{}def\}\allowbreak{}}{\textbackslash{}\allowbreak{}mathrm\{\allowbreak{}def\}\allowbreak{}}{TYPESETTING:\allowbreak{}7;\allowbreak{}\allowbreak{} NOTATION\_\allowbreak{}DER:\allowbreak{}17}
\DERapparatusitem{ENV0157}{55;\allowbreak{}\allowbreak{} ms89,\allowbreak{}\allowbreak{} final\allowbreak{} parenthesis}{\textbackslash{}\allowbreak{}subseteq}{\textbackslash{}\allowbreak{}in}{DER\_\allowbreak{}II-b:\allowbreak{}53}
\DERapparatusitem{ENV0158}{56;\allowbreak{}\allowbreak{} ms90\allowbreak{} Remark7.\allowbreak{}2}{The\allowbreak{} \textbackslash{}\allowbreak{}textit\{\allowbreak{}étale\}\allowbreak{}\allowbreak{} morphisms\allowbreak{} are\allowbreak{} locally\allowbreak{} filtering.\allowbreak{}}{The\allowbreak{} \textbackslash{}\allowbreak{}textit\{\allowbreak{}étale\}\allowbreak{}\allowbreak{} morphisms\allowbreak{} are\allowbreak{} locally\allowbreak{} cofiltering.\allowbreak{}}{DER\_\allowbreak{}II-b:\allowbreak{}56}
\DERapparatusitem{ENV0159}{57;\allowbreak{}\allowbreak{} ms92\allowbreak{} Remark7.\allowbreak{}4(\allowbreak{}3)\allowbreak{}}{In\allowbreak{} particular,\allowbreak{}\allowbreak{} it\allowbreak{} is\allowbreak{} \$\allowbreak{}W\_\allowbreak{}0\$\allowbreak{}-aspheric}{In\allowbreak{} particular,\allowbreak{}\allowbreak{} it\allowbreak{} is\allowbreak{} \$\allowbreak{}W\_\allowbreak{}0\$\allowbreak{}-coaspheric}{DER\_\allowbreak{}II-b:\allowbreak{}57}
\DERapparatusitem{ENV0160}{57;\allowbreak{}\allowbreak{} ms92\allowbreak{} Proposition7.\allowbreak{}4(\allowbreak{}b)\allowbreak{}}{f/\allowbreak{}x\allowbreak{} :\allowbreak{}\allowbreak{} X/\allowbreak{}x\allowbreak{} \textbackslash{}\allowbreak{}longrightarrow\allowbreak{} Y/\allowbreak{}x}{f/\allowbreak{}x\allowbreak{} :\allowbreak{}\allowbreak{} X/\allowbreak{}x\allowbreak{} \textbackslash{}\allowbreak{}longrightarrow\allowbreak{} Y/\allowbreak{}y}{DER\_\allowbreak{}II-b:\allowbreak{}58}
\DERapparatusitem{ENV0161}{57;\allowbreak{}\allowbreak{} ms92\allowbreak{} Proposition7.\allowbreak{}4(\allowbreak{}c)\allowbreak{},\allowbreak{}\allowbreak{} definition\allowbreak{} of\allowbreak{} V}{f/\allowbreak{}x}{(\allowbreak{}f/\allowbreak{}x)\allowbreak{}|\_\allowbreak{}U}{DER\_\allowbreak{}II-b:\allowbreak{}59}
\DERapparatusitem{ENV0162}{57;\allowbreak{}\allowbreak{} ms92\allowbreak{} Proposition7.\allowbreak{}4(\allowbreak{}c)\allowbreak{},\allowbreak{}\allowbreak{} definition\allowbreak{} of\allowbreak{} V}{\textbackslash{}\allowbreak{}backslash(\allowbreak{}X/\allowbreak{}x)\allowbreak{}}{\textbackslash{}\allowbreak{}backslash\allowbreak{} U}{DER\_\allowbreak{}II-b:\allowbreak{}59}
\DERapparatusitem{ENV0163}{57;\allowbreak{}\allowbreak{} ms92-93\allowbreak{} Proposition7.\allowbreak{}4(\allowbreak{}c)\allowbreak{},\allowbreak{}\allowbreak{} letter\allowbreak{} consistency}{q}{g}{DER\_\allowbreak{}II-b:\allowbreak{}61}
\DERapparatusitem{ENV0164}{58;\allowbreak{}\allowbreak{} running\allowbreak{} text\allowbreak{} or\allowbreak{} formula}{\textbackslash{}\allowbreak{}text\{\allowbreak{}def\}\allowbreak{}}{\textbackslash{}\allowbreak{}mathrm\{\allowbreak{}def\}\allowbreak{}}{TYPESETTING:\allowbreak{}7;\allowbreak{}\allowbreak{} NOTATION\_\allowbreak{}DER:\allowbreak{}17}
\DERapparatusitem{ENV0165}{58;\allowbreak{}\allowbreak{} ms94\allowbreak{} second\allowbreak{} diagram}{y'\textbackslash{}\allowbreak{}backslash\allowbreak{} X\allowbreak{} \textbackslash{}\allowbreak{}arrow[\allowbreak{}r]\allowbreak{}\allowbreak{} \textbackslash{}\allowbreak{}arrow[\allowbreak{}d,\allowbreak{}\allowbreak{} "\textbackslash{}\allowbreak{}text\{\allowbreak{}ét.\allowbreak{}\}\allowbreak{}"]\allowbreak{}\allowbreak{} \textbackslash{}\allowbreak{}arrow[\allowbreak{}dr,\allowbreak{}\allowbreak{} phantom,\allowbreak{}\allowbreak{} "\textbackslash{}\allowbreak{}textit\{\allowbreak{}cart.\allowbreak{}\}\allowbreak{}"]\allowbreak{}\allowbreak{} \&\allowbreak{}\allowbreak{} y\textbackslash{}\allowbreak{}backslash\allowbreak{} X\allowbreak{} \textbackslash{}\allowbreak{}arrow[\allowbreak{}r]\allowbreak{}\allowbreak{} \textbackslash{}\allowbreak{}arrow[\allowbreak{}d,\allowbreak{}\allowbreak{} "\textbackslash{}\allowbreak{}text\{\allowbreak{}ét.\allowbreak{}\}\allowbreak{}"]\allowbreak{}\allowbreak{} \textbackslash{}\allowbreak{}arrow[\allowbreak{}dr,\allowbreak{}\allowbreak{} phantom,\allowbreak{}\allowbreak{} "\textbackslash{}\allowbreak{}textit\{\allowbreak{}cart.\allowbreak{}\}\allowbreak{}"]\allowbreak{}\allowbreak{} \&\allowbreak{}\allowbreak{} X\allowbreak{} \textbackslash{}\allowbreak{}arrow[\allowbreak{}d,\allowbreak{}\allowbreak{} "\textbackslash{}\allowbreak{}text\{\allowbreak{}ét.\allowbreak{}\}\allowbreak{}"]\allowbreak{}}{y'\textbackslash{}\allowbreak{}backslash\allowbreak{} X\allowbreak{} \textbackslash{}\allowbreak{}arrow[\allowbreak{}r]\allowbreak{}\allowbreak{} \textbackslash{}\allowbreak{}arrow[\allowbreak{}d]\allowbreak{}\allowbreak{} \textbackslash{}\allowbreak{}arrow[\allowbreak{}dr,\allowbreak{}\allowbreak{} phantom,\allowbreak{}\allowbreak{} "\textbackslash{}\allowbreak{}textit\{\allowbreak{}cart.\allowbreak{}\}\allowbreak{}"]\allowbreak{}\allowbreak{} \&\allowbreak{}\allowbreak{} y\textbackslash{}\allowbreak{}backslash\allowbreak{} X\allowbreak{} \textbackslash{}\allowbreak{}arrow[\allowbreak{}r]\allowbreak{}\allowbreak{} \textbackslash{}\allowbreak{}arrow[\allowbreak{}d]\allowbreak{}\allowbreak{} \textbackslash{}\allowbreak{}arrow[\allowbreak{}dr,\allowbreak{}\allowbreak{} phantom,\allowbreak{}\allowbreak{} "\textbackslash{}\allowbreak{}textit\{\allowbreak{}cart.\allowbreak{}\}\allowbreak{}"]\allowbreak{}\allowbreak{} \&\allowbreak{}\allowbreak{} X\allowbreak{} \textbackslash{}\allowbreak{}arrow[\allowbreak{}d]\allowbreak{}}{DER\_\allowbreak{}II-b:\allowbreak{}64}
\DERapparatusitem{ENV0166}{59;\allowbreak{}\allowbreak{} ms94\allowbreak{} continuation\allowbreak{} after\allowbreak{} second\allowbreak{} diagram}{}{'}{DER\_\allowbreak{}II-b:\allowbreak{}65}
\DERapparatusitem{ENV0167}{59;\allowbreak{}\allowbreak{} ms94\allowbreak{} end\allowbreak{} of\allowbreak{} proof\allowbreak{} of(\allowbreak{}a)\allowbreak{}}{This\allowbreak{} proves\allowbreak{} d)\allowbreak{}.\allowbreak{}\space{}\allowbreak{}\space{}\allowbreak{}Applying\allowbreak{} a)\allowbreak{}}{This\allowbreak{} proves\allowbreak{} a)\allowbreak{}.\allowbreak{}\allowbreak{} Applying\allowbreak{} a)\allowbreak{}}{DER\_\allowbreak{}II-b:\allowbreak{}66}
\DERapparatusitem{ENV0168}{59;\allowbreak{}\allowbreak{} ms95\allowbreak{} Remark7.\allowbreak{}5(\allowbreak{}1)\allowbreak{},\allowbreak{}\allowbreak{} definition\allowbreak{} of\allowbreak{} V}{X}{U}{DER\_\allowbreak{}II-b:\allowbreak{}67}
\DERapparatusitem{ENV0169}{60;\allowbreak{}\allowbreak{} ms97\allowbreak{} locally\allowbreak{} cofiltered\allowbreak{} node}{(\allowbreak{}X\allowbreak{} \textbackslash{}\allowbreak{}textit\{\allowbreak{}locally\allowbreak{} cofiltering\}\allowbreak{})\allowbreak{}}{(\allowbreak{}\$\allowbreak{}X\$\allowbreak{}\allowbreak{} \textbackslash{}\allowbreak{}textit\{\allowbreak{}locally\allowbreak{} cofiltering\}\allowbreak{})\allowbreak{}}{DER\_\allowbreak{}II-b:\allowbreak{}69}
\DERapparatusitem{ENV0170}{61;\allowbreak{}\allowbreak{} ms98\allowbreak{} abbreviated\allowbreak{} total-connectedness\allowbreak{} nodes}{\$\allowbreak{}X\textasciicircum{}\allowbreak{}o\allowbreak{} \textbackslash{}\allowbreak{}longrightarrow\allowbreak{} X\textasciicircum{}\allowbreak{}o\allowbreak{} \textbackslash{}\allowbreak{}times\allowbreak{} X\textasciicircum{}\allowbreak{}o\$\allowbreak{}\allowbreak{} cofinal}{\$\allowbreak{}X\textbackslash{}\allowbreak{}neq\textbackslash{}\allowbreak{}emptyset\$\allowbreak{}\allowbreak{} and\allowbreak{} \$\allowbreak{}X\textasciicircum{}\allowbreak{}o\allowbreak{} \textbackslash{}\allowbreak{}longrightarrow\allowbreak{} X\textasciicircum{}\allowbreak{}o\allowbreak{} \textbackslash{}\allowbreak{}times\allowbreak{} X\textasciicircum{}\allowbreak{}o\$\allowbreak{}\allowbreak{} cofinal}{DER\_\allowbreak{}II-b:\allowbreak{}72}
\DERapparatusitem{ENV0171}{61;\allowbreak{}\allowbreak{} ms98\allowbreak{} abbreviated\allowbreak{} total-connectedness\allowbreak{} nodes}{\$\allowbreak{}X\textasciicircum{}\allowbreak{}o\allowbreak{} \textbackslash{}\allowbreak{}longrightarrow\allowbreak{} X\textasciicircum{}\allowbreak{}o\allowbreak{} \textbackslash{}\allowbreak{}times\allowbreak{} X\textasciicircum{}\allowbreak{}o\$\allowbreak{}\allowbreak{} cofinal}{\$\allowbreak{}X\textbackslash{}\allowbreak{}neq\textbackslash{}\allowbreak{}emptyset\$\allowbreak{}\allowbreak{} and\allowbreak{} \$\allowbreak{}X\textasciicircum{}\allowbreak{}o\allowbreak{} \textbackslash{}\allowbreak{}longrightarrow\allowbreak{} X\textasciicircum{}\allowbreak{}o\allowbreak{} \textbackslash{}\allowbreak{}times\allowbreak{} X\textasciicircum{}\allowbreak{}o\$\allowbreak{}\allowbreak{} cofinal}{DER\_\allowbreak{}II-b:\allowbreak{}72}
\DERapparatusitem{ENV0172}{62;\allowbreak{}\allowbreak{} ms98\allowbreak{} membership\allowbreak{} table}{Likewise\allowbreak{} for\allowbreak{} \&\allowbreak{}}{The\allowbreak{} same\allowbreak{} exclusion\allowbreak{} from\allowbreak{} the\allowbreak{} four\allowbreak{} connected\allowbreak{} classes\allowbreak{} for\allowbreak{} \&\allowbreak{}}{DER\_\allowbreak{}II-b:\allowbreak{}74}
\DERapparatusitem{ENV0173}{62;\allowbreak{}\allowbreak{} ms99\allowbreak{} first\allowbreak{} diagram,\allowbreak{}\allowbreak{} PS*1\allowbreak{} arrow\allowbreak{} label}{(\allowbreak{}the\allowbreak{} \$\allowbreak{}x\textbackslash{}\allowbreak{}backslash\allowbreak{} X\$\allowbreak{}\allowbreak{} satisfy}{(\allowbreak{}the\allowbreak{} \$\allowbreak{}X/\allowbreak{}x\$\allowbreak{}\allowbreak{} satisfy}{DER\_\allowbreak{}II-b:\allowbreak{}77}
\DERapparatusitem{ENV0174}{62;\allowbreak{}\allowbreak{} ms98\allowbreak{} abbreviated\allowbreak{} total-connectedness\allowbreak{} nodes}{\$\allowbreak{}X\textasciicircum{}\allowbreak{}o\allowbreak{} \textbackslash{}\allowbreak{}longrightarrow\allowbreak{} X\textasciicircum{}\allowbreak{}o\allowbreak{} \textbackslash{}\allowbreak{}times\allowbreak{} X\textasciicircum{}\allowbreak{}o\$\allowbreak{}\allowbreak{} cofinal}{\$\allowbreak{}X\textbackslash{}\allowbreak{}neq\textbackslash{}\allowbreak{}emptyset\$\allowbreak{}\allowbreak{} and\allowbreak{} \$\allowbreak{}X\textasciicircum{}\allowbreak{}o\allowbreak{} \textbackslash{}\allowbreak{}longrightarrow\allowbreak{} X\textasciicircum{}\allowbreak{}o\allowbreak{} \textbackslash{}\allowbreak{}times\allowbreak{} X\textasciicircum{}\allowbreak{}o\$\allowbreak{}\allowbreak{} cofinal}{DER\_\allowbreak{}II-b:\allowbreak{}72;\allowbreak{}\allowbreak{} DER\_\allowbreak{}II-b:\allowbreak{}80}
\DERapparatusitem{ENV0175}{62;\allowbreak{}\allowbreak{} ms99\allowbreak{} first\allowbreak{} diagram,\allowbreak{}\allowbreak{} local\allowbreak{} cofiltered\allowbreak{} parenthesis}{connected\allowbreak{} components\allowbreak{} filtering}{connected\allowbreak{} components\allowbreak{} cofiltering}{DER\_\allowbreak{}II-b:\allowbreak{}76}
\DERapparatusitem{ENV0176}{62;\allowbreak{}\allowbreak{} ms99\allowbreak{} first\allowbreak{} diagram,\allowbreak{}\allowbreak{} local\allowbreak{} irreducibility\allowbreak{} parenthesis}{pseudo-filtering)\allowbreak{}\textbackslash{}\allowbreak{}\textbackslash{}\allowbreak{}(\allowbreak{}\$\allowbreak{}\textbackslash{}\allowbreak{}PS\textasciicircum{}\allowbreak{}*1\$\allowbreak{})\allowbreak{}}{pseudo-cofiltering)\allowbreak{}\textbackslash{}\allowbreak{}\textbackslash{}\allowbreak{}(\allowbreak{}\$\allowbreak{}\textbackslash{}\allowbreak{}PS\textasciicircum{}\allowbreak{}*1\$\allowbreak{})\allowbreak{}}{DER\_\allowbreak{}II-b:\allowbreak{}78}
\DERapparatusitem{ENV0177}{62;\allowbreak{}\allowbreak{} ms99\allowbreak{} dual\allowbreak{} diagram\allowbreak{} lower-right\allowbreak{} node}{\textbackslash{}\allowbreak{}text\{\allowbreak{}\$\allowbreak{}X\$\allowbreak{}\allowbreak{} pseudo-cofiltering\}\allowbreak{}}{\textbackslash{}\allowbreak{}text\{\allowbreak{}\$\allowbreak{}X\$\allowbreak{}\allowbreak{} pseudo-filtering\}\allowbreak{}}{DER\_\allowbreak{}II-b:\allowbreak{}79}
\subsection*{Chapter III}
\DERapparatusitem{ENV0178}{1;\allowbreak{}\allowbreak{} III\allowbreak{} ms1,\allowbreak{}\allowbreak{} equation\allowbreak{} (\allowbreak{}2)\allowbreak{}\allowbreak{} and\allowbreak{} following\allowbreak{} paragraph}{\textbackslash{}\allowbreak{}mathcal\{\allowbreak{}A\}\allowbreak{},\allowbreak{}\allowbreak{} \textbackslash{}\allowbreak{}mathcal\{\allowbreak{}A\}\allowbreak{}}{\textbackslash{}\allowbreak{}D,\allowbreak{}\allowbreak{} \textbackslash{}\allowbreak{}D}{DER\_\allowbreak{}III-V:\allowbreak{}1}
\DERapparatusitem{ENV0179}{1;\allowbreak{}\allowbreak{} III\allowbreak{} ms1,\allowbreak{}\allowbreak{} equation\allowbreak{} (\allowbreak{}2)\allowbreak{}\allowbreak{} and\allowbreak{} following\allowbreak{} paragraph}{\textbackslash{}\allowbreak{}uHom}{\textbackslash{}\allowbreak{}hm}{DER\_\allowbreak{}III-V:\allowbreak{}1}
\DERapparatusitem{ENV0180}{2;\allowbreak{}\allowbreak{} III\allowbreak{} ms2,\allowbreak{}\allowbreak{} introduction\allowbreak{} to\allowbreak{} equation\allowbreak{} (\allowbreak{}3)\allowbreak{}}{F}{M}{DER\_\allowbreak{}III-V:\allowbreak{}2}
\DERapparatusitem{ENV0181}{2;\allowbreak{}\allowbreak{} III\allowbreak{} ms3,\allowbreak{}\allowbreak{} last\allowbreak{} line}{F}{M}{DER\_\allowbreak{}III-V:\allowbreak{}3}
\DERapparatusitem{ENV0182}{3;\allowbreak{}\allowbreak{} III\allowbreak{} ms4,\allowbreak{}\allowbreak{} localization\allowbreak{} paragraph}{Thus\allowbreak{} there\allowbreak{} is\allowbreak{} ,\allowbreak{}}{Thus\allowbreak{} there\allowbreak{} is,\allowbreak{}}{DER\_\allowbreak{}III-V:\allowbreak{}4}
\DERapparatusitem{ENV0183}{3;\allowbreak{}\allowbreak{} III\allowbreak{} ms4,\allowbreak{}\allowbreak{} equation\allowbreak{} (\allowbreak{}7)\allowbreak{}}{D}{\textbackslash{}\allowbreak{}D}{DER\_\allowbreak{}III-V:\allowbreak{}5}
\DERapparatusitem{ENV0184}{preamble;\allowbreak{}\allowbreak{} local\allowbreak{} macro\allowbreak{} definition}{\textbackslash{}\allowbreak{}text\{\allowbreak{}def\}\allowbreak{}}{\textbackslash{}\allowbreak{}mathrm\{\allowbreak{}def\}\allowbreak{}}{TYPESETTING:\allowbreak{}7;\allowbreak{}\allowbreak{} NOTATION\_\allowbreak{}DER:\allowbreak{}17}
\DERapparatusitem{ENV0185}{5;\allowbreak{}\allowbreak{} III\allowbreak{} ms7,\allowbreak{}\allowbreak{} paragraph\allowbreak{} introducing\allowbreak{} Hom\_\allowbreak{}xi}{\textbackslash{}\allowbreak{}xi\textbackslash{}\allowbreak{}in(\allowbreak{}\textbackslash{}\allowbreak{}mathrm\{\allowbreak{}hot\}\allowbreak{})\allowbreak{}}{\textbackslash{}\allowbreak{}xi\textbackslash{}\allowbreak{}in\textbackslash{}\allowbreak{}Hot}{DER\_\allowbreak{}III-V:\allowbreak{}8}
\subsection*{Chapter IV}
\DERapparatusitem{ENV0186}{1;\allowbreak{}\allowbreak{} IV\allowbreak{} ms73,\allowbreak{}\allowbreak{} skeleton\allowbreak{} functor\allowbreak{} display}{\textbackslash{}\allowbreak{}mathcal\{\allowbreak{}D\}\allowbreak{}}{\textbackslash{}\allowbreak{}D}{DER\_\allowbreak{}III-V:\allowbreak{}13}
\DERapparatusitem{ENV0187}{1;\allowbreak{}\allowbreak{} IV\allowbreak{} ms74,\allowbreak{}\allowbreak{} walking-arrow\allowbreak{} paragraph}{0}{1}{DER\_\allowbreak{}III-V:\allowbreak{}14}
\DERapparatusitem{ENV0188}{1;\allowbreak{}\allowbreak{} IV\allowbreak{} ms74,\allowbreak{}\allowbreak{} walking-arrow\allowbreak{} paragraph}{0}{1}{DER\_\allowbreak{}III-V:\allowbreak{}14}
\DERapparatusitem{ENV0189}{1;\allowbreak{}\allowbreak{} IV\allowbreak{} ms74,\allowbreak{}\allowbreak{} sentence\allowbreak{} before\allowbreak{} coefficient\allowbreak{} square}{the\allowbreak{} substantial\allowbreak{} diagrams\allowbreak{} of\allowbreak{} type\allowbreak{} \$\allowbreak{}I\$\allowbreak{}\allowbreak{} are\allowbreak{} written}{the\allowbreak{} substantial\allowbreak{} diagrams\allowbreak{} of\allowbreak{} type\allowbreak{} \$\allowbreak{}I\textasciicircum{}\allowbreak{}o\$\allowbreak{}\allowbreak{} are\allowbreak{} written}{DER\_\allowbreak{}III-V:\allowbreak{}15}
\DERapparatusitem{ENV0190}{3;\allowbreak{}\allowbreak{} IV\allowbreak{} ms77,\allowbreak{}\allowbreak{} quasi-inverse\allowbreak{} description}{U}{I}{DER\_\allowbreak{}III-V:\allowbreak{}16}
\DERapparatusitem{ENV0191}{4;\allowbreak{}\allowbreak{} IV\allowbreak{} ms78,\allowbreak{}\allowbreak{} introduction\allowbreak{} to\allowbreak{} relative\allowbreak{} morphism\allowbreak{} diagram}{the\allowbreak{} situation\allowbreak{} opposite\allowbreak{} \textbackslash{}\allowbreak{}texttt\{\allowbreak{}[\allowbreak{}below]\allowbreak{}\}\allowbreak{}}{the\allowbreak{} situation\allowbreak{} below\allowbreak{} \textbackslash{}\allowbreak{}texttt\{\allowbreak{}[\allowbreak{}below]\allowbreak{}\}\allowbreak{}}{DER\_\allowbreak{}III-V:\allowbreak{}17}
\DERapparatusitem{ENV0192}{4;\allowbreak{}\allowbreak{} running\allowbreak{} text\allowbreak{} or\allowbreak{} formula}{\textbackslash{}\allowbreak{}text\{\allowbreak{}def\}\allowbreak{}}{\textbackslash{}\allowbreak{}mathrm\{\allowbreak{}def\}\allowbreak{}}{TYPESETTING:\allowbreak{}7;\allowbreak{}\allowbreak{} NOTATION\_\allowbreak{}DER:\allowbreak{}17}
\DERapparatusitem{ENV0193}{5;\allowbreak{}\allowbreak{} IV\allowbreak{} ms79,\allowbreak{}\allowbreak{} first\allowbreak{} paragraph}{are\allowbreak{} now\allowbreak{} fleshed\allowbreak{} out)\allowbreak{}}{are\allowbreak{} now\allowbreak{} fleshed\allowbreak{} out}{DER\_\allowbreak{}III-V:\allowbreak{}18}
\DERapparatusitem{ENV0194}{5;\allowbreak{}\allowbreak{} IV\allowbreak{} ms79,\allowbreak{}\allowbreak{} object\allowbreak{} set\allowbreak{} of\allowbreak{} K(\allowbreak{}f)\allowbreak{}}{K(\allowbreak{}I)\allowbreak{}\allowbreak{} \textbackslash{}\allowbreak{}amalg\allowbreak{} \textbackslash{}\allowbreak{}Ob\allowbreak{} K(\allowbreak{}J)\allowbreak{}}{I\allowbreak{} \textbackslash{}\allowbreak{}amalg\allowbreak{} \textbackslash{}\allowbreak{}Ob\allowbreak{} J}{DER\_\allowbreak{}III-V:\allowbreak{}19}
\DERapparatusitem{ENV0195}{6;\allowbreak{}\allowbreak{} running\allowbreak{} text\allowbreak{} or\allowbreak{} formula}{\textbackslash{}\allowbreak{}text\{\allowbreak{}def\}\allowbreak{}}{\textbackslash{}\allowbreak{}mathrm\{\allowbreak{}def\}\allowbreak{}}{TYPESETTING:\allowbreak{}7;\allowbreak{}\allowbreak{} NOTATION\_\allowbreak{}DER:\allowbreak{}17}
\DERapparatusitem{ENV0196}{6;\allowbreak{}\allowbreak{} IV\allowbreak{} ms81,\allowbreak{}\allowbreak{} inline\allowbreak{} diagram\allowbreak{} of\allowbreak{} structural\allowbreak{} maps}{I\allowbreak{} \textbackslash{}\allowbreak{}arrow[\allowbreak{}r,\allowbreak{}\allowbreak{} hook,\allowbreak{}\allowbreak{} shift\allowbreak{} left,\allowbreak{}\allowbreak{} "j"]\allowbreak{}}{J\allowbreak{} \textbackslash{}\allowbreak{}arrow[\allowbreak{}r,\allowbreak{}\allowbreak{} hook,\allowbreak{}\allowbreak{} shift\allowbreak{} left,\allowbreak{}\allowbreak{} "j"]\allowbreak{}}{DER\_\allowbreak{}III-V:\allowbreak{}20}
\DERapparatusitem{ENV0197}{7;\allowbreak{}\allowbreak{} IV\allowbreak{} ms82,\allowbreak{}\allowbreak{} universal\allowbreak{} property}{\textbackslash{}\allowbreak{}\{\allowbreak{}F,\allowbreak{}\textbackslash{}\allowbreak{}\allowbreak{} G,\allowbreak{}\textbackslash{}\allowbreak{}\allowbreak{} u\textbackslash{}\allowbreak{}\}\allowbreak{}\textasciicircum{}\allowbreak{}o(\allowbreak{}\textbackslash{}\allowbreak{}mathcal\{\allowbreak{}M\}\allowbreak{})\allowbreak{}}{\textbackslash{}\allowbreak{}\{\allowbreak{}F,\allowbreak{}\textbackslash{}\allowbreak{}\allowbreak{} G,\allowbreak{}\textbackslash{}\allowbreak{}\allowbreak{} u\textbackslash{}\allowbreak{}\}\allowbreak{}(\allowbreak{}\textbackslash{}\allowbreak{}mathcal\{\allowbreak{}M\}\allowbreak{})\allowbreak{}}{DER\_\allowbreak{}III-V:\allowbreak{}21}
\subsection*{Chapter V}
\DERapparatusitem{ENV0222}{1;\allowbreak{}\allowbreak{} Chapter\allowbreak{} V,\allowbreak{}\allowbreak{} ms.\allowbreak{}\allowbreak{} 1:\allowbreak{}\allowbreak{} odd\allowbreak{} generator\allowbreak{} target}{\textbackslash{}\allowbreak{}alpha-1}{\textbackslash{}\allowbreak{}alpha}{DER\_\allowbreak{}III-V:\allowbreak{}24}
\DERapparatusitem{ENV0223}{1;\allowbreak{}\allowbreak{} Chapter\allowbreak{} V,\allowbreak{}\allowbreak{} ms.\allowbreak{}\allowbreak{} 1:\allowbreak{}\allowbreak{} same\allowbreak{} composition\allowbreak{} sentence}{if\allowbreak{} two\allowbreak{} arrows\allowbreak{} are\allowbreak{} composable,\allowbreak{}\allowbreak{} they\allowbreak{} are\allowbreak{} the\allowbreak{} identity}{if\allowbreak{} two\allowbreak{} arrows\allowbreak{} are\allowbreak{} composable,\allowbreak{}\allowbreak{} one\allowbreak{} of\allowbreak{} them\allowbreak{} is\allowbreak{} the\allowbreak{} identity}{DER\_\allowbreak{}III-V:\allowbreak{}26}
\DERapparatusitem{ENV0224}{1;\allowbreak{}\allowbreak{} Chapter\allowbreak{} V,\allowbreak{}\allowbreak{} ms.\allowbreak{}\allowbreak{} 1:\allowbreak{}\allowbreak{} strict\allowbreak{} order\allowbreak{} formula}{\$\allowbreak{}\textbackslash{}\allowbreak{}alpha\$\allowbreak{}\allowbreak{} even,\allowbreak{}\allowbreak{} \$\allowbreak{}\textbackslash{}\allowbreak{}beta=\allowbreak{}\textbackslash{}\allowbreak{}alpha+1\$\allowbreak{}\allowbreak{} }{\$\allowbreak{}\textbackslash{}\allowbreak{}alpha\$\allowbreak{}\allowbreak{} even,\allowbreak{}\allowbreak{} \$\allowbreak{}\textbackslash{}\allowbreak{}beta=\allowbreak{}\textbackslash{}\allowbreak{}alpha+1\$\allowbreak{}\allowbreak{} or\allowbreak{} \$\allowbreak{}\textbackslash{}\allowbreak{}beta=\allowbreak{}\textbackslash{}\allowbreak{}alpha-1\$\allowbreak{}\allowbreak{} }{DER\_\allowbreak{}III-V:\allowbreak{}27}
\DERapparatusitem{ENV0225}{1;\allowbreak{}\allowbreak{} Chapter\allowbreak{} V,\allowbreak{}\allowbreak{} ms.\allowbreak{}\allowbreak{} 1:\allowbreak{}\allowbreak{} last\allowbreak{} paragraph}{\textbackslash{}\allowbreak{}Psi}{\textbackslash{}\allowbreak{}Phi}{DER\_\allowbreak{}III-V:\allowbreak{}28}
\DERapparatusitem{ENV0226}{2;\allowbreak{}\allowbreak{} Chapter\allowbreak{} V,\allowbreak{}\allowbreak{} ms.\allowbreak{}\allowbreak{} 2:\allowbreak{}\allowbreak{} clamping\allowbreak{} display}{\textbackslash{}\allowbreak{}[\allowbreak{}\space{}\allowbreak{}\textbackslash{}\allowbreak{}left\textbackslash{}\allowbreak{}\{\allowbreak{}\space{}\allowbreak{}\textbackslash{}\allowbreak{}begin\{\allowbreak{}array\}\allowbreak{}\{\allowbreak{}l\}\allowbreak{}\space{}\allowbreak{}\textbackslash{}\allowbreak{}varphi\_\allowbreak{}\{\allowbreak{}n,\allowbreak{}m\}\allowbreak{}:\allowbreak{}\textbackslash{}\allowbreak{}Phi\_\allowbreak{}m\textbackslash{}\allowbreak{}longrightarrow\textbackslash{}\allowbreak{}Phi\_\allowbreak{}n\textbackslash{}\allowbreak{}subset\textbackslash{}\allowbreak{}Phi\_\allowbreak{}m\textbackslash{}\allowbreak{}\textbackslash{}\allowbreak{}[\allowbreak{}1ex]\allowbreak{}\space{}\allowbreak{}\textbackslash{}\allowbreak{}varphi\_\allowbreak{}\{\allowbreak{}n,\allowbreak{}m\}\allowbreak{}(\allowbreak{}\textbackslash{}\allowbreak{}\{\allowbreak{}\textbackslash{}\allowbreak{}alpha\_\allowbreak{}i\textbackslash{}\allowbreak{}\}\allowbreak{})\allowbreak{}=\allowbreak{}\space{}\allowbreak{}\textbackslash{}\allowbreak{}left\textbackslash{}\allowbreak{}\{\allowbreak{}\space{}\allowbreak{}\textbackslash{}\allowbreak{}begin\{\allowbreak{}array\}\allowbreak{}\{\allowbreak{}rcll@\{\allowbreak{}\textbackslash{}\allowbreak{}quad\}\allowbreak{}l@\{\allowbreak{}\textbackslash{}\allowbreak{}\allowbreak{} \textbackslash{}\allowbreak{}\allowbreak{} \}\allowbreak{}r\}\allowbreak{}\space{}\allowbreak{}\textbackslash{}\allowbreak{}alpha\_\allowbreak{}i\allowbreak{} \&\allowbreak{}\allowbreak{} =\allowbreak{}\allowbreak{} \&\allowbreak{}\allowbreak{} i\allowbreak{} \&\allowbreak{}\allowbreak{} \textbackslash{}\allowbreak{}in\textbackslash{}\allowbreak{}Ob\textbackslash{}\allowbreak{}Phi\_\allowbreak{}n,\allowbreak{}\allowbreak{} \&\allowbreak{}\allowbreak{} \textbackslash{}\allowbreak{}text\{\allowbreak{}i.\allowbreak{}e.\allowbreak{}\}\allowbreak{}\allowbreak{} \&\allowbreak{}\allowbreak{} -n\textbackslash{}\allowbreak{}le\allowbreak{} i\textbackslash{}\allowbreak{}le\allowbreak{} n\textbackslash{}\allowbreak{}\textbackslash{}\allowbreak{}[\allowbreak{}1ex]\allowbreak{}\space{}\allowbreak{}\textbackslash{}\allowbreak{}alpha\_\allowbreak{}i\allowbreak{} \&\allowbreak{}\allowbreak{} =\allowbreak{}\allowbreak{} \&\allowbreak{}\allowbreak{} n\allowbreak{} \&\allowbreak{}\allowbreak{} \&\allowbreak{}\allowbreak{} \textbackslash{}\allowbreak{}text\{\allowbreak{}if\}\allowbreak{}\allowbreak{} \&\allowbreak{}\allowbreak{} n\textbackslash{}\allowbreak{}le\allowbreak{} i\textbackslash{}\allowbreak{}le\allowbreak{} m\textbackslash{}\allowbreak{}\textbackslash{}\allowbreak{}[\allowbreak{}1ex]\allowbreak{}\space{}\allowbreak{}\textbackslash{}\allowbreak{}alpha\_\allowbreak{}i\allowbreak{} \&\allowbreak{}\allowbreak{} =\allowbreak{}\allowbreak{} \&\allowbreak{}\allowbreak{} -n\allowbreak{} \&\allowbreak{}\allowbreak{} \&\allowbreak{}\allowbreak{} \textbackslash{}\allowbreak{}text\{\allowbreak{}if\}\allowbreak{}\allowbreak{} \&\allowbreak{}\allowbreak{} -m\textbackslash{}\allowbreak{}le\allowbreak{} i\textbackslash{}\allowbreak{}le\allowbreak{} -n\space{}\allowbreak{}\textbackslash{}\allowbreak{}end\{\allowbreak{}array\}\allowbreak{}\space{}\allowbreak{}\textbackslash{}\allowbreak{}right.\allowbreak{}\space{}\allowbreak{}\textbackslash{}\allowbreak{}end\{\allowbreak{}array\}\allowbreak{}\space{}\allowbreak{}\textbackslash{}\allowbreak{}right.\allowbreak{}\space{}\allowbreak{}\textbackslash{}\allowbreak{}]\allowbreak{}}{\textbackslash{}\allowbreak{}[\allowbreak{}\space{}\allowbreak{}\textbackslash{}\allowbreak{}begin\{\allowbreak{}array\}\allowbreak{}\{\allowbreak{}rcl\}\allowbreak{}\space{}\allowbreak{}\textbackslash{}\allowbreak{}varphi\_\allowbreak{}\{\allowbreak{}n,\allowbreak{}m\}\allowbreak{}:\allowbreak{}\textbackslash{}\allowbreak{}Phi\_\allowbreak{}m\allowbreak{} \&\allowbreak{}\allowbreak{} \textbackslash{}\allowbreak{}longrightarrow\allowbreak{} \&\allowbreak{}\allowbreak{} \textbackslash{}\allowbreak{}Phi\_\allowbreak{}n\textbackslash{}\allowbreak{}subset\textbackslash{}\allowbreak{}Phi\_\allowbreak{}m\textbackslash{}\allowbreak{}\textbackslash{}\allowbreak{}[\allowbreak{}1ex]\allowbreak{}\space{}\allowbreak{}\textbackslash{}\allowbreak{}varphi\_\allowbreak{}\{\allowbreak{}n,\allowbreak{}m\}\allowbreak{}(\allowbreak{}\textbackslash{}\allowbreak{}\{\allowbreak{}i\textbackslash{}\allowbreak{}\}\allowbreak{})\allowbreak{}\allowbreak{} \&\allowbreak{}\allowbreak{} =\allowbreak{}\allowbreak{} \&\allowbreak{}\space{}\allowbreak{}\textbackslash{}\allowbreak{}left\textbackslash{}\allowbreak{}\{\allowbreak{}\space{}\allowbreak{}\textbackslash{}\allowbreak{}begin\{\allowbreak{}array\}\allowbreak{}\{\allowbreak{}l@\{\allowbreak{}\textbackslash{}\allowbreak{}qquad\}\allowbreak{}l@\{\allowbreak{}\textbackslash{}\allowbreak{}quad\}\allowbreak{}l@\{\allowbreak{}\textbackslash{}\allowbreak{}\allowbreak{} \textbackslash{}\allowbreak{}\allowbreak{} \}\allowbreak{}r\}\allowbreak{}\space{}\allowbreak{}\textbackslash{}\allowbreak{}\{\allowbreak{}i\textbackslash{}\allowbreak{}\}\allowbreak{}\allowbreak{} \&\allowbreak{}\allowbreak{} \textbackslash{}\allowbreak{}text\{\allowbreak{}if\}\allowbreak{}\textbackslash{}\allowbreak{}\allowbreak{} i\textbackslash{}\allowbreak{}in\textbackslash{}\allowbreak{}Ob\textbackslash{}\allowbreak{}Phi\_\allowbreak{}n,\allowbreak{}\allowbreak{} \&\allowbreak{}\allowbreak{} \textbackslash{}\allowbreak{}text\{\allowbreak{}i.\allowbreak{}e\}\allowbreak{}\allowbreak{} \&\allowbreak{}\allowbreak{} -n\textbackslash{}\allowbreak{}le\allowbreak{} i\textbackslash{}\allowbreak{}le\allowbreak{} n\textbackslash{}\allowbreak{}\textbackslash{}\allowbreak{}[\allowbreak{}1ex]\allowbreak{}\space{}\allowbreak{}\textbackslash{}\allowbreak{}\{\allowbreak{}n\textbackslash{}\allowbreak{}\}\allowbreak{}\allowbreak{} \&\allowbreak{}\allowbreak{} \&\allowbreak{}\allowbreak{} \textbackslash{}\allowbreak{}text\{\allowbreak{}if\}\allowbreak{}\allowbreak{} \&\allowbreak{}\allowbreak{} n\textbackslash{}\allowbreak{}le\allowbreak{} i\textbackslash{}\allowbreak{}le\allowbreak{} m\textbackslash{}\allowbreak{}\textbackslash{}\allowbreak{}[\allowbreak{}1ex]\allowbreak{}\space{}\allowbreak{}\textbackslash{}\allowbreak{}\{\allowbreak{}-n\textbackslash{}\allowbreak{}\}\allowbreak{}\allowbreak{} \&\allowbreak{}\allowbreak{} \&\allowbreak{}\allowbreak{} \textbackslash{}\allowbreak{}text\{\allowbreak{}if\}\allowbreak{}\allowbreak{} \&\allowbreak{}\allowbreak{} -m\textbackslash{}\allowbreak{}le\allowbreak{} i\textbackslash{}\allowbreak{}le\allowbreak{} -n\space{}\allowbreak{}\textbackslash{}\allowbreak{}end\{\allowbreak{}array\}\allowbreak{}\space{}\allowbreak{}\textbackslash{}\allowbreak{}right.\allowbreak{}\space{}\allowbreak{}\textbackslash{}\allowbreak{}end\{\allowbreak{}array\}\allowbreak{}\space{}\allowbreak{}\textbackslash{}\allowbreak{}]\allowbreak{}}{DER\_\allowbreak{}III-V:\allowbreak{}30}
\DERapparatusitem{ENV0227}{2;\allowbreak{}\allowbreak{} running\allowbreak{} text\allowbreak{} or\allowbreak{} formula}{\textbackslash{}\allowbreak{}text\{\allowbreak{}def\}\allowbreak{}}{\textbackslash{}\allowbreak{}mathrm\{\allowbreak{}def\}\allowbreak{}}{TYPESETTING:\allowbreak{}7;\allowbreak{}\allowbreak{} NOTATION\_\allowbreak{}DER:\allowbreak{}17}
\DERapparatusitem{ENV0228}{2;\allowbreak{}\allowbreak{} Chapter\allowbreak{} V,\allowbreak{}\allowbreak{} ms.\allowbreak{}\allowbreak{} 2:\allowbreak{}\allowbreak{} edge\allowbreak{} range}{}{-1}{DER\_\allowbreak{}III-V:\allowbreak{}31}
\DERapparatusitem{ENV0229}{2;\allowbreak{}\allowbreak{} Chapter\allowbreak{} V,\allowbreak{}\allowbreak{} ms.\allowbreak{}\allowbreak{} 2:\allowbreak{}\allowbreak{} target\allowbreak{} of\allowbreak{} padding\allowbreak{} functor}{Y}{X}{DER\_\allowbreak{}III-V:\allowbreak{}32}
\DERapparatusitem{ENV0230}{3;\allowbreak{}\allowbreak{} Chapter\allowbreak{} V,\allowbreak{}\allowbreak{} ms.\allowbreak{}\allowbreak{} 3:\allowbreak{}\allowbreak{} full\allowbreak{} faithfulness\allowbreak{} sentence}{is\allowbreak{} surjective\allowbreak{} on\allowbreak{} objects\allowbreak{} and\allowbreak{} on\allowbreak{} arrows.\allowbreak{}}{is\allowbreak{} surjective\allowbreak{} on\allowbreak{} objects\allowbreak{} and\allowbreak{} on\allowbreak{} arrows)\allowbreak{}.\allowbreak{}}{DER\_\allowbreak{}III-V:\allowbreak{}34}
\DERapparatusitem{ENV0231}{3;\allowbreak{}\allowbreak{} running\allowbreak{} text\allowbreak{} or\allowbreak{} formula}{\textbackslash{}\allowbreak{}text\{\allowbreak{}def\}\allowbreak{}}{\textbackslash{}\allowbreak{}mathrm\{\allowbreak{}def\}\allowbreak{}}{TYPESETTING:\allowbreak{}7;\allowbreak{}\allowbreak{} NOTATION\_\allowbreak{}DER:\allowbreak{}17}
\DERapparatusitem{ENV0232}{3;\allowbreak{}\allowbreak{} Chapter\allowbreak{} V,\allowbreak{}\allowbreak{} ms.\allowbreak{}\allowbreak{} 4:\allowbreak{}\allowbreak{} automorphism\allowbreak{} group}{The\allowbreak{} group\allowbreak{} of\allowbreak{} automorphisms\allowbreak{} of\allowbreak{} the\allowbreak{} \$\allowbreak{}\textbackslash{}\allowbreak{}Phi\_\allowbreak{}n\$\allowbreak{}\allowbreak{} is}{The\allowbreak{} group\allowbreak{} of\allowbreak{} automorphisms\allowbreak{} of\allowbreak{} the\allowbreak{} \$\allowbreak{}\textbackslash{}\allowbreak{}Phi\_\allowbreak{}n\$\allowbreak{},\allowbreak{}\allowbreak{} for\allowbreak{} \$\allowbreak{}n>0\$\allowbreak{},\allowbreak{}\allowbreak{} is}{DER\_\allowbreak{}III-V:\allowbreak{}37}
\DERapparatusitem{ENV0233}{3;\allowbreak{}\allowbreak{} Chapter\allowbreak{} V,\allowbreak{}\allowbreak{} ms.\allowbreak{}\allowbreak{} 4:\allowbreak{}\allowbreak{} reflection\allowbreak{} parenthesis}{symmetry\allowbreak{} \$\allowbreak{}\textbackslash{}\allowbreak{}\{\allowbreak{}i\textbackslash{}\allowbreak{}\}\allowbreak{}\textbackslash{}\allowbreak{}longrightarrow\textbackslash{}\allowbreak{}\{\allowbreak{}-i\textbackslash{}\allowbreak{}\}\allowbreak{}\$\allowbreak{}.\allowbreak{}\allowbreak{} So}{symmetry\allowbreak{} \$\allowbreak{}\textbackslash{}\allowbreak{}\{\allowbreak{}i\textbackslash{}\allowbreak{}\}\allowbreak{}\textbackslash{}\allowbreak{}longrightarrow\textbackslash{}\allowbreak{}\{\allowbreak{}-i\textbackslash{}\allowbreak{}\}\allowbreak{}\$\allowbreak{})\allowbreak{}.\allowbreak{}\allowbreak{} So}{DER\_\allowbreak{}III-V:\allowbreak{}36}
\DERapparatusitem{ENV0234}{4;\allowbreak{}\allowbreak{} Chapter\allowbreak{} V,\allowbreak{}\allowbreak{} ms.\allowbreak{}\allowbreak{} 4:\allowbreak{}\allowbreak{} endpoint\allowbreak{} maps\allowbreak{} called\allowbreak{} inclusions}{define\allowbreak{} inclusion\allowbreak{} functors\allowbreak{} (\allowbreak{}fully\allowbreak{} faithful)\allowbreak{}}{define\allowbreak{} functors\allowbreak{} (\allowbreak{}inclusions\allowbreak{} and\allowbreak{} fully\allowbreak{} faithful\allowbreak{} for\allowbreak{} \$\allowbreak{}n>0\$\allowbreak{})\allowbreak{}}{DER\_\allowbreak{}III-V:\allowbreak{}39}
\DERapparatusitem{ENV0235}{4;\allowbreak{}\allowbreak{} Chapter\allowbreak{} V,\allowbreak{}\allowbreak{} ms.\allowbreak{}\allowbreak{} 5:\allowbreak{}\allowbreak{} source-target\allowbreak{} variable}{x)\allowbreak{},\allowbreak{}\textbackslash{}\allowbreak{}beta(\allowbreak{}x}{c)\allowbreak{},\allowbreak{}\textbackslash{}\allowbreak{}beta(\allowbreak{}c}{DER\_\allowbreak{}III-V:\allowbreak{}40}
\DERapparatusitem{ENV0236}{4;\allowbreak{}\allowbreak{} Chapter\allowbreak{} V,\allowbreak{}\allowbreak{} ms.\allowbreak{}\allowbreak{} 5:\allowbreak{}\allowbreak{} eventual\allowbreak{} endpoint\allowbreak{} formulas}{\textbackslash{}\allowbreak{}sigma(\allowbreak{}x)\allowbreak{}\allowbreak{} \&\allowbreak{}=\allowbreak{}\allowbreak{} A\_\allowbreak{}i\textbackslash{}\allowbreak{}quad\textbackslash{}\allowbreak{}text\{\allowbreak{}for\allowbreak{} \}\allowbreak{}\allowbreak{} i\textbackslash{}\allowbreak{}leq\allowbreak{} -\textbackslash{}\allowbreak{}mathrm\{\allowbreak{}long\}\allowbreak{}(\allowbreak{}x)\allowbreak{}}{\textbackslash{}\allowbreak{}sigma(\allowbreak{}c)\allowbreak{}\allowbreak{} \&\allowbreak{}=\allowbreak{}\allowbreak{} A\_\allowbreak{}i\textbackslash{}\allowbreak{}quad\textbackslash{}\allowbreak{}text\{\allowbreak{}for\allowbreak{} \}\allowbreak{}\allowbreak{} i\textbackslash{}\allowbreak{}leq\allowbreak{} -\textbackslash{}\allowbreak{}mathrm\{\allowbreak{}long\}\allowbreak{}(\allowbreak{}c)\allowbreak{}}{DER\_\allowbreak{}III-V:\allowbreak{}41}
\DERapparatusitem{ENV0237}{4;\allowbreak{}\allowbreak{} Chapter\allowbreak{} V,\allowbreak{}\allowbreak{} ms.\allowbreak{}\allowbreak{} 5:\allowbreak{}\allowbreak{} eventual\allowbreak{} endpoint\allowbreak{} formulas}{\textbackslash{}\allowbreak{}beta(\allowbreak{}x)\allowbreak{}\allowbreak{} \&\allowbreak{}=\allowbreak{}\allowbreak{} A\_\allowbreak{}i\textbackslash{}\allowbreak{}quad\textbackslash{}\allowbreak{}text\{\allowbreak{}for\allowbreak{} \}\allowbreak{}\allowbreak{} i\textbackslash{}\allowbreak{}geq\allowbreak{} \textbackslash{}\allowbreak{}mathrm\{\allowbreak{}long\}\allowbreak{}(\allowbreak{}x)\allowbreak{}}{\textbackslash{}\allowbreak{}beta(\allowbreak{}c)\allowbreak{}\allowbreak{} \&\allowbreak{}=\allowbreak{}\allowbreak{} A\_\allowbreak{}i\textbackslash{}\allowbreak{}quad\textbackslash{}\allowbreak{}text\{\allowbreak{}for\allowbreak{} \}\allowbreak{}\allowbreak{} i\textbackslash{}\allowbreak{}geq\allowbreak{} \textbackslash{}\allowbreak{}mathrm\{\allowbreak{}long\}\allowbreak{}(\allowbreak{}c)\allowbreak{}}{DER\_\allowbreak{}III-V:\allowbreak{}41}
\DERapparatusitem{ENV0238}{5;\allowbreak{}\allowbreak{} Chapter\allowbreak{} V,\allowbreak{}\allowbreak{} ms.\allowbreak{}\allowbreak{} 6:\allowbreak{}\allowbreak{} footnote\allowbreak{} 3}{The\allowbreak{} set\allowbreak{} of\allowbreak{} its\allowbreak{} objects\allowbreak{} is\allowbreak{} denoted\allowbreak{} \$\allowbreak{}\textbackslash{}\allowbreak{}mathrm\{\allowbreak{}Ch\}\allowbreak{}(\allowbreak{}X)\allowbreak{}\$\allowbreak{}.\allowbreak{}}{The\allowbreak{} set\allowbreak{} of\allowbreak{} its\allowbreak{} objects\allowbreak{} is\allowbreak{} denoted\allowbreak{} \$\allowbreak{}\textbackslash{}\allowbreak{}mathrm\{\allowbreak{}Ch\}\allowbreak{}\_\allowbreak{}\textbackslash{}\allowbreak{}Sigma(\allowbreak{}X)\allowbreak{}\$\allowbreak{}.\allowbreak{}}{DER\_\allowbreak{}III-V:\allowbreak{}43}
\DERapparatusitem{ENV0239}{5;\allowbreak{}\allowbreak{} Chapter\allowbreak{} V,\allowbreak{}\allowbreak{} ms.\allowbreak{}\allowbreak{} 6:\allowbreak{}\allowbreak{} footnote\allowbreak{} 3}{Its\allowbreak{} elements\allowbreak{} are\allowbreak{} the\allowbreak{} \$\allowbreak{}\textbackslash{}\allowbreak{}Sigma\$\allowbreak{}-paths\allowbreak{} from\allowbreak{} \$\allowbreak{}x\$\allowbreak{}\allowbreak{} to\allowbreak{} \$\allowbreak{}y\$\allowbreak{}}{Its\allowbreak{} elements\allowbreak{} are\allowbreak{} the\allowbreak{} \$\allowbreak{}\textbackslash{}\allowbreak{}Sigma\$\allowbreak{}-paths}{DER\_\allowbreak{}III-V:\allowbreak{}43}
\DERapparatusitem{ENV0240}{5;\allowbreak{}\allowbreak{} Chapter\allowbreak{} V,\allowbreak{}\allowbreak{} ms.\allowbreak{}\allowbreak{} 6:\allowbreak{}\allowbreak{} retrograde\allowbreak{} parity}{For\allowbreak{} \$\allowbreak{}i\$\allowbreak{}\allowbreak{} odd,\allowbreak{}}{For\allowbreak{} \$\allowbreak{}i\$\allowbreak{}\allowbreak{} even,\allowbreak{}}{DER\_\allowbreak{}III-V:\allowbreak{}42}
\DERapparatusitem{ENV0241}{6;\allowbreak{}\allowbreak{} Chapter\allowbreak{} V,\allowbreak{}\allowbreak{} ms.\allowbreak{}\allowbreak{} 8--9:\allowbreak{}\allowbreak{} initial\allowbreak{} parity\allowbreak{} assertion}{\$\allowbreak{}m\$\allowbreak{},\allowbreak{}\allowbreak{} \$\allowbreak{}n\$\allowbreak{}\allowbreak{} have\allowbreak{} \textbackslash{}\allowbreak{}textit\{\allowbreak{}the\allowbreak{} same\allowbreak{} parity,\allowbreak{}\}\allowbreak{}}{\$\allowbreak{}m\$\allowbreak{},\allowbreak{}\allowbreak{} \$\allowbreak{}n\$\allowbreak{}\allowbreak{} are\allowbreak{} \textbackslash{}\allowbreak{}textit\{\allowbreak{}even,\allowbreak{}\}\allowbreak{}}{DER\_\allowbreak{}III-V:\allowbreak{}47}
\DERapparatusitem{ENV0242}{6;\allowbreak{}\allowbreak{} Chapter\allowbreak{} V,\allowbreak{}\allowbreak{} ms.\allowbreak{}\allowbreak{} 9:\allowbreak{}\allowbreak{} interval\allowbreak{} embeddings}{n\allowbreak{} \&\allowbreak{}\textbackslash{}\allowbreak{}xrightarrow\{\allowbreak{}\textbackslash{}\allowbreak{}alpha\}\allowbreak{}\allowbreak{} \textbackslash{}\allowbreak{}Phi\_\allowbreak{}\{\allowbreak{}n+m\}\allowbreak{}\textbackslash{}\allowbreak{}quad\allowbreak{} i\textbackslash{}\allowbreak{}longmapsto\allowbreak{} (\allowbreak{}i-m)\allowbreak{}\textbackslash{}\allowbreak{}\textbackslash{}\allowbreak{}\space{}\allowbreak{}\textbackslash{}\allowbreak{}Phi\_\allowbreak{}m\allowbreak{} \&\allowbreak{}\textbackslash{}\allowbreak{}xrightarrow\{\allowbreak{}\textbackslash{}\allowbreak{}beta\}\allowbreak{}\allowbreak{} \textbackslash{}\allowbreak{}Phi\_\allowbreak{}\{\allowbreak{}n+m\}\allowbreak{}\textbackslash{}\allowbreak{}quad\allowbreak{} i\textbackslash{}\allowbreak{}longmapsto\allowbreak{} (\allowbreak{}i+n}{m\allowbreak{} \&\allowbreak{}\textbackslash{}\allowbreak{}xrightarrow\{\allowbreak{}\textbackslash{}\allowbreak{}alpha\}\allowbreak{}\allowbreak{} \textbackslash{}\allowbreak{}Phi\_\allowbreak{}\{\allowbreak{}m+n\}\allowbreak{}\textbackslash{}\allowbreak{}quad\allowbreak{} i\textbackslash{}\allowbreak{}longmapsto\allowbreak{} (\allowbreak{}i-n)\allowbreak{}\textbackslash{}\allowbreak{}\textbackslash{}\allowbreak{}\space{}\allowbreak{}\textbackslash{}\allowbreak{}Phi\_\allowbreak{}n\allowbreak{} \&\allowbreak{}\textbackslash{}\allowbreak{}xrightarrow\{\allowbreak{}\textbackslash{}\allowbreak{}beta\}\allowbreak{}\allowbreak{} \textbackslash{}\allowbreak{}Phi\_\allowbreak{}\{\allowbreak{}m+n\}\allowbreak{}\textbackslash{}\allowbreak{}quad\allowbreak{} i\textbackslash{}\allowbreak{}longmapsto\allowbreak{} (\allowbreak{}i+m}{DER\_\allowbreak{}III-V:\allowbreak{}48}
\DERapparatusitem{ENV0243}{7;\allowbreak{}\allowbreak{} Chapter\allowbreak{} V,\allowbreak{}\allowbreak{} ms.\allowbreak{}\allowbreak{} 9:\allowbreak{}\allowbreak{} interval\allowbreak{} embeddings}{\textbackslash{}\allowbreak{}\{\allowbreak{}n\textbackslash{}\allowbreak{}\}\allowbreak{})\allowbreak{}=\allowbreak{}\textbackslash{}\allowbreak{}beta(\allowbreak{}\textbackslash{}\allowbreak{}\{\allowbreak{}-m\textbackslash{}\allowbreak{}\}\allowbreak{})\allowbreak{}=\allowbreak{}n-m}{\textbackslash{}\allowbreak{}\{\allowbreak{}m\textbackslash{}\allowbreak{}\}\allowbreak{})\allowbreak{}=\allowbreak{}\textbackslash{}\allowbreak{}beta(\allowbreak{}\textbackslash{}\allowbreak{}\{\allowbreak{}-n\textbackslash{}\allowbreak{}\}\allowbreak{})\allowbreak{}=\allowbreak{}m-n}{DER\_\allowbreak{}III-V:\allowbreak{}48}
\DERapparatusitem{ENV0244}{7;\allowbreak{}\allowbreak{} Chapter\allowbreak{} V,\allowbreak{}\allowbreak{} ms.\allowbreak{}\allowbreak{} 9:\allowbreak{}\allowbreak{} interval\allowbreak{} embeddings}{\{\allowbreak{}n+m\}\allowbreak{}\textbackslash{}\allowbreak{}simeq\textbackslash{}\allowbreak{}Phi\_\allowbreak{}n\textbackslash{}\allowbreak{}sqcup\_\allowbreak{}e\textbackslash{}\allowbreak{}Phi\_\allowbreak{}m}{\{\allowbreak{}m+n\}\allowbreak{}\textbackslash{}\allowbreak{}simeq\textbackslash{}\allowbreak{}Phi\_\allowbreak{}m\textbackslash{}\allowbreak{}sqcup\_\allowbreak{}e\textbackslash{}\allowbreak{}Phi\_\allowbreak{}n}{DER\_\allowbreak{}III-V:\allowbreak{}48}
\DERapparatusitem{ENV0245}{7;\allowbreak{}\allowbreak{} Chapter\allowbreak{} V,\allowbreak{}\allowbreak{} ms.\allowbreak{}\allowbreak{} 9:\allowbreak{}\allowbreak{} interval\allowbreak{} embeddings}{i.\allowbreak{}e.\allowbreak{}\allowbreak{} \$\allowbreak{}\textbackslash{}\allowbreak{}Phi\_\allowbreak{}\{\allowbreak{}n+m\}\allowbreak{}\$\allowbreak{}\allowbreak{} is\allowbreak{} the\allowbreak{} ``amalgamated\allowbreak{} sum''\allowbreak{} in\allowbreak{} (\allowbreak{}Cat)\allowbreak{},\allowbreak{}\allowbreak{} obtained\allowbreak{} by\allowbreak{} ``gluing''\allowbreak{} \$\allowbreak{}\textbackslash{}\allowbreak{}Phi\_\allowbreak{}n\$\allowbreak{}\allowbreak{} and\allowbreak{} \$\allowbreak{}\textbackslash{}\allowbreak{}Phi\_\allowbreak{}m\$\allowbreak{},\allowbreak{}\allowbreak{} by\allowbreak{} identification\allowbreak{} of\allowbreak{} the\allowbreak{} ``target\allowbreak{} element''\allowbreak{} of\allowbreak{} \$\allowbreak{}\textbackslash{}\allowbreak{}Phi\_\allowbreak{}n\$\allowbreak{}\allowbreak{} with\allowbreak{} the\allowbreak{} ``source\allowbreak{} element''\allowbreak{} of\allowbreak{} \$\allowbreak{}\textbackslash{}\allowbreak{}Phi\_\allowbreak{}m\$\allowbreak{}.\allowbreak{}}{i.\allowbreak{}e.\allowbreak{}\allowbreak{} \$\allowbreak{}\textbackslash{}\allowbreak{}Phi\_\allowbreak{}\{\allowbreak{}m+n\}\allowbreak{}\$\allowbreak{}\allowbreak{} is\allowbreak{} the\allowbreak{} ``amalgamated\allowbreak{} sum''\allowbreak{} in\allowbreak{} (\allowbreak{}Cat)\allowbreak{},\allowbreak{}\allowbreak{} obtained\allowbreak{} by\allowbreak{} ``gluing''\allowbreak{} \$\allowbreak{}\textbackslash{}\allowbreak{}Phi\_\allowbreak{}m\$\allowbreak{}\allowbreak{} and\allowbreak{} \$\allowbreak{}\textbackslash{}\allowbreak{}Phi\_\allowbreak{}n\$\allowbreak{},\allowbreak{}\allowbreak{} by\allowbreak{} identification\allowbreak{} of\allowbreak{} the\allowbreak{} ``target\allowbreak{} element''\allowbreak{} of\allowbreak{} \$\allowbreak{}\textbackslash{}\allowbreak{}Phi\_\allowbreak{}m\$\allowbreak{}\allowbreak{} with\allowbreak{} the\allowbreak{} ``source\allowbreak{} element''\allowbreak{} of\allowbreak{} \$\allowbreak{}\textbackslash{}\allowbreak{}Phi\_\allowbreak{}n\$\allowbreak{}.\allowbreak{}}{DER\_\allowbreak{}III-V:\allowbreak{}48}
\DERapparatusitem{ENV0246}{7;\allowbreak{}\allowbreak{} running\allowbreak{} text\allowbreak{} or\allowbreak{} formula}{\textbackslash{}\allowbreak{}text\{\allowbreak{}def\}\allowbreak{}}{\textbackslash{}\allowbreak{}mathrm\{\allowbreak{}def\}\allowbreak{}}{TYPESETTING:\allowbreak{}7;\allowbreak{}\allowbreak{} NOTATION\_\allowbreak{}DER:\allowbreak{}17}
\DERapparatusitem{ENV0247}{7;\allowbreak{}\allowbreak{} Chapter\allowbreak{} V,\allowbreak{}\allowbreak{} ms.\allowbreak{}\allowbreak{} 9:\allowbreak{}\allowbreak{} definition\allowbreak{} of\allowbreak{} fixed-stage\allowbreak{} components}{m}{n}{DER\_\allowbreak{}III-V:\allowbreak{}49}
\DERapparatusitem{ENV0248}{7;\allowbreak{}\allowbreak{} Chapter\allowbreak{} V,\allowbreak{}\allowbreak{} ms.\allowbreak{}\allowbreak{} 10:\allowbreak{}\allowbreak{} right-hand\allowbreak{} vertices\allowbreak{} of\allowbreak{} compatibility\allowbreak{} square}{y}{z}{DER\_\allowbreak{}III-V:\allowbreak{}50}
\DERapparatusitem{ENV0249}{7;\allowbreak{}\allowbreak{} Chapter\allowbreak{} V,\allowbreak{}\allowbreak{} ms.\allowbreak{}\allowbreak{} 10:\allowbreak{}\allowbreak{} right-hand\allowbreak{} vertices\allowbreak{} of\allowbreak{} compatibility\allowbreak{} square}{y}{z}{DER\_\allowbreak{}III-V:\allowbreak{}50}
\DERapparatusitem{ENV0250}{7;\allowbreak{}\allowbreak{} Chapter\allowbreak{} V,\allowbreak{}\allowbreak{} ms.\allowbreak{}\allowbreak{} 10:\allowbreak{}\allowbreak{} length\allowbreak{} convention}{of\allowbreak{} length\allowbreak{} \$\allowbreak{}m\$\allowbreak{}\allowbreak{} and\allowbreak{} \$\allowbreak{}n\$\allowbreak{}\allowbreak{} respectively}{of\allowbreak{} length\allowbreak{} \$\allowbreak{}2m\$\allowbreak{}\allowbreak{} and\allowbreak{} \$\allowbreak{}2n\$\allowbreak{}\allowbreak{} respectively}{DER\_\allowbreak{}III-V:\allowbreak{}51}
\DERapparatusitem{ENV0251}{7;\allowbreak{}\allowbreak{} Chapter\allowbreak{} V,\allowbreak{}\allowbreak{} ms.\allowbreak{}\allowbreak{} 10:\allowbreak{}\allowbreak{} length\allowbreak{} convention}{of\allowbreak{} length\allowbreak{} \$\allowbreak{}m'\$\allowbreak{},\allowbreak{}\allowbreak{} \$\allowbreak{}n'\$\allowbreak{}}{of\allowbreak{} length\allowbreak{} \$\allowbreak{}2m'\$\allowbreak{},\allowbreak{}\allowbreak{} \$\allowbreak{}2n'\$\allowbreak{}}{DER\_\allowbreak{}III-V:\allowbreak{}51}
\DERapparatusitem{ENV0252}{8;\allowbreak{}\allowbreak{} Chapter\allowbreak{} V,\allowbreak{}\allowbreak{} ms.\allowbreak{}\allowbreak{} 11:\allowbreak{}\allowbreak{} homotopy\allowbreak{} ambient\allowbreak{} category}{,\allowbreak{}}{+}{DER\_\allowbreak{}III-V:\allowbreak{}53}
\DERapparatusitem{ENV0253}{8;\allowbreak{}\allowbreak{} Chapter\allowbreak{} V,\allowbreak{}\allowbreak{} ms.\allowbreak{}\allowbreak{} 11:\allowbreak{}\allowbreak{} composite\allowbreak{} functor\allowbreak{} subscripts}{+}{,\allowbreak{}}{DER\_\allowbreak{}III-V:\allowbreak{}54}
\DERapparatusitem{ENV0254}{8;\allowbreak{}\allowbreak{} Chapter\allowbreak{} V,\allowbreak{}\allowbreak{} ms.\allowbreak{}\allowbreak{} 11:\allowbreak{}\allowbreak{} composite\allowbreak{} functor\allowbreak{} subscripts}{+}{,\allowbreak{}}{DER\_\allowbreak{}III-V:\allowbreak{}54}
\DERapparatusitem{ENV0255}{8;\allowbreak{}\allowbreak{} running\allowbreak{} text\allowbreak{} or\allowbreak{} formula}{\textbackslash{}\allowbreak{}text\{\allowbreak{}def\}\allowbreak{}}{\textbackslash{}\allowbreak{}mathrm\{\allowbreak{}def\}\allowbreak{}}{TYPESETTING:\allowbreak{}7;\allowbreak{}\allowbreak{} NOTATION\_\allowbreak{}DER:\allowbreak{}17}
\DERapparatusitem{ENV0256}{9;\allowbreak{}\allowbreak{} Chapter\allowbreak{} V,\allowbreak{}\allowbreak{} ms.\allowbreak{}\allowbreak{} 13:\allowbreak{}\allowbreak{} transition\allowbreak{} reduction}{1,\allowbreak{}n+1}{2,\allowbreak{}n)\allowbreak{}\$\allowbreak{}\allowbreak{} ou\allowbreak{} \$\allowbreak{}(\allowbreak{}m',\allowbreak{}n')\allowbreak{}=\allowbreak{}(\allowbreak{}m,\allowbreak{}n+2}{DER\_\allowbreak{}III-V:\allowbreak{}55}
\DERapparatusitem{ENV0257}{9;\allowbreak{}\allowbreak{} Chapter\allowbreak{} V,\allowbreak{}\allowbreak{} ms.\allowbreak{}\allowbreak{} 14:\allowbreak{}\allowbreak{} wrong\allowbreak{} shape\allowbreak{} indices}{a\allowbreak{} diagram\allowbreak{} of\allowbreak{} type\allowbreak{} \$\allowbreak{}\textbackslash{}\allowbreak{}Phi\_\allowbreak{}n\$\allowbreak{},\allowbreak{}\allowbreak{} but\allowbreak{} a\allowbreak{} diagram\allowbreak{} of\allowbreak{} \textbackslash{}\allowbreak{}textit\{\allowbreak{}dual\}\allowbreak{}\allowbreak{} type\allowbreak{} \$\allowbreak{}\textbackslash{}\allowbreak{}Phi\_\allowbreak{}n\textasciicircum{}\allowbreak{}o=\allowbreak{}\textbackslash{}\allowbreak{}Psi\_\allowbreak{}n\$\allowbreak{}}{a\allowbreak{} diagram\allowbreak{} of\allowbreak{} type\allowbreak{} \$\allowbreak{}\textbackslash{}\allowbreak{}Phi\_\allowbreak{}\{\allowbreak{}m+n\}\allowbreak{}\$\allowbreak{},\allowbreak{}\allowbreak{} but\allowbreak{} a\allowbreak{} diagram\allowbreak{} of\allowbreak{} \textbackslash{}\allowbreak{}textit\{\allowbreak{}dual\}\allowbreak{}\allowbreak{} type\allowbreak{} \$\allowbreak{}\textbackslash{}\allowbreak{}Phi\_\allowbreak{}\{\allowbreak{}m+n\}\allowbreak{}\textasciicircum{}\allowbreak{}o=\allowbreak{}\textbackslash{}\allowbreak{}Psi\_\allowbreak{}\{\allowbreak{}m+n\}\allowbreak{}\$\allowbreak{}}{DER\_\allowbreak{}III-V:\allowbreak{}56}
\DERapparatusitem{ENV0258}{9;\allowbreak{}\allowbreak{} Chapter\allowbreak{} V,\allowbreak{}\allowbreak{} ms.\allowbreak{}\allowbreak{} 14:\allowbreak{}\allowbreak{} unfinished\allowbreak{} target\allowbreak{} category}{\$\allowbreak{}\textbackslash{}\allowbreak{}Ch\_\allowbreak{}\{\allowbreak{}\textbackslash{}\allowbreak{}Sigma,\allowbreak{}n+2\}\allowbreak{}\$\allowbreak{}}{\$\allowbreak{}\textbackslash{}\allowbreak{}Ch\_\allowbreak{}\{\allowbreak{}\textbackslash{}\allowbreak{}Sigma,\allowbreak{}n+2\}\allowbreak{}(\allowbreak{}X;\allowbreak{}y,\allowbreak{}z)\allowbreak{}\$\allowbreak{}}{DER\_\allowbreak{}III-V:\allowbreak{}57}
\DERapparatusitem{ENV0259}{9;\allowbreak{}\allowbreak{} Chapter\allowbreak{} V,\allowbreak{}\allowbreak{} ms.\allowbreak{}\allowbreak{} 14:\allowbreak{}\allowbreak{} parenthetical\allowbreak{} parity\allowbreak{} variable}{n}{m}{DER\_\allowbreak{}III-V:\allowbreak{}58}
\DERapparatusitem{ENV0260}{10;\allowbreak{}\allowbreak{} Chapter\allowbreak{} V,\allowbreak{}\allowbreak{} ms.\allowbreak{}\allowbreak{} 15:\allowbreak{}\allowbreak{} transposed\allowbreak{} Hom\allowbreak{} diagram}{,\allowbreak{}}{+}{DER\_\allowbreak{}III-V:\allowbreak{}59}
\DERapparatusitem{ENV0261}{10;\allowbreak{}\allowbreak{} Chapter\allowbreak{} V,\allowbreak{}\allowbreak{} ms.\allowbreak{}\allowbreak{} 15:\allowbreak{}\allowbreak{} transposed\allowbreak{} Hom\allowbreak{} diagram}{,\allowbreak{}}{+}{DER\_\allowbreak{}III-V:\allowbreak{}59}
\DERapparatusitem{ENV0262}{10;\allowbreak{}\allowbreak{} Chapter\allowbreak{} V,\allowbreak{}\allowbreak{} ms.\allowbreak{}\allowbreak{} 15--16:\allowbreak{}\allowbreak{} ambient\allowbreak{} category\allowbreak{} of\allowbreak{} B}{+n}{'+n'}{DER\_\allowbreak{}III-V:\allowbreak{}60}
\DERapparatusitem{ENV0263}{10;\allowbreak{}\allowbreak{} Chapter\allowbreak{} V,\allowbreak{}\allowbreak{} ms.\allowbreak{}\allowbreak{} 16:\allowbreak{}\allowbreak{} values\allowbreak{} at\allowbreak{} the\allowbreak{} distinguished\allowbreak{} vertices}{y,\allowbreak{}\textbackslash{}\allowbreak{}\allowbreak{} c\textbackslash{}\allowbreak{}longrightarrow\allowbreak{} z}{z,\allowbreak{}\textbackslash{}\allowbreak{}\allowbreak{} c\textbackslash{}\allowbreak{}longrightarrow\allowbreak{} y}{DER\_\allowbreak{}III-V:\allowbreak{}62}
\DERapparatusitem{ENV0264}{11;\allowbreak{}\allowbreak{} Chapter\allowbreak{} V,\allowbreak{}\allowbreak{} ms.\allowbreak{}\allowbreak{} 16:\allowbreak{}\allowbreak{} reparametrization\allowbreak{} Hom\allowbreak{} order}{\textbackslash{}\allowbreak{}uHom(\allowbreak{}\textbackslash{}\allowbreak{}Phi\_\allowbreak{}\{\allowbreak{}m+n\}\allowbreak{},\allowbreak{}\textbackslash{}\allowbreak{}Phi\_\allowbreak{}\{\allowbreak{}m'+n'\}\allowbreak{};\allowbreak{}\textbackslash{}\allowbreak{}alpha)\allowbreak{}}{\textbackslash{}\allowbreak{}uHom(\allowbreak{}\textbackslash{}\allowbreak{}Phi\_\allowbreak{}\{\allowbreak{}m'+n'\}\allowbreak{},\allowbreak{}\textbackslash{}\allowbreak{}Phi\_\allowbreak{}\{\allowbreak{}m+n\}\allowbreak{};\allowbreak{}\textbackslash{}\allowbreak{}alpha')\allowbreak{}}{DER\_\allowbreak{}III-V:\allowbreak{}67}
\DERapparatusitem{ENV0265}{11;\allowbreak{}\allowbreak{} Chapter\allowbreak{} V,\allowbreak{}\allowbreak{} ms.\allowbreak{}\allowbreak{} 17:\allowbreak{}\allowbreak{} universal\allowbreak{} test\allowbreak{} variables}{for\allowbreak{} arbitrary\allowbreak{} \$\allowbreak{}X\$\allowbreak{},\allowbreak{}\allowbreak{} \$\allowbreak{}x\$\allowbreak{},\allowbreak{}\allowbreak{} \$\allowbreak{}y\$\allowbreak{},\allowbreak{}\allowbreak{} \$\allowbreak{}z\$\allowbreak{},\allowbreak{}}{for\allowbreak{} arbitrary\allowbreak{} \$\allowbreak{}X\$\allowbreak{},\allowbreak{}\allowbreak{} \$\allowbreak{}x\$\allowbreak{},\allowbreak{}\allowbreak{} \$\allowbreak{}z\$\allowbreak{},\allowbreak{}}{DER\_\allowbreak{}III-V:\allowbreak{}69}
\DERapparatusitem{ENV0266}{11;\allowbreak{}\allowbreak{} Chapter\allowbreak{} V,\allowbreak{}\allowbreak{} ms.\allowbreak{}\allowbreak{} 17:\allowbreak{}\allowbreak{} universal\allowbreak{} test\allowbreak{} variables}{on\allowbreak{} \$\allowbreak{}\textbackslash{}\allowbreak{}varphi\$\allowbreak{},\allowbreak{}\allowbreak{} \$\allowbreak{}\textbackslash{}\allowbreak{}psi\$\allowbreak{},\allowbreak{}\allowbreak{} \$\allowbreak{}\textbackslash{}\allowbreak{}alpha\$\allowbreak{}\allowbreak{} }{on\allowbreak{} \$\allowbreak{}\textbackslash{}\allowbreak{}varphi\$\allowbreak{},\allowbreak{}\allowbreak{} \$\allowbreak{}\textbackslash{}\allowbreak{}psi\$\allowbreak{},\allowbreak{}\allowbreak{} \$\allowbreak{}\textbackslash{}\allowbreak{}alpha'\$\allowbreak{}\allowbreak{} }{DER\_\allowbreak{}III-V:\allowbreak{}69}
\DERapparatusitem{ENV0267}{11;\allowbreak{}\allowbreak{} Chapter\allowbreak{} V,\allowbreak{}\allowbreak{} ms.\allowbreak{}\allowbreak{} 17:\allowbreak{}\allowbreak{} universal\allowbreak{} test\allowbreak{} variables}{\$\allowbreak{}x=\allowbreak{}a\$\allowbreak{},\allowbreak{}\allowbreak{} \$\allowbreak{}y=\allowbreak{}c\$\allowbreak{},\allowbreak{}\allowbreak{} \$\allowbreak{}z=\allowbreak{}b\$\allowbreak{}}{\$\allowbreak{}x=\allowbreak{}a\$\allowbreak{},\allowbreak{}\allowbreak{} \$\allowbreak{}z=\allowbreak{}b\$\allowbreak{}}{DER\_\allowbreak{}III-V:\allowbreak{}69}
\DERapparatusitem{ENV0268}{11;\allowbreak{}\allowbreak{} Chapter\allowbreak{} V,\allowbreak{}\allowbreak{} ms.\allowbreak{}\allowbreak{} 17:\allowbreak{}\allowbreak{} monotonicity\allowbreak{} omission}{\textbackslash{}\allowbreak{}uHom}{\textbackslash{}\allowbreak{}uHom\_\allowbreak{}\{\allowbreak{}cr\}\allowbreak{}}{DER\_\allowbreak{}III-V:\allowbreak{}70}
\DERapparatusitem{ENV0269}{12;\allowbreak{}\allowbreak{} Chapter\allowbreak{} V,\allowbreak{}\allowbreak{} ms.\allowbreak{}\allowbreak{} 18:\allowbreak{}\allowbreak{} unclosed\allowbreak{} parenthesis}{(\allowbreak{}independently\allowbreak{} of\allowbreak{} a\allowbreak{} \$\allowbreak{}\textbackslash{}\allowbreak{}varphi\$\allowbreak{}\allowbreak{} which\allowbreak{} would\allowbreak{} give\allowbreak{} rise\allowbreak{} to\allowbreak{} it\allowbreak{} ensure}{(\allowbreak{}independently\allowbreak{} of\allowbreak{} a\allowbreak{} \$\allowbreak{}\textbackslash{}\allowbreak{}varphi\$\allowbreak{}\allowbreak{} which\allowbreak{} would\allowbreak{} give\allowbreak{} rise\allowbreak{} to\allowbreak{} it)\allowbreak{}\allowbreak{} ensure}{DER\_\allowbreak{}III-V:\allowbreak{}76}
\DERapparatusitem{ENV0270}{12;\allowbreak{}\allowbreak{} Chapter\allowbreak{} V,\allowbreak{}\allowbreak{} ms.\allowbreak{}\allowbreak{} 19:\allowbreak{}\allowbreak{} terminal\allowbreak{} parity\allowbreak{} deduction}{because\allowbreak{} of\allowbreak{} a)\allowbreak{}\allowbreak{} \textbackslash{}\allowbreak{}texttt\{\allowbreak{}[\allowbreak{}rather\allowbreak{} a)\allowbreak{}\allowbreak{} and\allowbreak{} b)\allowbreak{}]\allowbreak{}\}\allowbreak{}}{because\allowbreak{} of\allowbreak{} a)\allowbreak{}\allowbreak{} and\allowbreak{} b)\allowbreak{}\allowbreak{} \textbackslash{}\allowbreak{}texttt\{\allowbreak{}[\allowbreak{}rather\allowbreak{} a)\allowbreak{}\allowbreak{} and\allowbreak{} b)\allowbreak{}]\allowbreak{}\}\allowbreak{}}{DER\_\allowbreak{}III-V:\allowbreak{}78}
\DERapparatusitem{ENV0271}{12;\allowbreak{}\allowbreak{} Chapter\allowbreak{} V,\allowbreak{}\allowbreak{} ms.\allowbreak{}\allowbreak{} 19:\allowbreak{}\allowbreak{} particular\allowbreak{} erased-edge\allowbreak{} set}{k-1,\allowbreak{}}{2k-1,\allowbreak{}2}{DER\_\allowbreak{}III-V:\allowbreak{}80}
\DERapparatusitem{ENV0272}{12;\allowbreak{}\allowbreak{} Chapter\allowbreak{} V,\allowbreak{}\allowbreak{} ms.\allowbreak{}\allowbreak{} 20:\allowbreak{}\allowbreak{} component\allowbreak{} count\allowbreak{} and\allowbreak{} closing\allowbreak{} parenthesis}{in\allowbreak{} one\allowbreak{} case,\allowbreak{}\allowbreak{} three}{in\allowbreak{} one\allowbreak{} case,\allowbreak{}\allowbreak{} two}{DER\_\allowbreak{}III-V:\allowbreak{}81}
\DERapparatusitem{ENV0273}{12;\allowbreak{}\allowbreak{} Chapter\allowbreak{} V,\allowbreak{}\allowbreak{} ms.\allowbreak{}\allowbreak{} 20:\allowbreak{}\allowbreak{} component\allowbreak{} count\allowbreak{} and\allowbreak{} closing\allowbreak{} parenthesis}{in\allowbreak{} the\allowbreak{} other.\allowbreak{})\allowbreak{}}{in\allowbreak{} the\allowbreak{} other.\allowbreak{}}{DER\_\allowbreak{}III-V:\allowbreak{}81}
\DERapparatusitem{ENV0274}{13;\allowbreak{}\allowbreak{} Chapter\allowbreak{} V,\allowbreak{}\allowbreak{} ms.\allowbreak{}\allowbreak{} 20:\allowbreak{}\allowbreak{} dual\allowbreak{} prefix/\allowbreak{}suffix\allowbreak{} lemma}{y}{u\_\allowbreak{}j}{DER\_\allowbreak{}III-V:\allowbreak{}82}
\DERapparatusitem{ENV0275}{13;\allowbreak{}\allowbreak{} Chapter\allowbreak{} V,\allowbreak{}\allowbreak{} ms.\allowbreak{}\allowbreak{} 20:\allowbreak{}\allowbreak{} cardinality\allowbreak{} of\allowbreak{} a\allowbreak{} block}{(\allowbreak{}i.\allowbreak{}e.\allowbreak{}\allowbreak{} \$\allowbreak{}s-r\$\allowbreak{}\allowbreak{} is\allowbreak{} even)\allowbreak{}}{(\allowbreak{}i.\allowbreak{}e.\allowbreak{}\allowbreak{} \$\allowbreak{}s-r\$\allowbreak{}\allowbreak{} is\allowbreak{} odd)\allowbreak{}}{DER\_\allowbreak{}III-V:\allowbreak{}83}
\DERapparatusitem{ENV0276}{14;\allowbreak{}\allowbreak{} Chapter\allowbreak{} V,\allowbreak{}\allowbreak{} ms.\allowbreak{}\allowbreak{} 22:\allowbreak{}\allowbreak{} components\allowbreak{} of\allowbreak{} path\allowbreak{} category}{\textbackslash{}\allowbreak{}mathrm\{\allowbreak{}Ch\}\allowbreak{}}{\textbackslash{}\allowbreak{}uCh}{DER\_\allowbreak{}III-V:\allowbreak{}88}
\DERapparatusitem{ENV0277}{15;\allowbreak{}\allowbreak{} Chapter\allowbreak{} V,\allowbreak{}\allowbreak{} ms.\allowbreak{}\allowbreak{} 23:\allowbreak{}\allowbreak{} path\allowbreak{} categories\allowbreak{} named\allowbreak{} Phi}{\textbackslash{}\allowbreak{}Phi\_\allowbreak{}2(\allowbreak{}}{\textbackslash{}\allowbreak{}uCh\_\allowbreak{}2(\allowbreak{}X;\allowbreak{}}{DER\_\allowbreak{}III-V:\allowbreak{}89}
\DERapparatusitem{ENV0278}{15;\allowbreak{}\allowbreak{} Chapter\allowbreak{} V,\allowbreak{}\allowbreak{} ms.\allowbreak{}\allowbreak{} 23:\allowbreak{}\allowbreak{} path\allowbreak{} categories\allowbreak{} named\allowbreak{} Phi}{\textbackslash{}\allowbreak{}Phi\_\allowbreak{}\{\allowbreak{}\textbackslash{}\allowbreak{}Sigma,\allowbreak{}0\}\allowbreak{}(\allowbreak{}}{\textbackslash{}\allowbreak{}uCh\_\allowbreak{}\{\allowbreak{}\textbackslash{}\allowbreak{}Sigma,\allowbreak{}0\}\allowbreak{}(\allowbreak{}X;\allowbreak{}}{DER\_\allowbreak{}III-V:\allowbreak{}89}
\DERapparatusitem{ENV0279}{15;\allowbreak{}\allowbreak{} Chapter\allowbreak{} V,\allowbreak{}\allowbreak{} ms.\allowbreak{}\allowbreak{} 24:\allowbreak{}\allowbreak{} gluing\allowbreak{} two\allowbreak{} arrows}{\textbackslash{}\allowbreak{}sqcup}{\textbackslash{}\allowbreak{}sqcup\_\allowbreak{}\{\allowbreak{}\textbackslash{}\allowbreak{}bDelta\_\allowbreak{}0\}\allowbreak{}}{DER\_\allowbreak{}III-V:\allowbreak{}90}
\DERapparatusitem{ENV0280}{16;\allowbreak{}\allowbreak{} Chapter\allowbreak{} V,\allowbreak{}\allowbreak{} ms.\allowbreak{}\allowbreak{} 25:\allowbreak{}\allowbreak{} adjacent\allowbreak{} interval\allowbreak{} vertices}{one\allowbreak{} arrow\allowbreak{} connecting\allowbreak{} \$\allowbreak{}i\$\allowbreak{},\allowbreak{}\allowbreak{} \$\allowbreak{}i+1\$\allowbreak{}}{one\allowbreak{} arrow\allowbreak{} connecting\allowbreak{} \$\allowbreak{}i-1\$\allowbreak{},\allowbreak{}\allowbreak{} \$\allowbreak{}i\$\allowbreak{}}{DER\_\allowbreak{}III-V:\allowbreak{}93}
\DERapparatusitem{ENV0281}{16;\allowbreak{}\allowbreak{} Chapter\allowbreak{} V,\allowbreak{}\allowbreak{} ms.\allowbreak{}\allowbreak{} 25:\allowbreak{}\allowbreak{} negative\allowbreak{} orientation\allowbreak{} sign}{i}{-}{DER\_\allowbreak{}III-V:\allowbreak{}94}
\DERapparatusitem{ENV0282}{17;\allowbreak{}\allowbreak{} Chapter\allowbreak{} V,\allowbreak{}\allowbreak{} ms.\allowbreak{}\allowbreak{} 27:\allowbreak{}\allowbreak{} last\allowbreak{} Hom\allowbreak{} display}{}{(\allowbreak{}x,\allowbreak{}y)\allowbreak{}}{DER\_\allowbreak{}III-V:\allowbreak{}98}
\DERapparatusitem{ENV0283}{18;\allowbreak{}\allowbreak{} V\allowbreak{} ms29\allowbreak{} boxed\allowbreak{} localization\allowbreak{} formula}{x,\allowbreak{}\allowbreak{} b\_\allowbreak{}x}{X,\allowbreak{}\allowbreak{} b\_\allowbreak{}X}{DER\_\allowbreak{}III-V:\allowbreak{}100}
\DERapparatusitem{ENV0284}{preamble;\allowbreak{}\allowbreak{} local\allowbreak{} macro\allowbreak{} definition}{\textbackslash{}\allowbreak{}text\{\allowbreak{}def\}\allowbreak{}}{\textbackslash{}\allowbreak{}mathrm\{\allowbreak{}def\}\allowbreak{}}{TYPESETTING:\allowbreak{}7;\allowbreak{}\allowbreak{} NOTATION\_\allowbreak{}DER:\allowbreak{}17}
\DERapparatusitem{ENV0285}{20;\allowbreak{}\allowbreak{} V\allowbreak{} ms33\allowbreak{} paragraph\allowbreak{} on\allowbreak{} lifts}{\textbackslash{}\allowbreak{}sigma\_\allowbreak{}\{\allowbreak{}\textbackslash{}\allowbreak{}cM\}\allowbreak{},\allowbreak{}\textbackslash{}\allowbreak{}beta\_\allowbreak{}\{\allowbreak{}\textbackslash{}\allowbreak{}cM\}\allowbreak{}}{F\_\allowbreak{}0,\allowbreak{}F\_\allowbreak{}1}{DER\_\allowbreak{}III-V:\allowbreak{}103}
\DERapparatusitem{ENV0286}{20;\allowbreak{}\allowbreak{} V\allowbreak{} ms33\allowbreak{} boxed\allowbreak{} endpoint\allowbreak{} isomorphism}{X}{\textbackslash{}\allowbreak{}cM(\allowbreak{}X)\allowbreak{}}{DER\_\allowbreak{}III-V:\allowbreak{}104}
\DERapparatusitem{ENV0287}{21;\allowbreak{}\allowbreak{} V\allowbreak{} ms34\allowbreak{} first\allowbreak{} paragraph}{K}{\textbackslash{}\allowbreak{}KK}{DER\_\allowbreak{}III-V:\allowbreak{}105}
\DERapparatusitem{ENV0288}{21;\allowbreak{}\allowbreak{} V\allowbreak{} ms35\allowbreak{} restriction\allowbreak{} display}{\textbackslash{}\allowbreak{}cM\_\allowbreak{}X}{\textbackslash{}\allowbreak{}cM(\allowbreak{}X)\allowbreak{}}{DER\_\allowbreak{}III-V:\allowbreak{}106}
\DERapparatusitem{ENV0289}{23;\allowbreak{}\allowbreak{} V\allowbreak{} ms37\allowbreak{} footnote6}{a\allowbreak{} retraction\allowbreak{} of\allowbreak{} \$\allowbreak{}\textbackslash{}\allowbreak{}psi\textbackslash{}\allowbreak{}varphi\allowbreak{} X\$\allowbreak{}\allowbreak{} onto\allowbreak{} \$\allowbreak{}X\$\allowbreak{}}{a\allowbreak{} retraction\allowbreak{} of\allowbreak{} \$\allowbreak{}\textbackslash{}\allowbreak{}psi\textbackslash{}\allowbreak{}varphi\allowbreak{} K\$\allowbreak{}\allowbreak{} onto\allowbreak{} \$\allowbreak{}K\$\allowbreak{}}{DER\_\allowbreak{}III-V:\allowbreak{}110}
\DERapparatusitem{ENV0290}{23;\allowbreak{}\allowbreak{} V\allowbreak{} ms38\allowbreak{} triangular\allowbreak{} Hom\allowbreak{} diagram\allowbreak{} upper\allowbreak{} right}{\textbackslash{}\allowbreak{}mathrm\{\allowbreak{}Hom\}\allowbreak{}\_\allowbreak{}C(\allowbreak{}\textbackslash{}\allowbreak{}varphi\allowbreak{} X,\allowbreak{}K')\allowbreak{}}{\textbackslash{}\allowbreak{}mathrm\{\allowbreak{}Hom\}\allowbreak{}\_\allowbreak{}\{\allowbreak{}C'\}\allowbreak{}(\allowbreak{}\textbackslash{}\allowbreak{}varphi\allowbreak{} X,\allowbreak{}K')\allowbreak{}}{DER\_\allowbreak{}III-V:\allowbreak{}111}
\DERapparatusitem{ENV0291}{24;\allowbreak{}\allowbreak{} V\allowbreak{} ms38\allowbreak{} implication\allowbreak{} chain\allowbreak{} explanation}{bearing\allowbreak{} in\allowbreak{} mind\allowbreak{} for\allowbreak{} (\allowbreak{}i)\allowbreak{}\allowbreak{} \$\allowbreak{}\textbackslash{}\allowbreak{}Longrightarrow\$\allowbreak{}\allowbreak{} (\allowbreak{}i)\allowbreak{}\allowbreak{} }{bearing\allowbreak{} in\allowbreak{} mind\allowbreak{} for\allowbreak{} (\allowbreak{}i)\allowbreak{}\allowbreak{} \$\allowbreak{}\textbackslash{}\allowbreak{}Longrightarrow\$\allowbreak{}\allowbreak{} (\allowbreak{}i\$\allowbreak{}'\$\allowbreak{})\allowbreak{}\allowbreak{} }{DER\_\allowbreak{}III-V:\allowbreak{}112}
\DERapparatusitem{ENV0292}{25;\allowbreak{}\allowbreak{} running\allowbreak{} text\allowbreak{} or\allowbreak{} formula}{\textbackslash{}\allowbreak{}text\{\allowbreak{}def\}\allowbreak{}}{\textbackslash{}\allowbreak{}mathrm\{\allowbreak{}def\}\allowbreak{}}{TYPESETTING:\allowbreak{}7;\allowbreak{}\allowbreak{} NOTATION\_\allowbreak{}DER:\allowbreak{}17}
\DERapparatusitem{ENV0293}{25;\allowbreak{}\allowbreak{} V\allowbreak{} ms39\allowbreak{} sentence\allowbreak{} about\allowbreak{} phi(\allowbreak{}q)\allowbreak{}}{i}{p}{DER\_\allowbreak{}III-V:\allowbreak{}113}
\DERapparatusitem{ENV0294}{25;\allowbreak{}\allowbreak{} V\allowbreak{} ms39\allowbreak{} proof\allowbreak{} conclusion}{'}{}{DER\_\allowbreak{}III-V:\allowbreak{}114}
\DERapparatusitem{ENV0295}{26;\allowbreak{}\allowbreak{} V\allowbreak{} ms40\allowbreak{} representability\allowbreak{} corollary\allowbreak{} conclusion}{}{'}{DER\_\allowbreak{}III-V:\allowbreak{}115}
\DERapparatusitem{ENV0296}{26;\allowbreak{}\allowbreak{} V\allowbreak{} ms40\allowbreak{} diagram\allowbreak{} of\allowbreak{} induced\allowbreak{} ind-functors}{}{leftarrow,\allowbreak{}}{DER\_\allowbreak{}III-V:\allowbreak{}116}
\DERapparatusitem{ENV0297}{26;\allowbreak{}\allowbreak{} running\allowbreak{} text\allowbreak{} or\allowbreak{} formula}{\textbackslash{}\allowbreak{}text\{\allowbreak{}def\}\allowbreak{}}{\textbackslash{}\allowbreak{}mathrm\{\allowbreak{}def\}\allowbreak{}}{TYPESETTING:\allowbreak{}7;\allowbreak{}\allowbreak{} NOTATION\_\allowbreak{}DER:\allowbreak{}17}
\DERapparatusitem{ENV0298}{27;\allowbreak{}\allowbreak{} V\allowbreak{} ms41\allowbreak{} Proposition2\allowbreak{} condition\allowbreak{} ii\allowbreak{} prime}{\$\allowbreak{}\textbackslash{}\allowbreak{}Ind\allowbreak{} C\$\allowbreak{}.\allowbreak{})\allowbreak{}}{\$\allowbreak{}\textbackslash{}\allowbreak{}Ind\allowbreak{} C\$\allowbreak{}.\allowbreak{}}{DER\_\allowbreak{}III-V:\allowbreak{}117}
\DERapparatusitem{ENV0299}{27;\allowbreak{}\allowbreak{} V\allowbreak{} ms41\allowbreak{} Proposition2\allowbreak{} condition\allowbreak{} ii\allowbreak{} double-prime}{C(\allowbreak{}X,\allowbreak{}K)\allowbreak{}\allowbreak{} \textbackslash{}\allowbreak{}longrightarrow\allowbreak{} \textbackslash{}\allowbreak{}Hom\_\allowbreak{}\{\allowbreak{}C'\}\allowbreak{}(\allowbreak{}\textbackslash{}\allowbreak{}varphi(\allowbreak{}X)\allowbreak{},\allowbreak{}\textbackslash{}\allowbreak{}varphi}{\{\allowbreak{}\textbackslash{}\allowbreak{}mathrm\{\allowbreak{}Ind\}\allowbreak{}(\allowbreak{}C)\allowbreak{}\}\allowbreak{}(\allowbreak{}X,\allowbreak{}K)\allowbreak{}\allowbreak{} \textbackslash{}\allowbreak{}longrightarrow\allowbreak{} \textbackslash{}\allowbreak{}Hom\_\allowbreak{}\{\allowbreak{}\textbackslash{}\allowbreak{}mathrm\{\allowbreak{}Ind\}\allowbreak{}(\allowbreak{}C')\allowbreak{}\}\allowbreak{}(\allowbreak{}\textbackslash{}\allowbreak{}Phi(\allowbreak{}X)\allowbreak{},\allowbreak{}\textbackslash{}\allowbreak{}Phi}{DER\_\allowbreak{}III-V:\allowbreak{}118}
\DERapparatusitem{ENV0300}{27;\allowbreak{}\allowbreak{} V\allowbreak{} ms42\allowbreak{} Proposition2\allowbreak{} proof}{and\allowbreak{} (\allowbreak{}i\$\allowbreak{}''\$\allowbreak{})\allowbreak{}\allowbreak{} (\allowbreak{}ii\$\allowbreak{}'\$\allowbreak{})\allowbreak{}.\allowbreak{}}{and\allowbreak{} (\allowbreak{}i\$\allowbreak{}''\$\allowbreak{})\allowbreak{}\allowbreak{} \$\allowbreak{}\textbackslash{}\allowbreak{}Longrightarrow\$\allowbreak{}\allowbreak{} (\allowbreak{}ii\$\allowbreak{}'\$\allowbreak{})\allowbreak{}.\allowbreak{}}{DER\_\allowbreak{}III-V:\allowbreak{}119}
\DERapparatusitem{ENV0301}{28;\allowbreak{}\allowbreak{} running\allowbreak{} text\allowbreak{} or\allowbreak{} formula}{\textbackslash{}\allowbreak{}text\{\allowbreak{}def\}\allowbreak{}}{\textbackslash{}\allowbreak{}mathrm\{\allowbreak{}def\}\allowbreak{}}{TYPESETTING:\allowbreak{}7;\allowbreak{}\allowbreak{} NOTATION\_\allowbreak{}DER:\allowbreak{}17}
\DERapparatusitem{ENV0302}{preamble;\allowbreak{}\allowbreak{} local\allowbreak{} macro\allowbreak{} definition}{\textbackslash{}\allowbreak{}text\{\allowbreak{}def\}\allowbreak{}}{\textbackslash{}\allowbreak{}mathrm\{\allowbreak{}def\}\allowbreak{}}{TYPESETTING:\allowbreak{}7;\allowbreak{}\allowbreak{} NOTATION\_\allowbreak{}DER:\allowbreak{}17}
\DERapparatusitem{ENV0303}{29;\allowbreak{}\allowbreak{} V\allowbreak{} ms45\allowbreak{} final\allowbreak{} display}{Y}{X}{DER\_\allowbreak{}III-V:\allowbreak{}126}
\subsection*{Chapter VI}
\DERapparatusitem{ENV0304}{preamble;\allowbreak{}\allowbreak{} local\allowbreak{} macro\allowbreak{} definition}{\textbackslash{}\allowbreak{}text\{\allowbreak{}def\}\allowbreak{}}{\textbackslash{}\allowbreak{}mathrm\{\allowbreak{}def\}\allowbreak{}}{TYPESETTING:\allowbreak{}7;\allowbreak{}\allowbreak{} NOTATION\_\allowbreak{}DER:\allowbreak{}17}
\DERapparatusitem{ENV0305}{8;\allowbreak{}\allowbreak{} running\allowbreak{} text\allowbreak{} or\allowbreak{} formula}{\textbackslash{}\allowbreak{}text\{\allowbreak{}def\}\allowbreak{}}{\textbackslash{}\allowbreak{}mathrm\{\allowbreak{}def\}\allowbreak{}}{TYPESETTING:\allowbreak{}7;\allowbreak{}\allowbreak{} NOTATION\_\allowbreak{}DER:\allowbreak{}17}
\DERapparatusitem{ENV0306}{17;\allowbreak{}\allowbreak{} ms\allowbreak{} 28,\allowbreak{}\allowbreak{} TeX\allowbreak{} line\allowbreak{} 728}{\textbackslash{}\allowbreak{}mathrm\{\allowbreak{}Cat/\allowbreak{}S\}\allowbreak{}}{\textbackslash{}\allowbreak{}mathrm\{\allowbreak{}Cat\}\allowbreak{}/\allowbreak{}S}{DER\_\allowbreak{}VI-a:\allowbreak{}12}
\DERapparatusitem{ENV0307}{19;\allowbreak{}\allowbreak{} ms\allowbreak{} 30-31:\allowbreak{}\allowbreak{} covariant\allowbreak{} square\allowbreak{} in\allowbreak{} a\allowbreak{} contravariant\allowbreak{} discussion}{X\_\allowbreak{}s\allowbreak{} \textbackslash{}\allowbreak{}arrow[\allowbreak{}r,\allowbreak{}\allowbreak{} "u\_\allowbreak{}*\textasciicircum{}\allowbreak{}X"]\allowbreak{}\allowbreak{} \textbackslash{}\allowbreak{}arrow[\allowbreak{}d,\allowbreak{}\allowbreak{} "i\_\allowbreak{}s"]\allowbreak{}\allowbreak{} \&\allowbreak{}\allowbreak{} X\_\allowbreak{}t\allowbreak{} \textbackslash{}\allowbreak{}arrow[\allowbreak{}d,\allowbreak{}\allowbreak{} "i\_\allowbreak{}t"]\allowbreak{}\allowbreak{} \textbackslash{}\allowbreak{}\textbackslash{}\allowbreak{}\space{}\allowbreak{}X/\allowbreak{}s\allowbreak{} \textbackslash{}\allowbreak{}arrow[\allowbreak{}r,\allowbreak{}\allowbreak{} "(\allowbreak{}\textbackslash{}\allowbreak{}mathrm\{\allowbreak{}id\}\allowbreak{}\_\allowbreak{}X)\allowbreak{}/\allowbreak{}u",\allowbreak{}\allowbreak{} "=\allowbreak{}\allowbreak{} u\_\allowbreak{}*\textasciicircum{}\allowbreak{}\{\allowbreak{}\textbackslash{}\allowbreak{}Phi(\allowbreak{}X)\allowbreak{}\}\allowbreak{}"']\allowbreak{}\allowbreak{} \textbackslash{}\allowbreak{}arrow[\allowbreak{}ur,\allowbreak{}\allowbreak{} bend\allowbreak{} left=\allowbreak{}30,\allowbreak{}\allowbreak{} shorten\allowbreak{} <=\allowbreak{}28pt,\allowbreak{}\allowbreak{} shorten\allowbreak{} >=\allowbreak{}12pt,\allowbreak{}\allowbreak{} "\textbackslash{}\allowbreak{}alpha"]\allowbreak{}\allowbreak{} \&\allowbreak{}\allowbreak{} X/\allowbreak{}t}{X\_\allowbreak{}s\allowbreak{} \textbackslash{}\allowbreak{}arrow[\allowbreak{}d,\allowbreak{}\allowbreak{} "i\_\allowbreak{}s"]\allowbreak{}\allowbreak{} \&\allowbreak{}\allowbreak{} X\_\allowbreak{}t\allowbreak{} \textbackslash{}\allowbreak{}arrow[\allowbreak{}l,\allowbreak{}\allowbreak{} "u\_\allowbreak{}X\textasciicircum{}\allowbreak{}*"']\allowbreak{}\allowbreak{} \textbackslash{}\allowbreak{}arrow[\allowbreak{}d,\allowbreak{}\allowbreak{} "i\_\allowbreak{}t"]\allowbreak{}\allowbreak{} \textbackslash{}\allowbreak{}\textbackslash{}\allowbreak{}\space{}\allowbreak{}s\textbackslash{}\allowbreak{}backslash\allowbreak{} X\allowbreak{} \&\allowbreak{}\allowbreak{} t\textbackslash{}\allowbreak{}backslash\allowbreak{} X\allowbreak{} \textbackslash{}\allowbreak{}arrow[\allowbreak{}l,\allowbreak{}\allowbreak{} "u\textbackslash{}\allowbreak{}backslash\textbackslash{}\allowbreak{}mathrm\{\allowbreak{}id\}\allowbreak{}\_\allowbreak{}X\allowbreak{} =\allowbreak{}\allowbreak{} u\_\allowbreak{}\{\allowbreak{}\textbackslash{}\allowbreak{}Phi(\allowbreak{}X)\allowbreak{}\}\allowbreak{}\textasciicircum{}\allowbreak{}*"']\allowbreak{}\space{}\allowbreak{}\textbackslash{}\allowbreak{}arrow[\allowbreak{}from=\allowbreak{}1-2,\allowbreak{}\allowbreak{} to=\allowbreak{}2-1,\allowbreak{}\allowbreak{} phantom,\allowbreak{}\allowbreak{} "\{\allowbreak{}\textbackslash{}\allowbreak{}Downarrow\textbackslash{}\allowbreak{}alpha\}\allowbreak{}"\allowbreak{} description]\allowbreak{}}{DER\_\allowbreak{}VI-a:\allowbreak{}45}
\DERapparatusitem{ENV0308}{19;\allowbreak{}\allowbreak{} ms\allowbreak{} 31,\allowbreak{}\allowbreak{} TeX\allowbreak{} line\allowbreak{} 793}{\textbackslash{}\allowbreak{}mathrm\{\allowbreak{}id\}\allowbreak{}\_\allowbreak{}\{\allowbreak{}\textbackslash{}\allowbreak{}uu\{\allowbreak{}Fib\}\allowbreak{}\_\allowbreak{}0\}\allowbreak{}\allowbreak{} S}{\textbackslash{}\allowbreak{}mathrm\{\allowbreak{}id\}\allowbreak{}\_\allowbreak{}\{\allowbreak{}\textbackslash{}\allowbreak{}uu\{\allowbreak{}Fib\}\allowbreak{}\_\allowbreak{}0\allowbreak{} S\}\allowbreak{}}{DER\_\allowbreak{}VI-a:\allowbreak{}14}
\DERapparatusitem{ENV0309}{20;\allowbreak{}\allowbreak{} ms\allowbreak{} 33:\allowbreak{}\allowbreak{} extra\allowbreak{} outer\allowbreak{} opposite\allowbreak{} in\allowbreak{} covariant\allowbreak{} formula}{\$\allowbreak{}(\allowbreak{}\textbackslash{}\allowbreak{}uu\{\allowbreak{}Fibsc\}\allowbreak{}(\allowbreak{}S\textasciicircum{}\allowbreak{}o)\allowbreak{})\allowbreak{}\textasciicircum{}\allowbreak{}o\$\allowbreak{}}{\$\allowbreak{}\textbackslash{}\allowbreak{}uu\{\allowbreak{}Fibsc\}\allowbreak{}(\allowbreak{}S\textasciicircum{}\allowbreak{}o)\allowbreak{}\$\allowbreak{}}{DER\_\allowbreak{}VI-a:\allowbreak{}40}
\DERapparatusitem{ENV0310}{20;\allowbreak{}\allowbreak{} ms\allowbreak{} 33:\allowbreak{}\allowbreak{} extra\allowbreak{} outer\allowbreak{} opposite\allowbreak{} in\allowbreak{} covariant\allowbreak{} formula}{\$\allowbreak{}(\allowbreak{}\textbackslash{}\allowbreak{}mathrm\{\allowbreak{}HOT\}\allowbreak{}\_\allowbreak{}W(\allowbreak{}S\textasciicircum{}\allowbreak{}o)\allowbreak{})\allowbreak{}\textasciicircum{}\allowbreak{}o\$\allowbreak{}}{\$\allowbreak{}\textbackslash{}\allowbreak{}mathrm\{\allowbreak{}HOT\}\allowbreak{}\_\allowbreak{}W(\allowbreak{}S\textasciicircum{}\allowbreak{}o)\allowbreak{}\$\allowbreak{}}{DER\_\allowbreak{}VI-a:\allowbreak{}40}
\DERapparatusitem{ENV0311}{20;\allowbreak{}\allowbreak{} ms\allowbreak{} 33:\allowbreak{}\allowbreak{} extra\allowbreak{} outer\allowbreak{} opposite\allowbreak{} in\allowbreak{} covariant\allowbreak{} formula}{\$\allowbreak{}\textbackslash{}\allowbreak{}mathrm\{\allowbreak{}HOT\}\allowbreak{}\_\allowbreak{}W(\allowbreak{}S\textasciicircum{}\allowbreak{}o)\allowbreak{}\textasciicircum{}\allowbreak{}o\$\allowbreak{}}{\$\allowbreak{}\textbackslash{}\allowbreak{}mathrm\{\allowbreak{}HOT\}\allowbreak{}\_\allowbreak{}W(\allowbreak{}S\textasciicircum{}\allowbreak{}o)\allowbreak{}\$\allowbreak{}}{DER\_\allowbreak{}VI-a:\allowbreak{}40}
\DERapparatusitem{ENV0312}{21;\allowbreak{}\allowbreak{} ms\allowbreak{} 36,\allowbreak{}\allowbreak{} TeX\allowbreak{} line\allowbreak{} 883}{which\allowbreak{} is\allowbreak{} quite\allowbreak{} remarkable.\allowbreak{}}{which\allowbreak{} is\allowbreak{} quite\allowbreak{} remarkable)\allowbreak{}.\allowbreak{}}{DER\_\allowbreak{}VI-a:\allowbreak{}16}
\DERapparatusitem{ENV0313}{22;\allowbreak{}\allowbreak{} ms\allowbreak{} 38:\allowbreak{}\allowbreak{} inference\allowbreak{} from\allowbreak{} one-sided\allowbreak{} inverse\allowbreak{} to\allowbreak{} fullness}{induces\allowbreak{} \textbackslash{}\allowbreak{}textit\{\allowbreak{}surjections\}\allowbreak{}\allowbreak{} on\allowbreak{} the\allowbreak{} \$\allowbreak{}\textbackslash{}\allowbreak{}mathrm\{\allowbreak{}Hom\}\allowbreak{}(\allowbreak{}-,\allowbreak{}-)\allowbreak{}\$\allowbreak{},\allowbreak{}\allowbreak{} while}{induces\allowbreak{} \textbackslash{}\allowbreak{}textit\{\allowbreak{}surjections\}\allowbreak{}\allowbreak{} on\allowbreak{} the\allowbreak{} \$\allowbreak{}\textbackslash{}\allowbreak{}mathrm\{\allowbreak{}Hom\}\allowbreak{}(\allowbreak{}-,\allowbreak{}-)\allowbreak{}\$\allowbreak{}\allowbreak{} between\allowbreak{} objects\allowbreak{} in\allowbreak{} the\allowbreak{} image\allowbreak{} of\allowbreak{} \$\allowbreak{}\textbackslash{}\allowbreak{}bar\{\allowbreak{}\textbackslash{}\allowbreak{}psi\}\allowbreak{}\$\allowbreak{},\allowbreak{}\allowbreak{} while}{DER\_\allowbreak{}VI-a:\allowbreak{}39}
\DERapparatusitem{ENV0314}{22;\allowbreak{}\allowbreak{} ms\allowbreak{} 38:\allowbreak{}\allowbreak{} inference\allowbreak{} from\allowbreak{} one-sided\allowbreak{} inverse\allowbreak{} to\allowbreak{} fullness}{i.\allowbreak{}e.\allowbreak{}\allowbreak{} \$\allowbreak{}\textbackslash{}\allowbreak{}bar\{\allowbreak{}k\}\allowbreak{}\$\allowbreak{}\allowbreak{} not\allowbreak{} even\allowbreak{} faithful}{which\allowbreak{} implies\allowbreak{} that\allowbreak{} \$\allowbreak{}\textbackslash{}\allowbreak{}bar\{\allowbreak{}k\}\allowbreak{}\$\allowbreak{}\allowbreak{} is\allowbreak{} not\allowbreak{} even\allowbreak{} faithful}{DER\_\allowbreak{}VI-a:\allowbreak{}39}
\DERapparatusitem{ENV0315}{preamble;\allowbreak{}\allowbreak{} local\allowbreak{} macro\allowbreak{} definition}{\textbackslash{}\allowbreak{}text\{\allowbreak{}def\}\allowbreak{}}{\textbackslash{}\allowbreak{}mathrm\{\allowbreak{}def\}\allowbreak{}}{TYPESETTING:\allowbreak{}7;\allowbreak{}\allowbreak{} NOTATION\_\allowbreak{}DER:\allowbreak{}17}
\DERapparatusitem{ENV0316}{30;\allowbreak{}\allowbreak{} ms\allowbreak{} 47:\allowbreak{}\allowbreak{} literal\allowbreak{} equality\allowbreak{} of\allowbreak{} coarse\allowbreak{} and\allowbreak{} fine\allowbreak{} weak\allowbreak{} equivalences}{of\allowbreak{} the\allowbreak{} diagram,\allowbreak{}\allowbreak{} it\allowbreak{} coincides\allowbreak{} with\allowbreak{} right\allowbreak{} fine\allowbreak{} \$\allowbreak{}W\$\allowbreak{}-equivalence}{of\allowbreak{} the\allowbreak{} diagram,\allowbreak{}\allowbreak{} fiberwise\allowbreak{} \$\allowbreak{}W\$\allowbreak{}-equivalence\allowbreak{} coincides\allowbreak{} with\allowbreak{} right\allowbreak{} fine\allowbreak{} \$\allowbreak{}W\$\allowbreak{}-equivalence}{DER\_\allowbreak{}VI-a:\allowbreak{}43}
\DERapparatusitem{ENV0317}{30;\allowbreak{}\allowbreak{} ms\allowbreak{} 47:\allowbreak{}\allowbreak{} literal\allowbreak{} equality\allowbreak{} of\allowbreak{} coarse\allowbreak{} and\allowbreak{} fine\allowbreak{} weak\allowbreak{} equivalences}{it\allowbreak{} coincides\allowbreak{} with\allowbreak{} left\allowbreak{} fine\allowbreak{} \$\allowbreak{}W\$\allowbreak{}-equivalence}{fiberwise\allowbreak{} \$\allowbreak{}W\$\allowbreak{}-equivalence\allowbreak{} coincides\allowbreak{} with\allowbreak{} left\allowbreak{} fine\allowbreak{} \$\allowbreak{}W\$\allowbreak{}-equivalence}{DER\_\allowbreak{}VI-a:\allowbreak{}43}
\DERapparatusitem{ENV0318}{31;\allowbreak{}\allowbreak{} ms\allowbreak{} 50:\allowbreak{}\allowbreak{} extra\allowbreak{} outer\allowbreak{} opposite\allowbreak{} in\allowbreak{} covariant\allowbreak{} category}{\textbackslash{}\allowbreak{}HOT\_\allowbreak{}W(\allowbreak{}S\textasciicircum{}\allowbreak{}o)\allowbreak{}\textasciicircum{}\allowbreak{}o}{\textbackslash{}\allowbreak{}HOT\_\allowbreak{}W(\allowbreak{}S\textasciicircum{}\allowbreak{}o)\allowbreak{}}{DER\_\allowbreak{}VI-a:\allowbreak{}44}
\DERapparatusitem{ENV0319}{32;\allowbreak{}\allowbreak{} ms\allowbreak{} 53,\allowbreak{}\allowbreak{} TeX\allowbreak{} line\allowbreak{} 1286}{\textbackslash{}\allowbreak{}psi}{\textbackslash{}\allowbreak{}Psi}{DER\_\allowbreak{}VI-a:\allowbreak{}27}
\DERapparatusitem{ENV0320}{34;\allowbreak{}\allowbreak{} ms\allowbreak{} 55,\allowbreak{}\allowbreak{} TeX\allowbreak{} line\allowbreak{} 1335}{X}{S}{DER\_\allowbreak{}VI-a:\allowbreak{}28}
\DERapparatusitem{ENV0321}{34;\allowbreak{}\allowbreak{} running\allowbreak{} text\allowbreak{} or\allowbreak{} formula}{\textbackslash{}\allowbreak{}text\{\allowbreak{}def\}\allowbreak{}}{\textbackslash{}\allowbreak{}mathrm\{\allowbreak{}def\}\allowbreak{}}{TYPESETTING:\allowbreak{}7;\allowbreak{}\allowbreak{} NOTATION\_\allowbreak{}DER:\allowbreak{}17}
\DERapparatusitem{ENV0322}{41;\allowbreak{}\allowbreak{} ms\allowbreak{} 68,\allowbreak{}\allowbreak{} TeX\allowbreak{} line\allowbreak{} 1645}{defined\allowbreak{} in\allowbreak{} \$\allowbreak{}A\textbackslash{}\allowbreak{}Sigma\textasciicircum{}\allowbreak{}\{\allowbreak{}-1\}\allowbreak{}\$\allowbreak{}}{defined\allowbreak{} in\allowbreak{} \$\allowbreak{}\textbackslash{}\allowbreak{}mathcal\{\allowbreak{}A\}\allowbreak{}\textbackslash{}\allowbreak{}Sigma\textasciicircum{}\allowbreak{}\{\allowbreak{}-1\}\allowbreak{}\$\allowbreak{}}{DER\_\allowbreak{}VI-a:\allowbreak{}34}
\DERapparatusitem{ENV0323}{42;\allowbreak{}\allowbreak{} ms\allowbreak{} 69,\allowbreak{}\allowbreak{} TeX\allowbreak{} line\allowbreak{} 1702}{\$\allowbreak{}C\_\allowbreak{}\{\allowbreak{}i(\allowbreak{}X)\allowbreak{}\}\allowbreak{}\$\allowbreak{}}{\$\allowbreak{}C(\allowbreak{}i\_\allowbreak{}X)\allowbreak{}\$\allowbreak{}}{DER\_\allowbreak{}VI-a:\allowbreak{}36}
\DERapparatusitem{ENV0324}{43;\allowbreak{}\allowbreak{} ms\allowbreak{} 69,\allowbreak{}\allowbreak{} TeX\allowbreak{} line\allowbreak{} 1711}{the\allowbreak{} functor\allowbreak{} C\allowbreak{} transforms}{the\allowbreak{} functor\allowbreak{} \$\allowbreak{}C\$\allowbreak{}\allowbreak{} transforms}{DER\_\allowbreak{}VI-a:\allowbreak{}37}
\DERapparatusitem{ENV0325}{46;\allowbreak{}\allowbreak{} ms74,\allowbreak{}\allowbreak{} first\allowbreak{} displayed\allowbreak{} dual\allowbreak{} diagram}{s\textbackslash{}\allowbreak{}longmapsto\textbackslash{}\allowbreak{}hot\_\allowbreak{}W(\allowbreak{}s\textbackslash{}\allowbreak{}backslash\allowbreak{} X)\allowbreak{}}{s\textbackslash{}\allowbreak{}longmapsto\textbackslash{}\allowbreak{}hot\_\allowbreak{}W(\allowbreak{}X/\allowbreak{}s)\allowbreak{}}{DER\_\allowbreak{}VI-b:\allowbreak{}8}
\DERapparatusitem{ENV0326}{50;\allowbreak{}\allowbreak{} running\allowbreak{} text\allowbreak{} or\allowbreak{} formula}{\textbackslash{}\allowbreak{}text\{\allowbreak{}def\}\allowbreak{}}{\textbackslash{}\allowbreak{}mathrm\{\allowbreak{}def\}\allowbreak{}}{TYPESETTING:\allowbreak{}7;\allowbreak{}\allowbreak{} NOTATION\_\allowbreak{}DER:\allowbreak{}17}
\DERapparatusitem{ENV0327}{50;\allowbreak{}\allowbreak{} running\allowbreak{} text\allowbreak{} or\allowbreak{} formula}{\textbackslash{}\allowbreak{}text\{\allowbreak{}def\}\allowbreak{}}{\textbackslash{}\allowbreak{}mathrm\{\allowbreak{}def\}\allowbreak{}}{TYPESETTING:\allowbreak{}7;\allowbreak{}\allowbreak{} NOTATION\_\allowbreak{}DER:\allowbreak{}17}
\DERapparatusitem{ENV0328}{50;\allowbreak{}\allowbreak{} running\allowbreak{} text\allowbreak{} or\allowbreak{} formula}{\textbackslash{}\allowbreak{}text\{\allowbreak{}def\}\allowbreak{}}{\textbackslash{}\allowbreak{}mathrm\{\allowbreak{}def\}\allowbreak{}}{TYPESETTING:\allowbreak{}7;\allowbreak{}\allowbreak{} NOTATION\_\allowbreak{}DER:\allowbreak{}17}
\DERapparatusitem{ENV0329}{52;\allowbreak{}\allowbreak{} ms.\allowbreak{}82,\allowbreak{}\allowbreak{} last\allowbreak{} complete\allowbreak{} paragraph\allowbreak{} of\allowbreak{} printed52,\allowbreak{}\allowbreak{} local/\allowbreak{}colocal\allowbreak{} tests}{local\allowbreak{} isomorphisms\allowbreak{} (\allowbreak{}to\allowbreak{} describe\allowbreak{} \$\allowbreak{}W\_\allowbreak{}S\textasciicircum{}\allowbreak{}\{\allowbreak{}\textbackslash{}\allowbreak{}mathrm\{\allowbreak{}d\}\allowbreak{}\}\allowbreak{}\$\allowbreak{})\allowbreak{},\allowbreak{}\allowbreak{} or\allowbreak{} that\allowbreak{} of\allowbreak{} colocal\allowbreak{} isomorphisms\allowbreak{} (\allowbreak{}to\allowbreak{} describe\allowbreak{} \$\allowbreak{}W\_\allowbreak{}S\textasciicircum{}\allowbreak{}\{\allowbreak{}\textbackslash{}\allowbreak{}mathrm\{\allowbreak{}g\}\allowbreak{}\}\allowbreak{}\$\allowbreak{})\allowbreak{}}{local\allowbreak{} isomorphisms\allowbreak{} (\allowbreak{}to\allowbreak{} describe\allowbreak{} \$\allowbreak{}W\_\allowbreak{}S\textasciicircum{}\allowbreak{}\{\allowbreak{}\textbackslash{}\allowbreak{}mathrm\{\allowbreak{}g\}\allowbreak{}\}\allowbreak{}\$\allowbreak{})\allowbreak{},\allowbreak{}\allowbreak{} or\allowbreak{} that\allowbreak{} of\allowbreak{} colocal\allowbreak{} isomorphisms\allowbreak{} (\allowbreak{}to\allowbreak{} describe\allowbreak{} \$\allowbreak{}W\_\allowbreak{}S\textasciicircum{}\allowbreak{}\{\allowbreak{}\textbackslash{}\allowbreak{}mathrm\{\allowbreak{}d\}\allowbreak{}\}\allowbreak{}\$\allowbreak{})\allowbreak{}}{DER\_\allowbreak{}VI-b:\allowbreak{}16}
\DERapparatusitem{ENV0330}{53;\allowbreak{}\allowbreak{} ms84,\allowbreak{}\allowbreak{} first\allowbreak{} sentence}{now\allowbreak{} convinces\allowbreak{} only\allowbreak{} that}{now\allowbreak{} convinces\allowbreak{} me\allowbreak{} that}{DER\_\allowbreak{}VI-b:\allowbreak{}18}
\DERapparatusitem{ENV0331}{54;\allowbreak{}\allowbreak{} ms86,\allowbreak{}\allowbreak{} Quillen\allowbreak{} lifting\allowbreak{} terminology}{the\allowbreak{} right\allowbreak{} lifting\space{}\allowbreak{}property\allowbreak{} with\allowbreak{} respect\allowbreak{} to\allowbreak{} \$\allowbreak{}F\$\allowbreak{}}{the\allowbreak{} left\allowbreak{} lifting\allowbreak{} property\allowbreak{} with\allowbreak{} respect\allowbreak{} to\allowbreak{} \$\allowbreak{}F\$\allowbreak{}}{DER\_\allowbreak{}VI-b:\allowbreak{}21}
\DERapparatusitem{ENV0332}{54;\allowbreak{}\allowbreak{} running\allowbreak{} text\allowbreak{} or\allowbreak{} formula}{\textbackslash{}\allowbreak{}text\{\allowbreak{}def\}\allowbreak{}}{\textbackslash{}\allowbreak{}mathrm\{\allowbreak{}def\}\allowbreak{}}{TYPESETTING:\allowbreak{}7;\allowbreak{}\allowbreak{} NOTATION\_\allowbreak{}DER:\allowbreak{}17}
\DERapparatusitem{ENV0333}{56;\allowbreak{}\allowbreak{} running\allowbreak{} text\allowbreak{} or\allowbreak{} formula}{\textbackslash{}\allowbreak{}text\{\allowbreak{}def\}\allowbreak{}}{\textbackslash{}\allowbreak{}mathrm\{\allowbreak{}def\}\allowbreak{}}{TYPESETTING:\allowbreak{}7;\allowbreak{}\allowbreak{} NOTATION\_\allowbreak{}DER:\allowbreak{}17}
\DERapparatusitem{ENV0334}{58;\allowbreak{}\allowbreak{} ms92,\allowbreak{}\allowbreak{} first\allowbreak{} displayed\allowbreak{} product}{\textbackslash{}\allowbreak{}times}{\textbackslash{}\allowbreak{}times\_\allowbreak{}S}{DER\_\allowbreak{}VI-b:\allowbreak{}25}
\DERapparatusitem{ENV0335}{preamble;\allowbreak{}\allowbreak{} local\allowbreak{} macro\allowbreak{} definition}{\textbackslash{}\allowbreak{}text\{\allowbreak{}def\}\allowbreak{}}{\textbackslash{}\allowbreak{}mathrm\{\allowbreak{}def\}\allowbreak{}}{TYPESETTING:\allowbreak{}7;\allowbreak{}\allowbreak{} NOTATION\_\allowbreak{}DER:\allowbreak{}17}
\DERapparatusitem{ENV0336}{59;\allowbreak{}\allowbreak{} running\allowbreak{} text\allowbreak{} or\allowbreak{} formula}{\textbackslash{}\allowbreak{}text\{\allowbreak{}def\}\allowbreak{}}{\textbackslash{}\allowbreak{}mathrm\{\allowbreak{}def\}\allowbreak{}}{TYPESETTING:\allowbreak{}7;\allowbreak{}\allowbreak{} NOTATION\_\allowbreak{}DER:\allowbreak{}17}
\DERapparatusitem{ENV0337}{59;\allowbreak{}\allowbreak{} ms95,\allowbreak{}\allowbreak{} target\allowbreak{} label\allowbreak{} of\allowbreak{} rho}{\textbackslash{}\allowbreak{}mathcal\{\allowbreak{}L\}\allowbreak{}}{}{DER\_\allowbreak{}VI-b:\allowbreak{}26}
\DERapparatusitem{ENV0338}{60;\allowbreak{}\allowbreak{} ms96,\allowbreak{}\allowbreak{} factorization\allowbreak{} of\allowbreak{} the\allowbreak{} base\allowbreak{} map\allowbreak{} g}{i\allowbreak{} \textbackslash{}\allowbreak{}in\allowbreak{} W,\allowbreak{}\textbackslash{}\allowbreak{}\allowbreak{} g\allowbreak{} \textbackslash{}\allowbreak{}in\allowbreak{} \textbackslash{}\allowbreak{}mathrm\{\allowbreak{}Fib\}\allowbreak{}}{i\allowbreak{} \textbackslash{}\allowbreak{}in\allowbreak{} W,\allowbreak{}\textbackslash{}\allowbreak{}\allowbreak{} \textbackslash{}\allowbreak{}bar\{\allowbreak{}g\}\allowbreak{}\allowbreak{} \textbackslash{}\allowbreak{}in\allowbreak{} \textbackslash{}\allowbreak{}mathrm\{\allowbreak{}Fib\}\allowbreak{}}{DER\_\allowbreak{}VI-b:\allowbreak{}28}
\DERapparatusitem{ENV0339}{63;\allowbreak{}\allowbreak{} ms100,\allowbreak{}\allowbreak{} parenthetical\allowbreak{} classification\allowbreak{} of\allowbreak{} the\allowbreak{} first\allowbreak{} square}{\textbackslash{}\allowbreak{}Fibu}{\textbackslash{}\allowbreak{}mathcal\allowbreak{} M/\allowbreak{}}{DER\_\allowbreak{}VI-b:\allowbreak{}29}
\DERapparatusitem{ENV0340}{64;\allowbreak{}\allowbreak{} ms102,\allowbreak{}\allowbreak{} formula\allowbreak{} for\allowbreak{} sigma\_\allowbreak{}S}{}{\_\allowbreak{}S}{DER\_\allowbreak{}VI-b:\allowbreak{}30}
\DERapparatusitem{ENV0341}{65;\allowbreak{}\allowbreak{} ms103,\allowbreak{}\allowbreak{} first\allowbreak{} displayed\allowbreak{} projection\allowbreak{} C(\allowbreak{}X)\allowbreak{}\allowbreak{} to\allowbreak{} X}{\textbackslash{}\allowbreak{}pi}{p}{DER\_\allowbreak{}VI-b:\allowbreak{}31}
\DERapparatusitem{ENV0342}{65;\allowbreak{}\allowbreak{} ms103,\allowbreak{}\allowbreak{} membership\allowbreak{} asserted\allowbreak{} for\allowbreak{} the\allowbreak{} section}{i\_\allowbreak{}X}{W}{DER\_\allowbreak{}VI-b:\allowbreak{}32}
\DERapparatusitem{ENV0343}{66;\allowbreak{}\allowbreak{} ms105,\allowbreak{}\allowbreak{} attribution\allowbreak{} of\allowbreak{} the\allowbreak{} two\allowbreak{} triples\allowbreak{} after\allowbreak{} respectivement}{the\allowbreak{} maps\allowbreak{} \$\allowbreak{}i\_\allowbreak{}\{\allowbreak{}C(\allowbreak{}X)\allowbreak{}\}\allowbreak{}\$\allowbreak{}\allowbreak{} and\allowbreak{} \$\allowbreak{}C(\allowbreak{}i\_\allowbreak{}X)\allowbreak{}\$\allowbreak{}}{the\allowbreak{} maps\allowbreak{} \$\allowbreak{}C(\allowbreak{}i\_\allowbreak{}X)\allowbreak{}\$\allowbreak{}\allowbreak{} and\allowbreak{} \$\allowbreak{}i\_\allowbreak{}\{\allowbreak{}C(\allowbreak{}X)\allowbreak{}\}\allowbreak{}\$\allowbreak{}}{DER\_\allowbreak{}VI-b:\allowbreak{}34}
\DERapparatusitem{ENV0344}{66;\allowbreak{}\allowbreak{} ms105,\allowbreak{}\allowbreak{} constant\allowbreak{} arrow\allowbreak{} in\allowbreak{} the\allowbreak{} first\allowbreak{} displayed\allowbreak{} path\allowbreak{} triple}{x}{f(\allowbreak{}x)\allowbreak{}}{DER\_\allowbreak{}VI-b:\allowbreak{}35}
\DERapparatusitem{ENV0345}{66;\allowbreak{}\allowbreak{} ms105,\allowbreak{}\allowbreak{} unit\allowbreak{} equation\allowbreak{} for\allowbreak{} a\allowbreak{} base\allowbreak{} path\allowbreak{} c}{c\textbackslash{}\allowbreak{}circ\allowbreak{} 1\_\allowbreak{}x=\allowbreak{}1\_\allowbreak{}s\textbackslash{}\allowbreak{}circ\allowbreak{} c=\allowbreak{}c}{c\textbackslash{}\allowbreak{}circ\allowbreak{} 1\_\allowbreak{}\{\allowbreak{}f(\allowbreak{}x)\allowbreak{}\}\allowbreak{}=\allowbreak{}1\_\allowbreak{}s\textbackslash{}\allowbreak{}circ\allowbreak{} c=\allowbreak{}c}{DER\_\allowbreak{}VI-b:\allowbreak{}36}
\DERapparatusitem{ENV0346}{67;\allowbreak{}\allowbreak{} ms106,\allowbreak{}\allowbreak{} internal-category\allowbreak{} paragraph\allowbreak{} immediately\allowbreak{} before\allowbreak{} the\allowbreak{} ms107\allowbreak{} marker}{}{X)\allowbreak{}(\allowbreak{}}{DER\_\allowbreak{}VI-b:\allowbreak{}37}
\DERapparatusitem{ENV0347}{67;\allowbreak{}\allowbreak{} ms108,\allowbreak{}\allowbreak{} footnote59}{These\allowbreak{} stability\allowbreak{} hypotheses\allowbreak{} on\allowbreak{} Fib\allowbreak{} will\allowbreak{} be\allowbreak{} used\allowbreak{} only\allowbreak{} in\allowbreak{} II,\allowbreak{}\allowbreak{} and\allowbreak{} even\allowbreak{} then\allowbreak{} only\allowbreak{} in\allowbreak{} case\allowbreak{} A)\allowbreak{},\allowbreak{}\allowbreak{} taken\allowbreak{} up\allowbreak{} again\allowbreak{} in\allowbreak{} toto\allowbreak{} in\allowbreak{} the\allowbreak{} form\allowbreak{} A\$\allowbreak{}'\$\allowbreak{}\allowbreak{} below\allowbreak{} (\allowbreak{}p.\allowbreak{}\textasciitilde{}112)\allowbreak{}.\allowbreak{}}{These\allowbreak{} stability\allowbreak{} hypotheses\allowbreak{} on\allowbreak{} Fib\allowbreak{} are\allowbreak{} also\allowbreak{} used\allowbreak{} in\allowbreak{} I\allowbreak{} under\allowbreak{} hypothesis\allowbreak{} B,\allowbreak{}\allowbreak{} pp.\allowbreak{}\textasciitilde{}92--93,\allowbreak{}\allowbreak{} and\allowbreak{} in\allowbreak{} II,\allowbreak{}\allowbreak{} in\allowbreak{} case\allowbreak{} A)\allowbreak{},\allowbreak{}\allowbreak{} taken\allowbreak{} up\allowbreak{} again\allowbreak{} in\allowbreak{} toto\allowbreak{} in\allowbreak{} the\allowbreak{} form\allowbreak{} A\$\allowbreak{}'\$\allowbreak{}\allowbreak{} below\allowbreak{} (\allowbreak{}p.\allowbreak{}\textasciitilde{}112)\allowbreak{}.\allowbreak{}}{DER\_\allowbreak{}VI-b:\allowbreak{}40}
\DERapparatusitem{ENV0348}{68;\allowbreak{}\allowbreak{} ms108,\allowbreak{}\allowbreak{} hypothesis\allowbreak{} B}{right\allowbreak{} lifting\allowbreak{} property\allowbreak{} \textbackslash{}\allowbreak{}texttt\{\allowbreak{}[\allowbreak{}rather\allowbreak{} left]\allowbreak{}\}\allowbreak{}}{left\allowbreak{} lifting\allowbreak{} property\allowbreak{} \textbackslash{}\allowbreak{}texttt\{\allowbreak{}[\allowbreak{}rather\allowbreak{} left]\allowbreak{}\}\allowbreak{}}{DER\_\allowbreak{}VI-b:\allowbreak{}41}
\DERapparatusitem{ENV0349}{69;\allowbreak{}\allowbreak{} running\allowbreak{} text\allowbreak{} or\allowbreak{} formula}{\textbackslash{}\allowbreak{}text\{\allowbreak{}def\}\allowbreak{}}{\textbackslash{}\allowbreak{}mathrm\{\allowbreak{}def\}\allowbreak{}}{TYPESETTING:\allowbreak{}7;\allowbreak{}\allowbreak{} NOTATION\_\allowbreak{}DER:\allowbreak{}17}
\DERapparatusitem{ENV0350}{70;\allowbreak{}\allowbreak{} ms111,\allowbreak{}\allowbreak{} sentence\allowbreak{} about\allowbreak{} images\allowbreak{} of\allowbreak{} W\_\allowbreak{}S\allowbreak{} under\allowbreak{} C}{o}{0}{DER\_\allowbreak{}VI-b:\allowbreak{}44}
\DERapparatusitem{ENV0351}{74;\allowbreak{}\allowbreak{} ms.\allowbreak{}120,\allowbreak{}\allowbreak{} sentence\allowbreak{} before\allowbreak{} Corollary2}{We\allowbreak{} have\allowbreak{} the\allowbreak{} choice\allowbreak{} between\allowbreak{} \$\allowbreak{}\textbackslash{}\allowbreak{}Phi\$\allowbreak{}\allowbreak{} and\allowbreak{} \$\allowbreak{}C\$\allowbreak{}\allowbreak{} to\allowbreak{} improve\allowbreak{} the\allowbreak{} situation!\allowbreak{} We\allowbreak{} know\allowbreak{} that\allowbreak{} \$\allowbreak{}X\$\allowbreak{}\allowbreak{} defines\allowbreak{} the\allowbreak{} same\allowbreak{} class\allowbreak{} in\allowbreak{} \$\allowbreak{}\textbackslash{}\allowbreak{}HOT\_\allowbreak{}W(\allowbreak{}S)\allowbreak{}\$\allowbreak{}\allowbreak{} as\allowbreak{} \$\allowbreak{}\textbackslash{}\allowbreak{}Phi(\allowbreak{}X)\allowbreak{}\$\allowbreak{},\allowbreak{}\allowbreak{} \$\allowbreak{}C(\allowbreak{}X)\allowbreak{}\$\allowbreak{},\allowbreak{}\allowbreak{} which\allowbreak{} are,\allowbreak{}\allowbreak{} the\allowbreak{} one\allowbreak{} smooth\allowbreak{} (\allowbreak{}and\allowbreak{} even\allowbreak{} split\allowbreak{} fibered)\allowbreak{},\allowbreak{}\allowbreak{} the\allowbreak{} other\allowbreak{} Hot-fibered\allowbreak{} (\allowbreak{}and\allowbreak{} even\allowbreak{} perfect)\allowbreak{}\allowbreak{} over\allowbreak{} \$\allowbreak{}S\$\allowbreak{}.\allowbreak{}}{We\allowbreak{} can\allowbreak{} use\allowbreak{} \$\allowbreak{}\textbackslash{}\allowbreak{}Phi\$\allowbreak{}\allowbreak{} to\allowbreak{} improve\allowbreak{} the\allowbreak{} situation!\allowbreak{} We\allowbreak{} know\allowbreak{} that\allowbreak{} \$\allowbreak{}X\$\allowbreak{}\allowbreak{} defines\allowbreak{} the\allowbreak{} same\allowbreak{} class\allowbreak{} in\allowbreak{} \$\allowbreak{}\textbackslash{}\allowbreak{}HOT\_\allowbreak{}W(\allowbreak{}S)\allowbreak{}\$\allowbreak{}\allowbreak{} as\allowbreak{} \$\allowbreak{}\textbackslash{}\allowbreak{}Phi(\allowbreak{}X)\allowbreak{}\$\allowbreak{},\allowbreak{}\allowbreak{} while\allowbreak{} \$\allowbreak{}C(\allowbreak{}X)\allowbreak{}\$\allowbreak{}\allowbreak{} represents\allowbreak{} the\allowbreak{} image\allowbreak{} of\allowbreak{} this\allowbreak{} class\allowbreak{} under\allowbreak{} \$\allowbreak{}i\textbackslash{}\allowbreak{}rho\$\allowbreak{}.\allowbreak{}\allowbreak{} These\allowbreak{} two\allowbreak{} objects\allowbreak{} are,\allowbreak{}\allowbreak{} the\allowbreak{} one\allowbreak{} smooth\allowbreak{} (\allowbreak{}and\allowbreak{} even\allowbreak{} split\allowbreak{} fibered)\allowbreak{},\allowbreak{}\allowbreak{} the\allowbreak{} other\allowbreak{} Hot-fibered\allowbreak{} (\allowbreak{}and\allowbreak{} even\allowbreak{} perfect)\allowbreak{}\allowbreak{} over\allowbreak{} \$\allowbreak{}S\$\allowbreak{}.\allowbreak{}}{DER\_\allowbreak{}VI-b:\allowbreak{}49}
\DERapparatusitem{ENV0352}{74;\allowbreak{}\allowbreak{} ms.\allowbreak{}120,\allowbreak{}\allowbreak{} Corollary2,\allowbreak{}\allowbreak{} second\allowbreak{} displayed\allowbreak{} isomorphism}{f\textasciicircum{}\allowbreak{}*(\allowbreak{}\textbackslash{}\allowbreak{}hot\_\allowbreak{}\{\allowbreak{}W,\allowbreak{}S\}\allowbreak{}(\allowbreak{}X)\allowbreak{})\allowbreak{}\allowbreak{} \textbackslash{}\allowbreak{}simeq\allowbreak{} \textbackslash{}\allowbreak{}hot\_\allowbreak{}\{\allowbreak{}W,\allowbreak{}S'\}\allowbreak{}(\allowbreak{}\textbackslash{}\allowbreak{}Phi(\allowbreak{}X)\allowbreak{}\allowbreak{} \textbackslash{}\allowbreak{}times\_\allowbreak{}S\allowbreak{} S')\allowbreak{}\allowbreak{} \textbackslash{}\allowbreak{}simeq\allowbreak{} \textbackslash{}\allowbreak{}hot\_\allowbreak{}\{\allowbreak{}W,\allowbreak{}S'\}\allowbreak{}(\allowbreak{}C(\allowbreak{}X)\allowbreak{}\allowbreak{} \textbackslash{}\allowbreak{}times\_\allowbreak{}S\allowbreak{} S')\allowbreak{}.\allowbreak{}}{f\textasciicircum{}\allowbreak{}*(\allowbreak{}\textbackslash{}\allowbreak{}hot\_\allowbreak{}\{\allowbreak{}W,\allowbreak{}S\}\allowbreak{}(\allowbreak{}X)\allowbreak{})\allowbreak{}\allowbreak{} \textbackslash{}\allowbreak{}simeq\allowbreak{} \textbackslash{}\allowbreak{}hot\_\allowbreak{}\{\allowbreak{}W,\allowbreak{}S'\}\allowbreak{}(\allowbreak{}\textbackslash{}\allowbreak{}Phi(\allowbreak{}X)\allowbreak{}\allowbreak{} \textbackslash{}\allowbreak{}times\_\allowbreak{}S\allowbreak{} S')\allowbreak{},\allowbreak{}\textbackslash{}\allowbreak{}qquad\allowbreak{} f\textasciicircum{}\allowbreak{}*i\textbackslash{}\allowbreak{}rho(\allowbreak{}\textbackslash{}\allowbreak{}hot\_\allowbreak{}\{\allowbreak{}W,\allowbreak{}S\}\allowbreak{}(\allowbreak{}X)\allowbreak{})\allowbreak{}\allowbreak{} \textbackslash{}\allowbreak{}simeq\allowbreak{} \textbackslash{}\allowbreak{}hot\_\allowbreak{}\{\allowbreak{}W,\allowbreak{}S'\}\allowbreak{}(\allowbreak{}C(\allowbreak{}X)\allowbreak{}\allowbreak{} \textbackslash{}\allowbreak{}times\_\allowbreak{}S\allowbreak{} S')\allowbreak{}.\allowbreak{}}{DER\_\allowbreak{}VI-b:\allowbreak{}50}
\DERapparatusitem{ENV0353}{75;\allowbreak{}\allowbreak{} running\allowbreak{} text\allowbreak{} or\allowbreak{} formula}{\textbackslash{}\allowbreak{}text\{\allowbreak{}def\}\allowbreak{}}{\textbackslash{}\allowbreak{}mathrm\{\allowbreak{}def\}\allowbreak{}}{TYPESETTING:\allowbreak{}7;\allowbreak{}\allowbreak{} NOTATION\_\allowbreak{}DER:\allowbreak{}17}
\DERapparatusitem{ENV0354}{75;\allowbreak{}\allowbreak{} running\allowbreak{} text\allowbreak{} or\allowbreak{} formula}{\textbackslash{}\allowbreak{}text\{\allowbreak{}def\}\allowbreak{}}{\textbackslash{}\allowbreak{}mathrm\{\allowbreak{}def\}\allowbreak{}}{TYPESETTING:\allowbreak{}7;\allowbreak{}\allowbreak{} NOTATION\_\allowbreak{}DER:\allowbreak{}17}
\DERapparatusitem{ENV0355}{75;\allowbreak{}\allowbreak{} ms.\allowbreak{}121,\allowbreak{}\allowbreak{} introduction\allowbreak{} to\allowbreak{} the\allowbreak{} list\allowbreak{} a–d}{one\allowbreak{} of\allowbreak{} the\allowbreak{} three\allowbreak{} types}{one\allowbreak{} of\allowbreak{} the\allowbreak{} four\allowbreak{} types}{DER\_\allowbreak{}VI-b:\allowbreak{}54}
\DERapparatusitem{ENV0356}{75;\allowbreak{}\allowbreak{} ms.\allowbreak{}121–122,\allowbreak{}\allowbreak{} sentence\allowbreak{} before\allowbreak{} Conjecture}{We\allowbreak{} can\allowbreak{} apply\allowbreak{} these\allowbreak{} results\allowbreak{} only\allowbreak{} to\allowbreak{} the\allowbreak{} derivator}{We\allowbreak{} cannot\allowbreak{} yet\allowbreak{} apply\allowbreak{} these\allowbreak{} results\allowbreak{} to\allowbreak{} the\allowbreak{} derivator}{DER\_\allowbreak{}VI-b:\allowbreak{}55}
\DERapparatusitem{ENV0357}{78;\allowbreak{}\allowbreak{} ms.\allowbreak{}127,\allowbreak{}\allowbreak{} definition\allowbreak{} of\allowbreak{} X-prime\allowbreak{} in\allowbreak{} the\allowbreak{} Hot-fibration\allowbreak{} paragraph}{\$\allowbreak{}X'\allowbreak{} =\allowbreak{}\allowbreak{} S\allowbreak{} \textbackslash{}\allowbreak{}times\_\allowbreak{}S\allowbreak{} S'\$\allowbreak{}}{\$\allowbreak{}X'\allowbreak{} =\allowbreak{}\allowbreak{} X\allowbreak{} \textbackslash{}\allowbreak{}times\_\allowbreak{}S\allowbreak{} S'\$\allowbreak{}}{DER\_\allowbreak{}VI-b:\allowbreak{}61}
\DERapparatusitem{ENV0358}{83;\allowbreak{}\allowbreak{} ms.\allowbreak{}136,\allowbreak{}\allowbreak{} first\allowbreak{} paragraph,\allowbreak{}\allowbreak{} definition\allowbreak{} by\allowbreak{} outgoing\allowbreak{} coslices}{\$\allowbreak{}\textbackslash{}\allowbreak{}overline\{\allowbreak{}S\}\allowbreak{}\textbackslash{}\allowbreak{}longrightarrow\allowbreak{} S\$\allowbreak{}\allowbreak{} is\allowbreak{} a\allowbreak{} local\allowbreak{} isomorphism}{\$\allowbreak{}\textbackslash{}\allowbreak{}overline\{\allowbreak{}S\}\allowbreak{}\textbackslash{}\allowbreak{}longrightarrow\allowbreak{} S\$\allowbreak{}\allowbreak{} is\allowbreak{} a\allowbreak{} local\allowbreak{} co-isomorphism}{DER\_\allowbreak{}VI-b:\allowbreak{}68}
\DERapparatusitem{ENV0359}{84;\allowbreak{}\allowbreak{} ms.\allowbreak{}139,\allowbreak{}\allowbreak{} square\allowbreak{} at\allowbreak{} the\allowbreak{} bottom\allowbreak{} of\allowbreak{} the\allowbreak{} leaf}{\&\allowbreak{}\allowbreak{} \textbackslash{}\allowbreak{}Phi\_\allowbreak{}S(\allowbreak{}X'/\allowbreak{}S')\allowbreak{}\allowbreak{} \textbackslash{}\allowbreak{}arrow[\allowbreak{}d]\allowbreak{}}{\&\allowbreak{}\allowbreak{} \textbackslash{}\allowbreak{}Phi\_\allowbreak{}S(\allowbreak{}X'/\allowbreak{}S)\allowbreak{}\allowbreak{} \textbackslash{}\allowbreak{}arrow[\allowbreak{}d]\allowbreak{}}{DER\_\allowbreak{}VI-b:\allowbreak{}69}
\DERapparatusitem{ENV0360}{88;\allowbreak{}\allowbreak{} ms.\allowbreak{}144,\allowbreak{}\allowbreak{} displayed\allowbreak{} comparison\allowbreak{} i\_\allowbreak{}X×id\_\allowbreak{}S′}{X\textbackslash{}\allowbreak{}times\_\allowbreak{}S\allowbreak{} S'\allowbreak{} =\allowbreak{}\allowbreak{} X'\allowbreak{} \textbackslash{}\allowbreak{}xrightarrow\{\allowbreak{}i\_\allowbreak{}X\textbackslash{}\allowbreak{}times\_\allowbreak{}S\textbackslash{}\allowbreak{}id\_\allowbreak{}\{\allowbreak{}S'\}\allowbreak{}\}\allowbreak{}\allowbreak{} \textbackslash{}\allowbreak{}Phi\_\allowbreak{}S(\allowbreak{}X')\allowbreak{}\textbackslash{}\allowbreak{}times\_\allowbreak{}S\allowbreak{} S'}{X\textbackslash{}\allowbreak{}times\_\allowbreak{}S\allowbreak{} S'\allowbreak{} =\allowbreak{}\allowbreak{} X'\allowbreak{} \textbackslash{}\allowbreak{}xrightarrow\{\allowbreak{}i\_\allowbreak{}X\textbackslash{}\allowbreak{}times\_\allowbreak{}S\textbackslash{}\allowbreak{}id\_\allowbreak{}\{\allowbreak{}S'\}\allowbreak{}\}\allowbreak{}\allowbreak{} \textbackslash{}\allowbreak{}Phi\_\allowbreak{}S(\allowbreak{}X)\allowbreak{}\textbackslash{}\allowbreak{}times\_\allowbreak{}S\allowbreak{} S'}{DER\_\allowbreak{}VI-b:\allowbreak{}72}
\DERapparatusitem{ENV0361}{89;\allowbreak{}\allowbreak{} ms145;\allowbreak{}\allowbreak{} line280;\allowbreak{}\allowbreak{} middle\allowbreak{} vertical\allowbreak{} arrow}{}{\textasciicircum{}\allowbreak{}}{DER\_\allowbreak{}VI-c:\allowbreak{}1}
\DERapparatusitem{ENV0362}{91;\allowbreak{}\allowbreak{} ms148-149;\allowbreak{}\allowbreak{} square(\allowbreak{}I)\allowbreak{},\allowbreak{}\allowbreak{} bottom\allowbreak{} arrow}{\textbackslash{}\allowbreak{}cA\allowbreak{} \textbackslash{}\allowbreak{}arrow[\allowbreak{}r]\allowbreak{}\allowbreak{} \&\allowbreak{}\allowbreak{} \textbackslash{}\allowbreak{}cA\_\allowbreak{}0.\allowbreak{}}{\textbackslash{}\allowbreak{}cA\allowbreak{} \textbackslash{}\allowbreak{}arrow[\allowbreak{}r,\allowbreak{}\allowbreak{} "\textbackslash{}\allowbreak{}rho"]\allowbreak{}\allowbreak{} \&\allowbreak{}\allowbreak{} \textbackslash{}\allowbreak{}cA\_\allowbreak{}0.\allowbreak{}}{DER\_\allowbreak{}VI-c:\allowbreak{}2}
\DERapparatusitem{ENV0363}{91;\allowbreak{}\allowbreak{} ms149;\allowbreak{}\allowbreak{} adjunction\allowbreak{} display\allowbreak{} after\allowbreak{} square(\allowbreak{}II)\allowbreak{}}{\textbackslash{}\allowbreak{}cA\_\allowbreak{}0\allowbreak{} \textbackslash{}\allowbreak{}adj\{\allowbreak{}f\_\allowbreak{}!\textasciicircum{}\allowbreak{}0\}\allowbreak{}\{\allowbreak{}f\_\allowbreak{}0\textasciicircum{}\allowbreak{}*\}\allowbreak{}\allowbreak{} \textbackslash{}\allowbreak{}cA'\_\allowbreak{}0,\allowbreak{}}{\textbackslash{}\allowbreak{}cA'\_\allowbreak{}0\allowbreak{} \textbackslash{}\allowbreak{}adj\{\allowbreak{}f\_\allowbreak{}!\textasciicircum{}\allowbreak{}0\}\allowbreak{}\{\allowbreak{}f\_\allowbreak{}0\textasciicircum{}\allowbreak{}*\}\allowbreak{}\allowbreak{} \textbackslash{}\allowbreak{}cA\_\allowbreak{}0,\allowbreak{}}{DER\_\allowbreak{}VI-c:\allowbreak{}3}
\DERapparatusitem{ENV0364}{96;\allowbreak{}\allowbreak{} ms157;\allowbreak{}\allowbreak{} introductory\allowbreak{} wording\allowbreak{} of\allowbreak{} first\allowbreak{} comma\allowbreak{} square}{relative\allowbreak{} to\allowbreak{} the\allowbreak{} commutative\allowbreak{} diagram\allowbreak{} in\allowbreak{} Cat}{relative\allowbreak{} to\allowbreak{} the\allowbreak{} diagram\allowbreak{} in\allowbreak{} Cat\allowbreak{} equipped\allowbreak{} with\allowbreak{} the\allowbreak{} natural\allowbreak{} transformation\allowbreak{} \$\allowbreak{}pf'\textbackslash{}\allowbreak{}Longrightarrow\allowbreak{} fq\$\allowbreak{}}{DER\_\allowbreak{}VI-c:\allowbreak{}9}
\DERapparatusitem{ENV0365}{96;\allowbreak{}\allowbreak{} ms157;\allowbreak{}\allowbreak{} displayed\allowbreak{} base-change\allowbreak{} arrow}{p\textasciicircum{}\allowbreak{}*(\allowbreak{}f\_\allowbreak{}!(\allowbreak{}\textbackslash{}\allowbreak{}xi')\allowbreak{})\allowbreak{}\textbackslash{}\allowbreak{}longrightarrow\allowbreak{} f'\_\allowbreak{}!(\allowbreak{}q\textasciicircum{}\allowbreak{}*(\allowbreak{}\textbackslash{}\allowbreak{}xi')\allowbreak{})\allowbreak{}}{f'\_\allowbreak{}!(\allowbreak{}q\textasciicircum{}\allowbreak{}*(\allowbreak{}\textbackslash{}\allowbreak{}xi')\allowbreak{})\allowbreak{}\textbackslash{}\allowbreak{}longrightarrow\allowbreak{} p\textasciicircum{}\allowbreak{}*(\allowbreak{}f\_\allowbreak{}!(\allowbreak{}\textbackslash{}\allowbreak{}xi')\allowbreak{})\allowbreak{}}{DER\_\allowbreak{}VI-c:\allowbreak{}8}
\DERapparatusitem{ENV0366}{97;\allowbreak{}\allowbreak{} ms158;\allowbreak{}\allowbreak{} line613;\allowbreak{}\allowbreak{} first\allowbreak{} display}{S}{e}{DER\_\allowbreak{}VI-c:\allowbreak{}10}
\DERapparatusitem{ENV0367}{97;\allowbreak{}\allowbreak{} running\allowbreak{} text\allowbreak{} or\allowbreak{} formula}{\textbackslash{}\allowbreak{}text\{\allowbreak{}def\}\allowbreak{}}{\textbackslash{}\allowbreak{}mathrm\{\allowbreak{}def\}\allowbreak{}}{TYPESETTING:\allowbreak{}7;\allowbreak{}\allowbreak{} NOTATION\_\allowbreak{}DER:\allowbreak{}17}
\DERapparatusitem{ENV0368}{97;\allowbreak{}\allowbreak{} ms158;\allowbreak{}\allowbreak{} line625;\allowbreak{}\allowbreak{} constant-coefficient\allowbreak{} comparison}{\textbackslash{}\allowbreak{}times\_\allowbreak{}S\textbackslash{}\allowbreak{}xi\textbackslash{}\allowbreak{}longrightarrow\textbackslash{}\allowbreak{}eta\textbackslash{}\allowbreak{}times\_\allowbreak{}S}{\textbackslash{}\allowbreak{}times\textbackslash{}\allowbreak{}xi\textbackslash{}\allowbreak{}longrightarrow\textbackslash{}\allowbreak{}eta\textbackslash{}\allowbreak{}times}{DER\_\allowbreak{}VI-c:\allowbreak{}11}
\DERapparatusitem{ENV0369}{97;\allowbreak{}\allowbreak{} running\allowbreak{} text\allowbreak{} or\allowbreak{} formula}{\textbackslash{}\allowbreak{}text\{\allowbreak{}def\}\allowbreak{}}{\textbackslash{}\allowbreak{}mathrm\{\allowbreak{}def\}\allowbreak{}}{TYPESETTING:\allowbreak{}7;\allowbreak{}\allowbreak{} NOTATION\_\allowbreak{}DER:\allowbreak{}17}
\DERapparatusitem{ENV0370}{97;\allowbreak{}\allowbreak{} running\allowbreak{} text\allowbreak{} or\allowbreak{} formula}{\textbackslash{}\allowbreak{}text\{\allowbreak{}def\}\allowbreak{}}{\textbackslash{}\allowbreak{}mathrm\{\allowbreak{}def\}\allowbreak{}}{TYPESETTING:\allowbreak{}7;\allowbreak{}\allowbreak{} NOTATION\_\allowbreak{}DER:\allowbreak{}17}
\DERapparatusitem{ENV0371}{98;\allowbreak{}\allowbreak{} printed\allowbreak{} p.\allowbreak{}98;\allowbreak{}\allowbreak{} ms160\allowbreak{} item1}{\textbackslash{}\allowbreak{}HOT}{\textbackslash{}\allowbreak{}HOT\textasciicircum{}\allowbreak{}\{\allowbreak{}\textbackslash{}\allowbreak{}mathrm\{\allowbreak{}lc\}\allowbreak{}\}\allowbreak{}}{DER\_\allowbreak{}VI-c:\allowbreak{}13}
\DERapparatusitem{ENV0372}{98;\allowbreak{}\allowbreak{} printed\allowbreak{} p.\allowbreak{}98;\allowbreak{}\allowbreak{} ms161\allowbreak{} item2}{\textbackslash{}\allowbreak{}HOT}{\textbackslash{}\allowbreak{}HOT\textasciicircum{}\allowbreak{}\{\allowbreak{}\textbackslash{}\allowbreak{}mathrm\{\allowbreak{}lc\}\allowbreak{}\}\allowbreak{}}{DER\_\allowbreak{}VI-c:\allowbreak{}14}
\DERapparatusitem{ENV0373}{98;\allowbreak{}\allowbreak{} running\allowbreak{} text\allowbreak{} or\allowbreak{} formula}{\textbackslash{}\allowbreak{}text\{\allowbreak{}def\}\allowbreak{}}{\textbackslash{}\allowbreak{}mathrm\{\allowbreak{}def\}\allowbreak{}}{TYPESETTING:\allowbreak{}7;\allowbreak{}\allowbreak{} NOTATION\_\allowbreak{}DER:\allowbreak{}17}
\DERapparatusitem{ENV0374}{98;\allowbreak{}\allowbreak{} printed\allowbreak{} p.\allowbreak{}98;\allowbreak{}\allowbreak{} ms161\allowbreak{} parenthesis\allowbreak{} defining\allowbreak{} pullback}{\{\allowbreak{}=\allowbreak{}\}\allowbreak{}X\textbackslash{}\allowbreak{}times\_\allowbreak{}S\allowbreak{} S'\$\allowbreak{}}{\{\allowbreak{}=\allowbreak{}\}\allowbreak{}X'\textbackslash{}\allowbreak{}times\_\allowbreak{}S\allowbreak{} S'\$\allowbreak{}}{DER\_\allowbreak{}VI-c:\allowbreak{}15}
\DERapparatusitem{ENV0375}{98;\allowbreak{}\allowbreak{} printed\allowbreak{} pp.\allowbreak{}98–99;\allowbreak{}\allowbreak{} ms161\allowbreak{} unit\allowbreak{} arrow\allowbreak{} and\allowbreak{} its\allowbreak{} diagram}{\textbackslash{}\allowbreak{}beta}{\textbackslash{}\allowbreak{}alpha}{DER\_\allowbreak{}VI-c:\allowbreak{}16}
\DERapparatusitem{ENV0376}{100;\allowbreak{}\allowbreak{} running\allowbreak{} text\allowbreak{} or\allowbreak{} formula}{\allowbreak{} (\allowbreak{}gear)\allowbreak{}}{}{GLOSSARY\_\allowbreak{}DER:\allowbreak{}164}
\DERapparatusitem{ENV0377}{101;\allowbreak{}\allowbreak{} printed\allowbreak{} p.\allowbreak{}101;\allowbreak{}\allowbreak{} ms164\allowbreak{} displayed\allowbreak{} Hom\allowbreak{} map\allowbreak{} after\allowbreak{} adjunction}{\textbackslash{}\allowbreak{}HOT\_\allowbreak{}W(\allowbreak{}S)\allowbreak{}\}\allowbreak{}(\allowbreak{}p\_\allowbreak{}\{\allowbreak{}S!\}\allowbreak{}(\allowbreak{}a\_\allowbreak{}S)\allowbreak{},\allowbreak{}\allowbreak{} b)\allowbreak{}\allowbreak{} \textbackslash{}\allowbreak{}longrightarrow\allowbreak{} \textbackslash{}\allowbreak{}Hom\_\allowbreak{}\{\allowbreak{}\textbackslash{}\allowbreak{}HOT\_\allowbreak{}W(\allowbreak{}S')\allowbreak{}}{\textbackslash{}\allowbreak{}Hot\_\allowbreak{}W\}\allowbreak{}(\allowbreak{}p\_\allowbreak{}\{\allowbreak{}S!\}\allowbreak{}(\allowbreak{}a\_\allowbreak{}S)\allowbreak{},\allowbreak{}b)\allowbreak{}\textbackslash{}\allowbreak{}longrightarrow\textbackslash{}\allowbreak{}Hom\_\allowbreak{}\{\allowbreak{}\textbackslash{}\allowbreak{}Hot\_\allowbreak{}W}{DER\_\allowbreak{}VI-c:\allowbreak{}19}
\DERapparatusitem{ENV0378}{101;\allowbreak{}\allowbreak{} running\allowbreak{} text\allowbreak{} or\allowbreak{} formula}{\textbackslash{}\allowbreak{}text\{\allowbreak{}def\}\allowbreak{}}{\textbackslash{}\allowbreak{}mathrm\{\allowbreak{}def\}\allowbreak{}}{TYPESETTING:\allowbreak{}7;\allowbreak{}\allowbreak{} NOTATION\_\allowbreak{}DER:\allowbreak{}17}
\DERapparatusitem{ENV0379}{111;\allowbreak{}\allowbreak{} printed\allowbreak{} p.\allowbreak{}111;\allowbreak{}\allowbreak{} ms181\allowbreak{} paragraph\allowbreak{} before\allowbreak{} section4}{for\allowbreak{} lack\allowbreak{} of\allowbreak{} axim\allowbreak{} W(\allowbreak{}4b)\allowbreak{}}{for\allowbreak{} lack\allowbreak{} of\allowbreak{} axiom\allowbreak{} W(\allowbreak{}4b)\allowbreak{}}{DER\_\allowbreak{}VI-c:\allowbreak{}27}
\DERapparatusitem{ENV0380}{112;\allowbreak{}\allowbreak{} running\allowbreak{} text\allowbreak{} or\allowbreak{} formula}{\textbackslash{}\allowbreak{}text\{\allowbreak{}def\}\allowbreak{}}{\textbackslash{}\allowbreak{}mathrm\{\allowbreak{}def\}\allowbreak{}}{TYPESETTING:\allowbreak{}7;\allowbreak{}\allowbreak{} NOTATION\_\allowbreak{}DER:\allowbreak{}17}
\DERapparatusitem{ENV0381}{117;\allowbreak{}\allowbreak{} printed117;\allowbreak{}\allowbreak{} manuscript190;\allowbreak{}\allowbreak{} paragraph\allowbreak{} constructing\allowbreak{} the\allowbreak{} retraction}{Now\allowbreak{} \$\allowbreak{}X''\$\allowbreak{}\allowbreak{} is\allowbreak{} obtained\allowbreak{} from\allowbreak{} \$\allowbreak{}X\$\allowbreak{}\allowbreak{} by\allowbreak{} adding\allowbreak{} a\allowbreak{} single\allowbreak{} element}{Now\allowbreak{} \$\allowbreak{}X\$\allowbreak{}\allowbreak{} is\allowbreak{} obtained\allowbreak{} from\allowbreak{} \$\allowbreak{}X''\$\allowbreak{}\allowbreak{} by\allowbreak{} adding\allowbreak{} a\allowbreak{} single\allowbreak{} element}{DER\_\allowbreak{}VI-c:\allowbreak{}34}
\DERapparatusitem{ENV0382}{117;\allowbreak{}\allowbreak{} printed117;\allowbreak{}\allowbreak{} manuscript190;\allowbreak{}\allowbreak{} final\allowbreak{} parenthetical\allowbreak{} cohomology\allowbreak{} comparison}{hence\allowbreak{} that\allowbreak{} of\allowbreak{} \$\allowbreak{}X\_\allowbreak{}0\$\allowbreak{}}{hence\allowbreak{} that\allowbreak{} of\allowbreak{} \$\allowbreak{}X\_\allowbreak{}a\$\allowbreak{}}{DER\_\allowbreak{}VI-c:\allowbreak{}33}
\DERapparatusitem{ENV0383}{118;\allowbreak{}\allowbreak{} printed118;\allowbreak{}\allowbreak{} manuscript191;\allowbreak{}\allowbreak{} centre\allowbreak{} of\allowbreak{} the\allowbreak{} contravariant\allowbreak{} fibre\allowbreak{} diagram\allowbreak{} of\allowbreak{} Phi(\allowbreak{}X)\allowbreak{}}{X\_\allowbreak{}0\allowbreak{} \&\allowbreak{}\allowbreak{} \textbackslash{}\allowbreak{}\textbackslash{}\allowbreak{}}{X\allowbreak{} \&\allowbreak{}\allowbreak{} \textbackslash{}\allowbreak{}\textbackslash{}\allowbreak{}}{DER\_\allowbreak{}VI-c:\allowbreak{}35}
\DERapparatusitem{ENV0384}{119;\allowbreak{}\allowbreak{} printed119;\allowbreak{}\allowbreak{} manuscript193;\allowbreak{}\allowbreak{} closure\allowbreak{} statement\allowbreak{} after\allowbreak{} W(\allowbreak{}6bis)\allowbreak{}}{a\allowbreak{} filtering\allowbreak{} direct\allowbreak{} limit\allowbreak{} of\allowbreak{} morphisms}{a\allowbreak{} sequential\allowbreak{} filtering\allowbreak{} direct\allowbreak{} limit\allowbreak{} of\allowbreak{} morphisms}{DER\_\allowbreak{}VI-c:\allowbreak{}36}
\DERapparatusitem{ENV0385}{120;\allowbreak{}\allowbreak{} printed120;\allowbreak{}\allowbreak{} manuscript195;\allowbreak{}\allowbreak{} formula\allowbreak{} for\allowbreak{} the\allowbreak{} minimal\allowbreak{} sequential\allowbreak{} localizer}{i\_\allowbreak{}x\allowbreak{} \textbackslash{}\allowbreak{}in\allowbreak{} W''\_\allowbreak{}\textbackslash{}\allowbreak{}omega}{i\_\allowbreak{}X\allowbreak{} \textbackslash{}\allowbreak{}in\allowbreak{} W''\_\allowbreak{}\textbackslash{}\allowbreak{}omega}{DER\_\allowbreak{}VI-c:\allowbreak{}38}
\DERapparatusitem{ENV0386}{120;\allowbreak{}\allowbreak{} running\allowbreak{} text\allowbreak{} or\allowbreak{} formula}{\textbackslash{}\allowbreak{}text\{\allowbreak{}def\}\allowbreak{}}{\textbackslash{}\allowbreak{}mathrm\{\allowbreak{}def\}\allowbreak{}}{TYPESETTING:\allowbreak{}7;\allowbreak{}\allowbreak{} NOTATION\_\allowbreak{}DER:\allowbreak{}17}
\DERapparatusitem{ENV0387}{124;\allowbreak{}\allowbreak{} printed124;\allowbreak{}\allowbreak{} manuscript200;\allowbreak{}\allowbreak{} opening\allowbreak{} pair\allowbreak{} of\allowbreak{} localizer\allowbreak{} symbols}{\$\allowbreak{}W\_\allowbreak{}S\textasciicircum{}\allowbreak{}\{\allowbreak{}\textbackslash{}\allowbreak{}mathrm\{\allowbreak{}d\}\allowbreak{}\}\allowbreak{}\$\allowbreak{},\allowbreak{}\allowbreak{} \$\allowbreak{}W\_\allowbreak{}S\textasciicircum{}\allowbreak{}\{\allowbreak{}\textbackslash{}\allowbreak{}mathrm\{\allowbreak{}d\}\allowbreak{}\}\allowbreak{}\$\allowbreak{}}{\$\allowbreak{}W\_\allowbreak{}S\textasciicircum{}\allowbreak{}\{\allowbreak{}\textbackslash{}\allowbreak{}mathrm\{\allowbreak{}d\}\allowbreak{}\}\allowbreak{}\$\allowbreak{},\allowbreak{}\allowbreak{} \$\allowbreak{}W\_\allowbreak{}X\textasciicircum{}\allowbreak{}\{\allowbreak{}\textbackslash{}\allowbreak{}mathrm\{\allowbreak{}d\}\allowbreak{}\}\allowbreak{}\$\allowbreak{}}{DER\_\allowbreak{}VI-c:\allowbreak{}43}
\DERapparatusitem{ENV0388}{124;\allowbreak{}\allowbreak{} running\allowbreak{} text\allowbreak{} or\allowbreak{} formula}{\textbackslash{}\allowbreak{}text\{\allowbreak{}def\}\allowbreak{}}{\textbackslash{}\allowbreak{}mathrm\{\allowbreak{}def\}\allowbreak{}}{TYPESETTING:\allowbreak{}7;\allowbreak{}\allowbreak{} NOTATION\_\allowbreak{}DER:\allowbreak{}17}
\DERapparatusitem{ENV0389}{124;\allowbreak{}\allowbreak{} printed124;\allowbreak{}\allowbreak{} manuscript201;\allowbreak{}\allowbreak{} last\allowbreak{} line\allowbreak{} of\allowbreak{} formula\allowbreak{} (\allowbreak{}**)\allowbreak{}}{\textbackslash{}\allowbreak{}mathrm\{\allowbreak{}hot\}\allowbreak{}\_\allowbreak{}\{\allowbreak{}S/\allowbreak{}W\}\allowbreak{}(\allowbreak{}X/\allowbreak{}S)\allowbreak{}}{\textbackslash{}\allowbreak{}mathrm\{\allowbreak{}hot\}\allowbreak{}\_\allowbreak{}\{\allowbreak{}S,\allowbreak{}W\}\allowbreak{}(\allowbreak{}X/\allowbreak{}S)\allowbreak{}}{DER\_\allowbreak{}VI-c:\allowbreak{}44}
\subsection*{Chapter VII}
\DERapparatusitem{ENV0390}{2;\allowbreak{}\allowbreak{} edition\allowbreak{} p2;\allowbreak{}\allowbreak{} ms3;\allowbreak{}\allowbreak{} displayed\allowbreak{} contiguous\allowbreak{} transition\allowbreak{} set}{}{,\allowbreak{}}{DER\_\allowbreak{}VII-a:\allowbreak{}7}
\DERapparatusitem{ENV0391}{3;\allowbreak{}\allowbreak{} edition\allowbreak{} p3;\allowbreak{}\allowbreak{} ms4–5;\allowbreak{}\allowbreak{} simplices\allowbreak{} and\allowbreak{} structural\allowbreak{} chains}{the\allowbreak{} full\allowbreak{} subcategories\allowbreak{} of\allowbreak{} \$\allowbreak{}I\$\allowbreak{}\allowbreak{} that\allowbreak{} are\allowbreak{} isomorphic}{the\allowbreak{} full\allowbreak{} subcategories\allowbreak{} of\allowbreak{} \$\allowbreak{}I\$\allowbreak{}\allowbreak{} whose\allowbreak{} vertices\allowbreak{} are\allowbreak{} consecutive\allowbreak{} in\allowbreak{} the\allowbreak{} total\allowbreak{} order\allowbreak{} and\allowbreak{} that\allowbreak{} are\allowbreak{} isomorphic}{DER\_\allowbreak{}VII-a:\allowbreak{}8}
\DERapparatusitem{ENV0392}{3;\allowbreak{}\allowbreak{} edition\allowbreak{} p3;\allowbreak{}\allowbreak{} ms4–5;\allowbreak{}\allowbreak{} simplices\allowbreak{} and\allowbreak{} structural\allowbreak{} chains}{that\allowbreak{} are\allowbreak{} injective\allowbreak{} on\allowbreak{} the\allowbreak{} vertices.\allowbreak{}}{that\allowbreak{} are\allowbreak{} injective\allowbreak{} on\allowbreak{} the\allowbreak{} vertices\allowbreak{} and\allowbreak{} whose\allowbreak{} image\allowbreak{} is\allowbreak{} an\allowbreak{} entire\allowbreak{} structural\allowbreak{} interval.\allowbreak{}}{DER\_\allowbreak{}VII-a:\allowbreak{}8}
\DERapparatusitem{ENV0393}{3;\allowbreak{}\allowbreak{} edition\allowbreak{} p3;\allowbreak{}\allowbreak{} ms5;\allowbreak{}\allowbreak{} Proposition\allowbreak{} 1(\allowbreak{}b)\allowbreak{}}{Every\allowbreak{} functor\allowbreak{} \$\allowbreak{}I\textbackslash{}\allowbreak{}longrightarrow\allowbreak{} I\$\allowbreak{}\allowbreak{} that\allowbreak{} fixes}{Every\allowbreak{} automorphism\allowbreak{} \$\allowbreak{}I\textbackslash{}\allowbreak{}longrightarrow\allowbreak{} I\$\allowbreak{}\allowbreak{} that\allowbreak{} fixes}{DER\_\allowbreak{}VII-a:\allowbreak{}9}
\DERapparatusitem{ENV0394}{4;\allowbreak{}\allowbreak{} edition\allowbreak{} p4;\allowbreak{}\allowbreak{} ms6;\allowbreak{}\allowbreak{} fork\allowbreak{} diagram\allowbreak{} rightmost\allowbreak{} vertex}{r}{s}{DER\_\allowbreak{}VII-a:\allowbreak{}10}
\DERapparatusitem{ENV0395}{4;\allowbreak{}\allowbreak{} edition\allowbreak{} p4;\allowbreak{}\allowbreak{} ms6;\allowbreak{}\allowbreak{} intrinsic-order\allowbreak{} reconstruction}{of\allowbreak{} \$\allowbreak{}I\$\allowbreak{}\allowbreak{} (\allowbreak{}among\allowbreak{} these}{of\allowbreak{} \$\allowbreak{}I\$\allowbreak{}\allowbreak{} (\allowbreak{}among\allowbreak{} its}{DER\_\allowbreak{}VII-a:\allowbreak{}11}
\DERapparatusitem{ENV0396}{4;\allowbreak{}\allowbreak{} edition\allowbreak{} p4;\allowbreak{}\allowbreak{} ms7;\allowbreak{}\allowbreak{} reconstruction\allowbreak{} in\allowbreak{} the\allowbreak{} remaining\allowbreak{} tail}{\$\allowbreak{}i\textbackslash{}\allowbreak{}backslash\allowbreak{} I\$\allowbreak{},\allowbreak{}\allowbreak{} \$\allowbreak{}I/\allowbreak{}i\$\allowbreak{}}{\$\allowbreak{}s\_\allowbreak{}i\textbackslash{}\allowbreak{}backslash\allowbreak{} I'\$\allowbreak{},\allowbreak{}\allowbreak{} \$\allowbreak{}I'/\allowbreak{}s\_\allowbreak{}i\$\allowbreak{}}{DER\_\allowbreak{}VII-a:\allowbreak{}12}
\DERapparatusitem{ENV0397}{6;\allowbreak{}\allowbreak{} edition\allowbreak{} p6;\allowbreak{}\allowbreak{} ms10;\allowbreak{}\allowbreak{} order\allowbreak{} of\allowbreak{} structural\allowbreak{} intervals}{each\allowbreak{} element\allowbreak{} of\allowbreak{} \$\allowbreak{}I\$\allowbreak{}}{each\allowbreak{} element\allowbreak{} of\allowbreak{} \$\allowbreak{}I\_\allowbreak{}i\$\allowbreak{}}{DER\_\allowbreak{}VII-a:\allowbreak{}18}
\DERapparatusitem{ENV0398}{6;\allowbreak{}\allowbreak{} edition\allowbreak{} p6;\allowbreak{}\allowbreak{} ms10;\allowbreak{}\allowbreak{} Proposition\allowbreak{} 3}{}{(\allowbreak{}(\allowbreak{}}{DER\_\allowbreak{}VII-a:\allowbreak{}19}
\DERapparatusitem{ENV0399}{6;\allowbreak{}\allowbreak{} edition\allowbreak{} p6;\allowbreak{}\allowbreak{} ms10;\allowbreak{}\allowbreak{} Proposition\allowbreak{} 3}{}{/\allowbreak{}s)\allowbreak{}\textbackslash{}\allowbreak{}cup(\allowbreak{}s\textbackslash{}\allowbreak{}backslash\allowbreak{} I)\allowbreak{})\allowbreak{}}{DER\_\allowbreak{}VII-a:\allowbreak{}19}
\DERapparatusitem{ENV0400}{7;\allowbreak{}\allowbreak{} edition\allowbreak{} p7;\allowbreak{}\allowbreak{} ms11;\allowbreak{}\allowbreak{} footnote\allowbreak{} 4\allowbreak{} second\allowbreak{} sentence}{Let\allowbreak{} us\allowbreak{} note\allowbreak{} moreover\allowbreak{} that\allowbreak{} \textbackslash{}\allowbreak{}textit\{\allowbreak{}every\}\allowbreak{}\allowbreak{} subcategory\allowbreak{} \$\allowbreak{}I'\$\allowbreak{}\allowbreak{} of\allowbreak{} \$\allowbreak{}I\$\allowbreak{}\allowbreak{} which\allowbreak{} is\allowbreak{} an\allowbreak{} interval,\allowbreak{}\allowbreak{} is\allowbreak{} necessarily\allowbreak{} full}{This\allowbreak{} fullness\allowbreak{} condition\allowbreak{} is\allowbreak{} part\allowbreak{} of\allowbreak{} the\allowbreak{} definition:\allowbreak{}\allowbreak{} a\allowbreak{} subcategory\allowbreak{} \$\allowbreak{}I'\$\allowbreak{}\allowbreak{} of\allowbreak{} \$\allowbreak{}I\$\allowbreak{}\allowbreak{} which\allowbreak{} is\allowbreak{} an\allowbreak{} interval\allowbreak{} is\allowbreak{} not\allowbreak{} necessarily\allowbreak{} full}{DER\_\allowbreak{}VII-a:\allowbreak{}21}
\DERapparatusitem{ENV0401}{9;\allowbreak{}\allowbreak{} edition\allowbreak{} p9;\allowbreak{}\allowbreak{} ms13;\allowbreak{}\allowbreak{} standardizing\allowbreak{} the\allowbreak{} type}{with\allowbreak{} the\allowbreak{} of\allowbreak{} type\allowbreak{} \$\allowbreak{}I\_\allowbreak{}0\$\allowbreak{}}{with\allowbreak{} those\allowbreak{} of\allowbreak{} type\allowbreak{} \$\allowbreak{}I\_\allowbreak{}0\$\allowbreak{}}{DER\_\allowbreak{}VII-a:\allowbreak{}22}
\DERapparatusitem{ENV0402}{9;\allowbreak{}\allowbreak{} edition\allowbreak{} p9;\allowbreak{}\allowbreak{} ms13;\allowbreak{}\allowbreak{} length\allowbreak{} equality}{}{(\allowbreak{}\textbackslash{}\allowbreak{}operatorname\{\allowbreak{}Ob\}\allowbreak{}}{DER\_\allowbreak{}VII-a:\allowbreak{}24}
\DERapparatusitem{ENV0403}{9;\allowbreak{}\allowbreak{} edition\allowbreak{} p9;\allowbreak{}\allowbreak{} ms13;\allowbreak{}\allowbreak{} length\allowbreak{} equality}{}{-1)\allowbreak{}}{DER\_\allowbreak{}VII-a:\allowbreak{}24}
\DERapparatusitem{ENV0404}{11;\allowbreak{}\allowbreak{} Edition18,\allowbreak{}\allowbreak{} ms28,\allowbreak{}\allowbreak{} final\allowbreak{} two\allowbreak{} arrow\allowbreak{} labels}{c)\allowbreak{},\allowbreak{}\textbackslash{}\allowbreak{}beta(\allowbreak{}c}{\textbackslash{}\allowbreak{}varphi,\allowbreak{}v)\allowbreak{},\allowbreak{}\textbackslash{}\allowbreak{}beta(\allowbreak{}\textbackslash{}\allowbreak{}varphi,\allowbreak{}v}{DER\_\allowbreak{}VII-a:\allowbreak{}40}
\DERapparatusitem{ENV0405}{12;\allowbreak{}\allowbreak{} edition\allowbreak{} p12;\allowbreak{}\allowbreak{} ms18;\allowbreak{}\allowbreak{} endpoint\allowbreak{} map\allowbreak{} α}{where\allowbreak{} \$\allowbreak{}\textbackslash{}\allowbreak{}alpha\$\allowbreak{}\allowbreak{} is\allowbreak{} the\allowbreak{} inclusion\allowbreak{} functor}{where\allowbreak{} \$\allowbreak{}\textbackslash{}\allowbreak{}alpha\$\allowbreak{}\allowbreak{} is\allowbreak{} the\allowbreak{} endpoint\allowbreak{} functor\allowbreak{} (\allowbreak{}an\allowbreak{} inclusion\allowbreak{} when\allowbreak{} the\allowbreak{} length\allowbreak{} is\allowbreak{} positive)\allowbreak{}}{DER\_\allowbreak{}VII-a:\allowbreak{}32}
\DERapparatusitem{ENV0406}{15;\allowbreak{}\allowbreak{} Edition15,\allowbreak{}\allowbreak{} ms23,\allowbreak{}\allowbreak{} identities\allowbreak{} for\allowbreak{} L}{C}{c}{DER\_\allowbreak{}VII-a:\allowbreak{}35}
\DERapparatusitem{ENV0407}{15;\allowbreak{}\allowbreak{} Edition15,\allowbreak{}\allowbreak{} ms24,\allowbreak{}\allowbreak{} item\allowbreak{} b}{A\allowbreak{} description\allowbreak{} of\allowbreak{} the\allowbreak{} compositions\allowbreak{} of\allowbreak{} paths}{A\allowbreak{} description\allowbreak{} of\allowbreak{} the\allowbreak{} composition\allowbreak{} of\allowbreak{} paths}{DER\_\allowbreak{}VII-a:\allowbreak{}38}
\DERapparatusitem{ENV0408}{17;\allowbreak{}\allowbreak{} running\allowbreak{} text\allowbreak{} or\allowbreak{} formula}{\textbackslash{}\allowbreak{}text\{\allowbreak{}def\}\allowbreak{}}{\textbackslash{}\allowbreak{}mathrm\{\allowbreak{}def\}\allowbreak{}}{TYPESETTING:\allowbreak{}7;\allowbreak{}\allowbreak{} NOTATION\_\allowbreak{}DER:\allowbreak{}17}
\DERapparatusitem{ENV0409}{17;\allowbreak{}\allowbreak{} running\allowbreak{} text\allowbreak{} or\allowbreak{} formula}{\textbackslash{}\allowbreak{}text\{\allowbreak{}def\}\allowbreak{}}{\textbackslash{}\allowbreak{}mathrm\{\allowbreak{}def\}\allowbreak{}}{TYPESETTING:\allowbreak{}7;\allowbreak{}\allowbreak{} NOTATION\_\allowbreak{}DER:\allowbreak{}17}
\DERapparatusitem{ENV0410}{17;\allowbreak{}\allowbreak{} running\allowbreak{} text\allowbreak{} or\allowbreak{} formula}{\textbackslash{}\allowbreak{}text\{\allowbreak{}def\}\allowbreak{}}{\textbackslash{}\allowbreak{}mathrm\{\allowbreak{}def\}\allowbreak{}}{TYPESETTING:\allowbreak{}7;\allowbreak{}\allowbreak{} NOTATION\_\allowbreak{}DER:\allowbreak{}17}
\DERapparatusitem{ENV0411}{19;\allowbreak{}\allowbreak{} Edition19,\allowbreak{}\allowbreak{} ms29,\allowbreak{}\allowbreak{} endpoint\allowbreak{} structure\allowbreak{} diagram}{hook',\allowbreak{}}{}{DER\_\allowbreak{}VII-a:\allowbreak{}43}
\DERapparatusitem{ENV0412}{19;\allowbreak{}\allowbreak{} Edition19,\allowbreak{}\allowbreak{} ms29,\allowbreak{}\allowbreak{} endpoint\allowbreak{} structure\allowbreak{} diagram}{hook,\allowbreak{}}{}{DER\_\allowbreak{}VII-a:\allowbreak{}43}
\DERapparatusitem{ENV0413}{22;\allowbreak{}\allowbreak{} Edition22,\allowbreak{}\allowbreak{} ms33,\allowbreak{}\allowbreak{} footnote11,\allowbreak{}\allowbreak{} comparable-subset\allowbreak{} lemma}{If\allowbreak{} \$\allowbreak{}I\_\allowbreak{}\{\allowbreak{}\textbackslash{}\allowbreak{}tau'\}\allowbreak{}\$\allowbreak{}\allowbreak{} is\allowbreak{} an\allowbreak{} interval,\allowbreak{}\allowbreak{} every\allowbreak{} subset}{If\allowbreak{} \$\allowbreak{}I\_\allowbreak{}\{\allowbreak{}\textbackslash{}\allowbreak{}tau'\}\allowbreak{}\$\allowbreak{}\allowbreak{} is\allowbreak{} an\allowbreak{} interval\allowbreak{} of\allowbreak{} strictly\allowbreak{} positive\allowbreak{} length,\allowbreak{}\allowbreak{} every\allowbreak{} subset}{DER\_\allowbreak{}VII-a:\allowbreak{}45}
\DERapparatusitem{ENV0414}{22;\allowbreak{}\allowbreak{} Edition22,\allowbreak{}\allowbreak{} ms33,\allowbreak{}\allowbreak{} paragraph\allowbreak{} before\allowbreak{} Proposition7}{if\allowbreak{} its\allowbreak{} value\allowbreak{} is\allowbreak{} not\allowbreak{} an\allowbreak{} endpoint\allowbreak{} of\allowbreak{} \$\allowbreak{}I\_\allowbreak{}\{\allowbreak{}\textbackslash{}\allowbreak{}tau'\}\allowbreak{}\$\allowbreak{}}{if\allowbreak{} its\allowbreak{} value\allowbreak{} is\allowbreak{} an\allowbreak{} interior\allowbreak{} critical\allowbreak{} vertex\allowbreak{} of\allowbreak{} \$\allowbreak{}I\_\allowbreak{}\{\allowbreak{}\textbackslash{}\allowbreak{}tau'\}\allowbreak{}\$\allowbreak{}}{DER\_\allowbreak{}VII-a:\allowbreak{}44}
\DERapparatusitem{ENV0415}{22;\allowbreak{}\allowbreak{} Edition22,\allowbreak{}\allowbreak{} ms34,\allowbreak{}\allowbreak{} paragraph\allowbreak{} after\allowbreak{} Proposition7}{this\allowbreak{} subinterval\allowbreak{} is\allowbreak{} \$\allowbreak{}[\allowbreak{}t\_\allowbreak{}0,\allowbreak{}t\_\allowbreak{}n]\allowbreak{}\$\allowbreak{}}{this\allowbreak{} subinterval\allowbreak{} has\allowbreak{} endpoints\allowbreak{} \$\allowbreak{}t\_\allowbreak{}0,\allowbreak{}t\_\allowbreak{}n\$\allowbreak{}\allowbreak{} and\allowbreak{} convex\allowbreak{} hull\allowbreak{} \$\allowbreak{}[\allowbreak{}t\_\allowbreak{}0,\allowbreak{}t\_\allowbreak{}n]\allowbreak{}\$\allowbreak{}}{DER\_\allowbreak{}VII-a:\allowbreak{}46}
\DERapparatusitem{ENV0416}{22;\allowbreak{}\allowbreak{} Edition22,\allowbreak{}\allowbreak{} ms34,\allowbreak{}\allowbreak{} first\allowbreak{} phrase\allowbreak{} of\allowbreak{} Corollary}{\textbackslash{}\allowbreak{}tau'}{\textbackslash{}\allowbreak{}tau}{DER\_\allowbreak{}VII-a:\allowbreak{}47}
\DERapparatusitem{ENV0417}{23;\allowbreak{}\allowbreak{} Edition23,\allowbreak{}\allowbreak{} ms35,\allowbreak{}\allowbreak{} Corollary2(\allowbreak{}a)\allowbreak{}}{'}{}{DER\_\allowbreak{}VII-a:\allowbreak{}49}
\DERapparatusitem{ENV0418}{23;\allowbreak{}\allowbreak{} Edition23,\allowbreak{}\allowbreak{} ms36,\allowbreak{}\allowbreak{} Corollary2(\allowbreak{}c)\allowbreak{}}{subseteq}{in}{DER\_\allowbreak{}VII-a:\allowbreak{}50}
\DERapparatusitem{ENV0419}{23;\allowbreak{}\allowbreak{} Edition23,\allowbreak{}\allowbreak{} ms36,\allowbreak{}\allowbreak{} parenthesis\allowbreak{} under\allowbreak{} alternative\allowbreak{} summation}{\textbackslash{}\allowbreak{}underbrace\{\allowbreak{}K\textbackslash{}\allowbreak{}in\textbackslash{}\allowbreak{}mathrm\{\allowbreak{}Cr\}\allowbreak{}(\allowbreak{}I\_\allowbreak{}\textbackslash{}\allowbreak{}tau)\allowbreak{},\allowbreak{}\textbackslash{}\allowbreak{}\allowbreak{} K\textbackslash{}\allowbreak{}subseteq\allowbreak{} J\}\allowbreak{}\_\allowbreak{}\{\allowbreak{}\textbackslash{}\allowbreak{}text\{\allowbreak{}(\allowbreak{}i.\allowbreak{}e.\allowbreak{}\allowbreak{} \}\allowbreak{}K\textbackslash{}\allowbreak{},\allowbreak{}\textbackslash{}\allowbreak{}in\textbackslash{}\allowbreak{},\allowbreak{}\textbackslash{}\allowbreak{}mathrm\{\allowbreak{}Cr\}\allowbreak{}(\allowbreak{}f\textasciicircum{}\allowbreak{}\{\allowbreak{}-1\}\allowbreak{}(\allowbreak{}J')\allowbreak{})\allowbreak{}\textbackslash{}\allowbreak{}text\{\allowbreak{})\allowbreak{}\}\allowbreak{}\}\allowbreak{}}{K\textbackslash{}\allowbreak{}in\textbackslash{}\allowbreak{}mathrm\{\allowbreak{}Cr\}\allowbreak{}(\allowbreak{}I\_\allowbreak{}\textbackslash{}\allowbreak{}tau)\allowbreak{},\allowbreak{}\textbackslash{}\allowbreak{}\allowbreak{} K\textbackslash{}\allowbreak{}subseteq\allowbreak{} J}{DER\_\allowbreak{}VII-a:\allowbreak{}51}
\DERapparatusitem{ENV0420}{24;\allowbreak{}\allowbreak{} Edition24,\allowbreak{}\allowbreak{} ms36-37,\allowbreak{}\allowbreak{} Corollary3\allowbreak{} display}{}{'}{DER\_\allowbreak{}VII-a:\allowbreak{}52}
\DERapparatusitem{ENV0421}{24;\allowbreak{}\allowbreak{} Edition24,\allowbreak{}\allowbreak{} ms36-37,\allowbreak{}\allowbreak{} Corollary3\allowbreak{} display}{ongrightarrow}{eq}{DER\_\allowbreak{}VII-a:\allowbreak{}52}
\DERapparatusitem{ENV0422}{25;\allowbreak{}\allowbreak{} ms39\allowbreak{} Proposition8\allowbreak{} target\allowbreak{} notation}{'\_\allowbreak{}\textbackslash{}\allowbreak{}tau}{\_\allowbreak{}\{\allowbreak{}\textbackslash{}\allowbreak{}tau\textasciicircum{}\allowbreak{}\textbackslash{}\allowbreak{}prime\}\allowbreak{}}{DER\_\allowbreak{}VII-a:\allowbreak{}54}
\DERapparatusitem{ENV0423}{26;\allowbreak{}\allowbreak{} ms41\allowbreak{} definition\allowbreak{} of\allowbreak{} set-theoretic\allowbreak{} part}{on\allowbreak{} \$\allowbreak{}u\$\allowbreak{},\allowbreak{}\allowbreak{} where\allowbreak{} of\allowbreak{} \$\allowbreak{}u\$\allowbreak{}\allowbreak{} for\allowbreak{} \$\allowbreak{}f\$\allowbreak{}}{on\allowbreak{} \$\allowbreak{}u\$\allowbreak{},\allowbreak{}\allowbreak{} or\allowbreak{} of\allowbreak{} \$\allowbreak{}u\$\allowbreak{}\allowbreak{} for\allowbreak{} \$\allowbreak{}f\$\allowbreak{}}{DER\_\allowbreak{}VII-a:\allowbreak{}55}
\DERapparatusitem{ENV0424}{26;\allowbreak{}\allowbreak{} ms41\allowbreak{} displayed\allowbreak{} list\allowbreak{} of\allowbreak{} surviving\allowbreak{} transitions}{n}{N}{DER\_\allowbreak{}VII-a:\allowbreak{}56}
\DERapparatusitem{ENV0425}{27;\allowbreak{}\allowbreak{} ms42\allowbreak{} unique\allowbreak{} target\allowbreak{} block}{=\allowbreak{}}{\textbackslash{}\allowbreak{}subseteq}{DER\_\allowbreak{}VII-a:\allowbreak{}73}
\DERapparatusitem{ENV0426}{27;\allowbreak{}\allowbreak{} ms42\allowbreak{} parenthetical\allowbreak{} part\allowbreak{} of\allowbreak{} f\allowbreak{} on\allowbreak{} v\_\allowbreak{}j}{(\allowbreak{}namely\allowbreak{} precisely\allowbreak{} \$\allowbreak{}[\allowbreak{}j-1,\allowbreak{}\allowbreak{} j]\allowbreak{}\$\allowbreak{})\allowbreak{}}{(\allowbreak{}namely\allowbreak{} precisely\allowbreak{} \$\allowbreak{}[\allowbreak{}t\_\allowbreak{}\{\allowbreak{}j-1\}\allowbreak{},\allowbreak{}\allowbreak{} t\_\allowbreak{}j]\allowbreak{}\$\allowbreak{})\allowbreak{}}{DER\_\allowbreak{}VII-a:\allowbreak{}74}
\DERapparatusitem{ENV0427}{29;\allowbreak{}\allowbreak{} ms45\allowbreak{} connected\allowbreak{} components\allowbreak{} in\allowbreak{} geometric\allowbreak{} reconstruction}{the\allowbreak{} ``connected\allowbreak{} components''\allowbreak{} (\allowbreak{}in\allowbreak{} the\allowbreak{} sense\allowbreak{} of\allowbreak{} the\allowbreak{} numbering\allowbreak{} of\allowbreak{} the\allowbreak{} vertices\allowbreak{} of\allowbreak{} \$\allowbreak{}I\_\allowbreak{}\textbackslash{}\allowbreak{}tau\$\allowbreak{})\allowbreak{}\allowbreak{} of\allowbreak{} the\allowbreak{} set\allowbreak{} \$\allowbreak{}E'\allowbreak{} =\allowbreak{}\allowbreak{} E'(\allowbreak{}f)\allowbreak{}\allowbreak{} \textbackslash{}\allowbreak{}subseteq\allowbreak{} I\_\allowbreak{}\textbackslash{}\allowbreak{}tau\allowbreak{} =\allowbreak{}\allowbreak{} [\allowbreak{}0,\allowbreak{}n]\allowbreak{}\$\allowbreak{}\allowbreak{} of\allowbreak{} vertices\allowbreak{} which\allowbreak{} are\allowbreak{} sources\allowbreak{} or\allowbreak{} targets\allowbreak{} of\allowbreak{} transition\allowbreak{} arrows\allowbreak{} degenerate\allowbreak{} for\allowbreak{} \$\allowbreak{}f\$\allowbreak{}}{the\allowbreak{} connected\allowbreak{} components\allowbreak{} of\allowbreak{} the\allowbreak{} graph\allowbreak{} \$\allowbreak{}E'\allowbreak{} =\allowbreak{}\allowbreak{} E'(\allowbreak{}f)\allowbreak{}\$\allowbreak{}\allowbreak{} whose\allowbreak{} vertices\allowbreak{} are\allowbreak{} all\allowbreak{} the\allowbreak{} vertices\allowbreak{} of\allowbreak{} \$\allowbreak{}I\_\allowbreak{}\textbackslash{}\allowbreak{}tau\allowbreak{} =\allowbreak{}\allowbreak{} [\allowbreak{}0,\allowbreak{}n]\allowbreak{}\$\allowbreak{}\allowbreak{} and\allowbreak{} whose\allowbreak{} edges\allowbreak{} are\allowbreak{} the\allowbreak{} transition\allowbreak{} arrows\allowbreak{} degenerate\allowbreak{} for\allowbreak{} \$\allowbreak{}f\$\allowbreak{}}{DER\_\allowbreak{}VII-a:\allowbreak{}75}
\DERapparatusitem{ENV0428}{29;\allowbreak{}\allowbreak{} ms45\allowbreak{} unique\allowbreak{} map\allowbreak{} to\allowbreak{} the\allowbreak{} target\allowbreak{} factor}{}{\{\allowbreak{}}{DER\_\allowbreak{}VII-a:\allowbreak{}76}
\DERapparatusitem{ENV0429}{29;\allowbreak{}\allowbreak{} ms45\allowbreak{} unique\allowbreak{} map\allowbreak{} to\allowbreak{} the\allowbreak{} target\allowbreak{} factor}{}{'\}\allowbreak{}}{DER\_\allowbreak{}VII-a:\allowbreak{}76}
\DERapparatusitem{ENV0430}{29;\allowbreak{}\allowbreak{} ms45\allowbreak{} definition\allowbreak{} of\allowbreak{} relative\allowbreak{} critical-image\allowbreak{} set\allowbreak{} at\allowbreak{} point\allowbreak{} target}{\$\allowbreak{}E'\$\allowbreak{}\allowbreak{} contains\allowbreak{} \$\allowbreak{}f(\allowbreak{}I\_\allowbreak{}\{\allowbreak{}\textbackslash{}\allowbreak{}tau\textbackslash{}\allowbreak{},\allowbreak{}\textbackslash{}\allowbreak{}text\{\allowbreak{}red\}\allowbreak{}\}\allowbreak{})\allowbreak{}\$\allowbreak{}}{When\allowbreak{} \$\allowbreak{}I\_\allowbreak{}\{\allowbreak{}\textbackslash{}\allowbreak{}tau'\}\allowbreak{}\$\allowbreak{}\allowbreak{} is\allowbreak{} a\allowbreak{} point,\allowbreak{}\allowbreak{} we\allowbreak{} set\allowbreak{} \$\allowbreak{}E'=\allowbreak{}\textbackslash{}\allowbreak{}\{\allowbreak{}0\textbackslash{}\allowbreak{}\}\allowbreak{}\$\allowbreak{},\allowbreak{}\allowbreak{} so\allowbreak{} that\allowbreak{} in\allowbreak{} every\allowbreak{} case\allowbreak{} \$\allowbreak{}E'=\allowbreak{}f(\allowbreak{}I\_\allowbreak{}\textbackslash{}\allowbreak{}tau)\allowbreak{}\$\allowbreak{}.\allowbreak{}\allowbreak{} \$\allowbreak{}E'\$\allowbreak{}\allowbreak{} contains\allowbreak{} \$\allowbreak{}f(\allowbreak{}I\_\allowbreak{}\{\allowbreak{}\textbackslash{}\allowbreak{}tau\textbackslash{}\allowbreak{},\allowbreak{}\textbackslash{}\allowbreak{}text\{\allowbreak{}red\}\allowbreak{}\}\allowbreak{})\allowbreak{}\$\allowbreak{}}{DER\_\allowbreak{}VII-a:\allowbreak{}77}
\DERapparatusitem{ENV0431}{30;\allowbreak{}\allowbreak{} ms46\allowbreak{} equivalent\allowbreak{} condition\allowbreak{} on\allowbreak{} E-prime}{(\allowbreak{}where\allowbreak{} what\allowbreak{} amounts\allowbreak{} to\allowbreak{} the\allowbreak{} same\allowbreak{} thing,\allowbreak{}}{(\allowbreak{}or\allowbreak{} what\allowbreak{} amounts\allowbreak{} to\allowbreak{} the\allowbreak{} same\allowbreak{} thing,\allowbreak{}}{DER\_\allowbreak{}VII-a:\allowbreak{}58}
\DERapparatusitem{ENV0432}{30;\allowbreak{}\allowbreak{} ms46\allowbreak{} enumeration\allowbreak{} of\allowbreak{} visited\allowbreak{} vertices\allowbreak{} before\allowbreak{} i\allowbreak{} in\allowbreak{} [\allowbreak{}1,\allowbreak{}N]\allowbreak{}}{r}{N}{DER\_\allowbreak{}VII-a:\allowbreak{}79}
\DERapparatusitem{ENV0433}{30;\allowbreak{}\allowbreak{} ms47\allowbreak{} sign\allowbreak{} codomain\allowbreak{} for\allowbreak{} arrow\allowbreak{} data}{[\allowbreak{}+1,\allowbreak{}-1]\allowbreak{}}{\textbackslash{}\allowbreak{}\{\allowbreak{}+1,\allowbreak{}-1\textbackslash{}\allowbreak{}\}\allowbreak{}}{DER\_\allowbreak{}VII-a:\allowbreak{}80}
\DERapparatusitem{ENV0434}{31;\allowbreak{}\allowbreak{} ms47\allowbreak{} explicit\allowbreak{} object-map\allowbreak{} table,\allowbreak{}\allowbreak{} third\allowbreak{} and\allowbreak{} general\allowbreak{} rows}{}{1+l\_\allowbreak{}}{DER\_\allowbreak{}VII-a:\allowbreak{}81}
\DERapparatusitem{ENV0435}{31;\allowbreak{}\allowbreak{} ms47\allowbreak{} explicit\allowbreak{} object-map\allowbreak{} table,\allowbreak{}\allowbreak{} third\allowbreak{} and\allowbreak{} general\allowbreak{} rows}{}{\textbackslash{}\allowbreak{}sum\_\allowbreak{}\{\allowbreak{}h=\allowbreak{}1\}\allowbreak{}\textasciicircum{}\allowbreak{}\{\allowbreak{}i+1\}\allowbreak{}}{DER\_\allowbreak{}VII-a:\allowbreak{}81}
\DERapparatusitem{ENV0436}{31;\allowbreak{}\allowbreak{} ms47\allowbreak{} explicit\allowbreak{} object-map\allowbreak{} table,\allowbreak{}\allowbreak{} third\allowbreak{} and\allowbreak{} general\allowbreak{} rows}{i}{h}{DER\_\allowbreak{}VII-a:\allowbreak{}81}
\DERapparatusitem{ENV0437}{33;\allowbreak{}\allowbreak{} ms50\allowbreak{} extension\allowbreak{} of\allowbreak{} concatenation\allowbreak{} to\allowbreak{} all\allowbreak{} Typ}{(\allowbreak{}In\allowbreak{} fact,\allowbreak{}\allowbreak{} we\allowbreak{} also\allowbreak{} have\allowbreak{} a\allowbreak{} bifunctor}{(\allowbreak{}In\allowbreak{} fact,\allowbreak{}\allowbreak{} restricting\allowbreak{} here\allowbreak{} the\allowbreak{} morphisms\allowbreak{} of\allowbreak{} \$\allowbreak{}\textbackslash{}\allowbreak{}uTyp\$\allowbreak{}\allowbreak{} to\allowbreak{} functors\allowbreak{} preserving\allowbreak{} the\allowbreak{} two\allowbreak{} endpoints,\allowbreak{}\allowbreak{} we\allowbreak{} also\allowbreak{} have\allowbreak{} a\allowbreak{} bifunctor}{DER\_\allowbreak{}VII-a:\allowbreak{}82}
\DERapparatusitem{ENV0438}{34;\allowbreak{}\allowbreak{} ms50-prime\allowbreak{} Cartesian-power\allowbreak{} parenthesis}{Cartesian\allowbreak{} power\allowbreak{} in\allowbreak{} Cat,\allowbreak{}\allowbreak{} where\allowbreak{} in\allowbreak{} \$\allowbreak{}\textbackslash{}\allowbreak{}uTy\$\allowbreak{}}{Cartesian\allowbreak{} power\allowbreak{} in\allowbreak{} Cat,\allowbreak{}\allowbreak{} or\allowbreak{} in\allowbreak{} \$\allowbreak{}\textbackslash{}\allowbreak{}uTy\$\allowbreak{}}{DER\_\allowbreak{}VII-a:\allowbreak{}62}
\DERapparatusitem{ENV0439}{36;\allowbreak{}\allowbreak{} ms52\allowbreak{} second\allowbreak{} map\allowbreak{} in\allowbreak{} composable\allowbreak{} pair}{q}{g}{DER\_\allowbreak{}VII-a:\allowbreak{}65}
\DERapparatusitem{ENV0440}{36;\allowbreak{}\allowbreak{} ms52\allowbreak{} Hom-set\allowbreak{} decomposition}{}{,\allowbreak{}\textbackslash{}\allowbreak{}mathbf\allowbreak{} N}{DER\_\allowbreak{}VII-a:\allowbreak{}66}
\DERapparatusitem{ENV0441}{36;\allowbreak{}\allowbreak{} ms52\allowbreak{} displayed\allowbreak{} two-step\allowbreak{} generator\allowbreak{} (\allowbreak{}1)\allowbreak{}}{1}{i}{DER\_\allowbreak{}VII-a:\allowbreak{}67}
\DERapparatusitem{ENV0442}{36;\allowbreak{}\allowbreak{} ms52\allowbreak{} displayed\allowbreak{} two-step\allowbreak{} generator\allowbreak{} (\allowbreak{}1)\allowbreak{}}{}{\_\allowbreak{}i}{DER\_\allowbreak{}VII-a:\allowbreak{}67}
\DERapparatusitem{ENV0443}{36;\allowbreak{}\allowbreak{} running\allowbreak{} text\allowbreak{} or\allowbreak{} formula}{\textbackslash{}\allowbreak{}text\{\allowbreak{}def\}\allowbreak{}}{\textbackslash{}\allowbreak{}mathrm\{\allowbreak{}def\}\allowbreak{}}{TYPESETTING:\allowbreak{}7;\allowbreak{}\allowbreak{} NOTATION\_\allowbreak{}DER:\allowbreak{}17}
\DERapparatusitem{ENV0444}{36;\allowbreak{}\allowbreak{} ms53\allowbreak{} final\allowbreak{} factor\allowbreak{} in\allowbreak{} unique\allowbreak{} two-step\allowbreak{} factorization}{}{,\allowbreak{}}{DER\_\allowbreak{}VII-a:\allowbreak{}68}
\DERapparatusitem{ENV0445}{37;\allowbreak{}\allowbreak{} ms54\allowbreak{} second\allowbreak{} factor\allowbreak{} g\_\allowbreak{}i}{l}{\textbackslash{}\allowbreak{}pi}{DER\_\allowbreak{}VII-a:\allowbreak{}69}
\DERapparatusitem{ENV0446}{38;\allowbreak{}\allowbreak{} ms55\allowbreak{} identity\allowbreak{} inserted\allowbreak{} before\allowbreak{} g}{The\allowbreak{} \$\allowbreak{}f=\allowbreak{}\textbackslash{}\allowbreak{}mathrm\{\allowbreak{}id\}\allowbreak{}\$\allowbreak{}\allowbreak{} is}{The\allowbreak{} case\allowbreak{} \$\allowbreak{}f=\allowbreak{}\textbackslash{}\allowbreak{}mathrm\{\allowbreak{}id\}\allowbreak{}\$\allowbreak{}\allowbreak{} is}{DER\_\allowbreak{}VII-a:\allowbreak{}71}
\DERapparatusitem{ENV0447}{38;\allowbreak{}\allowbreak{} running\allowbreak{} text\allowbreak{} or\allowbreak{} formula}{\textbackslash{}\allowbreak{}text\{\allowbreak{}def\}\allowbreak{}}{\textbackslash{}\allowbreak{}mathrm\{\allowbreak{}def\}\allowbreak{}}{TYPESETTING:\allowbreak{}7;\allowbreak{}\allowbreak{} NOTATION\_\allowbreak{}DER:\allowbreak{}17}
\DERapparatusitem{ENV0448}{38;\allowbreak{}\allowbreak{} ms56\allowbreak{} definition\allowbreak{} of\allowbreak{} length\allowbreak{} of\allowbreak{} pi\_\allowbreak{}(\allowbreak{}i-1)\allowbreak{}}{}{\{\allowbreak{}}{DER\_\allowbreak{}VII-a:\allowbreak{}72}
\DERapparatusitem{ENV0449}{38;\allowbreak{}\allowbreak{} ms56\allowbreak{} definition\allowbreak{} of\allowbreak{} length\allowbreak{} of\allowbreak{} pi\_\allowbreak{}(\allowbreak{}i-1)\allowbreak{}}{}{-1\}\allowbreak{}}{DER\_\allowbreak{}VII-a:\allowbreak{}72}
\DERapparatusitem{ENV0450}{38;\allowbreak{}\allowbreak{} ms56\allowbreak{} definition\allowbreak{} of\allowbreak{} length\allowbreak{} of\allowbreak{} pi\_\allowbreak{}(\allowbreak{}i-1)\allowbreak{}}{}{\{\allowbreak{}}{DER\_\allowbreak{}VII-a:\allowbreak{}72}
\DERapparatusitem{ENV0451}{38;\allowbreak{}\allowbreak{} ms56\allowbreak{} definition\allowbreak{} of\allowbreak{} length\allowbreak{} of\allowbreak{} pi\_\allowbreak{}(\allowbreak{}i-1)\allowbreak{}}{}{-1\}\allowbreak{}}{DER\_\allowbreak{}VII-a:\allowbreak{}72}
\DERapparatusitem{ENV0452}{38;\allowbreak{}\allowbreak{} ms57\allowbreak{} final\allowbreak{} boxed\allowbreak{} support\allowbreak{} range}{0}{1}{DER\_\allowbreak{}VII-a:\allowbreak{}83}
\DERapparatusitem{ENV0453}{39;\allowbreak{}\allowbreak{} ms57,\allowbreak{}\allowbreak{} zero-length\allowbreak{} propagation}{we\allowbreak{} must\allowbreak{} have\allowbreak{} \$\allowbreak{}\textbackslash{}\allowbreak{}mathrm\{\allowbreak{}supp\}\allowbreak{}(\allowbreak{}\textbackslash{}\allowbreak{}pi\_\allowbreak{}\{\allowbreak{}i+1\}\allowbreak{})\allowbreak{}\allowbreak{} =\allowbreak{}\allowbreak{} \textbackslash{}\allowbreak{}emptyset\$\allowbreak{},\allowbreak{}\allowbreak{} i.\allowbreak{}e.\allowbreak{}\allowbreak{} \$\allowbreak{}\textbackslash{}\allowbreak{}pi\_\allowbreak{}\{\allowbreak{}i+1\}\allowbreak{}\allowbreak{} =\allowbreak{}\allowbreak{} 0\$\allowbreak{}}{for\allowbreak{} \$\allowbreak{}i<r\$\allowbreak{},\allowbreak{}\allowbreak{} we\allowbreak{} must\allowbreak{} have\allowbreak{} \$\allowbreak{}\textbackslash{}\allowbreak{}mathrm\{\allowbreak{}supp\}\allowbreak{}(\allowbreak{}\textbackslash{}\allowbreak{}pi\_\allowbreak{}i)\allowbreak{}\allowbreak{} =\allowbreak{}\allowbreak{} \textbackslash{}\allowbreak{}emptyset\$\allowbreak{},\allowbreak{}\allowbreak{} i.\allowbreak{}e.\allowbreak{}\allowbreak{} \$\allowbreak{}\textbackslash{}\allowbreak{}pi\_\allowbreak{}i\allowbreak{} =\allowbreak{}\allowbreak{} 0\$\allowbreak{}}{DER\_\allowbreak{}VII-b:\allowbreak{}1}
\DERapparatusitem{ENV0454}{39;\allowbreak{}\allowbreak{} ms59,\allowbreak{}\allowbreak{} final\allowbreak{} sentence}{the\allowbreak{} calculation\allowbreak{} of\allowbreak{} \$\allowbreak{}\textbackslash{}\allowbreak{}pi\_\allowbreak{}0\$\allowbreak{}\allowbreak{} of\allowbreak{} an\allowbreak{} object\allowbreak{} of\allowbreak{} \$\allowbreak{}X\$\allowbreak{}}{the\allowbreak{} calculation\allowbreak{} of\allowbreak{} \$\allowbreak{}\textbackslash{}\allowbreak{}pi\_\allowbreak{}0\$\allowbreak{}\allowbreak{} of\allowbreak{} an\allowbreak{} object\allowbreak{} \$\allowbreak{}X\$\allowbreak{}\allowbreak{} of\allowbreak{} Cat}{DER\_\allowbreak{}VII-b:\allowbreak{}2}
\DERapparatusitem{ENV0455}{40;\allowbreak{}\allowbreak{} ms60,\allowbreak{}\allowbreak{} exception\allowbreak{} in\allowbreak{} the\allowbreak{} sentence\allowbreak{} about\allowbreak{} assigning\allowbreak{} an\allowbreak{} X-arrow}{of\allowbreak{} type\allowbreak{} \$\allowbreak{}\textbackslash{}\allowbreak{}Delta\_\allowbreak{}+\$\allowbreak{})\allowbreak{}}{of\allowbreak{} type\allowbreak{} \$\allowbreak{}\textbackslash{}\allowbreak{}Delta\_\allowbreak{}+\textasciicircum{}\allowbreak{}n\$\allowbreak{},\allowbreak{}\allowbreak{} or\allowbreak{} more\allowbreak{} generally\allowbreak{} if\allowbreak{} its\allowbreak{} retrograde\allowbreak{} arrows\allowbreak{} are\allowbreak{} invertible)\allowbreak{}}{DER\_\allowbreak{}VII-b:\allowbreak{}3}
\DERapparatusitem{ENV0456}{42;\allowbreak{}\allowbreak{} ms63,\allowbreak{}\allowbreak{} second\allowbreak{} cancellation\allowbreak{} diagram\allowbreak{} and\allowbreak{} following\allowbreak{} text}{\textbackslash{}\allowbreak{}begin\{\allowbreak{}tikzcd\}\allowbreak{}\space{}\allowbreak{}\allowbreak{} \&\allowbreak{}\allowbreak{} b\allowbreak{} \&\allowbreak{}\allowbreak{} a\allowbreak{} \textbackslash{}\allowbreak{}arrow[\allowbreak{}l,\allowbreak{}\allowbreak{} "u"']\allowbreak{}\allowbreak{} \textbackslash{}\allowbreak{}arrow[\allowbreak{}r,\allowbreak{}\allowbreak{} "u"]\allowbreak{}\allowbreak{} \&\allowbreak{}\allowbreak{} b\allowbreak{} \textbackslash{}\allowbreak{}\textbackslash{}\allowbreak{}\space{}\allowbreak{}\allowbreak{} \{\allowbreak{}\}\allowbreak{}\allowbreak{} \&\allowbreak{}\allowbreak{} a\allowbreak{} \textbackslash{}\allowbreak{}arrow[\allowbreak{}u,\allowbreak{}\allowbreak{} "u"']\allowbreak{}\allowbreak{} \&\allowbreak{}\allowbreak{} a\allowbreak{} \textbackslash{}\allowbreak{}arrow[\allowbreak{}u,\allowbreak{}\allowbreak{} "\textbackslash{}\allowbreak{}mathrm\{\allowbreak{}id\}\allowbreak{}"']\allowbreak{}\allowbreak{} \textbackslash{}\allowbreak{}arrow[\allowbreak{}l,\allowbreak{}\allowbreak{} "\textbackslash{}\allowbreak{}mathrm\{\allowbreak{}id\}\allowbreak{}"']\allowbreak{}\allowbreak{} \textbackslash{}\allowbreak{}arrow[\allowbreak{}r,\allowbreak{}\allowbreak{} "\textbackslash{}\allowbreak{}mathrm\{\allowbreak{}id\}\allowbreak{}"]\allowbreak{}\allowbreak{} \&\allowbreak{}\allowbreak{} a\allowbreak{} \textbackslash{}\allowbreak{}arrow[\allowbreak{}u,\allowbreak{}\allowbreak{} "u"']\allowbreak{}\allowbreak{} \textbackslash{}\allowbreak{}\textbackslash{}\allowbreak{}\space{}\allowbreak{}\allowbreak{} \{\allowbreak{}\}\allowbreak{}\allowbreak{} \textbackslash{}\allowbreak{}arrow[\allowbreak{}u]\allowbreak{}\allowbreak{} \&\allowbreak{}\allowbreak{} \&\allowbreak{}\allowbreak{} a\allowbreak{} \&\allowbreak{}\space{}\allowbreak{}\textbackslash{}\allowbreak{}end\{\allowbreak{}tikzcd\}\allowbreak{}}{\textbackslash{}\allowbreak{}begin\{\allowbreak{}tikzcd\}\allowbreak{}\space{}\allowbreak{}b\allowbreak{} \textbackslash{}\allowbreak{}arrow[\allowbreak{}d,\allowbreak{}\allowbreak{} "\textbackslash{}\allowbreak{}mathrm\{\allowbreak{}id\}\allowbreak{}"]\allowbreak{}\allowbreak{} \&\allowbreak{}\allowbreak{} a\allowbreak{} \textbackslash{}\allowbreak{}arrow[\allowbreak{}l,\allowbreak{}\allowbreak{} "u"']\allowbreak{}\allowbreak{} \textbackslash{}\allowbreak{}arrow[\allowbreak{}r,\allowbreak{}\allowbreak{} "u"]\allowbreak{}\allowbreak{} \textbackslash{}\allowbreak{}arrow[\allowbreak{}d,\allowbreak{}\allowbreak{} "u"]\allowbreak{}\allowbreak{} \&\allowbreak{}\allowbreak{} b\allowbreak{} \textbackslash{}\allowbreak{}arrow[\allowbreak{}d,\allowbreak{}\allowbreak{} "\textbackslash{}\allowbreak{}mathrm\{\allowbreak{}id\}\allowbreak{}"]\allowbreak{}\allowbreak{} \textbackslash{}\allowbreak{}\textbackslash{}\allowbreak{}\space{}\allowbreak{}b\allowbreak{} \&\allowbreak{}\allowbreak{} b\allowbreak{} \textbackslash{}\allowbreak{}arrow[\allowbreak{}l,\allowbreak{}\allowbreak{} "\textbackslash{}\allowbreak{}mathrm\{\allowbreak{}id\}\allowbreak{}"']\allowbreak{}\allowbreak{} \textbackslash{}\allowbreak{}arrow[\allowbreak{}r,\allowbreak{}\allowbreak{} "\textbackslash{}\allowbreak{}mathrm\{\allowbreak{}id\}\allowbreak{}"]\allowbreak{}\allowbreak{} \&\allowbreak{}\allowbreak{} b\space{}\allowbreak{}\textbackslash{}\allowbreak{}end\{\allowbreak{}tikzcd\}\allowbreak{}}{DER\_\allowbreak{}VII-b:\allowbreak{}5}
\DERapparatusitem{ENV0457}{42;\allowbreak{}\allowbreak{} ms63,\allowbreak{}\allowbreak{} second\allowbreak{} cancellation\allowbreak{} diagram\allowbreak{} and\allowbreak{} following\allowbreak{} text}{righ}{lef}{DER\_\allowbreak{}VII-b:\allowbreak{}5}
\DERapparatusitem{ENV0458}{42;\allowbreak{}\allowbreak{} ms63,\allowbreak{}\allowbreak{} superscript\allowbreak{} in\allowbreak{} original\allowbreak{} cancellation\allowbreak{} text}{0}{o}{DER\_\allowbreak{}VII-b:\allowbreak{}6}
\DERapparatusitem{ENV0459}{44;\allowbreak{}\allowbreak{} ms66,\allowbreak{}\allowbreak{} definition\allowbreak{} of\allowbreak{} Sigma-path}{X}{\textbackslash{}\allowbreak{}Sigma}{DER\_\allowbreak{}VII-b:\allowbreak{}8}
\DERapparatusitem{ENV0460}{44;\allowbreak{}\allowbreak{} ms66,\allowbreak{}\allowbreak{} last\allowbreak{} explanatory\allowbreak{} paragraph}{or\allowbreak{} arrows\allowbreak{} among\allowbreak{} those\allowbreak{} of\allowbreak{} \$\allowbreak{}\textbackslash{}\allowbreak{}tau\$\allowbreak{},\allowbreak{}\allowbreak{} or\allowbreak{} composites\allowbreak{} of\allowbreak{} a\allowbreak{} sequence\allowbreak{} of\allowbreak{} consecutive\allowbreak{} composable\allowbreak{} arrows\allowbreak{} of\allowbreak{} \$\allowbreak{}\textbackslash{}\allowbreak{}tau'\$\allowbreak{}}{or\allowbreak{} arrows\allowbreak{} among\allowbreak{} those\allowbreak{} of\allowbreak{} the\allowbreak{} path\allowbreak{} \$\allowbreak{}c'\$\allowbreak{}\allowbreak{} of\allowbreak{} type\allowbreak{} \$\allowbreak{}\textbackslash{}\allowbreak{}tau'\$\allowbreak{},\allowbreak{}\allowbreak{} or\allowbreak{} composites\allowbreak{} of\allowbreak{} a\allowbreak{} sequence\allowbreak{} of\allowbreak{} consecutive\allowbreak{} composable\allowbreak{} arrows\allowbreak{} of\allowbreak{} this\allowbreak{} path}{DER\_\allowbreak{}VII-b:\allowbreak{}10}
\DERapparatusitem{ENV0461}{45;\allowbreak{}\allowbreak{} ms68,\allowbreak{}\allowbreak{} final\allowbreak{} sentence}{bet}{alph}{DER\_\allowbreak{}VII-b:\allowbreak{}11}
\DERapparatusitem{ENV0462}{46;\allowbreak{}\allowbreak{} ms69,\allowbreak{}\allowbreak{} boxed\allowbreak{} equation\allowbreak{} (\allowbreak{}1)\allowbreak{}}{}{\_\allowbreak{}\textbackslash{}\allowbreak{}Sigma}{DER\_\allowbreak{}VII-b:\allowbreak{}15}
\DERapparatusitem{ENV0463}{46;\allowbreak{}\allowbreak{} ms69,\allowbreak{}\allowbreak{} equation\allowbreak{} (\allowbreak{}2)\allowbreak{}}{}{\_\allowbreak{}\textbackslash{}\allowbreak{}Sigma}{DER\_\allowbreak{}VII-b:\allowbreak{}16}
\DERapparatusitem{ENV0464}{48;\allowbreak{}\allowbreak{} ms72,\allowbreak{}\allowbreak{} uniqueness\allowbreak{} sentence}{when\allowbreak{} these\allowbreak{} inverses\allowbreak{} are\allowbreak{} invertible}{when\allowbreak{} the\allowbreak{} arrows\allowbreak{} under\allowbreak{} consideration\allowbreak{} are\allowbreak{} invertible}{DER\_\allowbreak{}VII-b:\allowbreak{}18}
\DERapparatusitem{ENV0465}{48;\allowbreak{}\allowbreak{} ms73,\allowbreak{}\allowbreak{} displayed\allowbreak{} map\allowbreak{} on\allowbreak{} components}{\textbackslash{}\allowbreak{}Sigma}{\textbackslash{}\allowbreak{}Sigma'}{DER\_\allowbreak{}VII-b:\allowbreak{}19}
\DERapparatusitem{ENV0466}{48;\allowbreak{}\allowbreak{} ms73,\allowbreak{}\allowbreak{} compatibility\allowbreak{} display}{}{'}{DER\_\allowbreak{}VII-b:\allowbreak{}20}
\DERapparatusitem{ENV0467}{48;\allowbreak{}\allowbreak{} ms73,\allowbreak{}\allowbreak{} compatibility\allowbreak{} display}{'}{}{DER\_\allowbreak{}VII-b:\allowbreak{}20}
\DERapparatusitem{ENV0468}{49;\allowbreak{}\allowbreak{} ms76,\allowbreak{}\allowbreak{} final\allowbreak{} sign\allowbreak{} in\allowbreak{} two\allowbreak{} displayed\allowbreak{} words}{}{-1}{DER\_\allowbreak{}VII-b:\allowbreak{}21}
\DERapparatusitem{ENV0469}{49;\allowbreak{}\allowbreak{} ms76,\allowbreak{}\allowbreak{} final\allowbreak{} sign\allowbreak{} in\allowbreak{} two\allowbreak{} displayed\allowbreak{} words}{}{-1}{DER\_\allowbreak{}VII-b:\allowbreak{}21}
\DERapparatusitem{ENV0470}{50;\allowbreak{}\allowbreak{} ms77,\allowbreak{}\allowbreak{} pushout\allowbreak{} square\allowbreak{} top\allowbreak{} object}{\textbackslash{}\allowbreak{}varepsilon\allowbreak{} =\allowbreak{}}{}{DER\_\allowbreak{}VII-b:\allowbreak{}23}
\DERapparatusitem{ENV0471}{50;\allowbreak{}\allowbreak{} ms77,\allowbreak{}\allowbreak{} pushout\allowbreak{} square\allowbreak{} top\allowbreak{} object}{}{1\_\allowbreak{}\{\allowbreak{}\textbackslash{}\allowbreak{}tau\}\allowbreak{}=\allowbreak{}}{DER\_\allowbreak{}VII-b:\allowbreak{}23}
\DERapparatusitem{ENV0472}{50;\allowbreak{}\allowbreak{} ms77,\allowbreak{}\allowbreak{} pushout\allowbreak{} square\allowbreak{} top\allowbreak{} object}{,\allowbreak{}1}{\_\allowbreak{}\{\allowbreak{}\textbackslash{}\allowbreak{}tau'\}\allowbreak{}}{DER\_\allowbreak{}VII-b:\allowbreak{}23}
\DERapparatusitem{ENV0473}{50;\allowbreak{}\allowbreak{} ms77,\allowbreak{}\allowbreak{} joining-vertex\allowbreak{} index}{+1}{}{DER\_\allowbreak{}VII-b:\allowbreak{}25}
\DERapparatusitem{ENV0474}{51;\allowbreak{}\allowbreak{} ms78,\allowbreak{}\allowbreak{} cartesian\allowbreak{} functor\allowbreak{} (\allowbreak{}**)\allowbreak{}\allowbreak{} indices}{}{'}{DER\_\allowbreak{}VII-b:\allowbreak{}28}
\DERapparatusitem{ENV0475}{51;\allowbreak{}\allowbreak{} ms78,\allowbreak{}\allowbreak{} cartesian\allowbreak{} functor\allowbreak{} (\allowbreak{}**)\allowbreak{}\allowbreak{} indices}{}{\textbackslash{}\allowbreak{}tau'}{DER\_\allowbreak{}VII-b:\allowbreak{}28}
\DERapparatusitem{ENV0476}{51;\allowbreak{}\allowbreak{} ms79,\allowbreak{}\allowbreak{} upper-right\allowbreak{} corner\allowbreak{} of\allowbreak{} compatibility\allowbreak{} square}{\textbackslash{}\allowbreak{}tau}{\textbackslash{}\allowbreak{}tau\textbackslash{}\allowbreak{}tau'}{DER\_\allowbreak{}VII-b:\allowbreak{}29}
\DERapparatusitem{ENV0477}{52;\allowbreak{}\allowbreak{} running\allowbreak{} text\allowbreak{} or\allowbreak{} formula}{\allowbreak{} (\allowbreak{}gear)\allowbreak{}}{}{GLOSSARY\_\allowbreak{}DER:\allowbreak{}164}
\DERapparatusitem{ENV0478}{53;\allowbreak{}\allowbreak{} ms82,\allowbreak{}\allowbreak{} variance\allowbreak{} of\allowbreak{} diagram\allowbreak{} category}{}{\textasciicircum{}\allowbreak{}\{\allowbreak{}o\}\allowbreak{}}{DER\_\allowbreak{}VII-b:\allowbreak{}31}
\DERapparatusitem{ENV0479}{54;\allowbreak{}\allowbreak{} ms83,\allowbreak{}\allowbreak{} endpoint\allowbreak{} diagram\allowbreak{} hooks}{,\allowbreak{}\allowbreak{} hook'}{}{DER\_\allowbreak{}VII-b:\allowbreak{}34}
\DERapparatusitem{ENV0480}{54;\allowbreak{}\allowbreak{} ms83,\allowbreak{}\allowbreak{} endpoint\allowbreak{} diagram\allowbreak{} hooks}{hook,\allowbreak{}}{}{DER\_\allowbreak{}VII-b:\allowbreak{}34}
\DERapparatusitem{ENV0481}{54;\allowbreak{}\allowbreak{} ms83,\allowbreak{}\allowbreak{} generating\allowbreak{} transition\allowbreak{} arrow}{the\allowbreak{} transition\allowbreak{} arrows\allowbreak{} \$\allowbreak{}\textbackslash{}\allowbreak{}ul\{\allowbreak{}u\}\allowbreak{}\_\allowbreak{}i:\allowbreak{}i-1\textbackslash{}\allowbreak{}longrightarrow\allowbreak{} i\$\allowbreak{}\allowbreak{} in\allowbreak{} \$\allowbreak{}I\_\allowbreak{}\textbackslash{}\allowbreak{}tau\$\allowbreak{}}{the\allowbreak{} transition\allowbreak{} arrows\allowbreak{} \$\allowbreak{}\textbackslash{}\allowbreak{}ul\{\allowbreak{}u\}\allowbreak{}\_\allowbreak{}i:\allowbreak{}i-1\textbackslash{}\allowbreak{}longrightarrow\allowbreak{} i\$\allowbreak{}\allowbreak{} if\allowbreak{} \$\allowbreak{}\textbackslash{}\allowbreak{}tau\_\allowbreak{}i=\allowbreak{}+\$\allowbreak{},\allowbreak{}\allowbreak{} or\allowbreak{} \$\allowbreak{}\textbackslash{}\allowbreak{}ul\{\allowbreak{}u\}\allowbreak{}\_\allowbreak{}i:\allowbreak{}i\textbackslash{}\allowbreak{}longrightarrow\allowbreak{} i-1\$\allowbreak{}\allowbreak{} if\allowbreak{} \$\allowbreak{}\textbackslash{}\allowbreak{}tau\_\allowbreak{}i=\allowbreak{}-\$\allowbreak{},\allowbreak{}\allowbreak{} in\allowbreak{} \$\allowbreak{}I\_\allowbreak{}\textbackslash{}\allowbreak{}tau\$\allowbreak{}}{DER\_\allowbreak{}VII-b:\allowbreak{}33}
\DERapparatusitem{ENV0482}{55;\allowbreak{}\allowbreak{} ms85,\allowbreak{}\allowbreak{} source\allowbreak{} of\allowbreak{} triv\_\allowbreak{}alpha}{\{\allowbreak{}}{}{DER\_\allowbreak{}VII-b:\allowbreak{}35}
\DERapparatusitem{ENV0483}{55;\allowbreak{}\allowbreak{} ms85,\allowbreak{}\allowbreak{} source\allowbreak{} of\allowbreak{} triv\_\allowbreak{}alpha}{\}\allowbreak{}}{}{DER\_\allowbreak{}VII-b:\allowbreak{}35}
\DERapparatusitem{ENV0484}{57;\allowbreak{}\allowbreak{} ms90,\allowbreak{}\allowbreak{} first\allowbreak{} inductive\allowbreak{} subinterval}{1}{0}{DER\_\allowbreak{}VII-b:\allowbreak{}36}
\DERapparatusitem{ENV0485}{57;\allowbreak{}\allowbreak{} ms90,\allowbreak{}\allowbreak{} length\allowbreak{} of\allowbreak{} a\allowbreak{} nonsingleton\allowbreak{} first\allowbreak{} run}{i.\allowbreak{}e.\allowbreak{}\allowbreak{} is\allowbreak{} of\allowbreak{} length\allowbreak{} 2,\allowbreak{}}{i.\allowbreak{}e.\allowbreak{}\allowbreak{} is\allowbreak{} of\allowbreak{} length\allowbreak{} at\allowbreak{} least\allowbreak{} 2,\allowbreak{}}{DER\_\allowbreak{}VII-b:\allowbreak{}37}
\DERapparatusitem{ENV0486}{58;\allowbreak{}\allowbreak{} ms91,\allowbreak{}\allowbreak{} opening\allowbreak{} sentence\allowbreak{} after\allowbreak{} first\allowbreak{} two\allowbreak{} cases}{a\allowbreak{} way\allowbreak{} to\allowbreak{} define\allowbreak{} \$\allowbreak{}F\$\allowbreak{}\allowbreak{} so\allowbreak{} as\allowbreak{} to\allowbreak{} satisfy\allowbreak{} a)\allowbreak{}}{a\allowbreak{} way\allowbreak{} to\allowbreak{} deform\allowbreak{} \$\allowbreak{}F\$\allowbreak{}\allowbreak{} so\allowbreak{} as\allowbreak{} to\allowbreak{} satisfy\allowbreak{} a)\allowbreak{}}{DER\_\allowbreak{}VII-b:\allowbreak{}38}
\DERapparatusitem{ENV0487}{58;\allowbreak{}\allowbreak{} ms92,\allowbreak{}\allowbreak{} second\allowbreak{} middle-vertex\allowbreak{} value}{or\allowbreak{} to\allowbreak{} 1\allowbreak{} \textbackslash{}\allowbreak{}texttt\{\allowbreak{}[\allowbreak{}rather\allowbreak{} to\allowbreak{} 0]\allowbreak{}\}\allowbreak{}}{or\allowbreak{} to\allowbreak{} 0\allowbreak{} \textbackslash{}\allowbreak{}texttt\{\allowbreak{}[\allowbreak{}rather\allowbreak{} to\allowbreak{} 0]\allowbreak{}\}\allowbreak{}}{DER\_\allowbreak{}VII-b:\allowbreak{}39}
\DERapparatusitem{ENV0488}{59;\allowbreak{}\allowbreak{} ms93,\allowbreak{}\allowbreak{} first\allowbreak{} reversal\allowbreak{} display}{'}{}{DER\_\allowbreak{}VII-b:\allowbreak{}42}
\DERapparatusitem{ENV0489}{59;\allowbreak{}\allowbreak{} ms93,\allowbreak{}\allowbreak{} first\allowbreak{} reversal\allowbreak{} display}{}{'}{DER\_\allowbreak{}VII-b:\allowbreak{}42}
\DERapparatusitem{ENV0490}{62;\allowbreak{}\allowbreak{} ms97,\allowbreak{}\allowbreak{} paragraph\allowbreak{} introducing\allowbreak{} endpoint\allowbreak{} projection}{categories\allowbreak{} \textbackslash{}\allowbreak{}textit\{\allowbreak{}over\allowbreak{} \$\allowbreak{}X\$\allowbreak{},\allowbreak{}\}\allowbreak{}}{categories\allowbreak{} \textbackslash{}\allowbreak{}textit\{\allowbreak{}over\allowbreak{} \$\allowbreak{}X\textbackslash{}\allowbreak{}times\allowbreak{} X\$\allowbreak{},\allowbreak{}\}\allowbreak{}}{DER\_\allowbreak{}VII-b:\allowbreak{}44}
\DERapparatusitem{ENV0491}{64;\allowbreak{}\allowbreak{} ms99,\allowbreak{}\allowbreak{} parenthetical\allowbreak{} characterization\allowbreak{} of\allowbreak{} ordered\allowbreak{} categories}{where\allowbreak{} every\allowbreak{} arrow\allowbreak{} is\allowbreak{} monomorphic}{where\allowbreak{} always\allowbreak{} \$\allowbreak{}\textbackslash{}\allowbreak{}\#\textbackslash{}\allowbreak{}Hom(\allowbreak{}a,\allowbreak{}b)\allowbreak{}\textbackslash{}\allowbreak{}le1\$\allowbreak{}}{DER\_\allowbreak{}VII-b:\allowbreak{}45}
\DERapparatusitem{ENV0492}{66;\allowbreak{}\allowbreak{} ms102,\allowbreak{}\allowbreak{} transition\allowbreak{} inside\allowbreak{} the\allowbreak{} fully\allowbreak{} faithful\allowbreak{} paragraph}{}{-1}{DER\_\allowbreak{}VII-b:\allowbreak{}46}
\DERapparatusitem{ENV0493}{66;\allowbreak{}\allowbreak{} ms102,\allowbreak{}\allowbreak{} transition\allowbreak{} inside\allowbreak{} the\allowbreak{} fully\allowbreak{} faithful\allowbreak{} paragraph}{-1}{}{DER\_\allowbreak{}VII-b:\allowbreak{}46}
\DERapparatusitem{ENV0494}{66;\allowbreak{}\allowbreak{} ms103,\allowbreak{}\allowbreak{} slice-preorder\allowbreak{} display}{}{\_\allowbreak{}\{\allowbreak{}0\}\allowbreak{}}{DER\_\allowbreak{}VII-b:\allowbreak{}51}
\DERapparatusitem{ENV0495}{66;\allowbreak{}\allowbreak{} ms103,\allowbreak{}\allowbreak{} slice-preorder\allowbreak{} display}{=\allowbreak{}}{\textbackslash{}\allowbreak{}longrightarrow}{DER\_\allowbreak{}VII-b:\allowbreak{}51}
\DERapparatusitem{ENV0496}{66;\allowbreak{}\allowbreak{} ms103,\allowbreak{}\allowbreak{} slice-preorder\allowbreak{} display}{}{\_\allowbreak{}\{\allowbreak{}0\}\allowbreak{}}{DER\_\allowbreak{}VII-b:\allowbreak{}51}
\DERapparatusitem{ENV0497}{73;\allowbreak{}\allowbreak{} ms116,\allowbreak{}\allowbreak{} defining\allowbreak{} H\_\allowbreak{}T\textasciicircum{}\allowbreak{}Y}{Cat\textbackslash{}\allowbreak{},\allowbreak{}hot}{\{\allowbreak{}Ind\}\allowbreak{}(\allowbreak{}\textbackslash{}\allowbreak{}mathrm\{\allowbreak{}Cat\textbackslash{}\allowbreak{},\allowbreak{}hot\}\allowbreak{})\allowbreak{}}{DER\_\allowbreak{}VII-b:\allowbreak{}57}
\DERapparatusitem{ENV0498}{73;\allowbreak{}\allowbreak{} running\allowbreak{} text\allowbreak{} or\allowbreak{} formula}{\textbackslash{}\allowbreak{}text\{\allowbreak{}def\}\allowbreak{}}{\textbackslash{}\allowbreak{}mathrm\{\allowbreak{}def\}\allowbreak{}}{TYPESETTING:\allowbreak{}7;\allowbreak{}\allowbreak{} NOTATION\_\allowbreak{}DER:\allowbreak{}17}
\DERapparatusitem{ENV0499}{73;\allowbreak{}\allowbreak{} ms117,\allowbreak{}\allowbreak{} final\allowbreak{} displayed\allowbreak{} localization\allowbreak{} formula}{H\_\allowbreak{}\{\allowbreak{}T\}\allowbreak{}(\allowbreak{}a,\allowbreak{}b)\allowbreak{}\allowbreak{} \textbackslash{}\allowbreak{}simeq\allowbreak{} \textbackslash{}\allowbreak{}mathrm\{\allowbreak{}Hom\}\allowbreak{}\_\allowbreak{}\{\allowbreak{}X(\allowbreak{}Y)\allowbreak{}\textbackslash{}\allowbreak{}Sigma\_\allowbreak{}\{\allowbreak{}Y\}\allowbreak{}\textasciicircum{}\allowbreak{}\{\allowbreak{}-1\}\allowbreak{}\}\allowbreak{}(\allowbreak{}a\_\allowbreak{}\{\allowbreak{}Y\}\allowbreak{},\allowbreak{}\allowbreak{} b\_\allowbreak{}\{\allowbreak{}Y\}\allowbreak{})\allowbreak{}}{H\_\allowbreak{}\{\allowbreak{}T\}\allowbreak{}\textasciicircum{}\allowbreak{}\{\allowbreak{}Y\}\allowbreak{}(\allowbreak{}a,\allowbreak{}b)\allowbreak{}\allowbreak{} \textbackslash{}\allowbreak{}simeq\allowbreak{} \textbackslash{}\allowbreak{}mathrm\{\allowbreak{}Hom\}\allowbreak{}\_\allowbreak{}\{\allowbreak{}X(\allowbreak{}Y)\allowbreak{}\textbackslash{}\allowbreak{}Sigma\_\allowbreak{}\{\allowbreak{}Y\}\allowbreak{}\textasciicircum{}\allowbreak{}\{\allowbreak{}-1\}\allowbreak{}\}\allowbreak{}(\allowbreak{}a\_\allowbreak{}\{\allowbreak{}Y\}\allowbreak{},\allowbreak{}\allowbreak{} b\_\allowbreak{}\{\allowbreak{}Y\}\allowbreak{})\allowbreak{}}{DER\_\allowbreak{}VII-b:\allowbreak{}59}
\DERapparatusitem{ENV0500}{76;\allowbreak{}\allowbreak{} ms121:\allowbreak{}\allowbreak{} ind-object\allowbreak{} comparison}{they\allowbreak{} define\allowbreak{} a\allowbreak{} homomorphism\allowbreak{} on\allowbreak{} the\allowbreak{} associated\allowbreak{} ind-objects}{they\allowbreak{} define\allowbreak{} an\allowbreak{} isomorphism\allowbreak{} on\allowbreak{} the\allowbreak{} associated\allowbreak{} ind-objects}{DER\_\allowbreak{}VII-c:\allowbreak{}3}
\DERapparatusitem{ENV0501}{79;\allowbreak{}\allowbreak{} ms126:\allowbreak{}\allowbreak{} common\allowbreak{} refinement\allowbreak{} of\allowbreak{} two\allowbreak{} interval\allowbreak{} types}{\textbackslash{}\allowbreak{}times}{*}{DER\_\allowbreak{}VII-c:\allowbreak{}5}
\DERapparatusitem{ENV0502}{80;\allowbreak{}\allowbreak{} ms127:\allowbreak{}\allowbreak{} printed\allowbreak{} misspelling}{the\allowbreak{} composititon\allowbreak{} defined\allowbreak{} here}{the\allowbreak{} composition\allowbreak{} defined\allowbreak{} here}{DER\_\allowbreak{}VII-c:\allowbreak{}8}
\DERapparatusitem{ENV0503}{81;\allowbreak{}\allowbreak{} ms.\allowbreak{}129,\allowbreak{}\allowbreak{} last\allowbreak{} display}{\textbackslash{}\allowbreak{}mathrm\{\allowbreak{}Hom\}\allowbreak{}\_\allowbreak{}\{\allowbreak{}\textbackslash{}\allowbreak{}mathrm\{\allowbreak{}Cat\textbackslash{}\allowbreak{},\allowbreak{}hot\}\allowbreak{}\}\allowbreak{}(\allowbreak{}Y,\allowbreak{}H(\allowbreak{}a,\allowbreak{}b)\allowbreak{})\allowbreak{}}{\textbackslash{}\allowbreak{}mathrm\{\allowbreak{}Hom\}\allowbreak{}\_\allowbreak{}\{\allowbreak{}\textbackslash{}\allowbreak{}mathrm\{\allowbreak{}IndCat\textbackslash{}\allowbreak{},\allowbreak{}hot\}\allowbreak{}\}\allowbreak{}(\allowbreak{}Y,\allowbreak{}H(\allowbreak{}a,\allowbreak{}b)\allowbreak{})\allowbreak{}}{DER\_\allowbreak{}VII-c:\allowbreak{}9}
\DERapparatusitem{ENV0504}{81;\allowbreak{}\allowbreak{} ms.\allowbreak{}129,\allowbreak{}\allowbreak{} underbrace\allowbreak{} in\allowbreak{} last\allowbreak{} display}{X}{Y}{DER\_\allowbreak{}VII-c:\allowbreak{}10}
\DERapparatusitem{ENV0505}{81;\allowbreak{}\allowbreak{} ms.\allowbreak{}129,\allowbreak{}\allowbreak{} underbrace\allowbreak{} in\allowbreak{} last\allowbreak{} display}{X}{Y}{DER\_\allowbreak{}VII-c:\allowbreak{}10}
\DERapparatusitem{ENV0506}{82;\allowbreak{}\allowbreak{} ms.\allowbreak{}131,\allowbreak{}\allowbreak{} boxed\allowbreak{} display}{X}{Y}{DER\_\allowbreak{}VII-c:\allowbreak{}12}
\DERapparatusitem{ENV0507}{82;\allowbreak{}\allowbreak{} ms.\allowbreak{}131,\allowbreak{}\allowbreak{} boxed\allowbreak{} display}{X}{Y}{DER\_\allowbreak{}VII-c:\allowbreak{}12}
\DERapparatusitem{ENV0508}{82;\allowbreak{}\allowbreak{} ms.\allowbreak{}131,\allowbreak{}\allowbreak{} adjunction\allowbreak{} targets}{cal}{bf}{DER\_\allowbreak{}VII-c:\allowbreak{}13}
\DERapparatusitem{ENV0509}{82;\allowbreak{}\allowbreak{} ms.\allowbreak{}131,\allowbreak{}\allowbreak{} adjunction\allowbreak{} targets}{M}{D}{DER\_\allowbreak{}VII-c:\allowbreak{}13}
\DERapparatusitem{ENV0510}{82;\allowbreak{}\allowbreak{} ms.\allowbreak{}131,\allowbreak{}\allowbreak{} adjunction\allowbreak{} targets}{}{(\allowbreak{}*)\allowbreak{}}{DER\_\allowbreak{}VII-c:\allowbreak{}13}
\DERapparatusitem{ENV0511}{82;\allowbreak{}\allowbreak{} ms.\allowbreak{}131,\allowbreak{}\allowbreak{} adjunction\allowbreak{} targets}{}{\textbackslash{}\allowbreak{}Sigma\textasciicircum{}\allowbreak{}\{\allowbreak{}-1\}\allowbreak{}}{DER\_\allowbreak{}VII-c:\allowbreak{}13}
\DERapparatusitem{ENV0512}{82;\allowbreak{}\allowbreak{} ms.\allowbreak{}131,\allowbreak{}\allowbreak{} adjunction\allowbreak{} targets}{cal}{bf}{DER\_\allowbreak{}VII-c:\allowbreak{}13}
\DERapparatusitem{ENV0513}{82;\allowbreak{}\allowbreak{} ms.\allowbreak{}131,\allowbreak{}\allowbreak{} adjunction\allowbreak{} targets}{M}{D}{DER\_\allowbreak{}VII-c:\allowbreak{}13}
\DERapparatusitem{ENV0514}{82;\allowbreak{}\allowbreak{} ms.\allowbreak{}131,\allowbreak{}\allowbreak{} adjunction\allowbreak{} targets}{}{(\allowbreak{}*)\allowbreak{}}{DER\_\allowbreak{}VII-c:\allowbreak{}13}
\DERapparatusitem{ENV0515}{82;\allowbreak{}\allowbreak{} ms.\allowbreak{}131,\allowbreak{}\allowbreak{} adjunction\allowbreak{} targets}{}{\textbackslash{}\allowbreak{}Sigma\textasciicircum{}\allowbreak{}\{\allowbreak{}-1\}\allowbreak{}}{DER\_\allowbreak{}VII-c:\allowbreak{}13}
\DERapparatusitem{ENV0516}{83;\allowbreak{}\allowbreak{} running\allowbreak{} text\allowbreak{} or\allowbreak{} formula}{\textbackslash{}\allowbreak{}text\{\allowbreak{}def\}\allowbreak{}}{\textbackslash{}\allowbreak{}mathrm\{\allowbreak{}def\}\allowbreak{}}{TYPESETTING:\allowbreak{}7;\allowbreak{}\allowbreak{} NOTATION\_\allowbreak{}DER:\allowbreak{}17}
\DERapparatusitem{ENV0517}{83;\allowbreak{}\allowbreak{} ms.\allowbreak{}133,\allowbreak{}\allowbreak{} bounded-degree\allowbreak{} sentence}{hence\allowbreak{} at\allowbreak{} least\allowbreak{} by\allowbreak{} 4}{hence\allowbreak{} at\allowbreak{} most\allowbreak{} by\allowbreak{} 4}{DER\_\allowbreak{}VII-c:\allowbreak{}14}
\DERapparatusitem{ENV0518}{83;\allowbreak{}\allowbreak{} ms.\allowbreak{}134,\allowbreak{}\allowbreak{} circle/\allowbreak{}sphere\allowbreak{} indexing}{=\allowbreak{}0}{\textbackslash{}\allowbreak{}geq2}{DER\_\allowbreak{}VII-c:\allowbreak{}16}
\DERapparatusitem{ENV0519}{86;\allowbreak{}\allowbreak{} ms.\allowbreak{}139,\allowbreak{}\allowbreak{} proposed\allowbreak{} factorization}{where\allowbreak{} \$\allowbreak{}\textbackslash{}\allowbreak{}bar\{\allowbreak{}f\}\allowbreak{}\$\allowbreak{}\allowbreak{} should\allowbreak{} be\allowbreak{} a\allowbreak{} \textbackslash{}\allowbreak{}textit\{\allowbreak{}``fibrant''\}\allowbreak{}\allowbreak{} arrow\allowbreak{} and\allowbreak{} a\allowbreak{} Hot-equivalence,\allowbreak{}\allowbreak{} \$\allowbreak{}\textbackslash{}\allowbreak{}alpha\$\allowbreak{}\allowbreak{} a\allowbreak{} \textbackslash{}\allowbreak{}textit\{\allowbreak{}``cofibrant''\}\allowbreak{}\allowbreak{} arrow\allowbreak{} in\allowbreak{} the\allowbreak{} sense\allowbreak{} of\allowbreak{} \textbackslash{}\allowbreak{}textsc\{\allowbreak{}Thomason\}\allowbreak{}.\allowbreak{}}{where\allowbreak{} \$\allowbreak{}\textbackslash{}\allowbreak{}bar\{\allowbreak{}f\}\allowbreak{}\$\allowbreak{}\allowbreak{} should\allowbreak{} be\allowbreak{} a\allowbreak{} \textbackslash{}\allowbreak{}textit\{\allowbreak{}``fibrant''\}\allowbreak{}\allowbreak{} arrow,\allowbreak{}\allowbreak{} \$\allowbreak{}\textbackslash{}\allowbreak{}alpha\$\allowbreak{}\allowbreak{} a\allowbreak{} Hot-equivalence\allowbreak{} and\allowbreak{} a\allowbreak{} \textbackslash{}\allowbreak{}textit\{\allowbreak{}``cofibrant''\}\allowbreak{}\allowbreak{} arrow\allowbreak{} in\allowbreak{} the\allowbreak{} sense\allowbreak{} of\allowbreak{} \textbackslash{}\allowbreak{}textsc\{\allowbreak{}Thomason\}\allowbreak{}.\allowbreak{}}{DER\_\allowbreak{}VII-c:\allowbreak{}18}
\DERapparatusitem{ENV0520}{87;\allowbreak{}\allowbreak{} ms.\allowbreak{}140,\allowbreak{}\allowbreak{} factorization\allowbreak{} cross-reference}{as\allowbreak{} on\allowbreak{} page\allowbreak{} 38\allowbreak{} to\allowbreak{} define}{as\allowbreak{} on\allowbreak{} page\allowbreak{} 138\allowbreak{} to\allowbreak{} define}{DER\_\allowbreak{}VII-c:\allowbreak{}20}
\DERapparatusitem{ENV0521}{87;\allowbreak{}\allowbreak{} ms.\allowbreak{}141,\allowbreak{}\allowbreak{} last\allowbreak{} line}{\textbackslash{}\allowbreak{}overline}{}{DER\_\allowbreak{}VII-c:\allowbreak{}22}
\DERapparatusitem{ENV0522}{90;\allowbreak{}\allowbreak{} ms.\allowbreak{}144,\allowbreak{}\allowbreak{} pair\allowbreak{} defining\allowbreak{} the\allowbreak{} lifting\allowbreak{} category}{;\allowbreak{}\allowbreak{} a',\allowbreak{}\allowbreak{} b'}{}{DER\_\allowbreak{}VII-c:\allowbreak{}24}
\DERapparatusitem{ENV0523}{90;\allowbreak{}\allowbreak{} ms.\allowbreak{}145,\allowbreak{}\allowbreak{} underbrace\allowbreak{} under\allowbreak{} t\_\allowbreak{}n}{}{2}{DER\_\allowbreak{}VII-c:\allowbreak{}26}
\DERapparatusitem{ENV0524}{91;\allowbreak{}\allowbreak{} ms.\allowbreak{}145,\allowbreak{}\allowbreak{} transition\allowbreak{} map}{o}{0}{DER\_\allowbreak{}VII-c:\allowbreak{}27}
\DERapparatusitem{ENV0525}{91;\allowbreak{}\allowbreak{} ms.\allowbreak{}145,\allowbreak{}\allowbreak{} support\allowbreak{} of\allowbreak{} Ch\_\allowbreak{}tn}{i}{n}{DER\_\allowbreak{}VII-c:\allowbreak{}28}
\DERapparatusitem{ENV0526}{91;\allowbreak{}\allowbreak{} ms.\allowbreak{}145,\allowbreak{}\allowbreak{} support\allowbreak{} of\allowbreak{} Ch\_\allowbreak{}tn}{i}{n}{DER\_\allowbreak{}VII-c:\allowbreak{}28}
\DERapparatusitem{ENV0527}{91;\allowbreak{}\allowbreak{} ms.\allowbreak{}145,\allowbreak{}\allowbreak{} support\allowbreak{} of\allowbreak{} Ch\_\allowbreak{}tn}{i}{n}{DER\_\allowbreak{}VII-c:\allowbreak{}28}
\DERapparatusitem{ENV0528}{91;\allowbreak{}\allowbreak{} ms.\allowbreak{}145,\allowbreak{}\allowbreak{} support\allowbreak{} of\allowbreak{} Ch\_\allowbreak{}tn}{i}{n}{DER\_\allowbreak{}VII-c:\allowbreak{}28}
\DERapparatusitem{ENV0529}{91;\allowbreak{}\allowbreak{} ms.\allowbreak{}145,\allowbreak{}\allowbreak{} support\allowbreak{} of\allowbreak{} Ch\_\allowbreak{}tn}{i}{n}{DER\_\allowbreak{}VII-c:\allowbreak{}28}
\DERapparatusitem{ENV0530}{91;\allowbreak{}\allowbreak{} ms.\allowbreak{}145,\allowbreak{}\allowbreak{} support\allowbreak{} of\allowbreak{} Ch\_\allowbreak{}tn}{i}{n}{DER\_\allowbreak{}VII-c:\allowbreak{}28}
\DERapparatusitem{ENV0531}{91;\allowbreak{}\allowbreak{} ms.\allowbreak{}146,\allowbreak{}\allowbreak{} support\allowbreak{} interval\allowbreak{} and\allowbreak{} labelled\allowbreak{} edge\allowbreak{} display}{}{+1}{DER\_\allowbreak{}VII-c:\allowbreak{}30}
\DERapparatusitem{ENV0532}{91;\allowbreak{}\allowbreak{} ms.\allowbreak{}146,\allowbreak{}\allowbreak{} support\allowbreak{} interval\allowbreak{} and\allowbreak{} labelled\allowbreak{} edge\allowbreak{} display}{-1}{}{DER\_\allowbreak{}VII-c:\allowbreak{}30}
\DERapparatusitem{ENV0533}{91;\allowbreak{}\allowbreak{} ms.\allowbreak{}146,\allowbreak{}\allowbreak{} support\allowbreak{} interval\allowbreak{} and\allowbreak{} labelled\allowbreak{} edge\allowbreak{} display}{}{\{\allowbreak{}}{DER\_\allowbreak{}VII-c:\allowbreak{}30}
\DERapparatusitem{ENV0534}{91;\allowbreak{}\allowbreak{} ms.\allowbreak{}146,\allowbreak{}\allowbreak{} support\allowbreak{} interval\allowbreak{} and\allowbreak{} labelled\allowbreak{} edge\allowbreak{} display}{}{+1\}\allowbreak{}}{DER\_\allowbreak{}VII-c:\allowbreak{}30}
\DERapparatusitem{ENV0535}{91;\allowbreak{}\allowbreak{} ms.\allowbreak{}146,\allowbreak{}\allowbreak{} support\allowbreak{} interval\allowbreak{} and\allowbreak{} labelled\allowbreak{} edge\allowbreak{} display}{\{\allowbreak{}}{}{DER\_\allowbreak{}VII-c:\allowbreak{}30}
\DERapparatusitem{ENV0536}{91;\allowbreak{}\allowbreak{} ms.\allowbreak{}146,\allowbreak{}\allowbreak{} support\allowbreak{} interval\allowbreak{} and\allowbreak{} labelled\allowbreak{} edge\allowbreak{} display}{-1\}\allowbreak{}}{}{DER\_\allowbreak{}VII-c:\allowbreak{}30}
\DERapparatusitem{ENV0537}{91;\allowbreak{}\allowbreak{} ms.\allowbreak{}146,\allowbreak{}\allowbreak{} n\allowbreak{} even\allowbreak{} m\allowbreak{} odd\allowbreak{} exponent}{(\allowbreak{}\textbackslash{}\allowbreak{}Delta\_\allowbreak{}+\allowbreak{} \textbackslash{}\allowbreak{}Delta\_\allowbreak{}-)\allowbreak{}\textasciicircum{}\allowbreak{}\{\allowbreak{}\textbackslash{}\allowbreak{}frac\{\allowbreak{}m-n\}\allowbreak{}\{\allowbreak{}2\}\allowbreak{}-1\}\allowbreak{}}{(\allowbreak{}\textbackslash{}\allowbreak{}Delta\_\allowbreak{}+\allowbreak{} \textbackslash{}\allowbreak{}Delta\_\allowbreak{}-)\allowbreak{}\textasciicircum{}\allowbreak{}\{\allowbreak{}\textbackslash{}\allowbreak{}frac\{\allowbreak{}m-n-1\}\allowbreak{}\{\allowbreak{}2\}\allowbreak{}\}\allowbreak{}}{DER\_\allowbreak{}VII-c:\allowbreak{}31}
\DERapparatusitem{ENV0538}{92;\allowbreak{}\allowbreak{} ms.\allowbreak{}147,\allowbreak{}\allowbreak{} n\allowbreak{} odd\allowbreak{} m\allowbreak{} even\allowbreak{} exponent}{(\allowbreak{}\textbackslash{}\allowbreak{}Delta\_\allowbreak{}-\allowbreak{} \textbackslash{}\allowbreak{}Delta\_\allowbreak{}+)\allowbreak{}\textasciicircum{}\allowbreak{}\{\allowbreak{}\textbackslash{}\allowbreak{}frac\{\allowbreak{}m-n\}\allowbreak{}\{\allowbreak{}2\}\allowbreak{}-1\}\allowbreak{}}{(\allowbreak{}\textbackslash{}\allowbreak{}Delta\_\allowbreak{}-\allowbreak{} \textbackslash{}\allowbreak{}Delta\_\allowbreak{}+)\allowbreak{}\textasciicircum{}\allowbreak{}\{\allowbreak{}\textbackslash{}\allowbreak{}frac\{\allowbreak{}m-n-1\}\allowbreak{}\{\allowbreak{}2\}\allowbreak{}\}\allowbreak{}}{DER\_\allowbreak{}VII-c:\allowbreak{}32}
\DERapparatusitem{ENV0539}{92;\allowbreak{}\allowbreak{} ms.\allowbreak{}147,\allowbreak{}\allowbreak{} terminal\allowbreak{} direction}{\$\allowbreak{}(\allowbreak{}-1)\allowbreak{}\textasciicircum{}\allowbreak{}m\$\allowbreak{}}{\$\allowbreak{}(\allowbreak{}-1)\allowbreak{}\textasciicircum{}\allowbreak{}\{\allowbreak{}m+1\}\allowbreak{}\$\allowbreak{}}{DER\_\allowbreak{}VII-c:\allowbreak{}33}
\DERapparatusitem{ENV0540}{96;\allowbreak{}\allowbreak{} ms155,\allowbreak{}\allowbreak{} paragraph\allowbreak{} introducing\allowbreak{} v\textasciicircum{}\allowbreak{}*}{Working\allowbreak{} with\allowbreak{} odd\allowbreak{} \$\allowbreak{}m\$\allowbreak{}}{Working\allowbreak{} with\allowbreak{} even\allowbreak{} \$\allowbreak{}m\$\allowbreak{}}{DER\_\allowbreak{}VII-c:\allowbreak{}37}
\DERapparatusitem{ENV0541}{98;\allowbreak{}\allowbreak{} ms157,\allowbreak{}\allowbreak{} definition\allowbreak{} of\allowbreak{} constant-path\allowbreak{} section}{X}{x}{DER\_\allowbreak{}VII-c:\allowbreak{}38}
\DERapparatusitem{ENV0542}{98;\allowbreak{}\allowbreak{} running\allowbreak{} text\allowbreak{} or\allowbreak{} formula}{projective\allowbreak{} system}{inverse\allowbreak{} system}{GLOSSARY\_\allowbreak{}DER:\allowbreak{}91}
\DERapparatusitem{ENV0543}{100;\allowbreak{}\allowbreak{} running\allowbreak{} text\allowbreak{} or\allowbreak{} formula}{\textbackslash{}\allowbreak{}text\{\allowbreak{}def\}\allowbreak{}}{\textbackslash{}\allowbreak{}mathrm\{\allowbreak{}def\}\allowbreak{}}{TYPESETTING:\allowbreak{}7;\allowbreak{}\allowbreak{} NOTATION\_\allowbreak{}DER:\allowbreak{}17}
\DERapparatusitem{ENV0544}{100;\allowbreak{}\allowbreak{} ms160–161,\allowbreak{}\allowbreak{} repeated\allowbreak{} verb\allowbreak{} across\allowbreak{} marker}{\textbackslash{}\allowbreak{}msp\{\allowbreak{}page\allowbreak{} 161\}\allowbreak{}\space{}\allowbreak{}is\allowbreak{} a\allowbreak{} Hot-equivalence.\allowbreak{}}{\textbackslash{}\allowbreak{}msp\{\allowbreak{}page\allowbreak{} 161\}\allowbreak{}\space{}\allowbreak{}a\allowbreak{} Hot-equivalence.\allowbreak{}}{DER\_\allowbreak{}VII-c:\allowbreak{}41}
\DERapparatusitem{ENV0545}{100;\allowbreak{}\allowbreak{} ms161,\allowbreak{}\allowbreak{} introduction\allowbreak{} to\allowbreak{} proof\allowbreak{} (\allowbreak{}2)\allowbreak{}}{the\allowbreak{} property\allowbreak{} will\allowbreak{} follow\allowbreak{} from\allowbreak{} it}{properness\allowbreak{} will\allowbreak{} follow\allowbreak{} from\allowbreak{} it}{DER\_\allowbreak{}VII-c:\allowbreak{}42}
\DERapparatusitem{ENV0546}{101;\allowbreak{}\allowbreak{} ms162,\allowbreak{}\allowbreak{} lower\allowbreak{} row\allowbreak{} of\allowbreak{} padding\allowbreak{} diagram}{'}{}{DER\_\allowbreak{}VII-c:\allowbreak{}43}
\DERapparatusitem{ENV0547}{101;\allowbreak{}\allowbreak{} ms162,\allowbreak{}\allowbreak{} lower\allowbreak{} row\allowbreak{} of\allowbreak{} padding\allowbreak{} diagram}{'}{}{DER\_\allowbreak{}VII-c:\allowbreak{}43}
\DERapparatusitem{ENV0548}{101;\allowbreak{}\allowbreak{} ms162,\allowbreak{}\allowbreak{} lower\allowbreak{} row\allowbreak{} of\allowbreak{} padding\allowbreak{} diagram}{'}{}{DER\_\allowbreak{}VII-c:\allowbreak{}43}
\DERapparatusitem{ENV0549}{101;\allowbreak{}\allowbreak{} ms163,\allowbreak{}\allowbreak{} equation\allowbreak{} defining\allowbreak{} the\allowbreak{} padded\allowbreak{} morphism}{\textbackslash{}\allowbreak{}tilde\{\allowbreak{}w\}\allowbreak{}\allowbreak{} =\allowbreak{}\allowbreak{} \textbackslash{}\allowbreak{}lambda(\allowbreak{}u)\allowbreak{}\allowbreak{} \textbackslash{}\allowbreak{}ast\allowbreak{} w'\allowbreak{} \textbackslash{}\allowbreak{}ast\allowbreak{} \textbackslash{}\allowbreak{}mu(\allowbreak{}v)\allowbreak{}}{\textbackslash{}\allowbreak{}tilde\{\allowbreak{}w\}\allowbreak{}'\allowbreak{} =\allowbreak{}\allowbreak{} \textbackslash{}\allowbreak{}lambda(\allowbreak{}u)\allowbreak{}\allowbreak{} \textbackslash{}\allowbreak{}ast\allowbreak{} w'\allowbreak{} \textbackslash{}\allowbreak{}ast\allowbreak{} \textbackslash{}\allowbreak{}mu(\allowbreak{}v)\allowbreak{}}{DER\_\allowbreak{}VII-c:\allowbreak{}44}
\DERapparatusitem{ENV0550}{102;\allowbreak{}\allowbreak{} ms164,\allowbreak{}\allowbreak{} lower\allowbreak{} brace\allowbreak{} in\allowbreak{} final\allowbreak{} diagram}{}{\textbackslash{}\allowbreak{}tilde\{\allowbreak{}}{DER\_\allowbreak{}VII-c:\allowbreak{}46}
\DERapparatusitem{ENV0551}{102;\allowbreak{}\allowbreak{} ms164,\allowbreak{}\allowbreak{} lower\allowbreak{} brace\allowbreak{} in\allowbreak{} final\allowbreak{} diagram}{}{\}\allowbreak{}}{DER\_\allowbreak{}VII-c:\allowbreak{}46}
\DERapparatusitem{ENV0552}{103;\allowbreak{}\allowbreak{} ms165,\allowbreak{}\allowbreak{} diagram\allowbreak{} of\allowbreak{} rho\_\allowbreak{}c,\allowbreak{}\allowbreak{} middle\allowbreak{} row\allowbreak{} final\allowbreak{} vertex\allowbreak{} and\allowbreak{} arrow}{'}{}{DER\_\allowbreak{}VII-c:\allowbreak{}47}
\DERapparatusitem{ENV0553}{103;\allowbreak{}\allowbreak{} ms165,\allowbreak{}\allowbreak{} diagram\allowbreak{} of\allowbreak{} rho\_\allowbreak{}c,\allowbreak{}\allowbreak{} middle\allowbreak{} row\allowbreak{} final\allowbreak{} vertex\allowbreak{} and\allowbreak{} arrow}{v}{\textbackslash{}\allowbreak{}mathrm\{\allowbreak{}id\}\allowbreak{}}{DER\_\allowbreak{}VII-c:\allowbreak{}47}
\DERapparatusitem{ENV0554}{104;\allowbreak{}\allowbreak{} ms166,\allowbreak{}\allowbreak{} end\allowbreak{} of\allowbreak{} proof\allowbreak{} (\allowbreak{}1′)\allowbreak{}}{Thus\allowbreak{} \$\allowbreak{}j'i'\$\allowbreak{}\allowbreak{} is\allowbreak{} aspheric,\allowbreak{}\allowbreak{} q.\allowbreak{}e.\allowbreak{}d.\allowbreak{}}{Thus\allowbreak{} \$\allowbreak{}j'i'\$\allowbreak{}\allowbreak{} is\allowbreak{} a\allowbreak{} Hot-equivalence,\allowbreak{}\allowbreak{} q.\allowbreak{}e.\allowbreak{}d.\allowbreak{}}{DER\_\allowbreak{}VII-c:\allowbreak{}48}
\DERapparatusitem{ENV0555}{104;\allowbreak{}\allowbreak{} ms167,\allowbreak{}\allowbreak{} exponent\allowbreak{} in\allowbreak{} parenthesis}{t\_\allowbreak{}\{\allowbreak{}n+1\}\allowbreak{}\allowbreak{} =\allowbreak{}\allowbreak{} (\allowbreak{}\textbackslash{}\allowbreak{}Delta\_\allowbreak{}+\textbackslash{}\allowbreak{}Delta\_\allowbreak{}-)\allowbreak{}\textasciicircum{}\allowbreak{}\{\allowbreak{}2n\}\allowbreak{}}{t\_\allowbreak{}\{\allowbreak{}n+1\}\allowbreak{}\allowbreak{} =\allowbreak{}\allowbreak{} (\allowbreak{}\textbackslash{}\allowbreak{}Delta\_\allowbreak{}+\textbackslash{}\allowbreak{}Delta\_\allowbreak{}-)\allowbreak{}\textasciicircum{}\allowbreak{}\{\allowbreak{}2n+2\}\allowbreak{}}{DER\_\allowbreak{}VII-c:\allowbreak{}49}
\DERapparatusitem{ENV0556}{105;\allowbreak{}\allowbreak{} ms168,\allowbreak{}\allowbreak{} equation\allowbreak{} for\allowbreak{} tilde\allowbreak{} w′}{}{'}{DER\_\allowbreak{}VII-c:\allowbreak{}50}
\DERapparatusitem{ENV0557}{105;\allowbreak{}\allowbreak{} ms168,\allowbreak{}\allowbreak{} second\allowbreak{} component\allowbreak{} transformation}{la}{}{DER\_\allowbreak{}VII-c:\allowbreak{}51}
\DERapparatusitem{ENV0558}{105;\allowbreak{}\allowbreak{} ms168,\allowbreak{}\allowbreak{} second\allowbreak{} component\allowbreak{} transformation}{bda}{u}{DER\_\allowbreak{}VII-c:\allowbreak{}51}
\DERapparatusitem{ENV0559}{105;\allowbreak{}\allowbreak{} ms169,\allowbreak{}\allowbreak{} lower\allowbreak{} row\allowbreak{} of\allowbreak{} final\allowbreak{} diagram,\allowbreak{}\allowbreak{} final\allowbreak{} vertex}{'}{}{DER\_\allowbreak{}VII-c:\allowbreak{}52}
\DERapparatusitem{ENV0560}{106;\allowbreak{}\allowbreak{} ms169,\allowbreak{}\allowbreak{} properness\allowbreak{} paragraph}{prove\allowbreak{} the\allowbreak{} smoothness\allowbreak{} of\allowbreak{} the\allowbreak{} case}{prove\allowbreak{} the\allowbreak{} smoothness\allowbreak{} in\allowbreak{} the\allowbreak{} case}{DER\_\allowbreak{}VII-c:\allowbreak{}53}
\DERapparatusitem{ENV0561}{107;\allowbreak{}\allowbreak{} running\allowbreak{} text\allowbreak{} or\allowbreak{} formula}{\textbackslash{}\allowbreak{}text\{\allowbreak{}def\}\allowbreak{}}{\textbackslash{}\allowbreak{}mathrm\{\allowbreak{}def\}\allowbreak{}}{TYPESETTING:\allowbreak{}7;\allowbreak{}\allowbreak{} NOTATION\_\allowbreak{}DER:\allowbreak{}17}
\subsection*{Chapter VIII}
\DERapparatusitem{ENV0562}{2;\allowbreak{}\allowbreak{} running\allowbreak{} text\allowbreak{} or\allowbreak{} formula}{\textbackslash{}\allowbreak{}text\{\allowbreak{}def\}\allowbreak{}}{\textbackslash{}\allowbreak{}mathrm\{\allowbreak{}def\}\allowbreak{}}{TYPESETTING:\allowbreak{}7;\allowbreak{}\allowbreak{} NOTATION\_\allowbreak{}DER:\allowbreak{}17}
\DERapparatusitem{ENV0563}{4;\allowbreak{}\allowbreak{} running\allowbreak{} text\allowbreak{} or\allowbreak{} formula}{\textbackslash{}\allowbreak{}text\{\allowbreak{}def\}\allowbreak{}}{\textbackslash{}\allowbreak{}mathrm\{\allowbreak{}def\}\allowbreak{}}{TYPESETTING:\allowbreak{}7;\allowbreak{}\allowbreak{} NOTATION\_\allowbreak{}DER:\allowbreak{}17}
\DERapparatusitem{ENV0564}{4;\allowbreak{}\allowbreak{} VIII\allowbreak{} ms.\allowbreak{}\allowbreak{} 4,\allowbreak{}\allowbreak{} unit,\allowbreak{}\allowbreak{} all\allowbreak{} three\allowbreak{} comma-category\allowbreak{} expressions}{s\textbackslash{}\allowbreak{}backslash\allowbreak{} G}{f(\allowbreak{}s)\allowbreak{}\textbackslash{}\allowbreak{}backslash\_\allowbreak{}T\allowbreak{} G}{DER\_\allowbreak{}VIII-X:\allowbreak{}1}
\DERapparatusitem{ENV0565}{4;\allowbreak{}\allowbreak{} VIII\allowbreak{} ms.\allowbreak{}\allowbreak{} 4,\allowbreak{}\allowbreak{} unit,\allowbreak{}\allowbreak{} all\allowbreak{} three\allowbreak{} comma-category\allowbreak{} expressions}{s\textbackslash{}\allowbreak{}backslash\allowbreak{} G}{f(\allowbreak{}s)\allowbreak{}\textbackslash{}\allowbreak{}backslash\_\allowbreak{}T\allowbreak{} G}{DER\_\allowbreak{}VIII-X:\allowbreak{}1}
\DERapparatusitem{ENV0566}{4;\allowbreak{}\allowbreak{} VIII\allowbreak{} ms.\allowbreak{}\allowbreak{} 4,\allowbreak{}\allowbreak{} unit,\allowbreak{}\allowbreak{} all\allowbreak{} three\allowbreak{} comma-category\allowbreak{} expressions}{s\textbackslash{}\allowbreak{}backslash\allowbreak{} G}{f(\allowbreak{}s)\allowbreak{}\textbackslash{}\allowbreak{}backslash\_\allowbreak{}T\allowbreak{} G}{DER\_\allowbreak{}VIII-X:\allowbreak{}1}
\DERapparatusitem{ENV0567}{4;\allowbreak{}\allowbreak{} running\allowbreak{} text\allowbreak{} or\allowbreak{} formula}{\textbackslash{}\allowbreak{}text\{\allowbreak{}def\}\allowbreak{}}{\textbackslash{}\allowbreak{}mathrm\{\allowbreak{}def\}\allowbreak{}}{TYPESETTING:\allowbreak{}7;\allowbreak{}\allowbreak{} NOTATION\_\allowbreak{}DER:\allowbreak{}17}
\DERapparatusitem{ENV0568}{4;\allowbreak{}\allowbreak{} VIII\allowbreak{} ms.\allowbreak{}\allowbreak{} 5,\allowbreak{}\allowbreak{} pullback\allowbreak{} square}{S\allowbreak{} T}{T\allowbreak{} S}{DER\_\allowbreak{}VIII-X:\allowbreak{}2}
\DERapparatusitem{ENV0569}{5;\allowbreak{}\allowbreak{} running\allowbreak{} text\allowbreak{} or\allowbreak{} formula}{\textbackslash{}\allowbreak{}text\{\allowbreak{}def\}\allowbreak{}}{\textbackslash{}\allowbreak{}mathrm\{\allowbreak{}def\}\allowbreak{}}{TYPESETTING:\allowbreak{}7;\allowbreak{}\allowbreak{} NOTATION\_\allowbreak{}DER:\allowbreak{}17}
\DERapparatusitem{ENV0570}{6;\allowbreak{}\allowbreak{} running\allowbreak{} text\allowbreak{} or\allowbreak{} formula}{\allowbreak{} (\allowbreak{}gear)\allowbreak{}}{}{GLOSSARY\_\allowbreak{}DER:\allowbreak{}164}
\DERapparatusitem{ENV0571}{7;\allowbreak{}\allowbreak{} VIII\allowbreak{} ms.\allowbreak{}\allowbreak{} 8,\allowbreak{}\allowbreak{} Theorem\allowbreak{} 3(\allowbreak{}2)\allowbreak{}\allowbreak{} parenthesis}{(\allowbreak{}\$\allowbreak{}f\$\allowbreak{}\allowbreak{} fully\allowbreak{} faithful)\allowbreak{}}{(\allowbreak{}\$\allowbreak{}f\textasciicircum{}\allowbreak{}*\$\allowbreak{}\allowbreak{} 2-faithful)\allowbreak{}}{DER\_\allowbreak{}VIII-X:\allowbreak{}3}
\DERapparatusitem{ENV0572}{8;\allowbreak{}\allowbreak{} VIII\allowbreak{} ms.\allowbreak{}\allowbreak{} 9,\allowbreak{}\allowbreak{} Theorem\allowbreak{} 3(\allowbreak{}3)\allowbreak{},\allowbreak{}\allowbreak{} final\allowbreak{} words}{on\allowbreak{} the\allowbreak{} fibers\allowbreak{} of\allowbreak{} \$\allowbreak{}G\$\allowbreak{}\allowbreak{} be\allowbreak{} a\allowbreak{} \$\allowbreak{}Z\$\allowbreak{}-equivalence}{on\allowbreak{} the\allowbreak{} fibers\allowbreak{} of\allowbreak{} \$\allowbreak{}G\$\allowbreak{}\allowbreak{} be\allowbreak{} an\allowbreak{} equivalence\allowbreak{} of\allowbreak{} groupoids}{DER\_\allowbreak{}VIII-X:\allowbreak{}4}
\subsection*{Chapter IX}
\DERapparatusitem{ENV0198}{1;\allowbreak{}\allowbreak{} running\allowbreak{} text\allowbreak{} or\allowbreak{} formula}{\textbackslash{}\allowbreak{}text\{\allowbreak{}def\}\allowbreak{}}{\textbackslash{}\allowbreak{}mathrm\{\allowbreak{}def\}\allowbreak{}}{TYPESETTING:\allowbreak{}7;\allowbreak{}\allowbreak{} NOTATION\_\allowbreak{}DER:\allowbreak{}17}
\DERapparatusitem{ENV0199}{2;\allowbreak{}\allowbreak{} IX\allowbreak{} manuscript\allowbreak{} page\allowbreak{} 3,\allowbreak{}\allowbreak{} covariant\allowbreak{} endpoint\allowbreak{} display}{\textbackslash{}\allowbreak{}Ch\_\allowbreak{}\textbackslash{}\allowbreak{}infty(\allowbreak{}a,\allowbreak{}b)\allowbreak{}\allowbreak{} \textbackslash{}\allowbreak{}longrightarrow\allowbreak{} \textbackslash{}\allowbreak{}Ch\_\allowbreak{}\textbackslash{}\allowbreak{}infty(\allowbreak{}b')\allowbreak{}}{\textbackslash{}\allowbreak{}Ch\_\allowbreak{}\textbackslash{}\allowbreak{}infty(\allowbreak{}a\_\allowbreak{}0,\allowbreak{}b)\allowbreak{}\allowbreak{} \textbackslash{}\allowbreak{}longrightarrow\allowbreak{} \textbackslash{}\allowbreak{}Ch\_\allowbreak{}\textbackslash{}\allowbreak{}infty(\allowbreak{}a\_\allowbreak{}0,\allowbreak{}b')\allowbreak{}}{DER\_\allowbreak{}VIII-X:\allowbreak{}6}
\DERapparatusitem{ENV0200}{3;\allowbreak{}\allowbreak{} IX\allowbreak{} manuscript\allowbreak{} page\allowbreak{} 4,\allowbreak{}\allowbreak{} corollary}{\$\allowbreak{}v\allowbreak{} :\allowbreak{}\allowbreak{} b\allowbreak{} \textbackslash{}\allowbreak{}longrightarrow\allowbreak{} b'\$\allowbreak{}\allowbreak{} are\allowbreak{} in\allowbreak{} \$\allowbreak{}\textbackslash{}\allowbreak{}mathcal\{\allowbreak{}M\}\allowbreak{}\_\allowbreak{}\textbackslash{}\allowbreak{}Pi\$\allowbreak{}}{\$\allowbreak{}v\allowbreak{} :\allowbreak{}\allowbreak{} b\allowbreak{} \textbackslash{}\allowbreak{}longrightarrow\allowbreak{} b'\$\allowbreak{}\allowbreak{} are\allowbreak{} invertible\allowbreak{} in\allowbreak{} \$\allowbreak{}\textbackslash{}\allowbreak{}mathcal\{\allowbreak{}M\}\allowbreak{}\_\allowbreak{}\textbackslash{}\allowbreak{}Pi\$\allowbreak{}}{DER\_\allowbreak{}VIII-X:\allowbreak{}7}
\DERapparatusitem{ENV0201}{3;\allowbreak{}\allowbreak{} IX\allowbreak{} manuscript\allowbreak{} page\allowbreak{} 4,\allowbreak{}\allowbreak{} final\allowbreak{} displayed\allowbreak{} bifunctor}{\textbackslash{}\allowbreak{}mathcal\{\allowbreak{}M\}\allowbreak{}\_\allowbreak{}\textbackslash{}\allowbreak{}Pi\allowbreak{} \textbackslash{}\allowbreak{}times\allowbreak{} \textbackslash{}\allowbreak{}mathcal\{\allowbreak{}M\}\allowbreak{}\_\allowbreak{}\textbackslash{}\allowbreak{}Pi\allowbreak{} \textbackslash{}\allowbreak{}longrightarrow\allowbreak{} \textbackslash{}\allowbreak{}mathrm\{\allowbreak{}Ind\}\allowbreak{}}{\textbackslash{}\allowbreak{}mathcal\{\allowbreak{}M\}\allowbreak{}\textasciicircum{}\allowbreak{}o\_\allowbreak{}\textbackslash{}\allowbreak{}Pi\allowbreak{} \textbackslash{}\allowbreak{}times\allowbreak{} \textbackslash{}\allowbreak{}mathcal\{\allowbreak{}M\}\allowbreak{}\_\allowbreak{}\textbackslash{}\allowbreak{}Pi\allowbreak{} \textbackslash{}\allowbreak{}longrightarrow\allowbreak{} \textbackslash{}\allowbreak{}mathrm\{\allowbreak{}Ind\}\allowbreak{}}{DER\_\allowbreak{}VIII-X:\allowbreak{}8}
\DERapparatusitem{ENV0202}{3;\allowbreak{}\allowbreak{} running\allowbreak{} text\allowbreak{} or\allowbreak{} formula}{\textbackslash{}\allowbreak{}text\{\allowbreak{}def\}\allowbreak{}}{\textbackslash{}\allowbreak{}mathrm\{\allowbreak{}def\}\allowbreak{}}{TYPESETTING:\allowbreak{}7;\allowbreak{}\allowbreak{} NOTATION\_\allowbreak{}DER:\allowbreak{}17}
\DERapparatusitem{ENV0203}{4;\allowbreak{}\allowbreak{} IX\allowbreak{} manuscript\allowbreak{} page\allowbreak{} 6,\allowbreak{}\allowbreak{} insertion\allowbreak{} elements}{\textbackslash{}\allowbreak{}beta\textbackslash{}\allowbreak{}in\allowbreak{} P\_\allowbreak{}1,\allowbreak{}\textbackslash{}\allowbreak{}\allowbreak{} \textbackslash{}\allowbreak{}gamma\textbackslash{}\allowbreak{}in\allowbreak{} P\_\allowbreak{}2}{\textbackslash{}\allowbreak{}beta\textbackslash{}\allowbreak{}in\textbackslash{}\allowbreak{}pi\_\allowbreak{}0(\allowbreak{}P\_\allowbreak{}1)\allowbreak{},\allowbreak{}\textbackslash{}\allowbreak{}\allowbreak{} \textbackslash{}\allowbreak{}gamma\textbackslash{}\allowbreak{}in\textbackslash{}\allowbreak{}pi\_\allowbreak{}0(\allowbreak{}P\_\allowbreak{}2)\allowbreak{}}{DER\_\allowbreak{}VIII-X:\allowbreak{}9}
\DERapparatusitem{ENV0204}{5;\allowbreak{}\allowbreak{} IX\allowbreak{} manuscript\allowbreak{} page\allowbreak{} 6,\allowbreak{}\allowbreak{} component\allowbreak{} pairing}{\textbackslash{}\allowbreak{}pi\_\allowbreak{}0(\allowbreak{}u\_\allowbreak{}2)\allowbreak{}\allowbreak{} :\allowbreak{}\allowbreak{} P\_\allowbreak{}1*P\_\allowbreak{}2\allowbreak{} \textbackslash{}\allowbreak{}longrightarrow\allowbreak{} Q\_\allowbreak{}1}{\textbackslash{}\allowbreak{}pi\_\allowbreak{}0(\allowbreak{}u\_\allowbreak{}2)\allowbreak{}\allowbreak{} :\allowbreak{}\allowbreak{} \textbackslash{}\allowbreak{}pi\_\allowbreak{}0(\allowbreak{}P\_\allowbreak{}1*P\_\allowbreak{}2)\allowbreak{}\allowbreak{} \textbackslash{}\allowbreak{}longrightarrow\allowbreak{} \textbackslash{}\allowbreak{}pi\_\allowbreak{}0(\allowbreak{}Q\_\allowbreak{}1)\allowbreak{}}{DER\_\allowbreak{}VIII-X:\allowbreak{}10}
\DERapparatusitem{ENV0205}{5;\allowbreak{}\allowbreak{} IX\allowbreak{} manuscript\allowbreak{} page\allowbreak{} 7,\allowbreak{}\allowbreak{} abstract\allowbreak{} composition\allowbreak{} and\allowbreak{} associativity\allowbreak{} displays}{\textbackslash{}\allowbreak{}times}{*}{DER\_\allowbreak{}VIII-X:\allowbreak{}11}
\DERapparatusitem{ENV0206}{5;\allowbreak{}\allowbreak{} IX\allowbreak{} manuscript\allowbreak{} page\allowbreak{} 7,\allowbreak{}\allowbreak{} abstract\allowbreak{} composition\allowbreak{} and\allowbreak{} associativity\allowbreak{} displays}{\textbackslash{}\allowbreak{}times}{*}{DER\_\allowbreak{}VIII-X:\allowbreak{}11}
\DERapparatusitem{ENV0207}{5;\allowbreak{}\allowbreak{} IX\allowbreak{} manuscript\allowbreak{} page\allowbreak{} 7,\allowbreak{}\allowbreak{} abstract\allowbreak{} composition\allowbreak{} and\allowbreak{} associativity\allowbreak{} displays}{\textbackslash{}\allowbreak{}times}{*}{DER\_\allowbreak{}VIII-X:\allowbreak{}11}
\DERapparatusitem{ENV0208}{5;\allowbreak{}\allowbreak{} IX\allowbreak{} manuscript\allowbreak{} page\allowbreak{} 7,\allowbreak{}\allowbreak{} abstract\allowbreak{} composition\allowbreak{} and\allowbreak{} associativity\allowbreak{} displays}{\textbackslash{}\allowbreak{}times}{*}{DER\_\allowbreak{}VIII-X:\allowbreak{}11}
\DERapparatusitem{ENV0209}{5;\allowbreak{}\allowbreak{} IX\allowbreak{} manuscript\allowbreak{} page\allowbreak{} 7,\allowbreak{}\allowbreak{} abstract\allowbreak{} composition\allowbreak{} and\allowbreak{} associativity\allowbreak{} displays}{\textbackslash{}\allowbreak{}times}{*}{DER\_\allowbreak{}VIII-X:\allowbreak{}11}
\DERapparatusitem{ENV0210}{9;\allowbreak{}\allowbreak{} IX\allowbreak{} manuscript\allowbreak{} page\allowbreak{} 14,\allowbreak{}\allowbreak{} first\allowbreak{} comparison\allowbreak{} with\allowbreak{} Quillen}{\textbackslash{}\allowbreak{}textsc\{\allowbreak{}Quilllen\}\allowbreak{}}{\textbackslash{}\allowbreak{}textsc\{\allowbreak{}Quillen\}\allowbreak{}}{DER\_\allowbreak{}VIII-X:\allowbreak{}18}
\DERapparatusitem{ENV0211}{10;\allowbreak{}\allowbreak{} IX\allowbreak{} manuscript\allowbreak{} page\allowbreak{} 16,\allowbreak{}\allowbreak{} two\allowbreak{} arrows\allowbreak{} between\allowbreak{} zero\allowbreak{} representatives}{But\allowbreak{} one\allowbreak{} has\allowbreak{} canonical\allowbreak{} arrows}{But\allowbreak{} one\allowbreak{} has\allowbreak{} a\allowbreak{} canonical\allowbreak{} arrow}{DER\_\allowbreak{}VIII-X:\allowbreak{}21}
\DERapparatusitem{ENV0212}{10;\allowbreak{}\allowbreak{} IX\allowbreak{} manuscript\allowbreak{} page\allowbreak{} 16,\allowbreak{}\allowbreak{} two\allowbreak{} arrows\allowbreak{} between\allowbreak{} zero\allowbreak{} representatives}{e\textasciicircum{}\allowbreak{}\{\allowbreak{}-\}\allowbreak{}\_\allowbreak{}\{\allowbreak{}a,\allowbreak{}b\}\allowbreak{}\textbackslash{}\allowbreak{}circ\allowbreak{} e\textasciicircum{}\allowbreak{}\{\allowbreak{}+\}\allowbreak{}\_\allowbreak{}\{\allowbreak{}b,\allowbreak{}c\}\allowbreak{}\allowbreak{} \textbackslash{}\allowbreak{};\allowbreak{}\textbackslash{}\allowbreak{};\allowbreak{}\textbackslash{}\allowbreak{}rightleftarrows\textbackslash{}\allowbreak{};\allowbreak{}\textbackslash{}\allowbreak{};\allowbreak{}\allowbreak{} e\_\allowbreak{}\{\allowbreak{}a,\allowbreak{}c\}\allowbreak{}}{e\textasciicircum{}\allowbreak{}\{\allowbreak{}-\}\allowbreak{}\_\allowbreak{}\{\allowbreak{}a,\allowbreak{}b\}\allowbreak{}\textbackslash{}\allowbreak{}circ\allowbreak{} e\textasciicircum{}\allowbreak{}\{\allowbreak{}+\}\allowbreak{}\_\allowbreak{}\{\allowbreak{}b,\allowbreak{}c\}\allowbreak{}\allowbreak{} \textbackslash{}\allowbreak{};\allowbreak{}\textbackslash{}\allowbreak{};\allowbreak{}\textbackslash{}\allowbreak{}longrightarrow\textbackslash{}\allowbreak{};\allowbreak{}\textbackslash{}\allowbreak{};\allowbreak{}\allowbreak{} e\_\allowbreak{}\{\allowbreak{}a,\allowbreak{}c\}\allowbreak{}}{DER\_\allowbreak{}VIII-X:\allowbreak{}21}
\DERapparatusitem{ENV0213}{10;\allowbreak{}\allowbreak{} IX\allowbreak{} manuscript\allowbreak{} page\allowbreak{} 16,\allowbreak{}\allowbreak{} two\allowbreak{} arrows\allowbreak{} between\allowbreak{} zero\allowbreak{} representatives}{which\allowbreak{} are\allowbreak{} the\allowbreak{} identity\allowbreak{} on\allowbreak{} all\allowbreak{} the\allowbreak{} vertices\allowbreak{} except\allowbreak{} the\allowbreak{} one\allowbreak{} of\allowbreak{} index\allowbreak{} 0.\allowbreak{}}{which\allowbreak{} is\allowbreak{} the\allowbreak{} identity\allowbreak{} on\allowbreak{} all\allowbreak{} the\allowbreak{} vertices\allowbreak{} except\allowbreak{} those\allowbreak{} of\allowbreak{} indices\allowbreak{} 0\allowbreak{} and\allowbreak{} 1,\allowbreak{}\allowbreak{} where\allowbreak{} its\allowbreak{} components\allowbreak{} are\allowbreak{} zero.\allowbreak{}}{DER\_\allowbreak{}VIII-X:\allowbreak{}21}
\DERapparatusitem{ENV0214}{10;\allowbreak{}\allowbreak{} IX\allowbreak{} manuscript\allowbreak{} page\allowbreak{} 16,\allowbreak{}\allowbreak{} pointed\allowbreak{} whiskering\allowbreak{} construction}{defines\allowbreak{} in\allowbreak{} Hot\allowbreak{} a\allowbreak{} canonical\allowbreak{} arrow}{defines\allowbreak{} in\allowbreak{} \$\allowbreak{}\textbackslash{}\allowbreak{}Hot\_\allowbreak{}\{\allowbreak{}\textbackslash{}\allowbreak{}bullet\}\allowbreak{}\$\allowbreak{}\allowbreak{} a\allowbreak{} canonical\allowbreak{} arrow}{DER\_\allowbreak{}VIII-X:\allowbreak{}20}
\DERapparatusitem{ENV0215}{10;\allowbreak{}\allowbreak{} IX\allowbreak{} manuscript\allowbreak{} page\allowbreak{} 17,\allowbreak{}\allowbreak{} formula\allowbreak{} for\allowbreak{} f-double-prime\allowbreak{} and\allowbreak{} c-double-prime}{c''\allowbreak{} :\allowbreak{}\allowbreak{} f''(\allowbreak{}e\_\allowbreak{}\{\allowbreak{}K\_\allowbreak{}\{\allowbreak{}\textbackslash{}\allowbreak{}bullet\}\allowbreak{}\}\allowbreak{})\allowbreak{}\allowbreak{} \textbackslash{}\allowbreak{}longrightarrow\allowbreak{} e\_\allowbreak{}\{\allowbreak{}M\_\allowbreak{}\{\allowbreak{}\textbackslash{}\allowbreak{}bullet\}\allowbreak{}\}\allowbreak{})\allowbreak{},\allowbreak{}}{c''\allowbreak{} :\allowbreak{}\allowbreak{} f''(\allowbreak{}e\_\allowbreak{}\{\allowbreak{}K\_\allowbreak{}\{\allowbreak{}\textbackslash{}\allowbreak{}bullet\}\allowbreak{}\}\allowbreak{})\allowbreak{}\allowbreak{} \textbackslash{}\allowbreak{}longrightarrow\allowbreak{} e\_\allowbreak{}\{\allowbreak{}M\_\allowbreak{}\{\allowbreak{}\textbackslash{}\allowbreak{}bullet\}\allowbreak{}\}\allowbreak{},\allowbreak{}}{DER\_\allowbreak{}VIII-X:\allowbreak{}23}
\DERapparatusitem{ENV0216}{10;\allowbreak{}\allowbreak{} IX\allowbreak{} manuscript\allowbreak{} page\allowbreak{} 17,\allowbreak{}\allowbreak{} monoidal\allowbreak{} assertions\allowbreak{} with\allowbreak{} footnotes\allowbreak{} 4\allowbreak{} and\allowbreak{} 5}{with\allowbreak{} the\allowbreak{} ordinary\allowbreak{} cartesian\allowbreak{} product}{with\allowbreak{} the\allowbreak{} product\allowbreak{} \$\allowbreak{}\textbackslash{}\allowbreak{}divideontimes\$\allowbreak{}\allowbreak{} with\allowbreak{} unit\allowbreak{} \$\allowbreak{}S\textasciicircum{}\allowbreak{}0\$\allowbreak{}}{DER\_\allowbreak{}VIII-X:\allowbreak{}24}
\DERapparatusitem{ENV0217}{10;\allowbreak{}\allowbreak{} IX\allowbreak{} manuscript\allowbreak{} page\allowbreak{} 17,\allowbreak{}\allowbreak{} monoidal\allowbreak{} assertions\allowbreak{} with\allowbreak{} footnotes\allowbreak{} 4\allowbreak{} and\allowbreak{} 5}{\$\allowbreak{}\textbackslash{}\allowbreak{}Hot\_\allowbreak{}\{\allowbreak{}\textbackslash{}\allowbreak{}bullet\}\allowbreak{}\allowbreak{} \textbackslash{}\allowbreak{}longrightarrow\allowbreak{} \textbackslash{}\allowbreak{}Hot\$\allowbreak{}\allowbreak{} is\allowbreak{} compatible\allowbreak{} with\allowbreak{} this\allowbreak{} product}{\$\allowbreak{}\textbackslash{}\allowbreak{}Hot\_\allowbreak{}\{\allowbreak{}\textbackslash{}\allowbreak{}bullet\}\allowbreak{}\allowbreak{} \textbackslash{}\allowbreak{}longrightarrow\allowbreak{} \textbackslash{}\allowbreak{}Hot\$\allowbreak{}\allowbreak{} is\allowbreak{} equipped\allowbreak{} with\allowbreak{} the\allowbreak{} lax\allowbreak{} monoidal\allowbreak{} comparison\allowbreak{} from\allowbreak{} this\allowbreak{} product}{DER\_\allowbreak{}VIII-X:\allowbreak{}24}
\DERapparatusitem{ENV0218}{11;\allowbreak{}\allowbreak{} IX\allowbreak{} manuscript\allowbreak{} page\allowbreak{} 18,\allowbreak{}\allowbreak{} truncated\allowbreak{} zero\allowbreak{} path}{a\allowbreak{} \textbackslash{}\allowbreak{}longrightarrow\allowbreak{} \textbackslash{}\allowbreak{}underbrace\{\allowbreak{}0\}\allowbreak{}\_\allowbreak{}\{\allowbreak{}=\allowbreak{}a\_\allowbreak{}0\}\allowbreak{}\allowbreak{} \textbackslash{}\allowbreak{}longleftarrow\allowbreak{} b}{a\allowbreak{} \textbackslash{}\allowbreak{}longrightarrow\allowbreak{} \textbackslash{}\allowbreak{}underbrace\{\allowbreak{}0\}\allowbreak{}\_\allowbreak{}\{\allowbreak{}=\allowbreak{}a\_\allowbreak{}0\}\allowbreak{}\allowbreak{} \textbackslash{}\allowbreak{}xleftarrow\{\allowbreak{}\textbackslash{}\allowbreak{}mathrm\{\allowbreak{}id\}\allowbreak{}\}\allowbreak{}\allowbreak{} 0\allowbreak{} \textbackslash{}\allowbreak{}longrightarrow\allowbreak{} b}{DER\_\allowbreak{}VIII-X:\allowbreak{}27}
\DERapparatusitem{ENV0219}{13;\allowbreak{}\allowbreak{} IX\allowbreak{} manuscript\allowbreak{} page\allowbreak{} 21,\allowbreak{}\allowbreak{} definition\allowbreak{} of\allowbreak{} pi\_\allowbreak{}1(\allowbreak{}a,\allowbreak{}b)\allowbreak{}}{\textbackslash{}\allowbreak{}pi\_\allowbreak{}1(\allowbreak{}a,\allowbreak{}b\textbackslash{}\allowbreak{},\allowbreak{};\allowbreak{}0\_\allowbreak{}a,\allowbreak{}0\_\allowbreak{}b)\allowbreak{}}{\textbackslash{}\allowbreak{}pi\_\allowbreak{}1(\allowbreak{}a,\allowbreak{}b\textbackslash{}\allowbreak{},\allowbreak{};\allowbreak{}0\_\allowbreak{}\{\allowbreak{}a,\allowbreak{}b\}\allowbreak{},\allowbreak{}0\_\allowbreak{}\{\allowbreak{}a,\allowbreak{}b\}\allowbreak{})\allowbreak{}}{DER\_\allowbreak{}VIII-X:\allowbreak{}30}
\DERapparatusitem{ENV0220}{16;\allowbreak{}\allowbreak{} running\allowbreak{} text\allowbreak{} or\allowbreak{} formula}{\textbackslash{}\allowbreak{}text\{\allowbreak{}def\}\allowbreak{}}{\textbackslash{}\allowbreak{}mathrm\{\allowbreak{}def\}\allowbreak{}}{TYPESETTING:\allowbreak{}7;\allowbreak{}\allowbreak{} NOTATION\_\allowbreak{}DER:\allowbreak{}17}
\DERapparatusitem{ENV0221}{16;\allowbreak{}\allowbreak{} running\allowbreak{} text\allowbreak{} or\allowbreak{} formula}{\textbackslash{}\allowbreak{}text\{\allowbreak{}def\}\allowbreak{}}{\textbackslash{}\allowbreak{}mathrm\{\allowbreak{}def\}\allowbreak{}}{TYPESETTING:\allowbreak{}7;\allowbreak{}\allowbreak{} NOTATION\_\allowbreak{}DER:\allowbreak{}17}
\subsection*{Chapter XI}
\DERapparatusitem{ENV0573}{1;\allowbreak{}\allowbreak{} XI\allowbreak{} ms127,\allowbreak{}\allowbreak{} footnote\allowbreak{} 2}{R\textasciicircum{}\allowbreak{}i\allowbreak{} f\_\allowbreak{}!(\allowbreak{}L)\allowbreak{}}{L\_\allowbreak{}i\allowbreak{} f\_\allowbreak{}!(\allowbreak{}L)\allowbreak{}}{DER\_\allowbreak{}XI-XIIa:\allowbreak{}1}
\DERapparatusitem{ENV0574}{2;\allowbreak{}\allowbreak{} XI\allowbreak{} ms128,\allowbreak{}\allowbreak{} torsor\allowbreak{} functor}{P\allowbreak{} \&\allowbreak{}\allowbreak{} \textbackslash{}\allowbreak{}longmapsto\allowbreak{} \&\allowbreak{}\allowbreak{} f\textasciicircum{}\allowbreak{}*P}{P\allowbreak{} \&\allowbreak{}\allowbreak{} \textbackslash{}\allowbreak{}longmapsto\allowbreak{} \&\allowbreak{}\allowbreak{} q\textasciicircum{}\allowbreak{}*P}{DER\_\allowbreak{}XI-XIIa:\allowbreak{}3}
\DERapparatusitem{ENV0575}{2;\allowbreak{}\allowbreak{} XI\allowbreak{} ms129,\allowbreak{}\allowbreak{} bottom\allowbreak{} left\allowbreak{} of\allowbreak{} nonabelian\allowbreak{} diagram}{H\textasciicircum{}\allowbreak{}1(\allowbreak{}Y',\allowbreak{}\allowbreak{} f'\_\allowbreak{}*F)\allowbreak{}}{H\textasciicircum{}\allowbreak{}1(\allowbreak{}Y',\allowbreak{}\allowbreak{} f'\_\allowbreak{}*F')\allowbreak{}}{DER\_\allowbreak{}XI-XIIa:\allowbreak{}4}
\DERapparatusitem{ENV0576}{3;\allowbreak{}\allowbreak{} XI\allowbreak{} ms130,\allowbreak{}\allowbreak{} equation\allowbreak{} (\allowbreak{}**)\allowbreak{}\allowbreak{} second\allowbreak{} term}{H\textasciicircum{}\allowbreak{}p(\allowbreak{}Y',\allowbreak{}\allowbreak{} R\textasciicircum{}\allowbreak{}q\allowbreak{} f'\_\allowbreak{}*(\allowbreak{}F)\allowbreak{})\allowbreak{}}{H\textasciicircum{}\allowbreak{}p(\allowbreak{}Y',\allowbreak{}\allowbreak{} R\textasciicircum{}\allowbreak{}q\allowbreak{} f'\_\allowbreak{}*(\allowbreak{}F')\allowbreak{})\allowbreak{}}{DER\_\allowbreak{}XI-XIIa:\allowbreak{}6}
\DERapparatusitem{ENV0577}{4;\allowbreak{}\allowbreak{} XI\allowbreak{} ms131,\allowbreak{}\allowbreak{} Proposition\allowbreak{} 1}{necessary\allowbreak{} and\allowbreak{} sufficient\allowbreak{} that\allowbreak{} \$\allowbreak{}f\$\allowbreak{}}{necessary\allowbreak{} and\allowbreak{} sufficient\allowbreak{} that\allowbreak{} \$\allowbreak{}f\textasciicircum{}\allowbreak{}\{\allowbreak{}o\}\allowbreak{}\$\allowbreak{}}{DER\_\allowbreak{}XI-XIIa:\allowbreak{}8}
\DERapparatusitem{ENV0578}{4;\allowbreak{}\allowbreak{} XI\allowbreak{} ms131,\allowbreak{}\allowbreak{} Theorem\allowbreak{} 2}{f\_\allowbreak{}s\allowbreak{} :\allowbreak{}\allowbreak{} X\_\allowbreak{}s\allowbreak{} \textbackslash{}\allowbreak{}longrightarrow\allowbreak{} Y\_\allowbreak{}s}{f\_\allowbreak{}s\allowbreak{} :\allowbreak{}\allowbreak{} X'\_\allowbreak{}s\allowbreak{} \textbackslash{}\allowbreak{}longrightarrow\allowbreak{} X\_\allowbreak{}s}{DER\_\allowbreak{}XI-XIIa:\allowbreak{}9}
\DERapparatusitem{ENV0579}{4;\allowbreak{}\allowbreak{} XI\allowbreak{} ms132,\allowbreak{}\allowbreak{} representative-point\allowbreak{} criterion}{for\allowbreak{} \$\allowbreak{}s\allowbreak{} \textbackslash{}\allowbreak{}in\allowbreak{} S\$\allowbreak{},\allowbreak{}\allowbreak{} then\allowbreak{} \$\allowbreak{}u\$\allowbreak{}}{for\allowbreak{} \$\allowbreak{}s\allowbreak{} \textbackslash{}\allowbreak{}in\allowbreak{} S\_\allowbreak{}0\$\allowbreak{},\allowbreak{}\allowbreak{} then\allowbreak{} \$\allowbreak{}u\$\allowbreak{}}{DER\_\allowbreak{}XI-XIIa:\allowbreak{}13}
\DERapparatusitem{ENV0580}{4;\allowbreak{}\allowbreak{} XI\allowbreak{} ms132,\allowbreak{}\allowbreak{} H0\allowbreak{} direct-image\allowbreak{} comparison}{H\textasciicircum{}\allowbreak{}0(\allowbreak{}Y,\allowbreak{}\allowbreak{} \textbackslash{}\allowbreak{}underbrace\{\allowbreak{}f\_\allowbreak{}*(\allowbreak{}F)\allowbreak{}\}\allowbreak{}}{H\textasciicircum{}\allowbreak{}0(\allowbreak{}S,\allowbreak{}\allowbreak{} \textbackslash{}\allowbreak{}underbrace\{\allowbreak{}p\_\allowbreak{}*(\allowbreak{}F)\allowbreak{}\}\allowbreak{}}{DER\_\allowbreak{}XI-XIIa:\allowbreak{}10}
\DERapparatusitem{ENV0581}{4;\allowbreak{}\allowbreak{} XI\allowbreak{} ms132,\allowbreak{}\allowbreak{} H0\allowbreak{} direct-image\allowbreak{} comparison}{H\textasciicircum{}\allowbreak{}0(\allowbreak{}Y,\allowbreak{}\allowbreak{} f'\_\allowbreak{}*(\allowbreak{}F')\allowbreak{})\allowbreak{}}{H\textasciicircum{}\allowbreak{}0(\allowbreak{}S,\allowbreak{}\allowbreak{} p'\_\allowbreak{}*(\allowbreak{}F')\allowbreak{})\allowbreak{}}{DER\_\allowbreak{}XI-XIIa:\allowbreak{}10}
\DERapparatusitem{ENV0582}{4;\allowbreak{}\allowbreak{} XI\allowbreak{} ms132,\allowbreak{}\allowbreak{} torsor\allowbreak{} comparison\allowbreak{} target}{(\allowbreak{}F\textbackslash{}\allowbreak{}text\{\allowbreak{}-torsors\allowbreak{} on\allowbreak{} \}\allowbreak{}\allowbreak{} X')\allowbreak{}}{(\allowbreak{}F'\textbackslash{}\allowbreak{}text\{\allowbreak{}-torsors\allowbreak{} on\allowbreak{} \}\allowbreak{}\allowbreak{} X')\allowbreak{}}{DER\_\allowbreak{}XI-XIIa:\allowbreak{}14}
\DERapparatusitem{ENV0583}{6;\allowbreak{}\allowbreak{} XI\allowbreak{} ms135,\allowbreak{}\allowbreak{} nonempty\allowbreak{} base}{Since\allowbreak{} \$\allowbreak{}S\allowbreak{} =\allowbreak{}\allowbreak{} \textbackslash{}\allowbreak{}emptyset\$\allowbreak{}}{Since\allowbreak{} \$\allowbreak{}S\allowbreak{} \textbackslash{}\allowbreak{}neq\allowbreak{} \textbackslash{}\allowbreak{}emptyset\$\allowbreak{}}{DER\_\allowbreak{}XI-XIIa:\allowbreak{}16}
\subsection*{Chapter XII, 1–4}
\DERapparatusitem{ENV0636}{1;\allowbreak{}\allowbreak{} XII\allowbreak{} ms1,\allowbreak{}\allowbreak{} first\allowbreak{} corollary}{y\textbackslash{}\allowbreak{}backslash\allowbreak{} X\allowbreak{} \textbackslash{}\allowbreak{}fl\allowbreak{} y\textbackslash{}\allowbreak{}backslash\allowbreak{} Y}{X/\allowbreak{}y\allowbreak{} \textbackslash{}\allowbreak{}fl\allowbreak{} Y/\allowbreak{}y}{DER\_\allowbreak{}XI-XIIa:\allowbreak{}18}
\DERapparatusitem{ENV0637}{3;\allowbreak{}\allowbreak{} XII\allowbreak{} ms3,\allowbreak{}\allowbreak{} Proposition2}{\textasciicircum{}\allowbreak{}\textbackslash{}\allowbreak{}circ}{}{DER\_\allowbreak{}XI-XIIa:\allowbreak{}19}
\DERapparatusitem{ENV0638}{4;\allowbreak{}\allowbreak{} XII\allowbreak{} ms4,\allowbreak{}\allowbreak{} proof\allowbreak{} of\allowbreak{} duality\allowbreak{} corollary}{For\allowbreak{} \$\allowbreak{}2'\$\allowbreak{}}{For\allowbreak{} \$\allowbreak{}3'\$\allowbreak{}}{DER\_\allowbreak{}XI-XIIa:\allowbreak{}20}
\DERapparatusitem{ENV0639}{4;\allowbreak{}\allowbreak{} XII\allowbreak{} ms4,\allowbreak{}\allowbreak{} proof\allowbreak{} of\allowbreak{} duality\allowbreak{} corollary}{For\allowbreak{} \$\allowbreak{}3'\$\allowbreak{}}{For\allowbreak{} \$\allowbreak{}2'\$\allowbreak{}}{DER\_\allowbreak{}XI-XIIa:\allowbreak{}20}
\DERapparatusitem{ENV0640}{5;\allowbreak{}\allowbreak{} XII\allowbreak{} ms5,\allowbreak{}\allowbreak{} recalled\allowbreak{} fibre\allowbreak{} criterion}{'}{\_\allowbreak{}1}{DER\_\allowbreak{}XI-XIIa:\allowbreak{}21}
\DERapparatusitem{ENV0641}{6;\allowbreak{}\allowbreak{} XII\allowbreak{} ms6,\allowbreak{}\allowbreak{} question\allowbreak{} after\allowbreak{} corollary}{\_\allowbreak{}S}{}{DER\_\allowbreak{}XI-XIIa:\allowbreak{}22}
\DERapparatusitem{ENV0642}{6;\allowbreak{}\allowbreak{} XII\allowbreak{} ms6-7,\allowbreak{}\allowbreak{} component\allowbreak{} of\allowbreak{} the\allowbreak{} base}{the\allowbreak{} connected\allowbreak{} component\allowbreak{} \$\allowbreak{}Y\_\allowbreak{}0\$\allowbreak{}\allowbreak{} of\allowbreak{} \$\allowbreak{}Y\$\allowbreak{}}{the\allowbreak{} connected\allowbreak{} component\allowbreak{} \$\allowbreak{}S\_\allowbreak{}0\$\allowbreak{}\allowbreak{} of\allowbreak{} \$\allowbreak{}S\$\allowbreak{}}{DER\_\allowbreak{}XI-XIIa:\allowbreak{}23}
\DERapparatusitem{ENV0643}{7;\allowbreak{}\allowbreak{} XII\allowbreak{} ms7,\allowbreak{}\allowbreak{} path-cover\allowbreak{} paragraph}{-apheric}{-aspheric}{DER\_\allowbreak{}XI-XIIa:\allowbreak{}24}
\DERapparatusitem{ENV0644}{7;\allowbreak{}\allowbreak{} XII\allowbreak{} ms8,\allowbreak{}\allowbreak{} constant-path\allowbreak{} comparison}{The\allowbreak{} fact\allowbreak{} that\allowbreak{} \$\allowbreak{}i\_\allowbreak{}S\allowbreak{} \textbackslash{}\allowbreak{}in\allowbreak{} W\$\allowbreak{}}{The\allowbreak{} fact\allowbreak{} that\allowbreak{} \$\allowbreak{}i\_\allowbreak{}S\allowbreak{} \textbackslash{}\allowbreak{}in\allowbreak{} W\_\allowbreak{}\textbackslash{}\allowbreak{}infty\$\allowbreak{}}{DER\_\allowbreak{}XI-XIIa:\allowbreak{}25}
\DERapparatusitem{ENV0645}{8;\allowbreak{}\allowbreak{} XII\allowbreak{} ms9,\allowbreak{}\allowbreak{} lifting\allowbreak{} data}{an\allowbreak{} arrow\allowbreak{} \$\allowbreak{}u\allowbreak{} \textbackslash{}\allowbreak{}in\allowbreak{} \textbackslash{}\allowbreak{}mathrm\{\allowbreak{}Fl\}\allowbreak{}(\allowbreak{}S\allowbreak{} \textbackslash{}\allowbreak{}times\allowbreak{} S)\allowbreak{}\$\allowbreak{}}{an\allowbreak{} arrow\allowbreak{} \$\allowbreak{}(\allowbreak{}u,\allowbreak{}v)\allowbreak{}\allowbreak{} \textbackslash{}\allowbreak{}in\allowbreak{} \textbackslash{}\allowbreak{}mathrm\{\allowbreak{}Fl\}\allowbreak{}(\allowbreak{}S\allowbreak{} \textbackslash{}\allowbreak{}times\allowbreak{} S)\allowbreak{}\$\allowbreak{}}{DER\_\allowbreak{}XI-XIIa:\allowbreak{}26}
\DERapparatusitem{ENV0646}{8;\allowbreak{}\allowbreak{} XII\allowbreak{} ms9,\allowbreak{}\allowbreak{} lifting-category\allowbreak{} variance}{liftings\allowbreak{} of\allowbreak{} \$\allowbreak{}u\$\allowbreak{}\allowbreak{} with\allowbreak{} source\allowbreak{} \$\allowbreak{}c\$\allowbreak{}}{liftings\allowbreak{} of\allowbreak{} \$\allowbreak{}(\allowbreak{}u,\allowbreak{}v)\allowbreak{}\$\allowbreak{}\allowbreak{} with\allowbreak{} target\allowbreak{} \$\allowbreak{}c\$\allowbreak{}}{DER\_\allowbreak{}XI-XIIa:\allowbreak{}27}
\DERapparatusitem{ENV0647}{10;\allowbreak{}\allowbreak{} running\allowbreak{} text\allowbreak{} or\allowbreak{} formula}{\textbackslash{}\allowbreak{}text\{\allowbreak{}def\}\allowbreak{}}{\textbackslash{}\allowbreak{}mathrm\{\allowbreak{}def\}\allowbreak{}}{TYPESETTING:\allowbreak{}7;\allowbreak{}\allowbreak{} NOTATION\_\allowbreak{}DER:\allowbreak{}17}
\DERapparatusitem{ENV0648}{11;\allowbreak{}\allowbreak{} XII\allowbreak{} ms13,\allowbreak{}\allowbreak{} W(\allowbreak{}4b')\allowbreak{}}{Let\allowbreak{} \$\allowbreak{}f\allowbreak{} :\allowbreak{}\allowbreak{} X\allowbreak{} \textbackslash{}\allowbreak{}longrightarrow\allowbreak{} Y\$\allowbreak{}\allowbreak{} be\allowbreak{} an\allowbreak{} arrow\allowbreak{} of\allowbreak{} Cat\$\allowbreak{}/\allowbreak{}S\$\allowbreak{}\allowbreak{} which\allowbreak{} is\allowbreak{} \$\allowbreak{}W\$\allowbreak{}-perfect\allowbreak{} and\allowbreak{} \$\allowbreak{}W\_\allowbreak{}\textbackslash{}\allowbreak{}infty\$\allowbreak{}-perfect,\allowbreak{}\allowbreak{} then\allowbreak{} \$\allowbreak{}f\allowbreak{} \textbackslash{}\allowbreak{}in\allowbreak{} W\$\allowbreak{}\allowbreak{} implies\allowbreak{} that\allowbreak{} the\allowbreak{} fibers\allowbreak{} of\allowbreak{} \$\allowbreak{}f\$\allowbreak{}\allowbreak{} are\allowbreak{} \$\allowbreak{}W\$\allowbreak{}-aspheric.\allowbreak{}\allowbreak{} (\allowbreak{}It\allowbreak{} is\allowbreak{} better\allowbreak{} to\allowbreak{} say:\allowbreak{}\allowbreak{} If\allowbreak{} \$\allowbreak{}f\$\allowbreak{}\allowbreak{} is\allowbreak{} \$\allowbreak{}W\_\allowbreak{}\textbackslash{}\allowbreak{}omega\$\allowbreak{}-perfect,\allowbreak{}\allowbreak{} it\allowbreak{} is\allowbreak{} an\allowbreak{} ultrastrong\allowbreak{} \$\allowbreak{}W\$\allowbreak{}-fibration.\allowbreak{})\allowbreak{}}{Let\allowbreak{} \$\allowbreak{}f\allowbreak{} :\allowbreak{}\allowbreak{} X\allowbreak{} \textbackslash{}\allowbreak{}longrightarrow\allowbreak{} Y\$\allowbreak{}\allowbreak{} be\allowbreak{} an\allowbreak{} arrow\allowbreak{} in\allowbreak{} Cat\$\allowbreak{}/\allowbreak{}S\$\allowbreak{}\allowbreak{} such\allowbreak{} that\allowbreak{} \$\allowbreak{}X\allowbreak{} \textbackslash{}\allowbreak{}longrightarrow\allowbreak{} S\$\allowbreak{}\allowbreak{} and\allowbreak{} \$\allowbreak{}Y\allowbreak{} \textbackslash{}\allowbreak{}longrightarrow\allowbreak{} S\$\allowbreak{}\allowbreak{} are\allowbreak{} \$\allowbreak{}W\_\allowbreak{}\textbackslash{}\allowbreak{}omega\$\allowbreak{}-perfect.\allowbreak{}\allowbreak{} If\allowbreak{} \$\allowbreak{}f\allowbreak{} \textbackslash{}\allowbreak{}in\allowbreak{} W\$\allowbreak{},\allowbreak{}\allowbreak{} then\allowbreak{} \$\allowbreak{}\textbackslash{}\allowbreak{}forall\textbackslash{}\allowbreak{},\allowbreak{}s\allowbreak{} \textbackslash{}\allowbreak{}in\allowbreak{} S\$\allowbreak{},\allowbreak{}\allowbreak{} \$\allowbreak{}f\_\allowbreak{}s\allowbreak{} :\allowbreak{}\allowbreak{} X\_\allowbreak{}s\allowbreak{} \textbackslash{}\allowbreak{}longrightarrow\allowbreak{} Y\_\allowbreak{}s\$\allowbreak{}\allowbreak{} is\allowbreak{} in\allowbreak{} \$\allowbreak{}W\$\allowbreak{}.\allowbreak{}}{DER\_\allowbreak{}XI-XIIa:\allowbreak{}32}
\DERapparatusitem{ENV0649}{11;\allowbreak{}\allowbreak{} XII\allowbreak{} ms14,\allowbreak{}\allowbreak{} conclusion\allowbreak{} of\allowbreak{} W(\allowbreak{}4b'')\allowbreak{}}{the\allowbreak{} fibers\allowbreak{} of\allowbreak{} \$\allowbreak{}f\$\allowbreak{}\allowbreak{} are\allowbreak{} \$\allowbreak{}W\_\allowbreak{}\textbackslash{}\allowbreak{}omega\$\allowbreak{}-aspheric}{the\allowbreak{} fibers\allowbreak{} of\allowbreak{} \$\allowbreak{}f\$\allowbreak{}\allowbreak{} are\allowbreak{} \$\allowbreak{}W\$\allowbreak{}-aspheric}{DER\_\allowbreak{}XI-XIIa:\allowbreak{}33}
\DERapparatusitem{ENV0650}{13;\allowbreak{}\allowbreak{} XII\allowbreak{} ms15,\allowbreak{}\allowbreak{} displayed\allowbreak{} factorization\allowbreak{} paragraph}{\textbackslash{}\allowbreak{}bar}{}{DER\_\allowbreak{}XI-XIIa:\allowbreak{}34}
\DERapparatusitem{ENV0651}{13;\allowbreak{}\allowbreak{} XII\allowbreak{} ms16,\allowbreak{}\allowbreak{} strong\allowbreak{} saturation\allowbreak{} and\allowbreak{} localization\allowbreak{} paragraph}{W\_\allowbreak{}i}{W\_\allowbreak{}n}{DER\_\allowbreak{}XI-XIIa:\allowbreak{}36}
\DERapparatusitem{ENV0652}{13;\allowbreak{}\allowbreak{} XII\allowbreak{} ms16,\allowbreak{}\allowbreak{} strong\allowbreak{} saturation\allowbreak{} and\allowbreak{} localization\allowbreak{} paragraph}{W\_\allowbreak{}i}{W\_\allowbreak{}n}{DER\_\allowbreak{}XI-XIIa:\allowbreak{}36}
\DERapparatusitem{ENV0653}{13;\allowbreak{}\allowbreak{} XII\allowbreak{} ms16,\allowbreak{}\allowbreak{} strong\allowbreak{} saturation\allowbreak{} and\allowbreak{} localization\allowbreak{} paragraph}{W\_\allowbreak{}i}{W\_\allowbreak{}n}{DER\_\allowbreak{}XI-XIIa:\allowbreak{}36}
\DERapparatusitem{ENV0654}{13;\allowbreak{}\allowbreak{} XII\allowbreak{} ms16,\allowbreak{}\allowbreak{} strong\allowbreak{} saturation\allowbreak{} and\allowbreak{} localization\allowbreak{} paragraph}{W\_\allowbreak{}i}{W\_\allowbreak{}n}{DER\_\allowbreak{}XI-XIIa:\allowbreak{}36}
\DERapparatusitem{ENV0655}{13;\allowbreak{}\allowbreak{} XII\allowbreak{} ms16,\allowbreak{}\allowbreak{} strong\allowbreak{} saturation\allowbreak{} and\allowbreak{} localization\allowbreak{} paragraph}{W\_\allowbreak{}i}{W\_\allowbreak{}n}{DER\_\allowbreak{}XI-XIIa:\allowbreak{}36}
\DERapparatusitem{ENV0656}{14;\allowbreak{}\allowbreak{} XII\allowbreak{} ms17,\allowbreak{}\allowbreak{} cokernel\allowbreak{} display}{\textbackslash{}\allowbreak{}pi\_\allowbreak{}n(\allowbreak{}X\_\allowbreak{}s)\allowbreak{}\allowbreak{} =\allowbreak{}\allowbreak{} \textbackslash{}\allowbreak{}mathrm\{\allowbreak{}Coker\}\allowbreak{}\textbackslash{}\allowbreak{}left(\allowbreak{}\textbackslash{}\allowbreak{}pi\_\allowbreak{}\{\allowbreak{}n+1\}\allowbreak{}(\allowbreak{}X)\allowbreak{}\allowbreak{} \textbackslash{}\allowbreak{}longrightarrow\allowbreak{} \textbackslash{}\allowbreak{}pi\_\allowbreak{}\{\allowbreak{}n+1\}\allowbreak{}(\allowbreak{}Y)\allowbreak{}\textbackslash{}\allowbreak{}right)\allowbreak{}}{\textbackslash{}\allowbreak{}pi\_\allowbreak{}n(\allowbreak{}X\_\allowbreak{}s)\allowbreak{}\allowbreak{} =\allowbreak{}\allowbreak{} \textbackslash{}\allowbreak{}begin\{\allowbreak{}cases\}\allowbreak{}\textbackslash{}\allowbreak{}mathrm\{\allowbreak{}Coker\}\allowbreak{}\textbackslash{}\allowbreak{}left(\allowbreak{}\textbackslash{}\allowbreak{}pi\_\allowbreak{}\{\allowbreak{}n+1\}\allowbreak{}(\allowbreak{}X)\allowbreak{}\allowbreak{} \textbackslash{}\allowbreak{}longrightarrow\allowbreak{} \textbackslash{}\allowbreak{}pi\_\allowbreak{}\{\allowbreak{}n+1\}\allowbreak{}(\allowbreak{}Y)\allowbreak{}\textbackslash{}\allowbreak{}right)\allowbreak{},\allowbreak{}\&\allowbreak{}n\textbackslash{}\allowbreak{}geq1,\allowbreak{}\textbackslash{}\allowbreak{}\textbackslash{}\allowbreak{}\allowbreak{} \textbackslash{}\allowbreak{}pi\_\allowbreak{}1(\allowbreak{}Y)\allowbreak{}/\allowbreak{}\textbackslash{}\allowbreak{}operatorname\{\allowbreak{}im\}\allowbreak{}\textbackslash{}\allowbreak{}!\textbackslash{}\allowbreak{}left(\allowbreak{}\textbackslash{}\allowbreak{}pi\_\allowbreak{}1(\allowbreak{}X)\allowbreak{}\textbackslash{}\allowbreak{}longrightarrow\textbackslash{}\allowbreak{}pi\_\allowbreak{}1(\allowbreak{}Y)\allowbreak{}\textbackslash{}\allowbreak{}right)\allowbreak{},\allowbreak{}\&\allowbreak{}n=\allowbreak{}0\textbackslash{}\allowbreak{}quad\textbackslash{}\allowbreak{}text\{\allowbreak{}(\allowbreak{}pointed\allowbreak{} set\allowbreak{} of\allowbreak{} cosets)\allowbreak{}\}\allowbreak{}.\allowbreak{}\textbackslash{}\allowbreak{}end\{\allowbreak{}cases\}\allowbreak{}}{DER\_\allowbreak{}XI-XIIa:\allowbreak{}37}
\DERapparatusitem{ENV0657}{22;\allowbreak{}\allowbreak{} XII:\allowbreak{}26,\allowbreak{}\allowbreak{} Proposition\allowbreak{} 3(\allowbreak{}d)\allowbreak{}}{which\allowbreak{} lift\allowbreak{} \$\allowbreak{}u\$\allowbreak{}\}\allowbreak{}}{which\allowbreak{} lift\allowbreak{} \$\allowbreak{}u\_\allowbreak{}0\$\allowbreak{}\}\allowbreak{}}{DER\_\allowbreak{}XI-XIIa:\allowbreak{}43}
\DERapparatusitem{ENV0658}{22;\allowbreak{}\allowbreak{} XII:\allowbreak{}27,\allowbreak{}\allowbreak{} after\allowbreak{} the\allowbreak{} lift-category\allowbreak{} identifications}{Hence\allowbreak{} the\allowbreak{} equivalence\allowbreak{} of\allowbreak{} a)\allowbreak{}\allowbreak{} b)\allowbreak{}\allowbreak{} c)\allowbreak{}.\allowbreak{}}{Hence\allowbreak{} the\allowbreak{} equivalence\allowbreak{} of\allowbreak{} a)\allowbreak{}\allowbreak{} b)\allowbreak{}\allowbreak{} d)\allowbreak{}.\allowbreak{}}{DER\_\allowbreak{}XI-XIIa:\allowbreak{}44}
\DERapparatusitem{ENV0659}{23;\allowbreak{}\allowbreak{} XII:\allowbreak{}28,\allowbreak{}\allowbreak{} Definition\allowbreak{} 2,\allowbreak{}\allowbreak{} dual\allowbreak{} inclusion}{the\allowbreak{} inclusion\allowbreak{} morphisms\allowbreak{} \$\allowbreak{}(\allowbreak{}*)\allowbreak{}\$\allowbreak{}}{the\allowbreak{} inclusion\allowbreak{} morphisms\allowbreak{} \$\allowbreak{}X\_\allowbreak{}y\allowbreak{} \textbackslash{}\allowbreak{}longrightarrow\allowbreak{} y\textbackslash{}\allowbreak{}backslash\allowbreak{} X\$\allowbreak{}}{DER\_\allowbreak{}XI-XIIa:\allowbreak{}47}
\DERapparatusitem{ENV0660}{24;\allowbreak{}\allowbreak{} XII:\allowbreak{}31,\allowbreak{}\allowbreak{} corollary\allowbreak{} proof\allowbreak{} opening}{Suppose\allowbreak{} e.\allowbreak{}g.\allowbreak{}\allowbreak{} \$\allowbreak{}f\$\allowbreak{}\allowbreak{} proper}{Suppose\allowbreak{} e.\allowbreak{}g.\allowbreak{}\allowbreak{} \$\allowbreak{}X\$\allowbreak{}\allowbreak{} and\allowbreak{} \$\allowbreak{}Y\$\allowbreak{}\allowbreak{} proper\allowbreak{} over\allowbreak{} \$\allowbreak{}S\$\allowbreak{}}{DER\_\allowbreak{}XI-XIIa:\allowbreak{}49}
\DERapparatusitem{ENV0661}{25;\allowbreak{}\allowbreak{} XII:\allowbreak{}31,\allowbreak{}\allowbreak{} reference\allowbreak{} to\allowbreak{} coslice\allowbreak{} lemma}{lemma\allowbreak{} p.\allowbreak{}\textasciitilde{}28}{lemma\allowbreak{} p.\allowbreak{}\textasciitilde{}29}{DER\_\allowbreak{}XI-XIIa:\allowbreak{}50}
\DERapparatusitem{ENV0662}{26;\allowbreak{}\allowbreak{} running\allowbreak{} text\allowbreak{} or\allowbreak{} formula}{\textbackslash{}\allowbreak{}text\{\allowbreak{}def\}\allowbreak{}}{\textbackslash{}\allowbreak{}mathrm\{\allowbreak{}def\}\allowbreak{}}{TYPESETTING:\allowbreak{}7;\allowbreak{}\allowbreak{} NOTATION\_\allowbreak{}DER:\allowbreak{}17}
\DERapparatusitem{ENV0663}{27;\allowbreak{}\allowbreak{} XII:\allowbreak{}34,\allowbreak{}\allowbreak{} transport\allowbreak{} base-change\allowbreak{} triangle\allowbreak{} labels}{\textbackslash{}\allowbreak{}hot\_\allowbreak{}W\textasciicircum{}\allowbreak{}\{\allowbreak{}X'/\allowbreak{}S'\}\allowbreak{}}{\textbackslash{}\allowbreak{}hot\_\allowbreak{}*\textasciicircum{}\allowbreak{}\{\allowbreak{}X'/\allowbreak{}S',\allowbreak{}W\}\allowbreak{}}{DER\_\allowbreak{}XI-XIIa:\allowbreak{}51}
\DERapparatusitem{ENV0664}{27;\allowbreak{}\allowbreak{} XII:\allowbreak{}34,\allowbreak{}\allowbreak{} transport\allowbreak{} base-change\allowbreak{} triangle\allowbreak{} labels}{\textbackslash{}\allowbreak{}hot\_\allowbreak{}W\textasciicircum{}\allowbreak{}\{\allowbreak{}X/\allowbreak{}S\}\allowbreak{}}{\textbackslash{}\allowbreak{}hot\_\allowbreak{}*\textasciicircum{}\allowbreak{}\{\allowbreak{}X/\allowbreak{}S,\allowbreak{}W\}\allowbreak{}}{DER\_\allowbreak{}XI-XIIa:\allowbreak{}51}
\DERapparatusitem{ENV0665}{27;\allowbreak{}\allowbreak{} XII:\allowbreak{}35,\allowbreak{}\allowbreak{} final\allowbreak{} restatement\allowbreak{} of\allowbreak{} Proposition\allowbreak{} 7}{c\allowbreak{} irc}{\textbackslash{}\allowbreak{}circ}{DER\_\allowbreak{}XI-XIIa:\allowbreak{}53}
\DERapparatusitem{ENV0666}{28;\allowbreak{}\allowbreak{} XII:\allowbreak{}37,\allowbreak{}\allowbreak{} name\allowbreak{} of\allowbreak{} the\allowbreak{} dual\allowbreak{} transport}{dually\allowbreak{} for\allowbreak{} \$\allowbreak{}\textbackslash{}\allowbreak{}lambda'\$\allowbreak{}}{dually\allowbreak{} for\allowbreak{} \$\allowbreak{}\textbackslash{}\allowbreak{}lambda\textasciicircum{}\allowbreak{}\textbackslash{}\allowbreak{}circ\$\allowbreak{}}{DER\_\allowbreak{}XI-XIIa:\allowbreak{}55}
\DERapparatusitem{ENV0667}{31;\allowbreak{}\allowbreak{} XII\allowbreak{} ms42,\allowbreak{}\allowbreak{} final\allowbreak{} arrow\allowbreak{} of\allowbreak{} Proposition\allowbreak{} 8\allowbreak{} proof}{\textbackslash{}\allowbreak{}text\{\allowbreak{}\allowbreak{} by\allowbreak{} c)\allowbreak{}\allowbreak{} \}\allowbreak{}}{\textbackslash{}\allowbreak{}text\{\allowbreak{}\allowbreak{} by\allowbreak{} b)\allowbreak{}\allowbreak{} \}\allowbreak{}}{DER\_\allowbreak{}XI-XIIa:\allowbreak{}58}
\DERapparatusitem{ENV0668}{32;\allowbreak{}\allowbreak{} XII\allowbreak{} ms43,\allowbreak{}\allowbreak{} opening\allowbreak{} hypothesis}{that\allowbreak{} \$\allowbreak{}f\$\allowbreak{}\allowbreak{} is\allowbreak{} \$\allowbreak{}W\$\allowbreak{}-smooth\allowbreak{} and\allowbreak{} a\allowbreak{} left\allowbreak{} \$\allowbreak{}W\$\allowbreak{}-fibration}{that\allowbreak{} \$\allowbreak{}f\$\allowbreak{}\allowbreak{} is\allowbreak{} \$\allowbreak{}W\$\allowbreak{}-presmooth\allowbreak{} and\allowbreak{} a\allowbreak{} \$\allowbreak{}W\$\allowbreak{}-pseudofibration}{DER\_\allowbreak{}XI-XIIa:\allowbreak{}60}
\DERapparatusitem{ENV0669}{33;\allowbreak{}\allowbreak{} running\allowbreak{} text\allowbreak{} or\allowbreak{} formula}{\textbackslash{}\allowbreak{}mathscr\{\allowbreak{}X\}\allowbreak{}}{\textbackslash{}\allowbreak{}mathcal\{\allowbreak{}X\}\allowbreak{}}{TYPESETTING:\allowbreak{}1}
\DERapparatusitem{ENV0670}{33;\allowbreak{}\allowbreak{} running\allowbreak{} text\allowbreak{} or\allowbreak{} formula}{\textbackslash{}\allowbreak{}mathscr\{\allowbreak{}X\}\allowbreak{}}{\textbackslash{}\allowbreak{}mathcal\{\allowbreak{}X\}\allowbreak{}}{TYPESETTING:\allowbreak{}1}
\DERapparatusitem{ENV0671}{33;\allowbreak{}\allowbreak{} running\allowbreak{} text\allowbreak{} or\allowbreak{} formula}{\textbackslash{}\allowbreak{}mathscr\{\allowbreak{}X\}\allowbreak{}}{\textbackslash{}\allowbreak{}mathcal\{\allowbreak{}X\}\allowbreak{}}{TYPESETTING:\allowbreak{}1}
\DERapparatusitem{ENV0672}{33;\allowbreak{}\allowbreak{} XII\allowbreak{} ms44,\allowbreak{}\allowbreak{} first\allowbreak{} displayed\allowbreak{} natural\allowbreak{} transformation}{f\textbackslash{}\allowbreak{}fl\textbackslash{}\allowbreak{}mathscr\{\allowbreak{}X\}\allowbreak{}\textbackslash{}\allowbreak{}circ\allowbreak{} g}{g\textasciicircum{}\allowbreak{}\{\allowbreak{}\textbackslash{}\allowbreak{}mathrm\{\allowbreak{}op\}\allowbreak{}\}\allowbreak{}\textbackslash{}\allowbreak{}longrightarrow\textbackslash{}\allowbreak{}mathcal\{\allowbreak{}X\}\allowbreak{}\textbackslash{}\allowbreak{}circ\allowbreak{} f\textasciicircum{}\allowbreak{}\{\allowbreak{}\textbackslash{}\allowbreak{}mathrm\{\allowbreak{}op\}\allowbreak{}\}\allowbreak{}}{DER\_\allowbreak{}XI-XIIa:\allowbreak{}61;\allowbreak{}\allowbreak{} TYPESETTING:\allowbreak{}1}
\DERapparatusitem{ENV0673}{33;\allowbreak{}\allowbreak{} running\allowbreak{} text\allowbreak{} or\allowbreak{} formula}{\textbackslash{}\allowbreak{}mathscr\{\allowbreak{}X\}\allowbreak{}}{\textbackslash{}\allowbreak{}mathcal\{\allowbreak{}X\}\allowbreak{}}{TYPESETTING:\allowbreak{}1}
\DERapparatusitem{ENV0674}{33;\allowbreak{}\allowbreak{} running\allowbreak{} text\allowbreak{} or\allowbreak{} formula}{\textbackslash{}\allowbreak{}mathscr\{\allowbreak{}X\}\allowbreak{}}{\textbackslash{}\allowbreak{}mathcal\{\allowbreak{}X\}\allowbreak{}}{TYPESETTING:\allowbreak{}1}
\DERapparatusitem{ENV0675}{33;\allowbreak{}\allowbreak{} running\allowbreak{} text\allowbreak{} or\allowbreak{} formula}{\textbackslash{}\allowbreak{}mathscr\{\allowbreak{}X\}\allowbreak{}}{\textbackslash{}\allowbreak{}mathcal\{\allowbreak{}X\}\allowbreak{}}{TYPESETTING:\allowbreak{}1}
\DERapparatusitem{ENV0676}{33;\allowbreak{}\allowbreak{} running\allowbreak{} text\allowbreak{} or\allowbreak{} formula}{\textbackslash{}\allowbreak{}mathscr\{\allowbreak{}X\}\allowbreak{}}{\textbackslash{}\allowbreak{}mathcal\{\allowbreak{}X\}\allowbreak{}}{TYPESETTING:\allowbreak{}1}
\DERapparatusitem{ENV0677}{33;\allowbreak{}\allowbreak{} running\allowbreak{} text\allowbreak{} or\allowbreak{} formula}{\textbackslash{}\allowbreak{}mathscr\{\allowbreak{}X\}\allowbreak{}}{\textbackslash{}\allowbreak{}mathcal\{\allowbreak{}X\}\allowbreak{}}{TYPESETTING:\allowbreak{}1}
\DERapparatusitem{ENV0678}{33;\allowbreak{}\allowbreak{} running\allowbreak{} text\allowbreak{} or\allowbreak{} formula}{\textbackslash{}\allowbreak{}mathscr\{\allowbreak{}X\}\allowbreak{}}{\textbackslash{}\allowbreak{}mathcal\{\allowbreak{}X\}\allowbreak{}}{TYPESETTING:\allowbreak{}1}
\DERapparatusitem{ENV0679}{33;\allowbreak{}\allowbreak{} running\allowbreak{} text\allowbreak{} or\allowbreak{} formula}{\textbackslash{}\allowbreak{}mathscr\{\allowbreak{}X\}\allowbreak{}}{\textbackslash{}\allowbreak{}mathcal\{\allowbreak{}X\}\allowbreak{}}{TYPESETTING:\allowbreak{}1}
\DERapparatusitem{ENV0680}{33;\allowbreak{}\allowbreak{} running\allowbreak{} text\allowbreak{} or\allowbreak{} formula}{\textbackslash{}\allowbreak{}mathscr\{\allowbreak{}X\}\allowbreak{}}{\textbackslash{}\allowbreak{}mathcal\{\allowbreak{}X\}\allowbreak{}}{TYPESETTING:\allowbreak{}1}
\DERapparatusitem{ENV0681}{33;\allowbreak{}\allowbreak{} XII\allowbreak{} ms45,\allowbreak{}\allowbreak{} sentence\allowbreak{} after\allowbreak{} the\allowbreak{} component\allowbreak{} formula}{\$\allowbreak{}\textbackslash{}\allowbreak{}mathscr\{\allowbreak{}X\}\allowbreak{}\$\allowbreak{}\allowbreak{} is\allowbreak{} in\allowbreak{} \$\allowbreak{}W\$\allowbreak{}\allowbreak{} fiber\allowbreak{} by\allowbreak{} fiber}{\$\allowbreak{}\textbackslash{}\allowbreak{}mathcal\{\allowbreak{}X\}\allowbreak{}*u\$\allowbreak{}\allowbreak{} is\allowbreak{} in\allowbreak{} \$\allowbreak{}W\$\allowbreak{}\allowbreak{} fiber\allowbreak{} by\allowbreak{} fiber}{DER\_\allowbreak{}XI-XIIa:\allowbreak{}62;\allowbreak{}\allowbreak{} TYPESETTING:\allowbreak{}1}
\DERapparatusitem{ENV0682}{33;\allowbreak{}\allowbreak{} running\allowbreak{} text\allowbreak{} or\allowbreak{} formula}{\textbackslash{}\allowbreak{}text\{\allowbreak{}def\}\allowbreak{}}{\textbackslash{}\allowbreak{}mathrm\{\allowbreak{}def\}\allowbreak{}}{TYPESETTING:\allowbreak{}7;\allowbreak{}\allowbreak{} NOTATION\_\allowbreak{}DER:\allowbreak{}17}
\DERapparatusitem{ENV0683}{37;\allowbreak{}\allowbreak{} XII\allowbreak{} ms51,\allowbreak{}\allowbreak{} Theorem\allowbreak{} 3(\allowbreak{}b1)\allowbreak{},\allowbreak{}\allowbreak{} parenthetical\allowbreak{} base-map\allowbreak{} explanation}{a\$\allowbreak{}'\$\allowbreak{}\allowbreak{} \textbackslash{}\allowbreak{}textup\{\allowbreak{}\textbackslash{}\allowbreak{}texttt\{\allowbreak{}[\allowbreak{}rather\allowbreak{} a]\allowbreak{}\}\allowbreak{}\}\allowbreak{}}{a\allowbreak{} \textbackslash{}\allowbreak{}textup\{\allowbreak{}\textbackslash{}\allowbreak{}texttt\{\allowbreak{}[\allowbreak{}rather\allowbreak{} a]\allowbreak{}\}\allowbreak{}\}\allowbreak{}}{DER\_\allowbreak{}XI-XIIa:\allowbreak{}67;\allowbreak{}\allowbreak{} DER\_\allowbreak{}XI-XIIa:\allowbreak{}68}
\DERapparatusitem{ENV0684}{37;\allowbreak{}\allowbreak{} XII\allowbreak{} ms51,\allowbreak{}\allowbreak{} Theorem\allowbreak{} 3(\allowbreak{}b1)\allowbreak{},\allowbreak{}\allowbreak{} parenthetical\allowbreak{} base-map\allowbreak{} explanation}{\$\allowbreak{}Y/\allowbreak{}s\_\allowbreak{}0\allowbreak{} \textbackslash{}\allowbreak{}longrightarrow\allowbreak{} Y/\allowbreak{}s\_\allowbreak{}1\$\allowbreak{}}{\$\allowbreak{}S/\allowbreak{}s\_\allowbreak{}0\allowbreak{} \textbackslash{}\allowbreak{}longrightarrow\allowbreak{} S/\allowbreak{}s\_\allowbreak{}1\$\allowbreak{}}{DER\_\allowbreak{}XI-XIIa:\allowbreak{}67}
\DERapparatusitem{ENV0685}{37;\allowbreak{}\allowbreak{} XII\allowbreak{} ms51,\allowbreak{}\allowbreak{} Theorem\allowbreak{} 3(\allowbreak{}b2)\allowbreak{},\allowbreak{}\allowbreak{} transport\allowbreak{} direction\allowbreak{} and\allowbreak{} base\allowbreak{} labels}{\$\allowbreak{}s\_\allowbreak{}0\textbackslash{}\allowbreak{}backslash\allowbreak{} X\allowbreak{} \textbackslash{}\allowbreak{}longrightarrow\allowbreak{} s\_\allowbreak{}1\textbackslash{}\allowbreak{}backslash\allowbreak{} X\$\allowbreak{}}{\$\allowbreak{}s\_\allowbreak{}1\textbackslash{}\allowbreak{}backslash\allowbreak{} X\allowbreak{} \textbackslash{}\allowbreak{}longrightarrow\allowbreak{} s\_\allowbreak{}0\textbackslash{}\allowbreak{}backslash\allowbreak{} X\$\allowbreak{}}{DER\_\allowbreak{}XI-XIIa:\allowbreak{}68}
\DERapparatusitem{ENV0686}{37;\allowbreak{}\allowbreak{} XII\allowbreak{} ms51,\allowbreak{}\allowbreak{} Theorem\allowbreak{} 3(\allowbreak{}b1)\allowbreak{},\allowbreak{}\allowbreak{} parenthetical\allowbreak{} base-map\allowbreak{} explanation}{a\$\allowbreak{}'\$\allowbreak{}\allowbreak{} \textbackslash{}\allowbreak{}textup\{\allowbreak{}\textbackslash{}\allowbreak{}texttt\{\allowbreak{}[\allowbreak{}rather\allowbreak{} a]\allowbreak{}\}\allowbreak{}\}\allowbreak{}}{a\allowbreak{} \textbackslash{}\allowbreak{}textup\{\allowbreak{}\textbackslash{}\allowbreak{}texttt\{\allowbreak{}[\allowbreak{}rather\allowbreak{} a]\allowbreak{}\}\allowbreak{}\}\allowbreak{}}{DER\_\allowbreak{}XI-XIIa:\allowbreak{}67;\allowbreak{}\allowbreak{} DER\_\allowbreak{}XI-XIIa:\allowbreak{}68}
\DERapparatusitem{ENV0687}{37;\allowbreak{}\allowbreak{} XII\allowbreak{} ms51,\allowbreak{}\allowbreak{} Theorem\allowbreak{} 3(\allowbreak{}b2)\allowbreak{},\allowbreak{}\allowbreak{} transport\allowbreak{} direction\allowbreak{} and\allowbreak{} base\allowbreak{} labels}{\$\allowbreak{}s\_\allowbreak{}1\textbackslash{}\allowbreak{}backslash\allowbreak{} Y\allowbreak{} \textbackslash{}\allowbreak{}longrightarrow\allowbreak{} s\_\allowbreak{}0\textbackslash{}\allowbreak{}backslash\allowbreak{} Y\$\allowbreak{}}{\$\allowbreak{}s\_\allowbreak{}1\textbackslash{}\allowbreak{}backslash\allowbreak{} S\allowbreak{} \textbackslash{}\allowbreak{}longrightarrow\allowbreak{} s\_\allowbreak{}0\textbackslash{}\allowbreak{}backslash\allowbreak{} S\$\allowbreak{}}{DER\_\allowbreak{}XI-XIIa:\allowbreak{}68}
\DERapparatusitem{ENV0688}{37;\allowbreak{}\allowbreak{} XII\allowbreak{} ms52,\allowbreak{}\allowbreak{} Definition\allowbreak{} 4\allowbreak{} cross-reference}{equivalent\allowbreak{} conditions\allowbreak{} a)\allowbreak{}\allowbreak{} to\allowbreak{} c\$\allowbreak{}\_\allowbreak{}2\$\allowbreak{})\allowbreak{}\allowbreak{} of\allowbreak{} proposition\allowbreak{} 9}{equivalent\allowbreak{} conditions\allowbreak{} a)\allowbreak{}\allowbreak{} to\allowbreak{} c\$\allowbreak{}\_\allowbreak{}2\$\allowbreak{})\allowbreak{}\allowbreak{} of\allowbreak{} theorem\allowbreak{} 3}{DER\_\allowbreak{}XI-XIIa:\allowbreak{}70}
\DERapparatusitem{ENV0689}{38;\allowbreak{}\allowbreak{} XII\allowbreak{} ms52,\allowbreak{}\allowbreak{} proof\allowbreak{} of\allowbreak{} Proposition\allowbreak{} 10}{it\allowbreak{} even\allowbreak{} suffices\allowbreak{} that\allowbreak{} \$\allowbreak{}X\$\allowbreak{}\allowbreak{} be\allowbreak{} preproper}{it\allowbreak{} even\allowbreak{} suffices\allowbreak{} that\allowbreak{} \$\allowbreak{}X\$\allowbreak{}\allowbreak{} be\allowbreak{} presmooth}{DER\_\allowbreak{}XI-XIIa:\allowbreak{}72}
\DERapparatusitem{ENV0690}{38;\allowbreak{}\allowbreak{} XII\allowbreak{} ms52,\allowbreak{}\allowbreak{} final\allowbreak{} reference\allowbreak{} in\allowbreak{} Proposition\allowbreak{} 10\allowbreak{} proof}{by\allowbreak{} \textbackslash{}\allowbreak{}texttt\{\allowbreak{}[\allowbreak{}the]\allowbreak{}\}\allowbreak{}\allowbreak{} proposition\allowbreak{} 5\allowbreak{} (\allowbreak{}page\allowbreak{} 29)\allowbreak{}}{by\allowbreak{} the\allowbreak{} \textbackslash{}\allowbreak{}texttt\{\allowbreak{}[\allowbreak{}the]\allowbreak{}\}\allowbreak{}\allowbreak{} corollary\allowbreak{} to\allowbreak{} proposition\allowbreak{} 5\allowbreak{} (\allowbreak{}page\allowbreak{} 30)\allowbreak{}}{DER\_\allowbreak{}XI-XIIa:\allowbreak{}71}
\DERapparatusitem{ENV0691}{39;\allowbreak{}\allowbreak{} XII\allowbreak{} ms53,\allowbreak{}\allowbreak{} W(\allowbreak{}4a)\allowbreak{}\allowbreak{} definition\allowbreak{} reference}{def.\allowbreak{}\allowbreak{} 3,\allowbreak{}\allowbreak{} p.\allowbreak{}\allowbreak{} 52}{def.\allowbreak{}\allowbreak{} 4,\allowbreak{}\allowbreak{} p.\allowbreak{}\allowbreak{} 52}{DER\_\allowbreak{}XI-XIIa:\allowbreak{}73}
\DERapparatusitem{ENV0692}{41;\allowbreak{}\allowbreak{} XII\allowbreak{} ms55,\allowbreak{}\allowbreak{} paragraph\allowbreak{} before\allowbreak{} the\allowbreak{} corollary}{Besides,\allowbreak{}\allowbreak{} since\allowbreak{} \$\allowbreak{}X\$\allowbreak{}\allowbreak{} is\allowbreak{} \$\allowbreak{}W\$\allowbreak{}-fibered\allowbreak{} over\allowbreak{} \$\allowbreak{}S\$\allowbreak{},\allowbreak{}\allowbreak{} practically\allowbreak{} by\allowbreak{} definition,\allowbreak{}\allowbreak{} \$\allowbreak{}\textbackslash{}\allowbreak{}xi\$\allowbreak{}\allowbreak{} is\allowbreak{} in\allowbreak{} \$\allowbreak{}\textbackslash{}\allowbreak{}HOT\textasciicircum{}\allowbreak{}\{\allowbreak{}\textbackslash{}\allowbreak{}mathrm\{\allowbreak{}lc\}\allowbreak{}\}\allowbreak{}\_\allowbreak{}W(\allowbreak{}S)\allowbreak{}\$\allowbreak{}}{Besides,\allowbreak{}\allowbreak{} since\allowbreak{} \$\allowbreak{}X\$\allowbreak{}\allowbreak{} is\allowbreak{} \$\allowbreak{}W\$\allowbreak{}-fibered\allowbreak{} over\allowbreak{} \$\allowbreak{}Y\$\allowbreak{},\allowbreak{}\allowbreak{} practically\allowbreak{} by\allowbreak{} definition,\allowbreak{}\allowbreak{} \$\allowbreak{}\textbackslash{}\allowbreak{}xi\$\allowbreak{}\allowbreak{} is\allowbreak{} in\allowbreak{} \$\allowbreak{}\textbackslash{}\allowbreak{}HOT\textasciicircum{}\allowbreak{}\{\allowbreak{}\textbackslash{}\allowbreak{}mathrm\{\allowbreak{}lc\}\allowbreak{}\}\allowbreak{}\_\allowbreak{}W(\allowbreak{}Y)\allowbreak{}\$\allowbreak{}}{DER\_\allowbreak{}XI-XIIa:\allowbreak{}75}
\DERapparatusitem{ENV0693}{41;\allowbreak{}\allowbreak{} XII\allowbreak{} ms56,\allowbreak{}\allowbreak{} constant\allowbreak{} coefficients\allowbreak{} in\allowbreak{} the\allowbreak{} corollary\allowbreak{} proof}{where\allowbreak{} \$\allowbreak{}\textbackslash{}\allowbreak{}eta\_\allowbreak{}S\allowbreak{} =\allowbreak{}\allowbreak{} p\textasciicircum{}\allowbreak{}*(\allowbreak{}\textbackslash{}\allowbreak{}eta)\allowbreak{}\$\allowbreak{},\allowbreak{}\allowbreak{} \$\allowbreak{}\textbackslash{}\allowbreak{}eta\_\allowbreak{}\{\allowbreak{}S'\}\allowbreak{}\allowbreak{} =\allowbreak{}\allowbreak{} p'\textasciicircum{}\allowbreak{}*(\allowbreak{}\textbackslash{}\allowbreak{}eta)\allowbreak{}\$\allowbreak{}}{where\allowbreak{} \$\allowbreak{}\textbackslash{}\allowbreak{}eta\_\allowbreak{}S\allowbreak{} =\allowbreak{}\allowbreak{} g\textasciicircum{}\allowbreak{}*(\allowbreak{}\textbackslash{}\allowbreak{}eta)\allowbreak{}\$\allowbreak{},\allowbreak{}\allowbreak{} \$\allowbreak{}\textbackslash{}\allowbreak{}eta\_\allowbreak{}\{\allowbreak{}S'\}\allowbreak{}\allowbreak{} =\allowbreak{}\allowbreak{} g'\textasciicircum{}\allowbreak{}*(\allowbreak{}\textbackslash{}\allowbreak{}eta)\allowbreak{}\$\allowbreak{}}{DER\_\allowbreak{}XI-XIIa:\allowbreak{}76}
\DERapparatusitem{ENV0694}{42;\allowbreak{}\allowbreak{} XII\allowbreak{} ms57,\allowbreak{}\allowbreak{} application\allowbreak{} of\allowbreak{} Proposition\allowbreak{} 2}{f\allowbreak{} :\allowbreak{}\allowbreak{} \textbackslash{}\allowbreak{}Phi(\allowbreak{}Y)\allowbreak{}\allowbreak{} \textbackslash{}\allowbreak{}longrightarrow\allowbreak{} S}{f\allowbreak{} :\allowbreak{}\allowbreak{} \textbackslash{}\allowbreak{}Phi(\allowbreak{}Z)\allowbreak{}\allowbreak{} \textbackslash{}\allowbreak{}longrightarrow\allowbreak{} S}{DER\_\allowbreak{}XI-XIIa:\allowbreak{}77}
\subsection*{Chapter XII, 5}
\DERapparatusitem{ENV0714}{43;\allowbreak{}\allowbreak{} ms59;\allowbreak{}\allowbreak{} printed43;\allowbreak{}\allowbreak{} displayed\allowbreak{} specialization\allowbreak{} relation}{x\allowbreak{} \textbackslash{}\allowbreak{}in\allowbreak{} \textbackslash{}\allowbreak{}overline\{\allowbreak{}\textbackslash{}\allowbreak{}\{\allowbreak{}y\textbackslash{}\allowbreak{}\}\allowbreak{}\}\allowbreak{}}{y\allowbreak{} \textbackslash{}\allowbreak{}in\allowbreak{} \textbackslash{}\allowbreak{}overline\{\allowbreak{}\textbackslash{}\allowbreak{}\{\allowbreak{}x\textbackslash{}\allowbreak{}\}\allowbreak{}\}\allowbreak{}}{DER\_\allowbreak{}XIIb-a:\allowbreak{}1}
\DERapparatusitem{ENV0715}{44;\allowbreak{}\allowbreak{} running\allowbreak{} text\allowbreak{} or\allowbreak{} formula}{\textbackslash{}\allowbreak{}text\{\allowbreak{}def\}\allowbreak{}}{\textbackslash{}\allowbreak{}mathrm\{\allowbreak{}def\}\allowbreak{}}{TYPESETTING:\allowbreak{}7;\allowbreak{}\allowbreak{} NOTATION\_\allowbreak{}DER:\allowbreak{}17}
\DERapparatusitem{ENV0716}{44;\allowbreak{}\allowbreak{} ms61;\allowbreak{}\allowbreak{} printed44;\allowbreak{}\allowbreak{} representability\allowbreak{} sentence\allowbreak{} and\allowbreak{} (\allowbreak{}*)\allowbreak{}}{(\allowbreak{}X'',\allowbreak{}U'',\allowbreak{}T'')\allowbreak{}\allowbreak{} \textbackslash{}\allowbreak{}longmapsto\allowbreak{} \textbackslash{}\allowbreak{}mathrm\{\allowbreak{}Hom\}\allowbreak{}\_\allowbreak{}\{\allowbreak{}\textbackslash{}\allowbreak{}mathrm\{\allowbreak{}Cat\}\allowbreak{}\}\allowbreak{}(\allowbreak{}U'',\allowbreak{}U)\allowbreak{}\allowbreak{} \textbackslash{}\allowbreak{}times\allowbreak{} \textbackslash{}\allowbreak{}mathrm\{\allowbreak{}Hom\}\allowbreak{}\_\allowbreak{}\{\allowbreak{}\textbackslash{}\allowbreak{}mathrm\{\allowbreak{}Cat\}\allowbreak{}\}\allowbreak{}(\allowbreak{}T'',\allowbreak{}T)\allowbreak{}}{\textbackslash{}\allowbreak{}begin\{\allowbreak{}aligned\}\allowbreak{}(\allowbreak{}X'',\allowbreak{}U'',\allowbreak{}T'')\allowbreak{}\textbackslash{}\allowbreak{}longmapsto\{\allowbreak{}\}\allowbreak{}\&\allowbreak{}\textbackslash{}\allowbreak{}mathrm\{\allowbreak{}Hom\}\allowbreak{}\_\allowbreak{}\{\allowbreak{}\textbackslash{}\allowbreak{}mathcal\{\allowbreak{}T\}\allowbreak{}\}\allowbreak{}(\allowbreak{}(\allowbreak{}X'',\allowbreak{}U'',\allowbreak{}T'')\allowbreak{},\allowbreak{}(\allowbreak{}X,\allowbreak{}U,\allowbreak{}T)\allowbreak{})\allowbreak{}\textbackslash{}\allowbreak{}\textbackslash{}\allowbreak{}\allowbreak{} \&\allowbreak{}\{\allowbreak{}\}\allowbreak{}\textbackslash{}\allowbreak{}times\_\allowbreak{}\{\allowbreak{}\textbackslash{}\allowbreak{}substack\{\allowbreak{}\textbackslash{}\allowbreak{}mathrm\{\allowbreak{}Hom\}\allowbreak{}\_\allowbreak{}\{\allowbreak{}\textbackslash{}\allowbreak{}mathrm\{\allowbreak{}Cat\}\allowbreak{}\}\allowbreak{}(\allowbreak{}U'',\allowbreak{}U)\allowbreak{}\textbackslash{}\allowbreak{}\textbackslash{}\allowbreak{}\{\allowbreak{}\}\allowbreak{}\textbackslash{}\allowbreak{}times\allowbreak{} \textbackslash{}\allowbreak{}mathrm\{\allowbreak{}Hom\}\allowbreak{}\_\allowbreak{}\{\allowbreak{}\textbackslash{}\allowbreak{}mathrm\{\allowbreak{}Cat\}\allowbreak{}\}\allowbreak{}(\allowbreak{}T'',\allowbreak{}T)\allowbreak{}\}\allowbreak{}\}\allowbreak{}\textbackslash{}\allowbreak{}bigl(\allowbreak{}\textbackslash{}\allowbreak{}mathrm\{\allowbreak{}Hom\}\allowbreak{}\_\allowbreak{}\{\allowbreak{}\textbackslash{}\allowbreak{}mathrm\{\allowbreak{}Cat\}\allowbreak{}\}\allowbreak{}(\allowbreak{}U'',\allowbreak{}U')\allowbreak{}\allowbreak{} \textbackslash{}\allowbreak{}times\allowbreak{} \textbackslash{}\allowbreak{}mathrm\{\allowbreak{}Hom\}\allowbreak{}\_\allowbreak{}\{\allowbreak{}\textbackslash{}\allowbreak{}mathrm\{\allowbreak{}Cat\}\allowbreak{}\}\allowbreak{}(\allowbreak{}T'',\allowbreak{}T')\allowbreak{}\textbackslash{}\allowbreak{}bigr)\allowbreak{}\textbackslash{}\allowbreak{}end\{\allowbreak{}aligned\}\allowbreak{}}{DER\_\allowbreak{}XIIb-a:\allowbreak{}4}
\DERapparatusitem{ENV0717}{44;\allowbreak{}\allowbreak{} ms61;\allowbreak{}\allowbreak{} printed44;\allowbreak{}\allowbreak{} representability\allowbreak{} sentence\allowbreak{} and\allowbreak{} (\allowbreak{}*)\allowbreak{}}{\textbackslash{}\allowbreak{}mathrm\{\allowbreak{}Hom\}\allowbreak{}\_\allowbreak{}\{\allowbreak{}\textbackslash{}\allowbreak{}mathcal\{\allowbreak{}T\}\allowbreak{}\}\allowbreak{}(\allowbreak{}(\allowbreak{}X'',\allowbreak{}U'',\allowbreak{}T'')\allowbreak{},\allowbreak{}(\allowbreak{}X',\allowbreak{}U',\allowbreak{}T')\allowbreak{})\allowbreak{}\allowbreak{} \textbackslash{}\allowbreak{}simeq\allowbreak{} \textbackslash{}\allowbreak{}mathrm\{\allowbreak{}Hom\}\allowbreak{}\_\allowbreak{}\{\allowbreak{}\textbackslash{}\allowbreak{}mathrm\{\allowbreak{}Cat\}\allowbreak{}\}\allowbreak{}(\allowbreak{}U'',\allowbreak{}U)\allowbreak{}\allowbreak{} \textbackslash{}\allowbreak{}times\allowbreak{} \textbackslash{}\allowbreak{}mathrm\{\allowbreak{}Hom\}\allowbreak{}\_\allowbreak{}\{\allowbreak{}\textbackslash{}\allowbreak{}mathrm\{\allowbreak{}Cat\}\allowbreak{}\}\allowbreak{}(\allowbreak{}T'',\allowbreak{}T)\allowbreak{}}{\textbackslash{}\allowbreak{}begin\{\allowbreak{}aligned\}\allowbreak{}\textbackslash{}\allowbreak{}mathrm\{\allowbreak{}Hom\}\allowbreak{}\_\allowbreak{}\{\allowbreak{}\textbackslash{}\allowbreak{}mathcal\{\allowbreak{}T\}\allowbreak{}\}\allowbreak{}(\allowbreak{}(\allowbreak{}X'',\allowbreak{}U'',\allowbreak{}T'')\allowbreak{},\allowbreak{}(\allowbreak{}X',\allowbreak{}U',\allowbreak{}T')\allowbreak{})\allowbreak{}\allowbreak{} \&\allowbreak{}\textbackslash{}\allowbreak{}simeq\allowbreak{} \textbackslash{}\allowbreak{}mathrm\{\allowbreak{}Hom\}\allowbreak{}\_\allowbreak{}\{\allowbreak{}\textbackslash{}\allowbreak{}mathcal\{\allowbreak{}T\}\allowbreak{}\}\allowbreak{}(\allowbreak{}(\allowbreak{}X'',\allowbreak{}U'',\allowbreak{}T'')\allowbreak{},\allowbreak{}(\allowbreak{}X,\allowbreak{}U,\allowbreak{}T)\allowbreak{})\allowbreak{}\textbackslash{}\allowbreak{}\textbackslash{}\allowbreak{}\allowbreak{} \&\allowbreak{}\textbackslash{}\allowbreak{}quad\{\allowbreak{}\}\allowbreak{}\textbackslash{}\allowbreak{}times\_\allowbreak{}\{\allowbreak{}\textbackslash{}\allowbreak{}substack\{\allowbreak{}\textbackslash{}\allowbreak{}mathrm\{\allowbreak{}Hom\}\allowbreak{}\_\allowbreak{}\{\allowbreak{}\textbackslash{}\allowbreak{}mathrm\{\allowbreak{}Cat\}\allowbreak{}\}\allowbreak{}(\allowbreak{}U'',\allowbreak{}U)\allowbreak{}\textbackslash{}\allowbreak{}\textbackslash{}\allowbreak{}\{\allowbreak{}\}\allowbreak{}\textbackslash{}\allowbreak{}times\allowbreak{} \textbackslash{}\allowbreak{}mathrm\{\allowbreak{}Hom\}\allowbreak{}\_\allowbreak{}\{\allowbreak{}\textbackslash{}\allowbreak{}mathrm\{\allowbreak{}Cat\}\allowbreak{}\}\allowbreak{}(\allowbreak{}T'',\allowbreak{}T)\allowbreak{}\}\allowbreak{}\}\allowbreak{}\textbackslash{}\allowbreak{}bigl(\allowbreak{}\textbackslash{}\allowbreak{}mathrm\{\allowbreak{}Hom\}\allowbreak{}\_\allowbreak{}\{\allowbreak{}\textbackslash{}\allowbreak{}mathrm\{\allowbreak{}Cat\}\allowbreak{}\}\allowbreak{}(\allowbreak{}U'',\allowbreak{}U')\allowbreak{}\allowbreak{} \textbackslash{}\allowbreak{}times\allowbreak{} \textbackslash{}\allowbreak{}mathrm\{\allowbreak{}Hom\}\allowbreak{}\_\allowbreak{}\{\allowbreak{}\textbackslash{}\allowbreak{}mathrm\{\allowbreak{}Cat\}\allowbreak{}\}\allowbreak{}(\allowbreak{}T'',\allowbreak{}T')\allowbreak{}\textbackslash{}\allowbreak{}bigr)\allowbreak{}\textbackslash{}\allowbreak{}end\{\allowbreak{}aligned\}\allowbreak{}}{DER\_\allowbreak{}XIIb-a:\allowbreak{}4}
\DERapparatusitem{ENV0718}{45;\allowbreak{}\allowbreak{} ms61;\allowbreak{}\allowbreak{} printed45;\allowbreak{}\allowbreak{} displayed\allowbreak{} (\allowbreak{}**)\allowbreak{}}{\textbackslash{}\allowbreak{}uHom\_\allowbreak{}\{\allowbreak{}\textbackslash{}\allowbreak{}mathcal\{\allowbreak{}T\}\allowbreak{}\}\allowbreak{}(\allowbreak{}(\allowbreak{}X'',\allowbreak{}U'',\allowbreak{}T'')\allowbreak{},\allowbreak{}(\allowbreak{}X',\allowbreak{}U',\allowbreak{}T')\allowbreak{})\allowbreak{}\allowbreak{} \textbackslash{}\allowbreak{}simeq\allowbreak{} \textbackslash{}\allowbreak{}uHom\_\allowbreak{}\{\allowbreak{}\textbackslash{}\allowbreak{}mathrm\{\allowbreak{}Cat\}\allowbreak{}\}\allowbreak{}(\allowbreak{}U'',\allowbreak{}U)\allowbreak{}\allowbreak{} \textbackslash{}\allowbreak{}times\allowbreak{} \textbackslash{}\allowbreak{}uHom\_\allowbreak{}\{\allowbreak{}\textbackslash{}\allowbreak{}mathrm\{\allowbreak{}Cat\}\allowbreak{}\}\allowbreak{}(\allowbreak{}T'',\allowbreak{}T)\allowbreak{}}{\textbackslash{}\allowbreak{}begin\{\allowbreak{}aligned\}\allowbreak{}\textbackslash{}\allowbreak{}uHom\_\allowbreak{}\{\allowbreak{}\textbackslash{}\allowbreak{}mathcal\{\allowbreak{}T\}\allowbreak{}\}\allowbreak{}(\allowbreak{}(\allowbreak{}X'',\allowbreak{}U'',\allowbreak{}T'')\allowbreak{},\allowbreak{}(\allowbreak{}X',\allowbreak{}U',\allowbreak{}T')\allowbreak{})\allowbreak{}\allowbreak{} \&\allowbreak{}\textbackslash{}\allowbreak{}simeq\allowbreak{} \textbackslash{}\allowbreak{}uHom\_\allowbreak{}\{\allowbreak{}\textbackslash{}\allowbreak{}mathcal\{\allowbreak{}T\}\allowbreak{}\}\allowbreak{}(\allowbreak{}(\allowbreak{}X'',\allowbreak{}U'',\allowbreak{}T'')\allowbreak{},\allowbreak{}(\allowbreak{}X,\allowbreak{}U,\allowbreak{}T)\allowbreak{})\allowbreak{}\textbackslash{}\allowbreak{}\textbackslash{}\allowbreak{}\allowbreak{} \&\allowbreak{}\textbackslash{}\allowbreak{}quad\{\allowbreak{}\}\allowbreak{}\textbackslash{}\allowbreak{}times\_\allowbreak{}\{\allowbreak{}\textbackslash{}\allowbreak{}substack\{\allowbreak{}\textbackslash{}\allowbreak{}uHom\_\allowbreak{}\{\allowbreak{}\textbackslash{}\allowbreak{}mathrm\{\allowbreak{}Cat\}\allowbreak{}\}\allowbreak{}(\allowbreak{}U'',\allowbreak{}U)\allowbreak{}\textbackslash{}\allowbreak{}\textbackslash{}\allowbreak{}\{\allowbreak{}\}\allowbreak{}\textbackslash{}\allowbreak{}times\allowbreak{} \textbackslash{}\allowbreak{}uHom\_\allowbreak{}\{\allowbreak{}\textbackslash{}\allowbreak{}mathrm\{\allowbreak{}Cat\}\allowbreak{}\}\allowbreak{}(\allowbreak{}T'',\allowbreak{}T)\allowbreak{}\}\allowbreak{}\}\allowbreak{}\textbackslash{}\allowbreak{}bigl(\allowbreak{}\textbackslash{}\allowbreak{}uHom\_\allowbreak{}\{\allowbreak{}\textbackslash{}\allowbreak{}mathrm\{\allowbreak{}Cat\}\allowbreak{}\}\allowbreak{}(\allowbreak{}U'',\allowbreak{}U')\allowbreak{}\allowbreak{} \textbackslash{}\allowbreak{}times\allowbreak{} \textbackslash{}\allowbreak{}uHom\_\allowbreak{}\{\allowbreak{}\textbackslash{}\allowbreak{}mathrm\{\allowbreak{}Cat\}\allowbreak{}\}\allowbreak{}(\allowbreak{}T'',\allowbreak{}T')\allowbreak{}\textbackslash{}\allowbreak{}bigr)\allowbreak{}\textbackslash{}\allowbreak{}end\{\allowbreak{}aligned\}\allowbreak{}}{DER\_\allowbreak{}XIIb-a:\allowbreak{}5}
\DERapparatusitem{ENV0719}{48;\allowbreak{}\allowbreak{} ms66;\allowbreak{}\allowbreak{} printed48;\allowbreak{}\allowbreak{} heading1.\allowbreak{}4}{an\allowbreak{} open\allowbreak{} immersion\allowbreak{} where\allowbreak{} closed}{an\allowbreak{} open\allowbreak{} immersion\allowbreak{} or\allowbreak{} closed\allowbreak{} immersion}{DER\_\allowbreak{}XIIb-a:\allowbreak{}7}
\DERapparatusitem{ENV0720}{50;\allowbreak{}\allowbreak{} ms70;\allowbreak{}\allowbreak{} printed50;\allowbreak{}\allowbreak{} sentence\allowbreak{} immediately\allowbreak{} before\allowbreak{} (\allowbreak{}*)\allowbreak{}}{Here\allowbreak{} \$\allowbreak{}F=\allowbreak{}f\_\allowbreak{}U\textasciicircum{}\allowbreak{}\{\allowbreak{}\textbackslash{}\allowbreak{}circ\}\allowbreak{}\textbackslash{}\allowbreak{}times\allowbreak{} f\_\allowbreak{}T\$\allowbreak{}\allowbreak{} is\allowbreak{} a\allowbreak{} covariant\allowbreak{} functor}{Here\allowbreak{} \$\allowbreak{}F=\allowbreak{}f\_\allowbreak{}U\textasciicircum{}\allowbreak{}\{\allowbreak{}\textbackslash{}\allowbreak{}circ\}\allowbreak{}\textbackslash{}\allowbreak{}times\allowbreak{} f\_\allowbreak{}T\$\allowbreak{},\allowbreak{}\allowbreak{} and\allowbreak{} \$\allowbreak{}H\$\allowbreak{}\allowbreak{} is\allowbreak{} a\allowbreak{} covariant\allowbreak{} functor}{DER\_\allowbreak{}XIIb-a:\allowbreak{}9}
\DERapparatusitem{ENV0721}{54;\allowbreak{}\allowbreak{} ms75;\allowbreak{}\allowbreak{} printed54;\allowbreak{}\allowbreak{} Proposition5}{\$\allowbreak{}X\$\allowbreak{},\allowbreak{}\allowbreak{} \$\allowbreak{}Y'\$\allowbreak{}\allowbreak{} of\allowbreak{} \$\allowbreak{}X\$\allowbreak{}\allowbreak{} having\allowbreak{} as\allowbreak{} union}{\$\allowbreak{}X\$\allowbreak{},\allowbreak{}\allowbreak{} \$\allowbreak{}Y'\$\allowbreak{}\allowbreak{} of\allowbreak{} \$\allowbreak{}X'\$\allowbreak{}\allowbreak{} having\allowbreak{} as\allowbreak{} union}{DER\_\allowbreak{}XIIb-a:\allowbreak{}11}
\DERapparatusitem{ENV0722}{54;\allowbreak{}\allowbreak{} ms75;\allowbreak{}\allowbreak{} printed54;\allowbreak{}\allowbreak{} Proposition5\allowbreak{} converse}{in\allowbreak{} such\allowbreak{} a\allowbreak{} situationm\allowbreak{} the\allowbreak{} corresponding\allowbreak{} square}{in\allowbreak{} such\allowbreak{} a\allowbreak{} situation,\allowbreak{}\allowbreak{} the\allowbreak{} corresponding\allowbreak{} square}{DER\_\allowbreak{}XIIb-a:\allowbreak{}12}
\DERapparatusitem{ENV0723}{60;\allowbreak{}\allowbreak{} ms.\allowbreak{}82,\allowbreak{}\allowbreak{} upper-right\allowbreak{} corner\allowbreak{} of\allowbreak{} the\allowbreak{} second\allowbreak{} displayed\allowbreak{} square}{\&\allowbreak{}\allowbreak{} T\allowbreak{} \textbackslash{}\allowbreak{}arrow[\allowbreak{}d,\allowbreak{}\allowbreak{} "j'"]\allowbreak{}}{\&\allowbreak{}\allowbreak{} T'\allowbreak{} \textbackslash{}\allowbreak{}arrow[\allowbreak{}d,\allowbreak{}\allowbreak{} "j'"]\allowbreak{}}{DER\_\allowbreak{}XIIb-a:\allowbreak{}17}
\DERapparatusitem{ENV0724}{61;\allowbreak{}\allowbreak{} ms.\allowbreak{}84,\allowbreak{}\allowbreak{} sentence\allowbreak{} preceding\allowbreak{} Proposition\allowbreak{} 7}{the\allowbreak{} inclusion\allowbreak{} arrow,\allowbreak{}\allowbreak{} Let\allowbreak{} us}{the\allowbreak{} inclusion\allowbreak{} arrow.\allowbreak{}\allowbreak{} Let\allowbreak{} us}{DER\_\allowbreak{}XIIb-a:\allowbreak{}18}
\DERapparatusitem{ENV0725}{preamble;\allowbreak{}\allowbreak{} local\allowbreak{} macro\allowbreak{} definition}{\textbackslash{}\allowbreak{}text\{\allowbreak{}def\}\allowbreak{}}{\textbackslash{}\allowbreak{}mathrm\{\allowbreak{}def\}\allowbreak{}}{TYPESETTING:\allowbreak{}7;\allowbreak{}\allowbreak{} NOTATION\_\allowbreak{}DER:\allowbreak{}17}
\DERapparatusitem{ENV0726}{63;\allowbreak{}\allowbreak{} ms.\allowbreak{}88,\allowbreak{}\allowbreak{} first\allowbreak{} formula\allowbreak{} on\allowbreak{} printed\allowbreak{} p.\allowbreak{}63}{}{'}{DER\_\allowbreak{}XIIb-a:\allowbreak{}19}
\DERapparatusitem{ENV0727}{63;\allowbreak{}\allowbreak{} ms.\allowbreak{}89,\allowbreak{}\allowbreak{} right\allowbreak{} vertical\allowbreak{} label\allowbreak{} in\allowbreak{} the\allowbreak{} last\allowbreak{} square}{F/\allowbreak{}V}{V/\allowbreak{}F}{DER\_\allowbreak{}XIIb-a:\allowbreak{}20}
\DERapparatusitem{ENV0728}{65;\allowbreak{}\allowbreak{} ms.\allowbreak{}92,\allowbreak{}\allowbreak{} factorization\allowbreak{} property\allowbreak{} a)\allowbreak{}}{}{'}{DER\_\allowbreak{}XIIb-a:\allowbreak{}21}
\DERapparatusitem{ENV0729}{65;\allowbreak{}\allowbreak{} ms.\allowbreak{}93,\allowbreak{}\allowbreak{} final\allowbreak{} displayed\allowbreak{} formula}{'/\allowbreak{}V}{/\allowbreak{}V'}{DER\_\allowbreak{}XIIb-a:\allowbreak{}23}
\DERapparatusitem{ENV0730}{69;\allowbreak{}\allowbreak{} ms.\allowbreak{}99,\allowbreak{}\allowbreak{} footnote\allowbreak{} 70}{\textbackslash{}\allowbreak{}Phi}{\textbackslash{}\allowbreak{}Psi}{DER\_\allowbreak{}XIIb-a:\allowbreak{}26}
\DERapparatusitem{ENV0731}{70;\allowbreak{}\allowbreak{} running\allowbreak{} text\allowbreak{} or\allowbreak{} formula}{\textbackslash{}\allowbreak{}text\{\allowbreak{}def\}\allowbreak{}}{\textbackslash{}\allowbreak{}mathrm\{\allowbreak{}def\}\allowbreak{}}{TYPESETTING:\allowbreak{}7;\allowbreak{}\allowbreak{} NOTATION\_\allowbreak{}DER:\allowbreak{}17}
\DERapparatusitem{ENV0732}{70;\allowbreak{}\allowbreak{} ms.\allowbreak{}100,\allowbreak{}\allowbreak{} paragraph\allowbreak{} after\allowbreak{} Corollary\allowbreak{} 2}{part\allowbreak{} b)\allowbreak{}\allowbreak{} (\allowbreak{}corrected)\allowbreak{}\allowbreak{} of\allowbreak{} page\allowbreak{} 9}{part\allowbreak{} b)\allowbreak{}\allowbreak{} (\allowbreak{}corrected)\allowbreak{}\allowbreak{} of\allowbreak{} proposition\allowbreak{} 9\allowbreak{} (\allowbreak{}ms.\allowbreak{}97--100)\allowbreak{}}{DER\_\allowbreak{}XIIb-a:\allowbreak{}27}
\DERapparatusitem{ENV0733}{71;\allowbreak{}\allowbreak{} running\allowbreak{} text\allowbreak{} or\allowbreak{} formula}{\textbackslash{}\allowbreak{}text\{\allowbreak{}def\}\allowbreak{}}{\textbackslash{}\allowbreak{}mathrm\{\allowbreak{}def\}\allowbreak{}}{TYPESETTING:\allowbreak{}7;\allowbreak{}\allowbreak{} NOTATION\_\allowbreak{}DER:\allowbreak{}17}
\DERapparatusitem{ENV0734}{72;\allowbreak{}\allowbreak{} ms.\allowbreak{}103,\allowbreak{}\allowbreak{} six-object\allowbreak{} diagram}{e\_\allowbreak{}1\textasciicircum{}\allowbreak{}+\allowbreak{} \textbackslash{}\allowbreak{}bullet\allowbreak{} \textbackslash{}\allowbreak{}arrow[\allowbreak{}ur]\allowbreak{}\allowbreak{} \&\allowbreak{}\allowbreak{} \textbackslash{}\allowbreak{}bullet\textbackslash{}\allowbreak{},\allowbreak{}\allowbreak{} e\_\allowbreak{}1\textasciicircum{}\allowbreak{}-\allowbreak{} \textbackslash{}\allowbreak{}arrow[\allowbreak{}ul]\allowbreak{}}{e\_\allowbreak{}1\textasciicircum{}\allowbreak{}+\allowbreak{} \textbackslash{}\allowbreak{}bullet\allowbreak{} \textbackslash{}\allowbreak{}arrow[\allowbreak{}ur]\allowbreak{}\allowbreak{} \textbackslash{}\allowbreak{}arrow[\allowbreak{}u]\allowbreak{}\allowbreak{} \&\allowbreak{}\allowbreak{} \textbackslash{}\allowbreak{}bullet\textbackslash{}\allowbreak{},\allowbreak{}\allowbreak{} e\_\allowbreak{}1\textasciicircum{}\allowbreak{}-\allowbreak{} \textbackslash{}\allowbreak{}arrow[\allowbreak{}ul]\allowbreak{}\allowbreak{} \textbackslash{}\allowbreak{}arrow[\allowbreak{}u]\allowbreak{}}{DER\_\allowbreak{}XIIb-a:\allowbreak{}28}
\DERapparatusitem{ENV0735}{72;\allowbreak{}\allowbreak{} ms.\allowbreak{}103,\allowbreak{}\allowbreak{} six-object\allowbreak{} diagram}{e\_\allowbreak{}0\textasciicircum{}\allowbreak{}+\allowbreak{} \textbackslash{}\allowbreak{}bullet\allowbreak{} \textbackslash{}\allowbreak{}arrow[\allowbreak{}ur]\allowbreak{}\allowbreak{} \&\allowbreak{}\allowbreak{} \textbackslash{}\allowbreak{}bullet\textbackslash{}\allowbreak{},\allowbreak{}\allowbreak{} e\_\allowbreak{}0\textasciicircum{}\allowbreak{}-\allowbreak{} \textbackslash{}\allowbreak{}arrow[\allowbreak{}ul]\allowbreak{}}{e\_\allowbreak{}0\textasciicircum{}\allowbreak{}+\allowbreak{} \textbackslash{}\allowbreak{}bullet\allowbreak{} \textbackslash{}\allowbreak{}arrow[\allowbreak{}ur]\allowbreak{}\allowbreak{} \textbackslash{}\allowbreak{}arrow[\allowbreak{}u]\allowbreak{}\allowbreak{} \&\allowbreak{}\allowbreak{} \textbackslash{}\allowbreak{}bullet\textbackslash{}\allowbreak{},\allowbreak{}\allowbreak{} e\_\allowbreak{}0\textasciicircum{}\allowbreak{}-\allowbreak{} \textbackslash{}\allowbreak{}arrow[\allowbreak{}ul]\allowbreak{}\allowbreak{} \textbackslash{}\allowbreak{}arrow[\allowbreak{}u]\allowbreak{}}{DER\_\allowbreak{}XIIb-a:\allowbreak{}28}
\DERapparatusitem{ENV0736}{73;\allowbreak{}\allowbreak{} running\allowbreak{} text\allowbreak{} or\allowbreak{} formula}{\textbackslash{}\allowbreak{}text\{\allowbreak{}def\}\allowbreak{}}{\textbackslash{}\allowbreak{}mathrm\{\allowbreak{}def\}\allowbreak{}}{TYPESETTING:\allowbreak{}7;\allowbreak{}\allowbreak{} NOTATION\_\allowbreak{}DER:\allowbreak{}17}
\DERapparatusitem{ENV0737}{74;\allowbreak{}\allowbreak{} ms.\allowbreak{}107,\allowbreak{}\allowbreak{} Proposition11\allowbreak{} cross-reference}{(\allowbreak{}p.\allowbreak{}\allowbreak{} 124\allowbreak{} \textbackslash{}\allowbreak{}texttt\{\allowbreak{}[\allowbreak{}rather\allowbreak{} 104]\allowbreak{}\}\allowbreak{})\allowbreak{}}{(\allowbreak{}p.\allowbreak{}\allowbreak{} 104\allowbreak{} \textbackslash{}\allowbreak{}texttt\{\allowbreak{}[\allowbreak{}rather\allowbreak{} 104]\allowbreak{}\}\allowbreak{})\allowbreak{}}{DER\_\allowbreak{}XIIb-a:\allowbreak{}29}
\DERapparatusitem{ENV0738}{75;\allowbreak{}\allowbreak{} ms.\allowbreak{}109,\allowbreak{}\allowbreak{} proof\allowbreak{} of\allowbreak{} Corollary1}{cobase\allowbreak{} change\allowbreak{} \$\allowbreak{}U\textbackslash{}\allowbreak{}xrightarrow\{\allowbreak{}f\}\allowbreak{}U'\$\allowbreak{}}{cobase\allowbreak{} change\allowbreak{} \$\allowbreak{}U\textbackslash{}\allowbreak{}xrightarrow\{\allowbreak{}p\}\allowbreak{}U'\$\allowbreak{}}{DER\_\allowbreak{}XIIb-a:\allowbreak{}32}
\DERapparatusitem{ENV0739}{76;\allowbreak{}\allowbreak{} ms.\allowbreak{}110,\allowbreak{}\allowbreak{} Definition2(\allowbreak{}d)\allowbreak{}}{X'=\allowbreak{}U'\textbackslash{}\allowbreak{}amalg\_\allowbreak{}U\allowbreak{} X'\textbackslash{}\allowbreak{}quad}{X'=\allowbreak{}U'\textbackslash{}\allowbreak{}amalg\_\allowbreak{}U\allowbreak{} X\textbackslash{}\allowbreak{}quad}{DER\_\allowbreak{}XIIb-a:\allowbreak{}34}
\DERapparatusitem{ENV0740}{77;\allowbreak{}\allowbreak{} ms.\allowbreak{}112,\allowbreak{}\allowbreak{} first\allowbreak{} sentence\allowbreak{} of\allowbreak{} proof}{aspheric)\allowbreak{}.\allowbreak{}\allowbreak{} then}{aspheric)\allowbreak{},\allowbreak{}\allowbreak{} then}{DER\_\allowbreak{}XIIb-a:\allowbreak{}36}
\DERapparatusitem{ENV0741}{77;\allowbreak{}\allowbreak{} ms.\allowbreak{}112,\allowbreak{}\allowbreak{} first\allowbreak{} sentence\allowbreak{} of\allowbreak{} proof}{it\allowbreak{} will\allowbreak{} follow\allowbreak{} that\allowbreak{} \$\allowbreak{}i'\$\allowbreak{}\allowbreak{} is\allowbreak{} \$\allowbreak{}W\$\allowbreak{}-aspheric}{it\allowbreak{} will\allowbreak{} follow\allowbreak{} that\allowbreak{} \$\allowbreak{}i'\_\allowbreak{}0\$\allowbreak{}\allowbreak{} is\allowbreak{} \$\allowbreak{}W\$\allowbreak{}-aspheric}{DER\_\allowbreak{}XIIb-a:\allowbreak{}37}
\DERapparatusitem{ENV0742}{77;\allowbreak{}\allowbreak{} ms.\allowbreak{}112,\allowbreak{}\allowbreak{} boundary\allowbreak{} weight\allowbreak{} in\allowbreak{} proof}{the\allowbreak{} element\allowbreak{} \$\allowbreak{}i\textasciicircum{}\allowbreak{}*(\allowbreak{}t')\allowbreak{}}{the\allowbreak{} element\allowbreak{} \$\allowbreak{}i'\textasciicircum{}\allowbreak{}*(\allowbreak{}t')\allowbreak{}}{DER\_\allowbreak{}XIIb-a:\allowbreak{}38}
\DERapparatusitem{ENV0743}{77;\allowbreak{}\allowbreak{} ms.\allowbreak{}112,\allowbreak{}\allowbreak{} displayed\allowbreak{} direct-image\allowbreak{} identity}{\textbackslash{}\allowbreak{}[\allowbreak{}\allowbreak{} i'\textasciicircum{}\allowbreak{}*(\allowbreak{}t)\allowbreak{}=\allowbreak{}p\_\allowbreak{}!(\allowbreak{}i\textasciicircum{}\allowbreak{}*(\allowbreak{}t)\allowbreak{})\allowbreak{}}{\textbackslash{}\allowbreak{}[\allowbreak{}\allowbreak{} i'\textasciicircum{}\allowbreak{}*(\allowbreak{}t')\allowbreak{}=\allowbreak{}p\_\allowbreak{}!(\allowbreak{}i\textasciicircum{}\allowbreak{}*(\allowbreak{}t)\allowbreak{})\allowbreak{}}{DER\_\allowbreak{}XIIb-a:\allowbreak{}39}
\DERapparatusitem{ENV0744}{77;\allowbreak{}\allowbreak{} ms.\allowbreak{}112,\allowbreak{}\allowbreak{} paragraph\allowbreak{} before\allowbreak{} lemma}{\$\allowbreak{}p\_\allowbreak{}!(\allowbreak{}F)\allowbreak{}\$\allowbreak{}\allowbreak{} is\allowbreak{} a\allowbreak{} \$\allowbreak{}W\$\allowbreak{}-aspheric\allowbreak{} element\allowbreak{} of\allowbreak{} \$\allowbreak{}U'\$\allowbreak{}}{\$\allowbreak{}p\_\allowbreak{}!(\allowbreak{}F)\allowbreak{}\$\allowbreak{}\allowbreak{} is\allowbreak{} a\allowbreak{} \$\allowbreak{}W\$\allowbreak{}-aspheric\allowbreak{} element\allowbreak{} of\allowbreak{} \$\allowbreak{}U'\textasciicircum{}\allowbreak{}\textbackslash{}\allowbreak{}wedge\$\allowbreak{}}{DER\_\allowbreak{}XIIb-a:\allowbreak{}41}
\DERapparatusitem{ENV0745}{77;\allowbreak{}\allowbreak{} ms.\allowbreak{}113,\allowbreak{}\allowbreak{} Corollary3}{X'=\allowbreak{}U'\textbackslash{}\allowbreak{}amalg\_\allowbreak{}U\allowbreak{} X'\$\allowbreak{}\}\allowbreak{}}{X'=\allowbreak{}U'\textbackslash{}\allowbreak{}amalg\_\allowbreak{}U\allowbreak{} X\$\allowbreak{}\}\allowbreak{}}{DER\_\allowbreak{}XIIb-a:\allowbreak{}42}
\DERapparatusitem{ENV0746}{78;\allowbreak{}\allowbreak{} ms.\allowbreak{}115,\allowbreak{}\allowbreak{} terminal-fibre\allowbreak{} inclusion}{of\allowbreak{} the\allowbreak{} fiber\allowbreak{} of\allowbreak{} \$\allowbreak{}X\$\allowbreak{}\allowbreak{} at\allowbreak{} \$\allowbreak{}i\_\allowbreak{}0\$\allowbreak{}\allowbreak{} is\allowbreak{} \$\allowbreak{}W\$\allowbreak{}-aspheric}{of\allowbreak{} the\allowbreak{} fiber\allowbreak{} of\allowbreak{} \$\allowbreak{}X\$\allowbreak{}\allowbreak{} at\allowbreak{} \$\allowbreak{}i\_\allowbreak{}0\$\allowbreak{}\allowbreak{} is\allowbreak{} \$\allowbreak{}W\$\allowbreak{}-coaspheric}{DER\_\allowbreak{}XIIb-a:\allowbreak{}44}
\DERapparatusitem{ENV0747}{79;\allowbreak{}\allowbreak{} ms.\allowbreak{}117,\allowbreak{}\allowbreak{} presheaf\allowbreak{} colimit}{\textbackslash{}\allowbreak{}varinjlim\_\allowbreak{}\{\allowbreak{}I\}\allowbreak{}\allowbreak{} G(\allowbreak{}i)\allowbreak{}}{\textbackslash{}\allowbreak{}varinjlim\_\allowbreak{}\{\allowbreak{}I\textasciicircum{}\allowbreak{}\textbackslash{}\allowbreak{}circ\}\allowbreak{}\allowbreak{} G(\allowbreak{}i)\allowbreak{}}{DER\_\allowbreak{}XIIb-a:\allowbreak{}46}
\DERapparatusitem{ENV0748}{81;\allowbreak{}\allowbreak{} ms.\allowbreak{}120,\allowbreak{}\allowbreak{} proof\allowbreak{} of\allowbreak{} Proposition11}{p\_\allowbreak{}\{\allowbreak{}i\_\allowbreak{}0\textbackslash{}\allowbreak{},\allowbreak{}!\}\allowbreak{}(\allowbreak{}X\_\allowbreak{}\{\allowbreak{}i\_\allowbreak{}0\}\allowbreak{})\allowbreak{}}{p\_\allowbreak{}\{\allowbreak{}i\_\allowbreak{}0\textbackslash{}\allowbreak{},\allowbreak{}!\}\allowbreak{}(\allowbreak{}e\_\allowbreak{}\{\allowbreak{}X\_\allowbreak{}\{\allowbreak{}i\_\allowbreak{}0\}\allowbreak{}\}\allowbreak{})\allowbreak{}}{DER\_\allowbreak{}XIIb-a:\allowbreak{}49}
\DERapparatusitem{ENV0749}{82;\allowbreak{}\allowbreak{} running\allowbreak{} text\allowbreak{} or\allowbreak{} formula}{\textbackslash{}\allowbreak{}mathscr\{\allowbreak{}X\}\allowbreak{}}{\textbackslash{}\allowbreak{}mathcal\{\allowbreak{}X\}\allowbreak{}}{TYPESETTING:\allowbreak{}1}
\DERapparatusitem{ENV0750}{82;\allowbreak{}\allowbreak{} running\allowbreak{} text\allowbreak{} or\allowbreak{} formula}{\textbackslash{}\allowbreak{}mathscr\{\allowbreak{}X\}\allowbreak{}}{\textbackslash{}\allowbreak{}mathcal\{\allowbreak{}X\}\allowbreak{}}{TYPESETTING:\allowbreak{}1}
\DERapparatusitem{ENV0751}{82;\allowbreak{}\allowbreak{} running\allowbreak{} text\allowbreak{} or\allowbreak{} formula}{\textbackslash{}\allowbreak{}mathscr\{\allowbreak{}X\}\allowbreak{}}{\textbackslash{}\allowbreak{}mathcal\{\allowbreak{}X\}\allowbreak{}}{TYPESETTING:\allowbreak{}1}
\DERapparatusitem{ENV0752}{82;\allowbreak{}\allowbreak{} running\allowbreak{} text\allowbreak{} or\allowbreak{} formula}{\textbackslash{}\allowbreak{}mathscr\{\allowbreak{}X\}\allowbreak{}}{\textbackslash{}\allowbreak{}mathcal\{\allowbreak{}X\}\allowbreak{}}{TYPESETTING:\allowbreak{}1}
\DERapparatusitem{ENV0753}{82;\allowbreak{}\allowbreak{} running\allowbreak{} text\allowbreak{} or\allowbreak{} formula}{\textbackslash{}\allowbreak{}mathscr\{\allowbreak{}X\}\allowbreak{}}{\textbackslash{}\allowbreak{}mathcal\{\allowbreak{}X\}\allowbreak{}}{TYPESETTING:\allowbreak{}1}
\DERapparatusitem{ENV0754}{83;\allowbreak{}\allowbreak{} running\allowbreak{} text\allowbreak{} or\allowbreak{} formula}{\textbackslash{}\allowbreak{}text\{\allowbreak{}def\}\allowbreak{}}{\textbackslash{}\allowbreak{}mathrm\{\allowbreak{}def\}\allowbreak{}}{TYPESETTING:\allowbreak{}7;\allowbreak{}\allowbreak{} NOTATION\_\allowbreak{}DER:\allowbreak{}17}
\DERapparatusitem{ENV0755}{85;\allowbreak{}\allowbreak{} ms126,\allowbreak{}\allowbreak{} parenthesis\allowbreak{} in\allowbreak{} cube\allowbreak{} corollary}{the\allowbreak{} diagram\allowbreak{} in\allowbreak{} solid\allowbreak{} lines\allowbreak{} opposite}{the\allowbreak{} diagram\allowbreak{} in\allowbreak{} solid\allowbreak{} lines\allowbreak{} below}{DER\_\allowbreak{}XIIb-b:\allowbreak{}3}
\DERapparatusitem{ENV0756}{86;\allowbreak{}\allowbreak{} ms127,\allowbreak{}\allowbreak{} replacement\allowbreak{} by\allowbreak{} closed\allowbreak{} pieces}{are\allowbreak{} replaced\allowbreak{} respectively\allowbreak{} not\allowbreak{} the\allowbreak{} closed\allowbreak{} subsets}{are\allowbreak{} replaced\allowbreak{} respectively\allowbreak{} by\allowbreak{} the\allowbreak{} closed\allowbreak{} subsets}{DER\_\allowbreak{}XIIb-b:\allowbreak{}4}
\DERapparatusitem{ENV0757}{86;\allowbreak{}\allowbreak{} ms127,\allowbreak{}\allowbreak{} description\allowbreak{} of\allowbreak{} replacement\allowbreak{} squares}{being\allowbreak{} inclusions\allowbreak{} of\allowbreak{} open\allowbreak{} subsets}{being\allowbreak{} inclusions\allowbreak{} of\allowbreak{} closed\allowbreak{} subsets}{DER\_\allowbreak{}XIIb-b:\allowbreak{}5}
\DERapparatusitem{ENV0758}{86;\allowbreak{}\allowbreak{} ms128,\allowbreak{}\allowbreak{} condition\allowbreak{} b\allowbreak{} prime}{f'\_\allowbreak{}1\allowbreak{} :\allowbreak{}\allowbreak{} \textbackslash{}\allowbreak{}cX/\allowbreak{}1\allowbreak{} \textbackslash{}\allowbreak{}longrightarrow\allowbreak{} \textbackslash{}\allowbreak{}cX'/\allowbreak{}1}{f'\_\allowbreak{}1\allowbreak{} :\allowbreak{}\allowbreak{} 1\textbackslash{}\allowbreak{}backslash\textbackslash{}\allowbreak{}cX\allowbreak{} \textbackslash{}\allowbreak{}longrightarrow\allowbreak{} 1\textbackslash{}\allowbreak{}backslash\textbackslash{}\allowbreak{}cX'}{DER\_\allowbreak{}XIIb-b:\allowbreak{}6}
\DERapparatusitem{ENV0759}{86;\allowbreak{}\allowbreak{} ms128,\allowbreak{}\allowbreak{} condition\allowbreak{} b\allowbreak{} prime}{f'\_\allowbreak{}2\allowbreak{} :\allowbreak{}\allowbreak{} \textbackslash{}\allowbreak{}cX/\allowbreak{}2\allowbreak{} \textbackslash{}\allowbreak{}longrightarrow\allowbreak{} \textbackslash{}\allowbreak{}cX'/\allowbreak{}2}{f'\_\allowbreak{}2\allowbreak{} :\allowbreak{}\allowbreak{} 2\textbackslash{}\allowbreak{}backslash\textbackslash{}\allowbreak{}cX\allowbreak{} \textbackslash{}\allowbreak{}longrightarrow\allowbreak{} 2\textbackslash{}\allowbreak{}backslash\textbackslash{}\allowbreak{}cX'}{DER\_\allowbreak{}XIIb-b:\allowbreak{}6}
\DERapparatusitem{ENV0760}{preamble;\allowbreak{}\allowbreak{} local\allowbreak{} macro\allowbreak{} definition}{\textbackslash{}\allowbreak{}mathscr\{\allowbreak{}X\}\allowbreak{}}{\textbackslash{}\allowbreak{}mathcal\{\allowbreak{}X\}\allowbreak{}}{TYPESETTING:\allowbreak{}1}
\DERapparatusitem{ENV0761}{preamble;\allowbreak{}\allowbreak{} local\allowbreak{} macro\allowbreak{} definition}{\textbackslash{}\allowbreak{}mathscr\{\allowbreak{}X\}\allowbreak{}}{\textbackslash{}\allowbreak{}mathcal\{\allowbreak{}X\}\allowbreak{}}{TYPESETTING:\allowbreak{}1}
\DERapparatusitem{ENV0762}{88;\allowbreak{}\allowbreak{} running\allowbreak{} text\allowbreak{} or\allowbreak{} formula}{\textbackslash{}\allowbreak{}text\{\allowbreak{}def\}\allowbreak{}}{\textbackslash{}\allowbreak{}mathrm\{\allowbreak{}def\}\allowbreak{}}{TYPESETTING:\allowbreak{}7;\allowbreak{}\allowbreak{} NOTATION\_\allowbreak{}DER:\allowbreak{}17}
\DERapparatusitem{ENV0763}{preamble;\allowbreak{}\allowbreak{} local\allowbreak{} macro\allowbreak{} definition}{\textbackslash{}\allowbreak{}mathscr\{\allowbreak{}X\}\allowbreak{}}{\textbackslash{}\allowbreak{}mathcal\{\allowbreak{}X\}\allowbreak{}}{TYPESETTING:\allowbreak{}1}
\DERapparatusitem{ENV0764}{preamble;\allowbreak{}\allowbreak{} local\allowbreak{} macro\allowbreak{} definition}{\textbackslash{}\allowbreak{}mathscr\{\allowbreak{}X\}\allowbreak{}}{\textbackslash{}\allowbreak{}mathcal\{\allowbreak{}X\}\allowbreak{}}{TYPESETTING:\allowbreak{}1}
\DERapparatusitem{ENV0765}{94;\allowbreak{}\allowbreak{} ms.\allowbreak{}139,\allowbreak{}\allowbreak{} parenthesis\allowbreak{} identifying\allowbreak{} Gamma;\allowbreak{}\allowbreak{} footnotes88–89}{\$\allowbreak{}\textbackslash{}\allowbreak{}boldsymbol\{\allowbreak{}\textbackslash{}\allowbreak{}Delta\}\allowbreak{}\textasciicircum{}\allowbreak{}1\textbackslash{}\allowbreak{}times\textbackslash{}\allowbreak{}boldsymbol\{\allowbreak{}\textbackslash{}\allowbreak{}Delta\}\allowbreak{}\textasciicircum{}\allowbreak{}1\$\allowbreak{}\allowbreak{} (\allowbreak{}\textbackslash{}\allowbreak{}footnotemark[\allowbreak{}88]\allowbreak{})\allowbreak{}}{the\allowbreak{} free\allowbreak{} diamond\allowbreak{} (\allowbreak{}\textbackslash{}\allowbreak{}footnotemark[\allowbreak{}88]\allowbreak{})\allowbreak{}}{DER\_\allowbreak{}XIIb-b:\allowbreak{}10}
\DERapparatusitem{ENV0766}{98;\allowbreak{}\allowbreak{} ms.\allowbreak{}142,\allowbreak{}\allowbreak{} paragraph\allowbreak{} beginning\allowbreak{} ce\allowbreak{} qui\allowbreak{} implique}{X\_\allowbreak{}*\allowbreak{} \textbackslash{}\allowbreak{}in\allowbreak{} \textbackslash{}\allowbreak{}mathrm\{\allowbreak{}Ob\}\allowbreak{}\textbackslash{}\allowbreak{},\allowbreak{}\textbackslash{}\allowbreak{}uHom(\allowbreak{}S,\allowbreak{}\textbackslash{}\allowbreak{}Cat)\allowbreak{}}{X\_\allowbreak{}*\allowbreak{} \textbackslash{}\allowbreak{}in\allowbreak{} \textbackslash{}\allowbreak{}mathrm\{\allowbreak{}Ob\}\allowbreak{}\textbackslash{}\allowbreak{},\allowbreak{}\textbackslash{}\allowbreak{}uHom(\allowbreak{}S\textasciicircum{}\allowbreak{}\textbackslash{}\allowbreak{}circ,\allowbreak{}\textbackslash{}\allowbreak{}Cat)\allowbreak{}}{DER\_\allowbreak{}XIIb-b:\allowbreak{}13}
\DERapparatusitem{ENV0767}{100;\allowbreak{}\allowbreak{} ms.\allowbreak{}146,\allowbreak{}\allowbreak{} paragraph\allowbreak{} completing\allowbreak{} the\allowbreak{} question\allowbreak{} of\allowbreak{} ms.\allowbreak{}145}{\$\allowbreak{}W\$\allowbreak{}-homotopically\allowbreak{} commutative\allowbreak{} \textbackslash{}\allowbreak{}texttt\{\allowbreak{}[\allowbreak{}rather\allowbreak{} cocartesian]\allowbreak{}\}\allowbreak{}}{\$\allowbreak{}W\$\allowbreak{}-homotopically\allowbreak{} cocartesian\allowbreak{} \textbackslash{}\allowbreak{}texttt\{\allowbreak{}[\allowbreak{}rather\allowbreak{} cocartesian]\allowbreak{}\}\allowbreak{}}{DER\_\allowbreak{}XIIb-b:\allowbreak{}15}
\DERapparatusitem{ENV0768}{101;\allowbreak{}\allowbreak{} ms.\allowbreak{}147,\allowbreak{}\allowbreak{} Corollary1,\allowbreak{}\allowbreak{} conclusion}{\$\allowbreak{}W\$\allowbreak{}-homotopically\allowbreak{} cartesian\}\allowbreak{}}{\$\allowbreak{}W\$\allowbreak{}-homotopically\allowbreak{} cocartesian\}\allowbreak{}}{DER\_\allowbreak{}XIIb-b:\allowbreak{}16}
\DERapparatusitem{ENV0769}{101;\allowbreak{}\allowbreak{} ms.\allowbreak{}147,\allowbreak{}\allowbreak{} Corollary1,\allowbreak{}\allowbreak{} parenthetical\allowbreak{} duality\allowbreak{} argument}{\$\allowbreak{}W\$\allowbreak{}-homotopically\allowbreak{} trivial\}\allowbreak{}}{\$\allowbreak{}W\$\allowbreak{}-homotopically\allowbreak{} cocartesian\}\allowbreak{}}{DER\_\allowbreak{}XIIb-b:\allowbreak{}17}
\DERapparatusitem{ENV0770}{103;\allowbreak{}\allowbreak{} running\allowbreak{} text\allowbreak{} or\allowbreak{} formula}{\textbackslash{}\allowbreak{}text\{\allowbreak{}def\}\allowbreak{}}{\textbackslash{}\allowbreak{}mathrm\{\allowbreak{}def\}\allowbreak{}}{TYPESETTING:\allowbreak{}7;\allowbreak{}\allowbreak{} NOTATION\_\allowbreak{}DER:\allowbreak{}17}
\DERapparatusitem{ENV0771}{103;\allowbreak{}\allowbreak{} running\allowbreak{} text\allowbreak{} or\allowbreak{} formula}{\textbackslash{}\allowbreak{}text\{\allowbreak{}def\}\allowbreak{}}{\textbackslash{}\allowbreak{}mathrm\{\allowbreak{}def\}\allowbreak{}}{TYPESETTING:\allowbreak{}7;\allowbreak{}\allowbreak{} NOTATION\_\allowbreak{}DER:\allowbreak{}17}
\DERapparatusitem{ENV0772}{107;\allowbreak{}\allowbreak{} Printed107,\allowbreak{}\allowbreak{} ms155,\allowbreak{}\allowbreak{} cartesian-arrow\allowbreak{} display}{x\_\allowbreak{}0\textbackslash{}\allowbreak{}longrightarrow\textbackslash{}\allowbreak{}alpha\_\allowbreak{}1(\allowbreak{}x\_\allowbreak{}0)\allowbreak{},\allowbreak{}\allowbreak{} \textbackslash{}\allowbreak{}qquad\allowbreak{} x\_\allowbreak{}0\textbackslash{}\allowbreak{}longrightarrow\textbackslash{}\allowbreak{}alpha\_\allowbreak{}2(\allowbreak{}x\_\allowbreak{}0)\allowbreak{}}{\textbackslash{}\allowbreak{}alpha\_\allowbreak{}1(\allowbreak{}x\_\allowbreak{}0)\allowbreak{}\textbackslash{}\allowbreak{}longrightarrow\allowbreak{} x\_\allowbreak{}0,\allowbreak{}\allowbreak{} \textbackslash{}\allowbreak{}qquad\allowbreak{} \textbackslash{}\allowbreak{}alpha\_\allowbreak{}2(\allowbreak{}x\_\allowbreak{}0)\allowbreak{}\textbackslash{}\allowbreak{}longrightarrow\allowbreak{} x\_\allowbreak{}0}{DER\_\allowbreak{}XIIb-b:\allowbreak{}18}
\DERapparatusitem{ENV0773}{107;\allowbreak{}\allowbreak{} running\allowbreak{} text\allowbreak{} or\allowbreak{} formula}{\textbackslash{}\allowbreak{}text\{\allowbreak{}def\}\allowbreak{}}{\textbackslash{}\allowbreak{}mathrm\{\allowbreak{}def\}\allowbreak{}}{TYPESETTING:\allowbreak{}7;\allowbreak{}\allowbreak{} NOTATION\_\allowbreak{}DER:\allowbreak{}17}
\DERapparatusitem{ENV0774}{108;\allowbreak{}\allowbreak{} Printed108,\allowbreak{}\allowbreak{} ms157,\allowbreak{}\allowbreak{} sentence\allowbreak{} before\allowbreak{} Corollary6}{is\allowbreak{} smooth\allowbreak{} at\allowbreak{} \$\allowbreak{}u\$\allowbreak{}\allowbreak{} \textbackslash{}\allowbreak{}texttt\{\allowbreak{}[\allowbreak{}rather\allowbreak{} at\allowbreak{} \$\allowbreak{}x\_\allowbreak{}0\$\allowbreak{}]\allowbreak{}\}\allowbreak{}}{is\allowbreak{} smooth\allowbreak{} at\allowbreak{} \$\allowbreak{}x\_\allowbreak{}0\$\allowbreak{}\allowbreak{} \textbackslash{}\allowbreak{}texttt\{\allowbreak{}[\allowbreak{}rather\allowbreak{} at\allowbreak{} \$\allowbreak{}x\_\allowbreak{}0\$\allowbreak{}]\allowbreak{}\}\allowbreak{}}{DER\_\allowbreak{}XIIb-b:\allowbreak{}21}
\DERapparatusitem{ENV0775}{108;\allowbreak{}\allowbreak{} Printed108,\allowbreak{}\allowbreak{} ms157,\allowbreak{}\allowbreak{} final\allowbreak{} sentence\allowbreak{} after\allowbreak{} Corollary6}{(\allowbreak{}\$\allowbreak{}i\allowbreak{} \textbackslash{}\allowbreak{}in\allowbreak{} \textbackslash{}\allowbreak{}\{\allowbreak{}0,\allowbreak{}1,\allowbreak{}2\textbackslash{}\allowbreak{}\}\allowbreak{}\$\allowbreak{})\allowbreak{}\allowbreak{} we\allowbreak{} find}{(\allowbreak{}\$\allowbreak{}i\allowbreak{} \textbackslash{}\allowbreak{}in\allowbreak{} \textbackslash{}\allowbreak{}\{\allowbreak{}0,\allowbreak{}1,\allowbreak{}2\textbackslash{}\allowbreak{}\}\allowbreak{}\$\allowbreak{})\allowbreak{},\allowbreak{}\allowbreak{} under\allowbreak{} the\allowbreak{} corresponding\allowbreak{} hypotheses\allowbreak{} of\allowbreak{} a)\allowbreak{}\allowbreak{} and\allowbreak{} b)\allowbreak{},\allowbreak{}\allowbreak{} we\allowbreak{} find}{DER\_\allowbreak{}XIIb-b:\allowbreak{}22}
\DERapparatusitem{ENV0776}{109;\allowbreak{}\allowbreak{} Printed109,\allowbreak{}\allowbreak{} ms159,\allowbreak{}\allowbreak{} Corollary1\allowbreak{} complement\allowbreak{} hypothesis}{suppose\allowbreak{} that\allowbreak{} \$\allowbreak{}T'\allowbreak{} \textbackslash{}\allowbreak{}longrightarrow\allowbreak{} T\$\allowbreak{}}{suppose\allowbreak{} that\allowbreak{} \$\allowbreak{}T\allowbreak{} \textbackslash{}\allowbreak{}longrightarrow\allowbreak{} T'\$\allowbreak{}}{DER\_\allowbreak{}XIIb-b:\allowbreak{}23}
\DERapparatusitem{ENV0777}{110;\allowbreak{}\allowbreak{} Printed110,\allowbreak{}\allowbreak{} ms161,\allowbreak{}\allowbreak{} Corollary2\allowbreak{} complement\allowbreak{} hypothesis}{suppose\allowbreak{} \$\allowbreak{}T'\allowbreak{} \textbackslash{}\allowbreak{}longrightarrow\allowbreak{} T\$\allowbreak{}\allowbreak{} an\allowbreak{} isomorphism}{suppose\allowbreak{} \$\allowbreak{}T\allowbreak{} \textbackslash{}\allowbreak{}longrightarrow\allowbreak{} T'\$\allowbreak{}\allowbreak{} an\allowbreak{} isomorphism}{DER\_\allowbreak{}XIIb-b:\allowbreak{}25}
\DERapparatusitem{ENV0778}{112;\allowbreak{}\allowbreak{} Printed112,\allowbreak{}\allowbreak{} ms166,\allowbreak{}\allowbreak{} last\allowbreak{} paragraph}{which\allowbreak{} do\allowbreak{} not\allowbreak{} factor\allowbreak{} through\allowbreak{} \$\allowbreak{}G\$\allowbreak{}}{which\allowbreak{} do\allowbreak{} not\allowbreak{} factor\allowbreak{} through\allowbreak{} \$\allowbreak{}F\$\allowbreak{}}{DER\_\allowbreak{}XIIb-b:\allowbreak{}29}
\DERapparatusitem{ENV0779}{114;\allowbreak{}\allowbreak{} Printed114,\allowbreak{}\allowbreak{} ms169,\allowbreak{}\allowbreak{} final\allowbreak{} cross-reference}{(\allowbreak{}pages\allowbreak{} 99,\allowbreak{}\allowbreak{} 100\allowbreak{} \textbackslash{}\allowbreak{}texttt\{\allowbreak{}[\allowbreak{}rather\allowbreak{} 109,\allowbreak{}\allowbreak{} 110]\allowbreak{}\}\allowbreak{})\allowbreak{}}{(\allowbreak{}pages\allowbreak{} 109,\allowbreak{}\allowbreak{} 110\allowbreak{} \textbackslash{}\allowbreak{}texttt\{\allowbreak{}[\allowbreak{}rather\allowbreak{} 109,\allowbreak{}\allowbreak{} 110]\allowbreak{}\}\allowbreak{})\allowbreak{}}{DER\_\allowbreak{}XIIb-b:\allowbreak{}30}
\subsection*{Chapter XII, 6}
\DERapparatusitem{ENV0780}{115;\allowbreak{}\allowbreak{} ms171,\allowbreak{}\allowbreak{} NB\allowbreak{} introducing\allowbreak{} the\allowbreak{} constant\allowbreak{} presheaf}{the\allowbreak{} constant\allowbreak{} sheaf\allowbreak{} of\allowbreak{} sets\allowbreak{} defined\allowbreak{} by\allowbreak{} \$\allowbreak{}F\$\allowbreak{}\allowbreak{} on\allowbreak{} \$\allowbreak{}Y\$\allowbreak{}}{the\allowbreak{} constant\allowbreak{} sheaf\allowbreak{} of\allowbreak{} sets\allowbreak{} defined\allowbreak{} by\allowbreak{} \$\allowbreak{}F\_\allowbreak{}0\$\allowbreak{}\allowbreak{} on\allowbreak{} \$\allowbreak{}Y\$\allowbreak{}}{DER\_\allowbreak{}XIIc-a:\allowbreak{}1}
\DERapparatusitem{ENV0781}{119;\allowbreak{}\allowbreak{} ms177\allowbreak{} continuation,\allowbreak{}\allowbreak{} arbitrary\allowbreak{} loop\allowbreak{} basepoint}{n}{x}{DER\_\allowbreak{}XIIc-a:\allowbreak{}4}
\DERapparatusitem{ENV0782}{119;\allowbreak{}\allowbreak{} ms179,\allowbreak{}\allowbreak{} definition\allowbreak{} of\allowbreak{} c\_\allowbreak{}n}{isomorphism\allowbreak{} for\allowbreak{} \$\allowbreak{}F\$\allowbreak{}\allowbreak{} a\allowbreak{} locally\allowbreak{} constant\allowbreak{} sheaf\allowbreak{} on\allowbreak{} \$\allowbreak{}Y\$\allowbreak{}\allowbreak{} (\allowbreak{}with\allowbreak{} \$\allowbreak{}F\$\allowbreak{}\allowbreak{} abelian\allowbreak{} if\allowbreak{} \$\allowbreak{}i\allowbreak{} \textbackslash{}\allowbreak{}geq\allowbreak{} 2\$\allowbreak{},\allowbreak{}\allowbreak{} \$\allowbreak{}F\$\allowbreak{}\allowbreak{} a\allowbreak{} sheaf\allowbreak{} of\allowbreak{} groups\allowbreak{} if\allowbreak{} \$\allowbreak{}i\allowbreak{} =\allowbreak{}\allowbreak{} 1\$\allowbreak{})\allowbreak{}}{isomorphism\allowbreak{} for\allowbreak{} \$\allowbreak{}0\textbackslash{}\allowbreak{}leq\allowbreak{} i\textbackslash{}\allowbreak{}leq\allowbreak{} n\$\allowbreak{}\allowbreak{} and\allowbreak{} \$\allowbreak{}F\$\allowbreak{}\allowbreak{} a\allowbreak{} locally\allowbreak{} constant\allowbreak{} sheaf\allowbreak{} on\allowbreak{} \$\allowbreak{}Y\$\allowbreak{}\allowbreak{} (\allowbreak{}with\allowbreak{} \$\allowbreak{}F\$\allowbreak{}\allowbreak{} abelian\allowbreak{} if\allowbreak{} \$\allowbreak{}i\allowbreak{} \textbackslash{}\allowbreak{}geq\allowbreak{} 2\$\allowbreak{},\allowbreak{}\allowbreak{} \$\allowbreak{}F\$\allowbreak{}\allowbreak{} a\allowbreak{} sheaf\allowbreak{} of\allowbreak{} groups\allowbreak{} if\allowbreak{} \$\allowbreak{}i\allowbreak{} =\allowbreak{}\allowbreak{} 1\$\allowbreak{},\allowbreak{}\allowbreak{} \$\allowbreak{}F\$\allowbreak{}\allowbreak{} a\allowbreak{} sheaf\allowbreak{} of\allowbreak{} sets\allowbreak{} if\allowbreak{} \$\allowbreak{}i=\allowbreak{}0\$\allowbreak{})\allowbreak{}}{DER\_\allowbreak{}XIIc-a:\allowbreak{}6}
\DERapparatusitem{ENV0783}{121;\allowbreak{}\allowbreak{} running\allowbreak{} text\allowbreak{} or\allowbreak{} formula}{\textbackslash{}\allowbreak{}text\{\allowbreak{}def\}\allowbreak{}}{\textbackslash{}\allowbreak{}mathrm\{\allowbreak{}def\}\allowbreak{}}{TYPESETTING:\allowbreak{}7;\allowbreak{}\allowbreak{} NOTATION\_\allowbreak{}DER:\allowbreak{}17}
\DERapparatusitem{ENV0784}{124;\allowbreak{}\allowbreak{} ms187,\allowbreak{}\allowbreak{} final\allowbreak{} sentence\allowbreak{} before\allowbreak{} section2}{is\allowbreak{} defined\allowbreak{} by\allowbreak{} the\allowbreak{} condition\allowbreak{} \$\allowbreak{}\textbackslash{}\allowbreak{}mathrm\{\allowbreak{}a\}\allowbreak{}\_\allowbreak{}n\$\allowbreak{}}{is\allowbreak{} defined\allowbreak{} by\allowbreak{} the\allowbreak{} condition\allowbreak{} \$\allowbreak{}\textbackslash{}\allowbreak{}mathrm\{\allowbreak{}b\}\allowbreak{}\_\allowbreak{}n\$\allowbreak{}}{DER\_\allowbreak{}XIIc-a:\allowbreak{}9}
\DERapparatusitem{ENV0785}{125;\allowbreak{}\allowbreak{} ms188,\allowbreak{}\allowbreak{} proof\allowbreak{} of\allowbreak{} Theorem2,\allowbreak{}\allowbreak{} initial\allowbreak{} backward\allowbreak{} page\allowbreak{} reference}{at\allowbreak{} the\allowbreak{} beginning\allowbreak{} (\allowbreak{}p.\allowbreak{}\textasciitilde{}71)\allowbreak{}}{at\allowbreak{} the\allowbreak{} beginning\allowbreak{} (\allowbreak{}p.\allowbreak{}\textasciitilde{}171)\allowbreak{}}{DER\_\allowbreak{}XIIc-a:\allowbreak{}10}
\DERapparatusitem{ENV0786}{129;\allowbreak{}\allowbreak{} ms197,\allowbreak{}\allowbreak{} prose\allowbreak{} below\allowbreak{} the\allowbreak{} diamond;\allowbreak{}\allowbreak{} existing\allowbreak{} sic\allowbreak{} flag}{}{'}{DER\_\allowbreak{}XIIc-a:\allowbreak{}16}
\DERapparatusitem{ENV0787}{131;\allowbreak{}\allowbreak{} ms200,\allowbreak{}\allowbreak{} both\allowbreak{} displays\allowbreak{} for\allowbreak{} passage\allowbreak{} to\allowbreak{} the\allowbreak{} ordinary\allowbreak{} local-system\allowbreak{} functor}{\textbackslash{}\allowbreak{}HOT\_\allowbreak{}W(\allowbreak{}S)\allowbreak{}\allowbreak{} \textbackslash{}\allowbreak{}arrow[\allowbreak{}rr]\allowbreak{}}{\textbackslash{}\allowbreak{}HOT\textasciicircum{}\allowbreak{}\{\allowbreak{}\textbackslash{}\allowbreak{}mathrm\{\allowbreak{}lc\}\allowbreak{}\}\allowbreak{}\_\allowbreak{}W(\allowbreak{}S)\allowbreak{}\allowbreak{} \textbackslash{}\allowbreak{}arrow[\allowbreak{}rr]\allowbreak{}}{DER\_\allowbreak{}XIIc-a:\allowbreak{}19}
\DERapparatusitem{ENV0788}{131;\allowbreak{}\allowbreak{} ms200,\allowbreak{}\allowbreak{} both\allowbreak{} displays\allowbreak{} for\allowbreak{} passage\allowbreak{} to\allowbreak{} the\allowbreak{} ordinary\allowbreak{} local-system\allowbreak{} functor}{\textbackslash{}\allowbreak{}Hom(\allowbreak{}\textbackslash{}\allowbreak{}Cat/\allowbreak{}S,\allowbreak{}\allowbreak{} \textbackslash{}\allowbreak{}HOT\_\allowbreak{}W(\allowbreak{}S)\allowbreak{})\allowbreak{}}{\textbackslash{}\allowbreak{}Hom(\allowbreak{}\textbackslash{}\allowbreak{}Cat/\allowbreak{}S,\allowbreak{}\allowbreak{} \textbackslash{}\allowbreak{}HOT\textasciicircum{}\allowbreak{}\{\allowbreak{}\textbackslash{}\allowbreak{}mathrm\{\allowbreak{}lc\}\allowbreak{}\}\allowbreak{}\_\allowbreak{}W(\allowbreak{}S)\allowbreak{})\allowbreak{}}{DER\_\allowbreak{}XIIc-a:\allowbreak{}19}
\DERapparatusitem{ENV0789}{132;\allowbreak{}\allowbreak{} ms201,\allowbreak{}\allowbreak{} first\allowbreak{} pair\allowbreak{} in\allowbreak{} common-refinement\allowbreak{} sentence;\allowbreak{}\allowbreak{} existing\allowbreak{} sic\allowbreak{} flag}{}{'}{DER\_\allowbreak{}XIIc-a:\allowbreak{}20}
\DERapparatusitem{ENV0790}{133;\allowbreak{}\allowbreak{} running\allowbreak{} text\allowbreak{} or\allowbreak{} formula}{\textbackslash{}\allowbreak{}text\{\allowbreak{}def\}\allowbreak{}}{\textbackslash{}\allowbreak{}mathrm\{\allowbreak{}def\}\allowbreak{}}{TYPESETTING:\allowbreak{}7;\allowbreak{}\allowbreak{} NOTATION\_\allowbreak{}DER:\allowbreak{}17}
\DERapparatusitem{ENV0791}{133;\allowbreak{}\allowbreak{} ms205,\allowbreak{}\allowbreak{} opening\allowbreak{} proof\allowbreak{} sentence}{The\allowbreak{} fact\allowbreak{} that\allowbreak{} \$\allowbreak{}(\allowbreak{}*)\allowbreak{}\$\allowbreak{}\allowbreak{} is\allowbreak{} an\allowbreak{} equivalence}{The\allowbreak{} fact\allowbreak{} that\allowbreak{} \$\allowbreak{}(\allowbreak{}**)\allowbreak{}\$\allowbreak{}\allowbreak{} is\allowbreak{} an\allowbreak{} equivalence}{DER\_\allowbreak{}XIIc-a:\allowbreak{}21}
\DERapparatusitem{ENV0792}{133;\allowbreak{}\allowbreak{} ms205,\allowbreak{}\allowbreak{} nonemptiness\allowbreak{} phrase;\allowbreak{}\allowbreak{} existing\allowbreak{} sic\allowbreak{} flag}{\$\allowbreak{}C\allowbreak{} \textbackslash{}\allowbreak{}neq\allowbreak{} \textbackslash{}\allowbreak{}emptyset\$\allowbreak{}}{\$\allowbreak{}\textbackslash{}\allowbreak{}mathcal\{\allowbreak{}C\}\allowbreak{}\allowbreak{} \textbackslash{}\allowbreak{}neq\allowbreak{} \textbackslash{}\allowbreak{}emptyset\$\allowbreak{}}{DER\_\allowbreak{}XIIc-a:\allowbreak{}22}
\DERapparatusitem{ENV0793}{134;\allowbreak{}\allowbreak{} running\allowbreak{} text\allowbreak{} or\allowbreak{} formula}{\allowbreak{} (\allowbreak{}gear)\allowbreak{}}{}{GLOSSARY\_\allowbreak{}DER:\allowbreak{}164}
\DERapparatusitem{ENV0794}{135;\allowbreak{}\allowbreak{} ms208,\allowbreak{}\allowbreak{} continuation\allowbreak{} of\allowbreak{} W-fibration\allowbreak{} definition}{takes\allowbreak{} \$\allowbreak{}W\_\allowbreak{}X\$\allowbreak{}\allowbreak{} into\allowbreak{} \$\allowbreak{}W\_\allowbreak{}Y\$\allowbreak{}}{takes\allowbreak{} \$\allowbreak{}W\_\allowbreak{}Y\$\allowbreak{}\allowbreak{} into\allowbreak{} \$\allowbreak{}W\_\allowbreak{}X\$\allowbreak{}}{DER\_\allowbreak{}XIIc-a:\allowbreak{}26}
\DERapparatusitem{ENV0795}{135;\allowbreak{}\allowbreak{} ms208,\allowbreak{}\allowbreak{} definition\allowbreak{} of\allowbreak{} W\_\allowbreak{}X;\allowbreak{}\allowbreak{} existing\allowbreak{} sic\allowbreak{} flag}{U}{u}{DER\_\allowbreak{}XIIc-a:\allowbreak{}27}
\DERapparatusitem{ENV0796}{136;\allowbreak{}\allowbreak{} running\allowbreak{} text\allowbreak{} or\allowbreak{} formula}{\textbackslash{}\allowbreak{}text\{\allowbreak{}def\}\allowbreak{}}{\textbackslash{}\allowbreak{}mathrm\{\allowbreak{}def\}\allowbreak{}}{TYPESETTING:\allowbreak{}7;\allowbreak{}\allowbreak{} NOTATION\_\allowbreak{}DER:\allowbreak{}17}
\DERapparatusitem{ENV0797}{136;\allowbreak{}\allowbreak{} running\allowbreak{} text\allowbreak{} or\allowbreak{} formula}{\textbackslash{}\allowbreak{}text\{\allowbreak{}def\}\allowbreak{}}{\textbackslash{}\allowbreak{}mathrm\{\allowbreak{}def\}\allowbreak{}}{TYPESETTING:\allowbreak{}7;\allowbreak{}\allowbreak{} NOTATION\_\allowbreak{}DER:\allowbreak{}17}
\DERapparatusitem{ENV0798}{136;\allowbreak{}\allowbreak{} ms211,\allowbreak{}\allowbreak{} after\allowbreak{} the\allowbreak{} displayed\allowbreak{} relative\allowbreak{} pullback\allowbreak{} functor}{and\allowbreak{} thus\allowbreak{} passes\allowbreak{} to\allowbreak{} the\allowbreak{} categories\allowbreak{} of\allowbreak{} fractions\allowbreak{} as\allowbreak{} an\allowbreak{} arrow\allowbreak{} \$\allowbreak{}(\allowbreak{}*)\allowbreak{}\$\allowbreak{}}{and\allowbreak{} thus\allowbreak{} passes\allowbreak{} to\allowbreak{} the\allowbreak{} categories\allowbreak{} of\allowbreak{} fractions\allowbreak{} as\allowbreak{} an\allowbreak{} arrow\allowbreak{} \$\allowbreak{}(\allowbreak{}**)\allowbreak{}\$\allowbreak{}}{DER\_\allowbreak{}XIIc-a:\allowbreak{}29}
\DERapparatusitem{ENV0799}{137;\allowbreak{}\allowbreak{} ms213,\allowbreak{}\allowbreak{} item\allowbreak{} 1°}{for\allowbreak{} arbitrary\allowbreak{} arrows\allowbreak{} \$\allowbreak{}S'\allowbreak{} \textbackslash{}\allowbreak{}longrightarrow\allowbreak{} S\$\allowbreak{}\allowbreak{} in\allowbreak{} Cat}{for\allowbreak{} arbitrary\allowbreak{} arrows\allowbreak{} \$\allowbreak{}S'\allowbreak{} \textbackslash{}\allowbreak{}longrightarrow\allowbreak{} S\$\allowbreak{}\allowbreak{} in\allowbreak{} \$\allowbreak{}\textbackslash{}\allowbreak{}cM\$\allowbreak{}}{DER\_\allowbreak{}XIIc-a:\allowbreak{}30}
\DERapparatusitem{ENV0800}{139;\allowbreak{}\allowbreak{} running\allowbreak{} text\allowbreak{} or\allowbreak{} formula}{\textbackslash{}\allowbreak{}text\{\allowbreak{}def\}\allowbreak{}}{\textbackslash{}\allowbreak{}mathrm\{\allowbreak{}def\}\allowbreak{}}{TYPESETTING:\allowbreak{}7;\allowbreak{}\allowbreak{} NOTATION\_\allowbreak{}DER:\allowbreak{}17}
\DERapparatusitem{ENV0801}{139;\allowbreak{}\allowbreak{} ms217,\allowbreak{}\allowbreak{} underbrace\allowbreak{} of\allowbreak{} slice\allowbreak{} localization;\allowbreak{}\allowbreak{} existing\allowbreak{} sic\allowbreak{} flag}{\textbackslash{}\allowbreak{}mathrm\{\allowbreak{}Hot\}\allowbreak{}\textasciicircum{}\allowbreak{}!\_\allowbreak{}W(\allowbreak{}S)\allowbreak{}}{\textbackslash{}\allowbreak{}mathrm\{\allowbreak{}Ho\}\allowbreak{}\textasciicircum{}\allowbreak{}!\_\allowbreak{}W(\allowbreak{}S)\allowbreak{}}{DER\_\allowbreak{}XIIc-a:\allowbreak{}32}
\DERapparatusitem{ENV0802}{preamble;\allowbreak{}\allowbreak{} local\allowbreak{} macro\allowbreak{} definition}{\textbackslash{}\allowbreak{}text\{\allowbreak{}def\}\allowbreak{}}{\textbackslash{}\allowbreak{}mathrm\{\allowbreak{}def\}\allowbreak{}}{TYPESETTING:\allowbreak{}7;\allowbreak{}\allowbreak{} NOTATION\_\allowbreak{}DER:\allowbreak{}17}
\DERapparatusitem{ENV0803}{141;\allowbreak{}\allowbreak{} ms222,\allowbreak{}\allowbreak{} first\allowbreak{} definition\allowbreak{} of\allowbreak{} saturation}{f\allowbreak{} :\allowbreak{}\allowbreak{} S\allowbreak{} \textbackslash{}\allowbreak{}longrightarrow\allowbreak{} S'}{f\allowbreak{} :\allowbreak{}\allowbreak{} S'\allowbreak{} \textbackslash{}\allowbreak{}longrightarrow\allowbreak{} S}{DER\_\allowbreak{}XIIc-a:\allowbreak{}35}
\DERapparatusitem{ENV0804}{142;\allowbreak{}\allowbreak{} ms223,\allowbreak{}\allowbreak{} sentence\allowbreak{} before\allowbreak{} full\allowbreak{} subcategory\allowbreak{} Prop\_\allowbreak{}W\allowbreak{} S}{whereas\allowbreak{} the\allowbreak{} arrows\allowbreak{} of\allowbreak{} \$\allowbreak{}\textbackslash{}\allowbreak{}mathrm\{\allowbreak{}Liss\}\allowbreak{}\_\allowbreak{}W\$\allowbreak{}}{whereas\allowbreak{} the\allowbreak{} arrows\allowbreak{} of\allowbreak{} \$\allowbreak{}\textbackslash{}\allowbreak{}mathrm\{\allowbreak{}Prop\}\allowbreak{}\_\allowbreak{}W\$\allowbreak{}}{DER\_\allowbreak{}XIIc-a:\allowbreak{}36}
\DERapparatusitem{ENV0805}{142;\allowbreak{}\allowbreak{} ms223,\allowbreak{}\allowbreak{} same\allowbreak{} underbrace:\allowbreak{}\allowbreak{} mathematical\allowbreak{} interpretation}{0}{\textbackslash{}\allowbreak{}circ}{DER\_\allowbreak{}XIIc-a:\allowbreak{}38}
\DERapparatusitem{ENV0806}{143;\allowbreak{}\allowbreak{} ms226,\allowbreak{}\allowbreak{} two\allowbreak{} displayed\allowbreak{} definitions\allowbreak{} of\allowbreak{} Wg\allowbreak{} and\allowbreak{} Wd}{\textbackslash{}\allowbreak{}mathcal\{\allowbreak{}F\}\allowbreak{}(\allowbreak{}\textbackslash{}\allowbreak{}mathcal\{\allowbreak{}M\}\allowbreak{}/\allowbreak{}S)\allowbreak{}}{\textbackslash{}\allowbreak{}mathrm\{\allowbreak{}Fl\}\allowbreak{}(\allowbreak{}\textbackslash{}\allowbreak{}mathcal\{\allowbreak{}M\}\allowbreak{}/\allowbreak{}S)\allowbreak{}}{DER\_\allowbreak{}XIIc-a:\allowbreak{}39}
\DERapparatusitem{ENV0807}{143;\allowbreak{}\allowbreak{} ms226,\allowbreak{}\allowbreak{} two\allowbreak{} displayed\allowbreak{} definitions\allowbreak{} of\allowbreak{} Wg\allowbreak{} and\allowbreak{} Wd}{\textbackslash{}\allowbreak{}mathcal\{\allowbreak{}F\}\allowbreak{}(\allowbreak{}\textbackslash{}\allowbreak{}mathcal\{\allowbreak{}M\}\allowbreak{}/\allowbreak{}S)\allowbreak{}}{\textbackslash{}\allowbreak{}mathrm\{\allowbreak{}Fl\}\allowbreak{}(\allowbreak{}\textbackslash{}\allowbreak{}mathcal\{\allowbreak{}M\}\allowbreak{}/\allowbreak{}S)\allowbreak{}}{DER\_\allowbreak{}XIIc-a:\allowbreak{}39}
\DERapparatusitem{ENV0808}{147;\allowbreak{}\allowbreak{} ms230,\allowbreak{}\allowbreak{} proof\allowbreak{} after\allowbreak{} equation\allowbreak{} for\allowbreak{} i\_\allowbreak{}f}{as\allowbreak{} a\allowbreak{} composite\allowbreak{} of\allowbreak{} \$\allowbreak{}Y\allowbreak{} \textbackslash{}\allowbreak{}longrightarrow\allowbreak{} Y\allowbreak{} \textbackslash{}\allowbreak{}amalg\allowbreak{} X\$\allowbreak{}}{as\allowbreak{} a\allowbreak{} composite\allowbreak{} of\allowbreak{} \$\allowbreak{}X\allowbreak{} \textbackslash{}\allowbreak{}longrightarrow\allowbreak{} Y\allowbreak{} \textbackslash{}\allowbreak{}amalg\allowbreak{} X\$\allowbreak{}}{DER\_\allowbreak{}XIIc-a:\allowbreak{}43}
\DERapparatusitem{ENV0809}{148;\allowbreak{}\allowbreak{} ms232,\allowbreak{}\allowbreak{} first\allowbreak{} paragraph;\allowbreak{}\allowbreak{} existing\allowbreak{} [\allowbreak{}sic]\allowbreak{}}{}{'}{DER\_\allowbreak{}XIIc-a:\allowbreak{}44}
\DERapparatusitem{ENV0810}{148;\allowbreak{}\allowbreak{} ms232,\allowbreak{}\allowbreak{} cobase-change\allowbreak{} question;\allowbreak{}\allowbreak{} existing\allowbreak{} [\allowbreak{}sic]\allowbreak{}}{'\allowbreak{} \textbackslash{}\allowbreak{}longrightarrow\allowbreak{} S}{\allowbreak{} \textbackslash{}\allowbreak{}longrightarrow\allowbreak{} S'}{DER\_\allowbreak{}XIIc-a:\allowbreak{}45}
\DERapparatusitem{ENV0811}{preamble;\allowbreak{}\allowbreak{} local\allowbreak{} macro\allowbreak{} definition}{\textbackslash{}\allowbreak{}mathscr\{\allowbreak{}X\}\allowbreak{}}{\textbackslash{}\allowbreak{}mathcal\{\allowbreak{}X\}\allowbreak{}}{TYPESETTING:\allowbreak{}1}
\DERapparatusitem{ENV0812}{152;\allowbreak{}\allowbreak{} ms.\allowbreak{}236,\allowbreak{}\allowbreak{} first\allowbreak{} paragraph\allowbreak{} following\allowbreak{} definition\allowbreak{} of\allowbreak{} I}{\textbackslash{}\allowbreak{}cX}{I}{DER\_\allowbreak{}XIIc-b:\allowbreak{}2}
\DERapparatusitem{ENV0813}{152;\allowbreak{}\allowbreak{} ms.\allowbreak{}237,\allowbreak{}\allowbreak{} third\allowbreak{} arrow\allowbreak{} in\allowbreak{} the\allowbreak{} three-row\allowbreak{} display}{\textbackslash{}\allowbreak{}cX\_\allowbreak{}\{\allowbreak{}020'\}\allowbreak{}\allowbreak{} \textbackslash{}\allowbreak{}quad\allowbreak{} \&\allowbreak{}\textbackslash{}\allowbreak{}xrightarrow\{\allowbreak{}\textbackslash{}\allowbreak{}varphi\_\allowbreak{}1\}\allowbreak{}}{\textbackslash{}\allowbreak{}cX\_\allowbreak{}\{\allowbreak{}020'\}\allowbreak{}\allowbreak{} \textbackslash{}\allowbreak{}quad\allowbreak{} \&\allowbreak{}\textbackslash{}\allowbreak{}xrightarrow\{\allowbreak{}\textbackslash{}\allowbreak{}varphi\_\allowbreak{}3\}\allowbreak{}}{DER\_\allowbreak{}XIIc-b:\allowbreak{}3}
\DERapparatusitem{ENV0814}{153;\allowbreak{}\allowbreak{} ms.\allowbreak{}238,\allowbreak{}\allowbreak{} third\allowbreak{} slice\allowbreak{} row}{\textbackslash{}\allowbreak{}Phi}{\textbackslash{}\allowbreak{}Delta\textasciicircum{}\allowbreak{}2}{DER\_\allowbreak{}XIIc-b:\allowbreak{}5}
\DERapparatusitem{ENV0815}{155;\allowbreak{}\allowbreak{} ms240,\allowbreak{}\allowbreak{} footnote129}{is\allowbreak{} a\allowbreak{} category\allowbreak{} with\allowbreak{} cofibrations\allowbreak{} of\allowbreak{} K.\allowbreak{}S.\allowbreak{}}{is\allowbreak{} a\allowbreak{} category\allowbreak{} with\allowbreak{} fibrations\allowbreak{} of\allowbreak{} K.\allowbreak{}S.\allowbreak{}}{DER\_\allowbreak{}XIIc-b:\allowbreak{}6}
\DERapparatusitem{ENV0816}{156;\allowbreak{}\allowbreak{} ms241,\allowbreak{}\allowbreak{} factorization\allowbreak{} condition}{}{,\allowbreak{}\textbackslash{}\allowbreak{}quad\allowbreak{} \textbackslash{}\allowbreak{}bar\allowbreak{} }{DER\_\allowbreak{}XIIc-b:\allowbreak{}7}
\DERapparatusitem{ENV0817}{156;\allowbreak{}\allowbreak{} ms241,\allowbreak{}\allowbreak{} diagram\allowbreak{} bottom-left\allowbreak{} arrow}{S\allowbreak{} \textbackslash{}\allowbreak{}\&\allowbreak{}\allowbreak{} \textbackslash{}\allowbreak{}bar\{\allowbreak{}X\}\allowbreak{}\_\allowbreak{}2\allowbreak{} \textbackslash{}\allowbreak{}arrow[\allowbreak{}l,\allowbreak{}\allowbreak{} "\textbackslash{}\allowbreak{}bar\{\allowbreak{}q\}\allowbreak{}\_\allowbreak{}2"']\allowbreak{}}{S\allowbreak{} \textbackslash{}\allowbreak{}\&\allowbreak{}\allowbreak{} \textbackslash{}\allowbreak{}bar\{\allowbreak{}X\}\allowbreak{}\_\allowbreak{}2\allowbreak{} \textbackslash{}\allowbreak{}arrow[\allowbreak{}l,\allowbreak{}\allowbreak{} "\textbackslash{}\allowbreak{}bar\{\allowbreak{}p\}\allowbreak{}\_\allowbreak{}2"']\allowbreak{}}{DER\_\allowbreak{}XIIc-b:\allowbreak{}8}
\DERapparatusitem{ENV0818}{157;\allowbreak{}\allowbreak{} running\allowbreak{} text\allowbreak{} or\allowbreak{} formula}{\textbackslash{}\allowbreak{}text\{\allowbreak{}def\}\allowbreak{}}{\textbackslash{}\allowbreak{}mathrm\{\allowbreak{}def\}\allowbreak{}}{TYPESETTING:\allowbreak{}7;\allowbreak{}\allowbreak{} NOTATION\_\allowbreak{}DER:\allowbreak{}17}
\DERapparatusitem{ENV0819}{158;\allowbreak{}\allowbreak{} ms244,\allowbreak{}\allowbreak{} display\allowbreak{} (\allowbreak{}*)\allowbreak{}}{\textbackslash{}\allowbreak{}Psi}{\textbackslash{}\allowbreak{}Phi}{DER\_\allowbreak{}XIIc-b:\allowbreak{}9}
\DERapparatusitem{ENV0820}{158;\allowbreak{}\allowbreak{} running\allowbreak{} text\allowbreak{} or\allowbreak{} formula}{\textbackslash{}\allowbreak{}text\{\allowbreak{}def\}\allowbreak{}}{\textbackslash{}\allowbreak{}mathrm\{\allowbreak{}def\}\allowbreak{}}{TYPESETTING:\allowbreak{}7;\allowbreak{}\allowbreak{} NOTATION\_\allowbreak{}DER:\allowbreak{}17}
\DERapparatusitem{ENV0821}{159;\allowbreak{}\allowbreak{} running\allowbreak{} text\allowbreak{} or\allowbreak{} formula}{\textbackslash{}\allowbreak{}text\{\allowbreak{}def\}\allowbreak{}}{\textbackslash{}\allowbreak{}mathrm\{\allowbreak{}def\}\allowbreak{}}{TYPESETTING:\allowbreak{}7;\allowbreak{}\allowbreak{} NOTATION\_\allowbreak{}DER:\allowbreak{}17}
\DERapparatusitem{ENV0822}{161;\allowbreak{}\allowbreak{} running\allowbreak{} text\allowbreak{} or\allowbreak{} formula}{\textbackslash{}\allowbreak{}text\{\allowbreak{}def\}\allowbreak{}}{\textbackslash{}\allowbreak{}mathrm\{\allowbreak{}def\}\allowbreak{}}{TYPESETTING:\allowbreak{}7;\allowbreak{}\allowbreak{} NOTATION\_\allowbreak{}DER:\allowbreak{}17}
\DERapparatusitem{ENV0823}{162;\allowbreak{}\allowbreak{} ms249,\allowbreak{}\allowbreak{} final\allowbreak{} codomain\allowbreak{} below\allowbreak{} cospan-completion\allowbreak{} display}{\textbackslash{}\allowbreak{}Psi}{\textbackslash{}\allowbreak{}tilde\textbackslash{}\allowbreak{}Psi}{DER\_\allowbreak{}XIIc-b:\allowbreak{}11}
\DERapparatusitem{ENV0824}{162;\allowbreak{}\allowbreak{} ms250,\allowbreak{}\allowbreak{} final\allowbreak{} parenthesis}{}{'}{DER\_\allowbreak{}XIIc-b:\allowbreak{}12}
\DERapparatusitem{ENV0825}{162;\allowbreak{}\allowbreak{} ms251,\allowbreak{}\allowbreak{} display\allowbreak{} (\allowbreak{}*)\allowbreak{}}{\textbackslash{}\allowbreak{}Phi}{\textbackslash{}\allowbreak{}Psi}{DER\_\allowbreak{}XIIc-b:\allowbreak{}13}
\DERapparatusitem{ENV0826}{162;\allowbreak{}\allowbreak{} ms251,\allowbreak{}\allowbreak{} adjunction\allowbreak{} display\allowbreak{} at\allowbreak{} foot\allowbreak{} of\allowbreak{} leaf}{\textbackslash{}\allowbreak{}Ho\_\allowbreak{}W(\allowbreak{}\textbackslash{}\allowbreak{}Phi)\allowbreak{}\allowbreak{} \textbackslash{}\allowbreak{}arrow[\allowbreak{}r,\allowbreak{}\allowbreak{} shift\allowbreak{} right,\allowbreak{}\allowbreak{} "i\_\allowbreak{}*"']\allowbreak{}\allowbreak{} \textbackslash{}\allowbreak{}\&\allowbreak{}\allowbreak{} \textbackslash{}\allowbreak{}Ho\_\allowbreak{}W(\allowbreak{}\textbackslash{}\allowbreak{}tilde\{\allowbreak{}\textbackslash{}\allowbreak{}Phi\}\allowbreak{})\allowbreak{}}{\textbackslash{}\allowbreak{}Ho\_\allowbreak{}W(\allowbreak{}\textbackslash{}\allowbreak{}Psi)\allowbreak{}\allowbreak{} \textbackslash{}\allowbreak{}arrow[\allowbreak{}r,\allowbreak{}\allowbreak{} shift\allowbreak{} right,\allowbreak{}\allowbreak{} "i\_\allowbreak{}*"']\allowbreak{}\allowbreak{} \textbackslash{}\allowbreak{}\&\allowbreak{}\allowbreak{} \textbackslash{}\allowbreak{}Ho\_\allowbreak{}W(\allowbreak{}\textbackslash{}\allowbreak{}tilde\{\allowbreak{}\textbackslash{}\allowbreak{}Psi\}\allowbreak{})\allowbreak{}}{DER\_\allowbreak{}XIIc-b:\allowbreak{}14}
\DERapparatusitem{ENV0827}{163;\allowbreak{}\allowbreak{} ms251,\allowbreak{}\allowbreak{} display\allowbreak{} (\allowbreak{}**)\allowbreak{}}{\textbackslash{}\allowbreak{}tilde\{\allowbreak{}\textbackslash{}\allowbreak{}Phi\}\allowbreak{}}{\textbackslash{}\allowbreak{}tilde\textbackslash{}\allowbreak{}Psi}{DER\_\allowbreak{}XIIc-b:\allowbreak{}15}
\DERapparatusitem{ENV0828}{163;\allowbreak{}\allowbreak{} ms251,\allowbreak{}\allowbreak{} Lemma2,\allowbreak{}\allowbreak{} every\allowbreak{} ambient\allowbreak{} square-shape\allowbreak{} occurrence}{\textbackslash{}\allowbreak{}tilde\{\allowbreak{}\textbackslash{}\allowbreak{}Phi\}\allowbreak{}}{\textbackslash{}\allowbreak{}tilde\{\allowbreak{}\textbackslash{}\allowbreak{}Psi\}\allowbreak{}}{DER\_\allowbreak{}XIIc-b:\allowbreak{}16}
\DERapparatusitem{ENV0829}{163;\allowbreak{}\allowbreak{} ms251,\allowbreak{}\allowbreak{} Lemma2,\allowbreak{}\allowbreak{} every\allowbreak{} ambient\allowbreak{} square-shape\allowbreak{} occurrence}{\textbackslash{}\allowbreak{}tilde\{\allowbreak{}\textbackslash{}\allowbreak{}Phi\}\allowbreak{}}{\textbackslash{}\allowbreak{}tilde\{\allowbreak{}\textbackslash{}\allowbreak{}Psi\}\allowbreak{}}{DER\_\allowbreak{}XIIc-b:\allowbreak{}16}
\DERapparatusitem{ENV0830}{163;\allowbreak{}\allowbreak{} ms251,\allowbreak{}\allowbreak{} Lemma2,\allowbreak{}\allowbreak{} every\allowbreak{} ambient\allowbreak{} square-shape\allowbreak{} occurrence}{\textbackslash{}\allowbreak{}tilde\{\allowbreak{}\textbackslash{}\allowbreak{}Phi\}\allowbreak{}}{\textbackslash{}\allowbreak{}tilde\{\allowbreak{}\textbackslash{}\allowbreak{}Psi\}\allowbreak{}}{DER\_\allowbreak{}XIIc-b:\allowbreak{}16}
\DERapparatusitem{ENV0831}{163;\allowbreak{}\allowbreak{} ms251,\allowbreak{}\allowbreak{} Lemma2,\allowbreak{}\allowbreak{} every\allowbreak{} ambient\allowbreak{} square-shape\allowbreak{} occurrence}{\textbackslash{}\allowbreak{}tilde\{\allowbreak{}\textbackslash{}\allowbreak{}Phi\}\allowbreak{}}{\textbackslash{}\allowbreak{}tilde\{\allowbreak{}\textbackslash{}\allowbreak{}Psi\}\allowbreak{}}{DER\_\allowbreak{}XIIc-b:\allowbreak{}16}
\DERapparatusitem{ENV0832}{163;\allowbreak{}\allowbreak{} ms251,\allowbreak{}\allowbreak{} Lemma2,\allowbreak{}\allowbreak{} every\allowbreak{} ambient\allowbreak{} square-shape\allowbreak{} occurrence}{\textbackslash{}\allowbreak{}tilde\{\allowbreak{}\textbackslash{}\allowbreak{}Phi\}\allowbreak{}}{\textbackslash{}\allowbreak{}tilde\{\allowbreak{}\textbackslash{}\allowbreak{}Psi\}\allowbreak{}}{DER\_\allowbreak{}XIIc-b:\allowbreak{}16}
\DERapparatusitem{ENV0833}{163;\allowbreak{}\allowbreak{} running\allowbreak{} text\allowbreak{} or\allowbreak{} formula}{\textbackslash{}\allowbreak{}text\{\allowbreak{}def\}\allowbreak{}}{\textbackslash{}\allowbreak{}mathrm\{\allowbreak{}def\}\allowbreak{}}{TYPESETTING:\allowbreak{}7;\allowbreak{}\allowbreak{} NOTATION\_\allowbreak{}DER:\allowbreak{}17}
\DERapparatusitem{ENV0834}{163;\allowbreak{}\allowbreak{} ms251,\allowbreak{}\allowbreak{} Lemma2,\allowbreak{}\allowbreak{} every\allowbreak{} ambient\allowbreak{} square-shape\allowbreak{} occurrence}{\textbackslash{}\allowbreak{}tilde\{\allowbreak{}\textbackslash{}\allowbreak{}Phi\}\allowbreak{}}{\textbackslash{}\allowbreak{}tilde\{\allowbreak{}\textbackslash{}\allowbreak{}Psi\}\allowbreak{}}{DER\_\allowbreak{}XIIc-b:\allowbreak{}16}
\DERapparatusitem{ENV0835}{165;\allowbreak{}\allowbreak{} ms254,\allowbreak{}\allowbreak{} concluding\allowbreak{} NB\allowbreak{} type\allowbreak{} of\allowbreak{} restriction}{\textbackslash{}\allowbreak{}Psi)\allowbreak{}\textbackslash{}\allowbreak{}longrightarrow\allowbreak{} \textbackslash{}\allowbreak{}Ho\_\allowbreak{}W(\allowbreak{}\textbackslash{}\allowbreak{}tilde}{\textbackslash{}\allowbreak{}tilde\textbackslash{}\allowbreak{}Psi)\allowbreak{}\textbackslash{}\allowbreak{}longrightarrow\textbackslash{}\allowbreak{}Ho\_\allowbreak{}W(\allowbreak{}}{DER\_\allowbreak{}XIIc-b:\allowbreak{}20}
\DERapparatusitem{ENV0836}{166;\allowbreak{}\allowbreak{} ms256,\allowbreak{}\allowbreak{} proof\allowbreak{} of\allowbreak{} Corollary1}{\textbackslash{}\allowbreak{}psi}{\textbackslash{}\allowbreak{}Psi}{DER\_\allowbreak{}XIIc-b:\allowbreak{}22}
\DERapparatusitem{ENV0837}{167;\allowbreak{}\allowbreak{} ms257,\allowbreak{}\allowbreak{} Lemma,\allowbreak{}\allowbreak{} direction\allowbreak{} of\allowbreak{} adjunction}{This\allowbreak{} functor\allowbreak{} admits\allowbreak{} a\allowbreak{} left\allowbreak{} adjoint}{This\allowbreak{} functor\allowbreak{} admits\allowbreak{} a\allowbreak{} right\allowbreak{} adjoint}{DER\_\allowbreak{}XIIc-b:\allowbreak{}23}
\DERapparatusitem{ENV0838}{167;\allowbreak{}\allowbreak{} ms258,\allowbreak{}\allowbreak{} Proposition19,\allowbreak{}\allowbreak{} formula\allowbreak{} for\allowbreak{} j2}{\textbackslash{}\allowbreak{}times}{\textbackslash{}\allowbreak{}times\_\allowbreak{}S}{DER\_\allowbreak{}XIIc-b:\allowbreak{}24}
\DERapparatusitem{ENV0839}{169;\allowbreak{}\allowbreak{} ms260,\allowbreak{}\allowbreak{} first\allowbreak{} paragraph\allowbreak{} of\allowbreak{} the\allowbreak{} pasting\allowbreak{} proof}{we\allowbreak{} can\allowbreak{} assume\allowbreak{} \$\allowbreak{}X\_\allowbreak{}0\textbackslash{}\allowbreak{}longrightarrow\allowbreak{} S\$\allowbreak{}\allowbreak{} to\allowbreak{} be\allowbreak{} in\allowbreak{} \$\allowbreak{}\textbackslash{}\allowbreak{}Fib\_\allowbreak{}W\allowbreak{} S\$\allowbreak{}}{we\allowbreak{} can\allowbreak{} assume\allowbreak{} \$\allowbreak{}X\_\allowbreak{}0'\textbackslash{}\allowbreak{}longrightarrow\allowbreak{} X\_\allowbreak{}0\$\allowbreak{}\allowbreak{} to\allowbreak{} be\allowbreak{} in\allowbreak{} \$\allowbreak{}\textbackslash{}\allowbreak{}Fib\_\allowbreak{}W\$\allowbreak{}}{DER\_\allowbreak{}XIIc-b:\allowbreak{}25}
\DERapparatusitem{ENV0840}{169;\allowbreak{}\allowbreak{} running\allowbreak{} text\allowbreak{} or\allowbreak{} formula}{\textbackslash{}\allowbreak{}text\{\allowbreak{}def\}\allowbreak{}}{\textbackslash{}\allowbreak{}mathrm\{\allowbreak{}def\}\allowbreak{}}{TYPESETTING:\allowbreak{}7;\allowbreak{}\allowbreak{} NOTATION\_\allowbreak{}DER:\allowbreak{}17}
\DERapparatusitem{ENV0841}{169;\allowbreak{}\allowbreak{} running\allowbreak{} text\allowbreak{} or\allowbreak{} formula}{\textbackslash{}\allowbreak{}text\{\allowbreak{}def\}\allowbreak{}}{\textbackslash{}\allowbreak{}mathrm\{\allowbreak{}def\}\allowbreak{}}{TYPESETTING:\allowbreak{}7;\allowbreak{}\allowbreak{} NOTATION\_\allowbreak{}DER:\allowbreak{}17}
\DERapparatusitem{ENV0842}{170;\allowbreak{}\allowbreak{} ms262,\allowbreak{}\allowbreak{} Proposition20,\allowbreak{}\allowbreak{} adjective}{For\allowbreak{} it\allowbreak{} to\allowbreak{} be\allowbreak{} \$\allowbreak{}W\$\allowbreak{}-cartesian}{For\allowbreak{} it\allowbreak{} to\allowbreak{} be\allowbreak{} \$\allowbreak{}W\$\allowbreak{}-cocartesian}{DER\_\allowbreak{}XIIc-b:\allowbreak{}29}
\DERapparatusitem{ENV0843}{170;\allowbreak{}\allowbreak{} ms262,\allowbreak{}\allowbreak{} Proposition20,\allowbreak{}\allowbreak{} integral\allowbreak{} diagram}{\textbackslash{}\allowbreak{}\&\allowbreak{}\allowbreak{} X\_\allowbreak{}1\allowbreak{} \textbackslash{}\allowbreak{}arrow[\allowbreak{}dl]\allowbreak{}\allowbreak{} \textbackslash{}\allowbreak{}\textbackslash{}\allowbreak{}\space{}\allowbreak{}\textbackslash{}\allowbreak{}int\allowbreak{} S\allowbreak{} \textbackslash{}\allowbreak{}\&\allowbreak{}\allowbreak{} \textbackslash{}\allowbreak{}\textbackslash{}\allowbreak{}\space{}\allowbreak{}\allowbreak{} \textbackslash{}\allowbreak{}\&\allowbreak{}\allowbreak{} X\_\allowbreak{}2\allowbreak{} \textbackslash{}\allowbreak{}arrow[\allowbreak{}ul]\allowbreak{}}{\allowbreak{} \textbackslash{}\allowbreak{}\&\allowbreak{}\allowbreak{} X\_\allowbreak{}1\allowbreak{} \textbackslash{}\allowbreak{}\textbackslash{}\allowbreak{}\space{}\allowbreak{}\textbackslash{}\allowbreak{}int\allowbreak{} S\allowbreak{} \textbackslash{}\allowbreak{}arrow[\allowbreak{}ur]\allowbreak{}\allowbreak{} \textbackslash{}\allowbreak{}arrow[\allowbreak{}dr]\allowbreak{}\allowbreak{} \textbackslash{}\allowbreak{}\&\allowbreak{}\allowbreak{} \textbackslash{}\allowbreak{}\textbackslash{}\allowbreak{}\space{}\allowbreak{}\allowbreak{} \textbackslash{}\allowbreak{}\&\allowbreak{}\allowbreak{} X\_\allowbreak{}2}{DER\_\allowbreak{}XIIc-b:\allowbreak{}30}
\DERapparatusitem{ENV0844}{171;\allowbreak{}\allowbreak{} ms263,\allowbreak{}\allowbreak{} factorization\allowbreak{} display}{}{\_\allowbreak{}1}{DER\_\allowbreak{}XIIc-b:\allowbreak{}31}
\DERapparatusitem{ENV0845}{171;\allowbreak{}\allowbreak{} ms263,\allowbreak{}\allowbreak{} triangle\allowbreak{} labelled\allowbreak{} mathcalX\_\allowbreak{}alpha=\allowbreak{}lambda}{\textbackslash{}\allowbreak{}mathcal\{\allowbreak{}X\}\allowbreak{}\_\allowbreak{}Q\allowbreak{} \textbackslash{}\allowbreak{}arrow[\allowbreak{}rr,\allowbreak{}\allowbreak{} "\textbackslash{}\allowbreak{}mathcal\{\allowbreak{}X\}\allowbreak{}\_\allowbreak{}\textbackslash{}\allowbreak{}alpha\allowbreak{} =\allowbreak{}\allowbreak{} \textbackslash{}\allowbreak{}lambda"]\allowbreak{}}{\textbackslash{}\allowbreak{}mathcal\{\allowbreak{}X\}\allowbreak{}\_\allowbreak{}Q\allowbreak{} \textbackslash{}\allowbreak{}arrow[\allowbreak{}rr,\allowbreak{}\allowbreak{} leftarrow,\allowbreak{}\allowbreak{} "\textbackslash{}\allowbreak{}mathcal\{\allowbreak{}X\}\allowbreak{}\_\allowbreak{}\textbackslash{}\allowbreak{}alpha\allowbreak{} =\allowbreak{}\allowbreak{} \textbackslash{}\allowbreak{}lambda"]\allowbreak{}}{DER\_\allowbreak{}XIIc-b:\allowbreak{}32}
\DERapparatusitem{ENV0846}{176;\allowbreak{}\allowbreak{} ms272,\allowbreak{}\allowbreak{} footnote137}{cf.\allowbreak{}\allowbreak{} XV\allowbreak{} \textbackslash{}\allowbreak{}S3}{cf.\allowbreak{}\allowbreak{} XVI\allowbreak{} \textbackslash{}\allowbreak{}S3}{DER\_\allowbreak{}XIIc-b:\allowbreak{}36}
\DERapparatusitem{ENV0847}{178;\allowbreak{}\allowbreak{} ms273,\allowbreak{}\allowbreak{} Cons1\allowbreak{} diagram}{S\allowbreak{} \textbackslash{}\allowbreak{}arrow[\allowbreak{}rr]\allowbreak{}\allowbreak{} \textbackslash{}\allowbreak{}arrow[\allowbreak{}dr]\allowbreak{}\allowbreak{} \textbackslash{}\allowbreak{}\&\allowbreak{}\allowbreak{} \textbackslash{}\allowbreak{}\&\allowbreak{}\allowbreak{} T}{S\allowbreak{} \textbackslash{}\allowbreak{}arrow[\allowbreak{}rr,\allowbreak{}\allowbreak{} "f"]\allowbreak{}\allowbreak{} \textbackslash{}\allowbreak{}arrow[\allowbreak{}dr]\allowbreak{}\allowbreak{} \textbackslash{}\allowbreak{}\&\allowbreak{}\allowbreak{} \textbackslash{}\allowbreak{}\&\allowbreak{}\allowbreak{} T}{DER\_\allowbreak{}XIIc-b:\allowbreak{}37}
\DERapparatusitem{ENV0848}{178;\allowbreak{}\allowbreak{} ms273,\allowbreak{}\allowbreak{} Cons2}{\textbackslash{}\allowbreak{}textbf\{\allowbreak{}Cons\allowbreak{} 2.\allowbreak{}\}\allowbreak{}\allowbreak{} \textbackslash{}\allowbreak{}textit\{\allowbreak{}If\allowbreak{} we\allowbreak{} have\allowbreak{} a\allowbreak{} cartesian\allowbreak{} square,\allowbreak{}\}\allowbreak{}}{\textbackslash{}\allowbreak{}textbf\{\allowbreak{}Cons\allowbreak{} 2.\allowbreak{}\}\allowbreak{}\allowbreak{} \textbackslash{}\allowbreak{}textit\{\allowbreak{}If\allowbreak{} we\allowbreak{} have\allowbreak{} a\allowbreak{} cocartesian\allowbreak{} square,\allowbreak{}\}\allowbreak{}}{DER\_\allowbreak{}XIIc-b:\allowbreak{}38}
\subsection*{Chapter XII, appendix}
\DERapparatusitem{ENV0695}{preamble;\allowbreak{}\allowbreak{} local\allowbreak{} macro\allowbreak{} definition}{\textbackslash{}\allowbreak{}text\{\allowbreak{}def\}\allowbreak{}}{\textbackslash{}\allowbreak{}mathrm\{\allowbreak{}def\}\allowbreak{}}{TYPESETTING:\allowbreak{}7;\allowbreak{}\allowbreak{} NOTATION\_\allowbreak{}DER:\allowbreak{}17}
\DERapparatusitem{ENV0696}{1;\allowbreak{}\allowbreak{} ms.\allowbreak{}1,\allowbreak{}\allowbreak{} last\allowbreak{} line}{\textbackslash{}\allowbreak{}alpha\_\allowbreak{}0\allowbreak{} :\allowbreak{}\allowbreak{} x\textbackslash{}\allowbreak{}longrightarrow\allowbreak{} a\_\allowbreak{}0}{\textbackslash{}\allowbreak{}alpha\_\allowbreak{}0\allowbreak{} :\allowbreak{}\allowbreak{} x\textbackslash{}\allowbreak{}longrightarrow\allowbreak{} a}{DER\_\allowbreak{}XIIapp:\allowbreak{}1}
\DERapparatusitem{ENV0697}{1;\allowbreak{}\allowbreak{} ms.\allowbreak{}2,\allowbreak{}\allowbreak{} first\allowbreak{} sentence}{}{)\allowbreak{}}{DER\_\allowbreak{}XIIapp:\allowbreak{}3}
\DERapparatusitem{ENV0698}{1;\allowbreak{}\allowbreak{} ms.\allowbreak{}2,\allowbreak{}\allowbreak{} first\allowbreak{} sentence}{u}{\textbackslash{}\allowbreak{}alpha}{DER\_\allowbreak{}XIIapp:\allowbreak{}2}
\DERapparatusitem{ENV0699}{2;\allowbreak{}\allowbreak{} ms.\allowbreak{}2,\allowbreak{}\allowbreak{} path-prefix\allowbreak{} indexing}{defined\allowbreak{} by\allowbreak{} the\allowbreak{} \$\allowbreak{}u\_\allowbreak{}i\$\allowbreak{}\allowbreak{} (\allowbreak{}\$\allowbreak{}i\textbackslash{}\allowbreak{}in[\allowbreak{}1,\allowbreak{}i]\allowbreak{}\$\allowbreak{})\allowbreak{}}{defined\allowbreak{} by\allowbreak{} the\allowbreak{} \$\allowbreak{}u\_\allowbreak{}j\$\allowbreak{}\allowbreak{} (\allowbreak{}\$\allowbreak{}j\textbackslash{}\allowbreak{}in[\allowbreak{}1,\allowbreak{}i]\allowbreak{}\$\allowbreak{})\allowbreak{}}{DER\_\allowbreak{}XIIapp:\allowbreak{}4}
\DERapparatusitem{ENV0700}{2;\allowbreak{}\allowbreak{} ms.\allowbreak{}3,\allowbreak{}\allowbreak{} definition\allowbreak{} of\allowbreak{} the\allowbreak{} two-object\allowbreak{} monoid\allowbreak{} category}{\$\allowbreak{}\textbackslash{}\allowbreak{}Hom\_\allowbreak{}X(\allowbreak{}b\textbackslash{}\allowbreak{}times\allowbreak{} b)\allowbreak{}\$\allowbreak{}}{\$\allowbreak{}\textbackslash{}\allowbreak{}Hom\_\allowbreak{}X(\allowbreak{}b\textbackslash{}\allowbreak{}times\allowbreak{} b,\allowbreak{}b)\allowbreak{}\$\allowbreak{}}{DER\_\allowbreak{}XIIapp:\allowbreak{}11}
\DERapparatusitem{ENV0701}{2;\allowbreak{}\allowbreak{} ms.\allowbreak{}3,\allowbreak{}\allowbreak{} cohomology\allowbreak{} display\allowbreak{} last\allowbreak{} line}{H\textasciicircum{}\allowbreak{}2(\allowbreak{}T,\allowbreak{}F)\allowbreak{}\allowbreak{} \&\allowbreak{}=\allowbreak{}\allowbreak{} \textbackslash{}\allowbreak{}Coker\textbackslash{}\allowbreak{}delta}{H\textasciicircum{}\allowbreak{}1(\allowbreak{}T,\allowbreak{}F)\allowbreak{}\allowbreak{} \&\allowbreak{}=\allowbreak{}\allowbreak{} \textbackslash{}\allowbreak{}Coker\textbackslash{}\allowbreak{}delta}{DER\_\allowbreak{}XIIapp:\allowbreak{}10}
\DERapparatusitem{ENV0702}{3;\allowbreak{}\allowbreak{} ms.\allowbreak{}4,\allowbreak{}\allowbreak{} spectral-sequence\allowbreak{} conclusion}{S}{T}{DER\_\allowbreak{}XIIapp:\allowbreak{}12}
\DERapparatusitem{ENV0703}{6;\allowbreak{}\allowbreak{} ms.\allowbreak{}9,\allowbreak{}\allowbreak{} represented\allowbreak{} pushout}{X}{V}{DER\_\allowbreak{}XIIapp:\allowbreak{}15}
\DERapparatusitem{ENV0704}{6;\allowbreak{}\allowbreak{} ms.\allowbreak{}9,\allowbreak{}\allowbreak{} asserted\allowbreak{} retraction\allowbreak{} paragraph}{X\_\allowbreak{}0\allowbreak{} \textbackslash{}\allowbreak{}hookrightarrow\allowbreak{} X}{V\_\allowbreak{}0\textbackslash{}\allowbreak{}hookrightarrow\allowbreak{} V}{DER\_\allowbreak{}XIIapp:\allowbreak{}16}
\DERapparatusitem{ENV0705}{7;\allowbreak{}\allowbreak{} ms.\allowbreak{}10,\allowbreak{}\allowbreak{} representable\allowbreak{} evaluation}{a(\allowbreak{}z)\allowbreak{}\allowbreak{} =\allowbreak{}\allowbreak{} \textbackslash{}\allowbreak{}Hom\_\allowbreak{}Z(\allowbreak{}a,\allowbreak{}z)\allowbreak{}}{a(\allowbreak{}z)\allowbreak{}\allowbreak{} =\allowbreak{}\allowbreak{} \textbackslash{}\allowbreak{}Hom\_\allowbreak{}Z(\allowbreak{}z,\allowbreak{}a)\allowbreak{}}{DER\_\allowbreak{}XIIapp:\allowbreak{}17}
\DERapparatusitem{ENV0706}{9;\allowbreak{}\allowbreak{} running\allowbreak{} text\allowbreak{} or\allowbreak{} formula}{\textbackslash{}\allowbreak{}text\{\allowbreak{}def\}\allowbreak{}}{\textbackslash{}\allowbreak{}mathrm\{\allowbreak{}def\}\allowbreak{}}{TYPESETTING:\allowbreak{}7;\allowbreak{}\allowbreak{} NOTATION\_\allowbreak{}DER:\allowbreak{}17}
\DERapparatusitem{ENV0707}{10;\allowbreak{}\allowbreak{} ms.\allowbreak{}13,\allowbreak{}\allowbreak{} first\allowbreak{} display\allowbreak{} of\allowbreak{} proof}{F}{V}{DER\_\allowbreak{}XIIapp:\allowbreak{}19}
\DERapparatusitem{ENV0708}{11;\allowbreak{}\allowbreak{} ms.\allowbreak{}14,\allowbreak{}\allowbreak{} sentence\allowbreak{} before\allowbreak{} p\_\allowbreak{}0}{of\allowbreak{} \$\allowbreak{}V\$\allowbreak{}\allowbreak{} into\allowbreak{} \$\allowbreak{}V'\$\allowbreak{}}{of\allowbreak{} \$\allowbreak{}V\$\allowbreak{}\allowbreak{} into\allowbreak{} \$\allowbreak{}V'\_\allowbreak{}0\$\allowbreak{}}{DER\_\allowbreak{}XIIapp:\allowbreak{}21}
\DERapparatusitem{ENV0709}{11;\allowbreak{}\allowbreak{} ms.\allowbreak{}15,\allowbreak{}\allowbreak{} last\allowbreak{} paragraph}{\$\allowbreak{}\textbackslash{}\allowbreak{}pi\_\allowbreak{}1(\allowbreak{}V'\_\allowbreak{}0)\allowbreak{}\allowbreak{} \textbackslash{}\allowbreak{}simeq\allowbreak{} \textbackslash{}\allowbreak{}mathbf\{\allowbreak{}Z\}\allowbreak{}\$\allowbreak{}}{\$\allowbreak{}\textbackslash{}\allowbreak{}pi\_\allowbreak{}1(\allowbreak{}V'\_\allowbreak{}0)\allowbreak{}=\allowbreak{}0\$\allowbreak{}\allowbreak{} and\allowbreak{} \$\allowbreak{}\textbackslash{}\allowbreak{}pi\_\allowbreak{}2(\allowbreak{}V'\_\allowbreak{}0)\allowbreak{}\allowbreak{} \textbackslash{}\allowbreak{}simeq\allowbreak{} \textbackslash{}\allowbreak{}mathbf\{\allowbreak{}Z\}\allowbreak{}\$\allowbreak{}}{DER\_\allowbreak{}XIIapp:\allowbreak{}22}
\DERapparatusitem{ENV0710}{13;\allowbreak{}\allowbreak{} ms.\allowbreak{}17,\allowbreak{}\allowbreak{} first\allowbreak{} paragraph}{or\allowbreak{} j\allowbreak{} \$\allowbreak{}W\$\allowbreak{}-coaspheric}{or\allowbreak{} \$\allowbreak{}j\$\allowbreak{}\allowbreak{} \$\allowbreak{}W\$\allowbreak{}-coaspheric}{DER\_\allowbreak{}XIIapp:\allowbreak{}23}
\DERapparatusitem{ENV0711}{13;\allowbreak{}\allowbreak{} ms.\allowbreak{}17,\allowbreak{}\allowbreak{} special-case\allowbreak{} parenthesis}{}{\_\allowbreak{}1}{DER\_\allowbreak{}XIIapp:\allowbreak{}24}
\DERapparatusitem{ENV0712}{13;\allowbreak{}\allowbreak{} running\allowbreak{} text\allowbreak{} or\allowbreak{} formula}{\textbackslash{}\allowbreak{}text\{\allowbreak{}def\}\allowbreak{}}{\textbackslash{}\allowbreak{}mathrm\{\allowbreak{}def\}\allowbreak{}}{TYPESETTING:\allowbreak{}7;\allowbreak{}\allowbreak{} NOTATION\_\allowbreak{}DER:\allowbreak{}17}
\DERapparatusitem{ENV0713}{13;\allowbreak{}\allowbreak{} ms.\allowbreak{}18,\allowbreak{}\allowbreak{} footnote10}{I\allowbreak{} do\allowbreak{} not\allowbreak{} have\allowbreak{} hardly\allowbreak{} any\allowbreak{} doubt}{I\allowbreak{} have\allowbreak{} hardly\allowbreak{} any\allowbreak{} doubt}{DER\_\allowbreak{}XIIapp:\allowbreak{}26}
\subsection*{Chapter XIII}
\DERapparatusitem{ENV0584}{5;\allowbreak{}\allowbreak{} ms.\allowbreak{}7–8,\allowbreak{}\allowbreak{} M\allowbreak{} 0}{and\allowbreak{} finite\allowbreak{} \$\allowbreak{}\textbackslash{}\allowbreak{}varprojlim\$\allowbreak{}\allowbreak{} \textbackslash{}\allowbreak{}texttt\{\allowbreak{}[\allowbreak{}rather\allowbreak{} \$\allowbreak{}\textbackslash{}\allowbreak{}varinjlim\$\allowbreak{}]\allowbreak{}\}\allowbreak{}}{and\allowbreak{} finite\allowbreak{} \$\allowbreak{}\textbackslash{}\allowbreak{}varinjlim\$\allowbreak{}\allowbreak{} \textbackslash{}\allowbreak{}texttt\{\allowbreak{}[\allowbreak{}rather\allowbreak{} \$\allowbreak{}\textbackslash{}\allowbreak{}varinjlim\$\allowbreak{}]\allowbreak{}\}\allowbreak{}}{DER\_\allowbreak{}XIII-a:\allowbreak{}3}
\DERapparatusitem{ENV0585}{6;\allowbreak{}\allowbreak{} ms.\allowbreak{}9,\allowbreak{}\allowbreak{} final\allowbreak{} M-double-prime\allowbreak{} in\allowbreak{} Proposition\allowbreak{} 3}{M\allowbreak{} 1\allowbreak{} to\allowbreak{} M\allowbreak{} 4,\allowbreak{}\allowbreak{} M\$\allowbreak{}'\$\allowbreak{},\allowbreak{}\allowbreak{} \$\allowbreak{}M''\$\allowbreak{}}{M\allowbreak{} 1\allowbreak{} to\allowbreak{} M\allowbreak{} 4,\allowbreak{}\allowbreak{} M\$\allowbreak{}'\$\allowbreak{},\allowbreak{}\allowbreak{} M\$\allowbreak{}''\$\allowbreak{}}{DER\_\allowbreak{}XIII-a:\allowbreak{}4}
\DERapparatusitem{ENV0586}{6;\allowbreak{}\allowbreak{} running\allowbreak{} text\allowbreak{} or\allowbreak{} formula}{\textbackslash{}\allowbreak{}text\{\allowbreak{}def\}\allowbreak{}}{\textbackslash{}\allowbreak{}mathrm\{\allowbreak{}def\}\allowbreak{}}{TYPESETTING:\allowbreak{}7;\allowbreak{}\allowbreak{} NOTATION\_\allowbreak{}DER:\allowbreak{}17}
\DERapparatusitem{ENV0587}{7;\allowbreak{}\allowbreak{} ms.\allowbreak{}11,\allowbreak{}\allowbreak{} Commentary\allowbreak{} on\allowbreak{} Proposition\allowbreak{} 3}{when\allowbreak{} only\allowbreak{} M\allowbreak{} 1,\allowbreak{}\allowbreak{} M2,\allowbreak{}\allowbreak{} M\$\allowbreak{}'\$\allowbreak{}}{when\allowbreak{} only\allowbreak{} M\allowbreak{} 1,\allowbreak{}\allowbreak{} M\allowbreak{} 2,\allowbreak{}\allowbreak{} M\$\allowbreak{}'\$\allowbreak{}}{DER\_\allowbreak{}XIII-a:\allowbreak{}5}
\DERapparatusitem{ENV0588}{8;\allowbreak{}\allowbreak{} ms.\allowbreak{}13,\allowbreak{}\allowbreak{} first\allowbreak{} paragraph\allowbreak{} of\allowbreak{} Scholie}{\textbackslash{}\allowbreak{}texttt\{\allowbreak{}[\allowbreak{}the]\allowbreak{}\}\allowbreak{}\allowbreak{} proposition\allowbreak{} 3}{\textbackslash{}\allowbreak{}texttt\{\allowbreak{}[\allowbreak{}the]\allowbreak{}\}\allowbreak{}\allowbreak{} proposition\allowbreak{} 4}{DER\_\allowbreak{}XIII-a:\allowbreak{}6}
\DERapparatusitem{ENV0589}{9;\allowbreak{}\allowbreak{} running\allowbreak{} text\allowbreak{} or\allowbreak{} formula}{\textbackslash{}\allowbreak{}text\{\allowbreak{}def\}\allowbreak{}}{\textbackslash{}\allowbreak{}mathrm\{\allowbreak{}def\}\allowbreak{}}{TYPESETTING:\allowbreak{}7;\allowbreak{}\allowbreak{} NOTATION\_\allowbreak{}DER:\allowbreak{}17}
\DERapparatusitem{ENV0590}{9;\allowbreak{}\allowbreak{} running\allowbreak{} text\allowbreak{} or\allowbreak{} formula}{\textbackslash{}\allowbreak{}text\{\allowbreak{}def\}\allowbreak{}}{\textbackslash{}\allowbreak{}mathrm\{\allowbreak{}def\}\allowbreak{}}{TYPESETTING:\allowbreak{}7;\allowbreak{}\allowbreak{} NOTATION\_\allowbreak{}DER:\allowbreak{}17}
\DERapparatusitem{ENV0591}{9;\allowbreak{}\allowbreak{} ms.\allowbreak{}16,\allowbreak{}\allowbreak{} paragraph\allowbreak{} introducing\allowbreak{} implication\allowbreak{} diagram}{tautological\allowbreak{} implication\allowbreak{} between}{tautological\allowbreak{} implications\allowbreak{} between}{DER\_\allowbreak{}XIII-a:\allowbreak{}8}
\DERapparatusitem{ENV0592}{10;\allowbreak{}\allowbreak{} ms.\allowbreak{}17,\allowbreak{}\allowbreak{} closures\allowbreak{} of\allowbreak{} generating\allowbreak{} classes}{\textbackslash{}\allowbreak{}overline\{\allowbreak{}TC\_\allowbreak{}0\}\allowbreak{}\allowbreak{} =\allowbreak{}\allowbreak{} TC,\allowbreak{}\allowbreak{} \textbackslash{}\allowbreak{}qquad\allowbreak{} \textbackslash{}\allowbreak{}overline\{\allowbreak{}C\_\allowbreak{}0\}\allowbreak{}\allowbreak{} =\allowbreak{}\allowbreak{} C}{\textbackslash{}\allowbreak{}widetilde\{\allowbreak{}TC\_\allowbreak{}0\}\allowbreak{}\allowbreak{} =\allowbreak{}\allowbreak{} TC,\allowbreak{}\allowbreak{} \textbackslash{}\allowbreak{}qquad\allowbreak{} \textbackslash{}\allowbreak{}widetilde\{\allowbreak{}C\_\allowbreak{}0\}\allowbreak{}\allowbreak{} =\allowbreak{}\allowbreak{} C}{DER\_\allowbreak{}XIII-a:\allowbreak{}9}
\DERapparatusitem{ENV0593}{11;\allowbreak{}\allowbreak{} ms.\allowbreak{}19,\allowbreak{}\allowbreak{} accessibility\allowbreak{} parenthesis}{the\allowbreak{} admissible\allowbreak{} \textbackslash{}\allowbreak{}texttt\{\allowbreak{}[\allowbreak{}rather\allowbreak{} accessible]\allowbreak{}\}\allowbreak{}\allowbreak{} subsets}{the\allowbreak{} accessible\allowbreak{} \textbackslash{}\allowbreak{}texttt\{\allowbreak{}[\allowbreak{}rather\allowbreak{} accessible]\allowbreak{}\}\allowbreak{}\allowbreak{} subsets}{DER\_\allowbreak{}XIII-a:\allowbreak{}10}
\DERapparatusitem{ENV0594}{11;\allowbreak{}\allowbreak{} ms.\allowbreak{}20,\allowbreak{}\allowbreak{} conjectural\allowbreak{} lemma\allowbreak{} (\allowbreak{}1)\allowbreak{},\allowbreak{}\allowbreak{} final\allowbreak{} parenthesis}{\textbackslash{}\allowbreak{}textit\{\allowbreak{}(\allowbreak{}resp.\allowbreak{}\}\allowbreak{}\allowbreak{} \$\allowbreak{}\textbackslash{}\allowbreak{}Phi\_\allowbreak{}0\allowbreak{} \textbackslash{}\allowbreak{}subset\allowbreak{} W\$\allowbreak{}}{\textbackslash{}\allowbreak{}textit\{\allowbreak{}(\allowbreak{}resp.\allowbreak{}\}\allowbreak{}\allowbreak{} \$\allowbreak{}(\allowbreak{}\textbackslash{}\allowbreak{}Phi\_\allowbreak{}0)\allowbreak{}\_\allowbreak{}*\allowbreak{} \textbackslash{}\allowbreak{}subset\allowbreak{} W\$\allowbreak{}}{DER\_\allowbreak{}XIII-a:\allowbreak{}11}
\DERapparatusitem{ENV0595}{14;\allowbreak{}\allowbreak{} manuscript\allowbreak{} p.\allowbreak{}25,\allowbreak{}\allowbreak{} Corollary(\allowbreak{}b)\allowbreak{},\allowbreak{}\allowbreak{} mathematical\allowbreak{} scope}{\textbackslash{}\allowbreak{}widetilde\{\allowbreak{}W\textasciicircum{}\allowbreak{}\{\allowbreak{}\textbackslash{}\allowbreak{}univ\}\allowbreak{}\allowbreak{} \textbackslash{}\allowbreak{}cup\allowbreak{} (\allowbreak{}W\_\allowbreak{}\{\allowbreak{}\textbackslash{}\allowbreak{}univ\}\allowbreak{})\allowbreak{}\}\allowbreak{}\textasciicircum{}\allowbreak{}\{\allowbreak{}*\}\allowbreak{}}{\textbackslash{}\allowbreak{}widetilde\{\allowbreak{}W\textasciicircum{}\allowbreak{}\{\allowbreak{}\textbackslash{}\allowbreak{}univ\}\allowbreak{}\textbackslash{}\allowbreak{}cup(\allowbreak{}W\_\allowbreak{}\{\allowbreak{}\textbackslash{}\allowbreak{}univ\}\allowbreak{})\allowbreak{}\textasciicircum{}\allowbreak{}*\}\allowbreak{}}{DER\_\allowbreak{}XIII-a:\allowbreak{}13}
\DERapparatusitem{ENV0596}{15;\allowbreak{}\allowbreak{} manuscript\allowbreak{} p.\allowbreak{}26,\allowbreak{}\allowbreak{} summary\allowbreak{} II(\allowbreak{}c)\allowbreak{}}{\_\allowbreak{}}{\textasciicircum{}\allowbreak{}}{DER\_\allowbreak{}XIII-a:\allowbreak{}14}
\DERapparatusitem{ENV0597}{16;\allowbreak{}\allowbreak{} manuscript\allowbreak{} p.\allowbreak{}26,\allowbreak{}\allowbreak{} parenthesis\allowbreak{} before\allowbreak{} the\allowbreak{} final\allowbreak{} inclusion}{with\allowbreak{} \$\allowbreak{}X\$\allowbreak{}\allowbreak{} discrete)\allowbreak{}}{with\allowbreak{} \$\allowbreak{}X\$\allowbreak{}\allowbreak{} nonempty\allowbreak{} and\allowbreak{} discrete)\allowbreak{}}{DER\_\allowbreak{}XIII-a:\allowbreak{}16}
\DERapparatusitem{ENV0598}{17;\allowbreak{}\allowbreak{} manuscript\allowbreak{} p.\allowbreak{}28,\allowbreak{}\allowbreak{} Remark(\allowbreak{}b)\allowbreak{},\allowbreak{}\allowbreak{} axiom\allowbreak{} references}{for\allowbreak{} the\allowbreak{} stabilities\allowbreak{} specified\allowbreak{} in\allowbreak{} M\allowbreak{} 2,\allowbreak{}\allowbreak{} M\allowbreak{} 3}{for\allowbreak{} the\allowbreak{} stabilities\allowbreak{} specified\allowbreak{} in\allowbreak{} M\allowbreak{} 3,\allowbreak{}\allowbreak{} M\allowbreak{} 4}{DER\_\allowbreak{}XIII-a:\allowbreak{}18}
\DERapparatusitem{ENV0599}{17;\allowbreak{}\allowbreak{} manuscript\allowbreak{} p.\allowbreak{}28,\allowbreak{}\allowbreak{} Remark(\allowbreak{}b)\allowbreak{},\allowbreak{}\allowbreak{} type\allowbreak{} of\allowbreak{} resulting\allowbreak{} structure}{we\allowbreak{} find\allowbreak{} a\allowbreak{} Quillen\allowbreak{} structure.\allowbreak{}}{we\allowbreak{} find\allowbreak{} a\allowbreak{} weak\allowbreak{} Quillen\allowbreak{} structure.\allowbreak{}}{DER\_\allowbreak{}XIII-a:\allowbreak{}19}
\DERapparatusitem{ENV0600}{18;\allowbreak{}\allowbreak{} running\allowbreak{} text\allowbreak{} or\allowbreak{} formula}{\textbackslash{}\allowbreak{}text\{\allowbreak{}def\}\allowbreak{}}{\textbackslash{}\allowbreak{}mathrm\{\allowbreak{}def\}\allowbreak{}}{TYPESETTING:\allowbreak{}7;\allowbreak{}\allowbreak{} NOTATION\_\allowbreak{}DER:\allowbreak{}17}
\DERapparatusitem{ENV0601}{19;\allowbreak{}\allowbreak{} manuscript\allowbreak{} p.\allowbreak{}31,\allowbreak{}\allowbreak{} localization\allowbreak{} display}{\textbackslash{}\allowbreak{}cM\_\allowbreak{}f}{\textbackslash{}\allowbreak{}cM\_\allowbreak{}\{\allowbreak{}\textbackslash{}\allowbreak{}mathrm\{\allowbreak{}f\}\allowbreak{}\}\allowbreak{}}{DER\_\allowbreak{}XIII-a:\allowbreak{}21}
\DERapparatusitem{ENV0602}{21;\allowbreak{}\allowbreak{} manuscript\allowbreak{} p.\allowbreak{}35,\allowbreak{}\allowbreak{} prose\allowbreak{} before\allowbreak{} equation(\allowbreak{}5)\allowbreak{}}{it\allowbreak{} remains\allowbreak{} to\allowbreak{} define\allowbreak{} \$\allowbreak{}C\$\allowbreak{}.\allowbreak{}}{it\allowbreak{} remains\allowbreak{} to\allowbreak{} define\allowbreak{} \$\allowbreak{}C(\allowbreak{}I)\allowbreak{}\$\allowbreak{}.\allowbreak{}}{DER\_\allowbreak{}XIII-a:\allowbreak{}23}
\DERapparatusitem{ENV0603}{27;\allowbreak{}\allowbreak{} ms.\allowbreak{}47,\allowbreak{}\allowbreak{} proof\allowbreak{} of\allowbreak{} Proposition\allowbreak{} 1(\allowbreak{}2)\allowbreak{}}{resp.\allowbreak{}\allowbreak{} \$\allowbreak{}\textbackslash{}\allowbreak{}Phi\allowbreak{} =\allowbreak{}\allowbreak{} C(\allowbreak{}I)\allowbreak{}\$\allowbreak{},\allowbreak{}\allowbreak{} \$\allowbreak{}\textbackslash{}\allowbreak{}Psi'\allowbreak{} =\allowbreak{}\allowbreak{} F(\allowbreak{}I')\allowbreak{}\$\allowbreak{}}{resp.\allowbreak{}\allowbreak{} \$\allowbreak{}\textbackslash{}\allowbreak{}Phi\allowbreak{} =\allowbreak{}\allowbreak{} C(\allowbreak{}I)\allowbreak{}\$\allowbreak{},\allowbreak{}\allowbreak{} \$\allowbreak{}\textbackslash{}\allowbreak{}Psi'\allowbreak{} =\allowbreak{}\allowbreak{} TF(\allowbreak{}I')\allowbreak{}\$\allowbreak{}}{DER\_\allowbreak{}XIII-a:\allowbreak{}26}
\DERapparatusitem{ENV0604}{27;\allowbreak{}\allowbreak{} running\allowbreak{} text\allowbreak{} or\allowbreak{} formula}{\textbackslash{}\allowbreak{}text\{\allowbreak{}def\}\allowbreak{}}{\textbackslash{}\allowbreak{}mathrm\{\allowbreak{}def\}\allowbreak{}}{TYPESETTING:\allowbreak{}7;\allowbreak{}\allowbreak{} NOTATION\_\allowbreak{}DER:\allowbreak{}17}
\DERapparatusitem{ENV0605}{29;\allowbreak{}\allowbreak{} ms.\allowbreak{}51,\allowbreak{}\allowbreak{} Example\allowbreak{} 2}{}{\textbackslash{}\allowbreak{}mathcal\allowbreak{} }{DER\_\allowbreak{}XIII-a:\allowbreak{}27}
\DERapparatusitem{ENV0606}{30;\allowbreak{}\allowbreak{} ms.\allowbreak{}52,\allowbreak{}\allowbreak{} nonnegative\allowbreak{} and\allowbreak{} nonpositive\allowbreak{} complexes}{with\allowbreak{} \$\allowbreak{}K\textasciicircum{}\allowbreak{}\{\allowbreak{}i\}\allowbreak{}\allowbreak{} =\allowbreak{}\allowbreak{} 0\$\allowbreak{}\allowbreak{} if\allowbreak{} \$\allowbreak{}i\allowbreak{} \textbackslash{}\allowbreak{}geq\allowbreak{} 0\$\allowbreak{},\allowbreak{}\allowbreak{} resp.\allowbreak{}\allowbreak{} if\allowbreak{} \$\allowbreak{}i\allowbreak{} \textbackslash{}\allowbreak{}leq\allowbreak{} 0\$\allowbreak{}}{with\allowbreak{} \$\allowbreak{}K\textasciicircum{}\allowbreak{}\{\allowbreak{}i\}\allowbreak{}\allowbreak{} =\allowbreak{}\allowbreak{} 0\$\allowbreak{}\allowbreak{} if\allowbreak{} \$\allowbreak{}i\allowbreak{} <\allowbreak{} 0\$\allowbreak{},\allowbreak{}\allowbreak{} resp.\allowbreak{}\allowbreak{} if\allowbreak{} \$\allowbreak{}i\allowbreak{} >\allowbreak{} 0\$\allowbreak{}}{DER\_\allowbreak{}XIII-a:\allowbreak{}29}
\DERapparatusitem{ENV0607}{31;\allowbreak{}\allowbreak{} running\allowbreak{} text\allowbreak{} or\allowbreak{} formula}{\textbackslash{}\allowbreak{}text\{\allowbreak{}def\}\allowbreak{}}{\textbackslash{}\allowbreak{}mathrm\{\allowbreak{}def\}\allowbreak{}}{TYPESETTING:\allowbreak{}7;\allowbreak{}\allowbreak{} NOTATION\_\allowbreak{}DER:\allowbreak{}17}
\DERapparatusitem{ENV0608}{32;\allowbreak{}\allowbreak{} ms.\allowbreak{}56,\allowbreak{}\allowbreak{} final\allowbreak{} paragraph\allowbreak{} before\allowbreak{} section\allowbreak{} 6}{\textbackslash{}\allowbreak{}star}{*}{DER\_\allowbreak{}XIII-a:\allowbreak{}32}
\DERapparatusitem{ENV0609}{32;\allowbreak{}\allowbreak{} ms.\allowbreak{}57,\allowbreak{}\allowbreak{} Definition\allowbreak{} 1(\allowbreak{}a)\allowbreak{},\allowbreak{}\allowbreak{} componentwise\allowbreak{} condition}{\&\allowbreak{}u\_\allowbreak{}\{\allowbreak{}s\}\allowbreak{}\allowbreak{} :\allowbreak{}\allowbreak{} X\_\allowbreak{}\{\allowbreak{}s\}\allowbreak{}\allowbreak{} \textbackslash{}\allowbreak{}longrightarrow\allowbreak{} Y\_\allowbreak{}\{\allowbreak{}s\}\allowbreak{}\allowbreak{} \textbackslash{}\allowbreak{}text\{\allowbreak{}\allowbreak{} is\allowbreak{} in\allowbreak{} \}\allowbreak{}\allowbreak{} C(\allowbreak{}I)\allowbreak{}\allowbreak{} \textbackslash{}\allowbreak{}text\{\allowbreak{}\allowbreak{} (\allowbreak{}resp.\allowbreak{}\allowbreak{} \}\allowbreak{}\allowbreak{} TC(\allowbreak{}I)\allowbreak{})\allowbreak{}}{\&\allowbreak{}u\_\allowbreak{}\{\allowbreak{}s\}\allowbreak{}\allowbreak{} :\allowbreak{}\allowbreak{} X\_\allowbreak{}\{\allowbreak{}s\}\allowbreak{}\allowbreak{} \textbackslash{}\allowbreak{}longrightarrow\allowbreak{} Y\_\allowbreak{}\{\allowbreak{}s\}\allowbreak{}\allowbreak{} \textbackslash{}\allowbreak{}text\{\allowbreak{}\allowbreak{} is\allowbreak{} in\allowbreak{} \}\allowbreak{}\allowbreak{} C\allowbreak{} \textbackslash{}\allowbreak{}text\{\allowbreak{}\allowbreak{} (\allowbreak{}resp.\allowbreak{}\allowbreak{} \}\allowbreak{}\allowbreak{} TC)\allowbreak{}}{DER\_\allowbreak{}XIII-a:\allowbreak{}33}
\DERapparatusitem{ENV0610}{33;\allowbreak{}\allowbreak{} ms.\allowbreak{}59,\allowbreak{}\allowbreak{} parenthetical\allowbreak{} justification\allowbreak{} of\allowbreak{} finite-index\allowbreak{} right\allowbreak{} Kan\allowbreak{} extensions}{\$\allowbreak{}\textbackslash{}\allowbreak{}mathcal\{\allowbreak{}M\}\allowbreak{}\$\allowbreak{}\allowbreak{} being\allowbreak{} assumed\allowbreak{} stable\allowbreak{} under\allowbreak{} finite\allowbreak{} \$\allowbreak{}\textbackslash{}\allowbreak{}varinjlim\$\allowbreak{}}{\$\allowbreak{}\textbackslash{}\allowbreak{}mathcal\{\allowbreak{}M\}\allowbreak{}\$\allowbreak{}\allowbreak{} being\allowbreak{} assumed\allowbreak{} stable\allowbreak{} under\allowbreak{} finite\allowbreak{} \$\allowbreak{}\textbackslash{}\allowbreak{}varprojlim\$\allowbreak{}}{DER\_\allowbreak{}XIII-a:\allowbreak{}34}
\DERapparatusitem{ENV0611}{34;\allowbreak{}\allowbreak{} running\allowbreak{} text\allowbreak{} or\allowbreak{} formula}{\textbackslash{}\allowbreak{}text\{\allowbreak{}def\}\allowbreak{}}{\textbackslash{}\allowbreak{}mathrm\{\allowbreak{}def\}\allowbreak{}}{TYPESETTING:\allowbreak{}7;\allowbreak{}\allowbreak{} NOTATION\_\allowbreak{}DER:\allowbreak{}17}
\DERapparatusitem{ENV0612}{38;\allowbreak{}\allowbreak{} ms.\allowbreak{}\allowbreak{} 68,\allowbreak{}\allowbreak{} right\allowbreak{} column\allowbreak{} of\allowbreak{} localization\allowbreak{} diagram}{\&\allowbreak{}\allowbreak{} \textbackslash{}\allowbreak{}cM(\allowbreak{}I')\allowbreak{}\_\allowbreak{}\{\allowbreak{}\textbackslash{}\allowbreak{}rf\}\allowbreak{}\allowbreak{} \textbackslash{}\allowbreak{}arrow[\allowbreak{}l,\allowbreak{}\allowbreak{} dotted,\allowbreak{}\allowbreak{} "\textbackslash{}\allowbreak{}sim"']\allowbreak{}}{\&\allowbreak{}\allowbreak{} \textbackslash{}\allowbreak{}cM(\allowbreak{}I')\allowbreak{}\_\allowbreak{}\{\allowbreak{}\textbackslash{}\allowbreak{}rf\}\allowbreak{}W(\allowbreak{}I')\allowbreak{}\_\allowbreak{}\{\allowbreak{}\textbackslash{}\allowbreak{}rf\}\allowbreak{}\textasciicircum{}\allowbreak{}\{\allowbreak{}-1\}\allowbreak{}\allowbreak{} \textbackslash{}\allowbreak{}arrow[\allowbreak{}l,\allowbreak{}\allowbreak{} dotted,\allowbreak{}\allowbreak{} "\textbackslash{}\allowbreak{}sim"']\allowbreak{}}{DER\_\allowbreak{}XIII-b:\allowbreak{}4}
\DERapparatusitem{ENV0613}{38;\allowbreak{}\allowbreak{} ms.\allowbreak{}\allowbreak{} 68,\allowbreak{}\allowbreak{} right\allowbreak{} column\allowbreak{} of\allowbreak{} localization\allowbreak{} diagram}{\&\allowbreak{}\allowbreak{} \textbackslash{}\allowbreak{}cM(\allowbreak{}I)\allowbreak{}\_\allowbreak{}\{\allowbreak{}\textbackslash{}\allowbreak{}rf\}\allowbreak{}\allowbreak{} \textbackslash{}\allowbreak{}arrow[\allowbreak{}l,\allowbreak{}\allowbreak{} dotted,\allowbreak{}\allowbreak{} "\textbackslash{}\allowbreak{}sim"']\allowbreak{}}{\&\allowbreak{}\allowbreak{} \textbackslash{}\allowbreak{}cM(\allowbreak{}I)\allowbreak{}\_\allowbreak{}\{\allowbreak{}\textbackslash{}\allowbreak{}rf\}\allowbreak{}W(\allowbreak{}I)\allowbreak{}\_\allowbreak{}\{\allowbreak{}\textbackslash{}\allowbreak{}rf\}\allowbreak{}\textasciicircum{}\allowbreak{}\{\allowbreak{}-1\}\allowbreak{}\allowbreak{} \textbackslash{}\allowbreak{}arrow[\allowbreak{}l,\allowbreak{}\allowbreak{} dotted,\allowbreak{}\allowbreak{} "\textbackslash{}\allowbreak{}sim"']\allowbreak{}}{DER\_\allowbreak{}XIII-b:\allowbreak{}4}
\DERapparatusitem{ENV0614}{38;\allowbreak{}\allowbreak{} ms.\allowbreak{}\allowbreak{} 69,\allowbreak{}\allowbreak{} label\allowbreak{} introducing\allowbreak{} the\allowbreak{} canonical\allowbreak{} comparison}{we\allowbreak{} define\allowbreak{} the\allowbreak{} arrow\allowbreak{} \$\allowbreak{}(\allowbreak{}*\$\allowbreak{}\allowbreak{} bis\$\allowbreak{})\allowbreak{}\$\allowbreak{}}{we\allowbreak{} define\allowbreak{} the\allowbreak{} arrow\allowbreak{} \$\allowbreak{}(\allowbreak{}***)\allowbreak{}\$\allowbreak{}}{DER\_\allowbreak{}XIII-b:\allowbreak{}6}
\DERapparatusitem{ENV0615}{39;\allowbreak{}\allowbreak{} ms.\allowbreak{}\allowbreak{} 70,\allowbreak{}\allowbreak{} isomorphism\allowbreak{} induced\allowbreak{} by\allowbreak{} λ\allowbreak{} [\allowbreak{}sic\allowbreak{} line\allowbreak{} 403]\allowbreak{}}{\textbackslash{}\allowbreak{}Hom\_\allowbreak{}W(\allowbreak{}u\textasciicircum{}\allowbreak{}*F',\allowbreak{}G')\allowbreak{}}{\textbackslash{}\allowbreak{}Hom\_\allowbreak{}W(\allowbreak{}u\textasciicircum{}\allowbreak{}*F,\allowbreak{}G')\allowbreak{}}{DER\_\allowbreak{}XIII-b:\allowbreak{}7}
\DERapparatusitem{ENV0616}{39;\allowbreak{}\allowbreak{} ms.\allowbreak{}\allowbreak{} 71,\allowbreak{}\allowbreak{} numbered\allowbreak{} conclusion\allowbreak{} 1)\allowbreak{}}{admits\allowbreak{} a\allowbreak{} left\allowbreak{} adjoint}{admits\allowbreak{} a\allowbreak{} right\allowbreak{} adjoint}{DER\_\allowbreak{}XIII-b:\allowbreak{}8;\allowbreak{}\allowbreak{} DER\_\allowbreak{}XIII-b:\allowbreak{}9}
\DERapparatusitem{ENV0617}{39;\allowbreak{}\allowbreak{} ms.\allowbreak{}\allowbreak{} 71,\allowbreak{}\allowbreak{} numbered\allowbreak{} conclusion\allowbreak{} 2)\allowbreak{}}{that\allowbreak{} this\allowbreak{} left\allowbreak{} adjoint\allowbreak{} is\allowbreak{} computed}{that\allowbreak{} this\allowbreak{} right\allowbreak{} adjoint\allowbreak{} is\allowbreak{} computed}{DER\_\allowbreak{}XIII-b:\allowbreak{}9}
\DERapparatusitem{ENV0618}{40;\allowbreak{}\allowbreak{} ms.\allowbreak{}\allowbreak{} 72,\allowbreak{}\allowbreak{} Commentaires\allowbreak{} 1)\allowbreak{}a)\allowbreak{},\allowbreak{}\allowbreak{} fibrancy\allowbreak{} shorthand}{fibrant,\allowbreak{}\allowbreak{} \textbackslash{}\allowbreak{}textit\{\allowbreak{}i.\allowbreak{}e.\allowbreak{}\}\allowbreak{}\allowbreak{} in\allowbreak{} \$\allowbreak{}F\$\allowbreak{}}{fibrant,\allowbreak{}\allowbreak{} \textbackslash{}\allowbreak{}textit\{\allowbreak{}i.\allowbreak{}e.\allowbreak{}\}\allowbreak{}\allowbreak{} such\allowbreak{} that\allowbreak{} \$\allowbreak{}X\_\allowbreak{}i\textbackslash{}\allowbreak{}longrightarrow\allowbreak{} e\$\allowbreak{},\allowbreak{}\allowbreak{} \$\allowbreak{}Y\_\allowbreak{}i\textbackslash{}\allowbreak{}longrightarrow\allowbreak{} e\$\allowbreak{}\allowbreak{} are\allowbreak{} in\allowbreak{} \$\allowbreak{}F\$\allowbreak{}}{DER\_\allowbreak{}XIII-b:\allowbreak{}10}
\DERapparatusitem{ENV0619}{40;\allowbreak{}\allowbreak{} ms.\allowbreak{}\allowbreak{} 73,\allowbreak{}\allowbreak{} Commentaires\allowbreak{} 1)\allowbreak{}b)\allowbreak{}\allowbreak{} [\allowbreak{}sic\allowbreak{} line\allowbreak{} 451]\allowbreak{}}{\textbackslash{}\allowbreak{}item[\allowbreak{}\textbackslash{}\allowbreak{}textit\{\allowbreak{}b\}\allowbreak{})\allowbreak{}]\allowbreak{}\allowbreak{} I\allowbreak{} the\allowbreak{} category}{\textbackslash{}\allowbreak{}item[\allowbreak{}\textbackslash{}\allowbreak{}textit\{\allowbreak{}b\}\allowbreak{})\allowbreak{}]\allowbreak{}\allowbreak{} \$\allowbreak{}I\$\allowbreak{}\allowbreak{} the\allowbreak{} category}{DER\_\allowbreak{}XIII-b:\allowbreak{}11}
\DERapparatusitem{ENV0620}{41;\allowbreak{}\allowbreak{} ms.\allowbreak{}\allowbreak{} 74,\allowbreak{}\allowbreak{} displayed\allowbreak{} right\allowbreak{} exactness\allowbreak{} condition}{C}{W}{DER\_\allowbreak{}XIII-b:\allowbreak{}14}
\DERapparatusitem{ENV0621}{41;\allowbreak{}\allowbreak{} ms.\allowbreak{}\allowbreak{} 75,\allowbreak{}\allowbreak{} footnote\allowbreak{} 38}{as\allowbreak{} right\allowbreak{} closed\allowbreak{} subsets\allowbreak{} \textbackslash{}\allowbreak{}texttt\{\allowbreak{}[\allowbreak{}rather\allowbreak{} on\allowbreak{} the\allowbreak{} left]\allowbreak{}\}\allowbreak{}}{as\allowbreak{} left\allowbreak{} closed\allowbreak{} subsets\allowbreak{} \textbackslash{}\allowbreak{}texttt\{\allowbreak{}[\allowbreak{}rather\allowbreak{} on\allowbreak{} the\allowbreak{} left]\allowbreak{}\}\allowbreak{}}{DER\_\allowbreak{}XIII-b:\allowbreak{}15}
\DERapparatusitem{ENV0622}{41;\allowbreak{}\allowbreak{} ms.\allowbreak{}\allowbreak{} 75,\allowbreak{}\allowbreak{} factorization\allowbreak{} sentence\allowbreak{} [\allowbreak{}sic\allowbreak{} line\allowbreak{} 495]\allowbreak{}}{with\allowbreak{} \$\allowbreak{}i\allowbreak{} \textbackslash{}\allowbreak{}in\allowbreak{} TC\$\allowbreak{},\allowbreak{}\allowbreak{} \$\allowbreak{}f\allowbreak{} \textbackslash{}\allowbreak{}in\allowbreak{} F\$\allowbreak{},\allowbreak{}\allowbreak{} and\allowbreak{} also\allowbreak{} with\allowbreak{} \$\allowbreak{}i\allowbreak{} \textbackslash{}\allowbreak{}in\allowbreak{} C\$\allowbreak{},\allowbreak{}\allowbreak{} \$\allowbreak{}f\allowbreak{} \textbackslash{}\allowbreak{}in\allowbreak{} TF\$\allowbreak{}}{with\allowbreak{} \$\allowbreak{}i\allowbreak{} \textbackslash{}\allowbreak{}in\allowbreak{} TC\$\allowbreak{},\allowbreak{}\allowbreak{} \$\allowbreak{}p\allowbreak{} \textbackslash{}\allowbreak{}in\allowbreak{} F\$\allowbreak{},\allowbreak{}\allowbreak{} and\allowbreak{} also\allowbreak{} with\allowbreak{} \$\allowbreak{}i\allowbreak{} \textbackslash{}\allowbreak{}in\allowbreak{} C\$\allowbreak{},\allowbreak{}\allowbreak{} \$\allowbreak{}p\allowbreak{} \textbackslash{}\allowbreak{}in\allowbreak{} TF\$\allowbreak{}}{DER\_\allowbreak{}XIII-b:\allowbreak{}16}
\DERapparatusitem{ENV0623}{42;\allowbreak{}\allowbreak{} ms.\allowbreak{}\allowbreak{} 77,\allowbreak{}\allowbreak{} NB\allowbreak{} under\allowbreak{} the\allowbreak{} generalized\allowbreak{} continuum\allowbreak{} hypothesis}{we\allowbreak{} shall\allowbreak{} always\allowbreak{} have\allowbreak{} \$\allowbreak{}\textbackslash{}\allowbreak{}pi'\textasciicircum{}\allowbreak{}\{\allowbreak{}\textbackslash{}\allowbreak{}pi\}\allowbreak{}\allowbreak{} =\allowbreak{}\allowbreak{} \textbackslash{}\allowbreak{}pi'\$\allowbreak{}}{we\allowbreak{} shall\allowbreak{} have\allowbreak{} \$\allowbreak{}\textbackslash{}\allowbreak{}pi'\textasciicircum{}\allowbreak{}\{\allowbreak{}\textbackslash{}\allowbreak{}pi\}\allowbreak{}\allowbreak{} =\allowbreak{}\allowbreak{} \textbackslash{}\allowbreak{}pi'\$\allowbreak{}\allowbreak{} for\allowbreak{} \$\allowbreak{}\textbackslash{}\allowbreak{}pi<\textbackslash{}\allowbreak{}operatorname\{\allowbreak{}cf\}\allowbreak{}(\allowbreak{}\textbackslash{}\allowbreak{}pi')\allowbreak{}\$\allowbreak{}}{DER\_\allowbreak{}XIII-b:\allowbreak{}17}
\DERapparatusitem{ENV0624}{46;\allowbreak{}\allowbreak{} ms.\allowbreak{}\allowbreak{} 85,\allowbreak{}\allowbreak{} display\allowbreak{} above\allowbreak{} Proposition\allowbreak{} 1\allowbreak{} bis\allowbreak{} [\allowbreak{}sic\allowbreak{} line\allowbreak{} 712]\allowbreak{}}{\textbackslash{}\allowbreak{}varprojlim}{\textbackslash{}\allowbreak{}varinjlim}{DER\_\allowbreak{}XIII-b:\allowbreak{}18}
\DERapparatusitem{ENV0625}{47;\allowbreak{}\allowbreak{} running\allowbreak{} text\allowbreak{} or\allowbreak{} formula}{\textbackslash{}\allowbreak{}text\{\allowbreak{}def\}\allowbreak{}}{\textbackslash{}\allowbreak{}mathrm\{\allowbreak{}def\}\allowbreak{}}{TYPESETTING:\allowbreak{}7;\allowbreak{}\allowbreak{} NOTATION\_\allowbreak{}DER:\allowbreak{}17}
\DERapparatusitem{ENV0626}{47;\allowbreak{}\allowbreak{} ms.\allowbreak{}\allowbreak{} 88,\allowbreak{}\allowbreak{} paragraph\allowbreak{} introducing\allowbreak{} I\_\allowbreak{}M}{For\allowbreak{} every\allowbreak{} direct\allowbreak{} system\allowbreak{} \$\allowbreak{}(\allowbreak{}X\_\allowbreak{}i)\allowbreak{}\_\allowbreak{}\{\allowbreak{}i\textbackslash{}\allowbreak{}in\allowbreak{} I\}\allowbreak{}\$\allowbreak{}\allowbreak{} in\allowbreak{} \$\allowbreak{}M\$\allowbreak{}}{For\allowbreak{} every\allowbreak{} direct\allowbreak{} system\allowbreak{} \$\allowbreak{}(\allowbreak{}X\_\allowbreak{}i)\allowbreak{}\_\allowbreak{}\{\allowbreak{}i\textbackslash{}\allowbreak{}in\allowbreak{} I\}\allowbreak{}\$\allowbreak{}\allowbreak{} in\allowbreak{} \$\allowbreak{}E\$\allowbreak{}}{DER\_\allowbreak{}XIII-b:\allowbreak{}20}
\DERapparatusitem{ENV0627}{48;\allowbreak{}\allowbreak{} ms.\allowbreak{}\allowbreak{} 88,\allowbreak{}\allowbreak{} condition\allowbreak{} b\_\allowbreak{}π}{M}{E}{DER\_\allowbreak{}XIII-b:\allowbreak{}21}
\DERapparatusitem{ENV0628}{49;\allowbreak{}\allowbreak{} ms.\allowbreak{}\allowbreak{} 91,\allowbreak{}\allowbreak{} pullback\allowbreak{} diagram\allowbreak{} [\allowbreak{}sic\allowbreak{} line\allowbreak{} 813]\allowbreak{}}{X}{i}{DER\_\allowbreak{}XIII-b:\allowbreak{}22}
\DERapparatusitem{ENV0629}{49;\allowbreak{}\allowbreak{} ms.\allowbreak{}\allowbreak{} 91,\allowbreak{}\allowbreak{} fibre\allowbreak{} of\allowbreak{} C→I}{at\allowbreak{} \$\allowbreak{}i\textbackslash{}\allowbreak{}in\allowbreak{} I\$\allowbreak{}\allowbreak{} is\allowbreak{} \$\allowbreak{}\textbackslash{}\allowbreak{}mathrm\{\allowbreak{}Filt\}\allowbreak{}\textasciicircum{}\allowbreak{}\textbackslash{}\allowbreak{}pi/\allowbreak{}X\_\allowbreak{}i\$\allowbreak{}}{at\allowbreak{} \$\allowbreak{}i\textbackslash{}\allowbreak{}in\allowbreak{} I\$\allowbreak{}\allowbreak{} is\allowbreak{} \$\allowbreak{}M\textasciicircum{}\allowbreak{}\textbackslash{}\allowbreak{}pi/\allowbreak{}X\_\allowbreak{}i\$\allowbreak{}}{DER\_\allowbreak{}XIII-b:\allowbreak{}23}
\DERapparatusitem{ENV0630}{49;\allowbreak{}\allowbreak{} ms.\allowbreak{}\allowbreak{} 92,\allowbreak{}\allowbreak{} paragraph\allowbreak{} after\allowbreak{} the\allowbreak{} two\allowbreak{} colimit\allowbreak{} expressions}{on\allowbreak{} \$\allowbreak{}I\$\allowbreak{}\allowbreak{} with\allowbreak{} values\allowbreak{} in\allowbreak{} \$\allowbreak{}X\$\allowbreak{},\allowbreak{}}{on\allowbreak{} \$\allowbreak{}I\$\allowbreak{}\allowbreak{} with\allowbreak{} values\allowbreak{} in\allowbreak{} \$\allowbreak{}E\$\allowbreak{},\allowbreak{}}{DER\_\allowbreak{}XIII-b:\allowbreak{}24}
\DERapparatusitem{ENV0631}{50;\allowbreak{}\allowbreak{} ms.\allowbreak{}\allowbreak{} 93,\allowbreak{}\allowbreak{} map\allowbreak{} of\allowbreak{} ordered\allowbreak{} sets}{\textbackslash{}\allowbreak{}mathcal\{\allowbreak{}C\}\allowbreak{}}{I}{DER\_\allowbreak{}XIII-b:\allowbreak{}26}
\DERapparatusitem{ENV0632}{52;\allowbreak{}\allowbreak{} ms.\allowbreak{}\allowbreak{} 97,\allowbreak{}\allowbreak{} final\allowbreak{} implication\allowbreak{} before\allowbreak{} ms.\allowbreak{}\allowbreak{} 98}{then\allowbreak{} \$\allowbreak{}a\_\allowbreak{}\textbackslash{}\allowbreak{}pi\textbackslash{}\allowbreak{}Longrightarrow\allowbreak{} b\_\allowbreak{}\textbackslash{}\allowbreak{}pi\$\allowbreak{}}{then\allowbreak{} \$\allowbreak{}a\_\allowbreak{}\textbackslash{}\allowbreak{}pi\textbackslash{}\allowbreak{}Longrightarrow\allowbreak{} b\_\allowbreak{}\{\allowbreak{}\textbackslash{}\allowbreak{}pi'\}\allowbreak{}\$\allowbreak{}}{DER\_\allowbreak{}XIII-b:\allowbreak{}30}
\DERapparatusitem{ENV0633}{62;\allowbreak{}\allowbreak{} Glossary\allowbreak{} Section\allowbreak{} 6\allowbreak{} [\allowbreak{}sic\allowbreak{} line\allowbreak{} 1261]\allowbreak{}}{weak\allowbreak{} Quillen\allowbreak{} strucure}{weak\allowbreak{} Quillen\allowbreak{} structure}{DER\_\allowbreak{}XIII-b:\allowbreak{}32}
\DERapparatusitem{ENV0634}{63;\allowbreak{}\allowbreak{} Glossary\allowbreak{} Section\allowbreak{} 7,\allowbreak{}\allowbreak{} strict\allowbreak{} cardinal\allowbreak{} filtration}{\$\allowbreak{}\textbackslash{}\allowbreak{}Filt\textasciicircum{}\allowbreak{}\{\allowbreak{}\textbackslash{}\allowbreak{}pi\}\allowbreak{}(\allowbreak{}E)\allowbreak{}\$\allowbreak{}\allowbreak{} is\allowbreak{} stable\allowbreak{} under\allowbreak{} filtering\allowbreak{} direct\allowbreak{} limits.\allowbreak{}}{\$\allowbreak{}\textbackslash{}\allowbreak{}Filt\textasciicircum{}\allowbreak{}\{\allowbreak{}\textbackslash{}\allowbreak{}pi\}\allowbreak{}(\allowbreak{}E)\allowbreak{}\$\allowbreak{}\allowbreak{} is\allowbreak{} stable\allowbreak{} under\allowbreak{} filtering\allowbreak{} direct\allowbreak{} limits\allowbreak{} of\allowbreak{} cardinality\allowbreak{} \$\allowbreak{}\textbackslash{}\allowbreak{}leq\textbackslash{}\allowbreak{}pi\$\allowbreak{}.\allowbreak{}}{DER\_\allowbreak{}XIII-b:\allowbreak{}33}
\DERapparatusitem{ENV0635}{63;\allowbreak{}\allowbreak{} Glossary\allowbreak{} Section\allowbreak{} 7,\allowbreak{}\allowbreak{} perfectly\allowbreak{} adapted\allowbreak{} representation}{a\allowbreak{} \textbackslash{}\allowbreak{}textit\{\allowbreak{}perfectly\allowbreak{} adapted\allowbreak{} representation\}\allowbreak{}\allowbreak{} of\allowbreak{} \$\allowbreak{}X\$\allowbreak{},\allowbreak{}\allowbreak{} in\allowbreak{} terms\allowbreak{} of\allowbreak{} objects\allowbreak{} of\allowbreak{} \$\allowbreak{}\textbackslash{}\allowbreak{}Filt\textasciicircum{}\allowbreak{}\{\allowbreak{}\textbackslash{}\allowbreak{}pi\}\allowbreak{}(\allowbreak{}E)\allowbreak{}\$\allowbreak{},\allowbreak{}\allowbreak{} is\allowbreak{} the\allowbreak{} datum\allowbreak{} of\allowbreak{} a\allowbreak{} well-adapted\allowbreak{} direct\allowbreak{} system}{a\allowbreak{} \textbackslash{}\allowbreak{}textit\{\allowbreak{}perfectly\allowbreak{} adapted\allowbreak{} representation\}\allowbreak{}\allowbreak{} of\allowbreak{} \$\allowbreak{}X\$\allowbreak{},\allowbreak{}\allowbreak{} in\allowbreak{} terms\allowbreak{} of\allowbreak{} objects\allowbreak{} of\allowbreak{} \$\allowbreak{}\textbackslash{}\allowbreak{}Filt\textasciicircum{}\allowbreak{}\{\allowbreak{}\textbackslash{}\allowbreak{}pi\}\allowbreak{}(\allowbreak{}E)\allowbreak{}\$\allowbreak{},\allowbreak{}\allowbreak{} is\allowbreak{} the\allowbreak{} datum\allowbreak{} of\allowbreak{} a\allowbreak{} perfectly\allowbreak{} adapted\allowbreak{} direct\allowbreak{} system}{DER\_\allowbreak{}XIII-b:\allowbreak{}34}
\subsection*{Chapter XV}
\DERapparatusitem{ENV0849}{4;\allowbreak{}\allowbreak{} running\allowbreak{} text\allowbreak{} or\allowbreak{} formula}{\textbackslash{}\allowbreak{}text\{\allowbreak{}def\}\allowbreak{}}{\textbackslash{}\allowbreak{}mathrm\{\allowbreak{}def\}\allowbreak{}}{TYPESETTING:\allowbreak{}7;\allowbreak{}\allowbreak{} NOTATION\_\allowbreak{}DER:\allowbreak{}17}
\DERapparatusitem{ENV0850}{6;\allowbreak{}\allowbreak{} ms.\allowbreak{}7,\allowbreak{}\allowbreak{} Corollary\allowbreak{} 2,\allowbreak{}\allowbreak{} ordinal\allowbreak{} hypothesis}{Let\allowbreak{} \$\allowbreak{}\textbackslash{}\allowbreak{}alpha\$\allowbreak{}\allowbreak{} be\allowbreak{} an\allowbreak{} ordinal\allowbreak{} large\allowbreak{} relative\allowbreak{} to\allowbreak{} \$\allowbreak{}\textbackslash{}\allowbreak{}pi\$\allowbreak{}}{Let\allowbreak{} \$\allowbreak{}\textbackslash{}\allowbreak{}alpha\$\allowbreak{}\allowbreak{} be\allowbreak{} a\allowbreak{} nonzero\allowbreak{} limit\allowbreak{} ordinal\allowbreak{} of\allowbreak{} cofinality\allowbreak{} strictly\allowbreak{} greater\allowbreak{} than\allowbreak{} \$\allowbreak{}\textbackslash{}\allowbreak{}pi\$\allowbreak{}}{DER\_\allowbreak{}XV-a:\allowbreak{}1}
\DERapparatusitem{ENV0851}{7;\allowbreak{}\allowbreak{} ms.\allowbreak{}7,\allowbreak{}\allowbreak{} proof\allowbreak{} of\allowbreak{} Corollary\allowbreak{} 2,\allowbreak{}\allowbreak{} diagram\allowbreak{} label}{\textbackslash{}\allowbreak{}arrow[\allowbreak{}d,\allowbreak{}\allowbreak{} "p\_\allowbreak{}\{\allowbreak{}\textbackslash{}\allowbreak{}alpha+1\}\allowbreak{}(\allowbreak{}f)\allowbreak{}"]\allowbreak{}\allowbreak{} \&\allowbreak{}\allowbreak{} X\_\allowbreak{}\textbackslash{}\allowbreak{}alpha(\allowbreak{}f)\allowbreak{}}{\textbackslash{}\allowbreak{}arrow[\allowbreak{}d,\allowbreak{}\allowbreak{} "p\_\allowbreak{}\{\allowbreak{}\textbackslash{}\allowbreak{}alpha\textasciicircum{}\allowbreak{}\{\allowbreak{}\textbackslash{}\allowbreak{}prime\}\allowbreak{}+1\}\allowbreak{}(\allowbreak{}f)\allowbreak{}"]\allowbreak{}\allowbreak{} \&\allowbreak{}\allowbreak{} X\_\allowbreak{}\textbackslash{}\allowbreak{}alpha(\allowbreak{}f)\allowbreak{}}{DER\_\allowbreak{}XV-a:\allowbreak{}3}
\DERapparatusitem{ENV0852}{7;\allowbreak{}\allowbreak{} ms.\allowbreak{}7,\allowbreak{}\allowbreak{} proof\allowbreak{} of\allowbreak{} Corollary\allowbreak{} 2,\allowbreak{}\allowbreak{} sentence\allowbreak{} beginning\allowbreak{} avec}{\textbackslash{}\allowbreak{}alpha}{\textbackslash{}\allowbreak{}alpha\textasciicircum{}\allowbreak{}\{\allowbreak{}\textbackslash{}\allowbreak{}prime\}\allowbreak{}}{DER\_\allowbreak{}XV-a:\allowbreak{}4}
\DERapparatusitem{ENV0853}{7;\allowbreak{}\allowbreak{} ms.\allowbreak{}8,\allowbreak{}\allowbreak{} definition\allowbreak{} (\allowbreak{}1.\allowbreak{}18)\allowbreak{},\allowbreak{}\allowbreak{} nonempty\allowbreak{} family\allowbreak{} restriction}{of\allowbreak{} a\allowbreak{} \textbackslash{}\allowbreak{}textit\{\allowbreak{}nonempty\}\allowbreak{}\allowbreak{} family\allowbreak{} of\allowbreak{} such\allowbreak{} arrows}{of\allowbreak{} a\allowbreak{} possibly\allowbreak{} empty\allowbreak{} family\allowbreak{} of\allowbreak{} such\allowbreak{} arrows}{DER\_\allowbreak{}XV-a:\allowbreak{}5}
\DERapparatusitem{ENV0854}{9;\allowbreak{}\allowbreak{} ms.\allowbreak{}11,\allowbreak{}\allowbreak{} Theorem\allowbreak{} 1(\allowbreak{}b)\allowbreak{},\allowbreak{}\allowbreak{} ordinal\allowbreak{} hypothesis}{for\allowbreak{} every\allowbreak{} ordinal\allowbreak{} number\allowbreak{} \$\allowbreak{}\textbackslash{}\allowbreak{}alpha\$\allowbreak{}\allowbreak{} large\allowbreak{} relative\allowbreak{} to\allowbreak{} \$\allowbreak{}\textbackslash{}\allowbreak{}pi\$\allowbreak{}}{for\allowbreak{} every\allowbreak{} nonzero\allowbreak{} limit\allowbreak{} ordinal\allowbreak{} \$\allowbreak{}\textbackslash{}\allowbreak{}alpha\$\allowbreak{}\allowbreak{} of\allowbreak{} cofinality\allowbreak{} strictly\allowbreak{} greater\allowbreak{} than\allowbreak{} \$\allowbreak{}\textbackslash{}\allowbreak{}pi\$\allowbreak{}}{DER\_\allowbreak{}XV-a:\allowbreak{}6}
\DERapparatusitem{ENV0855}{11;\allowbreak{}\allowbreak{} running\allowbreak{} text\allowbreak{} or\allowbreak{} formula}{\textbackslash{}\allowbreak{}mathscr\{\allowbreak{}X\}\allowbreak{}}{\textbackslash{}\allowbreak{}mathcal\{\allowbreak{}X\}\allowbreak{}}{TYPESETTING:\allowbreak{}1}
\DERapparatusitem{ENV0856}{11;\allowbreak{}\allowbreak{} running\allowbreak{} text\allowbreak{} or\allowbreak{} formula}{\textbackslash{}\allowbreak{}mathscr\{\allowbreak{}X\}\allowbreak{}}{\textbackslash{}\allowbreak{}mathcal\{\allowbreak{}X\}\allowbreak{}}{TYPESETTING:\allowbreak{}1}
\DERapparatusitem{ENV0857}{11;\allowbreak{}\allowbreak{} running\allowbreak{} text\allowbreak{} or\allowbreak{} formula}{\textbackslash{}\allowbreak{}mathscr\{\allowbreak{}X\}\allowbreak{}}{\textbackslash{}\allowbreak{}mathcal\{\allowbreak{}X\}\allowbreak{}}{TYPESETTING:\allowbreak{}1}
\DERapparatusitem{ENV0858}{11;\allowbreak{}\allowbreak{} running\allowbreak{} text\allowbreak{} or\allowbreak{} formula}{\textbackslash{}\allowbreak{}mathscr\{\allowbreak{}X\}\allowbreak{}}{\textbackslash{}\allowbreak{}mathcal\{\allowbreak{}X\}\allowbreak{}}{TYPESETTING:\allowbreak{}1}
\DERapparatusitem{ENV0859}{14;\allowbreak{}\allowbreak{} ms20,\allowbreak{}\allowbreak{} coloured-set\allowbreak{} example,\allowbreak{}\allowbreak{} bound\allowbreak{} on\allowbreak{} A}{\$\allowbreak{}\textbackslash{}\allowbreak{}card\allowbreak{} A\textbackslash{}\allowbreak{}geq\textbackslash{}\allowbreak{}pi\$\allowbreak{}}{\$\allowbreak{}\textbackslash{}\allowbreak{}card\allowbreak{} A\textbackslash{}\allowbreak{}leq\textbackslash{}\allowbreak{}pi\$\allowbreak{}}{DER\_\allowbreak{}XV-a:\allowbreak{}11}
\DERapparatusitem{ENV0860}{14;\allowbreak{}\allowbreak{} ms20,\allowbreak{}\allowbreak{} display\allowbreak{} bounding\allowbreak{} Hom(\allowbreak{}A\_\allowbreak{}i,\allowbreak{}X)\allowbreak{}}{\textbackslash{}\allowbreak{}Hom(\allowbreak{}A\_\allowbreak{}i,\allowbreak{}X)\allowbreak{}\textbackslash{}\allowbreak{}\allowbreak{} \textbackslash{}\allowbreak{}leq\textbackslash{}\allowbreak{}\allowbreak{} \textbackslash{}\allowbreak{}card(\allowbreak{}J)\allowbreak{}\textbackslash{}\allowbreak{}times\allowbreak{} C\_\allowbreak{}1}{\textbackslash{}\allowbreak{}card\textbackslash{}\allowbreak{}Hom(\allowbreak{}A\_\allowbreak{}i,\allowbreak{}X)\allowbreak{}\textbackslash{}\allowbreak{}\allowbreak{} \textbackslash{}\allowbreak{}leq\textbackslash{}\allowbreak{}\allowbreak{} \textbackslash{}\allowbreak{}card(\allowbreak{}J)\allowbreak{}\textbackslash{}\allowbreak{}times\allowbreak{} C\_\allowbreak{}1}{DER\_\allowbreak{}XV-a:\allowbreak{}12}
\DERapparatusitem{ENV0861}{14;\allowbreak{}\allowbreak{} ms20,\allowbreak{}\allowbreak{} inverse-limit\allowbreak{} index}{\textbackslash{}\allowbreak{}varprojlim\_\allowbreak{}\{\allowbreak{}j\textbackslash{}\allowbreak{}in\allowbreak{} I\}\allowbreak{}\textbackslash{}\allowbreak{}Hom(\allowbreak{}A\_\allowbreak{}i,\allowbreak{}X)\allowbreak{}}{\textbackslash{}\allowbreak{}varprojlim\_\allowbreak{}\{\allowbreak{}i\textbackslash{}\allowbreak{}in\allowbreak{} I\}\allowbreak{}\textbackslash{}\allowbreak{}Hom(\allowbreak{}A\_\allowbreak{}i,\allowbreak{}X)\allowbreak{}}{DER\_\allowbreak{}XV-a:\allowbreak{}13}
\DERapparatusitem{ENV0862}{15;\allowbreak{}\allowbreak{} ms21\allowbreak{} and\allowbreak{} diagram-count\allowbreak{} label\allowbreak{} on\allowbreak{} ms22}{\textbackslash{}\allowbreak{}textit\{\allowbreak{}objects\}\allowbreak{}\allowbreak{} \$\allowbreak{}A\_\allowbreak{}i\$\allowbreak{}\allowbreak{} in\allowbreak{} \$\allowbreak{}\textbackslash{}\allowbreak{}M\textasciicircum{}\allowbreak{}\textbackslash{}\allowbreak{}pi\$\allowbreak{}}{\textbackslash{}\allowbreak{}textit\{\allowbreak{}objects\}\allowbreak{}\allowbreak{} \$\allowbreak{}A\_\allowbreak{}i\$\allowbreak{}\allowbreak{} in\allowbreak{} \$\allowbreak{}\textbackslash{}\allowbreak{}M\textasciicircum{}\allowbreak{}\{\allowbreak{}\textbackslash{}\allowbreak{}pi\_\allowbreak{}0\}\allowbreak{}\$\allowbreak{}}{DER\_\allowbreak{}XV-a:\allowbreak{}14}
\DERapparatusitem{ENV0863}{15;\allowbreak{}\allowbreak{} ms21\allowbreak{} and\allowbreak{} diagram-count\allowbreak{} label\allowbreak{} on\allowbreak{} ms22}{in\allowbreak{} \$\allowbreak{}\textbackslash{}\allowbreak{}M\textasciicircum{}\allowbreak{}\textbackslash{}\allowbreak{}pi\$\allowbreak{},\allowbreak{}\allowbreak{} up\allowbreak{} to\textbackslash{}\allowbreak{}\textbackslash{}\allowbreak{}\allowbreak{} isomorphism}{in\allowbreak{} \$\allowbreak{}\textbackslash{}\allowbreak{}M\textasciicircum{}\allowbreak{}\{\allowbreak{}\textbackslash{}\allowbreak{}pi\_\allowbreak{}0\}\allowbreak{}\$\allowbreak{},\allowbreak{}\allowbreak{} up\allowbreak{} to\textbackslash{}\allowbreak{}\textbackslash{}\allowbreak{}\allowbreak{} isomorphism}{DER\_\allowbreak{}XV-a:\allowbreak{}14}
\DERapparatusitem{ENV0864}{15;\allowbreak{}\allowbreak{} ms22,\allowbreak{}\allowbreak{} sentence\allowbreak{} before\allowbreak{} the\allowbreak{} final\allowbreak{} slice-arrow\allowbreak{} count}{\$\allowbreak{}\textbackslash{}\allowbreak{}wavy\{\allowbreak{}\textbackslash{}\allowbreak{}uFl\textasciicircum{}\allowbreak{}r\textbackslash{}\allowbreak{},\allowbreak{}\textbackslash{}\allowbreak{}M\}\allowbreak{}\$\allowbreak{}\allowbreak{} \textbackslash{}\allowbreak{}\allowbreak{} (\allowbreak{}3.\allowbreak{}9)\allowbreak{}}{\$\allowbreak{}\textbackslash{}\allowbreak{}wavy\{\allowbreak{}\textbackslash{}\allowbreak{}uFl(\allowbreak{}\textbackslash{}\allowbreak{}M\textasciicircum{}\allowbreak{}r/\allowbreak{}X)\allowbreak{}\}\allowbreak{}\$\allowbreak{}\allowbreak{} \textbackslash{}\allowbreak{}\allowbreak{} (\allowbreak{}3.\allowbreak{}9)\allowbreak{}}{DER\_\allowbreak{}XV-a:\allowbreak{}15}
\DERapparatusitem{ENV0865}{16;\allowbreak{}\allowbreak{} ms23,\allowbreak{}\allowbreak{} Corollary\allowbreak{} 1}{\$\allowbreak{}A,\allowbreak{}X\textbackslash{}\allowbreak{}in\textbackslash{}\allowbreak{}Filt\textasciicircum{}\allowbreak{}\textbackslash{}\allowbreak{}pi\$\allowbreak{}}{\$\allowbreak{}A,\allowbreak{}X\textbackslash{}\allowbreak{}in\textbackslash{}\allowbreak{}Filt\textasciicircum{}\allowbreak{}r\$\allowbreak{}}{DER\_\allowbreak{}XV-a:\allowbreak{}16}
\DERapparatusitem{ENV0866}{18;\allowbreak{}\allowbreak{} ms27,\allowbreak{}\allowbreak{} definition\allowbreak{} of\allowbreak{} c\_\allowbreak{}1\allowbreak{} in\allowbreak{} Theorem\allowbreak{} 2}{c\_\allowbreak{}1\textbackslash{}\allowbreak{}\allowbreak{} =\allowbreak{}\textbackslash{}\allowbreak{}\allowbreak{} \textbackslash{}\allowbreak{}card\textbackslash{}\allowbreak{}Fl(\allowbreak{}\textbackslash{}\allowbreak{}M\textasciicircum{}\allowbreak{}\{\allowbreak{}\textbackslash{}\allowbreak{}pi\_\allowbreak{}1\}\allowbreak{})\allowbreak{}\textbackslash{}\allowbreak{}\allowbreak{} =\allowbreak{}\textbackslash{}\allowbreak{}\allowbreak{} \textbackslash{}\allowbreak{}card\textbackslash{}\allowbreak{}Fl\textasciicircum{}\allowbreak{}\{\allowbreak{}\textbackslash{}\allowbreak{}pi\_\allowbreak{}1\}\allowbreak{}(\allowbreak{}\textbackslash{}\allowbreak{}M)\allowbreak{}}{c\_\allowbreak{}1\textbackslash{}\allowbreak{}\allowbreak{} =\allowbreak{}\textbackslash{}\allowbreak{}\allowbreak{} \textbackslash{}\allowbreak{}card\textbackslash{}\allowbreak{}Fl(\allowbreak{}K\_\allowbreak{}1)\allowbreak{},\allowbreak{}\textbackslash{}\allowbreak{}qquad\allowbreak{} K\_\allowbreak{}1\textbackslash{}\allowbreak{}subset\textbackslash{}\allowbreak{}M\textasciicircum{}\allowbreak{}\{\allowbreak{}\textbackslash{}\allowbreak{}pi\_\allowbreak{}1\}\allowbreak{}\textbackslash{}\allowbreak{}\allowbreak{} \textbackslash{}\allowbreak{}text\{\allowbreak{}a\allowbreak{} chosen\allowbreak{} small\allowbreak{} full\allowbreak{} skeleton\}\allowbreak{}}{DER\_\allowbreak{}XV-a:\allowbreak{}18}
\DERapparatusitem{ENV0867}{20;\allowbreak{}\allowbreak{} ms30,\allowbreak{}\allowbreak{} proof\allowbreak{} of\allowbreak{} Proposition4.\allowbreak{}2,\allowbreak{}\allowbreak{} (\allowbreak{}b)\allowbreak{}⇒(\allowbreak{}d)\allowbreak{}}{the\allowbreak{} given\allowbreak{} diagram\allowbreak{} (\allowbreak{}4.\allowbreak{}2.\allowbreak{}2)\allowbreak{}\allowbreak{} (\allowbreak{}solid\allowbreak{} arrows)\allowbreak{}}{the\allowbreak{} given\allowbreak{} diagram\allowbreak{} (\allowbreak{}4.\allowbreak{}2.\allowbreak{}3)\allowbreak{}\allowbreak{} (\allowbreak{}solid\allowbreak{} arrows)\allowbreak{}}{DER\_\allowbreak{}XV-a:\allowbreak{}21}
\DERapparatusitem{ENV0868}{23;\allowbreak{}\allowbreak{} ms34,\allowbreak{}\allowbreak{} proof\allowbreak{} of\allowbreak{} Proposition4.\allowbreak{}3,\allowbreak{}\allowbreak{} paragraph\allowbreak{} after\allowbreak{} triangular\allowbreak{} diagram}{(\allowbreak{}prop.\allowbreak{}\allowbreak{} 3.\allowbreak{}2,\allowbreak{}}{(\allowbreak{}prop.\allowbreak{}\allowbreak{} 4.\allowbreak{}2,\allowbreak{}}{DER\_\allowbreak{}XV-a:\allowbreak{}22}
\DERapparatusitem{ENV0869}{preamble;\allowbreak{}\allowbreak{} local\allowbreak{} macro\allowbreak{} definition}{\textbackslash{}\allowbreak{}text\{\allowbreak{}def\}\allowbreak{}}{\textbackslash{}\allowbreak{}mathrm\{\allowbreak{}def\}\allowbreak{}}{TYPESETTING:\allowbreak{}7;\allowbreak{}\allowbreak{} NOTATION\_\allowbreak{}DER:\allowbreak{}17}
\DERapparatusitem{ENV0870}{24;\allowbreak{}\allowbreak{} ms35,\allowbreak{}\allowbreak{} proof\allowbreak{} of\allowbreak{} Proposition4.\allowbreak{}3\allowbreak{} Corollary2}{But\allowbreak{} since\allowbreak{} \$\allowbreak{}C\textbackslash{}\allowbreak{}cap\allowbreak{} W\textbackslash{}\allowbreak{}longleftrightarrow\allowbreak{} TC\$\allowbreak{}\allowbreak{} and\allowbreak{} \$\allowbreak{}C\textbackslash{}\allowbreak{}cap\allowbreak{} W\$\allowbreak{}\allowbreak{} is\allowbreak{} stable}{But\allowbreak{} since\allowbreak{} \$\allowbreak{}C\textbackslash{}\allowbreak{}cap\allowbreak{} W\textbackslash{}\allowbreak{}longleftrightarrow\allowbreak{} F\$\allowbreak{},\allowbreak{}\allowbreak{} \$\allowbreak{}(\allowbreak{}C\textbackslash{}\allowbreak{}cap\allowbreak{} W,\allowbreak{}F)\allowbreak{}\$\allowbreak{}\allowbreak{} is\allowbreak{} a\allowbreak{} Quillen\allowbreak{} precouple,\allowbreak{}\allowbreak{} and\allowbreak{} \$\allowbreak{}C\textbackslash{}\allowbreak{}cap\allowbreak{} W\$\allowbreak{}\allowbreak{} is\allowbreak{} stable}{DER\_\allowbreak{}XV-a:\allowbreak{}23}
\DERapparatusitem{ENV0871}{24;\allowbreak{}\allowbreak{} ms36,\allowbreak{}\allowbreak{} footnote13\allowbreak{} Corollary4}{two\allowbreak{} Quillen\allowbreak{} precouples\allowbreak{} such\allowbreak{} that}{two\allowbreak{} closed\allowbreak{} Quillen\allowbreak{} couples\allowbreak{} such\allowbreak{} that}{DER\_\allowbreak{}XV-a:\allowbreak{}24}
\DERapparatusitem{ENV0872}{25;\allowbreak{}\allowbreak{} ms37,\allowbreak{}\allowbreak{} first\allowbreak{} paragraph\allowbreak{} of4.\allowbreak{}4}{\textbackslash{}\allowbreak{}subset\textbackslash{}\allowbreak{}mathcal\{\allowbreak{}M\}\allowbreak{}}{\textbackslash{}\allowbreak{}subset\textbackslash{}\allowbreak{}Fl(\allowbreak{}\textbackslash{}\allowbreak{}mathcal\allowbreak{} M)\allowbreak{}}{DER\_\allowbreak{}XV-a:\allowbreak{}25}
\DERapparatusitem{ENV0873}{26;\allowbreak{}\allowbreak{} ms38,\allowbreak{}\allowbreak{} first\allowbreak{} paragraph}{by\allowbreak{} proposition\allowbreak{} 4.\allowbreak{}1,\allowbreak{}\allowbreak{} corollary\allowbreak{} 2}{by\allowbreak{} proposition\allowbreak{} 4.\allowbreak{}2,\allowbreak{}\allowbreak{} corollary\allowbreak{} 2}{DER\_\allowbreak{}XV-a:\allowbreak{}26}
\DERapparatusitem{ENV0874}{27;\allowbreak{}\allowbreak{} ms40,\allowbreak{}\allowbreak{} Corollary4.\allowbreak{}4.\allowbreak{}3\allowbreak{} hypothesis\allowbreak{} and\allowbreak{} equality}{\textbackslash{}\allowbreak{}textbf\{\allowbreak{}Corollary\allowbreak{} 4.\allowbreak{}4.\allowbreak{}3.\allowbreak{}\}\allowbreak{}\allowbreak{} \textbackslash{}\allowbreak{}textit\{\allowbreak{}Under\allowbreak{} the\allowbreak{} conditions\allowbreak{} of\allowbreak{} proposition\}\allowbreak{}\allowbreak{} 4.\allowbreak{}4.\allowbreak{}2,\allowbreak{}}{\textbackslash{}\allowbreak{}textbf\{\allowbreak{}Corollary\allowbreak{} 4.\allowbreak{}4.\allowbreak{}3.\allowbreak{}\}\allowbreak{}\allowbreak{} \textbackslash{}\allowbreak{}textit\{\allowbreak{}Under\allowbreak{} the\allowbreak{} conditions\allowbreak{} of\allowbreak{} proposition\}\allowbreak{}\allowbreak{} 4.\allowbreak{}4.\allowbreak{}2,\allowbreak{}\allowbreak{} \textbackslash{}\allowbreak{}textit\{\allowbreak{}and\allowbreak{} assuming\allowbreak{} that\allowbreak{} every\allowbreak{} arrow\allowbreak{} factors\allowbreak{} as\allowbreak{} a\allowbreak{} \$\allowbreak{}W\$\allowbreak{}-cofibration\allowbreak{} followed\allowbreak{} by\allowbreak{} a\allowbreak{} \$\allowbreak{}W\$\allowbreak{}-equivalence,\allowbreak{}\}\allowbreak{}}{DER\_\allowbreak{}XV-a:\allowbreak{}27}
\DERapparatusitem{ENV0875}{27;\allowbreak{}\allowbreak{} ms40,\allowbreak{}\allowbreak{} Corollary4.\allowbreak{}4.\allowbreak{}4\allowbreak{} display}{=\allowbreak{}}{\textbackslash{}\allowbreak{}subset}{DER\_\allowbreak{}XV-a:\allowbreak{}28}
\DERapparatusitem{ENV0876}{28;\allowbreak{}\allowbreak{} ms42,\allowbreak{}\allowbreak{} Proposition4.\allowbreak{}4.\allowbreak{}6\allowbreak{} before\allowbreak{} the\allowbreak{} first\allowbreak{} occurrence\allowbreak{} of\allowbreak{} F0}{and\allowbreak{} let\allowbreak{} \$\allowbreak{}TC\_\allowbreak{}0\$\allowbreak{}\allowbreak{} be\allowbreak{} a\allowbreak{} small\allowbreak{} subset\allowbreak{} of\allowbreak{} \$\allowbreak{}TC=\allowbreak{}C\textbackslash{}\allowbreak{}cap\allowbreak{} W\$\allowbreak{}.\allowbreak{}\allowbreak{} We\allowbreak{} then\allowbreak{} have}{and\allowbreak{} let\allowbreak{} \$\allowbreak{}TC\_\allowbreak{}0\$\allowbreak{}\allowbreak{} be\allowbreak{} a\allowbreak{} small\allowbreak{} subset\allowbreak{} of\allowbreak{} \$\allowbreak{}TC=\allowbreak{}C\textbackslash{}\allowbreak{}cap\allowbreak{} W\$\allowbreak{}.\allowbreak{}\allowbreak{} Put\allowbreak{} \$\allowbreak{}F\_\allowbreak{}0=\allowbreak{}(\allowbreak{}TC\_\allowbreak{}0)\allowbreak{}\_\allowbreak{}*\$\allowbreak{}.\allowbreak{}\allowbreak{} We\allowbreak{} then\allowbreak{} have}{DER\_\allowbreak{}XV-a:\allowbreak{}29}
\DERapparatusitem{ENV0877}{29;\allowbreak{}\allowbreak{} Lemma\allowbreak{} 5.\allowbreak{}3,\allowbreak{}\allowbreak{} ms45}{Then\allowbreak{} \$\allowbreak{}g\$\allowbreak{}\allowbreak{} is\allowbreak{} an\allowbreak{} epimorphism\allowbreak{} if\allowbreak{} and\allowbreak{} only\allowbreak{} if\allowbreak{} \$\allowbreak{}g\$\allowbreak{}\allowbreak{} does\allowbreak{} not\allowbreak{} factor}{Then\allowbreak{} \$\allowbreak{}f\$\allowbreak{}\allowbreak{} is\allowbreak{} an\allowbreak{} epimorphism\allowbreak{} if\allowbreak{} and\allowbreak{} only\allowbreak{} if\allowbreak{} \$\allowbreak{}f\$\allowbreak{}\allowbreak{} does\allowbreak{} not\allowbreak{} factor}{DER\_\allowbreak{}XV-a:\allowbreak{}31}
\DERapparatusitem{ENV0878}{30;\allowbreak{}\allowbreak{} Opening\allowbreak{} paragraph,\allowbreak{}\allowbreak{} ms46}{at\allowbreak{} issue\allowbreak{} in\allowbreak{} (\allowbreak{}\$\allowbreak{}d\_\allowbreak{}2\$\allowbreak{})\allowbreak{}.\allowbreak{}}{at\allowbreak{} issue\allowbreak{} in\allowbreak{} (\allowbreak{}\$\allowbreak{}d\_\allowbreak{}2\$\allowbreak{})\allowbreak{})\allowbreak{}.\allowbreak{}}{DER\_\allowbreak{}XV-a:\allowbreak{}32}
\DERapparatusitem{ENV0879}{30;\allowbreak{}\allowbreak{} Proof\allowbreak{} of\allowbreak{} Proposition\allowbreak{} 5.\allowbreak{}3,\allowbreak{}\allowbreak{} ms47}{we\allowbreak{} have\allowbreak{} \$\allowbreak{}Z\allowbreak{} \textbackslash{}\allowbreak{}in\allowbreak{} \textbackslash{}\allowbreak{}Filt\textasciicircum{}\allowbreak{}\{\allowbreak{}\textbackslash{}\allowbreak{}pi\}\allowbreak{}\$\allowbreak{}\allowbreak{} by\allowbreak{} (\allowbreak{}\$\allowbreak{}c\$\allowbreak{})\allowbreak{}}{we\allowbreak{} have\allowbreak{} \$\allowbreak{}Z\allowbreak{} \textbackslash{}\allowbreak{}in\allowbreak{} \textbackslash{}\allowbreak{}Filt\textasciicircum{}\allowbreak{}\{\allowbreak{}\textbackslash{}\allowbreak{}pi\}\allowbreak{}\$\allowbreak{}\allowbreak{} by\allowbreak{} (\allowbreak{}\$\allowbreak{}d\_\allowbreak{}0\$\allowbreak{})\allowbreak{}}{DER\_\allowbreak{}XV-a:\allowbreak{}33}
\DERapparatusitem{ENV0880}{30;\allowbreak{}\allowbreak{} Proof\allowbreak{} of\allowbreak{} Proposition\allowbreak{} 5.\allowbreak{}3,\allowbreak{}\allowbreak{} ms47}{which\allowbreak{} is\allowbreak{} in\allowbreak{} \$\allowbreak{}\textbackslash{}\allowbreak{}Filt\textasciicircum{}\allowbreak{}\{\allowbreak{}\textbackslash{}\allowbreak{}pi\}\allowbreak{}\$\allowbreak{}\allowbreak{} by\allowbreak{} (\allowbreak{}\$\allowbreak{}d\_\allowbreak{}1\$\allowbreak{})\allowbreak{}}{which\allowbreak{} is\allowbreak{} in\allowbreak{} \$\allowbreak{}\textbackslash{}\allowbreak{}Filt\textasciicircum{}\allowbreak{}\{\allowbreak{}\textbackslash{}\allowbreak{}pi\}\allowbreak{}\$\allowbreak{}\allowbreak{} by\allowbreak{} (\allowbreak{}\$\allowbreak{}d\_\allowbreak{}2\$\allowbreak{})\allowbreak{}}{DER\_\allowbreak{}XV-a:\allowbreak{}34}
\DERapparatusitem{ENV0881}{32;\allowbreak{}\allowbreak{} running\allowbreak{} text\allowbreak{} or\allowbreak{} formula}{\textbackslash{}\allowbreak{}text\{\allowbreak{}def\}\allowbreak{}}{\textbackslash{}\allowbreak{}mathrm\{\allowbreak{}def\}\allowbreak{}}{TYPESETTING:\allowbreak{}7;\allowbreak{}\allowbreak{} NOTATION\_\allowbreak{}DER:\allowbreak{}17}
\DERapparatusitem{ENV0882}{34;\allowbreak{}\allowbreak{} Proposition\allowbreak{} 5.\allowbreak{}7,\allowbreak{}\allowbreak{} opening,\allowbreak{}\allowbreak{} ms57}{\textbackslash{}\allowbreak{}M}{M}{DER\_\allowbreak{}XV-a:\allowbreak{}35;\allowbreak{}\allowbreak{} DER\_\allowbreak{}XV-b:\allowbreak{}1}
\DERapparatusitem{ENV0883}{34;\allowbreak{}\allowbreak{} Proposition\allowbreak{} 5.\allowbreak{}7,\allowbreak{}\allowbreak{} condition\allowbreak{} (\allowbreak{}c)\allowbreak{},\allowbreak{}\allowbreak{} ms57}{we\allowbreak{} have\allowbreak{} \$\allowbreak{}X\allowbreak{} \textbackslash{}\allowbreak{}in\allowbreak{} M\$\allowbreak{}\allowbreak{} if\allowbreak{} and\allowbreak{} only\allowbreak{} if\allowbreak{} the\allowbreak{} set\allowbreak{} \$\allowbreak{}J\_\allowbreak{}M\$\allowbreak{}}{putting\allowbreak{} \$\allowbreak{}J=\allowbreak{}\textbackslash{}\allowbreak{}Ss\_\allowbreak{}\textbackslash{}\allowbreak{}pi(\allowbreak{}X)\allowbreak{}\$\allowbreak{},\allowbreak{}\allowbreak{} we\allowbreak{} have\allowbreak{} \$\allowbreak{}X\allowbreak{} \textbackslash{}\allowbreak{}in\allowbreak{} M\$\allowbreak{}\allowbreak{} if\allowbreak{} and\allowbreak{} only\allowbreak{} if\allowbreak{} the\allowbreak{} set\allowbreak{} \$\allowbreak{}J\_\allowbreak{}M\$\allowbreak{}}{DER\_\allowbreak{}XV-a:\allowbreak{}36}
\DERapparatusitem{ENV0884}{35;\allowbreak{}\allowbreak{} Proposition\allowbreak{} 5.\allowbreak{}7,\allowbreak{}\allowbreak{} condition\allowbreak{} (\allowbreak{}cprime)\allowbreak{},\allowbreak{}\allowbreak{} ms58}{\$\allowbreak{}M\$\allowbreak{}\allowbreak{} is\allowbreak{} the\allowbreak{} smallest\allowbreak{} subset\allowbreak{} of\allowbreak{} \$\allowbreak{}\textbackslash{}\allowbreak{}M\$\allowbreak{}\allowbreak{} stable\allowbreak{} under\allowbreak{} filtering\allowbreak{} direct\allowbreak{} limits}{\$\allowbreak{}M\$\allowbreak{}\allowbreak{} is\allowbreak{} the\allowbreak{} smallest\allowbreak{} subset\allowbreak{} of\allowbreak{} \$\allowbreak{}\textbackslash{}\allowbreak{}M\$\allowbreak{}\allowbreak{} containing\allowbreak{} \$\allowbreak{}M\_\allowbreak{}0\$\allowbreak{}\allowbreak{} and\allowbreak{} stable\allowbreak{} under\allowbreak{} filtering\allowbreak{} direct\allowbreak{} limits}{DER\_\allowbreak{}XV-a:\allowbreak{}37;\allowbreak{}\allowbreak{} DER\_\allowbreak{}XV-b:\allowbreak{}2}
\DERapparatusitem{ENV0885}{35;\allowbreak{}\allowbreak{} Corollary\allowbreak{} 5.\allowbreak{}8,\allowbreak{}\allowbreak{} inverse\allowbreak{} construction\allowbreak{} parenthesis,\allowbreak{}\allowbreak{} ms58}{smallest\allowbreak{} subset\allowbreak{} of\allowbreak{} \$\allowbreak{}\textbackslash{}\allowbreak{}Ob\textbackslash{}\allowbreak{}M\$\allowbreak{}\allowbreak{} stable\allowbreak{} under\allowbreak{} filtering\allowbreak{} \$\allowbreak{}\textbackslash{}\allowbreak{}varinjlim\$\allowbreak{}}{smallest\allowbreak{} subset\allowbreak{} of\allowbreak{} \$\allowbreak{}\textbackslash{}\allowbreak{}Ob\textbackslash{}\allowbreak{}M\$\allowbreak{}\allowbreak{} containing\allowbreak{} \$\allowbreak{}M\textasciicircum{}\allowbreak{}\textbackslash{}\allowbreak{}pi\$\allowbreak{}\allowbreak{} and\allowbreak{} stable\allowbreak{} under\allowbreak{} filtering\allowbreak{} \$\allowbreak{}\textbackslash{}\allowbreak{}varinjlim\$\allowbreak{}}{DER\_\allowbreak{}XV-a:\allowbreak{}38;\allowbreak{}\allowbreak{} DER\_\allowbreak{}XV-b:\allowbreak{}3}
\DERapparatusitem{ENV0886}{35;\allowbreak{}\allowbreak{} Proof\allowbreak{} of\allowbreak{} Proposition\allowbreak{} 5.\allowbreak{}7,\allowbreak{}\allowbreak{} ms59}{}{,\allowbreak{}}{DER\_\allowbreak{}XV-a:\allowbreak{}39;\allowbreak{}\allowbreak{} DER\_\allowbreak{}XV-b:\allowbreak{}4}
\DERapparatusitem{ENV0887}{35;\allowbreak{}\allowbreak{} Proof\allowbreak{} of\allowbreak{} Proposition\allowbreak{} 5.\allowbreak{}7,\allowbreak{}\allowbreak{} ms59}{if\allowbreak{} all\allowbreak{} the\allowbreak{} objects\allowbreak{} of\allowbreak{} \$\allowbreak{}\textbackslash{}\allowbreak{}Filt\textasciicircum{}\allowbreak{}\textbackslash{}\allowbreak{}pi\$\allowbreak{}\allowbreak{} are\allowbreak{} accessible}{if\allowbreak{} all\allowbreak{} the\allowbreak{} objects\allowbreak{} of\allowbreak{} \$\allowbreak{}\textbackslash{}\allowbreak{}Filt\textasciicircum{}\allowbreak{}\textbackslash{}\allowbreak{}pi\$\allowbreak{}\allowbreak{} are\allowbreak{} \$\allowbreak{}\textbackslash{}\allowbreak{}pi\$\allowbreak{}-accessible}{DER\_\allowbreak{}XV-a:\allowbreak{}40;\allowbreak{}\allowbreak{} DER\_\allowbreak{}XV-b:\allowbreak{}5}
\DERapparatusitem{ENV0888}{36;\allowbreak{}\allowbreak{} running\allowbreak{} text\allowbreak{} or\allowbreak{} formula}{\textbackslash{}\allowbreak{}text\{\allowbreak{}def\}\allowbreak{}}{\textbackslash{}\allowbreak{}mathrm\{\allowbreak{}def\}\allowbreak{}}{TYPESETTING:\allowbreak{}7;\allowbreak{}\allowbreak{} NOTATION\_\allowbreak{}DER:\allowbreak{}17}
\DERapparatusitem{ENV0889}{37;\allowbreak{}\allowbreak{} Lemma\allowbreak{} 5.\allowbreak{}10\allowbreak{} and\allowbreak{} first\allowbreak{} display\allowbreak{} of\allowbreak{} proof,\allowbreak{}\allowbreak{} ms62–63}{and\allowbreak{} \$\allowbreak{}\textbackslash{}\allowbreak{}Ss\_\allowbreak{}\textbackslash{}\allowbreak{}pi(\allowbreak{}Y)\allowbreak{}\$\allowbreak{}\allowbreak{} is\allowbreak{} cofinal\allowbreak{} in\allowbreak{} it}{and\allowbreak{} \$\allowbreak{}\textbackslash{}\allowbreak{}Ss\_\allowbreak{}\textbackslash{}\allowbreak{}pi\textasciicircum{}\allowbreak{}\{\allowbreak{}M\textasciicircum{}\allowbreak{}\textbackslash{}\allowbreak{}pi\}\allowbreak{}(\allowbreak{}Y)\allowbreak{}\$\allowbreak{}\allowbreak{} is\allowbreak{} cofinal\allowbreak{} in\allowbreak{} it}{DER\_\allowbreak{}XV-a:\allowbreak{}42;\allowbreak{}\allowbreak{} DER\_\allowbreak{}XV-b:\allowbreak{}7}
\DERapparatusitem{ENV0890}{37;\allowbreak{}\allowbreak{} Lemma\allowbreak{} 5.\allowbreak{}10\allowbreak{} and\allowbreak{} first\allowbreak{} display\allowbreak{} of\allowbreak{} proof,\allowbreak{}\allowbreak{} ms62–63}{I\allowbreak{} \textbackslash{}\allowbreak{};\allowbreak{}=\allowbreak{}\textbackslash{}\allowbreak{};\allowbreak{}\allowbreak{} \textbackslash{}\allowbreak{}Ss\_\allowbreak{}\textbackslash{}\allowbreak{}pi(\allowbreak{}Y)\allowbreak{}\allowbreak{} \textbackslash{}\allowbreak{}xrightarrow}{I\allowbreak{} \textbackslash{}\allowbreak{};\allowbreak{}=\allowbreak{}\textbackslash{}\allowbreak{};\allowbreak{}\allowbreak{} \textbackslash{}\allowbreak{}Ss\_\allowbreak{}\textbackslash{}\allowbreak{}pi\textasciicircum{}\allowbreak{}\{\allowbreak{}M\textasciicircum{}\allowbreak{}\textbackslash{}\allowbreak{}pi\}\allowbreak{}(\allowbreak{}Y)\allowbreak{}\allowbreak{} \textbackslash{}\allowbreak{}xrightarrow}{DER\_\allowbreak{}XV-a:\allowbreak{}42;\allowbreak{}\allowbreak{} DER\_\allowbreak{}XV-b:\allowbreak{}8}
\DERapparatusitem{ENV0891}{37;\allowbreak{}\allowbreak{} Corollary\allowbreak{} 5.\allowbreak{}11,\allowbreak{}\allowbreak{} parenthetical\allowbreak{} adjective,\allowbreak{}\allowbreak{} ms64}{(\allowbreak{}\textbackslash{}\allowbreak{}textit\{\allowbreak{}fibrant\}\allowbreak{})\allowbreak{}}{(\allowbreak{}\textbackslash{}\allowbreak{}textit\{\allowbreak{}cofibrant\}\allowbreak{})\allowbreak{}}{DER\_\allowbreak{}XV-a:\allowbreak{}44;\allowbreak{}\allowbreak{} DER\_\allowbreak{}XV-b:\allowbreak{}10}
\DERapparatusitem{ENV0892}{40;\allowbreak{}\allowbreak{} ms71\allowbreak{} proof5.\allowbreak{}15,\allowbreak{}\allowbreak{} line\allowbreak{} before\allowbreak{} ms72\allowbreak{} marker}{\$\allowbreak{}i\_\allowbreak{}\{\allowbreak{}\textbackslash{}\allowbreak{}beta\}\allowbreak{}\$\allowbreak{}\allowbreak{} is\allowbreak{} in\allowbreak{} each\allowbreak{} \$\allowbreak{}I\_\allowbreak{}\{\allowbreak{}\textbackslash{}\allowbreak{}beta\_\allowbreak{}\{\allowbreak{}0\}\allowbreak{}\}\allowbreak{}\$\allowbreak{}}{\$\allowbreak{}i\_\allowbreak{}\{\allowbreak{}\textbackslash{}\allowbreak{}beta\}\allowbreak{}\$\allowbreak{}\allowbreak{} is\allowbreak{} in\allowbreak{} each\allowbreak{} \$\allowbreak{}J\_\allowbreak{}\{\allowbreak{}\textbackslash{}\allowbreak{}beta\_\allowbreak{}\{\allowbreak{}0\}\allowbreak{}\}\allowbreak{}\$\allowbreak{}}{DER\_\allowbreak{}XV-b:\allowbreak{}11}
\DERapparatusitem{ENV0893}{40;\allowbreak{}\allowbreak{} ms72\allowbreak{} proof\allowbreak{} finite-union\allowbreak{} assertion}{the\allowbreak{} relations\allowbreak{} \$\allowbreak{}(\allowbreak{}*)\allowbreak{}\$\allowbreak{}\allowbreak{} (\allowbreak{}p.\allowbreak{}\allowbreak{} 68)\allowbreak{}}{the\allowbreak{} relations\allowbreak{} \$\allowbreak{}(\allowbreak{}*)\allowbreak{}\$\allowbreak{}\allowbreak{} (\allowbreak{}p.\allowbreak{}\allowbreak{} 69)\allowbreak{}}{DER\_\allowbreak{}XV-b:\allowbreak{}12}
\DERapparatusitem{ENV0894}{40;\allowbreak{}\allowbreak{} ms72\allowbreak{} finite-union\allowbreak{} proof}{a\allowbreak{} union\allowbreak{} of\allowbreak{} the\allowbreak{} (\allowbreak{}filtering)\allowbreak{}\allowbreak{} subsets\allowbreak{} \$\allowbreak{}I\_\allowbreak{}\{\allowbreak{}M\_\allowbreak{}\{\allowbreak{}\textbackslash{}\allowbreak{}alpha\}\allowbreak{}\}\allowbreak{}\$\allowbreak{}}{a\allowbreak{} union\allowbreak{} of\allowbreak{} the\allowbreak{} subsets\allowbreak{} \$\allowbreak{}I\_\allowbreak{}\{\allowbreak{}M\_\allowbreak{}\{\allowbreak{}\textbackslash{}\allowbreak{}alpha\}\allowbreak{}\}\allowbreak{}\$\allowbreak{}}{DER\_\allowbreak{}XV-b:\allowbreak{}13}
\DERapparatusitem{ENV0895}{42;\allowbreak{}\allowbreak{} ms76\allowbreak{} Theorem6.\allowbreak{}1\allowbreak{} hypothesis(\allowbreak{}d)\allowbreak{}}{\textbackslash{}\allowbreak{}item[\allowbreak{}d)\allowbreak{}]\allowbreak{}\allowbreak{} \textbackslash{}\allowbreak{}textit\{\allowbreak{}The\}\allowbreak{}\allowbreak{} \$\allowbreak{}\textbackslash{}\allowbreak{}Filt\textasciicircum{}\allowbreak{}\{\allowbreak{}\textbackslash{}\allowbreak{}pi\}\allowbreak{}\$\allowbreak{}\allowbreak{} \textbackslash{}\allowbreak{}textit\{\allowbreak{}are\allowbreak{} stable\allowbreak{} under\}\allowbreak{}}{\textbackslash{}\allowbreak{}item[\allowbreak{}d)\allowbreak{}]\allowbreak{}\allowbreak{} \textbackslash{}\allowbreak{}textit\{\allowbreak{}The\}\allowbreak{}\allowbreak{} \$\allowbreak{}\textbackslash{}\allowbreak{}Filt\textasciicircum{}\allowbreak{}\{\allowbreak{}\textbackslash{}\allowbreak{}pi\}\allowbreak{}\$\allowbreak{}\allowbreak{} \textbackslash{}\allowbreak{}textit\{\allowbreak{}are\allowbreak{} stable\allowbreak{} under\allowbreak{} passage\allowbreak{} to\allowbreak{} subobjects\allowbreak{} and\allowbreak{} under\}\allowbreak{}}{DER\_\allowbreak{}XV-b:\allowbreak{}14}
\DERapparatusitem{ENV0896}{43;\allowbreak{}\allowbreak{} ms79\allowbreak{} end\allowbreak{} of\allowbreak{} transfinite\allowbreak{} termination\allowbreak{} argument}{if\allowbreak{} \$\allowbreak{}\textbackslash{}\allowbreak{}pi\$\allowbreak{}\allowbreak{} is\allowbreak{} the\allowbreak{} set\allowbreak{} of\allowbreak{} subobjects\allowbreak{} of\allowbreak{} \$\allowbreak{}B\$\allowbreak{}}{if\allowbreak{} \$\allowbreak{}\textbackslash{}\allowbreak{}mu\$\allowbreak{}\allowbreak{} is\allowbreak{} the\allowbreak{} cardinality\allowbreak{} of\allowbreak{} the\allowbreak{} set\allowbreak{} of\allowbreak{} subobjects\allowbreak{} of\allowbreak{} \$\allowbreak{}B\$\allowbreak{}}{DER\_\allowbreak{}XV-b:\allowbreak{}15}
\DERapparatusitem{ENV0897}{44;\allowbreak{}\allowbreak{} ms80\allowbreak{} first\allowbreak{} sentence}{\$\allowbreak{}\textbackslash{}\allowbreak{}mathrm\{\allowbreak{}card\}\allowbreak{}\textbackslash{}\allowbreak{},\allowbreak{}A\_\allowbreak{}\textbackslash{}\allowbreak{}alpha\allowbreak{} \textbackslash{}\allowbreak{}leq\allowbreak{} \textbackslash{}\allowbreak{}pi\$\allowbreak{}}{\$\allowbreak{}\textbackslash{}\allowbreak{}mathrm\{\allowbreak{}card\}\allowbreak{}\textbackslash{}\allowbreak{},\allowbreak{}\textbackslash{}\allowbreak{}alpha\allowbreak{} \textbackslash{}\allowbreak{}leq\allowbreak{} \textbackslash{}\allowbreak{}mu\$\allowbreak{}}{DER\_\allowbreak{}XV-b:\allowbreak{}16}
\DERapparatusitem{ENV0898}{44;\allowbreak{}\allowbreak{} ms81\allowbreak{} parenthesis\allowbreak{} after(\allowbreak{}e′)\allowbreak{}}{(\allowbreak{}It\allowbreak{} suffices\allowbreak{} it\allowbreak{} here}{(\allowbreak{}It\allowbreak{} suffices\allowbreak{} here}{DER\_\allowbreak{}XV-b:\allowbreak{}17}
\DERapparatusitem{ENV0899}{44;\allowbreak{}\allowbreak{} ms81\allowbreak{} condition(\allowbreak{}e0)\allowbreak{}}{The\allowbreak{} filtering\allowbreak{} \$\allowbreak{}\textbackslash{}\allowbreak{}varinjlim\$\allowbreak{}\allowbreak{} in\allowbreak{} \$\allowbreak{}I\$\allowbreak{}\allowbreak{} are\allowbreak{} left\allowbreak{} exact.\allowbreak{}}{The\allowbreak{} filtering\allowbreak{} \$\allowbreak{}\textbackslash{}\allowbreak{}varinjlim\$\allowbreak{}\allowbreak{} in\allowbreak{} \$\allowbreak{}\textbackslash{}\allowbreak{}M\$\allowbreak{}\allowbreak{} are\allowbreak{} left\allowbreak{} exact.\allowbreak{}}{DER\_\allowbreak{}XV-b:\allowbreak{}18}
\DERapparatusitem{ENV0900}{44;\allowbreak{}\allowbreak{} ms82\allowbreak{} use\allowbreak{} of\allowbreak{} bounded\allowbreak{} arrow\allowbreak{} filtration}{\textbackslash{}\allowbreak{}in\allowbreak{} \textbackslash{}\allowbreak{}Filt\textasciicircum{}\allowbreak{}\{\allowbreak{}\textbackslash{}\allowbreak{}pi\}\allowbreak{}}{\textbackslash{}\allowbreak{}in\textbackslash{}\allowbreak{}Fl\textasciicircum{}\allowbreak{}\textbackslash{}\allowbreak{}pi(\allowbreak{}\textbackslash{}\allowbreak{}mathcal\allowbreak{} M)\allowbreak{}}{DER\_\allowbreak{}XV-b:\allowbreak{}19}
\DERapparatusitem{ENV0901}{45;\allowbreak{}\allowbreak{} ms82\allowbreak{} final\allowbreak{} sentence\allowbreak{} before\allowbreak{} ms83}{the\allowbreak{} set\allowbreak{} of\allowbreak{} \$\allowbreak{}J\allowbreak{} \textbackslash{}\allowbreak{}subset\allowbreak{} \textbackslash{}\allowbreak{}Ss\_\allowbreak{}\{\allowbreak{}\textbackslash{}\allowbreak{}pi\}\allowbreak{}(\allowbreak{}B)\allowbreak{}\$\allowbreak{}}{the\allowbreak{} set\allowbreak{} \$\allowbreak{}J\allowbreak{} \textbackslash{}\allowbreak{}subset\allowbreak{} \textbackslash{}\allowbreak{}Ss\_\allowbreak{}\{\allowbreak{}\textbackslash{}\allowbreak{}pi\}\allowbreak{}(\allowbreak{}B)\allowbreak{}\$\allowbreak{}}{DER\_\allowbreak{}XV-b:\allowbreak{}20}
\DERapparatusitem{ENV0902}{49;\allowbreak{}\allowbreak{} running\allowbreak{} text\allowbreak{} or\allowbreak{} formula}{\textbackslash{}\allowbreak{}text\{\allowbreak{}def\}\allowbreak{}}{\textbackslash{}\allowbreak{}mathrm\{\allowbreak{}def\}\allowbreak{}}{TYPESETTING:\allowbreak{}7;\allowbreak{}\allowbreak{} NOTATION\_\allowbreak{}DER:\allowbreak{}17}
\DERapparatusitem{ENV0903}{49;\allowbreak{}\allowbreak{} ms90,\allowbreak{}\allowbreak{} flexible\allowbreak{} generator\allowbreak{} condition\allowbreak{} (\allowbreak{}7.\allowbreak{}2)\allowbreak{}}{For\allowbreak{} every\allowbreak{} monomorphism\allowbreak{} \$\allowbreak{}i\allowbreak{} :\allowbreak{}\allowbreak{} A\allowbreak{} \textbackslash{}\allowbreak{}longrightarrow\allowbreak{} B\$\allowbreak{}\allowbreak{} in\allowbreak{} \$\allowbreak{}\textbackslash{}\allowbreak{}mathcal\{\allowbreak{}M\}\allowbreak{}\$\allowbreak{},\allowbreak{}}{For\allowbreak{} every\allowbreak{} monomorphism\allowbreak{} \$\allowbreak{}i\allowbreak{} :\allowbreak{}\allowbreak{} A\allowbreak{} \textbackslash{}\allowbreak{}longrightarrow\allowbreak{} B\$\allowbreak{}\allowbreak{} in\allowbreak{} \$\allowbreak{}\textbackslash{}\allowbreak{}mathcal\{\allowbreak{}M\}\allowbreak{}\$\allowbreak{}\allowbreak{} which\allowbreak{} is\allowbreak{} not\allowbreak{} an\allowbreak{} isomorphism,\allowbreak{}}{DER\_\allowbreak{}XV-b:\allowbreak{}27}
\DERapparatusitem{ENV0904}{51;\allowbreak{}\allowbreak{} running\allowbreak{} text\allowbreak{} or\allowbreak{} formula}{\textbackslash{}\allowbreak{}text\{\allowbreak{}def\}\allowbreak{}}{\textbackslash{}\allowbreak{}mathrm\{\allowbreak{}def\}\allowbreak{}}{TYPESETTING:\allowbreak{}7;\allowbreak{}\allowbreak{} NOTATION\_\allowbreak{}DER:\allowbreak{}17}
\DERapparatusitem{ENV0905}{51;\allowbreak{}\allowbreak{} ms93,\allowbreak{}\allowbreak{} stopping-index\allowbreak{} sentence}{Thus\allowbreak{} there\allowbreak{} exists\allowbreak{} a\allowbreak{} cardinal\allowbreak{} \$\allowbreak{}\textbackslash{}\allowbreak{}alpha\$\allowbreak{}\allowbreak{} such\allowbreak{} that\allowbreak{} \$\allowbreak{}\textbackslash{}\allowbreak{}card\allowbreak{} \textbackslash{}\allowbreak{}alpha\allowbreak{} \textbackslash{}\allowbreak{}leq\allowbreak{} \textbackslash{}\allowbreak{}card\allowbreak{} \textbackslash{}\allowbreak{}Ss(\allowbreak{}B)\allowbreak{}\$\allowbreak{}}{Thus\allowbreak{} there\allowbreak{} exists\allowbreak{} an\allowbreak{} ordinal\allowbreak{} \$\allowbreak{}\textbackslash{}\allowbreak{}alpha\$\allowbreak{}\allowbreak{} such\allowbreak{} that\allowbreak{} \$\allowbreak{}\textbackslash{}\allowbreak{}card\allowbreak{} \textbackslash{}\allowbreak{}alpha\allowbreak{} \textbackslash{}\allowbreak{}leq\allowbreak{} \textbackslash{}\allowbreak{}card\allowbreak{} \textbackslash{}\allowbreak{}Ss(\allowbreak{}B)\allowbreak{}\$\allowbreak{}}{DER\_\allowbreak{}XV-b:\allowbreak{}28}
\DERapparatusitem{ENV0906}{53;\allowbreak{}\allowbreak{} ms96,\allowbreak{}\allowbreak{} cancellation\allowbreak{} reference}{by\allowbreak{} (\allowbreak{}7.\allowbreak{}1.\allowbreak{}4)\allowbreak{},\allowbreak{}\allowbreak{} (\allowbreak{}\textbackslash{}\allowbreak{}textit\{\allowbreak{}a\}\allowbreak{})\allowbreak{}}{by\allowbreak{} (\allowbreak{}7.\allowbreak{}1.\allowbreak{}4)\allowbreak{},\allowbreak{}\allowbreak{} (\allowbreak{}\textbackslash{}\allowbreak{}textit\{\allowbreak{}b\}\allowbreak{})\allowbreak{}}{DER\_\allowbreak{}XV-b:\allowbreak{}29}
\DERapparatusitem{ENV0907}{54;\allowbreak{}\allowbreak{} running\allowbreak{} text\allowbreak{} or\allowbreak{} formula}{\textbackslash{}\allowbreak{}text\{\allowbreak{}def\}\allowbreak{}}{\textbackslash{}\allowbreak{}mathrm\{\allowbreak{}def\}\allowbreak{}}{TYPESETTING:\allowbreak{}7;\allowbreak{}\allowbreak{} NOTATION\_\allowbreak{}DER:\allowbreak{}17}
\DERapparatusitem{ENV0908}{54;\allowbreak{}\allowbreak{} running\allowbreak{} text\allowbreak{} or\allowbreak{} formula}{\textbackslash{}\allowbreak{}text\{\allowbreak{}def\}\allowbreak{}}{\textbackslash{}\allowbreak{}mathrm\{\allowbreak{}def\}\allowbreak{}}{TYPESETTING:\allowbreak{}7;\allowbreak{}\allowbreak{} NOTATION\_\allowbreak{}DER:\allowbreak{}17}
\DERapparatusitem{ENV0909}{55;\allowbreak{}\allowbreak{} ms99,\allowbreak{}\allowbreak{} (\allowbreak{}7.\allowbreak{}4.\allowbreak{}3)\allowbreak{},\allowbreak{}\allowbreak{} fibration\allowbreak{} hypothesis\allowbreak{} below\allowbreak{} the\allowbreak{} square}{\textbackslash{}\allowbreak{}textit\{\allowbreak{}if\allowbreak{} \}\allowbreak{}\allowbreak{} f\textbackslash{}\allowbreak{}in\allowbreak{} F\allowbreak{} \textbackslash{}\allowbreak{}\allowbreak{} (\allowbreak{}\textbackslash{}\allowbreak{}text\{\allowbreak{}i.\allowbreak{}e.\allowbreak{}\allowbreak{} \}\allowbreak{}\allowbreak{} f\textbackslash{}\allowbreak{}in\allowbreak{} (\allowbreak{}C\textbackslash{}\allowbreak{}cap\allowbreak{} W)\allowbreak{}\_\allowbreak{}*)\allowbreak{}}{\textbackslash{}\allowbreak{}textit\{\allowbreak{}if\allowbreak{} \}\allowbreak{}\allowbreak{} f'\textbackslash{}\allowbreak{}in\allowbreak{} F\allowbreak{} \textbackslash{}\allowbreak{}\allowbreak{} (\allowbreak{}\textbackslash{}\allowbreak{}text\{\allowbreak{}i.\allowbreak{}e.\allowbreak{}\allowbreak{} \}\allowbreak{}\allowbreak{} f'\textbackslash{}\allowbreak{}in\allowbreak{} (\allowbreak{}C\textbackslash{}\allowbreak{}cap\allowbreak{} W)\allowbreak{}\_\allowbreak{}*)\allowbreak{}}{DER\_\allowbreak{}XV-b:\allowbreak{}30}
\DERapparatusitem{ENV0910}{55;\allowbreak{}\allowbreak{} ms100,\allowbreak{}\allowbreak{} presheaf\allowbreak{} category\allowbreak{} identification}{}{\textasciicircum{}\allowbreak{}\{\allowbreak{}\textbackslash{}\allowbreak{}mathrm\{\allowbreak{}op\}\allowbreak{}\}\allowbreak{}}{DER\_\allowbreak{}XV-b:\allowbreak{}31}
\DERapparatusitem{ENV0911}{57;\allowbreak{}\allowbreak{} ms104,\allowbreak{}\allowbreak{} formula\allowbreak{} (\allowbreak{}8.\allowbreak{}3.\allowbreak{}1)\allowbreak{}}{TC\allowbreak{} \textbackslash{}\allowbreak{};\allowbreak{}\&\allowbreak{}=\allowbreak{}\textbackslash{}\allowbreak{};\allowbreak{}\allowbreak{} C\textbackslash{}\allowbreak{}cap\allowbreak{} W\textbackslash{}\allowbreak{}\allowbreak{} ,\allowbreak{}}{TC\allowbreak{} \textbackslash{}\allowbreak{};\allowbreak{}\&\allowbreak{}=\allowbreak{}\textbackslash{}\allowbreak{};\allowbreak{}\allowbreak{} C\textbackslash{}\allowbreak{}cap\allowbreak{} W\_\allowbreak{}A\textbackslash{}\allowbreak{}\allowbreak{} ,\allowbreak{}}{DER\_\allowbreak{}XV-b:\allowbreak{}32}
\DERapparatusitem{ENV0912}{57;\allowbreak{}\allowbreak{} running\allowbreak{} text\allowbreak{} or\allowbreak{} formula}{\textbackslash{}\allowbreak{}text\{\allowbreak{}def\}\allowbreak{}}{\textbackslash{}\allowbreak{}mathrm\{\allowbreak{}def\}\allowbreak{}}{TYPESETTING:\allowbreak{}7;\allowbreak{}\allowbreak{} NOTATION\_\allowbreak{}DER:\allowbreak{}17}
\DERapparatusitem{ENV0913}{60;\allowbreak{}\allowbreak{} ms109,\allowbreak{}\allowbreak{} last\allowbreak{} sentence\allowbreak{} of\allowbreak{} proof8.\allowbreak{}6\allowbreak{} and\allowbreak{} existing\allowbreak{} [\allowbreak{}sic]\allowbreak{}\allowbreak{} note}{i}{f}{DER\_\allowbreak{}XV-b:\allowbreak{}33}
\DERapparatusitem{ENV0914}{64;\allowbreak{}\allowbreak{} ms.\allowbreak{}118,\allowbreak{}\allowbreak{} NB\allowbreak{} after\allowbreak{} Loc(\allowbreak{}8)\allowbreak{}}{Loc\textbackslash{}\allowbreak{},\allowbreak{}(\allowbreak{}7)\allowbreak{}\allowbreak{} tells\allowbreak{} us\allowbreak{} that\allowbreak{} \$\allowbreak{}W\textbackslash{}\allowbreak{}subset\textbackslash{}\allowbreak{}mathbf\{\allowbreak{}W\}\allowbreak{}\_\allowbreak{}0\$\allowbreak{}}{Loc\textbackslash{}\allowbreak{},\allowbreak{}(\allowbreak{}8)\allowbreak{}\allowbreak{} tells\allowbreak{} us\allowbreak{} that\allowbreak{} \$\allowbreak{}W\textbackslash{}\allowbreak{}subset\textbackslash{}\allowbreak{}mathbf\{\allowbreak{}W\}\allowbreak{}\_\allowbreak{}0\$\allowbreak{}}{DER\_\allowbreak{}XV-b:\allowbreak{}35}
\DERapparatusitem{ENV0915}{65;\allowbreak{}\allowbreak{} ms.\allowbreak{}120,\allowbreak{}\allowbreak{} top\allowbreak{} horizontal\allowbreak{} arrow\allowbreak{} of\allowbreak{} the\allowbreak{} two-square\allowbreak{} coproduct\allowbreak{} diagram}{X\_\allowbreak{}1\allowbreak{} \textbackslash{}\allowbreak{}arrow[\allowbreak{}r,\allowbreak{}\allowbreak{} "f"]\allowbreak{}\allowbreak{} \textbackslash{}\allowbreak{}arrow[\allowbreak{}d]\allowbreak{}\allowbreak{} \&\allowbreak{}\allowbreak{} Y\_\allowbreak{}1}{X\_\allowbreak{}1\allowbreak{} \textbackslash{}\allowbreak{}arrow[\allowbreak{}r,\allowbreak{}\allowbreak{} "f\_\allowbreak{}1"]\allowbreak{}\allowbreak{} \textbackslash{}\allowbreak{}arrow[\allowbreak{}d]\allowbreak{}\allowbreak{} \&\allowbreak{}\allowbreak{} Y\_\allowbreak{}1}{DER\_\allowbreak{}XV-b:\allowbreak{}37}
\DERapparatusitem{ENV0916}{66;\allowbreak{}\allowbreak{} ms.\allowbreak{}121,\allowbreak{}\allowbreak{} sentence\allowbreak{} introducing\allowbreak{} the\allowbreak{} perfect\allowbreak{} factorization}{by\allowbreak{} factoring\allowbreak{} \$\allowbreak{}X\_\allowbreak{}1\textbackslash{}\allowbreak{}longrightarrow\allowbreak{} Y\$\allowbreak{}\allowbreak{} as}{by\allowbreak{} factoring\allowbreak{} \$\allowbreak{}X\textbackslash{}\allowbreak{}longrightarrow\allowbreak{} Y\$\allowbreak{}\allowbreak{} as}{DER\_\allowbreak{}XV-b:\allowbreak{}38}
\DERapparatusitem{ENV0917}{66;\allowbreak{}\allowbreak{} ms.\allowbreak{}122,\allowbreak{}\allowbreak{} final\allowbreak{} assertion\allowbreak{} of\allowbreak{} Corollary\allowbreak{} 9.\allowbreak{}3}{\$\allowbreak{}(\allowbreak{}f\_\allowbreak{}0\allowbreak{} :\allowbreak{}\allowbreak{} X\_\allowbreak{}0\textbackslash{}\allowbreak{}longrightarrow\allowbreak{} Y\_\allowbreak{}0)\allowbreak{}\textbackslash{}\allowbreak{}in\allowbreak{} W\$\allowbreak{}}{\$\allowbreak{}(\allowbreak{}f\_\allowbreak{}0\allowbreak{} :\allowbreak{}\allowbreak{} X\_\allowbreak{}0\textbackslash{}\allowbreak{}longrightarrow\allowbreak{} X'\_\allowbreak{}0)\allowbreak{}\textbackslash{}\allowbreak{}in\allowbreak{} W\$\allowbreak{}}{DER\_\allowbreak{}XV-b:\allowbreak{}39}
\DERapparatusitem{ENV0918}{68;\allowbreak{}\allowbreak{} ms.\allowbreak{}125,\allowbreak{}\allowbreak{} final\allowbreak{} parenthetical\allowbreak{} sentence}{\textbackslash{}\allowbreak{}textit\{\allowbreak{}cf.\allowbreak{}\}\allowbreak{}\allowbreak{} Pursuing\allowbreak{} Stacks.\allowbreak{}}{\textbackslash{}\allowbreak{}textit\{\allowbreak{}cf.\allowbreak{}\}\allowbreak{}\allowbreak{} Pursuing\allowbreak{} Stacks.\allowbreak{})\allowbreak{}}{DER\_\allowbreak{}XV-b:\allowbreak{}40}
\DERapparatusitem{ENV0919}{68;\allowbreak{}\allowbreak{} ms.\allowbreak{}126,\allowbreak{}\allowbreak{} restriction\allowbreak{} to\allowbreak{} a\allowbreak{} missed\allowbreak{} connected\allowbreak{} component}{and\allowbreak{} by\allowbreak{} Loc\textbackslash{}\allowbreak{},\allowbreak{}(\allowbreak{}7)\allowbreak{},\allowbreak{}\allowbreak{} \$\allowbreak{}\textbackslash{}\allowbreak{}emptyset\textbackslash{}\allowbreak{}longrightarrow\allowbreak{} Y\_\allowbreak{}0\$\allowbreak{}}{and\allowbreak{} by\allowbreak{} Loc\textbackslash{}\allowbreak{},\allowbreak{}(\allowbreak{}5)\allowbreak{},\allowbreak{}\allowbreak{} \$\allowbreak{}\textbackslash{}\allowbreak{}emptyset\textbackslash{}\allowbreak{}longrightarrow\allowbreak{} Y\_\allowbreak{}0\$\allowbreak{}}{DER\_\allowbreak{}XV-b:\allowbreak{}41}
\DERapparatusitem{ENV0920}{68;\allowbreak{}\allowbreak{} ms.\allowbreak{}127,\allowbreak{}\allowbreak{} first\allowbreak{} declaration\allowbreak{} of\allowbreak{} the\allowbreak{} homotopy\allowbreak{} h}{\$\allowbreak{}h\textbackslash{}\allowbreak{}\allowbreak{} :\allowbreak{}\textbackslash{}\allowbreak{}\allowbreak{} X\textbackslash{}\allowbreak{}times\allowbreak{} I\textbackslash{}\allowbreak{}longrightarrow\allowbreak{} Y\$\allowbreak{}}{\$\allowbreak{}h\textbackslash{}\allowbreak{}\allowbreak{} :\allowbreak{}\textbackslash{}\allowbreak{}\allowbreak{} I\textbackslash{}\allowbreak{}times\allowbreak{} X\textbackslash{}\allowbreak{}longrightarrow\allowbreak{} Y\$\allowbreak{}}{DER\_\allowbreak{}XV-b:\allowbreak{}42}
\DERapparatusitem{ENV0921}{70;\allowbreak{}\allowbreak{} running\allowbreak{} text\allowbreak{} or\allowbreak{} formula}{\textbackslash{}\allowbreak{}text\{\allowbreak{}def\}\allowbreak{}}{\textbackslash{}\allowbreak{}mathrm\{\allowbreak{}def\}\allowbreak{}}{TYPESETTING:\allowbreak{}7;\allowbreak{}\allowbreak{} NOTATION\_\allowbreak{}DER:\allowbreak{}17}
\DERapparatusitem{ENV0922}{72;\allowbreak{}\allowbreak{} ms.\allowbreak{}133\allowbreak{} bis,\allowbreak{}\allowbreak{} prose\allowbreak{} after\allowbreak{} the\allowbreak{} descended\allowbreak{} component\allowbreak{} functor}{there\allowbreak{} is\allowbreak{} an\allowbreak{} inverse\allowbreak{} up\allowbreak{} to\allowbreak{} isomorphism,\allowbreak{}\allowbreak{} \$\allowbreak{}I\textbackslash{}\allowbreak{}longmapsto\textbackslash{}\allowbreak{}hot\_\allowbreak{}W(\allowbreak{}I)\allowbreak{}\$\allowbreak{}}{there\allowbreak{} is\allowbreak{} a\allowbreak{} right\allowbreak{} inverse\allowbreak{} up\allowbreak{} to\allowbreak{} isomorphism,\allowbreak{}\allowbreak{} \$\allowbreak{}I\textbackslash{}\allowbreak{}longmapsto\textbackslash{}\allowbreak{}hot\_\allowbreak{}W(\allowbreak{}I)\allowbreak{}\$\allowbreak{}}{DER\_\allowbreak{}XV-b:\allowbreak{}43}
\DERapparatusitem{ENV0923}{73;\allowbreak{}\allowbreak{} ms.\allowbreak{}133\allowbreak{} quinquies,\allowbreak{}\allowbreak{} final\allowbreak{} displayed\allowbreak{} copairing}{h\textbackslash{}\allowbreak{}\allowbreak{} =\allowbreak{}\textbackslash{}\allowbreak{}\allowbreak{} (\allowbreak{}f\_\allowbreak{}0,\allowbreak{}f\_\allowbreak{}1)\allowbreak{}\textbackslash{}\allowbreak{}\allowbreak{} :\allowbreak{}\textbackslash{}\allowbreak{}\allowbreak{} X\textbackslash{}\allowbreak{}amalg\allowbreak{} X\textbackslash{}\allowbreak{}\allowbreak{} \textbackslash{}\allowbreak{}longrightarrow\textbackslash{}\allowbreak{}\allowbreak{} X\textbackslash{}\allowbreak{}\allowbreak{} ,\allowbreak{}}{h\textbackslash{}\allowbreak{}\allowbreak{} =\allowbreak{}\textbackslash{}\allowbreak{}\allowbreak{} (\allowbreak{}f\_\allowbreak{}0,\allowbreak{}f\_\allowbreak{}1)\allowbreak{}\textbackslash{}\allowbreak{}\allowbreak{} :\allowbreak{}\textbackslash{}\allowbreak{}\allowbreak{} X\textbackslash{}\allowbreak{}amalg\allowbreak{} X\textbackslash{}\allowbreak{}\allowbreak{} \textbackslash{}\allowbreak{}longrightarrow\textbackslash{}\allowbreak{}\allowbreak{} Y\textbackslash{}\allowbreak{}\allowbreak{} ,\allowbreak{}}{DER\_\allowbreak{}XV-b:\allowbreak{}45}
\DERapparatusitem{ENV0924}{75;\allowbreak{}\allowbreak{} ms.\allowbreak{}136,\allowbreak{}\allowbreak{} hypothesis\allowbreak{} and\allowbreak{} explanatory\allowbreak{} parenthesis\allowbreak{} of\allowbreak{} 10.\allowbreak{}1(\allowbreak{}3b)\allowbreak{}}{Loc\textbackslash{}\allowbreak{},\allowbreak{}(\allowbreak{}1\textbackslash{}\allowbreak{},\allowbreak{}ter,\allowbreak{}\allowbreak{} 2,\allowbreak{}\allowbreak{} 3,\allowbreak{}\allowbreak{} 5\textbackslash{}\allowbreak{},\allowbreak{}b)\allowbreak{}}{Loc\textbackslash{}\allowbreak{},\allowbreak{}(\allowbreak{}1\textbackslash{}\allowbreak{},\allowbreak{}ter,\allowbreak{}\allowbreak{} 2,\allowbreak{}\allowbreak{} 3,\allowbreak{}\allowbreak{} 6\textbackslash{}\allowbreak{},\allowbreak{}b)\allowbreak{}}{DER\_\allowbreak{}XV-b:\allowbreak{}46}
\DERapparatusitem{ENV0925}{75;\allowbreak{}\allowbreak{} ms.\allowbreak{}136,\allowbreak{}\allowbreak{} hypothesis\allowbreak{} and\allowbreak{} explanatory\allowbreak{} parenthesis\allowbreak{} of\allowbreak{} 10.\allowbreak{}1(\allowbreak{}3b)\allowbreak{}}{Loc\textbackslash{}\allowbreak{},\allowbreak{}(\allowbreak{}5\textbackslash{}\allowbreak{},\allowbreak{}b)\allowbreak{}\allowbreak{} \textbackslash{}\allowbreak{}textit\{\allowbreak{}of\allowbreak{} \$\allowbreak{}W\$\allowbreak{}-fibrations}{Loc\textbackslash{}\allowbreak{},\allowbreak{}(\allowbreak{}6\textbackslash{}\allowbreak{},\allowbreak{}b)\allowbreak{}\allowbreak{} \textbackslash{}\allowbreak{}textit\{\allowbreak{}of\allowbreak{} \$\allowbreak{}W\$\allowbreak{}-fibrations}{DER\_\allowbreak{}XV-b:\allowbreak{}46}
\subsection*{Chapter XVI}
\DERapparatusitem{ENV0926}{2;\allowbreak{}\allowbreak{} Proposition1.\allowbreak{}3,\allowbreak{}\allowbreak{} ms3}{\textbackslash{}\allowbreak{}mathcal\{\allowbreak{}M\}\allowbreak{}}{\textbackslash{}\allowbreak{}Cat}{DER\_\allowbreak{}XVI:\allowbreak{}3}
\DERapparatusitem{ENV0927}{3;\allowbreak{}\allowbreak{} Proof\allowbreak{} of(\allowbreak{}3-prime\allowbreak{} b)\allowbreak{},\allowbreak{}\allowbreak{} final\allowbreak{} conclusion,\allowbreak{}\allowbreak{} ms5}{\textasciicircum{}\allowbreak{}\textbackslash{}\allowbreak{}circ\allowbreak{} \textbackslash{}\allowbreak{}in\allowbreak{} W\textasciicircum{}\allowbreak{}\textbackslash{}\allowbreak{}circ}{\textbackslash{}\allowbreak{}in\allowbreak{} W}{DER\_\allowbreak{}XVI:\allowbreak{}4}
\DERapparatusitem{ENV0928}{4;\allowbreak{}\allowbreak{} Proof\allowbreak{} of\allowbreak{} weak\allowbreak{} fibre\allowbreak{} criterion,\allowbreak{}\allowbreak{} ms7}{the\allowbreak{} fibers\allowbreak{} \$\allowbreak{}s\textbackslash{}\allowbreak{}backslash\allowbreak{} X\$\allowbreak{}\allowbreak{} of\allowbreak{} \$\allowbreak{}\textbackslash{}\allowbreak{}Phi\_\allowbreak{}S(\allowbreak{}X)\allowbreak{}\$\allowbreak{}\allowbreak{} are\allowbreak{} \$\allowbreak{}W\$\allowbreak{}-coaspheric}{the\allowbreak{} fibers\allowbreak{} \$\allowbreak{}s\textbackslash{}\allowbreak{}backslash\allowbreak{} X\$\allowbreak{}\allowbreak{} of\allowbreak{} \$\allowbreak{}\textbackslash{}\allowbreak{}Phi\_\allowbreak{}S(\allowbreak{}X)\allowbreak{}\$\allowbreak{}\allowbreak{} are\allowbreak{} \$\allowbreak{}W\$\allowbreak{}-aspheric}{DER\_\allowbreak{}XVI:\allowbreak{}5}
\DERapparatusitem{ENV0929}{5;\allowbreak{}\allowbreak{} Corollary2.\allowbreak{}2\allowbreak{} parenthesis,\allowbreak{}\allowbreak{} ms9}{It\allowbreak{} is\allowbreak{} equivalent\allowbreak{} to\}\allowbreak{}\allowbreak{} Loc\allowbreak{} (\allowbreak{}4)\allowbreak{}}{It\allowbreak{} is\allowbreak{} equivalent\allowbreak{} to\}\allowbreak{}\allowbreak{} Loc\allowbreak{} (\allowbreak{}4\allowbreak{} a)\allowbreak{}}{DER\_\allowbreak{}XVI:\allowbreak{}6}
\DERapparatusitem{ENV0930}{5;\allowbreak{}\allowbreak{} running\allowbreak{} text\allowbreak{} or\allowbreak{} formula}{\textbackslash{}\allowbreak{}textit\{\allowbreak{}def\}\allowbreak{}}{\textbackslash{}\allowbreak{}mathrm\{\allowbreak{}def\}\allowbreak{}}{TYPESETTING:\allowbreak{}7;\allowbreak{}\allowbreak{} NOTATION\_\allowbreak{}DER:\allowbreak{}17}
\DERapparatusitem{ENV0931}{7;\allowbreak{}\allowbreak{} Footnote5,\allowbreak{}\allowbreak{} ms12}{(\allowbreak{}condition\allowbreak{} Loc\allowbreak{} (\allowbreak{}1\allowbreak{} ter)\allowbreak{},\allowbreak{}}{(\allowbreak{}condition\allowbreak{} Loc\allowbreak{} (\allowbreak{}1\allowbreak{} ter)\allowbreak{})\allowbreak{},\allowbreak{}}{DER\_\allowbreak{}XVI:\allowbreak{}9}
\DERapparatusitem{ENV0932}{7;\allowbreak{}\allowbreak{} running\allowbreak{} text\allowbreak{} or\allowbreak{} formula}{\textbackslash{}\allowbreak{}text\{\allowbreak{}def\}\allowbreak{}}{\textbackslash{}\allowbreak{}mathrm\{\allowbreak{}def\}\allowbreak{}}{TYPESETTING:\allowbreak{}7;\allowbreak{}\allowbreak{} NOTATION\_\allowbreak{}DER:\allowbreak{}17}
\DERapparatusitem{ENV0933}{8;\allowbreak{}\allowbreak{} Commentary2.\allowbreak{}10,\allowbreak{}\allowbreak{} ms15}{(\allowbreak{}such\allowbreak{} as\allowbreak{} Cat,\allowbreak{}\allowbreak{} where\allowbreak{} \$\allowbreak{}\textbackslash{}\allowbreak{}mathrm\{\allowbreak{}Hot\}\allowbreak{}\_\allowbreak{}W\$\allowbreak{}}{(\allowbreak{}such\allowbreak{} as\allowbreak{} Cat,\allowbreak{}\allowbreak{} or\allowbreak{} \$\allowbreak{}\textbackslash{}\allowbreak{}mathrm\{\allowbreak{}Hot\}\allowbreak{}\_\allowbreak{}W\$\allowbreak{}}{DER\_\allowbreak{}XVI:\allowbreak{}11}
\DERapparatusitem{ENV0934}{8;\allowbreak{}\allowbreak{} running\allowbreak{} text\allowbreak{} or\allowbreak{} formula}{\textbackslash{}\allowbreak{}text\{\allowbreak{}def\}\allowbreak{}}{\textbackslash{}\allowbreak{}mathrm\{\allowbreak{}def\}\allowbreak{}}{TYPESETTING:\allowbreak{}7;\allowbreak{}\allowbreak{} NOTATION\_\allowbreak{}DER:\allowbreak{}17}
\DERapparatusitem{ENV0935}{9;\allowbreak{}\allowbreak{} Corollary2.\allowbreak{}11\allowbreak{} parenthesis,\allowbreak{}\allowbreak{} ms17\allowbreak{} (\allowbreak{}Ou)\allowbreak{}}{(\allowbreak{}\textbackslash{}\allowbreak{}textit\{\allowbreak{}Where\allowbreak{} again,\allowbreak{}\}\allowbreak{}}{(\allowbreak{}\textbackslash{}\allowbreak{}textit\{\allowbreak{}Or\allowbreak{} again,\allowbreak{}\}\allowbreak{}}{DER\_\allowbreak{}XVI:\allowbreak{}14}
\DERapparatusitem{ENV0936}{9;\allowbreak{}\allowbreak{} Corollary2.\allowbreak{}11\allowbreak{} parenthesis,\allowbreak{}\allowbreak{} ms17\allowbreak{} (\allowbreak{}de\allowbreak{} localisation)\allowbreak{}}{is\allowbreak{} a\allowbreak{} functor\allowbreak{} the\allowbreak{} localization}{is\allowbreak{} a\allowbreak{} localization\allowbreak{} functor}{DER\_\allowbreak{}XVI:\allowbreak{}15}
\DERapparatusitem{ENV0937}{10;\allowbreak{}\allowbreak{} Lemma2.\allowbreak{}14(\allowbreak{}a)\allowbreak{},\allowbreak{}\allowbreak{} second\allowbreak{} backwards-induction\allowbreak{} equality,\allowbreak{}\allowbreak{} ms13-prime}{1}{2}{DER\_\allowbreak{}XVI:\allowbreak{}17}
\DERapparatusitem{ENV0938}{11;\allowbreak{}\allowbreak{} Application\allowbreak{} of\allowbreak{} Lemma2.\allowbreak{}14,\allowbreak{}\allowbreak{} ms14-prime}{(\allowbreak{}this\allowbreak{} is\allowbreak{} essentially\allowbreak{} corollary\allowbreak{} 2.\allowbreak{}12.\allowbreak{})\allowbreak{}}{(\allowbreak{}this\allowbreak{} follows\allowbreak{} from\allowbreak{} Loc\allowbreak{} (\allowbreak{}5)\allowbreak{})\allowbreak{}}{DER\_\allowbreak{}XVI:\allowbreak{}20}
\DERapparatusitem{ENV0939}{14;\allowbreak{}\allowbreak{} Explanation\allowbreak{} immediately\allowbreak{} after\allowbreak{} cube\allowbreak{} Q-prime,\allowbreak{}\allowbreak{} ms18-prime}{1}{2}{DER\_\allowbreak{}XVI:\allowbreak{}28}
\DERapparatusitem{ENV0940}{18;\allowbreak{}\allowbreak{} Corollary3.\allowbreak{}7,\allowbreak{}\allowbreak{} ms26}{Loc\allowbreak{} (\allowbreak{}6\allowbreak{} B)\allowbreak{}\allowbreak{} \textbackslash{}\allowbreak{}textit\{\allowbreak{}above\}\allowbreak{})\allowbreak{},\allowbreak{}}{Loc\allowbreak{} (\allowbreak{}6\allowbreak{} B)\allowbreak{}\allowbreak{} \textbackslash{}\allowbreak{}textit\{\allowbreak{}above\}\allowbreak{},\allowbreak{}}{DER\_\allowbreak{}XVI:\allowbreak{}30}
\DERapparatusitem{ENV0941}{21;\allowbreak{}\allowbreak{} ms33\allowbreak{} §4.\allowbreak{}2}{}{\textasciicircum{}\allowbreak{}\textbackslash{}\allowbreak{}circ}{DER\_\allowbreak{}XVI:\allowbreak{}33}
\DERapparatusitem{ENV0942}{22;\allowbreak{}\allowbreak{} ms35\allowbreak{} before\allowbreak{} ms36\allowbreak{} marker}{except\allowbreak{} under\allowbreak{} special\allowbreak{} axioms\allowbreak{} on\allowbreak{} \$\allowbreak{}W\$\allowbreak{}\allowbreak{} or\allowbreak{} on\allowbreak{} \$\allowbreak{}W\_\allowbreak{}0\$\allowbreak{})\allowbreak{}.\allowbreak{}}{except\allowbreak{} under\allowbreak{} special\allowbreak{} axioms\allowbreak{} on\allowbreak{} \$\allowbreak{}W\$\allowbreak{}\allowbreak{} or\allowbreak{} on\allowbreak{} \$\allowbreak{}W\_\allowbreak{}0\$\allowbreak{}.\allowbreak{}}{DER\_\allowbreak{}XVI:\allowbreak{}35}
\DERapparatusitem{ENV0943}{23;\allowbreak{}\allowbreak{} ms38\allowbreak{} paragraph\allowbreak{} preceding\allowbreak{} Loc7}{(\allowbreak{}\$\allowbreak{}W\$\allowbreak{}\allowbreak{} strongly\allowbreak{} saturated,\allowbreak{}\allowbreak{} \textbackslash{}\allowbreak{}textit\{\allowbreak{}i.\allowbreak{}e.\allowbreak{}\}\allowbreak{}\allowbreak{} Loc\allowbreak{} (\allowbreak{}1\allowbreak{} bis)\allowbreak{})\allowbreak{}}{(\allowbreak{}\$\allowbreak{}W\$\allowbreak{}\allowbreak{} strongly\allowbreak{} saturated,\allowbreak{}\allowbreak{} \textbackslash{}\allowbreak{}textit\{\allowbreak{}i.\allowbreak{}e.\allowbreak{}\}\allowbreak{}\allowbreak{} Loc\allowbreak{} (\allowbreak{}1\allowbreak{} ter)\allowbreak{},\allowbreak{}\allowbreak{} and\allowbreak{} satisfying\allowbreak{} Loc\allowbreak{} (\allowbreak{}3\allowbreak{} bis)\allowbreak{})\allowbreak{}}{DER\_\allowbreak{}XVI:\allowbreak{}37}
\DERapparatusitem{ENV0944}{24;\allowbreak{}\allowbreak{} ms39\allowbreak{} Loc7F}{\$\allowbreak{}(\allowbreak{}W,\allowbreak{}\allowbreak{} \textbackslash{}\allowbreak{}Fib\_\allowbreak{}W)\allowbreak{}\$\allowbreak{}\allowbreak{} is\allowbreak{} a\allowbreak{} category\allowbreak{} with\allowbreak{} cofibrations}{\$\allowbreak{}(\allowbreak{}W,\allowbreak{}\allowbreak{} \textbackslash{}\allowbreak{}Fib\_\allowbreak{}W)\allowbreak{}\$\allowbreak{}\allowbreak{} is\allowbreak{} a\allowbreak{} category\allowbreak{} with\allowbreak{} fibrations}{DER\_\allowbreak{}XVI:\allowbreak{}38}
\DERapparatusitem{ENV0945}{24;\allowbreak{}\allowbreak{} ms39\allowbreak{} Loc7F}{\$\allowbreak{}X\allowbreak{} \textbackslash{}\allowbreak{}longrightarrow\allowbreak{} e\$\allowbreak{}\allowbreak{} is\allowbreak{} in\allowbreak{} \$\allowbreak{}\textbackslash{}\allowbreak{}Fib\_\allowbreak{}W\$\allowbreak{}\allowbreak{} for\allowbreak{} every\allowbreak{} \$\allowbreak{}X\$\allowbreak{}\allowbreak{} in\allowbreak{} Cat.\allowbreak{}}{\$\allowbreak{}X\allowbreak{} \textbackslash{}\allowbreak{}longrightarrow\allowbreak{} e\$\allowbreak{}\allowbreak{} is\allowbreak{} in\allowbreak{} \$\allowbreak{}\textbackslash{}\allowbreak{}Fib\_\allowbreak{}W\$\allowbreak{}\allowbreak{} for\allowbreak{} every\allowbreak{} \$\allowbreak{}X\$\allowbreak{}\allowbreak{} in\allowbreak{} Cat)\allowbreak{}.\allowbreak{}}{DER\_\allowbreak{}XVI:\allowbreak{}39}
\DERapparatusitem{ENV0946}{24;\allowbreak{}\allowbreak{} ms41\allowbreak{} paragraph\allowbreak{} beginning\allowbreak{} Nous}{(\allowbreak{}We\allowbreak{} will\allowbreak{} apply\allowbreak{} it\allowbreak{} in\allowbreak{} the\allowbreak{} following\allowbreak{} cases:\allowbreak{}}{We\allowbreak{} will\allowbreak{} apply\allowbreak{} it\allowbreak{} in\allowbreak{} the\allowbreak{} following\allowbreak{} cases:\allowbreak{}}{DER\_\allowbreak{}XVI:\allowbreak{}40}
\DERapparatusitem{ENV0947}{24;\allowbreak{}\allowbreak{} ms41\allowbreak{} footnote28}{cF}{Parf}{DER\_\allowbreak{}XVI:\allowbreak{}41}
\DERapparatusitem{ENV0948}{24;\allowbreak{}\allowbreak{} ms41\allowbreak{} footnote28}{cF}{Parf}{DER\_\allowbreak{}XVI:\allowbreak{}41}
\DERapparatusitem{ENV0949}{25;\allowbreak{}\allowbreak{} ms41\allowbreak{} Lemma4.\allowbreak{}4.\allowbreak{}1}{f}{p}{DER\_\allowbreak{}XVI:\allowbreak{}42}
\DERapparatusitem{ENV0950}{25;\allowbreak{}\allowbreak{} ms42\allowbreak{} underbrace\allowbreak{} in\allowbreak{} factorization}{\{\allowbreak{}}{}{DER\_\allowbreak{}XVI:\allowbreak{}44}
\DERapparatusitem{ENV0951}{25;\allowbreak{}\allowbreak{} ms42\allowbreak{} underbrace\allowbreak{} in\allowbreak{} factorization}{\_\allowbreak{}0}{}{DER\_\allowbreak{}XVI:\allowbreak{}44}
\DERapparatusitem{ENV0952}{25;\allowbreak{}\allowbreak{} ms42\allowbreak{} underbrace\allowbreak{} in\allowbreak{} factorization}{\}\allowbreak{}}{}{DER\_\allowbreak{}XVI:\allowbreak{}44}
\DERapparatusitem{ENV0953}{28;\allowbreak{}\allowbreak{} ms49\allowbreak{} §4.\allowbreak{}4.\allowbreak{}7}{two\allowbreak{} for\allowbreak{} the\allowbreak{} for\allowbreak{} the\allowbreak{} two\allowbreak{} triples}{two\allowbreak{} for\allowbreak{} the\allowbreak{} two\allowbreak{} triples}{DER\_\allowbreak{}XVI:\allowbreak{}46}
\DERapparatusitem{ENV0954}{28;\allowbreak{}\allowbreak{} ms49\allowbreak{} §4.\allowbreak{}4.\allowbreak{}7}{}{f}{DER\_\allowbreak{}XVI:\allowbreak{}47}
\DERapparatusitem{ENV0955}{28;\allowbreak{}\allowbreak{} ms49\allowbreak{} §4.\allowbreak{}4.\allowbreak{}7}{}{f}{DER\_\allowbreak{}XVI:\allowbreak{}47}
\DERapparatusitem{ENV0956}{29;\allowbreak{}\allowbreak{} ms49\allowbreak{} implication\allowbreak{} before\allowbreak{} ms50}{}{f}{DER\_\allowbreak{}XVI:\allowbreak{}48}
\DERapparatusitem{ENV0957}{29;\allowbreak{}\allowbreak{} ms50\allowbreak{} explanation\allowbreak{} of\allowbreak{} D=\allowbreak{}>F}{is\allowbreak{} a\allowbreak{} category\allowbreak{} with\allowbreak{} cofibrations\allowbreak{} for\allowbreak{} \$\allowbreak{}(\allowbreak{}W,\allowbreak{}\allowbreak{} \textbackslash{}\allowbreak{}Fib\_\allowbreak{}W)\allowbreak{}\$\allowbreak{}}{is\allowbreak{} a\allowbreak{} category\allowbreak{} with\allowbreak{} fibrations\allowbreak{} for\allowbreak{} \$\allowbreak{}(\allowbreak{}W,\allowbreak{}\allowbreak{} \textbackslash{}\allowbreak{}Fib\_\allowbreak{}W)\allowbreak{}\$\allowbreak{}}{DER\_\allowbreak{}XVI:\allowbreak{}49}
\DERapparatusitem{ENV0958}{29;\allowbreak{}\allowbreak{} ms50\allowbreak{} isolated\allowbreak{} display}{\textbackslash{}\allowbreak{}Parf\_\allowbreak{}\{\allowbreak{}W\_\allowbreak{}0\}\allowbreak{}\allowbreak{} \textbackslash{}\allowbreak{}subset\allowbreak{} W.\allowbreak{}}{\textbackslash{}\allowbreak{}Parf\_\allowbreak{}\{\allowbreak{}W\_\allowbreak{}0\}\allowbreak{}\allowbreak{} \textbackslash{}\allowbreak{}subset\allowbreak{} \textbackslash{}\allowbreak{}Fib\_\allowbreak{}W.\allowbreak{}}{DER\_\allowbreak{}XVI:\allowbreak{}50}
\DERapparatusitem{ENV0959}{29;\allowbreak{}\allowbreak{} ms50\allowbreak{} factorization\allowbreak{} display}{\textbackslash{}\allowbreak{}quad\allowbreak{} f\allowbreak{} \textbackslash{}\allowbreak{}in\allowbreak{} \textbackslash{}\allowbreak{}Parf\_\allowbreak{}\{\allowbreak{}W\_\allowbreak{}0\}\allowbreak{}}{\textbackslash{}\allowbreak{}quad\allowbreak{} p\allowbreak{} \textbackslash{}\allowbreak{}in\allowbreak{} \textbackslash{}\allowbreak{}Parf\_\allowbreak{}\{\allowbreak{}W\_\allowbreak{}0\}\allowbreak{}}{DER\_\allowbreak{}XVI:\allowbreak{}51}
\DERapparatusitem{ENV0960}{29;\allowbreak{}\allowbreak{} ms50\allowbreak{} closing\allowbreak{} sentence}{we\allowbreak{} find\allowbreak{} the\allowbreak{} factorization\allowbreak{} with\allowbreak{} \$\allowbreak{}i\allowbreak{} \textbackslash{}\allowbreak{}in\allowbreak{} W\$\allowbreak{},\allowbreak{}\allowbreak{} \$\allowbreak{}f\allowbreak{} \textbackslash{}\allowbreak{}in\allowbreak{} \textbackslash{}\allowbreak{}Fib\_\allowbreak{}W\$\allowbreak{}.\allowbreak{}}{we\allowbreak{} find\allowbreak{} the\allowbreak{} factorization\allowbreak{} with\allowbreak{} \$\allowbreak{}i\allowbreak{} \textbackslash{}\allowbreak{}in\allowbreak{} W\$\allowbreak{},\allowbreak{}\allowbreak{} \$\allowbreak{}p\allowbreak{} \textbackslash{}\allowbreak{}in\allowbreak{} \textbackslash{}\allowbreak{}Fib\_\allowbreak{}W\$\allowbreak{}.\allowbreak{}}{DER\_\allowbreak{}XVI:\allowbreak{}52}
\DERapparatusitem{ENV0961}{30;\allowbreak{}\allowbreak{} ms52\allowbreak{} proof\allowbreak{} of\allowbreak{} Lemma4.\allowbreak{}4.\allowbreak{}9(\allowbreak{}C)\allowbreak{}}{factoring\allowbreak{} as\allowbreak{} \$\allowbreak{}f\allowbreak{} =\allowbreak{}\allowbreak{} pi\$\allowbreak{},\allowbreak{}\allowbreak{} with\allowbreak{} \$\allowbreak{}i\allowbreak{} \textbackslash{}\allowbreak{}in\allowbreak{} W\textasciicircum{}\allowbreak{}\{\allowbreak{}\textbackslash{}\allowbreak{}mathrm\{\allowbreak{}u\}\allowbreak{}\}\allowbreak{}\_\allowbreak{}S\$\allowbreak{},\allowbreak{}\allowbreak{} \$\allowbreak{}g\allowbreak{} \textbackslash{}\allowbreak{}in\allowbreak{} \textbackslash{}\allowbreak{}mathcal\{\allowbreak{}G\}\allowbreak{}\$\allowbreak{}}{factoring\allowbreak{} as\allowbreak{} \$\allowbreak{}f\allowbreak{} =\allowbreak{}\allowbreak{} gi\$\allowbreak{},\allowbreak{}\allowbreak{} with\allowbreak{} \$\allowbreak{}i\allowbreak{} \textbackslash{}\allowbreak{}in\allowbreak{} W\textasciicircum{}\allowbreak{}\{\allowbreak{}\textbackslash{}\allowbreak{}mathrm\{\allowbreak{}u\}\allowbreak{}\}\allowbreak{}\_\allowbreak{}S\$\allowbreak{},\allowbreak{}\allowbreak{} \$\allowbreak{}g\allowbreak{} \textbackslash{}\allowbreak{}in\allowbreak{} \textbackslash{}\allowbreak{}mathcal\{\allowbreak{}G\}\allowbreak{}\$\allowbreak{}}{DER\_\allowbreak{}XVI:\allowbreak{}53}
\DERapparatusitem{ENV0962}{31;\allowbreak{}\allowbreak{} running\allowbreak{} text\allowbreak{} or\allowbreak{} formula}{\textbackslash{}\allowbreak{}text\{\allowbreak{}def\}\allowbreak{}}{\textbackslash{}\allowbreak{}mathrm\{\allowbreak{}def\}\allowbreak{}}{TYPESETTING:\allowbreak{}7;\allowbreak{}\allowbreak{} NOTATION\_\allowbreak{}DER:\allowbreak{}17}
\DERapparatusitem{ENV0963}{31;\allowbreak{}\allowbreak{} running\allowbreak{} text\allowbreak{} or\allowbreak{} formula}{\textbackslash{}\allowbreak{}text\{\allowbreak{}def\}\allowbreak{}}{\textbackslash{}\allowbreak{}mathrm\{\allowbreak{}def\}\allowbreak{}}{TYPESETTING:\allowbreak{}7;\allowbreak{}\allowbreak{} NOTATION\_\allowbreak{}DER:\allowbreak{}17}
\DERapparatusitem{ENV0964}{32;\allowbreak{}\allowbreak{} ms56\allowbreak{} proof\allowbreak{} of\allowbreak{} Corollary4.\allowbreak{}4.\allowbreak{}11}{\$\allowbreak{}f\allowbreak{} \textbackslash{}\allowbreak{}in\allowbreak{} \textbackslash{}\allowbreak{}Fib\_\allowbreak{}\{\allowbreak{}W\_\allowbreak{}0\}\allowbreak{}\$\allowbreak{}}{\$\allowbreak{}f\allowbreak{} \textbackslash{}\allowbreak{}in\allowbreak{} \textbackslash{}\allowbreak{}Fibf\_\allowbreak{}\{\allowbreak{}W\_\allowbreak{}0\}\allowbreak{}\$\allowbreak{}}{DER\_\allowbreak{}XVI:\allowbreak{}55}
\DERapparatusitem{ENV0965}{32;\allowbreak{}\allowbreak{} ms56\allowbreak{} proof\allowbreak{} of\allowbreak{} Corollary4.\allowbreak{}4.\allowbreak{}11}{}{\{\allowbreak{}}{DER\_\allowbreak{}XVI:\allowbreak{}56}
\DERapparatusitem{ENV0966}{32;\allowbreak{}\allowbreak{} ms56\allowbreak{} proof\allowbreak{} of\allowbreak{} Corollary4.\allowbreak{}4.\allowbreak{}11}{}{\textbackslash{}\allowbreak{},\allowbreak{}S\}\allowbreak{}}{DER\_\allowbreak{}XVI:\allowbreak{}56}
\DERapparatusitem{ENV0967}{33;\allowbreak{}\allowbreak{} ms57\allowbreak{} proof\allowbreak{} of\allowbreak{} Corollary4.\allowbreak{}4.\allowbreak{}12}{And\allowbreak{} the\allowbreak{} arrow\allowbreak{} \$\allowbreak{}i\_\allowbreak{}\textbackslash{}\allowbreak{}infty\$\allowbreak{}\allowbreak{} is\allowbreak{} in\allowbreak{} \$\allowbreak{}W\textasciicircum{}\allowbreak{}\{\allowbreak{}\textbackslash{}\allowbreak{}mathrm\{\allowbreak{}u\}\allowbreak{}\}\allowbreak{}\_\allowbreak{}S\$\allowbreak{},\allowbreak{}\allowbreak{} for\allowbreak{} every\allowbreak{} localizer\allowbreak{} satisfying\allowbreak{} Loc\allowbreak{} (\allowbreak{}8\allowbreak{} bis)\allowbreak{},\allowbreak{}\allowbreak{} hence\allowbreak{} for\allowbreak{} \$\allowbreak{}W\_\allowbreak{}0\$\allowbreak{}.\allowbreak{}}{And\allowbreak{} the\allowbreak{} arrow\allowbreak{} \$\allowbreak{}i\_\allowbreak{}\textbackslash{}\allowbreak{}infty\$\allowbreak{}\allowbreak{} is\allowbreak{} in\allowbreak{} \$\allowbreak{}W\textasciicircum{}\allowbreak{}\{\allowbreak{}\textbackslash{}\allowbreak{}mathrm\{\allowbreak{}u\}\allowbreak{}\}\allowbreak{}\_\allowbreak{}S\$\allowbreak{},\allowbreak{}\allowbreak{} since\allowbreak{} the\allowbreak{} arrow\allowbreak{} \$\allowbreak{}f\$\allowbreak{}\allowbreak{} under\allowbreak{} consideration\allowbreak{} is\allowbreak{} in\allowbreak{} \$\allowbreak{}\textbackslash{}\allowbreak{}Fibf\_\allowbreak{}W\$\allowbreak{}.\allowbreak{}}{DER\_\allowbreak{}XVI:\allowbreak{}57}
\DERapparatusitem{ENV0968}{35;\allowbreak{}\allowbreak{} ms61\allowbreak{} end\allowbreak{} of\allowbreak{} theorem4.\allowbreak{}4.\allowbreak{}13(\allowbreak{}3)\allowbreak{}}{of\allowbreak{} the\allowbreak{} axiom\}\allowbreak{}\allowbreak{} Loc\allowbreak{} (\allowbreak{}6\allowbreak{} a)\allowbreak{})\allowbreak{}}{of\allowbreak{} the\allowbreak{} axiom\}\allowbreak{}\allowbreak{} Loc\allowbreak{} (\allowbreak{}7\allowbreak{} a)\allowbreak{})\allowbreak{}}{DER\_\allowbreak{}XVI:\allowbreak{}58}
\DERapparatusitem{ENV0969}{35;\allowbreak{}\allowbreak{} ms61\allowbreak{} theorem4.\allowbreak{}4.\allowbreak{}13(\allowbreak{}4a)\allowbreak{}}{If\allowbreak{} \$\allowbreak{}W\_\allowbreak{}0\$\allowbreak{}\allowbreak{} satisfies\allowbreak{} Loc\allowbreak{} (\allowbreak{}7\allowbreak{} F)\allowbreak{}}{If\allowbreak{} \$\allowbreak{}W\_\allowbreak{}0\$\allowbreak{}\allowbreak{} satisfies\allowbreak{} Loc\allowbreak{} (\allowbreak{}7)\allowbreak{}}{DER\_\allowbreak{}XVI:\allowbreak{}59}
\DERapparatusitem{ENV0970}{36;\allowbreak{}\allowbreak{} ms63\allowbreak{} diagram\allowbreak{} annotation}{}{,\allowbreak{}\allowbreak{} 3}{DER\_\allowbreak{}XVI:\allowbreak{}61}
\DERapparatusitem{ENV0971}{37;\allowbreak{}\allowbreak{} ms65\allowbreak{} prose}{one\allowbreak{} that\allowbreak{} it\allowbreak{} is\allowbreak{} almost\allowbreak{} obvious\allowbreak{} that\allowbreak{} it\allowbreak{} does}{of\allowbreak{} which\allowbreak{} it\allowbreak{} is\allowbreak{} almost\allowbreak{} obvious\allowbreak{} that\allowbreak{} it\allowbreak{} does}{DER\_\allowbreak{}XVI:\allowbreak{}65}
\DERapparatusitem{ENV0972}{41;\allowbreak{}\allowbreak{} running\allowbreak{} text\allowbreak{} or\allowbreak{} formula}{\textbackslash{}\allowbreak{}text\{\allowbreak{}def\}\allowbreak{}}{\textbackslash{}\allowbreak{}mathrm\{\allowbreak{}def\}\allowbreak{}}{TYPESETTING:\allowbreak{}7;\allowbreak{}\allowbreak{} NOTATION\_\allowbreak{}DER:\allowbreak{}17}
\DERapparatusitem{ENV0973}{42;\allowbreak{}\allowbreak{} running\allowbreak{} text\allowbreak{} or\allowbreak{} formula}{\textbackslash{}\allowbreak{}text\{\allowbreak{}def\}\allowbreak{}}{\textbackslash{}\allowbreak{}mathrm\{\allowbreak{}def\}\allowbreak{}}{TYPESETTING:\allowbreak{}7;\allowbreak{}\allowbreak{} NOTATION\_\allowbreak{}DER:\allowbreak{}17}
\DERapparatusitem{ENV0974}{42;\allowbreak{}\allowbreak{} ms78;\allowbreak{}\allowbreak{} order\allowbreak{} on\allowbreak{} pairs\allowbreak{} of\allowbreak{} subcategories}{Y}{X}{DER\_\allowbreak{}XVI:\allowbreak{}70}
\DERapparatusitem{ENV0975}{42;\allowbreak{}\allowbreak{} ms78;\allowbreak{}\allowbreak{} order\allowbreak{} on\allowbreak{} pairs\allowbreak{} of\allowbreak{} subcategories}{i}{j}{DER\_\allowbreak{}XVI:\allowbreak{}70}
\DERapparatusitem{ENV0976}{42;\allowbreak{}\allowbreak{} ms78;\allowbreak{}\allowbreak{} order\allowbreak{} on\allowbreak{} pairs\allowbreak{} of\allowbreak{} subcategories}{X}{Y}{DER\_\allowbreak{}XVI:\allowbreak{}70}
\DERapparatusitem{ENV0977}{42;\allowbreak{}\allowbreak{} ms78;\allowbreak{}\allowbreak{} order\allowbreak{} on\allowbreak{} pairs\allowbreak{} of\allowbreak{} subcategories}{j}{i}{DER\_\allowbreak{}XVI:\allowbreak{}70}
\DERapparatusitem{ENV0978}{42;\allowbreak{}\allowbreak{} ms78–79;\allowbreak{}\allowbreak{} Corollary5.\allowbreak{}3.\allowbreak{}2\allowbreak{} parenthesis}{(\allowbreak{}with\allowbreak{} \$\allowbreak{}f(\allowbreak{}X\_\allowbreak{}i)\allowbreak{}\allowbreak{} \textbackslash{}\allowbreak{}subset\allowbreak{} Y\_\allowbreak{}i\$\allowbreak{})\allowbreak{},\allowbreak{}\allowbreak{} but\allowbreak{} in\allowbreak{} fact}{(\allowbreak{}with\allowbreak{} \$\allowbreak{}f(\allowbreak{}X\_\allowbreak{}i)\allowbreak{}\allowbreak{} \textbackslash{}\allowbreak{}subset\allowbreak{} Y\_\allowbreak{}i\$\allowbreak{},\allowbreak{}\allowbreak{} but\allowbreak{} in\allowbreak{} fact}{DER\_\allowbreak{}XVI:\allowbreak{}71}
\DERapparatusitem{ENV0979}{44;\allowbreak{}\allowbreak{} ms82;\allowbreak{}\allowbreak{} direction\allowbreak{} of\allowbreak{} the\allowbreak{} chosen\allowbreak{} point}{Let\allowbreak{} \$\allowbreak{}x\$\allowbreak{}\allowbreak{} be\allowbreak{} an\allowbreak{} object\allowbreak{} (\allowbreak{}point)\allowbreak{}\allowbreak{} of\allowbreak{} \$\allowbreak{}X\$\allowbreak{},\allowbreak{}\allowbreak{} \$\allowbreak{}x\_\allowbreak{}i\$\allowbreak{}\allowbreak{} its\allowbreak{} image\allowbreak{} in\allowbreak{} \$\allowbreak{}X\_\allowbreak{}i\$\allowbreak{}.\allowbreak{}}{Let\allowbreak{} \$\allowbreak{}x\$\allowbreak{}\allowbreak{} be\allowbreak{} an\allowbreak{} object\allowbreak{} (\allowbreak{}point)\allowbreak{}\allowbreak{} of\allowbreak{} \$\allowbreak{}X\$\allowbreak{},\allowbreak{}\allowbreak{} represented\allowbreak{} by\allowbreak{} a\allowbreak{} point\allowbreak{} \$\allowbreak{}x\_\allowbreak{}\{\allowbreak{}i\_\allowbreak{}0\}\allowbreak{}\$\allowbreak{}\allowbreak{} of\allowbreak{} \$\allowbreak{}X\_\allowbreak{}\{\allowbreak{}i\_\allowbreak{}0\}\allowbreak{}\$\allowbreak{};\allowbreak{}\allowbreak{} after\allowbreak{} restricting\allowbreak{} to\allowbreak{} the\allowbreak{} cofinal\allowbreak{} category\allowbreak{} \$\allowbreak{}i\_\allowbreak{}0\textbackslash{}\allowbreak{}backslash\allowbreak{} I\$\allowbreak{},\allowbreak{}\allowbreak{} let\allowbreak{} \$\allowbreak{}x\_\allowbreak{}i\$\allowbreak{}\allowbreak{} be\allowbreak{} its\allowbreak{} image\allowbreak{} in\allowbreak{} \$\allowbreak{}X\_\allowbreak{}i\$\allowbreak{}.\allowbreak{}}{DER\_\allowbreak{}XVI:\allowbreak{}74}
\DERapparatusitem{ENV0980}{44;\allowbreak{}\allowbreak{} running\allowbreak{} text\allowbreak{} or\allowbreak{} formula}{\textbackslash{}\allowbreak{}text\{\allowbreak{}def\}\allowbreak{}}{\textbackslash{}\allowbreak{}mathrm\{\allowbreak{}def\}\allowbreak{}}{TYPESETTING:\allowbreak{}7;\allowbreak{}\allowbreak{} NOTATION\_\allowbreak{}DER:\allowbreak{}17}
\DERapparatusitem{ENV0981}{45;\allowbreak{}\allowbreak{} ms86;\allowbreak{}\allowbreak{} heading6.\allowbreak{}4\allowbreak{} agreement}{fibrations\allowbreak{} of\allowbreak{} the\allowbreak{} types\}\allowbreak{}}{fibrations\allowbreak{} of\allowbreak{} the\allowbreak{} type\}\allowbreak{}}{DER\_\allowbreak{}XVI:\allowbreak{}75}
\DERapparatusitem{ENV0982}{45;\allowbreak{}\allowbreak{} ms86;\allowbreak{}\allowbreak{} exception\allowbreak{} in6.\allowbreak{}4(\allowbreak{}a)\allowbreak{}}{D\$\allowbreak{}\_\allowbreak{}\{\allowbreak{}W\_\allowbreak{}\textbackslash{}\allowbreak{}omega\}\allowbreak{}\$\allowbreak{}\allowbreak{} itself}{\$\allowbreak{}\textbackslash{}\allowbreak{}mathbf\{\allowbreak{}W\}\allowbreak{}\_\allowbreak{}\textbackslash{}\allowbreak{}infty\$\allowbreak{}\allowbreak{} itself}{DER\_\allowbreak{}XVI:\allowbreak{}76}
\DERapparatusitem{ENV0983}{47;\allowbreak{}\allowbreak{} running\allowbreak{} text\allowbreak{} or\allowbreak{} formula}{\textbackslash{}\allowbreak{}text\{\allowbreak{}def\}\allowbreak{}}{\textbackslash{}\allowbreak{}mathrm\{\allowbreak{}def\}\allowbreak{}}{TYPESETTING:\allowbreak{}7;\allowbreak{}\allowbreak{} NOTATION\_\allowbreak{}DER:\allowbreak{}17}
\DERapparatusitem{ENV0984}{48;\allowbreak{}\allowbreak{} ms95;\allowbreak{}\allowbreak{} incomplete\allowbreak{} printed\allowbreak{} subdivision\allowbreak{} drawing}{\textbackslash{}\allowbreak{}draw\allowbreak{} (\allowbreak{}0,\allowbreak{}0)\allowbreak{}\allowbreak{} --\allowbreak{} (\allowbreak{}0.\allowbreak{}95,\allowbreak{}0.\allowbreak{}95)\allowbreak{}\allowbreak{} --\allowbreak{} (\allowbreak{}1.\allowbreak{}9,\allowbreak{}0)\allowbreak{};\allowbreak{}}{\textbackslash{}\allowbreak{}draw\allowbreak{} (\allowbreak{}0,\allowbreak{}0)\allowbreak{}\allowbreak{} --\allowbreak{} (\allowbreak{}0.\allowbreak{}95,\allowbreak{}0.\allowbreak{}95)\allowbreak{}\allowbreak{} --\allowbreak{} (\allowbreak{}1.\allowbreak{}9,\allowbreak{}0)\allowbreak{}\allowbreak{} --\allowbreak{} cycle;\allowbreak{}}{DER\_\allowbreak{}XVI:\allowbreak{}81}
\DERapparatusitem{ENV0985}{48;\allowbreak{}\allowbreak{} ms97;\allowbreak{}\allowbreak{} unmatched\allowbreak{} quotation\allowbreak{} mark}{a\allowbreak{} theory\allowbreak{} of\allowbreak{} homotopy’,\allowbreak{}}{a\allowbreak{} theory\allowbreak{} of\allowbreak{} homotopy,\allowbreak{}}{DER\_\allowbreak{}XVI:\allowbreak{}82}
\DERapparatusitem{ENV0986}{50;\allowbreak{}\allowbreak{} ms103;\allowbreak{}\allowbreak{} footnote\allowbreak{} marks48–49\allowbreak{} versus\allowbreak{} texts49–50}{\textbackslash{}\allowbreak{}text\{\allowbreak{}\textbackslash{}\allowbreak{}footnotemark[\allowbreak{}48]\allowbreak{}\textbackslash{}\allowbreak{},\allowbreak{}\textbackslash{}\allowbreak{}footnotemark[\allowbreak{}49]\allowbreak{}\}\allowbreak{}}{\textbackslash{}\allowbreak{}text\{\allowbreak{}\textbackslash{}\allowbreak{}footnotemark[\allowbreak{}49]\allowbreak{}\textbackslash{}\allowbreak{},\allowbreak{}\textbackslash{}\allowbreak{}footnotemark[\allowbreak{}50]\allowbreak{}\}\allowbreak{}}{DER\_\allowbreak{}XVI:\allowbreak{}86}
\DERapparatusitem{ENV0987}{54;\allowbreak{}\allowbreak{} ms114;\allowbreak{}\allowbreak{} relative\allowbreak{} category\allowbreak{} notation}{W}{\textbackslash{}\allowbreak{}Sigma}{DER\_\allowbreak{}XVI:\allowbreak{}92}
\DERapparatusitem{ENV0988}{54;\allowbreak{}\allowbreak{} ms114;\allowbreak{}\allowbreak{} derivator\allowbreak{} index\allowbreak{} typography}{}{\{\allowbreak{}\textbackslash{}\allowbreak{}mathbf\{\allowbreak{}}{DER\_\allowbreak{}XVI:\allowbreak{}91}
\DERapparatusitem{ENV0989}{54;\allowbreak{}\allowbreak{} ms114;\allowbreak{}\allowbreak{} derivator\allowbreak{} index\allowbreak{} typography}{}{\}\allowbreak{}\}\allowbreak{}}{DER\_\allowbreak{}XVI:\allowbreak{}91}
\DERapparatusitem{ENV0990}{54;\allowbreak{}\allowbreak{} ms114;\allowbreak{}\allowbreak{} axiom\allowbreak{} prefix}{\$\allowbreak{}\textbackslash{}\allowbreak{}mathbf\{\allowbreak{}D\}\allowbreak{}\$\allowbreak{}\allowbreak{} satisfies\allowbreak{} Loc\allowbreak{} (\allowbreak{}4\$\allowbreak{}\textasciicircum{}\allowbreak{}*\$\allowbreak{},\allowbreak{}\allowbreak{} 5\$\allowbreak{}\textasciicircum{}\allowbreak{}*\$\allowbreak{})\allowbreak{}}{\$\allowbreak{}\textbackslash{}\allowbreak{}mathbf\{\allowbreak{}D\}\allowbreak{}\$\allowbreak{}\allowbreak{} satisfies\allowbreak{} Der\allowbreak{} (\allowbreak{}4\$\allowbreak{}\textasciicircum{}\allowbreak{}*\$\allowbreak{},\allowbreak{}\allowbreak{} 5\$\allowbreak{}\textasciicircum{}\allowbreak{}*\$\allowbreak{})\allowbreak{}}{DER\_\allowbreak{}XVI:\allowbreak{}93}
\DERapparatusitem{ENV0991}{54;\allowbreak{}\allowbreak{} ms115;\allowbreak{}\allowbreak{} proof\allowbreak{} transition}{\textbackslash{}\allowbreak{}textit\{\allowbreak{}we\allowbreak{} may\allowbreak{} assume\}\allowbreak{}\allowbreak{} that\allowbreak{} if\allowbreak{} \$\allowbreak{}S\$\allowbreak{}}{\textbackslash{}\allowbreak{}textit\{\allowbreak{}we\allowbreak{} must\allowbreak{} prove\}\allowbreak{}\allowbreak{} that\allowbreak{} if\allowbreak{} \$\allowbreak{}S\$\allowbreak{}}{DER\_\allowbreak{}XVI:\allowbreak{}94}
\DERapparatusitem{ENV0992}{54;\allowbreak{}\allowbreak{} ms116;\allowbreak{}\allowbreak{} localization\allowbreak{} subscripts}{W}{\textbackslash{}\allowbreak{}Sigma}{DER\_\allowbreak{}XVI:\allowbreak{}95}
\DERapparatusitem{ENV0993}{54;\allowbreak{}\allowbreak{} ms116;\allowbreak{}\allowbreak{} localization\allowbreak{} subscripts}{W}{\textbackslash{}\allowbreak{}Sigma}{DER\_\allowbreak{}XVI:\allowbreak{}95}
\subsection*{Chapter XVII}
\DERapparatusitem{ENV0994}{1;\allowbreak{}\allowbreak{} ms\allowbreak{} p1,\allowbreak{}\allowbreak{} corollary,\allowbreak{}\allowbreak{} typing\allowbreak{} of\allowbreak{} W}{\textbackslash{}\allowbreak{}mathcal\{\allowbreak{}M\}\allowbreak{}}{\textbackslash{}\allowbreak{}operatorname\{\allowbreak{}Fl\}\allowbreak{}(\allowbreak{}\textbackslash{}\allowbreak{}mathcal\allowbreak{} M)\allowbreak{}}{DER\_\allowbreak{}XVII-a:\allowbreak{}2}
\DERapparatusitem{ENV0995}{1;\allowbreak{}\allowbreak{} running\allowbreak{} text\allowbreak{} or\allowbreak{} formula}{\textbackslash{}\allowbreak{}text\{\allowbreak{}def\}\allowbreak{}}{\textbackslash{}\allowbreak{}mathrm\{\allowbreak{}def\}\allowbreak{}}{TYPESETTING:\allowbreak{}7;\allowbreak{}\allowbreak{} NOTATION\_\allowbreak{}DER:\allowbreak{}17}
\DERapparatusitem{ENV0996}{2;\allowbreak{}\allowbreak{} ms\allowbreak{} p3,\allowbreak{}\allowbreak{} reduction\allowbreak{} before\allowbreak{} condition\allowbreak{} (\allowbreak{}c)\allowbreak{}}{C}{W}{DER\_\allowbreak{}XVII-a:\allowbreak{}3}
\DERapparatusitem{ENV0997}{2;\allowbreak{}\allowbreak{} running\allowbreak{} text\allowbreak{} or\allowbreak{} formula}{\textbackslash{}\allowbreak{}text\{\allowbreak{}def\}\allowbreak{}}{\textbackslash{}\allowbreak{}mathrm\{\allowbreak{}def\}\allowbreak{}}{TYPESETTING:\allowbreak{}7;\allowbreak{}\allowbreak{} NOTATION\_\allowbreak{}DER:\allowbreak{}17}
\DERapparatusitem{ENV0998}{2;\allowbreak{}\allowbreak{} running\allowbreak{} text\allowbreak{} or\allowbreak{} formula}{\textbackslash{}\allowbreak{}text\{\allowbreak{}def\}\allowbreak{}}{\textbackslash{}\allowbreak{}mathrm\{\allowbreak{}def\}\allowbreak{}}{TYPESETTING:\allowbreak{}7;\allowbreak{}\allowbreak{} NOTATION\_\allowbreak{}DER:\allowbreak{}17}
\DERapparatusitem{ENV0999}{2;\allowbreak{}\allowbreak{} running\allowbreak{} text\allowbreak{} or\allowbreak{} formula}{\textbackslash{}\allowbreak{}text\{\allowbreak{}def\}\allowbreak{}}{\textbackslash{}\allowbreak{}mathrm\{\allowbreak{}def\}\allowbreak{}}{TYPESETTING:\allowbreak{}7;\allowbreak{}\allowbreak{} NOTATION\_\allowbreak{}DER:\allowbreak{}17}
\DERapparatusitem{ENV1000}{3;\allowbreak{}\allowbreak{} ms\allowbreak{} p6,\allowbreak{}\allowbreak{} parenthetical\allowbreak{} successor\allowbreak{} index}{\textbackslash{}\allowbreak{}beta}{i}{DER\_\allowbreak{}XVII-a:\allowbreak{}4}
\DERapparatusitem{ENV1001}{3;\allowbreak{}\allowbreak{} ms\allowbreak{} p6,\allowbreak{}\allowbreak{} exceptional\allowbreak{} initial\allowbreak{} attachment}{0}{-1}{DER\_\allowbreak{}XVII-a:\allowbreak{}6}
\DERapparatusitem{ENV1002}{3;\allowbreak{}\allowbreak{} running\allowbreak{} text\allowbreak{} or\allowbreak{} formula}{\textbackslash{}\allowbreak{}text\{\allowbreak{}def\}\allowbreak{}}{\textbackslash{}\allowbreak{}mathrm\{\allowbreak{}def\}\allowbreak{}}{TYPESETTING:\allowbreak{}7;\allowbreak{}\allowbreak{} NOTATION\_\allowbreak{}DER:\allowbreak{}17}
\DERapparatusitem{ENV1003}{preamble;\allowbreak{}\allowbreak{} local\allowbreak{} macro\allowbreak{} definition}{\textbackslash{}\allowbreak{}textup\{\allowbreak{}def\}\allowbreak{}}{\textbackslash{}\allowbreak{}mathrm\{\allowbreak{}def\}\allowbreak{}}{TYPESETTING:\allowbreak{}7;\allowbreak{}\allowbreak{} NOTATION\_\allowbreak{}DER:\allowbreak{}17}
\DERapparatusitem{ENV1004}{4;\allowbreak{}\allowbreak{} ms\allowbreak{} p8,\allowbreak{}\allowbreak{} first\allowbreak{} corollary,\allowbreak{}\allowbreak{} condition\allowbreak{} 2°}{T}{C}{DER\_\allowbreak{}XVII-a:\allowbreak{}7}
\DERapparatusitem{ENV1005}{5;\allowbreak{}\allowbreak{} ms\allowbreak{} p8\allowbreak{} continuation,\allowbreak{}\allowbreak{} applying\allowbreak{} Proposition\allowbreak{} 2}{X}{C}{DER\_\allowbreak{}XVII-a:\allowbreak{}9}
\DERapparatusitem{ENV1006}{8;\allowbreak{}\allowbreak{} ms\allowbreak{} p13,\allowbreak{}\allowbreak{} first\allowbreak{} paragraph,\allowbreak{}\allowbreak{} objectwise\allowbreak{} weak-equivalence\allowbreak{} assumption}{0\allowbreak{} \textbackslash{}\allowbreak{}in\allowbreak{} W\_\allowbreak{}\{\allowbreak{}I\_\allowbreak{}0\}\allowbreak{}}{\textbackslash{}\allowbreak{}bullet\textbackslash{}\allowbreak{}in\allowbreak{} W\_\allowbreak{}\{\allowbreak{}I\textasciicircum{}\allowbreak{}\{\allowbreak{}\textbackslash{}\allowbreak{}circ\}\allowbreak{}\}\allowbreak{}}{DER\_\allowbreak{}XVII-a:\allowbreak{}10}
\DERapparatusitem{ENV1007}{10;\allowbreak{}\allowbreak{} ms\allowbreak{} p16,\allowbreak{}\allowbreak{} footnote\allowbreak{} 9}{diagram\allowbreak{} (\allowbreak{}6)\allowbreak{}\allowbreak{} above}{diagram\allowbreak{} (\allowbreak{}6)\allowbreak{}\allowbreak{} below}{DER\_\allowbreak{}XVII-a:\allowbreak{}15}
\DERapparatusitem{ENV1008}{11;\allowbreak{}\allowbreak{} running\allowbreak{} text\allowbreak{} or\allowbreak{} formula}{\textbackslash{}\allowbreak{}textup\{\allowbreak{}def\}\allowbreak{}}{\textbackslash{}\allowbreak{}mathrm\{\allowbreak{}def\}\allowbreak{}}{TYPESETTING:\allowbreak{}7;\allowbreak{}\allowbreak{} NOTATION\_\allowbreak{}DER:\allowbreak{}17}
\DERapparatusitem{ENV1009}{11;\allowbreak{}\allowbreak{} ms\allowbreak{} p19,\allowbreak{}\allowbreak{} induction\allowbreak{} branch}{'}{}{DER\_\allowbreak{}XVII-a:\allowbreak{}17}
\DERapparatusitem{ENV1010}{14;\allowbreak{}\allowbreak{} ms\allowbreak{} p24,\allowbreak{}\allowbreak{} second\allowbreak{} case\allowbreak{} in\allowbreak{} the\allowbreak{} lemma\allowbreak{} proof}{=\allowbreak{}}{\textbackslash{}\allowbreak{}neq}{DER\_\allowbreak{}XVII-a:\allowbreak{}20}
\DERapparatusitem{ENV1011}{14;\allowbreak{}\allowbreak{} running\allowbreak{} text\allowbreak{} or\allowbreak{} formula}{\textbackslash{}\allowbreak{}text\{\allowbreak{}def\}\allowbreak{}}{\textbackslash{}\allowbreak{}mathrm\{\allowbreak{}def\}\allowbreak{}}{TYPESETTING:\allowbreak{}7;\allowbreak{}\allowbreak{} NOTATION\_\allowbreak{}DER:\allowbreak{}17}
\DERapparatusitem{ENV1012}{14;\allowbreak{}\allowbreak{} ms\allowbreak{} p24,\allowbreak{}\allowbreak{} progressive\allowbreak{} condition\allowbreak{} in\allowbreak{} corollary\allowbreak{} proof}{/\allowbreak{}}{\_\allowbreak{}\{\allowbreak{}<}{DER\_\allowbreak{}XVII-a:\allowbreak{}22}
\DERapparatusitem{ENV1013}{14;\allowbreak{}\allowbreak{} ms\allowbreak{} p24,\allowbreak{}\allowbreak{} progressive\allowbreak{} condition\allowbreak{} in\allowbreak{} corollary\allowbreak{} proof}{}{\}\allowbreak{}}{DER\_\allowbreak{}XVII-a:\allowbreak{}22}
\DERapparatusitem{ENV1014}{14;\allowbreak{}\allowbreak{} ms\allowbreak{} p24,\allowbreak{}\allowbreak{} descending-chain\allowbreak{} argument}{0}{1}{DER\_\allowbreak{}XVII-a:\allowbreak{}23}
\DERapparatusitem{ENV1015}{15;\allowbreak{}\allowbreak{} ms26,\allowbreak{}\allowbreak{} Theorem2.\allowbreak{}1(\allowbreak{}1)\allowbreak{},\allowbreak{}\allowbreak{} reference}{\&\allowbreak{}\allowbreak{} Y\_\allowbreak{}J\allowbreak{} \textbackslash{}\allowbreak{}\allowbreak{} .\allowbreak{}\space{}\allowbreak{}\textbackslash{}\allowbreak{}end\{\allowbreak{}tikzcd\}\allowbreak{}\space{}\allowbreak{}\textbackslash{}\allowbreak{}]\allowbreak{}}{\&\allowbreak{}\allowbreak{} Y\_\allowbreak{}J\allowbreak{} \textbackslash{}\allowbreak{}\allowbreak{} .\allowbreak{}\space{}\allowbreak{}\textbackslash{}\allowbreak{}end\{\allowbreak{}tikzcd\}\allowbreak{}\space{}\allowbreak{}\textbackslash{}\allowbreak{}tag*\{\allowbreak{}(\allowbreak{}2.\allowbreak{}4)\allowbreak{}\}\allowbreak{}\space{}\allowbreak{}\textbackslash{}\allowbreak{}]\allowbreak{}}{DER\_\allowbreak{}XVII-a:\allowbreak{}24}
\DERapparatusitem{ENV1016}{15;\allowbreak{}\allowbreak{} ms26,\allowbreak{}\allowbreak{} Theorem2.\allowbreak{}1(\allowbreak{}2)\allowbreak{},\allowbreak{}\allowbreak{} parenthetical\allowbreak{} general\allowbreak{} pair}{J}{K}{DER\_\allowbreak{}XVII-a:\allowbreak{}25}
\DERapparatusitem{ENV1017}{15;\allowbreak{}\allowbreak{} ms26,\allowbreak{}\allowbreak{} Theorem2.\allowbreak{}1(\allowbreak{}2)\allowbreak{},\allowbreak{}\allowbreak{} parenthetical\allowbreak{} general\allowbreak{} pair}{K}{J}{DER\_\allowbreak{}XVII-a:\allowbreak{}25}
\DERapparatusitem{ENV1018}{15;\allowbreak{}\allowbreak{} running\allowbreak{} text\allowbreak{} or\allowbreak{} formula}{\textbackslash{}\allowbreak{}text\{\allowbreak{}def\}\allowbreak{}}{\textbackslash{}\allowbreak{}mathrm\{\allowbreak{}def\}\allowbreak{}}{TYPESETTING:\allowbreak{}7;\allowbreak{}\allowbreak{} NOTATION\_\allowbreak{}DER:\allowbreak{}17}
\DERapparatusitem{ENV1019}{15;\allowbreak{}\allowbreak{} running\allowbreak{} text\allowbreak{} or\allowbreak{} formula}{\textbackslash{}\allowbreak{}text\{\allowbreak{}def\}\allowbreak{}}{\textbackslash{}\allowbreak{}mathrm\{\allowbreak{}def\}\allowbreak{}}{TYPESETTING:\allowbreak{}7;\allowbreak{}\allowbreak{} NOTATION\_\allowbreak{}DER:\allowbreak{}17}
\DERapparatusitem{ENV1020}{17;\allowbreak{}\allowbreak{} ms31,\allowbreak{}\allowbreak{} sentence\allowbreak{} introducing\allowbreak{} Lemma2.\allowbreak{}6}{a\allowbreak{} lemma\allowbreak{} on\allowbreak{} homomorphisms\allowbreak{} of\allowbreak{} cartesian\allowbreak{} squares}{a\allowbreak{} lemma\allowbreak{} on\allowbreak{} homomorphisms\allowbreak{} of\allowbreak{} cocartesian\allowbreak{} squares}{DER\_\allowbreak{}XVII-a:\allowbreak{}26}
\DERapparatusitem{ENV1021}{18;\allowbreak{}\allowbreak{} ms33,\allowbreak{}\allowbreak{} same\allowbreak{} final\allowbreak{} sentence,\allowbreak{}\allowbreak{} retained\allowbreak{} source\allowbreak{} wording}{being\allowbreak{} cartesian,\allowbreak{}\allowbreak{} the\allowbreak{} same\allowbreak{} holds}{being\allowbreak{} cocartesian,\allowbreak{}\allowbreak{} the\allowbreak{} same\allowbreak{} holds}{DER\_\allowbreak{}XVII-a:\allowbreak{}29}
\DERapparatusitem{ENV1022}{18;\allowbreak{}\allowbreak{} ms33,\allowbreak{}\allowbreak{} Lemma2.\allowbreak{}7\allowbreak{} statement}{-1}{0}{DER\_\allowbreak{}XVII-a:\allowbreak{}30}
\DERapparatusitem{ENV1023}{18;\allowbreak{}\allowbreak{} ms33,\allowbreak{}\allowbreak{} Lemma2.\allowbreak{}7\allowbreak{} statement}{\allowbreak{} (\allowbreak{}and\allowbreak{} from\allowbreak{} \$\allowbreak{}u\_\allowbreak{}\{\allowbreak{}I\_\allowbreak{}0\}\allowbreak{}\allowbreak{} :\allowbreak{}\allowbreak{} X\_\allowbreak{}\{\allowbreak{}I\_\allowbreak{}0\}\allowbreak{}\allowbreak{} \textbackslash{}\allowbreak{}longrightarrow\allowbreak{} Y\_\allowbreak{}\{\allowbreak{}I\_\allowbreak{}0\}\allowbreak{}\$\allowbreak{}\allowbreak{} for\allowbreak{} \$\allowbreak{}\textbackslash{}\allowbreak{}alpha\allowbreak{} =\allowbreak{}\allowbreak{} 0\$\allowbreak{})\allowbreak{}}{}{DER\_\allowbreak{}XVII-a:\allowbreak{}30}
\DERapparatusitem{ENV1024}{19;\allowbreak{}\allowbreak{} ms33–34,\allowbreak{}\allowbreak{} endpoint\allowbreak{} identities}{\textbackslash{}\allowbreak{}alpha}{0}{DER\_\allowbreak{}XVII-a:\allowbreak{}31}
\DERapparatusitem{ENV1025}{19;\allowbreak{}\allowbreak{} ms33–34,\allowbreak{}\allowbreak{} endpoint\allowbreak{} identities}{}{\{\allowbreak{}}{DER\_\allowbreak{}XVII-a:\allowbreak{}31}
\DERapparatusitem{ENV1026}{19;\allowbreak{}\allowbreak{} ms33–34,\allowbreak{}\allowbreak{} endpoint\allowbreak{} identities}{}{\_\allowbreak{}0\}\allowbreak{}}{DER\_\allowbreak{}XVII-a:\allowbreak{}31}
\DERapparatusitem{ENV1027}{19;\allowbreak{}\allowbreak{} ms33–34,\allowbreak{}\allowbreak{} endpoint\allowbreak{} identities}{}{\allowbreak{} \textbackslash{}\allowbreak{}simeq\allowbreak{} Y\_\allowbreak{}\{\allowbreak{}I\_\allowbreak{}\{\allowbreak{}\textbackslash{}\allowbreak{}alpha\_\allowbreak{}0\}\allowbreak{}\}\allowbreak{}}{DER\_\allowbreak{}XVII-a:\allowbreak{}31}
\DERapparatusitem{ENV1028}{19;\allowbreak{}\allowbreak{} ms34,\allowbreak{}\allowbreak{} top-right\allowbreak{} label\allowbreak{} in\allowbreak{} successor\allowbreak{} diagram}{rlap\{\allowbreak{}\$\allowbreak{}\{\allowbreak{}\}\allowbreak{}=\allowbreak{}\allowbreak{} Z\_\allowbreak{}0\$\allowbreak{}\}\allowbreak{}\allowbreak{} \textbackslash{}\allowbreak{}}{}{DER\_\allowbreak{}XVII-a:\allowbreak{}32}
\DERapparatusitem{ENV1029}{19;\allowbreak{}\allowbreak{} ms35,\allowbreak{}\allowbreak{} first\allowbreak{} paragraph\allowbreak{} relative-latching\allowbreak{} subscript}{\textbackslash{}\allowbreak{}alpha}{i}{DER\_\allowbreak{}XVII-a:\allowbreak{}33}
\DERapparatusitem{ENV1030}{19;\allowbreak{}\allowbreak{} ms35,\allowbreak{}\allowbreak{} Lemma2.\allowbreak{}8\allowbreak{} wording}{such\allowbreak{} that\allowbreak{} for\allowbreak{} \$\allowbreak{}u\_\allowbreak{}\textbackslash{}\allowbreak{}bullet|I/\allowbreak{}i\_\allowbreak{}0\$\allowbreak{}\allowbreak{} to\allowbreak{} be\allowbreak{} in\allowbreak{} \$\allowbreak{}C\_\allowbreak{}\{\allowbreak{}I/\allowbreak{}i\_\allowbreak{}0\}\allowbreak{}\$\allowbreak{}}{such\allowbreak{} that\allowbreak{} \$\allowbreak{}u\_\allowbreak{}\textbackslash{}\allowbreak{}bullet|I/\allowbreak{}i\_\allowbreak{}0\$\allowbreak{}\allowbreak{} is\allowbreak{} in\allowbreak{} \$\allowbreak{}C\_\allowbreak{}\{\allowbreak{}I/\allowbreak{}i\_\allowbreak{}0\}\allowbreak{}\$\allowbreak{}}{DER\_\allowbreak{}XVII-a:\allowbreak{}34}
\DERapparatusitem{ENV1031}{19;\allowbreak{}\allowbreak{} running\allowbreak{} text\allowbreak{} or\allowbreak{} formula}{\textbackslash{}\allowbreak{}text\{\allowbreak{}def\}\allowbreak{}}{\textbackslash{}\allowbreak{}mathrm\{\allowbreak{}def\}\allowbreak{}}{TYPESETTING:\allowbreak{}7;\allowbreak{}\allowbreak{} NOTATION\_\allowbreak{}DER:\allowbreak{}17}
\DERapparatusitem{ENV1032}{19;\allowbreak{}\allowbreak{} ms35–36,\allowbreak{}\allowbreak{} noetherian\allowbreak{} continuation}{Let\allowbreak{} us\allowbreak{} note\allowbreak{} that\allowbreak{} \$\allowbreak{}\textbackslash{}\allowbreak{}mathrm\{\allowbreak{}Esp\}\allowbreak{}(\allowbreak{}I)\allowbreak{}\$\allowbreak{}\allowbreak{} is\allowbreak{} noetherian}{Let\allowbreak{} us\allowbreak{} note\allowbreak{} that\allowbreak{} if\allowbreak{} \$\allowbreak{}\textbackslash{}\allowbreak{}mathrm\{\allowbreak{}Esp\}\allowbreak{}(\allowbreak{}I)\allowbreak{}\$\allowbreak{}\allowbreak{} is\allowbreak{} noetherian}{DER\_\allowbreak{}XVII-a:\allowbreak{}35}
\DERapparatusitem{ENV1033}{20;\allowbreak{}\allowbreak{} ms36,\allowbreak{}\allowbreak{} reference\allowbreak{} after\allowbreak{} pointwise\allowbreak{} TC\allowbreak{} conclusion}{(\allowbreak{}cf.\allowbreak{}\allowbreak{} cor.\allowbreak{}\allowbreak{} 2.\allowbreak{}1\allowbreak{} a)\allowbreak{}}{(\allowbreak{}cf.\allowbreak{}\allowbreak{} cor.\allowbreak{}\allowbreak{} 2.\allowbreak{}2\allowbreak{} a)\allowbreak{}}{DER\_\allowbreak{}XVII-a:\allowbreak{}36}
\DERapparatusitem{ENV1034}{20;\allowbreak{}\allowbreak{} ms36,\allowbreak{}\allowbreak{} hypothesis\allowbreak{} in\allowbreak{} proof\allowbreak{} of\allowbreak{} b⇒a}{}{\_\allowbreak{}I}{DER\_\allowbreak{}XVII-a:\allowbreak{}37}
\DERapparatusitem{ENV1035}{20;\allowbreak{}\allowbreak{} ms37,\allowbreak{}\allowbreak{} reference\allowbreak{} already\allowbreak{} flagged\allowbreak{} \%\allowbreak{} [\allowbreak{}sic]\allowbreak{}}{it\allowbreak{} follows\allowbreak{} from\allowbreak{} 3.\allowbreak{}9}{it\allowbreak{} follows\allowbreak{} from\allowbreak{} 2.\allowbreak{}9}{DER\_\allowbreak{}XVII-a:\allowbreak{}38}
\DERapparatusitem{ENV1036}{20;\allowbreak{}\allowbreak{} ms38,\allowbreak{}\allowbreak{} reference\allowbreak{} to\allowbreak{} the\allowbreak{} weaker\allowbreak{} condition}{the\allowbreak{} condition\allowbreak{} considered\allowbreak{} in\allowbreak{} 2.\allowbreak{}1.\allowbreak{}\allowbreak{} 1°}{the\allowbreak{} condition\allowbreak{} considered\allowbreak{} in\allowbreak{} 2.\allowbreak{}1.\allowbreak{}\allowbreak{} 2°}{DER\_\allowbreak{}XVII-a:\allowbreak{}40}
\DERapparatusitem{ENV1037}{20;\allowbreak{}\allowbreak{} ms38,\allowbreak{}\allowbreak{} transformation-class\allowbreak{} notation}{\$\allowbreak{}u\_\allowbreak{}\textbackslash{}\allowbreak{}bullet\allowbreak{} :\allowbreak{}\allowbreak{} X\_\allowbreak{}\textbackslash{}\allowbreak{}bullet\allowbreak{} \textbackslash{}\allowbreak{}longrightarrow\allowbreak{} Y\_\allowbreak{}\textbackslash{}\allowbreak{}bullet\$\allowbreak{}\allowbreak{} is\allowbreak{} in\allowbreak{} \$\allowbreak{}C\$\allowbreak{}}{\$\allowbreak{}u\_\allowbreak{}\textbackslash{}\allowbreak{}bullet\allowbreak{} :\allowbreak{}\allowbreak{} X\_\allowbreak{}\textbackslash{}\allowbreak{}bullet\allowbreak{} \textbackslash{}\allowbreak{}longrightarrow\allowbreak{} Y\_\allowbreak{}\textbackslash{}\allowbreak{}bullet\$\allowbreak{}\allowbreak{} is\allowbreak{} in\allowbreak{} \$\allowbreak{}C\_\allowbreak{}I\$\allowbreak{}}{DER\_\allowbreak{}XVII-a:\allowbreak{}41}
\DERapparatusitem{ENV1038}{21;\allowbreak{}\allowbreak{} running\allowbreak{} text\allowbreak{} or\allowbreak{} formula}{\textbackslash{}\allowbreak{}text\{\allowbreak{}def\}\allowbreak{}}{\textbackslash{}\allowbreak{}mathrm\{\allowbreak{}def\}\allowbreak{}}{TYPESETTING:\allowbreak{}7;\allowbreak{}\allowbreak{} NOTATION\_\allowbreak{}DER:\allowbreak{}17}
\DERapparatusitem{ENV1039}{22;\allowbreak{}\allowbreak{} ms41,\allowbreak{}\allowbreak{} italic\allowbreak{} M\allowbreak{} already\allowbreak{} flagged\allowbreak{} \%\allowbreak{} [\allowbreak{}sic]\allowbreak{}}{M}{\textbackslash{}\allowbreak{}M}{DER\_\allowbreak{}XVII-a:\allowbreak{}42}
\DERapparatusitem{ENV1040}{24;\allowbreak{}\allowbreak{} ms45,\allowbreak{}\allowbreak{} number\allowbreak{} of\allowbreak{} data}{let\allowbreak{} us\allowbreak{} consider\allowbreak{} the\allowbreak{} triples}{let\allowbreak{} us\allowbreak{} consider\allowbreak{} the\allowbreak{} quadruples}{DER\_\allowbreak{}XVII-a:\allowbreak{}44}
\DERapparatusitem{ENV1041}{24;\allowbreak{}\allowbreak{} ms45,\allowbreak{}\allowbreak{} factorization\allowbreak{} display}{\textbackslash{}\allowbreak{}xleftarrow}{\textbackslash{}\allowbreak{}xrightarrow}{DER\_\allowbreak{}XVII-a:\allowbreak{}45}
\DERapparatusitem{ENV1042}{24;\allowbreak{}\allowbreak{} ms45,\allowbreak{}\allowbreak{} union\allowbreak{} already\allowbreak{} flagged\allowbreak{} \%\allowbreak{} [\allowbreak{}sic]\allowbreak{}}{I}{J}{DER\_\allowbreak{}XVII-a:\allowbreak{}46}
\DERapparatusitem{ENV1043}{24;\allowbreak{}\allowbreak{} ms46,\allowbreak{}\allowbreak{} extension\allowbreak{} symbols\allowbreak{} already\allowbreak{} flagged\allowbreak{} \%\allowbreak{} [\allowbreak{}sic]\allowbreak{}}{j}{i}{DER\_\allowbreak{}XVII-a:\allowbreak{}48}
\DERapparatusitem{ENV1044}{24;\allowbreak{}\allowbreak{} ms46,\allowbreak{}\allowbreak{} intersection\allowbreak{} already\allowbreak{} flagged\allowbreak{} \%\allowbreak{} [\allowbreak{}sic]\allowbreak{}}{}{'}{DER\_\allowbreak{}XVII-a:\allowbreak{}49}
\DERapparatusitem{ENV1045}{24;\allowbreak{}\allowbreak{} ms46,\allowbreak{}\allowbreak{} reversed\allowbreak{} extension\allowbreak{} sentence}{those\allowbreak{} of\allowbreak{} \$\allowbreak{}(\allowbreak{}Z|K',\allowbreak{}i|K',\allowbreak{}p|K')\allowbreak{}\$\allowbreak{}\allowbreak{} to\allowbreak{} \$\allowbreak{}K\$\allowbreak{}}{those\allowbreak{} of\allowbreak{} \$\allowbreak{}(\allowbreak{}Z|K,\allowbreak{}i|K,\allowbreak{}p|K)\allowbreak{}\$\allowbreak{}\allowbreak{} to\allowbreak{} \$\allowbreak{}K'\$\allowbreak{}}{DER\_\allowbreak{}XVII-a:\allowbreak{}50}
\DERapparatusitem{ENV1046}{24;\allowbreak{}\allowbreak{} ms46,\allowbreak{}\allowbreak{} final-object\allowbreak{} name\allowbreak{} already\allowbreak{} flagged\allowbreak{} \%\allowbreak{} [\allowbreak{}sic]\allowbreak{}}{i}{e}{DER\_\allowbreak{}XVII-a:\allowbreak{}51}
\DERapparatusitem{ENV1047}{24;\allowbreak{}\allowbreak{} ms46–47,\allowbreak{}\allowbreak{} W-index\allowbreak{} already\allowbreak{} flagged\allowbreak{} \%\allowbreak{} [\allowbreak{}sic]\allowbreak{}}{}{'}{DER\_\allowbreak{}XVII-a:\allowbreak{}52}
\DERapparatusitem{ENV1048}{25;\allowbreak{}\allowbreak{} ms47,\allowbreak{}\allowbreak{} old\allowbreak{} factorization\allowbreak{} domain}{J}{I}{DER\_\allowbreak{}XVII-a:\allowbreak{}53}
\DERapparatusitem{ENV1049}{25;\allowbreak{}\allowbreak{} ms49,\allowbreak{}\allowbreak{} p\allowbreak{} symbol\allowbreak{} already\allowbreak{} flagged\allowbreak{} \%\allowbreak{} [\allowbreak{}sic]\allowbreak{}}{}{'}{DER\_\allowbreak{}XVII-a:\allowbreak{}54}
\DERapparatusitem{ENV1050}{26;\allowbreak{}\allowbreak{} ms51,\allowbreak{}\allowbreak{} Proposition3.\allowbreak{}2(\allowbreak{}a)\allowbreak{},\allowbreak{}\allowbreak{} local\allowbreak{} cofilteredness}{the\allowbreak{} category\allowbreak{} \$\allowbreak{}i'\textbackslash{}\allowbreak{}backslash\allowbreak{} I\$\allowbreak{}\allowbreak{} is\allowbreak{} locally\allowbreak{} filtering}{the\allowbreak{} category\allowbreak{} \$\allowbreak{}i'\textbackslash{}\allowbreak{}backslash\allowbreak{} I\$\allowbreak{}\allowbreak{} is\allowbreak{} locally\allowbreak{} cofiltering}{DER\_\allowbreak{}XVII-a:\allowbreak{}56}
\DERapparatusitem{ENV1051}{28;\allowbreak{}\allowbreak{} ms55,\allowbreak{}\allowbreak{} proof3.\allowbreak{}3,\allowbreak{}\allowbreak{} diagram\allowbreak{} cofibration\allowbreak{} notation}{(\allowbreak{}I\textasciicircum{}\allowbreak{}\textbackslash{}\allowbreak{}circ)\allowbreak{}}{\_\allowbreak{}\{\allowbreak{}I\textasciicircum{}\allowbreak{}\textbackslash{}\allowbreak{}circ\}\allowbreak{}}{DER\_\allowbreak{}XVII-a:\allowbreak{}59}
\DERapparatusitem{ENV1052}{28;\allowbreak{}\allowbreak{} ms56,\allowbreak{}\allowbreak{} proof3.\allowbreak{}3,\allowbreak{}\allowbreak{} final\allowbreak{} two-out-of-three\allowbreak{} step}{i}{p}{DER\_\allowbreak{}XVII-a:\allowbreak{}60}
\DERapparatusitem{ENV1053}{30;\allowbreak{}\allowbreak{} running\allowbreak{} text\allowbreak{} or\allowbreak{} formula}{\textbackslash{}\allowbreak{}text\{\allowbreak{}def\}\allowbreak{}}{\textbackslash{}\allowbreak{}mathrm\{\allowbreak{}def\}\allowbreak{}}{TYPESETTING:\allowbreak{}7;\allowbreak{}\allowbreak{} NOTATION\_\allowbreak{}DER:\allowbreak{}17}
\DERapparatusitem{ENV1054}{31;\allowbreak{}\allowbreak{} ms63,\allowbreak{}\allowbreak{} target\allowbreak{} of\allowbreak{} w\_\allowbreak{}!}{}{J\allowbreak{} =\allowbreak{}\allowbreak{} }{DER\_\allowbreak{}XVII-a:\allowbreak{}65;\allowbreak{}\allowbreak{} DER\_\allowbreak{}XVII-a:\allowbreak{}66;\allowbreak{}\allowbreak{} DER\_\allowbreak{}XVII-b:\allowbreak{}5}
\DERapparatusitem{ENV1055}{31;\allowbreak{}\allowbreak{} ms63,\allowbreak{}\allowbreak{} target\allowbreak{} of\allowbreak{} w\_\allowbreak{}!}{}{'}{DER\_\allowbreak{}XVII-a:\allowbreak{}65;\allowbreak{}\allowbreak{} DER\_\allowbreak{}XVII-a:\allowbreak{}66;\allowbreak{}\allowbreak{} DER\_\allowbreak{}XVII-b:\allowbreak{}5}
\DERapparatusitem{ENV1056}{32;\allowbreak{}\allowbreak{} printed\allowbreak{} p.\allowbreak{}32;\allowbreak{}\allowbreak{} ms.\allowbreak{}64,\allowbreak{}\allowbreak{} Corollary\allowbreak{} 3.\allowbreak{}9\allowbreak{} diagram\allowbreak{} domain}{\textbackslash{}\allowbreak{}textit\{\allowbreak{}artinian\allowbreak{} and\allowbreak{} such\allowbreak{} that\}\allowbreak{}}{\textbackslash{}\allowbreak{}textit\{\allowbreak{}noetherian\allowbreak{} and\allowbreak{} such\allowbreak{} that\}\allowbreak{}}{DER\_\allowbreak{}XVII-b:\allowbreak{}11}
\DERapparatusitem{ENV1057}{32;\allowbreak{}\allowbreak{} printed\allowbreak{} p.\allowbreak{}32;\allowbreak{}\allowbreak{} ms.\allowbreak{}64,\allowbreak{}\allowbreak{} Corollary\allowbreak{} 3.\allowbreak{}9\allowbreak{} diagram\allowbreak{} domain}{}{\textasciicircum{}\allowbreak{}\textbackslash{}\allowbreak{}circ}{DER\_\allowbreak{}XVII-b:\allowbreak{}11}
\DERapparatusitem{ENV1058}{32;\allowbreak{}\allowbreak{} printed\allowbreak{} p.\allowbreak{}32;\allowbreak{}\allowbreak{} ms.\allowbreak{}64,\allowbreak{}\allowbreak{} Corollary\allowbreak{} 3.\allowbreak{}9\allowbreak{} diagram\allowbreak{} domain}{\textbackslash{}\allowbreak{}textit\{\allowbreak{}consisting\allowbreak{} of\allowbreak{} the\allowbreak{} artinian\allowbreak{} ordered\allowbreak{} categories\allowbreak{} \$\allowbreak{}I\$\allowbreak{},\allowbreak{}\allowbreak{} nothing\allowbreak{} more.\allowbreak{}\}\allowbreak{}}{\textbackslash{}\allowbreak{}textit\{\allowbreak{}consisting\allowbreak{} of\allowbreak{} the\allowbreak{} noetherian\allowbreak{} ordered\allowbreak{} categories\allowbreak{} \$\allowbreak{}I\$\allowbreak{},\allowbreak{}\allowbreak{} nothing\allowbreak{} more.\allowbreak{}\}\allowbreak{}}{DER\_\allowbreak{}XVII-b:\allowbreak{}11}
\DERapparatusitem{ENV1059}{32;\allowbreak{}\allowbreak{} ms66,\allowbreak{}\allowbreak{} opening\allowbreak{} map\allowbreak{} of\allowbreak{} arbitrary\allowbreak{} categories}{'\textbackslash{}\allowbreak{}longrightarrow\allowbreak{} I}{\textbackslash{}\allowbreak{}longrightarrow\allowbreak{} I'}{DER\_\allowbreak{}XVII-a:\allowbreak{}68;\allowbreak{}\allowbreak{} DER\_\allowbreak{}XVII-b:\allowbreak{}7}
\DERapparatusitem{ENV1060}{32;\allowbreak{}\allowbreak{} ms66,\allowbreak{}\allowbreak{} pointwise\allowbreak{} Kan-extension\allowbreak{} formula}{}{'}{DER\_\allowbreak{}XVII-a:\allowbreak{}69;\allowbreak{}\allowbreak{} DER\_\allowbreak{}XVII-b:\allowbreak{}8}
\DERapparatusitem{ENV1061}{33;\allowbreak{}\allowbreak{} ms68,\allowbreak{}\allowbreak{} second\allowbreak{} assertion\allowbreak{} of\allowbreak{} adjunction}{*\textasciicircum{}\allowbreak{}W}{W\textasciicircum{}\allowbreak{}*}{DER\_\allowbreak{}XVII-a:\allowbreak{}73;\allowbreak{}\allowbreak{} DER\_\allowbreak{}XVII-b:\allowbreak{}14}
\DERapparatusitem{ENV1062}{33;\allowbreak{}\allowbreak{} printed\allowbreak{} p.\allowbreak{}33;\allowbreak{}\allowbreak{} ms.\allowbreak{}69,\allowbreak{}\allowbreak{} 3.\allowbreak{}11\allowbreak{} diagram\allowbreak{} domain}{defined\allowbreak{} above\allowbreak{} (\allowbreak{}artinian\allowbreak{} ordered\allowbreak{} categories\allowbreak{} \$\allowbreak{}I\$\allowbreak{}}{defined\allowbreak{} above\allowbreak{} (\allowbreak{}noetherian\allowbreak{} ordered\allowbreak{} categories\allowbreak{} \$\allowbreak{}I\$\allowbreak{}}{DER\_\allowbreak{}XVII-b:\allowbreak{}18}
\DERapparatusitem{ENV1063}{33;\allowbreak{}\allowbreak{} printed\allowbreak{} p.\allowbreak{}33;\allowbreak{}\allowbreak{} ms.\allowbreak{}69,\allowbreak{}\allowbreak{} 3.\allowbreak{}11\allowbreak{} diagram\allowbreak{} domain}{}{\textasciicircum{}\allowbreak{}\textbackslash{}\allowbreak{}circ}{DER\_\allowbreak{}XVII-b:\allowbreak{}18}
\DERapparatusitem{ENV1064}{33;\allowbreak{}\allowbreak{} printed\allowbreak{} p.\allowbreak{}33;\allowbreak{}\allowbreak{} ms.\allowbreak{}69,\allowbreak{}\allowbreak{} 3.\allowbreak{}11\allowbreak{} diagram\allowbreak{} domain}{\$\allowbreak{}\textbackslash{}\allowbreak{}mathcal\{\allowbreak{}D\}\allowbreak{}\allowbreak{} =\allowbreak{}\allowbreak{} \textbackslash{}\allowbreak{}mathrm\{\allowbreak{}Diag\}\allowbreak{}\$\allowbreak{}\allowbreak{} (\allowbreak{}artinian\allowbreak{} ordered\allowbreak{} categories\allowbreak{} \$\allowbreak{}I\$\allowbreak{}}{\$\allowbreak{}\textbackslash{}\allowbreak{}mathcal\{\allowbreak{}D\}\allowbreak{}\allowbreak{} =\allowbreak{}\allowbreak{} \textbackslash{}\allowbreak{}mathrm\{\allowbreak{}Diag\}\allowbreak{}\$\allowbreak{}\allowbreak{} (\allowbreak{}noetherian\allowbreak{} ordered\allowbreak{} categories\allowbreak{} \$\allowbreak{}I\$\allowbreak{}}{DER\_\allowbreak{}XVII-b:\allowbreak{}18}
\DERapparatusitem{ENV1065}{preamble;\allowbreak{}\allowbreak{} local\allowbreak{} macro\allowbreak{} definition}{\textbackslash{}\allowbreak{}text\{\allowbreak{}def\}\allowbreak{}}{\textbackslash{}\allowbreak{}mathrm\{\allowbreak{}def\}\allowbreak{}}{TYPESETTING:\allowbreak{}7;\allowbreak{}\allowbreak{} NOTATION\_\allowbreak{}DER:\allowbreak{}17}
\DERapparatusitem{ENV1066}{34;\allowbreak{}\allowbreak{} printed\allowbreak{} p.\allowbreak{}34;\allowbreak{}\allowbreak{} ms.\allowbreak{}72,\allowbreak{}\allowbreak{} 4.\allowbreak{}2\allowbreak{} final\allowbreak{} sentence\allowbreak{} before\allowbreak{} page\allowbreak{} break}{)\allowbreak{},\allowbreak{}\allowbreak{} is\allowbreak{} not\allowbreak{} stable\allowbreak{} under\allowbreak{} \$\allowbreak{}\textbackslash{}\allowbreak{}varinjlim\$\allowbreak{}}{)\allowbreak{},\allowbreak{}\allowbreak{} is\allowbreak{} stable\allowbreak{} under\allowbreak{} \$\allowbreak{}\textbackslash{}\allowbreak{}varinjlim\$\allowbreak{}}{DER\_\allowbreak{}XVII-b:\allowbreak{}19}
\DERapparatusitem{ENV1067}{36;\allowbreak{}\allowbreak{} printed\allowbreak{} p.\allowbreak{}36;\allowbreak{}\allowbreak{} ms.\allowbreak{}75,\allowbreak{}\allowbreak{} last\allowbreak{} line}{'}{}{DER\_\allowbreak{}XVII-b:\allowbreak{}20}
\DERapparatusitem{ENV1068}{preamble;\allowbreak{}\allowbreak{} local\allowbreak{} macro\allowbreak{} definition}{\textbackslash{}\allowbreak{}text\{\allowbreak{}def\}\allowbreak{}}{\textbackslash{}\allowbreak{}mathrm\{\allowbreak{}def\}\allowbreak{}}{TYPESETTING:\allowbreak{}7;\allowbreak{}\allowbreak{} NOTATION\_\allowbreak{}DER:\allowbreak{}17}
\DERapparatusitem{ENV1069}{39;\allowbreak{}\allowbreak{} ms.\allowbreak{}81;\allowbreak{}\allowbreak{} Proposition\allowbreak{} 4.\allowbreak{}8,\allowbreak{}\allowbreak{} parenthetical\allowbreak{} reference}{(\allowbreak{}def.\allowbreak{}\allowbreak{} 4.\allowbreak{}4,\allowbreak{}\allowbreak{} page\allowbreak{} 94)\allowbreak{}}{(\allowbreak{}definition\allowbreak{} 4.\allowbreak{}4,\allowbreak{}\allowbreak{} page\allowbreak{} 74)\allowbreak{}}{DER\_\allowbreak{}XVII-b:\allowbreak{}22}
\DERapparatusitem{ENV1070}{40;\allowbreak{}\allowbreak{} ms.\allowbreak{}82;\allowbreak{}\allowbreak{} equation\allowbreak{} (\allowbreak{}4.\allowbreak{}8.\allowbreak{}1)\allowbreak{},\allowbreak{}\allowbreak{} both\allowbreak{} class\allowbreak{} arguments}{)\allowbreak{})\allowbreak{})\allowbreak{}\allowbreak{} \textbackslash{}\allowbreak{};\allowbreak{}\textbackslash{}\allowbreak{}subset\textbackslash{}\allowbreak{};\allowbreak{}\allowbreak{} TC(\allowbreak{}I}{\_\allowbreak{}0)\allowbreak{})\allowbreak{})\allowbreak{}\textbackslash{}\allowbreak{}subset\allowbreak{} TC(\allowbreak{}I\_\allowbreak{}0}{DER\_\allowbreak{}XVII-b:\allowbreak{}23}
\DERapparatusitem{ENV1071}{40;\allowbreak{}\allowbreak{} ms.\allowbreak{}82;\allowbreak{}\allowbreak{} footnote\allowbreak{} 23,\allowbreak{}\allowbreak{} first\allowbreak{} proposed\allowbreak{} equality}{These\allowbreak{} inclusions\allowbreak{} should\allowbreak{} be\allowbreak{} equalities}{These\allowbreak{} inclusions\allowbreak{} should\allowbreak{} give\allowbreak{} \$\allowbreak{}TC(\allowbreak{}I)\allowbreak{}=\allowbreak{}(\allowbreak{}v\_\allowbreak{}!(\allowbreak{}TC(\allowbreak{}I\_\allowbreak{}0)\allowbreak{})\allowbreak{})\allowbreak{}\textbackslash{}\allowbreak{}cl\$\allowbreak{}\allowbreak{} and\allowbreak{} \$\allowbreak{}F(\allowbreak{}I)\allowbreak{}=\allowbreak{}(\allowbreak{}v\textasciicircum{}\allowbreak{}*)\allowbreak{}\textasciicircum{}\allowbreak{}\{\allowbreak{}-1\}\allowbreak{}(\allowbreak{}F(\allowbreak{}I\_\allowbreak{}0)\allowbreak{})\allowbreak{}\$\allowbreak{}}{DER\_\allowbreak{}XVII-b:\allowbreak{}24}
\DERapparatusitem{ENV1072}{41;\allowbreak{}\allowbreak{} ms.\allowbreak{}84;\allowbreak{}\allowbreak{} equation\allowbreak{} (\allowbreak{}5.\allowbreak{}5)\allowbreak{},\allowbreak{}\allowbreak{} initial\allowbreak{} equality}{T\textbackslash{}\allowbreak{}tilde\{\allowbreak{}C\}\allowbreak{}\allowbreak{} =\allowbreak{}}{}{DER\_\allowbreak{}XVII-b:\allowbreak{}26}
\DERapparatusitem{ENV1073}{42;\allowbreak{}\allowbreak{} ms.\allowbreak{}85;\allowbreak{}\allowbreak{} first\allowbreak{} equality\allowbreak{} in\allowbreak{} proof\allowbreak{} of\allowbreak{} Proposition\allowbreak{} 5.\allowbreak{}1}{=\allowbreak{}}{\textbackslash{}\allowbreak{}subset}{DER\_\allowbreak{}XVII-b:\allowbreak{}27}
\DERapparatusitem{ENV1074}{42;\allowbreak{}\allowbreak{} ms.\allowbreak{}85;\allowbreak{}\allowbreak{} lifting\allowbreak{} sentence\allowbreak{} and\allowbreak{} following\allowbreak{} triangle\allowbreak{} label}{}{\textbackslash{}\allowbreak{}tilde\{\allowbreak{}}{DER\_\allowbreak{}XVII-b:\allowbreak{}28}
\DERapparatusitem{ENV1075}{42;\allowbreak{}\allowbreak{} ms.\allowbreak{}85;\allowbreak{}\allowbreak{} lifting\allowbreak{} sentence\allowbreak{} and\allowbreak{} following\allowbreak{} triangle\allowbreak{} label}{}{\}\allowbreak{}}{DER\_\allowbreak{}XVII-b:\allowbreak{}28}
\DERapparatusitem{ENV1076}{42;\allowbreak{}\allowbreak{} ms.\allowbreak{}85;\allowbreak{}\allowbreak{} lifting\allowbreak{} sentence\allowbreak{} and\allowbreak{} following\allowbreak{} triangle\allowbreak{} label}{}{\textbackslash{}\allowbreak{}tilde\{\allowbreak{}}{DER\_\allowbreak{}XVII-b:\allowbreak{}28}
\DERapparatusitem{ENV1077}{42;\allowbreak{}\allowbreak{} ms.\allowbreak{}85;\allowbreak{}\allowbreak{} lifting\allowbreak{} sentence\allowbreak{} and\allowbreak{} following\allowbreak{} triangle\allowbreak{} label}{}{\}\allowbreak{}}{DER\_\allowbreak{}XVII-b:\allowbreak{}28}
\DERapparatusitem{ENV1078}{preamble;\allowbreak{}\allowbreak{} local\allowbreak{} macro\allowbreak{} definition}{\textbackslash{}\allowbreak{}text\{\allowbreak{}def\}\allowbreak{}}{\textbackslash{}\allowbreak{}mathrm\{\allowbreak{}def\}\allowbreak{}}{TYPESETTING:\allowbreak{}7;\allowbreak{}\allowbreak{} NOTATION\_\allowbreak{}DER:\allowbreak{}17}
\DERapparatusitem{ENV1079}{44;\allowbreak{}\allowbreak{} ms.\allowbreak{}88;\allowbreak{}\allowbreak{} reference\allowbreak{} to\allowbreak{} XV\allowbreak{} corollary\allowbreak{} 4,\allowbreak{}\allowbreak{} formula\allowbreak{} in\allowbreak{} prose}{\$\allowbreak{}W\allowbreak{} =\allowbreak{}\allowbreak{} TF\allowbreak{} \textbackslash{}\allowbreak{}circ\allowbreak{} TW\$\allowbreak{}}{\$\allowbreak{}W\allowbreak{} =\allowbreak{}\allowbreak{} TF\allowbreak{} \textbackslash{}\allowbreak{}circ\allowbreak{} TC\$\allowbreak{}}{DER\_\allowbreak{}XVII-b:\allowbreak{}31}
\DERapparatusitem{ENV1080}{45;\allowbreak{}\allowbreak{} ms.\allowbreak{}92;\allowbreak{}\allowbreak{} parenthetical\allowbreak{} notation\allowbreak{} comparison\allowbreak{} in\allowbreak{} proof}{was\allowbreak{} denoted\allowbreak{} \$\allowbreak{}\textbackslash{}\allowbreak{}tilde\{\allowbreak{}C\}\allowbreak{}'\$\allowbreak{}\allowbreak{} in\allowbreak{} proposition\allowbreak{} 4.\allowbreak{}5}{was\allowbreak{} denoted\allowbreak{} \$\allowbreak{}\textbackslash{}\allowbreak{}tilde\{\allowbreak{}C\}\allowbreak{}'\$\allowbreak{}\allowbreak{} in\allowbreak{} proposition\allowbreak{} 4.\allowbreak{}6}{DER\_\allowbreak{}XVII-b:\allowbreak{}34}
\DERapparatusitem{ENV1081}{47;\allowbreak{}\allowbreak{} ms.\allowbreak{}95;\allowbreak{}\allowbreak{} Q-saturation\allowbreak{} paragraph;\allowbreak{}\allowbreak{} supplied\allowbreak{} [\allowbreak{}sic]\allowbreak{}}{}{C}{DER\_\allowbreak{}XVII-b:\allowbreak{}36}
\DERapparatusitem{ENV1082}{48;\allowbreak{}\allowbreak{} ms.\allowbreak{}96;\allowbreak{}\allowbreak{} first\allowbreak{} display;\allowbreak{}\allowbreak{} supplied\allowbreak{} [\allowbreak{}sic]\allowbreak{}}{}{\_\allowbreak{}*}{DER\_\allowbreak{}XVII-b:\allowbreak{}37}
\DERapparatusitem{ENV1083}{48;\allowbreak{}\allowbreak{} ms.\allowbreak{}96;\allowbreak{}\allowbreak{} reference\allowbreak{} to\allowbreak{} triangle\allowbreak{} (\allowbreak{}*)\allowbreak{}}{diagram\allowbreak{} \$\allowbreak{}(\allowbreak{}*)\allowbreak{}\$\allowbreak{},\allowbreak{}\allowbreak{} page\allowbreak{} 114}{diagram\allowbreak{} \$\allowbreak{}(\allowbreak{}*)\allowbreak{}\$\allowbreak{},\allowbreak{}\allowbreak{} page\allowbreak{} 94}{DER\_\allowbreak{}XVII-b:\allowbreak{}38}
\DERapparatusitem{ENV1084}{50;\allowbreak{}\allowbreak{} ms.\allowbreak{}101,\allowbreak{}\allowbreak{} Theorem6.\allowbreak{}1(\allowbreak{}d)\allowbreak{}}{when\allowbreak{} \$\allowbreak{}I\$\allowbreak{}\allowbreak{} is\allowbreak{} an\allowbreak{} ordered\allowbreak{} set\allowbreak{} that\allowbreak{} is\}\allowbreak{}\allowbreak{} Artinian.\allowbreak{}}{when\allowbreak{} \$\allowbreak{}I\$\allowbreak{}\allowbreak{} is\allowbreak{} an\allowbreak{} ordered\allowbreak{} set\allowbreak{} that\allowbreak{} is\}\allowbreak{}\allowbreak{} Noetherian.\allowbreak{}}{DER\_\allowbreak{}XVII-b:\allowbreak{}41}
\DERapparatusitem{ENV1085}{51;\allowbreak{}\allowbreak{} ms.\allowbreak{}103,\allowbreak{}\allowbreak{} display\allowbreak{} before(\allowbreak{}6.\allowbreak{}1.\allowbreak{}16)\allowbreak{},\allowbreak{}\allowbreak{} \%\allowbreak{} [\allowbreak{}sic]\allowbreak{}}{I'\_\allowbreak{}0)\allowbreak{})\allowbreak{}\textbackslash{}\allowbreak{}tl}{C(\allowbreak{}I'\_\allowbreak{}0)\allowbreak{})\allowbreak{})\allowbreak{}\textasciicircum{}\allowbreak{}\{\allowbreak{}\textbackslash{}\allowbreak{}sim\}\allowbreak{}}{DER\_\allowbreak{}XVII-b:\allowbreak{}42}
\DERapparatusitem{ENV1086}{52;\allowbreak{}\allowbreak{} ms.\allowbreak{}105,\allowbreak{}\allowbreak{} definitions\allowbreak{} in6.\allowbreak{}2,\allowbreak{}\allowbreak{} \%\allowbreak{} [\allowbreak{}sic]\allowbreak{}}{U}{V}{DER\_\allowbreak{}XVII-b:\allowbreak{}43}
\DERapparatusitem{ENV1087}{52;\allowbreak{}\allowbreak{} running\allowbreak{} text\allowbreak{} or\allowbreak{} formula}{\textbackslash{}\allowbreak{}text\{\allowbreak{}def\}\allowbreak{}}{\textbackslash{}\allowbreak{}mathrm\{\allowbreak{}def\}\allowbreak{}}{TYPESETTING:\allowbreak{}7;\allowbreak{}\allowbreak{} NOTATION\_\allowbreak{}DER:\allowbreak{}17}
\DERapparatusitem{ENV1088}{53;\allowbreak{}\allowbreak{} ms.\allowbreak{}106,\allowbreak{}\allowbreak{} diagram(\allowbreak{}6.\allowbreak{}2.\allowbreak{}3)\allowbreak{},\allowbreak{}\allowbreak{} edition\allowbreak{} label}{0}{\textbackslash{}\allowbreak{}partial}{DER\_\allowbreak{}XVII-b:\allowbreak{}46}
\DERapparatusitem{ENV1089}{54;\allowbreak{}\allowbreak{} ms.\allowbreak{}108,\allowbreak{}\allowbreak{} sentence\allowbreak{} before6.\allowbreak{}3}{the\allowbreak{} equality\allowbreak{} \$\allowbreak{}C(\allowbreak{}I)\allowbreak{}\allowbreak{} =\allowbreak{}\allowbreak{} C\_\allowbreak{}I\$\allowbreak{}\allowbreak{} for\allowbreak{} \$\allowbreak{}I\$\allowbreak{}\allowbreak{} artinian}{the\allowbreak{} equality\allowbreak{} \$\allowbreak{}C(\allowbreak{}I)\allowbreak{}\allowbreak{} =\allowbreak{}\allowbreak{} C\_\allowbreak{}I\$\allowbreak{}\allowbreak{} for\allowbreak{} \$\allowbreak{}I\$\allowbreak{}\allowbreak{} noetherian}{DER\_\allowbreak{}XVII-b:\allowbreak{}48}
\DERapparatusitem{ENV1090}{54;\allowbreak{}\allowbreak{} ms.\allowbreak{}109,\allowbreak{}\allowbreak{} opening\allowbreak{} proof\allowbreak{} of6.\allowbreak{}3,\allowbreak{}\allowbreak{} \%\allowbreak{} [\allowbreak{}sic]\allowbreak{}}{u}{i}{DER\_\allowbreak{}XVII-b:\allowbreak{}49}
\DERapparatusitem{ENV1091}{55;\allowbreak{}\allowbreak{} running\allowbreak{} text\allowbreak{} or\allowbreak{} formula}{\textbackslash{}\allowbreak{}text\{\allowbreak{}def\}\allowbreak{}}{\textbackslash{}\allowbreak{}mathrm\{\allowbreak{}def\}\allowbreak{}}{TYPESETTING:\allowbreak{}7;\allowbreak{}\allowbreak{} NOTATION\_\allowbreak{}DER:\allowbreak{}17}
\DERapparatusitem{ENV1092}{56;\allowbreak{}\allowbreak{} ms.\allowbreak{}112,\allowbreak{}\allowbreak{} (\allowbreak{}6.\allowbreak{}3.\allowbreak{}8)\allowbreak{},\allowbreak{}\allowbreak{} \%\allowbreak{} [\allowbreak{}sic]\allowbreak{}}{X}{I}{DER\_\allowbreak{}XVII-b:\allowbreak{}53}
\DERapparatusitem{ENV1093}{56;\allowbreak{}\allowbreak{} ms.\allowbreak{}113,\allowbreak{}\allowbreak{} sentence\allowbreak{} before(\allowbreak{}6.\allowbreak{}3.\allowbreak{}9)\allowbreak{},\allowbreak{}\allowbreak{} \%\allowbreak{} [\allowbreak{}sic]\allowbreak{}}{C(\allowbreak{}}{\textbackslash{}\allowbreak{}Ph}{DER\_\allowbreak{}XVII-b:\allowbreak{}54}
\DERapparatusitem{ENV1094}{56;\allowbreak{}\allowbreak{} ms.\allowbreak{}113,\allowbreak{}\allowbreak{} sentence\allowbreak{} before(\allowbreak{}6.\allowbreak{}3.\allowbreak{}9)\allowbreak{},\allowbreak{}\allowbreak{} \%\allowbreak{} [\allowbreak{}sic]\allowbreak{}}{)\allowbreak{}}{\_\allowbreak{}I}{DER\_\allowbreak{}XVII-b:\allowbreak{}54}
\DERapparatusitem{ENV1095}{56;\allowbreak{}\allowbreak{} ms.\allowbreak{}113,\allowbreak{}\allowbreak{} sentence\allowbreak{} before(\allowbreak{}6.\allowbreak{}3.\allowbreak{}9)\allowbreak{},\allowbreak{}\allowbreak{} \%\allowbreak{} [\allowbreak{}sic]\allowbreak{}}{C(\allowbreak{}}{\textbackslash{}\allowbreak{}Phi\_\allowbreak{}\{\allowbreak{}}{DER\_\allowbreak{}XVII-b:\allowbreak{}54}
\DERapparatusitem{ENV1096}{56;\allowbreak{}\allowbreak{} ms.\allowbreak{}113,\allowbreak{}\allowbreak{} sentence\allowbreak{} before(\allowbreak{}6.\allowbreak{}3.\allowbreak{}9)\allowbreak{},\allowbreak{}\allowbreak{} \%\allowbreak{} [\allowbreak{}sic]\allowbreak{}}{)\allowbreak{}}{\}\allowbreak{}}{DER\_\allowbreak{}XVII-b:\allowbreak{}54}
\DERapparatusitem{ENV1097}{57;\allowbreak{}\allowbreak{} running\allowbreak{} text\allowbreak{} or\allowbreak{} formula}{\textbackslash{}\allowbreak{}text\{\allowbreak{}def\}\allowbreak{}}{\textbackslash{}\allowbreak{}mathrm\{\allowbreak{}def\}\allowbreak{}}{TYPESETTING:\allowbreak{}7;\allowbreak{}\allowbreak{} NOTATION\_\allowbreak{}DER:\allowbreak{}17}
\DERapparatusitem{ENV1098}{59;\allowbreak{}\allowbreak{} ms.\allowbreak{}116,\allowbreak{}\allowbreak{} last\allowbreak{} lifting\allowbreak{} sentence}{\textbackslash{}\allowbreak{}in\allowbreak{} \textbackslash{}\allowbreak{}Phi(\allowbreak{}I)\allowbreak{}}{\textbackslash{}\allowbreak{}in\textbackslash{}\allowbreak{}Phi\_\allowbreak{}I}{DER\_\allowbreak{}XVII-b:\allowbreak{}57}
\DERapparatusitem{ENV1099}{59;\allowbreak{}\allowbreak{} ms.\allowbreak{}116,\allowbreak{}\allowbreak{} last\allowbreak{} lifting\allowbreak{} sentence,\allowbreak{}\allowbreak{} \%\allowbreak{} [\allowbreak{}sic]\allowbreak{}}{o}{0}{DER\_\allowbreak{}XVII-b:\allowbreak{}58}
\DERapparatusitem{ENV1100}{60;\allowbreak{}\allowbreak{} ms.\allowbreak{}118,\allowbreak{}\allowbreak{} final\allowbreak{} sentence}{}{(\allowbreak{}X)\allowbreak{}}{DER\_\allowbreak{}XVII-b:\allowbreak{}59}
\subsection*{Chapter XVIII}
\DERapparatusitem{ENV1101}{1;\allowbreak{}\allowbreak{} ms50–51,\allowbreak{}\allowbreak{} first\allowbreak{} parenthesis}{(\allowbreak{}\textbackslash{}\allowbreak{}footnote\{\allowbreak{}This\allowbreak{} is\allowbreak{} the\allowbreak{} case\allowbreak{} in\allowbreak{} particular\allowbreak{} if\allowbreak{} \$\allowbreak{}\textbackslash{}\allowbreak{}mathcal\{\allowbreak{}M\}\allowbreak{}\$\allowbreak{}\allowbreak{} is\allowbreak{} a\allowbreak{} topos.\allowbreak{}\}\allowbreak{})\allowbreak{}.\allowbreak{}}{(\allowbreak{}\textbackslash{}\allowbreak{}footnote\{\allowbreak{}This\allowbreak{} is\allowbreak{} the\allowbreak{} case\allowbreak{} in\allowbreak{} particular\allowbreak{} if\allowbreak{} \$\allowbreak{}\textbackslash{}\allowbreak{}mathcal\{\allowbreak{}M\}\allowbreak{}\$\allowbreak{}\allowbreak{} is\allowbreak{} a\allowbreak{} topos.\allowbreak{}\}\allowbreak{})\allowbreak{})\allowbreak{}.\allowbreak{}}{DER\_\allowbreak{}XVIII-a:\allowbreak{}2}
\DERapparatusitem{ENV1102}{1;\allowbreak{}\allowbreak{} ms52,\allowbreak{}\allowbreak{} last\allowbreak{} paragraph}{i.\allowbreak{}e.\allowbreak{}\allowbreak{} that\allowbreak{} it\allowbreak{} admits\allowbreak{} a\allowbreak{} right\allowbreak{} adjoint}{i.\allowbreak{}e.\allowbreak{}\allowbreak{} that\allowbreak{} it\allowbreak{} admits\allowbreak{} a\allowbreak{} left\allowbreak{} adjoint}{DER\_\allowbreak{}XVIII-a:\allowbreak{}3}
\DERapparatusitem{ENV1103}{2;\allowbreak{}\allowbreak{} ms54,\allowbreak{}\allowbreak{} equation4.\allowbreak{}1.\allowbreak{}8}{Q}{P}{DER\_\allowbreak{}XVIII-a:\allowbreak{}4}
\DERapparatusitem{ENV1104}{preamble;\allowbreak{}\allowbreak{} local\allowbreak{} macro\allowbreak{} definition}{\textbackslash{}\allowbreak{}text\{\allowbreak{}def\}\allowbreak{}}{\textbackslash{}\allowbreak{}mathrm\{\allowbreak{}def\}\allowbreak{}}{TYPESETTING:\allowbreak{}7;\allowbreak{}\allowbreak{} NOTATION\_\allowbreak{}DER:\allowbreak{}17}
\DERapparatusitem{ENV1105}{4;\allowbreak{}\allowbreak{} ms58,\allowbreak{}\allowbreak{} formula\allowbreak{} below4.\allowbreak{}2.\allowbreak{}3}{\textbackslash{}\allowbreak{}Hom(\allowbreak{}y,\allowbreak{}\textbackslash{}\allowbreak{}psi(\allowbreak{}x}{\textbackslash{}\allowbreak{}Hom\_\allowbreak{}\{\allowbreak{}\textbackslash{}\allowbreak{}mathcal\allowbreak{} M\}\allowbreak{}(\allowbreak{}x,\allowbreak{}\textbackslash{}\allowbreak{}psi(\allowbreak{}y}{DER\_\allowbreak{}XVIII-a:\allowbreak{}6}
\DERapparatusitem{ENV1106}{6;\allowbreak{}\allowbreak{} ms61,\allowbreak{}\allowbreak{} first\allowbreak{} display}{\textbackslash{}\allowbreak{}uHom(\allowbreak{}S\textasciicircum{}\allowbreak{}\{\allowbreak{}\textbackslash{}\allowbreak{}wedge\textbackslash{}\allowbreak{}circ\}\allowbreak{},\allowbreak{}\textbackslash{}\allowbreak{}cN)\allowbreak{}}{\textbackslash{}\allowbreak{}uHom\textasciicircum{}\allowbreak{}\{\allowbreak{}!\}\allowbreak{}(\allowbreak{}S\textasciicircum{}\allowbreak{}\{\allowbreak{}\textbackslash{}\allowbreak{}wedge\textbackslash{}\allowbreak{}circ\}\allowbreak{},\allowbreak{}\textbackslash{}\allowbreak{}cN)\allowbreak{}}{DER\_\allowbreak{}XVIII-a:\allowbreak{}7}
\DERapparatusitem{ENV1107}{6;\allowbreak{}\allowbreak{} ms61,\allowbreak{}\allowbreak{} equation4.\allowbreak{}2.\allowbreak{}1.\allowbreak{}1}{\textbackslash{}\allowbreak{}uHom}{\textbackslash{}\allowbreak{}uHom\textasciicircum{}\allowbreak{}\{\allowbreak{}!\}\allowbreak{}}{DER\_\allowbreak{}XVIII-a:\allowbreak{}8}
\DERapparatusitem{ENV1108}{6;\allowbreak{}\allowbreak{} ms62,\allowbreak{}\allowbreak{} equation4.\allowbreak{}2.\allowbreak{}2.\allowbreak{}2}{\textbackslash{}\allowbreak{}wedge}{\textbackslash{}\allowbreak{}circ}{DER\_\allowbreak{}XVIII-a:\allowbreak{}9}
\DERapparatusitem{ENV1109}{7;\allowbreak{}\allowbreak{} ms64,\allowbreak{}\allowbreak{} final\allowbreak{} paragraph}{to\allowbreak{} the\allowbreak{} initial\allowbreak{} object\allowbreak{} of\allowbreak{} \$\allowbreak{}\textbackslash{}\allowbreak{}cM\$\allowbreak{})\allowbreak{}\allowbreak{} into\allowbreak{} a\allowbreak{} final\allowbreak{} object}{to\allowbreak{} the\allowbreak{} initial\allowbreak{} object\allowbreak{} of\allowbreak{} \$\allowbreak{}\textbackslash{}\allowbreak{}cM\$\allowbreak{}\allowbreak{} into\allowbreak{} a\allowbreak{} final\allowbreak{} object}{DER\_\allowbreak{}XVIII-a:\allowbreak{}10}
\DERapparatusitem{ENV1110}{9;\allowbreak{}\allowbreak{} ms66,\allowbreak{}\allowbreak{} Lemma4.\allowbreak{}3.\allowbreak{}2\allowbreak{} hypothesis}{\$\allowbreak{}\textbackslash{}\allowbreak{}mathcal\allowbreak{} E\$\allowbreak{}\allowbreak{} stable\allowbreak{} under\allowbreak{} small\allowbreak{} \$\allowbreak{}\textbackslash{}\allowbreak{}varprojlim\$\allowbreak{}}{\$\allowbreak{}\textbackslash{}\allowbreak{}mathcal\allowbreak{} E\$\allowbreak{}\allowbreak{} stable\allowbreak{} under\allowbreak{} small\allowbreak{} \$\allowbreak{}\textbackslash{}\allowbreak{}varinjlim\$\allowbreak{}}{DER\_\allowbreak{}XVIII-a:\allowbreak{}13}
\DERapparatusitem{ENV1111}{9;\allowbreak{}\allowbreak{} ms66,\allowbreak{}\allowbreak{} Lemma4.\allowbreak{}3.\allowbreak{}2\allowbreak{} evaluation\allowbreak{} domain}{\textbackslash{}\allowbreak{}mathcal\allowbreak{} E\textasciicircum{}\allowbreak{}\textbackslash{}\allowbreak{}circ\allowbreak{} \&\allowbreak{}\allowbreak{} \textbackslash{}\allowbreak{}longrightarrow\allowbreak{} \&\allowbreak{}\allowbreak{} \textbackslash{}\allowbreak{}mathrm\{\allowbreak{}Ens\}\allowbreak{}}{\textbackslash{}\allowbreak{}mathcal\allowbreak{} E\allowbreak{} \&\allowbreak{}\allowbreak{} \textbackslash{}\allowbreak{}longrightarrow\allowbreak{} \&\allowbreak{}\allowbreak{} \textbackslash{}\allowbreak{}mathrm\{\allowbreak{}Ens\}\allowbreak{}}{DER\_\allowbreak{}XVIII-a:\allowbreak{}14}
\DERapparatusitem{ENV1112}{10;\allowbreak{}\allowbreak{} ms67,\allowbreak{}\allowbreak{} proof\allowbreak{} of\allowbreak{} Lemma4.\allowbreak{}3.\allowbreak{}2,\allowbreak{}\allowbreak{} second\allowbreak{} Hom\allowbreak{} expression}{(\allowbreak{}\textbackslash{}\allowbreak{}rho(\allowbreak{}F\_\allowbreak{}i)\allowbreak{},\allowbreak{}\textbackslash{}\allowbreak{}xi)\allowbreak{}}{\textbackslash{}\allowbreak{}rho(\allowbreak{}F\_\allowbreak{}i)\allowbreak{},\allowbreak{}\textbackslash{}\allowbreak{}xi}{DER\_\allowbreak{}XVIII-a:\allowbreak{}15}
\DERapparatusitem{ENV1113}{10;\allowbreak{}\allowbreak{} running\allowbreak{} text\allowbreak{} or\allowbreak{} formula}{\textbackslash{}\allowbreak{}text\{\allowbreak{}def\}\allowbreak{}}{\textbackslash{}\allowbreak{}mathrm\{\allowbreak{}def\}\allowbreak{}}{TYPESETTING:\allowbreak{}7;\allowbreak{}\allowbreak{} NOTATION\_\allowbreak{}DER:\allowbreak{}17}
\DERapparatusitem{ENV1114}{10;\allowbreak{}\allowbreak{} ms67,\allowbreak{}\allowbreak{} Lemma4.\allowbreak{}3.\allowbreak{}3(\allowbreak{}b)\allowbreak{},\allowbreak{}\allowbreak{} last\allowbreak{} target}{\$\allowbreak{}\textbackslash{}\allowbreak{}rho(\allowbreak{}U)\allowbreak{}\$\allowbreak{}\allowbreak{} is\allowbreak{} a\allowbreak{} small\allowbreak{} strictly\allowbreak{} generating\allowbreak{} subset\allowbreak{} in\allowbreak{} \$\allowbreak{}\textbackslash{}\allowbreak{}mathcal\allowbreak{} F\$\allowbreak{}}{\$\allowbreak{}\textbackslash{}\allowbreak{}rho(\allowbreak{}U)\allowbreak{}\$\allowbreak{}\allowbreak{} is\allowbreak{} a\allowbreak{} small\allowbreak{} strictly\allowbreak{} generating\allowbreak{} subset\allowbreak{} in\allowbreak{} \$\allowbreak{}\textbackslash{}\allowbreak{}mathcal\allowbreak{} E\$\allowbreak{}}{DER\_\allowbreak{}XVIII-a:\allowbreak{}16}
\DERapparatusitem{ENV1115}{10;\allowbreak{}\allowbreak{} ms68,\allowbreak{}\allowbreak{} footnote8,\allowbreak{}\allowbreak{} dense\allowbreak{} presentation\allowbreak{} of\allowbreak{} x}{\textbackslash{}\allowbreak{}varprojlim}{\textbackslash{}\allowbreak{}varinjlim}{DER\_\allowbreak{}XVIII-a:\allowbreak{}18}
\DERapparatusitem{ENV1116}{10;\allowbreak{}\allowbreak{} ms68,\allowbreak{}\allowbreak{} reference\allowbreak{} following\allowbreak{} Lemma4.\allowbreak{}3.\allowbreak{}3}{proving\allowbreak{} formula\allowbreak{} (\allowbreak{}4.\allowbreak{}1.\allowbreak{}3.\allowbreak{}4)\allowbreak{}}{proving\allowbreak{} formula\allowbreak{} (\allowbreak{}4.\allowbreak{}3.\allowbreak{}1.\allowbreak{}4)\allowbreak{}}{DER\_\allowbreak{}XVIII-a:\allowbreak{}19}
\DERapparatusitem{ENV1117}{13;\allowbreak{}\allowbreak{} ms74,\allowbreak{}\allowbreak{} condition(\allowbreak{}a)\allowbreak{},\allowbreak{}\allowbreak{} limit\allowbreak{} sign}{\textbackslash{}\allowbreak{}varprojlim}{\textbackslash{}\allowbreak{}varinjlim}{DER\_\allowbreak{}XVIII-a:\allowbreak{}21}
\DERapparatusitem{ENV1118}{13;\allowbreak{}\allowbreak{} ms74,\allowbreak{}\allowbreak{} condition(\allowbreak{}b)\allowbreak{},\allowbreak{}\allowbreak{} limit\allowbreak{} sign}{\textbackslash{}\allowbreak{}varprojlim}{\textbackslash{}\allowbreak{}varinjlim}{DER\_\allowbreak{}XVIII-a:\allowbreak{}22}
\DERapparatusitem{ENV1119}{13;\allowbreak{}\allowbreak{} running\allowbreak{} text\allowbreak{} or\allowbreak{} formula}{\textbackslash{}\allowbreak{}text\{\allowbreak{}def\}\allowbreak{}}{\textbackslash{}\allowbreak{}mathrm\{\allowbreak{}def\}\allowbreak{}}{TYPESETTING:\allowbreak{}7;\allowbreak{}\allowbreak{} NOTATION\_\allowbreak{}DER:\allowbreak{}17}
\DERapparatusitem{ENV1120}{14;\allowbreak{}\allowbreak{} running\allowbreak{} text\allowbreak{} or\allowbreak{} formula}{\textbackslash{}\allowbreak{}text\{\allowbreak{}def\}\allowbreak{}}{\textbackslash{}\allowbreak{}mathrm\{\allowbreak{}def\}\allowbreak{}}{TYPESETTING:\allowbreak{}7;\allowbreak{}\allowbreak{} NOTATION\_\allowbreak{}DER:\allowbreak{}17}
\DERapparatusitem{ENV1121}{16;\allowbreak{}\allowbreak{} ms80,\allowbreak{}\allowbreak{} proof4.\allowbreak{}4.\allowbreak{}2,\allowbreak{}\allowbreak{} first\allowbreak{} comparison}{\textbackslash{}\allowbreak{}varinjlim\allowbreak{} \textbackslash{}\allowbreak{}varphi\_\allowbreak{}\textbackslash{}\allowbreak{}alpha\allowbreak{} \textbackslash{}\allowbreak{}longrightarrow\allowbreak{} \textbackslash{}\allowbreak{}varinjlim\allowbreak{} \textbackslash{}\allowbreak{}psi\textbackslash{}\allowbreak{}varphi\_\allowbreak{}\textbackslash{}\allowbreak{}alpha}{\textbackslash{}\allowbreak{}varinjlim\_\allowbreak{}I\allowbreak{} \textbackslash{}\allowbreak{}psi\textbackslash{}\allowbreak{}varphi\_\allowbreak{}\textbackslash{}\allowbreak{}alpha\allowbreak{} \textbackslash{}\allowbreak{}longrightarrow\allowbreak{} \textbackslash{}\allowbreak{}varinjlim\_\allowbreak{}\{\allowbreak{}C/\allowbreak{}F\}\allowbreak{}\textbackslash{}\allowbreak{}psi}{DER\_\allowbreak{}XVIII-a:\allowbreak{}25}
\DERapparatusitem{ENV1122}{18;\allowbreak{}\allowbreak{} ms83,\allowbreak{}\allowbreak{} equation4.\allowbreak{}4.\allowbreak{}6.\allowbreak{}4,\allowbreak{}\allowbreak{} selected\allowbreak{} object}{e}{\textbackslash{}\allowbreak{}varepsilon}{DER\_\allowbreak{}XVIII-a:\allowbreak{}29}
\DERapparatusitem{ENV1123}{19;\allowbreak{}\allowbreak{} ms84,\allowbreak{}\allowbreak{} definition\allowbreak{} of\allowbreak{} I}{/\allowbreak{}}{,\allowbreak{}}{DER\_\allowbreak{}XVIII-a:\allowbreak{}32}
\DERapparatusitem{ENV1124}{19;\allowbreak{}\allowbreak{} running\allowbreak{} text\allowbreak{} or\allowbreak{} formula}{\textbackslash{}\allowbreak{}text\{\allowbreak{}def\}\allowbreak{}}{\textbackslash{}\allowbreak{}mathrm\{\allowbreak{}def\}\allowbreak{}}{TYPESETTING:\allowbreak{}7;\allowbreak{}\allowbreak{} NOTATION\_\allowbreak{}DER:\allowbreak{}17}
\DERapparatusitem{ENV1125}{19;\allowbreak{}\allowbreak{} ms85,\allowbreak{}\allowbreak{} kind\allowbreak{} of\allowbreak{} index\allowbreak{} I}{along\allowbreak{} the\allowbreak{} filtering\allowbreak{} ordered\allowbreak{} set\allowbreak{} \$\allowbreak{}I\$\allowbreak{}}{along\allowbreak{} the\allowbreak{} filtering\allowbreak{} category\allowbreak{} \$\allowbreak{}I\$\allowbreak{}}{DER\_\allowbreak{}XVIII-a:\allowbreak{}34}
\DERapparatusitem{ENV1126}{19;\allowbreak{}\allowbreak{} ms85,\allowbreak{}\allowbreak{} evaluation\allowbreak{} of\allowbreak{} phi}{\$\allowbreak{}\textbackslash{}\allowbreak{}varphi(\allowbreak{}j)\allowbreak{}\allowbreak{} =\allowbreak{}\allowbreak{} \textbackslash{}\allowbreak{}varepsilon\_\allowbreak{}j\textasciicircum{}\allowbreak{}*(\allowbreak{}\textbackslash{}\allowbreak{}varphi)\allowbreak{}\$\allowbreak{}}{\$\allowbreak{}\textbackslash{}\allowbreak{}varphi(\allowbreak{}j)\allowbreak{}\allowbreak{} =\allowbreak{}\allowbreak{} \textbackslash{}\allowbreak{}varepsilon\_\allowbreak{}j\textasciicircum{}\allowbreak{}\{\allowbreak{}*\textbackslash{}\allowbreak{}prime\}\allowbreak{}(\allowbreak{}\textbackslash{}\allowbreak{}varphi)\allowbreak{}\$\allowbreak{}}{DER\_\allowbreak{}XVIII-a:\allowbreak{}35}
\DERapparatusitem{ENV1127}{19;\allowbreak{}\allowbreak{} ms85,\allowbreak{}\allowbreak{} evaluation\allowbreak{} after\allowbreak{} alpha\textasciicircum{}\allowbreak{}J}{\$\allowbreak{}\textbackslash{}\allowbreak{}varepsilon\_\allowbreak{}j\textasciicircum{}\allowbreak{}*(\allowbreak{}\textbackslash{}\allowbreak{}alpha\textasciicircum{}\allowbreak{}J\textbackslash{}\allowbreak{}psi)\allowbreak{}\allowbreak{} =\allowbreak{}\allowbreak{} (\allowbreak{}\textbackslash{}\allowbreak{}alpha\textbackslash{}\allowbreak{}psi\_\allowbreak{}j)\allowbreak{}\textbackslash{}\allowbreak{}varepsilon\_\allowbreak{}j\textasciicircum{}\allowbreak{}\textbackslash{}\allowbreak{}varphi\$\allowbreak{}}{\$\allowbreak{}\textbackslash{}\allowbreak{}varepsilon\_\allowbreak{}j\textasciicircum{}\allowbreak{}\{\allowbreak{}*\textbackslash{}\allowbreak{}prime\}\allowbreak{}(\allowbreak{}\textbackslash{}\allowbreak{}alpha\textasciicircum{}\allowbreak{}J\textbackslash{}\allowbreak{}psi)\allowbreak{}\allowbreak{} =\allowbreak{}\allowbreak{} (\allowbreak{}\textbackslash{}\allowbreak{}alpha\textbackslash{}\allowbreak{}psi\_\allowbreak{}j)\allowbreak{}\textbackslash{}\allowbreak{}varepsilon\_\allowbreak{}j\textasciicircum{}\allowbreak{}\textbackslash{}\allowbreak{}varphi\$\allowbreak{}}{DER\_\allowbreak{}XVIII-a:\allowbreak{}36}
\DERapparatusitem{ENV1128}{19;\allowbreak{}\allowbreak{} ms86,\allowbreak{}\allowbreak{} curried\allowbreak{} diagram}{\textbackslash{}\allowbreak{}Psi}{\textbackslash{}\allowbreak{}psi}{DER\_\allowbreak{}XVIII-a:\allowbreak{}37}
\DERapparatusitem{ENV1129}{21;\allowbreak{}\allowbreak{} ms90,\allowbreak{}\allowbreak{} proof\allowbreak{} heading}{\textbackslash{}\allowbreak{}textsc\{\allowbreak{}Proof\allowbreak{} of\}\allowbreak{}\allowbreak{} 4.\allowbreak{}4.\allowbreak{}8.\allowbreak{}}{\textbackslash{}\allowbreak{}textsc\{\allowbreak{}Proof\allowbreak{} of\}\allowbreak{}\allowbreak{} 4.\allowbreak{}4.\allowbreak{}9.\allowbreak{}}{DER\_\allowbreak{}XVIII-a:\allowbreak{}40}
\DERapparatusitem{ENV1130}{21;\allowbreak{}\allowbreak{} ms91;\allowbreak{}\allowbreak{} last\allowbreak{} paragraph,\allowbreak{}\allowbreak{} continuation\allowbreak{} of4.\allowbreak{}4.\allowbreak{}9}{(\allowbreak{}4.\allowbreak{}4.\allowbreak{}4.\allowbreak{}2)\allowbreak{}\allowbreak{} is\allowbreak{} an\allowbreak{} isomorphism}{(\allowbreak{}4.\allowbreak{}4.\allowbreak{}9.\allowbreak{}2)\allowbreak{}\allowbreak{} is\allowbreak{} an\allowbreak{} isomorphism}{DER\_\allowbreak{}XVIII-a:\allowbreak{}41}
\DERapparatusitem{ENV1131}{26;\allowbreak{}\allowbreak{} ms100;\allowbreak{}\allowbreak{} conclusion\allowbreak{} after\allowbreak{} the\allowbreak{} PF\_\allowbreak{}pi2\allowbreak{} diagram}{completes\allowbreak{} the\allowbreak{} proof\allowbreak{} of\allowbreak{} 4.\allowbreak{}4.\allowbreak{}14}{completes\allowbreak{} the\allowbreak{} proof\allowbreak{} of\allowbreak{} 4.\allowbreak{}4.\allowbreak{}15}{DER\_\allowbreak{}XVIII-a:\allowbreak{}43}
\DERapparatusitem{ENV1132}{27;\allowbreak{}\allowbreak{} ms102–103;\allowbreak{}\allowbreak{} unrestricted\allowbreak{} restriction\allowbreak{} functor}{\textbackslash{}\allowbreak{}alpha\textasciicircum{}\allowbreak{}*\_\allowbreak{}\{\allowbreak{}\textbackslash{}\allowbreak{}mathcal\{\allowbreak{}E\}\allowbreak{}\}\allowbreak{}\textbackslash{}\allowbreak{};\allowbreak{}:\allowbreak{}\textbackslash{}\allowbreak{};\allowbreak{}\textbackslash{}\allowbreak{}uHom(\allowbreak{}\textbackslash{}\allowbreak{}mathcal\{\allowbreak{}M\}\allowbreak{},\allowbreak{}\textbackslash{}\allowbreak{}mathcal\{\allowbreak{}E\}\allowbreak{})\allowbreak{}}{\textbackslash{}\allowbreak{}alpha\textasciicircum{}\allowbreak{}*\_\allowbreak{}\{\allowbreak{}\textbackslash{}\allowbreak{}mathcal\{\allowbreak{}E\}\allowbreak{}\}\allowbreak{}\textbackslash{}\allowbreak{};\allowbreak{}:\allowbreak{}\textbackslash{}\allowbreak{};\allowbreak{}\textbackslash{}\allowbreak{}uHom\_\allowbreak{}\{\allowbreak{}!\}\allowbreak{}(\allowbreak{}\textbackslash{}\allowbreak{}mathcal\{\allowbreak{}M\}\allowbreak{},\allowbreak{}\textbackslash{}\allowbreak{}mathcal\{\allowbreak{}E\}\allowbreak{})\allowbreak{}}{DER\_\allowbreak{}XVIII-a:\allowbreak{}46}
\DERapparatusitem{ENV1133}{28;\allowbreak{}\allowbreak{} running\allowbreak{} text\allowbreak{} or\allowbreak{} formula}{\textbackslash{}\allowbreak{}text\{\allowbreak{}def\}\allowbreak{}}{\textbackslash{}\allowbreak{}mathrm\{\allowbreak{}def\}\allowbreak{}}{TYPESETTING:\allowbreak{}7;\allowbreak{}\allowbreak{} NOTATION\_\allowbreak{}DER:\allowbreak{}17}
\DERapparatusitem{ENV1134}{28;\allowbreak{}\allowbreak{} ms106;\allowbreak{}\allowbreak{} target\allowbreak{} of\allowbreak{} the\allowbreak{} arrow\allowbreak{} R\allowbreak{} in\allowbreak{} the\allowbreak{} final\allowbreak{} diagram}{\textbackslash{}\allowbreak{}uHom\_\allowbreak{}\{\allowbreak{}!\textbackslash{}\allowbreak{},\allowbreak{}\textbackslash{}\allowbreak{}pi\}\allowbreak{}(\allowbreak{}\textbackslash{}\allowbreak{}Ind\_\allowbreak{}\textbackslash{}\allowbreak{}pi(\allowbreak{}C)\allowbreak{},\allowbreak{}\textbackslash{}\allowbreak{}mathcal\{\allowbreak{}E\}\allowbreak{})\allowbreak{}}{\textbackslash{}\allowbreak{}uHom(\allowbreak{}\textbackslash{}\allowbreak{}Ind\_\allowbreak{}\textbackslash{}\allowbreak{}pi(\allowbreak{}C)\allowbreak{},\allowbreak{}\textbackslash{}\allowbreak{}mathcal\{\allowbreak{}E\}\allowbreak{})\allowbreak{}}{DER\_\allowbreak{}XVIII-a:\allowbreak{}49}
\DERapparatusitem{ENV1135}{29;\allowbreak{}\allowbreak{} running\allowbreak{} text\allowbreak{} or\allowbreak{} formula}{\textbackslash{}\allowbreak{}text\{\allowbreak{}def\}\allowbreak{}}{\textbackslash{}\allowbreak{}mathrm\{\allowbreak{}def\}\allowbreak{}}{TYPESETTING:\allowbreak{}7;\allowbreak{}\allowbreak{} NOTATION\_\allowbreak{}DER:\allowbreak{}17}
\DERapparatusitem{ENV1136}{30;\allowbreak{}\allowbreak{} ms108;\allowbreak{}\allowbreak{} formula4.\allowbreak{}5.\allowbreak{}2.\allowbreak{}2}{F\allowbreak{} \textbackslash{}\allowbreak{}in\allowbreak{} \textbackslash{}\allowbreak{}mathcal\{\allowbreak{}M\}\allowbreak{}\allowbreak{} \textbackslash{}\allowbreak{}text\{\allowbreak{}\allowbreak{} which\allowbreak{} are\allowbreak{} \}\allowbreak{}\allowbreak{} \textbackslash{}\allowbreak{}pi\textbackslash{}\allowbreak{}text\{\allowbreak{}-accessible.\allowbreak{}\}\allowbreak{}}{F\allowbreak{} \textbackslash{}\allowbreak{}in\allowbreak{} C\textasciicircum{}\allowbreak{}\{\allowbreak{}\textbackslash{}\allowbreak{}wedge\}\allowbreak{}\allowbreak{} \textbackslash{}\allowbreak{}text\{\allowbreak{}\allowbreak{} which\allowbreak{} are\allowbreak{} \}\allowbreak{}\allowbreak{} \textbackslash{}\allowbreak{}pi\textbackslash{}\allowbreak{}text\{\allowbreak{}-accessible\allowbreak{} in\allowbreak{} \}\allowbreak{}C\textasciicircum{}\allowbreak{}\{\allowbreak{}\textbackslash{}\allowbreak{}wedge\}\allowbreak{}.\allowbreak{}}{DER\_\allowbreak{}XVIII-a:\allowbreak{}51}
\DERapparatusitem{ENV1137}{31;\allowbreak{}\allowbreak{} ms110;\allowbreak{}\allowbreak{} proof\allowbreak{} after\allowbreak{} Corollary4.\allowbreak{}5.\allowbreak{}4}{already\allowbreak{} proved\allowbreak{} of\allowbreak{} of\allowbreak{} lemma}{already\allowbreak{} proved\allowbreak{} of\allowbreak{} lemma}{DER\_\allowbreak{}XVIII-a:\allowbreak{}54}
\DERapparatusitem{ENV1138}{31;\allowbreak{}\allowbreak{} ms110;\allowbreak{}\allowbreak{} last\allowbreak{} sentence\allowbreak{} before\allowbreak{} ms111}{thanks\allowbreak{} to\allowbreak{} 4.\allowbreak{}14.\allowbreak{}3.\allowbreak{}b}{thanks\allowbreak{} to\allowbreak{} 4.\allowbreak{}4.\allowbreak{}13\allowbreak{} b}{DER\_\allowbreak{}XVIII-a:\allowbreak{}55}
\DERapparatusitem{ENV1139}{31;\allowbreak{}\allowbreak{} ms111;\allowbreak{}\allowbreak{} inclusion\allowbreak{} after\allowbreak{} the\allowbreak{} displayed\allowbreak{} definition\allowbreak{} of\allowbreak{} I}{\textbackslash{}\allowbreak{}mathcal\{\allowbreak{}M\}\allowbreak{}\allowbreak{} =\allowbreak{}\allowbreak{} \textbackslash{}\allowbreak{}Ind'\_\allowbreak{}\{\allowbreak{}\textbackslash{}\allowbreak{}pi\}\allowbreak{}C}{\textbackslash{}\allowbreak{}mathcal\{\allowbreak{}M\}\allowbreak{}\allowbreak{} =\allowbreak{}\allowbreak{} \textbackslash{}\allowbreak{}Ind'\_\allowbreak{}\{\allowbreak{}\textbackslash{}\allowbreak{}pi\_\allowbreak{}0\}\allowbreak{}C}{DER\_\allowbreak{}XVIII-a:\allowbreak{}56}
\DERapparatusitem{ENV1140}{32;\allowbreak{}\allowbreak{} ms114;\allowbreak{}\allowbreak{} first\allowbreak{} paragraph}{\textbackslash{}\allowbreak{}card\textbackslash{}\allowbreak{}Fl\allowbreak{} J'\_\allowbreak{}0\allowbreak{} \textbackslash{}\allowbreak{}leqslant\allowbreak{} \textbackslash{}\allowbreak{}pi\_\allowbreak{}0}{\textbackslash{}\allowbreak{}card\textbackslash{}\allowbreak{}Fl\allowbreak{} J'\_\allowbreak{}0\allowbreak{} \textbackslash{}\allowbreak{}leqslant\allowbreak{} \textbackslash{}\allowbreak{}pi}{DER\_\allowbreak{}XVIII-a:\allowbreak{}57}
\DERapparatusitem{ENV1141}{34;\allowbreak{}\allowbreak{} ms117–118,\allowbreak{}proof4⇒5}{and\allowbreak{} \$\allowbreak{}\textbackslash{}\allowbreak{}Filt\textasciicircum{}\allowbreak{}\textbackslash{}\allowbreak{}pi\$\allowbreak{}\allowbreak{} being\allowbreak{} stable\allowbreak{} under\allowbreak{} \$\allowbreak{}\textbackslash{}\allowbreak{}varinjlim\$\allowbreak{}}{and\allowbreak{} \$\allowbreak{}\textbackslash{}\allowbreak{}Filt\textasciicircum{}\allowbreak{}\textbackslash{}\allowbreak{}pi\$\allowbreak{}\allowbreak{} being\allowbreak{} stable}{DER\_\allowbreak{}XVIII-a:\allowbreak{}60}
\DERapparatusitem{ENV1142}{34;\allowbreak{}\allowbreak{} ms117–118,\allowbreak{}proof4⇒5}{large\allowbreak{} relative\allowbreak{} to\allowbreak{} \$\allowbreak{}\textbackslash{}\allowbreak{}pi\_\allowbreak{}0\$\allowbreak{},\allowbreak{}\allowbreak{} this\allowbreak{} implies\allowbreak{} that\allowbreak{} \$\allowbreak{}\textbackslash{}\allowbreak{}Filt\textasciicircum{}\allowbreak{}\textbackslash{}\allowbreak{}pi\$\allowbreak{}\allowbreak{} is\allowbreak{} stable\allowbreak{} under\allowbreak{} direct\allowbreak{} factors}{under\allowbreak{} direct\allowbreak{} factors\allowbreak{} (\allowbreak{}argument\allowbreak{} by\allowbreak{} argument)\allowbreak{}}{DER\_\allowbreak{}XVIII-a:\allowbreak{}60}
\DERapparatusitem{ENV1143}{34;\allowbreak{}\allowbreak{} running\allowbreak{} text\allowbreak{} or\allowbreak{} formula}{\textbackslash{}\allowbreak{}text\{\allowbreak{}def\}\allowbreak{}}{\textbackslash{}\allowbreak{}mathrm\{\allowbreak{}def\}\allowbreak{}}{TYPESETTING:\allowbreak{}7;\allowbreak{}\allowbreak{} NOTATION\_\allowbreak{}DER:\allowbreak{}17}
\DERapparatusitem{ENV1144}{34;\allowbreak{}\allowbreak{} running\allowbreak{} text\allowbreak{} or\allowbreak{} formula}{\textbackslash{}\allowbreak{}text\{\allowbreak{}def\}\allowbreak{}}{\textbackslash{}\allowbreak{}mathrm\{\allowbreak{}def\}\allowbreak{}}{TYPESETTING:\allowbreak{}7;\allowbreak{}\allowbreak{} NOTATION\_\allowbreak{}DER:\allowbreak{}17}
\DERapparatusitem{ENV1145}{35;\allowbreak{}\allowbreak{} ms120,\allowbreak{}evaluation\allowbreak{} right\allowbreak{} adjoint}{x,\allowbreak{}y}{y,\allowbreak{}x}{DER\_\allowbreak{}XVIII-a:\allowbreak{}61}
\DERapparatusitem{ENV1146}{37;\allowbreak{}\allowbreak{} ms123,\allowbreak{}proof4.\allowbreak{}5.\allowbreak{}10,\allowbreak{}first\allowbreak{} sentence}{since\allowbreak{} for\allowbreak{} \$\allowbreak{}\textbackslash{}\allowbreak{}pi\allowbreak{} \textbackslash{}\allowbreak{}geqslant\allowbreak{} \textbackslash{}\allowbreak{}pi'\$\allowbreak{},\allowbreak{}\allowbreak{} \$\allowbreak{}\textbackslash{}\allowbreak{}Filt\textasciicircum{}\allowbreak{}\textbackslash{}\allowbreak{}pi\$\allowbreak{}\allowbreak{} is\allowbreak{} a\allowbreak{} full\allowbreak{} subcategory\allowbreak{} of\allowbreak{} \$\allowbreak{}C\textasciicircum{}\allowbreak{}\textbackslash{}\allowbreak{}wedge(\allowbreak{}\textbackslash{}\allowbreak{}pi')\allowbreak{}\$\allowbreak{}}{since\allowbreak{} for\allowbreak{} \$\allowbreak{}\textbackslash{}\allowbreak{}pi\allowbreak{} \textbackslash{}\allowbreak{}geqslant\allowbreak{} \textbackslash{}\allowbreak{}pi\_\allowbreak{}1\$\allowbreak{},\allowbreak{}\allowbreak{} \$\allowbreak{}\textbackslash{}\allowbreak{}Filt\textasciicircum{}\allowbreak{}\textbackslash{}\allowbreak{}pi\$\allowbreak{}\allowbreak{} is\allowbreak{} a\allowbreak{} full\allowbreak{} subcategory\allowbreak{} of\allowbreak{} \$\allowbreak{}C\textasciicircum{}\allowbreak{}\textbackslash{}\allowbreak{}wedge(\allowbreak{}\textbackslash{}\allowbreak{}pi)\allowbreak{}\$\allowbreak{}}{DER\_\allowbreak{}XVIII-a:\allowbreak{}63}
\DERapparatusitem{ENV1147}{37;\allowbreak{}\allowbreak{} ms123,\allowbreak{}proof4.\allowbreak{}5.\allowbreak{}10,\allowbreak{}remaining\allowbreak{} bound\allowbreak{} paragraph}{\textbackslash{}\allowbreak{}pi'}{\textbackslash{}\allowbreak{}pi}{DER\_\allowbreak{}XVIII-a:\allowbreak{}64}
\DERapparatusitem{ENV1148}{37;\allowbreak{}\allowbreak{} ms123,\allowbreak{}proof4.\allowbreak{}5.\allowbreak{}10,\allowbreak{}remaining\allowbreak{} bound\allowbreak{} paragraph}{\textbackslash{}\allowbreak{}pi'}{\textbackslash{}\allowbreak{}pi}{DER\_\allowbreak{}XVIII-a:\allowbreak{}64}
\DERapparatusitem{ENV1149}{37;\allowbreak{}\allowbreak{} ms123,\allowbreak{}proof4.\allowbreak{}5.\allowbreak{}10,\allowbreak{}remaining\allowbreak{} bound\allowbreak{} paragraph}{\textbackslash{}\allowbreak{}pi'}{\textbackslash{}\allowbreak{}pi}{DER\_\allowbreak{}XVIII-a:\allowbreak{}64}
\DERapparatusitem{ENV1150}{37;\allowbreak{}\allowbreak{} ms123,\allowbreak{}proof4.\allowbreak{}5.\allowbreak{}10,\allowbreak{}remaining\allowbreak{} bound\allowbreak{} paragraph}{\textbackslash{}\allowbreak{}pi'}{\textbackslash{}\allowbreak{}pi}{DER\_\allowbreak{}XVIII-a:\allowbreak{}64}
\DERapparatusitem{ENV1151}{37;\allowbreak{}\allowbreak{} ms123,\allowbreak{}proof4.\allowbreak{}5.\allowbreak{}10,\allowbreak{}remaining\allowbreak{} bound\allowbreak{} paragraph}{\textbackslash{}\allowbreak{}pi'}{\textbackslash{}\allowbreak{}pi}{DER\_\allowbreak{}XVIII-a:\allowbreak{}64}
\DERapparatusitem{ENV1152}{37;\allowbreak{}\allowbreak{} ms123,\allowbreak{}proof4.\allowbreak{}5.\allowbreak{}10,\allowbreak{}remaining\allowbreak{} bound\allowbreak{} paragraph}{\textbackslash{}\allowbreak{}pi'}{\textbackslash{}\allowbreak{}pi}{DER\_\allowbreak{}XVIII-a:\allowbreak{}64}
\DERapparatusitem{ENV1153}{37;\allowbreak{}\allowbreak{} ms123,\allowbreak{}proof4.\allowbreak{}5.\allowbreak{}10,\allowbreak{}remaining\allowbreak{} bound\allowbreak{} paragraph}{\textbackslash{}\allowbreak{}pi'}{\textbackslash{}\allowbreak{}pi}{DER\_\allowbreak{}XVIII-a:\allowbreak{}64}
\DERapparatusitem{ENV1154}{37;\allowbreak{}\allowbreak{} ms123,\allowbreak{}proof4.\allowbreak{}5.\allowbreak{}10,\allowbreak{}remaining\allowbreak{} bound\allowbreak{} paragraph}{\textbackslash{}\allowbreak{}pi\_\allowbreak{}0}{\textbackslash{}\allowbreak{}pi}{DER\_\allowbreak{}XVIII-a:\allowbreak{}64}
\DERapparatusitem{ENV1155}{38;\allowbreak{}\allowbreak{} ms125,\allowbreak{}4.\allowbreak{}5.\allowbreak{}13\allowbreak{} conditionc}{\$\allowbreak{}x\_\allowbreak{}\textbackslash{}\allowbreak{}alpha\allowbreak{} \textbackslash{}\allowbreak{}longrightarrow\allowbreak{} x\$\allowbreak{},\allowbreak{}\allowbreak{} indexed\allowbreak{} by\allowbreak{} a\allowbreak{} set\allowbreak{} \$\allowbreak{}A\$\allowbreak{}}{\$\allowbreak{}x\_\allowbreak{}\textbackslash{}\allowbreak{}alpha\allowbreak{} \textbackslash{}\allowbreak{}longrightarrow\allowbreak{} x\$\allowbreak{},\allowbreak{}\allowbreak{} with\allowbreak{} \$\allowbreak{}x\_\allowbreak{}\textbackslash{}\allowbreak{}alpha\textbackslash{}\allowbreak{}in\textbackslash{}\allowbreak{}Ob\allowbreak{} C\$\allowbreak{}\allowbreak{} for\allowbreak{} every\allowbreak{} \$\allowbreak{}\textbackslash{}\allowbreak{}alpha\textbackslash{}\allowbreak{}in\allowbreak{} A\$\allowbreak{},\allowbreak{}\allowbreak{} indexed\allowbreak{} by\allowbreak{} a\allowbreak{} set\allowbreak{} \$\allowbreak{}A\$\allowbreak{}}{DER\_\allowbreak{}XVIII-a:\allowbreak{}65}
\DERapparatusitem{ENV1156}{41;\allowbreak{}\allowbreak{} ms134,\allowbreak{}Lemma4.\allowbreak{}5.\allowbreak{}14\allowbreak{} implication}{\$\allowbreak{}\textbackslash{}\allowbreak{}mathcal\{\allowbreak{}M\}\allowbreak{}\_\allowbreak{}0\$\allowbreak{}\allowbreak{} a\allowbreak{} strictly\allowbreak{} full\allowbreak{} subcategory\allowbreak{} of\allowbreak{} \$\allowbreak{}\textbackslash{}\allowbreak{}mathcal\{\allowbreak{}M\}\allowbreak{}\$\allowbreak{}.\allowbreak{}\allowbreak{} Consider\allowbreak{} the\allowbreak{} conditions}{\$\allowbreak{}\textbackslash{}\allowbreak{}mathcal\{\allowbreak{}M\}\allowbreak{}\_\allowbreak{}0\$\allowbreak{}\allowbreak{} a\allowbreak{} strictly\allowbreak{} full\allowbreak{} subcategory\allowbreak{} of\allowbreak{} \$\allowbreak{}\textbackslash{}\allowbreak{}mathcal\{\allowbreak{}M\}\allowbreak{}\$\allowbreak{},\allowbreak{}\allowbreak{} stable\allowbreak{} under\allowbreak{} direct\allowbreak{} factors.\allowbreak{}\allowbreak{} Consider\allowbreak{} the\allowbreak{} conditions}{DER\_\allowbreak{}XVIII-a:\allowbreak{}69}
\DERapparatusitem{ENV1157}{42;\allowbreak{}\allowbreak{} ms135,\allowbreak{}footnote53}{but\allowbreak{} true\allowbreak{} in\allowbreak{} the\allowbreak{} particular\allowbreak{} case\allowbreak{} below.\allowbreak{}}{and\allowbreak{} also\allowbreak{} false\allowbreak{} in\allowbreak{} the\allowbreak{} particular\allowbreak{} case\allowbreak{} below\allowbreak{} for\allowbreak{} nontrivial\allowbreak{} \$\allowbreak{}H\$\allowbreak{}.\allowbreak{}}{DER\_\allowbreak{}XVIII-a:\allowbreak{}71}
\DERapparatusitem{ENV1158}{43;\allowbreak{}\allowbreak{} ms137,\allowbreak{}Corollary4.\allowbreak{}5.\allowbreak{}16\allowbreak{} and\allowbreak{} its\allowbreak{} proof}{a\allowbreak{} family\}\allowbreak{}\allowbreak{} epimorphic\allowbreak{} \textbackslash{}\allowbreak{}underline\{\allowbreak{}universally\}\allowbreak{}}{a\allowbreak{} family\}\allowbreak{}\allowbreak{} strictly\allowbreak{} epimorphic\allowbreak{} \textbackslash{}\allowbreak{}underline\{\allowbreak{}universally\}\allowbreak{}}{DER\_\allowbreak{}XVIII-a:\allowbreak{}72}
\DERapparatusitem{ENV1159}{43;\allowbreak{}\allowbreak{} ms138,\allowbreak{}remark\allowbreak{} after4.\allowbreak{}5.\allowbreak{}17}{Nor\allowbreak{} is\allowbreak{} it\allowbreak{} the\allowbreak{} category\allowbreak{} of\allowbreak{} the}{It\allowbreak{} is\allowbreak{} also\allowbreak{} the\allowbreak{} category\allowbreak{} of\allowbreak{} the}{DER\_\allowbreak{}XVIII-a:\allowbreak{}73}
\DERapparatusitem{ENV1160}{44;\allowbreak{}\allowbreak{} ms139,\allowbreak{}\allowbreak{} PTA\allowbreak{} implies\allowbreak{} d}{then\allowbreak{} exists\allowbreak{} an\allowbreak{} infinite\allowbreak{} cardinal}{then\allowbreak{} there\allowbreak{} exists\allowbreak{} an\allowbreak{} infinite\allowbreak{} cardinal}{DER\_\allowbreak{}XVIII-b:\allowbreak{}4}
\DERapparatusitem{ENV1161}{44;\allowbreak{}\allowbreak{} ms140,\allowbreak{}\allowbreak{} indexing\allowbreak{} the\allowbreak{} transfinite\allowbreak{} construction}{for\allowbreak{} every\allowbreak{} small\allowbreak{} cardinal\allowbreak{} \$\allowbreak{}\textbackslash{}\allowbreak{}alpha\$\allowbreak{}}{for\allowbreak{} every\allowbreak{} small\allowbreak{} ordinal\allowbreak{} \$\allowbreak{}\textbackslash{}\allowbreak{}alpha\$\allowbreak{}}{DER\_\allowbreak{}XVIII-b:\allowbreak{}7}
\DERapparatusitem{ENV1162}{45;\allowbreak{}\allowbreak{} ms141,\allowbreak{}\allowbreak{} ordinal\allowbreak{} index\allowbreak{} set}{of\allowbreak{} cardinals\allowbreak{} \$\allowbreak{}\textbackslash{}\allowbreak{}alpha<\textbackslash{}\allowbreak{}alpha\_\allowbreak{}0\$\allowbreak{}}{of\allowbreak{} ordinals\allowbreak{} \$\allowbreak{}\textbackslash{}\allowbreak{}alpha<\textbackslash{}\allowbreak{}alpha\_\allowbreak{}0\$\allowbreak{}}{DER\_\allowbreak{}XVIII-b:\allowbreak{}14}
\DERapparatusitem{ENV1163}{45;\allowbreak{}\allowbreak{} ms141,\allowbreak{}\allowbreak{} required\allowbreak{} closure}{\$\allowbreak{}C\_\allowbreak{}\{\allowbreak{}\textbackslash{}\allowbreak{}alpha\_\allowbreak{}0\}\allowbreak{}\$\allowbreak{}\allowbreak{} is\allowbreak{} stable\allowbreak{} under\allowbreak{} \$\allowbreak{}\textbackslash{}\allowbreak{}varinjlim\$\allowbreak{}\allowbreak{} of\allowbreak{} cardinality\allowbreak{} \$\allowbreak{}\textbackslash{}\allowbreak{}leqslant\textbackslash{}\allowbreak{}pi\$\allowbreak{}}{\$\allowbreak{}C\_\allowbreak{}\{\allowbreak{}\textbackslash{}\allowbreak{}alpha\_\allowbreak{}0\}\allowbreak{}\$\allowbreak{}\allowbreak{} is\allowbreak{} stable\allowbreak{} under\allowbreak{} \$\allowbreak{}\textbackslash{}\allowbreak{}varinjlim\$\allowbreak{}\allowbreak{} of\allowbreak{} cardinality\allowbreak{} \$\allowbreak{}\textbackslash{}\allowbreak{}leqslant\textbackslash{}\allowbreak{}pi\_\allowbreak{}0\$\allowbreak{}}{DER\_\allowbreak{}XVIII-b:\allowbreak{}9}
\DERapparatusitem{ENV1164}{45;\allowbreak{}\allowbreak{} ms143,\allowbreak{}\allowbreak{} successor\allowbreak{} step\allowbreak{} display}{\textbackslash{}\allowbreak{}cardred\allowbreak{} C\_\allowbreak{}\textbackslash{}\allowbreak{}alpha\allowbreak{} \textbackslash{}\allowbreak{}\allowbreak{} \textbackslash{}\allowbreak{}leqslant\textbackslash{}\allowbreak{}\allowbreak{} (\allowbreak{}\textbackslash{}\allowbreak{}pi\_\allowbreak{}1\textasciicircum{}\allowbreak{}\{\allowbreak{}\textbackslash{}\allowbreak{}pi\_\allowbreak{}0\}\allowbreak{})\allowbreak{}\textasciicircum{}\allowbreak{}\{\allowbreak{}\textbackslash{}\allowbreak{}pi\_\allowbreak{}0\}\allowbreak{}}{\textbackslash{}\allowbreak{}cardred\allowbreak{} C\_\allowbreak{}\{\allowbreak{}\textbackslash{}\allowbreak{}alpha+1\}\allowbreak{}\allowbreak{} \textbackslash{}\allowbreak{}\allowbreak{} \textbackslash{}\allowbreak{}leqslant\textbackslash{}\allowbreak{}\allowbreak{} (\allowbreak{}\textbackslash{}\allowbreak{}pi\_\allowbreak{}1\textasciicircum{}\allowbreak{}\{\allowbreak{}\textbackslash{}\allowbreak{}pi\_\allowbreak{}0\}\allowbreak{})\allowbreak{}\textasciicircum{}\allowbreak{}\{\allowbreak{}\textbackslash{}\allowbreak{}pi\_\allowbreak{}0\}\allowbreak{}}{DER\_\allowbreak{}XVIII-b:\allowbreak{}12}
\DERapparatusitem{ENV1165}{45;\allowbreak{}\allowbreak{} ms143,\allowbreak{}\allowbreak{} limit\allowbreak{} step\allowbreak{} display}{\textbackslash{}\allowbreak{}alpha\allowbreak{} \textbackslash{}\allowbreak{}\allowbreak{} \textbackslash{}\allowbreak{}longrightarrow}{\textbackslash{}\allowbreak{}alpha\textbackslash{}\allowbreak{}leqslant}{DER\_\allowbreak{}XVIII-b:\allowbreak{}13}
\DERapparatusitem{ENV1166}{46;\allowbreak{}\allowbreak{} ms144,\allowbreak{}\allowbreak{} reference\allowbreak{} in\allowbreak{} proof\allowbreak{} of\allowbreak{} 4.\allowbreak{}6.\allowbreak{}1(\allowbreak{}2)\allowbreak{}}{by\allowbreak{} virtue\allowbreak{} of\allowbreak{} 4.\allowbreak{}5.\allowbreak{}19\allowbreak{} a\allowbreak{} \$\allowbreak{}\textbackslash{}\allowbreak{}Longrightarrow\$\allowbreak{}\allowbreak{} c}{by\allowbreak{} virtue\allowbreak{} of\allowbreak{} 4.\allowbreak{}5.\allowbreak{}14\allowbreak{} a\allowbreak{} \$\allowbreak{}\textbackslash{}\allowbreak{}Longrightarrow\$\allowbreak{}\allowbreak{} c}{DER\_\allowbreak{}XVIII-b:\allowbreak{}16}
\DERapparatusitem{ENV1167}{46;\allowbreak{}\allowbreak{} ms144,\allowbreak{}\allowbreak{} full\allowbreak{} faithfulness\allowbreak{} sentence}{The\allowbreak{} fact\allowbreak{} that\allowbreak{} this\allowbreak{} factor\allowbreak{} is\allowbreak{} fully\allowbreak{} faithful}{The\allowbreak{} fact\allowbreak{} that\allowbreak{} this\allowbreak{} functor\allowbreak{} is\allowbreak{} fully\allowbreak{} faithful}{DER\_\allowbreak{}XVIII-b:\allowbreak{}15}
\DERapparatusitem{ENV1168}{46;\allowbreak{}\allowbreak{} ms144,\allowbreak{}\allowbreak{} final\allowbreak{} term\allowbreak{} of\allowbreak{} Hom\allowbreak{} comparison}{}{\_\allowbreak{}0}{DER\_\allowbreak{}XVIII-b:\allowbreak{}17}
\DERapparatusitem{ENV1169}{46;\allowbreak{}\allowbreak{} ms145,\allowbreak{}\allowbreak{} parenthesis\allowbreak{} in\allowbreak{} (\allowbreak{}4.\allowbreak{}6.\allowbreak{}2.\allowbreak{}1)\allowbreak{}}{(\allowbreak{}=\allowbreak{}2\textasciicircum{}\allowbreak{}\{\allowbreak{}\textbackslash{}\allowbreak{}pi\}\allowbreak{}\textbackslash{}\allowbreak{}\allowbreak{} \textbackslash{}\allowbreak{}textit\{\allowbreak{}if\}\allowbreak{}\textbackslash{}\allowbreak{}\allowbreak{} \textbackslash{}\allowbreak{}pi\textbackslash{}\allowbreak{}geqslant\textbackslash{}\allowbreak{}pi\_\allowbreak{}0)\allowbreak{}}{(\allowbreak{}=\allowbreak{}2\textasciicircum{}\allowbreak{}\{\allowbreak{}\textbackslash{}\allowbreak{}pi\}\allowbreak{}\textbackslash{}\allowbreak{}\allowbreak{} \textbackslash{}\allowbreak{}textit\{\allowbreak{}if\}\allowbreak{}\textbackslash{}\allowbreak{}\allowbreak{} \textbackslash{}\allowbreak{}pi\textbackslash{}\allowbreak{}geqslant\textbackslash{}\allowbreak{}pi\_\allowbreak{}1\textbackslash{}\allowbreak{}geqslant2)\allowbreak{}}{DER\_\allowbreak{}XVIII-b:\allowbreak{}18}
\DERapparatusitem{ENV1170}{47;\allowbreak{}\allowbreak{} ms148,\allowbreak{}\allowbreak{} conclusion\allowbreak{} after\allowbreak{} pointwise\allowbreak{} calculation}{This\allowbreak{} completes\allowbreak{} the\allowbreak{} proof\allowbreak{} of\allowbreak{} c\allowbreak{} \$\allowbreak{}\textbackslash{}\allowbreak{}Longrightarrow\$\allowbreak{}\allowbreak{} d\allowbreak{} in\allowbreak{} \$\allowbreak{}(\allowbreak{}*)\allowbreak{}\$\allowbreak{},\allowbreak{}\allowbreak{} page\allowbreak{} 139.\allowbreak{}}{This\allowbreak{} completes\allowbreak{} the\allowbreak{} proof\allowbreak{} of\allowbreak{} d\allowbreak{} \$\allowbreak{}\textbackslash{}\allowbreak{}Longrightarrow\$\allowbreak{}\allowbreak{} c\allowbreak{} in\allowbreak{} \$\allowbreak{}(\allowbreak{}*)\allowbreak{}\$\allowbreak{},\allowbreak{}\allowbreak{} page\allowbreak{} 139.\allowbreak{}}{DER\_\allowbreak{}XVIII-b:\allowbreak{}19}
\DERapparatusitem{ENV1171}{48;\allowbreak{}\allowbreak{} ms149,\allowbreak{}\allowbreak{} Theorem4.\allowbreak{}6.\allowbreak{}3(\allowbreak{}2)\allowbreak{},\allowbreak{}\allowbreak{} beta(\allowbreak{}pi0)\allowbreak{}}{For\allowbreak{} every\allowbreak{} arrow\allowbreak{} \$\allowbreak{}y\textbackslash{}\allowbreak{}leqslant\allowbreak{} x\$\allowbreak{}\allowbreak{} in}{For\allowbreak{} every\allowbreak{} arrow\allowbreak{} \$\allowbreak{}y\textbackslash{}\allowbreak{}longrightarrow\allowbreak{} x\$\allowbreak{}\allowbreak{} in}{DER\_\allowbreak{}XVIII-b:\allowbreak{}22}
\DERapparatusitem{ENV1172}{48;\allowbreak{}\allowbreak{} ms150,\allowbreak{}\allowbreak{} Theorem\allowbreak{} 4.\allowbreak{}6.\allowbreak{}3(\allowbreak{}3)\allowbreak{}c\allowbreak{} and\allowbreak{} c-prime}{a\allowbreak{} small\allowbreak{} category\allowbreak{} \$\allowbreak{}S\$\allowbreak{},\allowbreak{}\allowbreak{} and\allowbreak{} an\allowbreak{} equivalence\allowbreak{} \$\allowbreak{}\textbackslash{}\allowbreak{}mathcal\{\allowbreak{}M\}\allowbreak{}\textbackslash{}\allowbreak{}simeq\$\allowbreak{}\allowbreak{} \$\allowbreak{}\textbackslash{}\allowbreak{}Ind'\_\allowbreak{}\{\allowbreak{}\textbackslash{}\allowbreak{}pi\_\allowbreak{}0\}\allowbreak{}(\allowbreak{}S)\allowbreak{}\$\allowbreak{}}{a\allowbreak{} small\allowbreak{} category\allowbreak{} \$\allowbreak{}S\$\allowbreak{}\allowbreak{} admitting\allowbreak{} direct\allowbreak{} limits\allowbreak{} of\allowbreak{} categories\allowbreak{} \$\allowbreak{}I\$\allowbreak{}\allowbreak{} with\allowbreak{} \$\allowbreak{}\textbackslash{}\allowbreak{}operatorname\{\allowbreak{}card\}\allowbreak{}\textbackslash{}\allowbreak{}operatorname\{\allowbreak{}Fl\}\allowbreak{}I\textbackslash{}\allowbreak{}leqslant\textbackslash{}\allowbreak{}pi\_\allowbreak{}0\$\allowbreak{},\allowbreak{}\allowbreak{} and\allowbreak{} an\allowbreak{} equivalence\allowbreak{} \$\allowbreak{}\textbackslash{}\allowbreak{}mathcal\{\allowbreak{}M\}\allowbreak{}\textbackslash{}\allowbreak{}simeq\$\allowbreak{}\allowbreak{} \$\allowbreak{}\textbackslash{}\allowbreak{}Ind'\_\allowbreak{}\{\allowbreak{}\textbackslash{}\allowbreak{}pi\_\allowbreak{}0\}\allowbreak{}(\allowbreak{}S)\allowbreak{}\$\allowbreak{}}{DER\_\allowbreak{}XVIII-b:\allowbreak{}21}
\DERapparatusitem{ENV1173}{49;\allowbreak{}\allowbreak{} ms151,\allowbreak{}\allowbreak{} first\allowbreak{} sentence\allowbreak{} of\allowbreak{} proof}{The\allowbreak{} first\allowbreak{} assertions\allowbreak{} in\allowbreak{} 2°)\allowbreak{}}{The\allowbreak{} first\allowbreak{} assertion\allowbreak{} in\allowbreak{} 2°)\allowbreak{}}{DER\_\allowbreak{}XVIII-b:\allowbreak{}23}
\DERapparatusitem{ENV1174}{49;\allowbreak{}\allowbreak{} ms151,\allowbreak{}\allowbreak{} second\allowbreak{} displayed\allowbreak{} functor}{}{\textbackslash{}\allowbreak{}mathcal\allowbreak{} }{DER\_\allowbreak{}XVIII-b:\allowbreak{}25}
\DERapparatusitem{ENV1175}{50;\allowbreak{}\allowbreak{} running\allowbreak{} text\allowbreak{} or\allowbreak{} formula}{\textbackslash{}\allowbreak{}text\{\allowbreak{}def\}\allowbreak{}}{\textbackslash{}\allowbreak{}mathrm\{\allowbreak{}def\}\allowbreak{}}{TYPESETTING:\allowbreak{}7;\allowbreak{}\allowbreak{} NOTATION\_\allowbreak{}DER:\allowbreak{}17}
\DERapparatusitem{ENV1176}{51;\allowbreak{}\allowbreak{} ms157,\allowbreak{}\allowbreak{} Proposition\allowbreak{} 4.\allowbreak{}6.\allowbreak{}7}{d)\allowbreak{}\allowbreak{} below)\allowbreak{},\allowbreak{}\allowbreak{} is\allowbreak{} equivalent\allowbreak{} to\allowbreak{} the\allowbreak{} set}{d)\allowbreak{}\allowbreak{} below)\allowbreak{},\allowbreak{}\allowbreak{} are\allowbreak{} equivalent\allowbreak{} to\allowbreak{} the\allowbreak{} set}{DER\_\allowbreak{}XVIII-b:\allowbreak{}29}
\DERapparatusitem{ENV1177}{51;\allowbreak{}\allowbreak{} running\allowbreak{} text\allowbreak{} or\allowbreak{} formula}{\textbackslash{}\allowbreak{}text\{\allowbreak{}def\}\allowbreak{}}{\textbackslash{}\allowbreak{}mathrm\{\allowbreak{}def\}\allowbreak{}}{TYPESETTING:\allowbreak{}7;\allowbreak{}\allowbreak{} NOTATION\_\allowbreak{}DER:\allowbreak{}17}
\DERapparatusitem{ENV1178}{53;\allowbreak{}\allowbreak{} ms163,\allowbreak{}\allowbreak{} Corollary\allowbreak{} 4.\allowbreak{}7.\allowbreak{}2(\allowbreak{}1)\allowbreak{}}{\$\allowbreak{}\textbackslash{}\allowbreak{}mathrm\{\allowbreak{}Ob\}\allowbreak{}\textbackslash{}\allowbreak{},\allowbreak{}\textbackslash{}\allowbreak{}mathrm\{\allowbreak{}Filt\}\allowbreak{}\textasciicircum{}\allowbreak{}\{\allowbreak{}\textbackslash{}\allowbreak{}pi\}\allowbreak{}\$\allowbreak{}\allowbreak{} \textbackslash{}\allowbreak{}textit\{\allowbreak{}where\allowbreak{} \$\allowbreak{}F\allowbreak{} \textbackslash{}\allowbreak{}in\$\allowbreak{}\}\allowbreak{}}{\$\allowbreak{}\textbackslash{}\allowbreak{}mathrm\{\allowbreak{}Ob\}\allowbreak{}\textbackslash{}\allowbreak{},\allowbreak{}\textbackslash{}\allowbreak{}mathrm\{\allowbreak{}Filt\}\allowbreak{}\textasciicircum{}\allowbreak{}\{\allowbreak{}\textbackslash{}\allowbreak{}pi\}\allowbreak{}\$\allowbreak{}\allowbreak{} \textbackslash{}\allowbreak{}textit\{\allowbreak{}or\allowbreak{} \$\allowbreak{}F\allowbreak{} \textbackslash{}\allowbreak{}in\$\allowbreak{}\}\allowbreak{}}{DER\_\allowbreak{}XVIII-b:\allowbreak{}33}
\DERapparatusitem{ENV1179}{53;\allowbreak{}\allowbreak{} ms163,\allowbreak{}\allowbreak{} equality\allowbreak{} in\allowbreak{} Corollary\allowbreak{} 4.\allowbreak{}7.\allowbreak{}2(\allowbreak{}2)\allowbreak{}}{3.\allowbreak{}2}{2.\allowbreak{}3}{DER\_\allowbreak{}XVIII-b:\allowbreak{}34}
\DERapparatusitem{ENV1180}{54;\allowbreak{}\allowbreak{} ms164,\allowbreak{}\allowbreak{} identity\allowbreak{} of\allowbreak{} the\allowbreak{} ambient\allowbreak{} category}{\textbackslash{}\allowbreak{}underbrace\{\allowbreak{}\textbackslash{}\allowbreak{}Ind\_\allowbreak{}\{\allowbreak{}\textbackslash{}\allowbreak{}pi\}\allowbreak{}(\allowbreak{}C)\allowbreak{}\}\allowbreak{}\_\allowbreak{}\{\allowbreak{}\textbackslash{}\allowbreak{}simeq\allowbreak{} \textbackslash{}\allowbreak{}Ind\_\allowbreak{}\{\allowbreak{}\textbackslash{}\allowbreak{}pi\}\allowbreak{}\allowbreak{} C\_\allowbreak{}0\}\allowbreak{}}{\textbackslash{}\allowbreak{}underbrace\{\allowbreak{}\textbackslash{}\allowbreak{}Ind\_\allowbreak{}\{\allowbreak{}\textbackslash{}\allowbreak{}pi\_\allowbreak{}0\}\allowbreak{}(\allowbreak{}C)\allowbreak{}\}\allowbreak{}\_\allowbreak{}\{\allowbreak{}\textbackslash{}\allowbreak{}simeq\allowbreak{} \textbackslash{}\allowbreak{}Ind\_\allowbreak{}\{\allowbreak{}\textbackslash{}\allowbreak{}pi\_\allowbreak{}0\}\allowbreak{}\allowbreak{} C\_\allowbreak{}0\}\allowbreak{}}{DER\_\allowbreak{}XVIII-b:\allowbreak{}35}
\DERapparatusitem{ENV1181}{55;\allowbreak{}\allowbreak{} ms166,\allowbreak{}\allowbreak{} colimit\allowbreak{} comparison}{\textbackslash{}\allowbreak{}varinjlim\_\allowbreak{}\{\allowbreak{}j\allowbreak{} \textbackslash{}\allowbreak{}in\allowbreak{} J\}\allowbreak{}\allowbreak{} x\_\allowbreak{}j\allowbreak{} \textbackslash{}\allowbreak{};\allowbreak{}\textbackslash{}\allowbreak{}simeq\textbackslash{}\allowbreak{};\allowbreak{}\allowbreak{} x}{\textbackslash{}\allowbreak{}varinjlim\_\allowbreak{}\{\allowbreak{}j\allowbreak{} \textbackslash{}\allowbreak{}in\allowbreak{} J\}\allowbreak{}\allowbreak{} x\_\allowbreak{}j\allowbreak{} \textbackslash{}\allowbreak{};\allowbreak{}\textbackslash{}\allowbreak{}simeq\textbackslash{}\allowbreak{};\allowbreak{}\allowbreak{} F}{DER\_\allowbreak{}XVIII-b:\allowbreak{}37}
\DERapparatusitem{ENV1182}{55;\allowbreak{}\allowbreak{} ms166,\allowbreak{}\allowbreak{} beginning\allowbreak{} of\allowbreak{} cardinal-bound\allowbreak{} paragraph}{There\allowbreak{} remains\allowbreak{} the\allowbreak{} bound\allowbreak{} on\allowbreak{} cardinals\allowbreak{} for\allowbreak{} \$\allowbreak{}F\allowbreak{} \textbackslash{}\allowbreak{}in\allowbreak{} \textbackslash{}\allowbreak{}Ob\allowbreak{} \textbackslash{}\allowbreak{}Filt\textasciicircum{}\allowbreak{}\{\allowbreak{}\textbackslash{}\allowbreak{}pi\}\allowbreak{}\$\allowbreak{}.\allowbreak{}}{There\allowbreak{} remains\allowbreak{} the\allowbreak{} bound\allowbreak{} on\allowbreak{} cardinals\allowbreak{} for\allowbreak{} \$\allowbreak{}F\allowbreak{} \textbackslash{}\allowbreak{}in\allowbreak{} \textbackslash{}\allowbreak{}Ob\allowbreak{} \textbackslash{}\allowbreak{}Filt\textasciicircum{}\allowbreak{}\{\allowbreak{}\textbackslash{}\allowbreak{}pi'\}\allowbreak{}\$\allowbreak{}.\allowbreak{}}{DER\_\allowbreak{}XVIII-b:\allowbreak{}38}
\DERapparatusitem{ENV1183}{55;\allowbreak{}\allowbreak{} ms166,\allowbreak{}\allowbreak{} end\allowbreak{} of\allowbreak{} cardinal-bound\allowbreak{} paragraph}{\textbackslash{}\allowbreak{}card\allowbreak{} I\allowbreak{} \textbackslash{}\allowbreak{}leqslant\allowbreak{} \textbackslash{}\allowbreak{}pi'\textasciicircum{}\allowbreak{}\{\allowbreak{}\textbackslash{}\allowbreak{}pi\_\allowbreak{}0\}\allowbreak{}}{\textbackslash{}\allowbreak{}card\allowbreak{} I\allowbreak{} \textbackslash{}\allowbreak{}leqslant\allowbreak{} \textbackslash{}\allowbreak{}pi'\textasciicircum{}\allowbreak{}\{\allowbreak{}\textbackslash{}\allowbreak{}pi\}\allowbreak{}}{DER\_\allowbreak{}XVIII-b:\allowbreak{}39}
\DERapparatusitem{ENV1184}{55;\allowbreak{}\allowbreak{} ms166,\allowbreak{}\allowbreak{} referring\allowbreak{} to\allowbreak{} the\allowbreak{} presentation\allowbreak{} of\allowbreak{} F}{in\allowbreak{} the\allowbreak{} form\allowbreak{} (\allowbreak{}4.\allowbreak{}7.\allowbreak{}3.\allowbreak{}2)\allowbreak{}}{in\allowbreak{} the\allowbreak{} form\allowbreak{} (\allowbreak{}4.\allowbreak{}7.\allowbreak{}3.\allowbreak{}1)\allowbreak{}}{DER\_\allowbreak{}XVIII-b:\allowbreak{}40}
\DERapparatusitem{ENV1185}{55;\allowbreak{}\allowbreak{} ms167,\allowbreak{}\allowbreak{} end\allowbreak{} of\allowbreak{} proof\allowbreak{} of\allowbreak{} Corollary\allowbreak{} 4.\allowbreak{}7.\allowbreak{}4}{which\allowbreak{} proves\allowbreak{} 4.\allowbreak{}7.\allowbreak{}3}{which\allowbreak{} proves\allowbreak{} 4.\allowbreak{}7.\allowbreak{}4}{DER\_\allowbreak{}XVIII-b:\allowbreak{}41}
\DERapparatusitem{ENV1186}{56;\allowbreak{}\allowbreak{} ms169–170,\allowbreak{}\allowbreak{} counting\allowbreak{} the\allowbreak{} selected\allowbreak{} families}{A\allowbreak{} \textbackslash{}\allowbreak{}in\allowbreak{} \textbackslash{}\allowbreak{}mathfrak\{\allowbreak{}P\}\allowbreak{}\_\allowbreak{}\{\allowbreak{}\textbackslash{}\allowbreak{}leqslant\allowbreak{} \textbackslash{}\allowbreak{}pi\_\allowbreak{}0\}\allowbreak{}(\allowbreak{}\textbackslash{}\allowbreak{}Ob\allowbreak{} J)\allowbreak{}}{A\allowbreak{} \textbackslash{}\allowbreak{}in\allowbreak{} \textbackslash{}\allowbreak{}mathfrak\{\allowbreak{}P\}\allowbreak{}\_\allowbreak{}\{\allowbreak{}\textbackslash{}\allowbreak{}leqslant\allowbreak{} \textbackslash{}\allowbreak{}pi\_\allowbreak{}0\}\allowbreak{}(\allowbreak{}\textbackslash{}\allowbreak{}Ob\allowbreak{} K\_\allowbreak{}\textbackslash{}\allowbreak{}alpha)\allowbreak{}}{DER\_\allowbreak{}XVIII-b:\allowbreak{}43}
\DERapparatusitem{ENV1187}{56;\allowbreak{}\allowbreak{} ms169–170,\allowbreak{}\allowbreak{} counting\allowbreak{} the\allowbreak{} selected\allowbreak{} families}{B\allowbreak{} \textbackslash{}\allowbreak{}in\allowbreak{} \textbackslash{}\allowbreak{}mathfrak\{\allowbreak{}P\}\allowbreak{}\_\allowbreak{}\{\allowbreak{}\textbackslash{}\allowbreak{}pi\_\allowbreak{}0\}\allowbreak{}(\allowbreak{}\textbackslash{}\allowbreak{}mathrm\{\allowbreak{}Fl\}\allowbreak{}\textbackslash{}\allowbreak{},\allowbreak{}\allowbreak{} J)\allowbreak{}}{B\allowbreak{} \textbackslash{}\allowbreak{}in\allowbreak{} \textbackslash{}\allowbreak{}mathfrak\{\allowbreak{}P\}\allowbreak{}\_\allowbreak{}\{\allowbreak{}\textbackslash{}\allowbreak{}pi\_\allowbreak{}0\}\allowbreak{}(\allowbreak{}\textbackslash{}\allowbreak{}mathrm\{\allowbreak{}Fl\}\allowbreak{}\textbackslash{}\allowbreak{},\allowbreak{}\allowbreak{} K\_\allowbreak{}\textbackslash{}\allowbreak{}alpha)\allowbreak{}}{DER\_\allowbreak{}XVIII-b:\allowbreak{}43}
\DERapparatusitem{ENV1188}{56;\allowbreak{}\allowbreak{} ms170,\allowbreak{}\allowbreak{} last\allowbreak{} bound\allowbreak{} on\allowbreak{} successor\allowbreak{} stage}{generated\allowbreak{} by\allowbreak{} this\allowbreak{} set\allowbreak{} of\allowbreak{} arrows\allowbreak{} is\allowbreak{} of\allowbreak{} cardinality\allowbreak{} \$\allowbreak{}\textbackslash{}\allowbreak{}leqslant\allowbreak{} 1\$\allowbreak{}}{generated\allowbreak{} by\allowbreak{} this\allowbreak{} set\allowbreak{} of\allowbreak{} arrows\allowbreak{} is\allowbreak{} of\allowbreak{} cardinality\allowbreak{} \$\allowbreak{}\textbackslash{}\allowbreak{}leqslant\allowbreak{} \textbackslash{}\allowbreak{}pi\$\allowbreak{}}{DER\_\allowbreak{}XVIII-b:\allowbreak{}44}
\DERapparatusitem{ENV1189}{57;\allowbreak{}\allowbreak{} ms171,\allowbreak{}\allowbreak{} filtering\allowbreak{} criterion\allowbreak{} name}{\textbackslash{}\allowbreak{}pi}{\textbackslash{}\allowbreak{}pi\_\allowbreak{}0}{DER\_\allowbreak{}XVIII-b:\allowbreak{}45}
\DERapparatusitem{ENV1190}{57;\allowbreak{}\allowbreak{} ms171,\allowbreak{}\allowbreak{} (\allowbreak{}4.\allowbreak{}7.\allowbreak{}6.\allowbreak{}1)\allowbreak{},\allowbreak{}\allowbreak{} right-hand\allowbreak{} underbrace}{\textbackslash{}\allowbreak{}underbrace\{\allowbreak{}\textbackslash{}\allowbreak{}Filt\textasciicircum{}\allowbreak{}\{\allowbreak{}\textbackslash{}\allowbreak{}pi\textasciicircum{}\allowbreak{}\{\allowbreak{}\textbackslash{}\allowbreak{}pi\_\allowbreak{}0\}\allowbreak{}\}\allowbreak{}\}\allowbreak{}}{\textbackslash{}\allowbreak{}underbrace\{\allowbreak{}\textbackslash{}\allowbreak{}Filt'\textasciicircum{}\allowbreak{}\{\allowbreak{}\textbackslash{}\allowbreak{}pi\textasciicircum{}\allowbreak{}\{\allowbreak{}\textbackslash{}\allowbreak{}pi\_\allowbreak{}0\}\allowbreak{}\}\allowbreak{}\}\allowbreak{}}{DER\_\allowbreak{}XVIII-b:\allowbreak{}46}
\DERapparatusitem{ENV1191}{57;\allowbreak{}\allowbreak{} ms171,\allowbreak{}\allowbreak{} final\allowbreak{} term\allowbreak{} of\allowbreak{} exponent\allowbreak{} calculation}{\textbackslash{}\allowbreak{}pi\textasciicircum{}\allowbreak{}\{\allowbreak{}\textbackslash{}\allowbreak{}pi\_\allowbreak{}0\allowbreak{} \textbackslash{}\allowbreak{}cdot\allowbreak{} \textbackslash{}\allowbreak{}pi\_\allowbreak{}0\}\allowbreak{}\allowbreak{} =\allowbreak{}\allowbreak{} \textbackslash{}\allowbreak{}pi\textasciicircum{}\allowbreak{}\{\allowbreak{}\textbackslash{}\allowbreak{}pi\_\allowbreak{}0\}\allowbreak{}\allowbreak{} =\allowbreak{}\allowbreak{} \textbackslash{}\allowbreak{}pi\$\allowbreak{}}{\textbackslash{}\allowbreak{}pi\textasciicircum{}\allowbreak{}\{\allowbreak{}\textbackslash{}\allowbreak{}pi\_\allowbreak{}0\allowbreak{} \textbackslash{}\allowbreak{}cdot\allowbreak{} \textbackslash{}\allowbreak{}pi\_\allowbreak{}0\}\allowbreak{}\allowbreak{} =\allowbreak{}\allowbreak{} \textbackslash{}\allowbreak{}pi\textasciicircum{}\allowbreak{}\{\allowbreak{}\textbackslash{}\allowbreak{}pi\_\allowbreak{}0\}\allowbreak{}\allowbreak{} =\allowbreak{}\allowbreak{} \textbackslash{}\allowbreak{}pi'\$\allowbreak{}}{DER\_\allowbreak{}XVIII-b:\allowbreak{}48}
\DERapparatusitem{ENV1192}{57;\allowbreak{}\allowbreak{} ms172,\allowbreak{}\allowbreak{} Proposition\allowbreak{} 4.\allowbreak{}7.\allowbreak{}7(\allowbreak{}b)\allowbreak{}}{and\allowbreak{} of\allowbreak{} cardinality\allowbreak{} \$\allowbreak{}\textbackslash{}\allowbreak{}leqslant\allowbreak{} \textbackslash{}\allowbreak{}pi\_\allowbreak{}0\$\allowbreak{}\allowbreak{} [\allowbreak{}?]\allowbreak{}.\allowbreak{}\}\allowbreak{}}{and\allowbreak{} of\allowbreak{} cardinality\allowbreak{} \$\allowbreak{}\textbackslash{}\allowbreak{}leqslant\allowbreak{} \textbackslash{}\allowbreak{}pi\$\allowbreak{}\allowbreak{} [\allowbreak{}?]\allowbreak{}.\allowbreak{}\}\allowbreak{}}{DER\_\allowbreak{}XVIII-b:\allowbreak{}49}
\DERapparatusitem{ENV1193}{58;\allowbreak{}\allowbreak{} ms172,\allowbreak{}\allowbreak{} Proposition\allowbreak{} 4.\allowbreak{}7.\allowbreak{}7(\allowbreak{}c)\allowbreak{},\allowbreak{}\allowbreak{} filtering\allowbreak{} cardinal}{and\allowbreak{} \$\allowbreak{}I\$\allowbreak{}\allowbreak{} an\allowbreak{} ordered\allowbreak{} category\allowbreak{} large\allowbreak{} relative\allowbreak{} to\allowbreak{} \$\allowbreak{}\textbackslash{}\allowbreak{}pi\_\allowbreak{}0\$\allowbreak{}.\allowbreak{}}{and\allowbreak{} \$\allowbreak{}I\$\allowbreak{}\allowbreak{} an\allowbreak{} ordered\allowbreak{} category\allowbreak{} large\allowbreak{} relative\allowbreak{} to\allowbreak{} \$\allowbreak{}\textbackslash{}\allowbreak{}pi\$\allowbreak{}.\allowbreak{}}{DER\_\allowbreak{}XVIII-b:\allowbreak{}50}
\DERapparatusitem{ENV1194}{58;\allowbreak{}\allowbreak{} ms172,\allowbreak{}\allowbreak{} Proposition\allowbreak{} 4.\allowbreak{}7.\allowbreak{}7(\allowbreak{}c)\allowbreak{},\allowbreak{}\allowbreak{} membership\allowbreak{} notation}{\textbackslash{}\allowbreak{}in\allowbreak{} \textbackslash{}\allowbreak{}Ob}{\textbackslash{}\allowbreak{}in\textbackslash{}\allowbreak{}Ob\textbackslash{}\allowbreak{}mathcal}{DER\_\allowbreak{}XVIII-b:\allowbreak{}51}
\DERapparatusitem{ENV1195}{58;\allowbreak{}\allowbreak{} ms174,\allowbreak{}\allowbreak{} proof\allowbreak{} heading}{Let\allowbreak{} us\allowbreak{} prove\allowbreak{} 4.\allowbreak{}7.\allowbreak{}2\allowbreak{} b)\allowbreak{}}{Let\allowbreak{} us\allowbreak{} prove\allowbreak{} 4.\allowbreak{}7.\allowbreak{}7\allowbreak{} b)\allowbreak{}}{DER\_\allowbreak{}XVIII-b:\allowbreak{}52}
\DERapparatusitem{ENV1196}{59;\allowbreak{}\allowbreak{} ms174,\allowbreak{}\allowbreak{} second\allowbreak{} display\allowbreak{} in\allowbreak{} proof\allowbreak{} of\allowbreak{} 4.\allowbreak{}7.\allowbreak{}7(\allowbreak{}c)\allowbreak{}}{\textbackslash{}\allowbreak{}mathcal\{\allowbreak{}F\}\allowbreak{}\textasciicircum{}\allowbreak{}\{\allowbreak{}\textbackslash{}\allowbreak{}pi'\}\allowbreak{}\allowbreak{} \textbackslash{}\allowbreak{};\allowbreak{}\textbackslash{}\allowbreak{}subset\textbackslash{}\allowbreak{};\allowbreak{}\allowbreak{} \textbackslash{}\allowbreak{}Filt\textasciicircum{}\allowbreak{}\{\allowbreak{}\textbackslash{}\allowbreak{}pi'\textasciicircum{}\allowbreak{}\{\allowbreak{}\textbackslash{}\allowbreak{}pi\}\allowbreak{}\}\allowbreak{}}{\textbackslash{}\allowbreak{}mathcal\{\allowbreak{}F\}\allowbreak{}\textasciicircum{}\allowbreak{}\{\allowbreak{}\textbackslash{}\allowbreak{}pi'\}\allowbreak{}\allowbreak{} \textbackslash{}\allowbreak{};\allowbreak{}\textbackslash{}\allowbreak{}subset\textbackslash{}\allowbreak{};\allowbreak{}\allowbreak{} \textbackslash{}\allowbreak{}Filt'\textasciicircum{}\allowbreak{}\{\allowbreak{}\textbackslash{}\allowbreak{}pi'\textasciicircum{}\allowbreak{}\{\allowbreak{}\textbackslash{}\allowbreak{}pi\}\allowbreak{}\}\allowbreak{}}{DER\_\allowbreak{}XVIII-b:\allowbreak{}53}
\DERapparatusitem{ENV1197}{59;\allowbreak{}\allowbreak{} ms175,\allowbreak{}\allowbreak{} Remark\allowbreak{} 4.\allowbreak{}7.\allowbreak{}7.\allowbreak{}2,\allowbreak{}\allowbreak{} proviso}{then\allowbreak{} the\allowbreak{} condition\allowbreak{} \$\allowbreak{}\textbackslash{}\allowbreak{}pi\textasciicircum{}\allowbreak{}\{\allowbreak{}\textbackslash{}\allowbreak{}pi'\}\allowbreak{}\textbackslash{}\allowbreak{}geqslant\textbackslash{}\allowbreak{}pi\_\allowbreak{}1\$\allowbreak{}\allowbreak{} in\allowbreak{} c)\allowbreak{}}{then\allowbreak{} the\allowbreak{} condition\allowbreak{} \$\allowbreak{}\textbackslash{}\allowbreak{}pi'\textasciicircum{}\allowbreak{}\{\allowbreak{}\textbackslash{}\allowbreak{}pi\}\allowbreak{}\textbackslash{}\allowbreak{}geqslant\textbackslash{}\allowbreak{}pi\_\allowbreak{}1\$\allowbreak{}\allowbreak{} in\allowbreak{} c)\allowbreak{}}{DER\_\allowbreak{}XVIII-b:\allowbreak{}55}
\DERapparatusitem{ENV1198}{60;\allowbreak{}\allowbreak{} ms177,\allowbreak{}\allowbreak{} proof\allowbreak{} of\allowbreak{} Corollary\allowbreak{} 4.\allowbreak{}7.\allowbreak{}10}{\textbackslash{}\allowbreak{}pi=\allowbreak{}}{\textbackslash{}\allowbreak{}pi\textbackslash{}\allowbreak{}subseteq}{DER\_\allowbreak{}XVIII-b:\allowbreak{}57}
\DERapparatusitem{ENV1199}{62;\allowbreak{}\allowbreak{} running\allowbreak{} text\allowbreak{} or\allowbreak{} formula}{\textbackslash{}\allowbreak{}text\{\allowbreak{}def\}\allowbreak{}}{\textbackslash{}\allowbreak{}mathrm\{\allowbreak{}def\}\allowbreak{}}{TYPESETTING:\allowbreak{}7;\allowbreak{}\allowbreak{} NOTATION\_\allowbreak{}DER:\allowbreak{}17}
\DERapparatusitem{ENV1200}{62;\allowbreak{}\allowbreak{} ms182,\allowbreak{}\allowbreak{} cardinal-bound\allowbreak{} display}{\textbackslash{}\allowbreak{}longrightarrow}{\textbackslash{}\allowbreak{}leqslant}{DER\_\allowbreak{}XVIII-b:\allowbreak{}59}
\DERapparatusitem{ENV1201}{63;\allowbreak{}\allowbreak{} running\allowbreak{} text\allowbreak{} or\allowbreak{} formula}{\textbackslash{}\allowbreak{}text\{\allowbreak{}def\}\allowbreak{}}{\textbackslash{}\allowbreak{}mathrm\{\allowbreak{}def\}\allowbreak{}}{TYPESETTING:\allowbreak{}7;\allowbreak{}\allowbreak{} NOTATION\_\allowbreak{}DER:\allowbreak{}17}
\DERapparatusitem{ENV1202}{63;\allowbreak{}\allowbreak{} ms185,\allowbreak{}\allowbreak{} beginning\allowbreak{} of\allowbreak{} proof\allowbreak{} of\allowbreak{} 4.\allowbreak{}7.\allowbreak{}19}{the\allowbreak{} bounds\allowbreak{} (\allowbreak{}4.\allowbreak{}7.\allowbreak{}9.\allowbreak{}1)\allowbreak{}\allowbreak{} and\allowbreak{} (\allowbreak{}4.\allowbreak{}7.\allowbreak{}9.\allowbreak{}2)\allowbreak{}}{the\allowbreak{} bounds\allowbreak{} (\allowbreak{}4.\allowbreak{}7.\allowbreak{}19.\allowbreak{}1)\allowbreak{}\allowbreak{} and\allowbreak{} (\allowbreak{}4.\allowbreak{}7.\allowbreak{}19.\allowbreak{}2)\allowbreak{}}{DER\_\allowbreak{}XVIII-b:\allowbreak{}61}
\DERapparatusitem{ENV1203}{64;\allowbreak{}\allowbreak{} ms186,\allowbreak{}\allowbreak{} first\allowbreak{} display}{\textbackslash{}\allowbreak{}longrightarrow}{\textbackslash{}\allowbreak{}leqslant}{DER\_\allowbreak{}XVIII-b:\allowbreak{}62}
\DERapparatusitem{ENV1204}{64;\allowbreak{}\allowbreak{} ms187,\allowbreak{}\allowbreak{} definition\allowbreak{} of\allowbreak{} maximal\allowbreak{} filtration}{\textbackslash{}\allowbreak{}textit\{\allowbreak{}maximal\}\allowbreak{}\allowbreak{} (\allowbreak{}where\allowbreak{} \textbackslash{}\allowbreak{}textit\{\allowbreak{}strictly\allowbreak{} canonical\}\allowbreak{})\allowbreak{}}{\textbackslash{}\allowbreak{}textit\{\allowbreak{}maximal\}\allowbreak{}\allowbreak{} (\allowbreak{}or\allowbreak{} \textbackslash{}\allowbreak{}textit\{\allowbreak{}strictly\allowbreak{} canonical\}\allowbreak{})\allowbreak{}}{DER\_\allowbreak{}XVIII-b:\allowbreak{}63}
\DERapparatusitem{ENV1205}{67;\allowbreak{}\allowbreak{} ms193,\allowbreak{}\allowbreak{} footnote82,\allowbreak{}\allowbreak{} weaker\allowbreak{} hypothesis\allowbreak{} in4.\allowbreak{}8.\allowbreak{}2(\allowbreak{}3)\allowbreak{}}{the\allowbreak{} weaker\allowbreak{} condition\allowbreak{} c)\allowbreak{}\allowbreak{} of\allowbreak{} 1°)\allowbreak{}}{the\allowbreak{} weaker\allowbreak{} condition\allowbreak{} a)\allowbreak{}\allowbreak{} of\allowbreak{} 1°)\allowbreak{}}{DER\_\allowbreak{}XVIII-b:\allowbreak{}68}
\DERapparatusitem{ENV1206}{68;\allowbreak{}\allowbreak{} ms195,\allowbreak{}\allowbreak{} proof4.\allowbreak{}8.\allowbreak{}2(\allowbreak{}2)\allowbreak{},\allowbreak{}\allowbreak{} uniqueness\allowbreak{} parenthesis}{by\allowbreak{} property\allowbreak{} c)\allowbreak{}\allowbreak{} of\allowbreak{} strict\allowbreak{} epimorphicity}{by\allowbreak{} property\allowbreak{} a)\allowbreak{}\allowbreak{} of\allowbreak{} strict\allowbreak{} epimorphicity}{DER\_\allowbreak{}XVIII-b:\allowbreak{}69}
\DERapparatusitem{ENV1207}{68;\allowbreak{}\allowbreak{} ms195,\allowbreak{}\allowbreak{} line1137,\allowbreak{}\allowbreak{} second\allowbreak{} display\allowbreak{} labelled(\allowbreak{}*)\allowbreak{}}{up\_\allowbreak{}x\allowbreak{} \&\allowbreak{}=\allowbreak{}\allowbreak{} v\_\allowbreak{}x\allowbreak{} \textbackslash{}\allowbreak{}quad}{up\_\allowbreak{}x\allowbreak{} \&\allowbreak{}=\allowbreak{}\allowbreak{} u\_\allowbreak{}x\allowbreak{} \textbackslash{}\allowbreak{}quad}{DER\_\allowbreak{}XVIII-b:\allowbreak{}70}
\DERapparatusitem{ENV1208}{68;\allowbreak{}\allowbreak{} ms195,\allowbreak{}\allowbreak{} line1138,\allowbreak{}\allowbreak{} second\allowbreak{} display\allowbreak{} labelled(\allowbreak{}*)\allowbreak{}}{up\_\allowbreak{}\textbackslash{}\allowbreak{}alpha\allowbreak{} \&\allowbreak{}=\allowbreak{}\allowbreak{} v\_\allowbreak{}\textbackslash{}\allowbreak{}alpha\allowbreak{} \textbackslash{}\allowbreak{}quad}{up\_\allowbreak{}\textbackslash{}\allowbreak{}alpha\allowbreak{} \&\allowbreak{}=\allowbreak{}\allowbreak{} u\_\allowbreak{}\textbackslash{}\allowbreak{}alpha\allowbreak{} \textbackslash{}\allowbreak{}quad}{DER\_\allowbreak{}XVIII-b:\allowbreak{}71}
\DERapparatusitem{ENV1209}{68;\allowbreak{}\allowbreak{} ms196,\allowbreak{}\allowbreak{} line1144,\allowbreak{}\allowbreak{} applying\allowbreak{} strict\allowbreak{} epimorphicity\allowbreak{} at\allowbreak{} X\_\allowbreak{}alpha}{we\allowbreak{} want\allowbreak{} that\allowbreak{} the\allowbreak{} family\allowbreak{} of\allowbreak{} arrows}{we\allowbreak{} see\allowbreak{} that\allowbreak{} the\allowbreak{} family\allowbreak{} of\allowbreak{} arrows}{DER\_\allowbreak{}XVIII-b:\allowbreak{}72}
\DERapparatusitem{ENV1210}{69;\allowbreak{}\allowbreak{} ms197,\allowbreak{}\allowbreak{} proof4.\allowbreak{}8.\allowbreak{}1,\allowbreak{}\allowbreak{} implication\allowbreak{} c=\allowbreak{}>b}{the\allowbreak{} fact\allowbreak{} that\allowbreak{} \$\allowbreak{}\textbackslash{}\allowbreak{}pi\_\allowbreak{}0\$\allowbreak{}\allowbreak{} a\allowbreak{} regular\allowbreak{} cardinal\allowbreak{} for\allowbreak{} the\allowbreak{} fcga}{the\allowbreak{} fact\allowbreak{} that\allowbreak{} \$\allowbreak{}\textbackslash{}\allowbreak{}pi\_\allowbreak{}0\$\allowbreak{}\allowbreak{} a\allowbreak{} useful\allowbreak{} cardinal\allowbreak{} for\allowbreak{} the\allowbreak{} fcga}{DER\_\allowbreak{}XVIII-b:\allowbreak{}73}
\DERapparatusitem{ENV1211}{73;\allowbreak{}\allowbreak{} ms206,\allowbreak{}\allowbreak{} continuation\allowbreak{} of\allowbreak{} Proposition4.\allowbreak{}8.\allowbreak{}8\allowbreak{} discussion}{this\allowbreak{} is\allowbreak{} so\allowbreak{} if\allowbreak{} we\allowbreak{} assume\allowbreak{} \$\allowbreak{}C\$\allowbreak{}\allowbreak{} stable\allowbreak{} under\allowbreak{} small\allowbreak{} \$\allowbreak{}\textbackslash{}\allowbreak{}varinjlim\$\allowbreak{}}{this\allowbreak{} is\allowbreak{} so\allowbreak{} if\allowbreak{} we\allowbreak{} assume\allowbreak{} \$\allowbreak{}\textbackslash{}\allowbreak{}mathcal\{\allowbreak{}M\}\allowbreak{}\$\allowbreak{}\allowbreak{} stable\allowbreak{} under\allowbreak{} small\allowbreak{} \$\allowbreak{}\textbackslash{}\allowbreak{}varinjlim\$\allowbreak{}}{DER\_\allowbreak{}XVIII-b:\allowbreak{}75}
\DERapparatusitem{ENV1212}{preamble;\allowbreak{}\allowbreak{} local\allowbreak{} macro\allowbreak{} definition}{\textbackslash{}\allowbreak{}text\{\allowbreak{}def\}\allowbreak{}}{\textbackslash{}\allowbreak{}mathrm\{\allowbreak{}def\}\allowbreak{}}{TYPESETTING:\allowbreak{}7;\allowbreak{}\allowbreak{} NOTATION\_\allowbreak{}DER:\allowbreak{}17}
\DERapparatusitem{ENV1213}{74;\allowbreak{}\allowbreak{} ms208,\allowbreak{}\allowbreak{} line1335,\allowbreak{}\allowbreak{} Lemma4.\allowbreak{}8.\allowbreak{}10\allowbreak{} hypothesis}{Suppose\allowbreak{} \$\allowbreak{}\textbackslash{}\allowbreak{}Ind\_\allowbreak{}\{\allowbreak{}\textbackslash{}\allowbreak{}pi\_\allowbreak{}0\}\allowbreak{}\$\allowbreak{}\allowbreak{} is\allowbreak{} stable}{Suppose\allowbreak{} \$\allowbreak{}\textbackslash{}\allowbreak{}Ind\_\allowbreak{}\{\allowbreak{}\textbackslash{}\allowbreak{}pi\_\allowbreak{}0\}\allowbreak{}C\$\allowbreak{}\allowbreak{} is\allowbreak{} stable}{DER\_\allowbreak{}XVIII-b:\allowbreak{}78}
\DERapparatusitem{ENV1214}{74;\allowbreak{}\allowbreak{} ms209,\allowbreak{}\allowbreak{} finite-cardinality\allowbreak{} parenthesis\allowbreak{} and(\allowbreak{}4.\allowbreak{}8.\allowbreak{}10.\allowbreak{}1)\allowbreak{}}{\textbackslash{}\allowbreak{}pi\textasciicircum{}\allowbreak{}\{\allowbreak{}c\}\allowbreak{}}{\textbackslash{}\allowbreak{}pi\textasciicircum{}\allowbreak{}\{\allowbreak{}\textbackslash{}\allowbreak{}max(\allowbreak{}1,\allowbreak{}c)\allowbreak{}\}\allowbreak{}}{DER\_\allowbreak{}XVIII-b:\allowbreak{}81}
\DERapparatusitem{ENV1215}{74;\allowbreak{}\allowbreak{} ms209,\allowbreak{}\allowbreak{} line1342,\allowbreak{}\allowbreak{} definition\allowbreak{} of\allowbreak{} X}{X\allowbreak{} \textbackslash{}\allowbreak{}\allowbreak{} \textbackslash{}\allowbreak{}defeq\textbackslash{}\allowbreak{}\allowbreak{} \textbackslash{}\allowbreak{}varinjlim\_\allowbreak{}\{\allowbreak{}j\textbackslash{}\allowbreak{}in\allowbreak{} J\}\allowbreak{}\allowbreak{} x\_\allowbreak{}j}{X\allowbreak{} \textbackslash{}\allowbreak{}\allowbreak{} \textbackslash{}\allowbreak{}defeq\textbackslash{}\allowbreak{}\allowbreak{} \textbackslash{}\allowbreak{}varprojlim\_\allowbreak{}\{\allowbreak{}j\textbackslash{}\allowbreak{}in\allowbreak{} J\}\allowbreak{}\allowbreak{} x\_\allowbreak{}j}{DER\_\allowbreak{}XVIII-b:\allowbreak{}79}
\DERapparatusitem{ENV1216}{74;\allowbreak{}\allowbreak{} ms209,\allowbreak{}\allowbreak{} line1342,\allowbreak{}\allowbreak{} comparison\allowbreak{} with\allowbreak{} product}{\textbackslash{}\allowbreak{}\allowbreak{} =\allowbreak{}\allowbreak{} \textbackslash{}\allowbreak{}\allowbreak{} \textbackslash{}\allowbreak{}underbrace\{\allowbreak{}\textbackslash{}\allowbreak{}prod\_\allowbreak{}\{\allowbreak{}j\textbackslash{}\allowbreak{}in\textbackslash{}\allowbreak{}Ob\allowbreak{} J\}\allowbreak{}\allowbreak{} x\_\allowbreak{}j\}\allowbreak{}}{\textbackslash{}\allowbreak{}\allowbreak{} \textbackslash{}\allowbreak{}hookrightarrow\textbackslash{}\allowbreak{}\allowbreak{} \textbackslash{}\allowbreak{}underbrace\{\allowbreak{}\textbackslash{}\allowbreak{}prod\_\allowbreak{}\{\allowbreak{}j\textbackslash{}\allowbreak{}in\textbackslash{}\allowbreak{}Ob\allowbreak{} J\}\allowbreak{}\allowbreak{} x\_\allowbreak{}j\}\allowbreak{}}{DER\_\allowbreak{}XVIII-b:\allowbreak{}80}
\DERapparatusitem{ENV1217}{74;\allowbreak{}\allowbreak{} ms209,\allowbreak{}\allowbreak{} finite-cardinality\allowbreak{} parenthesis\allowbreak{} and(\allowbreak{}4.\allowbreak{}8.\allowbreak{}10.\allowbreak{}1)\allowbreak{}}{\textbackslash{}\allowbreak{}pi\textasciicircum{}\allowbreak{}\{\allowbreak{}c\}\allowbreak{}}{\textbackslash{}\allowbreak{}pi\textasciicircum{}\allowbreak{}\{\allowbreak{}\textbackslash{}\allowbreak{}max(\allowbreak{}1,\allowbreak{}c)\allowbreak{}\}\allowbreak{}}{DER\_\allowbreak{}XVIII-b:\allowbreak{}81}
\DERapparatusitem{ENV1218}{74;\allowbreak{}\allowbreak{} ms209,\allowbreak{}\allowbreak{} finite-cardinality\allowbreak{} parenthesis\allowbreak{} and(\allowbreak{}4.\allowbreak{}8.\allowbreak{}10.\allowbreak{}1)\allowbreak{}}{\textbackslash{}\allowbreak{}pi\textasciicircum{}\allowbreak{}\{\allowbreak{}c\}\allowbreak{}}{\textbackslash{}\allowbreak{}pi\textasciicircum{}\allowbreak{}\{\allowbreak{}\textbackslash{}\allowbreak{}max(\allowbreak{}1,\allowbreak{}c)\allowbreak{}\}\allowbreak{}}{DER\_\allowbreak{}XVIII-b:\allowbreak{}81}
\DERapparatusitem{ENV1219}{74;\allowbreak{}\allowbreak{} ms209,\allowbreak{}\allowbreak{} finite-cardinality\allowbreak{} parenthesis\allowbreak{} and(\allowbreak{}4.\allowbreak{}8.\allowbreak{}10.\allowbreak{}1)\allowbreak{}}{\textbackslash{}\allowbreak{}pi\textasciicircum{}\allowbreak{}\{\allowbreak{}c\}\allowbreak{}}{\textbackslash{}\allowbreak{}pi\textasciicircum{}\allowbreak{}\{\allowbreak{}\textbackslash{}\allowbreak{}max(\allowbreak{}1,\allowbreak{}c)\allowbreak{}\}\allowbreak{}}{DER\_\allowbreak{}XVIII-b:\allowbreak{}81}
\DERapparatusitem{ENV1220}{74;\allowbreak{}\allowbreak{} ms210,\allowbreak{}\allowbreak{} Lemma4.\allowbreak{}8.\allowbreak{}11\allowbreak{} first\allowbreak{} sentence}{For\allowbreak{} \$\allowbreak{}\textbackslash{}\allowbreak{}xi\$\allowbreak{}\allowbreak{} to\allowbreak{} be\allowbreak{} representable,\allowbreak{}\allowbreak{} it\allowbreak{} is\allowbreak{} necessary\allowbreak{} and\allowbreak{} sufficient\allowbreak{} that\allowbreak{} it\allowbreak{} be\allowbreak{} left\allowbreak{} exact\allowbreak{} and\allowbreak{} accessible.\allowbreak{}}{For\allowbreak{} \$\allowbreak{}\textbackslash{}\allowbreak{}xi\$\allowbreak{}\allowbreak{} to\allowbreak{} be\allowbreak{} prorepresentable,\allowbreak{}\allowbreak{} it\allowbreak{} is\allowbreak{} necessary\allowbreak{} and\allowbreak{} sufficient\allowbreak{} that\allowbreak{} it\allowbreak{} be\allowbreak{} left\allowbreak{} exact\allowbreak{} and\allowbreak{} accessible.\allowbreak{}}{DER\_\allowbreak{}XVIII-b:\allowbreak{}83}
\DERapparatusitem{ENV1221}{75;\allowbreak{}\allowbreak{} ms211,\allowbreak{}\allowbreak{} line1378,\allowbreak{}\allowbreak{} Lemma4.\allowbreak{}8.\allowbreak{}12\allowbreak{} initial\allowbreak{} object-name}{\_\allowbreak{}0}{}{DER\_\allowbreak{}XVIII-b:\allowbreak{}84}
\DERapparatusitem{ENV1222}{76;\allowbreak{}\allowbreak{} ms214,\allowbreak{}4.\allowbreak{}8.\allowbreak{}15,\allowbreak{}\allowbreak{} criterion\allowbreak{} for\allowbreak{} Filt\allowbreak{} membership}{by\allowbreak{} 4.\allowbreak{}7.\allowbreak{}7,\allowbreak{}}{by\allowbreak{} 4.\allowbreak{}7.\allowbreak{}1,\allowbreak{}}{DER\_\allowbreak{}XVIII-b:\allowbreak{}87}
\DERapparatusitem{ENV1223}{76;\allowbreak{}\allowbreak{} ms214,\allowbreak{}4.\allowbreak{}8.\allowbreak{}15,\allowbreak{}\allowbreak{} auxiliary\allowbreak{} cardinal}{J}{1}{DER\_\allowbreak{}XVIII-b:\allowbreak{}88}
\DERapparatusitem{ENV1224}{76;\allowbreak{}\allowbreak{} ms214,\allowbreak{}4.\allowbreak{}8.\allowbreak{}15,\allowbreak{}\allowbreak{} reduction\allowbreak{} to\allowbreak{} a\allowbreak{} skeleton}{=\allowbreak{}}{\textbackslash{}\allowbreak{}leqslant}{DER\_\allowbreak{}XVIII-b:\allowbreak{}89}
\DERapparatusitem{ENV1225}{76;\allowbreak{}\allowbreak{} ms215,\allowbreak{}4.\allowbreak{}8.\allowbreak{}15,\allowbreak{}\allowbreak{} cofinal\allowbreak{} fibre}{the\allowbreak{} set\allowbreak{} of\allowbreak{} the\allowbreak{} \$\allowbreak{}I\_\allowbreak{}\{\allowbreak{}\textbackslash{}\allowbreak{}xi\}\allowbreak{}\$\allowbreak{}\allowbreak{} is\allowbreak{} cofinal\allowbreak{} in\allowbreak{} \$\allowbreak{}I\$\allowbreak{}}{one\allowbreak{} of\allowbreak{} the\allowbreak{} \$\allowbreak{}I\_\allowbreak{}\{\allowbreak{}\textbackslash{}\allowbreak{}xi\}\allowbreak{}\$\allowbreak{}\allowbreak{} is\allowbreak{} cofinal\allowbreak{} in\allowbreak{} \$\allowbreak{}I\$\allowbreak{}}{DER\_\allowbreak{}XVIII-b:\allowbreak{}90}
\DERapparatusitem{ENV1226}{76;\allowbreak{}\allowbreak{} ms215,\allowbreak{}4.\allowbreak{}8.\allowbreak{}15,\allowbreak{}\allowbreak{} size\allowbreak{} of\allowbreak{} endomorphism\allowbreak{} monoid}{\textbackslash{}\allowbreak{}leqslant\textbackslash{}\allowbreak{}\allowbreak{} \textbackslash{}\allowbreak{}pi}{\textbackslash{}\allowbreak{}leqslant\textbackslash{}\allowbreak{}pi\_\allowbreak{}0}{DER\_\allowbreak{}XVIII-b:\allowbreak{}91}
\DERapparatusitem{ENV1227}{77;\allowbreak{}\allowbreak{} ms216,\allowbreak{}\allowbreak{} line1478,\allowbreak{}\allowbreak{} left-cancellation\allowbreak{} condition}{\textbackslash{}\allowbreak{}cM}{M}{DER\_\allowbreak{}XVIII-b:\allowbreak{}92}
\DERapparatusitem{ENV1228}{78;\allowbreak{}\allowbreak{} ms218,\allowbreak{}\allowbreak{} indexed\allowbreak{} representation\allowbreak{} paragraph}{the\allowbreak{} essential\allowbreak{} image\allowbreak{} of\allowbreak{} (\allowbreak{}4.\allowbreak{}4.\allowbreak{}3)\allowbreak{}}{the\allowbreak{} essential\allowbreak{} image\allowbreak{} of\allowbreak{} (\allowbreak{}4.\allowbreak{}9.\allowbreak{}3)\allowbreak{}}{DER\_\allowbreak{}XVIII-b:\allowbreak{}93}
\DERapparatusitem{ENV1229}{79;\allowbreak{}\allowbreak{} ms219,\allowbreak{}\allowbreak{} Theorem4.\allowbreak{}9.\allowbreak{}3(\allowbreak{}b)\allowbreak{},\allowbreak{}\allowbreak{} roles\allowbreak{} of\allowbreak{} the\allowbreak{} two\allowbreak{} cardinals}{general\allowbreak{} conditions\allowbreak{} of\allowbreak{} 4.\allowbreak{}9.\allowbreak{}1\allowbreak{} for\allowbreak{} \$\allowbreak{}C\$\allowbreak{},\allowbreak{}\allowbreak{} \$\allowbreak{}\textbackslash{}\allowbreak{}pi\_\allowbreak{}0\$\allowbreak{},\allowbreak{}\allowbreak{} \$\allowbreak{}J\$\allowbreak{}\allowbreak{} (\allowbreak{}in\allowbreak{} particular\}\allowbreak{}\allowbreak{} \$\allowbreak{}\textbackslash{}\allowbreak{}cardred\allowbreak{} J\allowbreak{} \textbackslash{}\allowbreak{}leqslant\allowbreak{} \textbackslash{}\allowbreak{}pi\_\allowbreak{}0\$\allowbreak{}}{general\allowbreak{} conditions\allowbreak{} of\allowbreak{} 4.\allowbreak{}9.\allowbreak{}1\allowbreak{} for\allowbreak{} \$\allowbreak{}C\$\allowbreak{},\allowbreak{}\allowbreak{} \$\allowbreak{}\textbackslash{}\allowbreak{}pi\$\allowbreak{},\allowbreak{}\allowbreak{} \$\allowbreak{}J\$\allowbreak{}\allowbreak{} (\allowbreak{}in\allowbreak{} particular\}\allowbreak{}\allowbreak{} \$\allowbreak{}\textbackslash{}\allowbreak{}cardred\allowbreak{} J\allowbreak{} \textbackslash{}\allowbreak{}leqslant\allowbreak{} \textbackslash{}\allowbreak{}pi\$\allowbreak{}}{DER\_\allowbreak{}XVIII-b:\allowbreak{}95}
\DERapparatusitem{ENV1230}{79;\allowbreak{}\allowbreak{} ms219,\allowbreak{}\allowbreak{} Theorem4.\allowbreak{}9.\allowbreak{}3(\allowbreak{}b)\allowbreak{},\allowbreak{}\allowbreak{} roles\allowbreak{} of\allowbreak{} the\allowbreak{} two\allowbreak{} cardinals}{\textbackslash{}\allowbreak{}item[\allowbreak{}b)\allowbreak{}]\allowbreak{}\allowbreak{} There\allowbreak{} exists\allowbreak{} a\allowbreak{} cardinal\allowbreak{} \$\allowbreak{}\textbackslash{}\allowbreak{}pi\_\allowbreak{}0\allowbreak{} <\allowbreak{} \textbackslash{}\allowbreak{}pi\$\allowbreak{},\allowbreak{}\allowbreak{} and\allowbreak{} a\allowbreak{} cardinal\allowbreak{} \$\allowbreak{}\textbackslash{}\allowbreak{}pi'\_\allowbreak{}0\allowbreak{} >\allowbreak{} \textbackslash{}\allowbreak{}pi\_\allowbreak{}0\$\allowbreak{}\allowbreak{} \textbackslash{}\allowbreak{}textup\{\allowbreak{}(\allowbreak{}\}\allowbreak{}\textbackslash{}\allowbreak{}textbf\{\allowbreak{}\textbackslash{}\allowbreak{}upshape\allowbreak{} N.\allowbreak{}B.\allowbreak{}\}\allowbreak{}\allowbreak{} We\allowbreak{} may\allowbreak{} suppose\allowbreak{} that\allowbreak{} \$\allowbreak{}\textbackslash{}\allowbreak{}pi'\_\allowbreak{}0\$\allowbreak{}\allowbreak{} is\allowbreak{} the\allowbreak{} successor\allowbreak{} of\allowbreak{} \$\allowbreak{}\textbackslash{}\allowbreak{}pi\_\allowbreak{}0\$\allowbreak{},\allowbreak{}\allowbreak{} hence\allowbreak{} \$\allowbreak{}\textbackslash{}\allowbreak{}pi'\_\allowbreak{}0\allowbreak{} \textbackslash{}\allowbreak{}leqslant\allowbreak{} \textbackslash{}\allowbreak{}pi\$\allowbreak{})\allowbreak{},\allowbreak{}\allowbreak{} such\allowbreak{} that\allowbreak{} \$\allowbreak{}C\$\allowbreak{}\allowbreak{} is\allowbreak{} stable\allowbreak{} under\allowbreak{} direct\allowbreak{} limits\allowbreak{} large\allowbreak{} relative\allowbreak{} to\allowbreak{} \$\allowbreak{}\textbackslash{}\allowbreak{}pi\_\allowbreak{}0\$\allowbreak{},\allowbreak{}\allowbreak{} and\allowbreak{} of\allowbreak{} cardinality\allowbreak{} \$\allowbreak{}\textbackslash{}\allowbreak{}pi'\_\allowbreak{}0\$\allowbreak{}.\allowbreak{}}{\textbackslash{}\allowbreak{}item[\allowbreak{}b)\allowbreak{}]\allowbreak{}\allowbreak{} There\allowbreak{} exists\allowbreak{} a\allowbreak{} cardinal\allowbreak{} \$\allowbreak{}\textbackslash{}\allowbreak{}mu\allowbreak{} <\allowbreak{} \textbackslash{}\allowbreak{}pi\$\allowbreak{},\allowbreak{}\allowbreak{} and\allowbreak{} a\allowbreak{} cardinal\allowbreak{} \$\allowbreak{}\textbackslash{}\allowbreak{}mu'\allowbreak{} >\allowbreak{} \textbackslash{}\allowbreak{}mu\$\allowbreak{}\allowbreak{} \textbackslash{}\allowbreak{}textup\{\allowbreak{}(\allowbreak{}\}\allowbreak{}\textbackslash{}\allowbreak{}textbf\{\allowbreak{}\textbackslash{}\allowbreak{}upshape\allowbreak{} N.\allowbreak{}B.\allowbreak{}\}\allowbreak{}\allowbreak{} We\allowbreak{} may\allowbreak{} suppose\allowbreak{} that\allowbreak{} \$\allowbreak{}\textbackslash{}\allowbreak{}mu'\$\allowbreak{}\allowbreak{} is\allowbreak{} the\allowbreak{} successor\allowbreak{} of\allowbreak{} \$\allowbreak{}\textbackslash{}\allowbreak{}mu\$\allowbreak{},\allowbreak{}\allowbreak{} hence\allowbreak{} \$\allowbreak{}\textbackslash{}\allowbreak{}mu'\allowbreak{} \textbackslash{}\allowbreak{}leqslant\allowbreak{} \textbackslash{}\allowbreak{}pi\$\allowbreak{})\allowbreak{},\allowbreak{}\allowbreak{} such\allowbreak{} that\allowbreak{} \$\allowbreak{}C\$\allowbreak{}\allowbreak{} is\allowbreak{} stable\allowbreak{} under\allowbreak{} direct\allowbreak{} limits\allowbreak{} large\allowbreak{} relative\allowbreak{} to\allowbreak{} \$\allowbreak{}\textbackslash{}\allowbreak{}mu\$\allowbreak{},\allowbreak{}\allowbreak{} and\allowbreak{} of\allowbreak{} cardinality\allowbreak{} \$\allowbreak{}\textbackslash{}\allowbreak{}mu'\$\allowbreak{}.\allowbreak{}}{DER\_\allowbreak{}XVIII-b:\allowbreak{}95}
\DERapparatusitem{ENV1231}{79;\allowbreak{}\allowbreak{} ms219,\allowbreak{}\allowbreak{} Theorem4.\allowbreak{}9.\allowbreak{}3(\allowbreak{}b)\allowbreak{},\allowbreak{}\allowbreak{} roles\allowbreak{} of\allowbreak{} the\allowbreak{} two\allowbreak{} cardinals}{1)\allowbreak{}\allowbreak{} \$\allowbreak{}\textbackslash{}\allowbreak{}card\allowbreak{} \textbackslash{}\allowbreak{}alpha\allowbreak{} \textbackslash{}\allowbreak{}leqslant\allowbreak{} \textbackslash{}\allowbreak{}pi\_\allowbreak{}0\$\allowbreak{}.\allowbreak{}\allowbreak{} 2)\allowbreak{}\allowbreak{} The\allowbreak{} direct\allowbreak{} limits}{1)\allowbreak{}\allowbreak{} \$\allowbreak{}\textbackslash{}\allowbreak{}card\allowbreak{} \textbackslash{}\allowbreak{}alpha\allowbreak{} \textbackslash{}\allowbreak{}leqslant\allowbreak{} \textbackslash{}\allowbreak{}pi\$\allowbreak{}.\allowbreak{}\allowbreak{} 2)\allowbreak{}\allowbreak{} The\allowbreak{} direct\allowbreak{} limits}{DER\_\allowbreak{}XVIII-b:\allowbreak{}95}
\DERapparatusitem{ENV1232}{80;\allowbreak{}\allowbreak{} ms222,\allowbreak{}\allowbreak{} proof4.\allowbreak{}9.\allowbreak{}6.\allowbreak{}1\allowbreak{} existing\allowbreak{} bracketed\allowbreak{} editorial\allowbreak{} correction}{'\allowbreak{} \textbackslash{}\allowbreak{}in\allowbreak{} \textbackslash{}\allowbreak{}Ob\allowbreak{} I'}{\allowbreak{} \textbackslash{}\allowbreak{}in\allowbreak{} \textbackslash{}\allowbreak{}Ob\allowbreak{} I}{DER\_\allowbreak{}XVIII-c:\allowbreak{}5}
\DERapparatusitem{ENV1233}{80;\allowbreak{}\allowbreak{} running\allowbreak{} text\allowbreak{} or\allowbreak{} formula}{\textbackslash{}\allowbreak{}text\{\allowbreak{}def\}\allowbreak{}}{\textbackslash{}\allowbreak{}mathrm\{\allowbreak{}def\}\allowbreak{}}{TYPESETTING:\allowbreak{}7;\allowbreak{}\allowbreak{} NOTATION\_\allowbreak{}DER:\allowbreak{}17}
\DERapparatusitem{ENV1234}{81;\allowbreak{}\allowbreak{} ms223,\allowbreak{}\allowbreak{} line1639,\allowbreak{}\allowbreak{} evaluation\allowbreak{} functor\allowbreak{} domain}{\textbackslash{}\allowbreak{}Phi}{\textbackslash{}\allowbreak{}varphi}{DER\_\allowbreak{}XVIII-b:\allowbreak{}100;\allowbreak{}\allowbreak{} DER\_\allowbreak{}XVIII-c:\allowbreak{}6}
\DERapparatusitem{ENV1235}{81;\allowbreak{}\allowbreak{} ms224,\allowbreak{}\allowbreak{} Lemma4.\allowbreak{}9.\allowbreak{}6.\allowbreak{}3(\allowbreak{}b-prime)\allowbreak{}}{\textbackslash{}\allowbreak{}textup\{\allowbreak{})\allowbreak{}\}\allowbreak{}\allowbreak{} (\allowbreak{}equalization\allowbreak{} of\allowbreak{} double\allowbreak{} arrows)\allowbreak{}.\allowbreak{}}{\textbackslash{}\allowbreak{}textup\{\allowbreak{})\allowbreak{}\}\allowbreak{}\allowbreak{} (\allowbreak{}equalization\allowbreak{} of\allowbreak{} double\allowbreak{} arrows)\allowbreak{},\allowbreak{}\allowbreak{} and\allowbreak{} the\allowbreak{} categories\allowbreak{} \$\allowbreak{}i'\textbackslash{}\allowbreak{}backslash\allowbreak{} I\$\allowbreak{}\allowbreak{} are\allowbreak{} filtering\allowbreak{} for\allowbreak{} every\allowbreak{} \$\allowbreak{}i'\textbackslash{}\allowbreak{}in\textbackslash{}\allowbreak{}Ob\allowbreak{} I'\$\allowbreak{}.\allowbreak{}}{DER\_\allowbreak{}XVIII-b:\allowbreak{}101;\allowbreak{}\allowbreak{} DER\_\allowbreak{}XVIII-c:\allowbreak{}8}
\DERapparatusitem{ENV1236}{82;\allowbreak{}\allowbreak{} ms225,\allowbreak{}\allowbreak{} line1690,\allowbreak{}\allowbreak{} b\_\allowbreak{}pi\allowbreak{} double-prime}{satisfy\allowbreak{} PS\allowbreak{} \$\allowbreak{}1\_\allowbreak{}\textbackslash{}\allowbreak{}pi\{\allowbreak{}\}\allowbreak{}'\$\allowbreak{}}{satisfy\allowbreak{} PF\allowbreak{} \$\allowbreak{}1\_\allowbreak{}\textbackslash{}\allowbreak{}pi\{\allowbreak{}\}\allowbreak{}'\$\allowbreak{}}{DER\_\allowbreak{}XVIII-b:\allowbreak{}105;\allowbreak{}\allowbreak{} DER\_\allowbreak{}XVIII-c:\allowbreak{}11}
\DERapparatusitem{ENV1237}{82;\allowbreak{}\allowbreak{} ms225,\allowbreak{}\allowbreak{} footnote102,\allowbreak{}\allowbreak{} proof\allowbreak{} locator}{4.\allowbreak{}9.\allowbreak{}6.\allowbreak{}7,\allowbreak{}\allowbreak{} page\allowbreak{} 227.\allowbreak{}}{4.\allowbreak{}9.\allowbreak{}6.\allowbreak{}7,\allowbreak{}\allowbreak{} page\allowbreak{} 228.\allowbreak{}}{DER\_\allowbreak{}XVIII-b:\allowbreak{}107}
\DERapparatusitem{ENV1238}{82;\allowbreak{}\allowbreak{} ms225,\allowbreak{}4.\allowbreak{}9.\allowbreak{}6.\allowbreak{}4,\allowbreak{}\allowbreak{} condition\allowbreak{} to\allowbreak{} be\allowbreak{} proved}{We\allowbreak{} shall\allowbreak{} prove\allowbreak{} it\allowbreak{} by\allowbreak{} taking\allowbreak{} \$\allowbreak{}\textbackslash{}\allowbreak{}mathrm\{\allowbreak{}b\}\allowbreak{}\_\allowbreak{}\textbackslash{}\allowbreak{}pi\$\allowbreak{}}{We\allowbreak{} shall\allowbreak{} prove\allowbreak{} it\allowbreak{} by\allowbreak{} taking\allowbreak{} \$\allowbreak{}\textbackslash{}\allowbreak{}mathrm\{\allowbreak{}c\}\allowbreak{}\_\allowbreak{}\textbackslash{}\allowbreak{}pi\$\allowbreak{}}{DER\_\allowbreak{}XVIII-b:\allowbreak{}108;\allowbreak{}\allowbreak{} DER\_\allowbreak{}XVIII-c:\allowbreak{}14}
\DERapparatusitem{ENV1239}{82;\allowbreak{}\allowbreak{} ms226,\allowbreak{}\allowbreak{} line1726,\allowbreak{}\allowbreak{} F3\_\allowbreak{}pi\allowbreak{} source\allowbreak{} diagrams}{C,\allowbreak{}J}{J,\allowbreak{}C}{DER\_\allowbreak{}XVIII-b:\allowbreak{}110;\allowbreak{}\allowbreak{} DER\_\allowbreak{}XVIII-c:\allowbreak{}13}
\DERapparatusitem{ENV1240}{83;\allowbreak{}\allowbreak{} ms226-227,\allowbreak{}\allowbreak{} ambient\allowbreak{} category\allowbreak{} of\allowbreak{} H->phi}{there\allowbreak{} exists\allowbreak{} \$\allowbreak{}H\allowbreak{} \textbackslash{}\allowbreak{}xrightarrow\{\allowbreak{}q\}\allowbreak{}\allowbreak{} \textbackslash{}\allowbreak{}varphi\$\allowbreak{}\allowbreak{} in\allowbreak{} \$\allowbreak{}\textbackslash{}\allowbreak{}uHom(\allowbreak{}J,\allowbreak{}C)\allowbreak{}\$\allowbreak{}}{there\allowbreak{} exists\allowbreak{} \$\allowbreak{}H\$\allowbreak{}\allowbreak{} in\allowbreak{} \$\allowbreak{}\textbackslash{}\allowbreak{}uHom(\allowbreak{}J,\allowbreak{}C)\allowbreak{}\$\allowbreak{}\allowbreak{} and\allowbreak{} \$\allowbreak{}H\allowbreak{} \textbackslash{}\allowbreak{}xrightarrow\{\allowbreak{}q\}\allowbreak{}\allowbreak{} \textbackslash{}\allowbreak{}varphi\$\allowbreak{}\allowbreak{} in\allowbreak{} \$\allowbreak{}\textbackslash{}\allowbreak{}uHom(\allowbreak{}J,\allowbreak{}\textbackslash{}\allowbreak{}cM)\allowbreak{}\$\allowbreak{}}{DER\_\allowbreak{}XVIII-c:\allowbreak{}19}
\DERapparatusitem{ENV1241}{83;\allowbreak{}\allowbreak{} ms227,\allowbreak{}\allowbreak{} line1748,\allowbreak{}\allowbreak{} indexing\allowbreak{} of\allowbreak{} components}{\textbackslash{}\allowbreak{}in}{=\allowbreak{}}{DER\_\allowbreak{}XVIII-b:\allowbreak{}111;\allowbreak{}\allowbreak{} DER\_\allowbreak{}XVIII-c:\allowbreak{}16}
\DERapparatusitem{ENV1242}{83;\allowbreak{}\allowbreak{} ms227,\allowbreak{}\allowbreak{} line1748,\allowbreak{}\allowbreak{} conclusion\allowbreak{} for\allowbreak{} w\_\allowbreak{}alpha}{comes\allowbreak{} from\allowbreak{} a\allowbreak{} \textbackslash{}\allowbreak{}textit\{\allowbreak{}functor\}\allowbreak{}\allowbreak{} \$\allowbreak{}w\_\allowbreak{}\textbackslash{}\allowbreak{}alpha}{comes\allowbreak{} from\allowbreak{} a\allowbreak{} \textbackslash{}\allowbreak{}textit\{\allowbreak{}natural\allowbreak{} transformation\}\allowbreak{}\allowbreak{} \$\allowbreak{}w\_\allowbreak{}\textbackslash{}\allowbreak{}alpha}{DER\_\allowbreak{}XVIII-b:\allowbreak{}112;\allowbreak{}\allowbreak{} DER\_\allowbreak{}XVIII-c:\allowbreak{}20}
\DERapparatusitem{ENV1243}{83;\allowbreak{}\allowbreak{} ms227,\allowbreak{}\allowbreak{} proof4.\allowbreak{}9.\allowbreak{}6.\allowbreak{}5,\allowbreak{}\allowbreak{} filtered\allowbreak{} comma\allowbreak{} categories}{by\allowbreak{} 4.\allowbreak{}9.\allowbreak{}6.\allowbreak{}3,\allowbreak{}\allowbreak{} 1°)\allowbreak{}\allowbreak{} a\allowbreak{} \$\allowbreak{}\textbackslash{}\allowbreak{}Longrightarrow\$\allowbreak{}\allowbreak{} c,\allowbreak{}}{by\allowbreak{} 4.\allowbreak{}9.\allowbreak{}6.\allowbreak{}3,\allowbreak{}\allowbreak{} 1°)\allowbreak{}\allowbreak{} a\allowbreak{} \$\allowbreak{}\textbackslash{}\allowbreak{}Longrightarrow\$\allowbreak{}\allowbreak{} b,\allowbreak{}}{DER\_\allowbreak{}XVIII-b:\allowbreak{}113;\allowbreak{}\allowbreak{} DER\_\allowbreak{}XVIII-c:\allowbreak{}21}
\DERapparatusitem{ENV1244}{83;\allowbreak{}\allowbreak{} ms228,\allowbreak{}\allowbreak{} line1760,\allowbreak{}\allowbreak{} product\allowbreak{} lemma\allowbreak{} reference}{(\allowbreak{}cf.\allowbreak{}\allowbreak{} 4.\allowbreak{}4.\allowbreak{}6.\allowbreak{}2)\allowbreak{}}{(\allowbreak{}cf.\allowbreak{}\allowbreak{} 4.\allowbreak{}9.\allowbreak{}6.\allowbreak{}2)\allowbreak{}}{DER\_\allowbreak{}XVIII-b:\allowbreak{}114;\allowbreak{}\allowbreak{} DER\_\allowbreak{}XVIII-c:\allowbreak{}17}
\DERapparatusitem{ENV1245}{83;\allowbreak{}\allowbreak{} ms228,\allowbreak{}\allowbreak{} line1762,\allowbreak{}\allowbreak{} corollary\allowbreak{} heading}{\textbackslash{}\allowbreak{}textbf\{\allowbreak{}Corollary\allowbreak{} 4.\allowbreak{}4.\allowbreak{}6.\allowbreak{}6.\allowbreak{}\}\allowbreak{}}{\textbackslash{}\allowbreak{}textbf\{\allowbreak{}Corollary\allowbreak{} 4.\allowbreak{}9.\allowbreak{}6.\allowbreak{}6.\allowbreak{}\}\allowbreak{}}{DER\_\allowbreak{}XVIII-b:\allowbreak{}115;\allowbreak{}\allowbreak{} DER\_\allowbreak{}XVIII-c:\allowbreak{}18}
\DERapparatusitem{ENV1246}{84;\allowbreak{}\allowbreak{} ms\allowbreak{} 229;\allowbreak{}\allowbreak{} existing\allowbreak{} [\allowbreak{}sic]\allowbreak{}\allowbreak{} flag}{'}{}{DER\_\allowbreak{}XVIII-c:\allowbreak{}22}
\DERapparatusitem{ENV1247}{85;\allowbreak{}\allowbreak{} ms230-231,\allowbreak{}\allowbreak{} parallel-pair\allowbreak{} names}{i\_\allowbreak{}0\allowbreak{} \textbackslash{}\allowbreak{}\allowbreak{} \textbackslash{}\allowbreak{}overset\{\allowbreak{}v\_\allowbreak{}1\}\allowbreak{}\{\allowbreak{}\textbackslash{}\allowbreak{}underset\{\allowbreak{}v\_\allowbreak{}2\}\allowbreak{}\{\allowbreak{}\textbackslash{}\allowbreak{}rightrightarrows\}\allowbreak{}\}\allowbreak{}\allowbreak{} \textbackslash{}\allowbreak{}\allowbreak{} i\_\allowbreak{}1}{i\_\allowbreak{}0\allowbreak{} \textbackslash{}\allowbreak{}\allowbreak{} \textbackslash{}\allowbreak{}overset\{\allowbreak{}u\_\allowbreak{}1\}\allowbreak{}\{\allowbreak{}\textbackslash{}\allowbreak{}underset\{\allowbreak{}u\_\allowbreak{}2\}\allowbreak{}\{\allowbreak{}\textbackslash{}\allowbreak{}rightrightarrows\}\allowbreak{}\}\allowbreak{}\allowbreak{} \textbackslash{}\allowbreak{}\allowbreak{} i\_\allowbreak{}1}{DER\_\allowbreak{}XVIII-c:\allowbreak{}25}
\DERapparatusitem{ENV1248}{85;\allowbreak{}\allowbreak{} ms231,\allowbreak{}\allowbreak{} lifting\allowbreak{} equation}{\$\allowbreak{}f(\allowbreak{}w)\allowbreak{}v'\allowbreak{} =\allowbreak{}\allowbreak{} f(\allowbreak{}v)\allowbreak{}\$\allowbreak{}}{\$\allowbreak{}f(\allowbreak{}v)\allowbreak{}v'\allowbreak{} =\allowbreak{}\allowbreak{} f(\allowbreak{}w)\allowbreak{}\$\allowbreak{}}{DER\_\allowbreak{}XVIII-c:\allowbreak{}26}
\DERapparatusitem{ENV1249}{85;\allowbreak{}\allowbreak{} ms\allowbreak{} 231;\allowbreak{}\allowbreak{} existing\allowbreak{} [\allowbreak{}sic]\allowbreak{}\allowbreak{} flag}{\$\allowbreak{}\textbackslash{}\allowbreak{}xi\_\allowbreak{}1\allowbreak{} =\allowbreak{}\allowbreak{} (\allowbreak{}i\_\allowbreak{}0,\allowbreak{}\allowbreak{} v')\allowbreak{}\$\allowbreak{}}{\$\allowbreak{}\textbackslash{}\allowbreak{}xi\_\allowbreak{}1\allowbreak{} =\allowbreak{}\allowbreak{} (\allowbreak{}i\_\allowbreak{}1,\allowbreak{}\allowbreak{} v')\allowbreak{}\$\allowbreak{}}{DER\_\allowbreak{}XVIII-c:\allowbreak{}23}
\DERapparatusitem{ENV1250}{86;\allowbreak{}\allowbreak{} ms232,\allowbreak{}\allowbreak{} identity\allowbreak{} in\allowbreak{} lifting\allowbreak{} equation}{i\_\allowbreak{}0}{f(\allowbreak{}i\_\allowbreak{}0)\allowbreak{}}{DER\_\allowbreak{}XVIII-c:\allowbreak{}30}
\DERapparatusitem{ENV1251}{86;\allowbreak{}\allowbreak{} ms\allowbreak{} 233;\allowbreak{}\allowbreak{} existing\allowbreak{} [\allowbreak{}sic]\allowbreak{}\allowbreak{} flag}{=\allowbreak{}\allowbreak{} f(\allowbreak{}i\_\allowbreak{}1}{\textbackslash{}\allowbreak{}rightrightarrows\allowbreak{} f(\allowbreak{}i\_\allowbreak{}2}{DER\_\allowbreak{}XVIII-c:\allowbreak{}29}
\DERapparatusitem{ENV1252}{87;\allowbreak{}\allowbreak{} ms235,\allowbreak{}\allowbreak{} indexing\allowbreak{} of\allowbreak{} commas}{}{'}{DER\_\allowbreak{}XVIII-c:\allowbreak{}31}
\DERapparatusitem{ENV1253}{88;\allowbreak{}\allowbreak{} ms235,\allowbreak{}\allowbreak{} printed\allowbreak{} editorial\allowbreak{} corrections}{these\allowbreak{} conditions\allowbreak{} c,\allowbreak{}\allowbreak{} c\$\allowbreak{}'\$\allowbreak{}\allowbreak{} }{these\allowbreak{} conditions\allowbreak{} c,\allowbreak{}\allowbreak{} b\$\allowbreak{}''\$\allowbreak{}\allowbreak{} }{DER\_\allowbreak{}XVIII-c:\allowbreak{}33}
\DERapparatusitem{ENV1254}{88;\allowbreak{}\allowbreak{} ms235,\allowbreak{}\allowbreak{} printed\allowbreak{} editorial\allowbreak{} corrections}{the\allowbreak{} five\allowbreak{} conditions\allowbreak{} a)\allowbreak{}\allowbreak{} b)\allowbreak{}\allowbreak{} b\$\allowbreak{}'\$\allowbreak{})\allowbreak{}\allowbreak{} c)\allowbreak{}\allowbreak{} c\$\allowbreak{}'\$\allowbreak{})\allowbreak{}}{the\allowbreak{} five\allowbreak{} conditions\allowbreak{} a)\allowbreak{}\allowbreak{} b)\allowbreak{}\allowbreak{} b\$\allowbreak{}'\$\allowbreak{})\allowbreak{}\allowbreak{} c)\allowbreak{}\allowbreak{} b\$\allowbreak{}''\$\allowbreak{})\allowbreak{}}{DER\_\allowbreak{}XVIII-c:\allowbreak{}33}
\DERapparatusitem{ENV1255}{88;\allowbreak{}\allowbreak{} ms237,\allowbreak{}\allowbreak{} diagonal\allowbreak{} label\allowbreak{} and\allowbreak{} its\allowbreak{} editorial\allowbreak{} correction}{\textbackslash{}\allowbreak{}delta\_\allowbreak{}\{\allowbreak{}A'\}\allowbreak{}}{\textbackslash{}\allowbreak{}delta'\_\allowbreak{}A}{DER\_\allowbreak{}XVIII-c:\allowbreak{}34}
\DERapparatusitem{ENV1256}{88;\allowbreak{}\allowbreak{} ms237,\allowbreak{}\allowbreak{} reference\allowbreak{} to\allowbreak{} product\allowbreak{} theorem}{also\allowbreak{} cofinal\allowbreak{} (\allowbreak{}4.\allowbreak{}9.\allowbreak{}8.\allowbreak{}5)\allowbreak{}}{also\allowbreak{} cofinal\allowbreak{} (\allowbreak{}4.\allowbreak{}9.\allowbreak{}6.\allowbreak{}5)\allowbreak{}}{DER\_\allowbreak{}XVIII-c:\allowbreak{}35}
\DERapparatusitem{ENV1257}{88;\allowbreak{}\allowbreak{} ms237,\allowbreak{}\allowbreak{} diagonal\allowbreak{} label\allowbreak{} and\allowbreak{} its\allowbreak{} editorial\allowbreak{} correction}{\textbackslash{}\allowbreak{}delta\_\allowbreak{}\{\allowbreak{}A'\}\allowbreak{}}{\textbackslash{}\allowbreak{}delta'\_\allowbreak{}A}{DER\_\allowbreak{}XVIII-c:\allowbreak{}34}
\DERapparatusitem{ENV1258}{88;\allowbreak{}\allowbreak{} ms237,\allowbreak{}\allowbreak{} diagonal\allowbreak{} label\allowbreak{} and\allowbreak{} its\allowbreak{} editorial\allowbreak{} correction}{\textbackslash{}\allowbreak{}delta\_\allowbreak{}\{\allowbreak{}A'\}\allowbreak{}}{\textbackslash{}\allowbreak{}delta'\_\allowbreak{}A}{DER\_\allowbreak{}XVIII-c:\allowbreak{}34}
\DERapparatusitem{ENV1259}{88;\allowbreak{}\allowbreak{} ms237,\allowbreak{}\allowbreak{} diagonal\allowbreak{} label\allowbreak{} and\allowbreak{} its\allowbreak{} editorial\allowbreak{} correction}{\textbackslash{}\allowbreak{}delta\_\allowbreak{}\{\allowbreak{}A'\}\allowbreak{}}{\textbackslash{}\allowbreak{}delta'\_\allowbreak{}A}{DER\_\allowbreak{}XVIII-c:\allowbreak{}34}
\DERapparatusitem{ENV1260}{88;\allowbreak{}\allowbreak{} ms237,\allowbreak{}\allowbreak{} connectedness\allowbreak{} of\allowbreak{} comma}{}{'}{DER\_\allowbreak{}XVIII-c:\allowbreak{}36}
\DERapparatusitem{ENV1261}{89;\allowbreak{}\allowbreak{} ms\allowbreak{} 240;\allowbreak{}\allowbreak{} existing\allowbreak{} [\allowbreak{}sic]\allowbreak{}\allowbreak{} flag}{}{\_\allowbreak{}\textbackslash{}\allowbreak{}alpha}{DER\_\allowbreak{}XVIII-c:\allowbreak{}38}
\DERapparatusitem{ENV1262}{89;\allowbreak{}\allowbreak{} ms240,\allowbreak{}\allowbreak{} first\allowbreak{} F3\allowbreak{} equation}{\_\allowbreak{}\textbackslash{}\allowbreak{}alpha)\allowbreak{}u'\_\allowbreak{}\textbackslash{}\allowbreak{}alpha\allowbreak{} \textbackslash{}\allowbreak{}\allowbreak{} =\allowbreak{}\allowbreak{} \textbackslash{}\allowbreak{}\allowbreak{} f(\allowbreak{}w}{\textasciicircum{}\allowbreak{}0\_\allowbreak{}\textbackslash{}\allowbreak{}alpha)\allowbreak{}u'\_\allowbreak{}\textbackslash{}\allowbreak{}alpha\allowbreak{} \textbackslash{}\allowbreak{}\allowbreak{} =\allowbreak{}\allowbreak{} \textbackslash{}\allowbreak{}\allowbreak{} f(\allowbreak{}w\textasciicircum{}\allowbreak{}0}{DER\_\allowbreak{}XVIII-c:\allowbreak{}39}
\DERapparatusitem{ENV1263}{90;\allowbreak{}\allowbreak{} ms240,\allowbreak{}\allowbreak{} labels\allowbreak{} and\allowbreak{} equality\allowbreak{} underbrace}{}{\textasciicircum{}\allowbreak{}0}{DER\_\allowbreak{}XVIII-c:\allowbreak{}41}
\DERapparatusitem{ENV1264}{90;\allowbreak{}\allowbreak{} ms\allowbreak{} 240;\allowbreak{}\allowbreak{} existing\allowbreak{} [\allowbreak{}sic]\allowbreak{}\allowbreak{} flag}{h}{j}{DER\_\allowbreak{}XVIII-c:\allowbreak{}40}
\DERapparatusitem{ENV1265}{90;\allowbreak{}\allowbreak{} ms240,\allowbreak{}\allowbreak{} labels\allowbreak{} and\allowbreak{} equality\allowbreak{} underbrace}{\_\allowbreak{}\textbackslash{}\allowbreak{}alpha)\allowbreak{}u'\_\allowbreak{}\textbackslash{}\allowbreak{}alpha\}\allowbreak{}\_\allowbreak{}\{\allowbreak{}f(\allowbreak{}w}{\textasciicircum{}\allowbreak{}0\_\allowbreak{}\textbackslash{}\allowbreak{}alpha)\allowbreak{}u'\_\allowbreak{}\textbackslash{}\allowbreak{}alpha\}\allowbreak{}\_\allowbreak{}\{\allowbreak{}f(\allowbreak{}w\textasciicircum{}\allowbreak{}0}{DER\_\allowbreak{}XVIII-c:\allowbreak{}41}
\DERapparatusitem{ENV1266}{90;\allowbreak{}\allowbreak{} ms241,\allowbreak{}\allowbreak{} paragraph\allowbreak{} proving\allowbreak{} b-double-prime\_\allowbreak{}pi\allowbreak{} implies\allowbreak{} b\_\allowbreak{}pi}{}{'}{DER\_\allowbreak{}XVIII-c:\allowbreak{}42}
\DERapparatusitem{ENV1267}{91;\allowbreak{}\allowbreak{} ms244;\allowbreak{}\allowbreak{} three\allowbreak{} references\allowbreak{} to\allowbreak{} the\allowbreak{} inclusion\allowbreak{} and\allowbreak{} its\allowbreak{} members}{and\allowbreak{} the\allowbreak{} inclusion\allowbreak{} (\allowbreak{}4.\allowbreak{}10.\allowbreak{}3)\allowbreak{}}{and\allowbreak{} the\allowbreak{} inclusion\allowbreak{} (\allowbreak{}4.\allowbreak{}10.\allowbreak{}4)\allowbreak{}}{DER\_\allowbreak{}XVIII-c:\allowbreak{}44}
\DERapparatusitem{ENV1268}{91;\allowbreak{}\allowbreak{} ms244;\allowbreak{}\allowbreak{} three\allowbreak{} references\allowbreak{} to\allowbreak{} the\allowbreak{} inclusion\allowbreak{} and\allowbreak{} its\allowbreak{} members}{The\allowbreak{} second\allowbreak{} member\allowbreak{} of\allowbreak{} (\allowbreak{}4.\allowbreak{}10.\allowbreak{}3)\allowbreak{}}{The\allowbreak{} second\allowbreak{} member\allowbreak{} of\allowbreak{} (\allowbreak{}4.\allowbreak{}10.\allowbreak{}4)\allowbreak{}}{DER\_\allowbreak{}XVIII-c:\allowbreak{}44}
\DERapparatusitem{ENV1269}{91;\allowbreak{}\allowbreak{} ms244;\allowbreak{}\allowbreak{} three\allowbreak{} references\allowbreak{} to\allowbreak{} the\allowbreak{} inclusion\allowbreak{} and\allowbreak{} its\allowbreak{} members}{The\allowbreak{} first\allowbreak{} member\allowbreak{} of\allowbreak{} (\allowbreak{}4.\allowbreak{}10.\allowbreak{}3)\allowbreak{}}{The\allowbreak{} first\allowbreak{} member\allowbreak{} of\allowbreak{} (\allowbreak{}4.\allowbreak{}10.\allowbreak{}4)\allowbreak{}}{DER\_\allowbreak{}XVIII-c:\allowbreak{}44}
\DERapparatusitem{ENV1270}{92;\allowbreak{}\allowbreak{} ms245;\allowbreak{}\allowbreak{} transport\allowbreak{} notation\allowbreak{} in\allowbreak{} Theorem4.\allowbreak{}10.\allowbreak{}1}{\textbackslash{}\allowbreak{}alpha\textasciicircum{}\allowbreak{}}{\textbackslash{}\allowbreak{}alpha\_\allowbreak{}}{DER\_\allowbreak{}XVIII-c:\allowbreak{}46}
\DERapparatusitem{ENV1271}{92;\allowbreak{}\allowbreak{} ms245;\allowbreak{}\allowbreak{} (\allowbreak{}4.\allowbreak{}10.\allowbreak{}2.\allowbreak{}2)\allowbreak{}}{C}{\textbackslash{}\allowbreak{}mathcal\{\allowbreak{}D\}\allowbreak{}}{DER\_\allowbreak{}XVIII-c:\allowbreak{}47}
\DERapparatusitem{ENV1272}{92;\allowbreak{}\allowbreak{} ms246;\allowbreak{}\allowbreak{} description\allowbreak{} of\allowbreak{} slice\allowbreak{} transport}{C\_\allowbreak{}j\allowbreak{} \textbackslash{}\allowbreak{}\allowbreak{} \textbackslash{}\allowbreak{}longrightarrow\allowbreak{} \textbackslash{}\allowbreak{}\allowbreak{} C}{\textbackslash{}\allowbreak{}mathcal\{\allowbreak{}D\}\allowbreak{}\_\allowbreak{}j\allowbreak{} \textbackslash{}\allowbreak{}\allowbreak{} \textbackslash{}\allowbreak{}longrightarrow\allowbreak{} \textbackslash{}\allowbreak{}\allowbreak{} \textbackslash{}\allowbreak{}mathcal\{\allowbreak{}D\}\allowbreak{}}{DER\_\allowbreak{}XVIII-c:\allowbreak{}48}
\DERapparatusitem{ENV1273}{93;\allowbreak{}\allowbreak{} ms247;\allowbreak{}\allowbreak{} slice-transport\allowbreak{} composite}{C\_\allowbreak{}j/\allowbreak{}\textbackslash{}\allowbreak{}varphi(\allowbreak{}j)\allowbreak{}\allowbreak{} \textbackslash{}\allowbreak{}\allowbreak{} \textbackslash{}\allowbreak{}longrightarrow\allowbreak{} \textbackslash{}\allowbreak{}\allowbreak{} C\_\allowbreak{}\{\allowbreak{}j'\}\allowbreak{}/\allowbreak{}\textbackslash{}\allowbreak{}alpha\textasciicircum{}\allowbreak{}*(\allowbreak{}\textbackslash{}\allowbreak{}varphi(\allowbreak{}j)\allowbreak{})\allowbreak{}\allowbreak{} \textbackslash{}\allowbreak{}\allowbreak{} \textbackslash{}\allowbreak{}longrightarrow\allowbreak{} \textbackslash{}\allowbreak{}\allowbreak{} C}{\textbackslash{}\allowbreak{}mathcal\{\allowbreak{}D\}\allowbreak{}\_\allowbreak{}j/\allowbreak{}\textbackslash{}\allowbreak{}varphi(\allowbreak{}j)\allowbreak{}\allowbreak{} \textbackslash{}\allowbreak{}\allowbreak{} \textbackslash{}\allowbreak{}longrightarrow\allowbreak{} \textbackslash{}\allowbreak{}\allowbreak{} \textbackslash{}\allowbreak{}mathcal\{\allowbreak{}D\}\allowbreak{}\_\allowbreak{}\{\allowbreak{}j'\}\allowbreak{}/\allowbreak{}\textbackslash{}\allowbreak{}alpha\_\allowbreak{}*(\allowbreak{}\textbackslash{}\allowbreak{}varphi(\allowbreak{}j)\allowbreak{})\allowbreak{}\allowbreak{} \textbackslash{}\allowbreak{}\allowbreak{} \textbackslash{}\allowbreak{}longrightarrow\allowbreak{} \textbackslash{}\allowbreak{}\allowbreak{} \textbackslash{}\allowbreak{}mathcal\{\allowbreak{}D\}\allowbreak{}}{DER\_\allowbreak{}XVIII-c:\allowbreak{}49}
\DERapparatusitem{ENV1274}{93;\allowbreak{}\allowbreak{} ms247;\allowbreak{}\allowbreak{} hypothesis4.\allowbreak{}10.\allowbreak{}3(\allowbreak{}c)\allowbreak{}}{There\allowbreak{} exists\allowbreak{} a\allowbreak{} limit\allowbreak{} ordinal\allowbreak{} such\allowbreak{} that:\allowbreak{}}{There\allowbreak{} exists\allowbreak{} a\allowbreak{} limit\allowbreak{} ordinal\allowbreak{} \$\allowbreak{}\textbackslash{}\allowbreak{}alpha\$\allowbreak{}\allowbreak{} such\allowbreak{} that:\allowbreak{}}{DER\_\allowbreak{}XVIII-c:\allowbreak{}50}
\DERapparatusitem{ENV1275}{93;\allowbreak{}\allowbreak{} ms248;\allowbreak{}\allowbreak{} parenthetical\allowbreak{} ordinal\allowbreak{} example}{\textbackslash{}\allowbreak{}leqslant}{<}{DER\_\allowbreak{}XVIII-c:\allowbreak{}51}
\DERapparatusitem{ENV1276}{93;\allowbreak{}\allowbreak{} running\allowbreak{} text\allowbreak{} or\allowbreak{} formula}{\textbackslash{}\allowbreak{}textup\{\allowbreak{}def\}\allowbreak{}}{\textbackslash{}\allowbreak{}mathrm\{\allowbreak{}def\}\allowbreak{}}{TYPESETTING:\allowbreak{}7;\allowbreak{}\allowbreak{} NOTATION\_\allowbreak{}DER:\allowbreak{}17}
\DERapparatusitem{ENV1277}{94;\allowbreak{}\allowbreak{} ms250;\allowbreak{}\allowbreak{} cofinality-criterion\allowbreak{} reference}{Applying\allowbreak{} 4.\allowbreak{}9.\allowbreak{}6\allowbreak{} 2°}{Applying\allowbreak{} 4.\allowbreak{}9.\allowbreak{}6.\allowbreak{}3,\allowbreak{}\allowbreak{} 2°}{DER\_\allowbreak{}XVIII-c:\allowbreak{}52}
\DERapparatusitem{ENV1278}{95;\allowbreak{}\allowbreak{} ms252;\allowbreak{}\allowbreak{} target\allowbreak{} of\allowbreak{} natural\allowbreak{} section\allowbreak{} map}{\_\allowbreak{}\textbackslash{}\allowbreak{}alpha}{}{DER\_\allowbreak{}XVIII-c:\allowbreak{}53}
\DERapparatusitem{ENV1279}{96;\allowbreak{}\allowbreak{} ms254;\allowbreak{}\allowbreak{} label\allowbreak{} of\allowbreak{} prescribed\allowbreak{} component\allowbreak{} maps}{6.\allowbreak{}10}{10.\allowbreak{}6}{DER\_\allowbreak{}XVIII-c:\allowbreak{}54}
\DERapparatusitem{ENV1280}{96;\allowbreak{}\allowbreak{} ms254;\allowbreak{}\allowbreak{} X(\allowbreak{}j)\allowbreak{}→G(\allowbreak{}j)\allowbreak{}\allowbreak{} label}{\textbackslash{}\allowbreak{}xrightarrow\{\allowbreak{}v(\allowbreak{}j)\allowbreak{}\}\allowbreak{}\allowbreak{} G(\allowbreak{}j)\allowbreak{}}{\textbackslash{}\allowbreak{}xrightarrow\{\allowbreak{}w(\allowbreak{}j)\allowbreak{}\}\allowbreak{}\allowbreak{} G(\allowbreak{}j)\allowbreak{}}{DER\_\allowbreak{}XVIII-c:\allowbreak{}55}
\DERapparatusitem{ENV1281}{96;\allowbreak{}\allowbreak{} ms255;\allowbreak{}\allowbreak{} source\allowbreak{} of\allowbreak{} (\allowbreak{}4.\allowbreak{}10.\allowbreak{}6.\allowbreak{}5)\allowbreak{}}{}{'}{DER\_\allowbreak{}XVIII-c:\allowbreak{}56}
\DERapparatusitem{ENV1282}{97;\allowbreak{}\allowbreak{} ms256;\allowbreak{}\allowbreak{} condition(\allowbreak{}i)\allowbreak{}}{\textbackslash{}\allowbreak{}mathrm\{\allowbreak{}id\}\allowbreak{}\_\allowbreak{}J\allowbreak{} :\allowbreak{}\allowbreak{} J\textbackslash{}\allowbreak{}longrightarrow\allowbreak{} J}{\textbackslash{}\allowbreak{}mathrm\{\allowbreak{}id\}\allowbreak{}\_\allowbreak{}j\allowbreak{} :\allowbreak{}\allowbreak{} j\textbackslash{}\allowbreak{}longrightarrow\allowbreak{} j}{DER\_\allowbreak{}XVIII-c:\allowbreak{}57}
\DERapparatusitem{ENV1283}{97;\allowbreak{}\allowbreak{} ms256;\allowbreak{}\allowbreak{} lower\allowbreak{} row\allowbreak{} of\allowbreak{} diagram(\allowbreak{}ii)\allowbreak{}}{')\allowbreak{}\allowbreak{} \textbackslash{}\allowbreak{}arrow[\allowbreak{}r]\allowbreak{}\allowbreak{} \&\allowbreak{}\allowbreak{} X\_\allowbreak{}\{\allowbreak{}\textbackslash{}\allowbreak{}alpha+1\}\allowbreak{}(\allowbreak{}j')\allowbreak{}\allowbreak{} \textbackslash{}\allowbreak{}arrow[\allowbreak{}r,\allowbreak{}"\textbackslash{}\allowbreak{}tau\_\allowbreak{}\{\allowbreak{}\textbackslash{}\allowbreak{}alpha+1\}\allowbreak{}(\allowbreak{}j')\allowbreak{}"]\allowbreak{}\allowbreak{} \&\allowbreak{}\allowbreak{} X\_\allowbreak{}\{\allowbreak{}\textbackslash{}\allowbreak{}alpha+2\}\allowbreak{}(\allowbreak{}j')\allowbreak{}\allowbreak{} \textbackslash{}\allowbreak{}arrow[\allowbreak{}r,\allowbreak{}"\textbackslash{}\allowbreak{}tau\_\allowbreak{}\{\allowbreak{}\textbackslash{}\allowbreak{}alpha+2\}\allowbreak{}(\allowbreak{}j')\allowbreak{}"]\allowbreak{}\allowbreak{} \&\allowbreak{}\allowbreak{} X\_\allowbreak{}\{\allowbreak{}\textbackslash{}\allowbreak{}alpha+3\}\allowbreak{}(\allowbreak{}j'}{)\allowbreak{}\allowbreak{} \textbackslash{}\allowbreak{}arrow[\allowbreak{}r]\allowbreak{}\allowbreak{} \&\allowbreak{}\allowbreak{} X\_\allowbreak{}\{\allowbreak{}\textbackslash{}\allowbreak{}alpha+1\}\allowbreak{}(\allowbreak{}j)\allowbreak{}\allowbreak{} \textbackslash{}\allowbreak{}arrow[\allowbreak{}r,\allowbreak{}"\textbackslash{}\allowbreak{}tau\_\allowbreak{}\{\allowbreak{}\textbackslash{}\allowbreak{}alpha+1\}\allowbreak{}(\allowbreak{}j)\allowbreak{}"]\allowbreak{}\allowbreak{} \&\allowbreak{}\allowbreak{} X\_\allowbreak{}\{\allowbreak{}\textbackslash{}\allowbreak{}alpha+2\}\allowbreak{}(\allowbreak{}j)\allowbreak{}\allowbreak{} \textbackslash{}\allowbreak{}arrow[\allowbreak{}r,\allowbreak{}"\textbackslash{}\allowbreak{}tau\_\allowbreak{}\{\allowbreak{}\textbackslash{}\allowbreak{}alpha+2\}\allowbreak{}(\allowbreak{}j)\allowbreak{}"]\allowbreak{}\allowbreak{} \&\allowbreak{}\allowbreak{} X\_\allowbreak{}\{\allowbreak{}\textbackslash{}\allowbreak{}alpha+3\}\allowbreak{}(\allowbreak{}j}{DER\_\allowbreak{}XVIII-c:\allowbreak{}58}
\DERapparatusitem{ENV1284}{98;\allowbreak{}\allowbreak{} ms259;\allowbreak{}\allowbreak{} sequence\allowbreak{} in\allowbreak{} Lemma4.\allowbreak{}10.\allowbreak{}7}{i}{j}{DER\_\allowbreak{}XVIII-c:\allowbreak{}59}
\DERapparatusitem{ENV1285}{98;\allowbreak{}\allowbreak{} ms259;\allowbreak{}\allowbreak{} arrow\allowbreak{} in\allowbreak{} Lemma4.\allowbreak{}10.\allowbreak{}7}{t\_\allowbreak{}\textbackslash{}\allowbreak{}alpha(\allowbreak{}\textbackslash{}\allowbreak{}lambda)\allowbreak{}\allowbreak{} :\allowbreak{}\allowbreak{} X\_\allowbreak{}\textbackslash{}\allowbreak{}alpha(\allowbreak{}j')\allowbreak{}\textbackslash{}\allowbreak{}longrightarrow\allowbreak{} X\_\allowbreak{}\textbackslash{}\allowbreak{}alpha(\allowbreak{}j)\allowbreak{}}{t\_\allowbreak{}\textbackslash{}\allowbreak{}alpha(\allowbreak{}\textbackslash{}\allowbreak{}lambda)\allowbreak{}\allowbreak{} :\allowbreak{}\allowbreak{} X\_\allowbreak{}\textbackslash{}\allowbreak{}alpha(\allowbreak{}j')\allowbreak{}\textbackslash{}\allowbreak{}longrightarrow\allowbreak{} X\_\allowbreak{}\{\allowbreak{}\textbackslash{}\allowbreak{}alpha+1\}\allowbreak{}(\allowbreak{}j)\allowbreak{}}{DER\_\allowbreak{}XVIII-c:\allowbreak{}60}
\DERapparatusitem{ENV1286}{98;\allowbreak{}\allowbreak{} ms260;\allowbreak{}\allowbreak{} count\allowbreak{} of\allowbreak{} optional\allowbreak{} naturality\allowbreak{} equations}{\textbackslash{}\allowbreak{}card\textbackslash{}\allowbreak{}Ob\textbackslash{}\allowbreak{}mathcal\{\allowbreak{}E\}\allowbreak{}/\allowbreak{}j}{\textbackslash{}\allowbreak{}card\textbackslash{}\allowbreak{}Ob(\allowbreak{}J/\allowbreak{}j)\allowbreak{}}{DER\_\allowbreak{}XVIII-c:\allowbreak{}61}
\DERapparatusitem{ENV1287}{98;\allowbreak{}\allowbreak{} ms261;\allowbreak{}\allowbreak{} display\allowbreak{} defining\allowbreak{} G(\allowbreak{}λ)\allowbreak{}}{)\allowbreak{}=\allowbreak{}\textbackslash{}\allowbreak{}varinjlim\allowbreak{} X\_\allowbreak{}\textbackslash{}\allowbreak{}beta(\allowbreak{}j)\allowbreak{}\allowbreak{} \textbackslash{}\allowbreak{}longrightarrow\allowbreak{} G(\allowbreak{}j')\allowbreak{}\textbackslash{}\allowbreak{}simeq\textbackslash{}\allowbreak{}varinjlim\allowbreak{} X\_\allowbreak{}\{\allowbreak{}\textbackslash{}\allowbreak{}beta+1\}\allowbreak{}(\allowbreak{}j'}{')\allowbreak{}=\allowbreak{}\textbackslash{}\allowbreak{}varinjlim\allowbreak{} X\_\allowbreak{}\textbackslash{}\allowbreak{}beta(\allowbreak{}j')\allowbreak{}\allowbreak{} \textbackslash{}\allowbreak{}longrightarrow\allowbreak{} G(\allowbreak{}j)\allowbreak{}\textbackslash{}\allowbreak{}simeq\textbackslash{}\allowbreak{}varinjlim\allowbreak{} X\_\allowbreak{}\{\allowbreak{}\textbackslash{}\allowbreak{}beta+1\}\allowbreak{}(\allowbreak{}j}{DER\_\allowbreak{}XVIII-c:\allowbreak{}62}
\DERapparatusitem{ENV1288}{98;\allowbreak{}\allowbreak{} ms262;\allowbreak{}\allowbreak{} type\allowbreak{} of\allowbreak{} systems}{a\allowbreak{} homomorphism\allowbreak{} of\allowbreak{} inverse\allowbreak{} systems}{a\allowbreak{} homomorphism\allowbreak{} of\allowbreak{} direct\allowbreak{} systems}{DER\_\allowbreak{}XVIII-c:\allowbreak{}63}
\DERapparatusitem{ENV1289}{preamble;\allowbreak{}\allowbreak{} local\allowbreak{} macro\allowbreak{} definition}{\textbackslash{}\allowbreak{}text\{\allowbreak{}\textbackslash{}\allowbreak{}upshape\allowbreak{} def\}\allowbreak{}}{\textbackslash{}\allowbreak{}mathrm\{\allowbreak{}def\}\allowbreak{}}{TYPESETTING:\allowbreak{}7;\allowbreak{}\allowbreak{} NOTATION\_\allowbreak{}DER:\allowbreak{}17}
\DERapparatusitem{ENV1290}{99;\allowbreak{}\allowbreak{} ms263;\allowbreak{}\allowbreak{} condition(\allowbreak{}b)\allowbreak{}\allowbreak{} of4.\allowbreak{}10.\allowbreak{}8}{\$\allowbreak{}\textbackslash{}\allowbreak{}lambda\_\allowbreak{}*\allowbreak{} :\allowbreak{}\allowbreak{} \textbackslash{}\allowbreak{}mathcal\{\allowbreak{}E\}\allowbreak{}\_\allowbreak{}j\allowbreak{} \textbackslash{}\allowbreak{}longrightarrow\allowbreak{} \textbackslash{}\allowbreak{}mathcal\{\allowbreak{}E\}\allowbreak{}\_\allowbreak{}\{\allowbreak{}j'\}\allowbreak{}\$\allowbreak{}\allowbreak{} \textbackslash{}\allowbreak{}textit\{\allowbreak{}maps\}\allowbreak{}\allowbreak{} \$\allowbreak{}(\allowbreak{}\textbackslash{}\allowbreak{}mathcal\{\allowbreak{}E\}\allowbreak{}\_\allowbreak{}j)\allowbreak{}\_\allowbreak{}\{\allowbreak{}\textbackslash{}\allowbreak{}pi\_\allowbreak{}0\}\allowbreak{}\$\allowbreak{}\allowbreak{} \textbackslash{}\allowbreak{}textit\{\allowbreak{}into\}\allowbreak{}\allowbreak{} \$\allowbreak{}(\allowbreak{}\textbackslash{}\allowbreak{}mathcal\{\allowbreak{}E\}\allowbreak{}\_\allowbreak{}\{\allowbreak{}j'\}\allowbreak{})\allowbreak{}\_\allowbreak{}\{\allowbreak{}\textbackslash{}\allowbreak{}pi\_\allowbreak{}0\}\allowbreak{}\$\allowbreak{}}{\$\allowbreak{}\textbackslash{}\allowbreak{}lambda\_\allowbreak{}*\allowbreak{} :\allowbreak{}\allowbreak{} \textbackslash{}\allowbreak{}mathcal\{\allowbreak{}N\}\allowbreak{}\_\allowbreak{}\{\allowbreak{}j'\}\allowbreak{}\allowbreak{} \textbackslash{}\allowbreak{}longrightarrow\allowbreak{} \textbackslash{}\allowbreak{}mathcal\{\allowbreak{}N\}\allowbreak{}\_\allowbreak{}j\$\allowbreak{}\allowbreak{} \textbackslash{}\allowbreak{}textit\{\allowbreak{}maps\}\allowbreak{}\allowbreak{} \$\allowbreak{}(\allowbreak{}\textbackslash{}\allowbreak{}mathcal\{\allowbreak{}N\}\allowbreak{}\_\allowbreak{}\{\allowbreak{}j'\}\allowbreak{})\allowbreak{}\_\allowbreak{}\{\allowbreak{}\textbackslash{}\allowbreak{}pi\_\allowbreak{}0\}\allowbreak{}\$\allowbreak{}\allowbreak{} \textbackslash{}\allowbreak{}textit\{\allowbreak{}into\}\allowbreak{}\allowbreak{} \$\allowbreak{}(\allowbreak{}\textbackslash{}\allowbreak{}mathcal\{\allowbreak{}N\}\allowbreak{}\_\allowbreak{}j)\allowbreak{}\_\allowbreak{}\{\allowbreak{}\textbackslash{}\allowbreak{}pi\_\allowbreak{}0\}\allowbreak{}\$\allowbreak{}}{DER\_\allowbreak{}XVIII-c:\allowbreak{}64}
\DERapparatusitem{ENV1291}{99;\allowbreak{}\allowbreak{} ms263-264;\allowbreak{}\allowbreak{} description\allowbreak{} of\allowbreak{} M\_\allowbreak{}π}{M}{N}{DER\_\allowbreak{}XVIII-c:\allowbreak{}65}
\DERapparatusitem{ENV1292}{99;\allowbreak{}\allowbreak{} ms263-264;\allowbreak{}\allowbreak{} description\allowbreak{} of\allowbreak{} M\_\allowbreak{}π}{E}{N}{DER\_\allowbreak{}XVIII-c:\allowbreak{}65}
\DERapparatusitem{ENV1293}{99;\allowbreak{}\allowbreak{} ms263-264;\allowbreak{}\allowbreak{} description\allowbreak{} of\allowbreak{} M\_\allowbreak{}π}{\$\allowbreak{}(\allowbreak{}\textbackslash{}\allowbreak{}mathcal\{\allowbreak{}E\}\allowbreak{}\_\allowbreak{}j)\allowbreak{}\_\allowbreak{}\textbackslash{}\allowbreak{}pi\$\allowbreak{}\allowbreak{} denoting\allowbreak{} the\allowbreak{} full\allowbreak{} subcategory\allowbreak{} of\allowbreak{} \$\allowbreak{}\textbackslash{}\allowbreak{}mathcal\{\allowbreak{}E\}\allowbreak{}\_\allowbreak{}j\$\allowbreak{}\allowbreak{} consisting\allowbreak{} of\allowbreak{} the\allowbreak{} \$\allowbreak{}\textbackslash{}\allowbreak{}pi\$\allowbreak{}-accessible\allowbreak{} objects\allowbreak{} of\allowbreak{} \$\allowbreak{}\textbackslash{}\allowbreak{}mathcal\{\allowbreak{}E\}\allowbreak{}\_\allowbreak{}j\$\allowbreak{}}{\$\allowbreak{}(\allowbreak{}\textbackslash{}\allowbreak{}mathcal\{\allowbreak{}N\}\allowbreak{}\_\allowbreak{}j)\allowbreak{}\_\allowbreak{}\textbackslash{}\allowbreak{}pi\$\allowbreak{}\allowbreak{} denoting\allowbreak{} the\allowbreak{} full\allowbreak{} subcategory\allowbreak{} of\allowbreak{} \$\allowbreak{}\textbackslash{}\allowbreak{}mathcal\{\allowbreak{}N\}\allowbreak{}\_\allowbreak{}j\$\allowbreak{}\allowbreak{} consisting\allowbreak{} of\allowbreak{} the\allowbreak{} \$\allowbreak{}\textbackslash{}\allowbreak{}pi\$\allowbreak{}-accessible\allowbreak{} objects\allowbreak{} of\allowbreak{} \$\allowbreak{}\textbackslash{}\allowbreak{}mathcal\{\allowbreak{}N\}\allowbreak{}\_\allowbreak{}j\$\allowbreak{}}{DER\_\allowbreak{}XVIII-c:\allowbreak{}65}
\DERapparatusitem{ENV1294}{100;\allowbreak{}\allowbreak{} ms266;\allowbreak{}\allowbreak{} variance\allowbreak{} in4.\allowbreak{}10.\allowbreak{}10\allowbreak{} proof}{the\allowbreak{} functor\allowbreak{} \$\allowbreak{}\textbackslash{}\allowbreak{}lambda\_\allowbreak{}*\$\allowbreak{}\allowbreak{} :\allowbreak{}\allowbreak{} \$\allowbreak{}\textbackslash{}\allowbreak{}mathcal\{\allowbreak{}N\}\allowbreak{}\_\allowbreak{}\{\allowbreak{}j'\}\allowbreak{}\allowbreak{} \textbackslash{}\allowbreak{}longrightarrow\allowbreak{} \textbackslash{}\allowbreak{}mathcal\{\allowbreak{}N\}\allowbreak{}\_\allowbreak{}j\$\allowbreak{}\allowbreak{} admits\allowbreak{} a\allowbreak{} left\allowbreak{} adjoint\allowbreak{} \$\allowbreak{}\textbackslash{}\allowbreak{}lambda\textasciicircum{}\allowbreak{}*\allowbreak{} :\allowbreak{}\allowbreak{} \textbackslash{}\allowbreak{}mathcal\{\allowbreak{}N\}\allowbreak{}\_\allowbreak{}j\allowbreak{} \textbackslash{}\allowbreak{}longrightarrow\allowbreak{} \textbackslash{}\allowbreak{}mathcal\{\allowbreak{}N\}\allowbreak{}\_\allowbreak{}\{\allowbreak{}j'\}\allowbreak{}\$\allowbreak{}}{the\allowbreak{} functor\allowbreak{} \$\allowbreak{}\textbackslash{}\allowbreak{}lambda\textasciicircum{}\allowbreak{}*\$\allowbreak{}\allowbreak{} :\allowbreak{}\allowbreak{} \$\allowbreak{}\textbackslash{}\allowbreak{}mathcal\{\allowbreak{}N\}\allowbreak{}\_\allowbreak{}j\allowbreak{} \textbackslash{}\allowbreak{}longrightarrow\allowbreak{} \textbackslash{}\allowbreak{}mathcal\{\allowbreak{}N\}\allowbreak{}\_\allowbreak{}\{\allowbreak{}j'\}\allowbreak{}\$\allowbreak{}\allowbreak{} admits\allowbreak{} a\allowbreak{} left\allowbreak{} adjoint\allowbreak{} \$\allowbreak{}\textbackslash{}\allowbreak{}lambda\_\allowbreak{}!\allowbreak{} :\allowbreak{}\allowbreak{} \textbackslash{}\allowbreak{}mathcal\{\allowbreak{}N\}\allowbreak{}\_\allowbreak{}\{\allowbreak{}j'\}\allowbreak{}\allowbreak{} \textbackslash{}\allowbreak{}longrightarrow\allowbreak{} \textbackslash{}\allowbreak{}mathcal\{\allowbreak{}N\}\allowbreak{}\_\allowbreak{}j\$\allowbreak{}}{DER\_\allowbreak{}XVIII-c:\allowbreak{}68}
\DERapparatusitem{ENV1295}{100;\allowbreak{}\allowbreak{} ms266;\allowbreak{}\allowbreak{} proof\allowbreak{} conclusion\allowbreak{} reference}{we\allowbreak{} can\allowbreak{} apply\allowbreak{} 4.\allowbreak{}8.\allowbreak{}9}{we\allowbreak{} can\allowbreak{} apply\allowbreak{} 4.\allowbreak{}10.\allowbreak{}9}{DER\_\allowbreak{}XVIII-c:\allowbreak{}69}
\DERapparatusitem{ENV1296}{100;\allowbreak{}\allowbreak{} ms267;\allowbreak{}\allowbreak{} definition\allowbreak{} of\allowbreak{} section\allowbreak{} category}{\textbackslash{}\allowbreak{}uGamma(\allowbreak{}\textbackslash{}\allowbreak{}mathcal\{\allowbreak{}N\}\allowbreak{}/\allowbreak{}J)\allowbreak{}\allowbreak{} \textbackslash{}\allowbreak{}defeq\allowbreak{} \textbackslash{}\allowbreak{}Hom(\allowbreak{}J,\allowbreak{}\textbackslash{}\allowbreak{}mathcal\{\allowbreak{}N\}\allowbreak{})\allowbreak{}}{\textbackslash{}\allowbreak{}uGamma(\allowbreak{}\textbackslash{}\allowbreak{}mathcal\{\allowbreak{}N\}\allowbreak{}/\allowbreak{}J)\allowbreak{}\allowbreak{} \textbackslash{}\allowbreak{}defeq\allowbreak{} \textbackslash{}\allowbreak{}Hom\_\allowbreak{}J(\allowbreak{}J,\allowbreak{}\textbackslash{}\allowbreak{}mathcal\{\allowbreak{}N\}\allowbreak{})\allowbreak{}}{DER\_\allowbreak{}XVIII-c:\allowbreak{}70}
\DERapparatusitem{ENV1297}{101;\allowbreak{}\allowbreak{} ms270;\allowbreak{}\allowbreak{} parenthetical\allowbreak{} in\allowbreak{} proof4.\allowbreak{}11.\allowbreak{}2}{}{)\allowbreak{}}{DER\_\allowbreak{}XVIII-c:\allowbreak{}71}
\DERapparatusitem{ENV1298}{102;\allowbreak{}\allowbreak{} ms273;\allowbreak{}\allowbreak{} first\allowbreak{} label\allowbreak{} in4.\allowbreak{}12.\allowbreak{}1}{}{.\allowbreak{}1}{DER\_\allowbreak{}XVIII-c:\allowbreak{}73}
\DERapparatusitem{ENV1299}{102;\allowbreak{}\allowbreak{} ms274;\allowbreak{}\allowbreak{} first\allowbreak{} proof\allowbreak{} display}{\textbackslash{}\allowbreak{}mathcal\{\allowbreak{}N\}\allowbreak{}\_\allowbreak{}1\allowbreak{} \textbackslash{}\allowbreak{}\allowbreak{} \textbackslash{}\allowbreak{}xleftarrow\{\allowbreak{}\textbackslash{}\allowbreak{}\allowbreak{} \textbackslash{}\allowbreak{}sim\textbackslash{}\allowbreak{}\allowbreak{} \}\allowbreak{}\allowbreak{} \textbackslash{}\allowbreak{}Ind\_\allowbreak{}\{\allowbreak{}\textbackslash{}\allowbreak{}pi\_\allowbreak{}0\}\allowbreak{}(\allowbreak{}\textbackslash{}\allowbreak{}mathcal\{\allowbreak{}M\}\allowbreak{}\_\allowbreak{}\{\allowbreak{}\textbackslash{}\allowbreak{}pi\_\allowbreak{}0\}\allowbreak{})\allowbreak{}}{\textbackslash{}\allowbreak{}mathcal\{\allowbreak{}M\}\allowbreak{}\allowbreak{} \textbackslash{}\allowbreak{}\allowbreak{} \textbackslash{}\allowbreak{}xleftarrow\{\allowbreak{}\textbackslash{}\allowbreak{}\allowbreak{} \textbackslash{}\allowbreak{}sim\textbackslash{}\allowbreak{}\allowbreak{} \}\allowbreak{}\allowbreak{} \textbackslash{}\allowbreak{}Ind\_\allowbreak{}\{\allowbreak{}\textbackslash{}\allowbreak{}pi\_\allowbreak{}0\}\allowbreak{}(\allowbreak{}\textbackslash{}\allowbreak{}mathcal\{\allowbreak{}M\}\allowbreak{}\_\allowbreak{}\{\allowbreak{}\textbackslash{}\allowbreak{}pi\_\allowbreak{}0\}\allowbreak{})\allowbreak{}}{DER\_\allowbreak{}XVIII-c:\allowbreak{}74}
\DERapparatusitem{ENV1300}{102;\allowbreak{}\allowbreak{} ms274;\allowbreak{}\allowbreak{} exponent\allowbreak{} identity\allowbreak{} with\allowbreak{} editorial\allowbreak{} correction}{\textbackslash{}\allowbreak{}pi\_\allowbreak{}0}{\textbackslash{}\allowbreak{}pi'}{DER\_\allowbreak{}XVIII-c:\allowbreak{}75}
\DERapparatusitem{ENV1301}{103;\allowbreak{}\allowbreak{} ms274;\allowbreak{}\allowbreak{} objects\allowbreak{} in\allowbreak{} bounded\allowbreak{} presentation}{N}{M}{DER\_\allowbreak{}XVIII-c:\allowbreak{}76}
\DERapparatusitem{ENV1302}{103;\allowbreak{}\allowbreak{} ms276;\allowbreak{}\allowbreak{} composite-functor\allowbreak{} reference}{The\allowbreak{} composite\allowbreak{} functor\allowbreak{} (\allowbreak{}4.\allowbreak{}11.\allowbreak{}3.\allowbreak{}2)\allowbreak{}}{The\allowbreak{} composite\allowbreak{} functor\allowbreak{} (\allowbreak{}4.\allowbreak{}12.\allowbreak{}3.\allowbreak{}2)\allowbreak{}}{DER\_\allowbreak{}XVIII-c:\allowbreak{}77}
\DERapparatusitem{ENV1303}{104;\allowbreak{}\allowbreak{} ms.\allowbreak{}\allowbreak{} 279;\allowbreak{}\allowbreak{} flagged\allowbreak{} source\allowbreak{} slip}{\textbackslash{}\allowbreak{}M}{\textbackslash{}\allowbreak{}N}{DER\_\allowbreak{}XVIII-c:\allowbreak{}78}
\DERapparatusitem{ENV1304}{105;\allowbreak{}\allowbreak{} ms.\allowbreak{}\allowbreak{} 282;\allowbreak{}\allowbreak{} flagged\allowbreak{} source\allowbreak{} slip}{\$\allowbreak{}\textbackslash{}\allowbreak{}uS\$\allowbreak{}\allowbreak{} (\allowbreak{}assumed\allowbreak{} accessible\allowbreak{} in\allowbreak{} \$\allowbreak{}\textbackslash{}\allowbreak{}M\$\allowbreak{})\allowbreak{}\allowbreak{} is\allowbreak{} \$\allowbreak{}\textbackslash{}\allowbreak{}pi\_\allowbreak{}0\$\allowbreak{}-\textbackslash{}\allowbreak{}textit\{\allowbreak{}accessible\}\allowbreak{}}{\$\allowbreak{}\textbackslash{}\allowbreak{}uS\$\allowbreak{}\allowbreak{} (\allowbreak{}assumed\allowbreak{} accessible\allowbreak{} in\allowbreak{} \$\allowbreak{}\textbackslash{}\allowbreak{}M\$\allowbreak{})\allowbreak{}\allowbreak{} is\allowbreak{} \$\allowbreak{}\textbackslash{}\allowbreak{}pi'\_\allowbreak{}0\$\allowbreak{}-\textbackslash{}\allowbreak{}textit\{\allowbreak{}accessible\}\allowbreak{}}{DER\_\allowbreak{}XVIII-c:\allowbreak{}80}
\DERapparatusitem{ENV1305}{105;\allowbreak{}\allowbreak{} ms.\allowbreak{}\allowbreak{} 282;\allowbreak{}\allowbreak{} flagged\allowbreak{} source\allowbreak{} slip}{\textbackslash{}\allowbreak{}M}{\textbackslash{}\allowbreak{}M\_\allowbreak{}\textbackslash{}\allowbreak{}pi}{DER\_\allowbreak{}XVIII-c:\allowbreak{}81}
\DERapparatusitem{ENV1306}{105;\allowbreak{}\allowbreak{} ms.\allowbreak{}\allowbreak{} 284;\allowbreak{}\allowbreak{} flagged\allowbreak{} source\allowbreak{} slip}{\textbackslash{}\allowbreak{}pi\_\allowbreak{}1}{\textbackslash{}\allowbreak{}pi}{DER\_\allowbreak{}XVIII-c:\allowbreak{}82}
\DERapparatusitem{ENV1307}{105;\allowbreak{}\allowbreak{} ms.\allowbreak{}\allowbreak{} 284;\allowbreak{}\allowbreak{} flagged\allowbreak{} source\allowbreak{} slip}{\textbackslash{}\allowbreak{}M}{\textbackslash{}\allowbreak{}M\_\allowbreak{}\textbackslash{}\allowbreak{}pi}{DER\_\allowbreak{}XVIII-c:\allowbreak{}83}
\DERapparatusitem{ENV1308}{106;\allowbreak{}\allowbreak{} ms.\allowbreak{}285,\allowbreak{}\allowbreak{} Proposition4.\allowbreak{}13.\allowbreak{}2(\allowbreak{}b)\allowbreak{}}{satisfies\allowbreak{} the\allowbreak{} conditions\allowbreak{} of\allowbreak{} c)\allowbreak{},\allowbreak{}\allowbreak{} taking}{satisfies\allowbreak{} the\allowbreak{} conditions\allowbreak{} of\allowbreak{} a)\allowbreak{},\allowbreak{}\allowbreak{} taking}{DER\_\allowbreak{}XVIII-c:\allowbreak{}86}
\DERapparatusitem{ENV1309}{106;\allowbreak{}\allowbreak{} ms.\allowbreak{}\allowbreak{} 286;\allowbreak{}\allowbreak{} flagged\allowbreak{} source\allowbreak{} slip}{the\allowbreak{} conditions\allowbreak{} of\allowbreak{} 4.\allowbreak{}13.\allowbreak{}3\allowbreak{} a)\allowbreak{}}{the\allowbreak{} conditions\allowbreak{} of\allowbreak{} 4.\allowbreak{}13.\allowbreak{}2\allowbreak{} a)\allowbreak{}}{DER\_\allowbreak{}XVIII-c:\allowbreak{}84}
\DERapparatusitem{ENV1310}{107;\allowbreak{}\allowbreak{} ms.\allowbreak{}288,\allowbreak{}\allowbreak{} Proposition4.\allowbreak{}13.\allowbreak{}4}{the\allowbreak{} non-fully\allowbreak{} faithful\allowbreak{} functor}{the\allowbreak{} fully\allowbreak{} faithful\allowbreak{} functor}{DER\_\allowbreak{}XVIII-c:\allowbreak{}87}
\DERapparatusitem{ENV1311}{108;\allowbreak{}\allowbreak{} ms.\allowbreak{}292,\allowbreak{}\allowbreak{} same\allowbreak{} exponent;\allowbreak{}\allowbreak{} mathematical\allowbreak{} reading}{2\textasciicircum{}\allowbreak{}\{\allowbreak{}(\allowbreak{}\textbackslash{}\allowbreak{}pi\_\allowbreak{}\textbackslash{}\allowbreak{}alpha\textasciicircum{}\allowbreak{}\{\allowbreak{}\textbackslash{}\allowbreak{}pi\_\allowbreak{}0\}\allowbreak{})\allowbreak{}\}\allowbreak{}}{2\textasciicircum{}\allowbreak{}\{\allowbreak{}(\allowbreak{}\{\allowbreak{}\textbackslash{}\allowbreak{}pi'\_\allowbreak{}\textbackslash{}\allowbreak{}alpha\}\allowbreak{}\textasciicircum{}\allowbreak{}\{\allowbreak{}\textbackslash{}\allowbreak{}pi\_\allowbreak{}0\}\allowbreak{})\allowbreak{}\}\allowbreak{}}{DER\_\allowbreak{}XVIII-c:\allowbreak{}89}
\DERapparatusitem{ENV1312}{108;\allowbreak{}\allowbreak{} ms.\allowbreak{}292,\allowbreak{}\allowbreak{} final\allowbreak{} cardinal\allowbreak{} comparison}{\textbackslash{}\allowbreak{}pi\allowbreak{} \textbackslash{}\allowbreak{}geqslant\allowbreak{} 2\textasciicircum{}\allowbreak{}\{\allowbreak{}\textbackslash{}\allowbreak{}pi\_\allowbreak{}0\}\allowbreak{}\allowbreak{} >\allowbreak{} \textbackslash{}\allowbreak{}pi'\_\allowbreak{}0}{\textbackslash{}\allowbreak{}pi\allowbreak{} \textbackslash{}\allowbreak{}geqslant\allowbreak{} 2\textasciicircum{}\allowbreak{}\{\allowbreak{}\textbackslash{}\allowbreak{}pi'\_\allowbreak{}0\}\allowbreak{}\allowbreak{} >\allowbreak{} \textbackslash{}\allowbreak{}pi'\_\allowbreak{}0}{DER\_\allowbreak{}XVIII-c:\allowbreak{}90}
\DERapparatusitem{ENV1313}{108;\allowbreak{}\allowbreak{} running\allowbreak{} text\allowbreak{} or\allowbreak{} formula}{\textbackslash{}\allowbreak{}text\{\allowbreak{}def\}\allowbreak{}}{\textbackslash{}\allowbreak{}mathrm\{\allowbreak{}def\}\allowbreak{}}{TYPESETTING:\allowbreak{}7;\allowbreak{}\allowbreak{} NOTATION\_\allowbreak{}DER:\allowbreak{}17}
\DERapparatusitem{ENV1314}{108;\allowbreak{}\allowbreak{} running\allowbreak{} text\allowbreak{} or\allowbreak{} formula}{\textbackslash{}\allowbreak{}text\{\allowbreak{}def\}\allowbreak{}}{\textbackslash{}\allowbreak{}mathrm\{\allowbreak{}def\}\allowbreak{}}{TYPESETTING:\allowbreak{}7;\allowbreak{}\allowbreak{} NOTATION\_\allowbreak{}DER:\allowbreak{}17}
\DERapparatusitem{ENV1315}{109;\allowbreak{}\allowbreak{} ms.\allowbreak{}\allowbreak{} 293;\allowbreak{}\allowbreak{} flagged\allowbreak{} source\allowbreak{} slip}{x}{\textbackslash{}\allowbreak{}pi}{DER\_\allowbreak{}XVIII-c:\allowbreak{}91}
\DERapparatusitem{ENV1316}{109;\allowbreak{}\allowbreak{} running\allowbreak{} text\allowbreak{} or\allowbreak{} formula}{\textbackslash{}\allowbreak{}text\{\allowbreak{}def\}\allowbreak{}}{\textbackslash{}\allowbreak{}mathrm\{\allowbreak{}def\}\allowbreak{}}{TYPESETTING:\allowbreak{}7;\allowbreak{}\allowbreak{} NOTATION\_\allowbreak{}DER:\allowbreak{}17}
\DERapparatusitem{ENV1317}{109;\allowbreak{}\allowbreak{} ms.\allowbreak{}\allowbreak{} 294;\allowbreak{}\allowbreak{} flagged\allowbreak{} source\allowbreak{} slip}{\textbackslash{}\allowbreak{}pi}{\textbackslash{}\allowbreak{}pi/\allowbreak{}x}{DER\_\allowbreak{}XVIII-c:\allowbreak{}92}
\DERapparatusitem{ENV1318}{109;\allowbreak{}\allowbreak{} running\allowbreak{} text\allowbreak{} or\allowbreak{} formula}{\textbackslash{}\allowbreak{}text\{\allowbreak{}def\}\allowbreak{}}{\textbackslash{}\allowbreak{}mathrm\{\allowbreak{}def\}\allowbreak{}}{TYPESETTING:\allowbreak{}7;\allowbreak{}\allowbreak{} NOTATION\_\allowbreak{}DER:\allowbreak{}17}
\DERapparatusitem{ENV1319}{111;\allowbreak{}\allowbreak{} ms.\allowbreak{}\allowbreak{} 298;\allowbreak{}\allowbreak{} flagged\allowbreak{} source\allowbreak{} slip;\allowbreak{}\allowbreak{} colimit\allowbreak{} display}{0}{\textbackslash{}\allowbreak{}pi}{DER\_\allowbreak{}XVIII-c:\allowbreak{}96;\allowbreak{}\allowbreak{} DER\_\allowbreak{}XVIII-c:\allowbreak{}97}
\DERapparatusitem{ENV1320}{111;\allowbreak{}\allowbreak{} ms.\allowbreak{}\allowbreak{} 298;\allowbreak{}\allowbreak{} flagged\allowbreak{} source\allowbreak{} slip;\allowbreak{}\allowbreak{} colimit\allowbreak{} display}{0}{\textbackslash{}\allowbreak{}pi}{DER\_\allowbreak{}XVIII-c:\allowbreak{}96;\allowbreak{}\allowbreak{} DER\_\allowbreak{}XVIII-c:\allowbreak{}97}
\DERapparatusitem{ENV1321}{111;\allowbreak{}\allowbreak{} ms.\allowbreak{}\allowbreak{} 298;\allowbreak{}\allowbreak{} flagged\allowbreak{} source\allowbreak{} slip}{In\allowbreak{} conditions\allowbreak{} b)\allowbreak{}\allowbreak{} c)\allowbreak{}\allowbreak{} of\allowbreak{} 4.\allowbreak{}3.\allowbreak{}10}{In\allowbreak{} conditions\allowbreak{} b)\allowbreak{}\allowbreak{} c)\allowbreak{}\allowbreak{} of\allowbreak{} 4.\allowbreak{}13.\allowbreak{}10}{DER\_\allowbreak{}XVIII-c:\allowbreak{}98}
\DERapparatusitem{ENV1322}{112;\allowbreak{}\allowbreak{} running\allowbreak{} text\allowbreak{} or\allowbreak{} formula}{\textbackslash{}\allowbreak{}text\{\allowbreak{}def\}\allowbreak{}}{\textbackslash{}\allowbreak{}mathrm\{\allowbreak{}def\}\allowbreak{}}{TYPESETTING:\allowbreak{}7;\allowbreak{}\allowbreak{} NOTATION\_\allowbreak{}DER:\allowbreak{}17}
\DERapparatusitem{ENV1323}{112;\allowbreak{}\allowbreak{} ms.\allowbreak{}\allowbreak{} 300;\allowbreak{}\allowbreak{} flagged\allowbreak{} source\allowbreak{} slip}{i}{j}{DER\_\allowbreak{}XVIII-c:\allowbreak{}101}
\DERapparatusitem{ENV1324}{112;\allowbreak{}\allowbreak{} ms.\allowbreak{}301,\allowbreak{}\allowbreak{} after\allowbreak{} equalizing\allowbreak{} arrows}{and\allowbreak{} \$\allowbreak{}i\textbackslash{}\allowbreak{}longrightarrow\allowbreak{} i''\$\allowbreak{}}{and\allowbreak{} \$\allowbreak{}i'\textbackslash{}\allowbreak{}longrightarrow\allowbreak{} i''\$\allowbreak{}}{DER\_\allowbreak{}XVIII-c:\allowbreak{}106}
\DERapparatusitem{ENV1325}{113;\allowbreak{}\allowbreak{} ms.\allowbreak{}\allowbreak{} 302;\allowbreak{}\allowbreak{} flagged\allowbreak{} source\allowbreak{} slip}{none\allowbreak{} other\allowbreak{} than\allowbreak{} d)\allowbreak{}\allowbreak{} and\allowbreak{} e)\allowbreak{}\allowbreak{} of\allowbreak{} 4.\allowbreak{}3.\allowbreak{}10}{none\allowbreak{} other\allowbreak{} than\allowbreak{} d)\allowbreak{}\allowbreak{} and\allowbreak{} e)\allowbreak{}\allowbreak{} of\allowbreak{} 4.\allowbreak{}13.\allowbreak{}10}{DER\_\allowbreak{}XVIII-c:\allowbreak{}107}
\DERapparatusitem{ENV1326}{113;\allowbreak{}\allowbreak{} ms.\allowbreak{}302,\allowbreak{}\allowbreak{} reformulated4.\allowbreak{}13.\allowbreak{}8\allowbreak{} preliminary\allowbreak{} exponent}{\textbackslash{}\allowbreak{}pi\textasciicircum{}\allowbreak{}\{\allowbreak{}\textbackslash{}\allowbreak{}pi'\_\allowbreak{}0\}\allowbreak{}=\allowbreak{}\textbackslash{}\allowbreak{}pi\textbackslash{}\allowbreak{}geqslant\textbackslash{}\allowbreak{}pi''}{\textbackslash{}\allowbreak{}pi\textasciicircum{}\allowbreak{}\{\allowbreak{}\textbackslash{}\allowbreak{}pi'\}\allowbreak{}=\allowbreak{}\textbackslash{}\allowbreak{}pi\textbackslash{}\allowbreak{}geqslant\textbackslash{}\allowbreak{}pi''}{DER\_\allowbreak{}XVIII-c:\allowbreak{}109}
\DERapparatusitem{ENV1327}{113;\allowbreak{}\allowbreak{} ms.\allowbreak{}302–302bis,\allowbreak{}\allowbreak{} slices\allowbreak{} in\allowbreak{} reformulation}{}{(\allowbreak{}\textbackslash{}\allowbreak{}pi)\allowbreak{}}{DER\_\allowbreak{}XVIII-c:\allowbreak{}110}
\DERapparatusitem{ENV1328}{113;\allowbreak{}\allowbreak{} ms.\allowbreak{}302–302bis,\allowbreak{}\allowbreak{} slices\allowbreak{} in\allowbreak{} reformulation}{}{(\allowbreak{}\textbackslash{}\allowbreak{}pi)\allowbreak{}}{DER\_\allowbreak{}XVIII-c:\allowbreak{}110}
\DERapparatusitem{ENV1329}{114;\allowbreak{}\allowbreak{} ms.\allowbreak{}304,\allowbreak{}\allowbreak{} supremum\allowbreak{} notation}{i\_\allowbreak{}J}{i\_\allowbreak{}\{\allowbreak{}J'\}\allowbreak{}}{DER\_\allowbreak{}XVIII-c:\allowbreak{}111}
\DERapparatusitem{ENV1330}{114;\allowbreak{}\allowbreak{} ms.\allowbreak{}304,\allowbreak{}\allowbreak{} supremum\allowbreak{} notation}{i\_\allowbreak{}J}{i\_\allowbreak{}\{\allowbreak{}J'\}\allowbreak{}}{DER\_\allowbreak{}XVIII-c:\allowbreak{}111}
\DERapparatusitem{ENV1331}{114;\allowbreak{}\allowbreak{} running\allowbreak{} text\allowbreak{} or\allowbreak{} formula}{\textbackslash{}\allowbreak{}text\{\allowbreak{}def\}\allowbreak{}}{\textbackslash{}\allowbreak{}mathrm\{\allowbreak{}def\}\allowbreak{}}{TYPESETTING:\allowbreak{}7;\allowbreak{}\allowbreak{} NOTATION\_\allowbreak{}DER:\allowbreak{}17}
\DERapparatusitem{ENV1332}{114;\allowbreak{}\allowbreak{} ms.\allowbreak{}305,\allowbreak{}\allowbreak{} union\allowbreak{} subcategory}{I}{J}{DER\_\allowbreak{}XVIII-c:\allowbreak{}112}
\DERapparatusitem{ENV1333}{116;\allowbreak{}\allowbreak{} ms.\allowbreak{}\allowbreak{} 308;\allowbreak{}\allowbreak{} flagged\allowbreak{} source\allowbreak{} slip}{\$\allowbreak{}S(\allowbreak{}\textbackslash{}\allowbreak{}pi)\allowbreak{}\$\allowbreak{}}{\$\allowbreak{}\textbackslash{}\allowbreak{}underline\{\allowbreak{}S\}\allowbreak{}(\allowbreak{}\textbackslash{}\allowbreak{}pi)\allowbreak{}\$\allowbreak{}}{DER\_\allowbreak{}XVIII-c:\allowbreak{}113}
\DERapparatusitem{ENV1334}{116;\allowbreak{}\allowbreak{} ms.\allowbreak{}310,\allowbreak{}\allowbreak{} first\allowbreak{} sentence}{\textbackslash{}\allowbreak{}uS\_\allowbreak{}\textbackslash{}\allowbreak{}pi}{\textbackslash{}\allowbreak{}pi}{DER\_\allowbreak{}XVIII-c:\allowbreak{}115}
\DERapparatusitem{ENV1335}{116;\allowbreak{}\allowbreak{} running\allowbreak{} text\allowbreak{} or\allowbreak{} formula}{\textbackslash{}\allowbreak{}text\{\allowbreak{}def\}\allowbreak{}}{\textbackslash{}\allowbreak{}mathrm\{\allowbreak{}def\}\allowbreak{}}{TYPESETTING:\allowbreak{}7;\allowbreak{}\allowbreak{} NOTATION\_\allowbreak{}DER:\allowbreak{}17}
\DERapparatusitem{ENV1336}{116;\allowbreak{}\allowbreak{} ms.\allowbreak{}\allowbreak{} 310;\allowbreak{}\allowbreak{} flagged\allowbreak{} source\allowbreak{} slip}{\$\allowbreak{}S\_\allowbreak{}\{\allowbreak{}\textbackslash{}\allowbreak{}pi'\_\allowbreak{}0\}\allowbreak{}\$\allowbreak{}}{\$\allowbreak{}\textbackslash{}\allowbreak{}underline\{\allowbreak{}S\}\allowbreak{}\_\allowbreak{}\{\allowbreak{}\textbackslash{}\allowbreak{}pi'\_\allowbreak{}0\}\allowbreak{}\$\allowbreak{}}{DER\_\allowbreak{}XVIII-c:\allowbreak{}114}
\DERapparatusitem{ENV1337}{117;\allowbreak{}\allowbreak{} ms.\allowbreak{}311,\allowbreak{}\allowbreak{} parenthetical\allowbreak{} exponent}{\textbackslash{}\allowbreak{}pi'}{\textbackslash{}\allowbreak{}pi}{DER\_\allowbreak{}XVIII-c:\allowbreak{}116}
\DERapparatusitem{ENV1338}{117;\allowbreak{}\allowbreak{} ms.\allowbreak{}311,\allowbreak{}\allowbreak{} chosen\allowbreak{} adapted\allowbreak{} presentation}{\textbackslash{}\allowbreak{}pi}{\textbackslash{}\allowbreak{}pi'}{DER\_\allowbreak{}XVIII-c:\allowbreak{}117}
\DERapparatusitem{ENV1339}{117;\allowbreak{}\allowbreak{} ms.\allowbreak{}312,\allowbreak{}\allowbreak{} footnote138}{the\allowbreak{} stability\allowbreak{} hypothesis\allowbreak{} made\allowbreak{} in\allowbreak{} b)\allowbreak{}\allowbreak{} comes\allowbreak{} in,\allowbreak{}\allowbreak{} for\allowbreak{} by\allowbreak{} b)\allowbreak{}}{the\allowbreak{} stability\allowbreak{} hypothesis\allowbreak{} made\allowbreak{} in\allowbreak{} a\$\allowbreak{}'\$\allowbreak{})\allowbreak{}\allowbreak{} comes\allowbreak{} in,\allowbreak{}\allowbreak{} for\allowbreak{} by\allowbreak{} a\$\allowbreak{}'\$\allowbreak{})\allowbreak{}}{DER\_\allowbreak{}XVIII-c:\allowbreak{}118}
\DERapparatusitem{ENV1340}{118;\allowbreak{}\allowbreak{} running\allowbreak{} text\allowbreak{} or\allowbreak{} formula}{\textbackslash{}\allowbreak{}text\{\allowbreak{}def\}\allowbreak{}}{\textbackslash{}\allowbreak{}mathrm\{\allowbreak{}def\}\allowbreak{}}{TYPESETTING:\allowbreak{}7;\allowbreak{}\allowbreak{} NOTATION\_\allowbreak{}DER:\allowbreak{}17}
\DERapparatusitem{ENV1341}{118;\allowbreak{}\allowbreak{} ms.\allowbreak{}\allowbreak{} 314;\allowbreak{}\allowbreak{} flagged\allowbreak{} source\allowbreak{} slip}{i}{\textbackslash{}\allowbreak{}alpha}{DER\_\allowbreak{}XVIII-c:\allowbreak{}119}
\DERapparatusitem{ENV1342}{118;\allowbreak{}\allowbreak{} ms.\allowbreak{}314–315,\allowbreak{}\allowbreak{} colimit\allowbreak{} object\allowbreak{} used\allowbreak{} repeatedly}{\textbackslash{}\allowbreak{}varphi(\allowbreak{}i)\allowbreak{}}{\textbackslash{}\allowbreak{}varphi(\allowbreak{}i')\allowbreak{}}{DER\_\allowbreak{}XVIII-c:\allowbreak{}122}
\DERapparatusitem{ENV1343}{118;\allowbreak{}\allowbreak{} ms.\allowbreak{}\allowbreak{} 314;\allowbreak{}\allowbreak{} flagged\allowbreak{} source\allowbreak{} slip}{\textbackslash{}\allowbreak{}alpha}{\textbackslash{}\allowbreak{}beta}{DER\_\allowbreak{}XVIII-c:\allowbreak{}120}
\DERapparatusitem{ENV1344}{118;\allowbreak{}\allowbreak{} ms.\allowbreak{}314–315,\allowbreak{}\allowbreak{} colimit\allowbreak{} object\allowbreak{} used\allowbreak{} repeatedly}{\textbackslash{}\allowbreak{}varphi(\allowbreak{}i)\allowbreak{}}{\textbackslash{}\allowbreak{}varphi(\allowbreak{}i')\allowbreak{}}{DER\_\allowbreak{}XVIII-c:\allowbreak{}122}
\DERapparatusitem{ENV1345}{118;\allowbreak{}\allowbreak{} running\allowbreak{} text\allowbreak{} or\allowbreak{} formula}{\textbackslash{}\allowbreak{}text\{\allowbreak{}def\}\allowbreak{}}{\textbackslash{}\allowbreak{}mathrm\{\allowbreak{}def\}\allowbreak{}}{TYPESETTING:\allowbreak{}7;\allowbreak{}\allowbreak{} NOTATION\_\allowbreak{}DER:\allowbreak{}17}
\DERapparatusitem{ENV1346}{118;\allowbreak{}\allowbreak{} ms.\allowbreak{}314–315,\allowbreak{}\allowbreak{} colimit\allowbreak{} object\allowbreak{} used\allowbreak{} repeatedly}{\textbackslash{}\allowbreak{}varphi(\allowbreak{}i)\allowbreak{}}{\textbackslash{}\allowbreak{}varphi(\allowbreak{}i')\allowbreak{}}{DER\_\allowbreak{}XVIII-c:\allowbreak{}122}
\DERapparatusitem{ENV1347}{118;\allowbreak{}\allowbreak{} ms.\allowbreak{}\allowbreak{} 315;\allowbreak{}\allowbreak{} flagged\allowbreak{} source\allowbreak{} slip}{J}{j}{DER\_\allowbreak{}XVIII-c:\allowbreak{}121}
\DERapparatusitem{ENV1348}{118;\allowbreak{}\allowbreak{} ms.\allowbreak{}314–315,\allowbreak{}\allowbreak{} colimit\allowbreak{} object\allowbreak{} used\allowbreak{} repeatedly}{\textbackslash{}\allowbreak{}varphi(\allowbreak{}i)\allowbreak{}}{\textbackslash{}\allowbreak{}varphi(\allowbreak{}i')\allowbreak{}}{DER\_\allowbreak{}XVIII-c:\allowbreak{}122}
\DERapparatusitem{ENV1349}{118;\allowbreak{}\allowbreak{} ms.\allowbreak{}314–315,\allowbreak{}\allowbreak{} colimit\allowbreak{} object\allowbreak{} used\allowbreak{} repeatedly}{\$\allowbreak{}i\allowbreak{} \textbackslash{}\allowbreak{}in\allowbreak{} I'\$\allowbreak{}}{\$\allowbreak{}i'\allowbreak{} \textbackslash{}\allowbreak{}in\allowbreak{} I'\$\allowbreak{}}{DER\_\allowbreak{}XVIII-c:\allowbreak{}122}
\DERapparatusitem{ENV1350}{120;\allowbreak{}\allowbreak{} ms.\allowbreak{}317,\allowbreak{}\allowbreak{} parenthetical\allowbreak{} old\allowbreak{} bound}{(\allowbreak{}instead\allowbreak{} of\allowbreak{} \$\allowbreak{}\textbackslash{}\allowbreak{}beta'<\textbackslash{}\allowbreak{}alpha\$\allowbreak{})\allowbreak{}\allowbreak{} in\allowbreak{} \textbackslash{}\allowbreak{}circled\{\allowbreak{}1\}\allowbreak{}\allowbreak{} and\allowbreak{} \textbackslash{}\allowbreak{}circled\{\allowbreak{}2\}\allowbreak{}}{(\allowbreak{}instead\allowbreak{} of\allowbreak{} \$\allowbreak{}\textbackslash{}\allowbreak{}beta+1<\textbackslash{}\allowbreak{}alpha\$\allowbreak{})\allowbreak{}\allowbreak{} in\allowbreak{} \textbackslash{}\allowbreak{}circled\{\allowbreak{}1\}\allowbreak{}\allowbreak{} and\allowbreak{} \textbackslash{}\allowbreak{}circled\{\allowbreak{}2\}\allowbreak{}}{DER\_\allowbreak{}XVIII-c:\allowbreak{}125;\allowbreak{}\allowbreak{} DER\_\allowbreak{}XVIII-d:\allowbreak{}1}
\DERapparatusitem{ENV1351}{120;\allowbreak{}\allowbreak{} ms.\allowbreak{}\allowbreak{} 318;\allowbreak{}\allowbreak{} flagged\allowbreak{} source\allowbreak{} slip}{\textbackslash{}\allowbreak{}varphi(\allowbreak{}i\_\allowbreak{}\textbackslash{}\allowbreak{}alpha)\allowbreak{}\allowbreak{} \textbackslash{}\allowbreak{}arrow[\allowbreak{}r,\allowbreak{}\allowbreak{} "u\_\allowbreak{}\textbackslash{}\allowbreak{}alpha"\allowbreak{} \{\allowbreak{}draw,\allowbreak{}\allowbreak{} circle,\allowbreak{}\allowbreak{} inner\allowbreak{} sep=\allowbreak{}1pt\}\allowbreak{}]\allowbreak{}\allowbreak{} \&\allowbreak{}\allowbreak{} j\_\allowbreak{}\textbackslash{}\allowbreak{}beta}{\textbackslash{}\allowbreak{}varphi(\allowbreak{}i\_\allowbreak{}\textbackslash{}\allowbreak{}alpha)\allowbreak{}\allowbreak{} \textbackslash{}\allowbreak{}arrow[\allowbreak{}r,\allowbreak{}\allowbreak{} "u\_\allowbreak{}\textbackslash{}\allowbreak{}alpha"\allowbreak{} \{\allowbreak{}draw,\allowbreak{}\allowbreak{} circle,\allowbreak{}\allowbreak{} inner\allowbreak{} sep=\allowbreak{}1pt\}\allowbreak{}]\allowbreak{}\allowbreak{} \&\allowbreak{}\allowbreak{} j\_\allowbreak{}\textbackslash{}\allowbreak{}alpha}{DER\_\allowbreak{}XVIII-c:\allowbreak{}124;\allowbreak{}\allowbreak{} DER\_\allowbreak{}XVIII-d:\allowbreak{}3}
\DERapparatusitem{ENV1352}{122;\allowbreak{}\allowbreak{} ms.\allowbreak{}\allowbreak{} 323;\allowbreak{}\allowbreak{} flagged\allowbreak{} source\allowbreak{} slip}{:\allowbreak{}\textbackslash{}\allowbreak{}varphi(\allowbreak{}i\_\allowbreak{}\textbackslash{}\allowbreak{}alpha)\allowbreak{}\textbackslash{}\allowbreak{}rightrightarrows}{:\allowbreak{}\textbackslash{}\allowbreak{}varphi(\allowbreak{}i\_\allowbreak{}\{\allowbreak{}\textbackslash{}\allowbreak{}alpha\_\allowbreak{}0\}\allowbreak{})\allowbreak{}\textbackslash{}\allowbreak{}rightrightarrows}{DER\_\allowbreak{}XVIII-c:\allowbreak{}127;\allowbreak{}\allowbreak{} DER\_\allowbreak{}XVIII-d:\allowbreak{}5}
\DERapparatusitem{ENV1353}{122;\allowbreak{}\allowbreak{} ms.\allowbreak{}\allowbreak{} 324;\allowbreak{}\allowbreak{} flagged\allowbreak{} source\allowbreak{} slip}{0}{\textbackslash{}\allowbreak{}alpha}{DER\_\allowbreak{}XVIII-c:\allowbreak{}128;\allowbreak{}\allowbreak{} DER\_\allowbreak{}XVIII-d:\allowbreak{}6}
\DERapparatusitem{ENV1354}{122;\allowbreak{}\allowbreak{} ms.\allowbreak{}\allowbreak{} 324;\allowbreak{}\allowbreak{} flagged\allowbreak{} source\allowbreak{} slip}{\textbackslash{}\allowbreak{}alpha\_\allowbreak{}0,\allowbreak{}\textbackslash{}\allowbreak{}alpha}{\textbackslash{}\allowbreak{}alpha,\allowbreak{}\textbackslash{}\allowbreak{}alpha\_\allowbreak{}0}{DER\_\allowbreak{}XVIII-c:\allowbreak{}129;\allowbreak{}\allowbreak{} DER\_\allowbreak{}XVIII-d:\allowbreak{}7}
\DERapparatusitem{ENV1355}{123;\allowbreak{}\allowbreak{} ms.\allowbreak{}327,\allowbreak{}\allowbreak{} first\allowbreak{} limit-stage\allowbreak{} transition\allowbreak{} indices}{\textbackslash{}\allowbreak{}beta,\allowbreak{}\textbackslash{}\allowbreak{}alpha}{\textbackslash{}\allowbreak{}alpha,\allowbreak{}\textbackslash{}\allowbreak{}beta}{DER\_\allowbreak{}XVIII-c:\allowbreak{}134;\allowbreak{}\allowbreak{} DER\_\allowbreak{}XVIII-d:\allowbreak{}8}
\DERapparatusitem{ENV1356}{123;\allowbreak{}\allowbreak{} ms.\allowbreak{}\allowbreak{} 327;\allowbreak{}\allowbreak{} flagged\allowbreak{} source\allowbreak{} slip}{J}{j}{DER\_\allowbreak{}XVIII-c:\allowbreak{}131;\allowbreak{}\allowbreak{} DER\_\allowbreak{}XVIII-d:\allowbreak{}9}
\DERapparatusitem{ENV1357}{123;\allowbreak{}\allowbreak{} ms.\allowbreak{}327,\allowbreak{}\allowbreak{} second\allowbreak{} limit-stage\allowbreak{} transition\allowbreak{} indices}{\textbackslash{}\allowbreak{}beta,\allowbreak{}\textbackslash{}\allowbreak{}alpha}{\textbackslash{}\allowbreak{}alpha,\allowbreak{}\textbackslash{}\allowbreak{}beta}{DER\_\allowbreak{}XVIII-c:\allowbreak{}135;\allowbreak{}\allowbreak{} DER\_\allowbreak{}XVIII-d:\allowbreak{}10}
\DERapparatusitem{ENV1358}{123;\allowbreak{}\allowbreak{} ms.\allowbreak{}\allowbreak{} 328;\allowbreak{}\allowbreak{} flagged\allowbreak{} source\allowbreak{} slip}{J}{j}{DER\_\allowbreak{}XVIII-c:\allowbreak{}132;\allowbreak{}\allowbreak{} DER\_\allowbreak{}XVIII-d:\allowbreak{}11}
\DERapparatusitem{ENV1359}{123;\allowbreak{}\allowbreak{} ms.\allowbreak{}\allowbreak{} 328;\allowbreak{}\allowbreak{} flagged\allowbreak{} source\allowbreak{} slip}{j\_\allowbreak{}\textbackslash{}\allowbreak{}alpha}{\textbackslash{}\allowbreak{}bar\{\allowbreak{}j\}\allowbreak{}\_\allowbreak{}\textbackslash{}\allowbreak{}alpha}{DER\_\allowbreak{}XVIII-c:\allowbreak{}133;\allowbreak{}\allowbreak{} DER\_\allowbreak{}XVIII-d:\allowbreak{}12}
\DERapparatusitem{ENV1360}{124;\allowbreak{}\allowbreak{} ms330;\allowbreak{}\allowbreak{} sentence\allowbreak{} after\allowbreak{} 4.\allowbreak{}13.\allowbreak{}17}{only\allowbreak{} that\allowbreak{} it\allowbreak{} was\allowbreak{} a\allowbreak{} limit\allowbreak{} cardinal}{only\allowbreak{} that\allowbreak{} it\allowbreak{} was\allowbreak{} a\allowbreak{} limit\allowbreak{} ordinal}{DER\_\allowbreak{}XVIII-d:\allowbreak{}13}
\DERapparatusitem{ENV1361}{125;\allowbreak{}\allowbreak{} ms331;\allowbreak{}\allowbreak{} first\allowbreak{} sentence}{the\allowbreak{} implication\allowbreak{} b\allowbreak{} \$\allowbreak{}\textbackslash{}\allowbreak{}Longrightarrow\$\allowbreak{}\allowbreak{} d}{the\allowbreak{} implication\allowbreak{} a\$\allowbreak{}'\$\allowbreak{}\allowbreak{} \$\allowbreak{}\textbackslash{}\allowbreak{}Longrightarrow\$\allowbreak{}\allowbreak{} c}{DER\_\allowbreak{}XVIII-d:\allowbreak{}14}
\DERapparatusitem{ENV1362}{125;\allowbreak{}\allowbreak{} ms332;\allowbreak{}\allowbreak{} equation\allowbreak{} (\allowbreak{}**)\allowbreak{}}{}{(\allowbreak{}\textbackslash{}\allowbreak{}pi)\allowbreak{}}{DER\_\allowbreak{}XVIII-d:\allowbreak{}15}
\DERapparatusitem{ENV1363}{126;\allowbreak{}\allowbreak{} ms333;\allowbreak{}\allowbreak{} target\allowbreak{} of\allowbreak{} the\allowbreak{} functor\allowbreak{} defining\allowbreak{} inverse\allowbreak{} image}{\textbackslash{}\allowbreak{}pi}{\textbackslash{}\allowbreak{}pi/\allowbreak{}x}{DER\_\allowbreak{}XVIII-d:\allowbreak{}16}
\DERapparatusitem{ENV1364}{127;\allowbreak{}\allowbreak{} running\allowbreak{} text\allowbreak{} or\allowbreak{} formula}{\textbackslash{}\allowbreak{}text\{\allowbreak{}def\}\allowbreak{}}{\textbackslash{}\allowbreak{}mathrm\{\allowbreak{}def\}\allowbreak{}}{TYPESETTING:\allowbreak{}7;\allowbreak{}\allowbreak{} NOTATION\_\allowbreak{}DER:\allowbreak{}17}
\DERapparatusitem{ENV1365}{129;\allowbreak{}\allowbreak{} ms340;\allowbreak{}\allowbreak{} opening\allowbreak{} of\allowbreak{} 4.\allowbreak{}13.\allowbreak{}19}{Let\allowbreak{} \$\allowbreak{}\textbackslash{}\allowbreak{}mathcal\{\allowbreak{}M\}\allowbreak{}\$\allowbreak{}\allowbreak{} be\allowbreak{} an\allowbreak{} accessible\allowbreak{} category\allowbreak{} in\allowbreak{} \$\allowbreak{}\textbackslash{}\allowbreak{}mathcal\{\allowbreak{}M\}\allowbreak{}\$\allowbreak{}}{Let\allowbreak{} \$\allowbreak{}\textbackslash{}\allowbreak{}mathcal\{\allowbreak{}M\}\allowbreak{}\$\allowbreak{}\allowbreak{} be\allowbreak{} an\allowbreak{} accessible\allowbreak{} category,\allowbreak{}\allowbreak{} and\allowbreak{} \$\allowbreak{}\textbackslash{}\allowbreak{}underline\{\allowbreak{}S\}\allowbreak{}\$\allowbreak{}\allowbreak{} an\allowbreak{} accessible\allowbreak{} subcategory\allowbreak{} in\allowbreak{} \$\allowbreak{}\textbackslash{}\allowbreak{}mathcal\{\allowbreak{}M\}\allowbreak{}\$\allowbreak{}}{DER\_\allowbreak{}XVIII-d:\allowbreak{}22}
\DERapparatusitem{ENV1366}{129;\allowbreak{}\allowbreak{} ms340;\allowbreak{}\allowbreak{} ordinal-colimit\allowbreak{} hypothesis\allowbreak{} in\allowbreak{} 4.\allowbreak{}13.\allowbreak{}19}{are\allowbreak{} stable\allowbreak{} under\allowbreak{} \$\allowbreak{}\textbackslash{}\allowbreak{}varinjlim\$\allowbreak{}\allowbreak{} of\allowbreak{} type\allowbreak{} \$\allowbreak{}I\_\allowbreak{}\textbackslash{}\allowbreak{}alpha\$\allowbreak{}\allowbreak{} are\allowbreak{} stable\allowbreak{} under\allowbreak{} \$\allowbreak{}\textbackslash{}\allowbreak{}varinjlim\$\allowbreak{}\allowbreak{} of\allowbreak{} type\allowbreak{} \$\allowbreak{}I\_\allowbreak{}\textbackslash{}\allowbreak{}alpha\$\allowbreak{}}{are\allowbreak{} stable\allowbreak{} under\allowbreak{} \$\allowbreak{}\textbackslash{}\allowbreak{}varinjlim\$\allowbreak{}\allowbreak{} of\allowbreak{} type\allowbreak{} \$\allowbreak{}I\_\allowbreak{}\textbackslash{}\allowbreak{}alpha\$\allowbreak{}}{DER\_\allowbreak{}XVIII-d:\allowbreak{}23}
\DERapparatusitem{ENV1367}{129;\allowbreak{}\allowbreak{} running\allowbreak{} text\allowbreak{} or\allowbreak{} formula}{\textbackslash{}\allowbreak{}textup\{\allowbreak{}def\}\allowbreak{}}{\textbackslash{}\allowbreak{}mathrm\{\allowbreak{}def\}\allowbreak{}}{TYPESETTING:\allowbreak{}7;\allowbreak{}\allowbreak{} NOTATION\_\allowbreak{}DER:\allowbreak{}17}
\DERapparatusitem{ENV1368}{129;\allowbreak{}\allowbreak{} running\allowbreak{} text\allowbreak{} or\allowbreak{} formula}{\textbackslash{}\allowbreak{}textup\{\allowbreak{}def\}\allowbreak{}}{\textbackslash{}\allowbreak{}mathrm\{\allowbreak{}def\}\allowbreak{}}{TYPESETTING:\allowbreak{}7;\allowbreak{}\allowbreak{} NOTATION\_\allowbreak{}DER:\allowbreak{}17}
\DERapparatusitem{ENV1369}{129;\allowbreak{}\allowbreak{} ms342;\allowbreak{}\allowbreak{} naming\allowbreak{} the\allowbreak{} closure\allowbreak{} condition}{\textbackslash{}\allowbreak{}varphi}{\textbackslash{}\allowbreak{}Phi}{DER\_\allowbreak{}XVIII-d:\allowbreak{}24}
\DERapparatusitem{ENV1370}{130;\allowbreak{}\allowbreak{} ms342\allowbreak{} continuation;\allowbreak{}\allowbreak{} first\allowbreak{} line}{(\allowbreak{}where\allowbreak{} what\allowbreak{} amounts\allowbreak{} to\allowbreak{} the\allowbreak{} same\allowbreak{} thing,\allowbreak{}}{(\allowbreak{}or\allowbreak{} what\allowbreak{} amounts\allowbreak{} to\allowbreak{} the\allowbreak{} same\allowbreak{} thing,\allowbreak{}}{DER\_\allowbreak{}XVIII-d:\allowbreak{}25}
\DERapparatusitem{ENV1371}{131;\allowbreak{}\allowbreak{} ms346,\allowbreak{}\allowbreak{} Proposition\allowbreak{} 4.\allowbreak{}13.\allowbreak{}25,\allowbreak{}\allowbreak{} parenthetical\allowbreak{} ordinal\allowbreak{} definition}{Ph}{p}{DER\_\allowbreak{}XVIII-d:\allowbreak{}26}
\DERapparatusitem{ENV1372}{132;\allowbreak{}\allowbreak{} ms348,\allowbreak{}\allowbreak{} Proposition\allowbreak{} 4.\allowbreak{}13.\allowbreak{}27\allowbreak{} hypothesis}{which\allowbreak{} is\allowbreak{} \$\allowbreak{}\textbackslash{}\allowbreak{}Phi\$\allowbreak{}-\$\allowbreak{}\textbackslash{}\allowbreak{}pi\$\allowbreak{}-adapted.\allowbreak{}}{which\allowbreak{} is\allowbreak{} \$\allowbreak{}\textbackslash{}\allowbreak{}Phi\$\allowbreak{}-\$\allowbreak{}\textbackslash{}\allowbreak{}pi\$\allowbreak{}-accessible.\allowbreak{}}{DER\_\allowbreak{}XVIII-d:\allowbreak{}27}
\DERapparatusitem{ENV1373}{132;\allowbreak{}\allowbreak{} ms349,\allowbreak{}\allowbreak{} sentence\allowbreak{} before\allowbreak{} Lemma\allowbreak{} 4.\allowbreak{}13.\allowbreak{}30}{(\allowbreak{}I\allowbreak{} do\allowbreak{} not\allowbreak{} spell\allowbreak{} out\allowbreak{} by\allowbreak{} the\allowbreak{} proof)\allowbreak{}}{(\allowbreak{}I\allowbreak{} do\allowbreak{} not\allowbreak{} spell\allowbreak{} out\allowbreak{} the\allowbreak{} proof)\allowbreak{}}{DER\_\allowbreak{}XVIII-d:\allowbreak{}28}
\DERapparatusitem{ENV1374}{133;\allowbreak{}\allowbreak{} ms351,\allowbreak{}\allowbreak{} Corollary\allowbreak{} 4.\allowbreak{}13.\allowbreak{}32\allowbreak{} conclusion}{N}{M}{DER\_\allowbreak{}XVIII-d:\allowbreak{}29}
\DERapparatusitem{ENV1375}{preamble;\allowbreak{}\allowbreak{} local\allowbreak{} macro\allowbreak{} definition}{\textbackslash{}\allowbreak{}text\{\allowbreak{}def\}\allowbreak{}}{\textbackslash{}\allowbreak{}mathrm\{\allowbreak{}def\}\allowbreak{}}{TYPESETTING:\allowbreak{}7;\allowbreak{}\allowbreak{} NOTATION\_\allowbreak{}DER:\allowbreak{}17}
\DERapparatusitem{ENV1376}{134;\allowbreak{}\allowbreak{} ms352,\allowbreak{}\allowbreak{} condition\allowbreak{} 2\allowbreak{} of\allowbreak{} adapted\allowbreak{} presentation}{large\allowbreak{} relative\allowbreak{} to\allowbreak{} \$\allowbreak{}\textbackslash{}\allowbreak{}pi\_\allowbreak{}0\$\allowbreak{}\allowbreak{} is\allowbreak{} of\allowbreak{} cardinality}{large\allowbreak{} relative\allowbreak{} to\allowbreak{} \$\allowbreak{}\textbackslash{}\allowbreak{}pi\_\allowbreak{}0\$\allowbreak{}\allowbreak{} and\allowbreak{} of\allowbreak{} cardinality}{DER\_\allowbreak{}XVIII-d:\allowbreak{}30}
\DERapparatusitem{ENV1377}{134;\allowbreak{}\allowbreak{} ms352,\allowbreak{}\allowbreak{} sentence\allowbreak{} beginning\allowbreak{} Comme\allowbreak{} f\allowbreak{} commute}{(\allowbreak{}cf.\allowbreak{}\textbackslash{}\allowbreak{},\allowbreak{}\textbackslash{}\allowbreak{}cercle\{\allowbreak{}1\}\allowbreak{})\allowbreak{}}{(\allowbreak{}cf.\allowbreak{}\textbackslash{}\allowbreak{},\allowbreak{}\textbackslash{}\allowbreak{}cercle\{\allowbreak{}d\}\allowbreak{})\allowbreak{}}{DER\_\allowbreak{}XVIII-d:\allowbreak{}31}
\DERapparatusitem{ENV1378}{134;\allowbreak{}\allowbreak{} ms353,\allowbreak{}\allowbreak{} tuple\allowbreak{} for\allowbreak{} applying\allowbreak{} 4.\allowbreak{}13.\allowbreak{}22}{S}{T}{DER\_\allowbreak{}XVIII-d:\allowbreak{}32}
\DERapparatusitem{ENV1379}{134;\allowbreak{}\allowbreak{} ms353,\allowbreak{}\allowbreak{} indices\allowbreak{} selected\allowbreak{} after\allowbreak{} 4.\allowbreak{}13.\allowbreak{}22}{S}{T}{DER\_\allowbreak{}XVIII-d:\allowbreak{}33}
\DERapparatusitem{ENV1380}{134;\allowbreak{}\allowbreak{} ms354,\allowbreak{}\allowbreak{} second\allowbreak{} conjecture}{a\allowbreak{} small\allowbreak{} full\allowbreak{} subcategories.\allowbreak{}}{a\allowbreak{} small\allowbreak{} full\allowbreak{} subcategory.\allowbreak{}}{DER\_\allowbreak{}XVIII-d:\allowbreak{}35}
\DERapparatusitem{ENV1381}{136;\allowbreak{}\allowbreak{} ms359,\allowbreak{}\allowbreak{} Proposition\allowbreak{} 4.\allowbreak{}13.\allowbreak{}36\allowbreak{} c-double-prime}{\textbackslash{}\allowbreak{}textit\{\allowbreak{}\$\allowbreak{}\textbackslash{}\allowbreak{}pi\$\allowbreak{}-adapted\allowbreak{} representation}{\textbackslash{}\allowbreak{}textit\{\allowbreak{}L-\$\allowbreak{}\textbackslash{}\allowbreak{}pi\$\allowbreak{}-adapted\allowbreak{} representation}{DER\_\allowbreak{}XVIII-d:\allowbreak{}37}
\DERapparatusitem{ENV1382}{136;\allowbreak{}\allowbreak{} ms360,\allowbreak{}\allowbreak{} final\allowbreak{} item\allowbreak{} of\allowbreak{} Lemma\allowbreak{} 4.\allowbreak{}13.\allowbreak{}37}{\textbackslash{}\allowbreak{}item[\allowbreak{}\textbackslash{}\allowbreak{}textit\{\allowbreak{}d)\allowbreak{}\}\allowbreak{}]\allowbreak{}\allowbreak{} \textbackslash{}\allowbreak{}textit\{\allowbreak{}\$\allowbreak{}x\_\allowbreak{}i\allowbreak{} \textbackslash{}\allowbreak{}in}{\textbackslash{}\allowbreak{}item[\allowbreak{}\textbackslash{}\allowbreak{}textit\{\allowbreak{}4)\allowbreak{}\}\allowbreak{}]\allowbreak{}\allowbreak{} \textbackslash{}\allowbreak{}textit\{\allowbreak{}\$\allowbreak{}x\_\allowbreak{}i\allowbreak{} \textbackslash{}\allowbreak{}in}{DER\_\allowbreak{}XVIII-d:\allowbreak{}39}
\DERapparatusitem{ENV1383}{138;\allowbreak{}\allowbreak{} ms365,\allowbreak{}\allowbreak{} normal\allowbreak{} form\allowbreak{} parenthesis}{and\allowbreak{} \$\allowbreak{}u\$\allowbreak{},\allowbreak{}\allowbreak{} \$\allowbreak{}v\$\allowbreak{}\allowbreak{} arrows\allowbreak{} in\allowbreak{} \$\allowbreak{}K'\$\allowbreak{}}{and\allowbreak{} \$\allowbreak{}u\$\allowbreak{},\allowbreak{}\allowbreak{} \$\allowbreak{}v\$\allowbreak{}\allowbreak{} arrows\allowbreak{} in\allowbreak{} \$\allowbreak{}K\$\allowbreak{}}{DER\_\allowbreak{}XVIII-d:\allowbreak{}40}
\DERapparatusitem{ENV1384}{138;\allowbreak{}\allowbreak{} ms365,\allowbreak{}\allowbreak{} normal\allowbreak{} form\allowbreak{} simplification}{\$\allowbreak{}pv\allowbreak{} =\allowbreak{}\allowbreak{} pu\$\allowbreak{}}{\$\allowbreak{}pv\allowbreak{} =\allowbreak{}\allowbreak{} p\$\allowbreak{}}{DER\_\allowbreak{}XVIII-d:\allowbreak{}41}
\DERapparatusitem{ENV1385}{138;\allowbreak{}\allowbreak{} ms366,\allowbreak{}\allowbreak{} recursive\allowbreak{} sequence\allowbreak{} of\allowbreak{} subcategories}{J}{K}{DER\_\allowbreak{}XVIII-d:\allowbreak{}42}
\DERapparatusitem{ENV1386}{139;\allowbreak{}\allowbreak{} ms367,\allowbreak{}\allowbreak{} Corollary\allowbreak{} 4.\allowbreak{}13.\allowbreak{}39,\allowbreak{}\allowbreak{} arrows\allowbreak{} in\allowbreak{} (\allowbreak{}*)\allowbreak{}}{J}{K}{DER\_\allowbreak{}XVIII-d:\allowbreak{}43}
\DERapparatusitem{ENV1387}{139;\allowbreak{}\allowbreak{} ms367,\allowbreak{}\allowbreak{} Corollary\allowbreak{} 4.\allowbreak{}13.\allowbreak{}39\allowbreak{} object\allowbreak{} set}{J}{K}{DER\_\allowbreak{}XVIII-d:\allowbreak{}44}
\DERapparatusitem{ENV1388}{139;\allowbreak{}\allowbreak{} ms367,\allowbreak{}\allowbreak{} pairs\allowbreak{} used\allowbreak{} to\allowbreak{} construct\allowbreak{} p\_\allowbreak{}i}{J}{K}{DER\_\allowbreak{}XVIII-d:\allowbreak{}45}
\DERapparatusitem{ENV1389}{140;\allowbreak{}\allowbreak{} ms370,\allowbreak{}\allowbreak{} display\allowbreak{} (\allowbreak{}4.\allowbreak{}14.\allowbreak{}1.\allowbreak{}1)\allowbreak{},\allowbreak{}\allowbreak{} phi\allowbreak{} domain}{}{\_\allowbreak{}\{\allowbreak{}\textbackslash{}\allowbreak{}bullet\}\allowbreak{}}{DER\_\allowbreak{}XVIII-d:\allowbreak{}47}
\DERapparatusitem{ENV1390}{140;\allowbreak{}\allowbreak{} ms370,\allowbreak{}\allowbreak{} display\allowbreak{} (\allowbreak{}4.\allowbreak{}14.\allowbreak{}1.\allowbreak{}1)\allowbreak{},\allowbreak{}\allowbreak{} psi\allowbreak{} domain}{\textbackslash{}\allowbreak{}psi\textbackslash{}\allowbreak{}\allowbreak{} :\allowbreak{}\textbackslash{}\allowbreak{}\allowbreak{} I\allowbreak{} \textbackslash{}\allowbreak{}longrightarrow\allowbreak{} \textbackslash{}\allowbreak{}mathcal\{\allowbreak{}C\}\allowbreak{}}{\textbackslash{}\allowbreak{}psi\textbackslash{}\allowbreak{}\allowbreak{} :\allowbreak{}\textbackslash{}\allowbreak{}\allowbreak{} \{\allowbreak{}\}\allowbreak{}\_\allowbreak{}\{\allowbreak{}\textbackslash{}\allowbreak{}bullet\}\allowbreak{}I\allowbreak{} \textbackslash{}\allowbreak{}longrightarrow\allowbreak{} \textbackslash{}\allowbreak{}mathcal\{\allowbreak{}C\}\allowbreak{}}{DER\_\allowbreak{}XVIII-d:\allowbreak{}48}
\DERapparatusitem{ENV1391}{143;\allowbreak{}\allowbreak{} ms377;\allowbreak{}\allowbreak{} proof\allowbreak{} of\allowbreak{} 4.\allowbreak{}14.\allowbreak{}6}{0}{1}{DER\_\allowbreak{}XVIII-d:\allowbreak{}49}
\DERapparatusitem{ENV1392}{143;\allowbreak{}\allowbreak{} ms378;\allowbreak{}\allowbreak{} proof\allowbreak{} of\allowbreak{} 4.\allowbreak{}14.\allowbreak{}8}{\textbackslash{}\allowbreak{}pi}{1}{DER\_\allowbreak{}XVIII-d:\allowbreak{}50}
\DERapparatusitem{ENV1393}{preamble;\allowbreak{}\allowbreak{} local\allowbreak{} macro\allowbreak{} definition}{\textbackslash{}\allowbreak{}textup\{\allowbreak{}def\}\allowbreak{}}{\textbackslash{}\allowbreak{}mathrm\{\allowbreak{}def\}\allowbreak{}}{TYPESETTING:\allowbreak{}7;\allowbreak{}\allowbreak{} NOTATION\_\allowbreak{}DER:\allowbreak{}17}
\DERapparatusitem{ENV1394}{144;\allowbreak{}\allowbreak{} ms380;\allowbreak{}\allowbreak{} paragraph\allowbreak{} after\allowbreak{} 4.\allowbreak{}14.\allowbreak{}11}{In\allowbreak{} the\allowbreak{} cases\allowbreak{} of\allowbreak{} a\allowbreak{} covariant\allowbreak{} functor}{In\allowbreak{} the\allowbreak{} case\allowbreak{} of\allowbreak{} a\allowbreak{} covariant\allowbreak{} functor}{DER\_\allowbreak{}XVIII-d:\allowbreak{}52}
\DERapparatusitem{ENV1395}{144;\allowbreak{}\allowbreak{} ms380;\allowbreak{}\allowbreak{} final\allowbreak{} SGA\allowbreak{} cross-reference}{cf.\allowbreak{}\allowbreak{} loc.\allowbreak{}\allowbreak{} cit.\allowbreak{}\allowbreak{} 8.\allowbreak{}3.\allowbreak{}12}{cf.\allowbreak{}\allowbreak{} loc.\allowbreak{}\allowbreak{} cit.\allowbreak{}\allowbreak{} 8.\allowbreak{}13.\allowbreak{}2}{DER\_\allowbreak{}XVIII-d:\allowbreak{}53}
\DERapparatusitem{ENV1396}{144;\allowbreak{}\allowbreak{} ms382;\allowbreak{}\allowbreak{} colimit\allowbreak{} functor\allowbreak{} display}{}{\textbackslash{}\allowbreak{}mathcal\allowbreak{} }{DER\_\allowbreak{}XVIII-d:\allowbreak{}54}
\DERapparatusitem{ENV1397}{145;\allowbreak{}\allowbreak{} ms383;\allowbreak{}\allowbreak{} inverse-image\allowbreak{} citation}{(\allowbreak{}4.\allowbreak{}3.\allowbreak{}31,\allowbreak{}\allowbreak{} page\allowbreak{} 350)\allowbreak{}}{(\allowbreak{}4.\allowbreak{}13.\allowbreak{}31,\allowbreak{}\allowbreak{} page\allowbreak{} 350)\allowbreak{}}{DER\_\allowbreak{}XVIII-d:\allowbreak{}55}
\DERapparatusitem{ENV1398}{145;\allowbreak{}\allowbreak{} ms383;\allowbreak{}\allowbreak{} small\allowbreak{} indexing\allowbreak{} family}{\textbackslash{}\allowbreak{}Phi}{\textbackslash{}\allowbreak{}Lambda\_\allowbreak{}1}{DER\_\allowbreak{}XVIII-d:\allowbreak{}56}
\DERapparatusitem{ENV1399}{145;\allowbreak{}\allowbreak{} ms383;\allowbreak{}\allowbreak{} second\allowbreak{} proof\allowbreak{} opening}{\textbackslash{}\allowbreak{}mathrm\{\allowbreak{}Hom\}\allowbreak{}}{\textbackslash{}\allowbreak{}uHom}{DER\_\allowbreak{}XVIII-d:\allowbreak{}57}
\DERapparatusitem{ENV1400}{145;\allowbreak{}\allowbreak{} ms383;\allowbreak{}\allowbreak{} limit\allowbreak{} functor\allowbreak{} display}{\textbackslash{}\allowbreak{}varinjlim\_\allowbreak{}\{\allowbreak{}I\}\allowbreak{}}{\textbackslash{}\allowbreak{}varprojlim\_\allowbreak{}I:\allowbreak{}}{DER\_\allowbreak{}XVIII-d:\allowbreak{}58}
\DERapparatusitem{ENV1401}{145;\allowbreak{}\allowbreak{} ms383;\allowbreak{}\allowbreak{} sentence\allowbreak{} following\allowbreak{} limit\allowbreak{} display}{\textbackslash{}\allowbreak{}pi}{\textbackslash{}\allowbreak{}pi\_\allowbreak{}0}{DER\_\allowbreak{}XVIII-d:\allowbreak{}59}
\DERapparatusitem{ENV1402}{146;\allowbreak{}\allowbreak{} ms385;\allowbreak{}\allowbreak{} induced\allowbreak{} arrow\allowbreak{} functor\allowbreak{} display}{\textbackslash{}\allowbreak{}xi}{x}{DER\_\allowbreak{}XVIII-d:\allowbreak{}60}
\DERapparatusitem{ENV1403}{146;\allowbreak{}\allowbreak{} ms385;\allowbreak{}\allowbreak{} later\allowbreak{} reference\allowbreak{} to\allowbreak{} induced\allowbreak{} functors}{\textbackslash{}\allowbreak{}xi}{x}{DER\_\allowbreak{}XVIII-d:\allowbreak{}61}
\DERapparatusitem{ENV1404}{146;\allowbreak{}\allowbreak{} ms387;\allowbreak{}\allowbreak{} question\allowbreak{} 4.\allowbreak{}14.\allowbreak{}15,\allowbreak{}\allowbreak{} ambient\allowbreak{} category}{\textbackslash{}\allowbreak{}mathcal\{\allowbreak{}N\}\allowbreak{}}{\textbackslash{}\allowbreak{}uFl(\allowbreak{}\textbackslash{}\allowbreak{}mathcal\allowbreak{} N)\allowbreak{}}{DER\_\allowbreak{}XVIII-d:\allowbreak{}62}
\DERapparatusitem{ENV1405}{146;\allowbreak{}\allowbreak{} ms387;\allowbreak{}\allowbreak{} heading\allowbreak{} 4.\allowbreak{}14.\allowbreak{}15.\allowbreak{}1}{(\allowbreak{}of\allowbreak{} 4.\allowbreak{}14.\allowbreak{}13)\allowbreak{}}{(\allowbreak{}of\allowbreak{} 4.\allowbreak{}14.\allowbreak{}12)\allowbreak{}}{DER\_\allowbreak{}XVIII-d:\allowbreak{}64}
\DERapparatusitem{ENV1406}{147;\allowbreak{}\allowbreak{} ms388;\allowbreak{}\allowbreak{} last\allowbreak{} sentence\allowbreak{} of\allowbreak{} pushout\allowbreak{} argument}{by\allowbreak{} 4.\allowbreak{}14.\allowbreak{}13\allowbreak{} the\allowbreak{} latter}{by\allowbreak{} 4.\allowbreak{}14.\allowbreak{}12\allowbreak{} the\allowbreak{} latter}{DER\_\allowbreak{}XVIII-d:\allowbreak{}65}
\DERapparatusitem{ENV1407}{147;\allowbreak{}\allowbreak{} ms390;\allowbreak{}\allowbreak{} two\allowbreak{} shape\allowbreak{} diagrams}{}{\_\allowbreak{}\textbackslash{}\allowbreak{}bullet}{DER\_\allowbreak{}XVIII-d:\allowbreak{}66}
\DERapparatusitem{ENV1408}{148;\allowbreak{}\allowbreak{} ms391;\allowbreak{}\allowbreak{} 4.\allowbreak{}14.\allowbreak{}18(\allowbreak{}b)\allowbreak{}}{\textbackslash{}\allowbreak{}Ob}{\textbackslash{}\allowbreak{}Ob\textbackslash{}\allowbreak{}mathcal}{DER\_\allowbreak{}XVIII-d:\allowbreak{}67;\allowbreak{}\allowbreak{} DER\_\allowbreak{}XVIII-d:\allowbreak{}69}
\DERapparatusitem{ENV1409}{148;\allowbreak{}\allowbreak{} ms391;\allowbreak{}\allowbreak{} 4.\allowbreak{}14.\allowbreak{}18(\allowbreak{}c)\allowbreak{}}{Suppose\allowbreak{} \$\allowbreak{}pi\_\allowbreak{}0\$\allowbreak{}\allowbreak{} chosen\allowbreak{} so\allowbreak{} as}{Suppose\allowbreak{} \$\allowbreak{}\textbackslash{}\allowbreak{}pi\_\allowbreak{}0\$\allowbreak{}\allowbreak{} chosen\allowbreak{} so\allowbreak{} as}{DER\_\allowbreak{}XVIII-d:\allowbreak{}68}
\DERapparatusitem{ENV1410}{148;\allowbreak{}\allowbreak{} ms392;\allowbreak{}\allowbreak{} bounded\allowbreak{} section\allowbreak{} equivalence}{\textbackslash{}\allowbreak{}uGamma(\allowbreak{}\textbackslash{}\allowbreak{}mathcal\{\allowbreak{}E\}\allowbreak{}(\allowbreak{}\textbackslash{}\allowbreak{}pi)\allowbreak{}/\allowbreak{}C)\allowbreak{}}{\textbackslash{}\allowbreak{}uGamma(\allowbreak{}\textbackslash{}\allowbreak{}mathcal\{\allowbreak{}E\}\allowbreak{}(\allowbreak{}\textbackslash{}\allowbreak{}pi)\allowbreak{}/\allowbreak{}\textbackslash{}\allowbreak{}mathcal\{\allowbreak{}C\}\allowbreak{})\allowbreak{}}{DER\_\allowbreak{}XVIII-d:\allowbreak{}70}
\DERapparatusitem{ENV1411}{148;\allowbreak{}\allowbreak{} ms392;\allowbreak{}\allowbreak{} parenthetical\allowbreak{} before\allowbreak{} Ind\allowbreak{} equivalence}{\$\allowbreak{}\textbackslash{}\allowbreak{}pi\$\allowbreak{}\allowbreak{} being\allowbreak{} useful\allowbreak{} for\allowbreak{} \$\allowbreak{}\textbackslash{}\allowbreak{}mathcal\{\allowbreak{}E\}\allowbreak{}\$\allowbreak{}}{\$\allowbreak{}\textbackslash{}\allowbreak{}pi\$\allowbreak{}\allowbreak{} being\allowbreak{} useful\allowbreak{} for\allowbreak{} \$\allowbreak{}\textbackslash{}\allowbreak{}mathcal\{\allowbreak{}M\}\allowbreak{}=\allowbreak{}\textbackslash{}\allowbreak{}uGamma(\allowbreak{}\textbackslash{}\allowbreak{}mathcal\{\allowbreak{}E\}\allowbreak{}/\allowbreak{}\textbackslash{}\allowbreak{}mathcal\{\allowbreak{}C\}\allowbreak{})\allowbreak{}\$\allowbreak{}}{DER\_\allowbreak{}XVIII-d:\allowbreak{}71}
\DERapparatusitem{ENV1412}{149;\allowbreak{}\allowbreak{} ms394;\allowbreak{}\allowbreak{} intersection\allowbreak{} formula}{\textbackslash{}\allowbreak{}cart}{\textbackslash{}\allowbreak{}cocart}{DER\_\allowbreak{}XVIII-d:\allowbreak{}72}
\DERapparatusitem{ENV1413}{149;\allowbreak{}\allowbreak{} running\allowbreak{} text\allowbreak{} or\allowbreak{} formula}{\textbackslash{}\allowbreak{}text\{\allowbreak{}def\}\allowbreak{}}{\textbackslash{}\allowbreak{}mathrm\{\allowbreak{}def\}\allowbreak{}}{TYPESETTING:\allowbreak{}7;\allowbreak{}\allowbreak{} NOTATION\_\allowbreak{}DER:\allowbreak{}17}
\DERapparatusitem{ENV1414}{149;\allowbreak{}\allowbreak{} ms394;\allowbreak{}\allowbreak{} intersection\allowbreak{} index}{\textbackslash{}\allowbreak{}lambda\textbackslash{}\allowbreak{}in\textbackslash{}\allowbreak{}Ob}{\textbackslash{}\allowbreak{}lambda\textbackslash{}\allowbreak{}in\textbackslash{}\allowbreak{}Fl}{DER\_\allowbreak{}XVIII-d:\allowbreak{}73}
\DERapparatusitem{ENV1415}{150;\allowbreak{}\allowbreak{} ms395;\allowbreak{}\allowbreak{} 4.\allowbreak{}14.\allowbreak{}20(\allowbreak{}c)\allowbreak{}\allowbreak{} cardinal\allowbreak{} threshold}{\textbackslash{}\allowbreak{}cardred\textbackslash{}\allowbreak{}mathcal\{\allowbreak{}M\}\allowbreak{}\_\allowbreak{}\textbackslash{}\allowbreak{}pi}{\textbackslash{}\allowbreak{}cardred\textbackslash{}\allowbreak{}mathcal\{\allowbreak{}M\}\allowbreak{}\_\allowbreak{}\{\allowbreak{}\textbackslash{}\allowbreak{}pi\_\allowbreak{}0\}\allowbreak{}}{DER\_\allowbreak{}XVIII-d:\allowbreak{}74}
\DERapparatusitem{ENV1416}{151;\allowbreak{}\allowbreak{} ms398;\allowbreak{}\allowbreak{} definition\allowbreak{} of\allowbreak{} the\allowbreak{} quintuple}{\$\allowbreak{}x\_\allowbreak{}i\allowbreak{} \textbackslash{}\allowbreak{}in\allowbreak{} \textbackslash{}\allowbreak{}Ob\textbackslash{}\allowbreak{}mathcal\{\allowbreak{}E\}\allowbreak{}\$\allowbreak{}\allowbreak{} for\allowbreak{} \$\allowbreak{}i}{\$\allowbreak{}x\_\allowbreak{}i\allowbreak{} \textbackslash{}\allowbreak{}in\allowbreak{} \textbackslash{}\allowbreak{}Ob\textbackslash{}\allowbreak{}mathcal\{\allowbreak{}E\}\allowbreak{}\_\allowbreak{}i\$\allowbreak{}\allowbreak{} for\allowbreak{} \$\allowbreak{}i}{DER\_\allowbreak{}XVIII-d:\allowbreak{}76}
\DERapparatusitem{ENV1417}{151;\allowbreak{}\allowbreak{} ms398;\allowbreak{}\allowbreak{} name\allowbreak{} of\allowbreak{} full\allowbreak{} section\allowbreak{} subcategory}{0}{1}{DER\_\allowbreak{}XVIII-d:\allowbreak{}77}
\DERapparatusitem{ENV1418}{151;\allowbreak{}\allowbreak{} ms398;\allowbreak{}\allowbreak{} ambient\allowbreak{} category\allowbreak{} for\allowbreak{} Gamma\_\allowbreak{}1}{\textbackslash{}\allowbreak{}mathcal\{\allowbreak{}M\}\allowbreak{}}{\textbackslash{}\allowbreak{}Gamma}{DER\_\allowbreak{}XVIII-d:\allowbreak{}78}
\DERapparatusitem{ENV1419}{151;\allowbreak{}\allowbreak{} ms399;\allowbreak{}\allowbreak{} tuple\allowbreak{} for\allowbreak{} comparison-arrow\allowbreak{} functor}{1}{0}{DER\_\allowbreak{}XVIII-d:\allowbreak{}79}
\DERapparatusitem{ENV1420}{151;\allowbreak{}\allowbreak{} ms397;\allowbreak{}\allowbreak{} 4.\allowbreak{}14.\allowbreak{}21(\allowbreak{}a)\allowbreak{},\allowbreak{}\allowbreak{} both\allowbreak{} projections}{}{\_\allowbreak{}0}{DER\_\allowbreak{}XVIII-d:\allowbreak{}75}
\DERapparatusitem{ENV1421}{151;\allowbreak{}\allowbreak{} ms399;\allowbreak{}\allowbreak{} conclusion\allowbreak{} about\allowbreak{} iso-comma\allowbreak{} subcategory}{}{\_\allowbreak{}0}{DER\_\allowbreak{}XVIII-d:\allowbreak{}80}
\DERapparatusitem{ENV1422}{151;\allowbreak{}\allowbreak{} ms399;\allowbreak{}\allowbreak{} ambient\allowbreak{} arrow\allowbreak{} category\allowbreak{} in\allowbreak{} conclusion}{}{\_\allowbreak{}0}{DER\_\allowbreak{}XVIII-d:\allowbreak{}81}
\DERapparatusitem{ENV1423}{151;\allowbreak{}\allowbreak{} ms399;\allowbreak{}\allowbreak{} object\allowbreak{} to\allowbreak{} interpret\allowbreak{} as\allowbreak{} cocartesian\allowbreak{} sections}{}{\_\allowbreak{}0}{DER\_\allowbreak{}XVIII-d:\allowbreak{}82}
\DERapparatusitem{ENV1424}{151;\allowbreak{}\allowbreak{} ms400;\allowbreak{}\allowbreak{} final\allowbreak{} pair\allowbreak{} of\allowbreak{} projections}{0,\allowbreak{}1}{1,\allowbreak{}2}{DER\_\allowbreak{}XVIII-d:\allowbreak{}83}
\DERapparatusitem{ENV1425}{152;\allowbreak{}\allowbreak{} ms\allowbreak{} 402,\allowbreak{}\allowbreak{} definition\allowbreak{} of\allowbreak{} G\_\allowbreak{}λ}{s}{t}{DER\_\allowbreak{}XVIII-d:\allowbreak{}84}
\DERapparatusitem{ENV1426}{152;\allowbreak{}\allowbreak{} ms\allowbreak{} 402,\allowbreak{}\allowbreak{} definition\allowbreak{} of\allowbreak{} G\_\allowbreak{}λ}{1}{0}{DER\_\allowbreak{}XVIII-d:\allowbreak{}84}
\DERapparatusitem{ENV1427}{152;\allowbreak{}\allowbreak{} ms\allowbreak{} 402,\allowbreak{}\allowbreak{} definition\allowbreak{} of\allowbreak{} G\_\allowbreak{}λ}{s}{t}{DER\_\allowbreak{}XVIII-d:\allowbreak{}84}
\DERapparatusitem{ENV1428}{152;\allowbreak{}\allowbreak{} ms\allowbreak{} 402,\allowbreak{}\allowbreak{} definition\allowbreak{} of\allowbreak{} G\_\allowbreak{}λ}{1}{0}{DER\_\allowbreak{}XVIII-d:\allowbreak{}84}
\DERapparatusitem{ENV1429}{152;\allowbreak{}\allowbreak{} ms\allowbreak{} 402,\allowbreak{}\allowbreak{} definition\allowbreak{} of\allowbreak{} G\_\allowbreak{}λ}{s}{t}{DER\_\allowbreak{}XVIII-d:\allowbreak{}84}
\DERapparatusitem{ENV1430}{152;\allowbreak{}\allowbreak{} ms\allowbreak{} 402,\allowbreak{}\allowbreak{} definition\allowbreak{} of\allowbreak{} G\_\allowbreak{}λ}{1}{0}{DER\_\allowbreak{}XVIII-d:\allowbreak{}84}
\DERapparatusitem{ENV1431}{153;\allowbreak{}\allowbreak{} ms\allowbreak{} 404,\allowbreak{}\allowbreak{} first\allowbreak{} paragraph}{associates\allowbreak{} to\allowbreak{} every\allowbreak{} sections\allowbreak{} \$\allowbreak{}f\$\allowbreak{}}{associates\allowbreak{} to\allowbreak{} every\allowbreak{} section\allowbreak{} \$\allowbreak{}f\$\allowbreak{}}{DER\_\allowbreak{}XVIII-d:\allowbreak{}85}
\DERapparatusitem{ENV1432}{154;\allowbreak{}\allowbreak{} ms\allowbreak{} 406,\allowbreak{}\allowbreak{} Lemma\allowbreak{} 4.\allowbreak{}14.\allowbreak{}22.\allowbreak{}1}{x)\allowbreak{}\textbackslash{}\allowbreak{}longrightarrow\allowbreak{} G(\allowbreak{}x}{X)\allowbreak{}\textbackslash{}\allowbreak{}longrightarrow\allowbreak{} G(\allowbreak{}X}{DER\_\allowbreak{}XVIII-d:\allowbreak{}87}
\DERapparatusitem{ENV1433}{154;\allowbreak{}\allowbreak{} ms\allowbreak{} 408,\allowbreak{}\allowbreak{} closing\allowbreak{} paragraph}{is\allowbreak{} reduced\allowbreak{} to\allowbreak{} the\allowbreak{} cofibrant\allowbreak{} cases}{is\allowbreak{} reduced\allowbreak{} to\allowbreak{} the\allowbreak{} cofibrant\allowbreak{} case}{DER\_\allowbreak{}XVIII-d:\allowbreak{}88}
\DERapparatusitem{ENV1434}{155;\allowbreak{}\allowbreak{} ms\allowbreak{} 411,\allowbreak{}\allowbreak{} necessity\allowbreak{} proof\allowbreak{} first\allowbreak{} sentence}{\_\allowbreak{}s}{}{DER\_\allowbreak{}XVIII-d:\allowbreak{}91}
\DERapparatusitem{ENV1435}{155;\allowbreak{}\allowbreak{} ms\allowbreak{} 411,\allowbreak{}\allowbreak{} opfibration\allowbreak{} adjunction}{admits\allowbreak{} a\allowbreak{} right\allowbreak{} adjoint\allowbreak{} \$\allowbreak{}\textbackslash{}\allowbreak{}bar\{\allowbreak{}\textbackslash{}\allowbreak{}varphi\}\allowbreak{}\$\allowbreak{}}{admits\allowbreak{} a\allowbreak{} left\allowbreak{} adjoint\allowbreak{} \$\allowbreak{}\textbackslash{}\allowbreak{}bar\{\allowbreak{}\textbackslash{}\allowbreak{}varphi\}\allowbreak{}\$\allowbreak{}}{DER\_\allowbreak{}XVIII-d:\allowbreak{}92}
\DERapparatusitem{ENV1436}{155;\allowbreak{}\allowbreak{} ms\allowbreak{} 411,\allowbreak{}\allowbreak{} fibration\allowbreak{} adjunction}{admits\allowbreak{} a\allowbreak{} left\allowbreak{} adjoint\allowbreak{} \$\allowbreak{}\textbackslash{}\allowbreak{}tilde\{\allowbreak{}\textbackslash{}\allowbreak{}psi\}\allowbreak{}\$\allowbreak{}}{admits\allowbreak{} a\allowbreak{} right\allowbreak{} adjoint\allowbreak{} \$\allowbreak{}\textbackslash{}\allowbreak{}tilde\{\allowbreak{}\textbackslash{}\allowbreak{}psi\}\allowbreak{}\$\allowbreak{}}{DER\_\allowbreak{}XVIII-d:\allowbreak{}93}
\DERapparatusitem{ENV1437}{155;\allowbreak{}\allowbreak{} ms\allowbreak{} 411,\allowbreak{}\allowbreak{} composite\allowbreak{} defining\allowbreak{} ψ}{1}{0}{DER\_\allowbreak{}XVIII-d:\allowbreak{}94}
\DERapparatusitem{ENV1438}{156;\allowbreak{}\allowbreak{} ms\allowbreak{} 412,\allowbreak{}\allowbreak{} first\allowbreak{} adjunction\allowbreak{} display}{L}{h}{DER\_\allowbreak{}XVIII-d:\allowbreak{}95}
\DERapparatusitem{ENV1439}{156;\allowbreak{}\allowbreak{} ms\allowbreak{} 412,\allowbreak{}\allowbreak{} first\allowbreak{} adjunction\allowbreak{} display}{L}{h}{DER\_\allowbreak{}XVIII-d:\allowbreak{}95;\allowbreak{}\allowbreak{} DER\_\allowbreak{}XVIII-d:\allowbreak{}96}
\DERapparatusitem{ENV1440}{157;\allowbreak{}\allowbreak{} ms\allowbreak{} 414,\allowbreak{}\allowbreak{} reference\allowbreak{} at\allowbreak{} start}{condition\allowbreak{} a)\allowbreak{}\allowbreak{} of\allowbreak{} 4.\allowbreak{}14.\allowbreak{}23.\allowbreak{}3}{condition\allowbreak{} a)\allowbreak{}\allowbreak{} of\allowbreak{} 4.\allowbreak{}14.\allowbreak{}24}{DER\_\allowbreak{}XVIII-d:\allowbreak{}97}
\DERapparatusitem{ENV1441}{157;\allowbreak{}\allowbreak{} ms\allowbreak{} 414,\allowbreak{}\allowbreak{} opfibration\allowbreak{} cross-Hom\allowbreak{} argument}{}{\textasciicircum{}\allowbreak{}\{\allowbreak{}\textbackslash{}\allowbreak{}circ\}\allowbreak{}}{DER\_\allowbreak{}XVIII-d:\allowbreak{}98}
\DERapparatusitem{ENV1442}{157;\allowbreak{}\allowbreak{} ms\allowbreak{} 414,\allowbreak{}\allowbreak{} implication\allowbreak{} b⇒c}{}{\_\allowbreak{}1}{DER\_\allowbreak{}XVIII-d:\allowbreak{}99}
\DERapparatusitem{ENV1443}{158;\allowbreak{}\allowbreak{} ms\allowbreak{} 416,\allowbreak{}\allowbreak{} first\allowbreak{} sentence\allowbreak{} of\allowbreak{} proof\allowbreak{} of\allowbreak{} 1°}{}{\_\allowbreak{}1}{DER\_\allowbreak{}XVIII-d:\allowbreak{}100}
\DERapparatusitem{ENV1444}{158;\allowbreak{}\allowbreak{} ms\allowbreak{} 418,\allowbreak{}\allowbreak{} first\allowbreak{} object\allowbreak{} declaration}{}{\textbackslash{}\allowbreak{}mathcal\{\allowbreak{}}{DER\_\allowbreak{}XVIII-d:\allowbreak{}102}
\DERapparatusitem{ENV1445}{158;\allowbreak{}\allowbreak{} ms\allowbreak{} 418,\allowbreak{}\allowbreak{} first\allowbreak{} object\allowbreak{} declaration}{}{\}\allowbreak{}}{DER\_\allowbreak{}XVIII-d:\allowbreak{}102}
\DERapparatusitem{ENV1446}{158;\allowbreak{}\allowbreak{} ms\allowbreak{} 418,\allowbreak{}\allowbreak{} upper-fibre\allowbreak{} case}{x}{y}{DER\_\allowbreak{}XVIII-d:\allowbreak{}103}
\DERapparatusitem{ENV1447}{158;\allowbreak{}\allowbreak{} ms\allowbreak{} 418,\allowbreak{}\allowbreak{} lower-fibre\allowbreak{} case,\allowbreak{}\allowbreak{} after\allowbreak{} Sinon}{0}{1}{DER\_\allowbreak{}XVIII-d:\allowbreak{}104}
\DERapparatusitem{ENV1448}{preamble;\allowbreak{}\allowbreak{} local\allowbreak{} macro\allowbreak{} definition}{\textbackslash{}\allowbreak{}text\{\allowbreak{}\textbackslash{}\allowbreak{}upshape\allowbreak{} def\}\allowbreak{}}{\textbackslash{}\allowbreak{}mathrm\{\allowbreak{}def\}\allowbreak{}}{TYPESETTING:\allowbreak{}7;\allowbreak{}\allowbreak{} NOTATION\_\allowbreak{}DER:\allowbreak{}17}
\DERapparatusitem{ENV1449}{159;\allowbreak{}\allowbreak{} ms\allowbreak{} 419,\allowbreak{}\allowbreak{} sentence\allowbreak{} before\allowbreak{} presentation}{But\allowbreak{} if\allowbreak{} \$\allowbreak{}x\$\allowbreak{}\allowbreak{} is\allowbreak{} in\allowbreak{} \$\allowbreak{}\textbackslash{}\allowbreak{}cE\_\allowbreak{}\{\allowbreak{}0\textbackslash{}\allowbreak{},\allowbreak{}\textbackslash{}\allowbreak{}pi\}\allowbreak{}\$\allowbreak{}}{But\allowbreak{} if\allowbreak{} \$\allowbreak{}x\_\allowbreak{}0\$\allowbreak{}\allowbreak{} is\allowbreak{} in\allowbreak{} \$\allowbreak{}\textbackslash{}\allowbreak{}cE\_\allowbreak{}\{\allowbreak{}0\textbackslash{}\allowbreak{},\allowbreak{}\textbackslash{}\allowbreak{}pi\}\allowbreak{}\$\allowbreak{}}{DER\_\allowbreak{}XVIII-d:\allowbreak{}105;\allowbreak{}\allowbreak{} DER\_\allowbreak{}XVIII-e:\allowbreak{}1}
\DERapparatusitem{ENV1450}{159;\allowbreak{}\allowbreak{} ms\allowbreak{} 419,\allowbreak{}\allowbreak{} parenthesis\allowbreak{} after\allowbreak{} presentation}{\textbackslash{}\allowbreak{}cardred\allowbreak{} \textbackslash{}\allowbreak{}cE\_\allowbreak{}\{\allowbreak{}\textbackslash{}\allowbreak{}pi\_\allowbreak{}0\}\allowbreak{}}{\textbackslash{}\allowbreak{}cardred\allowbreak{} \textbackslash{}\allowbreak{}cE\_\allowbreak{}\{\allowbreak{}0\textbackslash{}\allowbreak{},\allowbreak{}\textbackslash{}\allowbreak{}pi\_\allowbreak{}0\}\allowbreak{}}{DER\_\allowbreak{}XVIII-d:\allowbreak{}106;\allowbreak{}\allowbreak{} DER\_\allowbreak{}XVIII-e:\allowbreak{}2}
\DERapparatusitem{ENV1451}{159;\allowbreak{}\allowbreak{} ms\allowbreak{} 419,\allowbreak{}\allowbreak{} variance\allowbreak{} sentence}{x\_\allowbreak{}0\allowbreak{} \textbackslash{}\allowbreak{}longrightarrow\allowbreak{} H(\allowbreak{}x\_\allowbreak{}0,\allowbreak{}x\_\allowbreak{}1)\allowbreak{}\allowbreak{} :\allowbreak{}\allowbreak{} \textbackslash{}\allowbreak{}cE\_\allowbreak{}0\allowbreak{} \textbackslash{}\allowbreak{}longrightarrow\allowbreak{} \textbackslash{}\allowbreak{}Ens}{x\_\allowbreak{}0\allowbreak{} \textbackslash{}\allowbreak{}longmapsto\allowbreak{} H(\allowbreak{}x\_\allowbreak{}0,\allowbreak{}x\_\allowbreak{}1)\allowbreak{}\allowbreak{} :\allowbreak{}\allowbreak{} \textbackslash{}\allowbreak{}cE\_\allowbreak{}0\allowbreak{} \textbackslash{}\allowbreak{}longrightarrow\allowbreak{} \textbackslash{}\allowbreak{}Ens\textasciicircum{}\allowbreak{}\{\allowbreak{}\textbackslash{}\allowbreak{}circ\}\allowbreak{}}{DER\_\allowbreak{}XVIII-d:\allowbreak{}107;\allowbreak{}\allowbreak{} DER\_\allowbreak{}XVIII-e:\allowbreak{}3}
\DERapparatusitem{ENV1452}{159;\allowbreak{}\allowbreak{} ms\allowbreak{} 419,\allowbreak{}\allowbreak{} parenthetical\allowbreak{} definition}{}{\_\allowbreak{}0}{DER\_\allowbreak{}XVIII-d:\allowbreak{}108;\allowbreak{}\allowbreak{} DER\_\allowbreak{}XVIII-e:\allowbreak{}4}
\DERapparatusitem{ENV1453}{159;\allowbreak{}\allowbreak{} ms\allowbreak{} 420,\allowbreak{}\allowbreak{} Corollary\allowbreak{} 4.\allowbreak{}14.\allowbreak{}25(\allowbreak{}1)\allowbreak{},\allowbreak{}\allowbreak{} right-hand\allowbreak{} condition}{necessary\allowbreak{} and\allowbreak{} sufficient\allowbreak{} that\allowbreak{} \$\allowbreak{}\textbackslash{}\allowbreak{}cE\_\allowbreak{}0\$\allowbreak{}\allowbreak{} satisfy\allowbreak{} \$\allowbreak{}L\_\allowbreak{}\textbackslash{}\allowbreak{}pi\$\allowbreak{}}{necessary\allowbreak{} and\allowbreak{} sufficient\allowbreak{} that\allowbreak{} \$\allowbreak{}\textbackslash{}\allowbreak{}cE\_\allowbreak{}0\$\allowbreak{}\allowbreak{} and\allowbreak{} \$\allowbreak{}\textbackslash{}\allowbreak{}cE\_\allowbreak{}1\$\allowbreak{}\allowbreak{} satisfy\allowbreak{} \$\allowbreak{}L\_\allowbreak{}\textbackslash{}\allowbreak{}pi\$\allowbreak{}}{DER\_\allowbreak{}XVIII-d:\allowbreak{}109;\allowbreak{}\allowbreak{} DER\_\allowbreak{}XVIII-e:\allowbreak{}5}
\DERapparatusitem{ENV1454}{159;\allowbreak{}\allowbreak{} ms\allowbreak{} 421,\allowbreak{}\allowbreak{} Corollary\allowbreak{} 4.\allowbreak{}14.\allowbreak{}25(\allowbreak{}2)\allowbreak{},\allowbreak{}\allowbreak{} quantified\allowbreak{} objects}{x\_\allowbreak{}0\allowbreak{} \textbackslash{}\allowbreak{}in\allowbreak{} \textbackslash{}\allowbreak{}Ob\textbackslash{}\allowbreak{}cE\_\allowbreak{}\{\allowbreak{}\textbackslash{}\allowbreak{}pi\_\allowbreak{}0\}\allowbreak{}}{x\_\allowbreak{}0\allowbreak{} \textbackslash{}\allowbreak{}in\allowbreak{} \textbackslash{}\allowbreak{}Ob\textbackslash{}\allowbreak{}cE\_\allowbreak{}\{\allowbreak{}0\textbackslash{}\allowbreak{},\allowbreak{}\textbackslash{}\allowbreak{}pi\_\allowbreak{}0\}\allowbreak{}}{DER\_\allowbreak{}XVIII-d:\allowbreak{}110;\allowbreak{}\allowbreak{} DER\_\allowbreak{}XVIII-e:\allowbreak{}6}
\DERapparatusitem{ENV1455}{160;\allowbreak{}\allowbreak{} ms\allowbreak{} 422,\allowbreak{}\allowbreak{} parenthesis\allowbreak{} after\allowbreak{} Theorem\allowbreak{} 4.\allowbreak{}14.\allowbreak{}26}{(\allowbreak{}generalization\allowbreak{} of\allowbreak{} 4.\allowbreak{}14.\allowbreak{}13)\allowbreak{}}{(\allowbreak{}generalization\allowbreak{} of\allowbreak{} 4.\allowbreak{}14.\allowbreak{}17)\allowbreak{}}{DER\_\allowbreak{}XVIII-d:\allowbreak{}111}
\DERapparatusitem{ENV1456}{161;\allowbreak{}\allowbreak{} ms\allowbreak{} 424,\allowbreak{}\allowbreak{} definition\allowbreak{} of\allowbreak{} M′}{x)\allowbreak{}\allowbreak{} \textbackslash{}\allowbreak{}longrightarrow\allowbreak{} G(\allowbreak{}x}{X)\allowbreak{}\allowbreak{} \textbackslash{}\allowbreak{}longrightarrow\allowbreak{} G(\allowbreak{}X}{DER\_\allowbreak{}XVIII-d:\allowbreak{}112;\allowbreak{}\allowbreak{} DER\_\allowbreak{}XVIII-e:\allowbreak{}7}
\DERapparatusitem{ENV1457}{161;\allowbreak{}\allowbreak{} ms\allowbreak{} 424,\allowbreak{}\allowbreak{} following\allowbreak{} sentence}{\textbackslash{}\allowbreak{}pi}{\textbackslash{}\allowbreak{}Pi}{DER\_\allowbreak{}XVIII-d:\allowbreak{}113;\allowbreak{}\allowbreak{} DER\_\allowbreak{}XVIII-e:\allowbreak{}8}
\DERapparatusitem{ENV1458}{161;\allowbreak{}\allowbreak{} ms\allowbreak{} 424,\allowbreak{}\allowbreak{} last\allowbreak{} composite}{h}{l}{DER\_\allowbreak{}XVIII-d:\allowbreak{}114;\allowbreak{}\allowbreak{} DER\_\allowbreak{}XVIII-e:\allowbreak{}9}
\DERapparatusitem{ENV1459}{161;\allowbreak{}\allowbreak{} ms\allowbreak{} 424,\allowbreak{}\allowbreak{} last\allowbreak{} composite}{ok}{ng}{DER\_\allowbreak{}XVIII-d:\allowbreak{}114;\allowbreak{}\allowbreak{} DER\_\allowbreak{}XVIII-e:\allowbreak{}9}
\DERapparatusitem{ENV1460}{162;\allowbreak{}\allowbreak{} ms\allowbreak{} 426,\allowbreak{}\allowbreak{} Corollary\allowbreak{} 4.\allowbreak{}15.\allowbreak{}2(\allowbreak{}a)\allowbreak{}}{t}{s}{DER\_\allowbreak{}XVIII-d:\allowbreak{}115;\allowbreak{}\allowbreak{} DER\_\allowbreak{}XVIII-e:\allowbreak{}10}
\DERapparatusitem{ENV1461}{162;\allowbreak{}\allowbreak{} ms\allowbreak{} 426,\allowbreak{}\allowbreak{} Corollary\allowbreak{} 4.\allowbreak{}15.\allowbreak{}2(\allowbreak{}b)\allowbreak{}}{}{\_\allowbreak{}1}{DER\_\allowbreak{}XVIII-d:\allowbreak{}116;\allowbreak{}\allowbreak{} DER\_\allowbreak{}XVIII-e:\allowbreak{}11}
\DERapparatusitem{ENV1462}{162;\allowbreak{}\allowbreak{} ms\allowbreak{} 426,\allowbreak{}\allowbreak{} Corollary\allowbreak{} 4.\allowbreak{}15.\allowbreak{}2(\allowbreak{}c)\allowbreak{},\allowbreak{}\allowbreak{} condition\allowbreak{} (\allowbreak{}1)\allowbreak{}}{\$\allowbreak{}\textbackslash{}\allowbreak{}pi\_\allowbreak{}0\$\allowbreak{}\allowbreak{} is\allowbreak{} useful\allowbreak{} for\allowbreak{} \$\allowbreak{}\textbackslash{}\allowbreak{}cE\_\allowbreak{}s\$\allowbreak{}\allowbreak{} }{\$\allowbreak{}\textbackslash{}\allowbreak{}pi\_\allowbreak{}0\$\allowbreak{}\allowbreak{} is\allowbreak{} useful\allowbreak{} for\allowbreak{} \$\allowbreak{}\textbackslash{}\allowbreak{}cE\_\allowbreak{}s\$\allowbreak{}\allowbreak{} and\allowbreak{} \$\allowbreak{}\textbackslash{}\allowbreak{}pi\_\allowbreak{}0\textbackslash{}\allowbreak{}geqslant\textbackslash{}\allowbreak{}sup\_\allowbreak{}\{\allowbreak{}\textbackslash{}\allowbreak{}lambda\textbackslash{}\allowbreak{}in\textbackslash{}\allowbreak{}Fl\textbackslash{}\allowbreak{}cC\}\allowbreak{}\textbackslash{}\allowbreak{}pi\_\allowbreak{}\textbackslash{}\allowbreak{}lambda\$\allowbreak{}\allowbreak{} }{DER\_\allowbreak{}XVIII-d:\allowbreak{}120}
\DERapparatusitem{ENV1463}{162;\allowbreak{}\allowbreak{} ms\allowbreak{} 426,\allowbreak{}\allowbreak{} Corollary\allowbreak{} 4.\allowbreak{}15.\allowbreak{}2(\allowbreak{}c)\allowbreak{},\allowbreak{}\allowbreak{} sentence\allowbreak{} after\allowbreak{} (\allowbreak{}*)\allowbreak{}}{\textbackslash{}\allowbreak{}cE\_\allowbreak{}\textbackslash{}\allowbreak{}pi}{\textbackslash{}\allowbreak{}cE}{DER\_\allowbreak{}XVIII-d:\allowbreak{}117;\allowbreak{}\allowbreak{} DER\_\allowbreak{}XVIII-e:\allowbreak{}12}
\DERapparatusitem{ENV1464}{162;\allowbreak{}\allowbreak{} ms\allowbreak{} 426,\allowbreak{}\allowbreak{} Ind-equivalence\allowbreak{} in\allowbreak{} part\allowbreak{} c}{\_\allowbreak{}\{\allowbreak{}s\textbackslash{}\allowbreak{},\allowbreak{}\textbackslash{}\allowbreak{}pi\}\allowbreak{}}{(\allowbreak{}\textbackslash{}\allowbreak{}pi)\allowbreak{}}{DER\_\allowbreak{}XVIII-d:\allowbreak{}118;\allowbreak{}\allowbreak{} DER\_\allowbreak{}XVIII-e:\allowbreak{}13}
\DERapparatusitem{ENV1465}{162;\allowbreak{}\allowbreak{} ms\allowbreak{} 427,\allowbreak{}\allowbreak{} opening\allowbreak{} of\allowbreak{} proof\allowbreak{} a}{0}{1}{DER\_\allowbreak{}XVIII-d:\allowbreak{}119;\allowbreak{}\allowbreak{} DER\_\allowbreak{}XVIII-e:\allowbreak{}14}
\DERapparatusitem{ENV1466}{162;\allowbreak{}\allowbreak{} ms\allowbreak{} 427,\allowbreak{}\allowbreak{} opening\allowbreak{} of\allowbreak{} proof\allowbreak{} a}{0}{1}{DER\_\allowbreak{}XVIII-d:\allowbreak{}119;\allowbreak{}\allowbreak{} DER\_\allowbreak{}XVIII-e:\allowbreak{}14}
\DERapparatusitem{ENV1467}{163;\allowbreak{}\allowbreak{} ms\allowbreak{} 427\allowbreak{} continuation,\allowbreak{}\allowbreak{} base-arrow\allowbreak{} set}{}{f(\allowbreak{}}{DER\_\allowbreak{}XVIII-d:\allowbreak{}121}
\DERapparatusitem{ENV1468}{163;\allowbreak{}\allowbreak{} ms\allowbreak{} 427\allowbreak{} continuation,\allowbreak{}\allowbreak{} base-arrow\allowbreak{} set}{}{)\allowbreak{}}{DER\_\allowbreak{}XVIII-d:\allowbreak{}121}
\DERapparatusitem{ENV1469}{163;\allowbreak{}\allowbreak{} ms\allowbreak{} 427\allowbreak{} continuation,\allowbreak{}\allowbreak{} base-arrow\allowbreak{} set}{}{f(\allowbreak{}}{DER\_\allowbreak{}XVIII-d:\allowbreak{}121}
\DERapparatusitem{ENV1470}{163;\allowbreak{}\allowbreak{} ms\allowbreak{} 427\allowbreak{} continuation,\allowbreak{}\allowbreak{} base-arrow\allowbreak{} set}{}{)\allowbreak{}}{DER\_\allowbreak{}XVIII-d:\allowbreak{}121}
\DERapparatusitem{ENV1471}{163;\allowbreak{}\allowbreak{} ms\allowbreak{} 427\allowbreak{} continuation,\allowbreak{}\allowbreak{} concluding\allowbreak{} preservation\allowbreak{} statement}{0}{1}{DER\_\allowbreak{}XVIII-d:\allowbreak{}122;\allowbreak{}\allowbreak{} DER\_\allowbreak{}XVIII-e:\allowbreak{}15}
\DERapparatusitem{ENV1472}{163;\allowbreak{}\allowbreak{} ms\allowbreak{} 429,\allowbreak{}\allowbreak{} bound\allowbreak{} before\allowbreak{} final\allowbreak{} reduction}{\textbackslash{}\allowbreak{}sup(\allowbreak{}\textbackslash{}\allowbreak{}card\allowbreak{} S,\allowbreak{}\allowbreak{} \textbackslash{}\allowbreak{}sup\_\allowbreak{}\textbackslash{}\allowbreak{}lambda\allowbreak{} \textbackslash{}\allowbreak{}pi\_\allowbreak{}\textbackslash{}\allowbreak{}lambda)\allowbreak{}}{\textbackslash{}\allowbreak{}sup(\allowbreak{}\textbackslash{}\allowbreak{}card\textbackslash{}\allowbreak{}cC,\allowbreak{}\allowbreak{} \textbackslash{}\allowbreak{}sup\_\allowbreak{}\textbackslash{}\allowbreak{}lambda\allowbreak{} \textbackslash{}\allowbreak{}pi\_\allowbreak{}\textbackslash{}\allowbreak{}lambda)\allowbreak{}}{DER\_\allowbreak{}XVIII-d:\allowbreak{}123;\allowbreak{}\allowbreak{} DER\_\allowbreak{}XVIII-e:\allowbreak{}16}
\DERapparatusitem{ENV1473}{164;\allowbreak{}\allowbreak{} ms430\allowbreak{} small-generator\allowbreak{} case}{s}{x}{DER\_\allowbreak{}XVIII-e:\allowbreak{}17}
\DERapparatusitem{ENV1474}{164;\allowbreak{}\allowbreak{} ms430\allowbreak{} codomain\allowbreak{} of\allowbreak{} H\_\allowbreak{}lambda(\allowbreak{}-)\allowbreak{}}{s}{t}{DER\_\allowbreak{}XVIII-e:\allowbreak{}18}
\DERapparatusitem{ENV1475}{164;\allowbreak{}\allowbreak{} ms430\allowbreak{} displayed\allowbreak{} limit\allowbreak{} formula}{\_\allowbreak{}0}{}{DER\_\allowbreak{}XVIII-e:\allowbreak{}19}
\DERapparatusitem{ENV1476}{164;\allowbreak{}\allowbreak{} ms431\allowbreak{} first\allowbreak{} sentence}{}{\_\allowbreak{}\textbackslash{}\allowbreak{}lambda}{DER\_\allowbreak{}XVIII-e:\allowbreak{}20}
\DERapparatusitem{ENV1477}{164;\allowbreak{}\allowbreak{} ms431\allowbreak{} conclusion\allowbreak{} before\allowbreak{} inclusion}{(\allowbreak{}\textbackslash{}\allowbreak{}pi)\allowbreak{}}{\_\allowbreak{}\textbackslash{}\allowbreak{}pi}{DER\_\allowbreak{}XVIII-e:\allowbreak{}21}
\DERapparatusitem{ENV1478}{164;\allowbreak{}\allowbreak{} ms431\allowbreak{} usefulness\allowbreak{} clause}{\textbackslash{}\allowbreak{}pi}{\textbackslash{}\allowbreak{}cE\_\allowbreak{}s}{DER\_\allowbreak{}XVIII-e:\allowbreak{}22}
\DERapparatusitem{ENV1479}{165;\allowbreak{}\allowbreak{} ms431\allowbreak{} proof(\allowbreak{}d)\allowbreak{},\allowbreak{}\allowbreak{} projector\allowbreak{} object}{x}{y}{DER\_\allowbreak{}XVIII-e:\allowbreak{}23}
\DERapparatusitem{ENV1480}{165;\allowbreak{}\allowbreak{} running\allowbreak{} text\allowbreak{} or\allowbreak{} formula}{\textbackslash{}\allowbreak{}text\{\allowbreak{}def\}\allowbreak{}}{\textbackslash{}\allowbreak{}mathrm\{\allowbreak{}def\}\allowbreak{}}{TYPESETTING:\allowbreak{}7;\allowbreak{}\allowbreak{} NOTATION\_\allowbreak{}DER:\allowbreak{}17}
\DERapparatusitem{ENV1481}{165;\allowbreak{}\allowbreak{} ms432\allowbreak{} base-change\allowbreak{} argument}{s}{\textbackslash{}\allowbreak{}lambda}{DER\_\allowbreak{}XVIII-e:\allowbreak{}24}
\DERapparatusitem{ENV1482}{166;\allowbreak{}\allowbreak{} ms434\allowbreak{} footnote179}{with\allowbreak{} \$\allowbreak{}p\$\allowbreak{}\allowbreak{} a\allowbreak{} projector\allowbreak{} in\allowbreak{} \$\allowbreak{}p\$\allowbreak{}}{with\allowbreak{} \$\allowbreak{}p\$\allowbreak{}\allowbreak{} a\allowbreak{} projector\allowbreak{} on\allowbreak{} \$\allowbreak{}x\$\allowbreak{}}{DER\_\allowbreak{}XVIII-e:\allowbreak{}27}
\DERapparatusitem{ENV1483}{166;\allowbreak{}\allowbreak{} running\allowbreak{} text\allowbreak{} or\allowbreak{} formula}{\textbackslash{}\allowbreak{}text\{\allowbreak{}def\}\allowbreak{}}{\textbackslash{}\allowbreak{}mathrm\{\allowbreak{}def\}\allowbreak{}}{TYPESETTING:\allowbreak{}7;\allowbreak{}\allowbreak{} NOTATION\_\allowbreak{}DER:\allowbreak{}17}
\DERapparatusitem{ENV1484}{167;\allowbreak{}\allowbreak{} ms436\allowbreak{} numerical\allowbreak{} bound2}{s}{\{\allowbreak{}s\textbackslash{}\allowbreak{},\allowbreak{}\textbackslash{}\allowbreak{}pi\_\allowbreak{}0\}\allowbreak{}}{DER\_\allowbreak{}XVIII-e:\allowbreak{}29}
\DERapparatusitem{ENV1485}{167;\allowbreak{}\allowbreak{} ms436\allowbreak{} first\allowbreak{} boxed\allowbreak{} transport\allowbreak{} formula}{s}{t}{DER\_\allowbreak{}XVIII-e:\allowbreak{}30}
\DERapparatusitem{ENV1486}{167;\allowbreak{}\allowbreak{} ms436\allowbreak{} second\allowbreak{} boxed\allowbreak{} transport\allowbreak{} formula}{s}{t}{DER\_\allowbreak{}XVIII-e:\allowbreak{}31}
\DERapparatusitem{ENV1487}{167;\allowbreak{}\allowbreak{} ms437\allowbreak{} diagram-introduction\allowbreak{} sentence}{}{\{\allowbreak{}}{DER\_\allowbreak{}XVIII-e:\allowbreak{}32}
\DERapparatusitem{ENV1488}{167;\allowbreak{}\allowbreak{} ms437\allowbreak{} diagram-introduction\allowbreak{} sentence}{}{\textbackslash{}\allowbreak{},\allowbreak{}*\}\allowbreak{}}{DER\_\allowbreak{}XVIII-e:\allowbreak{}32}
\DERapparatusitem{ENV1489}{167;\allowbreak{}\allowbreak{} ms437\allowbreak{} diagram-introduction\allowbreak{} sentence}{=\allowbreak{}}{\textbackslash{}\allowbreak{}leqslant}{DER\_\allowbreak{}XVIII-e:\allowbreak{}32}
\DERapparatusitem{ENV1490}{167;\allowbreak{}\allowbreak{} ms437\allowbreak{} transition\allowbreak{} display}{j}{i}{DER\_\allowbreak{}XVIII-e:\allowbreak{}33}
\DERapparatusitem{ENV1491}{168;\allowbreak{}\allowbreak{} ms438\allowbreak{} underbrace\allowbreak{} below\allowbreak{} composite\allowbreak{} transfer}{\textbackslash{}\allowbreak{}lambda\_\allowbreak{}i}{\textbackslash{}\allowbreak{}lambda}{DER\_\allowbreak{}XVIII-e:\allowbreak{}34}
\DERapparatusitem{ENV1492}{preamble;\allowbreak{}\allowbreak{} local\allowbreak{} macro\allowbreak{} definition}{\textbackslash{}\allowbreak{}text\{\allowbreak{}\textbackslash{}\allowbreak{}upshape\allowbreak{} def\}\allowbreak{}}{\textbackslash{}\allowbreak{}mathrm\{\allowbreak{}def\}\allowbreak{}}{TYPESETTING:\allowbreak{}7;\allowbreak{}\allowbreak{} NOTATION\_\allowbreak{}DER:\allowbreak{}17}
\DERapparatusitem{ENV1493}{169;\allowbreak{}\allowbreak{} ms.\allowbreak{}441,\allowbreak{}\allowbreak{} final\allowbreak{} sentence\allowbreak{} of4.\allowbreak{}15.\allowbreak{}7.\allowbreak{}1}{Moreover,\allowbreak{}\allowbreak{} if\allowbreak{} we\allowbreak{} can\allowbreak{} suppose\allowbreak{} \$\allowbreak{}J\$\allowbreak{}}{Moreover,\allowbreak{}\allowbreak{} we\allowbreak{} can\allowbreak{} suppose\allowbreak{} \$\allowbreak{}J\$\allowbreak{}}{DER\_\allowbreak{}XVIII-e:\allowbreak{}36}
\DERapparatusitem{ENV1494}{169;\allowbreak{}\allowbreak{} ms.\allowbreak{}442,\allowbreak{}\allowbreak{} final\allowbreak{} paragraph}{j}{\textbackslash{}\allowbreak{}lambda}{DER\_\allowbreak{}XVIII-e:\allowbreak{}37}
\DERapparatusitem{ENV1495}{170;\allowbreak{}\allowbreak{} ms.\allowbreak{}444,\allowbreak{}\allowbreak{} object-colour\allowbreak{} pigeonhole\allowbreak{} hypothesis}{>}{\textbackslash{}\allowbreak{}geqslant}{DER\_\allowbreak{}XVIII-e:\allowbreak{}39}
\DERapparatusitem{ENV1496}{170;\allowbreak{}\allowbreak{} ms.\allowbreak{}444,\allowbreak{}\allowbreak{} second\allowbreak{} displayed\allowbreak{} absorption\allowbreak{} relation}{j'}{'\allowbreak{} j}{DER\_\allowbreak{}XVIII-e:\allowbreak{}40}
\DERapparatusitem{ENV1497}{171;\allowbreak{}\allowbreak{} ms.\allowbreak{}445,\allowbreak{}\allowbreak{} first\allowbreak{} line,\allowbreak{}\allowbreak{} parenthesis\allowbreak{} beginning\allowbreak{} défini}{}{)\allowbreak{}}{DER\_\allowbreak{}XVIII-e:\allowbreak{}41}
\DERapparatusitem{ENV1498}{171;\allowbreak{}\allowbreak{} ms.\allowbreak{}445,\allowbreak{}\allowbreak{} proof4.\allowbreak{}15.\allowbreak{}7.\allowbreak{}4}{the\allowbreak{} automorphisms\allowbreak{} \$\allowbreak{}\textbackslash{}\allowbreak{}tau\$\allowbreak{}\allowbreak{} of\allowbreak{} the\allowbreak{} set\allowbreak{} \$\allowbreak{}E\$\allowbreak{}}{the\allowbreak{} endomorphisms\allowbreak{} \$\allowbreak{}\textbackslash{}\allowbreak{}tau\$\allowbreak{}\allowbreak{} of\allowbreak{} the\allowbreak{} set\allowbreak{} \$\allowbreak{}E\$\allowbreak{}}{DER\_\allowbreak{}XVIII-e:\allowbreak{}42}
\DERapparatusitem{ENV1499}{171;\allowbreak{}\allowbreak{} ms.\allowbreak{}445,\allowbreak{}\allowbreak{} last\allowbreak{} implication\allowbreak{} in\allowbreak{} proof4.\allowbreak{}15.\allowbreak{}7.\allowbreak{}4}{u}{\textbackslash{}\allowbreak{}tau}{DER\_\allowbreak{}XVIII-e:\allowbreak{}43}
\DERapparatusitem{ENV1500}{171;\allowbreak{}\allowbreak{} ms.\allowbreak{}446,\allowbreak{}\allowbreak{} cardinal\allowbreak{} bound}{1}{0}{DER\_\allowbreak{}XVIII-e:\allowbreak{}44}
\DERapparatusitem{ENV1501}{171;\allowbreak{}\allowbreak{} ms.\allowbreak{}446,\allowbreak{}\allowbreak{} same\allowbreak{} sentence}{J/\allowbreak{}j\_\allowbreak{}0}{j\_\allowbreak{}0\textbackslash{}\allowbreak{}backslash\allowbreak{} J}{DER\_\allowbreak{}XVIII-e:\allowbreak{}45}
\DERapparatusitem{ENV1502}{171;\allowbreak{}\allowbreak{} ms.\allowbreak{}446,\allowbreak{}\allowbreak{} two\allowbreak{} cancellation\allowbreak{} displays}{}{\textasciicircum{}\allowbreak{}\textbackslash{}\allowbreak{}natural}{DER\_\allowbreak{}XVIII-e:\allowbreak{}46}
\DERapparatusitem{ENV1503}{171;\allowbreak{}\allowbreak{} ms.\allowbreak{}446,\allowbreak{}\allowbreak{} two\allowbreak{} cancellation\allowbreak{} displays}{}{\textasciicircum{}\allowbreak{}\textbackslash{}\allowbreak{}natural}{DER\_\allowbreak{}XVIII-e:\allowbreak{}46}
\DERapparatusitem{ENV1504}{171;\allowbreak{}\allowbreak{} ms.\allowbreak{}446,\allowbreak{}\allowbreak{} two\allowbreak{} cancellation\allowbreak{} displays}{}{\textasciicircum{}\allowbreak{}\textbackslash{}\allowbreak{}natural}{DER\_\allowbreak{}XVIII-e:\allowbreak{}46}
\DERapparatusitem{ENV1505}{171;\allowbreak{}\allowbreak{} ms.\allowbreak{}446,\allowbreak{}\allowbreak{} two\allowbreak{} cancellation\allowbreak{} displays}{}{\textasciicircum{}\allowbreak{}\textbackslash{}\allowbreak{}natural}{DER\_\allowbreak{}XVIII-e:\allowbreak{}46}
\DERapparatusitem{ENV1506}{172;\allowbreak{}\allowbreak{} ms.\allowbreak{}448,\allowbreak{}\allowbreak{} second\allowbreak{} line\allowbreak{} of\allowbreak{} brace\allowbreak{} display}{-}{=\allowbreak{}}{DER\_\allowbreak{}XVIII-e:\allowbreak{}48}
\DERapparatusitem{ENV1507}{172;\allowbreak{}\allowbreak{} ms.\allowbreak{}449,\allowbreak{}\allowbreak{} sentence\allowbreak{} beginning\allowbreak{} de\allowbreak{} x}{x}{\textbackslash{}\allowbreak{}alpha}{DER\_\allowbreak{}XVIII-e:\allowbreak{}50}
\DERapparatusitem{ENV1508}{172;\allowbreak{}\allowbreak{} ms.\allowbreak{}449,\allowbreak{}\allowbreak{} compatibility\allowbreak{} display}{i\textbackslash{}\allowbreak{}tau\_\allowbreak{}\{\allowbreak{}ji\}\allowbreak{}\allowbreak{} \textbackslash{}\allowbreak{};\allowbreak{}=\allowbreak{}\textbackslash{}\allowbreak{};\allowbreak{}\allowbreak{} \textbackslash{}\allowbreak{}beta\_\allowbreak{}j}{j\textbackslash{}\allowbreak{}tau\_\allowbreak{}\{\allowbreak{}ji\}\allowbreak{}=\allowbreak{}\textbackslash{}\allowbreak{}beta\_\allowbreak{}i}{DER\_\allowbreak{}XVIII-e:\allowbreak{}51}
\DERapparatusitem{ENV1509}{172;\allowbreak{}\allowbreak{} ms.\allowbreak{}449,\allowbreak{}\allowbreak{} final\allowbreak{} display}{\textbackslash{}\allowbreak{}alpha\textbackslash{}\allowbreak{}beta\allowbreak{} \textbackslash{}\allowbreak{};\allowbreak{}=\allowbreak{}\textbackslash{}\allowbreak{};\allowbreak{}\allowbreak{} \textbackslash{}\allowbreak{}id\_\allowbreak{}z\textbackslash{}\allowbreak{}\allowbreak{} ,\allowbreak{}\allowbreak{} \textbackslash{}\allowbreak{}qquad\allowbreak{} \textbackslash{}\allowbreak{}beta\textbackslash{}\allowbreak{}alpha}{\textbackslash{}\allowbreak{}beta\textbackslash{}\allowbreak{}alpha=\allowbreak{}\textbackslash{}\allowbreak{}mathrm\{\allowbreak{}id\}\allowbreak{}\_\allowbreak{}z,\allowbreak{}\textbackslash{}\allowbreak{}quad\textbackslash{}\allowbreak{}alpha\textbackslash{}\allowbreak{}beta}{DER\_\allowbreak{}XVIII-e:\allowbreak{}52}
\DERapparatusitem{ENV1510}{173;\allowbreak{}\allowbreak{} ms.\allowbreak{}449\allowbreak{} continuation,\allowbreak{}\allowbreak{} sentence\allowbreak{} above\allowbreak{} (\allowbreak{}*\_\allowbreak{}i0)\allowbreak{}}{\textbackslash{}\allowbreak{}alpha\textbackslash{}\allowbreak{}beta}{\textbackslash{}\allowbreak{}beta\textbackslash{}\allowbreak{}alpha}{DER\_\allowbreak{}XVIII-e:\allowbreak{}53}
\DERapparatusitem{ENV1511}{173;\allowbreak{}\allowbreak{} ms.\allowbreak{}450,\allowbreak{}\allowbreak{} paragraph\allowbreak{} Exprimons}{beta\textbackslash{}\allowbreak{}alph}{alpha\textbackslash{}\allowbreak{}bet}{DER\_\allowbreak{}XVIII-e:\allowbreak{}54}
\DERapparatusitem{ENV1512}{173;\allowbreak{}\allowbreak{} ms.\allowbreak{}450,\allowbreak{}\allowbreak{} paragraph\allowbreak{} Exprimons}{beta\textbackslash{}\allowbreak{}alph}{alpha\textbackslash{}\allowbreak{}bet}{DER\_\allowbreak{}XVIII-e:\allowbreak{}54}
\DERapparatusitem{ENV1513}{173;\allowbreak{}\allowbreak{} ms.\allowbreak{}450,\allowbreak{}\allowbreak{} naturality\allowbreak{} display}{\textbackslash{}\allowbreak{}lambda}{j}{DER\_\allowbreak{}XVIII-e:\allowbreak{}55}
\DERapparatusitem{ENV1514}{174;\allowbreak{}\allowbreak{} running\allowbreak{} text\allowbreak{} or\allowbreak{} formula}{\textbackslash{}\allowbreak{}text\{\allowbreak{}def\}\allowbreak{}}{\textbackslash{}\allowbreak{}mathrm\{\allowbreak{}def\}\allowbreak{}}{TYPESETTING:\allowbreak{}7;\allowbreak{}\allowbreak{} NOTATION\_\allowbreak{}DER:\allowbreak{}17}
\DERapparatusitem{ENV1515}{174;\allowbreak{}\allowbreak{} ms.\allowbreak{}452,\allowbreak{}\allowbreak{} diagram\allowbreak{} defining\allowbreak{} beta\_\allowbreak{}i}{j}{i}{DER\_\allowbreak{}XVIII-e:\allowbreak{}57}
\DERapparatusitem{ENV1516}{174;\allowbreak{}\allowbreak{} ms.\allowbreak{}452,\allowbreak{}\allowbreak{} diagram\allowbreak{} defining\allowbreak{} beta\_\allowbreak{}i}{j}{i}{DER\_\allowbreak{}XVIII-e:\allowbreak{}57}
\DERapparatusitem{ENV1517}{175;\allowbreak{}\allowbreak{} ms.\allowbreak{}454,\allowbreak{}\allowbreak{} final\allowbreak{} membership}{whence\allowbreak{} \$\allowbreak{}u\allowbreak{} \textbackslash{}\allowbreak{}in\allowbreak{} N\_\allowbreak{}i\textbackslash{}\allowbreak{}alpha\_\allowbreak{}\{\allowbreak{}i'i\}\allowbreak{}\$\allowbreak{}}{whence\allowbreak{} \$\allowbreak{}u\allowbreak{} \textbackslash{}\allowbreak{}in\allowbreak{} N\_\allowbreak{}\{\allowbreak{}i'\}\allowbreak{}\textbackslash{}\allowbreak{}alpha\_\allowbreak{}\{\allowbreak{}i'i\}\allowbreak{}\$\allowbreak{}}{DER\_\allowbreak{}XVIII-e:\allowbreak{}58}
\DERapparatusitem{ENV1518}{176;\allowbreak{}\allowbreak{} ms.\allowbreak{}455,\allowbreak{}\allowbreak{} pigeonhole\allowbreak{} sentence}{\textbackslash{}\allowbreak{}pi\_\allowbreak{}0}{I}{DER\_\allowbreak{}XVIII-e:\allowbreak{}60}
\DERapparatusitem{ENV1519}{176;\allowbreak{}\allowbreak{} ms.\allowbreak{}455,\allowbreak{}\allowbreak{} quotient-map\allowbreak{} display}{\textbackslash{}\allowbreak{}longmapsto}{\textbackslash{}\allowbreak{}longrightarrow}{DER\_\allowbreak{}XVIII-e:\allowbreak{}61}
\DERapparatusitem{ENV1520}{177;\allowbreak{}\allowbreak{} ms.\allowbreak{}457,\allowbreak{}\allowbreak{} statement\allowbreak{} of\allowbreak{} key\allowbreak{} lemma}{\_\allowbreak{}0,\allowbreak{}i}{,\allowbreak{}i\_\allowbreak{}0}{DER\_\allowbreak{}XVIII-e:\allowbreak{}62}
\DERapparatusitem{ENV1521}{177;\allowbreak{}\allowbreak{} running\allowbreak{} text\allowbreak{} or\allowbreak{} formula}{\textbackslash{}\allowbreak{}text\{\allowbreak{}def\}\allowbreak{}}{\textbackslash{}\allowbreak{}mathrm\{\allowbreak{}def\}\allowbreak{}}{TYPESETTING:\allowbreak{}7;\allowbreak{}\allowbreak{} NOTATION\_\allowbreak{}DER:\allowbreak{}17}
\DERapparatusitem{ENV1522}{177;\allowbreak{}\allowbreak{} running\allowbreak{} text\allowbreak{} or\allowbreak{} formula}{\textbackslash{}\allowbreak{}text\{\allowbreak{}def\}\allowbreak{}}{\textbackslash{}\allowbreak{}mathrm\{\allowbreak{}def\}\allowbreak{}}{TYPESETTING:\allowbreak{}7;\allowbreak{}\allowbreak{} NOTATION\_\allowbreak{}DER:\allowbreak{}17}
\DERapparatusitem{ENV1523}{177;\allowbreak{}\allowbreak{} ms.\allowbreak{}459,\allowbreak{}\allowbreak{} cardinality\allowbreak{} display}{\textbackslash{}\allowbreak{}alpha}{\textbackslash{}\allowbreak{}infty}{DER\_\allowbreak{}XVIII-e:\allowbreak{}64}
\DERapparatusitem{ENV1524}{178;\allowbreak{}\allowbreak{} ms.\allowbreak{}458,\allowbreak{}\allowbreak{} definition\allowbreak{} of\allowbreak{} I\_\allowbreak{}alpha0}{\textbackslash{}\allowbreak{}alpha}{\textbackslash{}\allowbreak{}alpha\_\allowbreak{}0}{DER\_\allowbreak{}XVIII-e:\allowbreak{}63}
\DERapparatusitem{ENV1525}{178;\allowbreak{}\allowbreak{} ms.\allowbreak{}460,\allowbreak{}\allowbreak{} corollary4.\allowbreak{}15.\allowbreak{}8.\allowbreak{}4}{\textbackslash{}\allowbreak{}alpha}{\textbackslash{}\allowbreak{}pi}{DER\_\allowbreak{}XVIII-e:\allowbreak{}65}
\DERapparatusitem{ENV1526}{178;\allowbreak{}\allowbreak{} running\allowbreak{} text\allowbreak{} or\allowbreak{} formula}{\textbackslash{}\allowbreak{}text\{\allowbreak{}def\}\allowbreak{}}{\textbackslash{}\allowbreak{}mathrm\{\allowbreak{}def\}\allowbreak{}}{TYPESETTING:\allowbreak{}7;\allowbreak{}\allowbreak{} NOTATION\_\allowbreak{}DER:\allowbreak{}17}
\DERapparatusitem{ENV1527}{178;\allowbreak{}\allowbreak{} ms.\allowbreak{}462,\allowbreak{}\allowbreak{} condition(\allowbreak{}b)\allowbreak{}}{that\allowbreak{} \$\allowbreak{}\textbackslash{}\allowbreak{}alpha\_\allowbreak{}\{\allowbreak{}ii\_\allowbreak{}0\}\allowbreak{}\allowbreak{} =\allowbreak{}\allowbreak{} 1\$\allowbreak{}\allowbreak{} for\allowbreak{} \$\allowbreak{}i\allowbreak{} \textbackslash{}\allowbreak{}geqslant\allowbreak{} i\_\allowbreak{}0\$\allowbreak{},\allowbreak{}\allowbreak{} i.\allowbreak{}e.\allowbreak{}\allowbreak{} \$\allowbreak{}\textbackslash{}\allowbreak{}alpha\_\allowbreak{}\{\allowbreak{}ji\}\allowbreak{}\allowbreak{} =\allowbreak{}\allowbreak{} 1\$\allowbreak{}\allowbreak{} for\allowbreak{} \$\allowbreak{}j\allowbreak{} \textbackslash{}\allowbreak{}geqslant\allowbreak{} i\allowbreak{} \textbackslash{}\allowbreak{}geqslant\allowbreak{} i\_\allowbreak{}0\$\allowbreak{}}{that\allowbreak{} \$\allowbreak{}\textbackslash{}\allowbreak{}alpha\_\allowbreak{}\{\allowbreak{}ii\_\allowbreak{}0\}\allowbreak{}\$\allowbreak{}\allowbreak{} is\allowbreak{} invertible\allowbreak{} for\allowbreak{} \$\allowbreak{}i\allowbreak{} \textbackslash{}\allowbreak{}geqslant\allowbreak{} i\_\allowbreak{}0\$\allowbreak{},\allowbreak{}\allowbreak{} i.\allowbreak{}e.\allowbreak{}\allowbreak{} \$\allowbreak{}\textbackslash{}\allowbreak{}alpha\_\allowbreak{}\{\allowbreak{}ji\}\allowbreak{}\$\allowbreak{}\allowbreak{} is\allowbreak{} invertible\allowbreak{} for\allowbreak{} \$\allowbreak{}j\allowbreak{} \textbackslash{}\allowbreak{}geqslant\allowbreak{} i\allowbreak{} \textbackslash{}\allowbreak{}geqslant\allowbreak{} i\_\allowbreak{}0\$\allowbreak{}}{DER\_\allowbreak{}XVIII-e:\allowbreak{}66}
\DERapparatusitem{ENV1528}{preamble;\allowbreak{}\allowbreak{} local\allowbreak{} macro\allowbreak{} definition}{\textbackslash{}\allowbreak{}text\{\allowbreak{}def\}\allowbreak{}}{\textbackslash{}\allowbreak{}mathrm\{\allowbreak{}def\}\allowbreak{}}{TYPESETTING:\allowbreak{}7;\allowbreak{}\allowbreak{} NOTATION\_\allowbreak{}DER:\allowbreak{}17}
\DERapparatusitem{ENV1529}{179;\allowbreak{}\allowbreak{} ms.\allowbreak{}463,\allowbreak{}\allowbreak{} final\allowbreak{} contradiction\allowbreak{} display}{=\allowbreak{}\textbackslash{}\allowbreak{};\allowbreak{}\allowbreak{} \textbackslash{}\allowbreak{}card\allowbreak{} M\_\allowbreak{}\{\allowbreak{}\textbackslash{}\allowbreak{}alpha\_\allowbreak{}1\}\allowbreak{}\allowbreak{} \textbackslash{}\allowbreak{};\allowbreak{}\textbackslash{}\allowbreak{}leqslant\textbackslash{}\allowbreak{};\allowbreak{}\allowbreak{} \textbackslash{}\allowbreak{}card\allowbreak{} M\allowbreak{} \textbackslash{}\allowbreak{};\allowbreak{}=\allowbreak{}}{\textbackslash{}\allowbreak{}leqslant\textbackslash{}\allowbreak{}operatorname\{\allowbreak{}card\}\allowbreak{}M\_\allowbreak{}\{\allowbreak{}\textbackslash{}\allowbreak{}alpha\_\allowbreak{}1\}\allowbreak{}\textbackslash{}\allowbreak{}leqslant\textbackslash{}\allowbreak{}operatorname\{\allowbreak{}card\}\allowbreak{}M\textbackslash{}\allowbreak{}leqslant}{DER\_\allowbreak{}XVIII-e:\allowbreak{}67}
\DERapparatusitem{ENV1530}{180;\allowbreak{}\allowbreak{} ms467,\allowbreak{}representable\allowbreak{} functor}{}{\textasciicircum{}\allowbreak{}\{\allowbreak{}\textbackslash{}\allowbreak{}circ\}\allowbreak{}}{DER\_\allowbreak{}XVIII-e:\allowbreak{}72}
\DERapparatusitem{ENV1531}{181;\allowbreak{}\allowbreak{} ms467,\allowbreak{}presentation\allowbreak{} clause}{\_\allowbreak{}0}{}{DER\_\allowbreak{}XVIII-e:\allowbreak{}73}
\DERapparatusitem{ENV1532}{181;\allowbreak{}\allowbreak{} ms468,\allowbreak{}4.\allowbreak{}16.\allowbreak{}2\allowbreak{} stability\allowbreak{} hypothesis}{\$\allowbreak{}\textbackslash{}\allowbreak{}mathcal\{\allowbreak{}N\}\allowbreak{}\$\allowbreak{}\allowbreak{} stable\allowbreak{} under\allowbreak{} \$\allowbreak{}\textbackslash{}\allowbreak{}varprojlim\$\allowbreak{}\allowbreak{} large\allowbreak{} relative\allowbreak{} to\allowbreak{} \$\allowbreak{}\textbackslash{}\allowbreak{}pi\_\allowbreak{}1\$\allowbreak{})\allowbreak{}.\allowbreak{}}{\$\allowbreak{}\textbackslash{}\allowbreak{}mathcal\{\allowbreak{}N\}\allowbreak{}\$\allowbreak{}\allowbreak{} stable\allowbreak{} under\allowbreak{} \$\allowbreak{}\textbackslash{}\allowbreak{}varprojlim\$\allowbreak{}\allowbreak{} large\allowbreak{} relative\allowbreak{} to\allowbreak{} \$\allowbreak{}\textbackslash{}\allowbreak{}pi\$\allowbreak{})\allowbreak{},\allowbreak{}\allowbreak{} and\allowbreak{} \$\allowbreak{}\textbackslash{}\allowbreak{}underline\{\allowbreak{}T\}\allowbreak{}\$\allowbreak{}\allowbreak{} is\allowbreak{} stable\allowbreak{} in\allowbreak{} \$\allowbreak{}\textbackslash{}\allowbreak{}mathcal\{\allowbreak{}N\}\allowbreak{}\$\allowbreak{}\allowbreak{} under\allowbreak{} these\allowbreak{} same\allowbreak{} limits.\allowbreak{}}{DER\_\allowbreak{}XVIII-e:\allowbreak{}76;\allowbreak{}\allowbreak{} DER\_\allowbreak{}XVIII-e:\allowbreak{}77}
\DERapparatusitem{ENV1533}{181;\allowbreak{}\allowbreak{} ms468,\allowbreak{}footnote191}{I}{\textbackslash{}\allowbreak{}underline\allowbreak{} T}{DER\_\allowbreak{}XVIII-e:\allowbreak{}82}
\DERapparatusitem{ENV1534}{181;\allowbreak{}\allowbreak{} ms468-469,\allowbreak{}4.\allowbreak{}16.\allowbreak{}2\allowbreak{} test\allowbreak{} cases}{\textbackslash{}\allowbreak{}mathcal\{\allowbreak{}N\}\allowbreak{}}{\textbackslash{}\allowbreak{}underline\allowbreak{} T}{DER\_\allowbreak{}XVIII-e:\allowbreak{}78}
\DERapparatusitem{ENV1535}{181;\allowbreak{}\allowbreak{} ms468-469,\allowbreak{}4.\allowbreak{}16.\allowbreak{}2\allowbreak{} test\allowbreak{} cases}{\textbackslash{}\allowbreak{}mathcal\{\allowbreak{}N\}\allowbreak{}}{\textbackslash{}\allowbreak{}underline\allowbreak{} T}{DER\_\allowbreak{}XVIII-e:\allowbreak{}78}
\DERapparatusitem{ENV1536}{181;\allowbreak{}\allowbreak{} ms469,\allowbreak{}4.\allowbreak{}16.\allowbreak{}2\allowbreak{} NB}{They\allowbreak{} are\allowbreak{} to\allowbreak{} ensure\allowbreak{} that}{It\allowbreak{} serves\allowbreak{} to\allowbreak{} ensure\allowbreak{} that}{DER\_\allowbreak{}XVIII-e:\allowbreak{}79}
\DERapparatusitem{ENV1537}{181;\allowbreak{}\allowbreak{} ms469,\allowbreak{}4.\allowbreak{}16.\allowbreak{}3\allowbreak{} heading}{Related\allowbreak{} questions:\allowbreak{}}{Related\allowbreak{} question:\allowbreak{}}{DER\_\allowbreak{}XVIII-e:\allowbreak{}80}
\DERapparatusitem{ENV1538}{181;\allowbreak{}\allowbreak{} ms469,\allowbreak{}4.\allowbreak{}16.\allowbreak{}3\allowbreak{} image\allowbreak{} ambient\allowbreak{} category}{N}{M}{DER\_\allowbreak{}XVIII-e:\allowbreak{}81}
\DERapparatusitem{ENV1539}{183;\allowbreak{}\allowbreak{} ms473,\allowbreak{}second\allowbreak{} diagram\allowbreak{} label}{}{\_\allowbreak{}\textbackslash{}\allowbreak{}bullet}{DER\_\allowbreak{}XVIII-e:\allowbreak{}84}
\DERapparatusitem{ENV1540}{183;\allowbreak{}\allowbreak{} ms473,\allowbreak{}Proposition4.\allowbreak{}16.\allowbreak{}6.\allowbreak{}1\allowbreak{} reference}{of\allowbreak{} 4.\allowbreak{}15.\allowbreak{}1\allowbreak{} (\allowbreak{}page\allowbreak{} 415)\allowbreak{}}{of\allowbreak{} 4.\allowbreak{}15.\allowbreak{}2\allowbreak{} (\allowbreak{}page\allowbreak{} 415)\allowbreak{}}{DER\_\allowbreak{}XVIII-e:\allowbreak{}85}
\DERapparatusitem{ENV1541}{184;\allowbreak{}\allowbreak{} ms475,\allowbreak{}Corollary4.\allowbreak{}16.\allowbreak{}6.\allowbreak{}2(\allowbreak{}b)\allowbreak{}}{defined\allowbreak{} in\allowbreak{} 4.\allowbreak{}15.\allowbreak{}6.\allowbreak{}1\allowbreak{} above}{defined\allowbreak{} in\allowbreak{} 4.\allowbreak{}16.\allowbreak{}6.\allowbreak{}1\allowbreak{} above}{DER\_\allowbreak{}XVIII-e:\allowbreak{}87}
\DERapparatusitem{ENV1542}{184;\allowbreak{}\allowbreak{} ms476,\allowbreak{}4.\allowbreak{}16.\allowbreak{}6.\allowbreak{}5\allowbreak{} reference}{conditions\allowbreak{} of\allowbreak{} 4.\allowbreak{}16.\allowbreak{}6.\allowbreak{}1\allowbreak{} --}{conditions\allowbreak{} of\allowbreak{} 4.\allowbreak{}16.\allowbreak{}6.\allowbreak{}2\allowbreak{} --}{DER\_\allowbreak{}XVIII-e:\allowbreak{}88}
\DERapparatusitem{ENV1543}{185;\allowbreak{}\allowbreak{} ms477,\allowbreak{}union\allowbreak{} formula}{\textbackslash{}\allowbreak{}uS}{\textbackslash{}\allowbreak{}Ob\textbackslash{}\allowbreak{}underline\allowbreak{} S}{DER\_\allowbreak{}XVIII-e:\allowbreak{}92}
\DERapparatusitem{ENV1544}{187;\allowbreak{}\allowbreak{} ms483,\allowbreak{}first\allowbreak{} parenthesis}{(\allowbreak{}card\allowbreak{} \$\allowbreak{}x\_\allowbreak{}\textbackslash{}\allowbreak{}alpha\allowbreak{} >}{(\allowbreak{}since\allowbreak{} \$\allowbreak{}x\_\allowbreak{}\textbackslash{}\allowbreak{}alpha\allowbreak{} >}{DER\_\allowbreak{}XVIII-e:\allowbreak{}96}
\DERapparatusitem{ENV1545}{187;\allowbreak{}\allowbreak{} ms483,\allowbreak{}after\allowbreak{} display(\allowbreak{}*)\allowbreak{}}{s}{x}{DER\_\allowbreak{}XVIII-e:\allowbreak{}97}
\DERapparatusitem{ENV1546}{187;\allowbreak{}\allowbreak{} ms483,\allowbreak{}Proposition4.\allowbreak{}16.\allowbreak{}7.\allowbreak{}5\allowbreak{} universe}{\textbackslash{}\allowbreak{}mathcal}{\textbackslash{}\allowbreak{}mathbf}{DER\_\allowbreak{}XVIII-e:\allowbreak{}98}
\DERapparatusitem{ENV1547}{187;\allowbreak{}\allowbreak{} ms483,\allowbreak{}Proposition4.\allowbreak{}16.\allowbreak{}7.\allowbreak{}5\allowbreak{} first\allowbreak{} label}{\textbackslash{}\allowbreak{}item[\allowbreak{}1)\allowbreak{}]\allowbreak{}\allowbreak{} There\allowbreak{} exists\allowbreak{} a}{\textbackslash{}\allowbreak{}item[\allowbreak{}a)\allowbreak{}]\allowbreak{}\allowbreak{} There\allowbreak{} exists\allowbreak{} a}{DER\_\allowbreak{}XVIII-e:\allowbreak{}99}
\DERapparatusitem{ENV1548}{187;\allowbreak{}\allowbreak{} running\allowbreak{} text\allowbreak{} or\allowbreak{} formula}{\textbackslash{}\allowbreak{}text\{\allowbreak{}\textbackslash{}\allowbreak{}upshape\allowbreak{} def\}\allowbreak{}}{\textbackslash{}\allowbreak{}mathrm\{\allowbreak{}def\}\allowbreak{}}{TYPESETTING:\allowbreak{}7;\allowbreak{}\allowbreak{} NOTATION\_\allowbreak{}DER:\allowbreak{}17}
\DERapparatusitem{ENV1549}{187;\allowbreak{}\allowbreak{} ms484,\allowbreak{}presheaf\allowbreak{} formula}{}{\textasciicircum{}\allowbreak{}\{\allowbreak{}\textbackslash{}\allowbreak{}circ\}\allowbreak{}}{DER\_\allowbreak{}XVIII-e:\allowbreak{}100}
\DERapparatusitem{ENV1550}{188;\allowbreak{}\allowbreak{} running\allowbreak{} text\allowbreak{} or\allowbreak{} formula}{\textbackslash{}\allowbreak{}text\{\allowbreak{}\textbackslash{}\allowbreak{}upshape\allowbreak{} def\}\allowbreak{}}{\textbackslash{}\allowbreak{}mathrm\{\allowbreak{}def\}\allowbreak{}}{TYPESETTING:\allowbreak{}7;\allowbreak{}\allowbreak{} NOTATION\_\allowbreak{}DER:\allowbreak{}17}
\DERapparatusitem{ENV1551}{preamble;\allowbreak{}\allowbreak{} local\allowbreak{} macro\allowbreak{} definition}{\textbackslash{}\allowbreak{}text\{\allowbreak{}def\}\allowbreak{}}{\textbackslash{}\allowbreak{}mathrm\{\allowbreak{}def\}\allowbreak{}}{TYPESETTING:\allowbreak{}7;\allowbreak{}\allowbreak{} NOTATION\_\allowbreak{}DER:\allowbreak{}17}
\DERapparatusitem{ENV1552}{189;\allowbreak{}\allowbreak{} ms489,\allowbreak{}Lemma4.\allowbreak{}16.\allowbreak{}7.\allowbreak{}11(\allowbreak{}a)\allowbreak{}}{\textbackslash{}\allowbreak{}psi}{\textbackslash{}\allowbreak{}varphi}{DER\_\allowbreak{}XVIII-e:\allowbreak{}102}
\DERapparatusitem{ENV1553}{189;\allowbreak{}\allowbreak{} ms489,\allowbreak{}proof\allowbreak{} first\allowbreak{} sentence}{\textbackslash{}\allowbreak{}underline\{\allowbreak{}S\}\allowbreak{}}{\textbackslash{}\allowbreak{}varphi}{DER\_\allowbreak{}XVIII-e:\allowbreak{}103}
\DERapparatusitem{ENV1554}{189;\allowbreak{}\allowbreak{} ms489,\allowbreak{}proof\allowbreak{} Ind\allowbreak{} formula}{}{\_\allowbreak{}\textbackslash{}\allowbreak{}pi}{DER\_\allowbreak{}XVIII-e:\allowbreak{}104}
\DERapparatusitem{ENV1555}{190;\allowbreak{}\allowbreak{} ms490,\allowbreak{}cofinal\allowbreak{} colimit\allowbreak{} denominator}{\textbackslash{}\allowbreak{}pi}{x}{DER\_\allowbreak{}XVIII-e:\allowbreak{}105}
\DERapparatusitem{ENV1556}{190;\allowbreak{}\allowbreak{} ms490,\allowbreak{}following\allowbreak{} sentence}{\textbackslash{}\allowbreak{}pi}{\textbackslash{}\allowbreak{}pi/\allowbreak{}x}{DER\_\allowbreak{}XVIII-e:\allowbreak{}106}
\DERapparatusitem{ENV1557}{190;\allowbreak{}\allowbreak{} ms491,\allowbreak{}ambient\allowbreak{} category\allowbreak{} of\allowbreak{} cosieve}{\textbackslash{}\allowbreak{}underline\{\allowbreak{}S\}\allowbreak{}(\allowbreak{}\textbackslash{}\allowbreak{}pi)\allowbreak{}}{\textbackslash{}\allowbreak{}mathcal\allowbreak{} M\_\allowbreak{}\textbackslash{}\allowbreak{}pi/\allowbreak{}x}{DER\_\allowbreak{}XVIII-e:\allowbreak{}107}
\DERapparatusitem{ENV1558}{190;\allowbreak{}\allowbreak{} ms491,\allowbreak{}second\allowbreak{} underbrace}{\$\allowbreak{}\textbackslash{}\allowbreak{}underline\{\allowbreak{}T\}\allowbreak{}(\allowbreak{}\textbackslash{}\allowbreak{}pi)\allowbreak{}/\allowbreak{}x\$\allowbreak{}\allowbreak{} cofinal\allowbreak{} in\allowbreak{} \$\allowbreak{}\textbackslash{}\allowbreak{}cM\_\allowbreak{}\textbackslash{}\allowbreak{}pi\$\allowbreak{},\allowbreak{}}{\$\allowbreak{}\textbackslash{}\allowbreak{}underline\{\allowbreak{}T\}\allowbreak{}(\allowbreak{}\textbackslash{}\allowbreak{}pi)\allowbreak{}/\allowbreak{}x\$\allowbreak{}\allowbreak{} cofinal\allowbreak{} in\allowbreak{} \$\allowbreak{}\textbackslash{}\allowbreak{}cM\_\allowbreak{}\textbackslash{}\allowbreak{}pi/\allowbreak{}x\$\allowbreak{},\allowbreak{}}{DER\_\allowbreak{}XVIII-e:\allowbreak{}108}
\DERapparatusitem{ENV1559}{190;\allowbreak{}\allowbreak{} ms491,\allowbreak{}final\allowbreak{} paragraph\allowbreak{} membership\allowbreak{} criterion}{}{/\allowbreak{}x}{DER\_\allowbreak{}XVIII-e:\allowbreak{}109}
\DERapparatusitem{ENV1560}{190;\allowbreak{}\allowbreak{} ms491,\allowbreak{}final\allowbreak{} reference}{none\allowbreak{} other\allowbreak{} than\allowbreak{} 4.\allowbreak{}6.\allowbreak{}7.\allowbreak{}6}{none\allowbreak{} other\allowbreak{} than\allowbreak{} 4.\allowbreak{}16.\allowbreak{}7.\allowbreak{}6}{DER\_\allowbreak{}XVIII-e:\allowbreak{}110}
\DERapparatusitem{ENV1561}{192;\allowbreak{}\allowbreak{} Appendix\allowbreak{} ms.\allowbreak{}3,\allowbreak{}\allowbreak{} sufficiency\allowbreak{} in\allowbreak{} Proposition1}{\textbackslash{}\allowbreak{}neq\textbackslash{}\allowbreak{}emptyset}{\textbackslash{}\allowbreak{}neq0}{DER\_\allowbreak{}XVIII-e:\allowbreak{}114}
\DERapparatusitem{ENV1562}{192;\allowbreak{}\allowbreak{} Appendix\allowbreak{} ms.\allowbreak{}3,\allowbreak{}\allowbreak{} last\allowbreak{} sentence}{\textbackslash{}\allowbreak{}neq\textbackslash{}\allowbreak{}emptyset}{\textbackslash{}\allowbreak{}neq0}{DER\_\allowbreak{}XVIII-e:\allowbreak{}115}
\DERapparatusitem{ENV1563}{192;\allowbreak{}\allowbreak{} Appendix\allowbreak{} ms.\allowbreak{}4,\allowbreak{}\allowbreak{} definition\allowbreak{} of\allowbreak{} A}{\textbackslash{}\allowbreak{}emptyset}{0}{DER\_\allowbreak{}XVIII-e:\allowbreak{}116}
\DERapparatusitem{ENV1564}{192;\allowbreak{}\allowbreak{} Appendix\allowbreak{} ms.\allowbreak{}4,\allowbreak{}\allowbreak{} definition\allowbreak{} of\allowbreak{} B}{\textbackslash{}\allowbreak{}neq\textbackslash{}\allowbreak{}emptyset}{=\allowbreak{}0}{DER\_\allowbreak{}XVIII-e:\allowbreak{}117}
\DERapparatusitem{ENV1565}{192;\allowbreak{}\allowbreak{} Appendix\allowbreak{} ms.\allowbreak{}4,\allowbreak{}\allowbreak{} definition\allowbreak{} of\allowbreak{} g\allowbreak{} prime}{f}{g}{DER\_\allowbreak{}XVIII-e:\allowbreak{}118}
\DERapparatusitem{ENV1566}{192;\allowbreak{}\allowbreak{} Appendix\allowbreak{} ms.\allowbreak{}4,\allowbreak{}\allowbreak{} inverse\allowbreak{} identity}{=\allowbreak{}}{}{DER\_\allowbreak{}XVIII-e:\allowbreak{}119}
\DERapparatusitem{ENV1567}{192;\allowbreak{}\allowbreak{} Appendix\allowbreak{} ms.\allowbreak{}4,\allowbreak{}\allowbreak{} final\allowbreak{} sentence}{\textbackslash{}\allowbreak{}Phi\textbackslash{}\allowbreak{}in\textbackslash{}\allowbreak{}cF}{\textbackslash{}\allowbreak{}Phi\textbackslash{}\allowbreak{}in\textbackslash{}\allowbreak{}cF\_\allowbreak{}0}{DER\_\allowbreak{}XVIII-e:\allowbreak{}120}
\DERapparatusitem{ENV1568}{193;\allowbreak{}\allowbreak{} Appendix\allowbreak{} ms.\allowbreak{}6,\allowbreak{}\allowbreak{} verification\allowbreak{} a)\allowbreak{}}{k\_\allowbreak{}\textbackslash{}\allowbreak{}Phi}{\textbackslash{}\allowbreak{}kappa(\allowbreak{}\textbackslash{}\allowbreak{}mathfrak\allowbreak{} m)\allowbreak{}=\allowbreak{}k\_\allowbreak{}E/\allowbreak{}\textbackslash{}\allowbreak{}mathfrak\allowbreak{} m}{DER\_\allowbreak{}XVIII-e:\allowbreak{}121}
\DERapparatusitem{ENV1569}{193;\allowbreak{}\allowbreak{} Appendix\allowbreak{} ms.\allowbreak{}6,\allowbreak{}\allowbreak{} following\allowbreak{} verification\allowbreak{} d)\allowbreak{}}{A}{k\_\allowbreak{}E}{DER\_\allowbreak{}XVIII-e:\allowbreak{}122}
\DERapparatusitem{ENV1570}{193;\allowbreak{}\allowbreak{} Appendix\allowbreak{} ms.\allowbreak{}6,\allowbreak{}\allowbreak{} following\allowbreak{} verification\allowbreak{} d)\allowbreak{}}{A}{k\_\allowbreak{}E}{DER\_\allowbreak{}XVIII-e:\allowbreak{}122}
\DERapparatusitem{ENV1571}{194;\allowbreak{}\allowbreak{} Appendix\allowbreak{} ms.\allowbreak{}7,\allowbreak{}\allowbreak{} support-idempotent\allowbreak{} lemma}{A}{B}{DER\_\allowbreak{}XVIII-e:\allowbreak{}123}
\DERapparatusitem{ENV1572}{194;\allowbreak{}\allowbreak{} Appendix\allowbreak{} ms.\allowbreak{}8,\allowbreak{}\allowbreak{} after\allowbreak{} factorization\allowbreak{} diagram}{}{mathfrak\{\allowbreak{}m\}\allowbreak{}\_\allowbreak{}\textbackslash{}\allowbreak{}}{DER\_\allowbreak{}XVIII-e:\allowbreak{}124}
\DERapparatusitem{ENV1573}{195;\allowbreak{}\allowbreak{} Appendix\allowbreak{} ms.\allowbreak{}10,\allowbreak{}\allowbreak{} proof\allowbreak{} of\allowbreak{} Proposition3}{f}{k}{DER\_\allowbreak{}XVIII-e:\allowbreak{}125}
\DERapparatusitem{ENV1574}{195;\allowbreak{}\allowbreak{} Appendix\allowbreak{} ms.\allowbreak{}11,\allowbreak{}\allowbreak{} first\allowbreak{} sentence}{E}{\textbackslash{}\allowbreak{}Phi}{DER\_\allowbreak{}XVIII-e:\allowbreak{}126}
\DERapparatusitem{ENV1575}{196;\allowbreak{}\allowbreak{} running\allowbreak{} text\allowbreak{} or\allowbreak{} formula}{\textbackslash{}\allowbreak{}text\{\allowbreak{}def\}\allowbreak{}}{\textbackslash{}\allowbreak{}mathrm\{\allowbreak{}def\}\allowbreak{}}{TYPESETTING:\allowbreak{}7;\allowbreak{}\allowbreak{} NOTATION\_\allowbreak{}DER:\allowbreak{}17}
\DERapparatusitem{ENV1576}{196;\allowbreak{}\allowbreak{} running\allowbreak{} text\allowbreak{} or\allowbreak{} formula}{\textbackslash{}\allowbreak{}text\{\allowbreak{}def\}\allowbreak{}}{\textbackslash{}\allowbreak{}mathrm\{\allowbreak{}def\}\allowbreak{}}{TYPESETTING:\allowbreak{}7;\allowbreak{}\allowbreak{} NOTATION\_\allowbreak{}DER:\allowbreak{}17}
\DERapparatusitem{ENV1577}{197;\allowbreak{}\allowbreak{} Appendix\allowbreak{} ms.\allowbreak{}14,\allowbreak{}\allowbreak{} definition\allowbreak{} of\allowbreak{} the\allowbreak{} complementary\allowbreak{} open\allowbreak{} category}{=\allowbreak{}}{\textbackslash{}\allowbreak{}neq}{DER\_\allowbreak{}XVIII-e:\allowbreak{}128}
\DERapparatusitem{ENV1578}{197;\allowbreak{}\allowbreak{} Appendix\allowbreak{} ms.\allowbreak{}15,\allowbreak{}\allowbreak{} displayed\allowbreak{} formula\allowbreak{} (\allowbreak{}*)\allowbreak{}}{\{\allowbreak{}\textbackslash{}\allowbreak{}xi\textbackslash{}\allowbreak{}\}\allowbreak{}=\allowbreak{}}{xi\textbackslash{}\allowbreak{}in}{DER\_\allowbreak{}XVIII-e:\allowbreak{}129}
\DERapparatusitem{ENV1579}{197;\allowbreak{}\allowbreak{} Appendix\allowbreak{} ms.\allowbreak{}15,\allowbreak{}\allowbreak{} sentence\allowbreak{} following\allowbreak{} (\allowbreak{}*)\allowbreak{}}{\textbackslash{}\allowbreak{}\{\allowbreak{}\textbackslash{}\allowbreak{}xi\textbackslash{}\allowbreak{}\}\allowbreak{}}{\textbackslash{}\allowbreak{}xi}{DER\_\allowbreak{}XVIII-e:\allowbreak{}130}
\DERapparatusitem{ENV1580}{198;\allowbreak{}\allowbreak{} Appendix\allowbreak{} ms.\allowbreak{}17,\allowbreak{}\allowbreak{} definition\allowbreak{} of\allowbreak{} I\_\allowbreak{}\{\allowbreak{}x,\allowbreak{}S\allowbreak{} prime\}\allowbreak{}}{}{'}{DER\_\allowbreak{}XVIII-e:\allowbreak{}133}
\subsection*{Chapter XIX}
\DERapparatusitem{ENV1581}{1;\allowbreak{}\allowbreak{} ms.\allowbreak{}\allowbreak{} 1,\allowbreak{}\allowbreak{} prose\allowbreak{} after\allowbreak{} (\allowbreak{}1.\allowbreak{}3)\allowbreak{}}{These\allowbreak{} three\allowbreak{} cases\allowbreak{} of\allowbreak{} strengthened\allowbreak{} Q-orthogonality\allowbreak{} are\allowbreak{} thus\allowbreak{} those\allowbreak{} where\allowbreak{} we\allowbreak{} have\allowbreak{} in\allowbreak{} addition}{These\allowbreak{} three\allowbreak{} cases\allowbreak{} of\allowbreak{} strengthened\allowbreak{} Q-orthogonality\allowbreak{} thus\allowbreak{} respectively\allowbreak{} imply}{DER\_\allowbreak{}XIX:\allowbreak{}4}
\DERapparatusitem{ENV1582}{1;\allowbreak{}\allowbreak{} ms.\allowbreak{}\allowbreak{} 2,\allowbreak{}\allowbreak{} second\allowbreak{} displayed\allowbreak{} group\allowbreak{} of\allowbreak{} factorization\allowbreak{} arrows}{3}{5}{DER\_\allowbreak{}XIX:\allowbreak{}7}
\DERapparatusitem{ENV1583}{3;\allowbreak{}\allowbreak{} running\allowbreak{} text\allowbreak{} or\allowbreak{} formula}{\textbackslash{}\allowbreak{}text\{\allowbreak{}def\}\allowbreak{}}{\textbackslash{}\allowbreak{}mathrm\{\allowbreak{}def\}\allowbreak{}}{TYPESETTING:\allowbreak{}7;\allowbreak{}\allowbreak{} NOTATION\_\allowbreak{}DER:\allowbreak{}17}
\DERapparatusitem{ENV1584}{3;\allowbreak{}\allowbreak{} running\allowbreak{} text\allowbreak{} or\allowbreak{} formula}{\textbackslash{}\allowbreak{}text\{\allowbreak{}def\}\allowbreak{}}{\textbackslash{}\allowbreak{}mathrm\{\allowbreak{}def\}\allowbreak{}}{TYPESETTING:\allowbreak{}7;\allowbreak{}\allowbreak{} NOTATION\_\allowbreak{}DER:\allowbreak{}17}
\DERapparatusitem{ENV1585}{5;\allowbreak{}\allowbreak{} ms.\allowbreak{}\allowbreak{} 7,\allowbreak{}\allowbreak{} Corollary\allowbreak{} 1.\allowbreak{}7\allowbreak{} wording}{}{rm\{\allowbreak{}Fl\}\allowbreak{}(\allowbreak{}\textbackslash{}\allowbreak{}math}{DER\_\allowbreak{}XIX:\allowbreak{}10}
\DERapparatusitem{ENV1586}{5;\allowbreak{}\allowbreak{} ms.\allowbreak{}\allowbreak{} 7,\allowbreak{}\allowbreak{} Corollary\allowbreak{} 1.\allowbreak{}7\allowbreak{} wording}{}{)\allowbreak{}}{DER\_\allowbreak{}XIX:\allowbreak{}10}
\DERapparatusitem{ENV1587}{5;\allowbreak{}\allowbreak{} ms.\allowbreak{}\allowbreak{} 7,\allowbreak{}\allowbreak{} proof\allowbreak{} of\allowbreak{} Corollary\allowbreak{} 1.\allowbreak{}7}{\textbackslash{}\allowbreak{}in\allowbreak{} \textbackslash{}\allowbreak{}Psi}{\textbackslash{}\allowbreak{}in\textbackslash{}\allowbreak{}Psi'}{DER\_\allowbreak{}XIX:\allowbreak{}11}
\DERapparatusitem{ENV1588}{5;\allowbreak{}\allowbreak{} running\allowbreak{} text\allowbreak{} or\allowbreak{} formula}{\textbackslash{}\allowbreak{}text\{\allowbreak{}def\}\allowbreak{}}{\textbackslash{}\allowbreak{}mathrm\{\allowbreak{}def\}\allowbreak{}}{TYPESETTING:\allowbreak{}7;\allowbreak{}\allowbreak{} NOTATION\_\allowbreak{}DER:\allowbreak{}17}
\DERapparatusitem{ENV1589}{6;\allowbreak{}\allowbreak{} running\allowbreak{} text\allowbreak{} or\allowbreak{} formula}{\textbackslash{}\allowbreak{}text\{\allowbreak{}def\}\allowbreak{}}{\textbackslash{}\allowbreak{}mathrm\{\allowbreak{}def\}\allowbreak{}}{TYPESETTING:\allowbreak{}7;\allowbreak{}\allowbreak{} NOTATION\_\allowbreak{}DER:\allowbreak{}17}
\DERapparatusitem{ENV1590}{6;\allowbreak{}\allowbreak{} ms.\allowbreak{}\allowbreak{} 10,\allowbreak{}\allowbreak{} equation\allowbreak{} (\allowbreak{}1.\allowbreak{}14)\allowbreak{}}{Suppose\allowbreak{} now\allowbreak{} that\allowbreak{} one\allowbreak{} has}{Suppose\allowbreak{} now\allowbreak{} that\allowbreak{} \$\allowbreak{}\textbackslash{}\allowbreak{}Phi\$\allowbreak{}\allowbreak{} contains\allowbreak{} the\allowbreak{} identities\allowbreak{} and\allowbreak{} that\allowbreak{} one\allowbreak{} has}{DER\_\allowbreak{}XIX:\allowbreak{}12}
\DERapparatusitem{ENV1591}{6;\allowbreak{}\allowbreak{} running\allowbreak{} text\allowbreak{} or\allowbreak{} formula}{\textbackslash{}\allowbreak{}text\{\allowbreak{}def\}\allowbreak{}}{\textbackslash{}\allowbreak{}mathrm\{\allowbreak{}def\}\allowbreak{}}{TYPESETTING:\allowbreak{}7;\allowbreak{}\allowbreak{} NOTATION\_\allowbreak{}DER:\allowbreak{}17}
\DERapparatusitem{ENV1592}{8;\allowbreak{}\allowbreak{} ms.\allowbreak{}\allowbreak{} 12,\allowbreak{}\allowbreak{} formula\allowbreak{} (\allowbreak{}1.\allowbreak{}13.\allowbreak{}2)\allowbreak{}}{\textbackslash{}\allowbreak{}Phi(\allowbreak{}I')\allowbreak{})\allowbreak{}\allowbreak{} \textbackslash{}\allowbreak{}subset\allowbreak{} \textbackslash{}\allowbreak{}Phi\textbackslash{}\allowbreak{}langle\allowbreak{} I'}{\textbackslash{}\allowbreak{}Phi\textbackslash{}\allowbreak{}langle\allowbreak{} I'\textbackslash{}\allowbreak{}rangle)\allowbreak{}\textbackslash{}\allowbreak{}subset\textbackslash{}\allowbreak{}Phi\textbackslash{}\allowbreak{}langle\allowbreak{} I}{DER\_\allowbreak{}XIX:\allowbreak{}13}
\DERapparatusitem{ENV1593}{9;\allowbreak{}\allowbreak{} ms.\allowbreak{}\allowbreak{} 14,\allowbreak{}\allowbreak{} dual\allowbreak{} clause\allowbreak{} of\allowbreak{} (\allowbreak{}2.\allowbreak{}5)\allowbreak{}}{W}{C}{DER\_\allowbreak{}XIX:\allowbreak{}14}
\DERapparatusitem{ENV1594}{9;\allowbreak{}\allowbreak{} ms.\allowbreak{}\allowbreak{} 14,\allowbreak{}\allowbreak{} dual\allowbreak{} clause\allowbreak{} of\allowbreak{} (\allowbreak{}2.\allowbreak{}5)\allowbreak{}}{f}{p}{DER\_\allowbreak{}XIX:\allowbreak{}14}
\DERapparatusitem{ENV1595}{9;\allowbreak{}\allowbreak{} ms.\allowbreak{}\allowbreak{} 14,\allowbreak{}\allowbreak{} dual\allowbreak{} clause\allowbreak{} of\allowbreak{} (\allowbreak{}2.\allowbreak{}5)\allowbreak{}}{C}{W}{DER\_\allowbreak{}XIX:\allowbreak{}14}
\DERapparatusitem{ENV1596}{9;\allowbreak{}\allowbreak{} ms.\allowbreak{}\allowbreak{} 14,\allowbreak{}\allowbreak{} definition\allowbreak{} of\allowbreak{} special\allowbreak{} triple}{the\allowbreak{} two\allowbreak{} Quillen\allowbreak{} couples\allowbreak{} (\allowbreak{}2.\allowbreak{}3)\allowbreak{}\allowbreak{} are\allowbreak{} special}{the\allowbreak{} two\allowbreak{} Quillen\allowbreak{} couples\allowbreak{} (\allowbreak{}2.\allowbreak{}2)\allowbreak{}\allowbreak{} are\allowbreak{} special}{DER\_\allowbreak{}XIX:\allowbreak{}15}
\DERapparatusitem{ENV1597}{9;\allowbreak{}\allowbreak{} running\allowbreak{} text\allowbreak{} or\allowbreak{} formula}{\textbackslash{}\allowbreak{}text\{\allowbreak{}def\}\allowbreak{}}{\textbackslash{}\allowbreak{}mathrm\{\allowbreak{}def\}\allowbreak{}}{TYPESETTING:\allowbreak{}7;\allowbreak{}\allowbreak{} NOTATION\_\allowbreak{}DER:\allowbreak{}17}
\DERapparatusitem{ENV1598}{10;\allowbreak{}\allowbreak{} ms.\allowbreak{}\allowbreak{} 15,\allowbreak{}\allowbreak{} adjunction\allowbreak{} label\allowbreak{} preceding\allowbreak{} (\allowbreak{}2.\allowbreak{}1.\allowbreak{}2)\allowbreak{}}{}{.\allowbreak{}}{DER\_\allowbreak{}XIX:\allowbreak{}17}
\DERapparatusitem{ENV1599}{11;\allowbreak{}\allowbreak{} ms.\allowbreak{}\allowbreak{} 17,\allowbreak{}\allowbreak{} proof\allowbreak{} of\allowbreak{} Corollary\allowbreak{} 2.\allowbreak{}2,\allowbreak{}\allowbreak{} condition\allowbreak{} reference}{(\allowbreak{}2.\allowbreak{}1.\allowbreak{}6.\allowbreak{}\allowbreak{} b\$\allowbreak{}''\$\allowbreak{})\allowbreak{}}{(\allowbreak{}2.\allowbreak{}1.\allowbreak{}4.\allowbreak{}\allowbreak{} a)\allowbreak{}}{DER\_\allowbreak{}XIX:\allowbreak{}20}
\DERapparatusitem{ENV1600}{11;\allowbreak{}\allowbreak{} ms.\allowbreak{}\allowbreak{} 17,\allowbreak{}\allowbreak{} proof\allowbreak{} of\allowbreak{} Corollary\allowbreak{} 2.\allowbreak{}2,\allowbreak{}\allowbreak{} first\allowbreak{} pair}{TC}{TF}{DER\_\allowbreak{}XIX:\allowbreak{}21}
\DERapparatusitem{ENV1601}{11;\allowbreak{}\allowbreak{} ms.\allowbreak{}\allowbreak{} 17,\allowbreak{}\allowbreak{} first\allowbreak{} display\allowbreak{} in\allowbreak{} Case\allowbreak{} A}{'}{}{DER\_\allowbreak{}XIX:\allowbreak{}22}
\DERapparatusitem{ENV1602}{11;\allowbreak{}\allowbreak{} running\allowbreak{} text\allowbreak{} or\allowbreak{} formula}{\textbackslash{}\allowbreak{}text\{\allowbreak{}def\}\allowbreak{}}{\textbackslash{}\allowbreak{}mathrm\{\allowbreak{}def\}\allowbreak{}}{TYPESETTING:\allowbreak{}7;\allowbreak{}\allowbreak{} NOTATION\_\allowbreak{}DER:\allowbreak{}17}
\DERapparatusitem{ENV1603}{12;\allowbreak{}\allowbreak{} ms.\allowbreak{}19,\allowbreak{}\allowbreak{} Theorem\allowbreak{} 2.\allowbreak{}3,\allowbreak{}\allowbreak{} definition\allowbreak{} of\allowbreak{} C<I>}{}{T}{DER\_\allowbreak{}XIX:\allowbreak{}25}
\DERapparatusitem{ENV1604}{13;\allowbreak{}\allowbreak{} ms.\allowbreak{}20,\allowbreak{}\allowbreak{} end\allowbreak{} of\allowbreak{} Theorem\allowbreak{} 2.\allowbreak{}3(\allowbreak{}2)\allowbreak{}}{\textbackslash{}\allowbreak{}textit\{\allowbreak{}induced\allowbreak{} by\allowbreak{} \$\allowbreak{}u\textasciicircum{}\allowbreak{}*\$\allowbreak{}\}\allowbreak{}\allowbreak{} (\allowbreak{}3.\allowbreak{}3.\allowbreak{}1)\allowbreak{}}{\textbackslash{}\allowbreak{}textit\{\allowbreak{}induced\allowbreak{} by\allowbreak{} \$\allowbreak{}u\textasciicircum{}\allowbreak{}*\$\allowbreak{}\}\allowbreak{}\allowbreak{} (\allowbreak{}2.\allowbreak{}3.\allowbreak{}1)\allowbreak{}}{DER\_\allowbreak{}XIX:\allowbreak{}27}
\DERapparatusitem{ENV1605}{preamble;\allowbreak{}\allowbreak{} local\allowbreak{} macro\allowbreak{} definition}{\textbackslash{}\allowbreak{}text\{\allowbreak{}def\}\allowbreak{}}{\textbackslash{}\allowbreak{}mathrm\{\allowbreak{}def\}\allowbreak{}}{TYPESETTING:\allowbreak{}7;\allowbreak{}\allowbreak{} NOTATION\_\allowbreak{}DER:\allowbreak{}17}
\DERapparatusitem{ENV1606}{14;\allowbreak{}\allowbreak{} ms.\allowbreak{}23,\allowbreak{}\allowbreak{} first\allowbreak{} proof\allowbreak{} paragraph}{The\allowbreak{} left\allowbreak{} exactness\allowbreak{} axiom\allowbreak{} Der\allowbreak{} 5\allowbreak{} bis}{The\allowbreak{} left\allowbreak{} exactness\allowbreak{} axiom\allowbreak{} Der\allowbreak{} 6\allowbreak{} bis}{DER\_\allowbreak{}XIX:\allowbreak{}29}
\DERapparatusitem{ENV1607}{15;\allowbreak{}\allowbreak{} running\allowbreak{} text\allowbreak{} or\allowbreak{} formula}{Filtered\allowbreak{} inductive\allowbreak{} limits}{Filtering\allowbreak{} direct\allowbreak{} limits}{GLOSSARY\_\allowbreak{}DER:\allowbreak{}89;\allowbreak{}\allowbreak{} GLOSSARY\_\allowbreak{}DER:\allowbreak{}92}
\DERapparatusitem{ENV1608}{15;\allowbreak{}\allowbreak{} ms.\allowbreak{}25,\allowbreak{}\allowbreak{} (\allowbreak{}2.\allowbreak{}6.\allowbreak{}1d)\allowbreak{}}{pro}{in}{DER\_\allowbreak{}XIX:\allowbreak{}30}
\DERapparatusitem{ENV1609}{15;\allowbreak{}\allowbreak{} running\allowbreak{} text\allowbreak{} or\allowbreak{} formula}{filtered}{filtering}{GLOSSARY\_\allowbreak{}DER:\allowbreak{}92}
\DERapparatusitem{ENV1610}{15;\allowbreak{}\allowbreak{} ms.\allowbreak{}25,\allowbreak{}\allowbreak{} (\allowbreak{}2.\allowbreak{}6.\allowbreak{}3b)\allowbreak{}}{}{\_\allowbreak{}0}{DER\_\allowbreak{}XIX:\allowbreak{}31}
\DERapparatusitem{ENV1611}{15;\allowbreak{}\allowbreak{} running\allowbreak{} text\allowbreak{} or\allowbreak{} formula}{filtered}{filtering}{GLOSSARY\_\allowbreak{}DER:\allowbreak{}92}
\DERapparatusitem{ENV1612}{17;\allowbreak{}\allowbreak{} ms.\allowbreak{}27,\allowbreak{}\allowbreak{} definition\allowbreak{} 2.\allowbreak{}8\allowbreak{} opening}{as\allowbreak{} in\allowbreak{} theorem\allowbreak{} 2.\allowbreak{}1.\allowbreak{},\allowbreak{}}{as\allowbreak{} in\allowbreak{} theorem\allowbreak{} 2.\allowbreak{}6,\allowbreak{}}{DER\_\allowbreak{}XIX:\allowbreak{}35}
\DERapparatusitem{ENV1613}{17;\allowbreak{}\allowbreak{} running\allowbreak{} text\allowbreak{} or\allowbreak{} formula}{filtered}{filtering}{GLOSSARY\_\allowbreak{}DER:\allowbreak{}92}
\DERapparatusitem{ENV1614}{18;\allowbreak{}\allowbreak{} ms.\allowbreak{}29,\allowbreak{}\allowbreak{} (\allowbreak{}2.\allowbreak{}9.\allowbreak{}6)\allowbreak{}}{it}{rm}{DER\_\allowbreak{}XIX:\allowbreak{}36}
\DERapparatusitem{ENV1615}{18;\allowbreak{}\allowbreak{} ms.\allowbreak{}29,\allowbreak{}\allowbreak{} (\allowbreak{}2.\allowbreak{}9.\allowbreak{}7)\allowbreak{}}{it}{rm}{DER\_\allowbreak{}XIX:\allowbreak{}37}
\DERapparatusitem{ENV1616}{18;\allowbreak{}\allowbreak{} ms.\allowbreak{}30,\allowbreak{}\allowbreak{} stronger\allowbreak{} claim\allowbreak{} concerning\allowbreak{} TF[\allowbreak{}I]\allowbreak{}}{=\allowbreak{}\allowbreak{} \textbackslash{}\allowbreak{}\allowbreak{} v\textasciicircum{}\allowbreak{}*v\_\allowbreak{}*(\allowbreak{}TF(\allowbreak{}I\_\allowbreak{}0)\allowbreak{})\allowbreak{}\allowbreak{} \textbackslash{}\allowbreak{}}{}{DER\_\allowbreak{}XIX:\allowbreak{}39}
\DERapparatusitem{ENV1617}{preamble;\allowbreak{}\allowbreak{} local\allowbreak{} macro\allowbreak{} definition}{\textbackslash{}\allowbreak{}text\{\allowbreak{}def\}\allowbreak{}}{\textbackslash{}\allowbreak{}mathrm\{\allowbreak{}def\}\allowbreak{}}{TYPESETTING:\allowbreak{}7;\allowbreak{}\allowbreak{} NOTATION\_\allowbreak{}DER:\allowbreak{}17}
\DERapparatusitem{ENV1618}{21;\allowbreak{}\allowbreak{} ms38,\allowbreak{}\allowbreak{} Proposition3.\allowbreak{}5\allowbreak{} and\allowbreak{} equation3.\allowbreak{}5.\allowbreak{}1}{If\allowbreak{} \$\allowbreak{}\textbackslash{}\allowbreak{}Phi=\allowbreak{}\textbackslash{}\allowbreak{}Psi\textasciicircum{}\allowbreak{}*\$\allowbreak{}\allowbreak{} is\allowbreak{} a\allowbreak{} left}{If\allowbreak{} \$\allowbreak{}\textbackslash{}\allowbreak{}Phi=\allowbreak{}\textbackslash{}\allowbreak{}Psi\textasciicircum{}\allowbreak{}*\$\allowbreak{},\allowbreak{}\allowbreak{} with\allowbreak{} \$\allowbreak{}\textbackslash{}\allowbreak{}Psi=\allowbreak{}\textbackslash{}\allowbreak{}Phi\_\allowbreak{}*\$\allowbreak{},\allowbreak{}\allowbreak{} is\allowbreak{} a\allowbreak{} left}{DER\_\allowbreak{}XIX:\allowbreak{}41}
\DERapparatusitem{ENV1619}{23;\allowbreak{}\allowbreak{} ms40,\allowbreak{}\allowbreak{} first\allowbreak{} line\allowbreak{} of\allowbreak{} leaf}{4}{3}{DER\_\allowbreak{}XIX:\allowbreak{}43}
\DERapparatusitem{ENV1620}{23;\allowbreak{}\allowbreak{} ms41,\allowbreak{}\allowbreak{} proof\allowbreak{} of\allowbreak{} Proposition3.\allowbreak{}7,\allowbreak{}\allowbreak{} section\allowbreak{} of\allowbreak{} p}{p}{q}{DER\_\allowbreak{}XIX:\allowbreak{}45}
\DERapparatusitem{ENV1621}{23;\allowbreak{}\allowbreak{} ms42,\allowbreak{}\allowbreak{} question3.\allowbreak{}9(\allowbreak{}a)\allowbreak{}}{does\allowbreak{} not\allowbreak{} exist\allowbreak{} by\allowbreak{} always}{does\allowbreak{} not\allowbreak{} always\allowbreak{} exist}{DER\_\allowbreak{}XIX:\allowbreak{}46}
\DERapparatusitem{ENV1622}{25;\allowbreak{}\allowbreak{} ms46,\allowbreak{}\allowbreak{} equation3.\allowbreak{}12.\allowbreak{}2}{\textbackslash{}\allowbreak{}mathrm\{\allowbreak{}HO\}\allowbreak{}\_\allowbreak{}W(\allowbreak{}I)\allowbreak{}}{\textbackslash{}\allowbreak{}mathrm\{\allowbreak{}Ho\}\allowbreak{}\_\allowbreak{}W(\allowbreak{}I)\allowbreak{}}{DER\_\allowbreak{}XIX:\allowbreak{}49}
\DERapparatusitem{ENV1623}{25;\allowbreak{}\allowbreak{} running\allowbreak{} text\allowbreak{} or\allowbreak{} formula}{\textbackslash{}\allowbreak{}text\{\allowbreak{}def\}\allowbreak{}}{\textbackslash{}\allowbreak{}mathrm\{\allowbreak{}def\}\allowbreak{}}{TYPESETTING:\allowbreak{}7;\allowbreak{}\allowbreak{} NOTATION\_\allowbreak{}DER:\allowbreak{}17}
\end{document}
